\documentclass[11pt]{amsart}
\usepackage{amsmath,amssymb,amsthm}
\usepackage[margin=1in]{geometry}
\setlength{\emergencystretch}{2em}
\usepackage[hidelinks]{hyperref}

\theoremstyle{plain}
\newtheorem{theorem}{Theorem}[section]
\newtheorem{proposition}[theorem]{Proposition}
\newtheorem{lemma}[theorem]{Lemma}
\newtheorem{corollary}[theorem]{Corollary}
\theoremstyle{definition}
\newtheorem{definition}[theorem]{Definition}
\newtheorem{example}[theorem]{Example}
\newtheorem{hypothesis}[theorem]{Hypothesis}
\newtheorem{kstatement}{Statement}[section]
\theoremstyle{remark}
\newtheorem{remark}[theorem]{Remark}

\numberwithin{equation}{section}

\newcommand{\R}{\mathbb{R}}
\newcommand{\Z}{\mathbb{Z}}
\newcommand{\N}{\mathbb{N}}
\newcommand{\Q}{\mathbb{Q}}
\newcommand{\C}{\mathbb{C}}
\newcommand{\E}{\mathbb{E}}
\newcommand{\Tor}{\mathbb{T}}
\newcommand{\Oh}{O_h}
\newcommand{\Ot}{\mathrm{O}(3)}
\newcommand{\SO}{\mathrm{SO}}
\newcommand{\fW}{\mathfrak{W}}
\newcommand{\Wwin}{\mathfrak{W}}
\newcommand{\fF}{\mathfrak{F}}
\newcommand{\cA}{\mathfrak{A}}
\newcommand{\cAc}{\mathcal{A}}
\newcommand{\cB}{\mathcal{B}}
\newcommand{\cD}{\mathcal{D}}
\newcommand{\cF}{\mathcal{F}}
\newcommand{\cH}{\mathcal{H}}
\newcommand{\cI}{\mathcal{I}}
\newcommand{\cN}{\mathcal{N}}
\newcommand{\cS}{\mathcal{S}}
\newcommand{\cT}{\mathcal{T}}
\newcommand{\cV}{\mathcal{V}}
\newcommand{\Sig}{\Sigma}
\newcommand{\fl}{\flat}
\newcommand{\bg}{\mathbf{g}}
\newcommand{\eps}{\varepsilon}
\newcommand{\Nbar}{\bar N}
\newcommand{\nbar}{\bar n}
\newcommand{\XR}{X_{\mathrm R}}
\newcommand{\Xr}{X_{R}}
\newcommand{\Tub}{D_{\sigma_+,H}}
\newcommand{\abs}[1]{\lvert #1\rvert}
\newcommand{\norm}[1]{\lVert #1\rVert}
\newcommand{\ip}[2]{\langle #1,\,#2\rangle}
\newcommand{\ipj}[3]{\langle #2,\,#3\rangle_{#1}}
\newcommand{\ev}[1]{\langle #1\rangle}
\newcommand{\wick}[1]{{:}#1{:}}
\newcommand{\rRe}{\operatorname{Re}}
\newcommand{\Hp}[2]{H^{+}_{#1}(#2)}
\newcommand{\Hm}[2]{H^{-}_{#1}(#2)}
\DeclareMathOperator{\supp}{supp}
\DeclareMathOperator{\dist}{dist}
\DeclareMathOperator{\tr}{tr}
\newcommand{\Ker}{\operatorname{Ker}}

\begin{document}

\title[Symmetric phase of the abelian Higgs model]{The symmetric phase of the three-dimensional abelian Higgs model:\\
torus limits and the Osterwalder--Schrader axioms with cubic symmetry}
\author{}
\date{}
\subjclass[2020]{Primary 81T08; Secondary 81T13, 81T17, 81T25}
\keywords{abelian Higgs model, scalar quantum electrodynamics, constructive field theory, renormalisation group, Osterwalder--Schrader axioms, reflection positivity, cluster property}

\begin{abstract}
We study the Euclidean abelian Higgs model (scalar quantum electrodynamics) in three dimensions, with a non-compact $U(1)$ gauge field and a complex scalar field of positive mass squared $m^2$ and quartic coupling $\lambda>0$. The model is regularised on the lattice $\varepsilon\Z^3$ on periodic cubic tori, with a photon mass $\mu_0>0$ in Feynman form and King's mass counterterm computed on the torus at a fixed reference photon mass. For $0<\lambda<0.1m$ and $0<e^2<0.1m$ we prove that the generating functionals of the curvature $F(g)$ and of the renormalised Wick square $\wick{\abs{\varphi}^2}(h)$ are bounded on a complex tube, uniformly in the lattice spacing, in the size of the torus and in $\mu_0$, and also in a massless state in which the constant mode of the gauge field is averaged over its period. Along subsequences the spacing goes to zero and the torus to $\R^3$, and then $\mu_0$ goes to zero along a countable set. The limit states satisfy the Osterwalder--Schrader axioms OS0, OS1 for the translations and the cubic group, OS2 about every coordinate plane, and OS3. Their laws are integrals of translation-ergodic, reflection-positive components, each of which clusters in the mean and under separations in which all three coordinates diverge, and has a unique joint vacuum in each coordinate direction; observables linear in the curvature, with any square-integrable dependence on the Wick square, cluster along each coordinate axis. We also prove sign inequalities for truncated correlations, and that the curvature two-point function of every limit state is dominated by that of the free massless photon and has no massless component at positive photon mass. The ultraviolet input is Ba{\l}aban's renormalisation group, with three repairs of printed proofs; the proof of the ultraviolet bounds also uses seven technical estimates, stated in Appendix~K, whose proofs are not included in this version. Invariance under the full rotation group, the cluster property along a single axis for all observables, uniqueness of the limits, non-triviality and the masslessness of the photon are left open.
\end{abstract}

\maketitle
\setcounter{tocdepth}{1}
\tableofcontents

\section{Introduction}\label{sec:intro}

\subsection{The model and the question}\label{ssec:question}

The abelian Higgs model is the Euclidean field theory of a real vector field $A_\mu$ and a complex scalar field $\varphi$ on $\R^3$ with formal action
\begin{equation}\label{eq:action}
S(A,\varphi)=\int_{\R^3}\Big[\tfrac14 F_{\mu\nu}F_{\mu\nu}+\tfrac12\abs{(\partial_\mu+ieA_\mu)\varphi}^2+\tfrac12m^2\abs{\varphi}^2+\lambda\abs{\varphi}^4\Big]\,dx ,\qquad F_{\mu\nu}=\partial_\mu A_\nu-\partial_\nu A_\mu .
\end{equation}
This is King's normalisation \cite[(2.1)--(2.4), p.~652]{K1}, the one of the lattice action~\eqref{eq:matteraction} below, in which the coupling range of this paper is stated; in terms of $\psi:=\varphi/\sqrt2$ the action reads $\abs{(\partial_\mu+ieA_\mu)\psi}^2+m^2\abs{\psi}^2+4\lambda\abs{\psi}^4$, so in that normalisation the quartic coupling is $4\lambda$. The sign of $ieA_\mu$ is the one of the lattice covariant difference $\varepsilon^{-1}(e^{ie\varepsilon A_{x,\mu}}\varphi(x+\varepsilon e_\mu)-\varphi(x))$, and the gauge transformations are $\varphi\mapsto e^{i\chi}\varphi$, $A\mapsto A-e^{-1}d\chi$.

In three dimensions $e^2$ and $\lambda$ have the dimension of a mass and the model is superrenormalisable. For $m^2<0$ the scalar condenses and the photon acquires a mass by the Higgs mechanism. For $m^2>0$, the case treated here, the scalar is massive and the photon is expected to stay massless: this is the symmetric, or Coulomb, phase. This picture comes from perturbation theory, and at very small $\lambda$ it is expected to fail: one-loop perturbation theory predicts a first-order transition at positive $m^2$, on a surface $\lambda m^2\approx e^6/(72\pi^2)$ that lies inside the coupling range of our theorems, with a massive photon below it (Section~\ref{sec:disc}). None of our proofs uses the phase; where this paper speaks of the symmetric or Coulomb phase, it refers to couplings above that surface. The question we address is whether the Schwinger functions of the local gauge-invariant fields, the curvature $F_{\mu\nu}$ and the Wick square $\wick{\abs{\varphi}^2}$, can be constructed as limits of a lattice regularisation and satisfy the Osterwalder--Schrader axioms \cite{OS73,OS75}: regularity (OS0), Euclidean invariance (OS1), reflection positivity (OS2), symmetry (OS3) and clustering (OS4).

The gauge field is non-compact: on the lattice, $A$ is real-valued on the bonds and the gauge action is the quadratic form of the lattice curl, as in Ba{\l}aban \cite{Ba1,Ba2,Ba3,Ba4} and King \cite{K1,K2}. The exact lattice relation between the two formulations is that, for observables of integer charge on a finite lattice with free boundary conditions (trivial cohomology), the non-compact model is the compact Villain model with every monopole removed \cite[Sect.~3, Thm.~3.1, p.~286]{BBIJ}. This choice matters in the Coulomb phase: in the compact pure gauge theory in three dimensions monopoles produce a mass gap and confinement for all couplings \cite{GM82}, while the non-compact pure gauge theory is the free field, whose photon is massless. For the non-compact model with matter studied here, masslessness of the photon is expected for couplings above the first-order surface, but it is not proved (Section~\ref{sec:disc}).

\subsection{Background}\label{ssec:background}

Reflection positivity of lattice gauge theories goes back to Osterwalder and Seiler \cite{OS78} for compact gauge groups, and to \cite[Thm.~5.3, p.~272]{BFS1} for the non-compact abelian field; Seiler's lectures \cite{Sei82} contain the general framework. In two dimensions the model, with a non-compact continuum gauge field and lattice matter, was constructed by Brydges, Fr\"ohlich and Seiler \cite{BFS1,BFS2,BFS3}, who also introduced the correlation inequalities that make the generating functionals monotone in the photon covariance. The ultraviolet problem in three dimensions was solved by Ba{\l}aban \cite{Ba1,Ba2,Ba3,Ba4}, who proved two-sided bounds
\begin{equation}\label{eq:K0}
e^{-E_-\abs{T}}\ \le\ Z^\varepsilon(T)\ \le\ e^{E_+\abs{T}}
\end{equation}
on the partition function of the model with a photon mass on a periodic torus $T$, with constants independent of the lattice spacing $\varepsilon$ and of $T$ \cite[Thm.~(1.14), p.~606]{Ba1}, \cite[Thm.~(1.3), p.~556]{Ba2}. Ba{\l}aban's smallness conditions are conditions on the running couplings $et^{1/2}$, $\lambda t$ and $m^2t^2$ at short scales $t\le\varepsilon_0$, not on the couplings themselves.

King \cite{K1,K2} turned these stability bounds into limits. He proved convergence of the generating functional of $F(g)$ and $\wick{\abs{\varphi}^2}(h)$ as $\varepsilon\to0$ on tori and on cubes, constructed the infinite-volume limit through cubes with Neumann conditions for $A$ and Dirichlet conditions for $\varphi$ using the monotonicity in the volume given by the correlation inequalities, and proved translation invariance and reflection positivity of the resulting Schwinger functions. King kept the photon mass $\mu_0>0$ throughout, and listed what remained: ergodicity, rotation invariance and regularity \cite[p.~326]{K2}. For rotations he compared, on a fixed torus, the lattice with its copy rotated by the angle $\theta_0=\arctan(3/4)$ \cite[Thm.~2.4, pp.~326--327]{K2}.

The Higgs phase was treated on the unit lattice by Ba{\l}aban, Brydges, Imbrie and Jaffe \cite{BBIJ}, and in the continuum by Ba{\l}aban, Imbrie and Jaffe \cite{BIJ}, who constructed the effective action down to the unit scale. The ultraviolet renormalisation of scalar quantum electrodynamics in three dimensions was revisited by Dimock \cite{Di15}, who later treated the fermionic model \cite{Di18}. Continuum limits of the compact pure $U(1)$ theory were obtained by Gross \cite{Gr83} in three dimensions and Driver \cite{Dr87} in four. Gauge theories with matter in three dimensions have recently been approached through stochastic quantisation \cite{CCHS,CS25}.

\subsection{Main results}\label{ssec:mainresults}

We do not use King's cube states. We work throughout on periodic cubic tori $T_j$ of side $j!\,\ell_0$, on which the lattice model is exactly invariant under the translations of the lattice and under the cubic group, and is reflection positive about lattice planes. The photon mass is the only infrared regulator, and the scalar mass counterterm is King's counterterm computed on the torus itself at a fixed reference photon mass $\mu_*=m$; we call this choice the \emph{torus counterterm scheme} (Definition~\ref{def:counterterm}). Besides the states at photon mass $\mu_0\in(0,\mu_*]$ we use the state $\fl$, in which the photon has no mass and the constant mode of the gauge field is averaged over its period; at fixed lattice it is the limit $\mu_0\downarrow0$ (Lemma~\ref{lem:flat}). The objects in the following statements are defined in Section~\ref{sec:model}.

\medskip\noindent\textbf{Theorem A} (existence and OS0--OS3). \emph{Let $m>0$, $0<\lambda<0.1m$ and $0<e^2<0.1m$, and consider the lattice states $\langle\cdot\rangle_{\varepsilon,j,\mu_0}$ of Definition~\ref{def:states} in the torus counterterm scheme with reference photon mass $\mu_*=m$, on the tori $T_j$ of side $\ell_j=j!\,\ell_0$, at the spacings $\varepsilon=2^{-K}\varepsilon_0$, at the photon masses $\mu_0\in(0,m]$ and in the state $\fl$. There are $\varepsilon_0=2^{-n_0}$, $\ell_0\ge183/m$, $K_0\in\N$, $\sigma_+>0$ and $C<\infty$ such that the following holds, with $H:=m^2$.}
\begin{enumerate}
\item[(i)] \emph{(Uniform tube bound.) For all $K\ge K_0$, $j\ge1$ and $\mu_0\in(0,m]\cup\{\fl\}$, every finite family of real test functions supported in a ball of radius $R$ with $\ell_j>2R+2$, and every $\zeta$ in the tube $D_{\sigma_+,H}$, the generating functional of the curvature and the renormalised Wick square satisfies}
\[
\abs{G_{\varepsilon,j,\mu_0}(\zeta)}\ \le\ G_{\varepsilon,j,\mu_0}(\rRe\zeta)\ \le\ \exp\Big(\tfrac12\Big\|\sum_i(\rRe z_i)g_i\Big\|_{2,\varepsilon}^2+C\norm{k(\zeta)}_{L^1,\varepsilon}\Big).
\]
\item[(ii)] \emph{(Limits.) Let $Q_0\subset(0,m]$ be countable with $0$ as an accumulation point. Every sequence $(K_l,j_l)$ with $K_l,j_l\to\infty$ has a subsequence along which the lattice Schwinger functions at every $\mu_0\in Q_0$ and in the state $\fl$ converge, to limit states $S^{\mu_0}$ and $S^\fl$; along this subsequence $S^0:=\lim_{\mu_0\downarrow0,\,\mu_0\in Q_0}S^{\mu_0}$ exists. Each of these limit states is, on test functions in $C_c^\infty$, the family of moments of a Borel probability measure on $H^{-s'}_{\mathrm{loc}}(\R^3;\Lambda^2\R^3\oplus\R)$ for every $s'>3$.}
\item[(iii)] \emph{(Axioms.) Each limit state satisfies OS0 with the constants $\sigma_+$, $H$, $C$ of (i), OS1 for the group $\R^3\rtimes\Oh$ generated by the translations and the cubic group, OS2 about every coordinate plane, and OS3 (Definition~\ref{def:OS}). Its Schwinger functions satisfy $\abs{S_{n,p}(g;h)}\le e^{C\sigma_0}\mathrm{e}^n(n!)^{1/2}(\mathrm{e}/\sigma_0)^p\,p!\prod_i\norm{g_i}_2\prod_kN(h_k)$, $\sigma_0:=\min(\sigma_+,H)/2$.}
\item[(iv)] \emph{(Monotonicity.) For every generating functional $\Theta$ of the class $\cD$ of Section~\ref{ssec:observables}, $\mu_0\mapsto S^{\mu_0}(\Theta)$ is non-increasing on $Q_0$, and $S^{\mu_0}(\Theta)\le S^\fl(\Theta)$, $S^0(\Theta)\le S^\fl(\Theta)$.}
\end{enumerate}

\medskip
Theorem~A is Theorem~\ref{thm:main}, which also gives the convergence of the generating functionals on the tube. Euclidean invariance is established for the translations of $\R^3$ and the cubic group $\Oh$, the symmetry group of the lattice; invariance under the full rotation group is not claimed (Remark~\ref{rem:cubic}). Theorem~\ref{thm:mainL3} gives parts (ii)--(iv) also for the lattices $3^{-K}\varepsilon_0\Z^3$ with blocking factor $3$, under the hypothesis that the uniform tube bound holds for that family; the interest of blocking factor $3$ lies in the rotations (Section~\ref{sec:disc}).

\medskip\noindent\textbf{Theorem B} (weak clustering). \emph{Let $S$ be a limit state of Theorem~A. Its law is the only Borel probability measure $\mu$ with these moments on a Hermite--Sobolev space of distributions, and $\mu=\int\mu_\xi\,d\nu(\xi)$ with components $\mu_\xi$ such that, for $\nu$-almost every $\xi$: $\mu_\xi$ is invariant and ergodic under the translations of $\R^3$, invariant under the three coordinate reflections, and reflection positive about every coordinate plane; for all $G,G'\in L^2(\mu_\xi)$, $\E_{\mu_\xi}[G'\,U(a)G]\to\E_{\mu_\xi}[G']\,\E_{\mu_\xi}[G]$ in the mean over cubes and as $\min_j\abs{a_j}\to\infty$, where $U(a)$ is the translation by $a$; and for each coordinate direction $e_j$ the Osterwalder--Schrader reconstruction of $\mu_\xi$ with $e_j$ as Euclidean time has a positive Hamiltonian $H_j$ and a unique joint vacuum: the vectors fixed by the semigroup $e^{-tH_j}$, $t\ge0$, and by the translations orthogonal to $e_j$ form one line $\C\Omega_j$.}

\medskip
Theorem~B is Theorem~\ref{thm:cluster}. In the terms of Definition~\ref{add:OS4} it says that the components satisfy OS4 in the weak form, ergodicity under the translation group $\R^3$; equivalently, they cluster under every separation in which all three coordinates diverge. It is not the cluster property OS4 of Definition~\ref{def:OS4}, which asks for clustering as the separation runs along a single coordinate axis: single-axis clustering does not follow from the structure that Theorem~B uses, namely invariance under translations and reflections, reflection positivity, and ergodicity under $\R^3$ (Proposition~\ref{prop:axisobs}); it would follow from invariance under all rotations (Corollary~\ref{p:cor:rotcor}). For part of the observables we prove it.

\medskip\noindent\textbf{Theorem C} (single-axis clustering, partial). \emph{Let $S$ be a limit state of Theorem~A, $\mu$ its law, and $\mathcal M$ the $\sigma$-algebra generated by the Wick-square fields $\wick{\abs\varphi^2}(h)$.}
\begin{enumerate}
\item[(a)] \emph{For every real 2-form $g$ with components in $C_c^\infty(\R^3)$ and every $Y\in L^2(\mathcal M,\mu)$, the observable $F(g)Y$ clusters against every $G'\in L^2(\mu)$: $\E_\mu[G'\,U(a)(F(g)Y)]\to0=\E_\mu[G']\,\E_\mu[F(g)Y]$ as $\abs a\to\infty$ in any manner, in particular along each coordinate axis.}
\item[(b)] \emph{If the latent density $\Delta_S$ of the Wick square vanishes (Definition~\ref{def:os4:pressure}; this holds if some limit, along the defining sequence of $S$, of the torus pressures of the Wick square is differentiable at zero source), the same holds, with limit $\E_\mu[G']\,\E_\mu[G]$, for every $G$ in the closed span of $L^2(\mathcal M,\mu)$ and of the products in (a).}
\item[(c)] \emph{If the block pressures of the bounded local observables $\cos\Phi(f)$, $\sin\Phi(f)$, with $\Phi(f)$ a smeared field, are differentiable at zero (hypothesis (D) of Corollary~\ref{cor:os4:OS4}), then $\mu$ is ergodic, $S$ satisfies OS4, and every pair $G,G'\in L^2(\mu)$ clusters along each coordinate axis.}
\end{enumerate}

\medskip\noindent\textbf{Theorem D} (signs and the curvature two-point function). \emph{Let $S$ be a limit state of Theorem~A.}
\begin{enumerate}
\item[(a)] \emph{For real 2-forms $g,g'$ and non-negative $h,h'$ with components in $C_c(\R^3)$: $\langle F(g)^2;\wick{\abs\varphi^2}(h)\rangle_S\le0$, $\langle\wick{\abs\varphi^2}(h);\wick{\abs\varphi^2}(h')\rangle_S\ge0$, and $S^T_{4,0}(g,g,g',g')+2S_{2,0}(g,g')^2\ge0$, where $\langle\,\cdot\,;\cdot\,\rangle_S$ and $S^T$ denote truncated functions. If the first inequality is strict for one pair $(g,h)$, then $S$ is not Gaussian, and the curvature and the Wick square are not independent under its law.}
\item[(b)] \emph{The two-point function of the smeared magnetic field is a Fourier multiplier $m_S(p)$ with $0\le m_S(p)\le P_T(p)$, where $P_T(p)$ is the orthogonal projection of $\C^3$ onto $p^\perp$; the multiplier $P_T$ is that of the free massless photon. At positive photon mass, $m_S(p)\le\min(1,\abs p^2/\mu_0^2)P_T(p)$, so $S^{\mu_0}$ has no massless component.}
\item[(c)] \emph{If $S^\fl$ or $S^0$ is massless in the curvature sector (Definition~\ref{def:M-massless}), its curvature correlations are not integrable, and in every coordinate direction the curvature sector of its reconstruction has no mass gap.}
\end{enumerate}

\medskip
Theorems~C and~D are Theorems~\ref{thm:os4main} and~\ref{thm:NTMmain}. Theorem~D does not show that the limit states are non-Gaussian, nor that the photon is massless at $\mu_0=0$. In Theorem~\ref{thm:NT-cond} we show that the first inequality of (a) is strict at small $e^2+\lambda$, under two hypotheses that we formulate and do not prove: a remainder bound (R) for the effective action with sources, and a one-loop statement (C) (Section~\ref{sec:nontriv}).

The ultraviolet input is taken from the printed papers of Ba{\l}aban \cite{Ba1,Ba2,Ba3,Ba4} and King \cite{K1,K2}, with three repairs of printed steps. First, the printed proof of \cite[Lemma~2.4, p.~572]{Ba2} does not cover the lattice points near the edges of the region on which the lemma is applied; we prove the statement of the lemma, under explicit conditions on the regulator data, by an argument that covers these points (Lemma~\ref{lem:lem24}). Second, the proof of \cite[Prop.~1, pp.~420--445]{Ba3} contains sketched steps; we write out each of them, with the depth inside the localisation domain that it needs, because our application uses the proposition at the limit of what its printed nesting provides (Proposition~\ref{prop:prop1}). Third, the printed proof of the combinatorial estimate \cite[(3.42), p.~592]{Ba2}, the last step of the proof of Ba{\l}aban's upper vacuum bound \cite[Thm.~(1.3)]{Ba2}, is incomplete: the step to \cite[(3.47), p.~593]{Ba2}, introduced by ``after easy transformations'', is not written out; the cube count \cite[(3.45), p.~592]{Ba2} fails for blocks near the boundary of the region; and the closing estimates \cite[(3.51)--(3.54), p.~594]{Ba2} bound only the large-field sets $\Lambda_0^{(k)c}$, while (3.42) also contains collar and shell sets. We bound these sets by families of cubes, pay the entropy of the sum over regions, and close (3.54) at $p=6$ with an explicit upper bound on $\varepsilon_0$ (Proposition~\ref{app:G1:I6}). This completes the printed proof of the upper vacuum bound at this step; the other estimates that its printed proof omits are discussed in Remark~\ref{rem:UV-blockD}.

\subsection{Method}\label{ssec:method}

The infrared problem of a constructive field theory is usually controlled by decay: a mass gap makes a cluster expansion converge, the expansion gives a unique infinite-volume limit, and uniqueness gives the symmetries. In the Coulomb phase no such expansion is available. We replace it by three facts that hold uniformly in every regulator, on a single torus at a time, and are therefore inherited by every limit.

\emph{(a) A uniform bound on a tube.} The central estimate (Theorem~A(i)) bounds $\abs{\langle e^{\sum z_iF(g_i)+\sum w_k\XR(h_k)}\rangle}$ on a tube $\{\abs{\rRe w}\ \text{small}\}$, with all of $\C$ allowed for the coefficients $z_i$ of the curvature, uniformly in the spacing, the torus and $\mu_0$. It has three factors. The curvature factor $e^{\frac12\norm{g}_2^2}$ is Gaussian domination of the photon by the free field, which on a torus holds with the sharp constant because the componentwise lattice Laplacian is the Hodge Laplacian (Lemma~\ref{lem:C-Fsector}). The factor for positive $\wick{\abs{\varphi}^2}$-sources comes from a convexity bound for constant sources on the same torus (Proposition~\ref{prop:ch-Pplus}), spread over arbitrary sources by the chessboard estimate of reflection positivity (Proposition~\ref{prop:ch-T7}); its input is a positive gap between King's counterterms at two scalar masses on that torus (Theorem~\ref{thm:ch-W}), which we verify by an explicit computation of the lattice graphs for the couplings of the window~\eqref{eq:window} below. The factor for negative sources (Proposition~\ref{prop:ch-Tflat}) comes from a vacuum bound for the model with a deformed photon covariance (Theorem~\ref{thm:T6}), which we derive from Ba{\l}aban's renormalisation group.

\emph{(b) Compactness and exact symmetries.} The uniform bound gives compactness of the lattice Schwinger functions and tightness of the lattice laws. Limits exist along subsequences; we do not prove uniqueness, and do not need it: on cubic tori the translations of $\varepsilon\Z^3$ and the cubic group are exact symmetries of every lattice state, and reflection positivity holds about pairs of lattice planes, so these properties are closed under limits. Continuity, which the bound supplies, extends the translations to $\R^3$.

\emph{(c) Monotonicity in the photon mass.} At fixed lattice and torus, the Brydges--Fr\"ohlich--Seiler inequalities make the generating functionals of the class $e^{F(g)-\XR(q)}$, $q\ge0$, non-decreasing in the photon covariance \cite[Cor.~6.3(1), p.~281]{BFS1} (those of $e^{+\XR(q)}$, $q\ge0$, are non-increasing). Taking the lattice limit along one subsequence common to a countable set of photon masses, the monotonicity passes to the limit and gives the massless limit $S^0$ and the inequality $S^0\le S^\fl$.

\emph{The weak cluster property} (Theorem~B) uses only translations, the three coordinate reflections and reflection positivity. A function invariant under translations along a coordinate axis is fixed by the reflection in a plane orthogonal to that axis; hence the invariant $\sigma$-algebra is fixed by the reflections, the ergodic components inherit reflection positivity, and on each component the Osterwalder--Schrader reconstruction in each coordinate direction identifies the zero-energy vectors with the functions invariant under the translations along that direction.

\emph{Restriction on the couplings.} The window
\begin{equation}\label{eq:window}
\fW:=\{(e,\lambda,m,\eta,m_1,m_2):\ m>0,\ 0<\lambda<0.1m,\ 0<e^2<0.1m,\ 0<\eta<1,\ 2m<m_1<m_2\}
\end{equation}
enters through one input of our own: the positive gap of (a) between King's counterterms at the scalar masses $0.7m$ and $m$, uniform in the spacing and the torus (Theorem~\ref{thm:ch-W}). The other parameters $\eta,m_1,m_2$ of $\fW$ are auxiliary (anchor masses of the vacuum bound for the deformed model), and can be chosen for every $(e,\lambda,m)$ with $m>0$, $0<\lambda<0.1m$, $0<e^2<0.1m$; Lemma~\ref{lem:window} shows that all conditions on the regulator data can be met at every point of $\fW$.

\subsection{What is not proved here}\label{ssec:notin}

We do not prove invariance under rotations outside $\Oh$; OS4 in the sense of Definition~\ref{def:OS4}, for the limit states or for their ergodic components, beyond the parts given by Theorem~C; uniqueness of the limits, the equality $S^0=S^\fl$, or the identification with King's cube limits; non-triviality of the limit states; masslessness of the photon at $\mu_0=0$; extremality of the limit states, or the existence of Wightman fields. The proofs do not use the phase, and one-loop perturbation theory predicts a first-order transition inside the coupling range of Theorem~A. Section~\ref{sec:disc} states these problems, says where the methods of this paper stop for each of them, and explains what is known about rotations, including the role of blocking factor $3$.

\subsection{Organisation and conventions}\label{ssec:org}

Section~\ref{sec:model} fixes the regularised model and the objects of the main results, which are stated in Section~\ref{sec:main}. Sections~\ref{sec:UV}--\ref{sec:chain} contain the proof of the uniform tube bound: Section~\ref{sec:UV} collects the results of Ba{\l}aban and King that we use; Section~\ref{sec:T6} states the vacuum bound for the model with a deformed photon covariance (Theorem~\ref{thm:T6}), describes the architecture of its proof, and contains the three repairs; Section~\ref{sec:chain} derives the uniform tube bound (Theorem~\ref{thm:ch-T9}) from it. Sections~\ref{sec:C-lattice}--\ref{sec:C-OS} prove parts (ii)--(iv) of Theorem~A: lattice statements on one torus at a time (Section~\ref{sec:C-lattice}), limits along subsequences (Section~\ref{sec:C-limits}), and the axioms OS0--OS3 together with what is known about rotations of infinite order (Section~\ref{sec:C-OS}). Section~\ref{sec:OS4} proves Theorems~B and~C, Section~\ref{sec:recon} describes the reconstructed theory, Section~\ref{sec:nontriv} proves the correlation signs and non-triviality under hypotheses (R) and (C), and Section~\ref{sec:massless} studies the curvature two-point function. Section~\ref{sec:disc} is a list of open problems.

The appendices contain the longer proofs; each begins by saying which result of the body it proves. Appendix~\ref{app:B} fixes the regulator data, the order in which they are chosen, and the conditions they satisfy, and proves that these conditions can be met at every point of $\fW$ (Lemma~\ref{lem:window}). Appendix~B contains supplementary proofs for the lattice model, Appendix~C the identification of King's counterterms and vacuum constants with Ba{\l}aban's, Appendix~D the form of \cite[Prop.~1]{Ba3} on the torus, Appendices~E and~F the ingredients and the assembly of the proof of Theorem~\ref{thm:T6}, Appendix~G the counterterm gaps, Appendix~H the bounds on c-number differences, Appendix~I the bound for negative sources and the choice of the auxiliary parameters, and Appendix~\ref{app:A} the measure-theoretic facts used in Section~\ref{sec:OS4}.

The logical structure is as follows. Theorem~A(i) is Theorem~\ref{thm:ch-T9}, whose proof occupies Sections~\ref{sec:UV}--\ref{sec:chain} and Appendices~\ref{app:B} and C--I. Sections~\ref{sec:C-lattice}--\ref{sec:C-OS} use from Sections~\ref{sec:UV}--\ref{sec:chain} only Theorem~\ref{thm:ch-T9}, and Sections~\ref{sec:OS4}--\ref{sec:massless} use Theorem~A together with estimates on the lattice tori from Sections~\ref{sec:chain} and~\ref{sec:C-lattice}. A reader interested in the passage from the uniform bound to the axioms may take Theorem~\ref{thm:ch-T9} as given and read Sections~\ref{sec:model}, \ref{sec:main} and~\ref{sec:C-lattice}--\ref{sec:recon}.

A formalisation of parts of the argument in the Lean proof assistant is in preparation.

\emph{Conventions.} Throughout, $d=3$ and $\N=\{1,2,3,\dots\}$. The letter $e$ always denotes the charge; exponentials are written $e^x$ or $\exp(x)$, and Euler's number, where it occurs as a constant factor, is written $\mathrm{e}$. Masses and lengths are measured in the units of \cite{K1}, in which the coarsest renormalisation-group scale is $1$. Equations are numbered within sections, and theorems, propositions, lemmas, corollaries, definitions, examples and remarks share one counter in each section.

\section{The regularised model}\label{sec:model}

Throughout, $m>0$, $\lambda>0$ and $e\neq0$ are fixed, and $\mu_*:=m$ is the reference photon mass.

\subsection{Lattices and tori}\label{ssec:tori}

Fix the blocking factor $L\in\{2,3\}$ ($L=2$ in Theorem~\ref{thm:main}, $L=3$ in Theorem~\ref{thm:mainL3}), the block-size parameter $M_b\in\N$ of Ba{\l}aban's renormalisation group (King's $M$ \cite{K1}), and $n_0\in\N$. Put
\[
\varepsilon_0:=L^{-n_0},\qquad \varepsilon=\varepsilon_K:=L^{-K}\varepsilon_0\quad(K=0,1,2,\dots).
\]
Fix $\ell_0\in c_LM_b\N$, with $c_2=2$ and $c_3=18$. For $j\ge1$ let $\ell_j:=j!\,\ell_0$, let $T_j:=(\R/\ell_j\Z)^3$ be the periodic cube of side $\ell_j$, and $T_{j,\varepsilon}:=T_j\cap\varepsilon\Z^3$ the lattice torus. A site is a point $x\in T_{j,\varepsilon}$; a (positively oriented) bond is a pair $b=(x,x+\varepsilon e_\mu)$, written $(x,\mu)$; a plaquette is $p=(x;\mu,\nu)$ with $\mu<\nu$, the square with corners $x$, $x+\varepsilon e_\mu$, $x+\varepsilon e_\nu$, $x+\varepsilon e_\mu+\varepsilon e_\nu$. We write $\abs{T_j}=\ell_j^3$.

\begin{lemma}[torus sizes]\label{lem:sizes}
For all $j\ge1$ and $K\ge0$: (i) the number of sites per side $N=\ell_j/\varepsilon_K$ is an even integer; (ii) for every $k\in\{0,1,\dots,K+n_0\}$, $\ell_j/(M_bL^k\varepsilon_K)\in2\N$; (iii) $\ell_j\mid\ell_{j+1}$ and $\ell_j\to\infty$; (iv) $\varepsilon_{K'}\Z^3\supset\varepsilon_K\Z^3$ for $K'\ge K$.
\end{lemma}
\begin{proof}
$\ell_j/(M_bL^k\varepsilon_K)=j!\,(\ell_0/M_b)\,L^{K+n_0-k}$, and $\ell_0/M_b\in c_L\N\subset2\N$; this is (ii), and (i) is the case $k=0$ multiplied by $M_b$. (iii) and (iv) are immediate.
\end{proof}

\subsection{Fields, action and counterterm}\label{ssec:action}

The gauge field is a function $A$ on bonds, $A_{x,\mu}\in\R$ (non-compact); the Higgs field is a function $\varphi$ on sites, $\varphi(x)\in\C$, with two real components and $\abs{\varphi}^2=(\rRe\varphi)^2+(\operatorname{Im}\varphi)^2$. Functions on sites, bonds and plaquettes carry the inner products $\ip{u}{v}_\varepsilon:=\sum\varepsilon^3u\,v$, and $\norm u_{2,\varepsilon}^2:=\ip uu_\varepsilon$. The lattice exterior derivative is $(d\chi)_{x,\mu}:=\varepsilon^{-1}(\chi(x+\varepsilon e_\mu)-\chi(x))$ on functions and
\[
(dA)(x;\mu,\nu):=\varepsilon^{-1}\big(A_{x,\mu}+A_{x+\varepsilon e_\mu,\nu}-A_{x+\varepsilon e_\nu,\mu}-A_{x,\nu}\big)
\]
on bond functions; $d^*$ denotes the adjoint. We write $-\Delta_\varepsilon:=d^*d+dd^*$ on bond functions. On the periodic lattice this operator acts on each component $A_{\cdot,\mu}$, viewed as a function of $x$, as the lattice Laplacian $-\Delta_\varepsilon$ on functions on sites (Lemma~\ref{app:G1:X}); in particular $-\Delta_\varepsilon\ge0$, with kernel the constant bond functions.

For a photon mass $\mu_0>0$ and a real bare coefficient $a$ the action is
\begin{equation}\label{eq:latticeaction}
S_{\mu_0,a}(A,\varphi):=\tfrac12\ip{A}{(-\Delta_\varepsilon+\mu_0^2)A}_\varepsilon+S^{\mathrm m}_a(A,\varphi),
\end{equation}
\begin{equation}\label{eq:matteraction}
S^{\mathrm m}_a(A,\varphi):=\tfrac12\sum_{x,\mu}\varepsilon^3\,\varepsilon^{-2}\abs{e^{ie\varepsilon A_{x,\mu}}\varphi(x+\varepsilon e_\mu)-\varphi(x)}^2+\sum_x\varepsilon^3\big(\tfrac12a\abs{\varphi(x)}^2+\lambda\abs{\varphi(x)}^4\big).
\end{equation}
The gauge part is the Feynman form of \cite[(2.4), p.~652]{K1}; it equals $\frac12\norm{dA}_{2,\varepsilon}^2+\frac12\norm{d^*A}_{2,\varepsilon}^2+\frac12\mu_0^2\norm{A}_{2,\varepsilon}^2$. The matter part is invariant under the lattice gauge transformations $\varphi\mapsto e^{i\chi}\varphi$, $A\mapsto A-e^{-1}d\chi$, $\chi$ real on sites (Lemma~\ref{app:G1:gauge}).

\begin{definition}[King's counterterm on a torus; the torus counterterm scheme]\label{def:counterterm}
For $M>0$ let $\delta m^2_{K,T}(\varepsilon;M)$ be the mass counterterm of \cite[(3.32)--(3.33), p.~659]{K1} computed on the lattice torus $T_{j,\varepsilon}$, with scalar mass $M$, scalar propagator $(-\Delta_\varepsilon+M^2)^{-1}$ and photon propagator $(-\Delta_\varepsilon+\mu_*^2)^{-1}$, both on $T_{j,\varepsilon}$ with periodic boundary conditions and without block averaging, so that every line commutes with the translations by $\varepsilon\Z^3$. It is the polynomial in $(e,\lambda)$ over the orders $2\le\alpha+2\beta\le4$ of \cite[(3.32)]{K1} whose coefficients are fixed, order by order, by the condition that the amputated one-particle-irreducible two-point function of $\varphi$ at zero momentum vanish at the orders $e^\alpha\lambda^\beta$, $2\le\alpha+2\beta\le4$. Put $a_{K,T}(M):=M^2+\delta m^2_{K,T}(\varepsilon;M)$. The bare coefficient of the model on $T_j$ at spacing $\varepsilon$ is
\[
b_{T_j}(\varepsilon):=a_{K,T_j}(m),
\]
the same for every photon mass $\mu_0\in(0,\mu_*]$ and for the state $\fl$ below. We call this choice the \emph{torus counterterm scheme}.
\end{definition}
The odd orders in $e$ of $\delta m^2_{K,T}$ vanish, because the substitution $A\mapsto-A$, $e\mapsto-e$ fixes $\varphi$ and leaves the matter action~\eqref{eq:matteraction} and the centred Gaussian photon measure with covariance $(-\Delta_\varepsilon+\mu_*^2)^{-1}$ invariant (Lemma~\ref{app:G1:oddorders}); so the monomials are $e^2,\lambda,e^4,e^2\lambda,\lambda^2$.

\begin{remark}\label{rem:counterterm}
(i) On $T_{j,\varepsilon}$ the counterterm of Definition~\ref{def:counterterm} coincides with Ba{\l}aban's \cite[p.~417]{Ba3}: both are the zero-momentum value of the amputated one-particle-irreducible two-point function on the same lattice torus, and the derivative parts have zero sum over a torus (Lemma~\ref{app:G1:ctid}).
(ii) The coefficient is computed at the reference photon mass $\mu_*$, not at $\mu_0$. This is what makes the limit $\mu_0\downarrow0$ at fixed lattice possible (Lemma~\ref{lem:flat}): a counterterm computed at photon mass $\mu_0$ on a torus contains the contribution of the constant mode of the photon, which has no limit. At order $e^2$ it is exactly $-3e^2/(\mu_0^2\abs{T})$: the constant mode $c$ has variance $(\mu_0^2\abs{T})^{-1}$ in each of the $d=3$ components (Definition~\ref{def:states}), the seagull term $\tfrac12e^2\sum_{x,\mu}\varepsilon^3A_{x,\mu}^2\,\rRe\big(\bar\varphi(x)\varphi(x+\varepsilon e_\mu)\big)$ of \eqref{eq:matteraction} turns it into the mass shift $3e^2/(\mu_0^2\abs{T})$ at zero momentum, and the exchange graph gets no constant-mode contribution there, because its vertex at zero photon momentum is proportional to $\sin(\varepsilon p_\mu)$ and $p=0$. The orders $e^4$ and $e^2\lambda$ contain further such terms.
(iii) The additive vacuum constants of \cite{K1,Ba1} play no role for normalised expectations and are omitted.
\end{remark}

\subsection{The states}\label{ssec:states}\label{ssec:C-flat}

Write $\Pi_0$ for the orthogonal projection onto the constant bond functions $A_{x,\mu}=c_\mu$, $c\in\R^3$, and $\Pi^\perp:=1-\Pi_0$; every bond function decomposes uniquely as $A=A^\perp+c$. Put $P:=2\pi/(\abs{e}\,\ell_j)>0$; the coupling $e\ne0$ may have either sign, since the window constrains only $e^2$.

\begin{definition}[observables at fixed lattice]\label{def:algebra}
$\cA_\varepsilon$ is the set of complex-valued Borel functions of the plaquette variables $\big((dA)(p)\big)_p$ and of the site variables $\big(\abs{\varphi(x)}^2\big)_x$. For $F\in\cA_\varepsilon$ with $F\ge0$ or $F$ bounded, and for a real function $k$ on sites, put
\[
\Psi_F(A):=\int F(A,\varphi)\,e^{-S^{\mathrm m}_b(A,\varphi)}\prod_xd^2\varphi(x),\qquad
W_k(A):=\int e^{-S^{\mathrm m}_b(A,\varphi)+\sum_x\varepsilon^3k(x)\abs{\varphi(x)}^2}\prod_xd^2\varphi(x),
\]
with $b=b_{T_j}(\varepsilon)$.
\end{definition}

An integer-valued function $n$ on bonds is called a \emph{current}; it is \emph{divergence-free} if $\sum_\mu\big(n_{x,\mu}-n_{x-\varepsilon e_\mu,\mu}\big)=0$ at every site $x$.

\begin{lemma}[positive type; current representation]\label{lem:postype}\label{lem:C-current}
For every real function $k$ on sites, $W_k$ is finite and
\begin{equation}\label{eq:C-current}
W_k(A)=\sum_{n}w_k(n)\,e^{ie\varepsilon\sum_bn_bA_b},\qquad w_k(n)=w_k(-n)\ge0,\qquad \sum_nw_k(n)=W_k(0)<\infty,
\end{equation}
the sum running over the divergence-free currents $n$. In particular $0<W_k(A)\le W_k(0)$. For a divergence-free current $n$ and $\mu\in\{1,2,3\}$, the flux $\omega_\mu(n):=\sum_{x:\,x_\mu=s}n_{x,\mu}$ of $n$ through the bonds in direction $\mu$ that start in the lattice plane $\{x_\mu=s\}$ is an integer that does not depend on $s$ (the winding of $n$ in direction $\mu$), and for every constant $c$
\[
e\varepsilon\sum_bn_bc_{\mu(b)}=e\ell_j\sum_\mu c_\mu\,\omega_\mu(n),
\]
where $\mu(b)$ is the direction of $b$; hence $c\mapsto W_k(A^\perp+c)$ is periodic with period $P$ in each coordinate. If $F\in\cA_\varepsilon$ is bounded, then $\int F\,e^{-S^{\mathrm m}_b+\sum_x\varepsilon^3k(x)\abs{\varphi(x)}^2}\prod_xd^2\varphi(x)=\sum_nw_{k,F}(n)\,e^{ie\varepsilon\sum_bn_bA_b}$, with complex weights $w_{k,F}(n)$ that depend on $A$ only through $dA$ and satisfy $\abs{w_{k,F}(n)}\le\norm F_\infty\,w_k(n)$; in particular this integral is $P$-periodic in each $c_\mu$.
\end{lemma}
The representation~\eqref{eq:C-current} is the current representation of \cite[Thm.~4.1(2), p.~262]{BFS1}, written here on the lattice torus.
\begin{proof}
In polar coordinates $\varphi(x)=r_xe^{i\vartheta_x}$, $\abs{e^{ie\varepsilon A_{x,\mu}}\varphi(y)-\varphi(x)}^2=r_x^2+r_y^2-2r_xr_y\cos(\vartheta_y-\vartheta_x+e\varepsilon A_{x,\mu})$ for $y=x+\varepsilon e_\mu$. Every site is an endpoint of six bonds, hence
\[
e^{-S^{\mathrm m}_b+\sum\varepsilon^3kr^2}=\prod_xv_x(r_x)\prod_{(x,\mu)}\exp\big(\varepsilon r_xr_y\cos(\vartheta_y-\vartheta_x+e\varepsilon A_{x,\mu})\big),\qquad v_x(r):=e^{-3\varepsilon r^2-\varepsilon^3(\frac12b\,r^2+\lambda r^4-k(x)r^2)} .
\]
Expand each bond factor as $e^{t\cos\psi}=\sum_{\nu\in\Z}I_\nu(t)e^{i\nu\psi}$, where $I_\nu=I_{-\nu}\ge0$ on $[0,\infty)$ are the modified Bessel functions and $\sum_\nu I_\nu(t)=e^t$. The resulting multiple series converges absolutely after integration, because $\prod_{(x,\mu)}e^{\varepsilon r_xr_y}\le\prod_xe^{3\varepsilon r_x^2}$ (by $r_xr_y\le\frac12(r_x^2+r_y^2)$) and $\int_0^\infty e^{-\varepsilon^3(\frac12b\,r^2+\lambda r^4-k(x)r^2)}r\,dr<\infty$ for $\lambda>0$; the same holds with a bounded factor $F\in\cA_\varepsilon$ in the integrand, and $F$ does not depend on the angles $\vartheta_x$. So sum and integral may be interchanged. Integrating over the angles, $\int\prod_{(x,\mu)}e^{in_{x,\mu}(\vartheta_{x+\varepsilon e_\mu}-\vartheta_x)}\prod_xd\vartheta_x$ equals $(2\pi)^{\abs{T_{j,\varepsilon}}}$ if $n$ is divergence-free and $0$ otherwise. This gives~\eqref{eq:C-current} with
\[
w_k(n)=(2\pi)^{\abs{T_{j,\varepsilon}}}\int\prod_xv_x(r_x)\,r_x\,dr_x\prod_{(x,\mu)}I_{n_{x,\mu}}(\varepsilon r_xr_{x+\varepsilon e_\mu})\ \ge0 ,
\]
and the same formula with the factor $F$ under the integral gives $w_{k,F}(n)$; since $F$ depends on $A$ only through $dA$, so does $w_{k,F}(n)$, and $\abs{w_{k,F}(n)}\le\norm F_\infty w_k(n)$. Moreover $w_k(-n)=w_k(n)$ because $I_{-\nu}=I_\nu$, and $\sum_nw_k(n)=W_k(0)$ by~\eqref{eq:C-current} at $A=0$. Positivity of $W_k(A)$ is clear from its definition, and $\abs{W_k(A)}\le\sum_nw_k(n)$. For a divergence-free current $n$, the difference of the fluxes through the planes $\{x_\mu=s+\varepsilon\}$ and $\{x_\mu=s\}$ is $-\sum_{x:\,x_\mu=s+\varepsilon}\sum_{\nu\ne\mu}(n_{x,\nu}-n_{x-\varepsilon e_\nu,\nu})$, which vanishes because the inner sums telescope over the periodic plane; so $\omega_\mu(n)$ does not depend on $s$. There are $\ell_j/\varepsilon$ such planes, so $\sum_{b\parallel e_\mu}n_b=(\ell_j/\varepsilon)\,\omega_\mu(n)$ and $e\varepsilon\sum_bn_bc_{\mu(b)}=e\ell_j\sum_\mu c_\mu\omega_\mu(n)$, which is $P$-periodic in each $c_\mu$. The weights $w_{k,F}(n)$ depend on $A$ through $dA=dA^\perp$ only, so the same conclusion holds for the integral with the factor $F$.
\end{proof}

\begin{lemma}[the constant mode]\label{lem:holonomy}
For every $F\in\cA_\varepsilon$ (bounded or non-negative), $c\mapsto\Psi_F(A^\perp+c)$ is periodic with period $P$ in each coordinate. If $F$ is bounded, $\abs{\Psi_F(A)}\le\norm{F}_\infty W_0(0)$.
\end{lemma}
\begin{proof}
Let $\kappa:=eP=\pm2\pi/\ell_j$ (the sign of $e$) and $\varphi'(x):=e^{-i\kappa x_\mu}\varphi(x)$, with $x_\mu\in[0,\ell_j)$; since $e^{-i\kappa\ell_j}=1$ this is a well-defined function on $T_{j,\varepsilon}$. For the bonds in direction $\mu$, $e^{ie\varepsilon(A_{x,\mu}+P)}\varphi'(x+\varepsilon e_\mu)-\varphi'(x)=e^{-i\kappa x_\mu}\big(e^{ie\varepsilon A_{x,\mu}}\varphi(x+\varepsilon e_\mu)-\varphi(x)\big)$, including the bonds that wrap around the torus; the bonds in the other directions and the single-site terms are unchanged. Hence $S^{\mathrm m}_b(A+Pe_\mu,\varphi')=S^{\mathrm m}_b(A,\varphi)$, where $A+Pe_\mu$ adds $P$ to every bond in direction $\mu$. Since $F$ depends on $\varphi$ only through $\abs{\varphi}$, and on $A$ only through $dA$, $F(A+Pe_\mu,\varphi')=F(A,\varphi)$; the change of variables $\varphi\mapsto\varphi'$ preserves Lebesgue measure, which gives the periodicity. The bound follows from $\abs{\Psi_F(A)}\le\norm{F}_\infty W_0(A)$ and Lemma~\ref{lem:postype}.
\end{proof}

\begin{definition}[the states]\label{def:states}\label{def:C-flat}
Let $\gamma_{\mu_0}$ be the centred Gaussian measure on bond functions with covariance $C_{\mu_0}:=(-\Delta_\varepsilon+\mu_0^2)^{-1}$ for the pairing $\ip{\cdot}{\cdot}_\varepsilon$, that is, $\E[\ip fA_\varepsilon\ip gA_\varepsilon]=\ip f{C_{\mu_0}g}_\varepsilon$. Let $\gamma^\fl$ be the product of the centred Gaussian measure with covariance $(-\Delta_\varepsilon)^{-1}\Pi^\perp$ on the range of $\Pi^\perp$ and the uniform probability measure on the period cell $[0,P)^3$ for $c$. For $F\in\cA_\varepsilon$ with $\E_{\gamma}[\Psi_{\abs F}]<\infty$, $\gamma\in\{\gamma_{\mu_0},\gamma^\fl\}$, put
\[
\langle F\rangle_{\varepsilon,j,\mu_0}:=\frac{\E_{\gamma_{\mu_0}}[\Psi_F]}{\E_{\gamma_{\mu_0}}[W_0]}\quad(0<\mu_0\le\mu_*),\qquad
\langle F\rangle_{\varepsilon,j,\fl}:=\frac{\E_{\gamma^\fl}[\Psi_F]}{\E_{\gamma^\fl}[W_0]} .
\]
\end{definition}
Under $\gamma_{\mu_0}$, $A^\perp$ and $c$ are independent and $c$ is centred Gaussian with covariance $(\mu_0^2\abs{T_j})^{-1}\mathbf 1_3$ (Lemma~\ref{app:G1:zeromode}). The state $\langle\cdot\rangle_{\varepsilon,j,\mu_0}$ is the expectation in the probability measure $Z^{-1}e^{-S_{\mu_0,b}}\prod_bdA_b\prod_xd^2\varphi(x)$ on $\R^{\text{bonds}}\times\C^{\text{sites}}$, restricted to $\cA_\varepsilon$ (Lemma~\ref{app:G1:gibbs}); it is King's torus model at photon mass $\mu_0$ with the counterterm of the reference photon mass. The denominators are positive and finite by Lemma~\ref{lem:postype}. In the state $\fl$ the non-constant photon modes are massless, and the constant mode $c$ (the holonomy) is uniformly distributed on a period cell. By Lemma~\ref{lem:holonomy} the integrands $\Psi_F$ and $W_0$ are $P$-periodic in $c$, so $\langle F\rangle_{\varepsilon,j,\fl}$ does not depend on the choice of the period cell; $F\mapsto\langle F\rangle_{\varepsilon,j,\fl}$ is the expectation of a probability measure on the $\sigma$-algebra generated by $\cA_\varepsilon$, namely of the normalised measure $\gamma^\fl(dA)\,e^{-S^{\mathrm m}_b(A,\varphi)}\prod_xd^2\varphi(x)$.

\begin{lemma}[$\fl$ is the massless limit at fixed lattice]\label{lem:flat}
Fix $\varepsilon$ and $j$. For every bounded $F\in\cA_\varepsilon$, $\lim_{\mu_0\downarrow0}\langle F\rangle_{\varepsilon,j,\mu_0}=\langle F\rangle_{\varepsilon,j,\fl}$. The same holds for every $F\in\cA_\varepsilon$ with $\sup_{0<\mu_0\le\mu_*}\langle\abs F^{1+\delta}\rangle_{\varepsilon,j,\mu_0}<\infty$ for some $\delta>0$.
\end{lemma}
\begin{proof}
We write $\langle\cdot\rangle_{\mu_0}$ and $\langle\cdot\rangle_\fl$ for the states at fixed $\varepsilon$, $j$.
(i) The law of $A^\perp$ under $\gamma_{\mu_0}$ is the centred Gaussian with covariance $(-\Delta_\varepsilon+\mu_0^2)^{-1}\Pi^\perp$ on the finite-dimensional range of $\Pi^\perp$, on which $-\Delta_\varepsilon$ is positive definite; as $\mu_0\downarrow0$ it converges in total variation to the centred Gaussian with covariance $(-\Delta_\varepsilon)^{-1}\Pi^\perp$ (the densities converge pointwise, and Scheff\'e's lemma applies). (ii) By Lemmas~\ref{lem:postype} and~\ref{lem:holonomy}, $\Psi_F(A^\perp+c)$ and $W_0(A^\perp+c)$ are bounded and $P$-periodic in each $c_\mu$, so their $\gamma_{\mu_0}$-expectations only involve the law of $c$ modulo $P$, the wrapped Gaussian of variance $s=(\mu_0^2\abs{T_j})^{-1}$ in each coordinate. By Poisson summation its density on $[0,P)$ is $P^{-1}\big(1+2\sum_{n\ge1}e^{-2\pi^2n^2s/P^2}\cos(2\pi nx/P)\big)$, which converges uniformly to $P^{-1}$ as $s\to\infty$. (iii) $A^\perp$ and $c$ are independent under every $\gamma_{\mu_0}$ and under $\gamma^\fl$, so the joint laws converge in total variation, and numerator and denominator of $\langle F\rangle_{\mu_0}$ converge to those of $\langle F\rangle_\fl$; the limiting denominator is positive because $W_0>0$. For unbounded $F$ put $F_n:=F\mathbf 1_{\{\abs F\le n\}}\in\cA_\varepsilon$. Then $\abs{\langle F\rangle_{\mu_0}-\langle F_n\rangle_{\mu_0}}\le n^{-\delta}\sup_{\mu_0}\langle\abs F^{1+\delta}\rangle_{\mu_0}$ uniformly in $\mu_0$. The same bound holds for $\fl$: by the bounded case $\langle\min(\abs F,M)^{1+\delta}\rangle_\fl\le\sup_{\mu_0}\langle\abs F^{1+\delta}\rangle_{\mu_0}$ for every $M$, and monotone convergence applies. The claim follows from the bounded case.
\end{proof}

\subsection{Observables and generating functionals}\label{ssec:observables}

\paragraph{Wick constants.} Put $\hat p^2:=\sum_\mu4\varepsilon^{-2}\sin^2(\varepsilon p_\mu/2)$ and
\[
W_\varepsilon:=2\int_{[-\pi/\varepsilon,\pi/\varepsilon)^3}\frac{d^3p}{(2\pi)^3}\,\frac1{\hat p^2+m^2},
\]
the free-field value of $\E\abs{\varphi(0)}^2$ on $\varepsilon\Z^3$ at mass $m$. Let $X_\varepsilon(m)$ be the coefficient of $e^2$ in the perturbative expansion of $\langle\abs{\varphi(0)}^2\rangle$ for the model on $\varepsilon\Z^3$ at scalar mass $m$ and photon mass $\mu_*$: the seagull and exchange graphs, with the order-$e^2$ part of the counterterm on $\varepsilon\Z^3$ inserted, that is, King's \cite[(3.33)]{K1} with $\Omega=\varepsilon\Z^3$. (It is not the torus counterterm of Definition~\ref{def:counterterm}, whose photon zero mode $e^2/(\mu_*^2\abs{T})$ would make $z_\varepsilon$ depend on the torus.) Put $z_\varepsilon:=e^2X_\varepsilon(m)$. The same $W_\varepsilon$ and $z_\varepsilon$ are used on every torus, at every photon mass and in the state $\fl$.

\begin{definition}[test functions and sampling]\label{def:C-sampling}
Scalar test functions are real $h\in C_c(\R^3)$; curvature test functions are real 2-forms $g=\sum_{\mu<\nu}g_{\mu\nu}\,dx_\mu\wedge dx_\nu$ with $g_{\mu\nu}\in C_c(\R^3)$. Let $\supp h,\supp g\subset B_R$ and $\ell_j>2R+2$, and identify $T_j$ with the cell $[-\ell_j/2,\ell_j/2)^3$. The test functions are placed on $T_{j,\varepsilon}$ by sampling: $h$ by its values $h(x)$ at the sites of the cell, and $g$ by the value $g(p):=g_{\mu\nu}\big(x+\tfrac\varepsilon2(e_\mu+e_\nu)\big)$ at the centre of the plaquette $p=(x;\mu,\nu)$; both are extended periodically. The lattice norms are those of the sampled functions, $\norm h_{L^1,\varepsilon}:=\sum_x\varepsilon^3\abs{h(x)}$, $N_\varepsilon(h):=\max(\norm h_{L^1,\varepsilon},\norm h_\infty)$ and $\norm g^2_{2,\varepsilon}:=\sum_p\varepsilon^3g(p)^2$; the continuum norms are $\norm h_1$, $\norm h_\infty$, $N(h):=\max(\norm h_1,\norm h_\infty)$ and $\norm g_2^2:=\sum_{\mu<\nu}\norm{g_{\mu\nu}}^2_{L^2}$.
\end{definition}

\paragraph{Fields.} The lattice fields are
\[
F(g):=\ip{g}{dA}_\varepsilon=\sum_p\varepsilon^3g(p)\,(dA)(p),\qquad \XR(h)=\wick{\abs\varphi^2}_{\mathrm R}(h):=\sum_x\varepsilon^3h(x)\big(\abs{\varphi(x)}^2-W_\varepsilon-z_\varepsilon\big),
\]
for real lattice 2-forms $g$ and real lattice functions $h$, in particular for sampled test functions. Both are in $\cA_\varepsilon$, and so are their polynomials and exponentials. Since $\ip{d^*g}{\Pi_0A}_\varepsilon=\ip{g}{d\,\Pi_0A}_\varepsilon=0$, $F(g)=\ip{d^*g}{A}_\varepsilon$ does not depend on the constant mode $c$. For a real lattice function $k$,
\[
\int e^{-S^{\mathrm m}_b(A,\varphi)+\XR(k)}\prod_xd^2\varphi(x)=e^{-\sum_x\varepsilon^3k(x)(W_\varepsilon+z_\varepsilon)}\,W_k(A),
\]
so Lemma~\ref{lem:postype} applies to this weight, with $w_k(n)$ multiplied by a positive constant. Wilson loops are not among the observables: in three dimensions the sharp loop is not a tempered field, and smeared loops are functions of $F(g)$ by Stokes' theorem.

\paragraph{Generating functionals and the tube.} For real test functions $g_1,\dots,g_n$, $h_1,\dots,h_p$ and $\zeta=(z,w)\in\C^{n+p}$ put
\[
\Theta(\zeta):=\exp\Big[\sum_iz_iF(g_i)+\sum_kw_k\XR(h_k)\Big],\qquad G_{\varepsilon,j,\mu_0}(\zeta):=\langle\Theta(\zeta)\rangle_{\varepsilon,j,\mu_0},\quad \mu_0\in(0,\mu_*]\cup\{\fl\}.
\]
At fixed $\varepsilon$ and $j$ this is finite for every $\zeta$: $\abs{\Theta(\zeta)}=\Theta(\rRe\zeta)$, the matter integral of $\Theta(\rRe\zeta)$ is at most a constant times $e^{F(g)}W_k(0)$ with $g=\sum_i(\rRe z_i)g_i$ and $k=\sum_k(\rRe w_k)h_k$ (Lemma~\ref{lem:postype}), and $e^{F(g)}=e^{\ip{d^*g}{A}_\varepsilon}$ is integrable for the Gaussian measures of Definition~\ref{def:states}; since on a compact set of $\zeta$ the function $\Theta(\rRe\zeta)$ is bounded, by convexity, by the sum of its values at the corners of a box, dominated convergence shows that $G_{\varepsilon,j,\mu_0}$ is entire. Let $k(\zeta):=\sum_k(\rRe w_k)h_k=k_+-k_-$, with $k_\pm\ge0$ its positive and negative parts. For $\sigma_+,H>0$ the tube is
\[
D_{\sigma_+,H}:=\big\{\zeta\in\C^{n+p}:\ \norm{k_+(\zeta)}_\infty<\sigma_+/2,\ \norm{k_-(\zeta)}_\infty<H/2\big\},
\]
with no restriction on $z$. Its base, the set of the admissible $\rRe w$, is open and convex, so $D_{\sigma_+,H}$ is connected. The lattice Schwinger functions are
\[
S^{\varepsilon,j,\mu_0}_{n,p}(g;h):=\Big\langle\prod_{i=1}^nF(g_i)\prod_{k=1}^p\XR(h_k)\Big\rangle_{\varepsilon,j,\mu_0}=\partial_{z_1}\cdots\partial_{z_n}\partial_{w_1}\cdots\partial_{w_p}G_{\varepsilon,j,\mu_0}(0).
\]

\paragraph{The class $\cD$.} A generating functional of the class $\cD$ is $\Theta=\exp[F(g)-\XR(q)]$ with $g$ a real 2-form and $q\in C_c(\R^3)$, $q\ge0$, $\norm q_\infty<H/2$; it is the value of $\Theta(\zeta)$ at a real point $\zeta$ of the tube. At fixed lattice and torus, $\langle\Theta\rangle_{\varepsilon,j,\mu_0}$ is non-decreasing in the photon covariance for $\Theta\in\cD$ \cite[Cor.~6.3(1), p.~281]{BFS1} (on the lattice torus: Lemma~\ref{lem:C-cov}), which is the source of the monotonicity in Theorem~\ref{thm:main}(iv). We write $E^\cD$ for the set of real points $\zeta$ of the tube at which $\Theta(\zeta)$ belongs to $\cD$, that is, $E^\cD=\{\zeta\in D_{\sigma_+,H}\cap\R^{n+p}:\ \sum_kw_kh_k\le0\}$.

\subsection{Symmetries}\label{ssec:symmetries}

The cubic group $\Oh\subset O(3)$ is the group of the $48$ signed permutation matrices. The Euclidean group acts on test functions by
\[
(\tau_ah)(x):=h(x-a),\qquad (R_*h)(x):=h(R^{-1}x),\qquad (R_*g)_x(v,w):=g_{R^{-1}x}(R^{-1}v,R^{-1}w),
\]
for $a\in\R^3$ and $R\in O(3)$,
that is, $h$ transforms as a scalar and $g$ as a 2-form (push-forward, signs included; for an improper $R$ no extra sign appears). For $\nu\in\{1,2,3\}$ and $c\in\R$ let $\rho_{\nu,c}(x):=x-2(x_\nu-c)e_\nu$ be the reflection in the plane $\{x_\nu=c\}$; on test functions it acts by push-forward, so that $(\rho_{\nu,c})_*g$ changes the sign of the components $g_{\mu\nu}$, $g_{\nu\mu}$ and reflects the argument. On the lattice, $\Oh$ acts about the site $0$ on $\varphi$ as on scalars and on $A$ as on 1-forms ($A_{x,\mu}\mapsto\pm A$ on the image bond, with the sign $-1$ when the image bond is reversed), and charge conjugation acts by $(A,\varphi)\mapsto(-A,\bar\varphi)$. Invariance of the lattice states under these maps is Lemma~\ref{lem:C-sym} (Section~\ref{sec:C-lattice}).

\subsection{Schwinger families and the axioms}\label{ssec:axioms}

A \emph{Schwinger family} is a collection $S=(S_{n,p})_{n,p\ge0}$, $S_{0,0}=1$, of real multilinear functionals $S_{n,p}(g_1,\dots,g_n;h_1,\dots,h_p)$ of real curvature and scalar test functions in $C_c$, symmetric within each type. For a finite family of test functions, its generating functional on a tube $D$ is a holomorphic function $G^S$ on $D$ whose Taylor coefficients at $0$ are the $S_{n,p}$; when it exists we write $S(\Theta(\zeta)):=G^S(\zeta)$. For a polynomial $Q=\sum c_\alpha\prod F(g_{\alpha,i})\prod\XR(h_{\alpha,k})$ in the smeared fields with complex coefficients, $S(Q)$ is defined by linearity. We use the following forms of the axioms.

\begin{definition}[OS0--OS3]\label{def:OS}
Let $S$ be a Schwinger family.
\begin{enumerate}
\item[OS0] (regularity) There are $\sigma_+,H>0$ and $C<\infty$ such that for every finite family of real test functions $G^S$ exists on $D_{\sigma_+,H}$ and $\abs{G^S(\zeta)}\le\exp\big(\tfrac12\norm{\sum_i(\rRe z_i)g_i}_2^2+C\norm{k(\zeta)}_{L^1}\big)$ (\emph{tube form}); and, with $\abs f_4:=\sup_{x\in\R^3,\abs\alpha\le4}\abs{(1+\abs x^2)^2D^\alpha f(x)}$, every $S_{n,p}$ evaluated on the components of the test functions satisfies $\abs{S_{n,p}}\le\sigma_{n+p}\prod_i\abs{g_i}_4\prod_k\abs{h_k}_4$ with $\sigma_N\le\alpha_\beta(N!)^\beta$ for every $\beta>1$, i.e.\ the condition E0$''$ of \cite[(4.2), p.~287]{OS75} with $s=4$.
\item[OS1$(\mathcal G)$] (invariance under a subgroup $\mathcal G$ of the Euclidean group $\R^3\rtimes O(3)$) $S_{n,p}(\gamma_*g_1,\dots;\gamma_*h_1,\dots)=S_{n,p}(g_1,\dots;h_1,\dots)$ for all $\gamma\in\mathcal G$, with the actions of Section~\ref{ssec:symmetries}.
\item[OS2] (reflection positivity about coordinate planes) For every $\nu\in\{1,2,3\}$ and $c\in\R$, with $\theta:=\rho_{\nu,c}$ acting on test functions as in Section~\ref{ssec:symmetries} and antilinearly on coefficients: $S(\theta(Q)\,Q)\ge0$ for every polynomial $Q$ in the fields smeared with test functions supported in $\{x_\nu>c\}$; and $\sum_{i,l}\bar c_ic_l\,S(\theta(\Theta_i)\Theta_l)\ge0$ for all $c\in\C^N$ and all exponentials $\Theta_i=\Theta(\zeta_i)$ of such fields whose $\wick{\abs\varphi^2}$-amplitudes satisfy $\sum\abs{\rRe w}\norm h_\infty<\min(\sigma_+,H)/2$.
\item[OS3] (symmetry) $S_{n,p}$ is real for real test functions and symmetric under permutations within each field type; moreover $S_{n,p}=0$ for odd $n$.
\end{enumerate}
\end{definition}

The cluster property refers to a probability law of the fields. Let $X$ be a space of distributions on $\R^3$ with values in $\Lambda^2\R^3\oplus\R$ (the curvature and the Wick square), carrying the $\sigma$-algebra generated by the evaluations $\Phi\mapsto\Phi(f)$, $f\in C_c^\infty(\R^3;\Lambda^2\R^3\oplus\R)$, and invariant under the translations $(\tau_a\Phi)(f):=\Phi(f(\cdot+a))$, $a\in\R^3$. Let $\mu$ be a probability measure on $X$ invariant under every $\tau_a$, and let $U(a)G:=G\circ\tau_{-a}$ be the induced unitary representation of $\R^3$ on $L^2(\mu)$; thus $U(a)[\Phi(f)]=\Phi(f(\cdot-a))$.

\begin{definition}[the cluster property OS4]\label{def:OS4}
For $j\in\{1,2,3\}$ consider the two conditions:
\begin{enumerate}
\item[(OS4$_j$)] (ergodicity of the translations in the direction $e_j$) every $G\in L^2(\mu)$ with $U(te_j)G=G$ for all $t\in\R$ is $\mu$-almost surely constant;
\item[(OS4$'_j$)] (clustering along the axis $e_j$) $\E_\mu[G'\,U(te_j)G]\to\E_\mu[G']\,\E_\mu[G]$ as $\abs t\to\infty$, for all $G,G'\in L^2(\mu)$.
\end{enumerate}
The measure $\mu$ satisfies \emph{OS4} if (OS4$_j$) holds for $j=1,2,3$. A Schwinger family satisfies OS4 if it is the family of moments of exactly one translation-invariant probability measure $\mu$ on $X$, and this $\mu$ satisfies OS4.
\end{definition}
\begin{definition}[forms of the cluster axiom; the weak form of OS4]\label{add:OS4}
Let $X$, $\mu$ and $U$ be as in Definition~\ref{def:OS4}, and let $j\in\{1,2,3\}$ be the direction chosen as Euclidean time. In (i) assume moreover that $X$ and $\mu$ are invariant under the map $\theta_j$ of $X$ induced by the reflection $\rho_{j,0}$ in the plane $\{x_j=0\}$ (Section~\ref{ssec:symmetries}), and let $\mathcal P_j^+$ be the set of polynomials in the fields $\Phi(f)$, $f\in C_c^\infty(\{x_j>0\};\Lambda^2\R^3\oplus\R)$ real, that belong to $L^2(\mu)$.
\begin{enumerate}
\item[(i)] (cluster property, \cite[E4, p.~88]{OS73}) For all $F,G\in\mathcal P_j^+$ and every $a\in\R^3$ with $a\perp e_j$,
\[
\lim_{s\to\infty}\E_\mu\big[\overline{F\circ\theta_j}\;U(sa)G\big]=\E_\mu\big[\overline F\,\big]\,\E_\mu[G].
\]
\item[(ii)] (ergodicity of the time translations, \cite[(6.1.10), p.~92]{GJ}) Condition (OS4$_j$) of Definition~\ref{def:OS4} holds.
\item[(iii)] (weak form) $\mu$ is ergodic for the action of the translation group $\R^3$: every $G\in L^2(\mu)$ with $U(a)G=G$ for all $a\in\R^3$ is $\mu$-almost surely constant.
\end{enumerate}
We call (i) and (ii) the single-axis forms of OS4, and say that $\mu$ satisfies \emph{OS4 in the weak form} if (iii) holds.
\end{definition}

\begin{remark}\label{p:rem:OS4forms}
(a) With $e_j$ as Euclidean time, (OS4$_j$) is the form of the cluster axiom in \cite[Sect.~6.1]{GJ}, ergodicity of the time translations. For $k\ne j$, condition (OS4$'_k$) contains the cluster property E4 of \cite[p.~88]{OS73} for the spatial directions $a=\pm e_k$ and test functions in $C_c^\infty$: take $G'=\overline{F\circ\theta}$, with $\theta$ the map of $X$ induced by the reflection in $\{x_j=0\}$ (Section~\ref{ssec:symmetries}) and $F,G$ square-integrable polynomials in fields smeared with test functions supported in $\{x_j>0\}$. The axiom E4 asks for all spatial directions $a\perp e_j$.
(b) If $\mu$ is invariant under the coordinate permutations, for instance under $\Oh$ as the laws of the limit states of Theorem~\ref{thm:main} are, then (OS4$_1$), (OS4$_2$) and (OS4$_3$) are equivalent; in general they are three statements. With invariance under all rotations, the direction $e_j$ could be replaced by any unit vector.
(c) If $\mu$ satisfies the hypotheses (H1)--(H3) of Section~\ref{sec:OS4}, as the laws of the limit states do (Proposition~\ref{prop:discharge}), then for each $j$ the conditions (OS4$_j$) and (OS4$'_j$) are equivalent: by Proposition~\ref{prop:D:subgroup} with $J=\{j\}$, $\E_\mu[G'\,U(te_j)G]\to\E_\mu[G'\,P_jG]$ as $\abs t\to\infty$, where $P_j$ is the orthogonal projection onto the $U(te_j)$-invariant vectors; (OS4$_j$) says that $P_jG=\E_\mu G$ for every $G$, and the equivalence follows as in Lemma~\ref{lem:os4:red}.
(d) (OS4$_j$) for one $j$ implies that $\mu$ is ergodic for the action of $\R^3$, since an $\R^3$-invariant function is invariant under the translations along $e_j$. The converse fails: the law of Proposition~\ref{prop:axisobs} (realised on the present field space in Example~\ref{ex:D:axis}) satisfies (H1)--(H3), is $\Oh$-invariant and ergodic for $\R^3$, and its two-point function along $e_1$ has a positive limit, so (OS4$'_1$), and hence (OS4$_1$), fails. Theorem~\ref{thm:cluster} gives the weak form of Definition~\ref{add:OS4}, ergodicity under $\R^3$, for the ergodic components of the limit states, and Theorem~\ref{thm:os4main} gives (OS4$'_j$) for part of the observables; OS4 itself is discussed in Section~\ref{sec:disc}.
(e) For $\mu$ satisfying (H1)--(H3), Proposition~\ref{add:OS4rel} in Section~\ref{sec:OS4} shows that (OS4$_j$) holds if and only if $\Ker H_j=\C\Omega_j$ in the reconstruction of Proposition~\ref{prop:D:recon} in direction $j$; that the weak form (iii) of Definition~\ref{add:OS4} is equivalent to the uniqueness of the joint vacuum in one (equivalently, every) coordinate reconstruction, to clustering under diagonal separation and to clustering in the mean over cubes; and that, for the decomposition of Theorem~\ref{thm:cluster}, almost every component satisfies (OS4$_j$) if and only if $P_jG=\E[G\mid\cI]$ for every $G\in L^2(\mu)$, where $P_j$ is the orthogonal projection onto the $U(te_j)$-invariant vectors.
\end{remark}

\section{Main results}\label{sec:main}

The objects in the statements below are defined in Section~\ref{sec:model}. Throughout this section $m>0$, $0<\lambda<0.1m$, $0<e^2<0.1m$ and $\mu_*:=m$; the lattice states $\langle\cdot\rangle_{\varepsilon,j,\mu_0}$ are those of Definition~\ref{def:states} in the torus counterterm scheme of Definition~\ref{def:counterterm}.

\subsection{Existence and OS0--OS3 at blocking factor 2}\label{ssec:mainL2}

\begin{theorem}[existence and OS0--OS3, $L=2$]\label{thm:main}
Let $m>0$, $0<\lambda<0.1m$, $0<e^2<0.1m$, and $\mu_*:=m$. There are $M_b,n_0\in\N$, $\varepsilon_0=2^{-n_0}$, $\ell_0\in2M_b\N$ with $\ell_0\ge183/m$, $K_0\in\N$, $\sigma_+>0$ and $C<\infty$ such that the following holds for the states $\langle\cdot\rangle_{\varepsilon,j,\mu_0}$ of Definition~\ref{def:states} with $L=2$, $\varepsilon=2^{-K}\varepsilon_0$, $K\ge K_0$, $j\ge1$, and with $H:=m^2$.
\begin{enumerate}
\item[(i)] (uniform tube bound) For every finite family of real test functions, every $\zeta\in D_{\sigma_+,H}$, every $\mu_0\in(0,\mu_*]\cup\{\fl\}$, every $K\ge K_0$ and every $j$ with $\ell_j$ larger than twice the radius of the supports plus $2$,
\[
\abs{G_{\varepsilon,j,\mu_0}(\zeta)}\ \le\ G_{\varepsilon,j,\mu_0}(\rRe\zeta)\ \le\ \exp\Big(\tfrac12\Big\|\sum_i(\rRe z_i)g_i\Big\|_{2,\varepsilon}^2+C\norm{k(\zeta)}_{L^1,\varepsilon}\Big).
\]
\item[(ii)] (limits) Let $Q_0\subset(0,\mu_*]$ be countable and accumulate at $0$. Every sequence $(K_l,j_l)$ with $K_l\to\infty$ and $j_l\to\infty$ has a subsequence along which, for every $\mu_0\in Q_0$ and for $\mu_0=\fl$, the lattice Schwinger functions $S^{\varepsilon_{K_l},j_l,\mu_0}_{n,p}$ converge for all $n,p$ and all test functions in $C_c$, and the generating functionals converge locally uniformly on $D_{\sigma_+,H}$. Call the limits $S^{\mu_0}$ ($\mu_0\in Q_0$) and $S^\fl$. Along this subsequence, $S^0:=\lim_{\mu_0\downarrow0,\ \mu_0\in Q_0}S^{\mu_0}$ exists, with the same limit along every sequence in $Q_0$ decreasing to $0$, in the sense of Schwinger functions and locally uniformly on $D_{\sigma_+,H}$. Each $S\in\{S^{\mu_0},S^\fl,S^0\}$ is, on test functions in $C_c^\infty$, the family of moments of a Borel probability measure on $H^{-s'}_{\mathrm{loc}}(\R^3;\Lambda^2\R^3\oplus\R)$ for every $s'>3$, whose characteristic functional is $G^S$ at imaginary arguments.
\item[(iii)] (axioms) Each $S\in\{S^{\mu_0},S^\fl,S^0\}$ satisfies OS0 with the constants $\sigma_+$, $H=m^2$, $C$ of (i); OS1$(\R^3\rtimes\Oh)$, i.e.\ invariance under all translations of $\R^3$ and under $\Oh$ about every point of $\R^3$, $g$ transforming as a 2-form and $h$ as a scalar; OS2 about every coordinate plane; and OS3. Moreover
\[
\abs{S_{n,p}(g;h)}\le \mathrm{e}^{C\sigma_0}\,\mathrm{e}^n(n!)^{1/2}(\mathrm{e}/\sigma_0)^p\,p!\prod_i\norm{g_i}_2\prod_kN(h_k),\qquad\sigma_0:=\min(\sigma_+,H)/2 .
\]
\item[(iv)] (monotonicity) For $\Theta\in\cD$, $\mu_0\mapsto S^{\mu_0}(\Theta)$ is non-increasing on $Q_0$, $S^{\mu_0}(\Theta)\le S^\fl(\Theta)$, and $S^0(\Theta)\le S^\fl(\Theta)$.
\end{enumerate}
\end{theorem}
The constants of (i) are $\sigma_+=g/2$, where $g>0$ is the convexity constant of Proposition~\ref{prop:ch-Pplus} (not a test function), and $C=\mathsf K=\max(\kappa_+,2S^B_\fl)$, with $\kappa_+$ of Proposition~\ref{prop:ch-Pplus} and $S^B_\fl$ of Proposition~\ref{prop:ch-Tflat}, at the auxiliary parameters $(\eta,m_1,m_2)=(\frac12,3m,4m)$ (Theorem~\ref{thm:ch-T9} and Lemma~\ref{app:G1:window-apply}).

\begin{remark}[cubic symmetry, not rotations]\label{rem:cubic}
The group in OS1 is $\R^3\rtimes\Oh$: the translations of $\R^3$ and the $48$ signed permutations about every point. On the lattice, translations by $\varepsilon\Z^3$ and the cubic group about a site are exact symmetries of every state on a cubic torus; they pass to the limits, and continuity extends the translations to $\R^3$. Rotations outside $\Oh$ are not symmetries of any regularised state, and Theorem~\ref{thm:main} makes no claim about them; Section~\ref{ssec:introrot} discusses them. In the reconstruction of Section~\ref{sec:recon} only the subgroup $D_4$ of $\Oh$ fixing the reconstruction direction acts on the Hilbert space, together with the translations orthogonal to it.
\end{remark}

\begin{remark}[order of limits]\label{rem:order}
(a) In (ii) the spacing and the torus go to their limits jointly. Every iterated limit, first $\varepsilon\to0$ at fixed $T_j$ and then $j\to\infty$, at fixed $\mu_0$, is a limit of the joint family in the sense of (ii), and every limit $\mu_k\downarrow0$ of such iterated limits is a limit point of the joint family $(\varepsilon,T_j,\mu_0)\to(0,\R^3,0)$; so (iii) holds in every order (Proposition~\ref{prop:C-order}).
(b) The photon mass is removed last, along the countable set $Q_0$ and along one subsequence of $(\varepsilon,T_j)$ common to all $\mu_0\in Q_0$. The limit $S^0$ is monotone in $\mu_0$ on $\cD$ and does not depend on the sequence in $Q_0$, but it may depend on the common subsequence; we do not prove that it is independent of it. Continuity of $\mu_0\mapsto S^{\mu_0}$ is not claimed.
(c) The limit $S^\fl$ is taken in the other order: $\mu_0\downarrow0$ first, at fixed lattice (Lemma~\ref{lem:flat}), then the lattice limit. Part (iv) is the only relation between the two orders that we prove.
\end{remark}

\begin{remark}[the regulator data]\label{rem:scope}
(a) The constants $\varepsilon_0$, $\ell_0$ and $K_0$ come from Ba{\l}aban's renormalisation group with blocking factor $L=2$, small-field parameter $p=6$, perturbative order $\nbar=13$ with order cap $\Nbar=169$ and localisation radius $R\ge R^{**}$, and from the conditions on these data listed in Appendix~\ref{app:B}. The data are chosen in stages, in an order in which each condition depends only on data chosen at the same or at earlier stages, once the constants that depend on the torus side are evaluated at the torus floor (Lemma~\ref{app:G1:acyc}). The bound $\ell_0\ge183/m$ is the torus floor $128/M_1=182.857\ldots/m$ at the smallest scalar mass $M_1=0.7m$ used in the proof, rounded up (Lemma~\ref{app:G1:floor}).
(b) Lemma~\ref{lem:window} shows that all these conditions can be met at every $(e,\lambda,m)$ of Theorem~\ref{thm:main}, with the auxiliary parameters $(\eta,m_1,m_2)=(\frac12,3m,4m)$, which lie in the window~\eqref{eq:window} for every such $(e,\lambda,m)$ (Lemma~\ref{app:G1:window-apply}). An explicit admissible point is described in Section~\ref{p:appA:example}; nothing depends on it.
(c) Among the conditions are the condition (E.edge) on $\varepsilon_0$, produced by the proof of Lemma~\ref{lem:lem24}, and the torus-size condition (T.per), used in Lemma~\ref{lem:IMG}.
(d) Four ingredients that are first established at the anchor masses $m_1,m_2$ are used at all masses $M\in(0,m_2]$; Lemma~\ref{lem:B:massrange} supplies the extensions.
\end{remark}

\subsection{Blocking factor 3}\label{ssec:mainL3}

For blocking factor $L=3$ the lattices and tori of Section~\ref{ssec:tori} are $\varepsilon=3^{-K}\varepsilon_0$, $\varepsilon_0=3^{-n_0}$, and $\ell_j=j!\,\ell_0$ with $\ell_0\in18M_b\N$; Lemma~\ref{lem:sizes} holds with $c_3=18$. The states of Definition~\ref{def:states} are defined on these tori in the same way.

\begin{theorem}[blocking factor $3$, under a uniform tube bound]\label{thm:mainL3}
Let $m>0$, $0<\lambda<0.1m$, $0<e^2<0.1m$, $\mu_*:=m$ and $H:=m^2$, and consider the states $\langle\cdot\rangle_{\varepsilon,j,\mu_0}$ of Definition~\ref{def:states} with $L=3$, $\varepsilon=3^{-K}\varepsilon_0$, $\varepsilon_0=3^{-n_0}$, $\ell_0\in18M_b\N$ and $\ell_0\ge183/m$. Assume that for some $M_b,n_0\in\N$, some such $\ell_0$, some $K_0\in\N$, $\sigma_+>0$ and $C<\infty$ the following lattice tube bound holds: for every $K\ge K_0$, every $j\ge1$, every $\mu_0\in(0,\mu_*]\cup\{\fl\}$, all real lattice test functions $\bg_1,\dots,\bg_n$, $h_1,\dots,h_p$ on $T_{j,\varepsilon}$ and every $\zeta$ in the lattice tube $\mathcal D_{\sigma_+,H}$ of Section~\ref{ssec:ch-T9},
\begin{equation}\label{p:eq:T3}
\abs{G_{\varepsilon,T_j,\mu_0}(\zeta)}\le G_{\varepsilon,T_j,\mu_0}(\rRe\zeta)\le\exp\Big(\tfrac12\Big\|\sum_i(\rRe z_i)\bg_i\Big\|_{2,\varepsilon}^2+C\norm{k(\zeta)}_{L^1,\varepsilon}\Big).
\end{equation}
Then statements (ii)--(iv) of Theorem~\ref{thm:main} hold for this family, with these constants $\sigma_+$, $C$ and $H=m^2$.
\end{theorem}
\begin{proof}
Sections~\ref{sec:C-lattice}--\ref{sec:C-OS} derive (ii)--(iv) of Theorem~\ref{thm:main} from the lattice tube bound of Theorem~\ref{thm:ch-T9} and from statements on one lattice torus at a time. Apart from the tube bound they use only the properties (P1)--(P4) of Section~\ref{ssec:C-standing}. Of these, (P1)--(P3) do not involve the blocking factor, and (P4), the nesting of the lattices and the evenness of the number of sites per side, is Lemma~\ref{lem:sizes} with $c_3=18$. They also use the density in $\R^3$ of $\bigcup_K\varepsilon_K\Z^3$, here $\Z[1/3]^3\varepsilon_0$. With \eqref{p:eq:T3} in place of Theorem~\ref{thm:ch-T9}, the same proofs apply.
\end{proof}
\noindent Theorems~\ref{thm:cluster}, \ref{thm:os4main}(a),(c) and~\ref{thm:NTMmain} below, and the results of Sections~\ref{sec:OS4}--\ref{sec:massless} on which they rest, hold for the limit states of Theorem~\ref{thm:mainL3}, with the same proofs. Theorem~\ref{thm:os4main}(b) and Proposition~\ref{prop:os4:latent} are not asserted for these states: their proofs use Propositions~\ref{prop:ch-T7} and~\ref{prop:ch-Tflat}, not only for the bound $\Delta_S\le\kappa_++2S^B_\fl$ but already for the existence of the limit pressures, on which the definition of $\Delta_S$ rests. We expect that \eqref{p:eq:T3} can be obtained by carrying out Sections~\ref{sec:UV}--\ref{sec:chain} and Appendix~\ref{app:B} with blocking factor $3$; this is not done here. The interest of blocking factor $3$ lies in the rotations (Section~\ref{ssec:introrot}).

\subsection{Weak clustering}\label{ssec:clustering}

Let $\cV:=\Lambda^2\R^3\oplus\R\cong\R^4$ (the curvature and the Wick square). For $r>0$ let $\mathcal H^{-r}:=\mathcal H^{-r}(\R^3;\cV)$ be the Hermite--Sobolev space of order $-r$ of $\cV$-valued distributions: with $\mathcal A:=-\Delta+\abs x^2+1$, acting componentwise, and $\norm f_r:=\norm{\mathcal A^{r/2}f}_{L^2}$, $\mathcal H^{-r}$ is the dual of the completion of $\cS(\R^3;\cV)$ in $\norm\cdot_r$. It is a separable Hilbert space. Translations act on $\mathcal H^{-r}$ by $(\tau_a\Phi)(f):=\Phi(f(\cdot+a))$, and the reflection $\rho_{j,c}$ in the plane $\{x_j=c\}$ by the map $\theta_{j,c}$ induced by the push-forward of Section~\ref{ssec:symmetries}, with $\theta_j:=\theta_{j,0}$ (on the curvature components this is the reflection twisted by the signs of the 2-form). For an open $U\subset\R^3$ let $\Sig(U)$ be the $\sigma$-algebra generated by the evaluations $\Phi\mapsto\Phi(f)$, $f\in C_c^\infty(U;\cV)$, and write $\Hp jc:=\{x:x_j>c\}$, $U(a)G:=G\circ\tau_{-a}$, and $\cI$ for the $\sigma$-algebra of translation-invariant Borel sets.

\begin{theorem}[weak clustering]\label{thm:cluster}
Let $S$ be one of the limit states $S^{\mu_0}$ ($\mu_0\in Q_0$), $S^\fl$, $S^0$ of Theorem~\ref{thm:main}. For every even $r\ge14$ there is exactly one Borel probability measure $\mu$ on $\mathcal H^{-r}$ whose mixed moments on test functions in $C_c^\infty(\R^3;\cV)$ are the Schwinger functions of $S$, and there are a probability space $(\Xi,\nu)$ and Borel probability measures $\mu_\xi$ on $\mathcal H^{-r}$, with $\xi\mapsto\mu_\xi(B)$ measurable for every Borel set $B$, such that $\mu=\int\mu_\xi\,d\nu(\xi)$ and, for $\nu$-almost every $\xi$:
\begin{enumerate}
\item[(a)] $\mu_\xi$ is invariant under all translations and ergodic for the action of $\R^3$, and invariant under $\theta_1,\theta_2,\theta_3$;
\item[(b)] $\mu_\xi$ is reflection positive: $\E_{\mu_\xi}[\overline{F\circ\theta_{j,c}}\,F]\ge0$ for every $j$, every $c\in\R$ and every $F\in L^2(\Sig(\Hp jc),\mu_\xi)$;
\item[(c)] (clustering in the mean) for all $G,G'\in L^2(\mu_\xi)$, $\lim_{T\to\infty}(2T)^{-3}\int_{[-T,T]^3}\E_{\mu_\xi}[G'\,U(a)G]\,da=\E_{\mu_\xi}[G']\,\E_{\mu_\xi}[G]$;
\item[(d)] (clustering under diagonal separation) for all $G,G'\in L^2(\mu_\xi)$, $\E_{\mu_\xi}[G'\,U(a)G]\to\E_{\mu_\xi}[G']\,\E_{\mu_\xi}[G]$ as $\min_j\abs{a_j}\to\infty$;
\item[(e)] (reconstruction in each coordinate direction) for each $j$, the Osterwalder--Schrader reconstruction from $\mu_\xi$ in direction $j$ gives a Hilbert space $\cH_j$, a unit vector $\Omega_j$, a self-adjoint operator $H_j\ge0$ and a commuting unitary representation of the translations orthogonal to $e_j$; the joint fixed space of the operators $\mathrm{e}^{-tH_j}$ ($t\ge0$) and of these translations is $\C\Omega_j$ (a unique joint vacuum).
\end{enumerate}
For $\mu$ itself, (c) and (d) hold with the right side replaced by $\E_\mu\big[\E[G'\mid\cI]\,\E[G\mid\cI]\big]=\int\E_{\mu_\xi}[G']\,\E_{\mu_\xi}[G]\,d\nu(\xi)$.
\begin{enumerate}
\item[(f)] (decomposition along $\R^3\rtimes\Oh$) The decomposition may also be chosen with cubic-invariant components: there are a probability space $(\Xi',\nu')$ and Borel probability measures $\mu'_\zeta$ on $\mathcal H^{-r}$, with $\zeta\mapsto\mu'_\zeta(B)$ measurable for every Borel set $B$, such that $\mu=\int\mu'_\zeta\,d\nu'(\zeta)$ and, for $\nu'$-almost every $\zeta$, the measure $\mu'_\zeta$ is invariant under $\R^3\rtimes\Oh$, is reflection positive in the sense of (b), and is the uniform mixture $\mu'_\zeta=k^{-1}\sum_{i=1}^km_i$ of the $k$ distinct measures $m_1,\dots,m_k$ of one orbit of the group $S_3$ of coordinate permutations acting on the translation-invariant measures that are ergodic for $\R^3$; thus $k=k(\zeta)\in\{1,2,3,6\}$, and for each $j$ the joint fixed space of the operators $\mathrm{e}^{-tH_j}$ ($t\ge0$) and of the transverse translations in the reconstruction from $\mu'_\zeta$ in direction $j$, defined as in (e), has dimension $k$ (Theorem~\ref{add:Ohdecomp}).
\end{enumerate}
\end{theorem}
\noindent In the terms of Definition~\ref{add:OS4}, part (a) says that almost every component $\mu_\xi$ satisfies OS4 in the weak form.

\begin{remark}[the components]\label{rem:components}
For $\nu$-almost every $\xi$ the moments of $\mu_\xi$ exist for all Schwartz test functions and satisfy a bound $\abs{\E_{\mu_\xi}\prod_{i=1}^N\Phi(f_i)}\le\sigma^\xi_N\prod_i\abs{f_i}_k$ for a fixed Schwartz norm $\abs\cdot_k$, with $\sigma^\xi_N\le\alpha_\beta(\xi)(N!)^\beta$ for every $\beta>1$ (Corollary~\ref{cor:D:schwinger}(iv)); the constants depend on $\xi$. So the components satisfy a growth condition of the type of the second part of OS0, with the norm $\abs\cdot_k$ in place of $\abs\cdot_4$ and with $\xi$-dependent constants; we do not claim OS0 for them, in particular not the tube bound of Theorem~\ref{thm:main}(i). The components are invariant under the translations and the three coordinate reflections (part (a) of Theorem~\ref{thm:cluster}), so they satisfy OS1 for the group generated by these. A coordinate permutation maps a component to a component (Remark~\ref{rem:D:Ohcomp}), and nothing forces a component to be invariant under the permutations; grouping the components along the orbits of $\Oh$ gives the $\Oh$-invariant components of part (f) of Theorem~\ref{thm:cluster}, whose joint vacuum has dimension $k\in\{1,2,3,6\}$. The moments of a component are symmetric and real, and those of degree one in the curvature vanish (Theorem~\ref{thm:os4:FH}); the vanishing of its moments of odd degree three or more in the curvature, which holds for $S$ (Theorem~\ref{thm:C-OS3}), is not known for a single component.
\end{remark}

\begin{proposition}[the axis obstruction]\label{prop:axisobs}
For $r$ large there is a centred Gaussian measure on $\mathcal H^{-r}(\R^3)$ that satisfies the hypotheses (H1)--(H3) of Section~\ref{ss:abstract} with the untwisted reflections $\Phi\mapsto\Phi\circ\rho_j$, is invariant under translations and under $\Oh$, is ergodic for the action of $\R^3$, and whose two-point function does not tend to zero along any coordinate axis: $\E[\Phi(f)\,U(te_1)\Phi(f)]$ has a positive limit as $t\to\infty$ for every $f\in C_c^\infty(\R^3)$ whose integral over $x_1$ is not identically zero.
\end{proposition}
\noindent The example is $\Phi=\psi_1(x_2,x_3)+\psi_2(x_1,x_3)+\psi_3(x_1,x_2)$ with independent two-dimensional massive free fields $\psi_j$ (Example~\ref{ex:D:axis}, where it is realised on the field space $\cV$ with vanishing curvature component, so that the twisted and the untwisted reflections coincide). So single-axis clustering, even of the two-point function, is not a consequence of translation invariance, cubic invariance, ergodicity and reflection positivity.

\begin{proposition}[what closes single-axis clustering]\label{prop:closeaxis}
Let $\mu$ satisfy (H1)--(H3), let $\mathcal F\subset L^2(\mu)$ be a set whose linear span is dense, and let $P_j$ be the orthogonal projection onto the functions invariant under the translations along $e_j$. The following are equivalent: (i) $P_jG=\E[G\mid\cI]$ for all $G\in\mathcal F$ and $j=1,2,3$; (ii) for all $G\in\mathcal F$, $G'\in L^2(\mu)$ and each $j$, $\E[G'\,U(te_j)G]\to\E\big[\E[G'\mid\cI]\,\E[G\mid\cI]\big]$ as $\abs t\to\infty$; (iii) the same limit holds as $\abs a\to\infty$ in any manner. If $\mu$ is a centred Gaussian measure with covariance $\widehat C$ in Fourier space, then (i) holds for the linear fields $\Phi(f)$ (a statement about (i) alone; their linear span is not dense) if and only if $\widehat C(\{p:p_j=0\}\setminus\{0\})=0$ for $j=1,2,3$.
\end{proposition}

\begin{theorem}[single-axis clustering]\label{thm:os4main}
Let $S$ be one of the limit states $S^{\mu_0}$ ($\mu_0\in Q_0$), $S^\fl$, $S^0$ of Theorem~\ref{thm:main}, $\mu$ its law on $\mathcal H^{-r}$ (Theorem~\ref{thm:cluster}), $P_j$ the orthogonal projection of $L^2(\mu)$ onto the $e_j$-invariant vectors, and $\mathcal M$ the $\sigma$-algebra generated by the Wick-square fields $\XR(h)$, $h\in C_c^\infty(\R^3)$ real, completed by $\mu$-null sets. (Section~\ref{ss:os4axis} works with the law of $S$ on $H^{-s}_{\mathrm{loc}}(\R^3;\cV)$; Lemma~\ref{lem:D:hermite} carries every statement to the law on $\mathcal H^{-r}$, with $\mathcal M$ replaced by its preimage under the restriction map.)
\begin{enumerate}
\item[(a)] For every real 2-form $\bg$ with components in $C_c^\infty(\R^3)$ and every $Y\in L^2(\mathcal M)$: $F(\bg)Y\in L^2(\mu)$, $\E_\mu[F(\bg)^2\abs Y^2]\le\norm{\bg}_2^2\,\E_\mu\abs Y^2$, $P_j(F(\bg)Y)=0$ for $j=1,2,3$, and for every $G'\in L^2(\mu)$
\[
\lim_{\abs a\to\infty}\E_\mu\big[G'\,U(a)\big(F(\bg)Y\big)\big]=0=\E_\mu[G']\,\E_\mu[F(\bg)Y],
\]
the limit taken in any manner, in particular along each coordinate axis.
\item[(b)] Let $\Delta_S\ge0$ be the latent density of the Wick square in $S$ (Definition~\ref{def:os4:pressure}). For every $h\ge0$, $\abs{P_j\XR(h)-\E_\mu\XR(h)}\le\Delta_S\norm h_1$ $\mu$-almost surely, and $\Delta_S\le\kappa_++2S^B_\fl$ with the constants of Theorem~\ref{thm:ch-T9}. If $\Delta_S=0$, the conclusion of (a) holds for every $G$ in the closed span of $L^2(\mathcal M)$ and of the products $F(\bg)Y$, with limit $\E_\mu[G']\,\E_\mu[G]$.
\item[(c)] If hypothesis (D) of Corollary~\ref{cor:os4:OS4} holds, then $\mu$ is ergodic, $S$ satisfies OS4 (Definition~\ref{def:OS4}), and the conclusion of (a) holds for all $G,G'\in L^2(\mu)$, with limit $\E_\mu[G']\,\E_\mu[G]$.
\end{enumerate}
\end{theorem}

\subsection{Correlation signs and the curvature two-point function}\label{ssec:NTmain}

For a limit state $S$ let $S^T_{n,p}$ be its truncated functions (Definition~\ref{def:NT-trunc}) and $M(g;h):=S^T_{2,1}(g,g;h)$, the truncated correlation of $F(g)^2$ and $\XR(h)$. For a real vector field $G\in C_c(\R^3;\R^3)$ let $B(G):=F(\star^{-1}G)$ be the smeared magnetic field, $\mathfrak b_S(G,G')$ its two-point function, $P_T(p)$ the orthogonal projection of $\C^3$ onto $p^\perp$, and $\mathsf Z_S(G):=\limsup_{R\to\infty}R^3\mathfrak b_S(G_R,G_R)$ with $G_R:=R^{-3}G(\cdot/R)$ (Section~\ref{sec:massless}); $S$ is \emph{massless in the curvature sector} if $\mathsf Z_S(G)>0$ for some $G$ (Definition~\ref{def:M-massless}).

\begin{theorem}[signs, and the curvature two-point function]\label{thm:NTMmain}
Let $S\in\{S^{\mu_0}\ (\mu_0\in Q_0),S^\fl,S^0\}$ be a limit state of Theorem~\ref{thm:main}.
\begin{enumerate}
\item[(a)] (signs) For all real 2-forms $g,g'$ with components in $C_c(\R^3)$ and all $h,h'\in C_c(\R^3)$ with $h,h'\ge0$,
\[
M(g;h)\le0,\qquad S^T_{0,2}(h,h')\ge0,\qquad S^T_{4,0}(g,g,g',g')+2\,S_{2,0}(g,g')^2\ge0 .
\]
If $M(g;h)\neq0$ for one such pair, $S$ is not decoupled (Definition~\ref{def:NT-trunc}); hence it is not Gaussian, and the curvature and the Wick square are not independent under its law.
\item[(b)] (the curvature two-point function) There is a measurable $m_S$ with $0\le m_S(p)\le P_T(p)$ for a.e.\ $p$ and $\mathfrak b_S(G,G')=(2\pi)^{-3}\int\overline{\hat G(p)}^{\mathsf T}m_S(p)\hat G'(p)\,dp$; in particular $0\le\mathsf Z_S(G)\le\norm{\Pi_TG}_2^2$, with $\Pi_T$ the Fourier multiplier by $P_T$; the right side is the value for the free massless photon. If $S=S^{\mu_0}$ with $\mu_0>0$, then $m_S(p)\le\min(1,\abs p^2/\mu_0^2)P_T(p)$ and $\mathsf Z_S\equiv0$: the states at positive photon mass have no massless component. For $\mu_0\in Q_0$, $0=\mathsf Z_{S^{\mu_0}}\le\mathsf Z_{S^0}\le\mathsf Z_{S^\fl}$.
\item[(c)] (what masslessness would give) If $S\in\{S^\fl,S^0\}$ is massless in the curvature sector, its curvature correlations are not integrable, in particular not exponentially decaying, and in every coordinate direction the curvature sector of its reconstruction has no mass gap: the generator $H_F$ of Proposition~\ref{prop:M-gap}(a) has $\inf\operatorname{spec}H_F=0$.
\end{enumerate}
\end{theorem}
\noindent Theorem~\ref{thm:NT-cond} shows, under a remainder bound (R) comparing $e^{-2}M$ with its one-loop value and under Hypothesis~\ref{hyp:NT-C}, neither of which is proved here, that for every real 2-form $g$ with $d^*g\not\equiv0$ and every $0\le h\not\equiv0$ in $C_c^\infty(\R^3)$ one has $M(g;h)<0$ for $e^2+\lambda$ small enough depending on $(g,h)$; hence these limit states are not decoupled (Section~\ref{ss:nontriv}).

\begin{remark}[what is not claimed]\label{rem:notclaimed}
The theorems above do not assert uniqueness of the limits, the equality $S^0=S^\fl$ (part (iv) of Theorem~\ref{thm:main} gives one inequality on $\cD$), or the coincidence of the limits with King's cube limits. They do not assert that the limit states are non-Gaussian, that they or their ergodic components satisfy OS4 in the single-axis forms (i), (ii) of Definition~\ref{add:OS4} (Theorem~\ref{thm:cluster} gives the weak form (iii) for the ergodic components, and Theorem~\ref{thm:os4main} the single-axis forms for part of the observables), that the translation-ergodic components are invariant under $\Oh$, that the photon is massless at $\mu_0=0$, that the states lie in the symmetric phase (the proofs do not use the phase, and one-loop perturbation theory predicts a first-order transition surface inside the range of couplings of Theorem~\ref{thm:main}, Remark~\ref{rem:disc:firstorder}), or invariance under rotations outside $\Oh$. Section~\ref{sec:disc} discusses each of these questions.
\end{remark}

\subsection{Where the results are proved}\label{ssec:proofmap}

The proofs of Theorems~\ref{thm:main} and \ref{thm:T6} use the statements collected in Appendix~\ref{app:K}, whose proofs are not included in this version.

\emph{Theorem~\ref{thm:main}.} Part (i) is Theorem~\ref{thm:ch-T9} of Section~\ref{sec:chain}. Its curvature factor is the Gaussian domination of the photon on the torus (Lemma~\ref{lem:C-Fsector}). The factor for positive $\wick{\abs\varphi^2}$-sources is the convexity-and-chessboard bound on the same torus (Propositions~\ref{prop:ch-Pplus} and~\ref{prop:ch-T7}); its input is Ba{\l}aban's two-sided bound~\eqref{eq:K0} at the scalar masses $0.7m$ and $m$ on every $T_j$ (Proposition~\ref{prop:ch-K0}), together with the counterterm gap of Theorem~\ref{thm:ch-W}. The factor for negative sources (Proposition~\ref{prop:ch-Tflat}) follows from Theorem~\ref{thm:T6}, the lower vacuum bound of Theorem~\ref{thm:UV-K0minus} at $m_1$, and the counterterm gap between $m_1$ and $m_2$. The proofs of these statements use the results of Section~\ref{sec:UV} and Appendices~\ref{app:B} and~\ref{app:G1:area-UV}--\ref{app:G1:part-D6}. Part (ii) is compactness (Montel's theorem on the tube and a diagonal argument), tightness on $H^{-s'}_{\mathrm{loc}}$, and, for $S^0$, the monotonicity on the class $\cD$ at fixed lattice passed to the limit along the common subsequence (Section~\ref{sec:C-limits}, with Lemma~\ref{app:G1:S0measure} for the law of $S^0$). Part (iii) is Section~\ref{sec:C-OS}: the bound passes to the limit, and the lattice symmetries and lattice reflection positivity on tori (Section~\ref{sec:C-lattice}) are preserved under the limits. Part (iv) is the monotonicity in part (ii).

\emph{Theorem~\ref{thm:cluster} and Propositions~\ref{prop:axisobs}, \ref{prop:closeaxis}.} These are proved in Section~\ref{sec:OS4}: the ergodic decomposition with inheritance of the reflections and of reflection positivity, which rests on the fact that a function invariant under translations along $e_j$ is fixed by $\theta_j$ (Section~\ref{ss:decomp}); the reconstruction in a coordinate direction and the identification of its zero-energy vectors with the $e_j$-invariant functions (Section~\ref{ss:recon}); clustering in the mean from the mean ergodic theorem, and clustering under diagonal separation from the telescoping identity $1=P_1P_2P_3+(1-P_1)+P_1(1-P_2)+P_1P_2(1-P_3)$ and the spectral theorem for each $H_j$ (Section~\ref{ss:cluster}); the verification of (H1)--(H3) for the limit states (Section~\ref{ss:discharge}); and Corollary~\ref{cor:D:schwinger} with Lemma~\ref{lem:D:hermite}. Proposition~\ref{prop:axisobs} is Example~\ref{ex:D:axis}, carried to $\mathcal H^{-r}$ by the carrier argument of Lemma~\ref{lem:D:hermite}(b), and Proposition~\ref{prop:closeaxis} is Proposition~\ref{prop:D:axis}.

\emph{Theorems~\ref{thm:os4main} and~\ref{thm:NTMmain}.} Theorem~\ref{thm:os4main} is Theorem~\ref{thm:os4:summary}(a)--(c) of Section~\ref{ss:os4axis}. Part (a) is Theorem~\ref{thm:os4:FH}, whose input is a domination of the curvature conditional on the matter on every lattice torus (Lemma~\ref{lem:os4:condlat}); part (b) is Proposition~\ref{prop:os4:latent} and Theorem~\ref{thm:os4:matter}, from the chessboard estimate at the photon mass of the state and the association of the radial field (Lemma~\ref{lem:os4:fkg}); part (c) is Corollary~\ref{cor:os4:OS4}, from the multiple-reflection bound in infinite volume (Theorem~\ref{thm:os4:block}). Theorem~\ref{thm:NTMmain}(a) is Theorem~\ref{thm:NT-sign}(b) and Corollary~\ref{cor:NT-crit}; (b) is Theorem~\ref{thm:M-density} and Theorem~\ref{thm:M-dichotomy}(i)--(iii); (c) is Corollary~\ref{cor:M-summary}.

\section{Ultraviolet input from the literature}\label{sec:UV}

This section lists the results from the literature that the construction of Theorem~\ref{thm:main} uses, in the form in which they are used, with the page where each is printed: Ba{\l}aban's three papers on the $(\mathrm{Higgs})_{2,3}$ model in a finite volume \cite{Ba1,Ba2,Ba3}, his paper on lattice Green's functions \cite{Ba4}, and King's two papers on the $U(1)$ Higgs model \cite{K1,K2}, together with the classical results (correlation inequalities, power counting) of Subsection~\ref{ssec:UV-other}. Three printed proofs are incomplete at places that the construction uses; Section~\ref{sec:T6} completes them (Lemma~\ref{lem:lem24}, Proposition~\ref{prop:prop1} and Proposition~\ref{app:G1:I6}).

\subsection{Ba{\l}aban's torus and scales}\label{ssec:UV-torus}
Ba{\l}aban works on the lattice torus
\[
T_\eps=\{x\in\eps\Z^3:\ -L_\mu\le x_\mu<L_\mu,\ \mu=1,2,3\},\qquad \abs{x-y}=\max_\mu\min\{\abs{x_\mu-y_\mu},\,2L_\mu-\abs{x_\mu-y_\mu}\},
\]
with sides $2L_\mu$ such that $\eps^{-1}L_\mu\in L^K M_b\N$ \cite[(1.2)--(1.3), p.~604]{Ba1}; here $L\ge2$ is the block factor, $M_b$ the side of a big (``large'') block, and the distance is the torus sup-distance. All distances in Sections~\ref{sec:UV}--\ref{sec:chain} are measured in this metric. For the lattices of Theorem~\ref{thm:main}, $\eps=L^{-K}\eps_0$ with $L=2$, so the $k$-th renormalisation scale is $t_k:=L^k\eps$ and the last one is $t_K=\eps_0$. In the $k$-th step all objects are rescaled by $t_k$; scale-$k$ blocks then have side $1$ and the lattice spacing is $\eta_k:=L^{-k}$ (Ba{\l}aban writes $\eta$; we reserve $\eta$ for the parameter of the window $\Wwin$ of \eqref{eq:window}). The large-field regions of the $k$-th step are organised in nested unions of big blocks $\Lambda_0\supset\Lambda_1\supset\cdots$, where $\Lambda_{i+1}^c$ is the union of the big blocks within distance $r(t)$ of $\Lambda_i^c$, and
\begin{equation}\label{eq:UV-r}
r(t)=R\,(1+\log t^{-1})^{\varrho},\qquad \varrho>1,\quad R>R_0,
\end{equation}
\cite[(2.7)--(2.8), p.~558]{Ba2}; we take $\varrho=2$. Ba{\l}aban's small-field parameter is $p(t):=b_0(1+\log t^{-1})^{p}$; we take $p=6$. The constants $R$ and $b_0$ are fixed with the regulator data (Appendix~\ref{app:B}), and $R_0$ is the constant of Proposition~\ref{prop:UV-BaProp21}. At step $k$ the sets of scale $k-1$ are written $\Lambda_i^{(k-1)\prime}$, and in scale-$k$ units consecutive layers are at distance at least $\rho_k:=(r(t_{k-1})-M_b)/L$ \cite[p.~558, and Sect.~B, pp.~566\,ff.]{Ba2}.

The tori of the construction are the cubic tori $T_j$ of Subsection~\ref{ssec:tori}, of side $\ell_j=j!\,\ell_0$ with $\ell_0\in 2M_b\N$; the regulator data have $\ell_0\ge183/m$ (Appendix~\ref{app:B}). By Lemma~\ref{lem:sizes}, for every spacing $\eps=2^{-K}\eps_0$ the lattice torus $T_{j,\eps}$ is a torus of the form \cite[(1.2)]{Ba1} (part (ii) with $k=K$), and the number $N=\ell_j/\eps$ of sites per side is even (part (i)). We write $K_0'\in\N$ for the index from which the torus-uniform bounds of Theorem~\ref{thm:T6} hold; it is fixed with the regulator data (Appendix~\ref{app:B}). The spacing threshold of Theorem~\ref{thm:main} is
\[
K_0:=\max(K_0',\,n_0+1).
\]

\subsection{Ba{\l}aban's vacuum bounds}\label{ssec:UV-K0}
For a Higgs mass $M>0$ let
\[
Z_K(T;M)=\int dA\,d\varphi\,\exp[-S^\eps(T,A,\varphi;M)],
\]
where $S^\eps$ is King's periodic action \cite[(2.4), p.~652]{K1} on $\Omega=T_\eps$ with $m^2=M^2$, photon mass $\mu_0=\mu_*$, the mass counterterm $\delta m^2_{K,T}(M)$ and the vacuum constant $E_0(M)+E_1(M)$ of \cite[(2.6), (3.35)]{K1} (Subsection~\ref{ssec:UV-King}). By Proposition~\ref{prop:UV-B1e} this is the action of \cite[(1.11)]{Ba1}.

\begin{theorem}[{lower vacuum bound; \cite[Thm.~(1.14), p.~606]{Ba1}}]\label{thm:UV-K0minus}
Let $\mu_*>0$, $\lambda>0$, $M>0$, the vector field in the Feynman gauge, and the charge as in \cite[p.~605]{Ba1}. There is $E_-(M)<\infty$, independent of $\eps$ and of the torus, with $Z_K(T;M)\ge \mathrm{e}^{-E_-(M)\abs{T}}$.
\end{theorem}
\noindent In this paper Theorem~\ref{thm:UV-K0minus} is used at $M\in\{m,m_1,m_2\}$.

\begin{theorem}[{upper vacuum bound; \cite[Thm.~(1.3), p.~556]{Ba2}}]\label{thm:UV-K0plus}
Under the hypotheses of Theorem~\ref{thm:UV-K0minus} there is $E_+(M)<\infty$, independent of $\eps$ and of the torus, with $Z_K(T;M)\le \mathrm{e}^{E_+(M)\abs{T}}$.
\end{theorem}
\noindent Remark~\ref{rem:UV-blockD} describes how the estimates that its printed proof leaves out are supplied in this paper. Theorem~\ref{thm:UV-K0plus} is used at $M_1=0.7m$ (on the way to the tube bound, through Proposition~\ref{prop:ch-Pplus}) and at $m_1,m_2$.

\begin{lemma}[the hypotheses of Theorems~\ref{thm:UV-K0minus} and~\ref{thm:UV-K0plus} hold for $Z_K$]\label{rem:UV-K0hyp}
For every $M>0$, every torus $T_j$ and every spacing $\eps$, $Z_K(T;M)$ is Ba{\l}aban's Feynman-gauge integral \cite[(1.10)]{Ba1} with photon mass $\mu_0=\mu_*>0$, and its charge matrix is $q=\begin{pmatrix}0&-1\\1&0\end{pmatrix}$, the case ``$N=2$ and $q^2=-1$'' of \cite[p.~605]{Ba1}. In particular Theorems~\ref{thm:UV-K0minus} and~\ref{thm:UV-K0plus} apply to $Z_K(T;M)$.
\end{lemma}
\begin{proof}
(a) \emph{Feynman gauge.} In \cite[p.~605]{Ba1} the Feynman gauge is introduced by replacing the integral (1.9), with the action (1.8), by the integral (1.10), with the action (1.11); the vector-field part of (1.11) is $\frac12\langle A,(-\Delta^\eps+\mu_0^2)A\rangle$, with $-\Delta^\eps$ the Laplace operator on the torus $T_\eps$. The photon part of the action of $Z_K(T;M)$ is King's Feynman form \cite[(2.4), p.~652]{K1} at $\mu_0=\mu_*$, and by Proposition~\ref{prop:UV-B1e} the whole action is that of \cite[(1.11)]{Ba1}. So $Z_K(T;M)$ is Ba{\l}aban's Feynman-gauge integral \cite[(1.10)]{Ba1} with $\mu_0=\mu_*>0$.

(b) \emph{Charge.} \cite[p.~605]{Ba1} writes $U(A)=\exp(q\eps eA)$ with $q$ an antisymmetric $N\times N$ matrix and $\norm q\le1$. Here $\varphi(x)\in\C\cong\R^2$, so $N=2$, and the covariant difference in \eqref{eq:matteraction} multiplies by $\mathrm{e}^{ie\eps A_{x,\mu}}$, which acts on $\R^2$ as the rotation by the angle $e\eps A_{x,\mu}$. So $q=\begin{pmatrix}0&-1\\1&0\end{pmatrix}$, which is antisymmetric with $\norm q=1$ and $q^2=-1$. This is the case ``$N=2$ and $q^2=-1$'' of \cite[p.~605]{Ba1}, in which $\delta m^2$ is a number.
\end{proof}

\begin{remark}[on the printed proof of Theorem~\ref{thm:UV-K0plus}]\label{rem:UV-blockD}
The abstract of \cite{Ba2} (p.~555) states that the proof is completed there with the exception of some technical estimates, and \cite[p.~411]{Ba3} supplies them through \cite[Prop.~1]{Ba3}. Some estimates of \cite{Ba2,Ba3} are needed by the proofs of Theorems~\ref{thm:UV-K0plus} and~\ref{thm:T6} but are not written out in print. They are the statements $\mathcal I$ of Definition~\ref{app:G1:itemlist}. They are proved, at the two points $(u,M)\in\{(1,m_2),(-\eta,m_1)\}$ of Theorem~\ref{thm:T6}, in Lemma~\ref{app:G1:items} and the results it points to. Lemma~\ref{app:G1:items} also uses the margins of the two-domain bounds at the first step (Lemma~\ref{app:K:L18-1}) and the identification of the bad vacuum localisations at the last scale (Lemma~\ref{app:K:L18-2}). The transfer, Lemma~\ref{app:G1:RMG}, also uses the uniformity in $\eps_0$ of Ba{\l}aban's per-step constants, Statement~\ref{app:K:SL}, and some of the results below use the other statements collected in Appendix~\ref{app:K}; the proofs of those statements are not included in this version. The omitted estimates are proved, under these statements, in the following places:
\begin{itemize}
\item the depth at which each sketched step of \cite[Prop.~1]{Ba3} is applied: Proposition~\ref{prop:prop1}, Lemma~\ref{lem:Omegafree} and Corollary~\ref{cor:prop1R4};
\item \cite[Lemma~2.4]{Ba2} at edge points: Lemma~\ref{lem:lem24}, under the cap \eqref{eq:lem24-cap};
\item the comparison of torus and infinite-lattice propagators: Lemma~\ref{lem:IMG};
\item the combinatorial estimate \cite[(3.42)--(3.54), pp.~592--594]{Ba2}, which ends the proof of Theorem~\ref{thm:UV-K0plus}: Proposition~\ref{app:G1:I6};
\item the ingredients of Table~\ref{tab:T6-inputs}: Propositions~\ref{app:G1:I1}, \ref{app:G1:I2}, \ref{app:G1:I3}, \ref{app:G1:I4}, \ref{app:G1:I5}, \ref{app:G1:I7}, \ref{app:G1:I8} and~\ref{app:G1:I9};
\item the statements that make up Proposition~\ref{app:G1:I3}, listed in Table~\ref{app:G1:I3tab}: Subsection~\ref{app:G1:part-TI3}, Lemmas~\ref{app:G1:Rnest}, \ref{app:G1:Rthetab}, \ref{app:G1:Rloc7}, \ref{app:G1:AA}, \ref{app:G1:AA4}, \ref{app:G1:I3a-onebg}, \ref{app:G1:I3a-eta1}, \ref{app:G1:I3a-Bloc}, \ref{app:G1:I3a-margins}, \ref{app:G1:I3a-rate} and Propositions~\ref{app:G1:twodomain}, \ref{app:G1:T579};
\item the order cap $\Nbar=169$ for \cite[(1.33)]{Ba3}: Proposition~\ref{app:G1:I5};
\item the final step \cite[(2.116)--(2.118), p.~582]{Ba2} of the upper bound: Lemma~\ref{app:G1:N10}, with its hypothesis (F) supplied by Proposition~\ref{app:G1:assembly}.
\end{itemize}
The transfer of these estimates to the points $(u,M)=(0,M)$, $M\in\{M_1,m_1,m_2\}$, at which Theorem~\ref{thm:UV-K0plus} is used, is Lemma~\ref{app:G1:RMG}. The point $u=0$ is Ba{\l}aban's own model, on which these estimates carry no $u$-weight. Each estimate depends on the Higgs mass in one of two ways. Either the mass enters only from above, through $M^2t^2\le c$ (imposed at $m_2\ge M$). Or it enters through powers of $M$ of either sign and through logarithms such as $\log_+(1/(M\eps_0))$, in quantities that are finite at every fixed $M>0$. These quantities enter only three things: downward-closed upper bounds on $\eps_0$, constants evaluated at the four masses $M_1,m,m_1,m_2$, and the lower bound $\ell_0\ge128/M_1$ on the torus. None of the estimates orders $M$ against $\mu_*$ or uses $m_1>2m$. The estimate of Proposition~\ref{app:G1:I6} is combinatorial: it does not involve the fields or $u$. Four ingredients of these estimates hold at every $M\in(0,m_2]$ (Lemma~\ref{lem:B:massrange}): the torus-uniform final c-numbers, the $A\cdot A$ vertices, the tail beyond order $\nbar$ and the lower bound on the averaged mass operator.
\end{remark}

\subsection{Ba{\l}aban's propagator bounds}\label{ssec:UV-prop}\label{app:G1:ssec-BaProp21}
Let $\Omega=B^k(\Omega^{(k)})$, where $\Omega^{(k)}$ is a union of big blocks of the rescaled torus, let $A$ be a vector field regular on $\Omega$ in the sense of \cite[(2.23)]{Ba1}, and let $G_k(\Omega,A)$ be the Green's function of the operator $-\Delta^{\eta_k,N}_{A,\Omega}+m^2(L^k\eps)^2+a_kP_k(A)$ \cite[(2.22), p.~610]{Ba1} (Neumann conditions on $\partial\Omega$, $P_k=Q_k^*Q_k$ the block averaging, \cite[(2.11), p.~609]{Ba1}).

\begin{proposition}[{\cite[Props.~2.1--2.3, pp.~610--611]{Ba1}, proved in \cite{Ba4}}]\label{prop:UV-BaProp21}
If the running charge $e(L^k\eps):=e\,(L^k\eps)^{(4-d)/2}$ is sufficiently small and the H\"older exponent $\alpha$ of \cite[(2.24)]{Ba1} satisfies $\alpha<1$, there are $\delta_0,c_0,R_0>0$, depending only on $d$, $a$ and $M_b$ ($c_0$ also on $\alpha$; independent of $A$, $k$, $\Omega$ and $\eps$), such that:
\begin{enumerate}
\item[(2.24)--(2.25)] $\abs{(G_k(\Omega,A)f)(x)}$ and $\abs{(D^{\eta_k}G_k(\Omega,A)f)(x)}$, and the H\"older quotients of \cite[(2.24)]{Ba1}, are bounded by $c_0\mathrm{e}^{-\delta_0\dist(x,\supp f)}\norm f_\infty$ for $x\in\Omega$ with $\dist(x,\Omega^c)\ge R_0$ (for (2.24): both points at distance $\ge R_0$ from $\Omega^c$);
\item[(2.26)] for $\Omega\subset\Omega_0$, $\delta G_k(\Omega,\Omega_0,A):=G_k(\Omega,A)-G_k(\Omega_0,A)$ obeys the same bounds with the additional factor $\exp[-\delta_0\dist(\supp f,\Omega^c)-\delta_0\dist(\{x,x'\},\Omega^c)]$, $\Omega^c$ the complement of the smaller set;
\item[(p.~611)] if $\Omega$ is a rectangular parallelepiped, the bounds hold for all $x,x'\in\Omega$, without the restriction $\dist\ge R_0$;
\item[(2.34)] the block covariances $C^{(k)}_\Lambda(\Omega,A)$ obey exponential bounds without restriction.
\end{enumerate}
The operators are positive uniformly in the lattice spacing, $-\Delta^{\eta_k,N}_{A,\Omega}+aP_k(A)\ge\gamma_0$ \cite[(1.8), p.~573]{Ba4}, and their Green's functions decay exponentially \cite[Cor.~2.3, (2.30), p.~580]{Ba4}.
\end{proposition}
\noindent Two features of Proposition~\ref{prop:UV-BaProp21} are used repeatedly below. The restriction $R_0$ concerns the evaluation points only: in \cite[(2.24)--(2.25), p.~610]{Ba1} the condition $\dist(\cdot,\Omega^c)\ge R_0$ is imposed on $x$ (and $x'$), and the support of $f$ may touch $\partial\Omega$. And the constants are uniform over all unions of big blocks (``independent of $A$, $k$, $\Omega$'', \cite[p.~610]{Ba1}). The constants also do not depend on the side of the torus:

\begin{lemma}[independence of the torus side]\label{app:G1:torusuniform}
Let $T_{\eta_k}$ be a lattice torus whose side in direction $\mu$, in scale-$k$ units, is $n_\mu M_b$ with $n_\mu\in\N$. Then Proposition~\ref{prop:UV-BaProp21} and \cite[(1.8)]{Ba4} hold on $T_{\eta_k}$ with the constants $\gamma_0,\delta_0,c_0,R_0$ of \cite{Ba4}. These constants do not depend on $(n_\mu)$. Distances are those of \cite[(1.3)]{Ba1}.
\end{lemma}
\begin{proof}
\emph{Case $\min_\mu n_\mu\ge3$.} We follow the proof in \cite[Sect.~2, pp.~574--586]{Ba4}. Let $J=\prod_\mu\Z/n_\mu\Z$ index the big-block vertices $M_bj$, and let $C_j:=M_bj+[-M_b,M_b]^d$.

(1) Since $2M_b<n_\mu M_b$, the projection from $\eta_k\Z^d$ to $T_{\eta_k}$ is injective on $C_j$. It preserves bonds, unit blocks and the contours of $Q_k$ inside blocks. So $\Pi_j=\Omega\cap C_j$ is a union of at most $2^d$ big blocks in an embedded cube, and the constant-$A_0$ gauge $\lambda_0(x)=A_0\cdot(x-M_bj)$ is defined on $C_j$. Hence \cite[Lemmas~2.1, 2.2]{Ba4} apply to $G_k(\Pi_j,A_j)$ with their constants, which depend on $d,M_b,\alpha,c$ only. The same holds for \cite[Lemma~2.4]{Ba4}, which is proved on parallelepipeds of $\xi\Z^d$ \cite[pp.~584--586]{Ba4}. The field $A_j$ is built from $A|_{C_j}$ alone.

(2) Let $h\in C_0^\infty((-\tfrac23,\tfrac23))$ with $h=1$ on $[-\tfrac13,\tfrac13]$ and $\sum_{j\in\Z}h(\cdot-j)^2=1$, and put $h_j(x)=\prod_\mu h(x_\mu/M_b-j_\mu)$, $j\in J$ \cite[p.~575]{Ba4}. Since the sides are multiples of $M_b$, the identity $\sum_{j\in J}h_j^2=1$ holds on $T_{\eta_k}$, and $\supp h_j\subset C_j$. Hence \cite[(2.2)--(2.9)]{Ba4} hold as on $\eta_k\Z^d$. If $\Omega=T_{\eta_k}$, then $\partial\Omega=\emptyset$ and \cite[(2.6)]{Ba4} is trivial.

(3) For each $j\in J$ we have $\#\{j':\abs{j-j'}\le1\}\le3^d$, and each point lies in at most $2^d$ of the sets $\supp h_j$. Hence $\norm R_{2,2}\le c_d\,O(M_b^{-1})<1$ as in \cite[p.~577]{Ba4}, and the bound $2^d(3^d)^{n-1}2^d$ on the number of walks of length $n$ \cite[p.~579]{Ba4} holds unchanged.

(4) The metric \cite[(1.3)]{Ba1} is the quotient metric: $\dist_T(x,y)=\min_\gamma\abs{\tilde x-\tilde y-\gamma}_\infty$, the minimum taken over the deck translations $\gamma$. Consecutive cubes of a walk have centres at distance at most $M_b$. So $\dist_T(\{x,x'\},\supp f)\le(n+2)M_b$, i.e.\ $n\ge M_b^{-1}\dist_T-2$, and the sum \cite[(2.22)]{Ba4} gives $\delta_0,c_0$ as on $\eta_k\Z^d$. Two further steps use only (4): the definition \cite[(2.19)]{Ba4} of $R_0$, and the implication ``$\dist(\{x,x'\},\Omega^c)\ge R_0$ implies that the first $n_0$ cubes of the walk are contained in $\Omega$''. The cancellation of the interior-cube terms for $\delta G_k(\Omega,\Omega_0,A)$ uses only (1)--(2). If $\Omega=T_{\eta_k}$, every $\Pi_j=C_j$ is a cube, so no restriction on $x,x'$ is needed.

(5) For \cite[(1.8)]{Ba4}: $P_k(A)$ is block diagonal, so dropping the bonds between unit blocks gives $\ip{\varphi}{(-\Delta^{N}_{A,\Omega}+m^2(L^k\eps)^2+a_kP_k(A))\varphi}\ge\sum_{\Delta}\ip{\varphi}{(-\Delta^N_{A,\Delta}+a_kP_k(A))\varphi}_\Delta$, a sum over unit blocks $\Delta$. On each unit block, the gauge of (1) and the expansion in $A'$ reduce the bound to $A=0$, as in \cite[(2.26$'$)--(2.27), p.~580]{Ba4}. There $P_k$ is the orthogonal projection onto the constants. The Neumann gap of a block with $N\ge2$ sites per side is $4N^2\sin^2(\pi/2N)\ge8$. So the lower bound is $\min\{8,a_k\}-o(1)$, and it does not involve the torus.

(6) The block covariances \cite[(2.34)]{Ba1}, i.e.\ \cite[Prop.~2.3]{Ba1}, are derived in \cite[Sect.~3, pp.~586--589]{Ba4} from two ingredients. The first is \cite[Cor.~2.3]{Ba4}, which is covered by (1)--(4); these use only the $L^2$ bound of \cite[Lemma~2.1]{Ba4}. The second is a family of Gaussian integrals that factorise over the blocks $B^{k+1}(z)$. These are local, and their block lower bound is (5) with $L$-blocks. Hence \cite[(2.34)]{Ba1} holds on $T_{\eta_k}$ with constants independent of $(n_\mu)$.

(7) Let $\Omega$ be a rectangular parallelepiped that does not wrap around $T_{\eta_k}$. Then every $\Pi_j$ is a parallelepiped in an embedded cube, so \cite[Lemma~2.2]{Ba4} applies to every factor of the walk expansion of \cite[p.~579]{Ba4}. The bounds then hold for all $x,x'\in\Omega$, as in \cite[p.~579]{Ba4}.

No other constant in \cite[Sect.~2]{Ba4} involves the global geometry.

\emph{Case $\min_\mu n_\mu\le2$ (covering step).} Let $\tilde T$ be the torus of sides $3n_\mu M_b$, and let $\pi:\tilde T\to T$ be the covering map, with deck group $\Gamma\cong(\Z/3\Z)^d$. Put $\tilde\Omega:=\pi^{-1}\Omega$, a union of big blocks, and $\tilde A:=A\circ\pi$, which is regular with the same constants. Choose the $Q_k$-contours of $\tilde T$ as the lifts of those of $T$. Then $\tilde H:=-\Delta^N_{\tilde A,\tilde\Omega}+m^2(L^k\eps)^2+a_kP_k(\tilde A)$ satisfies $\tilde H(g\circ\pi)=(Hg)\circ\pi$, and $\tilde H$ is invertible by the first case. Hence $G_k(\tilde\Omega,\tilde A)(f\circ\pi)=(G_k(\Omega,A)f)\circ\pi$. The same holds for $D^{\eta_k}$, for the H\"older quotients (lift $x'$ next to $\tilde x$; the shortest contours lift) and for $\delta G_k$. By (4), $\dist_{\tilde T}(\tilde x,\pi^{-1}S)=\dist_T(x,S)$ for every set $S$, and $\norm{f\circ\pi}_\infty=\norm f_\infty$. So the bounds on $\tilde T$, whose sides are at least $3M_b$, transfer to $T$ with the same constants. Positivity on $T$ is positivity on the $\Gamma$-invariant subspace.
\end{proof}
\noindent\label{app:G1:torusuniform-use}On the tori $T_j$ at the steps $k\le K$, Lemma~\ref{lem:sizes}(ii) gives $n_\mu=j!\,(\ell_0/M_b)L^{K+n_0-k}\ge2L^{n_0}\ge4$ big blocks per side, so only the first case of the proof is needed there. Lemma~\ref{app:G1:torusuniform} is used in step S19 of the proof of Proposition~\ref{prop:prop1}, for an auxiliary torus $T'\supset\square$ whose side lies in $2M_b\N$, is at least $\max\{3M_b,3r_\square+2s_b+2\}$, and has a big-block grid compatible with $\square$. It also makes the constants of every application of Proposition~\ref{prop:UV-BaProp21} on $T_{j,\eps}$ the same for all $j$.

\subsection{Ba{\l}aban's Proposition~1 of Higgs III}\label{ssec:UV-prop1}
In the $k$-th step of \cite{Ba2} the torus is replaced by $\Omega:=B^k(\Lambda_2^{(k-1)\prime})$, a union of big blocks, and one sets $\Omega_1:=B^k(\Lambda_7^{(k-1)\prime})$; a second domain $\Omega_2\subset\Omega$ ``satisfying the same conditions as $\Omega$'' is allowed \cite[pp.~411--414]{Ba3}. The background fields $\tilde A,\tilde B$ satisfy $\abs{\tilde A},\abs{\partial^{\eta_k}\tilde A},\abs{\partial^{\eta_k}\tilde B}\le O(1)p(t_k)$ with the corresponding H\"older bounds, and the support hypothesis
\begin{equation}\label{eq:UV-R4}
\text{(R4)}\qquad \dist(\supp\tilde A,\partial\Omega)>2r(t_k)
\end{equation}
\cite[p.~412]{Ba3}. \cite[Prop.~1, pp.~420--421]{Ba3} bounds the localised renormalised perturbative expressions $E(G_{\rm ren},\{\square(v)\},\Phi_{\rm ext},A_{\rm ext})$ of \cite[(1.33)]{Ba3}, with constants ``independent of $\eps$, $k$, the domains $\Omega,\Omega_1,\Omega_2$, the vector field $\tilde B$'' \cite[p.~421]{Ba3}. Its proof occupies \cite[pp.~422--445]{Ba3} and contains sketched steps; Proposition~\ref{prop:prop1} writes each of them out and determines how far from $\partial\Omega$ it is applied.

\subsection{King's counterterms and their identification with Ba{\l}aban's}\label{ssec:UV-King}
For a Higgs mass $M$ and a torus $T$, King's mass counterterm on $\Omega=T$ at the reference photon mass $\mu_*$ is a finite sum of lattice graphs of orders $\lambda$, $e^2$, $\lambda^2$, $e^2\lambda$, $e^4$ \cite[(3.32)--(3.33)]{K1} (Definition~\ref{def:counterterm}); we write
\[
a_{K,T}(M):=M^2+\delta m^2_{K,T}(M)
\]
for the resulting bare quadratic coefficient, and $E_0(M)+E_1(M)$ for King's vacuum constant \cite[(2.6), (3.35)]{K1}.

\begin{proposition}[King's counterterm and vacuum constant are Ba{\l}aban's]\label{prop:UV-B1e}
$a_{K,T}(M)$ is the coefficient at which Ba{\l}aban's theorems are stated on $T_\eps$ (\cite[p.~417]{Ba3}, \cite[(1.11)--(1.13), pp.~605--606]{Ba1}), and King's vacuum constant equals Ba{\l}aban's $E=E_0+E_1$ at the truncation order $\nbar=13$ of the construction.
\end{proposition}
\begin{proof}
The statement on the mass counterterm is Lemma~\ref{app:G1:ctid}; the statement on the vacuum constant is Lemma~\ref{app:G1:vac}.
\end{proof}
\noindent In the slice models of Section~\ref{sec:T6} the order-zero vacuum constant is the scalar factor $E_{0,s}(M)$ of \cite[(1.12), p.~606]{Ba1} only, because the photon measure there is normalised.

\subsection{Other classical inputs}\label{ssec:UV-other}
\begin{enumerate}
\item \emph{Current representation and correlation inequalities} \cite{BFS1}: for every positive covariance $\Gamma$ of the photon and every real single-site source, the matter weight is of positive type in $A$, so that $Z(\Gamma;a)=\sum_nw(n)\,\mathrm{e}^{-\frac12 e^2\eps^2\langle n,\Gamma n\rangle}$ with $w(n)\ge0$, the sum running over divergence-free integer currents $n$ \cite[Thm.~4.1(2), p.~262]{BFS1} (Lemma~\ref{lem:postype}); the two classes of \cite[Sect.~6]{BFS1} move in opposite directions. For a site function $g\ge0$, $r^2(g):=\sum_xg_xr_x^2$, and a real linear functional $A(f)$ of the bond variables, $\langle \mathrm{e}^{r^2(g)}\mathrm{e}^{iA(f)}\rangle$ is non-increasing and $\langle \mathrm{e}^{-r^2(g)}\mathrm{e}^{A(f)}\rangle$ is non-decreasing in the photon covariance \cite[Cor.~6.3(1), p.~281]{BFS1} (on the torus graph, Lemma~\ref{lem:C-cov}); in particular, for $h\ge0$, $\langle \mathrm{e}^{\wick{\varphi^2}(h)}\rangle$ does not increase and $\langle \mathrm{e}^{-\wick{\varphi^2}(h)}\rangle$ does not decrease when the covariance increases. In the bare coefficient, \cite[Thm.~6.1, p.~276]{BFS1} with $P,Q$ polynomials in $r$ with positive coefficients gives $\langle P;r_x^2\rangle\ge0$, so $\langle P\rangle$ (for instance $\langle r_y^2\rangle$, or $\langle \mathrm{e}^{\wick{\varphi^2}(h)}\rangle$ with $h\ge0$) does not increase when the coefficient is increased at a site. In the hopping couplings the directions are reversed \cite[Thm.~6.2(2),(4), p.~280; Cor.~6.3(2), p.~281]{BFS1}: $\langle \mathrm{e}^{r^2(g)}\mathrm{e}^{iA(f)}\rangle$ does not decrease and $\langle \mathrm{e}^{-r^2(g)}\mathrm{e}^{A(f)}\rangle$ does not increase when a coupling is increased.
\item \emph{Chessboard estimate} for reflection-positive lattice measures on a torus with an even number of sites per side, for reflections in the planes halfway between lattice planes \cite[Sect.~4]{FILS78}; in the lattice-gauge setting it is \cite[Thm.~2.10, p.~252]{BFS1} and \cite[Thm.~2.2, p.~16]{Sei82}, both derived from reflection positivity alone. On the tori $T_j$ it rests on lattice reflection positivity about such planes, Proposition~\ref{prop:C-RP} with the plane pairs of case~(ii) of Lemma~\ref{lem:C-RPP}.
\item \emph{Lattice power counting} \cite[Thm.~1, p.~90]{Rei88} and lattice renormalisation \cite{Rei88b}; both are stated for four-dimensional lattices, and \cite[p.~86]{Rei88} notes that the arguments extend to other dimensions. The three-dimensional form is Proposition~\ref{app:G1:reisz3d}; it is used in the proof of Proposition~\ref{prop:ch-CN}.
\end{enumerate}
\noindent Sections~\ref{sec:UV}--\ref{sec:chain} use neither King's cube states \cite[Thms.~2.1--2.3]{K2}, nor any statement on the infinite lattice $\eps\Z^3$ with boundary conditions, nor any infrared renormalisation group; in particular, no bound uniform in the photon mass in a fixed cube is used.

\section{The vacuum bound for the slice models, and three repairs of printed proofs}\label{sec:T6}

\subsection{The slice models and the statement}\label{ssec:T6-statement}
Fix $j\ge1$ and a spacing $\eps=2^{-K}\eps_0$, and write $T:=T_j$. Apart from the statement of Theorem~\ref{thm:T6}, the torus is fixed in this section, and the letter $j$ indexes the fluctuation scales (slices) $H_j$, $\delta_j$, $G^{\eta_k}_{(j)}$; it is not the torus index. Let $-\Delta_\eps$ act on each component of the vector field (on the periodic lattice it is the Hodge Laplacian $d^*d+dd^*$, Lemma~\ref{app:G1:X}), and put $C_\mu:=(-\Delta_\eps+\mu^2)^{-1}$. Let $G_k:=(-\Delta_\eps+\mu_*^2+a_kt_k^{-2}P_k)^{-1}$ be King's block-spin covariances \cite[(2.13), p.~653]{K1} (vector version), and let $H_j:=G_{j+1}-G_j$ denote King's fluctuation covariances \cite[(2.17)]{K1}, with $G_0:=0$ (so $H_0=G_1$, King's $C^{(0)}$ \cite[(2.16), p.~653]{K1}); $H_j\ge0$ by Lemma~\ref{app:G1:slices}. Put
\begin{equation}\label{eq:T6-GammaB}
\Gamma^B:=\sum_{j<K}\delta_jH_j,\qquad \delta_j:=\kappa\mu_*^2t_{j+1}^2\le\bar\delta:=\kappa\mu_*^2\eps_0^2 .
\end{equation}
Here $\kappa>0$ is the constant of the domination inequality $D^<\le\Gamma^B\Pi^\perp$ on the torus, with $D^<$ the operator of Remark~\ref{rem:T6-routes} (Theorem~\ref{app:G1:DOM-nat} and Lemma~\ref{app:G1:struct}(b)); the inequality holds for every $\kappa\ge\kappa_1\kappa_3$, where $\kappa_1\kappa_3$ is the constant of Theorem~\ref{app:G1:DOM-nat}. At $L=2$, $a=1$ and under the cap $\mu_*\eps_0\le0.2$ (condition (E.dom) of Appendix~\ref{app:B}) one has $\kappa_1\kappa_3<7.5$ for every $K$ (Remark~\ref{app:G1:DOM-kappa}), and $\kappa\le7.5$ is fixed with the regulator data, independently of $K$. Since $\sum_{j<K}t_{j+1}^2H_j\le\eps_0^2G_K$, also $\Gamma^B\le\kappa\mu_*^2\eps_0^2\,G_K\le\bar\delta\, C_{\mu_*}$, for every $\kappa>0$ (Lemma~\ref{app:G1:GammaBdom}). For $u>-1/\bar\delta$ the operator $C_B(u):=C_{\mu_*}+u\Gamma^B$ is positive (Lemma~\ref{app:G1:CBpos}), and we call the lattice model with photon covariance $C_B(u)$, scalar coefficient $a_{K,T}(M)$, King's hopping and quartic term, the \emph{slice model} at $(u,M)$. Put
\[
\begin{gathered}
\Phi_M(u):=\log Z\big(C_B(u);a_{K,T}(M)\big),\\
Z(\Gamma;a):=\mathbb E_{N(0,\Gamma)}\big[Z_\varphi(A;a)\big],\qquad Z_\varphi(A;a):=\int \mathrm{e}^{-S^{\mathrm m}_a(A,\varphi)}\prod_{x}d^2\varphi(x),
\end{gathered}
\]
with $S^{\mathrm m}_a$ the matter action \eqref{eq:matteraction}: the photon measure is normalised and the scalar integral is not. Let $\mathcal E(u;M):=T_{\nbar}[\Phi_M(u)]$ be the Taylor polynomial of total order $\nbar=13$ in $(e,\lambda)$ at $(0,0)$, at fixed $M$ and with the coefficient $a_{K,T}(M)$ inside; this coefficient is computed with the photon covariance $C_{\mu_*}$ and does not depend on $u$ (the organisation of \cite[(1.12)--(1.13)]{Ba1}). The order-zero term of $\mathcal E(u;M)$ is $E_{0,s}(M)$. Set $e(\eps;u;M):=-\abs{T}^{-1}\mathcal E(u;M)$.

\begin{remark}[the direction of the deformation]\label{rem:T6-routes}
The translation-invariant direction $D^<:=\int_0^{\eps_0^2}\mathrm{e}^{t\Delta_\eps}(1-\mathrm{e}^{-t\mu_*^2})\,dt\,\Pi^\perp$ satisfies $D^<\le\Gamma^B\Pi^\perp$ (Lemma~\ref{app:G1:struct}(b)), and not conversely. An upper bound on the vacuum in one direction does not imply one in the other; Theorem~\ref{thm:T6} is stated and used with $\Gamma^B$.
\end{remark}

The function $u\mapsto\Phi_M(u)$ is convex on $(-1/\bar\delta,\infty)$: by the current representation (Subsection~\ref{ssec:UV-other}, item~1), $\Phi_M(u)=\log\sum_nw(n)\mathrm{e}^{-\frac12e^2\eps^2\langle n,C_{\mu_*}n\rangle}\mathrm{e}^{-\frac u2e^2\eps^2\langle n,\Gamma^Bn\rangle}$ with $w(n)\ge0$, a logarithmic Laplace transform. The polynomial $e(\eps;u;M)$ has degree at most $\lfloor\nbar/2\rfloor$ in $u$ (Lemma~\ref{app:G1:edeg}).

Recall the parameter window \eqref{eq:window},
\[
\Wwin=\{(e,\lambda,m,\eta,m_1,m_2):\ m>0,\ 0<\lambda<0.1m,\ 0<e^2<0.1m,\ 0<\eta<1,\ 2m<m_1<m_2\}.
\]

\begin{theorem}[the vacuum bound for the slice models]\label{thm:T6}
Let $(e,\lambda,m,\eta,m_1,m_2)\in\Wwin$ and $\mu_*:=m$, and let the regulator data ($L=2$, $p=6$, $\nbar=13$, $\Nbar=169$, $M_b$, $R$, $b_0$, $\eps_0$, $\ell_0$) satisfy the conditions of Appendix~\ref{app:B}. Then there are constants $C_+^B(1;m_2),C_+^B(-\eta;m_1)<\infty$ such that, for every $K\ge K_0$ and every $j\ge1$, with $T=T_j$ and $\eps=2^{-K}\eps_0$,
\begin{equation}\label{eq:T6}
\abs{T}^{-1}\Phi_M(u)+e(\eps;u;M)\le C_+^B(u;M)\qquad\text{at }(u,M)\in\{(1,m_2),(-\eta,m_1)\}.
\end{equation}
\end{theorem}
\noindent The proof of Theorem~\ref{thm:T6} is Proposition~\ref{app:G1:assembly}, which assembles \eqref{eq:T6} from the ingredients of Table~\ref{tab:T6-inputs}; its structure is described after Lemma~\ref{lem:T6-rows}.

\begin{remark}[conventions in Theorem~\ref{thm:T6}]\label{rem:T6-textfixes}
(i) The order-zero term of $\mathcal E(u;M)$ is the scalar factor $E_{0,s}(M)$, because the photon measure of the slice models is normalised. With the Lebesgue-normalised photon of Propositions~\ref{prop:ch-CN} and~\ref{prop:ch-Pplus} it is King's full constant. The two differ by a term that does not depend on $M$ (Lemma~\ref{app:G1:vac}(c)), and it cancels in all differences used.
(ii) As in \cite[(1.30), p.~419]{Ba3}, Ba{\l}aban's final and partial vacuum counterterms are global c-numbers: they carry no localisation function and do not depend on the domains.
(iii) The support hypothesis is used in the form (R4$'''$) of Corollary~\ref{cor:prop1R4}, whose proof covers the components anchored at $\Omega_1$ as well as those anchored at $\supp\tilde A$. The argument of that proof needs $L\le3$: the vertices anchored at $\Omega_1$ have depth about $5\rho_k\approx5r/L$, which exceeds the depth $\frac32r$ needed by the printed step P12 of Subsection~\ref{ssec:T6-prop1} only for $L\le3$; at $L\ge4$ it does not apply.
\end{remark}

\begin{lemma}[two consequences of the conditions]\label{lem:T6-rows}
Let the regulator data satisfy the conditions of Appendix~\ref{app:B} at $L=2$, and let $0<\eta<1$. Then {\rm(a)} $\eta\bar\delta\le\eta/12<1$, and {\rm(b)} $\eps_0\le\min(1/m_2,1/\mu_*,L)$.
\end{lemma}
\begin{proof}
(a) Condition (E.$\delta$) of stage 10a of Appendix~\ref{app:B} is $\bar\delta=\kappa\mu_*^2\eps_0^2\le1/12$; multiply by $\eta<1$. (b) Condition (E.m2), $m_2^2\eps_0^2\le1$, gives $\eps_0\le1/m_2$; condition (E.dom), $\mu_*\eps_0\le0.2$, gives $\eps_0\le1/\mu_*$; and $\eps_0\le1<2=L$ (stage 10a of Appendix~\ref{app:B}).
\end{proof}

\paragraph{Structure of the proof of Theorem~\ref{thm:T6}.} The bound \eqref{eq:T6} is obtained by running Ba{\l}aban's inductive upper-bound argument of \cite[Sects.~2--3]{Ba2} on the torus for the slice model, keeping track of the dependence on $u$: the weights $1+u\delta_j$ act on the photon fluctuations only, and the scalar side carries no $u$-weight. The argument uses the nine ingredients of Table~\ref{tab:T6-inputs}, together with Lemma~\ref{lem:lem24}, Proposition~\ref{prop:prop1}, Corollary~\ref{cor:prop1R4} and Lemma~\ref{lem:IMG}. One of the ingredients, Proposition~\ref{app:G1:I6}, completes Ba{\l}aban's combinatorial estimate \cite[(3.42)--(3.54), pp.~592--594]{Ba2} (Subsection~\ref{app:G1:ssec-I6}). The bounds of all ingredients are uniform in $\eps$, $K$ and $\ell_T$, and finite at $u=-\eta$ and $M=m_1$.

\begin{table}[ht]
\centering\footnotesize
\begin{tabular}{@{}p{0.2\textwidth}p{0.74\textwidth}@{}}
\hline
ingredient & content \\
\hline
Proposition~\ref{app:G1:I1} & the moved vacuum pieces at the scales $1\le k\le K-1$ (good, bad and outside localisations), including their $u$-dependent parts \\
Proposition~\ref{app:G1:I2} & the Ward--Takahashi identities on the torus at orders $\ge2$, through the periodisation of the propagators \\
Proposition~\ref{app:G1:I3} & the conditions and operations of \cite[Prop.~1]{Ba3} as they are used in \cite{Ba2}, including the $\square$-remainder of \cite[p.~433]{Ba3}, the separation step, the two-domain bounds for $\delta G_k$ and the support hypothesis in the form (R4$'''$) \\
Proposition~\ref{app:G1:I4} & in the bounds near Neumann faces, only real background fields are used \\
Proposition~\ref{app:G1:I5} & the order cap $\Nbar=169$ for the expressions \cite[(1.33)]{Ba3}; the tails of \cite[pp.~579, 582]{Ba2} are Propositions~\ref{app:G1:T579} and~\ref{app:G1:I5-T582}, which use it \\
Proposition~\ref{app:G1:I6} & the combinatorial estimate \cite[(3.42)--(3.54), pp.~592--594]{Ba2}: the step to (3.47) with the region entropy, the cube count (3.45) near the boundary, and the collar and shell sets of (3.42); it is also used in the proof of Theorem~\ref{thm:UV-K0plus} \\
Proposition~\ref{app:G1:I7} & the hypothesis $m^2\le O(1)$ of \cite[p.~589]{Ba2}, read as a bound on the rescaled mass; condition (E.md) \\
Proposition~\ref{app:G1:I8} & torus-uniform bounds on the final c-numbers $\delta m^2_{\rm fin}$, $E_{\rm fin}$, under $\eps_0\le\min(1/m_2,1/\mu_*,L)$ (Lemma~\ref{lem:T6-rows}(b)), with $\eps_0$ a power of $L$ (Lemma~\ref{app:G1:N9prem}) \\
Proposition~\ref{app:G1:I9} & the vacuum classes at the last scale \\
\hline
\end{tabular}
\caption{The ingredients of the proof of Theorem~\ref{thm:T6}. All their bounds are uniform in $\eps$, $K$ and $\ell_T$, and finite at $u=-\eta$ and $M=m_1$.}\label{tab:T6-inputs}
\end{table}

\subsection{Repair 1: Ba{\l}aban's Lemma 2.4 of Higgs II at the edge of \texorpdfstring{$\Lambda_7$}{Lambda7}}\label{ssec:T6-lem24}
We work in the $k$-th step of \cite[Chap.~2]{Ba2}, in scale-$k$ units, on the rescaled torus. The fields satisfy the small-field restrictions \cite[(2.55), p.~570]{Ba2} on $\Lambda_{-1}^{(k-1)\prime}$:
\[
\abs{\partial A}\le c_1p(t_{k-1}),\quad \abs{D_{\bar A^{(k)}}\varphi}\le c_1p(t_{k-1}),\quad \abs{\varphi(x)}\le\Upsilon_k:=c_1\big(\lambda t_{k-1}\big)^{-1/4}p(t_{k-1}),
\]
(and the bound on $A$ itself), and the background scalar field of the step is
\begin{equation}\label{eq:lem24-phik}
\varphi^{(k)}=a_kG_k\big(B^k(\Lambda_2^{(k-1)\prime}),A^{(k)}\big)Q_k^*(A^{(k)})\,\Lambda_6^{(k-1)\prime}\varphi
\end{equation}
\cite[(2.56)]{Ba2}: Green's function on $\Omega:=B^k(\Lambda_2^{(k-1)\prime})$, source cut to $\Lambda_6^{(k-1)\prime}$. In this subsection $\Lambda_i$ stands for $\Lambda_i^{(k-1)\prime}$, $r_k:=r(t_k)$, and $c_d:=\sqrt d\,M_b$; the letter $\kappa$ denotes Ba{\l}aban's exponent in the error $O(t_k^\kappa)$ of \cite[(2.67)]{Ba2}, not the constant $\kappa$ of Subsection~\ref{ssec:T6-statement}. From the nesting \cite[(2.8)]{Ba2}, $\dist(\Lambda_7,\Lambda_6^c)\ge\rho_k$, $\dist(\Lambda_6,\Lambda_2^c)\ge4\rho_k$ and $\dist(\Lambda_7,\Lambda_2^c)\ge5\rho_k$.

\cite[Lemma~2.4, p.~572]{Ba2} states: under \cite[(2.55)]{Ba2},
\begin{align}
\varphi^{(k)}(x)&=U\big(A^{(k)}(\Gamma^{(k)}_{x,y})\big)\varphi(y)+O(p(t_k))=(Q_k^*(A^{(k)})\varphi)(x)+O(p(t_k)),\qquad x\in B^k(y),\ y\in\Lambda_7,\label{eq:lem24-265}\\
(D^{\eta_k}_{A^{(k)}}\varphi^{(k)})(b)&=O(p(t_k)),\qquad b\subset B^k(\Lambda_7).\label{eq:lem24-266}
\end{align}
Its only consumer in \cite{Ba2} is the quartic convergence factor \cite[(2.78)--(2.79), p.~574]{Ba2}, a sum over $y\in\Lambda_7^{(k-1)\prime}\cap\Lambda_7^{(k)c}$; this set contains every point of $\Lambda_7$ within $2r_k$ of $\Lambda_7^c$.

\begin{remark}[the printed proof does not cover edge points]\label{rem:lem24-gap}
The printed proof defines $\square_1,\square_2$ as the unions of big blocks \emph{contained in $\Lambda_7$} at distance less than $2r_k$, $4r_k$ from $y$, puts $\square:=B^k(\square_2)$, and asserts \cite[(2.67)]{Ba2}: $\varphi^{(k)}(x)=a_k(G_k(\square,A^{(k)})Q_k^*\square_1\varphi)(x)+O(t_k^\kappa)$ by \cite[Prop.~2.2]{Ba2} and \eqref{eq:lem24-phik}. For $y\in\Lambda_7$ at distance $O(1)$ from $\Lambda_7^c$ this step is not justified as written. (i) The dropped source $\Lambda_6\varphi-\square_1\varphi$ contains $\varphi$ on $\Lambda_6\setminus\Lambda_7$ at distance $O(1)$ from $x$; the kernel $a_kG_kQ_k^*$ is $O(1)$ at distance $O(1)$ (it reproduces constants, \cite[(2.75)]{Ba2}), so this error is $O(\Upsilon_k)$, not $O(t_k^\kappa)$. (ii) The domain change from $\Omega$ to $\square$ happens at distance $O(1)$ from $x$, so the factor of \cite[(2.26)]{Ba1} gives no smallness. (iii) $x$ may lie within $R_0$ of $\partial\square$, where \cite[(2.24)--(2.25)]{Ba1} are not available unless $\square$ is a parallelepiped, and $\Lambda_7\cap(\text{cube})$ has re-entrant edges in general. The two errors in (i) and (ii) cancel to $O(p)$ for slowly varying $\varphi$, so the statement \eqref{eq:lem24-265}--\eqref{eq:lem24-266} of the lemma is not contradicted; it is the step \cite[(2.67)]{Ba2} of the printed proof that is not justified for such $y$. The gap is the same at every $L\ge2$. A box fitting around every $y\in\Lambda_7$ without any cut does not exist: a full $4r_k$-box inside $\Lambda_2$ would need $4r_k\le5\rho_k\approx5r_{k-1}/L$, which is false for small $\eps_0$, since $r_{k-1}/r_k\to1$.
\end{remark}

The printed proof of \cite[Lemma~2.4, p.~572]{Ba2} does not cover the points of $\Lambda_7^{(k-1)\prime}$ near its boundary (Remark~\ref{rem:lem24-gap}); Lemma~\ref{lem:lem24} gives a proof that does.

\begin{lemma}[Ba{\l}aban's Lemma 2.4 of \cite{Ba2} with a complete proof]\label{lem:lem24}
Assume \cite[(2.55)]{Ba2}, the conditions
\begin{equation}\label{eq:lem24-R0row}
\text{(R.edge)}\qquad 5\big(r(\eps_0)-M_b\big)/L\ge R_0+c_d+2M_b,\qquad 4r(\eps_0)\ge R_0+c_d,\qquad R\ge10M_b+10,
\end{equation}
$\eps_0\le\eps(L,\varrho)$ (so that $r(t_{k-1})/r(t_k)\le2L$), the torus-size condition (T.nw) of Appendix~\ref{app:B}, the mass-defect condition (E.md) of stage 10a (used in Step~5 at the Higgs mass $M$ of the step; it is imposed at $m_2$ and holds at every $M\le m_2$, since its left sides increase in $M$), and the upper bound on $\eps_0$
\begin{equation}\label{eq:lem24-cap}
\text{(E.edge)}\qquad C\,a_k\,c_1(\lambda t)^{-1/4}p(t)\,\big(1+r(t)\big)^{d-1}\exp\Big(-\delta_0\Big[\frac{r(t)-M_b}{L}-c_d\Big]\Big)\le(Lt)^\kappa\quad\text{for all }t\le\eps_0/L,
\end{equation}
with $C=C(c_0,\delta_0,d)$. Then \eqref{eq:lem24-265} and \eqref{eq:lem24-266} hold for every $y\in\Lambda_7^{(k-1)\prime}$, every $x\in B^k(y)$ and every bond $b\subset B^k(\Lambda_7^{(k-1)\prime})$, with constants independent of $\eps$, $k$, the torus and the position of $y$ relative to $\partial\Lambda_7^{(k-1)\prime}$. The condition (E.edge) holds for all small $\eps_0$ at every $R>0$.
\end{lemma}

\begin{proof}
Fix $y\in\Lambda_7$ and $x\in B^k(y)$. Put $s_k:=\min\{4r_k,\,5\rho_k-2M_b\}$, let $\square_2$ be the cube of big blocks of sup-radius $s_k$ around the big block of $y$, $\square:=B^k(\square_2)$, and let $\square_1\varphi:=1_{Q(y,s_k/2)}\,\Lambda_6\varphi$, where $Q(y,s)$ is the sup-ball of radius $s$. Since $\dist(\Lambda_7,\Lambda_2^c)\ge5\rho_k$, $\square\subset\Omega=B^k(\Lambda_2)$; $\square$ is a rectangular parallelepiped and a union of big blocks. By \eqref{eq:lem24-R0row} and since $r$ decreases in $t$, $s_k-c_d\ge R_0$ at every $k$.

\emph{Step 1 (source cut).} $\Lambda_6\varphi-\square_1\varphi$ is supported at distance $\ge s_k/2-c_d$ from $x$, and $\dist(x,\Omega^c)\ge5\rho_k-c_d\ge R_0$. By \cite[(2.25)]{Ba1} on $\Omega$, with the sup bound $\Upsilon_k$ on $\varphi$,
\[
\big|a_k\big(G_k(\Omega,A^{(k)})Q_k^*(\Lambda_6\varphi-\square_1\varphi)\big)(x)\big|\le C a_k\,\mathrm{e}^{-\delta_0(s_k/2-c_d)}\,\Upsilon_k .
\]
\emph{Step 2 (domain change).} $G_k(\square,A)=G_k(\Omega,A)+\delta G_k(\square,\Omega,A)$, and by \cite[(2.26)]{Ba1} the second term applied to $Q_k^*\square_1\varphi$ at $x$ carries the factor $\mathrm{e}^{-\delta_0(\dist(Q(y,s_k/2),\square^c)+\dist(x,\square^c))}\le \mathrm{e}^{-\delta_0(3s_k/2-2c_d)}$. The same two estimates hold for $D^{\eta_k}$ by \cite[Prop.~2.2, (2.58)]{Ba2}. Hence
\begin{equation}\label{eq:lem24-267}
\varphi^{(k)}(x)=a_k\big(G_k(\square,A^{(k)})Q_k^*(A^{(k)})\square_1\varphi\big)(x)+O\big(a_k\Upsilon_k \mathrm{e}^{-\delta_0(s_k/2-c_d)}\big),
\end{equation}
and likewise for the covariant derivative: this is \cite[(2.67)]{Ba2} with the boxes chosen above.

\emph{Step 3 (constant background).} Let $A_0:=A^{(k)}(y)$. Joining points of the convex box $\square$ by straight lattice paths and using \cite[(2.60)]{Ba2} on $B^k(\Lambda_2)\supset\square$, $\abs{A^{(k)}(z)-A_0}\le O(p)\sqrt d\,s_k\le O(p\,r_k)$ on $\square$. The expansion \cite[(3.44)]{Ba1} and \cite[Prop.~2.2]{Ba1} on the parallelepiped $\square$ replace $A^{(k)}$ by $A_0$ with an error $O(t_k^{\kappa_0})$ (Ba{\l}aban's exponent in \cite[(2.68)]{Ba2}), as in \cite[(2.68)]{Ba2}.

\emph{Step 4 (gauge-out).} The constant field is a gradient on the box, $A_0=\partial^{\eta_k}\lambda_0$ with $\lambda_0(z)=A_0\cdot(z-y)$, so $U(A_0(\Gamma_{z,y}))=\mathrm{e}^{ieA_0\cdot(z-y)}$ for every contour $\Gamma_{z,y}\subset\square$, and $G_k(\square,\partial\lambda_0)=U(\lambda_0)G_k(\square,0)U(\lambda_0)^*$, $Q_k(\partial\lambda_0)=U(\lambda_0)_{\rm coarse}Q_k(0)U(\lambda_0)^*$ (Neumann Laplacian on bonds of $\square$; contours inside blocks contained in $\square$). This is legitimate because $\square$ does not wrap around the torus: its diameter is at most $2s_k+M_b\le8r_k+M_b$, below the side of the torus by (T.nw). Write $\varphi'(z):=U(\lambda_0)(z)^{-1}\varphi(z)$.

\emph{Step 5 (reproduction of constants).} By \cite[(2.75)]{Ba2}, $a_kG_k(\square,0)Q_k^*1=1-m^2t_k^2G_k(\square,0)Q_k^*1$. Hence
\begin{multline*}
a_k\big(G_k(\square,0)Q_k^*\square_1\varphi'\big)(x)-\varphi'(y)=a_k\big(G_k(\square,0)Q_k^*[\square_1(\varphi'-\varphi'(y))]\big)(x)\\
-\varphi'(y)\,a_k\big(G_k(\square,0)Q_k^*1_{\square\setminus\square_1}\big)(x)-\varphi'(y)\,m^2t_k^2\big(G_k(\square,0)Q_k^*1\big)(x),
\end{multline*}
where $1_{\square\setminus\square_1}$ is the indicator of $\square\setminus(Q(y,s_k/2)\cap\Lambda_6)$. The first term is $O(p)$: on the convex cube $Q(y,s_k/2)\subset\Lambda_{-1}$ the restriction $\abs{D\varphi}\le c_1p$ of \cite[(2.55)]{Ba2} gives $\abs{\varphi'(z)-\varphi'(y)}\le O(p)\abs{z-y}$, and the kernel decays exponentially (the same reasoning as in \cite[Lemma~2.3]{Ba2}). The second term: $\square\setminus\square_1\subset(Q(y,s_k)\setminus Q(y,s_k/2))\cup(Q(y,s_k/2)\setminus\Lambda_6)$ lies at distance $\ge\min\{s_k/2,\rho_k\}-c_d$ from $x$, so it is $\le c_0\mathrm{e}^{-\delta_0(\min\{s_k/2,\rho_k\}-c_d)}\Upsilon_k$ by Proposition~\ref{prop:UV-BaProp21} on the parallelepiped $\square$. The third term is the mass-defect term of the printed proof; it is $O(p(t_k))$ under (E.md), as in the printed proof.

\emph{Step 6 (derivatives).} For $b\subset B^k(y)$, $b$ lies at distance $\ge s_k-c_d\ge R_0$ inside $\square$ and $\ge5\rho_k-c_d$ inside $\Omega$, so the printed derivative bounds apply at $b$ and Steps 1--5 give \eqref{eq:lem24-266} in the form \cite[(2.77)]{Ba2}. No re-entrant edge is approached.

\emph{Collection of the errors.} The inequality $s_k/2\ge\rho_k$ is equivalent to $\rho_k\le2r_k$ and $3\rho_k\ge2M_b$; the first holds because $r(t_{k-1})/r(t_k)\le2L$, the second because $R\ge10M_b+10$. So every error above is at most $C a_k\Upsilon_k(1+r_k)^{d-1}\mathrm{e}^{-\delta_0(\rho_k-c_d)}$, the polynomial factor coming from the sums of exponentially decaying kernels over sources at distance $D=O(r_k)$. With $t=t_{k-1}$ and $\rho_k=(r(t)-M_b)/L$, this is $\le(Lt)^\kappa=t_k^\kappa$ by (E.edge), which is the size of the error in \cite[(2.67)]{Ba2}; the rest of the printed proof, \cite[(2.69)--(2.77)]{Ba2}, now runs on the box $\square$. Finally, the logarithm of the prefactor in \eqref{eq:lem24-cap} is $O(\log t^{-1})+O(\log\log t^{-1})$ while the exponent is $\delta_0R(1+\log t^{-1})^\varrho/L$ with $\varrho>1$, so \eqref{eq:lem24-cap} holds for $t$ small, at every $R>0$.
\end{proof}

\begin{remark}[where the repair is used]\label{rem:lem24-uses}
In the proof of Theorem~\ref{thm:T6}, Lemma~\ref{lem:lem24} is used in the quartic lower bound \cite[(2.78)]{Ba2} at points of $\Lambda_7^{(k-1)\prime}$ near its boundary, whenever a large-field region is present. The gauge-out on a full cube used in \cite[p.~433]{Ba3} and the later lemma \cite[Lemma~2.7]{Ba2} do not use the cut at $\Lambda_7$ and are not affected. The condition (E.edge) depends only on $L$, $R$ and data fixed before them, not on $u$, $M$, $\eps$ or $\ell_T$; it is downward closed in $\eps_0$, and it is a condition of stage 10b of Appendix~\ref{app:B}. Lemma~\ref{lem:B:rows}(a) lists the data on which it depends, and shows that the conditions (R.edge) follow from the lower bound $R\ge R^*$ at $L=2$ and that $r(t_{k-1})/r(t_k)\le2L$ holds for every $\eps_0\le1$.
\end{remark}

\subsection{Repair 2: the sketched steps of Ba{\l}aban's Proposition~1 of Higgs III}\label{ssec:T6-prop1}
We use the setting of Subsection~\ref{ssec:UV-prop1}, in scale-$k$ units, with $r:=r(t_k)$, $s_b\le2M_b$ the rounding of a set to big blocks, $\ell_{\Nbar}:=2V(\Nbar)$ the side of a cube containing the localisations of a connected class with at most $\Nbar$ vertices (\cite[p.~432]{Ba3}), and $R_0$ the constant of Proposition~\ref{prop:UV-BaProp21}. The \emph{depth} of a point $x$ is $\dist(x,\Omega^c)$. At \cite[p.~433]{Ba3} the proof takes a cube $\square_1$ of size $r_\square$ (printed: $r_\square=r$) containing the localisations of the old vertices of a class and the concentric cube $\square$ of size $3r_\square$, and splits
\begin{equation}\label{eq:prop1-P12}
G_k(\Omega,\tilde B)=G_k(\square,\tilde B)+\delta G_k(\Omega,\square,\tilde B)\qquad[\text{or the same with }\Omega_2\text{ in place of }\Omega],
\end{equation}
treating $\delta G_k(\Omega,\square,\tilde B)$ as an external field. We call this printed step P12. Its precondition $\square\subset\Omega$ (respectively $\square\subset\Omega_2$) is not written in \cite{Ba3}.

The proof of \cite[Prop.~1, pp.~420--445]{Ba3} contains sketched steps; Proposition~\ref{prop:prop1} writes each of them out and determines the depth inside the localisation domain at which it is applied.

\begin{proposition}[the depth of each step of Ba{\l}aban's Proposition~1]\label{prop:prop1}
\begin{enumerate}
\item[(a)] \emph{(Inventory.)} The proof of \cite[Prop.~1]{Ba3}, \cite[pp.~422--445]{Ba3}, contains twenty sketched steps, S1--S11 and S13--S21 of Table~\ref{tab:prop1}, and the printed step P12 with its unwritten precondition. Besides these, a page-by-page reading of \cite[pp.~414--445]{Ba3} finds seven further sketched places (the $\Omega_2$ branch of P12; the identification of torus and infinite-lattice propagators in the counterterms; the repetition on $\square$ of S1 and S9 with $B'$ in place of $\tilde A$; three phrases of power counting and bookkeeping; the generation of all Ward--Takahashi identities); none of them uses depth beyond P12 or $R_0+O(1)$.
\item[(b)] \emph{(Anchored steps.)} The steps S1, S2, S3, S4, S9 apply Proposition~\ref{prop:UV-BaProp21} only at points of depth at least $R_0+12$ near the localisation of a vertex ($R_0+13$ at the first step $k=1$, where $\eta_k=\frac12$; S2: depth measured in $\Omega_2$), or at points of $U_c(\supp\tilde A)$ of depth $\ge R_0+2\sqrt3+1$ (S1).
\item[(c)] \emph{(The cube step.)} If an old vertex $x_v$ of the class has depth $>\frac32r_\square+2s_b+\ell_{\Nbar}$, then $\square\subset\Omega$ and every point of $\square_1$ lies at distance $\ge r_\square$ inside $\square$. The same holds with $\Omega_2$ in place of $\Omega$.
\item[(d)] \emph{(Depth-free steps.)} Once $\square\subset\Omega$, every identity and inequality used after the split \eqref{eq:prop1-P12}, on \cite[pp.~433--445]{Ba3}, is independent of $\dist(\cdot,\Omega^c)$ (Lemma~\ref{lem:Omegafree}); in particular S5--S8, S10, S11, S13--S21 need no depth.
\end{enumerate}
Consequently every sketched step of the proof needs depth at most $O(R_0+\ell_{\Nbar}+M_b)$, with a constant independent of $r$, and the only step needing depth of order $r$ is the printed step P12, with need $\frac32r_\square+2s_b+\ell_{\Nbar}$.
\end{proposition}

\begin{table}[ht]
\centering\footnotesize
\begin{tabular}{@{}p{0.05\textwidth}p{0.08\textwidth}p{0.47\textwidth}p{0.32\textwidth}@{}}
\hline
step & page & omitted argument & depth needed \\
\hline
S1 & 414 & kernel bound of $[G_k(\Omega,\tilde B)V_k]^nG_k(\Omega,\tilde A+\tilde B)[V_kG_k]^{n'}$, \cite[(1.16)]{Ba3} & $R_0+2\sqrt3+1$ on $U_c(\supp\tilde A)$; $R_0+O(1)$ at vertices \\
S2 & 414, 424 & $\delta G_k(\Omega,\Omega_2,\tilde B)$ as an external field, \cite[(2.5)]{Ba3} & $R_0+12$ inside $\Omega_2$ ($R_0+13$ at $k=1$) \\
S3 & 424, 426 & slice decomposition \cite[(2.6)]{Ba3} and slice bounds (2.10)--(2.12) & $R_0+1$ at line ends \\
S4 & 425 & summation by parts \cite[(2.8)]{Ba3} & $R_0+1$ \\
S5--S7 & 428--429 & induction (2.16); generalised graphs (Prop.~2.2); power counting, Cor.~2.3 & none \\
S8 & 431 & Ward--Takahashi identities for free boundary conditions as limits of periodic ones & none \\
S9 & 432 & bound \cite[(3.1)]{Ba3} between separated localisations & $R_0+5$ at vertex cubes \\
S10, S11 & 432 & regrouping of classes; localisation of a component in a cube of side $\ell_{\Nbar}$ & none \\
P12 & 433 & printed split \eqref{eq:prop1-P12}; precondition $\square\subset\Omega$ (or $\Omega_2$) unwritten & $\frac32r_\square+2s_b+\ell_{\Nbar}$ at an old vertex \\
S13 & 433 & constant field $B_0$ with $\abs{\tilde B-B_0}\le O(r_\square p)$ on $\square$ & as for P12 \\
S14--S16 & 433 & $B'$-expansion; old vertices in $\square_1$; chains by Prop.~2.1 on $\square$ & none beyond $\square\subset\Omega$ \\
S17, S18 & 433 & gauge covariance on $\square$; gauge-out of $B_0$ & none \\
S19 & 433 & bound on $\delta G_k(\square,\eta_k\Z^d,0)$ & none \\
S20 & 434 & comparison of the norms (1.32) at $B_0$ and at $\tilde B$ & none (a size factor, see Remark~\ref{rem:prop1-cited}) \\
S21 & 434--445 & renormalisation of the divergent classes & none \\
\hline
\end{tabular}
\caption{The sketched places in the proof of \cite[Prop.~1]{Ba3}. Pages are journal pages of \cite{Ba3}.}\label{tab:prop1}
\end{table}

\begin{proof}[Proof of Proposition~\ref{prop:prop1}]
\emph{(b), S3.} The slice decomposition \cite[(2.6)]{Ba3} is the iterated composition formula \cite[(2.43)]{Ba1}, an identity on any union of scale-$k$ blocks, with no depth. For \cite[(2.10)--(2.11)]{Ba3}, rescale to the $L^{-j}$ lattice; $\Omega$ is a union of scale-$j$ big blocks and $\tilde B$ is regular at scale $j$. For a coarse point $y_0$, $f:=Q_j^*\delta_{y_0}$ has $\norm f_\infty\le1$ and support $B^j(y_0)$, so \cite[(2.25)]{Ba1} at an evaluation point $y$ gives $\abs{(G_jQ_j^*)(y,y_0)}+\abs{(DG_jQ_j^*)(y,y_0)}\le c_0\mathrm{e}^{-\delta_0\dist(y,B(y_0))}$ provided $\dist(y,\Omega^c)\ge R_0$ in scale-$j$ units, for every $y_0$, including blocks touching $\partial\Omega$. The adjoint factor is bounded in the same way at the other end, and $C^{(j)}$ by \cite[(2.34)]{Ba1} without restriction. The sums over the coarse points converge with constants independent of $\eps$ and $r$. Rescaling back, the restriction is $\dist(\{x,x'\},\Omega^c)\ge R_0L^{j-k}\le R_0$ at the two line ends. The vector propagator lives on the torus \cite[p.~414]{Ba3} and needs no depth. Allowing one block for the contour legs of (2.12), the need is $R_0+1$.

\emph{(b), S4.} In \cite[(2.8)]{Ba3} the derivative is moved by summation by parts; no boundary term arises because the moved leg carries the localisation function $g$, supported in $\Omega_1$, at distance $\ge1$ from $\Omega^c$. After the move the derivative sits on the vector line (torus, no depth), on $g$, or on a scalar line end at the vertex: need $R_0+1$.

\emph{(b), S9.} \cite[(3.1)]{Ba3} bounds $\norm{hG_k(\Omega,\tilde B)h'}_{1,\alpha}$ for $\dist(\square(v),\square(v'))\ge1$. Summing \cite[(2.10)--(2.11)]{Ba3} over $j$ with one derivative at each end, the factors $(L^j\eta_k)^{-d-\alpha}\mathrm{e}^{-\delta_1\abs{x-x'}/(L^j\eta_k)}$ are summable because $\abs{x-x'}\ge1$, and the mixed derivative is pointwise bounded since the points are separated. The evaluation points lie within a collar of width $<\frac12$ around the vertex cubes: need $\le R_0+5$.

\emph{(b), S2.} The two-variable mixed norm of $\delta G_k(\Omega,\Omega_2,\tilde B)$ in \cite[(2.5)]{Ba3}, where the two ends may coincide, is bounded by the two-domain theorem, Proposition~\ref{app:G1:twodomain} (two cut-offs and a weighted $\ell^2$ bound for the middle factor), valid at points at distance $\ge R_0+11$ from $\Omega_2^c$ and supports at distance $\ge R_0+12$ when $\eta_k\le\frac14$ ($k\ge2$), and at $R_0+12$ and $R_0+13$ at $k=1$, $\eta_1=\frac12$, where the offsets of the cut-offs below one block use up the slack. The factor $\mathrm{e}^{-\delta_0\dist(\Omega_2,\partial\Omega)}$ of \cite[(2.5)]{Ba3} is obtained in the reading in which the distance is that of the vertex localisations to $\Omega_2^c$; read literally, as the distance between the sets $\Omega_2$ and $\partial\Omega$, it is not implied, and it fails for localisations near $\partial\Omega_2$ deep inside $\Omega$, where $\delta G_k$ is of order one. No consumer uses more than a bound of order one. In the proof of Theorem~\ref{thm:T6} these lines occur only in terms all of whose vertices lie in $\Lambda_7^{(k)}$ \cite[p.~580, Prop.~2.6]{Ba2}, at depth $\ge5r-O(1)$ in $\Omega_2$.

\emph{(b), S1.} The interaction $V_k(\tilde A,\tilde B)$ is supported in $U_c(\supp\tilde A)\times U_c(\supp\tilde A)$, $c\le2\sqrt3+1$ (nearest-neighbour covariant Laplacian; $P_k$ with contours inside unit blocks). All internal summation points therefore have depth $\ge\dist(\supp\tilde A,\Omega^c)-c$. After summation by parts each middle factor is of the form $\mathfrak DG_k\mathfrak D^*$, whose kernel has a diagonal singularity of order $\eta_k^{-d}$, so a pointwise sum is not uniform in $\eta_k$. Proposition~\ref{app:G1:sandwich} proves the bound by an $L^2$ sandwich with a weighted $\ell^2$ estimate valid for any boundary condition (Lemma~\ref{app:G1:WK-W}); it requires depth $\ge R_0+2\sqrt3+1$ at the internal points and $R_0+O(1)$ at the outer points. The power $(e\,t_k^{1-\alpha})^{n+n'}$ of \cite[p.~414]{Ba3} is a size bound ($\abs{e\tilde A}\le O(e\,p)$) and needs no depth.

\emph{(c).} The localisations of the old vertices lie in a cube of side $\ell_{\Nbar}$ containing $x_v$ (S11), and $\square_1$, $\square$ are concentric cubes of sides $r_\square$, $3r_\square$, rounded outward to big blocks. In the sup metric every point of $\square$ is within $\frac32r_\square+s_b$ of the common centre, and the centre is within $\ell_{\Nbar}+s_b$ of $x_v$. If the depth of $x_v$ exceeds $\frac32r_\square+2s_b+\ell_{\Nbar}$, the triangle inequality gives $\square\subset\Omega$; the points of $\square_1$ are then at distance $\ge r_\square$ from $\square^c$. With $\square\subset\Omega$, \cite[(2.26)]{Ba1} for the pair $(\square,\Omega)$, and the two-variable form of Proposition~\ref{app:G1:twodomain} at points at distance $\ge R_0+11$ inside $\square$ ($R_0+12$ at $k=1$), bound $\delta G_k(\Omega,\square,\tilde B)$ at points of $\square_1$; since $\square$ is a parallelepiped there is no further restriction, and \cite[(3.2)]{Ba3} uses only this boundedness. The argument for $\Omega_2$ is identical.

\emph{(d), S13, S17--S19.} S13: for $x\in\square$ and $x_c$ the centre, $\tilde B_\mu(x)-\tilde B_\mu(x_c)$ is a sum of $\eta_k\partial^{\eta_k}\tilde B_\mu$ along a coordinate path in the convex box, so $\abs{\tilde B-B_0}\le d(\frac32r_\square+s_b)\sup_\square\abs{\partial^{\eta_k}\tilde B}=O(r_\square p)$ with $B_0:=\tilde B(x_c)$. S17: for $\lambda:\square\to\R$, the Neumann covariant Laplacian on bonds of $\square$ and the averaging $Q_k$ (contours in blocks contained in $\square$) are covariant, hence $G_k(\square,B+\partial\lambda)=U(\lambda)G_k(\square,B)U(\lambda)^*$, and every class is invariant when the external scalar fields are transformed simultaneously. S18: $B_0$ constant gives $B_0=\partial^{\eta_k}\lambda_0$, $\lambda_0(x)=B_0\cdot(x-y)$, so the gauge-out of \cite[(2.68)]{Ba2} is exact on $\square$ provided $\square$ does not wrap around the torus (a condition on the size of the torus). S19: \cite[p.~433]{Ba3} writes $G_k(\square,0)=G_k(0)+\delta G_k(\square,\eta_k\Z^d,0)$ and treats $\delta G_k$ as an external field. Since $\square$ does not wrap (S18), it is a rectangular parallelepiped of $\eta_k\Z^d$ and a union of big blocks. \cite[p.~433]{Ba3} bounds $\delta G_k(\square,\eta_k\Z^d,0)$ by \cite[(1.11)--(1.12), p.~573]{Ba4} with $\Omega=\square\subset\Omega_0=\eta_k\Z^d$, the setting of \cite[p.~572]{Ba4}; since the proof in \cite[Sect.~2]{Ba4} is written for a finite $\Omega_0$, we argue on finite tori. Let $T'\supset\square$ be a torus whose side $\ell'$ is a multiple of $2M_b$, at least $\max\{3M_b,3r_\square+2s_b+2\}$, with big-block grid compatible with $\square$. At $A=0$ the operator is translation invariant and $G_k(T',0)(x,x')=\sum_{n\in\Z^d}G_k(\eta_k\Z^d,0)(x,x'+n\ell')$ (uniqueness of the periodic inverse, as in Lemma~\ref{lem:IMG}(i)), so
\[
\delta G_k(\square,\eta_k\Z^d,0)=\delta G_k(\square,T',0)+\sum_{n\ne0}G_k(\eta_k\Z^d,0)(\cdot,\cdot+n\ell').
\]
The first term is bounded by \cite[(2.26)]{Ba1} on $T'$, with constants independent of the side of $T'$ (Lemma~\ref{app:G1:torusuniform}): it, its derivatives and its H\"older quotients are at most $c_0\exp[-\delta_0(\dist(\{x,x'\},\supp f)+\dist(\supp f,\square^c)+\dist(\{x,x'\},\square^c))]\norm f_\infty$ for all $x,x'\in\square$, with no restriction on $x,x'$ because $\square$ is a parallelepiped. The second term tends to zero as $\ell'\to\infty$, by the exponential decay of $G_k(0)$ (the bound (H$_\infty$) of Subsection~\ref{ssec:T6-IMG}). Hence the same bound holds for $\delta G_k(\square,\eta_k\Z^d,0)$. At points of $\square_1$, at distance $\ge r_\square$ from $\square^c$, $\delta G_k(\square,\eta_k\Z^d,0)$ is therefore an external field of $(1,\alpha)$-norm $\le c_0\mathrm{e}^{-\delta_0r_\square}$. The remaining steps are covered by Lemma~\ref{lem:Omegafree}.

\emph{Conclusion.} Collecting (b)--(d) with Table~\ref{tab:prop1}: the sketched steps need at most $R_0+12$ ($R_0+13$ at $k=1$); P12 needs $\frac32r_\square+2s_b+\ell_{\Nbar}$; nothing else depends on $\dist(\cdot,\Omega^c)$. The inventory (a) is the result of a page-by-page reading of \cite[pp.~414--445]{Ba3}; the list, with the seven further places, is Lemma~\ref{app:G1:inventory}.
\end{proof}

\begin{lemma}[$\Omega$-freeness after the cube split]\label{lem:Omegafree}
Assume $\square\subset\Omega$ for the class. Then every identity and inequality on \cite[pp.~433--445]{Ba3} after the split \eqref{eq:prop1-P12} is a statement about a kernel built from: (i) vertex coefficients localised by the functions $g_k$, smooth partitions of unity with supports in cubes of side $2$, and the unit-cube partition of $\Omega_1$; (ii) counterterms with the same localisation and torus propagators; (iii) $G_k(0)$ (on $\eta_k\Z^d$ or the torus) and its slices, the vector propagator on the torus, and $C^\xi$; (iv) the cubes $\Delta(v)$ of side $L^j\eta_k$; and of external fields whose $(1,\alpha)$-norms at $\square_1$ enter linearly. None of these statements involves $\dist(\cdot,\Omega^c)$. The dependence on $\Omega_1$ is used only through Lipschitz bounds of the smooth localisation functions and through exact matching of graphs and counterterms.
\end{lemma}
\begin{proof}
After the split, every old vertex lies in $\square_1\subset\square\subset\Omega$ and the new vertices of the $B'$-expansion lie in $\square$, so the characteristic function of $\Omega$ is identically $1$ wherever a vertex or a line end occurs. Every object that depends on $\Omega$ ($\delta G_k(\Omega,\Omega_2,\tilde B)$, the operator \cite[(1.16)]{Ba3}, $\delta G_k(\Omega,\square,\tilde B)$, $\delta G_k(\square,\eta_k\Z^d,0)$, the gauge factors) has been moved into the external fields. Inspecting \cite[pp.~433--445]{Ba3}: every propagator is $G_k(0)$, a slice $G^{\eta_k}_{(j)}(0)$, $G^{\eta_k}_{j_0}(0)$, $G^\xi$, $C^\xi$ or $\Pi$; every sum runs over $\eta_k\Z^d$, $\xi\Z^d$ or the cubes $\Delta(v)$; Taylor's formula \cite[(3.10)]{Ba3} is applied to legs and to the smooth factors $g,g'$ (\cite[(3.23), (3.26), (3.31), (3.34), (3.36)]{Ba3}); and the localisation of the $\lambda$-vertices is matched exactly against counterterms carrying the same localisation (\cite[(3.19)]{Ba3}, p.~444, p.~445). None of these operations refers to a distance to $\Omega^c$.
\end{proof}
\noindent The kernels of Lemma~\ref{lem:Omegafree} do depend on $\Omega_1$ and on the torus; what they do not depend on is the depth.

\begin{corollary}[Proposition~1 of \cite{Ba3} under a weakened support hypothesis]\label{cor:prop1R4}
Let $L=2$ and $R\ge R^{**}$ (stage 5 of Appendix~\ref{app:B}); in particular $R$, and $r(t_k)\ge R$ at every scale, satisfy the condition (R.e) of \eqref{eq:eppp} below and the condition $R\ge3R_0+15$, which is part of $R\ge R^*$. Then \cite[Prop.~1]{Ba3} holds with (R4) replaced by
\begin{equation}\label{eq:R4ppp}
\text{(R4$'''$)}\qquad \dist(\supp\tilde A,\partial\Omega)>\tfrac74r+2s_b+\ell_{\Nbar}+1 .
\end{equation}
Moreover (R4$'''$) holds at every use of \cite[Prop.~1]{Ba3} in the proof of Theorem~\ref{thm:T6} under the condition
\begin{equation}\label{eq:eppp}
\text{(R.e)}\qquad R\ge4\ell_{\Nbar}+32M_b+25 .
\end{equation}
\end{corollary}

\begin{proof}
By Proposition~\ref{prop:prop1} the only consumer of depth of order $r$ is P12, with need $\frac32r+2s_b+\ell_{\Nbar}$ at an old vertex $x_v$ (printed cube, $r_\square=r$); every other step needs at most $R_0+13$ near the vertices or near $\supp\tilde A$ ($R_0+12$ for $k\ge2$). Every vertex of a class is anchored: it lies within distance $1$ of $\supp\tilde A$ or of $\supp g_k\subset\Omega_1$ (Lemma~\ref{app:G1:anchor}). A vertex anchored at $\supp g_k$ has depth at least $\dist(\Omega_1,\Omega^c)-1\ge5\rho_k-O(M_b)\ge\frac52(r(t_{k-1})-M_b-1)-4$ at $L=2$. A vertex anchored at $\supp\tilde A$ has, under (R4$'''$), depth $>\frac74r+2s_b+\ell_{\Nbar}$. Both exceed the need $\frac32r+2s_b+\ell_{\Nbar}$ of P12 (for the first, by \eqref{eq:eppp}), and both exceed $R_0+2\sqrt3+13$, since $r\ge R\ge3R_0+15$; the internal points of S1 have depth $>\frac74r-2\sqrt3-1\ge R_0+2\sqrt3+1$. At the uses in the proof of Theorem~\ref{thm:T6} the supply is at least $2r-M_\theta-2M_b-5$ with $M_\theta\le LM_b$ (Lemma~\ref{app:G1:supply}), which exceeds the right side of \eqref{eq:R4ppp} if $r/4>\ell_{\Nbar}+2s_b+M_\theta+2M_b+6$; with $s_b\le2M_b$ and $L=2$ this follows from \eqref{eq:eppp}, since $r(t)\ge R$.
\end{proof}

\begin{remark}[what is used as printed, and Ba{\l}aban's margin in Higgs II]\label{rem:prop1-cited}
(i) The renormalisation of the divergent classes on \cite[pp.~434--445]{Ba3} (S21) is used as printed; by Lemma~\ref{lem:Omegafree} it does not interact with the weakened hypothesis. The step at \cite[(3.35)]{Ba3} is used as printed; its torus form is Theorem~\ref{app:G1:VSR-T335}. (ii) The norm comparison S20 costs a factor $(1+O(r_\square p))^{\#\text{legs}}$, not $O(1)$ as printed (``exactly''); this is a size factor, polylogarithmic in $t_k$ at $r_\square=r$, absorbed by the slack $(L^k\eps)^{(4-d)/4-\beta}$ of the applications \cite[p.~421]{Ba3}. (iii) In \cite[p.~579]{Ba2} Ba{\l}aban applies Proposition~1 at the margin $2r-M-O(1)$; this margin is what his nesting supplies, not what the argument needs: the need, through P12, is $\frac32r+O(\ell_{\Nbar}+M_b)$. (iv) The identification of the torus propagators of the counterterms with the infinite-lattice propagators of the classes is Lemma~\ref{lem:IMG}. (v) Besides the steps of Table~\ref{tab:prop1}, the following printed statements are used as they are printed: \cite[Prop.~2.1, (2.57)]{Ba2}, whose printed proof refers to \cite[Prop.~1]{Ba3}; Ba{\l}aban's bound \cite[(3.3)]{Ba3} with its proof \cite[pp.~433--445]{Ba3}, in the coloured form of Theorem~\ref{app:G1:WP-WP1}; the comparison \cite[(2.68)]{Ba2} in the proof of \cite[Lemma~2.4]{Ba2}; the derivative bounds stated after \cite[(2.90) and (2.91), p.~575]{Ba2}; and the list of vertices \cite[(1.6)--(1.10)]{Ba3} at zero background.
\end{remark}

\subsection{Repair 3: the region entropy in Ba{\l}aban's estimate (3.42) of Higgs II}\label{app:G1:ssec-I6}
The printed proof of the combinatorial estimate \cite[(3.42), p.~592]{Ba2}, the last step of the proof of Theorem~\ref{thm:UV-K0plus}, is incomplete: the step to (3.47) is not written out, the cube count (3.45) fails near the boundary of the region, and (3.51)--(3.54) bound only the large-field sets while (3.42) also contains collar and shell sets; Proposition~\ref{app:G1:I6} completes it.

\emph{What is not written in \cite{Ba2}.} \cite[p.~592]{Ba2} reduces (2.117) to the combinatorial inequality (3.42), with $\zeta''_{\Lambda_0^{(k)}}$ the sum (3.41) over admissible minimal families. On p.~593 he bounds $\zeta''$ by (3.46) and then writes ``After easy transformations we get'' the bound (3.47), which is a supremum over the region sequences. Two steps are not written there: replacing the sum over region sequences by a supremum, and gathering the collar and shell sets of (3.42) into $O(1)(1+\log t_k^{-1})\abs{\Lambda_0^{(k)c}}$. Moreover the cube count (3.45) is false near $\partial\Lambda_7^{(k-1)\prime}$ (Remark~\ref{app:G1:I6-345}).

\emph{Conventions.} In this subsection $d=3$ (the arguments hold for every $d$); $T$ is Ba{\l}aban's periodic lattice torus of side $\ell_T$, $\abs T=\ell_T^d$; $L=2$; $\eps=2^{-K}\eps_0$, $t_k:=L^k\eps$ ($t_K=\eps_0$); $\hat\ell_k:=1+\ln t_k^{-1}$; $r_k:=r(t_k)=R\hat\ell_k^{\varrho}$, $\varrho=2$; $p(t)=b_0(1+\ln t^{-1})^{p}$, $p=6$; $M_b$ is Ba{\l}aban's block parameter. The letters $k$ and $j$ are scale indices.

\emph{Notation.} $T^{(k)}$ is the unit lattice at scale $k$, a discrete torus with $N_k:=\abs T t_k^{-d}$ points. Large blocks of $T^{(k)}$ have side $s_b:=LM_b$ lattice units (\cite[p.~607]{Ba1}), and $\vartheta:=s_b/R$, so $s_b\le\vartheta r_k$. The prime operation $X\mapsto X'$ of \cite[(1.16), p.~606]{Ba1} commutes with complements on unions of $L$-blocks. $B_\infty(y,h)$ is the $\ell^\infty$-ball and $N_w(X)$ the $\ell^\infty$-neighbourhood of width $w$ on the torus; we only use $\dist_\infty\le\dist\le\sqrt d\,\dist_\infty$. A family $F_k=(P_v^{(k)},\dots,R_s^{(k)})$ has bad union $\mathcal S_k:=B(P_v^{(k)})\cup\dots\cup R_s^{(k)}$ and size $\abs{F_k}:=\abs{P_v^{(k)}}+\dots+\abs{R_s^{(k)}}$. We write $X_k:=\Lambda_0^{(k)c}$, $\Lambda_i^{(-1)\prime}:=T^{(0)}$, $\mathrm{Col}_k:=\Lambda_7^{(k-1)\prime}\cap\Lambda_7^{(k)c}$ and $\mathrm{Sh}_k:=\Lambda_5^{(k-1)\prime}\cap\Lambda_5^{(k)c}$. Finally $c_7,c_5\ge0$ denote the two $O(1)$'s on the left of (3.42), $\kappa_0>0$ the exponent there, and $c_\square:=\min\{\tfrac18a3^{-d},\tfrac14\gamma_03^{-d},3^{-d}\}$ is Ba{\l}aban's constant $c_0$ of \cite[p.~593]{Ba2} (renamed to avoid a clash with the constant $c_0$ of Proposition~\ref{prop:UV-BaProp21}). The collar exponent is $\kappa_0=\frac15$ (Lemma~\ref{app:G1:kappa0}).

\emph{Definitions, as printed.} (D0) \cite[(2.7)--(2.8), p.~558]{Ba2}: $\Lambda_0^c$ is the union of the large blocks at distance $<r$ from the bad union, and $\Lambda_{i+1}^c$ is the union of the large blocks at distance $\le r$ from $\Lambda_i^c$. (Dk) \cite[p.~570]{Ba2}: $\Lambda_0^{(k)}$ is the union of those large blocks of $T^{(k)}$ that are contained in $\Lambda_7^{(k-1)\prime}$ and are at distance $>r_k$ from $\mathcal S_k$. \cite[p.~566]{Ba2}: these sets are unions of big blocks, $\Lambda_0^{(j+1)}\subset\Lambda_7^{(j)\prime}$, and $\Lambda_i^{(j)}$ is defined as $\Lambda_i$ with $r(t_j)$. (B-C): $\ell_T/(M_bt_K)\in\N$, so every $T^{(k)}$ with $k\le K-1$ is a union of large blocks. On the tori $T_j$, (B-C) holds by Lemma~\ref{lem:sizes}(ii).

\begin{lemma}[elementary consequences]\label{app:G1:I6-E}
For every admissible sequence and $k\le K-1$:
(a) every point of an element of an admissible $F_k$ lies in $X_k$;
(b) $\Lambda_i^{(k)c}\subset N_{i(r_k+s_b)}(X_k)$ for $i=0,\dots,7$, and $X_k\subset\Lambda_5^{(k)c}\subset\Lambda_7^{(k)c}$;
(c) put $I_k:=\{z:\dist_\infty(z,T^{(k)}\setminus\Lambda_7^{(k-1)\prime})\ge s_b\}$ ($I_0:=T^{(0)}$). If $z\in I_k\cap X_k$, the large block $\beta\ni z$ lies in $\Lambda_7^{(k-1)\prime}$ and $\dist(\beta,\mathcal S_k)\le r_k$;
(d) $T^{(k)}\setminus\Lambda_7^{(k-1)\prime}=(\Lambda_7^{(k-1)c})'$, hence $T^{(k)}\setminus I_k=N_{s_b-1}\big((\Lambda_7^{(k-1)c})'\big)$.
\end{lemma}
\begin{proof}
(a) Let $y$ lie in an element of $F_k$ and let $\beta\ni y$ be its large block. Then $\dist(\beta,\mathcal S_k)=0\le r_k$. If $\beta\subset\Lambda_7^{(k-1)\prime}$, (Dk) excludes $\beta$ from $\Lambda_0^{(k)}$. Otherwise $\beta$ is not a candidate block. Large blocks partition $T^{(k)}$, so $\beta\cap\Lambda_0^{(k)}=\emptyset$ in both cases.
(b) Each step of (2.8) moves at most $r_k$, plus at most $s_b-1$ inside a block.
(c) $\beta$ has $\ell^\infty$-diameter $s_b-1$, so $\beta\subset\Lambda_7^{(k-1)\prime}$. Since $\beta\not\subset\Lambda_0^{(k)}$, (Dk) gives $\dist(\beta,\mathcal S_k)\le r_k$.
(d) $\Lambda_7^{(k-1)}$ is a union of $L$-blocks, and the prime operation commutes with complements on such unions.
\end{proof}

\begin{lemma}[number of families]\label{app:G1:I6-N}
For admissible $\Lambda_0^{(k)}$, the number of six-tuples $F_k$ admissible for it is at most $2^{n_d\abs{X_k}}$, where $n_d:=4+2\max\{d,d(d-1)/2\}$ ($n_3=10$).
\end{lemma}
\begin{proof}
By Lemma~\ref{app:G1:I6-E}(a) every element has all its points in $X_k$. Charge each element to one point: an element of $P$-type to a point of its $L$-block, a bond or plaquette to its lowest corner. A point of $X_k$ is charged by at most one element of each of the types $P_v,R_v,P_s,R_s$, and by at most $\max\{d,d(d-1)/2\}$ elements of each of the types $Q_v,Q_s$. So a six-tuple is a subset of at most $n_d\abs{X_k}$ candidates.
\end{proof}

\begin{lemma}[cube count on the interior; repair of (3.45)]\label{app:G1:I6-C}
Partition each coordinate cycle of $T^{(k)}$ (a cycle of $n_k/s_b$ large blocks, $n_k:=\ell_T/t_k$) into consecutive arcs of $q_k,\dots,2q_k-1$ large blocks, where $q_k:=\lceil(r_k+s_b)/s_b\rceil$; if there are fewer than $q_k$ blocks, use one arc. Let $\mathcal Q_k$ be the product partition. Its boxes are unions of large blocks with sides in $[\min\{r_k+s_b,n_k\},\,2r_k+3s_b]$, and $2r_k+3s_b\le(2+3\vartheta)r_k$. Put $\mathcal C_k:=\{Q\in\mathcal Q_k:Q\cap X_k\cap I_k\ne\emptyset\}$. Then $\abs{\mathcal C_k}\le3^d\abs{F_k}$ for every family $F_k$ admissible for $\Lambda_0^{(k)}$.
\end{lemma}
\begin{proof}
(i) If $\dist_\infty(z,x)\le r_k+s_b$, then in each coordinate the cyclic distance is at most the minimal arc length (or there are at most two arcs). So the boxes $Q(z)$ and $Q(x)$ are equal or adjacent, and a box has at most $3^d$ equal or adjacent boxes.
(ii) Let $Q\in\mathcal C_k$, $z\in Q\cap X_k\cap I_k$, and $\beta\ni z$ its large block ($\beta\subset Q$). By Lemma~\ref{app:G1:I6-E}(c) there are $z'\in\beta$ and a point $y$ of an element $\mathfrak e$ of $F_k$ with $\dist_\infty(z',y)\le r_k$. The charge point $x_{\mathfrak e}$ is within $\max\{1,L-1\}$ of $y$. Hence $\dist_\infty(z',x_{\mathfrak e})\le r_k+L-1\le r_k+s_b$, and $Q=Q(z')$ is equal or adjacent to $Q(x_{\mathfrak e})$. Therefore $\mathcal C_k$ is contained in the union over $\mathfrak e\in F_k$ of the boxes equal or adjacent to $Q(x_{\mathfrak e})$.
\end{proof}

\begin{remark}[the printed (3.45)]\label{app:G1:I6-345}
\cite[p.~592]{Ba2} defines $\mathcal C_k$ through $\Lambda_7^{(k-1)\prime}\cap\Lambda_0^{(k)c}$ instead of the interior $I_k$. Now $\Lambda_7^{(k-1)\prime}$ is a union of $M_b$-cubes of $T^{(k)}$, not of scale-$k$ large blocks. So a large block straddling its boundary lies in $\Lambda_0^{(k)c}$ and can be far from $\mathcal S_k$. For $F_k=\emptyset$, the term $1$ of every $\zeta''$, the printed $\mathcal C_k$ is non-empty while $3^d\abs{F_k}=0$. The repair (the count on the interior, with the boundary layer covered by the inherited boxes, Lemma~\ref{app:G1:I6-D}) changes constants only.
\end{remark}

\begin{lemma}[bound on $\zeta''$]\label{app:G1:I6-Z}
For every admissible $\Lambda_0^{(k)}$ with $k\le K-1$: $\zeta''_{\Lambda_0^{(k)}}\le2^{n_d\abs{X_k}}\exp(-c_\square p(t_k)^2\abs{\mathcal C_k})$.
\end{lemma}
\begin{proof}
Each of the six coefficients $\tfrac18a,\tfrac14\gamma_0,\tfrac14\gamma_0,\tfrac18a,\tfrac14\gamma_0,1$ of (3.41) is at least $3^dc_\square$. So each term of (3.41) is at most $\exp(-3^dc_\square p(t_k)^2\abs{F_k})$, which is at most $\exp(-c_\square p(t_k)^2\abs{\mathcal C_k})$ by Lemma~\ref{app:G1:I6-C}. The number of terms is bounded by Lemma~\ref{app:G1:I6-N}.
\end{proof}

\begin{lemma}[weighted region entropy]\label{app:G1:I6-W}
For $k\le K-1$ and $\theta\ge0$: $\sum_X\mathrm{e}^{-\theta\abs X}=(1+\mathrm{e}^{-\theta s_b^d})^{N_k/s_b^d}\le\exp(N_k\mathrm{e}^{-\theta})$, where the sum runs over all unions $X$ of large blocks of $T^{(k)}$. With $\theta_k:=(d+1)\hat\ell_k$,
$\sum_{k<K}N_k\mathrm{e}^{-\theta_k}=\mathrm{e}^{-(d+1)}\abs T\sum_{k<K}t_k\le \mathrm{e}^{-(d+1)}\eps_0\abs T=\mathrm{e}^{-(d+1)}\eps_0^{d+1}\abs{T^{(K)}}$.
\end{lemma}
\begin{proof}
The large blocks partition $T^{(k)}$ into $N_k/s_b^d$ blocks (by (B-C)), and $X$ is any subset of them; use the product formula, then $1+x\le \mathrm{e}^x$ and $s_b\ge1$. Also $\mathrm{e}^{-\theta_k}=\mathrm{e}^{-(d+1)}t_k^{d+1}$, and $\sum_{k<K}t_k=\eps_0\sum_{m\ge1}2^{-m}$.
\end{proof}
\noindent A price $\theta$ independent of $k$ would give $\sum_kN_k\mathrm{e}^{-\theta}\asymp \mathrm{e}^{-\theta}\abs T\eps^{-d}\to\infty$. The logarithm in the second exponent of (3.47) is this price.

\begin{proposition}[from the sum to the supremum]\label{app:G1:regionsup}
Let $F(\mathbf\Lambda)$ be the summand of (3.42), and $C_{\rm fin}:=c_7\eps_0^{\kappa_0}+c_5+\mathrm{e}^{-(d+1)}\eps_0^{d+1}$. Then
\[
\sum_{\mathbf\Lambda\ \rm adm}F(\mathbf\Lambda)\le \mathrm{e}^{C_{\rm fin}\abs{T^{(K)}}}\sup_{\mathbf\Lambda\ \rm adm}\exp\Big(-\sum_{k<K}c_\square p(t_k)^2\abs{\mathcal C_k}+\Xi(\mathbf\Lambda)\Big),
\]
where $\Xi(\mathbf\Lambda):=\sum_{k<K}\big[(n_d\ln2+\theta_k)\abs{X_k}+c_7t_k^{\kappa_0}\abs{\mathrm{Col}_k}+c_5\abs{\mathrm{Sh}_k}\big]$.
\end{proposition}
\begin{proof}
Lemma~\ref{app:G1:I6-Z} bounds each $\zeta''$ by a function of the region alone. The terms with $k=K$ are at most $(c_7\eps_0^{\kappa_0}+c_5)\abs{T^{(K)}}$, since $\mathrm{Col}_K,\mathrm{Sh}_K\subset T^{(K)}$. Multiply and divide by $\mathrm{e}^{\sum_{k<K}\theta_k\abs{X_k}}$. Bound the resulting function $\Psi(\mathbf\Lambda)$ by its supremum and sum $\mathrm{e}^{-\sum_k\theta_k\abs{X_k}}$ over all sequences $(X_0,\dots,X_{K-1})$ of unions of large blocks: the map $\mathbf\Lambda\mapsto(X_k)_k$ is injective and all terms are non-negative. Finally apply Lemma~\ref{app:G1:I6-W}.
\end{proof}

\begin{lemma}[covering families; (3.48)--(3.52) repaired]\label{app:G1:I6-D}
Put $\nu:=14+20\vartheta$, $\nu':=16+23\vartheta$, $S_{\varrho,L}:=\sum_{m\ge0}L^{-m}(1+m\ln L)^{\varrho}$ ($S_{2,2}=2+4\ln2+6\ln^22=7.6553$), $C_D^0:=\nu'(S_{\varrho,L}-1)$ and $C_D:=C_D^0+14(1+\vartheta)$ ($C_D=144.35$ at $\vartheta=1/7$). Set $\mathcal D_0:=\mathcal C_0$. For $k\ge1$ set $\mathcal D_k:=\{\mathcal P_k(Q):Q\in\mathcal D_{k-1}\}\cup\mathcal C_k$, where $\mathcal P_k(Q)$ is the prime of the union of the $L$-blocks meeting $N_{w_{k-1}}(Q)$, with $w_{k-1}:=7(r_{k-1}+s_b)+Ls_b$. Then for $k\le K-1$:
(i) $X_k\subset\bigcup\mathcal D_k$;
(ii) $\abs{\mathcal D_k}\le\sum_{j\le k}\abs{\mathcal C_j}$;
(iii) every box of $\mathcal D_k$ has side at most $C_D^0r_k$;
(iv) $\abs{\Lambda_i^{(k)c}}\le(C_Dr_k)^d\sum_{j\le k}\abs{\mathcal C_j}$ for $i=0,\dots,7$.
\end{lemma}
\begin{proof}
(ii) is immediate.
(iii) A box born at scale $j$ has side $s_j\le(2+3\vartheta)r_j$ (Lemma~\ref{app:G1:I6-C}). One step gives $s_{i+1}\le(s_i+2w_i+2(L-1))/L$, and $2w_i+2L-2\le14r_i+20s_b\le\nu r_i$. Iterating, $s_k\le\nu'\sum_{i=j}^{k-1}L^{-(k-i)}r_i$. Since $\hat\ell_i=\hat\ell_k+(k-i)\ln L\le(1+(k-i)\ln L)\hat\ell_k$, we have $r_i\le(1+(k-i)\ln L)^{\varrho}r_k$, and summing gives $s_k\le C_D^0r_k$.
(i) We induct on $k$. The case $k=0$ holds because $I_0=T^{(0)}$, and $X_k\cap I_k\subset\bigcup\mathcal C_k$ by definition. Let $z\in X_k\setminus I_k$. By Lemma~\ref{app:G1:I6-E}(d) there is $z_1\in(\Lambda_7^{(k-1)c})'$ with $\dist_\infty(z,z_1)\le s_b-1$. Then $B(z_1)\subset\Lambda_7^{(k-1)c}\subset N_{7(r_{k-1}+s_b)}(X_{k-1})$ by Lemma~\ref{app:G1:I6-E}(b), and this is contained in $\bigcup_{Q\in\mathcal D_{k-1}}N_{7(r_{k-1}+s_b)}(Q)$ by induction. In scale-$(k-1)$ units the points of $B(z)$ are within $Ls_b-1$ of points of $B(z_1)$. So $B(z)$ meets some $N_{w_{k-1}}(Q)$, i.e.\ $z\in\mathcal P_k(Q)$.
(iv) By Lemma~\ref{app:G1:I6-E}(b) and (i), $\Lambda_i^{(k)c}\subset\bigcup_{Q\in\mathcal D_k}N_{7(r_k+s_b)}(Q)$. Each of these sets has at most $(C_D^0r_k+14(r_k+s_b))^d\le(C_Dr_k)^d$ points; now use (ii).
\end{proof}

\begin{lemma}[fatness of the large-field complement]\label{app:G1:I6-fat}
Assume $\vartheta\le1/7$ and $R\ge12\sqrt d$ (lower bounds on $R$ of stage 5 of Appendix~\ref{app:B}). Then for every admissible sequence and $k\le K-1$: $\abs{\Lambda_7^{(k)c}}\le C_F\abs{\Lambda_0^{(k)c}}$, where $C_F:=(4\sqrt d(8.434+9\vartheta)+2)^d$ ($=3.334\cdot10^5$ at $d=3$, $\vartheta=1/7$).
\end{lemma}
\begin{proof}
Let $h_k:=\lfloor r_k/(2\sqrt d)\rfloor-2$. Then $h_k\ge r_k/(4\sqrt d)\ge1$, because $r_k\ge R\ge12\sqrt d$.
(a) \emph{Centres.} For each $z\in X_k$ there is $y$ with $B_\infty(y,h_k)\subset X_k$ and $\dist_\infty(z,y)\le\omega_k:=(1.434+2\vartheta)r_k$.
If $z\in I_k$: as in Lemma~\ref{app:G1:I6-C}(ii), $\dist_\infty(z,x_{\mathfrak e})\le r_k+2s_b\le(1+2\vartheta)r_k$. Every $w$ with $\dist_\infty(w,x_{\mathfrak e})\le\lfloor r_k/\sqrt d\rfloor$ is at distance at most $r_k$ from $\mathcal S_k$, so its large block lies in $X_k$ by the argument of Lemma~\ref{app:G1:I6-E}(a). Take $y:=x_{\mathfrak e}$.
If $z\notin I_k$: by Lemma~\ref{app:G1:I6-E}(d) take $z_1$ as in the proof of Lemma~\ref{app:G1:I6-D}(i). The scale-$(k-1)$ large block $\beta_1\supset B(z_1)$ lies in $\Lambda_7^{(k-1)c}$. By (2.8) there are $z_2\in\beta_1$ and $v\in\Lambda_6^{(k-1)c}$ with $\dist(z_2,v)\le r_{k-1}$, and then $B_\infty(v,\lfloor r_{k-1}/\sqrt d\rfloor)\subset\Lambda_7^{(k-1)c}$. Let $u_v$ be the coarse point whose $L$-block contains $v$. Then $B_\infty(u_v,h'')\subset(\Lambda_7^{(k-1)c})'\subset X_k$ with $h''\ge r_{k-1}/(2\sqrt d)-2\ge h_k$. Moreover $\dist_\infty(z,u_v)\le r_{k-1}/2+1.5s_b\le((1+\ln2)^2/2+1.5\vartheta)r_k\le\omega_k$. Take $y:=u_v$.
(b) \emph{Reach.} By Lemma~\ref{app:G1:I6-E}(b), $\Lambda_7^{(k)c}$ lies within $\tau_k:=7(r_k+s_b)+\omega_k\le(8.434+9\vartheta)r_k$ of the set $Y_k$ of centres.
(c) \emph{Packing.} Let $Y'\subset Y_k$ be maximal with pairwise $\dist_\infty>2h_k$. The boxes $B_\infty(y',h_k)$, $y'\in Y'$, are disjoint and contained in $X_k$, and $\Lambda_7^{(k)c}$ lies within $\tau_k+2h_k$ of $Y'$. If $2h_k+1\ge n_k$, then $X_k=T^{(k)}$ and there is nothing to prove. Otherwise the ratio of the two cardinalities is at most $((2\tau_k+4h_k+1)/(2h_k+1))^d\le(\tau_k/h_k+2)^d\le C_F$.
\end{proof}

\begin{proposition}[the estimate (3.42) on the torus]\label{app:G1:I6}
Let $\eps_0\le \mathrm{e}^{-1}$ and $2p>\varrho d+2$ (at $p=6$, $\varrho=2$, $d=3$ the margin is $\hat\ell^4$). Put $\Theta:=C_1C_D^dR^d/(c_\square b_0^2\ln2)$ with $C_1:=n_d\ln2+d+1+c_7+c_5$, and assume
\begin{equation}\label{p:eq-ent}
\text{(E.ent)}\qquad (1+\ln(2/\eps_0))^{2p-2-\varrho d}\ge\Theta,\qquad\text{equivalently}\quad \eps_0\le2\exp\big(1-\Theta^{1/(2p-2-\varrho d)}\big).
\end{equation}
Then the left side of \cite[(3.42)]{Ba2} is at most $\exp(C_{\rm fin}\abs{T^{(K)}})=\exp(C_{\rm fin}\eps_0^{-d}\abs T)$. This holds uniformly in $K$, in $k$ and in the torus $T$, for every torus whose side $\ell_T$ satisfies $\ell_T/(M_bt_K)\in\N$ (Lemma~\ref{lem:sizes}(ii)). If in addition $\vartheta\le1/7$ and $R\ge12\sqrt d$, the printed bound \cite[(3.47)]{Ba2} holds with second $O(1)$ equal to $C_2:=n_d\ln2+(d+1)+(c_7+c_5)C_F$ and third $O(1)$ equal to $C_{\rm fin}$.
\end{proposition}
\begin{proof}
By Proposition~\ref{app:G1:regionsup} it suffices to show $\Xi(\mathbf\Lambda)\le\sum_{j<K}c_\square p(t_j)^2\abs{\mathcal C_j}$. The sets $X_k$, $\mathrm{Col}_k$ and $\mathrm{Sh}_k$ lie in $\Lambda_7^{(k)c}$ (Lemma~\ref{app:G1:I6-E}(b)), and $t_k^{\kappa_0}\le1\le\hat\ell_k$. So each coefficient in $\Xi$ is at most $C_1\hat\ell_k$. By Lemma~\ref{app:G1:I6-D}(iv), after swapping the sums,
$\Xi\le C_1C_D^d\sum_{j<K}\abs{\mathcal C_j}\sum_{k=j}^{K-1}\hat\ell_kr_k^d$. Here $\hat\ell_kr_k^d=R^d\hat\ell_k^{1+\varrho d}$ decreases in $k$, and the number of terms is $K-j=\log_L(\eps_0/t_j)\le\log_2(\eps_0/t_j)\le\hat\ell_j/\ln2$. Hence $\Xi\le c_\square b_0^2\Theta\sum_j\hat\ell_j^{2+\varrho d}\abs{\mathcal C_j}$. Since $p(t_j)^2=b_0^2\hat\ell_j^{2p}$ and $\hat\ell_j\ge1+\ln(2/\eps_0)$ for $j\le K-1$, (E.ent) gives the claim. For the printed form: by Lemma~\ref{app:G1:I6-fat}, $\abs{\mathrm{Col}_k},\abs{\mathrm{Sh}_k}\le C_F\abs{X_k}$, and $n_d\ln2+(d+1)\hat\ell_k+(c_7+c_5)C_F\le C_2\hat\ell_k$.
\end{proof}
\noindent The bound of Proposition~\ref{app:G1:I6} holds for arbitrary constants $c_5,c_7\ge0$. Let $\eps_1$ be the cap of stage 10a of Appendix~\ref{app:B}. If $c_5$ and $c_7$ are bounded for $\eps_0\in(0,\eps_1]$ (for $c_7$ this uses $\kappa_0<\frac14$, which $\kappa_0=\frac15$ meets), then (E.ent), with $\Theta$ evaluated at these bounds, is a downward-closed upper bound on $\eps_0$. The constants $c_5,c_7$ are bounded on $(0,\eps_1]$ by Statement~\ref{app:K:SL}. Hence (E.ent) is a condition of stage 10b of Appendix~\ref{app:B}.

\subsection{Torus and infinite-lattice propagators}\label{ssec:T6-IMG}
At step $k\le K$ the torus $T_\eps$, rescaled, has sides $S_k^{(\mu)}=2L_\mu/t_k$ in units of scale-$k$ blocks, with deck group $\Gamma_k:=\bigoplus_\mu S_k^{(\mu)}\Z e_\mu$; put $S_k:=\min_\mu S_k^{(\mu)}=2L_{\min}/t_k$, so $S_k=L^{K-k}S_K$ and $S_K=2L_{\min}/\eps_0$. Let $\mathcal K$ be the set of kernels $G_k(0)=(-\Delta^{\eta_k}+m^2t_k^2+a_kP_k)^{-1}$, its slices $G^{\eta_k}_{(j)}(0)$, $0\le j\le k$, and $C^\xi=(-\Delta^\xi+1)^{-1}$, with their derivatives of order $\le2$, each defined on $X\in\{T_{\eta_k},\eta_k\Z^d\}$ (written $K_T$, $K_\infty$). By Proposition~\ref{prop:UV-BaProp21} and \cite[(2.10)--(2.12), p.~426, and p.~437]{Ba3}, with constants depending on $d,a,M_b,L$ only:
\[
\text{(H}_\infty\text{)}\qquad \abs{K_\infty(x,x')}+\norm{K_\infty}_{1,\alpha;x,x'}\le C_0\mathrm{e}^{-\delta_0\abs{x-x'}}\ (\abs{x-x'}\ge1),\qquad \sum_{x'}\eta_k^d\abs{K_\infty(x,x')}\le C_0 .
\]
\noindent In (H$_\infty$) and in the definition of $\mathcal K$, a derivative of a kernel means at most one lattice difference in each of the two variables. Mixed second differences $\partial_x\partial_{x'}K$ are therefore included, and pure second differences in one variable are not. The norm $\norm{K}_{1,\alpha;x,x'}$ consists of the pointwise values of these kernels together with the H\"older quotients of their first differences. The summability bound $\sum_{x'}\eta_k^d\abs{K_\infty(x,x')}\le C_0$ is required only for kernels with at most one difference in total. For a mixed second difference of $G_i$ the sum is bounded only by a constant times $i$ (Proposition~\ref{app:K:L19-1}(iii)), so this bound is not uniform in $k$. On the tori $T_{\eta_k}$ the kernels $K_T$, $K\in\mathcal K$, satisfy the pointwise and H\"older bounds of (H$_\infty$) with constants $C_0,\delta_0$ that do not depend on the torus, by Proposition~\ref{prop:UV-BaProp21}, \cite[(2.10)--(2.12), p.~426]{Ba3} and Lemma~\ref{app:G1:torusuniform}. The constants obtained in this way are those of the kernels at the photon mass $\mu_*$. For the scalar kernels $G_k(0)$ and their slices at a Higgs mass $M\in(0,m_2]$ with $M\eps_0\le1$, (H$_\infty$) in the above form is Proposition~\ref{app:K:L19-1}(ii), with constants depending only on $a$, $L$ and $\alpha$. By (E.m2) this covers $M\in\{m_1,m_2\}$, and Lemma~\ref{lem:IMG} is applied at these masses only through that statement.

\begin{lemma}[periodisation]\label{lem:IMG}
Assume (H$_\infty$).
\begin{enumerate}
\item[(i)] $K_T(x,x')=\sum_{n\in\Gamma_k}K_\infty(\tilde x,\tilde x'+n)$ for every $K\in\mathcal K$ and all lifts $\tilde x,\tilde x'$.
\item[(ii)] For $F_T=\operatorname{Per}F_\infty$ as in (i), every $\Gamma_k$-periodic kernel $V_T$ and every weight $g$,
\[
\sum_{x'\in T_{\eta_k}}\eta_k^dF_T(x,x')V_T(x,x')g(x')=\sum_{y\in\eta_k\Z^d}\eta_k^dF_\infty(\tilde x,y)V_T(x,\pi y)g(\pi y),
\]
with $\pi$ the projection to the torus.
\item[(iii)] If $\tilde x,\tilde x'$ lie in one lifted cube of diameter $D$ and $S_k\ge\max\{2D,2/\delta_0\}$, then $\norm{K_T-K_\infty}_{1,\alpha}\le C_1\mathrm{e}^{-\delta_0S_k/2}$ on that cube, with $C_1=c_dC_0$.
\item[(iv)] Assume the condition (E.r), $\eps_0<\mathrm{e}^{1-\varrho}$ (stage 10b of Appendix~\ref{app:B}), and the condition
\begin{equation}\label{eq:GIMG}
\text{(T.per)}\qquad 2L_{\min}\ge\eps_0\max\{6\sqrt d\,r(\eps_0)+2,\ 2/\delta_0\}.
\end{equation}
Then $\mathrm{e}^{-\delta_0S_k/2}\le \mathrm{e}^{-3\delta_0r(t_k)}$ at every $k\le K$, and $\sum_{k\le K}\mathrm{e}^{-\delta_0S_k/2}\le2\mathrm{e}^{-\delta_0S_K/2}$.
\item[(v)] Transform every counterterm graph of \cite[(1.26)--(1.31)]{Ba3} as the classes are transformed on \cite[pp.~433, 445]{Ba3}. Then the near scalar lines of each connected component lie in one cube $\square$ of side $3r(t_k)$, and $\square$ does not wrap around the torus. On $\square\times\square$ write $G_k(T_{\eta_k},0)=G_k(0)+\delta G_k(T_{\eta_k},\eta_k\Z^d,0)$, and the same for each slice. The term with only $G_k(0)$ is the counterterm in the form used in \cite[(3.7)--(3.9), p.~435, (3.19), p.~438, and pp.~444--445]{Ba3}. Every other term contains an external scalar field of $(1,\alpha)$-norm $\le C_1\mathrm{e}^{-3\delta_0r(t_k)}$ and is bounded by \cite[(3.3)]{Ba3}.
\end{enumerate}
\end{lemma}

\begin{proof}
(i) By the pointwise bound of (H$_\infty$) the series converges absolutely, and it is $\Gamma_k$-periodic in both variables. The operator $H=-\Delta^{\eta_k}+m^2t_k^2+a_kP_k$ is nearest-neighbour plus block averaging. Every block of scale $j\le k$ has side $L^{j-k}\le1$, and the sides $S_k^{(\mu)}$ are multiples of the block side, so $H$ commutes with $\Gamma_k$. Unfolding one sum gives $H_T\operatorname{Per}K_\infty=\operatorname{Per}(H_\infty K_\infty)=\operatorname{Per}\delta=\delta_T$. $H_T$ is positive on a finite-dimensional space, hence invertible, and $\operatorname{Per}K_\infty=H_T^{-1}=K_T$.

For $1\le i\le k+1$ let $G_i$ be $H^{-1}$ with $a_kP_k$ replaced by $a_iL^{2(k-i)}P_i$, where $P_i$ is the projection onto the fields constant on the blocks of side $L^{i-k}$. On each lattice the slices satisfy $G_{(0)}=G_1$ and $G_{(j)}=G_{j+1}-G_j$ (Lemma~\ref{app:K:L19-2}(f), that is, part (b) with $\mu^2$ replaced by $m^2t_k^2$ and $\R^3$ by $\R^2$, together with the identification of these differences with the zero-background slices $G^{\eta_k}_{(j)}(0)$; at zero background $Q_k(0)$ is the plain block average). Each $G_i$ is periodised by the argument above. For $i\le k$ this uses blocks of side $\le1$. The kernel $G_{k+1}$, which enters $G_{(k)}$, uses blocks of side $L$ and the divisibility $L\mid S_k^{(\mu)}$. This divisibility holds for $k<K$ since $S_k^{(\mu)}=LS_{k+1}^{(\mu)}$, and at $k=K$ on the tori $T_j$ by the count of big blocks in the paragraph after Lemma~\ref{app:G1:torusuniform}. Derivatives commute with $\operatorname{Per}$.

(ii) Split the sum over lifts of $x'$; the remaining factors are periodic.

(iii) For $n\in\Gamma_k\setminus\{0\}$ we have $\abs n\ge S_k$ and $\abs{\tilde x-\tilde x'-n}\ge\abs n-D\ge\abs n/2\ge1$. So $\sum_{n\ne0}C_0\mathrm{e}^{-\delta_0\abs n/2}\le C_0\sum_{m\in\Z^d\setminus0}\mathrm{e}^{-\delta_0S_k\abs m/2}\le c_dC_0\mathrm{e}^{-\delta_0S_k/2}$ once $\delta_0S_k/2\ge1$. The same argument applies to the H\"older seminorms. Only the pointwise and H\"older bounds of (H$_\infty$) are used here.

(iv) $S_k/r(t_k)$ is decreasing in $t_k$ on $(0,\eps_0]$ for $\eps_0<\mathrm{e}^{1-\varrho}$ (condition (E.r); Lemma~\ref{lem:B:rows}(b)). So $S_K\ge6\sqrt d\,r(\eps_0)+2$ gives $S_k\ge6\sqrt d\,r(t_k)+2$ for all $k\le K$, hence $\mathrm{e}^{-\delta_0S_k/2}\le \mathrm{e}^{-3\sqrt d\,\delta_0r(t_k)}\le \mathrm{e}^{-3\delta_0r(t_k)}$. Moreover $\sum_k\mathrm{e}^{-\delta_0L^{K-k}S_K/2}\le \mathrm{e}^{-\delta_0S_K/2}\sum_{n\ge0}\mathrm{e}^{-(L^n-1)}\le2\mathrm{e}^{-\delta_0S_K/2}$, since $\delta_0S_K/2\ge1$.

(v) By S18 in the proof of Proposition~\ref{prop:prop1}, the near vertices of the class and of the counterterm lie in one cube $\square$ of side $3r(t_k)$, of diameter $D\le3\sqrt d\,r(t_k)$.

\emph{The cube does not wrap.} By (T.per) and (E.r), Lemma~\ref{lem:B:rows}(b) gives $S_k\ge\max\{6\sqrt d\,r(t_k)+2,\,2/\delta_0\}$ at every $k\le K$, so $3r(t_k)<S_k/2$. Choose a lift $\tilde\square\subset\eta_k\Z^d$ of $\square$ whose coordinate differences are at most $3r(t_k)$. Let $n\in\Gamma_k$ and let $\mu$ be a coordinate with $n_\mu\ne0$. Then $\abs{n_\mu}\ge S_k$, so $\abs{\tilde x_\mu-\tilde x'_\mu-n_\mu}>S_k/2>\abs{\tilde x_\mu-\tilde x'_\mu}$ for $\tilde x,\tilde x'\in\tilde\square$. Hence $\pi$ is a bijection from $\tilde\square$ onto $\square$, and $\abs{\tilde x-\tilde x'-n}\ge\max\{\abs{\tilde x-\tilde x'},S_k/2\}$ for $n\ne0$. So $\square$ does not wrap around the torus, and the torus distance on $\square$ is the distance on $\tilde\square$.

\emph{The expansion.} Since $S_k\ge2D$, (i), (iii) and (iv) give, for every $K\in\mathcal K$, $K_T=K_\infty+\delta K$ on $\square\times\square$ with $\norm{\delta K}_{1,\alpha}\le C_1\mathrm{e}^{-\delta_0S_k/2}\le C_1\mathrm{e}^{-3\delta_0r(t_k)}$. The derivatives in the norm are again kernels of $\mathcal K$, so no margin around $\square$ is needed. We argue by induction on the order. Lower-order counterterms inserted in the graph have smaller order \cite[p.~417]{Ba3} and are treated first.

Insert this splitting for every scalar line and expand the product; the expansion of the product of at most $\Nbar$ lines is finite. The vertices, the far lines (which became external fields in the same way on both sides) and the vector lines (the same periodic kernels in the class and in the counterterm) are not changed. So the term with $G_k(0)$ on every scalar line is the counterterm in the class form of \cite[(3.7)--(3.9), (3.19), pp.~444--445]{Ba3}, under the identification $\square\cong\tilde\square$.

Every other term has a nonempty set of lines carrying $\delta G_k$. These lines are replaced by the two-variable external fields $\delta G_k|_{\square\times\square}$. This is the operation that \cite[p.~433]{Ba3} performs with $\delta G_k(\square,\eta_k\Z^d,0)$, and the resulting decomposition into renormalised classes depends only on which lines are dropped \cite[p.~432]{Ba3}. Hence the term is of the type treated in \cite[(3.3)]{Ba3}, which is multilinear in the external-field norms.
\end{proof}
\noindent Here $L_{\min}$ is the half-side of Ba{\l}aban's torus $\{-L_\mu\le x_\mu<L_\mu\}$ \cite[(1.2), pp.~603--604]{Ba1}, so $2L_{\min}=\ell_j$ on $T_j$, and (T.per) is the lower bound $\ell_0\ge\eps_0\max\{6\sqrt3\,r(\eps_0)+2,\,2/\delta_0\}$ of stage 11 of Appendix~\ref{app:B}; it does not depend on $\eps$. By Lemma~\ref{lem:B:rows}(b), the passage in (iv) from $k=K$ to every $k\le K$ needs only the condition (E.r), $\eps_0<\mathrm{e}^{1-\varrho}$.

\section{From the vacuum bound to the uniform tube bound}\label{sec:chain}
This section proves Theorem~\ref{thm:ch-T9}, a bound on the generating functionals of the torus states that is uniform in $\eps$, in the torus and in the photon mass, from Theorem~\ref{thm:T6} and the results of Section~\ref{sec:UV}. The bound is a product of three factors. The factor for the curvature is the Gaussian domination of Lemma~\ref{lem:C-Fsector}. The factor for positive sources of the Wick square (Proposition~\ref{prop:ch-T7}) comes from a convexity bound for constant sources (Proposition~\ref{prop:ch-Pplus}), which uses Ba{\l}aban's vacuum bounds (Proposition~\ref{prop:ch-K0}), the c-number differences of Proposition~\ref{prop:ch-CN} and the identification of the counterterms (Proposition~\ref{prop:UV-B1e}). The factor for negative sources (Proposition~\ref{prop:ch-Tflat}) comes from the slice vacuum bound (Proposition~\ref{prop:ch-T5}), which follows from Theorem~\ref{thm:T6} at $(u,M)=(1,m_2)$ and $(-\eta,m_1)$. Every estimate is proved on one torus at a time; no comparison between tori of different sizes and no statement about the model on $\eps\Z^3$ is used.

\subsection{The scheme}\label{ssec:ch-scheme}
The states of Theorem~\ref{thm:main} are built in the torus counterterm scheme of Definition~\ref{def:counterterm}: on each $(\eps,T_j)$ the bare quadratic coefficient is $b_T:=b_{T_j}(\eps)=a_{K,T_j}(m)$, King's counterterm on the torus at Higgs mass $m$ and reference photon mass $\mu_*=m$, the same for every photon mass $\mu_0\in(0,\mu_*]$ and for the state $\flat$. The renormalised source is $\XR(h)=\wick{\abs\varphi^2}_{\rm R}(h)=\sum_x\eps^3h(x)(r_x^2-W_\eps-z_\eps)$ with $r_x:=\abs{\varphi(x)}$. Here $W_\eps=2\int_{[-\pi/\eps,\pi/\eps)^3}(\hat p^2+m^2)^{-1}\,d^3p/(2\pi)^3$ is the free value of $\mathbb E\abs{\varphi(0)}^2$ on the infinite lattice $\eps\Z^3$ (two real components, hence the factor $2$), which does not depend on the torus, and $z_\eps$ is the order-$e^2$ Wick constant of Subsection~\ref{ssec:observables}. The same $b_T$, $W_\eps$ and $z_\eps$ enter the states, the observables $\XR(h)$ and every bound of this section. The norms are $\norm h_{L^1,\eps}=\sum_x\eps^3\abs{h(x)}$, $\norm h_\infty$ and $\norm{\bg}_{2,\eps}$ of Subsection~\ref{ssec:observables}. On the periodic lattice, $-\Delta_\eps$ acting on each component of a bond function is the Hodge Laplacian $d^*d+dd^*$ (Lemma~\ref{app:G1:X}), and $C_\mu=(-\Delta_\eps+\mu^2)^{-1}$.

Two objects of this section are defined here. For $s\in(0,\infty)$ put
\[
C^\flat_s:=(-\Delta_\eps)^{-1}\Pi^\perp+s\abs T\,\Pi_0 ,
\]
with $(-\Delta_\eps)^{-1}$ the inverse on the range of $\Pi^\perp$: the photon covariance whose non-constant modes are massless and whose constant mode has variance $s$ in each coordinate $c_\mu$ (in the inner product of Subsection~\ref{ssec:action}, the constant bond function $c\mathbf 1_\mu$ has squared norm $\abs T\,c^2$). Then $C_{\mu_0}=(-\Delta_\eps+\mu_0^2)^{-1}\Pi^\perp+s(\mu_0)\abs T\,\Pi_0$ with $s(\mu_0)=(\mu_0^2\abs T)^{-1}$, and $C^\flat_s\ge C_{\mu_*}$ as forms if and only if $s\ge s(\mu_*)$, because $(-\Delta_\eps)^{-1}\ge(-\Delta_\eps+\mu_*^2)^{-1}$ on the range of $\Pi^\perp$. Let $\langle\cdot\rangle_{\eps,T,s}$ be the expectation of Definition~\ref{def:states} with $C_{\mu_0}$ replaced by $C^\flat_s$. For bounded $F\in\cA_\eps$, $\langle F\rangle_{\eps,T,s}\to\langle F\rangle_{\eps,T,\flat}$ as $s\to\infty$: step~(ii) of the proof of Lemma~\ref{lem:flat} uses only that the law of $A^\perp$ is fixed and the variance of the constant mode tends to infinity. Second, the proof of Proposition~\ref{prop:ch-Tflat} uses the slice-modified covariance
\[
C^\flat_{B,s}:=(C_{\mu_*}+\Gamma^B+\Gamma_\ge)\Pi^\perp+s\abs T\,\Pi_0\quad(0<s<\infty),\qquad C^\flat_B:=C^\flat_{B,\infty}\ \text{(Haar-distributed constant mode)},
\]
with $\Gamma^B$ of \eqref{eq:T6-GammaB} and the operator $\Gamma_\ge$ of Proposition~\ref{prop:ch-Tflat} (Definition~\ref{app:G1:Tflatdef}); it satisfies $C_{\mu_0}\le C^\flat_{B,s}$ for $\mu_0\in(0,\mu_*]$ and $s\ge s(\mu_0)$ (Lemma~\ref{app:G1:HBdom}). The class $\cD$, the set $E^\cD$ and the lattice norms are those of Subsection~\ref{ssec:observables}.

We write $\mathfrak C$ for the family of pairs $(\eps,T_j)$ with $\eps=2^{-K}\eps_0$, $K\ge K_0$, and $j\ge1$. Suprema over $\eps$ and $T$ below run over $\mathfrak C$ unless said otherwise.

\subsection{The counterterm gaps}\label{ssec:ch-window}
\begin{theorem}[counterterm gaps]\label{thm:ch-W}
Let $m>0$, $\mu_*:=m$, $M_1:=0.7m$, $H:=m^2$, and $(e,\lambda,m,\eta,m_1,m_2)\in\Wwin$. Then, uniformly in $\eps\le1/m$ and in every periodic torus $T$ of side $\ell_T\ge96/m$:
\begin{enumerate}
\item[(a)] $a_{K,T}(m)-a_{K,T}(M_1)\ge0.319\,m^2$;
\item[(b)] $g_{21}:=\inf_{\eps,T}[a_{K,T}(m_2)-a_{K,T}(m_1)]\ge(m_2-m_1)(m_2+m_1-0.3631m)>0$, the infimum running over the same $\eps$ and $T$;
\item[(c)] $a_{K,T}(m_1)\ge a_{K,T}(m)+2H$;
\item[(d)] $m<m_1<m_2$, $m_1\ge\sqrt2\mu_*$, $M_1<m$.
\end{enumerate}
\end{theorem}

\begin{proof}
\emph{Scaling.} Every c-number $a_{K,T}(M)$ is a finite sum of lattice graphs of orders $\lambda,e^2,\lambda^2,e^2\lambda,e^4$. Each graph is homogeneous under the change of units in which lengths ($x$, $\eps$, $\ell_T$) are multiplied by $m$, masses ($M$, $\mu$) are divided by $m$, and $\lambda\mapsto\lambda/m$, $e^2\mapsto e^2/m$, $a\mapsto a/m^2$ (in $d=3$, $\lambda$ and $e^2$ have the dimension of a mass). So it suffices to take $m=\mu_*=1$, $\eps\le1$, $\ell_T\ge96$.

(a) Proposition~\ref{app:G1:G2} gives, for $0<M_1<M_2$, $\eps\le1$ and every torus of side at least a given floor,
\[
a_{K,T}(M_2)-a_{K,T}(M_1)\ge(M_2^2-M_1^2)+16\lambda h_1-e^2\bar j-\lambda^2\bar s-e^2\lambda\bar B_3-e^4\bar B_4^-,
\]
with non-negative constants independent of the couplings. At $(M_1,M_2)=(0.7,1)$ and for tori of side $\ge96$, Corollary~\ref{app:G1:Wa} gives $\bar j\le0.0886293$, $\bar s\le7.179064$, $\bar B_3\le9.330427$, $\bar B_4^-\le1.703784$ and $h_1\ge0$. Drop the $\lambda$-term and use the monotonicity in $\lambda,e^2\ge0$. On $[0,0.1]^2$ the bound is then at least $0.51-0.00886293-0.07179064-0.09330427-0.01703784=0.31900432$.

(b) Proposition~\ref{app:G1:R} gives $\partial_Ma_{K,T}(M)\ge2M-p(M)$, where $p(M)$ is the majorant of that proposition (not the small-field parameter $p(t)$). It is a sum of the monomials $e^2,\lambda^2,e^2\lambda,e^4$ times non-negative coupling-free functions. Its supremum over $M\ge2$ is at most $e^2s_1+\lambda^2s_2+e^2\lambda s_3+e^4s_4$, with $s_1\le0.3178$, $s_2\le9.6941$, $s_3\le17.37755$, $s_4\le5.96785$ for tori of side $\ge64$ (Corollary~\ref{app:G1:Wb}). On the box $[0,0.1]^2$ this is at most $0.03178+0.096941+0.1737755+0.0596785=0.362175<0.3622\le0.3631$. Integrating from $m_1$ to $m_2$ gives (b).

(c) Corollary~\ref{app:G1:CorC} gives $a_{K,T}(m_1)-a_{K,T}(m)\ge(m_1-m)(m_1+m-S(m))$, where $S(m)\ge\sup_{M\ge m}p(M)$. On the box one may take $S(m)\le0.612m$ (Corollary~\ref{app:G1:Wc}). So (c) holds for $m_1\ge m_1^c:=\frac12[S(m)+((S(m)-2m)^2+8H)^{1/2}]$, and $m_1^c\le1.882m<2m<m_1$.

(d) holds by construction: $m<2m<m_1<m_2$, $m_1>2m>\sqrt2\mu_*$, and $M_1=0.7m<m$.

The ingredients of (a)--(c) are proved in Appendix~\ref{app:G1:part-D8}. The exact two-loop form of $a_{K,T}$ (Lemma~\ref{app:G1:poly}) rests on Lemma~\ref{app:G1:PD-low}, whose order $e^2\lambda$ uses the list of one-particle-irreducible graphs in Lemma~\ref{app:G1:PD-list21}, and on the order-$e^4$ coefficient of Corollary~\ref{app:G1:PD-e4}. Apart from the closed forms $s_1,s_2$ and those of $p$ in (c), the numerical constants are bounded by the cell upper sums~\eqref{app:G1:eq-cell} (proof of Corollary~\ref{app:G1:Wa}). For (b) and (c) these are combined with the bound for large masses of Lemma~\ref{app:K:L15-1}.
\end{proof}
\noindent That the conditions of Appendix~\ref{app:B} can be met at every point of $\Wwin$, with $\eps\le\eps_0\le1/m_2<1/m$ for every spacing of the construction and $\ell_0\ge96/m$, so that Theorem~\ref{thm:ch-W} applies to all spacings and tori of the construction, is Lemma~\ref{lem:window}.

\subsection{Ba{\l}aban's bounds at the masses used}\label{ssec:ch-anchors}
\begin{proposition}[Ba{\l}aban's vacuum bounds at the masses used]\label{prop:ch-K0}
Let $(e,\lambda,m,\eta,m_1,m_2)\in\Wwin$, $\mu_*:=m$, and let the regulator data satisfy the conditions of Appendix~\ref{app:B}. For the upper bound, assume Statement~\ref{app:K:SL}. For every $(\eps,T_j)\in\mathfrak C$, with $T=T_j$,
\[
Z_K(T;M)\ge \mathrm{e}^{-E_-(M)\abs T}\quad(M\in\{M_1,m,m_1,m_2\}),\qquad Z_K(T;M)\le \mathrm{e}^{E_+(M)\abs T}\quad(M\in\{M_1,m_1,m_2\}),
\]
with $E_\pm(M)$ independent of $\eps$ and $j$.
\end{proposition}
\begin{proof}
The hypotheses of Theorems~\ref{thm:UV-K0minus} and~\ref{thm:UV-K0plus} hold for $Z_K(T;M)$ by Lemma~\ref{rem:UV-K0hyp}, its action is Ba{\l}aban's by Proposition~\ref{prop:UV-B1e}, and every $T_{j,\eps}$ is a torus of the form \cite[(1.2)]{Ba1} (Subsection~\ref{ssec:UV-torus}). The lower bound is Theorem~\ref{thm:UV-K0minus}. The upper bound is Theorem~\ref{thm:UV-K0plus}; at the masses $M_1,m_1,m_2$ its printed proof is completed by Proposition~\ref{app:G1:I6} and Lemma~\ref{app:G1:RMG} (Remark~\ref{rem:UV-blockD}).
\end{proof}

\begin{proposition}[c-number differences]\label{prop:ch-CN}
Let $f_T(a):=\abs T^{-1}\log Z^F_T(a)+\frac12a(W_\eps+z_\eps)$, where $Z^F_T(a):=\int \mathrm{e}^{-S_{\mu_*,a}(A,\varphi)}\,dA\,d\varphi$ is the unnormalised integral (Lebesgue measure in $A$ and $\varphi$) with the action \eqref{eq:latticeaction} at photon mass $\mu_*$ and bare coefficient $a$, and put $\gamma_T(M):=f_T(a_{K,T}(M))-\abs T^{-1}\log Z_K(T;M)$. Then $\sup_{K,\,\ell_T\ge\ell_0}\abs{\gamma_T(M)-\gamma_T(M')}<\infty$ for $\{M,M'\}=\{M_1,m\}$ and for $\{M,M'\}=\{m_1,m_2\}$.
\end{proposition}
\begin{proof}
This is Proposition~\ref{app:G1:CN}; by Lemma~\ref{app:G1:vacnorm}(a) the $\gamma_T$ defined here is the one there. Its proof uses lattice power counting in $d=3$ (Proposition~\ref{app:G1:reisz3d}).
\end{proof}

\subsection{The factor for positive sources}\label{ssec:ch-plus}
For a bare coefficient $a$ put
\[
\alpha_T(\sigma;a):=\abs T^{-1}\log\langle\exp(\sigma \XR(1_T))\rangle_{\eps,T,\mu_*;a},
\]
the expectation in the model at photon mass $\mu_*$ with bare coefficient $a$.

\begin{proposition}[convexity bound for constant positive sources]\label{prop:ch-Pplus}
Let $(e,\lambda,m,\eta,m_1,m_2)\in\Wwin$, $\mu_*:=m$, $M_1:=0.7m$, and let the regulator data satisfy the conditions of Appendix~\ref{app:B}. Put
\[
\begin{gathered}
g:=\inf_{K,\,\ell_T\ge\ell_0}[a_{K,T}(m)-a_{K,T}(M_1)],\qquad \sigma_+:=g/2,\\
\bar\Gamma:=\sup_{K,\,\ell_T\ge\ell_0}[\gamma_T(M_1)-\gamma_T(m)],\qquad \kappa_+:=2[E_+(M_1)+E_-(m)+\bar\Gamma]_+/g .
\end{gathered}
\]
Then $g\ge0.319m^2$, $\bar\Gamma<\infty$, and for every $(\eps,T_j)\in\mathfrak C$ and $0\le\sigma\le\sigma_+$,
\[
\alpha_T(\sigma;b_T)\le\kappa_+\sigma,\qquad T=T_j .
\]
\end{proposition}

\begin{proof}
Every spacing of the family satisfies $\eps\le\eps_0\le1/m_2<1/m$ (Lemma~\ref{lem:T6-rows}(b)), and $\ell_0\ge96/m$ (Appendix~\ref{app:B}), so $g\ge0.319m^2$ by Theorem~\ref{thm:ch-W}(a); $\bar\Gamma<\infty$ by Proposition~\ref{prop:ch-CN}. Write $a_1:=a_{K,T}(M_1)$, $a_2:=b_T=a_{K,T}(m)$. The single-site weight of the model is $\mathrm{e}^{-\eps^3(\frac12ar_x^2+\lambda r_x^4)}$, so $\langle \mathrm{e}^{\sigma\sum_x\eps^3r_x^2}\rangle_a=Z^F_T(a-2\sigma)/Z^F_T(a)$, and therefore $\alpha_T(\sigma;a)=f_T(a-2\sigma)-f_T(a)$. The function $a\mapsto\log Z^F_T(a)$ is convex (its second derivative is the variance of $\frac12\sum_x\eps^3r_x^2$), hence so is $f_T$. For $0\le\sigma\le\sigma_+\le(a_2-a_1)/2$ write $a_2-2\sigma=(1-\theta)a_2+\theta a_1$ with $\theta=2\sigma/(a_2-a_1)\in[0,1]$; convexity gives
\begin{multline*}
\alpha_T(\sigma;a_2)\le\theta\,[f_T(a_1)-f_T(a_2)]=\frac{2\sigma}{a_2-a_1}\Big[\frac{\log Z_K(T;M_1)-\log Z_K(T;m)}{\abs T}+\gamma_T(M_1)-\gamma_T(m)\Big]\\
\le\frac{2\sigma}{g}\big[E_+(M_1)+E_-(m)+\bar\Gamma\big]_+ ,
\end{multline*}
by Proposition~\ref{prop:ch-K0} and $a_2-a_1\ge g$.
\end{proof}

\begin{proposition}[the factor for positive sources]\label{prop:ch-T7}
Under the hypotheses of Proposition~\ref{prop:ch-Pplus}, for every $(\eps,T_j)\in\mathfrak C$, every $\mu_0\in(0,\mu_*]\cup\{\flat\}$ and every lattice function $q$ with $0\le q\le\sigma_+$,
\[
\log\langle \mathrm{e}^{\XR(q)}\rangle_{\eps,T_j,\mu_0}\le\log\langle \mathrm{e}^{\XR(q)}\rangle_{\eps,T_j,\mu_*}\le\kappa_+\norm q_{L^1,\eps},
\]
all states at the coefficient $b_{T_j}$.
\end{proposition}

\begin{proof}
(1) For $q\ge0$, $\mathrm{e}^{\XR(q)}$ is a positive constant times $\mathrm{e}^{r^2(\eps^3q)}$, with $r^2(\eps^3q)=\sum_x\eps^3q(x)r_x^2$ in the notation of Subsection~\ref{ssec:UV-other}, item~1; its expectation does not increase when the photon covariance increases \cite[Cor.~6.3(1), p.~281]{BFS1} (on the torus graph, Lemma~\ref{lem:C-cov}). Since $b_T$, $W_\eps$, $z_\eps$ do not depend on $\mu_0$ and $C_{\mu_0}\ge C_{\mu_*}$ as forms (zero mode included), the first inequality holds for $\mu_0\in(0,\mu_*]$; this step uses only $q\ge0$. For $\flat$ the observable $F_q:=\mathrm{e}^{\XR(q)}$ is unbounded, so the limit $s\to\infty$ of Subsection~\ref{ssec:ch-scheme}, which concerns bounded observables, does not apply to it directly. Instead, $F_q$ is an elementary observable of Lemma~\ref{lem:C-flatlimit}: for $\delta>0$, the case just proved, applied to $(1+\delta)q\ge0$, gives $\sup_{0<\mu_0\le\mu_*}\langle\abs{F_q}^{1+\delta}\rangle_{\eps,T_j,\mu_0}\le\langle \mathrm{e}^{\XR((1+\delta)q)}\rangle_{\eps,T_j,\mu_*}<\infty$. Hence $\langle F_q\rangle_{\eps,T_j,\flat}=\lim_{\mu_0\downarrow0}\langle F_q\rangle_{\eps,T_j,\mu_0}\le\langle F_q\rangle_{\eps,T_j,\mu_*}$. (The proof of Lemma~\ref{lem:C-flatlimit} uses Lemma~\ref{lem:C-cov}, not this proposition.)

(2) The measure at photon mass $\mu_*$ on $T_j$ is reflection positive, on gauge-invariant observables, about the coordinate planes halfway between lattice planes (Proposition~\ref{prop:C-RP}(a), with the plane pairs of case~(ii) of Lemma~\ref{lem:C-RPP}), and $N$ is even, so the chessboard estimate (Subsection~\ref{ssec:UV-other}, item~2) applies to the site factors $F_x=\exp(\eps^3q(x)(r_x^2-W_\eps-z_\eps))$. The product of the reflected copies of $F_x$ over all sites is $\exp(q(x)\XR(1_T))$ in the same model, so
\[
\langle \mathrm{e}^{\XR(q)}\rangle_{\eps,T_j,\mu_*}\le\prod_x\langle \mathrm{e}^{q(x)\XR(1_T)}\rangle_{\eps,T_j,\mu_*}^{\eps^3/\abs T}=\exp\Big(\sum_x\eps^3\alpha_T(q(x);b_T)\Big).
\]
(3) By Proposition~\ref{prop:ch-Pplus}, $\alpha_T(q(x);b_T)\le\kappa_+q(x)$. Sum over $x$.
\end{proof}

\subsection{The factor for negative sources}\label{ssec:ch-minus}
With $a_i:=a_{K,T}(m_i)$ put
\[
\Delta_1^B:=\sup_{(\eps,T_j)\in\mathfrak C,\ T=T_j}\Big\{\abs T^{-1}\log\frac{Z(C_1^B;a_2)}{Z(C_1^B;a_1)}+\tfrac12(a_2-a_1)(W_\eps+z_\eps)\Big\},\qquad C_1^B:=C_B(1)=C_{\mu_*}+\Gamma^B .
\]
We call the statement $\Delta_1^B<\infty$ the \emph{slice vacuum bound}.

\begin{proposition}[the slice vacuum bound]\label{prop:ch-T5}
Assume Theorem~\ref{thm:T6} at $(1,m_2)$ and at $(-\eta,m_1)$. Assume the lower bound $\abs T^{-1}\Phi_{m_1}(0)+e(\eps;0;m_1)\ge-E_-(m_1)$ on every torus of $\mathfrak C$; this is Theorem~\ref{thm:UV-K0minus} at $m_1$, by Lemma~\ref{app:G1:vacnorm}(b) and Proposition~\ref{prop:ch-K0}. Assume also the two c-number bounds
\begin{align}
&\sup_{\eps,T}\ \sup_{u\in[-\eta,1]}\abs{\partial_u^2e(\eps;u;M)}\le E_2(M),\qquad M\in\{m_1,m_2\},\label{p:ch-E2B}\\
&\bar\Gamma^{(1)}_{21}:=\sup_{\eps,T}\big[e(\eps;1;m_1)-e(\eps;1;m_2)+\tfrac12(a_2-a_1)(W_\eps+z_\eps)\big]<\infty .\label{p:ch-CNu}
\end{align}
Then
\[
\Delta_1^B\le C_+^B(1;m_2)+\eta^{-1}C_+^B(-\eta;m_1)+(1+\eta^{-1})E_-(m_1)+\tfrac12(1+\eta)E_2(m_1)+\bar\Gamma^{(1)}_{21}<\infty .
\]
\end{proposition}

\begin{proof}
Since $\Phi_{m_1}$ is convex and $-\eta<0<1$,
\[
\Phi_{m_1}(0)\le\frac1{1+\eta}\Phi_{m_1}(-\eta)+\frac\eta{1+\eta}\Phi_{m_1}(1),
\]
that is, $\Phi_{m_1}(1)\ge(1+\eta^{-1})\Phi_{m_1}(0)-\eta^{-1}\Phi_{m_1}(-\eta)$. Insert the upper bounds \eqref{eq:T6} for $\Phi_{m_2}(1)$ and $\Phi_{m_1}(-\eta)$, and the lower bound for $\Phi_{m_1}(0)$, into $\abs T^{-1}[\Phi_{m_2}(1)-\Phi_{m_1}(1)]+\frac12(a_2-a_1)(W_\eps+z_\eps)$. Write $e(u;M)$ for $e(\eps;u;M)$. The c-numbers regroup as
\[
\big[e(1;m_1)-e(1;m_2)+\tfrac12(a_2-a_1)(W_\eps+z_\eps)\big]-\big[e(1;m_1)-(1+\eta^{-1})e(0;m_1)+\eta^{-1}e(-\eta;m_1)\big].
\]
The first bracket is at most $\bar\Gamma^{(1)}_{21}$. The second is a second divided difference of the polynomial $u\mapsto e(\eps;u;m_1)$ on $[-\eta,1]$, so its absolute value is at most $\frac12(1+\eta)\sup_{[-\eta,1]}\abs{\partial_u^2e}$.
\end{proof}
\noindent Assume the hypotheses of Theorem~\ref{thm:T6}, together with the two further hypotheses of Lemma~\ref{app:G1:E2B}: $2\nbar+1\le3R$, and (H$_\infty$) for the kernels at the masses $M$ and $\mu_*$. Then the bound \eqref{p:ch-E2B} at $M\in\{m_1,m_2\}$ is Lemma~\ref{app:G1:E2B}. Its proof in Appendix~\ref{app:G1:part-D5} rests on Theorem~\ref{app:K:L17-1}, and the inputs of that step are described after the proof. The bound \eqref{p:ch-CNu} is Proposition~\ref{app:G1:CNuB}. Its one-loop part is Lemma~\ref{app:G1:CU1}, which uses King's slice at a point (Lemma~\ref{app:G1:WP-RBdiag}). That lemma in turn rests on the block Poincar\'e inequality and the variational form of $\Delta^{(j)}$ (Lemmas~\ref{app:G1:DOM-poinc}, \ref{app:G1:Dk-var} and~\ref{app:G1:Cj-bound}), with the constant $c_1$ of Lemma~\ref{app:G1:diagC}.

\begin{proposition}[the factor for negative sources]\label{prop:ch-Tflat}
Assume $\Delta_1^B<\infty$ and the anchor condition
\begin{equation}\label{p:ch-anchor}
b_T+2H\le a_1\quad\text{on every torus of }\mathfrak C,\qquad g_{21}>0
\end{equation}
(Theorem~\ref{thm:ch-W}(b), (c)). For every $0\le h\le H$, every $\mu_0\in(0,\mu_*]\cup\{\flat\}$ and every $(\eps,T_j)\in\mathfrak C$, with $T=T_j$,
\[
\log\langle \mathrm{e}^{-\XR(h)}\rangle_{\eps,T,\mu_0}\le2\norm h_{L^1,\eps}\,S_\flat^B,\qquad S^B_\flat:=[\Delta_1^B+J_{\rm IR}+\bar Z]_+/g_{21}.
\]
Here $\gamma_\ge(u):=\int_{\eps_0^2}^\infty \mathrm{e}^{-tu}(1-\mathrm{e}^{-t\mu_*^2})\,dt$ for $u>0$, so that the operator $\Gamma_\ge$ of Subsection~\ref{ssec:ch-scheme} is $\gamma_\ge(-\Delta_\eps)\Pi^\perp$, and $\hat k^2$ runs over the nonzero eigenvalues of $-\Delta_\eps$ on one component. With these,
\[
J_{\rm IR}:=\sup_{K,\,\ell_T\ge\ell_0}\tfrac32\abs T^{-1}\sum_{\hat k^2\ne0}\gamma_\ge(\hat k^2)\,(\hat k^2+\mu_*^2),\qquad Z_T:=\abs T^{-1}\Big[\tfrac32\log_+\frac{4\pi^2\mu_*^2\ell_T}{e^2}+3\delta_P\Big],\quad\delta_P:=2\sum_{n\ge1}\mathrm{e}^{-2\pi^2n^2},
\]
and $\bar Z:=\sup_{j\ge1}Z_{T_j}$. The constant $J_{\rm IR}$ is finite and depends on $\mu_*,\eps_0,\ell_0$ only, not on $\eps$, $e$, $\lambda$ or $\ell_T$. The constant $Z_T$ depends on the side $\ell_T$ (and on $\mu_*$, $e$) and tends to $0$ as $\ell_T\to\infty$. The constant $\bar Z$ depends on $\mu_*,e,\ell_0$ only, and
\begin{equation}\label{eq:ch-Zbar}
\bar Z\le\frac{6\pi^2\mu_*^2}{e^2\ell_0^2}+\frac{3\delta_P}{\ell_0^3}<\infty,
\end{equation}
since $\log_+x\le x$ for $x>0$ and $\ell_{T_j}=j!\,\ell_0\ge\ell_0$.
\end{proposition}

\begin{proof}
The full proof is Proposition~\ref{app:G1:Tflat}; here we describe its steps. Let $\rho^B(y;b)$ be the one-site density $\langle r_y^2-W_\eps-z_\eps\rangle$ in the model with photon covariance $C^\flat_B$ (the slice-modified massless photon with Haar-distributed constant mode) and constant coefficient $b$, and let $\bar\rho^B(b)$ be its average over a cell.

(1) We use three facts. The first is the translation-averaging lemma on the torus (Lemma~\ref{app:G1:TA}), which uses that the model of the states is invariant under $\eps\Z^3$ and the comparison model under $\eps_0\Z^3$. The second is the monotonicity of $\langle \mathrm{e}^{-\XR(h)}\rangle$ in the photon covariance \cite[Cor.~6.3(1), p.~281]{BFS1} (Lemma~\ref{app:G1:HBdom}). The third is its monotonicity in the coefficient: by \cite[Thm.~6.1, p.~276]{BFS1} with $P,Q$ polynomials in $r$, $\langle r_x^2;r_y^2\rangle\ge0$ in every model of the family $b\mapsto b+2th$, so $t\mapsto\langle r_y^2\rangle_t$ is non-increasing. Together they give $\log\langle \mathrm{e}^{-\XR(h)}\rangle_{\mu_0}\le-\norm h_{L^1,\eps}\,\bar\rho^B(b_T+2H)\le-\norm h_{L^1,\eps}\,\bar\rho^B(a_1)$, the second inequality by \eqref{p:ch-anchor}.

(2) $f^B(b):=\abs T^{-1}\log Z(C^\flat_B;b)+\frac12b(W_\eps+z_\eps)$ is convex with $f^{B\prime}=-\frac12\bar\rho^B$, so $-\bar\rho^B(a_1)\le2[f^B(a_2)-f^B(a_1)]/(a_2-a_1)$.

(3) Comparing $C^\flat_B$ with $C_1^B$ from above and below costs $J_{\rm IR}+Z_T$ per unit volume. This follows from the monotonicity of $Z(\Gamma;a)$ in $\Gamma$ (current representation) and a determinant comparison for the infrared part (Lemma~\ref{app:G1:IR}).

(4) Hence, on $T=T_j$, $f^B(a_2)-f^B(a_1)\le\Delta_1^B+J_{\rm IR}+Z_{T_j}\le\Delta_1^B+J_{\rm IR}+\bar Z$. Divide by $a_2-a_1\ge g_{21}$ and keep the positive part. By \eqref{eq:ch-Zbar} the constant does not depend on $\eps$, $j$ or $\mu_0$.

(5) The case $\mu_0=\flat$. For $h\ge0$ we have $\mathrm{e}^{-\XR(h)}\le \mathrm{e}^{(W_\eps+z_\eps)\norm h_{L^1,\eps}}$, since $r_x^2\ge0$, so this observable is bounded, and it lies in the algebra $\fF$ of Definition~\ref{def:C-algebra}. By Lemma~\ref{lem:C-flatlimit}, $\langle \mathrm{e}^{-\XR(h)}\rangle_{\eps,T,\flat}=\lim_{\mu_0\downarrow0}\langle \mathrm{e}^{-\XR(h)}\rangle_{\eps,T,\mu_0}$, and the bound, which is uniform in $\mu_0\in(0,\mu_*]$, passes to the limit.

\emph{The domination of the infrared part.} The comparison $C_{\mu_0}\le C^\flat_{B,s}$ of Lemma~\ref{app:G1:HBdom} is used in step (1), and the slice covariance $\Gamma^B$ enters $C^\flat_B$ in step (3). Both rest on Corollary~\ref{app:G1:DOM-torus}(a): $D^<\le\Gamma^B\Pi^\perp$ on every torus $T_j$ and for every $K\ge1$, where $\Gamma^B=\kappa\mu_*^2D$ and $\kappa=\frac{15}2$ is the fixed constant of \eqref{eq:T6-GammaB}. This needs only $\kappa\ge\bar\kappa$, with $\bar\kappa$ the $K$-independent constant of Theorem~\ref{app:G1:DOM-nat}. At $L=2$ and $a=1$, Theorem~\ref{app:G1:DOM-nat} gives $\kappa_1\kappa_3\le\bar\kappa\le\frac{42176}{5625}<\frac{15}2$ for every $K\ge1$ under (E.dom), $\mu_*\eps_0\le0.2$. The condition (E.$\delta$), $\kappa\mu_*^2\eps_0^2\le\frac1{12}$, gives $\mu_*\eps_0\le90^{-1/2}<0.2$ and so implies (E.dom). The block Poincar\'e inequality behind Theorem~\ref{app:G1:DOM-nat} holds on every torus of $\mathfrak C$. The reason is that for every $k\le K$ the blocks of scale $t_k$ contain no wrap-around bond, because $\ell_j\ge2t_k$ (Lemma~\ref{app:G1:DOM-geom}(a), (c)). The domination constant does not depend on $\eps$, $K$, $\ell_j$ or $\mu_0$.
\end{proof}

\subsection{The tube bound}\label{ssec:ch-T9}
For real lattice 2-forms $\bg_1,\dots,\bg_n$, real lattice functions $h_1,\dots,h_p$ on $T_\eps$ and $\zeta=(z,w)\in\C^{n+p}$ let
\[
G_{\eps,T,\mu_0}(\zeta):=\Big\langle\exp\Big[\sum_iz_iF(\bg_i)+\sum_kw_k\XR(h_k)\Big]\Big\rangle_{\eps,T,\mu_0},\qquad k(\zeta):=\sum_k(\rRe w_k)h_k=k_+-k_-,
\]
with $k_\pm\ge0$ the pointwise positive and negative parts, and, for $\sigma_+,H>0$, let $\mathcal D_{\sigma_+,H}:=\{\zeta:\norm{k_+(\zeta)}_\infty<\sigma_+/2,\ \norm{k_-(\zeta)}_\infty<H/2\}$ (no restriction on $z$). These are the objects of Subsection~\ref{ssec:observables} with lattice test functions in place of sampled continuum ones; for sampled test functions the continuum tube lies in the lattice tube (Lemma~\ref{lem:C-tubeincl}).

\begin{theorem}[the uniform tube bound]\label{thm:ch-T9}\label{thm:C-T9}
Let $m>0$, $0<\lambda<0.1m$ and $0<e^2<0.1m$, as in Theorem~\ref{thm:main}, put $\mu_*:=m$ and $(\eta,m_1,m_2):=(\frac12,3m,4m)$, so that $(e,\lambda,m,\eta,m_1,m_2)\in\Wwin$ (any point of $\Wwin$ with these $(e,\lambda,m)$ may be used instead), and let the regulator data be those of Lemma~\ref{lem:window} at this point. Put $\sigma_+:=g/2$ and $\kappa_+$ as in Proposition~\ref{prop:ch-Pplus}, $S^B_\flat$ as in Proposition~\ref{prop:ch-Tflat}, $\mathsf K:=\max(\kappa_+,2S^B_\flat)$ and $H:=m^2$. Then for every $K\ge K_0$, every $j\ge1$, every $\mu_0\in(0,\mu_*]\cup\{\flat\}$, all real lattice test functions and every $\zeta\in\mathcal D_{\sigma_+,H}$,
\begin{equation}\label{eq:C-T9}
\begin{aligned}
\abs{G_{\eps,T_j,\mu_0}(\zeta)}&\le G_{\eps,T_j,\mu_0}(\rRe\zeta)\le\exp\Big(\tfrac12\Big\|\sum_i(\rRe z_i)\bg_i\Big\|_{2,\eps}^2+\kappa_+\norm{k_+}_{L^1,\eps}+2S^B_\flat\norm{k_-}_{L^1,\eps}\Big)\\
&\le\exp\Big(\tfrac12\Big\|\sum_i(\rRe z_i)\bg_i\Big\|_{2,\eps}^2+\mathsf K\norm k_{L^1,\eps}\Big).
\end{aligned}
\end{equation}
The constants $\sigma_+$, $\kappa_+$, $S^B_\flat$ and $\mathsf K$ do not depend on $\eps$, $j$ or $\mu_0$.
\end{theorem}
\noindent The constant $S^B_\flat$ contains $\bar Z$, not $Z_T$, so it does not depend on the torus, and neither does $\mathsf K$ (Corollary~\ref{app:G1:T9const}). The constants $\sigma_+$ and $\mathsf K$ are the constants $\sigma_+$ and $C$ of Theorem~\ref{thm:main}(i) (Lemma~\ref{app:G1:window-apply}).

\begin{proof}
By Lemma~\ref{app:G1:window-apply} and Proposition~\ref{app:G1:VACT}, with these data the hypotheses of Propositions~\ref{prop:ch-T7} and~\ref{prop:ch-Tflat} hold on every torus of $\mathfrak C$, and Lemma~\ref{lem:C-Fsector} holds there at the coefficient $b_{T_j}$. (1) For real test functions $F(\bg_i)$ and $\XR(h_k)$ are real, so $\abs{G(\zeta)}\le G(\rRe\zeta)$. (2) Put $\bg:=\sum_i(\rRe z_i)\bg_i$. By Lemma~\ref{lem:C-Fsector}, which holds on every torus for every real source of $\wick{\abs\varphi^2}$ and every bare coefficient, $G(\rRe\zeta)\le \mathrm{e}^{\frac12\norm{\bg}^2_{2,\eps}}\langle \mathrm{e}^{\XR(k)}\rangle$. (3) $\XR(k)=\XR(k_+)-\XR(k_-)$, and by the Cauchy--Schwarz inequality $\langle \mathrm{e}^{\XR(k)}\rangle\le\langle \mathrm{e}^{\XR(2k_+)}\rangle^{1/2}\langle \mathrm{e}^{-\XR(2k_-)}\rangle^{1/2}$. (4) Since $0\le2k_+<\sigma_+$, Proposition~\ref{prop:ch-T7} gives $\langle \mathrm{e}^{\XR(2k_+)}\rangle^{1/2}\le \mathrm{e}^{\kappa_+\norm{k_+}_{L^1,\eps}}$. (5) Since $0\le2k_-<H$, Proposition~\ref{prop:ch-Tflat} gives $\langle \mathrm{e}^{-\XR(2k_-)}\rangle^{1/2}\le \mathrm{e}^{2S^B_\flat\norm{k_-}_{L^1,\eps}}$. (6) $k_+$ and $k_-$ have disjoint supports, so $\norm{k_+}_{L^1,\eps}+\norm{k_-}_{L^1,\eps}=\norm k_{L^1,\eps}$. Every step is on one torus at the coefficient $b_{T_j}$, and the constants do not depend on $\eps$, $j$ or $\mu_0$.
\end{proof}

\begin{remark}[what the proof uses]\label{rem:ch-summary}
The proof of Theorem~\ref{thm:ch-T9} uses Ba{\l}aban's vacuum bounds at four Higgs masses, his propagator bounds and his Proposition~1 (Section~\ref{sec:UV}, with the repairs of Section~\ref{sec:T6}), King's counterterms, the correlation inequalities and the current representation of \cite{BFS1}, lattice reflection positivity on the tori and the chessboard estimate, and the c-number and counterterm-gap bounds of Proposition~\ref{prop:ch-CN} and Theorem~\ref{thm:ch-W}. It does not use King's cube states, a cluster expansion, or any bound uniform in the photon mass in a fixed cube; the uniformity in $\mu_0$ comes from the monotonicity in the photon covariance (Propositions~\ref{prop:ch-T7} and~\ref{prop:ch-Tflat}).
\end{remark}

\section{Lattice statements on the tori}\label{sec:C-lattice}

Sections~\ref{sec:C-lattice}--\ref{sec:C-OS} prove parts (ii)--(iv) of Theorem~\ref{thm:main} from part (i), the uniform tube bound, which is Theorem~\ref{thm:ch-T9}. Everything else in these three sections is either a statement about one lattice model on one torus, or an abstract compactness, holomorphy or positivity argument about the resulting family of functionals. No cube, no boundary condition other than periodic, and no comparison between two different tori is used in Sections~\ref{sec:C-lattice}--\ref{sec:C-OS}.

\subsection{Standing assumptions}\label{ssec:C-standing}

\emph{Parameters and regulators.} Throughout Sections~\ref{sec:C-lattice}--\ref{sec:C-OS}, $m>0$, $0<\lambda<0.1\,m$ and $0<e^2<0.1\,m$; the reference photon mass is $\mu_*:=m$, and $H:=m^2$. The regulator data (blocking factor $L=2$, $\nbar=13$, order cap $\Nbar=\nbar^2=169$, $M_b$, $R$, $b_0$, $\eps_0=2^{-n_0}$, and $\ell_0\in2M_b\N$ with $\ell_0\ge183/m$) are those of Theorem~\ref{thm:ch-T9}, chosen as in Lemma~\ref{lem:window} at the point $(e,\lambda,m,\frac12,3m,4m)$ of the window $\Wwin$ of Section~\ref{sec:T6}. Sections~\ref{sec:C-lattice}--\ref{sec:C-OS} use the regulator data only through Theorem~\ref{thm:ch-T9} and through the properties (P1)--(P4) below.

\emph{Lattices and tori.} The lattice spacings are $\eps=\eps_K:=2^{-K}\eps_0$ with $K\ge K_0=\max(K_0',n_0+1)$ (Subsection~\ref{ssec:UV-torus}); they are nested: $\eps_{K'}\Z^3\supset\eps_K\Z^3$ for $K'\ge K$. The tori are $T_j=(\R/\ell_j\Z)^3$, $\ell_j=j!\,\ell_0$, $j\ge1$, with the lattice tori $T_{j,\eps}=T_j\cap\eps\Z^3$ and $N:=\ell_j/\eps$ sites per side. In Sections~\ref{sec:C-lattice}--\ref{sec:C-OS} we write $T$ for a torus $T_j$, $\ell_T:=\ell_j$, $\abs T:=\ell_T^3$, $T_\eps:=T_{j,\eps}$, and $\ev{\cdot}_{\eps,T,\mu_0}:=\ev{\cdot}_{\eps,j,\mu_0}$ for the states of Definition~\ref{def:states}, $\mu_0\in(0,\mu_*]\cup\{\fl\}$.

\emph{Fields, forms and the action.} We use the notation of Subsection~\ref{ssec:action}: $A$ is a real function on the bonds of $T_\eps$, $\varphi$ is a function from the sites to $\R^2\cong\C$, and $r_x:=\abs{\varphi_x}$. Sites, bonds and plaquettes carry the weight $\eps^3$ in $\ip\cdot\cdot_\eps$, $\norm f_{2,\eps}^2:=\ip ff_\eps$, and $d$, $d^*$ are the lattice exterior derivative and its adjoint. At photon mass $\mu_0>0$, the action \eqref{eq:latticeaction} with the coefficient $a=b_T:=b_{T_j}(\eps)=a_{K,T}(m)$ of the torus counterterm scheme (Definition~\ref{def:counterterm}) reads, in real coordinates,
\begin{equation}\label{eq:C-action}
S_{\eps,T,\mu_0}(A,\varphi)=\tfrac12\ip{A}{(-\Delta_\eps+\mu_0^2)A}_\eps+S_m(A,\varphi),\qquad
S_m=\tfrac12\sum_b\eps^3\abs{(D_A\varphi)(b)}^2+\sum_x\eps^3\big(\tfrac12 b_T r_x^2+\lambda r_x^4\big),
\end{equation}
with $(D_A\varphi)(b)=\eps^{-1}\big(U(\eps eA_b)\varphi(b_+)-\varphi(b_-)\big)$ for the bond $b$ from $b_-$ to $b_+$, where $U(s)$ is the rotation of $\R^2$ by the angle $s$; thus $S_m=S^{\mathrm m}_{b_T}$ in the notation of \eqref{eq:matteraction}. The photon part is King's Feynman form. The observables are the smeared curvature and the renormalised square of the Higgs field of Subsection~\ref{ssec:observables},
\begin{equation}\label{eq:C-obs}
F(\bg)=\ip{\bg}{dA}_\eps=\ip{d^*\bg}{A}_\eps,\qquad \XR(h)=\sum_x\eps^3h(x)\big(r_x^2-W_\eps-z_\eps\big),
\end{equation}
for real lattice 2-forms $\bg$ and real lattice functions $h$; $W_\eps$ is the free Wick constant and $z_\eps$ the order-$e^2$ c-number of Subsection~\ref{ssec:observables}.

The arguments of Sections~\ref{sec:C-lattice}--\ref{sec:C-OS} use the following properties of the regularised model, and nothing else about $b_T$, $W_\eps$, $z_\eps$:
\begin{itemize}
\item[(P1)] $b_T$, $W_\eps$, $z_\eps$ are real numbers that depend on $(\eps,T)$ only: they do not depend on the site, on the direction, or on $\mu_0$ (Lemma~\ref{app:G1:P1}(a)).
\item[(P2)] $\lambda>0$, so that $\varphi\mapsto e^{-\lambda\eps^3\abs\varphi^4+a\abs\varphi^2}$ is integrable on $\R^2$ for every $a\in\R$ (Lemma~\ref{app:G1:P2}).
\item[(P3)] The same $b_T$, $W_\eps$, $z_\eps$ are used in the state, in the observables \eqref{eq:C-obs}, and in Theorem~\ref{thm:ch-T9} (Lemma~\ref{app:G1:P3}).
\item[(P4)] $N=\ell_T/\eps$ is even, the lattices are nested, and $\ell_j\to\infty$ (Lemma~\ref{lem:sizes}).
\end{itemize}
Apart from Theorem~\ref{thm:ch-T9}, Sections~\ref{sec:C-lattice}--\ref{sec:C-OS} use the blocking factor only through (P4) and through the density in $\R^3$ of the group $\bigcup_K\eps_K\Z^3=\Z[1/2]^3\eps_0$ (proof of Theorem~\ref{thm:C-OS1}). Both have the same form for blocking factor $3$ (Lemma~\ref{lem:sizes}, and the density of $\Z[1/3]^3\eps_0$); this is how Sections~\ref{sec:C-lattice}--\ref{sec:C-OS} enter the proof of Theorem~\ref{thm:mainL3}.

\emph{States.} For a positive semidefinite operator $C$ on the bond functions, symmetric for $\ip\cdot\cdot_\eps$, we write $\gamma_C$ for the centred Gaussian measure with $\E_{\gamma_C}[\ip fA_\eps\ip{f'}A_\eps]=\ip f{Cf'}_\eps$. For $\mu_0\in(0,\mu_*]$ the state $\ev{\cdot}_{\eps,T,\mu_0}$ is the normalised expectation with weight $e^{-S_{\eps,T,\mu_0}}\,dA\,d\varphi$, Lebesgue measure on bonds and sites (Lemma~\ref{app:G1:gibbs}); it is a probability measure by (P2) and $\mu_0>0$. Its photon part is $\gamma_{C_{\mu_0}}$ with $C_{\mu_0}=(-\Delta_\eps+\mu_0^2)^{-1}$, so that
\begin{equation}\label{eq:C-state}
\ev{G}_{\eps,T,\mu_0}=\frac{\int d\gamma_{C_{\mu_0}}(A)\int d\varphi\,G\,e^{-S_m(A,\varphi)}}{\int d\gamma_{C_{\mu_0}}(A)\int d\varphi\,e^{-S_m(A,\varphi)}} .
\end{equation}
Write $A=A^\perp+c$, where $c=\Pi_0A=\sum_{\mu=1}^3c_\mu\mathbf 1_\mu$ is the projection onto the constant 1-forms ($\mathbf 1_\mu(b)=1$ for the bonds $b$ in direction $e_\mu$ and $0$ for the others) and $A^\perp=\Pi^\perp A$. Under $\gamma_{C_{\mu_0}}$ the components $A^\perp$ and $c$ are independent; $A^\perp$ has covariance $(-\Delta_\eps+\mu_0^2)^{-1}\Pi^\perp$, and each $c_\mu$ is centred normal with variance $s(\mu_0)=1/(\mu_0^2\abs T)$ (Lemma~\ref{app:G1:zeromode}). In the state $\fl$ (Definition~\ref{def:states}) $A^\perp$ has the law $\gamma_{(-\Delta_\eps)^{-1}\Pi^\perp}$, and $c$ is independent of $A^\perp$ and uniformly distributed on the period cell $[0,P)^3$, $P=2\pi/(\abs e\,\ell_T)$.

\begin{definition}[the observable algebra $\fF$]\label{def:C-algebra}
$\fF$ is the algebra of complex Borel functions of the variables $(r_x)_{x}$ and $((dA)_p)_p$; it is the class $\cA_\eps$ of Definition~\ref{def:algebra}. For a region $\Lambda$ of the torus, $\fF(\Lambda)$ is the subalgebra of the functions of the $r_x$ with $x\in\Lambda$ and of the $(dA)_p$ with all bonds of $p$ in $\Lambda$. Every polynomial and every exponential $\exp[\sum_iz_iF(\bg_i)+\sum_kw_k\XR(h_k)]$, $z,w\in\C$, lies in $\fF$.
\end{definition}
Every element of $\fF$ is gauge invariant: under $(A,\varphi)\mapsto(A+d\chi,\,U(-e\chi)\varphi)$, $\chi$ a real function on the sites, $r_x$ and $dA$ do not change. Elements of $\fF$ are also unchanged by $A\mapsto A+t\mathbf 1_\mu$ for every $t\in\R$, since $d\mathbf 1_\mu=0$.

\begin{lemma}[Hodge facts on the periodic lattice]\label{p:lem:C-hodge}
On $T_\eps$ write $\partial_\nu f(x)=\eps^{-1}(f(x+\eps e_\nu)-f(x))$, identify a 1-form $A$ with its components $A_\mu(x):=A_{\langle x,x+\eps e_\mu\rangle}$, and let $\partial_\mu^*$ be the adjoint of $\partial_\mu$ for $\ip\cdot\cdot_\eps$ (minus the backward difference). Then $(dA)_{\mu\kappa}=\partial_\mu A_\kappa-\partial_\kappa A_\mu$, $d^*A=\sum_\mu\partial_\mu^*A_\mu$, and:
\begin{enumerate}
\item[(H-a)] $\norm{dA}_{2,\eps}^2+\norm{d^*A}_{2,\eps}^2=\sum_{\mu,\nu}\norm{\partial_\nu A_\mu}_{2,\eps}^2$; that is, $d^*d+dd^*$ is the componentwise lattice Laplacian.
\item[(H-b)] $\ker d$ on functions is the space of constants, and $\ker(d^*d+dd^*)$ on 1-forms is spanned by $\mathbf 1_1,\mathbf 1_2,\mathbf 1_3$.
\item[(H-c)] $\mathrm{ran}\,d^*$ (on 1-forms) is orthogonal to the constants and invariant under $-\Delta_\eps$, and $-\Delta_\eps=d^*d$ there; this restriction is injective.
\item[(H-d)] $\Pi:=d(d^*d)^+d^*$, with $(\cdot)^+$ the Moore--Penrose inverse, is the orthogonal projection onto $\mathrm{ran}\,d$ (on 2-forms).
\end{enumerate}
\end{lemma}
\begin{proof}
(H-a) All forward and backward differences commute on the torus, and $\norm{\partial_\mu^*f}=\norm{\partial_\mu f}$. Hence $\ip{\partial_\mu^*A_\mu}{\partial_\nu^*A_\nu}=\ip{A_\mu}{\partial_\mu\partial_\nu^*A_\nu}=\ip{\partial_\nu A_\mu}{\partial_\mu A_\nu}$. Expanding the squares, the mixed terms of $\norm{dA}^2$ (summed over unordered pairs $\mu<\kappa$, as the weight $\eps^3$ per unordered plaquette prescribes) and of $\norm{d^*A}^2$ cancel. (H-b) follows from (H-a), because a function with vanishing forward differences is constant on the connected torus. (H-c) $\ip{d^*\bg}{\mathbf 1_\mu}=\ip{\bg}{d\mathbf 1_\mu}=0$, and $-\Delta_\eps d^*\bg=d^*dd^*\bg+dd^*d^*\bg=d^*(dd^*\bg)$, since $d^*d^*=0$. The restriction of $d^*d$ to $\mathrm{ran}\,d^*=(\ker d)^\perp$ is injective, because $d^*dv=0$ implies $\norm{dv}=0$. (H-d) $\Pi$ is self-adjoint. Since $(d^*d)^+(d^*d)$ is the orthogonal projection onto $(\ker d)^\perp$, $\Pi d=d$; hence $\Pi^2=\Pi$ and $\mathrm{ran}\,d\subset\mathrm{ran}\,\Pi\subset\mathrm{ran}\,d$.
\end{proof}
Parts (H-a) and (H-b) are also contained in Lemma~\ref{app:G1:X}, which identifies $-\Delta_\eps$ with the lattice Laplacian on each component and gives the spectral gap $-\Delta_\eps\ge16/\ell_T^2$ on $\mathrm{ran}\,\Pi^\perp$.

\begin{lemma}[Gaussian domination of the curvature on tori]\label{lem:C-Fsector}
For every real lattice 2-form $\bg$, every real lattice function $k$ and every $\mu_0\in(0,\mu_*]\cup\{\fl\}$,
\[
\ev{e^{F(\bg)}e^{\XR(k)}}_{\eps,T,\mu_0}\le e^{\frac12V(\bg)}\ev{e^{\XR(k)}}_{\eps,T,\mu_0},\qquad V(\bg):=\ip{d^*\bg}{C\,d^*\bg}_\eps\le\norm{\Pi_{\mathrm{ran}\,d}\bg}_{2,\eps}^2\le\norm{\bg}_{2,\eps}^2 ,
\]
with $C=C_{\mu_0}$, or $C=(-\Delta_\eps)^{-1}\Pi^\perp$ in the state $\fl$, and $\Pi_{\mathrm{ran}\,d}$ the orthogonal projection onto the exact 2-forms.
\end{lemma}
\begin{proof}
Put $f:=d^*\bg$, so that $F(\bg)=\ip{f}{A}_\eps$. Since $\ip{f}{\mathbf 1_\mu}_\eps=\ip{\bg}{d\mathbf 1_\mu}_\eps=0$, $f$ is orthogonal to the constants and $F(\bg)=\ip{f}{A^\perp}_\eps$. By Subsection~\ref{ssec:observables}, $\int d\varphi\,e^{-S_m+\XR(k)}=c_k\,W_k(A)$ with the positive constant $c_k=e^{-\sum_x\eps^3k(x)(W_\eps+z_\eps)}$, which cancels in the ratios below, and $W_k$ has the current representation \eqref{eq:C-current} of Lemma~\ref{lem:postype}. Write its phase as $e\eps\sum_bn_bA_b=e\ip{j_n}{A}_\eps$ with $j_n:=\eps^{-2}n$.

\emph{Case $\mu_0>0$.} By \eqref{eq:C-state} and Lemma~\ref{lem:postype} (the series converges absolutely, so Fubini applies), with $C=C_{\mu_0}$,
\[
\E_{\gamma_C}\big[e^{\ip{f}{A}_\eps}W_k(A)\big]=\sum_nw_k(n)\,\E_{\gamma_C}\big[e^{\ip{f+iej_n}{A}_\eps}\big]=e^{\frac12\ip{f}{Cf}_\eps}\sum_nw_k(n)\,e^{-\frac12e^2\ip{j_n}{Cj_n}_\eps}\cos\big(e\ip{j_n}{Cf}_\eps\big),
\]
where the sine terms cancel in pairs because $w_k(n)=w_k(-n)$. The cosine is at most $1$, and $\sum_nw_k(n)e^{-\frac12e^2\ip{j_n}{Cj_n}_\eps}=\E_{\gamma_C}[W_k(A)]$ (the same identity at $f=0$). Divide by $\E_{\gamma_C}[W_0]$: the left side becomes $\ev{e^{F(\bg)}e^{\XR(k)}}$ and the right side is at most $e^{\frac12\ip{f}{Cf}}\ev{e^{\XR(k)}}$.

\emph{Case $\fl$.} Averaging \eqref{eq:C-current} over $c\in[0,P)^3$ with the uniform law replaces $w_k(n)$ by $w_k(n)\mathbf 1\{\omega(n)=0\}$ (Lemma~\ref{lem:postype}: the phase of the constant mode is $e\ell_T\sum_\mu c_\mu\omega_\mu(n)$ and $\abs e\ell_TP=2\pi$), which is still nonnegative and even in $n$. The Gaussian step is the same, with $A^\perp$ of covariance $C=(-\Delta_\eps)^{-1}\Pi^\perp$ and $f\in\mathrm{ran}\,\Pi^\perp$.

\emph{The constant.} By Lemma~\ref{p:lem:C-hodge}(H-c), $\mathrm{ran}\,d^*$ (on 1-forms) reduces $-\Delta_\eps$, which equals $d^*d$ there and is injective there. Hence on $\mathrm{ran}\,d^*$ we have $C_{\mu_0}=(d^*d+\mu_0^2)^{-1}\le(d^*d)^{-1}$, and $(-\Delta_\eps)^{-1}\Pi^\perp=(d^*d)^{-1}$, since $\mathrm{ran}\,d^*$ is orthogonal to the constants. Therefore $V(\bg)\le\ip{d^*\bg}{(d^*d)^{+}d^*\bg}_\eps=\ip{\bg}{d(d^*d)^{+}d^*\bg}_\eps=\norm{\Pi_{\mathrm{ran}\,d}\bg}_{2,\eps}^2$ by (H-d).
\end{proof}

\begin{remark}\label{p:rem:C-Fsector}
On the torus the constant is $\tfrac12\norm{\bg}^2_{2,\eps}$ at every lattice spacing and for every $\mu_0\ge0$, the state $\fl$ included. It needs no infinite-volume input and no infrared bound in the sense of \cite{FSS76}: the Hodge decomposition is exact on the periodic lattice.
\end{remark}

\begin{lemma}[monotonicity in the photon covariance on one torus]\label{lem:C-cov}
Fix $(\eps,T)$ and the matter part of \eqref{eq:C-action}, and let $\ev{\cdot}_{C}$ denote \eqref{eq:C-state} with $\gamma_{C_{\mu_0}}$ replaced by $\gamma_C$ for an arbitrary positive semidefinite $C$ on the bond functions. Let $g\ge0$ be a site function, $r^2(g):=\sum_xg_xr_x^2$, and $\ell$ a real linear functional of $A$. If $0\le C_0\le C_1$, then
\[
\ev{e^{r^2(g)}e^{i\ell(A)}}_{C_1}\le\ev{e^{r^2(g)}e^{i\ell(A)}}_{C_0},\qquad\ev{e^{-r^2(g)}e^{\ell(A)}}_{C_1}\ge\ev{e^{-r^2(g)}e^{\ell(A)}}_{C_0}.
\]
That is, the class $\cI$ of the functions $e^{r^2(g)}e^{i\ell(A)}$ is non-increasing in the covariance, and the class $\cD$ of the functions $e^{-r^2(g)}e^{\ell(A)}$ is non-decreasing.
\end{lemma}
This is the torus form of \cite[Cor.~6.3(1), p.~281]{BFS1}. The inequalities \cite[Thm.~6.2(1),(3), p.~280]{BFS1} that it uses are proved in \cite[Sect.~6.1, pp.~275--276]{BFS1} for an arbitrary finite lattice and an arbitrary centred Gaussian measure $\gamma_C$, $C\ge0$, with Dirichlet conditions, with the remark (p.~276) that periodic conditions on rectangles are allowed. Lemma~\ref{app:G1:BFScov} and Proposition~\ref{app:G1:BFStorus} prove them on every finite graph, the lattice torus with its wrap-around bonds included.
\begin{proof}
First let $C_0,C_1$ be positive definite. Put $C_t=(1-t)C_0+tC_1$ and $D=C_1-C_0\ge0$. Differentiating the Gaussian density, $\partial_t\ev{G}_{C_t}=\tfrac12\sum_\alpha d_\alpha\,\ev{G;\,k_\alpha(A)^2}_{C_t}$, where $D=\sum_\alpha d_\alpha v_\alpha v_\alpha^{\mathsf T}$ with $d_\alpha\ge0$, $k_\alpha(A)=v_\alpha^{\mathsf T}C_t^{-1}A$, and $\ev{G;K}:=\ev{GK}-\ev{G}\ev{K}$; the normalisation constant cancels. By \cite[Thm.~6.2(1)]{BFS1} (Lemma~\ref{app:G1:BFS62}(1) on the torus graph) the truncated correlation is $\le0$ for $G\in\cI$, and by \cite[Thm.~6.2(3)]{BFS1} (Lemma~\ref{app:G1:BFS62}(3)) it is $\ge0$ for $G\in\cD$; integrate over $t\in[0,1]$. For positive semidefinite $C_i$ replace $C_i$ by $C_i+\eta$ and let $\eta\downarrow0$. Here $\int G\,d\gamma_C=\E[G(C^{1/2}\xi)]$, with $\xi$ standard normal, is continuous in $C$ by dominated convergence, since for the two classes the $\varphi$-integral of $\abs{G}e^{-S_m}$ is bounded by a constant times $e^{a\abs{A}}$, by (P2) and $\abs{\cos}\le1$.
\end{proof}

\begin{lemma}[$\mu_0\downarrow0$ at fixed lattice; the elementary observables]\label{lem:C-flatlimit}
Let $G\in\fF$, and let $\delta>0$ with
\[
\sup_{0<\mu_0\le\mu_*}\ev{\abs{G}^{1+\delta}}_{\eps,T,\mu_0}<\infty .
\]
Then $\ev{\abs{G}}_{\eps,T,\fl}<\infty$ and $\lim_{\mu_0\downarrow0}\ev{G}_{\eps,T,\mu_0}=\ev{G}_{\eps,T,\fl}$, along every sequence. The hypothesis holds, with a bound that does not depend on $\mu_0$, for every finite product of: bounded elements of $\fF$; $e^{tF(\bg)}$, $t\in\R$; $e^{\XR(q)}$ with $q\ge0$; $e^{-\XR(q)}$ with $q\ge0$; polynomials in the $F(\bg_i)$ and $\XR(h_k)$; and the moduli of the exponentials $\exp[\sum_iz_iF(\bg_i)+\sum_kw_k\XR(h_k)]$. We call these the \emph{elementary observables}.
\end{lemma}
\begin{proof}
The first statement is Lemma~\ref{lem:flat}, since $\fF=\cA_\eps$; its proof gives $\ev{\abs G^{1+\delta}}_{\eps,T,\fl}\le\sup_{\mu_0}\ev{\abs G^{1+\delta}}_{\eps,T,\mu_0}$, hence $\ev{\abs G}_{\eps,T,\fl}<\infty$. For the elementary observables, by H\"older's inequality it suffices to bound each factor, raised to a power $s\ge1$, uniformly in $\mu_0$. Bounded factors are trivial. For $e^{tF(\bg)}$, Lemma~\ref{lem:C-Fsector} with $k=0$ gives $\ev{e^{stF(\bg)}}_{\mu_0}\le e^{\frac12s^2t^2\norm{\bg}_{2,\eps}^2}$. For $e^{\XR(q)}$ with $q\ge0$: $e^{s\XR(q)}$ is a positive constant times an element of the class $\cI$ of Lemma~\ref{lem:C-cov} ($g=s\eps^3q$, $\ell=0$), and $C_{\mu_0}\ge C_{\mu_*}$ as forms for $\mu_0\le\mu_*$, the zero mode included (Lemma~\ref{app:G1:mu-mono}(b)). So Lemma~\ref{lem:C-cov} gives $\ev{e^{s\XR(q)}}_{\mu_0}\le\ev{e^{s\XR(q)}}_{\mu_*}<\infty$, the same $b_T$ being used on both sides (P1). For $e^{-\XR(q)}$ with $q\ge0$ the pointwise bound $e^{-\XR(q)}\le e^{(W_\eps+z_\eps)\sum_x\eps^3q(x)}$ holds, since $r^2\ge0$. Polynomials are bounded by sums of products of such exponentials: $\abs{x}^k\le k!(e^x+e^{-x})$, and $\abs{\XR(h)}\le\XR(\abs h)+2\abs{W_\eps+z_\eps}\sum_x\eps^3\abs h$. The modulus of a complex exponential is the real exponential with the real parts $\rRe z_i$, $\rRe w_k$, which is $e^{F(\bg)}e^{\XR(k_+)}e^{-\XR(k_-)}$ with $\bg=\sum_i(\rRe z_i)\bg_i$ and $k_\pm$ the positive and negative parts of $\sum_k(\rRe w_k)h_k$, a product of factors of the previous types.
\end{proof}
Consequently every closed property of the states at $\mu_0>0$ on elementary observables also holds for the state $\fl$: positivity (Proposition~\ref{prop:C-RP}) and the symmetries (Lemma~\ref{lem:C-sym}). A Feynman or Proca photon measure at $\mu_0=0$ does not exist on the torus, because the constant mode is not normalisable; the state $\fl$ is not such a measure, but the limit of Lemma~\ref{lem:flat}.

\subsection{Lattice reflection positivity on tori}\label{ssec:C-RP}

\emph{Scope.} The statements of this subsection hold for every even $N$ (P4), every direction $\nu$, both types of plane pair (through sites, and between lattice planes), every photon mass $\mu_0\in(0,\mu_*]$, every value of the constant bare coefficient $b_T$ (P1), and, through Lemma~\ref{lem:C-flatlimit}, for the state $\fl$. In particular planes through sites are covered (Lemma~\ref{lem:C-RPP}(i)); this is the case used in Section~\ref{sec:C-OS}, because the reflection $x_\nu\mapsto-x_\nu$ in a coordinate plane through the origin maps every lattice $\eps\Z^3$ of the family to itself. Positivity is proved for the Proca form of the photon measure, in the manner of \cite{OS78}, and transferred to King's Feynman form on gauge-invariant observables by Lemma~\ref{lem:C-FP}; the Feynman measure itself is not reflection positive on functions of $A$ that are not gauge invariant (Remark~\ref{p:rem:C-Feynman}). Three routine points are used below without further comment. (1) The reduction to $c=0$, resp.\ $c=\eps/2$, in the proof of Lemma~\ref{lem:C-RPP} is the invariance of the Proca measure under the lattice translations: its action is a sum of translation-covariant terms whose coefficients do not depend on the site, by (P1). (2) In case (i) of Lemma~\ref{lem:C-RPP} the integrals are finite: with the decomposition $S^{\mathrm P}=\bar S_+(\omega_0,\omega_+)+\bar S_+(\omega_0,\omega_-^\rho)$ of the proof, $\int d\omega_0\big(\int d\omega_+\,\abs G\,e^{-\bar S_+(\omega_0,\omega_+)}\big)^2\le\norm G_\infty^2\int d\omega_0\big(\int d\omega_+\,e^{-\bar S_+(\omega_0,\omega_+)}\big)^2=\norm G_\infty^2Z^{\mathrm P}<\infty$. (3) In case (ii) the coupling constants of the expansion are read off \eqref{eq:C-action}: a crossing plaquette contributes $\tfrac12\eps^3(dA)_p^2=\tfrac\eps2(u_p(\xi)-u_p(\xi'))^2$, and a crossing bond contributes $-\eps\,\varphi(b_-)\cdot U(\eps eA_b)\varphi(b_+)$ to the action, besides the single-site terms $\tfrac\eps2(r_{b_-}^2+r_{b_+}^2)$. So $\gamma_\iota=\eps>0$ in both cases; the hopping term enters $-S$ with the positive sign, which is the only sign condition the expansion uses.

\emph{Reflections.} Fix a direction $\nu$ and $c\in\tfrac\eps2\Z$, and let $\rho=\rho_c$ be the reflection $x_\nu\mapsto2c-x_\nu$ of $T$. Since $N$ is even (P4), $\rho$ fixes the pair of planes $\{x_\nu=c\}$ and $\{x_\nu=c+\ell_T/2\}$. These are planes between lattice planes when $c\in\tfrac\eps2+\eps\Z$, and planes through sites when $c\in\eps\Z$. The induced map on fields is $(\hat\rho\varphi)_x=\varphi_{\rho x}$ and $(\hat\rho A)_b=\pm A_{\rho b}$: the sign is $-$ exactly for the bonds parallel to $e_\nu$, whose orientation $\rho$ reverses. $\hat\rho$ is an involution and a signed permutation of the coordinates, so it preserves Lebesgue measure. The action \eqref{eq:C-action} is $\hat\rho$-invariant, by (P1) and because the plaquette, gradient, mass and hopping terms are mapped to terms of the same kind. $\hat\rho$ maps gauge transformations to gauge transformations: $\hat\rho(A+d\chi,U(-e\chi)\varphi)=(\hat\rho A+d(\chi\circ\rho),U(-e\chi\circ\rho)\hat\rho\varphi)$. Put $\theta G:=\overline{G\circ\hat\rho}$. On curvatures, $(d\hat\rho A)_p=\mp(dA)_{\rho p}$, with $-$ exactly for the plaquettes containing the direction $\nu$; hence $F(\bg)\circ\hat\rho=F(\rho^*\bg)$, where $\rho^*\bg$ is the signed reflected form, and $\theta$ maps $\fF$ into itself. Let $\Lambda_+$ be the open half $\{c<x_\nu<c+\ell_T/2\}$ for planes between lattice planes, or the closed half $\{c\le x_\nu\le c+\ell_T/2\}$ for planes through sites. Let $\mathfrak B_+$ be the algebra of bounded Borel functions of $\varphi|_{\Lambda_+}$ and of $A$ on the bonds with both ends in $\Lambda_+$, and $\fF_+:=\fF(\Lambda_+)$.

\begin{lemma}[Feynman and Proca forms agree on gauge-invariant observables]\label{lem:C-FP}
Let $\mu>0$, and let $\ev{\cdot}^{\mathrm P}$ be \eqref{eq:C-state} with the photon form $\tfrac12\ip{A}{(-\Delta_\eps+\mu^2)A}_\eps$ replaced by the Proca form $\tfrac12\ip{A}{(d^*d+\mu^2)A}_\eps$. If $G$ is gauge invariant and $\abs{G}\le c\,e^{a\abs{A}}$ for some constants $a,c$, then $\ev{G}_{\eps,T,\mu}=\ev{G}^{\mathrm P}$.
\end{lemma}
This is the torus form of \cite[Thm.~5.1, p.~270]{BFS1}.
\begin{proof}
On the connected torus, the space of bond functions is $\ker d^*\oplus\mathrm{ran}\,d$, orthogonally, and $\ker d$ on functions is the constants. So $A=Z+d\chi$ uniquely, with $d^*Z=0$ and $\chi\perp1$, and $(Z,\chi)\mapsto A$ is linear and bijective, with constant Jacobian; moreover $dA=dZ$. Then $\ip{A}{d^*dA}+\mu^2\ip AA=\ip{dZ}{dZ}+\mu^2\ip ZZ+\mu^2\ip{d\chi}{d\chi}$. The Feynman form adds $\ip{d^*A}{d^*A}=\ip{d^*d\chi}{d^*d\chi}$, which again depends on $\chi$ only. Both $\chi$-forms are positive definite on $1^\perp$. Since $Ge^{-S_m}$ is gauge invariant, $\int d\varphi\,(Ge^{-S_m})(Z+d\chi,\varphi)=\int d\varphi\,(Ge^{-S_m})(Z,U(e\chi)\varphi)=\int d\varphi\,(Ge^{-S_m})(Z,\varphi)$, by the rotation invariance of $d\varphi$. So the integral factorises into a $Z$-integral and a Gaussian $\chi$-integral. The $\chi$-factor is the same for $G$ and for $G\equiv1$, and cancels in the normalised expectation. Fubini's theorem applies by the bound on $G$, (P2) and $\mu>0$.
\end{proof}

\begin{lemma}[reflection positivity of the Proca torus measure]\label{lem:C-RPP}
Let $\mu>0$ and $N$ even. For every $G\in\mathfrak B_+$ (complex valued, not necessarily gauge invariant), $\ev{\theta(G)G}^{\mathrm P}\ge0$, in each of the two cases: (i) planes through sites; (ii) planes between lattice planes. Hence $(B,C)\mapsto\ev{\theta(C)B}^{\mathrm P}$ is a positive semidefinite Hermitian form on $\mathfrak B_+$.
\end{lemma}
\begin{proof}
Translate so that $c=0$ in case (i) and $c=\eps/2$ in case (ii).

(i) Every bond either has both endpoints in $\Lambda_+$ or is the $\hat\rho$-image of such a bond; the bonds lying in the two fixed planes are fixed by $\hat\rho$. Every plaquette lies in the closure of one half. Split each term that lives in a fixed plane (plaquettes, masses and hopping inside the plane, the single-site term at the plane sites) into two equal halves. The Proca action has no term that couples the two sides other than through the plane variables $\omega_0$; the gauge-fixing term $\ip{d^*A}{d^*A}$, which would, is absent. Hence $S^{\mathrm P}=\bar S_+(\omega_0,\omega_+)+\bar S_+(\omega_0,\omega_-^\rho)$, where $\omega_-^\rho$ denotes the variables of the other half transported by $\hat\rho$. Therefore $Z^{\mathrm P}\ev{\theta(G)G}^{\mathrm P}=\int d\omega_0\,\big|\int d\omega_+\,G(\omega_0,\omega_+)e^{-\bar S_+(\omega_0,\omega_+)}\big|^2\ge0$; the integrals are finite by (P2) and $\mu>0$.

(ii) The crossing bonds $b\in B_0$ (parallel to $e_\nu$, joining the two halves across one of the two planes) are mapped to themselves, with $A_b\mapsto-A_b$. Write each crossing variable as $a_b=\alpha_b-\alpha'_b$ with $\alpha_b,\alpha'_b$ independent centred normal of variance $1/(2M)$, $M=\mu^2\eps^3$. This uses up the crossing mass term $\tfrac12Ma_b^2$, and the reflection becomes the swap $\alpha\leftrightarrow\alpha'$. Put $\xi=(\omega_+,\alpha)$ and $\xi'=(\omega_-^\rho,\alpha')$; their a priori laws are identical. Every remaining term of $S^{\mathrm P}$ is of one of four kinds: a term of $S_+(\xi)$; its mirror image $S_+(\xi')$; a crossing plaquette; or a crossing hopping term. For a crossing plaquette $p$, $(dA)_p=\eps^{-1}(u_p(\xi)-u_p(\xi'))$ with the same real linear function $u_p$ on both copies, so its contribution $\tfrac\eps2(u_p(\xi)-u_p(\xi'))^2$ splits into a term in $\xi$, a term in $\xi'$, and $-\eps u_p(\xi)u_p(\xi')$. For a crossing hopping term, $\varphi(b_-)\cdot U(\eps e(\alpha_b-\alpha'_b))\varphi(b_+)=\psi_b(\xi)\cdot\psi_b(\xi')$ with the same $\R^2$-valued function $\psi_b$ on both copies. Hence $e^{-S'}=\mathcal G(\xi)\mathcal G(\xi')\exp\big(\sum_\iota\gamma_\iota Y_\iota(\xi)Y_\iota(\xi')\big)$ with $\gamma_\iota\ge0$ and real $Y_\iota$. Expanding the last exponential,
\[
Z^{\mathrm P}\ev{\theta(G)G}^{\mathrm P}=K\sum_{n}\prod_\iota\frac{\gamma_\iota^{n_\iota}}{n_\iota!}\Big|\int d\xi\,G(\xi)\mathcal G(\xi)\prod_\iota Y_\iota(\xi)^{n_\iota}\Big|^2\ge0,\qquad K>0 .
\]
The interchange of sum and integral is justified by $\sum_n\prod_\iota\gamma_\iota^{n_\iota}\abs{Y_\iota Y_\iota'}^{n_\iota}/n_\iota!\le\exp(\tfrac12\sum_\iota\gamma_\iota(Y_\iota^2+Y_\iota'^2))$ and by $\int\abs{G}\mathcal G\,e^{\frac12\sum_\iota\gamma_\iota Y_\iota^2}d\xi<\infty$: the plaquette parts cancel exactly, $\abs{\psi_b(\xi)}$ is a radius at a site next to a plane, and $\lambda>0$. The statement on the form follows by sesquilinearity and polarisation.
\end{proof}

\begin{proposition}[reflection positivity of the torus states]\label{prop:C-RP}
Let $N$ be even, $\nu\in\{1,2,3\}$, and let the plane pair be of either type.
\begin{enumerate}
\item[(a)] For $\mu_0\in(0,\mu_*]$ and every gauge-invariant Borel function $G$ of the $\Lambda_+$-variables with $\ev{\abs G^2}_{\eps,T,\mu_0}<\infty$: $\ev{\theta(G)G}_{\eps,T,\mu_0}\ge0$. Moreover $\sum_{i,k}\bar\gamma_i\gamma_k\ev{\theta(G_i)G_k}_{\eps,T,\mu_0}\ge0$ for every finite family of such $G_i$ and every $\gamma\in\C^n$.
\item[(b)] For the state $\fl$ and every $G\in\fF_+$ with $\ev{\abs G^2}_{\eps,T,\fl}<\infty$: $\ev{\theta(G)G}_{\eps,T,\fl}\ge0$, with the matrix form as in (a).
\item[(c)] Every exponential $\Theta=\exp[\sum_iz_iF(\bg_i)+\sum_kw_k\XR(h_k)]$, with $z,w\in\C$ and real sources, lies in $\fF$ and in $L^2$ of every state. If all sources are supported in $\Lambda_+$ (for each $\bg_i$: the plaquettes on which it does not vanish have all their bonds in $\Lambda_+$), then $\Theta\in\fF_+$, so (a) and (b) apply to finite families of such $\Theta$.
\end{enumerate}
\end{proposition}
\begin{proof}
(a) For bounded gauge-invariant $G\in\mathfrak B_+$, $\theta(G)G$ is bounded and gauge invariant, so $\ev{\theta(G)G}_{\eps,T,\mu_0}=\ev{\theta(G)G}^{\mathrm P}\ge0$ by Lemmas~\ref{lem:C-FP} and \ref{lem:C-RPP}. For unbounded $G$ put $G_n=G\mathbf 1\{\abs G\le n\}$: it is bounded, gauge invariant and measurable with respect to the $\Lambda_+$-variables, so $\ev{\theta(G_n)G_n}\ge0$, and $\abs{\theta(G_n)G_n}\le\abs{\theta G}\abs G$. By the Cauchy--Schwarz inequality and the $\hat\rho$-invariance of the measure, $\ev{\abs{\theta G}\abs G}\le\ev{\abs G^2}$, and dominated convergence gives the claim. The matrix form is (a) applied to $G=\sum_k\gamma_kG_k$.
(b) For bounded $G\in\fF_+$, $\theta(G)G$ is a bounded element of $\fF$, so $\ev{\theta(G)G}_\fl=\lim_{\mu_0\downarrow0}\ev{\theta(G)G}_{\mu_0}\ge0$ by Lemma~\ref{lem:C-flatlimit} and (a). The state $\fl$ is $\hat\rho$-invariant on $\fF$: $\ev{K\circ\hat\rho}_\fl=\lim\ev{K\circ\hat\rho}_{\mu_0}=\lim\ev{K}_{\mu_0}=\ev K_\fl$ for bounded $K\in\fF$. So truncation and dominated convergence work in the state $\fl$ exactly as in (a).
(c) $\abs\Theta^2=\exp[\sum_i2\rRe z_i\,F(\bg_i)+\sum_k2\rRe w_k\,\XR(h_k)]$ is integrable at fixed lattice: the Gaussian dominates $\abs{F}$, and (P2) dominates $\abs{w}\abs{h}r^2$. It is an elementary observable of Lemma~\ref{lem:C-flatlimit}, so it is in $L^2$ of the state $\fl$ too.
\end{proof}

\begin{remark}\label{p:rem:C-Feynman}
The Gaussian photon measure $\gamma_{C_{\mu_0}}$ in Feynman form is not reflection positive on functions of $A$ that are not gauge invariant. For instance, let $b$ be a bond in direction $e_\nu$ with both ends in $\Lambda_+$. Then $\theta(A_b)=-A_{\rho b}$, and $\E_{\gamma_{C_{\mu_0}}}[\theta(A_b)A_b]=-\E_{\gamma_{C_{\mu_0}}}[A_{\rho b}A_b]<0$, because by Lemma~\ref{app:G1:X}(i) the covariance $C_{\mu_0}$ acts on the $e_\nu$-components of $A$ as the resolvent $(-\Delta_\eps+\mu_0^2)^{-1}$ of the lattice Laplacian on functions, whose matrix elements are all positive. Proposition~\ref{prop:C-RP} uses only gauge-invariant observables, and that is all the paper needs.
\end{remark}

\subsection{Exact lattice symmetries}\label{ssec:C-sym}

\begin{definition}[cells, lattice forms, push-forward]\label{def:C-cells}
An \emph{oriented $k$-cell} of $T_\eps$, $k\in\{0,1,2\}$, is a pair $c=(x;v_1\wedge\dots\wedge v_k)$. Here $x$ is a site, the $v_i$ are distinct signed basis vectors $\pm e_\mu$, and the cell is the cube $x+\eps\sum_i[0,1]v_i$ with the orientation $v_1\wedge\dots\wedge v_k$. Two descriptions of the same cube with orientations of opposite sign give the cells $c$ and $-c$. A \emph{lattice $k$-form} is a real function $\omega$ on oriented $k$-cells with $\omega(-c)=-\omega(c)$. Sites are $0$-cells, bonds $1$-cells, plaquettes $2$-cells, and the centre of $c$ is $z(c):=x+\frac\eps2\sum_iv_i$. A signed permutation matrix $R\in\Oh$ acts on cells about the site $0$ by $Rc:=(Rx;Rv_1\wedge\dots\wedge Rv_k)$, and on forms by the push-forward $(R_*\omega)(c):=\omega(R^{-1}c)$. A continuum test function is \emph{sampled} by $h_\eps(x):=h(x)$ on sites, and a continuum 2-form $\bg$ by $\bg_\eps(c):=\bg_{z(c)}(v_1,v_2)$ on plaquettes $c=(x;v_1\wedge v_2)$, after identifying $T$ with the cell $[-\ell_T/2,\ell_T/2)^3$. For $c=(x;e_\mu\wedge e_\kappa)$ this is the sampling of $g_{\mu\kappa}$ at the plaquette centre of Definition~\ref{def:C-sampling}, and $\bg_\eps$ is odd under reversal of the orientation because $\bg_z$ is alternating.
\end{definition}

\begin{lemma}[push-forward commutes with $d$, with the pairing and with sampling]\label{lem:C-push}
Let $T$ be cubic, so that every $R\in\Oh$ maps $T_\eps$ onto itself (about the site $0$), and let $R\in\Oh$.
\begin{enumerate}
\item[(i)] $R$ maps oriented $k$-cells bijectively onto oriented $k$-cells, $R(-c)=-Rc$, and $R\,\partial c=\partial(Rc)$ for the oriented boundary. Hence $dR_*=R_*d$, $\ip{R_*\omega}{R_*\omega'}_\eps=\ip\omega{\omega'}_\eps$, and $d^*R_*=R_*d^*$.
\item[(ii)] With $(\hat RA,\hat R\varphi):=(R_*A,R_*\varphi)$: $F(\bg)\circ\hat R=F(R^{-1}_*\bg)$ and $\XR(h)\circ\hat R=\XR(R^{-1}_*h)$, where $R^{-1}_*h=h\circ R$.
\item[(iii)] For continuum test functions supported in $B_\rho$ with $\ell_T>2\rho+2$: $(R_*h)_\eps=R_*(h_\eps)$ and $(R_*\bg)_\eps=R_*(\bg_\eps)$, where $(R_*\bg)_z(v,w):=\bg_{R^{-1}z}(R^{-1}v,R^{-1}w)$ is the continuum push-forward of Subsection~\ref{ssec:symmetries}. There is no extra sign for improper $R$.
\item[(iv)] The action \eqref{eq:C-action} and Lebesgue measure are $\hat R$-invariant. Hence $\ev{G\circ\hat R}_{\eps,T,\mu_0}=\ev G_{\eps,T,\mu_0}$ for every integrable $G$ and every $\mu_0\in(0,\mu_*]$, and for every bounded or nonnegative $G\in\fF$ when $\mu_0=\fl$.
\end{enumerate}
\end{lemma}
\begin{proof}
(i) $R$ is linear, permutes the signed basis vectors, and maps $\eps\Z^3$ to itself. It also preserves the period lattice $\ell_T\Z^3$, so it acts on $T=(\R/\ell_T\Z)^3$ and on $T_\eps$, because $\ell_T/\eps\in\N$. The image of the cube $x+\eps\sum[0,1]v_i$ is $Rx+\eps\sum[0,1]Rv_i$, with orientation $Rv_1\wedge\dots\wedge Rv_k$. Orientation reversal commutes with $R$, and so does taking the oriented boundary, because the boundary of a cube is determined by its vertices and its orientation, both of which $R$ transports. So $(dR_*\omega)(c)=\omega(R^{-1}\partial c)=\omega(\partial R^{-1}c)=(R_*d\omega)(c)$, and $R_*$ permutes the unoriented cells, which all carry the weight $\eps^3$. This gives the pairing identity, and $d^*R_*=R_*d^*$ follows by adjunction, since $R_*^{-1}=(R^{-1})_*$.

(ii) $F(\bg)(R_*A)=\ip\bg{dR_*A}_\eps=\ip\bg{R_*dA}_\eps=\ip{R^{-1}_*\bg}{dA}_\eps$ by (i). The second identity is the change of summation variable $x\mapsto R^{-1}x$.

(iii) For scalars, $(R_*h)_\eps(x)=h(R^{-1}x)=(R_*h_\eps)(x)$. For a plaquette $c=(x;v_1\wedge v_2)$, $R^{-1}c=(R^{-1}x;R^{-1}v_1\wedge R^{-1}v_2)$ has centre $R^{-1}z(c)$. So $(R_*\bg_\eps)(c)=\bg_\eps(R^{-1}c)=\bg_{R^{-1}z(c)}(R^{-1}v_1,R^{-1}v_2)=(R_*\bg)_{z(c)}(v_1,v_2)=(R_*\bg)_\eps(c)$. Both sides use only the values of $\bg$ on signed basis vectors, so no determinant appears. Supports in $B_\rho$ with $\ell_T>2\rho+2$ keep the sampled functions away from the faces of the cell, so the identification of $T$ with the cell commutes with $R$ on them.

(iv) By Lemma~\ref{p:lem:C-hodge}(H-a) and (i), $\tfrac12\ip A{(-\Delta_\eps+\mu_0^2)A}_\eps=\tfrac12(\norm{dA}^2+\norm{d^*A}^2+\mu_0^2\norm A^2)$ is invariant under $A\mapsto R_*A$. For a bond $b$ with endpoints $b_\pm$, $(D_{R_*A}R_*\varphi)(b)=\eps^{-1}(U(\eps eA_{R^{-1}b})\varphi(R^{-1}b_+)-\varphi(R^{-1}b_-))=(D_A\varphi)(R^{-1}b)$. The modulus $\abs{D_A\varphi(b)}$ does not depend on the orientation of $b$, because $U(-s)=U(s)^{-1}$ and $U$ is orthogonal. So the hopping term is invariant. The single-site terms are invariant by (P1). $\hat R$ is a signed permutation of the coordinates, so it preserves Lebesgue measure. For $\fl$: $R_*$ is unitary and leaves the constants invariant, so it commutes with $\Pi^\perp$; it maps the constants $\mathbf 1_\mu$ to $\pm\mathbf 1_{\pi(\mu)}$, $\pi$ the permutation part of $R$. The Gaussian measure with covariance $(-\Delta_\eps)^{-1}\Pi^\perp$ is invariant by (i). The cube $[0,P)^3$ is mapped onto another period cell of $P\Z^3$, because the period $P=2\pi/(\abs{e}\,\ell_T)$ is the same in the three directions of a cubic torus. By Lemma~\ref{lem:holonomy} the integrand is $P$-periodic in $c$, so the uniform average does not change.
\end{proof}
\noindent\emph{The symmetry statement.} Lemma~\ref{lem:C-push}(ii)--(iv) give, for $R\in\Oh$ about a site and lattice sources,
\begin{equation}\label{p:eq:C-symm}
\Big\langle\prod_iF(R_*\bg_i)\prod_k\XR(R_*h_k)\Big\rangle_{\eps,T,\mu_0}=\Big\langle\prod_iF(\bg_i)\prod_k\XR(h_k)\Big\rangle_{\eps,T,\mu_0}
\end{equation}
(apply (iv) to $G:=\prod_iF(R_*\bg_i)\prod_k\XR(R_*h_k)$, for which $G\circ\hat R=\prod_iF(\bg_i)\prod_k\XR(h_k)$ by (ii)). With (iii), the same identity holds for sampled continuum sources. This is the input of Theorem~\ref{thm:C-OS1}(b). About an arbitrary point $y\in\R^3$ the cubic group acts by $R_y:=\tau_yR\tau_{-y}$. In the continuum this gives invariance about every point once translation invariance holds (Theorem~\ref{thm:C-OS1}(a)); on the lattice it is exact for $y\in\eps\Z^3$. The curvature $F$ transforms as a 2-form, also under the reflections in $\Oh$; the magnetic field $\star F$ would be a pseudovector. The reflection $\rho_\nu\in\Oh$ acts on $F$ exactly as the map $\hat\rho$ used for reflection positivity (Subsection~\ref{ssec:C-RP}), with the sign $-$ on the components $g_{\mu\nu}$.

\begin{lemma}[translations, the cubic group, charge conjugation]\label{lem:C-sym}
Let $T$ be cubic and $\mu_0\in(0,\mu_*]\cup\{\fl\}$. The state $\ev{\cdot}_{\eps,T,\mu_0}$ is invariant under:
\begin{enumerate}
\item[(i)] the lattice translations $\tau_a$, $a\in\eps\Z^3$ (mod $T$);
\item[(ii)] the cubic group $\Oh$, the $48$ signed permutations of the coordinates, acting about any site. An element $R\in\Oh$ acts on $\varphi$ as on a scalar and on $A$ as on a 1-form: $A_b\mapsto\pm A_{Rb}$, with the sign $-$ when $R$ reverses the orientation of the bond. Hence it acts by $h\mapsto h\circ R^{-1}$ on $\XR(h)$ and by $\bg\mapsto R_*\bg$ (signed push-forward of 2-forms) on $F(\bg)$;
\item[(iii)] charge conjugation $(A,\varphi)\mapsto(-A,\bar\varphi)$, which maps $F(\bg)\mapsto-F(\bg)$ and $\XR(h)\mapsto\XR(h)$.
\end{enumerate}
\end{lemma}
\begin{proof}
For $\mu_0>0$: the Feynman form $\tfrac12\ip{A}{(-\Delta_\eps+\mu_0^2)A}_\eps$ is the componentwise lattice Laplacian plus a mass, and it is invariant under translations and under signed permutations of the components. The hopping term $\abs{D_A\varphi}^2$ is invariant under (i) and (ii); when the orientation of a bond is reversed, $U(-s)=U(s)^{\mathsf T}$ exchanges the roles of the endpoints. The single-site term is invariant because $b_T$ is site independent (P1). $\Oh$ maps $T_\eps$ to itself because $T$ is cubic and the centre is a site; about another site, conjugate by the translations of (i). Lebesgue measure is invariant under all these maps. For (iii), the Gaussian measure is even and $\abs{D_{-A}\bar\varphi}=\abs{D_A\varphi}$. For $\fl$: the statements are closed under the limit of Lemma~\ref{lem:C-flatlimit} on $\fF$. Directly: the constant mode $c$ is translation invariant; $\Oh$ permutes and negates its components; charge conjugation negates it; and the uniform law on $[0,P)^3$ is invariant under all of these modulo $P$, because the period $P=2\pi/(\abs{e}\,\ell_T)$ is the same in the three directions of a cubic torus.
\end{proof}
For the limits of King's cube states \cite{K2}, $\Oh$-invariance would follow from the uniqueness of the infinite-volume limit, which is not known (Section~\ref{sec:disc}); on the cubic tori the cubic group is an exact symmetry of every lattice state.

\subsection{The tube inclusion}\label{ssec:C-T9}

For real lattice sources $\bg_1,\dots,\bg_n$, $h_1,\dots,h_p$ on $T_\eps$, the lattice generating functional $G_{\eps,T,\mu_0}$ and the lattice tube $\mathcal D_{\sigma_+,H}$ are those of Subsection~\ref{ssec:ch-T9}; $G_{\eps,T,\mu_0}$ is an entire function of $\zeta\in\C^{n+p}$ (Proposition~\ref{prop:C-RP}(c) and dominated convergence; Subsection~\ref{ssec:observables}). Put
\[
\sigma_0:=\tfrac12\min(\sigma_+,H),\qquad D_{\sigma_0}:=\Big\{\zeta\in\C^{n+p}:\ \sum_k\abs{\rRe w_k}\norm{h_k}_\infty<\sigma_0\Big\}.
\]
Then $D_{\sigma_0}\subset\mathcal D_{\sigma_+,H}$, since $\norm{k_\pm(\zeta)}_\infty\le\norm{k(\zeta)}_\infty$. The base of the tube is open and convex, since $s\mapsto\norm{(\sum_ks_kh_k)_\pm}_\infty$ is continuous and convex; so the tube is connected. The same holds for the tube $D_{\sigma_+,H}$ of Subsection~\ref{ssec:observables}, built from continuum sources, and for the set $D_{\sigma_0}$ defined with the continuum norms.

\begin{lemma}[the continuum tube lies in every sampled tube]\label{lem:C-tubeincl}
Let $\bg_1,\dots,\bg_n$, $h_1,\dots,h_p$ be real continuum test functions supported in $B_R$, and let $\zeta\in\C^{n+p}$ lie in the tube $D_{\sigma_+,H}$ of Subsection~\ref{ssec:observables}, which is built from the continuum function $k(\zeta)=\sum_k(\rRe w_k)h_k$. Then for every $(\eps,T)$ with $\ell_T>2R+2$, $\zeta$ lies in the lattice tube $\mathcal D_{\sigma_+,H}$ built from the sampled sources.
\end{lemma}
\begin{proof}
Sampling (Definition~\ref{def:C-sampling}) is evaluation at points, so it is linear and commutes with taking pointwise positive and negative parts: the sampled $k(\zeta)$ is $\sum_k(\rRe w_k)(h_k)_\eps$, and its positive and negative parts are the sampled $k_\pm(\zeta)$. A supremum over the sample points is at most the supremum over $\R^3$, so $\norm{(k_\pm)_\eps}_\infty\le\norm{k_\pm}_\infty$.
\end{proof}
With Lemma~\ref{lem:C-tubeincl}, Theorem~\ref{thm:ch-T9} applies at every point of the tube $D_{\sigma_+,H}$ of Theorem~\ref{thm:main}(i), on every $(\eps,T)$ with $\ell_T>2R+2$. Theorem~\ref{thm:ch-T9} is the only input from Sections~\ref{sec:UV}--\ref{sec:chain} used in Sections~\ref{sec:C-lattice}--\ref{sec:C-OS}.

\section{Limits along subsequences}\label{sec:C-limits}

Throughout this section, ``uniformly'' means uniformly in $\eps$ ($K\ge K_0$), in $j\ge1$, and in $\mu_0\in(0,\mu_*]\cup\{\fl\}$. Every statement is proved from Theorem~\ref{thm:ch-T9} and the lattice facts of Section~\ref{sec:C-lattice}. In the torus counterterm scheme the bare coefficient $b_T$ changes with $(\eps,T)$, so a sequence of lattice states is a sequence of different models. None of the arguments below uses a relation between two members of the sequence. Each uses either a statement about one model, or a uniform bound passed to a limit, or an abstract compactness or holomorphy argument.

\subsection{Uniform bounds, compactness, tightness}\label{ssec:C-compact}\label{ssec:C-test}

Continuum test functions are placed on the lattice by sampling (Definition~\ref{def:C-sampling}), with the norms defined there. For $T=T_j$ the lattice Schwinger functions $S^{\eps,T,\mu_0}_{n,p}:=S^{\eps,j,\mu_0}_{n,p}$ are the mixed moments of Subsection~\ref{ssec:observables}, the Taylor coefficients at $0$ of $G_{\eps,T,\mu_0}$. For real test functions supported in $B_R$ with $\ell_T>2R+2$ they are finite, real, linear in each argument and symmetric within each field type (Lemma~\ref{app:G1:multilin}(a)).

\begin{proposition}[uniform multilinear bounds]\label{prop:C-bounds}
Uniformly, for all real lattice 2-forms $\bg_i$ and lattice functions $h_k$ (in particular for sampled test functions),
\[
\abs{S^{\eps,T,\mu_0}_{n,p}(\bg;h)}\le e^{\mathsf K\sigma_0}\,\mathrm{e}^n(n!)^{1/2}\,(\mathrm{e}/\sigma_0)^p\,p!\,\prod_i\norm{\bg_i}_{2,\eps}\prod_kN_\eps(h_k),
\]
where $\mathrm{e}$ is Euler's number and $N_\eps(h)=\max(\norm h_{L^1,\eps},\norm h_\infty)$.
\end{proposition}
\begin{proof}
Fix $\vartheta\in(0,1)$ and consider the polydisc $\abs{z_i}=\vartheta(\sqrt n\norm{\bg_i}_{2,\eps})^{-1}$, $\abs{w_k}=\vartheta\sigma_0/(pN_\eps(h_k))$. On it $\norm{\sum_i(\rRe z_i)\bg_i}_{2,\eps}\le\sqrt n$, $\norm{k(\zeta)}_\infty\le\vartheta\sigma_0<\sigma_0$ and $\norm{k(\zeta)}_{L^1,\eps}\le\vartheta\sigma_0$. So the polydisc lies in $D_{\sigma_0}\subset\mathcal D_{\sigma_+,H}$, and by \eqref{eq:C-T9} $\abs G\le e^{n/2+\mathsf K\sigma_0}$ there. The Cauchy estimate on the polydisc gives $\abs{S_{n,p}}\le e^{n/2+\mathsf K\sigma_0}\vartheta^{-n-p}n^{n/2}(p/\sigma_0)^p\prod\norm{\bg_i}_{2,\eps}\prod N_\eps(h_k)$. Let $\vartheta\uparrow1$ and use $e^{n/2}n^{n/2}\le\mathrm{e}^n(n!)^{1/2}$ and $p^p\le\mathrm{e}^pp!$.
\end{proof}

\begin{proposition}[compactness; identification of the tube functional]\label{prop:C-compact}
Let $(\eps_i,T_{j_i},\mu_i)$ be a sequence with $K_i\to\infty$, $\mu_i\in(0,\mu_*]\cup\{\fl\}$, and either $\ell_{T_{j_i}}\to\infty$ or $T_{j_i}\equiv T$ fixed. Write $S^{(i)}_{n,p}$ and $G_i$ for the corresponding lattice Schwinger functions and generating functionals. There is a subsequence along which:
\begin{enumerate}
\item[(a)] $S^{(i)}_{n,p}(\bg;h)$ converges for all $n,p$ and all test functions in $C_c$ (if $T_{j_i}\equiv T$ is fixed: all test functions supported in some ball $B_R$ with $2R+2<\ell_T$). The limit $S_{n,p}$ is multilinear and symmetric within each field type, and obeys the bound of Proposition~\ref{prop:C-bounds} with the continuum norms $\norm{\bg_i}_2$, $N(h_k)$;
\item[(b)] for every finite family of real test functions (supported in a ball $B_R$ with $2R+2<\ell_T$ if $T$ is fixed), $G_i\to G$ locally uniformly on $D_{\sigma_+,H}$. $G$ is holomorphic there, its Taylor coefficients at $0$ are the $S_{n,p}$, $\abs{G(\zeta)}\le G(\rRe\zeta)$, and $G$ obeys \eqref{eq:C-T9} with continuum norms.
\end{enumerate}
\end{proposition}
\begin{proof}
(i) \emph{Equicontinuity.} For supports in $B_R$ and $\eps\le1$: $\norm h_{L^1,\eps}\le\abs{B_{R+1}}\norm h_\infty$ and $\norm\bg_{2,\eps}\le(3\abs{B_{R+1}})^{1/2}\norm\bg_\infty$. Let $C_c^{(R)}$ be the space of continuous test functions with support in $\bar B_R$, with the supremum norm. By Proposition~\ref{prop:C-bounds} and Lemma~\ref{app:G1:multilin}(a), each $S^{(i)}_{n,p}$, restricted to $C_c^{(R)}$, is a multilinear form bounded uniformly in $i$ (for every $i$ for which sampling is defined, $\ell_{T_{j_i}}>2R+2$), hence uniformly Lipschitz on bounded sets. (ii) \emph{Diagonal argument.} $C_c^{(R)}$ is separable, as a subspace of the separable metric space $C(\bar B_R)$. Choose countable dense subsets $\mathcal Q_R\subset C_c^{(R)}$, $R\in\N$, and extract a diagonal subsequence along which $S^{(i)}_{n,p}$ converges on all tuples from $\bigcup_R\mathcal Q_R$, for all $(n,p)$. Given a tuple with entries in $C_c^{(R)}$ and $\delta>0$, choose entries of $\mathcal Q_R$ within $\delta$ in the supremum norm; by the uniform Lipschitz bound of (i) on $C_c^{(R)}$, the values of $S^{(i)}_{n,p}$ at the two tuples differ by $O(\delta)$ uniformly in $i$, so the values at the given tuple form a Cauchy sequence. Since every test function lies in some $C_c^{(R)}$, the convergence extends to all tuples (for fixed $T$, to all tuples with $2R+2<\ell_T$). The identities expressing multilinearity and symmetry pass to the limit (Lemma~\ref{app:G1:multilin}(b)). (iii) \emph{Norms.} Riemann sums give $\norm h_{L^1,\eps}\to\norm h_1$ and $\norm\bg_{2,\eps}\to\norm\bg_2$; lattice supremum norms are at most continuum supremum norms. (iv) \emph{Tube.} For a fixed finite family, \eqref{eq:C-T9} bounds $\{G_i\}$ uniformly on each compact subset of $D_{\sigma_+,H}$, because the tube condition is on the supremum norms of the sampled sources, which are at most the continuum ones (Lemma~\ref{lem:C-tubeincl}), and the lattice norms on the right side of \eqref{eq:C-T9} converge by (iii). By Montel's theorem in several variables, every subsequence has a further subsequence converging locally uniformly to a holomorphic function $\tilde G$. By Weierstrass's theorem its Taylor coefficients at $0$ are the limits of those of $G_i$, which are lattice Schwinger functions of the family and converge to the $S_{n,p}$ by (a). So $\tilde G$ is determined near $0$, hence on the connected tube by the identity theorem, and the whole subsequence converges. The inequality $\abs{G(\zeta)}\le G(\rRe\zeta)$ and the bound pass to the limit.
\end{proof}
Compactness replaces King's monotone limits: no monotonicity in the volume and no theorem of King about sources is used. The price is that the limits exist along subsequences and are not shown to be unique. For a fixed torus ($\eps\to0$ only) the same argument gives continuum states with the same bounds, defined on the tuples of test functions supported in a ball $B_R$ with $2R+2<\ell_T$ (not on periodic test functions on $(\R/\ell_T\Z)^3$, which would need a sampling of functions on the torus).

\begin{proposition}[tightness; the limits are probability laws on distributions]\label{prop:C-tight}
Let $\Phi_i=(F_i,X_i)$ be the lattice random fields $\bg\mapsto F(\bg)$, $h\mapsto\XR(h)$ along a subsequence as in Proposition~\ref{prop:C-compact} with $\ell_{T_{j_i}}\to\infty$. Their laws $\nu_i$ are tight on $H^{-s'}_{\mathrm{loc}}(\R^3;\Lambda^2\R^3\oplus\R)$ for every $s'>3$. Every weak limit point $\nu$ has the moments $\E_\nu[\prod_l\Phi(f_l)]=S_{n,p}(f)$ for $f_l\in C_c^\infty$, and its characteristic functional is $G$ at imaginary arguments.
\end{proposition}
\begin{proof}
Each $\Phi_i$ is a finite sum of weighted point masses (at sites, and at plaquette centres with a 2-form value), hence lies in $H^{-s}_{\mathrm{loc}}$ for $s>3/2$. Fix $R$, a cutoff $\chi\in C_c^\infty(B_{R+1})$ with $0\le\chi\le1$ and $\chi=1$ on $B_R$, the box $Q=[-R-1,R+1]^3$, and the orthonormal Fourier basis $(e_\alpha)_{\alpha\in\Z^3}$ of $L^2(Q)$. Then $N(\chi e_\alpha)\le C_R$ uniformly in $\alpha$, and Proposition~\ref{prop:C-bounds}, applied to real and imaginary parts with $n+p=2$, gives $\E\abs{\Phi_i(\chi e_\alpha)}^2\le C'_R$ for each component. Hence $\E\norm{\chi\Phi_i}^2_{H^{-s}_{\mathrm{per}}(Q)}=\sum_\alpha(1+\abs\alpha^2)^{-s}\E\abs{\Phi_i(\chi e_\alpha)}^2\le C'_R\sum_\alpha(1+\abs\alpha^2)^{-s}<\infty$ for $s>3/2$, uniformly in $i$. The embedding $H^{-s}_{\mathrm{per}}(Q)\hookrightarrow H^{-s'}_{\mathrm{per}}(Q)$ is compact for $s'>s$. By Chebyshev's inequality the laws of $\chi\Phi_i$ are tight for each $R$, and a diagonal choice over $R\in\N$ gives tightness on $H^{-s'}_{\mathrm{loc}}$. Prokhorov's theorem gives weak limit points. All moments $\E\abs{\Phi_i(f)}^{2q}$ are bounded uniformly (Proposition~\ref{prop:C-bounds}), so the moments converge by uniform integrability. The exponential moments at small real parameters are bounded uniformly (Theorem~\ref{thm:ch-T9}), so the characteristic functional of $\nu$ is $G(i\,\cdot)$.
\end{proof}

\begin{remark}\label{rem:C-tight}
The proof gives tightness on $H^{-s'}_{\mathrm{loc}}$ for every $s'>3/2$; the statement with $s'>3$ is the one used in Theorem~\ref{thm:main}(ii).
\end{remark}

\subsection{Order of limits}\label{ssec:C-order}

\begin{lemma}[diagonal choice]\label{lem:C-diag}
Let $\mathcal Q\subset C_c$ be a countable set of test functions such that, for every $R\in\N$, the elements of $\mathcal Q$ supported in $B_{R+1}$ are dense, for the supremum norm, in the set of continuous functions supported in $B_R$. Let $\mathcal T=\{t_1,t_2,\dots\}$ enumerate the tuples $t=(n,p;\bg_1,\dots,\bg_n;h_1,\dots,h_p)$ with entries in $\mathcal Q$, and let $R(t)$ be the least $R\in\N$ with all supports in $B_R$. For $l\ge1$ let $\Sigma^{(l)}$ be a Schwinger family that is the limit, on every tuple $t$ with $2R(t)+2<\varsigma_l$ (on all tuples if $\varsigma_l=\infty$), of the lattice Schwinger families $\Sigma^{(l,i)}:=S^{\eps_{l,i},T_{l,i},\mu_{l,i}}$ as $i\to\infty$, where $K_{l,i}\to\infty$, $\mu_{l,i}\in(0,\mu_*]\cup\{\fl\}$, and $\liminf_i\ell_{T_{l,i}}\ge\varsigma_l\in(0,\infty]$ with $\varsigma_l\to\infty$. Suppose $\Sigma^{(l)}\to\Sigma$ on all tuples. Then there are indices $i(l)$ such that $K_{l,i(l)}\to\infty$, $\ell_{T_{l,i(l)}}\to\infty$, the lattice families $\Sigma^{(l,i(l))}$ converge to $\Sigma$ on all tuples of test functions in $C_c$, and, for every finite family of real test functions, their generating functionals converge to the generating functional of $\Sigma$ locally uniformly on $D_{\sigma_+,H}$.
\end{lemma}
\begin{proof}
Put $\varsigma_l':=\min(\varsigma_l-1,l)$. For each $l$ choose $i(l)$ with $K_{l,i(l)}\ge l$, $\ell_{T_{l,i(l)}}\ge\varsigma_l'$, and $\abs{\Sigma^{(l,i(l))}(t_q)-\Sigma^{(l)}(t_q)}\le1/l$ for every $q\le l$ with $2R(t_q)+2<\varsigma_l'$. These are finitely many conditions, each of which holds for all large $i$ (the lattice values at $t_q$ are defined once $\ell_T>2R(t_q)+2$, and they converge to $\Sigma^{(l)}(t_q)$). Fix $q$. Since $\varsigma_l'\to\infty$, the condition at $t_q$ is imposed for all large $l$, and then $\abs{\Sigma^{(l,i(l))}(t_q)-\Sigma(t_q)}\le1/l+\abs{\Sigma^{(l)}(t_q)-\Sigma(t_q)}\to0$. Convergence on $\mathcal T$ extends to all tuples in $C_c$: for entries supported in $B_R$, approximated in the supremum norm by elements of $\mathcal Q$ supported in $B_{R+1}$, and for $\eps\le1$, Proposition~\ref{prop:C-bounds} and step (i) of the proof of Proposition~\ref{prop:C-compact} show that the lattice families are uniformly Lipschitz on bounded subsets of the functions supported in $\bar B_{R+1}$, uniformly in the $(l,i)$ with $\ell_{T_{l,i}}>2R+4$; $\Sigma$ obeys the same bound; and a $\delta/3$ argument with the density of $\mathcal Q$ applies. Finally, the sequence $(\eps_{l,i(l)},T_{l,i(l)},\mu_{l,i(l)})$ satisfies the hypotheses of Proposition~\ref{prop:C-compact} (with $\ell_T\to\infty$ and varying photon masses), so every subsequence of the generating functionals has a further subsequence converging locally uniformly on $D_{\sigma_+,H}$ to a holomorphic function whose Taylor coefficients at $0$ are those of $\Sigma$; this function is the same for all subsequences, so the whole sequence converges.
\end{proof}

\begin{proposition}[iterated limits are joint limit points]\label{prop:C-order}
Give the set of families $\{S_{n,p}\}$ that obey the bound of Proposition~\ref{prop:C-bounds} with continuum norms the topology of convergence on a fixed countable dense set of test-function tuples. On this uniformly bounded set the topology is metrisable and, by equicontinuity, equivalent to convergence on all tuples. Then:
\begin{enumerate}
\item[(i)] every iterated limit (first $\eps\to0$ at fixed $T_j$, then $j\to\infty$, at fixed $\mu_0$) is a limit point of the joint family $(\eps,T_j)\to(0,\infty)$ at that $\mu_0$;
\item[(ii)] every limit, as $\mu_k\downarrow0$, of such limit points at $\mu_0=\mu_k$ is a limit point of the joint family $(\eps,T_j,\mu_0)\to(0,\infty,0)$.
\end{enumerate}
\end{proposition}
\begin{proof}
(i) Apply Lemma~\ref{lem:C-diag} with $\Sigma^{(l)}$ the fixed-torus continuum limit on $T_{j_l}$ and $\varsigma_l=\ell_{j_l}$; each $\Sigma^{(l)}$ exists by Proposition~\ref{prop:C-compact} with $T$ fixed, since at fixed $T_j$ every spacing $\eps_K$, $K\ge K_0$, is admissible (P4). (ii) Apply Lemma~\ref{lem:C-diag} with $\Sigma^{(l)}$ the joint limit points at $\mu_0=\mu_l$ and $\varsigma_l=\infty$. In both cases the limit is a limit of a sequence of lattice states of the joint family, which is all that Section~\ref{sec:C-OS} uses.
\end{proof}
The paper uses the limits in the following order. (1) At fixed lattice, $\mu_0\downarrow0$ is a full limit, and it gives the state $\fl$ (Lemmas~\ref{lem:flat} and~\ref{lem:C-flatlimit}). (2) At fixed $\mu_0\in(0,\mu_*]$, and for $\fl$, the joint limit $(\eps,T_j)\to(0,\R^3)$ is taken along subsequences (Proposition~\ref{prop:C-compact}); by Proposition~\ref{prop:C-order}(i) the iterated order gives nothing new. (3) Afterwards $\mu_0\downarrow0$ is taken through a countable set $Q_0$, along one common $(\eps,T)$-subsequence (Proposition~\ref{prop:C-mu0}). So $S^\fl$ is the ``$\mu_0$ first'' order and $S^0$ the ``$\mu_0$ last'' order. Proposition~\ref{prop:C-mu0}(c) compares them by one inequality on $E^\cD$; no interchange of limits is claimed. Every property proved below is closed under the convergence of Proposition~\ref{prop:C-order}, so it holds for every limit point in every one of these orders.

\subsection{The limit \texorpdfstring{$\mu_0\downarrow0$}{mu0 to 0}}\label{p:ssec:C-mu0}

\emph{A uniqueness set.} Fix real test functions $\bg_1,\dots,\bg_n$, $h_1,\dots,h_p$. Write each scalar source as $h_k=h_k^1-h_k^2$ with continuous $h_k^1,h_k^2\ge0$ of compact support (for instance the positive and negative parts of $h_k$), and consider the generating functional of the larger family $(\bg_1,\dots,\bg_n;h^1_1,\dots,h^1_p,h^2_1,\dots,h^2_p)$ in the variables $(z,u,v)\in\C^{n+2p}$, that is, of
\[
\Theta(z,u,v)=\exp\Big[\sum_iz_iF(\bg_i)+\sum_k\big(u_k\XR(h_k^1)-v_k\XR(h_k^2)\big)\Big].
\]
The generating functional of the original family at $(z,w)$ is its value at $(z,w,w)$, and the map $(z,w)\mapsto(z,w,w)$ maps the tube of the original family into the tube of the larger one, since both points have the same function $k$. For $\delta'>0$ put
\[
E^\cD_{\delta'}:=\R^n\times(-\delta',0)^p\times(0,\delta')^p .
\]
If $\delta'\sum_k(\norm{h^1_k}_\infty+\norm{h^2_k}_\infty)<H/2$, then $E^\cD_{\delta'}$ is a nonempty open subset of the real points of the tube of the larger family, and on it $\Theta=e^{F(\sum_ix_i\bg_i)}e^{-\XR(q)}$ with $q=\sum_k(\abs{u_k}h^1_k+v_kh^2_k)\ge0$ and $\norm q_\infty<H/2$; so $E^\cD_{\delta'}$ is contained in the set $E^\cD$ of Subsection~\ref{ssec:observables} for the larger family. A holomorphic function on the connected tube that vanishes on $E^\cD_{\delta'}$ vanishes identically: its Taylor coefficients at a point of $E^\cD_{\delta'}$ are real derivatives, hence zero.

\begin{proposition}[the limit $\mu_0\downarrow0$ on the tube]\label{prop:C-mu0}\label{app:G1:mu0}
Fix a countable set $Q_0\subset(0,\mu_*]$ that accumulates at $0$ (other accumulation points are allowed), for instance $Q_0=\{\mu_*2^{-k}:k\ge0\}$ or a countable dense subset of $(0,\mu_*]$, and let $(K_l,j_l)$ be any sequence with $K_l\to\infty$ and $j_l\to\infty$. There is a subsequence $(\eps_i,T_{j_i})$ of $(\eps_{K_l},T_{j_l})$ along which $S^\mu:=\lim_i\ev{\cdot}_{\eps_i,T_{j_i},\mu}$ exists for every $\mu\in Q_0$, and $S^\fl:=\lim_i\ev{\cdot}_{\eps_i,T_{j_i},\fl}$ exists, in the sense of Proposition~\ref{prop:C-compact}. Along it, for every finite family of real test functions:
\begin{enumerate}
\item[(a)] for $\zeta\in E^\cD$, $\mu\mapsto G^\mu(\zeta)$ is non-increasing on $Q_0$, and $G^\mu(\zeta)\le G^\fl(\zeta)$;
\item[(b)] $S^0:=\lim_{\mu\downarrow0,\,\mu\in Q_0}S^\mu$ exists as a full limit within $Q_0$, that is, the same along every sequence in $Q_0$ tending to $0$. The convergence of the generating functionals is locally uniform on $D_{\sigma_+,H}$, the Schwinger functions converge, and $S^0$ obeys Proposition~\ref{prop:C-bounds} and \eqref{eq:C-T9} with continuum norms;
\item[(c)] $G^0(\zeta)\le G^\fl(\zeta)$ for $\zeta\in E^\cD$.
\end{enumerate}
\end{proposition}
\begin{proof}
\emph{Common subsequence.} Since $j_l\to\infty$, $\ell_{j_l}=j_l!\,\ell_0\to\infty$. So for each fixed $\mu\in Q_0\cup\{\fl\}$ the sequence $(\eps_{K_l},T_{j_l},\mu)$ satisfies the hypotheses of Proposition~\ref{prop:C-compact}. Enumerate $Q_0\cup\{\fl\}=\{\mu^{(1)},\mu^{(2)},\dots\}$. Apply Proposition~\ref{prop:C-compact} to the given sequence with $\mu\equiv\mu^{(1)}$, then to the resulting subsequence with $\mu\equiv\mu^{(2)}$, and so on. The diagonal subsequence is a subsequence of each of these, so the limits exist along it for every $\mu\in Q_0\cup\{\fl\}$.

(a) Fix $\zeta=(x,w)\in E^\cD$ and a member $(\eps,T)=(\eps_i,T_{j_i})$ of the subsequence; put $\bg:=\sum_ix_i\bg_i$ and $q:=-k(\zeta)\ge0$. By the linearity of sampling and of the fields, on $T_\eps$
\[
\Theta(\zeta)=e^{F(\bg_\eps)}e^{-\XR(q_\eps)}=c_\eps\,e^{-r^2(g)}e^{\ell(A)},\qquad g_x:=\eps^3q_\eps(x)\ge0,\quad\ell(A):=\ip{d^*\bg_\eps}{A}_\eps,
\]
with $c_\eps:=\exp\big((W_\eps+z_\eps)\sum_x\eps^3q_\eps(x)\big)>0$. By (P1), $c_\eps$ and the matter part $S_m$ of \eqref{eq:C-action} do not depend on the photon mass. The covariance $C_\mu$ is non-increasing in $\mu$ as a form, the zero mode included (Lemma~\ref{app:G1:mu-mono}(b)). By Lemma~\ref{lem:C-cov} (class $\cD$), $\ev{\Theta(\zeta)}_{\eps,T,\mu}\le\ev{\Theta(\zeta)}_{\eps,T,\mu'}$ for $\mu'\le\mu$. Letting $\mu'\downarrow0$ and using Lemma~\ref{lem:C-flatlimit} ($\Theta(\zeta)$ is elementary) gives $\ev{\Theta(\zeta)}_{\eps,T,\mu}\le\ev{\Theta(\zeta)}_{\eps,T,\fl}$. Since $E^\cD\subset D_{\sigma_+,H}$, Proposition~\ref{prop:C-compact}(b) gives convergence of both sides along the subsequence, and the inequalities pass to the limit. (The details are Lemma~\ref{app:G1:classD}.)

(b) First let $\zeta\in E^\cD$. By (a), $\mu\mapsto G^\mu(\zeta)$ is non-increasing on $Q_0$, and it is bounded above by \eqref{eq:C-T9} with continuum norms (Proposition~\ref{prop:C-compact}(b)), uniformly in $\mu$. Put $L(\zeta):=\sup_{\mu\in Q_0}G^\mu(\zeta)$. Given $\eta>0$, pick $\mu_\eta\in Q_0$ with $G^{\mu_\eta}(\zeta)>L(\zeta)-\eta$. Then $L(\zeta)-\eta<G^\mu(\zeta)\le L(\zeta)$ for every $\mu\in Q_0$ with $\mu\le\mu_\eta$. Every sequence in $Q_0$ tending to $0$ eventually lies below $\mu_\eta$, so it gives the limit $L(\zeta)$; such sequences exist because $0$ is an accumulation point of $Q_0$. This is the only place where a property of $Q_0$ is used. Now apply this to the larger family of the uniqueness set above, at the points of $E^\cD_{\delta'}\subset E^\cD$. The family $\{G^\mu\}_{\mu\in Q_0}$ of its generating functionals is locally bounded on its tube, by \eqref{eq:C-T9} with continuum norms. By Montel's theorem and the identity theorem with the open real uniqueness set $E^\cD_{\delta'}$ (the Vitali argument), every sequence $\mu_k\to0$ in $Q_0$ has a subsequence along which $G^{\mu_k}$ converges locally uniformly on the tube to a holomorphic function equal to the monotone limit on $E^\cD_{\delta'}$; this function does not depend on the sequence. Hence the whole family converges locally uniformly as $\mu\to0$ in $Q_0$, and so does its restriction to the tube of the original family under $(z,w)\mapsto(z,w,w)$. By Weierstrass's theorem the Schwinger functions converge. The bounds of Proposition~\ref{prop:C-bounds} and \eqref{eq:C-T9} hold for every $S^\mu$ with continuum norms and constants independent of $\mu$ (Proposition~\ref{prop:C-compact}), so they pass to $S^0$.

(c) follows from (a) and (b).
\end{proof}
For $\Theta=\exp[F(\bg)-\XR(q)]\in\cD$ put $S^{\mu_0}(\Theta):=G^{\mu_0}(\zeta)$, for any finite family and any $\zeta\in E^\cD$ with $\Theta(\zeta)=\Theta$; this does not depend on the choice, because the lattice values depend only on $(\bg,q)$, by linearity, and hence so do their limits (Lemma~\ref{app:G1:classD}). With this reading, Theorem~\ref{thm:main}(iv) is Proposition~\ref{prop:C-mu0}(a) and (c).

\begin{remark}[$\mu_0$ outside $Q_0$]\label{rem:C-mu0out}\label{p:rem:C-mu0out}
Let $\mu\in(0,\mu_*]\setminus Q_0$, and let $S^\mu_\star$ be any limit point at photon mass $\mu$ along a further subsequence (Proposition~\ref{prop:C-compact}). Then every selection of such limit points $S^\mu_\star$ converges to $S^0$ as $\mu\downarrow0$ through $(0,\mu_*]$, in the sense of Proposition~\ref{prop:C-mu0}(b). Continuity of $\mu\mapsto S^\mu$ for $\mu>0$ is not claimed.
\end{remark}
\begin{proof}
The lattice inequality in the proof of Proposition~\ref{prop:C-mu0}(a) holds for any three photon masses $\mu'<\mu<\mu''$, and the limits at $\mu',\mu''\in Q_0$ are the same along every further subsequence. Hence $G^{\mu''}(\zeta)\le G^\mu_\star(\zeta)\le G^{\mu'}(\zeta)\le G^0(\zeta)$ for $\zeta\in E^\cD$ and $\mu'<\mu<\mu''$ in $Q_0$, the last inequality by monotonicity. Given $\eta>0$, choose $\mu''\in Q_0$ with $G^{\mu''}(\zeta)>G^0(\zeta)-\eta$ (Proposition~\ref{prop:C-mu0}(b)); for $\mu<\mu''$ there is $\mu'\in Q_0$ with $\mu'<\mu$, since $Q_0$ accumulates at $0$, and then $\abs{G^\mu_\star(\zeta)-G^0(\zeta)}<\eta$. So $G^\mu_\star\to G^0$ on $E^\cD$, in particular on $E^\cD_{\delta'}$ for the larger family, and the Vitali argument of the proof of Proposition~\ref{prop:C-mu0}(b) gives locally uniform convergence on the tube and the convergence of the Schwinger functions.
\end{proof}

\begin{lemma}[the measure for $S^0$]\label{app:G1:S0measure}
Along the subsequence of Proposition~\ref{prop:C-mu0}, let $s'>3$. There is a Borel probability measure $\nu^0$ on $H^{-s'}_{\mathrm{loc}}(\R^3;\Lambda^2\R^3\oplus\R)$ such that $\E_{\nu^0}[\prod_l\Phi(f_l)]=S^0_{n,p}(f)$ for test functions in $C_c^\infty$, where $\Phi(f)=F(\bg)$ for a 2-form $f=\bg$ and $\Phi(f)=\XR(h)$ for a scalar $f=h$, and such that the characteristic functional of $\nu^0$ is $G^0$ at imaginary arguments.
\end{lemma}
Together with Proposition~\ref{prop:C-tight} for $S^\mu$ ($\mu\in Q_0$) and $S^\fl$, Lemma~\ref{app:G1:S0measure} gives the last sentence of Theorem~\ref{thm:main}(ii).
\begin{proof}
\emph{The laws at $\mu>0$.} For $\mu\in Q_0$, Proposition~\ref{prop:C-tight}, applied along the subsequence of Proposition~\ref{prop:C-mu0}, gives a Borel probability measure $\nu^\mu$ on $H^{-s'}_{\mathrm{loc}}$. Its moments are $S^\mu$ on $C_c^\infty$ test functions, and its characteristic functional is $G^\mu(i\,\cdot)$.

\emph{Uniform second moments.} Fix $R$, let $Q=[-R-1,R+1]^3$ and $(e_\alpha)_{\alpha\in\Z^3}$ be as in the proof of Proposition~\ref{prop:C-tight}, with the cutoff $\chi$ there, $0\le\chi\le1$. Let $f$ be one of the real functions $\chi\,\rRe\,e_\alpha$, $\chi\,\mathrm{Im}\,e_\alpha\in C_c^\infty$, placed in one field component. Then $\norm f_2\le1$, $\norm f_\infty\le\abs Q^{-1/2}$ and $\norm f_1\le\abs Q^{1/2}$, so $N(f)\le\abs Q^{1/2}$, since $\abs Q\ge1$. By Proposition~\ref{prop:C-compact}(a), $S^\mu$ obeys the bound of Proposition~\ref{prop:C-bounds} with continuum norms; its constant $e^{\mathsf K\sigma_0}\mathrm{e}^n(n!)^{1/2}(\mathrm{e}/\sigma_0)^pp!$ depends only on $\mathsf K$ and $\sigma_0$, which by Theorem~\ref{thm:ch-T9} do not depend on $\mu_0\in(0,\mu_*]\cup\{\fl\}$. With the moment identity of Proposition~\ref{prop:C-tight} this gives, for every $\mu\in Q_0$ and every $\alpha$, $\E_{\nu^\mu}\Phi(f)^2=S^\mu_2(f,f)\le C'_R:=e^{\mathsf K\sigma_0}\max\big(\sqrt2\,\mathrm{e}^2,\,2\mathrm{e}^2\sigma_0^{-2}\abs Q\big)$, using $S^\mu_{2,0}(\bg,\bg)\le\sqrt2\,e^{\mathsf K\sigma_0}\mathrm{e}^2\norm\bg_2^2$ for a 2-form component and $S^\mu_{0,2}(h,h)\le2e^{\mathsf K\sigma_0}\mathrm{e}^2\sigma_0^{-2}N(h)^2$ for the scalar component. Hence $\E_{\nu^\mu}\abs{\Phi(\chi e_\alpha)}^2\le2C'_R$ for each component, uniformly in $\alpha\in\Z^3$ and $\mu\in Q_0$. Fix $s\in(3/2,s')$. The map $u\mapsto\norm{\chi u}^2_{H^{-s}_{\mathrm{per}}(Q)}=\sum_\alpha(1+\abs\alpha^2)^{-s}\abs{u(\chi e_\alpha)}^2\in[0,\infty]$ is Borel on $H^{-s'}_{\mathrm{loc}}$, as a countable sum of continuous functions, and by Tonelli's theorem
\[
\E_{\nu^\mu}\norm{\chi\Phi}^2_{H^{-s}_{\mathrm{per}}(Q)}=\sum_\alpha(1+\abs\alpha^2)^{-s}\,\E_{\nu^\mu}\abs{\Phi(\chi e_\alpha)}^2\le 2C'_R\sum_\alpha(1+\abs\alpha^2)^{-s}<\infty ,
\]
uniformly in $\mu\in Q_0$.

\emph{Tightness and a limit point.} The embedding $H^{-s}_{\mathrm{per}}(Q)\hookrightarrow H^{-s'}_{\mathrm{per}}(Q)$ is compact, and Chebyshev's inequality bounds $\nu^\mu(\norm{\chi\Phi}_{H^{-s}_{\mathrm{per}}(Q)}>M)$ by $2C'_R\sum_\alpha(1+\abs\alpha^2)^{-s}/M^2$. The diagonal choice over $R\in\N$ of the proof of Proposition~\ref{prop:C-tight} shows that $\{\nu^\mu\}_{\mu\in Q_0}$ is tight on $H^{-s'}_{\mathrm{loc}}$. Choose $\mu_k\in Q_0$ with $\mu_k\downarrow0$; this is possible because $Q_0$ accumulates at $0$. By Prokhorov's theorem a subsequence of $(\nu^{\mu_k})$ converges weakly to a probability measure $\nu^0$.

\emph{Moments.} For $f_l\in C_c^\infty$ the map $u\mapsto\prod_{l=1}^Nu(f_l)$ is continuous on $H^{-s'}_{\mathrm{loc}}$. Also $\E_{\nu^\mu}\abs{\Phi(f)}^{2N}=S^\mu_{2N}(f,\dots,f)$ is bounded uniformly in $\mu$ (Proposition~\ref{prop:C-bounds}, continuum form). H\"older's inequality then bounds $\E_{\nu^\mu}\abs{\prod_l\Phi(f_l)}^2$ uniformly, so $\prod_l\Phi(f_l)$ is uniformly integrable along $(\nu^{\mu_k})$. Hence $\E_{\nu^0}[\prod_l\Phi(f_l)]=\lim_k S^{\mu_k}(f)=S^0(f)$, the last equality by Proposition~\ref{prop:C-mu0}(b).

\emph{Characteristic functional.} For real $t$ the point $it$ lies in $D_{\sigma_+,H}$, since $\rRe w=0$, and $u\mapsto e^{i\sum_lt_lu(f_l)}$ is bounded and continuous. So $\E_{\nu^0}e^{i\sum_lt_l\Phi(f_l)}=\lim_kG^{\mu_k}(it)=G^0(it)$, again by Proposition~\ref{prop:C-mu0}(b).
\end{proof}
Alternatively, by Proposition~\ref{prop:C-order}(ii) and Lemma~\ref{lem:C-diag}, $S^0$ is the limit of a sequence of lattice states $(\eps_{l,i(l)},T_{l,i(l)},\mu_l)$ with $\ell_T\to\infty$ and varying photon masses; that sequence satisfies the hypotheses of Proposition~\ref{prop:C-compact}, and the proof of Proposition~\ref{prop:C-tight} applies to it.

\begin{proposition}[what is identified]\label{prop:C-ident}
\begin{enumerate}
\item[(a)] The generating functional of every limit point is determined by its Schwinger functions (Proposition~\ref{prop:C-compact}(b)). $S^0$ is a monotone $\mu_0$-limit along the fixed $(\eps,T)$-subsequence, the same for every sequence in $Q_0$ tending to $0$ (Proposition~\ref{prop:C-mu0}(b)). At fixed lattice the state $\fl$ is the full limit $\mu_0\downarrow0$ (Lemma~\ref{lem:flat}).
\item[(b)] $G^0\le G^\fl$ on $E^\cD$ (Proposition~\ref{prop:C-mu0}(c)). If, for every finite family of real test functions and every $\zeta\in E^\cD$,
\[
(\mathrm U\text{-}\mu)_T:\qquad\lim_{\mu_0\downarrow0}\ \sup_{\eps,j}\ \abs{\ev{\Theta(\zeta)}_{\eps,T_j,\fl}-\ev{\Theta(\zeta)}_{\eps,T_j,\mu_0}}=0,
\]
the supremum running over the spacings and tori of the family on which the sources can be sampled, then $G^0(\zeta)=G^\fl(\zeta)$ for every $\zeta\in E^\cD$; that is, $\mu_0\downarrow0$ commutes with the $(\eps,T)$-limit on the class $\cD$.
\end{enumerate}
\end{proposition}
\begin{proof}
(a) is a summary of the results cited. (b) For $\mu\in Q_0$, $G^\fl(\zeta)-G^\mu(\zeta)$ is the limit along the subsequence of $\ev{\Theta(\zeta)}_{\eps_i,T_{j_i},\fl}-\ev{\Theta(\zeta)}_{\eps_i,T_{j_i},\mu}$ (Proposition~\ref{prop:C-compact}(b)), so $\abs{G^\fl(\zeta)-G^\mu(\zeta)}$ is at most the supremum in $(\mathrm U\text{-}\mu)_T$ at $\mu_0=\mu$. Let $\mu\to0$ in $Q_0$ and use Proposition~\ref{prop:C-mu0}(b).
\end{proof}
Neither $(\mathrm U\text{-}\mu)_T$ nor the equality $S^0=S^\fl$ is needed for Theorem~\ref{thm:main}; both are discussed in Section~\ref{sec:disc}.

\section{OS0, OS1 for translations and the cubic group, OS2, OS3}\label{sec:C-OS}

\emph{Limit points.} In this section a \emph{limit point} is any of: $S^{\mu_0}$ for $\mu_0\in Q_0$; $S^\fl$; $S^0$ (Proposition~\ref{prop:C-mu0}); a limit point at any $\mu_0\in(0,\mu_*]$ along a further subsequence (Proposition~\ref{prop:C-compact}); in every order of limits (Proposition~\ref{prop:C-order}). The fixed-torus continuum states are included where stated. Every limit point other than the fixed-torus states is a Schwinger family in the sense of Subsection~\ref{ssec:axioms}: $S_{0,0}=1$, because the lattice states are normalised; the $S_{n,p}$ are defined on all test functions in $C_c$, multilinear and symmetric within each field type (Proposition~\ref{prop:C-compact}(a)); and they are real, because the lattice moments of the real fields $F(\bg)$, $\XR(h)$ are real. For $S^0$ these identities pass to the limit of Proposition~\ref{prop:C-mu0}(b). The generating functional of a limit point on $D_{\sigma_+,H}$ is the limit $G$ of Proposition~\ref{prop:C-compact}(b), or of Proposition~\ref{prop:C-mu0}(b) for $S^0$.

\subsection{OS0}\label{p:ssec:C-OS0}

\begin{theorem}[OS0]\label{thm:C-OS0}
Every limit point $S$ satisfies:
\begin{enumerate}
\item[(a)] (tube form) for every finite family of real test functions, $G$ is holomorphic on $D_{\sigma_+,H}$, which is all of $\C$ in the coefficients $z_i$ of the curvature, and
\[
\abs{G(\zeta)}\le\exp\Big(\tfrac12\Big\|\sum_i(\rRe z_i)\bg_i\Big\|_2^2+\kappa_+\norm{k_+(\zeta)}_1+2S^B_\fl\norm{k_-(\zeta)}_1\Big);
\]
since $k_+(\zeta)$ and $k_-(\zeta)$ have disjoint supports and $\kappa_+,2S^B_\fl\le\mathsf K$, the right side is at most $\exp\big(\tfrac12\|\sum_i(\rRe z_i)\bg_i\|_2^2+\mathsf K\norm{k(\zeta)}_1\big)$, which is the tube form of OS0 in Definition~\ref{def:OS} with $C=\mathsf K$;
\item[(b)] (Schwinger functions) $\abs{S_{n,p}(\bg;h)}\le e^{\mathsf K\sigma_0}\mathrm{e}^n(n!)^{1/2}(\mathrm{e}/\sigma_0)^pp!\prod_i\norm{\bg_i}_2\prod_kN(h_k)$, with $\mathrm{e}$ Euler's number;
\item[(c)] (regularity) let $\abs f_4:=\sup_{x\in\R^3,\abs\alpha\le4}\abs{(1+x^2)^2D^\alpha f(x)}$ for scalar $f$, and $\abs\bg_4:=(\sum_{\mu<\kappa}\abs{g_{\mu\kappa}}_4^2)^{1/2}$ for 2-forms. Then
\[
\abs{S_{n,p}(\bg;h)}\le\sigma_{n+p}\prod_i\abs{\bg_i}_4\prod_k\abs{h_k}_4,\qquad\sigma_M:=e^{\mathsf K\sigma_0}C_0^MM!,\quad C_0:=\max\big(\mathrm{e}\pi/(2\sqrt2),\,\mathrm{e}\pi^2/\sigma_0\big).
\]
$S_{n,p}$ extends uniquely from $C_c^\infty$ to Schwartz test functions, complex test functions included, and defines a tempered distribution on $\R^{3(n+p)}$. For every $\beta>1$ there is $\alpha_\beta$ with $\sigma_M\le\alpha_\beta(M!)^\beta$. This is the regularity condition E0$''$ of \cite[(4.2), p.~287]{OS75} with $s=4$ and any $\beta>1$. For a 2-form with one component, $\bg=f\,dx_\mu\wedge dx_\kappa$, one has $\abs\bg_4=\abs f_4$; so the bound, read on such test functions, is E0$''$ for the components $F_{\mu\kappa}(f):=F(f\,dx_\mu\wedge dx_\kappa)$ of the curvature together with $\XR$, which is the form in which Definition~\ref{def:OS} states it.
\end{enumerate}
The constants are the same for $S^{\mu_0}$, $S^\fl$ and $S^0$.
\end{theorem}
\begin{proof}
(a) Proposition~\ref{prop:C-compact}(b) and Proposition~\ref{prop:C-mu0}(b), with the three-term form of \eqref{eq:C-T9}. (b) Proposition~\ref{prop:C-bounds}, passed to the limit (Proposition~\ref{prop:C-compact}(a)). (c) In $d=3$, $\int(1+x^2)^{-2}dx=\pi^2$ and $\int(1+x^2)^{-4}dx=\pi^2/8$. Hence $\norm f_1\le\pi^2\abs f_4$, $\norm f_2\le\frac{\pi}{2\sqrt2}\abs f_4$ and $\norm f_\infty\le\abs f_4$; so $N(h)\le\pi^2\abs h_4$, and $\norm\bg_2\le\frac\pi{2\sqrt2}\abs\bg_4$ by summing over the three components. Insert these into (b) and use $(n!)^{1/2}p!\le(n+p)!$. The seminorm $\abs\cdot_4$ is continuous on $\cS(\R^3)$ and $C_c^\infty$ is dense in $\cS$, so $S_{n,p}$ extends by continuity. Complex test functions are handled by multilinearity, at the cost of a factor at most $2^{n+p}$, and the kernel theorem gives the distribution. Finally $C_0^M/(M!)^{\beta-1}\to0$.
\end{proof}
OS0 rests on Theorem~\ref{thm:ch-T9} only. It does not use King's cube theorems, a decoupling into unit cubes, a comparison of a cube with a torus, or any monotonicity in the volume. In the tube, only the small amplitude $\sigma_+$ for positive sources is a limitation. Regularity in the stronger form of \cite{GJ}, for all complex $w$, would need a bound for positive sources of arbitrary size, which is not proved here.

\subsection{OS1}\label{p:ssec:C-OS1}

\begin{theorem}[OS1: translations and the cubic group]\label{thm:C-OS1}
\begin{enumerate}
\item[(a)] Every limit point is invariant under all translations $a\in\R^3$: $S_{n,p}(\tau_a\bg;\tau_ah)=S_{n,p}(\bg;h)$. The fixed-torus continuum states are defined on the tuples of test functions supported in a ball $B_R$ with $2R+2<\ell_T$ (Proposition~\ref{prop:C-compact}). They are locally translation invariant: $S_{n,p}(\tau_a\bg;\tau_ah)=S_{n,p}(\bg;h)$ for all such tuples and all $a\in\R^3$ with $2(R+\abs a)+2<\ell_T$.
\item[(b)] Every limit point is invariant under the cubic group $\Oh$ ($48$ elements) about every point of $\R^3$, acting by $h\mapsto h\circ R^{-1}$ and $\bg\mapsto R_*\bg$ (push-forward of 2-forms, signs included).
\end{enumerate}
So every limit point is invariant under the group $\R^3\rtimes\Oh$, with $F$ transforming as a 2-form and $\XR$ as a scalar.
\end{theorem}
\begin{proof}
(a) Let $a\in\eps_K\Z^3=2^{-K}\eps_0\Z^3$. For every $i$ with $K_i\ge K$ we have $a\in\eps_i\Z^3$, because the lattices are nested (P4). Sampling commutes with $\tau_a$ for test functions supported in $B_R$ once $\ell_{T_{j_i}}>2(R+\abs a)+2$, since no wrap-around occurs. By Lemma~\ref{lem:C-sym}(i), $S^{(i)}(\tau_af)=S^{(i)}(f)$ for these $i$. Passing to the limit gives invariance under the dense group $\Z[1/2]^3\eps_0$. For continuity, Proposition~\ref{prop:C-compact}(a) and multilinearity give $\abs{S(\tau_af)-S(f)}\le C_{n,p}\sum_l\norm{\tau_af_l-f_l}\prod_{l'\ne l}\norm{f_{l'}}$, with the norms $\norm\cdot_2$ and $N(\cdot)$. Each of these norms is at most $C_R\norm\cdot_\infty$ on functions supported in $B_{R+1}$, and $\norm{\tau_af-f}_\infty\to0$ as $a\to0$ for $f\in C_c$. Hence the limit points are invariant under $\R^3$. At a fixed torus the same argument applies to tuples supported in $B_R$ and to $a$ with $2(R+\abs a)+2<\ell_T$. In the continuity step, take $R'>R+\abs a$ with $2R'+2<\ell_T$ and dyadic $a_k\to a$ with $\abs{a_k}<R'-R$, so that all translates are supported in $B_{R'}$. The details are in Lemma~\ref{app:G1:OS1torus}. By Corollary~\ref{app:G1:OS1-chain} the identity holds more generally whenever both $f$ and $\tau_af$ are admissible. Invariance under $\R^3/\ell_T\Z^3$ holds for the periodic extension of Proposition~\ref{app:G1:OS1-per}, not for the states themselves, whose domain the torus group does not preserve (Remark~\ref{app:G1:OS1torus-rem}). For $S^0$, invariance passes to the limit of Proposition~\ref{prop:C-mu0}(b).
(b) A cubic symmetry $R\in\Oh$ about the origin maps $\eps_i\Z^3$ and the cubic torus $T_{j_i}$ to themselves, and the origin is a site of every $\eps_i\Z^3$. $R$ maps sites to sites and plaquette centres to plaquette centres, with the orientation sign of $R_*$, so it commutes with sampling (Lemma~\ref{lem:C-push}(iii)). Lemma~\ref{lem:C-sym}(ii), in the form \eqref{p:eq:C-symm}, and the limit give invariance about the origin. About any other point, conjugate by the translations of (a).
\end{proof}

\begin{remark}[cubic symmetry only]\label{rem:C-cubic}
Theorem~\ref{thm:C-OS1} is all the Euclidean invariance proved in this paper: the group is $\R^3\rtimes\Oh$, and reflection positivity is proved only about coordinate planes (Theorem~\ref{thm:C-OS2}). What is known about rotations of infinite order is Subsection~\ref{p:ssec:C-rot}, and the missing step is discussed in Section~\ref{sec:disc}.
\end{remark}

\subsection{OS2}\label{p:ssec:C-OS2}

\begin{theorem}[OS2: reflection positivity about every coordinate plane]\label{thm:C-OS2}
Let $S$ be a limit point, $\nu\in\{1,2,3\}$, $c\in\R$, and $\theta_c$ the reflection in the plane $\{x_\nu=c\}$. It acts on scalar test functions by $h\mapsto h\circ\rho_{\nu,c}$ and on 2-forms by the signed reflection $\bg\mapsto(\rho_{\nu,c})_*\bg$ (the components $g_{\mu\kappa}$ with $\nu\in\{\mu,\kappa\}$ change sign), and antilinearly on coefficients.
\begin{enumerate}
\item[(a)] (matrix form for exponentials) Let $\Theta_a=\exp[\sum_lz_{al}F(\bg_{al})+\sum_lw_{al}\XR(h_{al})]$, $a=1,\dots,N'$, with all sources supported in $\{x_\nu>c\}$, $z_{al},w_{al}\in\C$, and $\sum_l\abs{\rRe w_{al}}\norm{h_{al}}_\infty<\sigma_0$ for every $a$. Then $\sum_{a,b}\bar\gamma_a\gamma_bS(\theta_c(\Theta_a)\Theta_b)\ge0$ for all $\gamma\in\C^{N'}$.
\item[(b)] (Schwinger functions) $S(\theta_c(P)P)\ge0$ for every polynomial $P$ in the $F(\bg)$ and $\XR(h)$ with sources supported in $\{x_\nu>c\}$.
\end{enumerate}
Both hold for $S^{\mu_0}$ (every $\mu_0\in(0,\mu_*]$), for $S^\fl$ and for $S^0$.
\end{theorem}
\begin{proof}
(a) By Theorem~\ref{thm:C-OS1}(a), $S(\theta_c(\cdot)\,\cdot)$ with sources in $\{x_\nu>c\}$ equals $S(\theta_0(\cdot)\,\cdot)$ with the sources translated by $-ce_\nu$, so we may take $c=0$. Let all sources be supported in $\{x_\nu>\delta\}\cap B_R$ for some $\delta,R>0$. On the lattice use the pair of planes $\{x_\nu=0\}$, $\{x_\nu=\ell_T/2\}$ through sites (P4); $\rho_{\nu,0}$ maps $\eps_i\Z^3$ to itself, so sampling commutes with it. For $\eps_i<\delta$ and $\ell_{T_{j_i}}>2R+2$, every sampled scalar source lies in $\Lambda_+$, and every plaquette that carries a sampled 2-form has all its bonds at $x_\nu\ge\delta-\eps_i/2>0$. So $\Theta_a\in\fF_+$. Proposition~\ref{prop:C-RP}(a)--(c) gives $\sum\bar\gamma_a\gamma_b\ev{\theta(\Theta_a)\Theta_b}_i\ge0$ for every $i$, with $\mu_i\in(0,\mu_*]$ or $\fl$. Each entry $\ev{\theta(\Theta_a)\Theta_b}$ is the generating functional of the finite joint family $\{(\rho_{\nu,0})_*\bg_{al},h_{al}\circ\rho_{\nu,0},\bg_{bl},h_{bl}\}$ at the point $\zeta^*_{ab}=(\bar z_a,z_b,\bar w_a,w_b)$. The reflected and the unreflected sources have disjoint supports, so the supremum norms of $k_\pm(\zeta^*_{ab})$ are at most the larger of the two amplitudes, hence $<\sigma_0=\tfrac12\min(\sigma_+,H)$; that is, $\zeta^*_{ab}$ lies in the tube of that family. Proposition~\ref{prop:C-compact}(b), applied to the joint family along the same subsequence, gives the convergence of each entry: all Schwinger functions converge on all $C_c$ tuples, and the generating functional is determined by them. Positive semidefiniteness of a finite Hermitian matrix is a closed condition. For $S^0$, the matrices at $\mu\in Q_0$ converge entrywise by Proposition~\ref{prop:C-mu0}(b).
(b) $P$ is in $L^2$ of every lattice state and lies in $\fF_+$ for the indices $i$ above, so $\ev{\theta(P)P}_i\ge0$ by Proposition~\ref{prop:C-RP}. $\theta(P)P$ is a finite combination of mixed moments with reflected test functions, which converge by Proposition~\ref{prop:C-compact}(a) and, for $S^0$, by Proposition~\ref{prop:C-mu0}(b).
\end{proof}

For the Schwartz form of OS2 let $\Phi=(F_{12},F_{13},F_{23},\XR)$ be the four real field components, $\cV:=\Lambda^2\R^3\oplus\R$ and $\cV_\C$ its complexification, so that a test function of $n$ variables is a map $\R^{3n}\to\cV_\C^{\otimes n}\cong\C^{4^n}$. For $\nu\in\{1,2,3\}$ put $U_+:=\{x\in\R^3:x_\nu>0\}$ and
\[
\Theta f(x_1,\dots,x_n):=\overline{(\vartheta^{\otimes n}f)(\rho_{\nu,0}x_1,\dots,\rho_{\nu,0}x_n)},
\]
where $\vartheta$ is the diagonal matrix of the signs of the reflection on the four field components.

\begin{lemma}[density and continuity for the Schwartz form of OS2]\label{lem:C-tensordense}
Let $U\subset\R^3$ be open, $n\ge1$, and $V$ a finite-dimensional complex vector space.
\begin{enumerate}
\item[(a)] For every $g\in C^\infty_c(U^n;V^{\otimes n})$ there are a compact set $K'\subset U^n$ and finite sums $g_N$ of tensor products $f_1\otimes\cdots\otimes f_n$, $f_l\in C^\infty_c(U;V)$, with $\supp g_N\subset K'$ and $D^\alpha g_N\to D^\alpha g$ uniformly for every multi-index $\alpha$; in particular $g_N\to g$ in $\cS(\R^{3n};V^{\otimes n})$.
\item[(b)] Let $U=U_+$, and let $f\in\cS(\R^{3n};V^{\otimes n})$ vanish outside $U_+^n$. Then $f$ is a limit in $\cS$ of elements of $C^\infty_c(U_+^n;V^{\otimes n})$.
\item[(c)] Let $S$ be a limit point. For every $N$ there is a unique tempered distribution $\mathsf S_N\in\cS'(\R^{3N};(\cV_\C^{\otimes N})^*)$ with $\mathsf S_N(f_1\otimes\cdots\otimes f_N)=S_N(f_1,\dots,f_N)$ for $f_l\in C^\infty_c(\R^3;\cV_\C)$, and the form $Q(f_0,\dots,f_{n_0}):=\sum_{n,n'\le n_0}\mathsf S_{n+n'}(\Theta f_n\otimes f_{n'})$ is continuous on $\bigoplus_{n\le n_0}\cS(\R^{3n};\cV_\C^{\otimes n})$. Here $S_N(f_1,\dots,f_N)$ denotes the moments of the four-component field $\Phi$, finite sums of the mixed Schwinger functions $S_{n,p}$ with $n+p=N$.
\end{enumerate}
\end{lemma}
\begin{proof}
(a) The projections of $\supp g$ to the $n$ factors are compact subsets of $U$; let $K_1\subset U$ be a compact neighbourhood of their union, choose $\psi\in C^\infty_c(U)$ with $\psi=1$ on $K_1$, and put $K':=(\supp\psi)^n$. Then $\psi^{\otimes n}=1$ on $\supp g$, so $g=\psi^{\otimes n}g$. Let $Q$ be an open cube of side $\ell$ containing $\supp\psi$, and regard $g$ as a smooth $\ell\Z^{3n}$-periodic function. Its Fourier coefficients decay faster than every power, so the partial sums $s_N$ of its Fourier series converge to $g$ with all derivatives, uniformly on $Q^n$. Each term is $e^{2\pi ik\cdot x/\ell}v=\prod_le^{2\pi ik_l\cdot x_l/\ell}\,v$ with $v\in V^{\otimes n}$, a finite sum of elementary tensors $v_1\otimes\cdots\otimes v_n$; hence $g_N:=\psi^{\otimes n}s_N$ is a finite sum of tensor products of functions $\psi(x_l)e^{2\pi ik_l\cdot x_l/\ell}v_l\in C^\infty_c(U;V)$, with $\supp g_N\subset K'$, and $D^\alpha(g_N-g)=D^\alpha(\psi^{\otimes n}(s_N-g))\to0$ uniformly by the Leibniz rule. Uniform convergence of all derivatives with supports in a fixed compact set implies convergence in $\cS$.
(b) For $\delta>0$ let $f_\delta(x_1,\dots,x_n):=f(x_1-\delta e_\nu,\dots,x_n-\delta e_\nu)$. Translations act continuously on $\cS$, so $f_\delta\to f$ as $\delta\downarrow0$. Since $f$ vanishes outside $U_+^n$, $\supp f_\delta\subset\{x:x_{l,\nu}\ge\delta\ \forall l\}\subset U_+^n$. With $\chi\in C^\infty_c(\R^{3n})$ equal to $1$ on the unit ball and $\chi_R:=\chi(\cdot/R)$, $\chi_Rf_\delta\to f_\delta$ in $\cS$ as $R\to\infty$, and $\supp\chi_Rf_\delta$ is a compact subset of $U_+^n$.
(c) By Theorem~\ref{thm:C-OS0}(c), $(f_1,\dots,f_N)\mapsto S_N(f_1,\dots,f_N)$ extends, by multilinearity in the real and imaginary parts and by continuity, to a multilinear form on $\cS(\R^3;\cV_\C)^N$ that is bounded by $4^N\sigma_N\prod_l\abs{f_l}_4$, with $\abs{(\bg,h)}_4:=\max(\abs\bg_4,\abs h_4)$ (a factor $2^N$ from the curvature and scalar parts of each $f_l$, and $2^N$ from the real and imaginary parts), hence continuous. By the Schwartz kernel theorem \cite[\S V.3]{RS1} it is represented by a unique $\mathsf S_N$ as stated. The map $f\mapsto\Theta f$ (composition with the linear map $\rho_{\nu,0}^{\times n}$, multiplication by the constant sign matrix $\vartheta^{\otimes n}$, and complex conjugation) is a continuous antilinear map of $\cS$ to itself, and $(f,g)\mapsto f\otimes g$ is a separately continuous, hence continuous, bilinear map $\cS(\R^{3n})\times\cS(\R^{3n'})\to\cS(\R^{3(n+n')})$ of Fr\'echet spaces. So $Q$ is continuous.
\end{proof}

\begin{proposition}[OS2 for positive-time Schwartz test functions]\label{prop:C-OS2dens}
Let $\cS_+(\R^{3n})$ denote the Schwartz functions $f:\R^{3n}\to\cV_\C^{\otimes n}$ that vanish outside $U_+^n=\{x:x_{1,\nu}>0,\dots,x_{n,\nu}>0\}$; it contains the Schwartz functions with support in $U_+^n$. Then for every limit point $S$ and every finite sequence $f_n\in\cS_+(\R^{3n})$, $n=0,\dots,n_0$,
\[
\sum_{n,n'}\mathsf S_{n+n'}(\Theta f_n\otimes f_{n'})\ge0 .
\]
\end{proposition}
\begin{proof}
By Theorem~\ref{thm:C-OS2}(b), applied to complex test functions by multilinearity, the inequality holds when every $f_n$ is a finite sum of tensor products $f^{(1)}\otimes\cdots\otimes f^{(n)}$ with $f^{(l)}\in C_c^\infty(U_+;\cV_\C)$. Indeed, for such $f_n$ the sum equals $S(\theta_0(P)P)$ with $P=\sum_n\Phi(f_n)$, where $\Phi(f^{(1)}\otimes\cdots\otimes f^{(n)})=\prod_l\Phi(f^{(l)})$. Now let $f_n\in\cS_+(\R^{3n})$. By Lemma~\ref{lem:C-tensordense}(b), $f_n$ is a limit in $\cS$ of elements of $C_c^\infty(U_+^n;\cV_\C^{\otimes n})$, and by Lemma~\ref{lem:C-tensordense}(a) each of these is a limit in $\cS$ of finite sums of tensor products of functions in $C_c^\infty(U_+;\cV_\C)$. By Lemma~\ref{lem:C-tensordense}(c) the quadratic form is continuous, so positivity on these finite sums passes to $\cS_+$. The time-ordered support condition of \cite{OS73,OS75} is contained in the condition defining $\cS_+$. In \cite{OS73,OS75} the reflected function also reverses the order of the arguments; since each $S_N$ is symmetric under all permutations of its arguments (Remark~\ref{rem:C-OS3form}), this does not change the form, and the positivity of \cite{OS75} follows.
\end{proof}

\subsection{OS3}\label{p:ssec:C-OS3}

\begin{lemma}[symmetry and charge conjugation on the lattice]\label{lem:C-OS3lat}
Fix $(\eps,T)$ and $\mu_0\in(0,\mu_*]\cup\{\fl\}$. For a real test function $f=(\bg,h)$, $\bg$ a 2-form and $h$ a scalar, put $\Phi(f):=F(\bg)+\XR(h)$ (sampled on $T_\eps$), and $\mathsf Cf:=(-\bg,h)$. For $N\ge0$ and real $f_1,\dots,f_N$ the moments $M_N(f_1,\dots,f_N):=\ev{\prod_{l=1}^N\Phi(f_l)}_{\eps,T,\mu_0}$ are finite and real, multilinear and symmetric under all permutations of $(f_1,\dots,f_N)$, and
\[
M_N(\mathsf Cf_1,\dots,\mathsf Cf_N)=M_N(f_1,\dots,f_N).
\]
In particular $\ev{\prod_{i=1}^nF(\bg_i)\prod_{k=1}^p\XR(h_k)}_{\eps,T,\mu_0}=0$ for odd $n$, and the lattice generating functional satisfies $G_{\eps,T,\mu_0}(-z,w)=G_{\eps,T,\mu_0}(z,w)$ for all $(z,w)\in\C^{n+p}$.
\end{lemma}
\begin{proof}
$\prod_l\Phi(f_l)$ is a polynomial in the $F(\bg_l)$ and $\XR(h_l)$, hence an elementary observable, integrable in every state, the state $\fl$ included (Lemma~\ref{lem:C-flatlimit}). It is real, and the product of real random variables is commutative; sampling is linear, so $\Phi$ is linear in $f$ and $M_N$ is multilinear. By Lemma~\ref{lem:C-sym}(iii), every state, $\fl$ included, is invariant under charge conjugation $\mathcal C:(A,\varphi)\mapsto(-A,\bar\varphi)$: for $\mu_0>0$ the Gaussian measure is even and $S_m(-A,\bar\varphi)=S_m(A,\varphi)$; in the state $\fl$, $\mathcal C$ maps $A^\perp\mapsto-A^\perp$ and $c\mapsto-c$, and the uniform law on the period cell is invariant under $c\mapsto-c$ modulo $P$. So $\ev{K\circ\mathcal C}=\ev K$ for every integrable $K$. Since $F(\bg)\circ\mathcal C=-F(\bg)=F(-\bg)$ and $\XR(h)\circ\mathcal C=\XR(h)$, we get $\Phi(f)\circ\mathcal C=\Phi(\mathsf Cf)$, and the displayed identity follows. Taking $f_i=(\bg_i,0)$ for $i\le n$ and $f_{n+k}=(0,h_k)$ shows that the mixed moment equals $(-1)^n$ times itself. Finally, $\exp[\sum_iz_iF(\bg_i)+\sum_kw_k\XR(h_k)]$ is integrable (Proposition~\ref{prop:C-RP}(c)), and its composition with $\mathcal C$ is the same exponential with $z$ replaced by $-z$.
\end{proof}

\begin{theorem}[OS3; charge conjugation]\label{thm:C-OS3}
For every limit point, $S_{n,p}$ is symmetric under permutations of the (field component, test function) pairs within each field type, real for real test functions, and $S_{n,p}=0$ for odd $n$.
\end{theorem}
\begin{proof}
Lattice moments of commuting real random variables are symmetric and real. By Lemma~\ref{lem:C-sym}(iii) every lattice state, $\fl$ included, is invariant under charge conjugation, which maps $F(\bg)\mapsto-F(\bg)$ and $\XR(h)\mapsto\XR(h)$. So $S^{(i)}_{n,p}=(-1)^nS^{(i)}_{n,p}$. All three properties are closed under the limits of Propositions~\ref{prop:C-compact} and~\ref{prop:C-mu0}.
\end{proof}

\begin{remark}[OS3 in the form of \cite{OS75}]\label{rem:C-OS3form}
Let $S$ be a limit point, and let $S_N(f_1,\dots,f_N)$ be the moments of the four-component field $\Phi=(F_{12},F_{13},F_{23},\XR)$. Each $S_N$ is a finite sum of the mixed Schwinger functions $S_{n,p}$, $n+p=N$. By Theorem~\ref{thm:C-OS3} and Lemma~\ref{lem:C-OS3lat}, passed to the limit (Propositions~\ref{prop:C-compact}(a) and~\ref{prop:C-mu0}(b)), $S_N$ is real on real test functions in $C_c$, symmetric under all permutations of its arguments (mixed field types included), and invariant under $f_l\mapsto\mathsf Cf_l$; and $G^S(-z,w)=G^S(z,w)$ on the tube, since the tube is invariant under $z\mapsto-z$ (it restricts only $\rRe w$) and $G^S$ is the locally uniform limit of the lattice functionals there (Proposition~\ref{prop:C-compact}(b)). By Theorem~\ref{thm:C-OS0}(c), $S_N$ extends to a continuous multilinear form on $\cS(\R^3;\C^4)^{N}$. Symmetry and $\mathsf C$-invariance pass to the extension by continuity, and the reality condition takes the form $S_N(\bar f_1,\dots,\bar f_N)=\overline{S_N(f_1,\dots,f_N)}$, by expanding each $f_l$ into its real and imaginary parts. This is the symmetry axiom of \cite{OS73,OS75} for a boson field with several components; the extension of \cite{OS75} from one scalar field to several fields is \cite[Remark~5, p.~283]{OS75}.
\end{remark}

\subsection{Rotations of infinite order}\label{p:ssec:C-rot}

Theorem~\ref{thm:C-OS1} gives invariance under $\R^3\rtimes\Oh$ only. This subsection collects three elementary facts about the full orthogonal group $\Ot$, which are used in Section~\ref{sec:disc}.

\begin{lemma}[rational rotations]\label{lem:D:rational}
\begin{enumerate}
\item[(i)] If $R\in\Ot$ and $2^jR$ has integer entries for some $j\ge0$, then $R\in\Oh$.
\item[(ii)] The matrix
\[
R_3:=\frac13\begin{pmatrix}1&2&2\\2&1&-2\\-2&2&-1\end{pmatrix}
\]
lies in $\SO(3)$, $3R_3$ is integral, $R_3(3\eps\Z^3)\subset\eps\Z^3$ with index $27$, $R_3$ is the rotation about the axis $(1,1,0)$ by the angle $\vartheta$ with $\cos\vartheta=-\tfrac13$, and $R_3$ has infinite order. The same holds for the rotation $R_{12}(\theta_0)$ by the angle $\theta_0=\arctan(3/4)$ in the $(1,2)$-plane (King's rotation \cite[p.~326]{K2}), with $5R_{12}(\theta_0)$ integral and $\cos\theta_0=\tfrac45$.
\item[(iii)] No lattice $\Gamma\subset\R^3$ satisfies $R\Gamma=\Gamma$ for a rotation $R$ of infinite order.
\end{enumerate}
\end{lemma}
\begin{proof}
(i) Each column $v$ of $2^jR$ is an integer vector with $\abs v^2=4^j$. If $j\ge1$, then $v_1^2+v_2^2+v_3^2\equiv0\pmod4$; squares are $\equiv0$ or $1\pmod 4$, so all $v_i$ are even and $v/2$ is an integer vector of norm $2^{j-1}$. Descending, the columns of $R$ are integer unit vectors $\pm e_i$, so $R$ is a signed permutation matrix, that is, $R\in\Oh$.
(ii) The rows of $3R_3$ have squared norm $9$ and are pairwise orthogonal, and $\det(3R_3)=27$, so $R_3\in\SO(3)$. $\tr R_3=\tfrac13=1+2\cos\vartheta$ gives $\cos\vartheta=-\tfrac13$, and $(R_3-I)(1,1,0)^{\mathsf T}=0$. The index of $(3R_3)\Z^3$ in $\Z^3$ is $\abs{\det 3R_3}=27$. If $R_3^n=I$ for some $n\ge1$, the eigenvalues $e^{\pm i\vartheta}$ are roots of unity, hence algebraic integers, so $2\cos\vartheta=e^{i\vartheta}+e^{-i\vartheta}$ is an algebraic integer; being rational, it is an integer; but $2\cos\vartheta=-\tfrac23$. The same argument applies to $R_{12}(\theta_0)$, with $2\cos\theta_0=\tfrac85$.
(iii) If $R\Gamma=\Gamma$, the matrix of $R$ in a basis of $\Gamma$ is an integer matrix, and so are all its powers. These lie in the intersection of $GL_3(\Z)$ (conjugated by the basis matrix) with the compact group $\Ot$, which is discrete and compact, hence finite. So $R$ has finite order.
\end{proof}

\begin{remark}\label{p:rem:D:dyadic}
Part (i) shows that a comparison with rotated lattices cannot be made inside the dyadic family $\eps=2^{-K}\eps_0$ of Theorem~\ref{thm:main} in King's nested form: no rotation outside $\Oh$ maps a dyadic lattice into a dyadic refinement of itself. For blocking factor $3$, the family of Theorem~\ref{thm:mainL3}, the rotation $R_3$ maps $3\eps_K\Z^3=\eps_{K-1}\Z^3$ into $\eps_K\Z^3$, so the rotated lattice stays in the family.
\end{remark}

\begin{lemma}[continuity and closed stabilisers]\label{lem:D:cont}
Let $S$ be a limit point. For all $n,p$ and all test functions, the map $g\mapsto S_{n,p}(g_*\bg_1,\dots,g_*\bg_n;g_*h_1,\dots,g_*h_p)$ is continuous on $\Ot$. Hence $\mathrm{Stab}(S):=\{g\in\Ot:S_{n,p}(g_*\,\cdot)=S_{n,p}\ \text{for all }n,p\}$ is a closed subgroup of $\Ot$.
\end{lemma}
\begin{proof}
By Theorem~\ref{thm:C-OS0}(b), $S_{n,p}$ is multilinear and bounded by $C_{n,p}\prod_i\norm{\bg_i}_2\prod_kN(h_k)$, $N(h)=\max(\norm h_1,\norm h_\infty)$. For test functions supported in $B_R$, all these norms are bounded by $C_R\norm{\cdot}_\infty$ on functions supported in $B_{R+1}$, and $\norm{g_*f-f}_\infty\to0$ as $g\to1$ by the uniform continuity of $f$ (for 2-forms, $(g_*\bg)(x)=\Lambda^2g\,\bg(g^{-1}x)$). Telescoping over the arguments gives continuity at $1$, and continuity everywhere follows by applying this to the tuple $g_{0*}f$, since $g_*$ preserves the norms and the supports in $B_R$. The stabiliser is a subgroup and, by continuity, closed.
\end{proof}

\begin{lemma}[one rotation of infinite order suffices]\label{lem:D:onerot}
\begin{enumerate}
\item[(i)] Every closed subgroup of $\Ot$ that contains $\Oh$ and is infinite equals $\Ot$.
\item[(ii)] For every rotation $R$ of infinite order, the closure of the group generated by $R$ and $\Oh$ is $\Ot$. In particular this holds for $R_3$ and for $R_{12}(\theta_0)$.
\item[(iii)] If a limit point $S$ is invariant under one rotation $R$ of infinite order, it is invariant under $\Ot$.
\end{enumerate}
\end{lemma}
\noindent\emph{Inputs of the proof.} Part (i) uses two standard facts. The first is the closed-subgroup theorem: a closed subgroup of a Lie group is an embedded Lie subgroup \cite[Thm.~3.42]{War83}. The second is the identity $h(a\times b)=(\det h)\,(ha\times hb)$ for $h\in\Ot$ and $a,b\in\R^3$. To prove it, note that for every $c\in\R^3$,
\[
\ip{h(a\times b)}{hc}=\ip{a\times b}{c}=\det(a,b,c)=(\det h)^{-1}\det(ha,hb,hc)=(\det h)\ip{ha\times hb}{hc},
\]
since $(\det h)^{-1}=\det h$; and $hc$ runs over $\R^3$. Hence $hX_\omega h^{-1}v=h(\omega\times h^{-1}v)=(\det h)\,(h\omega\times v)$, that is, $hX_\omega h^{-1}=X_{(\det h)h\omega}$, where $X_\omega v:=\omega\times v$. The proof of (i) applies this only to rotations in $\Oh$, for which $\det h=1$. The infinite order of $R_3$ and of $R_{12}(\theta_0)$ is Lemma~\ref{lem:D:rational}(ii), and the continuity used in (iii) is Lemma~\ref{lem:D:cont}.
\begin{proof}
(i) Let $G\subset\Ot$ be closed and infinite, with $\Oh\subset G$. By the closed-subgroup theorem \cite[Thm.~3.42]{War83}, $G$ is an embedded Lie subgroup; it is compact, and a compact Lie group of dimension $0$ is finite, so its Lie algebra $\mathfrak g\subset\mathfrak{so}(3)$ is non-zero. $\mathfrak g$ is invariant under $\mathrm{Ad}(h)X=hXh^{-1}$, $h\in G$. Identify $\mathfrak{so}(3)$ with $\R^3$ by $\omega\mapsto X_\omega$; then $hX_\omega h^{-1}=X_{(\det h)h\omega}$. Let $0\ne\omega\in\mathfrak g$, with $\omega_i\ne0$. The rotations $\mathrm{diag}(1,-1,-1)$, $\mathrm{diag}(-1,1,-1)$, $\mathrm{diag}(-1,-1,1)$ lie in $\Oh$ and have determinant $1$; adding $\omega$ to its image under the $i$-th of them gives $2\omega_ie_i\in\mathfrak g$, so $e_i\in\mathfrak g$. The cyclic permutation of the coordinates is a rotation in $\Oh$ and maps $e_1\mapsto e_2\mapsto e_3$, so $\mathfrak g=\R^3=\mathfrak{so}(3)$. Hence the identity component of $G$ is the connected subgroup with Lie algebra $\mathfrak{so}(3)$, that is, $\SO(3)$, and $-I\in\Oh$ gives $G=\Ot$.
(ii) The closure is a closed subgroup that contains $\Oh$ and the infinite set $\{R^n:n\in\Z\}$; apply (i).
(iii) $\mathrm{Stab}(S)$ is closed (Lemma~\ref{lem:D:cont}), contains $\Oh$ (Theorem~\ref{thm:C-OS1}(b)) and $R$; apply (ii).
\end{proof}

\section{OS4: ergodic decomposition and clustering}\label{sec:OS4}

This section proves Theorems~\ref{thm:cluster} and~\ref{thm:os4main}. The argument for Theorem~\ref{thm:cluster} has two layers. The first (Sections~\ref{ss:abstract} and~\ref{ss:proj}--\ref{ss:axis}) is abstract: it uses a probability law $\mu$ on a Polish space of fields with three properties, (H1)--(H3) below, namely translation and coordinate-reflection invariance, a local structure, and reflection positivity about the coordinate planes. The cubic rotations are not used, and no input about the infrared behaviour of the model is used. The second layer (Section~\ref{ss:discharge}) shows that the laws of the limit states of Theorem~\ref{thm:main} have these three properties. Section~\ref{ss:os4axis} uses, in addition, estimates on the lattice tori, and proves single-axis clustering for part of the observables.

\begin{definition}[forms of clustering]\label{def:ladder}
Let $\mu$ be a translation-invariant probability law of a field and $U(a)$ the induced translations on $L^2(\mu)$. We distinguish:
(C1) \emph{clustering in the mean}: Ces\`aro averages of $\E[G'\,U(a)G]$ over cubes converge to $\E[G']\E[G]$;
(C2) \emph{clustering under diagonal separation}: $\E[G'\,U(a)G]\to\E[G']\E[G]$ as $\min_j\abs{a_j}\to\infty$;
(C3) clustering at a power-law rate;
(C4) clustering at an exponential rate (the Euclidean form of a mass gap).
\end{definition}
We prove (C1) and (C2) on ergodic components, and nothing about (C3) and (C4). Clustering as the separation runs along a single coordinate axis is the cluster property OS4 of Definition~\ref{def:OS4}; for laws with the properties (H1)--(H3) below it implies (C2) (Lemma~\ref{lem:os4:red}), and Section~\ref{ss:os4axis} proves it for part of the observables. An exponential rate is not expected in the symmetric phase, where the photon is expected to be massless (Section~\ref{sec:disc}).

\subsection{The abstract setting}\label{ss:abstract}

\emph{Notation.} $e_1,e_2,e_3$ is the standard basis of $\R^3$. For $J\subset\{1,2,3\}$ put $V_J=\operatorname{span}\{e_j:j\in J\}$ and $J^c=\{1,2,3\}\setminus J$. For $j\in\{1,2,3\}$ and $c\in\R$ the open half spaces are $\Hp{j}{c}=\{x:x_j>c\}$ and $\Hm{j}{c}=\{x:x_j<c\}$. The reflection in the plane $\{x_j=0\}$ is $\rho_j x=x-2x_je_j$, and $\rho_{j,c}$ is the reflection in $\{x_j=c\}$. A bounded open set is denoted $K$; $K\Subset U$ means that $\overline K\subset U$ is compact.

\begin{enumerate}
\item[(H1)] (\emph{Space, actions, measure.}) $X$ is a Polish space with a jointly continuous action $(a,x)\mapsto\tau_ax$ of $\R^3$ and three involutive homeomorphisms $\theta_1,\theta_2,\theta_3$ of $X$ with $\theta_j\tau_a=\tau_{\rho_ja}\theta_j$. $\mu$ is a Borel probability measure on $X$ invariant under every $\tau_a$ and every $\theta_j$.
\item[(H2)] (\emph{Local structure.}) To every open $U\subset\R^3$ is assigned a sub-$\sigma$-algebra $\Sig(U)$ of the Borel $\sigma$-algebra $\cB$ of $X$, such that: $\Sig(U)\subset\Sig(U')$ if $U\subset U'$; $\Sig(U)=\bigvee_{K\Subset U}\Sig(K)$; $\bigvee_{K}\Sig(K)=\cB$; each $\Sig(U)$ is generated by a countable family of bounded continuous functions on $X$; and, with
\begin{equation}\label{eq:D:cov}
U(a)G:=G\circ\tau_{-a},\qquad \Theta_jG:=G\circ\theta_j,
\end{equation}
$U(a)$ maps $\Sig(U)$-measurable functions to $\Sig(U+a)$-measurable ones and $\Theta_j$ maps $\Sig(U)$-measurable functions to $\Sig(\rho_jU)$-measurable ones.
\item[(H3)] (\emph{Coordinate reflection positivity.}) For every $j$, every $c\in\R$ and every $F\in L^2(\Sig(\Hp{j}{c}),\mu)$,
\begin{equation}\label{eq:D:RP}
\E_\mu\big[\,\overline{F\circ\theta_{j,c}}\;F\,\big]\ \ge\ 0,\qquad \theta_{j,c}:=\tau_{ce_j}\theta_j\tau_{-ce_j}.
\end{equation}
\end{enumerate}

By (H1), $a\mapsto U(a)$ is a unitary representation of $\R^3$ on $L^2(\mu)$, $\Theta_j$ is a unitary involution, and
\begin{equation}\label{eq:D:commute}
\Theta_jU(a)=U(\rho_ja)\Theta_j;\quad\text{in particular } \Theta_jU(te_j)=U(-te_j)\Theta_j,\ \ \Theta_jU(y)=U(y)\Theta_j\ (y\perp e_j).
\end{equation}
By translation invariance, \eqref{eq:D:RP} for $c=0$ implies it for every $c$: if $F\in L^2(\Sig(\Hp{j}{c}))$ then $U(-ce_j)F\in L^2(\Sig(\Hp{j}{0}))$ and $\E[\overline{F\circ\theta_{j,c}}\,F]=\E[\overline{(U(-ce_j)F)\circ\theta_j}\;U(-ce_j)F]$.

\begin{remark}[conventions]\label{rem:D:conv}
For a space of distributions we use $(\tau_a\Phi)(f):=\Phi(f(\cdot+a))$, the push-forward of $\Phi$ by $x\mapsto x+a$; then $U(a)[\Phi(f)]=\Phi(f(\cdot-a))$, whose test function is supported in $\supp f+a$, as (H2) requires.
\end{remark}

\subsection{The limit states satisfy (H1)--(H3)}\label{ss:discharge}

Lemma~\ref{lem:D:realize} realises each limit state as a probability law on a space $X$ of distributions, and Proposition~\ref{prop:discharge} verifies (H1)--(H3) there. Lemma~\ref{lem:D:hermite} carries the result to the Hermite--Sobolev space on which Theorem~\ref{thm:cluster} is stated, and proves that the law is the only one with the moments $S$.

Fix one of the limit states $S\in\{S^{\mu_0}\ (\mu_0\in Q_0),\,S^\fl,\,S^0\}$ of Theorem~\ref{thm:main}. Let $\cV:=\Lambda^2\R^3\oplus\R\cong\R^4$. A test function is $f=(\bg,h)\in C_c^\infty(\R^3;\cV)$, and the field is $\Phi(f):=F(\bg)+\XR(h)$. We write $G_S(\zeta)$ for the generating functional of $S$ on the tube (Section~\ref{sec:C-limits}), $S_{n,p}$ for its Schwinger functions, and $c_S:=e^{C\sigma_0}$ with $C$ and $\sigma_0$ the constants of Theorem~\ref{thm:main}(iii).

\emph{Actions on test functions.} For $a\in\R^3$, $(T_af)(x):=f(x-a)$. For $g\in\Ot$, $g_*f:=((g^{-1})^*\bg,\ h\circ g^{-1})$, the push-forward of a 2-form and of a scalar. For a coordinate reflection, $(\rho_{j*}\bg)_{\mu\nu}(x)=(-1)^{[\mu=j]+[\nu=j]}\,g_{\mu\nu}(\rho_jx)$: the components carrying the index $j$ change sign. This is the continuum form of the gauge-covariant lattice reflection used for OS2 (Theorem~\ref{thm:C-OS2}).

\emph{The space.} Fix $s>3$ and let $X:=H^{-s}_{\mathrm{loc}}(\R^3;\cV)$, the space of distributions $\Phi$ with $\chi\Phi\in H^{-s}(\R^3;\cV)$ for every $\chi\in C_c^\infty(\R^3)$, with the seminorms $\Phi\mapsto\norm{\chi_R\Phi}_{H^{-s}}$, $R\in\N$, for a fixed family of cut-offs $\chi_R$ equal to $1$ on the ball $B_R$. It is a separable Fr\'echet space, hence Polish. The actions are
\[
(\tau_a\Phi)(f):=\Phi(T_{-a}f),\qquad (g\Phi)(f):=\Phi(g^{-1}_*f)\ (g\in\Oh),\qquad (\theta_j\Phi)(f):=\Phi(\rho_{j*}f).
\]
For an open $U\subset\R^3$ let $\Sig(U)$ be the $\sigma$-algebra generated by the maps $\Phi\mapsto\Phi(f)$, $f\in C_c^\infty(U;\cV)$ real.

\begin{lemma}[realisation]\label{lem:D:realize}
For each $S$ as above there is exactly one Borel probability measure $\mu_S$ on $X$ such that
\begin{equation}\label{eq:D:charfun}
\E_{\mu_S}\big[e^{i\Phi(f)}\big]=G_S(if)\qquad\text{for every real }f\in C_c^\infty(\R^3;\cV),
\end{equation}
where $G_S(if)$ denotes the tube functional of the one-element family $\{f\}$ at the purely imaginary point. Its mixed moments are the Schwinger functions of $S$. For $S=S^{\mu_0}$ and $S=S^\flat$, $\mu_S$ is the weak limit on $X$ of the lattice laws along the subsequence of Theorem~\ref{thm:main}; for $S=S^0$ it is the weak limit of $\mu_{S^{\mu}}$ as $\mu\downarrow0$ in $Q_0$.
\end{lemma}
\begin{proof}
\emph{Uniqueness.} Let $(f_n)$ be a countable subset of real $C_c^\infty(\R^3;\cV)$, dense in it for the inductive-limit topology. The map $\Phi\mapsto(\Phi(f_n))_n$ from $X$ to $\R^{\N}$ is continuous and injective, so by the Lusin--Souslin theorem (Lemma~\ref{lem:A:LS}) the Borel $\sigma$-algebra of $X$ is generated by the evaluations $\Phi\mapsto\Phi(f_n)$. The sets $\{\Phi:(\Phi(f_1),\dots,\Phi(f_N))\in B\}$ form a $\pi$-system generating $\cB$, and the law of each vector $(\Phi(f_1),\dots,\Phi(f_N))$ is determined by its characteristic function $t\mapsto\E\exp(i\sum_kt_k\Phi(f_k))=G_S(i\sum_kt_kf_k)$ (linearity of $\Phi$). Hence \eqref{eq:D:charfun} determines $\mu_S$.

\emph{Existence for $S^{\mu_0}$, $S^\flat$.} The lattice laws are tight on $X$ along the subsequence, and every weak limit point has the Schwinger functions as moments and $G_S(i\,\cdot)$ as characteristic functional (Proposition~\ref{prop:C-tight}; the imaginary points lie in the tube, whose condition bounds only real parts of the coefficients of the Wick square). By uniqueness all limit points coincide, so the whole subsequence converges.

\emph{Existence for $S^0$.} By Theorem~\ref{thm:C-OS0}(b), whose constants are the same for $S^{\mu_0}$, $S^\flat$ and $S^0$, the bound of the proof of Proposition~\ref{prop:C-tight}, $\E_{\mu_{S^\mu}}\abs{\Phi(\chi_Re_\alpha)}^2\le C'_R$ for the Fourier basis $e_\alpha$ of a box containing $B_{R+1}$, holds with $C'_R$ independent of $\mu\in Q_0$. The Hilbert--Schmidt argument of Proposition~\ref{prop:C-tight} gives tightness of $\{\mu_{S^\mu}\}_{\mu\in Q_0}$ on $X$. Let $\mu'$ be a weak limit point as $\mu\downarrow0$. The functions $\Phi\mapsto e^{i\Phi(f)}$ are bounded and continuous on $X$, so $\E_{\mu'}e^{i\Phi(f)}=\lim G_{S^\mu}(if)=G_{S^0}(if)$ by the locally uniform convergence on the tube (Proposition~\ref{prop:C-mu0}(b)). By uniqueness $\mu'=\mu_{S^0}$, and the whole family converges. Moments converge by uniform integrability, since all moments are bounded uniformly in $\mu$.
\end{proof}

\begin{proposition}[(H1)--(H3) for the limit states]\label{prop:discharge}
For every $S$ as above, $(X,\tau,\theta,\mu_S,\Sig)$ satisfies (H1), (H2) and (H3). Moreover $\mu_S$ is invariant under the action of $\Oh$ about every point, and $\mu_S$ is the only Borel probability measure on $X$ whose mixed moments on $C_c^\infty(\R^3;\cV)$ are the Schwinger functions of $S$.
\end{proposition}
\begin{proof}
\emph{(H1).} $X$ is Polish. The map $(a,\Phi)\mapsto\tau_a\Phi$ is jointly continuous on $\R^3\times X$: translations act on $H^{-s}_{\mathrm{loc}}$ by a family of operators that is equicontinuous for $a$ in compacts and strongly continuous, and these two properties give joint continuity. Each $\theta_j$ is a linear homeomorphism of $X$ with $\theta_j^2=1$, and $\theta_j\tau_a=\tau_{\rho_ja}\theta_j$ because $\rho_j$ is linear. Invariance of $\mu_S$: by \eqref{eq:D:charfun} it suffices that $G_S(iT_af)=G_S(if)$ and $G_S(i\rho_{j*}f)=G_S(if)$. By Theorem~\ref{thm:C-OS1} every Schwinger function of $S$ is invariant under translations and under $\Oh$ (which contains $\rho_j$), and the tube functional is determined on the connected tube by its Taylor coefficients at $0$ (Proposition~\ref{prop:C-compact}(b)), so $G_S$ is invariant; the same argument gives invariance under all of $\Oh$ about every point.

\emph{(H2).} Monotonicity in $U$ is clear, and $\Sig(U)=\bigvee_{K\Subset U}\Sig(K)$ because each $f\in C_c^\infty(U)$ has support in some $K\Subset U$. The Lusin--Souslin argument in the uniqueness part of Lemma~\ref{lem:D:realize} gives $\bigvee_K\Sig(K)=\cB$. $\Sig(U)$ is generated by the countably many bounded continuous functions $\arctan\Phi(f_n)$, where $(f_n)$ is a countable subset of $C_c^\infty(U;\cV)$ that is dense in each $H^s$-norm on each compact subset of $U$; indeed $\Phi(f)$ is the pointwise limit of $\Phi(f_{n_k})$ whenever $f_{n_k}\to f$ in $H^s$ with supports in a fixed compact. The covariance: $U(a)[\Phi(f)]=\Phi(T_af)$ and $\supp T_af=\supp f+a$; $\Theta_j[\Phi(f)]=\Phi(\rho_{j*}f)$ and $\supp\rho_{j*}f=\rho_j\supp f$.

\emph{(H3).} Let $c=0$ (translation invariance, as noted after \eqref{eq:D:commute}). Let $\mathcal E_+$ be the linear span of the functions $e^{i\Phi(f)}$, $f$ real in $C_c^\infty(\Hp{j}{0};\cV)$. For $E_k=e^{i\Phi(f_k)}$ one has $\overline{E_k\circ\theta_j}=e^{-i\Phi(\rho_{j*}f_k)}$, which is the antilinear reflection $\theta(E_k)$ of Theorem~\ref{thm:C-OS2}: the coefficient $i$ is conjugated and the test function is reflected, with the sign of the $j$-components of the 2-form part. Theorem~\ref{thm:C-OS2}(a) (matrix reflection positivity of tube exponentials) applies to these $E_k$: their curvature coefficients are arbitrary complex numbers, and their Wick-square coefficients are purely imaginary, so the amplitude condition, which bounds only real parts, holds. Its matrix entries are expectations under $\mu_S$: by \eqref{eq:D:charfun} applied to $f_l-\rho_{j*}f_k$, $\E_{\mu_S}[\overline{E_k\circ\theta_j}\,E_l]=G_S(i(f_l-\rho_{j*}f_k))$, which is the entry $S(\theta_j(E_k)E_l)$ of the matrix of Theorem~\ref{thm:C-OS2}(a), both being limits of the same lattice expectations (Propositions~\ref{prop:C-compact}(b) and~\ref{prop:C-mu0}(b)). Hence $\E_{\mu_S}[\overline{F\circ\theta_j}\,F]\ge0$ for every $F\in\mathcal E_+$. $\mathcal E_+$ is an algebra of bounded continuous functions, closed under complex conjugation, containing the constants and generating $\Sig(\Hp{j}{0})$, so it is dense in $L^2(\Sig(\Hp{j}{0}),\mu_S)$ (Lemma~\ref{lem:A:mc}). If $F_n\to F$ in $L^2$ with $F_n\in\mathcal E_+$, then $\overline{F_n\circ\theta_j}F_n\to\overline{F\circ\theta_j}F$ in $L^1(\mu_S)$, because $\theta_j$ preserves $\mu_S$. This gives \eqref{eq:D:RP}.

\emph{Determination by $S$.} By Lemma~\ref{lem:D:realize}, $\mu_S$ is determined by $G_S(i\,\cdot)$, which is determined by the Schwinger functions (Proposition~\ref{prop:C-compact}(b)). If $\mu'$ is any Borel probability measure on $X$ with the same mixed moments, the argument of Lemma~\ref{lem:D:hermite}(c) applies verbatim on $X$, with the $\pi$-system of the proof of Lemma~\ref{lem:D:realize}, and gives $\mu'=\mu_S$.
\end{proof}
\noindent The field space is $\cV=\Lambda^2\R^3\oplus\R$: the three components of the curvature and the renormalised Wick square. Nothing above uses more of the field content than the transformation rules of the two summands.

\emph{The Hermite--Sobolev space.} $\mathcal A:=-\Delta+\abs x^2+1$ acts componentwise on $L^2(\R^3;\cV)$. Its eigenfunctions are $h_{\alpha,c}:=h_\alpha e_c$, with $h_\alpha$ the Hermite functions, $\alpha\in(\N\cup\{0\})^3$, and $(e_c)_{c\le4}$ a basis of $\cV$. The eigenvalues are $\lambda_\alpha=2\abs\alpha+4$. For $r\in2\N$ put $\norm f_r:=\norm{\mathcal A^{r/2}f}_{L^2}$, let $\mathcal H^r$ be the completion of $\cS(\R^3;\cV)$ in this norm, and let $\mathcal H^{-r}$ be its dual. So $\norm\Psi_{-r}^2=\sum_{\alpha,c}\lambda_\alpha^{-r}\abs{\Psi(h_{\alpha,c})}^2$, and $\mathcal H^{-r}$ is a separable Hilbert space. The weighted norm of OS0 is $\abs f_4:=\sup_{x,\abs\gamma\le4}(1+\abs x^2)^2\abs{D^\gamma f(x)}$.

\begin{lemma}[Hermite realisation and determinacy]\label{lem:D:hermite}
Let $S$ be as above, and let $\mu_S$ be the measure of Lemma~\ref{lem:D:realize} on $X=H^{-s}_{\mathrm{loc}}(\R^3;\cV)$, with $s:=r$.
\begin{enumerate}
\item[(a)] For $k\in2\N$ and $\abs\delta+\abs\gamma\le k$, $\norm{x^\delta\partial^\gamma f}_{L^2}\le k^{k/2}\norm f_k$. Hence $\abs f_4\le C\norm f_{10}$ for $f\in\cS$, and on $\cS$ the norm $\norm\cdot_{k}$ is equivalent to $\sum_{\abs\delta+\abs\gamma\le k}\norm{x^\delta\partial^\gamma\cdot}_{L^2}$.
\item[(b)] Let $r\ge14$ be even. There is exactly one Borel probability measure $\mu^H_S$ on $\mathcal H^{-r}$ whose image under the restriction map $\iota:\mathcal H^{-r}\to X$, $\iota\Psi:=\Psi|_{C_c^\infty}$, is $\mu_S$. The map $\iota$ is continuous, linear and injective, and $\mu_S(\iota\mathcal H^{-r})=1$. Moreover $\E_{\mu^H_S}e^{i\Psi(f)}=G_S(if)$, and $\E_{\mu^H_S}\prod_l\Psi(f_l)=S(f_1,\dots,f_N)$ for $f_l\in C_c^\infty$.
\item[(c)] $\mu^H_S$ is the only Borel probability measure on $\mathcal H^{-r}$ whose mixed moments on $C_c^\infty(\R^3;\cV)$ exist and equal those of $S$.
\item[(d)] Define the translations, the coordinate reflections and $\Oh$ on $\mathcal H^{-r}$ by the formulas used on $X$, and let $\Sig^H(U):=\iota^{-1}\Sig(U)$, the $\sigma$-algebra generated by $\Psi\mapsto\Psi(f)$, $f\in C_c^\infty(U;\cV)$. If $(X,\tau,\theta,\mu_S,\Sig)$ satisfies (H1)--(H3), so does $(\mathcal H^{-r},\tau,\theta,\mu^H_S,\Sig^H)$, and $\mu^H_S$ is $\Oh$-invariant about every point.
\end{enumerate}
\end{lemma}
\begin{proof}
(a) In each coordinate, $x_j=(a_j+a_j^*)/\sqrt2$ and $\partial_j=(a_j-a_j^*)/\sqrt2$, with $a_jh_\alpha=\sqrt{\alpha_j}\,h_{\alpha-e_j}$ and $a_j^*h_\alpha=\sqrt{\alpha_j+1}\,h_{\alpha+e_j}$. So $x^\delta\partial^\gamma$ is a sum of $2^{\abs\delta+\abs\gamma}$ words of length $\ell:=\abs\delta+\abs\gamma$ in the $a_j^\#$, each with coefficient $2^{-\ell/2}$. A word $w$ shifts $\alpha$ by a fixed vector, so it maps distinct basis vectors to distinct basis vectors or to $0$. Its coefficients are at most $(\abs\alpha+\ell)^{\ell/2}\le(\ell\lambda_\alpha/2)^{\ell/2}$. Hence $\norm{wf}^2\le(\ell/2)^\ell\sum_\alpha\lambda_\alpha^\ell\abs{\ip{h_\alpha}f}^2\le(\ell/2)^{\ell}\norm f_k^2$ for $\ell\le k$, since $\lambda_\alpha\ge1$, and $\norm{x^\delta\partial^\gamma f}\le2^{\ell/2}(\ell/2)^{\ell/2}\norm f_k\le k^{k/2}\norm f_k$. For the second claim, put $u:=(1+\abs x^2)^2D^\gamma f$ with $\abs\gamma\le4$. Sobolev's embedding in $\R^3$ gives $\norm u_\infty\le C\sum_{\abs\beta\le2}\norm{\partial^\beta u}_{L^2}$. Each $\partial^\beta u$ is a finite sum of terms $x^\delta\partial^{\gamma'}f$ with $\abs\delta\le4$ and $\abs{\gamma'}\le6$. For the equivalence, $\mathcal A^{k/2}$ is a differential operator with polynomial coefficients of total degree $\le k$, which gives the reverse inequality.

(b) \emph{A two-point bound on $\cS$.} For $f=(\bg,h)\in C_c^\infty$, Theorem~\ref{thm:main}(iii) and OS3 ($S_{1,1}=0$) give $\E_{\mu_S}\Phi(f)^2=S_{2,0}(\bg,\bg)+S_{0,2}(h,h)\le c_2(\norm\bg_2^2+N(h)^2)$. Moreover $\norm\bg_2+N(h)\le C\abs f_4$, since $\int(1+\abs x^2)^{-2}dx=\pi^2$ and $\int(1+\abs x^2)^{-4}dx=\pi^2/8$. With (a), $\E_{\mu_S}\Phi(f)^2\le C_2\norm f_{10}^2$. For $f\in\cS$ and cut-offs $\chi_R$ ($\chi_R=1$ on $B_R$, derivatives bounded uniformly in $R$), $\chi_Rf\to f$ in every seminorm $\norm{x^\delta\partial^\gamma\cdot}_{L^2}$, hence in $\norm\cdot_{10}$. So $\Phi(\chi_Rf)$ is Cauchy in $L^2(\mu_S)$. Call its limit $\tilde\Phi(f)$. It is linear in $f$ and $\E\tilde\Phi(f)^2\le C_2\norm f_{10}^2$.

\emph{The carrier.} Fix Borel versions $Z_{\alpha,c}$ of $\tilde\Phi(h_{\alpha,c})$, and put $M:=\sum_{\alpha,c}\lambda_\alpha^{-r}Z_{\alpha,c}^2$. Then $\E M\le4C_2\sum_\alpha\lambda_\alpha^{10-r}<\infty$, since there are $(n+1)(n+2)/2$ multi-indices with $\abs\alpha=n$ and $r-10>3$. On the Borel set $\{M<\infty\}$, which has full measure, define $\Psi(f):=\sum_{\alpha,c}\ip{h_{\alpha,c}}{f}Z_{\alpha,c}$. By the Cauchy--Schwarz inequality this is a continuous linear functional on $\mathcal H^r$ with $\norm\Psi_{-r}^2=M$. The map $\Phi\mapsto\Psi$ is Borel, because its coordinates are. For $f\in C_c^\infty$ the partial sums of $\Psi(f)$ converge to $\tilde\Phi(f)=\Phi(f)$ in $L^2(\mu_S)$, since the Hermite expansion of $f$ converges in $\norm\cdot_{10}$. So $\Psi(f)=\Phi(f)$ a.s. For each $n$ let $\mathcal D_n$ be a countable subset of $C_c^\infty(B_n;\cV)$ that is dense in the Sobolev norm $H^r$. Almost surely $\Psi(f)=\Phi(f)$ for all $f\in\bigcup_n\mathcal D_n$. Both sides are continuous on $C_c^\infty(B_n)$ for the $H^r$-norm. For $\Phi$ this holds because $\abs{\Phi(f)}\le\norm{\chi_{n+1}\Phi}_{H^{-s}}\norm f_{H^s}$ and $s=r$. For $\Psi$ it holds because $\norm{f}_r\le C_n\norm f_{H^r}$ for $\supp f\subset B_n$, by (a) and the boundedness of $x^\delta$ on $B_n$. Hence $\iota\Psi=\Phi$ a.s.

\emph{The map $\iota$.} For $\chi\in C_c^\infty$, $\abs{\Psi(\chi f)}\le\norm\Psi_{-r}\norm{\chi f}_r\le C_\chi\norm\Psi_{-r}\norm f_{H^r}$, so $\chi\,\iota\Psi\in H^{-r}=H^{-s}$, and $\iota$ is continuous into $X$. It is injective, because $C_c^\infty$ is dense in $\cS$ and $\cS$ is dense in $\mathcal H^r$. Put $\mu^H_S:=\Psi_*\mu_S$. Then $\iota_*\mu^H_S=\mu_S$. By Lemma~\ref{lem:A:LS} (applied to $\iota$ composed with countably many evaluations), $\iota$ is a Borel isomorphism onto the Borel set $\iota\mathcal H^{-r}$, which has full $\mu_S$-measure. This gives uniqueness of $\mu^H_S$ and the formulas for the characteristic functional and the moments, which hold for $\mu_S$ by Theorem~\ref{thm:main}(ii).

(c) Let $\mu'$ be such a measure, and fix real $f_1,\dots,f_N\in C_c^\infty$ and $t\in\R^N$. Put $f:=\sum_lt_lf_l=(\bg,h)$. The moments of the real random variable $Z:=\Psi(f)=\sum_lt_l\Psi(f_l)$ under $\mu'$ are $m_n=\sum_{k\le n}\binom nkS_{k,n-k}(\bg,\dots;h,\dots)$, the same as under $\mu^H_S$. By Theorem~\ref{thm:main}(iii) and $\binom nk(k!)^{1/2}(n-k)!\le n!$, they satisfy $\abs{m_n}\le c_S(n+1)K^n\,n!$ with $K:=\max(\mathrm e\norm\bg_2,\mathrm eN(h)/\sigma_0)$. The even moments are absolute moments, and the Cauchy--Schwarz inequality $\E\abs Z^{2n+1}\le(m_{2n}m_{2n+2})^{1/2}$ gives $\E\abs Z^n\le c_S(n+2)K^n\,n!$ for all $n$, under $\mu'$ and under $\mu^H_S$. Let $\phi(s):=\E e^{isZ}$. Then $\phi^{(k)}(s)=\E[(iZ)^ke^{isZ}]$, and from $\abs{e^{iy}-\sum_{k<n}(iy)^k/k!}\le\abs y^n/n!$,
\[
\Big\lvert\phi(s+u)-\sum_{k<n}\phi^{(k)}(s)\frac{u^k}{k!}\Big\rvert\le\frac{\abs u^n\,\E\abs Z^n}{n!}\le c_S(n+2)(K\abs u)^n\to0\qquad(\abs u<1/K).
\]
So on every interval $(s-\rho,s+\rho)$, $\rho<1/K$, $\phi$ is the sum of its Taylor series at $s$. The characteristic functions of $Z$ under $\mu'$ and $\mu^H_S$ have the same derivatives $i^km_k$ at $0$, so they agree on $(-\rho,\rho)$; there all their derivatives agree, so they agree near $\pm\rho/2$, and by induction on $\R$. Hence $\E_{\mu'}e^{i\sum_lt_l\Psi(f_l)}=\E_{\mu^H_S}e^{i\sum_lt_l\Psi(f_l)}$ for all $t\in\R^N$, and by uniqueness of characteristic functions on $\R^N$ the laws of $(\Psi(f_1),\dots,\Psi(f_N))$ agree. These vectors, over all finite families, generate a $\pi$-system that generates the Borel $\sigma$-algebra of $\mathcal H^{-r}$ (Lemma~\ref{lem:A:LS}; the evaluations at $C_c^\infty$ separate points). Hence $\mu'=\mu^H_S$.

(d) \emph{(H1).} $\mathcal H^{-r}$ is a separable Hilbert space, hence Polish. By (a), with $\norm\cdot_r$ replaced by the equivalent polynomial norm, $\norm{f(\cdot-a)}_r\le C_r(1+\abs a)^r\norm f_r$. So $\tau_a$, the adjoint of $f\mapsto f(\cdot+a)$, is bounded on $\mathcal H^{-r}$ with $\norm{\tau_a}\le C_r(1+\abs a)^r$. On the dense subspace $\cS\subset\mathcal H^{-r}$, embedded by the $L^2$ pairing, $\norm{\tau_a\psi-\psi}_{-r}\le\norm{\psi(\cdot-a)-\psi}_{L^2}\to0$. Local boundedness and an $\epsilon/3$ argument give strong continuity. Joint continuity follows from $\norm{\tau_a\Psi-\tau_b\Psi'}_{-r}\le C_r(1+\abs a)^r\norm{\Psi-\Psi'}_{-r}+\norm{\tau_a\Psi'-\tau_b\Psi'}_{-r}$. Every $g\in O(3)$ commutes with $\mathcal A$ and acts on $\cV$ by an orthogonal matrix, so $\theta_j$ and $\Oh$ act unitarily. The relations $\theta_j\tau_a=\tau_{\rho_ja}\theta_j$ are identities between the formulas on test functions. The actions commute with $\iota$, and $\iota\mathcal H^{-r}$ is invariant. Since $\iota_*\mu^H_S=\mu_S$ and $\iota$ is a Borel isomorphism onto its image, every invariance of $\mu_S$ is an invariance of $\mu^H_S$. \emph{(H2).} $\Sig^H(U)$ is generated by the bounded continuous functions $\arctan\Psi(f_n)$, $f_n$ in a countable subset of $C_c^\infty(U)$ dense in every $C_c^\infty(K)$, $K\Subset U$, for the $H^r$-norm. The covariance and monotonicity properties are those of $\Sig$, and $\bigvee_K\Sig^H(K)$ is the Borel $\sigma$-algebra by Lemma~\ref{lem:A:LS}. \emph{(H3).} A $\Sig^H(\Hp j0)$-measurable $F$ has the form $F=\phi((\Psi(f_n))_n)$ with $f_n\in C_c^\infty(\Hp j0)$ and $\phi$ Borel. So $F=F'\circ\iota$ with $F':=\phi((\Phi(f_n))_n)$, which is $\Sig(\Hp j0)$-measurable on $X$. Then $\E_{\mu^H_S}[\overline{F\circ\theta_j}F]=\E_{\mu_S}[\overline{F'\circ\theta_j}F']\ge0$.
\end{proof}

\noindent By Proposition~\ref{prop:discharge}, the hypothesis of part (d) holds. Hence Theorem~\ref{thm:cluster} may be stated, as it is, on $\mathcal H^{-r}$ for every even $r\ge14$, with $\mu=\mu^H_S$; its uniqueness clause is part (c). The abstract results of Sections~\ref{ss:proj}--\ref{ss:os4axis} apply on either space, and the two decompositions correspond under $\iota$.

\subsection{Translations and the mean-ergodic projections}\label{ss:proj}

Throughout Sections~\ref{ss:proj}--\ref{ss:axis}, $\mu$ is any measure satisfying (H1)--(H3).

\begin{lemma}\label{lem:D:strongcont}
$a\mapsto U(a)$ is strongly continuous on $L^2(\mu)$.
\end{lemma}
\begin{proof}
For bounded continuous $G$, $G(\tau_{-a}x)\to G(x)$ for every $x$ as $a\to0$ by continuity of the action, and dominated convergence gives $U(a)G\to G$ in $L^2$. Bounded continuous functions are dense in $L^2(\mu)$ for a Borel probability measure on a metric space, and $\norm{U(a)}=1$.
\end{proof}

For $J\subset\{1,2,3\}$ let $P_J$ be the orthogonal projection of $L^2(\mu)$ onto $\{G:U(a)G=G\ \forall a\in V_J\}$, with $P_\emptyset=1$ and $P_j:=P_{\{j\}}$. Let $\cI$ be the $\sigma$-algebra of $\R^3$-invariant Borel sets, completed by $\mu$-null sets.

\begin{lemma}\label{lem:D:proj}
(i) For $\Lambda^J_T:=[-T,T]^3\cap V_J$ and every $G\in L^2(\mu)$, $\abs{\Lambda^J_T}^{-1}\int_{\Lambda^J_T}U(a)G\,da\to P_JG$ in $L^2(\mu)$ as $T\to\infty$. The same holds for one-sided means over $[0,T]$ in one coordinate of $J$.
(ii) $P_J$ commutes with every $U(a)$ and every $\Theta_j$; $P_JP_{J'}=P_{J\cup J'}$; and $P_J=\E[\,\cdot\mid\cI_J]$, where $\cI_J$ is the $\sigma$-algebra of $V_J$-invariant Borel sets completed by $\mu$-null sets; in particular $P_{\{1,2,3\}}=\E[\,\cdot\mid\cI]$.
\end{lemma}
\begin{proof}
(i) For one coordinate $j$, Stone's theorem gives $U(te_j)=e^{itA_j}$ with $A_j$ self-adjoint, and $(2T)^{-1}\int_{-T}^Te^{itA_j}dt=\phi_T(A_j)$ with $\phi_T(\lambda)=\sin(T\lambda)/(T\lambda)$, $\phi_T(0)=1$. Since $\abs{\phi_T}\le1$ and $\phi_T\to\mathbf 1_{\{0\}}$ pointwise, the spectral theorem and dominated convergence give strong convergence to the spectral projection of $\{0\}$, which is the projection onto $\Ker A_j=\{G:U(te_j)G=G\ \forall t\}$. For one-sided means use $\phi_T(\lambda)=(e^{iT\lambda}-1)/(iT\lambda)$. The mean over $\Lambda^J_T$ is the product of the commuting one-dimensional means; products of commuting, uniformly bounded, strongly convergent families converge strongly to the product of the limits; and the product of commuting projections is the projection onto the intersection of their ranges, which is the $V_J$-fixed space.

(ii) The means commute with $U(a)$, and with $\Theta_j$ by \eqref{eq:D:commute} and $\rho_jV_J=V_J$; strong limits of commuting operators commute. $P_JP_{J'}=P_{J\cup J'}$ by the intersection statement in (i). It remains to show that the range of $P_J$ is $L^2(\cI_J,\mu)$. Let $G\in L^2(\mu)$ satisfy $U(a)G=G$ for all $a\in V_J$, and fix a Borel version of $G$. For each $a\in V_J$ the set $\{x:G(\tau_{-a}x)\neq G(x)\}$ is $\mu$-null, so by Fubini's theorem on $V_J\times X$ the set $X_0:=\{x:G(\tau_{-a}x)=G(x)\text{ for a.e.\ }a\in V_J\}$ is Borel and of full measure. Let $X_1$ be the set of $x$ for which $a\mapsto G(\tau_{-a}x)$ is a.e.\ constant on $V_J$, that is,
\[
\int_{\Lambda^J_n}\int_{\Lambda^J_n}\min\big(1,\abs{G(\tau_{-a}x)-G(\tau_{-b}x)}\big)\,da\,db=0\qquad\text{for all }n\in\N .
\]
By Tonelli's theorem $X_1$ is Borel; $X_0\subset X_1$; and $X_1$ is $V_J$-invariant, because $G(\tau_{-a}\tau_cx)=G(\tau_{-(a-c)}x)$ and a.e.\ constancy on $V_J$ is preserved by shifts of the variable. For $x\in X_1$ let $\tilde G(x)$ be the a.e.\ value of $a\mapsto G(\tau_{-a}x)$, that is, $\tilde G(x)=\lim_{N\to\infty}\abs{\Lambda^J_1}^{-1}\int_{\Lambda^J_1}t_N\big(G(\tau_{-a}x)\big)\,da$, where $t_N$ truncates real and imaginary parts at level $N$; put $\tilde G:=0$ off $X_1$. Then $\tilde G$ is Borel, $\tilde G\circ\tau_c=\tilde G$ everywhere for $c\in V_J$, and $\tilde G=G$ on $X_0$, hence $\mu$-a.e. Thus $G$ is $\cI_J$-measurable. Conversely, an $\cI_J$-measurable function is a.e.\ equal to a $V_J$-invariant Borel function and is therefore fixed by every $U(a)$, $a\in V_J$.
\end{proof}

The next lemma is the only place where reflection positivity enters the decomposition: functions invariant under translations in a direction live at infinity in that direction and are fixed by the corresponding reflection.

\begin{lemma}[localisation of invariant functions]\label{lem:D:invloc}
Let $j\in J\subset\{1,2,3\}$ and $W\in L^2(\mu)$ with $U(a)W=W$ for all $a\in V_J$. Then:
(i) $W\in L^2(\Sig(\Hp{j}{c}),\mu)\cap L^2(\Sig(\Hm{j}{c}),\mu)$ for every $c\in\R$;
(ii) $\Theta_jW=W$;
(iii) if $G\in L^2(\Sig(K),\mu)$, then $P_JG\in L^2(\Sig(K+V_J),\mu)$.
\end{lemma}
\begin{proof}
(i) Given $\epsilon>0$, by (H2) there are a bounded $K$ and $G\in L^2(\Sig(K))$ with $\norm{W-G}_2<\epsilon$ (an increasing union of $\sigma$-algebras generating $\cB$ has $L^2$ spaces whose union is dense). For every $t$, $\norm{W-U(te_j)G}_2=\norm{U(te_j)(W-G)}_2<\epsilon$, and $U(te_j)G\in L^2(\Sig(K+te_j))\subset L^2(\Sig(\Hp{j}{c}))$ for $t$ large. That space is closed, so it contains $W$. Letting $t\to-\infty$ gives the other half space.

(ii) $\Theta_jW$ is $V_J$-invariant by \eqref{eq:D:commute}, so by (i) both $W$ and $\Theta_jW$ lie in $L^2(\Sig(\Hp{j}{0}))$. Put $F=W-\Theta_jW$. Then $\Theta_jF=-F$, and \eqref{eq:D:RP} gives $0\le\E[\overline{\Theta_jF}\,F]=-\E\abs{F}^2$, so $F=0$.

(iii) The means of Lemma~\ref{lem:D:proj}(i) applied to $G$ lie in the closed subspace $L^2(\Sig(K+V_J))$, since $U(a)G\in L^2(\Sig(K+a))$ and $K+a\subset K+V_J$ for $a\in V_J$.
\end{proof}

\subsection{Ergodic decomposition and inheritance}\label{ss:decomp}

\begin{lemma}[conditional probabilities]\label{lem:D:rcp}
Let $X$ be Polish, $m$ a Borel probability measure on $X$ and $\mathcal F_0\subset\cB$ a $\sigma$-algebra. There are Borel probability measures $m_\xi$, $\xi\in X$, such that $\xi\mapsto m_\xi(B)$ is $\mathcal F_0$-measurable and $\int_Cm_\xi(B)\,dm(\xi)=m(B\cap C)$ for all $B\in\cB$, $C\in\mathcal F_0$.
\end{lemma}
\begin{proof}
Let $d\le1$ be a complete metric on $X$ and $(x_n)$ a dense sequence. Then $\jmath(x):=(d(x,x_n))_n$ is a continuous injection into the compact space $Q:=[0,1]^\N$, and by Lemma~\ref{lem:A:LS} it is a Borel isomorphism onto the Borel set $\jmath X$. Let $\mathcal C_0$ be a countable $\Q$-linear dense subspace of $C(Q)$ containing $1$, and choose $\mathcal F_0$-measurable versions $L_\xi(G)$ of $\E_m[G\circ\jmath\mid\mathcal F_0]$, $G\in\mathcal C_0$. Off a set $N_0\in\mathcal F_0$ with $m(N_0)=0$, $L_\xi$ is $\Q$-linear and positive, and $L_\xi(1)=1$. So $\abs{L_\xi(G)}\le\norm G_\infty$, and $L_\xi$ extends to a positive normalised functional on $C(Q)$, given by a probability measure $\kappa_\xi$ (Riesz representation theorem). On $N_0$ put $\kappa_\xi:=\jmath_*m$. The bounded Borel $G$ on $Q$ for which $\kappa_\xi(G)$ is a $\mathcal F_0$-measurable version of $\E_m[G\circ\jmath\mid\mathcal F_0]$ form a vector space that contains $C(Q)$ and is closed under bounded monotone limits, so it contains all bounded Borel $G$. In particular $\kappa_\xi(\jmath X)=1$ for $\xi\notin N_1$, $N_1\in\mathcal F_0$ null. Put $m_\xi:=(\jmath^{-1})_*\kappa_\xi$ off $N_0\cup N_1$ and $m_\xi:=m$ on it.
\end{proof}

\begin{theorem}[ergodic decomposition with inheritance]\label{thm:D:decomp}
Let $\cI_0\subset\cB$ be the $\sigma$-algebra of Borel sets $B$ with $\tau_aB=B$ for all $a\in\R^3$ (its $\mu$-completion is $\cI$), let $\mu_\xi$, $\xi\in X$, be the measures of Lemma~\ref{lem:D:rcp} for $m=\mu$ and $\mathcal F_0=\cI_0$, a regular conditional probability of $\mu$ given $\cI$, and put $(\Xi,\nu):=(X,\mu)$. Then $\xi\mapsto\mu_\xi(B)$ is $\cI$-measurable for every Borel set $B$, and:
\begin{enumerate}
\item[(a)] $\mu=\int_X\mu_\xi\,d\mu(\xi)$; for $G\in L^1(\mu)$ one has $G\in L^1(\mu_\xi)$ and $\E[G\mid\cI](\xi)=\E_{\mu_\xi}[G]$ for $\mu$-a.e.\ $\xi$; in particular $\E_\mu[WG]=\int W(\xi)\,\E_{\mu_\xi}[G]\,d\mu(\xi)$ for bounded $\cI$-measurable $W$;
\item[(b)] for $\nu$-a.e.\ $\xi$, $\mu_\xi$ is invariant under every $\tau_a$ and ergodic for the $\R^3$-action;
\item[(c)] for $\nu$-a.e.\ $\xi$, $\mu_\xi$ is invariant under $\theta_1,\theta_2,\theta_3$;
\item[(d)] for $\nu$-a.e.\ $\xi$, $\mu_\xi$ satisfies (H3);
\item[(e)] every $\mu_\xi$ satisfies (H2) with the same $\Sig(U)$, and every $G\in L^2(\mu)$ lies in $L^2(\mu_\xi)$ for $\nu$-a.e.\ $\xi$.
\end{enumerate}
\end{theorem}
\begin{proof}
(a) By Lemma~\ref{lem:D:rcp}, (a) holds for bounded $G$, and for $G\in L^1(\mu)$ by monotone convergence. Fix a complete metric $d\le1$ on $X$ and a dense sequence $(x_n)$, and let $\cD$ be the countable $\Q[i]$-algebra generated by $1$ and the functions $x\mapsto d(x,x_n)$. It consists of bounded continuous functions, generates $\cB$, and is dense in $L^2(m)$ for every Borel probability measure $m$ (Lemma~\ref{lem:A:mc}); so two Borel probability measures that agree on $\cD$ are equal.

(b) \emph{Invariance.} For $B\in\cI_0$, $\E_\mu[\mathbf 1_B\,G\circ\tau_a]=\E_\mu[(\mathbf 1_BG)\circ\tau_a]=\E_\mu[\mathbf 1_BG]$. Hence there is a null set off which $\mu_\xi(G\circ\tau_a)=\mu_\xi(G)$ for all $G\in\cD$ and $a\in\Q^3$. Since $a\mapsto\mu_\xi(G\circ\tau_a)$ is continuous (joint continuity of the action and dominated convergence), this holds for all $a$, so $\mu_\xi$ is $\tau$-invariant. \emph{Ergodicity.} For $G\in\cD$ let $q_G(\xi):=\mu_\xi(G)$, a bounded $\cI_0$-measurable version of $\E_\mu[G\mid\cI]$. Then $\mu_\xi(q_G)=q_G(\xi)$ and $\mu_\xi(\abs{q_G}^2)=\abs{q_G(\xi)}^2$ for a.e.\ $\xi$, so $\mu_\xi(\abs{q_G-q_G(\xi)}^2)=0$: for a.e.\ $\xi$, $q_G$ is $\mu_\xi$-a.s.\ the constant $\mu_\xi(G)$. Let $A_TG(x):=\abs{\Lambda_T}^{-1}\int_{\Lambda_T}G(\tau_{-a}x)\,da$. By Lemma~\ref{lem:D:proj}, $\int\mu_\xi(\abs{A_TG-q_G}^2)\,d\mu(\xi)=\E_\mu\abs{A_TG-q_G}^2\to0$. Choose $T_k\to\infty$ such that this is $\le2^{-k}$ for the first $k$ elements of an enumeration of $\cD$. Then for a.e.\ $\xi$, $A_{T_k}G\to\mu_\xi(G)$ in $L^2(\mu_\xi)$ for every $G\in\cD$. The proof of Lemma~\ref{lem:D:proj}(i) uses only the invariance of the measure, so $A_{T_k}G$ also tends in $L^2(\mu_\xi)$ to the projection of $G$ onto the $\tau$-fixed vectors. By density of $\cD$, that projection is $G\mapsto\mu_\xi(G)\mathbf 1$. If $B\in\cI_0$, $\mathbf 1_B$ is fixed, hence a.s.\ constant, so $\mu_\xi(B)\in\{0,1\}$.

(e) (H2) is a condition on $X$ and $\Sig$ alone and does not involve the measure, so it holds for every $\mu_\xi$. Also $\int\E_{\mu_\xi}\abs{G}^2\,d\nu=\E_\mu\abs{G}^2<\infty$.

(c) Let $W\ge0$ be bounded and $\cI$-measurable and $G$ bounded Borel. By Lemma~\ref{lem:D:invloc}(ii), $\Theta_jW=W$, so $\E_\mu[W\,\Theta_jG]=\E_\mu[\Theta_j(WG)]=\E_\mu[WG]$ by $\theta_j$-invariance of $\mu$. By (a), $\int W(\xi)\big(\E_{\mu_\xi}[\Theta_jG]-\E_{\mu_\xi}[G]\big)d\nu=0$ for all such $W$, hence $\E_{\mu_\xi}[\Theta_jG]=\E_{\mu_\xi}[G]$ for $\nu$-a.e.\ $\xi$. Taking $G$ in a countable algebra of bounded continuous functions generating $\cB$ makes the exceptional set independent of $G$, and gives $\theta_j$-invariance of $\mu_\xi$ for a.e.\ $\xi$.

(d) By translation invariance of $\mu_\xi$ it suffices to take $c=0$. Let $W\ge0$ be bounded and $\cI$-measurable and $F\in L^2(\Sig(\Hp{j}{0}),\mu)$. By Lemma~\ref{lem:D:invloc}(i)--(ii), $\sqrt W$ has a $\Sig(\Hp{j}{0})$-measurable version and $\Theta_j\sqrt W=\sqrt W$. Hence $\sqrt WF\in L^2(\Sig(\Hp{j}{0}))$ and
\[
\E_\mu\big[W\,\overline{F\circ\theta_j}\,F\big]=\E_\mu\big[\overline{(\sqrt WF)\circ\theta_j}\,(\sqrt WF)\big]\ge0 .
\]
With (a), $\E_{\mu_\xi}[\overline{F\circ\theta_j}F]\ge0$ for $\nu$-a.e.\ $\xi$. Let $\mathcal L$ be the countable $\Q[i]$-algebra, closed under conjugation, generated by a countable family of bounded continuous functions generating $\Sig(\Hp{j}{0})$ (H2). The exceptional set can be taken uniform over $F\in\mathcal L$. By Lemma~\ref{lem:A:mc}, $\mathcal L$ is dense in $L^2(\Sig(\Hp{j}{0}),m)$ for every Borel probability measure $m$, and \eqref{eq:D:RP} extends from $\mathcal L$ to $L^2(\Sig(\Hp{j}{0}),\mu_\xi)$ by the bound $\abs{\E[\overline{F\circ\theta_j}G]}\le\norm{F}_2\norm{G}_2$, valid because $\mu_\xi$ is $\theta_j$-invariant by (c).
\end{proof}

\begin{remark}\label{rem:D:Ohcomp}
The three coordinate reflections are inherited by the components because they act trivially on $\cI$ (Lemma~\ref{lem:D:invloc}(ii)). A coordinate permutation maps components to components (equation~\eqref{add:eq:equivariant} below), and the mixture $\mu$ is $\Oh$-invariant. After grouping the components along the orbits of the permutation group $S_3=\Oh/\langle\rho_1,\rho_2,\rho_3\rangle$ the components may be chosen $\Oh$-invariant, each a uniform mixture of at most six translation-ergodic ones (Theorem~\ref{add:Ohdecomp}). $\Oh$-invariance of the translation-ergodic components themselves does not follow from (H1)--(H3) and $\Oh$-invariance of $\mu$ (Example~\ref{add:Ohorbit}).
\end{remark}

Assume, in addition to (H1)--(H3), that $\Oh$ acts on $X$ by homeomorphisms $x\mapsto gx$, with $\rho_jx=\theta_jx$ for the coordinate reflections, $g\,\tau_a\,g^{-1}=\tau_{ga}$, $\mu$ invariant under every $g$, and $G\mapsto G\circ g^{-1}$ mapping $\Sig(U)$-measurable functions to $\Sig(gU)$-measurable ones. For the states of Theorem~\ref{thm:main} this holds on $X=H^{-s}_{\mathrm{loc}}(\R^3;\cV)$ with $(g\Phi)(f)=\Phi(g^{-1}_*f)$ (Proposition~\ref{prop:discharge}). Let $\mathcal G:=\R^3\rtimes\Oh$, let $N\cong(\Z/2)^3$ be the subgroup of $\Oh$ generated by $\rho_1,\rho_2,\rho_3$, so that $\Oh/N\cong S_3$, and let $\cI_{\mathcal G}\subset\cI$ be the $\sigma$-algebra of $\mathcal G$-invariant Borel sets, completed by $\mu$-null sets.

\begin{theorem}[decomposition with cubic symmetry]\label{add:Ohdecomp}
There are Borel probability measures $\mu'_x$ on $X$, $x\in X$, with $x\mapsto\mu'_x(B)$ $\cI_{\mathcal G}$-measurable for every Borel $B$ and $\mu=\int\mu'_x\,d\mu(x)$, such that for $\mu$-almost every $x$:
\begin{enumerate}
\item[(a)] $\mu'_x$ is invariant under $\mathcal G$ and ergodic for its action, and satisfies (H1)--(H3);
\item[(b)] there are an integer $k=k(x)\in\{1,2,3,6\}$ and pairwise distinct measures $m_1,\dots,m_k$, each invariant under the translations and the coordinate reflections, ergodic for $\R^3$ and satisfying (H1)--(H3), such that $\Oh$ permutes $\{m_1,\dots,m_k\}$ transitively and $\mu'_x=k^{-1}\sum_{i=1}^km_i$; the $m_i$ are components of the decomposition of Theorem~\ref{thm:D:decomp};
\item[(c)] in the reconstruction from $\mu'_x$ in any coordinate direction $j$ (Proposition~\ref{prop:D:recon}), the joint fixed space of $\{\mathrm{e}^{-tH_j}\}_{t\ge0}$ and $\{V_j(y)\}_{y\perp e_j}$ has dimension $k$;
\item[(d)] for all $G,G'\in L^2(\mu'_x)$, the Ces\`aro means over cubes of $\E_{\mu'_x}[G'\,U(a)G]$, and $\E_{\mu'_x}[G'\,U(a)G]$ itself as $\min_l\abs{a_l}\to\infty$, converge to $k^{-1}\sum_{i=1}^k\E_{m_i}[G']\,\E_{m_i}[G]$.
\end{enumerate}
Moreover $k(x)=1$ for $\mu$-almost every $x$ if and only if almost every component of Theorem~\ref{thm:D:decomp} is $\Oh$-invariant.
\end{theorem}
\begin{proof}
\emph{Step 1: $N$ acts trivially on $\cI$, and $\Oh$ preserves $\cI$.} A bounded $\cI$-measurable $W$ is invariant along every coordinate axis, so $W\circ\theta_j=W$ a.s.\ for $j=1,2,3$ by Lemma~\ref{lem:D:invloc}(ii); hence $W\circ n=W$ for $n\in N$. For $g\in\Oh$ and a translation-invariant set $B$, $\tau_a(gB)=g\,\tau_{g^{-1}a}B=gB$, so $g$ maps $\cI$ onto itself.

\emph{Step 2: conditional expectations.} Let $\mu_x$ be the components of Theorem~\ref{thm:D:decomp} and $\mu'_x$ a regular conditional probability of $\mu$ given $\cI_{\mathcal G}$ (Lemma~\ref{lem:D:rcp}, applied to the $\sigma$-algebra of $\mathcal G$-invariant Borel sets, whose $\mu$-completion is $\cI_{\mathcal G}$). For bounded Borel $F$ and $g\in\Oh$, $g$ preserves $\mu$ and maps $\cI$ onto itself, so $\E[F\circ g\mid\cI]=\E[F\mid\cI]\circ g$ a.s., that is
\begin{equation}\label{add:eq:equivariant}
\mu_{gx}=g_*\mu_x\qquad\text{for $\mu$-a.e.\ $x$,}
\end{equation}
first for $F$ in a countable algebra of bounded continuous functions generating $\cB$ and then as measures. By Step~1 and \eqref{add:eq:equivariant} with $g\in N$, $n_*\mu_x=\mu_x$ for $n\in N$, so $S_3=\Oh/N$ acts on the components by $\sigma\cdot\mu_x:=g_*\mu_x$, $g\in\sigma$. For bounded $\cI$-measurable $h$ put $\bar h:=\abs{\Oh}^{-1}\sum_{g\in\Oh}h\circ g$. It is $\cI$-measurable and $\Oh$-invariant, hence $\cI_{\mathcal G}$-measurable, and for bounded $\cI_{\mathcal G}$-measurable $W$ (which satisfies $W\circ g=W$ a.s.) $\E[W\bar h]=\abs{\Oh}^{-1}\sum_g\E[(W\circ g)(h\circ g)]=\E[Wh]$. So $\E[h\mid\cI_{\mathcal G}]=\bar h$. Applied to $h=\E[F\mid\cI]$, with $\E[F\mid\cI_{\mathcal G}]=\E[\E[F\mid\cI]\mid\cI_{\mathcal G}]$ and \eqref{add:eq:equivariant},
\[
\mu'_x=\frac1{48}\sum_{g\in\Oh}\mu_{gx}=\frac1{48}\sum_{g\in\Oh}g_*\mu_x=\frac16\sum_{\sigma\in S_3}\sigma\cdot\mu_x\qquad\text{for $\mu$-a.e.\ $x$.}
\]
Grouping equal terms, $\mu'_x=k^{-1}\sum_{i=1}^km_i$, where $\{m_1,\dots,m_k\}$ is the $S_3$-orbit of $\mu_x$, so $k=6/\abs{\mathrm{Stab}(\mu_x)}\in\{1,2,3,6\}$. For a.e.\ $x$ every $\mu_{gx}$ is a component with the properties of Theorem~\ref{thm:D:decomp}(b)--(e); this gives (b). Since $k=1$ exactly when $\mu_x$ is fixed by $S_3$, that is, $\Oh$-invariant, the last assertion follows.

\emph{Step 3: (a).} $\mu'_x$ is invariant under the translations and under $\Oh$, as an average over an $\Oh$-orbit of translation-invariant measures. If $B$ is a $\mathcal G$-invariant Borel set, then $m_i(B)\in\{0,1\}$ by ergodicity, and $(g_*m_i)(B)=m_i(g^{-1}B)=m_i(B)$, so all $m_i(B)$ are equal by transitivity and $\mu'_x(B)\in\{0,1\}$. Each $m_i$ satisfies (H3), and (H3) is preserved under convex combinations. (H2) holds as in Theorem~\ref{thm:D:decomp}(e).

\emph{Step 4: invariant sets separating the $m_i$.} Let $m\ne m'$ be two of the $m_i$, and $F$ bounded continuous with $m(F)\ne m'(F)$. With $A_TF:=\abs{\Lambda_T}^{-1}\int_{\Lambda_T}U(a)F\,da$, Lemma~\ref{lem:D:proj}(i) and ergodicity give $A_TF\to m(F)$ in $L^2(m)$ and $A_TF\to m'(F)$ in $L^2(m')$; choose $T_n\to\infty$ with a.s.\ convergence for both. Since $\sup_x\abs{A_TF(\tau_bx)-A_TF(x)}\le2\norm F_\infty\abs{(\Lambda_T+b)\triangle\Lambda_T}/\abs{\Lambda_T}\to0$ for each $b$, the sets $\{x:\lim_nA_{T_n}F(x)=m(F)\}$ and $\{x:\lim_nA_{T_n}F(x)=m'(F)\}$ are disjoint translation-invariant Borel sets of full $m$-measure and full $m'$-measure respectively. Doing this for every pair and intersecting gives pairwise disjoint translation-invariant Borel sets $B_1,\dots,B_k$ with $m_i(B_i)=1$.

\emph{Step 5: (c) and (d).} Let $\mu':=\mu'_x$ and $\cI'$ the $\sigma$-algebra of translation-invariant sets for $\mu'$. An $\cI'$-measurable function is $m_i$-a.s.\ constant for each $i$, so $L^2(\cI',\mu')$ is spanned by $\mathbf 1_{B_1},\dots,\mathbf 1_{B_k}$, has dimension $k$, and $\E_{\mu'}[G\mid\cI']=\sum_i\E_{m_i}[G]\,\mathbf 1_{B_i}$. By Proposition~\ref{prop:D:fixed}(ii) for $\mu'$, the projection onto the joint fixed space is $[G]\mapsto[\E_{\mu'}[G\mid\cI']]$, and $W\mapsto[W]$ is isometric on $L^2(\cI',\mu')$, because $\norm{[W]}^2=\E_{\mu'}[\overline{\Theta_jW}\,W]=\E_{\mu'}\abs W^2$ by Lemma~\ref{lem:D:invloc}(ii); the $\mathbf 1_{B_i}$ lie in $L^2_+$ by Lemma~\ref{lem:D:invloc}(i). So the joint fixed space is spanned by the $k$ orthogonal vectors $[\mathbf 1_{B_i}]$, which is (c). Theorems~\ref{thm:D:mean} and~\ref{thm:D:diag} for $\mu'$ give the limit $\E_{\mu'}\big[\E[G'\mid\cI']\,\E[G\mid\cI']\big]=k^{-1}\sum_i\E_{m_i}[G']\,\E_{m_i}[G]$, which is (d).
\end{proof}

\begin{remark}\label{add:Ohdecomp-rem}
(1) The components $\mu'_x$ satisfy OS0 in the form E0$''$ with $x$-dependent constants: the proof of Corollary~\ref{cor:D:schwinger}(iv) applies verbatim with the decomposition $\mu=\int\mu'_x\,d\mu(x)$. As for the components of Theorem~\ref{thm:D:decomp}, the tube form of OS0 and the clause of OS2 for exponentials are not claimed for them. They satisfy OS1 for $\R^3\rtimes\Oh$, OS2 about every coordinate plane and the symmetry and reality conditions of OS3; they are ergodic for $\mathcal G$, and their joint vacuum has dimension $k(x)\le6$.
(2) Charge conjugation $\mathsf C$ ($F\mapsto-F$, $\XR\mapsto\XR$) commutes with $\mathcal G$ and preserves $\mu$ (Theorem~\ref{thm:C-OS3} and the uniqueness in Proposition~\ref{prop:discharge}). The same argument with $\Oh\times\{1,\mathsf C\}$ in place of $\Oh$ gives components that are moreover $\mathsf C$-invariant, so that their moments of odd degree in the curvature vanish, with $k\in\{1,2,3,4,6,12\}$ (the group acts through the quotient $S_3\times\Z/2$, whose subgroups have index $1,2,3,4,6$ or $12$).
(3) In the reconstruction from $\mu'_x$ in direction $j$, the stabiliser $D_4$ of $e_j$ acts unitarily as in Theorem~\ref{thm:D:recon}(ii). It permutes the $k$ joint vacua $[\mathbf 1_{B_i}]$ through its image in $S_3$, which is generated by the transposition of the two coordinates transverse to $e_j$: if $g_*m_i=m_{i'}$, then $\mathbf 1_{gB_i}=\mathbf 1_{B_{i'}}$ $\mu'_x$-a.s.
\end{remark}

\begin{example}[cubic symmetry is not inherited by the components]\label{add:Ohorbit}
Let $m_1$ be the law of the field $\Phi(f):=\psi(f^{(1)})$, with $\psi$ the massive free field on $\R^2$ and $f^{(1)}$ as in Example~\ref{ex:D:axis}, realised on $\cV$ with vanishing curvature component; this is the summand $\psi_1(x_2,x_3)$ of that example alone. Let $m_2,m_3$ be its images under the transpositions of the coordinates $(x_1,x_2)$ and $(x_1,x_3)$. The proofs of Example~\ref{ex:D:axis}(ii)--(iii), with one plane instead of three, show that each $m_i$ satisfies (H1)--(H3) and is ergodic for $\R^3$; in the direction transverse to its plane reflection positivity holds because every function measurable for a half space is invariant under the reflection. The measure $m_1$ is invariant under the coordinate reflections and under the transposition of $(x_2,x_3)$, but not under that of $(x_1,x_2)$. So $\mu:=\frac13(m_1+m_2+m_3)$ satisfies (H1)--(H3), is invariant under $\R^3\rtimes\Oh$, has the translation-ergodic components $m_1,m_2,m_3$, none of which is $\Oh$-invariant, and $k=3$ in Theorem~\ref{add:Ohdecomp}. Thus $\Oh$-invariance of the translation-ergodic components does not follow from (H1)--(H3) and $\Oh$-invariance of the law.
\end{example}

\begin{remark}\label{add:Ohdecomp-B}
The components $\mu_\xi$ of Theorem~\ref{thm:cluster} are invariant under the translations and the three coordinate reflections, and a coordinate permutation maps a component to a component. Grouping the components along the orbits of $\Oh$ gives a decomposition into $\Oh$-invariant, reflection-positive components, each a uniform mixture of at most $6$ translation-ergodic ones, with a joint vacuum of dimension at most $6$ in each coordinate direction (Theorem~\ref{add:Ohdecomp}). Whether almost every translation-ergodic component is itself $\Oh$-invariant is not known; it holds if and only if $k=1$ almost everywhere in Theorem~\ref{add:Ohdecomp}.
\end{remark}

\subsection{Reconstruction in a coordinate direction}\label{ss:recon}

Fix $\mu$ satisfying (H1)--(H3) and a direction $j$. Write $L^2_+:=L^2(\Sig(\Hp{j}{0}),\mu)$ and define the form
\[
\ipj{j}{F}{G}:=\E_\mu\big[\overline{\Theta_jF}\,G\big],\qquad F,G\in L^2_+ .
\]
It is positive semidefinite by (H3), Hermitian because $\Theta_j$ is a unitary involution preserving $\mu$, and bounded: $\abs{\ipj{j}{F}{G}}\le\norm{F}_2\norm{G}_2$.

\begin{proposition}[reconstruction in direction $j$]\label{prop:D:recon}
Let $\cN=\{F\in L^2_+:\ipj{j}{F}{F}=0\}$, $\cH_j$ the completion of $L^2_+/\cN$, $[\,\cdot\,]:L^2_+\to\cH_j$ the quotient map, and $\Omega_j=[1]$.
\begin{enumerate}
\item[(i)] For $t\ge0$ and $y\perp e_j$, $T_j(t)[G]:=[U(te_j)G]$ and $V_j(y)[G]:=[U(y)G]$ are well defined. $T_j$ is a strongly continuous semigroup of self-adjoint contractions, so $T_j(t)=e^{-tH_j}$ with $H_j\ge0$ self-adjoint. $V_j$ is a strongly continuous unitary representation of $e_j^\perp\cong\R^2$. All these operators commute, and $T_j(t)\Omega_j=V_j(y)\Omega_j=\Omega_j$.
\item[(ii)] For $G\in L^2_+$ and $a=a_je_j+a_\perp$ with $a_j\ge0$: $[U(a)G]=T_j(a_j)V_j(a_\perp)[G]$.
\item[(iii)] For $F\in L^2(\Sig(\Hm{j}{0}),\mu)$ and $G\in L^2_+$: $\overline{\Theta_jF}\in L^2_+$ and $\E_\mu[FG]=\ip{[\overline{\Theta_jF}]}{[G]}_{\cH_j}$.
\end{enumerate}
\end{proposition}
\begin{proof}
This is the Osterwalder--Schrader construction \cite{OS73,OS75}, \cite[Thm.~6.1.3, p.~94]{GJ}; we write out the identities that drive it. For $G\in L^2_+$ and $t\ge0$, $U(te_j)G\in L^2(\Sig(\Hp{j}{t}))\subset L^2_+$ by (H2). By \eqref{eq:D:commute} and invariance of $\mu$,
\[
\ipj{j}{U(te_j)F}{U(te_j)G}=\E\big[\overline{U(-te_j)\Theta_jF}\;U(te_j)G\big]=\ipj{j}{F}{U(2te_j)G}.
\]
So the form is symmetric for $U(te_j)$, and by the Schwarz inequality for the form, $\ipj{j}{U(te_j)G}{U(te_j)G}\le\ipj{j}{G}{G}^{1/2}\ipj{j}{U(2te_j)G}{U(2te_j)G}^{1/2}$. Iterating $n$ times and using $\ipj{j}{U(se_j)G}{U(se_j)G}\le\norm{G}_2^2$ gives $\ipj{j}{U(te_j)G}{U(te_j)G}\le\ipj{j}{G}{G}^{1-2^{-n}}\norm{G}_2^{2^{1-n}}$, and $n\to\infty$ gives $\ipj{j}{U(te_j)G}{U(te_j)G}\le\ipj{j}{G}{G}$. So $T_j(t)$ maps $\cN$ into $\cN$ and is a symmetric contraction on the quotient, hence a self-adjoint contraction on $\cH_j$; the semigroup law is inherited from $U$. Strong continuity follows from $\norm{T_j(t)[G]-[G]}\le\norm{U(te_j)G-G}_2$ and Lemma~\ref{lem:D:strongcont}. For $y\perp e_j$, $\Theta_j$ commutes with $U(y)$, so $\ipj{j}{U(y)F}{U(y)G}=\ipj{j}{F}{G}$: $V_j(y)$ is isometric with isometric inverse $V_j(-y)$. Commutation is inherited from $U$, and (ii) is the definition. (iii): $\overline{\Theta_jF}$ is $\Sig(\Hp{j}{0})$-measurable by (H2), and $\ipj{j}{\overline{\Theta_jF}}{G}=\E\big[\overline{\Theta_j\overline{\Theta_jF}}\,G\big]=\E[FG]$.
\end{proof}

\begin{proposition}[zero-energy vectors and invariant functions]\label{prop:D:fixed}
In the setting of Proposition~\ref{prop:D:recon}, for every $G\in L^2_+$:
\begin{enumerate}
\item[(i)] $P_jG\in L^2_+$, $P_{\Ker H_j}[G]=[P_jG]$, and $\norm{[P_jG]}^2_{\cH_j}=\E_\mu\abs{P_jG}^2$;
\item[(ii)] $\E[G\mid\cI]\in L^2_+$ and $Q_j[G]=[\E[G\mid\cI]]$, where $Q_j$ is the projection onto the joint fixed space of $\{T_j(t)\}_{t\ge0}\cup\{V_j(y)\}_{y\perp e_j}$;
\item[(iii)] if $\mu$ is ergodic for the $\R^3$-action, the joint fixed space is $\C\Omega_j$.
\end{enumerate}
\end{proposition}
\begin{proof}
(i) $P_jG\in L^2_+$ by Lemma~\ref{lem:D:invloc}(i). The quotient map is bounded, so it commutes with Bochner integrals: $T^{-1}\int_0^TT_j(t)[G]\,dt=[T^{-1}\int_0^TU(te_j)G\,dt]$. As $T\to\infty$ the left side tends to $P_{\Ker H_j}[G]$ (spectral theorem, $T^{-1}\int_0^Te^{-t\lambda}dt\to\mathbf 1_{\{\lambda=0\}}$ boundedly on $[0,\infty)$) and the right side to $[P_jG]$ (one-sided means, Lemma~\ref{lem:D:proj}(i)). Finally $\norm{[P_jG]}^2=\E[\overline{\Theta_jP_jG}\,P_jG]=\E\abs{P_jG}^2$ by Lemma~\ref{lem:D:invloc}(ii).

(ii) $\E[G\mid\cI]=P_{\{1,2,3\}}G\in L^2_+$ by Lemma~\ref{lem:D:invloc}(i). The means of $W(a):=T_j(a_j)V_j(a_\perp)$ over $a\in[0,T]\times[-T,T]^2$ factor into three commuting means and converge strongly to the product of the projections onto the three fixed spaces, which is $Q_j$. Applying the same means to $[U(a)G]=W(a)[G]$ and using the argument of (i), the limit is $[P_{\{1,2,3\}}G]$.

(iii) If $\mu$ is ergodic, $\E[G\mid\cI]=\E[G]$ a.s., so $Q_j[G]=\E[G]\,\Omega_j=\ip{\Omega_j}{[G]}\Omega_j$. Since $[L^2_+]$ is dense in $\cH_j$, $Q_j$ is the projection onto $\C\Omega_j$.
\end{proof}

\begin{remark}\label{rem:D:kerH}
$\Ker H_j$ alone can be larger than $\C\Omega_j$ on an ergodic component: by Proposition~\ref{prop:D:fixed}(i) it is the image of the $e_j$-invariant functions, which need not be constant (Example~\ref{ex:D:axis}). Ergodicity controls only the joint fixed space.
\end{remark}

\subsection{Clustering}\label{ss:cluster}

``On an ergodic component'' means: for $\mu_\xi$, $\nu$-a.e.\ $\xi$, in the decomposition of Theorem~\ref{thm:D:decomp}. Each such $\mu_\xi$ satisfies (H1)--(H3) and is ergodic.

\begin{theorem}[clustering in the mean]\label{thm:D:mean}
For all $G,G'\in L^2(\mu)$ and $\Lambda_T=[-T,T]^3$,
\[
\lim_{T\to\infty}\frac1{\abs{\Lambda_T}}\int_{\Lambda_T}\E_\mu\big[G'\,U(a)G\big]\,da=\E_\mu\big[\E[G'\mid\cI]\,\E[G\mid\cI]\big]=\int_\Xi\E_{\mu_\xi}[G']\,\E_{\mu_\xi}[G]\,d\nu(\xi).
\]
On an ergodic component the right side is $\E_{\mu_\xi}[G']\,\E_{\mu_\xi}[G]$. For $\emptyset\neq J\subset\{1,2,3\}$ the first equality holds with $\Lambda_T$ replaced by $\Lambda^J_T$ and $\cI$ by $\cI_J$ (Lemma~\ref{lem:D:proj}(ii)). For $J\neq\{1,2,3\}$ the right side need not factorise on an ergodic component (Example~\ref{ex:D:axis}).
\end{theorem}
\begin{proof}
By Lemma~\ref{lem:D:proj}(i)--(ii) the means of $U(a)G$ converge in $L^2$ to $\E[G\mid\cI]$, and $\E[G'\,\E[G\mid\cI]]=\E[\E[G'\mid\cI]\,\E[G\mid\cI]]$. The integral form is Theorem~\ref{thm:D:decomp}(a).
\end{proof}

\begin{theorem}[clustering under diagonal separation]\label{thm:D:diag}
For all $G,G'\in L^2(\mu)$,
\[
\lim_{\min_j\abs{a_j}\to\infty}\E_\mu\big[G'\,U(a)G\big]=\E_\mu\big[\E[G'\mid\cI]\,\E[G\mid\cI]\big].
\]
On an ergodic component the limit is $\E_{\mu_\xi}[G']\,\E_{\mu_\xi}[G]$.
\end{theorem}
\begin{proof}
\emph{Step 1 (reductions).} Since $\abs{\E[G'U(a)G]-\E[G'_nU(a)G_n]}\le\norm{G'-G'_n}_2\norm{G}_2+\norm{G'_n}_2\norm{G-G_n}_2$ uniformly in $a$, and the right side of the claim is continuous in $(G,G')$ in the same way, it suffices to take $G\in L^2(\Sig(K))$ and $G'\in L^2(\Sig(K'))$ with $K,K'$ bounded (H2). A sequence $a^{(n)}$ with $\min_j\abs{a^{(n)}_j}\to\infty$ splits into finitely many subsequences on which the signs $\sigma_j=\operatorname{sgn}a^{(n)}_j$ are fixed. On each, replace $(G',G,a)$ by $(\Theta_jG',\Theta_jG,\rho_ja)$ for those $j$ with $\sigma_j=-1$; this leaves $\E[G'U(a)G]$ unchanged by $\theta_j$-invariance and \eqref{eq:D:commute}, and reflects the localisation regions. So assume $a_j\to+\infty$ for every $j$. By translation invariance, $\E[G'U(a)G]=\E[(U(b)G')\,U(a)(U(b)G)]$; choosing $b$ we may assume that $K'$ lies in the open negative octant $\bigcap_j\Hm{j}{0}$ and that $K\subset\{x:\abs{x_j}<R\ \forall j\}$ for some $R$.

\emph{Step 2 (decomposition).} For the ordering $(1,2,3)$,
\begin{equation}\label{eq:D:telescope}
1=P_{\{1,2,3\}}+(1-P_1)+P_1(1-P_2)+P_1P_2(1-P_3).
\end{equation}
The first term contributes $\E[G'\,U(a)\E[G\mid\cI]]=\E[\E[G'\mid\cI]\,\E[G\mid\cI]]$ for every $a$, since $\E[G\mid\cI]$ is invariant. It remains to show that each of the other three terms tends to $0$ as $a_\ell\to\infty$, uniformly in the other coordinates of $a$, where $\ell$ is the index of the factor $(1-P_\ell)$.

\emph{Step 3 (one term).} Fix $\ell$, let $J=\{1,\dots,\ell-1\}$ and $\tilde G=P_JG$ ($\tilde G=G$ for $\ell=1$). The term is $P_J(1-P_\ell)G=(1-P_\ell)\tilde G$. By Lemma~\ref{lem:D:invloc}(iii), $\tilde G\in L^2(\Sig(K+V_J))$, and $K+V_J\subset\{x_\ell>-R\}$ because $e_\ell\notin V_J$. Put $s=2R$, $a_\ell\ge s$, and $Y:=U(se_\ell)\tilde G\in L^2(\Sig(\Hp{\ell}{R}))\subset L^2_+$ (half spaces for direction $\ell$). Since $P_\ell$ commutes with translations,
\[
U(a)(1-P_\ell)\tilde G=U\big((a_\ell-s)e_\ell+a_\perp\big)(1-P_\ell)Y,
\]
and $Y$, $P_\ell Y$ lie in $L^2_+$ (Lemma~\ref{lem:D:invloc}(i)). Since $G'\in L^2(\Sig(\Hm{\ell}{0}))$, Proposition~\ref{prop:D:recon}(ii)--(iii) and Proposition~\ref{prop:D:fixed}(i) give, with $F:=\overline{\Theta_\ell G'}\in L^2_+$,
\[
\E\big[G'\,U(a)(1-P_\ell)\tilde G\big]=\Big\langle[F],\ V_\ell(a_\perp)\,T_\ell(a_\ell-s)\,(1-P_{\Ker H_\ell})[Y]\Big\rangle_{\cH_\ell}.
\]
Its modulus is at most $\norm{[F]}\,\norm{e^{-(a_\ell-s)H_\ell}(1-P_{\Ker H_\ell})[Y]}$, which tends to $0$ as $a_\ell\to\infty$ by the spectral theorem and dominated convergence, uniformly in $a_\perp$.
\end{proof}

\begin{proposition}[separation along a coordinate subgroup]\label{prop:D:subgroup}
Let $\emptyset\ne J\subset\{1,2,3\}$ and $G,G'\in L^2(\mu)$. Then, writing $a=a_J+a_{J^c}$ with $a_J\in V_J$,
\[
\E_\mu\big[G'\,U(a)G\big]-\E_\mu\big[G'\,U(a_{J^c})\,P_JG\big]\to0\qquad\text{as }\min_{j\in J}\abs{a_j}\to\infty,
\]
uniformly in $a_{J^c}$. For $j\in\{1,2,3\}$ and $G\in L^2(\Sig(\Hp{j}{0}),\mu)$, the function $t\mapsto\E_\mu[\overline{\Theta_jG}\,U(te_j)G]=\ip{[G]}{e^{-tH_j}[G]}_{\cH_j}$, $t\ge0$, is completely monotone and decreases to $\E_\mu\abs{P_jG}^2$; it tends to $0$ if and only if $P_jG=0$.
\end{proposition}
\begin{proof}
Apply \eqref{eq:D:telescope} with $\{1,2,3\}$ replaced by $J$ in any ordering; Step~3 of the proof of Theorem~\ref{thm:D:diag} applies verbatim to each term carrying a factor $(1-P_\ell)$, $\ell\in J$, after the reductions of Step~1, which do not involve $a_{J^c}$. The two-point statement is Proposition~\ref{prop:D:recon}(ii) and Proposition~\ref{prop:D:fixed}(i), with the spectral measure of $H_j$ in the vector $[G]$.
\end{proof}

\subsection{Single-axis clustering: the obstruction and what removes it}\label{ss:axis}

\begin{example}[the axis obstruction]\label{ex:D:axis}
Let $m>0$ and let $\psi_1,\psi_2,\psi_3$ be independent copies of the massive free field on $\R^2$, the centred Gaussian law on $\cS'(\R^2)$ with covariance $(-\Delta_2+m^2)^{-1}$. For $f\in\cS(\R^3)$ let $f^{(j)}\in\cS(\R^2)$ be the integral of $f$ over the coordinate $x_j$, and put
\[
\Phi(f):=\psi_1(f^{(1)})+\psi_2(f^{(2)})+\psi_3(f^{(3)}),
\]
the random distribution $\psi_1(x_2,x_3)+\psi_2(x_1,x_3)+\psi_3(x_1,x_2)$. Its law $\mu$ is the centred Gaussian measure with covariance $C=C_1+C_2+C_3$, $\widehat{C_j}(p)=2\pi\,\delta(p_j)\,(\abs{p_{\perp j}}^2+m^2)^{-1}$, where $p_{\perp j}$ is the projection of $p$ onto $e_j^\perp$. Realise it on the field space $\cV=\Lambda^2\R^3\oplus\R$ by putting the 2-form component equal to $0$ and the scalar component equal to $\Phi$, on $X=H^{-s}_{\mathrm{loc}}(\R^3;\cV)$, $s>1$. By the carrier argument of Lemma~\ref{lem:D:hermite}(b), with $\E\Phi(f)^2\le m^{-2}\sum_j\norm{f^{(j)}}^2_{L^2(\R^2)}\le c\norm f_2^2$ (Cauchy--Schwarz inequality in $x_j$ with weight $1+x_j^2$), $\mu$ is also carried by $\mathcal H^{-r}$ for $r>5$, as Proposition~\ref{prop:axisobs} states. Then:
\begin{enumerate}
\item[(i)] $\mu$ is invariant under translations and under $\Oh$, and has Gaussian moments, so every moment bound of the form $\abs{\E\prod_{i\le n}\Phi(f_i)}\le C^n(n!)^{1/2}\prod_i\norm{f_i}$ holds with the norm $\sum_j\norm{f^{(j)}}_{L^2(\R^2)}$. This norm is not bounded by $\norm f_{L^2(\R^3)}$; a bound $\E\Phi(f)^2\le c\norm f^2_{L^2(\R^3)}$ would give $\E e^{\Phi(f)}\le e^{c\norm f_2^2/2}$ for this Gaussian law, and Lemma~\ref{lem:os4:dom}, with $\varpi$ equal to $c$ times Lebesgue measure on the scalar component, would force $P_1\Phi(f)=\E[\Phi(f)\mid\cI]=0$ ($\mu$ is ergodic by (iii)), contradicting (iv).
\item[(ii)] $\mu$ satisfies (H1)--(H3) with the reflections of Section~\ref{ss:discharge}. On the scalar component these are the plain reflections $h\mapsto h\circ\rho_j$, and the 2-form component vanishes, so the twisted and the plain reflections coincide.
\item[(iii)] $\mu$ is ergodic for the $\R^3$-action.
\item[(iv)] For $f\in C_c^\infty(\R^3)$ with $f^{(1)}\ne0$, $\E[\Phi(f)\,U(te_1)\Phi(f)]\to\ip{f^{(1)}}{(-\Delta_2+m^2)^{-1}f^{(1)}}>0$ as $t\to\infty$. Thus $P_1\Phi(f)=\psi_1(f^{(1)})\neq0$, and the two-point function does not cluster along the coordinate axes, although it clusters under diagonal separation (Theorem~\ref{thm:D:diag}).
\end{enumerate}
\end{example}
\begin{proof}
$\Phi$ takes values in $X$: $\chi(x)\psi_1(x_2,x_3)$ lies in $H^{-s}(\R^3)$ for $\chi\in C_c^\infty(\R^3)$ when $\psi_1\in H^{-s}_{\mathrm{loc}}(\R^2)$, which holds a.s.\ for $s>1$, by a Fourier-space bound in which the weight $(1+\abs p^2)^{-s}\le(1+\abs{p_\perp}^2)^{-s}$ is used.

(i) Translations act on $\psi_j$ through the translations of its plane, under which the free field is invariant. Signed permutations of the coordinates permute the three planes and act on each by an element of $O(2)$; the $\psi_j$ are identically distributed and $O(2)$-invariant. Moments are Gaussian with $\ip{f}{C_jf}=\ip{f^{(j)}}{(-\Delta_2+m^2)^{-1}f^{(j)}}\le m^{-2}\norm{f^{(j)}}^2_{L^2(\R^2)}$.

(ii) (H1)--(H2) as in the proof of Proposition~\ref{prop:discharge}. For (H3) it suffices (by Lemma~\ref{lem:A:gaussRP}) that $\ip{\rho_j f}{Cf}\ge0$ for real $f\in C_c^\infty(\Hp{j}{0})$. For $C_j$: $(\rho_jf)^{(j)}=f^{(j)}$, so $\ip{\rho_jf}{C_jf}=\ip{f^{(j)}}{(-\Delta_2+m^2)^{-1}f^{(j)}}\ge0$. For $k\ne j$: $(\rho_jf)^{(k)}$ is the reflection of $f^{(k)}$ in the coordinate $x_j$ of the plane $e_k^\perp$, and $f^{(k)}$ is supported in $\{x_j>0\}$, so $\ip{\rho_jf}{C_kf}\ge0$ is reflection positivity of the two-dimensional massive free field \cite[Prop.~6.2.5, p.~103]{GJ} (the passage from the covariance to the Gaussian measure is \cite[Thm.~6.2.2, p.~101]{GJ} and Lemma~\ref{lem:A:gaussRP}). Sums of positive forms are positive.

(iii) By the Wiener--It\^o chaos decomposition, $L^2(\mu)=\bigoplus_{n\ge0}\cH^{:n:}$, and $\cH^{:n:}$ is unitarily equivalent to the symmetric subspace of $L^2((\R^3)^n,\widehat C^{\otimes n})$, with $U(a)$ acting by multiplication by $e^{-ia\cdot(p_1+\dots+p_n)}$ \cite[Thm.~2.6 and Ch.~4]{Jan97}. A vector fixed by all $U(a)$ in $\cH^{:n:}$, $n\ge1$, is supported on $Z_n=\{p_1+\dots+p_n=0\}$. The measure $\widehat C$ is a sum of three measures, the $j$-th supported on the plane $\pi_j=\{p_j=0\}$ and absolutely continuous with respect to its two-dimensional Lebesgue measure. For each choice $(j_1,\dots,j_n)$ the map $(p_1,\dots,p_n)\mapsto\sum_ip_i$ from $\pi_{j_1}\times\dots\times\pi_{j_n}$ to $\R^3$ is linear and onto a subspace of dimension $\ge2$, so its kernel $Z_n\cap(\pi_{j_1}\times\cdots\times\pi_{j_n})$ has codimension $\ge2$ there and is null for the product of the planar Lebesgue measures. Hence $\widehat C^{\otimes n}(Z_n)=0$, the only invariant vectors are the constants, and $\mu$ is ergodic.

(iv) $\E[\Phi(f)\,U(te_1)\Phi(f)]=\sum_k\ip{f}{C_kT_{te_1}f}$. For $k=1$, $(T_{te_1}f)^{(1)}=f^{(1)}$, giving the constant $\ip{f^{(1)}}{(-\Delta_2+m^2)^{-1}f^{(1)}}$, which is positive when $f^{(1)}\ne0$. For $k=2,3$, $(T_{te_1}f)^{(k)}$ is $f^{(k)}$ translated by $t$ in the plane, and the term tends to $0$ because the two-dimensional massive Green function decays. The statement on $P_1$ follows from Proposition~\ref{prop:D:subgroup}.
\end{proof}

\begin{proposition}[what removes the obstruction]\label{prop:D:axis}
Let $\mu$ satisfy (H1)--(H3) and let $\cF\subset L^2(\mu)$ be a set whose linear span is dense. The following are equivalent:
\begin{enumerate}
\item[(i)] $P_jG=\E[G\mid\cI]$ for all $G\in\cF$ and $j=1,2,3$;
\item[(ii)] for all $G\in\cF$, $G'\in L^2(\mu)$ and each $j$: $\E[G'\,U(te_j)G]\to\E[\E[G'\mid\cI]\,\E[G\mid\cI]]$ as $\abs t\to\infty$;
\item[(iii)] for all $G\in\cF$, $G'\in L^2(\mu)$: $\E[G'\,U(a)G]\to\E[\E[G'\mid\cI]\,\E[G\mid\cI]]$ as $\abs a\to\infty$ in any manner.
\end{enumerate}
If $\mu$ is a centred Gaussian law with covariance $C$ and $\cF$ is the set of linear fields $\Phi(f)$, then (i) holds if and only if $\widehat C(\{p:p_j=0\}\setminus\{0\})=0$ for $j=1,2,3$.
\end{proposition}
\begin{proof}
(i)$\Rightarrow$(iii): given $\abs{a^{(n)}}\to\infty$, every subsequence has a further subsequence along which each coordinate converges or diverges in modulus, at least one, say the $j$-th, diverging. Proposition~\ref{prop:D:subgroup} with $J=\{j\}$ reduces the limit to that of $\E[G'U(a_{J^c})P_jG]=\E[G'\,\E[G\mid\cI]]$. (iii)$\Rightarrow$(ii) is trivial. (ii)$\Rightarrow$(i): by Proposition~\ref{prop:D:subgroup} the limit in (ii) is $\E[G'P_jG]$; so $\E[G'(P_jG-\E[G\mid\cI])]=0$ for all $G'\in L^2$, i.e.\ $P_jG=\E[G\mid\cI]$. The Gaussian clause: the first chaos $\cH^{:1:}$ is invariant under $U$ and unitarily equivalent to the closure of $\{\hat f\}$ in $L^2(\R^3,\widehat C)$, with $U(te_j)$ acting by multiplication by $e^{-itp_j}$. So $P_j\Phi(f)$ corresponds to $\mathbf 1_{\{p_j=0\}}\hat f$ and $\E[\Phi(f)\mid\cI]$ to $\mathbf 1_{\{p=0\}}\hat f$; they agree for all $f$ exactly when $\widehat C$ gives no mass to $\{p_j=0\}\setminus\{0\}$.
\end{proof}

\subsection{Single-axis clustering}\label{ss:os4axis}

Throughout this subsection $S$ is one of the limit states $S^{\mu_0}$ ($\mu_0\in Q_0$), $S^\fl$, $S^0$ of Theorem~\ref{thm:main}, and $\mu=\mu_S$ is its law on $X=H^{-s}_{\mathrm{loc}}(\R^3;\cV)$ (Lemma~\ref{lem:D:realize}); everything holds verbatim on $\mathcal H^{-r}$ (Lemma~\ref{lem:D:hermite}). We write $F(\bg)$ for the curvature field and $\XR(h)$ for the Wick square, $h\in C_c^\infty(\R^3)$ real; $\mathcal M$ is the $\sigma$-algebra generated by the $\XR(h)$ and completed by $\mu$-null sets, and $L^2(\mathcal M)\subset L^2(\mu)$ is the corresponding closed subspace. $P_j$ is the projection onto the $e_j$-invariant vectors (Section~\ref{ss:proj}). We say that a pair $(G',G)$ in $L^2(\mu)$ \emph{clusters along the axes} if $\E_\mu[G'\,U(a)G]\to\E_\mu[G']\,\E_\mu[G]$ as $\abs a\to\infty$ in any manner; in particular as $a=te_j$, $\abs t\to\infty$, for each $j$. If every pair in $L^2(\mu)$ clusters along the axes, then $\mu$ satisfies (OS4$'_j$) of Definition~\ref{def:OS4} for $j=1,2,3$; under (H1)--(H3) this is equivalent to OS4 (Lemma~\ref{lem:os4:red} below, and Remark~\ref{p:rem:OS4forms}(c)).

The single tool that turns invariant parts into clustering is the following reformulation of Propositions~\ref{prop:D:subgroup} and~\ref{prop:D:axis}.

\begin{lemma}[clustering from vanishing invariant parts]\label{lem:os4:red}
Let $\mu$ satisfy (H1)--(H3) and $G\in L^2(\mu)$. If $P_jG=\E_\mu G$ for $j=1,2,3$, then $(G',G)$ clusters along the axes for every $G'\in L^2(\mu)$. Conversely, let $\mathcal S\subset L^2(\mu)$ be a closed subspace, closed under complex conjugation, containing $1$ and $P_jG$ (for instance $\mathcal S=L^2(\mu)$). If $\E[G'\,U(te_j)G]\to\E[G']\E[G]$ as $t\to\infty$ for every $G'\in\mathcal S$, then $P_jG=\E_\mu G$.
\end{lemma}
\begin{proof}
Let $\abs{a^{(n)}}\to\infty$. Every subsequence has a further subsequence along which some coordinate, say $a^{(n)}_j$, diverges in modulus. Proposition~\ref{prop:D:subgroup} with $J=\{j\}$ gives $\E[G'U(a^{(n)})G]-\E[G'U(a^{(n)}_{J^c})P_jG]\to0$, and $\E[G'U(a_{J^c})P_jG]=\E[G']\E[G]$ because $P_jG$ is a constant. So every subsequence has a further subsequence along which the limit is $\E G'\,\E G$. Conversely, by Proposition~\ref{prop:D:subgroup} the limit of $\E[G'U(te_j)G]$ is $\E[G'P_jG]$; so $\E[G'(P_jG-\E G)]=0$ for every $G'\in\mathcal S$. Taking $G'=\overline{P_jG-\E G}\in\mathcal S$ gives $P_jG=\E G$.
\end{proof}

\subsubsection{The curvature sector, and curvature times Wick square}

The first lemma makes precise what a Gaussian domination of the linear field gives in general.

\begin{lemma}[domination of the linear field]\label{lem:os4:dom}
Let $\mu$ satisfy (H1)--(H2), and let $f\mapsto\Phi(f)\in L^2(\mu)$, $f\in C_c^\infty(\R^3;\cV)$ real, be linear with $U(a)\Phi(f)=\Phi(f(\cdot-a))$. Suppose $\E_\mu e^{\Phi(f)}\le e^{\frac12C_0(f,f)}$ for all real $f$, where $C_0(f,f)=\int\hat f(p)^*\varpi(dp)\hat f(p)$ for a tempered measure $\varpi$ on $\R^3$ with values in the positive semidefinite $4\times4$ matrices. If $\varpi(\{p:p_j=0\}\setminus\{0\})=0$, then $\E\Phi(f)=0$ and $P_j\Phi(f)=\E[\Phi(f)\mid\cI]$ for every $f$.
\end{lemma}
\begin{proof}
For real $t$, $\E e^{t\Phi(f)}\le e^{\frac12t^2C_0(f,f)}$; both sides equal $1$ at $t=0$ and are differentiable there (all exponential moments of $\Phi(f)$ are finite), so comparing derivatives at $t=0$ gives $\E\Phi(f)=0$ and $\E\Phi(f)^2\le C_0(f,f)$. Since $\varpi$ is tempered, $C_0(f,f)\le C_K\norm{f}_{C^k}^2$ for $\supp f$ in a fixed compact $K$ and some $k$; hence $f\mapsto\Phi(f)$ is continuous from $C_c^\infty(K)$ with the $C^k$ norm into $L^2(\mu)$. Put $f_T:=T^{-1}\int_0^Tf(\cdot-te_j)\,dt$ and $f^\square_T:=(2T)^{-3}\int_{[-T,T]^3}f(\cdot-a)\,da$, both in $C_c^\infty$. Riemann sums of these integrals converge in $C^k$ with supports in a fixed compact, so $\Phi(f_T)=T^{-1}\int_0^TU(te_j)\Phi(f)\,dt$ and $\Phi(f^\square_T)=(2T)^{-3}\int U(a)\Phi(f)\,da$ in $L^2(\mu)$, and by Lemma~\ref{lem:D:proj} they converge to $P_j\Phi(f)$ and to $P_{\{1,2,3\}}\Phi(f)=\E[\Phi(f)\mid\cI]$. Now $\widehat{f_T}(p)=m_T(p_j)\hat f(p)$ and $\widehat{f^\square_T}(p)=m^\square_T(p)\hat f(p)$ with $m_T(s)=T^{-1}\int_0^Te^{-its}dt$ and $m^\square_T(p)=\prod_k\frac{\sin Tp_k}{Tp_k}$. Hence
\[
\norm{P_j\Phi(f)-\E[\Phi(f)\mid\cI]}_2^2=\lim_{T\to\infty}\E\,\Phi(f_T-f^\square_T)^2\le\liminf_{T\to\infty}\int\abs{m_T(p_j)-m^\square_T(p)}^2\,\hat f(p)^*\varpi(dp)\hat f(p).
\]
The integrand is bounded by $4\hat f^*\varpi\hat f$, a finite measure, and tends to $0$ at $p=0$ (both multipliers equal $1$), at every $p$ with $p_j\ne0$ (both tend to $0$), and to $1$ on $\{p_j=0\}\setminus\{0\}$. By dominated convergence the limit is $\int_{\{p_j=0\}\setminus\{0\}}\hat f^*\varpi\hat f=0$.
\end{proof}

For the limit states the curvature satisfies more than Lemma~\ref{lem:os4:dom}: its domination holds conditionally on the matter, on every lattice torus.

\begin{lemma}[conditional curvature domination on the torus]\label{lem:os4:condlat}
Fix $(\varepsilon,T_j)$ and $\mu_0\in(0,\mu_*]\cup\{\fl\}$, and write $\ev{\cdot}$ for $\ev{\cdot}_{\varepsilon,T_j,\mu_0}$ and $\mathcal R$ for the $\sigma$-algebra generated by the radial variables $(r_x)_x$, $r_x:=\abs{\varphi(x)}$. For every real lattice 2-form $\bg$ and every bounded Borel function $Y\ge0$ of $(r_x)_x$,
\begin{equation}\label{eq:os4:condmgf}
\ev{e^{F(\bg)}Y}\le e^{\frac12V(\bg)}\ev{Y},\qquad V(\bg)\le\norm{\bg}_{2,\varepsilon}^2,
\end{equation}
with $V$ as in Lemma~\ref{lem:C-Fsector}. Consequently, for real lattice 2-forms $\bg_1,\dots,\bg_m$ and bounded real $\mathcal R$-measurable $Y_1,\dots,Y_m$,
\begin{equation}\label{eq:os4:condvar}
\Big\langle\Big(\sum_{i}Y_iF(\bg_i)\Big)^2\Big\rangle\le\sum_{i,k}\ev{Y_iY_k}\,\ip{\bg_i}{\bg_k}_\varepsilon ,\qquad \ev{Y_1F(\bg_1)}=0 .
\end{equation}
\end{lemma}
\begin{proof}
\emph{(\ref{eq:os4:condmgf}).} In Lemma~\ref{lem:C-current} replace the factor $e^{\XR(k)}$ by $Y$: since $Y$ depends on the radial variables only, the expansion in the proof of Lemma~\ref{lem:C-current} gives $\int\prod_xd^2\varphi(x)\,e^{-S^{\mathrm m}_b(A,\varphi)}\,Y=\sum_nw_Y(n)e^{ie\varepsilon\sum_bn_bA_b}$ with
\[
w_Y(n)=(2\pi)^{\abs{T_{j,\varepsilon}}}\int\prod_xv_x(r_x)\,r_x\,dr_x\;Y(r)\prod_bI_{n_b}(\varepsilon r_{b_-}r_{b_+})\ \ge0,\qquad w_Y(-n)=w_Y(n),
\]
where $v_x$ is the single-site weight of that proof with $k=0$, and with absolute convergence because $Y$ is bounded. The proof of Lemma~\ref{lem:C-Fsector} uses only these two properties of the weight (in the $\flat$ case together with the winding selection, which keeps them). It gives $\ev{e^{F(\bg)}Y}\le e^{\frac12V(\bg)}\ev{Y}$, and $V(\bg)\le\norm{\bg}^2_{2,\varepsilon}$ is the bound of that lemma.

\emph{(\ref{eq:os4:condvar}).} Let $\mathbf F:=(F(\bg_1),\dots,F(\bg_m))$ and let $\pi(r,dy)$ be a regular conditional law of $\mathbf F$ given $\mathcal R$. Put $Q(c):=V(\sum_ic_i\bg_i)$, a positive semidefinite quadratic form on $\R^m$. Applying \eqref{eq:os4:condmgf} to $\sum_ic_i\bg_i$ and all bounded $Y\ge0$ gives $\int e^{c\cdot y}\pi(r,dy)\le e^{\frac12Q(c)}$ for almost every $r$, first for all $c\in\Q^m$ off one null set, then for all $c\in\R^m$ by Fatou's lemma. Fix such an $r$ and $c$. The function $\phi(t)=\int e^{tc\cdot y}\pi(r,dy)$ is finite for all $t$, hence smooth, and $\phi(t)\le e^{\frac12t^2Q(c)}$ with equality at $t=0$. So $\phi'(0)=0$ and $\phi''(0)\le Q(c)$: the conditional mean of $\mathbf F$ vanishes and its conditional covariance matrix $\Sigma(r)$ satisfies $c^{\mathsf T}\Sigma(r)c\le Q(c)$ for all $c$. Given $\mathcal R$, the $Y_i$ are constants, so
\[
\Big\langle\Big(\sum_iY_iF(\bg_i)\Big)^2\Big\rangle=\Big\langle\sum_{i,k}Y_iY_k\Sigma_{ik}(r)\Big\rangle\le\Big\langle V\Big(\sum_iY_i\bg_i\Big)\Big\rangle\le\Big\langle\Big\|\sum_iY_i\bg_i\Big\|^2_{2,\varepsilon}\Big\rangle=\sum_{i,k}\ev{Y_iY_k}\ip{\bg_i}{\bg_k}_\varepsilon,
\]
and $\ev{Y_1F(\bg_1)}=\ev{Y_1\int y_1\pi(r,dy)}=0$.
\end{proof}

\begin{theorem}[curvature times Wick square: no invariant part]\label{thm:os4:FH}\label{prop:D:Fcluster}
Let $S$ and $\mu=\mu_S$ be as above. For every real $\bg\in C_c^\infty(\R^3;\Lambda^2)$ and every $Y\in L^2(\mathcal M)$, the product $F(\bg)Y$ lies in $L^2(\mu)$, with
\begin{equation}\label{eq:os4:FHnorm}
\E_\mu[F(\bg)^2\abs Y^2]\le\norm\bg_2^2\,\E_\mu\abs Y^2,\qquad \E_\mu[F(\bg)Y]=0 ,
\end{equation}
and $P_j\big(F(\bg)Y\big)=0$ for $j=1,2,3$. Hence every pair $(G',F(\bg)Y)$, $G'\in L^2(\mu)$, clusters along the axes:
\[
\lim_{\abs a\to\infty}\E_\mu\big[G'\,U(a)\big(F(\bg)Y\big)\big]=0=\E_\mu[G']\,\E_\mu[F(\bg)Y].
\]
In particular, for every polynomial $P$ in the smeared fields, $S\big(P\cdot F(T_a\bg)\big)\to0$ as $\abs a\to\infty$. On $\nu$-almost every ergodic component $\mu_\xi$ of Theorem~\ref{thm:D:decomp}, for all $\bg$ and all $\mathcal M_0$-measurable $Y$ with $F(\bg)Y\in L^2(\mu_\xi)$ simultaneously, where $\mathcal M_0:=\sigma(\XR(h):h\in C_c^\infty(\R^3)\text{ real})$ (not completed): $P^\xi_j(F(\bg)Y)=0$ for $j=1,2,3$, with $P^\xi_j$ the projection for $\mu_\xi$, $\E_{\mu_\xi}[F(\bg)Y]=0$, and every pair $(G',F(\bg)Y)$, $G'\in L^2(\mu_\xi)$, clusters along the axes for $\mu_\xi$. The bound in \eqref{eq:os4:FHnorm} is not asserted for the components.
\end{theorem}
\begin{proof}
\emph{Step 1: the inequality \eqref{eq:os4:condvar} in the limit.} Let $\bg_1,\dots,\bg_m\in C_c^\infty(\R^3;\Lambda^2)$ be real and $Y_i=\psi_i(\XR(h_{i1}),\dots,\XR(h_{iq}))$ with $\psi_i$ bounded continuous on $\R^q$ and $h_{il}\in C_c^\infty(\R^3)$ real. We claim
\begin{equation}\label{eq:os4:contvar}
\E_\mu\Big[\Big(\sum_iY_iF(\bg_i)\Big)^2\Big]\le\sum_{i,k}\E_\mu[Y_iY_k]\,\ip{\bg_i}{\bg_k}_{L^2},\qquad\E_\mu[Y_1F(\bg_1)]=0 .
\end{equation}
For $S=S^{\mu_0}$ or $S^\fl$, let $\ev{\cdot}_n$ be the lattice states of the defining subsequence and $Y_i^{(n)}$, $F_n(\bg_i)$ the lattice observables with sampled test functions; the $Y^{(n)}_i$ are bounded functions of $(r_x)_x$, so \eqref{eq:os4:condvar} holds for them. The characteristic functionals of the finite vectors $(F_n(\bg_i),\XR(h_{il}))_{i,l}$ converge to those under $\mu$ (Proposition~\ref{prop:C-compact}(b) at imaginary points, Theorem~\ref{thm:main}(ii)), so these vectors converge in law; the fourth moments of $F_n(\bg_i)$ are bounded uniformly in $n$ (Proposition~\ref{prop:C-bounds}), so the products $F_n(\bg_i)F_n(\bg_k)Y^{(n)}_iY^{(n)}_k$ are uniformly integrable. Hence both sides of \eqref{eq:os4:condvar} converge to the two sides of \eqref{eq:os4:contvar}, using $\ip{\bg_i}{\bg_k}_\varepsilon\to\ip{\bg_i}{\bg_k}_{L^2}$ (Riemann sums). For $S=S^0$, \eqref{eq:os4:contvar} holds for every $S^{\mu}$, $\mu\in Q_0$, and passes to $\mu\downarrow0$ in the same way, since $\mu_{S^\mu}\to\mu_{S^0}$ weakly with uniformly bounded moments (Lemma~\ref{lem:D:realize}).

\emph{Step 2: \eqref{eq:os4:FHnorm}.} Bounded continuous cylinder functions $Y$ of the Wick square are dense in $L^2(\mathcal M)$ (Lemma~\ref{lem:A:mc}). For them \eqref{eq:os4:contvar} with $m=1$ is \eqref{eq:os4:FHnorm}. For complex $Y$ apply this to $\rRe Y$ and $\operatorname{Im}Y$ and add. If $Y_n\to Y$ in $L^2(\mathcal M)$, then $F(\bg)Y_n$ is Cauchy in $L^2(\mu)$ by \eqref{eq:os4:FHnorm}, and along a subsequence $F(\bg)Y_n\to F(\bg)Y$ almost surely; so $F(\bg)Y\in L^2(\mu)$ and \eqref{eq:os4:FHnorm} holds for $Y$. The map $Y\mapsto F(\bg)Y$ is bounded from $L^2(\mathcal M)$ to $L^2(\mu)$, with norm at most $\norm\bg_2$.

\emph{Step 3: $P_j(F(\bg)Y)=0$.} By Step~2 and the boundedness of $P_j$ it suffices to take $Y$ a bounded continuous cylinder function. Put $\bg_t:=\bg(\cdot-te_j)$, $Y_t:=U(te_j)Y$ (again such a cylinder function, of the translated test functions) and $D:=\operatorname{diam}\supp\bg$. Then $U(te_j)(F(\bg)Y)=F(\bg_t)Y_t$, and $t\mapsto F(\bg_t)Y_t$ is continuous in $L^2(\mu)$ (Lemma~\ref{lem:D:strongcont}). So $\mathsf A_T:=T^{-1}\int_0^TU(te_j)(F(\bg)Y)\,dt$ is the $L^2$-limit of the Riemann sums $N^{-1}\sum_{i<N}F(\bg_{t_i})Y_{t_i}$, $t_i=iT/N$. By \eqref{eq:os4:contvar}, since $\ip{\bg_t}{\bg_s}=0$ for $\abs{t-s}>D$ and $\abs{\ip{\bg_t}{\bg_s}}\le\norm\bg_2^2$,
\[
\E\Big[\Big(\frac1N\sum_{i<N}F(\bg_{t_i})Y_{t_i}\Big)^2\Big]\le\frac{\norm Y_\infty^2}{N^2}\sum_{i,k<N}\abs{\ip{\bg_{t_i}}{\bg_{t_k}}}\le\norm Y_\infty^2\norm\bg_2^2\Big(\frac{2D}{T}+\frac1N\Big).
\]
Letting $N\to\infty$, $\norm{\mathsf A_T}_2^2\le2D\norm Y_\infty^2\norm\bg_2^2/T$ for real $Y$ (\eqref{eq:os4:contvar} is stated for real $Y_i$). For complex $Y$ apply this to $\rRe Y$ and $\operatorname{Im}Y$, whose one-sided means are the real and imaginary parts of $\mathsf A_T$; this gives $\norm{\mathsf A_T}_2^2\le4D\norm Y_\infty^2\norm\bg_2^2/T$, the bound used in Step~4. By Lemma~\ref{lem:D:proj}(i) (one-sided means), $\mathsf A_T\to P_j(F(\bg)Y)$; hence $P_j(F(\bg)Y)=0$. The clustering statement is Lemma~\ref{lem:os4:red}, with $\E[F(\bg)Y]=0$ from \eqref{eq:os4:FHnorm}. For the statement on polynomials take $Y=1$ and $G'=P$, which lies in $L^2(\mu)$ because all moments of $\mu$ are finite; $U(a)F(\bg)=F(T_a\bg)$.

\emph{Step 4: components.} Fix an even $r\ge14$ and take $s=r$, so that Lemma~\ref{lem:D:hermite}(b) and Corollary~\ref{cor:D:schwinger}(iv) hold with the same $r$. Let $\mathcal L_2$ be a countable set of real $\bg\in C_c^\infty$ dense in $\cS$ for the norm $\norm\cdot_r$, and $\mathcal L_Y$ a countable $\Q$-algebra of real bounded continuous cylinder functions of the Wick square, containing $1$ and generating $\mathcal M_0$. For $\bg\in\mathcal L_2$, $Y\in\mathcal L_Y$ and $T\in\N$, Step~3 and Theorem~\ref{thm:D:decomp}(a) give $\int\norm{\mathsf A_T}^2_{L^2(\mu_\xi)}d\nu(\xi)=\norm{\mathsf A_T}^2_{L^2(\mu)}\le cT^{-1}$, where $\mathsf A_T$ is the same function on $X$ and, for $\nu$-a.e.\ $\xi$, is the one-sided mean in $L^2(\mu_\xi)$ (the Riemann sums converge in $L^2(\mu)$, hence in $L^2(\mu_\xi)$ along a subsequence, for a.e.\ $\xi$). Along $T=k^2$, $\sum_k\int\norm{\mathsf A_{k^2}}^2_{L^2(\mu_\xi)}d\nu<\infty$, so $\norm{\mathsf A_{k^2}}_{L^2(\mu_\xi)}\to0$ for a.e.\ $\xi$, and $P^\xi_j(F(\bg)Y)=0$ there (Lemma~\ref{lem:D:proj}(i) for $\mu_\xi$). Off one $\nu$-null set this holds for all $\bg\in\mathcal L_2$, $Y\in\mathcal L_Y$ and $j$, and $\E_{\mu_\xi}M<\infty$ with $M$ as in Corollary~\ref{cor:D:schwinger}(iv). Fix such $\xi$. For fixed $Y\in\mathcal L_Y$, $\bg\mapsto F(\bg)Y$ is continuous into $L^2(\mu_\xi)$ for $\norm\cdot_r$, since $\abs{F(\bg)}\le M^{1/2}\norm\bg_r$ for all $\bg$ simultaneously off one $\mu_S$-null set, hence $\mu_\xi$-a.s.\ for $\nu$-a.e.\ $\xi$ (the identification of $F(\bg)$ with its Hermite series holds for all test functions simultaneously, by the density argument in the proof of Lemma~\ref{lem:D:hermite}(b)); so $P^\xi_j(F(\bg)Y)=0$ for all $\bg$. For fixed $\bg$, the set of bounded real $\mathcal M_0$-measurable $Y$ with $P^\xi_j(F(\bg)Y)=0$ is a vector space containing $\mathcal L_Y$ and closed under bounded $\mu_\xi$-a.e.\ convergence (dominated convergence in $L^2(\mu_\xi)$, as $F(\bg)\in L^2(\mu_\xi)$); by the monotone class theorem it contains every bounded real $\mathcal M_0$-measurable $Y$, and by linearity every bounded complex one. Then $\E_{\mu_\xi}[F(\bg)Y]=\E_{\mu_\xi}[P^\xi_j(F(\bg)Y)]=0$, because $P^\xi_j$ is self-adjoint and fixes $1$, and Lemma~\ref{lem:os4:red} for $\mu_\xi$ (which satisfies (H1)--(H3), Theorem~\ref{thm:D:decomp}) gives the clustering. Finally, let $Y$ be $\mathcal M_0$-measurable with $F(\bg)Y\in L^2(\mu_\xi)$, and put $Y_n:=Y\mathbf 1_{\{\abs Y\le n\}}$. Then $F(\bg)Y_n\to F(\bg)Y$ in $L^2(\mu_\xi)$ by dominated convergence, so $P^\xi_j(F(\bg)Y)=0$; as before $\E_{\mu_\xi}[F(\bg)Y]=0$ and Lemma~\ref{lem:os4:red} gives the clustering. By continuity of $P^\xi_j$ the same holds on the closed span of such products in $L^2(\mu_\xi)$. This applies to every polynomial $Y$ in the $\XR(h)$, since all moments of $\mu_\xi$ are finite (Corollary~\ref{cor:D:schwinger}(iv)).
\end{proof}

\begin{remark}\label{rem:D:Fcluster}
For $Y=1$ the conclusion on $\mu$ also follows from the tube bound of OS0 alone. Let $G(z)$ be the tube functional of the one-element family $\{\bg\}$ with no source for the Wick square. It is entire and real on $\R$, $G(0)=1$, $G'(0)=S_{1,0}(\bg)=0$ by OS3 (Theorem~\ref{thm:C-OS3}), $G''(0)=\E_\mu F(\bg)^2$ by Lemma~\ref{lem:D:realize}, and $G(t)\le e^{t^2\norm{\bg}_2^2/2}$ for real $t$ (Theorem~\ref{thm:C-OS0}(a)); comparing second derivatives at $0$ gives $\E_\mu F(\bg)^2\le\norm\bg_2^2$, so $\bg\mapsto F(\bg)$ extends to a bounded linear map from $L^2(\R^3;\Lambda^2)$ to $L^2(\mu)$. With $\bg_T:=T^{-1}\int_0^TT_{te_j}\bg\,dt$, Riemann sums, this bound and the strong continuity of $U$ give $F(\bg_T)=T^{-1}\int_0^TU(te_j)F(\bg)\,dt$, which tends to $P_jF(\bg)$ (Lemma~\ref{lem:D:proj}(i)); and
\[
\norm{P_jF(\bg)}^2_{L^2(\mu)}=\lim_{T\to\infty}\E_\mu F(\bg_T)^2\le\liminf_{T\to\infty}\int_{\R^3}\abs{\hat{\bg}(p)}^2\,\Big|\frac1T\int_0^Te^{-itp_j}dt\Big|^2\frac{dp}{(2\pi)^3}=0
\]
by dominated convergence, since $\{p_j=0\}$ is Lebesgue-null; with $\E_\mu F(\bg)=0$ (OS3), Lemma~\ref{lem:os4:red} gives the clustering. Theorem~\ref{thm:os4:FH} extends this, by a domination of the curvature conditional on the matter on every lattice torus, to all products $F(\bg)Y$ with $Y$ a square-integrable function of the Wick square. By itself the OS0 bound says nothing about the Wick square: its part for the Wick square is an $L^1$-type bound and gives no $L^2$ control of the invariant part; the Wick-square sector is treated below (Theorem~\ref{thm:os4:matter}, under hypothesis (L)). Observables of degree at least two in the curvature are not covered (Remark~\ref{rem:os4:Fopen}). Example~\ref{ex:D:axis} shows that for a scalar field without a domination bound the conclusion can fail, and Example~\ref{ex:os4:telegraph} that a domination bound on the linear field does not control quadratic observables. With full rotation invariance the distinction between diagonal and single-axis separation would disappear for every observable, since a rotation maps an axis to a diagonal direction.
\end{remark}

The domination of the linear field does not control observables of higher degree in the field. The next example satisfies (H1)--(H3), is invariant under translations and $\Oh$, and its linear field is dominated by white noise, which charges no hyperplane; yet a quadratic observable does not cluster along an axis.

\begin{example}[domination does not control quadratic observables]\label{ex:os4:telegraph}
Let $0<b_-<b_+\le1$, $\gamma>0$, and let $\alpha_1,\alpha_2,\alpha_3$ be independent copies of the stationary telegraph process: the right-continuous Markov process on $\{b_-,b_+\}$ that jumps from either value to the other at rate $\gamma$, started in its stationary law (uniform on $\{b_-,b_+\}$). Let $W$ be a white noise on $\R^3$, independent of the $\alpha_k$. Put $\beta(x):=\alpha_1(x_1)\alpha_2(x_2)\alpha_3(x_3)$ and $\Phi(f):=\int\beta f\,dW$ (a Wiener integral, conditionally on the $\alpha_k$), for real $f\in C_c^\infty(\R^3)$. Realise $\Phi$ on $\cV=\Lambda^2\R^3\oplus\R$ with vanishing 2-form component, as in Example~\ref{ex:D:axis}. Then:
\begin{enumerate}
\item[(i)] the law $\mu$ satisfies (H1)--(H3) with the reflections of Section~\ref{ss:discharge} and is invariant under $\Oh$;
\item[(ii)] $\E e^{\Phi(f)}\le e^{\frac12\norm f_2^2}$, so Lemma~\ref{lem:os4:dom} applies with $\varpi$ the Lebesgue measure, and every pair $(G',\Phi(f))$ clusters along the axes;
\item[(iii)] for $G:=\Phi(f)^2$, $f\ne0$: $\E[G\mid\cI]=c_2^3\norm f_2^2$ is constant, $c_2:=\E\alpha_1(0)^2$, but $P_1G=c_2\int\alpha_2(y_2)^2\alpha_3(y_3)^2f(y)^2\,dy$ is not; so $\lim_{t\to\infty}\E[G\,U(te_1)G]=\E[(P_1G)^2]>(\E G)^2$, and the pair $(G,G)$ does not cluster along $e_1$.
\end{enumerate}
\end{example}
\begin{proof}
\emph{Facts on $\alpha$.} The two-state chain with equal rates is reversible: its stationary law is uniform and satisfies detailed balance, so $(\alpha(-t))_t$ (made right-continuous) has the law of $(\alpha(t))_t$. For $u:=\alpha^2$, $\E u(0)=c_2$ and $\operatorname{Cov}(u(s),u(t))=v\,e^{-2\gamma\abs{t-s}}$ with $v:=\big((b_+^2-b_-^2)/2\big)^2>0$, since the probability that the chain is at the same value after time $t$ is $(1+e^{-2\gamma t})/2$. Hence $\operatorname{Var}\big(T^{-1}\int_0^Tu(y+t)\,dt\big)\le v/(\gamma T)\to0$, uniformly in $y$.

(i) Conditionally on $\alpha=(\alpha_k)_k$, $\Phi$ is a centred Gaussian linear process with $\E[\Phi(f)^2\mid\alpha]=\norm{\beta f}_2^2\le\norm f_2^2$. By the carrier argument of Lemma~\ref{lem:D:hermite}(b) (with $\E\Phi(f)^2\le\norm f_2^2\le\norm f_r^2$) $\Phi$ has a version with values in $\mathcal H^{-r}\subset X$. (H1)--(H2) hold as in the proof of Proposition~\ref{prop:discharge}. Invariance: translations act by shifting each $\alpha_k$ and $W$, whose laws are stationary; the coordinate reflections act by reversing one $\alpha_k$ and reflecting $W$, which preserves the laws (reversibility); permutations permute the identically distributed $\alpha_k$. On the scalar component the reflections of Section~\ref{ss:discharge} are the plain ones. (H3) for $j$, at $c=0$: let $\mathcal G$ be generated by $\alpha_k$ ($k\ne j$) and $\alpha_j(0)$, and $\mathcal F_\pm$ by $W$ restricted to $\{\pm x_j>0\}$ and $\alpha_j$ restricted to $\pm t\ge0$. A $\Sig(\Hp j0)$-measurable $F$ is a function of $\Phi(f)$, $\supp f\subset\{x_j>0\}$, hence $\mathcal F_+\vee\mathcal G$-measurable, and $F\circ\theta_j$ is $\mathcal F_-\vee\mathcal G$-measurable. Given $\mathcal G$, $\mathcal F_+$ and $\mathcal F_-$ are independent (independence of $W$ on disjoint sets, independence of $W$ and $\alpha$, the Markov property of $\alpha_j$ at time $0$). The map of the underlying variables that implements $\theta_j$ preserves their law and fixes every $\mathcal G$-measurable function ($\alpha_j$ is a.s.\ continuous at $0$), so $\E[F\circ\theta_j\mid\mathcal G]=\E[F\mid\mathcal G]$. Hence $\E[\overline{F\circ\theta_j}\,F]=\E\abs{\E[F\mid\mathcal G]}^2\ge0$.

(ii) $\E e^{\Phi(f)}=\E\,e^{\frac12\norm{\beta f}_2^2}\le e^{\frac12\norm f_2^2}$ since $\abs\beta\le b_+^3\le1$. So Lemma~\ref{lem:os4:dom} applies with $\varpi$ the Lebesgue measure and gives $P_j\Phi(f)=\E[\Phi(f)\mid\cI]$; since $\varpi$ has no atom at $0$, the computation in its proof with the multiplier $m^\square_T$ alone gives $\E[\Phi(f)\mid\cI]=0$. Lemma~\ref{lem:os4:red} gives the clustering.

(iii) Write $f_a:=f(\cdot-a)$ and $G_a:=U(a)G=\Phi(f_a)^2=\norm{\beta f_a}_2^2+R_a$, with $R_a:=\Phi(f_a)^2-\norm{\beta f_a}_2^2$. Conditionally on $\alpha$ the pair $(\Phi(f_a),\Phi(f_{a'}))$ is centred Gaussian, so $\E[R_aR_{a'}\mid\alpha]=2\ip{\beta f_a}{\beta f_{a'}}^2$, which vanishes when $\abs{a-a'}>D:=\operatorname{diam}\supp f$ and is at most $2\norm f_2^4$. Hence the one-sided means of $R_{te_1}$ and the cube means of $R_a$ tend to $0$ in $L^2$. Next, $\norm{\beta f_{te_1}}_2^2=\int u_1(y_1+t)u_2(y_2)u_3(y_3)f(y)^2\,dy$ with $u_k=\alpha_k^2$. By the facts on $\alpha$ and Minkowski's integral inequality, $T^{-1}\int_0^T\norm{\beta f_{te_1}}_2^2dt\to c_2\int u_2(y_2)u_3(y_3)f(y)^2dy=:A$ in $L^2$, and the cube means converge to $c_2^3\norm f_2^2$ (the cube mean factorises into the three one-dimensional means). By Lemma~\ref{lem:D:proj}, $P_1G=A$ and $\E[G\mid\cI]=P_{\{1,2,3\}}G=c_2^3\norm f_2^2$. Finally
\[
\operatorname{Var}A=c_2^2\iint f(y)^2f(y')^2\Big[\big(c_2^2+ve^{-2\gamma\abs{y_2-y_2'}}\big)\big(c_2^2+ve^{-2\gamma\abs{y_3-y_3'}}\big)-c_2^4\Big]dy\,dy'>0 ,
\]
since the bracket is positive and $f\ne0$. By Proposition~\ref{prop:D:subgroup}, $\E[G\,U(te_1)G]\to\E[GP_1G]=\E A^2>(\E A)^2=(\E G)^2$.
\end{proof}

\begin{remark}[what the curvature bound leaves open]\label{rem:os4:Fopen}
Theorem~\ref{thm:os4:FH} covers every observable of degree one in the curvature (the closed span of the products $F(\bg)Y$), with an arbitrary $L^2$ dependence on the Wick square; observables of degree zero, the Wick-square sector, are treated next (Theorem~\ref{thm:os4:matter}). It says nothing about $F(\bg)^2$, $\cos F(\bg)$ or $F(\bg)F(\bg')Y$. On the lattice, the conditional law of the curvature given the radial field is the free Gaussian field shifted by the imaginary field $ie\,d\,Cj_n$ of a random current $n$ (the representation behind Lemma~\ref{lem:C-current}); the conditional covariance of the curvature is the free one minus $e^2$ times a covariance of these currents. The $e_j$-invariant part of a quadratic observable is therefore governed by the clustering of the current sector, on which no estimate is available, and Example~\ref{ex:os4:telegraph} shows that the domination \eqref{eq:os4:condmgf} cannot replace such an estimate (Section~\ref{p:ss:OS4open}).
\end{remark}

\subsubsection{The Wick-square sector}

\begin{lemma}[the radial field is associated]\label{lem:os4:fkg}
Fix $(\varepsilon,T_j)$ and $\mu_0\in(0,\mu_*]$. The law of $(r_x)_{x\in T_{j,\varepsilon}}$ under $\ev{\cdot}_{\varepsilon,T_j,\mu_0}$ is associated: $\ev{\phi\psi}\ge\ev\phi\ev\psi$ for all coordinatewise non-decreasing functions $\phi,\psi$ of $(r_x)_x$ with $\ev{\phi^2},\ev{\psi^2}<\infty$.
\end{lemma}
\begin{proof}
The proof combines \cite[Thm.~6.1, p.~276]{BFS1} with the FKG inequality for densities on products of totally ordered spaces \cite[Sect.~4.4]{GJ}, \cite{Pre74}. By the proof of Lemma~\ref{lem:C-current} and Fubini's theorem, the law of $r=(r_x)_x$ has the density $p(r)=Z(r)\prod_xr_xv_x(r_x)$ (up to normalisation, with $v_x$ as there at $k=0$) on $(0,\infty)^{T_{j,\varepsilon}}$, where
\[
Z(r)=\int\mu_C(dA)\int\prod_xd\vartheta_x\,\prod_b\exp\big(J_b\cos u_b\big),\qquad J_b:=\varepsilon r_{b_-}r_{b_+},\quad u_b:=\vartheta_{b_+}-\vartheta_{b_-}+\varepsilon eA_b ,
\]
and $C=C_{\mu_0}>0$. $Z$ is smooth and positive. Let $\ev{\cdot}_r$ be the probability measure on $(A,\vartheta)$ with density $\prod_be^{J_b\cos u_b}/Z(r)$ with respect to $\mu_C\otimes d\vartheta$. For $x\ne y$,
\[
\partial_{r_x}\partial_{r_y}\log Z=\sum_b(\partial_{r_x}\partial_{r_y}J_b)\ev{\cos u_b}_r+\sum_{b,b'}(\partial_{r_x}J_b)(\partial_{r_y}J_{b'})\big(\ev{\cos u_b\cos u_{b'}}_r-\ev{\cos u_b}_r\ev{\cos u_{b'}}_r\big).
\]
All derivatives of the $J_b$ are $\ge0$. Also $\ev{\cos u_b}_r\ge0$: expand every factor $e^{J\cos u}=\sum_\nu I_\nu(J)e^{i\nu u}$ with $I_\nu\ge0$; after the angular integration and the Gaussian integral $\E_C e^{ie\varepsilon\sum_bk_bA_b}=e^{-\frac12e^2\varepsilon^2\ip{k}{Ck}}>0$, $Z(r)\ev{e^{iu_b}}_r$ is a sum of nonnegative terms. The truncated correlation of $\cos u_b$ and $\cos u_{b'}$ is $\ge0$ by \cite[Thm.~6.1]{BFS1} on the torus graph, wrap-around bonds included (\cite{BFS1} allow periodic conditions on rectangles, and the theorem holds on every finite graph: Lemma~\ref{app:G1:ginibre} and Proposition~\ref{app:G1:BFStorus}), applied with the radial measures $d\rho_x=\delta_{r_x}$ (finite measures with all Gaussian moments, as allowed there), hopping coefficients $\varepsilon$, the Gaussian measure taken to be the law of the bond variables $(e\varepsilon A_b)_b$, which is centred with covariance $e^2\varepsilon^2C_{\mu_0}>0$ (the bond variable of \cite[Sect.~6.1]{BFS1} already contains this factor), $P=Q=1$, $n,m$ the coboundaries of the bonds $b,b'$ and $f,g$ the corresponding bond indicators. So $\partial_{r_x}\partial_{r_y}\log p\ge0$ for $x\ne y$. Let $u:=r\wedge r'$, $a:=r-u$, $b:=r'-u$. Then $a,b\ge0$ have disjoint supports and $r\vee r'=u+a+b$, so
\[
\log p(r\vee r')-\log p(r)-\log p(r')+\log p(r\wedge r')=\int_0^1\!\!\int_0^1\sum_{x\in\supp a,\ y\in\supp b}a_xb_y\,\partial_{r_x}\partial_{r_y}\log p(u+sa+tb)\,ds\,dt\ge0 ,
\]
since only $x\ne y$ occurs. For bounded non-decreasing $\phi$ and $\psi$, with $\psi>0$ after adding a constant, the densities $f_2:=p$ and $f_1:=\psi p/\int\psi p$ satisfy $f_1(r\vee r')f_2(r\wedge r')\ge f_1(r)f_2(r')$. Preston's theorem \cite[Thm.~3, p.~234]{Pre74}, applied on $(0,\infty)^{T_{j,\varepsilon}}$ with Lebesgue measure, gives $\int\phi f_1\ge\int\phi f_2$, that is, $\ev{\phi\psi}\ge\ev\phi\ev\psi$. Truncation $\phi\mapsto\max(-M,\min(M,\phi))$ and dominated convergence extend this to $L^2$.
\end{proof}

\begin{lemma}[a covariance inequality for associated variables]\label{lem:os4:newman}
Let $Z$ be an associated random vector and $Y_1,\dots,Y_m$, $Y'_1,\dots,Y'_{m'}$ coordinatewise non-decreasing functions of $Z$ in $L^2$. Let $\phi\in C^1(\R^m)$, $\psi\in C^1(\R^{m'})$ have bounded partial derivatives, $a_i:=\norm{\partial_i\phi}_\infty$, $b_l:=\norm{\partial_l\psi}_\infty$. Then
\[
\big\lvert\operatorname{Cov}\big(\phi(Y),\psi(Y')\big)\big\rvert\le\sum_{i,l}a_ib_l\operatorname{Cov}(Y_i,Y'_l).
\]
\end{lemma}
\noindent The inequality is Newman's \cite{New80}; we include the short proof.

\begin{proof}
The functions $\sum_ia_iy_i\pm\phi(y)$ and $\sum_lb_ly'_l\pm\psi(y')$ are non-decreasing in each variable, so, composed with $Y$ and $Y'$, they are non-decreasing functions of $Z$ in $L^2$. Association gives $\operatorname{Cov}(\sum a_iY_i+\epsilon\phi(Y),\sum b_lY'_l+\epsilon'\psi(Y'))\ge0$ for all four sign choices $\epsilon,\epsilon'\in\{\pm1\}$. Adding the choices $(+,-)$ and $(-,+)$ gives $\operatorname{Cov}(\phi(Y),\psi(Y'))\le\sum a_ib_l\operatorname{Cov}(Y_i,Y'_l)$; adding $(+,+)$ and $(-,-)$ gives the bound for $-\operatorname{Cov}(\phi(Y),\psi(Y'))$.
\end{proof}

\begin{theorem}[the Wick-square sector reduces to its linear field]\label{thm:os4:matter}
Let $S$ and $\mu=\mu_S$ be as above. Assume
\begin{equation}\label{eq:os4:L}
\text{(L)}\qquad P_j\XR(h)=\E_\mu\XR(h)\quad\text{for all real }h\in C_c^\infty(\R^3)\text{ and }j=1,2,3 .
\end{equation}
Then $P_jG=\E_\mu G$ for every $G\in L^2(\mathcal M)$ and every $j$, and every pair $(G',G)$ with $G'\in L^2(\mu)$, $G\in L^2(\mathcal M)$ clusters along the axes. In particular the restriction of $\mu$ to $\mathcal M$ is mixing, and for $\nu$-a.e.\ ergodic component, $\mu_\xi=\mu$ on $\mathcal M_0$.
\end{theorem}
\begin{proof}
\emph{Step 1: the linear field.} Let $h,h'\in C_c^\infty$. By (L) and Lemma~\ref{lem:os4:red}, $\operatorname{Cov}_\mu(\XR(h'),U(a)\XR(h))\to0$ as $\abs a\to\infty$ in any manner.

\emph{Step 2: a covariance bound in the limit.} Let $h_1,\dots,h_m,h'_1,\dots,h'_{m'}\in C_c^\infty$ be nonnegative, $\phi,\psi$ as in Lemma~\ref{lem:os4:newman} and also bounded, and $a\in\R^3$. We claim
\begin{equation}\label{eq:os4:newmanlim}
\big\lvert\operatorname{Cov}_\mu\big(\psi(\XR(h'_1),\dots),\,\phi(\XR(h_1(\cdot-a)),\dots)\big)\big\rvert\le\sum_{i,l}a_ib_l\operatorname{Cov}_\mu\big(\XR(h'_l),\XR(h_i(\cdot-a))\big).
\end{equation}
For $\mu_0\in(0,\mu_*]$ and every lattice in the defining sequence, $\XR(k)=\sum_x\varepsilon^3k(x)(r_x^2-W_\varepsilon-z_\varepsilon)$ is a non-decreasing function of $(r_x)_x$ when $k\ge0$, so Lemmas~\ref{lem:os4:fkg} and~\ref{lem:os4:newman} give \eqref{eq:os4:newmanlim} for the lattice state, with the sampled test functions. At fixed lattice it passes to $\mu_0\downarrow0$ by Lemma~\ref{lem:C-flatlimit} (both sides are elementary observables), so it holds for $\flat$ too. Along the defining subsequence the left side converges (convergence in law of finite vectors of smeared fields, as in Step~1 of the proof of Theorem~\ref{thm:os4:FH}), and the right side converges by Proposition~\ref{prop:C-compact}(a). For $S^0$, pass to $\mu\downarrow0$ in $Q_0$ as there.

\emph{Step 3: cylinder functions.} Steps 1--2 give $\operatorname{Cov}_\mu(G',U(a)G)\to0$ as $\abs a\to\infty$ for $G=\phi(\XR(h_1),\dots,\XR(h_m))$, $G'=\psi(\XR(h'_1),\dots)$ with nonnegative $h_i,h'_l$, because $U(a)G=\phi(\XR(h_1(\cdot-a)),\dots)$. Every real $h\in C_c^\infty$ is $h=(h+c\chi)-c\chi$ with $\chi\in C_c^\infty$, $0\le\chi$, $\chi=1$ on $\supp h$ and $c=\norm h_\infty$, so the functions $e^{i\sum_kt_k\XR(h_k)}$ with $h_k\ge0$ include all $e^{i\XR(h)}$; their real and imaginary parts are bounded with bounded derivatives. Their linear span is dense in $L^2(\mathcal M)$ (Lemma~\ref{lem:A:mc}). Since $\abs{\operatorname{Cov}(G',U(a)G)-\operatorname{Cov}(G'_n,U(a)G_n)}\le2(\norm{G'-G'_n}_2\norm G_2+\norm{G'_n}_2\norm{G-G_n}_2)$ uniformly in $a$, every pair in $L^2(\mathcal M)$ clusters along the axes.

\emph{Step 4: invariant parts and general $G'$.} For $G\in L^2(\mathcal M)$, $P_jG\in L^2(\mathcal M)$ (the means stay in this closed subspace), and Step~3 and the converse part of Lemma~\ref{lem:os4:red} with $\mathcal S=L^2(\mathcal M)$ give $P_jG=\E G$. Lemma~\ref{lem:os4:red} then gives clustering against every $G'\in L^2(\mu)$. Mixing on $\mathcal M$ is Step~3. For the components: $\E[G\mid\cI]=P_{\{1,2,3\}}P_jG=\E_\mu G$ (Lemma~\ref{lem:D:proj}(ii)), so $\E_{\mu_\xi}G=\E_\mu G$ for $\nu$-a.e.\ $\xi$ by Theorem~\ref{thm:D:decomp}(a); taking $G$ in a countable algebra of bounded continuous cylinder functions generating $\mathcal M_0$ gives $\mu_\xi=\mu$ on $\mathcal M_0$.
\end{proof}

Hypothesis (L) is a statement about the axis-invariant part of the linear Wick-square field. The torus pressures of the Wick square bound that part in $L^\infty$, and remove it when they have no corner at the origin.

\begin{definition}[pressures of the Wick square; latent density]\label{def:os4:pressure}
For a lattice state $\ev{\cdot}_n=\ev{\cdot}_{\varepsilon_n,T_n,\mu_0}$ of the defining sequence of $S\in\{S^{\mu_0},S^\fl\}$ put $\alpha_n(\sigma):=\abs{T_n}^{-1}\log\ev{e^{\sigma\XR(1_{T_n})}}_n$, $\sigma\in\R$. Let $\mathfrak P(S)$ be the set of pointwise limits on $I:=(-H/2,\sigma_+/2)$ of subsequences of $(\alpha_n)$; here $\sigma_+$ and $H=m^2$ are the constants of Theorem~\ref{thm:main}(i), and this constant $H$ is unrelated to the functions $Y\in L^2(\mathcal M)$ of Theorem~\ref{thm:os4:FH}. For $S^0$, let $\mathfrak P(S^0)$ be the set of pointwise limits on $I$ of sequences $\alpha^{(i)}\in\mathfrak P(S^{\mu_i})$, $\mu_i\in Q_0$, $\mu_i\downarrow0$ (all $S^{\mu_i}$ along the common subsequence of Theorem~\ref{thm:main}(ii)). For $\alpha\in\mathfrak P(S)$ let $\Delta(\alpha):=\alpha'_+(0)-\alpha'_-(0)$, and $\Delta_S:=\inf_{\alpha\in\mathfrak P(S)}\Delta(\alpha)$, the \emph{latent density} of the Wick square in $S$. Since the limits are not known to be unique, $\mathfrak P(S)$ and $\Delta_S$ are defined through the defining sequence of $S$ (the subsequence of Theorem~\ref{thm:main}(ii)) and may depend on it; ``the latent density in $S$'' always refers to this choice.
\end{definition}

\begin{proposition}[the invariant part of the Wick square]\label{prop:os4:latent}
$\mathfrak P(S)$ is not empty. Every $\alpha\in\mathfrak P(S)$ is convex on $I$, $\alpha(0)=0$, $\abs{\alpha(\sigma)}\le\mathsf K\abs\sigma$, and $0\le\Delta(\alpha)\le\kappa_++2S^B_\fl$, with $\kappa_+$, $S^B_\fl$ and $\mathsf K=\max(\kappa_+,2S^B_\fl)$ as in Theorem~\ref{thm:ch-T9}. For every $\alpha\in\mathfrak P(S)$, every nonnegative $h\in C_c^\infty(\R^3)$ and every $j$,
\begin{equation}\label{eq:os4:latent}
\alpha'_-(0)\norm h_1\le P_j\XR(h)\le\alpha'_+(0)\norm h_1\quad\mu\text{-a.s.},\qquad \alpha'_-(0)\norm h_1\le\E_\mu\XR(h)\le\alpha'_+(0)\norm h_1 .
\end{equation}
Hence $\abs{P_j\XR(h)-\E_\mu\XR(h)}\le\Delta_S\norm h_1$ a.s.\ for $h\ge0$, and (L) holds if $\Delta_S=0$, in particular if some $\alpha\in\mathfrak P(S)$ is differentiable at $0$.
\end{proposition}
\begin{proof}
\emph{Properties of $\alpha_n$.} $\alpha_n$ is finite (by property (P2) of Section~\ref{ssec:C-standing}) and convex (a logarithmic moment generating function), $\alpha_n(0)=0$ and $\alpha_n'(0)=\ev{r_0^2-W_\varepsilon-z_\varepsilon}_n$ by translation invariance (Lemma~\ref{lem:C-sym}). Proposition~\ref{prop:ch-T7} with $q\equiv\sigma$ gives $\alpha_n(\sigma)\le\kappa_+\sigma$ for $0\le\sigma\le\sigma_+$, and Proposition~\ref{prop:ch-Tflat} with $h\equiv p$ gives $\alpha_n(-p)\le2S^B_\fl p$ for $0\le p\le H$, for every $\mu_0\in(0,\mu_*]\cup\{\fl\}$. By convexity $\alpha_n(\sigma)\ge\alpha'_n(0)\sigma$, and evaluating at $\sigma=\sigma_+$ and $\sigma=-H$ gives $-2S^B_\fl\le\alpha_n'(0)\le\kappa_+$. So $\abs{\alpha_n(\sigma)}\le\mathsf K\abs\sigma$ on $[-H,\sigma_+]$, and the $\alpha_n$ are uniformly Lipschitz on $\bar I$ (a convex function bounded by $M$ on an interval is Lipschitz with constant $2M/\delta$ at distance $\delta$ from its ends). By the Arzel\`a--Ascoli theorem a subsequence converges uniformly on $\bar I$, so $\mathfrak P(S)\ne\emptyset$; for $S^0$ choose $\alpha^{(i)}\in\mathfrak P(S^{\mu_i})$, which obey the same bounds, and apply the Arzel\`a--Ascoli theorem once more. Convexity, $\alpha(0)=0$ and the bounds pass to pointwise limits, and give $\alpha'_+(0)\le\kappa_+$, $\alpha'_-(0)\ge-2S^B_\fl$.

\emph{The chessboard bound.} Fix $n$ and a real lattice function $q$. The state $\ev{\cdot}_n$ is reflection positive about the planes halfway between lattice planes (Proposition~\ref{prop:C-RP}(a) for $\mu_0>0$, (b) for $\flat$), invariant under these reflections (Lemma~\ref{lem:C-sym}), and the number of sites per side is even; so the chessboard estimate (Section~\ref{ssec:UV-other}, item~2) applies to the positive site factors $\exp(\varepsilon^3q(x)(r_x^2-W_\varepsilon-z_\varepsilon))$, exactly as in step~(2) of the proof of Proposition~\ref{prop:ch-T7}, but at the photon mass of the state:
\begin{equation}\label{eq:os4:chess}
\ev{e^{\XR(q)}}_n\le\exp\Big(\sum_x\varepsilon^3\alpha_n(q(x))\Big).
\end{equation}

\emph{Upper bound in \eqref{eq:os4:latent}.} Let $S\in\{S^{\mu_0},S^\fl\}$, $\alpha=\lim\alpha_n$ along a subsequence $\Lambda$, $h\ge0$, $s>0$, and $h_T:=T^{-1}\int_0^Th(\cdot-te_j)\,dt\in C_c^\infty$. Then $\norm{h_T}_\infty\le T^{-1}\sup_{x}\int_\R h(x+te_j)\,dt\to0$, and $\norm{h_T}_1=\norm h_1$. Fix $T$ so large that $\sigma_T:=s\norm{h_T}_\infty<\sigma_+/2$. For $0\le\sigma\le\sigma_T$ convexity and $\alpha_n(0)=0$ give $\alpha_n(\sigma)\le\sigma\,\alpha_n(\sigma_T)/\sigma_T$. With $q=sh_T$ (sampled) in \eqref{eq:os4:chess},
\[
\log\ev{e^{s\XR(h_T)}}_n\le s\sum_x\varepsilon^3h_T(x)\,\frac{\alpha_n(\sigma_T)}{\sigma_T}\ \longrightarrow\ s\norm h_1\frac{\alpha(\sigma_T)}{\sigma_T}\qquad(n\to\infty,\ n\in\Lambda).
\]
The law of $\XR(h_T)$ under $\ev{\cdot}_n$ converges to that of $\XR(h_T)$ under $\mu$ (characteristic functionals, Theorem~\ref{thm:main}(ii)); since $y\mapsto\min(e^{sy},L)$ is bounded and continuous for every $L>0$, $\E_\mu e^{s\XR(h_T)}\le\liminf_n\ev{e^{s\XR(h_T)}}_n\le\exp(s\norm h_1\alpha(\sigma_T)/\sigma_T)$. Next, $\XR(h_T)=T^{-1}\int_0^TU(te_j)\XR(h)\,dt$ in $L^2(\mu)$ (Riemann sums of $h_T$ converge uniformly with supports in a fixed compact, and $\E\XR(k)^2\le C'N(k)^2$ by Theorem~\ref{thm:main}(iii)), so $\XR(h_T)\to P_j\XR(h)$ in $L^2(\mu)$ (Lemma~\ref{lem:D:proj}(i)), and almost surely along a sequence $T_k\to\infty$. By Fatou's lemma and $\alpha(\sigma)/\sigma\downarrow\alpha'_+(0)$ as $\sigma\downarrow0$,
\[
\E_\mu e^{sP_j\XR(h)}\le\liminf_k\E_\mu e^{s\XR(h_{T_k})}\le e^{s\norm h_1\alpha'_+(0)}\qquad(s>0).
\]
Chebyshev's inequality gives $\mu(P_j\XR(h)>\alpha'_+(0)\norm h_1+\delta)\le e^{-s\delta}$ for every $s>0$ and $\delta>0$, hence the upper bound.

\emph{Lower bound.} The same argument with $q=-sh_T$, using $\alpha_n(-\sigma)\le\sigma\,\alpha_n(-\sigma_T)/\sigma_T$ for $0\le\sigma\le\sigma_T<H/2$ and $\alpha(-\sigma)/\sigma\downarrow-\alpha'_-(0)$, gives $\E_\mu e^{-sP_j\XR(h)}\le e^{-s\norm h_1\alpha'_-(0)}$, hence $P_j\XR(h)\ge\alpha'_-(0)\norm h_1$.

\emph{The mean.} $\E_\mu\XR(h)=\lim_{n\in\Lambda}\alpha'_n(0)\sum_x\varepsilon^3h(x)$, and for $\sigma\in(0,\sigma_+/2)$, $\alpha'_n(0)\le\alpha_n(\sigma)/\sigma\to\alpha(\sigma)/\sigma$ and $\alpha'_n(0)\ge\alpha_n(-\sigma)/(-\sigma)\to\alpha(-\sigma)/(-\sigma)$; letting $\sigma\downarrow0$ gives the second part of \eqref{eq:os4:latent}.

\emph{The state $S^0$.} Let $\alpha=\lim_i\alpha^{(i)}$ as in Definition~\ref{def:os4:pressure}. The bound $\E_{\mu_{S^{\mu_i}}}e^{s\XR(h_T)}\le\exp(s\norm h_1\alpha^{(i)}(\sigma_T)/\sigma_T)$ holds for each $i$ by the above, and passes to $i\to\infty$ by the weak convergence $\mu_{S^{\mu_i}}\to\mu_{S^0}$ (Lemma~\ref{lem:D:realize}) and the same truncation; the rest is unchanged, the mean included ($\E_{S^{\mu_i}}\XR(h)/\norm h_1\le\alpha^{(i)}(\sigma)/\sigma\to\alpha(\sigma)/\sigma$).

\emph{Conclusion.} For $h\ge0$ both $P_j\XR(h)$ and $\E\XR(h)$ lie in $[\alpha'_-(0),\alpha'_+(0)]\norm h_1$ for every $\alpha\in\mathfrak P(S)$, so they differ by at most $\Delta_S\norm h_1$. If $\Delta_S=0$, then $P_j\XR(h)=\E\XR(h)$ for $h\ge0$, and for every real $h\in C_c^\infty$ by the decomposition $h=(h+c\chi)-c\chi$ of the proof of Theorem~\ref{thm:os4:matter}.
\end{proof}

\begin{remark}[what (L) needs]\label{rem:os4:Lopen}
$\Delta(\alpha)>0$ means that the limit pressure has a corner at the reference value of the bare coefficient: a first-order transition in the mass, with the density $\E\,\XR$ jumping; one-loop perturbation theory predicts such a transition on a surface inside the coupling range of Theorem~\ref{thm:main} (Remark~\ref{rem:disc:firstorder}). A corner is excluded if the torus susceptibilities $\alpha_n''(\sigma)=\abs{T_n}^{-1}\operatorname{Var}_{n,\sigma}\XR(1_{T_n})$ in the tilted states are bounded uniformly in $n$ for $\abs\sigma$ small, which by Lemma~\ref{lem:os4:fkg} is a uniform integrability statement for the nonnegative truncated two-point function of the Wick square; the estimates of this paper are linear in $\sigma$ and give no such bound (Section~\ref{p:ss:OS4open}).
\end{remark}

\subsubsection{The full algebra under a differentiability hypothesis}

For the whole algebra we give a criterion in terms of the state alone. Fix a length $R>0$, the block $Q:=(0,R)^3$ and, for $b\in\Z^3$, the block $Q+Rb$. Assume, in addition to (H1)--(H3),
\begin{itemize}
\item[(H4)] $\theta_j\theta_k=\theta_k\theta_j$ for $j\ne k$.
\end{itemize}
It holds for the laws of the limit states, since the push-forwards $\rho_{j*}$ of Section~\ref{ss:discharge} commute. Put $\Theta_{j,c}G:=G\circ\theta_{j,c}$. From \eqref{eq:D:commute}, $\Theta_{j,c}=U(2ce_j)\Theta_j$, hence $\Theta_{j,c}\Theta_{j,c'}=U(2(c-c')e_j)$, $\Theta_{j,c}U(te_j)=U(-te_j)\Theta_{j,c}$, and $\Theta_{j,c}$ commutes with $U(y)$, $y\perp e_j$, and (by (H4)) with $\Theta_{k,c'}$, $k\ne j$. For $m\in\Z$ let $\vartheta^{(j)}_{2m}:=U(2mRe_j)$ and $\vartheta^{(j)}_{2m+1}:=U(2mRe_j)\Theta_{j,R}$, and for $b\in\Z^3$ let $\vartheta_b:=\vartheta^{(1)}_{b_1}\vartheta^{(2)}_{b_2}\vartheta^{(3)}_{b_3}$ (the factors commute). By (H2), $\vartheta_b$ maps $\Sig(Q)$-measurable functions to $\Sig(Q+Rb)$-measurable ones, and it preserves $\mu$. For a bounded real $\Sig(Q)$-measurable $\Gamma$ and $B_k:=\{0,\dots,2^k-1\}^3$ define the \emph{block pressure}
\[
\Pi_\Gamma(u):=\liminf_{k\to\infty}\abs{B_k}^{-1}\log\E_\mu\exp\Big(u\sum_{b\in B_k}\vartheta_b\Gamma\Big),\qquad u\in\R .
\]
By Jensen's inequality and the invariance of $\mu$, $\Pi_\Gamma(u)\ge u\,\E_\mu\Gamma$, and $\abs{\Pi_\Gamma(u)}\le\abs u\norm\Gamma_\infty$. We say that $\Gamma$ satisfies (D$_\Gamma$) if $\Pi_\Gamma(u)=u\,\E_\mu\Gamma+o(u)$ as $u\to0$.

\begin{theorem}[single-axis clustering from differentiable block pressures]\label{thm:os4:block}
Let $\mu$ satisfy (H1)--(H4), let $\Gamma$ be bounded, real and $\Sig(Q)$-measurable, and assume (D$_\Gamma$). Then $P_j\Gamma=\E_\mu\Gamma$ for $j=1,2,3$, and $(G',\Gamma)$ clusters along the axes for every $G'\in L^2(\mu)$.
\end{theorem}
\noindent Steps 1--2 of the proof are a variant, in the present setting, of the multiple-reflection bound of \cite[Prop.~10.5.3, p.~239]{GJ}.
\begin{proof}
\emph{Step 1: one doubling.} Let $W\ge0$ be bounded and $\Sig(U)$-measurable with $U\subset\Hm jc$. Then $\Theta_{j,c}W\in L^2(\Sig(\Hp jc))$ by (H2), and by (H3) and the Schwarz inequality for the positive form $\ip FG_c:=\E[\overline{\Theta_{j,c}F}\,G]$ on $L^2(\Sig(\Hp jc))$,
\[
(\E W)^2=\ip1{\Theta_{j,c}W}_c^2\le\ip11_c\,\ip{\Theta_{j,c}W}{\Theta_{j,c}W}_c=\E\big[W\cdot\Theta_{j,c}W\big].
\]

\emph{Step 2: from a line to a cube.} For a finite $B\subset\Z^3$ put $Z_B:=\prod_{b\in B}\vartheta_b(e^{u\Gamma})$. Let $N=2^p$ and $L_N:=\{0,\dots,N-1\}\times\{0\}\times\{0\}$. If $B$ is a box $\{0,\dots,2^{q_1}-1\}\times\{0,\dots,2^{q_2}-1\}\times\{0,\dots,2^{q_3}-1\}$, then $Z_B$ is measurable for the open box of side lengths $2^{q_i}R$, which lies in $\Hm j{2^{q_j}R}$, and $\Theta_{j,2^{q_j}R}Z_B=Z_{B'}$, where $B'$ is $B$ reflected in the plane $x_j=2^{q_j}R$ of block coordinates, i.e.\ $b_j\mapsto2^{q_j+1}-1-b_j$. Indeed, for $c=MR$ and $l=2m$, $\Theta_{j,c}\vartheta^{(j)}_l=U(-2mRe_j)\Theta_{j,c}=U(2(M-1-m)Re_j)\Theta_{j,R}=\vartheta^{(j)}_{2M-1-l}$, and for $l=2m+1$, $\Theta_{j,c}\vartheta^{(j)}_l=U(-2mRe_j)\Theta_{j,c}\Theta_{j,R}=U(2(M-1-m)Re_j)=\vartheta^{(j)}_{2M-1-l}$; the other factors of $\vartheta_b$ commute with $\Theta_{j,c}$. So Step~1 gives $\E Z_B\le(\E Z_{B\cup B'})^{1/2}$, and $B\cup B'$ is the box with $q_j$ replaced by $q_j+1$. Starting from $L_N$ and doubling $k-p$ times in direction $1$ and $k$ times in directions $2$ and $3$ ($k\ge p$),
\begin{equation}\label{eq:os4:line}
\log\E\exp\Big(u\sum_{l<N}\vartheta^{(1)}_l\Gamma\Big)=\log\E Z_{L_N}\le\frac N{\abs{B_k}}\log\E Z_{B_k},\qquad\text{hence}\qquad\log\E\exp\Big(u\sum_{l<N}\vartheta^{(1)}_l\Gamma\Big)\le N\,\Pi_\Gamma(u).
\end{equation}

\emph{Step 3: the line mean.} Put $\tilde\Gamma:=\Theta_{1,R}\Gamma$ and $\mathsf A_N:=N^{-1}\sum_{l<N}\vartheta^{(1)}_l\Gamma=N^{-1}\sum_{m<N/2}U(2mRe_1)(\Gamma+\tilde\Gamma)$. By the mean ergodic theorem for the unitary $U(2Re_1)$, $\mathsf A_N\to\mathsf A:=\frac12P'(\Gamma+\tilde\Gamma)$ in $L^2$, where $P'$ is the projection onto the $U(2Re_1)$-invariant vectors; and $\E\mathsf A_N=\E\Gamma$. By \eqref{eq:os4:line} with $u=s/N$ and (D$_\Gamma$), for every $s\in\R$, $\limsup_N\log\E e^{s\mathsf A_N}\le\limsup_NN\Pi_\Gamma(s/N)=s\E\Gamma$. Along a subsequence $\mathsf A_N\to\mathsf A$ a.s., and Fatou's lemma gives $\E e^{s(\mathsf A-\E\Gamma)}\le1$ for all $s\in\R$. Chebyshev's inequality, with $s\to\pm\infty$, gives $\mathsf A=\E\Gamma$ a.s.

\emph{Step 4: conclusion.} The range of $P_1$ is contained in that of $P'$, so $P_1P'=P_1$. $P_1$ commutes with $U$ and $\Theta_1$, and $\Theta_1$ fixes the range of $P_1$ (Lemma~\ref{lem:D:invloc}(ii)), so $P_1\tilde\Gamma=P_1U(2Re_1)\Theta_1\Gamma=P_1\Gamma$. Hence $\E\Gamma=P_1\mathsf A=\frac12(P_1\Gamma+P_1\tilde\Gamma)=P_1\Gamma$. The cube $B_k$ and $\Pi_\Gamma$ do not depend on the direction, so the same argument with lines along $e_2$, $e_3$ gives $P_j\Gamma=\E\Gamma$ for all $j$, and Lemma~\ref{lem:os4:red} gives the clustering.
\end{proof}

\begin{corollary}[OS4 for the limit state under (D)]\label{cor:os4:OS4}
Let $S$, $\mu=\mu_S$ be as above, and assume
\begin{itemize}
\item[(D)] for every $R>0$ and every real $f\in C_c^\infty(Q;\cV)$, $Q=(0,R)^3$, the functions $\Gamma=\cos\Phi(f)$ and $\Gamma=\sin\Phi(f)$ satisfy (D$_\Gamma$).
\end{itemize}
Then $P_j$ is the projection onto the constants for $j=1,2,3$: $\mu$ is ergodic, $\mu$ satisfies OS4 (Definition~\ref{def:OS4}), and every pair $(G',G)$ in $L^2(\mu)$ clusters along the axes. The same holds for the law $\mu^H_S$ on $\mathcal H^{-r}$, which is the only probability measure on $\mathcal H^{-r}$ with the moments $S$ (Lemma~\ref{lem:D:hermite}); so $S$ satisfies OS4. In particular the Schwinger functions satisfy $\lim_{\abs a\to\infty}S\big(P'\cdot P(\cdot-a)\big)=S(P')\,S(P)$ for all polynomials $P,P'$ in the smeared fields.
\end{corollary}
\begin{proof}
For real $f\in C_c^\infty(\R^3;\cV)$ choose $a$ and $R$ with $\supp f(\cdot-a)\subset Q$; Theorem~\ref{thm:os4:block} gives $P_je^{i\Phi(f(\cdot-a))}=\E e^{i\Phi(f(\cdot-a))}$, and translating back ($P_j$ commutes with $U(a)$, and $\mu$ is invariant) $P_je^{i\Phi(f)}=\E e^{i\Phi(f)}$. The span of the $e^{i\Phi(f)}$ is dense in $L^2(\mu)$ (Lemma~\ref{lem:A:mc}), so $P_jG=\E G$ for all $G\in L^2(\mu)$; this is (OS4$_j$), and in particular $\E[G\mid\cI]=P_{\{1,2,3\}}G=P_{\{1,2,3\}}P_1G=\E G$, so $\mu$ is ergodic. Lemma~\ref{lem:os4:red} gives the clustering; polynomials lie in $L^2(\mu)$. The map $\iota$ of Lemma~\ref{lem:D:hermite}(b) is a Borel isomorphism onto a set of full measure that intertwines the translations, so the statements carry over to $\mu^H_S$.
\end{proof}

\begin{remark}[on hypothesis (D)]\label{rem:os4:D}
If $\mu=\frac12(\mu_1+\mu_2)$ with $\E_{\mu_1}\Gamma\ne\E_{\mu_2}\Gamma$ and both $\mu_i$ invariant under the $\vartheta_b$, then $\log\E_\mu e^{uN_k}\ge\log\E_{\mu_i}e^{uN_k}-\log2$ (with $N_k:=\sum_{b\in B_k}\vartheta_b\Gamma$) and Jensen's inequality in $\mu_i$ give $\Pi_\Gamma(u)\ge\max_iu\,\E_{\mu_i}\Gamma$, so (D$_\Gamma$) fails: (D) excludes in particular a non-trivial ergodic decomposition, and it is a uniqueness-type statement. Theorem~\ref{thm:os4:block} applies to each component $\mu_\xi$ with its own block pressure, but (D) for a component is no easier to verify than for $\mu$. For the linear Wick-square field the same mechanism, run on the tori instead of in the limit state, is Proposition~\ref{prop:os4:latent}, where the corresponding pressure is controlled by the uniform estimates of this paper up to its corner at the origin.
\end{remark}

\subsubsection{Summary: which form of clustering holds for which state}

\begin{theorem}[clustering for the limit states and their components]\label{thm:os4:summary}
Let $S\in\{S^{\mu_0},S^\fl,S^0\}$ be a limit state of Theorem~\ref{thm:main}, $\mu=\mu_S$, $\mu=\int\mu_\xi\,d\nu$ its ergodic decomposition, and $\mathcal K_1$ the closed span in $L^2(\mu)$ of the products $F(\bg)Y$, $\bg\in C_c^\infty(\R^3;\Lambda^2)$ real, $Y\in L^2(\mathcal M)$.
\begin{enumerate}
\item[(a)] For every $G\in\mathcal K_1$ and $G'\in L^2(\mu)$, $\E_\mu[G'U(a)G]\to0=\E_\mu G'\,\E_\mu G$ as $\abs a\to\infty$ in any manner; the same holds on $\nu$-a.e.\ component for $G=F(\bg)Y$, $Y$ bounded and $\mathcal M_0$-measurable.
\item[(b)] For $h\ge0$, $\abs{P_j\XR(h)-\E\XR(h)}\le\Delta_S\norm h_1$ with $\Delta_S\le\kappa_++2S^B_\fl$. If $\Delta_S=0$, then (a) holds for every $G$ in the closed span of $L^2(\mathcal M)\cup\mathcal K_1$, with limit $\E G'\,\E G$, and all components have the same law on $\mathcal M_0$.
\item[(c)] If (D) holds, then $\mu$ is ergodic and satisfies OS4, and every pair $(G',G)$ in $L^2(\mu)$ clusters along the axes.
\item[(d)] On $\nu$-a.e.\ component: Theorem~\ref{thm:cluster}(c)--(e), that is, (C1), (C2) and a unique joint vacuum in each coordinate direction. In the reconstruction of Proposition~\ref{prop:D:recon} from $\mu_\xi$, the vectors $[F(\bg)Y]$, $F(\bg)Y\in L^2_+$ with $Y$ bounded and $\mathcal M_0$-measurable, are orthogonal to $\Ker H_j$; if $\Delta_S=0$, so are the vectors $[G-\E_{\mu_\xi}G]$, $G\in L^2_+$ bounded and $\mathcal M_0$-measurable.
\end{enumerate}
\end{theorem}
\begin{proof}
(a) For $G$ in the linear span of the products $F(\bg)Y$, $P_jG=0=\E_\mu G$ by Theorem~\ref{thm:os4:FH}. Both sides are continuous on $L^2(\mu)$, so this holds on $\mathcal K_1$, and Lemma~\ref{lem:os4:red} gives the clustering. On the components the limit is $\E_{\mu_\xi}G'\,\E_{\mu_\xi}G=0$ (Theorem~\ref{thm:os4:FH}). (b) Proposition~\ref{prop:os4:latent} and Theorem~\ref{thm:os4:matter}; for linear combinations use that $P_j$ is linear. (c) Corollary~\ref{cor:os4:OS4}. (d) Theorem~\ref{thm:cluster}; for the zero-energy statement, Proposition~\ref{prop:D:fixed}(i) for $\mu_\xi$ identifies $P_{\Ker H_j}[G]$ with $[P^\xi_jG]$, which vanishes for $G=F(\bg)Y$ (Theorem~\ref{thm:os4:FH}, components) and equals $\E_{\mu_\xi}G\,\Omega_j$ for bounded $\mathcal M_0$-measurable $G$ when $\Delta_S=0$ (Proposition~\ref{prop:os4:latent} and Theorem~\ref{thm:os4:matter}, components). For the last point: $\mu_\xi=\mu$ on $\mathcal M_0$ (Theorem~\ref{thm:os4:matter}); the map $(t,x)\mapsto G(\tau_{-te_j}x)$ is $\mathcal B(\R)\otimes\mathcal M_0$-measurable (for bounded continuous cylinder functions of the Wick square it is continuous in $t$ and $\mathcal M_0$-measurable in $x$; the monotone class theorem gives the rest), so the one-sided means $\mathsf A_T:=T^{-1}\int_0^TU(te_j)G\,dt$ are $\mathcal M_0$-measurable and $\norm{\mathsf A_T-\E_\mu G}_{L^2(\mu_\xi)}=\norm{\mathsf A_T-\E_\mu G}_{L^2(\mu)}\to\norm{P_jG-\E_\mu G}_{L^2(\mu)}=0$; hence $P^\xi_jG=\E_\mu G=\E_{\mu_\xi}G$ (Lemma~\ref{lem:D:proj}(i) for $\mu_\xi$).
\end{proof}
\noindent The theorem does not decide whether $\Ker H_j=\C\Omega_j$ on all of $\cH_j$ for a component; by Proposition~\ref{add:OS4rel}(a) below this is (OS4$_j$) for the component (Section~\ref{p:ss:OS4open}).

The forms of the cluster property are related as follows. Here (OS4$_j$) and (OS4$'_j$) are the conditions of Definition~\ref{def:OS4}: ergodicity of the translations $\{\tau_{te_j}\}_{t\in\R}$, and clustering along $e_j$; the weak form of OS4 is ergodicity of $\mu$ under $\R^3$ (Definition~\ref{add:OS4}). Under (H1)--(H3), (OS4$_j$) and (OS4$'_j$) are equivalent by Lemma~\ref{lem:os4:red} and Lemma~\ref{lem:D:proj}(ii).

\begin{proposition}[relations between the forms of OS4]\label{add:OS4rel}
Let $\mu$ satisfy (H1)--(H3), and let $j\in\{1,2,3\}$.
\begin{enumerate}
\item[(a)] (OS4$_j$) holds if and only if, in the reconstruction of Proposition~\ref{prop:D:recon} in direction $j$, $\Ker H_j=\C\Omega_j$.
\item[(b)] $\mu$ is ergodic under $\R^3$ if and only if, in the reconstruction in some (equivalently, every) coordinate direction $j$, the joint fixed space of $\{\mathrm{e}^{-tH_j}\}_{t\ge0}$ and $\{V_j(y)\}_{y\perp e_j}$ is $\C\Omega_j$; if and only if (C2) holds in the form $\E_\mu[G'\,U(a)G]\to\E_\mu[G']\,\E_\mu[G]$ as $\min_k\abs{a_k}\to\infty$, for all $G,G'\in L^2(\mu)$; and if and only if the Ces\`aro means over cubes of $\E_\mu[G'\,U(a)G]$ converge to $\E_\mu[G']\,\E_\mu[G]$ for all $G,G'\in L^2(\mu)$.
\item[(c)] (OS4$_j$) for one $j$ implies ergodicity under $\R^3$. Let $X$ be a space of distributions as in Section~\ref{ss:discharge} and $\mathcal P_j^+$ the algebra of polynomials in the fields $\Phi(f)$, $f\in C_c^\infty(\Hp j0;\cV)$ real. If $\mathcal P_j^+$ is dense in $L^2(\Sig(\Hp j0),\mu)$ and the cluster property in the form of \cite[E4, p.~88]{OS73} holds, $\lim_{\lambda\to\infty}\E_\mu[\overline{\Theta_jF}\,U(\lambda a)G]=\E_\mu[\overline F]\,\E_\mu[G]$ for all $F,G\in\mathcal P_j^+$ and $a\perp e_j$, then $\mu$ is ergodic under $\R^3$.
\item[(d)] Let $\mu=\int\mu_\xi\,d\nu(\xi)$ be the decomposition of Theorem~\ref{thm:D:decomp}. Then $\nu$-almost every $\mu_\xi$ satisfies (OS4$_j$) if and only if $P_jG=\E[G\mid\cI]$ for every $G\in L^2(\mu)$.
\end{enumerate}
\end{proposition}
\begin{proof}
(a) Let $G\in L^2_+=L^2(\Sig(\Hp j0),\mu)$. By Proposition~\ref{prop:D:fixed}(i), $P_{\Ker H_j}[G]=[P_jG]$ and $\norm{[W]}_{\cH_j}^2=\E_\mu\abs W^2$ for every $e_j$-invariant $W\in L^2_+$. If (OS4$_j$) holds, $P_jG=\E_\mu G$ by Lemma~\ref{lem:D:proj}(ii), so $P_{\Ker H_j}[G]=\E_\mu G\,\Omega_j$; since $[L^2_+]$ is dense in $\cH_j$, $\Ker H_j=\C\Omega_j$. Conversely, if $\Ker H_j=\C\Omega_j$, then $[P_jG]=c\,\Omega_j$ with $c=\ip{\Omega_j}{[P_jG]}=\E_\mu P_jG=\E_\mu G$, and $W:=P_jG-\E_\mu G$ is $e_j$-invariant with $[W]=0$, so $\E_\mu\abs W^2=0$. Thus $P_jG=\E_\mu G$ for $G\in L^2_+$. For $G\in L^2(\Sig(K))$, $K$ bounded, $U(se_j)G\in L^2_+$ for $s$ large and $P_j$ commutes with $U(se_j)$, so $P_jG=\E_\mu G$; these $G$ are dense in $L^2(\mu)$ by (H2). Hence $P_j$ is the projection onto the constants, and for a $\tau_{te_j}$-invariant Borel set $B$, $\mathbf 1_B=P_j\mathbf 1_B=\mu(B)$, so $\mu(B)\in\{0,1\}$.

(b) If $\mu$ is ergodic, $\E[\,\cdot\mid\cI]=\E_\mu$, and the three statements are Proposition~\ref{prop:D:fixed}(iii), Theorem~\ref{thm:D:diag} and Theorem~\ref{thm:D:mean}. Conversely, if $B$ is an $\R^3$-invariant Borel set, then $G=G'=\mathbf 1_B$ satisfies $\E_\mu[G'U(a)G]=\mu(B)$ for every $a$, so either clustering statement gives $\mu(B)=\mu(B)^2$; the joint-fixed-space statement gives the same by the converse in the proof of Theorem~\ref{thm:D:recon}(iv), which uses only Lemma~\ref{lem:D:invloc} and Proposition~\ref{prop:D:fixed}(i).

(c) An $\R^3$-invariant set is invariant under the $\tau_{te_j}$, which gives the first claim. For the second, let $W$ be bounded and $\R^3$-invariant. By Lemma~\ref{lem:D:invloc}(i)--(ii), $W\in L^2_+$ and $\Theta_jW=W$. Approximate $W$ in $L^2_+$ by $G_n\in\mathcal P_j^+$; then $\abs{\E[\overline{\Theta_jW}U(\lambda a)W]-\E[\overline{\Theta_jG_n}U(\lambda a)G_n]}\le\norm{W-G_n}_2(\norm W_2+\norm{G_n}_2)$ uniformly in $\lambda$, and the cluster property gives, for any $a\perp e_j$, $a\ne0$,
\[
\E_\mu\abs W^2=\lim_{\lambda\to\infty}\E_\mu\big[\overline{\Theta_jW}\,U(\lambda a)W\big]=\abs{\E_\mu W}^2 .
\]
So $W$ is constant.

(d) Let $G\in L^2(\mu)$ and $Y_T:=T^{-1}\int_0^TU(te_j)G\,dt$. By Lemma~\ref{lem:D:proj}(i), $Y_T\to P_jG$ in $L^2(\mu)$, and $\int\norm{Y_T-P_jG}^2_{L^2(\mu_\xi)}\,d\nu(\xi)=\norm{Y_T-P_jG}^2_{L^2(\mu)}$ by Theorem~\ref{thm:D:decomp}(a). Along a sequence $T_k\to\infty$, therefore, $Y_{T_k}\to P_jG$ in $L^2(\mu_\xi)$ for $\nu$-a.e.\ $\xi$; since $\mu_\xi$ is translation invariant, Lemma~\ref{lem:D:proj}(i) for $\mu_\xi$ also gives $Y_{T_k}\to P^\xi_jG$, the projection for $\mu_\xi$. So $P_jG=P^\xi_jG$ $\mu_\xi$-a.s.\ for $\nu$-a.e.\ $\xi$, and likewise $\E[G\mid\cI]=\E_{\mu_\xi}G$ $\mu_\xi$-a.s.\ (Theorem~\ref{thm:D:decomp}(a)). Take $G$ in a countable set dense in every $L^2(\mu_\xi)$ (Lemma~\ref{lem:A:mc}); off one $\nu$-null set, $P_jG=\E[G\mid\cI]$ $\mu$-a.s.\ for all such $G$ if and only if $P^\xi_jG=\E_{\mu_\xi}G$ for all such $G$ and $\nu$-a.e.\ $\xi$, and both sides are continuous in $G$. The second statement is (OS4$_j$) for $\mu_\xi$ by Lemma~\ref{lem:D:proj}(ii) for $\mu_\xi$, as in the end of the proof of (a).
\end{proof}

\begin{remark}\label{add:OS4rot}
Without rotations the forms are not equivalent (Example~\ref{ex:D:axis}: the law there is ergodic under $\R^3$, but $P_1\Phi(f)\ne\E\Phi(f)$). With invariance under the full Euclidean group, ergodicity under $\R^3$ implies (OS4$_j$), (OS4$'_j$) and the cluster property of \cite[E4, p.~88]{OS73} for every time direction (Corollary~\ref{p:cor:rotcor}), so the forms are then equivalent. In Theorem~\ref{thm:cluster} the weak form is proved for the ergodic components; Theorem~\ref{thm:os4main} proves (OS4$_j$) and (OS4$'_j$) for part of the observables.
\end{remark}

\subsection{The limit states: proof of \texorpdfstring{Theorem~\ref{thm:cluster}}{the weak clustering theorem}}\label{ss:proofcluster}

\begin{corollary}[Theorem~\ref{thm:cluster}]\label{cor:D:schwinger}
Let $S$ be one of the limit states of Theorem~\ref{thm:main}, and $\mu_S$ its law on $X$ (Lemma~\ref{lem:D:realize}). Then $\mu_S=\int\mu_\xi\,d\nu(\xi)$, where for $\nu$-a.e.\ $\xi$:
\begin{enumerate}
\item[(i)] $\mu_\xi$ is translation invariant and ergodic, invariant under $\theta_1,\theta_2,\theta_3$, and reflection positive about every coordinate plane;
\item[(ii)] (C1) Theorem~\ref{thm:D:mean} holds on $\mu_\xi$; (C2) $\E_{\mu_\xi}[G'U(a)G]\to\E_{\mu_\xi}[G']\,\E_{\mu_\xi}[G]$ as $\min_j\abs{a_j}\to\infty$, for all $G,G'\in L^2(\mu_\xi)$, in particular for polynomials in the smeared fields;
\item[(iii)] in each coordinate direction $j$, the reconstruction of Proposition~\ref{prop:D:recon} on $\mu_\xi$ has a unique joint vacuum: the joint fixed space of the $e^{-tH_j}$, $t\ge0$, and of the transverse translations is $\C\Omega_j$;
\item[(iv)] the Schwinger functions $S^\xi_N(f_1,\dots,f_N):=\E_{\mu_\xi}\prod_i\Phi(f_i)$ exist for all Schwartz test functions and satisfy $\abs{S^\xi_N(f_1,\dots,f_N)}\le\sigma^\xi_N\prod_i\abs{f_i}_k$ for a fixed Schwartz norm $\abs{\cdot}_k$ and constants with $\sigma^\xi_N\le\alpha_\beta(\xi)(N!)^\beta$ for every $\beta>1$; the constants depend on $\xi$.
\end{enumerate}
The decomposition is determined by $S$: $\mu_S$ is determined by $S$ (Proposition~\ref{prop:discharge}), and its ergodic decomposition is unique in the following sense. If $\mu_S=\int m_\eta\,d\lambda(\eta)$, with $\eta\mapsto m_\eta(B)$ measurable for every Borel $B$ and $m_\eta$ translation invariant and ergodic for $\lambda$-a.e.\ $\eta$, then the images of $\nu$ under $\xi\mapsto\mu_\xi$ and of $\lambda$ under $\eta\mapsto m_\eta$ agree on the $\sigma$-algebra of the space of probability measures on $X$ generated by the maps $m\mapsto m(B)$.
\end{corollary}
\begin{proof}
(i)--(iii) are Theorem~\ref{thm:D:decomp}, Theorems~\ref{thm:D:mean}--\ref{thm:D:diag} and Proposition~\ref{prop:D:fixed}(iii) applied to $\mu_S$, which satisfies (H1)--(H3) by Proposition~\ref{prop:discharge}; polynomials in the smeared fields lie in $L^2(\mu_\xi)$ for a.e.\ $\xi$ by Theorem~\ref{thm:D:decomp}(e) and the finiteness of all moments of $\mu_S$.

(iv) Fix an even $r\ge14$. By the uniqueness in Lemma~\ref{lem:D:realize}, the laws $\mu_S$ on $H^{-s}_{\mathrm{loc}}$ for different $s>3$ are images of one another under the continuous inclusions $H^{-s}_{\mathrm{loc}}\subset H^{-s'}_{\mathrm{loc}}$ ($s<s'$), and their ergodic decompositions correspond; we may therefore take $s=r$. Let $Z_{\alpha,c}$, $M=\sum_{\alpha,c}\lambda_\alpha^{-r}Z_{\alpha,c}^2$ and $\Psi$ be as in the proof of Lemma~\ref{lem:D:hermite}(b), and let $(f_n)$ be a countable subset of $C_c^\infty(\R^3;\cV)$ containing, for every $R$, a subset of $C_c^\infty(B_R;\cV)$ dense in the $H^r$-norm. The set
\[
B:=\{M<\infty\}\cap\{\Psi(f_n)=\Phi(f_n)\ \text{for all }n\}
\]
is Borel, $\mu_S(B)=1$, and on $B$ one has $\Psi(f)=\Phi(f)$ for every $f\in C_c^\infty(\R^3;\cV)$, since both sides are continuous on $C_c^\infty(B_R;\cV)$ for the $H^r$-norm. For Schwartz $f$ put $\Phi(f):=\Psi(f)$ on $B$; then $\abs{\Phi(f)}\le M^{1/2}\norm f_r$. By Theorem~\ref{thm:C-OS0}(c), each $h_{\alpha,c}$ being supported in one summand of $\cV$, and Lemma~\ref{lem:D:hermite}(a) ($\abs f_4\le C\norm f_{10}$),
\[
\E_{\mu_S}Z_{\alpha,c}^{2n}\le c_SC_0^{2n}(2n)!\,\abs{h_{\alpha,c}}_4^{2n}\le c_S(CC_0)^{2n}(2n)!\,\lambda_\alpha^{10n},
\]
with $c_S$ and $C_0$ the constants of Theorem~\ref{thm:C-OS0}(c). With $w:=4\sum_\alpha\lambda_\alpha^{10-r}<\infty$ and Jensen's inequality for the probability weights $\lambda_\alpha^{10-r}/w$,
\[
\E_{\mu_S}M^n\le w^n\sup_{\alpha,c}\lambda_\alpha^{-10n}\,\E_{\mu_S}Z_{\alpha,c}^{2n}\le c_SC_1^n(2n)!,\qquad C_1:=w(CC_0)^2 .
\]
Put $M_n(\xi):=\E_{\mu_\xi}M^n$. Since $\int\mu_\xi(B)\,d\nu=1$ and $\int M_n\,d\nu=\E_{\mu_S}M^n$, Markov's inequality and the Borel--Cantelli lemma give, for $\nu$-a.e.\ $\xi$: $\mu_\xi(B)=1$, $M_n(\xi)<\infty$ for all $n$, and $M_n(\xi)\le n^2c_SC_1^n(2n)!$ for $n\ge n_0(\xi)$. For such $\xi$ and all Schwartz $f_i$, using $\E_{\mu_\xi}M^{N/2}\le M_N(\xi)^{1/2}$,
\[
\abs{S^\xi_N(f_1,\dots,f_N)}\le M_N(\xi)^{1/2}\prod_i\norm{f_i}_r .
\]
Finally let $\abs f_k:=\sup_{x,\abs\gamma\le k}(1+\abs x^2)^{k/2}\abs{D^\gamma f(x)}$. By Lemma~\ref{lem:D:hermite}(a), $\norm f_r\le C_r\sum_{\abs\delta+\abs\gamma\le r}\norm{x^\delta\partial^\gamma f}_{L^2}$, and $\norm{x^\delta\partial^\gamma f}_{L^2}\le\norm{(1+\abs x^2)^{-1}}_{L^2}\sup_x(1+\abs x^2)^{1+\abs\delta/2}\abs{\partial^\gamma f(x)}\le\pi\abs f_{r+2}$. Hence $\norm f_r\le C'\abs f_{r+2}$, and (iv) holds with $k=r+2$ and
\[
\sigma^\xi_N:=c(\xi)\,N\,c_S^{1/2}C_1^{N/2}C'^N((2N)!)^{1/2}\le\alpha_\beta(\xi)(N!)^\beta\quad(\beta>1),
\]
where $c(\xi)\ge1$ absorbs the finitely many $N<n_0(\xi)$ and we used $((2N)!)^{1/2}\le2^NN!$.

Uniqueness of the decomposition: let $\mu_S=\int m_\eta\,d\lambda(\eta)$ be as in the statement, and let $\cD$, $A_T$ and $q_G$ be as in the proof of Theorem~\ref{thm:D:decomp}. As there, $\int m_\eta(\abs{A_TG-q_G}^2)\,d\lambda=\E_{\mu_S}\abs{A_TG-q_G}^2\to0$. Under an ergodic $m_\eta$ the fixed vectors are constants (the proof of Lemma~\ref{lem:D:proj}(ii) gives a strictly invariant version), so $A_TG\to m_\eta(G)$ in $L^2(m_\eta)$. Along a subsequence, therefore, $\mu_x(G)=q_G(x)=m_\eta(G)$ for $m_\eta$-a.e.\ $x$, for all $G\in\cD$ and $\lambda$-a.e.\ $\eta$; that is, $\mu_x=m_\eta$ for $m_\eta$-a.e.\ $x$. Hence $\mu_S\{x:\mu_x\in\mathcal A\}=\lambda\{\eta:m_\eta\in\mathcal A\}$ for every set $\mathcal A$ in the $\sigma$-algebra generated by the maps $m\mapsto m(B)$. Moreover $\mu_S$ is determined by $S$ (Proposition~\ref{prop:discharge}).
\end{proof}
\noindent Theorem~\ref{thm:cluster} is Corollary~\ref{cor:D:schwinger}, carried to $\mathcal H^{-r}$ by Lemma~\ref{lem:D:hermite}; for $\mu$ itself, parts (c) and (d) of Theorem~\ref{thm:cluster} are Theorems~\ref{thm:D:mean} and~\ref{thm:D:diag}, and the uniqueness of $\mu$ is Lemma~\ref{lem:D:hermite}(c). Theorem~\ref{thm:os4main} is Theorem~\ref{thm:os4:FH} together with Theorem~\ref{thm:os4:summary}(a)--(c) and Proposition~\ref{prop:os4:latent}.

The obstruction of Proposition~\ref{prop:axisobs} disappears as soon as the law is invariant under all rotations: reflection positivity then holds about every plane, every invariant function is fixed by every reflection, and a rotation carries a coordinate axis to a direction off the coordinate planes, along which Theorem~\ref{thm:D:diag} applies. We state this as a conditional statement; it is not used elsewhere. For a unit vector $n\in\R^3$ and $c\in\R$ put $\Hp{n}{c}:=\{x:x\cdot n>c\}$, $\Hm{n}{c}:=\{x:x\cdot n<c\}$, and let $\rho_n:=1-2nn^{\mathsf T}$ be the reflection in the plane $n^\perp$.

\begin{corollary}[the cluster property under rotation invariance]\label{p:cor:rotcor}\label{add:rotcond}
For a unit vector $n\in\R^3$ and $c\in\R$ let $\Hp nc:=\{x:x\cdot n>c\}$ and $\Hm nc:=\{x:x\cdot n<c\}$, and let $\rho_n:=1-2nn^{\mathsf T}$ be the reflection in the plane $n^\perp$; for $n=e_j$ these are the half spaces and the reflection $\rho_j$ of Section~\ref{ss:abstract}.
Let $\mu$ satisfy (H1)--(H3). Assume that $\Ot$ acts on $X$ by homeomorphisms $x\mapsto Rx$, continuously in $R$ for each $x$, such that $\rho_jx=\theta_jx$ for the coordinate reflections, $R\,\tau_a\,R^{-1}=\tau_{Ra}$, $\Gamma_RG:=G\circ R^{-1}$ maps $\Sig(U)$-measurable functions to $\Sig(RU)$-measurable ones, and $\mu$ is invariant under every $R\in\Ot$. Let $\mu=\int\mu_\xi\,d\nu(\xi)$ be the decomposition of Theorem~\ref{thm:D:decomp}. Then for $\nu$-almost every $\xi$:
\begin{enumerate}
\item[(a)] $\mu_\xi$ is invariant under the Euclidean group $\R^3\rtimes\Ot$;
\item[(b)] $\E_{\mu_\xi}[\overline{F\circ\theta_{n,c}}\,F]\ge0$ for every unit $n$, every $c\in\R$ and every $F\in L^2(\Sig(\Hp{n}{c}),\mu_\xi)$, where $\theta_{n,c}:=\tau_{cn}\rho_n\tau_{-cn}$;
\item[(c)] for all $G,G'\in L^2(\mu_\xi)$, $\E_{\mu_\xi}[G'\,U(a)G]\to\E_{\mu_\xi}[G']\,\E_{\mu_\xi}[G]$ as $\abs a\to\infty$ in any manner.
\end{enumerate}
Consequently $\mu_\xi$ satisfies (OS4$_j$) and (OS4$'_j$) of Definition~\ref{def:OS4} for $j=1,2,3$ and OS4 in each of the forms of Definition~\ref{add:OS4}, for every choice of the time direction (also for non-coordinate directions $n$, with the half spaces $\Hp{n}{0}$); and in the reconstruction from $\mu_\xi$ in any direction $n$, $\Ker H_n=\C\Omega_n$.
\end{corollary}
\begin{proof}
\emph{Step 1: reflection positivity about every plane.} Given a unit vector $n$, choose $R\in\SO(3)$ with $Re_1=n$. Then $\rho_n=R\rho_1R^{-1}$ and $R\,\Hp10=\Hp{n}{0}$. Let $F\in L^2(\Sig(\Hp{n}{0}),\mu)$ and $F':=F\circ R$, which is $\Sig(\Hp10)$-measurable. Since $\Gamma$ is a representation, $F\circ\rho_n=\Gamma_R\Theta_1F'$, and by $R$-invariance of $\mu$ and (H3), $\E_\mu[\overline{F\circ\rho_n}\,F]=\E_\mu[\Gamma_R(\overline{\Theta_1F'}\,F')]=\E_\mu[\overline{\Theta_1F'}\,F']\ge0$. For $c\ne0$ conjugate by $\tau_{cn}$, as after \eqref{eq:D:commute}. Moreover $\rho_n\tau_a=\tau_{\rho_na}\rho_n$, and $\Gamma_{\rho_n}$ maps $\Sig(U)$-measurable functions to $\Sig(\rho_nU)$-measurable ones. So $\mu$ satisfies (H1)--(H3) with the coordinate frame replaced by any orthonormal frame containing $n$.

\emph{Step 2: invariant functions are fixed by every reflection.} Let $W\in L^2(\mu)$ be invariant under $U(tn)$, $t\in\R$. The proof of Lemma~\ref{lem:D:invloc}(i)--(ii) applies verbatim with $e_j$ replaced by $n$, the half spaces $\Hp jc$, $\Hm jc$ by $\Hp nc$, $\Hm nc$ and $\Theta_j$ by $\Gamma_{\rho_n}$, using Step~1: $\Gamma_{\rho_n}W=W$. A translation-invariant $W$ is invariant along every $n$, so $\Gamma_{\rho_n}W=W$ for every $n$; every element of $\Ot$ is a product of at most three reflections, so $\Gamma_RW=W$ for all $R\in\Ot$.

\emph{Step 3: (a) and (b).} Let $\mathcal R_0\subset\Ot$ be a countable dense subgroup. Exactly as in the proof of Theorem~\ref{thm:D:decomp}(c), with $\Gamma_R$ in place of $\Theta_j$ and Step~2 in place of Lemma~\ref{lem:D:invloc}(ii), for $\nu$-a.e.\ $\xi$ the measure $\mu_\xi$ is invariant under every $R\in\mathcal R_0$. For bounded continuous $G$ the map $R\mapsto\E_{\mu_\xi}[G\circ R^{-1}]$ is continuous, by continuity of the action and dominated convergence, so $\mu_\xi$ is invariant under $\Ot$; with Theorem~\ref{thm:D:decomp}(b) this is (a). By Theorem~\ref{thm:D:decomp}(d) and (a), Step~1 applies to $\mu_\xi$ and gives (b).

\emph{Step 4: (c).} Fix such a $\xi$ and write $m:=\mu_\xi$; it is ergodic for $\R^3$, $\Ot$-invariant, and satisfies (H1)--(H3). The representation $R\mapsto\Gamma_R$ on $L^2(m)$ is strongly continuous (dominated convergence for bounded continuous $G$, density and $\norm{\Gamma_R}=1$ in general). Put $n_0:=(1,1,1)/\sqrt3$. Let $\abs{a^{(l)}}\to\infty$, and choose $R_l$ with $R_la^{(l)}=\abs{a^{(l)}}\,n_0$. Since $\Gamma_RU(a)=U(Ra)\Gamma_R$ and $m$ is $R$-invariant, $\E_m[G'\,U(a^{(l)})G]=\E_m[(\Gamma_{R_l}G')\,U(\abs{a^{(l)}}n_0)\,(\Gamma_{R_l}G)]$. Every subsequence has a further subsequence along which $R_l\to R$; along it the right side differs by $o(1)$ from $\E_m[(\Gamma_RG')\,U(\abs{a^{(l)}}n_0)(\Gamma_RG)]$. Since $\min_k\abs{a^{(l)}}(n_0)_k=\abs{a^{(l)}}/\sqrt3\to\infty$, Theorem~\ref{thm:D:diag} for $m$ gives the limit $\E_m[\Gamma_RG']\,\E_m[\Gamma_RG]=\E_m[G']\,\E_m[G]$.

\emph{Step 5: the forms of OS4.} With $G'=\overline{\Theta_jF}$, (c) along $\lambda a$ gives the cluster property of \cite[E4, p.~88]{OS73} for every time direction. If a Borel set $B$ is invariant under all $\tau_{te}$ for a unit vector $e$, (c) with $G=G'=\mathbf 1_B$ and $a=te$ gives $m(B)=m(B)^2$; this is ergodicity along every direction, in particular (OS4$_j$), and (OS4$'_j$) is (c) with $a=te_j$. Ergodicity under $\R^3$ is Theorem~\ref{thm:D:decomp}(b). Finally $\Ker H_n=\C\Omega_n$ by Proposition~\ref{add:OS4rel}(a), applied in a frame containing $n$ (Step~1).
\end{proof}
\begin{remark}\label{add:rotcond-rem}
(1) For the states of Theorem~\ref{thm:main} the hypotheses of Corollary~\ref{p:cor:rotcor} other than $\Ot$-invariance hold on $X=H^{-s}_{\mathrm{loc}}(\R^3;\cV)$ with $(R\Phi)(f):=\Phi(R^{-1}_*f)$, the push-forward of Section~\ref{ssec:symmetries}: the action is by linear homeomorphisms, continuous in $R$, and covariant for the local $\sigma$-algebras. $\Ot$-invariance of the law is equivalent to invariance of the Schwinger functions, $S_{n,p}(R_*g_1,\dots;R_*h_1,\dots)=S_{n,p}(g_1,\dots;h_1,\dots)$ for $R\in\Ot$, by Lemma~\ref{lem:D:realize}.
(2) So, if a limit state of Theorem~\ref{thm:main} is invariant under rotations, then almost every component of its law satisfies OS0 in the form of Remark~\ref{rem:components}, OS1 for the full Euclidean group, OS2 about every plane, the symmetry and reality conditions of OS3, and OS4 in every form. Rotation invariance is the only missing input for these conclusions; it is discussed in Section~\ref{ssec:introrot}.
(3) Only reflection positivity about the coordinate planes and invariance under rotations are used. Invariance under $\SO(3)$ suffices, since $\Ot=\SO(3)\cup\rho_1\SO(3)$ and $\rho_1\in\Oh$; so does invariance under a dense subgroup together with continuity of the action. In particular invariance of $S$ under one rotation about a coordinate axis by an angle not in $\pi\Q$ suffices: together with $\Oh$ it generates a dense subgroup of $\Ot$, and the set of $R$ with $R_*\mu=\mu$ is closed.
(4) The conclusion does not include the vanishing of the moments of odd degree $\ge3$ in the curvature for the components: $F\mapsto-F$ is not induced by any element of $\Ot$, since $R\in\Ot$ acts on 2-forms by a matrix similar to $\det(R)\,R$, which is never $-1$.
\end{remark}

\section{The reconstructed theory}\label{sec:recon}

Theorem~\ref{thm:D:recon} below uses no growth condition of \cite{OS73,OS75}: its proof uses the translations, the stabiliser $D_4$ of a coordinate axis in $\Oh$, reflection positivity about one coordinate plane, and the results of Section~\ref{sec:OS4}. On the relation with \cite{OS75}: the limit states satisfy the growth condition E0$''$ of \cite[(4.2), p.~287]{OS75} with $s=4$ and every $\beta>1$ (Theorem~\ref{thm:C-OS0}(c)), and E0$''$ implies the linear growth condition E0$'$ of \cite[(4.1), p.~287]{OS75} by the theorem of the Appendix of \cite[p.~303]{OS75}, whose proof is a Hermite expansion and uses no Euclidean invariance. The reconstruction theorem E$'\to$R$'$ of \cite[p.~287]{OS75}, which produces a Wightman theory, assumes in addition the axiom E1, invariance under all translations and all proper rotations \cite[p.~88]{OS73}. For the limit states E1 holds only for $\R^3\rtimes\Oh$, and the cluster axiom E4 only in the forms of Section~\ref{sec:OS4}; so that theorem does not apply as stated (Section~\ref{p:ss:extremal}).

Fix a limit state $S$ of Theorem~\ref{thm:main} and its law $\mu=\mu_S$. Choose the coordinate direction $e_3$ as Euclidean time. (By $\Oh$-invariance of $\mu$, every coordinate direction gives a unitarily equivalent theory; see (v) below.) Let $\cH:=\cH_3$, $\Omega:=\Omega_3$, $H:=H_3$, and $V(y):=V_3(y)$, $y\in\R^2=e_3^\perp$, as in Proposition~\ref{prop:D:recon}. Let $D_4:=\{g\in\Oh:ge_3=e_3\}$, the stabiliser of the time axis: the eight signed permutations of $(x_1,x_2)$, the symmetry group of a square.

\begin{theorem}[reconstruction with translations and $D_4$]\label{thm:D:recon}
\begin{enumerate}
\item[(i)] $\cH$ is a Hilbert space with the unit vector $\Omega$, $H\ge0$ is self-adjoint with $e^{-tH}[G]=[U(te_3)G]$, $H\Omega=0$, and $V$ is a strongly continuous unitary representation of the spatial translations commuting with $H$, with $V(y)\Omega=\Omega$. By Stone's theorem in two variables, $V(y)=e^{i(y_1P_1+y_2P_2)}$ with commuting self-adjoint momenta $P_1,P_2$.
\item[(ii)] For $g\in D_4$, $W(g)[G]:=[G\circ g^{-1}]$ defines a unitary representation of $D_4$ on $\cH$ with $W(g)\Omega=\Omega$, $W(g)e^{-tH}W(g)^{-1}=e^{-tH}$ and $W(g)V(y)W(g)^{-1}=V(gy)$.
\item[(iii)] Elements of $\Oh\setminus D_4$ do not act on $\cH$ in this way: those with $ge_3=-e_3$ exchange the two half spaces, and those with $ge_3\perp e_3$ do not preserve $\Hp{3}{0}$.
\item[(iv)] The joint vacuum is unique, that is, $\C\Omega$ is the joint fixed space of the $e^{-tH}$, $t\ge0$, and of $V$, if and only if $\mu$ is ergodic. On every ergodic component of Corollary~\ref{cor:D:schwinger} the joint vacuum is unique.
\item[(v)] For $g\in\Oh$ with $ge_3=e_j$, $[G]\mapsto[G\circ g^{-1}]$ is a unitary map $\cH_3\to\cH_j$ intertwining the reconstructions in directions $3$ and $j$.
\end{enumerate}
\end{theorem}
\begin{proof}
(i) is Proposition~\ref{prop:D:recon} with $j=3$; $H\Omega=0$ because $U(te_3)1=1$. (ii) Let $g\in D_4$. Then $g$ preserves $\Hp{3}{0}$ and commutes with $\rho_3$ (an orthogonal map fixing $e_3$ preserves $e_3^\perp$). So $G\mapsto G\circ g^{-1}$ maps $L^2_+$ onto itself, and by $g$-invariance of $\mu$ (Proposition~\ref{prop:discharge}),
\[
\ip{[G\circ g^{-1}]}{[G'\circ g^{-1}]}=\E\big[(\overline{\Theta_3G}\,G')\circ g^{-1}\big]=\ip{[G]}{[G']}.
\]
So $W(g)$ preserves $\cN$ and is isometric on the quotient, with inverse $W(g^{-1})$. The relations with $U(te_3)$ and $U(y)$ follow from $g\tau_ag^{-1}=\tau_{ga}$ and $ge_3=e_3$. (iii) For $ge_3=-e_3$, $g$ maps $\Hp{3}{0}$ to $\Hm{3}{0}$; such $g$ equal $\rho_3$ times an element of $D_4$, and $\rho_3$ is used in the inner product itself, not as an operator on $\cH$. For $ge_3=\pm e_k$, $k\ne3$, $g\Hp{3}{0}$ is a half space in direction $k$. (iv) is Proposition~\ref{prop:D:fixed}(ii)--(iii) and its converse: if $\mu$ is not ergodic, there is a non-constant bounded invariant $W$; by Lemma~\ref{lem:D:invloc}(i) it lies in $L^2_+$, $[W]$ is fixed, and $\norm{[W]-\E[W]\Omega}^2=\E\abs{W-\E W}^2>0$ by Proposition~\ref{prop:D:fixed}(i). The statement for components is Corollary~\ref{cor:D:schwinger}(iii). (v) As (ii), using $g\rho_3g^{-1}=\rho_j$ and $g\Hp{3}{0}=\Hp{j}{0}$.
\end{proof}
\noindent The theorem gives a quantum-mechanical system with a positive Hamiltonian, unitary spatial translations and the point group $D_4$. It gives no Lorentz covariance, and no Wightman-type distributions are claimed (Section~\ref{p:ss:extremal}).

\section{Correlation signs and non-triviality under a remainder bound}\label{sec:nontriv}

We ask whether a limit state of Theorem~\ref{thm:main} can be Gaussian, or the law of a pair of independent fields, one for the curvature and one for the Wick square, as for the free Maxwell field together with the Wick square of a free charged scalar. We do not settle this. What we prove is the following. A Gaussian limit state is decoupled (Lemma~\ref{lem:NT-gauss}), so it suffices to show that one mixed truncated function is non-zero. The simplest one, $M(g;h)=S^T_{2,1}(g,g;h)$, has a sign in every limit state (Theorem~\ref{thm:NT-sign}). On every lattice torus it equals $-e^2$ times a covariance of the matter (Lemma~\ref{lem:NT-dict}), and its coefficient at $e=\lambda=0$ is strictly negative on every lattice torus for $\mu_0>0$ (Proposition~\ref{prop:NT-lead}); its continuum value is the subject of Hypothesis~\ref{hyp:NT-C}. Non-triviality follows at small $e^2+\lambda$ from that hypothesis and one remainder bound, uniform in the lattice spacing and the torus (Theorem~\ref{thm:NT-cond}). The remainder bound is not available from the uniform estimates of this paper or from the literature it rests on (Remark~\ref{rem:NT-open} and Section~\ref{ss:nontriv}).

\subsection{Truncated functions; Gaussian and decoupled states}\label{ssec:NT-defs}

\begin{definition}[truncated functions]\label{def:NT-trunc}
Let $S$ be a Schwinger family whose generating functional $G^S$ exists on a tube $\Tub$ for every finite family of real test functions $g_1,\dots,g_n$, $h_1,\dots,h_p$ (every limit state has this property, by OS0 in Theorem~\ref{thm:main}(iii)). Since $G^S(0)=1$, $\log G^S$ is holomorphic near $0$. Its Taylor coefficients at $0$, $\partial_{z_1}\cdots\partial_{z_n}\partial_{w_1}\cdots\partial_{w_p}\log G^S(0)$, define the truncated functions $S^T_{n,p}(g_1,\dots,g_n;h_1,\dots,h_p)$ (test functions may repeat). The family $S$ is \emph{Gaussian} if $S^T_{n,p}=0$ whenever $n+p\ge3$. It is \emph{decoupled} if $S^T_{n,p}=0$ whenever $n\ge1$ and $p\ge1$. We put
\[
M(g;h):=S^T_{2,1}(g,g;h).
\]
The same definitions apply to the lattice states $\ev{\cdot}_{\eps,j,\mu_0}$, with lattice test functions; we write $M_{\eps,j,\mu_0}(g;h)$.
\end{definition}
A family is decoupled if and only if $G^S(z,w)=G^S(z,0)\,G^S(0,w)$ near $0$ for every family of test functions, hence on the connected tube by the identity theorem. This is the case if the curvature and the Wick square are independent under the law of Theorem~\ref{thm:main}(ii), for instance for the free Maxwell field (of any photon mass) together with the Wick square of a free charged scalar (of any mass).

\begin{lemma}[Gaussian limit states are decoupled]\label{lem:NT-gauss}
Let $S$ be a Schwinger family as in Definition~\ref{def:NT-trunc} with $S_{n,p}=0$ for odd $n$. Then $S^T_{1,1}=0$, $S^T_{1,0}=0$ and $M(g;h)=S_{2,1}(g,g;h)-S_{2,0}(g,g)\,S_{0,1}(h)$. If $S$ is Gaussian, it is decoupled. Every limit state $S\in\{S^{\mu_0},S^\fl,S^0\}$ of Theorem~\ref{thm:main} satisfies the hypothesis.
\end{lemma}
\begin{proof}
Expanding the logarithm to third order, $S^T_{1,0}=S_{1,0}$, $S^T_{1,1}=S_{1,1}-S_{1,0}S_{0,1}$ and $S^T_{2,1}(g,g;h)=S_{2,1}(g,g;h)-S_{2,0}(g,g)S_{0,1}(h)-2S_{1,1}(g;h)S_{1,0}(g)+2S_{1,0}(g)^2S_{0,1}(h)$; the terms with an odd number of curvature arguments vanish. If $S$ is Gaussian, $\log G^S$ is a polynomial of degree at most $2$ in $(z,w)$, and its only mixed term is $\sum_{i,k}z_iw_kS^T_{1,1}(g_i;h_k)=0$. The limit states satisfy $S_{n,p}=0$ for odd $n$ by OS3 (Theorem~\ref{thm:main}(iii), Definition~\ref{def:OS}).
\end{proof}

\subsection{The photon dictionary with a scalar source}\label{ssec:NT-dict}

Fix a lattice torus $T_{j,\eps}$, a real $b_T$, $\lambda>0$, $e\in\R$ ($e\ne0$ in the state $\fl$, whose period $P=2\pi/(\abs e\ell_j)$ needs it), and $\mu_0\in(0,\mu_*]\cup\{\fl\}$; on a fixed lattice torus the definitions of Section~\ref{sec:model} make sense for all such $(e,\lambda)$, and the standing convention $e\ne0$ is relaxed only in this section. Let $C=C_{\mu_0}=(-\Delta_\eps+\mu_0^2)^{-1}$, or $C=(-\Delta_\eps)^{-1}\Pi^\perp$ in the state $\fl$ (Definition~\ref{def:C-flat}). For a real lattice 2-form $g$ put
\[
f:=d^*g,\qquad a:=Cf,\qquad q:=\ip{f}{Cf}_\eps,\qquad Y(n):=\eps\sum_bn_ba_b=\ip{j_n}{a}_\eps\quad(j_n:=\eps^{-2}n),
\]
for integer-valued divergence-free currents $n$ on the bonds; let $\cN$ be the set of these currents, and in the state $\fl$ the subset of currents with zero windings, $\omega(n)=0$ (Lemma~\ref{lem:C-current}). Let $V(r):=3\eps r^2+\eps^3(\tfrac12b_Tr^2+\lambda r^4)$, and let $\mathbb P$ be the probability measure on $\cN\times\R_+^{T_{j,\eps}}$ given by
\begin{equation}\label{eq:NT-P}
\mathbb P(n,dr):=Z_{\mathbb P}^{-1}\,e^{-\frac12e^2\ip{j_n}{Cj_n}_\eps}\prod_bI_{n_b}\big(\eps r_{b_-}r_{b_+}\big)\prod_xe^{-V(r_x)}r_x\,dr_x .
\end{equation}
For a real site function $k$ let $\mathbb P_k$ be $\mathbb P$ tilted by $e^{r^2(\eps^3k)}$, where $r^2(u):=\sum_xu(x)r_x^2$, and let $\ev{\cdot}_k:=\ev{\cdot\,e^{\XR(k)}}/\ev{e^{\XR(k)}}$ denote the lattice state tilted by $e^{\XR(k)}$.

\begin{lemma}[photon dictionary with a scalar source]\label{lem:NT-dict}
$Z_{\mathbb P}<\infty$, $Y\in L^4(\mathbb P_k)$, and:
\begin{enumerate}
\item[(a)] $\ev{F(g)^2}_k=q-e^2\,\E_{\mathbb P_k}[Y^2]$ for every real site function $k$;
\item[(b)] for every real site function $h$, $\ev{F(g)}=\ev{F(g)\XR(h)}=0$ and
\[
M_{\eps,j,\mu_0}(g;h)=\ev{F(g)^2\XR(h)}-\ev{F(g)^2}\ev{\XR(h)}=-e^2\,\mathrm{Cov}_{\mathbb P}\big(Y^2,\,r^2(\eps^3h)\big).
\]
\end{enumerate}
\end{lemma}
\begin{proof}
\emph{Integrability.} By $\sum_{\nu\in\Z}I_\nu(t)=e^t$, $\sum_\nu\nu^2I_\nu(t)=te^t$ and $\sum_\nu\nu^4I_\nu(t)=(t+3t^2)e^t$ (the moments of $\nu$ under the weights $I_\nu(t)e^{-t}$, read off from $\sum_\nu I_\nu(t)e^{i\nu\theta}=e^{t\cos\theta}$), and by $\abs{Y}\le\eps\norm a_\infty\sum_b\abs{n_b}$, the sum over $n\in\Z^{\mathrm{bonds}}$ of $\prod_bI_{n_b}(\eps r_{b_-}r_{b_+})(1+Y^4)$ is at most a polynomial in $r$ times $\prod_be^{\eps r_{b_-}r_{b_+}}\le\prod_xe^{3\eps r_x^2}$. The factor $e^{3\eps r_x^2}$ cancels against $V$, and $e^{-\eps^3(\frac12b_T-k(x))r^2-\lambda\eps^3r^4}$ has all moments since $\lambda>0$. So $Z_{\mathbb P}<\infty$ and $Y\in L^4(\mathbb P_k)$.

(a) By the proof of Lemma~\ref{lem:C-Fsector}, with $tf$ in place of $f$ ($t\in\R$),
\[
\E_C\big[e^{tF(g)}W_k(A)\big]=e^{\frac12t^2q}\sum_{n\in\cN}w_k(n)\,e^{-\frac12e^2\ip{j_n}{Cj_n}_\eps}\cos\big(teY(n)\big),
\]
where $\ip{j_n}{Cf}_\eps=Y(n)$, and $w_k(n)\ge0$, $w_k(-n)=w_k(n)$ are the weights of Lemma~\ref{lem:C-current}; in the state $\fl$ the Haar average over the zero mode restricts the sum to $\omega(n)=0$, as there. By the formula for $w_k(n)$ in the proof of Lemma~\ref{lem:C-current}, $w_k(n)e^{-\frac12e^2\ip{j_n}{Cj_n}}$ is a positive constant (independent of $n$) times the $n$-marginal of $Z_{\mathbb P}\,e^{r^2(\eps^3k)}\,\mathbb P(n,dr)$; so the probability $\nu_k$ on $\cN$ proportional to $w_k(n)e^{-\frac12e^2\ip{j_n}{Cj_n}}$ is the $n$-marginal of $\mathbb P_k$. Dividing by the value at $t=0$,
\[
\ev{e^{tF(g)}}_k=e^{\frac12t^2q}\,\phi(t),\qquad \phi(t):=\E_{\nu_k}\big[\cos(teY)\big].
\]
The left side is finite for all $t$ (Lemma~\ref{lem:C-Fsector} with $\pm t$) and smooth in $t$, since $F(g)$ has exponential moments of every order. So $\phi$, the characteristic function of the symmetric law of $eY$ under $\nu_k$, is $C^2$, and $\phi''(0)=-e^2\E_{\nu_k}[Y^2]$: indeed $(1-\cos(tx))/t^2\to x^2/2$ and $0\le(1-\cos(tx))/t^2\le x^2/2$, so Fatou's lemma gives $\E[(eY)^2]\le-\phi''(0)<\infty$, and then dominated convergence gives equality. Differentiating twice at $t=0$ gives $\ev{F(g)^2}_k=q+\phi''(0)=q-e^2\E_{\mathbb P_k}[Y^2]$.

(b) $\ev{F(g)}=\ev{F(g)\XR(h)}=0$ by charge conjugation (Lemma~\ref{lem:C-sym}(iii)), which reverses the sign of $F$ and fixes $\XR$. Put $k=sh$ with $s\in\R$ small. The derivative at $s=0$ of $\ev{F(g)^2}_{sh}$ is $\ev{F(g)^2\XR(h)}-\ev{F(g)^2}\ev{\XR(h)}$, and the derivative of $\E_{\mathbb P_{sh}}[Y^2]$ is $\mathrm{Cov}_{\mathbb P}(Y^2,r^2(\eps^3h))$; in both cases differentiation under the integral is justified by the integrability above (with $e^{\abs s\,r^2(\eps^3\abs h)}$ absorbed by $e^{-\lambda\eps^3r^4}$). Apply (a). By Lemma~\ref{lem:NT-gauss} (on the lattice, odd curvature moments vanish by charge conjugation) the left side is $M_{\eps,j,\mu_0}(g;h)$.
\end{proof}

\subsection{Signs in every limit state}\label{ssec:NT-sign}

\begin{theorem}[signs]\label{thm:NT-sign}
(a) \emph{Lattice.} Let $\mu_0\in(0,\mu_*]$, let $T_{j,\eps}$ be any lattice torus, $b_T\in\R$, $\lambda>0$, $e\in\R$. For all real lattice 2-forms $g,g'$ and all real site functions $h,h'\ge0$:
\begin{enumerate}
\item[(i)] $M_{\eps,j,\mu_0}(g;h)\le0$; equivalently, $\mathrm{Cov}_{\mathbb P}(Y^2,r^2(\eps^3h))\ge0$ when $e\ne0$;
\item[(ii)] $\ev{\XR(h);\XR(h')}\ge0$;
\item[(iii)] $\ev{F(g)^2;F(g')^2}\ge0$,
\end{enumerate}
where $\ev{U;V}:=\ev{UV}-\ev{U}\ev{V}$. The same holds in the state $\fl$, for $e\ne0$, on every lattice torus $T_{j,\eps}$.

(b) \emph{Limits.} Let $S\in\{S^{\mu_0}\ (\mu_0\in Q_0),\,S^\fl,\,S^0\}$ be a limit state of Theorem~\ref{thm:main}. For all real 2-forms $g,g'$ with components in $C_c(\R^3)$ and all $h,h'\in C_c(\R^3)$ with $h,h'\ge0$:
\[
M(g;h)\le0,\qquad S^T_{0,2}(h,h')\ge0,\qquad S^T_{4,0}(g,g,g',g')+2\,S_{2,0}(g,g')^2\ge0 .
\]
\end{theorem}
\begin{proof}
(a) \emph{The setting of \cite{BFS1}, \S6.} There $d\mu_C$ is an arbitrary centred Gaussian measure with covariance $C>0$ on real bond variables $\vartheta_b$, the matter action is $-\sum_bc_br_{b_-}r_{b_+}\cos(-\theta_{b_-}+\vartheta_b+\theta_{b_+})$ with $c_b\ge0$, and at each site there is a finite measure $d\rho_x(r)$ with $\int e^{\alpha r^2}d\rho_x<\infty$ for all $\alpha>0$; periodic conditions on a rectangular lattice are allowed. The lattice state $\ev{\cdot}_{\eps,j,\mu_0}$ is of this form, with $\vartheta_b:=e\eps A_b$ (centred Gaussian with covariance $(e\eps)^2C_{\mu_0}>0$ for $e\ne0$; for $e=0$, where this covariance vanishes, the curvature is independent of the matter, so (i) holds with $M_{\eps,j,\mu_0}(g;h)=0$ and (iii) is the Gaussian identity $\ev{F(g)^2;F(g')^2}=2\ev{F(g)F(g')}^2\ge0$, while (ii) follows from the case $e\ne0$ by dominated convergence in $e$ at fixed lattice and fixed $b_T$, since $\abs{\cos}\le1$ in the hopping term and $\lambda>0$), $c_b=\eps$, and $d\rho_x(r)=e^{-V(r)}r\,dr$, by the polar form of the hopping term in the proof of Lemma~\ref{lem:postype}; the integrability condition holds because $\lambda>0$. (BFS let the angles run over $[-2\pi,2\pi)$; for $2\pi$-periodic functions this does not change expectations.) For real $u$ on bonds, $A_{\mathrm{BFS}}(u):=\sum_bu_b\vartheta_b$; with $u_b:=\eps^2(d^*g)_b/e$ this is $F(g)=\sum_b\eps^3(d^*g)_bA_b$.

(i) By \cite[Thm.~6.2(1), p.~280]{BFS1}, applied with the tilt $e^{r^2(\gamma)}$, $\gamma:=\eps^3(t\delta_x+\eta)$ ($t,\eta>0$), $f=0$ and the test function $u$: $\ev{e^{r^2(\gamma)};F(g)^2}\le0$. Let $\eta\downarrow0$ (dominated convergence; all moments of $e^{c\,r_y^2}F(g)^2$ are finite for every $c$, since $\lambda>0$). Then $\ev{e^{t\eps^3r_x^2};F(g)^2}\le0$ for every $t>0$. Divide by $t$ and let $t\downarrow0$: $t^{-1}(e^{t\eps^3r_x^2}-1)\to\eps^3r_x^2$, dominated by $\eps^3r_x^2e^{\eps^3r_x^2}$ for $t\le1$. Hence $\ev{r_x^2;F(g)^2}\le0$ for every site $x$. Multiply by $\eps^3h(x)\ge0$ and sum; the constants $W_\eps+z_\eps$ in $\XR(h)$ do not change covariances. So $\ev{\XR(h);F(g)^2}\le0$, which is $M_{\eps,j,\mu_0}(g;h)\le0$ by Lemma~\ref{lem:NT-dict}(b). For $e\ne0$ the equivalence with $\mathrm{Cov}_{\mathbb P}\ge0$ is the same lemma.

(ii) By \cite[Thm.~6.1, p.~276]{BFS1} with $P(r)=r_x^2$, $Q(r)=r_y^2$ (polynomials with positive coefficients), without current factors and without test functions for $A$ (in the notation of \cite[Thm.~6.1]{BFS1}, $n=m=0$ and $f=g=0$; these $m$ and $g$ are not the mass and the 2-form of this section): $\ev{r_x^2;r_y^2}\ge0$. Multiply by $\eps^6h(x)h'(y)\ge0$ and sum.

(iii) By \cite[Thm.~6.2(1)]{BFS1} with the tilt $e^{r^2(\eta)}$, $\eta>0$ constant, $f=tu'$ ($u'$ the coefficient vector of $F(g')$) and the test function $u$, and $\eta\downarrow0$: $\ev{\cos(tF(g'));F(g)^2}\le0$ for all $t\in\R$. (In the proof of \cite[Thm.~6.2]{BFS1}, $e^{iA(f)}$ is replaced by $\cos A(f)$, by the symmetry $A\to-A$, $\theta\to-\theta$.) Since $\abs{\cos x-1+x^2/2}\le x^4/24$ and $F(g'),F(g)$ have all moments, $\ev{\cos(tF(g'));F(g)^2}=-\frac{t^2}2\ev{F(g')^2;F(g)^2}+O(t^4)$. Hence $\ev{F(g')^2;F(g)^2}\ge0$.

\emph{The state $\fl$.} At fixed lattice, $\ev{\cdot}_{\eps,j,\fl}=\lim_{\mu_0\downarrow0}\ev{\cdot}_{\eps,j,\mu_0}$ on every $G\in\cA_\eps$ (Definition~\ref{def:algebra}) with $\sup_{0<\mu_0\le\mu_*}\ev{\abs G^2}_{\eps,j,\mu_0}<\infty$ (Lemma~\ref{lem:flat} with $\delta=1$; it is stated for $b=b_{T_j}(\eps)$, but its proof does not use the value of $b$, so it holds for every real $b_T$). The products that occur in (i)--(iii) are of this kind. By H\"older's inequality it suffices to bound $\ev{F(g)^8}_{\mu_0}$ and $\ev{r_x^{8}}_{\mu_0}$ uniformly in $\mu_0$. The first is at most $8!\,(\ev{e^{F(g)}}+\ev{e^{-F(g)}})\le2\cdot8!\,e^{\frac12\norm{g}_{2,\eps}^2}$ by Lemma~\ref{lem:C-Fsector}. For the second, put $\Psi(A):=\Psi_{r_x^8}(A)=\int r_x^8\,e^{-S^{\mathrm m}_{b_T}(A,\varphi)}\prod_yd^2\varphi(y)$ (Definition~\ref{def:algebra}, with $b_T$ in place of $b_{T_j}(\eps)$). The proof of Lemma~\ref{lem:postype}, with the factor $r_x^8$ inside the radial integral, gives $\Psi(A)=\sum_nw'(n)e^{ie\eps\sum_bn_bA_b}$ with weights $w'(n)\ge0$, so $0\le\Psi(A)\le\Psi(0)$. Moreover, $\E_{\gamma_{\mu_0}}[W_0]=\sum_nw_0(n)e^{-\frac12e^2\ip{j_n}{C_{\mu_0}j_n}}\ge w_0(0)>0$ (the $n=0$ term). Hence $\ev{r_x^8}_{\mu_0}=\E_{\gamma_{\mu_0}}[\Psi]/\E_{\gamma_{\mu_0}}[W_0]\le\Psi(0)/w_0(0)$ for all $\mu_0$. The inequalities (i)--(iii) pass to the limit $\mu_0\downarrow0$.

(b) The truncated functions in (b) are polynomials in the Schwinger functions $S_{2,1},S_{2,0},S_{0,1},S_{0,2},S_{4,0}$. Along the subsequence of Theorem~\ref{thm:main}(ii) the lattice Schwinger functions at sampled test functions converge to those of $S^{\mu_0}$ and $S^\fl$, and the Schwinger functions of $S^{\mu_0}$ converge to those of $S^0$ as $\mu_0\downarrow0$ in $Q_0$. Sampling preserves $h,h'\ge0$. The quantities in (a) are, on the lattice, $M_{\eps,j,\mu_0}(g;h)$ (Lemma~\ref{lem:NT-dict}(b)), $\ev{\XR(h);\XR(h')}=S^{T}_{0,2}(h,h')$ and $\ev{F(g)^2;F(g')^2}=S_{4,0}(g,g,g',g')-S_{2,0}(g,g)S_{2,0}(g',g')=S^T_{4,0}(g,g,g',g')+2S_{2,0}(g,g')^2$; the last identity is the fourth-order expansion of the logarithm, in which the terms with an odd number of curvature arguments vanish by charge conjugation, as in the proof of Lemma~\ref{lem:NT-gauss}. The inequalities pass to the limits.
\end{proof}

\begin{corollary}[a criterion for non-triviality]\label{cor:NT-crit}
Let $S$ be a limit state as in Theorem~\ref{thm:NT-sign}(b). If $M(g;h)<0$ for one real 2-form $g$ and one $h\ge0$, then $S$ is not decoupled. Hence $S$ is not Gaussian, and the curvature and the Wick square are not independent under the law of $S$. In particular $S$ is not the free Maxwell field (of any photon mass) together with the Wick square of a free charged scalar. By Theorem~\ref{thm:NT-sign}, the hypothesis is equivalent to $M(g;h)\ne0$.
\end{corollary}

\subsection{The leading coefficient}\label{ssec:NT-lead}

On a fixed lattice torus regard the coefficient $b_T=a_{K,T}(m)$ of the torus counterterm scheme (Definition~\ref{def:counterterm}) as the polynomial in $(e^2,\lambda)$ given there; at $e=\lambda=0$ it equals $m^2$. Write $\mathbb P^{e,\lambda}$ for the measure~\eqref{eq:NT-P} at these couplings, and for $\lambda=0$ use~\eqref{eq:NT-P} with $\lambda=0$ and $b_T=m^2$. A closed walk $\omega=(\omega_0,\omega_1,\dots,\omega_N=\omega_0)$ on $T_{j,\eps}$ has $\abs\omega=N\ge1$ nearest-neighbour steps. For a function $a$ on oriented bonds ($a_{-b}=-a_b$) put $\Phi_\omega(a):=\sum_{i=0}^{N-1}a_{(\omega_i,\omega_{i+1})}$.

\begin{proposition}[the leading coefficient on a fixed lattice torus]\label{prop:NT-lead}
Fix $\eps$, $j$, $m>0$ and $\mu_0\in(0,\mu_*]\cup\{\fl\}$, a real lattice 2-form $g$ and a real site function $h$.
\begin{enumerate}
\item[(a)] For $\lambda>0$, $\lim_{e\to0}e^{-2}M^{(e,\lambda)}_{\eps,j,\mu_0}(g;h)=c_{\eps,j,\mu_0}(g,h;\lambda):=-\mathrm{Cov}_{\mathbb P^{0,\lambda}}\big(Y^2,r^2(\eps^3h)\big)$. If $h\ge0$, then $c_{\eps,j,\mu_0}(g,h;\lambda)\le0$.
\item[(b)] $\lim_{\lambda\downarrow0}c_{\eps,j,\mu_0}(g,h;\lambda)=c^{\mathrm{free}}_{\eps,j,\mu_0}(g,h):=-\mathrm{Cov}_{\mathbb P^{0,0}}\big(Y^2,r^2(\eps^3h)\big)$.
\item[(c)] For $\mu_0\in(0,\mu_*]$, with $a=C_{\mu_0}d^*g$, the sum below converges absolutely and
\begin{equation}\label{eq:NT-walk}
c^{\mathrm{free}}_{\eps,j,\mu_0}(g,h)=-\frac{2\eps^4}{6+\eps^2m^2}\sum_{\omega}\frac{(6+\eps^2m^2)^{-\abs\omega}}{\abs\omega}\,\Phi_\omega(a)^2\sum_{i=0}^{\abs\omega-1}h(\omega_i),
\end{equation}
the sum running over all closed walks on $T_{j,\eps}$.
\item[(d)] For $\mu_0\in(0,\mu_*]$: if $h\ge0$, $h\not\equiv0$ and $d^*g\not\equiv0$, then $c^{\mathrm{free}}_{\eps,j,\mu_0}(g,h)<0$.
\end{enumerate}
\end{proposition}
\begin{proof}
(a) By Lemma~\ref{lem:NT-dict}(b), $e^{-2}M^{(e,\lambda)}=-\mathrm{Cov}_{\mathbb P^{e,\lambda}}(Y^2,r^2(\eps^3h))$ for $e\ne0$. The density of $\mathbb P^{e,\lambda}$ depends on $e$ through $e^{-\frac12e^2\ip{j_n}{Cj_n}}\le1$ and through $b_T$, which is continuous in $e$. Let $b_{\min}:=\min_{\abs e\le1}b_T$. For $\abs e\le1$ the integrands of $Z_{\mathbb P}$, $\E_{\mathbb P}[Y^2r^2(\eps^3h)]$, $\E_{\mathbb P}[Y^2]$ and $\E_{\mathbb P}[r^2(\eps^3h)]$ are dominated by the $e$-independent integrable function of the proof of Lemma~\ref{lem:NT-dict}, with $b_T$ replaced by $b_{\min}$. So the covariance is continuous in $e$, and its value at $e=0$ is the limit. The sign follows from Theorem~\ref{thm:NT-sign}(a)(i).

(b) At $e=0$, $b_T=m^2+\lambda\beta_1+\lambda^2\beta_2$ with real $\beta_1,\beta_2$ depending on $(\eps,T)$ (Definition~\ref{def:counterterm}: at $e=0$ the monomials are $\lambda$ and $\lambda^2$). For $0\le\lambda\le\lambda_1$ small, $b_T\ge m^2/2$. With $\prod_bI_{n_b}(\eps r_{b_-}r_{b_+})\le$ (polynomial in $r$ and $n$) $\times\prod_be^{\eps r_{b_-}r_{b_+}}$ and $\sum_b\eps r_{b_-}r_{b_+}\le3\eps\sum_xr_x^2$, the integrands are dominated by a polynomial times $\prod_xe^{-\eps^3m^2r_x^2/4}$, uniformly in $\lambda\in[0,\lambda_1]$. Dominated convergence gives the limit.

(c) Let $\Theta$ be a real function on oriented bonds. By the Bessel expansion of Lemma~\ref{lem:C-current}, with the phase $e\eps A_b$ replaced by $\Theta_b$,
\[
\sum_nw_k(n)\,e^{i\sum_bn_b\Theta_b}=\int e^{-S^{\Theta}(\varphi)+\XR(k)}\,d\varphi ,
\]
where $S^\Theta$ is the matter action with the hopping $\abs{U(\Theta_b)\varphi(b_+)-\varphi(b_-)}^2$. At $\lambda=0$, $b_T=m^2$, the right side is a complex Gaussian integral: $S^\Theta-\sum_x\eps^3k(x)\abs{\varphi(x)}^2=\frac12\ip{\varphi}{(H_\Theta-2k)\varphi}_\eps$ with $H_\Theta=-\Delta_\Theta+m^2$ ($-\Delta_\Theta\ge0$ the covariant lattice Laplacian on $\C$-valued functions). Hence, for $\norm k_\infty<m^2/2$, it equals a $\Theta$-independent constant times $\det(H_\Theta-2k)^{-1}$. Take $\Theta=\tau\eps a$, so that $\sum_bn_b\Theta_b=\tau Y(n)$. Since the weights are even in $n$, the first $\tau$-derivative at $0$ vanishes, and
\[
\E_{\mathbb P^{0,0}_k}[Y^2]=-\partial_\tau^2\log\Big(\sum_nw_k(n)e^{i\tau Y(n)}\Big)\Big|_{\tau=0}=\partial_\tau^2\,\tr\log\big(H_{\tau\eps a}-2k\big)\Big|_{\tau=0}.
\]
Write $H_\Theta-2k=D_k(1-K_\Theta)$, with $D_k$ the diagonal operator $6\eps^{-2}+m^2-2k(x)$ and $K_\Theta$ the hopping operator, whose matrix elements are $D_k(x)^{-1}\eps^{-2}e^{\pm i\Theta_b}$ for $x,y$ joined by a bond $b$. Its norm is at most $6/(6+\eps^2(m^2-2\norm k_\infty))<1$. So $\tr\log(H_\Theta-2k)=\tr\log D_k-\sum_{N\ge1}N^{-1}\tr K_\Theta^N$, and $\tr K_\Theta^N=\sum_{\abs\omega=N}u_k(\omega)e^{\pm i\Phi_\omega(\Theta)}$, where $u_k(\omega):=\prod_{i=0}^{N-1}\big(6+\eps^2(m^2-2k(\omega_i))\big)^{-1}$; the global sign of the phase is a convention and plays no role. Since $\abs{\Phi_\omega(a)}\le\abs\omega\norm a_\infty$ and there are at most $\abs{T_{j,\eps}}6^N$ closed walks of length $N$, the series may be differentiated termwise. This gives
\[
\E_{\mathbb P^{0,0}_k}[Y^2]=\eps^2\sum_\omega\abs\omega^{-1}u_k(\omega)\,\Phi_\omega(a)^2 .
\]
Finally $c^{\mathrm{free}}=-\partial_s\E_{\mathbb P^{0,0}_{sh}}[Y^2]\big|_{s=0}$ (as in the proof of Lemma~\ref{lem:NT-dict}(b); at $\lambda=0$ the integrability used there is supplied by the domination of (b), with $b_T=m^2$: a polynomial times $\prod_xe^{-\eps^3m^2r_x^2/4}$, valid for $\abs s\,\norm h_\infty<m^2/4$), and $\partial_su_{sh}(\omega)|_{s=0}=u_0(\omega)\sum_i2\eps^2h(\omega_i)/(6+\eps^2m^2)$ with $u_0(\omega)=(6+\eps^2m^2)^{-\abs\omega}$. Termwise differentiation in $s$ is justified by the same bounds, uniformly for $\abs s\,\norm h_\infty<m^2/4$. This is~\eqref{eq:NT-walk}.

(d) Every term of~\eqref{eq:NT-walk} is $\le0$, so it suffices to find one closed walk $\omega$ with $\Phi_\omega(a)\ne0$ that visits a site where $h>0$. Since $C_{\mu_0}>0$ and $f=d^*g\ne0$, $\ip{g}{dC_{\mu_0}d^*g}_\eps=\ip{f}{C_{\mu_0}f}_\eps>0$, so $da=dC_{\mu_0}d^*g\ne0$. Choose a plaquette $p$ with $(da)(p)\ne0$, a site $x_0$ with $h(x_0)>0$, and a nearest-neighbour path $\gamma$ from $x_0$ to a corner of $p$. Let $\omega$ be: $\gamma$, then once around $\partial p$, then $\gamma$ backwards. The two passages along $\gamma$ cancel in $\Phi_\omega$, so $\Phi_\omega(a)=\pm\sum_{b\in\partial p}\pm a_b=\pm\eps(da)(p)\ne0$, and $\omega_0=x_0$.
\end{proof}

Proposition~\ref{prop:NT-lead} is a statement on one lattice torus. Its continuum counterpart is the following hypothesis, which we formulate and do not prove.

\begin{hypothesis}[the continuum value of the leading coefficient]\label{hyp:NT-C}
For every real 2-form $g$ with components in $C_c^\infty(\R^3)$ and every $0\le h\in C_c^\infty(\R^3)$, sampled as in Definition~\ref{def:C-sampling}, the following holds, with $a_{\mu_0}:=(-\Delta+\mu_0^2)^{-1}d^*g$ on $\R^3$ for $\mu_0\in[0,\mu_*]$. For every $\mu_0\in[0,\mu_*]$ the integral
\begin{equation}\label{eq:NT-cont}
c^{\mathrm{free}}_{\mu_0}(g,h):=-2\int_0^\infty\frac{dT}{T}\,\frac{e^{-m^2T}}{(4\pi T)^{3/2}}\int_{\R^3}dx\ \E^T_{x,x}\Big[\Big(\oint_\omega a_{\mu_0}\cdot d\omega\Big)^2\int_0^Th(\omega_t)\,dt\Big]
\end{equation}
converges, where $\E^T_{x,x}$ is the expectation for the Brownian bridge $\omega$ with generator $\Delta$ from $x$ to $x$ in time $T$, and $\oint$ is the Stratonovich line integral. As $\eps\to0$ and $\ell_j\to\infty$, $c^{\mathrm{free}}_{\eps,j,\mu_0}(g,h)\to c^{\mathrm{free}}_{\mu_0}(g,h)$ for $\mu_0\in(0,\mu_*]$, and $c^{\mathrm{free}}_{\eps,j,\fl}(g,h)\to c^{\mathrm{free}}_0(g,h)$. The function $\mu_0\mapsto c^{\mathrm{free}}_{\mu_0}(g,h)$ is continuous on $[0,\mu_*]$, and $c^{\mathrm{free}}_{\mu_0}(g,h)<0$ if $d^*g\not\equiv0$ and $h\not\equiv0$.
\end{hypothesis}

\begin{remark}[on Hypothesis (C)]\label{p:rem:NT-C}
We refer to Hypothesis~\ref{hyp:NT-C} as (C). It is a one-loop statement: by Proposition~\ref{prop:NT-lead}(c), $c^{\mathrm{free}}_{\eps,j,\mu_0}(g,h)$ is a weighted sum over closed lattice walks, and \eqref{eq:NT-cont} is its expected Brownian-bridge limit. A proof would combine the convergence of the lattice coefficient (by lattice power counting \cite{Rei88} in its three-dimensional form, Proposition~\ref{app:G1:reisz3d}, applied to the two one-loop graphs, the triangle with two current vertices and one insertion of $\abs\varphi^2$ and the bubble with one seagull vertex and one insertion, both of superficial degree $-1$ in $d=3$, or by an invariance principle for the weighted walks of~\eqref{eq:NT-walk}, with the discrete line integral converging to the Stratonovich integral); the Feynman--Kac--It\^o formula for $-(\nabla-i\tau a_{\mu_0})^2+m^2-2sh$ ($\tau,s\in\R$), differentiated twice in $\tau$ and once in $s$ at $\tau=s=0$, which gives~\eqref{eq:NT-cont}; the support theorem for the Brownian bridge, for the strict negativity; and, for the state $\fl$, the fact that the restriction of $\cN$ to currents with zero windings changes the coefficient by $o(1)$ as $\ell_j\to\infty$. We do not carry this out.
\end{remark}

\subsection{Non-triviality given a remainder bound}\label{ssec:NT-cond}

\begin{theorem}[non-triviality under (R) and (C)]\label{thm:NT-cond}
Assume:
\begin{enumerate}
\item[(R)] There are $\delta_R\in(0,0.1]$ and a function $\rho$ on $(0,\delta_R]$ with $\rho(t)\to0$ as $t\downarrow0$ such that the following holds. For every pair $g,h$ as in Hypothesis~\ref{hyp:NT-C} there is $\kappa(g,h)<\infty$ such that, for every $(e,\lambda)$ with $\lambda>0$, $e^2>0$ and $e^2+\lambda\le\delta_Rm$ (such $(e,\lambda)$ satisfy the hypotheses of Theorem~\ref{thm:main}, since $\delta_R\le0.1$), the lattice states of Theorem~\ref{thm:main} at $(e,\lambda)$, with the constants $\eps_0$ and $K_0$ of that theorem chosen at $(e,\lambda)$ and $\eps_K=2^{-K}\eps_0$, satisfy
\[
\limsup_{K\to\infty,\ j\to\infty}\ \sup_{\mu_0\in(0,\mu_*]\cup\{\fl\}}\ \Big|e^{-2}M_{\eps_K,j,\mu_0}(g;h)-c^{\mathrm{free}}_{\eps_K,j,\mu_0}(g,h)\Big|\ \le\ \rho\big((e^2+\lambda)/m\big)\,\kappa(g,h).
\]
\end{enumerate}
Assume also (C). Let $g,h$ be as in Hypothesis~\ref{hyp:NT-C}, with $d^*g\not\equiv0$ and $h\not\equiv0$, and put $c_0(g,h):=-\max_{\mu\in[0,\mu_*]}c^{\mathrm{free}}_\mu(g,h)>0$. If $(e,\lambda)$ is as in (R) and $\rho((e^2+\lambda)/m)\kappa(g,h)<c_0(g,h)/2$, then every limit state $S\in\{S^{\mu_0}\ (\mu_0\in Q_0),S^\fl,S^0\}$ of Theorem~\ref{thm:main} at $(e,\lambda)$ satisfies
\[
M(g;h)\le-\tfrac12e^2c_0(g,h)<0 .
\]
In particular, for $e^2+\lambda$ small enough (depending on $g,h$) every limit state is not decoupled and not Gaussian (Corollary~\ref{cor:NT-crit}).
\end{theorem}
\begin{proof}
$c_0>0$ by the strict negativity and the continuity in (C) ($[0,\mu_*]$ is compact). Let $S^{\mu_0}$ or $S^\fl$ be the limit along $(\eps_{K_l},T_{j_l})$. By (R) and (C), $\limsup_le^{-2}M_{\eps_{K_l},j_l,\mu_0}(g;h)\le\lim_lc^{\mathrm{free}}_{\eps_{K_l},j_l,\mu_0}(g,h)+c_0/2\le-c_0/2$, and the same for $\fl$. By Theorem~\ref{thm:main}(ii) the left side is $e^{-2}M(g;h)$ for $S^{\mu_0}$, resp.\ $S^\fl$. For $S^0$, $M$ is the limit of the values for $S^{\mu_0}$, $\mu_0\downarrow0$ in $Q_0$, each of which is $\le-e^2c_0/2$.
\end{proof}

\begin{remark}[what (R) asks]\label{rem:NT-open}
(i) \emph{What the uniform bounds give.} Theorem~\ref{thm:C-T9} and Proposition~\ref{prop:C-bounds} bound every Schwinger function, and so $\abs{M(g;h)}$, by a constant times $\norm g_2^2N(h)$. The constants $\sigma_+$, $C$ (and $\eps_0,\ell_0,K_0$) are produced at each point of $\fW$ by the conditions of Appendix~\ref{app:B} and Lemma~\ref{lem:window}, and nothing is stated about how they depend on $(e,\lambda)$. So the Cauchy estimates give neither a factor $e^2$ nor any lower bound: the tube bound is uniform in the sources and in the regulators, not in the couplings.
(ii) \emph{What (R) asks.} By Lemma~\ref{lem:NT-dict}, (R) says that the matter covariance $\mathrm{Cov}_{\mathbb P^{e,\lambda}}(Y^2,r^2(\eps^3h))$ stays within $\rho\kappa$ of its value at $(e,\lambda)=(0,0)$, uniformly in $\eps$, $j$ and $\mu_0$. Moving $\lambda$ from $0$ inserts $\sum_x\eps^3r_x^4$ together with the $\lambda$-dependent part of $b_T$; one-loop power counting suggests that each of the two insertions alone diverges as $\eps\to0$, like $\lambda\eps^{-1}$, and that only their sum is finite. Moving $e$ from $0$ inserts the induced current interaction $e^{-\frac12e^2\ip{j_n}{Cj_n}}$ and the $e^2$-dependent part of $b_T$. A bound uniform in $\eps$ on these insertions in a three-point function is a renormalised perturbative estimate with a remainder. As printed, \cite{Ba1,Ba2,Ba3,Ba4} give two-sided bounds on partition functions, \cite[Thm.~2.1, p.~654]{K1} gives the convergence of $Z(T,g,h)$ and $\abs{\ln Z(T,g,h)}\le C\abs T$ with $C$ depending on $g$, $h$ and $\abs T$, and \cite[Thm.~3.4, p.~656]{K1} compares effective actions at two spacings; none of them states asymptotics in $(e,\lambda)$ of $\ln Z(T,g,h)-\ln Z(T,0,0)$.
(iii) \emph{Signs do not suffice.} The proof of Theorem~\ref{thm:NT-cond} uses only the upper half of (R), the bound $M\le e^2(c^{\mathrm{free}}+\rho\kappa)$. The correlation inequalities of \cite{BFS1} give only non-strict signs, and the strict sign of Proposition~\ref{prop:NT-lead}(d) holds on each lattice but is not uniform. Lebowitz's inequality concerns the field $\varphi$, which is not gauge invariant and is not among our observables.
\end{remark}

\begin{remark}[the notion of non-triviality]\label{rem:NT-scope}
Here non-triviality means that the state is not decoupled (Definition~\ref{def:NT-trunc}), the conclusion of Corollary~\ref{cor:NT-crit}. This is the informative notion: the free Maxwell field together with the Wick square of a free charged scalar is decoupled, but as a joint family it is not Gaussian, because the Wick square of a free field is not Gaussian. No result of this paper excludes a Gaussian limit state: such a state would have $M\equiv0$, which is compatible with every sign in Theorem~\ref{thm:NT-sign}, and the Gaussian domination bound of Lemma~\ref{lem:C-Fsector} allows the free Maxwell value. Nothing here concerns non-Gaussianity of the curvature field alone (Theorem~\ref{thm:NT-sign}(b) allows $S^T_{4,0}=0$).
\end{remark}

\section{The curvature two-point function}\label{sec:massless}

This section collects what the methods of the paper give about the photon mass in the massless states $S^\fl$ and $S^0$. The answer is one-sided. We prove the Bianchi identity in every limit state; the curvature two-point function of every limit state is a Fourier multiplier dominated by that of the free massless photon; at positive photon mass it is suppressed at small momenta, so $S^{\mu_0}$ has no massless component. For $S^\fl$ and $S^0$ we prove that a massless photon, in the sense of Definition~\ref{def:M-massless}, excludes integrable, in particular exponentially decaying, curvature correlations and excludes a gap in the curvature sector of the reconstructed theory (Theorem~\ref{thm:M-dichotomy}, Proposition~\ref{prop:M-gap}, Corollary~\ref{cor:M-summary}). Whether $S^\fl$ or $S^0$ is massless is not decided here: it needs a lower bound on the curvature two-point function, which no tool of this paper gives (Remark~\ref{p:rem:M-soft} and Section~\ref{p:ss:photon}).

\emph{Standing assumptions.} $S$ is a limit point in the sense of Section~\ref{sec:C-OS}. That is, $S$ is $S^{\mu_0}$ with $\mu_0\in Q_0$, or a limit point at a photon mass $\mu_0\in(0,\mu_*]$ along a further subsequence, or $S^\fl$, or $S^0$; the fixed-torus continuum states, which Section~\ref{sec:C-OS} includes where stated, are not limit points here (they are not invariant under all translations of $\R^3$, which Lemma~\ref{lem:M-mult} needs). We say that $S$ is \emph{at photon mass $\mu_0>0$} in the first two cases. The lattice states are those of Definition~\ref{def:states}. Everything below uses only Theorem~\ref{thm:main}(i)--(iv), Lemmas~\ref{lem:C-current}, \ref{lem:C-Fsector} and \ref{lem:C-sym}, Definition~\ref{def:C-sampling}, Propositions~\ref{prop:C-bounds}, \ref{prop:C-compact} and~\ref{prop:C-tight}, and Theorems~\ref{thm:C-OS1}, \ref{thm:C-OS2} and~\ref{thm:C-OS3}. Theorem~\ref{thm:mainL3} gives the statements (i)--(iv) of Theorem~\ref{thm:main} at blocking factor $3$, under its hypothesis, and the lattice lemmas used here concern one torus and do not involve $L$. So every result of this section holds, under the same hypothesis, for the limit states of Theorem~\ref{thm:mainL3}.

\emph{Notation.} We write curvature test functions as vector fields. For a real $G\in C_c(\R^3;\R^3)$ let $\star^{-1}G$ be the 2-form with components $g_{23}=G_1$, $g_{13}=-G_2$, $g_{12}=G_3$. Then $\norm{\star^{-1}G}_2=\norm G_2$, and we put
\[
B(G):=F(\star^{-1}G),\qquad \mathfrak b_S(G,G'):=S_{2,0}(\star^{-1}G,\star^{-1}G'),
\]
the smeared magnetic field $B=\star F$ and its two-point function. A direct computation gives $(\delta\star^{-1}G)_\nu=(\operatorname{curl}G)_\nu$ for $G\in C^1_c$, where $\delta$ is the $L^2$-adjoint of $d$ on forms, $(\delta g)_\nu=-\sum_\mu\partial_\mu g_{\mu\nu}$. The Fourier transform is $\hat f(p)=\int e^{-ip\cdot x}f(x)\,dx$, so that $\norm f_2^2=(2\pi)^{-3}\int\abs{\hat f}^2dp$. For $p\neq0$ let $P_T(p):=\mathbf 1-pp^{\mathsf T}/\abs p^2$ be the orthogonal projection of $\C^3$ onto $p^\perp$, and let $\Pi_T$ be the Fourier multiplier by $P_T$ on $L^2(\R^3;\C^3)$. For $R>0$ put
\[
G_R(x):=R^{-3}G(x/R),\qquad \widehat{G_R}(p)=\hat G(Rp),\qquad \norm{G_R}_2^2=R^{-3}\norm G_2^2 .
\]
The \emph{free massless photon} is the centred Gaussian curvature field with $\mathfrak b(G,G')=\ip{\Pi_TG}{\Pi_TG'}_{L^2}$, the curvature of the Feynman-gauge massless vector field:
\[
\ip{\delta\star^{-1}G}{(-\Delta)^{-1}\delta\star^{-1}G}=(2\pi)^{-3}\int\abs{p\times\hat G}^2\abs p^{-2}dp=\norm{\Pi_TG}_2^2 .
\]

\subsection{Two lattice facts}

\begin{lemma}[free part minus screening, on one torus]\label{lem:M-lattice}
Fix $(\varepsilon,T_j)$, $\mu_0\in(0,\mu_*]\cup\{\fl\}$ and a real lattice 2-form $r$ on $T_{j,\varepsilon}$. Let $C$ be as in Lemma~\ref{lem:C-Fsector}, let $w_0(n)$ be the weights of Lemma~\ref{lem:C-current} at $k=0$, and put $j_n:=\varepsilon^{-2}n$. Let $\tilde\nu$ be the probability measure on divergence-free integer currents $n$ with $\tilde\nu(n)\propto w_0(n)\,e^{-\frac12e^2\ip{j_n}{Cj_n}_\varepsilon}$, multiplied by $\mathbf 1\{\omega(n)=0\}$ in the $\fl$ state, and $a_r(n):=e\ip{j_n}{C\,d^*r}_\varepsilon$. Then:
\begin{enumerate}
\item[(a)] for every real $t$, $\ev{e^{tF(r)}}_{\varepsilon,j,\mu_0}=e^{\frac12t^2V(r)}\sum_n\tilde\nu(n)\cos\big(t\,a_r(n)\big)$, with $V(r)=\ip{d^*r}{C\,d^*r}_\varepsilon$;
\item[(b)] $\ev{F(r)}_{\varepsilon,j,\mu_0}=0$, $\sum_n\tilde\nu(n)a_r(n)^2<\infty$, and
\[
\ev{F(r)^2}_{\varepsilon,j,\mu_0}=V(r)-\sum_n\tilde\nu(n)\,a_r(n)^2\ \le\ V(r)\ \le\ \norm r_{2,\varepsilon}^2 ;
\]
\item[(c)] if $\mu_0\in(0,\mu_*]$, then $V(r)\le\mu_0^{-2}\norm{d^*r}_{2,\varepsilon}^2$.
\end{enumerate}
\end{lemma}

\begin{proof}
(a) Put $f:=t\,d^*r$. The proof of Lemma~\ref{lem:C-Fsector} (case $\mu_0>0$, with $k=0$) gives
\[
\E_C\big[e^{\ip fA_\varepsilon}W_0(A)\big]=e^{\frac12\ip f{Cf}_\varepsilon}\sum_nw_0(n)\,e^{-\frac12e^2\ip{j_n}{Cj_n}_\varepsilon}\cos\big(e\ip{j_n}{Cf}_\varepsilon\big),
\]
and the same identity at $f=0$ gives $\E_C[W_0]=\sum_nw_0(n)e^{-\frac12e^2\ip{j_n}{Cj_n}_\varepsilon}$. Divide. In the $\fl$ state the average over the zero mode replaces $w_0(n)$ by $w_0(n)\mathbf 1\{\omega(n)=0\}$, as in that proof. For such $n$ the phase $e\ip{j_n}{c}_\varepsilon=e\ell_T\sum_\mu c_\mu\omega_\mu(n)$ of a constant $c$ vanishes (Lemma~\ref{lem:C-current}), so only $A^\perp$, with covariance $C=(-\Delta_\varepsilon)^{-1}\Pi^\perp$, enters, and the same Gaussian computation applies.
(b) Let $\phi(t):=\sum_n\tilde\nu(n)\cos(ta_r(n))$ and $f(t):=\ev{e^{tF(r)}}$. Then $f(t)=e^{\frac12t^2V}\phi(t)$ is finite for all real $t$ by (a). Hence $\ev{e^{\abs t\abs{F(r)}}}\le f(t)+f(-t)<\infty$, and $f$ is real-analytic with $f^{(k)}(0)=\ev{F(r)^k}$. Since $\phi$ is even, so is $f$; thus $\ev{F(r)}=0$, and $e^{-\frac12t^2V}f(t)=1+\tfrac12t^2(\ev{F(r)^2}-V)+O(t^4)$. Therefore
\[
\sum_n\tilde\nu(n)\,\frac{2(1-\cos(ta_r(n)))}{t^2}=\frac{2(1-\phi(t))}{t^2}\ \longrightarrow\ V-\ev{F(r)^2}\qquad(t\to0).
\]
The summands are non-negative, bounded by $a_r(n)^2$, and converge to $a_r(n)^2$. By Fatou's lemma $\sum\tilde\nu a_r^2\le V-\ev{F(r)^2}<\infty$. Dominated convergence then gives equality. The bound $V(r)\le\norm r^2_{2,\varepsilon}$ is Lemma~\ref{lem:C-Fsector}.
(c) $C_{\mu_0}=(-\Delta_\varepsilon+\mu_0^2)^{-1}\le\mu_0^{-2}$ as an operator.
\end{proof}

\begin{lemma}[sampling and the lattice $d^*$]\label{lem:M-sampling}
Let $\supp G\subset B_R$, $\ell_j>2R+2$, and let $g_\varepsilon$ be the sample of $\star^{-1}G$ (Definition~\ref{def:C-sampling}).
\begin{enumerate}
\item[(a)] If $G\in C^1_c$, then $\norm{d^*g_\varepsilon}_{2,\varepsilon}\to\norm{\operatorname{curl}G}_2$ as $\varepsilon\to0$, uniformly in $j$.
\item[(b)] Let $\psi\in C^1_c(\R^3)$, $\supp\psi\subset B_{R-1}$, $G=\nabla\psi$, and let $\chi_\varepsilon(y):=\psi\big(y+\tfrac\varepsilon2(e_1+e_2+e_3)\big)$ be the sample of $\psi$ at the centres of the elementary cubes, indexed by their lowest corner $y$. Define $d$ on lattice 2-forms by $(d\sigma)(y):=\partial_1\sigma_{23}(y)-\partial_2\sigma_{13}(y)+\partial_3\sigma_{12}(y)$ (forward differences; $\sigma_{\mu\nu}(x)$ the value on the plaquette $(x;\mu,\nu)$), and $d^*$ as its adjoint for $\ip\cdot\cdot_\varepsilon$. Then $dd=0$ on 1-forms, and $r_\varepsilon:=g_\varepsilon+d^*\chi_\varepsilon$ satisfies $\norm{r_\varepsilon}_{2,\varepsilon}\to0$, uniformly in $j$.
\end{enumerate}
\end{lemma}

\begin{proof}
Write $\partial_\mu$ for the forward difference and $\partial_\mu^*=-\partial^-_\mu$ for its adjoint, minus the backward difference. Then $(dA)_{\mu\nu}=\partial_\mu A_\nu-\partial_\nu A_\mu$ (Section~\ref{ssec:action}).
(a) From $\ip g{dA}_\varepsilon=\sum_{\mu<\nu}\ip{g_{\mu\nu}}{\partial_\mu A_\nu-\partial_\nu A_\mu}=\sum_{\mu,\nu}\ip{g_{\mu\nu}}{\partial_\mu A_\nu}$, with $g_{\nu\mu}:=-g_{\mu\nu}$, we get $(d^*g)_\nu(x)=\sum_\mu\partial_\mu^*g_{\mu\nu}(x)$ on the bond $(x,\nu)$. The plaquettes $(x;\mu,\nu)$ and $(x-\varepsilon e_\mu;\mu,\nu)$ have their centres at the centre of the bond $(x,\nu)$ shifted by $\pm\frac\varepsilon2e_\mu$. So $(d^*g_\varepsilon)_\nu(x)=-\sum_\mu\varepsilon^{-1}\big(g_{\mu\nu}(z+\tfrac\varepsilon2e_\mu)-g_{\mu\nu}(z-\tfrac\varepsilon2e_\mu)\big)$, with $z$ the bond centre. This is a centred difference quotient of a $C^1_c$ function. By the mean value theorem and uniform continuity of $\nabla g$ it equals $(\delta\star^{-1}G)_\nu(z)+o(1)=(\operatorname{curl}G)_\nu(z)+o(1)$, uniformly in $z$. It vanishes unless $z\in B_{R+1}$. Riemann sums of the continuous compactly supported function $\abs{\operatorname{curl}G}^2$ over the bond centres converge to its integral.
(b) $d(dA)=\partial_1(\partial_2A_3-\partial_3A_2)-\partial_2(\partial_1A_3-\partial_3A_1)+\partial_3(\partial_1A_2-\partial_2A_1)=0$, because forward differences commute on the torus. The adjoint is $(d^*\omega)_{23}=\partial_1^*\omega$, $(d^*\omega)_{13}=-\partial_2^*\omega$, $(d^*\omega)_{12}=\partial_3^*\omega$. The two cubes that contain the plaquette $(x;2,3)$ have lowest corners $x$ and $x-\varepsilon e_1$, and their centres are the centre $c_p$ of the plaquette shifted by $\pm\frac\varepsilon2e_1$; similarly in the other two directions. Hence $(d^*\chi_\varepsilon)_{23}(p)=-\varepsilon^{-1}\big(\psi(c_p+\tfrac\varepsilon2e_1)-\psi(c_p-\tfrac\varepsilon2e_1)\big)=-\partial_1\psi(c_p)+o(1)$. In the same way $(d^*\chi_\varepsilon)_{13}=\partial_2\psi+o(1)$ and $(d^*\chi_\varepsilon)_{12}=-\partial_3\psi+o(1)$, uniformly. The sample $g_\varepsilon$ has $g_{23}=\partial_1\psi$, $g_{13}=-\partial_2\psi$, $g_{12}=\partial_3\psi$ at $c_p$. So $\abs{r_\varepsilon(p)}\le\eta(\varepsilon):=\sup_{\abs{y-y'}\le\varepsilon}\abs{\nabla\psi(y)-\nabla\psi(y')}\to0$. Also $r_\varepsilon(p)=0$ unless $c_p\in B_R$. Hence $\norm{r_\varepsilon}^2_{2,\varepsilon}\le3\abs{B_{R+1}}\eta(\varepsilon)^2\to0$.
\end{proof}

\subsection{The Bianchi identity and the spectral density}

\begin{lemma}[Bianchi identity]\label{lem:M-bianchi}
Let $S$ be a limit point and $\psi\in C^1_c(\R^3)$. Then every Schwinger function of $S$ with $\star^{-1}\nabla\psi$ among its curvature arguments vanishes:
\[
S_{1+n,p}(\star^{-1}\nabla\psi,\bg_1,\dots,\bg_n;h_1,\dots,h_p)=0 .
\]
In particular $\mathfrak b_S(\nabla\psi,\nabla\psi)=0$, and for $\psi\in C_c^\infty$, $B(\nabla\psi)=0$ almost surely under the law of $S$ (Theorem~\ref{thm:main}(ii)); that is, $\operatorname{div}B=0$, or $dF=0$.
\end{lemma}

\begin{proof}
Let $Q:=\prod_{i\le n}F(\bg_{i,\varepsilon})\prod_{k\le p}\Xr(h_{k,\varepsilon})$ on the lattice, the test functions being sampled. By Lemma~\ref{lem:M-sampling}(b), $F(g_\varepsilon)=F(r_\varepsilon)-\ip{\chi_\varepsilon}{ddA}_\varepsilon=F(r_\varepsilon)$, where $g_\varepsilon$ is the sample of $\star^{-1}\nabla\psi$. By the Cauchy--Schwarz inequality and Lemma~\ref{lem:M-lattice}(b),
\[
\abs{\ev{F(g_\varepsilon)\,Q}_{\varepsilon,j,\mu_0}}\le\ev{F(r_\varepsilon)^2}^{1/2}\ev{Q^2}^{1/2}\le\norm{r_\varepsilon}_{2,\varepsilon}\,\ev{Q^2}^{1/2}.
\]
$\ev{Q^2}$ is a lattice Schwinger function of order $(2n,2p)$, bounded uniformly by Proposition~\ref{prop:C-bounds}, and $\norm{r_\varepsilon}_{2,\varepsilon}\to0$. The left side converges to the Schwinger function in the statement along the subsequence that defines $S$ (Theorem~\ref{thm:main}(ii); for $S^0$, first for $S^{\mu}$, $\mu\in Q_0$, then as $\mu\downarrow0$). So the limit is $0$. For $\psi\in C^\infty_c$, $\E[B(\nabla\psi)^2]=\mathfrak b_S(\nabla\psi,\nabla\psi)=0$ by the moment identity of Proposition~\ref{prop:C-tight}.
\end{proof}

\begin{lemma}[translation-invariant bounded forms]\label{lem:M-mult}
Let $\mathfrak b$ be a Hermitian form on $C_c(\R^3;\C^3)$ with $0\le\mathfrak b(f,f)\le\norm f_2^2$ and $\mathfrak b(\tau_af,\tau_af')=\mathfrak b(f,f')$ for all $a\in\R^3$. Then there is a measurable $m:\R^3\to\{\text{Hermitian }3\times3\text{ matrices}\}$ with $0\le m(p)\le\mathbf 1$ for a.e.\ $p$ and
\[
\mathfrak b(f,f')=(2\pi)^{-3}\int\overline{\hat f(p)}^{\mathsf T}m(p)\hat f'(p)\,dp\qquad(f,f'\in C_c(\R^3;\C^3)).
\]
Moreover: (i) if $f_0\in C_c$ and $\mathfrak b(f_0,f_0)=0$, then $m\hat f_0=0$ a.e.; (ii) if $a$ is a measurable Hermitian-matrix-valued function with $\abs{a(p)}\le c(1+\abs p^2)$ and $\mathfrak b(f,f)\le(2\pi)^{-3}\int\overline{\hat f}^{\mathsf T}a\hat f$ for all $f\in C^1_c(\R^3;\C^3)$, then $m\le a$ a.e.
\end{lemma}

\begin{proof}
By the Cauchy--Schwarz inequality for the positive form, $\abs{\mathfrak b(f,f')}\le\norm f_2\norm{f'}_2$. So $\mathfrak b$ extends to $L^2$, and $\mathfrak b(f,f')=\ip f{Kf'}$ for a bounded $K$ with $0\le K\le1$ that commutes with all $\tau_a$. Let $\mathcal F=(2\pi)^{-3/2}\,\hat{}$ be the unitary Fourier transform and $\hat K:=\mathcal FK\mathcal F^{-1}$. Then $\hat K$ commutes with multiplication $M_u$ by $u=e_a$, $e_a(p):=e^{-ia\cdot p}$. For fixed $\varphi,\chi\in L^2$, $u\mapsto\ip\varphi{[\hat K,M_u]\chi}=\int u\,\big(\overline{\hat K\varphi}^{\mathsf T}\chi-\overline\varphi^{\mathsf T}\hat K\chi\big)\,dp$ is integration against an $L^1$ function. So the set of $u\in L^\infty$ with $[\hat K,M_u]=0$ is a weak-$*$ closed subspace that contains every $e_a$. Its pre-annihilator in $L^1$ consists of $h\in L^1$ with $\hat h\equiv0$, hence is $\{0\}$, so the subspace is all of $L^\infty$. Let $w(p):=e^{-\abs p^2}$ and $m_{\alpha\beta}:=(\hat K(we_\beta))_\alpha/w$, with $e_\beta$ the unit vectors of $\C^3$. For $u_\beta\in L^\infty$, $\hat K(w\sum_\beta u_\beta e_\beta)=w\sum_\beta u_\beta m_{\cdot\beta}$. So $\hat K f=mf$ on the set of $f=w\sum_\beta u_\beta e_\beta$, which is dense in $L^2$: for $f\in L^2$, $f\mathbf 1\{\abs f\le nw\}$ is of this form and converges to $f$. Testing with $f=wuv$, $u\in L^\infty$ scalar and $v\in\C^3$, gives $0\le\int\abs u^2w^2\,\overline v^{\mathsf T}mv\le\int\abs u^2w^2\abs v^2$ for all $u$. Hence $0\le\overline v^{\mathsf T}m(p)v\le\abs v^2$ for a.e.\ $p$, simultaneously for all $v$ in the countable set $(\Q+i\Q)^3$, hence for all $v$. By polarisation $m(p)$ is Hermitian, so $0\le m\le\mathbf 1$ a.e. Then $M_m$ is bounded and agrees with $\hat K$ on a dense set, so $\hat K=M_m$. This is the representation.
(i) $\ip{f_0}{Kf_0}=0$ and $K\ge0$ give $Kf_0=0$, that is, $m\hat f_0=0$ a.e.
(ii) The form $q(f):=(2\pi)^{-3}\int\overline{\hat f}^{\mathsf T}(a-m)\hat f$ is continuous on the Sobolev space $H^1(\R^3;\C^3)$, and $q\ge0$ on the dense subspace $C^\infty_c$, hence on $H^1$. Take $\hat f=uv$ with $u\in L^\infty$ of compact support and $v\in\C^3$; then $f\in H^1$, and $\int\abs u^2\,\overline v^{\mathsf T}(a-m)v\ge0$ for all such $u$. So $a-m\ge0$ a.e., as before.
\end{proof}

\begin{theorem}[the curvature spectral density]\label{thm:M-density}
\begin{enumerate}
\item[(a)] For every limit point $S$ there is a measurable function $m_S$ with values in the Hermitian $3\times3$ matrices, with $0\le m_S(p)\le P_T(p)$ for a.e.\ $p$, and
\[
\mathfrak b_S(G,G')=(2\pi)^{-3}\int\overline{\hat G(p)}^{\mathsf T}m_S(p)\,\hat G'(p)\,dp\qquad(G,G'\in C_c(\R^3;\R^3)).
\]
In particular $0\le\mathfrak b_S(G,G)\le\norm{\Pi_TG}_2^2$: the curvature two-point function of every limit point is dominated, as a form, by that of the free massless photon.
\item[(b)] If $S$ is at photon mass $\mu_0>0$, then $m_S(p)\le\min\big(1,\abs p^2/\mu_0^2\big)\,P_T(p)$ for a.e.\ $p$.
\item[(c)] For $\mu'<\mu$ in $Q_0$, $m_{S^\mu}\le m_{S^{\mu'}}\le m_{S^0}\le m_{S^\fl}$ a.e.; and $m_{S^{\mu_k}}\to m_{S^0}$ a.e.\ along every sequence $\mu_k\downarrow0$ in $Q_0$.
\end{enumerate}
\end{theorem}

\begin{proof}
(a) For real $G,G'\in C_c$, $\mathfrak b_S(G,G')$ is real and symmetric. It is the limit of the lattice second moments $\ev{F(g_\varepsilon)F(g'_\varepsilon)}$ of the samples. By Lemma~\ref{lem:M-lattice}(b), $0\le\ev{F(g_\varepsilon)^2}\le\norm{g_\varepsilon}^2_{2,\varepsilon}$, and $\norm{g_\varepsilon}_{2,\varepsilon}\to\norm G_2$ (Proposition~\ref{prop:C-compact}, step (iii)). Hence $0\le\mathfrak b_S(G,G)\le\norm G_2^2$. For $S^0$ this passes through the limit $\mu\downarrow0$ (Theorem~\ref{thm:main}(ii)). Extend $\mathfrak b_S$ to complex $f=G_1+iG_2$ as a Hermitian form, $\mathfrak b_S(f,f'):=\mathfrak b_S(G_1,G_1')+\mathfrak b_S(G_2,G_2')+i\big(\mathfrak b_S(G_1,G_2')-\mathfrak b_S(G_2,G_1')\big)$. Then $\mathfrak b_S(f,f)=\mathfrak b_S(G_1,G_1)+\mathfrak b_S(G_2,G_2)\in[0,\norm f_2^2]$. It is translation invariant by Theorem~\ref{thm:C-OS1}(a). Lemma~\ref{lem:M-mult} gives $m_S$ with $0\le m_S\le\mathbf 1$. For $\psi\in C^\infty_c$ real, $\mathfrak b_S(\nabla\psi,\nabla\psi)=0$ by Lemma~\ref{lem:M-bianchi}, and $\widehat{\nabla\psi}(p)=ip\hat\psi(p)$. By Lemma~\ref{lem:M-mult}(i), $m_S(p)p\,\hat\psi(p)=0$ a.e. For $\psi\neq0$, $\hat\psi$ is the restriction of an entire function, so its zero set is Lebesgue-null. Hence $m_S(p)p=0$ a.e. A positive semidefinite Hermitian matrix that annihilates $p$ satisfies $m=P_TmP_T$. So $m_S\le P_T\mathbf 1P_T=P_T$. The last claim is $\ip{\hat G}{m_S\hat G}\le\ip{\hat G}{P_T\hat G}$, integrated.
(b) Let $G\in C^1_c$ be real. By Lemma~\ref{lem:M-lattice}(b),(c), $\ev{F(g_\varepsilon)^2}\le\mu_0^{-2}\norm{d^*g_\varepsilon}^2_{2,\varepsilon}$, and by Lemma~\ref{lem:M-sampling}(a) the right side tends to $\mu_0^{-2}\norm{\operatorname{curl}G}_2^2=(2\pi)^{-3}\mu_0^{-2}\int\abs p^2\,\overline{\hat G}^{\mathsf T}P_T\hat G\,dp$, since $\abs{p\times v}^2=\abs p^2\,\overline v^{\mathsf T}P_T(p)v$ for $v\in\C^3$. Both sides are Hermitian forms induced by real symmetric forms on real fields, so the inequality extends to complex $f=G_1+iG_2$ by adding the inequalities for $G_1$ and $G_2$. Lemma~\ref{lem:M-mult}(ii) with $a=\mu_0^{-2}\abs p^2P_T$ gives $m_S\le\mu_0^{-2}\abs p^2P_T$. Since $m_S=P_Tm_SP_T$ by (a), and $m_S$ is bounded on $\mathrm{ran}\,P_T$ both by $1$ and by $\abs p^2/\mu_0^2$, we get $m_S\le\min(1,\abs p^2/\mu_0^2)P_T$.
(c) For real $\bg$ and real $t$, $e^{tF(\bg)}=\exp[F(t\bg)-\Xr(0)]$ lies in the class $\cD$ ($p=0$). Theorem~\ref{thm:main}(iv) gives $\phi_\mu(t)\le\phi_{\mu'}(t)\le\phi_\fl(t)$ and $\phi_0(t)\le\phi_\fl(t)$, where $\phi_S(t):=G^S(t)$ is the generating functional of the one-element family $\{\bg\}$. Each $\phi_S$ is entire, $\phi_S(0)=1$, and $\phi_S$ is even because its odd Taylor coefficients $S_{n,0}$ vanish (Theorem~\ref{thm:C-OS3}). If $f_1\le f_2$ are smooth with $f_1(0)=f_2(0)$, then $0$ is a maximum of $f_1-f_2$, so $f_1''(0)\le f_2''(0)$. Hence $S^\mu_{2,0}(\bg,\bg)\le S^{\mu'}_{2,0}(\bg,\bg)\le S^\fl_{2,0}(\bg,\bg)$ and $S^0_{2,0}(\bg,\bg)\le S^\fl_{2,0}(\bg,\bg)$. Also $S^0_{2,0}=\lim_{\mu\downarrow0}S^\mu_{2,0}$ (Theorem~\ref{thm:main}(ii)), so $S^{\mu'}_{2,0}(\bg,\bg)\le S^0_{2,0}(\bg,\bg)$. Extend to complex fields as in (b) and apply Lemma~\ref{lem:M-mult}(ii) with the bounded functions $a=m_{S^{\mu'}}$, $m_{S^0}$, $m_{S^\fl}$. This gives the a.e.\ inequalities. Along $\mu_k\downarrow0$ the matrices $m_{S^{\mu_k}}(p)$ increase and are bounded by $\mathbf 1$, off a null set, so they converge a.e.\ to some $m_*\le m_{S^0}$. Let $f=G_1+iG_2\in C_c(\R^3;\C^3)$ with $G_1,G_2$ real. By the representation of Lemma~\ref{lem:M-mult} for the Hermitian extension, $(2\pi)^{-3}\int\overline{\hat f}^{\mathsf T}m_{S^{\mu_k}}\hat f=\mathfrak b_{S^{\mu_k}}(G_1,G_1)+\mathfrak b_{S^{\mu_k}}(G_2,G_2)$, which tends to $\mathfrak b_{S^0}(G_1,G_1)+\mathfrak b_{S^0}(G_2,G_2)=\mathfrak b_{S^0}(f,f)$ (Theorem~\ref{thm:main}(ii)). The integrands are bounded by $\abs{\hat f}^2$, so by dominated convergence the left side tends to $(2\pi)^{-3}\int\overline{\hat f}^{\mathsf T}m_*\hat f$. Hence $\mathfrak b_{S^0}(f,f)=(2\pi)^{-3}\int\overline{\hat f}^{\mathsf T}m_*\hat f$ for all $f\in C_c(\R^3;\C^3)$. Lemma~\ref{lem:M-mult}(ii), applied to $\mathfrak b_{S^0}$ with the bounded function $a=m_*$, gives $m_{S^0}\le m_*$ a.e. Hence $m_{S^0}=m_*$ a.e.
\end{proof}

\subsection{Long-distance strength of the curvature}

\begin{definition}[massless curvature sector]\label{def:M-massless}
For a limit point $S$ and a real $G\in C_c(\R^3;\R^3)$ put
\[
\mathsf Z_S(G):=\limsup_{R\to\infty}R^3\,\mathfrak b_S(G_R,G_R)\ \in[0,\infty).
\]
$S$ is \emph{massless in the curvature sector} if $\mathsf Z_S(G)>0$ for some $G$.
\end{definition}
Since $R^3\norm{G_R}_2^2=\norm G_2^2$, $\mathsf Z_S(G)/\norm G_2^2$ is the asymptotic fraction of the $L^2$ size of a test field at scale $R$ that the two-point function retains. For the free massless photon $R^3\norm{\Pi_TG_R}^2_2=\norm{\Pi_TG}_2^2$ for every $R$, because $P_T$ is homogeneous of degree $0$. For the curvature of a free vector field of mass $M>0$, $m(p)=\abs p^2(\abs p^2+M^2)^{-1}P_T(p)$, and $\mathsf Z\equiv0$ by dominated convergence. If $m_S(p)-z\,P_T(p)\to0$ as $p\to0$ outside a Lebesgue null set (an essential limit; $m_S$ is defined only almost everywhere) for a number $z$, then $\mathsf Z_S(G)=z\norm{\Pi_TG}_2^2$. So $\mathsf Z_S$ is the Euclidean form of the field-strength renormalisation of the photon at zero momentum. The definition is the strong form of masslessness: a density $m(p)=\min(1,\abs p^{\eta})P_T(p)$, $0<\eta<1$, has power-law correlations and $\mathsf Z\equiv0$. The converse of Corollary~\ref{cor:M-summary}(c), from the absence of a gap to $\mathsf Z_S>0$, is not claimed.

\begin{theorem}[the long-distance dichotomy]\label{thm:M-dichotomy}
Let $S$ be a limit point and $G\in C_c(\R^3;\R^3)$ real.
\begin{enumerate}
\item[(i)] For every $R>0$,
\[
R^3\,\mathfrak b_S(G_R,G_R)=(2\pi)^{-3}\int\overline{\hat G(q)}^{\mathsf T}m_S(q/R)\,\hat G(q)\,dq\ \le\ \norm{\Pi_TG}_2^2 .
\]
Hence $0\le\mathsf Z_S(G)\le\norm{\Pi_TG}_2^2$.
\item[(ii)] (no massless component at positive photon mass) If $S$ is at photon mass $\mu_0>0$, then $\lim_{R\to\infty}R^3\mathfrak b_S(G_R,G_R)=0$, so $\mathsf Z_S\equiv0$. For $G\in C^1_c$, quantitatively, $R^3\mathfrak b_S(G_R,G_R)\le\mu_0^{-2}R^{-2}\norm{\operatorname{curl}G}_2^2$.
\item[(iii)] (order) For $\mu\in Q_0$, $0=\mathsf Z_{S^\mu}(G)\le\mathsf Z_{S^0}(G)\le\mathsf Z_{S^\fl}(G)$. In particular, if $S^0$ is massless in the curvature sector, so is $S^\fl$.
\item[(iv)] (integrable correlations exclude masslessness) Suppose that for some real $\psi\in C_c(\R^3)$ with $\int\psi\neq0$ and all $\alpha,\beta\in\{1,2,3\}$, the functions $a\mapsto\mathfrak b_S(\psi e_\alpha,\tau_a\psi e_\beta)$ are in $L^1(\R^3)$. Then $\lim_{R\to\infty}R^3\mathfrak b_S(G_R,G_R)=0$ for every $G$, so $\mathsf Z_S\equiv0$. This holds in particular if $\abs{\mathfrak b_S(\psi e_\alpha,\tau_a\psi e_\beta)}\le Ce^{-M\abs a}$ for some $M>0$.
\end{enumerate}
\end{theorem}

\begin{proof}
(i) Insert $\widehat{G_R}(p)=\hat G(Rp)$ into Theorem~\ref{thm:M-density}(a) and substitute $q=Rp$. The bound is $m_S\le P_T$ and $P_T(q/R)=P_T(q)$.
(ii) By Theorem~\ref{thm:M-density}(b), the integrand in (i) is at most $\min(1,\abs q^2/(R\mu_0)^2)\,\abs{\hat G(q)}^2$, which tends to $0$ pointwise as $R\to\infty$ and is dominated by $\abs{\hat G}^2\in L^1$. For $G\in C^1_c$, use the bound $\abs q^2/(R\mu_0)^2$ and $(2\pi)^{-3}\int\abs q^2\overline{\hat G}^{\mathsf T}P_T\hat G=\norm{\operatorname{curl}G}^2_2$.
(iii) Theorem~\ref{thm:M-density}(c) and (i), together with (ii) for $S^\mu$.
(iv) Since $\widehat{\tau_a\psi}=e^{-ia\cdot p}\hat\psi$,
\[
\phi_{\alpha\beta}(a):=\mathfrak b_S(\psi e_\alpha,\tau_a\psi e_\beta)=\int e^{-ia\cdot p}h_{\alpha\beta}(p)\,dp,
\]
with $h_{\alpha\beta}:=(2\pi)^{-3}\abs{\hat\psi}^2(m_S)_{\alpha\beta}\in L^1$. If $\phi_{\alpha\beta}\in L^1$, the Fourier inversion theorem shows that $h_{\alpha\beta}$ agrees a.e.\ with a continuous function. Since $\hat\psi(0)=\int\psi\neq0$ and $\hat\psi$ is continuous, $\abs{\hat\psi}\ge c>0$ on a ball $B_\rho$. So on $B_\rho$, $m_S$ agrees a.e.\ with a continuous matrix function $\tilde m$. By Theorem~\ref{thm:M-density}(a), $m_S(p)p=0$ a.e.; continuity gives $\tilde m(p)p=0$ for all $p\in B_\rho$. For a unit vector $n$ and $0<t<\rho$, $\tilde m(tn)n=0$, and $t\downarrow0$ gives $\tilde m(0)n=0$. Hence $\tilde m(0)=0$. In (i), split the integral at $\abs q=\rho R$. On $\{\abs q<\rho R\}$ the integrand equals $\overline{\hat G}^{\mathsf T}\tilde m(q/R)\hat G$ a.e.; it tends to $\overline{\hat G}^{\mathsf T}\tilde m(0)\hat G=0$ pointwise and is dominated by $\abs{\hat G}^2$. On $\{\abs q\ge\rho R\}$ it is at most $\abs{\hat G}^2$, whose integral there tends to $0$. Exponential decay implies integrability.
\end{proof}

\subsection{The curvature sector of the reconstruction}

Fix a coordinate direction $j$. Let $\mathcal V^{(j)}_+$ be the complex 2-forms in $C_c$ supported in $\{x_j>0\}$, with $F(g_1+ig_2):=F(g_1)+iF(g_2)$ and $S(F(u)F(v))$ the complex-bilinear extension of $S_{2,0}$. Let $\rho_j\in\Oh$ be the reflection in $\{x_j=0\}$, $\rho_j^*$ its signed action on 2-forms (Theorem~\ref{thm:C-OS2}), and $\theta_j$ the antilinear reflection $\theta_ju:=\overline{\rho_j^*u}$. Put
\[
\ip uv_{+}:=S\big(F(\theta_ju)F(v)\big)\qquad(u,v\in\mathcal V^{(j)}_+).
\]

\begin{proposition}[no gap without integrable correlations]\label{prop:M-gap}
Let $S$ be a limit point.
\begin{enumerate}
\item[(a)] $\ip\cdot\cdot_+$ is a positive semidefinite Hermitian form on $\mathcal V^{(j)}_+$ with $\ip uu_+\le\norm u_2^2$. Let $\cH_F^{(j)}$ be the completion of the quotient by its null space, and $[u]$ the class of $u$. For $t\ge0$, $T(t)[u]:=[\tau_{te_j}u]$ is a well-defined strongly continuous semigroup of self-adjoint contractions, so $T(t)=e^{-tH_F}$ with $H_F\ge0$ self-adjoint. For $y\perp e_j$, $V(y)[u]:=[\tau_yu]$ is a strongly continuous unitary representation commuting with $T$. For $\kappa\in\Oh$ with $\kappa e_j=e_k$, $[u]\mapsto[\kappa_*u]$ is unitary from $\cH^{(j)}_F$ onto $\cH^{(k)}_F$ and intertwines the two semigroups.
\item[(b)] Suppose $H_F\ge M$ for some $M>0$, in one coordinate direction (hence, by (a), in all three). Then for all real 2-forms $\bg,\bg'\in C_c$ supported in $\bar B_r$ and all $a\in\R^3$,
\[
\abs{S_{2,0}(\bg,\tau_a\bg')}\le e^{2Mr}\,e^{-M\abs a/\sqrt3}\,\norm\bg_2\norm{\bg'}_2 .
\]
Consequently the hypothesis of Theorem~\ref{thm:M-dichotomy}(iv) holds, and $\mathsf Z_S\equiv0$.
\end{enumerate}
\end{proposition}

\begin{proof}
(a) For $u=\sum_ic_i\bg_i$ with real $\bg_i$, $\ip uu_+=S(\theta(P)P)$ with $P=\sum_ic_iF(\bg_i)$. This is $\ge0$ by Theorem~\ref{thm:C-OS2}(b). The form is Hermitian because $S$ is real, symmetric and invariant under the reflection $\rho_j\in\Oh$ (Theorem~\ref{thm:C-OS1}(b)). Writing $\mathfrak b(u,v):=S(F(\bar u)F(v))$, Theorem~\ref{thm:M-density}(a) gives $\abs{\mathfrak b(u,v)}\le\norm u_2\norm v_2$, hence $\abs{\ip uv_+}=\abs{\mathfrak b(\rho_j^*u,v)}\le\norm u_2\norm v_2$. For $t\ge0$, $\tau_{te_j}$ maps $\mathcal V^{(j)}_+$ into itself. Since $\theta_j\tau_{te_j}=\tau_{-te_j}\theta_j$ and $S$ is translation invariant, $\ip u{\tau_{te_j}v}_+=\ip{\tau_{te_j}u}v_+$. Let $\norm u_+^2:=\ip uu_+$. The Cauchy--Schwarz inequality gives $\norm{\tau_{te_j}u}^2_+=\ip u{\tau_{2te_j}u}_+\le\norm u_+\norm{\tau_{2te_j}u}_+$. Iterating, and using $\norm{\tau_su}_+\le\norm{\tau_su}_2=\norm u_2$,
\[
\norm{\tau_{te_j}u}_+\le\norm u_+^{1-2^{-n}}\,\norm{\tau_{2^nte_j}u}_+^{2^{-n}}\le\norm u_+^{1-2^{-n}}\,\norm u_2^{2^{-n}} .
\] Letting $n\to\infty$ gives $\norm{\tau_{te_j}u}_+\le\norm u_+$. So $T(t)$ preserves the null space and defines a self-adjoint contraction, with $T(t)T(s)=T(t+s)$. The function $t\mapsto\ip u{T(t)v}_+=\mathfrak b(\overline{\theta_ju},\tau_{te_j}v)$ is continuous, because translations act continuously on $L^2$ and $\mathfrak b$ is bounded. A uniformly bounded semigroup that is weakly continuous on a dense set is weakly continuous, and a weakly continuous contraction semigroup on a Hilbert space is strongly continuous. The spectral theorem gives $H_F$. The transverse translations commute with $\theta_j$, so they preserve the form; they are invertible, hence unitary, and strongly continuous by the same argument. For $\kappa\in\Oh$ with $\kappa e_j=e_k$: $\kappa$ maps $\{x_j>0\}$ onto $\{x_k>0\}$, $\kappa\rho_j\kappa^{-1}=\rho_k$ and $\kappa\tau_{te_j}\kappa^{-1}=\tau_{te_k}$. With the $\Oh$-invariance of $S$ (Theorem~\ref{thm:C-OS1}(b)) this gives $\ip{\kappa_*u}{\kappa_*v}_+=\ip uv_+$, the left side formed in direction $k$ and the right side in direction $j$, and the intertwining.
(b) By (a) we may take the direction $j$ with $\abs{a_j}=\max_i\abs{a_i}\ge\abs a/\sqrt3$, and by the symmetry $S_{2,0}(\bg,\tau_a\bg')=S_{2,0}(\bg',\tau_{-a}\bg)$ we may take $a_j>0$. If $a_j\le2r$, the bound follows from $\abs{S_{2,0}(\bg,\tau_a\bg')}\le\norm\bg_2\norm{\bg'}_2$ and $e^{2Mr}e^{-Ma_j}\ge1$. Let $a_j>2r$, fix $\eta\in(0,a_j/2-r)$, and put $c:=a_j/2$ and $s:=c-r-\eta>0$. By translation invariance $S_{2,0}(\bg,\tau_a\bg')=S_{2,0}(\tau_{-ce_j}\bg,\tau_{a-ce_j}\bg')$. Here $\tau_{-ce_j}\bg$ is supported in $\{x_j\le-s-\eta\}$ and $\tau_{a-ce_j}\bg'$ in $\{x_j\ge s+\eta\}$, so the forms $u,v$ below are supported in $\{x_j\ge\eta\}$ and lie in $\mathcal V^{(j)}_+$. Write $\tau_{-ce_j}\bg=\theta_j\tau_{se_j}u$ and $\tau_{a-ce_j}\bg'=\tau_{a_\perp}\tau_{se_j}v$, where $a_\perp:=a-a_je_j$ and $u,v\in\mathcal V_+^{(j)}$ are real, $\norm u_2=\norm\bg_2$, $\norm v_2=\norm{\bg'}_2$. Using $\theta_j\tau_{se_j}=\tau_{-se_j}\theta_j$ and translation invariance,
\[
S_{2,0}(\bg,\tau_a\bg')=S\big(F(\theta_ju)F(\tau_{2se_j}\tau_{a_\perp}v)\big)=\ip{[u]}{T(2s)V(a_\perp)[v]}_{\cH_F^{(j)}} .
\]
With $\norm{T(2s)}\le e^{-2sM}$, $\norm{[u]}\le\norm u_2$ and $2s=a_j-2r-2\eta\ge\abs a/\sqrt3-2r-2\eta$, we get $\abs{S_{2,0}(\bg,\tau_a\bg')}\le e^{2M(r+\eta)}e^{-M\abs a/\sqrt3}\norm\bg_2\norm{\bg'}_2$; let $\eta\downarrow0$. It is integrable in $a$, and Theorem~\ref{thm:M-dichotomy}(iv) applies with any real $\psi\in C_c$ with $\int\psi\ne0$: take $\bg=\star^{-1}(\psi e_\alpha)$ and $\bg'=\star^{-1}(\psi e_\beta)$.
\end{proof}
\noindent Under the hypotheses of Theorem~\ref{thm:D:recon}, with $\cH$, $H$ and $\Omega$ as there, $[\bg]\mapsto[F(\bg)]$ is an isometry of $\cH^{(3)}_F$ onto the closed span of the vectors $[F(\bg)]$ in $\cH$, and it intertwines $e^{-tH_F}$ with $e^{-tH}$. So $H_F\ge M$ says that $H$ has a gap $M$ on the curvature subspace, which is orthogonal to $\Omega$ because $S_{1,0}=0$. Part (a) of the proposition is the Osterwalder--Schrader construction \cite{OS73,OS75} restricted to the linear curvature fields. Of OS2 the proposition uses only the case of polynomials of degree one; it also uses OS1, OS3 and Theorem~\ref{thm:M-density}(a), but not the hypotheses (H1)--(H3) of Section~\ref{ss:abstract} and not Theorem~\ref{thm:D:recon}.

\begin{corollary}[what masslessness would give]\label{cor:M-summary}
Let $S\in\{S^\fl,S^0\}$ be a limit state of Theorem~\ref{thm:main} or, under its hypothesis, of Theorem~\ref{thm:mainL3}. If $S$ is massless in the curvature sector, then:
\begin{enumerate}
\item[(a)] for every real $\psi\in C_c$ with $\int\psi\neq0$, some $a\mapsto\mathfrak b_S(\psi e_\alpha,\tau_a\psi e_\beta)$ is not integrable on $\R^3$. In particular this bounded function is not $O(\abs a^{-3-\delta})$ for any $\delta>0$, and does not decay exponentially;
\item[(b)] there are $R_k\to\infty$ with $\mathfrak b_S(G_{R_k},G_{R_k})\ge\tfrac12\mathsf Z_S(G)\,R_k^{-3}$, a power-law lower bound for the variance of the block averages $B(G_R)$. A state with integrable curvature correlations has $R^3\mathfrak b(G_R,G_R)\to0$ instead (Theorem~\ref{thm:M-dichotomy}(iv));
\item[(c)] in every coordinate direction, $\inf\operatorname{spec}H_F=0$: the curvature sector of the reconstructed theory has no mass gap.
\end{enumerate}
If $S^0$ is massless, then so is $S^\fl$ (Theorem~\ref{thm:M-dichotomy}(iii)). The states at positive photon mass are not massless in the curvature sector (Theorem~\ref{thm:M-dichotomy}(ii)).
\end{corollary}

\begin{proof}
(a) is the contrapositive of Theorem~\ref{thm:M-dichotomy}(iv), and (c) that of Proposition~\ref{prop:M-gap}(b): no $M>0$ has $H_F\ge M$; in particular $\cH^{(j)}_F\neq\{0\}$, since otherwise $H_F\ge M$ would hold vacuously, and as $H_F\ge0$, $\inf\operatorname{spec}H_F=0$. (b) is the definition of the $\limsup$.
\end{proof}

\begin{remark}[what these properties cannot give]\label{p:rem:M-soft}
The states $S^{\mu_0}$ at positive photon mass satisfy every conclusion of Theorem~\ref{thm:main}(ii)--(iii), the Bianchi identity (Lemma~\ref{lem:M-bianchi}) and the domination $m_S\le P_T$ (Theorem~\ref{thm:M-density}(a)), and yet $\mathsf Z_{S^{\mu_0}}\equiv0$ (Theorem~\ref{thm:M-dichotomy}(ii)). So no argument that uses only these properties of $S^\fl$ or $S^0$ proves that they are massless in the curvature sector. What distinguishes $S^0$ and $S^\fl$ from $S^{\mu_0}$ in this section is the weaker upper bound $m_S\le P_T$ in place of $m_S\le\abs p^2\mu_0^{-2}P_T$, and the comparison $m_{S^\mu}\le m_{S^0}\le m_{S^\fl}$ of Theorem~\ref{thm:M-density}(c), whose lower side is a state at positive photon mass. Section~\ref{p:ss:photon} discusses the lower bound that would be needed.
\end{remark}
\section{Open problems}\label{sec:disc}

We list the questions that this paper leaves open and, for each, the step at which its methods stop. Where a result of this paper depends on one of them, the dependence is an explicit hypothesis of that result (Theorems~\ref{thm:mainL3}, \ref{thm:os4main}(b),(c) and~\ref{thm:NT-cond}).

\subsection{Rotations}\label{ssec:introrot}

The limit states are invariant under $\R^3\rtimes\Oh$ (Theorem~\ref{thm:main}(iii)); rotations outside $\Oh$ are not symmetries of any regularised state and would have to emerge in the limit. Three facts are known. A rotation $R$ with $2^jR$ integral for some $j\ge0$ lies in $\Oh$, so no rotation outside $\Oh$ maps the dyadic lattices of Theorem~\ref{thm:main} into one another (Lemma~\ref{lem:D:rational}(i)). The rotation
\[
R_3=\tfrac13\begin{pmatrix}1&2&2\\2&1&-2\\-2&2&-1\end{pmatrix}
\]
has $3R_3$ integral and infinite order (Lemma~\ref{lem:D:rational}(ii)). The closed subgroup of $\Ot$ generated by $\Oh$ and one rotation of infinite order is $\Ot$, so a limit state invariant under one such rotation is $\Ot$-invariant (Lemma~\ref{lem:D:onerot}).

Invariance under a rotation $R$ of infinite order splits into two steps. (i) \emph{An ultraviolet comparison:} on a fixed torus, at a fixed photon mass $\mu_0>0$, the continuum limits of the model on $\varepsilon\Z^3$ and of the model on the rotated lattice $R\varepsilon\Z^3$ agree, when the same continuum 2-form test function is paired with each lattice in its own frame. King \cite[Thm.~2.4, pp.~326--327]{K2} states a comparison of this kind for the rotation by $\arctan(3/4)$ in a coordinate plane, at blocking factor $5$, with one scalar test function smeared against the plaquettes of both lattices; we do not use it. (ii) \emph{An infrared statement.} By the isometry $x\mapsto R^{-1}x$, the rotated lattice on a cubic torus $T$ is the axis lattice on the tilted torus $R^{-1}T$, with the same action; so (i) says that the continuum theory on the cubic torus, rotated by $R$, is the continuum theory on the tilted torus. Given (i), and with the limits taken in King's order, first $\varepsilon\to0$ on each fixed torus and then the torus to $\R^3$, invariance under $R$ of the limit along cubic tori becomes the statement
\begin{center}
\emph{the infinite-volume limits along cubic tori and along tilted tori agree.}
\end{center}
This statement would follow from uniqueness of the infinite-volume limit, but it is weaker. It cannot be arranged by the choice of tori, since no lattice is invariant under a rotation of infinite order (Lemma~\ref{lem:D:rational}(iii)). Nor does it follow from invariance under translations and $\Oh$, reflection positivity and clustering alone: the centred Gaussian scalar field with covariance
\[
\widehat G(k)=\frac13\sum_{\nu=1}^3\frac1{\abs k^2+ck_\nu^2+m^2},\qquad c>0,
\]
is invariant under translations and $\Oh$, reflection positive about the coordinate planes (each summand is the covariance of a free field with a diagonal metric), and clusters exponentially, but it is not invariant under any rotation of infinite order. Otherwise, by Lemma~\ref{lem:D:onerot}(ii) and the continuity of $\widehat G$, it would be $\Ot$-invariant; but with $a:=1+m^2$,
\[
3\Big(\widehat G(e_1)-\widehat G\big(2^{-1/2}(e_1+e_2)\big)\Big)=\frac1a+\frac1{a+c}-\frac2{a+c/2}>0
\]
by strict convexity of $x\mapsto1/x$.

For $R_3$ the natural family has blocking factor $3$: $R_3$ maps $3\varepsilon\Z^3$ into $\varepsilon\Z^3$ (Lemma~\ref{lem:D:rational}(ii)), so the rotated lattices stay in the family $\varepsilon=3^{-K}\varepsilon_0$. Theorem~\ref{thm:mainL3} gives parts (ii)--(iv) of Theorem~\ref{thm:main} for that family under the hypothesis of a uniform tube bound; a proof of the hypothesis would need the analogues for blocking factor $3$ of Theorem~\ref{thm:T6} and of Lemma~\ref{lem:window}. Neither step (i) nor step (ii) is carried out in this paper.
Invariance under rotations is the only input that would still be missing for the axioms on the ergodic components: by Corollary~\ref{p:cor:rotcor}, if a limit state $S$ is invariant under $\Ot$, then almost every ergodic component of its law satisfies OS0--OS4, with OS0 in the form of Remark~\ref{rem:components}.

A different finite-volume regularisation behaves better in this respect. For the curvature field of the box states built with King's gauge-fixed auxiliary action, invariance under $\R^3\rtimes\Ot$ and clustering on the ergodic components are proved in the companion paper \cite{EIC26}. Whether those box states coincide with the torus states of this paper is open. The only result in this direction, in \cite{EIC26}, is a statement about two-point functions: for limits in which the volume is removed before the spacing, and for photon masses outside a countable set, the curvature two-point functions of the two constructions coincide, so that the curvature two-point function of such limits of the torus states is invariant under $\Ot$. In particular, invariance under rotations outside $\Oh$ is not proved here for the states of Theorem~\ref{thm:main}.

\subsection{Uniqueness, the two massless states, and King's cube states}\label{ss:uniq}\label{ss:S0flat}

The limits of Theorem~\ref{thm:main} are taken along subsequences, and different subsequences might give different states; no uniqueness statement is proved. There are two candidates for the theory with a massless photon: $S^0$, the limit $\mu_0\downarrow0$ of the infinite-volume states at positive photon mass, and $S^\fl$, the infinite-volume limit of the states in which the constant mode of the gauge field is Haar-distributed on its period cell; at fixed lattice and torus the latter is the limit $\mu_0\downarrow0$ of the states at positive photon mass on the gauge-invariant observables (Lemma~\ref{lem:flat}). Both satisfy OS0--OS3 with cubic symmetry, and $S^0\le S^\fl$ on $E^\cD$ (Theorem~\ref{thm:main}(iv)). Whether they coincide is the question whether $\mu_0\downarrow0$ commutes with the limits $\varepsilon\to0$, $T\to\R^3$. By Proposition~\ref{prop:C-ident}(b) the equality on $E^\cD$ would follow from the condition $(\mathrm U\text{-}\mu)_T$ there, that the torus expectations at photon mass $\mu_0$ converge to those of the state $\fl$ as $\mu_0\downarrow0$, uniformly in the spacing and the torus, on the class $\cD$; this uniformity is not known at $e\ne0$. On tori there is no monotonicity in the volume, on which the corresponding argument for cubes would rely.

King's states \cite{K1,K2} are limits in cubes with Neumann conditions for the gauge field and Dirichlet conditions for the scalar field. The correlation inequalities of \cite{BFS1} suggest one-sided comparisons between cube limits and torus limits, but such comparisons would need, besides the existence of King's limits, a decay condition on the difference between the torus coefficient $b_T$ of Definition~\ref{def:counterterm} and King's free-boundary coefficient $a_\natural$, of the type $\sup_K\abs{b_T-a_\natural}\,\abs T\to0$ as $\ell_T\to\infty$; we do not prove them. The identification of the torus limits with King's cube limits, and the cubic symmetry of the latter, would follow from uniqueness of the infinite-volume state. On tori the cubic symmetry is exact at every lattice spacing, which is one reason this paper works with tori. The cluster expansion that gives uniqueness when there is a mass gap is not available in the Coulomb phase.

\subsection{Non-triviality}\label{ss:nontriv}

In every limit state the truncated functions have the signs of Theorem~\ref{thm:NT-sign}, and on every lattice torus the leading coefficient of $M(g;h)=S^T_{2,1}(g,g;h)$ is strictly negative (Proposition~\ref{prop:NT-lead}). Theorem~\ref{thm:NT-cond} derives $M(g;h)<0$ in the limit states at small $e^2+\lambda$, so that they are not decoupled, from two hypotheses: the one-loop statement (C) (Hypothesis~\ref{hyp:NT-C}), a statement about a weighted closed-walk sum and its Brownian-bridge limit, and the remainder bound (R). The step that fails is (R). It would follow from a bound
\[
\big|\ln Z(T,g,h)-\ln Z(T,0,0)-\mathcal P_2(g,h)\big|\le e^2\rho\big((e^2+\lambda)/m\big)\,\kappa(g,h),
\]
uniform in $\varepsilon$, $T$ and $\mu_0\in(0,\mu_*]\cup\{\fl\}$, for sources of the size needed to extract $S^T_{2,1}$ by Cauchy estimates; here $Z(T,g,h)$ is the lattice partition function with the sources $g$ for the curvature and $h$ for the Wick square, and $\mathcal P_2$ is its explicit one-loop part, whose coefficient of order $(2,1)$ in the sources is $e^2c^{\mathrm{free}}_{\varepsilon,j,\mu_0}(g,h)$. This is a renormalisation-group remainder estimate for source-dependent effective actions, with every remainder controlled in terms of the couplings at the unit scale. Neither the uniform bounds of this paper nor \cite{Ba1,Ba2,Ba3,Ba4,K1} provide it: Ba{\l}aban's papers give two-sided bounds on partition functions, and \cite[Thm.~2.1, p.~654]{K1} gives $\abs{\ln Z(T,g,h)}\le C\abs T$ with $C$ depending on $g$, $h$ and $\abs T$, without asymptotics in $(e,\lambda)$. The correlation inequalities of \cite{BFS1} give only non-strict signs, and we know of no inequality of Griffiths--Hurst--Sherman type for the even variable $\abs{\varphi(x)}^2$. No result of this paper excludes a Gaussian limit state: such a state would have $M\equiv0$, which is compatible with every sign in Theorem~\ref{thm:NT-sign}.

\subsection{The massless photon}\label{p:ss:photon}

Above the first-order surface of Section~\ref{ss:higgs} the photon is expected to stay massless in $S^\fl$ and $S^0$: the curvature two-point function should decay like $\abs x^{-3}$, and the Hamiltonian of Theorem~\ref{thm:D:recon} should have no gap above the joint vacuum. In the notation of Section~\ref{sec:massless} this is $\mathsf Z_S>0$ (Definition~\ref{def:M-massless}), in the uniform form
\[
(\mathrm{NS})\qquad \mathfrak b_S(G_R,G_R)\ge\delta R^{-3}\quad\text{for all }R\ge R_0,
\]
for some real $G\in C_c(\R^3;\R^3)$, $\delta>0$ and $R_0$. Even $\mathfrak b_S\not\equiv0$ is not known. What is proved is one-sided (Theorem~\ref{thm:M-dichotomy}, Corollary~\ref{cor:M-summary}): the curvature two-point function of every limit state is dominated by that of the free massless photon, the states at positive photon mass have no massless component, and masslessness would exclude integrable curvature correlations and a mass gap in the curvature sector. Every estimate of this paper on the curvature is an upper bound, or a comparison with the states at positive photon mass, which have every property used in Section~\ref{sec:massless} and $\mathsf Z\equiv0$ (Theorem~\ref{thm:M-dichotomy}(ii)); so these properties alone cannot prove masslessness. What is needed is a lower bound on the curvature two-point function: at long wavelengths and uniformly in the scale for $\mathsf Z_S>0$, or at one fixed scale, uniformly in the spacing and the torus, for $\mathfrak b_S\not\equiv0$. By Lemma~\ref{lem:M-lattice}(b) such a bound says that the closed currents of the matter do not completely screen long-wavelength magnetic fields. For $S^0$, and with it $S^\fl$, a bound $R^3\mathfrak b_{S^\mu}(G_R,G_R)\ge\delta$ for $R_0\le R\le c/\mu$, uniform in $\mu\in Q_0$, would suffice (Theorem~\ref{thm:M-density}(c)); by the same comparison, $\mathfrak b_{S^\mu}\not\equiv0$ for one $\mu\in Q_0$ would give $\mathfrak b_{S^0}\not\equiv0$ and $\mathfrak b_{S^\fl}\not\equiv0$, and this is not known either. At small $e^2+\lambda$, $\mathfrak b_S\not\equiv0$ would also follow from (R) and (C): if $\mathfrak b_S(G,G)=0$, then $F(g)=0$ almost surely for $g=\star^{-1}G$, so $M(g;h)=0$ for every $h$, which Theorem~\ref{thm:NT-cond} excludes when $g$ has components in $C_c^\infty$, $d^*g\not\equiv0$, $0\le h\not\equiv0$ and $e^2+\lambda$ is small enough, depending on $(g,h)$.

\subsection{The rest of OS4}\label{p:ss:OS4open}

Theorem~\ref{thm:cluster} gives, on almost every ergodic component of the law of a limit state, ergodicity under $\R^3$, reflection positivity, clustering in the mean and under diagonal separation, and a unique joint vacuum in each coordinate direction. Theorem~\ref{thm:os4main} gives clustering along the axes for every observable linear in the curvature, with any square-integrable dependence on the Wick square; for the Wick-square sector if $\Delta_S=0$; and for all observables, with OS4, under hypothesis (D). OS4 itself (Definition~\ref{def:OS4}) is not proved, for the limit states or for their components. For a component it is the statement that $\Ker H_j=\C\Omega_j$ for $j=1,2,3$ in the reconstruction of Proposition~\ref{prop:D:recon} (by Proposition~\ref{prop:D:fixed}(i) and translation invariance), and by Proposition~\ref{prop:closeaxis} it suffices to show $P_jG=\E[G\mid\cI]$ for $G$ in a set with dense span. Proposition~\ref{prop:axisobs} shows that (H1)--(H3), $\Oh$-invariance and ergodicity alone do not give it. Three steps fail.

(a) \emph{The latent density.} $\Delta_S=0$, which gives hypothesis (L) of Theorem~\ref{thm:os4:matter}, needs a bound on the torus susceptibilities of the Wick square uniform in the regulators (Remark~\ref{rem:os4:Lopen}). It is expected to fail on the first-order surface of Remark~\ref{rem:disc:firstorder}, where the limit pressure of the Wick square should have a corner at the reference mass and a limit state may be a mixture of two coexisting phases; there only the bound $\Delta_S\le\kappa_++2S^B_\fl$ of Proposition~\ref{prop:os4:latent} remains, and its constants are not shown to be uniform as $\lambda\downarrow0$. Convexity suggests a way around this step away from exceptional points. A constant tilt $\sigma\XR(1_T)$ of a torus state is the same model with $b_T$ replaced by $b_T-2\sigma$, and its pressure is $\alpha(\sigma+\cdot)-\alpha(\sigma)$; since the convex function $\alpha$ is differentiable off a countable set, limits of the tilted states would have latent density zero for all but countably many $\sigma$, provided the analysis of Sections~\ref{sec:C-limits}--\ref{sec:OS4} is carried over to the tilted family, which we do not do; and $\sigma=0$ may be exceptional. The analogous argument in the mass would need, in addition, regulator data of Lemma~\ref{lem:window} chosen uniformly on compact sets of masses and two-sided Lipschitz bounds on the torus coefficient as a function of the mass, which we do not verify; its exceptional set depends on the subsequence and is expected to contain the mass of the first-order transition.

(b) \emph{Higher degree in the curvature.} For observables of degree at least two in the curvature the axis-invariant parts are governed by the clustering of the current sector, for which there is no estimate (Remark~\ref{rem:os4:Fopen}); a Gaussian domination of the linear field does not control them (Example~\ref{ex:os4:telegraph}).

(c) \emph{Hypothesis (D).} Hypothesis (D) of Corollary~\ref{cor:os4:OS4} excludes a non-trivial ergodic decomposition and is a uniqueness-type statement (Remark~\ref{rem:os4:D}).

Two problems concern the components themselves. Let $\mu=\int\mu_\xi\,d\nu(\xi)$ be the decomposition of Theorem~\ref{thm:cluster} of the law of a limit state $S$, let $P_j$ be the orthogonal projection of $L^2(\mu)$ onto the functions invariant under the translations along $e_j$, and $\cI$ the $\sigma$-algebra of translation-invariant Borel sets.

\emph{Single-axis ergodicity of the components.} Show that for $\nu$-almost every $\xi$ the translations along some coordinate direction $e_j$ act ergodically on $\mu_\xi$, that is, $\mu_\xi$ satisfies (OS4$_j$) of Definition~\ref{def:OS4}; equivalently (Proposition~\ref{add:OS4rel}), $\Ker H_j=\C\Omega_j$ in the reconstruction from $\mu_\xi$ in direction $j$, or $P_jG=\E[G\mid\cI]$ for all $G\in L^2(\mu)$. Show also the cluster property of Definition~\ref{add:OS4}(i) for the components. The identity $P_jG=\E[G\mid\cI]$ holds for the observables of degree one in the curvature, with any square-integrable dependence on the Wick square (Theorem~\ref{thm:os4main}(a)), and for the Wick-square fields if $\Delta_S=0$ (Theorem~\ref{thm:os4main}(b)). It does not follow from (H1)--(H3), cubic invariance and ergodicity under $\R^3$ alone (Proposition~\ref{prop:axisobs}), and it would follow from invariance of $S$ under rotations (Corollary~\ref{p:cor:rotcor}). The conditions $\Delta_S=0$ and (D) are sufficient but stronger than needed: (D) makes $\mu$ itself ergodic, and $\Delta_S=0$ makes all components agree on the Wick-square fields (Theorem~\ref{thm:os4:matter}); both are expected to fail on the first-order surface of Remark~\ref{rem:disc:firstorder}, where the components may still satisfy (OS4$_j$).

\emph{Cubic symmetry and charge conjugation of the components.} Show that $\nu$-almost every $\mu_\xi$ is invariant under $\Oh$; equivalently, that the integer $k$ of Theorem~\ref{thm:cluster}(f) equals $1$ for $\nu'$-almost every $\zeta$; equivalently, that every bounded translation-invariant function is invariant under the coordinate permutations. Show also that almost every $\mu_\xi$ is invariant under charge conjugation, which acts on $\mathcal H^{-r}$ by changing the sign of the curvature component, so that the moments of $\mu_\xi$ of odd degree in the curvature vanish; in degree one this holds (Theorem~\ref{thm:os4:FH}). The cubic symmetry of the components does not follow from (H1)--(H3) and the cubic symmetry of the law (Example~\ref{add:Ohorbit}). The coordinate reflections act trivially on translation-invariant functions, by reflection positivity (Lemma~\ref{lem:D:invloc}); for a coordinate permutation the same argument would need reflection positivity about the diagonal planes $\{x_k=x_l\}$, which is not available for these states. Invariance of $S$ under rotations (Corollary~\ref{p:cor:rotcor}) or ergodicity of $\mu$ would give the cubic symmetry of the components.

The translation-ergodic components are not known to be $\Oh$-invariant (Remark~\ref{rem:D:Ohcomp}); the decomposition may instead be chosen with components invariant under $\R^3\rtimes\Oh$, each a uniform mixture of the $k\in\{1,2,3,6\}$ translation-ergodic measures of one orbit of the coordinate permutations (Theorem~\ref{thm:cluster}(f)), and for a law satisfying (H1)--(H3) the cubic symmetry of the law alone does not force $k=1$ (Example~\ref{add:Ohorbit}). The translation-ergodic components satisfy OS0 only with constants that depend on the component (Corollary~\ref{cor:D:schwinger}(iv)), and the vanishing of the odd curvature moments, OS3 for the limit states, is known for a component only in degree one (Theorem~\ref{thm:os4:FH}). No rate of decay is proved for any correlation.

\subsection{Other phases and block sizes}\label{ss:higgs}

The proofs of this paper do not use the phase. Sections~\ref{sec:C-lattice}--\ref{sec:recon} use the uniform tube bound, the exact lattice symmetries of tori and inequalities that hold on every lattice torus, and the inputs that give the tube bound (Sections~\ref{sec:UV}--\ref{sec:chain}) refer to the coupling range of Theorem~\ref{thm:main}, not to a phase. One-loop perturbation theory predicts a first-order transition inside that range (Remark~\ref{rem:disc:firstorder}); ``symmetric phase'' in the title refers to the couplings above the predicted surface. In the Higgs phase $m^2<0$ the ultraviolet analysis down to the Higgs scale must be redone (the effective action of \cite{BIJ}), and a unit-scale cluster expansion in unitary gauge, of the kind that \cite{BBIJ} carried out on the unit lattice, would be expected to give uniqueness of the infinite-volume limit, exponential clustering and, through the agreement of cubic and tilted torus limits, full rotation invariance; none of this is attempted here. Blocking factors $L\ge4$, in particular King's $L=5$, are outside the scope of the proof: the argument by which Ba{\l}aban's Proposition~1 is applied under the weakened support hypothesis (R4$'''$) needs $L\le3$ (Remark~\ref{rem:T6-textfixes}(iii)).

\begin{remark}[a predicted first-order surface]\label{rem:disc:firstorder}
The following is one-loop perturbation theory, which we do not justify. In the normalisation~\eqref{eq:action} a constant background $\abs\varphi=\rho$ gives the photon the mass $\abs e\rho$; we take $e>0$, since $(A,e)\mapsto(-A,-e)$ leaves the action invariant. In three dimensions the two polarisations of a vector field of mass $M$ add $-M^3/(6\pi)$ to the effective potential (Landau gauge; the term linear in $M^2$ is absorbed by the mass renormalisation); the scalar loops are analytic in $\rho^2$ for $m>0$ and negligible here. So, to one loop,
\[
V(\rho)=\tfrac12m^2\rho^2-\frac{e^3}{6\pi}\rho^3+\lambda\rho^4 .
\]
$V$ has a minimum at some $\rho_*>0$ degenerate with the one at $\rho=0$ exactly when the quadratic polynomial $V(\rho)/\rho^2$ has a double root, $(e^3/6\pi)^2=2m^2\lambda$, that is on the surface
\[
\lambda m^2=\frac{e^6}{72\pi^2}\,\big(1+O(e^2/m)\big),
\]
where the correction stands for the counterterm scheme, the photon-mass regulator and higher loops. There $\rho_*=6\pi m^2/e^3$ and the photon mass is $e\rho_*=6\pi m^2/e^2$, so the loop parameter $e^2/(e\rho_*)=e^4/(6\pi m^2)$ is below $10^{-3}$ when $e^2<0.1m$. For every $m$ and every $e$ with $0<e^2<0.1m$ the point $\lambda=e^6/(72\pi^2m^2)$ of the surface lies in the coupling range of Theorem~\ref{thm:main}: $\lambda<10^{-3}m/(72\pi^2)\approx1.4\cdot10^{-6}\,m$, that is $\lambda/e^2<1/(7200\pi^2)$, deep in the type-I regime $\lambda\ll e^2$; so does the region below it, since the range bounds $\lambda$ from below only by $\lambda>0$. Equivalently, for fixed $e$ and $\lambda<e^2/(7200\pi^2)$ the transition is at the mass $m_c(e,\lambda)=\abs e^3/(6\pi\sqrt{2\lambda})$. One loop thus predicts the symmetric, Coulomb phase above the surface ($\lambda m^2>e^6/(72\pi^2)$); on it a first-order transition in the mass, with the density of the Wick square jumping by about $\rho_*^2=36\pi^2m^4/e^6$; and below it, still at $m^2>0$, a phase with a massive photon. The theorems of this paper are not affected, since they do not use the phase; their reading is.
\end{remark}

\subsection{Extremality and relativistic fields}\label{p:ss:extremal}

Extremality of the limit states is not proved. Theorem~\ref{thm:cluster} decomposes their laws into translation-ergodic components, which are invariant under the coordinate reflections but are not known to be invariant under $\Oh$ (Section~\ref{p:ss:OS4open}; Theorem~\ref{thm:cluster}(f) gives $\Oh$-invariant components whose joint vacuum has dimension $k\in\{1,2,3,6\}$), and under hypothesis (D) the law is ergodic (Theorem~\ref{thm:os4main}(c)); whether the law of a limit state is ergodic, equivalently whether its joint vacuum is unique (Theorem~\ref{thm:D:recon}(iv)), is a uniqueness-type question of the kind in Section~\ref{ss:uniq}. The reconstruction of Theorem~\ref{thm:D:recon} has a positive Hamiltonian, unitary spatial translations and the point group $D_4$, but no Lorentz covariance, which would require invariance under the Euclidean rotations in the planes that contain the time axis (Section~\ref{ssec:introrot}). OS0 holds in the form E0$''$ of \cite[(4.2), p.~287]{OS75} and implies E0$'$ by the Appendix of \cite{OS75}, which uses no rotations (Section~\ref{sec:recon}). Under E0$'$ together with E1 and E2, \cite{OS75} continues the Schwinger functions analytically to tempered Wightman-type distributions with the spectrum condition and translation covariance. Whether that argument can be run with the translations, $D_4$ and reflection positivity alone has not been examined, and no Wightman-type distributions are claimed for the limit states.

\appendix

\section{Regulator data and the order of choices}\label{app:B}

This appendix fixes the regulator data of the construction (the blocking factor, the parameters of Ba{\l}aban's renormalisation group, the spacing $\eps_0$ of the coarsest lattice and the side $\ell_0$ of the smallest torus), lists the conditions they have to satisfy, and proves that these conditions can be met at every point of the window $\fW$ of \eqref{eq:window} (Lemma~\ref{lem:window}). Theorem~\ref{thm:T6}, Theorem~\ref{thm:ch-T9} and the results of Appendices~\ref{app:G1:area-T6}--\ref{app:G1:part-D6} are proved under these conditions. The conditions are named by type: (P.$\cdot$) conditions on the point of the window, (R.$\cdot$) lower bounds on $R$, (b.$\cdot$) conditions on $b_0$, (E.$\cdot$) upper bounds on $\eps_0$, and (T.$\cdot$) lower bounds and divisibility conditions on $\ell_0$. Throughout, $L=2$, $d=3$, $\hat\ell:=1+\log\eps_0^{-1}$, $\ell(t):=1+\log t^{-1}$, $r(t)=R\,\ell(t)^{\varrho}$ with $\varrho=2$ is Ba{\l}aban's localisation radius at scale $t$ \cite[(2.7), p.~558]{Ba2}, $p(t)=b_0\ell(t)^{p}$ is his small-field parameter, $t_k=L^k\eps$ ($0\le k\le K$, so that $t_K=\eps_0$), $a_\flat:=a(1-L^{-2})$, $\hat x:=\abs e\,\eps_0^{1/2}$, and $\mathbb M_4:=\{M_1,m,m_1,m_2\}$ with $M_1=0.7m$. The constant $c_0$ is the constant of stage~3 below; the constant called $c_0$ in \cite[p.~593]{Ba2} is written $c_\square$.

\subsection{The window and the regulator data}\label{p:appA:window}

Let $\pi=(e,\lambda,m,\eta,m_1,m_2)\in\fW$, and put $\mu_*:=m$, $M_1:=0.7m$ and $H:=m^2$. The regulator data at $L=2$ are: $p=6$; $\nbar=13$ and the order cap $\Nbar=\nbar^2=169$, with some decay rate $\delta(\Nbar)>0$ of the families at the cap (this is not Ba{\l}aban's $\delta_0$, which depends on $d$ only \cite[p.~420]{Ba3}); the block size $M_b$; the localisation parameter $R\ge R^{**}$ of Table~\ref{p:tab:A2}, with $\delta_0R\ge4$; the small-field parameter $b_0$; the coarsest spacing $\eps_0=2^{-n_0}$; and the side $\ell_0\in2M_b\N$ of the smallest torus. The tori and spacings of the construction are then $T_j:=j!\,T_0$ (side $\ell_j=j!\,\ell_0$) and $\eps=2^{-K}\eps_0$ with $K\ge n_0+1$; this definition imposes no condition, and $\eps\le1/m$ for all of them follows from condition (E.m2) below. Three rules are part of the conditions:
\begin{itemize}
\item[(i)] \emph{the lower bound on $R$}: $R\ge R^{**}$, which dominates every other lower bound on $R$ of Table~\ref{p:tab:A2};
\item[(ii)] \emph{the four-mass rule}: every upper bound on $\eps_0$ is imposed at each Higgs mass $M\in\mathbb M_4$ at which it is used, and every constant or condition of the stages 8--11 below that depends on the Higgs mass is evaluated at each $M\in\mathbb M_4$, the four values being combined by the maximum for constants and for lower bounds, and by the minimum for upper bounds on $\eps_0$;
\item[(iii)] \emph{the torus floor}: $\ell_0\ge\max\{96/m,\,128/m,\,128/M_1\}$ and $\ell_0$ is above every lower bound of Table~\ref{p:tab:A5}; since $128/M_1=182.857\ldots/m$, the choice $\ell_0\ge183/m$ meets the first three bounds, and $M_1$ is among the masses of stage~0, where the floor $\ell_0^\natural:=\max\{8,128/m_1,128/M_1\}=128/M_1$ is fixed.
\end{itemize}
Two conditions are produced by Section~\ref{sec:T6}: the upper bound (E.edge) of \eqref{eq:lem24-cap} on $\eps_0$ (stage 10b), from Lemma~\ref{lem:lem24}, and the lower bound (T.per) of \eqref{eq:GIMG} on $\ell_0$ (stage 11), from Lemma~\ref{lem:IMG}.

\subsection{The stages of the choice}\label{p:appA:stages}

The regulator data are chosen in the order of Table~\ref{p:tab:stages}. Each stage depends only on data fixed in earlier stages, apart from steps inside one stage that depend on earlier data and are not used again inside that stage: the lower bounds on $R$ give $R$ (stage 5), the upper bounds on $\eps_0$ give $\eps_0$ (stage 10b), the lower bounds on $\ell_0$ give $\ell_0$ (stage 11), and inside the stages 0 ($\kappa$ and $\ell_0^\natural$ from the point), 6--8 (constants from $p$ and $\mathbb W_0$) and 10a (suprema over $(0,\eps_1]$ from $\eps_1$) the dependence goes one way.

\begin{table}[ht]
\small
\begin{tabular}{@{}p{0.08\textwidth}p{0.72\textwidth}p{0.12\textwidth}@{}}
\hline
stage & choice & depends on \\
\hline
point & $(e,\lambda,m)$; $\mu_*:=m$; the torus counterterm scheme of Definition~\ref{def:counterterm}, with one Wick constant $z_\eps$; the torus side $96/m$ of Theorem~\ref{thm:ch-W}; $M_1:=0.7m$ and $\sigma_+$; $H:=m^2$; the anchor masses $m_1,m_2$ and $\eta\in(0,1)$ & --- \\
0 & $\pi=(e,\lambda,\mu_*,\kappa,\eta,M_1,m_1,m_2)$, with $\kappa$ the domination constant of Theorem~\ref{app:G1:DOM-nat} under the a~priori bound $\mu_*\eps_0\le0.2$ (Remark~\ref{app:G1:DOM-kappa}); the floor $\ell_0^\natural:=\max\{8,128/m_1,128/M_1\}=128/M_1$ & point \\
1--2 & $a=1$, $L=2$; the coefficients $a_k$, through the $k$-uniform bound $a_\flat<a_k\le a$ & 0 \\
3 & $M_b$; the constants $\delta_0,R_0,c_0,c_1,\dots$ of \cite{Ba1,Ba2,Ba3} (functions of $d$, $a$, $L$, $M_b$) & 0--2 \\
4 & $\nbar=13$, $\Nbar=\nbar^2=169$, a decay rate $\delta(\Nbar)>0$ of the families at the order cap; the constants $C_{\mathrm{loc}}(\Nbar)$ and $C_{\Nbar}$ of the conditions (R.cap) and (T.cap); the width $\ell_{\Nbar}$ of Subsection~\ref{ssec:T6-prop1}; the collar exponent $\kappa_0:=\frac15$ (Lemma~\ref{app:G1:kappa0}) & 0--3 \\
5 & $R\ge R^{**}$, $\delta_0R\ge4$, and the condition (R.cap), $R\hat\ell^2\ge C_{\mathrm{loc}}(\Nbar)$ (met by $R\ge C_{\mathrm{loc}}(\Nbar)$, since $\hat\ell\ge1$): lower bounds on $R$ only & 0--4 \\
6--8 & $p=6$ (so that $2p>\varrho d+2$, as Proposition~\ref{app:G1:I6} needs), $p_1=1$; $\mathbb W_0=[1-\eta,2]$ and its constants; every constant that depends on the Higgs mass, at each $M\in\mathbb M_4$ (maximum) & 0--5 \\
9 & $b_1=1$; $b_0$ above its lower bounds (Table~\ref{p:tab:A3}), all of them free of $\eps_0$ & 0--8 \\
10a & $\eps_1:=\min\{1,\eps_{(\mathrm{E.x})},\eps_{(\mathrm{E.md})},\eps_{(\mathrm{E}.\delta)},\eps_{(\mathrm{E.dom})},\eps_{(\mathrm{E.14})}\}$, where $\eps_{(\cdot)}$ is the largest $\eps_0\le1$ at which the condition $(\cdot)$ of Table~\ref{p:tab:A4} holds; every constant that depends on $\eps_0$ is replaced by its supremum over $(0,\eps_1]$ & 0--9 \\
10b & $\eps_0:=$ the largest $2^{-n_0}\le\eps_1$ that meets every upper bound of Table~\ref{p:tab:A4} at each $M\in\mathbb M_4$ (minimum), among them (E.r): $\eps_0<\mathrm{e}^{1-\varrho}=\mathrm{e}^{-1}$, (E.ent), (E.edge), (E.5), (E.ksp1)--(E.ksp3), and (R.cap) if it is not met at stage 5 & 0--10a \\
11 & $\ell_0\in2M_b\N$ above every lower bound of Table~\ref{p:tab:A5}, (T.cap) and (T.per) included; $T_j:=j!\,T_0$; $\eps=2^{-K}\eps_0$ & 0--10b \\
\hline
\end{tabular}
\caption{The order of choices of the regulator data.}\label{p:tab:stages}
\end{table}

The lower bound $b_0^2\hat\ell^{12}\ge\max\{2,\bar\mu\lambda^{-1/2}\}$, which is part of the hypotheses of \cite[(2.12)]{Ba2}, is imposed at stage~9 in the $\eps_0$-free form (b.3) of Table~\ref{p:tab:A3}, $b_0^2\ge\max\{2,\bar\mu\lambda^{-1/2}\}$, which implies it since $\hat\ell\ge1$; the second part, $\lambda^{-1/2}\eps_0^{-1/2}b_0^2\hat\ell^{11}\ge C_{2.12}$, is the downward-closed upper bound (E.6) on $\eps_0$ at stage 10b. Imposed in the other way, the first part would make $b_0$ depend on $\eps_0$ and close the cycle $\eps_0\to b_0\to\text{(E.edge)}\to\eps_0$.

\subsection{The conditions}\label{p:appA:conditions}

\begin{definition}[the conditions on the regulator data]\label{app:G1:rowlist}
The \emph{conditions} of this appendix are those of Tables~\ref{p:tab:A1}--\ref{p:tab:A5}, together with the rules (i)--(iii) of Subsection~\ref{p:appA:window}. The last column of each table lists the stages of Table~\ref{p:tab:stages} whose data the condition depends on.
\end{definition}

In the tables, letters other than the regulator data, the point $\pi$ and the quantities defined in this appendix denote constants of the estimates of \cite{Ba1,Ba2,Ba3} in which the condition is used (for instance $c_v$, $e_*$, $C_{\mathrm{reg}}$, $C_{2.12}$, $C_{2.21}$, $C_H'$, $C_{H''}$, $\delta'$, $\gamma_1$, $C_Z$, $\mathrm{Err}$, $c_M$, $C_d$, $c_b$, $b_*$, $C_A^{R1}$, $\rho_1$, $C_\alpha$, $\alpha$, $\alpha_I$, $C_\delta(M)$, $c_0^{\mathrm{ker}}$, $\bar\mu$, $c_1^I$, $C_\partial$, $c_{\mathrm{reg}}$, $C_W$, $c$, $C$; a subscript such as $2.21$ names the display \cite[(2.21)]{Ba2}); each of them depends only on the data of the stages in the last column, and none depends on $\eps$, $K$, $\ell_T$ or $u$. The constants $\delta_1,c_1',c_2,\delta_h,\rho,M_\theta$ are further constants of \cite{Ba1,Ba2,Ba3} fixed at stage~3, with $c_1\le3$, $c_2\le9$, $M_\theta\le LM_b$, and $b_d$ is a number depending on $d$ only. The constants $C_H,\delta_H$ of (E.5) are defined in Lemma~\ref{app:G1:RW-E06}. We write $e(t):=\abs e\,t^{1/2}$.

\begin{table}[ht]
\small
\begin{tabular}{@{}p{0.08\textwidth}p{0.70\textwidth}p{0.14\textwidth}@{}}
\hline
name & condition & depends on \\
\hline
(P.1) & $\lambda>0$, $e\ne0$, $m>0$ & point \\
(P.2) & $\pi\in\fW$: $0<\lambda<0.1m$, $0<e^2<0.1m$, $0<\eta<1$, $2m<m_1<m_2$ & point \\
(P.3) & Theorem~\ref{thm:ch-W}(a) at $\mu_*=m$, $M_1=0.7m$, on tori of side $\ge96/m$ and for $\eps\le1/m$; in units $m=1$ it follows from $\Pi(\lambda,e^2):=0.51+0.256576\lambda-0.088629e^2-7.1791\lambda^2-9.3304e^2\lambda-1.7038e^4>0$, and on the closed box $[0,0.1m]^2$ it gives $2\sigma_+\ge0.319m^2$ & point \\
(P.4) & Theorem~\ref{thm:ch-W}(c): $m_1\ge m_1^c$, with $m_1^c\le1.882m$ (Corollaries~\ref{app:G1:CorC} and \ref{app:G1:Wc}) & point \\
(P.5) & Theorem~\ref{thm:ch-W}(b): $m_2>\max\{m_1,S(m_1)-m_1\}$; $g_{21}\ge(m_2-m_1)(m_2+m_1-0.3631m)>0$ & point \\
(P.6) & Theorem~\ref{thm:ch-W}(d): $m_1\ge\sqrt2\mu_*$; $m<m_1<m_2$; $M_1<m$; $\eta\in(0,1)$ fixed & point \\
\hline
\end{tabular}
\caption{Conditions on the point $\pi$. They hold at every $\pi\in\fW$ by Theorem~\ref{thm:ch-W}; $S(m_1)$ is any number $\ge\sup_{M\ge m_1}p(M)$, with $p(M)$ the majorant of Proposition~\ref{app:G1:R}.}\label{p:tab:A1}
\end{table}

\begin{table}[ht]
\small
\begin{tabular}{@{}p{0.08\textwidth}p{0.70\textwidth}p{0.14\textwidth}@{}}
\hline
name & condition & depends on \\
\hline
(R.1) & $R>R_0$; $\delta_0R\ge2L=4$ & 3 \\
(R.2) & $R\ge107$; $R\ge3R_0+15$ & 3 \\
(R.3) & $R\ge\ell_{\Nbar}+4\sqrt3+2$; $R\ge2\ell_{\Nbar}+8M_b+8$; $R\ge2\ell_{\Nbar}+10M_b+10$ & 3, 4 \\
(R.4) & $R\ge R_0+2\sqrt3+\delta_h+1+\rho$; $2r(\eps_0)\ge R_0+2\sqrt3+1$ & 3 \\
(R.5) & $R\ge C_{\mathrm{loc}}(\Nbar)$; $R\ge R_0+13+c_1'$; $R\ge M_b+2$, $R\ge c_1'$, $R\ge c_2/4$ & 3, 4 \\
(R.6) & $R\ge[L(M_\theta+6)+6]/5$ (implied by (R.3)) & 3 \\
(R.7) & $r(\eps_0)-2M_b\ge R_0+2\sqrt3+1$; $r(\eps_0)>3M_b$, $R\hat\ell^2/4\ge2M_b+2$ & 3, 5 \\
(R.8) & $R\ge14b_d+1$, $R\ge(11\sqrt3/2)M_b$, $R\ge\max\{12\sqrt3,14M_b\}$; $r(t)\ge\ell_{\Nbar}+4\sqrt3+2$ for $t\le\eps_0$ & 3, 4 \\
(R.e) & $R\ge4\ell_{\Nbar}+32M_b+25$ (Corollary~\ref{cor:prop1R4}) & 3, 4 \\
(R.edge) & the three conditions \eqref{eq:lem24-R0row}, imposed in the sufficient form $r(\eps_0)\ge R$, which they follow from (Lemma~\ref{app:G1:R0rows}) & 3, 5 \\
(R.cap) & $R\hat\ell^2\ge C_{\mathrm{loc}}(\Nbar)$, imposed in the sufficient form $R\ge C_{\mathrm{loc}}(\Nbar)$ (Lemma~\ref{app:G1:RW-B20}) & 4, 5 \\
(R.nbar) & $2\nbar+1\le3R$ (at $\nbar=13$ it follows from (R.2)); used for the non-wrapping of the lifted clusters at zero background (Theorem~\ref{app:K:L17-1}) & 4 \\
(R.Z) & $\delta_5R\ge d\log L+\log2$, with $\delta_5$ the decay rate of Proposition~\ref{app:K:L13-7}; used for the expansion of the normalisations $Z^{(j)}$ & 5 \\
\hline
\end{tabular}
\caption{Lower bounds on $R$ (stage 5). None depends on $\eps$, $k$, $\ell_T$, $u$ or the Higgs mass. In (R.4), (R.7), (R.8), (R.edge) and (R.cap) the quantities $r(\eps_0)$, $r(t)$ and $\hat\ell$ are replaced by their lower bounds $R$, $R$ and $1$ (valid for $\eps_0\le1$); the forms so obtained imply the printed ones. All of them hold if $R\ge R^{**}:=\max\{R^*,4\ell_{\Nbar}+32M_b+25\}$ with $R^*:=\max\{4/\delta_0,107,3R_0+15+c^\sharp,2\ell_{\Nbar}+10M_b+10,C_{\mathrm{loc}}(\Nbar)\}$, $c^\sharp:=\max\{\delta_h+\rho,c_1',c_2\}$, together with the finitely many explicit bounds of (R.4)--(R.8).}\label{p:tab:A2}
\end{table}

\begin{table}[ht]
\small
\begin{tabular}{@{}p{0.08\textwidth}p{0.70\textwidth}p{0.14\textwidth}@{}}
\hline
name & condition & depends on \\
\hline
(b.0) & $p=6$, $p_1=1\in(2/3,p-1]$ & --- \\
(b.1) & $b_1=1$, $b_0\ge b_0^{\min}:=\max\{4d,c_1^I\}c_w(1+\log2)/(p\log2)$, $c_w\le(13/12)^{1/2}$ & 3--8 \\
(b.2) & $b_0\ge R/C$ & 5 \\
(b.3) & $b_0^2\ge\max\{2,\bar\mu\lambda^{-1/2}\}$ & 0 \\
(b.4) & the constants of stage 8 that depend on the Higgs mass ($c_7$, $c_1^I$, $\gamma_1''(M)$, $C_b'(M)$, $C_b^R(M)$, $C_E^R(M)$, $C_{F1}(M)$, $\mathfrak E_k(M)$) are taken at their maximum over $M\in\mathbb M_4$; the constant $C_>:=\sup_{t\in(0,1]}(\dots)$ is evaluated at the floor $\ell_0^\natural$ & 0--8 \\
\hline
\end{tabular}
\caption{The parameters $p$, $b_0$ and the constants of the stages 6--9.}\label{p:tab:A3}
\end{table}

\begin{table}[ht]
\footnotesize
\begin{tabular}{@{}p{0.09\textwidth}p{0.69\textwidth}p{0.14\textwidth}@{}}
\hline
name & condition at $\eps_0$ & depends on \\
\hline
(E.x) & $\hat x\le\min\{c_v,e_*,1\}$, $\lambda\eps_0\le c_v$ & 0, 8 \\
(E.1) & $\hat x^{1-\beta'}b_0\hat\ell^6\le c/(C_{\mathrm{reg}}c_1)$, $\beta'=1/56$ & 0, 3, 9 \\
(E.2) & $\hat xb_0R\hat\ell^8\le c$; $r(\eps_0)\le Cp(\eps_0)$ & 0, 5, 9 \\
(E.3) & $\abs e\lambda^{-1/4}\eps_0^{1/4}b_0\hat\ell^6\le(c_1-1)/C_{2.21}$ & 0, 3, 9 \\
(E.4) & $(24/11)(1-\eta)^{-1}C_H'\gamma_1^{-1}\mu_*^{-2}\eps_0^{-2}\mathrm{e}^{-\delta'R\hat\ell^2}\le\frac12$ & 0, 3, 5 \\
(E.5) & $2C_H(\gamma_1''M^2)^{-1}\eps_0^{-2}\mathrm{e}^{-\delta_HR\hat\ell^2}\le1$ and $\eps_0\le\exp(1-1/(\delta_HR))$, with $\gamma_1''=a_\flat/(a+M^2\eps_0^2)$, at each $M\in\mathbb M_4$ (Lemma~\ref{app:G1:RW-E06}; the condition is strongest at $M_1$) & 0, 3, 5, 8 \\
(E.6) & $\lambda^{-1/2}\eps_0^{-1/2}b_0^2\hat\ell^{11}\ge C_{2.12}$ & 0, 9 \\
(E.7) & $p(t)^2\ge C_\delta(M)t^{1/2}\lambda^{-1/2}$ for $t\le\eps_0$, at each $M\in\mathbb M_4$ & 0, 8, 9 \\
(E.8) & $C_Z\hat xb_0\hat\ell^6\le\frac12$ & 0, 5, 9 \\
(E.9) & $\mathrm{Err}_M(t)\le p(t)$ for all $t\le\eps_0$, at each $M\in\mathbb M_4$ & 0, 3, 5, 6, 9 \\
(E.10) & $R\hat\ell^2/4\ge2M_b+2$; $R\hat\ell^2\ge2(c_M+M_b)$; $a\,c_0c_1C_d\eps_0^{-1}\mathrm{e}^{-\delta_0R\hat\ell^2/8}\le\mu_*/2$ & 0--5 \\
(E.11) & $c_bb_1\hat\ell^{p_1}\ge b_*$; $C_A^{R1}\eps_0^{13/56}(c_bb_0\hat\ell^6)^{\rho_1}\le1$ & 0, 7, 9 \\
(E.m2) & $m_2^2\eps_0^2\le1$ & 0 \\
(E.$\Phi$) & $\lambda^{-1/4}m_2^2\eps_0^{7/4}\le a_\flat/C_\varphi$, $C_\varphi:=c_1L^{3/4}(1+\log L)^p$ ($=39.622\,c_1$ at $L=2$, $p=6$) & 0, 3 \\
(E.md) & (E.m2) and (E.$\Phi$) & 0, 3 \\
(E.12) & $C_\alpha\eps_0^\alpha\le b_0[\hat\ell^6-(\hat\ell-\log2)^6]-4d\hat\ell$ & 3, 9 \\
(E.13) & $\eps_0\le\min\{\frac12,(15b_0\log2/c_2)^{1/\alpha_I}\}$ & 3, 8, 9 \\
(E.$\delta$) & $\bar\delta:=\kappa\mu_*^2\eps_0^2\le1/12$ with the slice constant $\kappa=\frac{15}{2}$ of \eqref{eq:T6-GammaB}, i.e.\ $\eps_0\le(90\mu_*^2)^{-1/2}$; it implies (E.dom) (Lemma~\ref{app:G1:RW-B17}) & 0 \\
(E.14) & $4\kappa c_0^{\mathrm{ker}}c_d(1+\mu_*^2\eps_0^2)(1-\eta)^{-3}\mathrm{e}^{-\delta_0R\hat\ell^2/4}/(1-\mathrm{e}^{-\delta_0R\hat\ell\ln2/2})\le1/12$ (Lemma~\ref{app:G1:RW-KSP}) & 0, 3, 5 \\
(E.r) & $\eps_0<\mathrm{e}^{1-\varrho}=\mathrm{e}^{-1}$ & 4 \\
(E.15) & $C\eps_0^{\alpha+1/4}\le\frac12C_1b_0\hat\ell^6\lambda^{-1/4}$ & 0, 8, 9 \\
(E.dom) & $\mu_*\eps_0\le0.2$ & 0 \\
(E.fin) & $\eps_0\le\min(1/m_2,1/\mu_*,L)$ & 0 \\
(E.charge) & $\eps_0\le\min\{1,e^{-4}\}$, $e$ the charge (stage 10a); used for the small-field bounds of the scalar step (Proposition~\ref{app:K:L13-9}) & 0 \\
(E.Breg) & $(1+3d)C_\partial\,p(t)\,e(t)^{1-\beta}\le c_{\mathrm{reg}}$ for all $t\le\eps_0$, with $\beta<1$ the exponent fixed at stage 4 (stage 10a); used for the regularity of the transported background in the scalar H\"older estimate (Lemma~\ref{app:G1:H3-BD}) & 0, 3, 4, 6, 9 \\
(E.W) & $4C_We(t)p(t)\le\gamma_0$ for all $t\le\eps_0$; in particular $4w_j\le\gamma_0$ at every scale $t_j\le\eps_0$, where $w_j:=C_We(t_j)p(t_j)$ (stage 10a); used for the expansion of Ba{\l}aban's operator $W^{(j)}$ (Proposition~\ref{app:K:L13-7}) & 0, 3, 6, 9 \\
(E.rf) & the conditions (Sm1)--(Sm3) of Subsection~\ref{app:G1:Arem-sub} for all $t\le\eps_0$ (stage 10a); used for the remainder field of the scalar restrictions (Lemma~\ref{app:K:L09-1}) & 0, 3, 5 \\
(E.ent) & $(1+\ln(2/\eps_0))^{2p-2-\varrho d}\ge\Theta$ (Proposition~\ref{app:G1:I6}, Lemma~\ref{app:G1:RW-BM1}) & 0--10a \\
(E.comb) & the cap of the region-sum lemma (Lemma~\ref{app:G1:I145-comb}), with $c_A$ replaced by $c_A+C_{\rm old}+2C_{\rm loc}$ (Proposition~\ref{app:G1:I1}(iii)) & 0--9 \\
(E.edge) & \eqref{eq:lem24-cap} with $a_k$ replaced by $a$ (Lemma~\ref{app:G1:R0rows}); there $\kappa$ is the exponent of \cite[Prop.~2.2, (2.67)]{Ba2}, not the slice constant of (E.$\delta$) & 0--9 \\
(E.ksp1)--(E.ksp3) & the conditions of Lemma~\ref{app:G1:RW-KSP} & 0--9 \\
(E.vac) & $\eps_0$ is small enough for Theorems~\ref{thm:UV-K0minus} and \ref{thm:UV-K0plus} at $M_1,m,m_1,m_2$ & 0--9 \\
\hline
\end{tabular}
\caption{Upper bounds on $\eps_0$ (stages 10a and 10b). Every one of them is downward closed. The conditions (E.5), (E.7), (E.9), (E.ent) and (E.md) are imposed at each $M\in\mathbb M_4$ (minimum). The scalar zero mode (Lemma~\ref{app:G1:RW-ZMs}) imposes no condition beyond (E.m2), and the logarithmic entries (Lemma~\ref{app:G1:RW-logp}) impose none beyond (E.fin) (Lemma~\ref{app:G1:RW-E06}(c)).}\label{p:tab:A4}
\end{table}

\begin{table}[ht]
\small
\begin{tabular}{@{}p{0.09\textwidth}p{0.69\textwidth}p{0.14\textwidth}@{}}
\hline
name & condition & depends on \\
\hline
(T.floor) & $\ell_0\ge\ell_0^\natural=128/M_1$ (so $\ell_0\ge183/m$ suffices), $\ell_0\ge96/m$, $\ell_0\ge128/m$; it implies $\ell_T\ge L^3=8$ and $M_1\ell_0\ge32$ & 0 \\
(T.nw) & $\ell_0\ge\eps_0[8r(\eps_0)+4M_b+2]$: the boxes of radius $O(r(t_k))$ of the proofs, among them the box of Lemma~\ref{lem:lem24}, do not wrap around the torus & 3, 5, 10b \\
(T.cap) & $\ell_0\ge3r(\eps_0)\eps_0+C_{\Nbar}\eps_0$ (Lemma~\ref{app:G1:RW-B20}) & 4, 5, 10b \\
(T.T) & $\ell_0\ge\ell_0^T$, the torus side of Theorem~\ref{app:G1:VSR-T335} & 0--10b \\
(T.div) & $\ell_0\in2M_b\N$, $\eps_0=2^{-n_0}$ & 3, 10b \\
(T.per) & $\ell_0\ge\eps_0\max\{6\sqrt3\,r(\eps_0)+2,\ 2/\delta_0\}$ (Lemma~\ref{lem:IMG}, Lemma~\ref{lem:B:rows}(b)) & 3, 5, 10b \\
\hline
\end{tabular}
\caption{Lower bounds and divisibility for $\ell_0$ (stage 11). Nothing bounds $\ell_T$ from above or ties it to $\eps$.}\label{p:tab:A5}
\end{table}

\noindent\emph{The conditions at the order cap.} Let $\nbar=13$ and $\Nbar=\nbar^2=169$, and let $C_{\mathrm{loc}}(\Nbar)$ and $C_{\Nbar}$ be the numbers fixed at stage 4: the side, in units of blocks of the current scale, of a cube that contains a connected component of a family of order at most $\Nbar$, and the corresponding additive margin \cite[pp.~432--433]{Ba3}. The two conditions of Lemma~\ref{app:G1:RW-B20} are
\[
\text{(R.cap)}\quad R\hat\ell^2\ge C_{\mathrm{loc}}(\Nbar),\qquad\qquad \text{(T.cap)}\quad \ell_0\ge3r(\eps_0)\eps_0+C_{\Nbar}\eps_0 .
\]
(R.cap) is a lower bound on $R$ at stage 5. It is imposed in the form $R\ge C_{\mathrm{loc}}(\Nbar)$, which is part of (R.5) and of $R\ge R^*$ and implies (R.cap) for every $\eps_0\le1$, because $\hat\ell\ge1$. (T.cap) is a lower bound on $\ell_0$ at stage 11. Together with (E.r) they give, for every $k\le K$ and every torus of side $\ell_T\ge\ell_0$,
\[
r(t_k)\ge C_{\mathrm{loc}}(\Nbar),\qquad \big(3r(t_k)+C_{\Nbar}\big)t_k\le\ell_T .
\]
The first holds because $1+\log t_k^{-1}\ge\hat\ell$ for $t_k\le\eps_0$. The second holds because $t\mapsto t(1+\log t^{-1})^2$ is increasing on $(0,\mathrm{e}^{-1}]$, so that $r(t_k)t_k\le r(\eps_0)\eps_0$ when $t_k\le\eps_0<\mathrm{e}^{-1}$, and because $C_{\Nbar}t_k\le C_{\Nbar}\eps_0$.

The inequalities of the conditions on the scalar zero mode, on the Hessian of \cite[(3.12)]{Ba2} and of the logarithmic mass entries are given in Lemmas~\ref{app:G1:RW-ZMs}, \ref{app:G1:RW-MASS}, \ref{app:G1:RW-E06} and \ref{app:G1:RW-logp}: the first two impose no condition beyond (E.m2), (E.5) and (E.fin), and the last is covered by (E.fin). Some routes through the estimates of \cite{Ba2} that the proof of Theorem~\ref{thm:T6} does not use carry further conditions ($\eps_0\le\mathrm{e}^{-4}$, $r(\eps_0)>3M_b$ and the lower bounds $6r(\eps_0)\eps_0$, $21\eps_0$, $M_b\eps_0a_K(a_K-1)$ on $\ell_0$ of the cube-partition routes; $\mu_*\eps_0\le0.1525$); they are not needed and are not imposed.

\subsection{\texorpdfstring{The two conditions from Section~\ref{sec:T6}; acyclicity}{The two conditions from the vacuum-bound section; acyclicity}}\label{p:appA:acyc}

\begin{lemma}[the conditions (E.edge) and (T.per) in the order of choices]\label{lem:B:rows}
Let $c_d=\sqrt3M_b$.
\begin{enumerate}
\item[(a)] The condition (E.edge) of \eqref{eq:lem24-cap}, with $a_k$ replaced by its bound $a$, depends only on data of the stages 0--9 ($L$, $a$; $\lambda$; Ba{\l}aban's exponent $\kappa$ of \eqref{eq:lem24-cap}, fixed with $\nbar$ at stage 4 and distinct from the domination constant $\kappa$ of stage 0; $M_b,\delta_0,c_0,c_1$; $R$ through $r(t)$; $p$; $b_0$ through $p(t)$), and does not depend on $M$, $u$, $\eps$, $K$ or $\ell_T$. It is downward closed in $\eps_0$ and holds for all sufficiently small $\eps_0$, at every $R>0$. The conditions (R.edge) of \eqref{eq:lem24-R0row} follow from the lower bound $R\ge R^*$ of stage 5, and the condition $r(t_{k-1})/r(t_k)\le2L$ of Lemma~\ref{lem:lem24} holds for every $\eps_0\le1$.
\item[(b)] On the cubic torus $T_j$ the half-side of Ba{\l}aban's torus $\{-L_\mu\le x_\mu<L_\mu\}$ \cite[(1.2), pp.~603--604]{Ba1} is $L_{\min}=\ell_j/2$. So the condition \eqref{eq:GIMG} reads $\ell_j\ge\eps_0\max\{6\sqrt3\,r(\eps_0)+2,\,2/\delta_0\}$, and it is imposed at stage 11 as
\[
\text{(T.per)}\qquad \ell_0\ \ge\ \eps_0\max\{6\sqrt3\,r(\eps_0)+2,\ 2/\delta_0\}.
\]
This condition depends only on $\delta_0$ (stage 3), $R$ (stage 5) and $\eps_0$ (stage 10b), and it is a lower bound on $\ell_0$. If it holds and $\eps_0<\mathrm{e}^{1-\varrho}$ (condition (E.r)), then for every $j\ge1$, every $K$ and every scale $k\le K$, the rescaled side $S_k=\ell_j/t_k$ of $T_j$ satisfies $S_k\ge\max\{6\sqrt3\,r(t_k)+2,\ 2/\delta_0\}$; this is what Lemma~\ref{lem:IMG}(iv) uses.
\item[(c)] With these two conditions, the order of choices of Table~\ref{p:tab:stages} is acyclic and every stage can be filled.
\end{enumerate}
\end{lemma}
\begin{proof}
(a) The data are read off \eqref{eq:lem24-cap}; the coefficient $a_k$ enters only through $a_k\le a$, so the condition does not depend on $k$, and the masses do not enter. The condition is imposed for all $t\le\eps_0/L$, so if it holds for $\eps_0$ it holds for every smaller $\eps_0$. It holds for small $\eps_0$ because, with $s=\log t^{-1}$, the logarithm of the left side minus that of the right side is $(\tfrac14+\kappa)s+O(\log s)-\delta_0R(1+s)^{\varrho}/L+O(1)\to-\infty$ as $s\to\infty$, since $\varrho=2>1$; this uses only $R>0$. For (R.edge) note $r(\eps_0)\ge R$ for $\eps_0\le1$. The first condition of \eqref{eq:lem24-R0row} follows from $r(\eps_0)\ge0.4R_0+\big(1+0.4(2+\sqrt3)\big)M_b$, which is $\le0.4R_0+2.5M_b$. Now $R^*\ge\max\{3R_0+15,\ 2\ell_{\Nbar}+10M_b+10\}$. If $R_0\le18.7M_b+25$, then $10M_b+10\ge0.4R_0+2.5M_b$; otherwise $R_0>0.97M_b$, and $3R_0+15\ge0.4R_0+2.5M_b$. The second condition, $4r(\eps_0)\ge R_0+\sqrt3M_b$, follows from $4R\ge2(3R_0+15)+2(10M_b+10)$, and the third, $R\ge10M_b+10$, is part of $R\ge R^*$. Finally $r(t/L)/r(t)=\big(1+\log L/(1+\log t^{-1})\big)^{\varrho}\le(1+\log L)^2$ for $t\le1$, and $(1+\log2)^2<2.87\le4$.

(b) The data are read off \eqref{eq:GIMG}. Since $\ell_j=2L_{\min}\ge\ell_0$, the bound at $\ell_0$ gives \eqref{eq:GIMG} on every $T_j$. Put $f(t):=t(1+\log t^{-1})^{\varrho}$. Then $f'(t)=(1+\log t^{-1})^{\varrho-1}(\log t^{-1}+1-\varrho)>0$ for $0<t<\mathrm{e}^{1-\varrho}$. Hence $t\mapsto t\,(6\sqrt3\,r(t)+2)=6\sqrt3Rf(t)+2t$ is increasing on $(0,\eps_0]$ when $\eps_0<\mathrm{e}^{1-\varrho}$, and for $k\le K$, since $t_k\le t_K=\eps_0$,
\[
t_k\big(6\sqrt3\,r(t_k)+2\big)\le\eps_0\big(6\sqrt3\,r(\eps_0)+2\big)\le\ell_0\le\ell_j,
\]
that is, $S_k\ge6\sqrt3\,r(t_k)+2$. Also $S_k\ge\ell_j/\eps_0\ge2/\delta_0$. The condition is not implied by the other lower bounds on $\ell_0$: (T.nw) gives $\ell_0\ge\eps_0[8r(\eps_0)+4M_b+2]$ and $r(\eps_0)\ge R\ge R^*\ge10M_b+10>3M_b$, so $8r(\eps_0)+4M_b+2<9.34\,r(\eps_0)+2<6\sqrt3\,r(\eps_0)+2$ and (T.nw) never gives the first term of (T.per); and (T.floor) in the form $\ell_0\ge183/m$ together with (E.dom) gives $\ell_0/\eps_0\ge915$, which covers $2/\delta_0$ only if $\delta_0\ge2/915$, while $\delta_0$ is a constant of \cite{Ba1} without a printed value.

(c) This is Lemma~\ref{app:G1:acyc}, which treats all conditions of Tables~\ref{p:tab:A1}--\ref{p:tab:A5} with (E.edge) and (T.per) among them. Since by (b) the condition (T.per) exceeds (T.nw), $\ell_0$ is in general larger than (T.nw) alone requires; this violates no other condition, since no condition depends on $\ell_0$ and the conditions of Table~\ref{p:tab:A1} hold on all tori of side at least $96/m$.
\end{proof}

\begin{lemma}[the conditions (R.edge) and (E.edge)]\label{app:G1:R0rows}
Let $d=3$, $c_d=\sqrt3M_b$, and assume $M_b\ge2$, $R_0\ge0$, $\ell_{\Nbar}\ge0$, $\delta_0>0$ and $0<\eps_0\le1$.
\begin{enumerate}
\item[(i)] Suppose $R\ge3R_0+15$ and $R\ge2\ell_{\Nbar}+10M_b+10$ (part of (R.2) and (R.3)); in particular this holds if $R\ge R^{**}$. Then the three conditions \eqref{eq:lem24-R0row} hold. Since $c_d\ge\sqrt3$, so do the weaker forms $5(r(\eps_0)-M_b)/L\ge R_0+\sqrt3+2M_b$ and $4r(\eps_0)\ge R_0+\sqrt3$.
\item[(ii)] For $0<t\le1$, $r(t/L)/r(t)\le(1+\log L)^2\le2L$. Hence, for every $k$, with $\rho_k:=(r(t_{k-1})-M_b)/L$, we have $\rho_k\le2r(t_k)$, and $3\rho_k\ge2M_b$ under the third condition of \eqref{eq:lem24-R0row}.
\item[(iii)] Let $(\mathrm{E.edge})_a$ denote the condition \eqref{eq:lem24-cap} with $a_k$ replaced by $a$. It implies \eqref{eq:lem24-cap} at every $k$. It depends only on $L,a$ (stages 1--2), $\lambda$ (stage 0), Ba{\l}aban's exponent $\kappa$ (fixed with $\nbar$ at stage 4), $M_b,\delta_0,c_0,c_1$ (stage 3), $R$ (stage 5), $p$ (stage 6), $b_0$ (stage 9) and $C=C(c_0,\delta_0,d)$. It does not involve $M$, $u$, $\eps$, $K$ or $\ell_T$. It is downward closed in $\eps_0$, and for every $R>0$ it holds for all $\eps_0$ below a positive threshold that depends only on these data.
\end{enumerate}
\end{lemma}
\begin{proof}
(i) Since $\eps_0\le1$ we have $\log\eps_0^{-1}\ge0$, so $r(\eps_0)\ge R$. The maximum of two numbers is at least their mean. The first condition reads $r(\eps_0)\ge\frac25(R_0+c_d+2M_b)+M_b=0.4R_0+\big(1+0.4(2+\sqrt3)\big)M_b$, and $1+0.4(2+\sqrt3)=2.4928\ldots<2.5$. By the two hypotheses,
\[
R\ \ge\ \tfrac12(3R_0+15)+\tfrac12(2\ell_{\Nbar}+10M_b+10)\ \ge\ 1.5R_0+5M_b+12.5\ \ge\ 0.4R_0+2.5M_b.
\]
For the second condition, $4r(\eps_0)\ge4R\ge2(3R_0+15)+2(2\ell_{\Nbar}+10M_b+10)\ge R_0+\sqrt3M_b$. The third is $R\ge2\ell_{\Nbar}+10M_b+10\ge10M_b+10$.

(ii) Put $s=\log t^{-1}\ge0$. Then
\[
\frac{r(t/L)}{r(t)}=\Big(\frac{1+\log L+s}{1+s}\Big)^{2}=\Big(1+\frac{\log L}{1+s}\Big)^2\le(1+\log L)^2,
\]
and $(1+\log2)^2=2.8667\ldots\le4$. With $t=t_k\le\eps_0\le1$ and $t_{k-1}=t_k/L$ this gives $\rho_k\le r(t_{k-1})/L\le2r(t_k)$. Since $r(t_{k-1})\ge R$ and $L=2$, $3\rho_k\ge3(R-M_b)/2\ge R-M_b\ge9M_b+10\ge2M_b$.

(iii) Since $0<a_k\le a$, the left side of \eqref{eq:lem24-cap} with $a$ is at least the one with $a_k$. So $(\mathrm{E.edge})_a$ implies \eqref{eq:lem24-cap} at every $k$. Its only dependence on $k$ and on the mass came through $a_k$, and that is now gone. The list of data is read off the display; $b_0$ enters through $p(t)$.

The condition is imposed for all $t\le\eps_0/L$. If $\eps_0'\le\eps_0$, then this range of $t$ shrinks, so the condition is downward closed. For the threshold, put $s=\log t^{-1}$. Then $\log p(t)=\log b_0+p\log(1+s)$ and $\log(1+r(t))=O(\log(1+s))$. So the logarithm of the left side minus that of $(Lt)^\kappa$ is
\[
\log\!\big(C\,a\,c_1\lambda^{-1/4}b_0\,\mathrm{e}^{\delta_0(M_b/L+c_d)}L^{-\kappa}\big)+\big(\tfrac14+\kappa\big)s+O(\log(1+s))-\frac{\delta_0R}{L}(1+s)^{2},
\]
which tends to $-\infty$ as $s\to\infty$, because $\delta_0R>0$ and $\varrho=2>1$. Hence there is an $s_0<\infty$, depending only on the listed data, such that the condition holds for every $t\le\mathrm{e}^{-s_0}$, that is, whenever $\eps_0\le L\mathrm{e}^{-s_0}$. The factor $(1+r(t))^{d-1}$ changes only $s_0$.
\end{proof}

\begin{lemma}[the order of choices is acyclic and every stage can be filled]\label{app:G1:acyc}\label{app:G1:dom-node}\label{app:G1:acyc-GIMG}
Consider the conditions of Definition~\ref{app:G1:rowlist}, read as follows:
(1) (b.3) in its $\eps_0$-free form, the other part of the hypothesis of \cite[(2.12)]{Ba2} being the upper bound (E.6) at stage 10b;
(2) $a_k$ in (E.edge) replaced by its bound $a$;
(3) (R.edge) in its sufficient stage-5 form (Lemma~\ref{app:G1:R0rows}(i)).
Assume also that the torus-dependent constants are evaluated at $\ell_0^\natural$ (Lemma~\ref{app:G1:floor}), and that the $\eps_0$-dependent constants entering stage 10b are replaced by their suprema over $(0,\eps_1]$ (Lemma~\ref{app:G1:L10asup}).
Form the directed graph whose vertices are the data of Table~\ref{p:tab:stages} and the conditions, with an edge from each datum to each condition that depends on it (as listed in the last columns of Tables~\ref{p:tab:A1}--\ref{p:tab:A5}) and an edge from each condition to the datum it constrains. Then:
\begin{enumerate}
\item[(i)] every edge goes from a stage to the same or a later stage;
\item[(ii)] the graph is acyclic;
\item[(iii)] every stage can be filled;
\item[(iv)] no condition is an upper bound on $R$, $b_0$ or $\ell_0$ or a lower bound on $\eps_0$, and none depends on $\eps$, $k$, $\ell_T$ or $K$.
\end{enumerate}
\end{lemma}
\begin{proof}
(i) This is read off the last columns of Tables~\ref{p:tab:A1}--\ref{p:tab:A5}, of Lemma~\ref{app:G1:R0rows}(iii) for (E.edge) and of Lemma~\ref{lem:B:rows}(b) for (T.per). The following potential edges against the order are excluded by the form of the conditions.
(a) (E.edge) and $R$: $R$ is fixed at stage 5 from data of the stages 2--4, and no lower bound on $R$ depends on a datum of stage 10.
(b) (E.edge) and $b_0$: the lower bounds on $b_0$ ((b.1), (b.2), the maximum in (b.4), and (b.3) read as in (1)) depend neither on $\eps_0$ nor on (E.edge).
(c) $\eps_1$ (stage 10a) depends on no datum of stage 10b.
(d) The non-wrapping of the box of Lemma~\ref{lem:lem24} is the condition (T.nw) at stage 11, which depends on $\eps_0$; the edge points forward.
(e) (E.edge) does not involve the mass (Lemma~\ref{app:G1:R0rows}(iii)), so the four-mass rule adds no instance of it.
(f) The floor $\ell_0^\natural$ is not a datum of (E.edge).
(g) (R.edge) follows from $R\ge R^*$, so no edge into $R$ from it is needed.
(h) (E.13) depends on $b_0$ (stage 9) and is imposed at stage 10b; (b.1) does not depend on $\eps_0$.
The condition (T.per) depends on $\delta_0$ (stage 3), $R$ (stage 5, through $r(\eps_0)$) and $\eps_0$ (stage 10b); its one edge out goes to $\ell_0$ (stage 11), and no condition depends on $\ell_0$.

We now go through the tables condition by condition.

\emph{Table~\ref{p:tab:A1}.} The conditions (P.1)--(P.6) depend on the point only and constrain the point. They are infima or suprema over all tori of side at least $96/m$ and all $\eps\le1/m$, so they depend neither on $\eps_0$ nor on $\ell_0$.

\emph{Table~\ref{p:tab:A2}.} Every condition constrains $R$ (stage 5), and its coefficients are data of the stages 3 and 4. The quantities $r(\eps_0)$, $r(t)$, $\hat\ell$ in (R.4), (R.7), (R.8), (R.edge) and (R.cap) would introduce an edge from $\eps_0$ (stage 10b). As stated in the caption of Table~\ref{p:tab:A2}, they are replaced by $R$, $R$ and $1$. Since $r(t)\ge R$ and $\hat\ell\ge1$ for $t,\eps_0\le1$, the forms so obtained imply the printed ones, and they do not depend on $\eps_0$. (R.nbar) depends on $\nbar$ (stage 4).

\emph{Table~\ref{p:tab:A3}.} (b.1)--(b.3) constrain $b_0$ (stage 9). Their coefficients lie in the stages 0--8 ((b.2) depends on $R$, stage 5), and (b.3) is free of $\eps_0$ by (1). (b.4) fixes stage-8 constants as maxima over $\mathbb M_4$, whose elements are data of stage 0, and at the floor $\ell_0^\natural$, also a datum of stage 0. The constants of the stages 7 and 8 depend on $\eps_0$ only through numbers that are guaranteed by conditions of stage 10a. These are: the bound $\gamma_1''\ge a_\flat/(a+1)$ for $M\le m_2$, guaranteed by (E.m2); the number $\frac1{12}$, guaranteed by (E.$\delta$) and (E.14); and one term proportional to $\eps_0$ in a constant of \cite{Ba2}, which is bounded at stage 10a by its value at $\eps_1$ (Lemma~\ref{app:G1:L10asup}). The conditions (E.m2), (E.$\delta$) and (E.14) depend only on data of the stages 0, 3 and 5. So each of these dependences goes one way, and none closes a cycle. The slice constant $\kappa=\frac{15}{2}$ is a number fixed at stage 0. The inequality $\kappa\ge\kappa_1\kappa_3$ that it must satisfy holds under (E.dom) (Remark~\ref{app:G1:DOM-kappa}), and (E.$\delta$) implies (E.dom) (Lemma~\ref{app:G1:RW-B17}). This again is a one-way dependence. The constants of the per-step remainders that depend on $b_0$ and $R$ appear in no condition and enter only the output constants.

\emph{Table~\ref{p:tab:A4}.} Every condition constrains $\eps_1$ (stage 10a) or $\eps_0$ (stage 10b). Its coefficients lie in the stages 0--9, or, through suprema over $(0,\eps_1]$, in stage 10a. Under (2), the coefficients of (E.edge) are $L,a,\lambda$, the exponent $\kappa$, $M_b,\delta_0,c_0,c_1,R_0$, $R$ through $r(t)$, $p,p_1$, $b_0$ through $p(t)$, and a constant depending on $(c_0,\delta_0,d)$. All of these are data of the stages $\le9$. (E.ent) depends on $\Theta$, which is formed from $b_0$, $R$ and the suprema of $c_5,c_7$ over $(0,\eps_1]$ (Lemma~\ref{app:G1:RW-BM1}), so on the stages 0--10a. (E.9) is imposed at each mass and depends on $p$ (stage 6), on constants of stage 3, on $R$ and on $b_0$; (E.13) depends on $c_2,\alpha_I$, which are constants of \cite{Ba1} at stage 8. The last columns of (E.9), (E.13) and (E.ent) list these stages. (E.5) depends on $C_H,\delta_H$ (stage 3), on $R$ and on the mass; it also contains $\eps_0$ through $\gamma_1''$, but $\eps_0$ is the quantity it constrains, not one of its coefficients. (E.charge), (E.Breg) and (E.W) depend on data of the stages 0, 3, 4, 6 and 9 (as listed in their last columns) and are imposed at stage 10a.

\emph{Table~\ref{p:tab:A5}.} Every condition constrains $\ell_0$ (stage 11), and its coefficients lie in the stages $\le$10b. No datum of a stage $\le$10b depends on $\ell_0$. The torus-dependent constants are evaluated at $\ell_0^\natural$ (Lemma~\ref{app:G1:floor}), and the conditions of Table~\ref{p:tab:A1} are uniform over all tori of side at least $96/m$.

\emph{Masses and output.} The four-mass rule adds only edges from the masses (stage 0) into the stages 8--11 (Lemma~\ref{app:G1:mass}(iii)). The tori $T_j$, $K$ and $\eps$ are chosen after stage 11, and by (iv) no condition depends on them.

(ii) Order the stages $0<1<\dots<11$. By (i) the stage index does not decrease along any edge. So along a directed cycle it is constant: every cycle lies inside a single stage. Inside the stages 5, 9, 10b and 11 the edges inside the stage go from the conditions to the quantity they constrain ($R$, $b_0$, $\eps_0$, $\ell_0$), and from that quantity to the constants defined from it (such as $r(t)$, $p(t)$ and $C_>$); no condition of the same stage depends on that quantity, and under (3) this includes (R.edge). A condition of stage 10b is an inequality for $\eps_0$, also when $\eps_0$ enters through $\gamma_1''$ as in (E.5); there $\eps_0$ is the quantity constrained, not a coefficient, so it gives no edge from $\eps_0$ into the condition. Inside the stages 0--4, 6--8 and 10a the dependences inside the stage are one-way: from the point to $\kappa$ and $\ell_0^\natural$; at the stages 3 and 4 along the order in which the constants are fixed (at stage 3, $M_b$ from $(d,a,L)$ and the other constants from $(d,a,L,M_b)$; at stage 4, $\beta$ after $\beta_U$); from $p$ and $\mathbb W_0$ to their constants; from $\eps_1$ to the suprema over $(0,\eps_1]$. So no stage contains a cycle, and the graph is acyclic. In particular (T.per), whose data lie in earlier stages and whose one edge out ends at $\ell_0$, on which no condition depends, lies on no cycle. It is imposed only as a lower bound on $\ell_0$ at stage 11, never as an upper bound on $\eps_0$ at a fixed torus, which would close the cycle $\eps_0\to\text{(T.nw)}\to\ell_0\to\text{(T.per)}\to\eps_0$.

(iii) At the stages 5, 9 and 11 there are finitely many lower bounds, each a finite function of data of earlier stages, so their maximum is finite. At stage 11 there is also the divisibility $\ell_0\in2M_b\N$, which every large multiple of $2M_b$ meets. At stage 10a the conditions are finitely many and each has a positive threshold, so $\eps_1>0$. Among them, (E.charge) is explicit. (E.Breg) and (E.W) have the form ``$F(t)\le c$ for all $t\le\eps_0$'' with $F(t)\to0$ as $t\to0$, because $p(t)$ grows only logarithmically and $e(t)^{1-\beta}=\abs e^{1-\beta}t^{(1-\beta)/2}$ with $\beta<1$; so they are of this type by Lemma~\ref{app:G1:mon}(c). At stage 10b there are finitely many upper bounds on $\eps_0$, the instances at the four masses of $\mathbb M_4$ included. Each of them is downward closed and holds for all dyadic $\eps_0$ below a positive threshold: for the conditions of Table~\ref{p:tab:A4} other than (E.edge) this is Lemma~\ref{app:G1:ord}, and for (E.edge) it is Lemma~\ref{app:G1:R0rows}(iii). The minimum of finitely many positive thresholds is positive, and some $2^{-n_0}\le\eps_1$ lies below it.

(iv) This is read off Tables~\ref{p:tab:A2}--\ref{p:tab:A5}: the conditions of the stages 5, 9 and 11 are lower bounds, and those of the stages 10a and 10b are upper bounds on $\eps_0$. None of them depends on $\eps$, $k$, $\ell_T$ or $K$. These are chosen after stage 11 and occur in no last column of Tables~\ref{p:tab:A1}--\ref{p:tab:A5}. The conditions of Table~\ref{p:tab:A1} are uniform over all tori of side at least $96/m$ and all $\eps\le1/m$, and the torus-dependent constants are evaluated at $\ell_0^\natural$.
\end{proof}

\subsection{The conditions at every point of the window}\label{p:appA:windowlemma}

\begin{lemma}[the conditions at every point of the window]\label{lem:window}
Let $\pi=(e,\lambda,m,\eta,m_1,m_2)\in\fW$, $\mu_*:=m$, $M_1:=0.7m$, $H:=m^2$, $\mathbb M_4:=\{M_1,m,m_1,m_2\}$; let $M_1$ be among the data of stage 0, with the torus floor $\ell_0^\natural:=\max\{8,128/m_1,128/M_1\}$ fixed there (lengths in the units of Section~\ref{sec:model}; $\ell_0^\natural=128/M_1=182.857\ldots/m$ whenever $M_1\le16$), and let every constant or condition of the stages 8--11 that depends on the Higgs mass be evaluated at each $M\in\mathbb M_4$ and combined by the maximum (constants, lower bounds) and by the minimum (upper bounds on $\eps_0$). Then there are $R$, $b_0$, $n_0$ and $\ell_0\in2M_b\N$ such that, with $\eps_0:=2^{-n_0}$, every condition of Tables~\ref{p:tab:A1}--\ref{p:tab:A5} holds at $\pi$. For these data every spacing $\eps=2^{-K}\eps_0$ satisfies $\eps\le\eps_0\le1/m_2<1/m$, every torus $T_j$ has side $\ell_j\ge\ell_0\ge96/m$, and statements (a)--(d) of Theorem~\ref{thm:ch-W} hold at $\pi$.
\end{lemma}
\begin{proof}
\emph{Step 1: the four types of conditions.} Evaluate the constants that depend on $\ell_T$ at $\ell_0^\natural$ (Lemma~\ref{app:G1:floor}), and replace the $\eps_0$-dependent constants entering stage 10b by their suprema over $(0,\eps_1]$ (Lemma~\ref{app:G1:L10asup}). By Lemma~\ref{app:G1:acyc}, read in this way, the order of choices is acyclic: every condition depends only on data of its own or of earlier stages. The condition (T.per) is a lower bound on $\ell_0$ at stage 11 that depends only on $\delta_0$ (stage 3), $R$ (stage 5) and $\eps_0$ (stage 10b), and no condition depends on $\ell_0$. Every condition is of one of the following types.
\begin{itemize}
\item[(i)] A condition on the point $\pi$ alone: the conditions of Table~\ref{p:tab:A1}, which are statements (a)--(d) of Theorem~\ref{thm:ch-W}.
\item[(ii)] A lower bound $x\ge f$ on $R$ (stage 5), on $b_0$ (stage 9) or on $\ell_0$ (stage 11), where $f$ is a finite function of data of earlier stages.
\item[(iii)] An upper bound on $\eps_0$ (stages 10a--10b) that is downward closed (if it holds at a dyadic $\eps_0$, it holds at every dyadic $\eps_0'\le\eps_0$) and that holds for all dyadic $\eps_0$ below a positive threshold depending on data of earlier stages.
\item[(iv)] The divisibility $\ell_0\in2M_b\N$.
\end{itemize}
That the conditions of Tables~\ref{p:tab:A2}--\ref{p:tab:A5} are of the types (ii)--(iv) at every point of $\fW$, with thresholds finite, respectively positive, at every point, is Lemma~\ref{app:G1:ord}, together with Lemma~\ref{app:G1:R0rows}(iii) for (E.edge) and Lemma~\ref{lem:B:rows}(b) for (T.per). The family $T_j=j!\,T_0$, $\eps=2^{-K}\eps_0$, $K\ge n_0+1$, is a definition, and $\eps\le1/m$ follows from (E.m2). Moreover no condition bounds $b_0$ from above before $\eps_0$ is chosen, no condition bounds $\ell_0$ from above, the only conditions that bound $R$ from above at fixed $\eps_0$ are imposed after $R$ (in the order of choices they are upper bounds of type (iii) on $\eps_0$ or lower bounds of type (ii) on $b_0$, both chosen after $R$), and no condition involves the spacing $\eps$, the index $K$ or the torus except through the floor $\ell_0^\natural$ of stage 0 and the lower bounds of stage 11 (Lemma~\ref{app:G1:ord}(iii)).

\emph{Step 2: the conditions that involve masses, $\eta$ and the couplings are of these types at every point of $\fW$.}
(i) $R^{**}=\max\{R^*,4\ell_{\Nbar}+32M_b+25\}$, and $R^*$, are functions of $(d,a,L,M_b,R_0,\delta_0,\allowbreak\nbar,\Nbar,\allowbreak V(\Nbar),C_{\mathrm{loc}}(\Nbar))$, of fixed $O(1)$ localisation widths and of $b_d$, all fixed at the stages 0--4 independently of $\pi$; no Higgs mass, no coupling and no photon mass enters. The remaining lower bounds on $R$ ($R_0+1$, $4/\delta_0$ and the bounds from (E.10), (E.14) and (R.cap)) are free of the Higgs mass, and the only parameters they contain, $\kappa$, $\mu_*$ and $(1-\eta)^{-3}$, are finite on $\fW$. So $R\ge R^{**}$ is of type (ii) (Lemma~\ref{app:G1:mass}(i)).
(ii) The four-mass rule. Every constant of the stages 8--9 that depends on the Higgs mass depends on $M$ either from above, through $M^2t^2$ with $t\le\eps_0$, which (E.m2) at the largest mass, $m_2^2\eps_0^2\le1$, bounds by $1$ for every $M\le m_2$ (the constants are computed under this a priori bound, whose threshold $1/m_2$ depends on no constant, so no cycle arises), or through powers of $M$ of either sign, logarithms and other quantities finite at every fixed $M>0$ (for instance $(M^2\ell_T^3)^{-1}$, $(M\ell_T)^{-10}$ and $\log_+(1/(\min(M,\mu_*)\eps_0))$ in the constant of the final c-numbers, the constant $C_\delta(M)$, and the conditioning constant $\rho_{\min}(M)^{-1}$ of the bounds on the vector slices). No entry orders $M$ against $\mu_*$, and none uses $m_1>2m$; these conditions occur only in Table~\ref{p:tab:A1} (Lemma~\ref{app:G1:mass}(ii)). A maximum of four finite numbers is finite, and the constants enter only lower bounds, so these conditions are of type (ii). The upper bounds on $\eps_0$ of stage 10b that depend on the Higgs mass ((E.5), (E.7), the conditions of Lemmas~\ref{app:G1:RW-ZMs}--\ref{app:G1:RW-E06}, the logarithmic entries of Lemma~\ref{app:G1:RW-logp}, (E.md) and (E.fin)) are downward closed with a positive threshold at each fixed $M>0$; the minimum of four positive thresholds is positive, so these conditions are of type (iii).
(iii) The condition (E.edge) reads $C\,a_kc_1\lambda^{-1/4}t^{-1/4}p(t)\big(1+r(t)\big)^{d-1}\exp\big(-\delta_0[(r(t)-M_b)/L-\sqrt3M_b]\big)\le(Lt)^\kappa$ for all $t\le\eps_0/L$, with $p$ polynomial in $\ln t^{-1}$ and $r(t)=R\,\ell(t)^2$, $\ell(t)^2/\ln t^{-1}\to\infty$ as $t\downarrow0$. The exponential therefore decays faster than every power of $t$, and the inequality holds for all $t$ below a positive threshold that depends on $L,R,M_b,\delta_0,c_0,R_0,\kappa,p,\lambda,a,c_1$, all fixed before stage 10b; it has no explicit Higgs-mass entry (it depends on $b_0$ through $p(t)$, and $b_0$ is fixed at stage 9 as a maximum over $\mathbb M_4$, so no cycle arises). The conditions \eqref{eq:lem24-R0row}, with $c_d=\sqrt3M_b$, follow from $R\ge R^*$ (Lemma~\ref{lem:B:rows}(a)). So (E.edge) is of type (iii).
(iv) The condition (T.per), $2L_{\min}\ge\eps_0\max\{6\sqrt d\,r(\eps_0)+2,\,2/\delta_0\}$, where $2L_{\min}$ is the side of the torus, is a lower bound on the side by a function of $\eps_0,R,\delta_0$; since every $T_j$ has side $\ge\ell_0$, it is a condition of type (ii) at stage 11.
(v) (T.floor) reads $\ell_0\ge\max\{96/m,\ 128/m,\ \ell_0^\natural\}$, of type (ii) at stage 11. The constants that are evaluated at the floor ($C_\delta(M)$, $c_7$, $C_>$ and the zero-mode terms of the bounds on the vector slices) are non-increasing in the side of the torus, so their values at $\ell_0^\natural$ dominate their values on every torus of side $\ge\ell_0\ge\ell_0^\natural$; raising the floor from $\max\{8,128/m_1\}$ to $\ell_0^\natural$ can only weaken the conditions in which they occur (Lemma~\ref{app:G1:floor}).
(vi) The couplings and $\eta$ enter the conditions only through quantities finite for $e\ne0$, $\lambda>0$, $0<\eta<1$ (for instance $\lambda^{-1/4}$ in (E.$\Phi$) and in (E.edge), powers of $(1-\eta)^{-1}$ in (E.4) and (E.14)); the convexity range $0<\eta<\eta_{\max}(\eps_0)$ is an upper bound on $\eps_0$ at fixed $\eta$. No condition of the types (ii)--(iv) is a condition on $e$ or $\lambda$: the couplings are restricted only by the definition of $\fW$, whose upper bounds $\lambda<0.1m$, $e^2<0.1m$ enter through the conditions of type (i) (the proof of Theorem~\ref{thm:ch-W}(a)--(c)); in particular no condition imposes a lower bound on $\lambda$ beyond $\lambda>0$. As $\lambda\downarrow0$ at fixed $e$ the thresholds move (for instance through the factors $\lambda^{-1/4}$ above), but at every $\lambda>0$ they stay finite, respectively positive. The smallness conditions of Ba{\l}aban's renormalisation group concern the running couplings at the scales $t\le\eps_0$ and are conditions of type (iii).

\emph{Step 3: the choice at $\pi$.} Choose the data stage by stage. At stage 0 the parameters and $\ell_0^\natural$; at the stages 1--4 the discrete data ($a=1$, $L=2$, $M_b$, $\nbar=13$, $\Nbar=169$ and the others), which do not depend on $\pi$; at stage 5, $R$ an integer above the finitely many lower bounds of type (ii), $R^{**}$ among them; at stage 6, $p=6$; at the stages 7--8 the constants, at the four masses; at stage 9, $b_0$ the maximum of the finitely many lower bounds; at stage 10a, $\eps_1$ is the minimum of finitely many positive thresholds, and the constants that depend on $\eps_0$ are taken at their suprema over $(0,\eps_1]$, which are finite by Lemma~\ref{app:G1:L10asup}; at the stages 10a--10b, the finitely many upper bounds of type (iii), at the four masses and including (E.edge), have a positive common threshold $\bar\eps(\pi)$, and every $\eps_0=2^{-n_0}\le\bar\eps(\pi)$ satisfies all of them by downward closedness; at stage 11, $\ell_0$ the least element of $2M_b\N$ above the finitely many lower bounds. Each choice uses only data of earlier stages, and by Step~1 no later choice can violate an earlier condition. The conditions of type (i) hold by Theorem~\ref{thm:ch-W}, which proves (a)--(d) at every point of $\fW$, uniformly in $\eps\le1/m$ and in tori of side $\ge96/m$. The condition (E.m2), $m_2^2\eps_0^2\le1$, gives $\eps\le\eps_0\le1/m_2<1/m$. And $\ell_j\ge\ell_0\ge\max\{96/m,\ell_0^\natural\}\ge96/m$. So (a)--(d) apply to every lattice of the family.
\end{proof}

\emph{Mass ranges.} Four ingredients of the proof of Theorem~\ref{thm:T6} are used in the transfer of Lemma~\ref{app:G1:RMG} at the mass $M_1=0.7m$, which is not one of the masses of Theorem~\ref{thm:T6}. The next lemma shows that each of them holds at every $M\in(0,m_2]$, by its own proof. Here $t=t_k=L^k\eps\le\eps_0$ is the scale of the $k$-th step, $a_k$ is Ba{\l}aban's averaging constant of that step, with $a_\flat\le a_k\le a$, and $\rho_{\min}(M):=\min(M,\mu_*)\min\{1,(L\max(M,\mu_*))^{-1}\}$.

\begin{lemma}[mass range of four ingredients]\label{lem:B:massrange}
Let $(e,\lambda,m,\eta,m_1,m_2)\in\fW$ and $\mu_*=m$. Let the conditions (E.m2), (E.$\Phi$), (E.fin) and $\ell_0\ge L^3$ hold. Let $M\in(0,m_2]$, with the upper bounds on $\eps_0$ and the constants of the stages 8--10b that depend on the Higgs mass imposed at $M$. Then:
\begin{enumerate}
\item[(a)] (\emph{the final c-numbers}) the torus-uniform bounds on $\delta m^2_{\rm fin}$ and on $E_{\rm fin}/\abs T$ of Propositions~\ref{app:G1:I3b} and \ref{app:G1:Efin} hold at Higgs mass $M$. They hold with $\bar\ell(M):=\min(1/M,1/\mu_*,L)$ in place of $\bar\ell(m_i)$ and with constants that are polynomials in $1/\rho_{\min}(M)$. Their dependence on $\ell_T$ is through $Z_f=t_{\bar K}/(m_f^2\ell_T^3)$, $m_s=M$, $m_v=\mu_*$, which is non-increasing in $\ell_T$;
\item[(b)] (\emph{the $A\cdot A$ vertices}) the bounds on the $A\cdot A$ vertices of \cite[(3.30)]{Ba3} hold at Higgs mass $M$ with the same constants $\Gamma_3,\Gamma_3'$, uniformly in $M\in(0,m_2]$, and the bound on the vertex of \cite[(3.32)]{Ba3} holds at $M$ with a constant $\Gamma_4(M)<\infty$ that depends on $M$ only through $m_{j_0}^2\le1$ and the bound of (a) (Lemmas~\ref{app:G1:AA} and \ref{app:G1:AA4});
\item[(c)] (\emph{the tail bound}) under the hypotheses of Proposition~\ref{app:G1:T579}, its two-sided tail bound holds at Higgs mass $M$, with a constant $C'_{579}(M)<\infty$ that depends on $M$ only through the $\delta m^2$ vertices, that is, through the bound of (a) at the last scale and a bound free of the mass at the intermediate scales, and through $M^2t^2\in[0,1]$, on which it is uniform;
\item[(d)] for every real background $A$, every domain $\Omega$ with Neumann Laplacian $-\Delta^N_{A,\Omega}\ge0$ and every $Q_k=Q_k(A)$ with $Q_kQ_k^*=I$, the operator $\Delta^{(k)}(\Omega,A):=a_k-a_k^2Q_kG_k(\Omega,A)Q_k^*$, $G_k:=(-\Delta^N_{A,\Omega}+M^2t^2+a_kQ_k^*Q_k)^{-1}$, satisfies
\[
\frac{a_kM^2t^2}{a_k+M^2t^2}\ \le\ \Delta^{(k)}(\Omega,A)\ \le\ a_k,
\]
and $\gamma_1:=a_k/(a_k+M^2t^2)\ge a_\flat/(a+1)$ under (E.m2).
\end{enumerate}
At the four masses used by the proof, $M\in\mathbb M_4$, the upper bounds and constants of the hypothesis are those of the four-mass rule.
\end{lemma}
\begin{proof}
(d) Since $-\Delta^N_{A,\Omega}\ge0$ and the inverse is operator monotone, $G_k\le(M^2t^2+a_kP)^{-1}$ with $P:=Q_k^*Q_k$. Since $Q_kQ_k^*=I$, $P$ is an orthogonal projection, and $(M^2t^2+a_kP)^{-1}=P/(M^2t^2+a_k)+(1-P)/(M^2t^2)$. Also $Q_kP=Q_k$, so $Q_k(1-P)=0$ and $Q_kG_kQ_k^*\le Q_kPQ_k^*/(M^2t^2+a_k)=(M^2t^2+a_k)^{-1}$. Hence $\Delta^{(k)}\ge a_k-a_k^2/(a_k+M^2t^2)=a_kM^2t^2/(a_k+M^2t^2)$. The upper bound follows from $G_k\ge0$. Finally $M^2t^2\le m_2^2\eps_0^2\le1$ by (E.m2) and $t\le\eps_0$, so $\gamma_1\ge a_\flat/(a+1)$. Only $M>0$ was used.

(a) The construction of Propositions~\ref{app:G1:I3b} and \ref{app:G1:Efin} slices both fields, scalar and photon, up to one common index $\bar K$ with $t_{\bar K}\le\bar\ell(M)$. Every renormalised subgraph then has internal lines of index below $\bar K$ only, and the ratio of the two masses enters only through the top cluster, as the number $\rho_{\min}(M)$. The proof nowhere uses the order of $M$ and $\mu_*$: the light field, which is the scalar when $M<\mu_*$, carries its long remainder line in the top cluster in the same way as the photon does when $M>\mu_*$. Its hypotheses are $\eps_0\le\bar\ell(M)$, which follows from (E.fin) because $\bar\ell(M)\ge\bar\ell(m_2)$ for $M\le m_2$, and $\ell_T\ge L^3$, which does not involve the mass. Its constants are of three kinds. $C_b$ does not depend on $M$. $C_b^R$ is a polynomial in $1/\rho_{\min}(M)$ and in $\kappa_{R,\min}^{-1}$, $\kappa_{R,\min}=\min(\rho_{\min},\kappa_G)/8$. The remaining factors are $\max(LM,L\mu_*,1)$, $\log_L^+(\bar\ell/\eps_0)$ and $Z_f$. Each is finite at every $M>0$, so the statement holds at $M$ with these constants. The line bounds that carry this argument are those of Lemma~\ref{app:G1:VSR-F} and Theorem~\ref{app:G1:VSR-VS}. The renormalised subgraphs are bounded by the renormalisation machinery of Lemmas~\ref{app:G1:VSR-A}, \ref{app:G1:VSR-Btop}, \ref{app:G1:VSR-C}, Theorem~\ref{app:G1:VSR-T335} and Proposition~\ref{app:G1:VSR-5b}, and their constants are among the three kinds just listed.

(b) The constants $\Gamma_3,\Gamma_3'$ depend on the Higgs mass only through the rescaled mass $m_{j_0}^2=M^2t_{j_0}^2$. For $M\le m_2$, $t_{j_0}\le\eps_0$ and (E.m2), this lies in $[0,1]$, and the constants are uniform there. Lemma~\ref{app:G1:AA} is stated for every $m_{j_0}^2\in[0,1]$, so it applies at every $M\in(0,m_2]$. The constant $\Gamma_4$ of \cite[(3.32)]{Ba3} is that of Lemma~\ref{app:G1:AA4}. It also depends on $M$ through the finite mass vertices of the classes, and by (a) these are bounded at every $M\in(0,m_2]$. Uniformity as $M\downarrow0$ is not claimed, and the proof uses $\Gamma_4$ only at the four masses of $\mathbb M_4$. The bound on $\Gamma_4(M)$, with this dependence on $M$, is proved in steps (iv)--(v) of the proof of Lemma~\ref{app:G1:AA4}.

(c) The constant $C'_{579}$ depends on $M$ through the $\delta m^2$ vertices of the step functionals and through $M^2t^2\in[0,1]$. The $\delta m^2$ vertices are bounded at $M$ by (a) at the last scale, and by a bound free of the mass at the intermediate scales (Proposition~\ref{app:G1:T579}). None of its other arguments ($\nbar$, $\Nbar$, $\alpha_0$, $d$, $L$, $a$, $M_b$, $R$, $r$, $e$, $\lambda$, $\mu_*$, $b_0$, $p_0$, $c_1$, $\delta_0$, $\eta$) depends on $M$. The tail bound is Proposition~\ref{app:G1:T579}. Its proof rests on the remainder bounds of Proposition~\ref{app:G1:VSR-K3}, the weighted bounds of Lemma~\ref{app:G1:VSR-K4}, and the family bound of Definition~\ref{app:G1:VSR-IH} at the order cap (Proposition~\ref{app:G1:VSR-CAP}). In these statements the propagators carry $M$ only through $M^2t^2\in[0,1]$, and their bounds are uniform there.

In (a)--(c) the dependence on $u$ is unchanged. The transfer uses $u=0$, which lies in $[-\eta,1]$. Every factor $\abs u$ of the constants at $u=0$ is at most its value at the points $(u,M)\in\{(1,m_2),(-\eta,m_1)\}$ of Theorem~\ref{thm:T6}.
\end{proof}

\subsection{\texorpdfstring{The lemmas used in the proof of Lemma~\ref{lem:window}}{The lemmas used in the proof of the window lemma}}\label{app:G1:part-D2}

Throughout this subsection $L=2$, $a=1$, $a_\flat=\tfrac34$, $p=6$, $\mu_*=m$, and $C_\varphi:=2^{3/4}(1+\log2)^6c_1=39.622\ldots\,c_1$, where $c_1$ is Ba{\l}aban's constant of stage 8.

\begin{lemma}[monotonicity]\label{app:G1:mon}
Let $a>0$ and $q\ge0$.
\begin{enumerate}
\item[(a)] The function $f(t):=t^a\ell(t)^q$ is increasing on $(0,\mathrm{e}^{1-q/a}]$. Hence, for $\eps_0\le\mathrm{e}^{1-q/a}$, ``$f(t)\le c$ for all $t\in(0,\eps_0]$'' holds if and only if $f(\eps_0)\le c$.
\item[(b)] Let $\varrho>1$ and $\delta R>0$, and put $g(t):=t^{-a}\exp(-\delta R\,\ell(t)^{\varrho})$. Then $g$ is increasing on $(0,t_g]$, where $t_g:=\exp\big(1-(a/(\delta R\varrho))^{1/(\varrho-1)}\big)$, and $g(t)\to0$ as $t\downarrow0$.
\item[(c)] Let $F:(0,1]\to[0,\infty)$ satisfy $F(t)\to0$ as $t\downarrow0$, and let $c>0$, where neither $F$ nor $c$ depends on $\eps_0$. Consider the condition ``$F(t)\le c$ for all $t\in(0,\eps_0]$''. If it holds at $\eps_0$, it holds at every $\eps_0'\le\eps_0$. Moreover $\theta:=\sup\{s\in(0,1]:F\le c\text{ on }(0,s]\}>0$, and the condition holds for every $\eps_0<\theta$. The same is true for the reciprocal form ``$G(t)\ge c$ for all $t\in(0,\eps_0]$'' with $G(t)\to\infty$.
\end{enumerate}
\end{lemma}
\begin{proof}
(a) Since $\ell'(t)=-1/t$, $f'(t)=t^{a-1}\ell(t)^{q-1}\big(a\ell(t)-q\big)$. This is $>0$ when $\ell(t)>q/a$, i.e.\ when $t<\mathrm{e}^{1-q/a}$.
(b) $\frac{d}{dt}\log g=\big(-a+\delta R\varrho\,\ell(t)^{\varrho-1}\big)/t$. This is $>0$ when $\ell(t)^{\varrho-1}>a/(\delta R\varrho)$, i.e.\ when $t<t_g$. Also $\log g=a(\ell(t)-1)-\delta R\ell(t)^{\varrho}\to-\infty$ as $\ell(t)\to\infty$, because $\varrho>1$.

(c) The condition at $\eps_0'\le\eps_0$ is a statement about the smaller set $(0,\eps_0']\subset(0,\eps_0]$. Since $F(t)\to0<c$, there is $s>0$ with $F\le c$ on $(0,s]$, so $\theta>0$. If $\eps_0<\theta$, then $\eps_0\le s$ for some $s$ in the set defining $\theta$. The reciprocal form is the same argument applied to $1/G$.
\end{proof}

\begin{lemma}[the types of the conditions]\label{app:G1:ord}
Let $\pi=(e,\lambda,m,\eta,m_1,m_2)\in\fW$, and choose the data in the order of the stages 0--11 of Table~\ref{p:tab:stages}. Every constant that depends on $\eps_0$ and enters a condition of stage 10b is replaced by its supremum over $(0,\eps_1]$ (Lemma~\ref{app:G1:L10asup}), and every constant that depends on the side of the torus is evaluated at $\ell_0^\natural$ (Lemma~\ref{app:G1:floor}). Then:
\begin{enumerate}
\item[(i)] Each of the conditions of Tables~\ref{p:tab:A2}--\ref{p:tab:A5} is of type (ii) (a finite lower bound on $R$, $b_0$ or $\ell_0$), type (iii) (a downward-closed upper bound on $\eps_0$ with a positive threshold) or type (iv) (divisibility) of the proof of Lemma~\ref{lem:window}. The bound or threshold depends only on data of earlier stages and is finite, respectively positive, at every point of $\fW$.
\item[(ii)] The definition of the family $T_j=j!T_0$, $\eps=2^{-K}\eps_0$, $K\ge n_0+1$, imposes no condition on the regulator data, and $\eps\le1/m$ follows from (E.m2).
\item[(iii)] No condition bounds $b_0$ from above before $\eps_0$ is chosen, and no condition bounds $\ell_0$ from above. The only conditions that bound $R$ from above at fixed $\eps_0$ are (E.2) and the second part of it, $r(\eps_0)\le Cp(\eps_0)$; (E.2) is an upper bound on $\eps_0$, and $r(\eps_0)\le Cp(\eps_0)$ follows from the lower bound (b.2), $b_0\ge R/C$; both are imposed after $R$. No condition involves $\eps$, $K$ or $\ell_T$ except through $\ell_0^\natural$ and the lower bounds of stage 11.
\end{enumerate}
\end{lemma}
\begin{proof}
\emph{Reduction to Lemma~\ref{app:G1:mon}(c).} After stage 10a, every constant that enters an upper bound on $\eps_0$ is a function of the data of the stages $\le9$, of $\pi$, of $\eta$ and of $\ell_0^\natural$, but not of $\eps_0$ (Lemma~\ref{app:G1:L10asup}). Every upper bound on $\eps_0$ of Table~\ref{p:tab:A4} is either an explicit bound $\eps_0\le\theta$, or a condition ``for all $t\in(0,\eps_0]$''. The second kind is reduced by Lemma~\ref{app:G1:mon}(a),(b) to its value at $t=\eps_0$ plus an explicit bound. Its left side is built from $t^a\ell^q$ and $t^{-a}\exp(-\delta R\ell^{\varrho})$, or is a finite product or sum of these, so it tends to $0$; in the reciprocal form the left side tends to $\infty$. By Lemma~\ref{app:G1:mon}(c), such a condition is of type (iii).

\emph{Condition by condition.} The thresholds of (E.fin), (E.m2), (E.$\Phi$), (E.dom), (E.x) and (E.$\delta$) are explicit and positive at every point of $\fW$, since $m,m_2,\lambda>0$ and $e\neq0$.
\begin{center}\small
\begin{tabular}{@{}p{0.2\textwidth}p{0.36\textwidth}p{0.36\textwidth}@{}}
\hline
conditions & form & type; threshold \\ \hline
(E.fin) & $\eps_0\le\min(1/m_2,1/\mu_*,L)$ & (iii); $\min(1/m_2,1/m,2)$\\
(E.m2) & $m_2^2\eps_0^2\le1$ & (iii); $1/m_2$\\
(E.$\Phi$) & $\lambda^{-1/4}m_2^2\eps_0^{7/4}\le a_\flat/C_\varphi$ & (iii); $(\lambda^{1/4}a_\flat/(C_\varphi m_2^2))^{4/7}$\\
(E.dom) & $\mu_*\eps_0\le0.2$ & (iii); $0.2/m$\\
(E.x) & $\abs e\eps_0^{1/2}\le\min\{c_v,e_*,1\}$, $\lambda\eps_0\le c_v$ & (iii); $\min\{(\min\{c_v,e_*,1\}/\abs e)^2,\ c_v/\lambda\}$\\
(E.$\delta$) & $\kappa\mu_*^2\eps_0^2\le1/12$ & (iii); $(12\kappa m^2)^{-1/2}$\\
Table~\ref{p:tab:A2} & lower bounds on $R$; with $R\ge R^{**}$ & (ii) at stage 5; data of the stages 0--4 only (Lemma~\ref{app:G1:mass}(i))\\
(E.4), (E.5), (E.10), (E.14), (R.cap) & left sides of the form $\eps_0^{-a}\mathrm{e}^{-\delta R\hat\ell^2}$, or $R\hat\ell^2\ge c$ & (iii) by Lemma~\ref{app:G1:mon}(b),(c); the upper bound in (E.10) is void at $\delta_0R\ge4$\\
(E.ksp1)--(E.ksp3) & the conditions of Lemma~\ref{app:G1:RW-KSP} & (iii); for $t\le\eps_0$ the left sides tend to $0$\\
(E.1)--(E.3), (E.6)--(E.9), (E.11)--(E.13), (E.15), (E.r), (E.comb), (E.vac) & as in Table~\ref{p:tab:A4}, at each $M\in\mathbb M_4$ & (iii) (Lemma~\ref{app:G1:mon}(a),(c)); (E.6) and (E.7) in the reciprocal form\\
(b.1)--(b.4) & lower bounds on $b_0$ & (ii) at stage 9\\
Table~\ref{p:tab:A5} & lower bounds on $\ell_0$; $\ell_0\in2M_b\N$ & (ii) and (iv) at stage 11\\ \hline
\end{tabular}
\end{center}
For (E.$\Phi$) note that the left side $\lambda^{-1/4}m_2^2\eps_0^{7/4}$ is increasing in $\eps_0$, so the condition holds exactly up to the displayed threshold. For (E.vac), Theorems~\ref{thm:UV-K0minus} and \ref{thm:UV-K0plus} hold at each fixed Higgs mass $M>0$ for all sufficiently small $\eps_0$; with the four masses fixed at stage 0, this gives four positive thresholds.

Every threshold in the table depends only on data of earlier stages: on $\pi$, $\kappa$ and $\ell_0^\natural$ (stage 0), on Ba{\l}aban's constants (stages 3--4), on $R$ (stage 5), on $p$ (stage 6), on the constants of the stages 7--8, and on $b_0$ (stage 9). This proves (i) and (ii).

(iii) (E.2), $\hat x b_0R\hat\ell^8\le c$ with $\hat x=\abs e\eps_0^{1/2}$, bounds $R$ from above at fixed $\eps_0$. It is imposed at stage 10b, after $R$, as an upper bound on $\eps_0$. Its second part, $r\le Cp$, follows from the lower bound $b_0\ge R/C$, which is imposed at stage 9, after $R$. All other conditions on $R$, $b_0$ and $\ell_0$ in the tables are lower bounds. The index $K$ and the spacing $\eps$ enter no condition, because every scale used has $t_k\in(0,\eps_0]$ and the conditions are statements for all such $t$. The side of the torus enters only through $\ell_0^\natural$ and the bounds of stage 11.
\end{proof}

\begin{lemma}[$R$ is free of the masses; the four-mass rule]\label{app:G1:mass}
Let $\pi\in\fW$.
\begin{enumerate}
\item[(i)] $R^*$, and hence $R^{**}=\max\{R^*,4\ell_{\Nbar}+32M_b+25\}$, are functions of $(d,a,L,M_b,R_0,\delta_0,\allowbreak\nbar,\Nbar,\allowbreak V(\Nbar),C_{\mathrm{loc}}(\Nbar))$, of the $O(1)$ localisation widths and of $b_d$; no Higgs mass, coupling or photon mass enters. The remaining lower bounds on $R$ ($R_0+1$, $4/\delta_0$, and those from (E.10), (E.14) and (R.cap)) are also free of the Higgs mass. The bound from (E.14) contains only $\kappa$, $\mu_*$ and $(1-\eta)^{-3}$, which are finite on $\fW$. So $R\ge R^{**}$ is of type (ii).
\item[(ii)] Each entry of the Higgs mass $M$ in a constant or condition of the stages 8--11 is of one of two types. It is of type ($\alpha$) if $M$ enters from above through $M^2t^2$ with $t\le\eps_0$; then (E.m2) gives $M^2t^2\le m_2^2\eps_0^2\le1$ for every $M\le m_2$. It is of type ($\beta$) if $M$ enters through powers of either sign and logarithms of $M$, and through quantities finite at every fixed $M>0$ (among them $C_\delta(M)$, $1/\rho_{\min}(M)$, $(M^2\ell_T^3)^{-1}$, $(M\ell_T)^{-10}$, $\log_+(1/(M\eps_0))$ and the factor $(\gamma_1''M^2)^{-1}$ of (E.5)); such entries occur only in upper bounds on $\eps_0$, in monotone constants of the stages 8--9 and in the floor. No entry orders $M$ against $\mu_*$, and none uses $m_1>2m$ or $m_1\ge\sqrt2\mu_*$.
\item[(iii)] Consequently the four-mass rule can be met without creating a cycle, and each condition keeps its type.
\end{enumerate}
\end{lemma}
\begin{proof}
(i) The data on which $R$ depends lie in the stages 0--4, and the symbol $M$ in the lower bounds on $R$ is Ba{\l}aban's block size $M_b$, not a mass. The bound $4\ell_{\Nbar}+32M_b+25$ is free of the masses by inspection. The bound from (E.14) has as its only parameters $\kappa$, $\mu_*$ and $(1-\eta)^{-3}$; these are finite because $\mu_*=m$ and $\eta<1$.

(ii) The ingredients of the proof of Theorem~\ref{thm:T6} fall into the class ($\alpha$), into the class ($\beta$), or among those free of the mass, as follows.
($\alpha$) or free of the mass: the one-variable and two-variable remainder bounds (Lemmas~\ref{app:G1:I3r-U1} and \ref{app:G1:I3r-U1m}); the cube remainder (Lemma~\ref{app:G1:I3r-cube}); the separation lemma (Lemma~\ref{app:G1:I3r-SEP}); the two-domain bounds (Lemma~\ref{app:G1:I3r-DGM}); Lemmas~\ref{app:G1:I3r-Wsq} and \ref{app:G1:I3r-W2R}; Lemmas~\ref{app:G1:Rnest}, \ref{app:G1:Rthetab}, \ref{app:G1:Rloc7} and the support hypothesis (R4$'''$); the $A\cdot A$ family, whose $\Gamma_3,\Gamma_3'$ depend on $M^2t_{j_0}^2\le1$ only; the order cap (Proposition~\ref{app:G1:I5}); the region entropy (Proposition~\ref{app:G1:I6}); and (E.md), whose left sides increase in $M$.
($\beta$): the tail bound (Proposition~\ref{app:G1:T579}), through the $\delta m^2$ vertices; Proposition~\ref{app:G1:I4}, which needs only $-\Delta_A+M^2>0$; the final c-numbers (Proposition~\ref{app:G1:I8}: $\bar\ell=\min(1/M,1/\mu_*,L)$, $\rho_{\min}(M)\ge\min(M,\mu_*)\min\{1,(L\max(M,\mu_*))^{-1}\}$, constants polynomial in $1/\rho_{\min}$, $Z_f=t_{\bar K}/(m_f^2\ell_T^3)$); Lemma~\ref{app:G1:RW-MASS} and the scalar gap paid by (E.5) at the smallest mass; the scalar zero mode of Lemma~\ref{app:G1:RW-ZMs}, with gap $\ge M^2(1+4M^2\eps_0^2/(3a))^{-1}$; $C_\delta(M)$, finite for every $M>0$; $C_b'$, which contains $M+\mu_*$, $(M^2\ell_T^3)^{-1}$, $(M\ell_T)^{-10}$ and $\log_+(1/(\min(M,\mu_*)\eps_0))$; and $C_{F1}$, $\mathfrak E_k$, which are inherited from the first inductive step at mass $M$.
No entry of a third type occurs: orderings of $M$ against $\mu_*$, and the conditions $m_1>2m$, $m_1\ge\sqrt2\mu_*$, occur only in the anchor condition \eqref{p:ch-anchor} of Proposition~\ref{prop:ch-Tflat}, that is, in the conditions of Table~\ref{p:tab:A1}. The one place where the ordering of the two fields changes is the bound of the final c-numbers at $M_1<\mu_*$. It is covered by the symmetric construction of Proposition~\ref{app:G1:I8} with a common top index, which holds at $M_1$ by Lemma~\ref{lem:B:massrange}(a); at $(M_1,\mu_*)=(0.7,1)m$ one gets $\bar\ell=1/m$ and $\rho_{\min}\ge0.35$, the condition $\eps_0\le\min(1/M_1,1/\mu_*,L)$ follows from (E.fin), and $\eps_0$ is a power of $L$ by (T.div).

(iii) Each listed constant is finite at each $M\in\mathbb M_4$ by (ii). The maximum of four finite numbers is finite, and the minimum of four positive thresholds is positive. A larger constant at the stages 8--9 only raises $b_0$ and $\Theta$, which are lower bounds, and no condition bounds $b_0$ from above before $\eps_0$ (Lemma~\ref{app:G1:ord}(iii)). Each condition of stage 10b stays a downward-closed upper bound at the fixed $(R^{**},b_0,p)$ (Lemma~\ref{app:G1:mon}). The conditions (E.m2) and (E.$\Phi$) at $m_2$ imply the same conditions at every smaller $M$, because their left sides increase in $M$. The masses are data of stage 0, so the four-mass rule adds only edges from stage 0 into the stages 8--11. Adding data with edges forward to an acyclic graph whose edges do not decrease the stage keeps it acyclic.
\end{proof}

\begin{lemma}[the torus floor]\label{app:G1:floor}
Fix $\ell_0^\natural=\max\{8,128/m_1,128/M_1\}$ at stage 0, and evaluate at $\ell_T=\ell_0^\natural$ every constant that depends on $\ell_T$. These are: $C_\delta(M)$ and its majorants $\bar E(M,\ell)$; the loss in the extension of \cite[Prop.~2.1]{Ba2} to $Q^{(k)}$ (Lemma~\ref{app:G1:RW-KSP}); the zero-mode terms $(M^2\ell_T^3)^{-1}$, $(\mu_*^2\ell_T^3)^{-1}$, $\ell_T^{-1}$, $(M\ell_T)^{-10}$, $(\mu_*\ell_T)^{-10}$ of the bound on $\delta m^2_{\rm fin}$; the photon zero mode and the tail of $Q_0$; the twist defect $\mathrm{e}^{-\delta\ell_T/t}$; the zero-mode constants $Z_s,Z_v$ of Proposition~\ref{app:G1:I8}; and, through these, $c_C^{(0)}$, $c_L$, $c_7$, $C_>$. Then the value at $\ell_0^\natural$ bounds the constant on every torus of side $\ell_T\ge\ell_0^\natural$. Hence, with $\ell_0\ge\ell_0^\natural$ chosen at stage 11, every upper bound on $\eps_0$ decided at stage 10 remains true on every torus $T_j$, and no upper bound on $\eps_0$ depends on $\ell_0$. Moreover:
\begin{enumerate}
\item[(a)] $\ell_0^\natural=128/M_1=182.857\ldots/m$ whenever $M_1\le16$ (in the units of Section~\ref{sec:model}). Replacing the floor $\max\{8,128/m_1\}$ by $\ell_0^\natural$ can only weaken the upper bounds on $\eps_0$.
\item[(b)] The constants that depend on $\ell_T$ are used only in the range $\eps\le1/m$. Every spacing of the construction satisfies this, by (E.m2).
\end{enumerate}
\end{lemma}
\begin{proof}
Each listed constant is a non-increasing function $f$ of $\ell_T$. For $C_\delta$ every majorant, including the zero-mode terms $64\mathrm{e}^{-\mu^2\ell^2/8}/(\mu^2\ell^3)$ and $a^{-2\alpha}b^{-2\beta}/\ell^3$, decreases in $\ell$. The terms of the bound on $\delta m^2_{\rm fin}$ decrease; its logarithmic term uses $\ell^*=\min(\ell_T,\max(1/M,1/\mu_*))$ and is bounded by $\log_+(1/(\min(M,\mu_*)\eps_0))$ for every $\ell_T$. The terms $(\mu_*^2\abs T)^{-1}$, $\mathrm{e}^{-c\abs T}$ and $\mathrm{e}^{-\delta\ell_T/t}$ decrease. Also $Z_f$ decreases, and $C_>$ depends on $\ell_T$ only through its $\delta m^2_{\rm fin}$ and $C_\delta$ entries. Hence $f(\ell_T)\le f(\ell_0^\natural)$ for $\ell_T\ge\ell_0^\natural$.

Every upper bound on $\eps_0$ has the form $G(\eps_0;\text{constants})\le c$ and is monotone in each constant in the direction ``larger constant, stronger condition''; this is read off Table~\ref{p:tab:A4}, and in particular holds for (E.7), (E.ent) and the conditions of Lemma~\ref{app:G1:RW-KSP}. So replacing the constants by upper bounds keeps the conditions sufficient. The tori used have sides $\ell_j=j!\ell_0\ge\ell_0\ge\ell_0^\natural$ ((T.floor)).

(a) Since $M_1<m_1$, we have $128/M_1>128/m_1$, and $128/M_1\ge8$ exactly when $M_1\le16$. So the new floor is at least the old one. By monotonicity, every constant evaluated at the floor is no larger at the new floor, and every upper bound decided with the old values stays true. The new floor also dominates the lower bounds on $\ell_0$ tied to the masses ($M\ell_0\ge32$ at $M_1$, i.e.\ $\ell_0\ge45.71/m$; $96/m$; $128/m$; $\ell_T\ge L^3=8$). It is a single number computed from data of stage 0, so it creates no new dependence of earlier data on later data.

(b) The only other condition of validity of the constants that depend on $\ell_T$ is $\eps\le\eps_{\max}=1/m$. (E.m2) gives $\eps\le\eps_0\le1/m_2<1/m$.
\end{proof}

\begin{lemma}[the bound $\eps_1$ and the suprema of stage 10a]\label{app:G1:L10asup}
Let $\pi\in\fW$, and put
\[
\eps_1:=\min\{1,\eps_{(\mathrm{E.x})},\eps_{(\mathrm{E.md})},\eps_{(\mathrm{E}.\delta)},\eps_{(\mathrm{E.dom})},\eps_{(\mathrm{E.14})}\}.
\]
Here $\eps_{(\mathrm{E.x})}=\min\{(\min\{c_v,e_*,1\}/\abs e)^2,c_v/\lambda\}$, $\eps_{(\mathrm{E}.\delta)}=(12\kappa\mu_*^2)^{-1/2}$, $\eps_{(\mathrm{E.dom})}=0.2/\mu_*$, $\eps_{(\mathrm{E.md})}:=\min\{1/m_2,(\lambda^{1/4}a_\flat/(C_\varphi m_2^2))^{4/7}\}$, the largest $\eps_0$ that satisfies (E.m2) and (E.$\Phi$), and $\eps_{(\mathrm{E.14})}:=\sup\{s\le1:\text{(E.14) holds at }s\}$. Then $\eps_1>0$ is a function of data of the stages $\le9$ and of $\pi$, and does not depend on $\eps_0$ or $K$. Every constant that depends on $\eps_0$ and enters a condition of stage 10b is bounded on $(0,\eps_1]$ by a finite number free of $\eps_0$. These constants are $C'_{2.3}$, $\gamma_1$, $\gamma_1'$, $\gamma_1''$, $\bar\delta$, $c_K$, $C_{14}^w$, $C_P^w$, the term $C_\delta(M)^2\eps_0/(16\lambda)$ of $c_L$, $C_{\mathrm{fin}}$, the constant $O_\pi(1)$ of \cite[(3.42)]{Ba2}, $C_6(e,\lambda)$, $\bar v$, $c_A^{LF}+\bar v$, and $C_{\rm old}$.
\end{lemma}
\begin{proof}
\emph{Positivity.} Each of $\eps_{(\mathrm{E.x})}$, $\eps_{(\mathrm{E.md})}$, $\eps_{(\mathrm{E}.\delta)}$, $\eps_{(\mathrm{E.dom})}$ is an explicit positive number, because $e_*,c_v,a_\flat,C_\varphi,\kappa>0$ and $\lambda,m,m_2>0$ on $\fW$. The left sides of (E.m2) and (E.$\Phi$) increase in $\eps_0$, so $\eps_{(\mathrm{E.md})}$ is the largest $\eps_0$ that satisfies both. The left side of (E.14) increases in $\eps_0$ and tends to $0$ as $\eps_0\downarrow0$, since $\hat\ell\to\infty$ and $(1-\eta)^{-3}<\infty$. Hence $\eps_{(\mathrm{E.14})}>0$. It depends only on $\kappa$, $\mu_*$, $c_0^{\mathrm{ker}}$, $c_d$, $\delta_0$, $R$, $\eta$, all of the stages 0--5. None of the data of $\eps_1$ is $\eps_0$ or $K$.

\emph{Bounds on $(0,\eps_1]$.} Each item is monotone in $\eps_0$, or a polynomial in $\eps_0$ with $\eps_0\le1$. In detail:
\begin{itemize}
\item $C'_{2.3}$: its term $c_1\mu_*L\eps_0/a_k$ is at most $\tfrac43c_1\mu_*L\eps_1$, because $a_k>a_\flat=\tfrac34$.
\item $\gamma_1\ge(1+\mu_*^2\eps_1^2/a_\flat)^{-1}$ and $\gamma_1'\ge(1+4\mu_*^2\eps_1^2/(3a))^{-1}$: these are decreasing functions of $\eps_0$, bounded at $\eps_1$.
\item $\gamma_1''=a_\flat/(a+M^2\eps_0^2)\ge a_\flat/(a+1)$ for every $M\le m_2$, by (E.m2).
\item $\bar\delta\le1/12$ and $c_K\le1/3$ are a~priori bounds (Lemma~\ref{app:G1:RW-B17}).
\item $C^w_{14}$, $C^w_P$ are polynomials $\sum_ia_i\eps_0^i$, so $\sup_{(0,\eps_1]}\abs{\,\cdot\,}\le\sum_i\abs{a_i}$ since $\eps_1\le1$.
\item $C_\delta(M)^2\eps_0/(16\lambda)\le C_\delta(M)^2\eps_1/(16\lambda)$, with $C_\delta(M)$ evaluated at $\ell_0^\natural$ and at each $M\in\mathbb M_4$ (Lemmas~\ref{app:G1:mass}, \ref{app:G1:floor}).
\item $C_{\mathrm{fin}}\le c_7\eps_1^{\kappa_0}+c_5+\mathrm{e}^{-4}$ (with $\kappa_0=\frac15$), where $c_5$ and $c_7$ are bounded on $(0,\eps_1]$ by Statement~\ref{app:K:SL}. (The form $c_7\eps_1$ would hold only for $\kappa_0\ge1$.)
\item Under (E.ent), $O_\pi(1)$ is the constant $C_{\mathrm{fin}}$ of Proposition~\ref{app:G1:I6}, and it is bounded by the previous item. $C_6(e,\lambda)$ is bounded on $(0,\eps_1]$, since its dependence on $\eps_0$ is through monotone quantities.
\item $\bar v$ and $c^{LF}_A+\bar v$ are bounded on $(0,\eps_1]$; they are used only in the alternative large-field route of Remark~\ref{app:G1:S05-altroutes}.
\item $C_{\rm old}\le\bar S\,C_{PC_e}(1)(1+L^2/12)^{\bar N_{\rm ph}}$, a bound obtained through (E.$\delta$), which is part of $\eps_1$.
\end{itemize}
After this replacement, every constant that enters a condition of stage 10b is a function of the stages $\le9$, of $\pi$, of $\eta$ and of $\ell_0^\natural$, and not of $\eps_0$. This is the hypothesis ``$F$, $c$ independent of $\eps_0$'' of Lemma~\ref{app:G1:mon}(c).
\end{proof}

\subsection{The conditions named in the tables}\label{app:G1:part-Trows}

A condition ``for all $t\in(0,\eps_0]$'' is reduced to $t=\eps_0$ by Lemma~\ref{app:G1:mon}(a),(b). In this subsection $M_1$ is always the anchor mass $0.7m$.

\begin{lemma}[the scalar zero mode]\label{app:G1:RW-ZMs}
Work in physical units on the torus $T_\eps$, with the divisibility of Lemma~\ref{lem:sizes}(ii). Put $\Delta_s^{(0)}:=-\Delta^\eps+M^2$ and $\beta_k:=a\,t_{k+1}^{-2}$, and define recursively
\[
C_s^{(k)}:=(\Delta_s^{(k)}+\beta_kQ^*Q)^{-1},\qquad \Delta_s^{(k+1)}:=\beta_k-\beta_k^2\,Q\,C_s^{(k)}Q^*,
\]
where $Q$ is Ba{\l}aban's scalar block average at background $A=0$. Then for every $k\le K$ and every $M>0$,
\[
\Delta_s^{(k)}\ \ge\ \Big(M^{-2}+a^{-1}\sum_{j<k}t_{j+1}^2\Big)^{-1}\ \ge\ M^2\Big(1+\frac{4M^2\eps_0^2}{3a}\Big)^{-1},
\]
uniformly in $\ell_T$ and $\eps$. The first bound is attained on the constant mode, i.e.\ it is the eigenvalue on the constants. At $L=2$, $4/(3a)=1/a_\flat$. (The photon form is Lemma~\ref{app:G1:S05-ZM}.)
\end{lemma}
\begin{proof}
By the Woodbury identity, $(Q(\Delta_s^{(k)})^{-1}Q^*+\beta_k^{-1})^{-1}=\beta_k-\beta_kQ(\Delta_s^{(k)}+\beta_kQ^*Q)^{-1}Q^*\beta_k=\Delta_s^{(k+1)}$. Hence $(\Delta_s^{(k+1)})^{-1}=Q(\Delta_s^{(k)})^{-1}Q^*+\beta_k^{-1}$ exactly.

$Q$ is an average of point values whose weights sum to $1$. On $T_\eps$ the blocks tile every fine site with total weight $L^{-d}$. The Schur test with weights $t_{k+1}^d=L^dt_k^d$ therefore gives $\norm Q\le1$ and $Q^*1=1$.

So $\norm{(\Delta_s^{(k+1)})^{-1}}\le\norm{(\Delta_s^{(k)})^{-1}}+\beta_k^{-1}$, with $\norm{(\Delta_s^{(0)})^{-1}}=M^{-2}$, the constant mode included. Iterating gives the first bound. For the second, $\sum_{j<k}t_{j+1}^2\le\frac43t_k^2\le\frac43\eps_0^2$ at $L=2$.

On constants, $Q1=1$, $Q^*1=1$ and $\Delta_s^{(0)}1=M^21$. By induction $\Delta_s^{(k)}1=\lambda_s(k)1$ with $1/\lambda_s(k+1)=1/\lambda_s(k)+1/\beta_k$, which is the claimed equality.
\end{proof}
\begin{remark}\label{app:G1:RW-ZMsrem}
(i) The gap is positive for every $\eps_0$, so the scalar zero mode by itself imposes no upper bound on $\eps_0$. Under (E.m2), $m_2^2\eps_0^2\le1$, every $M\in\mathbb M_4$ (all are $\le m_2$) has gap $\ge M^2(1+4/(3a))^{-1}$, which is $3M^2/7$ at $a=1$; in particular it is finite and positive at $M_1$.
(ii) Ba{\l}aban's papers \cite{Ba1,Ba2} are themselves set on the torus, with constants independent of $\eps$ and $T_\eps$ \cite[(1.14), p.~606]{Ba1}. Lemma~\ref{app:G1:RW-ZMs} gives, with an explicit gap free of $\ell_T$, what that setting already contains.
(iii) The Gaussian factor of the constant mode is of the form $(2\pi/(M^2\eps^d))^{N/2}$ times a Jacobian. It sits in $E_{0,s}(M)$ and cancels.
\end{remark}

\begin{lemma}[lower bound on the averaged mass operator for every real background]\label{app:G1:RW-MASS}
Work in units of scale $k$, $t=L^k\eps$. Let $A$ be any real background and $\Omega$ any domain, with Neumann operator $-\Delta^N_{A,\Omega}\ge0$. Let $Q_k=Q_k(A)$ satisfy $Q_kQ_k^*=I$. Put $G_k(\Omega,A):=(-\Delta^N_{A,\Omega}+M^2t^2+a_kQ_k^*Q_k)^{-1}$ and $\Delta^{(k)}(\Omega,A):=a_k-a_k^2Q_kG_k(\Omega,A)Q_k^*$. Then for every $M>0$
\[
\frac{a_kM^2t^2}{a_k+M^2t^2}\ \le\ \Delta^{(k)}(\Omega,A)\ \le\ a_k.
\]
Consequently, the scalar part of \cite[(3.11), p.~585]{Ba2} holds for every real external field, without any regularity assumption:
\[
G''_k\ \ge\ \gamma_1M^2t^2\ \text{on }\Lambda_5^{(k)},\qquad \gamma_1=\frac{a_k}{a_k+M^2t^2}\ \ge\ \frac{a_\flat}{a_\flat+M^2t^2}\ \ge\ \frac{a_\flat}{a+M^2\eps_0^2}=:\gamma_1''.
\]
\end{lemma}
\begin{proof}
The two-sided bound is Lemma~\ref{lem:B:massrange}(d). For the consequence: principal blocks inherit lower bounds, and the averaging terms of \cite[(3.9)]{Ba2} are $\ge0$. The chain for $\gamma_1$ uses that $x\mapsto x/(x+c)$ is increasing, $a_k\ge a_\flat$ and $t\le\eps_0$. Only $Q_kQ_k^*=I$ is used. So $Q_k$ may come from any real background, or be a composite of single steps at different backgrounds.
\end{proof}

\begin{lemma}[the bound on the Hessian and condition (E.5)]\label{app:G1:RW-E06}
Fix $M>0$, and let $k$ satisfy $t_k\le\eps_0$. Let $V_k$ be the space of scalar fields on $\Lambda_5^{(k)}$ with $N=2$ real components, so that $\dim V_k=N\abs{\Lambda_5^{(k)}}$. Let $-\frac12\ip{\phi_k}{(G''_k-H''_k)\phi_k}$ be the $\phi_k$-quadratic part of the exponent in \cite[(3.9), p.~585]{Ba2}. Here $G''_k$ is the form of the first four terms there, and $H''_k$ is the form built from the changes (2.47) and from the operator $H_k$ of (3.7) of \cite{Ba2}; both are taken symmetric. Assume $G''_k\ge\gamma_1''M^2t_k^2$ on $V_k$ with $\gamma_1''=a_\flat/(a+M^2\eps_0^2)$ (Lemma~\ref{app:G1:RW-MASS}, Lemma~\ref{lem:B:massrange}(d)). Assume also $\norm{H''_k}\le h(t_k)$, where
\[
h(t_k):=C_{H''}\mathrm{e}^{-\delta'r(t_{k-1})}+C_Cc_d(\delta/2)\mathrm{e}^{-\delta r(t_k)/2}
\]
is the Schur bound of \cite[(3.12)]{Ba2}; the first term is absent at $k=0$. Put
\[
C_H:=C_{H''}+C_Cc_d(\delta/2),\qquad \delta_H:=\min\{\delta',\delta/2\}.
\]
\begin{enumerate}
\item[(a)] Suppose that
\begin{equation}\label{p:eq:E06}
h(t_k)\ \le\ \tfrac12\gamma_1''M^2t_k^2\qquad\text{for every $k$ with }t_k\le\eps_0 .
\end{equation}
Put $\varrho''_k:=h(t_k)/(\gamma_1''M^2t_k^2)$ and $X''_k:=G''^{-1/2}_kH''_kG''^{-1/2}_k$. Then $G''_k-H''_k\ge\frac12\gamma_1''M^2t_k^2$, $\norm{X''_k}\le\varrho''_k\le\frac12$,
\[
[\det(G''_k-H''_k)]^{-1/2}\le(\det G''_k)^{-1/2}\exp\big((\ln2)\varrho''_kN\abs{\Lambda_5^{(k)}}\big),\qquad \abs{\ln\det(\mathbf 1-X''_k)}\le2N\varrho''_k\abs{\Lambda_5^{(k)}},
\]
and $\norm{(G''_k-H''_k)^{-1}-G''^{-1}_k}\le2\varrho''_k(\gamma_1''M^2t_k^2)^{-1}$. The first bound is \cite[(3.11)]{Ba2} for $G''_k-H''_k$ with $\frac12\gamma_1''$ in place of $\gamma_1''$, as asserted in \cite[p.~586]{Ba2}. The determinant bounds are the scalar case of \cite[(3.18)--(3.20), pp.~586--587]{Ba2}, where $\varrho''_k$ bounds the norm of $X''_k$. The last bound is the one behind the replacement of $(G''_k-H''_k)^{-1}$ by $G''^{-1}_k$ in \cite[p.~587]{Ba2}.
\item[(b)] The condition (E.5) at $M$ reads
\[
2C_H(\gamma_1''M^2)^{-1}\eps_0^{-2}\exp(-\delta_HR\hat\ell^2)\le1,\qquad \eps_0\le\exp\big(1-1/(\delta_HR)\big).
\]
It implies \eqref{p:eq:E06}. It is downward closed in $\eps_0$, and it holds for all $\eps_0$ below a positive threshold that depends only on $C_H,\delta_H,R,a,L,M$. Its left side decreases in $M$, so at each $\eps_0$ the condition is strongest at $M_1$.
\item[(c)] The scalar zero mode and the logarithmic entries of Lemma~\ref{app:G1:RW-logp} impose no condition on $\eps_0$ beyond (E.m2) and (E.fin).
\end{enumerate}
\end{lemma}
\begin{proof}
(a) Write $g:=\gamma_1''M^2t_k^2$ and $h:=h(t_k)\le g/2$. For $v\in V_k$, $\ip{v}{(G''_k-H''_k)v}\ge(g-h)\abs v^2\ge\frac g2\abs v^2$. The operator $X''_k$ is symmetric, and $\norm{X''_k}\le g^{-1}\norm{H''_k}\le h/g=\varrho''_k\le\frac12$. From $G''_k-H''_k=G''^{1/2}_k(\mathbf 1-X''_k)G''^{1/2}_k$ we get $\det(G''_k-H''_k)=\det G''_k\det(\mathbf 1-X''_k)$. Let $x_1,\dots,x_n$, $n=N\abs{\Lambda_5^{(k)}}$, be the eigenvalues of $X''_k$; they lie in $[-\frac12,\frac12]$. The function $\varphi(x):=-\ln(1-x)$ is convex with $\varphi(0)=0$. Hence $0\le\varphi(x)\le2(\ln2)x$ for $0\le x\le\frac12$, and $\abs{\varphi(x)}=\ln(1+\abs x)\le\abs x$ for $-\frac12\le x<0$. So $\abs{\ln\det(\mathbf 1-X''_k)}\le\sum_i\abs{\varphi(x_i)}\le(2\ln2)n\varrho''_k\le2n\varrho''_k$, and $[\det(\mathbf 1-X''_k)]^{-1/2}=\exp(\frac12\sum_i\varphi(x_i))\le\exp((\ln2)n\varrho''_k)$. Finally, $(G''_k-H''_k)^{-1}-G''^{-1}_k=G''^{-1}_kH''_k(G''_k-H''_k)^{-1}$ has norm at most $g^{-1}h\cdot2g^{-1}=2\varrho''_k/g$.

The normalisation of \eqref{p:eq:E06} is the one used in \cite[(3.11), (3.18)--(3.20)]{Ba2}. The lower bound $\frac12\gamma_1''M^2t_k^2$ requires $h\le g/2$, and the determinant bound requires $\norm{X''_k}$ to stay away from $1$; both follow from $\varrho''_k\le\frac12$. Measured against $G''_k-H''_k$ instead of $G''_k$, the norm of $H''_k$ is at most $h/(g/2)=2\varrho''_k\le1$. This is the same condition, so no stronger form such as $h\le\frac14 g$ (equivalently, (E.5) with $4C_H$) is needed.

(b) The function $r$ is decreasing on $(0,1]$, and $t_{k-1}<t_k$, so $r(t_{k-1})\ge r(t_k)$. Hence $h(t_k)\le C_H\mathrm{e}^{-\delta_Hr(t_k)}$. Put $\psi(t):=C_Ht^{-2}\mathrm{e}^{-\delta_HR\ell(t)^2}$ with $\ell(t)=1+\ln t^{-1}$, and put $s=\ln t^{-1}$. Then $\frac{d}{ds}\ln\psi=2-2\delta_HR(1+s)$, which is $\le0$ exactly when $t\le t^*:=\exp(1-1/(\delta_HR))$. So $\psi$ is non-decreasing on $(0,t^*]$. By the second part of (E.5), $\eps_0\le t^*$, so $\psi(t)\le\psi(\eps_0)$ for $t\le\eps_0$. The first part of (E.5) says $\psi(\eps_0)\le\frac12\gamma_1''M^2$, and $\gamma_1''$ does not depend on $t$. Hence $h(t_k)\le\psi(t_k)t_k^2\le\frac12\gamma_1''M^2t_k^2$, which is \eqref{p:eq:E06}.

Now let $\eps_0'\le\eps_0$ with $\eps_0$ satisfying (E.5). Then $\eps_0'\le t^*$, $\psi(\eps_0')\le\psi(\eps_0)$ and $\gamma_1''(M,\eps_0')\ge\gamma_1''(M,\eps_0)$, so (E.5) holds at $\eps_0'$. For $\eps_0\le1$ we have $\gamma_1''\ge a_\flat/(a+M^2)$. Moreover $\ln\psi(\eps_0)=\ln C_H+2s-\delta_HR(1+s)^2\to-\infty$ as $s=\ln\eps_0^{-1}\to\infty$. This gives the positive threshold. Finally, $(\gamma_1''M^2)^{-1}=a/(a_\flat M^2)+\eps_0^2/a_\flat$ is decreasing in $M$, and the other factors do not depend on $M$.

(c) By Lemma~\ref{app:G1:RW-ZMs}, the scalar gap satisfies $\Delta_s^{(k)}(0)\ge\gamma_s(M;\eps_0)M^2$ with $\gamma_s(M;\eps_0)=(1+4M^2\eps_0^2/(3a))^{-1}$, uniformly in $\ell_T$ and $\eps$. This is not a condition on $\eps_0$. Under (E.m2) and $a=1$ it gives $\gamma_s\ge\frac37$ for $M\le m_2$. By Lemma~\ref{app:G1:RW-logp}(b),(c), once $\eps_0\le\min(1/m_2,1/\mu_*,L)$, which is (E.fin), the product $t^2\abs{\delta m^2_{\rm fin}}$ is bounded for $t\le\eps_0$ by a function of $t$ alone that tends to $0$ with $t$. So the factors $\log_L^+(\bar\ell/\eps_0)$ and $\log_+(1/(\min(M,\mu_*)\eps_0))$ add no condition on $\eps_0$. Every condition in which $t^2\delta m^2_{\rm fin}$ enters has the form $\sup_{0<t\le\eps_0}F(t)\le\theta$ with $F(t)\to0$, and so it is downward closed with a positive threshold (Lemma~\ref{app:G1:mon}(c)).
\end{proof}

\begin{lemma}[the logarithmic entries]\label{app:G1:RW-logp}
\begin{enumerate}
\item[(a)] In the proof of Theorem~\ref{thm:T6}, the logarithmic mass entries are those of the bound on the final mass counterterm of Proposition~\ref{app:G1:I8}. Put $\bar\ell:=\min(1/M,1/\mu_*,L)$ and assume $\eps_0\le\bar\ell$, with $\eps_0$ a power of $L$. Then
\[
\abs{\delta m^2_{\rm fin}(x)}\le C_b\Big\{(e^2+\lambda)\tfrac{L^2}{(L-1)\eps_0}+(e^4+e^2\lambda+\lambda^2)\big(2+\log_{L,+}(\bar\ell/\eps_0)\big)\Big\}+C_b^R\big\{\dots\big\},
\]
where $C_b=C(\alpha_0,\nbar,d,L,a,M_b)$ and $C_b^R$ is a polynomial in $1/\rho_{\min}$ (masses only). A weaker form of this bound carries $\log_+(1/(\min(M,\mu_*)\eps_0))$ instead.
\item[(b)] These entries enter a condition or an error constant only through the rescaled vertex $t_k^2\delta m^2(\cdot)$ with $t_k\le\eps_0$.
\item[(c)] For $0<t\le\eps_0\le\bar\ell$ and every $M>0$,
\[
t^2\log_{L,+}(\bar\ell/\eps_0)\ \le\ \frac{\bar\ell}{\mathrm{e}\ln L}\,t\ \le\ \frac{L}{\mathrm{e}\ln L}\,t\quad\Big(=1.0615\,t\text{ at }L=2\Big),\qquad t^2\log_+\frac1{\min(M,\mu_*)t}\le\frac{t}{\mathrm{e}\min(M,\mu_*)}.
\]
\end{enumerate}
In particular the logarithmic entries impose no condition beyond (E.fin).
\end{lemma}
\noindent Part (a) is the bound of Proposition~\ref{app:G1:I3b}.
\begin{proof}[Proof of (c)]
$x\log(1/x)\le1/\mathrm{e}$ for $x\in(0,1]$. With $x=t/\bar\ell$ and $\eps_0\ge t$, this gives $t\log_{L,+}(\bar\ell/\eps_0)\le t\log(\bar\ell/t)/\ln L\le\bar\ell/(\mathrm{e}\ln L)$. Multiplying by $t$ and using $\bar\ell\le L$ gives the first bound. The second follows in the same way with $x=\min(M,\mu_*)t$; if $x>1$ the left side is $0$.
\end{proof}

\begin{lemma}[the conditions (E.ksp1)--(E.ksp3) and (E.14)]\label{app:G1:RW-KSP}
Put $\vartheta(t):=\abs e\,t^{1/2}p(t)$ and $P_1(t):=\abs e\,c_1C\mu_*^{-1}t^{-1/2}p(t)$. Let $c_{AX}:=\sqrt d/(\mathrm{e}\,\delta_1)$, with $\delta_1$ fixed at stage 4. Let $C_B$ be the constant of the bound on the block-averaging operators $Q^{(k)}$ at real backgrounds (constants of stage 3 and $c_1$ of stage 8). The extension of \cite[Prop.~2.1]{Ba2} to $Q^{(k)}$ \cite[p.~579]{Ba2} needs the following upper bounds on $\eps_0$, each for all $t\in(0,\eps_0]$ and each at stage 10b.
\begin{enumerate}
\item[(E.ksp1)] $4c_{AX}C_B\abs e\,b_0\eps_0^{1/2}\hat\ell^{p}\le1$, together with $\eps_0\le\mathrm{e}^{1-2p}=\mathrm{e}^{-11}$. (E.8) does not imply (E.ksp1): the constants are unrelated.
\item[(E.ksp2)] With the capped radius $R_y:=\min\{\dist_j(y,\Omega^c)/2,\ r(t_k)L^{k-j}/2\}$:
\[
\text{(i)}\ \ (2d+3)\,\mathrm{e}^{-\delta_1r(t)/4}P_1(t)\le\tfrac12;\qquad \text{(ii)}\ \ \sqrt d\,C_B\,\mu_*\,t\,r(t)\le c_1C.
\]
As an alternative to (ii), one may keep the weight $\mathrm{e}^{-\delta_1R_y/2}$ and replace the factor $4$ of (E.ksp1) by $12$.
\item[(E.ksp3)] $4C_B\abs e^{1-\beta}t^{(1-\beta)/2}p(t)\le c$ with $\beta=1/56$. Equivalently, (E.1) with $C_{\mathrm{reg}}$ replaced by $\max\{C_{\mathrm{reg}},4C_B/c_1\}$, together with the bound $\eps_0\le\mathrm{e}^{1-2p/(1-\beta)}=\mathrm{e}^{-11.22}$.
\end{enumerate}
The condition (E.14) is used in the form of Table~\ref{p:tab:A4},
\[
4\kappa c_0^{\rm ker}c_d(1+\mu_*^2\eps_0^2)(1-\eta)^{-3}\,\frac{\mathrm{e}^{-\delta_0R\hat\ell^2/4}}{1-\mathrm{e}^{-\delta_0R\hat\ell\ln2/2}}\le\frac1{12};
\]
it enters stage 10a through $\eps_{(\mathrm{E.14})}$ and is imposed at stage 10b. Each of these conditions has a left side that tends to $0$ as $\eps_0\to0$, is downward closed below its explicit bound, depends only on the stages $\le9$ ($R$ at stage 5), and does not bound $\ell_T$ from above. So the order of choices stays acyclic.
\end{lemma}
\begin{proof}
(E.ksp1): the extension of \cite[Prop.~2.1]{Ba2} needs $c_{AX}E_j\le\frac12$ with $E_j\le2C_B\vartheta(t_k)$. The function $t^{1/2}\ell^p$ is increasing on $(0,\mathrm{e}^{1-2p}]$, so $\sup_{t\le\eps_0}\vartheta=\abs e\,b_0\eps_0^{1/2}\hat\ell^p$.

(E.ksp2): the corollary of that extension needs only $R_y\ge r(t_k)L^{k-j}/4$, and this holds because $\dist_k(Y,\Omega^c)\ge r(t_k)/2$. Then (i) follows from $\mathrm{e}^{-\delta_1R_y/2}\le\mathrm{e}^{-\delta_1r(t_k)/4}$. For (ii), $\sqrt dR_yE_j\le\sqrt d\,(r(t_k)L^{k-j}/2)(2C_B\vartheta_kL^{-2(k-j)})\le\sqrt d\,C_Br(t_k)\vartheta(t_k)/2$. The sufficient condition $\sqrt dC_Br\vartheta\le P_1$ is equivalent to $\sqrt dC_B\mu_*t\,r(t)\le c_1C$. Since $t\,r(t)=Rt\ell(t)^2$ is increasing on $(0,\mathrm{e}^{-1})$ (condition (E.r)), (ii) has a left side that tends to $0$.

(E.ksp3): regularity of $s(\tilde B_j+\tau a_j)$ in the estimate of the averaged fields requires $\sup\abs{\mathcal F}\le2E_j\le ce_j^{\beta}$; since $L^{-2m}\le L^{-\beta m/2}$ this gives the displayed condition. At $t=\eps_0$ the left side equals $4C_B\hat x^{1-\beta}b_0\hat\ell^6$, which is the form of (E.1).

(E.14): the deletion distance at level $j$ is $r(t_{j-1})/L\ge r(t_j)/L$, which gives the factor $\mathrm{e}^{-\delta_0R\ell_j^2/4}$ at $\varrho=L=2$. The sum runs over $t_j=2^{-m}\eps_0$ with $\ell_j=\hat\ell+m\ln2$, and $(\hat\ell+m\ln2)^2\ge\hat\ell^2+2m\hat\ell\ln2$. Summing the geometric series gives the displayed quotient.
\end{proof}

\begin{lemma}[the condition (E.ent)]\label{app:G1:RW-BM1}
Let $c_5,c_7$ be the constants of \cite[(3.42)]{Ba2} as used in Proposition~\ref{app:G1:I6}, taken at their suprema over $\eps_0\in(0,\eps_1]$, with $c_7$ at its maximum over $\mathbb M_4$. Let $\Theta:=C_1C_D^dR^d/(c_\square b_0^2\ln2)$ with $C_1:=n_d\ln2+d+1+c_7+c_5$ ($n_3=10$), $C_D:=C_D^0+14(1+\vartheta)$, $C_D^0:=A'(S_{\varrho,L}-1)$, $A'=16+23\vartheta$, $\vartheta:=M_bL/R\le\frac17$, $S_{\varrho,L}:=\sum_{m\ge0}L^{-m}(1+m\ln L)^{\varrho}$ ($S_{2,2}=2+4\ln2+6\ln^22=7.6553\ldots$, so $C_D\le144.36$), and $c_\square:=\min\{\frac18a3^{-d},\frac14\gamma_0 3^{-d},3^{-d}\}$. The condition at stage 10b is
\begin{multline*}
\text{(E.ent)}\qquad (1+\ln(2/\eps_0))^{2p-2-\varrho d}\ge\Theta\iff\eps_0\le\eps_{\rm ent}:=2\exp\big(1-\Theta^{1/(2p-2-\varrho d)}\big)\\
=2\exp(1-\Theta^{1/4})\qquad (p=6,\ \varrho=2,\ d=3).
\end{multline*}
Its left side increases as $\eps_0$ decreases, so it is a downward-closed upper bound. If half of the gain must be kept for another use, $\Theta$ is replaced by $2\Theta$. Under (E.ent), Proposition~\ref{app:G1:I6} gives \cite[(3.42)]{Ba2} on the torus.
\end{lemma}
\begin{proof}
$(1+\ln(2/\eps_0))^{4}\ge\Theta$ holds iff $\ln(2/\eps_0)\ge\Theta^{1/4}-1$, i.e.\ iff $\eps_0\le2\exp(1-\Theta^{1/4})$. The base $1+\ln(2/\eps_0)$ is the minimum of $\ell(t_j)$ over $j\le K-1$, because $t_j\le\eps_0/2$ there. The bound $\vartheta\le\frac17$ follows from $R\ge R^{**}\ge32M_b>7M_bL$. Since $C_D$ is increasing in $\vartheta$, its value $144.352\ldots$ at $\vartheta=\frac17$ bounds it.

The suprema of $c_5$ and $c_7$ over $\eps_0\in(0,\eps_1]$ are finite by Statement~\ref{app:K:SL}. Hence $\Theta$ is a finite positive number determined by $b_0$ (stage 9), $R$ (stage 5), constants of the stages 3 and 8, and $\eps_1$ (stage 10a); it depends on no datum of the stages 10b or 11. So $\eps_{\rm ent}>0$, and (E.ent) is a downward-closed upper bound on $\eps_0$ of stage 10b with a positive threshold, as used in Lemma~\ref{app:G1:acyc}.
\end{proof}

\begin{lemma}[the condition (E.$\delta$)]\label{app:G1:RW-B17}
Let $\kappa$ be the domination constant of stage 0, a fixed number $\kappa\ge\kappa_1\kappa_3$ (Theorem~\ref{app:G1:DOM-nat}); under (E.dom), $\mu_*\eps_0\le0.2$, every $\kappa\ge126528/16875$, in particular $\kappa=7.5$, is admissible (Remark~\ref{app:G1:DOM-kappa}). The condition (E.$\delta$) is
\[
\bar\delta:=\kappa\mu_*^2\eps_0^2\le\tfrac1{12},\qquad\text{i.e.}\quad \eps_0\le\eps_{(\mathrm E.\delta)}=(12\kappa\mu_*^2)^{-1/2},
\]
and it enters the bound $\eps_1$ of stage 10a. Consequences:
(a) the slice weights satisfy $w_k=1-\eta\delta_k\in[1-\eta\bar\delta,1]\subset[11/12,1]$, so $\eta\bar\delta\le\eta/12<1$ (Lemma~\ref{lem:T6-rows}(a));
(b) the coefficient $c_K:=\kappa\mu_*^2L^2t_K^2$ of the top slice satisfies $c_K\le1/3$.
At $\mu_*=1$ and $\kappa=7.5$ the condition reads $\eps_0\le90^{-1/2}=0.10541$.
\end{lemma}
\begin{proof}
(a) $\delta_k=\kappa\mu_*^2t_{k+1}^2\le\kappa\mu_*^2\eps_0^2=\bar\delta\le\frac1{12}$, and $0<\eta<1$. (b) $t_K=\eps_0$ and $L=2$, so $c_K=L^2\bar\delta\le\frac4{12}$.
\end{proof}
\begin{remark}
The value $6.07$ is the limit of $\kappa_1\kappa_3$ as $\mu_*\eps_0\to0$ and $K\to\infty$; since $\kappa_1\kappa_3>164/27>6.07$ (Remark~\ref{app:G1:DOM-kappa}), it may not be used for $\kappa$.
\end{remark}

\begin{lemma}[the conditions (R.cap) and (T.cap) at the order cap $\Nbar$]\label{app:G1:RW-B20}
At the order cap $\Nbar=\nbar^2=169$ the conditions on the localisation cubes read
\[
\text{(R.cap)}\ \ R\hat\ell^2\ge C_{\rm loc}(\Nbar),\qquad \text{(T.cap)}\ \ \ell_T\ge3r(\eps_0)\eps_0+C_{\Nbar}\eps_0.
\]
Here $C_{\rm loc}(\Nbar)$ and $C_{\Nbar}$ are finite numbers fixed at stage 4 together with $\Nbar$. (R.cap) is met at stage 5 by $R\ge C_{\rm loc}(\Nbar)$, since $\hat\ell\ge1$, or else at stage 10b; (T.cap) is met at stage 11. The order of choices stays acyclic, free of $\eps$ and uniform over tori with $\ell_T\ge\ell_0$. If $R$ is raised to meet (R.cap), the condition (E.2) makes $\eps_0$ smaller (an effect that depends on $\Nbar$ but not on $\eps$).
\end{lemma}
\begin{remark}
The decay rate of the trees of \cite[(1.33)]{Ba3} at the cap is some $\delta(\Nbar)>0$ (or the rate $\delta_1/2$, free of the cap, extracted before the shrinking induction of Lemma~\ref{app:G1:VSR-A}). It is not Ba{\l}aban's $\delta_0$ of stage 3, which is the propagator rate used in $\delta_0R\ge4$.
\end{remark}

\subsection{An explicit point}\label{p:appA:example}

\begin{example}[an admissible point]\label{app:G1:pstar}\label{app:G1:Pstar2}\label{app:G1:Pstar2-full}
In units $m=1$, let $(e^2,\lambda)=(0.05,0.05)$, $\mu_*=1$, $M_1=0.7$, $H=1$, $\eta=\frac12$, $(m_1,m_2)=(3,4)$, so that $\mathbb M_4=\{0.7,1,3,4\}$, with $L=2$, $a=1$, $a_\flat=\frac34$, $p=6$, $\varrho=2$, $\nbar=13$, $\Nbar=169$. Choose, in the order of the stages:
\begin{itemize}
\item $R:=\lceil\max\{R^{**},R_0+1,4/\delta_0,R_{(\mathrm{R.4})},\dots,R_{(\mathrm{R.8})},R_{(\mathrm{E.10})},R_{(\mathrm{E.14})},R_{(\mathrm{R.cap})}\}\rceil$. Here $R_{(\mathrm{R.4})},\dots,R_{(\mathrm{R.8})}$ are the explicit lower bounds of (R.4)--(R.8) in the forms of Table~\ref{p:tab:A2}. For the other conditions, $R_{(\cdot)}$ is the lower bound on $R$ that makes the condition $(\cdot)$ hold for all $\eps_0\le1$, if there is one; the bounds $R_{(\mathrm{E.10})}$ and $R_{(\mathrm{E.14})}$ also depend on $\kappa$, $\mu_*$ and $\eta$;
\item $b_0:=\max\{b_0^{\min},(\max\{2,\bar\mu\lambda^{-1/2}\})^{1/2},R/C\}$, with every constant of the stages 8--9 that depends on the Higgs mass taken at its maximum over $\mathbb M_4$;
\item $\eps_0:=2^{-n_0}$, where $n_0$ is the least integer such that $2^{-n_0}$ lies strictly below $\eps_1$ and below the threshold of every upper bound of Table~\ref{p:tab:A4} at this $R$ and $b_0$, at each mass of $\mathbb M_4$, (E.edge) included;
\item $\ell_0$ the least element of $2M_b\N$ that is $\ge\max\{183,128,96\}$ and above the lower bounds (T.nw), (T.cap), (T.T) and (T.per);
\item $T_j=j!\,T_0$ and $\eps=2^{-K}\eps_0$ with $K\ge K_0=\max(K_0',n_0+1)$.
\end{itemize}
Then every condition of Tables~\ref{p:tab:A1}--\ref{p:tab:A5} holds. Moreover, for every $j\ge1$, every such $K$ and every $k\le K$, the side $S_k:=\ell_j/t_k$ of $T_j$ in units of the scale-$k$ blocks satisfies $S_k\ge\max\{6\sqrt3\,r(t_k)+2,\,2/\delta_0\}$, which is the hypothesis of Lemma~\ref{lem:IMG}.

\emph{Order.} By Lemma~\ref{app:G1:acyc}, with the constants that depend on $\ell_T$ evaluated at $\ell_0^\natural$ (Lemma~\ref{app:G1:floor}) and the $\eps_0$-dependent constants entering stage 10b replaced by their suprema over $(0,\eps_1]$ (Lemma~\ref{app:G1:L10asup}), every condition depends only on data of its own stage and of earlier stages, and it constrains only the quantity chosen at its stage. The choices above are made in increasing order of stages. So it suffices to check each condition at its stage.

\emph{The point.} (P.1) and (P.2) hold because $0<\lambda=e^2=0.05<0.1$, $0<\eta=\frac12<1$ and $2<3<4$. The other conditions of Table~\ref{p:tab:A1} involve $R$, $\eps_0$ and $\ell_0$ only through ``$\ell_T\ge96$, $\eps\le1$'', and they hold by Corollaries~\ref{app:G1:Wa}, \ref{app:G1:Wb} and \ref{app:G1:Wc}: $a_{K,T}(1)-a_{K,T}(0.7)\ge0.3190043$; $g_{21}\ge(4-3)(4+3-0.3631)=6.6369$ and $S(m_1)-m_1\le0.3622-3<0<m_2$; $m_1=3>2>1.882\ge m_1^c$; $m_1=3\ge\sqrt2\,\mu_*$, $m<m_1<m_2$, $M_1=0.7<1$; and $\Pi(0.05,0.05)=0.4728641>0$. They are infima over all tori of side $\ge96$ (respectively $\ge64$). Every torus of the family has side $\ell_j\ge\ell_0\ge183$, and every spacing satisfies $\eps\le\eps_0\le\frac14$. So these conditions hold on every $T_j$.

\emph{$R$ and $b_0$.} By the caption of Table~\ref{p:tab:A2}, $R\ge R^{**}$ together with the explicit bounds of (R.4)--(R.8) gives every condition of Table~\ref{p:tab:A2}, (R.edge) included (Lemma~\ref{app:G1:R0rows}(i)). (R.nbar) reads $27\le3R$ and holds since $R\ge107$. Since $\delta_0R\ge4$, the part of (E.10) that bounds $\eps_0$ imposes nothing for $\eps_0\le1$. The three bounds defining $b_0$ are finite and free of $\eps_0$. Since $\hat\ell\ge1$, the $\eps_0$-free form of (b.3) implies the original.

\emph{$\eps_0$.} Each condition of Table~\ref{p:tab:A4} is downward closed and has a positive threshold (Lemma~\ref{app:G1:ord}, Lemma~\ref{app:G1:R0rows}(iii), Lemma~\ref{app:G1:acyc}(iii)). There are finitely many, the four-mass instances included, so $n_0$ exists. Some thresholds can be evaluated by hand:
\begin{itemize}
\item (E.m2) is $16\eps_0^2\le1$, i.e.\ $\eps_0\le0.25$;
\item (E.fin) is $\eps_0\le\min(1/m_2,1/\mu_*,L)=0.25$;
\item (E.dom) is $\eps_0\le0.2$;
\item (E.$\delta$) at $\kappa=7.5$ is $\eps_0\le90^{-1/2}=0.10540\ldots$;
\item (E.r) follows from (E.m2);
\item (E.charge) reads $\eps_0\le\min\{1,400\}=1$;
\item (E.$\Phi$), that is $\lambda^{-1/4}m_2^2\eps_0^{7/4}\le a_\flat/C_\varphi$ with $C_\varphi=39.622\,c_1$, holds if $\eps_0\le(1.89\cdot10^{-2}\cdot0.05^{1/4}/16)^{4/7}c_1^{-4/7}=0.0138418\,c_1^{-4/7}$, because $a_\flat/C_\varphi\ge1.89\cdot10^{-2}/c_1$.
\end{itemize}
Hence $n_0\ge4$, and $n_0\ge7$ if $c_1\ge1$.

The threshold of (E.edge) is also explicit in its constants. With $a_k$ replaced by $a$, (E.edge) at $L=2$, $d=3$ requires $F(t)\le1$ for all $0<t\le\eps_0/2$, where
\[
F(t):=C\,a\,c_1(\lambda t)^{-1/4}p(t)\big(1+r(t)\big)^{2}\exp\Big(-\delta_0\Big[\frac{r(t)-M_b}{2}-\sqrt3M_b\Big]\Big)(2t)^{-\kappa},
\]
with $C,c_1,\delta_0,M_b$ and the exponent $\kappa$ as in \eqref{eq:lem24-cap}. Since $a_k=a(1-L^{-2})(1-L^{-2k})^{-1}\le a$ for $k\ge1$ \cite[(2.13), p.~653]{K1}, and the left side of \eqref{eq:lem24-cap} increases in $a_k$, this form implies \eqref{eq:lem24-cap} at every $k\ge1$. Put $A:=\log(Cac_1\lambda^{-1/4}b_02^{-\kappa})+\delta_0(\frac{M_b}2+\sqrt3M_b)+2\log(1+R)$, $A_+:=\max\{A,0\}$, $B:=\kappa+p+\frac{17}4$, and $\eps_{24}:=\min\{1,2\exp(1-2(A_++B)/(\delta_0R))\}$. Let $0<t\le\eps_{24}/2$ and $s:=\log t^{-1}>0$. From $1+R(1+s)^2\le(1+R)(1+s)^2$ and $\log(1+s)\le s$ we get
\[
\log F(t)\le A_++Bs-\tfrac{\delta_0R}{2}(1+s)^2 .
\]
Moreover $1+s\ge2(A_++B)/(\delta_0R)$, so $\frac{\delta_0R}2(1+s)^2\ge(A_++B)(1+s)\ge A_++Bs$, and $F(t)\le1$. So every $\eps_0\le\eps_{24}$ satisfies (E.edge). This threshold does not depend on the Higgs mass, on $\eps$, $K$ or $\ell_0$. Here $\lambda^{-1/4}=2.1147\ldots$, but the constants $C,\kappa,\delta_0,c_1,M_b,R_0$ have no printed values in \cite{Ba1,Ba2}, and so neither has $\eps_{24}$. The same holds for the other thresholds; only their positivity is used.

\emph{$\ell_0$.} Each lower bound in the definition of $\ell_0$ is a finite function of data of the stages $\le$10b, and $2M_b\N$ is unbounded, so $\ell_0$ exists. We have $\ell_0\ge183\ge128/M_1=182.857\ldots=\ell_0^\natural$, so (T.floor) holds, together with $\ell_T\ge8$ and $M_1\ell_0\ge128.1\ge32$. Since $\eps_0\le1$ we have $r(\eps_0)\ge R\ge32M_b$. Hence
\[
6\sqrt3\,r+2-(8r+4M_b+2)=(6\sqrt3-8)r-4M_b\ge2.39\cdot32M_b-4M_b>0 ,
\]
so (T.per) dominates (T.nw). No condition of an earlier stage depends on $\ell_0$ (Lemma~\ref{app:G1:acyc}).

\emph{All tori and all scales.} Every condition of Table~\ref{p:tab:A5} is a lower bound on the side, so it holds on $T_j$, whose side is $\ell_j=j!\,\ell_0\ge\ell_0$. The scales satisfy $t_k=2^k\eps\le t_K=\eps_0<\mathrm{e}^{-1}$. The function $t\mapsto t(6\sqrt3\,r(t)+2)=6\sqrt3\,Rt(1+\log t^{-1})^2+2t$ is increasing on $(0,\mathrm{e}^{-1})$. So (T.per) gives $S_k=\ell_j/t_k\ge6\sqrt3\,r(t_k)+2$, and $S_k\ge\ell_0/\eps_0\ge2/\delta_0$. Finally $\eps\le\eps_0\le\frac14<1=1/m$.
\end{example}

\section{The lattice model: supplementary proofs}\label{p:appB}

This appendix proves the statements about the regularised model of Section~\ref{sec:model} and the standing properties (P1)--(P3) of Section~\ref{sec:C-lattice} that are used without proof there, together with the lattice facts behind the continuum limits of Section~\ref{sec:C-limits}: the identification of the action with King's and Ba{\l}aban's, the vanishing of the odd orders in $e$ of the counterterm, the operator $-\Delta_\eps=d^*d+dd^*$, gauge invariance, the law of the zero mode and the Gibbs form of the states (Subsection~\ref{app:G1:area-model}); the correlation inequalities of \cite[Sect.~6]{BFS1} on the torus graph, monotonicity in the photon mass, (P1)--(P3), the zero-mode split and the uniformity of the constants of Theorem~\ref{thm:ch-T9} over the tori (Subsection~\ref{app:G1:area-lattice}); multilinearity of the limits, local translation invariance at a fixed torus and monotonicity on the class $\cD$ (Subsection~\ref{app:G1:area-cont}). The letter $e$ is the charge; $\mathrm{e}$ is Euler's number.

\subsection{Model and states}\label{app:G1:area-model}

\subsubsection{The action and the odd orders of the counterterm}\label{app:G1:part-A1}

Throughout this subsection, $T_{j,\eps}$ is a lattice torus as in Subsection~\ref{ssec:tori}, sites are $x\in T_{j,\eps}$ and bonds are the positively oriented pairs $(x,\mu)$, each counted once.
For a function $f$ on sites put
\begin{gather*}
(\partial_\mu f)(x):=\eps^{-1}\big(f(x+\eps e_\mu)-f(x)\big),\\
(\partial_\mu^{*}f)(x):=\eps^{-1}\big(f(x-\eps e_\mu)-f(x)\big),\\
-\Delta_\eps:=\sum_{\mu=1}^3\partial_\mu^{*}\partial_\mu.
\end{gather*}
On the periodic lattice $\partial_\mu^*$ is the adjoint of $\partial_\mu$ for $\ip{\cdot}{\cdot}$, and all the operators $\partial_\mu,\partial_\nu^*$ commute, since they are linear combinations of translations.
A bond function $A$ is written $A_\mu(x):=A_{x,\mu}$.

\paragraph{The Feynman form and King's action.}

\begin{lemma}[the action of Subsection~\ref{ssec:action} is King's and Ba{\l}aban's]\label{app:G1:K1form}
Let $\mu_0>0$, $a\in\R$, $\lambda>0$ and $e\in\R$.
\begin{enumerate}
\item[(a)] On bond functions, $-\Delta_\eps:=d^*d+dd^*$ acts componentwise as the lattice Laplacian on sites: $(-\Delta_\eps A)_\nu=-\Delta_\eps A_\nu$. Hence $-\Delta_\eps\ge0$ on bond functions, and its kernel consists of the constant bond functions.
\item[(b)] $\tfrac12\ip{A}{(-\Delta_\eps+\mu_0^2)A}=\tfrac12\norm{dA}^2+\tfrac12\norm{d^*A}^2+\tfrac12\mu_0^2\norm{A}^2=\tfrac12\sum_\nu\ip{A_\nu}{(-\Delta_\eps+\mu_0^2)A_\nu}$.
\item[(c)] The matter action \eqref{eq:matteraction} can be written as
\[
S^{\mathrm m}_a(A,\varphi)=\sum_x\Big(3\eps\abs{\varphi(x)}^2+\eps^3\big(\tfrac12a\abs{\varphi(x)}^2+\lambda\abs{\varphi(x)}^4\big)\Big)-\sum_{x,\mu}\eps\,\rRe\big(\overline{\varphi(x)}\,\mathrm{e}^{ie\eps A_{x,\mu}}\varphi(x+\eps e_\mu)\big).
\]
\item[(d)] Let $q$ be the antisymmetric real $2\times2$ matrix with $\norm{q}=1$ of \cite[p.~651]{K1} and \cite[(1.7), p.~605]{Ba1} ($N=2$). Let $\iota:\R^2\to\C$ be the isometric identification with $\iota(qv)=i\,\iota(v)$. Write $S^\eps(\Omega,A,\phi)$ for King's action \cite[(2.4), p.~652]{K1}, with $\Omega=T_{j,\eps}$ and periodic operators \cite[(2.21), p.~654]{K1}, scalar mass $m$ and photon mass $\mu_0$. Write $S^{\mathrm{Ba}}(A,\phi)$ for Ba{\l}aban's action \cite[(1.11), p.~605]{Ba1}. Then for all $(A,\varphi)\in\R^{\text{bonds}}\times\C^{\text{sites}}$, with $\phi:=\iota^{-1}\circ\varphi$,
\[
S_{\mu_0,a}(A,\varphi)=S^\eps(T_{j,\eps},A,\phi)\big|_{\delta m^2\equiv a-m^2}-E_0-E_1
=S^{\mathrm{Ba}}(A,\phi)\big|_{m_0^2=a}-E.
\]
\item[(e)] Consequently the Gibbs factors $\mathrm{e}^{-S_{\mu_0,a}}$ and $\mathrm{e}^{-S^\eps(T_{j,\eps},\cdot)}$ differ only by the constant factor $\mathrm{e}^{E_0+E_1}$, so the normalised measures coincide. With $a=b_{T_j}(\eps)=a_{K,T_j}(m)$ (Definition~\ref{def:counterterm}), the model of Definition~\ref{def:states} at photon mass $\mu_0$ is King's torus model with $\delta m^2:=\delta m^2_{K,T}(\eps;m)$ computed at the reference photon mass $\mu_*$. This is King's own torus counterterm only when $\mu_0=\mu_*$ (Remark~\ref{rem:counterterm}(ii)).
\end{enumerate}
Part (d) is the identification stated after \eqref{eq:matteraction} in Subsection~\ref{ssec:action}, used in Proposition~\ref{prop:ch-CN}; part (a) is the statement on $-\Delta_\eps$ made there.
\end{lemma}

\begin{proof}
\emph{(a), (b)} Part (a) and the second equality of (b) are Lemma~\ref{app:G1:X}(i),(ii) below, which also gives the kernel and the spectral gap. The first equality of (b) is the definition of the adjoint, $\ip{A}{(d^*d+dd^*)A}=\norm{dA}^2+\norm{d^*A}^2$.

\emph{(c)} For each bond,
$\tfrac12\eps^{-2}\abs{\mathrm{e}^{ie\eps A_{x,\mu}}\varphi(x+\eps e_\mu)-\varphi(x)}^2
=\tfrac12\eps^{-2}\big(\abs{\varphi(x+\eps e_\mu)}^2+\abs{\varphi(x)}^2\big)-\eps^{-2}\rRe\big(\overline{\varphi(x)}\mathrm{e}^{ie\eps A_{x,\mu}}\varphi(x+\eps e_\mu)\big)$,
because $\abs{\mathrm{e}^{ie\eps A}}=1$.
Multiply by $\eps^3$ and sum over the bonds. Every site is the initial point of $3$ bonds and the end point of $3$ bonds, so the diagonal terms add up to $\sum_x3\eps\abs{\varphi(x)}^2$.

\emph{(d)} \emph{The identification $\iota$.} An antisymmetric real $2\times2$ matrix is $q=\theta\left(\begin{smallmatrix}0&-1\\1&0\end{smallmatrix}\right)$, and $\norm{q}=1$ forces $\theta=\pm1$. Take $\iota(v_1,v_2)=v_1+i\theta v_2$. Then $\iota$ is an isometry, it preserves Lebesgue measure, and $\iota(qv)=i\iota(v)$.
Hence $\iota\circ\exp(tq)=\mathrm{e}^{it}\iota$, and King's and Ba{\l}aban's $U(A_b)=\exp(\eps eqA_b)$ \cite[p.~651]{K1}, \cite[p.~605]{Ba1} becomes multiplication by $\mathrm{e}^{ie\eps A_b}$.

\emph{Matter part.} With $b=(x,\mu)$, $b_-=x$, $b_+=x+\eps e_\mu$, King's covariant derivative is \cite[(2.1)]{K1}
$(D_A^\eps\phi)(b)=\eps^{-1}(U(A_b)\phi(b_+)-\phi(b_-))$.
Its image under $\iota$ is $\eps^{-1}(\mathrm{e}^{ie\eps A_{x,\mu}}\varphi(x+\eps e_\mu)-\varphi(x))$.
By \cite[(2.2)--(2.3)]{K1}, $\tfrac12\ip{\phi}{-\Delta_A^{\eps,\Omega}\phi}=\tfrac12\sum_b\eps^3\abs{(D^\eps_A\phi)(b)}^2$, which is the gradient term of \eqref{eq:matteraction}.
Here the sum runs over the bonds of $T_{j,\eps}$, each once. This agrees with \cite[(1.8)]{Ba1} ($b\subset T_\eps$) and with $-\Delta_A^\eps=D_A^{\eps*}D_A^\eps$ in \cite[(1.11)]{Ba1}, and it is the reading under which $-\Delta^{\eps,\Omega}_A$ at $A=0$ is the lattice Laplacian $-\Delta_\eps$.
The remaining scalar terms of \cite[(2.4)]{K1} are $\tfrac12\ip{\phi}{m^2\phi}+\sum_x\eps^3(\tfrac12\delta m^2\abs{\phi}^2+\lambda\abs{\phi}^4)$. With $\delta m^2\equiv a-m^2$ these equal $\sum_x\eps^3(\tfrac12a\abs{\varphi}^2+\lambda\abs{\varphi}^4)$.
On a torus $\delta m^2(x)$ is constant in $x$.

\emph{Gauge part.} King's inner product \cite[(2.2)]{K1} on $A:\Omega\to\R^3$ is $\sum_{i}\sum_x\eps^3A_i(x)B_i(x)=\ip{A}{B}$. King's $-\Delta^{\eps,\Omega}$ on $\Omega=T_{j,\eps}$ is the periodic Laplacian acting on each component \cite[(2.4), (2.21)]{K1}.
By (b), $\tfrac12\ip{A}{(-\Delta^{\eps,\Omega}+\mu_0^2)A}=\tfrac12\ip{A}{(-\Delta_\eps+\mu_0^2)A}$.

\emph{Constants.} King's action contains $+E_0+E_1$ \cite[(2.4), (2.6), (3.35)]{K1}, and Ba{\l}aban's contains $+E$ \cite[(1.11)]{Ba1}. These are field-independent constants, omitted in \eqref{eq:latticeaction} (Remark~\ref{rem:counterterm}(iii)).

\emph{Ba{\l}aban's form.} \cite[(1.11)]{Ba1} is $\tfrac12\ip{\phi}{-\Delta^\eps_A\phi}+\sum_{x\in T_\eps}\eps^3(\tfrac12m_0^2\abs{\phi}^2+\lambda\abs{\phi}^4)+\tfrac12\ip{A}{(-\Delta^\eps+\mu_0^2)A}+E$. Here $-\Delta^\eps=\partial^{\eps*}\partial^\eps$ is the torus Laplacian and $U$ is as above.
With $m_0^2=a$, the same three identifications give the second equality.

\emph{(e)} This follows from (d), because $E_0+E_1$ does not depend on $(A,\varphi)$ and $\iota$ preserves Lebesgue measure. The statement about the coefficient is Definition~\ref{def:counterterm} (photon propagator $(-\Delta_\eps+\mu_*^2)^{-1}$) together with Remark~\ref{rem:counterterm}(ii).
\end{proof}

\paragraph{The odd orders in $e$ of the counterterm.}

Fix $M>0$ and the torus $T_{j,\eps}$. Let $\nu_M$ be the centred Gaussian measure on $\C^{\text{sites}}$ with density proportional to $\exp(-\tfrac12\ip{\varphi}{(-\Delta_\eps+M^2)\varphi})$, and let $\E_0$ be the expectation in $\gamma_{\mu_*}\otimes\nu_M$ ($\gamma_{\mu_*}$ as in Definition~\ref{def:states}). This measure is non-degenerate since $\mu_*,M>0$, so all its polynomial moments are finite.

For a polynomial $\delta=\sum_{2\le\alpha+2\beta\le4}e^\alpha\lambda^\beta\delta_{\alpha\beta}$ with real coefficients, define the interaction
\[
V_\delta:=S^{\mathrm m}_{M^2+\delta}(A,\varphi)-\tfrac12\ip{\varphi}{(-\Delta_\eps+M^2)\varphi}
=\sum_{n\ge1}e^nH^{(n)}+\sum_x\eps^3\big(\tfrac12\delta\abs{\varphi(x)}^2+\lambda\abs{\varphi(x)}^4\big),
\]
\[
H^{(n)}(A,\varphi):=-\sum_{x,\mu}\eps\,\rRe\Big(\overline{\varphi(x)}\,\frac{(i\eps A_{x,\mu})^n}{n!}\,\varphi(x+\eps e_\mu)\Big).
\]
The second expression for $V_\delta$ is Lemma~\ref{app:G1:K1form}(c) with the $n=0$ term removed: the $n=0$ term together with the diagonal term is $\tfrac12\ip{\varphi}{-\Delta_\eps\varphi}$.
The perturbative two-point function is the formal power series in $(e,\lambda)$ of $\E_0[\varphi_i(x)\varphi_k(y)\mathrm{e}^{-V_\delta}]/\E_0[\mathrm{e}^{-V_\delta}]$. By Wick's theorem its coefficients are sums of Feynman graphs with the following vertices: hopping vertices $e^nH^{(n)}$; quartic vertices $\lambda\abs\varphi^4$; counterterm vertices $e^{\alpha'}\lambda^{\beta'}\tfrac12\delta_{\alpha'\beta'}\abs\varphi^2$. Photon lines are $(-\Delta_\eps+\mu_*^2)^{-1}$ and scalar lines are $(-\Delta_\eps+M^2)^{-1}$.
Following \cite[(3.32)--(3.33), p.~659]{K1} and Definition~\ref{def:counterterm}, write $\Gamma(x,y;\alpha,\beta)$ for the coefficient of $e^\alpha\lambda^\beta$ in the sum of the amputated one-particle-irreducible graphs with two external scalar legs, other than the single counterterm vertex. The counterterm is the solution of
\begin{equation}\label{app:G1:eq-recursion}
-\delta_{\alpha\beta}+\sum_y\eps^3\,\Gamma(x,y;\alpha,\beta)=0,\qquad 2\le\alpha+2\beta\le4,
\end{equation}
solved recursively; $\Gamma(\cdot;\alpha,\beta)$ contains only coefficients $\delta_{\alpha'\beta'}$ with $\alpha'\le\alpha$, $\beta'\le\beta$, $(\alpha',\beta')\ne(\alpha,\beta)$. With $M$ and photon mass $\mu_*$ this is $\delta m^2_{K,T}(\eps;M)$.

\begin{lemma}[odd orders vanish]\label{app:G1:oddorders}
In the solution of \eqref{app:G1:eq-recursion}, $\delta_{\alpha\beta}=0$ whenever $\alpha$ is odd, and $\Gamma(x,y;\alpha,\beta)=0$ for odd $\alpha$. Hence
\[
\delta m^2_{K,T}(\eps;M)=e^2\delta_{20}+\lambda\delta_{01}+e^4\delta_{40}+e^2\lambda\,\delta_{21}+\lambda^2\delta_{02},
\]
and $a_{K,T}(M)=M^2+\delta m^2_{K,T}(\eps;M)$ is a polynomial in $(e^2,\lambda)$.
This is the statement in Definition~\ref{def:counterterm} that the odd orders in $e$ vanish.
\end{lemma}

\begin{proof}
\emph{Step 1 (vertices).} $H^{(n)}$ is homogeneous of degree $n$ in $A$: it is a sum of terms $A_{x,\mu}^n$ times a quadratic form in $\varphi$. The quartic and counterterm vertices do not contain $A$.
So a hopping vertex of order $e^n$ has exactly $n$ photon legs, and every other vertex has none.
The photon covariance does not depend on $e$.

\emph{Step 2 (contractions).} Under $\gamma_{\mu_*}\otimes\nu_M$ the fields $A$ and $\varphi$ are independent, and $\gamma_{\mu_*}$ is centred, hence invariant under $A\mapsto-A$.
So in a Wick contraction the photon legs are paired with each other. There are no external photon legs, because the external legs are scalar.
A graph with an odd number of photon legs therefore has value $0$. Equivalently, $\E_0[P]=0$ for every polynomial $P$ of odd degree in $A$, since $\E_0[P(A,\varphi)]=\E_0[P(-A,\varphi)]=-\E_0[P]$.

\emph{Step 3 (induction).} Order the pairs $(\alpha,\beta)$, $2\le\alpha+2\beta\le4$, in the order in which \eqref{app:G1:eq-recursion} determines them. Assume $\delta_{\alpha'\beta'}=0$ for all earlier pairs with $\alpha'$ odd.
Take a graph contributing to $\Gamma(\cdot;\alpha,\beta)$ with nonzero value. Let it have hopping vertices of orders $n_1,\dots,n_r$, $\beta_0$ quartic vertices, and counterterm vertices of orders $(\alpha_1,\beta_1),\dots,(\alpha_s,\beta_s)$. Then
\[
\alpha=\sum_k n_k+\sum_i\alpha_i,\qquad \beta=\beta_0+\sum_i\beta_i.
\]
Every vertex carries positive order, and the graph is not the single counterterm vertex. Hence each $(\alpha_i,\beta_i)$ is an earlier pair, and by hypothesis $\alpha_i$ is even (otherwise $\delta_{\alpha_i\beta_i}=0$ and the graph vanishes).
By Steps 1 and 2, $\sum_kn_k$ is the number of photon legs, which is even.
So $\alpha$ is even. Thus $\Gamma(\cdot;\alpha,\beta)=0$ for odd $\alpha$, and \eqref{app:G1:eq-recursion} gives $\delta_{\alpha\beta}=0$.
The pairs with $2\le\alpha+2\beta\le4$ are $(2,0),(0,1),(3,0),(1,1),(4,0),(2,1),(0,2)$; removing those with odd $\alpha$ leaves the five monomials stated.
\end{proof}

\begin{remark}\label{app:G1:oddorders-sym}
(i) Lemma~\ref{app:G1:oddorders} is the precise content of the symmetry named in Definition~\ref{def:counterterm}. Suppose $\delta$ is even in $e$. Then $H^{(n)}(-A,\varphi)=(-1)^nH^{(n)}(A,\varphi)$ gives
$V_\delta(-e;-A,\varphi)=V_\delta(e;A,\varphi)$: the substitution $A\mapsto-A$, $e\mapsto-e$ fixes $\varphi$ and leaves $V_\delta$ and $\gamma_{\mu_*}$ invariant. Hence the coefficient of $e^\alpha\lambda^\beta$ in the expansion of $\mathrm{e}^{-V_\delta}$ has $A$-parity $(-1)^\alpha$, and its $\E_0$-expectation, and that of its product with any polynomial in $\varphi$, vanishes for odd $\alpha$ (Step~2). In this form evenness in $e$ is carried along the recursion by Step~3.
(ii) Charge conjugation $(A,\varphi)\mapsto(-A,\bar\varphi)$ at fixed $e$ is also a symmetry of \eqref{eq:matteraction}. But it does not relate $e$ to $-e$, so it does not by itself give the lemma. It gives the vanishing of functions with an odd number of external photon legs.
(iii) Steps~1--3 count photon legs only. They therefore hold for every decomposition of the propagators into sums, in particular for the renormalisation-group graphs of \cite[(3.34)]{K1}.
\end{remark}

\subsubsection{The operator $-\Delta_\eps$, gauge invariance, the zero mode and the Gibbs form}\label{app:G1:part-A2}

We prove four statements of Section~\ref{sec:model}. Throughout, $j\ge1$ and $\eps$ are fixed, $N=\ell_j/\eps$ is even (Lemma~\ref{lem:sizes}), and $n:=3N^3$ is the number of bonds of $T_{j,\eps}$. For $\mu\in\{1,2,3\}$ let $\mathbf 1_\mu$ be the bond function equal to $1$ on the bonds parallel to $e_\mu$ and $0$ on the others.

\emph{Conventions.} (C1) Sites, bonds and plaquettes all carry the weight $\eps^3$ in $\ip{\cdot}{\cdot}$. Hence the adjoint $d^*$ is the transpose of $d$ in the standard coordinates, and the matrix of $-\Delta_\eps=d^*d+dd^*$ in the coordinates $(A_b)_b$ is symmetric. (C2) ``Centred Gaussian measure with covariance $C$'' on bond functions means the covariance operator for $\ip{\cdot}{\cdot}$: $\E[\ip fA\ip gA]=\ip f{Cg}$ for all bond functions $f,g$. For non-degenerate $C$ this measure has a Lebesgue density proportional to $\exp(-\tfrac12\ip{A}{C^{-1}A})$. This is the convention stated with \eqref{eq:C-state}. The value $s(\mu_0)=(\mu_0^2\abs{T_j})^{-1}$ in Definition~\ref{def:states} holds in this convention (Lemma~\ref{app:G1:zeromode}).

\begin{lemma}[the operator $d^*d+dd^*$ is the componentwise Laplacian]\label{app:G1:X}
Let $-\Delta_\eps:=d^*d+dd^*$ on bond functions, and for a function $f$ on sites put $(-\Delta_\eps f)(x):=\eps^{-2}\sum_{\nu=1}^3\big(2f(x)-f(x+\eps e_\nu)-f(x-\eps e_\nu)\big)$ (periodic shifts).
\begin{enumerate}
\item[(i)] For every bond function $A$ and every $\mu$, $(-\Delta_\eps A)_{x,\mu}=(-\Delta_\eps A_{\cdot,\mu})(x)$.
\item[(ii)] $\ip{A}{-\Delta_\eps A}=\norm{dA}^2+\norm{d^*A}^2=\sum_{x,\mu,\nu}\eps\,(A_{x+\eps e_\nu,\mu}-A_{x,\mu})^2\ge0$. Moreover $\Ker\Delta_\eps=\mathrm{span}\{\mathbf 1_1,\mathbf 1_2,\mathbf 1_3\}=\mathrm{ran}\,\Pi_0$, and $\Delta_\eps\Pi_0=\Pi_0\Delta_\eps=0$.
\item[(iii)] On $\mathrm{ran}\,\Pi^\perp$, $-\Delta_\eps\ge4\eps^{-2}\sin^2(\pi\eps/\ell_j)\ge16/\ell_j^2$. In particular $-\Delta_\eps+\mu_0^2$ is positive definite for every $\mu_0>0$, and $-\Delta_\eps$ is invertible on $\mathrm{ran}\,\Pi^\perp$.
\end{enumerate}
\end{lemma}
\begin{proof}
(i) By (C1) and the definition of $d$,
\[
(d^*A)(x)=\eps^{-1}\sum_\nu\big(A_{x-\eps e_\nu,\nu}-A_{x,\nu}\big),
\]
since $\ip{d\chi}{A}=\sum_{x,\nu}\eps^3\eps^{-1}(\chi(x+\eps e_\nu)-\chi(x))A_{x,\nu}$, and the periodic shift $x\mapsto x-\eps e_\nu$ is a bijection of $T_{j,\eps}$.
For a plaquette function $B$, extend it antisymmetrically, $B(x;\nu,\mu):=-B(x;\mu,\nu)$; the definition of $dA$ in Subsection~\ref{ssec:action} is antisymmetric in $(\mu,\nu)$. Then $\ip{dA}{B}=\tfrac12\sum_x\sum_{\mu\ne\nu}\eps^3(dA)(x;\mu,\nu)B(x;\mu,\nu)$. Collect the coefficient of $A_{x,\mu}$: the four terms of $(dA)(y;\alpha,\beta)$ that contain it come in two antisymmetric pairs. This gives
\[
(d^*B)_{x,\mu}=\eps^{-1}\sum_{\nu\ne\mu}\big(B(x;\mu,\nu)-B(x-\eps e_\nu;\mu,\nu)\big).
\]
Insert $B=dA$, with $(dA)(x;\mu,\nu)=\eps^{-1}(A_{x,\mu}+A_{x+\eps e_\mu,\nu}-A_{x+\eps e_\nu,\mu}-A_{x,\nu})$. Then
\[
(d^*dA)_{x,\mu}=\eps^{-2}\sum_{\nu\ne\mu}\big(2A_{x,\mu}-A_{x+\eps e_\nu,\mu}-A_{x-\eps e_\nu,\mu}\big)
+\eps^{-2}\sum_{\nu\ne\mu}\big(A_{x+\eps e_\mu,\nu}-A_{x,\nu}-A_{x+\eps e_\mu-\eps e_\nu,\nu}+A_{x-\eps e_\nu,\nu}\big),
\]
\[
(dd^*A)_{x,\mu}=\eps^{-1}\big((d^*A)(x+\eps e_\mu)-(d^*A)(x)\big)=\eps^{-2}\sum_{\nu}\big(A_{x+\eps e_\mu-\eps e_\nu,\nu}-A_{x+\eps e_\mu,\nu}-A_{x-\eps e_\nu,\nu}+A_{x,\nu}\big).
\]
The terms $\nu\ne\mu$ of the second line cancel the second sum of the first line. The term $\nu=\mu$ of the second line is $\eps^{-2}(2A_{x,\mu}-A_{x+\eps e_\mu,\mu}-A_{x-\eps e_\mu,\mu})$. Adding the two lines gives $((d^*d+dd^*)A)_{x,\mu}=\eps^{-2}\sum_\nu(2A_{x,\mu}-A_{x+\eps e_\nu,\mu}-A_{x-\eps e_\nu,\mu})$, which is~(i).

(ii) $\ip{A}{(d^*d+dd^*)A}=\ip{dA}{dA}+\ip{d^*A}{d^*A}$ by the definition of the adjoint. By~(i), and because $\sum_xf(x)f(x-\eps e_\nu)=\sum_xf(x+\eps e_\nu)f(x)$ on the torus, $\sum_x\eps^3f(x)(-\Delta_\eps f)(x)=\sum_{x,\nu}\eps\,(f(x+\eps e_\nu)-f(x))^2$ for each component $f=A_{\cdot,\mu}$. This gives the identity and $-\Delta_\eps\ge0$.
Since $-\Delta_\eps$ is symmetric and non-negative, $\Delta_\eps A=0$ if and only if $\ip{A}{-\Delta_\eps A}=0$, that is, if and only if every component has vanishing forward differences in every direction. Every site of the torus is reached from a fixed site by forward steps, so this means every component is constant, i.e.\ $A\in\mathrm{span}\{\mathbf 1_\mu\}=\mathrm{ran}\,\Pi_0$. Finally $\Delta_\eps\Pi_0=0$, and $\Pi_0\Delta_\eps=(\Delta_\eps\Pi_0)^*=0$ because $\Delta_\eps$ and $\Pi_0$ are self-adjoint.

(iii) For $k\in\{0,\dots,N-1\}^3$ put $f_k(x):=\mathrm{e}^{2\pi ik\cdot x/\ell_j}$, $x\in T_{j,\eps}$. These $N^3$ functions form an orthogonal basis of $\C^{T_{j,\eps}}$, $f_0$ is the constant, and $f_k(x+\eps e_\nu)=\mathrm{e}^{2\pi ik_\nu/N}f_k(x)$. Hence
\[
-\Delta_\eps f_k=\eps^{-2}\sum_\nu\big(2-2\cos(2\pi k_\nu/N)\big)f_k=4\eps^{-2}\sum_\nu\sin^2(\pi k_\nu/N)\,f_k.
\]
If $k\ne0$, some $k_\nu\in\{1,\dots,N-1\}$, and $\pi k_\nu/N\in[\pi/N,\pi-\pi/N]$ gives $\sin^2(\pi k_\nu/N)\ge\sin^2(\pi/N)$. Since $N\ge2$, $\pi/N\le\pi/2$, and $\sin t\ge2t/\pi$ on $[0,\pi/2]$ gives $4\eps^{-2}\sin^2(\pi/N)\ge4\eps^{-2}(2/N)^2=16/\ell_j^2$. So the real symmetric operator $-\Delta_\eps$ is $\ge16/\ell_j^2$ on the functions with zero sum. Now $\Pi_0A=\sum_\mu c_\mu\mathbf 1_\mu$, where $c_\mu$ is the mean of $A_{\cdot,\mu}$ (Lemma~\ref{app:G1:zeromode}(a)). So $\mathrm{ran}\,\Pi^\perp$ consists of the bond functions all of whose components have zero sum, and~(i) gives the bound on $\mathrm{ran}\,\Pi^\perp$. The remaining claims follow from $-\Delta_\eps\ge0$.
\end{proof}

\begin{lemma}[gauge invariance of the matter action]\label{app:G1:gauge}
For a real function $\chi$ on sites put $G_\chi(A,\varphi):=(A-\mathrm{e}^{-1}d\chi,\,\mathrm{e}^{i\chi}\varphi)$. For every real $a$ and all $(A,\varphi)$,
\[
S^{\mathrm m}_a(G_\chi(A,\varphi))=S^{\mathrm m}_a(A,\varphi),\qquad d(A-\mathrm{e}^{-1}d\chi)=dA,\qquad \abs{\mathrm{e}^{i\chi(x)}\varphi(x)}=\abs{\varphi(x)}.
\]
Hence every $F\in\cA_\eps$ and every product $F\,\mathrm{e}^{-S^{\mathrm m}_a}$ is invariant under $G_\chi$. The Gaussian part of \eqref{eq:latticeaction} is not invariant: $\norm{d^*A}^2$ and $\norm A^2$ change.
\end{lemma}
\begin{proof}
Let $y=x+\eps e_\mu$. Since $\eps e\cdot \mathrm{e}^{-1}(d\chi)_{x,\mu}=\chi(y)-\chi(x)$,
\[
\mathrm{e}^{ie\eps(A_{x,\mu}-\mathrm{e}^{-1}(d\chi)_{x,\mu})}\mathrm{e}^{i\chi(y)}\varphi(y)-\mathrm{e}^{i\chi(x)}\varphi(x)=\mathrm{e}^{i\chi(x)}\big(\mathrm{e}^{ie\eps A_{x,\mu}}\varphi(y)-\varphi(x)\big),
\]
and this holds also for the bonds that wrap around the torus. So every gradient term of \eqref{eq:matteraction} keeps its modulus. The single-site terms depend on $\abs{\varphi(x)}$ only. Finally $d(d\chi)=0$: each plaquette contains the four values $\chi$ at its corners, each once with each sign. Elements of $\cA_\eps$ are functions of $dA$ and $\abs\varphi^2$ (Definition~\ref{def:algebra}). $\norm{d^*A}^2$ and $\norm A^2$ change by $d^*d\chi$ and $d\chi$ terms, which do not vanish for non-constant $\chi$.
\end{proof}
\noindent\emph{Other parametrisations.} Section~\ref{sec:C-lattice} writes the gauge transformations as $(A+d\chi,\,U(-e\chi)\varphi)$. This is $G_{-e\chi}$, and $U(s)$ is multiplication by $\mathrm{e}^{is}$ on $\R^2\cong\C$. The proof of Lemma~\ref{lem:C-FP} uses $(G\mathrm{e}^{-S_m})(Z+d\chi,\varphi)=(G\mathrm{e}^{-S_m})(Z,U(e\chi)\varphi)$. This is Lemma~\ref{app:G1:gauge} applied with $G_{e\chi}$, since $G_{e\chi}(Z+d\chi,\varphi)=(Z,\mathrm{e}^{ie\chi}\varphi)$. Here $e\ne0$ is used.

\begin{lemma}[the zero mode under $\gamma_{\mu_0}$]\label{app:G1:zeromode}
Let $\mu_0>0$.
\begin{enumerate}
\item[(a)] $\ip{\mathbf 1_\mu}{\mathbf 1_\nu}=\abs{T_j}\delta_{\mu\nu}$, and $\Pi_0A=\sum_\mu c_\mu\mathbf 1_\mu$ with $c_\mu:=\abs{T_j}^{-1}\ip{\mathbf 1_\mu}{A}$, the mean of $A_{\cdot,\mu}$.
\item[(b)] The centred Gaussian measure $\gamma_{\mu_0}$ with covariance $C_{\mu_0}=(-\Delta_\eps+\mu_0^2)^{-1}$ exists and is non-degenerate. Under $\gamma_{\mu_0}$, $A^\perp:=\Pi^\perp A$ and $c=(c_1,c_2,c_3)$ are independent. $A^\perp$ is the centred Gaussian on $\mathrm{ran}\,\Pi^\perp$ with covariance $(-\Delta_\eps+\mu_0^2)^{-1}\Pi^\perp$. $c$ is the centred Gaussian on $\R^3$ with covariance matrix $s(\mu_0)\mathbf 1_3$, where $s(\mu_0)=(\mu_0^2\abs{T_j})^{-1}$.
\item[(c)] $(-\Delta_\eps)^{-1}\Pi^\perp$ is a positive definite operator on $\mathrm{ran}\,\Pi^\perp$, so the measure $\gamma^\fl$ of Definition~\ref{def:states} is well defined.
\end{enumerate}
\end{lemma}
\begin{proof}
(a) There are $N^3$ bonds parallel to $e_\mu$, and $N^3\eps^3=\ell_j^3=\abs{T_j}$. The $\mathbf 1_\mu$ are orthogonal, so the orthogonal projection onto their span is $\sum_\mu\abs{T_j}^{-1}\ip{\mathbf 1_\mu}{A}\mathbf 1_\mu$.

(b) $-\Delta_\eps+\mu_0^2$ is positive definite by Lemma~\ref{app:G1:X}(iii), so $C_{\mu_0}$ is a positive definite operator and $\gamma_{\mu_0}$ exists (Lemma~\ref{app:G1:gibbs}(i) gives its density). $A^\perp$ and $\Pi_0A$ are linear images of the same Gaussian vector, so they are jointly Gaussian. By Lemma~\ref{app:G1:X}(ii), $\Pi_0$ commutes with $\Delta_\eps$, hence with $C_{\mu_0}$. So for all bond functions $f,g$,
\[
\E_{\gamma_{\mu_0}}\big[\ip f{\Pi^\perp A}\ip g{\Pi_0A}\big]=\ip{\Pi^\perp f}{C_{\mu_0}\Pi_0g}=\ip{f}{\Pi^\perp\Pi_0C_{\mu_0}g}=0.
\]
Uncorrelated jointly Gaussian vectors are independent, so $A^\perp$ and $c$ (a linear function of $\Pi_0A$, by (a)) are independent.
The covariance of $A^\perp$ is $\Pi^\perp C_{\mu_0}\Pi^\perp=(-\Delta_\eps+\mu_0^2)^{-1}\Pi^\perp$. Since $\Delta_\eps\mathbf 1_\nu=0$, $C_{\mu_0}\mathbf 1_\nu=\mu_0^{-2}\mathbf 1_\nu$. Hence, by~(a),
\[
\mathrm{Cov}(c_\mu,c_\nu)=\abs{T_j}^{-2}\ip{\mathbf 1_\mu}{C_{\mu_0}\mathbf 1_\nu}=\abs{T_j}^{-2}\mu_0^{-2}\ip{\mathbf 1_\mu}{\mathbf 1_\nu}=\frac{\delta_{\mu\nu}}{\mu_0^2\abs{T_j}}.
\]
Equivalently, in terms of the density: $\ip{A^\perp}{(-\Delta_\eps+\mu_0^2)c}=\mu_0^2\ip{A^\perp}{c}=0$, so
\[
\tfrac12\ip{A}{(-\Delta_\eps+\mu_0^2)A}=\tfrac12\ip{A^\perp}{(-\Delta_\eps+\mu_0^2)A^\perp}+\tfrac12\mu_0^2\abs{T_j}\sum_\mu c_\mu^2.
\]
The density factorises, and the $c$-factor is the $N(0,s(\mu_0))$ density in each coordinate.

(c) $-\Delta_\eps$ maps $\mathrm{ran}\,\Pi^\perp$ into itself (it commutes with $\Pi^\perp$) and is $\ge16/\ell_j^2$ there, by Lemma~\ref{app:G1:X}(ii),(iii). So its inverse on $\mathrm{ran}\,\Pi^\perp$ exists and is positive definite.
\end{proof}

\begin{lemma}[Gibbs form of the states]\label{app:G1:gibbs}
Let $\mu_0>0$ and $a\in\R$. Let $Q:=\eps^3(-\Delta_\eps+\mu_0^2)$ be the matrix of the quadratic form $A\mapsto\ip{A}{(-\Delta_\eps+\mu_0^2)A}$ in the coordinates $(A_b)_b$, and $Z_\gamma:=(2\pi)^{n/2}(\det Q)^{-1/2}$. Let $\Psi^a_F$ and $W^a_0=\Psi^a_1$ be defined as in Definition~\ref{def:algebra}, with the coefficient $a$ in place of $b$.
\begin{enumerate}
\item[(i)] $\gamma_{\mu_0}(dA)=Z_\gamma^{-1}\mathrm{e}^{-\frac12\ip{A}{(-\Delta_\eps+\mu_0^2)A}}\prod_bdA_b$.
\item[(ii)] For every Borel $F\ge0$ on $\R^{\text{bonds}}\times\C^{\text{sites}}$, $A\mapsto\Psi^a_F(A)\in[0,\infty]$ is Borel, and
\[
\int F\,\mathrm{e}^{-S_{\mu_0,a}}\prod_bdA_b\prod_xd^2\varphi(x)=Z_\gamma\,\E_{\gamma_{\mu_0}}[\Psi^a_F].
\]
For complex Borel $F$ with $\E_{\gamma_{\mu_0}}[\Psi^a_{\abs F}]<\infty$, both sides are finite and equal.
\item[(iii)] $0<W^a_0(A)\le W^a_0(0)<\infty$. Hence $Z:=\int \mathrm{e}^{-S_{\mu_0,a}}\prod_bdA_b\prod_xd^2\varphi(x)=Z_\gamma\E_{\gamma_{\mu_0}}[W^a_0]\in(0,\infty)$, and $Z^{-1}\mathrm{e}^{-S_{\mu_0,a}}\prod_bdA_b\prod_xd^2\varphi(x)$ is a probability measure. For $a=b_{T_j}(\eps)$ and every $F\in\cA_\eps$ with $\E_{\gamma_{\mu_0}}[\Psi_{\abs F}]<\infty$ (in particular every bounded $F$),
\[
\langle F\rangle_{\eps,j,\mu_0}=Z^{-1}\int F\,\mathrm{e}^{-S_{\mu_0,b}}\prod_bdA_b\prod_xd^2\varphi(x).
\]
\item[(iv)] $0<\E_{\gamma^\fl}[W_0]\le W_0(0)$. For $F\in\cA_\eps$ bounded or non-negative, $\E_{\gamma^\fl}[\Psi_F]$ does not change if the period cell $[0,P)^3$ in the definition of $\gamma^\fl$ is replaced by any translate $\prod_\mu[t_\mu,t_\mu+P)$.
\end{enumerate}
\end{lemma}
\begin{proof}
(i) By (C1), $\ip{A}{(-\Delta_\eps+\mu_0^2)A}=A\cdot QA$ with $Q$ symmetric, and $Q$ is positive definite by Lemma~\ref{app:G1:X}(iii). The Gaussian integral gives $\int \mathrm{e}^{-\frac12A\cdot QA}\prod_bdA_b=(2\pi)^{n/2}(\det Q)^{-1/2}=Z_\gamma$. The normalised measure is the centred Gaussian with coordinate covariance $Q^{-1}=\eps^{-3}(-\Delta_\eps+\mu_0^2)^{-1}$. So $\E[\ip fA\ip gA]=\eps^6f\cdot Q^{-1}g=\eps^3f\cdot(-\Delta_\eps+\mu_0^2)^{-1}g=\ip f{C_{\mu_0}g}$, which is $\gamma_{\mu_0}$ in the convention (C2).

(ii) $S_{\mu_0,a}=\tfrac12\ip{A}{(-\Delta_\eps+\mu_0^2)A}+S^{\mathrm m}_a$ is continuous, so $F\mathrm{e}^{-S_{\mu_0,a}}$ is Borel. Tonelli's theorem on $\R^{\text{bonds}}\times\C^{\text{sites}}$, integrating first over $\varphi$, gives the Borel measurability of $\Psi^a_F$ and
\[
\int F\mathrm{e}^{-S_{\mu_0,a}}=\int\prod_bdA_b\,\mathrm{e}^{-\frac12\ip{A}{(-\Delta_\eps+\mu_0^2)A}}\,\Psi^a_F(A)=Z_\gamma\E_{\gamma_{\mu_0}}[\Psi^a_F],
\]
using~(i). For complex $F$, apply this to $\abs F$: then $F\mathrm{e}^{-S_{\mu_0,a}}$ is integrable, and Fubini's theorem, applied to the positive and negative parts of $\rRe F$ and $\operatorname{Im}F$, gives the identity.

(iii) Lemma~\ref{lem:postype} with $k=0$ gives $0<W^a_0(A)\le W^a_0(0)<\infty$. Its proof uses the coefficient of $\abs\varphi^2$ only through $\int_0^\infty \mathrm{e}^{-\eps^3(\frac12a\,r^2+\lambda r^4)}r\,dr<\infty$, which holds for every real $a$ because $\lambda>0$. Hence $0<\E_{\gamma_{\mu_0}}[W^a_0]\le W^a_0(0)$, and~(ii) with $F=1$ gives $Z=Z_\gamma\E_{\gamma_{\mu_0}}[W^a_0]\in(0,\infty)$. For $F\in\cA_\eps$ with $\E[\Psi_{\abs F}]<\infty$, (ii) gives
\[
\langle F\rangle_{\eps,j,\mu_0}=\frac{\E_{\gamma_{\mu_0}}[\Psi_F]}{\E_{\gamma_{\mu_0}}[W_0]}=\frac{Z_\gamma\E_{\gamma_{\mu_0}}[\Psi_F]}{Z_\gamma\E_{\gamma_{\mu_0}}[W_0]}=Z^{-1}\int F\mathrm{e}^{-S_{\mu_0,b}}.
\]
For bounded $F$, $\Psi_{\abs F}\le\norm F_\infty W_0\le\norm F_\infty W_0(0)$.

(iv) The bound is Lemma~\ref{lem:postype} again, since $\gamma^\fl$ is a probability measure (Lemma~\ref{app:G1:zeromode}(c)). Let $F$ be bounded or non-negative. By Lemma~\ref{lem:holonomy}, $g(c):=\Psi_F(A^\perp+c)$ is Borel and $P$-periodic in each $c_\mu$. For a $P$-periodic $h$ on $\R$, bounded or non-negative, $\int_t^{t+P}h=\int_0^Ph$: split $[t,t+P)$ at the point of $P\Z$ inside it and shift one piece by $P$. Apply this in each coordinate (Tonelli/Fubini) and then integrate over $A^\perp$.
\end{proof}

\subsection{Lattice standing assumptions}\label{app:G1:area-lattice}

\subsubsection{The correlation inequalities on the torus graph; monotonicity in the photon mass}\label{app:G1:part-B1}

We prove the correlation inequalities of \cite[Sect.~6]{BFS1} on an arbitrary finite graph and verify that the torus model \eqref{eq:C-action} is an instance; this is the input of Lemma~\ref{lem:C-cov}. Then we prove the monotonicity in the photon mass used in Proposition~\ref{prop:C-mu0}(a) and Lemma~\ref{lem:C-flatlimit}.

\paragraph{The inequalities of BFS~I, Sect.~6 on a finite graph.}\label{app:G1:ssec-BFS}
\cite[Sect.~6.1, p.~275]{BFS1} work on ``an arbitrary finite lattice'' with Dirichlet conditions, and remark (p.~276) that more general boundary conditions, e.g.\ periodic ones on rectangles, can be accommodated. We write their duplicate-variable argument \cite[Sect.~6.2--6.3, pp.~277--281]{BFS1} for an abstract finite graph. Nothing below uses a boundary, the lattice geometry or translation invariance, so the torus with its wrap-around bonds is an instance (Proposition~\ref{app:G1:BFStorus}).

\begin{definition}[finite-graph setting]\label{app:G1:BFSsetting}
Let $\Lambda$ be a finite set and $E$ a finite set of \emph{bonds}. Each $b\in E$ has endpoints $b_-,b_+\in\Lambda$, several bonds may join the same pair, and each bond carries a coupling $c_b\ge0$. For each $x\in\Lambda$ let $\nu_x$ be a nonzero finite measure on $[0,\infty)$ with $\int \mathrm{e}^{ar^2}d\nu_x(r)<\infty$ for every $a\ge0$. Angles $\vartheta_x\in\R/2\pi\Z$ carry the normalised Lebesgue measure $d\vartheta_x/2\pi$. Let $I$ be a finite index set, $\mu_C$ the centred Gaussian measure on $\R^I$ with covariance matrix $C\ge0$ (degenerate $C$ allowed), and for each $b\in E$ let $\bar A_b$ be a real linear functional on $\R^I$. Put
\begin{gather*}
u_b:=\vartheta_{b_+}-\vartheta_{b_-}+\bar A_b(A),\\
W:=\exp\Big(\sum_{b\in E}c_b\,r_{b_-}r_{b_+}\cos u_b\Big),\\
d\Pi:=\prod_{x\in\Lambda}d\nu_x(r_x)\frac{d\vartheta_x}{2\pi}\otimes d\mu_C(A).
\end{gather*}
Write $\omega=(r,\vartheta,A)$, $Z:=\int W\,d\Pi$, $\ev{F}:=Z^{-1}\int FW\,d\Pi$ and $\ev{F;G}:=\ev{FG}-\ev{F}\ev{G}$. Let $\mathcal P_r$ be the cone of polynomials in $(r_x)_{x\in\Lambda}$ with nonnegative coefficients. For $n\in\Z^\Lambda$ put $n\cdot\vartheta:=\sum_xn_x\vartheta_x$, and for $g\in[0,\infty)^\Lambda$ put $r^2(g):=\sum_xg_xr_x^2$. The class $\cI$ consists of the functions $\mathrm{e}^{r^2(g)}\mathrm{e}^{i\ell(A)}$ and the class $\cD$ of the functions $\mathrm{e}^{-r^2(g)}\mathrm{e}^{\ell(A)}$, where $g\ge0$ and $\ell$ is a real linear functional on $\R^I$.
\end{definition}
Write $\deg_c(x):=\sum_{b:\,b_-=x}c_b+\sum_{b:\,b_+=x}c_b$. Then $0<W\le\prod_x\exp(\tfrac12\deg_c(x)r_x^2)$, because $r_{b_-}r_{b_+}\le\frac12(r_{b_-}^2+r_{b_+}^2)$. Hence $Z\in(0,\infty)$. Moreover $\abs{P(r)}\mathrm{e}^{r^2(g)}\mathrm{e}^{a\abs A}$ is $W\,d\Pi$-integrable for every polynomial $P$, every $g\ge0$ and every $a\ge0$. This is the only integrability used below.

\begin{lemma}[angle doubling]\label{app:G1:angle}
For bounded measurable $f$ on $(\R/2\pi\Z)^2$,\quad $\displaystyle\int_{(\R/2\pi\Z)^2}f(\vartheta,\vartheta')\,d\vartheta\,d\vartheta'=\int_{[-\pi,\pi)^2}f(\sigma-\tau,\sigma+\tau)\,d\sigma\,d\tau$.
\end{lemma}
\begin{proof}
Let $L(\sigma,\tau):=(\sigma-\tau,\sigma+\tau)$, so $\abs{\det L}=2$. The map $L$ followed by reduction mod $2\pi\Z^2$ is a surjective homomorphism $\R^2\to(\R/2\pi\Z)^2$. Its kernel is $\Gamma:=L^{-1}(2\pi\Z^2)=\{\pi(a,b):a,b\in\Z,\ a\equiv b\bmod2\}$. This lattice has covolume $2\pi^2$, so $[-\pi,\pi)^2$ (area $4\pi^2$) consists of two $\Gamma$-cells. Since $f\circ L$ is $\Gamma$-periodic, $\int_{[-\pi,\pi)^2}f\circ L=2\int_{\R^2/\Gamma}f\circ L=2\abs{\det L}^{-1}\int_{(\R/2\pi\Z)^2}f$.
\end{proof}

\begin{lemma}[{\cite[Thm.~6.1]{BFS1} on a finite graph}]\label{app:G1:ginibre}
In the setting of Definition~\ref{app:G1:BFSsetting}, let $P,Q\in\mathcal P_r$, $n,m\in\Z^\Lambda$, and let $\ell,\ell'$ be real linear functionals on $\R^I$. Then
\[
\ev{P(r)\cos(n\cdot\vartheta+\ell(A));\ Q(r)\cos(m\cdot\vartheta+\ell'(A))}\ \ge0.
\]
\end{lemma}
\begin{proof}
Let $F:=P(r)\cos(n\cdot\vartheta+\ell(A))$ and $G:=Q(r)\cos(m\cdot\vartheta+\ell'(A))$.

\emph{Step 0 (duplication).} Let $\omega'=(r',\vartheta',A')$ be an independent copy of $\omega$. Expanding the right-hand side gives $Z^{-1}\int FGW\,d\Pi-Z^{-2}\int FW\,d\Pi\int GW\,d\Pi$, hence
\begin{equation}\label{app:G1:eq-dup}
\ev{F;G}=Z^{-2}\iint F(\omega)\,[G(\omega)-G(\omega')]\,W(\omega)W(\omega')\,d\Pi(\omega)\,d\Pi(\omega').
\end{equation}

\emph{Step 1 (new variables).} \emph{Gauge field.} Put $\alpha:=\frac12(A+A')$ and $\beta:=\frac12(A'-A)$, so that $A=\alpha-\beta$ and $A'=\alpha+\beta$. For $u,v\in\R^I$,
\[
\E\,\mathrm{e}^{i\ip{u}{\alpha}+i\ip{v}{\beta}}=\E\,\mathrm{e}^{i\ip{(u-v)/2}{A}+i\ip{(u+v)/2}{A'}}=\mathrm{e}^{-\frac18\ip{u-v}{C(u-v)}-\frac18\ip{u+v}{C(u+v)}}=\mathrm{e}^{-\frac14\ip{u}{Cu}-\frac14\ip{v}{Cv}}.
\]
Hence $\mu_C\otimes\mu_C$ is mapped to $\mu_{C/2}\otimes\mu_{C/2}$; only the fact that $\mu_C$ is a centred Gaussian is used. By linearity, $\ell(A)=\ell(\alpha)-\ell(\beta)$ and $\ell(A')=\ell(\alpha)+\ell(\beta)$, and likewise for $\ell'$ and $\bar A_b$.
\emph{Radii.} Put $s_x:=\frac12(r_x+r'_x)\ge0$ and $t_x:=\frac12(r_x-r'_x)$, and let $M_x$ be the image of $\nu_x\otimes\nu_x$; let $M:=\prod_xM_x$. Since $t_x\mapsto-t_x$ is the swap $r_x\leftrightarrow r'_x$, $M_x$ is invariant under $t_x\mapsto-t_x$, and it is supported in $\{s_x\ge0\}$. It has all moments.
\emph{Angles.} All integrands below are $2\pi$-periodic in each $\vartheta_x$ and $\vartheta'_x$, because $n,m$ are integer vectors. By Lemma~\ref{app:G1:angle}, applied site by site with Fubini, we may write $\vartheta=\sigma-\tau$ and $\vartheta'=\sigma+\tau$ with $\sigma,\tau\in[-\pi,\pi)^\Lambda$ independent and uniform.
\emph{Net effect.} $d\Pi(\omega)\,d\Pi(\omega')$ becomes $dM(s,t)\otimes dK(\sigma,\alpha)\otimes dK(\tau,\beta)$ with $K:=\prod_x\frac{d\sigma_x}{2\pi}\otimes\mu_{C/2}$. Thus $(\sigma,\alpha)$ and $(\tau,\beta)$ have the \emph{same} law $K$, independently of each other and of $(s,t)$.

\emph{Step 2 (algebra).} Put $h_b(\sigma,\alpha):=\sigma_{b_+}-\sigma_{b_-}+\bar A_b(\alpha)$, $X_b:=h_b(\sigma,\alpha)$ and $Y_b:=h_b(\tau,\beta)$; the same function $h_b$ is used on both copies. Then $u_b(\omega)=X_b-Y_b$ and $u_b(\omega')=X_b+Y_b$. Put $\pi^+_b:=s_{b_-}s_{b_+}+t_{b_-}t_{b_+}$ and $\pi^-_b:=s_{b_-}t_{b_+}+t_{b_-}s_{b_+}$. Then $r_{b_-}r_{b_+}=\pi^+_b+\pi^-_b$ and $r'_{b_-}r'_{b_+}=\pi^+_b-\pi^-_b$. From $\cos(X\mp Y)=\cos X\cos Y\pm\sin X\sin Y$,
\begin{align}
r_{b_-}r_{b_+}\cos u_b(\omega)+r'_{b_-}r'_{b_+}\cos u_b(\omega')&=2\pi^+_b\cos X_b\cos Y_b+2\pi^-_b\sin X_b\sin Y_b,\label{app:G1:eq-614}\\
r_{b_-}r_{b_+}\cos u_b(\omega)-r'_{b_-}r'_{b_+}\cos u_b(\omega')&=2\pi^-_b\cos X_b\cos Y_b+2\pi^+_b\sin X_b\sin Y_b,\nonumber
\end{align}
so that $W(\omega)W(\omega')=\prod_b\exp\big(2c_b[\pi^+_b\cos X_b\cos Y_b+\pi^-_b\sin X_b\sin Y_b]\big)$.
For $P\in\mathcal P_r$ write $P(s+t)=P_e+P_o$ and $P(s-t)=P_e-P_o$, where $P_e$ ($P_o$) collects the monomials of even (odd) total degree in $t$. Expanding $(s_x+t_x)^k$ shows that $P_e,P_o$ lie in the cone $\mathcal P_{s,t}$ of polynomials in $(s,t)$ with nonnegative coefficients. With $X_F:=n\cdot\sigma+\ell(\alpha)$, $Y_F:=n\cdot\tau+\ell(\beta)$, $X_G:=m\cdot\sigma+\ell'(\alpha)$ and $Y_G:=m\cdot\tau+\ell'(\beta)$,
\[
F(\omega)=P(s+t)\,[\cos X_F\cos Y_F+\sin X_F\sin Y_F],\qquad G(\omega)-G(\omega')=2Q_o\cos X_G\cos Y_G+2Q_e\sin X_G\sin Y_G.
\]

\emph{Step 3 (the paired cone).} Let $\mathcal K$ be the set of finite sums $\sum_kp_k(s,t)\,\Phi_k(\sigma,\alpha)\,\Phi_k(\tau,\beta)$, where $p_k\in\mathcal P_{s,t}$ and each $\Phi_k$ is a real measurable function on $[-\pi,\pi)^\Lambda\times\R^I$ with $\abs{\Phi_k(\sigma,\alpha)}\le c_k\mathrm{e}^{a_k\abs\alpha}$. The \emph{same} $\Phi_k$ is evaluated on the two copies. $\mathcal K$ is closed under sums and products, since $(p_1\Phi_1\Phi_1')(p_2\Phi_2\Phi_2')=(p_1p_2)(\Phi_1\Phi_2)(\Phi_1\Phi_2)'$. For every finite positive measure $M'$ that has all moments, is a product over $x$, is invariant under each $t_x\mapsto-t_x$ and is supported in $\{s\ge0\}$, and every $K\in\mathcal K$,
\begin{equation}\label{app:G1:eq-cone}
\int K\,dM'\,dK\,dK=\sum_k\Big(\int p_k\,dM'\Big)\Big(\int\Phi_k\,dK\Big)^2\ \ge0.
\end{equation}
The two $\Phi_k$-factors are integrated against the same law $K$, independently, which gives the square. Moreover $\int p_k\,dM'\ge0$: $M'$ is a product measure, so the integral of a monomial $s^at^j$ factorises over $x$. Each factor vanishes if $j_x$ is odd, by the symmetry $t_x\mapsto-t_x$, and is $\ge0$ if $j_x$ is even, because $s_x\ge0$.

\emph{Step 4 (assembly).} By Steps 1--2 the integrand of \eqref{app:G1:eq-dup} is
\begin{multline*}
P(s+t)[\cos X_F\cos Y_F+\sin X_F\sin Y_F]\cdot[2Q_o\cos X_G\cos Y_G+2Q_e\sin X_G\sin Y_G]\\
\cdot\prod_b\exp\big(2c_b[\pi^+_b\cos X_b\cos Y_b+\pi^-_b\sin X_b\sin Y_b]\big).
\end{multline*}
Expanding each exponential in its power series gives a series of elements of $\mathcal K$, because $c_b\ge0$, $\pi^\pm_b,P_e,P_o,Q_e,Q_o\in\mathcal P_{s,t}$, and the $\Phi$'s are products of $\cos$ and $\sin$ of $X$'s, hence bounded. Put $R_x:=s_x+\abs{t_x}=\max(r_x,r'_x)$. Since $s\ge0$, $\abs{\pi^+_b}+\abs{\pi^-_b}\le R_{b_-}R_{b_+}$. Hence the series of absolute values is dominated by
\[
P(R)\cdot2Q(R)\cdot\exp\Big(\sum_b2c_bR_{b_-}R_{b_+}\Big)\le P(R)\cdot2Q(R)\cdot\prod_x\exp\big(\deg_c(x)(r_x^2+r_x'^2)\big),
\]
where we used $R_x^2\le r_x^2+r_x'^2$. The right side is $\nu_x\otimes\nu_x$-integrable by the Gaussian moments. By dominated convergence the integral in \eqref{app:G1:eq-dup} is the sum of the integrals of the terms, and each is $\ge0$ by \eqref{app:G1:eq-cone} with $M'=M$. Since $Z>0$, $\ev{F;G}\ge0$.
\end{proof}

\begin{lemma}[{\cite[Thm.~6.2]{BFS1} on a finite graph}]\label{app:G1:BFS62}
In the setting of Definition~\ref{app:G1:BFSsetting}, let $g\ge0$, let $\ell$ be a real linear functional and $k$ a complex linear functional on $\R^I$, and let $b\in E$. Then
\[
\text{(1) }\ev{\mathrm{e}^{r^2(g)}\mathrm{e}^{i\ell(A)};\abs{k(A)}^2}\le0,\quad\text{(2) }\ev{\mathrm{e}^{r^2(g)}\mathrm{e}^{i\ell(A)};r_{b_-}r_{b_+}\cos u_b}\ge0,
\]
\[
\text{(3) }\ev{\mathrm{e}^{-r^2(g)}\mathrm{e}^{\ell(A)};\abs{k(A)}^2}\ge0,\quad\text{(4) }\ev{\mathrm{e}^{-r^2(g)}\mathrm{e}^{\ell(A)};r_{b_-}r_{b_+}\cos u_b}\le0.
\]
The expectations of elements of $\cI$ are real.
\end{lemma}
\begin{proof}
\emph{Symmetry.} The map $(A,\vartheta)\mapsto(-A,-\vartheta)$ preserves $d\Pi$, because $\mu_C$ is centred and the angles are uniform. It preserves $W$, since $u_b\mapsto-u_b$, and it leaves $r^2(g)$, $\abs{k(A)}^2$ and $r_{b_-}r_{b_+}\cos u_b$ invariant, while $\ell(A)\mapsto-\ell(A)$. Hence the $\sin\ell(A)$-parts drop out: in (1)--(2), $\mathrm{e}^{i\ell(A)}$ may be replaced by $\cos\ell(A)$, and $\ev{\mathrm{e}^{r^2(g)}\mathrm{e}^{i\ell(A)}}$ is real.

(2) Apply Lemma~\ref{app:G1:ginibre} with $P_N:=\sum_{j\le N}r^2(g)^j/j!\in\mathcal P_r$, $n=0$ and $\ell$, and with $Q=r_{b_-}r_{b_+}$, $m=\mathbf e_{b_+}-\mathbf e_{b_-}$ and $\ell'=\bar A_b$; here $\mathbf e_x$ is the indicator vector of $x$, and $m\cdot\vartheta+\bar A_b(A)=u_b$. Then let $N\to\infty$ by dominated convergence, since $\abs{P_N}\le \mathrm{e}^{r^2(g)}$.

(1) Let first $k$ be real. Then $k(A)^2=\lim_{\eta\downarrow0}2\eta^{-2}(1-\cos\eta k(A))$ with $0\le2\eta^{-2}(1-\cos\eta k)\le k^2$. Since $\ev{F;1}=0$,
\[
\ev{F;k(A)^2}=-\lim_{\eta\downarrow0}2\eta^{-2}\ev{F;\cos\eta k(A)}\le0
\]
for $F=\mathrm{e}^{r^2(g)}\cos\ell(A)$, by Lemma~\ref{app:G1:ginibre} (with $Q=1$, $m=0$, $\ell'=\eta k$) and the limit $P_N\to \mathrm{e}^{r^2(g)}$. For complex $k=k_1+ik_2$ use $\abs{k}^2=k_1^2+k_2^2$ \cite[(6.24)--(6.25)]{BFS1}.

(3), (4) This is the exchange of the two copies \cite[p.~281]{BFS1}. Relabelling $\omega\leftrightarrow\omega'$ in \eqref{app:G1:eq-dup} gives
\[
\ev{F;G}=-Z^{-2}\iint F(\omega')\,[G(\omega)-G(\omega')]\,WW'\,d\Pi\,d\Pi.
\]
Take $F=\mathrm{e}^{-r^2(g)}\mathrm{e}^{\ell(A)}$. Since $r'=s-t$ and $A'=\alpha+\beta$,
\[
F(\omega')=\mathrm{e}^{-s^2(g)-t^2(g)}\,\mathrm{e}^{2(st)(g)}\,\mathrm{e}^{\ell(\alpha)}\mathrm{e}^{\ell(\beta)},\qquad(st)(g):=\textstyle\sum_xg_xs_xt_x.
\]
The measure $dM^g:=\mathrm{e}^{-s^2(g)-t^2(g)}dM$ is again a product, invariant under each $t_x\mapsto-t_x$, supported in $\{s\ge0\}$, and has all moments. The factor $\mathrm{e}^{2(st)(g)}$ is a series with terms in $\mathcal P_{s,t}$, since $g\ge0$. The factor $\mathrm{e}^{\ell(\alpha)}\mathrm{e}^{\ell(\beta)}$ is a paired product $\Phi(\alpha)\Phi(\beta)$ with $\abs\Phi\le \mathrm{e}^{\norm\ell\abs\alpha}$. Finally $G-G'$ is as in Step~2. So the integrand is a series of elements of $\mathcal K$. Its absolute series is dominated as in Step~4 times $\mathrm{e}^{\abs{\ell(\alpha)}+\abs{\ell(\beta)}}$, using
\[
\mathrm{e}^{2(s\abs t)(g)}\mathrm{e}^{-(s^2+t^2)(g)}=\mathrm{e}^{-(s-\abs t)^2(g)}\le1
\]
(the series of absolute values of $\mathrm{e}^{2(st)(g)}$ is $\mathrm{e}^{2(s\abs t)(g)}$). This is integrable. By \eqref{app:G1:eq-cone} with $M'=M^g$, $\ev{\mathrm{e}^{-r^2(g)}\mathrm{e}^{\ell(A)};Q(r)\cos(m\cdot\vartheta+\ell'(A))}\le0$ for all $Q\in\mathcal P_r$, $m\in\Z^\Lambda$ and real $\ell'$. The choice $Q=r_{b_-}r_{b_+}$, $m=\mathbf e_{b_+}-\mathbf e_{b_-}$, $\ell'=\bar A_b$ gives (4). The choice $Q=1$, $m=0$, $\ell'=\eta k$, followed by the limit of (1), gives (3).
\end{proof}

\begin{lemma}[{\cite[Cor.~6.3]{BFS1} on a finite graph}]\label{app:G1:BFScov}
In the setting of Definition~\ref{app:G1:BFSsetting}, let $g\ge0$ and let $\ell$ be real linear.
\begin{enumerate}
\item[(a)] If $0\le C_0\le C_1$ in the positive semidefinite order on $\R^I$, and everything else is fixed, then $\ev{\mathrm{e}^{r^2(g)}\mathrm{e}^{i\ell(A)}}_{C_1}\le\ev{\mathrm{e}^{r^2(g)}\mathrm{e}^{i\ell(A)}}_{C_0}$ and $\ev{\mathrm{e}^{-r^2(g)}\mathrm{e}^{\ell(A)}}_{C_1}\ge\ev{\mathrm{e}^{-r^2(g)}\mathrm{e}^{\ell(A)}}_{C_0}$.
\item[(b)] For each $b$, $c_b\mapsto\ev{\mathrm{e}^{r^2(g)}\mathrm{e}^{i\ell(A)}}$ is non-decreasing and $c_b\mapsto\ev{\mathrm{e}^{-r^2(g)}\mathrm{e}^{\ell(A)}}$ is non-increasing on $[0,\infty)$.
\item[(c)] If $E'\subset E$ and the system on $E'$ is obtained by setting $c_b=0$ for $b\in E\setminus E'$, then the $\cI$-expectations on $E'$ are $\le$ those on $E$, and the $\cD$-expectations are $\ge$.
\end{enumerate}
\end{lemma}
\begin{proof}
(a) Let first $C_0,C_1$ be positive definite. Put $C_\theta:=(1-\theta)C_0+\theta C_1$ for $\theta\in[0,1]$ and $D:=C_1-C_0\ge0$. With $c_-:=\min(\lambda_{\min}(C_0),\lambda_{\min}(C_1))>0$ and $c_+:=\max(\norm{C_0},\norm{C_1})$ we have $c_-\le C_\theta\le c_+$. Let $\varrho_\theta(A)=(2\pi)^{-\abs I/2}\det(C_\theta)^{-1/2}\mathrm{e}^{-\frac12\ip{A}{C_\theta^{-1}A}}$ be the density of $\mu_{C_\theta}$. Then
\[
\partial_\theta\log\varrho_\theta=-\tfrac12\tr(C_\theta^{-1}D)+\tfrac12\ip{A}{C_\theta^{-1}DC_\theta^{-1}A},
\]
and hence $\abs{\partial_\theta\varrho_\theta(A)}\le(2\pi c_-)^{-\abs I/2}\big(\tfrac12\abs I\norm D/c_-+\tfrac12\norm D\abs A^2/c_-^2\big)\mathrm{e}^{-\abs A^2/(2c_+)}$, uniformly in $\theta\in[0,1]$. For $F$ in either class, $\Phi_F(A):=\int F\,W\prod_xd\nu_x\frac{d\vartheta_x}{2\pi}$ satisfies $\abs{\Phi_F(A)}\le c_F\mathrm{e}^{\norm\ell\abs A}$: we have $\abs{\mathrm{e}^{i\ell}}=1$, $\mathrm{e}^{-r^2(g)}\le1$, and $W$ is bounded uniformly in $A$ as above. So numerator and denominator of $\ev F_{C_\theta}$ are $C^1$ in $\theta$, with derivatives taken under the integral. The constant part of $\partial_\theta\log\varrho_\theta$ cancels in the normalised expectation, and we obtain
\[
\partial_\theta\ev{F}_{C_\theta}=\tfrac12\ev{F;\ip{A}{C_\theta^{-1}DC_\theta^{-1}A}}_{C_\theta}=\tfrac12\sum_jd_j\ev{F;k_j(A)^2}_{C_\theta},
\]
where $D=\sum_jd_jv_jv_j^{\mathsf T}$ with $d_j\ge0$ and $v_j$ orthonormal, and $k_j(A):=v_j^{\mathsf T}C_\theta^{-1}A$ is real linear. By Lemma~\ref{app:G1:BFS62}(1), resp.\ (3), applied at covariance $C_\theta$, this is $\le0$ for $F\in\cI$ and $\ge0$ for $F\in\cD$. Integrate over $\theta\in[0,1]$.

If $C_0,C_1$ are only positive semidefinite, apply this to $C_i+\eta\mathbf 1$ and let $\eta\downarrow0$. The map $C\mapsto\int\Phi_F\,d\mu_C=\E[\Phi_F(C^{1/2}\xi)]$, with $\xi$ standard normal on $\R^I$, is continuous on positive semidefinite $C$. Indeed, $\Phi_F$ is continuous by dominated convergence, $C\mapsto C^{1/2}$ is continuous, and $\abs{\Phi_F(C^{1/2}\xi)}\le c_F\mathrm{e}^{\norm\ell\norm{C^{1/2}}\abs\xi}$ is uniformly integrable for $C$ in bounded sets. The same holds for $F=1$, whose integral is $>0$.

(b) $\partial_{c_b}W=r_{b_-}r_{b_+}\cos u_b\cdot W$. The derivative is dominated, for $c_b$ in a bounded interval, by $r_{b_-}r_{b_+}\prod_x \mathrm{e}^{\frac12(\deg_c(x)+1)r_x^2}\abs F$, which is integrable. Hence $\partial_{c_b}\ev F=\ev{F;r_{b_-}r_{b_+}\cos u_b}$, and the sign follows from Lemma~\ref{app:G1:BFS62}(2), resp.\ (4). At $c_b=0$, continuity follows by dominated convergence. (c) follows by iterating (b).
\end{proof}

\begin{proposition}[the torus model is an instance]\label{app:G1:BFStorus}
Fix $(\eps,T)$, a lattice torus $T_{j,\eps}$ with $N\ge2$ sites per side, and the matter action $S_m$ of \eqref{eq:C-action}. Let $\mu_C$ be a centred Gaussian on the bond variables. Here $C\ge0$ is its covariance operator for $\ip\cdot\cdot_\eps$, as in \eqref{eq:C-state}; the corresponding covariance matrix of the coordinates $(A_b)$ is $\eps^{-3}C$, so the two positive semidefinite orders agree. Then the state $\ev{\cdot}_C$ of Lemma~\ref{lem:C-cov}, i.e.\ \eqref{eq:C-state} with $\mu_{C_{\mu_0}}$ replaced by $\mu_C$, is the state $\ev\cdot$ of Definition~\ref{app:G1:BFSsetting} with the following data:
\begin{itemize}
\item $\Lambda=T_{j,\eps}$, and $E$ = all $3N^3$ oriented bonds, wrap-around bonds included;
\item $c_b=\eps$ and $\bar A_b(A)=\eps eA_b$;
\item $d\nu_x(r)=r\,\mathrm{e}^{-V(r)}dr$ with $V(r)=\eps^3(\tfrac12b_Tr^2+\lambda r^4)+3\eps r^2$, and $I$ = bonds.
\end{itemize}
The same holds with any of the $c_b$ lowered to a value in $[0,\eps]$. Consequently Lemmas~\ref{app:G1:ginibre}, \ref{app:G1:BFS62} and \ref{app:G1:BFScov} hold on every lattice torus of the paper, for every $C\ge0$, in particular for $C=C_{\mu_0}$ with $\mu_0\in(0,\mu_*]$.
\end{proposition}
\begin{proof}
Write $\varphi_x=r_x(\cos\vartheta_x,\sin\vartheta_x)$, so that $d\varphi_x=r_x\,dr_x\,d\vartheta_x$. Expanding $\frac12\eps^3\abs{D_A\varphi(b)}^2=\frac12\eps(r_{b_+}^2+r_{b_-}^2)-\eps\,r_{b_-}r_{b_+}\cos(\vartheta_{b_+}-\vartheta_{b_-}+\eps eA_b)$ gives
\[
\mathrm{e}^{-S_m}=\prod_x\mathrm{e}^{-V(r_x)}\prod_b\exp\big(\eps r_{b_-}r_{b_+}\cos u_b\big),\qquad u_b=\vartheta_{b_+}-\vartheta_{b_-}+\eps eA_b,
\]
as in the proof of Lemma~\ref{lem:C-current}. Here the six diagonal halves $\frac12\eps r_x^2$ at each site give the term $3\eps r^2$ in $V$. By (P2), $\lambda>0$, so $\nu_x$ is nonzero and has Gaussian moments of all orders. The factors $2\pi$ from $d\vartheta_x$ cancel in normalised expectations. Covariance operator versus covariance matrix: if $E[\ip{f}{A}_\eps\ip{f'}{A}_\eps]=\ip{f}{Cf'}_\eps$, then with $f=\eps^{-3}\delta_b$ and $f'=\eps^{-3}\delta_{b'}$ one gets $E[A_bA_{b'}]=\eps^{-3}C_{bb'}$. Nothing in Definition~\ref{app:G1:BFSsetting} or in the proofs of Lemmas~\ref{app:G1:ginibre}--\ref{app:G1:BFScov} refers to how $E$ sits in a lattice. Multiple bonds between one pair of sites (at $N=2$) and the wrap-around bonds are bonds like any other.
\end{proof}

Lemma~\ref{lem:C-cov} is Lemma~\ref{app:G1:BFScov}(a) in the instance of Proposition~\ref{app:G1:BFStorus}. In the proof of Lemma~\ref{lem:C-cov}, the citations \cite[Thm.~6.2(1)]{BFS1} and \cite[Thm.~6.2(3)]{BFS1} may be read as Lemma~\ref{app:G1:BFS62}(1) and (3) on the torus graph, which removes the reliance on the periodic-boundary remark of \cite[p.~276]{BFS1}. The differentiation in the interpolation parameter is justified by the domination written in the proof of Lemma~\ref{app:G1:BFScov}(a). The inequalities compare two systems on \emph{one} bond set and one index set; they give no comparison between tori of different sizes, and none is used in the paper.

\paragraph{Monotonicity of the photon covariance in the photon mass.}\label{app:G1:ssec-mono}

\begin{lemma}[operator monotonicity]\label{app:G1:mu-mono}\leavevmode
\begin{enumerate}
\item[(a)] Let $Q,Q'$ be symmetric operators on a finite-dimensional real inner-product space with $0<Q'\le Q$. Then $Q^{-1}\le Q'^{-1}$.
\item[(b)] Fix $(\eps,T)$. For $0<\mu'\le\mu$, the photon covariances $C_\mu=(-\Delta_\eps+\mu^2)^{-1}$ of \eqref{eq:C-state} satisfy $C_\mu\le C_{\mu'}$ as quadratic forms on $(\R^{\text{bonds}},\ip\cdot\cdot_\eps)$, the zero mode included. On the constant 1-forms the inequality reads $\mu^{-2}\le\mu'^{-2}$. Equivalently, the variances $s(\mu)=1/(\mu^2\abs T)$ of the zero-mode coordinates $c_\mu$ are non-increasing in $\mu$.
\end{enumerate}
Part (b) is the monotonicity of $C_\mu$ in $\mu$, the zero mode included, used in the proof of Proposition~\ref{prop:C-mu0}(a), and the inequality $C_{\mu_0}\ge C_{\mu_*}$ for $\mu_0\le\mu_*$ used in the proof of Lemma~\ref{lem:C-flatlimit}. In both places Lemma~\ref{lem:C-cov} is then applied with $C_0=C_\mu$ and $C_1=C_{\mu'}$.
\end{lemma}
\begin{proof}
(a) Since $Q>0$, $Q^{\pm1/2}$ exist. From $Q'\le Q$ we get $B:=Q^{-1/2}Q'Q^{-1/2}\le\mathbf 1$, and $B>0$. A positive symmetric $B\le\mathbf 1$ has its eigenvalues in $(0,1]$, so $B^{-1}=Q^{1/2}Q'^{-1}Q^{1/2}\ge\mathbf 1$. Conjugating by $Q^{-1/2}$ gives $Q'^{-1}\ge Q^{-1}$.

(b) $-\Delta_\eps=d^*d+dd^*$ is symmetric for $\ip\cdot\cdot_\eps$ and $\ge0$, by Lemma~\ref{app:G1:X}(i),(ii). So $0<-\Delta_\eps+\mu'^2\le-\Delta_\eps+\mu^2$, and (a) gives $C_\mu\le C_{\mu'}$. Explicitly, let $(v_k)$ be an $\ip\cdot\cdot_\eps$-orthonormal eigenbasis of $-\Delta_\eps$ with eigenvalues $\kappa_k\ge0$. Then
\[
\ip{f}{(C_{\mu'}-C_\mu)f}_\eps=\sum_k\big[(\kappa_k+\mu'^2)^{-1}-(\kappa_k+\mu^2)^{-1}\big]\ip{v_k}{f}_\eps^2\ \ge0,
\]
with the terms $\kappa_k=0$ included. The constant 1-forms $\mathbf 1_\mu$ span the eigenvalue $0$ (Lemma~\ref{app:G1:X}(ii)), where $C_\mu\mathbf 1_\mu=\mu^{-2}\mathbf 1_\mu$, and the variance of $c_\mu$ is $s(\mu)=1/(\mu^2\abs T)$ by Lemma~\ref{app:G1:zeromode}(b), in agreement with \eqref{eq:C-state}ff.
\end{proof}

\subsubsection{Properties (P1)--(P3), the zero-mode split and the constants of \texorpdfstring{Theorem~\ref{thm:ch-T9}}{the tube bound}}\label{app:G1:part-B2}

Throughout, the notation is that of Section~\ref{sec:C-lattice}: $T=T_j$, $\ell_T=\ell_j=j!\,\ell_0$, $\abs T=\ell_T^3$, $N=\ell_T/\eps$, $-\Delta_\eps=d^*d+dd^*$ the componentwise lattice Laplacian on 1-forms, $P_0$ the $\ip{\cdot}{\cdot}_\eps$-orthogonal projection onto the constant 1-forms, $P^\perp:=1-P_0$, and $\mathbf 1_\mu$ the constant 1-form of direction $\mu$ (Subsection~\ref{ssec:C-standing}).

\paragraph{Property (P1); the per-torus coefficient.}

\paragraph{Property (P1); the coefficient is used torus by torus.}

\emph{Notation.} Fix $j\ge1$ and $\eps=\eps_K$, $K\ge K_0$, and write $\Lambda:=T_{j,\eps}$. Let $\mathcal H:=(\R^2)^\Lambda$ with the inner product $\ip{u}{v}_\eps=\sum_{x\in\Lambda}\eps^3\,u(x)\cdot v(x)$. The Higgs field $\varphi$ is an element of $\mathcal H$ through $\varphi(x)=(\operatorname{Re}\varphi(x),\operatorname{Im}\varphi(x))$. For $a\in\eps\Z^3$ let $\mathsf T_a$ be the translation: $(\mathsf T_au)(x):=u(x-a)$ on $\mathcal H$, and $(\mathsf T_aA)_{x,\mu}:=A_{x-a,\mu}$ on bond functions. A real $2\times2$ matrix $Q$ acts on $\mathcal H$ pointwise, $(Qu)(x):=Qu(x)$. We use the rotations $U(\theta)$, the complex conjugation $\mathsf K:=\mathrm{diag}(1,-1)$, and $J:=U(\pi/2)$. For a real function $f$ on $\Lambda$, $\mathsf M_f$ is the multiplication operator $(\mathsf M_fu)(x)=f(x)u(x)$. For $v\in\R^2$, $\mathbf 1\otimes v\in\mathcal H$ is the constant function $x\mapsto v$. For a linear operator $Y$ on $\mathcal H$ and $x\in\Lambda$, the \emph{zero-momentum value} $\hat Y(x)$ is the $2\times2$ matrix $v\mapsto\big(Y(\mathbf 1\otimes v)\big)(x)$. If $(Yu)(x)=\sum_y\eps^3Y(x,y)u(y)$, then $\hat Y(x)=\sum_{y\in\Lambda}\eps^3Y(x,y)$, which is the sum in \cite[(3.33), p.~659]{K1}.

\begin{lemma}[the lines commute with the symmetries]\label{app:G1:P1-lines}
For $M>0$ put $C_M:=(-\Delta_\eps+M^2)^{-1}$, acting on $\mathcal H$ componentwise, and $C^A:=(X+\mu_*^2)^{-1}$ on bond functions. Then $C_M$ commutes with $\mathsf T_a$ ($a\in\eps\Z^3$), with $U(\theta)$ ($\theta\in\R$) and with $\mathsf K$, and $C^A$ commutes with $\mathsf T_a$.
\end{lemma}
\begin{proof}
$(\Delta_\eps u)(x)=\eps^{-2}\sum_{\mu=1}^3\big(u(x+\eps e_\mu)+u(x-\eps e_\mu)-2u(x)\big)$. The map $x\mapsto x-a$ is a bijection of $\Lambda$ that commutes with $x\mapsto x\pm\eps e_\mu$, so $\Delta_\eps\mathsf T_a=\mathsf T_a\Delta_\eps$. Since $\Delta_\eps$ acts on each real component separately, it commutes with every pointwise matrix. On the periodic lattice, $X$ acts on each component $A_{\cdot,\mu}$ as $-\Delta_\eps$ (Lemma~\ref{app:G1:X}), so $X$ commutes with $\mathsf T_a$. The operators $-\Delta_\eps+M^2\ge M^2$ and $X+\mu_*^2\ge\mu_*^2$ are invertible. If $Y$ commutes with an invertible $W$, then $W^{-1}Y=W^{-1}YWW^{-1}=W^{-1}WYW^{-1}=YW^{-1}$.
\end{proof}

\emph{Formal perturbation theory on $\Lambda$.} Let $I:=\{(\alpha,\beta)\in\N_0^2:2\le\alpha+2\beta\le4\}$, and give the monomial $e^\alpha\lambda^\beta$ the weight $w(\alpha,\beta):=\alpha+2\beta$. Let $\mathfrak d=(\delta m^2_{\gamma})_{\gamma\in I}$ be a family of real functions on $\Lambda$; a priori they may depend on the site, as in \cite[(3.32), p.~659]{K1}. Put $\delta m^2_\gamma:=0$ for $\gamma\notin I$, and $\delta m^2:=\sum_{\gamma\in I}e^\alpha\lambda^\beta\delta m^2_\gamma$. For $M>0$, summation by parts on the torus and the identity $\tfrac12\eps^{-2}\abs{\mathrm{e}^{is}u-v}^2=\tfrac12\eps^{-2}\abs{u-v}^2+\eps^{-2}\operatorname{Re}\big(\bar vu(1-\mathrm{e}^{is})\big)$ give
\[
S_{\mu_*,M^2+\delta m^2}(A,\varphi)=\tfrac12\ip{A}{(X+\mu_*^2)A}_\eps+\tfrac12\ip{\varphi}{(-\Delta_\eps+M^2)\varphi}_\eps+V(A,\varphi),
\]
with $S_{\mu_0,a}$ as in \eqref{eq:latticeaction}, $V:=V^e+\lambda V^4+V^{\mathfrak d}$, and
\[
V^e:=\sum_{x,\mu}\eps\operatorname{Re}\big(\bar\varphi(x)\varphi(x+\eps e_\mu)(1-\mathrm{e}^{ie\eps A_{x,\mu}})\big),\quad
V^4:=\sum_x\eps^3\abs{\varphi(x)}^4,\quad
V^{\mathfrak d}:=\tfrac12\ip{\varphi}{\mathsf M_{\delta m^2}\varphi}_\eps .
\]
$V^e$ is entire in $e$, and $V^e=\sum_{k\ge1}e^kV^e_k$ with the polynomials $V^e_k:=-\sum_{x,\mu}\eps\operatorname{Re}\big(\bar\varphi(x)\varphi(x+\eps e_\mu)(i\eps A_{x,\mu})^k/k!\big)$. So $V=\sum_{\gamma\ne(0,0)}e^\alpha\lambda^\beta V_\gamma$ with polynomials $V_\gamma$, and $V_{(0,0)}=0$. Let $\gamma_M$ be the centred Gaussian measure on $\R^{\mathrm{bonds}}\times\mathcal H$ under which $A$ and $\varphi$ are independent, with covariances $C^A$ and $C_M$ for the pairings $\ip\cdot\cdot_\eps$. For a polynomial $F$ in $(A,\varphi)$, the \emph{formal expectation} $\langle F\rangle^{\mathrm f}$ is the formal power series in $(e,\lambda)$ obtained from $\sum_{k\ge0}\frac1{k!}\E_{\gamma_M}[F\,(-V)^k]$ by expanding $V$ in monomials. Its coefficient of weight $n$ is a finite sum of Gaussian moments of polynomials: every $V_\gamma$ has weight $\ge1$, so only $k\le n$ contribute. These coefficients are those of the perturbation expansions that define the counterterms of \cite[(3.32)--(3.35), p.~659]{K1} and \cite[(1.21)--(1.24), pp.~416--417]{Ba3}. Wick's theorem writes them as sums of lattice Feynman graphs whose lines are the kernels of $C_M$ and $C^A$. On the periodic torus these kernels are King's full propagators $C_s$, $C_v$ \cite[p.~658]{K1} with photon mass $\mu_*$, and Ba{\l}aban's $C^\eps_0$, $C^\eps$ \cite[p.~417]{Ba3}. The block-averaged operators $G_k$, $G_k(0)$ of \cite{K1} enter only through the splitting $C_v=G_k+F_{k,v}$, $C_s=G_k(0)+F_{k,s}$ \cite[p.~658]{K1}, which leaves the sum of the graphs unchanged.

The \emph{two-point function} $G=G[\mathfrak d]$ is the formal power series, with coefficients in the linear operators on $\mathcal H$, given by $\ip{u}{Gv}_\eps:=\langle\ip u\varphi_\eps\ip v\varphi_\eps\rangle^{\mathrm f}/\langle1\rangle^{\mathrm f}$. Here $\langle1\rangle^{\mathrm f}$ has constant term $1$, and $G_{(0,0)}=C_M$. Since $C_M$ is invertible, $G$ is invertible in this ring of formal power series. We put
\[
\Sigma=\Sigma[\mathfrak d]:=C_M^{-1}-G[\mathfrak d]^{-1},\qquad\text{equivalently}\qquad G=\sum_{n\ge0}C_M\big(\Sigma\,C_M\big)^n .
\]
This is the full amputated one-particle-irreducible two-point function. The second identity is the form \cite[(1.21), p.~416]{Ba3} of the expansion of the two-point function, where $\Sigma=-\delta m^2+\Sigma^\eps+\partial^{\eps*}\Sigma^\eps_1+\Sigma^{\eps*}_1\partial^\eps+\partial^{\eps*}\Sigma^\eps_2\partial^\eps$. In the notation of \cite[p.~659]{K1}, $\Sigma=-\delta m^2+\Gamma^\eps$. The counterterm condition of \cite[(3.33), p.~659]{K1} is
\[
\textup{(K)}\qquad \hat\Sigma[\mathfrak d]_\gamma(x)=0\qquad\text{for all }\gamma\in I,\ x\in\Lambda .
\]

\begin{proposition}[the torus counterterm is one number]\label{app:G1:P1-ct}
For every $M>0$ there is exactly one family $\mathfrak d=(\delta m^2_\gamma)_{\gamma\in I}$ of real functions on $\Lambda$ that satisfies \textup{(K)}. Each $\delta m^2_\gamma$ is constant on $\Lambda$. Its value depends only on $(\eps,j,M)$ and on the reference mass $\mu_*$, and it vanishes when $\alpha$ is odd (Lemma~\ref{app:G1:oddorders}). Consequently the counterterm $\delta m^2_{K,T_j}(\eps;M)=\sum_{\gamma\in I}e^\alpha\lambda^\beta\delta m^2_\gamma$ of Definition~\ref{def:counterterm} and $a_{K,T_j}(M)=M^2+\delta m^2_{K,T_j}(\eps;M)$ are real numbers, not functions of the site.
\end{proposition}
\begin{proof}
\emph{Step 1: the top-order counterterm enters linearly and alone.} Let $\gamma=(\alpha,\beta)$ with $n:=w(\gamma)\ge1$. Every coefficient of weight $n$ of $\langle F\rangle^{\mathrm f}$ is a sum of terms $\frac{(-1)^k}{k!}\E_{\gamma_M}[F\,V_{\gamma_1}\cdots V_{\gamma_k}]$ with $\gamma_1+\dots+\gamma_k=\gamma$ and $w(\gamma_i)\ge1$. The function $\delta m^2_\gamma$ occurs only in $V_\gamma$, and at weight $n$ the factor $V_\gamma$ occurs only in the term with $k=1$. Every other term contains only factors $V_{\gamma'}$ with $w(\gamma')\le n-1$. For a symmetric operator $B$ on $\mathcal H$, Wick's theorem gives
\[
\E_{\gamma_M}\big[\ip u\varphi_\eps\ip v\varphi_\eps\,\tfrac12\ip\varphi{B\varphi}_\eps\big]-\ip u{C_Mv}_\eps\,\E_{\gamma_M}\big[\tfrac12\ip\varphi{B\varphi}_\eps\big]=\ip{u}{C_MBC_Mv}_\eps .
\]
Write $G=N/Z$ with $N_{(0,0)}=C_M$ and $Z_{(0,0)}=1$, and apply this identity with $B=\mathsf M_{\delta m^2_\gamma}$. This gives $G_\gamma=-C_M\mathsf M_{\delta m^2_\gamma}C_M+R_\gamma$, where $R_\gamma$ is a polynomial in the coefficients of $V$ of weight $\le n-1$. In particular $R_\gamma$ depends on $\mathfrak d$ only through the $\delta m^2_{\gamma'}$ with $w(\gamma')<n$. Next, $G^{-1}=\sum_{k\ge0}C_M^{-1}\big(-(G-C_M)C_M^{-1}\big)^k$. Its coefficient at $\gamma$ is $-C_M^{-1}G_\gamma C_M^{-1}$ plus products $G_{\gamma_1}\cdots G_{\gamma_k}$, with factors $C_M^{-1}$ in between, where $k\ge2$, $\sum\gamma_i=\gamma$ and $1\le w(\gamma_i)\le n-1$. Hence
\begin{equation}\label{eq:app:G1:P1-top}
\Sigma[\mathfrak d]_\gamma=-\mathsf M_{\delta m^2_\gamma}+\Gamma_\gamma ,
\end{equation}
where $\Gamma_\gamma$ depends only on the $\delta m^2_{\gamma'}$ with $w(\gamma')<n$. Let $\mathfrak d^{<n}$ be the family obtained from $\mathfrak d$ by setting $\delta m^2_{\gamma'}=0$ for $w(\gamma')\ge n$. Applied to $\mathfrak d^{<n}$, \eqref{eq:app:G1:P1-top} has the same $\Gamma_\gamma$ and a vanishing first term. Therefore $\Gamma_\gamma=\Sigma[\mathfrak d^{<n}]_\gamma$.

\emph{Step 2: symmetries.} Let $g$ be one of the linear maps $(A,\varphi)\mapsto(A,U(\theta)\varphi)$ or $(A,\varphi)\mapsto(-A,\mathsf K\varphi)$, or, if every $\delta m^2_{\gamma'}$ is constant on $\Lambda$, the map $(A,\varphi)\mapsto(\mathsf T_aA,\mathsf T_a\varphi)$. Each such $g$ is orthogonal for $\ip\cdot\cdot_\eps$. By Lemma~\ref{app:G1:P1-lines} it commutes with the covariance of $\gamma_M$ (for $A\mapsto-A$ this is trivial), so $g$ preserves $\gamma_M$. Moreover $V_{\gamma'}\circ g=V_{\gamma'}$ for every $\gamma'$, for the following reasons. $V^4$ and $V^{\mathfrak d}$ depend on $\varphi$ only through $\abs{\varphi(x)}^2$, and $\delta m^2$ is translation invariant when it is constant. $\bar\varphi(x)\varphi(y)$ is invariant under $U(\theta)$. Under $(A,\varphi)\mapsto(-A,\bar\varphi)$ the summand of $V^e$ becomes $\operatorname{Re}\big(\overline{\bar\varphi(x)\varphi(y)(1-\mathrm{e}^{ie\eps A})}\big)$ for every real $e$, hence each $V^e_k$ is invariant. A translation permutes the summands of $V^e$, $V^4$ and $V^{\mathfrak d}$. Hence $\langle F\circ g\rangle^{\mathrm f}=\langle F\rangle^{\mathrm f}$, coefficient by coefficient. With $F=\ip u\varphi_\eps\ip v\varphi_\eps$ this says that every coefficient of $G$ commutes with the action of $g$ on $\mathcal H$ (by $U(\theta)$, $\mathsf K$ or $\mathsf T_a$). The same then holds for $G^{-1}$. Since $C_M$ commutes too, it holds for every coefficient of $\Sigma$.

\emph{Step 3: zero-momentum values.} Let $Y$ be a linear operator on $\mathcal H$ that commutes with all $U(\theta)$ and with $\mathsf K$. Since $U(\theta)(\mathbf 1\otimes v)=\mathbf 1\otimes U(\theta)v$, we get $\hat Y(x)U(\theta)=U(\theta)\hat Y(x)$, and likewise $\hat Y(x)\mathsf K=\mathsf K\hat Y(x)$. A real $2\times2$ matrix that commutes with $J$ has the form $s\mathbf 1+tJ$, and $\mathsf KJ\mathsf K=-J$ forces $t=0$. So $\hat Y(x)=s(x)\mathbf 1$ with $s(x)\in\R$. Suppose in addition that $Y$ commutes with every $\mathsf T_a$. Then $\mathsf T_a(\mathbf 1\otimes v)=\mathbf 1\otimes v$ gives $Y(\mathbf 1\otimes v)=\mathsf T_aY(\mathbf 1\otimes v)$, so $s(x-a)=s(x)$ for all $a\in\eps\Z^3$. Since $\Lambda$ is a single orbit of the translations, $s$ is constant.

\emph{Step 4: induction on the weight.} For $n\in\{2,3,4,5\}$ let $(\mathrm H_n)$ be the statement: there is exactly one choice of the $\delta m^2_{\gamma'}$, $\gamma'\in I$, $w(\gamma')<n$, for which \textup{(K)} holds at all these $\gamma'$, and these functions are constant. $(\mathrm H_2)$ holds vacuously. Assume $(\mathrm H_n)$ with $n\le4$, and let $\gamma\in I$ with $w(\gamma)=n$. By \eqref{eq:app:G1:P1-top}, $\hat\Sigma[\mathfrak d]_\gamma(x)=-\delta m^2_\gamma(x)\mathbf 1+\hat\Gamma_\gamma(x)$, and $\Gamma_\gamma=\Sigma[\mathfrak d^{<n}]_\gamma$ is fixed by $(\mathrm H_n)$. By $(\mathrm H_n)$ the family $\mathfrak d^{<n}$ is constant on $\Lambda$. So by Step~2 every coefficient of $\Sigma[\mathfrak d^{<n}]$ commutes with $U(\theta)$, $\mathsf K$ and all $\mathsf T_a$. By Step~3, $\hat\Gamma_\gamma(x)=s_\gamma\mathbf 1$ with a real $s_\gamma$ independent of $x$. Hence \textup{(K)} holds at $\gamma$ if and only if $\delta m^2_\gamma(x)=s_\gamma$ for all $x$. The conditions at the different $\gamma$ of weight $n$ do not involve one another, because $\Gamma_\gamma$ depends only on weights $<n$. This proves $(\mathrm H_{n+1})$, and $(\mathrm H_5)$ is the claim of existence, uniqueness and constancy. The values are real because all the Gaussian moments involved are real. They depend on $(\eps,j,M,\mu_*)$ only through $\Lambda$, $C_M$ and $C^A$, and they do not involve the photon mass $\mu_0$ of any state. The vanishing for odd $\alpha$ is Lemma~\ref{app:G1:oddorders}.
\end{proof}

\begin{corollary}\label{app:G1:P1-Ba}
Let $\mathfrak d$ be the family of Proposition~\ref{app:G1:P1-ct}. Write the coefficient $\Sigma[\mathfrak d]_\gamma$ as $-\delta m^2_\gamma+\Sigma^\eps_\gamma+\partial^{\eps*}\Sigma^\eps_{1,\gamma}+\Sigma^{\eps*}_{1,\gamma}\partial^\eps+\partial^{\eps*}\Sigma^\eps_{2,\gamma}\partial^\eps$, grouping the graphs according to the derivatives at their external vertices as in \cite[(1.21), p.~416]{Ba3}. Then $\hat\Sigma[\mathfrak d]_\gamma=-\delta m^2_\gamma\mathbf 1+\hat\Sigma^\eps_\gamma$. Hence $\mathfrak d$ is also the unique solution of the equations $-\delta m^2_\gamma+\sum_{x}\eps^3\Sigma^\eps_\gamma(x)=0$, $\gamma\in I$, of \cite[p.~417]{Ba3}.
\end{corollary}
\begin{proof}
All $\delta m^2_{\gamma'}$ are constant. So a translation maps each of the three groups of graphs onto itself, and $\Sigma^\eps_{1,\gamma}$, $\Sigma^\eps_{2,\gamma}$ commute with all $\mathsf T_a$, as in Step~2. The operator $\partial^\eps$ annihilates constants, so the last two terms vanish on $\mathbf 1\otimes v$. The function $\Sigma^\eps_{1,\gamma}(\mathbf 1\otimes v)$ is translation invariant, hence constant in each direction, and $\partial^{\eps*}$ annihilates such bond functions on the torus. Uniqueness follows by repeating Step~4 with this condition in place of \textup{(K)}.
\end{proof}

\begin{remark}[translation-invariant lines]\label{app:G1:P1-rem}
Steps 2--4 use only one property of the lines: they commute with the lattice translations. On the periodic torus this holds for King's full propagators. So the counterterm of Definition~\ref{def:counterterm} is King's counterterm \cite[(3.32)--(3.33)]{K1} on $\Omega=T_{j,\eps}$, with the photon mass taken to be $\mu_*$. The notation $\delta m^2(x)$ of \cite{K1} allows for domains $\Omega\ne T_\eps$, on which King's operators carry Neumann conditions and are not translation invariant. If the lines are only invariant under a block lattice, the argument gives only invariance under that lattice. Ba{\l}aban's splitting $\delta m^2=\delta m^2_{\mathrm{fin}}(x)+\delta m^2_K(x)$ \cite[(1.26), p.~418]{Ba3} has $x$-dependent terms because $G^\eps_K$ and $G^\eps_K(0)$ are not translation invariant; their sum is the constant counterterm of \cite[(1.23), p.~417]{Ba3}.
\end{remark}

\begin{lemma}[property (P1); the per-torus coefficient may be used throughout]\label{app:G1:P1}
Let $K\ge K_0$ and $j\ge1$.
\begin{itemize}
\item[(a)] $W_\eps$ and $z_\eps$ are real numbers that depend on $\eps$ only (for fixed $e,\lambda,m,\mu_*$). $b_{T_j}(\eps)=a_{K,T_j}(m)$ is a real number that depends on $(\eps,j)$ only. None of the three depends on the site, on a direction or on $\mu_0$, and the same three numbers are used at every $\mu_0\in(0,\mu_*]$ and in the state $\fl$. This is property (P1) of Subsection~\ref{ssec:C-standing}.
\item[(b)] Every result of Sections~\ref{sec:chain} and \ref{sec:C-lattice}--\ref{sec:C-OS} that involves a bare coefficient, other than Proposition~\ref{prop:C-ident}(c), is of one of three kinds:
  \begin{itemize}
  \item[(i)] a statement at fixed $(\eps,T_j)$ about models on $T_{j,\eps}$ whose quadratic coefficient is $a+k$, where $a$ is a real constant (for instance $b_{T_j}(\eps)$, a value $a_{K,T_j}(M)$, or a point between two such values) and $k$ is a real function on the sites coming from a source;
  \item[(ii)] a bound whose constants are infima or suprema, over $(\eps,T_j)$ and over $\mu_0$, of quantities of kind~(i), combined with the numbers $W_\eps$ and $z_\eps$;
  \item[(iii)] a compactness, holomorphy, positivity or closedness argument for a family of functionals that obey bounds of kind~(ii).
  \end{itemize}
  No such result compares the coefficients of two different tori or uses a model on $\eps\Z^3$; the counterterm on $\eps\Z^3$ enters only through the number $z_\eps$. Consequently the coefficient $b_{T_j}(\eps)$ of Definition~\ref{def:counterterm} can be used throughout.
\end{itemize}
\end{lemma}
\begin{proof}
\emph{(a), $W_\eps$.} By its definition in Subsection~\ref{ssec:observables}, $W_\eps=2\int_{[-\pi/\eps,\pi/\eps)^3}(2\pi)^{-3}(\hat p^2+m^2)^{-1}d^3p$. The integrand lies in $(0,m^{-2}]$ and the domain is bounded, so $W_\eps$ is a finite positive number. It is defined on $\eps\Z^3$ and involves no site, no direction, no torus and no photon mass.

\emph{(a), $z_\eps$.} $z_\eps=e^2X_\eps(m)$, where $X_\eps(m)$ is the $e^2$-coefficient of $\langle\abs{\varphi(0)}^2\rangle$ for the model on $\eps\Z^3$ at scalar mass $m$ and photon mass $\mu_*$ (Subsection~\ref{ssec:observables}). It is one number, defined without reference to a torus, to $\mu_0$, or to any site other than the origin.

\emph{(a), $b_T$.} By Proposition~\ref{app:G1:P1-ct} with $M=m$, $b_{T_j}(\eps)=a_{K,T_j}(m)$ is one real number per $(\eps,j)$. It multiplies the scalar $\abs{\varphi(x)}^2$ in \eqref{eq:C-action} and carries no direction. Its definition uses the photon propagator $(X+\mu_*^2)^{-1}$ at the reference mass $\mu_*$. The photon mass $\mu_0$ of the state does not enter it, so the same number is the coefficient at every $\mu_0\in(0,\mu_*]$ and in $\fl$. Remark~\ref{rem:counterterm}(ii) explains why the reference mass is needed.

\emph{(b).} Kind (i):
\begin{itemize}
\item Proposition~\ref{prop:UV-B1e}, by Proposition~\ref{app:G1:P1-ct}, Corollary~\ref{app:G1:P1-Ba} and Remark~\ref{app:G1:P1-rem}.
\item Lemmas~\ref{lem:C-current}, \ref{lem:C-Fsector}, \ref{lem:C-cov}, \ref{lem:C-flatlimit}, \ref{lem:C-FP}, \ref{lem:C-RPP}, \ref{lem:C-sym}, \ref{lem:C-OS3lat} and Proposition~\ref{prop:C-RP}, which hold at fixed lattice for every real constant coefficient.
\item The convexity of $a\mapsto\log Z^F_T(a)$ on $[a_{K,T}(M_1),b_T]$ in the proof of Proposition~\ref{prop:ch-Pplus}.
\item The monotonicity in the photon covariance at $b_T$ and the chessboard estimate in the proof of Proposition~\ref{prop:ch-T7}.
\item The convexity in $u$ in the proof of Proposition~\ref{prop:ch-T5}.
\item Steps (1)--(4) of the proof of Proposition~\ref{prop:ch-Tflat}: the translation averaging, the coefficients $b_T+2th$, and the convexity of $f^B$ on $[a_{K,T}(m_1),a_{K,T}(m_2)]$.
\item The fixed-lattice monotonicity in $\mu_0$ at $b_T$ behind Proposition~\ref{prop:C-ident}(a),(b).
\end{itemize}
Kind (ii):
\begin{itemize}
\item Theorem~\ref{thm:T6} and Propositions~\ref{prop:ch-K0}, \ref{prop:ch-CN}.
\item Theorem~\ref{thm:ch-W}, which bounds differences $a_{K,T}(M)-a_{K,T}(M')$ on one torus.
\item The constants $g,\sigma_+,\bar\Gamma,\kappa_+$ of Proposition~\ref{prop:ch-Pplus}, the constants $E_2$ and $\bar\Gamma^{(1)}_{21}$ of Proposition~\ref{prop:ch-T5}, and the constants $J_{\rm IR},\bar Z,g_{21},S^B_\flat$ of Proposition~\ref{prop:ch-Tflat}.
\item Theorems~\ref{thm:ch-T9}, \ref{thm:C-T9} and Proposition~\ref{prop:C-bounds}.
\end{itemize}
Kind (iii):
\begin{itemize}
\item Propositions~\ref{prop:C-compact}, \ref{prop:C-tight}, \ref{prop:C-order}, \ref{prop:C-mu0} and Lemma~\ref{lem:C-diag}, applied to sequences $(\eps_i,T_{j_i},\mu_i)$ in which each model has its own coefficient.
\item Theorems~\ref{thm:C-OS0}--\ref{thm:C-OS3} and Proposition~\ref{prop:C-OS2dens}.
\item The passage to the limit in Proposition~\ref{prop:C-ident}(a),(b).
\end{itemize}
None of these results uses a convergence of $b_{T_{j_i}}(\eps_i)$. The number $z_\eps$ enters the observables and the c-number differences only as a constant. Replacing it by another sequence multiplies the generating functionals by nowhere-vanishing c-number factors, so no model on $\eps\Z^3$ is involved. The same $b_T$, and one $z_\eps$ and one $W_\eps$, are used in the state, in the observables and in both factors (Lemma~\ref{app:G1:P3}).
\end{proof}

\begin{remark}
Part (a) is used in the following places:
\begin{itemize}
\item Proposition~\ref{prop:ch-T7}: the chessboard step needs a site-independent coefficient, and the monotonicity in $\mu_0$ needs one that does not depend on $\mu_0$.
\item Step (1) of Proposition~\ref{prop:ch-Tflat}: the translation averaging needs the torus model to be invariant under $\eps\Z^3$.
\item Lemma~\ref{lem:C-sym}.
\item Lemma~\ref{lem:C-flatlimit} and Proposition~\ref{prop:C-mu0}: limits and monotonicity in $\mu_0$ at a fixed coefficient.
\end{itemize}
\end{remark}

\begin{remark}
Lemma~\ref{app:G1:P1}(a) is what Proposition~\ref{prop:ch-T7} (the chessboard step needs a site-independent coefficient, and the $\mu_0$-monotonicity needs a $\mu_0$-free one), Lemma~\ref{lem:C-sym} (translations and $\Oh$ map the model to itself) and Lemma~\ref{lem:C-flatlimit} (the limit $\mu_0\downarrow0$ at a fixed coefficient) use. Part (b) is used by Sections~\ref{sec:chain}--\ref{sec:C-OS} as a whole.
\end{remark}

\paragraph{Property (P2): integrability.}

\paragraph{Property (P2): integrability.}
Throughout this paragraph $\lambda>0$ and $(\eps,T)$ is fixed. $T_\eps$ is the lattice torus with $N\ge2$ sites per side, so $\abs{T_\eps}=N^3$. Bonds are $b=(x,\mu)$ with $b_-=x$ and $b_+=x+\eps e_\mu$. For $a\in\R$ let $S_m^{(a)}$ be the matter action $S_m$ of \eqref{eq:C-action} with $b_T$ replaced by $a$, so that $S_m=S_m^{(b_T)}$, and let $S^{(a)}_{\eps,T,\mu_0}$ be the corresponding full action. Write $\varphi_x=r_x(\cos\vartheta_x,\sin\vartheta_x)$ and $u_b:=\vartheta_{b_+}-\vartheta_{b_-}+\eps eA_b$. For a real function $k$ on the sites put
\[
\kappa_k:=\exp\Big(-\sum_x\eps^3k(x)(W_\eps+z_\eps)\Big),\qquad
W^{(a)}_k(A):=\int\prod_xd^2\varphi_x\,\mathrm{e}^{-S^{(a)}_m(A,\varphi)+\Xr(k)}\in(0,\infty],
\]
and $V_a(r):=\eps^3\big(\tfrac12ar^2+\lambda r^4\big)+3\eps r^2$. Thus $W_k=W^{(b_T)}_k$ is the weight of Lemma~\ref{lem:C-current}.

\begin{lemma}[(P2)]\label{app:G1:P2}
Let $\lambda>0$.
\begin{itemize}
\item[(a)] For $c>0$, $a\in\R$ and $p\ge0$, $\int_0^\infty r^p\mathrm{e}^{-cr^4+ar^2}\,dr<\infty$. Moreover
\[
\int_{\R^2}\mathrm{e}^{-c\abs\varphi^4+a\abs\varphi^2}\,d^2\varphi=\pi\int_0^\infty \mathrm{e}^{-cu^2+au}\,du\le\pi^{3/2}c^{-1/2}\mathrm{e}^{a^2/(4c)} .
\]
\item[(b)] For every $a\in\R$, every real $k$ and all $(A,\varphi)$,
\begin{equation}\label{eq:G1P2-dom}
0<\mathrm{e}^{-S^{(a)}_m(A,\varphi)+\Xr(k)}\le\kappa_k\prod_x\exp\big(-\eps^3\lambda r_x^4+\eps^3(k(x)-\tfrac a2)r_x^2\big).
\end{equation}
Moreover $\mathrm{e}^{-S^{(a)}_m(A,\varphi)}=\prod_x\mathrm{e}^{-V_a(r_x)}\prod_b\mathrm{e}^{\eps r_{b_-}r_{b_+}\cos u_b}$, and
\[
\sum_xV_a(r_x)-\sum_b\eps r_{b_-}r_{b_+}=\sum_x\eps^3\big(\tfrac12ar_x^2+\lambda r_x^4\big)+\tfrac12\eps\sum_b(r_{b_+}-r_{b_-})^2 .
\]
In particular, for all couplings $c_b\in[0,\eps]$,
\[
\prod_x\mathrm{e}^{-V_a(r_x)}\prod_b\mathrm{e}^{c_br_{b_-}r_{b_+}}\le\prod_x\mathrm{e}^{-\eps^3(\frac12ar_x^2+\lambda r_x^4)} .
\]
\item[(c)] $W^{(a)}_k$ is a Borel function of $A$, and for every $A$
\[
0<W^{(a)}_k(A)\le W^{(a)}_k(0)\le\kappa_k\prod_x\pi^{3/2}(\eps^3\lambda)^{-1/2}\exp\Big(\frac{\eps^3(k(x)-a/2)^2}{4\lambda}\Big)<\infty .
\]
\item[(d)] Let $\mu_0>0$, let $f$ be a real $1$-form and $k$ a real function on the sites. Then
\[
0<\int\prod_bdA_b\prod_xd^2\varphi_x\,\mathrm{e}^{-S^{(a)}_{\eps,T,\mu_0}(A,\varphi)+\ip fA_\eps+\Xr(k)}<\infty .
\]
In particular $\mathrm{e}^{-S_{\eps,T,\mu_0}}\,dA\,d\varphi$ is a finite nonzero measure, so \eqref{eq:C-state} is the expectation of a probability measure. For real $2$-forms $\bg_i$, real functions $h_l$ and $z_i,w_l\in\C$, the observable $\Theta=\exp[\sum_iz_iF(\bg_i)+\sum_lw_l\Xr(h_l)]$ satisfies $\ev{\abs\Theta}_{\eps,T,\mu_0}<\infty$. Every polynomial in the $F(\bg_i)$ and $\Xr(h_l)$ is integrable as well.
\item[(e)] Let $\gamma^\flat:=\mu_{(-\Delta_\eps)^{-1}P^\perp}\otimes P^{-3}dc|_{[0,P)^3}$ be the measure of Definition~\ref{def:C-flat}. Then
\[
0<\int d\gamma^\flat(A^\perp,c)\,W_0(A^\perp+c)\le W_0(0)<\infty ,
\]
so the denominator of the $\flat$ state is positive and finite. For real $\bg$ and $k$,
\[
\int d\gamma^\flat(A^\perp,c)\,\mathrm{e}^{F(\bg)(A^\perp)}W_k(A^\perp+c)\le W_k(0)\exp\big(\tfrac12\ip{d^*\bg}{(-\Delta_\eps)^{-1}P^\perp d^*\bg}_\eps\big)<\infty .
\]
Hence the observables $\Theta$ of (d) satisfy $\ev{\abs\Theta}_{\eps,T,\flat}<\infty$.
\item[(f)] For every $a\in\R$ the measure $d\nu(r):=r\,\mathrm{e}^{-V_a(r)}\,dr$ on $[0,\infty)$ is nonzero, and $\int \mathrm{e}^{sr^2}d\nu(r)<\infty$ for every $s\in\R$.
\end{itemize}
Property (P2) of Subsection~\ref{ssec:C-standing} is part (a) with $c=\eps^3\lambda$; the integrability holds with respect to $r^p\,dr$ for any $p\ge0$, and with respect to $d^2\varphi$.
\end{lemma}

\begin{proof}
(a) By the AM--GM inequality, $ar^2\le\tfrac c2r^4+\tfrac{a^2}{2c}$. Hence $r^p\mathrm{e}^{-cr^4+ar^2}\le \mathrm{e}^{a^2/(2c)}r^p\mathrm{e}^{-cr^4/2}$, which is integrable on $[0,\infty)$. In polar coordinates $d^2\varphi=r\,dr\,d\vartheta$, and the substitution $u=r^2$, $r\,dr=\tfrac12du$, gives the identity. Next, $-cu^2+au=-c\big(u-\tfrac a{2c}\big)^2+\tfrac{a^2}{4c}$. Extending the $u$-integral to $\R$ and using $\int_\R \mathrm{e}^{-cv^2}dv=(\pi/c)^{1/2}$ gives the bound.

(b) $U(s)$ is a rotation, so $\abs{U(\eps eA_b)\varphi_{b_+}}=r_{b_+}$ and $\big(U(\eps eA_b)\varphi_{b_+}\big)\cdot\varphi_{b_-}=r_{b_+}r_{b_-}\cos u_b$. Therefore
\[
\tfrac12\eps^3\abs{(D_A\varphi)(b)}^2=\tfrac12\eps\big(r_{b_+}^2+r_{b_-}^2\big)-\eps r_{b_-}r_{b_+}\cos u_b\ \ge0 .
\]
Each site $x$ is the initial point of exactly three bonds, $(x,\mu)$, and the final point of exactly three bonds, $(x-\eps e_\mu,\mu)$. Hence $\sum_b\tfrac12\eps(r_{b_+}^2+r_{b_-}^2)=3\eps\sum_xr_x^2$, which gives the product formula for $\mathrm{e}^{-S^{(a)}_m}$. Setting $\cos u_b=1$ in the same computation gives $3\eps\sum_xr_x^2-\sum_b\eps r_{b_-}r_{b_+}=\tfrac12\eps\sum_b(r_{b_+}-r_{b_-})^2\ge0$. For $0\le c_b\le\eps$ and $r\ge0$ we have $c_br_{b_-}r_{b_+}\le\eps r_{b_-}r_{b_+}$, which gives the last inequality of (b). For \eqref{eq:G1P2-dom}, drop the nonnegative gradient term from $S^{(a)}_m$ and write $\Xr(k)=\sum_x\eps^3k(x)r_x^2+\log\kappa_k$. Positivity is clear.

(c) The integrand $\mathrm{e}^{-S^{(a)}_m+\Xr(k)}$ is continuous and positive on $\R^{\text{bonds}}\times(\R^2)^{T_\eps}$. By Tonelli's theorem, $W^{(a)}_k$ is Borel and $W^{(a)}_k(A)>0$. Integrate \eqref{eq:G1P2-dom} over $\varphi$ and use (a) at each site, with $c=\eps^3\lambda$ and with $\eps^3(k(x)-a/2)$ in place of $a$. This bounds $W^{(a)}_k(A)$ by the stated product for every $A$, in particular for $A=0$.

It remains to show $W^{(a)}_k(A)\le W^{(a)}_k(0)$. For $t\ge0$ and $u\in\R$,
\[
\mathrm{e}^{t\cos u}=\sum_{q\ge0}\frac{(t/2)^q}{q!}\big(\mathrm{e}^{iu}+\mathrm{e}^{-iu}\big)^q=\sum_{\nu\in\Z}I_\nu(t)\mathrm{e}^{i\nu u},
\qquad
I_\nu(t):=\sum_{\substack{0\le j\le q,\ 2j-q=\nu}}\frac{(t/2)^q}{j!\,(q-j)!}\ \ge0 .
\]
All terms are nonnegative, so the rearrangement is legitimate, and $\sum_\nu I_\nu(t)=\mathrm{e}^t$. Fix $r\in[0,\infty)^{T_\eps}$ and put $t_b:=\eps r_{b_-}r_{b_+}\ge0$. Then $\prod_b\mathrm{e}^{t_b\cos u_b}=\sum_{n\in\Z^{\text{bonds}}}\prod_bI_{n_b}(t_b)\,\mathrm{e}^{i\sum_bn_bu_b}$. The series converges absolutely, uniformly in $\vartheta$, with $\sum_n\prod_bI_{n_b}(t_b)=\prod_b\mathrm{e}^{t_b}$. Integrate term by term over $\vartheta\in[0,2\pi)^{T_\eps}$. Since $\int\prod_xd\vartheta_x\,\mathrm{e}^{i\sum_bn_b(\vartheta_{b_+}-\vartheta_{b_-})}=(2\pi)^{\abs{T_\eps}}\mathbf 1\{\partial n=0\}$, with $(\partial n)(x):=\sum_{b_+=x}n_b-\sum_{b_-=x}n_b$, we get
\[
J_r(A):=\int\prod_xd\vartheta_x\prod_b\mathrm{e}^{t_b\cos u_b}=(2\pi)^{\abs{T_\eps}}\sum_{n:\,\partial n=0}\prod_bI_{n_b}(t_b)\,\mathrm{e}^{ie\eps\sum_bn_bA_b},
\]
and therefore $0<J_r(A)=\abs{J_r(A)}\le J_r(0)$. By (b) and Tonelli's theorem,
\[
W^{(a)}_k(A)=\kappa_k\int\prod_xr_x\,dr_x\,\mathrm{e}^{-V_a(r_x)+\eps^3k(x)r_x^2}\,J_r(A) .
\]
Integrating $J_r(A)\le J_r(0)$ against this positive measure gives $W^{(a)}_k(A)\le W^{(a)}_k(0)$.

(d) The integrand is continuous and positive, so the integral is positive. By (b), the integral is at most
\[
\int\prod_bdA_b\,\mathrm{e}^{-\frac12\ip A{(-\Delta_\eps+\mu_0^2)A}_\eps+\ip fA_\eps}\ \times\ \kappa_k\int\prod_xd^2\varphi_x\,\mathrm{e}^{-\eps^3\lambda r_x^4+\eps^3(k(x)-\frac a2)r_x^2} .
\]
The second factor is finite by (a). In the coordinates $(A_b)$, $\ip A{(-\Delta_\eps+\mu_0^2)A}_\eps=A\cdot QA$ with $Q=\eps^3(-\Delta_\eps+\mu_0^2)$, and $\ip fA_\eps=\eps^3f\cdot A$. Here $Q\ge\eps^3\mu_0^2>0$, because $-\Delta_\eps\ge0$ (Lemma~\ref{app:G1:X}(ii)). So the first factor equals $(2\pi)^{n/2}(\det Q)^{-1/2}\exp\big(\tfrac12\eps^6f\cdot Q^{-1}f\big)<\infty$, with $n=3N^3$. Tonelli's theorem gives the claim. With $f=0$, $k=0$ and $a=b_T$ this is the normalisability of the state. For $\Theta$: $F$ and $\Xr$ are real and linear in $\bg$ and $h$. Hence $\abs\Theta=\exp\big[F(\sum_i\operatorname{Re}z_i\,\bg_i)+\Xr(\sum_l\operatorname{Re}w_l\,h_l)\big]$, which is of the form just treated with $f=d^*\sum_i\operatorname{Re}z_i\,\bg_i$, since $F(\bg)=\ip{d^*\bg}{A}_\eps$. Polynomials are covered by $\abs x^q\le q!\,(\mathrm{e}^x+\mathrm{e}^{-x})$.

(e) $W_0$ is Borel, and $0<W_0\le W_0(0)<\infty$ by (c). Moreover $\gamma^\flat$ is a probability measure: the Gaussian factor exists because $-\Delta_\eps\ge16/\ell_T^2$ on the range of $P^\perp$ (Lemma~\ref{app:G1:X}(iii)). Hence the integral of $W_0$ lies in $(0,W_0(0)]$; it is positive because the sets $\{W_0>1/n\}$ increase to the whole space. For the second claim: a constant $1$-form $c$ satisfies $dc=0$, so $F(\bg)(A^\perp+c)=\ip{d^*\bg}{A^\perp}_\eps$. Under the centred Gaussian with covariance $(-\Delta_\eps)^{-1}P^\perp$, this is a centred normal variable with variance $\ip{d^*\bg}{(-\Delta_\eps)^{-1}P^\perp d^*\bg}_\eps$. Bound $W_k(A^\perp+c)\le W_k(0)$ by (c) and compute the Gaussian exponential moment; this gives the bound. The claim for $\Theta$ follows as in (d).

(f) $\int_0^\infty \mathrm{e}^{sr^2}r\,\mathrm{e}^{-V_a(r)}\,dr=\int_0^\infty r\,\mathrm{e}^{-\eps^3\lambda r^4+(s-\eps^3a/2-3\eps)r^2}\,dr$ is finite by (a) with $p=1$. It is positive because the integrand is positive on $(0,\infty)$.
\end{proof}

\begin{remark}
Lemma~\ref{app:G1:P2} supplies every use of (P2) in Sections~\ref{sec:C-lattice}--\ref{sec:C-OS} and in this appendix, and none of these uses needs more:
\begin{itemize}
\item the normalisability of \eqref{eq:C-state} is part (d);
\item the domination that justifies Fubini's theorem in the proof of Lemma~\ref{lem:C-current} is part (b), last display (with $c_b=\eps$), together with part (a);
\item the integrability of $\abs\Theta^2$ in Proposition~\ref{prop:C-RP}(c) is parts (d) and (e);
\item the finiteness of the logarithmic moment generating functions $\alpha_n$ in Section~\ref{sec:OS4} is parts (d) and (e) with $\bg=0$ and constant $k$;
\item the hypotheses on $\nu_x$ in Definition~\ref{app:G1:BFSsetting} and Proposition~\ref{app:G1:BFStorus} are part (f) and part (b), last display.
\end{itemize}
The coefficient $a$ is arbitrary, so the lemma also applies at the anchor coefficients $a_{K,T}(M)$ and in Lemma~\ref{app:G1:gibbs}.
\end{remark}

\paragraph{Property (P3): one scheme.}

\begin{lemma}[property (P3)]\label{app:G1:P3}
On each $T_j$ and at each $\eps$ (with $K\ge K_0$), the following coincide:
\begin{itemize}
\item[(a)] the bare coefficient of the state \eqref{eq:C-state}, at every $\mu_0\in(0,\mu_*]$ and in $\fl$;
\item[(b)] the coefficient at which the bounds for positive sources (Propositions~\ref{prop:ch-Pplus}, \ref{prop:ch-T7}) are proved;
\item[(c)] the coefficient $a_0$ of the factor for negative sources (Proposition~\ref{prop:ch-Tflat}, the left side of the anchor condition \eqref{p:ch-anchor}: $b_T+2H\le a_{K,T}(m_1)$).
\end{itemize}
All three equal $b_T=a_{K,T}(m)$. Every anchor coefficient in Section~\ref{sec:chain} ($a_{K,T}(M_1)$ and $a_{K,T}(m)$ in Proposition~\ref{prop:ch-Pplus}; $a_{K,T}(m_1)$ and $a_{K,T}(m_2)$ in Theorem~\ref{thm:T6}, Propositions~\ref{prop:ch-T5}, \ref{prop:ch-Tflat}, \ref{prop:ch-CN} and the bound \eqref{p:ch-CNu}) is a value of the single function $M\mapsto a_{K,T}(M)$ of Definition~\ref{def:counterterm} on the same torus. One $z_\eps$ and one $W_\eps$, those of Subsection~\ref{ssec:observables}, appear in the observables \eqref{eq:C-obs}, in $\alpha_T$ and $\bar\Gamma$ of Proposition~\ref{prop:ch-Pplus}, in $\rho^B$, $f^B$, $\Delta_1^B$ of Subsection~\ref{ssec:ch-minus}, and in $\bar\Gamma^{(1)}_{21}$. Theorem~\ref{thm:T6} contains no $z_\eps$. This is property (P3) of Subsection~\ref{ssec:C-standing}.
\end{lemma}

\begin{proof}
(1) \emph{One state coefficient.} The state is defined at $b_T$ (Definition~\ref{def:counterterm}, and Lemma~\ref{app:G1:P1}(a) for $\mu_0$ and $\fl$). Proposition~\ref{prop:ch-Pplus} is stated with upper anchor $M_2=m$, i.e.\ at the coefficient $a_{K,T}(m)=b_T$ on the same $T$. Proposition~\ref{prop:ch-T7} is stated for all states at the coefficient $b_{T_j}$. Proposition~\ref{prop:ch-Tflat} is stated at $a_0=b_T$ through \eqref{p:ch-anchor}.
(2) \emph{One anchor function.} Each anchor listed in the statement is, by its definition in Section~\ref{sec:chain}, $a_{K,T}(M)$ at some $M\in\{M_1,m,m_1,m_2\}$ with the photon $\mu_*$ on the same torus. In the slice models of Subsection~\ref{ssec:T6-statement} the counterterm is fixed at the photon covariance $C_{\mu_*}$ for every $u$, so the coefficients of Theorem~\ref{thm:T6} at $u=1$ and $u=-\eta$ are $a_{K,T}(m_2)$ and $a_{K,T}(m_1)$, not a $u$-dependent profile.
(3) \emph{One $z_\eps$, one $W_\eps$.} $z_\eps$ and $W_\eps$ are the numbers of Subsection~\ref{ssec:observables}, the same on every torus (Lemma~\ref{app:G1:P1}(a)). They enter each statement of the chain only through the c-number $-\sum_x\eps^3h(x)(W_\eps+z_\eps)$ in $\Xr(h)$, or through $\tfrac12a(W_\eps+z_\eps)$ in $f_T(a)$ and $\gamma_T(M)$. Each proof on the path is an identity in $z$: the $z$-parts of the two sides of every comparison are equal. The only inputs that depend on $z$ are the bounds on $\gamma$-differences uniform in $\eps$ (Proposition~\ref{prop:ch-CN}, the bound \eqref{p:ch-CNu}, $\bar\Gamma$), and these are proved for the graphical $z_\eps=e^2X_\eps(m)$. Theorem~\ref{thm:T6} is a vacuum bound at anchor coefficients and contains no $z_\eps$.
(4) \emph{Vacuum constants.} In each statement that uses a vacuum constant, the $E_0$ is the photon-normalised $E_{0,s}$ of that statement. This is a condition within each statement. It does not affect Sections~\ref{sec:C-lattice}--\ref{sec:C-OS}, because normalised expectations contain no vacuum constants (Remark~\ref{rem:counterterm}(iii)), and $E_{0,v}(\mu_*)$ cancels in every $\gamma$-difference.
Items (1)--(3) are the statement, and (4) shows that the vacuum constants do not interfere.
\end{proof}

\begin{remark}
Mixed schemes are excluded: if the ``$+$'' input were taken at one coefficient and the ``$-$'' input at another, the Cauchy--Schwarz step (3) of the proof of Theorem~\ref{thm:ch-T9} would combine bounds for two different states. The difference would be the volume-extensive shift $r_T(m)\sum_x\eps^3r_x^2$, not a c-number.
\end{remark}

\paragraph{The zero-mode split.}

\begin{lemma}[spectrum of $-\Delta_\eps$ and the zero mode]\label{app:G1:zeromode-split}
On the real 1-forms of $T_{j,\eps}$, with the inner product $\ip{\cdot}{\cdot}_\eps$:
\begin{itemize}
\item[(a)] $-\Delta_\eps$ is self-adjoint and $\ge0$. Its kernel is the span of $\mathbf 1_1,\mathbf 1_2,\mathbf 1_3$, and $-\Delta_\eps\ge16/\ell_T^2$ on the range of $P^\perp$.
\item[(b)] $P_0$ commutes with $-\Delta_\eps$. For $\mu_0>0$, $C_{\mu_0}=(-\Delta_\eps+\mu_0^2)^{-1}=\mu_0^{-2}P_0+(-\Delta_\eps+\mu_0^2)^{-1}P^\perp$.
\item[(c)] Under $\mu_{C_{\mu_0}}$ write $A=A^\perp+c$ with $A^\perp:=P^\perp A$ and $c:=P_0A=\sum_\mu c_\mu\mathbf 1_\mu$. Then $A^\perp$ and $c$ are independent. $A^\perp$ is centred Gaussian with covariance $(-\Delta_\eps+\mu_0^2)^{-1}P^\perp$. The $c_\mu$ are independent $N(0,s(\mu_0))$ with $s(\mu_0)=1/(\mu_0^2\abs T)$. Equivalently, the mass term restricted to the constants is $\tfrac12\mu_0^2\norm{c}_{2,\eps}^2=\tfrac12\mu_0^2\abs T\sum_\mu c_\mu^2$.
\end{itemize}
\end{lemma}

\begin{proof}
This is Lemmas~\ref{app:G1:X} and \ref{app:G1:zeromode} in the notation of Section~\ref{sec:C-lattice}, where $P_0=\Pi_0$ and $P^\perp=\Pi^\perp$. (a) is Lemma~\ref{app:G1:X}(ii),(iii). (b) $P_0$ is the spectral projection of $-\Delta_\eps$ onto its kernel, so it commutes with $-\Delta_\eps$, and $(-\Delta_\eps+\mu_0^2)^{-1}$ equals $\mu_0^{-2}$ on the range of $P_0$. (c) is Lemma~\ref{app:G1:zeromode}(a),(b); the identity $\ip{c}{c}_\eps=\abs T\sum_\mu c_\mu^2$ is Lemma~\ref{app:G1:zeromode}(a).
\end{proof}

\paragraph{The constants of Theorem~\ref{thm:ch-T9} are uniform over the tori.}

Let $Z_T$ be the zero-mode price of Proposition~\ref{prop:ch-Tflat}. It comes from comparing the Haar zero mode with the massive zero mode of variance $s_K:=1/(\mu_*^2\abs T)$ on the period cell $P=2\pi/(\abs e\ell_T)$:
\begin{gather*}
Z_T:=\frac1{\ell_T^3}\Big[\tfrac32\ln_+\!\Big(\frac{P^2}{s_K}\Big)+3\delta_P\Big]=\frac1{\ell_T^3}\Big[\tfrac32\ln_+\!\Big(\frac{4\pi^2\mu_*^2\ell_T}{e^2}\Big)+3\delta_P\Big],\\
\delta_P:=2\sum_{n\ge1}\mathrm{e}^{-2\pi^2n^2}\ (\approx5.35\cdot10^{-9}),
\end{gather*}
where $\ln_+:=\max(\ln,0)$.
Unlike $\Delta_1^B$ (a supremum over $\eps$ and $\ell_T\ge\ell_0$, Subsection~\ref{ssec:ch-minus}) and $J_{\rm IR}$ (likewise a supremum), $Z_T$ depends on the torus. The constant $S^B_\flat$ of Proposition~\ref{prop:ch-Tflat} is therefore formed with $\bar Z:=\sup_jZ_{T_j}$, and the next lemma shows that $\bar Z<\infty$.

\begin{lemma}[$\sup Z_T<\infty$]\label{app:G1:ZT}
Let $e\ne0$ and $\ell_0>0$, and put $q:=4\pi^2\mu_*^2/e^2$. Then
\[
\sup_{\ell_T\ge\ell_0}Z_T\le\frac1{\ell_0^3}\Big[\tfrac32\ln_+(q\ell_0)+\frac1{2\mathrm{e}}+3\delta_P\Big]<\infty.
\]
\end{lemma}

\begin{proof}
For $x,y>0$, $\ln_+(xy)\le\ln_+x+\ln_+y$: if $xy\le1$ the left side is $0$, and otherwise $\ln(xy)=\ln x+\ln y$. With $t:=\ell_T/\ell_0\ge1$ this gives $\ln_+(q\ell_T)\le\ln_+(q\ell_0)+\ln t$. Hence
\[
Z_T\le\frac{\tfrac32\ln_+(q\ell_0)+3\delta_P}{\ell_0^3t^3}+\frac{3}{2\ell_0^3}\,\frac{\ln t}{t^3}\le\frac1{\ell_0^3}\Big[\tfrac32\ln_+(q\ell_0)+3\delta_P\Big]+\frac{3}{2\ell_0^3}\cdot\frac1{3\mathrm{e}},
\]
because $t^{-3}\le1$ and $\max_{t\ge1}t^{-3}\ln t=1/(3\mathrm{e})$, attained at $t=\mathrm{e}^{1/3}$.
\end{proof}

\begin{corollary}[the constants of Theorem~\ref{thm:ch-T9}]\label{app:G1:T9const}
Put $S^B_\flat(T):=[\Delta_1^B+J_{\rm IR}+Z_T]_+/g_{21}$ for a torus $T$, and
\[
S^B_\flat:=\big[\Delta_1^B+J_{\rm IR}+\sup_{\ell_T\ge\ell_0}Z_T\big]_+\big/g_{21}.
\]
Then $S^B_\flat<\infty$, and $S^B_\flat(T)\le S^B_\flat$ for every torus with $\ell_T\ge\ell_0$, in particular for every $T_j$. Hence Proposition~\ref{prop:ch-Tflat} holds with the constant $S^B_\flat$, uniformly in $\eps$, $j$ and $\mu_0\in(0,\mu_*]\cup\{\fl\}$, and Theorem~\ref{thm:ch-T9} holds with $\mathsf K:=\max(\kappa_+,2S^B_\flat)$, which is independent of $\eps$, $j$ and $\mu_0$.
\end{corollary}

\begin{proof}
The inputs are: $\Delta_1^B<\infty$ (Proposition~\ref{prop:ch-T5}); $J_{\rm IR}<\infty$, a supremum over $\eps$ and $\ell_T\ge\ell_0$ with an explicit bound free of $\eps$ (Lemma~\ref{app:G1:IR}(iv)); $g_{21}>0$ (Theorem~\ref{thm:ch-W}(b)); and Lemma~\ref{app:G1:ZT}. Together they give $S^B_\flat<\infty$. None of the four quantities depends on $\eps$, $j$ or $\mu_0$.
For each $T$ with $\ell_T\ge\ell_0$ we have $Z_T\le\sup_{\ell'\ge\ell_0}Z_{T'}$. Since $x\mapsto[x]_+$ is nondecreasing and $g_{21}>0$, it follows that $S^B_\flat(T)\le S^B_\flat$. The right side of the bound of Proposition~\ref{prop:ch-Tflat} on $T$ is therefore at most $2\norm h_{L^1,\eps}S^B_\flat$. In steps (5)--(6) of the proof of Theorem~\ref{thm:ch-T9} the bound $\kappa_+\norm{k_+}+2S^B_\flat\norm{k_-}\le\max(\kappa_+,2S^B_\flat)\norm k$ follows from $\norm{k_+}_{L^1,\eps}+\norm{k_-}_{L^1,\eps}=\norm k_{L^1,\eps}$. The constant $\Delta_1^B$ is the one of the slice vacuum bound (Proposition~\ref{prop:ch-T5}), formed with the slice covariance $\Gamma^B$.
\end{proof}

\begin{remark}
(i) $S^B_\flat$ is not uniform as $e\to0$: $\sup_{\ell_T\ge\ell_0}Z_T$ grows like $\tfrac32\ln(1/e^2)/\ell_0^3$, because the period cell $2\pi/(\abs e\ell_T)$ becomes infinite. In $\fW$ the charge satisfies $e^2>0$ and is fixed, so this does not matter for Theorem~\ref{thm:main}. (ii) Since $Z_T\to0$ as $\ell_T\to\infty$, the zero-mode price disappears in infinite volume, but uniformity over the sequence $(T_j)$ needs only the supremum.
\end{remark}

\subsection{The continuum limits}\label{app:G1:area-cont}

\subsubsection{Multilinearity, fixed-torus translations and the class $\cD$}\label{app:G1:part-C1}

We prove three steps of Sections~\ref{sec:C-limits}--\ref{sec:C-OS}. Throughout, $C^{(R)}$ is the real vector space of continuous scalar functions or continuous 2-forms on $\R^3$ with support in the open ball $B_R$. For a 2-form $\bg=\sum_{\mu<\kappa}g_{\mu\kappa}\,dx_\mu\wedge dx_\kappa$ we put $\norm{\bg}_\infty:=\max_{\mu<\kappa}\norm{g_{\mu\kappa}}_\infty$. Translations act by $(\tau_af)(x):=f(x-a)$, componentwise on 2-forms. A tuple $f=(\bg_1,\dots,\bg_n;h_1,\dots,h_m)$ is \emph{$R$-admissible for $T$} if all its entries lie in $C^{(R)}$ and $2R+2<\ell_T$. By Definition~\ref{def:C-sampling} this is exactly when the entries can be sampled on $T_\eps$.

\paragraph{Multilinearity and symmetry of the limits.}

\begin{lemma}[multilinearity]\label{app:G1:multilin}
\begin{enumerate}
\item[(a)] Fix $(\eps,T)$ and $\mu_0\in(0,\mu_*]\cup\{\flat\}$, and let $R$ satisfy $2R+2<\ell_T$. On $(C^{(R)})^{n+m}$ (2-forms in the first $n$ arguments, scalar functions in the last $m$), $S^{\eps,T,\mu_0}_{n,m}$ is finite and real. It is linear in each argument separately, and it is invariant under permutations of $(\bg_1,\dots,\bg_n)$ and, separately, of $(h_1,\dots,h_m)$.
\item[(b)] Let $S_{n,m}$ be the limit of Proposition~\ref{prop:C-compact}(a). Then $S_{n,m}$ has the properties of (a) on $(C^{(R)})^{n+m}$ for every $R$; if $T_{j_i}\equiv T$, for every $R$ with $2R+2<\ell_T$. The same holds for $S^0$ of Proposition~\ref{prop:C-mu0}(b).
\end{enumerate}
In particular the limits of Proposition~\ref{prop:C-compact}(a) are multilinear and symmetric within each field type.
\end{lemma}
\begin{proof}
(a) \emph{Linearity of sampling.} Sampling evaluates a function at the sites of the cell $[-\ell_T/2,\ell_T/2)^3$, and a 2-form component at plaquette centres (Definition~\ref{def:C-sampling}). Evaluation at a point is linear, so the sampled function of $\alpha f+\beta f'$ is $\alpha$ times the sampled $f$ plus $\beta$ times the sampled $f'$. Since $C^{(R)}$ is a vector space, all three functions are admissible.
\emph{Linearity of the fields.} By \eqref{eq:C-obs}, $F(\bg)=\ip{\bg}{dA}_\eps$ is linear in the sampled $\bg$. Also $\Xr(h)=\sum_x\eps^3h(x)(r_x^2-W_\eps-z_\eps)$ is linear in $h$, because $W_\eps,z_\eps$ do not depend on $h$ (P1). So for fixed field configuration $(A,\varphi)$ the product $\prod_iF(\bg_i)\prod_k\Xr(h_k)$ is a real number, linear in each argument. It is invariant under permutations of the $\bg_i$ and of the $h_k$, since multiplication of real numbers is commutative.
\emph{Integrability.} The product is a polynomial in the $F(\bg_i)$ and $\Xr(h_k)$, hence an elementary observable. It is therefore integrable in every state, the $\flat$ state included (Lemma~\ref{lem:C-flatlimit}). Expectation is linear on integrable functions and preserves reality, which gives (a). This is the same argument as in the proof of Lemma~\ref{lem:C-OS3lat}.
(b) Each of the properties in (a) is an identity between the values of $S^{(i)}_{n,m}$ at finitely many tuples, all of them admissible. Examples are $S^{(i)}(\dots,\alpha f+\beta f',\dots)=\alpha S^{(i)}(\dots,f,\dots)+\beta S^{(i)}(\dots,f',\dots)$, $S^{(i)}(f_\pi)=S^{(i)}(f)$ for a permutation $\pi$ within one field type, and $\mathrm{Im}\,S^{(i)}(f)=0$. Proposition~\ref{prop:C-compact}(a) gives convergence at each of these tuples, so each identity passes to the limit $i\to\infty$. For $S^0$, use the convergence of the Schwinger functions in Proposition~\ref{prop:C-mu0}(b) in the same way. The bound with continuum norms in Proposition~\ref{prop:C-compact}(a) is step~(iii) of its proof. Riemann sums give $\norm{\bg}_{2,\eps}\to\norm{\bg}_2$ and $\norm h_{L^1,\eps}\to\norm h_1$, while the sampled supremum is at most $\norm h_\infty$; hence $\limsup_iN_{\eps_i}(h)\le N(h)$.
\end{proof}

\paragraph{Translations of the fixed-torus continuum states.}

\paragraph{Translations of the fixed-torus continuum states.}
Throughout this paragraph $T=T_j$ is a fixed cubic torus of side $\ell:=\ell_T$, identified with the cell $Q:=[-\ell/2,\ell/2)^3$ as in Definition~\ref{def:C-sampling}. For $y\in\R^3$ let $[y]\in Q$ be its representative modulo $\ell\Z^3$. For $\rho\ge0$, $C_c^{(\rho)}$ is the space of real continuous scalar functions, or of real continuous 2-forms, with support in the closed Euclidean ball $\bar B_\rho$. For a 2-form we put $\norm{\bg}_\infty:=\max_{\mu<\kappa}\norm{g_{\mu\kappa}}_\infty$. A tuple $f=(\bg_1,\dots,\bg_n;h_1,\dots,h_m)$ is \emph{admissible} if all its entries lie in $C_c^{(\rho)}$ for some $\rho$ with $2\rho+2<\ell$. Such entries can be sampled, since $\bar B_\rho\subset B_{\rho'}$ for some $\rho'>\rho$ with $2\rho'+2<\ell$. For $a\in\R^3$ put $(\tau_ah)(x):=h(x-a)$ and $(\tau_a\bg)_x:=\bg_{x-a}$ (componentwise translation), and let $\tau_af$ be the entrywise translate. Every lattice spacing satisfies $\eps=2^{-K}\eps_0\le1$.

For $a\in\eps\Z^3$, translation by $a$ maps $T_\eps$ onto itself, because $\ell/\eps\in\N$. It acts on oriented cells by $\tau_a(x;v_1\wedge\dots\wedge v_k):=(x+a;v_1\wedge\dots\wedge v_k)$, on lattice forms by $(\tau_a\omega)(c):=\omega(\tau_{-a}c)$, and on field configurations by $\hat\tau_a(A,\varphi):=(\tau_aA,\tau_a\varphi)$.

\begin{lemma}[lattice translations]\label{app:G1:OS1-lat}
Let $a\in\eps\Z^3$ and $\mu_0\in(0,\mu_*]\cup\{\flat\}$.
\begin{enumerate}
\item[(i)] $d\tau_a=\tau_ad$, $\ip{\tau_a\omega}{\tau_a\omega'}_\eps=\ip{\omega}{\omega'}_\eps$ and $d^*\tau_a=\tau_ad^*$.
\item[(ii)] For real lattice sources, $F(\bg)\circ\hat\tau_a=F(\tau_{-a}\bg)$ and $\Xr(h)\circ\hat\tau_a=\Xr(\tau_{-a}h)$.
\item[(iii)] For real lattice 2-forms $\bg_i$ and real lattice functions $h_k$,
\[
\Big\langle\prod_iF(\tau_a\bg_i)\prod_k\Xr(\tau_ah_k)\Big\rangle_{\eps,T,\mu_0}=\Big\langle\prod_iF(\bg_i)\prod_k\Xr(h_k)\Big\rangle_{\eps,T,\mu_0}.
\]
\end{enumerate}
\end{lemma}
\begin{proof}
(i) $\tau_a$ maps oriented $k$-cells bijectively onto oriented $k$-cells. It commutes with reversal of orientation and with the oriented boundary, $\partial(\tau_ac)=\tau_a\partial c$, since it moves all vertices of a cube by the same vector and keeps the orientation. Hence $(d\tau_a\omega)(c)=\eps^{-1}\omega(\tau_{-a}\partial c)=\eps^{-1}\omega(\partial\tau_{-a}c)=(\tau_ad\omega)(c)$. Also $\tau_a$ permutes the unoriented cells, which all carry the weight $\eps^3$, so the pairing is invariant. Since $\tau_a^{-1}=\tau_{-a}$, taking adjoints gives $d^*\tau_a=\tau_ad^*$.

(ii) By (i), $F(\bg)(\tau_aA)=\ip{\bg}{d\tau_aA}_\eps=\ip{\bg}{\tau_adA}_\eps=\ip{\tau_{-a}\bg}{dA}_\eps$. Next, $r_x\circ\hat\tau_a=r_{x-a}$, and $W_\eps,z_\eps$ do not depend on the site (Lemma~\ref{app:G1:P1}(a)). Changing the summation variable to $y=x-a$ gives $\Xr(h)\circ\hat\tau_a=\sum_y\eps^3h(y+a)(r_y^2-W_\eps-z_\eps)=\Xr(\tau_{-a}h)$.

(iii) Put $G:=\prod_iF(\tau_a\bg_i)\prod_k\Xr(\tau_ah_k)$. By (ii), $G\circ\hat\tau_a=\prod_iF(\bg_i)\prod_k\Xr(h_k)$. Both $G$ and $G\circ\hat\tau_a$ are polynomials in the $F(\cdot)$ and $\Xr(\cdot)$. They are therefore elementary observables, integrable in every state (Lemma~\ref{lem:C-flatlimit}). Let $\mu_0>0$. The map $\hat\tau_a$ permutes the bond variables (without change of sign, since translations preserve the orientation of bonds) and the site variables. So it preserves Lebesgue measure, and by Lemma~\ref{lem:C-sym}(i) it preserves the action. Hence the probability measure of $\ev{\cdot}_{\eps,T,\mu_0}$ is $\hat\tau_a$-invariant, and $\ev{G\circ\hat\tau_a}_{\eps,T,\mu_0}=\ev{G}_{\eps,T,\mu_0}$. For $\flat$, Lemma~\ref{lem:C-flatlimit} applies to both $G$ and $G\circ\hat\tau_a$. It gives $\ev{G\circ\hat\tau_a}_{\eps,T,\flat}=\lim_{\mu_0\downarrow0}\ev{G\circ\hat\tau_a}_{\eps,T,\mu_0}=\lim_{\mu_0\downarrow0}\ev{G}_{\eps,T,\mu_0}=\ev{G}_{\eps,T,\flat}$.
\end{proof}

\begin{lemma}[sampling commutes with translations]\label{app:G1:OS1-samp}
Let $\rho\ge0$ and $a\in\eps\Z^3$ with $2(\rho+\abs a)+2<\ell$. Let $h$ be a scalar and $\bg$ a 2-form in $C_c^{(\rho)}$. Then $\tau_ah,\tau_a\bg\in C_c^{(\rho+\abs a)}$ can be sampled, and
\[
(\tau_ah)_\eps=\tau_a(h_\eps),\qquad(\tau_a\bg)_\eps=\tau_a(\bg_\eps).
\]
Consequently $S^{\eps,T,\mu_0}_{n,m}(\tau_af)=S^{\eps,T,\mu_0}_{n,m}(f)$ for every tuple $f$ with entries in $C_c^{(\rho)}$ and every $\mu_0\in(0,\mu_*]\cup\{\flat\}$.
\end{lemma}
\begin{proof}
\emph{Scalars.} Let $x\in Q$ be a site. Then $(\tau_ah)_\eps(x)=h(x-a)$ and $\tau_a(h_\eps)(x)=h([x-a])$. If $x-a\in Q$, the two values agree. Otherwise some coordinate of $x-a$ lies outside $[-\ell/2,\ell/2)$, so $\abs{x-a}\ge\ell/2>\rho$ and $h(x-a)=0$. Suppose $h([x-a])\ne0$. Then $\abs{[x-a]}\le\rho$, so $\abs{[x-a]+a}\le\rho+\abs a<\ell/2-1$, and $[x-a]+a$ lies in the interior of $Q$. But $[x-a]+a\equiv x$ modulo $\ell\Z^3$ and both points lie in $Q$, so $[x-a]+a=x$, that is, $x-a\in Q$, a contradiction. So both values vanish.

\emph{2-forms.} Let $c=(x;v_1\wedge v_2)$ be a plaquette with $x\in Q$, and put $z:=x+\frac\eps2(v_1+v_2)$; then $\abs{z-x}=\eps/\sqrt2<1$. By Definition~\ref{def:C-cells}, $(\tau_a\bg)_\eps(c)=\bg_{z-a}(v_1,v_2)$. Since $\tau_{-a}c=([x-a];v_1\wedge v_2)$, also $\tau_a(\bg_\eps)(c)=\bg_{y+\frac\eps2(v_1+v_2)}(v_1,v_2)$ with $y=[x-a]$. If $x-a\in Q$, then $y=x-a$ and the two values agree. Otherwise $\abs{x-a}\ge\ell/2$, so $\abs{z-a}\ge\ell/2-1>\rho$ and the first value vanishes. If the second value were nonzero, we would have $\abs y<\rho+1$, so $\abs{y+a}<\rho+\abs a+1<\ell/2$, and $y+a=x$ as in the scalar case, again a contradiction.

The last assertion follows from Lemma~\ref{app:G1:OS1-lat}(iii), applied to the sampled sources.
\end{proof}

\begin{lemma}[fixed-torus translations, local form]\label{app:G1:OS1torus}
Let $(\eps_i,T,\mu_i)$ be a sequence with $K_i\to\infty$ and $\mu_i\in(0,\mu_*]\cup\{\flat\}$, and let $S$ be the limit along a subsequence given by Proposition~\ref{prop:C-compact} with $T_{j_i}\equiv T$. Let $R\ge0$ and $a\in\R^3$ satisfy $2(R+\abs a)+2<\ell$. Then for every tuple $f$ with entries in $C_c^{(R)}$,
\[
S_{n,m}(\tau_a\bg_1,\dots,\tau_a\bg_n;\tau_ah_1,\dots,\tau_ah_m)=S_{n,m}(\bg_1,\dots,\bg_n;h_1,\dots,h_m).
\]
\end{lemma}
\begin{proof}
Let $c_{n,m}$ be the constant of Proposition~\ref{prop:C-bounds}, and let $\mathcal D:=\bigcup_K\eps_K\Z^3=\eps_0\Z[1/2]^3$, which is dense in $\R^3$.

\emph{Step 1 (dyadic translations).} Let $b\in\mathcal D$ with $2(R+\abs b)+2<\ell$, say $b\in\eps_K\Z^3$. The lattices are nested (P4), so $b\in\eps_i\Z^3$ whenever $K_i\ge K$, which holds for all large $i$. For those $i$, Lemma~\ref{app:G1:OS1-samp} gives $S^{(i)}_{n,m}(\tau_bf)=S^{(i)}_{n,m}(f)$. The tuples $f$ and $\tau_bf$ are admissible, with entries in $C_c^{(R)}$ and $C_c^{(R+\abs b)}$ respectively. So by Proposition~\ref{prop:C-compact}(a) both sides converge, and $S_{n,m}(\tau_bf)=S_{n,m}(f)$.

\emph{Step 2 (continuity).} Choose $R'$ with $R+\abs a<R'$ and $2R'+2<\ell$. Choose $a_k\in\mathcal D$ with $a_k\to a$ and $\abs{a_k}<R'-R$; this is possible since $\mathcal D$ is dense and $\abs a<R'-R$. Then $2(R+\abs{a_k})+2<\ell$, so Step~1 gives $S_{n,m}(\tau_{a_k}f)=S_{n,m}(f)$. All entries of $\tau_af$ and $\tau_{a_k}f$ lie in $C_c^{(R')}$. On tuples from $C_c^{(R')}$, $S_{n,m}$ is multilinear (Lemma~\ref{app:G1:multilin}(b)) and satisfies $\abs{S_{n,m}(f)}\le c_{n,m}\prod_i\norm{\bg_i}_2\prod_kN(h_k)$ (Proposition~\ref{prop:C-compact}(a)). On $C_c^{(R')}$ we have $\norm{\bg}_2\le(3\abs{B_{R'}})^{1/2}\norm\bg_\infty$ and $N(h)\le\max(1,\abs{B_{R'}})\norm h_\infty$. Put $C_{R'}:=\max(1,\abs{B_{R'}},(3\abs{B_{R'}})^{1/2})$, and write $f_1,\dots,f_{n+m}$ for the entries of $f$. By multilinearity,
\[
S_{n,m}(\tau_{a_k}f)-S_{n,m}(\tau_af)=\sum_{l=1}^{n+m}S_{n,m}\big(\tau_af_1,\dots,\tau_af_{l-1},\,\tau_{a_k}f_l-\tau_af_l,\,\tau_{a_k}f_{l+1},\dots,\tau_{a_k}f_{n+m}\big).
\]
Since $\norm{\tau_bf_l}_\infty=\norm{f_l}_\infty$, this gives
\[
\abs{S_{n,m}(\tau_{a_k}f)-S_{n,m}(\tau_af)}\le c_{n,m}C_{R'}^{n+m}\sum_l\norm{\tau_{a_k-a}f_l-f_l}_\infty\prod_{l'\ne l}\norm{f_{l'}}_\infty .
\]
The right side tends to $0$ because each $f_l$ is uniformly continuous. Hence $S_{n,m}(\tau_af)=\lim_kS_{n,m}(\tau_{a_k}f)=S_{n,m}(f)$.
\end{proof}

\begin{corollary}[translation invariance on admissible tuples]\label{app:G1:OS1-chain}
Let $S$ be as in Lemma~\ref{app:G1:OS1torus}, and let $a\in\R^3$ and $f$ be such that $f$ and $\tau_af$ are both admissible. Then $S_{n,m}(\tau_af)=S_{n,m}(f)$.
\end{corollary}
\begin{proof}
Let the entries of $f$ lie in $C_c^{(\rho_1)}$ and those of $\tau_af$ in $C_c^{(\rho_2)}$, and put $\rho:=\max(\rho_1,\rho_2)$, so that $2\rho+2<\ell$. Let $t\in[0,1]$ and let $f_l$ be an entry of $f$. Every point of $\supp\tau_{ta}f_l=\supp f_l+ta$ has the form $(1-t)y+t(y+a)$ with $y\in\supp f_l\subset\bar B_\rho$ and $y+a\in\supp\tau_af_l\subset\bar B_\rho$. Since $\bar B_\rho$ is convex, all entries of $\tau_{ta}f$ lie in $C_c^{(\rho)}$. Choose $M\in\N$ with $2(\rho+\abs a/M)+2<\ell$. For $k=0,\dots,M-1$, Lemma~\ref{app:G1:OS1torus}, applied with $R=\rho$, the tuple $\tau_{ka/M}f$ and the vector $a/M$, gives $S_{n,m}(\tau_{(k+1)a/M}f)=S_{n,m}(\tau_{ka/M}f)$. Chaining these $M$ identities gives the claim.
\end{proof}

\emph{Periodic test functions.} Let $C(T)$ be the space of real continuous functions on $T$, and $C(T;\Lambda^2)$ the space of real continuous 2-forms on $T$, both with the sup norm. Their elements are sampled by $h_\eps(x):=h(x)$ for $x\in T_\eps$, and by $\bg_\eps(c):=\bg_{z(c)}(v_1,v_2)$ for plaquettes $c=(x;v_1\wedge v_2)$, where $z(c)=x+\frac\eps2(v_1+v_2)\in T$; no identification with $Q$ is needed. Write $S^{(i),\mathrm{per}}_{n,m}$ for the corresponding lattice Schwinger functions. Put $\norm h_{L^1(T)}:=\int_T\abs h$, $\norm\bg_{L^2(T)}^2:=\sum_{\mu<\kappa}\int_Tg_{\mu\kappa}^2$ and $N_T(h):=\max(\norm h_{L^1(T)},\norm h_\infty)$. The periodisation of $h\in C_c^{(\rho)}$ with $2\rho+2<\ell$ is $(\mathsf Ph)(x+\ell\Z^3):=\sum_{v\in\Z^3}h(x+\ell v)$, a finite sum; the periodisation of a 2-form is defined componentwise. Translations act on $C(T)$ and $C(T;\Lambda^2)$ by the same formulas as on $\R^3$, and $\tau_{a+\ell v}=\tau_a$ for $v\in\Z^3$, so $a\mapsto\tau_a$ is an action of $\R^3/\ell\Z^3$.

\begin{proposition}[periodic extension and invariance under $\R^3/\ell\Z^3$]\label{app:G1:OS1-per}
Let $S$ be the limit of Lemma~\ref{app:G1:OS1torus}, along the subsequence $(i)$ given there. There is a further subsequence along which:
\begin{enumerate}
\item[(a)] $S^{(i),\mathrm{per}}_{n,m}(f)$ converges for all $n,m$ and every tuple $f$ with entries in $C(T;\Lambda^2)$ (2-forms) and $C(T)$ (scalars). The limit $S^{\mathrm{per}}_{n,m}$ is real, multilinear and symmetric within each field type, and satisfies $\abs{S^{\mathrm{per}}_{n,m}(f)}\le c_{n,m}\prod_i\norm{\bg_i}_{L^2(T)}\prod_kN_T(h_k)$;
\item[(b)] $S^{\mathrm{per}}_{n,m}(\mathsf Pf)=S_{n,m}(f)$ for every admissible tuple $f$;
\item[(c)] $S^{\mathrm{per}}_{n,m}(\tau_af)=S^{\mathrm{per}}_{n,m}(f)$ for every $a\in\R^3$ and every periodic tuple $f$; that is, $S^{\mathrm{per}}$ is invariant under $\R^3/\ell\Z^3$.
\end{enumerate}
\end{proposition}
\begin{proof}
(a) The proof of Proposition~\ref{prop:C-bounds} uses only \eqref{eq:C-T9}, which holds for arbitrary real lattice sources. So $\abs{S^{(i),\mathrm{per}}_{n,m}(f)}\le c_{n,m}\prod_i\norm{(\bg_i)_{\eps_i}}_{2,\eps_i}\prod_kN_{\eps_i}((h_k)_{\eps_i})$. The torus has $(\ell/\eps)^3$ sites and $(\ell/\eps)^3$ plaquettes of each orientation. Hence $\norm{h_\eps}_{L^1,\eps}\le\ell^3\norm h_\infty$ and $\norm{\bg_\eps}_{2,\eps}\le(3\ell^3)^{1/2}\norm\bg_\infty$. So the lattice Schwinger functions are bounded by $c_{n,m}C_T^{n+m}\prod_l\norm{f_l}_\infty$, uniformly in $i$, with $C_T:=\max(1,\ell^3,(3\ell^3)^{1/2})$. They are multilinear and symmetric within each field type, by the argument of Lemma~\ref{app:G1:multilin}(a): sampling is linear, and $F$ and $\Xr$ are linear in the source. Hence they are uniformly Lipschitz on bounded sets of the product of the sup-normed spaces. Since $T$ is a compact metric space, $C(T)$ and $C(T;\Lambda^2)$ are separable. Choose countable dense subsets. A diagonal argument over these subsets and over all $(n,m)$ gives a subsequence along which $S^{(i),\mathrm{per}}_{n,m}$ converges on all tuples drawn from them. By the uniform Lipschitz bound, as in step~(ii) of the proof of Proposition~\ref{prop:C-compact}, the convergence extends to all periodic tuples. Reality, multilinearity and symmetry are identities between finitely many convergent values, so they pass to the limit. The functions are uniformly continuous, so the Riemann sums converge: $\norm{h_\eps}_{L^1,\eps}\to\norm h_{L^1(T)}$ and $\norm{\bg_\eps}_{2,\eps}\to\norm\bg_{L^2(T)}$ as $\eps\to0$. Moreover, lattice sup norms are at most the continuum sup norms. This gives the bound.

(b) Let the entries of $f$ lie in $C_c^{(\rho)}$ with $2\rho+2<\ell$. For a site $x\in Q$ and $0\ne v\in\Z^3$, some coordinate of $x+\ell v$ has modulus at least $\ell/2$, so $\abs{x+\ell v}\ge\ell/2>\rho$; hence $(\mathsf Ph)(x)=h(x)$. For a plaquette $c=(x;v_1\wedge v_2)$ with $x\in Q$ and $z=x+\frac\eps2(v_1+v_2)$, we have $\abs{z+\ell v}\ge\ell/2-1>\rho$ for $v\ne0$, so $(\mathsf P\bg)_{z}=\bg_{z}$. Thus the sampled sources of $\mathsf Pf$ and of $f$ coincide, and so do the lattice Schwinger functions, for every $i$. The subsequence of (a) is a subsequence of the one along which $S^{(i)}_{n,m}(f)\to S_{n,m}(f)$, and (b) follows.

(c) Let $b\in\eps_i\Z^3$. Then $(\tau_bh)_{\eps_i}(x)=h(x-b)=(\tau_bh_{\eps_i})(x)$ on $T_{\eps_i}$. Since $z(\tau_{-b}c)=z(c)-b$ in $T$, also $(\tau_b\bg)_{\eps_i}(c)=\bg_{z(c)-b}(v_1,v_2)=\bg_{\eps_i}(\tau_{-b}c)=(\tau_b\bg_{\eps_i})(c)$. By Lemma~\ref{app:G1:OS1-lat}(iii), $S^{(i),\mathrm{per}}_{n,m}(\tau_bf)=S^{(i),\mathrm{per}}_{n,m}(f)$. If $b\in\eps_K\Z^3$, this holds for all $i$ with $K_i\ge K$, by nesting (P4). Passing to the limit with (a) gives invariance under the dense set $\mathcal D=\eps_0\Z[1/2]^3$. For arbitrary $a$, take $a_k\in\mathcal D$ with $a_k\to a$. The telescoping estimate of Step~2 of the proof of Lemma~\ref{app:G1:OS1torus}, with $C_{R'}$ replaced by $C_T$ and with the bound of (a), gives
\[
\abs{S^{\mathrm{per}}_{n,m}(\tau_{a_k}f)-S^{\mathrm{per}}_{n,m}(\tau_af)}\le c_{n,m}C_T^{n+m}\sum_l\norm{\tau_{a_k-a}f_l-f_l}_\infty\prod_{l'\ne l}\norm{f_{l'}}_\infty\to0 ,
\]
because continuous functions on the compact torus are uniformly continuous.
\end{proof}

By (b) and (c), $S_{n,m}(\tau_af)=S^{\mathrm{per}}_{n,m}(\mathsf P\tau_af)=S^{\mathrm{per}}_{n,m}(\tau_a\mathsf Pf)=S^{\mathrm{per}}_{n,m}(\mathsf Pf)=S_{n,m}(f)$ whenever $f$ and $\tau_af$ are admissible. This gives a second proof of Corollary~\ref{app:G1:OS1-chain}. Lemma~\ref{app:G1:OS1torus} is the fixed-torus clause of Theorem~\ref{thm:C-OS1}(a).

\begin{remark}[the torus group]\label{app:G1:OS1torus-rem}
At a fixed torus, the continuum states of Proposition~\ref{prop:C-compact} are defined only on admissible tuples, so the group $\R^3/\ell\Z^3$ does not act on their domain. Indeed, let $f\ne0$ be admissible with entries in $C_c^{(\rho)}$, $2\rho+2<\ell$, and let $0\ne a\in\ell\Z^3$. Every point $y+a$ with $y\in\bar B_\rho$ satisfies $\abs{y+a}\ge\ell-\rho>\ell/2+1$, so $\tau_af$ is not admissible. For these states, Corollary~\ref{app:G1:OS1-chain} is therefore the complete statement of translation invariance: $S_{n,m}(\tau_af)=S_{n,m}(f)$ whenever both sides are defined. Invariance under the torus group itself is a statement about the periodic extension of Proposition~\ref{app:G1:OS1-per}. That extension exists along a further subsequence and agrees with $S$ on periodisations of admissible tuples.
\end{remark}

\paragraph{Parts (a) and (c) of Proposition~\ref{prop:C-mu0} on the class $\cD$.}

\begin{lemma}[class-$\cD$ monotonicity]\label{app:G1:classD}
Fix real test functions $\bg_1,\dots,\bg_n,h_1,\dots,h_p$, and let
\[
E^\cD_{\mathrm{obs}}:=\{\zeta=(x,w)\in\R^{n+p}:\ k(\zeta)\le0\ \text{pointwise},\ \norm{k(\zeta)}_\infty<H/2\},\qquad k(\zeta)=\sum_kw_kh_k.
\]
This is the set of real points of the tube at which $\Theta(\zeta)$ lies in the class $\cD$ of Subsection~\ref{ssec:observables}. Along the subsequence of Proposition~\ref{prop:C-mu0}, for every $\zeta\in E^\cD_{\mathrm{obs}}$: $\mu\mapsto G^\mu(\zeta)$ is non-increasing on $Q_0$; $G^\mu(\zeta)\le G^\flat(\zeta)$; and $G^0(\zeta)\le G^\flat(\zeta)$.
\end{lemma}
\begin{proof}
Note first that $E^\cD_{\mathrm{obs}}\subset\Tub$, since $k_+(\zeta)=0$ and $k_-(\zeta)=-k(\zeta)$. Fix $\zeta\in E^\cD_{\mathrm{obs}}$ and put $\bg:=\sum_ix_i\bg_i$ and $q:=-k(\zeta)\ge0$, with $\norm q_\infty<H/2$. Fix also a member $(\eps,T)=(\eps_i,T_{j_i})$ of the subsequence, and write $\bg_\eps,q_\eps$ for the sampled functions; then $q_\eps\ge0$. Linearity of sampling and of the fields (Lemma~\ref{app:G1:multilin}) gives
\begin{gather*}
\Theta(\zeta)=\mathrm{e}^{F(\bg_\eps)}\mathrm{e}^{-\Xr(q_\eps)}=c_\eps\,\mathrm{e}^{-r^2(\gamma)}\mathrm{e}^{\ell(A)},\\
\gamma_x:=\eps^3q_\eps(x)\ge0,\quad \ell(A):=\ip{d^*\bg_\eps}{A}_\eps,
\end{gather*}
with $c_\eps:=\exp\big((W_\eps+z_\eps)\sum_x\eps^3q_\eps(x)\big)>0$. By (P1), $c_\eps$ does not depend on $\mu_0$, and the matter part $S_m$ of \eqref{eq:C-action} does not either.
\emph{Covariances.} By \eqref{eq:C-state}, $C_\mu=(-\Delta_\eps+\mu^2)^{-1}$. For $0<\mu'\le\mu$, Lemma~\ref{app:G1:mu-mono}(b) gives $C_\mu\le C_{\mu'}$ as forms, the zero mode included; on the zero mode this is $s(\mu)=1/(\mu^2\abs T)\le s(\mu')$.
\emph{Lattice inequalities.} Lemma~\ref{lem:C-cov} (class $\cD$, second inequality) with $C_0=C_\mu\le C_1=C_{\mu'}$ gives $\ev{\Theta(\zeta)}_{\eps,T,\mu}\le\ev{\Theta(\zeta)}_{\eps,T,\mu'}$; the constant $c_\eps$ factors out of the normalised expectations. $\Theta(\zeta)$ is an elementary observable, being a product of $\mathrm{e}^{tF(\bg)}$ ($t=1$) and $\mathrm{e}^{-\Xr(q)}$ with $q\ge0$. So Lemma~\ref{lem:C-flatlimit} gives $\ev{\Theta(\zeta)}_{\eps,T,\flat}=\lim_{\mu'\downarrow0}\ev{\Theta(\zeta)}_{\eps,T,\mu'}\ge\ev{\Theta(\zeta)}_{\eps,T,\mu}$.
\emph{Limits.} Since $\zeta\in\Tub$, Proposition~\ref{prop:C-compact}(b) gives $G_i^\mu(\zeta)\to G^\mu(\zeta)$ for each $\mu\in Q_0$ and $G_i^\flat(\zeta)\to G^\flat(\zeta)$. Non-strict inequalities pass to these limits, which gives the first two claims. By Proposition~\ref{prop:C-mu0}(b), $G^\mu\to G^0$ locally uniformly on $\Tub$, in particular at $\zeta$. So $G^0(\zeta)=\lim_{\mu\downarrow0,\,\mu\in Q_0}G^\mu(\zeta)\le G^\flat(\zeta)$. Finally, the lattice value at $\zeta$ depends only on $(\bg,q)$, by linearity; hence so does the limit. Thus $S^{\mu_0}(\Theta):=G^{\mu_0}(\zeta)$ is well defined for $\Theta=\exp[F(\bg)-\Xr(q)]\in\cD$, which is the meaning of Theorem~\ref{thm:main}(iv). So parts (a) and (c) of Proposition~\ref{prop:C-mu0} hold on $E^\cD_{\mathrm{obs}}$, which gives Theorem~\ref{thm:main}(iv).
\end{proof}

\section{Counterterm and vacuum identifications}\label{app:G1:area-UV}

This appendix proves Proposition~\ref{prop:UV-B1e}: King's mass counterterm and vacuum constant on a lattice torus are Ba{\l}aban's (Lemmas~\ref{app:G1:ctid} and \ref{app:G1:vac}), and it states and proves Reisz's power-counting theorems in three dimensions, item~3 of Subsection~\ref{ssec:UV-other} (Proposition~\ref{app:G1:reisz3d}). The torus is $T=T_j$ with spacing $\eps$; $-\Delta_\eps=d^*d+dd^*$ on bond functions.

\paragraph{King's counterterm and vacuum constant are Ba{\l}aban's.}\label{app:G1:ssec-B1e}
Fix a torus $T=T_j$, a spacing $\eps$, a Higgs mass $M>0$ and the reference photon mass $\mu_*>0$. \emph{King's model} is the action \cite[(2.4), p.~652]{K1} on $\Omega=T_{j,\eps}$ with $m^2=M^2$ and $\mu_0=\mu_*$, together with its counterterm $\delta m^2_{K,T}(\eps;M)$ (Definition~\ref{def:counterterm}). \emph{Ba{\l}aban's model} is the action \cite[(1.11), p.~605]{Ba1} on the same lattice torus, in the Feynman gauge, with $m^2=M^2$, $\mu_0=\mu_*$ and $m_0^2=m^2+\delta m^2$. Write $\delta m^2_{\rm B}(\eps;M)$ for Ba{\l}aban's counterterm as defined explicitly in \cite[(1.19)--(1.23), pp.~416--417]{Ba3}. The paper \cite{Ba1} only names $\delta m^2$ (``it can be written down explicitly but we will not do it'', \cite[p.~605]{Ba1}); \cite{Ba3} completes the definition, and we read the $\delta m^2$ of \cite{Ba1} as this one.

\begin{lemma}[King's mass counterterm is Ba{\l}aban's]\label{app:G1:ctid}
On $T_{j,\eps}$, $\delta m^2_{K,T}(\eps;M)=\delta m^2_{\rm B}(\eps;M)$ coefficient by coefficient in $e^\alpha\lambda^\beta$, $2\le\alpha+2\beta\le4$, as an identity of finite lattice sums. Hence $a_{K,T}(M)=M^2+\delta m^2_{\rm B}(\eps;M)$ is the bare coefficient $m_0^2$ of \cite[(1.11)]{Ba1} on $T_{j,\eps}$, i.e.\ the coefficient at which Ba{\l}aban's theorems (Theorems~\ref{thm:UV-K0minus}, \ref{thm:UV-K0plus}; \cite[p.~417]{Ba3}) are stated. This is the mass part of Proposition~\ref{prop:UV-B1e} and Remark~\ref{rem:counterterm}(i).
\end{lemma}

\begin{proof}
\emph{Step 1 (the two actions agree term by term).} King's action is
$\frac12\ip{A}{(-\Delta^{\eps,\Omega}+\mu_0^2)A}+\frac12\ip{\varphi}{(-\Delta^{\eps,\Omega}_A+m^2)\varphi}+\sum_{x\in\Omega}\eps^3(\frac12\delta m^2(x)\abs{\varphi(x)}^2+\lambda\abs{\varphi(x)}^4)+E_0+E_1$ \cite[(2.4), p.~652]{K1}. Ba{\l}aban's is $\frac12\ip{\varphi}{(-\Delta^\eps_A)\varphi}+\sum_{x\in T_\eps}\eps^3(\frac12m_0^2\abs\varphi^2+\lambda\abs\varphi^4)+\frac12\ip{A}{(-\Delta^\eps+\mu_0^2)A}+E$ \cite[(1.11), p.~605]{Ba1}. Both use the same covariant difference operator, with $U(A)=\exp(q\eps eA)$ \cite[(1.7)]{Ba1}, \cite[p.~651]{K1}, and both photon terms are in the Feynman form. On $\Omega=T_{j,\eps}$, King's $\delta m^2(x)$ is constant by translation invariance. So the vertices are the same, and so are the propagators: the scalar propagator $C^\eps_0=(-\Delta^\eps_0+M^2)^{-1}$ and the photon propagator $C^\eps=(-\Delta^\eps+\mu_*^2)^{-1}$ (which is $(-\Delta_\eps+\mu_*^2)^{-1}$ on bond functions, Subsection~\ref{ssec:action}), both on $T_{j,\eps}$ \cite[p.~416]{Ba3}.

\emph{Step 2 (the two defining equations).} King writes $(-\delta m^2+\Gamma^\eps)$ for ``the full one particle irreducible graph with two (undifferentiated) external scalar field lines''. He sets $\delta m^2(x)=\sum_{2\le\alpha+2\beta\le4}e^\alpha\lambda^\beta\delta m^2(x;\alpha,\beta)$ and imposes $-\delta m^2(x;\alpha,\beta)+\sum_{y\in\Omega}\eps^d\Gamma^\eps(x,y;\alpha,\beta)=0$, ``solved recursively'' \cite[(3.32)--(3.33), p.~659]{K1}. Ba{\l}aban expands the two-point function on $T_\eps$ as $G^\eps=\sum_nC^\eps_0[(-\delta m^2+\Sigma^\eps+\partial^{\eps*}\Sigma^\eps_1+\Sigma_1^{\eps*}\partial^\eps+\partial^{\eps*}\Sigma^\eps_2\partial^\eps)C^\eps_0]^n$, with $\Sigma^\eps,\Sigma^\eps_1,\Sigma^\eps_2$ amputated one-particle-irreducible \cite[(1.19)--(1.21), p.~416]{Ba3}. He then imposes $-\delta m^2_{(\alpha,\beta)}+\sum_{x\in T_\eps}\eps^d\Sigma^\eps_{(\alpha,\beta)}(x)=0$ for $2\le\alpha+2\beta\le4$ \cite[p.~417]{Ba3}. Both definitions include all graphs of each order, with the lower-order counterterms inserted.

\emph{Step 3 (the parts with derivative legs have zero sum).} Let $K(x,y)$ be a translation-invariant kernel on $T_{j,\eps}$. If $K=B\circ d$, with $d$ the lattice derivative of Subsection~\ref{ssec:action} acting componentwise, then $\sum_y\eps^3K(x,y)=(B\,d1)(x)=0$, because $d1=0$. If $K=d^*\circ B'$ with $B'$ translation invariant, then $(B'1)_{x,\mu}=:c_\mu$ does not depend on $x$, and
\[
(d^*c)(x)=\sum_\mu\eps^{-1}\big(c_\mu-c_\mu\big)=0.
\]
The kernels $\Sigma_1^{\eps*}\partial^\eps$ and $\partial^{\eps*}\Sigma_2^\eps\partial^\eps$ are of the first type and $\partial^{\eps*}\Sigma_1^\eps$ is of the second. So the zero-momentum sum of the full amputated one-particle-irreducible two-point function equals $\sum_x\eps^3\Sigma^\eps(x)$. The same sum also does not depend on how a convention divides that function into undifferentiated and differentiated legs. By translation invariance, $\sum_y\eps^3\Gamma^\eps(x,y)$ does not depend on $x$. Hence both conditions of Step~2 say the same thing: at each order, $\delta m^2_{(\alpha,\beta)}$ is the zero-momentum value of the order-$(\alpha,\beta)$ amputated one-particle-irreducible two-point function on $T_{j,\eps}$, with the lower-order counterterms inserted.

\emph{Step 4 (induction on the order).} Order the coefficients by $w:=\alpha+2\beta\in\{2,3,4\}$. In a one-particle-irreducible two-point graph of order $(\alpha,\beta)$ that contains a counterterm insertion $\delta m^2_{(\alpha',\beta')}$, there is at least one other vertex, which carries a factor $e$, $e^2$ or $\lambda$; so $\alpha'+2\beta'<w$. Assume the counterterms of the two models agree at all orders with $\alpha'+2\beta'<w$. By Step~1, the order-$(\alpha,\beta)$ coefficients of the two two-point functions are then the same finite lattice sums, and by Step~3 the two equations of Step~2 give the same number. The odd orders in $e$ vanish in both models (Definition~\ref{def:counterterm}).
\end{proof}

\begin{lemma}[the vacuum constants]\label{app:G1:vac}
Let $\nbar=13$, and let $\delta m^2=\delta m^2_{K,T}(\eps;M)=\delta m^2_{\rm B}(\eps;M)$ (Lemma~\ref{app:G1:ctid}).
\begin{enumerate}
\item[(a)] King's $E_0(M)$ \cite[(2.6), p.~652]{K1} equals Ba{\l}aban's $E_0=E_{0,v}(\mu_*)+E_{0,s}(M)$ \cite[(1.12), p.~606]{Ba1}.
\item[(b)] King's $E_1(M)$ \cite[(3.35), p.~659]{K1} equals Ba{\l}aban's $E_1$ \cite[(1.13), p.~606]{Ba1}. Hence $Z_K(T;M)$ (Subsection~\ref{ssec:UV-K0}) equals Ba{\l}aban's $Z^\eps$ at $(m^2,\mu_0)=(M^2,\mu_*)$, with no correction term.
\item[(c)] With the photon measure normalised, $Z(\Gamma;a):=\E_{N(0,\Gamma)}[Z_\varphi(A;a)]$ as in Subsection~\ref{ssec:T6-statement}, one has $Z_K(T;M)=Z(C_{\mu_*};a_{K,T}(M))\,\mathrm{e}^{-E_{0,s}(M)-E_1(M)}$. The constant $\gamma_T(M)$ of Proposition~\ref{prop:ch-CN} formed in this pairing and the one formed in the Lebesgue pairing differ by $\abs T^{-1}E_{0,v}(\mu_*)$. This term does not depend on $M$ and cancels in every difference $\gamma_T(M)-\gamma_T(M')$.
\end{enumerate}
This is the vacuum-constant part of Proposition~\ref{prop:UV-B1e} and Remark~\ref{rem:T6-textfixes}(i).
\end{lemma}

\begin{proof}
Write $S_{\rm rest}$ for the common action of Step~1 of Lemma~\ref{app:G1:ctid} without its additive constants. It depends on $(e,\lambda)$, also through $\delta m^2$, in the same way in both models. Both partition functions use the Lebesgue measure on the site and bond variables (\cite[(2.5)]{K1}; ``natural Lebesgue measure'', \cite[p.~605]{Ba1}).

(a) King: $\mathrm{e}^{E_0}=\int(dA)(d\varphi)\exp[-\frac12\ip{A}{(-\Delta^{\eps,\Omega}+\mu_0^2)A}-\frac12\ip{\varphi}{(-\Delta^{\eps,\Omega}+m^2)\varphi}]$. Ba{\l}aban: $\int dA\,\mathrm{e}^{-\frac12\ip{A}{(-\Delta^\eps+\mu_0^2)A}}\int d\varphi\,\mathrm{e}^{-\frac12\ip{\varphi}{(-\Delta_0^\eps+m^2)\varphi}}=\mathrm{e}^{E_{0,v}}\mathrm{e}^{E_{0,s}}=\mathrm{e}^{E_0}$. These are the same Gaussian integral: the photon factor is included on both sides, the sign is the same, and both are taken at $m^2=M^2$, not at $m_0^2$.

(b) King's (3.35) takes the Taylor coefficients, of total order $1\le\alpha+\beta\le\nbar$ at $e=\lambda=0$, of $\ln\int\exp[-S^\eps+E_1]=\ln\int \mathrm{e}^{-S_{\rm rest}}-E_0$. Ba{\l}aban's (1.13) takes the same coefficients of $\log\int\exp(-S+E)=\log\int \mathrm{e}^{-S_{\rm rest}}$. The two functions differ by the constant $E_0$, which does not depend on $(e,\lambda)$ (it is defined at $e=\lambda=0$), so every coefficient with $\alpha+\beta\ge1$ is the same on both sides. The two functions are therefore not literally the same: they agree only up to this constant. Both sides use the same index set $1\le\alpha+\beta\le\nbar$. King fixes $\nbar=13$ \cite[p.~657]{K1}, and Ba{\l}aban's definition admits it (``$\nbar>12$'', \cite[p.~418]{Ba3}). Hence $E_0+E_1$ is the same, and $Z_K(T;M)=\int \mathrm{e}^{-S_{\rm rest}-E_0-E_1}=Z^\eps$.

(c) Since $\int dA\,\mathrm{e}^{-\frac12\ip{A}{(-\Delta_\eps+\mu_*^2)A}}=\mathrm{e}^{E_{0,v}(\mu_*)}$, the Lebesgue integral equals $\mathrm{e}^{E_{0,v}(\mu_*)}Z(C_{\mu_*};a_{K,T}(M))$. With (b) this gives $Z_K=Z(C_{\mu_*};a_{K,T}(M))\,\mathrm{e}^{E_{0,v}-E_0-E_1}=Z(C_{\mu_*};a_{K,T}(M))\,\mathrm{e}^{-E_{0,s}-E_1}$. At $e=\lambda=0$ one has $a_{K,T}(M)=M^2$, so the order-zero term of $\log Z(C_{\mu_*};a_{K,T}(M))$ is $E_{0,s}(M)$, as stated in Subsection~\ref{ssec:T6-statement}. Its Taylor coefficients of order at least $1$ are those of $E_1$. So the two pairings give $\gamma_T(M)$'s that differ by $\abs T^{-1}E_{0,v}(\mu_*)$, which does not depend on $M$. Taking the photon-normalised quantity together with Ba{\l}aban's full $E_0$ would be off by $|T|^{-1}E_{0,v}(\mu_*)$ per volume, and this diverges as $\eps\to0$.
\end{proof}

\paragraph{Reisz's power counting in three dimensions.}\label{app:G1:ssec-reisz}
We use Reisz's notation \cite{Rei88,Rei88b}, with $a$ the lattice spacing. $\mathcal F_{\rm R}$ is Reisz's class of lattice Feynman integrands $F=V/\prod_i(\eta_i(l_ia)/a^2+\mu_i^2)$ (\cite[Def.~2.5]{Rei88}). Their numerators $V=a^{-m}F(ua,wa)$ are $(2\pi/a)$-periodic in the loop momenta, $\mu_i^2>0$ for every line, and the line momenta form a \emph{natural} set (\cite[Def.~3.1]{Rei88}). For a Zimmermann subspace $H$ of the loop space, $\mathrm{degr}_H$ is Reisz's UV degree. $\mathrm{BZ}=[-\pi/a,\pi/a]^{3m}$ is the Brillouin zone of $m$ three-dimensional loop momenta, and $\hat l^2:=\sum_\mu(4/a^2)\sin^2(al_\mu/2)$ is King's lattice momentum square. Reisz states both theorems for $d=4$ and remarks that ``the following considerations can easily be extended to other dimensions'' \cite[p.~86]{Rei88}.

\begin{proposition}[Reisz's theorems in three dimensions]\label{app:G1:reisz3d}
Let $d=3$.
\begin{enumerate}
\item[(a)] Let $F\in\mathcal F_{\rm R}$ have three-dimensional momenta, positive masses and a natural set of line momenta, and let $3\dim H+\mathrm{degr}_HF<0$ for every Zimmermann subspace $H$. Then $\lim_{a\to0}\int_{\rm BZ}F$ exists, and there are $a_0>0$ and $C<\infty$ with $\int_{\rm BZ}\abs{F(k,q;a)}\,d^{3m}k\le C$ for all $a<a_0$, uniformly for external momenta $q$ in a bounded set.
\item[(b)] Let $\Gamma$ be a one-particle-irreducible diagram with massive lines, renormalised by Reisz's forest formula with subtraction degrees $\delta(\gamma)$ obeying \cite[(4-3)--(4-4)]{Rei88b}; in particular $\delta(\gamma)=\omega(\gamma)$, with $\omega$ computed with $3m(\gamma)$. Then the forest integrand $R_\Gamma$ lies in $\mathcal F_{\rm R}$ with $3h+\mathrm{degr}_HR_\Gamma\le-1$ for every $H$, so (a) applies to $R_\Gamma$.
\item[(c)] Suppose in addition that every vertex function and every coefficient of a subtraction operator is entire in its momentum arguments, and that every propagator is $(\hat l^2+\mu^2)^{-1}$. Let $\mu_{\min}$ be the smallest mass. Let $c_\Gamma$ be the largest value, over all line and subdiagram momenta of $R_\Gamma$, of the sum of the absolute values of their coefficients on the loop momenta. Fix $\eta_0>0$ with $1.087\,c_\Gamma^2\eta_0^2\le\mu_{\min}^2/2$ and $c_\Gamma\eta_0\le1$, and let $a_0\le1$. Then
\[
\sup_{a<a_0}\ \sup_{\eta\in\R^{3m},\ \max_j\abs{\eta_j}\le\eta_0}\int_{\rm BZ}\abs{R_\Gamma(k+i\eta,q;a)}\,d^{3m}k\le C'<\infty.
\]
\end{enumerate}
The constants depend on the graph, the masses, the bound on the external momenta and $\eta_0$; they do not depend on $\eps$ or on the torus, and they are not explicit. This is item~3 of Subsection~\ref{ssec:UV-other}, as used in Proposition~\ref{prop:ch-CN}.
\end{proposition}
\noindent For (a) and (b) the proofs are those printed in \cite{Rei88,Rei88b} for $d=4$; the proof below does not repeat them, but lists every place where the dimension enters and checks each one, and the bound form is read off the majorants of the printed proofs. Part (c) is deduced from (a) and (b).

\begin{proof}
\emph{Step 1 (what Reisz's proof bounds).} The integral over BZ is split by non-negative cut-off functions into finitely many sectors: lines near a pole and lines away from all poles \cite[(4.1)--(4.7)]{Rei88}. Naturalness allows the poles to be translated to $0$ \cite[App.~D]{Rei88}. Then \cite[(4.8)--(4.9)]{Rei88} give $1/(\eta(la)/a^2+\mu^2)\le\alpha/(l^2+\mu^2)$ near a pole and $\le\gamma a^2$ away from the poles. By \cite[(7.4)--(7.5b)]{Rei88}, each sector integral is bounded by integrals of non-negative majorants, and these are controlled by the continuum power-counting theorem \cite[Thm.~2, p.~96]{Rei88}. Hence the proof bounds $\int_{\rm BZ}\abs F$, not only $F$, uniformly for bounded $q$ \cite[pp.~96, 102]{Rei88}. Complex-valued numerators are allowed: $\abs V\le\abs{\rRe V}+\abs{\operatorname{Im}V}$, and taking real and imaginary parts does not raise degrees. Neither does the photon index \cite[Sect.~6, p.~100]{Rei88b}.

\emph{Step 2 (where the dimension enters).} The dimension enters only as the dimension of each loop momentum, and the following replacements are the only changes.
(1) The degree definitions: \cite[(3.8), (5.9), (7.9), (7.14), (9.3), (9.29)]{Rei88} and \cite[(2-28), (4-11)--(4-16), (5-5), (5-10)]{Rei88b}. Each has the form ``$4\times$ (dimension of the loop space)'', and $4\to3$.
(2) The counting of \cite[App.~C]{Rei88b}, which is a \emph{lower} bound, $M_U(\Gamma)\ge4h$ and $M_U(\gamma)\ge4\,\mathrm{rank}$ \cite[(C-1)--(C-2), p.~106]{Rei88b}. Its proof is the count $M_U=4\cdot\#$, so it becomes $M_U\ge3h$; then \cite[(5-13)]{Rei88b} gives $3h+\mathrm{degr}_HR_\Gamma\le-1$.
(3) The scaling step \cite[Sect.~9]{Rei88}. The outer variable is $t_1\in\R^3$ and the inner variables $(t_2,\dots,t_m)\in\R^{3(m-1)}$ are rescaled by $\norm{t_1}$ \cite[(9.16)]{Rei88}. The Jacobian $\norm{t_1}^{3(m-1)}$ gives the factor $\norm{t_1}^{\omega_0-3}$ with $\omega_0:=3m+\Delta-2\sum n_j$ \cite[(9.17a)]{Rei88}.
(4) Two constants that use the dimension implicitly. The choice of $\delta$ in \cite[(7.8), p.~102]{Rei88} needs $\sqrt d\le4$. The proof of \cite[Lemma~D.1, p.~125]{Rei88} uses $\norm{k_j}\le2\pi/a$ on $[-\pi/a,\pi/a]^d$, i.e.\ $\sqrt d\le2$. Both hold for $d=3$.
All the other ingredients are dimension-free: the propagator bounds (4.8)--(4.9), Theorem~3 and Lemmas 6.1--6.3, Lemmas 8.1--8.10, Appendices A--D of \cite{Rei88}, the forest formula, Lemma~3.1 and Lemmas 5.1--5.2 of \cite{Rei88b}, and the subtraction operators $t^0,t^1$ (with $\sin(\eps q)/\eps$) of \cite[App.~A]{Rei88b}. The IR hypothesis is only $\mu_i^2>0$, which enters through lower bounds that do not depend on the dimension. Massless lines are not covered, and no graph bound of the construction uses one.

\emph{Step 3 (the outer integral, which Reisz does not print).} Reisz writes ``the remaining integration can be done without problems'' \cite[p.~119]{Rei88}. Here is that integral. With $\lambda'=b(\lambda+d)/\norm{t_1}+e$ \cite[Lemma~9.1, p.~115]{Rei88}, polar coordinates turn $r^{2}\cdot r^{\omega_0-3}$ into $r^{\omega_0-1}$. The outer integral $\int_1^{\lambda+K}r^{\omega_0-1}f(\lambda')\,dr$, with $f$ from \cite[(9.24)]{Rei88}, therefore has the same exponents as for $d=4$. If $f=1$ and $\omega_0<0$, it is bounded. If $f=\lambda'^{-1}\log^{m-1}\lambda'$ and $\omega_0\le-1$, it is $\le c\lambda^{-1}\log^m\lambda$. If $f=\lambda'^{\bar\omega}\log^{m-1}\lambda'$, it is $\le c\lambda^{\max(\omega_0,\bar\omega)}\log^m\lambda$. This is \cite[(9.30)]{Rei88} in every dimension. Lemma~9.2 then applies the induction hypothesis to the inner integral, with $\omega(\tilde H)=\omega(H)$ \cite[(9.29)]{Rei88}. This proves (a), and (b) follows by Step~2(2).

\emph{Step 4 (proof of (c)).} Write $R_\Gamma=V/(B_1B_2)$ \cite[(5-1)]{Rei88b}, with $V=\sum_ia^{-m_i}F_i(ka,qa)$ and each $F_i$ entire.
\emph{Numerator.} Put $V_\eta(k):=V(k+i\eta)=\sum_ia^{-m_i}\tilde F_i(ka,qa,\eta a)$ with $\tilde F_i(x,y,z):=F_i(x+iz,y)$. These are smooth and $2\pi$-periodic in $x$, and $a^{-m_i}\tilde F_i(ka,qa,\eta a)\to P_i(k+i\eta,q)$. So $V_\eta$ is in Reisz's numerator class, with $(q,\eta)$ as external variables. For a Zimmermann subspace with internal variables $u$ and complementary variables $w$, the shift $k\mapsto k+i\eta$ acts as $(u,w)\mapsto(u+iA\eta,w+iB\eta)$ with fixed real matrices $A,B$. Each derivative $\partial_w^b\partial_\eta^c$ at $(w,\eta)=0$ is a combination of terms $\partial_u^{c_1}\partial_w^{b+c_2}F(u,0)$ with $\abs{c_1}+\abs{c_2}=\abs c$. If $\abs b+\abs c$ is below Reisz's vanishing order $r_u$, these vanish identically. Hence $\mathrm{degr}_uV_\eta\le\mathrm{degr}_uV$.
\emph{Denominator.} For $l=\kappa+i\zeta$ with $\abs{a\zeta_\mu}\le1$,
\[
\rRe\big(1-\cos(al_\mu)\big)=(1-\cos a\kappa_\mu)-\cos(a\kappa_\mu)(\cosh a\zeta_\mu-1)\ge(1-\cos a\kappa_\mu)-(\cosh a\zeta_\mu-1).
\]
Moreover $(2/a^2)(\cosh a\zeta-1)=\zeta^2\big(\sinh(a\zeta/2)/(a\zeta/2)\big)^2\le1.0862\,\zeta^2$. Hence $\rRe\hat l^2\ge\hat\kappa^2-1.087\abs\zeta^2$, and therefore $\abs{\hat l^2+\mu^2}\ge\hat\kappa^2+\mu^2/2$ whenever $1.087\abs\zeta^2\le\mu^2/2$. Every momentum in $B_1B_2$ is a line or subdiagram momentum, so $\abs\zeta\le c_\Gamma\max_j\abs{\eta_j}\le c_\Gamma\eta_0$. Since $a<a_0\le1$ and $c_\Gamma\eta_0\le1$, also $\abs{a\zeta_\mu}\le1$. Hence $\abs{B_1B_2(k+i\eta)}\ge\tilde B(k)$, where $\tilde B$ is $B_1B_2$ with every mass $\mu$ replaced by $\mu/\sqrt2$.
\emph{Conclusion.} $\tilde F:=V_\eta/\tilde B\in\mathcal F_{\rm R}$, with the same natural set of line momenta, and masses do not enter degrees. So $3\dim H+\mathrm{degr}_H\tilde F\le3\dim H+\mathrm{degr}_HV-\mathrm{degr}_H(B_1B_2)\le-1$ by the numerator estimate and (b). Since $\abs{R_\Gamma(k+i\eta)}\le\tilde F(k)$, (a) gives the bound, uniformly for $a<a_0$ and $\max_j\abs{\eta_j}\le\eta_0$.
\end{proof}

\begin{remark}\label{app:G1:reisz-scope}
(i) For King's lattice Feynman rules, the degree of a one-particle-irreducible (sub)diagram is $\omega=3-\frac12(E_\varphi+E_A)-o_V-2v_2$, and every vertex has Reisz vertex dimension below $3$. The model is therefore super-renormalisable, and only finitely many classes have $\omega\ge0$.
(ii) What Proposition~\ref{prop:ch-CN} uses on the tori $T_j$ goes beyond Proposition~\ref{app:G1:reisz3d}: the Poisson resummation with a strip shift, the torus Ward defect at $k=0$, the graph of order $e^2$ that is logarithmic overall, and the identification of King's counterterm insertions with Reisz's forests. These are treated in Appendix~\ref{app:G1:part-D5} (Proposition~\ref{app:G1:CN} and Lemmas~\ref{app:G1:WP-R3transfer}, \ref{app:G1:WP-q2}).
\end{remark}

\section{Ba{\l}aban's Proposition~1 on the torus}\label{app:G1:area-T6}

This appendix proves the statements about \cite[Prop.~1]{Ba3} that Subsection~\ref{ssec:T6-prop1} uses: the anchoring of the vertices and the depth of the anchored classes in the proof of Corollary~\ref{cor:prop1R4}, the depth of the field $\tilde A$ at the places where the proposition is applied, and the list of sketched places behind Proposition~\ref{prop:prop1}(a) (Subsection~\ref{app:G1:part-E1}); the bound on the kernel \cite[(1.16)]{Ba3} and on domain differences of propagators in the mixed norm used in steps S1 and S2 (Subsection~\ref{app:G1:part-E2}); and the weighted $\ell^2$ and kernel bounds on which they rest (Subsection~\ref{app:G1:part-TWK}). Throughout $L=2$, distances are in scale-$k$ units, and the notation is that of Subsections~\ref{ssec:UV-prop1} and \ref{ssec:T6-prop1}.

\subsection{Anchoring and depth in \texorpdfstring{Corollary~\ref{cor:prop1R4}}{the support corollary}; the sketched places}\label{app:G1:part-E1}

\paragraph{Conventions.}
Distances are in the $\ell^\infty$ torus metric \cite[(1.3), p.~604]{Ba1}, in scale-$k$ units; $r_j:=r(t_j)$, and $r:=r_k$ at step $k$. Since $r(t)$ decreases in $t$, $r_{k-1}\ge r_k$. We use three nesting facts for the layers $\Lambda_i^{(j)}$ of \cite[(2.7)--(2.8), p.~558; p.~566]{Ba2}.
\begin{itemize}
\item A point of $\Lambda_{i+s}^{(j)}$ has distance $>s(r_j-1)$ from $(\Lambda_i^{(j)})^c$, in scale-$j$ units.
\item Passing from block sets to the point sets $B^j(\cdot)$ costs at most $1$ per set.
\item One scale-$(k-1)$ unit is $1/L$ scale-$k$ units.
\end{itemize}
Under a set-distance convention other than the infimum, each nesting step costs at most $M_b$ more.
We write $M_\theta$ for the thickness, in scale-$k$ units, of the slice on which $\theta_{k+1}$ goes from $1$ to $0$ \cite[p.~567]{Ba2}; $M_\theta\le LM_b$.
For a vertex $v$ of a class, its \emph{vertex point} $x_v$ depends on the vertex type.
\begin{itemize}
\item \cite[(1.6), (1.7)]{Ba3}: the summation point.
\item \cite[(1.8)--(1.11)]{Ba3}: an end point of a bond on which a factor $g_kA'$ or $\tilde A$ is non-zero.
\item \cite[(1.13)--(1.15)]{Ba3}: a point of the block $B^k(y)$ of the vertex.
\end{itemize}
The applications of \cite[Prop.~1]{Ba3} in \cite{Ba2}, at which the cube step P12 is used in the proof of Theorem~\ref{thm:T6}, are the following; we list them with their outer domains $\Omega_{\rm out}$ and keep their names.
\begin{itemize}
\item U0/U0$'$: Proposition~1 at step $k$, $\Omega_{\rm out}=B^k(\Lambda_2^{(k-1)\prime})$ \cite[p.~411]{Ba3} (also \cite[p.~582, (2.116)]{Ba2}).
\item U1a: \cite[p.~579, (2.106)--(2.107)]{Ba2}, $\Omega_{\rm out}=B^k(\Lambda_0^{(k)})$, with $\supp g_k\subset\Lambda_7^{(k)}$.
\item U1b: \cite[p.~580, Prop.~2.6]{Ba2}, $\Omega_{\rm out}=B^{k-1}(\Lambda_2^{(k-1)})$, with the terms localised in $B^k(\Lambda_7^{(k)})$.
\item U2: \cite[p.~578, (2.104)]{Ba2}, at scale $j\le k-1$, $\Omega_{\rm out}=B^j(\Lambda_2^{(j)})$.
\end{itemize}
The supports of $\tilde A$ are examined at three sites: (P)$=$U0, (Q)$=$U1a with $\tilde A=\tau\theta_{k+1}A'^{(k)}$, and (Z)$=$U2.

\paragraph{Anchoring.}
\begin{lemma}[every vertex is anchored]\label{app:G1:anchor}
Read \cite[(1.4)]{Ba3} as the definition of $E_k$, and \cite[(1.8)--(1.15)]{Ba3} as its expansion. Then every vertex $v$ of every class in \cite[(1.33)]{Ba3} has a vertex point $x_v$ with
\[
x_v\in\Omega_1\qquad\text{or}\qquad \dist\big(x_v,\ \supp g_k\cup\supp\tilde A\big)\le1.
\]
Since $\supp g_k\subset\Omega_1$, every vertex point lies within distance $1$ of $\Omega_1\cup\supp\tilde A$. In particular, every connected component of a class, after the regrouping S10 and the localisation S11 of Table~\ref{tab:prop1}, consists of anchored vertices only. This is the statement ``every vertex of a class is anchored'' used in the proof of Corollary~\ref{cor:prop1R4}.
\end{lemma}
\noindent The reading of \cite[(1.4)]{Ba3} as the definition of $E_k$ is the convention R1 of Lemma~\ref{app:G1:I3r-R123}; as printed, \cite[(1.14), (1.15)]{Ba3} carry no factor $g_k$.

\begin{proof}
(1) \emph{Scalar vertices.} The vertices \cite[(1.6), (1.7)]{Ba3} are sums over $x\in\Omega_1$ (p.~413). At p.~419 the text says: ``We localize the vertices (1.6) and (1.7) restricting summations to the set $\Omega_1$''. So $x_v\in\Omega_1$.

(2) \emph{Gating of $A'$.} In \cite[(1.4)]{Ba3} the field $A'$ enters the integrand only through the combination $e'g_kA'+A^{(k)}$, in the kernel $T$ and in the covariant Laplacian. The counterterms $\delta m^2(e',g_k,\lambda',\Omega_1,x)$ and $E_1(e',g_k,\lambda',\Omega_1)$ contain no $A'$. Hence every $A'$ produced by $\partial_{e'}$ at $e'=0$ in \cite[(1.5)]{Ba3} carries the factor $g_k$, in the kernel $T$ as well. So the factors $A'(\Gamma^{(k)}_{y,x})^n$ of (1.14), (1.15), which are printed without $g_k$, are $(g_kA')(\Gamma^{(k)}_{y,x})^n$. The form in \cite{Ba2} of the same polynomial, $Q^{(k)}(B^k(\Lambda_0^{(k)});B^{(k+1),\eta},g_kA'^{(k)})$ \cite[(2.107)]{Ba2}, is gated in the same way.

(3) \emph{Bond vertices.} The vertices (1.8)--(1.11) (p.~413) are sums over bonds $b$ of terms carrying $(g_kA'_b)^n(\tilde A_b)^{n'}$ with $n+n'\ge1$. The $R$-vertices (1.9), (1.11) carry $(\tilde A_b)^{\nbar+1}$. A term vanishes unless $b$ meets $\supp g_k$ or $\supp\tilde A$. An end point of $b$ is at distance $\le\eta_k\le1$ from that support. The $n=0$, $n'\ge1$ vertices produced by the $\tilde A$-expansion of $G_k(\Omega,\tilde A+\tilde B)$ \cite[p.~414]{Ba3} carry $\tilde A$-legs and are covered by the same count.

(4) \emph{Block vertices.} By \cite[p.~414]{Ba3} the expressions (1.12)--(1.15) enter only as pairs at the same $y\in T_1^{(k)}\cap\Omega$, ``with the restriction that at least one of them has to be of the form (1.14) or (1.15)''; (1.13)--(1.15) are the vertices and (1.12) is an external field. A factor (1.14) with $n\ge1$ vanishes unless the contour $\Gamma^{(k)}_{y,x}$ has a bond in $\supp g_k$, by step (2). A factor with $n'\ge1$, and every factor (1.15), vanishes unless the contour has a bond in $\supp\tilde A$. The contours satisfy $\Gamma^{(k)}_{y,x}\subset B^k(y)$, bonds included (\cite[(2.1)--(2.2), p.~608]{Ba1}). The block $B^k(y)$ has $\ell^\infty$-diameter $<1$ in scale-$k$ units. Hence both members of the pair, and so the (1.13) partner as well, have their vertex points within distance $1$ of $\supp g_k\cup\supp\tilde A$.

(5) \emph{Components.} Steps (1)--(4) apply to every vertex, so they apply to every vertex of every connected component produced by S10 and S11.
\end{proof}

\paragraph{Depth of the $\Omega_1$-anchored classes.}
\begin{lemma}[$\Omega_1$-anchored classes at each use]\label{app:G1:Omega1}
Let $L=2$, $R\ge R^{**}$ (so \eqref{eq:eppp} holds), and $r(t)\ge R$. Consider a class at one of the uses U0/U0$'$, U1a, U1b, U2, with outer domain $\Omega_{\rm out}$ and the cube $\square$ of P12 with $r_\square=r$. Suppose the class has an old vertex point $x_v$ anchored at $\Omega_1$ or at $\supp g_k$ in the sense of Lemma~\ref{app:G1:anchor}. Then $m:=\dist(x_v,\Omega_{\rm out}^c)$ satisfies
\[
\begin{gathered}
\text{U0/U0}':\ m\ge2r_k-4;\qquad \text{U1a}:\ m\ge7r_k-10;\\
\text{U1b}:\ m\ge\tfrac5L(r_{k-1}-1)-2\ge\tfrac52r_k-5;\qquad \text{U2}:\ m\ge5r_j-7,
\end{gathered}
\]
where at U2 the scale is $j$ and $r=r_j$. In every case
\begin{equation}\label{app:G1:eq-need}
m>\tfrac74r+2s_b+\ell_{\Nbar},
\end{equation}
hence $\square\subset\Omega_{\rm out}$ and $\dist(\square,\Omega_{\rm out}^c)>r/4$. In particular $m$ exceeds the need $\frac32r+2s_b+\ell_{\Nbar}$ of P12. This is the case of a vertex anchored at $\supp g_k$ in the proof of Corollary~\ref{cor:prop1R4}.
\end{lemma}
\noindent The lemma is stated at $L=2$. At $L=5$ the margin at U1b is $r_{k-1}-3$, so $\square\subset\Omega_{\rm out}$ can fail for such classes.

\begin{proof}
\emph{Margins.} (i) U1a: $x_v$ is within $1$ of $\supp g_k\subset\Lambda_7^{(k)}$ \cite[p.~579]{Ba2}, and $\Omega_{\rm out}=B^k(\Lambda_0^{(k)})$. Seven nesting steps give $m\ge7(r_k-1)-2-1=7r_k-10$.
(ii) U1b: $x_v\in B^k(\Lambda_7^{(k)})$, since the terms of \cite[Prop.~2.6, p.~580]{Ba2} are those localised in $B^k(\Lambda_7^{(k)})$. Moreover $\Lambda_7^{(k)}\subset\Lambda_0^{(k)}\subset\Lambda_7^{(k-1)\prime}$ \cite[p.~566]{Ba2}. Five nesting steps at scale $k-1$, converted to scale $k$, give $m\ge\frac5L(r_{k-1}-1)-2\ge\frac52r_k-5$ at $L=2$, since $r_{k-1}\ge r_k$. The terms of Prop.~2.6 localised outside $\Lambda_7^{(k)}$ are estimated with the domain of U0 and are covered by U0.
(iii) U2: $x_v\in B^k(\Lambda_0^{(k)})\subset B^j(\Lambda_7^{(j)})$, by iterating the inclusions of \cite[p.~566]{Ba2}. Five nesting steps at scale $j$ give $m\ge5r_j-7$.
(iv) U0/U0$'$: here $\Omega_{\rm out}=B^k(\Lambda_2^{(k-1)\prime})$, and the nesting of the layers of \cite[p.~566]{Ba2} gives $m\ge2r_k-4$ for every vertex type. Its dependence on the support of $\tilde B^{(k)}$ is removed by $\supp\tilde B^{(k)}\subset\supp\theta_{k+1}$ \cite[(2.89), p.~575]{Ba2}.

\emph{Arithmetic.} By \eqref{eq:eppp} and $r\ge R$ we have $r/4\ge\ell_{\Nbar}+8M_b+\frac{25}4$, and $s_b\le2M_b$. Each case of \eqref{app:G1:eq-need} reduces to one inequality:
\begin{itemize}
\item U0: $2r-4>\frac74r+2s_b+\ell_{\Nbar}$ iff $\frac r4>\ell_{\Nbar}+2s_b+4$, which follows from $\frac r4>\ell_{\Nbar}+4M_b+4$.
\item U1a: $\frac{21}4r>\ell_{\Nbar}+2s_b+10$.
\item U1b: $\frac34r>\ell_{\Nbar}+2s_b+5$.
\item U2: $\frac{13}4r_j>\ell_{\Nbar}+2s_b+7$.
\end{itemize}
All four follow from the bound on $r/4$.
The form $\frac52(r(t_{k-1})-M_b-1)-4$ used in the proof of Corollary~\ref{cor:prop1R4} is the U0 configuration $\Omega_1=B^k(\Lambda_7^{(k-1)\prime})\subset\Omega=B^k(\Lambda_2^{(k-1)\prime})$. It also satisfies \eqref{app:G1:eq-need}: with $r(t_{k-1})\ge r$, this reduces to $\frac34r>\ell_{\Nbar}+2s_b+\frac52M_b+\frac{13}2$, which follows from \eqref{eq:eppp}.

\emph{Cube.} By Proposition~\ref{prop:prop1}(c) every $p\in\square$ satisfies $\dist(p,x_v)\le\frac32r+2s_b+\ell_{\Nbar}$. By \eqref{app:G1:eq-need}, $\dist(p,\Omega_{\rm out}^c)>m-\frac32r-2s_b-\ell_{\Nbar}>r/4$.
\end{proof}

\paragraph{Supply at the uses.}
\begin{lemma}[supply of $\tilde A$-depth at the uses]\label{app:G1:supply}
Let $L=2$. At each of the sites (P), (Q), (Z) of the proof of Theorem~\ref{thm:T6}, with $r$ the value of $r(\cdot)$ at the scale of the site,
\begin{equation}\label{app:G1:eq-supply}
\dist(\supp\tilde A,\partial\Omega)>2r-M_\theta-2M_b-5.
\end{equation}
If \eqref{eq:eppp} holds and $r\ge R$, the right side exceeds $\frac74r+2s_b+\ell_{\Nbar}+1$, so (R4$'''$) \eqref{eq:R4ppp} holds at every use. This is the last statement of Corollary~\ref{cor:prop1R4}. It also gives $\dist(\supp\tilde A,\partial\Omega)\ge r$, the depth needed at the changes of domain.
\end{lemma}

\begin{proof}
(Q) Here $\Omega=B^k(\Lambda_0^{(k)})$ and $\tilde A=\tau\theta_{k+1}A'^{(k)}$ \cite[p.~579]{Ba2}. By \cite[p.~567]{Ba2}, $\theta_{k+1}=1$ on $B^k(\Lambda_2^{(k)})$ and varies on a slice of thickness $<M_\theta$ outside it. So $\supp\tilde A$ lies within $M_\theta$ of $B^k(\Lambda_2^{(k)})$. Two nesting steps give $\dist(B^k(\Lambda_2^{(k)}),\Omega^c)\ge2r_k-4$. With one more lattice unit for the bond, $\dist(\supp\tilde A,\partial\Omega)>2r_k-M_\theta-5$ under the infimum convention. Under any set-distance convention it is $>2r_k-M_\theta-2M_b-5$, since two nesting steps cost at most $2M_b$ more.

(P) Here $\Omega_P=B^k(\Lambda_2^{(k-1)\prime})\supset B^k(\Lambda_0^{(k)})=\Omega_Q$, because $\Lambda_0^{(k)}\subset\Lambda_7^{(k-1)\prime}\subset\Lambda_2^{(k-1)\prime}$ \cite[p.~566]{Ba2}. The field is $\tilde A=\theta_{k+1}A'^{(k)}+\tilde B^{(k)}$ \cite[(2.93)]{Ba2}. Each of the three terms defining $\tilde B^{(k)}$ in \cite[(2.89), p.~575]{Ba2} carries the factor $\theta_{k+1}$, so $\supp\tilde A\subset\supp\theta_{k+1}$, the same set as at (Q). Since $\Omega_Q\subset\Omega_P$, the margin at (P) is at least the margin at (Q).

(Z) At scale $j$ the margin is $\ge5r_j-8$, by the nesting argument of part (iii) of the proof of Lemma~\ref{app:G1:Omega1} with one more lattice unit for the bond, and $\ge5r_j-8-7M_b$ under any convention. It exceeds $2r_j-M_\theta-2M_b-5$, because $3r_j\ge3+5M_b$.

\emph{Arithmetic.} We have $\frac74r+2s_b+\ell_{\Nbar}+1<2r-M_\theta-2M_b-5$ iff $\frac r4>\ell_{\Nbar}+2s_b+M_\theta+2M_b+6$. With $s_b\le2M_b$ and $M_\theta\le LM_b$, the right side is at most $\ell_{\Nbar}+(6+L)M_b+6=\ell_{\Nbar}+8M_b+6$. By \eqref{eq:eppp} and $r\ge R$, $\frac r4\ge\ell_{\Nbar}+8M_b+\frac{25}4$. Finally, $2r-M_\theta-2M_b-5\ge r$ iff $r\ge(2+L)M_b+5$, which also follows from \eqref{eq:eppp}.

\emph{Completeness of the list of sites.} The sites (P), (Q), (Z) are all the places in the proof of Theorem~\ref{thm:T6} at which \cite[Prop.~1]{Ba3} is applied with a field $\tilde A$; this is read off \cite[pp.~583--594]{Ba2} and \cite[pp.~412--445]{Ba3}.
\end{proof}

\begin{remark}\label{app:G1:footing}
The supply \eqref{app:G1:eq-supply} is the margin Ba{\l}aban's own nesting provides at \cite[p.~579]{Ba2} (Remark~\ref{rem:prop1-cited}(iii)). The need, through P12, is only $\frac32r+O(\ell_{\Nbar}+M_b)$. The margin is therefore available, not needed.
\end{remark}

\paragraph{The sketched places of Proposition~\ref{prop:prop1}(a).}
\begin{lemma}[the sketched places of \cite{Ba3}, Prop.~1]\label{app:G1:inventory}\leavevmode
\begin{enumerate}
\item[(i)] On pages 412, 414, 415 and 419--445 of \cite{Ba3} and pp.~610--611 of \cite{Ba1}, the places of the proof where the text calls an argument similar, omitted, standard, obvious, easy, the same as or done in the same way as another, repeated, analogous or straightforward, or leaves out details, are the places S1--S21 of Table~\ref{tab:prop1}, with the number 12 used by the printed step P12. That is twenty sketched places, S1--S11 and S13--S21, plus P12.
\item[(ii)] Besides these, pages 414--445 of \cite{Ba3} contain seven further sketched places:
 \begin{enumerate}
 \item[(1)] the $\Omega_2$ branch ``[or $G_k(\Omega_2,\tilde B)$]'' of P12 (p.~433);
 \item[(2)] the identification, unwritten in \cite{Ba3}, of the torus propagators of the counterterms (pp.~418--419) with the $\eta_k\Z^d$ propagators of the classes, which the exact cancellations of \cite[(3.7), (3.19)]{Ba3} and of pp.~444--445 need;
 \item[(3)] p.~433: ``We repeat some of the operations done before \dots'', i.e.\ S1 and S9 again on $\square$, with $B'$ in place of $\tilde A$;
 \item[(4)] pp.~430--431: ``The same applies to the remaining definitions'' and ``Taking other functions $F$, or differentiating (2.24) to higher order in $A$, we can get all necessary Ward--Takahashi identities'';
 \item[(5)--(7)] three phrases of power counting and bookkeeping: p.~427 ``Now it is easy to see that the remaining powers of $L^j\eta$ \dots'', p.~428 ``In the case $v=v'$ we make the same operations'', and p.~433 ``Considering each connected component of $G_{\rm ren}$ separately \dots We can extract also all coupling constants''.
 \end{enumerate}
\item[(iii)] \emph{Depth of the seven places.}
 \begin{itemize}
 \item (1) needs $\square\subset\Omega_2$, i.e.\ the need of P12 measured from $\Omega_2^c$ (Proposition~\ref{prop:prop1}(c) with $\Omega_2$ in place of $\Omega$). It has the same coefficient of $r$ as P12 and is not a new item.
 \item (2) is torus-size, not depth: the differences are torus-image terms (Lemma~\ref{lem:IMG}).
 \item (3) takes place on the parallelepiped $\square$, where \cite[p.~611]{Ba1} imposes no restriction. Its coupling power uses only $\abs{B'}\le O(r_\square p)$, a size bound.
 \item (4)--(7) are Gaussian lattice identities or bookkeeping on \cite[(2.14)]{Ba3}, with no domain.
 \end{itemize}
 So none of the seven needs depth beyond that of P12 or $R_0+O(1)$, and no sketched place needs depth in $\big(\frac74r+O(\ell_{\Nbar}+M_b),\,2r-M_\theta-2M_b-5\big]$.
\end{enumerate}
This is the list of sketched places and the depth claim of Proposition~\ref{prop:prop1}(a).
\end{lemma}
\noindent Parts (i) and (ii) are statements about the printed text of \cite{Ba3}, not mathematical statements; the completeness of the list is what Proposition~\ref{prop:prop1}(a) uses.

\begin{proof}
Parts (i) and (ii) are read off the printed text. For each place of (ii) the depth claim of (iii) is read off the cited page, as follows.
\begin{itemize}
\item (1): the cube step is printed for $G_k(\Omega_2,\tilde B)$ as well. By Proposition~\ref{prop:prop1}(c) with $\Omega_2$ in place of $\Omega$, it needs an old vertex at depth $>\frac32r_\square+2s_b+\ell_{\Nbar}$ in $\Omega_2$.
\item (2): \cite[(1.25)--(1.31), pp.~418--419]{Ba3} builds the counterterms from propagators on $T_\eps$ rescaled to $T_{\eta}$. After p.~433 the classes use $G_k(0)$ on $\eta_k\Z^d$. The difference consists of image terms of order $\mathrm{e}^{-\delta_0\cdot(\text{torus side})}$ and does not involve $\Omega$.
\item (3): on $\square$ the bounds of \cite[Prop.~2.1]{Ba1} hold for all points (\cite[p.~611]{Ba1}).
\item (4)--(7): the Ward--Takahashi identities are identities of Gaussian measures with no domain, and the remaining phrases act on the power counting of \cite[(2.14)]{Ba3}.
\end{itemize}
Every object on pp.~433--445 is independent of the distance to $\Omega^c$ once $\square\subset\Omega$ (Lemma~\ref{lem:Omegafree}). The places (2)--(7) therefore need no depth, and (1) needs that of P12 with $\Omega_2$ in place of $\Omega$.
\end{proof}

\begin{remark}[scope]\label{app:G1:anchor-missing}
(i) The renormalisation combinatorics S21 of \cite[pp.~434--445]{Ba3}, and the generalisation S6 (``easily seen'' in print), are used as printed; they do not depend on the depth in $\Omega$. (ii) Lemmas~\ref{app:G1:Omega1} and \ref{app:G1:supply} are stated at $L=2$; at $L=5$ the margin at U1b is too small.
\end{remark}

\subsection{The $L^2$ sandwich and the two-domain theorem}\label{app:G1:part-E2}

We prove two bounds used in the proof of Proposition~\ref{prop:prop1}: the bound for the kernel \cite[(1.16)]{Ba3} used in step S1 (Proposition~\ref{app:G1:sandwich}), and the bound for the two-variable mixed norm of a domain difference $\delta G_k$ used in step S2 and in part (c) (Proposition~\ref{app:G1:twodomain}).

\emph{Setting and notation.} We fix a scale $k\ge1$ and work in scale-$k$ units, as in Subsection~\ref{ssec:UV-prop}. We write $\eta:=\eta_k=L^{-k}\le\frac12$, $t:=t_k=L^k\eps$, $e_k:=e(t_k)$, and $d=3$.
\begin{itemize}
\item \emph{Domains.} $\mathcal X$ is a union of big blocks of the rescaled torus $T_\eta$ with Neumann conditions, or $\mathcal X=T_\eta$ itself. For real $A$ we put $H_{\mathcal X}(A):=-\Delta^{\eta,N}_{A,\mathcal X}+m^2t^2+a_kP_k(A)$ and $G_k(\mathcal X,A):=H_{\mathcal X}(A)^{-1}$. Here $P_k=Q_k^*Q_k$, and $Q_k$ averages over unit blocks along contours inside the block, of length $\le d$.
\item \emph{Kernels and derivatives.} Kernels are taken with respect to $\eta^3$, i.e.\ $(Yf)(x)=\sum_y\eta^3Y(x,y)f(y)$. $D=D^{\eta}_{\tilde B}$ is the covariant forward difference and $D^*$ its $\ell^2(\eta^3)$-adjoint, so that $D^{b*}$ acts on the second variable. We put $\mathfrak Df:=(f,D_1f,D_2f,D_3f)$.
\item \emph{Depth and neighbourhoods.} $\rho_{\mathcal X}(z):=\dist_\infty(z,\mathcal X^c)$ is the depth of $z$ in $\mathcal X$ (the $\ell^\infty$ torus distance, which is $1$-Lipschitz), and $U_s(Y)$ denotes the $\ell^\infty$-neighbourhood of radius $s$ of a set $Y$.
\item \emph{Kernel classes.} For $0<\kappa<3$ and $\delta>0$ let $K^\delta_\kappa(z):=\max(\abs z,\eta)^{\kappa-3}\mathrm{e}^{-\delta\abs z}$. Further, $K^\delta_3(z):=(1+\log_+(1/\max(\abs z,\eta)))\mathrm{e}^{-\delta\abs z}$, and $K^\delta_\kappa(z):=\mathrm{e}^{-\delta\abs z}$ for $\kappa>3$.
\item \emph{Smallness parameters.} $\sigma:=e_k(\norm{\tilde A}_\infty+\norm{\partial^\eta\tilde A}_\infty)$ and $\theta(A):=e_k\norm{\partial^\eta A}_\infty$.
\end{itemize}
Here $\tilde A$ and $\tilde B$ are the real background fields of Subsection~\ref{ssec:UV-prop1}. Under the hypotheses of \cite[p.~412]{Ba3} one has $\sigma\le2c_1e_kp(t)$, $\theta(\tilde B)\le c_1e_kp(t)$ and $\theta(\tilde A+\tilde B)\le2c_1e_kp(t)$.

\emph{Norm.} For a two-variable kernel $Y$, put $Y_{ab}:=D^aYD^{b*}$ ($a,b\in\{0,1\}$, every direction), and for $0<\alpha<1$, $\beta\ge0$ define
\[
\mathfrak N_{\alpha,\beta}(Y):=\max_{a,b}\sup_{x,x'}\mathrm{e}^{\beta\abs{x-x'}}\max\{\norm{Y_{ab}(x,x')},\,a\,H_x,\,b\,H_{x'},\,ab\,H_{xx'}\}.
\]
Here $H_x:=\sup_{0<\abs{x-\tilde x}\le1}\abs{x-\tilde x}^{-\alpha}\norm{U(\tilde B(\Gamma_{\tilde x,x}))Y_{ab}(x,x')-Y_{ab}(\tilde x,x')}$; $H_{x'}$ is the same quotient in $x'$, with the transport acting from the right; and $H_{xx'}$ is the mixed second difference divided by $\abs{x-\tilde x}^\alpha\abs{x'-\tilde x'}^\alpha$. $\norm{hYh'}_{1,\alpha}$ denotes the same quantity at $\beta=0$ for the kernel $h(x)Y(x,x')h'(x')$. This is \cite[(1.32)]{Ba3} ``extended in a natural way'' to two variables, with the mixed component $D_xD^*_{x'}$, which is the reading under which \cite[(2.5)]{Ba3} is used. Quotients with $\abs{x-\tilde x}>1$ are bounded by twice the corresponding sup and need no separate estimate.

\begin{lemma}[inputs]\label{app:G1:WKinputs}
There are constants $\gamma_0,\beta_*,K_W,c_K,c_K(\alpha),\delta_K,\vartheta_R>0$, depending only on $d,a,M_b,L$ (and on $\alpha$ for $c_K(\alpha)$), such that the following hold whenever $A$ is real with $\theta(A)\le\vartheta_R$ and $0\le m^2t^2\le c$:
\begin{enumerate}
\item[(H0)] $H_{\mathcal X}(A)\ge\gamma_0$ on every $\mathcal X$ as above \cite[(1.8), p.~573]{Ba4}.
\item[(W)] For every $1$-Lipschitz $\rho$ and $\abs\beta\le\beta_*$, the operators $G$, $DG$, $GD^*$ and $DGD^*$ with $G=G_k(\mathcal X,A)$, conjugated as $Y\mapsto \mathrm{e}^{\beta\rho}Y\mathrm{e}^{-\beta\rho}$, have $\ell^2(\eta^3)$ operator norms $\le2/\gamma_0$ (for $G$) and $\le K_W$ (the other three). These bounds hold uniformly in $\eta$, $k$, $\eps$, the shape of $\mathcal X$ (wrapping allowed) and the boundary condition.
\item[(K)] For $x,w$ with $\rho_{\mathcal X}(x),\rho_{\mathcal X}(w)\ge R_0$ (and H\"older partners $\abs{x-\tilde x}\le1$ with the same property), where $R_0$ is the constant of Proposition~\ref{prop:UV-BaProp21}:
\begin{itemize}
\item[(K0)] $\norm{G(x,w)}\le c_KK^{\delta_K}_2(x-w)$;
\item[(K1)] $\norm{D_xG(x,w)}\le c_KK^{\delta_K}_1(x-w)$;
\item[(K2)] $\norm{U(A(\Gamma_{\tilde x,x}))D_xG(x,w)-D_{\tilde x}G(\tilde x,w)}\le c_K(\alpha)\abs{x-\tilde x}^\alpha\big(K^{\delta_K}_{1-\alpha}(x-w)+K^{\delta_K}_{1-\alpha}(\tilde x-w)\big)$.
\end{itemize}
On $\mathcal X=T_\eta$ the restriction $\rho_{\mathcal X}\ge R_0$ is void.
\item[(S)] For real $A$, $G_k(\mathcal X,A)$ is real symmetric, so $(GD^*)(x,w)=(D_wG(w,x))^T$.
\end{enumerate}
\end{lemma}
(W) is the weighted form of \cite[Cor.~2.3, (2.30), p.~580]{Ba4}, which bounds $\abs{\langle f,DGD^*f'\rangle}$ by $c_0\mathrm{e}^{-\delta_0\dist(\supp f,\supp f')}\norm f_2\norm{f'}_2$ without restriction on the supports. (K) is \cite[(2.10)--(2.11), p.~426]{Ba3}, summed over the slices of \cite[(2.6)]{Ba3}; the derivation of (2.10)--(2.11) is part (b), step S3, of the proof of Proposition~\ref{prop:prop1}. (H0) is cited, (S) is immediate, and the proofs of (W) and of the summed form (K) are Lemmas~\ref{app:G1:WK-W} and \ref{app:G1:WK-K} of Subsection~\ref{app:G1:part-TWK}.

\paragraph{The kernel (1.16): an $L^2$ sandwich.}\label{app:G1:ssec-sandwich}
Let $G:=G_k(\mathcal X,\tilde B)$ and $G':=G_k(\mathcal X,\tilde A+\tilde B)$. Put $V:=G^{-1}-G'^{-1}$; this is Ba{\l}aban's $V_k(\tilde A,\tilde B)$, the bracket of \cite[(3.44), p.~619]{Ba1}. For $n,n'\ge0$, put $m:=n+n'$ and
\[
K_{n,n'}:=[GV]^nG'[VG]^{n'}=L_0VL_1V\cdots VL_m,\qquad L_n:=G',\ L_i:=G\ (i\ne n),
\]
which is the kernel \cite[(1.16), p.~414]{Ba3}. Our hypothesis is
\[
\text{(R$_\vartheta$)}\qquad \theta(\tilde B),\ \theta(\tilde A+\tilde B)\le\vartheta_R,\qquad \sigma\le1,\qquad 0\le m^2t^2\le c.
\]
Under the hypotheses of \cite[p.~412]{Ba3}, (R$_\vartheta$) follows from the regularity condition $2c_1e_kp(t)\le\min\{\vartheta_R,1\}$ on the coupling.

\begin{lemma}[the vertex operator]\label{app:G1:vertex}
Put $\phi_\mu(x):=\eta^{-1}(1-U_{x,\mu}(\tilde A)^{-1})$, $\psi_\mu:=\phi_\mu^*\phi_\mu$ and $\Pi:=P_k(\tilde A+\tilde B)-P_k(\tilde B)$. Then:
\begin{enumerate}
\item[(i)] $V=-\big[\sum_\mu(D_\mu^*\phi_\mu+\phi_\mu^*D_\mu+\psi_\mu)+a_k\Pi\big]=-\mathfrak D^*\mathfrak E\mathfrak D$, where $\mathfrak E$ is the $4\times4$ matrix of multiplication and block operators with entries $\mathfrak E_{00}=\sum_\mu\psi_\mu+a_k\Pi$, $\mathfrak E_{0\mu}=\phi_\mu^*$, $\mathfrak E_{\mu0}=\phi_\mu$, all other entries $0$.
\item[(ii)] (Right-pushed form.) $V=-(R_0+\sum_\mu R_{1,\mu}D_\mu)=:-\mathfrak r\mathfrak D$. Here $(R_{1,\mu}h)(x):=\phi_\mu(x)^*h(x)-\phi_\mu(x-\eta e_\mu)U_{x-\eta e_\mu,\mu}(\tilde B)^{-1}h(x-\eta e_\mu)$, and $R_0:=\sum_\mu(\psi_\mu-M_\mu)+a_k\Pi$, with $M_\mu$ the multiplication by $(\partial^\eta_\mu\phi_\mu)(x-\eta e_\mu)$.
\item[(iii)] (Left-pushed form.) $V=V^*=-\mathfrak D^*\mathfrak r^*$.
\item[(iv)] For $\sigma\le1$, in $\ell^p(\eta^3)$, $p\in\{1,2,\infty\}$: $\norm{\mathfrak r},\norm{\mathfrak E}\le c_V\sigma$ with $c_V:=4d+a(2d+d^2)$. Every kernel has range $\le\sqrt3+\eta$, so for $1$-Lipschitz $\rho$ and $\abs\beta\le1$ the conjugates by $\mathrm{e}^{\pm\beta\rho}$ are bounded by $\mathrm{e}^{\beta(\sqrt3+\eta)}c_V\sigma$. Moreover $V(y,\cdot)=0$ unless $y\in U_{\sqrt3+\eta}(\supp\tilde A)$.
\end{enumerate}
On $\mathcal X=\Omega$ with Neumann conditions, (i)--(iv) hold as stated provided $\dist(\supp\tilde A,\Omega^c)>\sqrt3+2\eta$, since then every bond and every block on which $\phi$ or $\Pi$ is non-zero lies in $\Omega$.
\end{lemma}

\begin{proof}
All transporters are functions of the matrix $q$ ($q^2=-1$), so they commute, and $U(\tilde A+\tilde B)=U(\tilde A)U(\tilde B)$. We use $D_{A,\mu}f(x)=\eta^{-1}(U_{x,\mu}(A)f(x+\eta e_\mu)-f(x))$ \cite[(1.3)]{Ba4}. Then
\[
D_{\tilde A+\tilde B,\mu}f=U_\mu(\tilde A)\big(D_\mu f+\phi_\mu f\big).
\]
Since $U$ is pointwise orthogonal, summing over the bonds of $\mathcal X$ (the same bonds for both backgrounds) gives
\[
-\Delta_{\tilde A+\tilde B}=\sum_\mu(D_\mu+\phi_\mu)^*(D_\mu+\phi_\mu)=-\Delta_{\tilde B}+\sum_\mu(D^*_\mu\phi_\mu+\phi^*_\mu D_\mu+\psi_\mu).
\]
Since $Q_k(\tilde A+\tilde B)=Q_k(\tilde B)+F_2$ with $F_2:=Q_k(\tilde A+\tilde B)-Q_k(\tilde B)$, we get $P_k(\tilde A+\tilde B)-P_k(\tilde B)=\Pi$. The terms $m^2t^2$ cancel, and (i) follows.
For (ii), use $D^*_\mu h(x)=\eta^{-1}(U^{-1}_{x-\eta e_\mu}h(x-\eta e_\mu)-h(x))$ and $[\phi,U]=0$:
\[
D^*_\mu(\phi_\mu f)(x)=-\phi_\mu(x-\eta e_\mu)U^{-1}_{x-\eta e_\mu}(D_\mu f)(x-\eta e_\mu)-(\partial^\eta_\mu\phi_\mu)(x-\eta e_\mu)f(x).
\]
Inserting this into (i) gives (ii). (iii) holds because both inverses are real symmetric for real backgrounds.
For (iv): $\abs{\mathrm{e}^{q\vartheta}-1}\le\abs\vartheta$, hence $\abs{\phi_\mu}\le\sigma$, $\abs{\psi_\mu}\le\sigma^2$ and $\abs{\partial^\eta_\mu\phi_\mu}\le e_k\abs{\partial^\eta_\mu\tilde A_\mu}\le\sigma$. Also $\abs{U(\tilde A(\Gamma))-1}\le d\sigma$ for the contours, $\norm{Q_k}\le1$, and so $\norm\Pi\le(2d+d^2)\sigma$. This gives $\norm{R_0}\le(2d+a(2d+d^2))\sigma$, $\sum_\mu\norm{R_{1,\mu}}\le2d\sigma$ and $\norm{\mathfrak E}\le c_V\sigma$. A kernel of range $r$ changes by at most $\mathrm{e}^{\abs\beta r}$ under conjugation by $\mathrm{e}^{\beta\rho}$. The support statement holds because the Laplacian part of $V$ lives on bonds where $\tilde A\ne0$, and $\Pi$ on unit blocks meeting $\supp\tilde A$.
\end{proof}

\begin{lemma}[convolution and smearing]\label{app:G1:conv}
(a) For $\kappa_1,\kappa_2\in(0,3)$ and $0<\delta'<\delta\le1$, on $T_\eta$ or $\eta\Z^3$:
\[
\sum_w\eta^3K^\delta_{\kappa_1}(z-w)K^\delta_{\kappa_2}(w)\le c(\kappa_1,\kappa_2)\big(1+(\delta-\delta')^{-3}\big)K^{\delta'}_{\kappa_1+\kappa_2}(z),
\]
with the conventions above for $\kappa_1+\kappa_2\ge3$.

(b) For $\delta''<\delta$: $K^\delta_\kappa\le c(1+(\delta-\delta'')^{-1})K^{\delta''}_{\kappa'}$ whenever $\kappa\ge\kappa'$.

(c) Let $R$ be a multiplication, a one-bond shift, or a bounded kernel supported on single unit blocks, with $\abs{R(z,z')}\le r_R$. Then $\big(\abs R\,K^\delta_\kappa(\cdot-y)\big)(z)\le c(\kappa)\mathrm{e}^{2\sqrt3\delta}r_RK^\delta_\kappa(z-y)$.
\end{lemma}
\begin{proof}
For (a), use $\mathrm{e}^{-\delta\abs{z-w}-\delta\abs w}\le \mathrm{e}^{-\delta'\abs z}\mathrm{e}^{-(\delta-\delta')(\abs{z-w}+\abs w)}$, put $r:=\max(\abs z,\eta)$, and split the sum into three regions: $\abs w\le r/2$; $\abs{z-w}\le r/2$; and both distances $\ge r/2$. Each region is bounded with the lattice-sum bound $\sum_{\abs w\le R}\eta^3\max(\abs w,\eta)^{\kappa-3}\le c(\kappa)R^\kappa$. The tail $\abs w>2r$ gives $c\int_{2r}^\infty s^{\kappa_1+\kappa_2-4}\mathrm{e}^{-(\delta-\delta')s}\,ds$. (b): for $r\le1$ the smaller exponent gives the larger kernel, and for $r>1$ the power is absorbed by $\mathrm{e}^{-(\delta-\delta'')r}$. (c): for a one-bond shift, $\max(\abs{z'-y},\eta)\ge\frac12\max(\abs{z-y},\eta)$. For a block kernel, $\abs{z'-y}\ge\abs{z-y}/2$ if $\abs{z-y}\ge2\sqrt3$, while the block sum is $\le c(\kappa)$ if $\abs{z-y}<2\sqrt3$.
\end{proof}

\begin{lemma}[end factors]\label{app:G1:ends}
Let $L_0,L_1,L_2\in\{G,G'\}$ and $a\in\{0,1\}$, and set $E^{(2)}_a:=D^aL_0VL_1$, $E^{(3)}_a:=D^aL_0VL_1VL_2$, $\delta':=\delta_K/2$, $0\le\beta\le\delta_K/4$. Assume that (K) applies at every pair of points used below (automatic on $T_\eta$). Then:
\begin{enumerate}
\item[(i)] $\norm{E^{(2)}_a(x,y)}\le c_E\sigma K^{\delta'}_2(x-y)$ and $\norm{E^{(3)}_a(x,y)}\le c_E'\sigma^2K^{\delta'}_3(x-y)$.
\item[(ii)] For $a=1$ and $\abs{x-\tilde x}\le1$, the transported difference of $E^{(2)}_1$ is bounded by $c_E(\alpha)\sigma\abs{x-\tilde x}^\alpha(K^{\delta'}_{2-\alpha}(x-y)+K^{\delta'}_{2-\alpha}(\tilde x-y))$, and that of $E^{(3)}_1$ by the same expression with $c'_E(\alpha)\sigma^2$ and $K_{3-\alpha}$.
\item[(iii)] Put $J_\kappa:=(\sum_y\eta^3K^{\delta'}_\kappa(y)^2\mathrm{e}^{2\beta\abs y})^{1/2}$. Then $\norm{\mathrm{e}^{\beta\abs{\cdot-x}}E^{(2)}_a(x,\cdot)}_2\le c_EJ_2\sigma$, the transported difference in (ii) has weighted $\ell^2$ norm $\le2\mathrm{e}^\beta c_E(\alpha)J_{2-\alpha}\sigma\abs{x-\tilde x}^\alpha$, and the same holds for $E^{(3)}$ with $J_3$, $J_{3-\alpha}$ and $\sigma^2$.
\end{enumerate}
$J_\kappa<\infty$ uniformly in $\eta$ and in the torus side if and only if $\kappa>\frac32$. Hence $J_{2-\alpha}<\infty$ if and only if $\alpha<\frac12$, and $J_{3-\alpha}<\infty$ for every $\alpha<1$. The constants depend only on $d,a,M_b,L,\alpha,\beta$.
\end{lemma}
\begin{proof}
(i) By Lemma~\ref{app:G1:vertex}(ii), $E^{(2)}_a=-(D^aL_0)R_0L_1-\sum_\mu(D^aL_0)R_{1,\mu}(D_\mu L_1)$. Thus every line carries exactly one derivative: $L_0$ the external one, $L_1$ the vertex derivative. By (K), $\norm{D^aL_0(x,z)}\le c_KK^{\delta_K}_{2-a}$ and $\norm{\mathfrak DL_1(z,y)}\le c_K(K_2+K_1)^{\delta_K}$. Lemma~\ref{app:G1:conv}(c) moves $R$ (with $r_R\le c_V\sigma$) onto the kernel, and Lemma~\ref{app:G1:conv}(a),(b) give $K_1*K_1\to K_2$. For $E^{(3)}_a$ both vertices are right-pushed, so $E^{(3)}_a=(D^aL_0)\mathfrak r(\mathfrak DL_1)\mathfrak r(\mathfrak DL_2)$ with one derivative per line, and $K_1*K_1*K_1\to K_3$.

(ii) The difference acts on the first factor only. Use (K2), then $K_{1-\alpha}*K_1\to K_{2-\alpha}$, respectively $K_{1-\alpha}*K_1*K_1\to K_{3-\alpha}$.

(iii) Square and sum, using $\abs{y-x}\le\abs{y-\tilde x}+1$ for the $\tilde x$-term. Near $0$, $\sum_y\eta^3\max(\abs y,\eta)^{2\kappa-6}$ is bounded uniformly in $\eta$ if and only if $2\kappa>3$. The torus sum is dominated by the $\Z^3$-sum because $2\delta'-2\beta\ge\delta_K/2$.
\end{proof}

\begin{lemma}[localisation of two-variable kernel bounds]\label{app:G1:leibniz-loc}
Let $\eta\le\frac12$, $0<\alpha<1$, $\beta\ge0$ and $M\ge0$. Let $D$ be the covariant forward difference, $(D_\mu f)(x)=\eta^{-1}(U_\mu(x)f(x+\eta e_\mu)-f(x))$, with bond transporters $U_\mu(x)$. In the H\"older quotients, let the transport $U(\Gamma_{\tilde x,x})$ be the product of the same bond transporters along a shortest lattice contour $\Gamma_{\tilde x,x}$ from $\tilde x$ to $x$, as in the definition of $\mathfrak N_{\alpha,\beta}$ in Lemma~\ref{app:G1:I3r-U1m}. For a kernel $Y$ put $Y_{ab}:=D^aYD^{b*}$, where $a,b\in\{0,1\}$ and all directions are allowed. Take lattice points $y,\tilde y,y',\tilde y'$ with $0<\abs{y-\tilde y}\le1$ and $0<\abs{y'-\tilde y'}\le1$. Let $H_y(\tilde y)$, $H_{y'}(\tilde y')$ and $H_{yy'}(\tilde y,\tilde y')$ be the $\alpha$-H\"older quotients of $Y_{ab}$ at these partners: in the first variable, in the second variable, and jointly. The joint quotient is the mixed second difference divided by $\abs{y-\tilde y}^\alpha\abs{y'-\tilde y'}^\alpha$. Let $\mathcal P$ be a set of lattice points such that, for all such $y,y',\tilde y,\tilde y'\in\mathcal P$, all $a,b$ and all directions,
\begin{equation}\label{app:G1:leib-hyp}
\mathrm{e}^{\beta\abs{y-y'}}\max\{\norm{Y_{ab}(y,y')},\,aH_y(\tilde y),\,bH_{y'}(\tilde y'),\,abH_{yy'}(\tilde y,\tilde y')\}\le M .
\end{equation}
Let $0\le h,h'\le1$ satisfy $\abs{\partial^\eta_\mu h}\le c_h$ and $\abs{\partial^\eta_\nu\partial^\eta_\mu h}\le c_h$ for all $\mu,\nu$, and likewise for $h'$. Assume that every lattice point at $\ell^\infty$-distance $\le1+\eta$ from $\supp h\cup\supp h'$ lies in $\mathcal P$. Then
\[
\norm{hYh'}_{1,\alpha}\le \mathrm{e}^{3\beta}(1+3\sqrt3c_h)^2M\,\mathrm{e}^{-\beta\dist(\supp h,\supp h')}\le \mathrm{e}^{5\beta}(1+6c_h)^2M\,\mathrm{e}^{-\beta\dist(\supp h,\supp h')} .
\]
Here $\norm Z_{1,\alpha}$ is the following maximum, taken over $a,b$, the directions, all lattice points $x,x'$ and all partners at distance $\le1$: the maximum of $\norm{Z_{ab}(x,x')}$ and of the quotients $aH_x$, $bH_{x'}$, $abH_{xx'}$ of $Z$. If the quotients over $\abs{x-\tilde x}>1$ and $\abs{x'-\tilde x'}>1$ are included, the bound holds with an additional factor $4$. If, as in \cite[(1.32)]{Ba3}, the nine terms of the norm (value, derivative and quotient in each variable) are moreover added instead of maximised, it holds with an additional factor $16$.
\end{lemma}
\begin{proof}
\emph{Step 1 (Leibniz identity).} Since $h$ is real, $D_\mu(hf)(x)=h(x+\eta e_\mu)(D_\mu f)(x)+(\partial^\eta_\mu h)(x)f(x)$. Put $Y^\dagger(x',x):=Y(x,x')^*$. Then $(ZD^*_\mu)(x,x')=((D_\mu Z^\dagger)(x',x))^*$. So the components and quotients of $hYh'$ in $x'$ are those of $h'Y^\dagger h$ in its first variable, and everything below applies to both variables.

\emph{Step 2 (differences of the coefficients).} Let $\abs{x-\tilde x}\le1$, and let $\Gamma=\Gamma_{\tilde x,x}=(y_0=\tilde x,\dots,y_n=x)$. Its length is $\eta n=\abs{x-\tilde x}_1\le\sqrt3\abs{x-\tilde x}$. It is monotone in each coordinate, so its points lie in the box spanned by $\tilde x$ and $x$. If $\abs{\partial^\eta_\nu\varphi}\le c$ for all $\nu$, summing along $\Gamma$ gives $\abs{\varphi(x)-\varphi(\tilde x)}\le\sqrt3c\abs{x-\tilde x}$. We apply this to $\varphi=h(\cdot+\eta e_\mu)$ and to $\varphi=\partial^\eta_\mu h$.

\emph{Step 3 (telescoping).} Let $U_i$ be the transporter of the bond $\{y_{i-1},y_i\}$ from $y_i$ to $y_{i-1}$, so that $U(\Gamma)=U_1\cdots U_n$. For a fibre-valued $W$,
\[
U(\Gamma)W(x)-W(\tilde x)=\sum_{i=1}^nU_1\cdots U_{i-1}\big[U_iW(y_i)-W(y_{i-1})\big].
\]
The bracket equals $\eta(D_\nu W)(y_{i-1})$ if $y_i=y_{i-1}+\eta e_\nu$, and $-\eta U_i(D_\nu W)(y_i)$ if $y_i=y_{i-1}-\eta e_\nu$. The transporters are unitary, so $\norm{U(\Gamma)W(x)-W(\tilde x)}\le\abs{x-\tilde x}_1\max_{z\in\Gamma,\nu}\norm{(D_\nu W)(z)}$.

\emph{Step 4 (one variable).} Suppress the second variable. Then $(hY)_0=hY_0$ and $(hY)_1(x)=F_1(x)Y_1(x)+F_2(x)Y_0(x)$, where $F_1:=h(\cdot+\eta e_\mu)$ and $F_2:=\partial^\eta_\mu h$. Hence $\norm{(hY)_1(x)}\le(1+c_h)\max_a\norm{Y_a(x)}$. With $U:=U(\Gamma_{\tilde x,x})$,
\begin{align*}
U(hY)_1(x)-(hY)_1(\tilde x)&=F_1(x)\big[UY_1(x)-Y_1(\tilde x)\big]+\big(F_1(x)-F_1(\tilde x)\big)Y_1(\tilde x)\\
&\quad+F_2(x)\big[UY_0(x)-Y_0(\tilde x)\big]+\big(F_2(x)-F_2(\tilde x)\big)Y_0(\tilde x).
\end{align*}
\begin{itemize}
\item The first term is at most $\abs{x-\tilde x}^\alpha$ times the quotient of $Y_1$ at $(x,\tilde x)$.
\item By Step~2, the second and fourth terms are at most $\sqrt3c_h\abs{x-\tilde x}$ times a component of $Y$ at $\tilde x$.
\item By Step~3, the third term is at most $\sqrt3c_h\abs{x-\tilde x}\max_{z\in\Gamma,\nu}\norm{(D_\nu Y_0)(z)}$, and $D_\nu Y_0$ is a component $Y_1$.
\end{itemize}
Since $\abs{x-\tilde x}\le\abs{x-\tilde x}^\alpha$, the quotient of $(hY)_1$ at $(x,\tilde x)$ is at most $(1+3\sqrt3c_h)$ times the largest component or quotient of $Y$ at points of $\{x,\tilde x\}\cup\Gamma$, with partners in the same set.

\emph{Step 5 (two variables, points, exponential).} Each expansion in Step~4 is a sum of terms. Each term is a coefficient depending only on the expanded variable, times a linear functional of $Y$ in that variable: a value, a transported difference, or a telescoped sum of $D_\nu Y$. The expansions in $x$ and in $x'$ act on different variables, and the transports act from opposite sides, so the two expansions commute. This gives the following bounds, in each case times the largest of the quantities in \eqref{app:G1:leib-hyp} at pairs of evaluated points with evaluated partners:
\begin{itemize}
\item every component of $hYh'$: $(1+c_h)^2$;
\item every one-variable quotient: $(1+3\sqrt3c_h)(1+c_h)$;
\item the mixed quotient: $(1+3\sqrt3c_h)^2$.
\end{itemize}
A term vanishes when its coefficient vanishes. For an $x$-coefficient to be non-zero, $x$ or $\tilde x$ must lie within distance $\eta$ of $\supp h$ along a single coordinate direction. Hence all evaluated points lie within $\ell^\infty$- and Euclidean distance $1+\eta$ of $\supp h$, and similarly for $x'$ and $h'$. So these points lie in $\mathcal P$, and \eqref{app:G1:leib-hyp} applies. Moreover $\mathrm{e}^{-\beta\abs{y-y'}}\le \mathrm{e}^{2(1+\eta)\beta}\mathrm{e}^{-\beta\dist(\supp h,\supp h')}\le \mathrm{e}^{3\beta}\mathrm{e}^{-\beta\dist(\supp h,\supp h')}$. Since $3\sqrt3\le6$, both bounds follow.

\emph{Step 6 (other forms of the norm).} A quotient over a distance $>1$ is at most twice the supremum of the quantity whose difference is taken. Since $1+c_h\le1+3\sqrt3c_h$, this costs a factor $2$ per variable. In each variable, the sum of the value, the derivative and the quotient over all distances is at most $1+(1+c_h)+\max\{1+3\sqrt3c_h,2+2c_h\}\le4(1+3\sqrt3c_h)$. This gives the factor $16$ for the summed form.
\end{proof}
\noindent The bound on second differences enters only through $\abs{F_2(x)-F_2(\tilde x)}$ in Step~4. With first differences alone, this term is only bounded by $2c_h\abs{x-\tilde x}^{-\alpha}\norm{Y_0(\tilde x)}$, which is of order $\eta^{-\alpha}$ at $\abs{x-\tilde x}=\eta$.

\begin{lemma}[the cut-offs of \cite{Ba3}, (1.33)]\label{app:G1:ramp-ch}
For $\nu\ne\mu$, the cut-off $h$ of Lemma~\ref{app:G1:Rloc7}(iii) satisfies
\[
\abs{\partial^\eta_\mu h}\le\tfrac{15}4,\qquad\abs{\partial^\eta_\mu\partial^\eta_\mu h}\le\tfrac{40}{\sqrt3},\qquad\abs{\partial^\eta_\nu\partial^\eta_\mu h}\le\big(\tfrac{15}4\big)^2 .
\]
Hence the hypothesis of Lemma~\ref{app:G1:leibniz-loc} holds with $c_h:=40/\sqrt3$, uniformly in $\eta$ and $\square(v)$.
\end{lemma}
\begin{proof}
In Lemma~\ref{app:G1:Rloc7}(iii), $h(x)=\prod_\mu g_\mu(x_\mu)$ with $g_\mu(u)=\chi(\dist(u,I_\mu))$. Here $I_\mu$ is the $\mu$-th side interval of $\square(v)$, and $\chi(t)=1$ for $t\le\frac12$, $\chi(t)=1-s(2t-1)$ for $\frac12\le t\le1$, and $\chi(t)=0$ for $t\ge1$, where $s(t)=10t^3-15t^4+6t^5$.

Since $\chi=1$ on $[0,\frac12]$, $g_\mu=1$ on the $\frac12$-neighbourhood of $I_\mu$. Outside $I_\mu$ the distance is affine with slope $\pm1$. Hence $g_\mu\in C^2(\R)$, with $\abs{g_\mu'}\le\sup\abs{\chi'}=2\max s'=\frac{15}4$ and $\abs{g_\mu''}\le\sup\abs{\chi''}=4\max\abs{s''}=\frac{40}{\sqrt3}$. Here $s'=30t^2(1-t)^2$ is maximal at $t=\frac12$, and $\abs{s''}=60t(1-t)\abs{1-2t}$ is maximal at $t=\frac{3\pm\sqrt3}6$, where it equals $10/\sqrt3$.

Since $0\le g_\mu\le1$, we get $\abs{\partial_\mu h}\le\frac{15}4$, $\abs{\partial_\mu^2h}\le\frac{40}{\sqrt3}$ and $\abs{\partial_\nu\partial_\mu h}\le(\frac{15}4)^2$. Lattice differences are averages of derivatives, including the case $\nu=\mu$:
\[
\partial^\eta_\mu h(x)=\eta^{-1}\int_0^\eta\partial_\mu h(x+se_\mu)\,ds,\qquad\partial^\eta_\nu\partial^\eta_\mu h(x)=\eta^{-2}\int_0^\eta\!\!\int_0^\eta\partial_\nu\partial_\mu h(x+se_\mu+te_\nu)\,ds\,dt .
\]
Finally $\frac{15}4\le(\frac{15}4)^2\le\frac{40}{\sqrt3}$.
\end{proof}

\begin{proposition}[the bound for \cite{Ba3}, (1.16), in the mixed norm]\label{app:G1:sandwich}
Assume (R$_\vartheta$), and let $0<\beta\le\beta_U:=\min\{\beta_*,\delta_K/4\}$. Suppose that either $m=n+n'\ge4$ and $0<\alpha<\frac12$, or $m\ge6$ and $0<\alpha<1$; in particular, $\alpha=\alpha_0=\frac12$ requires $m\ge6$.
\begin{enumerate}
\item[(a)] If $\mathcal X=T_\eta$, then $\mathfrak N_{\alpha,\beta}(K_{n,n'})\le C_U(\alpha)(C_V\sigma)^m$. Here $C_V:=\mathrm{e}^{\beta(\sqrt3+1)}c_V\max\{c_W,1\}$ and $c_W:=(1+\sigma)^2\big(2\gamma_0^{-1}+((1+d)^2-1)K_W\big)$.
\item[(b)] Let $\mathcal X=\Omega=B^k(\Omega^{(k)})$ with Neumann conditions, and assume $\dist(\supp\tilde A,\Omega^c)\ge R_0+2\sqrt3+1$. Then the same bound holds when the supremum in $\mathfrak N_{\alpha,\beta}$ is restricted as follows: the points $x,x'$, together with their bond partners, have depth $\ge R_0$, and the quotients are taken only at partners $\tilde x,\tilde x'$ with the same property.
\item[(c)] Let $0\le h,h'\le1$ with $\abs{\partial^\eta_\mu h},\abs{\partial^\eta_\nu\partial^\eta_\mu h}\le c_h$ for all $\mu,\nu$, and likewise for $h'$. In case (b), assume in addition that $\supp h\cup\supp h'$ has depth $\ge R_0+3$. Then
\[
\norm{hK_{n,n'}h'}_{1,\alpha}\le \mathrm{e}^{5\beta}(1+6c_h)^2C_U(\alpha)(2c_1C_V)^{n+n'}(e_kp(t))^{n+n'}\mathrm{e}^{-\beta\dist(\supp h,\supp h')} .
\]
This is the bound \cite[(2.5)]{Ba3} for the operator (1.16), with Ba{\l}aban's rate $\delta_0$ taken $\le\beta_U$.
\end{enumerate}
$C_U(\alpha)$ depends only on $d,a,M_b,L,\gamma_0,K_W,\alpha,\beta$. The bound is uniform in $\eta$ (hence in $k$ and $\eps$), in the torus side, in $0\le m^2t^2\le c$ and in the backgrounds within (R$_\vartheta$). It is also uniform in the parameters $(u,M)$ of Theorem~\ref{thm:T6}, since no photon line occurs in (1.16).
\end{proposition}
\noindent The depth hypothesis of (b) follows from $R>R_0$ and $R\ge12\sqrt d$ (conditions (R.1) and (R.8) of Table~\ref{p:tab:A2}): these give $2r(t)\ge R_0+2\sqrt3+1$ under (R4), and under (R4$'''$) of Corollary~\ref{cor:prop1R4} the hypothesis holds with room. In the construction $n=n'=6$, so $m=12$, and the H\"older exponent of \cite[(2.5)]{Ba3} is $\alpha_0=\frac12$. This is the $L^2$ sandwich used in part (b), step S1, of the proof of Proposition~\ref{prop:prop1}.

\begin{proof}
\emph{Step 1 (why ends are needed).} $K_{n,n'}$ has $m$ vertices and $m+1$ lines, and every vertex is first order on each side (Lemma~\ref{app:G1:vertex}(i)). So $D^aK_{n,n'}D^{b*}$ carries $m+2$ derivatives on $m+1$ lines, and some line carries two. That line is placed in a middle factor, which enters only through its weighted $\ell^2$ norm (W). A pointwise bound of $D_xGD^*_y$ is never used.

\emph{Step 2 (factorisation, two-line ends, $m\ge4$).} Write $Y:=D^aK_{n,n'}D^{b*}=E_L\mathcal NE_R$ with
\[
E_L:=D^aL_0VL_1,\qquad \mathcal N:=VL_2V\cdots VL_{m-2}V,\qquad E_R:=(D^bL_mVL_{m-1})^*.
\]
Put the first vertex of $\mathcal N$ in right-pushed form, the last in left-pushed form, and the others in symmetric form (Lemma~\ref{app:G1:vertex}(i)--(iii)):
\[
\mathcal N=(-1)^{m-2}\mathfrak r(\mathfrak DL_2\mathfrak D^*)\mathfrak E\cdots\mathfrak E(\mathfrak DL_{m-2}\mathfrak D^*)\mathfrak r^*.
\]
For $L=G'$, the identity $D_{\tilde B,\mu}=U_\mu(\tilde A)^{-1}D_{\tilde A+\tilde B,\mu}-\phi_\mu$ costs a factor $1+\sigma$ on each side. With (W) and Lemma~\ref{app:G1:vertex}(iv), for $\abs\beta\le\beta_U$,
\[
\norm{\mathrm{e}^{\mp\beta\rho}\mathcal N\mathrm{e}^{\pm\beta\rho}}\le(\mathrm{e}^{\beta(\sqrt3+1)}c_V\sigma)^{m-2}c_W^{m-3}.
\]
By (S) and $V^*=V$, Lemma~\ref{app:G1:ends} applies to $E_R$. In $E_L$ and $E_R$ each line carries exactly one derivative. The pointwise bounds (K) used in Lemma~\ref{app:G1:ends} are those of Lemma~\ref{app:G1:WK-K}, whose proof sums the slice bounds by Lemma~\ref{app:G1:WK-slicesum}; that lemma holds at every distance $z\ge0$.

\emph{Step 3 (sandwich).} Fix $x,x'$ and put $\rho:=\abs{\cdot-x}$. Then
\[
Y(x,x')=\big\langle \mathrm{e}^{\beta\rho}E_L(x,\cdot)^*,\ (\mathrm{e}^{-\beta\rho}\mathcal N\mathrm{e}^{\beta\rho})\,\mathrm{e}^{-\beta\rho}E_R(\cdot,x')\big\rangle_{\ell^2(\eta^3)}.
\]
Cauchy--Schwarz, $\mathrm{e}^{-\beta\abs{y'-x}}\le \mathrm{e}^{-\beta\abs{x-x'}}\mathrm{e}^{\beta\abs{y'-x'}}$ and Lemma~\ref{app:G1:ends}(iii) give
\[
\mathrm{e}^{\beta\abs{x-x'}}\norm{Y(x,x')}\le(c_EJ_2\sigma)^2(\mathrm{e}^{\beta(\sqrt3+1)}c_V\sigma)^{m-2}c_W^{m-3}.
\]
$Y$ is linear in $E_L(x,\cdot)$ and in $E_R(\cdot,x')$ separately. So $H_x,H_{x'},H_{xx'}$ obey the same bound with $c_EJ_2$ replaced by $2\mathrm{e}^\beta c_E(\alpha)J_{2-\alpha}$, which is finite for $\alpha<\frac12$. Hence $\mathfrak N_{\alpha,\beta}(K_{n,n'})\le C^{(2)}_U(\alpha)\sigma^2(\mathrm{e}^{\beta(\sqrt3+1)}c_V\sigma)^{m-2}c_W^{m-3}$, with $C^{(2)}_U(\alpha):=\max\{c_EJ_2,2\mathrm{e}^\beta c_E(\alpha)J_{2-\alpha}\}^2$.

\emph{Step 4 (three-line ends, $m\ge6$, any $\alpha<1$).} Repeat Steps 2--3 with $E_L:=E^{(3)}_a$, $\mathcal N:=VL_3V\cdots VL_{m-3}V$ and $E_R:=(D^bL_mVL_{m-1}VL_{m-2})^*$. Now $J_3,J_{3-\alpha}<\infty$ for every $\alpha<1$, and $\mathfrak N_{\alpha,\beta}\le C^{(3)}_U(\alpha)\sigma^4(\mathrm{e}^{\beta(\sqrt3+1)}c_V\sigma)^{m-4}c_W^{m-5}$, with $C^{(3)}_U(\alpha):=\max\{c'_EJ_3,2\mathrm{e}^\beta c'_E(\alpha)J_{3-\alpha}\}^2$. Since $c_V\ge1$, both bounds are $\le C_U(\alpha)(C_V\sigma)^m$, with $C_U:=C^{(2)}_Uc_V^{-2}$, respectively $C^{(3)}_Uc_V^{-4}$. This proves (a).

\emph{Step 5 (Neumann domains).} On $\Omega$, (H0) and (W) hold with the Neumann condition (Lemma~\ref{app:G1:WKinputs}). Lemma~\ref{app:G1:vertex}(ii) holds on $\Omega$, because every bond of a vertex lies in $\Omega$. Since $V(y,\cdot)=0$ unless $y\in U_{\sqrt3+\eta}(\supp\tilde A)$, every pointwise use of (K) in Lemma~\ref{app:G1:ends} is at one of two kinds of points:
\begin{itemize}
\item points of $U_{2\sqrt3+1}(\supp\tilde A)$, which have depth $\ge R_0$ by hypothesis;
\item the points $x,x'$, their bond partners, and, for the quotients, $\tilde x,\tilde x'$; these have depth $\ge R_0$ by the restriction in (b).
\end{itemize}
Steps 2--4 apply verbatim, which proves (b).

\emph{Step 6 (localisation, (c)).} Apply Lemma~\ref{app:G1:leibniz-loc} with $M:=C_U(\alpha)(C_V\sigma)^m$. In case (a) take $\mathcal P:=T_\eta$. In case (b) take $\mathcal P:=\{z:\text{depth of }z\ge R_0+1+\eta\}$. Points of this $\mathcal P$ and their bond partners have depth $\ge R_0$, so (b) gives \eqref{app:G1:leib-hyp} on $\mathcal P$. Since $\eta\le\frac12$, every point within $\ell^\infty$-distance $1+\eta$ of a set of depth $\ge R_0+3$ has depth $\ge R_0+2-\eta\ge R_0+1+\eta$, so it lies in $\mathcal P$. Finally use $\sigma\le2c_1e_kp(t)$. The bound \cite[(2.5)]{Ba3} is used with Ba{\l}aban's tree rate $\delta_0\le\beta_U$, which only shrinks a free constant.
\end{proof}

\begin{remark}
(i) At $m=0$ the mixed kernel is not uniform in $\eta$. With $g:=\eta^{-3}1_{\{x\}}$ and $\Lambda:=4d\eta^{-2}+m^2t^2+a\ge\norm{G'^{-1}}$, one has
\[
(D_\mu G'D^*_\mu)(x,x)=\langle D^*_\mu g,G'D^*_\mu g\rangle\ge\Lambda^{-1}\norm{D^*_\mu g}_2^2=2\Lambda^{-1}\eta^{-5}\ge\eta^{-3}/(6+O(\eta^2))
\]
per direction. This is why \cite[p.~414]{Ba3} needs ``$n,n'$ sufficiently large''. (ii) The operator form of the bound, in the norms of \cite[(2.24)--(2.25)]{Ba1}, which is the literal statement of \cite[p.~414]{Ba3}, holds for every $m\ge0$. It is the argument of \cite[(2.33)]{Ba4} applied to one term of the series. It is not needed for (2.5).
\end{remark}

\paragraph{Domain differences in the mixed norm: the two-domain theorem.}\label{app:G1:ssec-twodomain}
Let $\mathcal X_{\rm in}\subset \mathcal X_{\rm out}$ be unions of big blocks of $T_\eta$, where $\mathcal X_{\rm out}$ may be $T_\eta$ itself. Let $A$ be real with $\theta(A)\le\vartheta_R$ on $\mathcal X_{\rm out}$, and $0\le m^2t^2\le c$. Put $G_{\rm in}:=G_k(\mathcal X_{\rm in},A)$, $G_{\rm out}:=G_k(\mathcal X_{\rm out},A)$, $H_{\rm in},H_{\rm out}$ the corresponding operators, $\delta:=G_{\rm out}-G_{\rm in}$ (signs are irrelevant for the bounds), and $\rho:=\rho_{\mathcal X_{\rm in}}$. Then $\rho_{\mathcal X_{\rm out}}\ge\rho$.

\emph{Cut-offs.} Fix $f\in C^\infty(\R)$ with $0\le f\le1$, $f=0$ on $(-\infty,0]$ and $f=1$ on $[1,\infty)$, and $\varphi\in C^\infty_c([-\frac14,\frac14]^3)$ with $\varphi\ge0$ and $\int\varphi=1$. Let $\tilde\rho:=\rho^{\rm cont}*\varphi$, where $\rho^{\rm cont}$ is the continuum $\ell^\infty$ distance to the closed complement of $\mathcal X_{\rm in}$. Then $\abs{\tilde\rho-\rho}\le\frac14+\eta$ on lattice points. For $s\in\R$ put $\chi_s:=f(\tilde\rho-s)$. Then $\chi_s=0$ on $\{\rho\le s-\frac14-\eta\}$, $\chi_s=1$ on $\{\rho\ge s+\frac54+\eta\}$, and the first and second lattice differences of $\chi_s$ are bounded by absolute constants, uniformly in $\eta$ and $\mathcal X_{\rm in}$. We take
\[
\chi:=\chi_{R_0+3},\qquad \psi:=\chi_{R_0+7}.
\]

Let $X_{\rm in}\subset X_{\rm out}$ be unions of big blocks of $T_\eta$, $\eta=L^{-k}\le\frac12$; here $X_{\rm out}$ may be $T_\eta$, while $X_{\rm in}\ne T_\eta$. Let $A$ be real with $\theta(A)\le\vartheta_R$, and let $0\le m^2t^2\le c$. Put $H_{\rm in/out}:=H_{X_{\rm in/out}}(A)$, $G_{\rm in/out}:=H_{\rm in/out}^{-1}$, $\delta:=G_{\rm out}-G_{\rm in}$ and $\rho(z):=\dist_\infty(z,X_{\rm in}^c)$; then $\rho_{X_{\rm out}}\ge\rho$. For a matrix-valued function $F$ of one variable put $\norm F_2:=(\sum_z\eta^3\norm{F(z)}^2)^{1/2}$; for a row of four such functions, the four norms are added. The quotients $H_x,H_{x'},H_{xx'}$ are those of $\mathfrak N_{\alpha,\beta}$, with transport along a shortest lattice contour.

\begin{lemma}[collars]\label{app:G1:collar}
Fix $f\in C^\infty(\R)$ with $0\le f\le1$, $f=0$ on $(-\infty,0]$ and $f=1$ on $[1,\infty)$. Fix $\varphi\in C^\infty_c([-\frac14,\frac14]^3)$ with $\varphi\ge0$ and $\int\varphi=1$. For a unit block $b$ with lowest corner $y_b$ let $\hat b:=y_b+[0,1)^3$, and put $\hat X_{\rm in}:=\bigcup_{b\subset X_{\rm in}}\hat b$. Let $\rho^{\rm c}(u)$ be the continuum $\ell^\infty$ distance from $u$ to the complement of $\hat X_{\rm in}$. Put $\tilde\rho:=\rho^{\rm c}*\varphi$, and let $\chi_s:=f(\tilde\rho-s)$ on lattice points. Let $c_1':=\norm{f'}_\infty$ and $c_2':=\norm{f''}_\infty+\norm{f'}_\infty\max_\nu\norm{\partial_\nu\varphi}_{L^1}$. Then:
\begin{enumerate}
\item[(i)] $\tilde\rho$ is $1$-Lipschitz for $\abs\cdot_\infty$, and $\rho-\frac14-\eta\le\tilde\rho\le\rho+\frac14$ at lattice points;
\item[(ii)] $\chi_s=0$ on $\{\rho\le s-\frac14\}$ and $\chi_s=1$ on $\{\rho\ge s+\frac54+\eta\}$;
\item[(iii)] $\abs{\partial^\eta_\mu\chi_s}\le c_1'$ and $\abs{\partial^\eta_\mu\partial^\eta_\nu\chi_s}\le c_2'$, uniformly in $\eta$, $s$ and $X_{\rm in}$;
\item[(iv)] if $\chi_s(u)\ne\chi_s(u')$ and $\abs{u-u'}_\infty\le\ell$, then $\tilde\rho(u),\tilde\rho(u')\in(s-\ell,s+1+\ell)$.
\end{enumerate}
\end{lemma}
\begin{proof}
(i) The cubes $\hat b$ tile the torus, so a lattice point outside $X_{\rm in}$ lies outside $\hat X_{\rm in}$; hence $\rho^{\rm c}\le\rho$. Conversely, $\rho^{\rm c}(z)$ is attained at a point of a closed cube $\overline{\hat b'}$ with $b'\not\subset X_{\rm in}$. The lattice points of $b'$ are $\eta$-dense in $\overline{\hat b'}$ for $\abs\cdot_\infty$, so $\rho(z)\le\rho^{\rm c}(z)+\eta$. Finally, $\rho^{\rm c}$ is $1$-Lipschitz, and mollification by $\varphi$ preserves this and moves values by at most $\frac14$.

(ii) follows from (i).

(iii) By the mean value theorem, the lattice differences are bounded by the suprema of the continuum derivatives. Here $\abs{\partial_\mu\tilde\rho}\le1$ and $\partial_\mu\partial_\nu\tilde\rho=(\partial_\mu\rho^{\rm c})*\partial_\nu\varphi$.

(iv) If $\tilde\rho(u)\le s-\ell$, then $\tilde\rho(u')\le s$, and both values of $\chi_s$ are $0$. If $\tilde\rho(u)\ge s+1+\ell$, both values are $1$. The same holds with $u$ and $u'$ exchanged.
\end{proof}

\begin{lemma}[commutators]\label{app:G1:commutator}
Let $s\ge1$, $\chi:=\chi_s$ and $\phi_\mu:=\partial^\eta_\mu\chi$.
\begin{enumerate}
\item[(i)] $H_{\rm out}(\chi g)=H_{\rm in}(\chi g)$ for every $g$ on $X_{\rm in}$ (extended by zero).
\item[(ii)] $C_\chi:=H_{\rm out}\chi-\chi H_{\rm in}=[H_{\rm in},\chi]=-\mathfrak r_\chi\mathfrak D=\mathfrak D^*\mathfrak r_\chi^*$. Here $\mathfrak r_\chi=(\mathfrak r_0,\dots,\mathfrak r_3)$ is given by
\[
\mathfrak r_0:=\sum_\mu(\partial^\eta_\mu\phi_\mu)(\cdot-\eta e_\mu)-a_k[P_k(A),\chi],\qquad(\mathfrak r_\mu g)(x):=\phi_\mu(x)g(x)+\phi_\mu(x-\eta e_\mu)U_{x-\eta e_\mu,\mu}(A)^{-1}g(x-\eta e_\mu).
\]
\item[(iii)] $\sum_{i=0}^3\sup_z\sum_w\eta^3\norm{\mathfrak r_i(w,z)}\le c_{\mathfrak r}:=6c_1'+3c_2'+ac_1'$.
\item[(iv)] If $\mathfrak r_i(w,z)\ne0$, then $\abs{w-z}\le\sqrt3$ and $w,z\in S_s:=\{s-\frac54-\eta<\rho<s+\frac94+2\eta\}$.
\end{enumerate}
\end{lemma}
\begin{proof}
(i) If $\chi(y)\ne0$, then $\rho(y)>s-\frac14>\eta$, so every bond at $y$ lies in $X_{\rm in}$. The unit block of $y$ also lies in $X_{\rm in}$, since $X_{\rm in}$ is a union of unit blocks. A bond of $X_{\rm out}$ that is not a bond of $X_{\rm in}$ has both endpoints at depth $\le\eta$, hence outside $\supp\chi$. Since $P_k$ acts inside unit blocks, the two operators agree on $\chi g$.

(ii) The Leibniz rules $D_\mu(\chi f)=\chi(\cdot+\eta e_\mu)D_\mu f+\phi_\mu f$ and $D_\mu^*(\chi(\cdot+\eta e_\mu)g)=\chi D^*_\mu g-\phi_\mu g$ give $[D^*_\mu D_\mu,\chi]=D^*_\mu\phi_\mu-\phi_\mu D_\mu$. Also $[m^2t^2,\chi]=0$. By \cite[(2.32), p.~581]{Ba4},
\[
D^*_\mu(\phi_\mu f)(x)=-\phi_\mu(x-\eta e_\mu)U^{-1}_{x-\eta e_\mu,\mu}(D_\mu f)(x-\eta e_\mu)-(\partial^\eta_\mu\phi_\mu)(x-\eta e_\mu)f(x).
\]
Collecting terms gives $C_\chi=-\mathfrak r_\chi\mathfrak D$. Since $H_{\rm in}$ is self-adjoint and $\chi$ is real, $C_\chi^*=-C_\chi$, which gives the second form.

(iii) For each $\mu$, the two parts of $\mathfrak r_\mu$ have column sums $\le c_1'$ each, which gives $6c_1'$ in total. The multiplication part of $\mathfrak r_0$ is bounded by $3c_2'$. For the last part, $[P_k,\chi](w,z)=P_k(w,z)(\chi(z)-\chi(w))$. Here $\sum_w\eta^3\norm{P_k(w,z)}\le1$, $\abs{\chi(z)-\chi(w)}\le c_1'$ inside a unit block, and $a_k\le a$.

(iv) There are two cases.
\begin{itemize}
\item For $\mathfrak r_\mu$ and the first part of $\mathfrak r_0$: here $z\in\{w,w-\eta e_\mu\}$, and $\chi$ differs at two adjacent points of $\{w-\eta e_\mu,w,w+\eta e_\mu\}$. Lemma~\ref{app:G1:collar}(iv) with $\ell=\eta$ gives $\tilde\rho(w),\tilde\rho(z)\in(s-3\eta,s+1+3\eta)$.
\item For $[P_k,\chi]$: here $\abs{w-z}_\infty\le1-\eta$ and $\chi(w)\ne\chi(z)$, which gives $\tilde\rho(w),\tilde\rho(z)\in(s-1+\eta,s+2-\eta)$.
\end{itemize}
Since $\eta\le\frac12$, both $3\eta$ and $1-\eta$ are $\le1+\eta$, and Lemma~\ref{app:G1:collar}(i) gives $w,z\in S_s$. Finally $\abs{w-z}\le\sqrt3\max(\eta,1-\eta)\le\sqrt3$.
\end{proof}

We take $\chi:=\chi_{R_0+3}$ and $\psi:=\chi_{R_0+13/2}$, with $S_\chi:=S_{R_0+3}$ and $S_\psi:=S_{R_0+13/2}$. We also put $\mathcal P_D:=\{z:\rho(z)\ge R_0+10+2\eta\}$. Since $10+2\eta\le11$, $\mathcal P_D\supset\{\rho\ge R_0+11\}$ for every $k\ge1$.

\begin{proposition}[the two-domain theorem]\label{app:G1:twodomain}
Let $0<\alpha<1$, $0<\beta\le\beta_D:=\min\{\beta_*,\delta_K/4\}$, $a,b\in\{0,1\}$ and $Y:=D^a\delta D^{b*}$. Then for $x,x'\in\mathcal P_D$, with partners $\tilde x,\tilde x'\in\mathcal P_D$,
\[
\mathrm{e}^{\beta\abs{x-x'}}\max\{\norm{Y(x,x')},aH_x,bH_{x'},abH_{xx'}\}\le C_D(\alpha)\exp\Big(-\tfrac{\delta_K}2\big[(\rho(x)-R_0-7)+(\rho(x')-R_0-10)\big]\Big).
\]
The constants are
\begin{gather*}
C_D(\alpha):=c'_Wc_{\rm end}(\alpha)^2,\qquad c'_W:=2/\gamma_0+3K_W,\\
c_{\rm end}(\alpha):=c_{\mathfrak r}\mathrm{e}^{\delta_K(\sqrt3+1)}\max\{c_K,2c_K(\alpha)\}c_3^{1/2},\qquad c_3:=(\tfrac12+4\sqrt3/\delta_K)^3.
\end{gather*}
They depend only on $d,a,M_b,L,\gamma_0,K_W,\alpha$ and on $f,\varphi$. They do not depend on $\eta$, $k$, $\eps$, the side of the torus, $X_{\rm in}$, $X_{\rm out}$, $A$, $m^2t^2\in[0,c]$ or $R_0$.
\end{proposition}
\noindent The proof uses only (H0), (W), (K) and (S) of Lemma~\ref{app:G1:WKinputs}, in the forms of Lemmas~\ref{app:G1:WK-W} and \ref{app:G1:WK-K}.

\begin{proof}
\emph{Step 1 (collars).} Lemma~\ref{app:G1:commutator}(iv) gives
\[
S_\chi\subset\{R_0+\tfrac74-\eta<\rho<R_0+\tfrac{21}4+2\eta\},\qquad S_\psi\subset\{R_0+\tfrac{21}4-\eta<\rho<R_0+\tfrac{35}4+2\eta\}.
\]
We use the following consequences.
\begin{itemize}
\item $\psi=0$ on $\{\rho\le R_0+\frac{25}4\}\supset S_\chi$, since $\frac{21}4+2\eta\le\frac{25}4$ for $\eta\le\frac12$.
\item $S_\chi\cup S_\psi\subset\{\rho>R_0+\frac54\}$.
\item $\chi=1$ on $\{\rho\ge R_0+\frac{17}4+\eta\}$ and $\psi=1$ on $\{\rho\ge R_0+\frac{31}4+\eta\}$. Both sets contain $\mathcal P_D$ and all bond partners of points of $\mathcal P_D$.
\item Since $\rho$ is $1$-Lipschitz, every $z\in\mathcal P_D$ satisfies
\[
\dist_\infty(z,S_\chi)>\max\{\rho(z)-R_0-7,\tfrac{19}4\},\qquad\dist_\infty(z,S_\psi)>\max\{\rho(z)-R_0-10,\tfrac54\}.
\]
\end{itemize}

\emph{Step 2 (representation).} We have $G_{\rm out}C_\chi G_{\rm in}=\chi G_{\rm in}-G_{\rm out}\chi$, so $\delta(y,y')=-(G_{\rm out}C_\chi G_{\rm in})(y,y')$ when $\chi(y)=\chi(y')=1$. Likewise $G_{\rm in}C_\psi G_{\rm in}=\psi G_{\rm in}-G_{\rm in}\psi$, so $G_{\rm in}(z,y')=-(G_{\rm in}C_\psi G_{\rm in})(z,y')$ when $\psi(z)=0$ and $\psi(y')=1$. In the first identity only points $z\in S_\chi\subset\{\psi=0\}$ occur. Inserting the second identity into the first and using Lemma~\ref{app:G1:commutator}(ii) gives
\[
\delta=-G_{\rm out}\mathfrak r_\chi(\mathfrak DG_{\rm in}\mathfrak D^*)\mathfrak r_\psi^*G_{\rm in}
\]
at all points of $\mathcal P_D$ and their bond partners. By Step~1, these are the only points at which $D^a$ and $D^{b*}$ evaluate $\delta$. Hence
\[
Y(x,x')=-\langle E_L(x,\cdot),\mathfrak ME_R(\cdot,x')\rangle,\qquad E_L:=D^aG_{\rm out}\mathfrak r_\chi,\quad\mathfrak M:=\mathfrak DG_{\rm in}\mathfrak D^*,\quad E_R:=\mathfrak r_\psi^*G_{\rm in}D^{b*}.
\]

\emph{Step 3 (left end).} $E_L(x,z)=\sum_w\eta^3D^a_xG_{\rm out}(x,w)\mathfrak r_\chi(w,z)$, with $w,z\in S_\chi$ and $\abs{w-z}\le\sqrt3$. Here (K) applies on $X_{\rm out}$ at $(x,w)$, and $\abs{x-w}\ge1$. So $K^{\delta_K}_{2-a}(x-w)\le \mathrm{e}^{-\delta_K\abs{x-w}}\le \mathrm{e}^{\delta_K\sqrt3}\mathrm{e}^{-\delta_K\abs{x-z}}$, and
\[
\norm{E_L(x,z)}\le c_{\mathfrak r}c_K\mathrm{e}^{\delta_K\sqrt3}\mathrm{e}^{-\delta_K\abs{x-z}}1_{S_\chi}(z).
\]
Since $2(\delta_K-\beta)\ge\frac32\delta_K$,
\[
\norm{\mathrm{e}^{\beta\abs{\cdot-x}}E_L(x,\cdot)}_2\le c_{\mathfrak r}c_K\mathrm{e}^{\delta_K\sqrt3}\mathrm{e}^{-\frac{\delta_K}2\dist(x,S_\chi)}\Big(\sum_z\eta^3\mathrm{e}^{-\frac{\delta_K}2\abs{z-x}}\Big)^{1/2}.
\]
To bound the lattice sum by $c_3$, pass on the torus to shortest lifts and use $\abs z\ge\abs z_1/\sqrt3$ and $\coth(c/2)\le1+2/c$. With Step~1 this gives $\le c_{\rm end}(\alpha)\mathrm{e}^{-\frac{\delta_K}2(\rho(x)-R_0-7)}$.

For $a=1$ and a partner $\tilde x\in\mathcal P_D$, use (K2) and $\abs{\tilde x-w}\ge1$. The transported difference then has kernel $\le2c_K(\alpha)\mathrm{e}^{\delta_K}\abs{x-\tilde x}^\alpha \mathrm{e}^{-\delta_K\abs{x-w}}$, and the same computation applies.

\emph{Step 4 (right end).} By (S), $(G_{\rm in}D^{b*})(v,x')=(D^b_{x'}G_{\rm in}(x',v))^T$ with $v\in S_\psi$, and (K) applies on $X_{\rm in}$. As in Step~3,
\[
\norm{\mathrm{e}^{\beta\abs{\cdot-x'}}E_R(\cdot,x')}_2\le c_{\rm end}(\alpha)\mathrm{e}^{-\frac{\delta_K}2(\rho(x')-R_0-10)}.
\]
The same bound, times $\abs{x'-\tilde x'}^\alpha$, holds for the transported difference.

\emph{Step 5 (middle).} By Lemma~\ref{app:G1:WK-W} on $X_{\rm in}$ with weight $\abs{\cdot-x}$, $\norm{\mathrm{e}^{-\beta\abs{\cdot-x}}\mathfrak M\mathrm{e}^{\beta\abs{\cdot-x}}}\le c'_W$. Together with $\mathrm{e}^{-\beta\abs{u-x}}\le \mathrm{e}^{-\beta\abs{x-x'}}\mathrm{e}^{\beta\abs{u-x'}}$ and Cauchy--Schwarz, this gives the bound for $\norm{Y(x,x')}$. $Y$ is linear in $E_L(x,\cdot)$ and in $E_R(\cdot,x')$ separately, so the quotients follow by replacing one or both ends by their differences. No kernel is evaluated at points less than $1$ apart, and the only doubly differentiated line enters through its operator norm.
\end{proof}

\begin{corollary}[localised form]\label{app:G1:twodomain-loc}
Let $0\le h,h'\le1$ with $\abs{\partial^\eta_\mu h},\abs{\partial^\eta_\nu\partial^\eta_\mu h}\le c_h$ for all $\mu,\nu$, and likewise for $h'$. Assume $\rho\ge R_0+11+3\eta$ on $\supp h\cup\supp h'$. This holds if $\rho\ge R_0+12$ there when $k\ge2$, and if $\rho\ge R_0+13$ there when $k=1$. Then, at any mutual distance of the supports (overlapping supports included),
\[
\norm{h\,\delta\,h'}_{1,\alpha}\le C_{\rm loc}\,\mathrm{e}^{-\beta\dist(\supp h,\supp h')},\qquad C_{\rm loc}:=\mathrm{e}^{5\beta}(1+6c_h)^2C_D(\alpha).
\]
\end{corollary}
\begin{proof}
Apply Lemma~\ref{app:G1:leibniz-loc} with $\mathcal P:=\mathcal P_D$ and $M:=C_D(\alpha)$. On $\mathcal P_D$ both brackets in the exponent of Proposition~\ref{app:G1:twodomain} are nonnegative, so the exponential factor is $\le1$ and \eqref{app:G1:leib-hyp} holds. Since $\rho$ is $1$-Lipschitz for $\ell^\infty$, every point within $\ell^\infty$-distance $1+\eta$ of the supports has $\rho\ge R_0+10+2\eta$, so it lies in $\mathcal P_D$.
\end{proof}

\begin{lemma}[nesting]\label{app:G1:nesting}
Let $\rho_0\ge1$, and let $\Lambda_{i_0}\supset\Lambda_{i_0+1}\supset\cdots$ be sets of points of a cubic lattice of spacing $\le1$. Assume that every point at $\ell^\infty$-distance $\le\rho_0$ from $\Lambda_i^c$ lies in $\Lambda_{i+1}^c$. Then $\dist_\infty(\Lambda_{i_0+m},\Lambda_{i_0}^c)>m\rho_0-(m-1)$ for every $m\ge1$.
\end{lemma}
\begin{proof}
For $m=1$ this is the hypothesis. For the induction step, let $x\in\Lambda_{i_0+m}$ and $z\in\Lambda_{i_0}^c$; then $\abs{x-z}_\infty>\rho_0$. A king-move lattice geodesic from $z$ to $x$ advances $\abs{\cdot-z}_\infty$ by one lattice step ($\le1$) at a time. So it contains a point $w$ with $\rho_0-1\le\abs{w-z}_\infty\le\rho_0$ and $\abs{x-w}_\infty=\abs{x-z}_\infty-\abs{w-z}_\infty$. Then $w\in\Lambda_{i_0+1}^c$. The induction hypothesis for the chain starting at $i_0+1$ gives $\abs{x-w}_\infty>(m-1)\rho_0-(m-2)$, hence $\abs{x-z}_\infty>m\rho_0-(m-1)$.
\end{proof}

\begin{remark}[uses]\label{app:G1:twodomain-inst}
(I1) $X_{\rm out}=\Omega$, $X_{\rm in}=\Omega_2=B^k(\Lambda_2^{(k)})$, $A=\tilde B$: this is step S2 of Proposition~\ref{prop:prop1}(b). These lines occur only in terms whose vertices lie in $B^k(\Lambda_7^{(k)})$ \cite[p.~579; Prop.~2.6, p.~580]{Ba2}. Use Lemma~\ref{app:G1:nesting} with $\rho_0=r(t_k)$ and \cite[(2.8), p.~558]{Ba2}, with one block for the passage to $B^k(\cdot)$. Then these vertices have depth $\ge5r(t_k)-5$ in $\Omega_2$. The localisation cubes $\square(v)$ have side $\le2$ \cite[p.~420]{Ba3}, and the supports of the localisation functions lie within $4$ of $B^k(\Lambda_7^{(k)})$. So the supports have depth $\ge5r(t_k)-9\ge5R-9$. This is $\ge R_0+13$ when $R>R_0$ \cite[(2.7)]{Ba2} and $R\ge\frac{11}2$; both follow from the lower bounds on $R$ in Table~\ref{p:tab:A2}. Corollary~\ref{app:G1:twodomain-loc} therefore applies at every $k\ge1$.

(I2) $X_{\rm out}=\Omega$, $X_{\rm in}=\square$, $A=\tilde B$, given $\square\subset\Omega$ (Proposition~\ref{prop:prop1}(c)). The vertex localisations lie at depth $\ge r_\square-2$ in $\square$. So $r_\square\ge R_0+13+3\eta$ suffices. This follows from $r_\square\ge R_0+15$, which is implied by $r\ge3R_0+15$.

(I3) $X_{\rm out}=T_\eta$, $X_{\rm in}=\square$, $A=0$: this is the line $\delta G_k(\square,\eta_k\Z^3,0)$ of \cite[p.~433]{Ba3} on the torus. The margins are as in (I2).

On the printed factor: \cite[(2.5)]{Ba3} prints the factor $\mathrm{e}^{-\delta_0\dist(\Omega_2,\partial\Omega)}$, and the two-point decay multiplies both alternatives. Proposition~\ref{app:G1:twodomain} gives this factor when the distance is read as the distance of the vertex localisations to $\Omega_2^c$; this is the factor of \cite[(1.12), p.~573]{Ba4} for the inner domain. Read literally, as a distance between sets, the factor is not implied: for localisations near $\partial\Omega_2$ but deep in $\Omega$, $\delta$ is of order one. The uses (I1)--(I3) need only Corollary~\ref{app:G1:twodomain-loc}. Ba{\l}aban's rate $\delta_0$ is taken $\le\beta_D$.
\end{remark}

\begin{remark}[the printed factor; rates]\label{app:G1:twodomain-rem}
(i) \cite[(2.5)]{Ba3} prints, for $\delta G_k(\Omega,\Omega_2,\tilde B)$, the bound $O(\mathrm{e}^{-\delta_0\dist(\Omega_2,\partial\Omega)})\cdot\allowbreak\mathrm{e}^{-\delta_0\dist(\supp h,\supp h')}$; the two-point factor multiplies both alternatives of (2.5). Proposition~\ref{app:G1:twodomain} matches this factor only in the reading ``distance of the vertex localisations to $\Omega_2^c$'' (the factor of \cite[(2.26)]{Ba1}). In the literal set-distance reading it is not implied, and it is false in general: localisations near $\partial\Omega_2$ but deep in $\Omega$ give a $\delta G$ of order one. No step of the proof of Theorem~\ref{thm:T6} uses more than an $O(1)$ bound together with the two-point decay \cite[p.~432]{Ba3}, and both are proved.

(ii) The rates $\beta_U$ and $\beta_D$ enter \cite[(2.13)]{Ba3} as the rate of a line. One takes Ba{\l}aban's $\delta_0\le\min\{\beta_U,\beta_D\}$, a free shrink of a constant; the conditions that use $\delta_0$, such as $\delta_0R\ge4$, are imposed after it.
\end{remark}

\begin{remark}\label{app:G1:sandwich-missing}
Two cases are not treated, and neither is used: the remainder of the $B'$-expansion on the cube $\square$ of \cite[p.~433]{Ba3}, which concerns Neumann faces and is a different object; and the H\"older exponents $\frac12\le\alpha<1$ at $m\in\{4,5\}$ in Proposition~\ref{app:G1:sandwich}.
\end{remark}

\subsection{The inputs (W) and (K)}\label{app:G1:part-TWK}
\emph{Conventions.} $\mathcal X$ is a union of big blocks of $T_\eta$ with Neumann conditions, or $\mathcal X=T_\eta$. $\mathcal B_{\mathcal X}$ is the bond set: bonds with both ends in $\mathcal X$, or all bonds of $T_\eta$. Functions on sites and on bonds carry the measure $\eta^3$. For a bond $b=(b_-,b_+)$ in direction $\mu$, $(D\phi)(b)=\eta^{-1}(U_b\phi(b_+)-\phi(b_-))$, with $U_b$ orthogonal on $\R^2$. Let $\rho$ be real and $1$-Lipschitz for the Euclidean torus distance; every function that is $1$-Lipschitz for the $\ell^\infty$ torus distance qualifies. On sites put $Y^\beta:=\mathrm{e}^{\beta\rho}Y\mathrm{e}^{-\beta\rho}$. A bond function is weighted at $b_-$, i.e.\ by $\mathrm{e}^{\beta\rho(b_-)}$. Then $\abs{\rho(b_+)-\rho(b_-)}\le\eta$ on every bond, and $\abs{\rho(x)-\rho(x')}\le\sqrt3$ inside a unit block. Write $H_{\mathcal X}(A)=D^*D+m^2t^2+a_kP_k(A)$.

\begin{lemma}[real-part Combes--Thomas bound]\label{app:G1:WK-CT}
Let $Y=D^*D+v+aP$ on $\ell^2(\mathcal X,\eta^3;\R^2)$, where the bond set and the transporters $U_b$ are arbitrary, $v\ge0$ is a multiplication operator, $0\le a$, and $P\ge0$ satisfies $P(x,x')=0$ for $\abs{x-x'}>\sqrt3$ and $\sup_x\sum_{x'}\eta^3\abs{P(x,x')}\le1$. Assume $Y\ge\gamma>0$. Let $\abs\beta\le1$ with
\[
c(\beta):=3\beta^2\cosh\beta+a\big(\mathrm{e}^{\sqrt3\abs\beta}-1\big)\le\gamma/2.
\]
Then for every $F$ the vector $w:=\mathrm{e}^{\beta\rho}Y^{-1}\mathrm{e}^{-\beta\rho}F$ satisfies $\norm w\le2\gamma^{-1}\norm F$ and $\norm{Dw}^2+\ip w{vw}\le4\gamma^{-1}\norm F^2$. This holds uniformly in $\eta$, the domain, the boundary condition, $U$ and $v$.
\end{lemma}
\begin{proof}
Negative $\beta$ is the case $-\rho$, so let $\beta\ge0$. The vector $w$ solves $Y^\rho w=F$ with $Y^\rho:=\mathrm{e}^{\beta\rho}Y\mathrm{e}^{-\beta\rho}$.

\emph{Bond terms.} Fix a bond $b$, put $t:=\beta(\rho(b_+)-\rho(b_-))$, so $\abs t\le\beta\eta$, and put $z:=\bar w(b_-)U_bw(b_+)$. The contribution of $b$ to $\rRe\ip w{\mathrm{e}^{\beta\rho}D^*D\mathrm{e}^{-\beta\rho}w}$ is
\[
\eta^{1}\big[\abs{w(b_+)}^2+\abs{w(b_-)}^2-2\cosh t\,\rRe z\big]
=q_b(w)-2(\cosh t-1)\eta\,\rRe z.
\]
The odd part ($\sinh t$) is anti-Hermitian and drops out of the real part. Since $\cosh t-1\le\frac12\beta^2\eta^2\cosh\beta$, $2\abs{\rRe z}\le\abs{w(b_+)}^2+\abs{w(b_-)}^2$, and each site lies in at most $6$ bonds, the bond part is $\ge\norm{Dw}^2-3\beta^2\cosh\beta\norm w^2$.

\emph{Other terms.} $v$ commutes with $\mathrm{e}^{\beta\rho}$. The kernel of $P^\rho-P$ is $(\mathrm{e}^{\beta(\rho(x)-\rho(x'))}-1)P(x,x')$, so by Schur's test $\norm{P^\rho-P}\le \mathrm{e}^{\sqrt3\beta}-1$. Hence
\[
\rRe\ip w{Y^\rho w}\ge\ip w{Yw}-c(\beta)\norm w^2\ge\tfrac\gamma2\norm w^2.
\]

\emph{Conclusion.} Since $\rRe\ip w{Y^\rho w}=\rRe\ip wF\le\norm w\norm F$, the first bound follows. Then $\norm{Dw}^2+\ip w{vw}\le\ip w{Yw}\le\norm w\norm F+\frac\gamma2\norm w^2\le4\gamma^{-1}\norm F^2$.
\end{proof}

\begin{lemma}[input (W): weighted $\ell^2$ bounds]\label{app:G1:WK-W}
Let $A$ be real with $H_{\mathcal X}(A)-m^2t^2\ge\gamma_0$ (that is, (H0)), let $m^2t^2\ge0$, and let $G:=G_k(\mathcal X,A)$. Fix $0<\beta_*\le1$ with $c(\beta_*)\le\gamma_0/4$ (the $c$ of Lemma~\ref{app:G1:WK-CT} at $a=a_k$), and put $c':=\sqrt3\,\beta_*\mathrm{e}^{\beta_*}$. Then for every $1$-Lipschitz $\rho$ and $\abs\beta\le\beta_*$:
\[
\norm{G^\beta}\le\frac2{\gamma_0},\qquad \norm{(DG)^\beta},\ \norm{(GD^*)^\beta},\ \norm{(DGD^*)^\beta}\le K_W:=\max\Big\{\frac2{\gamma_0^{1/2}}+\frac{2c'}{\gamma_0},\ \frac{1+c'^2}{\min\{\frac12,\frac{\gamma_0}4\}}\Big\}.
\]
These hold uniformly in $\eta$, $k$, $\eps$, $m^2\ge0$, the shape of $\mathcal X$ (wrapping allowed) and the boundary condition. Consequently the block operator $\mathfrak DG\mathfrak D^*$ satisfies $\norm{(\mathfrak DG\mathfrak D^*)^\beta}\le2\gamma_0^{-1}+3K_W$.
\end{lemma}
\noindent The only input is (H0) \cite[(1.8), p.~573]{Ba4}. The factor $\sqrt3$ in $c'$ accounts for the three directions of the vector difference $D$; per direction the constant is $\mathrm{e}^{\beta_*}$.
\begin{proof}
\emph{(a) The bound on $G^\beta$.} Apply Lemma~\ref{app:G1:WK-CT} to $Y=H_{\mathcal X}(A)$, with $v=m^2t^2$, $P=P_k(A)$, $a=a_k$ and $\gamma=\gamma_0$. This is admissible because $c(\beta)\le\gamma_0/4$. The operator $G^\beta=(H_{\mathcal X}^\beta)^{-1}$ exists because $\rRe H_{\mathcal X}^\beta\ge\gamma_0/2$. The lemma gives $\norm{G^\beta}\le2/\gamma_0$ and $\norm{DG^\beta}\le2\gamma_0^{-1/2}$.
\smallskip

\emph{(b) The bound on $(DG)^\beta$.} With the bond weight at $b_-$,
\[
(DG)^\beta=D^\beta G^\beta,\qquad D^\beta=D+\beta K_\beta,\qquad (K_\beta\phi)(b)=(\beta\eta)^{-1}\big(\mathrm{e}^{-\beta(\rho(b_+)-\rho(b_-))}-1\big)U_b\phi(b_+).
\]
The coefficient is at most $\mathrm{e}^{\abs\beta}$ in modulus. The map $\phi\mapsto(\phi(b_+))_b$ has norm $1$ for each direction and $\sqrt3$ for the three directions together. Hence $\norm{\beta K_\beta}\le c'$, and $\norm{(DG)^\beta}\le2\gamma_0^{-1/2}+2c'\gamma_0^{-1}$.
\smallskip

\emph{(c) The bound on $(GD^*)^\beta$.} $G$ is real symmetric, so $(GD^*)^\beta=((DG)^{-\beta})^*$, and (b) at $-\beta$ applies.
\smallskip

\emph{(d) The bound on $(DGD^*)^\beta$.} Let $w:=G^\beta(D^*)^\beta g$, so that $H_{\mathcal X}^\beta w=(D^*)^\beta g$, and note $((D^*)^\beta)^*=D^{-\beta}$.
\begin{itemize}
\item From the computation in Lemma~\ref{app:G1:WK-CT}, together with $H_{\mathcal X}\ge D^*D$, $H_{\mathcal X}\ge\gamma_0$ and $c(\beta)\le\gamma_0/4$,
\[
\rRe\ip w{H_{\mathcal X}^\beta w}\ge\tfrac12\norm{Dw}^2+\tfrac{\gamma_0}4\norm w^2.
\]
\item On the other hand, $\rRe\ip w{H^\beta_Xw}=\rRe\ip{D^{-\beta}w}g\le(\norm{Dw}+c'\norm w)\norm g$.
\end{itemize}
Put $s:=(\norm{Dw}^2+\norm w^2)^{1/2}$. Then $\min\{\frac12,\frac{\gamma_0}4\}s^2\le(1+c'^2)^{1/2}s\norm g$. Hence
\[
\norm{(DGD^*)^\beta g}=\norm{D^\beta w}\le(\norm{Dw}+c'\norm w)\le(1+c'^2)\norm g/\min\{\tfrac12,\tfrac{\gamma_0}4\}.
\]
\smallskip

\emph{Uniformity.} Lemma~\ref{app:G1:WK-CT} uses no geometric property of the bond set or of the boundary condition. Its only inputs are $\abs{\rho(b_+)-\rho(b_-)}\le\eta$ and the range $\sqrt3$ of $P_k$; both hold on every union of whole unit blocks, including wrapping ones, and on $T_\eta$. The lower bound (H0) holds on every such $\mathcal X$, because the form $\ip f{H_Xf}$ dominates the sum over unit blocks of the intra-block Neumann form plus $a_k\abs{Q_kf}^2$, and each of these is $\ge\gamma_0\norm f^2_{\rm block}$ by \cite[(1.8)]{Ba4}.
\smallskip

\emph{Block operator.} The bound on $\mathfrak DG\mathfrak D^*$ sums the four blocks: $G$, the block $DG$, the block $GD^*$ and the block $DGD^*$, with the vector reading of $D$.
\end{proof}
\noindent At $\beta=0$ and for $\rho=\dist(\cdot,\supp f')$, (W) reproduces the literature bound \cite[Cor.~2.3, (2.30), p.~580]{Ba4}: $\abs{\ip f{DGD^*f'}}\le c_0\mathrm{e}^{-\delta_0\dist(\supp f,\supp f')}\norm f_2\norm{f'}_2$, with no restriction on the supports.

\begin{lemma}[slice sums]\label{app:G1:WK-slicesum}
Let $L\ge2$, $\eta=L^{-k}$ and $s_j:=L^j\eta$ for $0\le j\le k-1$. For $p>0$ and $\delta>0$ put
\[
c_\Sigma(p,\delta,L):=(1-L^{-p})^{-1}+\sup_{1\le u\le L}\sum_{m\ge0}(uL^m)^p\mathrm{e}^{-\delta uL^m/2}.
\]
Then for all $z\ge0$:
\[
\sum_{j=0}^{k-1}s_j^{-p}\mathrm{e}^{-\delta z/s_j}\le c_\Sigma(p,\delta,L)\max(z,\eta)^{-p}\mathrm{e}^{-\delta z/2}.
\]
\end{lemma}
\begin{proof}
Put $r:=\max(z,\eta)$ and note $s_j\le L^{-1}<1$.

\emph{Slices with $s_j\ge r$.} These contribute at most $r^{-p}\sum_{m\ge0}L^{-mp}\mathrm{e}^{-\delta z}$.

\emph{Slices with $s_j<r$.} These occur only if $r=z$. Put $u_j:=z/s_j>1$; the $u_j$ form a sequence $uL^m$ ($m\ge0$) with $u\in(1,L]$. Each term equals $z^{-p}u_j^p\mathrm{e}^{-\delta u_j}$, and $\mathrm{e}^{-\delta u_j}\le \mathrm{e}^{-\delta u_j/2}\mathrm{e}^{-\delta z/2}$ because $u_j\ge z$.

Adding the two parts gives the claim.
\end{proof}

\begin{lemma}[input (K): one-derivative pointwise kernel bounds]\label{app:G1:WK-K}
Let $A$ be real and regular on $\mathcal X$ in the sense of \cite[(2.23)]{Ba1}, with $e_k$ small as in Proposition~\ref{prop:UV-BaProp21}; under the hypothesis (R$_\vartheta$) of the $L^2$ sandwich of Subsection~\ref{app:G1:part-E2} this holds for $A\in\{\tilde B,\tilde A+\tilde B\}$. Let $0\le m^2t^2\le c$, $0<\alpha<1$ and $G:=G_k(\mathcal X,A)$. Let $\delta_0,c_0,R_0$ be the constants of Proposition~\ref{prop:UV-BaProp21}, and put $\delta_K:=\delta_0/4$. There are $c_K$ and $c_K(\alpha)$, depending only on $d,a,M_b,L$ (and on $\alpha$ for $c_K(\alpha)$), such that for all $x,w$ with $\rho_{\mathcal X}(x),\rho_{\mathcal X}(w)\ge R_0$:
\begin{itemize}
\item[(K0)] $\norm{G(x,w)}\le c_KK^{\delta_K}_2(x-w)$;
\item[(K1)] $\norm{D_xG(x,w)}\le c_KK^{\delta_K}_1(x-w)$, and by (S) $\norm{(GD^*)(w,x)}\le c_KK^{\delta_K}_1(x-w)$;
\item[(K2)] for $\abs{x-\tilde x}\le1$ with $\rho_{\mathcal X}(\tilde x)\ge R_0$:
\[
\norm{U(A(\Gamma_{\tilde x,x}))D_xG(x,w)-D_{\tilde x}G(\tilde x,w)}\le c_K(\alpha)\abs{x-\tilde x}^\alpha\big(K^{\delta_K}_{1-\alpha}(x-w)+K^{\delta_K}_{1-\alpha}(\tilde x-w)\big).
\]
\end{itemize}
On $\mathcal X=T_\eta$ the restriction $\rho_{\mathcal X}\ge R_0$ is void.
\end{lemma}
\noindent The lemma is used for scalar lines only; its inputs are Proposition~\ref{prop:UV-BaProp21} and Lemma~\ref{app:G1:WK-slicesum}.
\begin{proof}
\emph{Step 1 (slices).} By \cite[(2.6), p.~424]{Ba3},
\[
G=\sum_{j=0}^{k-1}G^\eta_{(j)},\qquad G^\eta_{(0)}=C^{(0),\eta},\qquad G^\eta_{(j)}=a_j^2(L^j\eta)^{-4}G^\eta_jQ_j^*C^{(j),L^j\eta}Q_jG^\eta_j\quad(j\ge1).
\]
This is the renormalisation-group identity \cite[(2.42)--(2.43)]{Ba1}, which is exact. Put $s_j:=L^j\eta\in[\eta,L^{-1}]$.

\emph{Step 2 (per-slice bounds).} Ba{\l}aban's slice bounds \cite[(2.10)--(2.11), p.~426]{Ba3} follow from \cite[Props.~2.1, 2.3]{Ba1} (Proposition~\ref{prop:UV-BaProp21}), with $\delta_1=\delta_0/2$, as follows. For each $j$, rescale to the $L^{-j}$-lattice, $x=s_jy$. Then:
\begin{itemize}
\item \cite[(2.25)]{Ba1} applied to $f=Q_j^*\delta_{y_0}$ bounds $(G_jQ_j^*)(y,y_0)$ and its derivative by $c_0\mathrm{e}^{-\delta_0\dist(y,B(y_0))}$; \cite[(2.24)]{Ba1} gives the same bound for the $\alpha$-H\"older quotient.
\item \cite[(2.34)]{Ba1} bounds $C^{(j)}$ by $c_0\mathrm{e}^{-\delta_0\abs{y_0-y_0'}}$.
\item The composition is a two-exponential sum, at most $c_d(\delta_0/2)^{-6}\mathrm{e}^{-(\delta_0/2)\abs{y-y'}}$.
\item Rescaling back multiplies the kernel by $s_j^{-1}$, each derivative by $s_j^{-1}$ and each H\"older quotient by $s_j^{-1-\alpha}$; the decay length becomes $s_j$.
\end{itemize}
Thus, for depths as in the statement and $n\in\{0,1\}$,
\[
\norm{D_x^nG^\eta_{(j)}(x,w)}\le c_1s_j^{-1-n}\mathrm{e}^{-\delta_1\abs{x-w}/s_j},
\]
and the H\"older quotient of $D_xG^\eta_{(j)}$ is at most $c_1s_j^{-2-\alpha}\mathrm{e}^{-\delta_1\dist(\{x,\tilde x\},w)/s_j}$. For $j=0$ the bound is \cite[(2.33)--(2.34)]{Ba1} at the $\eta$-scale. The constants are independent of $\eta$, $k$ and $\mathcal X$; they depend on the regularity of the background, which is inherited at scale $j$.

\emph{Depth.} \cite[Props.~2.1, 2.3]{Ba1} require depth $\ge R_0$ in scale-$j$ units, i.e.\ $\ge R_0s_j\le R_0/L$ in scale-$k$ units. So $\rho_{\mathcal X}\ge R_0$ at the two evaluation points (and at $\tilde x$) suffices for every $j$. Both points must be deep because the slice is evaluated as $(G_jQ_j^*)(x,\cdot)$ and, by (S), as $(G_jQ_j^*)(w,\cdot)^T$. On $\mathcal X=T_\eta$ the restriction is omitted \cite[p.~611]{Ba1}. All constants are uniform over unions of big blocks \cite[p.~411]{Ba3} and in the side of the torus (Lemma~\ref{app:G1:torusuniform}).

\emph{Step 3 (summation).}
\begin{itemize}
\item (K0), (K1): sum the per-slice bounds and apply Lemma~\ref{app:G1:WK-slicesum} with $p=1$ and $p=2$, $\delta=\delta_1$. The results are $K^{\delta_1/2}_2$ and $K^{\delta_1/2}_1$, and $\delta_1/2=\delta_K$.
\item (K2): fix $x$, $\tilde x$ and $j$.
\begin{itemize}
\item If $s_j\ge\abs{x-\tilde x}$, use the slice H\"older bound.
\item If $s_j<\abs{x-\tilde x}$, bound the two terms separately. $U$ is orthogonal and $\abs{x-\tilde x}^{-\alpha}<s_j^{-\alpha}$, so the two derivative bounds give $c_1s_j^{-2-\alpha}\big(\mathrm{e}^{-\delta_1\abs{x-w}/s_j}+\mathrm{e}^{-\delta_1\abs{\tilde x-w}/s_j}\big)$.
\end{itemize}
In both cases, using $\mathrm{e}^{-\delta_1\min(a,b)/s}\le \mathrm{e}^{-\delta_1a/s}+\mathrm{e}^{-\delta_1b/s}$, the quotient of slice $j$ is at most $c_1s_j^{-2-\alpha}\big(\mathrm{e}^{-\delta_1\abs{x-w}/s_j}+\mathrm{e}^{-\delta_1\abs{\tilde x-w}/s_j}\big)$. Summing with Lemma~\ref{app:G1:WK-slicesum} at $p=2+\alpha$ gives (K2), with $c_K(\alpha)=2c_1c_\Sigma(2+\alpha,\delta_1,L)$.
\end{itemize}
The split covers both readings of the range of the H\"older quotient in \cite[(2.24)]{Ba1}.

\emph{Step 4 (column bounds).} $G$ is real symmetric, so $(GD^*)(w,x)=(D_xG(x,w))^T$.
\end{proof}

\begin{remark}[sources and limits of (W) and (K)]\label{app:G1:WK-src}
\begin{itemize}
\item[(i)] (K) is used only for scalar lines. The vector-field slices of \cite[(2.6)]{Ba3} are not derived from \cite[Props.~2.1, 2.3]{Ba1}, and they are not used.
\item[(ii)] On Neumann domains (K) is stated only at points of depth $\ge R_0$. Near re-entrant edges a derivative bound uniform in $\eta$ is not expected (one expects a corner mode $r^{2/3}$). No kernel is evaluated there.

\item[(iii)] For $\delta G_k(\square,\eta\Z^3,0)$ the infinite lattice $\eta\Z^3$ enters only through the bounds (K) on $\eta\Z^3$ at $A=0$, that is, the bounds of \cite[Props.~2.1, 2.3]{Ba1} read on the infinite lattice.
\item[(iv)] (W) needs no depth, no shape condition and no regularity beyond (H0). It is uniform in $m^2\ge0$. (K) is uniform for $0\le m^2t^2\le c$.
\end{itemize}
\end{remark}

\subsection{The two-domain bounds at lattice spacing one half}\label{app:G1:ssec-k1margins}

Throughout this subsection $k=1$, so $\eta=\frac12$. Let $X_{\rm in}\subset X_{\rm out}$, where $X_{\rm in}$ is a union of unit blocks and $X_{\rm out}$ is a union of big blocks or $T_\eta$. Let $A$ be real with $\theta(A)\le\vartheta_R$. Assume that the inputs (W), (K) and (S) of Lemma~\ref{app:G1:WKinputs} hold on $X_{\rm in}$ and on $X_{\rm out}$; Lemma~\ref{app:G1:WKinputs} gives this when both are unions of big blocks or $T_\eta$. Let $0\le m^2t^2\le c$. Put $H_{\rm in/out}:=H_{X_{\rm in/out}}(A)$, $G_{\rm in/out}:=H_{\rm in/out}^{-1}$, $\delta:=G_{\rm out}-G_{\rm in}$ and $\rho:=\rho_{\rm in}=\dist_\infty(\cdot,X_{\rm in}^c)$. Since $X_{\rm out}^c\subset X_{\rm in}^c$, we have $\rho_{X_{\rm out}}\ge\rho$.

If $X_{\rm in}=T_\eta$, then $\delta=0$ and there is nothing to prove, so assume $X_{\rm in}\ne T_\eta$. We write $\abs\cdot$, $\abs\cdot_\infty$ and $\abs\cdot_1$ for the Euclidean, $\ell^\infty$ and $\ell^1$ torus distances. Kernels are taken with respect to $\eta^3$. For a matrix-valued function $F$ of one lattice variable put $\norm F_2:=(\sum_z\eta^3\norm{F(z)}^2)^{1/2}$; for a row or column of four such functions, the four norms are added.

For a kernel $Y$ put $Y_{ab}:=D^aYD^{b*}$, with $a,b\in\{0,1\}$ and every direction. The quotients $H_x$, $H_{x'}$ and $H_{xx'}$ are the $\alpha$-H\"older quotients of $Y_{1b}$ in $x$, of $Y_{a1}$ in $x'$, and the mixed one of $Y_{11}$. They are taken over partners $0<\abs{x-\tilde x}\le1$ (respectively $0<\abs{x'-\tilde x'}\le1$), with transport $U(A(\Gamma))$ along a shortest lattice contour. As in the definition of $\mathfrak N_{\alpha,\beta}$, $\norm{hYh'}_{1,\alpha}$ is the maximum over $a,b$ of the suprema of $\norm{(hYh')_{ab}}$, $aH_x$, $bH_{x'}$ and $abH_{xx'}$ for the kernel $h(x)Y(x,x')h'(x')$. This is the norm of the paper. It is the two-variable form of \cite[(1.32), p.~420]{Ba3}, with the sum there replaced by a maximum and the quotients restricted to $\abs{x-\tilde x}\le1$. Since $U$ is orthogonal, a quotient with $\abs{x-\tilde x}>1$ is at most twice the corresponding supremum, so the norm of \cite[(1.32)]{Ba3} is at most a fixed multiple of this one. When the points are restricted to a set, the partners are restricted to the same set. Put
\begin{gather*}
c_3:=\big(\tfrac12+4\sqrt3/\delta_K\big)^3,\qquad c_W:=2/\gamma_0+3K_W,\\
c_{\rm end}(\alpha):=c_{\mathfrak r}\,\mathrm{e}^{\delta_K(\sqrt3+1)}\max\{c_K,2c_K(\alpha)\}\,c_3^{1/2},\qquad C_D(\alpha):=c_Wc_{\rm end}(\alpha)^2,
\end{gather*}
where $c_{\mathfrak r}$ is defined in Lemma~\ref{app:G1:k1-collar} and $c_K,c_K(\alpha),\delta_K,\gamma_0,K_W,\beta_*$ are the constants of Lemma~\ref{app:G1:WKinputs}.

\begin{lemma}[collars and commutators at spacing one half]\label{app:G1:k1-collar}
Let $f$ and $\varphi$ be the profiles of Lemma~\ref{app:G1:collar}: $f\in C^\infty(\R)$ with $0\le f\le1$, $f=0$ on $(-\infty,0]$ and $f=1$ on $[1,\infty)$, and $\varphi\in C^\infty_c([-\frac14,\frac14]^3)$ with $\varphi\ge0$ and $\int\varphi=1$. With these profiles, $C_D(\alpha)$ is the constant of Lemma~\ref{app:G1:I3r-DGM}. For a unit block $b$ with lowest corner $y_b$ put $\hat b:=y_b+[0,1)^3$ and $\hat X_{\rm in}:=\bigcup_{b\subset X_{\rm in}}\hat b$. Let $\rho^{\rm c}(u):=\dist_\infty(u,\hat X_{\rm in}^c)$ on the continuum torus, $\tilde\rho:=\rho^{\rm c}*\varphi$, and $\chi_s:=f(\tilde\rho-s)$ on lattice points. Put $c_1':=\norm{f'}_\infty$, $c_2':=\norm{f''}_\infty+\norm{f'}_\infty\max_\nu\norm{\partial_\nu\varphi}_{L^1}$ and $c_{\mathfrak r}:=6c_1'+3c_2'+ac_1'$. Then:
\begin{enumerate}
\item[(i)] $\tilde\rho$ is $1$-Lipschitz for $\abs\cdot_\infty$, and $\rho-\frac34\le\tilde\rho\le\rho+\frac14$ at lattice points;
\item[(ii)] $\chi_s=0$ on $\{\rho\le s-\frac14\}$ and $\chi_s=1$ on $\{\rho\ge s+\frac74\}$. If $\chi_s(u)\ne\chi_s(u')$ and $\abs{u-u'}_\infty\le\ell$, then $\tilde\rho(u),\tilde\rho(u')\in(s-\ell,s+1+\ell)$;
\item[(iii)] $\abs{\partial^\eta_\mu\chi_s}\le c_1'$ and $\abs{\partial^\eta_\mu\partial^\eta_\mu\chi_s}\le c_2'$;
\item[(iv)] let $s\ge1$ and $\phi_\mu:=\partial^\eta_\mu\chi_s$. Then $H_{\rm out}(\chi_sg)=H_{\rm in}(\chi_sg)$ for every $g$ on $X_{\rm in}$, extended by zero. Hence $C_s:=H_{\rm out}\chi_s-\chi_sH_{\rm in}=[H_{\rm in},\chi_s]$, with kernel $C_s(w,z)=H_{\rm in}(w,z)(\chi_s(z)-\chi_s(w))$;
\item[(v)] $C_s=-\mathfrak r_s\mathfrak D=\mathfrak D^*\mathfrak r_s^*$, where $\mathfrak Dg=(g,D_1g,D_2g,D_3g)$ and $\mathfrak r_s=(\mathfrak r_0,\dots,\mathfrak r_3)$ is given by
\[
\mathfrak r_0:=\sum_\mu(\partial^\eta_\mu\phi_\mu)(\cdot-\eta e_\mu)-a_1[P_1(A),\chi_s],\qquad(\mathfrak r_\mu g)(x):=\phi_\mu(x)g(x)+\phi_\mu(x-\eta e_\mu)U_{x-\eta e_\mu,\mu}(A)^{-1}g(x-\eta e_\mu).
\]
With $c_{\mathfrak r,i}:=\sup_z\sum_w\eta^3\norm{\mathfrak r_i(w,z)}$ we have $\sum_ic_{\mathfrak r,i}\le c_{\mathfrak r}$. If $C_s(w,z)\ne0$, then $\rho(z)<s+\frac94$. If $\mathfrak r_i(w,z)\ne0$, then $\abs{w-z}\le\frac{\sqrt3}2$ and $w,z\in S_s:=\{s-\frac74<\rho<s+\frac{13}4\}$.
\end{enumerate}
\end{lemma}
\begin{proof}
(i) The cubes $\hat b$ tile the torus. So a lattice point outside $X_{\rm in}$ lies outside $\hat X_{\rm in}$, and $\rho^{\rm c}\le\rho$ at lattice points. The distance $\rho^{\rm c}(z)$ is attained at a point $y$ of some $\overline{\hat b'}$ with $b'\not\subset X_{\rm in}$, and $b'\cap X_{\rm in}=\emptyset$. The lattice points of $b'$ are $\eta$-dense in $\overline{\hat b'}$ for $\abs\cdot_\infty$, so $\rho\le\rho^{\rm c}+\frac12$. A distance function is $1$-Lipschitz. Mollification by $\varphi$ preserves this and moves values by at most $\frac14$.

(ii) The first claim follows from (i). For the second: if $\tilde\rho(u)\le s-\ell$, then $\tilde\rho(u')\le s$ and both values of $\chi_s$ vanish; if $\tilde\rho(u)\ge s+1+\ell$, both values equal $1$. The same holds with $u,u'$ exchanged.

(iii) The first bound is the mean value theorem, since $\tilde\rho$ is $1$-Lipschitz. For the second, put $g(t):=f(\tilde\rho(z+te_\mu)-s)$. Then $\eta^{-2}(g(2\eta)-2g(\eta)+g(0))=g''(\xi)$ for some $\xi$. Here $g''=f''(\partial_\mu\tilde\rho)^2+f'\partial_\mu^2\tilde\rho$, $\abs{\partial_\mu\tilde\rho}\le1$ and $\partial^2_\mu\tilde\rho=(\partial_\mu\rho^{\rm c})*\partial_\mu\varphi$, with $\abs{\partial_\mu\rho^{\rm c}}\le1$ almost everywhere.

(iv) If $\chi_s(y)\ne0$, then $\rho(y)>s-\frac14\ge\frac34>\eta$. So the six bonds at $y$ are bonds of $X_{\rm in}$, and the unit block of $y$ lies in $X_{\rm in}$. A bond of $X_{\rm out}$ that is not a bond of $X_{\rm in}$ has both endpoints at depth $\le\eta$, hence outside $\supp\chi_s$. Its term in $D^*D$ involves only the values at its endpoints, so it contributes nothing on $\chi_sg$. $P_1(A)$ acts inside unit blocks, and a unit block lies in $X_{\rm in}$ or is disjoint from it. Hence $H_{\rm out}(\chi_sg)=H_{\rm in}(\chi_sg)$, extended by zero, and the kernel formula follows.

(v) The Leibniz rules $D_\mu(\chi g)=\chi(\cdot+\eta e_\mu)D_\mu g+\phi_\mu g$ and $D^*_\mu(\chi(\cdot+\eta e_\mu)g)=\chi D^*_\mu g-\phi_\mu g$ give $[D^*_\mu D_\mu,\chi]=D^*_\mu\phi_\mu-\phi_\mu D_\mu$. From $D^*_\mu g(x)=\eta^{-1}(U^{-1}_{x-\eta e_\mu,\mu}g(x-\eta e_\mu)-g(x))$ and $U^{-1}_{x-\eta e_\mu,\mu}g(x-\eta e_\mu)=g(x)-\eta U^{-1}_{x-\eta e_\mu,\mu}(D_\mu g)(x-\eta e_\mu)$ we get
\[
D^*_\mu(\phi_\mu g)(x)=-\phi_\mu(x-\eta e_\mu)U^{-1}_{x-\eta e_\mu,\mu}(D_\mu g)(x-\eta e_\mu)-(\partial^\eta_\mu\phi_\mu)(x-\eta e_\mu)g(x).
\]
Since $[m^2t^2,\chi]=0$, collecting terms gives $C_s=-\mathfrak r_s\mathfrak D$. $H_{\rm in}$ is self-adjoint and $\chi_s$ is real, so $C_s^*=-C_s=\mathfrak D^*\mathfrak r_s^*$.

Column sums: each of the two parts of $\mathfrak r_\mu$ gives at most $c_1'$, and the multiplication part of $\mathfrak r_0$ at most $3c_2'$, by (iii). For the last part, $P_1(A)=Q_1(A)^*Q_1(A)$, where $Q_1(A)$ averages over unit blocks (volume $1$) with orthogonal transporters. So $P_1(A)(w,z)$ is orthogonal for $w,z$ in one unit block and zero otherwise. Moreover $\abs{\chi_s(z)-\chi_s(w)}\le c_1'\abs{z-w}_\infty\le c_1'$, and $a_1=a$.

Supports. If $C_s(w,z)\ne0$, then $\chi_s(w)\ne\chi_s(z)$ and $\abs{w-z}_\infty\le\frac12$. Then (ii) gives $\tilde\rho(z)<s+\frac32$, and (i) gives $\rho(z)<s+\frac94$.

Now let $\mathfrak r_i(w,z)\ne0$, and first consider $\mathfrak r_\mu$ or the first part of $\mathfrak r_0$. Then $z\in\{w,w-\eta e_\mu\}$, and $\chi_s$ differs at two adjacent points $u,u'\in\{w-\eta e_\mu,w,w+\eta e_\mu\}$. So $\tilde\rho(u)\in(s-\frac12,s+\frac32)$, and since $\abs{w-u}_\infty,\abs{z-u}_\infty\le1$, $\tilde\rho(w),\tilde\rho(z)\in(s-\frac32,s+\frac52)$. For the part $[P_1,\chi_s]$ we have $\abs{w-z}_\infty\le\frac12$ and $\chi_s(w)\ne\chi_s(z)$, so $\tilde\rho(w),\tilde\rho(z)\in(s-\frac12,s+\frac32)$. By (i), $\rho\in[\tilde\rho-\frac14,\tilde\rho+\frac34]$, which gives $w,z\in S_s$. In all cases $\abs{w-z}_\infty\le\frac12$, hence $\abs{w-z}\le\frac{\sqrt3}2$.
\end{proof}

\begin{proposition}[two-domain bound at spacing one half]\label{app:G1:k1-core}
Let $\mathcal P:=\{z:\rho(z)\ge R_0+11\}$, $0<\alpha<1$, $0<\beta\le\beta_D=\min\{\beta_*,\delta_K/4\}$ and $Y:=\delta$. For $x,x'\in\mathcal P$, with partners $\tilde x,\tilde x'\in\mathcal P$, and all $a,b$,
\[
\mathrm{e}^{\beta\abs{x-x'}}\max\{\norm{Y_{ab}(x,x')},aH_x,bH_{x'},abH_{xx'}\}\le C_D(\alpha)\exp\!\big(-\tfrac{\delta_K}2[(\rho(x)-R_0-7)+(\rho(x')-R_0-10)]\big).
\]
\end{proposition}
\begin{proof}
\emph{Collars.} Put $\chi:=\chi_{R_0+3}$ and $\psi:=\chi_{R_0+13/2}$, with commutators $C_\chi=-\mathfrak r_\chi\mathfrak D$ and $C_\psi:=[H_{\rm in},\psi]=\mathfrak D^*\mathfrak r_\psi^*$; the latter is Lemma~\ref{app:G1:k1-collar} with $X_{\rm out}=X_{\rm in}$. By Lemma~\ref{app:G1:k1-collar}, $S_\chi\subset\{R_0+\frac54<\rho<R_0+\frac{25}4\}$ and $S_\psi\subset\{R_0+\frac{19}4<\rho<R_0+\frac{39}4\}$. If $C_\chi(w,z)\ne0$, then $\rho(z)<R_0+\frac{21}4$, so $\psi(z)=0$. Moreover $\chi=\psi=1$ on $\{\rho\ge R_0+\frac{33}4\}$, and this set contains the points of $\mathcal P$ and their bond partners, which have $\rho\ge R_0+\frac{21}2$. Since $\rho$ is $1$-Lipschitz, for $z\in\mathcal P$:
\begin{gather*}
\dist(z,S_\chi)\ge\dist_\infty(z,S_\chi)>\rho(z)-R_0-\tfrac{25}4\ge\max\{\rho(z)-R_0-7,1\},\\
\dist(z,S_\psi)>\rho(z)-R_0-\tfrac{39}4\ge\max\{\rho(z)-R_0-10,1\}.
\end{gather*}
Every point of $S_\chi\cup S_\psi$ has $\rho\ge R_0$, so (K) applies at all pairs used below, on $X_{\rm out}$ and on $X_{\rm in}$.

\emph{Representation.} By Lemma~\ref{app:G1:k1-collar}(iv), $G_{\rm out}C_\chi G_{\rm in}=\chi G_{\rm in}-G_{\rm out}\chi$. So $\delta(y,y')=-(G_{\rm out}C_\chi G_{\rm in})(y,y')$ if $\chi(y)=\chi(y')=1$. Likewise $G_{\rm in}C_\psi G_{\rm in}=\psi G_{\rm in}-G_{\rm in}\psi$, so $G_{\rm in}(z,y')=-(G_{\rm in}C_\psi G_{\rm in})(z,y')$ if $\psi(z)=0$ and $\psi(y')=1$. In the first identity only $z$ with $\psi(z)=0$ occur. Hence
\[
\delta=G_{\rm out}C_\chi G_{\rm in}C_\psi G_{\rm in}=-G_{\rm out}\mathfrak r_\chi(\mathfrak DG_{\rm in}\mathfrak D^*)\mathfrak r_\psi^*G_{\rm in}
\]
at all pairs $(y,y')$ of points of $\mathcal P$ or their bond partners. These are the only points at which $D^a$ and $D^{b*}$ evaluate $\delta$. Therefore
\[
Y_{ab}(x,x')=-\textstyle\sum_{u,u'}\eta^6E_L(x,u)\mathfrak M(u,u')E_R(u',x'),\qquad E_L:=D^aG_{\rm out}\mathfrak r_\chi,\ \mathfrak M:=\mathfrak DG_{\rm in}\mathfrak D^*,\ E_R:=\mathfrak r_\psi^*G_{\rm in}D^{b*}.
\]

\emph{Left end.} $E_{L,i}(x,z)=\sum_w\eta^3D^a_xG_{\rm out}(x,w)\mathfrak r_{\chi,i}(w,z)$, with $w,z\in S_\chi$ and $\abs{w-z}\le\frac{\sqrt3}2$. Since $\abs{x-w}\ge1$, (K0)--(K1) give $\norm{D^a_xG_{\rm out}(x,w)}\le c_K\mathrm{e}^{-\delta_K\abs{x-w}}\le c_K\mathrm{e}^{\delta_K\sqrt3}\mathrm{e}^{-\delta_K\abs{x-z}}$. Hence $\norm{E_{L,i}(x,z)}\le c_{\mathfrak r,i}c_K\mathrm{e}^{\delta_K\sqrt3}\mathrm{e}^{-\delta_K\abs{x-z}}1_{S_\chi}(z)$. Since $2(\delta_K-\beta)\ge\frac32\delta_K$,
\[
\sum_i\norm{\mathrm{e}^{\beta\abs{\cdot-x}}E_{L,i}(x,\cdot)}_2\le c_{\mathfrak r}c_K\mathrm{e}^{\delta_K\sqrt3}\mathrm{e}^{-\frac{\delta_K}2\dist(x,S_\chi)}\Big(\sum_z\eta^3\mathrm{e}^{-\frac{\delta_K}2\abs{z-x}}\Big)^{1/2}.
\]
Passing to shortest lifts (an injection into $x+\eta\Z^3$) and using $\abs n\ge\abs n_1/\sqrt3$, the sum is at most $(\sum_{m\in\Z}\eta\,\mathrm{e}^{-\delta_K\eta\abs m/(2\sqrt3)})^3\le(\eta+4\sqrt3/\delta_K)^3=c_3$. So the left side is $\le c_{\rm end}(\alpha)\mathrm{e}^{-\frac{\delta_K}2(\rho(x)-R_0-7)}$.

For $a=1$ and a partner $\tilde x\in\mathcal P$, the transported difference $U(A(\Gamma_{\tilde x,x}))E_L(x,\cdot)-E_L(\tilde x,\cdot)$ has, by (K2) and $\abs{\tilde x-w}\ge1$, kernel factor at most $c_K(\alpha)\abs{x-\tilde x}^\alpha(\mathrm{e}^{-\delta_K\abs{x-w}}+\mathrm{e}^{-\delta_K\abs{\tilde x-w}})\le2c_K(\alpha)\mathrm{e}^{\delta_K}\abs{x-\tilde x}^\alpha\mathrm{e}^{-\delta_K\abs{x-w}}$. The same computation bounds its weighted norm by $c_{\rm end}(\alpha)\abs{x-\tilde x}^\alpha\mathrm{e}^{-\frac{\delta_K}2(\rho(x)-R_0-7)}$.

\emph{Right end.} By (S), $(G_{\rm in}D^{b*})(v,x')=(D^b_{x'}G_{\rm in}(x',v))^T$, and $E_{R,i}(z,x')=\sum_v\eta^3\mathfrak r_{\psi,i}(v,z)^T(G_{\rm in}D^{b*})(v,x')$, with $v,z\in S_\psi$ and $\abs{x'-v}\ge1$. The computation of the left end, with column sums over $v$, gives
\[
\sum_i\norm{\mathrm{e}^{\beta\abs{\cdot-x'}}E_{R,i}(\cdot,x')}_2\le c_{\rm end}(\alpha)\mathrm{e}^{-\frac{\delta_K}2(\rho(x')-R_0-10)}.
\]
The same bound times $\abs{x'-\tilde x'}^\alpha$ holds for the transported difference in $x'$, which is the transpose of the difference in (K2).

\emph{Middle.} Let $\xi,\xi'$ be unit vectors, $F(u):=E_L(x,u)^T\xi$ and $F'(u'):=E_R(u',x')\xi'$. By (W) on $X_{\rm in}$, with the $1$-Lipschitz weight $\abs{\cdot-x}$ and $\abs\beta\le\beta_*$, each block of $\mathrm{e}^{-\beta\abs{\cdot-x}}\mathfrak M\,\mathrm{e}^{\beta\abs{\cdot-x}}$ has norm $\le c_W$. Using $\mathrm{e}^{-\beta\abs{u'-x}}\le\mathrm{e}^{-\beta\abs{x-x'}}\mathrm{e}^{\beta\abs{u'-x'}}$ and Cauchy--Schwarz in each block,
\[
\abs{\xi^TY_{ab}(x,x')\xi'}\le\mathrm{e}^{-\beta\abs{x-x'}}\Big(\sum_i\norm{\mathrm{e}^{\beta\abs{\cdot-x}}E_{L,i}(x,\cdot)}_2\Big)c_W\Big(\sum_j\norm{\mathrm{e}^{\beta\abs{\cdot-x'}}E_{R,j}(\cdot,x')}_2\Big).
\]
This proves the bound for $\norm{Y_{ab}}$. $Y_{ab}(x,x')$ is linear in $E_L(x,\cdot)$ and in $E_R(\cdot,x')$ separately, and $\mathfrak M$ does not depend on $x,x'$. So $H_x$, $H_{x'}$ and $H_{xx'}$ follow by replacing one or both ends by their transported differences.
\end{proof}

\begin{lemma}[Leibniz step at spacing one half]\label{app:G1:k1-leib}
Let $0\le h\le1$ with $\abs{\partial^\eta_\mu h}\le c$ for all $\mu$, and put $N:=\{y:\dist(y,\supp h)\le\frac32\}$. Let $Z$ be a matrix-valued function, and let $M\ge0$ be such that the following hold:
\begin{itemize}
\item $\norm{Z(y)},\norm{D_\mu Z(y)}\le M$ for $y\in N$;
\item $\abs{y-\tilde y}^{-\alpha}\norm{U(A(\Gamma_{\tilde y,y}))D_\mu Z(y)-D_\mu Z(\tilde y)}\le M$ whenever $\dist(y,\supp h)\le\eta$ and $0<\abs{y-\tilde y}\le1$.
\end{itemize}
Then $\norm{hZ}\le M$ and $\norm{D_\mu(hZ)}\le(1+c)M$ everywhere. For $0<\abs{y-\tilde y}\le1$, the quotient $\abs{y-\tilde y}^{-\alpha}\norm{U\,D_\mu(hZ)(y)-D_\mu(hZ)(\tilde y)}$ is at most $2(1+2c)M$ if $0<\alpha\le\frac12$, and at most $2(1+3c)M$ if $0<\alpha<1$.
\end{lemma}
\begin{proof}
We have $D_\mu(hZ)=F_1\,D_\mu Z+F_2\,Z$ with $F_1:=h(\cdot+\eta e_\mu)$ and $F_2:=\partial^\eta_\mu h$. Each $F_t$ vanishes outside $\{\dist(\cdot,\supp h)\le\eta\}$. This gives the two sup bounds. For the quotient let $d:=y-\tilde y$. We may assume that some $F_t$ is non-zero at $y$ or at $\tilde y$; then $y,\tilde y\in N$. Write
\[
U\,F_1(y)D_\mu Z(y)-F_1(\tilde y)D_\mu Z(\tilde y)=F_1(y)\big[U\,D_\mu Z(y)-D_\mu Z(\tilde y)\big]+\big(F_1(y)-F_1(\tilde y)\big)D_\mu Z(\tilde y).
\]
The first term vanishes unless $\dist(y,\supp h)\le\eta$, and then it is $\le M\abs d^\alpha$. Along a lattice path of $\abs d_1/\eta$ steps, $\abs{F_1(y)-F_1(\tilde y)}\le\min\{1,c\abs d_1\}$. Since $U$ is orthogonal, $\norm{U\,F_2(y)Z(y)-F_2(\tilde y)Z(\tilde y)}\le2cM$. So the quotient is at most $M[1+(\min\{1,c\abs d_1\}+2c)\abs d^{-\alpha}]$.

Since $d\in\frac12\Z^3$, we have $\abs d^2\in\{\frac14,\frac12,\frac34,1\}$ and $\abs d_1\le\sqrt3\abs d$. Let $\alpha\le\frac12$.
\begin{itemize}
\item If $\abs d=\frac12$, the bracket is $\le1+\sqrt2(\frac c2+2c)\le1+4c$.
\item Otherwise $\abs d^{-\alpha}\le2^{1/4}$. If $c\le1/\sqrt3$, the bracket is $\le1+2^{1/4}(\sqrt3+2)c\le1+4.44c\le2+4c$. If $c>1/\sqrt3$, it is $\le1+2^{1/4}+2^{5/4}c\le2+4c$, because $2^{1/4}-1\le(4-2^{5/4})/\sqrt3$.
\end{itemize}
Now let $\alpha<1$.
\begin{itemize}
\item If $\abs d=\frac12$, the bracket is $\le1+2(\frac c2+2c)=1+5c\le2+6c$.
\item Otherwise $\abs d^{-\alpha}\le\sqrt2$, and the bracket is $\le1+\sqrt2(\sqrt3+2)c\le2+6c$ if $c\le1/\sqrt3$, and $\le1+\sqrt2+2\sqrt2c\le2+6c$ if $c>1/\sqrt3$.
\end{itemize}
\end{proof}

\begin{lemma}[margins of the two-domain bounds at the first step]\label{app:K:L18-1}
Assume the setting of this subsection, with $k=1$, so that $\eta_1=\frac12$. Let $0<\alpha<1$ and $0<\beta\le\beta_D$, and put $X_{ab}=D^a\delta D^{b*}$. Then:
\begin{enumerate}
\item[(a)] for all $x,x'$, and their H\"older partners, with $\rho_{\rm in}(x),\rho_{\rm in}(x')\ge R_0+12$,
\[
\mathrm{e}^{\beta\abs{x-x'}}\max\{\norm{X_{ab}(x,x')},aH_x,bH_{x'},abH_{xx'}\}\le C_D(\alpha)\exp\!\big(-\tfrac{\delta_K}2[(\rho_{\rm in}(x)-R_0-7)+(\rho_{\rm in}(x')-R_0-10)]\big);
\]
\item[(b)] let $0\le h,h'\le1$ satisfy $\abs{\partial^{\eta}h},\abs{\partial^{\eta}h'}\le c_h$ and $\dist_\infty(\supp h\cup\supp h',X_{\rm in}^c)\ge R_0+13$. If $0<\alpha_0\le\frac12$, then
\[
\norm{h\,\delta\,h'}_{1,\alpha_0}\le4(1+2c_h)^2C_D(\alpha_0)\,\mathrm{e}^{-\beta\dist(\supp h,\supp h')}.
\]
For $0<\alpha_0<1$ the same bound holds with $4(1+3c_h)^2$ in place of $4(1+2c_h)^2$.
\end{enumerate}
\end{lemma}
\begin{proof}
(a) The points lie in $\{\rho\ge R_0+12\}$. A partner satisfies $\abs{x-\tilde x}_\infty\le\abs{x-\tilde x}\le1$, and $\rho$ is $1$-Lipschitz for $\abs\cdot_\infty$, so $\rho(\tilde x)\ge R_0+11$. Hence points and partners lie in $\mathcal P$, whether or not the partners are themselves assumed to have depth $\ge R_0+12$, and (a) is Proposition~\ref{app:G1:k1-core}.

(b) Put $W:=\delta$, and let $N$, $N'$ be the sets of Lemma~\ref{app:G1:k1-leib} for $h$ and $h'$. Points of $N\cup N'$ lie within $\ell^\infty$-distance $\frac32$ of the supports, so they have $\rho\ge R_0+\frac{23}2$ and lie in $\mathcal P$. For $y,\tilde y\in N$ and $y',\tilde y'\in N'$, the two gains of Proposition~\ref{app:G1:k1-core} add up to at least $\frac{\delta_K}2(\frac92+\frac32)=3\delta_K$. Moreover $\abs{y-y'}\ge\dist(\supp h,\supp h')-3$, and $3\beta\le3\delta_K$. Hence every quantity $\norm{W_{a'b'}}$, $H_x$, $H_{x'}$, $H_{xx'}$ evaluated at such base points and partners is at most
\[
M_0:=C_D(\alpha_0)\,\mathrm{e}^{-\beta\dist(\supp h,\supp h')}.
\]

Fix the second variable, either a point $x'$ or a pair $(x',\tilde x')$. The proof of Lemma~\ref{app:G1:k1-leib}, applied to the right variable, writes each of the following as a finite sum $\sum_s\kappa_sV_s(x)$:
\begin{itemize}
\item the kernel $\Phi_b(x,x'):=(Wh'D^{b*})(x,x')$, so that $(hWh')_{ab}=D^a_x(h\Phi_b)$; by $(Wh'D^*_\nu)(x,x')=h'(x'+\eta e_\nu)(WD^*_\nu)(x,x')+(\partial^\eta_\nu h')(x')W(x,x')$ its coefficients are $h'(x')$, $h'(x'+\eta e_\nu)$ or $\partial^\eta_\nu h'(x')$;
\item for $b=1$, its transported quotient in $x'$.
\end{itemize}
The scalar coefficients $\kappa_s$ do not depend on $x$, and they satisfy $\sum_s\abs{\kappa_s}\le1+c_h$ for a component and $\sum_s\abs{\kappa_s}\le\Lambda$ for a quotient. Here $\Lambda:=2(1+2c_h)$ if $\alpha_0\le\frac12$ and $\Lambda:=2(1+3c_h)$ otherwise. Each $V_s(x)$ is one of the following, multiplied on the right by a fixed orthogonal matrix:
\begin{itemize}
\item $W_{a'b'}(x,y')$ with $y'\in\{x',\tilde x'\}\subset N'$;
\item the $x'$-quotient of $W_{a'1}(x,\cdot)$ at $(x',\tilde x')$.
\end{itemize}
The derivatives $D_\mu$ and the transported $x$-differences of $V_s$ are again quantities of the list above; for quotient-type $V_s$ they are of type $H_{x'}$ and $H_{xx'}$. So $Z:=\sum_s\kappa_sV_s$ satisfies the hypotheses of Lemma~\ref{app:G1:k1-leib} with $M:=(\sum_s\abs{\kappa_s})M_0$, by the triangle inequality. That lemma, in the variable $x$, gives:
\begin{itemize}
\item each component $(hWh')_{ab}$ is at most $(1+c_h)^2M_0$;
\item $H_x$ and $H_{x'}$ are at most $(1+c_h)\Lambda M_0$;
\item $H_{xx'}$ is at most $\Lambda^2M_0$.
\end{itemize}
Every term vanishes unless its points lie in $N\times N'$, where the bounds above were used. Since $1+c_h\le\Lambda$ and $\Lambda^2\le4(1+2c_h)^2$, respectively $\Lambda^2=4(1+3c_h)^2$, this proves (b).
\end{proof}

\section{\texorpdfstring{The ingredients of Theorem~\ref{thm:T6}}{The ingredients of the vacuum bound for the slice models}}\label{p:appE}

This appendix proves the ingredients of the proof of Theorem~\ref{thm:T6} listed in Table~\ref{tab:T6-inputs}, except the region entropy, which is Proposition~\ref{app:G1:I6}. Subsection~\ref{app:G1:part-E3} treats the slice photon: the domination \eqref{eq:T6-GammaB}, the positivity of $C_B(u)$ and the degree in $u$ of $e(\eps;u;M)$. Subsection~\ref{app:G1:part-D3} proves the bounds on the final c-numbers that are uniform in the torus, the bounds on the $A\cdot A$ vertices and the tail beyond the order $\nbar$, together with Ba{\l}aban's line bounds on the torus on which they rest. Subsection~\ref{app:G1:part-E4} proves the ingredients themselves (Propositions~\ref{app:G1:I1}, \ref{app:G1:I2}, \ref{app:G1:I4}, \ref{app:G1:I5}, \ref{app:G1:I3}, \ref{app:G1:I7}, \ref{app:G1:I8}, \ref{app:G1:I9}). The torus is fixed and the letter $j$ is a slice or scale index, as stated at the start of each subsection.

\subsection{The slice photon: domination, positivity and the degree in $u$}\label{app:G1:part-E3}

We prove three statements of Subsection~\ref{ssec:T6-statement}: the domination in \eqref{eq:T6-GammaB}, the positivity of $C_B(u)$ for $u>-1/\bar\delta$, and the bound on the $u$-degree of $e(\eps;u;M)$.

\emph{Conventions.} Throughout, $T$ is one of the tori $T_j$ of Subsection~\ref{ssec:tori}, of side $\ell_T$, and $T_\eps=T\cap\eps\Z^3$ with $\eps=L^{-K}\eps_0$, $K\ge1$ ($L=2$ in Theorem~\ref{thm:T6}). The index $j$ below is the slice index of \eqref{eq:T6-GammaB}, not the torus index. Put $t_k:=L^k\eps$, so $t_K=\eps_0$. Vector fields $A=(A_\mu)_{\mu=1,2,3}$ are triples of real site functions on $T_\eps$, with the inner product $\ip{\cdot}{\cdot}$ of Subsection~\ref{ssec:action}. Operator inequalities are in the sense of quadratic forms on this finite-dimensional real space. The lattice Laplacian $-\Delta_\eps$ of $T_\eps$ acts on each component (Lemma~\ref{app:G1:X}), $M_*:=-\Delta_\eps+\mu_*^2$ ($\mu_*>0$), and $C_{\mu_*}=M_*^{-1}$. In this subsection $a_k$ always denotes King's block constants defined below, not the scalar coefficients $a_1,a_2$ of Section~\ref{sec:chain}.

\emph{Blocks (King's averaging).} For $1\le k\le K$ and $y\in T\cap t_k\Z^3$ let $B^k(y):=\{x\in T_\eps:\ y_i\le x_i<y_i+t_k,\ i=1,2,3\}$. Since $\ell_T/t_k=(\ell_T/(M_bL^k\eps))\,M_b\in\N$ by Lemma~\ref{lem:sizes}(ii), these blocks tile $T_\eps$. Let $Q_k$ be the componentwise average $(Q_kA_\mu)(y):=L^{-3k}\sum_{x\in B^k(y)}A_\mu(x)$ \cite[(2.10), p.~653]{K1}, with adjoint $Q_k^*$ taken for the inner product $t_k^3\sum_y$ on the block lattice. Then $Q_k^*$ is the constant extension, $Q_kQ_k^*=1$, and $P_k:=Q_k^*Q_k$ is the orthogonal projection onto the fields that are constant on every scale-$k$ block.
Fix King's constant $a>0$ and put
\[
a_k:=a(1-L^{-2})(1-L^{-2k})^{-1},\qquad \varrho^{\mathrm K}_k:=a_kt_k^{-2}\qquad(1\le k\le K),
\]
\cite[(2.13), p.~653]{K1}. Then $G_k=(M_*+\varrho^{\mathrm K}_kP_k)^{-1}$ for $1\le k\le K$ is the operator of Subsection~\ref{ssec:T6-statement}. We set $G_0:=0$, which is King's convention $C^{(0)}=G_1$, $G_k=\sum_{j=0}^{k-1}G_{(j)}$ \cite[(2.16)--(2.17), p.~653]{K1}. Thus $H_0=G_1$, and $H_j=G_{j+1}-G_j$ ($0\le j<K$) is King's $G_{(j)}$.

\begin{lemma}[King's slices are nonnegative]\label{app:G1:slices}
For every $K\ge1$, $a>0$, $\mu_*>0$ and every torus $T$ as above:
\begin{enumerate}
\item[(a)] $0<G_1\le G_2\le\dots\le G_K\le C_{\mu_*}$; equivalently $H_j\ge0$ for $0\le j<K$ and $C_{\mu_*}-G_K\ge0$;
\item[(b)] $\sum_{j<K}H_j=G_K$, hence $C_{\mu_*}=\sum_{j<K}H_j+(C_{\mu_*}-G_K)$, a sum of nonnegative operators.
\end{enumerate}
\end{lemma}
\noindent In particular $H_j\ge0$, which is used in Subsection~\ref{ssec:T6-statement}.
\begin{proof}
\emph{Step 1 (inversion reverses order).} If $0<A\le B$ are self-adjoint, then $A^{-1/2}BA^{-1/2}\ge1$, so its inverse $A^{1/2}B^{-1}A^{1/2}$ is $\le1$, i.e.\ $B^{-1}\le A^{-1}$.

\emph{Step 2 (nested blocks).} For $1\le k<K$ each block $B^{k+1}(y)$ is the disjoint union of the $L^3$ blocks $B^k(y')$, $y'\in y+t_k\{0,\dots,L-1\}^3$. So a field constant on the scale-$(k+1)$ blocks is constant on the scale-$k$ blocks, $\operatorname{ran}P_{k+1}\subset\operatorname{ran}P_k$, and $P_k-P_{k+1}$ is the orthogonal projection onto $\operatorname{ran}P_k\ominus\operatorname{ran}P_{k+1}$. In particular $P_k-P_{k+1}\ge0$.

\emph{Step 3 ($\varrho^{\mathrm K}_k$ decreases).} Since $1-L^{-2k}$ increases with $k$, $a_k$ decreases, and $\varrho^{\mathrm K}_k/\varrho^{\mathrm K}_{k+1}=L^2a_k/a_{k+1}\ge L^2>1$.

\emph{Step 4 ($H_j\ge0$).} For $1\le j<K$,
\[
G_j^{-1}-G_{j+1}^{-1}=\varrho^{\mathrm K}_jP_j-\varrho^{\mathrm K}_{j+1}P_{j+1}=\varrho^{\mathrm K}_j(P_j-P_{j+1})+(\varrho^{\mathrm K}_j-\varrho^{\mathrm K}_{j+1})P_{j+1}\ge0
\]
by Steps 2 and 3. Both inverses are $\ge M_*>0$ ($\varrho^{\mathrm K}_k>0$, $P_k\ge0$), so Step~1 gives $G_j\le G_{j+1}$, i.e.\ $H_j\ge0$. For $j=0$, $H_0=G_1>0$.

\emph{Step 5 ($G_K\le C_{\mu_*}$).} $G_K^{-1}=M_*+\varrho^{\mathrm K}_KP_K\ge M_*>0$, so $G_K\le M_*^{-1}=C_{\mu_*}$ by Step~1. This proves (a).

\emph{Step 6.} The sum telescopes: $\sum_{j<K}(G_{j+1}-G_j)=G_K-G_0=G_K$. Together with (a) this gives (b).
\end{proof}

\begin{lemma}[the domination in \eqref{eq:T6-GammaB}]\label{app:G1:GammaBdom}
Let $\kappa>0$ and $\Gamma^B=\sum_{j<K}\delta_jH_j$, $\delta_j=\kappa\mu_*^2t_{j+1}^2$, as in \eqref{eq:T6-GammaB}. Then $0\le\delta_j\le\bar\delta=\kappa\mu_*^2\eps_0^2$ for $j<K$, and
\[
0\le\Gamma^B\le\kappa\mu_*^2\eps_0^2\,G_K=\bar\delta\,G_K\le\bar\delta\,C_{\mu_*}.
\]
\end{lemma}
\noindent This is the domination $\Gamma^B\le\kappa\mu_*^2\eps_0^2G_K\le\bar\delta C_{\mu_*}$ stated after \eqref{eq:T6-GammaB}; it is uniform in $\eps$, $K$, $\ell_T$.
\begin{proof}
For $j<K$, $0<t_{j+1}\le t_K=\eps_0$, hence $0\le\delta_j\le\bar\delta$. Since $H_j\ge0$ (Lemma~\ref{app:G1:slices}(a)), $\Gamma^B\ge0$ and
\[
\bar\delta\,G_K-\Gamma^B=\kappa\mu_*^2\sum_{j<K}\big(t_K^2-t_{j+1}^2\big)H_j\ \ge0,
\]
where $G_K=\sum_{j<K}H_j$ (Lemma~\ref{app:G1:slices}(b)) was used. Finally $G_K\le C_{\mu_*}$ (Lemma~\ref{app:G1:slices}(a)).
\end{proof}

\begin{remark}[the constant $\kappa$]\label{app:G1:kappa}
Lemma~\ref{app:G1:GammaBdom} holds for every $\kappa>0$, since $\kappa$ is only the coefficient in the definition of $\Gamma^B$. The numerical bound $\kappa_1\kappa_3\le7.5$ quoted after \eqref{eq:T6-GammaB} belongs to a different inequality, the domination $D^<\le\Gamma^B$ of the slice-modified photon (Remark~\ref{rem:T6-routes}, Theorem~\ref{app:G1:DOM-nat}). That inequality fixes $\kappa$ and is the one used in the proof of Proposition~\ref{prop:ch-Tflat}; the admissible values of $\kappa$ are in Remark~\ref{app:G1:DOM-kappa}.
\end{remark}

\begin{lemma}[positivity of the slice photon]\label{app:G1:CBpos}
For every $u>-1/\bar\delta$,
\[
C_B(u)=C_{\mu_*}+u\Gamma^B=\sum_{j<K}(1+u\delta_j)H_j+(C_{\mu_*}-G_K)\ \ge\ \min\{1,1+u\bar\delta\}\,C_{\mu_*}\ >0,
\]
and every step weight satisfies $1+u\delta_j\ge\min\{1,1+u\bar\delta\}>0$.
\end{lemma}
\noindent This is the positivity of $C_B(u)$ for $u>-1/\bar\delta$ stated in Subsection~\ref{ssec:T6-statement}. It is used for the existence of $N(0,C_B(u))$ in the definition of $\Phi_M(u)$, for the convexity of $\Phi_M$ on $(-1/\bar\delta,\infty)$, and, at $u=-\eta$ with $\eta\bar\delta<1$, in Theorem~\ref{thm:T6}.
\begin{proof}
The identity follows from Lemma~\ref{app:G1:slices}(b) and the definition of $\Gamma^B$. If $u\ge0$, then $u\Gamma^B\ge0$ (Lemma~\ref{app:G1:GammaBdom}), so $C_B(u)\ge C_{\mu_*}$. If $-1/\bar\delta<u<0$, then $u\Gamma^B\ge u\bar\delta\,C_{\mu_*}$ by Lemma~\ref{app:G1:GammaBdom}, so $C_B(u)\ge(1+u\bar\delta)C_{\mu_*}$ with $1+u\bar\delta>0$. In both cases $C_{\mu_*}=M_*^{-1}>0$ because $-\Delta_\eps\ge0$ and $\mu_*>0$. For the weights: $u\delta_j\ge0$ if $u\ge0$, and $u\delta_j\ge u\bar\delta$ if $u<0$, because $0\le\delta_j\le\bar\delta$.
\end{proof}

\emph{Setting for the degree lemma.} Write $Z_\varphi(A;a):=\int\exp\big(-S^{\mathrm m}_a(A,\varphi)\big)\prod_{x\in T_\eps}d^2\varphi(x)$ for the unnormalised scalar integral with the matter action \eqref{eq:matteraction}. This is the factor in $Z(\Gamma;a)=\mathbb E_{N(0,\Gamma)}[Z_\varphi(A;a)]$, whose photon measure is normalised and which contains no c-numbers. The coefficient $a=a_{K,T}(M)=M^2+\delta m^2_{K,T}(\eps;M)$ of Definition~\ref{def:counterterm} is computed with the photon propagator $C_{\mu_*}$. It does not depend on $u$, and it is a polynomial in $(e^2,\lambda)$ without constant term added to $M^2$ (monomials $e^2,\lambda,e^4,e^2\lambda,\lambda^2$). We read $\mathcal E(u;M)=T_{\nbar}[\Phi_M(u)]$ as the Taylor polynomial of total order $\nbar$ in $(e,\lambda)$ at $(0,0)$, with $a_{K,T}(M)$ inside, as in \cite[(1.13)]{Ba1}. The proof below uses only that the coefficient is a $u$-independent smooth function of $(e^2,\lambda)$ equal to $M^2$ at the origin.

\begin{lemma}[the $u$-degree of $e(\eps;u;M)$]\label{app:G1:edeg}
Fix $\eps$, $T$, $M>0$, $\lambda\ge0$ and $\kappa,\mu_*>0$. For $u>-1/\bar\delta$ the Taylor coefficients defining $\mathcal E(u;M)$ exist. The coefficient of $e^\alpha\lambda^\beta$ ($\alpha+\beta\le\nbar$) vanishes for odd $\alpha$, and for $\alpha=2k$ it is a polynomial in $u$ of degree at most $k$. Hence $\mathcal E(u;M)$ and $e(\eps;u;M)=-\abs T^{-1}\mathcal E(u;M)$ are polynomials in $u$ of degree at most $\lfloor\nbar/2\rfloor$. Their order-zero term $\log Z_\varphi(0;M^2)\big|_{e=\lambda=0}=E_{0,s}(M)$ does not depend on $u$.
\end{lemma}
\noindent This is the bound on the degree in $u$ stated in Subsection~\ref{ssec:T6-statement}. It is used in the proof of Proposition~\ref{prop:ch-T5}: the second divided difference of the polynomial $u\mapsto e(\eps;u;m_1)$, and the meaning of $\partial_u^2e$ in the bound \eqref{p:ch-E2B}.
\begin{proof}
\emph{Step 1 (gauge field enters only through $eA$).} The action \eqref{eq:matteraction} depends on $A$ only through $\mathrm{e}^{ie\eps A_{x,\mu}}$. Hence $Z_\varphi(A;a)=\mathcal Z(eA;\lambda,a)$, where $\mathcal Z(B;\lambda,a)$ is the same integral with $\mathrm{e}^{i\eps B_{x,\mu}}$ in place of $\mathrm{e}^{ie\eps A_{x,\mu}}$. If $A\sim N(0,\Gamma)$ then $eA$ has the law of $\sqrt s\,Y$, $s:=e^2$, $Y\sim N(0,\Gamma)$. Thus, with $\Gamma=C_B(u)>0$ (Lemma~\ref{app:G1:CBpos}),
\[
Z(C_B(u);a)=h_u(e^2,\lambda,a),\qquad h_u(s,\lambda,a):=\mathbb E\big[\mathcal Z(\sqrt s\,Y_u;\lambda,a)\big],\quad Y_u\sim N(0,C_B(u)).
\]

\emph{Step 2 (regularity of $\mathcal Z$).} Let $a\ge a_{\min}>0$ and $\lambda\ge0$. The gradient and quartic terms of \eqref{eq:matteraction} are nonnegative, so the integrand is bounded by $\exp(-\tfrac12a_{\min}\sum_x\eps^3\abs{\varphi(x)}^2)$. Every derivative in $B$, $\lambda$ or $a$ multiplies the integrand by a polynomial in $\abs\varphi$ whose coefficients are bounded uniformly in $B$, because $\abs{\mathrm{e}^{i\eps B}}=1$. By dominated convergence $\mathcal Z$ is $C^\infty$ in $(B,\lambda,a)\in\R^{n}\times[0,\infty)\times[a_{\min},\infty)$ (one-sided in $\lambda$ at $0$), and all its $B$-derivatives are bounded uniformly in $B$.

\emph{Step 3 (Gaussian heat equation).} Let $c_u(b,b'):=\mathbb E[Y_u(b)Y_u(b')]$ be the second-moment matrix of $Y_u$. It is affine in $u$ because $C_B(u)$ is. Put $L_u:=\sum_{b,b'}c_u(b,b')\,\partial_{B_b}\partial_{B_{b'}}$. For $F$ smooth with bounded derivatives, $w(s):=\mathbb E[F(\sqrt sY_u)]$ satisfies $w'(s)=\mathbb E[(\tfrac12L_uF)(\sqrt sY_u)]$ for $s>0$: differentiate the density of $N(0,sc_u)$, which solves $\partial_sp=\tfrac12L_up$, and integrate by parts. The right side is continuous on $[0,\infty)$ by dominated convergence, so $w\in C^\infty([0,\infty))$ with $w^{(i)}(0)=2^{-i}(L_u^iF)(0)$. Applied to $F=\partial_\lambda^{b}\partial_a^{c}\mathcal Z(\cdot;\lambda,a)$ (Step 2), $h_u$ is $C^\infty$ on $[0,\infty)\times[0,\infty)\times[a_{\min},\infty)$, and
\[
\partial_s^{i}\partial_\lambda^{b}\partial_a^{c}h_u(0,\lambda,a)=2^{-i}\big(L_u^{i}\,\partial_\lambda^{b}\partial_a^{c}\mathcal Z\big)(0;\lambda,a).
\]
The right side is a homogeneous polynomial of degree $i$ in the entries of $c_u$, hence a polynomial in $u$ of degree $\le i$. For $i=0$ it is $\partial_\lambda^b\partial_a^c\mathcal Z(0;\lambda,a)$, independent of $u$, and $\mathcal Z(0;\lambda,a)>0$.

\emph{Step 4 (composition and Taylor polynomial).} Write $a(s,\lambda):=M^2+\delta m^2_{K,T}(\eps;M)|_{e^2=s}$. This is a polynomial in $(s,\lambda)$, independent of $u$, with $a(0,0)=M^2$ (Definition~\ref{def:counterterm}). Choose a neighbourhood $U$ of $(0,0)$ in $[0,\infty)^2$ on which $a(s,\lambda)\ge M^2/2=:a_{\min}$. On $U$ the function $\psi_u(s,\lambda):=\log h_u(s,\lambda,a(s,\lambda))$ is $C^\infty$, and $\Phi_M(u)=\psi_u(e^2,\lambda)$. By the chain rule and the formula for derivatives of $\log$, $\partial_s^k\partial_\lambda^\beta\psi_u(0,0)$ is a finite sum, with $u$-independent coefficients, of terms
\begin{gather*}
h_u(0,0,M^2)^{-m}\prod_r\partial_s^{i_r}\partial_\lambda^{b_r}\partial_a^{c_r}h_u(0,0,M^2)\times\big(\text{derivatives of }a(s,\lambda)\text{ at }(0,0)\big),\\
\textstyle\sum_ri_r\le k.
\end{gather*}
By Step~3 each such term is a polynomial in $u$ of degree $\le k$, and $h_u(0,0,M^2)=\mathcal Z(0;0,M^2)>0$ does not depend on $u$. Taylor's formula for $\psi_u$ at $(0,0)$ and the substitution $s=e^2$ give
\[
\Phi_M(u)=\sum_{2k+\beta\le\nbar}\frac{\partial_s^k\partial_\lambda^\beta\psi_u(0,0)}{k!\,\beta!}\,e^{2k}\lambda^\beta+o\big(\abs{(e,\lambda)}^{\nbar}\big),
\]
and by uniqueness of Taylor polynomials this sum is $T_{\nbar}[\Phi_M(u)]=\mathcal E(u;M)$. So the coefficients of odd powers of $e$ vanish, and the coefficient of $e^{2k}\lambda^\beta$ has $u$-degree $\le k\le\lfloor\nbar/2\rfloor$, since $2k\le\nbar$. The order-zero term is $\psi_u(0,0)=\log\mathcal Z(0;0,M^2)$, the free scalar constant $E_{0,s}(M)$ of Remark~\ref{rem:T6-textfixes}(i), independent of $u$. Multiplying by $-\abs T^{-1}$ gives the claim for $e(\eps;u;M)$.
\end{proof}

\begin{remark}\label{app:G1:normalised}
Lemma~\ref{app:G1:edeg} uses that the photon measure in $Z(\Gamma;a)$ is normalised. With the Lebesgue-normalised photon of Propositions~\ref{prop:ch-CN} and \ref{prop:ch-Pplus}, the free energy would contain the additional term $\tfrac12\log\det C_B(u)$, which is not a polynomial in $u$. This is consistent with Remark~\ref{rem:T6-textfixes}(i).
\end{remark}

\subsection{Final c-numbers, the $A\cdot A$ vertices and the tail beyond the order $\nbar$}\label{app:G1:part-D3}

We prove the three statements behind Lemma~\ref{lem:B:massrange}(a)--(c), directly at a Higgs mass $M\in(0,m_2]$: bounds on the final c-numbers uniform in the torus (Propositions~\ref{app:G1:I3b} and \ref{app:G1:Efin}), bounds on the coefficients of the $A\cdot A$ vertices (Lemmas~\ref{app:G1:AA} and \ref{app:G1:AA4}), and the tail bound beyond the order $\nbar$ (Proposition~\ref{app:G1:T579}). The line bounds, the remainder line and the renormalisation machinery on which they rest are proved in Subsection~\ref{app:G1:part-TVSR}.

\emph{Standing notation.} $d=3$, $L=2$, $\eps=2^{-K}\eps_0$, $t_k:=L^k\eps$ (so $t_K=\eps_0$), $e(t)=e\,t^{1/2}$, $\lambda(t)=\lambda t$, $p(t)=b_0(1+\ln t^{-1})^{p_0}$. The scalar field $f=s$ has mass $m_s:=M$; the photon $f=v$ (Feynman gauge, per component) has mass $m_v:=\mu_*$. On the torus $T_\eps$ of side $\ell_T$, $C_f$ is the free covariance of field $f$, and $G_{k,f}(0)$ is Ba{\l}aban's block-spin propagator at zero background. Its slices are $H_{j,f}:=G_{j+1,f}(0)-G_{j,f}(0)$.

\paragraph{The final c-numbers.}\label{app:G1:ssec-N9}
The c-numbers are $\delta m^2_{\rm fin}:=\delta m^2-\delta m^2_K$ and $E_{\rm fin}:=E_1-E_K$ of \cite[(1.26), (1.30)]{Ba3}. Both are fixed numbers built from free lines, so any exact decomposition of the lines above $\eps_0$ may be used to bound them.

\emph{The common top index.} Put $\bar\ell(M):=\min(1/M,1/\mu_*,L)$, and let $\bar K$ be the largest $k$ with $t_k\le\bar\ell(M)$. Assume
\begin{equation}\label{app:G1:eq-hyp}
\eps_0\le\bar\ell(M),\qquad \ell_T\ge L^3.
\end{equation}
Then $\bar K\ge K$. For both fields telescoping gives the exact identity
\begin{equation}\label{app:G1:eq-N9split}
C_f=G_{K,f}(0)+\sum_{j=K}^{\bar K-1}H_{j,f}+R_{\bar K,f},\qquad R_{\bar K,f}:=C_f-G_{\bar K,f}(0),\qquad f\in\{s,v\}.
\end{equation}
Put $\rho_f:=m_ft_{\bar K}$, $\rho_{\min}:=\min_f\rho_f$, $Z_f:=t_{\bar K}/(m_f^2\ell_T^3)$, and $\kappa_{R,\min}:=\min(\rho_{\min},\kappa_G)/8$. Here $\kappa_G:=\min(c,1/16)$ is the rate defined before Lemma~\ref{app:G1:VSR-E}; it depends only on $(a,L)$.

The following properties are elementary.
(P1) $t_{\bar K}\in(\bar\ell/L,\bar\ell]$ is a power of $L$ that does not depend on $\eps$ or $\ell_T$.
(P2) $\rho_f\in(0,1]$, since $m_ft_{\bar K}\le m_f\bar\ell\le1$.
(P3) $t_{\bar K}^{-1}<L/\bar\ell=\max(LM,L\mu_*,1)$.
(P4) $\bar K-K\le1+\log_L^+(\bar\ell/\eps_0)$.
(P5) $\rho_{\min}>\min_fm_f\bar\ell/L\ge\min(M,\mu_*)\min\{1,(L\max(M,\mu_*))^{-1}\}$.
(P6) $Z_f\le(m_f^2\max(M,\mu_*)\ell_T^3)^{-1}$, since $t_{\bar K}\le\bar\ell\le1/\max(M,\mu_*)$; and $Z_f$ is strictly decreasing in $\ell_T$.
(P7) $\bar K\le\min_fK_f$, where $K_f$, the top index of the field $f$ alone, is the largest $k$ with $t_k\le\min(1/m_f,\ell_T/L^2)$. Indeed $t_{\bar K}\le\min(1/m_f,L)$, and $L\le\ell_T/L^2$ is $\ell_T\ge L^3$.

\emph{Inputs (proved in Subsection~\ref{app:G1:part-TVSR}).}
(VS) Let $t_{j+1}\le\min(1/m_f,L)$. Then the slices $H_{j,f}$ obey Ba{\l}aban's line bounds \cite[(2.10)--(2.12), p.~426]{Ba3} at zero background. The constants depend only on $(a,L)$ (and on $\alpha$ for the H\"older quotients), and the bounds are uniform in $\eps$, $j$, $\ell_T$, $\mu_*$ and $M$. This holds below and above $\eps_0$ (Theorem~\ref{app:G1:VSR-VS} with Corollary~\ref{app:G1:VSR-VSphys}).
(Rem) In scale-$\bar K$ units the remainder $\tilde R$ of $R_{\bar K,f}$ obeys $\abs{\tilde R(x,y)}\le B_R[\mathrm{e}^{-\kappa_Rd(x,y)}+Z_f]$ and $\eta^3\sum_y\abs{\tilde R(x,y)}\le B_R'\rho_f^{-2}$, with $\kappa_R=\min(\rho_f,\kappa_G)/8$. The same bounds hold for first derivatives in either variable, for mixed derivatives and for H\"older quotients, with $B_R''$ and $B_R(\alpha)$ (Proposition~\ref{app:G1:VSR-R}). The prefactors depend only on $(a,L,\alpha)$.
(S6) The renormalisation machinery of \cite[\S3]{Ba3} on the torus: Lemma~\ref{app:G1:VSR-A} (with its pinned form (ii)), Lemmas~\ref{app:G1:VSR-Btop} and \ref{app:G1:VSR-C}, Theorem~\ref{app:G1:VSR-T335} and Proposition~\ref{app:G1:VSR-5b}. It bounds every index assignment whose renormalised subgraphs all have lowest external index $j_0\le\min_fK_f$, by the power counting of \cite[\S2]{Ba3} at the top index.

\begin{lemma}[no intermediate subgraph scale]\label{app:G1:common}
Assume \eqref{app:G1:eq-hyp}. Expand every line of a counterterm graph of $\delta m^2$ or of $E_1$ by \eqref{app:G1:eq-N9split}, and fix an index assignment; the line $R_{\bar K,f}$ has index $\bar K$. Let $S$ be a connected proper subgraph all of whose internal lines have index $<j_0$, where $j_0$ is the lowest index of its external lines. Then $j_0\le\bar K$, and no internal line of $S$ is an $R$-line. For each field, the sum over internal indices $i<j_0$ (together with the lines below $\eps_0$) is $G_{j_0,f}(0)$, of the same form for $f=s$ and $f=v$. In particular $j_0\le\min_fK_f$.
\end{lemma}
\begin{proof}
Both fields are sliced up to the same index $\bar K$, and every line has an index in $\{<K\}\cup\{K,\dots,\bar K-1\}\cup\{\bar K\}$, the same set for both fields. Hence $j_0\le\bar K$. An internal line has index $<j_0\le\bar K$, so it is not an $R$-line. The internal lines of field $f$ are therefore slices $H_{i,f}$ with $i<j_0$, or lines below $\eps_0$. By telescoping, their sum is $G_{K,f}(0)+\sum_{K\le i<j_0}H_{i,f}=G_{j_0,f}(0)$, of the same algebraic form for both fields. Finally $j_0\le\bar K\le\min_fK_f$ by (P7).
\end{proof}
In the case $j_0\le\min_fK_f$, the replacement of $G_{j_0,f}(0)$ by $C^\xi$ carries the factor $\omega_f:=1-m_f^2t_{j_0}^2\in[0,1]$. No renormalised subgraph contains an $R$-line or a full $C_f$. So no zero mode, and no Ward identity beyond those of \cite[\S3]{Ba3}, enters a renormalisation.

\begin{lemma}[the top cluster]\label{app:G1:TB}
Assume \eqref{app:G1:eq-hyp}, and consider a term with top index $j^*=\bar K$, in scale-$\bar K$ units. The \emph{lower clusters} are its maximal connected subgraphs made of lines of index $<\bar K$, each with its renormalisations. Choose a spanning tree $\mathcal T$ of $R$-lines on the lower clusters, rooted at the cluster that contains the pinned point $x$. Then
\[
\abs{\text{term}}\le\prod_{\text{lower clusters}}(\text{bound of Lemmas~\ref{app:G1:VSR-A}(ii) and \ref{app:G1:VSR-C}, uniform in the pin})\cdot\prod_{l\in\mathcal T}N_1(l)\cdot\prod_{l\notin\mathcal T,\ R}N_\infty(l),
\]
with $N_1(l):=\sup_y\eta^3\sum_{y'}\abs{\partial^a\partial'^b\tilde R_f(y,y')}\le B_R'\rho_f^{-2}+B_R''J_0(\kappa_R)(1+\rho_f^{-2})$ and $N_\infty(l):=\sup\abs{\partial^a\partial'^b\tilde R_f}\le B_R''(1+Z_f)$, $a,b\in\{0,1\}$. Here $J_0(\kappa):=\sup_x\eta^3\sum_y\mathrm{e}^{-\kappa d(x,y)}$ is the lattice sum of Subsection~\ref{app:G1:part-TVSR}.
\end{lemma}
\begin{proof}
By definition the lower clusters are joined to one another only by $R$-lines. Sum the cluster positions from the leaves of $\mathcal T$ towards the root. Each tree line is summed over its far endpoint, which costs its $\ell^1$ norm $N_1$. The far cluster's internal positions, relative to that endpoint, are summed inside its own pinned bound of Lemma~\ref{app:G1:VSR-A}(ii). Each non-tree $R$-line is bounded by its supremum $N_\infty$. The bounds on $N_1$ and $N_\infty$ are (Rem): on the torus the $Z_f$-part of the pointwise bound sums to $Z_f\ell^3=\rho_f^{-2}$. Where Lemma~\ref{app:G1:VSR-C} needs the regularity of an external $R$-line of a lower cluster, (Rem) supplies the derivative and H\"older bounds with the factor $1+Z_f$.
\end{proof}

\begin{proposition}[a bound on $\delta m^2_{\rm fin}$ uniform in the torus]\label{app:G1:I3b}
Let $M\in(0,m_2]$ and assume \eqref{app:G1:eq-hyp}. Uniformly in $\eps=2^{-K}\eps_0$, $x\in T_\eps$ and $\ell_T\in\{\ell_j\}$,
\begin{multline*}
\abs{\delta m^2_{\rm fin}(x)}\le C_b\Big\{(e^2+\lambda)\frac{L^2}{(L-1)\eps_0}+(e^4+e^2\lambda+\lambda^2)\big(2+\log_L^+(\bar\ell/\eps_0)\big)\Big\}\\
+C_b^R\Big\{(e^2+\lambda)\max(LM,L\mu_*,1)(1+Z_s+Z_v)^2+(e^4+e^2\lambda+\lambda^2)(1+Z_s+Z_v)^5\Big\}.
\end{multline*}
Here $C_b=C(\alpha_0,\nbar,d,L,a,M_b)$, and $C_b^R=C(\alpha_0,\nbar,d,L,a,M_b)(1+\rho_{\min}^{-2}+\kappa_{R,\min}^{-3})^5$. The torus side enters only through $Z_s$ and $Z_v$, both non-increasing in $\ell_T$.
\end{proposition}
\begin{proof}
Split the index assignments of \eqref{app:G1:eq-N9split} by the top index $j^*\in\{K,\dots,\bar K\}$.

(i) $K\le j^*<\bar K$. All lines are slices with $t_{j+1}\le t_{\bar K}\le\min(1/m_f,L)$, so (VS) applies to both fields. By Lemma~\ref{app:G1:common} every renormalised subgraph is in the case $j_0\le\min_fK_f$, so (S6) applies at top index $j^*$ with no exceptional subgraph. The degrees are $D=-1$ at one loop and $D=0$ at two loops. Rescaling $\delta m^2=t_{j^*}^{-2}\cdot(\text{scale-}j^*\text{ coefficient})$ with $e(t)^2=e^2t$ and $\lambda(t)=\lambda t$, the one-loop terms are $\le C(e^2+\lambda)t_{j^*}^{-1}$. Their sum is $\sum_{j^*\ge K}t_{j^*}^{-1}=t_K^{-1}\sum_{n\ge0}L^{-n}=L/((L-1)\eps_0)\le L^2/((L-1)\eps_0)$. The two-loop terms give one $O(1)$ per $j^*$, and there are $\bar K-K+1\le2+\log_L^+(\bar\ell/\eps_0)$ values of $j^*$, by (P4).

(ii) $j^*=\bar K$. Apply Lemma~\ref{app:G1:TB}, with at most $2$ lines at one loop and at most $5$ lines ($I=V+1\le5$) at two loops. Each $R$-factor is at most $C(1+\rho_f^{-2}+\kappa_R^{-3})(1+Z_f)$. Rescaling from scale $\bar K$ gives $t_{\bar K}^{-2}e^2t_{\bar K}=e^2t_{\bar K}^{-1}<e^2\max(LM,L\mu_*,1)$ at one loop, by (P3), and no factor at two loops. The polynomial in $\rho_f^{-1}$ and $\kappa_R^{-1}$ is at most its value at $\rho_{\min}$.

The torus side enters (i) nowhere: the constants of (VS) are uniform in $\ell_T$, and no condition on the side of the torus is used. It enters (ii) only through $Z_f$, which is decreasing by (P6). Here (VS) is Theorem~\ref{app:G1:VSR-VS}, which rests on the free kernel bounds of Lemma~\ref{app:G1:VSR-F}, and (S6) consists of Lemmas~\ref{app:G1:VSR-A}, \ref{app:G1:VSR-Btop}, \ref{app:G1:VSR-C}, Theorem~\ref{app:G1:VSR-T335} and Proposition~\ref{app:G1:VSR-5b} (Remark~\ref{app:G1:VSR-S6use}).
\end{proof}

\begin{lemma}[the splitting of $\delta m^2$ and its inner counterterms]\label{app:G1:dm2grade}
Every graph of $\delta m^2$ \cite[(1.23), p.~417]{Ba3} either has no mass counterterm vertex, or is a one-loop two-leg graph $\Gamma$ whose internal scalar line carries one vertex \cite[(1.7)]{Ba3} with coefficient $\delta m^2_1$, the part of order $\alpha+2\beta=2$. Let $\Gamma'$ be $\Gamma$ with coefficient $1$ at that vertex. For a set $J$ of line indices let $\Gamma'_J$, $s_J$ be the parts of $\Gamma'$, $\delta m^2_1$ whose line indices lie in $J$, with the index $K$ standing for the regular remainder. Split $\delta m^2=\delta m^2_{\rm fin}+\delta m^2_{K,k}+\delta m^2_k$ as in \cite[(1.26)--(1.28)]{Ba3} in two ways: (O) grading all lines, including those of $\delta m^2_1$; (B) grading the lines of $\Gamma'$ only, with $\delta m^2_1$ kept as a number. Then:
\begin{enumerate}
\item[(a)] both splittings are exact, and for every graph $\Gamma$ as above
\[
\delta m^{2,{\rm O}}_{K,k}-\delta m^{2,{\rm B}}_{K,k}=\Gamma'_{<k}\,s_{[k,K)}-\Gamma'_{[k,K)}\,s_{\ge K},\qquad
\delta m^{2,{\rm O}}_{\rm fin}-\delta m^{2,{\rm B}}_{\rm fin}=\Gamma'_{<K}\,s_{\ge K};
\]
the pieces $\Delta_k:=(\delta m^2_{K,k}-\delta m^2_{K,k+1})t_k^2$ of the two splittings differ by $(\Gamma'_{<k}s_k-\Gamma'_ks_{>k})t_k^2$;
\item[(b)] $D(\Gamma')=D(\Gamma)+2=1$, whereas a closed scalar or vector loop carrying one vertex (1.7) has degree $-1$;
\item[(c)] $\abs{\Gamma'_{\{s\}}}\le C(e^2+\lambda)t_s$, $\abs{s_j}\le C(e^2+\lambda)t_j^{-1}$, and all the differences in (a), multiplied by $t_k^2$, are at most $C(e^2+\lambda)^2t_k^2\big(1+\sum_f(1+Z_f)^2\big)$.
\end{enumerate}
In particular the bounds on $\delta m^2_{K,k}$, on the pieces $\Delta_k$ and on $\delta m^2_{\rm fin}$ hold for both splittings, with the same form.
\end{lemma}
\begin{proof}
(a) Expand multilinearly. In (O) a term of $\Gamma'\cdot\delta m^2_1$ belongs to $\delta m^2_{K,k}$ if all its indices are $<K$ and one of them is $\ge k$; this gives $\Gamma'_{[k,K)}s_{<K}+\Gamma'_{<k}s_{[k,K)}$. In (B) it gives $\Gamma'_{[k,K)}(s_{<K}+s_{\ge K})$. Subtracting gives the first identity. The other two follow in the same way, using $s_{>k}=s_{[k+1,K)}+s_{\ge K}$.
(b) By \cite[(2.2) and p.~423(ii)]{Ba3}, the new vertex has $D_G(v)=d+2(-d+2)/2=2$, and no other vertex changes. One-loop two-leg graphs have $D=-1$, and a free loop has $D=-3$.
(c) By \cite[(2.10)--(2.12)]{Ba3} and Theorem~\ref{app:G1:VSR-VS}, every index tuple of $\Gamma'$ with largest index $s$ contributes at most $Ct_s$ times $(t_a/t_s)^{\theta}$, $\theta\ge1$, for each smaller index $a$; for instance $\min(t_a,t_s)^3/(t_at_s)=t_s(t_a/t_s)^2$. Summing gives $\abs{\Gamma'_{\{s\}}}\le C(e^2+\lambda)t_s$, $\abs{\Gamma'_{<k}}\le C(e^2+\lambda)t_k$ and $\abs{\Gamma'_{<K}}\le C(e^2+\lambda)\eps_0$. One-loop graphs have degree $-1$, so $\abs{s_j}\le C(e^2+\lambda)t_j^{-1}$. By Proposition~\ref{app:G1:I3b}, and since $\max(LM,L\mu_*,1)=L/\bar\ell\le L/\eps_0$ by \eqref{app:G1:eq-hyp}, $\abs{s_{\ge K}}\le C(e^2+\lambda)\eps_0^{-1}(1+Z_s+Z_v)^2$, with $C$ depending also on $\rho_{\min}$ and $\kappa_{R,\min}$. Since $t_k\le\eps_0$ and $(1+Z_s+Z_v)^2\le2\sum_f(1+Z_f)^2$, (c) follows.
\end{proof}

\begin{remark}\label{app:G1:dm2grade-rem}
Lemma~\ref{app:G1:dm2grade}(b) is the reason why the two gradings agree for the mass counterterms but not for the vacuum pieces (Remark~\ref{app:G1:vacgrade-rem}): a finite mass vertex makes every two-leg graph convergent, whereas on a vacuum loop it leaves degree $-1$.
\end{remark}

\begin{definition}[graded vacuum pieces]\label{app:G1:vacgrade}
A \emph{bare vacuum graph} is a connected graph without external legs built from the vertices \cite[(1.6), (1.7), (1.8), (1.10)]{Ba3}, with $n'=0$, $\tilde B=0$, $g_k=1$ \cite[p.~416]{Ba3}. Every vertex \cite[(1.7)]{Ba3} carries the number $\delta m^2$ of \cite[(1.19)--(1.23)]{Ba3}, a polynomial in $e,\lambda$ of orders $2\le\alpha+2\beta\le4$; the order of the vertex is the order of the term of $\delta m^2$ it carries. The \emph{lines} of the graph are the lines joining its vertices; the graphs defining $\delta m^2$ are not lines. By \cite[(1.24)]{Ba3}, $E_1$ is the sum of the values of all bare vacuum graphs of orders $1\le\alpha+\beta\le\nbar$. Expand every line of field $f$ as $C_f=\sum_{j=0}^{\bar K}C_{(j),f}$, where for $j<K$ the $C_{(j),f}$ are the terms of the decomposition \cite[(1.27)]{Ba3} of $G_{K,f}(0)$ (with $C_{(0),f}$ the term $C^{(0)}$), $C_{(j),f}:=H_{j,f}$ for $K\le j<\bar K$, and $C_{(\bar K),f}:=R_{\bar K,f}$, as in \eqref{app:G1:eq-N9split}. The integer $j$ is the \emph{index} of the line. Expanding multilinearly, with every $\delta m^2$ kept as a number, define
\begin{itemize}
\item $E_{\rm fin}$: the sum of the terms with at least one line of index $\ge K$;
\item $E_K$: the sum of the terms with all line indices $<K$;
\item $E_{K,k}$: the sum of the terms with all line indices $<K$ and at least one $\ge k$;
\item $E_k$: the sum of the terms with all line indices $<k$;
\item $\cV_k:=E_{K,k}-E_{K,k+1}$, for $0\le k<K$.
\end{itemize}
For estimates the coefficients may be split by \cite[(1.26), (1.28)]{Ba3} as $\delta m^2=\delta m^2_{\rm fin}+\delta m^2_{K,k+1}+\delta m^2_{k+1}$; the first two parts are counterterms of degree $0$, and the last has the degree of its graph \cite[p.~423]{Ba3}. We keep the notation of \cite[(1.24), (1.30)--(1.31)]{Ba3}: $E_1$ denotes the vacuum energy counterterm, the sum of all bare vacuum graphs of orders $1\le\alpha+\beta\le\nbar$, and is not the case $k=1$ of $E_k$; similarly $E_0$ of the lemma below is the case $k=0$ of $E_k$, which vanishes.
\end{definition}

\begin{remark}[a grading that opens the coefficients fails]\label{app:G1:vacgrade-rem}
One could instead attach to each vertex (1.7) the graph of its coefficient and grade by all lines, including the attached ones; this is the literal reading of \cite[p.~418]{Ba3}. It defines different pieces term by term, and neither $\cV_k$ nor $E_{\rm fin}$ is then bounded. Take $N$ real components, $\ev{\varphi_a\varphi_b}=G\delta_{ab}$, $V=\lambda\abs\varphi^4+\delta m^2\abs\varphi^2$, $\delta m^2=-2(N+2)\lambda G$; the order-$\lambda$ vacuum energy is $\lambda N(N+2)\sum_xG(x,x)^2$. Let $k\ge1$ and write $G=A+B+R$, where $A$ collects the slices of index $<k$, $B$ is the slice $k$ and $R$ is the rest, and put $R':=C_0-G_K(0)\ge0$. In Definition~\ref{app:G1:vacgrade} the index-$k$ piece is $\lambda N(N+2)\sum_x(B^2+2BR)(x,x)$ and the final piece is $\lambda N(N+2)\sum_xR'(x,x)^2$. In the opened grading they are $\lambda N(N+2)\sum_x(B^2+2AB)(x,x)$ and $\lambda N(N+2)\sum_x(R'^2+2G_K(0)R')(x,x)$. Here $A(x,x)\ge(-\Delta_\eps+m^2+a_kt_k^{-2})^{-1}(x,x)\ge c\,\eps^{-1}$. So the opened index-$k$ piece exceeds the size $P(t_k)$ of Proposition~\ref{app:G1:I9} per scale-$k$ block by a factor of order $t_k/\eps$, and the opened final piece is at least of order $\lambda\eps^{-1}\abs T$. In Definition~\ref{app:G1:vacgrade} the term $2\lambda N(N+2)\sum_xA(x,x)B(x,x)$ lies in the steps $j<k$ given by the indices of the slices of $A$, where the coefficient $-2(N+2)\lambda B$ is a part of the degree-$0$ vertex $\delta m^2_{K,j+1}$.
\end{remark}

\begin{lemma}[grading identities]\label{app:G1:vacgrade-id}
(i) $E_1=E_{\rm fin}+E_K$ and $E_K=E_{K,k}+E_k$ for $0\le k\le K$, with $E_{K,K}=E_0=0$. Hence $E_1=E_{\rm fin}+\sum_{k<K}\cV_k$. The number $E_{\rm fin}$ does not depend on how $C_f-G_{K,f}(0)$ is decomposed.
(ii) For every exact Gaussian staging, $\cV_k$ is the field-independent part of order $\le\nbar$ of the output of step $k$.
\end{lemma}
\begin{proof}
(i) Every value is multilinear in the lines and every coefficient is a number. Sort the terms by their set of line indices. A term lies in $E_{\rm fin}$ exactly when at least one line carries a piece of $C_f-G_{K,f}(0)$.
(ii) Lemma~\ref{app:G1:S05-TS} identifies the step-$k$ output with the terms of the bare vacuum graphs whose largest line index is $k$. By (i) these are the terms of $\cV_k$.
\end{proof}

\begin{lemma}[the finite vertex]\label{app:G1:Efin-fv}
Assume \eqref{app:G1:eq-hyp}. For $K\le j\le\bar K$ let $\delta m^2_{\langle j\rangle}(x)$ be the expression $\delta m^2(x)$ in which every propagator, including those of the inner insertions in \cite[(1.23)]{Ba3}, is replaced by $\sum_{i\le j}C_{(i),f}$, and put $\delta m^2_{\succ j}:=\delta m^2-\delta m^2_{\langle j\rangle}$. Then $\delta m^2_{\succ\bar K}=0$, and for $K\le j<\bar K$, order by order and uniformly in $\eps$, $x$ and $\ell_T$,
\begin{gather*}
\sup_x\abs{\delta m^2_{\succ j}(x)}\,t_j^2\le\Gamma_f\Big[e(t_j)^2+\lambda(t_j)+e(t_j)^4+e(t_j)^2\lambda(t_j)+\lambda(t_j)^2\Big],\\
\Gamma_f:=C_b\big(1+\log_L^+(\bar\ell/\eps_0)\big)+C_b^R(1+Z_s+Z_v)^5 .
\end{gather*}
\end{lemma}
\begin{proof}
By multilinearity, $\delta m^2_{\succ j}$ is the sum of the index assignments of the graphs of $\delta m^2$ whose largest index $j'$ (inner insertions included) is $>j$. Since $j\ge K$, these are terms of $\delta m^2_{\rm fin}$, and the proof of Proposition~\ref{app:G1:I3b} bounds them for each $j'$ separately: for $j<j'<\bar K$ by $C_b(e^2+\lambda)t_{j'}^{-1}$ at one loop and $C_b(e^4+e^2\lambda+\lambda^2)$ at two loops; for $j'=\bar K$ by $C_b^R(e^2+\lambda)t_{\bar K}^{-1}(1+Z_s+Z_v)^2$ and $C_b^R(e^4+e^2\lambda+\lambda^2)(1+Z_s+Z_v)^5$. Use $\sum_{j'>j}t_{j'}^{-1}\le t_j^{-1}/(L-1)$, $t_{\bar K}^{-1}\le t_j^{-1}$, and $\bar K-j-1\le\bar K-K-1\le\log_L^+(\bar\ell/\eps_0)$ by (P4). Multiplying by $t_j^2$ and using $e(t)^2=e^2t$, $\lambda(t)=\lambda t$ gives the claim.
\end{proof}

\begin{lemma}[degrees with finite vertices]\label{app:G1:Efin-deg}
Let $Q$ be a connected graph built from the vertices of Definition~\ref{app:G1:vacgrade}, in which each vertex \cite[(1.7)]{Ba3} carries either a graph of $\delta m^2_{\langle j\rangle}$ attached as in \cite[p.~417]{Ba3}, or the number $\delta m^2_{\succ j}$ (a \emph{finite vertex}). Let $m_f$ be the number of finite vertices, $o'$ the total order of the other vertices, $E$ the number of external legs and $n_\partial$ the number of differentiations acting on external legs. Then
\[
D(Q)=o'+2m_f-3+\tfrac12E+n_\partial .
\]
If $m_f\ge1$, $Q$ has at least one line, $E\ge1$, and $Q$ does not vanish identically, then $D(Q)\ge\frac12$.
\end{lemma}
\begin{proof}
By \cite[(2.1)]{Ba3} with $d=3$, a leg on an internal line contributes $-\frac12$ and a differentiation of an internal line $-1$; external legs and their differentiations contribute $0$. So $D_Q(v)=D(v)+\frac12(\#\text{external legs at }v)+(\#\text{external differentiations at }v)$. By \cite[p.~423, (i), (iii), (iv)]{Ba3}, $D(v)=o(v)$ for the bare vertices. An attached graph $G_0$ has $D(v)+D(G_0)=o(G_0)$ (proof of Lemma~\ref{app:G1:cutclasses}(b)). A finite vertex has $D(v)=2$ and counterterm degree $0$ \cite[p.~423]{Ba3}. Now \cite[(2.3)]{Ba3} gives the formula. Every vertex has an even number of scalar legs; if $E=1$, the external leg is therefore a vector leg, and $Q$ vanishes \cite[p.~434]{Ba3}, as in Lemma~\ref{app:G1:cutclasses}(c). So $E\ge2$. If $o'>0$ then $o'\ge\frac12$ and $D\ge\frac12+2-3+1$. If $o'=0$, then $Q$ has a line only if $m_f\ge2$, and then $D\ge4-3+1$.
\end{proof}

\paragraph{Torus moments of the cut classes.} Let $K\le j_0\le\bar K$, $\xi=L^{-j_0}$, and work in scale-$j_0$ units. The rescaled torus is $T^\xi=\xi\Z^3/\Gamma$, $\Gamma=\bigoplus_\mu\ell'_\mu\Z e_\mu$, $\ell'_\mu\in\xi\N$. By (P1) and \eqref{app:G1:eq-hyp}, $t_{j_0}\le\bar\ell\le L$ and $\ell_T\ge L^3$, so $\ell':=\min_\mu\ell'_\mu\ge L^2\ge4$. For $w\in T^\xi$ let $\bar w\in\xi\Z^3$ be the representative with $\abs{\bar w_\mu}\le\ell'_\mu/2$, $\abs w_T:=\abs{\bar w}$, and $\rho_X(w):=\abs w_X+\xi$ for $X\in\{\xi\Z^3,T^\xi\}$. Let $C^\xi_\infty$ and $C^\xi_T$ be the resolvents of $-\Delta^\xi+1$ on $\xi\Z^3$ and on $T^\xi$ (componentwise for the photon); $C^\xi_T$ is the kernel $C^\xi$ of \cite[p.~437]{Ba3} on the torus, for both fields. $\partial^a$ is any product of $a$ forward or backward differences. We use: (C) $C^\xi_\infty$ is even and reflection invariant, and $\abs{\partial^aC^\xi_\infty(z)}\le C_0\rho(z)^{-1-a}\mathrm{e}^{-\delta_0\abs z}$ for $a\le2$ \cite[(2.10)--(2.12), p.~426, p.~437]{Ba3} (see (H$_\infty$) in Subsection~\ref{ssec:T6-IMG}); (Per) $C^\xi_T(w)=\sum_{n\in\Gamma}C^\xi_\infty(\bar w+n)$, by the proof of Lemma~\ref{lem:IMG}(i), since $-\Delta^\xi+1$ commutes with $\Gamma$; (W) on $\xi\Z^3$, with every line equal to $C^\xi_\infty$, the zeroth and first moments of C3 vanish \cite[(3.27)--(3.29), p.~441]{Ba3}. Let $N_0$ be the largest number of kernel factors (partner coefficients included) in a term of a member of C1$'$, C3 or C4, and $r_0:=1/(8N_0)$.

\begin{lemma}[images]\label{app:G1:TW-img}
For $n\in\Gamma\setminus\{0\}$, $\abs{\bar w+n}\ge\frac12\abs{\bar w}+\frac14\abs n$. Consequently, with $\operatorname{Per}F(w):=\sum_{n\in\Gamma}F(\bar w+n)$ for a function $F$ on $\xi\Z^3$ and $w\in T^\xi$:
(a) if $\abs{F(z)}\le A\mathrm{e}^{-\delta\abs z}$ on $\xi\Z^3$, then $\abs{\operatorname{Per}F(w)-F(\bar w)}\le Ac_\delta\mathrm{e}^{-\delta\abs w_T/2}$, $c_\delta:=\sum_{m\in\Z^3\setminus0}\mathrm{e}^{-\delta\abs m}$;
(b) $\abs{\partial^aC^\xi_T(w)}\le C_0'\rho_T(w)^{-1-a}\mathrm{e}^{-\delta_0\abs w_T/4}$ for $a\le2$;
(c) for $\xi\le1$: $\sum_{w\in X}\xi^3\rho_X(w)^{-k}\mathrm{e}^{-\delta\abs w_X}\le C(k,\delta)$ for $0\le k<3$, and $\le C(\delta)\xi^{-1}$ for $k=4$.
\end{lemma}
\begin{proof}
If $n_\mu\neq0$ then $\abs{n_\mu}\ge\ell'_\mu\ge2\abs{\bar w_\mu}$, so $\abs{\bar w_\mu+n_\mu}\ge\max(\abs{\bar w_\mu},\frac12\abs{n_\mu})$; this is trivial if $n_\mu=0$. Hence $\abs{\bar w+n}\ge\max(\abs{\bar w},\frac12\abs n)$; average. (a) $\abs n\ge4\abs m$ for $n=(\ell'_\mu m_\mu)_\mu$, so $\mathrm{e}^{-\delta\abs{\bar w+n}}\le\mathrm{e}^{-\delta\abs{\bar w}/2}\mathrm{e}^{-\delta\abs m}$; sum over $m\ne0$. (b) The term $n=0$ is (C). For $n\ne0$, $\abs{\bar w+n}\ge2$, so the terms are $\le C_0\mathrm{e}^{-\delta_0\abs{\bar w+n}}$, and (a) gives $C_0c_{\delta_0}\mathrm{e}^{-\delta_0\abs{\bar w}/2}\le C\rho_T^{-1-a}\mathrm{e}^{-\delta_0\abs{\bar w}/4}$. (c) $w\mapsto\bar w$ is injective and $\abs w_T=\abs{\bar w}$; compare with $\int(\abs z+\xi)^{-k}\mathrm{e}^{-\delta\abs z/2}d^3z$, the term $z=0$ contributing $\xi^{3-k}$.
\end{proof}

\begin{lemma}[short and long parts]\label{app:G1:TW-split}
Fix $\chi_0\in C^\infty([0,\infty))$, $0\le\chi_0\le1$, $\chi_0=1$ on $[0,\frac12]$, $\chi_0=0$ on $[1,\infty)$; put $\chi(z):=\chi_0(\abs z/r_0)$, $C^<:=\chi C^\xi_\infty$, $C^>:=C^\xi_\infty-C^<$. These are even and reflection invariant, and for $a\le2$: (a) $\partial^aC^<(z)=0$ for $\abs z\ge r_0+2\xi$ and $\abs{\partial^aC^<(z)}\le C_1\rho(z)^{-1-a}$; (b) $\abs{\partial^aC^>(z)}\le C_1\mathrm{e}^{-\delta_0\abs z}$; $C_1=C_1(C_0,\delta_0,r_0,\chi_0)$. With $C^<_T:=\operatorname{Per}C^<$, $C^>_T:=\operatorname{Per}C^>$ we have $C^\xi_T=C^<_T+C^>_T$; if $\xi<r_0$ then $\partial^aC^<_T(w)=\partial^aC^<(\bar w)$; and $\abs{\partial^aC^>_T(w)}\le C_2\mathrm{e}^{-\delta_0\abs w_T/2}$, with $C_2:=C_1(1+c_{\delta_0})$.
\end{lemma}
\begin{proof}
By $\partial_e(fg)(z)=(\partial_ef)(z)g(z+\xi e)+f(z)(\partial_eg)(z)$, every term of $\partial^a(\chi C^\xi_\infty)$ is $(\partial^b\chi)(z+\sigma)(\partial^{a-b}C^\xi_\infty)(z+\sigma')$ with $\abs\sigma,\abs{\sigma'}\le2\xi$ and $\abs{\partial^b\chi}\le c_br_0^{-b}$; also $\rho(z+\sigma')\ge\rho(z)/3$. (a) On the support $\rho(z)\le r_0+5\xi\le6$, so $r_0^{-b}\le(6/r_0)^b\rho(z)^{-b}$, and (C) gives the bound. (b) A term is non-zero only if $\abs{z+\sigma}\ge r_0/2$ for some $\abs\sigma\le2\xi$; then $\rho(z+\sigma')\ge r_0/12$ (from $\rho\ge\xi$ if $\xi\ge r_0/12$, otherwise from $\abs{z+\sigma'}\ge r_0/2-4\xi$), and (C) gives $C(r_0)\mathrm{e}^{-\delta_0\abs z+4\delta_0\xi}$. If $\xi<r_0$, the range $r_0+2\xi<\frac14\ell'$ of $\partial^aC^<$ is smaller than $\abs{\bar w+n}\ge2$ for $n\ne0$, so only $n=0$ contributes. For $C^>_T$ use Lemma~\ref{app:G1:TW-img}(a). Finally $C^\xi_T=C^<_T+C^>_T$ by (Per).
\end{proof}

A term of a class is a product of kernels $\partial^aC^\xi_X(u-v)$, with $u,v$ vertex positions or their shifts by at most $\xi$, all positions except the base $y$ summed; the coefficient of the partner of a subgraph $H$ is written as the graph of $H$ with one leg moved \cite[p.~417]{Ba3}. Expand every line as $C^\xi_X=C^<_X+C^>_X$; for an assignment $\mathfrak a$ of a letter $<$ or $>$ to every line of the parent graph (a subgraph and its partner carrying the same letters) this gives $\Pi_X=\sum_{\mathfrak a}\Pi^{\mathfrak a}_X$ for each moment $\Pi$; $\mathfrak s$ denotes the assignment with only $<$.

\begin{lemma}[no wrapping]\label{app:G1:TW-nowrap}
Let $\xi<r_0$. Then $\Pi^{\mathfrak s}_{T^\xi}=\Pi^{\mathfrak s}_{\xi\Z^3}$ for the zeroth moment of every member with its partners, and for the first moment, with $y'-y$ the minimal displacement on $T^\xi$.
\end{lemma}
\begin{proof}
Each term is connected: every point is linked to $y$ by at most $N_0$ kernels. By Lemma~\ref{app:G1:TW-split}(a) a non-vanishing kernel links points at torus distance $<r_0+4\xi$, so every point $p$ of a non-vanishing configuration has $\abs{p-y}_T<R_0:=N_0(r_0+4\xi)<\frac58$. Fix a lift $\tilde y$ of $y$ and put $\tilde p:=\tilde y+\overline{p-y}$. Then $\abs{\tilde u-\tilde v}<2R_0<\ell'/2$, so $\overline{u-v}=\tilde u-\tilde v$ and $\partial^aC^<_T(u-v)=\partial^aC^<(\tilde u-\tilde v)$. Every non-vanishing configuration on $\xi\Z^3$ lies in the ball of radius $R_0$ about $\tilde y$, which projects injectively to $T^\xi$ since $\abs n\ge\ell'>2R_0$ for $n\in\Gamma\setminus0$. So both sums run over the same configurations with the same values.
\end{proof}

\begin{proposition}[torus moments]\label{app:G1:TW}
Replace every line of the classes C1$'$ and C3 of Lemma~\ref{app:G1:cutclasses}, and of the partners of their proper divergent subgraphs, by $C^\xi_T$. Let $\Pi^{(1')}$, $\Pi^{(3)}$ be the zeroth moments at a base point $y\in T^\xi$ and $\Pi^{(3)}_\nu$ the first moment of C3. There is $C$, depending only on $d$ and on the classes (through $N_0$), such that, uniformly in $\xi$, $y$ and $\ell_T\ge L^3$: (i) $\abs{\Pi^{(3)}}+\sum_\nu\abs{\Pi^{(3)}_\nu}\le C$; (ii) $\abs{\Pi^{(1')}}\le C$.
\end{proposition}
\begin{proof}
If $\xi\ge r_0$, then $\rho\ge r_0$ and every kernel is $\le C(r_0)\mathrm{e}^{-\delta_0\abs w_T/4}$ by Lemma~\ref{app:G1:TW-img}(b); bound the kernels of a spanning tree by this and the others by their supremum, and use Lemma~\ref{app:G1:TW-img}(c) with $k=0$ (also for the weight $\abs w_T$). Let now $\xi<r_0$.

(i) By \cite[(3.27), first line, and (3.29)]{Ba3}, written with $C^\xi_X$, $\Pi^{(3)}$ is a linear combination, with coefficients depending only on $d$ and on the charge matrix, of bubble terms $\sum_w\xi^3(D_1C^\xi_X)(w)(D_2C^\xi_X)(w)$, $\operatorname{ord}D_1+\operatorname{ord}D_2=2$, and local terms $(D_3C^\xi_X)(\sigma)$, $\operatorname{ord}D_3\le2$, $\abs\sigma\le2\xi$ (translation invariance and evenness were used). $\Pi^{(3)}_\nu$ is a combination of the weighted bubbles $\sum_w\xi^3\bar w_\nu(D_1C^\xi_X)(D_2C^\xi_X)(w)$. Expand $C^\xi_X=C^<_X+C^>_X$. Every $\mathfrak a\ne\mathfrak s$ has a factor $C^>_X$: a local term with it is $\le C_2$; a bubble with orders $(1,1)$, or with $C^>$ carrying the order-$2$ factor, is $\le C\sum\xi^3\rho_X^{-2}\mathrm{e}^{-\delta\abs w_X}\le C$ (Lemmas~\ref{app:G1:TW-img}(c), \ref{app:G1:TW-split}); if $C^>$ carries the order-$0$ factor, sum by parts once, which is exact on $T^\xi$ by periodicity and on $\xi\Z^3$ by finite range, and obtain the same bound; in the weighted bubbles $\abs{\bar w_\nu}\rho^{-3}\le\rho^{-2}$. By (W) and Lemma~\ref{app:G1:TW-nowrap}, $\Pi^{(3)}_T=\Pi^{(3)}_T-\Pi^{(3)}_{\xi\Z^3}=\sum_{\mathfrak a\ne\mathfrak s}(\Pi^{\mathfrak a}_T-\Pi^{\mathfrak a}_{\xi\Z^3})$, and likewise for $\Pi^{(3)}_\nu$; both sums are bounded.

(ii) By \cite[(3.21)--(3.24), pp.~438--439]{Ba3}, written with $C^\xi_T$, $\Pi^{(1')}$ is a linear combination of $E_\sigma:=\sum_w\xi^3C^\xi_T(w+\sigma)(\partial_\mu C^\xi_T)(w)$ and $\xi C^\xi_T(\sigma')$, $\abs\sigma,\abs{\sigma'}\le2\xi$. Translation invariance gives $E_0=-\frac\xi2\sum_w\xi^3\abs{\partial_\mu C^\xi_T(w)}^2$, which is $\le C\xi\cdot\xi^{-1}$ by Lemma~\ref{app:G1:TW-img}(b),(c). Also $\abs{E_\sigma-E_0}\le C\sum_w\xi^3\xi\rho_T^{-4}\mathrm{e}^{-\delta\abs w_T}\le C$ and $\xi\abs{C^\xi_T(\sigma')}\le C$.
\end{proof}

\begin{proposition}[a bound on $E_{\rm fin}$ uniform in the torus, at $u=0$]\label{app:G1:Efin}
Let $M\in(0,m_2]$ and assume \eqref{app:G1:eq-hyp}. Let $E_{\rm fin}$ be as in Definition~\ref{app:G1:vacgrade}. Put $o:=\alpha/2+\beta$, $P_E(\eps_0):=\sum_{2\le\alpha+2\beta\le6}\abs e^\alpha\lambda^\beta\eps_0^{o}$ and $\Phi:=(1+\Gamma_f)^{N_v}$, with $\Gamma_f$ as in Lemma~\ref{app:G1:Efin-fv} and $N_v=N_v(\nbar)$ the largest number of vertices of a bare vacuum graph of order $\le\nbar$. Then, uniformly in $\eps=2^{-K}\eps_0$ and in $\ell_T\ge\ell_0$,
\begin{multline*}
\abs{E_{\rm fin}}\le\abs{T_\eps}\Big\{C_E\,\Phi\Big[L^3\eps_0^{-3}P_E(\eps_0)+\sum_{\alpha+2\beta=6}\abs e^\alpha\lambda^\beta\big(2+\log_L^+(\bar\ell/\eps_0)\big)+\sum_{7\le\alpha+2\beta\le2\nbar}\frac{\abs e^\alpha\lambda^\beta t_{\bar K}^{o-3}}{1-L^{-(o-3)}}\Big]\\
+C_E^Rt_{\bar K}^{-3}\sum_{2\le\alpha+2\beta\le2\nbar}\abs e^\alpha\lambda^\beta t_{\bar K}^{o}(1+Z_s+Z_v)^{N_\ell}\Big\},
\end{multline*}
with $C_E=C(\alpha_0,\nbar,d,L,a,M_b)$, $C_E^R=C(\alpha_0,\nbar,d,L,a,M_b)(1+\rho_{\min}^{-2}+\kappa_{R,\min}^{-3})^{N_\ell}$, $N_\ell=N(\nbar)$ the maximal number of lines, and $t_{\bar K}^{-3}<\max(LM,L\mu_*,1)^3$. The torus side enters only through $Z_s$ and $Z_v$; $\Gamma_f$ and the whole right side are non-increasing in $\ell_T$.
\end{proposition}
\begin{proof}
\emph{Step 1 (top index and vertex split).} For a term of $E_{\rm fin}$ let $j^*\in\{K,\dots,\bar K\}$ be its largest line index; then $E_{\rm fin}=\sum_{j^*=K}^{\bar K}E^{[j^*]}$. In $E^{[j^*]}$ write every $\delta m^2$ as $\delta m^2_{\langle j^*\rangle}+\delta m^2_{\succ j^*}$ (Lemma~\ref{app:G1:Efin-fv}). Write the first part in its graph representation \cite[p.~417]{Ba3} (the vertex replaced by the graph with both legs localised at its position and the other leg summed); its attached lines have indices $\le j^*$. The second part is a finite vertex. By multilinearity, $E^{[j^*]}$ is the sum of the terms with all lines, attached or not, of index $\le j^*$ and at least one non-attached line of index $j^*$. Decompose by Ba{\l}aban's orderings \cite[(3.4), p.~424]{Ba3}, with ties broken by a fixed order of the lines of each parent graph in which attached lines precede the others, so that the last line $l(m)$ is a non-attached line of index $j^*$. A term in which a vertex carries a graph $G_0$ of $\delta m^2_{\langle j^*\rangle}$ is the counterterm member of the parent in which that vertex is replaced by $G_0$. By Lemma~\ref{app:G1:VSR-Btop}, with $j_{l(m)}=j^*$ fixed, every divergent proper subgraph renormalised within an ordering is a $G_i$ with $i<m$, and its partner has attached lines of index $<j_0\le j^*$; so the partner is present in $E^{[j^*]}$. The transformations \cite[(3.5), pp.~434--445]{Ba3} act inside the sub-sum, and no term is moved to another top index. By Lemma~\ref{app:G1:common}, every renormalised subgraph has $j_0\le\bar K\le\min_fK_f$, its internal lines resum to $G_{j_0,f}(0)$, and this is the case covered by (S6).

\emph{Step 2 (the cut).} The proof of Lemma~\ref{app:G1:cut} uses only the definition of $J(\bar l)$ and $l(m)\notin G_i$ for $i<m$, so it applies at top index $j^*$. Cutting $l(m)$, i.e.\ treating its kernel $K_{j^*}=C_{(j^*),f}$ as a two-variable external field \cite[p.~419, p.~422]{Ba3}, maps every class term by term onto two-leg classes of $G\setminus l(m)$. For $j^*<\bar K$, in scale-$j^*$ units, $\norm{hK_{j^*}h'}_{1,\alpha_0}\le O(1)\mathrm{e}^{-\delta_1\dist(\supp h,\supp h')}$ by (VS). By Lemma~\ref{app:G1:cutclasses}(c), $G\setminus l(m)$ is connected or the term vanishes. By Lemma~\ref{app:G1:cutclasses}(d) and Lemma~\ref{app:G1:Efin-deg}, if $D(G\setminus l(m))\le0$ then:
\begin{itemize}
\item[(a)] $G$ is a single finite vertex closed by $l(m)$; its value at $x$ is a multiple of $\delta m^2_{\succ j^*}(x)\tr K_{j^*}(x,x)$, which in scale-$j^*$ units is at most $O(1)$ times the bound of Lemma~\ref{app:G1:Efin-fv};
\item[(b)] $G\setminus l(m)$ is a one-particle-irreducible member of C1 or C2; its partner, the vertex carrying $G_0\cong G\setminus l(m)$ with lines of index $<j^*$ closed by $l(m)$, belongs to $\delta m^2_{\langle j^*\rangle}$, and the class is renormalised as a pinned two-leg class by Lemma~\ref{app:G1:VSR-C};
\item[(c)] $G\setminus l(m)$ is a one-particle-reducible member of C2; it has no partner, and its class is bounded by Lemma~\ref{app:K:L04-3} with $F=K_{j^*}$;
\item[(d)] $G\setminus l(m)$ is of type C1$'$, C3 or C4; its class is bounded by Lemma~\ref{app:K:L04-1} with $F=K_{j^*}$, based at the pinned endpoint. For C1$'$ no hypothesis is needed; for C3 the hypothesis of that statement is Proposition~\ref{app:G1:TW}(i); for C4 it is Lemma~\ref{app:K:L04-2}. On the torus these moments need not vanish: the identities \cite[(3.27)--(3.29)]{Ba3} hold on $\xi\Z^3$ only, and the one-loop zeroth moment on $T^\xi$ is a sum over the non-zero images.
\end{itemize}
For a function $F$ on $T_\eps\times T_\eps$, in the units of the scale considered, let $N(F)$ be the maximum over $a,b\in\{0,1\}$ of $\sup\abs{\partial_1^a\partial_2^bF}$ plus the H\"older-$\alpha_0$ quotients of $\partial_1^a\partial_2^bF$ in either variable; for a function of one variable $N$ is defined in the same way, and $N(F)\le N(\Phi_1)N(\Phi_2)$ if $F(y,z)=\Phi_1(y)\Phi_2(z)$. For $j^*<\bar K$, $N(K_{j^*})\le C$ by \cite[(2.10)--(2.11), p.~426]{Ba3} and (VS). Every other component of $G\setminus l(m)$ has $D\ge\frac12$. Every proper component containing a finite vertex has $D\ge\frac12$ or vanishes (Lemma~\ref{app:G1:Efin-deg}). Every other proper component with $D\le0$ is renormalised by (S6), as in part~(i) of the proof of Proposition~\ref{app:G1:I3b}.

\emph{Step 3 ($K\le j^*<\bar K$).} All lines are slices: those of index $<K$ obey \cite[(2.10)--(2.12), p.~426]{Ba3}, and those of index $\ge K$ obey them by (VS) (Theorem~\ref{app:G1:VSR-VS}, resting on Lemma~\ref{app:G1:VSR-F}). Pin an endpoint $x$ of $l(m)$ and work in scale-$j^*$ units, with couplings $e(t_{j^*})$, $\lambda(t_{j^*})$ and finite-vertex coefficients $\delta m^2_{\succ j^*}t_{j^*}^2$ \cite[p.~418]{Ba3}. By Step~2, after renormalisation every component $G_i^{(\alpha)}$, $i<m$, has degree $\ge\vartheta>0$. Lemma~\ref{app:G1:VSR-A}(ii) with the largest index fixed to $j^*$ (no sign condition on $D(G)$, and $L^{j^*}\eta=1$), Lemmas~\ref{app:K:L04-1} and~\ref{app:K:L04-3} for the classes (c), (d), and Lemma~\ref{app:G1:Efin-fv} bound the pinned value of each class by $C\abs e^\alpha\lambda^\beta t_{j^*}^{o}\Gamma_f^{m_f}$, uniformly in $x$, where $(\alpha,\beta)$ is the total order including the finite vertices. The number of graphs and orderings depends only on $\nbar$, and there are $\abs{T_\eps}t_{j^*}^{-3}$ pin blocks. Hence $\abs{E^{[j^*]}}\le C_E\Phi\abs{T_\eps}\sum_{\alpha,\beta}\abs e^\alpha\lambda^\beta t_{j^*}^{o-3}$.

\emph{Step 4 ($j^*=\bar K$).} Here $\delta m^2_{\succ\bar K}=0$, so there are no finite vertices. Lemma~\ref{app:G1:TB}, with one vertex pinned per scale-$\bar K$ block, gives $\abs{E^{[\bar K]}}\le C_E^R\abs{T_\eps}t_{\bar K}^{-3}\sum\abs e^\alpha\lambda^\beta t_{\bar K}^{o}(1+Z_s+Z_v)^{N_\ell}$. In it, every lower cluster of type C1$'$, C3, C4, or a one-particle-reducible member of C2, is bounded by Lemma~\ref{app:K:L04-1} or \ref{app:K:L04-3}; the field on its two legs is one $R$-line, or a product $\Phi_1(y)\Phi_2(z)$ of two $R$-lines with their far ends held fixed and summed as in Lemma~\ref{app:G1:TB}, so $N(F)\le C(1+Z_s+Z_v)^2$ by Proposition~\ref{app:G1:VSR-R}, and these factors are among those counted by $(1+Z_s+Z_v)^{N_\ell}$.

\emph{Step 5 (sum over $j^*$).} The integer $\alpha$ is even, so $o\in\Z$. For $o\le2$: $\sum_{j^*\ge K}t_{j^*}^{o-3}\le\eps_0^{o-3}/(1-L^{-1})\le L^3\eps_0^{o-3}$. For $o=3$: there are $\bar K-K\le1+\log_L^+(\bar\ell/\eps_0)$ values of $j^*<\bar K$, by (P4). For $o\ge4$: $\sum_{j^*<\bar K}t_{j^*}^{o-3}\le t_{\bar K}^{o-3}/(1-L^{-(o-3)})$. Finally $t_{\bar K}$, $C_E$, $C_E^R$ do not depend on $\ell_T$, while $Z_f$, and with it $\Gamma_f$, is non-increasing in $\ell_T$ by (P6).
\end{proof}

\begin{remark}\label{app:G1:Efin-rem}
In the grading of Definition~\ref{app:G1:vacgrade}, $E_{\rm fin}$ contains the terms whose lines lie above $\eps_0$ while their $\delta m^2$ coefficients contain lines below $\eps_0$; it does not contain the terms in which only a coefficient line lies above $\eps_0$, which belong to $E_K$ with a finite vertex. The factor $\Phi$ accounts for the finite vertices $\delta m^2_{\succ j^*}$. The number $\delta m^2_{\rm fin}:=\delta m^2-\delta m^2_K$ is a coefficient and keeps the meaning of \cite[(1.26)]{Ba3}; by Lemma~\ref{app:G1:dm2grade} its bounds are the same in both gradings. The form of the bound at $u\ne0$ is Proposition~\ref{app:G1:I8}.
\end{remark}

\begin{remark}[the mass range]\label{app:G1:N9-supply}
(i) For $M\le m_2$, $\bar\ell(M)\ge\bar\ell(m_2)$. So (E.fin) gives \eqref{app:G1:eq-hyp}, and (T.floor) gives $\ell_T\ge L^3$.
(ii) The construction slices both fields to one index $\bar K$. The order of $M$ and $\mu_*$ is used nowhere: when $M<\mu_*$, as at $M_1=0.7m$, the scalar's $R$-line is the long one and is bounded by (Rem) in the same way. At $M=\mu_*$ the two fields coincide in the construction.
(iii) $C_b^R$ and $C_E^R$ are non-increasing in $\rho_{\min}$, since $\kappa_{R,\min}$ is non-decreasing in it. So they may be evaluated at the lower bound (P5), which is the quantity $\rho_{\min}(M)$ of Lemma~\ref{lem:B:massrange}.
\end{remark}

\paragraph{The \texorpdfstring{$A\cdot A$}{A.A} vertices.}\label{app:G1:ssec-AA}
Work in step-$k$ units: $\eta=L^{-k}$, and $\ell_j:=L^j\eta\le1$ is the scale of a line of index $j$. The vertices of \cite[(3.30), p.~442]{Ba3} are
\[
e(t_k)^2\sum_x\eta^d\,g(x)A_\mu(x)\,c_{\mu\mu'}(x)\,g'(x)A'_{\mu'}(x)
\]
and the first-moment vertex, with coefficient $c_{\mu\mu'\nu}$. Here $j_0$ is the lowest $j$-index of the external legs. On the $\xi$-lattice, $\xi:=L^{-j_0}$, put $C^\xi:=(-\Delta^\xi+1)^{-1}$ and $G^\xi_{j_0}(0):=(-\Delta^\xi+m^2_{j_0}+a_{j_0}P_{j_0})^{-1}$. The rescaled Higgs mass is $m^2_{j_0}=M^2t_{j_0}^2$, and $P_{j_0}$ is the block projection. Write $\rho(y,y'):=\abs{y-y'}+\xi$.

\begin{lemma}[coefficient bounds for the $A\cdot A$ vertices]\label{app:G1:AA}
Let $0\le m^2_{j_0}\le1$. Then
\[
\abs{c_{\mu\mu'}(x)}\le\Gamma_3\,\ell_{j_0}^{-1},\qquad\abs{c_{\mu\mu'\nu}(x)}\le\Gamma_3'\,\ell_{j_0}^{0}.
\]
The constants $\Gamma_3,\Gamma_3'$ depend only on $d$, $a$ (through $a_\flat\le a_{j_0}\le a$) and the block size. They are uniform in $m^2_{j_0}\in[0,1]$, $\xi$, $j_0$, $k$, $\eps$ and $\ell_T$, and are free of $\mu_*$, $u$ and the photon zero mode. In particular, under (E.m2) and $t_{j_0}\le\eps_0$ the bounds hold with the same constants at every Higgs mass $M\in(0,m_2]$.
\end{lemma}
\begin{proof}
\emph{Step 1 (resolvent identity).} $G^\xi_{j_0}(0)-C^\xi=G^\xi_{j_0}(0)\big[(C^\xi)^{-1}-G^\xi_{j_0}(0)^{-1}\big]C^\xi=G^\xi_{j_0}(0)\,\Upsilon\,C^\xi$, with $\Upsilon:=1-m^2_{j_0}-a_{j_0}P_{j_0}$. Because $P_{j_0}$ is an orthogonal projection, $\Upsilon$ has spectrum in $\{1-m^2_{j_0},\,1-m^2_{j_0}-a_{j_0}\}$. So $\norm \Upsilon\le1+a$ for $m^2_{j_0}\in[0,1]$. Moreover $\Upsilon=(1-m^2_{j_0})\mathbf 1-a_{j_0}P_{j_0}$, where the kernel of $P_{j_0}$ is bounded and supported in $\abs{z-z'}\le2\sqrt d$.

\emph{Step 2 (the line bounds and the mass).} At zero background, $G^\xi_{j_0}(0)$ is the sum of its slices of index $<j_0$, together with the bottom propagator. By Theorem~\ref{app:G1:VSR-VS}, which covers the scalar half at mass $M$ for $t_{j+1}\le\min(1/M,L)$ (here $t_{j_0}\le\eps_0\le1/m_2\le1/M$), it obeys $\abs{G^\xi_{j_0}(0)(y,y')}\le O(1)\rho(y,y')^{-1}\mathrm{e}^{-\delta\abs{y-y'}}$ and $\abs{\partial G^\xi_{j_0}(0)(y,y')}\le O(1)\rho^{-2}\mathrm{e}^{-\delta\abs{y-y'}}$, where $\rho(y,y'):=\abs{y-y'}+\xi$. The $O(1)$ is uniform in $M$ and in $\xi$. The mass enters the quadratic form only through $+m^2_{j_0}\ge0$, so it does not weaken the coercivity.

\emph{Step 3 (the replacement difference is one dimension better).} Put $D:=G^\xi_{j_0}(0)-C^\xi=G^\xi_{j_0}(0)\Upsilon C^\xi$. By Steps 1--2,
\[
\abs{D(y,y')}\le O(1)\sum_{z,z'}\xi^{2d}\rho(y,z)^{-1}\mathrm{e}^{-\delta\abs{y-z}}\,1_{\abs{z-z'}\le2\sqrt d}\,\rho(z',y')^{-1}\mathrm{e}^{-\frac12\abs{z'-y'}}\le O(1)\mathrm{e}^{-\delta'\abs{y-y'}},
\]
because $\int d^3z\,\abs{y-z}^{-1}\abs{z-y'}^{-1}\mathrm{e}^{-\delta(\abs{y-z}+\abs{z-y'})}$ is bounded uniformly in $y,y'$ in $d=3$. The display is the contribution of $a_{j_0}P_{j_0}$. The contribution of $(1-m^2_{j_0})\mathbf 1$ is the same expression with $z'=z$ and $\xi^{d}$ in place of $\xi^{2d}$, and it obeys the same bound. With one difference at one end, $\rho^{-2}*\rho^{-1}$ in $d=3$ gives $O(1)(1+\log_+(1/\rho))\mathrm{e}^{-\delta'\abs{y-y'}}$.

With one difference at each end, $\abs{\partial_y\partial_{y'}D(y,y')}\le O(1)\rho(y,y')^{-1}\mathrm{e}^{-\delta'\abs{y-y'}}$, after decreasing $\delta'$. Indeed, $\partial_y\partial_{y'}D=(\partial G^\xi_{j_0}(0))\,\Upsilon\,(\partial C^\xi)$, and $\abs{\partial C^\xi(z,y')}\le c\rho(z,y')^{-2}\mathrm{e}^{-\abs{z-y'}/4}$ by Lemma~\ref{app:G1:VSR-F}.
\begin{itemize}
\item The part with $a_{j_0}P_{j_0}$ has a bounded kernel supported in $\abs{z-z'}\le2\sqrt3$. Since $\rho^{-2}$ is summable near the diagonal in $d=3$, this part is at most $c\,\mathrm{e}^{-\delta''\abs{y-y'}}\le c'\rho^{-1}\mathrm{e}^{-\delta''\abs{y-y'}/2}$.
\item In the part with $(1-m^2_{j_0})$, put $r:=\abs{y-y'}$ and $\delta_0:=\min(\delta,\frac14)$, so that $\mathrm{e}^{-\delta\abs{y-z}-\abs{z-y'}/4}\le \mathrm{e}^{-\delta_0r}$. On $\{\abs{z-y'}\le2r\}$, one of $\rho(y,z),\rho(z,y')$ is at least $\rho(y,y')/2$, and $\sum_{\abs w\le3r}\xi^3\rho(w)^{-2}\le c\rho(y,y')$; this gives $c\rho(y,y')^{-1}$. On $\{\abs{z-y'}>2r\}$ we have $\abs{y-z}\ge\abs{z-y'}/2$, and $\sum_{\abs w>2r}\xi^3\rho(w)^{-4}\le c\rho(y,y')^{-1}$.
\end{itemize}

\emph{Step 4 (exact zeros).} In $\xi$-units the order-$e^2$ class gives the zeroth-moment coefficient $\Pi_{\mu\mu'}$ (degree $-d+2=-1$) and the first-moment coefficient $\Pi_{\mu\mu'\nu}$ (degree $-d+3=0$). Their bubble terms are of size $\sum_{y'}\xi^d\rho^{-4}$ and $\sum_{y'}\xi^d\rho^{-4}\rho$, respectively. Replace every $G_{j_0}(0)$ by $C^\xi$. Then the $\Pi_{\mu\mu'\nu}$ combination is \cite[(3.27)]{Ba3}, which vanishes by the $\mu\leftrightarrow\mu'$ symmetry \cite[(3.28)]{Ba3}. The $\Pi_{\mu\mu'}$ combination \cite[(3.29)]{Ba3} vanishes by the Ward--Takahashi identity \cite[(2.26)]{Ba3}; see also Lemma~\ref{app:G1:AA4-WT} with $M^2=1$. (The ``$M^2=1$'' of \cite{Ba3} is the normalisation of $C^\xi$, not the Higgs mass.)

\emph{Step 5 (remainders).} Write $\nabla_\nu f(y)=\xi^{-1}(f(y+\xi e_\nu)-f(y))$ and $\nabla^-_\nu f(y)=\xi^{-1}(f(y)-f(y-\xi e_\nu))$. For absolutely summable products, $\sum_y\xi^3(\nabla_\nu f)h=-\sum_y\xi^3f\,\nabla^-_\nu h$, and $(\nabla^-_\nu h)(y)=(\nabla_\nu h)(y-\xi e_\nu)$.

\emph{Line bounds.} By Step 3 and Lemma~\ref{app:G1:VSR-F}, a line with $a\le2$ differences, at most one at each end, obeys
\[
\abs{\partial^aC^\xi}\le c\rho^{-1-a}\mathrm{e}^{-\abs{y-y'}/4},\qquad\abs{\partial^aD}\le c\,h_a(\rho)\mathrm{e}^{-\delta'\abs{y-y'}},
\]
with $h_0=1$, $h_1=1+\log_+(1/\rho)$ and $h_2=\rho^{-1}$. Expand every scalar line as $G^\xi_{j_0}(0)=C^\xi+D$. The terms with $C^\xi$-lines only are the combinations of Step~4, which vanish. Every other term contains a line $D$.

\emph{First moment.} Every term with a line $D$ is at most $c\rho^{-3}\cdot\rho\,\mathrm{e}^{-\delta'\abs{y-y'}}$, and its sum over $y'$ is $O(1)$.

\emph{Zeroth moment.} The class consists of two kinds of graph:
\begin{itemize}
\item the bubble: two vertices $\varphi\cdot q\nabla\varphi\,A$ at $x,x'$, joined by two scalar lines that carry $a$ and $2-a$ differences, at most one at each end;
\item the seagull tadpole: one vertex at $x$ with a scalar self-loop, contributing only at $x'=x$.
\end{itemize}
In the tadpole, $\abs{D(x,x)}\le c$. In the bubble, every product containing a line $D$, except $(\nabla^{(1)}_\mu\nabla^{(2)}_\nu C^\xi)\cdot D$, is at most $c\rho^{-2}(1+\log_+(1/\rho))\mathrm{e}^{-\delta'\abs{x-x'}}$, and its sum $\sum_{x'}\xi^3$ is $O(1)$.

In the exceptional product, both differences act on $C^\xi$, the difference of the vertex at $x'$ acting at the $x'$-end. The remaining factor at $x'$ is the constant external field. Sum by parts in $x'$. The difference of a constant is $0$, so the difference moves to the line $D$, and we get
\[
-\sum_{x'}\xi^3\,\nabla^{(1)}_\mu C^\xi(x,x')\,(\nabla^{(2)}_\nu D)(x,x'-\xi e_\nu).
\]
Since $\rho(x,x'-\xi e_\nu)\ge\rho(x,x')/2$, this is at most $c\sum_{x'}\xi^3\rho^{-2}(1+\log_+(1/\rho))\mathrm{e}^{-\delta'\abs{x-x'}}\le c'$. Without this summation by parts, the exceptional product is only bounded by $c\sum_{x'}\xi^3\rho^{-3}\mathrm{e}^{-\delta'\abs{x-x'}}=O(\log(1/\xi))$.

So both coefficients are $O(1)$ in $\xi$-units, uniformly in $\xi$ and in $m^2_{j_0}\in[0,1]$.

\emph{Step 6 (rescaling and the torus).} \cite[p.~441]{Ba3} gives $\Pi^{(\eta)}_{\mu\mu'}(x)=(L^{j_0}\eta)^{-d+2}\Pi^{(\xi)}_{\mu\mu'}(y)$ and $\Pi^{(\eta)}_{\mu\mu'\nu}(x)=(L^{j_0}\eta)^{-d+3}\Pi^{(\xi)}_{\mu\mu'\nu}(y)$. At $d=3$ these are the factors $\ell_{j_0}^{-1}$ and $\ell_{j_0}^0$. On $T_\eps$ the expression lives on the cube $\square$ of \cite[p.~433]{Ba3}, which does not wrap around the torus by (T.nw). The gauge-out of the background on $\square$ and the transfer to $\eta\Z^3$ are exact, and the image differences are external kernels of $\delta G$ type of \cite[(1.33)]{Ba3}. Only scalar lines occur, so $\mu_*$, $u$ and the photon zero mode do not enter.

\emph{Step 7 (the Higgs mass).} The Higgs mass entered only through $m^2_{j_0}$, in $\Upsilon$ (Step 1) and in the line bounds (Steps 2--3), and both are uniform on $[0,1]$. For $M\le m_2$, $t_{j_0}\le\eps_0$ and (E.m2), $m^2_{j_0}\le m_2^2\eps_0^2\le1$. No restriction to $M\in[m_1,m_2]$ is used.
\end{proof}

\paragraph{Setting for the vertex of \texorpdfstring{\cite[(3.32)]{Ba3}}{(3.32)}.}
We keep the notation of Lemma~\ref{app:G1:AA}:
\begin{itemize}
\item $\xi=L^{-j_0}$, $\ell:=\ell_{j_0}=L^{j_0}\eta\le1$ and $s_j:=L^{j-j_0}$ for $0\le j\le j_0$;
\item $C^\xi=(-\Delta^\xi+1)^{-1}$ and $G:=G^\xi_{j_0}(0)$;
\item $\Upsilon=1-m^2_{j_0}-a_{j_0}P_{j_0}$, and $D:=G-C^\xi=G\Upsilon C^\xi$, the \emph{regular line}.
\end{itemize}
On $\xi\Z^3$ we use $\ip fh=\sum_y\xi^3f(y)\cdot h(y)$ and write kernels with respect to this measure. The scalar field is a map $\varphi:\xi\Z^3\to\R^2$, and $q=\begin{pmatrix}0&-1\\1&0\end{pmatrix}$. For a real field $a$ on the bonds $b=(y,y+\xi e_\nu)$, the covariant Laplacian is given by
\[
\ip{\varphi}{-\Delta^\xi_a\varphi}=\sum_{b=(y,y+\xi e_\nu)}\xi^{3}\xi^{-2}\bigl|\mathrm{e}^{-q\xi a(b)}\varphi(y+\xi e_\nu)-\varphi(y)\bigr|^2 .
\]
On $\Lambda_R:=\xi\Z^3\cap[-R,R]^3$ the Dirichlet version $-\Delta^{\xi,R}_a$ is the same form on fields with $\varphi=0$ off $\Lambda_R$. A gauge function $\chi$ acts by $\varphi\mapsto \mathrm{e}^{q\chi}\varphi$ and $a\mapsto a+d\chi$, where $(d\chi)(b)=\xi^{-1}(\chi(y+\xi e_\nu)-\chi(y))$.

\emph{Labels} \cite[pp.~419--420]{Ba3}. The labels are as follows:
\begin{itemize}
\item every internal vector leg carries a factor $t_\square g_\square$, where $(g_\square)$ is the unit-scale partition of unity of \cite{Ba3}, $\supp g_\square$ lies in a cube of side $2$, $\abs g\le1$, and $\abs{\nabla g}+\abs{\nabla\nabla g}\le c_g$;
\item every $\lambda$-vertex carries the indicator of a unit cube and a parameter $s_\square$;
\item every \emph{fixed mass vertex} carries the indicator of a unit cube and a parameter $r_\square$. A fixed mass vertex is an insertion $\kappa(y)\abs{\varphi(y)}^2$ with an $a$-independent coefficient: the finite counterterms of \cite[(1.7), (1.29)]{Ba3}, and the differences between counterterm lines built on the torus and on the infinite lattice.
\end{itemize}
A \emph{label configuration} $\Lambda$ is a monomial $t^ms^nr^p$.

\emph{Counterterms in graph representation.} For a connected graph $H$ with two external scalar legs at $x,x'$, let $H_{\rm nn}$ be the sum of its terms in which no difference acts on an external leg. Following \cite[(1.21), p.~416; p.~417; (3.7)--(3.9), p.~435]{Ba3}, the mass counterterm generated by $H$ is written as the graph of $H_{\rm nn}$, with the same lines and indices, in which the external leg at $x'$ is attached at $x$; $x'$ is summed and the sign is $-1$. Within a graph, this counterterm graph is the \emph{counterterm partner} of the graph containing $H$. For a member $T$, let $\mathcal K(T)$ be the set formed by $T$ and its counterterm partners.

\emph{The class.} Fix $(\alpha,\beta)\in\{(4,0),(2,1)\}$ and $\Lambda$. Let $\Pi^\Lambda_{G,\mu\nu}(b,b')$ be the sum of all connected graphs on $\xi\Z^3$ (the \emph{members}) with the following properties:
\begin{itemize}
\item they have order $e_\xi^\alpha\lambda_\xi^\beta$, label configuration $\Lambda$, and two external vector legs, at the bond $b$ in direction $\mu$ and at the bond $b'$ in direction $\nu$; the external couplings and the localisations $g,g'$ are removed;
\item every internal scalar line, including those of the counterterms in graph representation, is $G$;
\item every internal vector line is the covariance $\mathcal C_{\rm ph}$ of a centred Gaussian field on the bonds that does not depend on the external field.
\end{itemize}
Let $\Pi^\Lambda_C$ be the same sum with every internal scalar line replaced by $C^\xi$. For a member $T$ and a set $S$ of its scalar lines, $T^S$ denotes $T$ with the lines of $S$ replaced by $D$; the other lines stay $G$.

\begin{lemma}[Dirichlet scalar propagator]\label{app:G1:AA4-Dir}
Let $M^2>0$, let $C^\xi_M:=(-\Delta^\xi_0+M^2)^{-1}$ on $\xi\Z^3$, and let $C^{\xi,R}_M:=(-\Delta^{\xi,R}_0+M^2)^{-1}$ on $\Lambda_R$. Then:
\begin{itemize}
\item $0\le C^{\xi,R}_M(y,y')\le C^{\xi,R'}_M(y,y')\le C^\xi_M(y,y')$ for $R\le R'$;
\item $C^{\xi,R}_M(y,y')\to C^\xi_M(y,y')$ as $R\to\infty$;
\item $C^\xi_M(y,y')\le c(\xi,M)\,\theta^{\abs{y-y'}_1/\xi}$, with $\theta:=6/(6+M^2\xi^2)<1$.
\end{itemize}
\end{lemma}
\begin{proof}
Write $-\Delta^\xi_0+M^2=(6\xi^{-2}+M^2)(1-T)$, where $T=(6+M^2\xi^2)^{-1}A$ and $A$ is the adjacency matrix of $\xi\Z^3$. So $T$ has nonnegative entries and row sums $\theta<1$, and $C^\xi_M=(6\xi^{-2}+M^2)^{-1}\xi^{-3}\sum_{n\ge0}T^n$ entrywise.

On $\Lambda_R$ the same holds with $1_{\Lambda_R}T1_{\Lambda_R}$ in place of $T$. Its powers are nonnegative, nondecreasing in $R$, and converge entrywise to those of $T$. Monotone convergence gives the first two claims.

Finally, $T^n(y,y')=0$ unless $n\ge\abs{y-y'}_1/\xi$, and $T^n(y,y')\le\theta^n$. Hence $C^\xi_M(y,y')\le(6\xi^{-2}+M^2)^{-1}\xi^{-3}(1-\theta)^{-1}\theta^{\abs{y-y'}_1/\xi}$.
\end{proof}

\begin{lemma}[Ward--Takahashi identity at fixed labels]\label{app:G1:AA4-WT}
Let $M^2>0$, and let $F(\varphi;t,s,r)$ be a polynomial in the parameters $t,s,r$ and in the variables $\abs{\varphi(y)}^2$, $y$ in a finite set, with $a$-independent coefficients. Let $\mathcal A$ be a centred Gaussian field on the bonds, independent of $a$, and let $B:=\sum_\square t_\square g_\square\mathcal A$, a finite sum. Define, as a formal power series in $(t,s,r)$,
\[
W_R(a):=\log\E_{\mathcal A}\int_{(\R^2)^{\Lambda_R}}d\varphi\,\exp\Bigl[-\tfrac12\ip{\varphi}{(-\Delta^{\xi,R}_{a+B}+M^2)\varphi}-F(\varphi)\Bigr].
\]
Let $\Pi^{R,\Lambda}_{\mu\nu}(b,b')$ be the second derivative in $a(b),a(b')$ at $a=0$ of the coefficient of the monomial $\Lambda$. Then:
\begin{itemize}
\item[(a)] $\sum_{b'\parallel e_\nu}\Pi^{R,\Lambda}_{\mu\nu}(b,b')=0$ for every $R$, every $b$ and every $\mu,\nu$;
\item[(b)] as $R\to\infty$, $\Pi^{R,\Lambda}(b,b')$ converges to $\Pi^\Lambda(b,b')$, the sum of the connected graphs of label configuration $\Lambda$ with scalar lines $C^\xi_M$ and vector lines the covariance of $\mathcal A$. Moreover, $\sum_{b'}\abs{\Pi^\Lambda(b,b')}<\infty$ and $\sum_{b'\parallel e_\nu}\Pi^\Lambda_{\mu\nu}(b,b')=0$.
\end{itemize}
\end{lemma}
\begin{proof}
(a) Each $\mathrm{e}^{q\chi(y)}$ is a rotation commuting with $\mathrm{e}^{-q\xi a(b)}$, so
\[
\mathrm{e}^{-q\xi(a+d\chi)(b)}\mathrm{e}^{q\chi(y+\xi e_\nu)}\varphi(y+\xi e_\nu)-\mathrm{e}^{q\chi(y)}\varphi(y)=\mathrm{e}^{q\chi(y)}\bigl[\mathrm{e}^{-q\xi a(b)}\varphi(y+\xi e_\nu)-\varphi(y)\bigr].
\]
Bonds leaving $\Lambda_R$ contribute only $\abs{\varphi(y)}^2$. Hence $\ip{\mathrm{e}^{q\chi}\varphi}{-\Delta^{\xi,R}_{a+B+d\chi}\mathrm{e}^{q\chi}\varphi}=\ip{\varphi}{-\Delta^{\xi,R}_{a+B}\varphi}$ for every $B$. Also $F(\mathrm{e}^{q\chi}\varphi)=F(\varphi)$, the map $\varphi\mapsto \mathrm{e}^{q\chi}\varphi$ is orthogonal, and the form is positive definite. So, coefficient by coefficient, $W_R(a+d\chi)=W_R(a)$.

Each coefficient is a smooth function of $a$, being a finite combination of Gaussian moments with covariance $(-\Delta^{\xi,R}_a+M^2)^{-1}$. Differentiate along $\chi$ at $\chi=0$ and in $a(b)$ at $a=0$; this gives the finite sum identity $\sum_{b'}\Pi^{R,\Lambda}(b,b')(d\chi)(b')=0$. Take $\chi(y)=e_\nu\cdot y$, so that $(d\chi)(b')=1$ for $b'\parallel e_\nu$ and $0$ otherwise.

(b) By Wick's theorem and the linked-cluster property, $\Pi^{R,\Lambda}$ is the finite sum of the connected graphs of configuration $\Lambda$ with scalar lines $C^{\xi,R}_M$ and lattice differences of them. Every vertex is labelled, and so confined to a bounded set $K_\Lambda$, or is the vertex at $b$ or at $b'$. If the vertex at $b'$ is unlabelled, it is joined to the rest by a scalar line. By Lemma~\ref{app:G1:AA4-Dir}, each member is bounded, uniformly in $R$, by $c(\xi,\Lambda)\theta^{\dist_1(b',K_\Lambda\cup\{b\})/\xi}$, and converges pointwise as $R\to\infty$. Dominated convergence in $b'$ and (a) conclude.
\end{proof}

\noindent The identity does not apply with $G$ in place of $C^\xi_M$: the term $a_{j_0}P_{j_0}$ is not gauge covariant without the vertices \cite[(1.13)--(1.15)]{Ba3}. The insertion $F$ may contain the fixed mass vertices and the counterterms in graph representation evaluated with $C^\xi$-lines, so the identity holds separately for each monomial in $r$. No decay of $\mathcal C_{\rm ph}$ is used.

\begin{proposition}[exact form of the coefficient of {\cite[(3.32)]{Ba3}}]\label{app:G1:AA4-red}
Let $0\le m^2_{j_0}\le1$. For a label configuration $\Lambda$, write the zeroth-moment term of \cite[(3.31)]{Ba3} for the class of order $(\alpha,\beta)\in\{(4,0),(2,1)\}$ as $e(t_k)^\alpha\lambda(t_k)^\beta\sum_x\eta^3g(x)A_\mu(x)\,c^{(4),\Lambda}_{\mu\mu'}(x)\,g'(x)A'_{\mu'}(x)$. Let $b_\mu(y)$ be the bond at $y$ in direction $\mu$, and put $y=\ell^{-1}x$.
\begin{itemize}
\item[(a)] $c^{(4),\Lambda}_{\mu\mu'}(x)=\hat c^{\,\Lambda}_{\mu\mu'}(y):=\sum_{b'\parallel e_{\mu'}}\xi^3\bigl[\Pi^\Lambda_{G}-\Pi^\Lambda_{C}\bigr]_{\mu\mu'}(b_\mu(y),b')$, and $\sum_{b'\parallel e_{\mu'}}\xi^3\Pi^\Lambda_{C,\mu\mu'}(b,b')=0$.
\item[(b)] $\hat c^{\,\Lambda}_{\mu\mu'}(y)=\sum_T\sum_{\emptyset\ne S}(-1)^{\abs S+1}\sum_{b'\parallel e_{\mu'}}\xi^3T^S_{\mu\mu'}(b_\mu(y),b')$, where $T$ runs over the members of $\Pi^\Lambda_G$ and $S$ over the non-empty sets of scalar lines of $T$, and all sums converge absolutely.
\end{itemize}
\end{proposition}
\begin{proof}
\emph{Rescaling.} Put $x=\ell y$ and $\eta=\ell\xi$. In $d=3$, $\varphi_\eta(x)=\ell^{-1/2}\varphi_\xi(y)$ and $A_\eta(x)=\ell^{-1/2}A_\xi(y)$. Then
\begin{itemize}
\item $\eta^3\lambda\abs{\varphi_\eta}^4=\xi^3(\lambda\ell)\abs{\varphi_\xi}^4$;
\item $\eta^3M^2t_k^2\abs{\varphi_\eta}^2=\xi^3(M^2t_{j_0}^2)\abs{\varphi_\xi}^2$, since $t_k\ell=t_{j_0}$;
\item $e\,\eta A_\eta=(e\ell^{1/2})\xi A_\xi$.
\end{itemize}
So $e_\xi^2=e(t_k)^2\ell$, $\lambda_\xi=\lambda(t_k)\ell$, and the rescaled mass is $m^2_{j_0}$. A local coefficient of $A\cdot A'$ satisfies $c_\eta=\ell^{-2}c_\xi$. Writing $c_\xi=e_\xi^\alpha\lambda_\xi^\beta\hat c$ gives $c_\eta=e(t_k)^\alpha\lambda(t_k)^\beta\ell^{\alpha/2+\beta-2}\hat c$. The exponent is $-1$ for $(2,0)$, which is the factor $(L^{j_0}\eta)^{-d+2}$ of \cite[p.~441]{Ba3}, and $0$ for $(4,0)$ and $(2,1)$.

\emph{Resummation.} By \cite[(3.31) and p.~442]{Ba3}, the coefficient is the zeroth moment of the class, summed over the orderings in which every internal line precedes the external legs, and over the corresponding indices. The transformations performed beforehand on proper divergent subgraphs \cite[pp.~444--445]{Ba3} are exact identities (Taylor's formula with remainder, and resummation of indices), and the orderings in which a given subgraph is the last divergent one partition all orderings. So the zeroth moment is that of the untransformed class. Summing the slice decomposition \cite[(2.6)]{Ba3} over the indices $<j_0$ turns every internal scalar line, including those of the counterterms in graph representation, into $G$. The remaining mass insertions are the fixed mass vertices.

\emph{Cancellation.} In $\Pi^\Lambda_C$, every counterterm in graph representation is an $a$-independent insertion $\kappa(x)\abs{\varphi(x)}^2$, with $\kappa$ polynomial in $t$, and so is every fixed mass vertex. Apply Lemma~\ref{app:G1:AA4-WT} with $M^2=1$, with $\mathcal A$ the centred Gaussian field of covariance $\mathcal C_{\rm ph}$, and with $F$ containing these insertions. It gives $\sum_{b'\parallel e_{\mu'}}\xi^3\Pi^\Lambda_{C,\mu\mu'}(b,b')=0$. This proves (a).

(b) The members of $\Pi^\Lambda_G$ and $\Pi^\Lambda_C$ correspond one to one. Since $C^\xi=G-D$, for each member $\prod_lG_l-\prod_lC^\xi_l=\sum_{\emptyset\ne S}(-1)^{\abs S+1}\prod_{l\in S}D_l\prod_{l\notin S}G_l$. At fixed $\xi$, $G$ and $D$ decay exponentially (Lemma~\ref{app:G1:AA}, Steps~2--3), so the sums converge absolutely, as in the proof of Lemma~\ref{app:G1:AA4-WT}(b).
\end{proof}

\begin{definition}[generalised graphs]\label{app:G1:R-def}
A \emph{generalised graph} $\Gamma$ is a finite connected multigraph $(V,\mathcal L)$ with $\mathcal L\ne\emptyset$ and a root $v_0\in V$; self-loops are allowed. Each line $l$ has endpoints $u_l,w_l$ and a dimension $\omega_l\in\R$, and each vertex $v$ has an exponent $p_v\ge0$. For $\emptyset\ne\mathcal L_H\subset\mathcal L$, let $V_H$ be the set of endpoints of the lines of $\mathcal L_H$, and $H=(V_H,\mathcal L_H)$. The \emph{degree} of $H$ is
\[
D(H):=\sum_{l\in\mathcal L_H}\omega_l+d(\abs{V_H}-1)+\sum_{v\in V_H}p_v.
\]
For $\delta>0$, $N\ge0$, $s_j:=L^{j-N}$ and $\mathbf j=(j_l)\in\{0,\dots,N\}^{\mathcal L}$, put, with $y_{v_0}$ fixed,
\[
E_{\mathbf j}(\Gamma):=\sum_{(y_v)_{v\ne v_0}\in(\xi\Z^d)^{V\setminus v_0}}\prod_{v\ne v_0}\xi^d\prod_v\xi^{p_v}\prod_ls_{j_l}^{\omega_l}\mathrm{e}^{-\delta\abs{y_{u_l}-y_{w_l}}/s_{j_l}},\qquad\xi=L^{-N}.
\]
Let $\Gamma_i(\mathbf j)$ be the subgraph of the lines with $j_l\le i$.
\end{definition}

\noindent To a member $T^S$ we associate the generalised graph $\Gamma_T$: the graph of $T$, rooted at the vertex of $b$, with $p_v$ the number of factors $\xi$ of $v$ besides $\xi^3$. A scalar or vector line with $a_l$ differences, not in $S$, has dimension $-1-a_l$; a line of $S$ has dimension $1-a_l$. For a connected subgraph $H$, let $n_D(H)$ be the number of lines of $S$ in $H$, and let $D^{\rm st}(H)$ be its degree with every line counted at dimension $-1-a_l$. Then $D(H)=D^{\rm st}(H)+2n_D(H)$. $H$ is \emph{critical} if $n_D(H)=0$ and $D^{\rm st}(H)\le0$. A \emph{chain} is a member of the form
\[
\sum_{z,z'}\xi^6\,\Pi^1_{\mu\nu}(b,z)g_1(z)\,\mathcal C_{{\rm ph},\nu\nu'}(z,z')\,g_2(z')\,\Pi^2_{\nu'\mu'}(z',b'),
\]
with $\Pi^1,\Pi^2$ in the order-one class $\mathcal P$ of the bubbles and the seagull tadpole (Lemma~\ref{app:G1:AA}, Step~5).

\begin{lemma}[domination by scale sums]\label{app:G1:R-dom}
Let $\rho=\abs{y-y'}+\xi$ and $0<\delta_1\le\delta'$. Put $f_\omega(\rho)=\rho^\omega$ if $\omega<0$, $f_0(\rho)=1+\log_+(1/\rho)$, and $f_\omega=1$ if $\omega>0$. Then
\[
f_\omega(\rho)\mathrm{e}^{-\delta'\abs{y-y'}}\le \mathrm{e}^{\delta_1}(1+\log L)L^{\abs\omega}\sum_{j=0}^{j_0}s_j^\omega \mathrm{e}^{-\delta_1\abs{y-y'}/s_j}.
\]
In particular, by Lemma~\ref{app:G1:AA}, Step~3, the regular line $D$ with $a\in\{0,1,2\}$ differences, at most one at each end, is dominated by a sum over $j$ of lines of dimension $1-a$, two more than a slice line with $a$ differences.
\end{lemma}
\begin{proof}
If $\rho>1$, then $f_\omega(\rho)\le1$ and the term $j=j_0$ is at least $\mathrm{e}^{-\delta'\abs{y-y'}}$.

Let $\rho\le1$, and let $J=\{j:s_j\ge\rho\}=\{j_*,\dots,j_0\}$. This set is non-empty, since $s_{j_0}=1$, and $s_{j_*}<L\rho$, since $s_0=\xi\le\rho$. For $j\in J$, $\mathrm{e}^{-\delta_1\abs{y-y'}/s_j}\ge \mathrm{e}^{-\delta_1}$.
\begin{itemize}
\item For $\omega>0$, use the term $j_0$.
\item For $\omega<0$, use $s_{j_*}^\omega\ge L^{-\abs\omega}\rho^\omega$.
\item For $\omega=0$, use $\abs J\ge\max(1,\log(1/\rho)/\log L)\ge(1+\log(1/\rho))/(1+\log L)$.
\end{itemize}
\end{proof}

\begin{lemma}[degrees in the class]\label{app:G1:R-deg}
Let $T$ be a connected graph of order $o_T=2$ (the order of $e_\xi$ is $\frac12$ and that of $\lambda_\xi$ is $1$) with two external vector legs. Suppose $T$ is built, in $\xi$-units at zero background, from the vertices
\begin{itemize}
\item $\xi^{n-1}\varphi\cdot q^n\partial\varphi\,a^n$, $n\ge1$ (order $n/2$, one difference);
\item $\xi^{n-2}\varphi\cdot q^n\varphi\,a^n$, $n\ge2$ even (order $n/2$);
\item $\lambda\abs\varphi^4$ (order $1$);
\item $c(y)\abs{\varphi(y)}^2$ with $\abs c\le\bar c$ (order $o_v\ge\frac12$).
\end{itemize}
Suppose also that its lines are proportional to $1_2$ in the charge index. For a connected subgraph $H$, let $E_H$ be its number of external legs and $o_H$ its order.
\begin{itemize}
\item[(a)] If some connected subgraph without external scalar legs has an odd number of external vector legs, then $T^S=0$ for every $S$.
\item[(b)] Otherwise every connected subgraph satisfies $D^{\rm st}(H)\ge o_H+E_H/2-3\ge-1$, and $D^{\rm st}(T)\ge0$.
\item[(c)] Hence $D(H)\ge1$ whenever $n_D(H)\ge1$, and $D(T)\ge2$ if $S\ne\emptyset$.
\end{itemize}
\end{lemma}
\begin{proof}
Let $\deg_v$ be the number of legs at $v$, and $\delta_v$ the number of differences at $v$ acting on lines of $H$. Since $2\abs{\mathcal L_H}=\sum_v\deg_v-E_H$,
\[
D^{\rm st}(H)=\sum_{v\in V_H}c_v+\tfrac{E_H}2-3,\qquad c_v:=3+p_v-\tfrac{\deg_v}2-\delta_v .
\]
For the four vertex types, $c_v=\frac n2+1-\delta_v$, $\frac n2$, $1$ and $2$, respectively. So $c_v\ge o_v$; for the fourth type $c_v=2>o_v$, since $T$ has a vertex with a vector leg besides it.

(a) Each vertex has an even number of scalar legs. The charge indices of $H$ run along closed scalar loops, and a loop through vertices with $n_v$ vector legs carries $\pm\tr q^m$, $m=\sum n_v$, times scalars. Since $\sum_v n_v\equiv E_H\pmod2$, some loop has $m$ odd, and $\tr q^m=0$.

(b) The first inequality follows from $c_v\ge o_v$. We have $E_H\ge1$, and $E_H=1$ is excluded by (a) and the parity of the scalar legs. If $E_H=2$ and $\abs{V_H}\ge2$, then $o_H\ge1$. If $E_H=2$ and $V_H=\{v\}$, then $\deg_v\ge4$, so $o_v\ge1$. If $E_H\ge3$, then $o_H\ge\frac12$. For $T$, every difference acts on internal legs, so $D^{\rm st}(T)\ge o_T-2=0$.

(c) Use $D=D^{\rm st}+2n_D$.
\end{proof}

\begin{lemma}[members without critical subgraphs]\label{app:G1:AA4-mem}
Let $T$ be a member built from the vertices of Lemma~\ref{app:G1:R-deg}, whose internal vector lines are sums of slices with weights in $[0,2]$, as for the photon lines of Lemma~\ref{app:G1:AA4}. Let $S\ne\emptyset$, and suppose $T^S$ has no critical subgraph. Then $\sum_{b'}\xi^3\abs{T^S(b,b')}\le C$. Here $C$ depends only on $d$, $a$, $L$, the block size, $\bar c$, the constants of the slice bounds of the scalar and vector lines, and the graph of $T$; it is uniform in $\xi,j_0,k,\eps,\ell_T$ and $m^2_{j_0}\in[0,1]$.
\end{lemma}
\begin{proof}
If $T$ falls under Lemma~\ref{app:G1:R-deg}(a), then $T^S=0$. Otherwise take absolute values after the charge sums, and bound each vertex factor by its power of $\xi$, using $\abs g\le1$ and $\abs c\le\bar c$. Decompose the lines not in $S$ into slices, which obey $\abs{\partial^aG_{(j)}(y,y')}\le c\,s_j^{-1-a}\mathrm{e}^{-\delta_1\abs{y-y'}/s_j}$ by Theorem~\ref{app:G1:VSR-VS} and Lemma~\ref{app:G1:VSR-V0} (the vector slices with weights in $[0,2]$). Dominate the lines of $S$ by Lemma~\ref{app:G1:R-dom}. This bounds the sum by a constant times $\sum_{\mathbf j}E_{\mathbf j}(\Gamma_T)$, with $\delta=\delta_1$, where the vertex of $b'$ is summed or coincides with the root. By Lemma~\ref{app:G1:R-deg}(c) and the hypothesis, every connected subgraph has positive degree. Theorem~\ref{app:K:L05-1}(a) concludes.
\end{proof}

\begin{lemma}[coefficient bound for the vertex of {\cite[(3.32)]{Ba3}}]\label{app:G1:AA4}
Let $0\le m^2_{j_0}\le1$, and let the internal photon lines be the weighted fluctuation lines $(1+u\delta_j)H_j$ of the slice model (Lemma~\ref{app:G1:CBpos}), with weights at most $2$. For a label configuration $\Lambda$, write the zeroth-moment term of \cite[(3.31)]{Ba3} for the class of order $e^4$ or $e^2\lambda$ as
\[
e(t_k)^{\alpha}\lambda(t_k)^{\beta}\sum_x\eta^d\,g(x)A_\mu(x)\,c^{(4),\Lambda}_{\mu\mu'}(x)\,g'(x)A'_{\mu'}(x),\qquad(\alpha,\beta)\in\{(4,0),(2,1)\}.
\]
Then $\abs{c^{(4),\Lambda}_{\mu\mu'}(x)}\le\Gamma_4\,\ell_{j_0}^{0}$. The constant $\Gamma_4$ depends only on:
\begin{itemize}
\item $d$, on $a$ (through $a_\flat$), and on the block size;
\item the constants of the slice bounds of the scalar lines and of the photon lines, including the H\"older bound of a differenced photon slice, and the bound on the first two differences of the label functions $g,g'$;
\item the bound $\bar c_{\rm fin}$ on the rescaled fixed mass vertices \cite[(1.7)]{Ba3} that occur in the class.
\end{itemize}
It is uniform in $\xi,j_0,k,\eps,\ell_T$, $\Lambda$ and $m^2_{j_0}\in[0,1]$, and it does not depend on $\mu_*$, $u$ or the photon zero mode.
\end{lemma}
\begin{proof}
\emph{(i) Rescaling.} By Proposition~\ref{app:G1:AA4-red}(a), $c^{(4),\Lambda}(x)=\hat c^{\,\Lambda}(y)$ with the factor $\ell_{j_0}^0$. The torus enters as in Lemma~\ref{app:G1:AA}, Step~6, and through the fixed mass vertices. So it suffices to show $\abs{\hat c^{\,\Lambda}}\le\Gamma_4$.

\emph{(ii) Exact form.} By Proposition~\ref{app:G1:AA4-red}(b), $\hat c^{\,\Lambda}$ is the alternating sum of the terms $\sum_{b'}\xi^3T^S$ with $S\ne\emptyset$. A counterterm partner carries the lines of its subgraph, so $S$ selects the same lines in a member and in its partners. Photon lines are not replaced.

\emph{(iii) The pure $C^\xi$ part.} The part in which every scalar line is $C^\xi$ has been subtracted exactly, by Lemma~\ref{app:G1:AA4-WT} at fixed labels with the fixed mass vertices and the counterterms in graph representation as gauge-invariant insertions.

\emph{(iv) The remainder terms.} Every term $T^S$ contains at least one regular line. By Lemma~\ref{app:G1:R-dom}, a regular line has dimension $2$ higher than a slice line, so it raises by $2$ the degree of every subgraph containing it. By Lemma~\ref{app:G1:R-deg}(c), every subgraph containing a regular line has degree $\ge1$, and the whole graph has degree $\ge2$.

If $T^S$ has no critical subgraph, Lemma~\ref{app:G1:AA4-mem} bounds it by $O(1)$ uniformly in $\xi$; this uses only the power-counting estimate of \cite[Props.~2.1--2.2]{Ba3} for generalised graphs (Theorem~\ref{app:K:L05-1}). The critical subgraphs contain no regular line; their lines are slices of $G^\xi_{j_0}(0)$ and of the photon line. They are renormalised by the inner transformations of \cite[pp.~425, 434--445]{Ba3}. The zeroth-moment steps among them, \cite[(3.31), (3.36)--(3.38)]{Ba3}, replace the scalar and the photon lines by $C^\xi$. With the photon lines kept, these steps do not apply literally, and they produce regular lines inside the transformed subgraph. These transformations are made inside the following groupings of the remainder:
\begin{itemize}
\item the chains are grouped over both order-one factors;
\item every other member is grouped with its counterterm partners at fixed $S$.
\end{itemize}
The resulting bounds are Statement~\ref{app:K:L05-3}, and the assembly of all terms into $\abs{\hat c^{\,\Lambda}}\le K_R$ is Theorem~\ref{app:K:L05-2}. The line bounds are uniform in $m^2_{j_0}\in[0,1]$ for the scalar lines, and independent of $\mu_*$ for the photon lines. Members containing a fixed mass vertex are multiplied by its rescaled coefficient, at most $\bar c_{\rm fin}$. Hence $\abs{\hat c^{\,\Lambda}}\le\Gamma_4:=K_R$.

\emph{(v) Higgs-mass dependence.} The Higgs mass $M$ enters only through $m^2_{j_0}=M^2t_{j_0}^2\le m_2^2\eps_0^2\le1$ (by (E.m2)) and through $\bar c_{\rm fin}$. So $\Gamma_4=\Gamma_4(M)$ may depend on $M$ through $\bar c_{\rm fin}$, which is finite at every $M\in(0,m_2]$ by Lemma~\ref{lem:B:massrange}(a). If the class contains no fixed mass vertex, $\Gamma_4$ is uniform in $M\in(0,m_2]$. Uniformity as $M\downarrow0$ is not claimed and is not used; only the finiteness at the four masses is used.
\end{proof}

\paragraph{The tail beyond the order $\nbar$.}\label{app:G1:ssec-T579}
\cite[p.~579]{Ba2} states that, in the expansions of ``the basic quadratic form in the $\psi$, $\phi$ fields and the polynomial $\mathcal P^{(k)}$'', ``the terms with $R$-vertices, the terms containing the configuration $\tilde B^{(k)}$ and the terms of order $>\nbar$'' are $O((L^k\eps)^\kappa)\abs{\Lambda_2^{(k)}}$. The log-determinant expansions are not on this list; they are cut at \cite[pp.~577--578]{Ba2}. We write the order-$>\nbar$ part.

\emph{Setting.} This is the $k$-th step, $0\le k\le K-1$, with $t:=t_k$, on the support of the small-field restrictions $\chi_k$ of \cite[(2.55)]{Ba2}. Here $\abs{A'}\le O(1)p$ on $\Lambda_1^{(k)}$, and the scalar fields and their covariant differences are $\le Cp\lambda(t)^{-1/4}$. Write:
\begin{itemize}
\item $\mathfrak B:=B^k(\Lambda_0^{(k)})\tilde B$ for the regular background of \cite[p.~579]{Ba2};
\item $\tilde A:=\theta_{k+1}A'^{(k)}+\tilde B^{(k)}$ for the expansion field;
\item $\sigma:=e(t)(\norm{\tilde A}_\infty+\norm{\partial^\eta\tilde A}_\infty)$, so that $\sigma\le C_\sigma t^{1/2}p$ with $C_\sigma=eC_A$;
\item $N_k:=\abs T/t_k^3$ for the number of $t_k$-blocks of $T_\eps$, so that $\abs{\Lambda_2^{(k)}}_k\le N_k$.
\end{itemize}
The field $A$ enters only through:
\begin{itemize}
\item $G_k(\Omega,\mathfrak B+\tilde A)$, on $\Omega=B^k(\Lambda_0^{(k)})$ for the basic quadratic form and on $\Omega'=B^{k-1}(\Lambda_2^{(k-1)})$ for the lines of $\mathcal P^{(k)}$, whose vertices lie in $\Omega_1=B^k(\Lambda_7^{(k-1)\prime})$;
\item $Q_k(\mathfrak B+\tilde A)$;
\item the transports and covariant differences in the vertices of $\mathcal P^{(k)}_w$.
\end{itemize}
Every coefficient is an entire function of $e(t)\tilde A$.

\emph{Inputs.} (K1)--(K4), (IH) and (CAP) are stated in Subsection~\ref{app:G1:part-TVSR}; (K3) is proved as Lemma~\ref{app:G1:I3r-U1m}. (REG), ($\tilde B$) and (PO) are bounds of \cite{Ba1,Ba2} that are used as stated.
\begin{itemize}
\item (K1) The bounds of \cite[(1.8), (1.10), (1.12)]{Ba4} for $G_k(\Omega,\mathfrak B)$ on a regular background: decay, first derivative, $\alpha$-H\"older quotient of the derivative, and decay of $G_kQ_k^*$. These hold with $0\le M^2t^2\le c$ in $G_k$ and constants uniform there.
\item (K2) The slice bounds \cite[(2.10)--(2.12)]{Ba3} for the lines of the families of $\mathcal P^{(k)}$.
\item (K3) The one- and two-variable remainder bounds (Proposition~\ref{app:G1:VSR-K3}). For $\mathcal K=[G(\mathfrak B)V]^nG(\mathfrak B+\tilde A)[VG(\mathfrak B)]^{n'}$ with $n+n'=m$, the sup, the derivative and the $\alpha$-H\"older quotient of $\mathcal Kf$ are $\le C(C_\infty\sigma)^m\mathrm{e}^{-\beta\dist(x,\supp f)}\norm f_\infty$. Here $V$ has the right-pushed form $R_0+\sum_\mu R_{1,\mu}D_\mu$ with $\norm R\le c_V\sigma$. On Neumann domains this needs internal points at distance $\ge R_0$ from $\Omega^c$ and $R\ge12\sqrt d$.
\item (K4) The weighted $\ell^2$ bounds of Lemma~\ref{app:G1:VSR-K4}.
\item (REG) $\mathfrak B$ is regular for \cite[Prop.~2.1]{Ba1}, under the upper bounds on $\eps_0$ of Table~\ref{p:tab:A4}, among them (E.1) and (E.ksp3), read with the constant $C_A$.
\item ($\tilde B$) $\norm{\tilde B^{(k)}}_\infty+\norm{\partial^\eta\tilde B^{(k)}}_\infty\le C_\kappa t^\kappa$ for every $\kappa$; and $\theta_{k+1}A'^{(k)}=O(p)$, with its derivative.
\item (IH) Every localised family of $\mathcal P^{(k)}_w$ obeys \cite[(1.33)]{Ba3} with the factor $(\abs uc_k)^m$, $c_k\le1/3$.
\item (CAP) The same bound holds for families of graded order $\le\Nbar=\nbar^2=169$, with $O(1)=O(\alpha_0,\Nbar)$ and decay $\delta(\Nbar)>0$.
\item (PO) An elementary vertex of coupling order $q_v$ in a connected graph with at least two vertices gives $\le Ct^{q_v/4}\cdot$polylog.
\end{itemize}

\begin{lemma}[regularity of the backgrounds]\label{app:G1:T579-reg}
Let $0\le k\le K-1$, $t:=t_k$, $\ell(t):=1+\log t^{-1}$, $p(t)=b_0\ell(t)^{p_0}$ with $p_0=6$, $e(t)=et^{1/2}$, and let $\mathfrak B:=B^k(\Lambda_0^{(k)})\tilde B$ and $\tilde A$ be the background and the expansion field of the setting above. Write $\partial^\eta$ for the lattice derivative of the scale-$k$ lattice, and $\theta(A):=e(t)\norm{\partial^\eta A}_\infty$. Work on the support of the small-field restrictions $\chi_k$.
\begin{enumerate}
\item[(a)] There is $C_{\rm reg}=C_{\rm reg}(d,L,a,M_b,R,c_1,\mu_*)$ such that $\abs{(\partial^\eta_\mu\mathfrak B)(x)}\le C_{\rm reg}p(t)$ and $\abs{(\partial^\eta_\mu(\mathfrak B+\tilde A))(x)}\le C_{\rm reg}p(t)$ at every point $x$ of the domain on which these configurations are defined.
\item[(b)] For every $s\in[0,1]$, $\abs{\partial^\eta_\mu(\mathfrak B+s\tilde A)}\le C_{\rm reg}p(t)$.
\item[(c)] Let $\beta=1/56$ and $c_{\rm reg}$ be the constants of \cite[(2.23), p.~610]{Ba1}, let $\vartheta_R$ be the constant of Lemma~\ref{app:G1:WKinputs}, and let $C_A$ be the constant of the setting above, so that $\sigma\le\abs e\,C_At^{1/2}p(t)$. Suppose that for all $t\in(0,\eps_0]$
\[
\text{(Reg1)}\quad C_{\rm reg}\,\abs e^{1-\beta}t^{(1-\beta)/2}p(t)\le c_{\rm reg},\qquad
\text{(Reg2)}\quad \abs e\max\{C_{\rm reg},C_A\}\,t^{1/2}p(t)\le\min\{\vartheta_R,1\}.
\]
Then every $\mathfrak B+s\tilde A$, $s\in[0,1]$, is regular in the sense of \cite[(2.23)]{Ba1}, and $\theta(\mathfrak B)\le\vartheta_R$, $\theta(\mathfrak B+\tilde A)\le\vartheta_R$, $\sigma\le1$.
\item[(d)] For $a\in(0,1)$ the function $F_a(t):=t^a\ell(t)^{p_0}$ is increasing on $(0,\exp(1-p_0/a)]$, and $F_a(t)\to0$ as $t\downarrow0$. Hence (Reg1) and (Reg2) hold on $(0,\eps_0]$ as soon as they hold at $t=\eps_0$ and $\eps_0\le\exp(1-2p_0/(1-\beta))$. The admissible $\eps_0$ form an interval $(0,\eps_{\rm reg}]$ with $\eps_{\rm reg}>0$, and $\eps_{\rm reg}$ does not depend on the Higgs mass.
\end{enumerate}
\end{lemma}
\begin{proof}
(a) For $\mathfrak B$ and for the configuration (2.93) of \cite[p.~576]{Ba2}, of which $\mathfrak B+\tilde A$ is the case $j=k$, this is \cite[(2.94)--(2.98), pp.~576--577]{Ba2}. There the derivative of each shell term is bounded through \cite[Lemma~2.3, p.~571]{Ba2}, and on the $\theta$-slices through (2.96)--(2.97). The bound is pointwise, so it holds on every subdomain. For $\mathfrak B$ it is stated at \cite[p.~579]{Ba2}.
(b) $\mathfrak B+s\tilde A=(1-s)\mathfrak B+s(\mathfrak B+\tilde A)$, and the bound in (a) is preserved under convex combinations.
(c) By (a), (b) and (Reg1), $\abs{\partial^\eta_\mu(\mathfrak B+s\tilde A)}\le C_{\rm reg}p(t)\le c_{\rm reg}(\abs e t^{1/2})^{\beta-1}=c_{\rm reg}\abs{e(t)}^{\beta-1}$, which is \cite[(2.23)]{Ba1}. By (a) and (Reg2), $\theta(\mathfrak B)\le\abs eC_{\rm reg}t^{1/2}p(t)\le\vartheta_R$, and likewise for $\mathfrak B+\tilde A$. By (Reg2), $\sigma\le\abs eC_At^{1/2}p(t)\le1$.
(d) Since $\frac{d}{dt}\ell(t)=-1/t$, $\frac{d}{dt}\log F_a(t)=\frac1t\big(a-\frac{p_0}{\ell(t)}\big)$, which is positive if and only if $\ell(t)>p_0/a$, that is, $t<\exp(1-p_0/a)$. Also $F_a(t)\to0$ as $t\downarrow0$. (Reg1) is $F_{(1-\beta)/2}\le c_{\rm reg}/(C_{\rm reg}\abs e^{1-\beta}b_0)$ and (Reg2) is $F_{1/2}\le\min\{\vartheta_R,1\}/(\abs e\max\{C_{\rm reg},C_A\}b_0)$. Since $\exp(1-2p_0/(1-\beta))\le\exp(1-2p_0)$, both functions are increasing on $(0,\eps_0]$ when $\eps_0\le\exp(1-2p_0/(1-\beta))$, so their suprema there are their values at $\eps_0$. None of the constants involves the Higgs mass.
\end{proof}

\begin{lemma}[power count per coupling order]\label{app:G1:T579-po}
Let $t=t_k\le\eps_0\le1$, $\hat p:=\max\{1,p(t)\}$, $\lambda(t):=\lambda t$, and for constants $c_\Phi,c_A>0$ put $X_s:=\max\{1,c_\Phi p(t)\lambda(t)^{-1/4}\}$ and $X_v:=\max\{1,c_Ap(t)\}$. Let $F$ be a connected localised family of \cite[Prop.~1]{Ba3} with at least two vertices. Its vertices are those of \cite[(1.6)--(1.11)]{Ba3} and the products \cite[(1.18)]{Ba3}, each product counted as one vertex. The expressions \cite[(1.12)--(1.15)]{Ba3} occur only inside these products, and \cite[(1.12)]{Ba3} is an external scalar field. A vertex $v$ has $e$-order $a_v$, $\lambda$-order $b_v$, $q_v:=a_v+b_v$, and $s_v$ external scalar and $n_v$ external vector legs. We call $v$ \emph{exceptional} if $a_v=1$, $b_v=0$ and $s_v=2$. Put $q:=\sum_vq_v$, and let $d(F)$ be the localisation length of \cite[(1.33)]{Ba3}. Assume:
\begin{itemize}
\item[(PC1)] $\abs{E(F)}\le C_F\,\mathrm{e}^{-\delta d(F)}\prod_v\Phi_v$, where $\Phi_v:=w_vX_s^{s_v}X_v^{n_v}$ with $w_v=e(t)^{a_v}\lambda(t)^{b_v}$ for a vertex that is neither a mass vertex nor exceptional;
\item[(PC2)] every external scalar field has norm \cite[(1.32)]{Ba3} at most $X_s$, and every external vector field at most $X_v$;
\item[(PC3)] for a mass vertex \cite[(1.7)]{Ba3}, the piece of $\delta m^2(\cdot,x)t^2$ of orders $(a,b)$ is bounded by $C_\delta e(t)^a\lambda(t)^b(1+K-k)$, and $\Phi_v$ is this bound times $X_s^{s_v}$;
\item[(X)] for an exceptional vertex, $\Phi_v\le C_{\rm ex}\,t^{1/4}\hat p^{\,2}$.
\end{itemize}
Then
\[
\abs{E(F)}\le C(q)\,C_F\,t^{q/4}\hat p^{\,4q}(1+K-k)^q\mathrm{e}^{-\delta d(F)},\qquad 1+K-k\le1+\frac{\ell(t)}{\log L},
\]
where $C(q)$ depends on $q,e,\lambda,c_\Phi,c_A,C_\delta,C_{\rm ex}$.
\end{lemma}
\begin{proof}
Every vertex has $q_v\ge1$. For a product \cite[(1.18)]{Ba3}, the factor \cite[(1.14)]{Ba3} or \cite[(1.15)]{Ba3} carries $e(t)^{n+n'}$ with $n+n'\ge1$, or $e(t)^{n+\nbar+1}$. For a mass vertex $q_v\ge1$ holds because every piece in (PC3) has orders $(a,b)\ne(0,0)$. Every vector leg carries one factor $e$, so $n_v\le a_v$. Since $F$ is connected and has at least two vertices, every vertex has an internal leg. Since $t\le1$, $X_s\le\max\{1,c_\Phi\lambda^{-1/4}\}\hat p\,t^{-1/4}$ and $X_v\le\max\{1,c_A\}\hat p$. For a vertex that is neither a mass vertex nor exceptional this gives $\Phi_v\le C_v\,t^{a_v/2+b_v-s_v/4}\hat p^{\,s_v+n_v}$. Moreover $a_v/2+b_v-s_v/4\ge q_v/4$ holds if and only if $a_v+3b_v\ge s_v$.
\begin{itemize}
\item For \cite[(1.6)]{Ba3}: $a_v=0$, $b_v=1$, and one of the four scalar legs is internal, so $s_v\le3$.
\item For a mass vertex: $s_v\le2$. A piece of orders $(a,b)$ with $b\ge1$ has $a+3b\ge3$. A piece with $b=0$ has $a\ge2$, because its $e$-order is the number of contracted vector legs, which is even. So $\Phi_v\le C\,t^{q_v/4}\hat p^{\,2}(1+K-k)$.
\item For \cite[(1.8)--(1.11)]{Ba3} and for the products \cite[(1.18)]{Ba3}: $b_v=0$ and $s_v\le2$. If $a_v\ge2$, then $s_v\le a_v$. If $a_v=1$ and $s_v\le1$, the criterion holds. The remaining case is the exceptional one, and (X) gives $t^{1/4}=t^{q_v/4}$.
\end{itemize}
The number of legs is $s_v+n_v\le4q_v$ in every case, so $\Phi_v\le C_v\,t^{q_v/4}\hat p^{\,4q_v}(1+K-k)^{q_v}$. Take the product over $v$. Finally, $K-k=\log_L(\eps_0/t)\le\ell(t)/\log L$.
\end{proof}

\begin{lemma}[the field-independent part]\label{app:G1:T579-vac}
For fixed domains, the vacuum term $-E_1(\cdot,\Omega_1;u)$ of $\mathcal P^{(k)}_w$ \cite[(1.4), (1.30)]{Ba3} is a constant whose coupling order is at most $\nbar$. It belongs to $\mathcal Q_k$ in Proposition~\ref{app:G1:T579}. Every term of $\Delta_k^{\rm pre}$ contains at least one factor that depends on $\tilde A$: an $\tilde A$-leg, a vertex $V^{[r]}$, $V^{>}$ or $V^R$, or a factor $\mathcal K_m$. Consequently, at every $k$, including $k=0$, the proof of Proposition~\ref{app:G1:T579} uses the family bound of Definition~\ref{app:G1:VSR-IH} only for families with at least one external field.
\end{lemma}
\begin{proof}
By \cite[(1.30)--(1.31), p.~419]{Ba3}, $E_1=E_{\rm fin}+E_{K,k}+E_k(\cdot,\Omega_1)$ is a sum of vacuum graphs. Their lines are the zero-background propagators $G_K(0)$, $G_{K,k}(0)$, $G_k(0)$, $C^\eps$, with each vector line replaced by its $u$-version, and their localising function $g_k$ is fixed by $\Omega_1$. So $E_1$ does not depend on $A'$, $\psi$, $\phi$, $\tilde A$ or $\mathfrak B$. It enters the exponent in \cite[(1.4)]{Ba3} additively, and by \cite[(1.5)]{Ba3} only its coefficients of orders $1\le\alpha+\beta\le\nbar$ enter $\mathcal P^{(k)}_w$. The expansion in the proof of Proposition~\ref{app:G1:T579} acts on a sum, not on its exponential, so this constant is never multiplied by another term. It is not expanded, and by the stop rule (E3) it belongs to $\mathcal Q_k$. Every term of $\mathcal P^{(k)}_w$ has order at most $\nbar$. A term of $\Delta_k^{\rm pre}$ has order larger than $\nbar$, so it contains a factor created by (E1) or (E2), and every such factor depends on $\tilde A$.
\end{proof}

\begin{proposition}[the tail beyond the order $\nbar$]\label{app:G1:T579}
Let $M\in(0,m_2]$ and $u\in\{1,-\eta\}$. Let the inputs above hold, with the power count (PO) in the form of Lemma~\ref{app:G1:T579-po}. Let the conditions (Reg1) and (Reg2) of Lemma~\ref{app:G1:T579-reg} hold, and assume the bound of Statement~\ref{app:K:L06-4}. At $k=0$, assume moreover Statement~\ref{app:K:X0}. Pointwise on $\supp\chi_k$, for $0\le k\le K-1$,
\begin{gather*}
(\text{basic quadratic form}+\mathcal P^{(k)}_w)(\mathfrak B+\tilde A)=\mathcal Q_k+\Delta_k^{\rm pre},\\
\abs{\Delta_k^{\rm pre}}\le C_{579}\,t_k^{7/2}\,\mathrm{polylog}(t_k)\,\abs{\Lambda_2^{(k)}}_k\le C'_{579}\,t_k^{16/5}N_k.
\end{gather*}
Here $\mathcal Q_k$ is a polynomial in $\theta_{k+1}A'^{(k)}$, $\psi$, $\phi$ of order $\le\nbar$ and degree $\le D(\nbar)$, and the identity is exact (a two-sided bound). The constant is $C'_{579}=C(\nbar,\Nbar,\alpha_0,d,L,a,M_b,R,r,e,\lambda,\mu_*,b_0,p_0,c_1,\delta_0,\eta)$. It does not depend on $\eps$, $k$ or $\ell_T$: where the final mass counterterm vertices occur, it is evaluated at the torus floor. It depends on $M$ through the $\delta m^2$ vertices of $\mathcal P^{(k)}_w$. These are bounded at the last scale by Proposition~\ref{app:G1:I3b} and at the intermediate scales by a bound free of the mass. It also depends on $M$ through $M^2t^2\in[0,1]$ in (K1) and (K3), with constants uniform there.
\end{proposition}
\begin{proof}
\emph{Step 0 (the configuration and the backgrounds).} For $1\le k\le K-1$, the field $\tilde A$ is defined through the decomposition \cite[(2.86)--(2.89), p.~575]{Ba2} of $\theta_{k+1}A^{(k)}$. For $L=2$ this decomposition holds identically in $A'$ and $B$ by Proposition~\ref{app:K:L06-1}. Its two conditions hold: $r(\eps_0)\ge8M_b+8\ge2M_b$ by (E.10), and $R\ge107>\sqrt3/\log2$.
For $k=0$ the expansion field is the field $A''$ of \cite[(2.19)--(2.22), p.~561]{Ba2}, on the background $\tilde B^{(1)}$. The cut-off $\zeta$ there equals $1$ for $\abs{x-y}\le\frac12r(\eps)$ and $0$ otherwise. Since the lattice spacing is $1$ at $k=0$, its first lattice differences are bounded by $1$ and its second by $2$. The second and third terms of $A''-A'$ have the forms $\theta_1\zeta\,C^{(0)}Q^*\,\delta C_{\Lambda_0}$ and $(1-\zeta)C^{(0)}Q^*$. The bounds ($\tilde B$) of the inputs hold for them by the same argument, with $G_kQ_k^*$ replaced by $C^{(0)}Q^*$.
By Lemma~\ref{app:G1:T579-reg}, every $\mathfrak B+s\tilde A$, $s\in[0,1]$, is regular in the sense of \cite[(2.23)]{Ba1}, so (K1) applies at these backgrounds. Moreover $\theta(\mathfrak B),\theta(\mathfrak B+\tilde A)\le\vartheta_R$ and $\sigma\le1$, which are the smallness hypotheses under which Proposition~\ref{app:G1:VSR-K3} is applied below.

\emph{Step 1 (expansion with a stop rule).}
\begin{itemize}
\item[(E1)] Taylor-expand every exponential $\mathrm{e}^{iz}$ (with $z=qe(t)\eta\tilde A(b)$ or $z=qe(t)\tilde A(\Gamma)$) to order $\nbar$. This splits the bracket $V_k$ of \cite[(3.44)]{Ba1} into three parts: pieces $V^{[r]}$ homogeneous of degree $r\le\nbar$ in $e(t)\tilde A$; polynomial pieces $V^{>}$ of degree $r\in(\nbar,2\nbar]$; and pieces $V^R$ that contain a remainder $R_{\nbar+1}$. The same split is applied to $Q_k(\mathfrak B+\tilde A)$ and to the $\mathcal P$-vertices.
\item[(E2)] Expand each $G':=G_k(\mathfrak B+\tilde A)$ about $G:=G_k(\mathfrak B)$ by the two-sided identity $G'=\sum_{j<m}(GV)^jG+\mathcal K_m$. It follows by applying $G'=G+GVG'$ $n$ times on the left and $G'=G+G'VG$ $n'$ times on the right. Here $\mathcal K_m=[G(\mathfrak B)V]^nG(\mathfrak B+\tilde A)[VG(\mathfrak B)]^{n'}$ with $n+n'=m$, $m:=\max(6,\nbar+1-q_0)$, and $q_0$ is the order already carried by the term.
\item[(E3)] A term is expanded further only while its total coupling order is $\le\nbar$. Every other term goes into $\Delta_k^{\rm pre}$, with its unexpanded factors $\mathcal K_m$, $V^{>}$ and $V^R$.
\end{itemize}
Each $V$-insertion raises the order by at least one. So the procedure stops after at most $\nbar+1$ rounds. The number of terms depends only on $\nbar$ and on the vertex list \cite[(1.6)--(1.11), (1.18)]{Ba3}. The polynomial terms of $\Delta_k^{\rm pre}$ have graded order at most $3\nbar=39$; families in which $\mathcal K_m$ replaces a line have order at most $\nbar+14=27$. Both are within the cap $\Nbar=169$.

\emph{Step 2 (chains of the basic quadratic form).} Consider $\ip{hf}{\mathcal C h'f'}$, where $\mathcal C$ is a chain $GV^{[r_1]}G\cdots V^{[r_n]}G$ or contains $\mathcal K_m$ or $V^{>}$. The proof of (K3) writes $\mathcal K=L_0(VL_1)\cdots(VL_m)$ and uses $\norm{VL}_{\infty,\beta}\le\norm{R_0}\norm L+\sum_\mu\norm{R_{1,\mu}}\norm{D_\mu L}$ together with submultiplicativity. The degree-$r$ Taylor piece of the bond function is again a function of $q$. So by the lattice Leibniz rule $V^{[r]}$, $V^{>}$ and $V^R$ have the same right-pushed form, with $\norm{\mathfrak r^{[r]}}\le c_V(C(d)\sigma)^r$. Hence
\begin{gather*}
\abs{\ip{hf}{\mathcal C h'f'}}\le C_X(C_\infty C_\sigma t^{1/2}p)^q\mathrm{e}^{-\beta\dist(\supp h,\supp h')}\norm f_\infty\norm{f'}_\infty ,\\
\norm f_\infty\le Cp\lambda^{-1/4}t^{-1/4}.
\end{gather*}
A chain of order $q$ is therefore $\le C(q)t^{(q-1)/2}\lambda^{-1/2}\,\mathrm{polylog}\,\mathrm{e}^{-\beta\dist}$, and at $q\ge\nbar+1=14$ this is $\le Ct^{13/2}\,\mathrm{polylog}$.

\emph{Step 3 (families of $\mathcal P^{(k)}_w$).} After (E1)--(E3) the pieces enter Ba{\l}aban's estimates as follows.
\begin{itemize}
\item The expanded external scalar fields are the zeroth-order fields $\Phi_0:=a_kG_k(\Omega,\mathfrak B)Q_k^*(\mathfrak B)\phi$ and $\phi(y)$, and chains $\mathcal Kf$ of order $m\ge1$, with $\norm{h\cdot}_{1,\alpha_0}\le C(C_\infty\sigma)^mp\lambda^{-1/4}t^{-1/4}$ by (K3).
\item The $\tilde A$-legs are external photon fields with $\norm{h'\tilde A}_{1,\alpha_0}\le Cp$. This uses ($\tilde B$), the H\"older clause of (K1) for $a_kG_kQ_k^*$, and the $C^2$ profile of the cut-offs $\theta_{k+1}$ and $\zeta$.
\item A line replaced by $\mathcal K_m$ is an external two-variable field of \cite[(2.5)]{Ba3}, with $\norm{h\mathcal K_mh'}_{1,\alpha_0}\le\mathrm{e}^{5\beta}(1+6c_h)^2C_U(\alpha_0)(C_VC_\sigma t^{1/2}p)^m\mathrm{e}^{-\beta\dist}$ by Proposition~\ref{app:G1:VSR-K3}(b) ($m\ge6$, $\alpha_0=\frac12$, $c_h$ the constant of the cut-offs of Lemma~\ref{app:G1:ramp-ch}; on a Neumann domain under the hypothesis there on $\supp\tilde A$, and with $\supp h\cup\supp h'$ of depth $\ge R_0+3$). No domain difference of a $\mathcal K_m$-operator is needed, because all lines of a given term live on one domain. The two-domain bounds $\delta G$ enter only inside the families, through Definition~\ref{app:G1:VSR-IH} and Proposition~\ref{app:G1:VSR-CAP}.
\end{itemize}
By Lemma~\ref{app:G1:T579-vac}, every family that occurs has at least one external field, at every $k$ including $k=0$. For such families the family bound of Definition~\ref{app:G1:VSR-IH} is available. So each resulting family obeys \cite[(1.33)]{Ba3} at graded order $\le\Nbar$, by Definition~\ref{app:G1:VSR-IH} and Proposition~\ref{app:G1:VSR-CAP}, with the factor $(\abs uc_k)^m\le3^{-m}$. This is (PC1) of Lemma~\ref{app:G1:T579-po}. The norm bounds above give (PC2), and the bounds on the $\delta m^2$ vertices recalled in the statement give (PC3). A family of graded order $q\ge14$ is connected and has at least two vertices. By Lemma~\ref{app:G1:T579-po}, with (X) supplied in Step 3a, it is $\le C(q)t^{q/4}\hat p^{\,4q}(1+K-k)^q\mathrm{e}^{-\delta(\Nbar)d}$, and $(1+K-k)^q\le(1+\ell(t)/\log L)^q$ is a polylog. Each order gains at least $t^{1/4}$ up to polylog, a $V^{[r]}$-vertex through $(C_\sigma t^{1/2}p)^r$ and a $\mathcal K_m$-factor through $t^{m/2}\,\mathrm{polylog}$.

\emph{Step 3a (the exceptional vertices).} Let $v$ be an exceptional vertex in the sense of Lemma~\ref{app:G1:T579-po}, so $v$ has coupling order one and both of its scalar legs are external. Since the family is connected and has at least two vertices, the vector leg of $v$ is an internal leg $A'$. Hence $v$ is \cite[(1.8)]{Ba3} with $n=1$, $n'=0$, or the product of \cite[(1.12)]{Ba3} or \cite[(1.13)]{Ba3} with \cite[(1.14)]{Ba3} at $n=1$, $n'=0$.
The bound (X) does not follow from (PC2). For example, if $\norm\Phi_{1,\alpha_0}=X_s$ and $\abs{D\Phi}\asymp X_s$, the vertex \cite[(1.8)]{Ba3} is of order $e(t)X_s^2\asymp p(t)^2$. (X) rests on the cancellations of \cite[p.~421]{Ba3}, which use \cite[Lemma~2.4]{Ba2} and the identity $\phi(y)\cdot q\phi(y)=0$. We distinguish three cases. Proposition~\ref{app:K:L06-3} and Statement~\ref{app:K:L06-4} are formulated for $1\le k\le K-1$, so they are used in the second and third cases only for these $k$. At $k=0$ those two cases are covered by Statement~\ref{app:K:X0}, which is assumed in Proposition~\ref{app:G1:T579}. The first case applies at every $k$.
\begin{itemize}
\item One scalar leg of $v$ is a chain of order $m\ge1$. By Lemma~\ref{app:K:L06-2}, $v$ together with the chain contributes at most $Ct^{(1+m)/4}p^{2+m}$. This is (X), with the order $m$ of the chain counted in $q$ and $\hat p^{\,2+m}\le\hat p^{\,4(1+m)}$.
\item Both scalar legs of $v$ are $\Phi_0$ or $\phi(y)$, and the vector leg of $v$ is contracted with a bond leg $(g_kA')_{b'}$ of a vertex of type \cite[(1.8)--(1.11)]{Ba3}. Then (X) is Proposition~\ref{app:K:L06-3}.
\item Both scalar legs of $v$ are $\Phi_0$ or $\phi(y)$, and the vector leg of $v$ is contracted with a contour leg of a factor \cite[(1.14)]{Ba3} or \cite[(1.15)]{Ba3}. Then (X) is the assumed bound of Statement~\ref{app:K:L06-4}.
\end{itemize}

\emph{Step 4 (domain hypothesis of (K3)).} On a Neumann domain, (K3) needs:
\begin{itemize}
\item internal points $y\in\supp V\subset(\supp\tilde A)^{+\sqrt3+L^{-k}}$ at distance $\ge R_0$ from the complement, which holds if $\dist(\supp\tilde A,\partial\Omega)\ge R_0+2\sqrt3+1$;
\item external points at distance $\ge R_0$.
\end{itemize}
Every term of $\tilde A$ carries $\theta_{k+1}$, so $\supp\tilde A\subset\supp\theta_{k+1}$, which lies at distance $\ge r(t)-2M_b$ from $\Lambda_0^{(k)c}$. Condition (E.10) gives $r(\eps_0)\ge8M_b+8$, and (R.2) gives $r(t)\ge R\ge3R_0+15$; moreover $r(t)\ge r(\eps_0)$ for $t\le\eps_0$. Hence
\[
r(t)\ge\tfrac14(8M_b+8)+\tfrac34(3R_0+15)=2M_b+\tfrac94R_0+13.25\ge2M_b+R_0+2\sqrt3+1.
\]
This settles the basic quadratic form on $\Omega$. For the lines of $\mathcal P^{(k)}$ on $\Omega'$: $\Lambda_0^{(k)}\subset\Lambda_7^{(k-1)\prime}$, so the internal margin is at least as large. The external points are vertex points in $\Omega_1$. This set lies several $r$-layers inside $\Omega'$, which is the setting of (K3) on Neumann domains. And $R\ge107>12\sqrt3$.

\emph{Step 5 (sum).} For $q\ge14$, $t^{q/4}\hat p^{\,4q}\le t^{7/2}\,\mathrm{polylog}$, and $\sup_{t\in(0,1]}t^{3/10}\big(b_0(1+\ln t^{-1})^{p_0}\big)^{4q}<\infty$. So every term is $\le Ct^{16/5}$, with room $t^{3/10}$; Step 2 has more room. The localisations through a fixed block of $\Lambda_2^{(k)}$, weighted by $\mathrm{e}^{-\min(\beta,\delta(\Nbar))d}$, sum to $S(\Nbar,d)<\infty$. There are finitely many terms, and $\abs{\Lambda_2^{(k)}}_k\le N_k$. No inequality of the type \cite[(3.47)]{Ba1} is used, so the bound is two-sided.
\end{proof}

\begin{remark}[the other two cuts, and the mass range]\label{app:G1:T579-rem}
(i) The terms that contain a Taylor remainder $R_{\nbar+1}$ obey $\le C_Rt^{13/2}\,\mathrm{polylog}\,\abs{\Lambda_2^{(k)}}_k$. Since $\tilde A$ is real, $\abs{R_{\nbar+1}(z)}\le\abs z^{\nbar+1}/(\nbar+1)!$. The $\norm{\cdot}_{1,\alpha_0}$-bound $c_R(Ct^{1/2}p)^{\nbar+1}$ of an $R$-factor also needs $\abs{R''}\abs{Dz}\abs{\Delta z}$ and $\abs{R'}\abs{\Delta Dz}$, where $[\partial\tilde A]_\alpha$ comes from the H\"older clause of (K1).
(ii) The terms that contain $\tilde B^{(k)}$ obey $\le C_{\tilde B}t^{9/2}\,\mathrm{polylog}\,\abs{\Lambda_2^{(k)}}_k$ at $\kappa=31$. The fourth term of $\tilde B^{(k)}$ is $a_{k+1}L^{-2}\theta_{k+1}(1-\zeta^{(k+1)})G^\eta_{k+1}Q^*_{k+1}\Lambda_0B$.
Both follow by the argument of Proposition~\ref{app:G1:T579}, with the bounds just stated in place of Step~2.
(iii) The three cuts use $M$ only as stated in Proposition~\ref{app:G1:T579}, so they hold at every $M\in(0,m_2]$. At $u=0$ the factor $(\abs uc_k)^m$ is at most its value at $u\in\{1,-\eta\}$.
\end{remark}

\subsection{Line bounds, the remainder line and the renormalisation machinery on the torus}\label{app:G1:part-TVSR}
We prove the inputs (VS), (Rem) and (S6) of the bounds on the final c-numbers of Subsection~\ref{app:G1:part-D3}, and state the inputs (K1)--(K4), (IH) and (CAP) of the tail bound.

\paragraph{The zero-background slices and the last remainder line}
\emph{Setting (scale-$j$ units).} Fix a slice index $j$ with $m_ft_j\le1$, and measure lengths in units of $t_j$. Let $\Lambda$ be the $\eta$-lattice torus of side $\ell$, where $\eta:=\eps/t_j=L^{-j}\in(0,1]$, $1/\eta\in\N$, and $\ell:=\ell_T/t_j\in L\N$ with $\ell\ge L$. The inner product is $\ip fg=\eta^3\sum_xf(x)\bar g(x)$, and a kernel acts by $(Kf)(x)=\eta^3\sum_yK(x,y)f(y)$. $d$ is the $\ell^1$ distance on the torus. $-\Delta$ is the periodic nearest-neighbour Laplacian, and $\nabla_\mu$ is the forward difference. For $s\in\{1,L\}$, $P_s$ is the orthogonal projection onto functions constant on blocks of side $s$. Put $\tilde\mu:=m_ft_j\in[0,1]$, $b:=a_j$ and $b':=a_{j+1}L^{-2}$, with $a_k\in[a_\flat,a]$ and $a_\flat:=a(1-L^{-2})$. Define
\[
\mathsf X:=-\Delta+\tilde\mu^2+bP_1,\quad \mathsf X':=-\Delta+\tilde\mu^2+b'P_L,\quad G:=\mathsf X^{-1},\quad G':=\mathsf X'^{-1},\quad H:=G'-G,\quad F:=(-\Delta+1)^{-1}.
\]
In physical units the slice $H_{j,f}=G_{j+1,f}(0)-G_{j,f}(0)$ has kernel $t_j^{-1}H(x/t_j,y/t_j)$ (canonical scaling, $d=3$).

\begin{lemma}[the Feynman photon is componentwise]\label{app:G1:VSR-V0}
On the hypercubic torus, $\norm{dA}^2+\norm{d^*A}^2=\sum_{\mu,\nu}\norm{\partial_\nu A_\mu}^2$. Hence $C_v^\eps=(-\Delta^\eps+\mu_*^2)^{-1}\delta_{\mu\nu}$. At $A=0$, $G_{k,v}(0)$, $H_{j,v}$ and $R_{\bar K,v}$ are each $\delta_{\mu\nu}$ times the scalar object of mass $\mu_*$.
\end{lemma}
\begin{proof}
$\norm{dA}^2=\sum_{\mu\ne\nu}\norm{\partial_\mu A_\nu}^2-\sum_{\mu\ne\nu}\ip{\partial_\mu A_\nu}{\partial_\nu A_\mu}$ and $\norm{d^*A}^2=\sum_\mu\norm{\partial_\mu A_\mu}^2+\sum_{\mu\ne\nu}\ip{\partial_\mu^*A_\mu}{\partial_\nu^*A_\nu}$. Forward differences and their adjoints commute on the torus, so $\ip{\partial_\mu^*A_\mu}{\partial_\nu^*A_\nu}=\ip{\partial_\nu A_\mu}{\partial_\mu A_\nu}$, and the cross terms cancel. Componentwise block averaging at $A=0$ (\cite[(2.10)--(2.17), p.~653]{K1}: ``setting $m^2=\mu_0^2$, $N=d$ and $A=0$'') gives the rest.
\end{proof}

\paragraph{The heat kernel of $\eta\Z^3$.}
For $\eta\in(0,1]$ let $\Delta$ be the nearest-neighbour Laplacian of $\eta\Z^3$, $\Delta f(x)=\eta^{-2}\sum_\mu[f(x+\eta e_\mu)+f(x-\eta e_\mu)-2f(x)]$, and let $\nabla_\mu f(x):=\eta^{-1}(f(x+\eta e_\mu)-f(x))$. Kernels are taken with respect to $\eta^3\sum_y$. For a translation-invariant operator we write its kernel as a function of $w=x-y$, and $\nabla_\mu$ acts on $w$. Put $\abs w_1:=\sum_\mu\abs{w_\mu}$, $\hat k_\mu:=\tfrac2\eta\sin(\eta k_\mu/2)$, $\hat k^2:=\sum_\mu\hat k_\mu^2$, $B_\eta:=[-\pi/\eta,\pi/\eta]^3$ and $a:=\sqrt\pi/4$. For $z\ge0$ and $n\in\Z$ define $I_n(z)$ by $\mathrm{e}^{(z/2)(t+t^{-1})}=\sum_{n\in\Z}I_n(z)t^n$ ($t\ne0$). Then $I_n=I_{-n}\ge0$, and setting $t=1$ gives $\sum_nI_n(z)=\mathrm{e}^z$.

\begin{lemma}[free kernel bounds]\label{app:G1:VSR-F}
On $\eta\Z^3$, $\eta\in(0,1]$, let $F^\infty:=(-\Delta+1)^{-1}$ and $r:=\abs w_1$. Then:
\begin{itemize}
\item[(F0)] $F^\infty(w)\le A_0\mathrm{e}^{-r/4}/(r+\eta)$, with $A_0:=4.22$;
\item[(F1)] $\abs{\nabla_\mu F^\infty(w)}\le A_1\mathrm{e}^{-r/4}/(r+\eta)^2$, with $A_1:=98$;
\item[(F2)] $\abs{\nabla_\mu\nabla_\nu F^\infty(w)}\le A_2\mathrm{e}^{-r/4}/(r+\eta)^3$, with $A_2<\infty$ absolute;
\item[(F3)] $\abs{\nabla_\mu(F^\infty)^2(w)}\le A_3\mathrm{e}^{-r/4}$, with $A_3:=A_0(1+2\mathrm{e}^{1/4})$;
\item[(F4)] for $\kappa\le1/8$, $K_0:=\eta^3\sum_w\mathrm{e}^{\kappa r}F^\infty(w)\le8.4$ and $K_1:=\eta^3\sum_w\mathrm{e}^{\kappa r}\abs{\nabla F^\infty(w)}\le49A_1$, where $\abs{\nabla F^\infty(w)}:=\max_\mu\abs{\nabla_\mu F^\infty(w)}$.
\end{itemize}
Here $(F^\infty)^2=(-\Delta+1)^{-2}$. For $(-\Delta+\rho^2)^{-1}$ with $\rho\in(0,1]$, (F0) holds with $\mathrm{e}^{-\rho r/4}$ and the same $A_0$, and $\eta^3\sum_w(-\Delta+\rho^2)^{-1}(w)=\rho^{-2}$.
\end{lemma}

The proof rests on two lemmas.

\begin{lemma}[the heat kernel]\label{app:G1:VSR-HK}
Let $p_s$ ($s\ge0$) be the kernel of $\mathrm{e}^{s\Delta}$ on $\eta\Z^3$.
\begin{itemize}
\item[(a)] $p_s(w)=\int_{B_\eta}\mathrm{e}^{ik\cdot w}\mathrm{e}^{-s\hat k^2}\,\frac{d^3k}{(2\pi)^3}=\prod_{\mu=1}^3q_s(w_\mu)$, where $q_s(n\eta)=\eta^{-1}\mathrm{e}^{-z}I_n(z)$ and $z:=2s/\eta^2$. In particular $p_s\ge0$.
\item[(b)] For every $\lambda\in\R^3$, $p_s(w)\le\mathrm{e}^{-\lambda\cdot w}\exp\big(\sum_\mu\tfrac{2s}{\eta^2}(\cosh\eta\lambda_\mu-1)\big)\,p_s(0)$.
\item[(c)] $p_s(0)\le\min(\eta^{-1},a s^{-1/2})^3$.
\item[(d)] For every $\bar\lambda\ge0$ and every $\lambda\in\R^3$ with $\abs{\lambda_\nu}\le\bar\lambda$ for all $\nu$,
\[
\abs{\nabla_\mu p_s(w)}\le\mathrm{e}^{-\lambda\cdot w}\exp\Big(\sum_\nu\tfrac{2s}{\eta^2}(\cosh\eta\lambda_\nu-1)\Big)\mathrm{e}^{\eta\bar\lambda}\int_{B_\eta}\big(\abs{\hat k_\mu}+\bar\lambda\big)\mathrm{e}^{-s\hat k^2}\,\frac{d^3k}{(2\pi)^3}.
\]
\item[(e)] For $s>0$, $\abs{\nabla_\mu p_s(w)}\le\frac{\abs w_1+\eta}{s}\max\big(p_s(w),p_s(w+\eta e_\mu)\big)$.
\end{itemize}
\end{lemma}
\begin{proof}
(a) In one dimension, $\mathrm{e}^{s\Delta}=\mathrm{e}^{-z}\exp\big(\tfrac z2(S+S^{-1})\big)$, where $S$ is the unit shift and $z=2s/\eta^2$. The coefficient of $S^n$ is $\mathrm{e}^{-z}I_n(z)$, and the factor $\eta^{-1}$ comes from the normalisation of kernels. The Fourier formula is the same identity, because $\tfrac{2}{\eta^2}(1-\cos\eta k)=\hat k^2$. The three-dimensional kernel is the product of the one-dimensional ones.

(b) The integrand $k\mapsto\mathrm{e}^{ik\cdot w}\mathrm{e}^{-s\hat k^2}$ is entire and $2\pi/\eta$-periodic in each $k_\mu$. So the shift $k_\mu\to k_\mu+i\lambda_\mu$ does not change the integral: by periodicity, the contributions of the vertical sides cancel. Moreover $\mathrm{e}^{i(k+i\lambda)\cdot w}=\mathrm{e}^{ik\cdot w}\mathrm{e}^{-\lambda\cdot w}$ and
\[
\rRe\big[1-\cos\eta(k_\mu+i\lambda_\mu)\big]=\cosh(\eta\lambda_\mu)(1-\cos\eta k_\mu)-(\cosh\eta\lambda_\mu-1)\ge(1-\cos\eta k_\mu)-(\cosh\eta\lambda_\mu-1).
\]
Hence the modulus of the shifted integrand is at most $\mathrm{e}^{-\lambda\cdot w}\exp(\sum_\mu\frac{2s}{\eta^2}(\cosh\eta\lambda_\mu-1))\mathrm{e}^{-s\hat k^2}$. Integrating gives (b), since $p_s(0)=\int_{B_\eta}\mathrm{e}^{-s\hat k^2}\frac{d^3k}{(2\pi)^3}$.

(c) We have $q_s(0)\le\eta^{-1}$, because $\eta\sum_nq_s(n\eta)=\mathrm{e}^{-z}\sum_nI_n(z)=1$ and $q_s\ge0$. On $B_\eta$, the inequality $\abs{\sin x}\ge2\abs x/\pi$ for $\abs x\le\pi/2$ gives $\hat k_\mu^2\ge4k_\mu^2/\pi^2$. Hence
\[
q_s(0)\le\tfrac1{2\pi}\int_\R \mathrm{e}^{-4sk^2/\pi^2}dk=\tfrac{\sqrt\pi}{4}s^{-1/2}.
\]

(d) We have $\nabla_\mu p_s(w)=\int_{B_\eta}\mathrm{e}^{ik\cdot w}m(k_\mu)\mathrm{e}^{-s\hat k^2}\frac{d^3k}{(2\pi)^3}$ with $m(k)=(\mathrm{e}^{i\eta k}-1)/\eta$. This integrand is again entire and periodic, so we shift as in (b). Then
\[
\abs{m(k+i\lambda_\mu)}\le\mathrm{e}^{-\eta\lambda_\mu}\abs{\mathrm{e}^{i\eta k}-1}/\eta+\abs{\mathrm{e}^{-\eta\lambda_\mu}-1}/\eta\le\mathrm{e}^{\eta\bar\lambda}(\abs{\hat k_\mu}+\bar\lambda),
\]
using $\abs{\mathrm{e}^{i\eta k}-1}=\eta\abs{\hat k_\mu}$ and $\abs{\mathrm{e}^{-y}-1}\le\abs y\mathrm{e}^{\abs y}$.

(e) First we show $I_n\ge I_{n+1}$ for $n\ge0$. Differentiating the generating function in $z$ gives $I_n'=\frac12(I_{n-1}+I_{n+1})$. Put $d_n:=I_n-I_{n+1}$ for $n\ge0$. Then $d_{-1}=I_1-I_0=-d_0$, so $d=(d_n)_{n\ge0}$ solves $d'=Ad$ with $d(0)=e_0$. Here $(Ad)_n=\frac12(d_{n-1}+d_{n+1})$ for $n\ge1$ and $(Ad)_0=\frac12d_1-\frac12d_0$. The operator $A$ is bounded on $\ell^\infty(\{0,1,2,\dots\})$, and $A+\frac12$ has nonnegative entries. So $\tilde d(z):=\mathrm{e}^{-z/2}\mathrm{e}^{z(A+1/2)}e_0$ has nonnegative entries and solves the same equation. Both $d$ and $\tilde d$ satisfy $f(z)=e_0+\int_0^zAf$ coordinatewise, and both are bounded in $\ell^\infty$ on compact $z$-intervals ($0\le I_n\le\mathrm{e}^z$). Gronwall's inequality applied to $\sup_n\abs{d_n-\tilde d_n}$ then gives $d=\tilde d\ge0$.

Second, differentiating the generating function in $t$ gives $\frac z2(I_{n-1}-I_{n+1})=nI_n$. Write $q(n):=q_s(n\eta)$. For $n\ge1$,
\[
0\le q(n)-q(n+1)\le q(n-1)-q(n+1)=\tfrac{2n}{z}q(n)=\tfrac{n\eta^2}{s}q(n).
\]
For $n=0$, $0\le q(0)-q(1)\le q(0)-q(2)=\frac{\eta^2}{s}q(1)$. Negative $n$ follow from the symmetry $q(n)=q(-n)$. Together, $\abs{q(n+1)-q(n)}\le\frac{(\abs n+1)\eta^2}{s}\max(q(n),q(n+1))$, that is, $\abs{\nabla q_s(x)}\le\frac{\abs x+\eta}{s}\max(q_s(x),q_s(x+\eta))$. Multiply by the other two factors of (a), and use $\abs{w_\mu}+\eta\le\abs w_1+\eta$.
\end{proof}

\begin{lemma}[lattice sums]\label{app:G1:VSR-Jsum}
For $\phi>0$, $\eta\in(0,1]$ and $p\in\{0,1,2\}$ put $J_p(\phi):=\eta^3\sum_{v\in\eta\Z^3}\mathrm{e}^{-\phi\abs v_1}(\abs v_1+\eta)^{-p}$. Then $J_0(\phi)\le(1+2/\phi)^3$, $J_1(\phi)\le1+6/\phi^2$ and $J_2(\phi)\le1+6/\phi$.
\end{lemma}
\begin{proof}
Put $x:=\mathrm{e}^{-\phi\eta}$ and $u:=\phi\eta$.
\begin{itemize}
\item $J_0(\phi)=\big(\eta\frac{1+x}{1-x}\big)^3=\big(\eta+\frac{2\eta}{\mathrm{e}^{u}-1}\big)^3\le(1+2/\phi)^3$, since $\mathrm{e}^u-1\ge u$.
\item For $n\ge1$, the number of $v\in\eta\Z^3$ with $\abs v_1=n\eta$ is $4n^2+2\le6n^2$. Hence
\[
J_1(\phi)\le\eta^2+6\eta^2\sum_{n\ge1}n\mathrm{e}^{-un}=\eta^2+\frac{6}{\phi^2}\frac{u^2\mathrm{e}^{-u}}{(1-\mathrm{e}^{-u})^2}\le1+\frac6{\phi^2},
\]
because $u\mathrm{e}^{-u/2}/(1-\mathrm{e}^{-u})=u/(2\sinh(u/2))\le1$.
\item $J_2(\phi)\le\eta+6\eta\sum_{n\ge1}\mathrm{e}^{-un}=\eta+\frac{6\eta}{\mathrm{e}^u-1}\le1+\frac6\phi$.
\end{itemize}
\end{proof}

\begin{proof}[Proof of Lemma~\ref{app:G1:VSR-F}]
We have $F^\infty=\int_0^\infty \mathrm{e}^{-s}p_s\,ds$, $(F^\infty)^2=\int_0^\infty s\mathrm{e}^{-s}p_s\,ds$ and $(-\Delta+\rho^2)^{-1}=\int_0^\infty \mathrm{e}^{-\rho^2s}p_s\,ds$. By Lemma~\ref{app:G1:VSR-HK}(a), all three kernels are $\ge0$. We use $\cosh y-1\le\frac{y^2}2\cosh y$, which holds term by term in the power series because $(2k)!\ge2(2k-2)!$. For $\abs y\le\frac12$ it gives $\cosh y-1\le\frac12\cosh(\frac12)\,y^2\le0.5639\,y^2$.

\emph{The choice of $\lambda$.} Put $\lambda_\mu:=\mathrm{sgn}(w_\mu)\bar\lambda$, with $\mathrm{sgn}(0):=1$ and $\bar\lambda:=\frac14+\frac1{4\max(\sqrt s,\eta)}$. Then $\lambda\cdot w=\bar\lambda r$ and $\eta\bar\lambda\le\frac12$. Also
\[
\sum_\mu\tfrac{2s}{\eta^2}(\cosh\eta\bar\lambda-1)\le3.384\,s\bar\lambda^2\le3.384\,s\Big(\tfrac18+\tfrac1{8\max(s,\eta^2)}\Big)\le0.423\,s+0.423 .
\]
By Lemma~\ref{app:G1:VSR-HK}(b), and since $\mathrm{e}^{0.423}\le1.527$,
\begin{equation}\label{app:G1:eq-HKconj}
p_s(w)\le1.527\,\mathrm{e}^{0.423s}p_s(0)\,\mathrm{e}^{-r/4}\mathrm{e}^{-r/(4\max(\sqrt s,\eta))}.
\end{equation}

(F0) On $s\le\eta^2$ we use $p_s(0)\le\eta^{-3}$ and $\mathrm{e}^{-s}\mathrm{e}^{0.423s}\le1$. The contribution is at most $1.527\,\eta^{-1}\mathrm{e}^{-r/(4\eta)}\mathrm{e}^{-r/4}$. With $u=r/\eta$ we have $\eta^{-1}\mathrm{e}^{-u/4}=(r+\eta)^{-1}(1+u)\mathrm{e}^{-u/4}$, and $\max_{u\ge0}(1+u)\mathrm{e}^{-u/4}=4\mathrm{e}^{-3/4}\le1.890$.

On $s\ge\eta^2$ we use $p_s(0)\le a^3s^{-3/2}$ and substitute $u=s^{-1/2}$:
\[
\int_{\eta^2}^\infty a^3s^{-3/2}\mathrm{e}^{-r/(4\sqrt s)}ds=2a^3\int_0^{1/\eta}\mathrm{e}^{-ru/4}du\le2a^3\min(4/r,1/\eta)\le10a^3/(r+\eta).
\]
For the last step, treat $r\le4\eta$ and $r\ge4\eta$ separately. Since $10a^3\le0.871$, the total is at most $1.527(1.890+0.871)\mathrm{e}^{-r/4}/(r+\eta)\le4.22\,\mathrm{e}^{-r/4}/(r+\eta)$.

For the mass $\rho\in(0,1]$ take $\bar\lambda:=\frac\rho4+\frac1{4\max(\sqrt s,\eta)}$. Then still $\eta\bar\lambda\le\frac12$, the exponent is at most $0.423\rho^2s+0.423$, and $\mathrm{e}^{-\rho^2s}\mathrm{e}^{0.423\rho^2s}\le1$. The same two integrals give $A_0\mathrm{e}^{-\rho r/4}/(r+\eta)$. In particular the kernel is summable. Hence $\eta^3\sum_w(-\Delta+\rho^2)^{-1}(w)$ equals its Fourier transform $(\rho^2+\hat k^2)^{-1}$ at $k=0$, that is, $\rho^{-2}$.

(F1) \emph{The range $s\ge\eta^2$.} Apply Lemma~\ref{app:G1:VSR-HK}(d) with the $\lambda$ above; here $\mathrm{e}^{\eta\bar\lambda}\le\mathrm{e}^{1/2}$. By Cauchy--Schwarz in $d^3k/(2\pi)^3$,
\[
\int_{B_\eta}\abs{\hat k_\mu}\mathrm{e}^{-s\hat k^2}\le\Big(\int_{B_\eta}\hat k_\mu^2\mathrm{e}^{-s\hat k^2}\Big)^{1/2}p_s(0)^{1/2}.
\]
In direction $\mu$ use $\hat k_\mu^2\le k_\mu^2$ and $\mathrm{e}^{-s\hat k_\mu^2}\le\mathrm{e}^{-4sk_\mu^2/\pi^2}$; in the other two directions use Lemma~\ref{app:G1:VSR-HK}(c). This gives $\int\hat k_\mu^2\mathrm{e}^{-s\hat k^2}\le\frac{\pi^{5/2}}{32}s^{-3/2}\cdot a^2s^{-1}$, hence $\int\abs{\hat k_\mu}\mathrm{e}^{-s\hat k^2}\le0.0967\,s^{-2}$. Also $\bar\lambda p_s(0)\le\frac{a^3}4(s^{-3/2}+s^{-2})$, with $a^3/4\le0.0218$. So
\[
\abs{\nabla_\mu p_s(w)}\le1.527\,\mathrm{e}^{0.423s}\mathrm{e}^{-r/4}\mathrm{e}^{-r/(4\sqrt s)}\big[0.196\,s^{-2}+0.036\,s^{-3/2}\big].
\]
Multiply by $\mathrm{e}^{-s}$ and use $\sqrt s\,\mathrm{e}^{-0.577s}\le0.565$. Then substitute $u=s^{-1/2}$:
\[
\int_{\eta^2}^\infty s^{-2}\mathrm{e}^{-r/(4\sqrt s)}ds=2\int_0^{1/\eta}u\mathrm{e}^{-ru/4}du\le2\min(16/r^2,1/(2\eta^2))\le44.4/(r+\eta)^2.
\]
Here we used that the maximum over $r$ of $\min(16/r^2,1/(2\eta^2))(r+\eta)^2$ is attained at $r^2=32\eta^2$, where it equals $(\sqrt{32}+1)^2/2\le22.16$. The range $s\ge\eta^2$ therefore contributes at most $14.7\,\mathrm{e}^{-r/4}/(r+\eta)^2$.

\emph{The range $s\le\eta^2$.} Use $\abs{\nabla_\mu p_s(w)}\le\eta^{-1}(p_s(w)+p_s(w+\eta e_\mu))$ and apply \eqref{app:G1:eq-HKconj} at both points, with $\bar\lambda=\frac14+\frac1{4\eta}$. Since $\abs{w+\eta e_\mu}_1\ge r-\eta$, the second point costs at most a factor $\mathrm{e}^{\eta\bar\lambda}\le\mathrm{e}^{1/2}$. With $p_s(0)\le\eta^{-3}$ and $\int_0^{\eta^2}ds=\eta^2$, the contribution is at most $1.527(1+\mathrm{e}^{1/2})\eta^{-2}\mathrm{e}^{-r/(4\eta)}\mathrm{e}^{-r/4}$. Finally $\eta^{-2}\mathrm{e}^{-u/4}=(r+\eta)^{-2}(1+u)^2\mathrm{e}^{-u/4}$ with $u=r/\eta$, and $\max_u(1+u)^2\mathrm{e}^{-u/4}=64\mathrm{e}^{-7/4}\le11.13$. This gives at most $45.1\,\mathrm{e}^{-r/4}/(r+\eta)^2$. Since $14.7+45.1\le98$, the value $A_1=98$ suffices.

(F2) $\nabla_\mu\nabla_\nu p_s(w)$ is the integral of $\mathrm{e}^{ik\cdot w}m(k_\mu)m(k_\nu)\mathrm{e}^{-s\hat k^2}$, again entire and periodic. After the shift of Lemma~\ref{app:G1:VSR-HK}(b), the product of the two multipliers is bounded by $\mathrm{e}^{2\eta\bar\lambda}(\abs{\hat k_\mu}+\bar\lambda)(\abs{\hat k_\nu}+\bar\lambda)$, with $\mathrm{e}^{2\eta\bar\lambda}\le\mathrm{e}$.
\begin{itemize}
\item By Cauchy--Schwarz and the one-dimensional Gaussian bounds used in (F1), for $s\ge\eta^2$ we have $\int\abs{\hat k_\mu\hat k_\nu}\mathrm{e}^{-s\hat k^2}\le Cs^{-5/2}$, $\bar\lambda\int\abs{\hat k_\mu}\mathrm{e}^{-s\hat k^2}\le C(s^{-2}+s^{-5/2})$ and $\bar\lambda^2p_s(0)\le C(s^{-3/2}+s^{-5/2})$. After multiplication by $\mathrm{e}^{-0.577s}$, each of these is at most $C's^{-5/2}$.
\item $\int_{\eta^2}^\infty s^{-5/2}\mathrm{e}^{-r/(4\sqrt s)}ds=2\int_0^{1/\eta}u^2\mathrm{e}^{-ru/4}du\le2\min(128/r^3,1/(3\eta^3))\le C/(r+\eta)^3$.
\item On $s\le\eta^2$, the second difference is $\eta^{-2}$ times a signed sum of four values of $p_s$, at points of $\ell^1$-norm $\ge r-2\eta$. By \eqref{app:G1:eq-HKconj}, $p_s(0)\le\eta^{-3}$ and $\int_0^{\eta^2}ds=\eta^2$, this gives at most $C\eta^{-3}\mathrm{e}^{-r/(4\eta)}\mathrm{e}^{-r/4}\le C'\mathrm{e}^{-r/4}/(r+\eta)^3$.
\end{itemize}
All constants here are absolute.

(F3) Here a cancellation is used: the factor $s$ of $(F^\infty)^2=\int_0^\infty s\mathrm{e}^{-s}p_s\,ds$ cancels the factor $s^{-1}$ of Lemma~\ref{app:G1:VSR-HK}(e). Thus
\[
\abs{\nabla_\mu(F^\infty)^2(w)}\le\int_0^\infty s\mathrm{e}^{-s}\tfrac{r+\eta}{s}\big(p_s(w)+p_s(w+\eta e_\mu)\big)ds=(r+\eta)\big(F^\infty(w)+F^\infty(w+\eta e_\mu)\big).
\]
Let $r':=\abs{w+\eta e_\mu}_1$. Then $r'\ge r-\eta$ and $r'+\eta\ge\max(r,\eta)\ge(r+\eta)/2$. So (F0) gives $(r+\eta)F^\infty(w)\le A_0\mathrm{e}^{-r/4}$ and $(r+\eta)F^\infty(w+\eta e_\mu)\le2\mathrm{e}^{1/4}A_0\mathrm{e}^{-r/4}$. Hence $A_3=A_0(1+2\mathrm{e}^{1/4})$ works.

(F4) We have $\mathrm{e}^{\kappa r}=\prod_\mu \mathrm{e}^{\kappa\abs{w_\mu}}\le\sum_{\sigma\in\{\pm1\}^3}\mathrm{e}^{\kappa\sigma\cdot w}$. For $\lambda\in\C^3$ with $\abs{\rRe\lambda_\mu}<\frac14$, the series $\Phi(\lambda):=\eta^3\sum_w\mathrm{e}^{\lambda\cdot w}F^\infty(w)$ converges absolutely and locally uniformly by (F0), so $\Phi$ is analytic there. Put $D(\lambda):=1-\sum_\mu\frac2{\eta^2}(\cosh\eta\lambda_\mu-1)$. For $\lambda=-ik$ with $k\in\R^3$, $\Phi(\lambda)=(1+\hat k^2)^{-1}$ and $D(\lambda)=1+\hat k^2$. Thus the analytic function $\Phi D-1$ vanishes on $i\R^3$, and so it vanishes identically (argue one variable at a time). For real $\lambda$ with $\abs{\lambda_\mu}\le\kappa\le\frac18$, the bound $\frac2{\eta^2}(\cosh\eta\kappa-1)\le\kappa^2\cosh\kappa$ gives $D(\lambda)\ge1-3\kappa^2\cosh\kappa>0$. Since $F^\infty\ge0$, this yields $K_0\le8/(1-3\kappa^2\cosh\kappa)\le8.4$. For $K_1$, (F1) gives $\mathrm{e}^{\kappa r}\abs{\nabla F^\infty(w)}\le A_1\mathrm{e}^{-(1/4-\kappa)r}/(r+\eta)^2$. Lemma~\ref{app:G1:VSR-Jsum} with $\frac14-\kappa\ge\frac18$ then gives $K_1\le A_1J_2(\frac14-\kappa)\le A_1(1+48)=49A_1$.
\end{proof}

\emph{Lattice sums.} For $\phi>0$ and $\eta\le1$: $J_0(\phi):=\eta^3\sum_{v\in\eta\Z^3}\mathrm{e}^{-\phi\abs v_1}\le(1+2/\phi)^3$, $J_1(\phi):=\eta^3\sum \mathrm{e}^{-\phi\abs v_1}/(\abs v_1+\eta)\le1+6/\phi^2$, and $J_2(\phi):=\eta^3\sum \mathrm{e}^{-\phi\abs v_1}/(\abs v_1+\eta)^2\le1+6/\phi$.

\begin{lemma}[torus images]\label{app:G1:VSR-IM}
Let $\abs{K^\infty(w)}\le A\mathrm{e}^{-\kappa\abs w_1}/(\abs w_1+\eta)^p$ with $p\in\{0,1,2,3\}$, and let $K^T(z):=\sum_{n\in\Z^3}K^\infty(z+n\ell)$ with $\ell\ge1$. Then
\[
\abs{K^T(z)}\le A\big(1+5^pS(\kappa\ell)\big)\mathrm{e}^{-(\kappa/2)d(z)}/(d(z)+\eta)^p,\qquad S(u):=\big(3+2/(\mathrm{e}^{u/2}-1)\big)^3,
\]
and $\eta^3\sum_{z\in\Lambda}\mathrm{e}^{\kappa'd(z)}\abs{K^T(z)}\le\eta^3\sum_{w\in\eta\Z^3}\mathrm{e}^{\kappa'\abs w_1}\abs{K^\infty(w)}$ for every $\kappa'\ge0$.
\end{lemma}
\begin{proof}
Take $z$ in the fundamental domain. For $n\ne0$, $\abs{z+n\ell}_1\ge\max(d(z),\ell/2)$, so $(\abs{z+n\ell}_1+\eta)^{-p}\le5^p(d(z)+\eta)^{-p}$. Moreover $\mathrm{e}^{-\kappa\abs{z+n\ell}_1}\le \mathrm{e}^{-(\kappa/2)d(z)}\mathrm{e}^{-(\kappa/2)\abs{z+n\ell}_1}$, and $\sum_n\mathrm{e}^{-(\kappa/2)\abs{z+n\ell}_1}\le S(\kappa\ell)$. For the $\ell^1$ statement, unfold $\sum_{z\in\Lambda}\sum_n=\sum_{w\in\eta\Z^3}$ and use $d(z)\le\abs{z+n\ell}_1$.
\end{proof}
The torus kernel of $(-\Delta+1)^{-1}$ is $F=F^T$. Put $S_L:=S(L/4)$, so that $S_L\le1013$ at $L=2$ (the value is $1012.54$). Put $A_p^T:=A_p(1+5^pS_L)$ for $p=0,1,2$, $A_3^T:=A_3(1+S_L)$ and $\phi:=1/8$. For $\kappa\le1/16$ the column norm is $c_F:=\sup_x\norm{\mathrm{e}^{\kappa d(x,\cdot)}F(x,\cdot)}_{\ell^2}\le7A_0^T$.

\begin{lemma}[spectral gap and Combes--Thomas]\label{app:G1:VSR-CT}
(i) For $s\in\{1,L\}$, $\beta>0$ and $\tilde\mu\ge0$: $-\Delta+\tilde\mu^2+\beta P_s\ge\min(4/s^2,\beta)$.
(ii) Let $\rho$ be $1$-Lipschitz for $d$, let $0\le\kappa\le\frac12$, and let $Y:=-\Delta+\tilde\mu^2+\beta P_s\ge\gamma$. If $3.4\kappa^2+\beta(\mathrm{e}^{3s\kappa}-1)\le\gamma/2$, then $\norm{\mathrm{e}^{\kappa\rho}Y^{-1}\mathrm{e}^{-\kappa\rho}}\le2/\gamma$.
\end{lemma}
\begin{proof}
(i) The torus Dirichlet form dominates the sum of the Neumann forms of the $s$-blocks. A block with $s/\eta$ points per side has Neumann gap $(4/\eta^2)\sin^2(\pi\eta/(2s))\ge4/s^2$. Hence $\ip f{Yf}\ge(4/s^2)\norm{(1-P_s)f}^2+\beta\norm{P_sf}^2$.
(ii) $\mathrm{e}^{\kappa\rho}(-\Delta)\mathrm{e}^{-\kappa\rho}+\Delta$ splits into a real antisymmetric $\sinh$ part, which has zero real part, and a symmetric $\cosh-1$ part, of norm $\le3.4\kappa^2$. The conjugated $P_s$ differs from $P_s$ by a kernel of Schur norm $\le \mathrm{e}^{3s\kappa}-1$. So $\rRe\ip f{\mathrm{e}^{\kappa\rho}Y\mathrm{e}^{-\kappa\rho}f}\ge(\gamma/2)\norm f^2$.
\end{proof}
\emph{Rates.} Put $\gamma_0:=\min(4,a_\flat)$, and let $c=c(a,L)\in(0,\frac12]$ be the largest number with $3.4c^2+a(\mathrm{e}^{3c}-1)\le\gamma_0/2$; at $a=1$, $L=2$ this gives $c=0.0981$. For $\mathsf X$, (ii) holds with $\kappa\le c$; for $\mathsf X'$ (gap $\ge\gamma_0L^{-2}$) with $\kappa\le c/L$. With $W:=(1-\tilde\mu^2)-bP_1$ and $W':=(1-\tilde\mu^2)-b'P_L$, the conjugated norms are $\le w_1:=1+a\mathrm{e}^{3c}$. Put $\kappa_G:=\min(c,1/16)$, $\kappa':=\min(c/L,1/16)$ and $\kappa_H:=\min(c/L,1/32)$.

\begin{lemma}[structure of $G$ and $G'$]\label{app:G1:VSR-E}
Put $E:=G-F$ and $E':=G'-F$. For all $x,y,z\in\Lambda$:
\begin{itemize}
\item[(E1)] $\abs{E(x,y)}\le A_E\mathrm{e}^{-\kappa_Gd(x,y)}$ with $A_E:=c_F^2w_1(1+2w_1/\gamma_0)$, and likewise for $E'$ with $(A_E',\kappa')$, where $A_E':=c_F^2w_1(1+2w_1L^2/\gamma_0)$;
\item[(E2)] $\abs{(P_1G)(z,y)}\le A_P\mathrm{e}^{-\kappa_Gd(z,y)}$ and $\abs{(P_LG)(z,y)}\le A_P\mathrm{e}^{3L\kappa_G}\mathrm{e}^{-\kappa_Gd(z,y)}$, with $A_P:=\mathrm{e}^{3\kappa_G}(J_1(\phi/2)A_0^T+A_E)$, and likewise for $G'$;
\item[(E3)] $\abs{\nabla_1E(x,y)}\le A_\nabla \mathrm{e}^{-\kappa_Gd(x,y)}$ with $A_\nabla:=A_3^T+K_1(A_E+aA_P)$, and likewise for $E'$;
\item[(E4)] $\abs{\nabla_y(P_1G)(z,y)}$ and $\abs{\nabla_y(P_LG')(z,y)}$ are $\le A_P^\nabla \mathrm{e}^{-\kappa'd(z,y)}$.
\end{itemize}
\end{lemma}
\begin{proof}
$\mathsf X=F^{-1}-W$ gives $G=F+FWG$, hence $E=FWF+FWGWF$. Weighted Cauchy--Schwarz with Lemma~\ref{app:G1:VSR-CT}(ii) at $\rho_y:=d(\cdot,y)$ gives (E1). (E2): average $F+E$ over the unit block. (E3): $\nabla_1E=(1-\tilde\mu^2)[\nabla(F^2)+\nabla F\cdot E]-b\nabla F\cdot P_1G$, bounded by (F3), Lemma~\ref{app:G1:VSR-IM} and (F4). (E4): $G$ is symmetric, so $\nabla_yG(w,y)=\nabla F(y-w)+\nabla_1E(y,w)$; average over the block.
\end{proof}

\begin{lemma}[$G'$ maps decaying data to decaying data]\label{app:G1:VSR-Bdec}
If $0<\kappa\le\kappa'$ and $\abs{f(z)}\le \mathrm{e}^{-2\kappa d(z,y)}$, then $\abs{(G'f)(x)}\le A_B(\kappa)\mathrm{e}^{-\kappa d(x,y)}$, with $A_B(\kappa):=K_0+c_Fw_1(2L^2/\gamma_0)(1+1/\kappa)^{3/2}$. The same holds for $G$, with $\gamma_0$ and $\kappa\le\kappa_G$.
\end{lemma}
\begin{proof}
$G'f=Ff+FW'G'f$. The first term is $\le K_0\mathrm{e}^{-\kappa d(x,y)}$. For the second, use weighted Cauchy--Schwarz together with $\norm{\mathrm{e}^{\kappa\rho_y}f}_2^2\le(1+1/\kappa)^3$.
\end{proof}

\begin{lemma}[H\"older bound for $\nabla F$]\label{app:G1:VSR-Hol}
Let $\abs{h(z)}\le \mathrm{e}^{-\kappa d(z,y)}+Z_0$ with $0\le\kappa\le1/32$ and $Z_0\ge0$, and let $\delta:=d(x,x')$ and $\alpha\in[0,1)$. Then
\[
\abs{(\nabla_\mu Fh)(x)-(\nabla_\mu Fh)(x')}\le C_{\rm H}(\alpha)\delta^\alpha\big[\mathrm{e}^{-\kappa\min(d(x,y),d(x',y))}+Z_0\big],
\]
where $C_{\rm H}(\alpha):=2K_1+\mathrm{e}^{1/8}(38A_1^T+48A_2^T)(4+\log16)(1+1/(e(1-\alpha)))$.
\end{lemma}
\begin{proof}
If $\delta\ge1$, bound each term separately by $K_1(\dots)$. If $\eta\le\delta<1$, split into the near region $d(z,x)\le2\delta$, where the contribution is $\le38A_1^T\delta$, and the far region, where the second differences give $\le8A_2^T\delta(d(x,z)+\eta)^{-3}\mathrm{e}^{-\phi(d(x,z)-\delta)}$ and the sum is $\le6(\log(16/\delta)+2)$. Finally $\delta^{1-\alpha}\log(1/\delta)\le1/(e(1-\alpha))$.
\end{proof}

\begin{theorem}[Ba{\l}aban's line bounds for the zero-background slices at every scale]\label{app:G1:VSR-VS}
In the setting above, for all $x,x',y\in\Lambda$, directions $\mu,\nu$ and $\alpha\in[0,1)$:
\begin{itemize}
\item[(S1)] $\abs{H(x,y)}\le C_0\mathrm{e}^{-\kappa_Hd(x,y)}$;
\item[(S2)] $\abs{\nabla_{x,\mu}H(x,y)}$ and $\abs{\nabla_{y,\mu}H(x,y)}$ are $\le C_1\mathrm{e}^{-\kappa_Hd(x,y)}$;
\item[(S3)] $\abs{\nabla_{x,\mu}\nabla_{y,\nu}H(x,y)}\le C_2\mathrm{e}^{-\kappa_Hd(x,y)}$;
\item[(S4)] $\abs{\nabla_{x,\mu}H(x,y)-\nabla_{x,\mu}H(x',y)}\le C_3(\alpha)d(x,x')^\alpha \mathrm{e}^{-\kappa_H\min(d(x,y),d(x',y))}$, and the same for $\nabla_{x,\mu}\nabla_{y,\nu}H$;
\item[(S5)] for every lattice path $\Gamma$ of length $\abs\Gamma$, $\abs{\sum_{b\in\Gamma}\eta H(b,y)}\le\abs\Gamma C_0\mathrm{e}^{-\kappa_H\dist(\Gamma,y)}$, and likewise with $\nabla_y$.
\end{itemize}
The constants are $C_0:=C_VA_B(\kappa_H)$ with $C_V:=aA_P(1+L^{-2}\mathrm{e}^{3L\kappa_G})$; $C_g:=C_0+aA_P+aL^{-2}\mathrm{e}^{3L\kappa'}A_P'$; $C_1:=K_1C_g$; $C_2:=K_1(C_1+2aA_P^\nabla)$; and $C_3(\alpha):=C_{\rm H}(\alpha)\max(C_g,C_1+2aA_P^\nabla)$. They and the rate $\kappa_H$ depend only on $(a,L)$ (and on $\alpha$). They are uniform in $\eta\in(0,1]$, $\tilde\mu\in[0,1]$, $\ell\in L\N$, $a_j\in[a_\flat,a]$ and $x,y$.
\end{theorem}
\begin{proof}
(S1) Resolvent identity: $H=G'(\mathsf X-\mathsf X')G=G'VG$ with $V:=bP_1-b'P_L$. By (E2), $f:=VG(\cdot,y)$ obeys $\abs{f(z)}\le C_V\mathrm{e}^{-\kappa_Gd(z,y)}$. Since $\kappa_G\ge2\kappa_H$ and $\kappa_H\le\kappa'$, Lemma~\ref{app:G1:VSR-Bdec} applies.
(S2) From $\mathsf XG=\mathsf X'G'=1$ we get $(-\Delta+1)G=1+(1-\tilde\mu^2)G-bP_1G$, and similarly for $G'$. Subtracting, the Kronecker deltas cancel exactly:
\[
H(\cdot,y)=Fg_y,\qquad g_y:=(1-\tilde\mu^2)H(\cdot,y)+bP_1G(\cdot,y)-b'P_LG'(\cdot,y),
\]
with $\abs{g_y(z)}\le C_g\mathrm{e}^{-\kappa_Hd(z,y)}$ by (S1) and (E2). Hence $\abs{\nabla_xH}\le K_1C_g\mathrm{e}^{-\kappa_Hd}$, by Lemma~\ref{app:G1:VSR-IM}. $H$ is symmetric, which gives $\nabla_y$.
(S3) $\nabla_{y,\nu}H(\cdot,y)=F(\nabla_{y,\nu}g_y)$, and this is bounded by (S2) and (E4).
(S4) Apply Lemma~\ref{app:G1:VSR-Hol} with $h=g_y/C_g$, respectively $h=\nabla_{y,\nu}g_y/(C_1+2aA_P^\nabla)$, and $Z_0=0$.
(S5) Sum (S1), or (S2), over the bonds of $\Gamma$.
\end{proof}
\begin{remark}\label{app:G1:VSR-VSrem}
No cancellation between slices of different fields or scales is used. The only cancellations are those in (F3) and in the identity $H=Fg_y$. There is no zero mode, because $\mathsf X,\mathsf X'\ge\gamma_0L^{-2}$ on all of $\ell^2(\Lambda)$, constants included. No large-block hypothesis enters.
\end{remark}

\begin{corollary}[(VS) in physical units]\label{app:G1:VSR-VSphys}
Let $f\in\{s,v\}$, with $v$ taken per component (Lemma~\ref{app:G1:VSR-V0}), and let $t_{j+1}\le\min(1/m_f,L)$ on a torus with $\ell_T/t_j\in L\N$. (This holds on the tori $T_j$ whenever $t_{j+1}\le L$.) Then:
\begin{itemize}
\item $\abs{H_{j,f}(x,x')}\le C_0t_j^{-1}\mathrm{e}^{-\kappa_Hd_T(x,x')/t_j}$;
\item each lattice difference in $x$ or $x'$ costs one extra factor $t_j^{-1}$;
\item the $\alpha$-H\"older quotient of a derivative is of total order $t_j^{-2-\alpha}$;
\item a contour sum along Ba{\l}aban's $\Gamma^{(j+1)}$ gains one factor $t_j$, with constant $C_0dL^2/(L-1)$.
\end{itemize}
These are \cite[(2.10)--(2.12), p.~426]{Ba3} at $A=0$, with $\delta_1=\kappa_H$, below and above $\eps_0$ alike, uniformly in $\eps$, $j$, $\ell_T$, $\mu_*$ and $M$.
\end{corollary}

\begin{proposition}[the last remainder line]\label{app:G1:VSR-R}
Work in scale-$\bar K$ units: $t_{\bar K}=1$, $\eta=\eps/t_{\bar K}$, $\ell:=\ell_T/t_{\bar K}\in\N$, $\rho:=m_ft_{\bar K}\in(0,1]$. Put $\tilde C:=(-\Delta+\rho^2)^{-1}$, $\tilde G:=(-\Delta+\rho^2+bP_1)^{-1}$ with $b=a_{\bar K}$, and $\tilde R:=\tilde C-\tilde G$. In physical units, $R_{\bar K,f}=C_f-G_{\bar K,f}(0)$ has kernel $t_{\bar K}^{-1}\tilde R(x/t_{\bar K},y/t_{\bar K})$. Let $\kappa_R:=\min(\rho,\kappa_G)/8$ and $Z:=(\rho^2\ell^3)^{-1}=t_{\bar K}/(m_f^2\ell_T^3)$. Then:
\begin{itemize}
\item[(R1)] $\abs{\tilde R(x,y)}\le B_R[\mathrm{e}^{-\kappa_Rd(x,y)}+Z]$;
\item[(R2)] $\eta^3\sum_y\abs{\tilde R(x,y)}\le B_R'\rho^{-2}$;
\item[(R3)] $\abs{\nabla_x\tilde R}$, $\abs{\nabla_y\tilde R}$ and $\abs{\nabla_x\nabla_y\tilde R}$ are $\le B_R''[\mathrm{e}^{-\kappa_Rd(x,y)}+Z]$; H\"older quotients obey (S4) with $B_R(\alpha)[\mathrm{e}^{-\kappa_R\min d}+Z]$; contour sums obey (S5).
\end{itemize}
The constants $B_R,B_R',B_R'',B_R(\alpha)$ depend only on $(a,L)$ (and on $\alpha$), and not on $\rho$, $\eta$ or $\ell$.
\end{proposition}
\begin{proof}
(R1) $\tilde R=\tilde C\,bP_1\tilde G$, and $f_y:=bP_1\tilde G(\cdot,y)$ obeys $\abs{f_y}\le aA_P\mathrm{e}^{-\kappa_Gd(\cdot,y)}$ by (E2) at $\tilde\mu=\rho$. Unfold the torus: $\abs{\tilde R(x,y)}\le aA_P\sum_nI(\hat x-\hat y-n\ell)$ with $I(D):=\eta^3\sum_v\tilde C^\infty(D-v)\mathrm{e}^{-\kappa_G\abs v_1}$. By Lemma~\ref{app:G1:VSR-F} at mass $\rho$, $I(D)\le A_I\mathrm{e}^{-\kappa_R\abs D_1}/(1+\abs D_1)$ with $A_I:=A_0(13+J_0(\kappa_G/2))(3+8/\kappa_G)$, which is free of $\rho$. For $n\ne0$, $\abs{\hat x-\hat y-n\ell}_1\ge\ell\abs n_1/2$ and $\#\{\abs n_1=k\}\le6k^2$, so $\sum_{n\ne0}I\le48A_I/(\kappa_R^2\ell^3)\le3072A_I\kappa_G^{-2}Z$, using $\kappa_R\ge\rho\kappa_G/8$. Hence $B_R:=aA_PA_I\max(1,3072\kappa_G^{-2})$.
(R2) $\eta^3\sum_y\abs{\tilde R}\le aA_PJ_0(\kappa_G)\,\eta^3\sum_z\tilde C(x-z)=aA_PJ_0(\kappa_G)\rho^{-2}$.
(R3) $\tilde R(\cdot,y)=F\tilde g_y$ with $\tilde g_y:=(1-\rho^2)\tilde R(\cdot,y)+bP_1\tilde G(\cdot,y)$, and $\abs{\tilde g_y}\le(B_R+aA_P)[\mathrm{e}^{-\kappa_Rd}+Z]$. Then argue as in (S2)--(S4); note $\kappa_R\le\kappa_G/8\le1/128$.
\end{proof}

\paragraph{The renormalisation machinery (S6).}
\emph{Conventions} \cite[(2.1)--(2.16), pp.~423--428]{Ba3}. A \emph{generalised graph} has two kinds of lines and vertices:
\begin{itemize}
\item lines with slices obeying \cite[(2.10)--(2.12)]{Ba3}, with an arbitrary line dimension $\delta_l$;
\item vertices with $\eta$-powers, H\"older factors $\abs{x-x'}^\gamma$ ($\gamma\ge0$) and external fields in $\norm\cdot_{1,\alpha_0}$.
\end{itemize}
Fix an ordering $\bar l=(l(1),\dots,l(m))$ of the lines, with nested subgraphs $G_1\subset\dots\subset G_m=G$. $D(\cdot)$ is the degree of \cite[(2.2)--(2.3)]{Ba3}, and $J(\bar l)$ is the set of index assignments compatible with $\bar l$. In the \emph{pinned} expression $E^{(x)}$, the vertex $v_0$ carrying a constant external leg (the pin $p$) is fixed at $x$: it has no factor $(L^{j(v_0)}\eta)^d$ and no localisation sum.

\begin{lemma}[the shrinking induction, with or without a pin]\label{app:G1:VSR-A}
Let $G$ be a connected generalised graph, and assume every connected component $G_i^{(\alpha)}$, $i<m$, has $D(G_i^{(\alpha)})\ge\vartheta>0$.
\begin{itemize}
\item[(i)] Unpinned: $\sum_{j\in J(\bar l)}\sum_{\{\Delta(v)\}}\tilde E(G(j),\{\Delta(v)\})\le C_A\sum_{j_{l(m)}}(L^{j_{l(m)}}\eta)^{D(G)}$. This is $\le C_A$ if $D(G)>0$, and $\le C_A(L^{j^*}\eta)^{D(G)}$ when the largest index is fixed to $j^*$, with no sign condition on $D(G)$.
\item[(ii)] Pinned: the same bounds hold, with the same $C_A$, when one vertex $v_0$ is fixed at a point $x$.
\end{itemize}
Here
\[
C_A\le c_d^m(4/\delta_1)^{dm}(4m+2)^{dm(m-1)/2}(1-L^{-\vartheta})^{-m},
\]
with $m\le m(\nbar)$ the maximal number of lines of a graph of order $\le\nbar$, and $\vartheta\ge\min(\alpha_0,\frac12)$. $C_A$ depends on $d,\delta_1,L,\alpha_0,\nbar$ and is independent of $\eta$, $k$, $\eps$, $x$, the localisations, $\ell_T$, $\mu_*$ and $m^2$.
\end{lemma}
\begin{proof}
Write $\tilde E$ as in \cite[(2.14)]{Ba3}:
\begin{itemize}
\item each vertex $v$ carries $(L^{j(v)}\eta)^d$ and a sum over $\Delta(v)\subset\square(v)$;
\item each leg of a line $l$ carries $(L^{j_l}\eta)^{\delta_l/2}$, and each derivative on $l$ carries $(L^{j_l}\eta)^{-1}$;
\item each line carries $\exp[-\delta^{(1)}(L^{j_l}\eta)^{-1}\dist(\Delta,\Delta')]$, with $\delta^{(1)}:=\delta_1/2$ after the tree factor has been extracted.
\end{itemize}
Induct over $i$: shrink $l:=l(i+1)$, with index $j$ and $s:=L^j\eta$, joining $w,w'$.

(a$'$) Suppose $w\ne w'$. Let $n'\le2m$ be the number of legs of other lines at $w'$, and put $\delta^{(i+1)}:=\delta^{(i)}/(2(n'+1))$. For each such line $l'$ (scale $s'\ge s$, other end $w''$),
\[
\mathrm{e}^{-\delta^{(i)}\dist(\Delta(w'),\Delta(w''))/s'}\le \mathrm{e}^{-\delta^{(i+1)}\dist(\Delta(w),\Delta(w''))/s'}\,\mathrm{e}^{+\delta^{(i+1)}\dist(\Delta(w),\Delta(w'))/s}.
\]
The $n'$ costs total at most $\mathrm{e}^{+\frac12\delta^{(i)}\dist/s}$. They are paid by $l$'s own factor, which leaves $\mathrm{e}^{-\frac12\delta^{(i)}\dist/s}$, and with it $\sum_{\Delta(w')}\le c_d(\delta^{(i)}/2)^{-d}$.
(b) Relocalise the merged vertex using $w$'s volume factor: $\sum_{\Delta(w)\subset\Delta(w\cup w')}(L^j\eta)^d=(L^{j(w\cup w')}\eta)^d$.
(c) The remaining powers of $L^j\eta$ multiply to $(L^j\eta)^{D(G_{i+1}^{(\beta)})}$, by additivity of the degree.
(d) Summing $j_{l(i+1)}\le j_{l(i+2)}$ gives the factor $(1-L^{-\vartheta})^{-1}$.
(e) If $w=w'$, only (c)--(d) are needed.

Since $\delta^{(i)}\ge(\delta_1/2)(4m+2)^{-(i-1)}$, the product of the factors $c_d(\delta^{(i)}/2)^{-d}(1-L^{-\vartheta})^{-1}$ is the displayed $C_A$.
(ii) In step (a) choose $w$ to be the vertex containing $v_0$. Then $w'$ never contains the pin, and its volume factor is used as before. The pinned merged vertex needs no volume factor, so the exponent attached to the component is again $D(G_{i+1}^{(\beta)})$. At the end the pin contributes $1$ in place of $\abs\square=1$.
\end{proof}

\begin{lemma}[largest index fixed]\label{app:G1:VSR-Btop}
In the decomposition \cite[(3.4)--(3.5)]{Ba3}, fix the index of the last line $l(m)$ to $j^*$. Then every divergent proper subgraph renormalised within an ordering is a $G_i$ with $i<m$, and contains only lines of index $\le j^*$. The transformations of \cite[(3.5)]{Ba3} act inside the sub-sum with $j_{l(m)}=j^*$, and Lemma~\ref{app:G1:VSR-A}(i) at fixed $j^*$ gives $(L^{j^*}\eta)^{D(G)}$, with no renormalisation of $G$ itself.
\end{lemma}
\begin{proof}
A connected proper subgraph containing $l(m)$ is not of the form $G_i$. The transformations of a $G_i$ sum only over internal indices below its lowest external index (\cite[pp.~437, 444]{Ba3}), which is $\le j^*$. The transformed graph has degree $\ge D(G_0)$, and $L^{j^*}\eta\le1$.
\end{proof}

\begin{theorem}[leg transport {\cite[(3.35)]{Ba3}} for every member of class C5, every ordering and every receiving vertex]\label{app:G1:VSR-T335}
Let $G_0$ be a C5 graph ($\phi\phi A$ class) with external legs $E$. Let $\bar l$ be an ordering with at most one outer divergent two-vertex component $H=\{a,b;h\}$ of type \cite[(2.4)]{Ba3} (true for every C5 member). Let $r$ be any vertex, and let the pin, if any, be $r$. Transport the legs by exact telescoping moves
\[
\textstyle\prod_{\ell\in S}F_\ell(x_u)=\prod_{\ell\in S}F_\ell(x_{u'})+\sum_i[F_{\ell_i}(x_u)-F_{\ell_i}(x_{u'})]\prod_{k<i}F_{\ell_k}(x_{u'})\prod_{k>i}F_{\ell_k}(x_u),
\]
keeping the factors $\abs{x_u-x_{u'}}^\alpha$ unsplit until the end and then splitting them along the path by subadditivity. Use the following schedule:
\begin{itemize}
\item if there is no $H$, move every leg to $r$ along any path;
\item if $r=b$, first move $a$'s legs to $b$ along $h$, then all others;
\item if $r\notin H$, first move $a$'s legs to $b$ along $h$, then the legs at $b$ and elsewhere to $r$.
\end{itemize}
Then every remainder either carries $\abs{x_a-x_b}^\alpha$ on $h$, or has $a$ leg-free with $a\ne r$. In the second case \cite[(2.8)]{Ba3} at $a$ acts on no external leg. After (2.8), every component has degree $>0$, with $D_R(H)\ge\alpha$, $D_R(K)\ge\frac12$ otherwise, and $D_R(G_0)\ge\alpha$. The local term has all legs at $r$, and is Ba{\l}aban's \cite[(3.36)]{Ba3}.
\end{theorem}
\begin{proof}
Moves along the single line $h$ give $+\alpha$ on $h$, so $D_R(H)=\alpha$. In the later moves no leg sits at $a$ and $a\notin\{c,b,r\}$, so $a$ is leg-free. Then (2.8) at $a$ differentiates only internal lines of $G_0$ or $g$. This never lowers a component degree, and it raises $D(H)$ to $1$. The chain piece adds $\alpha$. The components $K\ne H,G_0$ are not divergent, so $D(K)\ge\frac12$, since degrees are half-integers.
\end{proof}

\begin{proposition}[the pinned bound {\cite[(2.21e)]{Ba3}} with the pin at the external-photon vertex]\label{app:G1:VSR-5b}
Let $S$ be a (2.21e)-shaped component of a pinned family with $p=L\in S$. Here $L$ is an A1 vertex carrying leg 1 and the line $c$, and $x_1$ is an A2 vertex carrying a photon tadpole and an explicit $\eta$. Transport $S$'s external leg at $x_1$ to $p$ along $c$, as in Theorem~\ref{app:G1:VSR-T335} with no $H$. Then:
(i) every remainder carries $\abs{x_1-p}^\alpha$ on $c$ and has $D_R(S)=\alpha>0$;
(ii) the local term is (the fields of $S$ at $p$)$\times\Gamma_e(p)$, with $\Gamma_e(p):=\sum_{x_1}\eta^dK_c(p,x_1)\kappa(x_1)\prod g(x_1)$, where $K_c=(D_\mu)_p(D_\nu)_{x_1}c(p,x_1)$ carries one derivative at each end of $c$ and $\kappa(x_1)\propto\eta D^{(<j_0)}(x_1,x_1)$ is the photon tadpole;
\[
\abs{\Gamma_e(p)}\le C_e(d,L,a,M_b,\alpha_0,\norm g_{C^1}),\quad\text{uniformly in }\eta,j_0,k,p.
\]
\end{proposition}
\begin{proof}
(i) In the ordering $(\mathrm{loop},c)$ the component $\{\mathrm{loop}\}$ has $D=0$. Replace it first by its bounded value $\abs\kappa\le O(1)$, using $\eta\sum_jO((L^j\eta)^{-1})\le O(1)L/(L-1)$. Then apply Lemma~\ref{app:G1:VSR-A}(ii) to the remainders.
(ii) First extend the $g$-weighted $x_1$-sum to all of $\eta\Z^d$; the support of $\prod g$ lies inside the cube. Next resum $S$'s internal indices up to $j_0$; this is legitimate because the base $p$ is independent of the ordering and of the indices. Then rescale to $\xi=L^{-j_0}$, with $\rho:=\abs{y-y'}+\xi$. Write $c=C^\xi+G^\xi_{j_0}(0)\Upsilon C^\xi$, and $\kappa=\kappa_0+\kappa_1$ with $\kappa_0:=c_v\xi C_v^\xi(0)$, a constant.

The photon line is $D-C_v=-D(D^{-1}-C_v^{-1})C_v$. Here the middle factor $(\mu_*^2t^2-1)+aQ^*Q$ is bounded and of range $O(1)$: at $A=0$ Ba{\l}aban's and King's photon $G_{j_0}(0)$ is the componentwise plain block average, with no contours. With the vector line bounds this gives $(D-C_v)(y,y)=O(1)$, so $\kappa_1=O(\xi)$.

Split $\Gamma_e$ into four pieces:
\begin{itemize}
\item[(1)] the translation-invariant piece $\kappa_0\prod g(p)\sum_y\xi^dK^C(p,y)$ is exactly $0$, as a summed lattice difference of a summable kernel;
\item[(2)] the $g$-difference: $\abs{K^C}\le O(1)\rho^{-3}\mathrm{e}^{-\delta\rho}$ and $\abs{\prod g(y)-\prod g(p)}\le C\norm g_{C^1}\abs{y-p}$, so this piece is $O(1)$;
\item[(3)] the tadpole difference is $O(\xi)\sum\xi^d\rho^{-3}\mathrm{e}^{-\delta\rho}=O(\xi\log(1/\xi))$;
\item[(4)] the line difference: $\abs{K^\Upsilon}\le O(1)\rho^{-1}\mathrm{e}^{-\delta'\rho}$ (Step~3 of Lemma~\ref{app:G1:AA}) and $\abs\kappa\le O(1)$, so this piece is $O(1)$.
\end{itemize}
\end{proof}

\begin{lemma}[renormalisation with a pin]\label{app:G1:VSR-C}
Let $G_{\rm ren}$ be a family each of whose members carries the external leg 1 with $\Phi_{\rm ext}\equiv1$ at the pin $p$. Then every transformation of \cite[pp.~434--445]{Ba3} can be performed without summing over the position of $p$, and $E^{(x)}(G_{\rm ren},\dots)$ obeys \cite[(1.33)]{Ba3} with the same $O(1)$ and $\delta_0$.
\end{lemma}
\begin{proof}
(i) A subgraph containing $p$ has a scalar external leg at $p$, so it is never a $\Pi$/WT class. Those transformations act only on subgraphs not containing $p$.
(ii) A (2.4) subgraph through $p$ has $D=0$, which forces both derivatives onto its internal scalar line. Apply (2.8) at its other vertex $a\ne p$. At most one (2.4) contains $p$.
(iii) Scalar self-energies containing $p$: Ba{\l}aban's rule expands the leg with the smaller index \cite[pp.~435--436]{Ba3}. Leg 1 is unit-scale, so the base of the subtraction is $p$, and the counterterm partner is localised at $p$. The symmetry $\Sigma_S(y,z)_{ab}=\Sigma_S(z,y)_{ba}$ makes this choice legitimate.
(iv) (3.21)-classes containing $p$: at $\tilde B=0$, a derivative on leg 1 vanishes. At $\tilde B\ne0$, $D_{\tilde B}1=O(p(t))\ne0$, and the differentiated leg may be leg 1. Then (3.21)--(3.24) are based at $p$ by the rule of (iii), expanding the other leg.
(v) $\phi\phi A$ classes containing $p$: use Theorem~\ref{app:G1:VSR-T335} with $r=p$, and Proposition~\ref{app:G1:VSR-5b} for the pinned (2.21e) case. The vanishing \cite[(3.37)--(3.38)]{Ba3} is a statement about a translation-invariant $C^\xi$-expression.
(vi) Nested and one-particle-reducible cases: induct on the last divergent subgraph; the transformations used never sum over a vertex carrying an external field.
(vii) Line cuts and localisation \cite[pp.~432--433]{Ba3} act at fixed positions, and a pinned vertex is the localisation $\{x\}\subset\square(p)$. Cuts use the graph representation with partners of \emph{cut} divergent subgraphs.
Combining (i)--(vii) with Lemma~\ref{app:G1:VSR-A}(ii) gives the claim.
\end{proof}
\begin{remark}[how (S6) is used]\label{app:G1:VSR-S6use}
In Lemma~\ref{app:G1:common}, every renormalised subgraph has $j_0\le\bar K\le\min_fK_f$, and its internal lines resum to $G_{j_0,f}(0)$, of the same algebraic form for both fields. This is exactly the case in which Lemma~\ref{app:G1:VSR-C}, Step~3 of Lemma~\ref{app:G1:AA} and Theorem~\ref{app:G1:VSR-T335} apply. Lemma~\ref{app:G1:TB} uses Lemma~\ref{app:G1:VSR-A}(ii) in its pinned form, which is uniform in the position of the pin.
\end{remark}

\paragraph{The inputs (K1)--(K4), (IH) and (CAP) of the tail bound.}
\begin{itemize}
\item[(K1)] For the background $\mathfrak B$, regular in the sense of \cite[Prop.~2.1]{Ba1} (input (REG)), $G_k(\Omega,\mathfrak B)$ obeys \cite[(1.8), (1.10), (1.12)]{Ba4}: decay $c_0\mathrm{e}^{-\delta_0\abs{x-y}}$, the first derivative and the $\alpha$-H\"older quotient of $DG_kf$ ($\alpha<1$), uniformly in $\eta$; and $G_kQ_k^*$ decays. These hold for $0\le M^2t^2\le c$ in $G_k$, uniformly there.
\item[(K2)] The slice bounds \cite[(2.10)--(2.12)]{Ba3} for the lines of the families of $\mathcal P^{(k)}$, obtained in \cite[p.~426]{Ba3} from \cite[Props.~2.1, 2.3]{Ba1} as $a_j^2t_j^{-4}G_jQ_j^*C^{(j)}Q_jG_j$. At zero background, Corollary~\ref{app:G1:VSR-VSphys} gives an independent proof.
\end{itemize}

\begin{proposition}[(K3): the one- and two-variable remainder bounds]\label{app:G1:VSR-K3}
Let $V:=G_k(\Omega,\tilde B)^{-1}-G_k(\Omega,\tilde A+\tilde B)^{-1}$, $\mathcal K:=[G(\tilde B)V]^nG(\tilde A+\tilde B)[VG(\tilde B)]^{n'}$, $m:=n+n'$ and $\sigma:=e(t)(\norm{\tilde A}_\infty+\norm{\partial^\eta\tilde A}_\infty)$. Assume the hypothesis (R$_\vartheta$) of Subsection~\ref{app:G1:part-E2}: $\theta(\tilde B),\theta(\tilde A+\tilde B)\le\vartheta_R$ and $\sigma\le1$. This hypothesis follows from (E.1), read with the constant $C_A$ of $\sigma\le eC_At^{1/2}p$.
\begin{enumerate}
\item[(a)] For every $m\ge0$ and $\alpha<1$: $\abs{D^a\mathcal Kf(x)}$ and the $\alpha$-H\"older quotient of $D\mathcal Kf$ are $\le C(C_\infty\sigma)^m\mathrm{e}^{-\beta\dist(x,\supp f)}\norm f_\infty$. The same holds for $\mathcal K^*$.
\item[(b)] Let $0\le h,h'\le1$ with $\abs{\partial^\eta_\mu h},\abs{\partial^\eta_\nu\partial^\eta_\mu h}\le c_h$ for all $\mu,\nu$, and likewise for $h'$; the cut-offs of Lemma~\ref{app:G1:ramp-ch} qualify. Assume $m\ge6$ and $0<\alpha<1$, or $m\ge4$ and $0<\alpha<\frac12$. Then for $0<\beta\le\beta_U$, on $T_\eta$,
\[
\norm{h\mathcal Kh'}_{1,\alpha}\le \mathrm{e}^{5\beta}(1+6c_h)^2C_U(\alpha)(C_V\sigma)^m\mathrm{e}^{-\beta\dist(\supp h,\supp h')}
\]
in the mixed two-variable norm, including the $D_xD^*_{x'}$ quotients. The same bound holds on the Neumann domains of Proposition~\ref{app:G1:sandwich}(b), under its hypothesis on $\supp\tilde A$, if $\supp h\cup\supp h'$ has depth $\ge R_0+3$. If the quotients over distances $>1$ are included, multiply the constant by $4$.
\end{enumerate}
The bounds are uniform in $\eta$, $k$, $\eps$, $\ell_T$, and in $0\le m^2t^2\le c$. The constants depend on $d,a,M_b,L,\gamma_0,K_W,\alpha,\beta$. In \cite[(2.5)]{Ba3}, the tree rate of \cite[Prop.~1]{Ba3} is taken $\le\beta_U$.
\end{proposition}
\begin{proof}
(b) Apply Lemma~\ref{app:G1:leibniz-loc} with $M:=C_U(\alpha)(C_V\sigma)^m$.
\begin{itemize}
\item On $T_\eta$, take $M$ from Lemma~\ref{app:G1:I3r-U1m}(b) and $\mathcal P:=T_\eta$.
\item On the Neumann domains, take $M$ from Proposition~\ref{app:G1:sandwich}(b) and $\mathcal P:=\{z:\text{depth of }z\ge R_0+1+\eta\}$. Points of this $\mathcal P$ and their bond partners have depth $\ge R_0$, so Proposition~\ref{app:G1:sandwich}(b) gives \eqref{app:G1:leib-hyp} on $\mathcal P$. Since $\eta\le\frac12$, every point within $\ell^\infty$-distance $1+\eta$ of the supports lies in $\mathcal P$.
\end{itemize}

(a) Write $D^a\mathcal KD^{b*}=(\text{end})\cdot\mathcal N\cdot(\text{end})^*$. In each end, push every vertex derivative onto the line to its right \cite[(2.32)]{Ba4}, so that every end line carries one derivative. The end factors lie in weighted $\ell^2$ when the kernel exponent exceeds $3/2$. This is why two-line ends need $\alpha<\frac12$, while three-line ends allow every $\alpha<1$. The doubly differentiated lines sit in the middle factor $\mathcal N$, which is bounded in weighted $\ell^2$ by (K4). Cauchy--Schwarz at the two end factors converts $\ell^2$ to $\ell^\infty$, and no pointwise bound on $D_xD_yG$ is used. The details are those of the proof of Proposition~\ref{app:G1:sandwich}, with one end replaced by $f$, and of Lemmas~\ref{app:G1:I3r-U1} and \ref{app:G1:I3r-U1m}.
\end{proof}

\begin{lemma}[(K4): weighted $\ell^2$ bounds]\label{app:G1:VSR-K4}
Let $\rho(x):=\abs{x-y}$ and $O^\beta:=\mathrm{e}^{\beta\rho}O\mathrm{e}^{-\beta\rho}$. Assume:
\begin{itemize}
\item[(H0)] $\mathsf X(A)-m^2t_j^2-v\ge\gamma_0$ \cite[(1.8)]{Ba4};
\item[(M1)] $\gamma_0'\le M(A)\le a_j+aL^{-2}$ \cite[(1.15)]{Ba4};
\item[(M2)] $\abs{M(A)(y,y')}\le c_M\mathrm{e}^{-\delta\abs{y-y'}}$.
\end{itemize}
Then there are $\beta_*>0$ and $K_W=K_W(d,a,\gamma_0,L)$ such that for $\abs\beta\le\beta_*$, uniformly in $\eta,j,k,\eps,\Omega$, the boundary conditions, regular $A$, $v\ge0$ and $s\ge0$:
(i) $\norm{G^\beta}\le2/\gamma_0$, and $\norm{(\nabla_AG)^\beta}$, $\norm{(G\nabla_A^*)^\beta}$, $\norm{(\nabla_AG\nabla_A^*)^\beta}$ are $\le K_W$;
(ii) $\norm{Q_j(A)^\beta},\norm{(Q_j(A)^*)^\beta}\le \mathrm{e}^{\beta\sqrt d}$;
(iii) $\norm{((M+s)^{-1})^\beta}\le(\gamma_0'/2+s)^{-1}$.
\end{lemma}
\begin{proof}
(i) $G^\beta=(\mathsf X^\beta)^{-1}$. Lemma~\ref{app:G1:WK-CT} applies when $c(\beta):=d\beta^2\cosh\beta+a(\mathrm{e}^{\sqrt d\beta}-1)\le\gamma_0/4$, and gives $\norm w\le2\gamma_0^{-1}\norm F$ and $\norm{\nabla_Aw}^2\le4\gamma_0^{-1}\norm F^2$ for $w=G^\beta F$. Next, $(\nabla_AG)^\beta=\nabla_A^\beta G^\beta$ with $\nabla_A^\beta=\nabla_A+\beta K_\beta$ and $\norm{K_\beta}\le \mathrm{e}^{\beta_*}$ per direction, hence $\le\sqrt d\,\mathrm{e}^{\beta_*}$ for the vector difference (Lemma~\ref{app:G1:WK-W}(b)); and $(G\nabla^*)^\beta=((\nabla G)^{-\beta})^*$. For $\nabla G\nabla^*$, let $w:=G^\beta(\nabla^*)^\beta g$. Then $\rRe\ip w{\mathsf X^\beta w}\ge\frac12\norm{\nabla_Aw}^2+\frac{\gamma_0}4\norm w^2$ and $\rRe\ip w{\mathsf X^\beta w}\le(\norm{\nabla_Aw}+\mathrm{e}^{\beta_*}\beta\norm w)\norm g$.
(ii) Schur test, using $\abs{\rho(x)-\rho(y')}\le\sqrt d$ on the support of $Q_j$.
(iii) By (M2) and Schur, $\norm{M^\beta-M}\le\gamma_0'/2$; then use (M1).
\end{proof}

\begin{definition}[(IH): the family bound carried by $\mathcal P^{(k)}_w$]\label{app:G1:VSR-IH}
For every localised family $G^\natural_{\rm ren}$ of $\mathcal P^{(k)}_w$,
\[
\abs{E(G^\natural_{\rm ren},\{\square(v)\},\Phi_{\rm ext},A_{\rm ext})}\le O(1)(\abs uc_k)^m(e(t_k))^{d_v}(\lambda(t_k))^{d_s}\mathrm{e}^{-\delta_0d(\{\square(v)\})}\norm{h\Phi_{\rm ext}}_{1,\alpha_0}\norm{h'A_{\rm ext}}_{1,\alpha_0}.
\]
Here $\norm\cdot_{1,\alpha_0}$ is the covariant H\"older norm \cite[(1.32), p.~420]{Ba3}; $m$ is the number of D-lines; $c_k\le1/3$; $O(1)=O(\alpha_0,\nbar)$; and $\delta_0>0$ is the tree rate of \cite[(1.33)]{Ba3}. At the cap it is some $\delta(\Nbar)>0$ (Proposition~\ref{app:G1:VSR-CAP}); \cite[p.~420]{Ba3} says that it depends on $d$ only, which is not what the proof gives (the line-moving rates of \cite[p.~427]{Ba3} depend on $\nbar$). Everything is on the small-field restrictions $\chi_k$ of \cite[(2.55), p.~570]{Ba2}. On $\Lambda_{-1}^{(k-1)\prime}$ these are:
\begin{itemize}
\item $\abs{(\partial A)(b)}\le c_1p(t_{k-1})$;
\item $\abs{A(x)}\le c_1p(t_{k-1})/(\mu_*t_{k-1})$;
\item $\abs{(D_{\tilde A^{(k)}}\phi)(b)}\le c_1p(t_{k-1})$;
\item $\abs{\phi(x)}\le c_1\lambda(t_{k-1})^{-1/4}p(t_{k-1})$.
\end{itemize}
\end{definition}

\begin{proposition}[(CAP): the family bound at a finite order cap]\label{app:G1:VSR-CAP}
Fix $\Nbar=\nbar^2=169$. Every family of graded coupling order $q\le\Nbar$ obeys \cite[(1.33)]{Ba3}, and its D-line form of Definition~\ref{app:G1:VSR-IH}. The constant is $O(1)=O(\alpha_0,\Nbar;d,L,a,M_b)$, with $C_A$ of Lemma~\ref{app:G1:VSR-A} at $m\le m(\Nbar)$; the tree rate is some $\delta(\Nbar)>0$; the $u$-factor is $(1+\abs uc_k)^{N_v}\le5^{84}$ at $L=2$. These are independent of $\eps$, $k$, $x$ and the domains. Families containing $\delta m^2_{\rm fin}$ vertices are uniform in $\ell_T\ge\ell_0$, being evaluated at the torus floor. The conditions (R.cap) and (T.cap) of Appendix~\ref{app:B} are read at the cap: $R\hat\ell^2\ge C_{\rm loc}(169)$ and $\ell_T\ge3r(\eps_0)\eps_0+C_{169}\eps_0$.
\end{proposition}
\begin{proof}
Every step of (S6) and of Theorem~\ref{app:G1:WP-WP1} is a finite induction over one graph of fixed order, namely the sequence $G_1\subset\dots\subset G_m$ with $m\le m(q)$. There is no induction over the order and no series in the order. So the cap enters in four ways only:
\begin{itemize}
\item[(i)] finite counts: the number of graphs, $m(\Nbar)$ lines, $N_s(\Nbar)$ localisation labels, and $N_v\le\Nbar/2$ vector lines;
\item[(ii)] the line-moving rates $\delta^{(i)}\ge(\delta_1/2)(4m+2)^{-(i-1)}$ of Lemma~\ref{app:G1:VSR-A};
\item[(iii)] the constant $\sup_{t\in(0,1]}t^{(q-13)/4+1/20}p(t)^{4q}<\infty$;
\item[(iv)] the size $O(\nbar)$ of the component cube inside $\square_1\subset\square$ \cite[pp.~432--433]{Ba3}, which is where (R.cap) and (T.cap) tighten with the cap.
\end{itemize}
Renormalisation does not depend on the cap in $d=3$. With $\omega=3-E/2-\alpha/2-\beta$, every superficially divergent graph has $\alpha+2\beta\le6-E\le6<\nbar$, so no counterterm object is re-read at the cap. All new numbers are fixed at stage 4 of Appendix~\ref{app:B}, before $\eps_0$ is chosen, so the order of choices stays acyclic.
\end{proof}

\subsection{The ingredients}\label{app:G1:part-E4}

We prove the ingredients of the proof of Theorem~\ref{thm:T6} listed in Table~\ref{tab:T6-inputs}: the moved vacuum pieces (Proposition~\ref{app:G1:I1}), the Ward--Takahashi step on the torus (Proposition~\ref{app:G1:I2}), the use of real backgrounds only (Proposition~\ref{app:G1:I4}), the order cap (Proposition~\ref{app:G1:I5}), the conditions of \cite[Prop.~1]{Ba3} at its applications (Proposition~\ref{app:G1:I3}), the mass reading of \cite[p.~589]{Ba2} (Proposition~\ref{app:G1:I7}), the final c-numbers (Proposition~\ref{app:G1:I8}) and the vacuum classes at the last scale (Proposition~\ref{app:G1:I9}). The per-localisation bound, the geometric and region-sum lemmas, and \cite[Prop.~1]{Ba3} with its sketched steps written out are stated in Subsection~\ref{app:G1:part-TI145}.
Throughout: $d=3$, $L=2$, $\nbar=13$, $\Nbar=169$, $t_k=L^k\eps$ with $t_K=\eps_0$, $r(t)=R(1+\log t^{-1})^{\varrho}$ with $\varrho=2$ (Appendix~\ref{app:B}), and the two points $\pi_+=(u,M)=(1,m_2)$, $\pi_-=(-\eta,m_1)$ of \eqref{eq:T6}.

\paragraph{The moved vacuum pieces at every scale.}\label{app:G1:ssec-I1}
\emph{Notation at step $k$.} The scale-$k$ blocks are the unit cubes of the lattice rescaled by $t_k$. For a union $X$ of such blocks, $\abs X_k$ is their number. Let $\Lambda_7^{(k)}$ be Ba{\l}aban's small-field set at step $k$ \cite[pp.~557--558]{Ba2}, and put $W_k:=B^k(\Lambda_7^{(k)})$. The \emph{layer} is $\mathcal L_k:=\{\square\subset W_k:\dist(\square,W_k^c)\le r(t_k)\}$, which equals $B^k(\Lambda_7^{(k)}\cap\Lambda_8^{(k)c})$ of \cite[p.~422]{Ba3}. The \emph{deep part} is $\mathcal D_k:=W_k\setminus\mathcal L_k$. The \emph{collar} is $\mathrm{Col}_k:=\Lambda_7^{(k-1)\prime}\cap\Lambda_7^{(k)c}$, and $\mathrm{Old}_k:=\Lambda_7^{(k-1)\prime c}$. Since $\Lambda_7^{(k)}\subset\Lambda_7^{(k-1)\prime}$, we have $W_k^c=\mathrm{Col}_k\sqcup\mathrm{Old}_k=\Lambda_7^{(k)c}$.
Let $E_K(u)$ be the vacuum counterterm of the slice model, and $E_{K,k}(u)$ its part from graphs with operator indices in $[k,K)$ \cite[(1.27)--(1.28), p.~418]{Ba3}, with $E_{K,0}=E_K$ and $E_{K,K}=0$. \emph{Convention (G)} is Ba{\l}aban's \cite[(1.30), p.~419]{Ba3}: at step $k\ge1$, $E_{\rm fin}(u)$ and $E_{K,k}(u)$ are global c-numbers, with no $\Omega_1$ and no $g_k$. The set moved at step $k$ is
\[
\mathcal V_k(u):=E_{K,k}(u)-E_{K,k+1}(u),
\]
the terms whose indices are all $\le k$, at least one of them equal to $k$ \cite[p.~422]{Ba3}.

Fix a pair $(G,\tilde l)$ of a graph and one of Ba{\l}aban's orderings of its lines, whose last line has index $k$. Its terms are localised by the rule of \cite[p.~420]{Ba3}: there is one unit cube per vertex of type (1.6)/(1.7), and one support cube of side $2$ per vector leg of the other vertices. We write $\Lambda=(\square_\iota)_{\iota\in I(G)}$, $N:=\abs{I(G)}\le\bar N(\nbar)<\infty$. Labels at one vertex have overlapping supports (otherwise the term vanishes). So they are at distance $0$, and the tree length $d(\Lambda)$ of \cite[(1.33)]{Ba3} is the same over labels and over vertices. Order $I(G)$ so that the two ends of the last line come first. Then:
\begin{itemize}
\item $\Lambda$ is \emph{good} if $\Lambda\subset W_k$;
\item $\Lambda$ is \emph{bad} if it meets both $W_k$ and $W_k^c$; its anchor is the first index whose cube lies in $W_k$;
\item $\Lambda$ is \emph{outside} if $\Lambda\cap W_k=\emptyset$; its anchor is the first index.
\end{itemize}
Every $\Lambda$ has exactly one class and one anchor.
If the moved terms are localised at scale $t_{k+1}$ rather than $t_k$ (the two readings of ``passing from $k$ to $k+1$'', p.~422), block counts change by factors $L^{\pm3}$ and the constants below by at most $L^3$.

\begin{lemma}[the localisation sum]\label{app:G1:treesum}
Let $N\ge1$ and $\delta>0$. On the scale-$k$ block lattice $(\Z/n)^3$ of $T_\eps$, fix an anchor index $\iota_0$ and a cube $\square_0$. Then for every $D_0\ge0$,
\begin{gather*}
\sum_{\Lambda:\ \square_{\iota_0}=\square_0,\ d(\Lambda)\ge D_0}\mathrm{e}^{-\delta d(\Lambda)}\le \mathrm{e}^{-\delta D_0/2}S_N(\delta),\\
S_N(\delta):=N^{N-2}\big(\mathrm{e}^{\delta}\sigma(\delta)\big)^{N-1},
\end{gather*}
where $\sigma(\delta):=1+\sum_{r\ge1}(24r^2+2)\mathrm{e}^{-\delta r/2}$. The bound does not depend on $n$, i.e.\ on $\ell_T$.
\end{lemma}
\begin{proof}
On $\{d\ge D_0\}$ we have $\mathrm{e}^{-\delta d}\le \mathrm{e}^{-\delta D_0/2}\mathrm{e}^{-\delta d/2}$. Next, $d(\Lambda)=\min_{\mathcal T}L_{\mathcal T}(\Lambda)$ over the spanning trees $\mathcal T$ on the $N$ labelled indices, where $L_{\mathcal T}$ is the sum of the cube distances over the edges of $\mathcal T$. Hence $\mathrm{e}^{-\delta d/2}=\max_{\mathcal T}\mathrm{e}^{-\delta L_{\mathcal T}/2}\le\sum_{\mathcal T}\mathrm{e}^{-\delta L_{\mathcal T}/2}$, and by Cayley's formula there are $N^{N-2}$ trees.
Fix a tree, root it at $\iota_0$, and sum the free cubes from the leaves inwards. Cube centres $c,c'$ satisfy $\dist(\square,\square')\ge(\abs{c-c'}_\infty-2)_+$, so each edge contributes at most $\sum_{c'}\mathrm{e}^{-(\delta/2)(\abs{c-c'}_\infty-2)_+}\le \mathrm{e}^{\delta}\sum_{w}\mathrm{e}^{-(\delta/2)\abs w_\infty}$. The minimal-image map $(\Z/n)^3\to\Z^3$ is injective, so this is at most the $\Z^3$ sum. Since $\#\{w\in\Z^3:\abs w_\infty=r\}=(2r+1)^3-(2r-1)^3=24r^2+2$, the $\Z^3$ sum is $\sigma(\delta)$. There are $N-1$ edges.
\end{proof}

\emph{The per-localisation bound (Lemma~\ref{app:G1:I145-perloc}).} Put $\delta:=\min(\delta_T,\delta_1)$, with $\delta_T>0$ the tree rate of \cite[(1.33)]{Ba3} as given by Proposition~\ref{app:G1:I145-prop1base} (it depends on $\nbar$; it is not Ba{\l}aban's ``$\delta_0$ depends on $d$ only'') and $\delta_1$ the rate of \cite[(2.10), p.~426]{Ba3}. Define
\[
A_k:=C_P(1+\abs u\bar c)^{\bar N_{\rm ph}}P(t_k),\qquad P(t):=\sum_{2\le\alpha+2\beta\le2\nbar}\abs{e(t)}^\alpha\lambda(t)^\beta,
\]
where $e(t)^2=e^2t$ and $\lambda(t)=\lambda t$, $\bar c:=\kappa\mu_*^2\eps_0^2L^2$, $\bar N_{\rm ph}$ is the largest number of photon lines, and $C_P$ is the $O(1)$ of \cite[(1.33)]{Ba3} times that of \cite[(2.10)--(2.11)]{Ba3}. Lemma~\ref{app:G1:I145-perloc} states that for every $(G,\tilde l)$ other than the order-$e^2$ family with one $D$-line, and every $\Lambda$,
\begin{equation}\label{app:G1:eq-perloc}
\abs{\mathcal V_{G,\tilde l}(\Lambda)}\le A_k\,\mathrm{e}^{-\delta d(\Lambda)}.
\end{equation}
The proof cuts the line of index $k$ as a two-variable external field \cite[p.~419, (2.10)--(2.11) p.~426]{Ba3}. It then uses the list of the resulting two-leg classes in the bare grading, in which counterterm vertices are single vertices and are never opened. The one-particle-irreducible scalar two-leg classes get their $\delta m^2$ partners, which are already inside $E_1$. The one-particle-reducible members of C2 (21 of the 59 non-vanishing classes) have no partner; they are handled by the product mechanism of Step~3 of the proof of Lemma~\ref{app:G1:I145-perloc}. The photon two-leg classes C3, C4 need the Ward--Takahashi class, valid on $T_\eps$ after the transfer of \cite[p.~433]{Ba3} (Proposition~\ref{app:G1:I2}). In every non-vanishing term the cut graph stays connected (scalar lines are never bridges, and photon bridges give vanishing pairs), so (1.33) of the cut graph applies. For $\abs u\le1$, the $u$-part is \eqref{app:G1:eq-perloc} with $(1+\abs u\bar c)^{\bar N_{\rm ph}}$ replaced by $\bar N_{\rm ph}\abs u\bar c(1+\bar c)^{\bar N_{\rm ph}-1}$. The order-$e^2$ family with one $D$-line is bounded per anchor cube in Lemma~\ref{app:G1:I1c-e2}.
Since every vacuum graph has $o=\alpha/2+\beta\ge1$, $\abs{e(t)}^\alpha\lambda(t)^\beta=\abs e^\alpha\lambda^\beta t\cdot t^{o-1}\le\abs e^\alpha\lambda^\beta\eps_0^{o-1}t$ for $t\le\eps_0\le1$. Hence
\[
P(t_k)\le C_e t_k,\qquad C_e:=\sum\abs e^\alpha\lambda^\beta\eps_0^{o-1},
\]
and the rate exponent is $1$, not the collar exponent $\kappa_0=\frac15$ used elsewhere.

\begin{proposition}[bad and outside localisations]\label{app:G1:R9vac}
Assume \eqref{app:G1:eq-perloc}. Put $\bar S:=\sum_{(G,\tilde l)}S_{N(G)}(\delta)$, the sum over the graph--ordering pairs of order $\le\nbar$; it depends on $\nbar$, $d$, $\delta$ only. At both points and every $0\le k\le K-1$:
\begin{enumerate}
\item[(i)] $\big|\sum_{\Lambda\ \rm bad}\mathcal V_k(\Lambda)\big|\le\bar N\bar SA_k\big[\abs{\mathcal L_k}+\mathrm{e}^{3\delta}\mathrm{e}^{-\delta r(t_k)/2}\abs{\mathcal D_k}\big]$;
\item[(ii)] for every union $X$ of scale-$k$ blocks of $W_k^c$, $\big|\sum_{\Lambda\ \rm outside,\ anchor\in X}\mathcal V_k(\Lambda)\big|\le\bar SA_k\abs X_k$;
\item[(iii)] (i) and (ii) hold for $\mathcal V_k(u)-\mathcal V_k(0)$ with $A_k$ replaced by $A_k^u\le\bar N_{\rm ph}\abs u\bar c(1+\bar c)^{\bar N_{\rm ph}-1}C_PP(t_k)$.
\end{enumerate}
The bounds do not depend on $\eps$, $K$ or $\ell_T$, they are linear in $\abs u$ with no sign used, and they are finite at $M=m_1$.
\end{proposition}
\noindent The proof uses \cite[(1.33)]{Ba3} in the form of Proposition~\ref{app:G1:I5} and \cite[(2.10)--(2.11), p.~426]{Ba3}, through \eqref{app:G1:eq-perloc}. Lemma~\ref{app:G1:VS} is a second form of the same localisation sum.
\begin{proof}
(i) There are at most $\bar N$ choices of the anchor index. A bad $\Lambda$ has a cube outside $W_k$, so its tree joins $\square_0$ to $W_k^c$ and $d(\Lambda)\ge\dist(\square_0,W_k^c)-3$. Apply \eqref{app:G1:eq-perloc} and Lemma~\ref{app:G1:treesum} with $D_0=(\dist(\square_0,W_k^c)-3)_+$. For $\square_0\in\mathcal L_k$ use $\mathrm{e}^{-\delta D_0/2}\le1$. For $\square_0\in\mathcal D_k$ we have $\dist>r(t_k)$, so $\mathrm{e}^{-\delta D_0/2}\le \mathrm{e}^{3\delta/2}\mathrm{e}^{-\delta r(t_k)/2}\le \mathrm{e}^{3\delta}\mathrm{e}^{-\delta r(t_k)/2}$. Then sum over $(G,\tilde l)$.
(ii) Apply \eqref{app:G1:eq-perloc} and Lemma~\ref{app:G1:treesum} with $D_0=0$. The anchor is fixed and the anchor cube runs over $X$; the condition $\Lambda\cap W_k=\emptyset$ only removes terms from an absolute sum.
(iii) The same argument, with the $u$-part of \eqref{app:G1:eq-perloc}. For uniformity: $\bar S$ and $\bar N$ contain no sum over steps; the torus count is at most the $\Z^3$ count; and $C_P$ does not depend on the domains \cite[p.~421]{Ba3}.
\end{proof}

\emph{The region-sum lemma (Lemma~\ref{app:G1:I145-comb}, under its hypothesis (H-LF)).} The large-field region sum of the inductive bound of the proof of Theorem~\ref{thm:T6}, multiplied by $\exp\big(c_A\sum_{k<K}(1+\log t_k^{-1})\abs{\Lambda_7^{(k)c}}_k\big)$, is $\le\exp\big(2(d+2)\Sigma_*(\eps_0;\pi)\abs T\big)$ for any $c_A\ge0$. The hypothesis is the upper bound (E.comb) on $\eps_0$: $\hat\ell^{4}\ge2c_AC_VR^d/(c_*b_0^2\ln L)$, with $\hat\ell=1+\log\eps_0^{-1}$. Here $\Sigma_*(\eps_0;\pi)=\sum_{n\ge n_0}L^{3n}\exp(-\frac12c_*b_0^2(1+n\log L)^{2p})$ contains no $c_A$.

\begin{lemma}[absorption of the outside pieces]\label{app:G1:IF1abs}
Put $C_{\rm old}:=\bar SC_PC_e(1+\bar c)^{\bar N_{\rm ph}}$. At both points,
\[
\sum_{k<K}\bar SA_k\abs{\Lambda_7^{(k)c}}_k\le C_{\rm old}\sum_{k<K}(1+\log t_k^{-1})\abs{\Lambda_7^{(k)c}}_k.
\]
Hence, under (E.comb) read with $c_A+C_{\rm old}$, the outside pieces cost at most $\exp(2(d+2)\Sigma_*(\eps_0;\pi)\abs T)$ in the region sum. $C_{\rm old}$ does not depend on $\eps$, $K$, $k$, $\ell_T$, $b_0$, $R$ or $p$, and it is fixed at stage 10a of Appendix~\ref{app:B} with $\bar c$ and $C_e$ taken at their suprema over $(0,\eps_1]$. Along the collars the outside pieces cost $2C_{\rm old}\eps_0$ per block of $\mathrm{Col}_i$.
\end{lemma}

\begin{proof}
The outside-anchored set is $X=W_k^c=\mathrm{Col}_k\sqcup\mathrm{Old}_k=\Lambda_7^{(k)c}$. Proposition~\ref{app:G1:R9vac}(ii) gives $\bar SA_k\abs{\Lambda_7^{(k)c}}_k$. For $\abs u\le1$ we have $A_k\le C_P(1+\bar c)^{\bar N_{\rm ph}}C_et_k$, and $t_k\le\eps_0\le1\le1+\log t_k^{-1}$. This gives the display. Lemma~\ref{app:G1:I145-comb} with $c_A\mapsto c_A+C_{\rm old}$ gives the cost bound.
For the collar form: every scale-$k$ block of $\mathrm{Old}_k$ meets $\bigcup_{i<k}\mathrm{Col}_i$, and the blocks are dyadically nested, so $\abs{\mathrm{Old}_k}_k\le\sum_{i<k}\abs{\mathrm{Col}_i}_i$. Moreover $\sum_{k=i+1}^Kt_k\le\eps_0\sum_{n\ge0}L^{-n}=\frac L{L-1}\eps_0\le2\eps_0$.
\end{proof}

\begin{proposition}[the moved vacuum pieces]\label{app:G1:I1}
Assume convention (G) and $R\ge5$, and let $\mathcal V_k$ be as in Definition~\ref{app:G1:vacgrade}. Let $1\le k\le K-1$, and work at either of the points $\pi_\pm$. The moved set $\mathcal V_k(u)$ is localised once, at its own scale $k$. Each localisation is exactly one of good, bad or outside, and has one anchor, by the anchor rule of the paragraph containing Lemma~\ref{app:G1:treesum}. Write $\mathcal V_k=\mathcal V_k^{\rm good}+\mathcal V_k^{\rm bad}+\mathcal V_k^{\rm out}$ accordingly. Then:
\begin{enumerate}
\item[(i)] \emph{good:} $\mathcal V_k^{\rm good}$ is a sum of terms with all cubes in $W_k$. It equals the generated step-$k$ vacuum family and cancels against it.
\item[(ii)] \emph{bad:} $\abs{\mathcal V_k^{\rm bad}}\le\big(\bar N\bar SA_k+4C_{\rm loc}t_k\big)\abs{\mathcal L_k}_k+\big(\bar N\bar SA_k\mathrm{e}^{3\delta}\mathrm{e}^{-\delta r(t_k)/2}+2C_{\rm dp}t_k^3\mathrm{e}^{-2\delta_1(r(t_k)-4)}\big)\abs{\mathcal D_k}_k$.
\item[(iii)] \emph{outside:} for every union $X$ of scale-$k$ blocks of $W_k^c$, the part of $\mathcal V_k^{\rm out}$ with anchor in $X$ is at most $(\bar SA_k+2C_{\rm loc}t_k)\abs X_k$. In particular $\abs{\mathcal V_k^{\rm out}}\le C'_{\rm old}t_k\abs{\Lambda_7^{(k)c}}_k$, where $C'_{\rm old}:=C_{\rm old}+2C_{\rm loc}$. Suppose that the hypothesis (H-LF) of Lemma~\ref{app:G1:I145-comb} holds, and that condition (E.comb) holds with $c_A$ replaced by $c_A+C'_{\rm old}$. Then the factor $\exp\big(\sum_{k}\abs{\mathcal V_k^{\rm out}}\big)$ costs at most $\exp\big(2(d+2)\Sigma_*(\eps_0;\pi)\abs T\big)$ in the region sum of that lemma.
\item[(iv)] \emph{$u$-parts:} (ii) and (iii) hold for $\mathcal V_k(u)-\mathcal V_k(0)$, with $A_k$ replaced by $A_k^u$ of Proposition~\ref{app:G1:R9vac}(iii).
\item[(v)] No term is moved twice. $E_{\rm fin}(u)$ lies in no $\mathcal V_k$, and it remains a global factor that is paid once.
\item[(vi)] The sum over $k$ of the $\abs{\mathcal D_k}_k$-terms of (ii) is at most $C\abs T$, uniformly in $K$.
\end{enumerate}
Here $\delta$, $\delta_1$ and $A_k$ are as in Lemma~\ref{app:G1:I145-perloc}, and $\bar N$, $\bar S$ as in Proposition~\ref{app:G1:R9vac}. $C_{\rm loc}$ and $C_{\rm dp}$ are the constants of Theorem~\ref{app:G1:B5e2-thm} with $R_0=1$, $C_{\rm old}$ is the constant of Lemma~\ref{app:G1:IF1abs}, and $\Sigma_*$, $c_A$ are as in Lemma~\ref{app:G1:I145-comb}. All constants are bounded uniformly in $\eps$, $K$ and $\ell_T\ge\ell_0$, and they are finite at $u=-\eta$ and $M=m_1$.
\end{proposition}
\noindent The proof uses Lemma~\ref{app:G1:I145-perloc}, which rests on Proposition~\ref{app:G1:I145-prop1base} and on Proposition~\ref{app:G1:I2}. It also uses \cite[p.~433]{Ba3}, whose remainder on the cube $\square$ is treated in Proposition~\ref{app:G1:I3}, and Theorem~\ref{app:G1:B5e2-thm}. Throughout, convention (G) is in force.
\begin{proof}
(v) By Lemma~\ref{app:G1:vacgrade-id}(i), the sets $\mathcal V_k$ partition the terms of $E_{\rm vac}-E_{\rm fin}$. A term lies in the set $\mathcal V_k$ whose $k$ is the largest index of its lines, so the telescoping $E_K=E_{K,k}+\sum_{j<k}\mathcal V_j$ moves each term once, at its top scale. The terms of $E_{\rm fin}$ have a line of index $K$ \cite[(1.26), p.~418]{Ba3}, so they lie in no $\mathcal V_k$. $E_{\rm fin}(u)$ is bounded once, as a global factor: by Proposition~\ref{app:G1:Efin} at $u=0$, and by Proposition~\ref{app:K:L04-4} for $u\ne0$.

The three classes partition the localisations, and each localisation has one anchor, by the anchor rule.

(i) A good localisation has all its cubes in $W_k$. By Lemma~\ref{app:G1:vacgrade-id}(ii), $\mathcal V_k$ is the field-independent part of order $\le\nbar$ of the output of step $k$. The cancellation of $\mathcal V_k^{\rm good}$ against the vacuum output generated at step $k$ is Lemma~\ref{app:G1:S05-TS} together with item (C11) of Lemma~\ref{app:G1:S05-RF}. It is used in the closing step, Lemma~\ref{app:G1:S05-CL}.

(ii)--(iv) Consider first all graph--ordering pairs other than the order-$e^2$ family with one $D$-line. For them, Lemma~\ref{app:G1:I145-perloc} gives \eqref{app:G1:eq-perloc}, so Proposition~\ref{app:G1:R9vac}(i)--(iii) apply, since their proofs use only \eqref{app:G1:eq-perloc} and Lemma~\ref{app:G1:treesum}. The order-$e^2$ family with one $D$-line contributes the terms with $C_{\rm loc}$ and $C_{\rm dp}$, by Lemma~\ref{app:G1:I1c-e2}. That lemma applies with $R_0=1$, because $r(t)=R(1+\log t^{-1})^2\ge R\ge5$ for $t\le1$. For the second claim of (iii), take $X=W_k^c=\Lambda_7^{(k)c}$, and use $\bar SA_k\le C_{\rm old}t_k$ for $\abs u\le1$ (Lemma~\ref{app:G1:IF1abs}). For the absorption, use $t_k\le\eps_0\le1\le1+\log t_k^{-1}$. This gives $\sum_k\abs{\mathcal V_k^{\rm out}}\le C'_{\rm old}\sum_k(1+\log t_k^{-1})\abs{\Lambda_7^{(k)c}}_k$, which is the cost term of (H-LF) with $c_A$ replaced by $c_A+C'_{\rm old}$. Lemma~\ref{app:G1:I145-comb} with $c_A+C'_{\rm old}$ then gives the stated cost.

(vi) We have $\abs{\mathcal D_k}_k\le\abs Tt_k^{-3}$ and $A_k\le Ct_k$. Since $r(t)=R(1+\log t^{-1})^2$, for every $q>0$ there is a constant $C_q$ with $\mathrm{e}^{-\delta r(t)/2}+\mathrm{e}^{-2\delta_1(r(t)-4)}\le C_qt^q$. Take $q=3$. The deep sum is then at most $C\abs T\sum_kt_k\le2C\eps_0\abs T$.

Uniformity: $\bar S$ and $\bar N$ contain no sum over steps. On the torus, the count of Lemma~\ref{app:G1:treesum} is at most the count on $\Z^3$. $C_P$ is independent of the domains \cite[p.~421]{Ba3}. Finally, $C_{\rm loc}$ and $C_{\rm dp}$ are bounded uniformly in $\eps$, $k$ and $\ell_T\ge\ell_0$ by Theorem~\ref{app:G1:B5e2-thm}.
\end{proof}

\paragraph{The Ward--Takahashi step on the torus.}\label{app:G1:ssec-I2}
\begin{lemma}[a constant background is a pure gauge on a non-wrapping set]\label{app:G1:gaugetriv}
Let $\Omega\subset T_\eps$ be a set of sites with the bonds between them. Assume that $\Omega$ lifts isometrically: there is $\iota:\Omega\to\eps\Z^3$ with $\pi\circ\iota=\mathrm{id}$ and $\iota(x+\eps e_\mu)=\iota(x)+\eps e_\mu$ on every bond of $\Omega$. This holds if $\Omega$ lies in a box of side $<\ell_T$. Then for every $a\in\R^3$ the constant background on the bonds of $\Omega$ equals $\nabla\chi$ with $\chi:=a\cdot\iota$. Hence:
\begin{itemize}
\item every gauge-covariant quantity built only from fields and operators on $\Omega$ is independent of $a$, up to phases $\mathrm{e}^{\pm ie\chi}$ on its external scalar legs;
\item the localised zeroth-moment identity $\sum_{b'\subset\Omega}\Pi^{\rm loc}_\Omega(b,b')=0$ holds exactly on $T_\eps$ for every class supported in $\Omega$.
\end{itemize}
\end{lemma}
\noindent This is Ba{\l}aban's gauge-out of the constant field $B_0$ \cite[p.~433]{Ba3}, written for the torus.
\begin{proof}
$\chi$ is well defined because $\iota$ is injective, and $\nabla_\mu\chi(x)=a\cdot(\iota(x+\eps e_\mu)-\iota(x))/\eps=a_\mu$ on the bonds of $\Omega$. The gauge transformation $\varphi\mapsto \mathrm{e}^{ie\chi}\varphi$, $A\mapsto A-\nabla\chi$ on $\Omega$ removes $a$ from every object supported in $\Omega$. For a class supported in $\Omega$ the second $a$-derivative is therefore $0$. The second $a$-derivative of a localised two-photon class at zero momentum is $\sum_{b'\subset\Omega}\Pi^{\rm loc}(b,b')$.
\end{proof}

\begin{lemma}[the cube of {\cite[p.~433]{Ba3}} does not wrap]\label{app:G1:nowrap}
Assume the conditions (T.nw), $\ell_0\ge\eps_0[8r(\eps_0)+4M_b+2]$, and (E.r), $\eps_0<\mathrm{e}^{1-\varrho}=\mathrm{e}^{-1}$, of Appendix~\ref{app:B}, and let $\ell_T\ge\ell_0$. Then for every $k$ with $t_k\le\eps_0$, a cube of side $3r(t_k)$ in scale-$k$ units (physical side $3r(t_k)t_k$) lies in a box of side $<\ell_T$. In particular it lifts isometrically.
\end{lemma}
\begin{proof}
With $s=\log t^{-1}$, $\frac{d}{dt}\big[tR(1+s)^2\big]=R(\log^2t-1)>0$ for $t<\mathrm{e}^{-1}$. So $t\mapsto t\,r(t)$ increases on $(0,\eps_0]$, and $3t_kr(t_k)\le3\eps_0r(\eps_0)<8\eps_0r(\eps_0)\le\ell_0\le\ell_T$.
\end{proof}

\begin{proposition}[the Ward--Takahashi step on the torus]\label{app:G1:I2}
At every scale with $t_k\le\eps_0$, localise each renormalised class as in \cite[p.~433]{Ba3}: in a cube $\square_1$ of size $r(L^k\eps)$ containing its $O(\nbar)$-cluster, inside a cube $\square$ of size $3r(L^k\eps)$. Then, on $T_\eps$:
\begin{itemize}
\item the gauge-out of the constant field $B_0$ and the replacement $G_k(\square,0)=G_k(0)+\delta G_k(\square,\eta\Z^d,0)$ are valid;
\item every Ward--Takahashi identity used for the photon-cut classes C3, C4 holds exactly, with no torus image defect;
\item the torus enters only through the external fields $\delta G_k$ that \cite[(1.33)]{Ba3} already carries.
\end{itemize}
Ba{\l}aban's renormalisation \cite[pp.~434--445]{Ba3} and the vacuum case of Proposition~\ref{app:G1:R9vac} use the same torus input.
\end{proposition}
\noindent The proposition uses \cite[p.~433]{Ba3} and the condition (T.nw).
\begin{proof}
By Lemma~\ref{app:G1:nowrap}, $\square$ lifts isometrically. By Lemma~\ref{app:G1:gaugetriv}, $B_0$ is a gradient on $\square$, and its removal is the gauge transformation of \cite[p.~433]{Ba3}, after which ``we get the same expressions as above with $B_0=0$ only''. Since $\square$ lifts, it can be identified with a cube of $\eta\Z^d$. There the Ward--Takahashi identities of free boundary conditions hold exactly (the step S8 of Table~\ref{tab:prop1}), and the difference $\delta G_k(\square,\eta\Z^d,0)$ is treated as an external field, as in the source. So no Ward--Takahashi step runs on a wrapping domain.
\end{proof}

\paragraph{Only real gauge backgrounds are used.}\label{app:G1:ssec-I4}
A \emph{gauge background} is any vector-field or group configuration at which one of the following is evaluated: an operator $G_k(\Omega,\cdot)$, $C_k(\Omega,\cdot)$, a covariant Laplacian, a transporter $U(\cdot)$, a determinant or Gaussian measure, or an expansion in the background. It is \emph{real} if $A\in\R$ on every bond, equivalently $U(A)\in\SO(N)$. The \emph{allowed} complex objects are: the entire function $R_{\nbar+1}(z)$ evaluated at real $A$; momentum continuations at $A=0$ \cite[p.~586]{Ba4}; complex sources and Cauchy estimates in sources or in $u$; the complex structure of the charged scalar.

\begin{definition}[gauge backgrounds]\label{app:G1:I145-defbg}
In Definition~\ref{app:G1:I145-defbg}, Lemmas~\ref{app:G1:I145-closure}, \ref{app:G1:I145-realbg}, \ref{app:G1:I145-unitary} and Proposition~\ref{app:G1:I4}, $N=2$, $q$ is the real antisymmetric $2\times2$ matrix with $q^2=-1$ of \cite[(1.2), p.~572]{Ba4}, $e>0$ is the charge, $\eta=L^{-k}$ and $U(A)=\mathrm{e}^{qe\eta A}$. Vector fields are real-valued functions on bonds \cite[p.~572]{Ba4}. Scalar fields take values in $\R^2$, and $q$ acts on $\R^2$ as the complex structure.
\begin{itemize}
\item[(a)] A \emph{gauge background} is a configuration $A$ of complex numbers on the bonds of a subset of $T_\eta$, or the family $\{U(A_b)\}_b$ it defines, at which one of the following is evaluated: the transporters $U(A(\Gamma))$; $Q_k(A)$ and $P_k(A)=Q_k^*(A)Q_k(A)$ \cite[(1.4)--(1.5)]{Ba4}; the Neumann covariant Laplacian $-\Delta^{\eta,N}_{A,\Omega}$ \cite[(1.3)]{Ba4}; $G_k(\Omega,A)=(-\Delta^{\eta,N}_{A,\Omega}+m^2+aP_k(A))^{-1}$ \cite[(1.6)]{Ba4} and $C_k(\Omega,A)$ \cite[p.~586]{Ba4}; any operator obtained from these by sums, products, restrictions and inverses; a determinant or Gaussian measure defined by such an operator; or an expansion in powers of $A-B$ around a background $B$.
\item[(b)] A gauge background is \emph{real} if $A_b\in\R$ for every bond $b$.
\item[(c)] The following complex quantities are not gauge backgrounds: the value of the entire function $R_{n+1}(z):=(n+1)\int_0^1(1-t)^n\mathrm{e}^{tz}\,dt$ \cite[(3.14), p.~614]{Ba1} at a matrix $qe\eta A$ with $A$ real; complex momenta in Fourier representations of operators at $A=0$ or of Feynman integrands; complex test functions and complex parameters multiplying sources; the complex structure $q$.
\end{itemize}
\end{definition}

\begin{lemma}[realness is preserved by the operations of the construction]\label{app:G1:I145-closure}
Let $\Omega\subset T_\eta$ be a union of blocks and let $A,A_1,A_2$ be real vector fields.
\begin{itemize}
\item[(a)] $U(A(\Gamma))$, $Q_k(A)$ and $P_k(A)$ are real matrices. $X_\Omega(A):=-\Delta^{\eta,N}_{A,\Omega}+m^2+aP_k(A)$ is real symmetric and nonnegative, and whenever it is invertible (in particular under \cite[(1.8), p.~573]{Ba4}) $G_k(\Omega,A)$ is real symmetric. The same holds for every operator obtained from these by sums, products, restrictions to subsets of $\Omega$, adjoints and inverses, in particular for $C_k(\Omega,A)$.
\item[(b)] The following are real vector fields: $TA$ for a linear map $T$ with real matrix (for instance $Q^*$, $C^{(0)}Q^*$, block averages $\bar A_b=L^{-k}\sum_{\langle x,x'\rangle\subset b}A_{\langle x,x'\rangle}$); $(1-\theta)A_1+\theta A_2$ for real $\theta$; $A_1+A_2$; restrictions and extensions by zero; $A+\partial^\eta\lambda$ for real $\lambda$. Contour sums $A(\Gamma)$ are real.
\item[(c)] Let $\square$ be a set of sites, $\tilde B$ a real vector field on $\square$ and $B_0\in\C^d$ a constant field with $\max_\mu\sup_{x\in\square}\abs{\tilde B_\mu(x)-B_{0,\mu}}\le\delta$. Then $\rRe B_0\in\R^d$ satisfies the same inequality. If $\square$ does not wrap around the torus, that is, if it has an isometric lift $\iota$ to a subset of $\eta\Z^d$, the constant field $\rRe B_0$ equals $\partial^\eta\lambda$ on $\square$ for the real function $\lambda(x)=\sum_\mu(\rRe B_{0,\mu})\iota(x)_\mu$.
\item[(d)] If $P$ is a real symmetric positive definite operator on $(\R^n)^\Lambda$, $N(0,P^{-1})$ is a probability measure on $(\R^n)^\Lambda$. For the slice model at real $u>-1/\bar\delta$ the photon covariance $C_B(u)=C_{\mu_*}+u\Gamma^B$ of Subsection~\ref{ssec:T6-statement} is real symmetric and positive definite; this covers $u\in\{1,-\eta\}$, since $\eta\bar\delta<1$ (Lemma~\ref{lem:T6-rows}(a)). The weights $1+u\delta_j$ are real and positive there.
\end{itemize}
\end{lemma}
\begin{proof}
(a) $U(A(\Gamma))$ is real by Lemma~\ref{app:G1:I145-unitary}(c); hence so is $Q_k(A)$, a sum of terms $\eta^dU(A(\Gamma))$ \cite[(1.4)]{Ba4}, and $P_k(A)=Q_k^*Q_k$. The quadratic form $\sum_{b\subset\Omega}\eta^{d-2}\abs{U(A_b)\phi(b_+)-\phi(b_-)}^2$ of \cite[(1.3)]{Ba4} is real and nonnegative on real $\phi:\Omega\to\R^N$, so $-\Delta^{\eta,N}_{A,\Omega}$ and $X_\Omega(A)$ ($m^2\ge0$, $a>0$) are real symmetric and nonnegative. Sums, products, restrictions and adjoints of real matrices are real, and the inverse of an invertible real symmetric matrix is real symmetric.
(b) Each operation is a real-linear map with real coefficients.
(c) For real $t$ and complex $c$, $\abs{t-\rRe c}=\abs{\rRe(t-c)}\le\abs{t-c}$; apply this with $t=\tilde B_\mu(x)$ and $c=B_{0,\mu}$. With the lift $\iota$, $\partial^\eta_\mu\lambda=\rRe B_{0,\mu}$ on the bonds of $\square$ (Lemma~\ref{app:G1:gaugetriv}). A set contained in a box of side smaller than the side of the torus has such a lift (Lemma~\ref{app:G1:gaugetriv}), and so does the cube of \cite[p.~433]{Ba3} (Lemma~\ref{app:G1:nowrap}).
(d) $N(0,P^{-1})$ has the density $(\det P/(2\pi)^{n\abs\Lambda})^{1/2}\mathrm{e}^{-\frac12\langle z,Pz\rangle}$. $C_{\mu_*}$ and $\Gamma^B$ are real symmetric with $0\le\Gamma^B\le\bar\delta C_{\mu_*}$, so $C_B(u)\ge(1+\min(u,0)\bar\delta)C_{\mu_*}>0$ for $u>-1/\bar\delta$. Finally $1+u\delta_j\ge1-\eta\bar\delta>0$ at $u=-\eta$, since $0\le\delta_j\le\bar\delta$.
\end{proof}

\begin{lemma}[the backgrounds of the construction are real]\label{app:G1:I145-realbg}
Consider the inductive construction of \cite[Sects.~2--3]{Ba2} on the torus for the slice model at $(u,M)\in\{(1,m_2),(-\eta,m_1)\}$, as used in the proof of Theorem~\ref{thm:T6} (Subsection~\ref{app:G1:ssec-assembly}). At every step $k$ the following are real vector fields:
\begin{itemize}
\item[(1)] the photon field $A$, the fluctuation fields $A'$, $A'_j$, and the block-spin fields $B$, $B^{(k)}$;
\item[(2)] the background fields $B^{(1)}=aL^{-2}C^{(0)}Q^*B$ \cite[(3.10), p.~614]{Ba1}, their analogues $A^{(k)}$, $B^{(k+1)}$, $B^{(k+1),\eta}$ at general $k$, and the averages $\bar A^{(k)}$ \cite[p.~570]{Ba2};
\item[(3)] the interpolated fields $\tilde B=(1-\theta_{k+1})A^{(k)}+\theta_{k+1}B^{(k+1),\eta}$ \cite[p.~579]{Ba2};
\item[(4)] the arguments $\tilde A+\tilde B$ of \cite[(1.16), p.~414]{Ba3}, and $B'$, $B_0$ of the decomposition $\tilde B=B_0+B'$ \cite[p.~433]{Ba3}, with $B_0$ chosen real as in Lemma~\ref{app:G1:I145-closure}(c);
\item[(5)] the fields obtained from these by the gauge transformation that removes $B_0$ \cite[p.~433]{Ba3}, and the field $A^\eps$ of \cite[(3.37), p.~591]{Ba2}.
\end{itemize}
Every operator of Definition~\ref{app:G1:I145-defbg}(a) that the construction evaluates is evaluated at one of these fields or at $A=0$. Every Gaussian measure of the construction is $N(0,P^{-1})$ on a real space with $P$ real symmetric positive definite; for instance $d\mu_{C^{(0)}}(A')\,d\mu_{C^{(0)}(B^{(1)})}(\phi')$ \cite[(3.22), p.~616]{Ba1} and the measure of a positive operator in \cite[(2.28)--(2.31), p.~563]{Ba2}.
\end{lemma}
\begin{proof}
Induction on $k$. At the start the integration variables are the photon field, with law $N(0,C_B(u))$ and $C_B(u)$ real symmetric positive definite (Lemma~\ref{app:G1:I145-closure}(d)), and the $\R^2$-valued scalar field; the block-spin fields are their real block averages. If all fields of step $k-1$ are real, the fields of step $k$ arise by: the translations \cite[(3.10), (3.18), p.~614]{Ba1}, \cite[(2.56)]{Ba2}, which apply operators of Lemma~\ref{app:G1:I145-closure}(a) at real fields to real fields; Gaussian integrations with real symmetric positive covariances \cite[(3.11), (3.19)]{Ba1}, \cite[(2.28)]{Ba2}, where the weights $1+u\delta_j>0$ change only the photon covariances (Lemma~\ref{app:G1:I145-closure}(d)); averages and the interpolation with the real cut-off $\theta_{k+1}$ (Lemma~\ref{app:G1:I145-closure}(b)); the expansions \cite[(3.14)--(3.16)]{Ba1}, \cite[(1.15)--(1.16)]{Ba3} in a real field ($A'$, $\tilde A$ or $B'$) around a real field; the constant approximant and the gauge transformation of \cite[p.~433]{Ba3} (Lemma~\ref{app:G1:I145-closure}(c)). In \cite[(3.37)]{Ba2}, $A^\eps$ is the variable of the outer real integral $\int dB$, real by (1). This proves (1)--(5). The last assertion holds because \cite{Ba1,Ba2,Ba3,Ba4} define the operators of Definition~\ref{app:G1:I145-defbg}(a) only for vector fields, i.e.\ real configurations \cite[p.~572]{Ba4}; the only evaluations at other arguments are Fourier representations at $A=0$ (e.g.\ the contour shift in the proof of \cite[Lemma~2.4, p.~586]{Ba4}) and $R_{\nbar+1}$ at real $A$ \cite[p.~615]{Ba1}, \cite[p.~414]{Ba3}, both of the kind in Definition~\ref{app:G1:I145-defbg}(c).
\end{proof}

\begin{proposition}[only real backgrounds]\label{app:G1:I4}
In the proof of Theorem~\ref{thm:T6}:
\begin{enumerate}
\item[(i)] every gauge background at which the construction, or a statement of this appendix used in the proof, evaluates an operator of Definition~\ref{app:G1:I145-defbg}(a) is real;
\item[(ii)] every estimate taken from the literature is stated for real backgrounds, for orthogonal gauge configurations, or for real Gaussian variables and real polynomial interactions, and is applied at a real background. These estimates are: the propagator bounds of \cite{Ba4} and \cite[Props.~2.1--2.3]{Ba1}; \cite[(1.16), p.~414; (1.32)--(1.33), p.~420]{Ba3}; $\norm{R_{\nbar+1}(qA)}\le1$ \cite[p.~615]{Ba1}; the diamagnetic inequality \cite[(3.17), p.~257]{BFS1}; and the cumulant lemma \cite[p.~152]{BCGNOPS} as applied in \cite[(3.22)--(3.24), p.~616]{Ba1} and \cite[(2.28)--(2.31), p.~563]{Ba2};
\item[(iii)] the printed proofs of these estimates contain no non-real gauge background; their only complex quantities are those of Definition~\ref{app:G1:I145-defbg}(c);
\item[(iv)] the two steps of \cite{Ba3} whose proofs are not printed, the bound on the operator (1.16) on the Neumann domain $\Omega$ \cite[p.~414]{Ba3} and the cube remainder \cite[p.~433]{Ba3}, are stated at real backgrounds, and their proofs in Lemmas~\ref{app:G1:I3r-U1mOmega}, \ref{app:G1:I3r-SEP} and \ref{app:G1:I3r-cube} use real backgrounds only. The second presupposes the inclusion of sets $\square\subset\Omega$ (Lemma~\ref{app:G1:I3r-Wsq}) and the replacement of $G_k(\Omega_2,\tilde B)$ (Lemma~\ref{app:G1:I3r-W2R}), a kernel identity at the real background $\tilde B$ that does not use $\square\subset\Omega_2$.
\end{enumerate}
In particular no bound on $G_k(\Omega,A)$, $C_k(\Omega,A)$ or an operator built from them is used at a non-real $A$, near a Neumann face of $\Omega$ or elsewhere, and no coupling constant takes a non-real value.
\end{proposition}
\noindent Part (iv) uses the cube inclusion $\square\subset\Omega$ and, in place of the inclusion $\square\subset\Omega_2$, the replacement of Lemma~\ref{app:G1:I3r-W2R}; both are parts of Proposition~\ref{app:G1:I3}. A second proof is given after Lemma~\ref{app:G1:I145-unitary}.
\begin{proof}
\emph{(i).} The construction is covered by Lemma~\ref{app:G1:I145-realbg}. The statements of this appendix used in the proof of Theorem~\ref{thm:T6} (Subsection~\ref{app:G1:ssec-assembly}) concern the objects of that lemma; for instance Lemma~\ref{app:G1:I3r-SEP} is stated for $A\in\{B_0,B_0+B'\}$ in the class of \cite[Prop.~2.1]{Ba1}, and Lemma~\ref{app:G1:I3r-U1mOmega} for the backgrounds $\tilde A$, $\tilde B$ of \cite[p.~414]{Ba3}.

\emph{(ii).} \cite{Ba4} defines its operators for real vector fields \cite[p.~572]{Ba4}, and \cite[Props.~2.1--2.3]{Ba1} are stated for regular real vector fields. In \cite[p.~414]{Ba3} the operator (1.16) is built from $G_k(\Omega,\tilde B)$ and $G_k(\Omega,\tilde A+\tilde B)$, with $\tilde A$ a sum of the real Gaussian fields $A'_j$, and the norms (1.32) are H\"older norms of kernels at these backgrounds \cite[p.~420]{Ba3}. The inequality $\norm{R_{\nbar+1}(qA)}\le1$ ``holds for arbitrary real $A$'' \cite[p.~615]{Ba1} and fails for non-real $A$ (Lemma~\ref{app:G1:I145-unitary}(b)). \cite[(3.17)]{BFS1} is stated for configurations of a compact gauge group and proved from \cite[(3.7), p.~255]{BFS1}, whose proof uses $\norm{U^\Phi(g)}=1$; in \cite[p.~591]{Ba2} it is applied to (3.37) at the real field $A^\eps$ (Lemma~\ref{app:G1:I145-realbg}(5)), where $U(A^\eps)$ is orthogonal (Lemma~\ref{app:G1:I145-unitary}(a)). The lemma of \cite[p.~152]{BCGNOPS} concerns real Gaussian variables and a real polynomial; in \cite[(3.22)--(3.24)]{Ba1} it is applied to $d\mu_{C^{(0)}}(A')\,d\mu_{C^{(0)}(B^{(1)})}(\phi')$ and a polynomial with real coefficients at the real background $B^{(1)}$, and in \cite[(2.28)--(2.31)]{Ba2} to ``a probabilistic Gaussian measure'' of ``a positive operator''; in each case $\mu=N(0,P^{-1})$ on $(\R^n)^\Lambda$ with $P$ real symmetric positive. At the pages of \cite{K1} that are used (pp.~653--657, 659) operators are evaluated at real fields, and the only complex quantities are sources. By (i), each estimate is applied at a real background.

\emph{(iii).} In \cite{Ba4} the only complex quantity is the contour shift in the proof of Lemma~2.4 \cite[p.~586]{Ba4}, at $A=0$. In \cite{Ba1,Ba2,Ba3} the expansions in the background are in real fields around real fields (Lemma~\ref{app:G1:I145-realbg}), their remainders are controlled by $\norm{R_{\nbar+1}(qA)}\le1$ at real $A$, and determinants are expanded through $\log(1+\lambda)$ for real eigenvalues $\lambda$ of real symmetric operators. The proof of \cite[(3.7)]{BFS1} is the path expansion \cite[(3.6)]{BFS1} with $\norm{U^\Phi(g)}=1$. The proof of the lemma of \cite{BCGNOPS} uses real Gaussian covariances at real momenta \cite[App.~C, pp.~164--165]{BCGNOPS}.

\emph{(iv).} At \cite[p.~414]{Ba3} the backgrounds $\tilde B$, $\tilde A+\tilde B$ are real. At \cite[p.~433]{Ba3}, $\tilde B=B_0+B'$ with $\tilde B$ real; by Lemma~\ref{app:G1:I145-closure}(c), $B_0$ may be taken real, and then $B'$ is real. The proof of Lemma~\ref{app:G1:I3r-U1mOmega} uses \cite[Prop.~2.1]{Ba1} and the real-background Lemma~\ref{app:G1:I3r-U1m}. The proofs of Lemmas~\ref{app:G1:I3r-SEP} and \ref{app:G1:I3r-cube} use \cite[Prop.~2.1]{Ba1}, \cite[Lemma~2.1, Cor.~2.3, (2.30), pp.~580--581]{Ba4} and the identity $D_{B_0}=U(B')^{-1}D_A-\phi$, where $A=B_0+B'$ with $B_0,B'$ real and $\phi$ is a same-site multiplication operator. The decomposition $G_k(\Omega,\tilde B)=G_k(\square,\tilde B)+\delta G_k(\Omega,\square,\tilde B)$ needs $\square\subset\Omega$ (Lemma~\ref{app:G1:I3r-Wsq}); the replacement of $G_k(\Omega_2,\tilde B)$ is Lemma~\ref{app:G1:I3r-W2R}.

\emph{Last assertion.} By (i) and (ii), no pointwise or derivative bound on $G_k$ or $C_k$ is used at a non-real argument. The proofs written in this appendix that enter Theorem~\ref{thm:T6} contain only these complex quantities: momentum shifts in Feynman integrands at no gauge background (Lemma~\ref{app:G1:CNtorus}), $R_{\nbar+1}$ at real arguments, and $q$. The Taylor coefficients in $(e,\lambda)$ that enter are derivatives at real couplings, and Theorem~\ref{thm:T6} is evaluated at the two real points $u\in\{1,-\eta\}$.
\end{proof}

\begin{remark}[scope]\label{app:G1:I145-I4scope}
(a) Some results used in the proof of Theorem~\ref{thm:T6} rest on statements whose proofs are not given in this paper (the statements of Appendix~\ref{app:K}). They are used through their statements, which concern the real objects of Lemma~\ref{app:G1:I145-realbg}, so (i) and (ii) of Proposition~\ref{app:G1:I4} apply to them. Proposition~\ref{app:G1:I4} makes no assertion about proofs that are not written here.
(b) Two complex-analytic devices are used in no proof written here or printed in \cite{Ba1,Ba2,Ba3,Ba4}: Cauchy estimates in the gauge field, which would need bounds on $G_k(\Omega,A)$ for non-real $A$, and bounds at a complex charge $ze$ with $\abs z=1$. By Lemma~\ref{app:G1:I145-unitary}(a),(b), such bounds would not follow from \cite{Ba4}, from \cite[(3.17)]{BFS1} or from $\norm{R_{\nbar+1}(qA)}\le1$.
\end{remark}

\paragraph{Proposition~1 at the order cap $\Nbar=169$.}\label{app:G1:ssec-I5}
\emph{Notation.} Fix a step $k\to k+1$, $0\le k\le K-1$, and write $t:=t_k=L^k\eps\le\eps_0$, $e(t)^2=e^2t$, $\lambda(t)=\lambda t$, $p(t)=b_0(1+\log t^{-1})^{p}$ with $p=6$, $b_0\ge\sqrt2$, and $\mathsf d=3$. The \emph{coupling order} of a vertex $v$ is $q_v:=d_v(v)+d_s(v)$, where $d_v(v)$, $d_s(v)$ are its orders in $e$ and in $\lambda$ \cite[p.~420]{Ba3}; the order of a graph or family is $q:=\sum_vq_v$. Couplings inside the covariant kernels $U(\tilde B)$, $G_k(\Omega,\tilde B)$ are not counted. The fields satisfy the small-field restrictions of Definition~\ref{app:G1:VSR-IH}. Since $p(t_{k-1})\le(1+\log L)^pp(t_k)$ and $\lambda(t_{k-1})^{-1/4}=L^{1/4}\lambda(t_k)^{-1/4}$, there is $C_1=C_1(c_1,L,p)$ with
\begin{equation}\label{app:G1:I5-legs}
\norm{h\Phi_{\rm ext}}_{1,\alpha_0}\le C_1p(t)\lambda(t)^{-1/4},\qquad \norm{h'A_{\rm ext}}_{1,\alpha_0}\le C_1p(t)
\end{equation}
for each scalar and each vector leg; for the external scalar fields $a_kG_k(\Omega,\tilde B)Q_k^*(\tilde B)\phi$ the covariant derivative is $\le C_1p(t)$ (Lemma~\ref{lem:lem24}). Put $P(t):=p(t)^4(1+\log t^{-1})\ge1$, let $N_k:=\abs{T_\eps}/t_k^3$, and let $r(t)=R(1+\log t^{-1})^2$, $\hat\ell=1+\log\eps_0^{-1}$ be as in Appendix~\ref{app:B}.

\begin{definition}[the family bound at order $q$]\label{app:G1:I5-F}
For $q\in\N$, $(\mathrm F_q)$ is the statement: there are $C(q)=C(q,\alpha_0;\mathsf d,L,a,M_b)$ and $\delta(q)>0$ such that every localised renormalised family $G_{\rm ren}$ of order $\le q$, on $T_\eps$, at every $k\le K$, with domains $\Omega,\Omega_1,\Omega_2$ and background $\tilde B$ as in Corollary~\ref{cor:prop1R4}, satisfies
\[
\abs{E(G_{\rm ren},\{\square(v)\},\Phi_{\rm ext},A_{\rm ext})}\le C(q)\,e(t)^{d_v(G_{\rm ren})}\lambda(t)^{d_s(G_{\rm ren})}\,\mathrm{e}^{-\delta(q)d(\{\square(v)\})}\,\norm{h\Phi_{\rm ext}}_{1,\alpha_0}\norm{h'A_{\rm ext}}_{1,\alpha_0},
\]
and coloured families (Theorem~\ref{app:G1:WP-WP1}) satisfy the same bound with the extra factor $(\abs uc_k)^m$, $m$ the number of $D$-lines.
\end{definition}
\noindent \cite[(1.33), p.~420]{Ba3} has a rate $\delta_0$ that ``depends on the dimension $d$ only'' and a constant that ``depends on $\alpha_0,\nbar$'' \cite[p.~421]{Ba3}. $(\mathrm F_q)$ is weaker: its rate may depend on the order, since the rates used to move lines decrease with the number of lines (\cite[p.~427]{Ba3}; Lemma~\ref{app:G1:VSR-A}).

\begin{lemma}[divergent subgraphs do not depend on the cap]\label{app:G1:I5-pc}
In $\mathsf d=3$ a graph with $E$ external legs, $e$-order $\alpha$ and $\lambda$-order $\beta$ has superficial degree $\omega=3-E/2-\alpha/2-\beta$. Hence $\omega\ge0$ if and only if $\alpha+2\beta\le6-E$, and every superficially divergent graph has coupling order $\le6<\nbar$.
\end{lemma}
\begin{proof}
$\varphi$, $A$ and $e$ have dimension $\frac12$, $\lambda$ has dimension $1$. For $E=0$ this is $\alpha+2\beta\le6$, i.e.\ \cite[p.~418]{Ba3}, ``$2\le\alpha+2\beta\le6$''; there is no case with $E\ge6$; and $\alpha+\beta\le\alpha+2\beta$.
\end{proof}

\begin{lemma}[the cube conditions at every scale]\label{app:G1:I5-cube}
Let $C_{\rm loc}(\Nbar)$ and $C_{\Nbar}$ be the numbers fixed at stage 4 of Appendix~\ref{app:B}, and assume
\[
\text{(R.cap)}\ \ R\hat\ell^{2}\ge C_{\rm loc}(\Nbar),\qquad\text{(T.cap)}\ \ \ell_T\ge3r(\eps_0)\eps_0+C_{\Nbar}\eps_0,\qquad\text{(E.r)}\ \ \eps_0<\mathrm{e}^{-1}.
\]
Then for every $0\le k\le K$: (a) $r(t_k)\ge C_{\rm loc}(\Nbar)$; (b) $(3r(t_k)+C_{\Nbar})t_k\le\ell_T$. So at every scale the cube of \cite[p.~432]{Ba3} containing a connected component of a family of order $\le\Nbar$ fits into $\square_1$ (side $r(t_k)$ in units of $t_k$-blocks), and $\square$ (side $3r(t_k)$) together with this cube has physical size $\le\ell_T$, so it does not wrap.
\end{lemma}
\begin{proof}
(a) $t_k\le\eps_0$ gives $1+\log t_k^{-1}\ge\hat\ell$. (b) $f(t):=t(1+\log t^{-1})^2$ has $f'(t)=s^2-1$, $s=\log t^{-1}$, so $f$ increases on $(0,\mathrm{e}^{-1}]$; hence $r(t_k)t_k=Rf(t_k)\le r(\eps_0)\eps_0$ and $C_{\Nbar}t_k\le C_{\Nbar}\eps_0$. Apply (T.cap).
\end{proof}

\begin{proposition}[the order cap]\label{app:G1:I5}
Let $\nbar=13$, $\Nbar=\nbar^2=169$, $0<\alpha_0<1$, $L=2$, and assume (R.cap), (T.cap) and (E.r) of Lemma~\ref{app:G1:I5-cube}. Assume Ba{\l}aban's estimate \cite[(1.33), p.~420]{Ba3} in the form of Proposition~\ref{app:G1:I145-prop1base}, with the structure of its proof described in steps (1)--(7) below, together with its version for graphs with coloured photon lines ($D$-lines; Definition~\ref{app:G1:VSR-IH}, Theorem~\ref{app:G1:WP-WP1}). Then $(\mathrm F_q)$ holds for every finite $q$, in particular for $q=\Nbar$, with:
\begin{itemize}
\item[(i)] $C(\Nbar)$ a finite multiple of the constant $C_A\le c_{\mathsf d}^m(4/\delta_1)^{\mathsf dm}(4m+2)^{\mathsf dm(m-1)/2}(1-L^{-\vartheta})^{-m}$ of Lemma~\ref{app:G1:VSR-A} with $m\le m(\Nbar)$ lines, times the number of graphs of order $\le\Nbar$;
\item[(ii)] some rate $\delta(\Nbar)>0$, and on the coloured side the factor $(1+\abs uc_k)^{N_v}\le(1+L^2)^{N_v}\le5^{84}$ for all $u\in[-\eta,1]$, where $N_v\le\Nbar/2$ is the number of vector lines.
\end{itemize}
These constants do not depend on $\eps$, $k$, the position, $\Omega,\Omega_1,\Omega_2$ or $\tilde B$. They are uniform over tori with $\ell_T\ge\ell_0$, with values taken at the floor $\ell_0^\natural$ for families containing $\delta m^2_{\rm fin}$ vertices. No step requires a smallness condition on $e,\lambda,e(t),\lambda(t)$ that depends on the order. Every superficially divergent subgraph has $\alpha+2\beta\le6-E$ (Lemma~\ref{app:G1:I5-pc}), so the counterterms do not depend on the cap.
\end{proposition}
\noindent (R.cap) is met at stage 5 (by $R\ge C_{\rm loc}(\Nbar)$, since $\hat\ell\ge1$) or at stage 10b, (E.r) at stage 10b, and (T.cap) at stage 11 (Lemma~\ref{app:G1:RW-B20}).
\begin{proof}
Proposition~1 of \cite{Ba3} bounds one family at a time. Each step of its proof, with the sketched steps supplied as described in the paragraph \emph{The base} below, is a finite induction over the nested subgraphs $G_1\subset\dots\subset G_m=G$ of one graph of fixed order; nothing is induced over the order, and no series in the order is summed. The order enters as follows.
(1) \emph{Graph expansion} \cite[p.~420]{Ba3}: the order fixes only the range of the final finite sum over families, whose number depends only on $\Nbar$.
(2) \emph{Single lines} \cite[(2.10)--(2.12), p.~426]{Ba3} (Proposition~\ref{prop:UV-BaProp21}): no order enters; the condition ``$e(L^k\eps)$ small'' for $G_k(\Omega,\tilde B)$ is a property of the kernel.
(3) \emph{Tree factor} \cite[(2.13)--(2.14)]{Ba3}: before the induction each line keeps $\exp[-\frac12\delta_1\dist/(L^{j_l}\eta)]$, independently of the order.
(4) \emph{Moving lines} (Lemma~\ref{app:G1:VSR-A}, step (a$'$)): $\delta^{(i)}\ge\frac{\delta_1}2(4m+2)^{-(i-1)}$, $m\le m(\Nbar)$, the ``$\delta_2(\delta_1,\nbar)$'' of \cite[p.~427]{Ba3}; this gives $\delta(\Nbar)>0$ (not Ba{\l}aban's ``$\delta_0$ depends on $d$ only'' \cite[p.~420]{Ba3}). The sums over $\Delta(w')$ and over scale differences give $C_A$; the geometric sum uses $\vartheta\ge\min(\alpha_0,\frac12)$. In $\mathsf d=3$ the degrees grow with the number of vertices, so higher orders only enlarge the margins.
(5) \emph{Renormalisation} \cite[Cor.~2.3, pp.~429--432; pp.~434--445]{Ba3}: by Lemma~\ref{app:G1:I5-pc} every divergent subgraph has order $\le6$, so the divergent classes and the counterterms $\delta m^2$, $E_1$ with their partners are the same at every cap. By Lemma~\ref{app:G1:WP-deg}(d) every divergent subgraph is $C$-only; the hypothesis $\nbar>12$ of Theorem~\ref{app:G1:WP-WP1} is a lower bound, kept by $\Nbar$; the colour factor is counted by $N_v\le\Nbar/2\le84$.
(6) \emph{Localisation of a component} \cite[pp.~432--433]{Ba3}: each connected component lies in a cube whose size depends on its number of vertices; this cube must fit into $\square_1$, and $\square$ must not wrap. Lemma~\ref{app:G1:I5-cube} gives both at every scale, so Proposition~\ref{prop:prop1}(c) and Corollary~\ref{cor:prop1R4}, stated with $\ell_{\Nbar}$, apply as written.
(7) \emph{The constant}: \cite[p.~421]{Ba3} states that $O(1)$ ``depends on $\alpha_0,\nbar$''; here it is $C(\Nbar)$ of (i).
Only (4) and (6) depend on the cap, and neither imposes a smallness condition on the couplings. The torus enters only through the finite size of $T_\eps$ and (6); the only $\ell_T$-dependent constant is the zero-mode part of $\delta m^2_{\rm fin}$, which decreases in $\ell_T$ and is evaluated at $\ell_0^\natural$.

\emph{The base.} The structure of the proof of Proposition~\ref{app:G1:I145-prop1base} used in (1)--(7) is the proof of \cite[Prop.~1]{Ba3} on pp.~422--434, completed by Lemmas~\ref{app:G1:VSR-A}, \ref{app:G1:VSR-Btop}, \ref{app:G1:VSR-C}, Theorem~\ref{app:G1:VSR-T335}, Propositions~\ref{app:G1:VSR-5b}, \ref{prop:prop1}, Corollary~\ref{cor:prop1R4}, the graded counterterm bound Lemma~\ref{app:G1:I5-ct}, for coloured families Theorem~\ref{app:G1:WP-WP1}, and by the four steps that \cite{Ba3} only sketches: the closure of the renormalised classes under cutting \cite[p.~432]{Ba3} (Theorem~\ref{app:K:L09-3}); the bad localisations of the counterterms \cite[p.~422]{Ba3} (Proposition~\ref{app:K:L09-4}); the H\"older bounds of the background fields \cite[p.~412]{Ba3}, for vector fields Lemma~\ref{app:K:L09-5} and for the scalar external fields Lemma~\ref{app:G1:H3-HOs} below, whose hypothesis (N$_{\rm nn}$) is Proposition~\ref{app:K:L09-2}; and the nested and one-particle-reducible cases in step (vi) of Lemma~\ref{app:G1:VSR-C} (Theorem~\ref{app:K:L09-6}).

\emph{Tail constant.} The constant $C_>$ of the tails of order $>\nbar$ (Proposition~\ref{app:G1:I5-T582}) is $C_{582}\sup_{t\in(0,1]}t^{3/10}P(t)^{\Nbar}$: for $14\le q\le\Nbar$ and $t\le1$, $t^{q/4}P(t)^q\le t^{16/5}\cdot t^{3/10}P(t)^{\Nbar}$ (since $P\ge1$), and a positive power of $t$ beats every power of $1+\log t^{-1}$. So the supremum is finite and does not depend on $\eps_1$.
\end{proof}

\noindent\emph{Scalar H\"older bounds.} Let $d=3$, $\eta=L^{-k}$, $e(t)=et^{1/2}$, and let $G(X,A):=G_k(X,A)$ be the propagator \cite[(2.22)]{Ba1} on a union $X$ of big blocks, with Neumann conditions on $\partial X$ (none if $X=T_\eta$). Put $U(a)=\exp(\eta qe(t)a)$ for a bond variable $a$, $U(A(\Gamma))=\exp(qe(t)A(\Gamma))$ with $A(\Gamma)=\sum_{b\in\Gamma}\pm\eta A(b)$, and $D_{A,\mu}u(z)=\eta^{-1}(U(A_\mu(z))u(z+\eta e_\mu)-u(z))$. For $0<\abs{x-x'}\le1$ let $\Gamma_{x,x'}$ be a shortest lattice contour (length $\abs{x-x'}_1\le\sqrt d\abs{x-x'}$) and
\begin{gather*}
H^A_{x,x'}(u):=\max_\mu\abs{x-x'}^{-\alpha}\abs{U(A(\Gamma_{x,x'}))D_{A,\mu}u(x')-D_{A,\mu}u(x)},\\
\mathcal N^A_x(u):=\max\Big\{\max_\mu\abs{D_{A,\mu}u(x)},\sup_{0<\abs{x'-x}\le1}H^A_{x,x'}(u)\Big\}.
\end{gather*}
$c_0,\delta_0,R_0$ are the constants of Proposition~\ref{prop:UV-BaProp21} at the exponent $\alpha<1$, $M_b$ is the side of the big blocks, $c_r,\beta_r$ are the constants of \cite[(2.23)]{Ba1}, $\ell_Q\le d$ bounds the lengths of the contours in $Q_k$, and $\vartheta_R$ and $C_{\rm rem}$ are the constants of Proposition~\ref{app:G1:VSR-K3}(a) at $m=1$. The transporters are orthogonal and commute with each other, and $\abs{\mathrm{e}^{q\vartheta}-1}\le\abs\vartheta$ for real $\vartheta$. We assume $0\le m^2t^2\le c$ for a fixed constant $c$.

\begin{lemma}[change of the transporting field]\label{app:G1:H3-TC}
Let $A_0$ be a constant field, $B=A_0+\tilde A$, $x$ a site, $N_x:=\{z:\abs{z-x}\le\sqrt d+2\}$ and $s:=e(t)(\sup_{N_x}\abs{\tilde A}+\sup_{N_x}\abs{\partial^\eta\tilde A})\le1$. Then for every field $u$, every $\mu$ and $0<\abs{x'-x}\le1$: (a) $\abs{D_{B,\mu}u(x)}\le\abs{D_{A_0,\mu}u(x)}+s\abs{u(x)}$; (b) $H^B_{x,x'}(u)\le H^{A_0}_{x,x'}(u)+2\sqrt d\,s(\sup_{N_x}\abs{D_{A_0}u}+\sup_{N_x}\abs u)$.
\end{lemma}
\begin{proof}
With $\hat\phi_\mu:=\eta^{-1}(U(\tilde A_\mu)-1)$, $D_{B,\mu}u=U(\tilde A_\mu)D_{A_0,\mu}u+\hat\phi_\mu u$, $\abs{\hat\phi_\mu}\le e(t)\abs{\tilde A_\mu}$, $\abs{\hat\phi_\mu(z')-\hat\phi_\mu(z)}\le e(t)\abs{z'-z}_1\sup_{N_x}\abs{\partial^\eta\tilde A}$ (chaining in $N_x$); this gives (a). For (b), with $\Gamma=\Gamma_{x,x'}$,
\begin{align*}
U(B(\Gamma))D_{B,\mu}u(x')-D_{B,\mu}u(x)&=U(\tilde A_\mu(x))[U(A_0(\Gamma))D_{A_0,\mu}u(x')-D_{A_0,\mu}u(x)]\\
&\quad+U(A_0(\Gamma))[U(\tilde A(\Gamma))U(\tilde A_\mu(x'))-U(\tilde A_\mu(x))]D_{A_0,\mu}u(x')\\
&\quad+\hat\phi_\mu(x')[U(B(\Gamma))u(x')-u(x)]+(\hat\phi_\mu(x')-\hat\phi_\mu(x))u(x).
\end{align*}
The first term has norm $\abs{x-x'}^\alpha$ times the $\mu$-term of $H^{A_0}_{x,x'}(u)$. The second bracket is $\le\sqrt d\,s\abs{x-x'}$; telescoping along $\Gamma$ gives $\abs{U(B(\Gamma))u(x')-u(x)}\le\sqrt d\abs{x-x'}(\sup_{N_x}\abs{D_{A_0}u}+s\sup_{N_x}\abs u)$, and $\abs{\hat\phi_\mu(x')}\le s$; the last term is $\le\sqrt d\,s\abs{x-x'}\abs{u(x)}$. The last three terms are at most $\sqrt d\,s\abs{x-x'}\big(2\sup_{N_x}\abs{D_{A_0}u}+(1+s)\sup_{N_x}\abs u\big)\le2\sqrt d\,s\abs{x-x'}(\sup_{N_x}\abs{D_{A_0}u}+\sup_{N_x}\abs u)$, since $s\le1$. Divide by $\abs{x-x'}^\alpha$, use $\abs{x-x'}\le1$, and take the maximum over $\mu$.
\end{proof}

\begin{lemma}[background differences]\label{app:G1:H3-BD}
Let $\Omega$ be a union of big blocks and $\tilde B$ a real field on $\Omega$ with $\abs{\partial^\eta\tilde B}\le C_\partial p(t)$. Let $x$ lie in the unit block of the site $y$, $Q_s:=\{z:\abs{z-y}_\infty\le s\}$, and let $\rho\ge2$ be an integer with $Q_{2\rho+3}\subset\Omega$ not wrapping around the torus, $\dist(Q_2,\Omega^c)\ge R_0$ and $\rho'':=\rho-M_b-\sqrt d-4\ge R_0$. Let $\square$ be the union of the big blocks contained in $Q_\rho$ (a rectangular parallelepiped, $\square\supset Q_{\rho-M_b}$), $A_{0,\mu}:=\tilde B_\mu(y)$, $\tilde A:=\tilde B-A_0$, $\sigma_*:=3dC_\partial e(t)p(t)(\rho+3)$, and assume $\sigma_*\le\min\{\vartheta_R,1\}$ and $(1+3d)C_\partial p(t)\le c_re(t)^{\beta_r-1}$. For $\psi$ on the unit lattice of $\Omega$ let $\psi_\square:=\psi1_{\{y':B(y')\subset\square\}}$, $\phi:=a_kG(\Omega,\tilde B)Q_k^*(\tilde B)\psi$ and $\Psi:=a_kG(\square,A_0)Q_k^*(A_0)\psi_\square$. Then, with $\mathcal E_\psi:=C_D(\mathrm{e}^{-\delta_0\rho''}+\sigma_*)a_k\norm\psi_\infty$ and $C_D:=c_0(4+8\sqrt d+\ell_Q)+C_{\rm rem}$,
\[
\max_\mu\abs{D_{\tilde B,\mu}\phi(x)}\le\max_\mu\abs{D_{A_0,\mu}\Psi(x)}+\mathcal E_\psi,\qquad H^{\tilde B}_{x,x'}(\phi)\le H^{A_0}_{x,x'}(\Psi)+\mathcal E_\psi\quad(0<\abs{x'-x}\le1).
\]
\end{lemma}
\begin{proof}
The second assumption makes $\tilde B$ regular in the sense of \cite[(2.23)]{Ba1}, and chaining gives $\abs{\tilde A_\mu(z)}\le\abs{z-y}_1C_\partial p$ on $Q_{2\rho+3}$. Let $\chi(z):=\min\{1,\max\{0,2-\abs{z-y}_\infty/(\rho+1)\}\}$, $\tilde A^\sharp_\mu:=\chi\tilde A_\mu$ (zero off $Q_{2\rho+2}$) and $\tilde B^\sharp:=A_0+\tilde A^\sharp$ on $T_\eta$. Then $\norm{\tilde A^\sharp}_\infty\le d(2\rho+2)C_\partial p$, $\norm{\partial^\eta\tilde A^\sharp}_\infty\le(1+3d)C_\partial p$, so $\sigma^\sharp:=e(t)(\norm{\tilde A^\sharp}_\infty+\norm{\partial^\eta\tilde A^\sharp}_\infty)\le\sigma_*$; $\tilde B^\sharp$ is regular, $\theta(\tilde B^\sharp)\le\vartheta_R$, $\theta(A_0)=0$, and $\sigma_\square:=e(t)(\sup_{Q_{\rho+1}}\abs{\tilde A}+\sup_{Q_{\rho+1}}\abs{\partial^\eta\tilde A})\le\sigma_*$. Since $\tilde B^\sharp=\tilde B$ on the bonds based in $Q_{\rho+1}\supset\square$, $G(\square,\tilde B^\sharp)=G(\square,\tilde B)$, and $D_{\tilde B^\sharp}$, $H^{\tilde B^\sharp}$ at $x,x'$ equal those for $\tilde B$. Also $\dist(\{x,x'\},\square^c)\ge\rho''$, and $N_x\subset Q_{\rho+1}$ has distance $\ge\rho''\ge R_0$ from $\square^c$.
With $f_{\rm full}:=a_kQ_k^*(\tilde B)\psi$, $f:=a_kQ_k^*(\tilde B)\psi_\square$, $f_0:=a_kQ_k^*(A_0)\psi_\square$ (sup norms $\le a_k\norm\psi_\infty$), $\phi=\Psi+\sum_{i=1}^6R_i$ with $R_1:=G(\Omega,\tilde B)(f_{\rm full}-f)$, $R_2:=[G(\Omega,\tilde B)-G(\square,\tilde B)]f$, $R_3:=G(\square,\tilde B)(f-f_0)$, $R_4:=[G(\square,\tilde B^\sharp)-G(T_\eta,\tilde B^\sharp)]f_0$, $R_5:=[G(T_\eta,\tilde B^\sharp)-G(T_\eta,A_0)]f_0$, $R_6:=[G(T_\eta,A_0)-G(\square,A_0)]f_0$. By \cite[(2.24)--(2.26)]{Ba1} (Proposition~\ref{prop:UV-BaProp21}), $\mathcal N^{\tilde B}_x(R_i)\le c_0\mathrm{e}^{-\delta_0\rho''}a_k\norm\psi_\infty$ for $i=1,2,4$. Since $\abs{U(\tilde B(\Gamma))-U(A_0(\Gamma))}\le\ell_Q\sigma_\square$ for the contours of $Q_k$, $\mathcal N^{\tilde B}_x(R_3)\le c_0\ell_Q\sigma_*a_k\norm\psi_\infty$. By the resolvent identity $R_5=-G(T_\eta,\tilde B^\sharp)VG(T_\eta,A_0)f_0$ with $V:=G(T_\eta,\tilde B^\sharp)^{-1}-G(T_\eta,A_0)^{-1}$, and Proposition~\ref{app:G1:VSR-K3}(a) ($n=1$, $n'=0$, perturbation $-\tilde A^\sharp$) gives $\mathcal N^{\tilde B}_x(R_5)\le C_{\rm rem}\sigma^\sharp a_k\norm\psi_\infty$. Lemma~\ref{app:G1:H3-TC} with $s\le\sigma_*$, and (2.26), (2.25) for $G(\square,A_0)$, give $\mathcal N^{\tilde B}_x(R_6)\le c_0(1+4\sqrt d)\mathrm{e}^{-\delta_0\rho''}a_k\norm\psi_\infty$ and bound both differences between the $\tilde B$- and the $A_0$-quantities of $\Psi$ by $4\sqrt d\,c_0\sigma_*a_k\norm\psi_\infty$. Summing gives at most $\mathcal E_\psi$.
\end{proof}
\noindent Lemma~\ref{app:G1:H3-BD} measures $\tilde B-A_0$ on all of $Q_{2\rho+3}$. A pointwise kernel bound in which $\tilde B-A_0$ enters only through its values on $[x,x']$ and on the block of $x'$ is false already for $-\Delta_A+m^2$: its random-walk expansion has positive weights, so $(-\Delta_A+m^2)^{-1}(x,x)<(-\Delta_0+m^2)^{-1}(x,x)$ whenever $A$ is non-zero (modulo the period) on one bond away from $x$.

\begin{lemma}[scalar external fields]\label{app:G1:H3-HOs}
Let $\Omega$, $\tilde B$ be as in Lemma~\ref{app:G1:H3-BD} and $\phi^{(k)}=a_kG(\Omega,\tilde B)Q_k^*(\tilde B)\psi$. Assume
\begin{itemize}
\item[($\Psi$)] $\norm\psi_\infty\le C_\psi\lambda(t)^{-1/4}p(t)$;
\item[(N$_{\rm nn}$)] $\abs{U(\tilde B(\Gamma_{y_1,y_2}))\psi(y_2)-\psi(y_1)}\le C_{\rm nn}p(t)$ for nearest neighbours $y_1,y_2$ of the unit lattice in $\Omega$;
\item[(D)] $\abs{D_{\tilde B}\phi^{(k)}(z)}\le C_{2.66}p(t)$ for $\dist(z,\Omega^c)\ge r(t)$ (as in \cite[(2.66)]{Ba2}; Lemma~\ref{lem:lem24});
\end{itemize}
and let $t\le\eps_0$, where for $t\le\eps_0$: $r(t)\ge4\sqrt d(R_0+M_b+\sqrt d+6)$; $\sigma_*(t):=3dC_\partial e(t)p(t)(r(t)+3)\le\min\{\vartheta_R,1\}$; $(1+3d)C_\partial p(t)\le c_re(t)^{\beta_r-1}$; $\ell_T/t\ge8r(t)$. Then for $\dist(\{x,x'\},\Omega^c)\ge r(t)$ and $\alpha<1$,
\[
\abs{x-x'}^{-\alpha}\abs{U(\tilde B(\Gamma_{x,x'}))D_{\tilde B}\phi^{(k)}(x')-D_{\tilde B}\phi^{(k)}(x)}\le C^sp+\omega(t)p,
\]
with $S_d:=\sum_{z\in\Z^d}\abs z_1\mathrm{e}^{-\delta_0(\abs z-2\sqrt d-1)_+}$, $C^s:=\max\{2C_{2.66},c_0aS_dC_{\rm nn}\}$ and
\[
\omega(t):=aC_\psi\lambda(t)^{-1/4}\big[(c_0S_d+C_D)\sigma_*(t)+C_D\mathrm{e}^{-\delta_0(r(t)/(4\sqrt d)-M_b-\sqrt d-5)}\big]\to0\quad(t\to0).
\]
\end{lemma}
\begin{proof}
In this proof $e$, $\lambda$, $p$, $r$ stand for $e(t)$, $\lambda(t)$, $p(t)$, $r(t)$. For $\abs{x-x'}>1$, (D) bounds the quotient by $2C_{2.66}p$. Let $\abs{x-x'}\le1$, $y$ the block of $x$, $\rho:=\lfloor r(t)/(4\sqrt d)\rfloor$. Then $\rho\ge2$, $\rho''\ge r/(4\sqrt d)-M_b-\sqrt d-5\ge R_0+1$, $Q_{2\rho+3}$ lies at distance $\ge r/2-4\sqrt d\ge R_0$ from $\Omega^c$ and has side $<\ell_T/t$, and $3dC_\partial ep(\rho+3)\le\sigma_*(t)$; Lemma~\ref{app:G1:H3-BD} with $a_k\le a$ and ($\Psi$) gives $H^{\tilde B}_{x,x'}(\phi^{(k)})\le H^{A_0}_{x,x'}(\Psi)+aC_\psi\lambda^{-1/4}C_D(\mathrm{e}^{-\delta_0\rho''}+\sigma_*(t))p$. On the non-wrapping box $\square$ put $\Lambda(z):=\sum_\mu A_{0,\mu}(y-z)_\mu$ and $(\mathcal Uu)(z):=\exp(qe(t)\Lambda(z))u(z)$; then $\Psi=\mathcal Uv$ with $v:=a_kG(\square,0)Q_k^*\psi'$, $\psi'(y'):=\exp(-qe(t)\Lambda(y'))\psi_\square(y')$, and $H^{A_0}_{x,x'}(\Psi)=\max_\mu\abs{x-x'}^{-\alpha}\abs{\partial^\eta_\mu v(x')-\partial^\eta_\mu v(x)}$. Since $G(\square,0)Q_k^*1=(a_k+m^2t^2)^{-1}1_\square$, $\partial^\eta v=\sum_{y'}\partial^\eta G(\square,0)g_{y'}$, $g_{y'}:=a_kQ_k^*[(\psi'(y')-\psi'(y))1_{\{y'\}}]$. Along a monotone nearest-neighbour path from $y$ to $y'$ in $\square$, (N$_{\rm nn}$) and $\abs{U(A_0(\Gamma_i))-U(\tilde B(\Gamma_i))}\le\sigma_*(t)$ give $\norm{g_{y'}}_\infty\le a\abs{y'-y}_1(C_{\rm nn}+C_\psi\lambda^{-1/4}\sigma_*(t))p$. \cite[(2.24)]{Ba1} at zero field on $\square$, with $\dist(\{x,x'\},B(y'))\ge\abs{y'-y}-2\sqrt d-1$, gives $H^{A_0}_{x,x'}(\Psi)\le c_0aS_d(C_{\rm nn}+C_\psi\lambda^{-1/4}\sigma_*(t))p$. Adding the two bounds, the quotient is $\le c_0aS_dC_{\rm nn}p+\omega(t)p\le C^sp+\omega(t)p$. Finally $\lambda^{-1/4}\sigma_*(t)=O(t^{1/4}(1+\log t^{-1})^{c})$ and $\lambda^{-1/4}\mathrm{e}^{-\delta_0r/(4\sqrt d)}\le Ct^{-1/4}\mathrm{e}^{-c'(1+\log t^{-1})^2}$; both tend to $0$.
\end{proof}
\noindent By Proposition~\ref{app:K:L09-2}, applied at $t=t_k$ with $\Omega'$ a union of blocks of $B^k(\Lambda^{(k)}_0)$ that contains $\square$, and with (NN2) given by Lemma~\ref{app:K:L09-1}, (N$_{\rm nn}$) holds with $C_{\rm nn}=c_\psi(t)=4Ld+O(t^{1/4}p(t))$. Using it, one replaces $C_{\rm nn}$ by $4Ld$ in $C^s$ and adds $c_0aS_d(c_\psi(t)-4Ld)=O(t^{1/4}p(t))$ to $\omega(t)$; the conclusion is unchanged.

\begin{lemma}[graded bound on the mass counterterm]\label{app:G1:I5-ct}
Let $\eps=L^{-K}\eps_0$, $0\le k\le K$, $t=t_k$, and let $\delta m^2_{K,k}(x)$ be the part of $\delta m^2_K(x)$ formed by the index assignments whose largest line index is $\ge k$ \cite[(1.27)--(1.28), p.~418]{Ba3}. Assume $\ell_T\ge6r(\eps_0)\eps_0$ (implied by (T.nw)). Then for every $x\in T_\eps$
\[
\abs{\delta m^2_{K,k}(x)}t^2\le C_m\Big[(e^2+\lambda)t\tfrac{L}{L-1}+(e^4+e^2\lambda+\lambda^2)t^2(K-k)\Big],\qquad K-k=\log_L(\eps_0/t),
\]
with $C_m=C_m(\alpha_0,\nbar,\mathsf d,L,a,M)$ independent of $\eps,k,x,\ell_T$. In particular $\abs{\delta m^2_{K,k}}t^2=O(t)$, as asserted in \cite[p.~418]{Ba3}.
\end{lemma}
\begin{proof}
In units of scale $K$, $\delta m^2_{K,k}(x)t_K^2=\sum_{G_0}c^{[\ge k]}_{G_0}(x)$, a finite sum over the $C$-only scalar self-energy families of order $o:=\alpha/2+\beta\in\{1,2\}$ and the index assignments with largest index $j^*\ge k$. With the legs equal to $1$ and the vertex of leg 1 at $x$, Lemmas~\ref{app:G1:VSR-Btop}, \ref{app:G1:VSR-C} and \ref{app:G1:VSR-A}(ii) give $\abs{c^{[j^*]}_{G_0}(x)}\le O(1)e(t_K)^\alpha\lambda(t_K)^\beta(L^{j^*}\eta_K)^{D(G_0)}$, and $D(G_0)=o-2$ by Lemma~\ref{app:G1:WP-deg}(b). For $o=1$, $\sum_{j^*=k}^{K-1}L^{K-j^*}\le\frac{t_K}t\frac L{L-1}$; for $o=2$ there are $K-k$ terms. Undo the scaling with $e(t_K)^2=e^2t_K$, $\lambda(t_K)=\lambda t_K$ and multiply by $t^2$; finally $t\log_L(\eps_0/t)\le\eps_0/(\mathrm{e}\ln L)$. The torus condition makes the cubes of \cite[p.~433]{Ba3} fit at scale $K$.
\end{proof}

\begin{lemma}[gain per unit of coupling order]\label{app:G1:I5-po}
Let $G$ be a connected graph with at least two vertices, each one of \cite[(1.6)--(1.11), (1.14), (1.15), (1.18), pp.~413--415]{Ba3} or a mass counterterm vertex (1.7) bounded by Lemma~\ref{app:G1:I5-ct}; the coupling-free vertices (1.12), (1.13) are counted as part of the external scalar fields \cite[p.~421]{Ba3}. For a vertex of order $q_v\ge1$ let $F_v$ be its coupling factor times the bounds \eqref{app:G1:I5-legs} of its external legs. Then $F_v\le C(q_v)t^{q_v/4}P(t)^{q_v}$, hence $\prod_vF_v\le C(q)t^{q/4}P(t)^q$, with $C(q)$ depending on $q,e,\lambda,C_1,C_m,L$.
\end{lemma}
\begin{proof}
Every vertex has an internal leg, a vertex of order $q_v$ has at most $4q_v$ legs, a scalar leg costs $\le C_1p\lambda^{-1/4}t^{-1/4}$ and a vector leg $\le C_1p$. (1.6): $\lambda t(C_1p\lambda^{-1/4}t^{-1/4})^3\le C t^{1/4}p^3$. (1.7): the first term of Lemma~\ref{app:G1:I5-ct} with two scalar legs gives $Ct^{1/2}p^2$; the second gives $Ct^{3/2}(1+K-k)p^2\le Ct^{q_v/4}P(t)$, since $q_v\le6$ and $K-k\le\log_L(1/t)\le C(1+\log t^{-1})$. (1.8), $q_v\ge2$: $Ct^{q_v/2}p^{q_v+2}\lambda^{-1/2}t^{-1/2}$ and $\frac{q_v-1}2\ge\frac{q_v}4$; $q_v=1$: $\le Ct^{1/4}p^2$ with one external scalar leg, and with two the differentiated one is $\le C_1p$ (Lemma~\ref{lem:lem24}), giving $Ct^{1/4}p^3$. (1.10) as (1.8). (1.9), (1.11): factor $e(t)^{n+\nbar+1}$ with $q_v=n+\nbar+1\ge2$, and $\norm{R_{\nbar+1}}\le1$ at real argument (Lemma~\ref{app:G1:I145-unitary}(b)); the count is that of (1.8) and (1.10). (1.14), (1.15), (1.18): as (1.8), and the first two vertices of (1.18) with $n+n'=1$ are $O(t^{1/4}p^2)$ by Lemma~\ref{lem:lem24} and \cite[p.~421]{Ba3}. In all cases $F_v\le C(q_v)t^{q_v/4}p^{4q_v}(1+\log t^{-1})\le C(q_v)t^{q_v/4}P^{q_v}$.
\end{proof}

\begin{lemma}[sum over localisations]\label{app:G1:I5-locsum}
Let a family have $N$ vertices localised in cubes $\square(v)$ of diameter $\le c_0$ attached to points of a lattice of spacing $1$, with at most $n_0$ cube types per point \cite[p.~420]{Ba3}, and let $\delta>0$. The sum of $\mathrm{e}^{-\delta d(\{\square(v)\})}$ over the localisations with $\square(v_1)$ fixed is at most $S(N,\delta):=n_0^{N-1}N^{N-2}\big(\mathrm{e}^{\delta c_0}\sum_{z\in\Z^3}\mathrm{e}^{-\delta\abs z/2}\big)^{N-1}$.
\end{lemma}
\begin{proof}
A shortest Steiner tree has length at least half that of a minimal spanning tree for the edge lengths $\dist(\square(v),\square(w))$ (double, take an Euler tour, shortcut). So $\mathrm{e}^{-\delta d}\le\sum_T\prod_{\{v,w\}\in T}\mathrm{e}^{-\frac\delta2\dist(\square(v),\square(w))}$ over the $N^{N-2}$ labelled trees, and $\dist(\square(v),\square(w))\ge\abs{y_v-y_w}-2c_0$. Sum leaf by leaf from the fixed vertex.
\end{proof}

\begin{lemma}[cumulants are sums of connected graphs]\label{app:G1:I5-lc}
Let $\mu_C$ be a centred Gaussian measure and $V_1,\dots,V_m$ polynomials in the field, each a finite sum of vertices with legs. The joint cumulant $\kappa(V_1,\dots,V_m)$ is the sum over the complete pairings of the legs of $V_1,\dots,V_m$ by lines of covariance $C$ whose graph on $\{1,\dots,m\}$ is connected.
\end{lemma}
\begin{proof}
By Wick's theorem $\int\prod_{i\in S}V_i\,d\mu_C=\sum_\pi\prod_{B\in\pi}c(B)$, grouping pairings by the partition into connected components, $c(B)$ the connected sum. The cumulants satisfy the same moment--cumulant relation, so $\kappa(S)=c(S)$ by M\"obius inversion.
\end{proof}

\begin{proposition}[the tail of order $>\nbar$ after the fluctuation integral, {\cite[p.~582]{Ba2}}]\label{app:G1:I5-T582}
Let $u\in\{1,-\eta\}$, $M\in(0,m_2]$, assume the hypotheses of Propositions~\ref{app:G1:I5} and \ref{app:G1:T579}, and let $0\le k\le K-1$. After the cut of Proposition~\ref{app:G1:T579} the exponent under the second integral of \cite[(2.115), p.~582]{Ba2} is $V=\sum_{q=1}^{\nbar}V^{[q]}$, each $V^{[q]}$ a finite sum of localised families of coupling order $q$. With $\kappa$ the cumulants of the Gaussian fluctuation law with weighted covariance ($C$- and $D$-lines, Theorem~\ref{app:G1:WP-WP1}), let
\[
\Delta_k:=\sum_{m=2}^{\nbar}\frac1{m!}\sum_{1\le q_1,\dots,q_m\le\nbar,\ q_1+\dots+q_m>\nbar}\kappa(V^{[q_1]},\dots,V^{[q_m]}).
\]
Then, pointwise on the small-field set of step $k$, $\abs{\Delta_k}\le C_{582}t_k^{7/2}P(t_k)^{\Nbar}N_k\le C_>t_k^{16/5}N_k$ with $C_>:=C_{582}\sup_{t\in(0,1]}t^{3/10}P(t)^{\Nbar}<\infty$. $C_{582}$ and $C_>$ depend only on $\nbar,\Nbar,\alpha_0,\mathsf d,L,a,M_b,e,\lambda,\mu_*,M,b_0,p,\eta$; they do not depend on $\eps$, $k$, $\ell_T\ge\ell_0$ (at $\ell_0^\natural$ where $\delta m^2_{\rm fin}$ vertices occur) or $\eps_0$, and are finite at $u=-\eta$.
\end{proposition}
\begin{proof}
By Lemma~\ref{app:G1:I5-lc}, $\kappa(V^{[q_1]},\dots,V^{[q_m]})$ is a sum of connected graphs obtained by joining legs by fluctuation lines (\cite[p.~623]{Ba1}, \cite[p.~582]{Ba2}), of order $q=\sum q_i$; $q>\nbar$ forces $m\ge2$, and $q\le m\nbar\le\Nbar$. They regroup into renormalised families of fixed order (\cite[p.~420]{Ba3}, Lemma~\ref{app:G1:WP-fam}); every divergent subgraph has order $\le6$ (Lemma~\ref{app:G1:I5-pc}) and its counterterm partner, of order $\le4$, is a vertex of $V$, so the families close at every order $\le\Nbar$. When long lines are cut, they remain disjoint unions of renormalised classes (Theorem~\ref{app:K:L09-3}(c)). By Proposition~\ref{app:G1:I5} each family obeys $(\mathrm F_{\Nbar})$, with colour factor $(\abs uc_k)^m\le3^{-m}$ ($c_k\le\frac13$, Definition~\ref{app:G1:VSR-IH}). By Lemma~\ref{app:G1:I5-po} a family of order $q$ is bounded by $C(q)C(\Nbar)t^{q/4}P(t)^q\mathrm{e}^{-\delta(\Nbar)d}$. For $q\ge14$, $t^{q/4}\le t^{7/2}$ and $P^q\le P^{\Nbar}$; the number of families is finite; with one block fixed, Lemma~\ref{app:G1:I5-locsum} bounds the other localisations by $S(N,\delta(\Nbar))$, $N\le\Nbar$; the fixed block ranges over $\le N_k$ blocks. The second inequality is $t^{7/2}P^{\Nbar}=t^{16/5}\cdot t^{3/10}P^{\Nbar}$. Torus images only remove localisations.
\end{proof}

\begin{remark}[the other cuts of order $>\nbar$]\label{app:G1:I5-othercuts}
(i) The cut before the fluctuation integral \cite[p.~579]{Ba2} is Proposition~\ref{app:G1:T579} with Remark~\ref{app:G1:T579-rem}; it makes $V$ a polynomial of fixed degree, which Proposition~\ref{app:G1:I5-T582} presupposes. (ii) The terms of order $>\nbar$ of the determinants $Z^{(j)}$ \cite[p.~578]{Ba2} are one-loop graphs, estimated by \cite[Props.~2.1--2.3]{Ba1}; they are not covered by Lemma~\ref{app:G1:I5-po}. (iii) Cumulants of order $m>\nbar$ are not part of $\Delta_k$.
\end{remark}

\subsection{The per-localisation bound, the region sum and Ba{\l}aban's Proposition~1}\label{app:G1:part-TI145}

\paragraph{The per-localisation bound}\label{app:G1:I145-perlocsec}
We use the notation of the paragraph on the moved vacuum pieces in Subsection~\ref{app:G1:part-E4}: $\mathcal V_k(u)$, the pairs $(G,\tilde l)$, the localisations $\Lambda=(\square_\iota)_{\iota\in I(G)}$, $N=\abs{I(G)}\le\Nbar(\nbar)$, $d(\Lambda)$, $\bar c=\kappa\mu_*^2\eps_0^2L^2$, $\bar N_{\rm ph}$. Put $P(t):=\sum_{2\le\alpha+2\beta\le2\nbar}\abs{e(t)}^\alpha\lambda(t)^\beta$, where $e(t)^2=e^2t$ and $\lambda(t)=\lambda t$. Let $m$ be the number of D-coloured lines of a family.

\noindent In this paragraph, all graphs are Ba{\l}aban's graphs \cite[(1.6)--(1.15), pp.~413--414]{Ba3}, graded as in Definition~\ref{app:G1:vacgrade}. Each vector line carries one of two colours, according to the multilinear splitting of the photon covariance $C_{\mu_*}+u\Gamma^B$ of \eqref{eq:T6-GammaB}. A line of colour $C$ has Ba{\l}aban's kernel. A line of colour $D$ carries the slice photon; in the units of scale $k$, its kernel is $D^\eta_k=\sum_{j<k}(L^j\eta)^2G^\eta_{(j)}$ with the coefficient $uc_k$, where $0\le c_k\le\bar c:=\kappa\mu_*^2\eps_0^2L^2$. We write $m(G)$ for the number of $D$-lines of $G$. For a graph with coupling factor $e^\alpha\lambda^\beta$ we put $o(G):=\alpha/2+\beta$. Finally, $E_s$, $E_v$ and $E_\partial$ denote the numbers of external scalar legs, of external vector legs (background legs included), and of differentiations acting on external legs.

\begin{lemma}[per-localisation bound]\label{app:G1:I145-perloc}
Let $\delta_T>0$ be the tree rate of \cite[(1.33)]{Ba3} given by Proposition~\ref{app:G1:I145-prop1base}, and let $\delta_1$ be the rate of \cite[(2.10), p.~426]{Ba3}. Put $\delta:=\min(\delta_T,\delta_1)$. Work at either of the points $\pi_\pm$, and let $0\le k\le K-1$. Consider a pair $(G,\tilde l)$ of a bare vacuum graph $G$ and an ordering $\tilde l$ (Step~1 below) whose last line has index $k$, other than the order-$e^2$ family with one $D$-line ($m(G)=1$, $o(G)=1$). Then for every localisation $\Lambda$,
\begin{equation}\label{app:G1:I145-eq-perloc}
\abs{\mathcal V_{G,\tilde l}(\Lambda)}\le A_k\,\mathrm{e}^{-\delta d(\Lambda)},\qquad A_k:=C_P(1+\abs u\bar c)^{\bar N_{\rm ph}}P(t_k),
\end{equation}
with $P$ and $\bar N_{\rm ph}$ as in the paragraph containing Lemma~\ref{app:G1:treesum}. Here $C_P$ is the constant of \cite[(1.33)]{Ba3} times that of \cite[(2.10)--(2.11)]{Ba3}. If $\abs u\le1$, the $u$-part $\mathcal V_{G,\tilde l}(\Lambda;u)-\mathcal V_{G,\tilde l}(\Lambda;0)$ obeys the same bound with $(1+\abs u\bar c)^{\bar N_{\rm ph}}$ replaced by $\bar N_{\rm ph}\abs u\bar c(1+\bar c)^{\bar N_{\rm ph}-1}$. The order-$e^2$ family with one $D$-line is bounded per anchor cube in Lemma~\ref{app:G1:I1c-e2}.
\end{lemma}
\noindent The lemma rests on Proposition~\ref{app:G1:I145-prop1base} and \cite[(2.10)--(2.11), p.~426]{Ba3}. For the photon-cut classes it also rests on \cite[p.~433]{Ba3}, whose remainder on $\square$ is treated in Proposition~\ref{app:G1:I3}.
\begin{proof}
\emph{Step 1 (exact cut, in the grading of Definition~\ref{app:G1:vacgrade}).} Decompose the index sums by the orderings of \cite[(3.4), p.~424]{Ba3}, with ties broken by a fixed order of the lines; the classes are indexed by parent graphs. For an index assignment, $\tilde l=(l(1),\dots,l(n))$ lists the lines of the bare graph by increasing (index, fixed order). Every assignment belongs to exactly one ordering. Hence every term of $\mathcal V_k$ belongs to exactly one pair $(G,\tilde l)$. Its last line $l(n)$ is a line of the bare graph, and it has index $k$, because $k$ is the largest index on $\mathcal V_k$.

The sets $J(\tilde l)$ of \cite[p.~424]{Ba3} attach to $\tilde l$ the increasing chain of subgraphs $G_i$ spanned by $l(1),\dots,l(i)$, so $l(n)\notin G_i$ for $i<n$. Thus every counterterm operation active in $\tilde l$ acts on a divergent subgraph that does not contain $l(n)$, and it produces a coefficient in $\delta m^2_{k+1}$. Degrees are intrinsic \cite[(2.1)--(2.3), pp.~422--423]{Ba3}. So the divergent subgraphs of $G\setminus l(n)$ are exactly the divergent subgraphs of $G$ that do not contain $l(n)$.

Now cut $l(n)$, i.e.\ treat its kernel $K_k(x,x')$ as a two-variable external field \cite[p.~419]{Ba3}. This maps the class, term by term and with the same coefficients, onto renormalised two-leg classes of $G\setminus l(n)$. In these classes the vertices carrying $\delta m^2_{\rm fin}$ or $\delta m^2_{K,k+1}$ are finite vertices of degree $0$. The coefficients $\delta m^2$ are never opened.

\emph{Step 2 (the cut kernel).} In scale-$k$ units, $L^k\eta=1$. So \cite[(2.10)--(2.11)]{Ba3} at $j=k$ give $\norm{hK_kh'}_{1,\alpha_0}\le O(1)\mathrm{e}^{-\delta_1\dist(\supp h,\supp h')}$. This holds for scalar and vector lines alike (``the same inequalities hold for vector field propagators'', p.~426). It also holds for $D$-lines, because their slice $k$ carries the factor $(L^k\eta)^2=1$. The bound holds on $T_\eps$ (Corollary~\ref{app:G1:VSR-VSphys} at zero background; Proposition~\ref{prop:UV-BaProp21} with Lemma~\ref{app:G1:torusuniform} in general).

\emph{Step 3 (which two-leg classes occur; partners).} By Proposition~\ref{app:G1:I1c-table} (see also Lemma~\ref{app:G1:cutclasses}), the divergent components of $G\setminus l(n)$ are of the following kinds.
\begin{itemize}
\item One-particle-irreducible components of types C1 and C2. Their $\delta m^2$ partners are in $E_1$, which is computed with the full action \cite[(1.24), p.~417]{Ba3}.
\item One-particle-reducible chains of type C2 ($21$ of the $59$ non-vanishing isomorphism classes). They have no partner, since Ba{\l}aban's $\delta m^2$ is built from one-particle-irreducible graphs \cite[p.~416, p.~430]{Ba3}. They need no counterterm, by Lemma~\ref{app:G1:I1c-1PR}.
\item Components of type C1$'$. These are finite vertices \cite[(3.22)--(3.24)]{Ba3}.
\item Components of types C3 and C4. They are renormalised by the Ward--Takahashi class \cite[(3.25)--(3.32)]{Ba3}. On $T_\eps$ this class is used after the transfer of \cite[p.~433]{Ba3} to a non-wrapping cube, where it is exact by Proposition~\ref{app:G1:I2} (Lemmas~\ref{app:G1:gaugetriv} and \ref{app:G1:nowrap}); this needs condition (T.nw).
\item C6 pairs, which vanish.
\end{itemize}
All other components are convergent. In particular, no divergent component contains a vertex (1.13)--(1.15) (Lemma~\ref{app:K:L08-4}). So the cut family is a renormalised family in the sense of Proposition~\ref{app:G1:I145-prop1base}, extended by Lemma~\ref{app:G1:I1c-1PR}.

\emph{Step 4 (coloured families).} Consider a family with $m$ $D$-lines, where $m=0$, $m\ge2$, or $m=1$ and $o\ge2$. For it, the coloured form of (1.33) in Proposition~\ref{app:G1:I145-prop1base} (Theorem~\ref{app:G1:WP-WP1}(c)) gives the bound per localisation times $(\abs uc_k)^m$, with $c_k\le\bar c$. Summing over the colourings of at most $\bar N_{\rm ph}$ photon lines gives the factor $(1+\abs u\bar c)^{\bar N_{\rm ph}}$. The $u$-part is the sum over $m\ge1$, which is bounded by $(1+\abs u\bar c)^{\bar N_{\rm ph}}-1\le\bar N_{\rm ph}\abs u\bar c(1+\bar c)^{\bar N_{\rm ph}-1}$. This inequality holds for $\abs u\le1$ by the mean value theorem: $(1+x)^N-1=Nx(1+\xi)^{N-1}$ with $\xi\in[0,x]$ and $x=\abs u\bar c\le\bar c$. The remaining family, $m=1$ and $o=1$, is excluded from the lemma. It is bounded per anchor cube directly in Ba{\l}aban's localisation, without Proposition~1 (Lemma~\ref{app:G1:I1c-e2}), and this form suffices for Proposition~\ref{app:G1:I1}.

\emph{Step 5 (connectedness).} In a non-vanishing term, $G\setminus l(n)$ is connected. A scalar line is never a bridge, because every vertex has an even number of scalar legs, so the scalar lines form closed charge loops. A photon bridge would leave two pieces with one external vector leg each, i.e.\ a C6 pair. Each such piece has an odd number of vector legs on a closed scalar strand, so it vanishes ($\tr q^{\rm odd}=0$) after the gauge transformation of \cite[p.~433]{Ba3}. So the tree of (1.33) for $G\setminus l(n)$ joins all cubes of $\Lambda$. Labels at one vertex have overlapping supports, since the term vanishes otherwise. Hence $d$ over labels equals $d$ over vertices, and the tree length $d'$ of $G\setminus l(n)$ satisfies $d'\ge d(\Lambda)$.

\emph{Step 6 (assembly).} Apply Proposition~\ref{app:G1:I145-prop1base} to the cut class, with the external field of Step~2 and the colour factor of Step~4. This gives \eqref{app:G1:I145-eq-perloc}, where $e(t_k)^{d_v}\lambda(t_k)^{d_s}\le P(t_k)$ and $\mathrm{e}^{-\delta_Td'}\le\mathrm{e}^{-\delta d(\Lambda)}$. The quantity $\mathcal V$ is a c-number, so no field restriction and no factor $p(t)$ or $\lambda^{-1/4}$ enters.
\end{proof}

\noindent We name the divergent two-leg types by $(E_s,E_v,E_\partial,o)$, following \cite[Cor.~2.3, p.~429; (3.6), (3.18), (3.21), (3.25), (3.31)]{Ba3}:
\begin{itemize}
\item C1 $=(2,0,0,1)$ and C2 $=(2,0,0,2)$, the scalar self-energies;
\item C1$'=(2,0,1,1)$, the self-energy with a differentiated leg;
\item C3 $=(0,2,\cdot,1)$ and C4 $=(0,2,\cdot,2)$, the vacuum polarisations;
\item a C6 pair: two pieces with one external vector leg each, joined by a photon line.
\end{itemize}

\begin{proposition}[degrees and the cut table]\label{app:G1:I1c-table}
Let $d=3$, and let $D$ be Ba{\l}aban's degree \cite[(2.1)--(2.3)]{Ba3}. A mass vertex carrying $\delta m^2_{\rm fin}$ or $\delta m^2_{K,k+1}$ (a \emph{finite mass vertex}) has degree $0$. A mass vertex carrying the contribution of a graph $G_0$ to $\delta m^2_{k+1}$ has degree $D(G_0)$.
\begin{enumerate}
\item[(a)] Let $G$ be a connected graph built from the vertices (1.6), (1.8), (1.10) and from mass vertices of divergent self-energy graphs. Then $D(G)=o(G)+\frac12(E_s+E_v)+E_\partial-3$. If $D(G)\le0$ and $G$ has $V$ vertices, then $V+E_s+E_v+2E_\partial\le6$.
\item[(b)] Let $G$ be a vacuum graph with $D(G)\le0$, and let $\ell$ be a line of $G$. The components of $G\setminus\ell$ with $D\le0$ are of the following kinds only:
\begin{itemize}
\item if $\ell$ is a scalar line that is not differentiated: C1, C2, or a bare vertex;
\item if $\ell$ is a scalar line differentiated once: C1$'$;
\item if $\ell$ is a scalar line differentiated twice: none;
\item if $\ell$ is a photon line: C3, C4, or the two pieces of a C6 pair.
\end{itemize}
The one-particle-reducible members of C2 are chains $H_1$--$\ell'$--$H_2$ of two order-$1$ two-scalar-leg pieces joined by one scalar line. Vacuum graphs with $o\ge4$ give $D(G\setminus\ell)\ge2$. Vacuum graphs with a finite mass vertex give only bare vertices, convergent pieces, and vanishing pieces with an odd number of vector legs on a closed scalar strand.
\end{enumerate}
\end{proposition}
\begin{proof}
(a) By \cite[(2.1)]{Ba3}, a leg external to $G$ contributes $0$ in place of $-\frac12$, and a differentiation of an external leg contributes $0$ in place of $-1$. Summing over the vertices and using \cite[(2.3)]{Ba3} gives $D(G)=\sum_vD(v)+\frac12(E_s+E_v)+E_\partial-3+\sum\deg$. By the list on \cite[p.~423]{Ba3}, the vertex (1.6) has $D=1=o(v)$, and the vertices (1.8), (1.10) with $n$ vector legs have $D=n/2=o(v)$. For a mass vertex, $D(v)+\deg=2+D(G_0)=o(v)$, since $D(G_0)=o(G_0)-2$ by induction on the order. Every such vertex has $o(v)\ge\frac12$, so $0\ge D\ge\frac12V+\frac12(E_s+E_v)+E_\partial-3$.

(b) A vertex (1.9) or (1.11) has at least $\nbar+1>13$ vector legs, so a graph containing one has $D>0$ \cite[p.~423, (iii)--(iv)]{Ba3}. A connected graph containing a vertex (1.13)--(1.15) has $D\ge1$ by Lemma~\ref{app:K:L08-4}. A finite mass vertex has $D_G(v)\ge2$. Let $H$ be a component of $G\setminus\ell$ with $D(H)\le0$. If $H$ has a line, it is one of the divergent graphs of Lemma~\ref{app:G1:H3-class}(ii); otherwise it is a bare vertex. Since $G$ is a vacuum graph, the external legs of $H$ are the legs of $\ell$ that lie in $H$: two legs if both ends of $\ell$ lie in $H$, one otherwise, of the type of $\ell$ and carrying its differentiations. If $\ell$ is a scalar line without differentiation, then $E_v=0$, $E_s\in\{1,2\}$ and $E_\partial=0$; the table of Lemma~\ref{app:G1:H3-class}(ii) has no divergent graph with one scalar leg, so $H$ is of class C1 or C2, or a bare vertex. If $\ell$ is differentiated once, $E_\partial=1$, and the table leaves only C1$'$. If $\ell$ is differentiated twice, $E_\partial=2$, and the table contains no divergent graph. If $\ell$ is a photon line, then $E_s=0$ and $E_v\in\{1,2\}$, which gives C3, C4 or the class C6 of that lemma. All other components are convergent. A one-particle-reducible member of C2 has order $2$, and each piece of a decomposition at a separating line has order $\ge1$, so it consists of exactly two pieces of order $1$, each with two scalar legs, joined by one scalar line. Cutting a line adds two external legs, so for $o\ge4$ part (a) gives $D(G\setminus\ell)\ge o-3+1\ge2$.
\end{proof}

\begin{lemma}[one-particle-reducible chains need no counterterm]\label{app:G1:I1c-1PR}
Work after the transfer of \cite[p.~433]{Ba3} to a cube $\square$ of $\eta\Z^3$ (Proposition~\ref{app:G1:I2}). Let $S^{(1)}$ be the amputated order-$1$ one-particle-irreducible two-scalar-point operator at fixed localisation labels. It consists of the order-$e^2$ and order-$\lambda$ graphs and of the order-$1$ counterterm vertex, with derivative legs written as lattice differences. Write $S^{(1)}=\Sigma_{nn}+\Sigma_{nd}\partial+\partial^*\Sigma_{dn}+\partial^*\Sigma_{dd}\partial$ according to where the differences act \cite[(1.21)]{Ba3}. Then:
\begin{enumerate}
\item[(a)] the sum of the one-particle-reducible members of a C2 family at fixed labels is $\mathcal K:=(S^{(1)}_1-\mathrm{ct}_1)\,K_{\ell'}\,(S^{(1)}_2-\mathrm{ct}_2)$, where $K_{\ell'}$ is the kernel of the scalar line $\ell'$;
\item[(b)] $\sum_{x'}\mathcal K(x,x')=\sum_{y,y'}(S^{(1)}_1-\mathrm{ct}_1)(x,y)K_{\ell'}(y,y')(\partial^*_\nu c_\nu)(y')$, where $c_\nu(y'):=\sum_{x'}\Sigma_{dn,\nu}(y',x')$ is bounded and Lipschitz at the unit scale;
\item[(c)] \cite[(3.19)]{Ba3} applies to $\mathcal K$ with zero counterterm. Both resulting terms have positive degree, and all their proper divergent subgraphs are renormalised. Hence $\mathcal K$ obeys bounds of the type \eqref{app:G1:I145-eq-perloc}, by \cite[Props.~2.1--2.2, pp.~424--428]{Ba3}.
\end{enumerate}
\end{lemma}
\begin{proof}
(a) As formal power series, coefficient by coefficient, the amputated connected two-point function is $S+SK_{\ell'}S+\cdots$. At $o=2$, a one-particle-reducible two-scalar-leg graph is cut by a scalar bridge into two connected two-leg pieces, each of order $\ge1$. A photon bridge would leave a piece with one external vector leg, which vanishes by $\tr q^{\rm odd}=0$ after the gauge transformation of \cite[p.~433]{Ba3}. A two-leg piece of order $\frac12$ is a single vertex (1.8) with its vector leg external, and this is not of type C2. The subgraph counterterms are present in the family, and the labels factor because they are attached to the pieces.

(b) The counterterm is the zeroth moment of the $nn$ part \cite[p.~417, (3.9)]{Ba3}. So $\sum_{x'}(\Sigma_{nn}-\mathrm{ct})(y',x')=0$, and tadpole pieces cancel their counterterms exactly \cite[p.~435]{Ba3}. A term with a difference acting at $x'$ has zero sum, because a sum of differences of a summable function vanishes, on $\eta\Z^3$ and over a period. The remaining terms give $\partial^*_{y'}c(y')$. The propagator $G_k(0)$ is only block periodic \cite[p.~418]{Ba3}. After the replacement of \cite[p.~437]{Ba3}, however, the $C^\xi$ part is translation invariant and the remainder is regular. So $c$ depends on $y'$ only through the unit-scale localisation functions and the regular part, and it is bounded and Lipschitz at the unit scale \cite[(3.22)--(3.24)]{Ba3}.

(c) \cite[(3.19)]{Ba3} is the identity
\[
\varphi(x)\mathcal K(x,x')\varphi'(x')=\varphi(x)\mathcal K(x,x')\abs{x'-x}^\alpha\frac{\varphi'(x')-\varphi'(x)}{\abs{x'-x}^\alpha}+\varphi(x)\varphi'(x)\sum_{x'}\mathcal K(x,x').
\]
The first term has degree $\alpha>0$. By (b), in the second term one difference on the line is replaced by a unit-scale difference of localisation functions, so its degree is $\ge1$. The inner pieces are of types C1, C1$'$ or tadpole, with their counterterms present by (a). The torus counterterm $\delta m^2_k$ is built with the torus propagator \cite[(1.28)]{Ba3}, while the members after the transfer use $G_k(0)$ on $\eta\Z^3$. Their difference is a mass vertex with two scalar legs and no line, whose coefficient is exponentially small in $\ell_T/t_k$, by \cite[(2.10), p.~426]{Ba3}. Like $\delta m^2_{\rm fin}$ and $\delta m^2_{K,k+1}$, it is not the contribution of a graph, so it has degree $0$, and we count it among the finite mass vertices of Proposition~\ref{app:G1:I1c-table}. By the last clause of part (b) of that proposition, its vacuum cuts produce no divergent class.
\end{proof}

\paragraph{The order-$e^2$ family with one $D$-line.}\label{app:G1:part-B5e2}
\emph{Setting.} On $T_\eps$ we use $\ip fg=\sum_x\eps^3f(x)g(x)$, and kernels are taken relative to $\sum\eps^3$. We put $\nabla_\mu=\nabla^+_\mu$, $\Delta_\mu=\nabla^+_\mu\nabla^-_\mu$ and $-\Delta_\eps=-\sum_\mu\Delta_\mu$. For $z\in T_\eps$, $\abs z_1$ and $\abs z_\infty$ are the norms of the minimal representative, and $r:=\abs z_\infty$. Let $c$ be the kernel of $C=(-\Delta_\eps+M^2)^{-1}$. Throughout this paragraph $a=1$ and $t_K=\eps_0$ (Lemma~\ref{app:K:L08-3}), so the hypotheses of Lemma~\ref{app:K:L08-2}(c) hold. Let $G_k=G^M_k$ be King's scalar operators (Lemma~\ref{app:K:L08-1}(S4)), with slices $G_{(i)}:=G_{i+1}-G_i$, $G_0:=0$, and put $\mathcal C_k:=C-G_k\ge0$. The photon slices are $H_j\ge0$ (Lemma~\ref{app:K:L08-1}), $\delta_j:=\kappa\mu_*^2L^2t_j^2$, and $\Gamma^B=\sum_{j<K}\delta_jH_j$. For a two-variable function $f$, $\nabla f$ and $\nabla'f$ are the forward differences in the first and second variable.

For a real symmetric kernel $D\ge0$ and a real symmetric scalar covariance $S$ put
\[
\operatorname{Bub}(S,S)(x,y):=2\eps^{-2}\big[S(x+\eps e_\mu,y+\eps e_\mu)S(x,y)-S(x+\eps e_\mu,y)S(x,y+\eps e_\mu)\big],
\]
\[
V(D;S):=\sum_\mu\Big\{-e^2\sum_x\eps^3D(x,x)S(x+\eps e_\mu,x)+\frac{e^2}2\sum_{x,y}\eps^6D(x,y)\operatorname{Bub}(S,S)(x,y)\Big\}.
\]
This is the order-$e^2$ vacuum coefficient produced by a photon of covariance $\delta_{\mu\nu}D$; it consists of the seagull and the bubble whose two vector legs both end on the photon line. The order-$e^2$ part of the $u$-dependence of the vacuum terms is $u\sum_{j<K}\delta_jV(H_j;C)$. A graph has top index $\ge k$ exactly when its $D$-slice has index $\ge k$, or when that index is $<k$ and the graph contains a line of $\mathcal C_k$.

\emph{Leg partition} \cite[p.~420]{Ba3}. Fix $t\in\{t_i\}$. Put $\phi(s):=\cos^2(\pi s/2)\mathbf 1_{\abs s<1}$ and $h_a(x):=\prod_{i=1}^3\phi(x_i/t-a_i)$ for $a\in(\Z/(\ell_T/t))^3$. Then $\sum_ah_a=1$ and $0\le h_a\le1$, and $h_a$ is supported in the open cube $S_a$ of side $2t$ centred at $ta$. At most $N_{\rm ov}=8$ of the $h_a$ are non-zero at any point, and $\abs{h_a(x)-h_a(y)}\le\varpi(x-y):=\min(c_{h1}\abs{x-y}_1/t,1)$ with $c_{h1}=\pi/2$. Each vector leg is placed at the base point of its bond. Put
\[
F_{D,S}(a,a'):=\sum_\mu\Big\{-e^2\sum_x\eps^3h_ah_{a'}(x)D(x,x)S(x+\eps e_\mu,x)+\frac{e^2}2\sum_{x,y}\eps^6h_a(x)h_{a'}(y)D(x,y)\operatorname{Bub}(S,S)(x,y)\Big\}.
\]
With the partition at scale $t=t_k$, put
\[
F^{\ge k}:=u\Big[\sum_{k\le j<K}\delta_jF_{H_j,C}+\sum_{j<k}\delta_j\big(F_{H_j,C}-F_{H_j,G_k}\big)\Big],\qquad F^{=k}:=F^{\ge k}-F^{\ge k+1},
\]
where $F^{\ge k+1}$ uses the same scale-$t_k$ partition.

\begin{lemma}[the localised family and its Ward form]\label{app:G1:B5e2-sym}
For real symmetric $D$ and $S$, $F_{D,S}(a,a')=F_{D,S}(a',a)$ and $\sum_{a,a'}F_{D,S}(a,a')=V(D;S)$. Moreover
\[
F_{D,C}(a,a')=-\tfrac12\sum_\mu\sum_{x,y}\eps^6h_a(x)h_{a'}(y)D(x,y)\Pi_\mu(x-y),
\]
where $\Pi_\mu(z):=2e^2c(\eps e_\mu)\eps^{-3}\delta_{z,0}-B_\mu(z)$ and $B_\mu(z):=2e^2\eps^{-2}[c(z)^2-c(z+\eps e_\mu)c(z-\eps e_\mu)]$. If in addition $D\ge0$, let $A\sim N(0,D)$, and put $g:=h_a\otimes h_{a'}+h_{a'}\otimes h_a$, $\varphi:=gD$, $\Delta\varphi(x,y):=\varphi(x,x)+\varphi(y,y)-2\varphi(x,y)$ and $\widehat\Pi_\mu(0):=\sum_z\eps^3\Pi_\mu(z)$. Then
\[
F_{D,C}(a,a')=-\tfrac14\sum_\mu\Big[\sum_x\eps^3g(x,x)D(x,x)\widehat\Pi_\mu(0)+\tfrac12\sum_{x,y}\eps^6\Delta\varphi(x,y)B_\mu(x-y)\Big].
\]
Here $\Delta\varphi=g(x,y)\Delta D+D(x,y)\Delta g+(g(x,x)-g(x,y))(D(x,x)-D(x,y))+(g(y,y)-g(x,y))(D(y,y)-D(x,y))$, with $\Delta D(x,y)=\E(A(x)-A(y))^2$ and $\Delta g(x,y)=2(h_a(x)-h_a(y))(h_{a'}(x)-h_{a'}(y))$.
\end{lemma}
\begin{proof}
$\operatorname{Bub}(S,S)$ is symmetric for symmetric $S$. Exchanging $x$ and $y$ and using $D(y,x)=D(x,y)$ gives the symmetry; the seagull sum is symmetric in $(a,a')$. The total follows from $\sum_ah_a=1$. For $S=C$, we have $S(x+\eps e_\mu,x)=c(\eps e_\mu)$ and $\frac{e^2}2\operatorname{Bub}(C,C)(x,y)=\frac12B_\mu(x-y)$.

For the Ward form, symmetrise: $F=-\frac14\sum_\mu\sum\eps^6\varphi\Pi_\mu$. Then insert $\varphi(x,y)=\frac12\varphi(x,x)+\frac12\varphi(y,y)-\frac12\Delta\varphi(x,y)$. The first two terms give $\widehat\Pi_\mu(0)$, because $\Pi_\mu$ is even. In the third, $\Delta\varphi(x,x)=0$ and $\Pi_\mu=-B_\mu$ off the diagonal. The decomposition of $\Delta\varphi$ follows by writing $g(x,x)=g(x,y)+\alpha$, $g(y,y)=g(x,y)+\beta$, $D(x,x)=D(x,y)+\gamma$ and $D(y,y)=D(x,y)+\gamma'$.
\end{proof}

\begin{proposition}[per-anchor bound]\label{app:G1:B5e2-Wlocbd}
Let $w_T$, $\hat B$, $I=I(\ell_T)$ and the functionals $s_0(\cdot)$, $s_1(\cdot)$ be as in Lemma~\ref{app:K:L08-2}. Let $D\ge0$ be real symmetric, put $s_0:=s_0(D)$ and $s_1:=s_1(D)$, let the partition have any scale $t$, and put
\[
\hat S:=312\hat B+156I,\quad C_1:=4N_{\rm ov}c_{h1}(156\hat B+78I),\quad C_2:=104(N_{\rm ov}+1)\sqrt{18c_{h1}}\,\hat B,\quad C_3:=312(N_{\rm ov}+1)c_{h1}I.
\]
Then for every label $a$,
\[
\sum_{a'}\abs{F_{D,C}(a,a')}\le\tfrac32e^2(2t)^3\Big\{s_0w_T+\sqrt{s_0s_1}\,\hat S+(C_1+C_3)\frac{s_0}t+C_2s_0^{3/4}s_1^{1/4}t^{-1/2}\Big\}.
\]
\end{proposition}
\begin{proof}
Put $z=y-x$ and $q:=\min(\abs z_1^2s_1,4s_0)$. Write $A(y)-A(x)$ as a sum of $\abs z_1/\eps$ differences and apply Cauchy--Schwarz. This gives $\abs{D(x,y)}\le s_0$, $\Delta D\le q$ and $\abs{D(x,x)-D(x,y)}\le\sqrt{s_0q}$. Put $\delta h:=h(x)-h(y)$. At most $2N_{\rm ov}$ labels have $\delta h_{a'}\ne0$, so $\sum_{a'}\abs{\delta h_{a'}}\le2N_{\rm ov}m_1$, and $\sum_{a'}g(x,y)=h_a(x)+h_a(y)$. By Lemma~\ref{app:G1:B5e2-sym},
\[
\sum_{a'}\abs{\Delta\varphi}\le(h_a(x)+h_a(y))[q+2N_{\rm ov}\varpi\sqrt{s_0q}]+\abs{\delta h_a}[4N_{\rm ov}s_0\varpi+2\sqrt{s_0q}].
\]
We have $\abs{\delta h_a}\le\varpi(z)(\mathbf 1_{S_a}(x)+\mathbf 1_{S_a}(y))$, $\sum_{a'}g(x,x)=2h_a(x)$, and $\sum_x\eps^3h_a\le(2t)^3$. With Lemma~\ref{app:K:L08-2}(a), this yields
\[
\sum_{a'}\abs F\le\tfrac34(2t)^3[2e^2s_0w_T+\Sigma_A+4N_{\rm ov}s_0\Sigma_B+2(N_{\rm ov}+1)\Sigma_C],
\]
where $\Sigma_A,\Sigma_B,\Sigma_C$ are the suprema over $\mu$ of $\sum_z\eps^3\abs{B_\mu}w$ with the weights $w=q$, $\varpi^2$ and $\sqrt{s_0q}\,\varpi$, respectively.

On the fundamental cell, Lemma~\ref{app:K:L08-2}(b) gives $\abs{B_\mu}\le2e^2[(g_1+\rho_1)^2+(g_0+\rho_0)(g_2+\rho_2)]$ with $g_a\le\hat G_ar^{-a-1}$, and $\abs z_1\le3r$. The $\eps\Z^3$ parts are estimated with the shell sums of Lemma~\ref{app:K:L08-2}(e) and $\int_0^\infty\min(A,B/r^2)\,dr=2\sqrt{AB}$. The image and cross terms are estimated with $\min(9r^2s_1,4s_0)\le6r\sqrt{s_0s_1}$, $\varpi^2\le3c_{h1}r/t$ and $\sqrt{s_0q}\,\varpi\le6c_{h1}s_0r/t$. The result is
\[
\Sigma_A\le2e^2\sqrt{s_0s_1}\hat S,\qquad\Sigma_B\le2e^2c_{h1}(156\hat B+78I)/t,
\]
\[
\Sigma_C\le2e^2\big[52\sqrt{18c_{h1}}\hat Bs_0^{3/4}s_1^{1/4}t^{-1/2}+156c_{h1}s_0I/t\big].
\]
Collecting terms with $\frac34\cdot2e^2=\frac32e^2$ proves the bound. Every weight vanishes as $r\to0$, so no logarithm of $t/\eps$ appears.
\end{proof}

\begin{lemma}[members with a coarse scalar line]\label{app:G1:B5e2-CSloc}
Let $1\le k\le K$ and $0\le j<k$, and put $F^{\rm cs}_{j,k}:=F_{H_j,C}-F_{H_j,G_k}$. Let $s_0=s_0(H_j)$, $d_0=s_0(\mathcal C_k)$ and $d_1=s_1(\mathcal C_k)$, and let $\sigma_0,h,\gamma,\theta$ and $\Lambda_1$ be the constants of Lemma~\ref{app:K:L08-2}(c),(d). Then for every label $a$ and every partition scale $t$,
\[
\sum_{a'}\abs{F^{\rm cs}_{j,k}(a,a')}\le(2t)^3e^2\Big[5s_0d_0+3s_0d_1M^{-2}+6s_0\sqrt{d_0d_1}\Lambda_1+3d_1M^{-1}\sqrt{\norm{H_j}s_0d_0}+3\sqrt{d_0d_1}\sqrt{\norm{H_j}s_0d_0}\Big].
\]
Consequently $\abs u\sum_{j<k}\delta_j\sum_{a'}\abs{F^{\rm cs}_{j,k}(a,a')}\le(2t)^3\abs ue^2\kappa\mu_*^2L^2A_<t_k^{-2}$, where
\[
A_<:=5\sigma_0\gamma\eps_0+3\sigma_0\theta M^{-2}+6\sigma_0\sqrt{\gamma\theta}\Lambda_1\eps_0^{1/2}+\frac{3\theta\sqrt{h\sigma_0\gamma}\,\eps_0^{1/2}}{(L^{5/2}-1)M}+\frac{3\gamma\sqrt{\theta h\sigma_0}\,\eps_0}{L^{5/2}-1}.
\]
\end{lemma}
\begin{proof}
Write $D=H_j$ and $\mathcal C=\mathcal C_k$, so that $G_k=C-\mathcal C$. By bilinearity, $F^{\rm cs}_{j,k}(a,a')$ is
\[
\sum_\mu\Big\{-e^2\sum_x\eps^3h_ah_{a'}D(x,x)\mathcal C(x+\eps e_\mu,x)+e^2\sum_{x,y}\eps^6h_a(x)h_{a'}(y)D(x,y)J(x,y)\Big\},
\]
where $J:=C\nabla\nabla'\mathcal C+\mathcal C\nabla\nabla'C-\nabla C\nabla'\mathcal C-\nabla\mathcal C\nabla'C-\mathcal C\nabla\nabla'\mathcal C+\nabla\mathcal C\nabla'\mathcal C$. Call the seagull term $\tau_1$ and the six terms of $J$ $\tau_2$--$\tau_7$.

All terms except $\tau_3$ are bounded in absolute value, after summing over $a'$ with $\sum_{a'}h_{a'}=1$, by Lemma~\ref{app:K:L08-2}(c),(d):
\begin{itemize}
\item $\tau_1$ by $s_0d_0$, and $\tau_2$ by $d_1s_0M^{-2}$;
\item $\tau_4$ and $\tau_5$ by $s_0\sqrt{d_0d_1}\Lambda_1$ each;
\item $\tau_6$ by $d_1(\norm Ds_0)^{1/2}(M^{-2}d_0)^{1/2}$, and $\tau_7$ by $\sqrt{d_0d_1}(\norm Ds_0d_0)^{1/2}$.
\end{itemize}
Each of these carries a factor $3$ from the sum over $\mu$.

$\tau_3$ is summed over $\mu$ first. Since $\nabla_x\nabla'_yc(x-y)=-(\Delta_\mu c)(x-y)$ and $(-\Delta_\eps+M^2)c=\eps^{-3}\delta$, it equals $h_{a'}(x)D(x,x)\mathcal C(x,x)-M^2\sum_y\eps^3h_{a'}(y)D\mathcal C\,c(x-y)$, which is at most $2s_0d_0$ after the sum over $a'$. Finally $\sum_x\eps^3h_a\le(2t)^3$.

For the sum over $j$, insert $\delta_j\le\kappa\mu_*^2L^2t_j^2$, $s_0\le\sigma_0/t_j$, $\norm{H_j}\le ht_j^2$, $d_0\le\gamma t_k^{-2}$ and $d_1\le\theta t_k^{-3}$. Then use $\sum_{j<k}t_j\le t_k$, $\sum_{j<k}t_j^{5/2}\le t_k^{5/2}/(L^{5/2}-1)$ and $t_k\le\eps_0$.
\end{proof}

\begin{theorem}[order-$e^2$ anchor bounds]\label{app:G1:B5e2-thm}
Let $1\le k\le K-1$, $\abs u\le1$ and $M\in\{m_1,m_2\}$, and use the partition at scale $t=t_k$. With $\sigma_0,\theta_H$ as in Lemma~\ref{app:K:L08-2}(c) and $\hat S,C_1,C_2,C_3$ as in Proposition~\ref{app:G1:B5e2-Wlocbd}, put
\[
C_W:=12e^2\abs u\kappa\mu_*^2L^2\eps_0^2\big[\sigma_0w_T\eps_0+0.2654\sqrt{\sigma_0\theta_H}\hat S+(C_1+C_3)\sigma_0+2.415\,C_2\sigma_0^{3/4}\theta_H^{1/4}\big]
\]
and $C_{\rm loc}:=C_W+8\abs ue^2\kappa\mu_*^2L^2A_<$.
\begin{enumerate}
\item[(a)] $F^{=k}(a,a')=F^{=k}(a',a)$, and
\[
\sum_{a,a'}F^{=k}(a,a')=u\Big[\delta_kV(H_k;G_{k+1})+\sum_{j<k}\delta_j\big(V(H_j;G_{k+1})-V(H_j;G_k)\big)\Big].
\]
This is the part carried by graphs of top index exactly $k$; all its lines have index $\le k$.
\item[(b)] For every label $a$, the sums $\sum_{a'}\abs{F^{\ge k}(a,a')}$ and $\sum_{a'}\abs{F^{\ge k+1}(a,a')}$ are at most $C_{\rm loc}t_k$, and $\sum_{a'}\abs{F^{=k}(a,a')}\le2C_{\rm loc}t_k$.
\item[(c)] Let $c_G,c_H,\delta_1$ be as in Lemma~\ref{app:K:L08-3}, and fix $R_0>0$. Put $Z_p:=\sum_{n\ge0}2^{np}\mathrm{e}^{-\delta_1R_0(2^n-1)}$, $C_b:=2c_G^2(Z_1Z_3+Z_2^2)$ and $d_{aa'}:=\dist(S_a,S_{a'})$. If $d_{aa'}\ge R_0t_k$, then $\abs{F^{=k}(a,a')}\le C_{\rm pair}t_k^3\mathrm{e}^{-3\delta_1d_{aa'}/t_k}$, with $C_{\rm pair}:=384c_HC_be^2\abs u\kappa\mu_*^2L^2$. For every $R\ge R_0t_k$,
\[
\sum_{a':\,d_{aa'}\ge R}\abs{F^{=k}(a,a')}\le C_{\rm dp}t_k^3\mathrm{e}^{-2\delta_1R/t_k},\qquad C_{\rm dp}:=C_{\rm pair}\sum_{w\in\Z^3}\mathrm{e}^{-\delta_1(\abs w_\infty-2)_+}.
\]
If the scale-$k$ block containing the centre of $S_a$ lies in $\mathcal D_k$, and $S_{a'}$ meets $W_k^c$, then $d_{aa'}>(r(t_k)-3)t_k$.
\end{enumerate}
The constants $C_{\rm loc}$, $C_{\rm pair}$ and $C_{\rm dp}$ do not depend on $\eps$, $k$ or $K$. They are bounded for $\ell_T\ge\ell_0$, linear in $\abs u$ (no sign of $u$ is used), and finite at $M=m_1$.
\end{theorem}
\begin{proof}
(a) $F^{\ge k}$ and $F^{\ge k+1}$ are linear combinations of the $F_{H_j,S}$ with symmetric $S\in\{C,G_k,G_{k+1}\}$, so Lemma~\ref{app:G1:B5e2-sym} gives the symmetry and the totals. The difference of the totals is $u$ times
\[
\delta_kV(H_k;C)-\delta_k\big(V(H_k;C)-V(H_k;G_{k+1})\big)+\sum_{j<k}\delta_j\big(V(H_j;G_{k+1})-V(H_j;G_k)\big),
\]
and $G_{k+1}=\sum_{i\le k}G_{(i)}$.

(b) Let $k\le j<K$. Apply Proposition~\ref{app:G1:B5e2-Wlocbd} to $D=H_j$ at $t=t_k$, with $s_0\le\sigma_0/t_j$ and $s_1\le\theta_H/t_j^3$ (Lemma~\ref{app:K:L08-2}(c)), and multiply by $\delta_j$. The four terms become proportional to $t_j$, $1$, $t_j/t_k$ and $t_j^{1/2}t_k^{-1/2}$. Then use:
\begin{itemize}
\item $\sum_jt_j\le\eps_0$, $\#\{j\}\le\log_2(\eps_0/t_k)$ and $\sum_jt_j^{1/2}\le2.415\,\eps_0^{1/2}$;
\item the factor $(2t_k)^3=8t_k^3$, the bound $t_k\le\eps_0$, and $t_k^2\log_2(\eps_0/t_k)\le\eps_0^2/(2\mathrm{e}\ln2)\le0.2654\,\eps_0^2$.
\end{itemize}
This gives at most $C_Wt_k$. The members $j<k$ add at most $8t_k^3\abs ue^2\kappa\mu_*^2L^2A_<t_k^{-2}$, by Lemma~\ref{app:G1:B5e2-CSloc}. For $F^{\ge k+1}$ at the same scale $t_k$, the slices $j\ge k+1$ form a subset of those above. The coarse members carry $\mathcal C_{k+1}$, and Lemma~\ref{app:G1:B5e2-CSloc} at $k+1\le K$ gives $(2t_k)^3A_<t_{k+1}^{-2}\le8A_<t_k$, up to the common factor. The third bound follows from $F^{=k}=F^{\ge k}-F^{\ge k+1}$.

(c) If $d_{aa'}>0$, there is no seagull term, and
\[
F^{=k}(a,a')=u\frac{e^2}2\sum_\mu\sum_{x,y}\eps^6h_a(x)h_{a'}(y)\Big[\delta_kH_k\operatorname{Bub}_{k+1}+\sum_{j<k}\delta_jH_j(\operatorname{Bub}_{k+1}-\operatorname{Bub}_k)\Big](x,y),
\]
where $\operatorname{Bub}_{k'}:=\operatorname{Bub}(G_{k'},G_{k'})$. For $\abs{x-y}\ge d_{aa'}\ge R_0t_k$ we have $\sum_{i\le k}t_i^{-p}\mathrm{e}^{-\delta_1\abs{x-y}/t_i}\le Z_pt_k^{-p}\mathrm{e}^{-\delta_1d_{aa'}/t_k}$. With Lemma~\ref{app:K:L08-3}, this gives $\abs{\operatorname{Bub}_{k'}}\le C_bt_k^{-4}\mathrm{e}^{-2\delta_1d_{aa'}/t_k}$ for $k'\in\{k,k+1\}$. Moreover $\delta_j\abs{H_j}\le\kappa\mu_*^2L^2c_Ht_j\mathrm{e}^{-\delta_1d_{aa'}/t_k}$ for $j\le k$, and $\sum_{j\le k}t_j\le2t_k$. With $\sum_x\eps^3h_a\le(2t_k)^3$,
\[
\abs{F^{=k}}\le\tfrac32e^2\abs u(2t_k)^6\kappa\mu_*^2L^2c_H\cdot2t_k\cdot2C_bt_k^{-4}\mathrm{e}^{-3\delta_1d_{aa'}/t_k},
\]
and $\frac32\cdot64\cdot4=384$. For the partner sum, use $d_{aa'}\ge R$ and $d_{aa'}\ge(\abs{a'-a}_\infty-2)_+t_k$ in the minimal-image metric, and note that the minimal-image map from the label torus to $\Z^3$ is injective. For the last claim: every point of $S_a$ is within $\infty$-distance $t_k$ of the block, and every point of $S_{a'}$ is within $2t_k$ of $W_k^c$.

Uniformity: $w_T$, $\hat S$, $I$, $\Lambda_1$, $\sigma_0$, $\theta_H$, $\gamma$, $\theta$ and $h$ are free of $\eps$ and bounded for $\ell_T\ge\ell_0$ (Lemma~\ref{app:K:L08-2}). The only count of scales is absorbed by $t_k^2$. The factors $M^{-1}$ and $M^{-2}$ are finite at $m_1$.
\end{proof}
\noindent Each localisation of the seagull alone is of order $\eps^{-1}$. In Lemma~\ref{app:G1:B5e2-sym} this part cancels inside the single localisation $(a,a')$ against the zeroth moment of the bubble, and only the image sum $\widehat\Pi_\mu(0)$ remains.

\begin{lemma}[the order-$e^2$ family with one $D$-line]\label{app:G1:I1c-e2}
Let $1\le k\le K-1$, $\abs u\le1$ and $r(t_k)-4\ge R_0$. Let $C_{\rm loc}$ and $C_{\rm dp}$ be as in Theorem~\ref{app:G1:B5e2-thm}. For the order-$e^2$ vacuum family with one $D$-line at step $k$, the following hold.
\begin{enumerate}
\item[(i)] The sum over the bad localisations is at most $4C_{\rm loc}t_k\abs{\mathcal L_k}_k+2C_{\rm dp}t_k^3\mathrm{e}^{-2\delta_1(r(t_k)-4)}\abs{\mathcal D_k}_k$.
\item[(ii)] For every union $X$ of scale-$k$ blocks of $W_k^c$, the sum over the outside localisations with anchor in $X$ is at most $2C_{\rm loc}t_k\abs X_k$.
\end{enumerate}
Both bounds are proportional to $\abs u$.
\end{lemma}
\begin{proof}
Under convention (G), every moved piece of this family is localised once, at its own scale $k$, by the pair of labels of its two photon legs. This localisation is $F^{=k}(a,a')$ (Theorem~\ref{app:G1:B5e2-thm}(a)). A pair $(a,a')$ is bad if exactly one of $S_a,S_{a'}$ lies in $W_k$. By the symmetry in Theorem~\ref{app:G1:B5e2-thm}(a), the bad sum is at most $2\sum_{S_a\subset W_k}\sum_{S_{a'}\not\subset W_k}\abs{F^{=k}(a,a')}$. For $a$ with block in $\mathcal L_k$, use $\sum_{a'}\abs{F^{=k}}\le2C_{\rm loc}t_k$. For $a$ with block in $\mathcal D_k$, Theorem~\ref{app:G1:B5e2-thm}(c) gives $d_{aa'}>(r(t_k)-4)t_k\ge R_0t_k$, so its partner sum applies with $R=(r(t_k)-4)t_k$. For (ii), the anchor $a$ runs over $X$, and $\sum_{a'}\abs{F^{=k}}\le2C_{\rm loc}t_k$. The family carries exactly one $D$-line, so it is proportional to $u$.
\end{proof}

\paragraph{The region sum.}\label{app:G1:I145-combsec}
Notation at scale $k$ (Ba{\l}aban's large-field sets on $T_\eps$, all distances being torus distances):
\begin{itemize}
\item A \emph{tuple} $\mathfrak t_k=(P_v,Q_v,R_v,P_s,Q_s,R_s)$ is a $6$-tuple of subsets of $\Lambda_7^{(k-1)\prime}$ (scale-$(k+1)$ blocks, bonds and sites of scale $k$; $\Lambda_7^{(-1)\prime}:=T_1^{(0)}$), the set of violated restrictions. We write $\abs{\mathfrak t}$ for its total number of elements.
\item $\Lambda_0^{(k)}$ is the union of the large blocks contained in $\Lambda_7^{(k-1)\prime}$ at distance $\ge r(t_k)$ from the elements of $\mathfrak t_k$, and $\Lambda_{i+1}^{(k)c}$ is the union of the large blocks within distance $r(t_k)$ of $\Lambda_i^{(k)c}$ \cite[(2.7)--(2.8), p.~558; p.~570]{Ba2}.
\item Given $\Lambda_7^{(k-1)\prime}$, the set $\Lambda_0^{(k)}$ is a function of $\mathfrak t_k$. Hence a sequence $\mathfrak S=(\mathfrak t_0,\dots,\mathfrak t_{K-1})$ determines at most one admissible region sequence.
\item $\abs\cdot_k$ counts scale-$k$ points, and $N_i:=\abs{T_1^{(i)}}=\abs T/t_i^3$.
\end{itemize}

\begin{lemma}[geometric lemma]\label{app:G1:I145-geo}
Put $g(t):=8r(t)+8\sqrt d\,M_b+\sqrt d(L+1)$, $g_0:=8+(8\sqrt dM_b+\sqrt d(L+1))/R$, $C_r:=\sum_{n\ge0}L^{-n}(1+n\log L)^{\varrho}$ and $C_V:=(2g_0C_r+1)^d$. For every admissible region sequence generated by $(\mathfrak t_0,\dots,\mathfrak t_{K-1})$ and every $k<K$,
\[
\abs{\Lambda_7^{(k)c}}_k\le C_V\,r(t_k)^d\sum_{i\le k}\abs{\mathfrak t_i}.
\]
\end{lemma}
\begin{proof}
(i) The set $\Lambda_0^{(k)c}$ is contained in the union of the large blocks not contained in $\Lambda_7^{(k-1)\prime}$ and the large blocks at distance $<r(t_k)$ from an element of $\mathfrak t_k$. Moreover $\Lambda_7^{(k)c}$ lies in the $(7r(t_k)+7\sqrt dM_b)$-neighbourhood of $\Lambda_0^{(k)c}$ (scale-$k$ units). A point of the second kind lies within $r(t_k)+\sqrt dM_b+\sqrt dL$ of an element of $\mathfrak t_k$. A point $z$ of the first kind lies in a large block containing a point $z'\notin\Lambda_7^{(k-1)\prime}$, with $\abs{z-z'}\le\sqrt dM_b$; the unit block of $z'$ contains a scale-$(k-1)$ point of $\Lambda_7^{(k-1)c}$, at distance $\le\sqrt d$.
(ii) By induction on $k$, every point of $\Lambda_7^{(k)c}$ lies within physical distance $D(i,k):=\sum_{m=i}^kt_mg(t_m)$ of an element of some $\mathfrak t_i$, $i\le k$.
(iii) Since $r\ge R$, $g(t)\le g_0r(t)$. For $t_k\le1$, $1+\log t_{k-n}^{-1}=(1+\log t_k^{-1})+n\log L\le(1+\log t_k^{-1})(1+n\log L)$, so $r(t_{k-n})\le(1+n\log L)^{\varrho}r(t_k)$. With $t_{k-n}=L^{-n}t_k$ this gives $D(i,k)\le g_0C_rt_kr(t_k)$.
(iv) A physical ball of radius $D$ contains at most $(2D/t_k+1)^d$ scale-$k$ points of $\Z^3$, hence at most that many of $T_1^{(k)}$, since projection to the torus does not increase the count. As $r\ge R\ge1$, this is $\le C_Vr(t_k)^d$. A union bound over the elements of the tuples concludes.
\end{proof}

\begin{lemma}[region-sum lemma]\label{app:G1:I145-comb}
Hypothesis (H-LF) at a point $\pi$: the large-field region sum of the closing estimate (\cite[(2.117), p.~582; (3.40)--(3.42), pp.~591--592]{Ba2} for the slice model) is bounded by $\mathrm{e}^{C_0(\pi)\abs T}$ times
\[
\sum_{\mathfrak S=(\mathfrak t_0,\dots,\mathfrak t_{K-1})\ \rm admissible}\ \prod_{k<K}\mathrm{e}^{-c_*(\pi)p(t_k)^2\abs{\mathfrak t_k}}\ \exp\Big(c_A(\pi)\sum_{k<K}(1+\log t_k^{-1})\abs{\Lambda_7^{(k)c}}_k\Big),
\]
with $c_*(\pi)>0$, $c_A(\pi)<\infty$ and $C_0(\pi)$ independent of $\eps$, $\ell_T$ and $\mathfrak S$. Assume $2p>d\varrho+2$ and
\[
\text{(E.comb)}\qquad \hat\ell^{\,2p-d\varrho-2}\ge\frac{2c_A(\pi)C_VR^d}{c_*(\pi)b_0^2\log L},\qquad \hat\ell:=1+\log\eps_0^{-1}.
\]
Then, uniformly in $\eps=L^{-K'}$ and $\ell_T$, the sum is at most $\exp\big(2(d+2)\Sigma_*(\eps_0;\pi)\abs T\big)$, where
\[
\Sigma_*(\eps_0;\pi):=\sum_{n\ge n_0}L^{3n}\exp\big(-\tfrac12c_*(\pi)b_0^2(1+n\log L)^{2p}\big),\qquad n_0:=\log_L\eps_0^{-1}.
\]
$\Sigma_*$ is finite, contains no $c_A$, and decreases to $0$ as $\eps_0\downarrow0$. At $p=6$, $\varrho=2$, $d=3$ we have $2p-d\varrho-2=4$.
\end{lemma}
\noindent The constants $c_A,c_*,C_V$ are evaluated at the bound $\eps_1$ (stage 10a of Appendix~\ref{app:B}) before $\eps_0$ is chosen by (E.comb) at stage 10b, so there is no cycle. The hypothesis (H-LF) is the form of the region sum in the closing step of the proof of Theorem~\ref{thm:T6} (Lemma~\ref{app:G1:S05-CL}).
\begin{proof}
(1) \emph{Costs to elements.} Put $x_i:=\log t_i^{-1}$. For $k\ge i$ we have $t_k\ge t_i$, so $1+\log t_k^{-1}\le1+x_i$ and $r(t_k)\le R(1+x_i)^{\varrho}$. By Lemma~\ref{app:G1:I145-geo},
\[
\sum_{k<K}(1+\log t_k^{-1})\abs{\Lambda_7^{(k)c}}_k\le C_VR^d\sum_i\abs{\mathfrak t_i}(K-i)(1+x_i)^{1+d\varrho},
\]
and $K-i=\log_L(\eps_0/t_i)\le(1+x_i)/\log L$ since $\eps_0\le1$.
(2) \emph{Absorption.} The summand is therefore at most $\prod_i\exp\big(-[c_*b_0^2(1+x_i)^{2p}-c_AC_VR^d(\log L)^{-1}(1+x_i)^{d\varrho+2}]\abs{\mathfrak t_i}\big)$. Since $2p>d\varrho+2$, $x\mapsto(1+x)^{2p-d\varrho-2}$ is increasing, and $x_i\ge\log\eps_0^{-1}$. So (E.comb) makes the bracket at least $\frac12c_*b_0^2(1+x_i)^{2p}=\frac12c_*p(t_i)^2$.
(3) \emph{Entropy.} By determinism, each $\mathfrak S$ occurs at most once. The region expansion \cite[(2.9)]{Ba2} holds in the ``$\le$'' direction, which is the only one used, and dropping admissibility only adds non-negative terms. Hence, with $a_i:=\frac12c_*p(t_i)^2$,
\[
\sum_{\mathfrak S}\prod_i\mathrm{e}^{-a_i\abs{\mathfrak t_i}}\le\prod_i(1+\mathrm{e}^{-a_i})^{n_i}\le\exp\Big(\sum_in_i\mathrm{e}^{-a_i}\Big),
\]
where $n_i$ is the number of possible elements of a tuple at scale $i$: $P_v,P_s$ give at most $N_i/L^d$ each, $Q_v,Q_s$ at most $dN_i$ each, and $R_v,R_s$ at most $N_i$ each, so $n_i\le2(d+2)N_i$.
(4) \emph{Scales.} $t_i=L^{-n}$ with distinct integers $n\ge n_0$ for distinct $i$, and $N_i=L^{3n}\abs T$, so $\sum_in_i\mathrm{e}^{-a_i}\le2(d+2)\abs T\,\Sigma_*(\eps_0;\pi)$. $\Sigma_*$ converges because $2p>1$.
\end{proof}
\begin{remark}\label{app:G1:I145-combuse}
Lemma~\ref{app:G1:I145-comb} holds for every $c_A$. Lemma~\ref{app:G1:IF1abs} uses it with $c_A\mapsto c_A+C_{\rm old}$. Under convention (G), the cost of the vacuum counterterm on ${\rm Col}_k$ is contained in $C_{\rm old}$.
\end{remark}

\paragraph{Ba{\l}aban's Proposition~1 with its sketched steps written out.}\label{app:G1:I145-prop1basesec}
\begin{proposition}[Ba{\l}aban's Proposition~1 on the torus]\label{app:G1:I145-prop1base}
Let $\nbar$, $\alpha_0$, $k$, $\eta$, the domains $\Omega,\Omega_1,\Omega_2$, $\tilde B$, the renormalised families and the external fields be as in Proposition~\ref{app:K:L08-5}. Then the bound \cite[Prop.~1, (1.33), p.~420]{Ba3} holds on $T_\eps$, and on King's torus with photon mass $\mu_*>0$ and $D$-coloured lines. It holds with a tree rate $\delta_T>0$ and a constant $O(1)=O(\alpha_0,\nbar;d,L,a,M_b)$, both independent of $\eps$, $k$, the domains and $\tilde B$. The rate $\delta_T$ depends on $\nbar$ through the number of lines; it is not asserted to depend on $d$ only. The coloured families with $m=0$, $m\ge2$, or $m=1$ and $o\ge2$ carry the extra factor $(\abs uc_k)^m$. This proposition is Proposition~\ref{app:K:L08-5}.
\end{proposition}
\noindent The proposition rests on the written formulas of \cite{Ba3} (including p.~433) and on \cite{Ba4}. Two items are treated separately, in Proposition~\ref{app:G1:I3}:
\begin{itemize}
\item the two-variable form of the remainder kernel of \cite[(1.16), (2.5)]{Ba3} (Lemmas~\ref{app:G1:I3r-U1m} and \ref{app:G1:I3r-U1mOmega});
\item the remainder on the cube $\square$ of \cite[p.~433]{Ba3} (Lemma~\ref{app:G1:I3r-cube}).
\end{itemize}

\begin{remark}[the sketched steps of {\cite[Prop.~1]{Ba3}}]\label{app:G1:I145-prop1base-rem}
The steps of \cite[pp.~418--445]{Ba3} that are only sketched there fall into the items below. For each item we name the results of this paper that treat it, and, where there is one, the closest step of Table~\ref{tab:prop1}. The domain conditions of all items are Proposition~\ref{prop:prop1} and Corollary~\ref{cor:prop1R4}.
\begin{itemize}
\item \emph{I1} (p.~432, family combinatorics, ``inspecting all possible cases''; cf.\ S10). For the vacuum families of the moved pieces, the classes, their cut table and the one-particle-reducible C2 chains are given by Proposition~\ref{app:G1:I1c-table}, Lemma~\ref{app:G1:cutclasses} and Lemma~\ref{app:G1:I1c-1PR}. The closure of the classes under the operations of pp.~434--445 in general, including the nested cases, is part of Proposition~\ref{app:K:L08-5}.
\item \emph{I2} (p.~422, bad localisations). The vacuum part is Lemma~\ref{app:G1:I145-perloc} with Proposition~\ref{app:G1:R9vac}, together with Lemma~\ref{app:G1:I1c-e2} for the order-$e^2$ family with one $D$-line. The mass part is the same localisation sum for the two-leg classes (Lemma~\ref{app:G1:VS}).
\item \emph{I3} (p.~418, pointwise counterterm bounds). The final counterterm bounds are Propositions~\ref{app:G1:I3b} and \ref{app:G1:Efin}. The bounds at the intermediate scales are part of Proposition~\ref{app:K:L08-5}.
\item \emph{I4} (\cite[Props.~2.1--2.3]{Ba1}, behind (2.10)--(2.12)). This item is cited, as Proposition~\ref{prop:UV-BaProp21}. The slices at zero background are also Theorem~\ref{app:G1:VSR-VS}.
\item \emph{I5} (p.~412, H\"older norms of the background). The H\"older bounds of the vector and scalar backgrounds in the norm \cite[(1.32)]{Ba3} are part of Proposition~\ref{app:K:L08-5}.
\item \emph{(3.35)} (p.~443, leg transport; inside S21). This is Theorem~\ref{app:G1:VSR-T335}, with the pinned case \cite[(2.21e)]{Ba3} in Proposition~\ref{app:G1:VSR-5b}.
\item \emph{Rates and constants} (\cite[p.~427]{Ba3}). They are $\delta^{(i+1)}=\delta^{(i)}/(2(n'+1))$ and
\[
C_A\le c_d^m(4/\delta_1)^{dm}(4m+2)^{dm(m-1)/2}(1-L^{-\vartheta})^{-m},
\]
independent of $\eta$, $k$, $\eps$, $x$ and $\ell_T$ (Lemma~\ref{app:G1:VSR-A}).
\end{itemize}
\end{remark}

\paragraph{Real backgrounds.}\label{app:G1:I145-I4check}
\begin{lemma}[real backgrounds and unitarity]\label{app:G1:I145-unitary}
For $w\in\C$ put $U_w:=\mathrm{e}^{\eta qw}$, so that $U(A)=U_{eA}$, and a complex charge $ze$, $z\in\C$, gives $U_{zeA}$. Let $R_{n+1}$ be as in Definition~\ref{app:G1:I145-defbg}(c).
\begin{itemize}
\item[(a)] $\norm{U_w}=\mathrm{e}^{\eta\abs{\Im w}}$. Hence $U_w\in\SO(2)$ if and only if $w\in\R$; $U(A)$ is orthogonal if and only if $A\in\R$; and for real $A\ne0$ and $z=\mathrm{e}^{i\alpha}$ with $\sin\alpha\ne0$, $U_{zeA}$ is not orthogonal.
\item[(b)] For $n\ge0$ and $w\in\R$, $\norm{R_{n+1}(qw)}\le1$. For every $s>0$, $\sup_{\abs{\Im w}\le s}\norm{R_{n+1}(qw)}\ge2^{-(n+1)}\mathrm{e}^{s/2}$.
\item[(c)] The matrix entries of $U_w$ and $R_{n+1}(qw)$ are entire in $w$ with real Taylor coefficients; so $U_{\bar w}=\overline{U_w}$, $R_{n+1}(q\bar w)=\overline{R_{n+1}(qw)}$, and both are real for real $w$.
\end{itemize}
Consequently an estimate proved through $\norm U=1$, such as \cite[(3.6)--(3.7), p.~255; (3.17), p.~257]{BFS1}, or through $\norm{R_{\nbar+1}(qA)}\le1$ \cite[p.~615]{Ba1}, \cite[p.~414]{Ba3}, gives no information at a non-real background or at a complex charge.
\end{lemma}
\begin{proof}
With $S$ the unitary matrix with columns $(1,\mp i)/\sqrt2$, $q=S\,\mathrm{diag}(i,-i)S^{-1}$, so $f(qw)=S\,\mathrm{diag}(f(iw),f(-iw))S^{-1}$ is normal with norm $\max(\abs{f(iw)},\abs{f(-iw)})$ for every entire $f$. (a) $\abs{\mathrm{e}^{\pm i\eta w}}=\mathrm{e}^{\mp\eta\Im w}$; for real $w$, $U_w$ is the rotation by $\eta w$; for $w=zeA$, $\Im w=eA\sin\alpha$. (b) For real $w$, $\abs{R_{n+1}(\pm iw)}\le(n+1)\int_0^1(1-t)^n\,dt=1$. At $w=is$ the eigenvalue $R_{n+1}(s)$ is real and $R_{n+1}(s)\ge(n+1)\int_{1/2}^1(1-t)^n\mathrm{e}^{s/2}\,dt=2^{-(n+1)}\mathrm{e}^{s/2}$. (c) $U_w=\sum_j(\eta q)^jw^j/j!$ and $R_{n+1}(qw)=\sum_j\frac{(n+1)!}{(n+1+j)!}q^jw^j$, since $(n+1)\int_0^1(1-t)^nt^j\,dt=\frac{(n+1)!\,j!}{(n+1+j)!}$.
\end{proof}

\begin{proof}[Second proof of Proposition~\ref{app:G1:I4}]
\emph{Part 1 (the printed literature).} Each item below is a statement about printed pages.
\begin{itemize}
\item[(B)] \cite{Ba4}: vector fields are real-valued, $U(A)=\mathrm{e}^{qe\eta A}$, $e$ real \cite[p.~572]{Ba4}; the only complex object is the contour shift of the free $A=0$ propagator \cite[p.~586]{Ba4}. In \cite{Ba1,Ba2,Ba3} every background is real, and $R_{\nbar+1}$ is used at real $A$, where $\norm{R_{\nbar+1}(qA)}\le1$ (Lemma~\ref{app:G1:I145-unitary}(b)). \cite[(1.16), p.~414; (1.32)--(1.33), p.~420]{Ba3} are stated at real backgrounds. At the cited pages of \cite{K1}, \cite{K2} (\cite[pp.~653--657, 659]{K1}, \cite[pp.~334, 336]{K2}) the only complex objects are sources.
\item[(C)] The cumulant lemma \cite[p.~152]{BCGNOPS} concerns real Gaussian variables, with no gauge field; it is applied in \cite[p.~563, (2.28)--(2.31)]{Ba2} to ``a positive operator'' with ``a probabilistic Gaussian measure''. \cite[(3.17)]{BFS1} is proved via $\norm{U^\Phi(g)}=1$ for configurations $g$ of a compact gauge group \cite[(3.7), p.~255]{BFS1} and applied in \cite[p.~591, (3.37)]{Ba2} at the real field $A^\eps$, the variable of the outer real integral, where $U(A^\eps)$ is orthogonal (Lemma~\ref{app:G1:I145-unitary}(a)).
\item[(D)] The two unprinted steps are stated at real backgrounds: \cite[p.~414]{Ba3}, with $\tilde A$ the real Gaussian fluctuation, and \cite[p.~433]{Ba3}, with $\tilde B=B_0+B'$ real. Their proofs in this paper (Lemmas~\ref{app:G1:I3r-U1mOmega}, \ref{app:G1:I3r-SEP}, \ref{app:G1:I3r-cube}) use real backgrounds only. The p.~433 split presupposes the inclusion of sets $\square\subset\Omega$ (Proposition~\ref{prop:prop1}(c), Lemma~\ref{app:G1:I3r-Wsq}); the domain $\Omega_2$ enters only through the replacement of Lemma~\ref{app:G1:I3r-W2R}, at the real background $\tilde B$.
\end{itemize}
\emph{Part 2 (the results of this paper).} By Lemma~\ref{app:G1:I145-realbg}, every operator of Definition~\ref{app:G1:I145-defbg}(a) evaluated in the construction is evaluated at a real field or at $A=0$, and every Gaussian measure is a real Gaussian measure with real symmetric positive definite precision; by Lemma~\ref{app:G1:I145-closure}, the operations of the construction keep backgrounds real, and $B_0$ in (D) may be taken real. The statements of this appendix used in the proof of Theorem~\ref{thm:T6} concern these objects (proof of Proposition~\ref{app:G1:I4}(i)). Bounds at phase-rotated couplings $(ze,z\lambda)$ with $\abs z=1$, where $U$ is not orthogonal (Lemma~\ref{app:G1:I145-unitary}(a)), are not used: Theorem~\ref{thm:T6} is evaluated at real couplings and at $u\in\{1,-\eta\}$.
Part~1 (B), (C) and Part~2 give (i) and (ii); Part~1 (B), (C) give (iii); Part~1 (D) and Lemma~\ref{app:G1:I145-closure}(c) give (iv).
\end{proof}

\begin{table}[ht]
\centering\footnotesize
\begin{tabular}{@{}p{0.22\textwidth}p{0.45\textwidth}p{0.25\textwidth}@{}}
\hline
use & what is used & type \\
\hline
Cauchy estimates in the fluctuation field & the fluctuation bound (real polynomial $V$, real symmetric $P$), the norms \cite[(1.32)]{Ba3}, the bound \cite[(1.16)]{Ba3} & U0: no complex background \\
analyticity in the background & an analyticity lemma and its use & U0 + U1; not used \\
pointwise bounds on a route through Cauchy estimates & interior points & U1 / U2--U3 at interior points; not used \\
the large-field step & real expansions; an optional Cauchy estimate & U0; the optional part U2 at interior points, not used \\
conditions of the route through Cauchy estimates & upper bounds on $\eps_0$ of that route & not used \\
lists and renamings & --- & no estimate \\

Cauchy estimates in sources and in $(u,b)$ & --- & not a gauge background \\

\hline
\end{tabular}
\caption{The uses in the proof of Theorem~\ref{thm:T6} of bounds that involve complex parameters. U0: no complex background; U1: at the level of $L^2$ norms or forms; U2/U3: pointwise or derivative bounds; ``at interior points'': only at points of depth $\ge R_0$.}\label{app:G1:I145-ledger}
\end{table}

\subsection{The conditions of Ba{\l}aban's Proposition~1 at its applications}\label{app:G1:part-E5}

We prove (a) three conventions of Ba{\l}aban's papers on domains and supports, Lemmas~\ref{app:G1:Rnest} (nesting of the domains), \ref{app:G1:Rthetab} (the meaning of $\theta_b$) and \ref{app:G1:Rloc7} (localisation in $\Lambda_7$), and (b) Proposition~\ref{app:G1:I3}, which assembles the conditions of \cite[Prop.~1]{Ba3} at its applications in \cite{Ba2} from the statements of Table~\ref{app:G1:I3tab}. The support hypothesis in the form (R4$'''$) is Corollary~\ref{cor:prop1R4}. Lemma~\ref{app:G1:I3a-eta1} and the statements of the last row of the table are proved in this subsection, before Proposition~\ref{app:G1:I3}. Proposition~\ref{app:G1:twodomain} is proved in Subsection~\ref{app:G1:part-E2}, and Proposition~\ref{app:G1:T579} and Lemmas~\ref{app:G1:AA}, \ref{app:G1:AA4} in Subsection~\ref{app:G1:part-D3}. The remaining statements are proved in Subsection~\ref{app:G1:part-TI3}.

\paragraph{Setting.}\label{app:G1:ssec-I3setting}
We work in the $k$-th step, $k\ge1$, of \cite[Chap.~2]{Ba2}, $d=3$, on the torus $T_\eps$ of \cite[(1.2)]{Ba1}, at $L=2$. Distances are those of the torus sup-metric \cite[(1.3), p.~604]{Ba1}, $\abs{x-y}=\max_\mu\min\{\abs{x_\mu-y_\mu},2L_\mu-\abs{x_\mu-y_\mu}\}$, and the distance of two sets is the infimum of point distances (``inf convention''; only the point metric is printed). We write $r_j:=r(L^j\eps)$ and $M_b$ for Ba{\l}aban's large-block factor $M$. For $\Lambda\subset T_1^{(j)}$ (the unit lattice of scale $j$), $\Lambda':=\Lambda\cap L\,T_1^{(j)}$ and $B(y)$ is the block with base point $y$ \cite[(1.16)--(1.17), p.~606]{Ba1}. $B^k(y)$ is the scale-$k$ block in $T_\eta$, $\eta=L^{-k}$, and $B^k(X):=\bigcup_{y\in X}B^k(y)$. Large blocks are unions of blocks \cite[p.~605]{Ba1}. The sets $\Lambda_i^{(j)}\subset T_1^{(j)}$ are defined by \cite[(2.7)--(2.8), p.~558]{Ba2} with $r(\eps)$ replaced by $r_j$ \cite[p.~566]{Ba2}: $\Lambda_0^{(j)c}$ is a union of large blocks, and $\Lambda_{i+1}^{(j)c}$ is the union of all large blocks whose distance from $\Lambda_i^{(j)c}$ is at most $r_j$. The function $\theta_k$ on $T_\eta$ equals $1$ on $B^{k-1}(\Lambda_2^{(k-1)})$ and varies from $1$ to $0$ on a slice of thickness $<M_b$ around it \cite[pp.~566--567]{Ba2}. In \cite{Ba3}, $\Omega=B^k(\Lambda_2^{(k-1)\prime})$ \cite[p.~411]{Ba3}, $\Omega_1=B^k(\Lambda_7^{(k-1)\prime})$, and $g_k$ is smooth with $\supp g_k\subset\Omega_1$ \cite[p.~412]{Ba3}.

On the path of Theorem~\ref{thm:T6} the propagator is changed to the domain $\Omega_2:=B^k(\Lambda_2^{(k)})$ at three places.
\begin{itemize}
\item \emph{U1a}, in the polynomial $Q^{(k)}$ \cite[p.~579]{Ba2}: $G_k(B^k(\Lambda_0^{(k)}),\tilde B)\to G_k(B^k(\Lambda_2^{(k)}),B^{(k+1),\eta})$, where $\tilde B=(1-\theta_{k+1})A^{(k)}+\theta_{k+1}B^{(k+1),\eta}$ on $B^k(\Lambda_0^{(k)})$.
\item \emph{U1b}, in the interaction terms \cite[p.~580, Prop.~2.6]{Ba2}: $G_k(B^{k-1}(\Lambda_2^{(k-1)}),\tilde B)\to G_k(B^k(\Lambda_2^{(k)}),B^{(k+1),\eta})$.
\item \emph{U2}, at scale $j$ \cite[p.~578]{Ba2}.
\end{itemize}

\paragraph{Nesting of the domains.}
\begin{lemma}[nesting of the domains]\label{app:G1:Rnest}
Fix $k\ge1$ and an admissible sequence $(\Lambda_0^{(j)})_j$ of \cite[(2.43)]{Ba2}.
\begin{enumerate}
\item[(i)] For every $j$ and $i\ge0$: $\Lambda_{i+1}^{(j)}\subset\Lambda_i^{(j)}$. This holds under every block-to-set distance convention for which a block contained in a set has distance $0$ from it.
\item[(ii)] $\Lambda\mapsto\Lambda'$ and $X\mapsto B^k(X)$ are monotone, and $B^k(\Lambda')=B^{k-1}(\Lambda)$ for every $\Lambda\subset T_1^{(k-1)}$ that is a union of blocks.
\item[(iii)] At U1b: $\Omega_2=B^k(\Lambda_2^{(k)})\subset B^k(\Lambda_0^{(k)})\subset\Omega_1\subset\Omega$. Moreover $\Omega=B^{k-1}(\Lambda_2^{(k-1)})$, the old domain of U1b.
\item[(iv)] $B^k(\Lambda_2^{(k)})\subset B^l(\Lambda_2^{(l)})$ for every $l\le k$.
\item[(v)] (Margins, scale-$k$ units, inf convention.) If $x\in\Lambda_{i+s}^{(j)}$ then $\dist(x,\Lambda_i^{(j)c})\ge s\lfloor r_j\rfloor+1$. Hence $\dist(\Omega_2,\Omega_1^c)>2r_k-2$ and $\dist(\Omega_1,\Omega^c)>(5r_{k-1}-4)/L-1$. Under any other block-to-set convention (block centre, representative point) each step \cite[(2.8)]{Ba2} loses at most one large-block diameter $M_b$ of its scale, which gives $>2(r_k-M_b)-2$ and $>(5(r_{k-1}-M_b)-4)/L-1$.
\end{enumerate}
The constants depend on $L$ and $M_b$ only.
\end{lemma}
\noindent The lemma uses \cite[p.~566]{Ba2}. Parts (i)--(iv) do not depend on the convention for the distance between sets. Part (v) is stated under the infimum convention, with the loss of at most $M_b$ per step under any other convention; this loss is absorbed by $r\ge R\gg M_b$ and covered by \eqref{eq:eppp}.

\begin{proof}
(i) By \cite[(2.7)]{Ba2}, $\Lambda_0^{(j)c}$ is a union of large blocks, and by \cite[(2.8)]{Ba2} so is every $\Lambda_{i+1}^{(j)c}$. Suppose $\Lambda_i^{(j)c}$ is a union of large blocks. Each of these blocks has distance $0\le r_j$ from $\Lambda_i^{(j)c}$ (under the inf, sup or centre conventions alike), so it belongs to $\Lambda_{i+1}^{(j)c}$. Thus $\Lambda_i^{(j)c}\subset\Lambda_{i+1}^{(j)c}$, and induction on $i$ gives (i).

(ii) Monotonicity is immediate from the definitions. For $y\in L\,T_1^{(k-1)}$ the scale-$k$ block is $B^k(y)=\bigcup_{x\in B(y)}B^{k-1}(x)$ (iterated blocks, \cite[(1.17)]{Ba1}). If $\Lambda=\bigcup_{y\in\Lambda'}B(y)$, then $B^k(\Lambda')=\bigcup_{y\in\Lambda'}\bigcup_{x\in B(y)}B^{k-1}(x)=B^{k-1}(\Lambda)$.

(iii) The definition of ``admissible'' in \cite[p.~566]{Ba2} contains the clause ``In general $\Lambda_0^{(j+1)}\subset\Lambda_7^{(j)\prime}$, and either $\Lambda_0^{(j+1)}$ is a maximal set composed of big blocks and satisfying this inclusion, or \dots''. This is the general-$j$ analogue of the $j=0$ condition printed just before it, and it has the same either/or form. It is therefore a condition on every term of the expansion \cite[(2.43)]{Ba2}, not a hedge: if ``in general'' meant ``usually'', the either/or clause of the definition would also be optional. With $j=k-1$: $\Lambda_0^{(k)}\subset\Lambda_7^{(k-1)\prime}$. Now
\[
\Omega_2=B^k(\Lambda_2^{(k)})\overset{\rm (i),(ii)}{\subset}B^k(\Lambda_0^{(k)})\overset{\text{p.~566},\,{\rm (ii)}}{\subset}B^k(\Lambda_7^{(k-1)\prime})=\Omega_1\overset{\rm (i),(ii)}{\subset}B^k(\Lambda_2^{(k-1)\prime})=\Omega.
\]
Since $\Lambda_2^{(k-1)}$ is a union of large blocks, (ii) gives $\Omega=B^{k-1}(\Lambda_2^{(k-1)})$.

(iv) By the chain in (iii), $B^k(\Lambda_2^{(k)})\subset B^k(\Lambda_2^{(k-1)\prime})=B^{k-1}(\Lambda_2^{(k-1)})$. Iterating in $k$ gives (iv).

(v) Let $x\in\Lambda_{i+1}^{(j)}$. The large block containing $x$ is not in $\Lambda_{i+1}^{(j)c}$, so its distance from $\Lambda_i^{(j)c}$ is $>r_j$. Hence no lattice point $z$ with $\abs{z-x}\le\lfloor r_j\rfloor$ lies in $\Lambda_i^{(j)c}$; that is, the closed sup-ball of radius $\lfloor r_j\rfloor$ about $x$ lies in $\Lambda_i^{(j)}$. On the cyclic lattice, sup-balls of integer radii add coordinatewise. Iterating $s$ times therefore puts the ball of radius $s\lfloor r_j\rfloor$ about $x\in\Lambda_{i+s}^{(j)}$ inside $\Lambda_i^{(j)}$, and so $\dist(x,\Lambda_i^{(j)c})\ge s\lfloor r_j\rfloor+1$.

For $x\in B^k(y)$ and $z\in B^k(y')$, $\abs{x-z}>\abs{y-y'}-1$ in scale-$k$ units. With $y\in\Lambda_2^{(k)}$ and $y'\notin\Lambda_0^{(k)}$ ($s=2$), and using $\Omega_1\supset B^k(\Lambda_0^{(k)})$ by (iii), we get $\dist(\Omega_2,\Omega_1^c)\ge2\lfloor r_k\rfloor>2r_k-2$. Likewise, $\Omega_1=B^{k-1}(\Lambda_7^{(k-1)})$ and $\Omega=B^{k-1}(\Lambda_2^{(k-1)})$ by (ii). With $s=5$, in scale-$(k-1)$ units their distance is $\ge5\lfloor r_{k-1}\rfloor$, i.e.\ $\ge5\lfloor r_{k-1}\rfloor/L>(5r_{k-1}-5)/L\ge(5r_{k-1}-4)/L-1$ in scale-$k$ units, because $L\ge1$. Under a non-inf convention the block of $x$ may be declared at distance $\le r_j$ although its points are up to one large-block diameter $M_b$ farther away. Replacing $r_j$ by $r_j-M_b$ in each step gives the stated bounds.
\end{proof}

\begin{remark}[scope of the order $\Omega_1\supset\Omega_2$]\label{app:G1:Rnest-scope}
The order $\Omega\supset\Omega_1\supset\Omega_2$ holds at U1b. At U1a the new localisation has $\supp g_k\subset\Lambda_7^{(k)}$ \cite[p.~579]{Ba2}, so there the order is $\Omega=B^k(\Lambda_0^{(k)})\supset\Omega_2\supset B^k(\Lambda_7^{(k)})$. At U2 the vertices lie in $\supp\theta_{k+1}$ \cite[p.~578]{Ba2}. No use in the proof needs $\Omega_1\supset\Omega_2$ at U1a or U2; so the step S17 of Table~\ref{tab:prop1} is used only in the form U1b. Lemma~\ref{app:G1:Rnest} does not need the identification of the $\Omega_2$ of \cite{Ba3} with $B^k(\Lambda_2^{(k)})$: \cite[Prop.~1]{Ba3} is uniform over $\Omega_2\subset\Omega$ ``satisfying the same conditions as $\Omega$'' \cite[p.~414]{Ba3}, its constants are independent of $\Omega,\Omega_1,\Omega_2$ \cite[p.~421]{Ba3}, and $B^k(\Lambda_2^{(k)})$ is a union of big blocks.
\end{remark}

\paragraph{The meaning of $\theta_b$.}
\begin{lemma}[the meaning of $\theta_b$]\label{app:G1:Rthetab}
Let $\Omega_2=B^k(\Lambda_2^{(k)})$, and call a bond of $T_\eta$ \emph{interior} if both its end points lie in $\Omega_2$.
\begin{enumerate}
\item[(i)] $\theta_{k+1}=1$ at every point of $\Omega_2$.
\item[(ii)] For every vector field $A$ and every interior bond $b$, $(\theta_{k+1}A)_b=A_b$ and $((1-\theta_{k+1})A)_b=0$. This holds for the printed identification of vector fields with bond functions and for every local convention for the product of a point function with a vector field (one in which the value on $b$ is a function of $\theta(b_-),\theta(b_+)$ equal to $1$ when both are $1$). In particular $\tilde B_b=B^{(k+1),\eta}_b$ on interior bonds.
\item[(iii)] $G_k(\Omega_2,\tilde B)=G_k(\Omega_2,B^{(k+1),\eta})$ and $Q_k(\tilde B)=Q_k(B^{(k+1),\eta})$ on $\Omega_2$.
\item[(iv)] At U1b, the background summands $(1-\theta_{l+1})\theta_lA^{(l)}$, $l+1\le k+1$, of \cite[(2.93), p.~576]{Ba2} vanish on interior bonds. The vertex-borne field $B^{(k+1),\eta}$ of \cite[(2.107)]{Ba2} enters only on interior bonds.
\end{enumerate}
\end{lemma}
\noindent The lemma uses \cite[p.~604]{Ba1} and \cite[p.~567]{Ba2}, and the standard meaning, not spelled out in \cite{Ba1}, that the Neumann form on $\Omega_2$ involves only bonds with both end points in $\Omega_2$.

\begin{proof}
(i) This is the definition of $\theta_k$ \cite[pp.~566--567]{Ba2} with $k$ replaced by $k+1$: $\theta_{k+1}=1$ on $B^k(\Lambda_2^{(k)})$.

(ii) \cite[p.~604]{Ba1} identifies a vector field $A:T_\eps\to\R^d$ with a bond function by $A_{\langle x,x+\eps e_\mu\rangle}=A_\mu(x)$ for positively oriented bonds, with $g_{-b}=-g_b$. The product $\theta A$ of the point function $\theta$ \cite[p.~567]{Ba2} with a vector field is printed only as such a product. Through the $\R^d$-valued representation it is $(\theta A)_\mu(x)=\theta(x)A_\mu(x)$, i.e.\ $\theta$ is evaluated at the lower end point of the \emph{positively oriented} bond. On an interior bond both end points lie in $\Omega_2$, where $\theta_{k+1}=1$ by (i). So $(\theta_{k+1}A)_b=A_b$ under this convention and under every local convention. With the printed $\tilde B=(1-\theta_{k+1})A^{(k)}+\theta_{k+1}B^{(k+1),\eta}$ \cite[p.~579]{Ba2}, this gives $\tilde B_b=B^{(k+1),\eta}_b$.

(iii) $G_k(\Omega_2,B)$ is the inverse of the covariant Laplacian on $\Omega_2$ with Neumann boundary conditions \cite[p.~609]{Ba1}, plus the mass term, plus $a_kQ_k^*(B)Q_k(B)$ \cite[(2.20)--(2.22)]{Ba1}. It depends on $B$ only in two ways:
\begin{itemize}
\item through the transporters $U(B_b)$ of the bonds of the Neumann form, all of which are interior;
\item through $(Q_k(B)f)(y)=L^{-kd}\sum_{x\in B^k(y)}U(B(\Gamma^{(k)}_{y,x}))f(x)$ \cite[(2.11)]{Ba1}. The contours $\Gamma^{(k)}_{y,x}$ are coordinate-wise segments inside $B^k(y)$ (\cite[(2.1)--(2.2), p.~608]{Ba1}; convention R3 of Lemma~\ref{app:G1:I3r-R123}), and $B^k(y)\subset\Omega_2$ for $y\in\Lambda_2^{(k)}$, so these bonds are interior too.
\end{itemize}
By (ii), $\tilde B$ and $B^{(k+1),\eta}$ agree on all these bonds. This is the replacement Ba{\l}aban prints at U1a and U1b \cite[p.~579; p.~580, Prop.~2.6]{Ba2}.

(iv) First claim. By Lemma~\ref{app:G1:Rnest}(iv), $\Omega_2\subset B^l(\Lambda_2^{(l)})$ for $l\le k$, and $\theta_{l+1}=1$ on $B^l(\Lambda_2^{(l)})$ by (i) at index $l$. So $(1-\theta_{l+1})=0$ at both ends of every interior bond, and the summand vanishes there by (ii).

Second claim. By Lemma~\ref{app:G1:Rloc7}(i) the vertices of \cite[(2.107)]{Ba2} lie in $B^k(\Lambda_7^{(k)})$. By Lemma~\ref{app:G1:Rnest}(v) with $s=5$, these points have depth $\ge5\lfloor r_k\rfloor$ in $\Omega_2$. By Lemma~\ref{app:G1:Rloc7}(ii)--(iii), every bond used by a vertex lies within distance $c_1\le3$ of it. Since $5\lfloor r_k\rfloor>3+1$ (as $r_k\ge R\ge107$), all these bonds are interior, and (ii) applies.
\end{proof}

\paragraph{Localisation in $\Lambda_7$.}
\begin{lemma}[localisation in $\Lambda_7$]\label{app:G1:Rloc7}
At U1a and U1b:
\begin{enumerate}
\item[(i)] Every vertex of every term to which the propagator change is applied lies in $B^k(\Lambda_7^{(k)})$.
\item[(ii)] Every localisation cube $\square(v)$ of \cite[p.~420]{Ba3} has sup-side at most $2$ (scale-$k$ units).
\item[(iii)] The cut-off $h$ of \cite[(1.33)]{Ba3} can be chosen with $h=1$ on $\square(v)$ and its $\eta$-neighbours, $\supp h\subset U_1(\square(v))$ (the closed $1$-neighbourhood), and
\[
\norm{h\Phi}_{1,\alpha_0}\le C_d\,\norm{\Phi}_{1,\alpha_0;\,U_1(\square(v))},
\]
where $\norm{\cdot}_{1,\alpha_0}$ is the norm \cite[(1.32)]{Ba3}, the right side is the same norm with all suprema restricted to points of $U_1(\square(v))$ and their $\eta$-neighbours, and $C_d$ depends on $d$ only. The localisation width is therefore $c_1\le3$, an absolute constant.
\item[(iv)] Under $R\ge3R_0+15$ the margin $5r_k-c_2$ at the uses, $c_2=5+c_1\le9$, is at least $R_0+12$.
\end{enumerate}
\end{lemma}
\noindent The lemma uses \cite[pp.~579--580]{Ba2} and \cite[p.~420]{Ba3}; the localisation width is $c_1\le3$.

\begin{proof}
(i) At U1a, \cite[p.~579]{Ba2} prints: ``The terms containing the vertices localized in $B^k(\Lambda_7^{(k)c})$ can be estimated by $O((L^k\eps)^{\kappa_0})\abs{\Lambda_2^{(k)}\cap\Lambda_7^{(k)c}}$ \dots\ We choose the new localization with the help of a function $g_k$, \dots\ $\supp g_k\subset\Lambda_7^{(k)}$ and $g_k=1$ on $\Lambda_8^{(k)}$. Now the summations in the vertices are restricted to $\Lambda_7^{(k)}$''.

At U1b, \cite[p.~579]{Ba2} prints ``We estimate the terms localized in $B^k(\Lambda_7^{(k)c})$ by $O((L^k\eps)^{\kappa_0})\abs{\Lambda_2^{(k)}\cap\Lambda_7^{(k)c}}$''. \cite[Prop.~2.6, p.~580]{Ba2} acts on ``the interaction terms localized in $B^k(\Lambda_7^{(k)})$''. And \cite[p.~580]{Ba2} prints: ``If $x$ in $(aL^{-2}C^{(k)}_{\Lambda_4^{(k)}}(\dots)Q^*(\dots)\psi)(x)$ is a vertex variable, then $x\in\Lambda_7^{(k)}$''.

This covers all vertices, including the vertices with $n=0$, $n'\ge1$ that come from the $\tilde A$-expansion of $G_k(\Omega,\tilde A+\tilde B)$ \cite[p.~414]{Ba3}. Those vertices are not summands of the interaction \cite[(1.4)]{Ba3}, so an argument based on the local form of (1.4) alone would not cover them.

(ii) \cite[p.~420]{Ba3} prints: ``For the vertices (1.6) and (1.7) we localize simply by representing $\Omega_1$ as a sum of unit cubes of the $\eta$-lattice. For the remaining vertices we localize, taking for each leg of the vector field a smooth partition of unity satisfying the condition that a support of each function is contained in a cube with sides of length 2. \dots\ The expressions (1.12)--(1.15) are localized by fixing a point $y\in T_1^{(k)}\cap\Omega$, or the unit cube $B^k(y)$.'' Hence every $\square(v)$ has side $\le2$.

(iii) The only unprinted width is that of $h$ (``$h=1$ on $\square(v)$ and $h=0$ outside some neighborhood of $\square(v)$'', \cite[p.~420]{Ba3}). It is a choice of whoever applies \cite[Prop.~1]{Ba3}, for two reasons. First, the vertex depends on $\Phi_{\rm ext}$ only on $\square(v)$ and its $\eta$-neighbours, so the expression $E$ of \cite[(1.33)]{Ba3} satisfies $E(\Phi_{\rm ext})=E(h\Phi_{\rm ext})$ for any $h$ equal to $1$ there. Second, the constant of \cite[(1.33)]{Ba3} is independent of $\eps$, $k$ and the domains \cite[p.~421]{Ba3}, and $h$ enters only through $\norm{h\Phi_{\rm ext}}_{1,\alpha_0}$, $\norm{h'A_{\rm ext}}_{1,\alpha_0}$.

We now construct $h$. Let $s(t)=10t^3-15t^4+6t^5$ on $[0,1]$; then $s(0)=0$, $s(1)=1$, $s'=s''=0$ at both ends, $\max s'=15/8$ and $\max\abs{s''}=10/\sqrt3$. Put $\chi(t)=1$ for $t\le\frac12$, $\chi(t)=1-s(2t-1)$ for $\frac12\le t\le1$, and $\chi(t)=0$ for $t\ge1$. Then $\chi\in C^2$, $\abs{\chi'}\le15/4$ and $\abs{\chi''}\le40/\sqrt3<32$. Set $h(x)=\prod_\mu\chi(t_\mu(x))$, where $t_\mu(x)$ is the distance of $x_\mu$ from the $\mu$-th side interval of $\square(v)$. Since $\eta=L^{-k}\le\frac12$, $h=1$ on $\square(v)$ and its $\eta$-neighbours, and $\supp h\subset U_1(\square(v))$, a cube of side $\le4$. The lattice first and second differences of $h$ are bounded by constants depending on $d$ only, uniformly in $\eta$.

Write $D_\mu$ for the covariant difference of \cite[(1.32)]{Ba3}, with unimodular transporter $U_\mu(x)$ on the bond $\langle x,x+\eta e_\mu\rangle$. Then exactly
\[
D_\mu(hf)(x)=h(x+\eta e_\mu)\,(D_\mu f)(x)+(\partial^\eta_\mu h)(x)\,f(x).
\]
This gives the sup and first-difference parts of the norm with a constant depending on $d$ only.

For the H\"older quotient, put $F:=D_\mu(hf)$ and $U_\Gamma:=U(\tilde B(\Gamma_{x,x'}))$. Then
\begin{align*}
U_\Gamma F(x')-F(x)&=h(x'+\eta e_\mu)\big(U_\Gamma D_\mu f(x')-D_\mu f(x)\big)+\big(h(x'+\eta e_\mu)-h(x+\eta e_\mu)\big)D_\mu f(x)\\
&\quad+(\partial^\eta_\mu h)(x')\big(U_\Gamma f(x')-f(x)\big)+\big((\partial^\eta_\mu h)(x')-(\partial^\eta_\mu h)(x)\big)f(x).
\end{align*}
Here $\abs{U_\Gamma f(x')-f(x)}\le d\abs{x-x'}\sup\abs{Df}$, by telescoping along the shortest contour. The differences of $h$ and $\partial^\eta h$ are $\le C_d\abs{x-x'}$. For $\abs{x-x'}\le1$ this gives the bound $C_d\abs{x-x'}^{1-\alpha_0}\norm{f}_{1,\alpha_0}\le C_d\norm f_{1,\alpha_0}$ after division by $\abs{x-x'}^{\alpha_0}$; for $\abs{x-x'}>1$, bound each term by twice a supremum.

If $x$ lies outside $S:=\{x:\,x\text{ or }x+\eta e_\mu\in\supp h\}$, then $F(x)=0$ and the restricted norm sees only $F(x')$, $x'\in S$. Because $h$ and $\partial^\eta h$ vanish to first order at the edge of $\supp h$ ($\chi'=0$ at $t=1$), $\abs{F(x')}\le C_d\abs{x-x'}\sup_{U_1(\square(v))}(\abs f+\abs{Df})$, which is again bounded after division by $\abs{x-x'}^{\alpha_0}$. A merely linear ramp would not do: its difference jumps across one $\eta$-bond, and the quotient is then of order $\eta^{-\alpha_0}$.

So $c_1\le$ (side of $\square(v)$) $+$ (collar) $\le2+1=3$.

(iv) $5r_k-9\ge5R-9\ge R_0+12$, because $5R\ge15R_0+75\ge R_0+21$.
\end{proof}

\paragraph{Assembly.}
\paragraph{The applications of Proposition~1 of \cite{Ba3}.}
Throughout this paragraph $L=2$, $\nbar=13$ and $\Nbar=169$, and the regulator data are chosen in the order of Appendix~\ref{app:B}. In particular $R\ge R^{**}$, so that
\begin{gather*}
R\ge107,\qquad R\ge3R_0+15,\qquad R\ge R_0+13+c_1',\\
R\ge2\ell_{\Nbar}+10M_b+10,\qquad R\ge4\ell_{\Nbar}+32M_b+25,\qquad \delta_0R\ge4 .
\end{gather*}
Also $\eps_0<\mathrm{e}^{-1}$, the $\eps_0$-bounds of Lemma~\ref{app:G1:RW-KSP} hold, $\ell_0\ge\eps_0[8r(\eps_0)+4M_b+2]$ (condition (T.nw)) and $\ell_T/(M_bt_k)\in2\N$. Since $t\le\eps_0<\mathrm{e}^{-1}$, we have $1+\log t^{-1}>2$, and hence
\begin{equation}\label{eq:I3a-r4R}
r(t)=R(1+\log t^{-1})^2>4R .
\end{equation}
The notation is that of the setting paragraph of Subsection~\ref{app:G1:part-E5} and of Subsection~\ref{app:G1:part-TI3}. The constant $c_1'$ is the localisation width at scale $j$ of Lemma~\ref{app:G1:I3r-HOm2}.

\begin{definition}[the applications of Proposition~1 at step $k$]\label{app:G1:I3a-uses}
Fix $k\ge1$. In \cite{Ba2,Ba3}, \cite[Prop.~1]{Ba3} is applied, or the domain of a scalar propagator is changed in the way its proof requires, in the situations listed below. For each one we name the \emph{outer domain} $\Omega_{\rm out}$ and the \emph{new domain}.
\begin{itemize}
\item[(U0)] \cite[pp.~411--414, 432--433]{Ba3} itself, and its application at \cite[p.~582, (2.116)]{Ba2}: $\Omega_{\rm out}=\Omega=B^k(\Lambda_2^{(k-1)\prime})$; the new domain is $\Omega_2$.
\item[(U1a)] \cite[p.~579]{Ba2}: $\Omega_{\rm out}=B^k(\Lambda_0^{(k)})$; the new domain is $\Omega_2=B^k(\Lambda_2^{(k)})$; the background is $\tilde B=(1-\theta_{k+1})A^{(k)}+\theta_{k+1}B^{(k+1),\eta}$.
\item[(U1b)] \cite[p.~580, Prop.~2.6]{Ba2}: $\Omega_{\rm out}=B^{k-1}(\Lambda_2^{(k-1)})$, which equals $\Omega$ by Lemma~\ref{app:G1:Rnest}(iii); the new domain is $\Omega_2$.
\item[(U2)] \cite[p.~578]{Ba2}, at a scale $0\le j\le k-1$, on the lattice $T_{L^{-j}}$ of spacing $\eta_j=L^{-j}$: $\Omega_{\rm out}=B^j(\Lambda_2^{(j)})$; the new domain is $B^k(\Lambda_0^{(k)})$. The propagator $G_j(B^j(\Lambda_2^{(j)}),\tilde B)$ is replaced by $G_j(B^k(\Lambda_0^{(k)}),\tilde B)$.
\item[(U3)] Inside the proof of \cite[Prop.~1]{Ba3}, at any of the above: the substitution $G_k(\square,0)=G_k(0)+\delta G_k(\square,\eta\Z^d,0)$ of \cite[p.~433]{Ba3}.
\end{itemize}
An \emph{evaluation point} of a kernel at an application is any of the following: a line end, a bond partner, a block point of $Q_k$, a H\"older partner, or a point of the support of a localisation function $h,h'$ of \cite[(1.33)]{Ba3}.
\end{definition}

\begin{lemma}[one background at the change to $\Omega_2$]\label{app:G1:I3a-onebg}
At (U1a) and (U1b), let $\tilde B$ be the background and $\Omega_{\rm out}$ the outer domain of the application. Then $\Omega_2\subset\Omega_{\rm out}$, both sets are unions of big blocks, and, as kernels on $\Omega_2\times\Omega_2$,
\[
G_k(\Omega_{\rm out},\tilde B)-G_k(\Omega_2,B^{(k+1),\eta})=G_k(\Omega_{\rm out},\tilde B)-G_k(\Omega_2,\tilde B)=:\delta G_k(\Omega_{\rm out},\Omega_2,\tilde B).
\]
At (U2) the printed replacement is $\delta G_j(B^j(\Lambda_2^{(j)}),B^k(\Lambda_0^{(k)}),\tilde B)$. It is again taken at the single background $\tilde B$, and $B^k(\Lambda_0^{(k)})\subset B^j(\Lambda_2^{(j)})$. In each case the background is regular in the sense of hypothesis (R) of Subsection~\ref{app:G1:part-TI3}.
\end{lemma}
\begin{proof}
\emph{Nesting.} At (U1a), $\Omega_2=B^k(\Lambda_2^{(k)})\subset B^k(\Lambda_0^{(k)})$ by Lemma~\ref{app:G1:Rnest}(i),(ii). At (U1b), $\Omega_2\subset\Omega$ by Lemma~\ref{app:G1:Rnest}(iii). At (U2), Lemma~\ref{app:G1:Rnest}(i)--(iii) give $B^k(\Lambda_0^{(k)})\subset B^k(\Lambda_7^{(k-1)\prime})=B^{k-1}(\Lambda_7^{(k-1)})\subset B^{k-1}(\Lambda_2^{(k-1)})$. The argument of Lemma~\ref{app:G1:Rnest}(iv) then gives $B^{k-1}(\Lambda_2^{(k-1)})\subset B^j(\Lambda_2^{(j)})$ for $j\le k-1$. All these sets are unions of big blocks \cite[p.~566]{Ba2}, and $B^l$ of a union of big blocks is again a union of big blocks.

\emph{Identity.} By Lemma~\ref{app:G1:Rthetab}(ii)--(iv), on every bond with both end points in $\Omega_2$ the background $\tilde B$ of (U1a) equals $B^{(k+1),\eta}$. The same holds for the background of (U1b): all its summands other than $\theta_{k+1}B^{(k+1),\eta}$ vanish there. By Lemma~\ref{app:G1:Rthetab}(iii), $G_k(\Omega_2,\cdot)$ depends on its background only through such bonds. Hence $G_k(\Omega_2,\tilde B)=G_k(\Omega_2,B^{(k+1),\eta})$, which is the displayed identity. At (U2) the propagator keeps its background \cite[p.~578]{Ba2}.

\emph{Regularity.} At (U1a) the regularity of the background is \cite[(2.98), p.~579]{Ba2}. The backgrounds of (U1b) and (U2) are configurations of the form \cite[(2.93), p.~576]{Ba2}, and their regularity is \cite[(2.94)--(2.97), p.~576]{Ba2}. The smallness required in (R) is given by the $\eps_0$-bounds of Lemma~\ref{app:G1:RW-KSP}. The H\"older part for the $\tilde A$-legs is Proposition~\ref{app:G1:AAext}.
\end{proof}

\begin{lemma}[two-domain bound at lattice spacing $1$]\label{app:G1:I3a-eta1}
Let $\eta=1$. Let $X_{\rm in}\subset X_{\rm out}$ be unions of blocks of $T_1$ (or let $X_{\rm out}=T_1$), and let $A$ be a regular vector field in the sense of \cite[p.~573]{Ba4}. Put $\delta:=G(X_{\rm out},A)-G(X_{\rm in},A)$, where $G$ denotes the Green's functions of \cite[(1.6)]{Ba4} with Neumann conditions, and let $c_0,\delta_0,R_0$ be the constants of the theorem on \cite[p.~573]{Ba4}. Let $x,x'$ be points such that $x$, $x'$ and their lattice neighbours have distance $\ge R_0$ from $X_{\rm in}^c$. Write $D$ for the covariant difference and $D^*$ for its adjoint. Then
\[
\max\big\{\abs{\delta(x,x')},\abs{D_x\delta(x,x')},\abs{(\delta D^*)(x,x')},\abs{(D_x\delta D^*)(x,x')}\big\}\le2c_0\mathrm{e}^{\delta_0}\mathrm{e}^{-\delta_0\abs{x-x'}} .
\]
Every $\alpha$-H\"older quotient of these kernels over $0<\abs{x-\tilde x}\le1$ is at most twice the supremum of the kernel over the points involved if it is taken in one variable, and at most four times that supremum if it is taken jointly. Consequently, if $0\le h,h'\le1$ and every point of $\supp h\cup\supp h'$ has distance $\ge R_0+2$ from $X_{\rm in}^c$, then
\[
\norm{h\,\delta\,h'}_{1,\alpha}\le C(c_0,\delta_0)\,\mathrm{e}^{-\delta_0\dist(\supp h,\supp h')},
\]
with the norm of \cite[(1.32)]{Ba3} in both variables.
\end{lemma}
\begin{proof}
At $\eta=1$ kernels are taken with respect to counting measure. Fix $y$, and let $f$ be the point mass at $y$ with values in a unit vector. Then $\norm f_\infty=1$, and $(\delta f)(x)$ is the column $\delta(x,y)$. We apply \cite[(1.10)--(1.12), p.~573]{Ba4} with $\Omega=X_{\rm in}\subset\Omega_0=X_{\rm out}$; the additional factor (1.12) is at most $1$. This gives, for every $y$,
\[
\abs{\delta(x,y)}+\abs{D_x\delta(x,y)}\le c_0\mathrm{e}^{-\delta_0\abs{x-y}}\qquad\text{whenever }\dist(x,X_{\rm in}^c)\ge R_0 .
\]
If $y\notin X_{\rm in}$, then $\delta(\cdot,y)=G(X_{\rm out},A)(\cdot,y)$, and (1.10) for $X_{\rm out}$ applies, since $\dist(x,X_{\rm out}^c)\ge\dist(x,X_{\rm in}^c)$.

At $\eta=1$, $(D^*_\mu g)(x')$ is a combination of $g(x')$ and $g(x'-e_\mu)$ with coefficients of modulus $1$. Hence, for $K=\delta$ or $K=D_x\delta$,
\[
\abs{(KD^*)(x,x')}\le\abs{K(x,x')}+\abs{K(x,x'-e_\mu)}\le2c_0\mathrm{e}^{\delta_0}\mathrm{e}^{-\delta_0\abs{x-x'}} .
\]
This gives the first display.

Distinct lattice points with $\abs{x-\tilde x}\le1$ satisfy $\abs{x-\tilde x}=1$. So each H\"older quotient is at most $\abs{UF(\tilde x)}+\abs{F(x)}$: a sum of two kernel values (four, for the joint quotient) at points that again satisfy the depth hypothesis.

Finally, $D_\mu(hF)(x)=h(x+e_\mu)D_\mu F(x)+(\partial_\mu h)(x)F(x)$ with $\abs{\partial_\mu h}\le1$ at $\eta=1$, and similarly for $h'$. Every point at which a kernel is then evaluated lies within distance $2$ of $\supp h\cup\supp h'$, and so has depth $\ge R_0$. For such points, $\mathrm{e}^{-\delta_0\abs{x-x'}}\le\mathrm{e}^{2\delta_0}\mathrm{e}^{-\delta_0\dist(\supp h,\supp h')}$. Summing the finitely many terms gives the last bound.
\end{proof}

\begin{lemma}[localisation after the $B'$-expansion]\label{app:G1:I3a-Bloc}
At \cite[p.~433]{Ba3}, let the localisation cubes of a connected component lie in a cube $Q_0$ of side $\ell_{\Nbar}$. Let $\square_1$ be a union of big blocks of size $r=r(L^k\eps)$ that contains the cube of side $r$ concentric with $Q_0$. Suppose that the expansion in $B'$ of each scalar line \cite[(1.3.14), (3.44), (3.45)]{Ba1} produces chains with at most $n_{B'}$ new vertices. Suppose also that, as printed, every line and every operator \cite[(1.16)]{Ba3} whose end localisations satisfy $\dist(\square(v),\square(v'))\ge1$ is put into the external fields. If
\begin{equation}\label{eq:I3a-Bloc}
r\ge\ell_{\Nbar}+6n_{B'},
\end{equation}
then every vertex of every class that contains an old vertex is localised in $\square_1$. Under the standing hypotheses, \eqref{eq:I3a-Bloc} holds for $n_{B'}=\nbar=13$.
\end{lemma}
\begin{proof}
New vertices arise only inside the chains. Each chain replaces one line between two vertices of the expression before the expansion, and a new vertex is joined by lines only to its two neighbours in its chain. Let $w$ be a new vertex of a class that contains an old vertex. The class is connected and $w$ is joined to it only through its chain. Hence the class contains an uncut segment of the chain, of at most $n_{B'}$ lines, from $w$ to one of the two old end vertices.

Every line of the class joins localisation cubes at distance $<1$, and every cube has side $\le2$ (Lemma~\ref{app:G1:Rloc7}(ii)). In the sup metric, $\dist(Q,Q')<1$ and $\operatorname{side}(Q')\le2$ give $Q'\subset U_3(Q)$. By induction along the segment, $\square(w)\subset U_{3n_{B'}}(Q_0)$. This is the cube of side $\ell_{\Nbar}+6n_{B'}$ concentric with $Q_0$; it does not wrap the torus, since $\square$ does not (Lemma~\ref{app:G1:I3r-Wsq}). Under \eqref{eq:I3a-Bloc} this cube lies in the concentric cube of side $r$, and hence in $\square_1$.

For the last claim: by \eqref{eq:I3a-r4R}, $R\ge2\ell_{\Nbar}+10M_b+10$ and $R\ge107$, we get $r>2R+2R\ge\ell_{\Nbar}+214>\ell_{\Nbar}+78$.
\end{proof}

\begin{lemma}[depth supplied at the applications]\label{app:G1:I3a-margins}
Let $k\ge1$.
\begin{enumerate}
\item[(a)] At (U1a) and (U1b), every evaluation point of $G_k(\Omega_2,\cdot)$ and of $\delta G_k(\Omega_{\rm out},\Omega_2,\cdot)$ has distance $\ge5r_k-9\ge R_0+13$ from $\Omega_2^c$.
\item[(b)] At (U2), at scale $j$, every evaluation point has distance $>L^{k-j}r_k-c_1'\ge R_0+13$ from the complement of $B^k(\Lambda_0^{(k)})$, in scale-$j$ units.
\item[(c)] For a class with an old vertex, at any application, every point of the support of a localisation function of a vertex of the class has distance $\ge r-1\ge R_0+13$ from $\square^c$, and hence from $\Omega_{\rm out}^c$.
\item[(d)] Consequently, the support condition $R_0+11+3\eta\le R_0+13$ (for $\eta\le\frac12$) of Corollary~\ref{app:G1:twodomain-loc} holds in (a), in (c), and in (b) for $j\ge1$. The condition $R_0+2$ of Lemma~\ref{app:G1:I3a-eta1} holds in (b) for $j=0$.
\end{enumerate}
\end{lemma}
\begin{proof}
(a) Lemma~\ref{app:G1:I3r-HOm2} gives the lower bound $5r_k-5-c_1$ with $c_1\le4$. Moreover $5r_k-9\ge5R-9\ge15R_0+66$.

(b) Lemma~\ref{app:G1:I3r-HOm2} gives $>L^{k-j}r_k-c_1'\ge2R-c_1'$, and $R\ge R_0+13+c_1'$ gives $2R-c_1'\ge R+R_0+13$.

(c) By Lemma~\ref{app:G1:I3a-Bloc} the localisation cubes lie in $\square_1$. The localisation functions are supported in $U_1(\square(v))$ (Lemma~\ref{app:G1:Rloc7}(iii)). Every point of $\square_1$ has distance $\ge r$ from $\square^c$ (Proposition~\ref{prop:prop1}(c)). And $r-1>4R-1\ge12R_0+59$ by \eqref{eq:I3a-r4R}. Moreover $\square\subset\Omega_{\rm out}$ (Lemma~\ref{app:G1:I3r-Wsq}).

(d) For $\eta\le\frac12$ we have $R_0+\frac{45}4+3\eta\le R_0+\frac{51}4$. At $j=0$ the bound in (b) exceeds $R_0+2$.
\end{proof}

\begin{lemma}[lowering the tree rate]\label{app:G1:I3a-rate}
If a family satisfies \cite[(1.33)]{Ba3} with tree rate $\delta>0$, then it satisfies it with the same constant and with every tree rate $\delta'\in(0,\delta]$.
\end{lemma}
\begin{proof}
In \cite[(1.33), p.~420]{Ba3} the rate enters only through the factor $\exp[-\delta\,d(\{\square(v)\})]$, where $d\ge0$ is the length of a shortest tree. This factor does not decrease when $\delta$ decreases.
\end{proof}

\begin{proposition}[the conditions of Proposition~1 as used in \cite{Ba2}]\label{app:G1:I3}
Assume the hypotheses of Theorem~\ref{thm:T6} and the standing hypotheses of this paragraph. Let $k\ge1$, and consider any application (U0), (U1a), (U1b) or (U2) of Definition~\ref{app:G1:I3a-uses}. Then the following hold.
\begin{enumerate}
\item[(i)] \emph{(Hypotheses.)} The background is regular in the sense of (R). The support hypothesis holds in the form (R4$'''$) of \eqref{eq:R4ppp}, and \cite[Prop.~1]{Ba3} holds with (R4) replaced by (R4$'''$).
\item[(ii)] \emph{(Remainder kernels.)} The operator \cite[(1.16)]{Ba3} satisfies the two-variable bound of Lemma~\ref{app:G1:I3r-U1m}(b),(c) at $\alpha=\alpha_0=\frac12$ on $T_\eta$. It satisfies the bound of Lemma~\ref{app:G1:I3r-U1mOmega} at every evaluation point on the Neumann domain $\Omega$.
\item[(iii)] \emph{(Change of domain.)} At (U1a), (U1b) and (U2), the printed replacement is the subtraction of one kernel $\delta G$, at one regular background, for nested domains (Lemma~\ref{app:G1:I3a-onebg}). Moreover $\norm{h\,\delta G\,h'}_{1,\alpha_0}\le C\,\mathrm{e}^{-\beta\dist(\supp h,\supp h')}$ for all localisation functions $h,h'$ of the application. \cite[Chap.~3]{Ba2} contains no further change of domain of a scalar propagator (Lemma~\ref{app:G1:I3r-HOm2ch3}).
\item[(iv)] \emph{(The cube step of \cite[p.~433]{Ba3}.)} Let a class have an old vertex. Then:
\begin{itemize}
\item $\square\subset\Omega_{\rm out}$;
\item after the $B'$-expansion, every vertex of the class is localised in $\square_1$;
\item the kernels $\delta G_k(\Omega_{\rm out},\square,\tilde B)$, $\delta G_k(\Omega,\square,\tilde B)-\delta G_k(\Omega,\Omega_2,\tilde B)$ and $\delta G_k(\square,\eta\Z^d,0)$ are external fields whose $(1,\alpha_0)$-norm at the localisations of the class is at most $C\,\mathrm{e}^{-\beta\dist(\supp h,\supp h')}$;
\item the $B'$-remainder satisfies the bound of Lemma~\ref{app:G1:I3r-cube};
\item the inclusion $\square\subset\Omega_2$ is not used.
\end{itemize}
\item[(v)] \emph{(Mass entries.)} The conclusions of Lemma~\ref{app:G1:I3r-ME} hold. At the two points $(u,M)\in\{(1,m_2),(-\eta,m_1)\}$ of Theorem~\ref{thm:T6}, the $A\cdot A$ vertices of \cite[(3.30), (3.32)]{Ba3} satisfy (N1)--(N4) of Lemma~\ref{app:G1:I3r-AA}.
\item[(vi)] \emph{(The tails of \cite[p.~579]{Ba2}.)} The decomposition and the bounds of Lemma~\ref{app:G1:I3r-T579} hold.
\end{enumerate}
Here $\beta:=\beta_P:=\min\{\delta(\Nbar),\beta_U,\beta_D,\delta_0,\beta_\square\}$. In this definition, $\delta(\Nbar)$ is the rate of Proposition~\ref{app:G1:VSR-CAP}, $\beta_U$ and $\beta_D$ are the rates of Lemmas~\ref{app:G1:I3r-U1m} and~\ref{app:G1:I3r-DGM}, $\delta_0$ is the rate of \cite[p.~573]{Ba4}, and $\beta_\square$ is the rate of Lemma~\ref{app:G1:I3r-cube}.

Consequently, \cite[(1.33)]{Ba3} holds at each application, with tree rate $\beta_P$, for every localised family of order at most $\Nbar$. The rate $\beta_P$, the constants of (i)--(iv) and the constant of this bound depend only on $d$, $a$, $L$, $M_b$, $\Nbar$, $\alpha_0$, the constants of \cite{Ba1,Ba4} and the constants of Lemma~\ref{app:G1:WKinputs}. They do not depend on $\eps$, $k$, $\ell_T$, on the domains $\Omega,\Omega_1,\Omega_2$, or on the background within (R). They depend on the scalar mass only through the condition $0\le m^2t^2\le c$. The constants of (v) and (vi) are those of Lemmas~\ref{app:G1:I3r-AA} and~\ref{app:G1:I3r-T579}. They are independent of $\eps$, $k$ and $\ell_T$, and they are finite at the two points $(u,M)$ of Theorem~\ref{thm:T6}.
\end{proposition}

\begin{table}[ht]
\centering\footnotesize
\begin{tabular}{@{}p{0.05\textwidth}p{0.56\textwidth}p{0.33\textwidth}@{}}
\hline
item & statement & proved in \\
\hline
1 & one-variable bound on the kernel \cite[(1.16)]{Ba3} at points of depth $\ge R_0$ & Lemma~\ref{app:G1:I3r-U1} \\
2 & two-variable bound on \cite[(1.16)]{Ba3} on $T_\eta$ & Lemma~\ref{app:G1:I3r-U1m} \\
3 & the same on the Neumann domain $\Omega$, at every evaluation point & Lemma~\ref{app:G1:I3r-U1mOmega} \\
4 & no field-independent mass term other than the $\delta m^2$ insertions & Lemma~\ref{app:G1:I3r-ME} \\
5, 5$'$ & the abbreviated steps of \cite{Ba2,Ba3} that need more than the printed text & Lemma~\ref{app:G1:I3r-Pw} \\
6 & separation bound on a cube that is a union of big blocks & Lemma~\ref{app:G1:I3r-SEP} \\
7 & the $B'$-remainder of \cite[p.~433]{Ba3} & Lemma~\ref{app:G1:I3r-cube} \\
8 & two-domain bounds for nested unions of big blocks, for $0<\eta\le\frac12$ and for $\eta=1$ & Proposition~\ref{app:G1:twodomain}, Lemmas~\ref{app:G1:I3r-DGM}, \ref{app:G1:I3a-eta1} \\
9 & $\square\subset\Omega_{\rm out}$ at every application & Lemma~\ref{app:G1:I3r-Wsq} \\
10 & the cube step without $\square\subset\Omega_2$ & Lemma~\ref{app:G1:I3r-W2R} \\
11, 12 & depth of the evaluation points of $\Omega_2$-kernels in \cite[Chaps.~2, 3]{Ba2} & Lemmas~\ref{app:G1:I3r-HOm2}, \ref{app:G1:I3r-HOm2ch3}, \ref{app:G1:I3a-margins} \\
13 & nesting of the $\Lambda_i^{(j)}$, and $\Omega\supset\Omega_1\supset\Omega_2$ & Lemma~\ref{app:G1:Rnest} \\
14 & the support hypothesis in the form (R4$'''$) & Corollary~\ref{cor:prop1R4} \\
15 & gating by $g_k$, metric, contours inside blocks & Lemmas~\ref{app:G1:I3r-R123}, \ref{app:G1:Rthetab}, \ref{app:G1:Rloc7} \\
15$'$ & the list of applications is complete & Lemma~\ref{app:G1:I3r-Rexh} \\
16 & the $A\cdot A$ local terms & Lemma~\ref{app:G1:I3r-AA}, Propositions~\ref{app:G1:I3c-AA}, \ref{app:G1:AAext} \\
17 & the tails of \cite[p.~579]{Ba2} & Lemma~\ref{app:G1:I3r-T579}, Propositions~\ref{app:G1:T579}, \ref{app:G1:I3c-R}, \ref{app:G1:I3c-B} \\
-- & one background; localisation after the $B'$-expansion; depth at the applications; tree rate & Lemmas~\ref{app:G1:I3a-onebg}, \ref{app:G1:I3a-Bloc}, \ref{app:G1:I3a-margins}, \ref{app:G1:I3a-rate} \\
\hline
\end{tabular}
\caption{The statements that make up Proposition~\ref{app:G1:I3}.}\label{app:G1:I3tab}
\end{table}

\begin{proof}
\emph{(i).} At (U1a), (U1b) and (U2), regularity is the last part of Lemma~\ref{app:G1:I3a-onebg}. At (U0) it is the hypothesis of \cite[p.~412]{Ba3}, which is verified in the same way from \cite[(2.94)--(2.98)]{Ba2}. By Corollary~\ref{cor:prop1R4} and \eqref{eq:eppp}, the form (R4$'''$) holds at every application, and \cite[Prop.~1]{Ba3} holds under it.

\emph{(ii).} On $T_\eta$ this is Lemma~\ref{app:G1:I3r-U1m}(b),(c) with $m=n+n'\ge6$; this covers $\alpha_0=\frac12$. On $\Omega$ it is Lemma~\ref{app:G1:I3r-U1mOmega}. The depth that lemma needs at the vertex points, $R_0+2\sqrt3+1$, is supplied under (R4$'''$) by part~(3) of its proof.

\emph{(iii).} By Lemma~\ref{app:G1:I3a-onebg}, the replacement is one kernel $\delta G$, for a nested pair of unions of big blocks, at one regular background. For $\eta\le\frac12$ the corollary of Proposition~\ref{app:G1:twodomain} applies, with rate $\beta_D$, by Lemma~\ref{app:G1:I3a-margins}(a),(b),(d). At (U2) with $j=0$, Lemma~\ref{app:G1:I3a-eta1} applies, with rate $\delta_0$, by Lemma~\ref{app:G1:I3a-margins}(b),(d). Lemma~\ref{app:G1:I3r-HOm2} shows that (U1a), (U1b) and (U2) are all the changes of domain of a scalar propagator in \cite[Chap.~2]{Ba2}. Lemma~\ref{app:G1:I3r-HOm2ch3} shows that \cite[Chap.~3]{Ba2} adds none.

\emph{(iv).} The inclusion $\square\subset\Omega_{\rm out}$ is Lemma~\ref{app:G1:I3r-Wsq}. At (U0), its hypothesis on the support of $\tilde B^{(k)}$ holds for the following reason. $\tilde B^{(k)}$ is defined as the sum of the second, third and fourth terms of \cite[(2.89), p.~575]{Ba2}, and each of these terms has $\theta_{k+1}$ as its left factor; so $\supp\tilde B^{(k)}\subset\supp\theta_{k+1}$. The localisation in $\square_1$ is Lemma~\ref{app:G1:I3a-Bloc}. The bound on $\delta G_k(\Omega_{\rm out},\square,\tilde B)$ is the corollary of Proposition~\ref{app:G1:twodomain} for the pair $(\Omega_{\rm out},\square)$; it applies by Lemma~\ref{app:G1:I3a-margins}(c),(d). The cube step relative to $\Omega_2$ is Lemma~\ref{app:G1:I3r-W2R}. Its inputs are the two-domain bounds for $(\Omega,\square)$ and $(\Omega,\Omega_2)$ just obtained, Lemma~\ref{app:G1:I3r-Wsq}, and Lemmas~\ref{app:G1:I3r-HOm2} and~\ref{app:G1:I3r-HOm2ch3}; it never uses $\square\subset\Omega_2$. The bound on $\delta G_k(\square,\eta\Z^d,0)$ at points of $\square_1$ is step S19 in the proof of Proposition~\ref{prop:prop1}. The $B'$-remainder is Lemma~\ref{app:G1:I3r-cube}, with Lemma~\ref{app:G1:I3r-SEP}. It is applied with $\alpha=\alpha_0<1$, with $\beta=\beta_P\le\beta_\square$, at remainder orders $m\ge6$, and with $D$ and the transports taken with respect to the same background. Under these hypotheses, Lemma~\ref{app:G1:cube-margin} supplies the hypotheses on the localisation functions $h,h'$ of the class.

\emph{(v), (vi).} These are Lemmas~\ref{app:G1:I3r-ME}, \ref{app:G1:I3r-AA} and~\ref{app:G1:I3r-T579}.

\emph{The bound (1.33).} The proof of \cite[Prop.~1]{Ba3} on \cite[pp.~422--445]{Ba3} consists of printed steps and of the sketched steps listed in Table~\ref{tab:prop1}. By Proposition~\ref{prop:prop1}, these sketched steps are supplied as follows:
\begin{itemize}
\item S1 by (ii), and S2 by (iii);
\item P12 and S13--S19 by (iv);
\item S3, S4 and S9 by Proposition~\ref{prop:prop1}(b), whose depth requirements are met by Lemma~\ref{app:G1:I3a-margins};
\item the remaining steps need no depth (Lemma~\ref{lem:Omegafree}).
\end{itemize}
With the order cap, Proposition~\ref{app:G1:VSR-CAP} gives (1.33) with rate $\delta(\Nbar)$ for families of order at most $\Nbar$. The kernel bounds used along the way carry the rates $\beta_U$, $\beta_D$ and $\delta_0$. By Lemma~\ref{app:G1:I3a-rate} all of them hold with rate $\beta_P$, so (1.33) holds with tree rate $\beta_P$. The quantity $\beta_P$ depends only on constants that are fixed before $R$ is chosen (Appendix~\ref{app:B}), so the order of choices has no cycle (Lemma~\ref{app:G1:acyc}).

\emph{Order of use.} The derivation of (1.33) uses only (i)--(iv) and Lemma~\ref{app:G1:I3r-ME}. Lemmas~\ref{app:G1:I3r-AA} and~\ref{app:G1:I3r-T579} use (1.33) and the cube step (iv), and they are not used in the proofs of (i)--(iv). Hence there is no circularity.

\emph{Uniformity.} Each cited statement carries the stated dependence of its constants. The kernels in (ii)--(iv) are scalar Green's functions and do not involve the photon covariance. The uniformity of (v) and (vi) at the two points $(u,M)$ of Theorem~\ref{thm:T6} is part of Lemmas~\ref{app:G1:I3r-AA} and~\ref{app:G1:I3r-T579}.
\end{proof}

\begin{remark}[completeness of the list]\label{app:G1:I3a-complete}
The list of applications in Definition~\ref{app:G1:I3a-uses} is complete by Lemma~\ref{app:G1:I3r-Rexh} and Lemmas~\ref{app:G1:I3r-HOm2}, \ref{app:G1:I3r-HOm2ch3}. The sketched steps of the proof of \cite[Prop.~1]{Ba3} are those of Proposition~\ref{prop:prop1}(a). These completeness statements concern the printed text of \cite{Ba1,Ba2,Ba3}. They were obtained by reading the pages concerned, and they are not consequences of a mathematical argument. Lemma~\ref{app:G1:I3r-Pw} lists the abbreviated steps of \cite[pp.~555--594]{Ba2} and \cite[pp.~411--445]{Ba3} that need more than the printed text. That the remaining abbreviated steps on these pages hold under the hypotheses of Theorem~\ref{thm:T6} and the conditions of Appendix~\ref{app:B} is Statement~\ref{app:K:L10-1}.
\end{remark}

\begin{remark}[scope]\label{app:G1:I3-missing}
(i) Every clause of Proposition~\ref{app:G1:I3} is proved in this paper, with one exception. In the proof of Lemma~\ref{app:G1:I3r-AA}, the uniformity of the constant of \cite[Prop.~1]{Ba3} under a constant shift of the background is Proposition~\ref{app:K:L10-2}.
(ii) Lemma~\ref{app:G1:I3a-Bloc} is stated for a general order $n_{B'}$ of the $B'$-expansion; any $n_{B'}\le(r-\ell_{\Nbar})/6$ suffices. The value $n_{B'}=\nbar$ is the order to which \cite[p.~433]{Ba3} refers through \cite[(3.45)]{Ba1}.
(iii) Corollary~\ref{cor:prop1R4} is proved at $L=2$, and hence so is Proposition~\ref{app:G1:I3}.
\end{remark}

\subsection{\texorpdfstring{The statements of Table~\ref{app:G1:I3tab}}{The statements of the table of conditions}}\label{app:G1:part-TI3}
We prove the statements of Table~\ref{app:G1:I3tab} other than 13 and 14 (Lemma~\ref{app:G1:Rnest} and Corollary~\ref{cor:prop1R4}). Notation is that of the setting paragraph of Subsection~\ref{app:G1:part-E5}, together with the following:
\begin{itemize}
\item $t=t_k=L^k\eps$, $r=r(t)=R(1+\log t^{-1})^2$, $e_k=e(t)$, $\eta=L^{-k}$, and $\abs{\cdot}$ is the $\ell^\infty$ torus distance.
\item $G(\mathcal B)=G_k(\mathcal X,\mathcal B)=(-\Delta^{\eta,N}_{\mathcal B,\mathcal X}+m^2t^2+a_kP_k(\mathcal B))^{-1}$ on a block union $\mathcal X$, with Neumann conditions on $\partial\mathcal X$.
\item $D=D^\eta_{\tilde B}$, $D^*$ its $\ell^2(\eta^3)$-adjoint, and $\mathfrak Df:=(f,D_1f,D_2f,D_3f)$.
\item $\sigma:=e_k(\norm{\tilde A}_\infty+\norm{\partial^\eta\tilde A}_\infty)$ and $\theta(\mathcal B):=e_k\norm{\partial^\eta\mathcal B}_\infty$.
\end{itemize}
Kernels are taken with respect to $\eta^3$.

\paragraph{Inputs used throughout.}
\begin{itemize}
\item[(H0)] $-\Delta^{\eta,N}_{A,\mathcal X}+a_kP_k(A)\ge\gamma_0$ on every block union $\mathcal X$, for every b.c.\ and every regular $A$ \cite[(1.8), p.~573]{Ba4}.
\item[(W)] For $\abs\beta\le\beta_*$ and every $1$-Lipschitz $\rho$: $\norm{\mathrm{e}^{\beta\rho}G\mathrm{e}^{-\beta\rho}}\le2/\gamma_0$, and the weighted norms of $\nabla G$, $G\nabla^*$ and $\nabla G\nabla^*$ are $\le K_W$. This holds uniformly in $\eta$, the domain and the b.c.\ (Lemma~\ref{app:G1:WK-W}; literature form \cite[Cor.~2.3, (2.30), p.~580]{Ba4}).
\item[(K)] The pointwise bounds $\norm{G(x,y)}\le c_KK_2^{\delta_K}(x-y)$ and $\norm{D_xG(x,y)}\le c_KK_1^{\delta_K}(x-y)$, and for $\abs{x-\tilde x}\le1$ the $\alpha$-H\"older bound of $D_xG$ with $K_{1-\alpha}^{\delta_K}$. Here $K_\kappa^\delta(z):=\max(\abs z,\eta)^{\kappa-3}\mathrm{e}^{-\delta\abs z}$ for $0<\kappa<3$. These hold at points at distance $\ge R_0$ from $\mathcal X^c$ (on $T_\eta$ and on rectangular parallelepipeds at all points): \cite[(2.10)--(2.11), p.~426]{Ba3} summed via \cite[(2.6)]{Ba3}, from \cite[Props.~2.1, 2.3]{Ba1} (Lemma~\ref{app:G1:WK-K}).
\end{itemize}
The hypothesis (R$_\vartheta$) is: $\theta(\tilde B),\theta(\tilde A+\tilde B)\le\vartheta_R$ and $\sigma\le1$; it follows from (E.1) and (E.ksp3) of Appendix~\ref{app:B}.

\paragraph{The statement $\Sigma$ and the premises (statements 5, 5$'$).}
\begin{definition}[the statement $\Sigma$]\label{app:G1:I3r-KSigma}
Let $\mathcal R$ be the set of conditions of Tables~\ref{p:tab:A1}--\ref{p:tab:A5} other than (E.md), (E.ksp1)--(E.ksp3) and (E.14), and let $\mathcal R^+$ be $\mathcal R$ together with (E.md), (E.ksp1)--(E.ksp3) and (E.14). \emph{$\Sigma$} is the statement that every condition \emph{invoked} by the renormalisation step of \cite{Ba2} and \cite{Ba3} (including the pages of \cite{Ba1} they call) is a condition of $\mathcal R^+$ or is not a condition (a definition, a notation or a forward reference). Here each abbreviated step is replaced by its written argument. $\Sigma$ includes the statement that there is no other mass entry (Lemma~\ref{app:G1:I3r-ME}).
\end{definition}

\begin{lemma}[the abbreviated steps that need more than the printed text]\label{app:G1:I3r-Pw}
Assume the hypotheses of Theorem~\ref{thm:T6}, and let the regulator data satisfy the conditions of Definition~\ref{app:G1:rowlist}. Each step (a)--(i) below of \cite[pp.~555--594]{Ba2} and \cite[pp.~411--445]{Ba3} is stated there without a complete argument. Each of them holds. Its argument is written at the place named, and that argument uses, besides printed statements of \cite{Ba1,Ba2,Ba3,Ba4,K1}, only the conditions named.
\begin{enumerate}
\item[(a)] \cite[p.~573, (2.75)--(2.76)]{Ba2} (``so the same conclusion holds''): the scalar analogue on $\square$ of \cite[(2.62)--(2.63)]{Ba2}, including the mass defect left by (2.75). Place: Lemma~\ref{app:G1:M4D}. Condition: (E.md).
\item[(b)] \cite[Lemma~2.4]{Ba2} at the edge points not covered by its printed proof (Remark~\ref{rem:lem24-gap}). Place: Lemma~\ref{lem:lem24}. Conditions: (R.edge) and (E.edge) (Lemma~\ref{app:G1:R0rows}).
\item[(c)] \cite[p.~579]{Ba2}, the extension of \cite[Prop.~2.1]{Ba2} to $Q^{(k)}$. Place: Lemma~\ref{app:G1:RW-KSP}. Conditions: the $\eps_0$-bounds of that lemma.
\item[(d)] \cite[p.~579]{Ba2}, the cut of the expansion at the order $\nbar$. Place: Lemma~\ref{app:G1:I3r-T579}, with Propositions~\ref{app:G1:T579}, \ref{app:G1:I3c-R} and~\ref{app:G1:I3c-B}. The proof of Proposition~\ref{app:G1:T579} also uses Statement~\ref{app:K:L06-4}, which it assumes.
\item[(e)] \cite[p.~421]{Ba3}, the use of \cite[Lemma~2.4]{Ba2}. Places: as in (a) and (b). Conditions: (E.md), (R.edge) and (E.edge).
\item[(f)] \cite[(1.16), p.~414]{Ba3}, the two-variable remainder bound. Places: Lemmas~\ref{app:G1:I3r-U1m} and~\ref{app:G1:I3r-U1mOmega}.
\item[(g)] \cite[p.~433]{Ba3}, the remainder on the cube $\square$. Place: Lemma~\ref{app:G1:I3r-cube}, with Lemmas~\ref{app:G1:I3r-Wsq} and~\ref{app:G1:I3r-W2R}.
\item[(h)] The sketched steps S1--S21 and P12 of the proof of \cite[Prop.~1]{Ba3}. Places: Table~\ref{tab:prop1} and Proposition~\ref{prop:prop1}.
\item[(i)] \cite[(3.30), (3.32), p.~442]{Ba3}, the $A\cdot A$ local terms. Place: Lemma~\ref{app:G1:I3r-AA}, whose proof also uses Proposition~\ref{app:K:L10-2}.
\end{enumerate}
The conditions (E.md), (R.edge), (E.edge) and those of Lemma~\ref{app:G1:RW-KSP} are not printed in \cite{Ba1,Ba2,Ba3,Ba4}. In Definition~\ref{app:G1:rowlist}, (R.edge) is a lower bound on $R$, and the others are upper bounds on $\eps_0$.
\end{lemma}

\begin{proof}
Each item is the statement proved at the place named, and the conditions used in that proof are the ones named in the item. For (i), the proof at the place named also uses Proposition~\ref{app:K:L10-2}.
\end{proof}
\noindent That every other abbreviated step on these pages, and on the pages of \cite{Ba1} invoked there, holds under the hypotheses of Theorem~\ref{thm:T6} and the conditions of Definition~\ref{app:G1:rowlist} is Statement~\ref{app:K:L10-1}.

\paragraph{The remainder kernel (1.16): statements 1--3.}
For $n,n'\ge0$ and $m:=n+n'$, put $V:=G(\tilde B)^{-1}-G(\tilde A+\tilde B)^{-1}$. This is the bracket of \cite[(3.44)]{Ba1}, and the $m^2t^2$ terms cancel in it. Put
\[
K=K_{n,n'}:=[G(\tilde B)V]^nG(\tilde A+\tilde B)[VG(\tilde B)]^{n'}=L_0VL_1V\cdots VL_m,
\]
where $L_n=G(\tilde A+\tilde B)$ and $L_i=G(\tilde B)$ for $i\ne n$. This is the operator \cite[(1.16), p.~414]{Ba3}.

\begin{lemma}[vertex structure]\label{app:G1:I3r-V}
Put $\phi_\mu(x):=\eta^{-1}(1-U_{x,\mu}(\tilde A)^{-1})$, $\psi_\mu:=\phi_\mu^*\phi_\mu$, $Q:=Q_k(\tilde B)$, $F_2:=F_{2,k}(\tilde A,\tilde B)$ and $\Pi:=Q^*F_2+F_2^*Q+F_2^*F_2$. Then:
\begin{enumerate}
\item[(i)] $V=-\mathfrak D^*\mathfrak E\mathfrak D$, where $\mathfrak E$ has first row $(\sum_\mu\psi_\mu+a_k\Pi,\phi_1^*,\phi_2^*,\phi_3^*)$, first column $(\cdot,\phi_1,\phi_2,\phi_3)^T$ and zeros elsewhere.
\item[(ii)] (right-pushed form) $V=-\mathfrak r\mathfrak D$ with $\mathfrak r=R_0+\sum_\mu R_{1,\mu}$. Here $(R_{1,\mu}h)(x)=\phi_\mu(x)^*h(x)-\phi_\mu(x-\eta e_\mu)U_{x-\eta e_\mu,\mu}(\tilde B)^{-1}h(x-\eta e_\mu)$ and $R_0=\sum_\mu(\psi_\mu-M_\mu)+a_k\Pi$, where $M_\mu$ is multiplication by $(\partial^\eta_\mu\phi_\mu)(x-\eta e_\mu)$.
\item[(iii)] (left-pushed form) $V=V^*=-\mathfrak D^*\mathfrak r^*$.
\item[(iv)] For $\sigma\le1$, in $\ell^p(\eta^3)$ for $p\in\{1,2,\infty\}$: $\norm{\mathfrak r},\norm{\mathfrak E}\le c_V\sigma$ with $c_V=4d+a(2d+d^2)$. Every kernel has range $\le\sqrt3$, so $\norm{\mathfrak r^\beta},\norm{\mathfrak E^\beta}\le \mathrm{e}^{\beta\sqrt3}c_V\sigma$ for every $1$-Lipschitz weight $\rho$ and $\abs\beta\le1$. Moreover $V(x,\cdot)=0$ unless $\dist(x,\supp\tilde A)\le\sqrt3+\eta$.
\end{enumerate}
\end{lemma}
\noindent Ba{\l}aban's site form \cite[(3.44)]{Ba1} and the shifted form describe the same operator.
\begin{proof}
By \cite[p.~580]{Ba4}, $D_{\tilde A+\tilde B,\mu}=U_\mu(\tilde A)(D_{\tilde B,\mu}+\phi_\mu)$, since all transporters are functions of the generator $q$. $U$ is pointwise orthogonal, so bond by bond
\[
-\Delta_{\tilde A+\tilde B}=-\Delta_{\tilde B}+\sum_\mu(D_\mu^*\phi_\mu+\phi_\mu^*D_\mu+\psi_\mu).
\]
The same bonds are summed on $T_\eta$ and on a Neumann domain. Also $Q_k(\tilde A+\tilde B)=Q+F_2$. This gives (i).

For (ii): with $[\phi,U]=0$ one has $D_\mu^*(\phi_\mu f)(x)=-\phi_\mu(x-\eta e_\mu)U^{-1}(D_\mu f)(x-\eta e_\mu)-(\partial^\eta_\mu\phi_\mu)(x-\eta e_\mu)f(x)$. Inserting this into (i) gives (ii). (iii) holds because both inverses are real symmetric for real backgrounds.

For (iv): $\abs{\mathrm{e}^{q\vartheta}-1}\le\abs\vartheta$ gives $\abs{\phi_\mu}\le\sigma$ and $\abs{\partial^\eta\phi}\le\sigma$. Also $\abs{U(\tilde A(\Gamma))-1}\le d\sigma$, because the contours lie in $B^k(y)$ and have length $\le d$. Finally $\norm Q\le1$. Schur bounds give the constants. The support clause holds because each term carries $\phi$ or $F_{1}'$, which vanish off bonds or contours meeting $\supp\tilde A$.
\end{proof}

\begin{lemma}[one-variable remainder bound]\label{app:G1:I3r-U1}
Assume the hypotheses of \cite[p.~412]{Ba3} on $\tilde A,\tilde B$ and (R$_\vartheta$). Let $N:=n+n'\ge4$ and $\alpha\in(0,1)$, and write $K(x,y)$ for the kernel of the operator $K_{n,n'}$ above. Then for all $x,y$ at distance $\ge R_0$ from $\Omega^c$:
\[
\abs{K(x,y)}+\abs{D_xK(x,y)}+\sup_{\abs{x'-x}\le1}\abs{x-x'}^{-\alpha}\abs{U(\Gamma_{x,x'})D_xK(x',y)-D_xK(x,y)}\le C_N(C_*\eps_V)^N\mathrm{e}^{-\delta_N\abs{x-y}},
\]
and the same holds in the variable $y$ with $x$ fixed. Here $\eps_V=e_k\mathfrak a(1+e_k\mathfrak a)\mathrm{e}^{e_k\mathfrak a}\le Ce(t)p(t)$, $\mathfrak a=\sup(\abs{\tilde A}+\abs{\nabla\tilde A})$, $\delta_N=\delta_24^{-N}$ and $\delta_2=\delta_0/4$. The bound is per variable only; it contains no mixed derivative $D_xD_yK$.
\end{lemma}
\noindent The lemma does not give the two-variable form that \cite[(2.5)]{Ba3} is used for; that form is Lemma~\ref{app:G1:I3r-U1m}. The tree rate of Proposition~\ref{app:G1:I145-prop1base} is taken $\le\delta_N$, a choice made at stage 4 of Appendix~\ref{app:B}.
\begin{proof}
\emph{Inputs.} (I-G) is the pointwise kernel bound for $G$ and $D_{\mathcal B}G$ (decay $\mathrm{e}^{-\delta_2\abs{x-u}}$, singularities $\rho_u^{-1}$ and $\rho_u^{-2}$, where $\rho_u(x)=\abs{x-u}+\eta$). It is Lemma~\ref{app:G1:WK-K}. (I-R) is \cite[(2.24)--(2.25)]{Ba1}.

\emph{Convolution.} For $0<a,b<3$:
\[
\sum_u\eta^3\rho_x(u)^{-a}\rho_y(u)^{-b}\mathrm{e}^{-\delta(\abs{x-u}+\abs{u-y})}\le C\,\mathrm{e}^{-\delta\abs{x-y}/4}\Phi_{a+b}(\rho_y(x)),
\]
where $\Phi_s(\rho)=\rho^{3-s}$, $\rho^{-1/2}$ or $1$ according as $s>3$, $s=3$ or $s<3$. This is shown by splitting into the regions $\abs{u-x}\le r/2$, $\abs{u-y}\le r/2$ and the rest, as in Lemma~\ref{app:G1:conv}(a); it is uniform in $\eta$.

\emph{Reduction to $x$.} $V$ and both propagators are self-adjoint, so $K_{n,n'}(x,y)=K_{n',n}(y,x)^*$.

\emph{Chain.} Write $K=Y_N\cdots Y_1Y_0$ with $Y_0$ the rightmost propagator and $Y_i=G_iV$. Fix $y$ and set $u_0=G_0(\cdot,y)$ and $u_i=G_iVu_{i-1}$, so that $K(\cdot,y)=u_N$. Lemma~\ref{app:G1:I3r-V}(iv) and the convolution bound give successively $\abs{u_1}\lesssim\eps_V\rho_y^{-1/2}$, $\abs{u_2}\lesssim\eps_V^2$, and $u_3$, $Du_3$ bounded. From $i\ge4$ on, (I-R) applied to $f=Vu_{i-1}=\sum_\square f1_\square$ gives the sup, derivative and $\alpha$-H\"older bounds with decay. The decay rate is quartered at most once per step, which gives $\delta_N$.

\emph{Where the $R_0$ restriction is harmless.} Every point at which $Vu_{i-1}\ne0$ lies within $2\sqrt d+1$ of $\supp\tilde A$. The endpoints of (1.16) lines lie in $\Omega_1\cup U_{2\sqrt d+1}(\supp\tilde A)$. Moreover $\dist(\Omega_1,\Omega^c)\ge5r/L>R_0$ (Lemma~\ref{app:G1:Rnest}(v)), and $\dist(\supp\tilde A,\partial\Omega)$ exceeds the support margin, (R4$'''$) of Corollary~\ref{cor:prop1R4}. Hence (I-G) applies at every intermediate point, given $R>R_0$ and $R\ge12\sqrt d$ ((R.1) and (R.8)).
\end{proof}

\begin{lemma}[two-variable remainder bound on $T_\eta$]\label{app:G1:I3r-U1m}
Let $\Omega=T_\eta$, assume (R$_\vartheta$), and let $0<\beta\le\beta_U:=\min\{\beta_*,\delta_K/4\}$. For a kernel $Y$ and $a,b\in\{0,1\}$ put $Y_{ab}:=D^aYD^{b*}$. Define
\[
\mathfrak N_{\alpha,\beta}(Y):=\max_{a,b}\sup_{x,x'}\mathrm{e}^{\beta\abs{x-x'}}\max\{\norm{Y_{ab}(x,x')},\,a\,H_x,\,b\,H_{x'},\,ab\,H_{xx'}\}.
\]
Here $H_x$ and $H_{x'}$ are the $\alpha$-H\"older quotients in $x$ and in $x'$ over $0<\abs{x-\tilde x}\le1$. The transporter is $U(\tilde B(\Gamma_{\tilde x,x}))$, where $\Gamma_{\tilde x,x}$ is a shortest lattice contour from $\tilde x$ to $x$. $H_{xx'}$ is the mixed second difference divided by $\abs{x-\tilde x}^\alpha\abs{x'-\tilde x'}^\alpha$.
\begin{enumerate}
\item[(a)] (Operator form; every $m\ge0$, $\alpha<1$.)
\[
\abs{(D^aKf)(x)}\le c_KC_{(a)}(C_\infty\sigma)^m\mathrm{e}^{-\beta\dist(x,\supp f)}\norm f_\infty.
\]
The same holds for the $\alpha$-H\"older quotient of $DKf$ and for $K^*$.
\item[(b)] (Two-variable form.) If $m\ge4$ and $0<\alpha<\frac12$, or if $m\ge6$ and $0<\alpha<1$, then $\mathfrak N_{\alpha,\beta}(K)\le C_U(\alpha)(C_V\sigma)^m$, with $C_V:=\mathrm{e}^{\beta(\sqrt3+1)}c_V\max\{c_W,1\}$.
\item[(c)] (Form of \cite[(2.5)]{Ba3}.) Under the hypotheses of (b), let $0\le h,h'\le1$ with $\abs{\partial^\eta_\mu h},\abs{\partial^\eta_\nu\partial^\eta_\mu h}\le c_h$ for all $\mu,\nu$, and likewise for $h'$. Then
\[
\norm{hKh'}_{1,\alpha}\le \mathrm{e}^{5\beta}(1+6c_h)^2C_U(\alpha)(2c_1C_V)^m(e_kp(t))^m\mathrm{e}^{-\beta\dist(\supp h,\supp h')}.
\]
If the quotients over $\abs{x-\tilde x}>1$ are included, multiply the constant by $4$.
\end{enumerate}
Here $c_W:=(1+\sigma)^2(2\gamma_0^{-1}+((1+d)^2-1)K_W)$. All constants depend only on $d,a,M_b,L,\gamma_0,K_W,\alpha,\beta$. The bounds are uniform in $\eta$ (hence in $k$ and $\eps$), in $\ell_T$, in the backgrounds within (R$_\vartheta$), and in $0\le m^2t^2\le c$. In the construction $m=12$, so $\alpha_0=\frac12$ is covered.
\end{lemma}
\noindent The inputs are (H0) \cite[(1.8)]{Ba4}, (W) and (K); the tree rate of Proposition~\ref{app:G1:I145-prop1base} is taken $\le\beta_U$. Part (b) is Proposition~\ref{app:G1:sandwich}(a). We repeat the structure of its proof, and prove (a) and (c).
\begin{proof}
\emph{Step 1 (convolution and smearing).} By Lemma~\ref{app:G1:conv}(a), for $\kappa_i\in(0,3)$ and $\delta'<\delta$,
\[
\sum_w\eta^3K^\delta_{\kappa_1}(z-w)K^\delta_{\kappa_2}(w)\le c_CK^{\delta'}_{\kappa_1+\kappa_2}(z),
\]
with a logarithm at $\kappa_1+\kappa_2=3$ and boundedness above $3$. By Lemma~\ref{app:G1:conv}(c), a local operator of range $\le\sqrt3$ with entries $\le r_R$ maps $K_\kappa^\delta(\cdot-y)$ to $c_sr_RK_\kappa^\delta(\cdot-y)$.

\emph{Step 2 (end factors).} Define the two- and three-line ends $E^{(2)}_a:=D^aL_0VL_1$ and $E^{(3)}_a:=D^aL_0VL_1VL_2$. Write every vertex in right-pushed form (Lemma~\ref{app:G1:I3r-V}(ii)). Then each line of an end carries exactly one derivative. With $\delta'=\delta_K/2$, (K) together with Step~1 gives $\norm{E^{(2)}_a(x,y)}\le c_E\sigma K_2^{\delta'}(x-y)$ and $\norm{E^{(3)}_a(x,y)}\le c_E'\sigma^2K_3^{\delta'}(x-y)$. The H\"older differences give $K_{2-\alpha}$, respectively $K_{3-\alpha}$. In weighted $\ell^2$, the end factors are bounded by $J_\kappa:=(\sum_y\eta^3K_\kappa^{\delta'}(y)^2\mathrm{e}^{2\beta\abs y})^{1/2}$, and $J_\kappa<\infty$ uniformly in $\eta$ if and only if $\kappa>\frac32$. This explains the thresholds: $J_{2-\alpha}$ is finite if and only if $\alpha<\frac12$, and $J_{3-\alpha}$ is finite for every $\alpha<1$.

\emph{Step 3 (sandwich, $m\ge4$).} Write $K_{ab}=E_L\mathcal N E_R$ with $E_L=E^{(2)}_a$, $E_R=(D^bL_mVL_{m-1})^*$ and $\mathcal N=VL_2V\cdots VL_{m-2}V$. Put the first vertex of $\mathcal N$ in right-pushed form, the last in left-pushed form, and the others in symmetric form:
\[
\mathcal N=\pm\mathfrak r(\mathfrak DL_2\mathfrak D^*)\mathfrak E\cdots\mathfrak E(\mathfrak DL_{m-2}\mathfrak D^*)\mathfrak r^*.
\]
Every factor is bounded in weighted $\ell^2$; in particular $\norm{(\mathfrak DL\mathfrak D^*)^\beta}\le c_W$ by (W). For $L=G(\tilde A+\tilde B)$ use $D_{\tilde B,\mu}=U_\mu(\tilde A)^{-1}D_{\tilde A+\tilde B,\mu}-\phi_\mu$. With the vertex range of Lemma~\ref{app:G1:vertex}(iv), $\norm{\mathcal N^{\pm\beta}}\le(\mathrm{e}^{\beta(\sqrt3+1)}c_V\sigma)^{m-2}c_W^{m-3}$. Cauchy--Schwarz at the two end factors, with $\rho=\abs{\cdot-x}$ and $\mathrm{e}^{-\beta\abs{y'-x}}\le \mathrm{e}^{-\beta\abs{x-x'}}\mathrm{e}^{\beta\abs{y'-x'}}$, gives
\[
\mathrm{e}^{\beta\abs{x-x'}}\norm{K_{ab}(x,x')}\le(c_EJ_2\sigma)^2(\mathrm{e}^{\beta(\sqrt3+1)}c_V\sigma)^{m-2}c_W^{m-3}.
\]
No pointwise bound on $D_xD_yG$ is used. $K_{ab}$ is linear in each end, so the H\"older quotients in $x$, in $x'$ and jointly follow with $c_EJ_2$ replaced by $2\mathrm{e}^\beta c_E(\alpha)J_{2-\alpha}$. For $m\ge6$, repeat with three-line ends ($J_3$, $J_{3-\alpha}$).

\emph{Step 4 ((a) and (c)).} For (a): $\norm{VL}_{\infty,\beta}\le C_\infty\sigma$ by Lemma~\ref{app:G1:I3r-V}(ii) and (K), and the norm is submultiplicative. This is the argument of \cite[(2.33)]{Ba4} applied to one term.

For (c): apply Lemma~\ref{app:G1:leibniz-loc} with $\mathcal P=T_\eta$ and $M=C_U(\alpha)(C_V\sigma)^m$ from (b). The quotient of the undifferentiated component is telescoped along the shortest lattice contour that carries the transport, so no holonomy correction arises. Finally, $\sigma\le2c_1e_kp(t)$.

\emph{Remark (the threshold).} At $m=0$ the mixed diagonal kernel is of order $\eta^{-3}$ (see the remark after Proposition~\ref{app:G1:sandwich}). For $m\le3$, power counting suggests $\abs{x-x'}^{m-3}$. The case $\frac12\le\alpha<1$ with $m\in\{4,5\}$ is not treated and not needed.
\end{proof}

\begin{lemma}[two-variable remainder bound on Neumann $\Omega$, restricted form]\label{app:G1:I3r-U1mOmega}
Let $\Omega=B^k(\Lambda_2^{(k-1)\prime})$ with Neumann conditions. $\Omega$ may be any union of big blocks, including faces and re-entrant edges and corners, and may wrap the torus. Then the bounds of Lemma~\ref{app:G1:I3r-U1m}(b),(c) hold with the $T_\eta$ constants, uniformly in $\eta,\eps,\ell_T$ and $\Omega$, at every point that the path evaluates. These are the external points, all vertex points and their H\"older partners, and they lie at $\ell^\infty$-distance $\ge2R-O(1)$ from $\Omega^c$.
\end{lemma}
\noindent The inputs are Lemma~\ref{app:G1:I3r-U1m}, \cite[Prop.~2.1]{Ba1}, (W), the conditions (R.1) (with an $O(1)$ shift) and (R.8) ($R\ge12\sqrt d$) of Appendix~\ref{app:B}, the support hypothesis of \cite[p.~412]{Ba3} in the form (R4$'''$), (E.1) and (E.ksp3), and the conventions ($g_k$ gating, $\ell^\infty$ metric \cite[(1.3)]{Ba1}, contours inside blocks, the nesting of \cite[(2.8)]{Ba2}: statement 15 and Lemma~\ref{app:G1:Rnest}). The unrestricted form is not expected to be uniform near re-entrant edges (at the level of derivatives, of order $\eta^{-1/3}$); those points are not evaluated.
\begin{proof}
(1) \emph{Restriction to $\supp V$.} $\mathcal N=VL\cdots LV$ begins and ends with a vertex, and $V(y,\cdot)=0$ unless $y$ is within one bond or one unit block of $\supp\tilde A$ (Lemma~\ref{app:G1:I3r-V}(iv)). Hence $K_{ab}(x,x')=\langle E_L(x,\cdot)\chi,\mathcal N\chi E_R(\cdot,x')\rangle$ with $\chi=1_{\supp V}$. So only $y,y'\in\supp V\subset(\supp\tilde A)^{+\sqrt3+\eta}$ enter, and every pointwise evaluation of (K) in the end factors is at a vertex point. $V=-\mathfrak D^*\mathfrak E\mathfrak D$ is exact on $\Omega$ with no face term, because both operators have the same Neumann bond set and the same $m^2$, and the $Q$-contours lie inside unit blocks.

(2) \emph{Every endpoint is deep.} \cite[(2.5)]{Ba3} and \cite[(1.33)]{Ba3} take the $1,\alpha$ norm of $h(1.16)h'$ on \emph{every} line, so derivative-free legs must be traced as well. All endpoints lie in $U_{O(1)}(\Omega_1)\cup U_{O(1)}(\supp\tilde A)$. The block-spin legs (1.13)--(1.15) are gated by $g_kA'$ (from the $T^\eta_{\dots,e'g_kA'+A^{(k)}}$ of \cite[(1.4)]{Ba3}) or by $\tilde A$. The $n=0$ $\tilde A$-vertices sit at $\supp\tilde A$.

(3) \emph{Margin.} Under the literal hypothesis $\dist(\supp\tilde A,\partial\Omega)>2r$ the margin is $2R-O(1)$. Under (R4$'''$) of Corollary~\ref{cor:prop1R4} the vertex points satisfy
\[
\dist(\cdot,\Omega^c)>\tfrac74r-\sqrt3-1\ge R_0+2\sqrt3+1
\]
whenever $R\ge3R_0+15$. The proof uses only a margin of type $R_0$ at these points; the weaker margin $2r_k-M_b-5$ would also suffice.

(4) \emph{Middle factor.} The middle factor is uniform in the domain: Lemma~\ref{app:G1:WK-W} applies to a union of whole blocks with any b.c.\ and a $1$-Lipschitz weight. The end factors use \cite[Prop.~2.1]{Ba1} only at admissible pairs. With (1)--(3), the proof of Lemma~\ref{app:G1:I3r-U1m} goes through.
\end{proof}

\paragraph{No other mass entry: statement 4.}
\begin{lemma}[no other mass entry]\label{app:G1:I3r-ME}
A \emph{mass entry} is a term $c(x)\abs{\phi(x)}^2$ or $c(x)A_\mu(x)A_\nu(x)$ of an effective action, or of a quadratic form defining a covariance, whose coefficient $c$ does not depend on any field. Consider the step-$k$ expansion of \cite{Ba3}: vertices (1.6)--(1.15); lines $G_k(\Omega_2,\tilde B)$, $\delta G_k$ and (1.16); the vector $G_k$. Through the renormalisation of \cite[Prop.~1]{Ba3} (pp.~424--445) and the localisations (pp.~420, 422, 432--433):
\begin{enumerate}
\item[(i)] every field-independent coefficient of $\abs\phi^2$ is either a piece of $\delta m^2_{\rm fin}$, $\delta m^2_{K,k}$, $\delta m^2_k$ (or their localisations) or the covariance entry $m^2t^2$;
\item[(ii)] no field-independent coefficient of $A\cdot A$ is added to any covariance, and there is no photon mass counterterm;
\item[(iii)] every other field-independent quadratic local term is an interaction vertex of $\mathcal P$ bounded inside \cite[(1.33)]{Ba3} with at least one coupling factor, and its $M$-dependence is only through $m^2t^2$;
\item[(iv)] hence every occurrence of the Higgs mass in \cite{Ba3} is of type ($\alpha$) or ($\beta$) of Lemma~\ref{app:G1:mass}(ii), or free of the mass.
\end{enumerate}
\end{lemma}
\noindent The lemma uses the classification of the two-leg classes (Lemma~\ref{app:G1:cutclasses}) and the written formulas of \cite[pp.~412--445]{Ba3}. It does not depend on Lemma~\ref{app:G1:I3r-U1}.
\begin{proof}
The divergent two-leg classes are C1, C1$'$, C2 (scalar) and C3, C4 (photon). Every other two-leg graph has $D>0$ and is bounded as a non-local expression, so it produces no local term.
\begin{itemize}
\item \emph{C1} (pp.~435--437). The tadpoles are cancelled exactly by their counterterm insertions (3.7). For the exchange graph, the zeroth moment is cancelled exactly by the counterterm of (3.9), whose counterterm is the same graph with both legs at $x$. The first moment is the derivative vertex (3.15). The remainder (3.12) is non-local.
\item \emph{C1$'$} gives only first-derivative vertices.
\item \emph{C2}: the zeroth moment is cancelled exactly by (3.19).
\item \emph{Photon classes.} Under the $C^\xi$ replacement, (3.27)--(3.28) kill the $C^\xi$ part of the first moment $\Pi_{\mu\mu'\nu}$, and (3.29) (Ward--Takahashi identity at $M^2=1$) kills the zeroth moment $\Pi_{\mu\mu'}$.
\end{itemize}
The localisations are partitions of unity. The cuts of p.~432 replace lines by pairs of external fields. $\delta G_k(\Omega,\square,\tilde B)$ and $\delta G_k(\square,\eta\Z^d,0)$ are non-local two-variable external fields. The $B_0$ gauge-out is exact. The $B'$-expansion produces only vertices with $B'$ legs. Finally, \cite[(1.20), p.~416]{Ba3} has $\delta m^2$ for the scalar only. Together these give (i)--(iii), and (iv) follows from (i)--(iii) and the list of the mass occurrences in Lemma~\ref{app:G1:mass}(ii).

\emph{The $A\cdot A$ vertices.} The rewritings (3.30) and (3.32) leave field-independent local $A\cdot A$ (``$\sim\!\ast\!\sim$'') vertices of order $e(t)^2$. Their coefficient is $(L^{j_0}\eta)^{-d+2}=(L^{j_0}\eta)^{-1}$ times a convergent expression built from $\Upsilon=1-m^2t^2-aP$. They are mass-\emph{type} vertices of $\mathcal P$, not mass entries. They are harmless for Theorem~\ref{thm:T6} by Lemma~\ref{app:G1:I3r-AA}.
\end{proof}

\paragraph{The remainder on the cube $\square$ of \texorpdfstring{\cite[p.~433]{Ba3}}{p. 433}: statements 6, 7.}
Setting (p.~433). $\square$ is a cube of side $3r$ and $\square_1$ the concentric cube of side $r$; both are unions of big blocks, and $d(z):=\dist(z,\square^c)$. Put $G_0=G_k(\square,B_0)$ with $B_0$ constant and $G_1=G_k(\square,B_0+B')$, both Neumann. Put $\sigma':=e_k(\norm{B'}_\infty+\norm{\partial^\eta B'}_\infty)\le c_2'e_kp^2$. The remainder is $K^\square=L_0VL_1\cdots VL_m$ with $m\ge6$.

\begin{lemma}[separation with one interior point, $\ell^2$ in the other]\label{app:G1:I3r-SEP}
Let $A\in\{B_0,B_0+B'\}$ be in the class of \cite[Prop.~2.1]{Ba1} on $\square$, with $\square$ a union of big blocks. Let $L=G_k(\square,A)$, $a\in\{0,1\}$, $\rho\ge16$, and $z\in\square$ with $d(z)\ge\rho/2+R_0+4$. Let $S\subset\square$ be arbitrary ($S$ may meet $\partial\square$) with $\dist(z,S)\ge\rho$. Then for $0\le\beta\le\delta_0/8$:
\begin{enumerate}
\item[(i)] $\norm{\mathrm{e}^{\beta\abs{\cdot-z}}1_S(D^aL\mathfrak D^*)(z,\cdot)}_{\ell^2(\eta^3)}\le C_*\mathrm{e}^{-\delta_*\rho}$ with $\delta_*:=\delta_1/32$;
\item[(ii)] for $\abs{z-\tilde z}\le1$ and $a=1$, the $\alpha$-H\"older version of (i) holds with constant $C_*(\alpha)$;
\item[(iii)] the mirror statement holds in the first variable.
\end{enumerate}
The constants are uniform in $\eta$, $\ell_T$, $r(t)$, $m^2t^2$, $B_0$ and $B'$. No constancy of $A$ near $\partial\square$ is assumed.
\end{lemma}
\noindent A pointwise bound on $(DG_1\mathfrak D^*)(z,w)$ for $w$ on $\partial\square$ is not used.
\begin{proof}
\emph{Inputs.} (P1) is the following interior bound: for $u,z$ with $d(u),d(z)\ge R_0$ and $\abs{z-u}\ge2$, $\abs{(D^aLD_A^{b*})(z,u)}\le c_S\mathrm{e}^{-\delta_1\abs{z-u}/2}$, with H\"older quotient $\le c_S(\alpha)\mathrm{e}^{-\delta_1\abs{z-u}/4}$. This follows from \cite[(2.6), (2.10)--(2.11)]{Ba3} and \cite[Prop.~2.1]{Ba1}. (P2) is \cite[Cor.~2.3, (2.30), pp.~580--581]{Ba4}: the four forms $\langle f,Gf'\rangle$, $\langle f,DGf'\rangle$, $\langle f,GD^*f'\rangle$, $\langle f,DGD^*f'\rangle$ are $\le c_0\mathrm{e}^{-\delta_0\dist(\supp f,\supp f')}\norm f_2\norm{f'}_2$ for $\square$ a union of big blocks. This rests on \cite[Lemma~2.1]{Ba4} only; the constancy near $\partial\square$ is a hypothesis of Lemma~2.2 alone. (P3) is the conversion $D_{B_0}=U(B')^{-1}D_A-\phi$, a bounded same-site multiplication with $\abs\phi\le\sigma'$.

Step 1 (duality). $\norm{1_{S'}K(z,\cdot)}_2=\sup\{\abs{(Kf)(z)}:\norm f_2\le1,\ \supp f\subset S'\}$. Put $g=L\tilde{\mathfrak D}^*f$.

Step 2 (interior cut-off). Let $\chi=\varphi(\abs{\cdot-z})$, with $\varphi=1$ on $[0,\rho/4]$, $\varphi=0$ on $[\rho/2,\infty)$ and linear in between. Then $\chi L^{-1}g=0$, so $\chi g=L[L^{-1},\chi]g=L\tilde{\mathfrak D}^*h$. Here $h$ is supported in $A'=\{\rho/4-2\le\abs{y-z}\le\rho/2+2\}$, and $\norm h_2\le(c_d/\rho)(\norm{1_{A'}g}_2+\norm{1_{A'}D_Ag}_2)$. All links used are interior ($d\ge R_0+2$), so no boundary term appears.

Step 3. By (P1), $\abs{D^ag(z)}\le c_S\mathrm{e}^{\delta_1}(5\rho)^{3/2}\mathrm{e}^{-\delta_1\rho/8}\norm h_2$.

Step 4 (far field in $\ell^2$ only). By (P2), $\norm{1_{A'}g}_2+\norm{1_{A'}D_Ag}_2\le2(d+1)^2c_0\mathrm{e}^{-\delta_0\dist(A',S')}\norm f_2$.

Step 5. Decompose $S$ into dyadic shells $S_n=S\cap\{2^n\rho\le\abs{\cdot-z}<2^{n+1}\rho\}$ and sum, using $\beta\le\delta_0/8$. This gives (i) for $\tilde{\mathfrak D}$; (P3) gives (i) for $\mathfrak D$, with $C_*\mapsto4C_*$.

Step 6. The same steps with the H\"older line of (P1) give (ii) with $\delta_1/32$. (iii) follows from $L^*=L$. No pointwise value within $\rho/2$ of $\partial\square$ is ever taken.
\end{proof}

\begin{lemma}[the remainder on the cube $\square$]\label{app:G1:I3r-cube}
Let $m\ge6$, $0<\alpha<1$ and $0<\beta\le\beta_\square:=\min\{\beta_U,\delta_1/32,\delta_0/8\}$. Take $D$ and the transports in $\mathfrak N_\alpha$ with respect to the same background. Let $0\le h,h'\le1$ with $\abs{\partial^\eta_\mu h},\abs{\partial^\eta_\nu\partial^\eta_\mu h}\le c_h$ for all $\mu,\nu$, and likewise for $h'$. Assume $\dist_\infty(\supp h\cup\supp h',\square_1^c)\ge2$. Then
\[
\sup_{x,x'\in\square_1}\mathrm{e}^{\beta\abs{x-x'}}\cdot(\text{all components of }\mathfrak N_\alpha(K^\square))\le C_\square(\alpha)(C_\square\sigma')^m ,
\]
\[
\norm{hK^\square h'}_{1,\alpha}\le \mathrm{e}^{5\beta}(1+6c_h)^2C_\square(\alpha)(C_\square\sigma')^m\mathrm{e}^{-\beta\dist(\supp h,\supp h')} .
\]
Here $C_\square=2C_0$ and $C_0=2\mathrm{e}^{\beta\sqrt3}c_V'\max\{c_W,c_S,C_*,C_*(\alpha),c_K,1\}$, with $c_V'=c_V+2d$. The bounds are uniform in $\eta$, $\ell_T$, $r(t)$, $m^2t^2$, $B_0$ and $B'$. The $\eta^{-1}$ boundary term never enters.
\end{lemma}
\noindent The inputs are Lemmas~\ref{app:G1:I3r-U1m}, \ref{app:G1:I3r-Wsq} and \ref{app:G1:leibniz-loc}, (H0), (W), (E.1), (R.1), and the choice $r(t)\ge\max\{107,3R_0+15,\tfrac43(R_0+4),O(\nbar)+4\sqrt3+2\}$, which does not depend on $\eps$. For the cut-offs of Lemma~\ref{app:G1:Rloc7}(iii), the support hypothesis holds by Lemma~\ref{app:G1:cube-margin}.
\begin{proof}
(1) \emph{The defect.} Lemma~\ref{app:G1:I3r-V}(i) holds \emph{exactly} on the Neumann cube, boundary sites included: $V=-\mathfrak D^*\mathfrak E\mathfrak D$ with $\norm{\mathfrak E}\le c_V\sigma'$. The unpaired term $-\eta^{-1}\phi_\mu(z)f(z)$ appears only after a vertex touching the collar is right- or left-pushed. It is a property of the rewriting, not of $V$.

(2) \emph{Split.} Put $\theta(z)=\min\{1,\max\{0,3d(z)/r-1\}\}$ and
\[
V=V_\theta+V_c,\qquad V_\theta=-\mathfrak D^*\Theta\mathfrak E\Theta\mathfrak D,\qquad V_c=-\mathfrak D^*(\mathfrak E-\Theta\mathfrak E\Theta)\mathfrak D.
\]
$V_\theta$ satisfies Lemma~\ref{app:G1:I3r-V} with $c_V'$ and has no unpaired term. $V_c$ is supported in $\mathcal C\times\mathcal C$, where $\mathcal C=\{d<2r/3+2\sqrt3\}$, and $\norm{\mathfrak E_c}\le2c_V\sigma'$. For $x\in\square_1$ and $w\in\mathcal C$, $\abs{x-w}\ge s:=r/3-2\sqrt3\ge r/6$. Grouping by the first and last $c$-vertex gives $K^\square=K_\theta+\sum_{i_1\le i_2}K_{i_1,i_2}$, a sum of at most $m(m+1)/2+1$ terms.

(3) \emph{$K_\theta$.} The proof of Lemma~\ref{app:G1:I3r-U1m}(b) with three-line ends applies. Given $r\ge3(R_0+2+\sqrt3+1)$, every pointwise evaluation of (K) is between points with $d\ge R_0+2$: external points have $d\ge r$, and $V_\theta$-sites have $d\ge r/3-\sqrt3-\eta$. So only interior points are used, and the parallelepiped sentence of \cite[p.~611]{Ba1} is not needed.

(4) \emph{$K_{i_1,i_2}$.} Keep both $V_c$ and every full $V$ between them in symmetric form, and push the $V_\theta$ outwards. Then $D^aK_{i_1,i_2}D^{b*}(x,x')=\pm\langle\mathcal L_x,\mathcal N\mathcal R_{x'}\rangle$. Here $\mathcal N=\mathfrak E_c(\mathfrak DL\mathfrak D^*)\mathfrak E\cdots\mathfrak E_c$ is bounded in weighted $\ell^2$ by (W), for any boundary condition. The end factors $\mathcal L_x,\mathcal R_{x'}$ are needed only in weighted $\ell^2(\mathcal C)$. Split the points $z\in B_x=\{\abs{z-x}<s/2\}$ from those with $z\notin B_x$.
\begin{itemize}
\item On $B_x$, use Lemma~\ref{app:G1:I3r-SEP} with $\rho=s/2$. Its hypotheses hold if and only if $r\ge3(32+2\sqrt3)$ and $\frac34r\ge R_0+4$.
\item Off $B_x$, use (W).
\end{itemize}
A pointwise separation bound for $G_1$ near $\partial\square$ is not available, because the proof of \cite[p.~581]{Ba4} needs $A$ to be constant near $\partial\square$; Lemma~\ref{app:G1:I3r-SEP} avoids it. The H\"older bound in $x$ works for every $\alpha<1$. The assembly is that of Step~3 of the proof of Lemma~\ref{app:G1:I3r-U1m}.

(5) Summing, $m(m+1)/2+2\le2^m$. The $h,h'$ form follows from Lemma~\ref{app:G1:leibniz-loc} with $\mathcal P=\square_1$: by the support hypothesis, every lattice point within $\ell^\infty$-distance $1+\eta\le\frac32$ of $\supp h\cup\supp h'$ lies in $\square_1$.
\end{proof}

\begin{lemma}[support margin for the remainder on $\square$]\label{app:G1:cube-margin}
Let $d$ be the $\ell^\infty$ distance on the torus, $B_\varrho(c):=\{z:d(z,c)\le\varrho\}$, $s_b$ the side of a big block, and $r=r(t)=R(1+\log t^{-1})^2$ with $t\le\eps_0<\mathrm{e}^{-1}$. Let $K$ be a component of the step of \cite[p.~433]{Ba3} that contains an old vertex. Let $Q$ be a closed cube of side $\ell$ containing the localisation cubes $\square(u)$ of all old vertices $u$ of $K$, and let $c\in Q$. Put
\[
\square_1:=\bigcup\{b:\ b\cap B_{r/2}(c)\ne\emptyset\},\qquad\square:=\bigcup\{b:\ b\cap B_{3r/2+s_b}(c)\ne\emptyset\},
\]
where the unions run over big blocks $b$. Let $w\ge0$, let $v,v'$ be old vertices of $K$, and let $\supp h\subset U_w(\square(v))$ and $\supp h'\subset U_w(\square(v'))$, where $U_w$ is the closed $\ell^\infty$-neighbourhood of width $w$. If $r\ge2\ell+2w+4$, then:
\begin{enumerate}
\item[(i)] $Q\subset\square_1\subset\square$, and $d(x,y)>r$ and $\abs{x-y}>r$ for all $x\in\square_1$ and $y\notin\square$;
\item[(ii)] $\dist_\infty(\square(u),\square_1^c)>w+2$ for every old vertex $u$ of $K$;
\item[(iii)] $\dist_\infty(\supp h\cup\supp h',\square_1^c)>2$.
\end{enumerate}
For the cut-offs of Lemma~\ref{app:G1:Rloc7}(iii) one has $w=1$. With $\ell=\ell_{\Nbar}$, the condition $r\ge2\ell+6$ holds for every $t\le\eps_0$ under the lower bound $R\ge2\ell_{\Nbar}+10M_b+10$ of Lemma~\ref{app:G1:I3r-Wsq}. In that case, let $h,h'$ be these cut-offs of the old vertices $v,v'$ at the two ends of the line carrying $K^\square$. Assume the remaining hypotheses of Lemma~\ref{app:G1:I3r-cube} on the sets $\square_1\subset\square$ above: $m\ge6$, $0<\alpha<1$, $0<\beta\le\beta_\square$, and $D$ and the transports taken with respect to the same background. Then the second display of Lemma~\ref{app:G1:I3r-cube} holds with $c_h=40/\sqrt3$.
\end{lemma}
\begin{proof}
\emph{Step 1 (the hulls).} A big block has $\ell^\infty$-diameter at most $s_b$, and the big blocks cover the torus. Hence $x\in\square_1$ implies $d(x,c)\le r/2+s_b$, and $y\notin\square_1$ implies $d(y,c)>r/2$. Likewise, $y\notin\square$ implies $d(y,c)>\frac32r+s_b$.

\emph{Step 2 (part (i)).} Every $q\in Q$ satisfies $d(q,c)\le\ell\le r/2$, so $Q\subset B_{r/2}(c)\subset\square_1$. For $x\in\square_1$ and $y\notin\square$, $d(x,y)>\frac32r+s_b-\frac r2-s_b=r$. In particular $\square_1\subset\square$. Since the Euclidean distance dominates the $\ell^\infty$ distance, also $\abs{x-y}>r$.

\emph{Step 3 (part (ii)).} Since $\square(u)\subset Q$, for $q\in\square(u)$ and $y\notin\square_1$ we have $d(q,y)>\frac r2-\ell\ge w+2$.

\emph{Step 4 (part (iii)).} For $z\in\supp h$ there is $q\in\square(v)$ with $d(z,q)\le w$. So for $y\notin\square_1$, $d(z,y)>w+2-w=2$; the same holds for $h'$.

\emph{Step 5 (the cut-offs and $R$).} In Lemma~\ref{app:G1:Rloc7}(iii), $h(x)=\prod_\mu\chi(t_\mu(x))$ with $\chi=0$ on $[1,\infty)$. So $h(x)\ne0$ implies $\dist_\infty(x,\square(v))<1$, and $w=1$. For $t\le\eps_0$ we have $r(t)\ge R$. Since $M_b\ge1$, $R\ge2\ell_{\Nbar}+10M_b+10\ge2\ell_{\Nbar}+20>2\ell_{\Nbar}+6$.

Finally, by (i) the sets $\square_1\subset\square$ satisfy $\dist(\square_1,\square^c)\ge r$, and $\square$ is a box and a union of big blocks; this is the only geometry of $\square_1\subset\square$ used in the proof of Lemma~\ref{app:G1:I3r-cube}. By (iii) and Lemma~\ref{app:G1:ramp-ch}, the hypotheses on $h,h'$ of the second display of Lemma~\ref{app:G1:I3r-cube} hold with $c_h=40/\sqrt3$. Together with the remaining hypotheses assumed in the statement, this gives the second display.
\end{proof}

\paragraph{Two-domain bounds $\delta G$: statement 8.}
\begin{lemma}[two-domain bounds]\label{app:G1:I3r-DGM}
Let $X_{\rm in}\subset X_{\rm out}$ be unions of big blocks of $T_\eta$, $\eta=L^{-k}\le\frac12$, where $X_{\rm out}$ may be $T_\eta$ and $X_{\rm in}\ne T_\eta$. Let $A$ be real with $\theta(A)\le\vartheta_R$, let $0\le m^2t^2\le c$, and put $\delta:=G_k(X_{\rm out},A)-G_k(X_{\rm in},A)$ and $\rho_{\rm in}:=\dist_\infty(\cdot,X_{\rm in}^c)$. Let $0<\alpha<1$ and $0<\beta\le\beta_D:=\min\{\beta_*,\delta_K/4\}$. Then for $x,x'$, and their H\"older partners, with $\rho_{\rm in}\ge R_0+11$:
\[
\mathrm{e}^{\beta\abs{x-x'}}\max\{\norm{X_{ab}(x,x')},aH_x,bH_{x'},abH_{xx'}\}\le C_D(\alpha)\exp\!\big(-\tfrac{\delta_K}2[(\rho_{\rm in}(x)-R_0-7)+(\rho_{\rm in}(x')-R_0-10)]\big),
\]
where $X_{ab}=D^a\delta D^{b*}$. $C_D$ depends only on $d,a,M_b,L,\gamma_0,K_W,\alpha$ and on the profiles $f,\varphi$ of Lemma~\ref{app:G1:collar}.

\emph{(loc)} Let $0\le h,h'\le1$ satisfy $\abs{\partial^\eta_\mu h},\abs{\partial^\eta_\nu\partial^\eta_\mu h}\le c_h$ for all $\mu,\nu$, and likewise for $h'$. If $\dist_\infty(\supp h\cup\supp h',X_{\rm in}^c)\ge R_0+12$ (and $\ge R_0+13$ if $k=1$), then
\[
\norm{h\delta h'}_{1,\alpha_0}\le C_{\rm loc}\mathrm{e}^{-\beta\dist(\supp h,\supp h')},\qquad C_{\rm loc}=\mathrm{e}^{5\beta}(1+6c_h)^2C_D(\alpha_0).
\]
This applies to the pairs $(X_{\rm out},X_{\rm in})=(\Omega,\Omega_2),(\Omega,\square),(T_\eta,\square)$.
\end{lemma}
\noindent The inputs are (H0), (W), (K) and the pushing identity \cite[(2.32)]{Ba4}. At every $k\ge1$, the point margin is $R_0+10+2\eta\le R_0+11$, and the support margin is $R_0+11+3\eta$ (Proposition~\ref{app:G1:twodomain} and Corollary~\ref{app:G1:twodomain-loc}). At the uses the margins are met; see Remark~\ref{app:G1:twodomain-inst}.
\begin{proof}
The setting is that of Proposition~\ref{app:G1:twodomain}, with $\rho=\rho_{\rm in}$. Points with $\rho_{\rm in}\ge R_0+11$ lie in $\mathcal P_D=\{\rho\ge R_0+10+2\eta\}$, so the first display is Proposition~\ref{app:G1:twodomain}. Part (loc) is Corollary~\ref{app:G1:twodomain-loc} with $\alpha=\alpha_0$, since $R_0+12\ge R_0+11+3\eta$ for $k\ge2$ and $R_0+13\ge R_0+11+3\eta$ for $k=1$.

We recall the structure of the proof.
\begin{itemize}
\item The cut-offs are $\chi=\chi_{R_0+3}$ and $\psi=\chi_{R_0+13/2}$ of Lemma~\ref{app:G1:collar}, for every $\eta\le\frac12$.
\item The commutators $C_\chi=-\mathfrak r_\chi\mathfrak D=\mathfrak D^*\mathfrak r_\chi^*$ of Lemma~\ref{app:G1:commutator} are supported in $S_\chi\subset\{R_0+\frac74-\eta<\rho_{\rm in}<R_0+\frac{21}4+2\eta\}\subset\{\psi=0\}$ and $S_\psi\subset\{R_0+\frac{21}4-\eta<\rho_{\rm in}<R_0+\frac{35}4+2\eta\}$.
\item This gives the two-commutator representation $X_{ab}=-E_L\,\mathfrak M\,E_R$, with $E_L=D^aG_{\rm out}\mathfrak r_\chi$, $\mathfrak M=\mathfrak DG_{\rm in}\mathfrak D^*$ and $E_R=\mathfrak r_\psi^*G_{\rm in}D^{b*}$.
\item The ends are one-derivative kernels at pairs $\ge1$ apart and at depth $>R_0+\frac54$ in $X_{\rm in}$. So (K) bounds them in weighted $\ell^2$, with the gains $\mathrm{e}^{-(\delta_K/2)(\rho_{\rm in}-R_0-7)}$ and $\mathrm{e}^{-(\delta_K/2)(\rho_{\rm in}-R_0-10)}$.
\item The middle factor is bounded by (W), and Cauchy--Schwarz concludes.
\end{itemize}
No kernel is evaluated at points less than $1$ apart, or within $R_0+1$ of $X_{\rm in}^c$.
\end{proof}

\paragraph{The cube lies in the outer domain: statement 9.}
\begin{lemma}[the cube lies in the outer domain]\label{app:G1:I3r-Wsq}
Let $s_b\le2M_b$ be the big-block side and $\ell:=\ell_{\Nbar}=2V(\Nbar)$ with $\Nbar=169$, where $V(\Nbar)$ is the maximal number of localisation cubes of an expression. Assume the conditions (T.nw) ($\ell_0\ge\eps_0[8r(\eps_0)+4M_b+2]$), (E.r) ($\eps_0<\mathrm{e}^{-1}$), (T.div) (the block tiling is compatible with the torus) and
\[
\text{(W-R}'\text{)}\qquad R\ge2\ell_{\Nbar}+10M_b+10.
\]
(W-R$'$) is implied by $R\ge R^{**}$. Let $K$ be a component of the p.~433 step with an old vertex, let $Q$ be the closed cube of side $\ell$ containing its localisation cubes, $c\in Q$, and $\square$ the big-block hull of $B_{3r_\square/2+s_b}(c)$.
\begin{enumerate}
\item[(G)] (domain-generic) If some old vertex has a point $x_v$ with $d(x_v,\Omega_{\rm out}^c)>\rho_\square:=\frac32r_\square+2s_b+\ell$, then $\square\subset\Omega_{\rm out}$.
\item[(U0)] At \cite[pp.~411--414, 432--433]{Ba3}, with $\Omega_{\rm out}=B^k(\Lambda_2^{(k-1)\prime})$: $m_{\rm use}\ge2r_k-4$, so $\square\subset\Omega_{\rm out}$. The same holds at \cite[p.~582, (2.116)]{Ba2} and for the cap families $\Nbar=169$.
\item[(U1a)] At \cite[p.~579]{Ba2}, with $\Omega_{\rm out}=B^k(\Lambda_0^{(k)})$: $m_{\rm use}>2r_k-M_b-5$, so $\square\subset\Omega_{\rm out}$ under (W-R$'$).
\item[(U1b)] At \cite[Prop.~2.6]{Ba2}, with $\Omega_{\rm out}=B^{k-1}(\Lambda_2^{(k-1)})$: $m_{\rm use}\ge2.5r_k-5$.
\item[(U2)] At \cite[p.~578]{Ba2}, scale $j\le k-1$, with $\Omega_{\rm out}=B^j(\Lambda_2^{(j)})$: $m_{\rm use}\ge5r_j-7$.
\item[(U3)] $\square$ is a non-wrapping box, so $\delta G_k(\square,\eta\Z^d,0)$ is defined.
\end{enumerate}
\end{lemma}
\noindent The support hypothesis for $\tilde B^{(k)}$ at U0 follows from $\supp\tilde B^{(k)}\subset\supp\theta_{k+1}$, which holds because each of the three terms defining $\tilde B^{(k)}$ in \cite[(2.89), p.~575]{Ba2} carries $\theta_{k+1}$ as its left factor. The completeness of the list of uses is statement 15$'$.
\begin{proof}
(G) Since $x_v,c\in Q$, $d(c,x_v)\le\ell$. For $p\in\square$, $d(p,x_v)\le\frac32r_\square+2s_b+\ell<d(x_v,\Omega_{\rm out}^c)$, so $p\notin\Omega_{\rm out}^c$. The hull satisfies $B_{3r/2+s_b}(c)\subset\square\subset B_{3r/2+2s_b}(c)$, since a big block has diameter $\le s_b$ and the hull of a box is a box ((T.div)). Moreover $d(\square_1,\square^c)\ge r$.

The per-use margins use three facts. N1 is the nesting gap (Lemma~\ref{app:G1:Rnest}(v)). N2 is $\Lambda_0^{(i+1)}\subset\Lambda_7^{(i)\prime}$ \cite[p.~566]{Ba2}, iterated. N4 is that $\supp\theta_{k+1}$ lies in the $M_b$-neighbourhood of $B^k(\Lambda_2^{(k)})$ \cite[p.~567]{Ba2}.
\begin{itemize}
\item (U0) The old-vertex bonds have $d\ge2r_k-4$. (The claim is restricted to bonds occurring in vertices, since vertex sums run over $\Omega$.) The condition needed is $r_k/2>\ell+2s_b+4$.
\item (U1a) $g_k$-vertices: $7r_k-10$. $\theta_{k+1}A'$-vertices: $2r_k-4-M_b-1$. The condition needed is $r_k/2\ge\ell+2s_b+M_b+5$, which is (W-R$'$) at $s_b=2M_b$.
\item (U1b) $x_v\in B^k(\Lambda_7^{(k)})\subset\Lambda_7^{(k-1)}$-region, and N1 at scale $k-1$ with $s=5$ gives $(5/L)(r_{k-1}-1)-2\ge2.5r_k-5$.
\item (U2) $d(\supp\theta_{k+1},B^k(\Lambda_0^{(k)c}))>L^{k-j}r_k$ is printed at \cite[p.~578]{Ba2}, and N1 at scale $j$ gives $5r_j-7$.
\item (U3) $t(3r(t)+8M_b+1)\le\eps_0(3r(\eps_0)+8M_b+1)<\ell_0\le\ell_T$, by monotonicity on $t\le\eps_0<\mathrm{e}^{-1}$ and (T.nw).
\end{itemize}
In each case the condition is implied by (W-R$'$).
\end{proof}

\paragraph{The inner domain $\Omega_2$ is not needed: statement 10.}
\begin{lemma}[the inner domain $\Omega_2$ is not needed]\label{app:G1:I3r-W2R}
Let $\ell$ be a $G_k(\Omega_2,\tilde B)$ line of a p.~433 component with ends $v,v'\in\square_1$, and suppose $\square\subset\Omega_{\rm out}$ for the outer domain of the use (Lemma~\ref{app:G1:I3r-Wsq}). Put $X_\ell:=\delta G_k(\Omega,\square,\tilde B)-\delta G_k(\Omega,\Omega_2,\tilde B)$. Then:
\begin{enumerate}
\item[(i)] $G_k(\Omega_2,\tilde B)=G_k(\square,\tilde B)+X_\ell$ as kernels on $\supp h_v\times\supp h_{v'}$; if $\square\subset\Omega_2$, then $X_\ell=\delta G_k(\Omega_2,\square,\tilde B)$, which is p.~433 verbatim;
\item[(ii)] $\norm{hX_\ell h'}_{1,\alpha_0}\le C_X\mathrm{e}^{-\beta_D\dist(\supp h,\supp h')}$, where $C_X$ is built from the constants of Lemma~\ref{app:G1:I3r-DGM}(loc) for the two pairs $(\Omega,\square)$ and $(\Omega,\Omega_2)$;
\item[(iii)] the p.~433 transformation with (i) produces exactly the renormalised classes of p.~433. Every external field produced obeys the two-point $(3.1)/(2.5)$ form, and $\square\subset\Omega_2$ is never used. Hence \cite[(3.2) $\Rightarrow$ (1.33)]{Ba3} holds for these expressions with tree rate $\delta_P=\min\{\delta(\Nbar),\beta_U,\beta_D\}$ and constant $O(1)\max\{1,C_X\}^{N_\ell(\nbar)}$, uniformly in $\eps,\eta,\ell_T$.
\end{enumerate}
\end{lemma}
\noindent The proof composes Lemma~\ref{app:G1:I3r-DGM}(loc) for $(\Omega,\Omega_2)$ and $(\Omega,\square)$, Lemma~\ref{app:G1:I3r-Wsq} at the outer domain of each use, Lemmas~\ref{app:G1:I3r-HOm2} and \ref{app:G1:I3r-HOm2ch3}, the closure of the classes under cuts (Lemma~\ref{app:G1:cutclasses}) and the printed base of p.~433.
\begin{proof}
(i) $[G(\Omega)-G(\square)]-[G(\Omega)-G(\Omega_2)]=G(\Omega_2)-G(\square)$. This is pure algebra and uses no nesting between $\square$ and $\Omega_2$. $G(\Omega,\cdot)$ is evaluated at points admissible for $\Omega_2$; by monotonicity, $\dist(z,\Omega^c)\ge\dist(z,\Omega_2^c)$ when $\Omega_2\subset\Omega$.

(ii) Apply the triangle inequality and Lemma~\ref{app:G1:I3r-DGM}(loc) to each pair. The points of $\square_1$ are $\ge r\ge R_0+13$ from $\square^c$, and the line ends are $\ge R_0+12$ from $\Omega_2^c$ by Lemmas~\ref{app:G1:I3r-HOm2} and \ref{app:G1:I3r-HOm2ch3}; at $k=1$ they are $\ge R_0+13$ from $\Omega_2^c$ by Lemma~\ref{app:K:L18-1}.

(iii) $E(G;\ell)=E(G;\ell\mapsto G_k(\square))+E(G\setminus\ell;X_\ell)$, and $E$ is linear in the external field, so the graphs are those of p.~433. Relabelling the domain of a line leaves every degree \cite[(2.1), p.~422]{Ba3} unchanged. Cutting a line does not lower the degree \cite[p.~432]{Ba3}: if the cut leaves the graph connected, the two legs of $\ell$ and their differentiations become external and contribute $0$ in place of $-\frac12$ and $-1$ in \cite[(2.1)]{Ba3}, so the degree rises by $1+n_\partial(\ell)$, as in Lemma~\ref{app:G1:cutclasses}(a) for vacuum graphs. $X_\ell$ is $\tilde B$-covariant, so the $B_0$ gauge-out preserves its norm. \cite[(3.2), p.~432]{Ba3} requires no smallness of external fields, so the two-point decay with an $O(1)$ constant suffices.
\end{proof}

\paragraph{End points of $\Omega_2$-lines: statements 11, 12.}
\begin{lemma}[end points of $\Omega_2$-lines in {\cite[Chap.~2]{Ba2}}]\label{app:G1:I3r-HOm2}
Every evaluation point of a $G_k(\Omega_2,\cdot)$ or $\delta G_k(\Omega,\Omega_2,\cdot)$ line on the path of Theorem~\ref{thm:T6} is far from $\Omega_2^c$. The points are line ends, bond partners, block points of $Q_k$, H\"older partners and supports of localisation functions. The bounds are:
\begin{itemize}
\item at U1a and U1b, with $\Omega_2=B^k(\Lambda_2^{(k)})$: $\dist(\cdot,\Omega_2^c)\ge5r_k-5-c_1$ with $c_1\le4$;
\item at U2 (scale $j$): $\dist>L^{k-j}r_k-c_1'\ge2R-c_1'$.
\end{itemize}
Hence the margin is $\ge R_0$ under the printed $R>R_0$, and $\ge R_0+12$ under $R\ge\max(R_0,12+c_1')$. The list U1a, U1b, U2 contains all the $G_k(\Omega_2)$-type changes in \cite[Chap.~2]{Ba2}.
\end{lemma}
\noindent The lemma uses the conventions that ``localized in $B^k(\Lambda_7^{(k)})$'' means all vertices and that the localisation width is $O(1)$ (Lemma~\ref{app:G1:Rloc7}), and that $\theta_{k+1}=1$ on the bonds of $B^k(\Lambda_2^{(k)})$ (Lemma~\ref{app:G1:Rthetab}). The quadratic-form use \cite[(2.108)]{Ba2} needs no restriction to depth $R_0$ (\cite[Prop.~2.2]{Ba1} has none).
\begin{proof}
\emph{Nesting.} If $\Lambda_{i+1}^c$ contains every point within $\rho$ of $\Lambda_i^c$ for $m$ consecutive $i$, then $\dist(\Lambda_{i_0+m},\Lambda_{i_0}^c)>m\rho-(m-1)$. (Induction along an $\ell^\infty$ geodesic, which changes the distance by $\le1$ per step.) So $\dist(B^k(\Lambda_7^{(k)}),\Omega_2^c)\ge5r_k-5$.

\emph{U1a and U1b.} The vertex variables lie in $\Lambda_7^{(k)}$ \cite[pp.~579--580]{Ba2}. Evaluation points are offset by $0$ for $(1.6)$, $(1.7)$; by $\le\eta$ for bond legs; by $\le1$ for $Q_k$ contours; and by $\le2$ for H\"older partners and supports of $h$. The (1.16) lines are estimated away before Prop.~2.6, so no interior chain point of a $G_k(\Omega_2)$ line occurs.

\emph{U2.} \cite[p.~578]{Ba2} prints that the distance of $\supp\theta_{k+1}$ to $B^k(\Lambda_0^{(k)c})$, measured on $T_{L^{-j}}$, is $>L^{k-j}r(L^k\eps)>r(L^j\eps)$. The loop vertices are within $c_1'$ of $\supp\theta_{k+1}$.

The list of changes of the domain to $\Omega_2$ is read off the pages of \cite[Chap.~2]{Ba2}. The changes at the first step of \cite[Sect.~A]{Ba2} have margins $5r(\eps)-O(1)$ and $2r(\eps)-O(1)$ and do not affect the result.
\end{proof}

\begin{lemma}[end points of $\Omega_2$-lines in {\cite[Chap.~3]{Ba2}}]\label{app:G1:I3r-HOm2ch3}
\cite[Chap.~3]{Ba2} (pp.~583--594) is on the path. It contains no $G_k(\Omega_2)$ graph line and no p.~433-type cube step; the cubes of Sect.~3C are combinatorial. Its only $\Omega_2=B^k(\Lambda_2^{(k)})$ object is the form $\langle\Lambda_3^{(k)}\phi,\Delta^{(k)}(\Omega_2,\tilde A^{(k+1)})\Lambda_3^{(k)}\phi\rangle$ of (3.5). This form is turned into $\Omega$ by (2.108) run backwards ((3.7)--(3.8)), and its evaluation points satisfy $\dist(\cdot,\Omega_2^c)\ge r(t_k)-2\ge R_0$ under $R\ge R_0+2$.
\end{lemma}
\noindent The margin is measured in Ba{\l}aban's metric. With a Euclidean margin the requirement becomes $R\ge\sqrt3(R_0+2)$, which is implied by $R\ge R^{**}$. The identity $\tilde A^{(k)}=\tilde A^{(k+1)}$ on $\Omega_2$ is not used.
\begin{proof}
On the twelve pages of \cite[Chap.~3]{Ba2} none of the occurrences of ``replac'', ``propagat'', ``Proposition'' is a change of the domain of $G_k$, and ``interaction'' does not occur. In (3.5) the points are $z\in B(x)$, $z'\in B(x')$ with $x,x'\in\Lambda_3^{(k)}$, plus bond partners. By \cite[(2.8)]{Ba2}, every large block of $\Lambda_3^{(k)}$ is at distance $>r(t_k)$ from $(\Lambda_2^{(k)})^c$. Subtracting $1$ for the block and $1$ for the bond gives $r(t_k)-2$.
\end{proof}

\paragraph{Conventions and completeness: statements 15, 15$'$.}
\begin{lemma}[conventions R1--R3]\label{app:G1:I3r-R123}
\begin{enumerate}
\item[(R1)] In \cite[(1.13)--(1.15)]{Ba3} the $A'$-legs enter only as $g_kA'$, with $\supp g_k\subset\Omega_1$.
\item[(R2a)] $\Omega=B^k(\Lambda_2^{(k-1)\prime})$ and $\Omega_1=B^k(\Lambda_7^{(k-1)\prime})$, and $\Lambda_i^{(j)}$ is defined by \cite[(2.7)--(2.8)]{Ba2} with $r(L^j\eps)$.
\item[(R2b)] $\Lambda_{i+1}\subset\Lambda_i$, so $\Omega_1\subset\Omega$.
\item[(R2c)] $\dist(\Lambda_{i+s},\Lambda_i^c)>s(r-1)$.
\item[(R2d)] The metric is the $\ell^\infty$ torus distance \cite[(1.3)]{Ba1}.
\item[(R3)] The contours $\Gamma^{(k)}_{y,x}$ lie inside $B^k(y)$.
\end{enumerate}
\end{lemma}
\noindent The infimum convention for the distance between sets is not printed. Any other convention costs at most one large-block diameter per step, which $R\gg M_bL$ absorbs. The order $\Omega\supset\Omega_1\supset\Omega_2$ holds at U1b (Lemma~\ref{app:G1:Rnest}).
\begin{proof}
(R1) The vertices of (1.5) are built from $T^\eta_{\dots,e'g_kA'+A^{(k)}}$ \cite[(1.4)]{Ba3}, and p.~414 lists ``(1.14) with $n'=0$ and $\tilde B=A^{(k)}$'' among them. An ungated reading contradicts both. It is corroborated by $g_kA'^{(k)}$ in \cite[(2.107)]{Ba2}.
(R2a) is printed at \cite[pp.~411--412]{Ba3} and \cite[p.~566]{Ba2}.
(R2b) and (R2c) are Lemma~\ref{app:G1:Rnest}(i) and (v).
(R2d) \cite[(1.3)]{Ba1} is the only distance defined; \cite{Ba2} uses its symbol \cite[p.~566]{Ba2}, and \cite[p.~411]{Ba3} adopts the notation of \cite{Ba1,Ba2}.
(R3) follows from the staircase contours of \cite[(2.1)]{Ba1}, their concatenation (2.2), and box blocks \cite[(1.17)]{Ba1}. Long arcs are excluded by \cite[(1.2)]{Ba1}.
\end{proof}

\begin{lemma}[completeness of the lists of uses]\label{app:G1:I3r-Rexh}
The uses in \cite{Ba3}, in the proof of Theorem~\ref{thm:T6}, of the following are exactly those treated in Lemmas~\ref{app:G1:I3r-Wsq}--\ref{app:G1:I3r-HOm2} and Corollary~\ref{cor:prop1R4}, plus two further items (O1, O2): $\Omega_2$ or another domain; $G_k(\cdot)$ with a domain argument; $\delta G$; (2.5) and (1.16); cube constructions; and the hypothesis of p.~412 on $\supp\tilde A$. The consumers of the p.~412 hypothesis are p.~433, p.~422 and the domain changes (p.~414 and (2.5), pp.~424, 426, 432).
\begin{itemize}
\item O1 is \cite[(2.8)]{Ba3} on $T_\eta$, with margin $2$, at the stage of the cube $\square$;
\item O2 is the comparison on p.~437 of torus and $\Z^d$ integrals, which is the image kernel of Lemma~\ref{lem:IMG}.
\end{itemize}
\end{lemma}
\noindent The same holds for the uses in \cite{Ba2} (Lemmas~\ref{app:G1:I3r-HOm2}, \ref{app:G1:I3r-HOm2ch3}). The Ward--Takahashi identity of \cite[p.~442]{Ba3} with $\lambda$ linear is exact on $\zeta\Z^d$, and on a torus it holds up to the image term of Lemma~\ref{lem:IMG}.
\begin{proof}
Every occurrence of the six objects on the journal pages 411--445 of \cite{Ba3} is read in context. None is used with a margin outside (W-R$'$), outside $R\ge R_0+12+O(1)$, or outside depth $>r$ at the changes of domain.
\end{proof}

\paragraph{Statements 16 and 17.}
\begin{lemma}[$A\cdot A$ local terms]\label{app:G1:I3r-AA}
Fix $\mu_*>0$, and let $(u,M)=(1,m_2)$ or $(u,M)=(-\eta,m_1)$. Then, uniformly in $\eps$, $k$ and $\ell_T$, the $\sim\!\ast\!\sim$ vertices of \cite[(3.30), (3.32)]{Ba3} satisfy:
\begin{itemize}
\item[(N1)] they enter no covariance and no Gaussian weight;
\item[(N2)] every localised family containing them obeys $(1.33)^\natural$, that is, the bound \cite[(1.33)]{Ba3} times $(\abs uc_k)^m$, where $m$ is the number of coloured lines of the family. The tree rate of \cite[(1.33)]{Ba3} is the rate $\beta_P$ of Proposition~\ref{app:G1:I3}. This rate is at most the rate $\delta(\Nbar)>0$ at the cap (Proposition~\ref{app:G1:VSR-CAP}), and it is not a rate depending on $d$ only. The constant is independent of $\eps,k,\Omega,\Omega_1,\Omega_2,\tilde B,\ell_T$;
\item[(N3)] they are not iterated as a running coupling;
\item[(N4)] they create no new counterterm and no new vacuum entry.
\end{itemize}
\end{lemma}
\noindent The proof is given after Proposition~\ref{app:G1:I3c-AA} below.

\paragraph{The two-photon vertices of \texorpdfstring{\cite[(3.30), (3.32)]{Ba3}}{[Ba3, (3.30), (3.32)]}.}
We work in step-$k$ units. For $j\le k$, $\ell_j:=L^j\eta\le1$ is the scale of a line of index $j$; in particular $\ell_k=1$. We use the following two properties.
\begin{itemize}
\item[(LB)] (Line bounds.) Every line $G$ of index $j$ in the expansion of \cite[p.~414]{Ba3}, scalar or vector, satisfies $\abs{G(x,x')}\le K_\ell\,\ell_j^{-1}\mathrm{e}^{-a\abs{x-x'}/\ell_j}$. Each derivative at an end costs a factor $\ell_j^{-1}$, and an $\alpha$-H\"older quotient of a derivative costs a factor $\ell_j^{-1-\alpha}$; this holds in each variable separately and in both together \cite[(2.10)--(2.12), p.~426; p.~418]{Ba3}. The constant $K_\ell$ is independent of $\eps,k,\ell_T$. For the vector lines it is independent of $\mu_*$. The slice weights of \cite{K1} multiply it by at most $2$ per line, and coloured lines carry the factor $(\abs uc_k)$ per line.
\item[(VB)] (Vertex bounds.) In the notation of Lemma~\ref{app:G1:AA}, the two-photon vertex is
\[
e(t_k)^2\sum_z\eta^3g(z)A_\mu(z)c_{\mu\mu'}(z)g'(z)A'_{\mu'}(z),
\]
with $\abs{c_{\mu\mu'}}\le\Gamma_3\ell_{j_0}^{-1}$, where $j_0$ is the lowest index of the legs of the vertex. The first-moment vertex has $\abs{c_{\mu\mu'\nu}}\le\Gamma_3'$, and the vertex of \cite[(3.32)]{Ba3} has $\abs{c^{(4),\Lambda}}\le\Gamma_4$ for every label configuration $\Lambda$, with $\Gamma_4=\Gamma_4(M)$ finite at the Higgs mass considered (Lemmas~\ref{app:G1:AA}, \ref{app:G1:AA4} and~\ref{lem:B:massrange}(b)).
\end{itemize}
We say that a two-photon vertex has an \emph{unshared lowest-index line} if its line of index $j_0$ is neither the lowest-index line of another two-photon vertex nor the transported leg of a renormalised subgraph of the same class \cite[p.~440]{Ba3}.

\begin{lemma}[lattice sum]\label{app:G1:I3c-latsum}
For $0<\eta\le\ell$ and $b>0$, $\sum_{z\in\eta\Z^3}\eta^3\mathrm{e}^{-b\abs z/\ell}\le S(b)\ell^3$, where $S(b):=(2\sqrt3/b+1)^3$.
\end{lemma}
\begin{proof}
Since $\abs z\ge\abs z_1/\sqrt3$, the sum is at most $\big(\eta\sum_{n\in\Z}\mathrm{e}^{-c\abs n}\big)^3$ with $c=b\eta/(\sqrt3\ell)$. Now $\sum_n\mathrm{e}^{-c\abs n}=\coth(c/2)\le1+2/c$. Hence $\eta\coth(c/2)\le\eta+2\sqrt3\ell/b\le(1+2\sqrt3/b)\ell$.
\end{proof}

\begin{lemma}[insertion bound]\label{app:G1:I3c-INS}
Assume (LB), (VB) and $L=2$, and let a two-photon vertex at $z$ have an unshared lowest-index line.
\begin{enumerate}
\item[(a)] (Both legs on lines.) Let the legs be lines $G_0,G_1$ of indices $j_0\le j_1$, and let $0<\delta'<a$. Then $K(x,y):=\sum_z\eta^3G_0(x,z)c(z)G_1(z,y)$ obeys
\[
\abs{K(x,y)}\le K_\ell^2\Gamma_3S(\delta')\,\ell_{j_0}\,\ell_{j_1}^{-1}\mathrm{e}^{-(a-\delta')\abs{x-y}/\ell_{j_1}}.
\]
This is the bound (LB) of a line of index $j_1$, times $K_\ell\Gamma_3S(\delta')\ell_{j_0}$. A derivative at $y$ costs $\ell_{j_1}^{-1}$, and a contour leg at $x$ gives an extra factor $\ell_{j_0}\le\ell_{j_1}$. Moreover $\sum_{j_0\le j_1}\ell_{j_0}<2\ell_{j_1}$.
\item[(b)] (One external leg.) For an external field $\tilde A$ at $z$ and a line of index $j_0$ to $x$,
\[
\Big|\sum_z\eta^3G_0(x,z)c(z)g(z)\tilde A(z)\Big|\le K_\ell\Gamma_3S(a/2)\ell_{j_0}\norm{\tilde A}_\infty\mathrm{e}^{-\frac a2\dist(x,\supp g)},
\]
and $\sum_{j_0\le k}\ell_{j_0}<2$.
\item[(c)] (Both legs external; then $j_0=k$.) Per unit cube $\square$, $\big|\sum_{x\in\square}\eta^3g\tilde A_\mu c_{\mu\mu'}g'\tilde A'_{\mu'}\big|\le\Gamma_3\norm{h\tilde A}_{1,\alpha_0}\norm{h'\tilde A'}_{1,\alpha_0}$.
\item[(d)] (First-moment vertex.) With a derivative on the $j_1$-leg, the relative factor is $K_\ell\Gamma_3'S\,\ell_{j_0}^2/\ell_{j_1}\le K_\ell\Gamma_3'S\,\ell_{j_0}$. With a derivative on the $j_0$-leg, it is $K_\ell\Gamma_3'S\,\ell_{j_0}$.
\item[(e)] (Chains.) If $N_v$ insertions occur in series on one path, the choice $\delta'=a/(2N_v)$ at each insertion keeps the decay rate $\ge a/2$, at the price of a factor $S(a/(2N_v))^{N_v}$.
\item[(f)] (Transported legs.) Let a renormalised subgraph $H$ of internal scale $\ell_i$, $i<j_0$, transport its leg $G_0$ to order $1+\alpha$ \cite[(3.26), p.~440]{Ba3}. Then the composite of (a), with $H$ contracted to its local vertex at scale $j_0$, is bounded by the bound of (a) times $C(\ell_i/\ell_{j_0})^{\alpha}$, and $\sum_{i<j_0}(\ell_i/\ell_{j_0})^\alpha\le(1-L^{-\alpha})^{-1}$.
\end{enumerate}
\end{lemma}
\begin{proof}
(a) Since $\ell_{j_0}\le\ell_{j_1}$, we have $\mathrm{e}^{-a\abs{x-z}/\ell_{j_0}}\le\mathrm{e}^{-(a-\delta')\abs{x-z}/\ell_{j_1}}\mathrm{e}^{-\delta'\abs{x-z}/\ell_{j_0}}$ and $\mathrm{e}^{-a\abs{z-y}/\ell_{j_1}}\le\mathrm{e}^{-(a-\delta')\abs{z-y}/\ell_{j_1}}$. By the triangle inequality, the dependence on $y$ is bounded by $\mathrm{e}^{-(a-\delta')\abs{x-y}/\ell_{j_1}}$. By Lemma~\ref{app:G1:I3c-latsum}, the $z$-sum is at most $\ell_{j_0}^{-1}\Gamma_3\ell_{j_0}^{-1}S(\delta')\ell_{j_0}^3$. Finally, $\sum_{j_0\le j_1}\ell_{j_0}=\ell_{j_1}\sum_{n=0}^{j_1}2^{-n}<2\ell_{j_1}$.

(b) Since $\ell_{j_0}\le1$, $\mathrm{e}^{-a\abs{x-z}/\ell_{j_0}}\le\mathrm{e}^{-\frac a2\abs{x-z}/\ell_{j_0}}\mathrm{e}^{-\frac a2\abs{x-z}}$. Sum over $z$ as in (a).

(c) Here $\ell_k=1$, the cube has volume $1$, and $\norm{\cdot}_\infty\le\norm{\cdot}_{1,\alpha_0}$.

(d) As in (a), with one more factor $\ell^{-1}$ from the differentiated line.

(e) Apply (a) $N_v$ times.

(f) The remainder factor $\abs{x'-x}^{1+\alpha}$ is absorbed by the lines of $H$, which give $\le C\ell_i^{1+\alpha}$ \cite[p.~438]{Ba3}. By (LB), the quotient of order $1+\alpha$ of $G_0$ costs $\ell_{j_0}^{-1-\alpha}$ and keeps the decay $\mathrm{e}^{-a\abs{x-z}/\ell_{j_0}}$. So in the transported term the factor of $G_0$ relative to (LB) is $(C\ell_i)^{1+\alpha}\ell_{j_0}^{-1-\alpha}$. The subgraph $H$, of degree $-1$, contributes $\ell_i^{-1}$ in place of the local vertex factor $\ell_{j_0}^{-1}$. The product is $C(\ell_i/\ell_{j_0})^{\alpha}$ times the unmoved expression. The $z$-sum of (a) is unchanged, because the quotient keeps the decay of $G_0$, and the sum over $i$ is geometric.
\end{proof}

\begin{remark}[two configurations not covered by Lemma~\ref{app:G1:I3c-INS}]\label{app:G1:I3c-INSrem}
(i) \emph{Shared lines.} Let two two-photon vertices $V_1,V_2$ be joined by a line of index $j_0$ that lies below their other legs. A bound by absolute values gives $e(t_k)^4\Gamma_3^2\ell_{j_0}^{-3}S(a)\ell_{j_0}^3=e(t_k)^4\Gamma_3^2S(a)$ per unit cube for each $j_0$, and the sum over $j_0<k$ is of order $k$. Such families must therefore be bounded as whole classes (Lemma~\ref{app:G1:I3c-fact} and Proposition~\ref{app:G1:I3c-AA}).
(ii) \emph{Transport charged to the composite.} Suppose the transport of Lemma~\ref{app:G1:I3c-INS}(f) is charged to the composite line of index $j_1$ instead of to $G_0$. Then the ratio to the required bound is $L^{n\alpha}\ell_{j_1}/(1-L^{-\alpha})$ with $n=j_1-i$, and this ratio is unbounded at $j_1=k$. This is why Lemma~\ref{app:G1:I3c-INS}(f) charges the gain to $G_0$.
\end{remark}

\begin{lemma}[zeroth moment of a one-particle-reducible chain]\label{app:G1:I3c-fact}
Let $\Pi^1,\Pi^2$ be kernels with two vector legs and $C_\gamma$ a vector line on $\xi\Z^3$, all with absolutely summable products. Then
\[
\sum_{x'}(\Pi^1C_\gamma\Pi^2)(x,x')=\sum_{y,y'}\Pi^1(x,y)C_\gamma(y,y')\,m^{(2)}(y'),\qquad m^{(2)}(y'):=\sum_{x'}\Pi^2(y',x').
\]
Suppose in particular that every scalar line of $\Pi^2$ is the propagator $C^\xi$, and that $\Pi^2$ is the sum of the order-$e^2$ class with two vector legs at fixed labels. Then $m^{(2)}\equiv0$ by \cite[(3.29), p.~441]{Ba3}, that is, by the Ward--Takahashi identity \cite[(2.26)]{Ba3} with $M^2=1$. Hence the one-particle-reducible members of the order-$e^4$ class with two vector legs contribute nothing to the part of its zeroth moment in which all scalar lines are $C^\xi$.
\end{lemma}
\begin{proof}
Use Fubini's theorem, then \cite[(3.29)]{Ba3} at fixed labels.
\end{proof}

\paragraph{The external vector fields of the $A\cdot A$ vertices.}\label{app:G1:ssec-AAext}
Fix a step $k$ with $1\le k\le K-1$, put $t:=t_k$, and work pointwise on the support of the small-field restrictions of \cite[(2.55)]{Ba2} at scales $k$ and $k+1$. As in \cite[(2.80)--(2.89), pp.~574--575]{Ba2}, $A$ is the block field of scale $k$ on $T_1^{(k)}$, $B$ is the block field of scale $k+1$, and
\[
A=A'+aL^{-2}C^{(k)}_{\Lambda_0^{(k)}}Q^*B,\qquad
\theta_{k+1}A'^{(k)}=a_k\theta_{k+1}\zeta^{(k)}G_kQ_k^*A',\qquad
\tilde A:=\theta_{k+1}A'^{(k)}+\tilde B^{(k)}.
\]
Here $(\zeta^{(k)}G_kQ_k^*)(x,y)=\zeta^{(k)}(x,y)(G_kQ_k^*)(x,y)$, and $\tilde B^{(k)}$ is the sum of the second, third and fourth terms of \cite[(2.89)]{Ba2}. Everything else in $A^{(k)}$ is the background $\mathfrak B$. For a vector field $f$ on the $\eta$-lattice put
\[
\norm f_{1,\alpha}:=\sup_x\abs{f(x)}+\sup_{x,\mu}\abs{\partial_\mu f(x)}+[\partial f]_\alpha,\qquad
[\partial f]_\alpha:=\sup_{x\ne x',\mu}\frac{\abs{\partial_\mu f(x')-\partial_\mu f(x)}}{\abs{x-x'}^{\alpha}}.
\]
This is the norm \cite[(1.32)]{Ba3} for external vector fields (there with $\tilde B=0$). We use the following properties of the construction of \cite{Ba1,Ba2}.
\begin{itemize}
\item[(F1)] On $\Lambda_0^{(k)\prime}$, $\abs{\partial B}\le c_1p(t)$ and $\abs B\le c_1p(t)/(\mu_0t)$. For $x\in B(y)$ with $y\in\Lambda_1^{(k)\prime}$, $\abs{A(x)-(Q^*B)(x)}\le c_2p(t)$, and the three remainder terms of \cite[(2.82)]{Ba2} are at most $c_3p(t)$. These are the restrictions and the estimates of \cite[(2.81)--(2.86)]{Ba2}. They involve the values of $A$ and $B$ only through differences, or multiplied by exponentially small kernels.
\item[(F2)] Let $W\in\{a_kG_kQ_k^*,\;a_{k+1}L^{-2}G^\eta_{k+1}Q^*_{k+1}\}$. There are $c_0,\delta_0>0$, independent of $\eps,k,\ell_T$, such that $\abs{W(x,y)}+\abs{\partial_xW(x,y)}\le c_0\mathrm{e}^{-\delta_0\abs{x-y}}$, and, for $\abs{x-x'}\le1$,
\[
\abs{\partial_xW(x',y)-\partial_xW(x,y)}\le c_0\abs{x-x'}^{\alpha_0}\big(\mathrm{e}^{-\delta_0\abs{x-y}}+\mathrm{e}^{-\delta_0\abs{x'-y}}\big)
\]
\cite[Props.~2.1--2.3]{Ba1}, \cite[Prop.~2.2]{Ba2}.
\item[(F3)] $\theta_{k+1}$, $\zeta^{(k)}(\cdot,y)$ and $\zeta^{(k+1)}(\cdot,y)$ take values in $[0,1]$, and their first and second differences in $x$ are bounded by $c_\theta$.
\item[(F4)] If $\dist(x,\supp\theta_{k+1})\le2\eta$ and $\zeta^{(k)}(x,y)\ne0$, then $y\in\Lambda_1^{(k)}$.
\item[(F5)] Each term of $\tilde B^{(k)}$ has the form $\theta_{k+1}\sum_y\chi(\cdot,y)W(\cdot,y)F(y)$, with $\chi\in\{\zeta^{(k)},1-\zeta^{(k)},1-\zeta^{(k+1)}\}$, $W$ as in (F2), and $\mathrm{e}^{-\delta_0d_F/2}M_F\le C_\kappa t^\kappa$ for every $\kappa$. Here $M_F$ and $d_F$ are defined in Lemma~\ref{app:G1:kerHol}. That (F5) holds is shown in Lemma~\ref{app:G1:I3c-Bt}.
\end{itemize}

\begin{lemma}[product rule]\label{app:G1:prodHol}
Let $d=3$ and $0<\alpha<1$, and let $h$ be real with $\abs h\le1$ and with first and second differences bounded by $c_h$. Then $\norm{hf}_{1,\alpha}\le(3+5c_h)\norm f_{1,\alpha}$ for every vector field $f$.
\end{lemma}
\begin{proof}
Write $\norm\cdot$ for sup norms. We have $\partial_\mu(hf)=u\,\partial_\mu f+v\,f$ with $u:=h(\cdot+\eta e_\mu)$ and $v:=\partial_\mu h$. Hence $\norm{hf}\le\norm f$ and $\norm{\partial(hf)}\le\norm{\partial f}+c_h\norm f$.

Let $w$ be a lattice function with $\norm{\partial w}\le c$. Joining $x$ to $x'$ by a monotone lattice path gives $\abs{w(x')-w(x)}\le\min(2\norm w,\sqrt3\,c\abs{x-x'})$. Treating the cases $\abs{x-x'}\le1$ and $\abs{x-x'}>1$ separately, this is $\le\max(2\norm w,\sqrt3\,c)\abs{x-x'}^\alpha$.

Put $g:=\partial_\mu(hf)$ and write
\[
g(x')-g(x)=u(x')\big(\partial f(x')-\partial f(x)\big)+\partial f(x)\big(u(x')-u(x)\big)+v(x')\big(f(x')-f(x)\big)+f(x)\big(v(x')-v(x)\big).
\]
This gives $[g]_\alpha\le[\partial f]_\alpha+(2+\sqrt3c_h)\norm{\partial f}+c_h(2\norm f+\sqrt3\norm{\partial f})+2c_h\norm f$. Adding the three parts of the norm gives $\norm{hf}_{1,\alpha}\le(1+5c_h)\norm f+(3+2\sqrt3c_h)\norm{\partial f}+[\partial f]_\alpha$, which implies the claim.
\end{proof}

\begin{lemma}[cut-off kernels]\label{app:G1:kerHol}
Let $\theta,\chi$ satisfy (F3), let $W$ satisfy (F2), and let $F$ be a function on the unit lattice. Put $f(x):=\theta(x)\sum_y\chi(x,y)W(x,y)F(y)$. Let $Y$ be the set of $y$ such that $\chi(\hat x,y)\ne0$ for some $\hat x$ with $\dist(\hat x,\supp\theta)\le2\eta$. Put $M_F:=\sup_Y\abs F$ and
\[
d_F:=\inf\{\abs{\hat x-y}:\ \dist(\hat x,\supp\theta)\le2\eta,\ \chi(\hat x,y)\ne0\}.
\]
Then $\norm f_{1,\alpha_0}\le C_K\,\mathrm{e}^{-\delta_0d_F/2}M_F$, where $C_K$ depends only on $c_0,\delta_0,c_\theta$.
\end{lemma}
\begin{proof}
Let $K(x,y):=\theta(x)\chi(x,y)W(x,y)$. By the Leibniz rule,
\[
\partial_xK=\theta(\cdot+\eta e)\chi(\cdot+\eta e,y)\partial W+\theta(\cdot+\eta e)(\partial\chi)W+(\partial\theta)\chi W,
\]
so $\abs{\partial_xK(x,y)}\le(1+2c_\theta)c_0\mathrm{e}^{-\delta_0\abs{x-y}}$. In each term, every factor is of one of two kinds. Either it is bounded and Lipschitz in $x$, with constants depending on $c_\theta,c_0,\delta_0$ (for $W$ this follows from (F2), along a path of length $\le\sqrt3\abs{x-x'}$). Or it is $\alpha_0$-H\"older in $x$ for $\abs{x-x'}\le1$; this is the factor $\partial W$. Hence, for $\abs{x-x'}\le1$,
\[
\abs{\partial_xK(x',y)-\partial_xK(x,y)}\le C\abs{x-x'}^{\alpha_0}\big(\mathrm{e}^{-\delta_0\abs{x-y}}+\mathrm{e}^{-\delta_0\abs{x'-y}}\big).
\]
For $\abs{x-x'}>1$ we bound the difference by twice the supremum of $\partial f$.

Only $y\in Y$ contribute. For these, $\min(\abs{x-y},\abs{x'-y})\ge d_F-2\eta$ whenever the summand is non-zero. So $\mathrm{e}^{-\delta_0\abs{x-y}}\le\mathrm{e}^{3\delta_0}\mathrm{e}^{-\delta_0d_F/2}\mathrm{e}^{-\delta_0\abs{x-y}/2}$, and similarly at $x'$. Summing over $y$, with $S_0:=\sup_x\sum_y\mathrm{e}^{-\delta_0\abs{x-y}/2}<\infty$, gives the claim.
\end{proof}

\begin{lemma}[the translation removes the constant mode]\label{app:G1:Aprime-const}
Assume (F1). Then, for $x\in B(y)$ and $y\in\Lambda_1^{(k)\prime}$,
\[
\abs{A'(x)}\le C_1p(t),\qquad C_1:=c_2+c_3+a^{-1}L^2c_1\,\mu_0\eps_0 .
\]
The constant $C_1$ does not depend on the values of $A$ and $B$; in particular it does not depend on their constant parts.
\end{lemma}
\begin{proof}
By \cite[(2.82)]{Ba2}, $aL^{-2}(C^{(k)}_{\Lambda_0^{(k)}}Q^*B)(x)=aL^{-2}(C^{(k)}Q^*1)(x)\,B(y)+\mathcal R(x)$, where $\abs{\mathcal R(x)}\le c_3p(t)$ by (F1). By \cite[(2.83)--(2.85)]{Ba2}, $aL^{-2}C^{(k)}Q^*1=(1+X)^{-1}$ with
\[
0\le X=a^{-1}L^2a_k\mu_0^2t^2/(a_k+\mu_0^2t^2)\le a^{-1}L^2\mu_0^2t^2,
\]
so $\abs{aL^{-2}C^{(k)}Q^*1-1}\le a^{-1}L^2\mu_0^2t^2$. Since $(Q^*B)(x)=B(y)$,
\[
A'(x)=\big(A(x)-(Q^*B)(x)\big)-\big(aL^{-2}C^{(k)}Q^*1-1\big)B(y)-\mathcal R(x).
\]
The bound $\abs B\le c_1p(t)/(\mu_0t)$ bounds the middle term by $a^{-1}L^2c_1\mu_0t\,p(t)\le a^{-1}L^2c_1\mu_0\eps_0\,p(t)$.
\end{proof}

\begin{proposition}[size of the external vector fields]\label{app:G1:AAext}
Assume (F1)--(F5), and let $\eps_0$ be so small that $C_\kappa t^\kappa\le p(t)$ for $t\le\eps_0$, at one fixed $\kappa$. Then, for every localisation function $h'$ of \cite[p.~420]{Ba3} whose first and second differences are bounded by $c_h$,
\[
\norm{h'\tilde A}_{1,\alpha_0}\le(3+5c_h)\big(C_KC_1+3C_K\big)\,p(t).
\]
The constant depends only on $d,a,L,c_0,\delta_0,c_\theta,c_1,c_2,c_3,\mu_0\eps_0,c_h$. It does not depend on $\eps$, $k$, $\ell_T$, the domains $\Omega,\Omega_1,\Omega_2$, or the constant parts of $A$ and $B$. In particular, the hypothesis of \cite[p.~412]{Ba3} on the external vector fields holds uniformly in the photon zero mode and in the twist on the torus.
\end{proposition}
\begin{proof}
For $\theta_{k+1}A'^{(k)}$, apply Lemma~\ref{app:G1:kerHol} with $\chi=\zeta^{(k)}$, $W=a_kG_kQ_k^*$ and $F=A'$. By (F4), $Y\subset\Lambda_1^{(k)}$, so $M_F\le C_1p(t)$ by Lemma~\ref{app:G1:Aprime-const}; also $d_F\ge0$. Hence $\norm{\theta_{k+1}A'^{(k)}}_{1,\alpha_0}\le C_KC_1p(t)$. For each of the three terms of $\tilde B^{(k)}$, Lemma~\ref{app:G1:kerHol} and (F5) give $\norm{\cdot}_{1,\alpha_0}\le C_KC_\kappa t^\kappa\le C_Kp(t)$. Lemma~\ref{app:G1:prodHol} with $h=h'$ concludes.

The bound involves only three kinds of input: the vector operators $G_k,Q_k,C^{(k)}_{\Lambda_0^{(k)}},G_{k+1},Q_{k+1}$, which do not depend on any background field; the fixed cut-offs; and the bounds (F1), which see the values of $A$ and $B$ only through $\abs B\le c_1p/(\mu_0t)$. So the bound is the same for every constant part.
\end{proof}

\begin{remark}[where the constant mode goes]\label{app:G1:AAext-rem}
(i) By Lemma~\ref{app:G1:Aprime-const}, the constant part of the block fields, including the constant photon mode of the torus, is carried by the background $\mathfrak B$; only $O(\mu_0^2t^2)\abs B=O(\mu_0t\,p)$ of it reaches $A'$. Hence no vertex of \cite[(3.30), (3.32)]{Ba3} carries the torus constant on an external leg. For the background, \cite[p.~412]{Ba3} asks only for bounds on $\partial^\eta\mathfrak B$ and on its H\"older quotient, and these do not change under $\mathfrak B\mapsto\mathfrak B+c$ with $c$ constant. A constant $c'$ with $\abs{c'}\le\norm{\tilde A}_\infty$ can be moved from $\tilde A$ to the background by passing to $(\mathfrak B+\theta_{k+1}c',\ \tilde A-\theta_{k+1}c')$. This keeps $\supp\tilde A\subset\supp\theta_{k+1}$, and by (F3) it changes all the bounds above by at most a factor depending on $c_\theta$.
(ii) Proposition~\ref{app:G1:AAext} does not use the statement that the constant $O(1)$ of \cite[Prop.~1]{Ba3} is uniform under $\mathfrak B\mapsto\mathfrak B+c$ on $T_\eta$. That statement is Proposition~\ref{app:K:L10-2}.
(iii) At $L=3$, $\tilde B^{(k)}$ has a fourth term of the same form, (F5) holds for it, and the proof is unchanged. At $k=0$ the same argument applies to \cite[(2.19)--(2.22)]{Ba2}.
\end{remark}

\begin{proposition}[the two-photon vertices]\label{app:G1:I3c-AA}
Fix $\mu_*>0$ and $(u,M)\in\{(1,m_2),(-\eta,m_1)\}$. Assume (LB) and (VB). Assume that the family bounds of Definition~\ref{app:G1:VSR-IH} and Proposition~\ref{app:G1:VSR-CAP} hold for families of graded order $\le\Nbar$. Assume also the following two hypotheses on the frame of the expansion on $T_\eps$.
\begin{itemize}
\item[(ZM)] The constant mode $c$ of the photon field is carried as a constant background in $\tilde B$, equivalently as a twist $\vartheta=ec\ell_T$ of the scalar boundary condition. Every bound on scalar operators used in \cite[Prop.~1]{Ba3} is uniform in $\vartheta$, provided the localisation cubes do not wrap the torus.
\item[(EXT)] The external vector fields of the families satisfy $\abs{\tilde A},\norm{h'\tilde A}_{1,\alpha_0}\le O(1)p(t_k)$, with the constant mode removed as in (ZM). This is the hypothesis of \cite[p.~412]{Ba3}.
\end{itemize}
Then the vertices of \cite[(3.30), (3.32)]{Ba3} satisfy (N1)--(N4) of Lemma~\ref{app:G1:I3r-AA}, uniformly in $\eps,k,\ell_T$.
\end{proposition}
\begin{proof}
\emph{(N1).} The covariances of the expansion are fixed before any graph is formed: ``$A'_j$ are independent Gaussian fields with covariances $C^{(j),L^j\eta}$ \dots\ the field $A'$ has the covariance $G_k$. The basic propagator for the scalar field $\phi'$ is $G_k(\Omega,\tilde B)$'' \cite[p.~414]{Ba3}. Moreover, ``all the $A'$-legs are contracted, i.e.\ they are divided into pairs and each pair is replaced by the corresponding propagator'' (ibid.). The vertex list \cite[(1.6)--(1.15), pp.~413--414]{Ba3} contains no two-photon term with a coefficient proportional to $\ell_{j_0}^{-1}$. The vertices of (3.30) and (3.32) arise only on p.~442, where the values of subgraph classes are rewritten inside the estimate. So they enter no covariance and no Gaussian weight.

\emph{(N2).} Every renormalised family is bounded by absolute values \cite[p.~426]{Ba3}, using the counting (2.13)--(2.16). By \cite[Prop.~2.2, p.~428]{Ba3}, this counting needs only lines with representations (2.6) and bounds (2.10)--(2.12) (``lines with the same exponential factors but with arbitrary dimensions''), and vertices whose coefficients are bounded times powers of the scale. The two-photon vertex is the only vertex whose coefficient carries a negative power $\ell_{j_0}^{-1}$ of a line scale. There are two cases.
\begin{itemize}
\item \emph{The lowest-index line is unshared.} Lemma~\ref{app:G1:I3c-INS}(a),(d),(e),(f) replaces the vertex, together with that line, by a generalised line with an extra positive power $\ell_{j_0}$, which is summable over $j_0$. If one or both legs are external, Lemma~\ref{app:G1:I3c-INS}(b),(c) replaces the vertex by a bounded smeared external leg, respectively by a bounded local vertex. These steps use $\norm{\tilde A}_\infty\le\norm{h'\tilde A}_{1,\alpha_0}\le O(1)p$, which is (EXT). By (ZM), the constant mode is not a vector leg; it enters only through scalar operators, uniformly in $\vartheta$. So the family obeys $(1.33)^\natural$ with $O(1)$ replaced by $O(1)(1+C_{\rm INS})^{N_v}$, where
\[
C_{\rm INS}:=2K_\ell^2\max(\Gamma_3,\Gamma_3',\Gamma_4)(1-2^{-\alpha})^{-1}S(a/(2N_v))^{N_v}
\]
and $N_v\le N(\Nbar)$.
\item \emph{The lowest-index line is shared} (Remark~\ref{app:G1:I3c-INSrem}(i)). Then the family is a one-particle-reducible or nested member of the class of order $e^4$ (or $e^2\lambda$) with two external vector legs, and its lowest external index $j_{\rm ext}$ exceeds the index of the shared line. Such members are treated in \cite{Ba3} as whole graphs: ``we resum over proper orderings and $j$-indices, we replace the propagators $G_{j_0}(0)$, $G_{j_0}$ by $C^\xi$ \dots\ we apply a proper Ward--Takahashi identity'' \cite[p.~442]{Ba3}. Only the last divergent subgraph of an ordering is decomposed, and the remaining divergent graphs are the whole graphs, to which (3.31) applies \cite[p.~445]{Ba3}. So the index of the shared line is summed inside the resummed propagators at index $j_{\rm ext}$, and no sum over $j_0$ of terms of size $O(1)$ arises. By (3.31), the family is the sum of a generalised graph of degree $\alpha$ and a zeroth-moment term. The coefficient of the zeroth-moment term is bounded by Lemma~\ref{app:G1:AA4}, and its pure-$C^\xi$ part vanishes; for the one-particle-reducible part this is already Lemma~\ref{app:G1:I3c-fact}. The inner transformations are exact identities on the set of orderings in which the inner subgraph is the last divergent one, so the whole-class sums are unchanged.
\end{itemize}
In both cases the bound is uniform in $\eps,k,\ell_T$ and finite at both points $(u,M)$. The only dependence on $u$ is the factor $(\abs uc_k)^m$ of the coloured lines. The only dependence on $M$ is through $m^2_{j_0}=M^2t_{j_0}^2\le1$ in Lemmas~\ref{app:G1:AA} and~\ref{app:G1:AA4}, and, in Lemma~\ref{app:G1:AA4}, through the bound $\bar c_{\rm fin}$ on the fixed mass vertices, which is finite at every $M\in(0,m_2]$ by Lemma~\ref{lem:B:massrange}(a).

\emph{(N3).} The interaction after $k$ steps is given by formula (1.5) of \cite[p.~412]{Ba3}, built from the original vertices (1.6)--(1.15). By (N1) the two-photon vertices are not among them. So they are created anew inside the estimate of each step and are never carried to step $k+1$ as a coupling. Hence no factor $L^2$ per step can accumulate, and their weight per step is bounded through (N2).

\emph{(N4).} (3.30)--(3.32) are identities between values of classes \cite[p.~442]{Ba3}, so the total of a family is unchanged. The mass counterterm is defined from complete scalar classes (Lemma~\ref{app:G1:ctid}), and inserting these vertices into a scalar two-leg family changes neither the family nor its zeroth moment. The photon has no mass counterterm. Contracting the two legs of such a vertex with each other gives a vacuum graph, which belongs to the vacuum classes of Proposition~\ref{app:G1:I9}. The torus counterterm is built with the torus propagator \cite[(1.28)]{Ba3}. Its difference from the members on $\eta\Z^3$ after \cite[p.~433]{Ba3} is an image term of order $\mathrm{e}^{-\delta\ell_T/t}$. This term enters as a finite vertex of type \cite[(1.7)]{Ba3}, not as a new counterterm.
\end{proof}

\begin{proof}[Proof of Lemma~\ref{app:G1:I3r-AA}]
We apply Proposition~\ref{app:G1:I3c-AA} and verify its hypotheses.
\begin{itemize}
\item (LB) is \cite[(2.10)--(2.12), p.~426]{Ba3}, together with Theorem~\ref{app:G1:VSR-VS} for the vector lines.
\item (VB) is Lemmas~\ref{app:G1:AA} and~\ref{app:G1:AA4}.
\item The family bounds hold at step $k$ by Definition~\ref{app:G1:VSR-IH} and Proposition~\ref{app:G1:VSR-CAP}, through the bound (1.33) of Proposition~\ref{app:G1:I3}; that bound does not use the present lemma.
\item (EXT) follows from the proof of Proposition~\ref{app:G1:AAext}, whose hypothesis (F5) holds by Lemma~\ref{app:G1:I3c-Bt}. Its smallness condition on $\eps_0$ is not needed here. Fix $\kappa=1$. Since $t\le\eps_0<1$ and $p(t)\ge b_0$, the three terms of $\tilde B^{(k)}$ are bounded by $C_KC_\kappa t^\kappa\le C_K(C_\kappa/b_0)\,p(t)$. Hence $\norm{\tilde A}_{1,\alpha_0}\le C_K(C_1+3C_\kappa/b_0)\,p(t)$ and $\norm{h'\tilde A}_{1,\alpha_0}\le(3+5c_h)C_K(C_1+3C_\kappa/b_0)\,p(t)$, which give both bounds of (EXT).
\item In (ZM), the first part (the constant mode is carried by the background) is Lemma~\ref{app:G1:Aprime-const} and Remark~\ref{app:G1:AAext-rem}(i). The localisation cubes do not wrap the torus, by (T.nw) and Lemmas~\ref{app:G1:nowrap} and~\ref{app:G1:gaugetriv}. The uniformity in $\vartheta$ is Proposition~\ref{app:K:L10-2}.
\end{itemize}
\end{proof}

\begin{lemma}[{the tails of \cite[p.~579]{Ba2}}]\label{app:G1:I3r-T579}
Pointwise on the small-field support, at step $k$,
\[
(\text{basic quadratic form}+\mathcal P_w^{(k)})=(\text{polynomial in }\theta_{k+1}A'^{(k)},\psi,\phi\text{ of order}\le\nbar)+\Delta_k^{\rm pre}+\Delta_k^R+\Delta_k^{\tilde B},
\]
with the following two-sided bounds:
\begin{itemize}
\item $\abs{\Delta^{\rm pre}_k}\le Ct_k^{7/2}\,\mathrm{polylog}\,\abs{\Lambda_2^{(k)}}_k$ (the terms of order $>\nbar$, cut before integration);
\item $\abs{\Delta^R_k}\le Ct_k^{13/2}\,\mathrm{polylog}\,\abs{\Lambda_2^{(k)}}_k$ (the terms with Taylor remainders $R_{\nbar+1}$);
\item $\abs{\Delta^{\tilde B}_k}\le Ct_k^{9/2}\,\mathrm{polylog}\,\abs{\Lambda_2^{(k)}}_k$ at $\kappa=31$ (the terms with $\tilde B^{(k)}$, \cite[(2.90)]{Ba2}).
\end{itemize}
Each of these is $\le C't_k^{16/5}N_k$. The constants do not depend on $\eps$, $k$ or $\ell_T$. They are finite at $u=-\eta$ and at $M=m_1$. Where $\delta m^2_{\rm fin}$ vertices occur, they are evaluated at the torus floor $\ell_0$. The expansion of the logarithm of the determinant \cite[(3.46)--(3.47)]{Ba1} is not part of this cut; it is cut at \cite[pp.~577--578]{Ba2}.
\end{lemma}
\noindent The inputs are those of Proposition~\ref{app:G1:T579}:
\begin{itemize}
\item the kernels of \cite{Ba4} and Lemma~\ref{app:G1:WK-K};
\item Lemma~\ref{app:G1:I3r-U1m}(a), and Lemma~\ref{app:G1:I3r-U1mOmega} on Neumann domains under (R.1) and (R.8);
\item Lemma~\ref{app:G1:WK-W};
\item the regularity of the background under the conditions of Definition~\ref{app:G1:rowlist};
\item the cap $\Nbar=169$, with (R.cap) and (T.cap);
\item the per-order bound in the form of Lemma~\ref{app:G1:T579-po}, and Lemma~\ref{app:G1:I3c-PO} below;
\item the conditions (Reg1) and (Reg2) of Lemma~\ref{app:G1:T579-reg};
\item the bound of Statement~\ref{app:K:L06-4};
\item (E.10), and $R\ge R^{**}$.
\end{itemize}
Proposition~\ref{app:G1:I5} uses Proposition~\ref{app:G1:I3} only through item 7 of Table~\ref{app:G1:I3tab}, and item 7 does not use item 17, so there is no cycle.

\paragraph{The remainder and \texorpdfstring{$\tilde B$}{B-tilde} parts.}
The setting, the fields $\mathfrak B$ and $\tilde A$, the bound $\sigma\le C_\sigma t^{1/2}p$, the number $N_k$ and the inputs (K1)--(K4), (REG), ($\tilde B$), (IH), (CAP) are those of Proposition~\ref{app:G1:T579}. The expansion (E1)--(E3) is that of its proof; here $\nbar=13$ and $\Nbar=169$. We write $t=t_k$ and $p=p(t)=b_0(1+\ln t^{-1})^{p_0}$. ``Polylog'' means a factor bounded by $C(1+\ln t^{-1})^{c}$, with constants independent of $\eps,k,\ell_T$. The order of a term is $\alpha+\beta$ for a factor $e^\alpha\lambda^\beta$, and each factor $e_k\tilde A$ counts $1$. On $\supp\chi_k$, scalar external fields obey $\norm{h\Phi_{\rm ext}}_{1,\alpha_0}\le Cp\lambda(t)^{-1/4}$, and vector ones obey $\norm{h'A_{\rm ext}}_{1,\alpha_0}\le Cp$ (Proposition~\ref{app:G1:AAext}). The mass vertices satisfy
\[
\text{(}\delta m\text{)}\qquad\abs{\delta m^2_{K,k}}t^2\le O(1)\big[e(t)^2+\lambda(t)+(e(t)^4+e(t)^2\lambda(t)+\lambda(t)^2)(1+K-k)\big],
\]
with $K-k\le\log_2(\eps_0/t)$, $e(t)=et^{1/2}$ and $\lambda(t)=\lambda t$.

\begin{lemma}[per-order power count]\label{app:G1:I3c-PO}
Consider a connected family with at least two vertices. After its external legs are bounded as above, every elementary vertex of \cite[(1.6)--(1.15), (1.7), (1.18)]{Ba3} of order $q_v$ contributes at most $Ct^{q_v/4}$ times a polylog. Hence a family of total order $q$ is $\le C(q)t^{q/4}p^{4q}\mathrm{e}^{-\delta(\Nbar)d}$.
\end{lemma}
\begin{proof}
Every vertex has at least one internal leg. We go through the vertices.
\begin{itemize}
\item $\lambda\abs\phi^4$ (order $1$) has at most three external scalar legs, which gives $\lambda t\cdot t^{-3/4}=t^{1/4}$.
\item $e^2A^2\abs\phi^2$ (order $2$) with two external scalar legs gives $e^2t\cdot t^{-1/2}=t^{1/2}$.
\item $e\bar\phi A\cdot D\phi$ (order $1$) gives $O(t^{(4-d)/4}p^2)=O(t^{1/4}p^2)$ by \cite[p.~421]{Ba3}, via \cite[Lemma~2.4]{Ba2}.
\item For the mass vertex with two external scalar legs, ($\delta m$) gives $t^{1/2}$ for the $\lambda$-part (order $1$), $t^{1/2}$ for the $e^2$-part (order $2$), and $t^{3/2}(1+K-k)$, that is $t^{3/2}$ times a polylog, for the parts $e^4$, $e^2\lambda$, $\lambda^2$ (orders $4$, $3$, $2$). Since $t\le1$, each of these is at most $t^{q_v/4}$ times a polylog.
\item The vertices (1.18) with $n+n'\ge2$ give $t^{(n+n')/2-1/2}\le t^{(n+n')/4}$.
\item The remaining vertices are better \cite[p.~421]{Ba3}.
\end{itemize}
Internal lines and integrated fluctuations contribute to the $O(1)$ of $(1.33)^\natural$, by (IH) and (CAP). Multiplying over the vertices gives the claim.
\end{proof}

\begin{lemma}[Taylor-remainder factors]\label{app:G1:I3c-Rfac}
Let $R_{n+1}(z):=\mathrm{e}^{iz}-\sum_{r\le n}(iz)^r/r!$, and let $\tilde A$ be real. Then $\abs{R_{n+1}(z)}\le\abs z^{n+1}/(n+1)!$, $\abs{R_{n+1}'(z)}\le\abs z^n/n!$ and $\abs{R_{n+1}''(z)}\le\abs z^{n-1}/(n-1)!$. Assume that the cut-offs $\theta_{k+1},\zeta$ have bounded second differences and that $[\partial^\eta\tilde A]_\alpha\le Cp$ (the H\"older clause of (K1); see Proposition~\ref{app:G1:AAext}). Consider each factor $R_{\nbar+1}(z)$ with $z=qe_k\tilde A(\Gamma)$ for a contour of length $O(1)$, and each bond factor $\eta^{-1}R_{\nbar+1}(\eta qe_k\tilde A(b))$. As a coefficient or as an external field, each such factor satisfies
\[
\norm{\cdot}_{1,\alpha_0}\le c_R(Ct^{1/2}p)^{\nbar+1},\qquad c_R=c_R(\nbar,d,L).
\]
\end{lemma}
\begin{proof}
We have $R_{n+1}'=iR_n$ and $R_{n+1}''=-R_{n-1}$, and $\abs{R_m(z)}\le\abs z^m/m!$ for real $z$ by Taylor's formula with integral remainder. On $\supp\chi_k$, $\abs z\le Ce_k\norm{\tilde A}_\infty\le C'_\sigma t^{1/2}p$ and $\abs{\partial^\eta z}\le Ce_k\norm{\partial^\eta\tilde A}_\infty\le C'_\sigma t^{1/2}p$. The lattice difference of $R(z(\cdot))$ is $\int_0^1R'(z_s)\,ds\cdot\partial^\eta z$, where $z_s$ interpolates between the two values. Its H\"older difference over $\abs{x-\tilde x}\le1$ is bounded by
\[
\sup\abs{R''}\,\abs{z(\tilde x)-z(x)}\,\abs{\partial^\eta z}+\sup\abs{R'}\,\abs{\partial^\eta z(\tilde x)-\partial^\eta z(x)}\le\big[\abs z^{\nbar-1}(Ct^{1/2}p)^2+\abs z^{\nbar}e_k[\partial^\eta\tilde A]_\alpha\big]\abs{x-\tilde x}^\alpha ,
\]
where $\abs{x-\tilde x}\le\abs{x-\tilde x}^\alpha$ was used. This is $\le(Ct^{1/2}p)^{\nbar+1}\abs{x-\tilde x}^\alpha$. For the bond factor, each estimate carries an extra factor $\eta^{\nbar}\le1$.
\end{proof}

\begin{lemma}[$\tilde B^{(k)}$ in the norm (1.32)]\label{app:G1:I3c-Bt}
By \cite[(2.89), p.~575]{Ba2}, $\tilde B^{(k)}=\sum_{i=2}^4\theta_{k+1}\zeta_iu_i$, with
\begin{align*}
\zeta_2u_2&=a_kaL^{-2}\zeta^{(k)}G_kQ_k^*\delta C^{(k)}_{\Lambda_0^{(k)}}Q^*B, &
\zeta_3u_3&=-a_kaL^{-2}(1-\zeta^{(k)})G_kQ_k^*C^{(k)}Q^*\Lambda_0^{(k)\prime}B,\\
\zeta_4u_4&=a_{k+1}L^{-2}(1-\zeta^{(k+1)})G^\eta_{k+1}Q^*_{k+1}\Lambda_0^{(k)\prime}B .
\end{align*}
Assume that $\theta_{k+1},\zeta^{(k)},\zeta^{(k+1)}$ take values in $[0,1]$ and have first and second lattice differences $\le c_\theta$. Assume further that, on the support of $\theta_{k+1}\zeta_i$, for every $\kappa$,
\[
(\tilde B_0)\qquad\abs{u_i}+\abs{\partial^\eta u_i}+[\partial^\eta u_i]_{\alpha_0}\le C_\kappa t^\kappa .
\]
Then $\norm{h\tilde B^{(k)}}_{1,\alpha_0}\le C(c_\theta)C_\kappa t^\kappa$ for every $\kappa$.

Hypothesis $(\tilde B_0)$ holds, and so does (F5) of Subsection~\ref{app:G1:part-TI3}. The sup and derivative parts of $(\tilde B_0)$ are \cite[(2.90)]{Ba2} (``$\tilde B^{(k)}(x)=O((L^k\eps)^\kappa)$ for every $\kappa$, and similar estimates for the derivative''). For the remaining part: each $u_i$ is a kernel of $G_kQ_k^*$ (respectively $G^\eta_{k+1}Q^*_{k+1}$) applied to a source of size $\le c_1p/(\mu_0t)$ supported at distance $\ge r(t)/2-2M_b$ from $\supp(\theta_{k+1}\zeta_i)$. The decay, derivative and H\"older clauses of (K1) carry the same factor $\mathrm{e}^{-\delta_0(r(t)/2-2M_b)}$. Since $r(t)=R(1+\ln t^{-1})^2$, this factor is $\le C_\kappa t^{\kappa+1}$ for every $\kappa$.
\end{lemma}
\begin{proof}
$\partial^\eta(\theta\zeta u)$ is the sum of the three terms $(\partial^\eta\theta)\zeta u$, $\theta(\partial^\eta\zeta)u$ and $\theta\zeta\,\partial^\eta u$, with shifted arguments. Their H\"older quotients over $\abs{x-\tilde x}\le1$ are bounded using $(\tilde B_0)$ and the Lipschitz bounds on $\partial^\eta\theta$ and $\partial^\eta\zeta$, which follow from the bound $c_\theta$ on the second differences.
\end{proof}

\begin{proposition}[the Taylor-remainder cut]\label{app:G1:I3c-R}
In the expansion (E1)--(E3), let $\Delta^R_k$ be the sum of the terms that contain at least one factor $R_{\nbar+1}$: in the bracket $V_k$, in $Q_k(\mathfrak B+\tilde A)$, or in a transport of a vertex of $\mathcal P^{(k)}_w$. Then, pointwise on $\supp\chi_k$ and for $0\le k\le K-1$,
\[
\abs{\Delta^R_k}\le C_Rt_k^{13/2}\,\mathrm{polylog}\,\abs{\Lambda_2^{(k)}}_k\le C_R't_k^{16/5}N_k .
\]
The constant $C_R$ depends on $\nbar,\Nbar,d,L,a,M_b,e,\lambda,\mu_*,b_0,p_0,c_1,R,\delta_0,\eta$, and not on $\eps$, $k$ or $\ell_T$. It is finite at $u=-\eta$ and at $M=m_1$, with the same dependence on $u$ and $M$ as in Proposition~\ref{app:G1:T579}.
\end{proposition}
\begin{proof}
(1) Each $R$-factor obeys Lemma~\ref{app:G1:I3c-Rfac}.

(2) \emph{An $R$-factor in $V_k$.} By the lattice Leibniz rule, the piece $V^R$ has the right-pushed form $R_0+\sum_\mu R_{1,\mu}D_\mu$ of Lemma~\ref{app:G1:I3r-V}(ii), with coefficients $\le c_V(Ct^{1/2}p)^{\nbar+1}$. Consider a chain with one $V^R$-vertex and further vertices of orders $r_i$. As in Step~2 of the proof of Proposition~\ref{app:G1:T579}, it is bounded by a product of per-vertex factors, namely by $C(Ct^{1/2}p)^{\nbar+1+\sum r_i}\norm f_\infty\norm{f'}_\infty\mathrm{e}^{-\beta\dist}$. With end fields $\le(Cp\lambda^{-1/4}t^{-1/4})^2$, this is $\le Ct^{(\nbar+1)/2-1/2}=Ct^{13/2}$ times a polylog. Inside a line of a family of $\mathcal P^{(k)}_w$, the same chain is an external two-variable field of \cite[(2.5)]{Ba3}. It is bounded by (K3) (Lemmas~\ref{app:G1:I3r-U1m} and~\ref{app:G1:I3r-U1mOmega}) with one coefficient replaced, and the family is then bounded by (IH) and (CAP).

(3) \emph{An $R$-factor in a vertex transport.} $(1.33)^\natural$ is linear in the norm of each external field. So the family receives the factor $c_R(Ct^{1/2}p)^{\nbar+1}$, that is $t^7$ times a polylog. At most two scalar legs sit at that vertex, and they give $t^{-1/2}$. By Lemma~\ref{app:G1:I3c-PO}, each of the other vertices gives $\ge t^{q_v/4}$ times a polylog. The total is $\le Ct^{13/2}$ times a polylog times $\mathrm{e}^{-\delta(\Nbar)d}$.

(4) There are finitely many terms, each with at most $D(\nbar)$ factors. The localisations through a fixed block of $\Lambda_2^{(k)}$ sum to $\le S(\Nbar,d)$, and $\abs{\Lambda_2^{(k)}}_k\le N_k$. Finally, $\sup_{t\in(0,1]}t^{13/2-16/5}\,\mathrm{polylog}<\infty$. The bound is two-sided, because the expansion is an exact identity and every term is bounded in absolute value.
\end{proof}

\begin{proposition}[the $\tilde B$ cut]\label{app:G1:I3c-B}
In every polynomial term of order $\le\nbar$ of the expansion (E1)--(E3), split $\tilde A=\theta_{k+1}A'^{(k)}+\tilde B^{(k)}$ multilinearly. Terms of higher order are part of $\Delta^{\rm pre}_k$ in Proposition~\ref{app:G1:T579}, with $\tilde B^{(k)}$ included in $\tilde A$. Let $\Delta^{\tilde B}_k$ be the sum of the terms with at least one factor $\tilde B^{(k)}$. Let $s(\nbar)$ be the largest number of scalar legs of a term of order $\le\nbar$. Every elementary vertex of \cite[(1.6)--(1.15)]{Ba3} has order $\ge1$ and at most four scalar legs, so $s(\nbar)\le8\nbar$. Put $\kappa:=\lceil s(\nbar)/4\rceil+5$; this is $31$ when $s(\nbar)=8\nbar$. Then, pointwise on $\supp\chi_k$,
\[
\abs{\Delta^{\tilde B}_k}\le C_{\tilde B}t_k^{9/2}\,\mathrm{polylog}\,\abs{\Lambda_2^{(k)}}_k\le C_{\tilde B}'t_k^{16/5}N_k .
\]
The constant $C_{\tilde B}$ depends on $\nbar,\Nbar,\kappa,\mu_*,c_1,b_0,p_0,R,\delta_0,M_b,\alpha_0,e,\lambda,\eta,d,L,a$, and not on $\eps$, $k$ or $\ell_T$.
\end{proposition}
\begin{proof}
(1) By Lemma~\ref{app:G1:I3c-Bt}, $\norm{h\tilde B^{(k)}}_{1,\alpha_0}\le C'_\kappa t^\kappa$.

(2) In a term with a factor $\tilde B^{(k)}$, replace one vector-leg bound $Cp$ by $C'_\kappa t^\kappa$, and bound every other factor without using any gain. The factors $e_k\tilde A$ are $\le Ct^{1/2}p$, which is at most a polylog. The end fields are $\le(Cp\lambda^{-1/4}t^{-1/4})^2$. By (IH) and (CAP), the families obey $(1.33)^\natural$, since their graded order is $\le3\nbar\le\Nbar$. They have at most $s(\nbar)$ scalar legs, which contribute $\le t^{-s(\nbar)/4}$ times a polylog. Hence every such term is $\le C(\nbar)C'_\kappa t^{\kappa-s(\nbar)/4-1/2}\mathrm{e}^{-\delta d}$ times a polylog, and $\kappa-s(\nbar)/4-\frac12\ge\frac92$. A factor $\tilde B$ inside a $V$ of a line is covered by the same multilinear count, through the product structure of Step~2 of the proof of Proposition~\ref{app:G1:T579}.

(3) There are at most $2^{D(\nbar)}$ splittings per term. The localisation sum and the volume are treated as in Proposition~\ref{app:G1:I3c-R}. Finally, $\sup_tt^{9/2-16/5}\,\mathrm{polylog}<\infty$.
\end{proof}

\begin{remark}[scope of the three cuts]\label{app:G1:I3c-scope}
(i) Propositions~\ref{app:G1:T579}, \ref{app:G1:I3c-R} and~\ref{app:G1:I3c-B} do not cover two operations: the propagator change $G_k(B^k(\Lambda_0^{(k)}),\tilde B)\to G_k(B^k(\Lambda_2^{(k)}),B^{(k+1),\eta})$ at the end of \cite[p.~579]{Ba2} (\cite[Prop.~2.6]{Ba2}), and the cuts of the determinant expansions \cite[(2.103)--(2.104), p.~578]{Ba2}.
(ii) Through (IH) and (CAP), the three cuts depend on the two-domain bounds (Proposition~\ref{app:G1:twodomain}, Lemma~\ref{app:G1:I3r-DGM}) and on the $\square$-remainder of \cite[p.~433]{Ba3}, but only inside the families of \cite[Prop.~1]{Ba3}. No domain difference of an operator (1.16) occurs, because all lines of a given term live on one domain. So this dependence enters Proposition~\ref{app:G1:I3} once, through items 2--10 of Table~\ref{app:G1:I3tab}.
\end{remark}

\begin{proof}[Proof of Lemma~\ref{app:G1:I3r-T579}]
The decomposition is the exact expansion (E1)--(E3) of the proof of Proposition~\ref{app:G1:T579}, in which $\tilde A=\theta_{k+1}A'^{(k)}+\tilde B^{(k)}$ is split multilinearly in the terms of order $\le\nbar$ (Proposition~\ref{app:G1:I3c-B}).
\begin{itemize}
\item The bound on $\Delta^{\rm pre}_k$ is Proposition~\ref{app:G1:T579}, under its hypotheses; its per-order input is Lemma~\ref{app:G1:T579-po}, and it assumes the bound of Statement~\ref{app:K:L06-4}.
\item The bound on $\Delta^R_k$ is Proposition~\ref{app:G1:I3c-R}.
\item The bound on $\Delta^{\tilde B}_k$ is Proposition~\ref{app:G1:I3c-B}, at $\kappa=31$.
\end{itemize}
The hypotheses on $\tilde A$ and on $\tilde B^{(k)}$ used in Lemmas~\ref{app:G1:I3c-Rfac} and~\ref{app:G1:I3c-Bt} are supplied by Proposition~\ref{app:G1:AAext}, in the form without a smallness condition on $\eps_0$ obtained in the proof of Lemma~\ref{app:G1:I3r-AA}, and by Lemma~\ref{app:G1:I3c-Bt} itself. The two-domain bounds enter only inside the families (Remark~\ref{app:G1:I3c-scope}(ii)). They hold at every $k\ge1$, including $k=1$, where $\eta=\frac12$, by Lemma~\ref{app:G1:I3a-margins}(d). Each of the three bounds is two-sided, is $\le C't_k^{16/5}N_k$, and has the stated dependence of its constant.
\end{proof}

\subsection{The mass reading, the final c-numbers and the vacuum classes at the last scale}\label{app:G1:part-E6}

We prove Propositions~\ref{app:G1:I7}, \ref{app:G1:I8} and \ref{app:G1:I9} of Table~\ref{tab:T6-inputs}; the region entropy is Proposition~\ref{app:G1:I6}.
Throughout: $d=3$; Ba{\l}aban's periodic lattice torus $T$ of side $\ell_T$, $|T|=\ell_T^d$; $L=2$; $\eps=2^{-K}\eps_0$, $t_k:=L^k\eps$ ($t_K=\eps_0$); $\ell_k:=1+\ln t_k^{-1}$; $r_k:=r(t_k)=R\ell_k^{\varrho}$, $\varrho=2$; $p(t)=b_0(1+\ln t^{-1})^{p}$, $p=6$; $M_b$ is Ba{\l}aban's block parameter.

\paragraph{The mass bound \texorpdfstring{``$m^2\le O(1)$''}{m2 <= O(1)} of \texorpdfstring{\cite[p.~589]{Ba2}}{[Ba2, p. 589]}.}\label{app:G1:ssec-I7}
\cite[Prop.~3.1, p.~589]{Ba2} ends: ``We assume $m^2\le O(1)$ also.'' (read at the page). Write $\nu_k:=M^2t_k^2$ for the rescaled Higgs mass at scale $k$, $\lambda(t):=\lambda t^{4-d}$, $a_\flat:=a(1-L^{-2})<a_k\le a$ ($a_\flat=3/4$), and $c_1$ for the constant of \cite[(2.55)]{Ba2}. Since $t_k\le\eps_0$, $\sup_{k\le K}\nu_k\le M^2\eps_0^2$.

\begin{lemma}[where no mass bound is needed]\label{app:G1:M4A}
(i) \cite[Theorem (1.9)--(1.12), p.~573]{Ba4} (\,$=$\,\cite[Prop.~2.1]{Ba1}) holds for every $\nu\ge0$, with constants depending on $d,M_b$ ($c_0$ also on $\alpha$).
(ii) For $0\le\nu\le\nu'$: $\Delta_0\le\Delta_\nu\le\Delta_{\nu'}\le a_kI$, where $\Delta_\nu:=a_k-a_k^2Q_kG_\nu Q_k^*$ and $G_\nu:=(-\Delta^{\eta,N}_{A,\Omega}+\nu+a_kP_k(A))^{-1}$. Hence \cite[(2.33), p.~611]{Ba1} holds for every $\nu\ge0$, with its $\nu=0$ lower constant and upper constant $a(1+L^{-2})=5/4$.
\end{lemma}
\begin{proof}[Proof of (ii)]
$X_\nu\le X_{\nu'}$ as forms, and $X_0>0$ by \cite[(1.8)]{Ba4}. Operator monotonicity of the inverse gives $G_{\nu'}\le G_\nu$, and $G_\nu\ge0$ gives $\Delta_\nu\le a_k\le a$. Finally $\norm{P}\le1$.
\end{proof}

\begin{lemma}[Prop.~3.1 of \cite{Ba2} with explicit constant]\label{app:G1:M4B}
Let $0\le\nu\le c_m$ and $\gamma_0(c_m):=\tfrac12\min\{\gamma_0^{(0)},a_\flat/(a_\flat+c_m)\}$, where $\gamma_0^{(0)}$ is the constant of \cite[(1.22), p.~574]{Ba4} at $m^2=0$. Under the hypotheses of that proposition,
$\ip\varphi{\Delta_\nu\varphi}\ge\gamma_0(c_m)\big(\sum_{\langle x,x'\rangle}\abs{U(A)\varphi(x')-\varphi(x)}^2+\nu\sum_x\abs{\varphi(x)}^2\big)-\tfrac12O(1)\mathrm{e}^{2-\alpha}\sum_x\abs{\varphi(x)}^2$.
Consequently \cite[(3.26)]{Ba2} holds with $\gamma_0^{\rm II}:=\min\{1,\gamma_0(c_m)\}$ whenever $M^2\eps_0^2\le c_m$. We take $c_m:=1$.
\end{lemma}
\begin{proof}
Since $-\Delta^N\ge0$, $G_\nu\le(\nu+a_kP_k)^{-1}=P_k/(\nu+a_k)+(I-P_k)/\nu$. As $Q_kQ_k^*=I$, this gives $Q_kG_\nu Q_k^*\le(\nu+a_k)^{-1}$ and $\Delta_\nu\ge a_k\nu/(a_k+\nu)\ge[a_\flat/(a_\flat+c_m)]\nu$. Also $\Delta_\nu\ge\Delta_0$ (Lemma~\ref{app:G1:M4A}). Write $\Delta_\nu=\tfrac12\Delta_\nu+\tfrac12\Delta_\nu$ and apply \cite[(1.22)]{Ba4} at $m^2=0$ to $\Delta_0$. For (3.26): (3.27)--(3.28) decouple the scales by Neumann conditions for every $m$; the $k$-th term rescaled is (3.29), which is the display with $\nu=\nu_k$; and the $k=0$ term has coefficient $1$.
\end{proof}

\begin{lemma}[{the mass defect in \cite[Lemma~2.4, (2.75)--(2.76)]{Ba2}}]\label{app:G1:M4D}
On a cube $\square$ of large blocks with Neumann conditions and $A_0$ gauged out, $a_kG_k(\square,0)Q_k^*1=\frac{a_k}{a_k+\nu_k}1$. So the constant reproduction (2.76) carries the extra term $E_{\rm mass}:=-\frac{\nu_k}{a_k+\nu_k}\varphi'(y)$. With $C_\varphi:=c_1L^{3/4}(1+\log L)^p$ ($\approx39.6\,c_1$ at $L=2$, $p=6$) one has $\abs{\varphi'(y)}\le C_\varphi\lambda(t_k)^{-1/4}p(t_k)$, hence
$\abs{E_{\rm mass}}\le(C_\varphi/a_\flat)\lambda^{-1/4}M^2t_k^{7/4}p(t_k)$. Under the condition
\[
\text{(E.$\Phi$)}\qquad \lambda^{-1/4}m_2^2\eps_0^{7/4}\le a_\flat/C_\varphi\qquad(a_\flat/C_\varphi\approx1.89\cdot10^{-2}/c_1),
\]
$\abs{E_{\rm mass}}\le p(t_k)$ for all $k\le K$ and $M\le m_2$. Then (2.65) and (2.76) hold with Ba{\l}aban's $O(p)$-constant plus $1$, and (2.78) keeps a per-site error $O(t^{1/4}p^4)\le O(t^{\kappa_0})$ for $\kappa_0<1/4$, in particular for the value $\kappa_0=\frac15$ used here, with an $m$-free constant.
\end{lemma}
\begin{proof}
$-\Delta^{\eta,N}_{0,\square}1=0$ and $P_k1=1$ on $\square$, so $X_\nu1=(a_k+\nu_k)1$. Split $a_kGQ^*1_{\square_1}\varphi'$ into $\varphi'(y)\,a_kGQ^*1$, a cut term (exponentially small in $r(t)$, uniform in $\nu$ by Lemma~\ref{app:G1:M4A}(i)) and Ba{\l}aban's gradient term $O(p)$. The field bound comes from (2.55) at index $k-1$: $\lambda(t/L)^{-1/4}=L^{1/4}\lambda(t)^{-1/4}$ ($d=3$), the canonical factor is $L^{1/2}$, and $p(t/L)\le(1+\log L)^pp(t)$ because $1+\log L+\log t^{-1}\le(1+\log L)(1+\log t^{-1})$ for $t\le1$. Next, $\nu_k\lambda(t_k)^{-1/4}=\lambda^{-1/4}M^2t_k^{7/4}$; this increases in $t_k\le\eps_0$ and in $M\le m_2$. For (2.78): $\lambda(t)\abs{\abs{\varphi^{(k)}}^4-\abs{\varphi(y)}^4}\le4C'(C_\varphi+C')^3\lambda(t)^{1/4}p(t)^4$, with $C'$ Ba{\l}aban's constant plus $1$.
\end{proof}
\noindent\emph{Remark (role of (E.$\Phi$)).} (E.$\Phi$) is a normalisation. The condition (E.9) implies $6C_m\lambda^{-1/4}m_2^2\eps_0^{7/4}\le1$, so (2.65) holds with the increment $C_\varphi/(6a_\flat C_m)$, free of $\eps_0$ and $M$; (E.$\Phi$) makes the increment $1$, independent of $c_1$, so that no feedback between $c_1$ and \cite[Lemma~2.4]{Ba2} can arise.

\paragraph{The constants of Lemma~2.7 of \cite{Ba2} and of the counterterm coefficients.}\label{app:G1:ssec-I7const}
In this paragraph $d=3$, $L=2$, $a=1$ and $a_\flat=a(1-L^{-2})$. We write $t=t_k=L^k\eps\le\eps_0$ and $\nu=\nu_k=M^2t^2$, where $M\in(0,m_2]$ is the Higgs mass. We also write $e(t)=et^{1/2}$, $\lambda(t)=\lambda t$, $p(t)=b_0(1+\log t^{-1})^p$ and $r(t)=R(1+\log t^{-1})^{\varrho}$ with $\varrho>1$. $M_b$ is the large-block parameter of \cite{Ba2}. The constants $c_0,\delta_0,R_0$ are those of \cite[Theorem, (1.9)--(1.12), p.~573]{Ba4}, of \cite[Prop.~2.2, (2.58), pp.~570--571]{Ba2} and of \cite[(1.15)--(1.16), p.~574]{Ba4}, taken equal: we use the minimum of the $\delta_0$'s and the maximum of the $c_0$'s and of the $R_0$'s. By Lemma~\ref{app:G1:M4A}(i) and \cite[p.~572, (1.6)]{Ba4} they depend only on $d$ and $M_b$ (and $\alpha$), and they hold for every rescaled mass $\nu\ge0$. For $\delta>0$ put
\[
S_d(\delta):=\sum_{z\in\Z^d}\mathrm{e}^{-\delta\abs z},\qquad J_d(\delta):=\mathrm{e}^{\delta\sqrt d}\big(1+2\sqrt d/\delta\big)^d,\qquad S_F(\delta):=d\,\mathrm{e}^{\delta\sqrt d}S_d(\delta).
\]
These satisfy three elementary bounds:
\begin{itemize}
\item $S_d(\delta)\le(1+2\sqrt d/\delta)^d$;
\item for every $\eta\in(0,1]$ and every unit block $B$, $\eta^d\sum_{z\in\eta\Z^d}\mathrm{e}^{-\delta\dist(z,B)}\le J_d(\delta)$;
\item for every point $x'$, $\sum_F\mathrm{e}^{-\delta\dist(F,x')}\le S_F(\delta)$, where $F$ runs over the unit faces of the unit lattice.
\end{itemize}

\emph{Setting for Lemma~2.7 of \cite{Ba2}.} We work in scale-$k$ units on the fine lattice of spacing $\eta=L^{-k}$. The unit lattice is the lattice of unit blocks $B(x)$. Let:
\begin{itemize}
\item $\Omega:=B^k(\Lambda_2^{(k)})$, and $\hat B:=B^{(k+1),\eta}$ the background of \cite[(2.110)--(2.113), pp.~580--581]{Ba2};
\item $G_k:=(-\Delta^{\eta,N}_{\hat B,\Omega}+\nu+a_kP_k(\hat B))^{-1}$, where $P_k=Q_k^*Q_k$ and $Q_kQ_k^*=I$;
\item $\Delta:=a_k-a_k^2Q_kG_kQ_k^*$, acting on functions on $\Omega^{(k)}$;
\item $Q:=Q(\hat B)$ the averaging from the unit lattice to the lattice of $L$-blocks, so that $QQ^*=I$;
\item $\mathsf M:=\Delta+aL^{-2}Q^*Q$ and $C_{\Lambda_4}:=(\Lambda_4\mathsf M\Lambda_4)^{-1}$, as in \cite[(1.13)--(1.14)]{Ba4}.
\end{itemize}
Kernels of operators from the unit lattice carry unit weights, $(G_kQ_k^*f)(z)=\sum_y(G_kQ_k^*)(z,y)f(y)$, and fine-lattice pairings carry the weight $\eta^d$. We assume the following.
\begin{itemize}
\item[(G)] $\dist(\Lambda_5^{(k)},\Lambda_4^{(k)c})>r(t)$, $\dist(\Lambda_4^{(k)},\Omega^c)>2r(t)$ and $r(t)\ge R_0+1+\sqrt d$. The first two inequalities follow from the definition \cite[(2.8)]{Ba2} of the sets $\Lambda_i$. The third holds since $r(t)\ge3R_0+15$ under the conditions of Appendix~\ref{app:B}. Each $\Lambda_i^{(k)}$ is a union of large blocks.
\item[(Reg)] $\hat B$ is regular on $\Omega$: $\abs{\partial\hat B}\le c_Bp(t)$. Hence the conclusions of \cite[Prop.~2.2]{Ba2} hold for $\hat B$, by \cite[Lemma~2.3, (2.60)]{Ba2} at scale $k+1$.
\item[(Res)] The scalar block field $\psi$ and the scalar field $\phi$ satisfy three bounds:
\begin{itemize}
\item $\abs{\psi(y)}\le2\lambda(t)^{-1/4}p(t)$;
\item $\abs{\psi(y)-U(\hat B(\Gamma_{y,x}))\phi(x)}\le2Ldp(t)$ for $x\in B(y)$;
\item $\abs{U(\hat B(\Gamma_{y,x}\cup\langle x,x'\rangle\cup\Gamma_{x',y'}))\psi(y')-\psi(y)}\le4Ldp(t)$ for every bond $\langle x,x'\rangle$ with $x\in B(y)$ and $x'\in B(y')$.
\end{itemize}
These are the restrictions \cite[(2.16), p.~560]{Ba2} at scale $t$, except that there the transports are taken with the block field rather than with $\hat B$. Replacing one transport by the other changes the right-hand sides by terms of size $e(t)p(t)\abs\psi=O(t^{1/4}p(t)^2)$, and this affects only $\mathcal E(t)$ below.
\item[($\Gamma$)] The contours $\Gamma$ of \cite[\S1]{Ba1} define $Q$, $Q_k$ and their adjoints. For a fine bond $b$, consider the closed loop formed by $b$ and the contours that transport $\psi$ to the two endpoints of $b$. This loop bounds a lattice surface of area at most $c_\Gamma L^2$, and of area at most $c_\Gamma\eta$ if $b$ lies inside one unit block. Here $c_\Gamma$ depends only on $d$.
\end{itemize}

\begin{lemma}[Lemma~2.7 of \cite{Ba2} with the mass term explicit]\label{app:G1:I7-L27}
Assume (G), (Reg), (Res) and ($\Gamma$). Then for every $x\in\Lambda_5^{(k)}$
\[
aL^{-2}(C_{\Lambda_4}Q^*\psi)(x)=(Q^*\psi)(x)+\mathcal R(x),\qquad
\abs{\mathcal R(x)}\le\big(K_2^0+\mathcal E(t)+K_m\lambda^{-1/4}M^2t^{7/4}\big)\,p(t),
\]
where
\begin{gather*}
K_2^0:=4Ld\,a\,c_0^2\mathrm{e}^{\delta_0}S_d(\delta_0)S_F(\delta_0),\qquad K_m:=2a\,c_0^2S_d(\delta_0),\\
\mathcal E(t):=a\,c_0S_d(\delta_0)\,\mathcal E_1(t)+4a^2c_0^2\mathrm{e}^{\delta_0\sqrt d}S_d(\delta_0/2)^2\lambda^{-1/4}t^{-1/4}\mathrm{e}^{-\delta_0r(t)/2},\\
\mathcal E_1(t):=2\lambda^{-1/4}\Big[c_0c_Bc_\Gamma e\,\mathrm{e}^{\delta_0}\big(S_F(\delta_0)L^2+dJ_d(\delta_0)\big)t^{1/4}p(t)+c_0\mathrm{e}^{2\delta_0\sqrt d}S_F(\delta_0/2)\,t^{-1/4}\mathrm{e}^{-\delta_0r(t)/2}\Big].
\end{gather*}
$K_2^0$ and $K_m$ depend only on $d,a,L,M_b$. $\mathcal E(t)$ depends in addition on $e,\lambda,b_0,p,R,\varrho,c_B,c_\Gamma$, but not on $M$, $\eps$, $k$ or $\eps_0$. Moreover $\mathcal E(t)\to0$ as $t\to0$, and $\mathcal E^*:=\sup_{0<t\le1}\mathcal E(t)<\infty$. Consequently \cite[(2.114)]{Ba2} holds on $\Lambda_5^{(k)}$ in the form
\[
\abs{\phi'(x)}\le\big(2Ld+K_2^0+\mathcal E(t)+K_m\lambda^{-1/4}M^2t^{7/4}\big)p(t),
\]
where $\phi'=\phi-aL^{-2}C_{\Lambda_4}Q^*\psi$ \cite[(2.110)]{Ba2}.
\end{lemma}
\begin{proof}
Put $g:=Q^*\psi$. Then $\abs g\le2\lambda(t)^{-1/4}p$, and all transports are unimodular.

\emph{Decomposition.} Since $QQ^*=I$ we have $Q^*Qg=g$, and $\Lambda_4$ commutes with $Q^*Q$. Hence $\Lambda_4\mathsf M\Lambda_4g=aL^{-2}\Lambda_4g+\Lambda_4\Delta g-\Lambda_4\Delta\Lambda_4^cg$. Applying $C_{\Lambda_4}$ gives $aL^{-2}C_{\Lambda_4}\Lambda_4g=\Lambda_4g+C_{\Lambda_4}\Lambda_4\Delta\Lambda_4^cg-C_{\Lambda_4}\Lambda_4\Delta g$.

Let $S$ be the union of the unit blocks $B(y)\subset\Omega$ with $\dist(B(y),\Omega^c)\ge R_0+1$, and let $\chi_S$ be its indicator. By (G), $\Lambda_4\subset S$ and $\dist(\Lambda_4,S^c)\ge r(t)$. Therefore
\[
\mathcal R=C_{\Lambda_4}\Lambda_4\big[\Delta\Lambda_4^cg-\Delta\chi_{S^c}g\big]-C_{\Lambda_4}\Lambda_4\Delta(\chi_Sg),
\]
and $\norm{C_{\Lambda_4}}_{\infty\to\infty}\le c_0S_d(\delta_0)$ by \cite[(1.16)]{Ba4}.

\emph{Off-diagonal terms.} Let $f$ be supported in $\Lambda_4^c$ with $\abs f\le\abs g$, and let $x'\in\Lambda_4$. Then $(\Delta f)(x')=-a_k^2(Q_kG_kQ_k^*f)(x')$. Every point of $B(x')$ has distance $\ge R_0$ from $\Omega^c$, so \cite[Prop.~2.2]{Ba2} gives
\[
\abs{(\Delta f)(x')}\le a^2c_0\mathrm{e}^{\delta_0\sqrt d}\sum_{x''\notin\Lambda_4}\mathrm{e}^{-\delta_0\abs{x'-x''}}\abs{g(x'')}.
\]
For $x\in\Lambda_5$ we have $\abs{x-x'}+\abs{x'-x''}>r(t)$. Splitting the exponentials in halves gives $\abs{(C_{\Lambda_4}\Lambda_4\Delta f)(x)}\le2a^2c_0^2\mathrm{e}^{\delta_0\sqrt d}S_d(\delta_0/2)^2\mathrm{e}^{-\delta_0r(t)/2}\lambda^{-1/4}t^{-1/4}p$. This applies to $f=\Lambda_4^cg$ and to $f=\chi_{S^c}g$, and it gives the last term of $\mathcal E(t)$.

\emph{Key identity.} Put $u:=Q_k^*(\chi_Sg)$. Then $Q_k^*Q_ku=u$, because $\chi_S$ is constant on unit blocks. With $G_k^{-1}=(-\Delta_{\hat B}+\nu)+a_kQ_k^*Q_k$, it follows that $a_kG_kQ_k^*(\chi_Sg)=u-G_k(-\Delta_{\hat B}+\nu)u$. Applying $a_kQ_k$, and using that $G_k$ is self-adjoint,
\[
(\Delta\chi_Sg)(x')=a_k\ip{w_{x'}}{(-\Delta_{\hat B}+\nu)u},\qquad w_{x'}:=G_kQ_k^*\delta_{x'} .
\]

\emph{Mass term.} We have $a_k\nu\ip{w_{x'}}u=a_k\nu(Q_kG_ku)(x')$. By \cite[(1.10)]{Ba4}, $\abs{(G_ku)(z)}\le c_0\norm u_\infty$ on $B(x')$. Hence $\abs{a_k\nu\ip{w_{x'}}u}\le2ac_0\lambda^{-1/4}M^2t^{7/4}p$, and after $C_{\Lambda_4}$ this is the $K_m$-term. This is the only place where $\nu$ enters, apart from $c_0,\delta_0,R_0$.

\emph{Kinetic term.} Write $\ip w{-\Delta_{\hat B}u}=\sum_{b\subset\Omega}\eta^d\,\overline{(\eta^{-1}d_bw)}(\eta^{-1}d_bu)$, where $\eta^{-1}d_b$ is the covariant difference on the bond $b$. If $d_bu\ne0$, then $\dist(b,\Omega^c)\ge R_0$, and \cite[(2.58)]{Ba2} gives $\abs{\eta^{-1}d_bw_{x'}}\le c_0\mathrm{e}^{-\delta_0\dist(b,x')}$. The bonds with $d_bu\ne0$ are of three kinds.
\begin{itemize}
\item Bonds crossing a unit face, between two $L$-blocks in $S$ or inside one $L$-block. By (Res), ($\Gamma$), Stokes' formula, (Reg) and $\abs{\mathrm{e}^{i\theta}-1}\le\abs\theta$, such a bond has $\abs{d_bu}\le4Ldp+c_\Gamma L^2e(t)c_Bp\abs\psi$.
\item Bonds inside a unit block. Here $\abs{\eta^{-1}d_bu}\le c_\Gamma e(t)c_Bp\abs\psi$.
\item Bonds crossing a face of $\partial B^k(S)$. Here $\abs{d_bu}\le\abs\psi$, and the face has distance $\ge r(t)-\sqrt d$ from $x'$.
\end{itemize}
A unit face carries $\eta^{1-d}$ fine bonds of weight $\eta^{d-1}$. So a face $F$ contributes at most
\[
c_0\mathrm{e}^{\delta_0}\mathrm{e}^{-\delta_0\dist(F,x')}\sup_{b\in F}\abs{d_bu}.
\] Summing with $S_F$, $J_d$ and $S_F(\delta_0/2)$, and using $e(t)p\abs\psi\le2e\lambda^{-1/4}t^{1/4}p^2$ and $\abs\psi\le2\lambda^{-1/4}t^{-1/4}p$, we get
\[
\abs{a_k\ip{w_{x'}}{-\Delta_{\hat B}u}}\le a\big(4Ld\,c_0\mathrm{e}^{\delta_0}S_F(\delta_0)+\mathcal E_1(t)\big)p.
\]

\emph{Sum.} Multiply the mass and kinetic bounds by $c_0S_d(\delta_0)$ and add the off-diagonal terms. This gives the bound on $\mathcal R$. The function $t^{1/4}(1+\log t^{-1})^p$ is bounded on $(0,1]$ and tends to $0$ as $t\to0$. Since $\varrho>1$, the same holds for $t^{-1/4}\mathrm{e}^{-\delta_0R(1+\log t^{-1})^\varrho/2}$. For the last claim, $\abs{\phi'}\le\abs{\phi-Q^*\psi}+\abs{\mathcal R}$, and $\abs{\phi-Q^*\psi}\le2Ldp$ by (Res).
\end{proof}

\begin{corollary}[the mass defect of Lemma~2.7 of \cite{Ba2}]\label{app:G1:I7-U6}
Assume (E.$\Phi$): $\lambda^{-1/4}m_2^2\eps_0^{7/4}\le a_\flat/C_\varphi$, with $C_\varphi$ as in Lemma~\ref{app:G1:M4D}. Then for every $M\in(0,m_2]$, every $t\le\eps_0\le1$ and every $x\in\Lambda_5^{(k)}$,
\[
\abs{\mathcal R(x)}\le K_2^\natural\,p(t),\qquad K_2^\natural:=K_2^0+\mathcal E^*+K_ma_\flat/C_\varphi ,
\]
and $\abs{\phi'(x)}\le(2Ld+K_2^\natural)p(t)$. The constant $K_2^\natural$ does not depend on $M$, $\eps$, $k$ or $\eps_0$.
\end{corollary}
\begin{proof}
$\lambda^{-1/4}M^2t^{7/4}$ is increasing in $M$ and in $t$, so it is at most $\lambda^{-1/4}m_2^2\eps_0^{7/4}\le a_\flat/C_\varphi$. Now apply Lemma~\ref{app:G1:I7-L27}.
\end{proof}

In the next three lemmas, $\kappa_G$, $\kappa_H$, $A_0^T$, $A_1^T$, $S_L$, $K_1$ (from (F4)), $A_E$, $A_\nabla$ and the lattice sums $J_0,J_1,J_2$ are those of Lemmas~\ref{app:G1:VSR-F}, \ref{app:G1:VSR-IM} and~\ref{app:G1:VSR-E} and of Subsection~\ref{app:G1:part-TVSR}.

\begin{lemma}[the coefficient {\cite[(3.16)]{Ba3}} uniformly in the mass]\label{app:G1:I7-W}
Let $j''\ge1$ and $\xi=L^{-j''}$. Work on the $\xi$-lattice torus $\Lambda$ of side $\ell=\ell_T/t_{j''}\in L\N$, with $\ell\ge L$. Let $\nu\in[0,1]$; for the scalar field $\nu=M^2t_{j''}^2$, and for a photon component $\nu=\mu_*^2t_{j''}^2$. Let $b:=a_{j''}\in[a_\flat,a]$, and
\[
G:=G^\xi_{j''}(0)=(-\Delta^\xi+\nu+bP_1)^{-1},\qquad C:=C^\xi=(-\Delta^\xi+1)^{-1},\qquad W:=(1-\nu)-bP_1 ,
\]
so that $G=C+GWC$ \cite[(3.16), p.~437]{Ba3}. Let $g,g'$ be bounded functions and $\abs q\le1$. Take the coordinate differences $y'_\mu-y_\mu$ in the fundamental domain centred at $y$, so that $\abs{y'_\mu-y_\mu}\le d(y,y')$. Then the coefficient of \cite[(3.16)]{Ba3},
\begin{multline*}
\Xi(y):=-q^2\sum_{y'}\xi^dC(y-y')g(y)G(y,y')g'(y')(y'_\mu-y_\mu)\\
+q^2\sum_{y'}\xi^d\sum_{\nu'=1}^d\big(\partial^\xi_{\nu'}GWC\partial^{\xi*}_{\nu'}\big)(y,y')\,g(y)G(y,y')g'(y')(y'_\mu-y_\mu),
\end{multline*}
satisfies
\[
\abs{\Xi(y)}\le\norm g_\infty\norm{g'}_\infty A_G\big[1+d(1+a)K_1\big(K_1+A_\nabla J_0(\kappa_G)\big)\big]=:K_{3.16},\qquad A_G:=A_0^T+\frac{A_E}{\mathrm{e}\,\kappa_G}.
\]
$K_{3.16}$ depends only on $(a,L)$ and on $\norm g_\infty\norm{g'}_\infty$. It does not depend on $\nu\in[0,1]$, $\xi$, $j''$, $\ell$ or $M$.
\end{lemma}
\begin{proof}
By Lemma~\ref{app:G1:VSR-E}, $G=F+E$ with $F=C$, $\abs E\le A_E\mathrm{e}^{-\kappa_Gd}$ and $\abs{\nabla_1E}\le A_\nabla\mathrm{e}^{-\kappa_Gd}$. By Lemmas~\ref{app:G1:VSR-F} and~\ref{app:G1:VSR-IM}, $0\le F\le A_0^T\mathrm{e}^{-d/8}/(d+\xi)$, since $S(\ell/4)\le S_L$ for $\ell\ge L$. These constants are uniform in $\nu\in[0,1]$. Hence $\abs{G(y,y')}d(y,y')\le A_G$.

The first term of $\Xi$ is at most $\norm g_\infty\norm{g'}_\infty A_G$, because $C\ge0$ and $\xi^d\sum_{y'}C(y-y')=1$.

Put $h(y'):=G(y,y')g'(y')(y'_\mu-y_\mu)$, so that $\norm h_\infty\le A_G\norm{g'}_\infty$. The second term equals $q^2g(y)\sum_{\nu'}(\partial_{\nu'}GWC\partial^*_{\nu'}h)(y)$. Its absolute value is at most $d\norm g_\infty\norm{\partial G}_{\infty\to\infty}\norm W_{\infty\to\infty}\norm{C\partial^*}_{\infty\to\infty}\norm h_\infty$, and the three operator norms are bounded as follows:
\begin{itemize}
\item $\norm W_{\infty\to\infty}\le\abs{1-\nu}+b\le1+a$, since $P_1$ is an average;
\item $\norm{C\partial^*}_{\infty\to\infty}\le K_1$, by summation by parts, (F4) and Lemma~\ref{app:G1:VSR-IM};
\item $\norm{\partial G}_{\infty\to\infty}\le K_1+A_\nabla J_0(\kappa_G)$, since $\nabla_1G=\nabla F+\nabla_1E$ and a torus sum of $\mathrm{e}^{-\kappa d}$ is at most the lattice sum $J_0(\kappa)$.\qedhere
\end{itemize}
\end{proof}

\begin{lemma}[the coefficients of {\cite[p.~441]{Ba3}} uniformly in the mass]\label{app:G1:I7-W441}
Work in the setting of Lemma~\ref{app:G1:I7-W}, with $j_0\ge1$ in place of $j''$. Let $G:=G^\xi_{j_0}(0)$ and $E:=G-C=GWC$. Fix $y\in\Lambda$ and directions $\mu,\mu',\nu_0$. Let $m(y')$ be either $1$ or $y'_{\nu_0}-y_{\nu_0}$, with coordinates taken in the fundamental domain centred at $y$. For $K\in\{G,C\}$ put
\begin{gather*}
X_1[K]:=\xi^d\sum_{y'}(K\partial^{\xi*}_{\mu'})(y,y')(K\partial^{\xi*}_{\mu})(y',y)\,m(y'),\qquad
X_2[K]:=\xi^d\sum_{y'}K(y,y')(\partial^\xi_{\mu'}K\partial^{\xi*}_{\mu})(y',y)\,m(y'),\\
X_3[K]:=K(y,y),\qquad X_4[K]:=\xi\,(K\partial^{\xi*}_\mu)(y,y).
\end{gather*}
Then $\abs{X_i[G]-X_i[C]}\le K_{441}$ for $i=1,\dots,4$. The constant $K_{441}$ depends only on $(a,L)$, and not on $\nu\in[0,1]$, $\xi$, $j_0$, $\ell\ge L$, $y$ or the directions.
\end{lemma}
\begin{proof}
Write $r:=d(w,w')$ and $\kappa:=\kappa_G\le\frac1{16}$. By Lemmas~\ref{app:G1:VSR-F}, \ref{app:G1:VSR-IM} and~\ref{app:G1:VSR-E} (items (E1) and (E3)), together with the symmetry of $E$:
\begin{itemize}
\item $\abs{C}\le A_0^T\mathrm{e}^{-r/8}/(r+\xi)$ and $\abs{\nabla C}\le A_1^T\mathrm{e}^{-r/8}/(r+\xi)^2$;
\item $\abs E\le A_E\mathrm{e}^{-\kappa r}$, and every difference of $E$ is at most $A_\nabla\mathrm{e}^{-\kappa r}$;
\item hence $\abs{\nabla G}\le(A_1^T(r+\xi)^{-2}+A_\nabla)\mathrm{e}^{-\kappa r}$.
\end{itemize}
Shifting the point at which a difference is evaluated by $\xi$ changes these bounds by at most the factor $c_*:=4\mathrm{e}^{1/8}$. Moreover $\abs m\le1+r$, and every torus sum below is at most the corresponding sum over $\xi\Z^3$, that is, at most $J_0$, $J_1$ or $J_2$.

\emph{$X_3$ and $X_4$.} The differences are $E(y,y)$ and $\xi(E\partial^*_\mu)(y,y)$, which are at most $A_E$ and $A_\nabla$.

\emph{$X_1$.} The difference $X_1[G]-X_1[C]$ is the sum over $y'$ of $\xi^d\big[(E\partial^*_{\mu'})(y,y')(G\partial^*_\mu)(y',y)+(C\partial^*_{\mu'})(y,y')(E\partial^*_\mu)(y',y)\big]m(y')$. Each product has one factor at most $A_\nabla\mathrm{e}^{-\kappa r}$ and one factor at most $(A_1^T(r+\xi)^{-2}+A_\nabla)\mathrm{e}^{-\kappa r}$. Using $(1+r)(r+\xi)^{-2}\le(r+\xi)^{-2}+(r+\xi)^{-1}$ and $(1+r)\mathrm{e}^{-\kappa r}\le(1+2/(\mathrm{e}\kappa))\mathrm{e}^{-\kappa r/2}$, the sum is at most $A_\nabla[2A_1^T(J_2(\kappa)+J_1(\kappa))+A_\nabla(1+2/(\mathrm{e}\kappa))J_0(\kappa/2)]$.

\emph{$X_2$.} Put $f_K:=K(y,\cdot)m$ and $k_K:=(K\partial^*_\mu)(\cdot,y)$. Then $X_2[G]-X_2[C]=\xi^d\sum_{y'}[f_E\,\partial_{\mu'}k_G+f_C\,\partial_{\mu'}k_E](y')$. Sum by parts in $y'$ on the torus.
\begin{itemize}
\item If $m\equiv1$, there is no boundary term.
\item Otherwise $\partial^*_{\mu'}(K(y,\cdot)m)=(\partial^*_{\mu'}K(y,\cdot))\,m(\cdot-\xi e_{\mu'})+K(y,\cdot)\,\partial^*_{\mu'}m$. Here $\abs{\partial^*_{\mu'}m}\le1$ except on at most $(\ell/\xi)^{d-1}$ sites. On those sites $\abs{\partial^*m}\le\ell/\xi$ and $d(y,y')\ge\ell/2-1$, so they contribute at most $c_*^2\ell^d(A_E+A_0^T)(A_1^T+A_\nabla)\mathrm{e}^{-\kappa(\ell/2-1)}$, which is bounded uniformly in $\ell$.
\end{itemize}
On the remaining sites:
\begin{itemize}
\item $\abs{\partial^*f_E}\le c_*(A_\nabla(2+r)+A_E)\mathrm{e}^{-\kappa r}$ and $\abs{k_G}\le(A_1^T(r+\xi)^{-2}+A_\nabla)\mathrm{e}^{-\kappa r}$;
\item $\abs{\partial^*f_C}\le c_*(A_1^T(2+r)(r+\xi)^{-2}+A_0^T(r+\xi)^{-1})\mathrm{e}^{-r/8}$ and $\abs{k_E}\le A_\nabla\mathrm{e}^{-\kappa r}$.
\end{itemize}
In each product only one factor is singular, of order at most $(r+\xi)^{-2}$. So the sums are bounded by constant multiples of $J_0(\kappa/2)$, $J_1(\kappa)$ and $J_2(\kappa)$.
\end{proof}

\begin{lemma}[the counterterm lines at zero background, uniformly in the mass]\label{app:G1:I7-lines}
Let $0<M\le m_2$ and $\mu_*\ge0$, and assume $m_2\eps_0\le1$, $\mu_*\eps_0\le1$ and $\eps_0\le1$. By \cite[(1.27), p.~418]{Ba3},
\[
G_K(0)=\sum_{j=0}^{K-1}\Gamma_j,\qquad \Gamma_0=G_1(0)=C^{(0),\eps}(0),\qquad \Gamma_j=G_{j+1}(0)-G_j(0)\quad(j\ge1).
\]
For the scalar field (mass $M$), and for each component of the photon (mass $\mu_*$, Lemma~\ref{app:G1:VSR-V0}), every $\Gamma_j$ satisfies \cite[(2.10)--(2.12), p.~426]{Ba3} at background $0$ with index $j$. The constants, and the rate $\delta_1:=\min\{\kappa_H,\min(1/8,\kappa_G)/L\}$, depend only on $(a,L)$ and $\alpha$. They are uniform in $\eps$, $j$, $\ell_T$, $M$ and $\mu_*$.
\end{lemma}
\begin{proof}
\emph{The case $j\ge1$.} This is Corollary~\ref{app:G1:VSR-VSphys}, with $m_s=M$ and $m_v=\mu_*$. It applies because $t_{j+1}\le t_K\le\eps_0\le\min(1/m_f,1)<L$, so $m_ft_j\in[0,1]$.

\emph{The case $j=0$.} Work in scale-$1$ units, with spacing $1/L$. By Lemma~\ref{app:G1:VSR-E} at $j=1$, $G_1(0)=F+E$. Then (F0), (F1), (E1) and (E3) give two bounds:
\begin{itemize}
\item $\abs{G_1(0)}\le(LA_0^T+A_E)\mathrm{e}^{-\min(1/8,\kappa_G)d}$;
\item each lattice difference costs at most a factor $L$, with constant $L^2A_1^T+A_\nabla$.
\end{itemize}
In physical units these are (2.10) and its derivative form at index $0$, with rate $\min(1/8,\kappa_G)/L$. Distinct lattice points are at distance $\ge\eta$, so the H\"older quotient (2.11) is at most twice the supremum of the difference divided by $\eta^\alpha$. The contour sums (2.12) follow by summing (2.10) along the contour.
\end{proof}

\emph{The counterterm coefficients.} By \cite[(1.23)--(1.28), pp.~416--418]{Ba3}, $\delta m^2$ is a finite sum over one-particle-irreducible graphs $G_0$ with two scalar legs, including the lower-order $\delta m^2$-insertions. Their orders $e^\alpha\lambda^\beta$ satisfy $2\le\alpha+2\beta\le4$, so that $o:=\alpha/2+\beta\in\{1,2\}$. $\delta m^2_K(x)$ is the same sum with propagators $G_K(0)$. Expand each line by Lemma~\ref{app:G1:I7-lines}. Then $\delta m^2_{K,k}(x)$ is the sum over index assignments with largest index $\ge k$, and $\delta m^2_k(x)$ the sum over those with all indices $<k$ \cite[(1.28)]{Ba3}.

Work in scale-$K$ units, with $\eta_K:=\eps/t_K$. Then $L^j\eta_K=t_j/t_K\le1$, $e(t_K)^2=e^2t_K$ and $\lambda(t_K)=\lambda t_K$. For an amputated kernel $\Sigma_{G_0}(y_1,y_2)$ put $c_{G_0}(x):=\sum_{y_2}\eta_K^d\Sigma_{G_0}(x,y_2)$. Then $t_K^2\,\delta m^2_{K,k}(x)=\sum_{G_0}c^{[\ge k]}_{G_0}(x)$.

\begin{theorem}[pointwise bound on $\delta m^2_{K,k}$]\label{app:G1:I7-I3a}
Let $\eps=L^{-K}\eps_0$ with $K\in\N$, so that $t_K=\eps_0$. Let $0\le k\le K$ and $t:=t_k$, and assume the hypotheses of Lemma~\ref{app:G1:I7-lines}. Then for every point $x$ of the torus of spacing $\eps$,
\[
\abs{\delta m^2_{K,k}(x)}\le C_m\Big[(e^2+\lambda)\,t^{-1}\frac{L}{L-1}+(e^4+e^2\lambda+\lambda^2)(K-k)\Big],
\]
and $K-k=\log_L(\eps_0/t)$, so that $t^2\abs{\delta m^2_{K,k}(x)}=O(t)$. Here $C_m$ is the constant $O(1)$ of \cite[(1.33)]{Ba3} for the expressions of Lemma~\ref{app:G1:VSR-C} (it contains the constant $C_A$ of Lemma~\ref{app:G1:VSR-A}, with $\delta_1$ of Lemma~\ref{app:G1:I7-lines}), times the number of graphs. It does not depend on $\eps$, $k$, $x$ or $\ell_T$. For $k=K$ the left side is $0$.
\end{theorem}
\begin{proof}
Fix a graph $G_0$ and a value $j^*\in\{k,\dots,K-1\}$ of the largest index. By Lemma~\ref{app:G1:VSR-Btop}, all renormalisations act inside the sub-sum with top index $j^*$, and $G_0$ itself is not renormalised. By Lemma~\ref{app:G1:VSR-C}, the expression with the leg at $y_1$ fixed at $x$ obeys \cite[(1.33)]{Ba3}. The external fields are constant with norm $O(1)$. Lemma~\ref{app:G1:VSR-A}(ii) at fixed $j^*$ then gives
\[
\abs{c^{[j^*]}_{G_0}(x)}\le O(1)\,e(t_K)^\alpha\lambda(t_K)^\beta(t_{j^*}/t_K)^{D(G_0)},
\]
with the degree $D(G_0)$ of that lemma. At background $0$, a derivative on a constant external leg vanishes. Close the two legs by a scalar line $\ell$. Lemma~\ref{app:G1:cutclasses}(a),(b) with $n_\partial(\ell)=0$ give $D(G_0)\ge o-2$.

Now sum over $j^*$, using $t_{j^*}/t_K\le1$.
\begin{itemize}
\item For $o=1$: $\sum_{j^*=k}^{K-1}(t_{j^*}/t_K)^{-1}\le(t_K/t)L/(L-1)$, and $t_K^{-2}(e^2+\lambda)t_K(t_K/t)L/(L-1)=(e^2+\lambda)t^{-1}L/(L-1)$.
\item For $o=2$: there are $K-k$ terms, each at most $1$, which gives $(e^4+e^2\lambda+\lambda^2)(K-k)$.
\end{itemize}
Finally $K-k=\log_L(\eps_0/t)$, and $t\ln(\eps_0/t)\le\eps_0/\mathrm{e}$.
\end{proof}

\begin{lemma}[where the masses enter the proof of {\cite[(1.33)]{Ba3}}]\label{app:G1:I7-masses}
Consider the graphs of orders $2\le\alpha+2\beta\le4$ at zero background without operator lines. Fix $j_0\ge1$ and work in units of $t_{j_0}$ on the torus of spacing $\xi=L^{-j_0}$. Let $G_s$ be the scalar propagator $G^\xi_{j_0}(0)$ and $G_p$ each component of the photon propagator. Put $E_f:=G_f-C^\xi$ for $f\in\{s,p\}$.

In \cite[pp.~434--445]{Ba3}, the masses $M,\mu_*$ enter the constant $O(1)$ of \cite[(1.33)]{Ba3} only through the following quantities:
\begin{itemize}
\item[(i)] the line bounds \cite[(2.10)--(2.12)]{Ba3};
\item[(ii)] pointwise bounds on $G_s,G_p,E_s,E_p$ and on their differences, at pp.~437--444;
\item[(iii)] the coefficient of (3.16) on p.~437, and the coefficients of p.~441;
\item[(iv)] the quantity $\xi G_f(y,y)$ in (3.21) on p.~438;
\item[(v)] the difference between the coefficient of (3.23) on p.~439 and its $C^\xi$-version;
\item[(vi)] for the classes (2.19), (2.20) and (2.21c), the difference between the vertex coefficient $\sum\xi^{2d}\Gamma'_\mu$ of p.~444 and its $C^\xi$-version;
\item[(vii)] the expressions produced by the replacements in (iii), (v) and (vi), which \cite{Ba3} calls convergent.
\end{itemize}
In particular, p.~442 contributes no mass dependence at these orders.
\end{lemma}
\begin{proof}
We go through the pages in order.
\begin{itemize}
\item P.~434: (3.4) decomposes the propagators by (2.6) into slices, which are estimated by (2.10)--(2.12); this is (i). Graphs with an odd number of external vector legs vanish, since $\tr q^{2n+1}=0$.
\item Pp.~435--436: the last two terms of (3.7) cancel those of (3.6) identically. (3.9)--(3.14) are Taylor's formula followed by (2.10)--(2.11) for a line of index $j''$.
\item P.~437: the sum over orderings produces $G^\eta_{j''}(0)$ and the photon propagator. (3.16) uses $G=C^\xi+G(1-m_{j''}^2-a_{j''}P_{j''})C^\xi$, together with pointwise bounds on $C^\xi$, on the propagators and on their derivatives; this gives (ii) and (iii).
\item P.~438: (3.19) is an exact subtraction. The first graph of (3.21) is controlled by the boundedness of $\eta G_k(x,x)$, which in scale-$k$ units is $\xi G_f(y,y)$; this is (iv).
\item P.~439: (3.23) is treated by the method of (3.16), with both propagators replaced; this gives (v) and a convergent remainder, which is part of (vii). (3.24) contains $C^\xi$ only.
\item Pp.~440--441: the last bracket of (3.26) is estimated as in (3.13)--(3.14). The replacement of $G_{j_0}(0)$ is (iii). (3.27)--(3.28) contain $C^\xi$ only.
\item P.~442: (3.29) vanishes by the Ward--Takahashi identity (2.26) for $C^\xi$, whose mass is $1$. The graphs (3.31) have two external vector legs and order $\ge4$, so inside a graph of order $\le4$ with two external scalar legs they would give order $\ge6$; hence they do not occur. (2.21a) vanishes, (2.21e) carries an extra factor $\eta$, and (2.21d), (2.21f) cancel against (2.21g).
\item P.~443: (3.33)--(3.35) transport legs, and the result is estimated by line bounds.
\item P.~444: (2.21b), (2.21g) and (2.18) vanish exactly. The replacements of the propagators give (vi) and convergent remainders, while the $C^\xi$-versions (3.37), ``derivative of a constant'', and (3.38) vanish.
\item P.~445: the induction on the last divergent subgraph, and the cancellation of vacuum graphs, are combinatorial.\qedhere
\end{itemize}
\end{proof}

\begin{proposition}[independence of $C_m$ from the masses]\label{app:G1:I7-Cm}
Assume the hypotheses of Lemma~\ref{app:G1:I7-lines}; in particular (E.m2) $m_2^2\eps_0^2\le1$ and $\mu_*\eps_0\le1$. Then the constant $C_m$ of Theorem~\ref{app:G1:I7-I3a} can be chosen depending only on $\alpha_0,\nbar,d,L,a,M_b$. It is the same for all $M\in(0,m_2]$ and all such $\mu_*$, and it does not depend on $\ell_T$.
\end{proposition}
\begin{proof}
By Theorem~\ref{app:G1:I7-I3a}, $C_m$ is the constant $O(1)$ of \cite[(1.33)]{Ba3} for the expressions of Lemma~\ref{app:G1:VSR-C} times the number of graphs, and this $O(1)$ contains the constant $C_A$ of Lemma~\ref{app:G1:VSR-A}. We treat the three quantities in turn:
\begin{itemize}
\item the number of graphs depends only on $\nbar$ and $d$;
\item the constant $C_A$ contains no mass; its rate $\delta_1$ is free of $M$ and $\mu_*$ by Lemma~\ref{app:G1:I7-lines};
\item the constant $O(1)$ of \cite[(1.33)]{Ba3} for the expressions of Lemma~\ref{app:G1:VSR-C}.
\end{itemize}
For $j_0=0$ every quantity is a single slice, which is bounded by Lemma~\ref{app:G1:I7-lines}. For $j_0\ge1$ we have $M^2t_{j_0}^2\le1$ and $\mu_*^2t_{j_0}^2\le1$, since $t_{j_0}\le\eps_0$. By Lemma~\ref{app:G1:I7-masses}, the masses enter $O(1)$ only through (i)--(vii), and each of these is bounded uniformly in the masses.
\begin{itemize}
\item (i) is uniform by Lemma~\ref{app:G1:I7-lines}.
\item (ii) follows from Proposition~\ref{app:K:L11-1}(a).
\item (iii) follows from Lemmas~\ref{app:G1:I7-W} and~\ref{app:G1:I7-W441}.
\item (iv) and (v) follow from Proposition~\ref{app:K:L11-1}(b), items (C1)--(C3). The coefficient of \cite[(3.23)]{Ba3} has the form (C2). The remainder of \cite[(3.23)]{Ba3} that comes from $g'(x')-g'(x)$ has the form (C3). The scalar charge satisfies $\abs q\le1$ \cite[p.~426]{Ba3}.
\item (vi) for (2.19) and (2.20). Both terms of \cite[(3.37)]{Ba3} have the form (C2). Each term in the first line of \cite[(3.38)]{Ba3} has one twice-differenced factor and one undifferenced factor, and these two share a summed variable that carries one of the differences. The replacement is made line by line, so the same holds for its $C^\xi$-version. On the torus, $\xi^3\sum_z(\partial_\nu f)g=\xi^3\sum_zf(\partial^*_\nu g)$, and the kernel of $K\partial^*_\nu$ is the forward difference of $K$ in its second variable. So one summation by parts brings both versions to the form $T$ of (C4). Items (C2) and (C4) of Proposition~\ref{app:K:L11-1}(b) then apply.
\item (vi) for (2.21c) follows from Lemma~\ref{app:K:L11-2}(a) together with Proposition~\ref{app:K:L11-1}(c), with $K_1=K_2=G_s$ and $K_3=G_p$.
\item (vii): by Lemma~\ref{app:K:L11-2}(b), these expressions are the differences between the coefficients in (iii), (v), (vi) and their $C^\xi$-versions, and they enter only through bounds on their absolute values. The $C^\xi$-versions contain no mass, so these differences are bounded by the bounds just given plus mass-free quantities.
\end{itemize}
All constants depend only on $(a,L)$ and on the finitely many graphs. They are uniform in $\xi$, $j_0$ and $\ell\ge L$, and hence in $\ell_T$.
\end{proof}

\begin{proposition}[the mass reading of {\cite[p.~589]{Ba2}}]\label{app:G1:I7}
Assume (E.m2): $m_2^2\eps_0^2\le1$, and (E.$\Phi$): $\lambda^{-1/4}m_2^2\eps_0^{7/4}\le a_\flat/C_\varphi$. Equivalently, $\eps_0\le\eps_{(\mathrm{E.md})}:=\min\{1/m_2,(a_\flat\lambda^{1/4}/(C_\varphi m_2^2))^{4/7}\}$ (stage 10a of Appendix~\ref{app:B}). Assume moreover $\mu_*\eps_0\le1$, which holds under (E.dom), and $\eps_0\le1$. Then at every $M\in(0,m_2]$, in particular at $M=m_2$ and $M=m_1$, every estimate of \cite{Ba1,Ba2,Ba3,Ba4} that uses an upper bound on the mass holds with the following constants:
\begin{itemize}
\item $\gamma_0^{\rm II}=\min\{1,\tfrac12\min\{\gamma_0^{(0)},a_\flat/(a_\flat+1)\}\}$ in \cite[(3.26)]{Ba2}, with $\gamma_0^{(0)}$ as in Lemma~\ref{app:G1:M4B};
\item $\gamma_1=a_\flat/(a_\flat+1)=3/7$ in the bound $G''_k\ge\gamma_1\nu_k$ of \cite[(3.11)]{Ba2};
\item the constant of \cite[(2.65)]{Ba2} increased by $1$;
\item $K_2^\natural$ of Corollary~\ref{app:G1:I7-U6} in \cite[(2.113)]{Ba2}, and $2Ld+K_2^\natural$ in \cite[(2.114)]{Ba2};
\item $K_{3.16}$ of Lemma~\ref{app:G1:I7-W} for \cite[(3.16), p.~437]{Ba3}, and $K_{441}$ of Lemma~\ref{app:G1:I7-W441} for the coefficients of \cite[p.~441]{Ba3};
\item $C_m$ of Proposition~\ref{app:G1:I7-Cm} in the bound of Theorem~\ref{app:G1:I7-I3a} on $\delta m^2_{K,k}$.
\end{itemize}
None of these depends on $\eps$, $\eps_0$, $k$, $\ell_T$, $M$, or the deformation parameter $u$ of Subsection~\ref{ssec:T6-statement}. This proposition is the mass ingredient listed in Table~\ref{tab:T6-inputs}.
\end{proposition}
\begin{proof}
Write $\nu_k=M^2t_k^2$. Since $t_k\le\eps_0$, (E.m2) gives $\nu_k\le1$ for all $k\le K$. The estimates of \cite{Ba1,Ba2,Ba3,Ba4} in which the mass occurs are the following ten, which we call U1--U10.
\begin{itemize}
\item[(U1)] \cite[Theorem (1.9)--(1.12)]{Ba4} and \cite[(2.58)]{Ba2}. No mass bound is needed (Lemma~\ref{app:G1:M4A}(i)).
\item[(U2)] \cite[(2.33)--(2.35)]{Ba1} $=$ \cite[(1.15)--(1.20)]{Ba4}. No mass bound is needed (Lemma~\ref{app:G1:M4A}(ii)).
\item[(U3)] \cite[Prop.~3.1, (3.26)]{Ba2}, used through (3.29). This is Lemma~\ref{app:G1:M4B} with $c_m=1$.
\item[(U4)] \cite[(3.11)]{Ba2}. $G''_k$ is a Neumann $\Delta^{(k)}$-form plus non-negative terms. So the proof of Lemma~\ref{app:G1:M4B} gives $G''_k\ge a_k\nu_k/(a_k+\nu_k)\ge\tfrac37\nu_k$. The uses (3.17)--(3.20) need only $\norm{G''^{-1}_k}\le1/\nu_k+1/a_k$.
\item[(U5)] \cite[Lemma~2.4, (2.65), (2.76), (2.78)]{Ba2}. This is Lemma~\ref{app:G1:M4D} under (E.$\Phi$).
\item[(U6)] \cite[Lemma~2.7, (2.113)--(2.114)]{Ba2}. This is Lemma~\ref{app:G1:I7-L27} with Corollary~\ref{app:G1:I7-U6}, under (E.$\Phi$) and $\eps_0\le1$.
\item[(U7)] \cite[(3.16), p.~437]{Ba3} is Lemma~\ref{app:G1:I7-W}, and \cite[p.~441]{Ba3} is Lemma~\ref{app:G1:I7-W441}; both need only $\nu_j\le1$. The counterterm bound is Theorem~\ref{app:G1:I7-I3a} with Proposition~\ref{app:G1:I7-Cm}.
\item[(U8)] \cite[(1.1), (2.1), (2.11)--(2.13), (3.21)--(3.23), (3.37)--(3.38)]{Ba2} and \cite[(1.4), (1.19)--(1.22)]{Ba3}. Each of these steps is one of the following:
\begin{itemize}
\item a positive mass term dropped in an upper bound;
\item an exact Gaussian identity;
\item the diamagnetic inequality, which holds for every $M>0$;
\item a Gaussian measure whose propagator obeys (U1).
\end{itemize}
No mass bound is needed.
\item[(U9)] The counterterm graphs of \cite[p.~418]{Ba3}. For $\delta m^2_{K,k}$ they are covered by (U7). The final number $\delta m^2_{\rm fin}$ is bounded in Proposition~\ref{app:G1:I3b}, whose constants carry their own dependence on $M$ and $\mu_*$.
\item[(U10)] \cite[(3.15)--(3.20)]{Ba2} and the other bounds that use the mass from below. These are not upper-bound uses: they bind at $m_1>0$ and are separate conditions.
\end{itemize}
Three further occurrences need only $\nu_k\le1$:
\begin{itemize}
\item \cite[(3.52), p.~621]{Ba1}: in that paper all fields are cut at $\abs\cdot\le p$, so $\nu_k\abs{\psi^{(k+1)}}\le p$;
\item \cite[(A.1), (A.7)--(A.8), p.~676]{K1}: the target $Cp\lambda(t)^{-1/4}$ has the size of the field bound of (A.1);
\item \cite[p.~441]{Ba3}, which is Lemma~\ref{app:G1:I7-W441} in (U7).
\end{itemize}
The scalar operators carry no weight in $u$, so none of these bounds depends on $u$. Finally, every left side ($\nu_k$, and $\lambda^{-1/4}M^2t^{7/4}$) is increasing in $M$, and every constant listed is free of $M$. So the bounds at $m_2$ give those at every $M<m_2$, in particular at $m_1$. The lower bounds on the mass are strongest at $m_1$; they are separate conditions.
\end{proof}

\paragraph{Bounds on \texorpdfstring{$\delta m^2_{\rm fin}$}{dm2fin} and \texorpdfstring{$E_{\rm fin}$}{Efin} uniform in the torus.}\label{app:G1:ssec-I8}
Let $f\in\{s,v\}$ label the scalar and photon fields, with masses $m_s=M\in\{m_1,m_2\}$ and $m_v=\mu_*$. The number $\delta m^2_{\rm fin}:=\delta m^2-\delta m^2_K$ is the coefficient of \cite[(1.26)]{Ba3}, and $E_{\rm fin}$, $E_K$ are the vacuum pieces of Definition~\ref{app:G1:vacgrade}, so that $E_1=E_{\rm fin}+E_K$ (Lemma~\ref{app:G1:vacgrade-id}); these are the final c-numbers of \cite[(1.26), (1.30)]{Ba3}, in the grading of Definition~\ref{app:G1:vacgrade}. They are fixed numbers, so by Lemma~\ref{app:G1:vacgrade-id}(i) any exact decomposition of the lines above $\eps_0$ may be used to bound them. For $k>K$ put $t_k:=L^{k-K}\eps_0$. Let
\[
\bar\ell:=\min(1/M,1/\mu_*,L),\qquad \bar K:=\max\{k:t_k\le\bar\ell\}.
\]
For both fields we use the exact identity $C_f^\eps=G^\eps_{K,f}(0)+\sum_{j=K}^{\bar K-1}H_{j,f}+R_{\bar K,f}$, where the $H_{j,f}$ are King's fluctuation covariances and $R_{\bar K,f}:=C_f-G_{\bar K,f}(0)$.

The common top index, the lemma on intermediate subgraph scales, the top-cluster lemma and the two bounds are in Subsection~\ref{app:G1:part-D3}, at every Higgs mass $M\in(0,m_2]$: properties (P1)--(P7) there are the facts about the common top index (the inclusion $\bar K\le\min_fK_f$ is (P7)), Lemma~\ref{app:G1:common} is the lemma on intermediate scales, Lemma~\ref{app:G1:TB} the top-cluster lemma, and Propositions~\ref{app:G1:I3b} and \ref{app:G1:Efin} the bounds, with constants that keep $M_b$. The inputs (VS), (Rem) and (S6) are proved in Subsection~\ref{app:G1:part-TVSR}. The dependence on $u$ is treated next.

\emph{The photon lines at $u\ne0$.} Let $M\in(0,m_2]$ and assume \eqref{app:G1:eq-hyp}. For every index $j$ put $c_j:=\kappa\mu_*^2L^2t_j^2$. In the slice model at $(u,M)$ every vector line carries $C_B(u)=C_{\mu_*}+u\Gamma^B$ (see \eqref{eq:T6-GammaB}) and every scalar line carries $C_s$; the mass counterterm is the vertex \cite[(1.7)]{Ba3} with the number $\delta m^2_{\rm ref}$ as its coefficient, built from $C_{\mu_*}$-lines and the propagators at $u=0$. $E_1(u)$ is the vacuum coefficient of \cite[(1.26), (1.30)]{Ba3} in this model, $E_K(u)$ is the same expression with $G_K(u):=G_K+u\Gamma^B$ for the vector lines and $G_{K,s}(0)$ for the scalar lines, and
\begin{equation}\label{app:G1:eq-Efinu}
E_{\rm fin}(u):=E_1(u)-E_K(u),
\end{equation}
all graded as in Definition~\ref{app:G1:vacgrade}, with every $\delta m^2$ kept as a number. A vector line of index $i<K$ is $H_i+u\delta_iH_i$; we call $H_i$ its $C$-part and $u\delta_iH_i$ its $D$-part (a \emph{$D$-line}).

\begin{lemma}[colour expansion of $E_{\rm fin}(u)$]\label{app:G1:I8-colour}
For every $u\in\R$, $E_{\rm fin}(u)=\sum_{m=0}^{N_\ell}u^mE^{(m)}_{\rm fin}$ with $E^{(0)}_{\rm fin}=E_{\rm fin}(0)$, where $u^mE^{(m)}_{\rm fin}$ is the sum of the values of the coloured graphs with exactly $m$ $D$-lines. Each such graph has at least one line from $\{C_{\mu_*}-G_K,\ C_s-G_{K,s}(0)\}$; these lines do not depend on $u$, and all $D$-lines have index $<K$. Moreover $\delta m^2_{\rm ref}$, $\delta m^2_K$ and $\delta m^2_{\rm fin}$ do not depend on $u$.
\end{lemma}
\begin{proof}
Graph values are multilinear in the line kernels. Since $C_B(u)-G_K(u)=C_{\mu_*}-G_K$, the graphs of $E_1(u)$ that are not in $E_K(u)$ are exactly those with at least one line $C_{\mu_*}-G_K$ or $C_s-G_{K,s}(0)$. Write $G_K(u)=\sum_{i<K}(1+u\delta_i)H_i$ and expand every vector line of index $i<K$ into its $C$- and $D$-parts. A graph with $m$ $D$-lines carries $u^m$, and $m\le N_\ell$. Finally, $\delta m^2_{\rm ref}$ and its pieces $\delta m^2_K$, $\delta m^2_{\rm fin}$ are built from $C_{\mu_*}$-lines and from the $u=0$ propagators, so they carry no $u$.
\end{proof}

\begin{lemma}[the $D$-line in scale-$j$ units]\label{app:G1:I8-Dfac}
Let $i<K\le j\le\bar K$ and $\eta_j:=L^{-j}$, so that $L^i\eta_j=t_i/t_j\le L^{-1}$. Then
\[
\delta_i=c_j\,(L^i\eta_j)^2,\qquad c_j\le c_{\bar K}\le\kappa L^2(\mu_*\bar\ell)^2\le\kappa L^2 .
\]
The same identity holds for $i<j\le K$, and then $c_j\le c_K=\kappa\mu_*^2L^2\eps_0^2$. If $\eps_0$ is a power of $L$, then $c_{\bar K}$ does not depend on $\eps$, $\eps_0$ or $\ell_T$.
\end{lemma}
\begin{proof}
$\delta_i=\kappa\mu_*^2t_{i+1}^2=\kappa\mu_*^2L^2t_j^2(t_i/t_j)^2$, and $j\mapsto c_j$ is increasing. By (P1), $t_{\bar K}\le\bar\ell\le1/\mu_*$. If $\eps_0$ is a power of $L$, the $t_k$ run through powers of $L$, and since $\eps_0\le\bar\ell$, $t_{\bar K}$ is the largest power of $L$ that is $\le\bar\ell$.
\end{proof}

\begin{proposition}[the final c-numbers]\label{app:G1:I8}
Assume $\eps_0\le\min(1/m_2,1/\mu_*,L)$, with $\eps_0$ a power of $L$. Then at both anchor masses $M\in\{m_1,m_2\}$ the bounds of Propositions~\ref{app:G1:I3b} and \ref{app:G1:Efin} hold, uniformly in $\eps=2^{-K}\eps_0$, $x\in T_\eps$ and the tori $T_j$, and the bound on $\delta m^2_{\rm fin}$ does not depend on $u$. For every $u\in\R$, $\abs{E_{\rm fin}(u)}\le\bar\sigma_N(u)\hat B$, where $\hat B$ is as in Proposition~\ref{app:K:L04-4} and
\[
\bar\sigma_N(u):=\big(1+\abs u\,\kappa\mu_*^2L^2t_{\bar K}^2\big)^{N_\ell}\le\big(1+\abs u\,\kappa L^2\big)^{N_\ell};\qquad
\bar\sigma_N(1)\le(1+\kappa L^2)^{N_\ell},\quad \bar\sigma_N(-\eta)\le(1+\eta\kappa L^2)^{N_\ell}.
\]
These numbers do not depend on $\eps$, $\eps_0$ or $\ell_T$. The torus enters only through $Z_s$ and $Z_v$, which are non-increasing in $\ell_T$; hence the supremum over $\ell_T\ge\ell_0$ is at most the value at $\ell_0^\natural$, a number fixed at stage 0 of Appendix~\ref{app:B}, before $\eps_0$. In particular the one-loop constant of the bound on $\delta m^2_{\rm fin}$ depends on $\ell_0^\natural$.
\end{proposition}
\begin{proof}
For $M\le m_2$, $\bar\ell(M)\ge\bar\ell(m_2)$, so the hypothesis gives the first half of \eqref{app:G1:eq-hyp} at $M=m_1$ and $M=m_2$; the second half, $\ell_T\ge L^3$, holds on the tori of the paper, since $\ell_T\ge\ell_0\ge\ell_0^\natural\ge8=L^3$ (Remark~\ref{app:G1:N9-supply}(i)). Propositions~\ref{app:G1:I3b} and \ref{app:G1:Efin} then apply at both masses, and Proposition~\ref{app:K:L04-4} gives the bound at $u\ne0$. The $u$-independence of $\delta m^2_{\rm fin}$ is Lemma~\ref{app:G1:I8-colour}, the bounds on $\bar\sigma_N$ are Lemma~\ref{app:G1:I8-Dfac} (note $c_{\bar K}=\kappa\mu_*^2L^2t_{\bar K}^2$), and the monotonicity in $\ell_T$ is (P6), since $t_{\bar K}$, and with it $\bar\sigma_N$, does not depend on $\ell_T$.
\end{proof}

\begin{remark}[the factor $c_K$]\label{app:G1:I8-cK}
The argument does not give the smaller factor $(1+\abs u\,c_K)^{N_\ell}$. Bounding a $D$-line by the number $\abs{u\delta_i}\le\abs u\,c_K$ leaves the kernel $(t_i/\eps_0)^2H_i$, which has dimension $-d+4$ only relative to the scale $\eps_0$. A subgraph that contains a $D$-line, has no counterterm and has lowest external index $j_0>K$ needs the factor $(t_i/t_{j_0})^2$, which costs $c_{j_0}=c_K(t_{j_0}/\eps_0)^2$. When every index is $\le K$, as for the pieces of the flow below $\eps_0$, $c_{j_0}\le c_K$, and the factor $(1+\abs u\,c_K)^{N_\ell}$ is valid there by Theorem~\ref{app:G1:WP-WP1}(c),(d).
\end{remark}

\begin{lemma}[the hypothesis of Proposition~\ref{app:G1:I8} is a condition of Appendix~\ref{app:B}]\label{app:G1:N9prem}
Every $\eps_0$ admissible in the order of choices of Appendix~\ref{app:B} satisfies $\eps_0\le\min(1/m_2,1/\mu_*,L)$. At $L=2$ it is moreover a power of $L$.
\end{lemma}
\begin{proof}
By stage 10a, $\eps_0\le\eps_1\le1<L$ and $\eps_0\le\eps_{(\mathrm{E.md})}\le1/m_2$; (E.m2) reads $m_2^2\eps_0^2\le1$. The bound (E.dom), $\mu_*\eps_0\le0.2$ (stages 0 and 10a), gives $\eps_0\le0.2/\mu_*<1/\mu_*$. Since $m_1<m_2$, also $\eps_0\le\bar\ell$ at $M=m_1$. Finally $\eps_0=2^{-n_0}$ by (T.div).
\end{proof}

\paragraph{The vacuum classes at the last scale.}\label{app:G1:ssec-I9}
\cite[p.~422]{Ba3} treats, for the mass counterterms, the terms that pass from scale $k$ to $k+1$ by cutting the index-$k$ line, and adds: ``The corresponding considerations for vacuum energy counterterms are even simpler and we omit them.'' We write out the omitted vacuum case, in the grading of Definition~\ref{app:G1:vacgrade}; Remark~\ref{app:G1:vacgrade-rem} shows why the coefficients are not opened. Write $t_k:=L^k\eps$, $e(t)^2=e^2t$, $\lambda(t)=\lambda t$, $o:=\alpha/2+\beta$, and $P(t):=\sum_{2\le\alpha+2\beta\le2\nbar}\abs{e(t)}^\alpha\lambda(t)^\beta$. By Definition~\ref{app:G1:vacgrade}, $\cV_k$ is the sum of the terms of the bare vacuum graphs of orders $1\le\alpha+\beta\le\nbar$ in which every line has index $\le k$ and at least one has index $k$. In such a term every coefficient is split at index $k+1$ by \cite[(1.26), (1.28), (1.29)]{Ba3}, $\delta m^2=\delta m^2_{\rm fin}+\delta m^2_{K,k+1}+\delta m^2_{k+1}$. Following \cite[p.~423]{Ba3}, a vertex carrying $\delta m^2_{\rm fin}$ or $\delta m^2_{K,k+1}$ is a \emph{finite} vertex of degree $0$, localised by the unit cube of its position; a vertex carrying the part $\delta m^2_{G_0}$ of $\delta m^2_{k+1}$ belonging to a two-scalar-leg graph $G_0$ has degree $D(G_0)$. Localisations are those of \cite[p.~420]{Ba3}. In scale-$k$ units, $\delta m^2_{K,k}t_k^2=O(t_k)$ \cite[p.~418]{Ba3}.

\begin{lemma}[exact cut]\label{app:G1:cut}
Fix a total order of the lines of every bare vacuum graph, and decompose the index sums by Ba{\l}aban's orderings $\bar l=(l(1),\dots,l(m))$ \cite[(3.4), p.~424]{Ba3}, with ties broken by that order. Then every term of $\cV_k$ belongs to exactly one ordering, and: (i) the last line $l(m)$ has index $k$; (ii) the counterterm operations active in $\bar l$ act on divergent components of $G_1\subset\dots\subset G_{m-1}$, none of which contains $l(m)$, so every member of the vacuum class contributing to $\bar l$ contains $l(m)$ as a line of its bare graph. Cutting $l(m)$, i.e.\ treating its index-$k$ kernel $K_k(x,x')$ as a two-variable external field \cite[p.~419]{Ba3}, maps the class term by term, with the same coefficients, onto two-leg classes of $G\setminus l(m)$, in which every line has index $\le k$ and every counterterm vertex carries $\delta m^2_{k+1}$ or a finite coefficient.
\end{lemma}
\begin{proof}
(i) is the definition of $J(\bar l)$ \cite[p.~424]{Ba3}, since all indices are $\le k$ and one equals $k$. (ii) By construction $l(m)\notin G_i$ for $i<m$, and $G_m=G$. A counterterm operation replaces a divergent component $S\subset G_i$ by the vertex $\delta m^2_S$, or by the other members of its class \cite[pp.~429--432]{Ba3}, and does not touch the lines outside $S$; hence $l(m)$ remains a line. The lines of $S$ have indices $\le k$, so $\delta m^2_S$ is a part of $\delta m^2_{k+1}$. The parts of a coefficient with a line of index $>k$ form $\delta m^2_{K,k+1}$, a finite vertex.
\end{proof}

\begin{lemma}[cut classes]\label{app:G1:cutclasses}
Let $G$ be a bare vacuum graph, with degree $D$ as in \cite[(2.1)--(2.3), pp.~422--423]{Ba3}, let $\ell$ be one of its lines, and let $n_\partial(\ell)\in\{0,1,2\}$ be the number of differentiated legs among its two endpoints.
\begin{enumerate}
\item[(a)] If $G\setminus\ell$ is connected, then $D(G\setminus\ell)=D(G)+1+n_\partial(\ell)$.
\item[(b)] If $G$ has no finite vertex, then $o(G)\in\Z_{\ge1}$ and $D(G)=o(G)-3$. Otherwise $D(G)\ge-1$, with equality only if $G$ is one finite vertex closed by one line.
\item[(c)] If $G\setminus\ell$ is disconnected, then $\ell$ is a photon line and each component has exactly one external leg, a vector leg. Both components vanish \cite[p.~434]{Ba3}.
\item[(d)] Every component of $G\setminus\ell$ with $D\le0$ is one of the following. For a scalar $\ell$: C1 ($o=1$, $n_\partial=0$, $D=-1$), C2 ($o=2$, $n_\partial=0$, $D=0$), C1$'$ ($o=1$, $n_\partial=1$, $D=0$), or a bare finite vertex with two legs, when $G$ is that vertex closed by $\ell$. For a photon $\ell$: C3 ($o=1$, $D=-1$), C4 ($o=2$, $D=0$), or a vanishing pair as in (c). All other components have $D\ge1$; this includes every $G$ with $o\ge3$.
\end{enumerate}
\end{lemma}
\begin{proof}
(a) In \cite[(2.1)]{Ba3} the two endpoint legs of $\ell$ change from internal, of dimension $-(d-2)/2=-\frac12$, to external, of dimension $0$, and a differentiation acting on such a leg changes from $-1$ to $0$. The numbers of factors $\eta$ and the counterterm degrees in (2.3) are unchanged, and $-d$ is subtracted once, since $G\setminus\ell$ is connected.
(b) By \cite[p.~423, (i)--(iv)]{Ba3} with $d=3$: $D(v)=1=o(v)$ for (1.6), and $D(v)=n/2=o(v)$ for (1.8) and (1.10) with $n$ vector legs and no background leg. For a vertex carrying $\delta m^2_{G_0}$, $D(v)+\deg=2+D(G_0)$; by induction on the order, $D(G_0)=o(G_0)-2$ for the two-scalar-leg graph $G_0$, which has no differentiated external leg, because $\delta m^2$ is built from the underived part of \cite[(1.21)]{Ba3}. Hence $D(v)+\deg=o(v)$, and (2.3) gives $D(G)=o(G)-3$. Every vector line uses two vector legs, so $\alpha$ is even and $o\in\Z$; since $o(v)\ge\frac12$, $o\ge1$. A finite vertex has $D(v)=2$ and degree $0$, and every other vertex has $D(v)\ge\frac12$; so $D(G)\ge2-3$, with equality only for one finite vertex closed by one line.
(c) Every vertex has an even number of scalar legs, so a component of $G\setminus\ell$ has an even number of external scalar legs, and a scalar line is never a bridge. If $\ell$ is a photon bridge, each side carries one external vector leg and no external scalar leg; by \cite[p.~434]{Ba3} each side contains a closed scalar strand with an odd power of $q$, and $\tr q^{2n+1}=0$; this holds after the gauge transformation of \cite[p.~433]{Ba3}.
(d) By (a) and (b), a connected $G\setminus\ell$ without finite vertex has $D=o-2+n_\partial$, with $n_\partial=0$ for photon lines. With a finite vertex, $D(G)=o'-1$, where $o'\in\Z_{\ge0}$ is the order of the other vertices, so $D(G\setminus\ell)=o'+n_\partial\ge0$, with equality only for $o'=0$, $n_\partial=0$. Listing $D\le0$ over $o\in\Z_{\ge1}$ and $n_\partial\in\{0,1,2\}$ gives the stated classes.
\end{proof}

\emph{Cut classes.} By Lemma~\ref{app:G1:cutclasses}, a scalar cut gives the classes C1, C2 (with their $\delta m^2$-partners, Lemma~\ref{app:G1:I9-partners}, except the one-particle-reducible chains of C2, which have no partner and are treated by the product mechanism of Step~3 of the proof of Lemma~\ref{app:G1:I145-perloc}) or C1$'$, or, when $G$ is a finite vertex closed by the cut line, that bare vertex with two legs; a photon cut gives the Ward classes C3, C4, or vanishing pairs. The photon cut needs the Ward class, so the vacuum case is not simpler than the mass case.

\begin{lemma}[cuts inside opened counterterm vertices]\label{app:G1:cutopen}
Let $G$ be as in Lemma~\ref{app:G1:cutclasses}, without finite vertices, and let $v$ be a vertex \cite[(1.7)]{Ba3} of $G$ carrying the counterterm $\delta m^2_{G_0}$ of a graph $G_0$ of order $o_0\in\{1,2\}$. Let $\ell_1,\ell_2$ be the lines of $G$ at $v$ and $u_1,u_2$ the vertices of $G_0$ carrying its external legs. Let $G^T$ be obtained from $G$ by replacing $v$ by $G_0$, with $\ell_i$ attached at $u_i$, and $G^v$ from $G^T$ by moving the end of $\ell_2$ from $u_2$ to $u_1$ (the representation of \cite[p.~417]{Ba3}). Degrees are those of \cite[(2.1)--(2.3)]{Ba3}. Then:
(a) $G_0$ is one-particle irreducible, and neither of its external legs is differentiated;
(b) $D(G^v)=D(G^T)=D(G)=o(G)-3$;
(c) for every line $\ell$ of $G_0$, $G^v\setminus\ell$ and $G^T\setminus\ell$ are connected, and $D(G^v\setminus\ell)=D(G^T\setminus\ell)=o(G)-2+n_\partial(\ell)$;
(d) hence every such cut with $D\le0$ belongs to C1, C2, C1$'$, C3 or C4 with $o=o(G)\le2$; no photon bridge and no bare vertex occurs.
\end{lemma}
\begin{proof}
(a) By \cite[(1.21), p.~416]{Ba3}, $\Sigma^\eps,\Sigma^\eps_1,\Sigma^\eps_2$ are given by amputated one-particle-irreducible graphs, and $\delta m^2$ is defined from $\Sigma^\eps$, the part without differentiated external legs \cite[p.~417]{Ba3}.
(b) Each internal leg contributes $-\frac12$ to (2.1), so $D(H)=\sum_w\eta(w)-I(H)-\partial(H)-d+\sum\deg$, with $\eta(w)$ the number of factors $\eta$ at $w$, $I$ the number of lines, $\partial$ the number of differentiations on lines and the last sum over unopened counterterm vertices. In $G^T$ the vertex $v$ ($\eta(v)=d$, $\deg v=D(G_0)$) is replaced by the vertices of $G_0$; since $I(G^T)=I(G)+I(G_0)$ and $\partial(G^T)=\partial(G)+\partial(G_0)$ by (a), $D(G^T)=D(G)$. Moving an undifferentiated leg-end changes no term, so $D(G^v)=D(G^T)$. Finally $D(G)=o(G)-3$ by Lemma~\ref{app:G1:cutclasses}(b).
(c) $G_0\setminus\ell$ is connected by (a) and contains $u_1,u_2$, and every path of $G$ through $v$ passes through it. Removing $\ell$ lowers $I$ by $1$ and $\partial$ by $n_\partial(\ell)$.
(d) Combine (c) with the list in the proof of Lemma~\ref{app:G1:cutclasses}(d).
\end{proof}

\begin{remark}\label{app:G1:cutopen-rem}
In the grading of Definition~\ref{app:G1:vacgrade} the cut line $l(m)$ of Lemma~\ref{app:G1:cut} is a line of the bare graph, so it never lies inside a coefficient. Let $G$ have no finite vertex. If the coefficient of $v$ were opened and a line $\ell$ of $G_0$ were cut, the cut member would factor as $\sum_xR_v(x)c_{G_0,\ell}(x)$, where $R_v$ is the value of $G$ with coefficient $1$ at $v$ and $c_{G_0,\ell}$ is the coefficient graph with $\ell$ as external field. The vertex of $R_v$ at $v$ has degree $0$ in place of $D(G_0)=o_0-2$ (proof of Lemma~\ref{app:G1:cutclasses}(b)), so $R_v$ has degree $D(G)-D(G_0)=o(G)-o_0-1$. If $o(G)-o_0\le1$, which by Lemma~\ref{app:G1:cutopen}(d) is the case for every such cut with $D\le0$, the factor $R_v$ has degree $-1$ or $0$ and no partner in the cut family. This is the mechanism of Remark~\ref{app:G1:vacgrade-rem}.
\end{remark}

\begin{lemma}[the partners are present]\label{app:G1:I9-partners}
In the cut classes of Lemma~\ref{app:G1:cut}, every one-particle-irreducible component of type C1 or C2 occurs together with its mass counterterm, with coefficient in $\delta m^2_{k+1}$, and every component of type C3 or C4 occurs together with all connected two-vector-leg graphs of the same order whose lines have index $\le k$.
\end{lemma}
\begin{proof}
Let $H$ be a connected two-leg graph of order $o\le\nbar$ with lines of index $\le k$, and $\ell$ a line of index $k$ of the kind of its legs. Closing $H$ by $\ell$ gives a bare vacuum graph whose term lies in $\cV_k$, and cutting $\ell$ gives back $H$. If $H$ is a one-particle-irreducible two-scalar-leg graph, the vertex $\delta m^2_H$ (a part of $\delta m^2_{k+1}$) closed by $\ell$ is also a bare vacuum graph in $\cV_k$, and its cut is the counterterm member of $H$ in the representation of \cite[p.~417]{Ba3}. For two vector legs this applies to every connected two-vector-leg graph of the order, which is the class of \cite[p.~431]{Ba3}.
\end{proof}

\begin{lemma}[localisation sum]\label{app:G1:VS}
Let $F=\sum_\Lambda F(\Lambda)$ be a sum over localisations $\Lambda=(\square(v))_{v\in G}$, i.e.\ $N_s$ cubes of side $1$ or $2$ of the scale-$k$ lattice, each a union of unit cubes (the localisations of \cite[p.~420]{Ba3}). Assume $\abs{F(\Lambda)}\le A\,\mathrm{e}^{-\delta d(\Lambda)}$, where $d(\Lambda)$ is the length of a shortest tree through the cubes. Let $W$ be a union of unit blocks and let $F^{\rm bad}$ be the sum of $F(\Lambda)$ over the $\Lambda$ that have a cube meeting $W$ and a cube not contained in $W$. Then for every $r>0$
\[
\abs{F^{\rm bad}}\le 2^dA\,N_s\,S\Big[\#\{\square\subset W:\dist(\square,W^c)\le r\}+\mathrm{e}^{\delta(c_d+1)/2}\mathrm{e}^{-\delta r/2}\,\#\{\square\subset W:\dist(\square,W^c)>r\}\Big],
\]
where $S:=\sum_{D\ge0}(c''_d(D+2)^d)^{N_s}\mathrm{e}^{-\delta D/2}$, $c_d:=4\sqrt d$, and $c''_d$ is a lattice constant.
\end{lemma}
\begin{proof}
Attribute each bad $\Lambda$ to a unit cube $\square_0\subset W$ contained in the cube of the first vertex (in a fixed order) whose cube meets $W$; there are at most $2^dN_s$ choices of the vertex and of $\square_0$. Since $\Lambda$ also has a cube that is not contained in $W$, which meets $W^c$ and has diameter at most $2\sqrt d$, $d(\Lambda)\ge\dist(\square_0,W^c)-c_d$. With $\square_0$ fixed, every other cube lies within $d(\Lambda)+c_d$ of $\square_0$, so at most $(c''_d(D+2)^d)^{N_s}$ localisations have $d(\Lambda)\in[D,D+1)$, for $c''_d$ large enough. With $D_0:=(\dist(\square_0,W^c)-c_d-1)_+$ and $\mathrm{e}^{-\delta D}\le \mathrm{e}^{-\delta D_0/2}\mathrm{e}^{-\delta D/2}$ for $D\ge D_0$, the sum over $\Lambda$ with $\square_0=\square$ is at most $2^dAN_s\mathrm{e}^{-\delta D_0/2}S$. Use $\mathrm{e}^{-\delta D_0/2}\le1$ if $\dist\le r$, and $\mathrm{e}^{-\delta D_0/2}\le \mathrm{e}^{\delta(c_d+1)/2}\mathrm{e}^{-\delta r/2}$ if $\dist>r$. Sum over $\square\subset W$.
\end{proof}

\begin{lemma}[the layer and the deep part]\label{app:G1:I9-geo}
In scale-$k$ units let $W:=B^k(\Lambda_7^{(k)})$ and $r:=r(t_k)$. Every unit block $\square\subset W$ with $\dist(\square,W^c)\le r$ lies in $B^k(\Lambda_7^{(k)}\cap\Lambda_8^{(k)c})$, and every unit block of $B^k(\Lambda_8^{(k)})$ has $\dist(\square,W^c)>r$.
\end{lemma}
\begin{proof}
By \cite[(2.8), p.~558]{Ba2}, $\Lambda_8^{(k)c}$ is the union of the large blocks at distance $\le r$ from $\Lambda_7^{(k)c}$. If $\square\subset W$ and $\dist(\square,W^c)\le r$, the large block $b\supset\square$ has $\dist(b,W^c)\le r$, so $b\subset\Lambda_8^{(k)c}$. If $\square\subset B^k(\Lambda_8^{(k)})$, its large block is not within $r$ of $\Lambda_7^{(k)c}$.
\end{proof}

\begin{lemma}[degrees of open graphs]\label{app:G1:I9-deg}
Let $Q$ be a connected graph with $E=E_s+E_v\ge1$ external legs ($E_s$ scalar, $E_v$ vector), $n_\partial$ of them differentiated, $m_f$ finite vertices and $o'$ the order of its other vertices. Then
\begin{equation}\label{app:G1:I9-degf}
D(Q)=o'+2m_f-3+\tfrac12E+n_\partial .
\end{equation}
If $Q$ does not vanish and $D(Q)\le0$, then $Q$ is a single vertex; or the tree \cite[(2.4)]{Ba3} (two vertices \cite[(1.8)]{Ba3} with $n=1$ joined by a scalar line carrying both derivatives); or a member of C1, C1$'$, C2, C3, C4 of Lemma~\ref{app:G1:cutclasses}, or of the class C5 of the graphs with two scalar legs and one vector leg at $o=\frac32$. A connected graph with a finite vertex and $E\ge1$ has $D\ge0$, with equality only for a single finite vertex or for a vanishing scalar loop through one finite vertex and one vertex with a single, external, vector leg.
\end{lemma}
\begin{proof}
External legs and derivatives contribute $0$ in place of $-\frac12$ and $-1$ \cite[(2.1)--(2.3)]{Ba3}; summing the vertex degrees of the proof of Lemma~\ref{app:G1:Efin-deg} gives \eqref{app:G1:I9-degf}. Constraints: $E_s$ is even; $o\in\frac12E_v+\Z$; every non-finite vertex has $o(v)\ge\frac12$; a graph with an odd number of external vector legs and no external scalar leg vanishes, as in Lemma~\ref{app:G1:cutclasses}(c). If $m_f=0$, $D\le0$ means $o'+\frac12E+n_\partial\le3$: $E=1$ or $(E_s,E_v)=(0,3)$ vanish; $E=2$ gives $o'\in\{1,2\}$, $o'+n_\partial\le2$, i.e.\ C1, C1$'$, C2 ($E_s=2$) and C3, C4 ($E_v=2$); $(2,1)$ gives a single vertex ($o'=\frac12$) or C5 ($o'=\frac32$, $n_\partial=0$); $E=4$ needs $o'=1$, $n_\partial=0$: $(4,0)$ is a single $\lambda$ vertex (two $n=1$ vertices joined by a photon have $n_\partial\ge2$), $(2,2)$ is a single $n=2$ vertex or the tree \cite[(2.4)]{Ba3}, $(0,4)$ needs $o'\ge2$; $E\ge5$ needs $o'\le\frac12$, impossible. If $m_f\ge1$, $D\ge o'-1+\frac12E$: for $o'=0$, $D=2m_f-3+\frac12E\ge0$, with equality only for one vertex with $E=2$; for $o'=\frac12$ the only vector leg is external, $E=1$ gives $D=2m_f-2$, which is $0$ only for the vanishing loop, and $E\ge3$ gives $D\ge1$; for $o'\ge1$, $D\ge\frac12$.
\end{proof}

\begin{lemma}[image kernels at scale $k$]\label{app:G1:I9-imgk}
Assume (T.nw), (E.r) and (T.per). Let $0\le k<K$, $S_k:=\ell_T/t_k$ the side of $T_\eps$ in scale-$k$ units, $\square$ a cube of side $3r(t_k)$ in scale-$k$ units and $\iota$ its isometric lift (Lemma~\ref{app:G1:nowrap}). Let $F$ be the kernel of a line of index $\le k$ of a step-$k$ cut class, scalar or vector, or a kernel obtained from it by at most one lattice difference in each of the two variables, in scale-$k$ units. Then on $\square\times\square$
\[
F_T(x,x')=F_\infty(\iota x,\iota x')+\delta F(x,x'),\qquad \norm{\delta F}_{1,\alpha}\le C_1\mathrm{e}^{-\delta_0S_k/2}\le C_1\mathrm{e}^{-3\delta_0r(t_k)},
\]
with $C_1$ and $\delta_0$ those of Lemma~\ref{lem:IMG}.
\end{lemma}
\begin{proof}
For the scalar kernels, Lemma~\ref{lem:IMG}(i) gives $F_T=\operatorname{Per}F_\infty$; for the vector kernels this is Lemma~\ref{app:K:L04-5}. The diameter of $\square$ is $3\sqrt3\,r(t_k)$; by (T.per) and the proof of Lemma~\ref{lem:IMG}(iv), $S_k\ge6\sqrt3\,r(t_k)+2$ (twice the diameter) and $S_k\ge2/\delta_0$. Now Lemma~\ref{lem:IMG}(iii),(iv) give the bound.
\end{proof}

\begin{lemma}[image lines regroup into renormalised classes]\label{app:G1:I9-regroup}
Make the assumptions of Lemma~\ref{app:G1:I9-imgk} and $R\ge R^{**}$. Fix $k<K$, a pair $(G,\bar l)$ and a localisation $\Lambda$ on $T_\eps$. Cut $l(m)$ (Lemma~\ref{app:G1:cut}) and the lines between cubes at distance $\ge1$ \cite[(3.1), p.~432]{Ba3}; these are torus kernels bounded in torus distance and are not split. Regroup into classes \cite[p.~432]{Ba3}; each connected component $\Gamma$ is localised in a cube $\square_1$ of side $r(t_k)\ge\ell_{\Nbar}$, centred in a cube $\square$ of side $3r(t_k)$ \cite[p.~433]{Ba3}. A \emph{line position} of $\Gamma$ is a line of a member of $\Gamma$; in a member carrying $\delta m^2_S$, written in the form of \cite[p.~417]{Ba3}, the lines of $S$ are the same line positions as in the parent member, and the members of a Ward--Takahashi class (\cite[pp.~431, 440--445]{Ba3}, types C3, C4, C5) share their outer lines and have their own inner lines. Split each line position once, $K_s=K_s^\infty+\delta K_s$ (Lemma~\ref{app:G1:I9-imgk}), so that a counterterm member and its parent receive the same choice; for each Ward--Takahashi subclass $\Pi_Q$, from the outermost inwards, write $\Pi_Q=\Pi_Q^\infty+\Pi_Q^\delta$ with $\Pi_Q^\infty$ the terms whose inner lines are all $K^\infty$. Then $\Gamma=\Gamma^\infty+\sum_{J\ne\emptyset}\Gamma^J$, $J$ the set of line positions carrying $\delta K$, with at most $2^{N_{\rm lp}}$ terms, $N_{\rm lp}=N_{\rm lp}(\nbar)$, and:
(a) $\Gamma^\infty$ is the class lifted by $\iota$ to $\eta\Z^3$ with every propagator $G_k(0)$, i.e.\ the expression after the last substitution of \cite[p.~433]{Ba3}, bounded by \cite[(3.3)]{Ba3};
(b) in every $\Gamma^J$, replacing each line in $J$ by a pair of external fields carrying $\delta K_s$ gives a finite sum of renormalised classes on $\eta\Z^3$, each bounded by \cite[(3.3)]{Ba3} with the $\delta K_s$ among its external fields;
(c) the per-localisation bound of Lemma~\ref{app:G1:I145-perloc} holds on $T_\eps$ with $C_P$ replaced by $C_P^T:=C_P'(1+C_1)^{N_{\rm lp}}$, $C_P'=C_P'(\nbar,\alpha_0)$, and the part with $J\ne\emptyset$ is at most $C_P'N_{\rm lp}(1+C_1)^{N_{\rm lp}}\mathrm{e}^{-\delta_0S_k/2}P(t_k)\mathrm{e}^{-\hat\delta d(\Lambda)}$, with $d(\Lambda)$ the tree length in torus distance and $\hat\delta$ as in Proposition~\ref{app:G1:I9}.
\end{lemma}
\begin{proof}
$\square_1$ exists since $\ell_{\Nbar}\le r(t_k)$ by $R\ge R^{**}$, and $\square$ lifts by Lemma~\ref{app:G1:nowrap}. Every line position joins two points of $\square_1$, since the vertices of a counterterm member lie in the cubes of its lines \cite[p.~420]{Ba3}; so Lemma~\ref{app:G1:I9-imgk} applies. The vacuum terms are at zero background, so the expansion in $B'$ and the gauge transformation removing $B_0$ of \cite[p.~433]{Ba3} are trivial; this gives (a).
(b) Induction on the number of lines. Let $M$ be a member of $\Gamma^J$ and $Q$ a non-vanishing connected component of $M\setminus J$ with $D(Q)\le0$; Lemma~\ref{app:G1:I9-deg} applies. Single vertices and the tree \cite[(2.4)]{Ba3} are bounded as in \cite[(2.4), Props.~2.1--2.2]{Ba3}; a component with a finite vertex has $D\ge0$, with equality only for a single vertex or a vanishing loop; C1$'$ is the finite vertex of \cite[(3.22)--(3.24)]{Ba3}. Otherwise $Q$ is of type C1, C2, C3, C4 or C5; its degree is the same in $M$ and in $M\setminus J$, so $Q$ is a divergent subgraph of $M$ without line position in $J$, and its partners are members of $\Gamma$ \cite[p.~432]{Ba3}. For one-particle-irreducible C1, C2, the partner $M[Q\to\delta m^2_Q]$ has, in the form of p.~417, the line positions of $M$, hence the same choices, and the pair is Ba{\l}aban's renormalised pair on $\eta\Z^3$. The one-particle-reducible members of C2 are treated by Step~3 of the proof of Lemma~\ref{app:G1:I145-perloc}; their pieces share line positions, and if the joining line lies in $J$ the pieces are two C1 pairs. For C3, C4, C5 the term contains $\Pi_Q^\infty$, the lifted class on $\eta\Z^3$, where the Ward identities hold exactly (Lemma~\ref{app:G1:gaugetriv}, Proposition~\ref{app:G1:I2}) and \cite[(3.25)--(3.38)]{Ba3} apply as printed; a term with $\Pi_Q^\delta$ consists of single members with an inner image line, treated by the same case analysis with fewer lines. A counterterm member $M[S\to\delta m^2_S]$ with $S\cap J\ne\emptyset$, written in the one-vertex form of p.~417, carries $c^J_S(x)=\sum_{x'}\eta^3\Sigma^J_S(x,x')$; by the induction hypothesis for the class of $S$ (subgraph counterterms included in $\delta m^2_S$ by \cite[(1.23), p.~430]{Ba3}) and \cite[(3.3)]{Ba3}, $\abs{c^J_S}\le C\prod_{s\in J\cap S}\norm{\delta K_s}$ times the couplings of $S$. So it is a graph with one finite vertex of degree $0$; by Lemma~\ref{app:G1:I9-deg} its components through that vertex need no partner, and its other components are subgraphs of $M$ outside $S$.
(c) Each term is bounded by \cite[(3.3)]{Ba3}, linear in the norms of the external fields; the exponential factor comes from the cut kernels as in Lemma~\ref{app:G1:I145-perloc}. Sum over the at most $2^{N_{\rm lp}}$ terms with $\norm{\delta K_s}\le C_1\mathrm{e}^{-\delta_0S_k/2}$.
\end{proof}

\begin{proposition}[the vacuum classes at the last scale]\label{app:G1:I9}
Assume (T.nw), (E.r), (T.per), $R\ge R^{**}$ and $\delta_0R\ge4$, where $\delta_0$ is the constant of Lemma~\ref{lem:IMG}, and let $S_k:=\ell_T/t_k$. Let $\delta_T>0$ be the tree rate of \cite[(1.33)]{Ba3} as established in Proposition~\ref{app:G1:I145-prop1base}, $\delta_P$ the propagator rate of \cite[(2.10), p.~426]{Ba3}, and $\hat\delta:=\min(\delta_T,\delta_P)$. Call a localisation of a term of $\cV_k$ bad if it meets both $W=B^k(\Lambda_7^{(k)})$ and $W^c$, and let $\cV_k^{\rm bad}$ be the sum of the terms with bad localisations. At the step $k\to k+1$, on $T_\eps$ with all lines torus lines,
\[
\abs{\cV_k^{\rm bad}}\le C_9\,P(t_k)\Big[\abs{B^k(\Lambda_7^{(k)}\cap\Lambda_8^{(k)c})}+\mathrm{e}^{-\hat\delta r(t_k)/2}\abs{B^k(\Lambda_8^{(k)})}\Big],
\]
where $\abs\cdot$ counts scale-$k$ unit blocks, $C_9=C_P^T\,N_G\,N_s\,S_{\hat\delta}\,(1+\mathrm{e}^{\hat\delta(c_3+1)/2})$, $C_P^T$ is the constant of Lemma~\ref{app:G1:I9-regroup}(c), $N_G$ is the number of pairs (bare vacuum graph, ordering) of order $\le\nbar$, $N_s$ the largest number of localised vertices, $S_{\hat\delta}:=\sum_{D\ge0}(c''_3(D+2)^3)^{N_s}\mathrm{e}^{-\hat\delta D/2}$ and $c_3=2\sqrt3$. $C_9$ is independent of $\eps$, $k$, the domains and $\ell_T$. Let $\cV_k^{\rm img}$ be the part of $\cV_k$, over all localisations, in which some internal line carries its image kernel. Then, with constants $C$ independent of $\eps$, $k$, the domains and $\ell_T$,
\[
\abs{\cV_k^{\rm img}}\le C\,P(t_k)\,\mathrm{e}^{-\delta_0S_k/2}\,\abs T\,t_k^{-3},\qquad
\sum_{k<K}\abs{\cV_k^{\rm img}}\le C\,\abs T\,\eps_0^{-3}\mathrm{e}^{-\delta_0\ell_T/\eps_0}\le C\,\abs T\,\eps_0^{-3}\mathrm{e}^{-\delta_0\ell_0/\eps_0}.
\]
\end{proposition}
\begin{proof}
\emph{Per localisation.} By Lemma~\ref{app:G1:cut}, $\cV_k$ is a finite sum over pairs (graph, ordering) of two-leg classes of $G\setminus l(m)$ with external field $K_k$. By Lemma~\ref{app:G1:cutclasses}(c), in every non-vanishing term $G\setminus l(m)$ is connected, so the tree of \cite[(1.33)]{Ba3} joins all cubes of $\Lambda$. By Lemma~\ref{app:G1:cutclasses}(d), the components with $D\le0$ are: C1, C2 with their counterterms (Lemma~\ref{app:G1:I9-partners}), the one-particle-reducible C2 members by Step~3 of the proof of Lemma~\ref{app:G1:I145-perloc}; C1$'$, a finite vertex \cite[(3.22)--(3.24)]{Ba3}; C3, C4 with their Ward--Takahashi classes \cite[(3.25)--(3.32)]{Ba3}, run after the transfer of \cite[p.~433]{Ba3} to a non-wrapping cube (Lemmas~\ref{app:G1:gaugetriv}, \ref{app:G1:nowrap}), the local $AA$ vertex left by \cite[(3.30)]{Ba3} having coefficient $O(1)$ at $j_0=k$ and being multiplied by $K_k(x,x)=O(1)$; and a bare finite vertex, localised in one cube and never bad. In the grading of Definition~\ref{app:G1:vacgrade} the line $l(m)$ is never inside a coefficient (Remark~\ref{app:G1:cutopen-rem}). In scale-$k$ units, $\norm{hK_kh'}_{1,\alpha_0}\le O(1)\mathrm{e}^{-\delta_P\dist(\supp h,\supp h')}$ for scalar and vector lines \cite[(2.10)--(2.11), p.~426]{Ba3}. Proposition~\ref{app:G1:I145-prop1base}, applied with this external field \cite[p.~419]{Ba3} at scale $k+1$, and Lemma~\ref{app:G1:I9-regroup}(c) on the torus give $\abs{\cV_{G,\bar l}(\Lambda)}\le C_P^TP(t_k)\mathrm{e}^{-\hat\delta d(\Lambda)}$; summing over the $N_G$ pairs gives the hypothesis of Lemma~\ref{app:G1:VS} with $A=C_P^TN_GP(t_k)$ and $\delta=\hat\delta$.

\emph{Localisation sum.} Lemma~\ref{app:G1:VS} holds on $T_\eps$ with torus distances, since a torus ball contains at most as many unit cubes as the ball of the same radius in $\Z^3$. Apply it with $W$ as above and $r=r(t_k)$. By Lemma~\ref{app:G1:I9-geo}, which uses the torus distances of \cite[(2.8)]{Ba2}, the blocks of $W$ within $r$ of $W^c$ are at most $\abs{B^k(\Lambda_7^{(k)}\cap\Lambda_8^{(k)c})}$, and the other blocks of $W$ at most $\abs{B^k(\Lambda_7^{(k)}\cap\Lambda_8^{(k)c})}+\abs{B^k(\Lambda_8^{(k)})}$. With $\mathrm{e}^{-\hat\delta r/2}\le1$ this gives the first bound. Vacuum terms have no external legs, so the cut-off functions of external legs play no role.

\emph{Torus lines.} The cut kernels are torus kernels bounded in torus distance; the internal lines of each component are split on a non-wrapping cube by Lemma~\ref{app:G1:I9-imgk}, and by Lemma~\ref{app:G1:I9-regroup}(b) every resulting term is a sum of renormalised classes on $\eta\Z^3$ with the image kernels as external fields. For $\cV_k^{\rm img}$, sum the bound of Lemma~\ref{app:G1:I9-regroup}(c) over the localisations anchored at each of the $\abs Tt_k^{-3}$ unit blocks; with the anchor fixed the sum over $\Lambda$ is at most $N_sS_{\hat\delta}$, and $P(t_k)\le P(\eps_0)$. Put $a:=\delta_0\ell_T/(2\eps_0)$ and $t_k=\eps_0L^{-n}$, $n\ge1$; then $t_k^{-3}\mathrm{e}^{-\delta_0S_k/2}=(L^n/\eps_0)^3\mathrm{e}^{-aL^n}$. Since $r(\eps_0)\ge R$ and $\delta_0R\ge4$, (T.per) gives $a\ge3\sqrt3\,\delta_0r(\eps_0)\ge12\sqrt3>\log(2L^3)$ for $L\le3$, so consecutive terms decrease by a factor $\le L^3\mathrm{e}^{-a}\le\frac12$, and $\sum_{n\ge1}(L^n/\eps_0)^3\mathrm{e}^{-aL^n}\le2L^3\eps_0^{-3}\mathrm{e}^{-aL}\le2L^3\eps_0^{-3}\mathrm{e}^{-\delta_0\ell_T/\eps_0}$, since $L\ge2$. Finally $\ell_T\ge\ell_0$.
\end{proof}

\subsection{Kernel bounds at the Higgs masses and the vector slices}\label{app:G1:kernelbounds}

This subsection works in scale-$k$ units on $\eta\Z^3$, $\eta=\eta_k=L^{-k}$, or on a lattice torus of spacing $\eta$. Kernels are taken with respect to $\eta^3\sum$. Fields may have several real components, and all operators below act componentwise. Let $\Gamma=\bigoplus_\mu S^{(\mu)}\Z e_\mu$ with $S^{(\mu)}\in\eta\N$, let $T=\eta\Z^3/\Gamma$ and let $\pi$ be the projection. A kernel $K$ on $\eta\Z^3$ is \emph{$\Gamma$-invariant} if $K(x+n,y+n)=K(x,y)$ for $n\in\Gamma$. It is \emph{row-summable} if $\sup_x\eta^3\sum_y\abs{K(x,y)}<\infty$. For such $K$ put $(\operatorname{Per}K)(\pi x,\pi y):=\sum_{n\in\Gamma}K(x,y+n)$; this does not depend on the lifts. For a finite-range $\Gamma$-invariant operator $D$ on $\eta\Z^3$ (a lattice difference, a plaquette sum or a contour sum), its torus version is $D_Tf:=(D\tilde f)\circ\pi^{-1}$, where $\tilde f$ is the periodic lift of $f$. With this convention, a repeated neighbour on a small torus is counted with its multiplicity.

We read (H$_\infty$) of Subsection~\ref{ssec:T6-IMG} for a family of kernels that is closed under the differences considered, with constants $C_0,\delta_0$, as follows:
\begin{itemize}
\item every member satisfies $\abs{K(x,x')}\le C_0\mathrm{e}^{-\delta_0\abs{x-x'}}$ for $\abs{x-x'}\ge1$;
\item every member that carries a difference in $x$ satisfies $\abs{K(x_1,x')-K(x_2,x')}\le C_0\abs{x_1-x_2}^\alpha \mathrm{e}^{-\delta_0\min_r\abs{x_r-x'}}$ for $\min_r\abs{x_r-x'}\ge1$, and likewise in $x'$;
\item where it is claimed, $\eta^3\sum_{x'}\abs{K(x,x')}\le C_0$.
\end{itemize}
On a torus, $\abs{x-x'}$ is the torus distance.

\begin{lemma}[periodisation of products]\label{app:G1:per}
Let $K,K'$ be $\Gamma$-invariant and row-summable. Then $KK'$ is $\Gamma$-invariant and row-summable, and $\operatorname{Per}(KK')=(\operatorname{Per}K)(\operatorname{Per}K')$, with absolutely convergent sums. Moreover $\operatorname{Per}(\eta^{-3}\delta)=\eta^{-3}\delta_T$. A finite-range $\Gamma$-invariant operator $D$ satisfies $\operatorname{Per}D=D_T$, and $\operatorname{Per}(D^t)=(D_T)^t$.
\end{lemma}

\begin{proof}
First, $\eta^3\sum_w\sum_n\abs{K(x,w)}\abs{K'(w,y+n)}\le\eta^3\sum_w\abs{K(x,w)}\,\eta^{-3}\sup_w\eta^3\sum_v\abs{K'(w,v)}<\infty$, which justifies the rearrangements below. Using $K'(\tilde z,y+n)=K'(\tilde z+m,y+n+m)$ and $n':=n+m$, we get
\[
\eta^3\sum_{z\in T}\sum_{m}K(x,\tilde z+m)\sum_nK'(\tilde z,y+n)=\sum_{n'}\eta^3\sum_{w\in\eta\Z^3}K(x,w)K'(w,y+n')=\operatorname{Per}(KK')(\pi x,\pi y).
\]
The identity for $\eta^{-3}\delta$ is immediate. Next, $(\operatorname{Per}D)f(\pi x)=\eta^3\sum_{z\in T}\sum_nD(x,\tilde z+n)f(z)=\eta^3\sum_wD(x,w)\tilde f(w)=(D\tilde f)(x)$. Finally, by $\Gamma$-invariance, $\operatorname{Per}(D^t)(\pi x,\pi y)=\sum_nD(y+n,x)=\sum_nD(y,x-n)=(\operatorname{Per}D)(\pi y,\pi x)$.
\end{proof}

\begin{lemma}[a priori bounds]\label{app:G1:apriori}
Let $1\le i\le k+1$ and $\ell_i:=L^{i-k}$. Let $\Lambda$ be $\eta\Z^3$, or a torus of spacing $\eta$ whose sides are multiples of $\ell_i$. Let $P_i$ be the orthogonal projection onto the fields constant on the blocks $y+[0,\ell_i)^3$, $y\in\ell_i\Z^3$ (on a torus, their images). Then:
\begin{enumerate}
\item[(i)] $\ip\varphi{-\Delta\varphi}\ge4\ell_i^{-2}\norm{(1-P_i)\varphi}^2$;
\item[(ii)] for $c\ge0$ and $\beta>0$, the operator $H:=-\Delta+c+\beta P_i$ satisfies $H\ge\gamma:=\min(4\ell_i^{-2},\beta)$;
\item[(iii)] on $\eta\Z^3$ there is $\rho>0$ with $\abs{H^{-1}(x,y)}\le2\gamma^{-1}\eta^{-3}\mathrm{e}^{-\rho\abs{x-y}}$.
\end{enumerate}
\end{lemma}

\begin{proof}
(i) A block has $\ell_i/\eta=L^i\ge2$ sites per side. Dropping the bonds between different blocks (and, on a torus, any bond that closes a block around the torus) leaves the Neumann Laplacians of the blocks. Their kernel consists of the constants, and their gap is the gap of the one-dimensional Neumann Laplacian, namely $4\eta^{-2}\sin^2(\pi\eta/2\ell_i)\ge4\ell_i^{-2}$, since $\sin u\ge2u/\pi$ on $[0,\pi/2]$.

(ii) By (i), $\ip\varphi{H\varphi}\ge4\ell_i^{-2}\norm{(1-P_i)\varphi}^2+\beta\norm{P_i\varphi}^2\ge\gamma\norm\varphi^2$.

(iii) The kernel of $H$ vanishes for $\abs{x-y}>r:=2L$, and $h:=\sup_x\eta^3\sum_y\abs{H(x,y)}<\infty$. Choose $\rho>0$ with $h(\mathrm{e}^{\rho r}-1)\le\gamma/2$. For a unit vector $\nu$ and $R>0$ put $f(x):=\max(-R,\min(R,\rho\,\nu\cdot x))$; it is bounded and $\rho$-Lipschitz. By the Schur test, $E:=\mathrm{e}^fH\mathrm{e}^{-f}-H$ has $\norm E\le\gamma/2$. Hence $\rRe\ip\varphi{\mathrm{e}^fH\mathrm{e}^{-f}\varphi}\ge\frac\gamma2\norm\varphi^2$, and $\mathrm{e}^fH\mathrm{e}^{-f}$ is invertible with inverse $\mathrm{e}^fH^{-1}\mathrm{e}^{-f}$ of norm $\le2/\gamma$. Apply this to $\eta^{-3}\mathbf 1_y$, which has norm $\eta^{-3/2}$, and use $\abs{\psi(x)}\le\eta^{-3/2}\norm\psi$. This gives $\abs{H^{-1}(x,y)}\le2\gamma^{-1}\eta^{-3}\mathrm{e}^{-(f(x)-f(y))}$. Taking $\nu=(x-y)/\abs{x-y}$ and $R\ge\rho(\abs x+\abs y)$ gives the claim.
\end{proof}

\begin{proposition}[the bounds (H$_\infty$) for the scalar kernels at a mass $M$]\label{app:K:L19-1}\label{app:G1:Hinf-M}
Let $L\ge2$ be an integer, $a>0$, $a_\flat:=a(1-L^{-2})$, $a_j\in[a_\flat,a]$ for $j\ge0$, and $\alpha\in[0,1)$. Let $\eps_0>0$, $K\ge0$, $t_j:=L^{j-K}\eps_0$, $0\le k\le K$, and let $M\ge0$ with $M\eps_0\le1$. Work on $\eta\Z^3$, $\eta:=L^{-k}$, with $\partial_{u,\mu}f(u):=\eta^{-1}(f(u+\eta e_\mu)-f(u))$. For $1\le i\le k+1$ let $P_i$ be as in Lemma~\ref{app:G1:apriori}, and put
\[
G_i:=\big(-\Delta+M^2t_k^2+a_iL^{2(k-i)}P_i\big)^{-1},\qquad G_{(0)}:=G_1,\qquad G_{(j)}:=G_{j+1}-G_j\ (1\le j\le k),
\]
and $\ell_j:=L^{j-k}$ for $0\le j\le k$. For $\sigma,\tau\in\{0,e_1,e_2,e_3\}$ let $\partial^\sigma_u\partial^\tau_v$ be at most one difference in each variable, and put $n:=\abs\sigma+\abs\tau$. There are $\kappa>0$ and $C<\infty$, depending only on $a,L,\alpha$, such that the following hold for all $K,k,\eps_0,M,(a_j)$ as above, all $0\le j\le k$ and all $1\le i\le k+1$.
\begin{enumerate}
\item[(i)] For all $u,v$:
\[
\abs{\partial^\sigma_u\partial^\tau_vG_{(j)}(u,v)}\le C\ell_j^{-1-n}\mathrm{e}^{-\kappa\abs{u-v}_1/\ell_j},\qquad \eta^3\sum_v\abs{\partial^\sigma_u\partial^\tau_vG_{(j)}(u,v)}\le C\ell_j^{2-n}.
\]
If $\abs\sigma=1$, then for all $u_1\ne u_2$,
\[
\abs{\partial^\sigma\partial^\tau G_{(j)}(u_1,v)-\partial^\sigma\partial^\tau G_{(j)}(u_2,v)}\le C\ell_j^{-1-n-\alpha}\abs{u_1-u_2}^\alpha \mathrm{e}^{-\kappa\min_r\abs{u_r-v}_1/\ell_j}.
\]
The same H\"older bound holds in the second variable when $\abs\tau=1$.
\item[(ii)] Put $\delta_0:=\kappa/2$. There is $C_0=C_0(a,L,\alpha)$ such that for every $G\in\{G_i,G_{(j)}\}$:
\begin{itemize}
\item $\abs{\partial^\sigma_u\partial^\tau_vG(u,v)}\le C_0\mathrm{e}^{-\delta_0\abs{u-v}}$ if $\abs{u-v}\ge1$;
\item if $\abs\sigma=1$ and $\min_r\abs{u_r-v}\ge1$, the difference quotient of (i) is at most $C_0\abs{u_1-u_2}^\alpha \mathrm{e}^{-\delta_0\min_r\abs{u_r-v}}$, and likewise in $v$;
\item if $n\le1$, then $\eta^3\sum_v\abs{\partial^\sigma_u\partial^\tau_vG(u,v)}\le C_0$.
\end{itemize}
\item[(iii)] If $n=2$, then $\eta^3\sum_v\abs{\partial^\sigma_u\partial^\tau_vG_i(u,v)}\le C\,i$.
\end{enumerate}
\end{proposition}

\begin{proof}
\emph{Step 1: a priori decay.} By Lemma~\ref{app:G1:apriori}, every $G_i$ is bounded with $\abs{G_i(u,v)}\le c\,\mathrm{e}^{-\rho\abs{u-v}}$, where $c,\rho$ may depend on $k$ and $i$. Hence every kernel $\partial^\sigma_u\partial^\tau_vG$, $G\in\{G_i,G_{(j)}\}$, is row-summable. These constants are used only qualitatively. Since $G_i=G_1+\sum_{j=1}^{i-1}(G_{j+1}-G_j)$, we have $G_i=\sum_{j<i}G_{(j)}$.

\emph{Step 2: periodisation.} For $N\ge2$ let $\Gamma_N:=LN\Z^3$ and $T_N:=\eta\Z^3/\Gamma_N$. Since $LN/\ell_i=NL^{k+1-i}\in\N$, the blocks project to a partition of $T_N$, and a block meets each $\Gamma_N$-orbit at most once. So $P_i$ is $\Gamma_N$-invariant and $\operatorname{Per}P_i$ is the block projection of $T_N$. With Lemma~\ref{app:G1:per}, $\operatorname{Per}H_i=H_{i,N}$, the same operator on $T_N$ with the periodic Laplacian. Moreover $H_{i,N}\operatorname{Per}G_i=\operatorname{Per}(H_iG_i)=\eta^{-3}\delta_{T_N}$. By Lemma~\ref{app:G1:apriori}(ii) on $T_N$, $H_{i,N}$ is invertible. Hence $G_{i,N}=\operatorname{Per}G_i$, and $G_{(j),N}:=G_{j+1,N}-G_{j,N}=\operatorname{Per}G_{(j)}$. Differences commute with $\operatorname{Per}$, so
\[
\partial^\sigma\partial^\tau G_N(\pi u,\pi v)=\sum_{n\in\Gamma_N}\partial^\sigma\partial^\tau G(u,v+n).
\]

\emph{Step 3: scaling.} Fix $j$ and put $s:=\ell_j$ and $\eta_j:=\eta/s=L^{-j}$. The map $u\mapsto u/s$ sends $T_N$ to the torus $\Lambda$ of spacing $\eta_j$ and side $\ell:=L^{k-j}LN\in L\N$, and it sends the blocks of side $\ell_i$ to blocks of side $L^{i-j}$. With $(Uf)(x):=f(sx)$ one has $H_{i,N}=s^{-2}U^{-1}X_iU$, where
\[
X_i:=-\Delta_{\eta_j}+(Mt_j)^2+a_iL^{2(j-i)}P_{L^{i-j}} .
\]
Comparing $\eta^3\sum_v$ with $\eta_j^3\sum_y$ gives $G_{i,N}(u,v)=s^{-1}X_i^{-1}(u/s,v/s)$, and $\partial_u=s^{-1}\nabla_x$. Moreover $X_j=-\Delta+\tilde\mu^2+a_jP_1$ and $X_{j+1}=-\Delta+\tilde\mu^2+a_{j+1}L^{-2}P_L$, where $\tilde\mu:=Mt_j\le M\eps_0\le1$. These are the operators $X,X'$ of the setting of Theorem~\ref{app:G1:VSR-VS}, with $b=a_j$ and $b'=a_{j+1}L^{-2}$, spacing $\eta_j\in(0,1]$, $1/\eta_j\in\N$ and $\ell\in L\N$, $\ell\ge L$.

\emph{Step 4: slice bounds on $T_N$.} Let $d_N$ be the $\ell^1$ distance of $T_N$.

Case $j\ge1$. Here $G_{(j),N}(u,v)=s^{-1}H(u/s,v/s)$ with $H=X_{j+1}^{-1}-X_j^{-1}$. By (S1)--(S4) of Theorem~\ref{app:G1:VSR-VS}, whose constants and rate $\kappa_H$ depend only on $(a,L,\alpha)$ and are uniform in the data of Step 3, we get $\abs{\partial^\sigma\partial^\tau G_{(j),N}(u,v)}\le C_Vs^{-1-n}\mathrm{e}^{-\kappa_Hd_N(u,v)/s}$. For $\abs\sigma=1$ the difference quotient is at most $C_Vs^{-1-n}(d_N(u_1,u_2)/s)^\alpha \mathrm{e}^{-\kappa_H\min_rd_N(u_r,v)/s}$. Since $X_j,X_{j+1}$ are symmetric, $H(x,y)=H(y,x)$, which gives the bounds in the second variable.

Case $j=0$. Here $s=\eta$, $\eta_0=1$, and $G_{(0),N}=G_{1,N}$ has kernel $s^{-1}X'^{-1}(u/s,v/s)$, where $X'=X_1$. Put $\gamma_0:=\min(4,a_\flat)$, and let $c=c(a,L)\in(0,\frac12]$ be the largest number with $3.4c^2+a(\mathrm{e}^{3c}-1)\le\gamma_0/2$. Then $\kappa_H\le c/L$. By Lemma~\ref{app:G1:VSR-CT}(i), $X'\ge\gamma_0L^{-2}$. The condition of Lemma~\ref{app:G1:VSR-CT}(ii) with $s=L$, $\beta\le aL^{-2}$, $\gamma=\gamma_0L^{-2}$ and $\kappa=c/L$ is the defining inequality of $c$, divided by $L^2$. Hence $\norm{\mathrm{e}^{\kappa\rho}X'^{-1}\mathrm{e}^{-\kappa\rho}}\le2L^2/\gamma_0$ for $\rho=d(\cdot,y)$. At spacing $1$ a kernel value is a matrix element between unit vectors, so $\abs{X'^{-1}(x,y)}\le2L^2\gamma_0^{-1}\mathrm{e}^{-(c/L)d(x,y)}$. Each difference at spacing $1$ is bounded by two values at points within distance $1$, so $\abs{\nabla^\sigma\nabla^\tau X'^{-1}(x,y)}\le8L^2\gamma_0^{-1}\mathrm{e}^{2c/L}\mathrm{e}^{-(c/L)d(x,y)}$. Distinct points satisfy $d(x_1,x_2)\ge1$, so the difference quotients are at most $d(x_1,x_2)^\alpha$ times twice this bound, with $\min_r$ in the exponent. Since $\kappa_H\le c/L$, the bounds of the case $j\ge1$ hold for $j=0$ as well.

\emph{Step 5: passage to $\eta\Z^3$.} Fix the finitely many points that enter a given bound. By Steps 1 and 2, $\partial^\sigma\partial^\tau G_N(\pi u,\pi v)-\partial^\sigma\partial^\tau G(u,v)$ is a tail of an absolutely convergent series, and it tends to $0$ as $N\to\infty$. For $LN>2\abs{u-v}_1$ one has $d_N(\pi u,\pi v)=\abs{u-v}_1$. So the bounds of Step 4 hold on $\eta\Z^3$ with $d_N$ replaced by $\abs\cdot_1$. Using $\abs w\le\abs w_1\le\sqrt3\abs w$, this gives the pointwise and H\"older bounds of (i), with $\kappa:=\kappa_H$. For the summability, the substitution $v=sy$ gives $\eta^3\sum_v\mathrm{e}^{-\kappa\abs{u-v}_1/s}=s^3J_0(\kappa)$, where $J_0(\kappa)=\eta_j^3\sum_{y\in\eta_j\Z^3}\mathrm{e}^{-\kappa\abs y_1}\le(1+2/\kappa)^3$ because $\eta_j\le1$. Multiplying by $Cs^{-1-n}$ gives $Cs^{2-n}J_0(\kappa)$.

\emph{Step 6: summation.} Put $q_j:=1/\ell_j=L^{k-j}$; distinct $j$ give distinct powers of $L$. For $q,z\ge1$ one has $2qz-q-z=q(z-1)+z(q-1)\ge0$, hence $\mathrm{e}^{-\kappa qz}\le \mathrm{e}^{-\kappa z/2}\mathrm{e}^{-\kappa q/2}$. Let $\abs{u-v}\ge1$, so that $z:=\abs{u-v}_1\ge1$. By (i) and $G_i=\sum_{j<i}G_{(j)}$,
\[
\abs{\partial^\sigma\partial^\tau G_i(u,v)}\le C\sum_{j<i}q_j^{1+n}\mathrm{e}^{-\kappa q_jz}\le C\mathrm{e}^{-\kappa z/2}\Sigma,\qquad \Sigma:=\sum_{m\ge0}L^{4m}\mathrm{e}^{-\kappa L^m/2}<\infty .
\]
The H\"older bound follows in the same way, using $q_j^{1+n+\alpha}\le q_j^4$ and $z:=\min_r\abs{u_r-v}_1\ge1$. For a single slice, $q_j^{1+n+\alpha}\mathrm{e}^{-\kappa q_jz}\le \mathrm{e}^{-\kappa z/2}\Sigma$. For $n\le1$, summing (i) gives $\eta^3\sum_v\abs{\partial^\sigma\partial^\tau G_i(u,v)}\le CJ_0(\kappa)\sum_{j<i}\ell_j\le CJ_0(\kappa)L/(L-1)$; a single slice gives at most $CJ_0(\kappa)$. Finally, $\mathrm{e}^{-\kappa z/2}\le \mathrm{e}^{-\delta_0\abs{u-v}}$. This proves (ii) with $C_0:=C\max(\Sigma,J_0(\kappa)L/(L-1))$. For $n=2$ the same sum is at most $CJ_0(\kappa)\,i$, which is (iii).
\end{proof}

\begin{remark}\label{app:G1:Hinf-M-rem}
The bounds (i) are \cite[(2.10)--(2.11), p.~426]{Ba3} at zero background with $\delta_1=\kappa$, in scale-$k$ units (Corollary~\ref{app:G1:VSR-VSphys}). They are uniform in $k$, $\eps$ and $M\in[0,\eps_0^{-1}]$, in particular for $M\in\{m_1,m_2\}$ when $m_2\eps_0\le1$. Part (ii) is (H$_\infty$) for these kernels in the reading above. The summability clause holds for $n\le1$ only; for $n=2$ the bound (iii) grows with $i$. Pure second differences in a single variable are not covered.
\end{remark}

\begin{lemma}[vector propagators on the torus and on the lattice]\label{app:K:L19-2}\label{app:K:L04-5}
Let $0\le k\le K$ and $\eta:=\eta_k$. Let $T_\eta$ be the torus at step $k$ of Subsection~\ref{ssec:T6-IMG}, with sides $S_k^{(\mu)}$ and deck group $\Gamma_k$. Assume that every $S_k^{(\mu)}$ is a multiple of $L$; for $k<K$ this holds because $S_k^{(\mu)}=LS_{k+1}^{(\mu)}$. Let $\Lambda\in\{T_\eta,\eta\Z^3\}$, write $A_\mu(x):=A_{(x,x+\eta e_\mu)}$, and let $\cH:=\ell^2(\Lambda;\R^3)$ with $\ip fg=\sum_x\eta^3f(x)\cdot g(x)$.

For $1\le i\le k+1$ let $\ell_i:=L^{i-k}$ and $B_i(y):=y+[0,\ell_i)^3$, $y\in\ell_i\Z^3$ (on $T_\eta$, their images). Let $\cH_i$ be the fields on $\ell_i\Z^3\cap\Lambda$ with weight $\ell_i^3$. Let $(Q_if)(y):=\ell_i^{-3}\sum_{x\in B_i(y)}\eta^3f(x)$ and $P_i:=Q_i^*Q_i$. Let $Q\colon\cH_i\to\cH_{i+1}$ be the one-step average, so that $Q_{i+1}=QQ_i$, and put $P:=Q^*Q$.

Let $\mu_0\ge0$, $\mu^2:=\mu_0^2t_k^2$, $\alpha_i:=a_iL^{2(k-i)}$ and $\beta_i:=aL^{2(k-i)-2}$, where the $a_i\in[a_\flat,a]$ are the constants of \cite[(2.13)]{Ba1}, so that $\alpha_{i+1}=\alpha_i\beta_i/(\alpha_i+\beta_i)$. With $-\Delta$ the componentwise lattice Laplacian, put
\[
H_i:=-\Delta+\mu^2+\alpha_iP_i,\quad G_i:=H_i^{-1},\quad C^{(i)}:=(\alpha_i-\alpha_i^2Q_iG_iQ_i^*+\beta_iP)^{-1},
\]
\[
W_{(0)}:=G_1,\qquad W_{(i)}:=\alpha_i^2G_iQ_i^*C^{(i)}Q_iG_i\quad(1\le i\le k).
\]
Up to the rescaling \cite[(2.20)]{K1}, these are the vector propagators of \cite[(2.13)--(2.17)]{K1} and \cite[(2.20)--(2.22), (2.30)--(2.31), (2.42)--(2.43)]{Ba1} at $N=3$, $A=0$, $m=\mu_0$ \cite[p.~608]{Ba1}, \cite[p.~653]{K1}. The $W_{(j)}$ are the covariances of the fields $A'^{(j)}$ of \cite[(1.2), pp.~412, 414]{Ba3}, that is, the vector slices of index $j$ of \cite[p.~426]{Ba3}.

Let $\mathcal W^\sharp$ be the set of the kernels $G_i$ ($1\le i\le k+1$) and $W_{(i)}$ ($0\le i\le k$), and of their compositions, in either variable, with lattice derivatives of order $\le2$, with plaquette sums, and with the contour sums $f\mapsto\sum_{b\in\Gamma^{(j+1)}_{x_{j+1},x}}\eta f(b)$, $j+1\le k$, of \cite[(3.27)]{K1}. Let $\mathcal W\subset\mathcal W^\sharp$ consist of:
\begin{itemize}
\item the kernels $G_i$ and $W_{(i)}$;
\item their compositions, in either variable, with lattice derivatives of order $\le2$ and with plaquette sums;
\item the contour composites $W_{(j)}(\Gamma^{(j+1)}_{x_{j+1},x},b)=\sum_{b'\in\Gamma^{(j+1)}_{x_{j+1},x}}\eta\,W_{(j)}(b',b)$, $j+1\le k$, $x\in B_{j+1}(x_{j+1})$ \cite[(3.27)]{K1}, \cite[(1.3)]{Ba3}.
\end{itemize}
Let $\mathcal W'\subset\mathcal W$ consist of the elements that contain no derivative of order $2$ of a kernel $G_i$; here a plaquette sum counts as one derivative. Then:
\begin{enumerate}
\item[(a)] $H_i\ge\gamma_i:=\min(4\ell_i^{-2},\alpha_i)>0$ on $\cH$, and on $\eta\Z^3$ one has $\abs{G_i(x,y)}\le c_i\mathrm{e}^{-\rho_i\abs{x-y}}$ for some $c_i,\rho_i>0$.
\item[(b)] On each $\Lambda$, $C^{(i)}$ is bounded and $W_{(i)}=G_{i+1}-G_i$ for $1\le i\le k$.
\item[(c)] For every $W\in\mathcal W^\sharp$, $W_T(x,x')=\sum_{n\in\Gamma_k}W_\infty(\tilde x,\tilde x'+n)$, with absolute convergence.
\item[(d)] Let $\cT$ be a family of tori at step $k$ with spacing $\eta$, all of whose sides are multiples of $L$, with minimal sides tending to infinity. Let $\cV'\subset\cV\subset\mathcal W^\sharp$. Suppose that for every $W\in\cV$ the kernels $W_T$, $T\in\cT$, satisfy the pointwise and H\"older bounds of (H$_\infty$), and for every $W\in\cV'$ the summability bound, with constants $C_0,\delta_0$ independent of $T$. Then the kernels $W_\infty$ satisfy the same bounds. In particular this holds for $\cV=\mathcal W^\sharp$ and $\cV'=\mathcal W'$.
\item[(e)] Let $\delta_0$ be the rate of (H$_\infty$) entering condition (T.per) of \eqref{eq:GIMG}, and assume (E.r) and (T.per). Let $\square\subset T_\eta$ be a non-wrapping cube of side $3r(t_k)$, with lift map $\iota$. Then every $W\in\mathcal W$ satisfies $W_T=W_\infty\circ(\iota\times\iota)+\delta W$ on $\square\times\square$, with $\norm{\delta W}_{1,\alpha}\le C_1\mathrm{e}^{-\delta_0S_k/2}\le C_1\mathrm{e}^{-3\delta_0r(t_k)}$. Here $C_1=c_dC_0$, $c_d$ is absolute, and $C_0$ depends only on $a,L,M_b,\alpha$; in particular it is independent of $k$, of the torus and of $\eps_0$.
\item[(f)] Parts (a)--(c) hold with $\R^3$ replaced by $\R^2$ and $\mu^2$ by $m^2t_k^2$, $m\ge0$. Then $W_{(0)}=G_1$ and $W_{(j)}=G_{j+1}-G_j$ ($1\le j\le k$) are the zero-background scalar slices $G^{\eta_k}_{(j)}(0)$ of \cite[(1.27), (2.6)]{Ba3}.
\item[(g)] Let $0\le j\le k<K$. The kernel on $T_\eps$, in scale-$k$ units, of the vector slice $W_{(j)}$, and of each of its lattice derivatives of order $\le2$, is the periodisation of the corresponding kernel on $\eta\Z^3$. The latter satisfies (H$_\infty$) of Subsection~\ref{ssec:T6-IMG}, including the summability bound.
\end{enumerate}
\end{lemma}

\begin{proof}
\emph{Preliminaries.} $Q_i^*u$ is $u$ composed with the block map, so $Q_iQ_i^*=1$, and $P_i$ is the orthogonal projection onto the block-constant fields. Hence $\norm{Q_i}\le1$, $QQ^*=1$, $P^2=P$ and $Q_i^*PQ_i=Q_{i+1}^*Q_{i+1}=P_{i+1}$. Every $S_k^{(\mu)}$ is a multiple of $L$, hence of every $\ell_i$ with $i\le k+1$. So every block partition is $\Gamma_k$-invariant, and $H_i$ commutes with the translations by $\Gamma_k$.

(a) Apply Lemma~\ref{app:G1:apriori} componentwise with $c=\mu^2$ and $\beta=\alpha_i$.

(b) On $\cH\oplus\cH_i$ consider
\[
F(\varphi,\psi):=\ip\varphi{(-\Delta+\mu^2)\varphi}+\alpha_i\norm{\psi-Q_i\varphi}_i^2+\beta_i\norm{P\psi}_i^2 .
\]
Its bounded self-adjoint operator is
\[
\mathbf M=\begin{pmatrix}H_i&-\alpha_iQ_i^*\\-\alpha_iQ_i&\alpha_i+\beta_iP\end{pmatrix}.
\]
We first show that $F$ is coercive. From $\norm{PQ_i\varphi}\le\norm{Q_i\varphi-\psi}+\norm{P\psi}$ and Cauchy--Schwarz,
\[
\tfrac{\alpha_i}2\norm{\psi-Q_i\varphi}^2+\tfrac{\beta_i}2\norm{P\psi}^2\ge\tfrac{\alpha_{i+1}}2\norm{PQ_i\varphi}^2=\tfrac{\alpha_{i+1}}2\norm{P_{i+1}\varphi}^2 .
\]
With Lemma~\ref{app:G1:apriori}(i) at scale $i+1\le k+1$, this gives $F\ge c'\norm\varphi^2+\frac{\alpha_i}2\norm{\psi-Q_i\varphi}^2$, where $c'=\min(4\ell_{i+1}^{-2},\alpha_{i+1}/2)$. Since $\norm\psi^2\le2\norm{\psi-Q_i\varphi}^2+2\norm\varphi^2$, we get $F\ge c(\norm\varphi^2+\norm\psi^2)$ with $c:=\frac14\min(c',\alpha_i/2)>0$.

So $\mathbf M\ge c$, and its diagonal blocks are $\ge c$. The Schur complements
\[
S_1:=\alpha_i+\beta_iP-\alpha_i^2Q_iG_iQ_i^*=(C^{(i)})^{-1},\qquad S_2:=H_i-\alpha_i^2Q_i^*(\alpha_i+\beta_iP)^{-1}Q_i
\]
satisfy $\ip\psi{S_1\psi}=\min_\varphi F(\varphi,\psi)\ge c\norm\psi^2$ and $\ip\varphi{S_2\varphi}=\min_\psi F(\varphi,\psi)\ge c\norm\varphi^2$. In particular $C^{(i)}$ is bounded. Since $(\alpha_i+\beta_iP)^{-1}=\alpha_i^{-1}(1-P)+(\alpha_i+\beta_i)^{-1}P$, we have $\alpha_i-\alpha_i^2(\alpha_i+\beta_iP)^{-1}=\alpha_{i+1}P$. Hence $S_2=-\Delta+\mu^2+\alpha_{i+1}Q_i^*PQ_i=H_{i+1}$. The two block-inversion formulas for the upper left block of $\mathbf M^{-1}$ give $S_2^{-1}$ and $G_i+G_i(\alpha_iQ_i^*)S_1^{-1}(\alpha_iQ_i)G_i$. Hence $G_{i+1}=G_i+W_{(i)}$, which is \cite[(2.42)]{Ba1}.

(c) First let $W=G_i$. By (a), $G_{i,\infty}$ is row-summable, and it is $\Gamma_k$-invariant because $H_{i,\infty}$ is. By the preliminaries and Lemma~\ref{app:G1:per}, $\operatorname{Per}H_{i,\infty}=H_{i,T}$, and $H_{i,T}\operatorname{Per}G_{i,\infty}=\operatorname{Per}(H_{i,\infty}G_{i,\infty})=\eta^{-3}\delta_T$. By (a), $H_{i,T}$ is invertible, so $\operatorname{Per}G_{i,\infty}=G_{i,T}$. By (b) on both lattices, $W_{(i),T}=G_{i+1,T}-G_{i,T}=\operatorname{Per}W_{(i),\infty}$. The derivatives, plaquette sums and contour sums are finite-range operators that commute with $\Gamma_k$. For the contour sums this holds because a contour $\Gamma^{(j+1)}_{x_{j+1},x}$ lies in $B_{j+1}(x_{j+1})$ and is fixed relative to that block \cite[(2.1)--(2.2)]{Ba1}, \cite[(2.12)]{K1}, while $\Gamma_k\subset\ell_{j+1}\Z^3$. A composition in the first variable is a product $DW$, and one in the second variable is a product $WD^t$. So Lemma~\ref{app:G1:per} gives $\operatorname{Per}(DW)=D_TW_T$ and $\operatorname{Per}(WD^t)=W_T(D_T)^t$, and by induction this covers every element of $\mathcal W^\sharp$. Each element of $\mathcal W^\sharp$ is row-summable, so the sums converge absolutely.

(d) Take $T^{(N)}\in\cT$ with $\min_\mu S^{(N),\mu}\to\infty$, and fix $\tilde x,\tilde x'\in\eta\Z^3$. By (c),
\[
\abs{W_{T^{(N)}}(\pi\tilde x,\pi\tilde x')-W_\infty(\tilde x,\tilde x')}\le\sum_{n\in\Gamma^{(N)}\setminus0}\abs{W_\infty(\tilde x,\tilde x'+n)}\to0 .
\]
The same holds for the finitely many values that enter a H\"older quotient. For large $N$ the torus distance of the projected points equals $\abs{\tilde x-\tilde x'}$. So the pointwise and H\"older bounds pass to the limit. For a finite set $F\subset\eta\Z^3$, the map $\pi$ is injective on $F$ for large $N$, so $\eta^3\sum_{x'\in F}\abs{W_\infty(\tilde x,x')}=\lim_N\eta^3\sum_{x'\in F}\abs{W_{T^{(N)}}(\pi\tilde x,\pi x')}\le C_0$; this gives the summability bound.

(e) \emph{Torus bounds.} Let $\cT$ be the family of tori of spacing $\eta$ whose sides, in scale-$k$ units, lie in $LM_b\N$ and are at least $3M_b$. Let $T\in\cT$. By \cite[(2.10)--(2.12), p.~426]{Ba3}, which hold for vector propagators (loc.\ cit.) and are obtained there from \cite[Props.~2.1, 2.3]{Ba1} by rescaling, we have, in scale-$k$ units with $\ell_j=L^{j-k}\le1$,
\[
\abs{\partial^{a}W_{(j),T}(x,x')}\le C\,\ell_j^{-1-\abs a}\mathrm{e}^{-\delta_1\abs{x-x'}/\ell_j}\qquad(\abs a\le2,\ \text{differences in either variable}).
\]
The H\"older quotients carry the extra factor $\ell_j^{-\alpha}\abs{x_1-x_2}^\alpha$, with $\min_r$ in the exponent. The contour composites satisfy $\abs{W_{(j),T}(\Gamma^{(j+1)}_{x_{j+1},x},b)}\le C\mathrm{e}^{-\delta_1\dist(B_{j+1}(x_{j+1}),b)/\ell_j}\le C\mathrm{e}^{\sqrt3L\delta_1}\mathrm{e}^{-\delta_1\abs{x-b}/\ell_j}$, because $B_{j+1}(x_{j+1})$ has side $L\ell_j$. Here $C,\delta_1$ depend only on $a,L,M_b,\alpha$ (Proposition~\ref{prop:UV-BaProp21}; for $\Omega=T$ there is no restriction on the evaluation points \cite[p.~611]{Ba1}), and by Lemma~\ref{app:G1:torusuniform} they do not depend on the torus in $\cT$. They do not depend on $k$ or $\eps_0$ either.

For $\abs{x-x'}\ge1$ and $\varsigma:=\ell_j^{-1}\ge1$ we have $\varsigma\abs{x-x'}\ge(\varsigma+\abs{x-x'})/2$. Hence a single slice satisfies $\ell_j^{-1-\abs a-\alpha}\mathrm{e}^{-\delta_1\abs{x-x'}/\ell_j}\le \mathrm{e}^{-\delta_1\abs{x-x'}/2}\Sigma_1$, with $\Sigma_1:=\sum_{m\ge0}L^{4m}\mathrm{e}^{-\delta_1L^m/2}$. By (b), $G_{i,T}=\sum_{j<i}W_{(j),T}$, and the same estimate, summed over the distinct values $\varsigma=L^{k-j}$, gives the pointwise and H\"older bounds for the kernels $G_i$ and their composites, with rate $\delta_1/2$. For the summability, $\eta^3\sum_{x'}\ell_j^{-1-\abs a}\mathrm{e}^{-\delta_1\abs{x-x'}/\ell_j}\le\ell_j^{2-\abs a}J$, where $J:=(1+2\sqrt3/\delta_1)^3$. This is at most $J$ for a slice with $\abs a\le2$ and for a contour composite. For $G_i$ with $\abs a\le1$ it is at most $J\sum_j\ell_j\le JL/(L-1)$. So the kernels $W_T$, $W\in\mathcal W$, satisfy the pointwise and H\"older bounds of (H$_\infty$), and the elements of $\mathcal W'$ satisfy the summability bound, with rate $\delta_1/2$ and a constant $C_0$ that depends only on $a,L,M_b,\alpha$. These bounds hold for every smaller rate as well. The rate $\delta_0$ of (T.per) is at most $\delta_1/2$: it is the rate obtained in the same way from \cite[(2.10)--(2.12)]{Ba3} for the scalar kernels of Subsection~\ref{ssec:T6-IMG}, and the vector propagators satisfy the same inequalities \cite[p.~426]{Ba3}. So the bounds hold with the rate $\delta_0$.

\emph{Infinite lattice.} By (d) with $\cV=\mathcal W$ and $\cV'=\mathcal W'$, the kernels $W_\infty$, $W\in\mathcal W$, satisfy the same bounds.

\emph{Image bound.} Let $x,x'\in\square$ and $\tilde x=\iota x$, $\tilde x'=\iota x'$. By (c), $\delta W(x,x')=\sum_{n\in\Gamma_k\setminus0}W_\infty(\tilde x,\tilde x'+n)$. The same representation holds for each kernel that enters $\norm\cdot_{1,\alpha}$, because differences commute with $\operatorname{Per}$. The lifted cube has diameter $D=3\sqrt3\,r(t_k)$. By (E.r), (T.per) and Lemma~\ref{lem:B:rows}(b) (see the proof of Lemma~\ref{lem:IMG}(iv)), $S_k\ge6\sqrt3\,r(t_k)+2\ge2D+2$ and $S_k\ge2/\delta_0$. For $n\ne0$ we have $\abs n\ge S_k$, so $\abs{\tilde x-\tilde x'-n}\ge\abs n-D\ge\abs n/2\ge1$, and likewise with $\min_r$ for the H\"older quotients. So only the pointwise and H\"older bounds of $W_\infty$ are used. Write $n=\sum_\mu m_\mu S_k^{(\mu)}e_\mu$, so that $\abs n\ge S_k\abs m$, and put $t:=\delta_0S_k/2\ge1$. Then
\[
\sum_{n\ne0}C_0\mathrm{e}^{-\delta_0\abs n/2}\le C_0\sum_{m\in\Z^3\setminus0}\mathrm{e}^{-t\abs m}\le C_0\,\mathrm{e}^{-t}\sum_{m\ne0}\mathrm{e}^{-(\abs m-1)} .
\]
Counting the terms of $\norm\cdot_{1,\alpha}$ gives $\norm{\delta W}_{1,\alpha}\le c_dC_0\mathrm{e}^{-\delta_0S_k/2}$. Finally, $\delta_0S_k/2\ge3\sqrt3\,\delta_0r(t_k)\ge3\delta_0r(t_k)$.

(f) The proofs of (a)--(c) use neither the number of components nor the value of $\mu^2\ge0$, so they hold verbatim for $\R^2$ and $m^2t_k^2$. At zero background the block average $Q_k(0)$ is the plain block average. The decomposition \cite[(1.27)]{Ba3} of $G_k(0)$ into the slices $G^{\eta_k}_{(j)}(0)$ is the normalisation group equation \cite[(2.43)]{Ba1}, built from the recursion \cite[(2.42)]{Ba1}, which is (b). Hence $G^{\eta_k}_{(0)}(0)=G_1$ and $G^{\eta_k}_{(j)}(0)=G_{j+1}-G_j$, and these are the slices of \cite[(2.6)]{Ba3}.

(g) For $k<K$ the divisibility hypothesis holds, so (c) gives the periodisation. The torus bounds in the proof of (e), which do not use (E.r) or (T.per), give on every torus of the family $\cT$ defined there the pointwise, H\"older and summability bounds of (H$_\infty$) for the slice $W_{(j)}$ and its derivatives of order $\le2$; for a single slice, including $W_{(0)}=G_1$, the summability bound holds for every $\abs a\le2$. By (d), with $\cV=\mathcal W$ and $\cV'$ the set of these kernels, the same bounds hold on $\eta\Z^3$.
\end{proof}

\begin{remark}\label{app:G1:vec-rem}
For $W=\partial^2G_i$ the summability bound holds only with a constant $O(1+i)$: summing the slice bounds gives $O(i)$, and the kernel behaves like the lattice Green function in the range $\eta\le\abs{x-x'}\le\ell_i$. Part (e) uses only the pointwise and H\"older bounds. A contour sum of a kernel $G_i$, or of a slice of index different from $j$, along $\Gamma^{(j+1)}$ need not satisfy (H$_\infty$) uniformly in $k$; such kernels lie in $\mathcal W^\sharp$ but not in $\mathcal W$. The full propagators $(-\Delta+\mu_0^2t_k^2)^{-1}$ and $F_{k,v}$ of \cite[p.~658]{K1} are not in $\mathcal W^\sharp$.
\end{remark}

\subsection{The final c-numbers: two-leg classes above $\eps_0$ and the transfer to $u\neq0$}\label{app:G1:ssec-Efin-classes}

\paragraph{Setting.} Throughout this subsection, $K\le j^*\le\bar K$, and we work in scale-$j^*$ units. So the lattice spacing is $\eta=\eps/t_{j^*}$ and $L^{j^*}\eta=1$. Write $t:=t_{j^*}$ and $\rho(y,y'):=\abs{y-y'}+\eta$, where $\abs{\cdot}$ is the torus distance. Let $\partial$ be the forward lattice difference and $\partial^*$ its adjoint in $\ell^2$ with weight $\eta^d$, so that $-\partial^*$ is the backward lattice difference. For a function $F$ on $T_\eps\times T_\eps$, let $N(F)$ be the maximum over $a,b\in\{0,1\}$ of $\sup\abs{\partial_1^a\partial_2^bF}$ plus the H\"older-$\alpha_0$ quotients of $\partial_1^a\partial_2^bF$ in either variable. For a function of one variable, $N$ is defined in the same way. If $F(y,z)=\Phi_1(y)\Phi_2(z)$, then $N(F)\le N(\Phi_1)N(\Phi_2)$. By \cite[(2.10)--(2.11), p.~426]{Ba3} and Theorem~\ref{app:G1:VSR-VS}, $N(C_{(j^*),f})\le C$ for $j^*<\bar K$. By Proposition~\ref{app:G1:VSR-R}, $N(\tilde R_f)\le C(1+Z_f)$.

The lines of a class considered below are pieces $C_{(i),f}$ (Definition~\ref{app:G1:vacgrade}) of index $\le j^*$ and $<\bar K$. Sum over these indices and telescope, as in the proof of Lemma~\ref{app:G1:common}. Each such line then becomes $G_{j_0,f}(0)$, with $j_0\in\{j^*,j^*+1\}$ and $j_0\le\bar K$. We split $G_{j_0,f}(0)=C^\xi_f+R_f$, where $\xi=L^{-j_0}$ and $R_f:=G^\xi_{j_0,f}(0)X_fC^\xi_f$. Here $C^\xi_f$ is the resolvent of $-\Delta^\xi+1$ on the rescaled torus. It is translation invariant, and it is the periodisation of the resolvent on $\xi\Z^3$.

Lemma~\ref{app:G1:AA} is stated for $0\le m^2_{j_0}\le1$, and its consequence at every Higgs mass is drawn under $t_{j_0}\le\eps_0$. Steps~1--3 of its proof use only Lemma~\ref{app:G1:VSR-F}, which does not involve the mass, and Theorem~\ref{app:G1:VSR-VS} at $t_{j_0}\le\min(1/m_f,L)$ and $m_f^2t_{j_0}^2\le1$. Both conditions hold for $j_0\le\bar K$ by (P1)--(P2). Hence these steps, together with the proof of Proposition~\ref{app:G1:VSR-5b}(4), give, for both fields and with some $\delta,\delta'>0$,
\begin{equation}\label{app:G1:Efin-L}
\begin{aligned}
&\abs{\partial^{\le2}C^\xi_f(z)}\le C\rho^{-1-\#\partial}\mathrm{e}^{-\delta\abs z},\qquad \abs{R_f(y,y')}\le C\mathrm{e}^{-\delta'\abs{y-y'}},\\
&\abs{\partial_iR_f(y,y')}\le C\big(1+\log_+\tfrac1\rho\big)\mathrm{e}^{-\delta'\abs{y-y'}},\qquad \abs{\partial_1\partial_2^*R_f(y,y')}\le C\rho^{-1}\mathrm{e}^{-\delta'\abs{y-y'}} .
\end{aligned}
\end{equation}
By ``the lemmas of the renormalisation machinery'' we mean Lemmas~\ref{app:G1:VSR-A}, \ref{app:G1:VSR-Btop} and \ref{app:G1:VSR-C}, Theorem~\ref{app:G1:VSR-T335} and Proposition~\ref{app:G1:VSR-5b}.

\smallskip\noindent\textup{(TM)} \emph{Torus moments.} Let $K\le j_0\le\bar K$ and $\xi=L^{-j_0}$. Consider the classes C3 and C4 of Lemma~\ref{app:G1:cutclasses} with every line replaced by $C^\xi_f$, together with the partners of their proper divergent subgraphs. Their zeroth moments at a fixed base point, and for C3 also their first moments, are bounded by a constant $C_{\rm TW}$, uniformly in $\xi$, in the base point and in $\ell_T\ge L^3$.
\smallskip

\noindent On $\xi\Z^3$ the C3 moments vanish \cite[(3.27)--(3.29), p.~441]{Ba3}, and the C4 moment is finite \cite[(3.31)--(3.32), p.~442]{Ba3}. On the torus these identities are not available, because the one-loop zeroth moment equals a sum over the non-zero images of the torus, which does not vanish in general. For C3, (TM) is Proposition~\ref{app:G1:TW}(i). For C4 it is Lemma~\ref{app:K:L04-2}.

\begin{lemma}[smoothing by a line]\label{app:G1:Efin-smooth}
Let $f\in\{s,v\}$, let $\ell=\sum_{i\le j'}C_{(i),f}$ with $j'\le j^*+1$, and let $b$ be bounded. Write $[g]_{\alpha_0}:=\sup_{y\neq y''}\abs{g(y)-g(y'')}/\abs{y-y''}^{\alpha_0}$. Then
\[
\norm{\ell b}_\infty+\norm{\partial(\ell b)}_\infty+[\partial(\ell b)]_{\alpha_0}\le C\norm b_\infty,\qquad C=C(L,\alpha_0,a).
\]
\end{lemma}
\begin{proof}
For $i<K$ use \cite[(2.10)--(2.11), p.~426]{Ba3}, which hold for vector propagators as well. For $i\ge K$ use Theorem~\ref{app:G1:VSR-VS}, which covers both fields. Together they give $\abs{\partial^aC_{(i)}(y,z)}\le Cs_i^{-1-a}\mathrm{e}^{-\delta_1\abs{y-z}/s_i}$ for $a\le2$, with $s_i:=L^i\eta$. Hence the H\"older quotient of $\partial C_{(i)}$ in $y$ is at most $Cs_i^{-2-\alpha_0}$ times the sum of the same exponentials at the two points. Moreover $\sum_z\eta^d\mathrm{e}^{-\delta_1\abs{y-z}/s}\le Cs^d$ for $s\ge\eta$. So the three quantities for $C_{(i)}b$ are at most $C\norm b_\infty s_i^2$, $C\norm b_\infty s_i$ and $C\norm b_\infty s_i^{1-\alpha_0}$, respectively. Since $s_i=L^{i-j^*}\le L$, the sums over $i\le j'$ are geometric.
\end{proof}

\begin{lemma}[pinned two-leg classes without a partner, above $\eps_0$]\label{app:K:L04-1}\label{app:G1:Efin-W}
Assume \eqref{app:G1:eq-hyp}. Let $Q$ be a two-leg class of type C1$'$, C3 or C4 of Lemma~\ref{app:G1:cutclasses}, nested and one-particle-reducible members included, whose lines are pieces $C_{(i),f}$ of index $\le j^*$ and $<\bar K$. Let both external legs carry a two-variable field $F$ with $N(F)<\infty$, and let $e(t)^\alpha\lambda(t)^\beta$ be the coupling factor of $Q$. Fix the vertex $v_0$ of one external leg at $x$. Let $V_Q(x)$ be the sum of the members together with the partners of their proper divergent subgraphs, over the positions of all vertices other than $v_0$ and over the indices. If $Q$ is of type C3 or C4, assume \textup{(TM)}. Then, uniformly in $x$, $\eps$ and $\ell_T\ge L^3$,
\[
\abs{V_Q(x)}\le C\,e(t)^\alpha\lambda(t)^\beta N(F),\qquad C=C(\nbar,d,L,a,\alpha_0,M_b,C_{\rm TW}).
\]
\end{lemma}
\begin{proof}
(i) \emph{The base point.} In \cite[pp.~435--442]{Ba3} the field of one external leg is expanded about the position of the other leg. Since the two-leg expressions are symmetric in the two legs, either leg may be expanded \cite[p.~435]{Ba3}. The photon classes are treated by the same transformations. Every later sum is taken at the fixed base: over the second external vertex, the internal vertices and the indices below $j_0$. The result is a function of the base \cite[pp.~440--442]{Ba3}. We take the base to be $x$, and we write $x'-x$ for the minimal displacement on the torus.

\emph{Class C3.} Summing along a shortest lattice path gives
\[
F(x,x')=F(x,x)+\textstyle\sum_\nu(x'-x)_\nu(\partial_{2,\nu}F)(x,x)+r(x,x'),\qquad\abs{r(x,x')}\le C_dN(F)\abs{x'-x}^{1+\alpha_0}.
\]
Hence
\[
V_Q(x)=F(x,x)\,\Pi(x)+\textstyle\sum_\nu\partial_{2,\nu}F(x,x)\,\Pi_\nu(x)+(\text{remainder}),
\]
where $\Pi(x)$ and $\Pi_\nu(x)$ are the zeroth and first moments at base $x$. The remainder has degree $\alpha_0>0$. Rescale to $\xi=L^{-j_0}$ and split every line as $C^\xi_f+R_f$. The terms with at least one $R_f$ are bounded by \eqref{app:G1:Efin-L}, as in Step~5 of the proof of Lemma~\ref{app:G1:AA}. The parts of $\Pi$ and $\Pi_\nu$ that contain only $C^\xi$ lines are bounded by (TM).

\emph{Class C4.} Write $F(x,x')=F(x,x)+[F(x,x')-F(x,x)]$, where $\abs{F(x,x')-F(x,x)}\le N(F)\abs{x'-x}^{\alpha_0}$, as in \cite[(3.31)]{Ba3}. Then argue as for C3, using (TM) for the zeroth moment.

\emph{Class C1$'$.} The derivative on one leg acts on $F$. Expand the differentiated factor $\partial_iF(x,\cdot)$ about $x$. This is allowed because $N(F)$ contains the H\"older quotients in either variable. The local coefficient is the zeroth moment of the kernel at base $x$.
\begin{itemize}
\item For the one-vertex graph of \cite[(3.21)]{Ba3}, this coefficient is $\eta\,G_{j_0,v}(0)$ near the diagonal, which is bounded.
\item For the exchange graph, its $C^\xi$ part is a multiple of
\[
\textstyle\sum_z\xi^3(\partial_\mu C^\xi)(z)\,C^\xi(z)=-\frac\xi2\sum_z\xi^3\abs{(\partial_\mu C^\xi)(z)}^2 .
\]
This identity holds because $\sum_zC^\xi(z+\xi e_\mu)^2=\sum_zC^\xi(z)^2$ on the torus. The right-hand side is bounded uniformly in $\xi$, since $\sum_z\xi^3\rho^{-4}\mathrm{e}^{-2\delta\abs z}\le C\xi^{-1}$. The rest of the coefficient is bounded by \eqref{app:G1:Efin-L}.
\end{itemize}
If $x$ is the differentiated end, use the symmetry of Lemma~\ref{app:G1:VSR-C}(iii).

(ii) \emph{Inner subgraphs.} A divergent proper subgraph that does not contain $v_0$ is transformed as in \cite[\S3]{Ba3}. Let $S\ni v_0$ be a divergent proper subgraph. Then the leg at $x$ is an external leg of $S$, and $x$ stays fixed while $S$ is transformed.
\begin{itemize}
\item If $S$ is of type C3 (a one-particle-reducible member of C4), base it at $x$ as in (i). The field on its other leg is the photon bridge applied to a bounded function. By Lemma~\ref{app:G1:Efin-smooth}, this field is bounded together with its first difference and the H\"older quotient of that difference.
\item If $S$ is a $\phi\phi A$ class through the vertex that carries the leg at $x$, use Theorem~\ref{app:G1:VSR-T335} with $r=x$, and Proposition~\ref{app:G1:VSR-5b}.
\end{itemize}
Nested subgraphs are handled by induction on the last divergent subgraph, with base $x$ at every level that contains $v_0$.

(iii) After (i) and (ii), every component $G_i^{(\alpha)}$ with $i<m$ has degree $\ge\vartheta>0$. Lemma~\ref{app:G1:VSR-A}(ii), pinned at $x$ and with the largest index fixed (Lemma~\ref{app:G1:VSR-Btop}), gives the bound. The field enters only through $N(F)$.
\end{proof}

\paragraph{Factorisation at a bridge.} Let $q$ be a real antisymmetric $N\times N$ matrix. Let $(C_{(j)})_j$ and $(D_{(j)})_j$ be finite families of real symmetric kernels on the lattice torus $T_\eta$, the scalar and vector slices. No translation invariance is assumed. Let $\phi^{(j)}=(\phi^{(j)}_a)_{a\le N}$ and $A^{(j)}$ be independent centred Gaussian fields with covariances $\delta_{jj'}\delta_{ab}C_{(j)}$ and $\delta_{jj'}D_{(j)}$. Put $\phi:=\sum_j\phi^{(j)}$ and $A:=\sum_jA^{(j)}$.
\begin{itemize}
\item A \emph{scalar factor} is $\partial^t\phi^{(j)}_a(u)$ with $t\in\{0,1\}$, where $\partial^1=\partial_\mu$. A \emph{photon factor} is $A^{(j)}_b$.
\item The interaction is $\sum_rg^{(r)}m_r$, where each $g^{(r)}$ is $e^\alpha\lambda^\beta$ times a real number, possibly depending on the positions. The \emph{order} of a vertex is $o(r)=\alpha/2+\beta>0$.
\item By the Gaussian pairing formula, a labelled graph with $n$ vertices has weight
\[
\frac{(-1)^n}{n!}\prod g^{(r)}\prod_{\text{lines}}\E_0[\cdot\,\cdot].
\] A scalar line joining $\partial^t\phi(u)$ to $\partial^s\phi(w)$ has value $(\partial_1^t\partial_2^sC_{(j)})(u,w)$. Fixing the indices of all lines amounts to fixing the slice labels of all factors.
\end{itemize}
We assume:
\begin{itemize}
\item[(V1)] The scalar factors of each $m_r$ are grouped into bilinears $\Phi\cdot M\Phi'$, and every photon factor is assigned to exactly one bilinear. Either $M=q^m$, where $m$ is the number of photon factors assigned to that bilinear, or $M$ is an even polynomial in $q$ and no photon factor is assigned to the bilinear.
\item[(V2)] Every vertex of order $\frac12$ has exactly one photon factor.
\end{itemize}
These hold for the vertices (1.6)--(1.8) and (1.10) of \cite{Ba3} at zero background, where $\abs\phi^4=(\phi\cdot\phi)(\phi\cdot\phi)$, and for the counterterm vertices below.

An \emph{amputated two-leg piece} is a labelled graph with $n\ge1$ vertices in which two scalar factors, the in-leg and the out-leg, are left unpaired. For a set $h$ of pieces that is closed under relabelling, let $\Sigma^{tt'}_h(u,w)_{ab}$ be $\eta^{-2d}$ times the sum of the weights of the members of $h$ with in-leg $\partial^t\phi_a(u)$ and out-leg $\partial^{t'}\phi_b(w)$. Put
\[
S_h:=\Sigma^{00}_h+\Sigma^{01}_{h,\mu}\partial_\mu+\partial^*_\mu\Sigma^{10}_{h,\mu}+\partial^*_\mu\Sigma^{11}_{h,\mu\nu}\partial_\nu,
\]
with operators acting through $\sum_w\eta^d$. Let $\mathcal H$ be the order-one one-particle-irreducible pieces without counterterm vertices. They are grouped into classes $h$ that are closed under relabelling and under the exchange of the two legs. Put $\delta m^2_h(u):=\sum_w\eta^d\Sigma^{00}_h(u,w)$, built with the same lines and indices. The partner of $h$ is the vertex $\frac12\sum_u\eta^d\phi(u)\cdot\delta m^2_h(u)\phi(u)$, and $\delta m^2_1:=\sum_h\delta m^2_h$.

\begin{lemma}[factorisation at a bridge]\label{app:G1:sliced-dyson}
Let $F$ be a function of two variables, and fix $x$ and the indices of all lines. Let $V(x)$ be $\eta^{-d}$ times the sum of the weights of the connected amputated two-leg pieces of order $2$ that have a bridge and have their in-leg at $x$, the two legs being contracted with $F$.
\begin{enumerate}
\item[(a)] The pieces all of whose bridges leave both legs on one side contribute $0$, already at fixed positions.
\item[(b)] Every other piece has exactly one bridge $\ell$ that separates the two legs, and $\ell$ is a scalar line of some slice $C_{(j)}$. Removing $\ell$ leaves $Q_1\ni x$ and $Q_2$. Up to terms that vanish by the argument of (a), each of them is of order one and one-particle irreducible, or is a single partner vertex. With $S_i:=\sum_{h\in\mathcal H}S_h-\delta m^2_1$ acting on the legs of piece $i$,
\[
V(x)=(S_1\,C_{(j)}\,S_2\,F^{\top})(x,x),
\]
where the derivatives on the legs contracted with $F$ act on $F$. No combinatorial factor is left over.
\item[(c)] $\delta m^2_h(u)$ is a symmetric matrix. It equals $\sum_w\eta^d\Sigma^{00}_h(w,u)^{\top}$, so it does not matter which leg is kept at $u$. Moreover, exactly,
\[
\textstyle\sum_{x'}\eta^d(S_h-\delta m^2_h)(y',x')=(\partial^*_\mu c_{h,\mu})(y'),\qquad c_{h,\mu}(y'):=\sum_{x'}\eta^d\Sigma^{10}_{h,\mu}(y',x'),
\]
and the same holds with $S_2$ and $\delta m^2_1$ in place of $S_h$ and $\delta m^2_h$.
\end{enumerate}
\end{lemma}
\begin{proof}
\emph{Parity.} Every vertex has an even number of scalar factors, and the lines pair scalar factors with scalar factors. So a connected set of vertices has an even number of unpaired scalar factors.

(a) Let $e$ be a bridge with both legs on one side, and let $P_e$ be the other side. By parity, the end of $e$ in $P_e$ is not scalar, so $e$ is a photon line, and every scalar factor of $P_e$ is paired inside $P_e$. The bilinears of $P_e$, joined by scalar lines, therefore form disjoint cycles. Scalar lines are diagonal in the internal index, so at fixed positions each cycle gives $\tr(M_1^\pm\cdots M_p^\pm)$, with $M^\pm\in\{M,M^\top\}$ and $(q^m)^\top=(-1)^mq^m$. The number of photon factors of $P_e$ is $2\,\#\{\text{photon lines of }P_e\}+1$, which is odd. By (V1), some cycle then carries an odd number of them, and its trace is a combination of the numbers $\tr q^m$ with $m$ odd. These vanish, since $\tr q^m=\tr(q^m)^\top=(-1)^m\tr q^m$.

(b) Let $\ell$ be a bridge separating the legs. By parity applied to the side of the in-leg, $\ell$ is scalar. Removing $\ell$ leaves two components $Q_1\ni x$ and $Q_2$, and with the ends of $\ell$ as new legs both are amputated two-leg pieces. Every vertex has order $\ge\frac12$, so a piece of order $\frac12$ is a single vertex; by (V2) it would have one unpaired photon factor, which is impossible. So $o(Q_i)\ge1$, and since $o(Q_1)+o(Q_2)=2$, both orders equal $1$. A second separating bridge would give three pieces of order $\ge1$, so $\ell$ is unique. Every other bridge of $Q_i$ is of the kind in (a), and the argument of (a) shows that a piece with such a bridge contributes $0$ at fixed positions.

A piece of order one that contains a counterterm vertex is that vertex alone. The partner of $h$ has two choices of in-leg, and its weight as a piece is $-\frac12\eta^d(\delta m^2_h+\delta m^{2\,\top}_h)=-\eta^d\delta m^2_h$, by (c).

Conversely, take an ordered partition $(S_1,S_2)$ of $\{1,\dots,n\}$ into non-empty blocks. Put pieces on $S_1$ and on $S_2$, relabelled increasingly, and join the out-leg of the first to the in-leg of the second. This is a bijection that preserves all slice labels. With $n_i=\abs{S_i}$,
\[
w(Q)=\tfrac{n_1!\,n_2!}{n!}\,w(Q_1)\,w(Q_2)\,\E_0[\text{ends of }\ell],
\]
and there are $n!/(n_1!n_2!)$ partitions with these block sizes. Hence the weights factorise with no factor left over. The operator form follows from $\sum_{y'}(\partial_{2,\mu}C)(y,y')g(y')=(C\partial^*_\mu g)(y)$.

(c) The fields are real, so exchanging the two legs of a labelled piece preserves its weight. Since $h$ is closed under this exchange, $\Sigma^{00}_h(u,w)_{ab}=\Sigma^{00}_h(w,u)_{ba}$.

In a non-vanishing term, the scalar path from the in-leg to the out-leg passes through bilinears whose product is $\pm q^m$, where $m$ counts the photon factors assigned to them. All other cycles carry an even number of photon factors, by the proof of (a), and the total number of photon factors is even. So $m$ is even, $q^m$ is symmetric, and $\delta m^2_h(u)$ is symmetric. The same argument shows that the counterterm vertices satisfy (V1).

Finally, $\partial\mathbf 1=0$ on $T_\eta$. Hence $S_h\mathbf 1=\Sigma^{00}_h\mathbf 1+\partial^*_\mu(\Sigma^{10}_{h,\mu}\mathbf 1)$, and $\Sigma^{00}_h\mathbf 1=\delta m^2_h$. Summing over $h$ gives the statement for $S_2$.
\end{proof}

\begin{lemma}[the coefficient is Lipschitz]\label{app:G1:Efin-lip}
Let $H_2$ be a one-particle-irreducible scalar two-leg piece of order one, and let $c(y):=\sum_{x'}\eta^d\Sigma_{2,1}(y,x')$, where $\Sigma_{2,1}$ is the kernel of $H_2$ with a derivative on the in-leg, the leg at $y$, i.e.\ the kernel $\Sigma^{10}$ of Lemma~\ref{app:G1:sliced-dyson}(c). With the lines of $H_2$ summed over their indices,
\[
\norm{\partial c}_\infty+\norm{\partial^*c}_\infty\le Ce(t)^2,
\]
uniformly in $\eps$ and $\ell_T$.
\end{lemma}
\begin{proof}
Only the exchange graph and the one-vertex graph of \cite[(3.21)]{Ba3} contribute to $c$.

\emph{The exchange graph.} It gives a finite sum of terms $e(t)^2\sum_{y'}\eta^dU(y,y')W(y,y')$, where $U$ is a photon line and $W=\ell_s\partial_2^*$; lattice shifts by $\eta e_\nu$ do not affect the argument. Split $U=C_U+R_U$ and $W=W_C+W_R$ as in \eqref{app:G1:Efin-L}. The product $C_UW_C$ is translation invariant, so its sum over $y'$ does not depend on $y$. For the other products use $\partial(fg)=(\partial f)g+f(\cdot+\eta e)\partial g$ in $y$, at fixed $y'$, together with \eqref{app:G1:Efin-L}:
\begin{itemize}
\item the difference of $C_UW_R$ is at most $C[\rho^{-2}(1+\log\frac1\rho)+\rho^{-2}]$ times an exponential;
\item the difference of $R_UW_R$ is at most $C[(1+\log\frac1\rho)^2+\rho^{-1}]$ times an exponential.
\end{itemize}
Both are summable uniformly in $\eta$.

In $R_UW_C$, the term $(\partial_1R_U)W_C$ is at most $C(1+\log\frac1\rho)\rho^{-2}$. For the other term write $R_U(y+\eta e,y')=R_U(y+\eta e,y)+\Delta$, where $\abs\Delta\le C\rho(1+\log\frac1\rho)$ by summing $\partial_2R_U$ along a lattice path. The sum over the torus of a lattice difference of a summable function vanishes, so $\sum_{y'}\partial_y[W_C(y-y')]=0$. Moreover $\abs\Delta\abs{\partial_yW_C}\le C\rho^{-2}(1+\log\frac1\rho)$, which is summable.

\emph{The one-vertex graph.} It gives $e(t)^2$ times a numerical coefficient times $\eta\,G_{j_0,v}(0)(y,y+\sigma)$. Its $C^\xi$ part does not depend on $y$, and $\eta\abs{\partial R_v}\le C\eta(1+\log\frac1\eta)\le C$.

The bound for $\partial^*$ is the same.
\end{proof}

\begin{lemma}[one-particle-reducible chains above $\eps_0$]\label{app:K:L04-3}\label{app:G1:Efin-1PR}
Assume the setting of Lemma~\ref{app:K:L04-1}. Let the two external legs of a scalar two-leg class of type C2 whose members are one-particle reducible carry a two-variable field $F$ with $N(F)<\infty$, and fix the external vertex at $x$. Let $V(x)$ be the pinned value of the class: the members together with the partners of their divergent subgraphs. Then, uniformly in $x$, $\eps$ and $\ell_T\ge L^3$,
\[
\abs{V(x)}\le C\sum_{\alpha+2\beta=4}\abs{e(t)}^\alpha\lambda(t)^\beta\,N(F),\qquad C=C(\nbar,d,L,a,\alpha_0,M_b).
\]
\end{lemma}
\begin{proof}
\emph{Step 1 (structure).} Apply Lemma~\ref{app:G1:sliced-dyson} at fixed indices. The members have order $2$ and a bridge, so up to terms that vanish they have exactly one bridge $\ell$, a scalar line. Removing it leaves $H_1\ni x$ and $H_2$, each of order $1$, and each either one-particle irreducible or a partner vertex. Moreover
\[
V(x)=(S_1\,\ell\,S_2\,F^\top)(x,x),\qquad S_i=-\delta m^2_{H_i}+\Sigma_i+\partial^*\Sigma_{i,1}+\Sigma_{i,1}^*\partial+\partial^*\Sigma_{i,2}\partial,
\]
in the notation of \cite[(1.21)]{Ba3}, where the derivatives on the legs attached to $F$ act on $F$. The coefficient of the partner of $H_i$ is $\delta m^2_{H_i}(y)=\sum_z\eta^d\Sigma_{H_i}(y,z)$, built with the same lines, as in \cite[p.~417]{Ba3}. By Lemma~\ref{app:G1:sliced-dyson}(c), or equivalently by the symmetry in Lemma~\ref{app:G1:VSR-C}(iii), it does not matter which leg is kept at $y$.

\emph{Step 2 (zeroth moment).} For every index assignment we have $\Sigma_{2,1}^*\partial\mathbf 1=\partial^*\Sigma_{2,2}\partial\mathbf 1=0$ and $-\delta m^2_{H_2}+\Sigma_2\mathbf 1=0$. By Lemma~\ref{app:G1:sliced-dyson}(c), therefore,
\[
(S_2\mathbf 1)(y)=(\partial^*c)(y),\qquad c(y):=\textstyle\sum_{x'}\eta^d\Sigma_{2,1}(y,x').
\]
One-vertex pieces with both legs at a single vertex and no derivative, namely the $\lambda$-tadpole and the photon-tadpole seagull, cancel their partners identically.

\emph{Step 3 (splitting).} The terms of $V$ with a derivative on an external leg have degree $1$, and all their proper subgraphs are renormalised. They are bounded by the lemmas of the renormalisation machinery. The remaining terms are
\[
V_{nn}(x)=\textstyle\sum_{x'}\eta^d\mathcal K(x,x')F(x,x'),\qquad \mathcal K:=T_1\,\ell\,(-\delta m^2_{H_2}+\Sigma_2+\partial^*\Sigma_{2,1}),\qquad T_1:=-\delta m^2_{H_1}+\Sigma_1+\Sigma_{1,1}^*\partial .
\]
Write $F(x,x')=F(x,x)+[F(x,x')-F(x,x)]$.
\begin{itemize}
\item The second part is the first term of \cite[(3.19)]{Ba3}, an identity valid for any kernel. It has degree $\alpha_0$, with the factor $\abs{x'-x}^{\alpha_0}$ split along the chain \cite[p.~438]{Ba3}. It is bounded by the lemmas of the renormalisation machinery.
\item By Step~2, the first part is $F(x,x)\,(T_1u)(x)$ with $u:=\ell(\partial^*c)$.
\end{itemize}
Steps~2 and 3 hold for each index assignment, so we may now sum over the indices of $\ell$ and of the lines of $H_2$.

\emph{Step 4 (the local term).} By Lemma~\ref{app:G1:Efin-lip}, $\norm{\partial^*c}_\infty\le Ce(t)^2$. By Lemma~\ref{app:G1:Efin-smooth} with $f=s$, $\norm u_{1,\alpha_0}:=\norm u_\infty+\norm{\partial u}_\infty+[\partial u]_{\alpha_0}\le Ce(t)^2$. The expression $T_1u$ is the renormalised two-leg class of $H_1$, pinned at $x$, with field $\equiv1$ there and field $u$ on the other leg. Lemma~\ref{app:G1:VSR-C} gives $\abs{(T_1u)(x)}\le C(e(t)^2+\lambda(t))\norm u_{1,\alpha_0}$. Finally, $\abs{F(x,x)}\le N(F)$. Together with Step~3, this gives the bound.
\end{proof}

\begin{remark}[use in the proof of Proposition~\ref{app:G1:Efin}]\label{app:G1:Efin-use}
Cut the top line $l(m)$ of a term.
\begin{itemize}
\item If $G\setminus l(m)$ is a one-particle-reducible member of C2, the term has no partner, and its class is bounded by Lemma~\ref{app:K:L04-3} with $F=l(m)$.
\item If $G\setminus l(m)$ is of type C1$'$, C3 or C4, its class is bounded by Lemma~\ref{app:K:L04-1} with $F=l(m)$, based at the pinned endpoint.
\end{itemize}
For $j^*<\bar K$ we have $N(F)\le C$. For $j^*=\bar K$, the two lemmas are applied in Lemma~\ref{app:G1:TB} to every lower cluster of these types. The field on the two legs of such a cluster is either one line $\tilde R_f$, or a product $\Phi_1(y)\Phi_2(z)$ of two such lines whose far ends are held fixed and later summed as in Lemma~\ref{app:G1:TB}. In both cases $N(F)\le C(1+Z_s+Z_v)^2$, and these factors are among those counted by $(1+Z_s+Z_v)^{N_\ell}$.
\end{remark}

\paragraph{The transfer to $u\ne0$.} In this paragraph $M\in(0,m_2]$, \eqref{app:G1:eq-hyp} holds, $\eps_0$ is a power of $L$, and the notation is that of Subsection~\ref{app:G1:part-D3}, of \eqref{eq:T6-GammaB} and of \eqref{app:G1:eq-Efinu}. We write $t_k=L^k\eps$, $\eta_j:=L^{-j}$ and $c_j:=\kappa\mu_*^2L^2t_j^2$. The $C$-parts, the $D$-parts ($D$-lines) and the expansion $E_{\rm fin}(u)=\sum_mu^mE^{(m)}_{\rm fin}$ are those of Lemma~\ref{app:G1:I8-colour}. The mass counterterm enters as the vertex \cite[(1.7)]{Ba3} whose coefficient is the number $\delta m^2_{\rm ref}$, built from the propagators at $u=0$. For fixed $\kappa$, the number $c_{\bar K}$ does not depend on $\eps$, $\eps_0$ or $\ell_T$ (Lemma~\ref{app:G1:I8-Dfac}).

\begin{lemma}[orderings of the coloured vacuum graphs]\label{app:G1:I8-order}
Let $m\ge1$ and $K\le j^*\le\bar K$. Expand the lines of the coloured graphs of $E^{(m)}_{\rm fin}$ by \eqref{app:G1:eq-N9split}, and keep the terms whose largest line index is $j^*$. In them, write every coefficient $\delta m^2_{\rm ref}$ as $\delta m^2_{\langle j^*\rangle}+\delta m^2_{\succ j^*}$, the first part in its graph representation and the second as a finite vertex, as in Step~1 of the proof of Proposition~\ref{app:G1:Efin}. Denote the resulting sum by $E^{(m),[j^*]}_{\rm fin}$. Fix an index assignment and an ordering $\bar l=(l(1),\dots,l(n))$ of all lines, attached lines included, with non-decreasing indices and with ties broken so that attached lines precede the others.
\begin{enumerate}
\item[(a)] $l(n)$ is a non-attached $C$-line of index $j^*$. Every $D$-line is non-attached and has index $<K$.
\item[(b)] Every connected proper subgraph that contains a $D$-line has $D^\natural\ge1$. In particular, $D^\natural(G\setminus l(n))\ge1$ whenever the term does not vanish.
\item[(c)] Let $S$ be a connected proper $C$-only subgraph all of whose internal lines have index below the lowest index $j_0$ of its external lines. Then:
\begin{itemize}
\item $S$ contains neither $l(n)$ nor an $R$-line, $j_0\le\min_fK_f$, and the internal lines of each field sum to $G_{j_0,f}(0)$;
\item the partner of $S$, namely the vertex carrying the graph $G_0\cong S$ with attached lines of index $<j_0$, is a term of $E^{(m),[j^*]}_{\rm fin}$ with the same $D$-lines and the same line $l(n)$;
\item conversely, every vertex carrying a graph of $\delta m^2_{\langle j^*\rangle}$ is the counterterm member of such a pair;
\item every connected proper subgraph that contains a finite vertex has degree $\ge\frac12$ or vanishes.
\end{itemize}
\item[(d)] Let $m=1$ and $o=1$. Then every $C$-only divergent component $G_r^{(\beta)}$, $r<n$, of $\bar l$ is a tree of type \cite[(2.4)]{Ba3}. If the $D$-line is replaced by a pair of external fields by \cite[(3.1)]{Ba3}, the remaining graph has largest index $j^*$, and its proper components are trees of type \cite[(2.4)]{Ba3}.
\end{enumerate}
\end{lemma}
\begin{proof}
(a) Attached lines belong to $\delta m^2_{\rm ref}$ and are $C$-lines. By Lemma~\ref{app:G1:I8-colour}, the term has a non-attached line of index $\ge K$, while $D$-lines have index $<K$. By the tie-break, the last line is a non-attached line of the largest index.

(b) A connected proper subgraph has an external leg, so the first claim is Lemma~\ref{app:G1:WP-deg}(c). By Lemma~\ref{app:G1:cutclasses}(c), $G\setminus l(n)$ is connected unless the term vanishes, and it contains every $D$-line.

(c) The internal lines of $S$ have index $<j_0\le j^*$, so $l(n)\notin S$. The next claims are Lemma~\ref{app:G1:common}. By Lemma~\ref{app:G1:VSR-Btop} with $j_{l(n)}=j^*$ fixed, $S$ is some $G_i$ with $i<n$. Its partner has attached lines of index $<j_0\le j^*$, so it is a term of $\delta m^2_{\langle j^*\rangle}$. The converse statement is Step~1 of the proof of Proposition~\ref{app:G1:Efin}. The partner operation changes only $S$ and the vertex, all of whose lines are $C$-lines of index $<j^*$, so it leaves the $D$-lines and $l(n)$ unchanged. The last claim is Lemma~\ref{app:G1:Efin-deg}, since $D^\natural\ge D$.

(d) Since $o=1$, $(\alpha,\beta)\in\{(2,0),(0,1)\}$. A $\lambda$-vertex has no vector leg, and a counterterm vertex of order $e^2$ closed by a scalar line gives $m=0$. So $G$ is one of three graphs: the bubble (two vertices \cite[(1.8)]{Ba3}, $n=1$), the seagull (the vertex \cite[(1.10)]{Ba3}, $n=2$), or the self-contraction of \cite[(1.8)]{Ba3} with $n=2$. Its only vector line is the $D$-line, of index $i<K$. Hence some scalar line has index $\ge K>i$.

A component $G_r^{(\beta)}$ consists of the lines in the first $r$ positions of $\bar l$. If it is $C$-only, it lies before the $D$-line, so all its lines are scalar lines of index $\le i$.
\begin{itemize}
\item In the seagull and the self-contraction, the only scalar line has index $\ge K$, so there is no such component.
\item In the bubble, with scalar lines $a,b$ and $\max(a,b)\ge K$, the only possibility is $\{a\}$ with $a\le i$. This is a tree with one internal scalar line, two external vector legs and two external scalar legs, i.e.\ of type \cite[(2.4)]{Ba3}.
\end{itemize}
If the $D$-line of the bubble is replaced by external fields, the remaining graph is the scalar loop $\{a,b\}$. Its largest index is $j^*$, and its proper components $\{a\}$, $\{b\}$ are trees of type \cite[(2.4)]{Ba3}. In the seagull and the self-contraction the $D$-line joins a vertex to itself and is never replaced.
\end{proof}

\begin{lemma}[coloured families at a fixed largest index]\label{app:G1:I8-fixed}
Let $m\ge1$ and $K\le j^*\le\bar K$. Let $\mathcal F$ be a family of terms of $E^{(m),[j^*]}_{\rm fin}$ of order $(\alpha,\beta)$, i.e.\ a class of the partner relation of Lemma~\ref{app:G1:I8-order}(c). Fix one vertex in a block of side $t_{j^*}$. Then the sum of the terms of $\mathcal F$ is, in absolute value, at most
\[
C_{\mathcal F}\,\Phi\,\abs e^\alpha\lambda^\beta t_{j^*}^{\,o}\,(\abs u\,c_{j^*})^m\quad(j^*<\bar K),\qquad
C_{\mathcal F}^R\abs e^\alpha\lambda^\beta t_{\bar K}^{\,o}(1+Z_s+Z_v)^{N_\ell}(\abs u\,c_{\bar K})^m\quad(j^*=\bar K).
\]
Here $\Phi$ is as in Proposition~\ref{app:G1:Efin}, $C_{\mathcal F}=C(\mathcal F,\alpha_0,\nbar,d,L,a,M_b)$ and $C_{\mathcal F}^R=C(\mathcal F,\alpha_0,\nbar,d,L,a,M_b)(1+\rho_{\min}^{-2}+\kappa_{R,\min}^{-3})^{N_\ell}$. Both constants are independent of $\eps$, $\eps_0$, $\ell_T$, $u$, $\kappa$ and $\mu_*$.
\end{lemma}
\begin{proof}
Work in scale-$j^*$ units, $\eta=\eta_{j^*}$.

\emph{Lines and vertices.} Every $C$-line of index $j<\bar K$ is a slice of $G_{K,f}(0)$, or an $H_{j,f}$ with $t_{j+1}\le t_{\bar K}\le\min(1/m_f,L)$. So it obeys \cite[(2.10)--(2.12)]{Ba3} at zero background, by Theorem~\ref{app:G1:VSR-VS} and Corollary~\ref{app:G1:VSR-VSphys}. A line of index $\bar K$ obeys the bounds of Proposition~\ref{app:G1:VSR-R}. By Lemma~\ref{app:G1:I8-Dfac}, a $D$-line of index $i$ equals $u\,c_{j^*}$ times $(L^i\eta)^2H_i$. By Lemma~\ref{app:G1:WP-deg}(a), this kernel obeys \cite[(2.10)--(2.12)]{Ba3} with line dimension $-\mathsf d+4$. A finite vertex carries $\delta m^2_{\succ j^*}t_{j^*}^2$, which is bounded by Lemma~\ref{app:G1:Efin-fv}. After the factor $(uc_{j^*})^m$ is taken out, every term of $\mathcal F$ is a generalised graph in the sense of Lemma~\ref{app:G1:VSR-A}, of degree $D^\natural$.

\emph{Renormalisation.} The subgraphs of Lemma~\ref{app:G1:I8-order}(c) are renormalised together with their partners, which lie in $\mathcal F$. This is done by the lemmas of the renormalisation machinery, as in Step~1 and at the end of Step~2 of the proof of Proposition~\ref{app:G1:Efin}. These operations act on $C$-parts, whose index sums are those at $u=0$. A $D$-line meets them only as an external leg, which is regular by the derivative and H\"older bounds of Lemma~\ref{app:G1:WP-deg}(a).

Trees of type \cite[(2.4)]{Ba3}, including those with $D$-line legs, are integrated by parts by \cite[(2.8)]{Ba3}. A line, $C$ or $D$, between localisations at distance at least $1$ becomes a pair of external fields by \cite[(3.1)]{Ba3}. The components that remain are subgraphs of the original graph, as in steps (1)--(6) of the proof of Theorem~\ref{app:G1:WP-WP1}. In the case of Lemma~\ref{app:G1:I8-order}(d), the remaining graph keeps the largest index $j^*$.

By Lemma~\ref{app:G1:I8-order}(b)--(d), and since the transformations of \cite[\S3]{Ba3} do not lower the degree of a larger subgraph, every component $G_r^{(\beta)}$, $r<n$, of every admissible ordering has positive degree. No condition is imposed on the degree of the whole graph.

\emph{The bound for $j^*<\bar K$.} By Lemma~\ref{app:G1:VSR-Btop}, the terms form a sub-sum of a majorant sum. Lemma~\ref{app:G1:VSR-A}(ii), with the largest index fixed, bounds each of them by $C_A(L^{j^*}\eta)^{D^\natural}=C_A$, times $\Gamma_f^{m_f}\le\Phi$ for the $m_f$ finite vertices. Multiply by $\abs{e(t_{j^*})}^\alpha\lambda(t_{j^*})^\beta=\abs e^\alpha\lambda^\beta t_{j^*}^{\,o}$, by $(\abs uc_{j^*})^m$ and by the combinatorial factors of $\mathcal F$. This gives the first bound.

\emph{The bound for $j^*=\bar K$.} Here $\delta m^2_{\succ\bar K}=0$, so there are no finite vertices. The lines of index $\bar K$ do not depend on $u$, and every $D$-line lies in a lower cluster. Each lower cluster, with its renormalisations, has components of positive degree. So Lemma~\ref{app:G1:VSR-A}(ii) bounds it uniformly in the pin; the $C$-only lower clusters of the types of Lemmas~\ref{app:K:L04-1} and \ref{app:K:L04-3} are bounded as in Remark~\ref{app:G1:Efin-use}. Lemma~\ref{app:G1:TB} then gives the second bound, with the polynomial in $\rho_f^{-1},\kappa_R^{-1}$ evaluated at $\rho_{\min}$.
\end{proof}

\begin{proposition}[$E_{\rm fin}(u)$ at $u\ne0$]\label{app:K:L04-4}\label{app:G1:I8-transfer}
Let $M\in(0,m_2]$, assume \eqref{app:G1:eq-hyp}, let $E_{\rm fin}$ be graded as in Definition~\ref{app:G1:vacgrade}, and let $\eps_0$ be a power of $L$. Let $E_{\rm fin}(u)$ be as in \eqref{app:G1:eq-Efinu}, and put
\[
\bar\sigma_N(u):=\big(1+\abs u\,\kappa\mu_*^2L^2t_{\bar K}^2\big)^{N_\ell}\le\big(1+\abs u\,\kappa L^2\big)^{N_\ell}.
\]
Let $\hat B$ be the right-hand side of Proposition~\ref{app:G1:Efin}, including the factor $\Phi$, with $C_E,C_E^R$ replaced by constants $\hat C_E=C(\alpha_0,\nbar,d,L,a,M_b)\ge C_E$ and $\hat C_E^R=C(\alpha_0,\nbar,d,L,a,M_b)(1+\rho_{\min}^{-2}+\kappa_{R,\min}^{-3})^{N_\ell}\ge C_E^R$. Then, for every $u\in\R$ and every $\eps=2^{-K}\eps_0$,
\[
\abs{E_{\rm fin}(u)-E_{\rm fin}(0)}\le\big(\bar\sigma_N(u)-1\big)\hat B,\qquad \abs{E_{\rm fin}(u)}\le\bar\sigma_N(u)\,\hat B .
\]
For fixed $\kappa$, the factor $\bar\sigma_N(u)$ depends only on $u,\kappa,\mu_*,M,L$ and $\nbar$.
\end{proposition}
\begin{proof}
By Lemma~\ref{app:G1:I8-colour}, $E_{\rm fin}(u)-E_{\rm fin}(0)=\sum_{m=1}^{N_\ell}u^mE^{(m)}_{\rm fin}$, and $E^{(m)}_{\rm fin}=\sum_{j^*=K}^{\bar K}E^{(m),[j^*]}_{\rm fin}$. By Lemma~\ref{app:G1:I8-order}(c), each $E^{(m),[j^*]}_{\rm fin}$ is a disjoint union of families.

Let $n_{\alpha,\beta}$ be the number of graphs of order $(\alpha,\beta)$ with all lines uncoloured, each counterterm vertex either opened or finite; it depends only on $\nbar$. A coloured graph with $m$ $D$-lines is such a graph with $m$ of its at most $N_\ell$ non-attached lines marked. So for each $(\alpha,\beta,m,j^*)$ there are at most $\binom{N_\ell}mn_{\alpha,\beta}$ families.

There are $\abs{T_\eps}t_{j^*}^{-3}$ blocks of side $t_{j^*}$. By Lemma~\ref{app:G1:I8-fixed} and $c_{j^*}\le c_{\bar K}$ (Lemma~\ref{app:G1:I8-Dfac}), a family of order $(\alpha,\beta)$, summed over $j^*$, contributes at most
\[
\abs{T_\eps}(\abs uc_{\bar K})^m\abs e^\alpha\lambda^\beta\Big[C_{\mathcal F}\,\Phi\sum_{j^*=K}^{\bar K-1}t_{j^*}^{\,o-3}+C_{\mathcal F}^R\,t_{\bar K}^{-3}t_{\bar K}^{\,o}(1+Z_s+Z_v)^{N_\ell}\Big].
\]
The sum over $j^*$ is bounded as in the proof of Proposition~\ref{app:G1:Efin}:
\begin{itemize}
\item for $o\le2$ by $L^3\eps_0^{o-3}$;
\item for $o=3$ by $2+\log_L^+(\bar\ell/\eps_0)$;
\item for $o\ge4$ by $t_{\bar K}^{o-3}/(1-L^{-(o-3)})$.
\end{itemize}
Let $\hat C_E$ be the maximum of $C_E$ and $\max_{\alpha,\beta}n_{\alpha,\beta}\max_{\mathcal F}C_{\mathcal F}$. Define $\hat C_E^R$ in the same way, with $C_E^R$ and $C_{\mathcal F}^R$. Summing over families, orders and $m$ gives
\[
\abs{E_{\rm fin}(u)-E_{\rm fin}(0)}\le\sum_{m=1}^{N_\ell}\binom{N_\ell}m(\abs uc_{\bar K})^m\,\hat B=\big(\bar\sigma_N(u)-1\big)\hat B .
\]
The second inequality follows from $\abs{E_{\rm fin}(0)}\le\hat B$, which is Proposition~\ref{app:G1:Efin}. The bound on $\bar\sigma_N$ and the last sentence follow from Lemma~\ref{app:G1:I8-Dfac}.
\end{proof}

\subsection{The split of $\theta_{k+1}A^{(k)}$ and exceptional vertices with a chain leg}\label{app:G1:ssec-T579split}

This subsection proves two statements used in the proof of Proposition~\ref{app:G1:T579}. The first is the exact form of the split of $\theta_{k+1}A^{(k)}$ in \cite[(2.89)]{Ba2}. The second is the bound for an exceptional vertex one of whose scalar legs is a chain. Throughout, $\ell(t)=1+\log t^{-1}$, $r(t)=R\ell(t)^2$, $\hat\ell:=\ell(\eps_0)$ and $t_k=L^k\eps$, so that $t_K=\eps_0$. For a kernel $K(x,y)$ and a function $\zeta(x,y)$, the kernel product $\zeta K$ is $(\zeta K)(x,y):=\zeta(x,y)K(x,y)$. Operators to the right of a kernel product are composed with it in the usual way.

\begin{proposition}[exactness of the split of $\theta_{k+1}A^{(k)}$ at $L=2$]\label{app:K:L06-1}\label{app:G1:T579-zeta}
Let $L=2$ and $1\le k\le K-1$. Measure lengths in units of $T_1^{(k)}$, with the Euclidean distance of the torus. We use the following conventions of \cite{Ba2}.
\begin{itemize}
\item[(i)] The sets $\Lambda_i^{(k)}\subset T_1^{(k)}$ are built from $r(t_k)$ as in \cite[(2.7)--(2.8), p.~558; p.~566]{Ba2}: $(\Lambda_{i+1}^{(k)})^c$ is the union of the large blocks (of side $M_b$) at distance $\le r(t_k)$ from $(\Lambda_i^{(k)})^c$.
\item[(ii)] For a union $X$ of $L$-blocks, $X'$ is the set of $y\in T_L^{(k)}$ with $B(y)\subset X$. Every point of $B(y)$ lies within $L\sqrt d$ of $y$.
\item[(iii)] $\theta_{k+1}=1$ on $B^k(\Lambda_2^{(k)})$, and $\theta_{k+1}$ vanishes at distance $\ge LM_b$ from it.
\item[(iv)] $\zeta^{(k+1)}(x,y)=0$ for $\abs{x-y}\ge L\big(r(t_{k+1})-2M_b\big)$. Here the support radius of \cite[(2.44), p.~566]{Ba2} is measured in units of $T_1^{(k+1)}$.
\end{itemize}
Let $A'$, $B$, $\delta C^{(k)}_{\Lambda_0^{(k)}}=C^{(k)}_{\Lambda_0^{(k)}}-C^{(k)}$, $G_k$, $G^\eta_{k+1}$, $Q_k^*$, $Q^*_{k+1}$, $Q^*$, $a_k$, $a_{k+1}$ and $a$ be as in \cite[(2.80)--(2.89), pp.~574--575]{Ba2}. Assume
\begin{equation}\label{app:G1:eq-splitZ}
r(\eps_0)\ge 2M_b,\qquad M_b+2R\log 2\;\hat\ell\ \ge\ 2\sqrt d .
\end{equation}
The second inequality holds whenever $R\ge\sqrt d/\log2$, because $\hat\ell\ge1$. Then $\zeta^{(k+1)}(x,y)=0$ whenever $\theta_{k+1}(x)\neq0$ and $y\notin\Lambda_0^{(k)\prime}$. Moreover, identically in $A'$ and $B$,
\begin{equation}\label{app:G1:eq-split289}
\begin{aligned}
\theta_{k+1}A^{(k)}={}&a_k\theta_{k+1}\zeta^{(k)}G_kQ_k^*A'+a_kaL^{-2}\theta_{k+1}\zeta^{(k)}G_kQ_k^*\,\delta C^{(k)}_{\Lambda_0^{(k)}}Q^*B\\
&-a_kaL^{-2}\theta_{k+1}(1-\zeta^{(k)})G_kQ_k^*C^{(k)}Q^*\Lambda_0^{(k)\prime}B
+a_{k+1}L^{-2}\theta_{k+1}(1-\zeta^{(k+1)})G^\eta_{k+1}Q^*_{k+1}\Lambda_0^{(k)\prime}B\\
&+a_{k+1}L^{-2}\theta_{k+1}\zeta^{(k+1)}G^\eta_{k+1}Q^*_{k+1}B .
\end{aligned}
\end{equation}
This is \cite[(2.89)]{Ba2} as printed.
\end{proposition}

\begin{proof}
\emph{Step 1: a layer estimate.} Write $r:=r(t_k)$. We claim that $\abs{x_0-z}>2r-\sqrt d$ whenever $x_0\in\Lambda_2^{(k)}$ and $z\notin\Lambda_0^{(k)}$.

Every lattice point $w$ lies in a large block, and that block is at distance $\le\dist(w,(\Lambda_0^{(k)})^c)$ from $(\Lambda_0^{(k)})^c$. By (i), therefore, $\{w:\dist(w,(\Lambda_0^{(k)})^c)\le r\}\subset(\Lambda_1^{(k)})^c$. In the same way, $\{w:\dist(w,(\Lambda_1^{(k)})^c)\le r\}\subset(\Lambda_2^{(k)})^c$, and in particular $(\Lambda_1^{(k)})^c\subset(\Lambda_2^{(k)})^c$.

Let $D:=\abs{x_0-z}$ and suppose that $D\le2r-\sqrt d$. Choose lifts of $x_0$ and $z$ to the covering lattice at Euclidean distance $D$; distances on the torus are at most the distances between lifts.
\begin{itemize}
\item If $D\le r$, then $x_0\in(\Lambda_1^{(k)})^c\subset(\Lambda_2^{(k)})^c$, which contradicts $x_0\in\Lambda_2^{(k)}$.
\item Otherwise $r<D\le2r-\sqrt d$, so $r>\sqrt d$. Let $P$ be the point of the segment from $z$ to $x_0$ at distance $r-\frac{\sqrt d}2$ from $z$, and let $w$ be a lattice point with $\abs{w-P}\le\frac{\sqrt d}2$. Then $\abs{w-z}\le r$, so $w\in(\Lambda_1^{(k)})^c$. Also $\abs{w-x_0}\le D-r+\sqrt d\le r$, so $x_0\in(\Lambda_2^{(k)})^c$, which is again a contradiction.
\end{itemize}

\emph{Step 2: the supports.} Let $\theta_{k+1}(x)\neq0$ and $y\notin\Lambda_0^{(k)\prime}$.
\begin{itemize}
\item By (iii), $\dist(x,B^k(\Lambda_2^{(k)}))<LM_b$. Every point of $B^k(x_0)$ lies within $\sqrt d$ of $x_0$, so there is $x_0\in\Lambda_2^{(k)}$ with $\abs{x-x_0}<LM_b+\sqrt d$.
\item By (ii), $B(y)\not\subset\Lambda_0^{(k)}$, so there is $z\notin\Lambda_0^{(k)}$ with $\abs{y-z}\le L\sqrt d$.
\end{itemize}
By Step 1 and $L=2$,
\[
\abs{x-y}\ge\abs{x_0-z}-\abs{x-x_0}-\abs{y-z}>2r(t_k)-2M_b-4\sqrt d .
\]
Put $\ell_j:=\ell(t_j)$. Since $t_{k+1}=2t_k$, we have $\ell_{k+1}=\ell_k-\log2$. Since $t_k<t_{k+1}\le t_K=\eps_0$ and $\ell$ is decreasing, $\ell_k\ge\ell_{k+1}\ge\hat\ell$. Hence
\[
r(t_k)-r(t_{k+1})=R(\ell_k-\ell_{k+1})(\ell_k+\ell_{k+1})=R\log2\,(\ell_k+\ell_{k+1})\ge2R\log2\,\hat\ell .
\]
By the second inequality of \eqref{app:G1:eq-splitZ},
\[
2r(t_k)-2M_b-4\sqrt d\ \ge\ 2r(t_{k+1})+4R\log2\,\hat\ell-2M_b-4\sqrt d\ \ge\ 2r(t_{k+1})-4M_b=L\big(r(t_{k+1})-2M_b\big).
\]
This radius is nonnegative by the first inequality of \eqref{app:G1:eq-splitZ}, since $r(t_{k+1})\ge r(\eps_0)$. So $\abs{x-y}\ge L(r(t_{k+1})-2M_b)$, and (iv) gives $\zeta^{(k+1)}(x,y)=0$.

\emph{Step 3: the identity.} In units of $T_1^{(k)}$, \cite[(2.44), (2.88)]{Ba2} read $A^{(k)}=a_k\zeta^{(k)}G_kQ_k^*A$, and $A=A'+aL^{-2}C^{(k)}_{\Lambda_0^{(k)}}Q^*B$ by \cite[(2.80)]{Ba2}. The first equality of \cite[(2.82)]{Ba2} contains terms in $B(y)$. After these terms cancel, it is equivalent to
\[
C^{(k)}_{\Lambda_0^{(k)}}Q^*B=C^{(k)}Q^*\Lambda_0^{(k)\prime}B+\delta C^{(k)}_{\Lambda_0^{(k)}}Q^*B .
\]
Inserting this, we get
\[
\theta_{k+1}A^{(k)}=a_k\theta_{k+1}\zeta^{(k)}G_kQ_k^*A'+a_kaL^{-2}\theta_{k+1}\zeta^{(k)}G_kQ_k^*\,\delta C^{(k)}_{\Lambda_0^{(k)}}Q^*B+a_kaL^{-2}\theta_{k+1}\zeta^{(k)}G_kQ_k^*C^{(k)}Q^*\Lambda_0^{(k)\prime}B .
\]
In the last term write $\zeta^{(k)}=1-(1-\zeta^{(k)})$. The kernel product is linear in $\zeta^{(k)}$, so the last term becomes
\[
a_kaL^{-2}\theta_{k+1}G_kQ_k^*C^{(k)}Q^*\Lambda_0^{(k)\prime}B-a_kaL^{-2}\theta_{k+1}(1-\zeta^{(k)})G_kQ_k^*C^{(k)}Q^*\Lambda_0^{(k)\prime}B .
\]
Now use the composition identity
\[
a_kaL^{-2}G_kQ_k^*C^{(k)}Q^*=a_{k+1}L^{-2}G^\eta_{k+1}Q^*_{k+1}.
\]
This is the vector-field analogue of \cite[(2.41), p.~612]{Ba1}, taken in adjoint form, and it follows by the same computation (compare the use of \cite[(I.2.43)]{Ba1} for the vector field at \cite[p.~414]{Ba3}). The first term above therefore equals $a_{k+1}L^{-2}\theta_{k+1}G^\eta_{k+1}Q^*_{k+1}\Lambda_0^{(k)\prime}B$. Split $1=(1-\zeta^{(k+1)})+\zeta^{(k+1)}$ in it, and in the $\zeta^{(k+1)}$-part write $\Lambda_0^{(k)\prime}=1-(1-\Lambda_0^{(k)\prime})$. This gives
\[
\theta_{k+1}A^{(k)}=(\text{right side of }\eqref{app:G1:eq-split289})-E,\qquad
E:=a_{k+1}L^{-2}\theta_{k+1}\zeta^{(k+1)}G^\eta_{k+1}Q^*_{k+1}(1-\Lambda_0^{(k)\prime})B .
\]
The kernel of $E$ at $(x,y)$, with $y\in T_L^{(k)}$, is
\[
a_{k+1}L^{-2}\theta_{k+1}(x)\,\zeta^{(k+1)}(x,y)\,(G^\eta_{k+1}Q^*_{k+1})(x,y)\,\mathbf 1_{y\notin\Lambda_0^{(k)\prime}} .
\]
By Step 2 this vanishes for all $(x,y)$. Hence $E\equiv0$, identically in $A'$ and $B$.
\end{proof}

\medskip
\noindent\emph{Setting for Lemma~\ref{app:K:L06-2}.} Let $0\le k\le K-1$, $t=t_k$, $p=p(t)=b_0\ell(t)^{p_0}$, $e(t)=e\,t^{1/2}$ and $\lambda(t)=\lambda t$, and work on $\supp\chi_k$ with the notation of the proof of Proposition~\ref{app:G1:T579}:
\begin{itemize}
\item $\mathfrak B$ is the background and $\tilde A$ the expansion field;
\item $\sigma:=e(t)(\norm{\tilde A}_\infty+\norm{\partial^\eta\tilde A}_\infty)$ satisfies $\sigma\le C_\sigma t^{1/2}p$;
\item the scalar fields $f$ of the step satisfy $\norm f_\infty\le X'_s:=c_1'p\lambda(t)^{-1/4}$, with $c_1'=2^{1/4}(1+\log2)^6c_1$, by the restrictions \cite[(2.55)]{Ba2};
\item $\Phi_0:=a_kG_k(\Omega_*,\mathfrak B)Q_k^*(\mathfrak B)\phi$, where $\Omega_*$ is the domain of the scalar propagators of the family.
\end{itemize}
The norm of an external leg is the norm $\norm{h\,\cdot}_{1,\alpha_0}$ of \cite[(1.32)]{Ba3}.

\begin{lemma}[exceptional vertices with a chain leg]\label{app:K:L06-2}\label{app:G1:T579-chainleg}
In the setting above, assume the conclusions of Lemma~\ref{app:G1:T579-reg}(c). Let $v$ be a vertex of coupling order one of a connected localised family of $\mathcal P^{(k)}_w$ with at least two vertices, both of whose scalar legs are external fields. Let one of these legs be a chain $\mathcal Kf$ of order $m\ge1$, that is, a term with at least one insertion of the expansions (E1)--(E2), in the proof of Proposition~\ref{app:G1:T579}, of $G_k(\Omega_*,\mathfrak B+\tilde A)$ and $Q_k(\mathfrak B+\tilde A)$. Then $v$ together with this chain contributes to the bound \cite[(1.33)]{Ba3} at most
\[
C\,t^{(1+m)/4}p^{2+m},
\]
with $C$ independent of $\eps$, $k$ and $\ell_T$.
\end{lemma}

\begin{proof}
The bound \cite[(1.33)]{Ba3} carries the factor $e(t)=e\,t^{1/2}$ for the vertex $v$, and it is linear in the norm \cite[(1.32)]{Ba3} of each external leg of $v$. So it suffices to bound the norms of the two scalar legs.

\emph{The chain.} By Proposition~\ref{app:G1:VSR-K3}, the sup, the derivative and the $\alpha_0$-H\"older quotient of $\mathcal Kf$ are $\le C(C_\infty\sigma)^m\norm f_\infty$. With $\sigma\le C_\sigma t^{1/2}p$ and $\norm f_\infty\le X'_s=c_1'p\lambda^{-1/4}t^{-1/4}$, this gives
\[
\norm{h\mathcal Kf}_{1,\alpha_0}\le C\,(C_\infty C_\sigma t^{1/2}p)^m\,p\lambda^{-1/4}t^{-1/4}.
\]

\emph{The other leg.} This leg is $\Phi_0$, or $\phi(y)$, or a chain of order $m'\ge1$.
\begin{itemize}
\item For $\Phi_0$: by \cite[Prop.~2.2]{Ba2} and the H\"older bounds \cite[(1.8), (1.10), (1.12)]{Ba4} for $G_k(\Omega_*,\mathfrak B)Q_k^*(\mathfrak B)$, its norm is $\le CX'_s$.
\item For $\phi(y)$: its norm is $\le CX'_s$ by \cite[(2.55)]{Ba2}.
\item For a chain of order $m'$: by the bound just proved, its norm is $\le C(C_\infty C_\sigma t^{1/2}p)^{m'}p\lambda^{-1/4}t^{-1/4}$. Here $t^{1/2}p(t)\le\sup_{0<s\le1}b_0s^{1/2}\ell(s)^{p_0}<\infty$, and $m'$ is bounded by a number depending only on $\nbar$, by the stop rule of the expansion. So this norm is also $\le CX'_s$.
\end{itemize}

\emph{Collection.} The contribution of $v$ together with the chain is therefore at most
\[
e\,t^{1/2}\cdot CX'_s\cdot C(C_\infty C_\sigma t^{1/2}p)^mp\lambda^{-1/4}t^{-1/4}
\le C\lambda^{-1/2}\,t^{m/2}\,p^{2+m}\le C\,t^{(1+m)/4}p^{2+m}.
\]
The last step uses $m/2\ge(1+m)/4$, which holds for $m\ge1$, and $t\le1$.

The constants $C_\infty$, $C_\sigma$, $c_1'$, the constants of Proposition~\ref{app:G1:VSR-K3}, those of \cite[Prop.~2.2]{Ba2} and those of \cite{Ba4} are independent of $\eps$, $k$ and $\ell_T$, and $m$ is bounded in terms of $\nbar$. Hence $C$ is independent of $\eps$, $k$ and $\ell_T$.
\end{proof}

\subsection{King's slices and the order-$e^2$ anchor bounds}\label{app:G1:kingslices}

In this subsection $j$ and $k$ are slice indices, not torus indices. $T_\eps=\eps\Z^3/\ell_T\Z^3$ is one fixed torus, and we write $\ell:=\ell_T$.

\paragraph{Setting.} Let $L=2$, $a>0$, $\eps>0$, $K\ge2$ and $t_j:=L^j\eps$. Assume $\ell/t_K\in\N$, and put $T^{(j)}:=T_\eps\cap t_j\Z^3$ for $0\le j\le K$. All operators below act in the same way on each real field component, so we state everything for one real component. Following \cite[(2.2)]{K1}, we use the scalar products
\[
\ip fg=\sum_{x\in T_\eps}\eps^3f(x)g(x),\qquad \ipj jB{B'}=\sum_{y\in T^{(j)}}t_j^3B(y)B'(y),
\]
with norms $\norm\cdot$ and $\norm\cdot_j$. Put $\nabla^+_\mu f(x)=\eps^{-1}(f(x+\eps e_\mu)-f(x))$ and $\nabla^-_\mu f(x)=\eps^{-1}(f(x)-f(x-\eps e_\mu))$, and write $\nabla^{\pm*}_\mu$ for their adjoints. Then $\nabla^{+*}_\mu=-\nabla^-_\mu$ and $-\Delta_\eps=\sum_\mu\nabla^{+*}_\mu\nabla^+_\mu$. Hence $\ip f{-\Delta_\eps f}=\sum_x\eps^3\sum_\mu\abs{\nabla^+_\mu f(x)}^2$ \cite[(2.3)]{K1}.

For $m>0$ put $Y_m:=-\Delta_\eps+m^2$ and $C_m:=Y_m^{-1}$. With $\delta_x:=\eps^{-3}\mathbf 1_{\{x\}}$, the kernel of an operator $B$ is $B(x,y)=\ip{\delta_x}{B\delta_y}$. For $B\ge0$ we put $s_0(B):=\sup_xB(x,x)$.

The block of $y\in T^{(j)}$ is $\square_j(y)=\{x:y_i\le x_i<y_i+t_j\}$. We follow \cite[(2.10), (2.13)--(2.17)]{K1} at zero background, with photon mass $\mu_*>0$ (written $\mu_0$ in \cite{K1}). Define:
\begin{itemize}
\item the block averages $(Q_jf)(y)=t_j^{-3}\sum_{x\in\square_j(y)}\eps^3f(x)$ for $1\le j\le K$, and $Q$, the block average from $T^{(j)}$ to $T^{(j+1)}$;
\item the constants $a_j:=a(1-L^{-2})(1-L^{-2j})^{-1}$, $r_j:=a_jt_j^{-2}$ and $a_\flat:=a(1-L^{-2})$;
\item for $1\le j\le K-1$: $G_j:=(Y_{\mu_*}+r_jQ_j^*Q_j)^{-1}$, $\Delta^{(j)}:=r_j-r_j^2Q_jG_jQ_j^*$ and $C^{(j)}:=(\Delta^{(j)}+at_{j+1}^{-2}Q^*Q)^{-1}$ on $(\R^{T^{(j)}},\ipj j\cdot\cdot)$;
\item the slices $H_0:=G_1$ and $H_j:=r_j^2G_jQ_j^*C^{(j)}Q_jG_j$ for $1\le j\le K-1$;
\item for $M>0$ and $1\le k\le K$: $G^M_k:=(Y_M+r_kQ_k^*Q_k)^{-1}$.
\end{itemize}
By \cite[(2.17)]{K1} the $H_j$ are the slices of the photon propagator, and by \cite[p.~328]{K2} the same holds for the vector-field averaging used there. The operator $G^M_k$ is \cite[(2.13)]{K1} at zero background, with $m^2$ replaced by $M^2$.

\begin{lemma}[block Poincar\'e inequality]\label{app:G1:KS-poinc}
For $1\le j\le K$ and every $f$, $\norm{(1-P_j)f}^2\le\frac{t_j^2}8\ip f{-\Delta_\eps f}$, where $P_j$ replaces $f$ on each block $\square_j(y)$ by its average.
\end{lemma}
\begin{proof}
Let $n=t_j/\eps\ge2$. Drop from $\ip f{-\Delta_\eps f}$ all bonds that are not contained in a single block. What remains is the sum over blocks of the Neumann forms $q_\square(f)=\sum_{x\in\square}\eps^3\sum_\mu\abs{\nabla^+_\mu f(x)}^2$, each restricted to the bonds inside $\square$.

On the path $\{0,\dots,n-1\}$, the form $\sum_{i=0}^{n-2}(g_{i+1}-g_i)^2$ has the eigenvectors $\cos(\pi m(i+\frac12)/n)$ with eigenvalues $4\sin^2(\pi m/2n)$, $0\le m\le n-1$. On the cube the form is the Kronecker sum of three such forms. Its kernel consists of the constants, and its smallest non-zero eigenvalue is $4\sin^2(\pi/2n)$. Hence
\[
\sum_{x\in\square}\eps^3\abs{f-\operatorname{avg}_\square f}^2\le\frac{\eps^2}{4\sin^2(\pi/2n)}\,q_\square(f).
\]
The function $n\sin(\pi/2n)$ is increasing in $n$ and equals $\sqrt2$ at $n=2$. So $4\eps^{-2}\sin^2(\pi/2n)\ge8/t_j^2$. Summing over the blocks gives the claim.
\end{proof}

\begin{lemma}[resolvent facts]\label{app:G1:KS-res}
Let $Y>0$ on a finite-dimensional real Hilbert space $\mathcal H$, let $P$ be an orthogonal projection, $r>0$ and $G=(Y+rP)^{-1}$.
\begin{itemize}
\item[(i)] $0<G\le Y^{-1}$, and $Y^{-1}-G=Y^{-1}rPG=GrPY^{-1}$.
\item[(ii)] $0\le rPGP\le1$. Hence $\norm{PGv}\le\norm{PY^{-1}v}$ for every $v$.
\item[(iii)] Suppose $Y\ge B\ge0$ and $\norm{(1-P)f}^2\le\gamma^{-1}\ip f{Bf}$ for all $f$. Then $\norm G\le\max(\gamma^{-1},r^{-1})$.
\item[(iv)] Let $Q:\mathcal H\to\mathcal H'$ satisfy $QQ^*=1$ and $Q^*Q=P$. Then for every $B'\in\mathcal H'$,
\[
r\norm{B'}^2-r^2\ip{Q^*B'}{GQ^*B'}=\min_{A\in\mathcal H}\big[\ip A{YA}+r\norm{QA-B'}^2\big],
\]
and the minimum is attained at $A^*=rGQ^*B'$.
\end{itemize}
\end{lemma}
\begin{proof}
(i) We have $G^{-1}\ge Y>0$, and inversion is operator monotone. Moreover $Y^{-1}-G=Y^{-1}(G^{-1}-Y)G=G(G^{-1}-Y)Y^{-1}$.

(ii) Put $S=r^{1/2}PY^{-1/2}$. Then $Y+rP=Y^{1/2}(1+S^*S)Y^{1/2}$, so
\[
rPGP=S(1+S^*S)^{-1}S^*=SS^*(1+SS^*)^{-1}\in[0,1).
\]
By (i), $PGv=(1-rPGP)PY^{-1}v$, and $\norm{1-rPGP}\le1$.

(iii) $\ip f{(Y+rP)f}\ge\gamma\norm{(1-P)f}^2+r\norm{Pf}^2\ge\min(\gamma,r)\norm f^2$.

(iv) Since $\norm{QA}^2=\ip A{PA}$, the bracket equals
\[
\ip A{G^{-1}A}-2r\ip A{Q^*B'}+r\norm{B'}^2=\ip{A-A^*}{G^{-1}(A-A^*)}+r\norm{B'}^2-r^2\ip{Q^*B'}{GQ^*B'}.\qedhere
\]
\end{proof}

\begin{lemma}[lattice Green's function on the torus]\label{app:G1:KS-green}
Put $C_{\rm osc}^2:=(1+\frac{\sqrt3}2)^3\frac{\pi^2}8$. Let $\mu>0$ and $x,y\in T_\eps$, and let $z\in\eps\Z^3$ be any representative of $x-y$. Then
\[
C_\mu(x,x)\le\frac{13}{16\eps}+\frac1{\mu^2\ell^3},\qquad\norm{C_\mu(\delta_x-\delta_y)}^2\le C_{\rm osc}^2\abs z.
\]
\end{lemma}
\begin{proof}
Let $T^*=(2\pi/\ell)\Z^3\cap(-\pi/\eps,\pi/\eps]^3$ and $\lambda(k)=\sum_i4\eps^{-2}\sin^2(\eps k_i/2)$. Since $\abs{\sin u}\ge2\abs u/\pi$ for $\abs u\le\pi/2$, we have $\lambda(k)\ge\frac4{\pi^2}\abs k^2$ on $T^*$. Moreover $F(Y_\mu)(x,y)=\ell^{-3}\sum_{k\in T^*}F(\lambda(k)+\mu^2)\mathrm{e}^{ik\cdot(x-y)}$.

\emph{Diagonal.} Write $k=2\pi n/\ell$. The term $k=0$ is $\ell^{-3}\mu^{-2}$. For $n\ne0$ we have $\lambda(k)\ge16\abs n_\infty^2/\ell^2$, and $\#\{n:\abs n_\infty=r\}=24r^2+2\le26r^2$. Here $r\le N/2$ with $N=\ell/\eps$. So the other terms sum to at most $\ell^{-3}\frac{\ell^2}{16}\sum_{r=1}^{\lfloor N/2\rfloor}26\le\frac{13}{16\eps}$.

\emph{Oscillation.} We have
\[
\norm{C_\mu(\delta_x-\delta_y)}^2=\ell^{-3}\sum_{k\in T^*}\abs{\mathrm{e}^{ik\cdot z}-1}^2(\lambda(k)+\mu^2)^{-2}.
\]
The term $k=0$ vanishes. For $k\ne0$ the summand is at most $f(\abs k):=\frac{\pi^4}{16}\min(4\abs k^{-4},\abs z^2\abs k^{-2})$, a non-increasing function of $\abs k$. Extend the sum to all $k\in(2\pi/\ell)\Z^3\setminus0$, and put $\kappa=1+\frac{\sqrt3}2$. Let $Z_k$ be the cube of side $2\pi/\ell$ centred at $k$. For $k'\in Z_k$, $\abs{k'}\le\abs k+\sqrt3\pi/\ell\le\kappa\abs k$, so $f(\abs k)\le f(\abs{k'}/\kappa)$. The cubes $Z_k$ have disjoint interiors, hence
\[
\ell^{-3}\sum_{k\ne0}f(\abs k)\le\frac{\kappa^3}{8\pi^3}\int_{\R^3}f(\abs q)\,d^3q.
\]
Finally
\[
\int_{\R^3}f(\abs q)\,d^3q=\frac{\pi^5}4\Big[\int_0^{2/\abs z}\abs z^2dq+\int_{2/\abs z}^\infty4q^{-2}dq\Big]=\pi^5\abs z .
\]
So the bound is $\kappa^3\pi^2\abs z/8=C_{\rm osc}^2\abs z$.
\end{proof}

\begin{lemma}[King's slices]\label{app:K:L08-1}\label{app:G1:KS-slices}
In the setting above, let $P_j:=Q_j^*Q_j$. These operators, the operators $Q_j,Q$, the coefficients $a_j$, $r_j=a_jt_j^{-2}$, and the operators $G_j$, $C^{(j)}$, $F_j:=r_jG_jQ_j^*$, $H_j$ and $G^M_k$ are those of \cite[(2.13)--(2.17)]{K1}, with $0<a_\flat\le a$. Put $m_a:=\max(\frac18,\frac1{a_\flat})$, $c_C:=\frac4{a_\flat}+\frac{L^2}8+\frac{L^2}a$ and $C_H:=c_Ca^2(m_a+3^{1/4}C_{\rm osc})^2$, with $C_{\rm osc}$ as in Lemma~\ref{app:G1:KS-green}. At $L=2$, $a=1$ these are $m_a=\frac43$, $c_C=\frac{59}6$, $C_{\rm osc}^2<8.02$, $C_H<253.66$ and $c_Ca^2m_a^2=\frac{944}{54}<17.49$. Then:
\begin{itemize}
\item[(S1)] For $1\le j\le K$: $Q_jQ_j^*=1$, and $Q_j^*Q_j=P_j$ is the orthogonal projection onto the functions that are constant on the blocks of side $t_j$. Moreover $Q_{j+1}=QQ_j$, $QQ^*=1$ and $P_jP_{j+1}=P_{j+1}P_j=P_{j+1}$.
\item[(S2)] For $1\le j\le K-1$:
\begin{itemize}
\item $a_\flat\le a_j\le a$, so $0<r_j\le at_j^{-2}$;
\item $G_j=(Y_{\mu_*}+r_jP_j)^{-1}$ and $\norm{G_j}\le m_at_j^2$;
\item $C^{(j)}$ is well defined, and $0<C^{(j)}\le c_Ct_j^2$;
\item $H_j=F_jC^{(j)}F_j^*\ge0$.
\end{itemize}
\item[(S3)] For $1\le j\le K-1$, $s_0(H_j)\le C_H/t_j$. Moreover $0<H_0=G_1\le C_{\mu_*}$ and $s_0(H_0)\le\frac{13}{16\eps}+\frac1{\mu_*^2\ell_T^3}$.
\item[(S4)] For $1\le k\le K$ and $M>0$: $G^M_k=(Y_M+r_kP_k)^{-1}$ with $0<r_k\le at_k^{-2}$. Moreover $0<G^M_k\le C_M$ and $C_M-G^M_k=C_Mr_kP_kG^M_k$.
\end{itemize}
\end{lemma}
\begin{proof}
\emph{(S1).} By the choice of the weights, $\ipj j{Q_jf}B=\sum_x\eps^3f(x)B(y(x))$, where $y(x)$ labels the block containing $x$. So $Q_j^*B=B\circ y$ and $Q_jQ_j^*=1$. Hence $P_j$ is self-adjoint and idempotent, and it replaces $f$ by its block averages. Each block of scale $j+1$ is the union of $L^3$ blocks of scale $j$, so $Q_{j+1}=QQ_j$. The identity $QQ^*=1$ is proved as for $Q_j$. Since $\operatorname{ran}P_{j+1}\subset\operatorname{ran}P_j$, we get $P_jP_{j+1}=P_{j+1}$, and taking adjoints gives $P_{j+1}P_j=P_{j+1}$.

\emph{(S2).} The map $j\mapsto(1-L^{-2j})^{-1}$ decreases from $(1-L^{-2})^{-1}$ to $1$. Hence $a_\flat<a_j\le a_1=a$. By (S1), $G_j=(Y_{\mu_*}+r_jP_j)^{-1}$. Lemma~\ref{app:G1:KS-poinc} gives $\gamma=8t_j^{-2}$ with $B=-\Delta_\eps$. Then Lemma~\ref{app:G1:KS-res}(iii) gives $\norm{G_j}\le t_j^2\max(\frac18,\frac1{a_j})\le m_at_j^2$.

For the covariance, apply Lemma~\ref{app:G1:KS-res}(iv) with $Y=Y_{\mu_*}$, $Q=Q_j$ and $r=r_j$. Since $\ip{Q_j^*B}{G_jQ_j^*B}=\ipj jB{Q_jG_jQ_j^*B}$, this gives
\[
\ipj jB{\Delta^{(j)}B}=\ip{A^*}{Y_{\mu_*}A^*}+r_j\norm{B'-B}_j^2,\qquad B':=Q_jA^*,
\]
so in particular $\Delta^{(j)}\ge0$. We need three facts.
\begin{itemize}
\item $\norm{QD}_{j+1}\le\norm D_j$, because $Q^*Q$ is an orthogonal projection.
\item $\norm{Q_j^*D}=\norm D_j$ and $\norm{Q^*D}_j=\norm D_{j+1}$, because $Q_jQ_j^*=1$ and $QQ^*=1$.
\item By (S1), $\norm{(1-Q^*Q)B'}_j=\norm{(P_j-P_{j+1})A^*}=\norm{P_j(1-P_{j+1})A^*}\le\norm{(1-P_{j+1})A^*}$. By Lemma~\ref{app:G1:KS-poinc} at scale $j+1$, this is at most $\tau^{1/2}\ip{A^*}{Y_{\mu_*}A^*}^{1/2}$ with $\tau=t_{j+1}^2/8$.
\end{itemize}
Since $\norm B_j\le\norm{B-B'}_j+\norm{(1-Q^*Q)B'}_j+\norm{Q^*QB'}_j$, $\norm{Q^*QB'}_j=\norm{QB'}_{j+1}$ and $\norm{QB'}_{j+1}\le\norm{QB}_{j+1}+\norm{B'-B}_j$, we obtain
\[
\norm B_j\le2\norm{B-B'}_j+\tau^{1/2}\ip{A^*}{Y_{\mu_*}A^*}^{1/2}+\norm{QB}_{j+1}.
\]
Put $\beta=at_{j+1}^{-2}$. The Cauchy--Schwarz inequality with weights $(r_j,\tau^{-1},\beta)$ gives
\[
\norm B_j^2\le\Big(\frac4{r_j}+\tau+\frac1\beta\Big)\big[\ipj jB{\Delta^{(j)}B}+\beta\norm{QB}^2_{j+1}\big]=t_j^2\Big(\frac4{a_j}+\frac{L^2}8+\frac{L^2}a\Big)\ipj jB{(\Delta^{(j)}+\beta Q^*Q)B}.
\]
So $\Delta^{(j)}+\beta Q^*Q$ is positive definite, $C^{(j)}>0$ exists, and $C^{(j)}\le c_Ct_j^2$ because $a_j\ge a_\flat$. Finally $F_j^*=r_jQ_jG_j$, so $H_j=F_jC^{(j)}F_j^*\ge0$.

\emph{(S3).} Let $1\le j\le K-1$. By (S2) and $\norm{Q_jf}_j=\norm{P_jf}$,
\[
H_j(x,x)=\ipj j{F_j^*\delta_x}{C^{(j)}F_j^*\delta_x}\le c_Ct_j^2r_j^2\norm{P_jG_j\delta_x}^2 .
\]
Write $\delta_x=\rho_x+v$ with $\rho_x:=P_j\delta_x=t_j^{-3}\mathbf 1_{\square_x}$, where $\square_x$ is the block containing $x$. Then $\norm{\rho_x}=t_j^{-3/2}$.
\begin{itemize}
\item First, $\norm{P_jG_j\rho_x}\le m_at_j^{1/2}$.
\item Second, Lemma~\ref{app:G1:KS-res}(ii) gives $\norm{P_jG_jv}\le\norm{C_{\mu_*}v}=\sup_{\norm w\le1}\ip{C_{\mu_*}w}v$. With $\psi_w:=C_{\mu_*}w$ we have $\ip{\psi_w}v=\psi_w(x)-\operatorname{avg}_{\square_x}\psi_w$. Hence $\abs{\ip{\psi_w}v}\le\max_{y\in\square_x}\abs{\ip w{C_{\mu_*}(\delta_x-\delta_y)}}$. For $y\in\square_x$, $x-y$ has a representative $z$ with $\abs z\le\sqrt3t_j$. Lemma~\ref{app:G1:KS-green} then gives $\norm{P_jG_jv}\le3^{1/4}C_{\rm osc}t_j^{1/2}$.
\end{itemize}
Therefore $\norm{P_jG_j\delta_x}\le(m_a+3^{1/4}C_{\rm osc})t_j^{1/2}$. With $r_j\le at_j^{-2}$ this gives $H_j(x,x)\le C_H/t_j$.

For $j=0$: $H_0=G_1$, and $0<G_1\le C_{\mu_*}$ by Lemma~\ref{app:G1:KS-res}(i). Lemma~\ref{app:G1:KS-green} bounds $C_{\mu_*}(x,x)$.

\emph{(S4).} By (S1), $G^M_k=(Y_M+r_kP_k)^{-1}$, and $0<r_k\le at_k^{-2}$ as in (S2). The remaining claims are Lemma~\ref{app:G1:KS-res}(i) with $Y=Y_M$.

\emph{Constants.} At $L=2$, $a=1$ we have $a_\flat=\frac34$, so $m_a=\frac43$ and $c_C=\frac{16}3+\frac12+4=\frac{59}6$. Moreover $C_{\rm osc}^2=8.0160\ldots$. With $C_{\rm osc}^2\le8.1$ we get $C_H\le\frac{59}6(\frac43+3^{1/4}\sqrt{8.1})^2=253.657\ldots$. Finally $c_Cm_a^2=\frac{59}6\cdot\frac{16}9=\frac{944}{54}$.
\end{proof}

\paragraph{Kernels on $\eps\Z^3$ and on the torus.} For the rest of this subsection let $M>0$ and assume $\ell\ge16\eps$. We use the following notation.
\begin{itemize}
\item $\mathrm{BZ}=(-\pi/\eps,\pi/\eps]^3$, $\bar dp=d^3p/(2\pi)^3$, $\hat p^2=\sum_\mu4\eps^{-2}\sin^2(\eps p_\mu/2)$ and $\hat G=(\hat p^2+M^2)^{-1}$.
\item A kernel on $\eps\Z^3$ is written $\mathcal K(x)=\int_{\rm BZ}\hat{\mathcal K}(p)\mathrm{e}^{ip\cdot x}\bar dp$. We write $G=G_M$ for the kernel of $\hat G$ and $\mathcal G=\mathcal G_M$ for the kernel of $\hat G^2$.
\item $T^*=(2\pi/\ell)\Z^3\cap\mathrm{BZ}$ and $\abs T=\ell^3$.
\item $\Lambda_0=\{z\in\eps\Z^3:-\ell/2<z_i\le\ell/2\}$ is the fundamental cell, and $\Lambda_0^+=\{y:\abs y_\infty\le\ell/2+2\eps\}$.
\item $c$ is the kernel of $C=C_M$ on $T_\eps$, so $C(x,y)=c(x-y)$.
\item $c_N=\pi^3/8$, and $\eta(s):=\frac2\eps\operatorname{arcsinh}\frac{\eps\sqrt s}{2\sqrt2}$.
\end{itemize}

\begin{lemma}[two elementary tools]\label{app:G1:B5e2-tools}
\begin{itemize}
\item[(i)] If $F\ge0$ is non-increasing, then $\int_{\rm BZ}F(\hat p^2)\,d^3p\le c_N\int_{\R^3}F(\abs p^2)\,d^3p$.
\item[(ii)] If $\Phi$ is smooth and $2\pi/\eps$-periodic, then $\abs T^{-1}\sum_{p\in T^*}\Phi(p)=\sum_{n\in\Z^3}\check\Phi(n\ell)$, where $\check\Phi(x)=\int_{\rm BZ}\Phi(p)\mathrm{e}^{ip\cdot x}\bar dp$.
\end{itemize}
\end{lemma}
\begin{proof}
(i) On $\mathrm{BZ}$, $\abs{\hat p_\mu}\ge\frac2\pi\abs{p_\mu}$, so $F(\hat p^2)\le F(4\abs p^2/\pi^2)$. Extend the integral to $\R^3$ and substitute $q=2p/\pi$.

(ii) Insert the absolutely convergent series $\Phi(p)=\sum_{x\in\eps\Z^3}\eps^3\check\Phi(x)\mathrm{e}^{-ip\cdot x}$, and use $\abs T^{-1}\sum_{p\in T^*}\mathrm{e}^{-ip\cdot x}=\eps^{-3}\mathbf 1[x\in\ell\Z^3]$.
\end{proof}

\begin{lemma}[contour bound]\label{app:G1:B5e2-K}
Let $s>0$, $a\in\{0,1,2\}$, and let $\nabla^\alpha$ be a product of $a$ forward or backward differences. Put $b=s/2$ and $I(\beta',b)=\int_{\R^3}(k^2+b)^{-\beta'}\bar dk$. Let $z\ne0$, $r=\abs z_\infty$ and $\eta=\eta(s)$. Then the kernel $K_s$ of $(\hat p^2+s)^{-3}$ satisfies
\[
\abs{\nabla^\alpha K_s(z)}\le c_N\mathrm{e}^{-\eta(r-a\eps)}\mathcal J_a,
\]
where $\mathcal J_0=I(3,b)$ and $\mathcal J_a=2^{a-1}[I(3-\frac a2,b)+\eta^aI(3,b)]$ for $a\ge1$.
\end{lemma}
\begin{proof}
Let $\abs{z_\nu}=r$ and $\sigma=\operatorname{sign}z_\nu$. Let $m_\alpha(p)$ be the product of the factors $\pm(\mathrm{e}^{\pm i\eps p_\lambda}-1)/\eps$. The integrand $\mathrm{e}^{ip\cdot z}m_\alpha(p)(\hat p^2+s)^{-3}$ is analytic and $2\pi/\eps$-periodic in $p_\nu$. Shift $p_\nu\to p_\nu+i\sigma\eta$; the vertical sides of the contour cancel by periodicity. On the shifted contour three estimates hold.
\begin{itemize}
\item $\abs{\mathrm{e}^{ip\cdot z}}=\mathrm{e}^{-\eta r}$.
\item We have
\[
\rRe\,4\eps^{-2}\sin^2(\eps(p_\nu+i\sigma\eta)/2)=4\eps^{-2}\sin^2(\eps p_\nu/2)\cosh\eps\eta-4\eps^{-2}\sinh^2(\eps\eta/2),
\]
and $4\eps^{-2}\sinh^2(\eps\eta/2)=s/2$. Hence $\abs{\hat p^2+s}\ge\hat p^2+b$, with $\hat p^2$ taken at real $p$. In particular there is no pole in the strip.
\item A factor in the direction $\nu$ is at most $\mathrm{e}^{\eps\eta}\abs{\hat p_\nu}+(\mathrm{e}^{\eps\eta}-1)/\eps\le\mathrm{e}^{\eps\eta}(\abs{\hat p}+\eta)$. So $\abs{m_\alpha}\le\mathrm{e}^{a\eps\eta}2^{(a-1)_+}(\abs{\hat p}^a+\eta^a)$.
\end{itemize}
Moreover $\abs{\hat p}^a(\hat p^2+b)^{-3}\le(\hat p^2+b)^{a/2-3}$, and both $(\hat p^2+b)^{a/2-3}$ and $(\hat p^2+b)^{-3}$ are non-increasing functions of $\hat p^2$. Lemma~\ref{app:G1:B5e2-tools}(i) now gives the bound. All integrals are finite because $3-a/2>3/2$.
\end{proof}

\begin{corollary}[majorants]\label{app:G1:B5e2-K1}
Put $v_0=M/\sqrt2$, $\beta=\frac12\operatorname{arcsinh}2$, $c_0=\frac1{32\pi}$, $c_1=\frac1{6\pi^2}+\frac1{32\pi}$, $c_2=\frac5{16\pi}$, $\bar w=\sqrt{12}c_N/(2\pi^2)$, $\tau_0=\frac1{96}$, $\tau_1=\frac18$ and $\tau_2=2$. Let $z\ne0$, $r=\abs z_\infty$, and let $\nabla^\alpha$ be as in Lemma~\ref{app:G1:B5e2-K}.
\begin{itemize}
\item[(i)] $\abs{\nabla^\alpha G(z)}\le\kappa_a(r):=16c_Nc_a\int_{v_0}^\infty v^a\mathrm{e}^{-\eta(2v^2)(r-a\eps)}dv$, which is non-increasing in $r$. Likewise $\abs{\nabla^\alpha\mathcal G(z)}\le\kappa^{\mathcal G}_a(r):=8c_Nc_a\int_{v_0}^\infty v^{a-2}\mathrm{e}^{-\eta(2v^2)(r-a\eps)}dv$.
\item[(ii)] If $r\ge(a+2)\eps$, put $\rho:=r-a\eps$. Then
\[
\kappa_a(r)\le16c_Nc_a\big(a!(\beta\rho)^{-a-1}+\tau_a\rho^{-a-1}\big)\le16c_Nc_a(a!\beta^{-a-1}+\tau_a)\big(\tfrac{a+2}2\big)^{a+1}r^{-a-1}.
\]
More precisely, the part $\eps v\le4$ of the integral is at most $\int_{v_0}^\infty v^a\mathrm{e}^{-\beta v\rho}dv$, and the part $\eps v>4$ is at most $\rho^{-a-1}n^{a+1}4^{a+1-2n}/(2n-a-1)$, with $n=\rho/\eps$.
\item[(iii)] If $0<r<(a+2)\eps$, then $\abs{\nabla^\alpha G(z)}\le2^a\bar w\eps^{-a-1}\le2^a\bar w(a+2)^{a+1}r^{-a-1}$.
\end{itemize}
Hence $\abs{\nabla^\alpha G(z)}\le\hat G_ar^{-a-1}$ with $\hat G_0=1.361$, $\hat G_1=12.25$ and $\hat G_2=361.2$. Moreover:
\begin{itemize}
\item $0\le G(z)\le G(0)\le\bar w/\eps$;
\item $G(z)\le\hat A_0/\max(r,\eps)$ for all $z$, with $\hat A_0=0.86101$;
\item if $\rho:=r-\eps\ge2\eps$, then $\abs{\nabla^\pm_\mu G(z)}\le\hat G_1^*\rho^{-2}$ with $\hat G_1^*=3.402$.
\end{itemize}
All these constants are independent of $M$ and $\eps$.
\end{corollary}
\begin{proof}
\emph{Integral representation.} We have $\hat G=\int_{M^2}^\infty2(s-M^2)(\hat p^2+s)^{-3}ds$ and $\hat G^2=\int_{M^2}^\infty2(\hat p^2+s)^{-3}ds$. Apply Lemma~\ref{app:G1:B5e2-K} with $v=\sqrt{s/2}$. Since $\operatorname{arcsinh}x\le x$, we have $\eta(s)\le v$, so $\mathcal J_a\le c_av^{a-3}$. Here
\[
c_0=I(3,1),\qquad c_1=I(\tfrac52,1)+I(3,1),\qquad c_2=2I(2,1)+2I(3,1),
\]
from $I(\beta',b)=\Gamma(\beta'-\frac32)b^{3/2-\beta'}/((4\pi)^{3/2}\Gamma(\beta'))$. With $ds=4v\,dv$ this gives (i).

\emph{Far branch (ii).} For $\eps v\le4$, concavity gives $\operatorname{arcsinh}x\ge\beta x$ on $[0,2]$, so $\eta\ge\beta v$. That part of the integral is therefore at most $\int_{v_0}^\infty v^a\mathrm{e}^{-\beta v\rho}dv\le a!(\beta\rho)^{-a-1}$. For $\eps v>4$, $\operatorname{arcsinh}x\ge\ln2x$, so $\mathrm{e}^{-\eta\rho}\le(\eps v)^{-2n}$ with $n=\rho/\eps\ge2$. That part is at most
\[
\eps^{-a-1}4^{a+1-2n}/(2n-a-1)=\rho^{-a-1}n^{a+1}4^{a+1-2n}/(2n-a-1).
\]
This expression is decreasing in $n$, and its value at $n=2$ is $\tau_a\rho^{-a-1}$. Finally $\rho\ge2r/(a+2)$.

\emph{Near branch (iii).} We have $\abs{\nabla^\alpha f}\le2^a\eps^{-a}\max\abs f$, and $\max\abs G=G(0)$ because $\hat G\ge0$. Lemma~\ref{app:G1:B5e2-tools}(i) with $F(u)=u^{-1}\mathbf 1_{u\le12/\eps^2}$ gives
\[
\eps G(0)\le\eps\int_{\rm BZ}\hat p^{-2}\bar dp\le\eps c_N\int_{\abs p^2\le12/\eps^2}\abs p^{-2}\bar dp=\bar w.
\]
Finally $\eps>r/(a+2)$. For the positivity of $G$, let $(\mathcal Af)(x):=\eps^{-2}\sum_{\mu,\pm}f(x\pm\eps e_\mu)$ on $\ell^2(\eps\Z^3)$. Then $-\Delta_\eps+M^2=(6\eps^{-2}+M^2)-\mathcal A$ and $\norm{\mathcal A}\le6\eps^{-2}$, so $(-\Delta_\eps+M^2)^{-1}=\sum_{n\ge0}(6\eps^{-2}+M^2)^{-n-1}\mathcal A^n$ is a norm-convergent series of operators with non-negative kernels. Hence $G\ge0$.

\emph{Numerical values.} For $a=0,1,2$, the far and near branches give the constants $(0.861005,\,1.36035)$, $(7.6541,\,12.2431)$ and $(361.128,\,174.125)$. For $\hat A_0$: if $r\ge2\eps$, use (ii) at $a=0$ with $\rho=r$, which gives the constant $16c_Nc_0(\beta^{-1}+\tau_0)=0.8610045\ldots$; if $r<2\eps$, use $G(z)\le\bar w/\eps$ with $\bar w=0.6802\ldots$. For $\hat G_1^*$, use (ii) at $a=1$, which gives $16c_Nc_1(\beta^{-2}+\tau_1)=3.40183\ldots$.
\end{proof}

\begin{lemma}[images]\label{app:G1:B5e2-I}
For a kernel $\mathcal K$ with the non-increasing majorant $\kappa_a$ of Corollary~\ref{app:G1:B5e2-K1}, put
\[
\rho_a(\mathcal K):=\sum_{j\ge1}(24j^2+2)\,\kappa_a\big((j-\tfrac12)\ell-4\eps\big),
\]
with $\kappa^{\mathcal G}_a$ in place of $\kappa_a$ for $\mathcal K=\mathcal G$.
\begin{itemize}
\item[(i)] For all $z$, $c(z)=\sum_{n\in\Z^3}G(z+n\ell)$. The difference $R:=c-G$ on $\Lambda_0^+$ satisfies $\abs{\nabla^\alpha R(y)}\le\rho_a(G)$ for every $a$-fold stencil ($a\le2$) based at $y\in\Lambda_0^+$.
\item[(ii)] The same holds for $\mathcal G$ and the kernel of $C^2$ on $T_\eps$.
\item[(iii)] Each $\rho_a$ is finite and non-increasing in $\ell$. It is bounded uniformly in $\eps\le\eps_0$, and it tends to $0$ exponentially fast as $\ell\to\infty$, uniformly in $\eps\le\eps_0$.
\end{itemize}
\end{lemma}
\begin{proof}
For (i), apply Lemma~\ref{app:G1:B5e2-tools}(ii) to $\Phi(p)=\hat G(p)\mathrm{e}^{ip\cdot z}$. Suppose $\abs n_\infty=j\ge1$. Then every point of the stencil based at $y+n\ell$ satisfies $\abs\cdot_\infty\ge(j-\frac12)\ell-4\eps\ge4\eps\ge(a+2)\eps$. Moreover $\#\{n:\abs n_\infty=j\}=24j^2+2$. Now use Corollary~\ref{app:G1:B5e2-K1}(i) and the monotonicity of $\kappa_a$. Statement (ii) is proved in the same way.

(iii) Let $\rho$ be the argument of $\kappa_a$ minus $a\eps$. Since $\ell\ge16\eps$, $\rho\ge(j-\frac12)\ell-6\eps\ge j\ell/8$. By Corollary~\ref{app:G1:B5e2-K1}(ii), the part $\eps v\le4$ of the integral is at most $\int_{v_0}^\infty v^a\mathrm{e}^{-\beta v\rho}dv$, which decays like $\mathrm{e}^{-\beta v_0\rho}$. The part $\eps v>4$ is at most
\[
4^{a+1}\rho^{-a-1}n^{a+1}16^{-n}/(2n-a-1)\le64\,\rho^{-a-1}n^{a+1}16^{-n},\qquad n=\rho/\eps\ge\max(2,\rho/\eps_0).
\]
Since $n^{a+1}16^{-n}$ is decreasing for $n\ge2$, this is at most its value at $n=\max(2,\rho/\eps_0)$. Both bounds are independent of $\eps$ and decrease in $\rho$, and their sums against $24j^2+2$ converge and tend to $0$ exponentially fast as $\ell\to\infty$. For $\mathcal G$, use in addition $v^{a-2}\le v_0^{-2}v^a$. Monotonicity in $\ell$ holds because $\kappa_a$ is non-increasing and the arguments increase with $\ell$.
\end{proof}

\begin{lemma}[inputs of the order-$e^2$ anchor bounds]\label{app:K:L08-2}\label{app:G1:B5e2-inputs}
Let $L=2$, $0<\eps\le\eps_0$ and $M>0$, and let $\ell=\ell_T\ge\ell_0$ with $\ell_T\ge16\eps$. Let $G_M$, $\mathcal G_M$, $c$, $c_N$, $v_0$, $\eta$, $c_a$, $\beta$, $\bar w$, $\kappa_a$, $\kappa^{\mathcal G}_a$ and $\rho_a$ be as above (Corollary~\ref{app:G1:B5e2-K1}, Lemma~\ref{app:G1:B5e2-I}). For a positive semidefinite operator $\mathcal K$ on $T_\eps$ put $s_0(\mathcal K):=\sup_x\mathcal K(x,x)$ and $s_1(\mathcal K):=\sup_{x,\mu,\pm}\ip{\nabla^{\pm*}_\mu\delta_x}{\mathcal K\nabla^{\pm*}_\mu\delta_x}$. Put $\ell':=\ell_T/2+\eps$, and
\[
I(\ell_T):=\hat G_1\rho_1\ell'^2+\hat G_0\rho_2\ell'^3/3+\hat G_2\rho_0\ell'+(\rho_1^2+\rho_0\rho_2)\ell'^4/4,\qquad\rho_a=\rho_a(G_M).
\]
Then:
\begin{itemize}
\item[(a)] The quantity
\[
\widehat\Pi_\mu(0):=2e^2c(\eps e_\mu)-2e^2\eps^{-2}\sum_z\eps^3[c(z)^2-c(z+\eps e_\mu)c(z-\eps e_\mu)]
\]
satisfies $\abs{\widehat\Pi_\mu(0)}\le e^2w_T$, where $w_T:=2\rho_0(G_M)+4\rho_2(\mathcal G_M)$.
\item[(b)] For $z\ne0$ and a product $\nabla^\alpha$ of $a\le2$ forward or backward differences, $\abs{\nabla^\alpha G_M(z)}\le\hat G_a\abs z_\infty^{-a-1}$, with $\hat G_0=1.361$, $\hat G_1=12.25$ and $\hat G_2=361.2$. Moreover $\hat B:=\hat G_1^2+\hat G_0\hat G_2\le641.7$. On $\Lambda_0^+$, and in particular on the fundamental cell, $c=G_M+R$ with $\abs{\nabla^\alpha R}\le\rho_a(G_M)$.
\item[(c)] Let $a=1$, let the operators be those of Lemma~\ref{app:K:L08-1}, put $t_0=\eps$, and assume $t_K\le\eps_0$. Put
\begin{itemize}
\item $W_1:=4374\hat A_0^2+2S_4\hat G_1^{*2}$, with $\hat A_0=0.86101$, $\hat G_1^*=3.402$ and $S_4:=\sum_{m\ge2}(24m^2+2)(m-1)^{-4}\le125.32$, so that $W_1\le6144$;
\item $W_{\rm img}(m):=216\rho_0(G_m)^2\eps_0+2\ell_T^3\rho_1(G_m)^2$;
\item $\sigma_0:=C_H+\eps_0/(\mu_*^2\ell_0^3)$, $h:=c_Ca^2m_a^2\le17.49$ and $\theta_H:=c_Ca^2[W_1+\eps_0W_{\rm img}(\mu_*)]$;
\item $\gamma:=a[c_N/(8\pi M)+\rho_0(\mathcal G_M)]$ and $\theta:=a[W_1+\eps_0W_{\rm img}(M)]$,
\end{itemize}
with $c_C=\frac{59}6$. Then:
\begin{itemize}
\item for $0\le j\le K-1$: $s_0(H_j)\le\sigma_0/t_j$, $\norm{H_j}\le ht_j^2$ and $s_1(H_j)\le\theta_H/t_j^3$;
\item for $1\le k\le K$, the operator $\mathcal C_k:=C-G^M_k$ satisfies $0\le\mathcal C_k\le C$, $s_0(\mathcal C_k)\le\gamma t_k^{-2}$ and $s_1(\mathcal C_k)\le\theta t_k^{-3}$. Moreover $\sum_y\eps^3\sum_\mu\abs{\nabla^+_{\mu,y}\mathcal C_k(x,y)}^2\le\mathcal C_k(x,x)$ and $\norm{\mathcal C_k}\le M^{-2}$;
\item for every positive semidefinite $\mathcal K$: $\abs{\mathcal K(x,y)}\le s_0(\mathcal K)$, $\abs{\nabla\mathcal K}\le\sqrt{s_0(\mathcal K)s_1(\mathcal K)}$, $\abs{\nabla\nabla'\mathcal K}\le s_1(\mathcal K)$ and $\sum_y\eps^3\mathcal K(x,y)^2\le\norm{\mathcal K}\mathcal K(x,x)$. Here $\nabla$ and $\nabla'$ are single differences in the first and in the second variable.
\end{itemize}
\item[(d)] We have
\[
\sum_{z\in T_\eps}\eps^3\abs{\nabla^\pm_\mu c(z)}\le\Lambda_1:=26\cdot\tfrac94\cdot16c_Nc_1\Big[\frac{2\sqrt2}{\beta^3M}+0.1345\eps_0\Big]+250\bar w\eps_0+\ell_T^3\rho_1(G_M).
\]
\item[(e)] We have:
\begin{itemize}
\item $\#\{z\in\eps\Z^3:\abs z_\infty=j\eps\}=24j^2+2$;
\item if $f\ge0$ is non-increasing, then $\sum_{j\ge1}\eps f(j\eps)\le\int_0^\infty f$;
\item for $q\ge-2$, $\sum_{0<\abs z_\infty\le\ell_T/2}\eps^3\abs z_\infty^q\le26\int_0^{\ell'}r^{2+q}dr$.
\end{itemize}
\end{itemize}
Each $\rho_a$ is finite and non-increasing in $\ell_T$, it is bounded uniformly in $\eps\le\eps_0$, and it tends to $0$ exponentially fast as $\ell_T\to\infty$, uniformly in $\eps\le\eps_0$. All constants above are independent of $\eps\le\eps_0$, $j$ and $k$, and are bounded for $\ell_T\ge\ell_0$.
\end{lemma}
\begin{proof}
\emph{(e)} We have $(2j+1)^3-(2j-1)^3=24j^2+2\le26j^2$. For non-increasing $f$, $\eps f(j\eps)\le\int_{(j-1)\eps}^{j\eps}f$. For $q\ge-2$ the function $r^{2+q}$ is non-decreasing, so
\[
\sum_{0<\abs z_\infty\le\ell/2}\eps^3\abs z_\infty^q\le26\sum_{j\le\ell/(2\eps)}\eps(j\eps)^{2+q}\le26\int_0^{\ell/2+\eps}r^{2+q}dr.
\]

\emph{(b)} This is Corollary~\ref{app:G1:B5e2-K1} together with Lemma~\ref{app:G1:B5e2-I}(i). Numerically, $\hat G_1^2+\hat G_0\hat G_2=641.66\ldots$.

\emph{(a)} By Parseval on $T_\eps$:
\begin{itemize}
\item $c(\eps e_\mu)=\abs T^{-1}\sum_{T^*}\hat G\cos\eps p_\mu$;
\item $\sum_z\eps^3c(z)^2=\abs T^{-1}\sum_{T^*}\hat G^2$;
\item $\sum_z\eps^3c(z+\eps e_\mu)c(z-\eps e_\mu)=\abs T^{-1}\sum_{T^*}\hat G^2\cos2\eps p_\mu$.
\end{itemize}
Hence $\widehat\Pi_\mu(0)=e^2\abs T^{-1}\sum_{T^*}\Phi$, where
\[
\Phi=2\cos(\eps p_\mu)\hat G-4\eps^{-2}\sin^2(\eps p_\mu)\hat G^2=\partial^2_{p_\mu}\ln(\hat p^2+M^2),
\]
because $\partial_{p_\mu}\hat p^2=\frac2\eps\sin\eps p_\mu$. By Lemma~\ref{app:G1:B5e2-tools}(ii), $\widehat\Pi_\mu(0)=e^2\sum_n\check\Phi(n\ell)$. Here $\check\Phi(0)=\int_{\rm BZ}\partial^2_{p_\mu}\ln(\hat p^2+M^2)\bar dp=0$ by periodicity. Write $\Delta_\mu:=\nabla^+_\mu\nabla^-_\mu$. Using
\[
f(x+2\eps e_\mu)+f(x-2\eps e_\mu)-2f(x)=\eps^2[\Delta_\mu f(x+\eps e_\mu)+2\Delta_\mu f(x)+\Delta_\mu f(x-\eps e_\mu)],
\]
we get
\[
\check\Phi(x)=G(x+\eps e_\mu)+G(x-\eps e_\mu)+\Delta_\mu\mathcal G(x+\eps e_\mu)+2\Delta_\mu\mathcal G(x)+\Delta_\mu\mathcal G(x-\eps e_\mu).
\]
The base points $\pm\eps e_\mu$ and $0$ lie in $\Lambda_0^+$. Summing over $n\ne0$ with Lemma~\ref{app:G1:B5e2-I}(i),(ii) gives $\abs{\widehat\Pi_\mu(0)}\le e^2[2\rho_0(G)+4\rho_2(\mathcal G)]=e^2w_T$.

\emph{(d)} Sum over $z\in\Lambda_0$. There $\abs{\nabla^\pm c}\le\abs{\nabla^\pm G}+\rho_1$, and the $\rho_1$ part contributes at most $\ell^3\rho_1$. The $125$ sites with $\abs z_\infty\le2\eps$ contribute at most $125\eps^3\cdot2\bar w\eps^{-2}\le250\bar w\eps_0$, by Corollary~\ref{app:G1:B5e2-K1}.

Now let $\abs z_\infty=j\eps\ge3\eps$, and put $\rho=(j-1)\eps$ and $n=j-1\ge2$. By Corollary~\ref{app:G1:B5e2-K1}(ii) at $a=1$:
\begin{itemize}
\item the part $\eps v\le4$ is at most $16c_Nc_1\int_{v_0}^\infty v\mathrm{e}^{-\beta v\rho}dv=16c_Nc_1\rho^{-2}\beta^{-2}\mathrm{e}^{-x}(1+x)$, with $x=\beta v_0\rho$;
\item the part $\eps v>4$ is at most $16c_Nc_1\rho^{-2}n^24^{2-2n}/(2n-2)$.
\end{itemize}
Since $j/(j-1)\le\frac32$,
\[
(j\eps)^2\abs{\nabla^\pm G(z)}\le\tfrac94\cdot16c_Nc_1\big[\beta^{-2}\mathrm{e}^{-x}(1+x)+n^24^{2-2n}/(2n-2)\big].
\]
Multiply by $\eps^3(j\eps)^{-2}$ and sum over $z$, using $24j^2+2\le26j^2$. This leaves $26\cdot\frac94\cdot16c_Nc_1\sum_{n\ge2}\eps[\cdots]$. The function $\rho\mapsto\beta^{-2}\mathrm{e}^{-\beta v_0\rho}(1+\beta v_0\rho)$ is decreasing. By (e) its sum is at most its integral over $(0,\infty)$, which is $2/(\beta^3v_0)=2\sqrt2/(\beta^3M)$. Moreover $\sum_{n\ge2}n^24^{2-2n}/(2n-2)=0.13449\ldots\le0.1345$ and $\eps\le\eps_0$. Together these give $\Lambda_1$.

\emph{(c)} The operators satisfy (S1)--(S4) of Lemma~\ref{app:K:L08-1}.

\emph{Step c1 (block averages).} Let $t\in\{t_j:1\le j\le K\}$, so that $2\eps\le t\le\eps_0$. Let $m>0$, write $P_t$ for the corresponding $P_j$, and let $w=C_m\nabla^{+*}_\mu\delta_x$. We show
\[
\norm{P_tw}^2\le t^{-1}\big[W_1+\eps_0W_{\rm img}(m)\big].
\]
We have $w=-\nabla^-_\mu f$ with $f=c_m(\cdot-x)$, where $c_m$ is the torus kernel of $C_m$. Moreover $f\ge0$, by Lemma~\ref{app:G1:B5e2-I}(i) and $G_m\ge0$ (Corollary~\ref{app:G1:B5e2-K1}). Summing along $\mu$-lines in a block $\square$,
\[
\sum_{y\in\square}\eps^3\nabla^-_\mu f(y)=\eps^2\sum_{y_\perp}\big[f(\text{top face})-f(\text{face below }\square)\big],
\]
so $\abs{\operatorname{avg}_\square w}\le t^{-3}\eps^2\sum_{\text{two faces}}f$.

\emph{Near blocks.} On a face, write $f\le\hat A_0/\max(\abs{z_\perp}_\infty,\eps)+\rho_0$, by Corollary~\ref{app:G1:B5e2-K1} and Lemma~\ref{app:G1:B5e2-I}(i) for the mass $m$. Here $z$ is the representative of $y-x$ in $\Lambda_0$, and the in-plane offsets $z_\perp$ are distinct points of $\eps\Z^2$. A radially decreasing function summed over $(t/\eps)^2$ distinct points is at most its sum over $\{\abs v_\infty\le t/2\}$. That sum is $\eps^{-2}\hat A_0(\eps+4t)\le4.5\eps^{-2}\hat A_0t$. Hence, on the $27$ blocks at block offset $\le1$ from the block of $x$,
\[
\abs{\operatorname{avg}_\square w}\le9\hat A_0t^{-2}+2\rho_0t^{-1}.
\]

\emph{Far blocks.} At block offset $m'\ge2$ there are at most $24m'^2+2$ blocks. Every point $y$ of such a block has $\abs{y-x}_\infty-\eps\ge(m'-1)t\ge2\eps$, so $\abs w\le\hat G_1^*((m'-1)t)^{-2}+\rho_1$.

\emph{Sum.} Therefore
\[
\norm{P_tw}^2=\sum_\square t^3\abs{\operatorname{avg}_\square w}^2\le t^{-1}\big(4374\hat A_0^2+2S_4\hat G_1^{*2}\big)+216\rho_0^2t+2\ell^3\rho_1^2\le t^{-1}\big[W_1+\eps_0W_{\rm img}(m)\big],
\]
because $t\le\eps_0$. Backward differences are treated in the same way.

\emph{Step c2 (the slices).} Let $1\le j\le K-1$ and $v$ arbitrary. By (S2), $\norm{Q_jg}_j=\norm{P_jg}$ and Lemma~\ref{app:G1:KS-res}(ii),
\[
\ip v{H_jv}\le c_Ct_j^2r_j^2\norm{P_jG_jv}^2\le c_Ca^2t_j^{-2}\norm{P_jC_{\mu_*}v}^2 .
\]
With $v=\nabla^{\pm*}_\mu\delta_x$ and Step c1 at $m=\mu_*$, this gives $s_1(H_j)\le\theta_H/t_j^3$. Also $\norm{H_j}\le c_Ct_j^2\,a^2t_j^{-4}\,m_a^2t_j^4=ht_j^2$, and $s_0(H_j)\le C_H/t_j\le\sigma_0/t_j$ by (S3).

For $j=0$, we have $H_0=G_1\le C_{\mu_*}$. Hence
\[
s_1(H_0)\le\abs T^{-1}\sum_{q\in T^*}\hat q_\mu^2/(\hat q^2+\mu_*^2)\le\eps^{-3}\le\theta_H/t_0^3,
\]
because $\theta_H\ge1$. Moreover $\norm{H_0}\le m_aL^2\eps^2=\frac{16}3\eps^2\le h\eps^2$. Finally, by (S3), $s_0(H_0)\le t_0^{-1}[\frac{13}{16}+\eps/(\mu_*^2\ell^3)]\le\sigma_0/t_0$, because $\ell\ge\ell_0$, $\eps\le\eps_0$ and $\frac{13}{16}\le C_H$.

\emph{Step c3 (coarse scalar part).} By (S4), $0<G^M_k\le C$, so $0\le\mathcal C_k\le C$, and $\mathcal C_k=Cr_kP_kG^M_k$. With Lemma~\ref{app:G1:KS-res}(ii),
\[
\ip v{\mathcal C_kv}=r_k\ip{P_kCv}{P_kG^M_kv}\le r_k\norm{P_kCv}^2 .
\]
For $v=\delta_x$ we have $\norm{Cv}^2=(C^2)(x,x)\le\int_{\rm BZ}\hat G^2\bar dp+\rho_0(\mathcal G)\le c_N/(8\pi M)+\rho_0(\mathcal G)$, by Lemmas~\ref{app:G1:B5e2-tools}(i) and~\ref{app:G1:B5e2-I}(ii). Since $r_k\le at_k^{-2}$, this gives $s_0(\mathcal C_k)\le\gamma t_k^{-2}$. For $v=\nabla^{\pm*}_\mu\delta_x$, Step c1 with $m=M$ gives $s_1(\mathcal C_k)\le\theta t_k^{-3}$.

\emph{Step c4 (general kernels).} Let $\mathcal K\ge0$.
\begin{itemize}
\item A difference of $\mathcal K(x,y)$ in $x$ is $\ip{\nabla^{\pm*}_\mu\delta_x}{\mathcal K\delta_y}$ up to sign. So the Cauchy--Schwarz inequality for the form $\ip\cdot{\mathcal K\cdot}$ gives $\abs{\mathcal K(x,y)}\le s_0$, $\abs{\nabla\mathcal K}\le\sqrt{s_0s_1}$ and $\abs{\nabla\nabla'\mathcal K}\le s_1$.
\item $\sum_y\eps^3\mathcal K(x,y)^2=(\mathcal K^2)(x,x)\le\norm{\mathcal K}\mathcal K(x,x)$.
\item $\sum_y\eps^3\sum_\mu\abs{\nabla^+_{\mu,y}\mathcal C_k(x,y)}^2=\ip{\delta_x}{\mathcal C_k(-\Delta_\eps)\mathcal C_k\delta_x}\le\norm{(-\Delta_\eps)^{1/2}\mathcal C_k(-\Delta_\eps)^{1/2}}\,\mathcal C_k(x,x)\le\mathcal C_k(x,x)$. The last step holds because $(-\Delta_\eps)^{1/2}\mathcal C_k(-\Delta_\eps)^{1/2}\le(-\Delta_\eps)^{1/2}C(-\Delta_\eps)^{1/2}\le1$.
\item $\norm{\mathcal C_k}\le\norm C\le M^{-2}$.
\end{itemize}

\emph{Constants.} Numerically $S_4=125.3175\ldots$, $W_1\le6143.5$ and $h=\frac{944}{54}$. The properties of $\rho_a$ are Lemma~\ref{app:G1:B5e2-I}(iii). The constants are independent of $\eps$, $j$ and $k$. They are bounded for $\ell\ge\ell_0$, because $\ell^3\rho_1(\ell)$ is bounded for $\ell\ge\ell_0$ by the exponential bound of Lemma~\ref{app:G1:B5e2-I}(iii).
\end{proof}

\paragraph{Degrees of graphs with a vertex (1.13)--(1.15).} We use the vertices (1.6)--(1.15) of \cite[pp.~413--414]{Ba3}. Following \cite[p.~414]{Ba3}, the expressions (1.13)--(1.15) are vertices and (1.12) is an external field. Each expression (1.13)--(1.15) has exactly one scalar leg $\phi'(x)$ and $n\ge0$ vector legs $A'$, with $n=0$ for (1.13). The background factors $\tilde A$ are external fields. The $y$-summed scalar product of two such expressions is not a line. A line joins two scalar legs or two vector legs, and a graph has at least one line \cite[pp.~414--415]{Ba3}.

For a vertex $v$ of a graph $G$ we use the following notation.
\begin{itemize}
\item $n_\eta(v)$ is the exponent of the lattice factor $\eta$ in $v$.
\item $\ell_G(v)$ is the number of legs of $v$ on lines of $G$; a self-loop counts twice.
\item $\partial_G(v)$ is the number of differentiations of $v$ acting on legs on lines.
\item $q_G(v)$ is the number of vector legs of $v$ on lines if $v$ is of type (1.14) or (1.15), and $q_G(v)=0$ otherwise.
\item $e_G(v)$ is the number of legs of $v$ not on lines, and $\partial^{\rm ext}_G(v)$ is the number of differentiations of $v$ on such legs.
\end{itemize}
Ba{\l}aban's degree is defined in \cite[(2.1)--(2.3)]{Ba3} as follows:
\[
D_G(v)=n_\eta(v)-\tfrac12\ell_G(v)-\partial_G(v)+q_G(v),\qquad D(G)=\sum_v\tilde D(v)-3,\qquad\tilde D(v):=D_G(v)+\deg(v).
\]
Here $\deg(v)$ is the degree of the mass counterterm carried by the mass vertex (1.7), and $\deg(v)=0$ for the other vertices. A mass vertex carrying $\delta m^2_{\rm fin}$ or $\delta m^2_{K,k}$ has degree $0$ \cite[p.~423]{Ba3}. A mass vertex carrying the contribution $\delta m^2_{G_0}$ of a graph $G_0$ has degree $D(G_0)$.

\begin{lemma}[graphs with a vertex (1.13)--(1.15)]\label{app:K:L08-4}\label{app:G1:CUT-Q}
Let $d=3$, and let $D$ be the degree just defined. Then:
\begin{itemize}
\item every mass counterterm satisfies
\begin{equation}\label{app:G1:ctdeg}
\deg(\delta m^2_{G_0})=D(G_0)\ge-1;
\end{equation}
\item $D(G)\in\Z$ for every connected graph $G$;
\item every connected graph $G$ built from the vertices \cite[(1.6)--(1.15)]{Ba3}, with (1.12) an external field, that contains a vertex of type (1.13)--(1.15) satisfies $D(G)\ge1$.
\end{itemize}
The bound $D(G)\ge1$ is attained.
\end{lemma}
\begin{proof}
We prove the three statements together, by induction on the order. The counterterms that occur in a graph of a given order are of lower order \cite[p.~417]{Ba3}, so we may assume \eqref{app:G1:ctdeg} for them.

\emph{Step 1 (vertex degrees).} Reading $n_\eta(v)$ from (1.6)--(1.15) \cite[pp.~413--414]{Ba3} (cf.\ \cite[p.~423, (i)--(v)]{Ba3}) gives, for $d=3$:
\begin{itemize}
\item (1.6): $D_G(v)=1+\frac12e_G(v)$.
\item (1.8), (1.9): $D_G(v)=\frac n2+n'+\frac12e_G(v)+\partial^{\rm ext}_G(v)$, with $n+n'\ge1$; here $n'=\bar n+1$ for (1.9).
\item (1.10), (1.11): $D_G(v)=\frac n2+n'+\frac12e_G(v)$, with $n+n'\ge2$ for (1.10) and $n'=\bar n+1$ for (1.11).
\item (1.7): $D_G(v)=2+\frac12e_G(v)$.
\item (1.13)--(1.15): $D_G(v)=3-\frac12s+\frac12q$, where $s\in\{0,1\}$ indicates whether the scalar leg lies on a line and $q=q_G(v)$. Legs of $v$ that are not on lines do not contribute.
\end{itemize}

\emph{Step 2 (integrality).} Since $\sum_v\ell_G(v)=2\#\{\text{lines}\}$,
\[
D(G)=\sum_vn_\eta(v)-\#\{\text{lines}\}-\sum_v\partial_G(v)+\sum_vq_G(v)-3+\sum_v\deg(v).
\]
The degrees of $\delta m^2_{\rm fin}$ and $\delta m^2_{K,k}$ are $0$, and $\deg(\delta m^2_{G_0})=D(G_0)$ is an integer by induction. Hence $D(G)\in\Z$.

\emph{Step 3 (lower bounds).} By Step 1 and by \eqref{app:G1:ctdeg} for the counterterms of lower order (which gives $\tilde D(v)\ge1+\frac12e_G(v)$ at a mass vertex):
\begin{itemize}
\item every vertex $v$ that is not of type (1.13)--(1.15) satisfies $\tilde D(v)\ge\frac12+\frac12e_G(v)$, and even $\tilde D(v)\ge1+\frac12e_G(v)$ unless $v$ is a vertex (1.8) with $n=1$, $n'=0$;
\item a vertex of type (1.13)--(1.15) satisfies $\tilde D(v)\ge\frac52$;
\item in particular $\tilde D(v)\ge\frac12$ for every vertex, and $\tilde D(v)\ge1$ unless $v$ is a vertex (1.8) with $n=1$, $n'=0$ and all three legs on lines.
\end{itemize}

\emph{Step 4 (cases).} Let $v_0$ be a vertex of type (1.13)--(1.15).
\begin{itemize}
\item[(i)] \emph{$v_0$ is the only vertex.} Every line is a self-loop at $v_0$. A self-loop joins two legs of the same kind, and $v_0$ has only one scalar leg, so $s=0$. The graph has at least one line, so $q\ge2$. Then $D(G)=\frac q2\ge1$.
\item[(ii)] \emph{There are at least two vertices, and $q\ge1$.} Then $D_G(v_0)\ge3$, and every other vertex contributes at least $\frac12$. So $D(G)\ge\frac12$.
\item[(iii)] \emph{There are at least two vertices, and $q=0$.} By connectedness $s=1$, so $D_G(v_0)=\frac52$. Suppose $D(G)\le0$. Then $\sum_{v\ne v_0}\tilde D(v)\le\frac12$. So there is exactly one other vertex $w$, and $w$ is a vertex (1.8) with $n=1$, $n'=0$ and all three legs on lines. The only line at $v_0$ joins its scalar leg to a scalar leg of $w$. The remaining scalar leg and the vector leg of $w$ would then have to form a line together, which is impossible. Hence $D(G)\ge\frac12$.
\end{itemize}
In every case $D(G)>0$, so $D(G)\ge1$ by Step 2.

\emph{Proof of \eqref{app:G1:ctdeg}.} The graph $G_0$ is one-particle irreducible, with exactly two external legs, both scalar \cite[p.~417]{Ba3}. If $G_0$ contains a vertex of type (1.13)--(1.15), the argument above gives $D(G_0)\ge1$. Otherwise Step 3 gives $\tilde D(v)\ge\frac12+\frac12e_{G_0}(v)$ for every vertex, and $\sum_v\frac12e_{G_0}(v)=1$.
\begin{itemize}
\item If $G_0$ has a single vertex, this vertex is not a vertex (1.8) with $n=1$: its vector leg could be neither external nor on a line. So $\sum_v\tilde D(v)\ge1+1=2$.
\item If $G_0$ has at least two vertices, then $\sum_v\tilde D(v)\ge\frac12+\frac12+1=2$.
\end{itemize}
Hence $D(G_0)=\sum_v\tilde D(v)-3\ge-1$.

\emph{Attainment.} Take a vertex (1.13) whose scalar leg is joined by a line to the differentiated scalar leg of a vertex (1.8) with $n=1$, $n'=0$. Let the undifferentiated scalar leg and the vector leg of the latter be external. Then $D=\frac52+(3-\frac12-1)-3=1$.
\end{proof}

\subsection{Renormalised classes under cutting, bad localisations, H\"older bounds and the nearest-neighbour bound}\label{app:G1:cutHol}

Throughout $d=3$. We work on the $\eta$-lattice of scale $k$, rescaled so that $L^k\eta=1$, and write $t=L^k\eps$, $e(t)^2=e^2t$, $\lambda(t)=\lambda t$. The graphs are those of \cite[(1.17)--(1.18)]{Ba3}, built from the vertices \cite[(1.6)--(1.11), (1.13)--(1.15)]{Ba3}. We use the following names for the vertices:
\begin{itemize}
\item $\mathsf L$ is (1.6);
\item $\mathsf A_n$ is (1.8), with $n$ vector legs and the derivative $D^\eta$ on one scalar leg;
\item $\mathsf S_n$ is (1.10);
\item $\mathsf M_1,\mathsf M_2$ are the counterterm vertices (1.7) of a divergent self-energy of degree $-1$, $0$;
\item $\mathsf M_f$ is a finite counterterm vertex ($\delta m^2_{\rm fin}$, $\delta m^2_{K,k}$);
\item $\mathsf Q_n$ are (1.13)--(1.15), and $\mathsf R$ are (1.9), (1.11).
\end{itemize}
A graph with coupling factor $e^\alpha\lambda^\beta$ has \emph{order} $o:=\alpha/2+\beta$. A vertex $\mathsf M_1,\mathsf M_2$ carries the order of the graph defining its counterterm, and $\mathsf M_f$ carries the order $2$. We write $E_s,E_v$ for the numbers of external scalar and vector legs; background legs count as vector legs \cite[p.~423]{Ba3}. $E_\partial$ is the number of differentiations of (1.8) that act on external legs. $D$ is the degree of \cite[(2.1)--(2.3)]{Ba3}, and a graph is \emph{divergent} if $D\le0$.

\paragraph{Degrees and the divergent classes.}

\begin{lemma}[degree identity]\label{app:G1:H3-deg}
For every connected graph $G$ built from $\mathsf L,\mathsf A_n,\mathsf S_n,\mathsf M_1,\mathsf M_2,\mathsf M_f$,
\[
D(G)=o(G)+\tfrac12(E_s+E_v)+E_\partial-3 .
\]
\end{lemma}
\begin{proof}
By \cite[(2.1)]{Ba3}, a leg external to $G$ contributes $0$ instead of $-\frac12$ to $D_G(v)$, and a differentiation on an external leg contributes $0$ instead of $-1$. Hence by \cite[(2.2)--(2.3)]{Ba3}
\[
D(G)=\sum_vD(v)+\tfrac12(E_s+E_v)+E_\partial-3+\sum\deg ,
\]
where the last sum runs over the degrees of the counterterms at the vertices $\mathsf M_1,\mathsf M_2$. The vertex degrees are those of \cite[p.~423]{Ba3}:
\begin{itemize}
\item $D(\mathsf L)=1$;
\item $D(\mathsf A_n)=D(\mathsf S_n)=n/2$;
\item $D(\mathsf M)=2$.
\end{itemize}
Thus $D(v)=o(v)$ for $\mathsf L,\mathsf A_n,\mathsf S_n,\mathsf M_f$. Let $\mathsf M$ carry the counterterm $\delta m^2_{G_0}$. Then $G_0$ has two scalar legs and $E_\partial(G_0)=0$, so by induction on the order $D(G_0)=o(G_0)-2$. Hence $D(v)+\deg=o(G_0)=o(v)$, and summing over the vertices gives the claim.
\end{proof}

\begin{lemma}[the divergent classes]\label{app:G1:H3-class}
\begin{itemize}
\item[(i)] A graph that contains an $\mathsf R$-vertex, or a vertex (1.13)--(1.15), has $D>0$.
\item[(ii)] The divergent connected graphs with at least one line are exactly those of the following table, together with the divergent graphs of the leg types $(0,1)$ and $(0,3)$. The latter form the class $\mathsf C6$, and their graphs vanish.
\end{itemize}
\begin{center}\small
\begin{tabular}{llll}
$(E_s,E_v)$ & $o$ & $E_\partial$, $D$ & class \\\hline
$(0,0)$ & $1,2,3$ & $D=-2,-1,0$ & vacuum \\
$(2,0)$ & $1$ & $0$, $-1$ & $\mathsf C1$: self-energy \cite[(3.6)]{Ba3}\\
$(2,0)$ & $1$ & $1$, $0$ & $\mathsf C1'$: \cite[(3.21)]{Ba3}\\
$(2,0)$ & $2$ & $0$, $0$ & $\mathsf C2$: self-energy \cite[(3.18)]{Ba3}\\
$(0,2)$ & $1$, $2$ & $0$; $-1$, $0$ & $\mathsf C3$, $\mathsf C4$: vacuum polarisation\\
$(2,1)$ & $\frac32$ & $0$, $0$ & $\mathsf C5$: \cite[(2.18)--(2.22)]{Ba3}\\
$(2,2)$ & $1$ & $0$, $0$ & $\mathsf C7$: the tree \cite[(2.4)]{Ba3}
\end{tabular}
\end{center}
\begin{itemize}
\item[(iii)] Let $G$ have at least one line and contain an $\mathsf M_f$ vertex. Then $D(G)>0$, unless $G$ is a vacuum graph or vanishes.
\end{itemize}
\end{lemma}
\begin{proof}
(i) An $\mathsf R$-vertex has at least $\nbar+1>13$ vector legs, so its degree exceeds $6$, while every other vertex has nonnegative degree. Graphs with a vertex (1.13)--(1.15) have $D\ge1$ \cite[pp.~423, 429]{Ba3}.

(ii) Each photon line uses two vector legs, so $o\ge I_v+E_v/2$, where $I_v$ is the number of photon lines. Every vertex has an even number of scalar legs, so $E_s$ is even. The derivatives of (1.8) sit on scalar legs, so $E_\partial=0$ if $E_s=0$.

A graph with a line and $o=\frac12$ is a single $\mathsf A_1$ with a scalar self-line, of type $(0,1)$. Otherwise $o\ge1$, and $o\ge\frac32$ if $E_v\ge3$ or $(E_s,E_v)=(2,1)$.

By Lemma~\ref{app:G1:H3-deg}, $D\le0$ means $o+\frac12(E_s+E_v)+E_\partial\le3$. Going through the leg types:
\begin{itemize}
\item $E=0$ gives $o\le3$;
\item $(2,0)$ gives $(o,E_\partial)\in\{(1,0),(2,0),(1,1)\}$;
\item $(0,2)$ gives $o\in\{1,2\}$;
\item $(2,1)$ gives $o=\frac32$, $E_\partial=0$;
\item $(2,2)$ gives $o=1$, $E_\partial=0$, i.e.\ two $\mathsf A_1$ joined by a scalar line with internal derivatives, which is the tree (2.4);
\item $(0,1)$, $(0,3)$ are the remaining odd cases.
\end{itemize}
The case $(4,0)$ needs $o=1$. Order $1$ with a line means $\mathsf L$ or $\mathsf A_2$ with a self-line (at most two external scalar legs), or two $\mathsf A_1$ joined by a photon line ($E_\partial=2$). All other leg types have $o+\frac12E>3$.

The class $\mathsf C6$ vanishes: after the gauge transformation removing the constant part of the background \cite[p.~433]{Ba3}, all scalar propagators are proportional to $1_N$. Then every scalar strand without external scalar legs carries $\tr q^{m}$ with some odd $m$, and this trace vanishes because $q$ is antisymmetric.

(iii) Let $G$ be neither a vacuum graph nor vanishing. By (i) we may assume that $G$ has no $\mathsf R$-vertex and no vertex (1.13)--(1.15). A graph with a line whose only vertex is $\mathsf M_f$ is the vacuum tadpole, so $G$ has a vertex besides $\mathsf M_f$, and $o\ge\frac52$. Moreover $E_s$ is even and $(E_s,E_v)=(0,1)$ is of class $\mathsf C6$, so $E_s+E_v\ge2$. Hence $D\ge\frac52+1-3>0$ by Lemma~\ref{app:G1:H3-deg}.
\end{proof}

\emph{Partners.} The class $G_{\rm ren}$ of \cite[pp.~420, 429--432]{Ba3} is the smallest class closed under the following operations on divergent subgraphs $S$:
\begin{itemize}
\item[(P1)] if $S$ is a one-particle-irreducible $\mathsf C1$ or $\mathsf C2$, replace $S$ by its counterterm;
\item[(P2)] if $S\in\mathsf C3,\mathsf C4$, add all connected graphs of the same order with two external vector legs and the same gauge-invariant insertions;
\item[(P3)] if $S\in\mathsf C5$, add the members of (2.18)--(2.21) and the permuted graphs;
\item[(P4)] for vacuum graphs, add the vacuum counterterm.
\end{itemize}
One-particle-reducible two-scalar-leg graphs receive no counterterm \cite[(1.21), p.~430]{Ba3}.

\begin{definition}[graph form of the counterterm members]\label{app:G1:H3-R}
Let $S\subset T$ be a divergent two-scalar-leg subgraph whose external legs are the lines (or external fields) $\ell_1$ at $v_1$ and $\ell_2$ at $v_2$, where $\ell_2$ is the leg of lower index \cite[p.~435]{Ba3}. The partner $T_{\rm ct}=T[S\to\delta m^2_S]$ is the graph with the same vertices and labels as $T$. Its $S$-lines are replaced by the lines of the counterterm, every other line of $T$ is kept, and the end of $\ell_2$ is moved from $v_2$ to $v_1$:
\[
T_{\rm ct}=-\sum_{x,x'}\bigl[\text{rest of }T\text{ with }\ell_1,\ell_2\text{ at }x\bigr]\,\eta^{2d}\Sigma_S(x,x'),
\]
where $x,x'$ are the positions of $v_1,v_2$. Finite counterterm vertices remain one-vertex graphs, localised by the unit cube of their position.
\end{definition}
By $\delta m^2_S(x)=\sum_{x'}\eta^d\Sigma_S(x,x')$ \cite[p.~417]{Ba3} and the localisation \cite[(1.31), p.~420]{Ba3}, $T_{\rm ct}$ equals the one-vertex member of \cite{Ba3} at each label configuration. For uncut classes we return to the one-vertex form.

A \emph{label} is each vector leg (one factor $e$), each $\lambda$-vertex, and the cube of each finite counterterm vertex and of each vertex (1.13)--(1.15). A \emph{label configuration} $\Lambda$ assigns a localisation cube to every label and to every external leg; $N_s$ is the number of these. The support cubes (side $2$) of $\Lambda$ are grouped into \emph{clusters}, the components of the relation $\dist(\square,\square')<1$. A line of a member is \emph{cut} at $\Lambda$ if its end vertices lie in different clusters.

\begin{lemma}[labels and clusters]\label{app:G1:H3-slots}
\begin{itemize}
\item[(i)] All members of a class have the same labels: $\alpha$ vector labels, of which $E_v$ are external, and $\beta$ $\lambda$-labels.
\item[(ii)] In the form of Definition~\ref{app:G1:H3-R}, every vertex carries only labels whose supports contain its position.
\item[(iii)] Every cut line joins vertices $v,v'$ with $\dist(\square(v),\square(v'))\ge1$.
\item[(iv)] After cutting, every connected component lies in one cluster, and each cluster lies in a cube of side at most $3N_s$.
\item[(v)] For every member $T$, $d(\{\square(v)\})\le\sum_{\ell\text{ cut in }T}\dist(\square(v_\ell),\square(v'_\ell))+3N_s$.
\end{itemize}
\end{lemma}
\begin{proof}
(i) Every vertex with $n$ vector legs carries $e^n$, so $\alpha=2I_v+E_v$, and $\beta$ counts the $\lambda$-vertices. Partners carry the labels of the subgraph they replace; for counterterms this is Definition~\ref{app:G1:H3-R}. The finite insertions are common to all members of a smallest class.

(ii) The vertices of $T_{\rm ct}$ are those of $T$, each with its own labels, and a moved line carries no label.

(iii) Different clusters are at distance at least $1$.

(iv) An uncut line stays inside one cluster, and a chain of $m$ side-$2$ cubes with gaps less than $1$ has diameter less than $3m$.

(v) $T$ is connected. A spanning tree of the cubes is obtained from the cut lines together with a tree inside each cluster, and by (iv) these trees have total length at most $3N_s$.
\end{proof}

\begin{lemma}[degrees under cutting]\label{app:G1:H3-degcut}
\begin{itemize}
\item[(a)] For $S\subset G\setminus\ell$, $D(S)$ is the same in $G$ and in $G\setminus\ell$. Hence the divergent subgraphs of $G\setminus\ell$ are the divergent subgraphs of $G$ not containing $\ell$.
\item[(b)] Let $S$ be divergent, non-vacuum and non-vanishing, $\ell\in S$, and $S\setminus\ell$ connected. Then $D(S\setminus\ell)\ge0$, with equality only for the tree (2.4).
\end{itemize}
\end{lemma}
\begin{proof}
(a) By (2.1), $D(S)$ counts only the legs of the vertices of $S$ on lines of $S$.

(b) $S\setminus\ell$ has the vertices of $S$ and two more external legs, possibly with differentiations moved from internal to external. So $D(S\setminus\ell)\ge D(S)+1\ge0$, because $D(S)\ge-1$ in every non-vacuum class of Lemma~\ref{app:G1:H3-class}. By that table, a connected graph with $E(S)+2$ legs and order $o(S)$ has $D\le0$ only in the class $(2,2)$, $o=1$, which is the tree (2.4).
\end{proof}

\begin{theorem}[closure of the renormalised classes under cutting]\label{app:K:L09-3}\label{app:G1:H3-closure}
Let $G_{\rm ren}$ be a renormalised class of \cite[pp.~420, 429--432]{Ba3}, with counterterm members written as in Definition~\ref{app:G1:H3-R}, and let $\Lambda$, $N_s$ and the clusters be as above. Then:
\begin{itemize}
\item[(a)] at each $\Lambda$, either no member has a cut line, or every member that does not vanish at $\Lambda$ has one;
\item[(b)] every cut line joins vertices with $\dist(\square(v),\square(v'))\ge1$, every cluster lies in a cube of side $\le3N_s$, and
\[
d(\{\square(v)\})\le\sum_{\ell\ \rm cut}\dist(\square(v_\ell),\square(v'_\ell))+3N_s ;
\]
\item[(c)] the cut family, with cut lines read as two-variable external fields, is a disjoint union of renormalised classes, and each connected component of a member is localised in one cluster. Hence \cite[(3.2)]{Ba3} for these classes implies \cite[(1.33)]{Ba3}, with $\mathrm{e}^{-\delta_0d(\{\square(v)\})}$ recovered up to $\mathrm{e}^{3\delta_0N_s}$;
\item[(d)] for a divergent vacuum graph $G$ and a line $\ell$ of $G$, either the term vanishes ($\ell$ is a photon bridge), or $G\setminus\ell$ is a bare vertex, convergent, or a divergent two-leg graph whose partners are members of the vacuum class of $G$.
\end{itemize}
\end{theorem}
\begin{proof}
(a) By Lemma~\ref{app:G1:H3-slots}(i) every member uses all labels of $\Lambda$. If some member has no cut line, then by (ii) all its labels, hence all labels, lie in one cluster, and no member has a cut line. Moreover, let $T''$ arise from $T$ by a partner operation on a divergent $S'\subset T$ none of whose lines is cut. The labels of $S'$ lie in one cluster $K$, the partner carries the same labels, and every moved line joins two vertices of $K$. So $T''$ has exactly the cut lines of $T$.

(b) is Lemma~\ref{app:G1:H3-slots}(iii)--(v).

(c) By (a) a class is never split. The cut family is closed under the partner operations of its uncut divergent subgraphs, with the inherited coefficients. Now let $S\subset T$ be divergent with at least one cut line. There are three cases.
\begin{itemize}
\item[(i)] A one-vertex partner of $S$ (for instance the seagull tadpole partners of a cut $\mathsf C3$) has all its labels in one cluster, while those of $S$ lie in at least two. So it vanishes at $\Lambda$.
\item[(ii)] The counterterm partner has the same vertices, labels and $S$-lines as $T$. It is cut along the same lines of $S$, and its $S$-part splits into the same pieces as $S\setminus{\rm cut}$, with equal degrees. By Lemma~\ref{app:G1:H3-degcut} these degrees are $\ge0$, with equality only for (2.4), which is treated by \cite[(2.8)]{Ba3}.
\item[(iii)] Let $Q$ be a subgraph of $T_{\rm ct}\setminus{\rm cut}$ containing the moved line $\ell_2$. Write $Q=Q_S\cup Q_R$, where $Q_S=Q\cap S\ni v_1$ and $Q_R$ is the rest together with a point vertex $V_x$ carrying $\ell_1,\ell_2$. Counting the volume factor of $v_1$ once, (2.1)--(2.2) give $D(Q)=D(Q_S)+D(Q_R)$.
\begin{itemize}
\item If $v_2\in Q$: moving the end of a line between two vertices of $Q$ changes no term of (2.1), since the moved leg carries no derivative. The subgraph $Q^T$ of $T$ with the same lines is connected, because $v_2$ meets the rest of $T$ only through $\ell_2$. Hence $D(Q)=D(Q^T)$, and the partners of $Q$ are the images of those of $Q^T$.
\item If $v_2\notin Q$: $Q_R$ is a non-vacuum graph with a line ($\ell_2$) and a finite mass vertex, so $D(Q_R)>0$ unless $Q$ vanishes (Lemma~\ref{app:G1:H3-class}(iii)). Moreover $D(Q_S)\ge0$. Indeed, $S$ is one-particle irreducible, since only such $S$ receive a counterterm (P1), and $Q_S\subsetneq S$, since $v_2\in S\setminus Q$. So at least two lines of $S$ leave $Q_S$, $Q_S$ has at least three external legs, and the table gives $D(Q_S)\ge0$. Hence $D(Q)>0$.
\end{itemize}
\end{itemize}
If $\ell_2$ itself is cut, it becomes a two-variable external field, and the pieces are covered by (ii)--(iii).

The three uses of the cut in \cite[pp.~432--433]{Ba3} are supplied by (b): the kernel bound \cite[(3.1)]{Ba3} on cut lines, the localisation of each component in a cube of size $O(N_s)$, and the exponential factor. For the last one, by (b),
\[
\prod_{\ell\ \rm cut}\mathrm{e}^{-\delta_0\dist(\square(v_\ell),\square(v'_\ell))}\le\mathrm{e}^{3\delta_0N_s}\mathrm{e}^{-\delta_0d(\{\square(v)\})}.
\]
The same rule is used at the second cut of \cite[p.~433]{Ba3}. The expansion in the background field is multilinear over the whole class: its order-$n$ term is the full class with $n$ extra external legs, each with its own label. The splitting of the propagator differences acts line by line. The induction on (order, number of lines) is used for the first cut only. The pieces of the second cut are renormalised directly, by \cite[Prop.~2.1]{Ba3} and \cite[p.~433]{Ba3}.

(d) $E_s$ is even for every connected graph, so a scalar line is never a bridge. A photon bridge leaves two graphs with one vector leg each, of class $\mathsf C6$, so the term vanishes. Otherwise $G\setminus\ell$ is connected, has the order of $G$ and two external legs, and has $E_\partial$ equal to the number of derivatives on $\ell$. By Lemma~\ref{app:G1:H3-deg}, $D(G\setminus\ell)=o-2+E_\partial$. Reading off the table:
\begin{itemize}
\item for scalar $\ell$, $(o,E_\partial)=(1,0)$, $(2,0)$, $(1,1)$ give $\mathsf C1$, $\mathsf C2$, $\mathsf C1'$, and all other cases are convergent;
\item for a photon $\ell$, $o=1$ gives $\mathsf C3$, $o=2$ gives $\mathsf C4$, and $o=3$ is convergent;
\item the tadpole of $\mathsf M_f$ gives a bare vertex.
\end{itemize}
$G\setminus\ell$ is a divergent subgraph of $G$, so the partners it needs are generated inside the class of $G$. One-particle-reducible $\mathsf C2$ pieces need no partner (Theorem~\ref{app:K:L09-6}(b)).
\end{proof}

\paragraph{Bad localisations.}
Let $0\le k<K$, $W:=B^k(\Lambda_7^{(k)})$, $c_d:=2\sqrt d$ and $\delta_1:=\delta_0/2$. By \cite[(2.8)]{Ba2}, $\dist(x,W^c)>r(t)-c_d$ for $x\in B^k(\Lambda_8^{(k)})$. Let
\[
\Delta_k(x):=[\delta m^2_{K,k}(x)-\delta m^2_{K,k+1}(x)]t^2=\sum_{G_0}c^{[k]}_{G_0}(x)
\]
(notation of Lemma~\ref{app:G1:I5-ct}), where the sum runs over the counterterm graphs, with their partners in the form of Definition~\ref{app:G1:H3-R}, whose largest line index is exactly $k$. Localise as in \cite[p.~420]{Ba3}, $\Delta_k(x)=\sum_\Lambda\Delta_k(x;\Lambda)$, with the pinned cube $\square(p)\ni x$. A localisation is \emph{bad} if one of its cubes is not contained in $W$, and $\Delta^{\rm bad}_k(x)$ is the sum of the $\Delta_k(x;\Lambda)$ over the bad $\Lambda$.

\begin{lemma}[exact cut]\label{app:G1:H3-cut}
Decompose the index sums of the vacuum terms $E_{K,k}$ of step $k$ by orderings \cite[(3.4)]{Ba3}, breaking ties by a fixed order of the lines. Then:
\begin{itemize}
\item the last line $l(m)$ has index $k$;
\item the partner operations active in an ordering act on components of $G_1\subset\dots\subset G_{m-1}$, none of which contains $l(m)$.
\end{itemize}
Hence cutting $l(m)$, and reading its index-$k$ kernel $K_k$ as a two-variable external field, maps the vacuum class, term by term and with the same coefficients, to the renormalised two-leg classes of $G\setminus l(m)$ listed in Theorem~\ref{app:K:L09-3}(d).
\end{lemma}
\begin{proof}
The first statement is the definition of the index sets \cite[p.~424]{Ba3}. The second follows from $l(m)\notin G_i$ for $i<m$ and Lemma~\ref{app:G1:H3-degcut}(a). When $l(m)$ is a line of the graph attached to a $\delta m^2_K$ vertex, it is read in the form of Definition~\ref{app:G1:H3-R}.
\end{proof}

\begin{proposition}[bad localisations of the counterterms]\label{app:K:L09-4}\label{app:G1:H3-bad}
Assume $\ell_T/t\ge6r(\eps_0)$, $\mu_*t\le1$ and $m^2t^2\le1$, and put $P_m(t):=(e^2+\lambda)t+(e^4+e^2\lambda+\lambda^2)t^2$.
\begin{itemize}
\item[(i)] There is $C_d$, independent of $\eps,k,x$, the domains and $\ell_T$, with $\abs{\Delta_k(x)},\abs{\Delta^{\rm bad}_k(x)}\le C_dP_m(t)$, and
\[
\abs{\Delta^{\rm bad}_k(x)}\le C_dP_m(t)\mathrm{e}^{3\delta_1c_d}\mathrm{e}^{-\delta_1r(t)}\qquad(x\in B^k(\Lambda_8^{(k)})).
\]
\item[(ii)] Let $E_{K,k}(\Lambda)$ be the localised vacuum terms of step $k$ (largest line index $k$), anchored at the cube $\square(v_0)$ of a fixed vertex, and put $P_v(t):=\sum_{2\le\alpha+2\beta\le2\nbar}e(t)^\alpha\lambda(t)^\beta$. There is $C=O(\alpha_0,\nbar)$, independent of $\eps,k$ and the domains, with
\[
\sum_{\Lambda:\square(v_0)\subset X}\abs{E_{K,k}(\Lambda)}\le CP_v(t)\abs X
\]
for every union $X$ of unit cubes, and
\[
\sum_{\Lambda\ {\rm bad},\,\square(v_0)\subset B^k(\Lambda_8^{(k)})}\abs{E_{K,k}(\Lambda)}\le CP_v(t)\mathrm{e}^{3\delta_1c_d}\mathrm{e}^{-\delta_1r(t)}\abs{B^k(\Lambda_8^{(k)})}.
\]
\end{itemize}
One may take $C_d=C_PS$, where $C_P$ is the constant of the pinned bound times the number of counterterm graphs, and
\[
S:=\max_{N_s}\sum_{D\ge0}\bigl(c'_d(D+2)^d\bigr)^{N_s}\mathrm{e}^{-\delta_0D/2},
\]
the maximum being over the label numbers of these graphs.
\end{proposition}
\begin{proof}
\emph{Step 1 (one localisation).} Lemma~\ref{app:G1:VSR-Btop} at the top index $k$ keeps every renormalisation inside the sub-sum with top index $k$. Theorem~\ref{app:K:L09-3} turns long lines into external fields with the factor $\mathrm{e}^{-\delta_0\dist}$. Lemma~\ref{app:G1:VSR-A}(ii) and Lemma~\ref{app:G1:VSR-C} then give
\[
\abs{\Delta_k(x;\Lambda)}\le C_Pe(t)^\alpha\lambda(t)^\beta\mathrm{e}^{-\delta_0d(\Lambda)},
\]
since $(L^k\eta)^{D(G_0)}=1$. The hypotheses $\ell_T/t\ge6r(\eps_0)$, $\mu_*t\le1$, $m^2t^2\le1$ are those of these lemmas, and they make $C_P$ independent of $\eps,k,x$, the domains and $\ell_T$.

\emph{Step 2 (the sum over localisations).} With $\square(p)\ni x$ fixed, every cube of $\Lambda$ lies within distance $d(\Lambda)+2\sqrt d$ of $\square(p)$. Hence
\[
\#\{\Lambda:d(\Lambda)\in[D,D+1)\}\le\bigl(c'_d(D+2)^d\bigr)^{N_s},
\qquad
\sum_\Lambda\mathrm{e}^{-\delta_0d(\Lambda)}\le\sum_\Lambda\mathrm{e}^{-\delta_0d(\Lambda)/2}\le S .
\]
The counterterm graphs have order $1$ or $2$, so $\sum_{G_0}e(t)^\alpha\lambda(t)^\beta$ is at most the number of graphs times $P_m(t)$. This gives the first two bounds of (i). The sub-sum over bad $\Lambda$ obeys the same bound.

\emph{Step 3 (bad localisations deep inside).} Let $x\in B^k(\Lambda_8^{(k)})$, and let $\Lambda$ be bad, with a cube $Q\ni z$, $z\notin W$. The tree length $d(\Lambda)$ is at least the distance of any two of its cubes, and the cubes have diameter $c_d$. So
\[
d(\Lambda)\ge\dist(\square(p),Q)\ge\abs{x-z}-2c_d>r(t)-3c_d .
\]
Hence
\[
\sum_{\Lambda\ \rm bad}\mathrm{e}^{-\delta_0d(\Lambda)}\le\mathrm{e}^{-\delta_1(r(t)-3c_d)}\sum_\Lambda\mathrm{e}^{-\delta_0d(\Lambda)/2}\le\mathrm{e}^{3\delta_1c_d}\mathrm{e}^{-\delta_1r(t)}S,
\]
which gives the last bound of (i).

\emph{Step 4 (vacuum terms).} By Lemma~\ref{app:G1:H3-cut} and Theorem~\ref{app:K:L09-3}(d), $E_{K,k}$ is a sum of renormalised two-leg classes with external field $K_k$. In scale-$k$ units, $\norm{hK_kh'}_{1,\alpha_0}\le O(1)\mathrm{e}^{-\delta_1\dist(\supp h,\supp h')}$, by \cite[(2.10)--(2.11), p.~426]{Ba3} and Lemma~\ref{app:K:L09-5}. Then \cite[(1.33)]{Ba3}, which admits external fields of several variables \cite[p.~419]{Ba3}, bounds each localised class by $O(1)e(t)^\alpha\lambda(t)^\beta\mathrm{e}^{-\delta_0d(\Lambda)}$.

A photon cut produces $\mathsf C3$ or $\mathsf C4$ classes. They are renormalised by their Ward--Takahashi classes, which belong to the vacuum class, and the coefficient of the local vertex is $(L^{j_0}\eta)^{-d+2}=1$ at $j_0=k$.

Summing over $\Lambda$ with the anchor fixed, as in Steps~2--3, gives $\le CP_v(t)$ per anchor cube, and $\le CP_v(t)\mathrm{e}^{3\delta_1c_d}\mathrm{e}^{-\delta_1r(t)}$ per anchor cube in $B^k(\Lambda_8^{(k)})$ for the bad ones. Summing over the at most $\abs X$ anchor cubes, respectively $\abs{B^k(\Lambda_8^{(k)})}$ anchor cubes, gives (ii). The orders are bounded by $\nbar$, so $C=O(\alpha_0,\nbar)$.
\end{proof}

\paragraph{H\"older bound of the vector backgrounds.}

\begin{lemma}[H\"older bound of the vector backgrounds]\label{app:K:L09-5}\label{app:G1:H3-HOv}
Let $\Omega$ be a union of big blocks, $G_k:=G_k(\Omega,0)$ the vector propagator with Neumann boundary conditions, and $\tilde A:=a_kG_kQ_k^*B$. Here $B$ is a unit-lattice field with $\abs{B(y)-B(y')}\le c_1p\abs{y-y'}_1$, $c_1=3Ld$ \cite[(2.17)]{Ba2}. Let $\alpha<1$ and $x,x'\in\Omega$ with $\dist(\{x,x'\},\Omega^c)\ge R_0+1$. Then
\[
\abs{x-x'}^{-\alpha}\abs{\partial^\eta_\mu\tilde A(x')-\partial^\eta_\mu\tilde A(x)}\le C_{\rm H}p,
\]
\[
C_{\rm H}:=2c_0(\alpha)a_kc_1\mathrm{e}^{\delta_0(\sqrt d+1)}\sum_{z\in\Z^d}(1+\abs z_1)\mathrm{e}^{-\delta_0\abs z}+2\sup_\Omega\abs{\partial\tilde A}/p ,
\]
where $c_0(\alpha),\delta_0$ are the constants of \cite[(1.9)]{Ba4}. The same holds for $B^{(k+1),\eta}$.
\end{lemma}
\begin{proof}
(i) For a constant $c$, $Q_k^*c=c$, $Q_kc=c$ and $-\Delta^N_\Omega c=0$. Hence
\[
(-\Delta^N_\Omega+\mu_0^2t^2+a_kQ_k^*Q_k)\frac{ca_k}{a_k+\mu_0^2t^2}=a_kQ_k^*c ,
\]
so $a_kG_kQ_k^*c$ is constant, and $\partial\tilde A$ does not change under $B\mapsto B-c$.

(ii) If $\abs{x-x'}>1$, the quotient is at most $2\sup_\Omega\abs{\partial\tilde A}\le C_{\rm H}p$.

(iii) If $\abs{x-x'}\le1$, let $y_0$ be the block of $x$. By (i) we may replace $B$ by $B-B(y_0)$. Write $a_kQ_k^*(B-B(y_0))=\sum_yf_y$, where $f_y$ is supported in the block $B(y)$ and $\norm{f_y}_\infty\le a_kc_1p\abs{y-y_0}_1$. For $w\in B(y)$ we have $\abs{w-x}\ge\abs{y-y_0}-\sqrt d$, so $\dist(\{x,x'\},B(y))\ge\abs{y-y_0}-\sqrt d-1$. Under the distance hypothesis, \cite[(1.9)]{Ba4} applies to each $G_kf_y$ at $x,x'$ and gives
\[
\abs{\partial^\eta_\mu G_kf_y(x')-\partial^\eta_\mu G_kf_y(x)}\le c_0(\alpha)\abs{x-x'}^\alpha\mathrm{e}^{-\delta_0(\abs{y-y_0}-\sqrt d-1)}\norm{f_y}_\infty .
\]
Summing over $y=y_0+z$ gives the bound $c_0(\alpha)a_kc_1\mathrm{e}^{\delta_0(\sqrt d+1)}\sum_z\abs z_1\mathrm{e}^{-\delta_0\abs z}\,p\,\abs{x-x'}^\alpha\le C_{\rm H}p\abs{x-x'}^\alpha$.

The argument uses only the form $a_kG_kQ_k^*B$ with a field $B$ satisfying the Lipschitz bound. It therefore applies verbatim to $B^{(k+1),\eta}$, which is a field of this form with the block field of the restrictions \cite[(2.17)]{Ba2}.
\end{proof}

\paragraph{Nested and one-particle-reducible classes.}
Let $H$ be an inner divergent subgraph whose external legs carry indices at least those of its lines. The transformation \cite[(3.5)]{Ba3} for $H$ produces two kinds of terms:
\begin{itemize}
\item[(T1)] \emph{remainder terms}: a factor $\abs{x-x'}^s$ between two vertices of $H$, and an order-$s$ difference quotient on one leg $\ell$ of $H$ (on two legs if $H\in\mathsf C5$). The legs chosen are those of lowest index \cite[pp.~435, 443]{Ba3}, and $s=1+\alpha$ if $D(H)=-1$, $s=\alpha$ if $D(H)=0$;
\item[(T2)] \emph{local terms}: a local vertex at $x$, with the internal sums of $H$ resummed up to $j_0=j_\ell$ \cite[(3.15), (3.23), (3.30), (3.32), (3.36)]{Ba3}.
\end{itemize}

\begin{lemma}[differentiated legs are lines]\label{app:G1:H3-N1}
Let $H\subsetneq G_0$ be divergent and not of type (2.4), where $G_0$ is a non-vanishing member of a $\mathsf C2$, $\mathsf C4$ or $\mathsf C5$ class. Then the legs of $H$ that (T1) and (T2) differentiate are lines of $G_0$.
\end{lemma}
\begin{proof}
In the orderings admitted at \cite[p.~445]{Ba3}, every line of $G_0$ has index below every external leg of $G_0$. So it suffices that $H$ has at least one (for $H\in\mathsf C5$, two) legs that are lines of $G_0$.
\begin{itemize}
\item If $H$ has two legs and both were external legs of $G_0$, connectedness would force $G_0=H$.
\item If $H\in\mathsf C5$ and $G_0\in\mathsf C4$, the scalar legs of $H$ are lines of $G_0$.
\item If $H\in\mathsf C5$ and $G_0\in\mathsf C2$, the vector leg of $H$ is a line of $G_0$. If both scalar legs of $H$ were external legs of $G_0$, then $G_0\setminus H$ would be attached by one photon line, so it would be of class $\mathsf C6$ and $G_0$ would vanish.
\item $H\in\mathsf C5$ and $G_0\in\mathsf C5$ is impossible, since divergent $\mathsf C5$ graphs have at most three vertices.
\end{itemize}
\end{proof}

\begin{lemma}[degree bookkeeping]\label{app:G1:H3-N2}
Fix an ordering of $G_0$ in which $H$ is the last divergent proper subgraph, completed at step $i_H$. Replace $H$ by a term (T1) with leg $\ell$, splitting $\abs{x-x'}^s$ along a chain of lines of $H$ as in \cite[p.~438]{Ba3}. Then the new degree $D'(G_i)$ of a subgraph $G_i$ of the ordering is:
\begin{itemize}
\item $D'(G_i)=D(G_i)+s$ if $G_i\supseteq H$ and $\ell\notin G_i$;
\item $D'(G_i)=D(G_i)$ if $\ell\in G_i$, or if $G_i$ contains no line of $H$;
\item $D'(G_i)\ge D(G_i)$ if $G_i$ contains some but not all lines of $H$.
\end{itemize}
In particular $D'(H)>0$, $D'(G_0)=D(G_0)$, and every $G_i$ with $i_H<i<m$ keeps a positive degree. The term (T2) enters every $G_i\supseteq H$ as a shrunk vertex of \cite[(2.16)]{Ba3}, with no internal index sum left.
\end{lemma}
\begin{proof}
Since $\ell\notin G_{i_H}$, every $G_i\ni\ell$ contains $H$. Each chain piece adds $+s$ to the degree of every $G_i$ containing its line. By (2.1), the difference quotient on $\ell$ contributes $-s$ exactly when $\ell$ is internal to $G_i$. By Lemma~\ref{app:G1:H3-N1}, $\ell\in G_0$.

If $\ell$ is the transported leg of two inner subgraphs, it is differentiated at both ends. The line bounds \cite[(2.10)--(2.11)]{Ba3} hold in both variables, so the two counts add. The resulting bound is \cite[Prop.~2.2, p.~428]{Ba3}.
\end{proof}

\begin{lemma}[partners are transformed identically]\label{app:G1:H3-N3}
Group the pairs (graph, ordering) of the class of $G_0$ by outside structure, outside ordering and the class of the last divergent subgraph $H$. This grouping is a partition, and each group contains every member of the class of $H$. The inner transformation is an exact identity, so after the inner transformations every whole-graph sum equals the untransformed whole-graph sum.
\end{lemma}
\begin{proof}
Degrees depend only on order and leg structure, so ``last divergent'' is intrinsic to each pair, and the grouping is a partition. In each group the inner sum runs over the full class of $H$, so \cite[(3.5)]{Ba3} applies to $H$. It consists of Taylor's formula with remainder and the telescoping resummation of the indices of $H$, so the total is unchanged term by term.
\end{proof}

\begin{theorem}[nested and one-particle-reducible classes, with a pin]\label{app:K:L09-6}\label{app:G1:H3-nest}
\begin{itemize}
\item[(a)] For every renormalised class of \cite{Ba3} whose members contain a divergent proper subgraph, the decomposition \cite[(3.5)]{Ba3} holds, and hence so do \cite[(3.2)]{Ba3} and \cite[(1.33)]{Ba3}.
\item[(b)] At fixed labels, the one-particle-reducible members of a class of two-scalar-leg graphs of order $2$ sum to $K=S^{(1)}_1GS^{(1)}_2$. Here $S^{(1)}_i$ are amputated order-one one-particle-irreducible two-point operators (including $\delta m^2_1$) and $G$ is the bridging scalar line. \cite[(3.19)]{Ba3} applies to $K$ with zero counterterm and leaves a term of degree $\alpha$ and a term of degree $\ge1$.
\item[(c)] If every member carries leg $1$ with $\Phi_{\rm ext}\equiv1$ at a fixed pin $p$, then the transformations of (a) and (b) never sum over the position of $p$.
\end{itemize}
\end{theorem}
\begin{proof}
(b) A one-particle-reducible two-scalar-leg graph of order $2$ is separated by a scalar bridge into two two-leg pieces of order $1$; a photon bridge would give $\mathsf C6$. The labels factor over the pieces, so the sum is $K=S^{(1)}_1GS^{(1)}_2$.

Split $S=\Sigma_{nn}+\Sigma_{nd}\partial+\partial^*\Sigma_{dn}+\partial^*\Sigma_{dd}\partial$ according to where the derivatives sit. The counterterm is minus the zeroth moment of the $nn$-part \cite[(1.31), (3.9)]{Ba3}. Tadpoles cancel exactly against their counterterms \cite[p.~435]{Ba3}, and a derivative at $x'$ gives $\sum_{x'}(\partial f)(x')=0$. Hence $(S-{\rm ct})1=\partial^*c$, and
\[
\sum_{x'}K(x,x')=\sum_{y,y'}S^{(1)}_1(x,y)G(y,y')(\partial^*_\nu c_\nu)(y'),
\]
where $c_\nu(y')$ is the zeroth moment of the part of $S^{(1)}_2$ carrying a derivative on the leg on the side of the bridge.

After the replacement of \cite[p.~437]{Ba3} the kernel is translation invariant up to a regular remainder. So $c(y')=g(y')\sum_{x'}K_0(y'-x')g'(x')$ plus regular terms, and $\partial c$ is a unit-scale derivative of products of the cut-off functions, with $\abs{\partial c}\le O(1)\norm g_{C^2}\norm{g'}_{C^2}$ by \cite[(3.24)]{Ba3}. The difference between the torus and the $\eta\Z^d$ counterterms is an image term, booked as a finite vertex (1.7).

\cite[(3.19)]{Ba3} is an identity for every kernel, applied here with zero counterterm. The order-one difference on the leg costs $(L^{j_\ell}\eta)^{-1}$. Replacing it by a unit-scale derivative of the cut-off functions raises the degree by $1$, which gives the term of degree $\ge1$; the other term has degree $\alpha$.

(a) Induct on the order. Lemmas~\ref{app:G1:H3-N1}--\ref{app:G1:H3-N3} give the inner step, and disjoint inner subgraphs are treated one after another. The outer step applies the primitive formula to the transformed whole graph, all of whose proper subgraphs now have positive degree:
\begin{itemize}
\item For $\mathsf C2$, the identity
\[
\phi\Sigma\phi'-\phi(\textstyle\sum_{x'}\Sigma)\phi=\phi\Sigma\abs{x'-x}^\alpha[\phi'(x')-\phi'(x)]\abs{x'-x}^{-\alpha}
\]
holds for every kernel. The counterterm is fixed by the untransformed class (Lemma~\ref{app:G1:H3-N3}), so the result has degree $\alpha$. The one-particle-reducible members are treated by (b).
\item For $\mathsf C4$, the zeroth moment equals the untransformed one, and its part with the lines $C^\xi$ vanishes by Lemma~\ref{app:G1:AA4-WT}. A term containing a regular line $(1-m^2-aP)C^\xi$ is convergent by power counting, since such a line raises the degree of every subgraph containing it by at least $2$. The chain of two polarisation insertions joined by a photon line factorises through the order-$e^2$ identity.
\item For $\mathsf C5$, the zeroth moment of ``self-energy on a leg plus counterterm (2.21g)'' factorises through the zeroth moment of the self-energy and vanishes \cite[p.~444]{Ba3}. The members (2.21d) and (2.21f) cancel exactly against (2.21g).
\end{itemize}
\cite[Props.~2.1--2.2]{Ba3} bound the result. This proves \cite[(3.2)]{Ba3}, and \cite[(1.33)]{Ba3} follows by Theorem~\ref{app:K:L09-3}(c).

(c) There are three cases.
\begin{itemize}
\item Inner $H\not\ni p$: the sums of (T1) and (T2) run over vertices of $H$. The transported leg is a line of $G_0$ (Lemma~\ref{app:G1:H3-N1}), and its far end, possibly $p$, stays fixed.
\item Inner $H\ni p$: leg $1$ is external to $H$ and has the highest index, so the base point is $p$, by steps (iii)--(v) of Lemma~\ref{app:G1:VSR-C}.
\item Reducible chains through $p$: the identity of (b) is used with $x=p$ fixed.
\end{itemize}
Inside the one-particle-irreducible counterterm graphs, a proper self-energy through $p$ would create a bridge. So only local insertions at $p$ occur, and these cancel exactly. This is the statement used in step (vi) of Lemma~\ref{app:G1:VSR-C}.
\end{proof}

\paragraph{The nearest-neighbour bound.}
In this paragraph lengths are measured in the units of the block lattice of $\psi$. The blocks $B(y)$, $y\in\Z^d$, are unit cubes of fine sites, and the step lattice has spacing $1/L$. For $x$ in the block $B(y)$, $\Gamma_{y,x}\subset B(y)$ is the axis-ordered contour on the step lattice \cite[(2.1)]{Ba1}; it has at most $d(L-1)$ bonds and length $\le d$.

An oriented fine bond is $[z,z+\eta e_\mu]$, with $[z+\eta e_\mu,z]:=-[z,z+\eta e_\mu]$. A contour is identified with the integral $1$-chain of its bonds, and $\ell(\Gamma)$ is its length. Fluxes $A(\Gamma)$ are extended linearly to chains. A $2$-chain $S=\sum_Pn_PP$ of plaquettes has boundary $\partial S=\sum_Pn_P\partial P$ and area $\abs S:=\eta^2\sum_P\abs{n_P}$. A box is a set $\{z:a_\mu\le z_\mu\le b_\mu\}$ that does not wrap around the torus.

The transport is $U(a)=\exp(e(t)qa)$, with $q$ real antisymmetric and $\norm q\le1$. The matrices $U(a)$ are orthogonal, $U(a)U(a')=U(a+a')$, and
\[
\norm{U(a)-U(a')}=\norm{U(a-a')-1}\le e(t)\abs{a-a'} ,
\]
since $qa$ is normal with eigenvalues $i\theta$, $\abs\theta\le\abs a$. Since $\abs{U(a)v-w}=\abs{v-U(-a)w}$, the statements below hold verbatim in the opposite transport convention.

\begin{lemma}[lattice Stokes formula with area bound]\label{app:G1:N-stokes}
Let $c$ be a closed contour in a box $Q$, let $o$ be a vertex of $c$, and put $R(c):=\max_{z\in c}\abs{z-o}_1$. Then there is a $2$-chain $S$ in $Q$ with $\partial S=c$ and $\abs S\le\ell(c)R(c)$. If $\abs{\partial^\eta_\mu A_\nu}\le K_0$ on $Q$ for all $\mu,\nu$, then $\abs{A(c)}\le2K_0\ell(c)R(c)$.
\end{lemma}
\begin{proof}
For $z\in Q$ let $\Pi(z)$ be the axis-ordered path from $o$ to $z$: first along $e_1$, then $e_2$, then $e_3$. It lies in the box spanned by $o$ and $z$.

Let $b=[z,z']$ with $z'=z+s\eta e_\mu$, $s=\pm1$, and put $w:=(z_1,\dots,z_\mu,o_{\mu+1},\dots,o_d)$. The paths $\Pi(z)$ and $\Pi(z')$ coincide up to their $e_\mu$-legs, and as chains these legs differ by the bond $[w,w+s\eta e_\mu]$. The remaining parts are the path $\Pi_{>\mu}(z)$ from $w$ to $z$, which uses only directions $\nu>\mu$, and its translate by $s\eta e_\mu$. Hence
\[
\Pi(z)+b-\Pi(z')=\Pi_{>\mu}(z)+[z,z']-[w,w+s\eta e_\mu]-\tau_{s\eta e_\mu}\Pi_{>\mu}(z).
\]
This is the boundary of the ribbon $S(b)$ of the plaquettes spanned by the bonds of $\Pi_{>\mu}(z)$ and $s\eta e_\mu$, each with the sign of its bond: the inner rungs cancel. The ribbon lies in the box spanned by $o,z,z'$ and has at most $\eta^{-1}\abs{z-o}_1$ plaquettes, so $\abs{S(b)}\le\eta R(c)$.

For $c=\sum_{i<n}[z_i,z_{i+1}]$ with $z_n=z_0$, put $S:=\sum_iS([z_i,z_{i+1}])$. Then $\partial S=\sum_i(\Pi(z_i)+[z_i,z_{i+1}]-\Pi(z_{i+1}))=c$ and $\abs S\le n\eta R(c)=\ell(c)R(c)$. Finally,
\[
A(\partial P_{\mu\nu}(z))=\eta^2[\partial^\eta_\mu A_\nu(z)-\partial^\eta_\nu A_\mu(z)],
\]
so $\abs{A(c)}\le\sum_P\abs{n_P}\abs{A(\partial P)}\le2K_0\abs S$.
\end{proof}

\begin{lemma}[change of contour and of field in the transport]\label{app:G1:N-field}
\begin{itemize}
\item[(i)] Let $\Gamma,\Gamma'$ be contours from $y_1$ to $y_2$ in a box $Q$ on which $\abs{\partial^\eta_\mu\tilde B_\nu}\le C_\partial p$, and put $c:=\Gamma'-\Gamma$, $o:=y_1$. Then $\norm{U(\tilde B(\Gamma'))-U(\tilde B(\Gamma))}\le2e(t)C_\partial p\,\ell(c)R(c)$.
\item[(ii)] Let $\Gamma$ consist of $n$ bonds of the step lattice on which $A'':=\mathcal A-\tilde B$ satisfies $\abs{A''(\langle x,x'\rangle)}\le c''p$. Then $\norm{U(\mathcal A(\Gamma))-U(\tilde B(\Gamma))}\le e(t)nc''p$.
\end{itemize}
\end{lemma}
\begin{proof}
(i) The difference is $e(t)\abs{\tilde B(c)}$ at most, and Lemma~\ref{app:G1:N-stokes} bounds $\abs{\tilde B(c)}$. (ii) Fluxes are linear in the field, so $\mathcal A(\Gamma)-\tilde B(\Gamma)=A''(\Gamma)$, and $\abs{A''(\Gamma)}\le nc''p$.
\end{proof}

\begin{proposition}[nearest-neighbour bound with transport by the background]\label{app:K:L09-2}\label{app:G1:N-nn}
Let $\Omega'$ be a union of blocks, $\mathcal A$ the fine field of the step \cite[(2.54)]{Ba2}, $\tilde B$ the background, and $p=p(t)$. Assume the following:
\begin{itemize}
\item[(NN1)] (\cite[(2.16)]{Ba2} at scale $t$.) For blocks in $\Omega'$, $\abs{\psi(y)}\le2\lambda(t)^{-1/4}p$. For nearest neighbours $x\in B(y_1)$, $x'\in B(y_2)$ of the step lattice with $y_1\ne y_2$,
\[
\abs{U(\mathcal A(\Gamma_{y_1,x,x',y_2}))\psi(y_2)-\psi(y_1)}\le4Ld\,p,\qquad \Gamma_{y_1,x,x',y_2}:=\Gamma_{y_1,x}+\langle x,x'\rangle-\Gamma_{y_2,x'}.
\]
\item[(NN2)] (Lemma~\ref{app:K:L09-1}, when $t=t_k$ with $0\le k\le K-1$, $\tilde B$ is the background of that lemma, its hypotheses hold, and $\Omega'\subset B^k(\Lambda^{(k)}_0)$.) $A'':=\mathcal A-\tilde B$ satisfies $\abs{A''(\langle x,x'\rangle)}\le c''p$ for every bond of the step lattice in $\Omega'$.
\item[(NN3)] $\abs{\partial^\eta_\mu\tilde B_\nu}\le C_\partial p$ on $\Omega'$ for all $\mu,\nu$.
\end{itemize}
Assume also $\ell_T/t>2$, and put $\kappa(t):=\abs e\lambda^{-1/4}t^{1/4}p(t)$. Then for nearest neighbours $y_1,y_2$ of the unit lattice with $B(y_1),B(y_2)\subset\Omega'$, and $\Gamma_{y_1,y_2}$ the straight contour,
\[
\abs{U(\tilde B(\Gamma_{y_1,y_2}))\psi(y_2)-\psi(y_1)}\le c_\psi(t)p(t),\qquad c_\psi(t):=4Ld+\big(2(2dL-2d+1)c''+8(d+1)^2C_\partial\big)\kappa(t).
\]
\end{proposition}
\begin{proof}
\emph{The box.} Since $B(y_1),B(y_2)$ are adjacent unit cubes of fine sites, $Q:=B(y_1)\cup B(y_2)$ is a box. It does not wrap since $\ell_T/t>2$, and $Q\subset\Omega'$. Choose nearest neighbours $x\in B(y_1)$, $x'\in B(y_2)$ of the step lattice, and put $\Gamma':=\Gamma_{y_1,x,x',y_2}$. The contours $\Gamma_{y_1,y_2}$ and $\Gamma'$ run from $y_1$ to $y_2$ and lie in $Q$.

\emph{Change of field.} $\Gamma'$ consists of at most $2d(L-1)+1$ bonds of the step lattice. By (NN1), (NN2) and Lemma~\ref{app:G1:N-field}(ii),
\[
\abs{U(\tilde B(\Gamma'))\psi(y_2)-\psi(y_1)}\le4Ldp+e(t)(2dL-2d+1)c''p\cdot2\lambda(t)^{-1/4}p=C_0p,
\]
with $C_0:=4Ld+2(2dL-2d+1)c''\kappa(t)$, since $e(t)\lambda(t)^{-1/4}=\abs e\lambda^{-1/4}t^{1/4}$.

\emph{Geometry of the closed contour.} Let $c:=\Gamma'-\Gamma_{y_1,y_2}$ and $o:=y_1$, and write $\ell_Q:=d$, which bounds the lengths of $\Gamma_{y_1,x}$ and $\Gamma_{y_2,x'}$. Every point $z$ of $c$ satisfies one of the following:
\begin{itemize}
\item $z\in\Gamma_{y_1,x}$, and then $\abs{z-y_1}_1\le\ell_Q$;
\item $z\in\langle x,x'\rangle\cup\Gamma_{y_2,x'}$, and then $\abs{z-y_1}_1\le\abs{z-y_2}_1+1\le\ell_Q+1$;
\item $z\in\Gamma_{y_1,y_2}$, and then $\abs{z-y_1}_1\le1$.
\end{itemize}
Hence $R(c)\le\ell_Q+1$ and $\ell(c)\le2\ell_Q+2$, so $\ell(c)R(c)\le2(d+1)^2$.

\emph{Conclusion.} By Lemma~\ref{app:G1:N-field}(i) and (NN3), $\norm{U(\tilde B(\Gamma_{y_1,y_2}))-U(\tilde B(\Gamma'))}\le4e(t)C_\partial(d+1)^2p$. With $\abs{\psi(y_2)}\le2\lambda(t)^{-1/4}p$,
\[
\abs{U(\tilde B(\Gamma_{y_1,y_2}))\psi(y_2)-\psi(y_1)}\le C_0p+8(d+1)^2C_\partial\kappa(t)p=c_\psi(t)p .
\]
\end{proof}

\subsection{Kernel and coefficient bounds for the mass reading}\label{app:G1:I7-kernels}

\emph{Setting.} Let $a=1$, $L=2$, $a_\flat=a(1-L^{-2})$, $j_0\ge1$ and $\xi=L^{-j_0}\le\frac12$. Let $\Lambda=\xi\Z^3/\ell\Z^3$ with $\ell\in L\N$ and $\ell\ge L$. Let $d$ be the $\ell^1$ torus distance and $\rho(w,w'):=d(w,w')+\xi$. Kernels are taken with respect to $\xi^3\sum$. $P_1$ is the orthogonal projection onto functions that are constant on the unit blocks. Since every unit block has volume $1$, $(P_1u)(z)=\xi^3\sum_{z'\in B(z)}u(z')$, where $B(z)$ is the unit block containing $z$; any two points of a block satisfy $d(z,z')\le3$.

Put $F:=(-\Delta^\xi+1)^{-1}$. Its kernel is $F(x,z)=f(x-z)$ with $f$ even. Let $\mathcal K$ be the set of operators $(-\Delta^\xi+\nu+bP_1)^{-1}$ with $\nu\in[0,1]$ and $b\in[a_\flat,a]$. For $K\in\mathcal K$ put $E_K:=K-F$. The operators $K$ and $F$ are real and self-adjoint, so the kernels $K$, $F$ and $E_K$ are symmetric.

The forward difference in direction $e_\mu$ is $(g(w+\xi e_\mu)-g(w))/\xi$, and the backward difference is $(g(w)-g(w-\xi e_\mu))/\xi$. A single difference is either of these, in any direction. $D^{(1)}$ and $D^{(2)}$ denote a single difference acting in the first and in the second variable. Put $c_*:=4\mathrm{e}^{1/8}$.

We use the following inputs.
\begin{itemize}
\item[(F)] By Lemmas~\ref{app:G1:VSR-F} and~\ref{app:G1:VSR-IM}, $0\le F(w,w')\le A_0^T\rho^{-1}\mathrm{e}^{-d/8}$. For a forward difference $D$ in the first variable, $\abs{D^{(1)}F(w,w')}\le A_1^T\rho^{-2}\mathrm{e}^{-d/8}$.
\item[(E)] By Lemma~\ref{app:G1:VSR-E}, for every $K\in\mathcal K$ we have $\abs{E_K(w,w')}\le A_E\mathrm{e}^{-\kappa_Gd}$. For a forward difference in the first variable, $\abs{D^{(1)}E_K(w,w')}\le A_\nabla\mathrm{e}^{-\kappa_Gd}$.
\end{itemize}
Here $d=d(w,w')$ and $\rho=\rho(w,w')$. The rate $\kappa_G\in(0,\frac1{16}]$ and the constants $A_0^T,A_1^T,A_E,A_\nabla$ depend only on $(a,L)$. They are uniform in $\nu\in[0,1]$, $b\in[a_\flat,a]$, $\xi$, $j_0$ and $\ell$.

For $\phi>0$ and $n\in\{0,1,2\}$ let
\[
J_n(\phi):=\xi^3\sum_{v\in\xi\Z^3}\mathrm{e}^{-\phi\abs v_1}(\abs v_1+\xi)^{-n}
\]
be the lattice sums of Lemma~\ref{app:G1:VSR-Jsum}. Each $J_n$ is decreasing in $\phi$. For $m\ge1$ the set $\{v\in\xi\Z^3:\abs v_1=m\xi\}$ has $4m^2+2$ points. With $u:=\phi\xi$, and using $4m^2+2\le4m(m+1)\le4(m+1)^2$, $\coth x\le1+1/x$, $\sum_{m\ge1}m\mathrm{e}^{-um}\le u^{-2}$ and $(\mathrm{e}^u-1)^{-1}\le u^{-1}$, we get
\[
J_0(\phi)\le(1+2/\phi)^3,\qquad J_1(\phi)\le1+4/\phi^2,\qquad J_2(\phi)\le1+4/\phi .
\]
In particular these bounds are uniform in $\xi\le1$.

\begin{lemma}[elementary kernel bounds]\label{app:G1:I7-K}
Let $K\in\mathcal K$, $E:=E_K$, and let $D$ be a single difference in either variable. Put $B_*:=A_1^T+A_\nabla\sup_{s\ge0}(1+s)^2\mathrm{e}^{-\kappa_Gs/2}$. Then for all $w,w',x,y\in\Lambda$:
\begin{itemize}
\item[(i)] $\abs{K(w,w')}\le(A_0^T\rho^{-1}+A_E)\mathrm{e}^{-\kappa_Gd}$ and $\abs{E(w,w')}\le A_E\mathrm{e}^{-\kappa_Gd}$.
\item[(ii)] $\abs{DF(w,w')}\le c_*A_1^T\rho^{-2}\mathrm{e}^{-\kappa_Gd}$ and $\abs{DE(w,w')}\le c_*A_\nabla\mathrm{e}^{-\kappa_Gd}$. Hence
\[
\abs{DK(w,w')}\le c_*(A_1^T\rho^{-2}+A_\nabla)\mathrm{e}^{-\kappa_Gd}\le c_*B_*\rho^{-2}\mathrm{e}^{-\kappa_Gd/2}.
\]
\item[(iii)] For every $g\ge0$ on $[0,\infty)$, $\xi^3\sum_{z\in\Lambda}g(d(z,y))\le\xi^3\sum_{v\in\xi\Z^3}g(\abs v_1)$.
\item[(iv)] We have
\[
\xi^3\sum_{d(z,y)\le R}\rho(z,y)^{-2}\le4(R+\xi),\qquad \xi^3\sum_{d(z,y)\ge R}\rho(z,y)^{-4}\le\frac8{R+\xi}\qquad(R\ge0),
\]
and
\[
\xi^3\sum_z\rho(x,z)^{-2}\rho(z,y)^{-2}\le176\,\rho(x,y)^{-1}.
\]
\end{itemize}
\end{lemma}

\begin{proof}
(i) follows from (F) and (E), since $\kappa_G\le\frac1{16}\le\frac18$.

(ii) For forward differences in the first variable this is (F) and (E), with constant $1\le c_*$. For the second variable we use the symmetry of $F$ and $E$. A backward difference at $w$ is a forward difference at $w-\xi e_\mu$. The shifted distance $d'$ satisfies $d'\ge d-\xi$. Hence $d'+\xi\ge\max(d,\xi)\ge\rho/2$, and $\mathrm{e}^{-\kappa'd'}\le\mathrm{e}^{\kappa'\xi}\mathrm{e}^{-\kappa'd}\le\mathrm{e}^{1/8}\mathrm{e}^{-\kappa'd}$ for $\kappa'\le\frac18$. Together with $\rho'^{-2}\le4\rho^{-2}$ this gives the factor $c_*$. The bound on $DK=DF+DE$ follows. The last inequality uses $\rho\le1+d$, so that $A_\nabla\mathrm{e}^{-\kappa_Gd}\le A_\nabla\rho^{-2}(1+d)^2\mathrm{e}^{-\kappa_Gd/2}\,\mathrm{e}^{-\kappa_Gd/2}$.

(iii) Send each $z$ to a representative $v\in\xi\Z^3$ of $z-y$ with $\abs v_1=d(z,y)$. This map is injective.

(iv) We use (iii) and the count of points on the $\ell^1$-spheres. Each sphere of radius $m\xi$, $m\ge1$, contributes at most $\xi^3(4m^2+2)\xi^{-2}(m+1)^{-2}\le4\xi$ to the first sum. There are at most $R/\xi$ such spheres, and the centre contributes $\xi$; this gives the first bound.

In the second sum, each sphere with $m\ge1$ contributes at most $4\xi^{-1}(m+1)^{-2}$.
\begin{itemize}
\item If $R\ge\xi$, the sum over $m\ge R/\xi$ is at most $4/R\le8/(R+\xi)$.
\item If $R<\xi$, the whole sum, centre included, is at most $\xi^{-1}(1+4(\pi^2/6-1))\le4/\xi\le8/(R+\xi)$.
\end{itemize}

For the third bound put $D_*:=d(x,y)$.
\begin{itemize}
\item On $\{d(z,y)\le D_*/2\}$ we have $\rho(x,z)\ge D_*/2+\xi\ge\rho(x,y)/2$. By the first bound with $R=D_*/2$, the contribution is at most $4\rho(x,y)^{-2}\cdot4\rho(x,y)=16\rho(x,y)^{-1}$.
\item The same holds on $\{d(x,z)\le D_*/2\}$.
\item On the remaining set, $d(z,y)\ge\max(d(x,z)-D_*,D_*/2)\ge d(x,z)/3$. So $\rho(z,y)\ge\rho(x,z)/3$, and the summand is at most $9\rho(x,z)^{-4}$. By the second bound with $R=D_*/2$, the contribution is at most $72/(D_*/2+\xi)\le144\rho(x,y)^{-1}$.
\end{itemize}
Adding the three contributions gives $176\rho(x,y)^{-1}$.
\end{proof}

\begin{lemma}[mixed second differences]\label{app:G1:I7-M}
Let $K\in\mathcal K$ and let $D,D'$ be single differences. Put $\sigma_\kappa:=\sup_{s\ge0}(s+4)\mathrm{e}^{-\kappa_Gs/4}$ and $\sigma'_\kappa:=\sup_{s\ge0}(1+s)\mathrm{e}^{-\kappa_Gs/8}$. Then for all $x,y\in\Lambda$
\[
\abs{D^{(1)}D'^{(2)}E_K(x,y)}\le A_M\,\rho(x,y)^{-1}\mathrm{e}^{-\kappa_Gd(x,y)/8},
\]
where
\[
A_M:=c_*^2A_1^TB_*\big[176+4a\,\mathrm{e}^{3\kappa_G/2}\sigma_\kappa\sigma'_\kappa J_2(3\kappa_G/4)\big].
\]
\end{lemma}

\begin{proof}
Write $K=(-\Delta^\xi+\nu+bP_1)^{-1}$ and $W:=(1-\nu)-bP_1$. Then $K^{-1}=F^{-1}-W$, so $E_K=F(F^{-1}-K^{-1})K=FWK$. In this product $D^{(1)}$ acts on $F$ and $D'^{(2)}$ acts on $K$. With $u_y:=(D'^{(2)}K)(\cdot,y)$ this gives
\[
D^{(1)}D'^{(2)}E_K(x,y)=(1-\nu)\,\xi^3\sum_z(D^{(1)}F)(x,z)u_y(z)-b\,\xi^3\sum_z(D^{(1)}F)(x,z)(P_1u_y)(z).
\]
By Lemma~\ref{app:G1:I7-K}(ii), $\abs{D^{(1)}F(x,z)}\le c_*A_1^T\rho(x,z)^{-2}\mathrm{e}^{-\kappa_Gd(x,z)}$ and $\abs{u_y(z')}\le c_*B_*\rho(z',y)^{-2}\mathrm{e}^{-\kappa_Gd(z',y)/2}$.

\emph{First term.} We have $\mathrm{e}^{-\kappa_Gd(x,z)-\kappa_Gd(z,y)/2}\le\mathrm{e}^{-\kappa_Gd(x,y)/2}$. By Lemma~\ref{app:G1:I7-K}(iv) the first term is at most $176c_*^2A_1^TB_*\rho(x,y)^{-1}\mathrm{e}^{-\kappa_Gd(x,y)/2}$.

\emph{Second term.} Put $D_0:=d(z,y)$. Every $z'\in B(z)$ satisfies $d(z,z')\le3$. Hence $\mathrm{e}^{-\kappa_Gd(z',y)/2}\le\mathrm{e}^{3\kappa_G/2}\mathrm{e}^{-\kappa_GD_0/2}$ and $B(z)\subset\{d(z',y)\le D_0+3\}$. By Lemma~\ref{app:G1:I7-K}(iv),
\[
\abs{(P_1u_y)(z)}\le4c_*B_*\mathrm{e}^{3\kappa_G/2}(D_0+4)\mathrm{e}^{-\kappa_GD_0/2}\le4c_*B_*\mathrm{e}^{3\kappa_G/2}\sigma_\kappa\mathrm{e}^{-\kappa_GD_0/4}.
\]
We have $\mathrm{e}^{-\kappa_Gd(x,z)-\kappa_Gd(z,y)/4}\le\mathrm{e}^{-\kappa_Gd(x,y)/4}\mathrm{e}^{-3\kappa_Gd(x,z)/4}$. By Lemma~\ref{app:G1:I7-K}(iii) and $b\le a$, the second term is at most $4a\,c_*^2A_1^TB_*\mathrm{e}^{3\kappa_G/2}\sigma_\kappa J_2(3\kappa_G/4)\,\mathrm{e}^{-\kappa_Gd(x,y)/4}$.

\emph{Conclusion.} Since $\rho\le1+d$, we have $\mathrm{e}^{-\kappa_Gd/4}\le\sigma'_\kappa\rho^{-1}\mathrm{e}^{-\kappa_Gd/8}$. Also $\mathrm{e}^{-\kappa_Gd/2}\le\mathrm{e}^{-\kappa_Gd/8}$ and $\abs{1-\nu}\le1$. Adding the two terms gives the claim.
\end{proof}

\begin{lemma}[coefficient bounds]\label{app:G1:I7-C}
Let $K_1,K_2,K_3\in\mathcal K$, let $D,D_1,D_2,D_3$ be single differences, fix $y\in\Lambda$, and put $\kappa:=\kappa_G$. Then:
\begin{itemize}
\item[(C1)] $\xi\abs{K_1(y,y)}\le A_0^T+A_E$.
\item[(C2)] $\big|\xi^3\sum_{y'}[(DK_1)(y,y')K_2(y,y')-(DF)(y,y')F(y,y')]\big|\le K_{\rm C}$, where
\[
K_{\rm C}:=c_*[A_\nabla(A_0^TJ_1(\kappa)+A_EJ_0(\kappa))+A_1^TA_EJ_2(\kappa)].
\]
\item[(C3)] If $\abs{h(y')}\le\norm h\,d(y,y')$ for all $y'$, then $\big|\xi^3\sum_{y'}(DK_1)(y,y')K_2(y,y')h(y')\big|\le K_{\rm C}'\norm h$, where
\[
K_{\rm C}':=c_*(A_1^T+A_\nabla)(A_0^T+A_E)[2J_2(\kappa)+(3+1/(\mathrm{e}\kappa))J_0(\kappa)].
\]
\item[(C4)] Put $T[K_1,K_2,K_3]:=\xi^6\sum_{y',y''}(D_1K_1)(y',y)(D_2K_2)(y,y'')(D_3K_3)(y',y'')$. Then $\abs{T[K_1,K_2,K_3]-T[F,F,F]}\le K_\triangle$, where
\[
K_\triangle:=c_*^3A_\nabla\big[176B_*^2J_1(\kappa)+A_1^TB_*(8+J_0(\kappa))J_2(\kappa)+(A_1^T)^2J_2(\kappa)^2\big].
\]
\end{itemize}
\end{lemma}

\begin{proof}
All sums are reduced to sums over $\xi\Z^3$ by Lemma~\ref{app:G1:I7-K}(iii). Write $E_i:=K_i-F$.

(C1) is Lemma~\ref{app:G1:I7-K}(i) at $d=0$, where $\rho=\xi$, together with $\xi\le1$.

(C2) We have $(DK_1)K_2-(DF)F=(DE_1)K_2+(DF)E_2$. By Lemma~\ref{app:G1:I7-K}(i),(ii), the first product is at most $c_*A_\nabla(A_0^T\rho^{-1}+A_E)\mathrm{e}^{-2\kappa d}$ and the second at most $c_*A_1^TA_E\rho^{-2}\mathrm{e}^{-2\kappa d}$. Summing, and using $J_n(2\kappa)\le J_n(\kappa)$, gives $K_{\rm C}$.

(C3) By Lemma~\ref{app:G1:I7-K}(i),(ii), $\abs{DK_1}\le c_*(A_1^T+A_\nabla)(\rho^{-2}+1)\mathrm{e}^{-\kappa d}$ and $\abs{K_2}\le(A_0^T+A_E)(\rho^{-1}+1)\mathrm{e}^{-\kappa d}$. Also $\abs h\le\norm h\rho$. So the summand is at most
\[
c_*(A_1^T+A_\nabla)(A_0^T+A_E)(\rho^{-2}+1)(1+\rho)\mathrm{e}^{-2\kappa d}\norm h .
\]
Since $\rho^{-1}\le\rho^{-2}+1$ and $\rho\le1+d$, we have $(\rho^{-2}+1)(1+\rho)\le2\rho^{-2}+3+d$. Moreover $d\mathrm{e}^{-2\kappa d}\le(\mathrm{e}\kappa)^{-1}\mathrm{e}^{-\kappa d}$. Summing gives $K_{\rm C}'$.

(C4) By multilinearity,
\[
T[K_1,K_2,K_3]-T[F,F,F]=T[E_1,K_2,K_3]+T[F,E_2,K_3]+T[F,F,E_3].
\]
\begin{itemize}
\item In the first term, $\abs{D_1E_1(y',y)}\le c_*A_\nabla\mathrm{e}^{-\kappa d(y,y')}$. The other two factors are at most $c_*B_*\rho^{-2}$ each. By Lemma~\ref{app:G1:I7-K}(iv) the sum over $y''$ is at most $176c_*^2B_*^2\rho(y,y')^{-1}$, and the sum over $y'$ then gives $J_1(\kappa)$.
\item In the second term, split the sum over $y''$ at $d(y',y'')=1$ and use Lemma~\ref{app:G1:I7-K}(iv) on the inner part and $\rho\ge1$ on the outer part. This gives $\xi^3\sum_{y''}\mathrm{e}^{-\kappa d(y,y'')}\rho(y',y'')^{-2}\le8+J_0(\kappa)$. Together with $\abs{D_2E_2}\le c_*A_\nabla\mathrm{e}^{-\kappa d(y,y'')}$, $\abs{D_3K_3}\le c_*B_*\rho(y',y'')^{-2}$ and $\abs{D_1F(y',y)}\le c_*A_1^T\rho^{-2}\mathrm{e}^{-\kappa d}$, this gives $c_*^3A_\nabla A_1^TB_*(8+J_0(\kappa))J_2(\kappa)$.
\item In the third term, $\abs{D_3E_3}\le c_*A_\nabla$. The two remaining sums factorise, and each is at most $c_*A_1^TJ_2(\kappa)$.\qedhere
\end{itemize}
\end{proof}

\begin{lemma}[the doubly differenced coefficient]\label{app:G1:I7-C21c}
Let $K_1,K_2,K_3\in\mathcal K$, let $D_1,D_2,D_3$ be single differences, fix $y'\in\Lambda$, and put
\[
T^{\rm c}[K_1,K_2,K_3]:=\xi^6\sum_{y,y''}(D_1^{(1)}D_2^{(2)}K_1)(y',y)\,(D_3^{(2)}K_2)(y,y'')\,K_3(y,y'').
\]
Then $T^{\rm c}[F,F,F]=0$ and $\abs{T^{\rm c}[K_1,K_2,K_3]}\le K_{\rm c}$, where, with $\kappa:=\kappa_G$,
\[
K_{\rm c}:=A_MQ_\infty J_1(\kappa/8)+c_*A_1^TJ_2(\kappa)(Q_1+Q_2),\qquad Q_\infty:=4c_*^2(A_1^T)^2+K_{\rm C},
\]
\[
Q_1:=\mathrm{e}^{1/16}\big[A_M\big(2A_0^TJ_2(\kappa/8)+A_EJ_1(\kappa/8)\big)+c_*^2A_\nabla B_*J_2(\kappa/2)\big],\qquad Q_2:=2c_*^2A_\nabla A_1^TJ_2(\kappa).
\]
\end{lemma}

\begin{proof}
All sums are finite. Write $E_i:=K_i-F$ and $g(y):=\xi^3\sum_{y''}(D_3^{(2)}K_2)(y,y'')K_3(y,y'')$, so that
\[
T^{\rm c}=\xi^3\sum_y(D_1^{(1)}D_2^{(2)}K_1)(y',y)g(y).
\]
Decompose $g=g_F+g_1+g_2$ with
\[
g_F:=\xi^3\sum_{y''}(D_3^{(2)}F)F,\qquad g_1:=\xi^3\sum_{y''}(D_3^{(2)}E_2)K_3,\qquad g_2:=\xi^3\sum_{y''}(D_3^{(2)}F)E_3 .
\]

\emph{Step 1: $g_F$ is a constant with $\abs{g_F}\le4c_*^2(A_1^T)^2$.} Let $D_3$ act in direction $e_\mu$. With $w=y-y''$ we have $(D_3^{(2)}F)(y,y'')=\sigma\xi^{-1}[f(w+\tau\xi e_\mu)-f(w)]$ for some $\sigma,\tau\in\{\pm1\}$. Hence $g_F$ does not depend on $y$. Since $\sum_wf(w+\tau\xi e_\mu)^2=\sum_wf(w)^2$, the identity $(p-q)q=-\frac12(p-q)^2+\frac12(p^2-q^2)$ gives
\[
g_F=-\frac{\sigma}{2\xi}\,\xi^3\sum_w\big[f(w+\tau\xi e_\mu)-f(w)\big]^2,\qquad\text{so}\qquad\abs{g_F}\le\frac\xi2\,\xi^3\sum_w\big(c_*A_1^T\rho(w,0)^{-2}\big)^2
\]
by Lemma~\ref{app:G1:I7-K}(ii). By Lemma~\ref{app:G1:I7-K}(iv) with $R=0$, $\xi^3\sum_w\rho(w,0)^{-4}\le8/\xi$.

\emph{Step 2: $\sup\abs g\le Q_\infty$.} This is Step 1 together with (C2) of Lemma~\ref{app:G1:I7-C}, applied to $K_2,K_3$ and $D=D_3^{(2)}$.

\emph{Step 3: the $E_1$-part.} By Lemma~\ref{app:G1:I7-M} and Step 2, $T_E:=\xi^3\sum_y(D_1^{(1)}D_2^{(2)}E_1)(y',y)g(y)$ satisfies $\abs{T_E}\le A_MQ_\infty J_1(\kappa/8)$.

\emph{Step 4: the $F$-part.} Put $\phi(y):=(D_1^{(1)}F)(y',y)$, so that $\abs\phi\le c_*A_1^T\rho(y',y)^{-2}\mathrm{e}^{-\kappa d(y',y)}$ and $(D_1^{(1)}D_2^{(2)}F)(y',y)=(D_2\phi)(y)$. Let $\bar D_2$ be the backward difference if $D_2$ is forward, and vice versa, in the same direction. On the torus $\xi^3\sum_y(D_2\phi)v=-\xi^3\sum_y\phi\,(\bar D_2v)$ for every $v$. Since $\bar D_2g_F=0$,
\[
T_F:=\xi^3\sum_y(D_2\phi)(y)g(y)=-\xi^3\sum_y\phi(y)\big(\bar D_2g_1+\bar D_2g_2\big)(y).
\]
We show $\abs{\bar D_2g_1}\le Q_1$ and $\abs{\bar D_2g_2}\le Q_2$.
\begin{itemize}
\item[(a)] Keep $y''$ fixed and use the discrete product rule $\bar D(uv)(y)=(\bar Du)(y)v(\hat y)+u(\check y)(\bar Dv)(y)$ with $\hat y,\check y\in\{y,y\pm\xi e_\mu\}$. Take $u=(D_3^{(2)}E_2)(\cdot,y'')$ and $v=K_3(\cdot,y'')$.
\begin{itemize}
\item By Lemma~\ref{app:G1:I7-M}, $\abs{\bar D^{(1)}D_3^{(2)}E_2}\le A_M\rho^{-1}\mathrm{e}^{-\kappa d/8}$.
\item Let $d',\rho'$ denote $d,\rho$ with $y$ replaced by a point of $\{y,y\pm\xi e_\mu\}$. Then $d'\ge d-\xi$, hence $\rho'\ge\max(d,\xi)\ge\rho/2$ and $\mathrm{e}^{-\kappa d'}\le\mathrm{e}^{1/16}\mathrm{e}^{-\kappa d}$, since $\kappa\xi\le\frac1{16}$. So Lemma~\ref{app:G1:I7-K}(i) gives $\abs{v(\hat y)}\le\mathrm{e}^{1/16}(2A_0^T\rho^{-1}+A_E)\mathrm{e}^{-\kappa d}$.
\item Likewise Lemma~\ref{app:G1:I7-K}(ii) gives $\abs{u(\check y)}\le c_*A_\nabla\mathrm{e}^{1/16}\mathrm{e}^{-\kappa d}$ and $\abs{\bar D^{(1)}K_3}\le c_*B_*\rho^{-2}\mathrm{e}^{-\kappa d/2}$.
\end{itemize}
Here $d=d(y,y'')$ and $\rho=\rho(y,y'')$. Summing over $y''$ with Lemma~\ref{app:G1:I7-K}(iii) and the monotonicity of $J_n$ gives $Q_1$.
\item[(b)] Write $g_2(y)=\xi^3\sum_w\psi(w)E_3(y,y-w)$ with $\psi(w):=(D_3^{(2)}F)(y,y-w)$. This $\psi$ does not depend on $y$, and $\abs\psi\le c_*A_1^T\rho(w,0)^{-2}\mathrm{e}^{-\kappa d(w,0)}$. A difference in $y$ now shifts both arguments of $E_3$. For instance
\[
\xi^{-1}\big[E_3(y+\xi e_\mu,y+\xi e_\mu-w)-E_3(y,y-w)\big]=(D^{(1)}E_3)(y,y+\xi e_\mu-w)+(D^{(2)}E_3)(y,y-w),
\]
and similarly for backward differences. Each term is at most $c_*A_\nabla$ by Lemma~\ref{app:G1:I7-K}(ii), which gives $Q_2$.
\end{itemize}
Hence $\abs{T_F}\le c_*A_1^TJ_2(\kappa)(Q_1+Q_2)$. Since $T^{\rm c}=T_E+T_F$, the bound follows.

If all $K_i=F$, then $E_i=0$ and $g=g_F$ is constant. So $T^{\rm c}[F,F,F]=g_F\,\xi^3\sum_y(D_2\phi)(y)=0$, because the sum over the torus of a difference vanishes.
\end{proof}

\begin{proposition}[uniform kernel and coefficient bounds on the torus]\label{app:K:L11-1}\label{app:G1:I7-KC}
In the setting of this subsection there are constants $\kappa\in(0,\frac1{16}]$ and $C<\infty$, depending only on $a$ and $L$, such that the following hold. They hold for all $K,K_1,K_2,K_3\in\mathcal K$, all single differences $D,D',D_1,D_2,D_3$, all $w,w',x,y,y'\in\Lambda$, and all $\xi$, $j_0$ and $\ell$ as above.
\begin{itemize}
\item[(a)] $\abs{K(w,w')}\le C\rho^{-1}\mathrm{e}^{-\kappa d}$ and $\abs{(K-F)(w,w')}\le C\mathrm{e}^{-\kappa d}$. In either variable, $\abs{DK(w,w')}\le C\rho^{-2}\mathrm{e}^{-\kappa d}$ and $\abs{D(K-F)(w,w')}\le C\mathrm{e}^{-\kappa d}$. Moreover $\abs{D^{(1)}D'^{(2)}(K-F)(x,y)}\le C\rho(x,y)^{-1}\mathrm{e}^{-\kappa d(x,y)/8}$.
\item[(b)] The bounds (C1), (C2) and (C4) of Lemma~\ref{app:G1:I7-C} hold with right-hand side $C$, and (C3) holds with right-hand side $C\norm h$.
\item[(c)] (M1) With $T^{\rm c}$ as in Lemma~\ref{app:G1:I7-C21c}, $T^{\rm c}[F,F,F]=0$ and $\abs{T^{\rm c}[K_1,K_2,K_3]}\le C$.
\end{itemize}
\end{proposition}

\begin{proof}
Put $\kappa:=\kappa_G/2\in(0,\frac1{32}]$ and $\sigma_0:=\sup_{s\ge0}(1+s)\mathrm{e}^{-\kappa_Gs/2}$.

(a) By Lemma~\ref{app:G1:I7-K}(i) and $\rho\le1+d$, we have $\abs K\le(A_0^T+A_E\sigma_0)\rho^{-1}\mathrm{e}^{-\kappa d}$ and $\abs{K-F}\le A_E\mathrm{e}^{-\kappa d}$. By Lemma~\ref{app:G1:I7-K}(ii), $\abs{DK}\le c_*B_*\rho^{-2}\mathrm{e}^{-\kappa d}$ and $\abs{D(K-F)}\le c_*A_\nabla\mathrm{e}^{-\kappa d}$. By Lemma~\ref{app:G1:I7-M}, the mixed difference is at most $A_M\rho^{-1}\mathrm{e}^{-\kappa_Gd/8}\le A_M\rho^{-1}\mathrm{e}^{-\kappa d/8}$.

(b) is Lemma~\ref{app:G1:I7-C}, and (c) is Lemma~\ref{app:G1:I7-C21c}.

Take
\[
C:=\max\{A_0^T+A_E\sigma_0,\ c_*B_*,\ c_*A_\nabla,\ A_M,\ K_{\rm C},\ K_{\rm C}',\ K_\triangle,\ K_{\rm c}\}.
\]
Every quantity entering $C$ depends only on $(a,L)$, by the inputs (F), (E) and the bounds on $J_n$.
\end{proof}

\begin{lemma}[the vertex coefficient of the class (2.21c)]\label{app:G1:I7-C21c-coef}
Let $j_0\ge1$, $Mt_{j_0}\le1$ and $\mu_*t_{j_0}\le1$, and use the notation of Lemma~\ref{app:G1:I7-masses}. Consider a graph of order $2\le\alpha+2\beta\le4$ at zero background without operator lines, in the proof of \cite[(1.33)]{Ba3}. Let $\Gamma^{\rm c}$ be the vertex coefficient $\sum\xi^{2d}\Gamma'_\mu$ of \cite[p.~444]{Ba3} for the class \cite[(2.21c)]{Ba3}, summed over orderings and indices, and let $\Gamma^{\rm c}_F$ be its $C^\xi$-version. Assume the representation of Lemma~\ref{app:K:L11-2}(a), with $N$ terms. Then $\Gamma^{\rm c}_F=0$ and
\[
\abs{\Gamma^{\rm c}-\Gamma^{\rm c}_F}\le N\,C,
\]
with $C$ the constant of Proposition~\ref{app:K:L11-1}. In particular the difference is bounded by a constant that depends only on $(a,L)$ and on the graph, uniformly in $M$, $\mu_*$, $\xi$, $j_0$ and $\ell$.
\end{lemma}

\begin{proof}
$G_s=(-\Delta^\xi+M^2t_{j_0}^2+a_{j_0}P_1)^{-1}$ and $G_p=(-\Delta^\xi+\mu_*^2t_{j_0}^2+a_{j_0}P_1)^{-1}$. Here $M^2t_{j_0}^2,\mu_*^2t_{j_0}^2\in[0,1]$, and $a_{j_0}=a(1-L^{-2})(1-L^{-2j_0})^{-1}\in[a_\flat,a]$. Hence $G_s,G_p\in\mathcal K$, and $C^\xi=F$.

By Lemma~\ref{app:K:L11-2}(a),
\[
\Gamma^{\rm c}=\sum_{i=1}^Nc_i\,T^{\rm c}_i[G_s,G_s,G_p],\qquad\Gamma^{\rm c}_F=\sum_{i=1}^Nc_i\,T^{\rm c}_i[F,F,F],\qquad\abs{c_i}\le1,
\]
where each $T^{\rm c}_i$ has the form of Proposition~\ref{app:K:L11-1}(c) for some choice of single differences. By Proposition~\ref{app:K:L11-1}(c), each $T^{\rm c}_i[F,F,F]=0$ and $\abs{T^{\rm c}_i[G_s,G_s,G_p]}\le C$.
\end{proof}

\subsection{Power counting for generalised graphs}\label{app:G1:ssec-PC}

In this subsection we use the notation of Definition~\ref{app:G1:R-def}. Thus $L\ge2$, $N\in\N$, $\xi=L^{-N}$ and $s_j=L^{j-N}$ for $0\le j\le N$, so that $\xi=s_0\le s_j\le s_N=1$. The arguments work in every dimension $d\ge1$. For an index assignment $\mathbf j=(j_l)_{l\in\mathcal L}$ we write $\min\mathbf j:=\min_lj_l$.

A \emph{connected subgraph} of $\Gamma$ is a subgraph $H=(V_H,\mathcal L_H)$ as in Definition~\ref{app:G1:R-def}: it is given by a non-empty set of lines $\mathcal L_H$, $V_H$ is the set of their endpoints, and $H$ is connected. Every connected component of $\Gamma_i(\mathbf j)$ is a connected subgraph in this sense.

For the torus version, let $\ell\in\xi\N$. On $\xi\Z^d/\ell\Z^d$ put
\[
d_\ell(y,y'):=\min\{\abs z:\ z\in\xi\Z^d,\ z\equiv y-y'\ \text{mod}\ \ell\Z^d\}.
\]
Let $E^{\ell}_{\mathbf j}(\Gamma)$ be defined like $E_{\mathbf j}(\Gamma)$, except that the $y_v$ range over $\xi\Z^d/\ell\Z^d$ and $\abs{y_{u_l}-y_{w_l}}$ is replaced by $d_\ell(y_{u_l},y_{w_l})$.

\begin{remark}[the degree]\label{app:G1:PC-deg}
Give each line the dimension $-d+2$ minus the number of its differentiations, and let $p_v$ be the number of factors $\eta$ of $v$ other than the volume factor $\eta^d$. Then $D(H)$ is the degree \cite[(2.1)--(2.2)]{Ba3} of a graph without the vertices \cite[(1.7)]{Ba3} and \cite[(1.14)--(1.15)]{Ba3}. Indeed, write $D_H(v)$ for the degree of the vertex $v$ in $H$ in the sense of \cite[(2.1)--(2.2)]{Ba3}. Each line has two legs of dimension $-(d-2)/2$, so
\[
\sum_{v\in V_H}D_H(v)-d=d(\abs{V_H}-1)+\sum_{v\in V_H}p_v+\sum_{l\in\mathcal L_H}\omega_l.
\]
Theorem~\ref{app:K:L05-1} below is therefore the estimate \cite[(2.15)]{Ba3} for the generalised graphs of \cite[Prop.~2.2]{Ba3}, in the form used in this paper.
\end{remark}

\begin{lemma}[a spanning tree adapted to the scales]\label{app:G1:PC-tree}
Let $\mathbf j\in\{0,\dots,N\}^{\mathcal L}$. There is a spanning tree $T\subset\mathcal L$ of $\Gamma$ with the following property. For every $0\le i\le N$ and every connected component $H$ of $\Gamma_i(\mathbf j)$, the set $T\cap\mathcal L_H$ is a spanning tree of $H$; in particular it has exactly $\abs{V_H}-1$ lines.
\end{lemma}

\begin{proof}
Enumerate the lines as $l(1),\dots,l(\abs{\mathcal L})$ so that $j_{l(1)}\le j_{l(2)}\le\dots$; ties are broken arbitrarily. For each $p$ let $\Gamma^{(p)}$ be the graph with vertex set $V$ and lines $l(1),\dots,l(p)$. Let $T$ be the set of lines $l(p+1)$ whose endpoints lie in different components of $\Gamma^{(p)}$. Put $T_p:=T\cap\{l(1),\dots,l(p)\}$.

We show by induction on $p$ that $(V,T_p)$ is a forest with the same components (as vertex sets) as $\Gamma^{(p)}$. For $p=0$ this is clear. For the step, let $l=l(p+1)$.
\begin{itemize}
\item If the endpoints of $l$ lie in different components of $\Gamma^{(p)}$, these are also different components of $(V,T_p)$. Adding $l$ to both graphs merges the same two components, and creates no cycle in $T_{p+1}=T_p\cup\{l\}$.
\item Otherwise $l\notin T$, and adding $l$ to $\Gamma^{(p)}$ does not change its components. (Self-loops are in this case.)
\end{itemize}

Take $p=\abs{\mathcal L}$. Since $\Gamma$ is connected, $T$ is a spanning tree of $\Gamma$.

Now fix $i$. Because the enumeration is non-decreasing in $j$, the lines with $j_l\le i$ are exactly $l(1),\dots,l(p)$ for some $p$. So $\Gamma_i(\mathbf j)$ consists of the components of $\Gamma^{(p)}$ that contain a line. By the claim, every line of $T_p$ lies in one such component $H$. The lines of $T_p$ in $H$ form an acyclic connected graph on $V_H$, that is, a spanning tree of $H$ with $\abs{V_H}-1$ lines. Finally, $T\cap\mathcal L_H=T_p\cap\mathcal L_H$, since $\mathcal L_H\subset\{l(1),\dots,l(p)\}$.
\end{proof}

\begin{lemma}[one summation]\label{app:G1:PC-sum}
Let $\delta>0$ and $\xi\le s$. Then
\[
\sum_{y\in\xi\Z^d}\xi^d\,\mathrm{e}^{-\delta\abs y/s}\le C_\delta s^d,\qquad C_\delta=(1+2\sqrt d/\delta)^d .
\]
Moreover, for $\ell\in\xi\N$ and every $y'$,
\[
\sum_{y\in\xi\Z^d/\ell\Z^d}\xi^d\,\mathrm{e}^{-\delta d_\ell(y,y')/s}\le C_\delta s^d .
\]
\end{lemma}

\begin{proof}
By the Cauchy--Schwarz inequality, $\abs y\ge(\abs{y_1}+\dots+\abs{y_d})/\sqrt d$. So the first sum is at most $\big(\sum_{n\in\Z}\xi\,\mathrm{e}^{-a\abs n}\big)^d$ with $a:=\delta\xi/(s\sqrt d)>0$. Using $\mathrm{e}^a-1\ge a$ and $\xi\le s$,
\[
\sum_{n\in\Z}\xi\,\mathrm{e}^{-a\abs n}=\xi\Big(1+\frac2{\mathrm{e}^a-1}\Big)\le\xi+\frac{2\xi}a=\xi+\frac{2s\sqrt d}\delta\le s\Big(1+\frac{2\sqrt d}\delta\Big).
\]
For the torus, choose for each $y$ some $z(y)\in\xi\Z^d$ with $z(y)\equiv y-y'$ and $\abs{z(y)}=d_\ell(y,y')$. The map $y\mapsto z(y)$ is injective, because $z(y)$ determines $y$ modulo $\ell\Z^d$. Hence the torus sum is at most the sum over $\xi\Z^d$.
\end{proof}

\begin{theorem}[power counting for generalised graphs]\label{app:K:L05-1}\label{app:G1:R-PC}
Let $\Gamma=(V,\mathcal L)$ be a generalised graph (Definition~\ref{app:G1:R-def}), $L\ge2$, $N\ge0$, $\delta>0$, $\kappa>0$ and $J\subset\{0,\dots,N\}^{\mathcal L}$. Assume that for every $\mathbf j\in J$, every $i$ with $\min\mathbf j\le i<N$, and every connected component $H$ of $\Gamma_i(\mathbf j)$, we have $D(H)\ge\kappa$. Then
\begin{equation}\label{app:G1:PC-main}
\sum_{\mathbf j\in J}E_{\mathbf j}(\Gamma)\le C_\delta^{\abs V-1}\sum_{n\ge0}(n+1)^{\abs{\mathcal L}}L^{-\kappa n}<\infty,\qquad C_\delta:=(1+2\sqrt d/\delta)^d,
\end{equation}
uniformly in $N$ and $y_{v_0}$. The same bound holds for $E^\ell_{\mathbf j}(\Gamma)$, for every $\ell\in\xi\N$, uniformly in $\ell$. Consequently, for $E_{\mathbf j}$ and for $E^\ell_{\mathbf j}$:
\begin{itemize}
\item[(a)] if every connected subgraph of $\Gamma$ has positive degree, then $\sum_{\mathbf j}E_{\mathbf j}(\Gamma)\le C$, where the sum runs over all assignments $\mathbf j\in\{0,\dots,N\}^{\mathcal L}$;
\item[(b)] fix an ordering $l(1),\dots,l(\abs{\mathcal L})$ of $\mathcal L$, let $J=\{\mathbf j:j_{l(1)}\le\dots\le j_{l(\abs{\mathcal L})}\}$ and $G_p=\{l(1),\dots,l(p)\}$; if every component of every $G_p$ has positive degree, then $\sum_{\mathbf j\in J}E_{\mathbf j}(\Gamma)\le C$.
\end{itemize}
In (a) and (b), $C$ depends only on $d,\delta,L,\abs V,\abs{\mathcal L}$, the dimensions $\omega_l$ and the exponents $p_v$. In particular, $C$ does not depend on $N$, $y_{v_0}$, $\ell$, or on which vertices the lines join.
\end{theorem}

\begin{proof}
\emph{Step 1: a bound at fixed $\mathbf j$.} Fix $\mathbf j$, and let $T$ be the spanning tree of Lemma~\ref{app:G1:PC-tree}. Bound the exponential of every line not in $T$ by $1$. The remaining integrand is
\[
c_{\mathbf j}\prod_{l\in T}\mathrm{e}^{-\delta\abs{y_{u_l}-y_{w_l}}/s_{j_l}},\qquad c_{\mathbf j}:=\prod_v\xi^{p_v}\prod_ls_{j_l}^{\omega_l}.
\]
Root $T$ at $v_0$. For $v\ne v_0$ let $l_v\in T$ be the line from $v$ to its parent; then $T=\{l_v:v\ne v_0\}$. We claim that for every tree on $V$ rooted at $v_0$, every choice of scales on its lines and every $y_{v_0}$,
\[
\sum_{(y_v)_{v\ne v_0}}\prod_{v\ne v_0}\xi^d\prod_{l\in T}\mathrm{e}^{-\delta\abs{y_{u_l}-y_{w_l}}/s_{j_l}}\le\prod_{v\ne v_0}C_\delta s_{j_{l_v}}^d .
\]
We prove this by induction on $\abs V$; for $\abs V=1$ both sides equal $1$. If $\abs V\ge2$, the tree has a leaf $v\ne v_0$, and $y_v$ occurs only in the factor of $l_v$. By Lemma~\ref{app:G1:PC-sum}, the sum over $y_v$ is at most $C_\delta s_{j_{l_v}}^d$, uniformly in the position of the parent of $v$. What remains is the same expression for the tree $T\setminus\{l_v\}$ on $V\setminus\{v\}$, and the induction hypothesis applies.

Let $m_v:=\min\{j_l:\ l\text{ has endpoint }v\}$. This is well defined, since $\Gamma$ is connected and $\mathcal L\ne\emptyset$. Since $\xi\le s_{m_v}$ and $p_v\ge0$, we have $\xi^{p_v}\le s_{m_v}^{p_v}$. Hence
\begin{equation}\label{app:G1:PC-fixed}
E_{\mathbf j}(\Gamma)\le C_\delta^{\abs V-1}\prod_vs_{m_v}^{p_v}\prod_ls_{j_l}^{\omega_l}\prod_{l\in T}s_{j_l}^d .
\end{equation}

\emph{Step 2: the scale decomposition.} Write $s_j^a=\prod_{i=j}^{N-1}L^{-a}$; the product is empty for $j=N$. Then the right side of \eqref{app:G1:PC-fixed} equals $C_\delta^{\abs V-1}\prod_{i=0}^{N-1}L^{-A_i}$, where
\[
A_i:=\sum_{l:\,j_l\le i}\big(\omega_l+d\,1_T(l)\big)+\sum_{v:\,m_v\le i}p_v .
\]
The vertices with $m_v\le i$ are exactly the vertices of $\Gamma_i(\mathbf j)$. The vertex sets and line sets of the components $H$ of $\Gamma_i(\mathbf j)$ partition those of $\Gamma_i(\mathbf j)$, and by Lemma~\ref{app:G1:PC-tree}, $\abs{T\cap\mathcal L_H}=\abs{V_H}-1$. Therefore
\[
A_i=\sum_H\Big(\sum_{l\in\mathcal L_H}\omega_l+d(\abs{V_H}-1)+\sum_{v\in V_H}p_v\Big)=\sum_HD(H),
\]
where the sum runs over the components $H$ of $\Gamma_i(\mathbf j)$. If $i<\min\mathbf j$, then $\Gamma_i(\mathbf j)$ is empty and $A_i=0$.

\emph{Step 3: the sum over $J$.} Let $\mathbf j\in J$ and $\min\mathbf j\le i<N$. Then $\Gamma_i(\mathbf j)$ has at least one component, and each component has degree at least $\kappa$; so $A_i\ge\kappa$. Hence
\[
E_{\mathbf j}(\Gamma)\le C_\delta^{\abs V-1}L^{-\kappa(N-\min\mathbf j)}.
\]
An assignment with $\min\mathbf j=N-n$, $0\le n\le N$, has all $j_l\in\{N-n,\dots,N\}$, so there are at most $(n+1)^{\abs{\mathcal L}}$ of them. Summing over $n$ gives \eqref{app:G1:PC-main}. The series converges, since $\kappa>0$ and $L\ge2$. Nothing in the argument depends on $N$ or $y_{v_0}$.

\emph{Step 4: the torus.} Lemma~\ref{app:G1:PC-tree} and Step~2 do not involve the positions. In Step~1 use the torus part of Lemma~\ref{app:G1:PC-sum}, which is uniform in $\ell$ and in the parent's position. Step~3 is unchanged.

\emph{Step 5: (a) and (b).} Let $\mathcal F$ be the finite set of numbers
\[
\sum_{l\in A}\omega_l+dk+\sum_{v\in B}p_v,\qquad A\subset\mathcal L,\ B\subset V,\ k\in\{0,\dots,\abs V-1\}.
\]
Under the hypotheses of (a) or (b), $\mathcal F$ contains a positive number: the degree of $\Gamma$ itself, or of a component of $G_{\abs{\mathcal L}}=\mathcal L$. Let $\kappa_*$ be the least positive element of $\mathcal F$. It depends only on $d$, $\abs V$, the $\omega_l$ and the $p_v$. Every connected subgraph $H$ has $D(H)\in\mathcal F$ (take $A=\mathcal L_H$, $B=V_H$, $k=\abs{V_H}-1$). So $D(H)>0$ implies $D(H)\ge\kappa_*$.
\begin{itemize}
\item[(a)] Every component of every $\Gamma_i(\mathbf j)$ is a connected subgraph, so it has degree at least $\kappa_*$. Apply \eqref{app:G1:PC-main} with $J=\{0,\dots,N\}^{\mathcal L}$ and $\kappa=\kappa_*$.
\item[(b)] Let $\mathbf j\in J$ and $\min\mathbf j\le i<N$. Put $p:=\max\{q:\ j_{l(q)}\le i\}$; then $p\ge1$, because $j_{l(1)}=\min\mathbf j$. Since $\mathbf j$ is non-decreasing along the ordering, $j_{l(q)}\le i$ holds exactly for $q\le p$, so $\Gamma_i(\mathbf j)=G_p$. Its components have positive degree, hence degree at least $\kappa_*$. Apply \eqref{app:G1:PC-main} with $\kappa=\kappa_*$.
\end{itemize}
In both cases $C:=C_\delta^{\abs V-1}\sum_{n\ge0}(n+1)^{\abs{\mathcal L}}L^{-\kappa_*n}$ works, for $E_{\mathbf j}$ and for $E^\ell_{\mathbf j}$. It depends only on the listed data.
\end{proof}

\subsection{The remainder of the $A\cdot A$ vertex and inner renormalisation}\label{app:G1:AArem}

Throughout this subsection $0\le m^2_{j_0}\le1$, $(\alpha,\beta)\in\{(4,0),(2,1)\}$, $\Lambda$ is a label configuration, $N:=j_0$, and we work in $\xi$-units on $\xi\Z^3$ or on a lattice torus. The members $T$ of $\Pi^\Lambda_G$, the bonds $b=b_\mu(y)$ and $b'$, and the graphs $T^S$ are as in Proposition~\ref{app:G1:AA4-red}. For a set $S$ of scalar lines of $T$, the graph $T^S$ is $T$ with the lines of $S$ equal to the regular line $D:=G-C^\xi$ and the other scalar lines equal to $G=G^\xi_{j_0}(0)$. The vertex of $b'$ is summed, so it carries the constant external field. The members are built from the vertices
(i) $\xi^{n-1}\varphi\cdot q^n\nabla\varphi\,a^n$, $n\ge1$, of order $n/2$;
(ii) $\xi^{n-2}\varphi\cdot q^n\varphi\,a^n$, $n\ge2$ even, of order $n/2$;
(iii) $\lambda\abs\varphi^4$, of order $1$;
(iv) $c(y)\abs{\varphi(y)}^2$ with $\abs c\le\bar c_{\rm fin}$, of order $o_v\ge\frac12$.
All lines are proportional to $1_2$ in the charge index. A vertex of type (i) carries one lattice difference, on one of its two scalar legs. Generalised graphs, $E_{\mathbf j}$ and the degree $D$ are as in Definition~\ref{app:G1:R-def}. The standard degree $D^{\rm st}$, the number $n_D$ of regular lines, the order $o_H$ and the number $E_H$ of external legs of a connected subgraph $H$ are as in Lemma~\ref{app:G1:R-deg}, and $D=D^{\rm st}+2n_D$. For $H$ connected put $c_v:=3+p_v-\deg_v/2-\delta_v$, where $\delta_v$ is the number of differences of $v$ acting on lines of $H$. Every line of $H$ has standard dimension $-1$ minus its number of differences, and $2\abs{\mathcal L_H}=\sum_{v\in V_H}\deg_v-E_H$. Hence
\begin{equation}\label{app:G1:Rm-degf}
D^{\rm st}(H)=\sum_{v\in V_H}c_v+\tfrac{E_H}2-3 .
\end{equation}
Here $c_v=\frac n2+1-\delta_v$, $\frac n2$, $1$, $2$ for the types (i)--(iv). Thus $c_v\ge o_v$, and $c_v=o_v+1$ for a vertex of type (i) whose difference acts on an external leg of $H$. For a vertex of type (iv) we have $c_v=2\ge o_v+\frac12$, because $T$ contains a vertex with a vector leg besides it.

\paragraph{Lines.} Every internal scalar line $G$ is the sum of its slices of index $1\le j<N$ and of the bottom propagator, counted as the slice of index $0$. A slice with $\tau\in\{0,1,2\}$ differences, at most one at each end, satisfies $\abs{\partial^\tau G_{(j)}(y,y')}\le c\,s_j^{-1-\tau}\mathrm{e}^{-\delta_1\abs{y-y'}/s_j}$, by Theorem~\ref{app:G1:VSR-VS} and Lemma~\ref{app:G1:VSR-V0}, rescaled. For the photon lines and the label functions we use:
\begin{itemize}
\item[(V1)] every photon line $\gamma$ is a sum of slices $\gamma_{(j)}$ with weights in $[0,2]$, satisfying the same slice bounds, with at most one difference at each end;
\item[(V2)] for some $\vartheta\in(0,1)$,
$\abs{\nabla^{(1)}\gamma_{(j)}(w,z')-\nabla^{(1)}\gamma_{(j)}(w',z')}\le c\abs{w-w'}^\vartheta s_j^{-2-\vartheta}\big(\mathrm{e}^{-\delta_1\abs{w-z'}/s_j}+\mathrm{e}^{-\delta_1\abs{w'-z'}/s_j}\big)$ (compare \cite[(2.10)--(2.11)]{Ba3});
\item[(V3)] the label functions $g$ satisfy $\abs g\le1$ and $\abs{\nabla g}+\abs{\nabla\nabla g}\le c_g$.
\end{itemize}
Summing slices gives $\abs{\partial^\tau G(y,y')}+\abs{\partial^\tau\gamma(y,y')}\le c\rho^{-1-\tau}\mathrm{e}^{-\delta\abs{y-y'}}$ for $\tau\le2$, with at most one difference at each end. The regular line satisfies $\abs{\partial^\tau D(y,y')}\le c\,f_{1-\tau}(\rho)\mathrm{e}^{-\delta'\abs{y-y'}}$ for $\tau\le2$, where $\delta'$ is decreased to $\min(\delta',\delta'')$: for $\tau\le1$ by Lemma~\ref{app:G1:AA}, Step~3, and for $\tau=2$ by Lemma~\ref{app:G1:Rm-reg}. Shifting an endpoint by $\xi\mathrm{e}_\nu$ multiplies the constant of a slice bound by at most $\mathrm{e}^{\delta_1}$, since $\xi\le s_j$. It multiplies the constant of a bound of the regular line by at most $2\mathrm{e}^{\delta'}$, since $\xi\le1$, $\rho/2\le\rho(y-\xi\mathrm{e}_\nu,w)\le2\rho(y,w)$ and $1+\log_+(2/\rho)\le2(1+\log_+(1/\rho))$. Put $c_{\rm sh}:=2\mathrm{e}^{\max(\delta_1,\delta')}$. By Lemma~\ref{app:G1:R-dom} it is therefore dominated by a sum over $j\in\{0,\dots,N\}$ of lines of dimension $1-\tau$. A slice line with $\tau$ differences has dimension $-1-\tau$. All constants are uniform in $\xi,j_0,k,\eps,\ell_T$ and $m^2_{j_0}\in[0,1]$.

\begin{lemma}[the regular line with two differences]\label{app:G1:Rm-reg}
$\abs{\partial_y\partial_{y'}D(y,y')}\le c\,\rho(y,y')^{-1}\mathrm{e}^{-\delta''\abs{y-y'}}$, with $c,\delta''>0$ depending only on the constants of Lemma~\ref{app:G1:AA}.
\end{lemma}
\begin{proof}
We have $D=G\,X\,C^\xi$ with $X=(1-m^2_{j_0})\mathbf 1-a_{j_0}P_{j_0}$. By Lemma~\ref{app:G1:AA}, Steps~1--2, and Lemma~\ref{app:G1:VSR-F}(F1), $\abs{\partial G(y,z)}\le c\rho(y,z)^{-2}\mathrm{e}^{-\delta\abs{y-z}}$ and $\abs{\partial C^\xi(z,y')}\le c\rho(z,y')^{-2}\mathrm{e}^{-\abs{z-y'}/4}$. The kernel of $P_{j_0}$ is bounded and supported in $\abs{z-z'}\le2\sqrt3$. Put $r:=\abs{y-y'}$ and $\delta_0:=\min(\delta,\frac14)$.

The $P_{j_0}$ part is at most $c\,\mathrm{e}^{-\delta_0r/2}\big(\sum_z\xi^3\rho(y,z)^{-2}\mathrm{e}^{-\delta_0\abs{y-z}/2}\big)\sup_{z}\sum_{z'}\xi^3\rho(z',y')^{-2}1_{\abs{z-z'}\le2\sqrt3}\le c\,\mathrm{e}^{-\delta_0r/2}$. Since $\rho(y,y')\le1+r$, this is at most $c'\rho^{-1}\mathrm{e}^{-\delta_0r/4}$.

For the identity part, $\mathrm{e}^{-\delta\abs{y-z}-\abs{z-y'}/4}\le\mathrm{e}^{-\delta_0r}$. On $\{\abs{z-y'}\le2r\}$ we have $\abs{y-z}\le3r$, and $\rho(y,z)\ge\rho(y,y')/2$ or $\rho(z,y')\ge\rho(y,y')/2$. Moreover $\sum_{\abs w\le3r}\xi^3\rho(w)^{-2}\le c\rho(y,y')$. This part is therefore at most $c\rho^{-1}\mathrm{e}^{-\delta_0r}$. On $\{\abs{z-y'}>2r\}$ we have $\abs{y-z}\ge\abs{z-y'}/2$, hence $\rho(y,z)\ge\rho(z,y')/2$, and $\sum_{\abs w>2r}\xi^3\rho(w)^{-4}\le c\rho(y,y')^{-1}$. So this part is also at most $c\rho^{-1}\mathrm{e}^{-\delta_0r}$, and the lemma holds with $\delta''=\delta_0/4$.
\end{proof}

\begin{definition}[critical subgraphs; the sets $\mathcal K(T)$]\label{app:G1:Rm-crit}
Let $T$ be a member and $S\ne\emptyset$. A \emph{critical subgraph} of $T^S$ is a connected subgraph $H$ of $T$ (a non-empty set of lines with their endpoints) that contains no line of $S$ and has $D^{\rm st}(H)\le0$. It is of
type 1 if it is a single scalar line with two differences joining two distinct vertices $\varphi\cdot q\nabla\varphi\,a$;
type 2 if $E_H=2$, both external legs are scalar, and no difference of a vertex of $H$ acts on them;
type 3 if $E_H=2$, both external legs are scalar, and a difference acts on exactly one of them;
type 4 if $o_H=1$ and $H$ has exactly two external legs, both vector;
type 5 if $E_H=3$, with two scalar legs and one vector leg.
$\mathcal K(T)$ denotes the set formed by $T$ and its counterterm partners. A \emph{chain} is a member of the form of Proposition~\ref{app:G1:Rm-chain}(i).
\end{definition}

\begin{lemma}[the critical subgraphs]\label{app:G1:Rm-types}
Let $T$ be a member and $S\ne\emptyset$, and suppose $T$ does not fall under Lemma~\ref{app:G1:R-deg}(a). Then every critical subgraph of $T^S$ is of one of the types 1--5.
\end{lemma}
\begin{proof}
Let $H$ be critical. Since $H$ contains no line of $S$, $H\ne T$.

We first show $E_H\ge1$. If $V_H$ is the whole vertex set, a line of $T$ outside $H$ contributes two legs. Otherwise connectedness of $T$ gives a line from $V_H$ to the outside. The number of external scalar legs of $H$ is even, because each vertex has an even number of scalar legs. By Lemma~\ref{app:G1:R-deg}(a), $H$ does not have an odd number of external legs all of which are vector. Hence $E_H\ne1$.

Let $n_H$ be the number of vertices of type (i) in $H$ whose difference acts on an external leg of $H$. By \eqref{app:G1:Rm-degf} and the values of $c_v$, $0\ge D^{\rm st}(H)\ge o_H+E_H/2-3+n_H$.

If $E_H\ge5$, then $o_H\le\frac12$. So $V_H$ is one vertex of order $\frac12$ with self-loops, and $E_H\le1$, a contradiction.

If $E_H=4$, then $o_H\le1$. A single vertex with self-loops has $E_H\le2$. So $V_H$ consists of two vertices of order $\frac12$. A vertex of type (iv) has two legs, so with it $E_H\le3$. Hence both vertices are of type (i) with $n=1$ and $\abs{\mathcal L_H}=(6-4)/2=1$. A vector line carries no difference and has $D^{\rm st}=-1+3=2$. A scalar line with $\tau$ differences has $D^{\rm st}=2-\tau$. So the line is scalar with $\tau=2$, which is type 1.

If $E_H=3$, the scalar external legs are even in number and not all legs are vector, which is type 5.

Let $E_H=2$. If $\abs{V_H}=1$, the lines are self-loops and $\deg_v\ge4$, so $o_v\ge1$ by the list. If $\abs{V_H}\ge2$, then $o_H\ge1$. Next, $o_H=2$ is impossible: then $V_H$ contains all vertices and $E_H\ge2+2$. Suppose $o_H=\frac32$. Then exactly one vertex $v$, of order $\frac12$, is outside $V_H$; it has two scalar legs and at most one vector leg. Now $E_H=2$ is the sum of the external legs of $T$ at $V_H$, twice the number of lines of $T\setminus H$ inside $V_H$, and the number of lines joining $V_H$ to $v$, which is at least $1$. If two lines join $V_H$ to $v$, both external vector legs of $T$ sit at $v$, which has at most one leg left. If one line joins them, one external vector leg of $T$ sits at $v$, and the last leg of $v$ can neither be external nor form a line. Both cases are impossible. Hence $o_H=1$ and $n_H\le1$. Two scalar legs give type 2 or 3. Two vector legs give type 4. One leg of each kind contradicts parity.
\end{proof}

\paragraph{Counterterm partners.}\label{app:G1:Rm-ctpar}
For a connected graph $H$ with two external scalar legs, write $H_{\rm nn}$ for the sum of its terms in which no difference acts on an external leg. Following \cite[(1.21), p.~416; p.~417; (3.7)--(3.9), p.~435]{Ba3}, the mass counterterm generated by $H$ is $-\sum_{x'}\xi^3H_{\rm nn}(x,x')$, multiplying $\varphi(x)\cdot\varphi(x)$. In graph representation it is the graph of $H_{\rm nn}$, with the same lines and indices, in which the external leg at $x'$ is attached at $x$. A partner carries the lines of its subgraph, so $S$ determines the same set of lines in a member and in each of its partners.

\begin{lemma}[value of the counterterm]\label{app:G1:Rm-ctval}
Let $\Gamma=\Gamma_{\rm nn}+\sum_\mu\Gamma^\mu_{\rm nd}\nabla_\mu+\sum_\mu\nabla^*_\mu\Gamma^\mu_{\rm dn}+\sum_{\mu,\nu}\nabla^*_\mu\Gamma^{\mu\nu}_{\rm dd}\nabla_\nu$ be a kernel on $\xi\Z^3$ or on a lattice torus, with translation-invariant summable kernels. Then $\sum_{x'}\Gamma(x,x')=\sum_{x'}\Gamma_{\rm nn}(x,x')$ for every $x$.
\end{lemma}
\begin{proof}
Apply $\Gamma$ to the constant function $1$. Since $\nabla_\nu1=0$, the second and fourth terms vanish. The third term gives $\nabla^*_\mu c_\mu$ with $c_\mu(x)=\sum_{x'}\Gamma^\mu_{\rm dn}(x,x')$, which is constant by translation invariance. So $\nabla^*_\mu c_\mu=0$.
\end{proof}

Hence, for translation-invariant lines, the zero-momentum value of the amputated one-particle-irreducible function is represented by the parts $H_{\rm nn}$ alone, and terms with a difference on an external leg have no partner. In a member, a partner arises only from an order-one subgraph with two scalar legs. Such a subgraph is one of the following:
\begin{itemize}
\item a vertex $\lambda\abs\varphi^4$ with a scalar self-loop;
\item a vertex $\varphi\cdot q^2\varphi\,a^2$ with a vector self-loop;
\item two vertices $\varphi\cdot q\nabla\varphi\,a$ joined by one vector line and one scalar line (the \emph{exchange}).
\end{itemize}
Indeed, two vertices of order $\frac12$ must have their vector legs joined. If their scalar line were a self-loop, the vertex carrying it would form a subgraph with one vector leg and no scalar leg, and $T^S=0$ by Lemma~\ref{app:G1:R-deg}(a). A counterterm of order $\ge2$ leaves no vertex for the external vector legs. For a single-vertex subgraph, moving a leg from $x$ to $x$ changes nothing: the partner coincides with the member line by line, with the opposite sign \cite[(3.6)--(3.7), p.~435]{Ba3}.

\begin{lemma}[degrees in counterterm partners]\label{app:G1:Rm-ct}
Let $T$ be a member that does not fall under Lemma~\ref{app:G1:R-deg}(a), with lines of slice or regular type; or let $T$ be a relocated integrand of such a member (Definition~\ref{app:G1:Rm-reloc}) in which $x$ and $x'$ below are not relocated. Let $H\subset T$ be an exchange in its nn part: vertices $x\ne x'$ joined by a scalar line $s$, which carries the differences of $x$ and $x'$, and by a vector line $p$. Let $\ell,\ell'$ be the other scalar lines at $x,x'$, and let $w$ be the other endpoint of $\ell'$. Let $T^c$ be the graph with the same lines, dimensions and exponents in which $\ell'$ joins $w$ to $x$. Write $D$ and $D^c$ for degrees in $T$ and in $T^c$.
\begin{itemize}
\item[(a)] If $K$ is a connected subgraph of $T^c$ with $\ell'\in K$ and $\{s,p\}\not\subset K$, then $D^c(K)\ge1+2n_D(K)$.
\item[(b)] Let $G$ be a set of lines and $K$ a connected component of $G$ in $T^c$ with $\ell'\notin K$ or $\{s,p\}\subset K$. Then $K$ is a connected component of $G$ in $T$, with the same vertex set, and $D^c(K)=D(K)$.
\item[(c)] Every connected subgraph $K$ of $T^c$ satisfies $D^c(K)\ge-1+2n_D(K)$. In every ordering, the components of nonpositive degree of the initial segments of $T^c$ are components of the corresponding initial segments of $T$, with the same degree.
\end{itemize}
\end{lemma}
\begin{proof}
\emph{Preliminaries.} We have $p_x=p_{x'}=0$. Also $\ell\ne\ell'$, since otherwise $T=H\cup\{\ell\}$ would have no external vector leg. Next, $w\ne x'$, because $x'$ has only one scalar leg besides $s$. And $w\ne x$, because otherwise $\ell'=\ell$. In $T^c$ the lines at $x'$ are $s,p$, and those at $x$ are $\ell,\ell',s,p$; all other incidences agree. At standard dimensions $\omega^{\rm st}_s=-3$ and $\omega^{\rm st}_p=-1$, so $D^{\rm st}(H)=-1$. Lemma~\ref{app:G1:R-deg}(b) holds for relocated integrands with the same proof, since every vertex still carries at most its own difference, so that $c_v\ge o_v$.

(a) Put $K_H:=K\cap\{s,p\}$ and $M:=K\cup\{s,p\}$. Then $M$ is connected in $T$: a path in $K$ using $\ell'$ at $x$ in $T^c$ uses $\ell'$ at $x'$ in $T$, and $x,x'$ are joined by $s,p$. The vertex sets of $M$ in $T$ and of $K$ in $T^c$ agree, except that $x'\in V^c(K)$ if and only if $K_H\ne\emptyset$, while $x'\in V(M)$ always. The line sets differ by $\{s,p\}\setminus K_H$. Hence $D^c(K)=D(M)-\sum_{l\in\{s,p\}\setminus K_H}\omega_l-3\cdot1_{K_H=\emptyset}$. Writing $D=D^{\rm st}+2n_D$ on both sides gives $D^c(K)-2n_D(K)=D^{\rm st}(M)+c(K_H)$, with $c(\emptyset)=3+1-3=1$, $c(\{s\})=1$ and $c(\{p\})=3$.

It remains to show $D^{\rm st}(M)\ge0$. If $M=T$, this is Lemma~\ref{app:G1:R-deg}(b). Otherwise $D^{\rm st}(M)\ge o_M+E_M/2-3$ with $E_M\ge2$, and $o_M\ge\frac32$ because $M$ contains $x,x',w$. If $M$ contains a vertex of type (iv), then $c_v\ge o_v+\frac12$ gives $D^{\rm st}(M)\ge0$. Otherwise suppose $o_M=\frac32$. Then $2o_M$ is congruent modulo $2$ to the number of vector legs of $M$, which is therefore odd, while the number of scalar legs is even. So $E_M\ge3$. In all cases $D^{\rm st}(M)\ge0$.

(b) If $\{s,p\}\subset K$, the vertex sets, lines and dimensions coincide, and $\ell'$ joins $w$ to a vertex of $K$ in both graphs. Now let $\ell'\notin K$. Then $K$ is connected in $T$, with the same vertex set and degree. For maximality, suppose $\ell'\in G$. If $x'\in V(K)$, then $K$ contains $s$ or $p$, hence $x$, and $\ell'$ meets $K$ in $T^c$, which is excluded. If $x\in V(K)$, then $\ell'$ meets $K$ in $T^c$, again excluded. So the lines of $G$ meeting $V(K)$ are the same in $T$ and $T^c$.

(c) In case (a) the bound is $\ge1+2n_D$. In case (b) $D^c(K)=D(K)\ge-1+2n_D(K)$, by Lemma~\ref{app:G1:R-deg}(b),(c).
\end{proof}

\begin{corollary}\label{app:G1:Rm-ctcor}
Fix an ordering of the lines of $T^c$, and let $Q$ be a component of nonpositive degree of an initial segment $G_q$. Suppose a transformation replaces the graph by a sum of graphs. In each of them one line of $Q$ has its dimension raised by $s>0$, and one leg, at a vertex of $Q$, of a line coming after all lines of $Q$ has its dimension lowered by $s$. Then no component of any $G_{q'}$ decreases its degree. In particular, the bound of Lemma~\ref{app:G1:Rm-ct}(a) persists for every component of every $G_{q'}$.
\end{corollary}
\begin{proof}
A component containing neither the raised line nor the lowered leg is unchanged. A component containing only the raised line gains $s$. Let $K$ be a component of $G_{q'}$ containing the lowered line. Then all lines of $Q$ lie in $G_{q'}$, and $Q\subset K$, because the lowered line meets a vertex of $Q$. So the degree of $K$ changes by $s-s=0$.
\end{proof}

\paragraph{Integration by parts.}
Put $\nabla_\nu f(y)=\xi^{-1}(f(y+\xi\mathrm{e}_\nu)-f(y))$ and $\nabla^-_\nu f(y)=\xi^{-1}(f(y)-f(y-\xi\mathrm{e}_\nu))$. For absolutely summable products, $\sum_y\xi^3(\nabla_\nu f)h=-\sum_y\xi^3f\nabla^-_\nu h$. Moreover $\nabla^-_\nu(h_1h_2h_3)(y)=(\nabla^-_\nu h_1)(y)h_2(y)h_3(y)+h_1(y-\xi\mathrm{e}_\nu)(\nabla^-_\nu h_2)(y)h_3(y)+h_1h_2(y-\xi\mathrm{e}_\nu)(\nabla^-_\nu h_3)(y)$.

\begin{definition}[relocations]\label{app:G1:Rm-reloc}
Let $I$ be $T^S$, or the graph $(T^c)^S$ of a partner, with its non-regular lines decomposed into slices. A \emph{relocation} at a vertex $v$ of type (i) with $n=1$ replaces the difference $\nabla_\nu$, acting on the leg of a line $l$ at $v$, by $-\nabla^-_\nu$ acting on one other factor at $v$: the other scalar leg, the vector leg or the constant external field, or the label. The factors to its left are evaluated at $y_v-\xi\mathrm{e}_\nu$. A \emph{relocated integrand} is obtained by relocations at distinct vertices. Its generalised graph has the same lines, with $\omega_l=-1-\tau_l$ for non-regular lines and $\omega_l=1-\tau_l$ for regular lines, where $\tau_l$ is the number of differences on $l$.
\end{definition}

\begin{lemma}[{integration by parts, \cite[(2.8)]{Ba3}}]\label{app:G1:Rm-ibp}
Let $I$ be a relocated integrand, and fix an ordering of its non-regular lines with initial segments $G^\natural_q$. Let $l$ be a non-regular scalar line with $\tau_l=2$, joining vertices $x\ne x'$ of type (i) with $n=1$, neither of them relocated. Assume that no line preceding $l$ meets $x$ or $x'$, and that $x'$ is not the vertex of $b$. Then:
\begin{itemize}
\item[(a)] $\sum_{y_{x'}}\xi^3I=\sum_kI_k$, where the $I_k$, at most three, are the relocations of $I$ at $x'$.
\item[(b)] In $I_k$ the dimension of $l$ rises by $1$, and at most one other line $l^*_k$, a line at $x'$, loses $1$. If $l^*_k$ is non-regular, it follows $l$.
\item[(c)] The degree of a component $K$ of $G^\natural_q$ in $I_k$ is $D(K)+1$ if $l\in K\not\ni l^*_k$, and $D(K)$ otherwise.
\item[(d)] $I_k$ satisfies the line bounds, with constants enlarged by the factors $c_{\rm sh}$ and $c_g$.
\end{itemize}
\end{lemma}
\begin{proof}
At $x'$ the summand is $(\nabla_\nu f)(y)h_1h_2h_3(y)$, with $y=y_{x'}$ and $f$ the slice of $l$ as a function of $y$. Here $h_1$ is the other scalar line, $h_2$ the vector line or the constant external field, and $h_3$ the label. At fixed $\xi$ everything decays exponentially, so summation by parts and the Leibniz rule give (a); the term in which $\nabla^-$ falls on the constant field vanishes. In each term $f$ carries one difference fewer, and $\nabla^-h_i(y)=(\nabla h_i)(y-\xi\mathrm{e}_\nu)$. So the line $h_i$ gains a difference at its $x'$-end, where it had none, because $x'$ differentiated $l$. Since $l$ is not a self-loop and the lines at $x'$ other than $l$ do not precede $l$, (b) follows. Part (d) follows from the remark on shifts and (V3). For (c): a non-regular $l^*_k$ follows $l$ and meets $x'$, so every $G^\natural_q$ containing $l^*_k$ contains $l$ in the same component, and the net change there is $0$. A regular $l^*_k$ lies in no $G^\natural_q$.
\end{proof}

\begin{lemma}[members whose critical subgraphs are single lines]\label{app:G1:Rm-N1}
Let $T$ be a member that is not a chain, $S\ne\emptyset$, and suppose every critical subgraph of $T^S$ is of type 1 (possibly there is none). Then for every $T'\in\mathcal K(T)$,
\[
\Bigl|\sum_{b'}\xi^3T'^S(b,b')\Bigr|\le C,
\]
where $C$ depends only on $d,a,L$, the block size, $\bar c_{\rm fin}$, the constants of the line bounds and of (V1)--(V3), and the graph of $T$. In particular $\abs{\sum_{T'\in\mathcal K(T)}\sum_{b'}\xi^3T'^S(b,b')}\le\abs{\mathcal K(T)}\,C$.
\end{lemma}
\begin{proof}
If $T$ falls under Lemma~\ref{app:G1:R-deg}(a), then $T^S=0$; let $T$ not do so.

\emph{The elements of $\mathcal K(T)$.} If the subgraph of a partner is a single vertex, the partner is $-T^S$ line by line. If it is an exchange $\{s,p\}$, then $D^{\rm st}(\{s,p\})=-1$. By hypothesis $\{s,p\}$ is not critical, so $s\in S$, and the partner is $-(T^c)^S$ in the notation of Lemma~\ref{app:G1:Rm-ct}. A type-1 line $l$ carries the differences of both its endpoints, while the differences of $x,x'$ act on $s\in S$. Hence $l$ meets neither $x$ nor $x'$, and its endpoints have the same lines in $T$ and $T^c$. It thus suffices to bound $\sum_{b'}\xi^3I$ for $I=T^S$ and for $I=(T^c)^S$.

\emph{Transformation.} Decompose the non-regular lines into slices and split the indices by orderings of the non-regular lines, as in \cite[(2.7)]{Ba3}; fix an ordering. Let $Q$ be a component of some $G^\natural_q$ with $D(Q)\le0$. If $I=T^S$, then $Q$ is critical, hence a single type-1 line. If $I=(T^c)^S$, then by Lemma~\ref{app:G1:Rm-ct}(c) $Q$ is a component of the corresponding segment of $T$ with the same degree. So again $Q=\{l\}$ with $l$ a type-1 line of $T$, and $l$ has the same endpoints in $I$. In both cases no line preceding $l$ meets its endpoints. One endpoint is not the vertex of $b$, because $l$ is not a self-loop. Two such lines are vertex-disjoint, since the earlier one would meet an endpoint of the later one. Neither is a line $l^*_k$ of the other, since $l^*_k$ follows the transformed line and meets its endpoint. Apply Lemma~\ref{app:G1:Rm-ibp} to all of them in turn. By part (c), degrees never decrease, and each such $\{l\}$ ends with degree $1$. Since degrees are integers, every component of every $G^\natural_q$ now has positive degree.

\emph{Power counting.} Take absolute values, dominate the regular lines by Lemma~\ref{app:G1:R-dom}, and split the full index set into the orderings of all lines that extend the fixed one. A component of an initial segment $G_p$ without regular lines is a component of some $G^\natural_q$, so it has positive degree. A component $K$ with a regular line has degree $\ge1$. For $T$ this follows from Lemma~\ref{app:G1:R-deg}(c), which holds for relocated integrands with the same proof. For $T^c$ it follows from Lemma~\ref{app:G1:Rm-ct}(c), which applies to the relocated integrands, because the relocations take place at vertices other than $x,x'$. Theorem~\ref{app:K:L05-1}, in its form for a fixed ordering of the lines, bounds each term. The numbers of orderings and of terms $I_k$ depend only on the graph of $T$.
\end{proof}

\begin{lemma}[zeroth moment of the order-one vacuum polarisation]\label{app:G1:Rm-zm}
Let $\mathcal P$ be the order-one class with two external vector legs: the bubbles and the seagull tadpole. Then $\sup_x\abs{\sum_{\Pi\in\mathcal P}\sum_{x'}\xi^3\Pi_{\mu\mu'}(G)(x,x')}\le C_0$, uniformly in $\xi,j_0,k,\eps,\ell_T$ and $m^2_{j_0}\in[0,1]$.
\end{lemma}
\begin{proof}
By Lemma~\ref{app:G1:AA4-WT}, the corresponding sum with all scalar lines equal to $C^\xi$ vanishes. Write $G=C^\xi+D$ and expand; every remaining term has a line $D$. In the tadpole, $\abs{D(x,x)}\le c$. In a bubble, the two lines carry $\tau$ and $2-\tau$ differences, at most one at each end. All products except $(\nabla^{(1)}_\mu\nabla^{(2)}_\nu C^\xi)D$ are at most $c\rho^{-2}(1+\log_+\rho^{-1})\mathrm{e}^{-\delta''\abs{x-x'}}$, by Lemma~\ref{app:G1:AA}, Step~3, and Lemma~\ref{app:G1:Rm-reg}; their sums over $x'$ are bounded. In the exceptional product, the vertex at $x'$ carries the constant external field. Summation by parts in $x'$ gives $-\sum_{x'}\xi^3\nabla^{(1)}_\mu C^\xi(x,x')(\nabla^{(2)}_\nu D)(x,x'-\xi\mathrm{e}_\nu)$, which is at most $c\sum_{x'}\xi^3\rho^{-2}(1+\log_+\rho^{-1})\mathrm{e}^{-\delta''\abs{x-x'}}\le c'$.
\end{proof}

\begin{proposition}[the chains]\label{app:G1:Rm-chain}
\textup{(i)} If a member $T$ contains a connected subgraph $H$ of order one with two vector legs and $n_D(H)=0$, then
\[
T_{\mu\mu'}(b,b')=\sum_{z,z'}\xi^6\,\Pi^1_{\mu\nu}(b,z)\,g_1(z)\,\gamma_{\nu\nu'}(z,z')\,g_2(z')\,\Pi^2_{\nu'\mu'}(z',b'),\qquad\Pi^1,\Pi^2\in\mathcal P,
\]
with labels $g_1,g_2$, and $T$ has no counterterm partner and no vertex of type (iv). In particular, if $T^S$ has a critical subgraph of type 4, then $T$ is a chain.

\textup{(ii)} Let $\mathcal K_{\rm ch}$ be the set of chains of $\Pi^\Lambda_G$. Then
\[
\Bigl|\sum_{T\in\mathcal K_{\rm ch}}\sum_{S\ne\emptyset}(-1)^{\abs S+1}\sum_{b'}\xi^3T^S(b,b')\Bigr|\le C_{\rm ch},
\]
uniformly in $\Lambda,b,\xi,j_0,k,\eps,\ell_T$ and $m^2_{j_0}\in[0,1]$. Here $C_{\rm ch}$ depends only on $d,a,L$, the block size and the constants of the line bounds and of (V1)--(V3).
\end{proposition}
\begin{proof}
(i) If both legs of $H$ were lines of $T$, the rest of $T$ would have order one and four vector legs. So one leg of $H$ is an external leg of $T$, and the other is a photon line whose far side has order one, one external vector leg and no other leg; that far side lies in $\mathcal P$. A vertex of type (iv) has order $\ge\frac12$ and would raise the order above $2$. A chain contains no order-one subgraph with two scalar legs, so by Lemma~\ref{app:G1:Rm-ctval} and the list following it, it has no partner. A type-4 critical subgraph is such an $H$.

(ii) $T$ is multilinear in its scalar lines, so $\sum_{S}(-1)^{\abs S}T^S=T(C^\xi)$, the photon lines being kept. For fixed $\Pi^1$ and fixed $\emptyset\ne S_1\subset\mathcal L_s(\Pi^1)$, the sum over $\Pi^2$ and $S_2\subset\mathcal L_s(\Pi^2)$ of $(-1)^{\abs{S_1}+\abs{S_2}+1}\sum_{b'}\xi^3T^{S_1\cup S_2}$ therefore equals $(-1)^{\abs{S_1}+1}\Pi^{1,S_1}g_1\gamma g_2\sum_{\Pi^2}\sum_{b'}\xi^3\Pi^2(C^\xi)=0$, by Lemma~\ref{app:G1:AA4-WT}. All sums converge absolutely at fixed $\xi$. What remains is the part with $S_1=\emptyset$. By multilinearity it equals $\sum_{\Pi^1,\Pi^2}\sum_{b'}\xi^3\Pi^1(G)g_1\gamma g_2\big(\Pi^2(G)-\Pi^2(C^\xi)\big)$. The term with $\Pi^2(C^\xi)$ vanishes by Lemma~\ref{app:G1:AA4-WT}, so this part is $\sum_{\Pi^1}\sum_z\xi^3\Pi^1(G)(b,z)h(z)$, where
\[
h(z):=g_1(z)\sum_{z'}\xi^3\gamma(z,z')g_2(z')c_2(z'),\qquad c_2(z'):=\sum_{\Pi^2}\sum_{b'}\xi^3\Pi^2(G)(z',b').
\]
By Lemma~\ref{app:G1:Rm-zm}, $\abs{c_2}\le C_0$. By (V1)--(V3), $\abs h+\abs{\nabla h}+[\nabla h]_\vartheta\le c\,C_0\sum_j(s_j^2+s_j+s_j^{1-\vartheta})\le C_1$. A lattice path from $b$ to $z$, taken one coordinate at a time, gives $h(z)=h(b)+\sum_\nu(z-b)_\nu\nabla_\nu h(b)+R(b,z)$ with $\abs{R(b,z)}\le3\abs{z-b}^{1+\vartheta}[\nabla h]_\vartheta$. The zeroth moment $\sum_{\Pi^1}\sum_z\xi^3\Pi^1(G)(b,z)$ is bounded by Lemma~\ref{app:G1:Rm-zm}, and the first moment by Lemma~\ref{app:G1:AA}, Steps~4--5. The tadpole is local, so its remainder vanishes. For the bubbles, $\abs{\Pi^1(G)(b,z)}\le c\rho^{-4}\mathrm{e}^{-\delta\abs{b-z}}$, so their remainder is at most $cC_1\sum_z\xi^3\rho^{-4}\abs{z-b}^{1+\vartheta}\mathrm{e}^{-\delta\abs{b-z}}\le C$.
\end{proof}

\noindent The chains must be grouped as in (ii). For a single chain with $S\cap\mathcal L_s(\Pi^2)=\emptyset$, the factor $\sum_{b'}\Pi^2(G)$ is the zeroth moment of a single graph of degree $-1$, and the argument does not apply.

\begin{theorem}[{the remainder of \cite[(3.31)]{Ba3} reduced to the inner renormalisation}]\label{app:K:L05-2}\label{app:G1:Rm-main}
Let $0\le m^2_{j_0}\le1$ and $(\alpha,\beta)\in\{(4,0),(2,1)\}$. Let the photon lines be as in Lemma~\ref{app:G1:AA4}, satisfying (V1)--(V3), and assume the bounds (N.2), (N.3) and (N.5) of Statement~\ref{app:K:L05-3}. Then there is $K_R<\infty$ with $\abs{\hat c^{\,\Lambda}_{\mu\mu'}(y)}\le K_R$ for every label configuration $\Lambda$, every $y$, $\xi\in L^{-\N}$, $j_0$, $k$, $\eps$, $\ell_T$, and every $m^2_{j_0}\in[0,1]$, where $\hat c^{\,\Lambda}$ is as in Proposition~\ref{app:G1:AA4-red}. The constant $K_R$ depends only on $d$, $a$, $L$ and the block size; on the constants of the slice bounds of the scalar lines and of the photon lines (the latter independent of $\mu_*$, with weights in $[0,2]$); on the constants of (V2), (V3); on $C_N$; and on $\bar c_{\rm fin}$, with $\bar c_{\rm fin}:=0$ if the class contains no vertex of type (iv). The bound is obtained by arranging the absolutely convergent sum of Proposition~\ref{app:G1:AA4-red}(b) as follows. All chains are summed together, over all $S\ne\emptyset$ and over $\Pi^1,\Pi^2\in\mathcal P$. Every other member $T$ is summed at fixed $S$ over $\mathcal K(T)$.
\end{theorem}
\begin{proof}
By Proposition~\ref{app:G1:AA4-red}(b), $\hat c^{\,\Lambda}_{\mu\mu'}(y)=\sum_T\sum_{S\ne\emptyset}(-1)^{\abs S+1}\sum_{b'}\xi^3T^S(b,b')$ with $b=b_\mu(y)$, and the series converges absolutely; it may therefore be rearranged. By Proposition~\ref{app:G1:Rm-chain}(i) the chains have no partners. The remaining terms are grouped into the sets $\mathcal K(T)$, $T$ a member that is not a chain and not a counterterm, each counterterm belonging to the set of the member it is a partner of. At fixed $S$ the lines of $S$ are the same in all elements of $\mathcal K(T)$.

The chains contribute at most $C_{\rm ch}$, by Proposition~\ref{app:G1:Rm-chain}(ii). Let $T$ be a member that is not a chain, and let $S\ne\emptyset$. If $T$ falls under Lemma~\ref{app:G1:R-deg}(a), all elements of $\mathcal K(T)$ vanish at $S$. Otherwise, by Lemma~\ref{app:G1:Rm-types}, every critical subgraph of $T^S$ is of one of the types 1--5. Type 4 is excluded by Proposition~\ref{app:G1:Rm-chain}(i). If a critical subgraph of type 2, 3 or 5 exists, the sum over $\mathcal K(T)$ is bounded by $C_N$, by (N.2), (N.3) or (N.5). Otherwise every critical subgraph is of type 1, and Lemma~\ref{app:G1:Rm-N1} bounds the sum.

At fixed $\Lambda$, the number of members, the size of each $\mathcal K(T)$ and the number of patterns $S$ are bounded by numbers depending only on the order. The graphs of the members range over a finite set. Hence $\abs{\hat c^{\,\Lambda}_{\mu\mu'}(y)}\le K_R$ with $K_R$ as stated.
\end{proof}

\subsection{The torus moment of the class C4}\label{app:G1:ssec-C4moment}

\paragraph{Setting.} Let $K\le j_0\le\bar K$, $\xi=L^{-j_0}\le1$, and work in scale-$j_0$ units. Then $T^\xi=\xi\Z^3/\Gamma$ with $\Gamma=\bigoplus_\mu\ell'_\mu\Z e_\mu$ and $\ell':=\min_\mu\ell'_\mu\ge L^2\ge4$, as in the setting of Lemma~\ref{app:G1:TW-img}. Let $C:=(-\Delta^\xi+1)^{-1}$ on $\xi\Z^3$, and let $C_T$ be the same resolvent on $T^\xi$. Every line, scalar or photon, is $C_T$; the photon line is $C_T\,\delta_{\mu\mu'}$. For $v\in\xi\Z^3$ we write $C_T(v)$ for the value at the projection of $v$. Put $\rho(v):=\abs v+\xi$ and $\rho_T(v):=\dist(v,\Gamma)+\xi$. The vertices are those of \cite[(1.6), (1.7), (1.8), (1.10), p.~413]{Ba3} with $g_k\equiv1$. Here $q$ is the real antisymmetric $2\times2$ charge matrix of \cite[(1.8)]{Ba3}, so $q^2=-1$ and $\exp(\theta q)=\cos\theta+q\sin\theta$. Every symbol $\partial^a$ denotes a product of $a$ forward or backward differences, with arguments shifted by at most $2\xi$. We use:
\begin{itemize}
\item[(K1)] $\abs{\partial^aC(w)}\le A\rho(w)^{-1-a}\mathrm{e}^{-\abs w/4}$ for $a\le2$ (Lemma~\ref{app:G1:VSR-F}, (F0)--(F2), with $\abs w_1\ge\abs w$);
\item[(K2)] $\abs{\partial^aC_T(v)}\le A\rho_T(v)^{-1-a}$ for $a\le1$ (Lemma~\ref{app:G1:TW-img}(b));
\item[(K3)] $\sum_{w\in\xi\Z^3}\xi^3\rho(w)^{-k}\mathrm{e}^{-\kappa\abs w}\le C(k,\kappa)$ for $0\le k<3$, and $\le C(\kappa)\xi^{-1}$ for $k=4$ (Lemma~\ref{app:G1:TW-img}(c)); moreover $\sum_{v\in\xi\Z^3}\xi^3\mathrm{e}^{-\kappa\abs v}\rho_T(v)^{-1}\le C(\kappa)$;
\item[(K4)] $\sum_{W\in\Gamma\setminus0}\abs W^2\mathrm{e}^{-\kappa\abs W}\le C(\kappa)$.
\end{itemize}
The constants in (K3), (K4) do not depend on $\xi\le1$ or on $\ell'\ge4$. To prove the second part of (K3), let $m\in\Gamma$ be a point of $\Gamma$ nearest to $v$ and put $w:=v-m$. Since $0\in\Gamma$ we have $\abs w\le\abs v$ and $\abs m\le2\abs v$, so $\mathrm{e}^{-\kappa\abs v}\le\mathrm{e}^{-\kappa\abs m/4}\mathrm{e}^{-\kappa\abs w/2}$ and $\rho_T(v)=\rho(w)$. Summing over $w$ with the first part of (K3), and then over $m$, gives the bound, since $\sum_{m\in\Gamma}\mathrm{e}^{-\kappa\abs m/4}$ is bounded for $\ell'\ge4$. For (K4), write $W=(\ell'_\mu m_\mu)_\mu$ with $m\in\Z^3\setminus0$. Then $\abs W\ge4\abs m_\infty$, and $\abs W^2\mathrm{e}^{-\kappa\abs W}\le C\mathrm{e}^{-\kappa\abs W/2}\le C\mathrm{e}^{-2\kappa\abs m_\infty}$.

\paragraph{The vacuum functional.} Covariant Laplacians $\Delta_A$ on $T^\xi$ or $\xi\Z^3$ are built with the parallel transporters $\exp(s\xi qA_b)$, where $s\in\{1,-1\}$ is the sign convention of \cite[(1.8)]{Ba3}. For $\hat c\in\R^3$, $e\in\R$ and $\lambda\ge0$ put
\[
G(\hat c;e,\lambda):=-\log\int d\mu_{C_T}(a)\int d\varphi\;\exp\Big(-\tfrac12\ip{\varphi}{(-\Delta_{\hat c+ea}+1)\varphi}-\lambda\sum_x\xi^3\abs{\varphi(x)}^4-\tfrac12\sum_x\xi^3\delta m_1^2\abs{\varphi(x)}^2\Big).
\]
Here $\delta m_1^2=e^2s_e+\lambda s_\lambda$ is the part of order one of the mass counterterm \cite[(1.23), p.~417]{Ba3}, built with the lines $C_T$. Thus $\delta m^2_G=\sum_{x'}\xi^3\Sigma_G(x-x')$, where $\Sigma_G$ is the self-energy with undifferentiated external legs. The constant $\delta m_1^2$ does not depend on $\hat c$. Since $-\Delta_A\ge0$, $G$ is smooth in $(\hat c,e)$ and one-sidedly smooth in $\lambda\ge0$. Its Taylor coefficients at $e=\lambda=0$ are obtained by differentiating under the integral and applying Wick's theorem.

\begin{lemma}[the moment as a curvature in a constant field]\label{app:G1:C4-vac}
Let $\Pi^{(3)}$ and $\Pi^{(4)}$ be the zeroth moments of the classes C3 and C4 at a base point $y$, with every line, including the lines of the partners, equal to $C_T$. Then $\Pi^{(3)}_{\mu\nu}=\abs{T^\xi}^{-1}\partial_{\hat c_\mu}\partial_{\hat c_\nu}G(0;0,0)$. The part of $\Pi^{(4)}_{\mu\nu}$ of order $e^4$ is the coefficient of $e^2$ in $\abs{T^\xi}^{-1}\partial_{\hat c_\mu}\partial_{\hat c_\nu}G(0;e,0)$, and its part of order $e^2\lambda$ is the coefficient of $\lambda$ in $\abs{T^\xi}^{-1}\partial_{\hat c_\mu}\partial_{\hat c_\nu}G(0;0,\lambda)$. In particular, $\Pi^{(4)}$ does not depend on $y$.
\end{lemma}
\begin{proof}
By \cite[(1.4), (1.5), pp.~412--413]{Ba3}, the vertices (1.8) and (1.10) are the Taylor coefficients of the covariant lattice Laplacian in the total field $e(g_kA'+\tilde A)$, and (1.6), (1.7) are the remaining terms of the exponent. The graphs of \cite[p.~413]{Ba3} are the Wick terms of this expansion. Replacing every line by $C_T$ therefore gives the Wick terms of $G$, with $\delta m^2_1$ built from the same lines. A two-leg class with vector legs is the corresponding part of the second derivative in the total field. The two legs enter the same way, since $A'$ and $\tilde A$ occur only through $g_kA'+\tilde A$. Put the field equal to a constant $\tilde A\equiv c$ and set $\hat c=ec$. The quadratic term becomes $\frac12\abs{T^\xi}\,c\cdot\Pi c$, where $\Pi$ is the sum over the second leg, i.e.\ the zeroth moment, which is independent of the base by translation invariance. Each derivative $\partial_c=e\partial_{\hat c}$ carries one factor $e$. The internal order of C3 is $0$, and that of C4 is $e^2$ or $\lambda$, because the classes are defined by their order \cite[(2.1)--(2.3)]{Ba3} (Lemma~\ref{app:G1:cutclasses}(d)). Members containing a finite vertex do not occur in C4: such a member is connected, non-vanishing, has a line and two external legs, so $D\ge\frac12$ by Lemma~\ref{app:G1:Efin-deg}, while C4 has $D=0$. The partners of the C4 members are the counterterm terms of order one. When $\partial^2_{\hat c}$ is applied to $\frac12\delta m_1^2\sum\abs\varphi^2$ closed by a line, the result is exactly the members in which a self-energy subgraph without external legs is replaced by its local part.
\end{proof}

\begin{lemma}[constant fields and windings]\label{app:G1:C4-twist}
Let $K$ be one of the kernels $C$, or the kernel of a vertex at zero field, on $\xi\Z^3$, and let $K_{\hat c}$ be the corresponding kernel at the constant field $\hat c$ on $T^\xi$. Then
\[
K_{\hat c}(x,y)=\sum_{n\in\Gamma}\exp\big(sq\,\hat c\cdot(\tilde x-\tilde y+n)\big)K(\tilde x-\tilde y+n)
\]
for any lifts $\tilde x,\tilde y$. Consider a closed scalar loop of a Wick term, with kernels $K^{(i)}_{\hat c}(x_i,x_{i+1})$, and expand each kernel in this way with image $n_i$. The term with images $(n_i)$ equals $\operatorname{tr}[\exp(sq\,\hat c\cdot W)M]$, where $W:=\sum_in_i$ and $M$ is the product of the kernels at $\hat c=0$, which commutes with $q$. In particular,
\[
C_{T,\hat c}(x,x)=a(\hat c)\,1,\qquad a(\hat c)=\sum_{n\in\Gamma}\cos(\hat c\cdot n)\,C(n),\qquad\nabla a(0)=0 .
\]
\end{lemma}
\begin{proof}
On $\xi\Z^3$ the map $(g\varphi)(x):=\exp(sq\,\hat c\cdot x)\varphi(x)$ satisfies $\Delta_{\hat c+ea}=g\Delta_{ea}g^{-1}$ for every $a$. If the orientation convention of \cite[(1.4), (1.8)]{Ba3} gives $\Delta_{\hat c+ea}=g^{-1}\Delta_{ea}g$ instead, replace $\hat c$ by $-\hat c$ throughout. Every later use is a second derivative at $\hat c=0$ of $\cos(\hat c\cdot W)$, of $a(\hat c)$, or of an expression quadratic in $j$, and none of these changes. Hence the kernel at $\hat c$ is $\exp(sq\,\hat c\cdot(x-y))K(x-y)$; the factors commute, since each $K$ is a scalar times a power of $q$. The torus operator lifts to the $\Gamma$-periodic functions on $\xi\Z^3$. Its resolvent kernel is therefore the periodisation, which converges absolutely by (K1); the vertex kernels have range $\le\xi<\ell'/2$. Around a loop the arguments add up to $\sum_i(\tilde x_i-\tilde x_{i+1}+n_i)=W$. For the last claim, $\sum_n\sin(\hat c\cdot n)C(n)=0$ and $\sum_n n\,C(n)=0$, since $C$ is even.
\end{proof}

\begin{lemma}[the order $e^2\lambda$]\label{app:G1:C4-lam}
The part of $\Pi^{(4)}$ of order $e^2\lambda$ vanishes.
\end{lemma}
\begin{proof}
By Wick's theorem and Lemma~\ref{app:G1:C4-twist}, the coefficient of $\lambda$ in $G(\hat c;0,\lambda)/\abs{T^\xi}$ is $c_1a(\hat c)^2+c_2a_0a(\hat c)+c_3$, with $a_0:=a(0)=C_T(0)$ and constants $c_1,c_2,c_3$. The term $c_2a_0a(\hat c)$ is the counterterm. Its partner property, $\delta m^2_G=\sum_{x'}\xi^3\Sigma_G$ \cite[p.~417]{Ba3}, says that it cancels the terms of $\partial_{\hat c}^2(c_1a^2)$ in which both derivatives act on one loop and the other loop is the tadpole subgraph: $c_2a_0\partial^2a=-2c_1a_0\partial^2a$. (With the constants of \cite[(1.23)]{Ba3}, $N=2$: $c_1=8$, $c_2=-16$.) Hence the expression is $c_1(a-a_0)^2+\mathrm{const}$. Its second derivative at $0$ is $2c_1[\nabla a\otimes\nabla a+(a-a_0)\nabla^2a](0)=0$, by Lemma~\ref{app:G1:C4-twist}. Apply Lemma~\ref{app:G1:C4-vac}.
\end{proof}

\paragraph{The order $e^2$.} Write $-\Delta_{\hat c+ea}=-\Delta_{\hat c}+eM_1(a)+e^2M_2(a)+O(e^3)$. Expanding $-\log$ of the expectation gives, up to a constant, the order-$e^2$ part of $G(\hat c;e,0)$:
\[
\tfrac12E_a\tr(M_2C_{T,\hat c})-\tfrac14E_a\tr(M_1C_{T,\hat c}M_1C_{T,\hat c})-\tfrac18E_a\big(\tr M_1C_{T,\hat c}\big)^2+\tfrac12 s_e\tr C_{T,\hat c}.
\]
The kernels $M_1$ and $M_2$ come from the vertex (1.8) with $(n,n')=(0,1)$, and from (1.8), (1.10) with $(n,n')=(0,2)$. The vertex (1.8) has the form $[(\partial_\mu\varphi)(b)\cdot q\varphi(b_-)]\,a_b$ \cite[(1.8)]{Ba3}. Its leg $\varphi(b_-)$ is the undifferentiated position, and its leg $(\partial\varphi)(b)$ the differentiated position.

\begin{lemma}[the photon bridge]\label{app:G1:C4-bridge}
The third term contributes $-(\Pi^{(3)}\Pi^{(3)})_{\mu\nu}$ to the part of $\Pi^{(4)}_{\mu\nu}$ of order $e^4$.
\end{lemma}
\begin{proof}
We have $\tr(M_1(a)C_{T,\hat c})=\sum_b\xi^3a_bj_b(\hat c)$, and by translation invariance $j_b$ depends only on the direction of $b$. Moreover $j(0)=0$, because $M_1\propto q$, $C_T\propto1$ and $\tr q=0$. By Lemma~\ref{app:G1:C4-vac} at order zero, $\partial_{\hat c_\nu}j_{\mu}(0)=2\Pi^{(3)}_{\mu\nu}$, since $j_\mu=2\abs{T^\xi}^{-1}\partial_{\hat c_\mu}G(\hat c;0,0)$. Together with $\sum_{b'}\xi^3C_T(b,b')=1$ (because $(-\Delta+1)1=1$), the second derivative at $0$ of $-\frac18E_a(\cdot)^2/\abs{T^\xi}$ equals $-\frac14\sum_\rho\partial_\mu j_\rho\,\partial_\nu j_\rho=-(\Pi^{(3)}\Pi^{(3)})_{\mu\nu}$.
\end{proof}

By Lemma~\ref{app:G1:C4-twist}, every other term has a single scalar loop. Its charge trace is $\pm\tr\exp(sq\,\hat c\cdot W)=\pm2\cos(\hat c\cdot W)$. Hence these terms, divided by $\abs{T^\xi}$, equal $\sum_{W\in\Gamma}\cos(\hat c\cdot W)\,\tau_W$, with real $\tau_W$ independent of $\hat c$, and their second derivative at $0$ is $-\sum_{W\ne0}W_\mu W_\nu\tau_W$. The term $W=0$ drops out exactly.

\begin{lemma}[single loops with a winding]\label{app:G1:C4-loop}
There are constants $C,\kappa>0$, depending only on the constants in (K1)--(K3), such that $\abs{\tau_W}\le C\mathrm{e}^{-\kappa\abs W}$ for every $W\in\Gamma\setminus0$, uniformly in $\xi\le1$ and $\ell'\ge4$.
\end{lemma}
\begin{proof}
\emph{Unfolding.} Fix the reference point $u=b_-$ of one vertex at $0$, and let $u'=b'_-$ be that of the other. Lift the two scalar lines of the exchange term so that their displacements are $v\in\xi\Z^3$ and $W-v$. The photon line joins the midpoints of $b$ and $b'$; it is diagonal, so both bonds have the same direction $\mu$, and the photon factor is $C_T(v)=C_T(W-v)$. Let $S_W(v)$ be the resulting summand, summed over $\mu$ and over the four position patterns. By cyclicity of the trace, translation invariance, the reflection invariance of the vertex polynomial $\varphi(b_+)\cdot q\varphi(b_-)a_b$, and $C_T(v)=C_T(W-v)$, we have $S_W(v)=S_W(W-v)$. In the pattern where both lines leave from the undifferentiated positions of the self-energy, the term is $\sigma(v)C(W-v)$, with
\[
\sigma(v)=\kappa_\sigma\sum_\mu(\partial_\mu\partial^*_\mu C)(v)\,C_T(v),\qquad\sigma(-v)=\sigma(v),
\]
because $\partial_{1,\mu}\partial_{2,\mu}C(u,u')=-(\partial_\mu\partial^*_\mu C)(u'-u)$ is even. This is the self-energy of \cite[(1.23)]{Ba3}, so the exchange part of $s_e$ is $\sum_v\xi^3\sigma(v)$. If the internal line takes exactly one differentiated position, its factor is $\pm g(v+t)$ with $g:=\partial_\mu C$, $t\in\{0,-\xi e_\mu\}$, and $g(-w-\xi e_\mu)=-g(w)$; the closing line then carries one difference. If the internal line takes no differentiated position, its factor is $C(v)$, and the closing line carries two differences.

\emph{The case $\xi\ge r_0:=1/64$.} Here $\rho,\rho_T\ge r_0$. By (K1), every scalar kernel with at most two differences at an argument $p$ is at most $C(r_0)\mathrm{e}^{-\abs p/4}$, and shifts by at most $2\xi\le2$ cost a factor $\mathrm{e}^{1/2}$. By (K2), the photon kernel is bounded: $\abs{C_T}\le A/r_0$. It is periodic, so it does not decay. Hence $\abs{C_T(0)}\le A/r_0$, $\xi\abs{C_T(0)}\le A$, and $\abs{s_e}\le C\sum_v\xi^3\mathrm{e}^{-\abs v/4}+C\le C$ by (K3) with $k=0$. The single-vertex and counterterm terms are at most $C\mathrm{e}^{-\abs W/4}$. The exchange term is at most $C\sum_v\xi^3\mathrm{e}^{-(\abs v+\abs{W-v})/4}\le C\mathrm{e}^{-\abs W/8}\sum_v\xi^3\mathrm{e}^{-\abs v/8}\le C\mathrm{e}^{-\abs W/16}$.

\emph{The case $\xi<r_0$.} Let $h(v):=\chi_0(8\abs v)$, with $\chi_0$ smooth, $0\le\chi_0\le1$, $\chi_0=1$ on $[0,\frac12]$ and $\chi_0=0$ on $[1,\infty)$. Since $\abs W\ge4$, $h(v)h(W-v)=0$. By the symmetry $S_W(v)=S_W(W-v)$,
\[
-\tfrac14\sum_v\xi^3S_W=-\tfrac12\sum_v\xi^3h(v)S_W(v)-\tfrac14\sum_v\xi^3\big(1-h(v)-h(W-v)\big)S_W(v).
\]
Put $F(v):=C(W-v)$. For $\abs p\le\frac18+2\xi$ we have $\abs{W-p}\ge\abs W/2\ge2$, so (K1) gives $\abs{\partial^aF(p)}\le C\mathrm{e}^{-\abs W/8}$ for $a\le2$. For $\abs v\le\frac18$ we have $\rho_T(v)=\rho(v)$, so $\abs{\sigma(v)}\le C\rho(v)^{-4}$.

(a) \emph{Undifferentiated positions and counterterm.} The counterterm term is $\frac12\sum_v\xi^3\sigma(v)\,F(0)$, with the normalisation of $-\frac14\tr(\cdot)$ above. Together with the term $-\frac12\sum h\sigma F$, it gives $-\frac12$ times
\[
\sum_v\xi^3h(v)\sigma(v)\big[F(v)-F(0)-v\cdot\partial F(0)\big]-F(0)\sum_v\xi^3(1-h(v))\sigma(v).
\]
The linear term may be inserted because $h\sigma$ is even. Summing along a monotone lattice path from $0$ to $v$ gives $\abs{F(v)-F(0)-v\cdot\partial F(0)}\le3\rho(v)^2\max_{\abs p\le\abs v+\xi}\abs{\partial^2F(p)}$. Since $\sum_{\abs v\le1/8}\xi^3\rho^{-2}\le C$, the first sum is at most $C\mathrm{e}^{-\abs W/8}$. For $\abs v\ge\frac1{16}$, (K1), (K2) give $\abs\sigma\le C\mathrm{e}^{-\abs v/4}\rho_T(v)^{-1}$, so the second sum is at most $C\mathrm{e}^{-\abs W/4}$ by (K3).

(b) \emph{One differentiated position on the internal line.} The term is $\sum_vh(v)C_T(v)(\pm g(v+t))\,\partial F(v+t')$ with $\abs{t'}\le2\xi$. Replacing $\partial F(v+t')$ by $\partial F(0)$ costs at most $C\mathrm{e}^{-\abs W/8}\sum_{\abs v\le1/8}\xi^3\rho^{-3}\rho\le C\mathrm{e}^{-\abs W/8}$. For the remaining term put $k(w):=(hC_T)(w-t)$. The involution $w\mapsto-w-\xi e_\mu$ and the oddness of $g$ give
\[
\sum_w\xi^3g(w)k(w)=\tfrac12\sum_w\xi^3g(w)\big[(hC_T)(w-t)-(hC_T)(w+\xi e_\mu+t)\big].
\]
The two arguments differ by one lattice step. By (K2) and $\sup\abs{\nabla h}\le C$, the bracket is at most $C\xi\rho(w)^{-2}$ on its support $\abs w\le\frac14$. With $\abs g\le A\rho^{-2}$ and (K3) for $k=4$, the sum is at most $C\xi\cdot\xi^{-1}=C$; it is multiplied by $\abs{\partial F(0)}\le C\mathrm{e}^{-\abs W/8}$.

(c) \emph{No differentiated position on the internal line.} The term is at most $C\mathrm{e}^{-\abs W/8}\sum_{\abs v\le1/8}\xi^3\rho^{-2}\le C\mathrm{e}^{-\abs W/8}$.

(d) \emph{Both lines long.} On the support of $1-h(v)-h(W-v)$, every scalar argument $p$ satisfies $\abs p\ge\frac1{16}-2\xi\ge\frac1{32}$. So by (K1) the scalar factors are at most $C\mathrm{e}^{-(\abs v+\abs{W-v})/4}\le C\mathrm{e}^{-\abs W/8}\mathrm{e}^{-\abs v/8}$, and the photon is at most $A\rho_T(v)^{-1}$. By (K3) the sum is at most $C\mathrm{e}^{-\abs W/8}$.

(e) \emph{The terms with $M_2$.} The vertex (1.10) with $(n,n')=(0,2)$ is local at $b_-$, with undifferentiated legs. Its loop term $\propto C_T(0)C(W)$ is cancelled exactly by its part of the counterterm, since the local part of a local kernel is the kernel itself. The vertex (1.8) with $(n,n')=(0,2)$ carries the factor $\xi$ \cite[(1.8)]{Ba3} and one differentiated position, so it has no counterterm. Its loop term is at most $\xi\abs{C_T(0)}\,\abs{\partial C(W+O(\xi))}\le A\cdot C\mathrm{e}^{-\abs W/8}$ by (K2), (K1).

Adding (a)--(e) gives the claim with $\kappa=1/16$.
\end{proof}

\begin{lemma}[the torus moment of the class C4]\label{app:K:L04-2}
Assume \eqref{app:G1:eq-hyp}. Let $K\le j_0\le\bar K$ and $\xi=L^{-j_0}$, and replace every line of the members of the class C4 of Lemma~\ref{app:G1:cutclasses}, and of the partners of their proper divergent subgraphs, by the resolvent $C^\xi_T$ of $-\Delta^\xi+1$ on the rescaled torus $T^\xi$. Then the zeroth moment $\Pi^{(4)}(y)$ at a base point $y\in T^\xi$ does not depend on $y$ and satisfies $\abs{\Pi^{(4)}_{\mu\nu}(y)}\le C$, uniformly in $\xi$, $y$ and $\ell_T\ge L^3$. The constant $C$ depends only on $d$, on the classes (through Proposition~\ref{app:G1:TW}(i)), on $\chi_0$ and on the constants in (K1)--(K3); in particular $C=C(\nbar,d,L,a)$. More precisely,
\[
\Pi^{(4)}_{\mu\nu}=-(\Pi^{(3)}\Pi^{(3)})_{\mu\nu}-\sum_{W\in\Gamma\setminus0}W_\mu W_\nu\,\tau_W
\]
at order $e^4$, and $\Pi^{(4)}=0$ at order $e^2\lambda$.
\end{lemma}
\begin{proof}
By Lemma~\ref{app:G1:C4-vac}, $\Pi^{(4)}$ is independent of $y$ and is a second derivative in $\hat c$ at $0$. At order $e^2\lambda$ it vanishes by Lemma~\ref{app:G1:C4-lam}. At order $e^4$ the displayed formula follows from Lemma~\ref{app:G1:C4-bridge} and the winding expansion above. By Proposition~\ref{app:G1:TW}(i), $\abs{\Pi^{(3)}}\le C$. By Lemma~\ref{app:G1:C4-loop} and (K4), $\sum_{W\ne0}\abs W^2\abs{\tau_W}\le C\sum_{W\ne0}\abs W^2\mathrm{e}^{-\abs W/16}\le C$. All constants are independent of $\xi\le1$ and of $\ell'\ge4$, and $\ell_T\ge L^3$ gives $\ell'\ge4$.
\end{proof}

\subsection{Exceptional vertices with a bond-type partner}\label{app:G1:exc-bond}

We work in scale-$k$ units, with $\eta=L^{-k}$, $L=2$, $d=3$, $t=t_k$, $e(t)=t^{1/2}e$, $\lambda(t)=\lambda t$ and $p=p(t)=b_0(1+\log t^{-1})^{p_0}$. The scale-$j$ length is $s_j:=L^j\eta\le L^{-1}$ for $0\le j\le k-1$. The configuration $\mathfrak B$ is real, so every transport $U(\mathfrak B(\Gamma))=\exp(e(t)q\,\mathfrak B(\Gamma))$ is a rotation of $\R^2$. Here $q$ is antisymmetric with $q^2=-1$, so $w\cdot qw=0$ and $\tr q^{2i+1}=(-1)^i\tr q=0$. Contours satisfy $\Gamma_{x,y}=-\Gamma_{y,x}$ \cite[p.~572]{Ba2}, and $\sum_{x\in B^k(y)}\eta^d=1$. In (a), $\phi$ is a field on the coarse points of $\Omega_*$. Bound (a) holds for all coarse nearest neighbours in $\Omega_*$, and the vertex points are those in $\Lambda_7^{(k-1)\prime}$. In (b), $\mathfrak B$ is in addition regular on $\Omega_*$ in the sense of \cite[(2.23)]{Ba1}. In the proof of Proposition~\ref{app:G1:T579} both hold: $\phi$ is the full scalar field, which on $\supp\chi_k$ obeys the restrictions of Definition~\ref{app:G1:VSR-IH}, and the regularity is Lemma~\ref{app:G1:T579-reg}. Since $M\le m_2$, condition (E.m2) of Appendix~\ref{app:B} gives $M^2t^2\le1$. Following \cite[(2.8), p.~425]{Ba3}, we write $g_k(b)=g_k(b_-)$, and $\supp g_k\subset\Omega_1=B^k(\Lambda_7^{(k-1)\prime})$ \cite[p.~412]{Ba3}.

\begin{lemma}[the zeroth-order field near a vertex point]\label{app:G1:exc-L24}
Assume (a)--(c) of Proposition~\ref{app:K:L06-3}. There is a constant $C_{24}$ such that, for every vertex point $y$, every $x\in B^k(y)$ and every bond $b$ with $b_-\in B^k(y)$,
\[
\abs{\Phi_0(x)-U(\mathfrak B(\Gamma_{x,y}))\phi(y)}\le C_{24}p,\qquad\abs{(D_{\mathfrak B}\Phi_0)(b)}\le C_{24}p .
\]
In particular $\abs{\Phi_0(x)}\le X'_s+C_{24}p$. The constant $C_{24}$ depends only on $d,L,a,M_b,R,\varrho,e,\lambda,c_1',C_{\rm reg},m_2,b_0,p_0$ and on the constants of \cite[Props.~2.1--2.2]{Ba1}. It is independent of $\eps$, $k$ and $\ell_T$.
\end{lemma}

\begin{proof}
Condition (c) is the nesting used in the proof of Lemma~\ref{lem:lem24}, with $\Omega_*$ in place of $B^k(\Lambda_2^{(k-1)\prime})$ and the vertex points in place of $\Lambda_7^{(k-1)\prime}$. So we may take the objects defined there:
\begin{itemize}
\item $s:=\min\{4r(t_k),5\rho_k-2M_b\}$;
\item the cube $\square_2$ of big blocks of sup-radius $s$ around the big block of $y$, and the Neumann parallelepiped $\square:=B^k(\square_2)\subset\Omega_*$;
\item the sup-ball $Q(y,s)$, and the field $\square_1\phi:=\mathbf 1_{Q(y,s/2)}\phi$.
\end{itemize}
They satisfy $s/2\ge\rho_k$ and $s-c_d\ge R_0$, where $c_d=\sqrt d M_b$. Unlike in Lemma~\ref{lem:lem24}, the source is not cut, since $\phi$ is defined on all coarse points of $Q(y,s/2)\subset\Omega_*$.

\emph{Step 1 (box).} We follow Steps 1--2 of the proof of Lemma~\ref{lem:lem24}, which use \cite[(2.25), (2.26)]{Ba1} and \cite[Prop.~2.2]{Ba2} on the regular background $\mathfrak B$. They give
\[
\Phi_0(x)=a_k\big(G_k(\square,\mathfrak B)Q_k^*(\mathfrak B)\square_1\phi\big)(x)+E_1,\qquad \abs{E_1}\le Ca\,X'_s(1+r(t_k))^{d-1}\mathrm{e}^{-\delta_0(\rho_k-c_d)},
\]
and the same holds for $D_{\mathfrak B}$ at $b$. Since $p$ and $t^{-1/4}$ decrease in $t$ and $t_k=Lt_{k-1}$, we have $X'_s\le(c_1'/c_1)\,c_1(\lambda t_{k-1})^{-1/4}p(t_{k-1})$, and $r(t_k)\le r(t_{k-1})$. So \eqref{eq:lem24-cap} gives $\abs{E_1}\le C(c_1'/c_1)t_k^\kappa$.

\emph{Step 2 (constant background).} Let $A_0$ be the constant configuration with $A_{0,\mu}:=\mathfrak B_\mu(b_\mu)$, for fixed bonds $b_\mu\subset B^k(y)$. Join bonds by straight lattice paths in the convex set $\square$. By (b), $\abs{\mathfrak B-A_0}\le C\,r(t_k)\,p$ on $\square$, so $\sigma:=\abs e\,t^{1/2}\sup_\square\abs{\mathfrak B-A_0}\le Ct^{1/2}r(t_k)p$. The expansion \cite[(3.44)]{Ba1} and \cite[Prop.~2.2]{Ba1} on $\square$, used exactly as in \cite[(2.68), p.~572]{Ba2}, replace $(\mathfrak B,\mathfrak B)$ by $(A_0,A_0)$ in $G_k(\square,\cdot)Q_k^*(\cdot)$. The error is $\le C\sigma X'_s\le Cc_1'\lambda^{-1/4}t^{1/4}r(t)p^2$, both for the value at $x$ and for the derivative at $b$.

\emph{Step 3 (gauge).} By (T.nw) the box $\square$ does not wrap around the torus. So $A_0(\Gamma)$, for $\Gamma\subset\square$, depends only on the endpoints, by the Stokes argument of \cite[(2.69)--(2.72), pp.~572--573]{Ba2}. Put $\phi'(y'):=U(A_0(\Gamma_{y',y}))^{-1}\phi(y')$. By \cite[(2.74)]{Ba2},
\[
a_k\big(G_k(\square,A_0)Q_k^*(A_0)\square_1\phi\big)(x)=U(A_0(\Gamma_{x,y}))\big(a_kG_k(\square,0)Q_k^*\square_1\phi'\big)(x),
\]
and by \cite[(2.70)]{Ba2}, $D_{A_0}\big(U(A_0(\Gamma_{\cdot,y}))g\big)(b)=U(A_0(\Gamma_{b_-,y}))(\partial^\eta g)(b)$.

For coarse neighbours $y',y''\in Q(y,s/2)$ we have $\abs{\phi'(y'')-\phi'(y')}=\abs{U(A_0(\langle y',y''\rangle))\phi(y'')-\phi(y')}$ \cite[p.~573]{Ba2}. The coarse configuration $\bar{\mathfrak B}(\langle y',y''\rangle)$ is a linear average of $\mathfrak B$ over contours in $B^k(y')\cup B^k(y'')$ that reproduces constant configurations; this is the property used at \cite[p.~573]{Ba2}. Hence $\abs{\bar{\mathfrak B}(\langle y',y''\rangle)-A_0(\langle y',y''\rangle)}\le Cr(t)p$. By (a),
\[
\abs{\phi'(y'')-\phi'(y')}\le c_1'p+\abs e\,t^{1/2}Cr(t)p\,X'_s\le C_\phi p,\qquad C_\phi:=c_1'+C\abs ec_1'\lambda^{-1/4}\sup_{0<t\le1}t^{1/4}r(t)p(t)<\infty .
\]
Coarse straight paths in the cube $Q(y,s/2)$ give $\abs{\phi'(y')-\phi(y)}\le C_\phi p\,\abs{y'-y}_1$. Also $\phi'(y)=\phi(y)$.

\emph{Step 4 (zero background).} By \cite[(2.75)]{Ba2}, $a_kG_k(\square,0)Q_k^*1=1-M^2t^2G_k(\square,0)Q_k^*1$. Hence
\[
a_k\big(G_k(\square,0)Q_k^*\square_1\phi'\big)(x)-\phi(y)=T_1-T_2-T_3,
\]
where
\begin{gather*}
T_1:=a_k\big(G_k(\square,0)Q_k^*[\square_1(\phi'-\phi(y))]\big)(x),\qquad
T_2:=\phi(y)\,a_k\big(G_k(\square,0)Q_k^*\mathbf 1_{\square\setminus Q(y,s/2)}\big)(x),\\
T_3:=\phi(y)M^2t^2\big(G_k(\square,0)Q_k^*1\big)(x).
\end{gather*}
The kernel of $a_kG_k(\square,0)Q_k^*$ decays exponentially, uniformly for $M^2t^2\in[0,1]$ (Proposition~\ref{prop:UV-BaProp21}, on the parallelepiped $\square$). With Step 3 this gives:
\begin{itemize}
\item $\abs{T_1}\le C\,C_\phi p$;
\item $\abs{T_2}\le C\mathrm{e}^{-\delta_0(s/2-c_d)}X'_s$, which is bounded as in Step 1;
\item $\abs{T_3}\le CM^2t^2X'_s\le Cm_2^2c_1'\lambda^{-1/4}t^{7/4}p\le Cm_2^2c_1'\lambda^{-1/4}p$.
\end{itemize}
The same three bounds hold for $\partial^\eta$ at $b$, by \cite[(2.77)]{Ba2}, the derivative bounds of Proposition~\ref{prop:UV-BaProp21}, and $\partial^\eta(a_kG_kQ_k^*1)=-M^2t^2\partial^\eta G_kQ_k^*1$. These derivative bounds apply for the following reason. The point $y$ lies at distance $\ge s$ from $\square^c$, and $b$ lies within distance $\sqrt d+\eta$ of $y$. So $b$ lies at depth $\ge s-\sqrt d-\eta\ge s-c_d\ge R_0$ in $\square$. Here $c_d=\sqrt dM_b\ge\sqrt d+\eta$, because the big blocks are unions of $L$-blocks, so $M_b\ge L=2$. This is the depth used in Step~6 of the proof of Lemma~\ref{lem:lem24}.

\emph{Step 5 (back to $\mathfrak B$).} The contour $\Gamma_{x,y}$ has length $\le C(d,L)$ \cite[(2.2)]{Ba1}. So
\[
\abs{U(A_0(\Gamma_{x,y}))-U(\mathfrak B(\Gamma_{x,y}))}\abs{\phi(y)}\le\abs e\,t^{1/2}C\sup_{B^k(y)^{+1}}\abs{\mathfrak B-A_0}\,X'_s\le Ct^{1/4}p^2 .
\]
Also $\abs{(D_{\mathfrak B}-D_{A_0})\Phi_0(b)}\le\abs e\,t^{1/2}\abs{\mathfrak B_b-A_{0}}\abs{\Phi_0(b_+)}$. Steps 1--4 at the point $b_+$, which lies within distance $\eta$ of $B^k(y)$, give $\abs{\Phi_0(b_+)}\le X'_s+Cp$. So this term is $\le Ct^{1/4}p^2$ as well.

\emph{Collection.} The rotations $U$ preserve norms. Every error is bounded by one of $Ct_k^\kappa$, $C\,t^{1/4}r(t)p^2$, $Cp$, $Ct^{1/4}p^2$. Since $p\ge b_0$ and $\sup_{0<t\le1}t^{1/4}r(t)p(t)<\infty$ (by \eqref{eq:UV-r}, $r$ grows like a power of $1+\log t^{-1}$), each of them is $\le C'p$.
\end{proof}

\begin{lemma}[the contracted source]\label{app:G1:exc-src}
Let $\ell(A')$ be one of
\begin{gather*}
\text{(I)}\ \ \sum_b\kappa(b)\cdot A'_b,\quad \abs{\kappa(b)}\le\eta^dK,\quad\supp\kappa\subset\square_v;\\
\text{(II)}\ \ \sum_{x\in B^k(y)}\eta^dj(x)A'(\Gamma^{(k)}_{y,x}),\quad\abs{j}\le K,\quad\square_v:=B^k(y).
\end{gather*}
In case (I), $\square_v$ is a cube of side $2$. Let the line from $\ell$ to the leg $A'_{b'}$ have the kernel $\sum_{j<k}w_jG^\eta_{(j)}$, with $\abs{w_j}\le c_\ell$. This covers two cases:
\begin{itemize}
\item an ordinary vector line \cite[(2.6), p.~424]{Ba3}, with $w_j=1$ and $c_\ell=1$;
\item a D-line of Definition~\ref{app:G1:VSR-IH}, with $w_j=uc_ks_j^2$ and $c_\ell=\abs uc_k$, as in Theorem~\ref{app:G1:WP-WP1}.
\end{itemize}
Let $F(b')$ be the contraction of $\ell$ with $A'_{b'}$. Let $h'$ have bounded first and second differences, and let $S:=(\supp h')^{+\sqrt d}$. Then
\[
\norm{h'g_kF}_{1,\alpha_0}\le C_Fc_\ell K\,\mathrm{e}^{-\delta_{\rm l}\dist(S,\square_v)/2},
\]
where $\delta_{\rm l}$ is the rate in \cite[(2.10)--(2.12), p.~426]{Ba3} and $C_F$ depends only on $d$, $L$, $\alpha_0$, the constants there, and the $C^2$ bounds of $h'g_k$.
\end{lemma}

\begin{proof}
By \cite[p.~426]{Ba3}, the bounds (2.10)--(2.11) hold for the vector propagators $G^\eta_{(j)}$, with $\tilde B=0$. They are:
\begin{itemize}
\item $\abs{\partial^nG^\eta_{(j)}(b,b')}\le c\,s_j^{-d+2-n}\mathrm{e}^{-\delta_{\rm l}\abs{b-b'}/s_j}$ for $n\in\{0,1\}$;
\item the H\"older bound $c\,s_j^{-d+1-\alpha}\abs{b'_1-b'_2}^\alpha\mathrm{e}^{-\delta_{\rm l}\dist(\{b'_1,b'_2\},b)/s_j}$, for $\alpha=\alpha_0$.
\end{itemize}
The slices are symmetric, so the derivative may act on either variable. For contour legs, \cite[(2.12)]{Ba3} and the sentence before it give the same bounds with one more factor $s_j$, and with $\dist(B^j(x),b')\ge\abs{x-b'}-\sqrt ds_j$ in the exponent. By \cite[(1.3), p.~412]{Ba3}, $A'(\Gamma^{(k)}_{y,x})=\sum_jA'^{(j),\eta}(\Gamma^{(j+1)}_{x_{j+1},x})$. So in case (II),
\[
F(b')=\sum_jw_j\sum_x\eta^dj(x)G^\eta_{(j)}(\Gamma^{(j+1)}_{x_{j+1},x},b'),
\]
and in case (I), $F(b')=\sum_jw_j\sum_b\kappa(b)G^\eta_{(j)}(b,b')$.

For $\eta\le s\le1$, $c>0$ and any $z_0$,
\[
\sum_{z\in\eta\Z^d}\eta^d\mathrm{e}^{-c\abs{z-z_0}/s}\le\prod_{i=1}^d\Big(\eta\sum_{n\in\Z}\mathrm{e}^{-c\eta\abs{n-a_i}/(\sqrt d s)}\Big)\le\big(2\eta+2\sqrt d s/c\big)^d\le C_cs^d ,
\]
where $z_0=\eta a$, and we used $\abs w\ge\abs w_1/\sqrt d$ and $\sum_{n\in\Z}\mathrm{e}^{-\beta\abs{n-a_i}}\le2/(1-\mathrm{e}^{-\beta})\le2+2/\beta$, together with $\eta\le s$; so $C_c=(2+2\sqrt d/c)^d$.
\[
\]
Bonds add a factor $d$. Let $D$ be the distance of the evaluation points to $\square_v$. Since $s_j\le1$, we have $\mathrm{e}^{-\delta_{\rm l}\abs{z-b'}/s_j}\le\mathrm{e}^{-\delta_{\rm l}D/2}\mathrm{e}^{-\delta_{\rm l}\abs{z-b'}/(2s_j)}$ for $z\in\square_v$. This gives, in case (I),
\[
\abs{\partial^nF(b')}\le Cc_\ell K\mathrm{e}^{-\delta_{\rm l}D/2}\sum_js_j^{2-n},
\]
and in case (II) the same with $s_j^{3-n}$ and an extra factor $\mathrm{e}^{\sqrt d\delta_{\rm l}}$. The H\"older quotient of $\partial F$ is bounded in the same way, by $Cc_\ell K\mathrm{e}^{-\delta_{\rm l}\min_iD(b'_i)/2}\sum_js_j^{1-\alpha_0}$. Moreover $\sum_js_j^{2-n}\le\sum_{i\ge1}2^{-i}=1$ and $\sum_js_j^{1-\alpha_0}\le(2^{1-\alpha_0}-1)^{-1}$.

Finally, let $\psi:=h'g_k$, with $\abs\psi,\abs{\partial\psi},\abs{\partial^2\psi}\le c_2$. Then:
\begin{itemize}
\item $\partial_\mu(\psi F)(x)=\psi(x+\eta e_\mu)\partial_\mu F(x)+(\partial_\mu\psi)(x)F(x)$;
\item $\abs{\psi(x)-\psi(x')}$ and $\abs{\partial\psi(x)-\partial\psi(x')}$ are $\le2c_2\min\{1,\abs{x-x'}\}\le2c_2\abs{x-x'}^{\alpha_0}$;
\item $\psi F$ and $\partial(\psi F)$ vanish at points outside $S$. For $\abs{x-x'}\ge1$ the H\"older quotients of $\psi F$ and $\partial(\psi F)$ are therefore at most $2\sup_S\abs{\psi F}$ and $2\sup_S\abs{\partial(\psi F)}$. For $\abs{x-x'}<1$ with one point in $\supp h'$, both points lie in $S$. A lattice path inside $S$ then gives $\abs{F(x)-F(x')}\le\sqrt d\sup_S\abs{\partial F}\abs{x-x'}^{\alpha_0}$, and the H\"older quotient of $\partial F$ on $S$ is the one bounded above.
\end{itemize}
Hence $\norm{\psi F}_{1,\alpha_0}$ in the norm \cite[(1.32), p.~420]{Ba3}, with $\tilde B=0$, is at most $Cc_2$ times the suprema over $S$ found above.
\end{proof}

\begin{lemma}[mass counterterms through the exceptional vertex vanish]\label{app:G1:exc-odd}
Let $v$ have coupling order one, both scalar legs external, and a single internal vector leg. Let $H\ni v$ be a subgraph whose only external legs are the two scalar legs of $v$. Then the value of $H$ at zero background ($\tilde B=0$, $\tilde A=0$) vanishes.
\end{lemma}

\begin{proof}
Let $H_1:=H\setminus\{v\}$. The external legs of $H_1$ reduce to the one vector leg joined to $v$. So $H_1$ contains no $\tilde A$-leg and no external field (1.12) of \cite{Ba3}, and every scalar leg of $H_1$ is contracted inside $H_1$. The value of $H$ is $\sum_{b,b''}(\text{kernel of }v)(b)\,(\text{line})(b,b'')\,T_{H_1}(b'')$. It suffices to show $T_{H_1}\equiv0$.

At zero background every scalar line is $X\otimes\mathbf 1_{\R^2}$, with $X$ a real kernel, and every transport is $\mathbf 1$ (compare \cite[(1.31), p.~419]{Ba3}). The vertices of \cite[(1.6)--(1.8), (1.10), (1.13)--(1.14)]{Ba3} that can occur, and their products \cite[p.~414]{Ba3}, are then sums of products of terms $\phi_1\cdot q^m\phi_2$, with $m$ equal to the number of their $A'$-legs and with $(q^m)^T=\pm q^m$. The $R$-vertices carry $\tilde A^{\nbar+1}$ and do not occur. Contracting the scalar legs gives closed loops, each loop $\ell$ contributing $\pm\tr q^{m_\ell}$ times kernels. Here $\sum_\ell m_\ell$ equals the number of $A'$-legs of $H_1$, which is $2\#\{\text{lines of }H_1\}+1$, an odd number. So some $m_\ell$ is odd, and then $\tr q^{m_\ell}=0$.
\end{proof}

\begin{proposition}[exceptional vertices with zeroth-order legs: bond-type partner]\label{app:K:L06-3}
Let $1\le k\le K-1$, $t=t_k$, $p=p(t)$, $\lambda(t)=\lambda t$, $M\le m_2$, and $\Omega_*=B^k(\Lambda_2^{(k-1)\prime})$, with vertex points $y\in\Lambda_7^{(k-1)\prime}$. Let $\phi$ be a scalar field on the coarse points of $\Omega_*$, and put $\Phi_0:=a_kG_k(\Omega_*,\mathfrak B)Q_k^*(\mathfrak B)\phi$. Assume:
\begin{itemize}
\item[(a)] $\abs{\phi}\le X'_s:=c_1'p\lambda(t)^{-1/4}$, and $\abs{U(\bar{\mathfrak B}(\langle y',y''\rangle))\phi(y'')-\phi(y')}\le c_1'p$ for coarse nearest neighbours $y',y''$ in $\Omega_*$;
\item[(b)] $\abs{\partial^\eta_\mu\mathfrak B}\le C_{\rm reg}p$ on $\Omega_*$, and $\mathfrak B$ is regular there in the sense of \cite[(2.23)]{Ba1};
\item[(c)] $\dist(y,\Omega_*^c)\ge5\rho_k$ for every vertex point $y$, with $\rho_k:=(r(t_{k-1})-M_b)/L$, and \eqref{eq:lem24-R0row}, \eqref{eq:lem24-cap}, $\eps_0\le\eps(L,\varrho)$ and (T.nw) hold.
\end{itemize}
Then:
\begin{itemize}
\item[(i)] the bounds of Lemma~\ref{app:G1:exc-L24} hold, with $C_{24}$ independent of $\eps$, $k$ and $\ell_T$;
\item[(ii)] a vertex $v$ of coupling order one whose scalar legs are $\Phi_0$ or $\phi(y)$ is, as a function of its internal leg, of one of the following forms:
\begin{itemize}
\item $-e(t)\sum_b\eta^dJ(b)g_kA'_b$, for \cite[(1.8)]{Ba3};
\item $\sum_y\sum_{x\in B^k(y)}\eta^dA'(\Gamma_{y,x})j_y(x)$, for the products of \cite[(1.12)]{Ba3} or \cite[(1.13)]{Ba3} with \cite[(1.14)]{Ba3}.
\end{itemize}
Here $\sup_{b:\,g_k(b)\ne0}\abs{e(t)J(b)}\le C_{\rm ex}t^{1/4}p^2$ and $\sup_{y,x}\abs{j_y(x)}\le C_{\rm ex}t^{1/4}p^2$, with $C_{\rm ex}:=\max\{1,a\}\abs eC_{24}(2c_1'\lambda^{-1/4}+C_{24})$;
\item[(iii)] assume the counterterm rule of \cite[p.~420]{Ba3}, with mass counterterm coefficients equal to their zero-background values \cite[(1.29)--(1.31)]{Ba3}. Assume also the family bound of Definition~\ref{app:G1:VSR-IH} at graded order $\le\Nbar$ (Proposition~\ref{app:G1:VSR-CAP}), with constant $C_{\rm fam}$ and rate $\delta_{\rm fam}>0$, for families with external vector fields of finite norm \cite[(1.32)]{Ba3}. Let $F$ be a connected localised family of $\mathcal P^{(k)}_w$ with at least two vertices that contains such a vertex $v$, whose vector leg is contracted with a bond leg $(g_kA')_{b'}$ of a vertex $w$ of type \cite[(1.8)--(1.11)]{Ba3}. Then $F$ obeys \cite[(1.33)]{Ba3} with the factor $(\abs uc_k)^m$ and the rate $\delta_1:=\min\{\delta_{\rm fam},\delta_{\rm l}/2\}$, where the factor $e(t)\norm{h\Phi_1}_{1,\alpha_0}\norm{h\Phi_2}_{1,\alpha_0}$ of $v$ is replaced by $C\,C_{\rm ex}t^{1/4}p^2$. The constant $C$ is independent of $\eps$, $k$ and $\ell_T$. Only the suprema of $J$ and $j_y$ enter.
\end{itemize}
\end{proposition}

\begin{proof}
(i) This is Lemma~\ref{app:G1:exc-L24}.

(ii) By \cite[(1.8), p.~413]{Ba3} with $n=1$, $n'=0$, the vertex is $-e(t)\sum_b\eta^dJ(b)g_kA'_b$ with $J(b)=(D_{\mathfrak B}\Phi_0)(b)\cdot q\Phi_0(b_-)$. If $g_k(b)\ne0$, then $b_-\in\Omega_1$, so $b_-\in B^k(y)$ for some vertex point $y$. By (i),
\[
\abs{e(t)J(b)}\le\abs et^{1/2}C_{24}p\big(c_1'p\lambda^{-1/4}t^{-1/4}+C_{24}p\big)\le C_{\rm ex}t^{1/4}p^2 .
\]
By \cite[(1.14)]{Ba3} with $n=1$, $n'=0$, and the product rule of \cite[p.~414]{Ba3} (factor $-a_k$), the products with \cite[(1.12)]{Ba3} and \cite[(1.13)]{Ba3} give
\[
j_y(x)=a_ke(t)\phi(y)\cdot qU(\mathfrak B(\Gamma_{y,x}))\Phi_0(x),\qquad j_y(x)=-a_ke(t)(Q_k(\mathfrak B)\Phi_0)(y)\cdot qU(\mathfrak B(\Gamma_{y,x}))\Phi_0(x).
\]
Put $s_y(x):=U(\mathfrak B(\Gamma_{y,x}))\Phi_0(x)-\phi(y)$. Since $\Gamma_{y,x}=-\Gamma_{x,y}$, (i) gives $\abs{s_y}\le C_{24}p$. Also $(Q_k(\mathfrak B)\Phi_0)(y)=\phi(y)+r_y$, with $r_y:=\sum_{x'\in B^k(y)}\eta^ds_y(x')$ and $\abs{r_y}\le C_{24}p$. Since $\phi(y)\cdot q\phi(y)=0$ (as at \cite[p.~421]{Ba3}), the two coefficients equal
\[
a_ke(t)\phi(y)\cdot qs_y(x)\qquad\text{and}\qquad -a_ke(t)\big(\phi(y)\cdot qs_y+r_y\cdot q\phi(y)+r_y\cdot qs_y\big).
\]
With $a_k\le a$, the bounds are $a\abs et^{1/4}C_{24}p^2c_1'\lambda^{-1/4}$ and $a\abs et^{1/4}C_{24}p^2(2c_1'\lambda^{-1/4}+C_{24}t^{1/4})$. Both are $\le C_{\rm ex}t^{1/4}p^2$.

(iii) \emph{Every nonzero term keeps $v$ and its line.} Under the counterterm rule, a subgraph is replaced only if it has exactly two external scalar legs and no other external legs, or no external legs at all. Consider a subgraph containing exactly one of $v,w$. It has the line $vw$ as an external vector leg, so it is not replaced. A subgraph containing $v$ has the external scalar legs of $v$, so it is not a vacuum subgraph. A subgraph containing both $v$ and $w$ that is of mass type has exactly the two scalar legs of $v$ as its external legs. Its coefficient vanishes by Lemma~\ref{app:G1:exc-odd}.

\emph{Reduction.} Fix a graph $G$ of $F$ with nonzero contribution. The scale fields $A'^{(j),\eta}$ are independent and centred \cite[p.~414]{Ba3}. So the pair (leg of $v$, leg of $w$) contributes $g_k(b')F_v(b')$, where $F_v$ is the field of Lemma~\ref{app:G1:exc-src} with $K=K_v$, $K_v:=\sup\abs{e(t)J}$ in case (I) or $\sup\abs{j_y}$ in case (II). In case (I), $\kappa=-e(t)\eta^dJg_kh_v$. Hence $E(G)=E(G')$. Here $G'$ is $G$ with $v$ deleted and the leg of $w$ replaced by the external vector field $g_kF_v$. Moreover $d_v(G')=d_v(G)-1$, because $e(t)$ is contained in $K_v$, and $d_s(G')=d_s(G)$.

The subgraphs of $G'$ are the subgraphs of $G$ not containing $v$. They have the same external legs, since the leg of $w$ is external in both. So the counterterm rule acts on $G'$ exactly as on $G$, and $F'=\{G'\}$, with the same coefficients, is a localised family of graded order $\le\Nbar$ with one more external vector field. By the assumed family bound, $\abs{E(F)}=\abs{E(F')}$ is at most
\[
C_{\rm fam}(\abs uc_k)^{m'}e(t)^{d_v-1}\lambda(t)^{d_s}\mathrm{e}^{-\delta_{\rm fam}d'}\norm{h'g_kF_v}_{1,\alpha_0}\prod\norm{\cdot}_{1,\alpha_0},
\]
where the product runs over the other external fields, $d':=d(\{\square(u)\}_{u\ne v})$, and $m'=m-1$ if the line $vw$ is a D-line and $m'=m$ otherwise.

\emph{The new field.} Lemma~\ref{app:G1:exc-src} and part (ii) give $\norm{h'g_kF_v}_{1,\alpha_0}\le C_Fc_\ell C_{\rm ex}t^{1/4}p^2\mathrm{e}^{-\delta_{\rm l}D/2}$. Here $c_\ell c_k^{\,m'}\abs u^{m'}=(\abs uc_k)^m$, and $D:=\dist(S,\square(v))\ge\dist(\square(w),\square(v))-c_h$, where $c_h$ is the width of the support of $h'$ around $\square(w)$ plus $\sqrt d$.

\emph{The tree.} The localisation cubes have side $\le2$. Adding to a shortest tree for $\{\square(u)\}_{u\ne v}$ one edge from $\square(w)$ to $\square(v)$ gives $d(\{\square(u)\}_{u\in G})\le d'+\dist(\square(w),\square(v))+4\sqrt d$. Therefore
\[
\mathrm{e}^{-\delta_{\rm fam}d'}\mathrm{e}^{-\delta_{\rm l}D/2}\le\mathrm{e}^{\delta_1(4\sqrt d+c_h)}\mathrm{e}^{-\delta_1d(\{\square(u)\}_{u\in G})}.
\]
This is \cite[(1.33)]{Ba3} for $F$, with rate $\delta_1$, the factor $(\abs uc_k)^m$, and $C:=C_{\rm fam}C_F\mathrm{e}^{\delta_1(4\sqrt d+c_h)}$ in place of the factor $e(t)\norm{h\Phi_1}\norm{h\Phi_2}$ of $v$. No constant depends on $\eps$, $k$ or $\ell_T$, and $J$ and $j_y$ enter only through $K_v$.
\end{proof}

\subsection{Decay of King's slices}\label{app:G1:KS-decaysec}

In this subsection $i$ and $j$ are slice indices. Kernels are taken with respect to $\sum_y\eps^3$, as in Lemma~\ref{app:K:L08-1}. $d_1(x,y)$ is the $\ell^1$ distance on $T_\eps$, that is, the minimum over all representatives. The Euclidean torus distance and the $\ell^\infty$ torus distance are both at most $d_1$. So every bound below that decays in $d_1$ also holds with $\abs{x-y}$ in place of $d_1(x,y)$.

\begin{lemma}[King's slices are differences of block Green's functions]\label{app:G1:KS-rec}
In the setting of Lemma~\ref{app:K:L08-1}, put $\beta_j:=at_{j+1}^{-2}$. For $1\le j\le K-1$ we have $r_{j+1}=r_j\beta_j/(r_j+\beta_j)$ and
\[
H_j=G_{j+1}-G_j,\qquad G_j=(Y_{\mu_*}+r_jP_j)^{-1}.
\]
Together with $H_0=G_1$, this says that $H_i=G^{\mu_*}_{i+1}-G^{\mu_*}_i$ for $0\le i\le K-1$, with $G^{\mu_*}_0:=0$.
\end{lemma}

\begin{proof}
\emph{The recursion.} The identity $1/r_{j+1}=1/r_j+1/\beta_j$ is the same as $t_{j+1}^2/a_{j+1}=t_j^2/a_j+t_{j+1}^2/a$, that is, $1/a_{j+1}=L^{-2}/a_j+1/a$. Insert $1/a_j=(1-L^{-2j})/(a(1-L^{-2}))$. The right side becomes
\[
\frac{(1-L^{-2})+L^{-2}-L^{-2j-2}}{a(1-L^{-2})}=\frac{1-L^{-2(j+1)}}{a(1-L^{-2})}=\frac1{a_{j+1}} .
\]
\emph{The identity.} Write $r=r_j$ and $\beta=\beta_j$. By (S1) of Lemma~\ref{app:K:L08-1}, $Q_{j+1}=QQ_j$, so
\[
G_{j+1}^{-1}=Y_{\mu_*}+r_{j+1}Q_j^*Q^*QQ_j=G_j^{-1}-Q_j^*DQ_j,\qquad D:=r-r_{j+1}Q^*Q .
\]
Since $QQ^*=1$, $Q^*Q$ is an orthogonal projection. By the recursion, $r-r_{j+1}=r^2/(r+\beta)>0$. Hence $D>0$, and
\[
D^{-1}=r^{-1}(1-Q^*Q)+\frac{r+\beta}{r^2}Q^*Q=r^{-1}+\beta r^{-2}Q^*Q .
\]
Therefore
\[
S:=D^{-1}-Q_jG_jQ_j^*=r^{-2}\big(r-r^2Q_jG_jQ_j^*+\beta Q^*Q\big)=r^{-2}(C^{(j)})^{-1},
\]
and $S$ is invertible by (S2) of Lemma~\ref{app:K:L08-1}.

For invertible $A$ and $D$ with $S=D^{-1}-U^*A^{-1}U$ invertible, a direct multiplication shows that
\[
(A-UDU^*)\big(A^{-1}+A^{-1}US^{-1}U^*A^{-1}\big)=1+U\big[(1-DU^*A^{-1}U)S^{-1}-D\big]U^*A^{-1}=1,
\]
because $1-DU^*A^{-1}U=DS$. Apply this with $A=G_j^{-1}$ and $U=Q_j^*$, in finite dimension. It gives
\[
G_{j+1}=G_j+r_j^2G_jQ_j^*C^{(j)}Q_jG_j=G_j+F_jC^{(j)}F_j^*=G_j+H_j .
\]
The last claim follows from $H_0=G_1$ and $G_j=G^{\mu_*}_j$.
\end{proof}

\begin{lemma}[decay of King's slices]\label{app:K:L08-3}\label{app:G1:KS-decay}
In the setting of Lemma~\ref{app:K:L08-1}, let $M\in\{m_1,m_2\}$ and assume
\[
\max(M,\mu_*)\,t_{K-1}\le1 .
\]
Let $G_{(i)}:=G^M_{i+1}-G^M_i$ for $0\le i\le K-1$, with $G^M_0:=0$. Let $\gamma_0=\min(4,a_\flat)$, $c=c(a,L)$ and $\kappa_H$ be the constants defined after Lemma~\ref{app:G1:VSR-CT}, and let $C_0,C_1,C_2$ be the constants of Theorem~\ref{app:G1:VSR-VS}. Put
\[
\delta_1:=\kappa_H,\qquad c_G:=\max\big(C_0,C_1,C_2,8L^2\gamma_0^{-1}\mathrm{e}^{2c/L}\big),\qquad c_H:=\max\big(C_0,2L^2\gamma_0^{-1}\big).
\]
Then for all $x,y\in T_\eps$ and $0\le i\le K-1$,
\[
\abs{\nabla^\alpha G_{(i)}(x,y)}\le c_Gt_i^{-1-\abs\alpha}\mathrm{e}^{-\delta_1d_1(x,y)/t_i},\qquad\abs{H_i(x,y)}\le c_Ht_i^{-1}\mathrm{e}^{-\delta_1d_1(x,y)/t_i}.
\]
The first bound holds for all $\abs\alpha\le2$ with at most one forward difference in each variable. Both bounds hold with $\abs{x-y}$ in place of $d_1(x,y)$.

The constants $c_G$, $c_H$ and $\delta_1$ depend only on $a$ and $L$. In particular they are independent of $\eps$, $i$, $K$, $M$, $\mu_*$ and $\ell_T$. At zero background, \cite[(2.10)--(2.11), p.~426]{Ba3} hold with the rate $\delta_1$ (Corollary~\ref{app:G1:VSR-VSphys}).

In the setting of Lemma~\ref{app:K:L08-1} we have $t_K=\eps_0$, so $t_{K-1}=\eps_0/L$. The assumption therefore holds because $m_1<m_2$, $m_2\eps_0\le1$ and $\mu_*\eps_0\le0.2$ (Appendix~\ref{app:B}, Table~\ref{p:tab:A4}).
\end{lemma}

\begin{proof}
\emph{Scaling.} Fix $0\le i\le K-1$, put $\eta:=\eps/t_i=L^{-i}$, and let $(Uf)(u):=f(t_iu)$. This map sends the functions on $T_\eps$ to the functions on the $\eta$-lattice torus $\Lambda$ of side $\ell:=\ell_T/t_i$. Since $\ell_T/t_K\in\N$ and $i+1\le K$, we have $\ell\in L\N$ and $\ell\ge L$. Moreover:
\begin{itemize}
\item $-\Delta_\eps=t_i^{-2}U^{-1}(-\Delta_\eta)U$;
\item the forward difference satisfies $\nabla_{\eps,\mu}=t_i^{-1}U^{-1}\nabla_{\eta,\mu}U$;
\item for $k\ge i$, $P_k$ becomes the projection $P_{L^{k-i}}$ onto the functions constant on the blocks of side $L^{k-i}$.
\end{itemize}
By the definitions in Lemma~\ref{app:K:L08-1} and its part (S4), $G^M_k=(Y_M+r_kP_k)^{-1}$ with $r_k=a_kt_k^{-2}$ for $1\le k\le K$. Hence
\[
Y_M+r_iP_i=t_i^{-2}U^{-1}XU,\qquad Y_M+r_{i+1}P_{i+1}=t_i^{-2}U^{-1}X'U,
\]
with
\[
X:=-\Delta_\eta+\tilde\mu^2+a_iP_1,\qquad X':=-\Delta_\eta+\tilde\mu^2+a_{i+1}L^{-2}P_L,\qquad\tilde\mu:=Mt_i .
\]
Let $B$ be an operator on $\Lambda$ with kernel taken with respect to $\eta^3\sum_v$. Since $\eps^3=t_i^3\eta^3$, the operator $t_i^2U^{-1}BU$ has the kernel $t_i^{-1}B(x/t_i,y/t_i)$ with respect to $\eps^3\sum_y$. Each forward difference contributes a factor $t_i^{-1}$, because $\nabla_{\eta,\mu}U=t_iU\nabla_{\eps,\mu}$, and $d(x/t_i,y/t_i)=d_1(x,y)/t_i$ for the $\ell^1$ torus distance $d$ of $\Lambda$. Also, $\tilde\mu=Mt_i\le Mt_{K-1}\le1$. By (S2) of Lemma~\ref{app:K:L08-1}, $a_\flat\le a_k\le a$ for $1\le k\le K-1$. The formula $1/a_k=(1-L^{-2k})/(a(1-L^{-2}))$, used in the proof of Lemma~\ref{app:G1:KS-rec}, holds for every $k\ge1$, including $k=K$. It gives $a(1-L^{-2})<a_k\le a$, that is, $a_\flat<a_k\le a$ for every $k\ge1$, since $a_\flat=a(1-L^{-2})$.

\emph{The case $1\le i\le K-1$.} Here $\eta\in(0,1]$, $1/\eta\in\N$, $\ell\in L\N$, $\tilde\mu\in[0,1]$, $b=a_i$ and $b'=a_{i+1}L^{-2}$ with $a_i,a_{i+1}\in[a_\flat,a]$. This is the setting of Theorem~\ref{app:G1:VSR-VS} with slice index $i$ and mass $M$. There $H:=X'^{-1}-X^{-1}$, so $G_{(i)}$ has the kernel $t_i^{-1}H(x/t_i,y/t_i)$. Parts (S1), (S2) and (S3) of that theorem bound $H$, a single difference in either variable, and one difference in each variable, by $C_0$, $C_1$ and $C_2$ times $\mathrm{e}^{-\kappa_Hd}$. After scaling, this is the first bound with constant $\max(C_0,C_1,C_2)$.

\emph{The case $i=0$.} Here $\eta=1$, $t_0=\eps$, $a_1=a$, and $G_{(0)}=G^M_1$ has the kernel $\eps^{-1}X'^{-1}(x/\eps,y/\eps)$ with $X'=-\Delta_1+\tilde\mu^2+aL^{-2}P_L$.
\begin{itemize}
\item By Lemma~\ref{app:G1:VSR-CT}(i) with $s=L$, $X'\ge\min(4L^{-2},aL^{-2})\ge\gamma_0L^{-2}$.
\item Apply Lemma~\ref{app:G1:VSR-CT}(ii) with $\gamma=\gamma_0L^{-2}$, $\beta=aL^{-2}$, $\kappa=c/L\le\frac12$ and $\rho=d(\cdot,y)$, which is $1$-Lipschitz. Its condition reads $L^{-2}\big[3.4c^2+a(\mathrm{e}^{3c}-1)\big]\le L^{-2}\gamma_0/2$, which is the definition of $c$. Hence $\norm{\mathrm{e}^{\kappa\rho}X'^{-1}\mathrm{e}^{-\kappa\rho}}\le2L^2/\gamma_0$.
\item At spacing $1$ the vectors $\delta_x$ are unit vectors and $\rho(y)=0$. So $\abs{X'^{-1}(x,y)}\le2L^2\gamma_0^{-1}\mathrm{e}^{-(c/L)d(x,y)}$.
\item A forward difference at spacing $1$ is the difference of two values at points whose distance changes by at most $1$. So one difference in each variable costs at most the factor $4\mathrm{e}^{2c/L}$.
\end{itemize}
Since $\kappa_H\le c/L$, after scaling we get $\abs{\nabla^\alpha G_{(0)}(x,y)}\le8L^2\gamma_0^{-1}\mathrm{e}^{2c/L}\eps^{-1-\abs\alpha}\mathrm{e}^{-\kappa_Hd_1(x,y)/\eps}$, and the factor is $2L^2\gamma_0^{-1}$ when $\alpha=0$.

\emph{The slices $H_i$.} By Lemma~\ref{app:G1:KS-rec}, $H_i=G^{\mu_*}_{i+1}-G^{\mu_*}_i$ for $0\le i\le K-1$. Since $\mu_*t_{K-1}\le1$, the two cases above apply with $M$ replaced by $\mu_*$ and $\alpha=0$. This gives the second bound with $c_H=\max(C_0,2L^2\gamma_0^{-1})$.

\emph{Uniformity.} The constants of Theorem~\ref{app:G1:VSR-VS} and of Lemma~\ref{app:G1:VSR-CT} depend only on $(a,L)$. They are uniform in $\eta\in(0,1]$, $\tilde\mu\in[0,1]$, $\ell\in L\N$ and $a_k\in[a_\flat,a]$. The torus side enters only through $\ell_T/t_K\in\N$.
\end{proof}

\subsection{The remainder field of a step}\label{app:G1:Arem-sub}
In this subsection we use the units of step $k$ of \cite{Ba2}: $t=t_k$, the unit lattice $T^{(k)}_1$ carries the field $A$, the $L$-lattice carries the new block variable $B$, and $\eta=L^{-k}$. For $y$ in the $L$-lattice, $\cB(y)$ is its block of $L^d$ unit sites, and $y(x)$ is the block point with $x\in\cB(y(x))$. $\Lambda_0=\Lambda^{(k)}_0$ is a union of big blocks, and $\Lambda_0'$ is the set of $y$ with $\cB(y)\subset\Lambda_0$. We write $r=r(t)$, $r'=r(Lt)\le r$, $\abs\cdot$ for the Euclidean norm, and $\mu_*$ for the photon mass of \cite{Ba2}. By \cite[p.~608]{Ba1}, the operators for vector fields are the scalar ones with $N=d$ and external field $0$, with mass $\mu_*$. Let $c_0,\delta_0\le1$ be constants for which the zero-field bounds \cite[(2.25), (2.34), (2.36)]{Ba1} of Proposition~\ref{prop:UV-BaProp21} hold on $T_\eta$ and $T^{(k)}_1$ (no restriction near $\partial\Omega$, since $\Omega$ is the torus \cite[p.~611]{Ba1}). They depend only on $d,a,M_b$, and they do not depend on the torus side. Put
\[
S(\delta):=\sum_{z\in\Z^d}(1+\abs z_1)\mathrm{e}^{-\delta\abs z},\qquad \Theta:=(1+d)\mathrm{e}^{\delta_0d}S(\delta_0/4),\qquad K_*:=(1+a)^2(1+c_0)^2\mathrm{e}^{\delta_0d}\Theta^2(1+dL).
\]

\begin{lemma}[lattice tails]\label{app:G1:Arem-tail}
Let $\delta\in\{\delta_0,\delta_0/2\}$, $x\in\R^d$ and $\rho\ge0$. Then $\sum_{z\in\Z^d,\abs{x-z}\ge\rho}(1+\abs{x-z}_1)\mathrm{e}^{-\delta\abs{x-z}}\le\Theta\,\mathrm{e}^{-\delta\rho/2}$. On a torus the same holds for the torus distance. If $\abs{K_1(x,w)}\le\kappa_1\mathrm{e}^{-\delta_0\abs{x-w}}$ and $\abs{K_2(w,u)}\le\kappa_2\mathrm{e}^{-\delta_0\abs{w-u}}$, then $\abs{\sum_wK_1(x,w)K_2(w,u)}\le\kappa_1\kappa_2\Theta\,\mathrm{e}^{-\delta_0\abs{x-u}/2}$.
\end{lemma}
\begin{proof}
Take $y_x\in\Z^d$ with $x-y_x\in[0,1)^d$, so that $\abs{x-y_x}\le d$. Then $1+\abs{x-z}_1\le(1+d)(1+\abs{y_x-z}_1)$. For $\abs{x-z}\ge\rho$,
\[
\mathrm{e}^{-\delta\abs{x-z}}\le\mathrm{e}^{-\delta\rho/2}\mathrm{e}^{\delta d/2}\mathrm{e}^{-\delta\abs{y_x-z}/2}.
\]
Summing over $z$ gives the bound $(1+d)\mathrm{e}^{\delta d/2}S(\delta/2)\mathrm{e}^{-\delta\rho/2}\le\Theta\mathrm{e}^{-\delta\rho/2}$, because $S$ is decreasing. On a torus, choose for every point a lift at minimal distance; these lifts are distinct points of $\Z^d$. For the composition, $\abs{x-w}+\abs{w-u}\ge\abs{x-u}$ gives $\mathrm{e}^{-\delta_0(\abs{x-w}+\abs{w-u})}\le\mathrm{e}^{-\delta_0\abs{x-u}/2}\mathrm{e}^{-\delta_0\abs{x-w}/2}$. Then apply the first part with $\delta=\delta_0/2$ and $\rho=0$.
\end{proof}

\noindent\emph{Hypotheses.} On the small-field set of step $k$ we have, for $y,y'\in\Lambda_0'$ and $x\in\cB(y)$ (this is \cite[(2.17)]{Ba2} with $\eps$ replaced by $t$, \cite[pp.~560, 574]{Ba2}):
\begin{itemize}
\item[(Rs)] $\abs{B(y)-A(x)}\le2Ldp$, $\abs{B(y)}\le2p/(\mu_*t)$, and $\abs{B(y)-B(y')}\le3Ldp$ if $y,y'$ are nearest neighbours.
\end{itemize}
The geometry \cite[(2.7)--(2.8), p.~558; pp.~560, 566--567]{Ba2} gives:
\begin{itemize}
\item[(Gm)] $\Lambda_0$ is the union of the blocks $\cB(y)$, $y\in\Lambda_0'$. Fine points of $B^k(\Lambda^{(k)}_2)$ lie at distance $>2r-2\sqrt d$ from $\Lambda_0^c$ (proved below). $0\le\theta_{k+1}\le1$, and $\theta_{k+1}$ vanishes at distance $\ge M_b$ from $B^k(\Lambda^{(k)}_2)$. $\zeta^{(k)}(x,y)=1$ for $\abs{x-y}\le r/2$, $\zeta^{(k)}(x,y)=0$ for $\abs{x-y}\ge r-2M_b$, and $\abs{\zeta^{(k)}(x,y)-\zeta^{(k)}(x',y)}\le\eta$ for fine nearest neighbours $x,x'$. For $\zeta^{(k+1)}$, in the units of step $k$: $\zeta^{(k+1)}(x,y)=1$ for $\abs{x-y}\le r'/2$, and $\abs{\zeta^{(k+1)}(x,y)-\zeta^{(k+1)}(x',y)}\le\eta$ for fine nearest neighbours. Indeed, the scale-$(k+1)$ properties give the cut-off $Lr'/2\ge r'/2$ and increments $\le\eta/L$.
\end{itemize}
\noindent\emph{Proof of the distance bound.} Let $x\in\Lambda_2^{(k)}$ and $u\in\Lambda_0^c$, and recall $r>\sqrt d$ by (Sm1). Every unit site $w$ with $\dist(w,\Lambda_0^c)\le r$ lies in a large block at distance $\le r$ from $\Lambda_0^c$, hence $w\in\Lambda_1^c$ by \cite[(2.8)]{Ba2}. Since $x$ lies in a large block at distance $>r$ from $\Lambda_1^c$, we get $\abs{x-w}>r$ for every such $w$, and in particular $\abs{x-u}>r$. Let $p$ be the point of a shortest segment from $x$ to $u$ with $\abs{p-u}=r-\sqrt d/2$, and let $w$ be a unit site with $\abs{w-p}\le\sqrt d/2$. Then $\abs{w-u}\le r$, so $\abs{x-u}=\abs{x-p}+\abs{p-u}\ge\abs{x-w}-\sqrt d/2+r-\sqrt d/2>2r-\sqrt d$. A fine point of $B^k(\Lambda_2^{(k)})$ lies within $\sqrt d$ of a point of $\Lambda_2^{(k)}$.

Since $\zeta^{(k)}(y,y)=1$, summing the increments along a fine path from $y$ to $x$ gives
\begin{equation}\label{app:G1:Arem-zeta}
\abs{\zeta^{(k)}(x,y)}\le1+\abs{x-y}_1,\qquad\abs{1-\zeta^{(k)}(x,y)}\le\abs{x-y}_1,
\end{equation}
and the same holds for $\zeta^{(k+1)}$. We also assume:
\begin{itemize}
\item[(Sm1)] $r'\ge8d(L+M_b)$;
\item[(Sm2)] $2a^{-1}L^2\mu_*t\le1$;
\item[(Sm3)] $12K_*\mathrm{e}^{-\delta_0r'/(16d)}\le\mu_*t$.
\end{itemize}
Since $r(s)=R(1+\log s^{-1})^{\varrho}$ with $\varrho>1$ \cite[(2.7)]{Ba2}, (Sm1)--(Sm3) hold for all $t\le\eps_0$ once $\eps_0$ is small enough.

\begin{lemma}[the fluctuation field]\label{app:G1:Arem-fluct}
Assume (Rs), (Gm), (Sm1)--(Sm3). Let $A'=A-aL^{-2}C^{(k)}_{\Lambda_0}Q^*B$ be the field of \cite[(2.80)]{Ba2}, and let $x$ be a unit site with $\dist(x,\Lambda_0^c)\ge r/2$. Then $\abs{A'(x)}\le c_{A'}p$, where $c_{A'}:=2Ld+\frac32+3ad(1+d)c_0\Theta$.
\end{lemma}
\begin{proof}
Write $C:=C^{(k)}$ and $\delta C:=C^{(k)}_{\Lambda_0}-C$ on $\Lambda_0$ \cite[(2.31)--(2.32), (2.35)]{Ba1}, and $y:=y(x)$. At zero field $(Q^*B)(x')=B(y(x'))$. As in \cite[(2.82)]{Ba2}, this gives the exact decomposition
\begin{multline*}
aL^{-2}(C^{(k)}_{\Lambda_0}Q^*B)(x)=aL^{-2}(C1)(x)B(y)-aL^{-2}\!\!\sum_{x'\notin\Lambda_0}\!\!C(x,x')B(y)\\
+aL^{-2}\!\!\sum_{x'\in\Lambda_0}\!\!C(x,x')(B(y(x'))-B(y))+aL^{-2}(\delta C\,Q^*B)(x).
\end{multline*}
\emph{First term.} $P1=1$, and $\Delta^{(k)}1=\delta_\mu1$ with $\delta_\mu=a_k\mu_*^2t^2/(a_k+\mu_*^2t^2)\le\mu_*^2t^2$ \cite[(2.84)]{Ba2}. Hence $C1=(aL^{-2}+\delta_\mu)^{-1}$, as in \cite[(2.83), (2.85)]{Ba2}, and so $0\le1-aL^{-2}C1\le a^{-1}L^2\mu_*^2t^2$. With (Rs) and (Sm2),
\[
\abs{(1-aL^{-2}C1)B(y)}\le2a^{-1}L^2\mu_*tp\le p .
\]
\emph{Tails.} By \cite[(2.34), (2.36)]{Ba1} and Lemma~\ref{app:G1:Arem-tail}, three contributions are each $\le ac_0\Theta\,\mathrm{e}^{-\delta_0r/(8d)}\cdot4p/(\mu_*t)$:
\begin{itemize}
\item the second term, since $\abs{x-x'}\ge r/2$ there (tail with $\rho=r/2$, and $\mathrm{e}^{-\delta_0r/4}\le\mathrm{e}^{-\delta_0r/(8d)}$);
\item the last term, since $\dist(x,\Lambda_0^c)\ge r/2$ and $\abs{Q^*B}\le2p/(\mu_*t)$ on $\Lambda_0$, which gives the factor $c_0\mathrm{e}^{-\delta_0r/2}S(\delta_0)\le c_0\Theta\mathrm{e}^{-\delta_0r/(8d)}$;
\item the part of the third term with $\abs{x-x'}>\rho:=r/(4d)$, where we use $\abs{B(y(x'))-B(y)}\le4p/(\mu_*t)$.
\end{itemize}
By (Sm3), each of them is $\le p/6$.
\emph{Near part.} Let $\abs{x-x'}\le\rho$ and $y'=y(x')$. Every point $w$ of the box spanned by $\cB(y)\cup\cB(y')$ satisfies $\abs{w-x}_\infty\le\rho+L$. Hence $\abs{w-x}\le d(\rho+L)=r/4+dL<r/2$ by (Sm1), so $w\in\Lambda_0$. A monotone nearest-neighbour path from $y$ to $y'$ therefore runs in $\Lambda_0'$, and it has $\abs{y'-y}_1/L\le(\abs{x'-x}_1+dL)/L$ steps. By (Rs),
\[
aL^{-2}\abs{B(y')-B(y)}\le3adL^{-2}(\abs{x'-x}_1+dL)p\le3ad(1+d)(1+\abs{x-x'}_1)p .
\]
Summing against $c_0\mathrm{e}^{-\delta_0\abs{x-x'}}$ gives at most $3ad(1+d)c_0\Theta p$.
\emph{Conclusion.} $A'(x)=(A(x)-B(y))+(B(y)-aL^{-2}(C^{(k)}_{\Lambda_0}Q^*B)(x))$. With (Rs) the bound follows: $\abs{A'(x)}\le2Ldp+p+\frac p2+3ad(1+d)c_0\Theta p$.
\end{proof}

\begin{lemma}[smallness of the remainder field]\label{app:K:L09-1}
Let $0\le k\le K-1$ and $t=t_k$. Let $\mathcal A=A^{(k)}=a_k\zeta^{(k)}G_kQ_k^*A$ be the fine field \cite[(2.54)]{Ba2}, whose fluxes $\bar A^{(k)}_b=L^{-k}\sum_{\langle x,x'\rangle\subset b}A^{(k)}_{\langle x,x'\rangle}$ define the scalar restrictions \cite[(2.16), (2.55), p.~570]{Ba2}. Let $\tilde B$ be the real interpolated background of Lemma~\ref{app:G1:I145-realbg}(3), that is, $\tilde B=(1-\theta_{k+1})A^{(k)}+\theta_{k+1}B^{(k+1),\eta}$ on $B^k(\Lambda^{(k)}_0)$ \cite[p.~579]{Ba2}. Let $A'':=\mathcal A-\tilde B$, and assume (Rs), (Gm), (Sm1)--(Sm3). Put
\[
c'':=ac_0\mathrm{e}^{\delta_0d}\Theta\,c_{A'}+\tfrac12 ,
\]
which depends only on $d,L,a,M_b$. Then $\abs{A''_\mu(z)}\le c''p(t)$ at every fine site $z$ of $B^k(\Lambda^{(k)}_0)$ and for every $\mu$. Consequently $\abs{A''(b)}\le c''p(t)$ for every bond $b$ of the step lattice in $\Omega'\subset B^k(\Lambda^{(k)}_0)$.
\end{lemma}
\begin{proof}
\emph{The identity.} Pointwise, $A''=A^{(k)}-(1-\theta_{k+1})A^{(k)}-\theta_{k+1}B^{(k+1),\eta}=\theta_{k+1}A^{(k)}-\theta_{k+1}B^{(k+1),\eta}$. By \cite[(2.88)--(2.89), p.~575]{Ba2},
\[
A''=\theta_{k+1}A'^{(k)}+T_2+T_3+T_4,
\]
where $\theta_{k+1}A'^{(k)}=a_k\theta_{k+1}\zeta^{(k)}G_kQ_k^*A'$ and $T_2,T_3,T_4$ are the second, third and fourth terms of \cite[(2.89)]{Ba2}, whose sum is $\tilde B^{(k)}$. If $\theta_{k+1}(z)=0$, then $A''(z)=0$. So let $\theta_{k+1}(z)>0$. Then $\dist(z,\Lambda_0^c)>2r-2\sqrt d-M_b$ by (Gm).

\emph{Kernels.} At zero field $(Q_k^*f)(z)=f(y')$ for $z$ in the unit block of $y'$, and $a_k\le a$. By \cite[(2.25)]{Ba1}, $\abs{(a_kG_kQ_k^*)(z,y')}\le ac_0\mathrm{e}^{\delta_0d}\mathrm{e}^{-\delta_0\abs{z-y'}}$.

\emph{The main term.} If $\zeta^{(k)}(z,y')\ne0$, then $\abs{z-y'}<r-2M_b$, so $\dist(y',\Lambda_0^c)>r+M_b-2\sqrt d\ge r/2$, because $r\ge4\sqrt d$ by (Sm1). Lemma~\ref{app:G1:Arem-fluct} then gives $\abs{A'(y')}\le c_{A'}p$. By \eqref{app:G1:Arem-zeta}, $0\le\theta_{k+1}\le1$ and Lemma~\ref{app:G1:Arem-tail} with $\rho=0$,
\[
\abs{\theta_{k+1}A'^{(k)}(z)}\le ac_0\mathrm{e}^{\delta_0d}\Theta c_{A'}p .
\]

\emph{$T_2=a_kaL^{-2}\theta_{k+1}\zeta^{(k)}G_kQ_k^*\delta C\,Q^*B$.} For the intermediate unit site $w$ with $\zeta^{(k)}(z,w)\ne0$ we have $\abs{z-w}<r-2M_b$, hence $\dist(w,\Lambda_0^c)\ge r/2$ as in the main term. So by \cite[(2.36)]{Ba1}, (Rs) and Lemma~\ref{app:G1:Arem-tail},
\[
\abs{T_2(z)}\le a^2c_0^2\mathrm{e}^{\delta_0d}\Theta^2\mathrm{e}^{-\delta_0r/2}\cdot2p/(\mu_*t).
\]

\emph{$T_3=-a_kaL^{-2}\theta_{k+1}(1-\zeta^{(k)})G_kQ_k^*CQ^*\Lambda_0'B$.} By \eqref{app:G1:Arem-zeta}, $\abs{1-\zeta^{(k)}(z,x')}\le\abs{z-x'}_1$, and it vanishes for $\abs{z-x'}\le r/2$. So $\abs{T_3(z)}\le a^2c_0^2\mathrm{e}^{\delta_0d}\Theta^2\mathrm{e}^{-\delta_0r/4}\cdot2p/(\mu_*t)$.

\emph{$T_4=a_{k+1}L^{-2}\theta_{k+1}(1-\zeta^{(k+1)})G^\eta_{k+1}Q^*_{k+1}\Lambda_0'B$.} Take the adjoint of \cite[(2.41)]{Ba1} at zero field and rescale it as in \cite[(2.89)]{Ba2}. This gives $a_{k+1}L^{-2}G^\eta_{k+1}Q^*_{k+1}=a_kaL^{-2}G_kQ_k^*C^{(k)}Q^*$. By Lemma~\ref{app:G1:Arem-tail}, the kernel of $G_kQ_k^*C^{(k)}$ is bounded by $c_0^2\mathrm{e}^{\delta_0d}\Theta\,\mathrm{e}^{-\delta_0\abs{z-x''}/2}$. For $x''\in\cB(y)$ we have $\abs{1-\zeta^{(k+1)}(z,y)}\le\abs{z-y}_1\le(1+dL)(1+\abs{z-x''}_1)$. This vanishes unless $\abs{z-x''}>r'/2-dL\ge r'/4$, by (Sm1). Hence
\[
\abs{T_4(z)}\le a^2c_0^2\mathrm{e}^{\delta_0d}\Theta^2(1+dL)\,\mathrm{e}^{-\delta_0r'/16}\cdot2p/(\mu_*t).
\]

\emph{Conclusion.} Each of $T_2,T_3,T_4$ is $\le K_*\mathrm{e}^{-\delta_0r'/(16d)}\cdot2p/(\mu_*t)\le p/6$ by (Sm3). So $\abs{A''_\mu(z)}\le c''p$. A bond $b$ of the step lattice consists of $L^k$ fine bonds of length $\eta$, so $\abs{A''(b)}\le\max_z\abs{A''_\mu(z)}\le c''p$.
\end{proof}

\subsection{The domain of the step and the bad vacuum localisations}\label{app:G1:L18-2sec}
Fix $1\le k<K$ and work in scale-$k$ units on $T_\eps$. Here $\cV_k$ is the sum of the terms of top index $k$ of Definition~\ref{app:G1:vacgrade}. These are the vacuum counterterm pieces $E_{K,k}-E_{K,k+1}$ that are moved at the step from $k$ to $k+1$ \cite[(1.27)--(1.28), p.~418; p.~422]{Ba3}.

We localise every vertex as in \cite[p.~420]{Ba3}. A vertex (1.6) or (1.7) is a sum over unit cubes of its position. For every other vertex, each vector leg carries a function from a smooth partition of unity whose support lies in a cube of side $2$. For each function of the partition we fix once and for all one unit cube $\square$ of the scale-$k$ lattice that meets its support, and call it the label of the function. The label depends only on the function, not on the term. A \emph{vertex piece} is a vertex together with labels for all of its localised legs. A localisation $\Lambda$ of a term is the family of labels of its vertex pieces. We write $\Lambda\subset W$ if all labels lie in $W$, and $\Lambda\cap W\ne\emptyset$ if some label does. By multilinearity, $\cV_k=\sum_\Lambda F(\Lambda)$, where $F(\Lambda)$ is the sum of the terms of $\cV_k$ with localisation $\Lambda$. The tree length $d(\Lambda)$ of \cite[(1.33)]{Ba3} is computed over the labels.

\begin{lemma}[the domain of the step]\label{app:G1:L18-2dom}
Put $\Omega_1:=B^k(\Lambda_7^{(k-1)\prime})$ \cite[(1.4), p.~412]{Ba3} and $W:=B^k(\Lambda_7^{(k)})$. Then:
\begin{enumerate}
\item[(a)] $W\subset\Omega_1$, and $\Omega_1\setminus B^k(\Lambda_7^{(k-1)\prime}\cap\Lambda_7^{(k)c})=W$;
\item[(b)] if $x$ lies in the support of a localisation function whose label lies in $W$, then $x\in\Omega_1$ and $\dist(x,\partial\Omega_1)\ge r(t_k)-4\sqrt3-1$;
\item[(c)] if $R\ge R^{**}$, the function $g_k$ of \cite[p.~412]{Ba3} equals $1$ on all such supports.
\end{enumerate}
\end{lemma}
\begin{proof}
(a) Let $\Lambda_0^{(k)},\dots,\Lambda_7^{(k)}$ be the sets of \cite[(2.7)--(2.8), p.~558]{Ba2}, built with $r(t_k)$ \cite[p.~566]{Ba2}. By (2.8), $\Lambda_{i+1}^{(k)c}$ is the union of the large blocks at distance $\le r(t_k)$ from $\Lambda_i^{(k)c}$. The set $\Lambda_0^{(k)}$ is a union of large blocks \cite[p.~570]{Ba2}. By induction on $i$, every $\Lambda_i^{(k)c}$ is a union of large blocks, each at distance $0$ from $\Lambda_i^{(k)c}$. So the union contains $\Lambda_i^{(k)c}$, and $\Lambda_{i+1}^{(k)}\subset\Lambda_i^{(k)}$. Admissibility gives $\Lambda_0^{(k)}\subset\Lambda_7^{(k-1)\prime}$ \cite[p.~566, p.~570]{Ba2}. Hence $\Lambda_7^{(k)}\subset\Lambda_6^{(k)}\subset\Lambda_0^{(k)}\subset\Lambda_7^{(k-1)\prime}$. The operation $B^k$ maps disjoint sets of points of $T_1^{(k)}$ to disjoint unions of unit cubes. Therefore
\[
\Omega_1\setminus B^k(\Lambda_7^{(k-1)\prime}\cap\Lambda_7^{(k)c})=B^k(\Lambda_7^{(k-1)\prime}\cap\Lambda_7^{(k)})=W .
\]
(b) Let $\square=B^k(y)\subset W$ with $y\in\Lambda_7^{(k)}$. Let $y$ lie in the large block $b\subset\Lambda_7^{(k)}$. Let $x'\notin\Omega_1$, say $x'\in B^k(y')$ with $y'\notin\Lambda_7^{(k-1)\prime}$. By (a), $y'\in\Lambda_0^{(k)c}\subset\Lambda_6^{(k)c}$. By (2.8), $\dist(b,\Lambda_6^{(k)c})>r(t_k)$, so $\abs{y-y'}>r(t_k)$. A unit cube has diameter $\sqrt3$. A point $x$ of a side-$2$ cube that meets $\square$ satisfies $\abs{x-y}\le\sqrt3+2\sqrt3$. Also $\abs{x'-y'}\le\sqrt3$. Hence $\abs{x-x'}>r(t_k)-4\sqrt3$. Every point of $\partial\Omega_1$ is within one lattice spacing $\eta\le1$ of $\Omega_1^c$, and this gives (b). The same argument holds with torus distances.
(c) By \cite[p.~412]{Ba3}, $g_k(x)=1$ whenever $\dist(x,\partial\Omega_1)\ge M_b$, where $M_b$ is the block-size constant $M$. Since $t_k\le\eps_0\le1$, we have $r(t_k)\ge R$. Moreover $R^{**}\ge4\ell_{\Nbar}+32M_b+25$, so $r(t_k)-4\sqrt3-1\ge M_b$. Now apply (b).
\end{proof}

By \cite[p.~412]{Ba3} and \cite[(2.57), p.~570]{Ba2}, the first operation of the step from $k$ to $k+1$ removes the interaction from $B^k(\Lambda_7^{(k-1)\prime}\cap\Lambda_7^{(k)c})$. That is, it deletes every vertex piece with a label in this set. What remains is $\mathcal P^{(k)}(\Lambda_7^{(k)},\cdot)$, whose vertex sums are restricted to $W$ by Lemma~\ref{app:G1:L18-2dom}(a). By Lemma~\ref{app:G1:L18-2dom}(c), the coupling $e\,g_k$ of a remaining piece equals $e$ on all of its legs. The fluctuation integrated at this step is the operator of index $k$ in \cite[(1.27), p.~418]{Ba3}.

\begin{definition}[the generated partners]\label{app:G1:L18-2def}
Let $U$ be the bare interaction of Definition~\ref{app:G1:vacgrade}, with the coefficients $\delta m^2$ kept as numbers. Let $U_W$ be the sum of its vertex pieces all of whose labels lie in $W$. For a covariance $\Gamma$ let $c_\Gamma(U_W)$ be the part of $\log\E_{\mu_\Gamma}\mathrm{e}^{-U_W}$ of orders $1\le\alpha+\beta\le\nbar$, with the order convention of Definition~\ref{app:G1:vacgrade}. Let $\Gamma_{<k}$ and $\Gamma_{\le k}$ be the sums of the slices of index $<k$ and $\le k$ of the lines of $\cV_k$. The \emph{generated partners} at the step from $k$ to $k+1$ are
\[
G_{k,W}:=c_{\Gamma_{\le k}}(U_W)-c_{\Gamma_{<k}}(U_W).
\]
\end{definition}

\begin{lemma}[the bad vacuum localisations at the last scale]\label{app:K:L18-2}
\label{app:G1:L18-2id}
Let $1\le k<K$, $R\ge R^{**}$, $W=B^k(\Lambda_7^{(k)})$, and $\cV_k=\sum_\Lambda F(\Lambda)$ as above. Then:
\begin{enumerate}
\item[(a)] $G_{k,W}=\sum_{\Lambda\subset W}F(\Lambda)$;
\item[(b)] $\cV_k-G_{k,W}=E^{\rm bad}_{K,k}+E^{\rm out}_{K,k}$, where
\[
E^{\rm bad}_{K,k}:=\sum_{\Lambda:\ \Lambda\cap W\ne\emptyset,\ \Lambda\not\subset W}F(\Lambda),\qquad E^{\rm out}_{K,k}:=\sum_{\Lambda:\ \Lambda\cap W=\emptyset}F(\Lambda);
\]
\item[(c)] the part of the localisations of $\cV_k$ meeting $W$ that is not cancelled by the generated partners,
\[
\sum_{\Lambda\cap W\ne\emptyset}F(\Lambda)-G_{k,W},
\]
equals $E^{\rm bad}_{K,k}$.
\end{enumerate}
\end{lemma}
\begin{proof}
(a) The proof of Lemma~\ref{app:G1:vacgrade-id}(ii) uses only one property of the interaction: it is a polynomial in the fields whose coefficients have order $\ge1$. This holds for $U_W$. So $c_{\Gamma_{\le k}}(U_W)$ is the sum of the connected vacuum graphs of $U_W$ of order $\le\nbar$ with covariance $\Gamma_{\le k}$, and similarly for $\Gamma_{<k}$. We have $\Gamma_{\le k}=\Gamma_{<k}+C_k$, where $C_k$ is the slice of index $k$, and $\Gamma_{<k}$ is the sum of the slices of index $<k$. The value of a graph is multilinear in its line kernels. Hence the difference consists of the graphs with every line of index $\le k$ and at least one line of index $k$. These are the graphs of $\cV_k$. The value of a graph is also multilinear in its vertex factors. Every vertex of $U$ is the sum of its vertex pieces \cite[p.~420]{Ba3}, and $U_W$ keeps exactly the pieces with all labels in $W$. So each graph of $U_W$ is the sum, over the localisations $\Lambda\subset W$, of the corresponding terms of the same graph of $U$. Summing over graphs and index assignments gives $\sum_{\Lambda\subset W}F(\Lambda)$.

(b) Every $\Lambda$ falls into exactly one of three classes: $\Lambda\subset W$; $\Lambda\cap W\ne\emptyset$ and $\Lambda\not\subset W$; or $\Lambda\cap W=\emptyset$. Subtract (a) from $\cV_k=\sum_\Lambda F(\Lambda)$.

(c) A localisation is non-empty, so $\{\Lambda\subset W\}\subset\{\Lambda\cap W\ne\emptyset\}$. Now apply (a).
\end{proof}

\begin{remark}[how the lemma is used]\label{app:G1:L18-2rem}
Here is how the lemma connects to the step of \cite{Ba2,Ba3}.
\begin{itemize}
\item By Lemma~\ref{app:G1:L18-2dom} and the paragraph after it, $U_W$ is the vertex content of the interaction on which the step acts. Lemma~\ref{app:K:L18-2} therefore says that the set $W$ of Lemma~\ref{app:G1:VS} in the proof of Proposition~\ref{app:G1:I9} is $B^k(\Lambda_7^{(k)})$, the set named in \cite[p.~422]{Ba3}. The bad terms of Proposition~\ref{app:G1:I9} are exactly $E^{\rm bad}_{K,k}$.
\item The localisations in $E^{\rm out}_{K,k}$ have no label in the domain of the step. They form the outside class of Proposition~\ref{app:G1:R9vac}(ii), which Lemma~\ref{app:G1:IF1abs} absorbs.
\item The actual step uses lines with domain-dependent operators and mass vertices whose coefficients are localised \cite[(1.29), p.~419]{Ba3}. These differ from $G_{k,W}$ by terms of two kinds. The first are the differences of operators, treated as external fields \cite[p.~419]{Ba3}. The second are the mass counterterms with bad localisations \cite[p.~422]{Ba3} (Proposition~\ref{app:G1:H3-bad}). Neither kind is a term of $\cV_k$ in the grading of Definition~\ref{app:G1:vacgrade}, so neither enters $E^{\rm bad}_{K,k}$. Both are field-independent contributions of the step, and they are bounded in the results cited, not by this lemma.
\end{itemize}
\end{remark}

\subsection{A constant shift of the background}\label{app:G1:cshift}

Fix $k\ge1$, put $t:=t_k$ and $\eta:=L^{-k}$, and work in scale-$k$ units on the torus $T_\eta$ of side $\ell$. For a vector field $A$ and a lattice contour $\Gamma$ put $A(\Gamma):=\sum_{b\subset\Gamma}\eta A_b$, as in \cite[p.~412]{Ba3}. Let $U(\cdot)$ be the unitary representation used in the covariant differences of \cite[(1.32)]{Ba3}, so that $(D^\eta_{A,\mu}f)(x)=\eta^{-1}\big(U(A(b))f(x+\eta e_\mu)-f(x)\big)$ for the bond $b=\langle x,x+\eta e_\mu\rangle$. A constant $c\in\R^3$ is the vector field with value $c_\mu$ on every bond in direction $\mu$. A box of side $s$ is a set $\prod_\mu\{a_\mu,a_\mu+\eta,\dots,a_\mu+s-\eta\}$, taken together with the bonds between its sites. Put $\theta(A):=e_k\norm{\partial^\eta A}_\infty$ and $\sigma:=e_k(\norm{\tilde A}_\infty+\norm{\partial^\eta\tilde A}_\infty)$. For a background $\tilde B$ and an external vector field $\tilde A$ consider the conditions
\begin{itemize}
\item[(CS1)] $\abs{\tilde A}$, $\abs{\partial^\eta\tilde A}$, $\abs{\partial^\eta\tilde B}$ and the $\alpha_0$-H\"older quotients of $\partial^\eta\tilde A$ and $\partial^\eta\tilde B$ are $\le O(1)p(t)$, and $\dist(\supp\tilde A,\partial\Omega)>2r(t)$ \cite[p.~412]{Ba3};
\item[(CS2)] $\tilde B$ and $\tilde A$ are real, $\theta(\tilde B),\theta(\tilde A+\tilde B)\le\vartheta_R$, $\sigma\le1$ and $0\le m^2t^2\le c$. This is the regularity hypothesis of Lemma~\ref{app:G1:WKinputs} and Propositions~\ref{app:G1:sandwich} and~\ref{app:G1:twodomain}.
\end{itemize}

\begin{lemma}[the conditions are invariant under constant shifts]\label{app:G1:cshift-hyp}
If $(\tilde B,\tilde A)$ satisfies (CS1) and (CS2), then so does $(\tilde B+c,\tilde A)$ for every $c\in\R^3$, with the same constants.
\end{lemma}
\begin{proof}
$\partial^\eta_\nu c=0$, so $\partial^\eta(\tilde B+c)=\partial^\eta\tilde B$. Hence the H\"older quotients of $\partial^\eta(\tilde B+c)$ are those of $\partial^\eta\tilde B$, and $\theta(\tilde B+c)=\theta(\tilde B)$, $\theta(\tilde A+\tilde B+c)=\theta(\tilde A+\tilde B)$. The quantities $\tilde A$, $\sigma$, $\Omega$ and $m^2t^2$ do not change. Neither (CS1) nor (CS2) contains $\tilde B$ in any other way.
\end{proof}

\begin{lemma}[the shift as a twist]\label{app:G1:cshift-twist}
Let $Y\subset T_\eta$ be a box of side $<\ell$. Then there is a fundamental domain $F=\prod_\mu\{a_\mu,\dots,a_\mu+\ell-\eta\}\supset Y$ whose cut bonds $\langle a_\mu+\ell-\eta,a_\mu\rangle$ are not bonds of $Y$. Put $g_c(x):=U(c\cdot(x-a))$ for $x\in F$ and $\vartheta_\mu:=U(\ell c_\mu)=U(c(\Gamma_\mu))$, where $\Gamma_\mu$ is the closed contour that winds once in direction $\mu$. For a background $\tilde B$, let $\mathcal O^\vartheta(\tilde B)$ denote the operator $\mathcal O(\tilde B)$ built with the transporter $U(\tilde B(b))\vartheta_\mu$ on every cut bond $b$ in direction $\mu$. Let $\mathcal O$ be any operator obtained by sums, products, adjoints and inverses from the bond transports $f\mapsto U(\tilde B(b))f(\cdot+\eta e_\mu)$, from multiplication operators by scalar functions or by values of $U$, and from the evaluation $E$ of functions on $T_\eta$ at the points of a sublattice of $T_\eta$ (in particular the unit lattice of the blocks), together with its adjoint $E^*$. This includes $D^\eta_{\tilde B}$, $\Delta^\eta_{\tilde B}$, $Q_k(\tilde B)$, $P_k(\tilde B)$, $G_k(\Omega,\tilde B)$ and $G_k(\Omega,\tilde B)Q_k^*(\tilde B)$. Here and below $M_{g_c}$ acts on functions on a sublattice as multiplication by the restriction of $g_c$. Then
\[
\mathcal O(\tilde B+c)=M_{\bar g_c}\,\mathcal O^\vartheta(\tilde B)\,M_{g_c},
\]
where $M_g$ is multiplication by $g$. Moreover, on $Y$ we have $g_c(x')\overline{g_c(x)}=U(c(\Gamma))$ for every contour $\Gamma\subset Y$ from $x$ to $x'$.
\end{lemma}
\begin{proof}
The projection of $Y$ on the axis $\mu$ consists of fewer than $\ell/\eta$ consecutive sites. Let $a_\mu$ be the first of them. Then $a_\mu+\ell-\eta\notin Y$, so the cut bond in direction $\mu$ is not a bond of $Y$.

Let $b=\langle x,x+\eta e_\mu\rangle$ be a bond. If $b$ is not cut, then $g_c(x+\eta e_\mu)\overline{g_c(x)}=U(\eta c_\mu)=U(c(b))$. If $b$ is cut, then the $\mu$-th coordinate of $x-a$ is $\ell-\eta$ and that of $x+\eta e_\mu-a$ is $0$, so $g_c(x+\eta e_\mu)\overline{g_c(x)}=U(-(\ell-\eta)c_\mu)=U(c(b))\overline{\vartheta_\mu}$. Since $U$ is abelian and unimodular, in both cases
\[
U((\tilde B+c)(b))f(x+\eta e_\mu)=\overline{g_c(x)}\,U(\tilde B(b))\vartheta_\mu^{[b\ \mathrm{cut}]}\,(g_cf)(x+\eta e_\mu).
\]
So every bond transport at $\tilde B+c$ equals $M_{\bar g_c}(\cdot)^\vartheta M_{g_c}$ of the bond transport at $\tilde B$. Multiplication operators by scalar functions or by values of $U$ commute with $M_{g_c}$, because $U$ is abelian. Moreover $EM_{g_c}=M_{g_c}E$ and $M_{g_c}E^*=E^*M_{g_c}$, with $M_{g_c}$ on the sublattice as in the statement. Finally, $M_{g_c}$ is unitary. So conjugation by $M_{g_c}$ respects sums, products, adjoints and inverses, and it commutes with $E$ and $E^*$. On a contour in $Y$ no bond is cut, and the product of the factors $U(c(b))$ telescopes.
\end{proof}

\begin{proposition}[uniformity of Proposition~1 of \cite{Ba3} in a constant background]\label{app:K:L10-2}\label{app:G1:cshift-main}
Let $k\ge1$. Let $\mathfrak B$ be a real background on $T_\eta$ at which the hypotheses of \cite[Prop.~1]{Ba3} hold in the form of Proposition~\ref{app:G1:I3}(i). That is, (CS1) and (CS2) hold with $\tilde B=\mathfrak B$, together with conditions that do not involve the background (domains, couplings, masses, and the conditions of Appendix~\ref{app:B}, among them (T.nw) and (E.r)). Let $c\in\R^3$. Then:
\begin{enumerate}
\item[(i)] these hypotheses hold at $\mathfrak B+c$ with the same constants;
\item[(ii)] let a localised family of order at most $\Nbar$ of the expansion of \cite[p.~414]{Ba3} at the background $\mathfrak B+c$ have all its localisation cubes in a box $Y$ of side $<\ell$. Then \cite[(1.33)]{Ba3} holds for it, with the norms of the external scalar fields taken at $\mathfrak B+c$ as in \cite[(1.32)]{Ba3}, and with the same tree rate $\beta_P$ and the same constant $O(1)$ as at $\mathfrak B$. In particular, $\beta_P$ and $O(1)$ do not depend on $c$;
\item[(iii)] by Lemma~\ref{app:G1:cshift-twist}, the family in (ii) equals the family at $\mathfrak B$ with every scalar operator replaced by its $\vartheta$-twisted form, with the kernels conjugated by $g_c$ and with the external scalar fields multiplied by $g_c$. So (ii) is the same statement for the scalar boundary condition twisted by $\vartheta$.
\end{enumerate}
\end{proposition}
\begin{proof}
(i) is Lemma~\ref{app:G1:cshift-hyp}. (iii) follows from Lemma~\ref{app:G1:cshift-twist} applied to every line, vertex and transporter of the family. The vector lines and $\tilde A$ do not depend on the background.

(ii) Write $\tilde B:=\mathfrak B$, and recall how \cite[Sect.~3]{Ba3} proves (1.33) at a background $\tilde B$. We check each step at $\tilde B+c$.

\emph{Step 1 (vertices).} The vertices of \cite[(1.6)--(1.18)]{Ba3} depend on the scalar background only through bond transporters and through the operators $G_k(\Omega,\cdot)$ and $Q_k(\cdot)$. The counterterms $\delta m^2(e',g_k,\lambda',\Omega_1,x)$ and $E_1(e',g_k,\lambda',\Omega_1)$ of \cite[(1.4), p.~412]{Ba3} contain no background at all.

\emph{Step 2 (lines that become external fields).} By \cite[p.~432]{Ba3}, the operators $\delta G_k$ and (1.16), and every line joining vertices with $\dist(\square(v),\square(v'))\ge1$, are treated as external fields. They are bounded by (3.1), $\norm{hG_k(\Omega,\tilde B)h'}_{1,\alpha}\le O(1)\exp(-\delta_0\dist(\square(v),\square(v')))$, and by \cite[(2.5)]{Ba3}. Both reduce (1.33) to (3.2), and ``the exponential factor in (1.33) is obtained from the estimates of external fields in (3.2)'' (ibid.). On $T_\eta$ these bounds are given by Lemma~\ref{app:G1:WKinputs} and Propositions~\ref{app:G1:sandwich} and~\ref{app:G1:twodomain}, and by the bounds derived from them. Their constants depend on the background only through (CS2), and wrapping of $\Omega$ is allowed. By (i), they hold at $\tilde B+c$ with the same constants and the same decay rate. So the rate $\beta_P$ and the factors that these bounds contribute to $O(1)$ are the same at $\tilde B$ and at $\tilde B+c$.

\emph{Step 3 (the cube $\square$).} Each connected component is localised in a cube of size $O(\nbar)$. Around it one takes cubes $\square_1$ and $\square$ of sizes $r(t)$ and $3r(t)$, which are unions of big blocks, and writes $G_k(\Omega,\tilde B)=G_k(\square,\tilde B)+\delta G_k(\Omega,\square,\tilde B)$ \cite[p.~433]{Ba3}. The external field $\delta G_k(\Omega,\square,\tilde B+c)$ obeys the bound of Proposition~\ref{app:G1:twodomain} with the same constant, by (i). By Lemma~\ref{app:G1:nowrap}, which uses (T.nw) and (E.r), the cube $\square$ lies in a box of side $<\ell$, so it lifts isometrically.

\emph{Step 4 (constant part and gauge transformation).} Ba{\l}aban chooses a constant field $B_0$ with $\abs{\tilde B-B_0}\le O(r(t)p(t))\le O(p(t)^2)$ on $\square$. This uses only $\norm{\partial^\eta\tilde B}_{0,\alpha}\le O(p(t))$. He puts $B':=\tilde B-B_0$ and expands in $B'$ \cite[p.~433]{Ba3}. At $\tilde B+c$ we choose the constant $B_0+c$. Then $B'$ is the same field, with the same bound. The expressions of the renormalised classes on $\square$ are gauge invariant ``with respect to gauge transformations of the field $B_0$, if the external scalar fields are simultaneously transformed''. The gauge transformation that removes $B_0$ gives ``the same expressions as above with $B_0=0$ only'' (ibid.). This step uses only that $B_0$ is constant on the lifted cube $\square$. By Lemma~\ref{app:G1:gaugetriv}, the constant $B_0+c$ is a gradient on $\square$, so the same transformation, now with $B_0+c$, gives exactly the same expressions in $B'$ and in the external fields. The external scalar fields are transformed by the phase of $B_0+c$ instead of that of $B_0$. Moreover, for every constant $B_0$ the covariant norm \cite[(1.32)]{Ba3} at $B_0+B'$ of a field on $\square$ is the norm at $B'$ of the gauge-transformed field, by the last assertion of Lemma~\ref{app:G1:cshift-twist} with $c$ replaced by $B_0$. So these norms enter at $\tilde B+c$ in exactly the way the norms at $\tilde B$ enter at $c=0$.

\emph{Step 5 (the background-free estimate).} After the substitution $G_k(\square,0)=G_k(0)+\delta G_k(\square,\eta\Z^d,0)$ on the lifted cube, the propagator is $G_k(0)$ on $\eta\Z^d$. The remaining expressions contain no scalar background other than the external vector field $B'$. They are bounded by \cite[(3.3)]{Ba3}, $\abs{E(G_{\rm ren},\{\square(v)\},\Phi_{\rm ext},A_{\rm ext})}\le O(1)\norm{\Phi_{\rm ext}}_{1,\alpha_0}\norm{A_{\rm ext}}_{1,\alpha_0}$, which is proved on \cite[pp.~433--445]{Ba3}. Neither these expressions nor this constant depend on $c$. The torus enters only through the external fields of Steps 2 and 3 (Proposition~\ref{app:G1:I2}).

\emph{Step 4$'$ (classes without old vertices).} The $B'$-expansion also creates classes that contain no old vertex. These are chains of propagators $G_k(\square,B_0)$ with vertices (1.8)--(1.11), (1.13)--(1.15) and external vector field $B'$ only, and for them (3.2) is ``a consequence of Proposition~2.1'' \cite[p.~433]{Ba3}. At $\tilde B+c$ the chain is built from $G_k(\square,B_0+c)$ with the same $B'$. The cube $\square$ lifts (Step 3), and $B_0+c$ and $B_0$ are gradients on $\square$ (Lemma~\ref{app:G1:gaugetriv}). So both chains are gauge-equivalent to the chain at $B_0=0$; only the phases of their external scalar legs differ. By the norm identity of Step 4, the bound (3.2) at $\tilde B+c$ is the bound at $\tilde B$, with the same constant.

Steps 1--5 and 4$'$ are the whole route from (3.3) to (1.33) \cite[pp.~432--433]{Ba3}, and they are the route of Proposition~\ref{app:G1:I145-prop1base} on $T_\eta$. At each step, the data at $\tilde B+c$ either coincide with the data at $\tilde B$ (vertices, $B'$, the expressions of Step 5) or obey the same bounds with the same constants (Steps 2 and 3). So (3.2), and hence (1.33), holds at $\tilde B+c$ with the same $\beta_P$ and the same $O(1)$. This agrees with \cite[p.~421]{Ba3}: the constant of (1.33) ``is independent of $\eps$, $k$, the domains $\Omega,\Omega_1,\Omega_2$, the vector field $\tilde B$ (if they satisfy the conditions mentioned previously)''. The box $Y$ is used only in (iii).
\end{proof}

\begin{remark}\label{app:G1:cshift-rem}
Together with Proposition~\ref{app:G1:AAext}, Proposition~\ref{app:K:L10-2} gives hypothesis (ZM) of Proposition~\ref{app:G1:I3c-AA}. The constant mode of the photon field is a constant $c$ added to the background. By Proposition~\ref{app:K:L10-2}(ii),(iii), every bound on scalar operators used in \cite[Prop.~1]{Ba3} is uniform in $c$, and hence in the twist $\vartheta$, for families whose localisation cubes lie in a box of side $<\ell$.
\end{remark}

\subsection{Ba{\l}aban's Proposition~1 with coloured vector lines}\label{app:G1:prop1col}

We use the notation of \cite[pp.~411--424]{Ba3}, with $\eta=L^{-k}$. By \cite[p.~414]{Ba3} the vector field $A'$ has covariance $G_k$. This covariance lives on the torus and does not depend on $\Omega$ or on $\tilde B$. By \cite[(2.6), p.~424]{Ba3} and the sentence after it, $G_k=\sum_{j<k}G^\eta_{(j)}$ with vector slices $G^\eta_{(j)}$. By \cite[p.~426]{Ba3}, these slices obey \cite[(2.10)--(2.12)]{Ba3}.

On the torus with photon covariance $C_{\mu_*}+u\Gamma^B$, we use the colour expansion of the paragraph before Lemma~\ref{app:G1:WP-deg}. A line of colour $C$ carries $G_k$. In scale-$k$ units, a line of colour $D$ carries $uc_kD^\eta_k$, where $D^\eta_k:=\sum_{j<k}(L^j\eta)^2G^\eta_{(j)}$. The mass counterterm is built from graphs with $C$-lines only, and the partner operations act only on divergent subgraphs with $C$-lines only. A coloured renormalised family is a family in the sense of Lemma~\ref{app:G1:WP-fam}, so all its members have the same $(m,d_v,d_s)$. On $T_\eps$ every line has colour $C$, i.e.\ $m=0$. We write $E^\natural(F,\dots)$ for the value of a family $F$ in which every $D$-line carries $D^\eta_k$ instead of $uc_kD^\eta_k$.

\begin{lemma}[the colour factor]\label{app:G1:prop1col-exp}
For every coloured renormalised family $F$ with $m$ $D$-lines, every localisation and all external fields,
$E(F,\{\square(v)\},\Phi_{\rm ext},A_{\rm ext})=(uc_k)^m\,E^\natural(F,\{\square(v)\},\Phi_{\rm ext},A_{\rm ext})$.
\end{lemma}
\begin{proof}
At fixed localisation, the value of a graph is a finite sum over the vertex positions of a product of vertex factors, external fields and line kernels. It is therefore linear in the kernel of each line separately. Replacing the kernel $uc_kD^\eta_k$ of each $D$-line by $D^\eta_k$ multiplies the value by $uc_k$ once per $D$-line. By Lemma~\ref{app:G1:WP-fam}, every member of $F$ has the same number $m$ of $D$-lines, and the coefficients of the members do not depend on $u$. Summing over the members gives the claim.
\end{proof}

\begin{lemma}[$D$-lines]\label{app:G1:prop1col-deg}
(a) The kernels $(L^j\eta)^2G^\eta_{(j)}$, $j<k$, obey \cite[(2.10)--(2.12), p.~426]{Ba3} with the powers $(L^j\eta)^{-d+2}$, $(L^j\eta)^{-d+2-1-\alpha}$, $(L^j\eta)^{-d+3}$ multiplied by $(L^j\eta)^2\le1$. The exponential factors are the same, and the bounds hold at all points of the torus.

(b) Let $D^\natural$ be the degree of \cite[(2.1)--(2.3)]{Ba3}, computed with dimension $(-d+4)/2$ for each leg of a $D$-line, and let $D$ be the degree with the colours forgotten. Then $D^\natural=D+2m$. Let $G$ be a connected coloured graph built from the vertices \cite[(1.6)--(1.8), (1.10)]{Ba3} and from mass vertices of divergent self-energies. If $G$ is a vacuum graph with $m\ge2$, or with $m=1$ and $o\ge2$, then $D^\natural(G)\ge1$. If $G$ has an external leg and $m\ge1$, then $D^\natural(G)\ge1$.
\end{lemma}
\begin{proof}
(a) Multiply the slice bounds of \cite[p.~426]{Ba3} by $(L^j\eta)^2$. Since $j<k$ and $\eta=L^{-k}$, we have $L^j\eta\le L^{-1}<1$. By \cite[p.~414]{Ba3}, the fields $A'_j$ are independent Gaussian fields with covariances $C^{(j),L^j\eta}$, and these are defined without reference to $\Omega$ or $\tilde B$. The vector slices of \cite[(2.6), p.~424]{Ba3} are built from them, and \cite[p.~426]{Ba3} states their bounds (``the same inequalities hold for vector field propagators'') with no condition on the distance to $\Omega^c$.

(b) Both legs of a line carry the same index $j$, so by (a) each $D$-leg has dimension $(-d+4)/2=(-d+2)/2+1$. This gives $D^\natural=D+2m$. For the stated vertex types, Lemma~\ref{app:G1:H3-deg} gives $D=o+\frac12(E_s+E_v)+E_\partial-3$. Each internal vector line uses two factors $e$, so $o\ge I_v\ge m$.

For a vacuum graph, $D^\natural=o-3+2m$. This is $\ge1$ if $m=1$ and $o\ge2$, and $\ge3m-3\ge3$ if $m\ge2$.

Now let $G$ have an external leg and $m\ge1$, so $o\ge1$ and $E_s+E_v\ge1$. If $o=1$, then $\alpha=2$, $I_v=1$ and $E_v=0$. Since $E_s$ is even, $E_s\ge2$ and $D\ge-1$. If $o\ge\frac32$, then $D\ge\frac32+\frac12-3=-1$. In all cases $D^\natural\ge1$. The remaining vertex types are covered by Lemma~\ref{app:G1:WP-deg}(c).
\end{proof}

\begin{lemma}[the weakened support hypothesis for coloured families]\label{app:G1:prop1col-depth}
Theorem~\ref{app:G1:WP-WP1}(c) holds, with a rate $\delta_0'>0$ and a constant $O(\alpha_0,\nbar;d,L,a,M_b)$ independent of $\eps$, $k$, the domains, $\tilde B$, $u$ and $\kappa$, if the hypotheses of \cite[Prop.~1]{Ba3} on $\Omega,\Omega_1,\Omega_2,\tilde B$ are replaced by those of Corollary~\ref{cor:prop1R4}.
\end{lemma}
\begin{proof}
The proof of Theorem~\ref{app:G1:WP-WP1}(c) is the proof of \cite[pp.~422--445]{Ba3}, with the $D$-lines added as vector lines of dimension $-d+4$ and with \cite[Prop.~2.2, p.~428]{Ba3} in place of \cite[Prop.~2.1]{Ba3}. Proposition~2.2 admits lines ``with the same exponential factors but with arbitrary dimensions instead of $-d+2$'' whose propagators have representations corresponding to \cite[(2.6)]{Ba3} with the estimates \cite[(2.10)--(2.12)]{Ba3}; Lemma~\ref{app:G1:prop1col-deg}(a) supplies them.

The proof of Corollary~\ref{cor:prop1R4} uses the hypothesis on $\Omega$ only through the depth needs listed in Proposition~\ref{prop:prop1} and Table~\ref{tab:prop1}. These needs involve only the following:
\begin{itemize}
\item the localisations of the vertices, through their anchoring at $\supp\tilde A$ or at $\supp g_k$;
\item the scalar propagators $G_k(\Omega,\tilde B)$, $\delta G_k(\Omega,\Omega_2,\tilde B)$ and the operator \cite[(1.16)]{Ba3};
\item the cube split \eqref{eq:prop1-P12}, which acts on scalar propagators only.
\end{itemize}
After the split, Lemma~\ref{lem:Omegafree} applies.

Colouring changes none of the vertices, localisations or scalar lines. It changes only the kernels of vector lines. By Lemma~\ref{app:G1:prop1col-deg}(a), the $D$-kernels are built from the torus vector slices and obey their bounds with no depth condition. In particular, they are among the depth-free objects of item (iii) of Lemma~\ref{lem:Omegafree}, up to the scalar weights $(L^j\eta)^2$. Hence every depth requirement of the coloured proof is a requirement of the uncoloured proof, and the argument of Corollary~\ref{cor:prop1R4} applies word for word.
\end{proof}

\begin{proposition}[Ba{\l}aban's Proposition~1 on the torus]\label{app:K:L08-5}
Fix $\nbar>12$, $0<\alpha_0<1$, $k\le K$ and $\eta=L^{-k}$, and let $\Omega,\Omega_1,\Omega_2,\tilde B$ satisfy the hypotheses of \cite[Prop.~1]{Ba3} in the form of Corollary~\ref{cor:prop1R4}. Work on $T_\eps$, or on the torus with photon covariance $C_{\mu_*}+u\Gamma^B$ ($\mu_*>0$, $\abs u\le1$), with the colours and families described above; in scale-$k$ units a $D$-line carries the coefficient $uc_k$, $0\le c_k\le\kappa\mu_*^2\eps_0^2L^2$. Then every renormalised family $G_{\rm ren}$ with $m=0$ satisfies
\[
\abs{E(G_{\rm ren},\{\square(v)\},\Phi_{\rm ext},A_{\rm ext})}\le O(1)\,e(L^k\eps)^{d_v}\lambda(L^k\eps)^{d_s}\mathrm{e}^{-\delta_Td(\{\square(v)\})}\norm{h\Phi_{\rm ext}}_{1,\alpha_0}\norm{h'A_{\rm ext}}_{1,\alpha_0},
\]
with $d$, $d_v$, $d_s$, $h$, $h'$ as in \cite[(1.33), p.~420]{Ba3}. Every family with $m\ge2$, or with $m=1$ and $o\ge2$, satisfies the same bound times $(\abs uc_k)^m$. The constant is $O(1)=O(\alpha_0,\nbar;d,L,a,M_b)$, and the rate $\delta_T>0$ may depend on $\nbar$. Both are independent of $\eps$, $k$, the domains and $\tilde B$.
\end{proposition}
\begin{proof}
\emph{Step 1: families with $m=0$.} On $T_\eps$ the families are those of \cite[pp.~418--421]{Ba3}. Ba{\l}aban's setting is the torus $T^{(k)}_1$ \cite[p.~420]{Ba3}, which is $T_\eps$ in scale-$k$ units. On the second torus, a family with $m=0$ has only $C$-lines, which carry $G_k$, and its vertices, scalar lines and counterterms are Ba{\l}aban's. So in both cases it is one of Ba{\l}aban's families. Corollary~\ref{cor:prop1R4} gives \cite[(1.33)]{Ba3} for it, with a rate $\delta_0>0$ and a constant $O(1)$. By Corollary~\ref{cor:prop1R4}, which carries over the uniformity stated in \cite[pp.~420--421]{Ba3} for domains ``satisfying the conditions mentioned previously'' to its weakened hypotheses, this constant depends on $\alpha_0$, $\nbar$ and the fixed data $d,L,a,M_b$ (Ba{\l}aban's $M$), and is independent of $\eps$, $k$, $\Omega,\Omega_1,\Omega_2$ and $\tilde B$.

\emph{Step 2: reduction of the coloured families.} Let $F$ have $m\ge1$ $D$-lines. By Lemma~\ref{app:G1:prop1col-exp}, $\abs{E(F,\dots)}=(\abs uc_k)^m\abs{E^\natural(F,\dots)}$, since $c_k\ge0$.

\emph{Step 3: the bound for $E^\natural$.} If $m\ge2$, or $m=1$ and $o\ge2$, then $F$ is either a non-vacuum family or a vacuum family with $m\ge2$ or with $m=1$, $o\ge2$. Hence it lies in the scope of Theorem~\ref{app:G1:WP-WP1}(c). This is consistent with Lemma~\ref{app:G1:prop1col-deg}(b): every vacuum graph with these values of $(m,o)$ is convergent, and only $C$-only divergent subgraphs need partners (Lemma~\ref{app:G1:WP-deg}(d)). By Lemma~\ref{app:G1:prop1col-depth}, Theorem~\ref{app:G1:WP-WP1}(c) applies under the hypotheses of Corollary~\ref{cor:prop1R4}. It bounds $\abs{E^\natural(F,\dots)}$ by the right side of the display, with a rate $\delta_0'>0$ and a constant $O(\alpha_0,\nbar)$. That constant is independent of $\eps$, $k$, the domains, $\tilde B$, $u$ and $\kappa$.

\emph{Step 4: rate and constant.} Put $\delta_T:=\min(\delta_0,\delta_0')>0$, and take the larger of the two constants. Since $d(\{\square(v)\})\ge0$, lowering the rate weakens both bounds, so Steps 1--3 give the claim. Only $\delta_T>0$ is used; no dependence of $\delta_T$ on $d$ alone is claimed.
\end{proof}

\subsection{The coefficient of the class (2.21c) and the convergent expressions of {\cite[\S3]{Ba3}}}\label{app:G1:I7-read-sec}

Here we use the notation of Lemma~\ref{app:G1:I7-masses}. All kernels live on the torus $\Lambda$ of spacing $\xi=L^{-j_0}$. We put $(\partial_\nu f)(y):=\xi^{-1}(f(y+\xi e_\nu)-f(y))$. For a kernel $K$, $\partial^{(1)}_\nu K$ and $\partial^{(2)}_\nu K$ denote this difference taken in the first and in the second variable. Recall $F=C^\xi$, and $W_f:=1-\nu_f-a_{j_0}P_1$ with $\nu_s=M^2t_{j_0}^2$ and $\nu_p=\mu_*^2t_{j_0}^2$. For $y'\in\Lambda$ and directions $\mu,\nu$ let $T^{\rm c}_{y';\mu,\nu}[K_1,K_2,K_3]$ be the expression $T^{\rm c}$ of Lemma~\ref{app:G1:I7-C21c} with the point $y'$, with $D_1$ the forward difference in direction $\mu$ and with $D_2=D_3$ the forward difference in direction $\nu$. Thus $D_1^{(1)}=\partial^{(1)}_\mu$ and $D_2^{(2)}=D_3^{(2)}=\partial^{(2)}_\nu$.

\begin{lemma}[two identities]\label{app:G1:I7-read-id}
\begin{itemize}
\item[(i)] For every function $u$ on $\Lambda$ and every single difference $D$ we have $\xi^3\sum_{y\in\Lambda}(Du)(y)=0$.
\item[(ii)] $G_f=F+G_fW_fF$ for $f\in\{s,p\}$. Let $\Phi[K_1,\dots,K_n]$ be linear in each of its kernel arguments. Write $\Phi_F:=\Phi[F,\dots,F]$, and for $S\subset\{1,\dots,n\}$ let $X^S_i:=K_iW_{f_i}F$ if $i\in S$ and $X^S_i:=F$ otherwise. If $K_i=G_{f_i}$ for all $i$, then $\Phi[K_1,\dots,K_n]-\Phi_F=\sum_{S\ne\emptyset}\Phi[X^S]$.
\end{itemize}
\end{lemma}

\begin{proof}
(i) The map $y\mapsto y\pm\xi e_\nu$ is a bijection of $\Lambda$, so $\sum_yu(y\pm\xi e_\nu)=\sum_yu(y)$.

(ii) We have $F^{-1}=-\Delta^\xi+1$ and $G_f^{-1}=-\Delta^\xi+\nu_f+a_{j_0}P_1$, so $F^{-1}-G_f^{-1}=W_f$. Hence $G_f-F=G_f(F^{-1}-G_f^{-1})F=G_fW_fF$. Insert $K_i=F+K_iW_{f_i}F$ in every argument and expand by linearity. The term $S=\emptyset$ is $\Phi_F$.
\end{proof}

\begin{lemma}[the vertex coefficient of the class (2.21c) and the convergent expressions of {\cite[\S3]{Ba3}}]\label{app:K:L11-2}\label{app:G1:I7-read}
Consider the graphs of orders $2\le\alpha+2\beta\le4$ at zero background without operator lines, in the proof of \cite[(1.33)]{Ba3} at a scale $j_0\ge1$.
\begin{itemize}
\item[(a)] Fix the following data:
\begin{itemize}
\item a member of the class \cite[(2.21c)]{Ba3}, that is, one of the two graphs that differ by the exchange of the external scalar legs: the graph (2.21c) drawn on \cite[p.~430]{Ba3} and its permuted graph (the classes are symmetric under permutation of the external scalar legs, \cite[p.~429]{Ba3});
\item the direction $\mu$ of its external vector leg;
\item components $\beta,\delta$ of the external scalar legs at the vertex $y'$ of the vector leg and at the other external vertex;
\item the receiving vertex of the transport \cite[(3.34)]{Ba3}.
\end{itemize}
Let $\Gamma$ be the corresponding entry of the coefficient $\sum\xi^{2d}\Gamma'_\mu$ of \cite[p.~444]{Ba3}, summed over orderings and indices and rescaled to the spacing $\xi$. Let $\Gamma_F$ be its $C^\xi$-version. If the receiving vertex is $y'$, then
\[
\Gamma=\sum_{\nu=1}^dc_\nu\,T^{\rm c}_{y';\mu,\nu}[G_s,G_s,G_p],\qquad\Gamma_F=\sum_{\nu=1}^dc_\nu\,T^{\rm c}_{y';\mu,\nu}[F,F,F],
\]
where $c_\nu=\sigma\,(q^{\rm T}qq)_{\beta\delta}$, the sign $\sigma\in\{\pm1\}$ does not depend on $\nu$, and $\abs{c_\nu}\le1$. If the receiving vertex is the other external vertex, then $\Gamma=\Gamma_F=0$. Hence, summed over the members of the class, the coefficient is a sum of at most $2d$ expressions $c\,T^{\rm c}[G_s,G_s,G_p]$ with $\abs c\le1$, and its $C^\xi$-version is the same sum with $T^{\rm c}[F,F,F]$. In these expressions:
\begin{itemize}
\item $D_1,D_2$ are the differences at the two ends of the scalar line from $y'$ to the inner vertex;
\item $D_3$ is the difference on the inner scalar line at the vertex of the external scalar leg.
\end{itemize}
\item[(b)] Consider the rescaled coefficients listed below, and write $\Phi_F$ for the version of a coefficient $\Phi$ in which every $G_s,G_p$ is replaced by $C^\xi$:
\begin{itemize}
\item the coefficient of (3.23);
\item a coefficient of p.~441;
\item the coefficient $\sum\xi^{2d}\Gamma'_\mu$ of p.~444 for the classes (2.19), (2.20) and (2.21c).
\end{itemize}
The expressions that \cite{Ba3} calls convergent after the replacements $G=C^\xi+GWC^\xi$ at pp.~439, 441 and~444 are the terms other than $\Phi_F$ produced from such a $\Phi$. Their sum is $\Phi-\Phi_F$. At p.~437 the replacement gives the identity (3.16) for the rescaled coefficient of (3.15) itself, and nothing is set aside there. The coefficients and these sums enter the constant $O(1)$ of \cite[(1.33)]{Ba3} only through upper bounds on their absolute values.
\end{itemize}
\end{lemma}

\begin{proof}
(a) \emph{The graph.} At zero background the vector-scalar vertices of order $e$ are the vertices \cite[(1.8)]{Ba3} with $n+n'=1$ (p.~413):
\[
e(L^k\eps)\,\eta^d\sum_b\big[(D^\eta_b\phi')(b)\cdot q\,\phi'(b_-)\big]\,g_kA'_b,
\]
up to a sign. Here $[u\cdot qv]=\sum_{\alpha\gamma}u_\alpha q_{\alpha\gamma}v_\gamma$. The vertex has one differentiated scalar position, carrying the forward difference in the direction of $b$, and one undifferentiated scalar position. Both positions and the vector leg sit at $b_-$. In the graphical notation of \cite{Ba3} an arrow at an end of a scalar line marks a difference at that end. This can be seen from (1.22) on p.~416: the graph with arrows at both ends of the scalar line is $e^2\sum_\mu q(\partial_\mu C\partial^*_\mu)(x-x')qC(x-x')$, and the kernel of $C\partial^*_\mu$ is $\partial^{(2)}_\mu C$. The graph (2.21c) on p.~430 has three such vertices:
\begin{itemize}
\item the vertex $y'$, which carries the external vector leg and one external scalar leg;
\item an inner vertex $y$;
\item a vertex $y''$, which carries the other external scalar leg.
\end{itemize}
It has three lines:
\begin{itemize}
\item a scalar line from $y'$ to $y$ with arrows at both ends;
\item a scalar line from $y$ to $y''$ with an arrow at $y''$;
\item a photon line from $y$ to $y''$.
\end{itemize}
The external scalar legs therefore occupy the undifferentiated positions at $y'$ and $y''$. The first line occupies the differentiated positions at $y'$ and $y$. The second line occupies the undifferentiated position at $y$ and the differentiated position at $y''$.

\emph{The propagators.} By \cite[p.~444]{Ba3}, the summation of (3.36) over orderings and indices gives the same expression with the propagators $G^\eta_{j_0}(0)$ and $G^\eta_{j_0}$. The summations are unrestricted and there are no localisation functions, since these were moved to the receiving vertex. Since $\sum\eta^{2d}\Gamma'_\mu$ has degree $0$, the rescaling to $\xi$ is exact (p.~444), and it turns these propagators into $G_s$ and the photon propagator. By Lemma~\ref{app:G1:VSR-V0} the photon propagator is $\delta_{\nu\nu'}G_p$. The scalar propagator is diagonal in the components. A line whose end at a point is attached to a differentiated position of direction $\nu$ carries $\partial_\nu$ in the variable of that point.

\emph{The contraction.} Let $\nu$ be the direction of the bond at $y$ and $\nu'$ the direction of the bond at $y''$. Contracting the positions gives:
\begin{itemize}
\item the factor $(\partial^{(1)}_\mu\partial^{(2)}_\nu G_s)(y',y)$ from the first line;
\item the factor $(\partial^{(2)}_{\nu'}G_s)(y,y'')$ from the second line;
\item the factor $\delta_{\nu\nu'}G_p(y,y'')$ from the photon line;
\item the charge factor $\sum_{\alpha\gamma}q_{\alpha\beta}q_{\alpha\gamma}q_{\gamma\delta}=(q^{\rm T}qq)_{\beta\delta}$.
\end{itemize}
The couplings $e(L^k\eps)$ are extracted as in \cite[(2.13)]{Ba3}. The combinatorial factor is $1$, for two reasons. First, the factor $1/3!$ of the third-order term is cancelled by the $3!$ assignments of the three distinct roles to the vertices. Second, the graph has no non-trivial automorphism, since $y'$ and $y''$ are fixed by their external legs, the two positions of each vertex are distinguishable, and the three lines have different end positions. The remaining signs combine to a single $\sigma\in\{\pm1\}$. By \cite[p.~426]{Ba3}, $\norm q\le1$, so $\abs{(q^{\rm T}qq)_{\beta\delta}}\le\norm q^3\le1$.

\emph{The receiving vertex.} In \cite[(3.33)]{Ba3} the kernel $\Gamma_\mu(x,x',x'')$ is paired with $A_\mu(x)\phi'(x')\phi''(x'')$. The coefficient in (3.34), (3.36) is the sum over the two points that do not carry the receiving leg. Two external legs sit at $y'$, so the kernel contains $\eta^{-d}$ times a Kronecker delta that identifies their two points.

If the receiving leg sits at $y'$, the delta removes one of the two summations. What remains is the sum over $y,y''$ with $y'$ fixed:
\[
\Gamma=\sigma(q^{\rm T}qq)_{\beta\delta}\sum_{\nu=1}^d\xi^6\sum_{y,y''}(\partial^{(1)}_\mu\partial^{(2)}_\nu G_s)(y',y)\,(\partial^{(2)}_\nu G_s)(y,y'')\,G_p(y,y''),
\]
which is the stated sum of terms $T^{\rm c}_{y';\mu,\nu}$.

If the receiving leg sits at $y''$, then $y''$ is fixed and the summation runs over $y'$ and $y$. For fixed $y$, the function $y'\mapsto(\partial^{(1)}_\mu\partial^{(2)}_\nu G_s)(y',y)$ is the difference $\partial_\mu u$ of $u:=(\partial^{(2)}_\nu G_s)(\cdot,y)$. Its sum over $y'$ vanishes by Lemma~\ref{app:G1:I7-read-id}(i), so $\Gamma=0$.

\emph{The $C^\xi$-version.} The same computation with every propagator equal to $F$ gives the formulas for $\Gamma_F$. The class has two members and each contributes at most $d$ terms, which gives the count $2d$.

(b) \emph{The page p.~437.} The replacement there acts only on $G^\xi_{j''}(0)$ in the twice-differenced factor. The resulting $C^\xi$-term is evaluated by $\sum_\nu\partial^\xi_\nu C^\xi\partial^{\xi*}_\nu=\delta^\xi-C^\xi$ and $\delta^\xi(y'-y)(y'_\mu-y_\mu)=0$, and the outcome is (3.16). So (3.16) is an exact identity for the rescaled coefficient, and no term of it is called convergent.

\emph{The pages pp.~439, 441 and~444.} Each coefficient $\Phi$ listed in (b) is linear in each of its propagators: it is a finite sum, over lattice points and indices, of products of these propagators and their differences. At p.~441 the convergent expressions are, in print, the terms in which at least one $G_{j_0}(0)$ is replaced by $G_{j_0}(0)(1-m^2_{j_0}-a_{j_0}P_{j_0})C^\xi$. These are the terms $\Phi[X^S]$ with $S\ne\emptyset$ of Lemma~\ref{app:G1:I7-read-id}(ii). At pp.~439 and~444 the replacement is printed as $\Phi=(\text{convergent expressions})+\Phi_F$ (p.~439: ``we get some convergent expressions plus'' the $C^\xi$-expression; p.~444: ``We get a convergent expression plus $\sum\xi^{2d}\Gamma''_\mu$''). Each single replacement is the identity of Lemma~\ref{app:G1:I7-read-id}(ii). In every case, therefore, the convergent expressions add up to $\Phi-\Phi_F$, whatever the order in which the replacements are made.

\emph{How the coefficients enter (1.33).} Each $\Phi$ is the coefficient function of one local vertex: the vertex of (3.15), the last vertex in (3.22), the vertices of p.~441, and the vertex of (3.36). The expressions that \cite{Ba3} calls convergent are produced inside the bracket of this coefficient. They therefore multiply the same vertex, and may be added before any estimate is made. By \cite[p.~432]{Ba3}, the bounds (3.2), together with (2.5) and (3.1), imply (1.33). On p.~433, (3.2) is reduced to (3.3). By p.~434, (3.3) follows from the representation (3.5), in which these vertices occur in generalised expressions, together with \cite[Prop.~2.2, p.~428]{Ba3}. Proposition~2.2 is Proposition~2.1 for generalised expressions. The proof of Proposition~2.1 rests on the estimate (2.13), which is obtained by ``taking absolute values of all factors'' \cite[p.~426]{Ba3}. A vertex coefficient therefore enters only through an upper bound on its absolute value.
\end{proof}

\begin{remark}\label{app:G1:I7-read-rem}
In Lemma~\ref{app:K:L11-2}(b), $\Phi_F=0$ for the classes (2.19), (2.20) and (2.21c), by (3.37), (3.38) and Lemma~\ref{app:K:L11-2}(a) together with Proposition~\ref{app:K:L11-1}(c). For the coefficient of (3.23), $\Phi_F=q^2\,\xi^3\sum_{x'}(\partial_\mu C^\xi)(x-x')C^\xi(x-x')$ (p.~439) contains no mass. Each of its entries is bounded on $\Lambda$ by $4c_*^2(A_1^T)^2$, uniformly in $\xi$, $j_0$ and $\ell$, by Step~1 of the proof of Lemma~\ref{app:G1:I7-C21c} and $\norm q\le1$; (3.24) is the infinite-lattice form of $\Phi_F$. So in each case the vertex coefficient that enters (1.33) is bounded by $\abs{\Phi-\Phi_F}+\abs{\Phi_F}$, where $\Phi_F$ contains no mass.
\end{remark}

\section{\texorpdfstring{Assembly of Theorem~\ref{thm:T6}}{Assembly of the vacuum bound for the slice models}}\label{app:G1:part-E7}

This appendix assembles the proof of Theorem~\ref{thm:T6} from its ingredients (Appendix~\ref{p:appE} and Proposition~\ref{app:G1:I6}) and the step lemmas of Ba{\l}aban's inductive argument, which are stated in Subsection~\ref{app:G1:part-TS05}. The collar exponent is $\kappa_0=\frac15$ (Lemma~\ref{app:G1:kappa0}); it is not carried to other exponents, such as $\kappa_B=\frac{16}5$ and the exponents of (E.11), which are separate quantities.

\subsection{The assembly}\label{app:G1:ssec-assembly}
The ingredients are those of Table~\ref{tab:T6-inputs}: Propositions~\ref{app:G1:I1}, \ref{app:G1:I2}, \ref{app:G1:I4}, \ref{app:G1:I5}, \ref{app:G1:I3}, \ref{app:G1:I6}, \ref{app:G1:I7}, \ref{app:G1:I8} and \ref{app:G1:I9}. The argument is Ba{\l}aban's inductive upper-bound argument of \cite[Sects.~2--3]{Ba2}, run on the torus for the slice model with the photon weights $w_j:=1+u\delta_j$ ($\delta_j$ as in \eqref{eq:T6-GammaB}). Besides the ingredients it uses six groups of step results, (S0)--(S5), stated in Subsection~\ref{app:G1:part-TS05}.

\paragraph{Notation.} $L=2$, $\eps=2^{-K}\eps_0$, and $t_k:=L^k\eps$ ($0\le k\le K$), so $t_0=\eps$ and $t_K=\eps_0$. For $X\subset T$, $\abs{X}_k$ is the number of blocks of side $t_k$ in $X$, and $N_k:=\abs{T}t_k^{-3}$. Put $\kappa_0:=\tfrac15$ and $\kappa_B:=\tfrac{16}{5}$. (Every $\kappa_0<\frac14$ works; $\kappa_0=\frac14$ does not, because some collar terms are $t^{1/4}$ times a power of $\log t^{-1}$, so that $C_{\kappa_0}=\infty$ at $\frac14$.) The admissible large-field data $\boldsymbol\Lambda_k=(\Lambda_0^{(j)})_{j<k}$, the regions $\Lambda_i^{(k-1)\prime}$, $\Omega_1:=B^k(\Lambda_7^{(k-1)\prime})$, and the collars $\mathrm{Col}_j:=\Omega_1^{(j)}\setminus\Omega_1^{(j+1)}$ are those of \cite[(2.7)--(2.8), (2.43), pp.~558, 566--570]{Ba2}, with $\mathrm{Col}_0:=\Lambda_7^{(0)c}$. $\mathcal Z_u$ is the integral defining the partition function of the slice model at $(u,M)$, taken with the \emph{Lebesgue} photon measure. That is, $d\mu_{N(0,C_B(u))}(A)=N_w\,\mathrm{e}^{-\frac12\ip{A}{C_B(u)^{-1}A}}dA$ and $\mathcal Z_u:=\int dA\,\mathrm{e}^{-\frac12\ip{A}{C_B(u)^{-1}A}}Z_\varphi(A;a_{K,T}(M))$. Put $E_0^w:=\log\mathcal Z_u|_{e=\lambda=0}=-\log N_w+E_{0,s}(M)$ and $E_1(u):=T_{\nbar}[\Phi_M(u)]-E_{0,s}(M)$. We write $E_1(u)=E_{\rm fin}(u)+E_K(u)$. Here $E_{\rm fin}(u)$ is the sum of the vacuum graphs with at least one line of index $\ge K$. It is a $c$-number, independent of the field, of $\boldsymbol\Lambda$ and of $k$.

\paragraph{The collar exponent.}
\begin{lemma}[the collar exponent]\label{app:G1:kappa0}
Let $p(t)=b_0(1+\log t^{-1})^{p}$ with $p=6$, and let $0<\kappa_0<\frac14$. For $a>0$ and $n\ge0$ put $S(a,n):=\sup_{0<t\le1}t^{a}(1+\log t^{-1})^{n}$. Then $S(a,n)=\mathrm{e}^{a-n}(n/a)^{n}$ if $n\ge a$, and $S(a,n)=1$ otherwise. Hence, for $t\le1$ and $q\le4$,
\[
t^{1/4}p(t)^q\le b_0^{q}\,S(\tfrac14-\kappa_0,6q)\,t^{\kappa_0},\qquad
t^{1/2}p(t)^2\le b_0^{2}\,S(\tfrac12-\kappa_0,12)\,t^{\kappa_0},
\]
and $t\le t^{\kappa_0}$, $t^3(1+\log_2(\eps_0/t))\le C(\kappa_0)\,t^{\kappa_0}$ for $t\le\eps_0\le1$. For $\kappa_0=\frac14$ and $q\ge1$ no constant $C$ satisfies $t^{1/4}p(t)^q\le C\,t^{1/4}$ for all $t\in(0,1]$. At $\kappa_0=\frac15$, $\log S(\frac1{20},24)=124.22\ldots$.
\end{lemma}
\begin{proof}
Put $s=\log t^{-1}\ge0$ and $g(s)=\mathrm{e}^{-as}(1+s)^n$. Then $g'(s)=g(s)\,(n/(1+s)-a)$, so $g$ increases for $1+s<n/a$ and decreases afterwards. If $n\ge a$, the maximum is at $1+s=n/a$, where $g=\mathrm{e}^{-a(n/a-1)}(n/a)^n$. If $n<a$, the maximum is $g(0)=1$. The displayed bounds follow from $t^{1/4}p^q=b_0^q\,t^{\kappa_0}\cdot t^{1/4-\kappa_0}(1+\log t^{-1})^{6q}$ and the analogous identity for $t^{1/2}p^2$. Also $t\le t^{\kappa_0}$ since $\kappa_0\le1$, and $t^{3-\kappa_0}(1+\log_2(\eps_0/t))\le S(3-\kappa_0,1)\max(1,1/\log2)$. At $\kappa_0=\frac14$ the ratio $t^{1/4}p(t)^q/t^{1/4}=p(t)^q$ tends to $\infty$ as $t\downarrow0$.
\end{proof}
\noindent Sizes $t^{1/4}p^q$ occur in \cite[(2.22), p.~561]{Ba2} ($q=3$), \cite[(2.78), p.~574]{Ba2} ($q=4$) and the count of \cite[p.~421]{Ba3} ($q=2,3$). By Lemma~\ref{app:G1:kappa0} they are at most $C(\kappa_0)t^{\kappa_0}$ per collar block, with $C(\kappa_0)<\infty$ at $\kappa_0=\frac15$. Every other use of $\kappa_0$ needs only $0<\kappa_0\le1$.

\begin{proposition}[assembly of Theorem~\ref{thm:T6}]\label{app:G1:assembly}
Assume the hypotheses of Theorem~\ref{thm:T6}; in particular (E.fin), $\eps_0\le\min(1/m_2,1/\mu_*,L)$ (Lemma~\ref{app:G1:N9prem}), and condition (E.comb) of Lemma~\ref{app:G1:I145-comb} with $c_A$ replaced by $c_A+C'_{\rm old}$, $C'_{\rm old}=C_{\rm old}+2C_{\rm loc}$ (Proposition~\ref{app:G1:I1}(iii)). Assume the ingredients (Propositions~\ref{app:G1:I1}--\ref{app:G1:I9}) and the step results:
\begin{itemize}
\item[(S1)] the scalar small-field step (Lemma~\ref{app:G1:S05-SC});
\item[(S2)] the photon small-field step (Lemma~\ref{app:G1:S05-PH});
\item[(S3)] the large-field part of the step (Lemma~\ref{app:G1:S05-LF}), with the region sum of Lemmas~\ref{app:G1:I145-geo} and~\ref{app:G1:I145-comb};
\item[(S4)] the renormalisation flow (Lemmas~\ref{app:G1:S05-RF}, \ref{app:G1:S05-TS}, Propositions~\ref{app:G1:I1}, \ref{app:G1:R9vac} and Lemma~\ref{app:G1:VS});
\item[(S5)] the closing step (Lemma~\ref{app:G1:S05-CL}), with Lemmas~\ref{app:G1:S05-F3}, \ref{app:G1:S05-N0}, \ref{app:G1:S05-FM}, \ref{app:G1:S05-ZM}, \ref{app:G1:IF1abs} and the order of choices (Lemma~\ref{app:G1:S05-ACY}).
\end{itemize}
The base identity (S0) is Lemma~\ref{app:G1:S05-B0}; it rests on Lemma~\ref{app:G1:S05-T}, Theorem~\ref{app:G1:S05-HA} and the identification of the counterterms (Proposition~\ref{prop:UV-B1e}).
Then \eqref{eq:T6} holds for every $K\ge0$ and every $j\ge1$, at $(u,M)\in\{(1,m_2),(-\eta,m_1)\}$, with
\begin{equation}\label{app:G1:eq-Cplus}
C_+^B(u;M)=C_{F1}\eps_0^{\kappa_0-3}+B_E^{\rm fin}(u,M)+B_{\mathfrak E}+B_{GS}+\tfrac12B_{F3}+B_{\rm comb},
\end{equation}
where $\eps_0^{\kappa_0-3}=\eps_0^{-14/5}$ and the six terms are defined in Steps~4 and~5 of the proof. None of them depends on $\eps$, $K$ or $\ell_T$, and none blows up at $u=-\eta$ or at $M=m_1$.
\end{proposition}
\noindent Together with Propositions~\ref{app:G1:I1}--\ref{app:G1:I9} and the step lemmas (S1)--(S5) of Subsection~\ref{app:G1:part-TS05}, Proposition~\ref{app:G1:assembly} gives Theorem~\ref{thm:T6}. The identification of the photon quadratic form and of $N_w$ is made in the Notation paragraph above and in Theorem~\ref{app:G1:S05-HA}.

\begin{proof}
The proof composes the step lemmas. Steps~1 and~3 use the convention (G) of Remark~\ref{rem:T6-textfixes}(ii): the vacuum counterterms are global c-numbers \cite[(1.30), p.~419]{Ba3}. We write $E_{\rm fin}(u)$ and $E_{K,k}(u)$ for the parts of the vacuum series $E_1(u)$ defined in Lemma~\ref{app:G1:S05-TS}. Here $E_{\rm fin}(u)$ is the sum of the monomials with at least one line of index $K$, and $E_{K,k}(u)$ is the sum of the monomials with at least one line index in $[k,K)$ and none equal to $K$. Thus $E_{K,0}(u)=E_K(u)=E_1(u)-E_{\rm fin}(u)$ and $E_{K,K}(u)=0$.

\emph{Step 0 (reduction).} By definition $e(\eps;u;M)=-\abs{T}^{-1}T_{\nbar}[\Phi_M(u)]$, so the left side of \eqref{eq:T6} equals $\abs{T}^{-1}(\Phi_M(u)-T_{\nbar}[\Phi_M(u)])$. The photon measure in $\Phi_M$ is normalised, so $\Phi_M(u)=\log(N_w\mathcal Z_u)$ by Theorem~\ref{app:G1:S05-HA}(a),(b) (see the proof of Lemma~\ref{app:G1:S05-B0}). The order-zero term of $T_{\nbar}[\Phi_M(u)]$ is $E_{0,s}(M)$ (Remark~\ref{rem:T6-textfixes}(i)). Hence
\[
\mathrm{e}^{\Phi_M(u)-T_{\nbar}[\Phi_M(u)]}=\mathrm{e}^{-E_0^w-E_1(u)}\mathcal Z_u,
\]
and \eqref{eq:T6} is equivalent to $\mathrm{e}^{-E_0^w-E_1(u)}\mathcal Z_u\le \mathrm{e}^{C_+^B(u;M)\abs{T}}$. The statement has to use the normalised photon, because the secant argument of Proposition~\ref{prop:ch-T5} needs convexity in $u$, while $\log\det C_B(u)$ is concave. Finally, $\bar\delta=\kappa\mu_*^2\eps_0^2\le\frac1{12}$ by (E.$\delta$), so $\eta\bar\delta\le\eta/12<1$ and $C_B(-\eta)>0$. Hence both points lie in the domain of $\Phi_M$.

\emph{Step 1 (inductive hypothesis).} For $0\le k\le K$ and both points, (IH$_k$) is the statement
\begin{equation}\label{app:G1:eq-IH}
\begin{split}
\mathrm{e}^{-E_1(u)}\mathcal Z_u\le{}& \mathrm{e}^{-E_{\rm fin}(u)-E_{K,k}(u)}\sum_{\boldsymbol\Lambda_k}\int dA_k\,d\varphi_k\int dA_{k+1}\cdots dA_K\prod_{j=k}^{K-1}\mathcal K^w_j(A_j,A_{j+1})\\
&\times\rho^{(k)}_w(\boldsymbol\Lambda_k;\cdot)\,\exp\big[\mathcal P^{(k)}_w(\Omega_1;\cdot)+\mathfrak E_k(\boldsymbol\Lambda_k)+\mathfrak O_k(\boldsymbol\Lambda_k)\big],
\end{split}
\end{equation}
with equality at $k=0$ and an empty coarse product at $k=K$. In \eqref{app:G1:eq-IH}:
\begin{itemize}
\item $\mathcal K^w_j$ are the (unnormalised) kernels of the photon chain of Lemma~\ref{app:G1:S05-T} with weights $w_j$. The fields of the coarse chain that are not yet integrated are kept.
\item $\rho^{(k)}_w$ is Ba{\l}aban's local Gaussian part and outer density \cite[(2.49)--(2.52), (2.55)]{Ba2}, with the coefficient $w_k^{-1}$ on the vector part and the small-field restrictions of \cite[(2.55), p.~570]{Ba2}.
\item $\mathcal P^{(k)}_w$ is the polynomial part of order $\le\nbar$ \cite[(1.4)--(1.5)]{Ba3} with the weighted lines $G_k(u)=G_k+uc_kD_k$. It contains no part of $E_{\rm fin}(u)$ or of $E_{K,k}(u)$. Its field-independent part is the vacuum remainder $\mathcal R^{\rm vac}_k$ of Lemma~\ref{app:G1:S05-RF}, with $\mathcal R^{\rm vac}_0=0$.
\item The remainders are
\begin{gather*}
\mathfrak E_k(\boldsymbol\Lambda_k):=\sum_{j<k}C_{\kappa_0}t_j^{\kappa_0}\abs{\mathrm{Col}_j}_j+\sum_{j<k}C_\kappa t_j^{\kappa_B}N_j+\sum_{j=1}^{k}\nu(t_j)\abs{\Lambda_7^{(j-1)\prime}}_j,\qquad
\mathfrak O_k(\boldsymbol\Lambda_k):=\sum_{j<k}\bar S A_j\abs{\Lambda_7^{(j)c}}_j,\\
\mathfrak E_0:=\mathfrak O_0:=0.
\end{gather*}
Here $\bar S$ and $A_j$ are as in Proposition~\ref{app:G1:R9vac}, and $\bar SA_j\le C_{\rm old}t_j$ with the constant $C_{\rm old}$ of Lemma~\ref{app:G1:S05-RF}. The function $\nu$ gives the field-independent part $\nu(t_{k+1})\abs{\Lambda_7^{(k)\prime}}_{k+1}$ of the cost of deleting the spectator cross term in the photon small-field step (Lemma~\ref{app:G1:S05-PH}(d)).
\end{itemize}
The constants $C_{\kappa_0},C_\kappa,C_{\rm old}$ and the function $\nu$, which satisfies $\nu(t)=O(t^{s})$ for every $s>0$, depend only on the parameters, on the window $[1-\eta,2]$ of the weights, and on the regulator data. They depend on $\eps_0$ only through quantities evaluated at the cap $\eps_1$, because Ba{\l}aban's per-step constants are bounded uniformly in $\eps_0\in(0,\eps_1]$ (Statement~\ref{app:K:SL}). They do not depend on $\eps$, $k$, $\ell_T$ or $\boldsymbol\Lambda_k$.

The constant $C_\kappa$ contains the factor $\sup_{t\le\eps_1}t^{1/20}(1+\log t^{-1})^q<\infty$, which converts volume terms of size $t^{13/4}$ times a power of $\log t^{-1}$ into $t^{16/5}$. Likewise $C_{\kappa_0}$ contains $\sup_{t\le\eps_1}t^{1/20}(1+\log t^{-1})^{q}<\infty$ ($q\le24$). This factor converts the collar terms of size $t^{1/4}$ times a power of $\log t^{-1}$ of \cite[(2.22), (2.57), (2.78)]{Ba2} and \cite[p.~421]{Ba3} into $t^{1/5}$ (Lemma~\ref{app:G1:kappa0}); at $\kappa_0=\frac14$ it would be infinite. $C_\kappa$ also contains the tail constants of \cite[pp.~579, 582]{Ba2} (Propositions~\ref{app:G1:T579} and~\ref{app:G1:I5-T582}), taken as suprema over $t\in(0,1]$; these tails are volume terms of size $t^{7/2}$ times a power of $\log t^{-1}$, per block of scale $t$.

\emph{Step 2 (base case).} By Lemma~\ref{app:G1:S05-B0}, (IH$_0$) holds with equality: the front factor is $\mathrm{e}^{-E_{\rm fin}(u)-E_{K,0}(u)}$ with $E_{K,0}(u)=E_K(u)$, $\mathcal P^{(0)}_w=-V$ and $\mathfrak E_0=\mathfrak O_0=0$. Here $V(\varphi)=\sum_x\eps^3(\lambda\abs{\varphi(x)}^4+\frac12\delta m^2_{K,T}(M)\abs{\varphi(x)}^2)$. The mass counterterm is the torus-built $\delta m^2_{K,T}(M)$, which is Ba{\l}aban's $\delta m^2$ \cite[(1.19)--(1.23), pp.~416--417]{Ba3} by Proposition~\ref{prop:UV-B1e}. The integrand is integrable because the whole chain is ($HA>0$, Lemma~\ref{app:G1:S05-T}), although at $u=1$ single kernels grow.

\emph{Step 3 (inductive step).} For $0\le k\le K-1$ and both points, (IH$_k$) implies (IH$_{k+1}$) with
\[
\mathfrak E_{k+1}=\mathfrak E_k+C_{\kappa_0}t_k^{\kappa_0}\abs{\mathrm{Col}_k}_k+C_\kappa t_k^{\kappa_B}N_k+\nu(t_{k+1})\abs{\Lambda_7^{(k)\prime}}_{k+1},\qquad
\mathfrak O_{k+1}=\mathfrak O_k+\bar SA_k\abs{\Lambda_7^{(k)c}}_k .
\]
The operations follow Ba{\l}aban's order \cite[pp.~569--582]{Ba2}: regions; collar removal; the vector translation and the fluctuation law; the scalar translation and $Z^{(j)}$; the joint fluctuation integral; the new effective action and its vacuum part; and finally the closing of the photon spectators.
\begin{enumerate}
\item[(a)] \emph{Large fields and collars} (S3). The new regions are a function of the past and of a minimal family. The region expansion is needed only in the direction ``$\le$'', since the integrand is $\ge0$. For $k\ge1$, on the collar $\mathrm{Col}_k=\Lambda_7^{(k-1)\prime}\cap\Lambda_7^{(k)c}$ the whole interaction, including the step-$k$ mass pieces, is removed under the step-$(k-1)$ restrictions \cite[Prop.~2.1, (2.57), p.~570]{Ba2}. The cost is at most $c_C^{(k)}t_k^{\kappa_0}\abs{\mathrm{Col}_k}_k$ with $c_C^{(k)}\le\sigma_+2^{N_\ell}c_C^{u=0}$ (Lemma~\ref{app:G1:S05-LF}(b), with the sizes bounded by Lemma~\ref{app:G1:kappa0}). At $k=0$, where $\mathrm{Col}_0=\Lambda_7^{(0)c}$, the cost is at most $C_\delta(M)^2\eps/(16\lambda)\le\frac{C_\delta(M)^2}{16\lambda}t_0^{\kappa_0}$ per site (Lemma~\ref{app:G1:S05-LF}(b)). We take $C_{\kappa_0}\ge\max\{\sigma_+2^{N_\ell}c_C^{u=0},C_\delta(M)^2/(16\lambda)\}$, so every collar costs at most $C_{\kappa_0}t_k^{\kappa_0}\abs{\mathrm{Col}_k}_k$. No vacuum c-number is removed on a collar.
\item[(b)] \emph{Scalar and photon small-field steps} (S1), (S2). Every localised family of $\mathcal P_w$ obeys the tree bound \cite[(1.33)]{Ba3} in weighted form.
\begin{itemize}
\item At each $k$ this bound is given by Proposition~\ref{app:G1:I5} (order cap $\Nbar=169$). It is applied under the conditions and operations of \cite[Prop.~1]{Ba3} of Proposition~\ref{app:G1:I3}. These include the support hypothesis in the form (R4$'''$) of Corollary~\ref{cor:prop1R4} and the repairs of Lemmas~\ref{lem:lem24} and~\ref{lem:IMG}.
\item The joint fluctuation integral uses no bound at a complex background near a Neumann face (Proposition~\ref{app:G1:I4}).
\item The torus defect of the zeroth moment at orders $\ge2$ is the image term of the Ward--Takahashi step (Proposition~\ref{app:G1:I2}).
\item The uniformity in $\ell_T$ of the constant $C_b$ of the $\delta m^2_{\rm fin}$ vertex is Proposition~\ref{app:G1:I8}.
\item The weights enter the vector part only through a factor $w_k$ on Ba{\l}aban's Dirichlet covariance and through the spectator cross term. Deleting that term costs $\nu(t_{k+1})\abs{\Lambda_7^{(k)\prime}}_{k+1}+\xi_{k+1}\norm{A_{k+1}|_{\Lambda_6^{(k)\prime c}}}^2$ (Lemma~\ref{app:G1:S05-PH}(d)). The second term depends on $A_{k+1}$ only through $\Lambda_6^{(k)\prime c}A_{k+1}$.
\end{itemize}
\item[(c)] \emph{Renormalisation flow} (S4). The mass counterterm is $u$-independent and is the same number at both values of $u$ for the same $M$ (Lemma~\ref{app:G1:S05-RF}). Write $\mathrm{e}^{-E_{K,k}(u)}=\mathrm{e}^{-E_{K,k+1}(u)}\mathrm{e}^{-\cV_k(u)}$ with $\cV_k(u):=E_{K,k}(u)-E_{K,k+1}(u)$. Only $\cV_k(u)$ is localised at step $k$. It is localised at scale $k$ by the rule of \cite[p.~420]{Ba3}, applied to whole renormalised classes with one external field each; the anchor of a class is its first leg, the same in every member. Put $W_k:=B^k(\Lambda_7^{(k)})$, $\mathcal L_k:=B^k(\Lambda_7^{(k)}\cap\Lambda_8^{(k)c})$ and $\mathcal D_k:=B^k(\Lambda_8^{(k)})$. A localisation $\Lambda$ is \emph{good} if $\Lambda\subset W_k$, \emph{bad} if it meets both $W_k$ and $W_k^c$, and \emph{outside} if $\Lambda\cap W_k=\emptyset$.
\begin{itemize}
\item \emph{Good localisations.} By (a), every vertex of a family generated at step $k$ lies in $W_k$. By Lemma~\ref{app:G1:S05-TS}, the field-independent output of step $k$, truncated at order $\nbar$ and computed with the free index-$k$ lines, is $\cV_k(u)$, as the same sum of index-assigned Wick graphs. A generated graph and the corresponding counterterm graph have the same vertex set, hence the same localisation. So the good localisations of $\cV_k(u)$ cancel the generated output localisation by localisation.
\item \emph{What remains in $\mathcal R^{\rm vac}_{k+1}-\mathcal R^{\rm vac}_k$} is the bad localisations and the domain differences (Lemma~\ref{app:G1:S05-RF}).
\item \emph{Bad localisations.} By Proposition~\ref{app:G1:I1}(ii) they are bounded by $(\bar N\bar SA_k+4C_{\rm loc}t_k)\abs{\mathcal L_k}_k+(\bar N\bar SA_k\mathrm{e}^{3\delta}\mathrm{e}^{-\delta r(t_k)/2}+2C_{\rm dp}t_k^3\mathrm{e}^{-2\delta_1(r(t_k)-4)})\abs{\mathcal D_k}_k$. The terms with $A_k$ come from Proposition~\ref{app:G1:R9vac}. The terms with $C_{\rm loc}$ and $C_{\rm dp}$ come from the order-$e^2$ vacuum family with one $D$-line (Lemma~\ref{app:G1:I1c-e2}). We have $A_k\le Ct_k$, $\delta\le\delta_1$ and $t_k\le1$. So the coefficient of $\abs{\mathcal L_k}_k$ is at most $Ct_k$, and the coefficient of $\abs{\mathcal D_k}_k$ is at most $Ct_k\mathrm{e}^{-\delta r(t_k)/2}$. For the families with $C$-lines only and $\delta m^2$ insertions, Proposition~\ref{app:G1:I9} applies, via Lemma~\ref{app:G1:VS}.
\item \emph{The layer.} Let $\beta$ be a block of $\Lambda_7^{(k)}\cap\Lambda_8^{(k)c}$. By \cite[(2.8)]{Ba2} it lies within $r(t_k)$ of a point of $\Lambda_7^{(k)c}$; let $b$ be the unit block that contains this point. If $b\subset\Lambda_0^{(k)c}$, then $\beta\subset\Lambda_1^{(k)c}\subset\Lambda_7^{(k)c}$ by \cite[(2.8)]{Ba2}, which is impossible. Hence $b\subset\Lambda_0^{(k)}\cap\Lambda_7^{(k)c}\subset\mathrm{Col}_k$. Counting gives $\abs{\mathcal L_k}_k\le(2r(t_k)+2\sqrt3M_b+1)^3\abs{\mathrm{Col}_k}_k$. The coefficient $\bar N\bar SA_k+4C_{\rm loc}t_k$ of $\abs{\mathcal L_k}_k$ is at most $Ct_k$; moreover $r(t)=R(1+\log t^{-1})^2$ and $\kappa_0<1$. Hence the layer costs are at most $C(\kappa_0)t_k^{\kappa_0}\abs{\mathrm{Col}_k}_k$. They are part of the collar term of the increment of $\mathfrak E$.
\item \emph{Deep part and domain differences.} The deep part is at most $Ct_k\mathrm{e}^{-\delta r(t_k)/2}N_k$, and the domain differences are at most $C_\kappa t_k^{\kappa_B}N_k$ with $\kappa_B=\frac{16}5$ (Lemma~\ref{app:G1:S05-RF}). Both are counted in the term $C_\kappa t_k^{\kappa_B}N_k$.
\item \emph{Outside localisations.} They are at most $\bar SA_k\abs{\Lambda_7^{(k)c}}_k\le C_{\rm old}t_k\abs{\Lambda_7^{(k)c}}_k$ (Proposition~\ref{app:G1:R9vac}). This c-number is not part of $\mathcal P_w$; it is the increment of $\mathfrak O$.
\end{itemize}
Each term of $E_K(u)=E_{K,k}(u)+\sum_{j<k}\cV_j(u)$ leaves the front factor exactly once, at its top index, and $E_{\rm fin}(u)$ is never moved. The global sum of the order-$e^2$ family with one $D$-line, which is of size $O(K)\abs T$, is never bounded. It occurs only inside $E_1(u)$, on both sides of the identity of Step~0.
\end{enumerate}
The constants are the same at both points. The data of the stages up to 9 of Appendix~\ref{app:B} do not depend on the point, and the conditions that do are imposed at the point where they are strongest. A single $\eps_0$ (the minimum over the two points) therefore gives the same $K$ at both points.

\emph{Step 4 (closing: volume terms).} Apply (IH$_K$); its front factor is $\mathrm{e}^{-E_{\rm fin}(u)}$, since $E_{K,K}(u)=0$. Three bounds follow.
\begin{itemize}
\item (C1) On the final small-field set, $\mathcal P^{(K)}_w(\Lambda_7^{(K-1)\prime})\le C_{F1}\eps_0^{\kappa_0}\abs{\Lambda_7^{(K-1)\prime}}_K\le C_{F1}\eps_0^{\kappa_0-3}\abs{T}$, using $\abs{\cdot}_K\le N_K=\abs T\eps_0^{-3}$. This is \cite[(2.116)]{Ba2} in weighted form. It is read with $\abs{\Lambda_7^{(K-1)}}$ in place of the printed $\abs{\Lambda_7^{(K)}}$, since no region of scale $K$ exists. The bound rests on Proposition~\ref{app:G1:I5}, and uses Propositions~\ref{app:G1:I2}, \ref{app:G1:I4} and, for the supremum over $\ell_T$, \ref{app:G1:I8}.
\item (C2) We have $\abs{\Lambda_7^{(j-1)\prime}}_j\le N_j$, $\nu(t)\le C_\nu t^{\kappa_B}$, $t_j^{\kappa_B}N_j=\abs T t_j^{\kappa_B-3}$ and $t_j=2^{-(K-j)}\eps_0$. Hence
\[
\sum_{j<K}C_\kappa t_j^{\kappa_B}N_j+\sum_{j\le K}\nu(t_j)\abs{\Lambda_7^{(j-1)\prime}}_j\le(C_\kappa+C_\nu)\abs T\sum_{m\ge0}(2^{-m}\eps_0)^{1/5}=:B_{\mathfrak E}\abs T,
\]
\[
B_{\mathfrak E}=(1-2^{-1/5})^{-1}(C_\kappa+C_\nu)\eps_0^{1/5}\le7.73\,(C_\kappa+C_\nu)\eps_0^{1/5}.
\]
Only $\kappa_B>3$ is used, and this keeps the sum uniform in $\eps$. The collar terms and $\mathfrak O_K$ stay in the summand and are paid in Step~5.
\item (C3) The backward replacements of \cite[pp.~583--587]{Ba2} carry the weight $w_{k-1}^{-1}$. Their determinants are at most $\exp(d\varrho^{\det}_k\abs{\Lambda_5^{(k)}})$, with $\varrho^{\det}_k\le\frac12$ (condition (E.4)) and $\varrho^{\det}_k\le C_4't_k^4$ (Lemma~\ref{app:G1:S05-CL}). Hence $\sum_kd\varrho^{\det}_kN_k\le B_{GS}\abs T$ with $B_{GS}:=2dC_4'\eps_0$.
\end{itemize}

\emph{Step 5 (closing: Gaussian integrals, region sum, $E_{\rm fin}$).}
\begin{itemize}
\item (C4) \emph{Scalar Gaussian integral.} The scalar field is integrated first, at fixed background \cite[Prop.~3.1, (3.30)--(3.38)]{Ba2}. The diamagnetic bound is used on the plain operator only, at every holonomy, and gives at most $\exp(E_{0,s}(M))$. The Gaussian mass is $M^2$. The reading of ``$m^2\le O(1)$'' \cite[p.~589]{Ba2} is Proposition~\ref{app:G1:I7}, with (E.m2) and (E.$\Phi$).
\item (C4, continued) \emph{Vector Gaussian integral and the normalisation defect.} On the vector side the chain rule gives, for every pattern $\boldsymbol\Lambda$, exactly $(2\pi)^{n/2}\det HA_{\rm del}(\boldsymbol\Lambda)^{-1/2}$, where $n$ is the number of vector variables. The factor $\mathrm{e}^{+E_0^w}$ produced in this way cancels the $\mathrm{e}^{-E_0^w}$ of Step~0 identically for every $\boldsymbol\Lambda$ (Lemma~\ref{app:G1:S05-N0}). By Lemma~\ref{app:G1:S05-F3}, which needs $\delta_0R\ge4$, $\abs{\log\det HA_{\rm del}(\boldsymbol\Lambda)-\log\det HA}\le B_{F3}\abs T$ with
\[
B_{F3}:=12\abs u\kappa C_{\ell^1}(1-\hat\xi_{\rm tot})^{-1}w_{\min}^{-2}w_{\rm lo}^{-1}\gamma_1^{-1}\eps_0.
\]
Here $\hat\xi_{\rm tot}\le\frac1{12}$ is the deletion budget ((E.14)), $w_{\min}=\min_jw_j\ge\frac{11}{12}$, $w_{\rm lo}$ is as in Theorem~\ref{app:G1:S05-HA}(c), $\gamma_1$ is the photon gap constant of Lemma~\ref{app:G1:S05-ZM}, and $C_{\ell^1}$ is the constant of Lemma~\ref{app:G1:S05-F3}. None of them depends on $\eps$, $K$ or $\ell_T$, and $B_{F3}\propto\abs u$.
\item (C4, continued) \emph{Zero modes.} The zero modes are massive, with gaps $\ge\gamma_1\mu_*^2$ (photon) and $\ge M^2(1+4M^2\eps_0^2/(3a))^{-1}$ (scalar), both free of $\ell_T$.
\item (C5) \emph{Region sum.} After the cancellation, consider the sum over $\boldsymbol\Lambda_K$ of the large-field factors, multiplied by $\exp$ of the collar terms, the shell terms, the count of kept coarse variables and $\mathfrak O_K(\boldsymbol\Lambda_K)$. This sum is at most $\exp(B_{\rm comb}\abs T)$. It is the region entropy of \cite[p.~593]{Ba2} (Proposition~\ref{app:G1:I6}), obtained via Lemmas~\ref{app:G1:I145-geo} and \ref{app:G1:I145-comb} under $2p=12>\varrho d+2$ and (E.comb). The constant is evaluated at $c_A+C_{\rm old}$ at $\eps_1$; the term $\mathfrak O_K$ is absorbed by Lemma~\ref{app:G1:IF1abs}. The constants of \cite[(3.42)]{Ba2} that enter (E.ent) are bounded on $(0,\eps_1]$ by Statement~\ref{app:K:SL}. With $n_0$ defined by $\eps_0=2^{-n_0}$, and $c_{\rm fin}\eps_0^{-3}\abs T$ the bound for the factor $\exp(O(1)\abs{T^{(K)}})$ of the kept variables at the final scale \cite[(3.47)]{Ba2}, the constant is
\[
B_{\rm comb}=c_{\rm fin}\eps_0^{-3}+2(d+2)\Sigma_*,\qquad \Sigma_*=\sum_{n\ge n_0}L^{3n}\exp\big(-\tfrac12c_0'^{\,w}b_0^2(1+n\log L)^{2p}\big).
\]
\item ($E_{\rm fin}$) The factor $\mathrm{e}^{-E_{\rm fin}(u)}$ in front of (IH$_K$) is paid once: $\abs{E_{\rm fin}(u)}\le B_E^{\rm fin}(u,M)\abs T$. The bound is uniform in $\eps$ and in $\ell_T$ by Propositions~\ref{app:G1:Efin} and \ref{app:G1:I8}. This is the only place where $B_E^{\rm fin}$ enters $C_+^B$.
\end{itemize}
Collecting Steps 0--5: (IH$_K$), (C1) and (C2) reduce the claim to \cite[(2.117)]{Ba2} in weighted form. For each $\boldsymbol\Lambda$ we apply (C4) and pay $\mathrm{e}^{(B_{GS}+\frac12B_{F3})\abs T}$. We then sum over $\boldsymbol\Lambda$ by (C5), and finally pay $\mathrm{e}^{-E_{\rm fin}(u)}$. This gives $\mathrm{e}^{-E_0^w-E_1(u)}\mathcal Z_u\le \mathrm{e}^{C_+^B(u;M)\abs T}$ with \eqref{app:G1:eq-Cplus}, which is \eqref{eq:T6} by Step~0.

\emph{Step 6 (uniformity and order of choices).} The regulator data are chosen in the order of the stages 0--11 of Appendix~\ref{app:B}.
\begin{itemize}
\item Every lower bound on $R$ depends only on data of the stages 2--4.
\item The upper bounds on $\eps_0$ of stage 10b, among them (E.fin), (E.edge), (E.$\delta$), (E.14) and (E.comb), are downward closed.
\item The lower bounds on $\ell_0$ of stage 11, among them (T.cap), (T.div) and (T.per), are fixed after their data.
\end{itemize}
So the order is acyclic and the conditions can be met (Lemmas~\ref{app:G1:acyc}, \ref{app:G1:S05-ACY} and \ref{lem:window}). The constants of the upper bounds on $\eps_0$ that depend on $\ell_T$ are evaluated at the floor $\ell_0^\natural=\max\{8,128/m_1,128/M_1\}$ (Lemma~\ref{app:G1:floor}). The majorants needed for $\delta m^2_{\rm fin}$ are those of Proposition~\ref{app:G1:I8}. The Ward identity of \cite[p.~442]{Ba3} with $\lambda$ linear in $y'$ is exact on $\eta\Z^d$. On the torus it holds up to the image term of Lemma~\ref{lem:IMG}; this is a $\delta G$-type external field and costs at most one more lower bound on $\ell_0$ at stage 11.

Every constant in Steps 1--5 is therefore free of $\eps$, $K$ and $\ell_T$. At $u=-\eta$ only $(1-\eta)^{-1}$, $(1-\eta)^{-3}$, $w_{\min}\ge\frac{11}{12}$ and $\vartheta_-\ge\frac{11}{24}$ enter. At $M=m_1$ only $C_\delta(m_1)$, $m_1^{-2}$, $\log_+(1/(m_1\eps_0))$, the scalar gap and the floor $128/m_1$ enter, and all of these are finite. Since the bound holds for every $K\ge0$, it holds in particular for $K\ge K_0$.
\end{proof}

\subsection{The step lemmas (S0)--(S5)}\label{app:G1:part-TS05}

\paragraph{(S0) The base case.}

\paragraph{The photon chain.}
\emph{Setting.} In this paragraph $j$ and $k$ are scale indices. Fix a torus $T$, $L=2$, $\eps=2^{-K}\eps_0$ and $t_k=L^k\eps$.

\emph{Spaces and block averages.} For $0\le k\le K$ let $\mathbb V_k$ be the space of real functions on the blocks of side $t_k$; every operator below acts componentwise on the vector field. Put $\ip fg_k:=t_k^3\sum fg$ and $\norm f_k^2=\ip ff_k$, and let $dA_k$ be the Lebesgue measure of $(\mathbb V_k,\ip\cdot\cdot_k)$. Let $Q_k:\mathbb V_0\to \mathbb V_k$ be the block average and $Q:=Q_{k\to k+1}:\mathbb V_k\to \mathbb V_{k+1}$ the one-step average. Adjoints are taken with respect to these products, so that $QQ^*=1$, $Q_kQ_k^*=1$, $Q_{k+1}=QQ_k$, and we put $P_k:=Q_k^*Q_k$.

\emph{Operators.} Let $\mathbb M:=-\Delta_\eps+\mu_*^2$ and $C_1:=\mathbb M^{-1}=C_{\mu_*}$. Let $a>0$ be the averaging constant of the block-spin transformation \cite[(2.13)--(2.18), p.~653]{K1}. Put $a_k=a(1-L^{-2})(1-L^{-2k})^{-1}$, $r_k:=a_kt_k^{-2}$ for $k\ge1$ (not to be confused with the function $r(t)$ of \eqref{eq:UV-r}), $r_0^{-1}:=0$, and $\beta_k:=at_{k+1}^{-2}$. Then
\begin{equation}\label{app:G1:eq-rrec}
r_{k+1}^{-1}=r_k^{-1}+\beta_k^{-1}\qquad(0\le k<K),
\end{equation}
since $1/a_{k+1}=L^{-2}/a_k+1/a$ and $a_1=a$. Put
\begin{itemize}
\item $G_k:=(\mathbb M+r_kP_k)^{-1}$ ($k\ge1$), $G_0:=0$, and $H_k:=G_{k+1}-G_k$;
\item $F_k:=r_kG_kQ_k^*$ ($k\ge1$) and $F_0:=1$;
\item $\Delta^{(0)}:=\mathbb M$ and $\Delta^{(k)}:=r_k-r_k^2Q_kG_kQ_k^*$ ($k\ge1$) \cite[(2.13)--(2.18), p.~653]{K1};
\item $C^{(k)}:=(\Delta^{(k)}+\beta_kQ^*Q)^{-1}$, $\mathcal H_k:=\beta_kC^{(k)}Q^*$ and $D_k(B):=\ip B{\Delta^{(k)}B}_k$.
\end{itemize}

\emph{Weights.} For weights $w_0,\dots,w_{K-1}>0$ put $w_K:=1$ and $\omega_k:=w_k^{-1}-w_{k-1}^{-1}$ ($1\le k\le K$). At the two points, $w_j=1+u\delta_j$ with $\delta_j=\kappa\mu_*^2t_{j+1}^2$ as in \eqref{eq:T6-GammaB}. A centred Gaussian $N(0,\Sigma)$ on $(\mathbb V_k,\ip\cdot\cdot_k)$ has a density with respect to $dA_k$ proportional to $\mathrm{e}^{-\frac12\ip x{\Sigma^{-1}x}_k}$.

\begin{lemma}[Gaussian conditioning]\label{app:G1:S05-gauss}
Let $X\sim N(0,\Sigma)$ with $\Sigma>0$ on a finite-dimensional inner product space $V$, let $S:V\to W$ be linear, and let $Y=SX+\xi$ with $\xi\sim N(0,\Lambda)$ independent of $X$ and $\Lambda>0$. Then:
\begin{itemize}
\item $Y\sim N(0,S\Sigma S^*+\Lambda)$;
\item the conditional law of $X$ given $Y$ is $N(\Pi Y,\Sigma_{X|Y})$, with $\Sigma_{X|Y}=(\Sigma^{-1}+S^*\Lambda^{-1}S)^{-1}$ and $\Pi=\Sigma_{X|Y}S^*\Lambda^{-1}$;
\item $\Sigma=\Sigma_{X|Y}+\Pi(S\Sigma S^*+\Lambda)\Pi^*$;
\item $(\Lambda+S\Sigma S^*)^{-1}=\Lambda^{-1}-\Lambda^{-1}S(\Sigma^{-1}+S^*\Lambda^{-1}S)^{-1}S^*\Lambda^{-1}$.
\end{itemize}
\end{lemma}
\begin{proof}
The joint density of $(X,Y)$ is proportional to $\exp(-\frac12\ip x{\Sigma^{-1}x}-\frac12\ip{y-Sx}{\Lambda^{-1}(y-Sx)})$. Completing the square in $x$ gives the conditional law. The covariance identity is the law of total covariance. The last identity (Woodbury) follows by multiplying the right side by $\Lambda+S\Sigma S^*$.
\end{proof}

\begin{lemma}[block actions and lifts]\label{app:G1:S05-block}
For $0\le k<K$:
\begin{itemize}
\item[(i)] $(\Delta^{(k)})^{-1}=Q_kC_1Q_k^*+r_k^{-1}$; in particular $\Delta^{(k)}>0$.
\item[(ii)] $\Delta^{(k+1)}=\beta_k-\beta_k^2QC^{(k)}Q^*$ \cite[(3.41), p.~619]{Ba1}.
\item[(iii)] $F_k\mathcal H_k=F_{k+1}$ \cite[(3.42), p.~619]{Ba1}.
\item[(iv)] $C_1-G_k=F_k(\Delta^{(k)})^{-1}F_k^*$ and $H_k=F_kC^{(k)}F_k^*$ \cite[(2.17), p.~653]{K1}.
\item[(v)] For $A\in \mathbb V_k$ and $B\in \mathbb V_{k+1}$,
\[
D_k(A)+\beta_k\norm{QA-B}_{k+1}^2=\ip{A-\mathcal H_kB}{(C^{(k)})^{-1}(A-\mathcal H_kB)}_k+D_{k+1}(B).
\]
\end{itemize}
\end{lemma}
\begin{proof}
(i) For $k=0$ this is $\mathbb M^{-1}=C_1$. For $k\ge1$ apply Woodbury's identity of Lemma~\ref{app:G1:S05-gauss} with $\Lambda=r_k^{-1}$, $S=Q_k$ and $\Sigma=C_1$: $(r_k^{-1}+Q_kC_1Q_k^*)^{-1}=r_k-r_k^2Q_k(\mathbb M+r_kP_k)^{-1}Q_k^*=\Delta^{(k)}$.

(ii) By (i) at $k+1$, $Q_{k+1}=QQ_k$, $QQ^*=1$ and \eqref{app:G1:eq-rrec},
\[
(\Delta^{(k+1)})^{-1}=Q\big(Q_kC_1Q_k^*+r_k^{-1}\big)Q^*+\beta_k^{-1}=Q(\Delta^{(k)})^{-1}Q^*+\beta_k^{-1}.
\]
Woodbury's identity with $\Lambda=\beta_k^{-1}$, $S=Q$ and $\Sigma=(\Delta^{(k)})^{-1}$ gives $\beta_k-\beta_k^2Q(\Delta^{(k)}+\beta_kQ^*Q)^{-1}Q^*$.

(iii), (iv) Let $A\sim N(0,C_1)$, $B=Q_kA+r_k^{-1/2}\xi$ and $B'=QB+\beta_k^{-1/2}\xi'$, with $\xi,\xi'$ standard and independent (at $k=0$, $B=A$). By Lemma~\ref{app:G1:S05-gauss} and (i), $B\sim N(0,Q_kC_1Q_k^*+r_k^{-1})=N(0,(\Delta^{(k)})^{-1})$.
\begin{itemize}
\item The law of $A$ given $B$ is $N(F_kB,G_k)$, since the posterior precision is $\mathbb M+r_kP_k$.
\item The law of $B$ given $B'$ is $N(\mathcal H_kB',C^{(k)})$, since the posterior precision is $\Delta^{(k)}+\beta_kQ^*Q$.
\item Given $B$, the fields $A$ and $B'$ are independent, so $\E[A\,|\,B']=\E[F_kB\,|\,B']=F_k\mathcal H_kB'$.
\item On the other hand $B'=Q_{k+1}A+\zeta$, where $\zeta=r_k^{-1/2}Q\xi+\beta_k^{-1/2}\xi'$ is independent of $A$ and has covariance $r_k^{-1}+\beta_k^{-1}=r_{k+1}^{-1}$. So $\E[A\,|\,B']=F_{k+1}B'$.
\end{itemize}
Since $B'$ has a non-degenerate law, $F_k\mathcal H_k=F_{k+1}$. The law of total covariance for $(A,B)$ gives $C_1=G_k+F_k(\Delta^{(k)})^{-1}F_k^*$. For $(B,B')$ it gives $(\Delta^{(k)})^{-1}=C^{(k)}+\mathcal H_k(\Delta^{(k+1)})^{-1}\mathcal H_k^*$. Hence, by (iii),
\[
H_k=(C_1-G_k)-(C_1-G_{k+1})=F_k\big[(\Delta^{(k)})^{-1}-\mathcal H_k(\Delta^{(k+1)})^{-1}\mathcal H_k^*\big]F_k^*=F_kC^{(k)}F_k^*.
\]

(v) Expand the right side. Use $(C^{(k)})^{-1}=\Delta^{(k)}+\beta_kQ^*Q$ and $(C^{(k)})^{-1}\mathcal H_k=\beta_kQ^*$. Also $\ip{\mathcal H_kB}{(C^{(k)})^{-1}\mathcal H_kB}_k=\beta_k^2\ip B{QC^{(k)}Q^*B}_{k+1}=\ip B{(\beta_k-\Delta^{(k+1)})B}_{k+1}$, by (ii).
\end{proof}

\begin{lemma}[the quadratic form of the photon chain]\label{app:G1:S05-T}
Let $w_0,\dots,w_{K-1}>0$, and let $HA$ be the quadratic form on $\mathbb V_0\times\cdots\times \mathbb V_K$ for which $N_w\mathrm{e}^{-\frac12HA}$ is the joint density of the photon chain $(\hat B_0,\dots,\hat B_K)$ of Theorem~\ref{app:G1:S05-HA}. Let $I_k(A_k,A_{k+1}):=\ip{A_k-\mathcal H_kA_{k+1}}{(C^{(k)})^{-1}(A_k-\mathcal H_kA_{k+1})}_k$ and $D_K(B):=\ip{B}{\Delta^{(K)}B}_K$. Then
\[
HA=\sum_{k<K}w_k^{-1}I_k(A_k,A_{k+1})+D_K(A_K).
\]
Equivalently, $\mathrm{e}^{-\frac12HA}=\mathrm{e}^{-\frac12w_0^{-1}D_0(A_0)}\prod_{j=0}^{K-1}\mathcal K^w_j(A_j,A_{j+1})$ with
\[
\mathcal K^w_j(A_j,A_{j+1})=\exp\big(-\tfrac12w_j^{-1}\beta_j\norm{QA_j-A_{j+1}}_{j+1}^2-\tfrac12\omega_{j+1}D_{j+1}(A_{j+1})\big).
\]
Each term $\omega_{j+1}D_{j+1}$ occurs in exactly one kernel. $HA>0$, in particular at both points, where $w_j=1+u\delta_j\ge1-\eta/12>0$ by (E.$\delta$).
\end{lemma}

\begin{theorem}[hierarchical form of the slice photon]\label{app:G1:S05-HA}
Let $w_0,\dots,w_{K-1}>0$. Let $\hat S_k\sim N(0,w_kC^{(k)})$ on $\mathbb V_k$ ($k<K$) and $\hat B_K\sim N(0,(\Delta^{(K)})^{-1})$ be independent, and put $\hat B_k:=\hat S_k+\mathcal H_k\hat B_{k+1}$ for $k=K-1,\dots,0$ (the photon chain).
\begin{itemize}
\item[(a)] $A:=\hat B_0$ has covariance $\sum_{j<K}w_jH_j+(C_1-G_K)$. At $w_j=1+u\delta_j$ this is $C_B(u)=C_{\mu_*}+u\Gamma^B$.
\item[(b)] The joint law of $(\hat B_0,\dots,\hat B_K)$ has density $N_w \mathrm{e}^{-\frac12HA}$ with respect to $dA_0\cdots dA_K$, where
\[
HA=\sum_{k<K}w_k^{-1}I_k(\hat B_k,\hat B_{k+1})+D_K(\hat B_K)
=w_0^{-1}D_0(\hat B_0)+\sum_{k<K}w_k^{-1}\beta_k\norm{Q\hat B_k-\hat B_{k+1}}_{k+1}^2+\sum_{k=1}^{K}\omega_kD_k(\hat B_k).
\]
Here $I_k\ge0$ is as in Lemma~\ref{app:G1:S05-T}, $n_k=\dim \mathbb V_k$, and $N_w^{-1}=\prod_{k<K}(2\pi)^{n_k/2}\det(w_kC^{(k)})^{1/2}\cdot(2\pi)^{n_K/2}\det(\Delta^{(K)})^{-1/2}$.
\item[(c)] Put $w_{\rm lo}:=\min(1,\min_kw_k^{-1})$, $w_{\rm hi}:=\max(1,\max_kw_k^{-1})$ and $S^1:=HA|_{w\equiv1}$. Then $w_{\rm lo}S^1\le HA\le w_{\rm hi}S^1$ as quadratic forms. The same inequalities hold after minimisation over any set of variables and on every diagonal block.
\end{itemize}
\end{theorem}
\begin{proof}
(a) Unfolding the recursion gives $\hat B_0=\sum_{j<K}\mathcal H_0\cdots\mathcal H_{j-1}\hat S_j+\mathcal H_0\cdots\mathcal H_{K-1}\hat B_K$. By Lemma~\ref{app:G1:S05-block}(iii), iterated from $F_0=1$, $\mathcal H_0\cdots\mathcal H_{j-1}=F_j$. By independence and Lemma~\ref{app:G1:S05-block}(iv),
\[
\mathrm{Cov}(A)=\sum_jw_jF_jC^{(j)}F_j^*+F_K(\Delta^{(K)})^{-1}F_K^*=\sum_jw_jH_j+C_1-G_K .
\]
Since $\sum_{j<K}H_j=G_K$, at $w_j=1+u\delta_j$ this equals $C_1+u\sum_j\delta_jH_j=C_B(u)$ by \eqref{eq:T6-GammaB}.

(b) $\hat B_{k+1}$ is a function of $(\hat S_j)_{j>k}$ and $\hat B_K$, so it is independent of $\hat S_k$. Hence the law of $\hat B_k$ given $(\hat B_{k+1},\dots,\hat B_K)$ is $N(\mathcal H_k\hat B_{k+1},w_kC^{(k)})$, whose normalisation does not depend on the condition. The joint density is the product of the density of $\hat B_K$ and these conditional densities; this gives the first expression for $HA$ and the value of $N_w$. For the second expression substitute Lemma~\ref{app:G1:S05-block}(v): $I_k=D_k(\hat B_k)+\beta_k\norm{Q\hat B_k-\hat B_{k+1}}_{k+1}^2-D_{k+1}(\hat B_{k+1})$. The coefficient of $D_0$ is then $w_0^{-1}$, that of $D_k$ ($1\le k<K$) is $w_k^{-1}-w_{k-1}^{-1}=\omega_k$, and that of $D_K$ is $1-w_{K-1}^{-1}=\omega_K$.

(c) $HA-w_{\rm lo}S^1=\sum_k(w_k^{-1}-w_{\rm lo})I_k+(1-w_{\rm lo})D_K\ge0$, since $I_k\ge0$ and $\Delta^{(K)}>0$; the bound with $w_{\rm hi}$ is proved in the same way. Pointwise inequalities between quadratic forms are preserved by minimisation over a set of variables and by restriction to a coordinate subspace.
\end{proof}

\begin{proof}[Proof of Lemma~\ref{app:G1:S05-T}]
The first formula for $HA$ is Theorem~\ref{app:G1:S05-HA}(b). Grouping the terms of the second expression of Theorem~\ref{app:G1:S05-HA}(b) by $j=k$ for the $\beta$-terms and $j=k-1$ for the $\omega_k$-terms gives the product formula with the kernels $\mathcal K^w_j$; each $\omega_{j+1}D_{j+1}$ occurs in exactly one kernel. $HA$ is the precision of a non-degenerate Gaussian, so $HA>0$. At the two points, $0\le\delta_j\le\kappa\mu_*^2\eps_0^2\le\frac1{12}$ by (E.$\delta$), since $t_{j+1}\le t_K=\eps_0$. So $w_j\ge1-\eta/12>0$, and the weights are admissible.
\end{proof}

\begin{lemma}[the base identity]\label{app:G1:S05-B0}
At both points $(u,M)\in\{(1,m_2),(-\eta,m_1)\}$, with $\boldsymbol\Lambda_0=\emptyset$, $\Lambda_i^{(-1)\prime}=T_\eps$, $\Omega=\Omega_1=T_\eps$, $A^{(0)}=A_0$,
\[
\mathrm{e}^{-E_1(u)}\mathcal Z_u=\mathrm{e}^{-E_{\rm fin}(u)}\int dA_0\,d\varphi_0\int dA_1\cdots dA_K\prod_{j=0}^{K-1}\mathcal K^w_j(A_j,A_{j+1})\,\mathrm{e}^{-\mathcal G^{(0)}_w(A_0,\varphi_0)}\exp[-V(\varphi_0)-E_K(u)].
\]
Here
\[
\mathcal G^{(0)}_w=\tfrac12w_0^{-1}\ip{A_0}{\mathbb M A_0}_0+\tfrac12\ip{\varphi_0}{(-\Delta^\eps_{A_0}+M^2)\varphi_0},\qquad
V(\varphi)=\sum_{x\in T_\eps}\eps^3\big(\lambda\abs{\varphi(x)}^4+\tfrac12\delta m^2_{K,T}(M)\abs{\varphi(x)}^2\big),
\]
and $E_K(u)=E_1(u)-E_{\rm fin}(u)=E_{K,0}(u)$ (Lemma~\ref{app:G1:S05-TS}). Hence $\rho'\equiv1$, $\chi_0\equiv1$, $\mathfrak E_0=\mathfrak O_0=0$ and $HA_{\rm del}^{(0)}=HA$. In the form \eqref{app:G1:eq-IH} at $k=0$ the front factor is $\mathrm{e}^{-E_{\rm fin}(u)-E_{K,0}(u)}$, and
\[
\mathcal P^{(0)}_w=-V .
\]
\end{lemma}

\begin{proof}
\emph{The photon.} The photon of the slice model is the normalised Gaussian $\mu_{C_B(u)}$. By Theorem~\ref{app:G1:S05-HA}(a),(b), $\mu_{C_B(u)}$ is the $A_0$-marginal of the law $N_w\mathrm{e}^{-\frac12HA}dA_0\cdots dA_K$. Let $S_\varphi(A_0,\varphi_0)$ be the scalar action of the slice model, and put $Z_\varphi(A_0):=\int d\varphi_0\,\mathrm{e}^{-S_\varphi(A_0,\varphi_0)}$. The integrand is nonnegative, so Tonelli's theorem gives
\[
\int d\mu_{C_B(u)}(A)\,Z_\varphi(A)=N_w\int dA_0\cdots dA_K\,\mathrm{e}^{-\frac12HA}Z_\varphi(A_0)=N_w\mathcal Z_u .
\]
By Lemma~\ref{app:G1:S05-T}, $\mathrm{e}^{-\frac12HA}=\mathrm{e}^{-\frac12w_0^{-1}D_0(A_0)}\prod_{j=0}^{K-1}\mathcal K^w_j$, and $D_0(A_0)=\ip{A_0}{\mathbb M A_0}_0$ since $\Delta^{(0)}=\mathbb M$.

\emph{The scalar.} King's torus scalar action equals Ba{\l}aban's at $m^2=M^2$, $\mu_0=\mu_*$, $\delta m^2=\delta m^2_{K,T}(M)$ (Lemma~\ref{app:G1:ctid}; Proposition~\ref{prop:UV-B1e}). Hence $S_\varphi=\frac12\ip{\varphi}{(-\Delta^\eps_{A_0}+M^2)\varphi}+V$. Multiplying by $\mathrm{e}^{-E_1(u)}=\mathrm{e}^{-E_{\rm fin}(u)}\mathrm{e}^{-E_K(u)}$ gives the displayed identity.

\emph{The form \eqref{app:G1:eq-IH} at $k=0$.} In convention (G) of Remark~\ref{rem:T6-textfixes}(ii), both c-numbers $E_{\rm fin}(u)$ and $E_{K,0}(u)=E_K(u)$ stand in front, so the polynomial part is $\mathcal P^{(0)}_w=-V$. It has order $\le\nbar$ in the couplings, since $\lambda$ has order $1$ and $\delta m^2_{K,T}$ has orders $1\le\alpha+\beta\le4\le\nbar$; its field-independent part is $0=\mathcal R^{\rm vac}_0$. No region has been chosen, so there are no restrictions and no deleted variables.

\emph{Integrability.} $HA>0$ (Lemma~\ref{app:G1:S05-T}). At $u=1$ single kernels grow ($\omega_k<0$ for $1\le k<K$); only the whole chain is integrable.
\end{proof}

\paragraph{(S1), (S2) The small-field steps.}

\begin{lemma}[scalar small-field step]\label{app:G1:S05-SC}
Fix the regulator data of Theorem~\ref{thm:T6} ($L=2$, $p=6$, $p_1=b_1=1$, $\nbar=13$, $R\ge R^{**}$, $\eps_0$, $\ell_0$) and put $\kappa_0=\frac15$. Let $(u,M)\in\{(1,m_2),(-\eta,m_1)\}$, $\eps=2^{-K}\eps_0$, $\ell_T\in\{j!\ell_0\}$ and $0\le k\le K-1$. We work in scale-$k$ units ($\eta=L^{-k}$, $t=t_k$). $\abs X$ is the number of scale-$k$ blocks in $X$, $\mathrm{Col}_k:=\Lambda_7^{(k-1)\prime}\cap\Lambda_7^{(k)c}$, and $N_\ell$ is the largest number of vector lines of a graph of order $\le\nbar$. Assume:
\begin{itemize}
\item[(SC1)] (IH$_k$) of \eqref{app:G1:eq-IH}. In particular: the scalar Gaussian part of $\rho^{(k)}_w$ is
\[
-\tfrac12\ip{\Lambda_6^{(k-1)\prime}\varphi}{\Delta^{(k)}(B^k(\Lambda_2^{(k-1)\prime}),A^{(k),\eps})\Lambda_6^{(k-1)\prime}\varphi}
\]
\cite[(2.49), p.~568]{Ba2}, with no weight; $\rho^{(k)}_w$ depends on the configuration on $B^k(\Lambda_7^{(k-1)\prime})$ only through $\prod_{j<k}Z^{(j)}_{\Lambda_5^{(j)}}(B^j(\Lambda_2^{(j)}),A^\eps)$ (ibid.); the restrictions \cite[(2.55), p.~570]{Ba2} hold on $\Lambda_{-1}^{(k-1)\prime}$, among them $\abs{(D_{\tilde A^{(k)}}\varphi)(b)}\le c_1p(t_{k-1})$ and $\abs{\varphi(x)}\le c_1\lambda(t_{k-1})^{-1/4}p(t_{k-1})$; $\theta_kA^{(k)}$ is an independent external field that occurs only in scalar expressions; the scalar normalisations are carried exactly.
\item[(SC2)] Part (a) of Lemma~\ref{app:G1:S05-PH}; the regularity of $\tilde B=(1-\theta_{k+1})\theta_kA^{(k)}+\theta_{k+1}B^{(k+1),\eta}$ and of $B^{(k+1),\eta}$ required in \cite[Prop.~2.1]{Ba1}, as supplied by \cite[(2.98), p.~577]{Ba2}; and the kernel bounds \cite[(2.10)--(2.12)]{Ba3} for the $D$-line, with dimension $-d+4$.
\item[(SC3)] Part (a) of Lemma~\ref{app:G1:S05-LF}.
\item[(SC4)] The mass counterterm is $\delta m^2_{K,T}(M)$ (Proposition~\ref{prop:UV-B1e}), with the scale split of \cite[(1.26)--(1.29)]{Ba3}.
\end{itemize}
Then the following hold. The constants depend only on the parameters, $\eta$, the regulator data and $\ell_0^\natural$; they do not depend on $\eps$, $k$, $\ell_T$, $\boldsymbol\Lambda_k$, the background, or the sign of $u$. Every $u$-dependent constant at $u=-\eta$ is at most its value at $u=1$.
\begin{itemize}
\item[(C1)] Under \cite[(2.55)]{Ba2},
\[
\mathcal P^{(k)}_w(\Lambda_7^{(k-1)\prime},\theta_kA^{(k)},\varphi)=-\lambda(t)\!\!\sum_{x\in B^k(\mathrm{Col}_k)}\!\!\eta^d\abs{\varphi^{(k)}(x)}^4+\mathcal P^{(k)}_w(\Lambda_7^{(k)},\theta_kA^{(k)},\varphi)+\mathcal R_1,
\]
with $\abs{\mathcal R_1}\le C_1(1+2^{N_\ell}\abs u c_k)\,t^{\kappa_0}\abs{\mathrm{Col}_k}$. Here $\varphi^{(k)}=a_kG_k(B^k(\Lambda_2^{(k-1)\prime}),A^{(k)})Q_k^*(A^{(k)})\Lambda_6^{(k-1)\prime}\varphi$ and $c_k=\kappa\mu_*^2L^2t_k^2\le\frac13$. At $k=0$, moreover, $\lambda\eps\abs\varphi^4+\frac12\delta m^2_{K,T}(M)\eps^2\abs\varphi^2\ge-C_\delta(M)^2\eps/(16\lambda)$ at every site.
\item[(C2)] Replacing the scalar quadratic forms by their localised versions \cite[(2.47)--(2.48), (2.108)]{Ba2} costs $O(t^\kappa)\abs{\Lambda_6^{(k-1)}}$ and $O(t^\kappa)\abs{\Lambda_3^{(k)}}$ respectively, for every $\kappa$. The difference operator $H_k$ satisfies $\abs{h_k(x,x')}\le O(1)\mathrm{e}^{-\delta_1r(t)}\mathrm{e}^{-\delta_0\abs{x-x'}}$ \cite[(2.109)]{Ba2}.
\item[(C3)] We have
\[
\prod_{j<k}Z^{(j)}(A+B)\le\prod_{j<k}Z^{(j)}(\tilde B)\exp\big(\mathcal Q^{(k)}(B^k(\Lambda_2^{(k)});B^{(k+1),\eta},g_kA'^{(k)})+C(\kappa_0)t^{\kappa_0}\abs{\mathrm{Col}_k}+O(t^\kappa)\abs{\Lambda_7^{(k)}}\big).
\]
The collar constant carries no factor $k$.
\item[(C4)] In the interaction localised in $B^k(\Lambda_7^{(k)})$, replace $G_k(B^{k-1}(\Lambda_2^{(k-1)}),\tilde B)$ by $G_k(B^k(\Lambda_2^{(k)}),B^{(k+1),\eta})$. All terms containing a difference are $O(t^\kappa)\abs{\Lambda_7^{(k)}}$. After the translation \cite[(2.110)]{Ba2}, the identity \cite[(2.111)--(2.112)]{Ba2} holds up to $O(t^\kappa)$.
\item[(C5)] For $x\in\Lambda_5^{(k)}$, \cite[(2.113)]{Ba2} holds, and $\abs{\varphi'(x)}\le O(1)p(t)$.
\item[(C6)] With the vector half from Lemma~\ref{app:G1:S05-PH}(b),
\[
\int d\mu_{w_kC^{(k)}_{\Lambda_5}}(A')\int d\mu_{C^{(k)}_{\Lambda_5}(B^k(\Lambda_2^{(k)}),B^{(k+1),\eta})}(\varphi')\,\chi(A')\chi(\varphi')\,\mathrm{e}^{V^{(k)}}\le\exp\Big(\sum_{m\le\nbar}\frac{\kappa_m(V^{(k)})}{m!}+O(t^\kappa)\abs{\Lambda_5^{(k)}}\Big).
\]
\item[(C7)] The cumulant sum, after the replacements of \cite[p.~582]{Ba2}, is $\mathcal P^{(k+1)}_w(\Lambda_7^{(k)},B^{(k+1),\eta},\psi)$, with vector lines $G_{k+1}(u)=G_{k+1}+uc_{k+1}D_{k+1}$. Every localised family except the order-$e^2$ vacuum family with one $D$-line obeys \cite[(1.33)]{Ba3} (Proposition~\ref{app:G1:I5}), the $D$-coloured ones with the factor $(\abs uc_{k+1})^m$; that family is part of the vacuum part of (C8) (Lemma~\ref{app:G1:S05-PH}(c)). The step-$(k+1)$ mass coefficients obey the pointwise bounds of Theorem~\ref{app:G1:I7-I3a}, with the constant $C_m$ of Proposition~\ref{app:G1:I7-Cm}.
\item[(C8)] The vacuum part is as stated in Lemma~\ref{app:G1:S05-RF}(V), (C11), with the $u$-part from Proposition~\ref{app:G1:I1} and the $C$-only bad localisations from Proposition~\ref{app:G1:I9}.
\item[(C9)] \cite[Prop.~3.1, (3.26)]{Ba2} holds at every real $\tilde A$ satisfying the vector restrictions, with $\gamma_0=\gamma_0(d,a)$ and error $\sum_kO(t_k^{\kappa_0})\abs{\Lambda_k}$; $\Delta^{(k)}(\Omega,A)\ge a_kM^2t^2/(a_k+M^2t^2)$ for every real $A$ and every $\Omega$; the difference forms of \cite[(3.7)--(3.8)]{Ba2} obey the kernel bound of (C2).
\item[(C10)] The scalar items of (IH$_{k+1}$) hold, with $\mathfrak E_{k+1}-\mathfrak E_k$ given by the scalar contributions to Step~3 of the proof of Proposition~\ref{app:G1:assembly}.
\end{itemize}
Every scalar condition is free of $w$; $u$ enters only through the vector background law and the $D$-lines.
\end{lemma}
\noindent The constant $C_b$ uniform in $\ell_T$ is that of Proposition~\ref{app:G1:I8}, and by Proposition~\ref{app:G1:I4} no bound at a complex background is used.

\begin{lemma}\label{app:G1:SC-quad}
Let $\lambda,\eps>0$, $\abs\delta\le C_\delta/\eps$ and $x\ge0$. Then $\lambda\eps x^2+\frac12\delta\eps^2x\ge-C_\delta^2\eps/(16\lambda)$.
\end{lemma}
\begin{proof}
The quadratic in $x$ is minimised at $x=-\delta\eps/(4\lambda)$, with value $-\delta^2\eps^3/(16\lambda)$. If $\delta\ge0$, the minimum over $x\ge0$ is $0$. In both cases the bound follows from $\delta^2\eps^2\le C_\delta^2$.
\end{proof}

\begin{lemma}\label{app:G1:SC-diff}
Let $\mathcal P(\Omega)=\sum_{X\subset\Omega}E(X)$, the sum running over connected unions of blocks, and assume $\sum_{X\ni\square}\abs{E(X)}\le a$ for every block $\square$. Then $\abs{\mathcal P(\Omega)-\mathcal P(\Omega')}\le a\abs{\Omega\setminus\Omega'}$ for $\Omega'\subset\Omega$.
\end{lemma}
\begin{proof}
$\mathcal P(\Omega)-\mathcal P(\Omega')=\sum_{X\subset\Omega,\,X\not\subset\Omega'}E(X)$, and every such $X$ contains a block of $\Omega\setminus\Omega'$. Summing over that block first gives the bound.
\end{proof}

\begin{lemma}\label{app:G1:SC-pow}
Let $0<t\le\eps_0\le1$, $0<\kappa_0<\frac14$, $\ell(t)=1+\log t^{-1}$ and $n_t:=\log_2(\eps_0/t)$. There is $C(\kappa_0)<\infty$ such that each of
\[
t^{1/4}p(t)^3,\quad t^{1/2}p(t)^q\ (q\le6),\quad t\,(1+t\,n_t),\quad t\,\ell(t)^{2p}(r(t)+M_b)^3,\quad r(t)^3t^3(1+n_t)
\]
is at most $C(\kappa_0)\,t^{\kappa_0}$.
\end{lemma}
\begin{proof}
Each quantity is $t^{\kappa_0}$ times $t^{a}\ell(t)^n$ with $a>0$, times bounded factors. Here $r(t)=R\ell(t)^2$, $p(t)=b_0\ell(t)^6$, $t\,n_t\le\eps_0/(\mathrm{e}\log2)$ and $n_t\le\log_2(\mathrm{e})\,\ell(t)$. Now apply $S(a,n)=\sup_{0<t\le1}t^a\ell(t)^n<\infty$ (Lemma~\ref{app:G1:kappa0}).
\end{proof}

\begin{lemma}\label{app:G1:SC-col}
Let $G$ be a graph of order $\le\nbar$ with at most $N_\ell$ vector lines. Assume each colouring of these lines with $m\ge1$ $D$-lines is bounded by $(\abs uc)^m$ times the bound $B_0$ of the $C$-only colouring, where $0\le\abs uc\le1$. Then the sum over all colourings with $m\ge1$ is at most $2^{N_\ell}\abs uc\,B_0$.
\end{lemma}
\begin{proof}
$\sum_{m\ge1}\binom{N_\ell}{m}(\abs uc)^m\le\abs uc\sum_{m\ge1}\binom{N_\ell}{m}\le2^{N_\ell}\abs uc$.
\end{proof}

\begin{proof}[Proof of Lemma~\ref{app:G1:S05-SC}]
We follow the order of operations of \cite[pp.~569--582]{Ba2}.

\emph{Step 0 (the weights act only on the photon law).} By Lemma~\ref{app:G1:S05-T}, the weights $w_j$ enter only the kernels $\mathcal K^w_j$ and, through Lemma~\ref{app:G1:S05-PH}(a), the precision $w_k^{-1}(C^{(k)}_{\Lambda_5})^{-1}$ of the photon fluctuation, whose minimiser is $w$-free. All scalar objects are deterministic functions of the vector configuration and the scalar fields, with no $w$ in them: $G_k(\Omega,B)$, $\Delta^{(k)}$, $Z^{(j)}$, $\chi_k$, $\varphi^{(k)}$, the scalar fluctuation law, and the statements \cite[Prop.~2.6, Lemma~2.7]{Ba2}. Hence every scalar bound below is Ba{\l}aban's, with his constants. The parameter $u$ enters only through the restriction constants (Proposition~\ref{app:K:L13-11}(iv), in the form of Corollary~\ref{app:G1:restr-slice-use}, whose premise is Lemma~\ref{app:K:L13-10}(ii)) and through the $D$-lines of $\mathcal P_w$, in both cases only via $\abs u\le1$ and $w_j\in[1-\eta,2]$.

\emph{Step 1, (C1).} \emph{The $u=0$ part.} Write $\mathcal P^{(k)}$ as a sum over localisations. The vertex \cite[(1.6)]{Ba3} with all four scalar legs external is kept explicitly as $-\lambda(t)\sum\eta^d\abs{\varphi^{(k)}}^4$. Every other localised term is bounded by \cite[(1.33)]{Ba3} (Proposition~\ref{app:G1:I5}), with external fields controlled by \cite[(2.55)]{Ba2}. By \cite[p.~421]{Ba3}, the worst term carries $O(\lambda(t)^{1/4}p(t)^3)$, i.e.\ $O(t^{1/4}\ell(t)^{18})$ in $d=3$. Summed over localisations containing a fixed block, this is $\le C\,t^{1/4}p(t)^3\le C(\kappa_0)t^{\kappa_0}$ by Lemma~\ref{app:G1:SC-pow}. Lemma~\ref{app:G1:SC-diff} with $\Omega=\Lambda_7^{(k-1)\prime}$, $\Omega'=\Lambda_7^{(k)}$ gives the $u=0$ part of $\mathcal R_1$.
\emph{The step-$k$ mass-counterterm pieces.} At $k\ge1$ these pieces leave $\mathrm{Col}_k$ together with the rest of the interaction. With $\abs{\varphi^{(k)}}^2\le O(1)c_1^2\lambda^{-1/2}t^{-1/2}p^2$ and the pointwise bound $\abs{\delta m^2_{K,k}}t^2\le C_m[(e^2+\lambda)t\frac{L}{L-1}+(e^4+e^2\lambda+\lambda^2)t^2(K-k)]$ (Theorem~\ref{app:G1:I7-I3a}, with $C_m$ of Proposition~\ref{app:G1:I7-Cm}), each vertex is $\le C\,t^{1/2}p^2(1+t\log_2(\eps_0/t))\le C(\kappa_0)t^{\kappa_0}$ per site (Lemma~\ref{app:G1:SC-pow}). The finite piece $\delta m^2_{\rm fin}$ gives $O(C_bt^{3/2}p^2)$ per site, with $C_b$ uniform in $\ell_T$ (Proposition~\ref{app:G1:I8}).
\emph{The $D$-coloured part.} By the coloured form of \cite[(1.33)]{Ba3} (Proposition~\ref{app:G1:I5}) and Lemma~\ref{app:G1:SC-col}, it is bounded by $2^{N_\ell}\abs uc_k$ times the $u=0$ bound, with $c_k\le4\kappa\mu_*^2\eps_0^2\le\frac13$ by condition (E.$\delta$).
\emph{The case $k=0$.} Combine Lemma~\ref{app:G1:SC-quad} with $\eps\abs{\delta m^2_{K,T}(M)}\le C_\delta(M)$ (Proposition~\ref{app:K:L13-12}). The absorption of this loss by half of the step-$0$ gain is not used here.
The localisation cubes of side $3r(t)$ do not wrap, by Lemma~\ref{app:G1:nowrap}.

\emph{Step 2, (C2).} The decomposition \cite[(2.47)--(2.48)]{Ba2} and the identity \cite[(2.108)]{Ba2} are domain changes at a single background. By Proposition~\ref{app:K:L13-8} and \cite[(2.55)]{Ba2}, they cost $O(t^\kappa)$ per block of $\Lambda_6^{(k-1)}$, resp.\ of $\Lambda_3^{(k)}$, and \cite[(2.109)]{Ba2} holds. Since $M^2t^2\le M^2\eps_0^2\le1$ (condition (E.m2); Proposition~\ref{app:G1:I7}), the constants are uniform in $M\le m_2$.

\emph{Step 3, (C3).} By \cite[(2.105)]{Ba2}, $\prod_{j<k}Z^{(j)}=Z_k$. The fluctuation expansion and the re-localisation are Proposition~\ref{app:K:L13-7}, whose inputs are Lemma~\ref{app:K:L13-6}. Part (d) of that proposition also assumes $4w_j\le\gamma_0\le1$ for all $j<k$ and $\delta_5r(t_k)\ge d\log L+\log2$; the first is condition (E.W) of Table~\ref{p:tab:A4}, and the second is condition (R.Z) of Table~\ref{p:tab:A2}, since $r(t_k)\ge R$. The collar terms produced by the extension of \cite[Prop.~2.1]{Ba1} to the polynomial $\mathcal Q^{(k)}$ \cite[p.~579]{Ba2} are bounded by Lemma~\ref{app:K:L13-4} and Lemma~\ref{app:G1:SC-pow}. The final replacement $G_k(B^k(\Lambda_0^{(k)}),\tilde B)\to G_k(B^k(\Lambda_2^{(k)}),B^{(k+1),\eta})$ is taken at one field on $B^k(\Lambda_2^{(k)})$ (Proposition~\ref{app:K:L13-11}(ii)), so it is a domain change at a single real background, at distance $\ge r(t)$, and it costs $O(t^\kappa)\abs{\Lambda_7^{(k)}}$ (Lemma~\ref{app:G1:I3r-DGM}). By the admissibility of $\boldsymbol\Lambda_{k+1}$ \cite[p.~566]{Ba2}, $\Lambda_2^{(k)}\cap\Lambda_7^{(k)c}\subset\mathrm{Col}_k$. The $Z^{(j)}$ and $\mathcal Q^{(k)}$ are evaluated only through $\ell^2$-bounds at real backgrounds (Proposition~\ref{app:G1:I4}).

\emph{Step 4, (C4).} By Proposition~\ref{app:K:L13-11}(ii), the two propagators in \cite[Prop.~2.6]{Ba2} are taken at the same field on $B^k(\Lambda_2^{(k)})$, so the replacement is a domain change. Vertices localised in $B^k(\Lambda_7^{(k)})$ lie at distance $>5r(t)$ from $\Lambda_2^{(k)c}$ \cite[(2.8)]{Ba2}. Lemma~\ref{app:G1:I3r-DGM} bounds each propagator difference by $O(1)\mathrm{e}^{-c\,r(t)}=O(t^\kappa)$ for every $\kappa$, and the graph sums are those of \cite[(1.33)]{Ba3} (Proposition~\ref{app:G1:I5}). The translation \cite[(2.110)]{Ba2} is Gaussian algebra at $B^{(k+1),\eta}$; the identity \cite[(2.111)--(2.112)]{Ba2} follows from \cite[(2.41)]{Ba1} and the domain-change bounds for the scalar covariances of Proposition~\ref{app:K:L13-8}, since the covariance replaced there is a fluctuation covariance and not a propagator $G_k$.

\emph{Step 5, (C5).} This is Lemma~\ref{app:K:L13-5}, together with the restrictions on $\psi,\varphi$ \cite[(2.114)]{Ba2}. Consequently the cut-off $\chi(\varphi')$ of \cite[(2.114)]{Ba2} on $\Lambda_5^{(k)}$ may be inserted. All other scalar characteristic functions on $\Lambda_5^{(k)}$ except $\chi_{k+1}$ are bounded by $1$ \cite[p.~581]{Ba2}.

\emph{Step 6, (C6).} Apply Theorem~\ref{app:K:L13-1} once to the product measure of (C6), with the data of the proof of Lemma~\ref{app:G1:S05-PH}(b): $z_x=(A'(x),\varphi'(x))$, $K$ the product of two unit balls ($c_K=\sqrt2$), $J=\Lambda_7^{(k)}$, the precision in the class $\mathcal P_{\mathbb W_0}$ of Lemma~\ref{app:G1:PH-CL}, the scalar covariance bounds of \cite[(1.9)--(1.20)]{Ba4} (uniform on the torus and in the zero mode, Proposition~\ref{app:K:L13-8}), $\norm{V^{(k)}}_\kappa\le C\,t^{13/56}$ by \cite[(3.58)]{Ba1} and Proposition~\ref{app:G1:I5}, and $b\ge b_*$, $\norm{V^{(k)}}_\kappa b^{\rho_1}\le1$ by Lemma~\ref{lem:window}. The scalar half is $w$-free, and no complex gauge background is used (Proposition~\ref{app:G1:I4}). The two remainders are bounded by Lemma~\ref{app:K:L13-13}.

\emph{Step 7, (C7).} We follow \cite[p.~582]{Ba2}: terms of order $>\nbar$ are estimated, Dirichlet covariances are replaced by free ones at a cost of $O(t^\kappa)\abs{\Lambda_7^{(k)}}$ (Corollary~\ref{app:K:L13-2}, with the domain differences of the step covariances of Proposition~\ref{app:K:L13-8} and Lemma~\ref{app:G1:PH-O2}), and some are composed by \cite[(2.42)]{Ba1}. This yields $\mathcal P^{(k+1)}_w$. Graph values are multilinear in the lines, so the vector lines compose to $\sum_{j\le k}w_jG_{(j)}=G_{k+1}+uc_{k+1}D_{k+1}$ (proof of Lemma~\ref{app:G1:S05-PH}(c)). The bounds are those of Proposition~\ref{app:G1:I5}, applied under the conditions of Proposition~\ref{app:G1:I3}. The torus enters only through $\delta G$-type external fields (Proposition~\ref{app:G1:I2}); the constant of the $\delta m^2_{\rm fin}$ vertex is uniform in $\ell_T$ (Proposition~\ref{app:G1:I8}).

\emph{Step 8, (C8).} The field-independent output of order $\le\nbar$ is handled by Lemmas~\ref{app:G1:S05-TS} and \ref{app:G1:S05-RF}(V), (C11), completed by the restriction tails of Lemma~\ref{app:K:L13-13}(d). That lemma is applied with $\kappa_B=\frac{16}5>3$. The hypothesis of its part (c) holds with $\kappa_*=\frac{13}{56}$ and $\nbar=13$, since $14\cdot\frac{13}{56}=\frac{13}4>\frac{16}5$ (Step~6), and, by its part (d), the bound of Lemma~\ref{app:G1:S05-RF}(V) holds with $C_\kappa$ replaced by $C_\kappa+2^{1/5}\big(C_{\rm res}(\frac{16}5)+C_{\rm tr}\big)$; in $\mathfrak E_k$, $C_\kappa$ denotes this enlarged constant. Its order-$e^2$ part with one $D$-line is Proposition~\ref{app:G1:I1}; the bad localisations of the $C$-only families with $\delta m^2$ insertions are Proposition~\ref{app:G1:I9} with Lemma~\ref{app:G1:VS}. The collar shares are at most $C(\kappa_0)t^{\kappa_0}$ per block of $\mathrm{Col}_k$: the layer has at most $(2r(t)+2\sqrt dM_b+1)^3\abs{\mathrm{Col}_k}$ blocks, and Lemma~\ref{app:G1:SC-pow} applies.

\emph{Step 9, (C9).} The first item is Proposition~\ref{app:K:L13-9}; its error term is already $C(\kappa_0)t_k^{\kappa_0}$ per block of $S_k$. The second is Lemma~\ref{lem:B:massrange}(d). The third is (C2).

\emph{Step 10, (C10) and uniformity.} We follow \cite[p.~582]{Ba2}: rescale to the $L^{k+1}\eps$-lattice and collect the estimates into \cite[(2.43)]{Ba2} at $k+1$. The new scalar form on $\Lambda_6^{(k)\prime}$ at the background $A^{(k+1),\eps}$ is \cite[(2.47)]{Ba2} at $k+1$, obtained from \cite[(2.110)--(2.112)]{Ba2}. The factor $Z^{(k)}$ enters the density, $\chi_{k+1}$ and $\theta_{k+1}A^{(k+1)}$ are carried, and $\mathcal P^{(k+1)}_w$ is given by (C7). The errors of (C1)--(C8) are of the two types $C\,t^{\kappa_0}\abs{\mathrm{Col}_k}$ and $C\,t^{\kappa}N_k$, which is the scalar part of Step~3 of Proposition~\ref{app:G1:assembly}. The only factors growing with $k$ are $K-k=\log_2(\eps_0/t)$, and they always multiply $t^2$ or $t^3$ (Lemma~\ref{app:G1:SC-pow}). Ba{\l}aban's per-step constants are bounded uniformly in $\eps_0\in(0,\eps_1]$ (Statement~\ref{app:K:SL}). Dependence on $\ell_T$ enters only through lower bounds on $\ell_0$, through Lemma~\ref{app:K:L13-4}, and through $C_b$ (Proposition~\ref{app:G1:I8}). Masses enter through $M^2t^2\le1$ and the mass floor, which binds at $m_1>0$.
\end{proof}

\begin{remark}
(a) The bound in (C1) needs $\kappa_0<\frac14$: by \cite[p.~421]{Ba3}, the family of the vertex \cite[(1.6)]{Ba3} with three external scalar legs is of order $t^{1/4}p^3$. The counterterm and $\abs\varphi^6$ terms alone would allow any $\kappa_0<\frac12$.
(b) In (C4) we use the printed volume $\abs{\Lambda_7^{(k)}}$ of \cite[Prop.~2.6]{Ba2}.
(c) The scalar step does not use a block-by-block bound on the order-$e^2$ vacuum family with one $D$-line; that family is treated in Proposition~\ref{app:G1:I1}. A direct treatment cutting the index-$k$ lines \cite[p.~422]{Ba3} also works, with three amendments: the per-slice factor is $L^{-(k-j)}$, not $L^{-2(k-j)}$, so the constant gains a factor $L(1-L^{-1})^{-1}$; the torus lines of index $<k$ differ from the $\eta\Z^d$ lines by image fields of size $O(\mathrm{e}^{-\delta\ell_0/(2\eps_0)})$ for some $\delta>0$; the dumbbell graph vanishes, and the $e^nA^n$ vertices with $n\ge3$ first contribute at higher order.
\end{remark}

\begin{lemma}[photon small-field step]\label{app:G1:S05-PH}
Under the same hypotheses:
\begin{itemize}
\item[(a)] The vector translation \cite[(2.80)]{Ba2} and the lift are $w$-free. After the translation, $A'|_{\Lambda_5^{(k)}}\sim N(0,w_kC^{(k)}_{\Lambda_5})$ in the $A'$-gauge. The small-field restriction $\chi(A')$ is $w$-free, and $\abs{A'}\le O(1)p$ on $\Lambda_1^{(k)}$.
\item[(b)] The joint fluctuation integral is bounded by the cumulant bound of Theorem~\ref{app:K:L13-1}, used once on $\mathcal P_{\mathbb W_0}$.
\item[(c)] The new photon line is $G_{k+1}(u)=G_{k+1}+uc_{k+1}D^\varsigma_{k+1}$ exactly, with $c_k=\kappa\mu_*^2L^2t_k^2\le\frac13$. It satisfies the weighted tree bound \cite[(1.33)]{Ba3} block by block. No new counterterm and no new Ward requirement arise. The order-$e^2$ vacuum family with one $D$-line is \emph{not} bounded here; it is handed to Lemma~\ref{app:G1:S05-RF}.
\item[(d)] Closing at $k+1$: $w_k^{-1}+\omega_{k+1}=w_{k+1}^{-1}$. The replacements \cite[(2.47)--(2.49)]{Ba2} cost at most $\sigma_+=(1-\eta)^{-1}$ times Ba{\l}aban's. With $X_{k+1}(B):=\ip{\Lambda_7^{(k)\prime}B}{\Delta^{(k+1)}\Lambda_6^{(k)\prime c}B}$, the deletion of $\omega_{k+1}X_{k+1}$ costs $\nu(t_{k+1})\abs{\Lambda_7^{(k)\prime}}_{k+1}+\xi_{k+1}\norm{A_{k+1}|_{\Lambda_6^{(k)\prime c}}}^2$, with
\[
\nu(t)=\tfrac12\abs u\kappa L^2c_0c_dc_1^2w_{\min}^{-2}p(t)^2\mathrm{e}^{-\delta_0r(t)/2},\qquad \xi_{k+1}=\tfrac12\abs{\omega_{k+1}}c_0c_d\mathrm{e}^{-\delta_0r(t_{k+1})/2},
\]
where $c_0=c_0^\Delta$, $\delta_0$ are the constants of Lemma~\ref{app:G1:PH-L1}, $c_d=\sum_{z\in\Z^3}\mathrm{e}^{-\delta_0\abs z/2}$ and $c_1$ is the constant of \cite[(2.55)]{Ba2}. Put $HA^{(k)}_{\rm del}:=HA-2\sum_{1\le j\le k}\omega_jX_j(A_j)$, $\hat\xi_j:=\abs u\kappa L^2c_0c_dw_{\min}^{-2}(w_{\rm lo}\gamma_1)^{-1}\mathrm{e}^{-\delta_0r(t_j)/2}$ and $\hat\xi_{\rm tot}:=\sum_{1\le j\le K}\hat\xi_j$. Then $HA_{\rm del}^{(k+1)}=HA_{\rm del}^{(k)}-2\omega_{k+1}X_{k+1}$ and $(1-\hat\xi_{\rm tot})HA\le HA_{\rm del}^{(k+1)}\le(1+\hat\xi_{\rm tot})HA$.
\item[(e)] The backward objects satisfy $H'^w_{k+1}=w_k^{-1}H'_{k+1}\ge0$ and $\varrho_{k+1}\le\frac12$ ((E.4)). Lemma~\ref{app:G1:S05-ZM} holds (photon zero mode; $\mu_*>0$ is necessary and sufficient).
\end{itemize}
\end{lemma}
\noindent The bounds at complex backgrounds are excluded by Proposition~\ref{app:G1:I4}; the Ward--Takahashi part of the defect of the zeroth moment at orders $\ge2$ is Proposition~\ref{app:G1:I2}.

\smallskip\noindent\emph{Setting for the proof.} Throughout, $d=3$, $L=2$, $t_k=L^k\eps$, $\ell(t)=1+\log t^{-1}$, $r(t)=R\,\ell(t)^2$, $p(t)=b_0\ell(t)^{6}$. The weights are $w_j=1+u\delta_j$, with $\delta_j=\kappa\mu_*^2t_{j+1}^2\le\frac1{12}$ for $j<K$ (condition (E.$\delta$)) and $\delta_K:=0$, so that $w_K=1$; further $\omega_k=w_k^{-1}-w_{k-1}^{-1}$, $w_{\min}=\min_jw_j$, $w_{\max}=\max_jw_j$, $w_{\rm lo}=\min(1,\min_kw_k^{-1})$, $\sigma_-=\frac12$, $\sigma_+=(1-\eta)^{-1}$. In scale-$k$ units $M_k=\Delta^{(k)}+\hat\beta Q^*Q$ with $\hat\beta=aL^{-2}$, $C^{(k)}=M_k^{-1}$, $C^{(k)}_\Lambda=(\Lambda M_k\Lambda)^{-1}$, $\Delta^{(k+1)}=\beta_k-\beta_k^2QC^{(k)}Q^*$, $\mathcal H_k=\beta_kC^{(k)}Q^*$, $D_k(B)=\ip{B}{\Delta^{(k)}B}$. Distances are torus distances; sums over a torus are dominated by sums over $\Z^d$. The torus satisfies Lemma~\ref{lem:sizes}(ii) and (T.nw).

\begin{lemma}[kernel decay and lower bound of the step operators]\label{app:G1:PH-L1}
There are $c_0^\Delta,\delta_0>0$, depending on $d$, $M_b$, $a$, $\mu_*\eps_0$ only, such that for $0\le k\le K$
\[
\abs{\Delta^{(k)}(x,y)}\le c_0^\Delta\mathrm{e}^{-\delta_0\abs{x-y}},\qquad \abs{M_k(x,y)}\le c_M\mathrm{e}^{-\delta_0\abs{x-y}},\quad c_M:=c_0^\Delta+\hat\beta\mathrm{e}^{\delta_0L\sqrt d},
\]
$0\le\Delta^{(k)}\le a_k$ for $k\ge1$, and $\norm{\Delta^{(0)}}\le4d+\mu_*^2\eps_0^2$. Moreover $\Lambda M_k\Lambda\ge c_C^{-1}$ for every set $\Lambda$ (in particular $\Lambda=T^{(k)}$), with $c_C=4/a_\flat+L^2/8+L^2/a$, $a_\flat=a(1-L^{-2})$. All constants are independent of $k$, $\eps$, $\ell_T$ and $\Lambda$, and the constant $c_C$ of the lower bound is also independent of $\mu_*\ge0$.
\end{lemma}
\begin{proof}
Let $k\ge1$. By \cite[(1.14), p.~573]{Ba4} at $A=0$, $\Delta^{(k)}(y,y')_{ij}=a_k\delta_{yy'}\delta_{ij}-a_k^2\ip{q_k(y)_i}{G_kq_k(y')_j}$, where $q_k(y)$ has unit norm and support in $B^k(y)$ \cite[(3.7), p.~588]{Ba4}. By \cite[Cor.~2.3, (2.30), p.~580]{Ba4}, valid on the torus \cite[pp.~572, 575]{Ba4},
$\abs{\Delta^{(k)}(y,y')}\le a_k\delta_{yy'}+a_k^2c_0\mathrm{e}^{-\delta_0\dist(B^k(y),B^k(y'))}\le(a+a^2c_0\mathrm{e}^{\delta_0\sqrt d})\mathrm{e}^{-\delta_0\abs{y-y'}}$.
Put $Y=a_k^{1/2}Q_k$ and $X=-\Delta+m^2\ge0$. Then $a_kQ_kG_kQ_k^*=Y(X+Y^*Y)^{-1}Y^*\le I$, as the limit $\delta\downarrow0$ of $Y(X+\delta+Y^*Y)^{-1}Y^*=I-(I+Y(X+\delta)^{-1}Y^*)^{-1}$. Hence $0\le\Delta^{(k)}\le a_k$; this also covers tori too small for the cubes of \cite{Ba4}, at the cost of a factor $\mathrm{e}^{\delta_0D}$, $D$ the diameter. For $k=0$, $\Delta^{(0)}=-\Delta+\mu_*^2$ is nearest-neighbour with entries at most $2d+\mu_*^2\eps_0^2$, so the kernel bound holds with the constant $(2d+\mu_*^2\eps_0^2)\mathrm{e}^{\delta_0}$; take $c_0^\Delta$ to be the larger of this constant and $a+a^2c_0\mathrm{e}^{\delta_0\sqrt d}$. The kernel of $Q^*Q$ is bounded by $1$ and vanishes for $\abs{x-y}>L\sqrt d$, which gives $c_M$. The lower bound is step~(a) of the proof of Lemma~\ref{app:G1:WP-RBdiag} ($C^{(k)}\le c_Ct_k^2$ in physical units); for $k=0$ the same chain holds with the term $4/r_j$ absent. A principal block of an operator $\ge c_C^{-1}$ is $\ge c_C^{-1}$.
\end{proof}

\begin{lemma}[class closure]\label{app:G1:PH-CL}
Let $\mathcal P(\gamma_0,\gamma_1,c_0,\delta_0)$ be the set of real symmetric block operators $P$ with $\gamma_0\le P\le\gamma_1$ and $\norm{P(x,y)}\le c_0\mathrm{e}^{-\delta_0\abs{x-y}}$. If $P_A,P_\varphi\in\mathcal P(\gamma_0,\gamma_1,c_0,\delta_0)$ and $s\in[\sigma_-,\sigma_+]$, then $sP_A\oplus P_\varphi\in\mathcal P_{\mathbb W_0}:=\mathcal P(\sigma_-\gamma_0,\sigma_+\gamma_1,\sigma_+c_0,\delta_0)$.
\end{lemma}
\begin{proof}
Multiply the inequalities for $P_A$ by $s$; the scalar block is unchanged.
\end{proof}

\begin{lemma}[weighted lines]\label{app:G1:PH-WL}
Let $\mathcal L=\sum_i\mathcal L_i$ be a fluctuation vector line in the sense of Lemma~\ref{app:K:L13-14}, obtained from the step covariances $C^{(i)}_{\Lambda_5^{(i)}}$, written by the step of origin $i$, by the operations listed in Lemma~\ref{app:K:L13-14}. In the weighted chain the line at the same place is $\sum_iw_i\mathcal L_i$. Its $D$-part is $u\sum_i\delta_i\mathcal L_i$, and every bound used for $\mathcal L_i$ holds for the $D$-part with the extra factor $\abs u\delta_i\le\abs uc_k$.
\end{lemma}
\begin{proof}
By Lemma~\ref{app:G1:S05-PH}(a), the step-$i$ fluctuation has the law $N(0,w_iC^{(i)}_{\Lambda_5^{(i)}})$, and the translations, the lift and $\chi(A')$ are $w$-free. The conditional means are $w$-free by Lemma~\ref{app:G1:VL-gauss}(i). A line is a covariance pushed forward by such linear maps $\mathcal M,\mathcal M'$, and $\mathcal M(w_iC)\mathcal M'^*=w_i\,\mathcal MC\mathcal M'^*$. Restriction to domains, Dirichlet-to-free replacement, localisation cut-offs and differences are linear in the covariance and act step by step. Vector lines do not depend on the background, so the same holds for the $\delta G$-type differences, which concern scalar lines only.
\end{proof}

\begin{lemma}[per-block bounds of the $D$-coloured part]\label{app:G1:PH-U3}
Let $\abs uc_k\le1$, and let $N_\ell$ be the largest number of vector lines of a graph of order $\le\nbar$. Then the $D$-coloured part of $\mathcal P^{(k)}_w$ obeys the localised bounds \cite[(2.57), (2.116)]{Ba2} with constant at most $2^{N_\ell}\abs uc_k$ times the $u=0$ constant.
\end{lemma}
\begin{proof}
At $u=0$, (2.57) and (2.116) follow by summing the right side of \cite[(1.33)]{Ba3} over the families of order $\le\nbar$ with a vertex in a given block, using only the factors $e_k^{d_v}\lambda_k^{d_s}$, the decay $\mathrm{e}^{-\delta_0d}$, external-field norms bounded by the restrictions, and the finiteness of the set of graphs. By Theorem~\ref{app:G1:WP-WP1}(c), a coloured family has the same right side times $\prod_{D\text{-lines}}\abs uc_k(L^j\varsigma)^2\le(\abs uc_k)^m$. Colourings multiply the number of graphs by at most $2^{N_\ell}$.
\end{proof}

\begin{lemma}[step covariances: decay and domain differences]\label{app:G1:PH-O2}
(i) Choose $\delta'\in(0,\delta_0/2]$ with $c_M\sum_{z\in\Z^d}(\mathrm{e}^{\delta'\abs z}-1)\mathrm{e}^{-\delta_0\abs z}\le(2c_C)^{-1}$. Then $\abs{C^{(k)}_\Lambda(x,y)}\le2c_C\mathrm{e}^{-\delta'\abs{x-y}}$ for every set $\Lambda$.
(ii) For $x,x'\in\Lambda$, $(C^{(k)}-C^{(k)}_\Lambda)(x,x')=-\sum_{z\in\Lambda,w\notin\Lambda}C^{(k)}_\Lambda(x,z)M_k(z,w)C^{(k)}(w,x')$.
(iii) With $d_x=\dist(x,\Lambda^c)$, $c_d(s)=\sum_{z\in\Z^d}\mathrm{e}^{-s\abs z}$ and $K_1=4c_C^2c_Mc_d(\delta'/2)^2$, $\abs{(C^{(k)}-C^{(k)}_\Lambda)(x,x')}\le K_1\mathrm{e}^{-(\delta'/4)(\abs{x-x'}+d_x+d_{x'})}$.
\end{lemma}
\begin{proof}
(i) Let $\rho$ be $1$-Lipschitz and $M^\rho=\mathrm{e}^{\delta'\rho}\Lambda M_k\Lambda\mathrm{e}^{-\delta'\rho}=\Lambda M_k\Lambda+R$. Then $\abs{R(x,y)}\le(\mathrm{e}^{\delta'\abs{x-y}}-1)\abs{M_k(x,y)}$, so $\norm R\le(2c_C)^{-1}$ by the Schur test and Lemma~\ref{app:G1:PH-L1}, and $\norm{(M^\rho)^{-1}}\le(c_C^{-1}-\norm R)^{-1}\le2c_C$. Take $\rho=\dist(\cdot,y)$.
(ii) The rows of $M_kC^{(k)}=I$ in $\Lambda$ give $(\Lambda M_k\Lambda)C_{\Lambda\Lambda}+M_{\Lambda\Lambda^c}C_{\Lambda^c\Lambda}=I_\Lambda$; multiply by $C_\Lambda$.
(iii) By (ii), (i) and Lemma~\ref{app:G1:PH-L1}, the left side is at most $4c_C^2c_M\sum_{z,w}\mathrm{e}^{-\delta'E}$, $E=\abs{x-z}+\abs{z-w}+\abs{w-x'}$. Since $w\notin\Lambda$, $E\ge d_x+d_{x'}$, and $E\ge\abs{x-x'}$; so $\mathrm{e}^{-\delta'E}\le\mathrm{e}^{-(\delta'/4)(\abs{x-x'}+d_x+d_{x'})}\mathrm{e}^{-\delta'E/2}$, and $\sum_{z,w}\mathrm{e}^{-\delta'E/2}\le c_d(\delta'/2)^2$.
\end{proof}

\begin{lemma}[zero modes]\label{app:G1:PH-ZM}
(a) For every $\Lambda\subset T^{(k)}$, including $\Lambda=T^{(k)}$, $w_kC^{(k)}_\Lambda\le w_kc_C$ in scale-$k$ units, independently of $\mu_*\ge0$ and $\ell_T$.
(b) $HA$ is positive definite if and only if $\mu_*>0$. Then its only candidate flat direction, $A_K$ constant with the finer fields at their conditional means, has precision $\ip{c1}{\Delta^{(K)}c1}\ge\gamma_1\mu_*^2\abs T\abs c^2$, and the mean $\bar A$ of $A$ has variance in $[\min(1,w_{\min}),\max(1,w_{\max})]\,(\mu_*^2\abs T)^{-1}$ per component.
(c) The constant mode $c$ of a real background enters no scalar operator bound used in the step: the bounds of \cite[(1.8), (1.9)--(1.12), p.~573; (1.15)--(1.22), p.~574]{Ba4} are stated for subsets of a torus with periodic conditions \cite[p.~572]{Ba4}, uniformly in the background, and since $\partial^\eta c=0$, $c+A''$ is regular iff $A''$ is, with the same constants.
(d) The event $\abs{\bar A}\ge c_1p(t)/(2\mu_*t)$, for one scale $t\le\eps_0$, has probability at most $6\exp(-c_1^2p(t)^2\abs T/(24w_{\max}t^2))$.
\end{lemma}
\begin{proof}
(a) Lemma~\ref{app:G1:PH-L1}.
(b) By Lemma~\ref{app:G1:S05-T}, $HA=\sum_{k<K}w_k^{-1}\ip{\hat S_k}{M_k\hat S_k}+D_K(A_K)$ with $\hat S_k=A_k-\mathcal H_kA_{k+1}$, a unit-triangular change of variables. Since $M_k\ge c_C^{-1}$, $HA>0$ iff $\Delta^{(K)}>0$. If $\mu_*>0$, $\Delta^{(K)}\ge\gamma_1\mu_*^2$ by Lemma~\ref{app:G1:S05-ZM}. If $\mu_*=0$, then $\Delta^{(0)}1=0$, $M_01=\beta_01$ gives $\Delta^{(1)}1=0$, and inductively $\Delta^{(K)}1=0$, so $\mathrm{e}^{-\frac12HA}$ is not integrable along the constants. For the variance, $C_B(u)=\sum_{j<K}w_jH_j+G_K$ and $C_{\mu_*}=\sum_{j<K}H_j+G_K$ with $H_j,G_K\ge0$, so $\min(1,w_{\min})C_{\mu_*}\le C_B(u)\le\max(1,w_{\max})C_{\mu_*}$; moreover $\operatorname{Var}\bar A=\abs T^{-2}\ip1{C_B1}$ and $C_{\mu_*}1=\mu_*^{-2}1$.
(c) The cited statements, together with $\partial^\eta c=0$.
(d) By (b), each of the three components of $\bar A$ is centred Gaussian with variance at most $w_{\max}(\mu_*^2\abs T)^{-1}$; if $\abs{\bar A}\ge s$, some component is at least $s/\sqrt3$. Use the Gaussian tail with $s=c_1p(t)/(2\mu_*t)$.
\end{proof}

\begin{lemma}[backward replacement, vector sector]\label{app:G1:PH-D}
Let $\Lambda_5:=\Lambda_5^{(k+1)}\subset\Lambda_0^{(k+1)}\subset\Lambda_7^{(k)\prime}$ and, in scale-$(k+1)$ units,
$H'_{k+1}:=\Lambda_5(\Delta^{(k+1)}_{\Lambda_5^{(k)}}-\Delta^{(k+1)})\Lambda_5=\beta_k^2\Lambda_5Q(C^{(k)}-C^{(k)}_{\Lambda_5^{(k)}})Q^*\Lambda_5$.
(i) $H'_{k+1}\ge0$ and $\norm{H'_{k+1}}\le C_{H'}\mathrm{e}^{-\delta'r(t_k)}$, $C_{H'}=\hat\beta^2K_1\mathrm{e}^{\delta'L\sqrt d}c_d(\delta'L/4)$.
(ii) In the weighted chain the difference between the exact and the replaced Hessian is $H'^w_{k+1}=w_k^{-1}H'_{k+1}$, and the linear part is $F'^w_{k+1}=w_k^{-1}F'_{k+1}$.
(iii) The exact Hessian $G^w_{k+1}$ satisfies $G^w_{k+1}\ge G^w_{k+1}-H'^w_{k+1}\ge(1-\hat\xi_{\rm tot})w_{\rm lo}\gamma_1\mu_*^2t_{k+1}^2$.
(iv) $X^w:=G^{w,-1/2}H'^wG^{w,-1/2}$ satisfies $0\le X^w\le\varrho^{\det}_{k+1}$ with $\varrho^{\det}_{k+1}\le\frac{24}{11}(1-\eta)^{-1}C_{H'}\gamma_1^{-1}\mu_*^{-2}t_{k+1}^{-2}\mathrm{e}^{-\delta'r(t_k)}\le\frac12$ under (E.4) and $\eps_0\le\exp(1-1/(\delta'R))$; and $1\le[\det(I-X^w)]^{-1/2}\le\exp(d\varrho^{\det}_{k+1}\abs{\Lambda_5})$.
\end{lemma}
\begin{proof}
(i) The identity is the definition of $\Delta^{(k+1)}$ and of its conditional version. Positivity: $\ip f{C_\Lambda f}=\sup_{g\subset\Lambda}[2\ip fg-\ip g{M_kg}]\le\ip f{C^{(k)}f}$. The fine sites of the blocks of $\Lambda_5$ satisfy $d_x\ge2r(t_k)-L\sqrt d$ by \cite[(2.8)]{Ba2}; Lemma~\ref{app:G1:PH-O2}(iii), the averaging weights of $Q$ and the Schur test give the bound.
(ii) The variables $A_{k+1}|_{\Lambda_5}$ occur only in the carried form $w_{k+1}^{-1}\Lambda_6^{(k)\prime}\Delta^{(k+1)}\Lambda_6^{(k)\prime}$, in the exact composed kernels of later steps, and in scalar forms through the external field $\tilde A$; the rest of the density depends on $A_{k+1}$ only through $\Lambda_7^{(k)\prime c}A_{k+1}$ (step (d1) of the proof of Lemma~\ref{app:G1:S05-PH} below). On $\Lambda_5^{(k)}$ the $A'$-Gaussian is $w_k^{-1}$ times Ba{\l}aban's (Lemma~\ref{app:G1:S05-PH}(a)), and minimisation is homogeneous. Hence the exact step-$k$ output, as a form in $A_{k+1}$, is $w_k^{-1}$ times Ba{\l}aban's exact form plus $\omega_{k+1}(D_{k+1}-2X_{k+1})$, and the replaced one is $w_k^{-1}$ times his replaced form plus the same term. Subtract. The blocks of $X_{k+1}$ lie in $\Lambda_7^{(k)\prime}\times\Lambda_6^{(k)\prime c}$ and its transpose, so $X_{k+1}$ has no block on $\Lambda_5\times\Lambda_5$.
(iii) The exact Hessian is the principal $\Lambda_5$-block of the Hessian of the carried form of the chain from level $k+1$, minimised over the coarser variables. By Lemma~\ref{app:G1:S05-T} and the sandwich of Lemma~\ref{app:G1:S05-PH}(d) (step (d2) of its proof below, which does not use this lemma), this form is at least $(1-\hat\xi_{\rm tot})w_{\rm lo}$ times the unweighted chain $\sum_{l>k}I_l+D_K$. The identity $D_l(B)+\beta_l\norm{QB-B'}^2=I_l(B,B')+D_{l+1}(B')$, minimised over $B'$, shows inductively that the minimum of the unweighted chain over the coarser variables is $D_{k+1}(A_{k+1})$. Inequalities between forms pass to minima and principal blocks, and $\lambda_{\min}(\Delta^{(k+1)})\ge\gamma_1\mu_*^2t_{k+1}^2$ (Lemma~\ref{app:G1:S05-ZM}). Finally $G\ge G-H'$ by (i).
(iv) $\norm{X^w}\le w_k^{-1}\norm{H'_{k+1}}/((1-\hat\xi_{\rm tot})w_{\rm lo}\gamma_1\mu_*^2t_{k+1}^2)$, with $w_k^{-1}\le\sigma_+$, $w_{\rm lo}\ge\sigma_-$ and $(1-\hat\xi_{\rm tot})^{-1}\le\frac{12}{11}$ (condition (E.14)). Since $r(t_k)\ge r(t_{k+1})$, the right side is at most $g(t_{k+1})$, $g(t)=\mathrm{const}\cdot t^{-2}\mathrm{e}^{-\delta'R\ell(t)^2}$, which is increasing on $(0,\exp(1-1/(\delta'R))]$; so $g(t_{k+1})\le g(\eps_0)\le\frac12$ by (E.4). Finally $-\frac12\log\det(I-X)=\frac12\sum_{n\ge1}\tr X^n/n\le\frac12\tr X/(1-\varrho)\le d\abs{\Lambda_5}\varrho$, and the ratio is $\ge1$ since $X\ge0$.
\end{proof}
\noindent The linear terms are Ba{\l}aban's estimates \cite[(3.13)--(3.17), pp.~585--586]{Ba2} with $\gamma_1$ replaced by $(1-\hat\xi_{\rm tot})w_{\rm lo}\gamma_1$ and $\abs{F'^w}\le\sigma_+\abs{F'}$: $\frac12\ip\Phi{[(G-H')^{-1}-G^{-1}]\Phi}=\frac12\ip\Phi{G^{-1}H'(G-H')^{-1}\Phi}=O(t_{k+1}^\kappa)\abs{\Lambda_5}$ for every $\kappa$, and similarly for the $F'$ terms.

\begin{proof}[Proof of Lemma~\ref{app:G1:S05-PH}]
\emph{(a).} By Lemma~\ref{app:G1:S05-T} and (IH$_k$), the terms of the vector exponent that depend on $A_k|_{\Lambda_0^{(k)}}$ are $\frac12w_k^{-1}\ip{\Lambda_6A_k}{\Delta^{(k)}\Lambda_6A_k}$ and $\frac12w_k^{-1}\beta_k\norm{QA_k-A_{k+1}}^2$; the spectator $\frac12\omega_{k+1}D_{k+1}(A_{k+1})$ does not involve $A_k$. A common positive factor of the quadratic and linear parts does not move the minimiser, so the translation \cite[(2.80)]{Ba2} is Ba{\l}aban's and on $\Lambda_5^{(k)}$ the form in $A'$ is $w_k^{-1}$ times $\frac12\ip{\Lambda_5^{(k)}A'}{(C^{(k)}_{\Lambda_5^{(k)}})^{-1}\Lambda_5^{(k)}A'}$ \cite[p.~575]{Ba2}. The restriction \cite[(2.86)]{Ba2}, which gives $\abs{A'}\le O(1)p$ on $\Lambda_1^{(k)}$, and the split \cite[(2.88)--(2.91)]{Ba2} are statements about the configuration and the lift, hence $w$-free.

\emph{(b).} In scale-$k$ units $P_A:=(C^{(k)}_{\Lambda_5^{(k)}})^{-1}=\Lambda_5^{(k)}M_k\Lambda_5^{(k)}$ satisfies $c_C^{-1}\le P_A\le\gamma_1^A:=\max(a,4d+\mu_*^2\eps_0^2)+\hat\beta\norm Q^2$ and $\abs{P_A(x,y)}\le c_M\mathrm{e}^{-\delta_0\abs{x-y}}$, uniformly in $\Lambda_5^{(k)}$, $k$, $\eps$, $\ell_T$, $\mu_*$ (Lemma~\ref{app:G1:PH-L1}). The joint precision is $w_k^{-1}P_A\oplus P_\varphi$, with $P_\varphi$ the scalar precision of Lemma~\ref{app:G1:S05-SC}(C6), which depends on $B^{(k+1)}$ and not on $A'$. Since $0\le\delta_k\le\frac1{12}$ and $u\in\{1,-\eta\}$, we have $w_k^{-1}\in[\frac{12}{13},(1-\frac\eta{12})^{-1}]\subset[\sigma_-,\sigma_+]$, and Lemma~\ref{app:G1:PH-CL} puts it in $\mathcal P_{\mathbb W_0}$ with $\gamma_0=\min(c_C^{-1},\gamma_0^\varphi)$, $\gamma_1=\max(\gamma_1^A,\gamma_1^\varphi)$, $c_0=\max(c_M,c_0^\varphi)$, $\delta_0=\min(\delta_0,\delta_0^\varphi)$. The interaction $V^{(k)}_w$ is a real polynomial in $(A',\varphi')$ of degree $\le n(\nbar)$ \cite[Prop.~3.2, (3.57), p.~622]{Ba1}, \cite[p.~581]{Ba2}, with coefficients depending on the real background only; they obey \cite[(3.58)]{Ba1} with rate $\kappa_*=\frac14-\frac1{56}=\frac{13}{56}$, times the colouring factor $(1+\abs uc_k)^{N_v}$ of Proposition~\ref{app:G1:I5}, which is at most $(4/3)^{N_v}$ since $c_k\le\frac13$ by condition (E.$\delta$). So $\mathsf A_w:=\norm{V^{(k)}_w}_\kappa\le C_At_k^{13/56}$. The restrictions $\chi(A')\chi(\varphi')$ of \cite[(2.86), (2.114), p.~581]{Ba2} are $w$-free and of the form $\chi_b$ with $b=c_bp(t_k)$, $K$ the product of two unit balls. Theorem~\ref{app:K:L13-1}, applied once in $\mathcal P_{\mathbb W_0}$ with $t=\nbar=13$, $b\ge b_*$ and $\mathsf A_wb^{\rho_1}\le1$ (Lemma~\ref{lem:window}), gives
\[
\int d\mu_{w_kC^{(k)}_{\Lambda_5^{(k)}}}(A')\int d\mu_{C^{(k)}_{\Lambda_5^{(k)}}(\cdot)}(\varphi')\,\chi(A')\chi(\varphi')\mathrm{e}^{V^{(k)}_w}\le\exp\Big\{\sum_{m\le\nbar}\frac{\kappa(V^{(k)}_w;m)}{m!}+S\abs{\Lambda_7^{(k)}}\big[(\mathsf A_wb^{\rho_1})^{14}+\mathrm{e}^{-\rho_3b^{3/2}}\big]\Big\}.
\]
Now $(\mathsf A_wb^{\rho_1})^{14}\le C\,t_k^{13/4}\ell(t_k)^{84\rho_1}$ with $\sup_{t\le\eps_0}t^{1/20}\ell(t)^{84\rho_1}<\infty$, and $\mathrm{e}^{-\rho_3b^{3/2}}=\exp(-\rho_3(c_bb_0)^{3/2}\ell(t_k)^9)\le C_Nt_k^N$. So the remainder is at most $C_{\rm R}t_k^{16/5}\abs{\Lambda_7^{(k)}}$, with $C_{\rm R}$ depending only on the parameters and $\eta$ (see also Lemma~\ref{app:K:L13-13}). Only this upper bound is used, and every background at which a bound is invoked is real (Proposition~\ref{app:G1:I4}).

\emph{(c).} \emph{The line.} Each $A'A'$ contraction carries $w_kC^{(k)}_{\Lambda_5^{(k)}}$. Corollary~\ref{app:K:L13-2}, with collar $\ge2r(t_k)$, replaces it by $w_kC^{(k)}$ at a cost $O(t_k^{\kappa})\abs{\Lambda_7^{(k)}}$ for every $\kappa$. The interaction sees $A'$ only through the $w$-free lift $F_k$, so the new line is $w_kF_kC^{(k)}F_k^*=w_kG_{(k)}$. Vector lines are background-free, so the composition \cite[(2.42)]{Ba1} (cf.\ \cite[p.~582]{Ba2}) is linear:
$\sum_{j\le k}w_jG_{(j)}=G_{k+1}+u\sum_{j\le k}\delta_jG_{(j)}=G_{k+1}+u\,c_{k+1}D^\varsigma_{k+1}$,
because $c_{k+1}(L^j\varsigma')^2=\kappa\mu_*^2L^{2j+2}\eps^2=\delta_j$ with $\varsigma'=L^{-k-1}$. Every graph value is multilinear in its lines, so $\mathcal P^{(k+1)}_w=\sum_m(uc_{k+1})^m\mathcal P^{(k+1)}_m$ exactly.
\emph{Bounds.} Each coloured family obeys Theorem~\ref{app:G1:WP-WP1}(c), with Ba{\l}aban's rate and constants independent of $\eps$, $k$, the domains, the background, $u$ and $\kappa$. By Lemma~\ref{app:G1:PH-U3} with $c_{k+1}\le\frac13$, the bounds \cite[(2.57), (2.116)]{Ba2} hold for the $D$-coloured part with constant $\le2^{N_\ell}\abs uc_{k+1}$ times the $u=0$ constant, and for all of $\mathcal P^{(k+1)}_w$ with $(4/3)^{N_v}$ times it. In the remaining weighted uses (the domain changes \cite[(2.32)--(2.42)]{Ba2}, \cite[Props.~3.1--3.2]{Ba1}, the output of \cite[p.~582]{Ba2} and \cite[(2.117)--(2.118)]{Ba2}) $u$ enters only through line weights (Lemma~\ref{app:G1:PH-WL} and Lemma~\ref{app:K:L13-14}) and family factors. The expansions of $Z^{(j)}$ \cite[(2.92)--(2.107)]{Ba2} contain no vector line.
\emph{Counterterms.} The vertex list of $\mathcal P^{(k+1)}_w$ is Ba{\l}aban's $u=0$ list, with the scale split of the fixed $\delta m^2_{\rm ref}$ \cite[(1.26)--(1.28)]{Ba3} at $k+1$; passing from $k$ to $k+1$ moves $C$-only, $u$-independent pieces between the terms of the split. By Lemma~\ref{app:G1:WP-deg}(c),(d), every connected coloured subgraph with an external leg and $m\ge1$ has $D^\natural\ge1$, and every divergent subgraph is $C$-only and one of Ba{\l}aban's classes; by Lemma~\ref{app:G1:WP-fam} these are renormalised by $\delta m^2_{\rm ref}$ and their Ward--Takahashi and $\varphi\varphi A$ partners exactly as at $u=0$. No $u$-dependent vertex arises. Transversality is needed only for $C$-only vacuum-polarisation subgraphs \cite[(3.25)--(3.32)]{Ba3}, and on $T_\eps$ it holds exactly by Proposition~\ref{app:G1:I2}; a vacuum-polarisation subgraph with an internal $D$-line has $D^\natural\ge2$.
\emph{The excluded family.} The only $D$-coloured family with $D^\natural=0$ is the vacuum family with one $D$-line at order $e^2$. At a fixed largest index it is $O(e^2\abs u\kappa\mu_*^2L^2t_k^3)$ per block before cancellation, so booked on every block it would sum to $O(\abs T\log(\eps_0/\eps))$. It is not bounded here; it cancels identically in $E_1(u)$, and step by step up to the bad localisations (Lemmas~\ref{app:G1:S05-RF} and~\ref{app:G1:S05-TS}).

\emph{(d).} Write $B:=A_{k+1}$, $\Lambda_6':=\Lambda_6^{(k)\prime}$, $\Lambda_7':=\Lambda_7^{(k)\prime}$; $\Lambda_7'\subset\Lambda_6'$ and $\dist(\Lambda_7',\Lambda_6'^c)\ge r(t_{k+1})$ \cite[(2.8)]{Ba2}. The identity $w_k^{-1}+\omega_{k+1}=w_{k+1}^{-1}$ is the definition of $\omega_{k+1}$.
\emph{(d1) Absorption.} With $\Sigma_{k+1}(B):=\ip{\Lambda_6'^cB}{\Delta^{(k+1)}\Lambda_6'^cB}+2\ip{(\Lambda_6'\setminus\Lambda_7')B}{\Delta^{(k+1)}\Lambda_6'^cB}$ we have $D_{k+1}(B)=\ip{\Lambda_6'B}{\Delta^{(k+1)}\Lambda_6'B}+\Sigma_{k+1}(B)+2X_{k+1}(B)$. Adding the spectator $\frac12\omega_{k+1}D_{k+1}(B)$ of $\mathcal K^w_k$ (Lemma~\ref{app:G1:S05-T}) to the carried form gives exactly
$\frac12w_k^{-1}\ip{\Lambda_6'B}{\Delta^{(k+1)}\Lambda_6'B}+\frac12\omega_{k+1}D_{k+1}(B)=\frac12w_{k+1}^{-1}\ip{\Lambda_6'B}{\Delta^{(k+1)}\Lambda_6'B}+\frac12\omega_{k+1}\Sigma_{k+1}(B)+\omega_{k+1}X_{k+1}(B)$.
The first term is the carried form at $k+1$. $\Sigma_{k+1}$ depends on $B$ only through $\Lambda_7'^cB$, which is never integrated forward ($\Lambda_5^{(k+1)}\subset\Lambda_0^{(k+1)}\subset\Lambda_7'$); it joins the vector part of $\rho^{(k+1)}_w$.
\emph{(d2) Deletion.} By Lemma~\ref{app:G1:PH-L1}, $\abs{x-y}\ge r(t_{k+1})$ for $x\in\Lambda_7'$, $y\in\Lambda_6'^c$, and the Schur test,
$\abs{\omega_{k+1}X_{k+1}(B)}\le\frac12\abs{\omega_{k+1}}c_0c_d\mathrm{e}^{-\delta_0r(t_{k+1})/2}(\norm{B|_{\Lambda_7'}}^2+\norm{B|_{\Lambda_6'^c}}^2)$.
Since $\abs{\omega_{k+1}}=\abs u\abs{\delta_{k+1}-\delta_k}/(w_kw_{k+1})\le\abs u\kappa L^2\mu_*^2t_{k+1}^2w_{\min}^{-2}$ (also for $k+1=K$, where $\delta_K=0$) and $\abs B\le c_1p(t_{k+1})/(\mu_*t_{k+1})$ on $\Lambda_7'$ by the restrictions \cite[(2.55)]{Ba2} at scale $k+1$, the first part is $\le\nu(t_{k+1})\abs{\Lambda_7'}_{k+1}$ and the second is $\xi_{k+1}\norm{B|_{\Lambda_6'^c}}^2$. The term $\omega_{k+1}X_{k+1}$ is deleted from the carried form, so $HA^{(k+1)}_{\rm del}=HA^{(k)}_{\rm del}-2\omega_{k+1}X_{k+1}$. For the sandwich, the same Schur test gives $\abs{2\omega_jX_j(A_j)}\le\abs{\omega_j}c_0c_d\mathrm{e}^{-\delta_0r(t_j)/2}\norm{A_j}^2$. By Lemma~\ref{app:G1:S05-T}, $HA$ is at least $w_{\rm lo}$ times the unweighted chain, whose minimum over all variables other than $A_j$ is at least $D_j(A_j)\ge\gamma_1\mu_*^2t_j^2\norm{A_j}^2$ (the minimisation identity of the proof of Lemma~\ref{app:G1:PH-D}(iii), and Lemma~\ref{app:G1:S05-ZM}). Hence $\abs{2\omega_jX_j(A_j)}\le\hat\xi_jHA$, the factor $\mu_*^2t_j^2$ cancelling, and $\abs{HA^{(k+1)}_{\rm del}-HA}\le\hat\xi_{\rm tot}HA$ as forms.
\emph{(d3) Replacement.} After the step-$k$ integration the minimised form is $w_k^{-1}$ times Ba{\l}aban's (part (a)), and his replacement \cite[(2.47)--(2.49)]{Ba2} of it, with Neumann conditions not introduced \cite[p.~568]{Ba2}, costs $w_k^{-1}O(t_k^\kappa)\abs{\Lambda_5^{(k)}}\le\sigma_+$ times his cost: by Lemma~\ref{app:G1:PH-D}(ii),(iv) and the bound on the linear terms after it.

\emph{(e).} The identities and bounds for $H'^w_{k+1}$, the floor and $\varrho^{\det}_{k+1}\le\frac12$ are Lemma~\ref{app:G1:PH-D}; the zero-mode statements are Lemma~\ref{app:G1:PH-ZM}. Inside the step the constant mode enters only in four places: (i) through $\lambda_{\min}(\Delta^{(k+1)})$, where in the deletion coefficients of (d) the factor $\mu_*^2$ cancels against $\abs{\omega_{k+1}}$ and survives only in $\varrho^{\det}$ (condition (E.4)); (ii) through the restriction $\abs A\le c_1p/(\mu_*t)$ of \cite[(2.55), p.~570]{Ba2}, Lemma~\ref{app:G1:PH-ZM}(d) being only a size estimate; (iii) through the scalar conditions, which are uniform in it by Lemma~\ref{app:G1:PH-ZM}(c); (iv) through c-numbers that do not increase with $\ell_T$, bounded at the floor by Lemma~\ref{app:G1:floor}.
\end{proof}

\paragraph{(S3) The large-field part of the step and the region sum.}

\paragraph{The large-field step: notation.}\label{app:G1:part-LF}
In this paragraph and in Lemmas~\ref{app:G1:S05-LF}--\ref{app:G1:LF-sigma}, the letters $j$ and $k$ are scale indices. We have $d=3$ and $L=2$. Put $t_k=L^k\eps$, $\ell(t)=1+\log t^{-1}$, $\ell_k=\ell(t_k)$, $\hat\ell=\ell(\eps_0)$, $r_k=r(t_k)=R\ell_k^{\varrho}$ with $\varrho=2$, and $p_k=p(t_k)=b_0\ell_k^{6}$. $T^{(k)}$ is the unit lattice of scale $k$. It has $N_k=\abs Tt_k^{-3}$ points, and $\abs X_k$ is the number of points of $X\subset T^{(k)}$. Distances are $\ell^\infty$ distances on the torus, and $\mathcal N_s(X)$ is the closed $s$-neighbourhood of $X$. For a union $X$ of $L$-blocks of $T^{(k-1)}$, $X'\subset T^{(k)}$ is its image at scale $k$ \cite[p.~566]{Ba2}. For a set $X$ at scale $k$, $B^j(X)$ is the corresponding set at scale $k-j$, and $B(x):=B^1(\{x\})$. A large block of $T^{(k)}$ is a cube of side $b:=M_bL$. By Lemma~\ref{lem:sizes}, $\ell_T/(M_bt_k)\in2\N$ for $k\le K$, so $T^{(k)}$ is a union of large blocks. We assume $\vartheta_M:=M_bL/R\le\frac17$. Since $r_k\ge R$, this gives $L\le b\le\vartheta_Mr_k$.

A \emph{family} at scale $k$ is a six-tuple $\mathfrak t=(P_v,Q_v,R_v,P_s,Q_s,R_s)$ whose entries are sets of coarse points, of bonds or plaquettes, and of points, as in \cite[(2.2)--(2.5), p.~557]{Ba2} at scale $t_k$. Put $\abs{\mathfrak t}:=\abs{P_v}+\dots+\abs{R_s}$ and $S(\mathfrak t):=B(P_v)\cup Q_v\cup R_v\cup B(P_s)\cup Q_s\cup R_s$. The sets $\Lambda_i^{(k)}$, $i\le8$, are built from $\Lambda_0^{(k)}$ by \cite[(2.8)]{Ba2} at radius $r_k$ \cite[p.~566]{Ba2}, and $\Lambda_i^{(-1)\prime}:=T^{(0)}$. We write $X_k:=\Lambda_0^{(k)c}$, $\mathrm{Col}_k:=\Lambda_7^{(k-1)\prime}\cap\Lambda_7^{(k)c}$ (the collar), $\mathrm{Sh}_k:=\Lambda_5^{(k-1)\prime}\cap\Lambda_5^{(k)c}$ (the shell) and $W_k:=B^k(\Lambda_7^{(k)})$. Finally $n_d:=4+2\max\{d,d(d-1)/2\}=10$.

\begin{lemma}[the large-field part of the step]\label{app:G1:S05-LF}
Assume the same hypotheses as above, that is, (IH$_k$) of \eqref{app:G1:eq-IH} with $0\le k\le K-1$, at either point $(u,M)\in\{(1,m_2),(-\eta,m_1)\}$.
\begin{itemize}
\item[(a)] The scale-$t_k$ restrictions \cite[(2.2)--(2.6), (2.55)]{Ba2} are conditions on field values in $A'$-coordinates. They do not involve the weights. $\Lambda_0^{(k)}$ is the union of the large blocks of $T^{(k)}$ that are contained in $\Lambda_7^{(k-1)\prime}$ and lie at distance $>r_k$ from $S(\mathfrak t_k)$, for a minimal admissible family $\mathfrak t_k$ (Lemma~\ref{app:G1:LF-fam}). So it is a function of the past and of $\mathfrak t_k$. The region expansion \cite[(2.9)--(2.10)]{Ba2} holds in the direction ``$\le$''.
\item[(b)] On $\mathrm{Col}_k$ the interaction is removed under the restrictions \cite[(2.55)]{Ba2}. This includes the $D$-coloured part and the mass counterterm pieces of step $k$. No vacuum constant is removed on $\mathrm{Col}_k$. The cost is $c_C^{(k)}t_k^{\kappa_0}\abs{\mathrm{Col}_k}_k$. For $k\ge1$, $c_C^{(k)}\le\sigma_+2^{N_\ell}c_C^{u=0}$, where $c_C^{u=0}$ is the collar constant of \cite[Prop.~2.1, (2.57)]{Ba2} at $u=0$ and $\sigma_+$ comes from Lemma~\ref{app:G1:S05-PH}(d). At $k=0$ the cost is at most $C_\delta(M)^2\eps/(16\lambda)$ per site, so $c_C^{(0)}:=C_\delta(M)^2/(16\lambda)$ may be taken. The moved vacuum pieces localised outside $W_k$ cost at most $C_{\rm old}t_k$ per block of $\Lambda_7^{(k)c}$, with $C_{\rm old}$ as in Lemma~\ref{app:G1:S05-RF}(V).
\item[(c)] The forward clause carries the characteristic functions $\chi^c$, which do not involve the weights, and the quartic factor. The quartic factor produces $\exp(-p_k^2\abs{R_s})$. No vector gain is carried forward.
\item[(d)] The scale-$k$ factor delivered by the backward integration is at most $\zeta''_w(\Lambda_0^{(k)})$. Here $W_k^w$ is as in Proposition~\ref{app:G1:LF-sigma}, $\vartheta=\frac12(1-\hat\xi_{\rm tot})w_{\rm lo}$, $\vartheta_+\ge\frac{11}{26}$ at $(1,m_2)$ and $\vartheta_-\ge\frac{11}{24}$ at $(-\eta,m_1)$. Each family element therefore gains $\mathrm{e}^{-c_0'^{\,w}p_k^2}$, with $c_0'^{\,w}=\vartheta\min\{\frac18a,\frac14\gamma_0\}$ and $\gamma_0$ as in Proposition~\ref{app:G1:LF-V}. Each shell site costs at most $\mathrm{e}^{c_B^w}$, where $c_B^w=c_B+\frac d2\log3\cdot c_{\rm sh}$, $c_B$ is the shell constant of \cite[(3.42)]{Ba2} at $u=0$, and $c_{\rm sh}$ is the number of kept vector variables per shell site.
\item[(e)] The region sum is bounded by Proposition~\ref{app:G1:LF-sigma} under \eqref{app:G1:eq-SigmaW}, or by Lemma~\ref{app:G1:I145-comb} under (E.comb) with $c_A+C_{\rm old}$.
\end{itemize}
All constants are independent of $\eps$, $k$, $\boldsymbol\Lambda$ and $\ell_T\ge\ell_0$, and one $\eps_0$ serves both points.
\end{lemma}
\noindent The region expansion is used in the direction ``$\le$'' only. It is printed with ``$=$'' in \cite[(2.9)--(2.10)]{Ba2}, and equality can fail (Lemma~\ref{app:G1:LF-fam}(iii)).

\begin{lemma}[regions are functions of their families]\label{app:G1:LF-fam}
For $Y\subset T^{(k)}$ and a family $\mathfrak t$, let $\mathfrak L_k(Y,\mathfrak t)$ be the union of the large blocks of $T^{(k)}$ that are contained in $Y$ and lie at distance $>r_k$ from $S(\mathfrak t)$. Call $\mathfrak t$ admissible for $\Lambda_0^{(k)}$ if $\Lambda_0^{(k)}=\mathfrak L_k(\Lambda_7^{(k-1)\prime},\mathfrak t)$ \cite[p.~570]{Ba2}. At $k=0$ the distance condition is read as in \cite[(2.7)]{Ba2}.
\begin{itemize}
\item[(i)] Define $\boldsymbol\Lambda(\mathfrak t)$ recursively by $\Lambda_0^{(k)}:=\mathfrak L_k(\Lambda_7^{(k-1)\prime},\mathfrak t_k)$. This assigns to each tuple $(\mathfrak t_0,\dots,\mathfrak t_{K-1})$ at most one region sequence. Every pair $(\boldsymbol\Lambda,(\mathfrak t_k)_k)$ in which each $\mathfrak t_k$ is admissible for $\Lambda_0^{(k)}$ is of the form $(\boldsymbol\Lambda(\mathfrak t),\mathfrak t)$.
\item[(ii)] For all weights $\omega_k\ge0$ on families and every $\Phi\ge0$,
\[
\sum_{\boldsymbol\Lambda}\prod_{k<K}\Big(\sum_{\mathfrak t\ \text{adm.\ min.\ for}\ \Lambda_0^{(k)}}\omega_k(\mathfrak t)\Big)\Phi(\boldsymbol\Lambda)\le\sum_{(\mathfrak t_k)_{k<K}:\ \boldsymbol\Lambda(\mathfrak t)\ \text{adm.}}\ \prod_{k<K}\omega_k(\mathfrak t_k)\,\Phi(\boldsymbol\Lambda(\mathfrak t)).
\]
\item[(iii)] The resummation \cite[(2.9)]{Ba2} holds with ``$\le$'' in place of ``$=$''.
\end{itemize}
\end{lemma}
\begin{proof}
(i) By \cite[p.~570]{Ba2}, $\Lambda_0^{(k)}$ is determined by $\Lambda_7^{(k-1)\prime}$ and $S(\mathfrak t_k)$. By \cite[(2.8), p.~566]{Ba2}, $\Lambda_7^{(k-1)\prime}$ is determined by $\Lambda_0^{(k-1)}$. Now induct on $k$.

(ii) Expanding the left side gives a sum over pairs $(\boldsymbol\Lambda,\mathfrak t)$. By (i), the map $(\boldsymbol\Lambda,\mathfrak t)\mapsto\mathfrak t$ is injective, and enlarging the index set adds only non-negative terms.

(iii) Fix a configuration whose violation set generates $\Lambda_0$. Then the left side of \cite[(2.9)]{Ba2} equals $1$. The right side counts the minimal admissible subfamilies of the violation set. There is at least one such subfamily, and its summand equals $1$. Two violations within distance $r(\eps)$ of the same large block can give two minimal subfamilies, so equality fails in general.
\end{proof}

\begin{lemma}[elementary geometry]\label{app:G1:LF-geo}
Let $\mathfrak t$ be admissible for $\Lambda_0^{(k)}$, and put $I_k:=\{z:\dist(z,T^{(k)}\setminus\Lambda_7^{(k-1)\prime})\ge b\}$, with $I_0:=T^{(0)}$.
\begin{itemize}
\item[(E1)] Every point of every element of $\mathfrak t$ lies in $X_k$.
\item[(E2)] $\Lambda_i^{(k)c}\subset\mathcal N_{i(r_k+b)}(X_k)$ for $i\le7$, and $X_k\subset\Lambda_5^{(k)c}\subset\Lambda_7^{(k)c}$.
\item[(E3)] Let $z\in X_k\cap I_k$, and let $\beta$ be the large block containing $z$. Then $\beta\subset\Lambda_7^{(k-1)\prime}$ and $\dist(\beta,S(\mathfrak t))\le r_k$.
\item[(E4)] $T^{(k)}\setminus I_k=\mathcal N_{b-1}\big((\Lambda_7^{(k-1)c})'\big)$.
\end{itemize}
\end{lemma}
\begin{proof}
(E1) Let $y$ be a point of an element of $\mathfrak t$, and let $\beta\ni y$ be its large block. Then $\dist(\beta,S(\mathfrak t))=0$, so $\beta$ is not one of the blocks that form $\Lambda_0^{(k)}$. Large blocks partition $T^{(k)}$, hence $y\notin\Lambda_0^{(k)}$.

(E2) By \cite[(2.8)]{Ba2}, a point of $\Lambda_{i+1}^{(k)c}$ lies in a block at distance $\le r_k$ from $\Lambda_i^{(k)c}$. The block has diameter $b-1$, so the point lies within $r_k+b$ of $\Lambda_i^{(k)c}$. Induct on $i$.

(E3) $\beta$ has diameter $b-1<b$ and contains $z\in I_k$, so $\beta\subset\Lambda_7^{(k-1)\prime}$. Since $\beta\not\subset\Lambda_0^{(k)}$, the definition of $\mathfrak L_k$ gives $\dist(\beta,S(\mathfrak t))\le r_k$.

(E4) On unions of $L$-blocks the operation $'$ commutes with complements, so $T^{(k)}\setminus\Lambda_7^{(k-1)\prime}=(\Lambda_7^{(k-1)c})'$.
\end{proof}

\begin{lemma}[cube count]\label{app:G1:LF-cubes}
Put $a_k:=\lceil(r_k+b)/b\rceil$. In each coordinate direction, cut the cycle of large blocks into consecutive arcs of between $a_k$ and $2a_k-1$ blocks; if there are fewer than $a_k$ blocks, take a single arc. Let $\mathcal Q_k$ be the product partition, and put $\mathcal C_k:=\{Q\in\mathcal Q_k:Q\cap X_k\cap I_k\ne\emptyset\}$. Then every box of $\mathcal Q_k$ has sides $\le2r_k+3b$, and for every family $\mathfrak t$ admissible for $\Lambda_0^{(k)}$
\[
\abs{\mathcal C_k}\le3^d\abs{\mathfrak t}.
\]
\end{lemma}
\begin{proof}
The side bound is $(2a_k-1)b\le2r_k+3b$.

Suppose $\dist(z,x)\le r_k+b$. In each coordinate the cyclic distance between $z$ and $x$ is then at most the minimal arc length $a_kb$, so the boxes $Q(z)$ and $Q(x)$ are equal or adjacent. Each box has at most $3^d$ boxes equal or adjacent to it.

Let $Q\in\mathcal C_k$ and $z\in Q\cap X_k\cap I_k$. By (E3) there are $z'\in\beta(z)\subset Q$ and a point $y$ of an element $e'\in\mathfrak t$ with $\dist(z',y)\le r_k$. Charge $e'$ to one of its points $x_{e'}$. Then $\dist(y,x_{e'})\le L-1$, hence $\dist(z',x_{e'})\le r_k+b$, and $Q$ is equal or adjacent to $Q(x_{e'})$.
\end{proof}
The literal family of \cite[(3.43)--(3.45)]{Ba2} uses $\Lambda_7^{(k-1)\prime}$ in place of $I_k$, and it can violate (3.45). Take $\mathfrak t=\emptyset$ and $\partial\Lambda_7^{(k-1)\prime}$ not aligned with the scale-$k$ large blocks: the literal family is non-empty, while $3^d\abs{\mathfrak t}=0$. This is why $I_k$ is used.

\begin{lemma}[cube families and the collar volume]\label{app:G1:LF-D}
Put $A_1:=14+20\vartheta_M$, $A_2:=16+23\vartheta_M$, $S_\varrho:=\sum_{m\ge0}L^{-m}(1+m\log L)^{\varrho}$ ($=7.6553\ldots$), $C_D^0:=A_2(S_\varrho-1)$ and $C_D:=C_D^0+14(1+\vartheta_M)$. At $\vartheta_M=\frac17$, $C_D=144.35\ldots$. For every admissible $\boldsymbol\Lambda$, every $k\le K-1$ and every $i\le7$,
\[
\abs{\Lambda_i^{(k)c}}_k\le(C_Dr_k)^d\textstyle\sum_{j\le k}\abs{\mathcal C_j}.
\]
In particular $\abs{\mathrm{Col}_k}_k$ and $\abs{\mathrm{Sh}_k}_k$ obey the same bound.
\end{lemma}
\begin{proof}
\emph{Construction.} Put $\mathcal E_0:=\mathcal C_0$. For a box $Q$ of $T^{(k-1)}$, let $\Phi_k(Q)$ be the image under $'$ of the union of the $L$-blocks that meet $\mathcal N_{w_{k-1}}(Q)$, where $w_{k-1}:=7(r_{k-1}+b)+Lb$. Put $\mathcal E_k:=\{\Phi_k(Q):Q\in\mathcal E_{k-1}\}\cup\mathcal C_k$. Then $\abs{\mathcal E_k}\le\sum_{j\le k}\abs{\mathcal C_j}$.

\emph{Sides.} Let $s_j$ be the side of a box at the scale $j$ where it is born. Then $s_j\le(2+3\vartheta_M)r_j$. One step gives $s_{i+1}\le(s_i+2w_i+2L-2)/L$, and $2w_i+2L-2\le14r_i+20b\le A_1r_i$. Since $\ell_i=\ell_k+(k-i)\log L\le(1+(k-i)\log L)\ell_k$, iteration gives
\[
s_k\le A_2\textstyle\sum_{i=j}^{k-1}L^{-(k-i)}r_i\le A_2r_k\sum_{m\ge1}L^{-m}(1+m\log L)^{\varrho}=C_D^0r_k .
\]

\emph{Covering.} We show $X_k\subset\bigcup\mathcal E_k$ by induction on $k$. Points of $X_k\cap I_k$ lie in $\bigcup\mathcal C_k$. Let $z\in X_k\setminus I_k$. By (E4) there is $z_1\in(\Lambda_7^{(k-1)c})'$ with $\dist(z,z_1)\le b-1$. By (E2) at scale $k-1$ and the induction hypothesis,
\[
B(z_1)\subset\mathcal N_{7(r_{k-1}+b)}(X_{k-1})\subset\textstyle\bigcup_{Q\in\mathcal E_{k-1}}\mathcal N_{7(r_{k-1}+b)}(Q).
\]
At scale $k-1$, $B(z)$ lies within $Lb-1$ of $B(z_1)$. Hence $z\in\Phi_k(Q)$ for some $Q\in\mathcal E_{k-1}$.

\emph{Count.} By (E2), $\Lambda_i^{(k)c}\subset\bigcup_{Q\in\mathcal E_k}\mathcal N_{7(r_k+b)}(Q)$. Each of these neighbourhoods has at most $(C_D^0r_k+14(r_k+b))^d\le(C_Dr_k)^d$ points. Finally, $\mathrm{Col}_k$ and $\mathrm{Sh}_k$ are contained in $\Lambda_7^{(k)c}$.
\end{proof}

\begin{proposition}[vector form on the kept variables]\label{app:G1:LF-V}
Let $\Psi=(A_0|_{S_0},A_k|_{S_k})_{1\le k\le K}$ with $S_0:=\Lambda_5^{(0)c}$ and $S_k:=\mathrm{Sh}_k$. Let $Q_k$ be the $k$-fold block average, $a_k\in[a_\flat,a]$ with $a_\flat=a(1-L^{-2})$, and
\[
\ip\Psi{\Delta_v\Psi}:=\min_{A_0|_{\Lambda_5^{(0)}}}\Big[\ip{A_0}{(-\Delta+\mu_0^2)A_0}+\sum_{k\ge1}a_k t_k^{d-2}\sum_{x\in S_k}\abs{A_k(x)-(Q_kA_0)(x)}^2\Big].
\]
If $\mu_0\eps_0\le1$, then
\[
\ip\Psi{\Delta_v\Psi}\ge\gamma_0\sum_{k=0}^K\Big[\sum_{\langle x,x'\rangle\subset S_k}t_k^{d-2}\abs{A_k(x')-A_k(x)}^2+\sum_{x\in S_k}t_k^d\mu_0^2\abs{A_k(x)}^2\Big],
\]
with $\gamma_0^{-1}:=\max\{1+6dC_P,(6d+2)/a_\flat,3\}$. Here $C_P\le\frac34$ is the two-block Poincar\'e constant: $n^{d-2}\abs{\mathrm{avg}_{B_1}f-\mathrm{avg}_{B_2}f}^2\le C_P\sum_{b'\subset B_1\cup B_2}\abs{\nabla_{b'}f}^2$ for adjacent blocks $B_1,B_2$ of side $n$. At $d=3$, $a=1$ this gives $\gamma_0=3/80$.
\end{proposition}
\begin{proof}
Fix $A_0$ and write $\bar A_k:=Q_kA_0$.
(a) For $k=0$, the bond terms and the mass terms on $S_0$ are terms of $\ip{A_0}{(-\Delta+\mu_0^2)A_0}$.
(b) For $k\ge1$ and a bond $\langle x,x'\rangle\subset S_k$,
\[
\abs{A_k(x')-A_k(x)}^2\le3\abs{A_k(x')-\bar A_k(x')}^2+3\abs{A_k(x)-\bar A_k(x)}^2+3\abs{\bar A_k(x')-\bar A_k(x)}^2.
\]
Apply the Poincar\'e inequality with $n=L^k$ and multiply by $\eps^{d-2}$. The last term then contributes at most $C_P\sum\eps^{d-2}\abs{\nabla A_0}^2$ over the double block. Each fine bond lies in at most $2d$ double blocks per level. The fine regions $B^k(S_k)$ are disjoint in $k$, because $S_{k'}\subset\Lambda_5^{(k)}$ for $k'>k$. So the Poincar\'e terms sum to at most $6dC_P\ip{A_0}{-\Delta A_0}$. Each $A_k(x)$ lies in at most $2d$ shell bonds, so the self terms sum to at most $(6d/a_k)\,a_kt_k^{d-2}\abs{A_k-\bar A_k}^2$.
(c) By Jensen's inequality, $t_k^d\mu_0^2\abs{A_k}^2\le2\mu_0^2t_k^2\cdot t_k^{d-2}\abs{A_k-\bar A_k}^2+2\mu_0^2\sum_{z\in B^k(x)}\eps^d\abs{A_0(z)}^2$, and $2\mu_0^2t_k^2\le2$.
Adding (a)--(c), the right side with $\gamma_0=1$ is at most $\gamma_0^{-1}$ times the bracket. Finally minimise over $A_0|_{\Lambda_5^{(0)}}$.
\end{proof}

\begin{corollary}[weighted vector gains]\label{app:G1:LF-Vw}
Let $S^1$ be the quadratic form of the photon chain of Lemma~\ref{app:G1:S05-T} with all weights equal to $1$. By Lemmas~\ref{app:K:L14-1} and~\ref{app:K:L14-2}, $\Delta_v=\operatorname{Sch}_\Psi S^1$. Let $HA_{\rm del}$ be one of the photon forms $HA^{(k)}_{\rm del}$, $0\le k\le K$, of Lemma~\ref{app:G1:S05-PH}(d), and put $\Delta^w_v:=\operatorname{Sch}_\Psi HA_{\rm del}$. If $HA_{\rm del}\ge(1-\hat\xi_{\rm tot})w_{\rm lo}S^1$, then $\Delta_v^w\ge(1-\hat\xi_{\rm tot})w_{\rm lo}\Delta_v$.
\end{corollary}
\begin{proof}
$\ip\Psi{\operatorname{Sch}_\Psi(H)\Psi}=\min_\chi\ip{(\Psi,\chi)}{H(\Psi,\chi)}$ is monotone and positively homogeneous in $H$.
\end{proof}

\begin{proposition}[region sum]\label{app:G1:LF-sigma}
For a family $\mathfrak t$ at scale $k$, put
\[
W_k^w(\mathfrak t):=\exp\Big(-p_k^2\big[\vartheta(\tfrac a8\abs{P_v}+\tfrac{\gamma_0}4\abs{Q_v}+\tfrac{\gamma_0}4\abs{R_v})+\tfrac a8\abs{P_s}+\tfrac{\gamma_0}4\abs{Q_s}+\abs{R_s}\big]\Big),
\qquad
\zeta''_w(\Lambda_0^{(k)}):=\sum_{\mathfrak t\ \text{adm.\ min.}}W_k^w(\mathfrak t),
\]
with $0<\vartheta\le1$ and $\gamma_0\le\frac13$. Put $c_0'^{\,w}:=\vartheta\min\{\frac a8,\frac{\gamma_0}4\}$ and $c_L:=\sup_kc_C^{(k)}t_k^{\kappa_0}+c_B^w+C_{\rm old}\eps_0$, and assume
\begin{equation}\label{app:G1:eq-SigmaW}
\tfrac12c_0'^{\,w}p(t)^2\ge3^dc_LC_D^dR^d\ell(t)^{1+\varrho d}/\log2\qquad(t\le\eps_0).
\end{equation}
At $d=3$, $\varrho=2$ and $p(t)=b_0\ell(t)^6$, condition \eqref{app:G1:eq-SigmaW} reads $\hat\ell^{5}\ge54c_LC_D^3R^3/(c_0'^{\,w}b_0^2\log2)$. Then, uniformly in $\eps=2^{-K}\eps_0$, $\ell_T$ and $K$,
\[
\sum_{\boldsymbol\Lambda}\prod_{k<K}\zeta''_w(\Lambda_0^{(k)})\exp\Big(\sum_{k=0}^K\big[c_C^{(k)}t_k^{\kappa_0}\abs{\mathrm{Col}_k}_k+c_B^w\abs{\mathrm{Sh}_k}_k\big]+\sum_{k<K}C_{\rm old}t_k\abs{\Lambda_7^{(k)c}}_k\Big)\le\mathrm{e}^{C_\Sigma\abs T},
\]
\[
C_\Sigma:=c_L\eps_0^{-3}+n_d\sum_{m\ge1}(2^{-m}\eps_0)^{-3}\exp\big(-\tfrac12c_0'^{\,w}b_0^2(\hat\ell+m\log2)^{12}\big).
\]
\end{proposition}
\begin{proof}
Since $\vartheta\le1$ and $\min\{\frac a8,\frac{\gamma_0}4\}\le1$, each of the six coefficients in $W_k^w$ is at least $c_0'^{\,w}$, so $W_k^w(\mathfrak t)\le\mathrm{e}^{-c_0'^{\,w}p_k^2\abs{\mathfrak t}}$. The terms with $k=K$ are at most $c_L\abs{T^{(K)}}=c_L\eps_0^{-3}\abs T$. For $k<K$, the sets $\mathrm{Col}_k$, $\mathrm{Sh}_k$ and $\Lambda_7^{(k)c}$ all lie in $\Lambda_7^{(k)c}$, and $C_{\rm old}t_k\le C_{\rm old}\eps_0$.

Fix a tuple $\mathfrak t$ that generates $\boldsymbol\Lambda(\mathfrak t)$. The radii $r_k$ decrease in $k$, and $K-j=\log_2(\eps_0/t_j)\le\ell_j/\log2$. So Lemmas~\ref{app:G1:LF-D} and~\ref{app:G1:LF-cubes} give
\[
\sum_{k<K}c_L\abs{\Lambda_7^{(k)c}}_k\le3^dc_LC_D^d\sum_j\abs{\mathfrak t_j}\sum_{k=j}^{K-1}r_k^d\le3^dc_LC_D^dR^d\sum_j\abs{\mathfrak t_j}\ell_j^{1+\varrho d}/\log2 .
\]
By Lemma~\ref{app:G1:LF-fam}(ii) and \eqref{app:G1:eq-SigmaW}, the left side is at most $\mathrm{e}^{c_L\eps_0^{-3}\abs T}\prod_{j<K}\sum_{\mathfrak t_j}\mathrm{e}^{-\frac12c_0'^{\,w}p_j^2\abs{\mathfrak t_j}}$. Charge each element to one point (Lemma~\ref{app:G1:LF-geo}(E1)). Each point carries at most $n_d$ candidate elements, so $\sum_{\mathfrak t_j}\mathrm{e}^{-x\abs{\mathfrak t_j}}\le(1+\mathrm{e}^{-x})^{n_dN_j}\le\exp(n_dN_j\mathrm{e}^{-x})$. Finally insert $N_j=\abs T(2^{-m}\eps_0)^{-3}$ and $\ell_j=\hat\ell+m\log2$, with $m=K-j\ge1$.
\end{proof}

\begin{proof}[Proof of Lemma~\ref{app:G1:S05-LF}]
(a) The partition of unity \cite[(2.4)--(2.6)]{Ba2} is an identity for every choice of thresholds. The integrand is $\ge0$, so dropping restrictions gives an upper bound. Then apply Lemma~\ref{app:G1:LF-fam}.

(b) Use \cite[Prop.~2.1, (2.57), p.~570]{Ba2} relative to \cite[Prop.~1]{Ba3}, as in Lemma~\ref{app:G1:S05-SC}(C1). The factor $\sigma_+$ comes from the replacements of Lemma~\ref{app:G1:S05-PH}(d). The factor $2^{N_\ell}$ comes from the $D$-coloured lines, with $c_k\le\frac13$. At $k=0$, the bound $\lambda\eps x^2+\frac12\delta\eps^2x\ge-\delta^2\eps^3/(16\lambda)$ with $\eps\abs\delta\le C_\delta(M)$ gives the per-site cost. Since $\eps\le\eps^{\kappa_0}$, this is at most $c_C^{(0)}t_0^{\kappa_0}$. Under the convention that the vacuum counterterms are global $c$-numbers \cite[(1.30)]{Ba3}, no vacuum constant is localised on a collar (Lemma~\ref{app:G1:S05-RF}(V)). The outside pieces are bounded by Lemma~\ref{app:G1:S05-RF}(V), Proposition~\ref{app:G1:R9vac}(ii) and Lemma~\ref{app:G1:IF1abs}.

(c) Use the first term of \cite[(2.57)]{Ba2}, as in \cite[(2.13)]{Ba2}.

(d) \emph{Scalar entries.} These come from \cite[(3.22)--(3.39)]{Ba2} at fixed external $\tilde A$, and they do not involve the weights. For $\tilde A\ne0$, \cite[Prop.~3.1]{Ba2} is supplied by \cite[Prop.~3.1$'$, (1.22), p.~574]{Ba4}. The hypothesis ``$m^2\le O(1)$'' is read as (E.m2) (Proposition~\ref{app:G1:I7}).

\emph{Vector entries.} Lemma~\ref{app:G1:S05-PH}(d) gives $HA_{\rm del}\ge(1-\hat\xi_{\rm tot})HA$, and $HA\ge w_{\rm lo}S^1$ by Lemma~\ref{app:G1:S05-T}. Together they give the hypothesis of Corollary~\ref{app:G1:LF-Vw}. Combine that corollary with Proposition~\ref{app:G1:LF-V} (where $\mu_0\eps_0\le1$, that is $\mu_*\eps_0\le1$, which holds under (E.dom)) and with the split $\frac14+\frac14$ of \cite[(3.30), (3.40)]{Ba2}, in which the $P_v$ term and the bound of the type of \cite[Prop.~3.1]{Ba2} share one quarter. A violated gradient restriction on $Q_v$ gives $\frac14(1-\hat\xi_{\rm tot})w_{\rm lo}\gamma_0\cdot\frac12p_k^2=\vartheta\frac{\gamma_0}4p_k^2$ per element. The same holds on $R_v$, since $\mu_0^2t_k^2\abs A^2>p_k^2$ there. The $P_v$ term \cite[(3.33)]{Ba2} gives $\frac1{16}(1-\hat\xi_{\rm tot})w_{\rm lo}aL^{d-2}p_k^2\ge\vartheta\frac a8p_k^2$.

\emph{Arithmetic.} We have $\hat\xi_{\rm tot}\le\frac1{12}$ and $\bar\delta=\kappa\mu_*^2\eps_0^2\le\frac1{12}$, the latter by (E.$\delta$). At $(1,m_2)$, $w_{\rm lo}\ge(1+\bar\delta)^{-1}\ge\frac{12}{13}$; at $(-\eta,m_1)$, $w_{\rm lo}=1$. Hence $\vartheta_+\ge\frac12\cdot\frac{11}{12}\cdot\frac{12}{13}=\frac{11}{26}$ and $\vartheta_-\ge\frac{11}{24}$.

\emph{Kept variables.} The Gaussian cost of each kept vector variable is at most $3^{1/2}$. This gives the term $\frac d2\log3\cdot c_{\rm sh}$ in $c_B^w$.

(e) Proposition~\ref{app:G1:LF-sigma}, or Lemma~\ref{app:G1:I145-comb}.

\emph{Uniformity.} None of the constants depends on $\eps$, $k$, $\boldsymbol\Lambda$ or $\ell_T\ge\ell_0$. By Statement~\ref{app:K:SL}, $c_C^{u=0}$ and $c_B$ are bounded uniformly in $\eps_0\in(0,\eps_1]$. Hence so are the constants $c_5=c_B^w$ and $c_7=\max\{\sigma_+2^{N_\ell}c_C^{u=0},C_\delta(M)^2/(16\lambda)\}$ of condition (E.ent), and so is $c_L$. The thresholds on $\eps_0$ are imposed at the binding point, so one $\eps_0$ serves both points.
\end{proof}

\emph{The region sum} is Lemmas~\ref{app:G1:I145-geo} and~\ref{app:G1:I145-comb}. The constants $c_A,c_*,C_V$ are evaluated at $\eps_1$, and (E.comb) is then one more downward-closed upper bound on $\eps_0$; in it $c_A$ is replaced by $c_A+C_{\rm old}$ (Lemma~\ref{app:G1:IF1abs}).

\begin{remark}[what $c_A$ dominates; other routes]\label{app:G1:S05-altroutes}
(i) In the closing step, (H-LF) holds with $c_*(\pi)=c_0'^{\,w}$ and $C_0(\pi)=c_{\rm fin}\eps_0^{-3}$; the latter collects the final-scale kept variables at $O(1)$ each. The constant $c_A(\pi)$ dominates every per-site cost on a subset of $\Lambda_7^{(k)c}$: the collar costs of Lemma~\ref{app:G1:S05-LF}(b), the vacuum collar shares, the shells, the forward replacement errors near large fields ($\times w_{\rm hi}$), and Ba{\l}aban's $(1+\log)\abs{\Lambda_0^{(k)c}}$. It also covers the kept vector variables, including the coarse variables over earlier large-field sets. All of these are finite at both points.
(ii) The region sum can also be bounded through the cube count of \cite[(3.45)]{Ba2}, as in Proposition~\ref{app:G1:I6}; such routes need the repair of Remark~\ref{app:G1:I6-345}, because the printed count fails for large blocks straddling $\partial\Lambda_7^{(k-1)\prime}$. The route above uses neither (3.45) nor cubes.
\end{remark}

\paragraph{(S4) The renormalisation flow.}

\begin{lemma}[the vacuum output of a step]\label{app:G1:S05-TS}
For any exact Gaussian staging, in particular King's photon chain (Lemma~\ref{app:G1:S05-T}) and Ba{\l}aban's scalar staging, the following holds. The field-independent output of order $\le\nbar$ at step $k$ of the $T_{\nbar}$-truncated flow is the sum of the terms of the bare vacuum graphs of Definition~\ref{app:G1:vacgrade} (counterterm vertices included, each coefficient $\delta m^2$ kept as a number) all of whose lines have index $\le k$ and at least one of which has index $k$; that is, it equals $\cV_k$. Consequently $E_1(u)=\sum_{k<K}\cV_k(u)+E_{\rm fin}(u)$, where $\cV_k(u):=E_{K,k}(u)-E_{K,k+1}(u)$, the pieces $E_{K,k}(u)$ are those of Definition~\ref{app:G1:vacgrade} built with the lines of the slice model at $u$, and $E_{\rm fin}(u)$ is as in \eqref{app:G1:eq-Efinu}; and $\cV_k(u)$ equals the step-$k$ output computed with the free index-$k$ covariances.
\end{lemma}
\begin{proof}
First let the covariances of the staging not depend on the fields, as for King's photon chain. Let $W_0$ be minus the interaction, including $\delta m^2\abs\varphi^2$, which has order $\ge1$, and let $W_{k+1}=\log\E_{\mu_k}\mathrm{e}^{W_k}$. Since $\log\E\mathrm{e}^W=\sum_m\kappa_m(W)/m!$ with the $m$-th cumulant $m$-linear and of order $\ge m$, the part of order $\le\nbar$ of $\log\E\mathrm{e}^{W}$ depends only on $T_{\nbar}W$; so the truncated flow agrees with the exact one up to order $\nbar$, and constants pass through later steps unchanged, $\log\E\mathrm{e}^{c+U}=c+\log\E\mathrm{e}^U$. In the exact flow, the field-independent part of $W_{k+1}$ is the sum of the connected vacuum graphs of the bare interaction whose lines carry $\sum_{j\le k}C_{(j),f}$. Subtracting the same sum with $\sum_{j<k}C_{(j),f}$ leaves, by multilinearity, the terms with all line indices $\le k$ and at least one equal to $k$, that is, $\cV_k$. The identity for $E_1(u)$ follows by the multilinear sorting of the proof of Lemma~\ref{app:G1:vacgrade-id}(i), applied with the lines of the slice model at $u$. For Ba{\l}aban's scalar staging, whose step covariances depend on the gauge background, the same conclusion is Statement~\ref{app:K:L04-6}.
\end{proof}

\emph{The localisation sum} is Lemma~\ref{app:G1:VS}. Its anchor is the first vertex lying in $W$. When the cut kernel supplies the tree link, the rate is $\delta_1$ in place of $\delta$.

\begin{lemma}[mass and vacuum flow along the step]\label{app:G1:S05-RF}
Let $(u,M)\in\{(1,m_2),(-\eta,m_1)\}$, $\eps=2^{-K}\eps_0$ and $0\le k\le K$, and put
\[
P_K(t):=2(e^2+\lambda)t+(e^4+e^2\lambda+\lambda^2)t^2\big(1+\log_2(\eps_0/t)\big),\qquad P_m(t):=(e^2+\lambda)t+(e^4+e^2\lambda+\lambda^2)t^2 .
\]
Put $\mathcal L_k:=B^k(\Lambda_7^{(k)}\cap\Lambda_8^{(k)c})$ (the layer) and $\mathcal D_k:=B^k(\Lambda_8^{(k)})$, so that $W_k=\mathcal L_k\sqcup\mathcal D_k$. At both points and at every step the following hold.
\begin{itemize}
\item[(M)] $a_{K,T}(M)=M^2+\delta m^2_{K,T}(M)$. Here $\delta m^2_{K,T}(M)$ is, coefficient by coefficient in $e^\alpha\lambda^\beta$ with $2\le\alpha+2\beta\le4$, the counterterm of \cite[(1.19)--(1.23)]{Ba3} on $T_\eps$. It has $C$-lines only, so it is the same number at both values of $u$ for the same $M$. For every $k$ it splits as $\delta m^2_{\rm fin}+\delta m^2_{K,k}+\delta m^2_k$ \cite[(1.26)--(1.29)]{Ba3}. In this split $\delta m^2_{\rm fin}$ and $\delta m^2_{K,k}$ are $c$-numbers, invariant under the translations of $t_K\Z^3$ and of $t_k\Z^3$ respectively, and $\delta m^2_k$ is a localised graph sum. There is no $u$-dependent mass vertex. We have $\abs{\delta m^2_{K,T}(M)}\le C_\delta(M)\eps^{-1}$, and per site
\[
\lambda(\eps)\abs\varphi^4+\tfrac12\delta m^2_{K,T}(M)\eps^2\abs\varphi^2\ge-C_\delta(M)^2\eps/(16\lambda).
\]
Moreover $\sup_x\abs{\delta m^2_{K,k}(x)}t_k^2\le C_mP_K(t_k)$, and the step differences obey Lemma~\ref{app:K:L14-3}(b). At $k\ge1$ the mass pieces of step $k$ leave $\mathrm{Col}_k$ together with the whole interaction, at a cost contained in $O(t_k^{\kappa_0})\abs{\mathrm{Col}_k}_k$; at $k\ge1$ no mass term is absorbed site by site. A connected coloured family with at least one $D$-line and an external leg is convergent and receives no counterterm. Per block, it is at most $C_{U1}e^2\abs u\kappa\mu_*^2L^2t_{k+1}^3\sup\norm{h\Phi_{\rm ext}}_{1,\alpha_0}$.
\item[(V)] $E_0^w$ is cancelled exactly for each $\boldsymbol\Lambda$ (Lemma~\ref{app:G1:S05-N0}). Further, $E_1(u)=E_{\rm fin}(u)+E_K(u)$ with $E_K(u)=E_{K,k}(u)+\sum_{j<k}\cV_j(u)$ and $\cV_j(u):=E_{K,j}(u)-E_{K,j+1}(u)$. The terms $E_{\rm fin}(u)$ and $E_{K,k}(u)$ stand as global factors $\mathrm{e}^{-E_{\rm fin}(u)-E_{K,k}(u)}$ in front of the inductive bound. At the step $k\to k+1$ only $\cV_k(u)$ is localised. This is done at scale $k$, by the rule of \cite[p.~420]{Ba3} applied to whole renormalised classes with one external field each. The anchor of a class is its first leg, which is the same in every member. A localisation $\Lambda$ is \emph{good} if $\Lambda\subset W_k$, \emph{bad} if it meets both $W_k$ and $W_k^c$, and \emph{outside} if $\Lambda\cap W_k=\emptyset$. Let $A_k$, $\bar N$, $\bar S$, $\delta$ be as in Proposition~\ref{app:G1:R9vac}. Then:
\begin{itemize}
\item[(V1)] The sum of $\cV_k(\Lambda)$ over the good $\Lambda$ equals the field-independent output of order $\le\nbar$ of step $k$, computed with the free index-$k$ covariances and localised by the same rule. Hence the field-independent part of $\mathcal P^{(k)}_w$ is a remainder $\mathcal R^{\rm vac}_k$, and $\mathcal R^{\rm vac}_{k+1}-\mathcal R^{\rm vac}_k$ consists of the bad localisations of $\cV_k$ and of domain differences. Moreover
\[
\abs{\mathcal R^{\rm vac}_{k+1}-\mathcal R^{\rm vac}_k}\le\bar N\bar SA_k\big[\abs{\mathcal L_k}_k+\mathrm{e}^{3\delta}\mathrm{e}^{-\delta r(t_k)/2}\abs{\mathcal D_k}_k\big]+C_\kappa t_k^{\kappa_B}N_k,\qquad\kappa_B=\tfrac{16}5 .
\]
\item[(V2)] $\abs{\sum_{\Lambda\ \rm outside}\cV_k(\Lambda)}\le\bar SA_k\abs{\Lambda_7^{(k)c}}_k\le C_{\rm old}t_k\abs{\Lambda_7^{(k)c}}_k$. This $c$-number is not part of $\mathcal P_w$. It is booked as $\mathfrak O_{k+1}:=\mathfrak O_k+\bar SA_k\abs{\Lambda_7^{(k)c}}_k$, with $\mathfrak O_0:=0$.
\item[(V3)] No vacuum $c$-number is removed on a collar, no term of $E_K(u)$ is moved twice, and $E_{\rm fin}(u)$ is not part of $\mathcal P^{(k)}_w$. The layer terms of (V1) are at most $C(\kappa_0)t_k^{\kappa_0}\abs{\mathrm{Col}_k}_k$. Summed over $k$, the deep terms and the domain terms are at most $B\abs T$, with $B$ independent of $\eps$, $K$ and $\ell_T$.
\end{itemize}
\item[(C11)] Let $F_{11}(u)=\frac u2\sum_{i<K}\delta_i\ip{H_i}{\Pi_2}$ be the order-$e^2$ vacuum family with one $D$-line (the bubble and the seagull closed by one $D$-line). Split $F_{11}=\sum_{k<K}F^{[k]}_{11}+F^{\rm fin}_{11}$ by the top index of the lines, where $F^{\rm fin}_{11}$ collects the terms with a regular scalar line. Then $F^{[k]}_{11}$ is the order-$e^2$, one-$D$-line part of $\cV_k(u)$, and $F^{\rm fin}_{11}$ is part of $E_{\rm fin}(u)$. Moreover:
\begin{itemize}
\item[(a)] for every localisation $\Lambda\subset W_k$, the generated free vacuum output and the moved piece agree: $\mathrm{Gen}^{\rm free}_k(\Lambda)=F^{[k]}_{11}(\Lambda)$;
\item[(b)] the localisations that straddle $\partial W_k$ satisfy
\[
\Big|\sum_{\Lambda\ \rm bad}F^{[k]}_{11}(\Lambda)\Big|\le3S_3C_{\rm pl}e^2\abs u\kappa\mu_*^2L^3t_k^3\Big[\abs{\mathcal L_k}_k+\mathrm{e}^{\delta_1(c_d+1)/2}\mathrm{e}^{-\delta_1r(t_k)/2}\abs{\mathcal D_k}_k\Big],
\]
where $S_3:=\sum_{D\ge0}(c''_d(D+2)^d)^3\mathrm{e}^{-\delta_1D/2}$ and $C_{\rm pl}$ is the constant of the order-$e^2$ anchor bounds for this family (Theorem~\ref{app:G1:B5e2-thm}). The layer part is at most $C(\kappa_0)t_k^{\kappa_0}\abs{\mathrm{Col}_k}_k$. The deep part, summed over $k$, is at most $C\abs T$, uniformly in $K$;
\item[(c)] the outside localisations of the $u$-part of $\cV_k$, in particular those of $F^{[k]}_{11}$, are bounded per block of $\Lambda_7^{(k)c}$ by $\bar S\bar N_{\rm ph}\abs u\bar c(1+\bar c)^{\bar N_{\rm ph}-1}C_PP(t_k)+2C_{\rm loc}t_k$. Here the notation is that of Proposition~\ref{app:G1:R9vac}(iii), and $C_{\rm loc}$ is as in Proposition~\ref{app:G1:I1}. The term $2C_{\rm loc}t_k$ is the share of the order-$e^2$ family with one $D$-line (Proposition~\ref{app:G1:I1}(iii), (iv) and Lemma~\ref{app:G1:I1c-e2}). They are booked in $\mathfrak O_{k+1}$ and carry no factor $1+K-k$.
\end{itemize}
The global sum $F_{11}(u)$ is of size $O(K)\abs T$ and is never bounded. It occurs only inside $E_1(u)$, both in front of the inductive bound and in $T_{\nbar}[\Phi_M(u)]$, and as the sum of the generated constants; these two occurrences are equal.
\end{itemize}
The constants $C_\delta(M)$, $C_m$, $C_{U1}$, $C_\kappa$, $C_{\rm old}$, $C_{\rm pl}$ and $C(\kappa_0)$ do not depend on $\eps$, $k$, $x$, $\boldsymbol\Lambda$ or $u$. They depend on $\ell_T$ only through terms that do not increase with $\ell_T$, evaluated at $\ell_0$, and they are finite at $M=m_1$.
\end{lemma}
\noindent The cancellation of the good localisations in (C11) is part of Proposition~\ref{app:G1:I1}, which is proved under the convention that the vacuum counterterms are global $c$-numbers \cite[(1.30)]{Ba3}.

\begin{lemma}[geometry of the layer]\label{app:G1:layer}
Every block of $\mathcal L_k$ lies within distance $r(t_k)+b$ of a block of $\mathrm{Col}_k$, and
\[
\abs{\mathcal L_k}_k\le\big(2r(t_k)+2b+1\big)^d\,\abs{\mathrm{Col}_k}_k .
\]
\end{lemma}
\begin{proof}
Let $\beta$ be a large block of $\Lambda_7^{(k)}\cap\Lambda_8^{(k)c}$. By \cite[(2.8)]{Ba2} with $i=7$, $\beta$ lies within $r(t_k)$ of a point of $\Lambda_7^{(k)c}$. Let $\square$ be the unit block containing that point. If $\square\subset\Lambda_0^{(k)c}$, then $\beta$ lies within $r(t_k)$ of $\Lambda_0^{(k)c}$. By (2.8) with $i=0$ this gives $\beta\subset\Lambda_1^{(k)c}\subset\Lambda_7^{(k)c}$, which contradicts $\beta\subset\Lambda_7^{(k)}$. Therefore $\square\subset\Lambda_0^{(k)}\cap\Lambda_7^{(k)c}\subset\Lambda_7^{(k-1)\prime}\cap\Lambda_7^{(k)c}=\mathrm{Col}_k$. Since $\beta$ has diameter $b-1$, every unit block of $\beta$ lies within $r(t_k)+b$ of $\square$. Now count the unit blocks within $r(t_k)+b$ of each such $\square$.
\end{proof}

\begin{proof}[Proof of Lemma~\ref{app:G1:S05-RF}]
\emph{(M).} The identity $a_{K,T}(M)=M^2+\delta m^2_{K,T}(M)$ is Proposition~\ref{prop:UV-B1e}. The coefficients are defined on $T_\eps$ by $-\delta m^2_{(\alpha,\beta)}+\sum_{x\in T_\eps}\eps^d\Sigma^\eps_{(\alpha,\beta)}(x)=0$ \cite[p.~417]{Ba3}. The self-energy graphs entering them carry $C$-lines only, so $u$ does not enter.

The split is the definition \cite[(1.25)--(1.28), p.~418]{Ba3}. On $T_\eps$ the propagators $G^\eps_K$ and $G^\eps_K(0)$ commute with the translations of $t_K\Z^3$, and $G^\eps_{K,k}$ commutes with those of $t_k\Z^3$. The coefficients inherit this periodicity. The localisation of $\delta m^2_k$ is that of \cite[p.~420]{Ba3}.

The bound $\abs{\delta m^2_{K,T}(M)}\le C_\delta(M)\eps^{-1}$ is the size statement $\delta m^2=O(\eps^{-1})$ of \cite[p.~559]{Ba2}. Its explicit constant $C_\delta(M)$ is that of Proposition~\ref{app:K:L14-4}, with the majorant constants of its hypothesis (Bd) evaluated at $\ell_T=\ell_0$; they do not increase with $\ell_T$. Put $s=\abs\varphi^2\ge0$ and $\delta=\delta m^2_{K,T}(M)$. Then $\lambda\eps s^2+\frac12\delta\eps^2s\ge-\delta^2\eps^3/(16\lambda)\ge-C_\delta(M)^2\eps/(16\lambda)$.

The pointwise bound on $\delta m^2_{K,k}$ and the bounds on the step differences are Lemma~\ref{app:K:L14-3}.

The removal from $\mathrm{Col}_k$ is \cite[Prop.~2.1, (2.57), p.~570]{Ba2}. Under the restrictions (2.55), the interaction is removed from $B^k(\mathrm{Col}_k)$, leaving only $-\lambda(t_k)\sum\eta^d\abs{\phi^{(k)}}^4$, at the cost $O(t_k^{\kappa_0})\abs{\mathrm{Col}_k}$. With $\kappa_0=\frac15$ this is Lemma~\ref{app:G1:kappa0}.

The statement on coloured families is Lemma~\ref{app:K:L14-5}(b),(c) at scale $n=k+1$. Its part (c) also bounds the share anchored in a scale-$k$ sub-block. The transfer of the line bounds to the $D$-parts is Lemma~\ref{app:K:L14-5}(d), whose hypothesis is Lemma~\ref{app:K:L14-6}.

\emph{(V1), the identity.} By \cite[p.~570]{Ba2}, the whole interaction is removed from $\mathrm{Col}_k$ before step $k$. So every vertex of a family generated at step $k$ lies in $W_k$, and the generated output has only localisations $\Lambda\subset W_k$. By Lemma~\ref{app:G1:S05-TS}, that output, computed with the free index-$k$ covariances, is $\cV_k(u)$ as a truncated power series. In fact this is an identity of the same index-assigned Wick graphs: a contraction with at least one index-$k$ line and all other lines of index $<k$ is exactly an assignment of top index $k$. A generated graph and the corresponding counterterm graph have the same vertex set. The anchor is the first leg of the whole class, the same in every member, so the rule of \cite[p.~420]{Ba3} gives both graphs the same $\Lambda$. The identity therefore holds localisation by localisation for $\Lambda\subset W_k$. The bad and outside localisations of $\cV_k$ have no generated partner.

\emph{(V1), the bounds.} By \cite[(2.8)]{Ba2}, $W_k$ lies $2r(t_k)$ inside $\Lambda_5^{(k)}$. Replace the restricted lines of step $k$ by the free ones, and $g_{k+1}$ by $1$. This costs $O((L^j\eps)^\kappa)\abs{B^{k-j}(\Lambda_2^{(k)})}$ for every $\kappa>d$ \cite[p.~578]{Ba2}, since $g_{k+1}=1$ at distance $\ge M_b$ from $\partial\Omega_1$ \cite[p.~412]{Ba3}. Take $\kappa=\kappa_B$. The bad localisations are bounded by Proposition~\ref{app:G1:R9vac}(i). Every vacuum graph has $o=\alpha/2+\beta\ge1$. So, for $\abs u\le1$, $A_k\le C_P(1+\bar c)^{\bar N_{\rm ph}}C_et_k$ with the constants of Proposition~\ref{app:G1:R9vac}, and $\bar SA_k\le C_{\rm old}t_k$.

\emph{(V2)} is Proposition~\ref{app:G1:R9vac}(ii) with $X=W_k^c=\Lambda_7^{(k)c}$, counted in scale-$k$ blocks.

\emph{(V3).} By Lemma~\ref{app:G1:layer}, and since $r(t)=R(1+\log t^{-1})^2$ and $\kappa_0<1$,
\[
A_k\abs{\mathcal L_k}_k\le Ct_k\big(2r(t_k)+2b+1\big)^d\abs{\mathrm{Col}_k}_k\le C(\kappa_0)t_k^{\kappa_0}\abs{\mathrm{Col}_k}_k .
\]
Next, $A_k\mathrm{e}^{-\delta r(t_k)/2}N_k\le C\abs Tt_k^{-2}\mathrm{e}^{-\delta r(t_k)/2}$, and $\sum_kt_k^{-2}\mathrm{e}^{-\delta r(t_k)/2}$ converges uniformly in $K$. Finally $\sum_kt_k^{\kappa_B}N_k=\abs T\sum_kt_k^{1/5}\le(1-2^{-1/5})^{-1}\eps_0^{1/5}\abs T$. The outside costs are paid in the region sum by Lemma~\ref{app:G1:IF1abs}. Each term of $E_K(u)=E_{K,k}(u)+\sum_{j<k}\cV_j(u)$ leaves the global factor exactly once, at its top index, and $E_{\rm fin}(u)$ is never moved.

\emph{(C11).} The split is multilinearity in $C_0=\sum_{j<K}S_j+R_s$ and in $D=\sum_{i<K}\delta_iH_i$; $D$ has no regular remainder. The set of index assignments with all indices $\le k$ and at least one equal to $k$ is the set moved at $k\to k+1$ \cite[p.~422]{Ba3}. Hence $F^{[k]}_{11}$ is the order-$e^2$, one-$D$-line part of $\cV_k(u)$.

(a) This is the identity of (V1) restricted to this family, since the generated family and $F^{[k]}_{11}$ consist of the same index-assigned graphs with the same anchors. At order $e^2$ no truncation enters.

(b) After the cut at top index $k$, each localisation is at most $C_{\rm pl}e^2\abs u\kappa\mu_*^2L^3t_k^3\mathrm{e}^{-\delta_1d(\Lambda)}$. This is the size $e(t_{k+1})^2\abs u\delta_k$ with $\delta_k=\kappa\mu_*^2t_{k+1}^2$. If only the $D$-line has index $k$, the remainder is a $C$-only class with its Ward partner. If a scalar line has index $k$, the remainder is convergent. In the case of ties, the cut kernels supply the tree link at rate $\delta_1$. There are $N_s\le3$ localised vertices. Apply Lemma~\ref{app:G1:VS} with $A=C_{\rm pl}e^2\abs u\kappa\mu_*^2L^3t_k^3$, $\delta=\delta_1$ and $r=r(t_k)$. By \cite[(2.8)]{Ba2}, the blocks within $r(t_k)$ of $W_k^c$ lie in $\mathcal L_k$, and the others lie in $\mathcal D_k$. For the layer part, Lemma~\ref{app:G1:layer} gives $t_k^3(2r(t_k)+2b+1)^d\le C(\kappa_0)t_k^{\kappa_0}$. The deep part sums to $C\abs T\sum_k\mathrm{e}^{-\delta_1r(t_k)/2}$, which is finite uniformly in $K$.

(c) This is Proposition~\ref{app:G1:R9vac}(iii) with $X=\Lambda_7^{(k)c}$. Only top index $k$ is localised at step $k$, so there is no factor $1+K-k$.

For the final statement, group $E_1(u)$ by largest index and apply Lemma~\ref{app:G1:S05-TS}: the two occurrences of $\sum_kF^{[k]}_{11}$ coincide, so no bound on $F_{11}(u)$ is used.
\end{proof}

\begin{remark}[collar shares under a localised convention]\label{app:G1:FAii}
Suppose instead that the parts of $E_{K,j}$ not yet moved were removed on the collars. Weight each class by the volume fraction of its scale-$j^*$ anchor block that lies in the removal set, where $j^*$ is the top index. For a scale-$j$ block $\square$ this fraction is $(t_j/t_{j^*})^3$. Each weighted class is bounded per anchor block by Lemma~\ref{app:K:L14-8}. Summing the geometric series in $j^*$ gives
\[
\abs{E_{K,j}(u)\!\restriction_\square}\le C_AC_F\sigma_N(u)\sum_oc_o\,t_j^3\sum_{j^*=j}^{K-1}t_{j^*}^{o-3}\le C^{\rm FA}_{\kappa_0}t_j^{\kappa_0},
\]
with $C^{\rm FA}_{\kappa_0}=C_AC_F\sigma_N(u)C_{\nbar}(e,\lambda)\,\eps_0^{1-\kappa_0}\max\{1,((1-\kappa_0)\log2)^{-1}\}$. Under that convention the order-$e^2$, one-$D$-line part is $\le Ce^2\abs u\kappa\mu_*^2L^2t_j^3(1+K-j)$. The convention of Lemma~\ref{app:G1:S05-RF}(V) does not use these bounds.
\end{remark}

\paragraph{(S5) The closing step and the order of choices.}

\begin{lemma}[constants do not enter the remainder]\label{app:G1:S05-N0}
If $\Phi'=\Phi+c$ with $c$ independent of $(e,\lambda)$, then $\Phi'-T_{\nbar}[\Phi']=\Phi-T_{\nbar}[\Phi]$.
\end{lemma}

\begin{proof}
$T_{\nbar}$ is linear and $T_{\nbar}[c]=c$.
\end{proof}

\begin{lemma}[the final-scale photon law]\label{app:G1:S05-FM}
Under the law $\mathrm{e}^{-\frac12HA}$, at both points, the final-scale field $A_K$ has law $N(0,(\Delta^{(K)})^{-1})$, and
\[
\int dA_0\cdots dA_{K-1}\,\mathrm{e}^{-\frac12HA}=\prod_{k<K}(2\pi)^{n_k/2}\det(w_kC^{(k)})^{1/2}\,\mathrm{e}^{-\frac12D_K(A_K)}.
\]
\end{lemma}

\begin{proof}
By Lemma~\ref{app:G1:S05-T}, $A_0$ occurs only in $I_0$. Integrate out $A_0$, then $A_1$, and so on up to $A_{K-1}$. Each step is a translated Gaussian integral $\int dA_k\,\mathrm{e}^{-\frac12w_k^{-1}I_k}=(2\pi)^{n_k/2}\det(w_kC^{(k)})^{1/2}$, which does not depend on $A_{k+1}$.
\end{proof}

\begin{lemma}[zero modes]\label{app:G1:S05-ZM}
Work in physical units, with $\beta_k=at_{k+1}^{-2}$, $\Delta^{(0)}=-\Delta^\eps+\mu_*^2$ and $\Delta^{(k+1)}=\beta_k-\beta_k^2QC^{(k)}Q^*$. Then for every $k\le K$, uniformly in $\ell_T$ and $\eps$,
\[
\Delta^{(k)}\ge\lambda_*(k):=\Big(\mu_*^{-2}+a^{-1}\sum_{j<k}t_{j+1}^2\Big)^{-1}\ge\gamma_1\mu_*^2,\qquad\gamma_1:=\big(1+4\mu_*^2\eps_0^2/(3a)\big)^{-1},
\]
with equality on the constant mode. The scalar Gaussian of the closing step has mass $M^2$, and $\Delta_s^{(k)}(0)\ge M^2(1+4M^2\eps_0^2/(3a))^{-1}$, with equality on constants.
\end{lemma}

\begin{proof}[Proof of the photon bound]
$(\Delta^{(k+1)})^{-1}$ is the covariance of $A_{k+1}=QA_k+\xi_k$, where $\xi_k\sim N(0,\beta_k^{-1})$ is independent of $A_k\sim N(0,(\Delta^{(k)})^{-1})$. Hence $(\Delta^{(k+1)})^{-1}=Q(\Delta^{(k)})^{-1}Q^*+\beta_k^{-1}\le\norm{(\Delta^{(k)})^{-1}}+\beta_k^{-1}$, using $\norm Q\le1$, and $\norm{(\Delta^{(0)})^{-1}}=\mu_*^{-2}$. Iterating gives the first bound. Since $L=2$, $\sum_{j<k}t_{j+1}^2\le\frac43\eps_0^2$. On constants, $Q1=1$, $Q^*1=1$ and $\Delta^{(0)}1=\mu_*^21$, so by induction $\Delta^{(k)}1=\lambda_*(k)1$. The scalar free bound has the same proof with $\Delta_s^{(0)}=-\Delta^\eps+M^2$ (Lemma~\ref{app:G1:RW-ZMs}).
\end{proof}

\begin{lemma}[the normalisation defect]\label{app:G1:S05-F3}
Assume $\delta_0R\ge4$ and $\hat\xi_{\rm tot}\le\frac1{12}$ ((E.14)). Then for every admissible $\boldsymbol\Lambda$ at both points,
\[
\abs{\log\det HA_{\rm del}(\boldsymbol\Lambda)-\log\det HA}\le B_{F3}\abs T,\qquad B_{F3}:=12\abs u\kappa C_{\ell^1}(1-\hat\xi_{\rm tot})^{-1}w_{\min}^{-2}w_{\rm lo}^{-1}\gamma_1^{-1}\eps_0.
\]
Here $C_{\ell^1},\delta_0$ are the kernel-decay constants of $\Delta^{(j)}$ on $T_\eps$, $w_{\min}=\min_jw_j$, $w_{\rm lo}:=\min(1,\min_kw_k^{-1})$, and $\gamma_1$ is the constant of Lemma~\ref{app:G1:S05-ZM}. Without the condition $\delta_0R\ge4$, the same holds with $B_{F3}$ multiplied by $C_*:=\exp(\max\{0,16/(\delta_0R)-4\})$.
\end{lemma}
\noindent The proof uses $\delta_0R\ge4$ in place of the printed $\delta_0R\ge2$.
\begin{proof}
Work in rescaled units at scale $j$.
(i) $HA_{\rm del}-HA=-2\sum_j\omega_j\hat X_j$. Put $Y:=HA^{-1/2}(HA_{\rm del}-HA)HA^{-1/2}$; then $\norm Y\le\hat\xi_{\rm tot}$.
(ii) $\abs{\tr\log(I+Y)}\le(1-\norm Y)^{-1}\norm{Y}_1$.
(iii) $\norm{HA^{-1/2}P_j\hat X_jP_jHA^{-1/2}}_1\le\norm{P_jHA^{-1}P_j}\,\norm{\hat X_j}_1$. Since $HA\ge w_{\rm lo}S^1$ (Lemma~\ref{app:G1:S05-T}), $P_jHA^{-1}P_j\le w_{\rm lo}^{-1}(\Delta^{(j)})^{-1}\le(w_{\rm lo}\gamma_1\mu_*^2t_j^2)^{-1}$ by Lemma~\ref{app:G1:S05-ZM}.
(iv) On level $j$, $\dist(\Lambda_7^{(j-1)\prime},\Lambda_6^{(j-1)\prime c})>r(t_{j-1})/L\ge r(t_j)/2$, so $\norm{\hat X_j}_1\le C_{\ell^1}\mathrm{e}^{-\delta_0r(t_j)/4}N_j$; the periodic images are summed into $C_{\ell^1}$.
(v) $2\abs{\omega_j}\le6\abs u\kappa\mu_*^2t_j^2w_{\min}^{-2}$, where $\delta_j-\delta_{j-1}=3\kappa\mu_*^2t_j^2$ at $L=2$; this also holds for $j=K$. The factor $\mu_*^2t_j^2$ cancels against (iii).
(vi) With $x=\log t_j^{-1}$ and $(1+x)^2\ge4x$: $\mathrm{e}^{-\delta_0R(1+x)^2/4}\le t_j^{\delta_0R}\le t_j^4$. Hence $\sum_j\mathrm{e}^{-\delta_0r(t_j)/4}N_j\le\abs T\sum_jt_j\le2\eps_0\abs T$. Collecting (i)--(vi) gives the claim. For general $R$, bound $\mathrm{e}^{-\delta_0R(1+x)^2/4}\le C_*t^4$ instead.
\end{proof}

\begin{lemma}[the closing step]\label{app:G1:S05-CL}
Assume (IH$_K$) at both points, and the hypotheses of Lemmas~\ref{app:G1:I145-comb} and \ref{app:G1:S05-F3}. Then
\[
\mathrm{e}^{-E_0^w-E_1(u)}\mathcal Z_u\le\exp\big(C_+^B(u;M)\abs T\big),\qquad C_+^B(u;M)=C_{F1}\eps_0^{\kappa_0-3}+B_E^{\rm fin}(u,M)+B_{\mathfrak E}+B_{GS}+\tfrac12B_{F3}+B_{\rm comb},
\]
with $B_{\mathfrak E}=(1-2^{-1/5})^{-1}(C_\kappa+C_\nu)\eps_0^{1/5}$, $B_{GS}=2dC_4'\eps_0$, $B_{F3}$ as in Lemma~\ref{app:G1:S05-F3}, and
\[
B_{\rm comb}=c_{\rm fin}\eps_0^{-3}+2(d+2)\Sigma_*(\eps_0;\pi)
\]
(Lemma~\ref{app:G1:I145-comb} with $c_*=c_0'^{\,w}$, $2p=12$).
\end{lemma}
\noindent The proof is Steps~4--5 of the proof of Proposition~\ref{app:G1:assembly} (C1--C5): $E_0^w$ cancels by Lemma~\ref{app:G1:S05-N0}, the zero modes are handled by Lemmas~\ref{app:G1:S05-FM} and \ref{app:G1:S05-ZM}, the normalisation defect by Lemma~\ref{app:G1:S05-F3}, and the region sum by Lemma~\ref{app:G1:I145-comb}. The factor $\mathrm{e}^{-E_{\rm fin}(u)}$ is paid once, by Propositions~\ref{app:G1:Efin} and \ref{app:G1:I8}.

\begin{lemma}[the order of choices]\label{app:G1:S05-ACY}
(i) In the order of the stages 0--11 of Appendix~\ref{app:B}, every constant depends only on data of earlier stages and on data of its own stage that are fixed before it. Dependences within one stage occur only in two ways, and both are one-way. First, at stage 0 ($\kappa$ and $\ell_0^\natural$ are computed from the point), at the stages 6--8 (constants computed from $p$ and $\mathbb W_0$) and at stage 10a (suprema over $(0,\eps_1]$ computed from $\eps_1$), a quantity is computed from data fixed before it in the same stage. Second, at the stages 5, 10b and 11 the conditions of the stage determine $R$, $\eps_0$ and $\ell_0$, respectively, and no constant entering a condition of that stage depends on the quantity it determines. In particular, the upper bounds on $\eps_0$ depend, besides $\eps_0$, on data of the stages $\le9$ together with stage 10a, and the lower bounds on $\ell_0$ (stage 11) depend on data of the stages $\le$ 10b.
(ii) Fix a floor $\ell_0^\natural>0$ at stage 0, and evaluate every constant that depends on $\ell_T$ and enters an upper bound on $\eps_0$ at $\ell_T=\ell_0^\natural$. Then each value bounds the constant for all $\ell_T\ge\ell_0^\natural$. Hence, with $\ell_0:=\max\{\ell_0^\natural,\ell_0^*(\eps_0)\}$ at stage 11, every upper bound on $\eps_0$ decided at stage 10 holds for every torus $T_j$, and no upper bound on $\eps_0$ depends on $\ell_0$.
\end{lemma}
\noindent Part (i) is Lemma~\ref{app:G1:acyc}; part (ii) is Lemma~\ref{app:G1:floor} for an arbitrary fixed floor.
\begin{proof}
(i) Evaluate the constants that depend on $\ell_T$ as in (ii) at the floor $\ell_0^\natural$ of Lemma~\ref{app:G1:floor}, and replace the $\eps_0$-dependent constants entering stage 10b by their suprema over $(0,\eps_1]$ (Lemma~\ref{app:G1:L10asup}). Then, by Lemma~\ref{app:G1:acyc}, every dependence goes from a stage to the same or a later stage, and the dependences within one stage are those listed in (i). The latter are read off the ``depends on'' columns of Table~\ref{p:tab:stages} and Tables~\ref{p:tab:A2}--\ref{p:tab:A5}. The upper bounds on $\eps_0$ (Table~\ref{p:tab:A4}) involve, besides $\eps_0$, only data of the stages $\le9$ and the suprema of stage 10a. The lower bounds on $\ell_0$ (Table~\ref{p:tab:A5}) involve only data of the stages $\le$ 10b. In particular, the constants $c_5$ and $c_7$ of condition (E.ent) are bounded uniformly in $\eps_0\in(0,\eps_1]$ by Statement~\ref{app:K:SL}, so (E.ent) depends on no datum of stage 10b other than $\eps_0$.

(ii) This is Lemma~\ref{app:G1:floor}. The $\ell_T$-dependent constants are the following, and all are non-increasing in $\ell_T$: $C_\delta(M)$, through its torus zero-mode majorants (Lemma~\ref{app:G1:floor}); the loss in the extension of \cite[Prop.~2.1]{Ba2} to $Q^{(k)}$ through $C_\delta$; the terms $(M^2\ell_T^3)^{-1}$, $(\mu_*^2\ell_T^3)^{-1}$, $\ell_T^{-1}$, $(M\ell_T)^{-10}$, $(\mu_*\ell_T)^{-10}$ of the bound on $\delta m^2_{\rm fin}$, whose logarithmic term is bounded by $\log_+(1/(\min(M,\mu_*)\eps_0))$ for every $\ell_T$; the photon zero mode and tail terms; and the twist defect $\mathrm{e}^{-\delta\ell_T/t}$. For non-increasing $f$, $f(\ell_T)\le f(\ell_0^\natural)$ when $\ell_T\ge\ell_0^\natural$. Every upper bound on $\eps_0$ is monotone in each constant in the direction ``larger constant $\Rightarrow$ stronger condition''. So upper bounds on the constants keep the conditions sufficient, and $\ell_0\ge\ell_0^\natural$ holds by (T.floor). For $\delta m^2_{\rm fin}$ the majorant is the explicit bound of Proposition~\ref{app:G1:I3b}.
\end{proof}

\subsection{The cumulant bound, the Dirichlet-to-free comparison and the collars}\label{app:G1:cumcoll}

\paragraph{Setting of the cumulant bound.}\label{app:G1:part-R1}
Let $\mathbb L=\prod_{a=1}^d\Z/N_a\Z$ carry the torus Euclidean distance $\abs{x-y}$, and let $\Lambda\subset\mathbb L$. Fields are $z=(z_x)_{x\in\Lambda}$ with $z_x\in\R^n$, and $\abs{\cdot}$ is the Euclidean norm on $\R^n$. Let $\mu=N(0,C)$ with $C=P^{-1}$, where $P$ is real symmetric and
\[
\text{(Pa)}\quad P\ge\gamma_0,\qquad\qquad \text{(Pb)}\quad \norm{P(x,y)}\le c_0\mathrm{e}^{-\delta_0\abs{x-y}}
\]
(norms of $n\times n$ blocks). Every $P$ of the class $\mathcal P(\gamma_0,\gamma_1,c_0,\delta_0)$ of Lemma~\ref{app:G1:PH-CL} satisfies (Pa) and (Pb). A polynomial $F=\sum_YF_Y$ is a sum of real polynomials $F_Y$ in $(z_y)_{y\in Y}$. If every monomial is replaced by the set of sites it actually depends on, the norms below do not increase. We therefore assume that every monomial of $F_Y$ depends on every $z_y$, $y\in Y$. Let $\abs{F_Y}_1$ be the sum of the absolute values of the coefficients. Let $\ell(Y)$ be the minimal total length of a connected graph whose vertex set contains $Y$; then $\ell(Y')\le\ell(Y)$ for $Y'\subset Y$, and $\abs{x-y}\le\ell(Y)$ for $x,y\in Y$. Put $\norm F_\theta:=\sup_x\sum_{Y\ni x}\mathrm{e}^{\theta\ell(Y)}\abs{F_Y}_1$. Let $V=\sum_{Y\subset J}V_Y$ with $J\subset\Lambda$, of degree $\le D$ and without constant term, and put $\mathsf A:=\norm V_\kappa$ and $V_R:=\sum_{Y\subset R}V_Y$ for $R\subset\Lambda$. Let $K\subset\R^n$ be convex with $B(1)\subset K\subset B(c_K)$, and put $\chi^X_\beta:=\prod_{x\in X}\mathbf 1\{z_x\in\beta K\}$ and $\chi_\beta:=\chi^\Lambda_\beta$. We write $\kappa_\nu(F_1,\dots,F_k)$ for the joint cumulant under $\nu$, and $\kappa_\nu(F;k)$ for the $k$-th cumulant.

\emph{Constants.} Put $\tilde c(\theta):=\sum_{z\in\Z^d}\mathrm{e}^{-\theta\abs z}$, $K_d:=\sum_z\abs z\mathrm{e}^{-\delta_0\abs z/2}$, $\mu_{\rm c}:=\min(\delta_0/2,\gamma_0/(2c_0K_d))$, $g_0:=2/\gamma_0$, $g:=\max(1,g_0)$, $M_1:=g_0c_0\tilde c(\mu_{\rm c}/2)^2$, $c_{\rm D}:=g_0^2c_0\tilde c(\mu_{\rm c}/2)^2$, $\gamma_K:=\min\{\frac12,(2c_KM_1)^{-1}\}$ and $\bar\mu:=\min(\mu_{\rm c},\kappa)$. For $g'\ge1$ put $K_k(\theta;g'):=(D+1)^k(kD)^{kD/2}g'^{\,kD/2}\tilde c(\theta/2)^{k-1}$, and $K_k(\theta):=K_k(\theta;g)$. For fixed $x$, the map sending $y\in\mathbb L$ to the representative of minimal norm of $y-x$ is injective into $\Z^d$. Hence $\sum_{y\in\mathbb L}f(\abs{x-y})\le\sum_{z\in\Z^d}f(\abs z)$ for every $f\ge0$; this is used without comment.

\begin{theorem}[cumulant bound]\label{app:K:L13-1}\label{app:G1:PH-R1}
Let $t\in\N$, $\rho_1:=D+2d$ and $\rho_3:=\bar\mu/16$. There are $b_*,S<\infty$, depending only on $(d,n,D,t,\kappa,\gamma_0,c_0,\delta_0,c_K)$, such that for all $\mathbb L$, $\Lambda$, $J$, $V$ and $P$ as above, all $b\ge b_*$ and $\mathsf Ab^{\rho_1}\le1$,
\[
\int\chi_b\mathrm{e}^V\,d\mu\le\exp\Big\{\sum_{m=1}^t\frac{\kappa_\mu(V;m)}{m!}+S\abs J\big[(\mathsf Ab^{\rho_1})^{t+1}+\mathrm{e}^{-\rho_3b^{3/2}}\big]\Big\}.
\]
\end{theorem}

\begin{lemma}[covariances]\label{app:G1:R1-cov}
Let $X\subset\Lambda$, $Y:=\Lambda\setminus X$ and $G_X:=(P_{XX})^{-1}$.
(a) $\norm{G_X(x,y)}\le g_0\mathrm{e}^{-\mu_{\rm c}\abs{x-y}}$; in particular this holds for $C=G_\Lambda$.
(b) $G_X\le C_{XX}\le\gamma_0^{-1}$.
(c) For every $z_Y$, $h:=-G_XP_{XY}z_Y$ satisfies $\abs{h(x)}\le M_1\mathrm{e}^{-(\mu_{\rm c}/2)\dist(x,Y)}\sup_{v\in Y}\abs{z_v}$.
(d) $\norm{(C-G_X)(x,y)}\le c_{\rm D}\mathrm{e}^{-(\mu_{\rm c}/8)(\dist(x,Y)+\dist(y,Y))-(\mu_{\rm c}/4)\abs{x-y}}$ for $x,y\in X$.
\end{lemma}
\begin{proof}
(a) Let $\rho$ be $1$-Lipschitz and put $P^\rho:=\mathrm{e}^{\mu_{\rm c}\rho}P_{XX}\mathrm{e}^{-\mu_{\rm c}\rho}=P_{XX}+E$. Since $\abs{\mathrm{e}^s-1}\le\abs s\mathrm{e}^{\abs s}$ and $\mu_{\rm c}\le\delta_0/2$, we have $\norm{E(x,y)}\le c_0\mu_{\rm c}\abs{x-y}\mathrm{e}^{-\delta_0\abs{x-y}/2}$. The Schur test gives $\norm E\le c_0\mu_{\rm c}K_d\le\gamma_0/2$. Since $P_{XX}\ge\gamma_0$, we get $\ip f{P^\rho f}\ge\frac{\gamma_0}2\norm f^2$ for real $f$, hence $\norm{(P^\rho)^{-1}}\le g_0$. Take $\rho=\abs{\cdot-y}$ and use $G_X(x,y)=\mathrm{e}^{-\mu_{\rm c}\rho(x)}(P^\rho)^{-1}(x,y)$.
(b) By block inversion, $C_{XX}=(P_{XX}-P_{XY}P_{YY}^{-1}P_{YX})^{-1}$. The subtracted term is $\ge0$, and inversion reverses the order of positive operators, so $G_X\le C_{XX}$. Moreover $C_{XX}$ is a compression of $C\le\gamma_0^{-1}$.
(c) For $u\in X$ and $v\in Y$ we have $\norm{G_X(x,u)}\norm{P(u,v)}\le g_0c_0\mathrm{e}^{-\mu_{\rm c}\abs{x-u}-2\mu_{\rm c}\abs{u-v}}$ and $\abs{x-u}+\abs{u-v}\ge\dist(x,Y)$. Hence $\abs{h(x)}\le g_0c_0\mathrm{e}^{-(\mu_{\rm c}/2)\dist(x,Y)}\sum_{u,v}\mathrm{e}^{-(\mu_{\rm c}/2)(\abs{x-u}+\abs{u-v})}\sup_{v\in Y}\abs{z_v}$, and the double sum is at most $\tilde c(\mu_{\rm c}/2)^2$.
(d) The $X$-rows of $PC=I$ give $C_{XX}-G_X=-G_XP_{XY}C_{YX}$. By (a) and (Pb), the kernel is bounded by $g_0^2c_0\sum_{u\in X,v\in Y}\mathrm{e}^{-\mu_{\rm c}L}$ with $L=\abs{x-u}+\abs{u-v}+\abs{v-y}$. Here $L\ge\abs{x-y}$, and $L\ge\max(\dist(x,Y),\dist(y,Y))\ge\frac12(\dist(x,Y)+\dist(y,Y))$. Write $\mathrm{e}^{-\mu_{\rm c}L}=\mathrm{e}^{-\mu_{\rm c}L/4}\mathrm{e}^{-\mu_{\rm c}L/4}\mathrm{e}^{-\mu_{\rm c}L/2}$ and sum the last factor over $u,v$.
\end{proof}

\begin{lemma}[bounded variables, tails]\label{app:G1:R1-bdd}
Let $\nu$ be a probability measure.
(a) If $\abs{X_j}\le M_j$, then $\abs{\kappa_\nu(X_1,\dots,X_k)}\le k!\,2^{k-1}\prod_jM_j$. If $\abs X\le M$, then $\abs{\log\int\mathrm{e}^Xd\nu-\sum_{m\le t}\kappa_\nu(X;m)/m!}\le(2M)^{t+1}$.
(b) If $\chi$ is an indicator with $\nu(1-\chi)\le\varepsilon\le\frac12$ and $\abs W\le M$, then $\int\chi\mathrm{e}^Wd\nu\le\int\mathrm{e}^{\chi W}d\nu\le\mathrm{e}^{2\mathrm{e}^M\varepsilon}\int\chi\mathrm{e}^Wd\nu$.
(c) If $\nu=N(m,G)$ on $(\R^n)^X$ with $\abs{m_x}\le\beta/2$ and $G\le\gamma_0^{-1}$, then $\nu(\exists x:\abs{z_x}>\beta)\le\varepsilon(\abs X,\beta):=2n\abs X\mathrm{e}^{-\gamma_0\beta^2/(8n)}$, and $\nu(\abs{z_x-m_x}>s)\le2n\mathrm{e}^{-\gamma_0s^2/(2n)}$.
\end{lemma}
\begin{proof}
(a) By M{\"o}bius inversion, $\kappa=\sum_\pi(-1)^{\abs\pi-1}(\abs\pi-1)!\prod_{B\in\pi}\E\prod_{j\in B}X_j$, and $\sum_\pi(\abs\pi-1)!$ is at most the number of ordered set partitions, which is $\le k!\,2^{k-1}$ (cut a permutation into consecutive blocks). For the second claim, $f(s)=\log\int\mathrm{e}^{sX}d\nu$ has $f(0)=0$ and $f^{(m)}(s)=\kappa_{\nu_s}(X;m)$ with $\nu_s\propto\mathrm{e}^{sX}\nu$. Use Taylor's formula with the Lagrange remainder and the first claim.
(b) $\chi\mathrm{e}^W\le\mathrm{e}^{\chi W}$. Also $\int\mathrm{e}^{\chi W}=\int\chi\mathrm{e}^W+\nu(1-\chi)$ and $\int\chi\mathrm{e}^W\ge\mathrm{e}^{-M}/2$, so the middle integral is at most $(1+2\varepsilon\mathrm{e}^M)\int\chi\mathrm{e}^W$.
(c) If $\abs{z_x}>\beta$, then some component of $z_x-m_x$ exceeds $\beta/(2\sqrt n)$ in modulus. Its variance is $\le\gamma_0^{-1}$, so the Gaussian tail bound and a union bound give the first claim. The second claim follows in the same way, with $s/\sqrt n$ in place of $\beta/(2\sqrt n)$.
\end{proof}

\begin{lemma}[connected Wick bound]\label{app:G1:R1-conn}
(a) Let $\zeta$ be centred Gaussian with covariance $G$, and let $F_1,\dots,F_k$ be monomials, each a multiset of coordinate legs. Then $\kappa(F_1,\dots,F_k)=\sum_\pi\prod_{\{u,v\}\in\pi}G(u,v)$. The sum runs over the pairings $\pi$ of all legs for which the following graph on $\{1,\dots,k\}$ is connected: it has an edge for each pair joining two factors.
(b) Let $\nu=N(m,G)$ on $(\R^n)^X$ with $\norm{G(x,y)}\le g'\mathrm{e}^{-\theta\abs{x-y}}$ for some $g'\ge1$, and $\abs{m_x}\le\beta$. Let each $F_j$ have degree $\le D$ and $\norm{F_j}_{\theta'}\le a_j$, where $\theta'\le\theta$. Suppose every support of $F_1$ meets $S_1$ and every support of $F_2$ lies in $S_2$ (for $k=1$ put $S_2:=\mathbb L$). Then
\[
\abs{\kappa_\nu(F_1,\dots,F_k)}\le\abs{S_1\cap X}\,K_k(\theta';g')(1+\beta)^{kD}\prod_ja_j\,\mathrm{e}^{-(\theta'/2)\dist(S_1,S_2)}.
\]
(c) For $\beta=0$, the bound of (b) also holds for every sum over the connected pairings of (a) in which each pair $\{u,v\}$ carries a matrix entry of a kernel $w_{\{u,v\}}$ with $\norm{w_{\{u,v\}}(x,y)}\le g'\mathrm{e}^{-\theta\abs{x-y}}$ (the kernels may differ from pair to pair).
\end{lemma}
\begin{proof}
(a) By Isserlis' theorem, $\E\prod_{j\in B}F_j$ is the sum over all pairings of the legs of $B$. Group the pairings by the partition of $B$ into connected components. This gives $\E\prod_BF=\sum_\sigma\prod_{b\in\sigma}\mathcal K(b)$, where $\mathcal K(b)$ denotes the connected sums. The moment--cumulant relations determine the cumulants recursively, so $\mathcal K=\kappa$.
(b) Substitute $z=m+\zeta$ and keep every new monomial under its original support. Coefficient sums grow by at most $\prod(1+\abs{m_y^a})^{\alpha_{y,a}}\le(1+\beta)^D$. For $k\ge2$, constants do not change cumulants and are dropped. Then apply (a) term by term. With $d_j\le D$ legs per factor, there are at most $(D+1)^k$ degree patterns and at most $(kD)^{kD/2}$ pairings. Each pair contributes at most $g'\mathrm{e}^{-\theta\abs{u-v}}$, and there are at most $kD/2$ pairs.
Choose a spanning tree of the factor graph rooted at factor $1$, and bound the non-tree pairs by $g'$. Sum the root over $a\in S_1\cap X$ and $Y_1\ni a$. Sum each child $j$ over the site $v\in Y_j$ of its tree leg and over $Y_j\ni v$. Write $\abs{F_{j,Y}}_1=\mathrm{e}^{-\theta'\ell(Y)}\big(\mathrm{e}^{\theta'\ell(Y)}\abs{F_{j,Y}}_1\big)$ and $\mathrm{e}^{-\theta\abs{u-v}}\le\mathrm{e}^{-(\theta'/2)\abs{u-v}}\mathrm{e}^{-(\theta'/2)\abs{u-v}}$. Follow the tree path from factor $1$ to factor $2$. It goes from $a$ to a point of $Y_2\subset S_2$, through tree pairs and through segments joining two points of the same $Y_j$, each of length $\le\ell(Y_j)$. Hence $\dist(S_1,S_2)\le\sum_j\ell(Y_j)+\sum_{\rm tree}\abs{u-v}$, which gives the factor $\mathrm{e}^{-(\theta'/2)\dist(S_1,S_2)}$. The remaining factors give $\tilde c(\theta'/2)$ per tree pair and $\norm{F_j}_{\theta'}$ per factor.
(c) The argument of (b) uses only the bound on each pair.
\end{proof}

\begin{lemma}[removing $\chi$; replacing the law]\label{app:G1:R1-chi}
Let $F_1,\dots,F_k$ ($k\le t$) have degree $\le D$, supports in a set $X$ with $\abs X\le v$, and $\norm{F_j}_0\le a_j$.
(a) Let $\nu=N(m,G)$ with $\abs{m_x}\le\beta/2$, $G\le\gamma_0^{-1}$ and $\beta\ge2$. Let $\chi\ge\prod_{x\in X}\mathbf 1\{\abs{z_x}\le\beta\}$ be an indicator, and put $\varepsilon:=\nu(1-\chi)$. Assume
\[
\text{(MC)}\qquad 2^qv(2n)^{1/2}\big(n^q(2q-1)!!\big)^{1/2}\gamma_0^{-q/2}\mathrm{e}^{-\gamma_0\beta^2/(4n)}\le(2\beta)^q\quad\text{for all }q\le2tD .
\]
Then $\abs{\kappa_\nu(\chi F_1,\dots,\chi F_k)-\kappa_\nu(F_1,\dots,F_k)}\le k\,k!\,4^k\prod_j(a_jv)(2\beta)^{kD}\varepsilon^{1/2}$.
(b) Let $\nu,\tilde\nu$ be Gaussian laws of $z_X$ with means bounded by $R$, block covariances of norm $\le\gamma_0^{-1}$, $\sup\abs{m-\tilde m}\le\eta_1$ and $\sup\norm{G-\tilde G}\le\eta_2$. Then
\[
\abs{\kappa_\nu(F)-\kappa_{\tilde\nu}(F)}\le k\,k!\,4^k(kD)^{kD+1}(1+R+\gamma_0^{-1})^{kD}\prod_j(a_jv)\max(\eta_1,\eta_2).
\]
\end{lemma}
\begin{proof}
(a) Put $Z:=\max(1,\max_{x\in X}\abs{z_x})$, $\zeta:=z-m$ and $A:=\{\max_x\abs{\zeta_x}>\beta\}$. Off $A$ we have $Z\le2\beta$, and on $A$ we have $Z\le2\max_x\abs{\zeta_x}$. By Jensen's inequality and the Gaussian moments, $\E\abs{\zeta_x}^{2q}\le n^q(2q-1)!!\gamma_0^{-q}$, so $\E\max_x\abs{\zeta_x}^{2q}\le\abs X\,n^q(2q-1)!!\gamma_0^{-q}$. By Lemma~\ref{app:G1:R1-bdd}(c), $\nu(A)\le2n\abs X\mathrm{e}^{-\gamma_0\beta^2/(2n)}$. The Cauchy--Schwarz inequality gives
\[
\E\big[Z^q\mathbf 1_A\big]\le2^q\big(\E\max_x\abs{\zeta_x}^{2q}\big)^{1/2}\nu(A)^{1/2}\le2^qv(2n)^{1/2}\big(n^q(2q-1)!!\big)^{1/2}\gamma_0^{-q/2}\mathrm{e}^{-\gamma_0\beta^2/(4n)},
\]
which is at most $(2\beta)^q$ by (MC). Hence $\E Z^q\le2(2\beta)^q$ for $q\le2tD$. Also $\abs{F_j}\le a_jvZ^D$. In the M{\"o}bius formula $\chi^{\abs B}=\chi$. Each block moment changes by $\abs{\E(1-\chi)\prod_BF}\le\varepsilon^{1/2}(\E\prod_BF^2)^{1/2}\le\sqrt2\,\varepsilon^{1/2}\prod_B(a_jv)(2\beta)^{\abs BD}$, and each block moment is at most $2\prod_B(a_jv)(2\beta)^{\abs BD}$. Telescope over the at most $k$ blocks and use the count of Lemma~\ref{app:G1:R1-bdd}(a).
(b) A moment of $L\le kD$ coordinates is a sum over at most $L!\le L^L$ partial pairings of products of at most $L$ factors (means and covariance entries), each at most $1+R+\gamma_0^{-1}$. Replacing one factor at a time changes such a product by at most $L\max(\eta_1,\eta_2)(1+R+\gamma_0^{-1})^{L}$. Coefficient sums are at most $\prod(a_jv)$. Then conclude as in (a).
\end{proof}

\begin{lemma}[decoupling]\label{app:G1:R1-dec}
Let $\Lambda=\Gamma\sqcup\Delta$ with $\Delta=\bigsqcup_i\Delta_i$ and $\dist(\Delta_i,\Delta_j)\ge w_s$ for $i\ne j$. Put $\omega:=c_0\tilde c(\delta_0/2)\mathrm{e}^{-\delta_0w_s/2}$ and assume $\omega\le\gamma_0/2$. Let $\nu_{\rm loc}(\cdot|z_\Gamma):=\bigotimes_iN(-G_{\Delta_i}P_{\Delta_i\Gamma}z_\Gamma,G_{\Delta_i})$. If $\abs{z_x}\le R$ for all $x\in\Lambda$, then
\[
\Big|\log\frac{d\mu(\cdot|z_\Gamma)}{d\nu_{\rm loc}(\cdot|z_\Gamma)}(z_\Delta)\Big|\le\omega_{\rm tot}:=\omega\abs\Delta\Big[\frac{R^2}2+\frac n{\gamma_0}+\frac{(c_0\tilde c(\delta_0)R)^2}{2\gamma_0^2}\Big].
\]
Hence, if $\abs{z_x}\le R$ on $\Gamma$, then $\mathrm{e}^{-\omega_{\rm tot}}\int F\,d\nu_{\rm loc}\le\int F\,d\mu(\cdot|z_\Gamma)\le\mathrm{e}^{\omega_{\rm tot}}\int F\,d\nu_{\rm loc}$ for every $F\ge0$ that vanishes unless $\abs{z_x}\le R$ on $\Delta$.
\end{lemma}
\begin{proof}
Let $B:=\bigoplus_iP_{\Delta_i\Delta_i}$, $O:=P_{\Delta\Delta}-B$ and $u:=P_{\Delta\Gamma}z_\Gamma$. Both densities are proportional to $\mathrm{e}^{-\frac12\ip{z_\Delta}{(\cdot)z_\Delta}-\ip{z_\Delta}u}$, with $P_{\Delta\Delta}$ and $B$ respectively. Comparing the normalisations gives
\[
\log\frac{d\mu}{d\nu_{\rm loc}}=-\tfrac12\ip{z_\Delta}{Oz_\Delta}+\tfrac12\log\det(P_{\Delta\Delta}B^{-1})-\tfrac12\ip u{(P_{\Delta\Delta}^{-1}-B^{-1})u}.
\]
The kernel of $O$ vanishes unless $\abs{x-y}\ge w_s$, so the Schur test gives $\norm O\le\omega$. Since $B,P_{\Delta\Delta}\ge\gamma_0$, the $n\abs\Delta$ eigenvalues of $B^{-1/2}OB^{-1/2}$ are at most $\omega/\gamma_0\le\frac12$ in modulus, and $\abs{\log(1+\lambda)}\le2\abs\lambda$ there. Moreover $\norm{P_{\Delta\Delta}^{-1}-B^{-1}}=\norm{P_{\Delta\Delta}^{-1}OB^{-1}}\le\omega/\gamma_0^2$, and $\abs{u_x}\le c_0\tilde c(\delta_0)R$.
\end{proof}

\begin{lemma}[pavements]\label{app:G1:R1-geo}
Let $b\ge b_{\rm geo}:=\max\{16,36(d+1)^2\}$, $\ell_b:=\lceil b^2\rceil$, $w:=\lceil b^{3/2}\rceil$ and $s:=\lfloor\ell_b/(d+1)\rfloor$; then $s\ge5w+2$. Call a coordinate $a$ cut if $N_a\ge\ell_b$. For each cut $a$, choose half-integer cuts $B^0_a\subset\Z/N_a\Z+\frac12$ whose cyclic gaps all lie in $[\ell_b,2\ell_b)$, and put $B^i_a:=B^0_a+is$ for $i=0,\dots,d$. The cells of $\mathcal P^i$ are the products of the arcs between consecutive cuts (whole cycles in the coordinates that are not cut). Let $\lambda_i(x):=\min_{a\ \rm cut}\dist(x_a,B^i_a)$, with $\lambda_i\equiv+\infty$ if no coordinate is cut. Then:
(G1) $\lambda_i$ is $1$-Lipschitz;
(G2) if $x,y$ lie in different cells of $\mathcal P^i$, then $\abs{x-y}\ge\lambda_i(x)+\lambda_i(y)$;
(G3) $\bigcap_{i=0}^d\{\lambda_i<2.5w+1\}=\emptyset$;
(G4) every cell has at most $v_b:=(4b^2)^d$ sites.
\end{lemma}
\begin{proof}
The gaps can be chosen: with $m=\lfloor N/\ell_b\rfloor\ge1$ cuts, take one gap $\ell_b+(N-m\ell_b)$ and $m-1$ gaps $\ell_b$. Next, $s\ge b^2/(d+1)-1\ge6b^{3/2}-1\ge5b^{3/2}+7\ge5w+2$, by $b^{1/2}\ge6(d+1)$ and $b\ge16$.
(G1) holds coordinatewise.
(G2) Some cut coordinate $a$ separates $x$ and $y$. The shorter arc between $x_a$ and $y_a$ contains a cut, so $\abs{x-y}\ge\abs{x_a-y_a}\ge\dist(x_a,B^i_a)+\dist(y_a,B^i_a)$.
(G3) Two cuts of $B^i_a$ and $B^{i'}_a$ with $i\ne i'$ are at least $s$ apart. If they come from the same cut, they differ by $\abs{i-i'}s\in[s,ds]$, and $N_a\ge(d+1)s$. Otherwise their distance is at least $\ell_b-ds\ge s$. Suppose $x$ lay in all $d+1$ sets. Then two indices would use the same coordinate, giving two cuts less than $5w+2\le s$ apart.
(G4) The sides are at most $2\ell_b-1\le4b^2$.
\end{proof}

\paragraph{One stage.}\label{app:G1:R1-onestage}
Fix $i$, write $\lambda=\lambda_i$, and intersect all sets with $\Lambda$. For a cell $\square$ put $\Gamma_1=\square\cap\{\lambda<w\}$, $\Gamma_2=\square\cap\{w\le\lambda<2w\}$, $\Gamma_3=\square\cap\{1.5w\le\lambda<2w\}$, $\square'=\square\cap\{\lambda\ge2w\}$, $\Gamma_4=\square\cap\{2w\le\lambda<2.5w\}$, $\mathrm{Dp}(\square)=\square\cap\{\lambda\ge2.5w\}$ and $\Delta_\square=\square\cap\{\lambda\ge w\}$. By (G1)--(G2):
\begin{itemize}
\item distinct $\Delta_\square$ are at least $2w$ apart;
\item $\mathrm{Dp}(\square)$ is at distance $\ge\frac w2$ from $\square\cap\{\lambda<2w\}$ and from $\{\lambda<2.5w\}\setminus\Gamma_4(\square)$, and at distance $\ge2.5w$ from the other cells;
\item $\square\cap\{\lambda\ge1.5w\}$ is at distance $\ge\frac w2$ from $\Lambda\setminus\Delta_\square$.
\end{itemize}
Let $J'\subset J$. Let $\mathcal B'$ be the set of cells meeting $J'$, put $N':=\abs{\mathcal B'}\le\abs{J'}$, $\Delta:=\bigcup_{\mathcal B'}\Delta_\square$, $\Gamma:=\Lambda\setminus\Delta$ and $J'':=J'\cap\{\lambda<2.5w\}$. We classify the $V_Y$ with $Y\subset J'$; here $\square\in\mathcal B'$:
\begin{itemize}
\item $V_\Gamma:=\sum_{Y\subset\Gamma\cap J'}V_Y$;
\item $\Psi_1$: $Y\subset\square\cap\{\lambda\ge1.5w\}$ and $Y\cap\square'\ne\emptyset$; it splits into $\Psi_1'$ ($Y\subset\Gamma_3\cup\Gamma_4$) and $\Psi_1''$ ($Y\cap\mathrm{Dp}\ne\emptyset$);
\item $\Psi_2$: $Y\subset\Gamma_1\cup\Gamma_2$ and $Y\cap\Gamma_2\ne\emptyset$;
\item $\Psi_3$: $Y\subset\Delta_\square$, and $Y$ meets both $\square'$ and $\Gamma_2\setminus\Gamma_3$.
\end{itemize}
Put $\Psi_\square:=\Psi_1+\Psi_2+\Psi_3$, $W_\square:=\Psi_1'+\Psi_2$, $R_0:=V_{J'}-V_\Gamma-\sum\Psi_\square$ and $R_1:=V_{J''}-V_\Gamma-\sum W_\square$. These classes are disjoint. $V_\Gamma$ and the $W_\square$ are parts of $V_{J''}$, because points of $J'\cap\Gamma$ have $\lambda<w$.

\emph{Length facts.} A support of $R_0$ meets some $\Delta_\square$ and has $\ell\ge w$. A support of $R_1$ meets some $\Delta_\square$ and has $\ell\ge\frac w2$. A support of $\Psi_3$ has $\ell\ge\frac w2$.
To see the first, note that $Y\not\subset\Gamma$, so $Y$ meets some $\Delta_\square$. If $Y\not\subset\square$, then $\ell\ge w$ by (G2). If $Y\subset\square$, then $Y$ meets $\square'$ (otherwise it is in $\Psi_2$), meets $\{\lambda<1.5w\}$ (otherwise it is in $\Psi_1$), and meets $\Gamma_1$ (otherwise it is in $\Psi_3$); so $\ell\ge w$ by (G1). For $R_1$: if $Y\subset\square\cap J''$, then $Y$ meets $\square'$ (otherwise it is in $\Psi_2$) and meets $\{\lambda<1.5w\}$ (otherwise $Y\subset\Gamma_3\cup\Gamma_4$, i.e.\ it is in $\Psi_1'$).
We have $\sum_{Y\ni x,\ell(Y)\ge r}\abs{V_Y}_1\le\mathsf A\mathrm{e}^{-\kappa r}$. Hence, on $\{\abs{z_x}\le R\ \forall x\}$ with $R\ge1$: $\abs{\Psi_\square},\abs{W_\square}\le M:=\mathsf Av_bR^D$, $\abs{\Psi_3}\le M\mathrm{e}^{-\kappa w/2}$, $\abs{R_0}\le N'M\mathrm{e}^{-\kappa w}$ and $\abs{R_1}\le N'M\mathrm{e}^{-\kappa w/2}$.

\emph{Parameters.} Put $R:=c_K\gamma_K^{-(d+1)}b$ and $m_*:=4^d(c_K\gamma_K^{-(d+1)})^D$, so that $M=m_*\mathsf Ab^{\rho_1}$. Put $\varepsilon_b:=\varepsilon(v_b,b)$ and $\omega:=c_0\tilde c(\delta_0/2)\mathrm{e}^{-\delta_0w}$. Fix a number $b_*\ge b_{\rm geo}$ such that for all $b\ge b_*$: $b\ge2$, $\omega\le\gamma_0/2$, $\mathrm{e}^{2m_*}\varepsilon_b\le\frac12$, and (MC) holds with $v=v_b$ and $\beta=b$. Such a number exists, because each condition holds for all sufficiently large $b$, and it depends only on $(d,n,D,t,\kappa,\gamma_0,c_0,\delta_0,c_K)$. Define the error per box
\[
\begin{aligned}
E_{\rm box}:={}&2M\mathrm{e}^{-\kappa w/2}+2v_b\omega\Big[\tfrac{R^2}2+\tfrac n{\gamma_0}+\tfrac{(c_0\tilde c(\delta_0)R)^2}{2\gamma_0^2}\Big]+2(2M)^{t+1}+2\mathrm{e}^M\varepsilon_b+\sum_{k\le t}(8M)^k\mathrm{e}^{-\kappa w/2}\\
&+\sum_{k\le t}\Big[k16^k(\mathsf Av_bR^D)^k(2R)^{kD}\varepsilon_b^{1/2}+\tfrac{3^k}{k!}v_bK_k(\bar\mu)(2R)^{kD}(\mathsf AR^D)^k\mathrm{e}^{-\bar\mu w/4}\\
&\qquad+k8^k(kD)^{kD+1}(1+R+\gamma_0^{-1})^{kD}(\mathsf Av_b)^k(M_1R+c_{\rm D})\mathrm{e}^{-\mu_{\rm c}w/8}+\tfrac{2\cdot6^k}{k!}v_bK_k(\bar\mu/2)\mathsf A^k\mathrm{e}^{-\bar\mu w/4}\Big],
\end{aligned}
\]
and let $\Phi(\beta;J'):=\int\chi_\beta\mathrm{e}^{V_{J'}}d\mu$.

\begin{lemma}[stage]\label{app:G1:R1-stage}
Let $b\ge b_*$, $\mathsf Ab^{\rho_1}\le1$, $\beta\in[b,\gamma_K^{-d}b]$ and $\beta':=\beta/\gamma_K$. Then
\[
\Phi(\beta;J')\le\Phi(\beta';J'')\exp\Big\{\sum_{k\le t}\frac{\kappa_\mu(V_{J'};k)-\kappa_\mu(V_{J''};k)}{k!}+N'E_{\rm box}\Big\}.
\]
\end{lemma}
\begin{proof}
On $\supp\chi_{\beta'}$ we have $\abs{z_x}\le c_K\beta'\le R$.

\emph{Step 1.} On $\supp\chi_\beta$ we have $\mathrm{e}^{V_{J'}}\le\mathrm{e}^{N'M\mathrm{e}^{-\kappa w}}\mathrm{e}^{V_\Gamma+\sum\Psi_\square}$, and $\chi_\beta\le\chi^\Gamma_\beta\chi^\Delta_{\beta'}$ because $\beta K\subset\beta'K$. Disintegrate $\mu=\mu_\Gamma\otimes\mu(\cdot|z_\Gamma)$. For $z_\Gamma\in\supp\chi^\Gamma_\beta$, the inner integrand $\prod_\square\chi^{\Delta_\square}_{\beta'}\mathrm{e}^{\Psi_\square}$ is $\ge0$ and factorises over the boxes. Lemma~\ref{app:G1:R1-dec} with $w_s=2w$ and $\abs\Delta\le N'v_b$ gives
\[
\int\mu(dz_\Delta|z_\Gamma)\prod_\square\chi^{\Delta_\square}_{\beta'}\mathrm{e}^{\Psi_\square}\le\mathrm{e}^{\omega_{\rm tot}}\prod_\square I_\square,\qquad I_\square:=\int\chi^{\Delta_\square}_{\beta'}\mathrm{e}^{\Psi_\square}d\nu_\square,
\]
with $\nu_\square=N(h_\square,G_{\Delta_\square})$, $h_\square=-G_{\Delta_\square}P_{\Delta_\square\Gamma}z_\Gamma$ and $\omega_{\rm tot}\le N'v_b\omega[\dots]$ as in $E_{\rm box}$.

\emph{Step 2: one box, $z_\Gamma$ fixed.} Apply Lemma~\ref{app:G1:R1-cov}(c) with $X=\Delta_\square$, $Y=\Lambda\setminus\Delta_\square$, and $z_Y$ equal to $z_\Gamma$ on $\Gamma$ and to $0$ on the other boxes. This gives $\abs{h_\square}\le M_1c_K\beta=M_1c_K\gamma_K\beta'\le\beta'/2$. Part (b) gives $G_{\Delta_\square}\le\gamma_0^{-1}$. Put $\chi:=\chi^{\Delta_\square}_{\beta'}$. Since $B(1)\subset K$, we have $1-\chi\le\mathbf 1\{\exists x\in\Delta_\square:\abs{z_x}>\beta'\}$. Since $\varepsilon(\cdot,\cdot)$ increases in its first argument and decreases in its second, Lemma~\ref{app:G1:R1-bdd}(c) gives $\nu_\square(1-\chi)\le\varepsilon_b$. Since $\chi\mathrm{e}^\Psi\le\mathrm{e}^{\Psi\chi}$, Lemma~\ref{app:G1:R1-bdd}(a) gives $\log I_\square\le\sum_k\kappa_{\nu_\square}(\Psi\chi;k)/k!+(2M)^{t+1}$.
Expand multilinearly in the letters $\Psi_1',\Psi_1'',\Psi_2,\Psi_3$; there are at most $4^k$ words.
(i) Words containing $\Psi_3$ contribute at most $\sum_k(8M)^k\mathrm{e}^{-\kappa w/2}$, by Lemma~\ref{app:G1:R1-bdd}(a).
(ii) Words without $\Psi_3$, with both $\Psi_1''$ and $\Psi_2$. View $\Psi_2$ as a polynomial in $z_{\Delta_\square}$ with $z_\Gamma$-dependent coefficients. It has no constant term, its supports lie in $\Gamma_2(\square)$, and $\norm{\Psi_2}_{\bar\mu}\le\mathsf AR^D$, since supports only shrink. First remove $\chi$ by Lemma~\ref{app:G1:R1-chi}(a), with $v=v_b$; (MC) at $(v_b,b)$ implies (MC) at $(\abs{\Delta_\square},\beta')$, and $2\beta'\le2R$. Then apply Lemma~\ref{app:G1:R1-conn}(b) with $S_1=\mathrm{Dp}(\square)$, $S_2=\square\cap\{\lambda<2w\}$, $\theta=\mu_{\rm c}$, $\theta'=\bar\mu$, $g'=g$ and mean bound $R$.
(iii) Words in the letters $\Psi_1',\Psi_1''$ only, with at least one $\Psi_1''$. Remove $\chi$ in the same way. Then apply Lemma~\ref{app:G1:R1-chi}(b) on $S=\square\cap\{\lambda\ge1.5w\}$, comparing the marginals of $\nu_\square$ and of $\mu$. Here $\eta_1\le M_1R\mathrm{e}^{-\mu_{\rm c}w/4}$ by Lemma~\ref{app:G1:R1-cov}(c), and $\eta_2\le c_{\rm D}\mathrm{e}^{-\mu_{\rm c}w/8}$ by part (d). Under $\mu$ these words sum exactly to $k!$ times the $k$-th term of
\[
c_\square:=\sum_{k\le t}\tfrac1{k!}\big[\kappa_\mu(\Psi_1'+\Psi_1'';k)-\kappa_\mu(\Psi_1';k)\big],
\]
which does not depend on $z_\Gamma$.
(iv) The words in $\Psi_1',\Psi_2$ only sum to $\kappa_{\nu_\square}(W_\square\chi;k)$. By Lemma~\ref{app:G1:R1-bdd}(a),(b),
\[
\sum_k\kappa_{\nu_\square}(W_\square\chi;k)/k!\le\log\int\chi\mathrm{e}^{W_\square}d\nu_\square+2\mathrm{e}^M\varepsilon_b+(2M)^{t+1}.
\]
Hence $I_\square\le\mathrm{e}^{c_\square+e_\square}\int\chi\mathrm{e}^{W_\square}d\nu_\square$, where $e_\square$ is the sum of the terms of $E_{\rm box}$ from $2(2M)^{t+1}$ up to the $(M_1R+c_{\rm D})$-term.

\emph{Step 3.} We have $\prod_\square\int\chi\mathrm{e}^{W_\square}d\nu_\square=\int\prod_\square\chi^{\Delta_\square}_{\beta'}\mathrm{e}^{W_\square}d\nu_{\rm loc}$, which by Lemma~\ref{app:G1:R1-dec} is at most $\mathrm{e}^{\omega_{\rm tot}}\int\mu(dz_\Delta|z_\Gamma)\chi^\Delta_{\beta'}\mathrm{e}^{\sum W_\square}$. Integrate against $\chi^\Gamma_\beta\mathrm{e}^{V_\Gamma}d\mu_\Gamma$, and use $\chi^\Gamma_\beta\le\chi^\Gamma_{\beta'}$ and $V_\Gamma+\sum W_\square=V_{J''}-R_1$. This gives at most $\mathrm{e}^{N'M\mathrm{e}^{-\kappa w/2}}\Phi(\beta';J'')$ times the factors collected so far.

\emph{Step 4: bookkeeping.} We show that
\[
\Big|\sum_{\mathcal B'}c_\square-\sum_k\tfrac1{k!}[\kappa_\mu(V_{J'};k)-\kappa_\mu(V_{J''};k)]\Big|\le N'\sum_k\tfrac{2\cdot6^k}{k!}v_bK_k(\bar\mu/2)\mathsf A^k\mathrm{e}^{-\bar\mu w/4}.
\]
Write $V_{J'}=V_{J''}+X+R_2$ with $X:=\sum_\square X_\square$ and $X_\square:=\Psi_1''(\square)$. A support of $R_2$ meets some $\mathrm{Dp}(\square)$ and is not contained in $\square\cap\{\lambda\ge1.5w\}$ (otherwise it is in $\Psi_1''$). Hence $\ell\ge w$, and $\norm{R_2}_{\bar\mu/2}\le\mathsf A\mathrm{e}^{-\bar\mu w/2}$. Expand multilinearly. The measure $\mu$ is centred and $\norm{C(x,y)}\le g_0\mathrm{e}^{-\mu_{\rm c}\abs{x-y}}$ (Lemma~\ref{app:G1:R1-cov}(a)); so Lemma~\ref{app:G1:R1-conn}(b) applies with $\beta=0$, $\theta=\mu_{\rm c}$ and $g'=g$.
Words containing $R_2$: anchor at $\bigcup\mathrm{Dp}(\square)$, which has at most $N'v_b$ sites, with $\theta'=\bar\mu/2$.
Other words containing an $X$: split the first $X$ as $\sum_\square X_\square$. Given $\square$, split the other $X$'s as $X_\square+X_{\ne\square}$, and split $V_{J''}=\Psi_1'(\square)+Y^{\rm far}+Y^{\rm long}$. Here $Y^{\rm long}$ consists of the supports meeting $\Gamma_4(\square)$ that are not in $\Psi_1'(\square)$; they have $\ell\ge\frac w2$. The supports of $Y^{\rm far}$ lie in $J''\setminus\Gamma_4(\square)$. Then:
\begin{itemize}
\item words with $X_{\ne\square}$ or $Y^{\rm far}$: Lemma~\ref{app:G1:R1-conn}(b) with $S_1=\mathrm{Dp}(\square)$, at distance $\ge\frac w2$ from $S_2$, and $\theta'=\bar\mu$;
\item words with $Y^{\rm long}$: the anchored bound with $\theta'=\bar\mu/2$ and $\norm{Y^{\rm long}}_{\bar\mu/2}\le\mathsf A\mathrm{e}^{-\bar\mu w/4}$;
\item words in $X_\square,\Psi_1'(\square)$ only: these sum exactly to $k!$ times the $k$-th term of $c_\square$.
\end{itemize}
There are at most $3^k+5^k\le2\cdot6^k$ words, and $K_k(\bar\mu)\le K_k(\bar\mu/2)$.
Steps 1--4 together prove the lemma.
\end{proof}

\begin{proof}[Proof of Theorem~\ref{app:K:L13-1}]
If no coordinate is cut, then $\lambda\equiv\infty$, $\Lambda$ is one cell with $\abs\Lambda\le v_b$, and a single stage suffices. Otherwise apply Lemma~\ref{app:G1:R1-stage} for $i=0,\dots,d$ with $J_0=J$, $J_{i+1}:=J_i''$ (with respect to $\mathcal P^i$) and $\beta_i=\gamma_K^{-i}b$; then $N_i\le\abs J$. By (G3), $J_{d+1}=\emptyset$, so $V_{J_{d+1}}=0$. The cumulant increments telescope to $\sum_k\kappa_\mu(V_J;k)/k!$, and $\Phi(\beta_{d+1};\emptyset)=\mu(\chi_{\beta_{d+1}})\le1$.
Next, $2(2M)^{t+1}=2(2m_*)^{t+1}(\mathsf Ab^{\rho_1})^{t+1}$. All other terms of $E_{\rm box}$ are nondecreasing in $\mathsf A$; let $Q(b)$ be their sum at $\mathsf A=b^{-\rho_1}$. Each of them is a polynomial in $b$ times $\mathrm{e}^{-\bar\mu w/8}$ or $\varepsilon_b^{1/2}$, since $\kappa/2,\delta_0,\bar\mu/4\ge\bar\mu/8$ and $\mu_{\rm c}\ge\bar\mu$. Since $w\ge b^{3/2}$, $S:=(d+1)\max\{2(2m_*)^{t+1},\sup_{b\ge b_*}\mathrm{e}^{\rho_3b^{3/2}}Q(b)\}<\infty$, and the claim follows.
\end{proof}

\begin{corollary}[Dirichlet to free cumulants]\label{app:K:L13-2}\label{app:G1:PH-DF}
Let $P$ on $\Lambda_{\rm big}\subset\mathbb L$ satisfy (Pa)--(Pb) (for instance $P\in\mathcal P(\gamma_0,\gamma_1,c_0,\delta_0)$), let $\Lambda\subset\Lambda_{\rm big}$, and let $V$ be as in the setting with $J\subset\Lambda$. Put $r:=\dist(J,\Lambda_{\rm big}\setminus\Lambda)$ and $\theta':=\min(\mu_{\rm c}/4,\kappa)$. Then for every $k\ge1$,
\[
\abs{\kappa_{N(0,(P_{\Lambda\Lambda})^{-1})}(V;k)-\kappa_{N(0,P^{-1})}(V;k)}\le C_k\abs J\,\mathsf A^k\mathrm{e}^{-\mu_{\rm c}r/4},\qquad C_k:=\tfrac{kD}2\,K_k(\theta';\max(g,c_{\rm D})).
\]
\end{corollary}
\begin{proof}
Both laws are centred, and only covariances between points of $J$ enter: $G_\Lambda:=(P_{\Lambda\Lambda})^{-1}$ for the first law and $C:=P^{-1}$ for the second. Expand both cumulants by Lemma~\ref{app:G1:R1-conn}(a). For a fixed connected pairing, telescope the at most $kD/2$ lines one at a time. Apply Lemma~\ref{app:G1:R1-cov} with $\Lambda_{\rm big}$ in place of $\Lambda$ and $X=\Lambda$. By part (d), the replaced line obeys $\norm{(C-G_\Lambda)(x,y)}\le c_{\rm D}\mathrm{e}^{-\mu_{\rm c}r/4}\mathrm{e}^{-(\mu_{\rm c}/4)\abs{x-y}}$ for $x,y\in J$. By part (a), the other lines obey $\le g_0\mathrm{e}^{-\mu_{\rm c}\abs{x-y}}$. So every line is bounded by $\max(g,c_{\rm D})\mathrm{e}^{-(\mu_{\rm c}/4)\abs{x-y}}$, with the extra factor $\mathrm{e}^{-\mu_{\rm c}r/4}$ on the replaced one. Finally apply Lemma~\ref{app:G1:R1-conn}(c) with $\theta=\mu_{\rm c}/4$, the given $\theta'$, $g'=\max(g,c_{\rm D})$, $S_1=J$ and $S_2=\mathbb L$; note that $\norm V_{\theta'}\le\mathsf A$ because $\theta'\le\kappa$.
\end{proof}

\begin{remark}[application]\label{app:G1:R1-appl}
Theorem~\ref{app:K:L13-1} is applied once, to the joint fluctuation integral of \cite[(2.115)]{Ba2}, in Lemma~\ref{app:G1:S05-SC}, step (C6), and in Lemma~\ref{app:G1:S05-PH}(b). The data are: $n=d+2$, $z_x=(A'(x),\varphi'(x))\in\R^d\times\R^2$, $K=B(1)\times B(1)$ (so $c_K=\sqrt2$), $\Lambda=\Lambda_5^{(k)}$ and $J=\Lambda_7^{(k)}$. The measure $\mu$ is the product Gaussian measure with covariance $w_kC^{(k)}_{\Lambda_5}\oplus C^{(k)}_{\Lambda_5}(\cdot)$, and $V=V^{(k)}$ is the real polynomial of \cite[(3.57)--(3.58), p.~622]{Ba1} at real backgrounds. The integrand is nonnegative and vanishes unless $\abs{A'(x)}\le b$ and $\abs{\varphi'(x)}\le b$ on $\Lambda_5^{(k)}$, with $b=c_bp(t_k)$ (Lemma~\ref{app:G1:S05-PH}(a) and \cite[Lemma~2.7]{Ba2}); so it is bounded by $\chi_b\mathrm{e}^V$. The bound \cite[(3.58)]{Ba1} gives $\mathsf A\le Ct_k^{\kappa_0}$ with $\kappa=\delta_0/2$, while $b$ grows like a power of $\log t_k^{-1}$. Hence $b\ge b_*$ and $\mathsf Ab^{\rho_1}\le1$ hold for $\eps_0$ small enough, uniformly in $k$, and $\mathrm{e}^{-\rho_3b^{3/2}}=O(t_k^N)$ for every $N$. Corollary~\ref{app:K:L13-2} is used with $\Lambda_{\rm big}=T_1^{(k)}$, $\Lambda=\Lambda_5^{(k)}$ and $J=\Lambda_7^{(k)}$, when the Dirichlet covariances in the cumulants are replaced by the free ones \cite[p.~582]{Ba2} (Lemma~\ref{app:G1:S05-SC}, step (C7), and Lemma~\ref{app:G1:S05-PH}(c)). If $r$ is at least a fixed multiple of $r(t_k)$, then $\mathrm{e}^{-\mu_{\rm c}r/4}=O(t_k^N)$ for every $N$.
\end{remark}

\paragraph{Collars: setting.}\label{app:G1:Q-setting}
Let $d=3$, $L=2$, $0\le k\le K-1$ and $t=t_k\le\eps_0$. Fix a scale $j<k$ and work in scale-$j$ units, with $\eta=L^{-j}$, on the domain $\Omega$ of the polynomial $\mathcal Q^{(k)}$ of \cite[(2.105)--(2.106), p.~579]{Ba2}. A real field enters through its density $\cAc(b)=e_jA(b)$, and fields are identified with their densities. Put $U_b=\mathrm{e}^{\mathrm{i}\eta\cAc(b)}$, and $d\lambda(b):=\eta^{-1}(\lambda(b_+)-\lambda(b_-))$ for functions $\lambda$ on the fine sites. Let $\mathcal F$ be the curvature of $\cAc$, so that the flux of a plaquette $p$ is $\eta^2\mathcal F(p)$. Put $(\nabla_A\phi)(b)=\eta^{-1}(U_b\phi(b_+)-\phi(b_-))$, with bonds inside $\Omega$. Let $Q_j(A)$ and $Q(A)$ be the averaging operators with transporters along contours of length $\le d$, resp.\ $\le dL$, ending at the unit-lattice point. Put $X=\nabla_A^*\nabla_A+m^2t_j^2+a_jQ_j^*Q_j$ with $0<a_j\le a$, $G=X^{-1}$, $\Delta^{(j)}=a_j-a_j^2Q_jGQ_j^*$, $\mathsf M(A)=aL^{-2}Q^*Q+\Delta^{(j)}$ and $f_y(A):=\tr_{\R^2}(\log\mathsf M(A))(y,y)$. Then $-\frac12\sum_yf_y$ is the field-dependent part of $\log Z^{(j)}$, and
\[
\mathcal Q^{(k)}(\Omega;\tilde B,\alpha)=\sum_{j<k}\sum_yq_{j,y}(\alpha),\qquad q_{j,y}(\alpha)=-\tfrac12\sum_{n=1}^{\nbar}\tfrac1{n!}D^nf_y(\tilde B_j)[\alpha_j,\dots,\alpha_j],
\]
where $\tilde B_j$ and $\alpha_j$ are the scale-$j$ forms of $\tilde B$ and $\alpha$. (The field argument of $\mathcal Q^{(k)}$ is written $\alpha$, because $a>0$ and $a_j$ denote the constants of the averaging terms above.) We assume $\nbar\ge6$, and (Q0), (Q1) for regular real $A$ in the sense of \cite[(2.23)]{Ba1}:
(Q0) $X-m^2t_j^2\ge\gamma_0$ \cite[(1.8)]{Ba4};
(Q1) $\gamma_0'\le\mathsf M(A)\le a_j+aL^{-2}$ \cite[(1.15)]{Ba4}.
By Lemma~\ref{app:G1:Q-W}(iii) below, these imply, for regular real $A$,
(Q2) $\abs{\mathsf M(A)(y,y')}\le c_M\mathrm{e}^{-\delta\abs{y-y'}}$, with the numbers $\delta$ and $c_M$ of that lemma, which depend only on $(d,L,a,\gamma_0)$.

\begin{lemma}[gauge invariance and evenness]\label{app:G1:Q-GC}
(a) $f_y(A+d\lambda)=f_y(A)$ for real $\lambda$. Hence $D^nf_y(A)[h_1+d\nu,h_2,\dots]=D^nf_y(A)[h_1,h_2,\dots]$, and the same holds in every argument.
(b) $f_y$ is real and $f_y(-A)=f_y(A)$. Hence $Df_y(B)[h]=\int_0^1D^2f_y(sB)[B,h]\,ds$.
\end{lemma}
\begin{proof}
(a) Every building block is conjugated by the sitewise phase $\mathrm{e}^{\mathrm{i}\lambda}$ (the contours end at the unit-lattice point). Hence $\log\mathsf M(A+d\lambda)=U_\lambda\log\mathsf M(A)U_\lambda^*$ with a diagonal unitary $U_\lambda$, and the trace of the $(y,y)$ block is unchanged. For the derivatives, differentiate $f_y(A+\sum_is_ih_i+d(s_1\nu))=f_y(A+\sum_is_ih_i)$ at $s=0$.
(b) Complex conjugation $\sigma$ satisfies $\sigma U_b(A)\sigma=U_b(-A)$, and the $A$-independent pieces are real; so $\mathsf M(-A)=\sigma\mathsf M(A)\sigma$, and $f_y$ is real and even. Evenness gives $Df_y(0)=0$, and $Df_y(B)[h]=Df_y(0)[h]+\int_0^1D^2f_y(sB)[B,h]ds$.
\end{proof}

\begin{lemma}[cut-off axial gauge]\label{app:G1:Q-AX}
Let $Q_y$ be the box of half-side $R_y$ about $y$, with $Q_y\subset\Omega$ not wrapping, and let $Q_y'$ be the concentric box of half-side $R_y/2$. For $x\in Q_y$ let $\lambda_y(x):=-\sum_{b\in\Gamma^{\rm ax}_{y\to x}}\eta\cAc(b)$ along the axis-ordered path, which stays in $Q_y$. Let $\chi:\Omega\to[0,1]$ satisfy $\chi=1$ on $Q_y'$, $\chi=0$ off $Q_y$, and $\abs{\chi(x)-\chi(x')}\le2\eta/R_y$ on bonds. Put $\cAc^g:=\cAc+d(\chi\lambda_y)$, where $\chi\lambda_y:=0$ off $Q_y$. If $\sup\abs{\cAc}\le P_1$, $\sup_{Q_y}\abs{\mathcal F}\le E$ and $\sqrt d\,R_yE\le P_1$, then
\[
\abs{\cAc^g(b)}\le\norm{b_--y}_1E\quad(b\subset Q_y'),\qquad
\abs{\cAc^g(b)}\le(2d+\sqrt d)P_1\le(2d+3)P_1\quad\text{otherwise}.
\]
\end{lemma}
\begin{proof}
For $b\subset Q_y$, $\eta(\cAc+d\lambda_y)(b)$ is minus the flux through the surface spanned by $\Gamma^{\rm ax}_{y\to b_-}\circ b\circ(\Gamma^{\rm ax}_{y\to b_+})^{-1}$. This surface lies in $Q_y$ and consists of at most $\norm{b_--y}_1/\eta$ plaquettes, each with flux at most $\eta^2E$. Hence $\abs{\cAc+d\lambda_y}(b)\le\norm{b_--y}_1E\le dR_yE\le\sqrt d\,P_1$. Since $\chi\equiv1$ on $Q_y'$, this gives the first bound. Also $\abs{\lambda_y(x)}\le\norm{x-y}_1P_1\le dR_yP_1$ on $Q_y$.
If $b\subset Q_y$, then
\[
\cAc^g(b)=(1-\chi(b_+))\cAc(b)+\chi(b_+)(\cAc+d\lambda_y)(b)+\lambda_y(b_-)\,\frac{\chi(b_+)-\chi(b_-)}\eta,
\]
so $\abs{\cAc^g(b)}\le\max(P_1,\sqrt dP_1)+dR_yP_1\cdot2/R_y$. If exactly one endpoint $x$ of $b$ lies in $Q_y$, then $\chi$ vanishes at the other endpoint. So $\chi(x)\le2\eta/R_y$ and $\abs{\cAc^g(b)}\le P_1+\chi(x)\abs{\lambda_y(x)}/\eta\le(2d+1)P_1$. If $b\cap Q_y=\emptyset$, then $\cAc^g(b)=\cAc(b)$. Finally, $\sqrt d\le3$.
\end{proof}

\begin{lemma}[weighted bounds]\label{app:G1:Q-W}
Let $\rho=\abs{\cdot-y}$ and $O^\beta=\mathrm{e}^{\beta\rho}O\mathrm{e}^{-\beta\rho}$; on bond functions, $\rho$ is evaluated at $b_-$. There is $\beta_0=\beta_0(d,a,\gamma_0)\in(0,1]$ such that the following hold for $\abs\beta\le\beta_0$, uniformly in $\eta,j,k,\eps,\Omega$, the boundary conditions and regular $A$:
(i) $\norm{G^\beta}\le2/\gamma_0$, and $(\nabla G)^\beta$, $(G\nabla^*)^\beta$, $(\nabla G\nabla^*)^\beta$ have norm $\le K_W(d,\gamma_0)$;
(ii) $\norm{Q_j^\beta},\norm{(Q_j^*)^\beta}\le\mathrm{e}^{\abs\beta\sqrt d}$ and $\norm{Q^\beta},\norm{(Q^*)^\beta}\le\mathrm{e}^{\abs\beta\sqrt dL}$.
(iii) for regular real $A$, (Q2) holds with $\delta:=\beta_0$ and $c_M:=a+aL^{-2}\mathrm{e}^{2\beta_0\sqrt dL}+2a^2\gamma_0^{-1}\mathrm{e}^{2\beta_0\sqrt d}$.
Moreover, there is $\beta_*=\beta_*(d,L,a,\gamma_0,\gamma_0')\in(0,\beta_0/2]$ such that the following holds for $\abs\beta\le\beta_*$ and $s\ge0$:
(iv) $\norm{((\mathsf M+s)^{-1})^\beta}\le(\gamma_0'/2+s)^{-1}$.
\end{lemma}
\begin{proof}
(i) Put $c(\beta):=d\beta^2\cosh\beta+a(\mathrm{e}^{\sqrt d\abs\beta}-1)$, and choose $\beta_0\le1$ with $c(\beta_0)\le\gamma_0/4$. For a bond $b$ let $\tau_b=\beta(\rho(b_+)-\rho(b_-))$, so $\abs{\tau_b}\le\abs\beta\eta$. Write $u=U_bw(b_+)$ and $v=w(b_-)$. The real part of the conjugated bond form is $\eta^{d-2}\big[\abs{u-v}^2-2(\cosh\tau_b-1)\rRe\ip uv\big]\ge\eta^{d-2}\big[\abs{u-v}^2-(\cosh\tau_b-1)(\abs u^2+\abs v^2)\big]$. Moreover $\norm{(Q_j^*Q_j)^\beta-Q_j^*Q_j}\le\mathrm{e}^{\sqrt d\abs\beta}-1$ by Schur's test, and $a_j\le a$. Hence $\rRe\ip w{X^\beta w}\ge\ip w{Xw}-c(\beta)\norm w^2$. Since $X\ge\gamma_0$ by (Q0) and $X\ge\nabla^*\nabla$, this gives
\[
\rRe\ip w{X^\beta w}\ge\tfrac34\gamma_0\norm w^2,\qquad \rRe\ip w{X^\beta w}\ge\tfrac12\norm{\nabla w}^2+\tfrac14\gamma_0\norm w^2 .
\]
With $X^\beta w=F$, the first inequality gives $\norm w\le\frac4{3\gamma_0}\norm F$, so $\norm{G^\beta}\le2/\gamma_0$. The second gives $\norm{\nabla w}^2\le2\norm w\norm F$, so $\norm{\nabla w}\le(8/(3\gamma_0))^{1/2}\norm F$. Next, $\nabla^\beta=\nabla+K_\beta$ with $(K_\beta\phi)(b)=\eta^{-1}(\mathrm{e}^{-\tau_b}-1)U_b\phi(b_+)$, and $\norm{K_\beta}\le\sqrt d\,\mathrm{e}\abs\beta$. Hence $(\nabla G)^\beta=\nabla^\beta G^\beta$ is bounded, and $(G\nabla^*)^\beta=((\nabla G)^{-\beta})^*$. For $(\nabla G\nabla^*)^\beta g=\nabla^\beta w$, let $w:=G^\beta(\nabla^*)^\beta g$. Then
\[
\tfrac12\norm{\nabla w}^2+\tfrac14\gamma_0\norm w^2\le\rRe\ip{\nabla^{-\beta}w}g\le(\norm{\nabla w}+\sqrt d\,\mathrm{e}\abs\beta\norm w)\norm g,
\]
so $\norm{\nabla w}+\norm w\le C(d,\gamma_0)\norm g$.
(ii) $\abs{Q_j(y',x)}\le\eta^d\mathbf 1_{x\in B(y')}$, with $\abs{\rho(x)-\rho(y')}\le\sqrt d$ on the support, and $Q_jQ_j^*=I$; apply Schur's test. The same argument works for $Q$, with $\sqrt dL$.
(iii) For an operator $O$ on the unit lattice, $\rho(y)=0$ gives $O^{-\beta}(y,y')=\mathrm{e}^{\beta\abs{y-y'}}O(y,y')$, so $\norm{O(y,y')}\le\norm{O^{-\beta_0}}\mathrm{e}^{-\beta_0\abs{y-y'}}$. Apply this to the three terms of $\mathsf M=aL^{-2}Q^*Q+a_j-a_j^2Q_jGQ_j^*$, using (i), (ii) and $a_j\le a$: $\norm{(Q^*Q)^{-\beta_0}}\le\mathrm{e}^{2\beta_0\sqrt dL}$ and $\norm{(Q_jGQ_j^*)^{-\beta_0}}\le2\gamma_0^{-1}\mathrm{e}^{2\beta_0\sqrt d}$.
(iv) Choose $\beta_*\le\beta_0/2$ with $c_M\sum_{z\in\Z^d}\mathrm{e}^{-\beta_0\abs z}(\mathrm{e}^{\beta_*\abs z}-1)\le\gamma_0'/2$. By (iii) and Schur's test, $\norm{\mathsf M^\beta-\mathsf M}\le\gamma_0'/2$ for $\abs\beta\le\beta_*$. Hence $\rRe(\mathsf M^\beta+s)\ge\gamma_0'/2+s$ by (Q1).
\end{proof}

\begin{lemma}[sensitivity]\label{app:G1:Q-S}
Let $\delta_1:=\beta_*/(\nbar+1)$ and $\norm h_y:=\sup_b\mathrm{e}^{-\delta_1\abs{b_--y}}\abs{h(b)}$. For real regular $A$, all $y$ and $1\le n\le\nbar+1$,
\[
\abs{D^nf_y(A)[h_1,\dots,h_n]}\le C_n\prod_i\norm{h_i}_y,
\]
with $C_n$ depending only on $(d,L,a,\gamma_0,\gamma_0',c_M,\delta,\nbar)$. No condition on $\dist(y,\partial\Omega)$ is needed.
\end{lemma}
\begin{proof}
Write $\log\mathsf M=\int_0^\infty[(1+s)^{-1}-(\mathsf M+s)^{-1}]\,ds$. Then $D^n(\mathsf M+s)^{-1}$ is a sum of at most $2^{n-1}n!$ chains $\pm R\,D^{\pi_1}\mathsf M\,R\cdots D^{\pi_r}\mathsf M\,R$, with $R=(\mathsf M+s)^{-1}$ and $r+1\ge2$ factors $R$. Expand $D^\pi\mathsf M$, $D^\pi G$ and $D^\pi X$ by the Leibniz rule. We have
\[
D^\pi\nabla_A=\mathrm{i}^{\abs\pi}\eta^{\abs\pi-1}S_{\prod_{i\in\pi}h_i},\qquad (S_g\phi)(b)=g(b)U_b\phi(b_+),\qquad \norm{S_g}\le\sqrt d\sup\abs g .
\]
The kernel of $D^\pi Q_j$ is $\prod_{i\in\pi}(\mathrm{i}\Phi_{h_i}(\Gamma_{y',x}))\,Q_j(y',x)$, with $\Phi_h(\Gamma)=\sum_{b\in\Gamma}\eta h(b)$. Write $h_i=\mathrm{e}^{\delta_1\rho}\hat h_i$ with $\sup\abs{\hat h_i}=\norm{h_i}_y$. Then $\abs{\Phi_{h_i}(\Gamma)}\le\mathrm{e}^{\delta_1\rho(y')}\,d\,\mathrm{e}^{\delta_1\sqrt d}\norm{h_i}_y$, and the same holds for $Q$ with $dL$ and $\sqrt dL$. Each $G$ carries at most one $\nabla$ on each side. So every chain is a product of operators from $\{R,G,\nabla G,G\nabla^*,\nabla G\nabla^*,Q_j^{(*)},Q^{(*)},S_{\hat g}\}$ and exactly $n$ multiplications by $\mathrm{e}^{\delta_1\rho}$, up to a factor $\prod_i(\sqrt d+d\mathrm{e}^{\delta_1\sqrt dL}L)\norm{h_i}_y$. Conjugate from the left with weights $0\le\beta_1\le\dots\le n\delta_1\le\beta_*$, raising the weight by $\delta_1$ at each multiplication by $\mathrm{e}^{\delta_1\rho}$. Since $\rho(y)=0$, we have $\mathrm{e}^{n\delta_1\rho}\delta_y=\delta_y$. Every conjugated factor is bounded by Lemma~\ref{app:G1:Q-W}, and $\int_0^\infty(\gamma_0'/2+s)^{-r-1}ds\le(2/\gamma_0')^r/r$.
\end{proof}

\paragraph{The fields.}\label{app:G1:Q-fields}
Let $\vartheta(t):=\abs e\,t^{1/2}p(t)$ and $E_j:=2C_B\vartheta(t)L^{-2(k-j)}$, where $C_B$ is the constant of Lemma~\ref{app:K:L13-3}. Let $P_1(t)$ be the bound on the fields at scale $j$ that Lemma~\ref{app:K:L13-3}(a) gives after the rescaling of Remark~\ref{app:G1:bg-rem}(ii). We use that statement (under (IH$_k$)) in the following form, on $\Omega$:
(Fd1) $\sup\abs{\mathcal F(\tilde B_j)},\ \sup\abs{\mathcal F(\alpha_j)}\le E_j$;
(Fd2) $\sup\abs{\tilde B_j},\ \sup\abs{\alpha_j}\le P_1(t)$, where $P_1$ is bounded by a polynomial in $t^{-1/2}\ell(t)$, $\ell(t)=1+\log t^{-1}$;
(Fd3) $s\tilde B_j$ is regular for $s\in[0,1]$.
For $t\le\eps_0$, conditions (E.ksp1), (E.ksp2) give (Lemma~\ref{app:G1:RW-KSP}):
(Sm1) $2c_{AX}C_B\vartheta(t)\le\frac12$, where $c_{AX}:=\sqrt d/(\mathrm{e}\delta_1)$;
(Sm2) $(2d+3)\mathrm{e}^{-\delta_1r(t)/4}P_1(t)\le\frac12$ and $\sqrt d\,C_Br(t)\vartheta(t)/2\le P_1(t)$, where $r(t)=R\ell(t)^2$.
\emph{Radius.} Assume $\Omega^c\ne\emptyset$, and put $R_y:=\min\{\dist_j(y,\Omega^c)/2,\ r(t)L^{k-j}/2\}$. Then $\sqrt dR_y<\dist_j(y,\Omega^c)$, so $Q_y\subset\Omega$. Moreover $R_y$ is at most $\sqrt d/4$ times the side of the scale-$j$ torus, so $Q_y$ does not wrap. Finally $\sqrt dR_yE_j\le\sqrt d\,C_Br(t)\vartheta(t)/2\le P_1(t)$ by (Sm2).

\begin{theorem}[site bound]\label{app:G1:Q-site}
Let $Q_y\subset\Omega$ be a non-wrapping box of half-side $R_y$ with $\sqrt dR_yE_j\le P_1(t)$, and assume (Fd1)--(Fd3). Let $T_{j,y}:=(2d+3)\mathrm{e}^{-\delta_1R_y/2}P_1(t)$, $\tau:=c_{AX}E_j+T_{j,y}$ and $C_Q:=\sum_{n\le\nbar+1}C_n/n!$. If $\tau\le1$, then $\abs{q_{j,y}(\alpha)}\le C_Q\tau^2$. The constants $\delta_1$ and $C_Q$ depend only on $(d,L,a,\gamma_0,\gamma_0',c_M,\delta,\nbar)$.
\end{theorem}
\begin{proof}
Lemma~\ref{app:G1:Q-AX} applies to $\tilde B_j$ and to $\alpha_j$, with $E=E_j$ and $P_1=P_1(t)$; denote the gauge-fixed fields by $\tilde B^g$ and $\alpha^g$. For $b\subset Q_y'$ we have $\mathrm{e}^{-\delta_1\abs{b_--y}}\sqrt d\abs{b_--y}E_j\le c_{AX}E_j$, since $\sup_{u\ge0}u\mathrm{e}^{-\delta_1u}=1/(\mathrm{e}\delta_1)$. If $b\not\subset Q_y'$, then $\abs{b_--y}\ge R_y/2-\eta$, so the weight is at most $\mathrm{e}^{\delta_1}\mathrm{e}^{-\delta_1R_y/2}$. Since $\delta_1=\beta_*/(\nbar+1)\le1/7$, we have $(2d+\sqrt d)\mathrm{e}^{\delta_1}\le2d+3$. Hence $\norm{\tilde B^g}_y,\norm{\alpha^g}_y\le\tau$.
By Lemma~\ref{app:G1:Q-GC}(a), applied in each argument at the regular base point $\tilde B_j$, $D^nf_y(\tilde B_j)[\alpha_j,\dots,\alpha_j]=D^nf_y(\tilde B_j)[\alpha^g,\dots,\alpha^g]$. For $n\ge2$ this is at most $C_n\tau^n\le C_n\tau^2$ by Lemma~\ref{app:G1:Q-S}. For $n=1$, Lemma~\ref{app:G1:Q-GC}(b) and then (a) in both arguments give $Df_y(\tilde B_j)[\alpha_j]=\int_0^1D^2f_y(s\tilde B_j)[\tilde B^g,\alpha^g]\,ds$. This is at most $C_2\tau^2$, because $s\tilde B_j$ is regular by (Fd3). Summing over $n$ with the factor $\frac12$ gives the claim, since $\frac12\big(C_2+\sum_{2\le n\le\nbar}C_n/n!\big)\le C_Q$.
\end{proof}

\begin{corollary}[sum over the sites]\label{app:G1:Q-sum}
Assume (Fd1)--(Fd3), (Sm1), (Sm2), $\Omega^c\ne\emptyset$ and the radius $R_y$ above. If $\dist_k(Y,\Omega^c)\ge r(t)/2$, then for every $N$
\[
\sum_{j<k}\sum_{y\in B^{k-j}(Y)}\abs{q_{j,y}(\alpha)}\le\Big[\frac{8C_Qc_{AX}^2C_B^2\vartheta(t)^2}{L-1}+C_Nt^N\Big]\abs Y ,
\]
where $C_N$ depends only on $N$, $R$, $\delta_1$, $C_Q$ and the polynomial bounding $P_1$.
\end{corollary}
\begin{proof}
Let $m=k-j\ge1$. Each scale-$k$ block holds $L^{3m}$ sites $y$, each with $\dist_j(y,\Omega^c)\ge r(t)L^m/2$, hence with $R_y\ge r(t)L^m/4$. Then $c_{AX}E_j\le2c_{AX}C_B\vartheta\le\frac12$ by (Sm1), and $T_{j,y}\le(2d+3)\mathrm{e}^{-\delta_1r(t)L^m/8}P_1(t)\le\frac12$ by (Sm2), since $L=2$. So $\tau\le1$, and Theorem~\ref{app:G1:Q-site} applies. Use $\tau^2\le2c_{AX}^2E_j^2+2T_{j,y}^2$. The first part contributes at most $C_Q\sum_{m\ge1}L^{3m}\,8c_{AX}^2C_B^2\vartheta^2L^{-4m}\abs Y=8C_Qc_{AX}^2C_B^2\vartheta^2\abs Y/(L-1)$. The second contributes at most $2C_Q(2d+3)^2P_1(t)^2\sum_{m\ge1}L^{3m}\mathrm{e}^{-\delta_1r(t)L^m/4}\abs Y\le C_Nt^N\abs Y$. This holds because $r(t)=R\ell(t)^2$ and $P_1$ is polynomial in $t^{-1/2}\ell(t)$.
\end{proof}

\subsection{The counterterm bound and the restriction tails}\label{app:G1:ct-tails}

\paragraph{Notation for the counterterm bound.} Let $X\in\{\eps\Z^3,T\}$, where $T$ is a lattice torus of side $\ell\ge\ell_0$ with $\ell/\eps\in2\N$, and let $\eps_{\max}$ be a fixed upper bound for $\eps$. Write $\int_X$ for the Brillouin-zone integral $\int_{(-\pi/\eps,\pi/\eps]^3}\frac{d^3p}{(2\pi)^3}$ on $\eps\Z^3$, and for $\ell^{-3}\sum_{p\in T^*}$ on $T$, where $T^*=\frac{2\pi}{\ell}\Z^3\cap(-\pi/\eps,\pi/\eps]^3$. Put $\hat x_\nu:=\frac2\eps\sin\frac{\eps x_\nu}2$ componentwise, $\hat p^2=\sum_\nu\hat p_\nu^2$, $\hat p=(\hat p^2)^{1/2}$, $\hat G=(\hat p^2+M^2)^{-1}$ and $\hat D=(\hat p^2+\mu_*^2)^{-1}$, and let $G_X,D_X$ be their kernels. Define
\[
D_X:=D_X(0)=\int_X\hat D,\quad G_X(0)=\int_X\hat G,\quad \Psi_X:=\int_X\hat D\hat G,\quad \delta_X:=\int_X\hat D^2,\quad \tau_X:=\int_X\hat G^2,
\]
and $J_X:=M^2\Psi_X$, $Q_X:=\sum_z\eps^3G_X(z)^3$. Write $\delta m^2_{K,X}(M)$ for the counterterm of Definition~\ref{def:counterterm}, computed on $X$. By Proposition~\ref{app:G1:P1-ct} it is one number on each $X$, so it does not depend on the position.

\begin{lemma}[representation of the counterterm]\label{app:G1:ct-rep}
Let $\Sigma_X(p):=e^2D_X\sum_\mu\cos(\eps p_\mu)-e^2\int_{X,p'}\hat D(p-p')\,\widehat{(p+p')}^2\hat G(p')$ and $\tilde\Pi_X(p):=\Sigma_X(p)-\Sigma_X(0)$. Then
\[
\delta m^2_{K,X}(M)=-16\lambda G_X(0)-e^2(2D_X+J_X)+128\lambda^2Q_X+16\lambda\int_X\hat G^2\tilde\Pi_X+\delta m^2_{e^4,X},
\]
\[
\delta m^2_{e^4,X}=\tfrac34\eps^2e^4D_X^2+(A)+(C)+(D1)+(D2)+(D3),
\]
where $s_\mu(q):=\frac2\eps\sin\frac{\eps q_\mu}2$ and
\begin{align*}
(A)&=e^2M^2\int_X\hat D\hat G^2\tilde\Pi_X,\qquad (C)=-\tfrac12\eps^2e^4D_X(D_X-J_X),\\
(D1)&=\tfrac{e^4}2\sum_\mu\sum_z\eps^3\big[D_X(z)^2+D_X(z+\eps e_\mu)^2\big]\big[G_X(z)+G_X(z+\eps e_\mu)\big],\\
(D2)&=-e^4\sum_\mu\int_{X,k}s_\mu(k)\hat G(k)\hat D(k)S_\mu(k),\qquad S_\mu(k):=\int_{X,P}\big[s_\mu(k)+s_\mu(2P+k)\big]\hat G(P)\hat D(P-k),\\
(D3)&=e^2\int_{X,k}\hat D(k)^2\sum_\mu\Pi_{\mu\mu}(k),\qquad \Pi_{\mu\mu}(k)/e^2=2G_X(\eps e_\mu)-\int_{X,q}\widehat{(2q+k)}_\mu^{\,2}\hat G(q)\hat G(q+k).
\end{align*}
\end{lemma}
\begin{proof}
The coefficients of $\lambda$, $e^2$, $\lambda^2$ and $e^2\lambda$, and the formula for $\Sigma_X$, are Lemma~\ref{app:G1:PD-low}(i)--(iv), with $C_X(M)=G_X(0)$ and $F_X(M)=e^{-2}\int_X\hat G^2\tilde\Pi_X$. The coefficient of $e^4$ is given by Corollary~\ref{app:G1:PD-e4}, which is derived from the Ward form of the zero-momentum condition at $\lambda=0$ (Proposition~\ref{app:G1:PD-ward}). This coefficient does not involve $\lambda$, so it is the coefficient of $e^4$ for every $\lambda$. The terms $\frac34\eps^2e^4D_X^2$, $(A)$, $(C)$, $(D1)$, $(D2)$, $(D3)$ are $e^4$ times the terms of Corollary~\ref{app:G1:PD-e4}; in $(D2)$ the $P$-integral is written as $S_\mu(k)$.
\end{proof}

\begin{lemma}[one-loop sums]\label{app:G1:ct-oneloop}
Let $w:=\int_{[-\pi,\pi]^3}\frac{d^3q}{(2\pi)^3}\big(\sum_\nu4\sin^2\frac{q_\nu}2\big)^{-1}=0.25273\ldots$, and let $\bar E$ be as in Lemma~\ref{app:G1:PD-img}(a). For a mass $a>0$ put $\bar\omega_a:=w+\eps_{\max}\bar E(a,\ell_0)$ on tori and $\bar\omega_a:=w$ on $\eps\Z^3$. Then $0\le G_T(0)-G_{\eps\Z^3}(0)\le\bar E(M,\ell)$, and for $\eps\le\eps_{\max}$
\[
\eps D_X\le\bar\omega_{\mu_*},\qquad \eps G_X(0)\le\bar\omega_M .
\]
Let $\bar\psi(M)$, $\bar\delta$ and $\bar\tau(M)$ be the right side of the inequality of Lemma~\ref{app:G1:PD-riem} at $\ell=\ell_0$ for $F(u)=(u+M^2)^{-1}(u+\mu_*^2)^{-1}$, $F(u)=(u+\mu_*^2)^{-2}$ and $F(u)=(u+M^2)^{-2}$, respectively. These numbers are finite, independent of $\eps$ and non-increasing in $\ell_0$. Moreover $\Psi_X\le\bar\psi(M)$, $\delta_X\le\bar\delta$ and $\tau_X\le\bar\tau(M)$ for every $\eps$, on $\eps\Z^3$ and on every torus with $\ell\ge\ell_0$, $\ell/\eps\in2\N$.
\end{lemma}
\begin{proof}
By Lemma~\ref{app:G1:PD-pos}(b), $G_T(0)-G_{\eps\Z^3}(0)=E_T(M)\ge0$. By Lemma~\ref{app:G1:PD-img}(a), $E_T(M)\le\bar E(M,\ell)\le\bar E(M,\ell_0)$; the same holds at the mass $\mu_*$ for $D$. On $\eps\Z^3$, scaling gives $\eps D_{\eps\Z^3}\le\eps\int\hat p^{-2}=w$ and likewise $\eps G_{\eps\Z^3}(0)\le w$. Adding $\eps E_T\le\eps_{\max}\bar E(\cdot,\ell_0)$ gives the torus bounds.

Each $F$ above is non-negative and non-increasing on $[0,\infty)$, and $F(u)=O(u^{-2})$. Hence the four terms of Lemma~\ref{app:G1:PD-riem} are finite, and each of them is non-increasing in $\ell$. This gives the torus bounds and the monotonicity in $\ell_0$. On $\eps\Z^3$ the first term alone is a bound, by Lemma~\ref{app:G1:PD-dos}.
\end{proof}

\begin{lemma}[elementary facts]\label{app:G1:ct-elem}
\begin{itemize}
\item[(a)] For $s,t>0$ the kernels of $\hat G^s$ and $\hat D^t$ on $X$ are pointwise non-negative, and $0\le G_X(z)\le G_X(0)$.
\item[(b)] For kernels $K,K'$ that are pointwise non-negative, with real even Fourier transforms, and for every $q$: $\abs{\int_{X,P}\hat K(P)\hat K'(P-q)}\le\int_X\hat K\hat K'$. In particular $\int_{X,P}\hat G(P)\hat D(P-k)\le\Psi_X$ and $\int_{X,P}\hat D(P-k)\hat D(P+k)\le\delta_X$.
\item[(c)] $\abs{\widehat{(a+b)}_\nu}\le\abs{\hat a_\nu}+\abs{\hat b_\nu}$; hence $\widehat{(a+b)}^2\le2\hat a^2+2\hat b^2$.
\item[(d)] $\widehat{(P+k)}^2-\widehat{(P-k)}^2=\sum_\nu\frac4{\eps^2}\sin(\eps P_\nu)\sin(\eps k_\nu)$, and $\abs{\frac2\eps\sin(\eps x)}\le2\abs{\hat x}$. Hence $\abs{\hat D(P-k)-\hat D(P+k)}\le4\hat P\hat k\,\hat D(P-k)\hat D(P+k)$.
\end{itemize}
\end{lemma}
\begin{proof}
(a) We have $(-\Delta_\eps+M^2)^{-s}=\Gamma(s)^{-1}\int_0^\infty\tau^{s-1}\mathrm{e}^{-\tau M^2}\mathrm{e}^{\tau\Delta_\eps}\,d\tau$. The off-diagonal entries of $\Delta_\eps$ are non-negative, so $\mathrm{e}^{\tau\Delta_\eps}$ has non-negative entries; this holds on $\eps\Z^3$ and on $T$. The same argument applies to $\hat D^t$. Since $\hat G>0$, $G_X(z)=\int_X\hat G(p)\mathrm{e}^{ipz}\le\int_X\hat G=G_X(0)$.

(b) By Parseval, the left side equals $\abs{\sum_x\eps^3K(x)K'(x)\mathrm{e}^{-iqx}}\le\sum_x\eps^3K(x)K'(x)=\int_X\hat K\hat K'$. For the second instance substitute $P\mapsto P+k$, which is a symmetry of the momentum set.

(c) This follows from $\abs{\sin(x+y)}\le\abs{\sin x}+\abs{\sin y}$.

(d) The identity is $\sin^2\alpha-\sin^2\beta=\sin(\alpha+\beta)\sin(\alpha-\beta)$, and the inequality is $\abs{\frac2\eps\sin\eps x}=\abs{\frac4\eps\sin\frac{\eps x}2\cos\frac{\eps x}2}$. Together they give $\abs{\widehat{(P+k)}^2-\widehat{(P-k)}^2}\le4\sum_\nu\abs{\hat P_\nu}\abs{\hat k_\nu}\le4\hat P\hat k$, by Cauchy--Schwarz. Finally $\hat D(P-k)-\hat D(P+k)=\big(\widehat{(P+k)}^2-\widehat{(P-k)}^2\big)\hat D(P-k)\hat D(P+k)$.
\end{proof}

\begin{lemma}[one-loop self-energies and the vertex]\label{app:G1:ct-se}
For all $p,k$:
\begin{itemize}
\item[(i)] $\abs{\Sigma_X(p)}\le e^2[5D_X+2\hat p^2\Psi_X]$, hence $\abs{\tilde\Pi_X(p)}\le e^2[10D_X+2\hat p^2\Psi_X]$;
\item[(ii)] $\abs{\sum_\mu\Pi_{\mu\mu}(k)}\le10e^2G_X(0)$;
\item[(iii)] $\sum_\mu\abs{s_\mu(k)}\abs{S_\mu(k)}\le\hat k^2\,(2\Psi_X+4\delta_X)$.
\end{itemize}
\end{lemma}
\begin{proof}
(i) The first term of $\Sigma_X$ has modulus at most $3e^2D_X$. In the second term, Lemma~\ref{app:G1:ct-elem}(c) gives $\widehat{(p+p')}^2\le2\hat p^2+2\hat p'^2$, and $\hat p'^2\hat G(p')\le1$. Using also $\int_{X,p'}\hat D(p-p')=D_X$ and Lemma~\ref{app:G1:ct-elem}(b), the second term is at most $e^2[2\hat p^2\Psi_X+2D_X]$. The bound on $\tilde\Pi_X$ follows from $\abs{\tilde\Pi_X(p)}\le\abs{\Sigma_X(p)}+\abs{\Sigma_X(0)}$.

(ii) By Lemma~\ref{app:G1:ct-elem}(a), $0\le2\sum_\mu G_X(\eps e_\mu)\le6G_X(0)$. By Lemma~\ref{app:G1:ct-elem}(c), $\sum_\mu\widehat{(2q+k)}_\mu^{\,2}\le2\hat q^2+2\widehat{(q+k)}^2$. With $\hat q^2\hat G(q)\le1$ the non-negative bubble term is therefore at most $2\int\hat G(q+k)+2\int\hat G(q)=4G_X(0)$. A difference of two numbers in $[0,6G_X(0)]$ and $[0,4G_X(0)]$ has modulus at most $10G_X(0)$.

(iii) Expand $s_\mu(2P+k)=\frac2\eps[\sin(\eps P_\mu)\cos\frac{\eps k_\mu}2+\cos(\eps P_\mu)\sin\frac{\eps k_\mu}2]$ and use $\frac2\eps\sin\frac{\eps k_\mu}2=s_\mu(k)$:
\[
S_\mu(k)=s_\mu(k)\int\hat G(P)\hat D(P-k)+\cos\tfrac{\eps k_\mu}2\,T_\mu(k)+s_\mu(k)\int\cos(\eps P_\mu)\hat G(P)\hat D(P-k),
\]
where $T_\mu(k):=\int\frac2\eps\sin(\eps P_\mu)\hat G(P)\hat D(P-k)$. The first and the last terms each have modulus at most $\abs{s_\mu(k)}\Psi_X$, by Lemma~\ref{app:G1:ct-elem}(b).

The substitution $P\mapsto-P$ is a symmetry of the momentum set, and $\hat G,\hat D$ are even. It gives $T_\mu=\frac12\int\frac2\eps\sin(\eps P_\mu)\hat G(P)[\hat D(P-k)-\hat D(P+k)]$. Lemma~\ref{app:G1:ct-elem}(d) then gives $\abs{T_\mu}\le4\hat k\int\abs{\hat P_\mu}\hat P\hat G(P)\hat D(P-k)\hat D(P+k)$. Since $\abs{s_\mu(k)}=\abs{\hat k_\mu}$, $\sum_\mu\abs{\hat k_\mu}\abs{\hat P_\mu}\le\hat k\hat P$ and $\hat P^2\hat G(P)\le1$,
\[
\sum_\mu\abs{s_\mu}\abs{S_\mu}\le2\hat k^2\Psi_X+4\hat k^2\int\hat D(P-k)\hat D(P+k)\le\hat k^2(2\Psi_X+4\delta_X),
\]
where the last step is Lemma~\ref{app:G1:ct-elem}(b).
\end{proof}

\begin{proposition}[size of the counterterm]\label{app:K:L13-12}\label{app:G1:ct-size}
For $X\in\{\eps\Z^3,T\}$ (with $\ell\ge\ell_0$), every $\eps\le\eps_{\max}$ and every $M>0$,
\[
\eps\,\abs{\delta m^2_{K,X}(M)}\le C_\delta(M):=(16\lambda+3e^2)\bar\omega_*+128\lambda^2\bar\omega_M\bar\tau(M)+16\lambda e^2\big[10\bar\omega_{\mu_*}\bar\tau(M)+2\bar\psi(M)\bar\omega_M\big]+e^4C_4(M),
\]
where $\bar\omega_*:=\max(\bar\omega_{\mu_*},\bar\omega_M)$ and
\[
C_4(M):=\tfrac54\eps_{\max}\bar\omega_{\mu_*}^2+16\bar\omega_M\bar\delta+14\bar\psi(M)\bar\omega_{\mu_*}+4\bar\delta\bar\omega_{\mu_*},
\]
with the constants of Lemma~\ref{app:G1:ct-oneloop}. The constant $C_\delta(M)$ depends only on $e,\lambda,\mu_*,M,\ell_0$ and $\eps_{\max}$. It is independent of $\eps\le\eps_{\max}$, of $\ell\ge\ell_0$ and of the position, and it is non-increasing in $\ell_0$.
\end{proposition}
\begin{proof}
We bound the terms of Lemma~\ref{app:G1:ct-rep} one by one, using Lemma~\ref{app:G1:ct-oneloop} and $1\le\eps_{\max}/\eps$ throughout.

\emph{One loop.} $16\lambda G_X(0)\le16\lambda\bar\omega_M/\eps$. Since $M^2\hat G\le1$, we have $0\le J_X\le D_X$, so $e^2(2D_X+J_X)\le3e^2\bar\omega_{\mu_*}/\eps$.

\emph{Order $\lambda^2$.} By Lemma~\ref{app:G1:ct-elem}(a) and Parseval, $0\le Q_X\le G_X(0)\sum_z\eps^3G_X(z)^2=G_X(0)\tau_X\le\bar\omega_M\bar\tau(M)/\eps$.

\emph{Order $e^2\lambda$.} By Lemma~\ref{app:G1:ct-se}(i) and $\hat p^2\hat G^2\le\hat G$, we get $\abs{16\lambda\int\hat G^2\tilde\Pi_X}\le16\lambda e^2[10D_X\tau_X+2\Psi_XG_X(0)]$, which is at most $16\lambda e^2[10\bar\omega_{\mu_*}\bar\tau(M)+2\bar\psi(M)\bar\omega_M]/\eps$.

\emph{Order $e^4$.}
\begin{itemize}
\item $\frac34\eps^2e^4D_X^2\le\frac34e^4\bar\omega_{\mu_*}^2\le\frac34e^4\bar\omega_{\mu_*}^2\eps_{\max}/\eps$.
\item $\abs{(C)}\le\frac12\eps^2e^4D_X^2\le\frac12e^4\bar\omega_{\mu_*}^2\eps_{\max}/\eps$, since $0\le D_X-J_X\le D_X$.
\item $0\le(D1)\le\frac{e^4}2\cdot3\cdot2\delta_X\cdot2G_X(0)\le6e^4\bar\delta\bar\omega_M/\eps$. Here we used $D_X,G_X\ge0$, $G_X(z)\le G_X(0)$, and $\sum_z\eps^3D_X(z)^2=\sum_z\eps^3D_X(z+\eps e_\mu)^2=\delta_X$.
\item $\abs{(D3)}\le10e^4\delta_XG_X(0)\le10e^4\bar\delta\bar\omega_M/\eps$, by Lemma~\ref{app:G1:ct-se}(ii).
\item By $M^2\hat G\le1$, Lemma~\ref{app:G1:ct-se}(i) and $\int\hat D\hat G\hat p^2\le D_X$, we get $\abs{(A)}\le e^4[10D_X\Psi_X+2\Psi_XD_X]\le12e^4\bar\psi(M)\bar\omega_{\mu_*}/\eps$.
\item By Lemma~\ref{app:G1:ct-se}(iii) and $\hat k^2\hat G\le1$, $\abs{(D2)}\le e^4(2\Psi_X+4\delta_X)D_X\le e^4(2\bar\psi(M)+4\bar\delta)\bar\omega_{\mu_*}/\eps$.
\end{itemize}
Summing gives $\abs{\delta m^2_{e^4,X}}\le e^4C_4(M)/\eps$, and adding the lower orders gives the claim. Every majorant used is independent of $\eps$ and of $\ell\ge\ell_0$, and non-increasing in $\ell_0$, by Lemma~\ref{app:G1:ct-oneloop}. Independence of the position holds because $\delta m^2_{K,X}(M)$ is one number on $X$.
\end{proof}

\begin{remark}
Take $M=\mu_*=1$ on $\eps\Z^3$ with $\eps_{\max}=1$. Lemma~\ref{app:G1:PD-dos} gives $\bar\psi=\bar\delta=\pi^2/64$, and $\bar\omega_{\mu_*}=\bar\omega_M=w$, so $C_4\le1.405$. Lemma~\ref{app:G1:ct-se}(ii) holds with $6$ in place of $10$; we keep $10$.
\end{remark}

\paragraph{Setting for the restriction tails.} Fix a step $0\le k\le K-1$ and work in scale-$k$ units, so that sites of the unit lattice are scale-$k$ blocks. As in Lemma~\ref{app:G1:kappa0}, $t_k=2^k\eps$ and $\eps=2^{-K}\eps_0$ with $\eps_0\le1$; $N_k=\abs{T}t_k^{-3}$ is the number of scale-$k$ blocks of $T_\eps$; and $\abs{Y}_k$ is the number of scale-$k$ blocks in $Y\subset T$. Put $t=t_k$ and $\ell(t):=1+\log t^{-1}$; here $\ell(t)$ is a function of the scale, not the side of a torus.

Let $z_x=(A'(x),\varphi'(x))\in\R^3\times\R^2$, $x\in\Lambda_5^{(k)}$, be the joint fluctuation variables of (C6) in Lemma~\ref{app:G1:S05-SC}, and let $\mu_k$ be their centred Gaussian law. Let $V=V^{(k)}_w$ be the real polynomial of (C6), of degree at most $D$ and localised in $J=\Lambda_7^{(k)}$. Put $\mathsf A:=\norm{V}_\theta$, the norm of Theorem~\ref{app:K:L13-1} with decay rate $\theta>0$; we write $\theta$ for that rate, which Theorem~\ref{app:K:L13-1} writes $\kappa$, because $\kappa$ is used below for an exponent. For $b>0$ let
\[
\chi_b:=\prod_{x\in\Lambda_5^{(k)}}1\{\abs{A'(x)}\le b\}\,1\{\abs{\varphi'(x)}\le b\},\qquad \chi^{\rm fl}_k:=\chi_{b_k},\qquad c_-\ell(t)\le b_k\le c_+\ell(t)^q,
\]
with fixed $c_\pm>0$ and $q\ge1$.

In the upper bound of \cite[(2.115), p.~581]{Ba2}, the fluctuation cut-offs are those of \cite[(2.86), (2.114)]{Ba2} on $\Lambda_5^{(k)}$, at level $O(1)p(t_k)$ with $p(t)=b_0\ell(t)^{p}$; this is the case $q=p$, $c_+=O(1)b_0$. The cut-off $p_1(t)=b_1\ell(t)^{p_1}$ of \cite[(3.50), (3.54), p.~622]{Ba1} belongs to the lower bound, where $p_1$ is restricted only by $p-1\ge p_1$.

The cumulant bound of Theorem~\ref{app:K:L13-1} enters through the following property, which we assume.

\smallskip\noindent\textbf{(CB)} There are constants $S_R,\rho_3>0$ and $b_*\ge1$ with the following property. They depend only on the covariance class of $\mu_k$ and on $\nbar,D,d,\theta$, and not on $\eps$, $k$, $\ell_T$, the domains or the background. Put $\rho_1:=D+2d$. Then for every $b\ge b_*$ with $\mathsf Ab^{\rho_1}\le1$,
\[
\int d\mu_k\,\chi_b\,\mathrm{e}^{V}\le\exp\Big(\sum_{m=1}^{\nbar}\frac{\kappa_{\mu_k}(V;m)}{m!}+S_R\abs{J}_k\big[(\mathsf Ab^{\rho_1})^{\nbar+1}+\mathrm{e}^{-\rho_3b^{3/2}}\big]\Big).
\]

\begin{lemma}[restriction tails]\label{app:K:L13-13}\label{app:G1:restr-tail}
Assume (CB), $c_-\ell(\eps_0)\ge b_*$ and $\mathsf Ab_k^{\rho_1}\le1$. Then:
\begin{itemize}
\item[(a)] 
\[
\int d\mu_k\,\chi^{\rm fl}_k\,\mathrm{e}^{V}\le\exp\Big(\sum_{m=1}^{\nbar}\frac{\kappa_{\mu_k}(V;m)}{m!}+\mathcal R^{\rm tr}_k+\mathcal R^{\rm res}_k\Big),
\]
with $\mathcal R^{\rm tr}_k:=S_R\abs{J}_k(\mathsf Ab_k^{\rho_1})^{\nbar+1}\ge0$ and $\mathcal R^{\rm res}_k:=S_R\abs{J}_k\,\mathrm{e}^{-\rho_3b_k^{3/2}}\ge0$. Moreover $\abs{J}_k\le N_k$.
\item[(b)] For every $\kappa>0$, $\mathcal R^{\rm res}_k\le C_{\rm res}(\kappa)\,t_k^{\kappa}N_k$, with $C_{\rm res}(\kappa):=S_R\exp\big(\frac{4\kappa^3}{27\rho_3^2c_-^3}-\kappa\big)$.
\item[(c)] If $\mathsf A\le C_At_k^{\kappa_*}$ with $(\nbar+1)\kappa_*>\kappa_B$, then $\mathcal R^{\rm tr}_k\le C_{\rm tr}\,t_k^{\kappa_B}N_k$, with
\[
C_{\rm tr}:=S_R\,(C_Ac_+^{\rho_1})^{\nbar+1}\,S\big((\nbar+1)\kappa_*-\kappa_B,\ (\nbar+1)q\rho_1\big),
\]
where $S(\cdot,\cdot)$ is as in Lemma~\ref{app:G1:kappa0}.
\item[(d)] Assume also the hypothesis of (c), and $\kappa_B>3$. The cumulants in (a) are those of the unrestricted Gaussian $\mu_k$. In the upper bound of step $k$, the indicator functions of the construction change the exponent by at most the non-negative term
\[
\mathcal R^{\rm tr}_k+\mathcal R^{\rm res}_k\le\big(C_{\rm res}(\kappa_B)+C_{\rm tr}\big)t_k^{\kappa_B}N_k .
\]
The change is measured relative to the computation with the unrestricted fluctuation integral, to which Lemmas~\ref{app:G1:S05-TS} and~\ref{app:G1:S05-RF} refer. In particular, the field-independent part of the cumulant sum in (a) is the step-$k$ vacuum output of Lemma~\ref{app:G1:S05-TS}. In the upper bound, the term $\mathcal R^{\rm vac}_k$ of Lemma~\ref{app:G1:S05-RF}(V) is enlarged by $\mathcal R^{\rm tr}_k+\mathcal R^{\rm res}_k$, and the bound there persists with $C_\kappa$ replaced by $C_\kappa+2^{\kappa_B-3}\big(C_{\rm res}(\kappa_B)+C_{\rm tr}\big)$. No lower bound is asserted.
\item[(e)] Let $\kappa_B>3$, and assume the hypotheses of (c) at every step $k<K$, with the same constants. Then, uniformly in $\eps$, $K$ and $\ell_T$,
\[
\sum_{k<K}\big(\mathcal R^{\rm tr}_k+\mathcal R^{\rm res}_k\big)\le\frac{C_{\rm res}(\kappa_B)+C_{\rm tr}}{1-2^{-(\kappa_B-3)}}\,\eps_0^{\kappa_B-3}\abs{T}.
\]
At $\kappa_B=\frac{16}5$ the prefactor is $(1-2^{-1/5})^{-1}=7.725\ldots$.
\end{itemize}
All constants depend only on the constants of (CB), on $\kappa$, $\kappa_*$, $\kappa_B$, $c_\pm$, $q$ and $C_A$. If the constants of (CB) and $C_A$ do not depend on the slice parameter $u$ of Subsection~\ref{ssec:T6-statement}, then neither do the bounds (a)--(e).
\end{lemma}

\begin{proof}
(a) Apply (CB) with $b=b_k$. Since $\ell$ is decreasing and $t_k\le\eps_0$, we have $b_k\ge c_-\ell(t_k)\ge c_-\ell(\eps_0)\ge b_*$; and $\mathsf Ab_k^{\rho_1}\le1$ by assumption. Since $J=\Lambda_7^{(k)}\subset T$, $\abs{J}_k\le N_k$.

(b) Write $\ell=\ell(t)\ge1$, so that $t^\kappa=\mathrm{e}^{-\kappa(\ell-1)}$, and put $\rho:=\rho_3c_-^{3/2}$. Since $b_k\ge c_-\ell$,
\[
\mathrm{e}^{-\rho_3b_k^{3/2}}t^{-\kappa}\le\exp\big(\kappa\ell-\rho\ell^{3/2}-\kappa\big)\le\exp\Big(\sup_{x\ge0}(\kappa x-\rho x^{3/2})-\kappa\Big)=\exp\Big(\frac{4\kappa^3}{27\rho^2}-\kappa\Big).
\]
The supremum is attained at $x^{1/2}=2\kappa/(3\rho)$, and $\rho^2=\rho_3^2c_-^3$. Multiply by $S_R\abs{J}_k\le S_RN_k$.

(c) Since $b_k\le c_+\ell^q$,
\[
(\mathsf Ab_k^{\rho_1})^{\nbar+1}\le(C_Ac_+^{\rho_1})^{\nbar+1}t^{\kappa_B}\cdot t^{(\nbar+1)\kappa_*-\kappa_B}\ell(t)^{(\nbar+1)q\rho_1}.
\]
Since $0<t\le\eps_0\le1$, Lemma~\ref{app:G1:kappa0} bounds the last two factors by $S((\nbar+1)\kappa_*-\kappa_B,(\nbar+1)q\rho_1)$. Multiply by $S_R\abs{J}_k\le S_RN_k$.

(d) The integrand of the step is non-negative. The indicator functions and the large-field factors $\zeta_{\Lambda_0^{(j)}}$ of the construction take values in $[0,1]$. Following \cite[p.~581]{Ba2}, they are treated as follows:
\begin{itemize}
\item the cut-off $\chi^{\rm fl}_k$ is inserted in the joint fluctuation integral;
\item the factor $\chi_{k+1}$, a function of the fields of scale $k+1$ only, is kept outside that integral;
\item all remaining indicator functions of the fields on $\Lambda_5^{(k)}$, together with the factors $\zeta$, are bounded by $1$.
\end{itemize}
This is legitimate in an upper bound and changes no Gaussian cumulant. The restrictions \cite[(2.55)]{Ba2} on the current fields enter only as conditions on the external fields, under which (a) holds pointwise. The complementary large-field terms are those of Lemma~\ref{app:G1:S05-LF}; they are not part of the small-field output.

By (a), the only change of the exponent relative to the unrestricted cumulant sum is $\mathcal R^{\rm tr}_k+\mathcal R^{\rm res}_k$, which (b) and (c) at $\kappa=\kappa_B$ bound by $(C_{\rm res}(\kappa_B)+C_{\rm tr})t_k^{\kappa_B}N_k$. Lemma~\ref{app:G1:S05-RF}(V) bounds increments by $C_\kappa t_{k-1}^{\kappa_B}N_{k-1}$. Put $a_k:=\mathcal R^{\rm tr}_k+\mathcal R^{\rm res}_k\ge0$; then $\abs{a_k-a_{k-1}}\le\max(a_k,a_{k-1})$. Moreover $t_k^{\kappa_B}N_k=\abs{T}t_k^{\kappa_B-3}=2^{\kappa_B-3}t_{k-1}^{\kappa_B}N_{k-1}$, because $t_k=2t_{k-1}$; since $\kappa_B>3$, $t^{\kappa_B}N$ increases with the scale. Hence the enlarged increment is at most $\big(C_\kappa+2^{\kappa_B-3}(C_{\rm res}(\kappa_B)+C_{\rm tr})\big)t_{k-1}^{\kappa_B}N_{k-1}$ plus the layer and deep terms of Lemma~\ref{app:G1:S05-RF}(V1). This is the stated enlargement of $C_\kappa$.

(e) Since $t_k^{\kappa_B}N_k=\abs{T}\,t_k^{\kappa_B-3}$ and $t_k=2^{-(K-k)}\eps_0$, (b) and (c) bound the sum by
\[
(C_{\rm res}(\kappa_B)+C_{\rm tr})\abs{T}\eps_0^{\kappa_B-3}\sum_{m\ge1}2^{-m(\kappa_B-3)}\le(C_{\rm res}(\kappa_B)+C_{\rm tr})\abs{T}\eps_0^{\kappa_B-3}\big(1-2^{-(\kappa_B-3)}\big)^{-1}.
\]
The bound does not depend on $\eps$, $K$ or $\ell_T$.
\end{proof}

\begin{remark}
Since $b_k\ge c_-\ell(t)$, the restriction tail $\mathcal R^{\rm res}_k$ is $O(t^\kappa N_k)$ for every $\kappa$. The slowest volume term of the step is the truncation tail, of order $t^{(\nbar+1)\kappa_*}$ times a power of $\ell(t)$. Take $\kappa_*=\frac14-\beta'$, which is the exponent $\frac{4-d}4-\beta$ of \cite[p.~421]{Ba3}, with $\beta'=\frac1{56}$ and $\nbar=13$. Then the truncation tail is $t^{13/4}$ up to logarithms. The logarithmic factor excludes $\kappa_B=\frac{13}4$, which is why $\kappa_B=\frac{16}5\in(3,\frac{13}4)$. Only the upper half of the cumulant bound is used.
\end{remark}

\subsection{The Schur complement, the mass and vacuum flow, and the anchor blocks}\label{app:G1:SchurFlow}

Throughout, $t_j=L^j\eps$ and $\eps=L^{-K}\eps_0$. For $y\in T^{(k')}$ and $0\le n\le k'$, $B^n(y)\subset T^{(k'-n)}$ is the set of points of scale $k'-n$ in the block of $y$. For a set $Y$ we put $B^n(Y):=\bigcup_{y\in Y}B^n(y)$. Since blocks are nested, $B^{n+n'}(Y)=B^n(B^{n'}(Y))$.

\paragraph{Gaussian marginals.}

\begin{lemma}[Gaussian marginals]\label{app:G1:SchurMarg}
Let $q$ be a positive definite quadratic form on $\R^V\oplus\R^W$, and put $\operatorname{Sch}_Vq(v):=\min_wq(v,w)$. Then:
\begin{itemize}
\item[(a)] $\int_{\R^W}\mathrm{e}^{-\frac12q(v,w)}\,dw=(2\pi)^{\abs W/2}\det(q_{WW})^{-1/2}\,\mathrm{e}^{-\frac12\operatorname{Sch}_Vq(v)}$;
\item[(b)] $\operatorname{Sch}_{V_1}\operatorname{Sch}_Vq=\operatorname{Sch}_{V_1}q$ for $V_1\subset V$;
\item[(c)] $q\le q'$ implies $\operatorname{Sch}_Vq\le\operatorname{Sch}_Vq'$, and $\operatorname{Sch}_V(cq)=c\operatorname{Sch}_Vq$ for $c>0$;
\item[(d)] if $D$ is an invertible linear map of $\R^V\oplus\R^W$ that maps $\R^V$ onto $\R^V$ and $\R^W$ onto $\R^W$, then $\operatorname{Sch}_V(q\circ D)=(\operatorname{Sch}_Vq)\circ D|_{\R^V}$.
\end{itemize}
Moreover, let $q_1$ and $q_2$ be quadratic forms such that $\mathrm{e}^{-\frac12q_1}$ and $\mathrm{e}^{-\frac12q_2}$ are proportional as functions of $v$. Then $q_1=q_2$.
\end{lemma}

\begin{proof}
Write $q(v,w)=\ip{(v,w)}{\mathcal Q(v,w)}$ with blocks $q_{VV},q_{VW},q_{WV},q_{WW}$. Here $q_{WW}>0$. Put $w_*(v):=-q_{WW}^{-1}q_{WV}v$. Completing the square gives
\[
q(v,w)=\operatorname{Sch}_Vq(v)+\ip{w-w_*(v)}{q_{WW}(w-w_*(v))},
\]
so the minimum over $w$ is attained at $w_*(v)$. Statement (a) follows from the Gaussian integral over $w-w_*(v)$.

Statement (b) holds because minimising in two stages equals minimising jointly. Statement (c) holds because pointwise inequalities and positive factors survive minimisation. Statement (d) holds because $D$ maps the variables in $\R^W$ bijectively onto themselves, so the minimisation over them commutes with $D$.

For the last claim: $q_1-q_2$ is constant, and it vanishes at $v=0$.
\end{proof}

\paragraph{The all-scale photon form.}
The photon field is treated in the Feynman form, with external field $0$. The form below is a sum of identical forms over the $d$ components, so we treat one component, and all fields are real valued.

On every lattice we use the unweighted inner product $\ip fg=\sum_xf(x)g(x)$, and $Q^*$ is the adjoint for these products. By \cite[(2.7), p.~608]{Ba1} at zero external field, $(QB)(y)=L^{-d}\sum_{x\in B(y)}B(x)$. Hence $QQ^*=L^{-d}I$, and by induction $Q^n(Q^n)^*=L^{-dn}I$.

Let $B_0$ be a field on $T^{(0)}$ and $B_j$ a field on $T^{(j)}$, $1\le j\le K$. Put $c_j:=a\,t_j^{d-2}$, the coefficient of the kernel \cite[(2.6), p.~608]{Ba1}. Let $\Delta^{(0)}$ be the positive definite operator of the $B_0$-term of the photon chain of Lemma~\ref{app:G1:S05-T}. Put
\[
S^1(B_0,\dots,B_K):=\ip{B_0}{\Delta^{(0)}B_0}+\sum_{j=1}^Kc_j\sum_{y\in T^{(j)}}\abs{B_j(y)-(QB_{j-1})(y)}^2 .
\]
This is the quadratic form of the photon chain of Lemma~\ref{app:G1:S05-T} with all weights equal to $1$.

$S^1$ is positive definite. Indeed, if $S^1(B)=0$ then $B_0=0$, and inductively $B_j=QB_{j-1}=0$. Only the positive definiteness of $\Delta^{(0)}$ is used below, so everything holds for $\Delta^{(0)}=-\Delta+\mu_0^2$ as well as for $\Delta^{(0)}=-\Delta+\mu_*^2$. By Lemma~\ref{app:G1:SchurMarg}(d), it also holds in any diagonally rescaled units, for instance the rescaled units of \cite{Ba2}.

\begin{lemma}[composed kernels are a Schur complement]\label{app:K:L14-1}\label{app:G1:SchurComp}
Put $Q_k:=Q^k$ and $a_k:=a(1-L^{-2})/(1-L^{-2k})$. Let $Y_0\subset T^{(0)}$ and $Y_k\subset T^{(k)}$, $1\le k\le K$, satisfy the shadow condition
\[
\text{(Sh)}\qquad Y_k\cap B^{k'-k}(Y_{k'})=\emptyset\qquad(1\le k<k'\le K).
\]
Put $\Psi:=(B_0|_{Y_0},B_k|_{Y_k})_{1\le k\le K}$. Then $\operatorname{Sch}_\Psi S^1=\Delta_v$, where
\[
\ip\Psi{\Delta_v\Psi}=\min_{B_0|_{T^{(0)}\setminus Y_0}}\Big[\ip{B_0}{\Delta^{(0)}B_0}+\sum_{k=1}^Ka_kt_k^{d-2}\sum_{x\in Y_k}\abs{B_k(x)-(Q_kB_0)(x)}^2\Big].
\]
Equivalently, integrating $\mathrm{e}^{-\frac12S^1}$ over all variables other than $\Psi$ gives a constant times $\mathrm{e}^{-\frac12\ip\Psi{\Delta_v\Psi}}$. For the shells $Y_k=\mathrm{Sh}_k$ of Lemma~\ref{app:K:L14-2}, this is the vector form of \cite[(3.23), p.~588]{Ba2}.
\end{lemma}

\begin{proof}
\emph{Step 1: a probabilistic description.} Normalise $\mathrm{e}^{-\frac12S^1}$ to a probability density. We have
\[
\mathrm{e}^{-\frac12S^1}=\mathrm{e}^{-\frac12\ip{B_0}{\Delta^{(0)}B_0}}\prod_{j=1}^K\mathrm{e}^{-\frac12c_j\norm{B_j-QB_{j-1}}^2}.
\]
For every $B_{j-1}$, the $j$-th factor has integral $(2\pi/c_j)^{n_j/2}$ in $B_j$, where $n_j$ is the number of real variables of $B_j$ \cite[(2.8), p.~609]{Ba1}. Integrating successively over $B_K,\dots,B_1$ shows the following:
\begin{itemize}
\item $B_0$ has density proportional to $\mathrm{e}^{-\frac12\ip{B_0}{\Delta^{(0)}B_0}}$;
\item $B_j=QB_{j-1}+\xi_j$, where the $\xi_j\sim N(0,c_j^{-1}I)$ are independent of each other and of $B_0$.
\end{itemize}
Thus $B_k=Q_kB_0+\sum_{j\le k}Q^{k-j}\xi_j$.

\emph{Step 2: the conditional covariance.} Condition on $B_0$, and let $k\le k'$, $x\in T^{(k)}$, $x'\in T^{(k')}$. Using $Q^{k'-j}=Q^{k'-k}Q^{k-j}$ and $Q^{k-j}(Q^{k-j})^*=L^{-d(k-j)}I$,
\[
\operatorname{Cov}\big(B_k(x),B_{k'}(x')\big)=\sum_{j=1}^kc_j^{-1}\big(Q^{k-j}(Q^{k-j})^*(Q^{k'-k})^*\big)(x,x')=\Big(\sum_{j=1}^kc_j^{-1}L^{-d(k-j)}\Big)(Q^{k'-k})(x',x).
\]
Since $t_j=L^{j-k}t_k$, we have $c_j^{-1}L^{-d(k-j)}=a^{-1}t_k^{2-d}L^{-2(k-j)}$. Hence the bracket equals
\[
a^{-1}t_k^{2-d}\sum_{i=0}^{k-1}L^{-2i}=(a_kt_k^{d-2})^{-1};
\]
this is the solution \cite[(2.15), p.~609]{Ba1} of the recursion \cite[(2.13)]{Ba1}.

Moreover $(Q^{k'-k})(x',x)=L^{-d(k'-k)}\mathbf 1\{x\in B^{k'-k}(x')\}$. So the covariance is $(a_kt_k^{d-2})^{-1}\delta_{xx'}$ if $k=k'$. It vanishes if $k<k'$, $x\in Y_k$ and $x'\in Y_{k'}$, by (Sh).

\emph{Step 3: the marginal.} Jointly Gaussian variables with zero covariance are independent. So, given $B_0$, the variables $B_k(x)$ ($k\ge1$, $x\in Y_k$) are independent with law $N\big((Q_kB_0)(x),(a_kt_k^{d-2})^{-1}\big)$. Hence the marginal density of $(B_0,(B_k|_{Y_k})_{k\ge1})$ is proportional to $\mathrm{e}^{-\frac12q_\flat}$, where $q_\flat$ is the bracket in the statement.

By Lemma~\ref{app:G1:SchurMarg}(a), this marginal is also proportional to $\mathrm{e}^{-\frac12\operatorname{Sch}S^1}$, the Schur complement being taken onto the same variables. By the last clause of Lemma~\ref{app:G1:SchurMarg}, $q_\flat$ equals that Schur complement. Minimising over $B_0|_{T^{(0)}\setminus Y_0}$ and using Lemma~\ref{app:G1:SchurMarg}(b) gives $\operatorname{Sch}_\Psi S^1=\Delta_v$. The equivalent formulation follows from Lemma~\ref{app:G1:SchurMarg}(a).
\end{proof}

\begin{lemma}[the shells satisfy the shadow condition]\label{app:K:L14-2}\label{app:G1:SchurShad}
For every admissible region sequence $\boldsymbol\Lambda$, put $\mathrm{Sh}_0:=\Lambda_5^{(0)c}$ and $\mathrm{Sh}_k:=\Lambda_5^{(k-1)\prime}\cap\Lambda_5^{(k)c}$ for $1\le k\le K$, with $\Lambda_5^{(K)}:=\emptyset$ as in \cite[(3.23), p.~588]{Ba2}. These sets satisfy condition (Sh) of Lemma~\ref{app:K:L14-1}.
\end{lemma}

\begin{proof}
The prime operation is $\Lambda'=\Lambda\cap L\eta\Z^d$ \cite[(1.16)--(1.17), p.~606]{Ba1}, and every point $y$ lies in its own block $B(y)$. Hence, if $\Lambda$ is a union of $L$-blocks, then $Y\subset\Lambda'$ implies $B^1(Y)\subset\Lambda$.

All the sets $\Lambda_i^{(j)}$ are unions of large blocks \cite[p.~566; (2.7)--(2.8), p.~558]{Ba2}, hence unions of $L$-blocks. By \cite[(2.8)]{Ba2}, $\Lambda_{i+1}^{(j)}\subset\Lambda_i^{(j)}$. Admissibility gives $\Lambda_0^{(j)}\subset\Lambda_7^{(j-1)\prime}$ for $j\ge1$ \cite[p.~566]{Ba2}. Hence
\[
\Lambda_5^{(j)}\subset\Lambda_0^{(j)}\subset\Lambda_7^{(j-1)\prime}\subset\Lambda_5^{(j-1)\prime},
\]
that is, $B^1(\Lambda_5^{(j)})\subset\Lambda_5^{(j-1)}$ for $j\ge1$.

Let $1\le k<k'\le K$. Since $\mathrm{Sh}_{k'}\subset\Lambda_5^{(k'-1)\prime}$, we get $B^1(\mathrm{Sh}_{k'})\subset\Lambda_5^{(k'-1)}$. Using $B^{n+1}=B^n\circ B^1$ and the previous inclusion $k'-k-1$ times gives $B^{k'-k}(\mathrm{Sh}_{k'})\subset\Lambda_5^{(k)}$. This set is disjoint from $\mathrm{Sh}_k\subset\Lambda_5^{(k)c}$.
\end{proof}

\paragraph{The mass flow.}

\begin{lemma}[pinned mass coefficients along the flow]\label{app:K:L14-3}\label{app:G1:RFpin}
Let $L=2$, $\eps=2^{-K}\eps_0$, $t_j=2^j\eps$ (so $t_K=\eps_0$), $e(t)=et^{1/2}$ and $\lambda(t)=\lambda t$. As in Lemma~\ref{app:G1:S05-RF}, put
\[
P_m(t):=(e^2+\lambda)t+(e^4+e^2\lambda+\lambda^2)t^2,\qquad P_K(t):=2(e^2+\lambda)t+(e^4+e^2\lambda+\lambda^2)t^2\big(1+\log_2(\eps_0/t)\big).
\]
Consider the counterterm graphs $G_0$ of \cite[(1.23)]{Ba3} of order $e^\alpha\lambda^\beta$ with $2\le\alpha+2\beta\le4$, and put $o:=\alpha/2+\beta$. Lower-order $\delta m^2$ insertions are kept as the partners of their self-energy subgraphs. Then:
\begin{itemize}
\item[(a)] for all $x\in T_\eps$ and $0\le k\le K$,
\[
\abs{\delta m^2_{K,k}(x)}t_k^2\le C_m\big[2(e^2+\lambda)t_k+(e^4+e^2\lambda+\lambda^2)t_k^2\log_2(\eps_0/t_k)\big]\le C_mP_K(t_k);
\]
\item[(b)] work in scale-$k$ units ($L^k\eta=1$), and put $t=t_k$, $W:=B^k(\Lambda_7^{(k)})$, $\mathcal L_k:=B^k(\Lambda_7^{(k)}\cap\Lambda_8^{(k)c})$ and $\mathcal D_k:=B^k(\Lambda_8^{(k)})$, as in Lemma~\ref{app:G1:S05-RF}.

Let $\Delta_k(x):=[\delta m^2_{K,k}(x)-\delta m^2_{K,k+1}(x)]t^2$. This is the sum over $G_0$ of the scale-$k$ coefficients with largest index $k$, leg~1 fixed at $x$, leg~2 summed, external fields $\equiv1$ and background $0$. Localise it as in \cite[p.~420]{Ba3}, $\Delta_k(x)=\sum_\Lambda\Delta_k(x;\Lambda)$, with $\square(p)\ni x$ for the vertex $p$ of leg~1. The localisation cubes are cubes of side $1$ or $2$ with corners in the unit lattice, and $c_d$ is twice their largest diameter.

Call $\Lambda$ bad if some cube of $\Lambda$ is not contained in $W$, and let $\Delta_k^{\rm bad}(x)$ be the sum over the bad $\Lambda$. Then for all $x$
\[
\abs{\Delta_k(x)},\ \abs{\Delta_k^{\rm bad}(x)}\le C_dP_m(t),\qquad \abs{\Delta^{\rm bad}_k(x)}\le C_dP_m(t)\,\mathrm{e}^{\delta_0c_d}\mathrm{e}^{-\delta_0r(t)/2}\quad(x\in\mathcal D_k).
\]
If $\abs{\varphi(x)}\le c_1p(t)\lambda(t)^{-1/4}$ on $W$, then
\[
\Big|\sum_{x\in W}\eta^d\Delta^{\rm bad}_k(x)\abs{\varphi(x)}^2\Big|\le c_1^2C_dP_m(t)\lambda(t)^{-1/2}p(t)^2\big[\abs{\mathcal L_k}+\mathrm{e}^{\delta_0c_d}\mathrm{e}^{-\delta_0r(t)/2}\abs{\mathcal D_k}\big].
\]
\end{itemize}
Here $C_m=C_PN_{\rm ct}$ and $C_d=C_PN_{\rm ct}S_0$. In these formulas:
\begin{itemize}
\item $C_P$ is the constant of the pinned bound of Lemmas~\ref{app:G1:VSR-A}, \ref{app:G1:VSR-Btop} and~\ref{app:G1:VSR-C}, see (P) below;
\item $N_{\rm ct}$ is the number of counterterm graphs of order $\le4$;
\item $S_0:=\sum_{D\ge0}(c'_d(D+2)^d)^{N_s}\mathrm{e}^{-\delta_0D/2}$, where $N_s$ is the maximal number of localised vertices of a graph of order $\le\nbar$ and $c'_d\ge1$ depends only on $d$ and $c_d$.
\end{itemize}
The constants are independent of $\eps$, $k$, $x$, the domains $\Lambda_i^{(k)}$ and $\ell_T$.
\end{lemma}

\begin{proof}
For $k\le j^*<K$, let $c^{[j^*]}_{G_0}(x)$ be the scale-$K$ coefficient of $G_0$ with leg~1 fixed at $x$, leg~2 summed, external fields $\equiv1$ and background $0$, summed over the index assignments \cite[(1.27)]{Ba3} with largest index $j^*$. By \cite[(1.27)--(1.28)]{Ba3} and multilinearity,
\[
\delta m^2_{K,k}(x)t_K^2=\sum_{G_0}\sum_{j^*=k}^{K-1}c^{[j^*]}_{G_0}(x).
\]

\emph{Pinned bound.} Fix $G_0$ and $j^*$. By Lemma~\ref{app:G1:VSR-Btop}, every renormalisation acting in the sub-sum with largest index $j^*$ concerns proper subgraphs with lines of index $\le j^*$, and $G_0$ itself is not renormalised. By Lemma~\ref{app:G1:VSR-C}, all these transformations can be done with the pin fixed. Lemma~\ref{app:G1:VSR-A}(ii) at fixed largest index then gives, in scale-$K$ units,
\[
\abs{c^{[j^*]}_{G_0}(x)}\le C_P\,e(t_K)^\alpha\lambda(t_K)^\beta\,(L^{j^*}\eta_K)^{D(G_0)}.\tag{P}
\]

Now keep the localisation $\Lambda$. Extract the tree factor of \cite[(2.13)--(2.14), pp.~426--427]{Ba3} before the induction, together with the cuts of \cite[pp.~432--433]{Ba3} in the form of Lemma~\ref{app:G1:VSR-C}(vii). Then (P) holds for each $\Lambda$ with the extra factor $\mathrm{e}^{-\delta_0d(\Lambda)}$. Here $d(\Lambda)$ is the length, in units of the scale $j^*$, of a shortest tree through the cubes of $\Lambda$.

On $T_\eps$ all distances are torus distances. The slices obey \cite[(2.10)--(2.12)]{Ba3} with torus distance, and every localisation count on $T_\eps$ is at most the corresponding count on $\eta\Z^d$.

\emph{Degree.} $G_0$ has two scalar legs and no external vector leg. It has no external derivative leg, because a difference of a constant external field vanishes at zero background. In a representative without insertions, every vector line joins two vector legs, each carrying one factor $e$. Hence $\alpha=2I_v$, with $I_v$ the number of vector lines; in particular $\alpha$ is even, and $o\in\{1,2\}$. Also the number of scalar vertices of degree four is $\beta$.

Lemma~\ref{app:G1:WP-deg}(b), extended to counterterm insertions by \cite[(2.3)]{Ba3}, gives $D(G_0)=I_v+\beta+1-3=o-2$.

\emph{(a).} We have $e(t_K)^\alpha\lambda(t_K)^\beta=\abs e^\alpha\lambda^\beta t_K^{o}$ and $L^{j^*}\eta_K=t_{j^*}/t_K$.
\begin{itemize}
\item For $o=1$: $\sum_{j^*=k}^{K-1}t_K/t_{j^*}=\sum_{n=1}^{K-k}2^n\le2t_K/t_k$.
\item For $o=2$: the sum has $K-k=\log_2(\eps_0/t_k)$ terms, each equal to $1$.
\end{itemize}
At $o=1$, $\abs e^\alpha\lambda^\beta\le e^2+\lambda$; at $o=2$, $\abs e^\alpha\lambda^\beta\le e^4+e^2\lambda+\lambda^2$. Hence
\[
\abs{\delta m^2_{K,k}(x)}t_k^2\le C_PN_{\rm ct}\Big[(e^2+\lambda)\frac{t_k^2}{t_K^2}\,t_K\,\frac{2t_K}{t_k}+(e^4+e^2\lambda+\lambda^2)\,t_k^2\log_2\frac{\eps_0}{t_k}\Big],
\]
which is the first inequality. The second holds because $\log_2(\eps_0/t)\le1+\log_2(\eps_0/t)$. For $k=K$ the sum is empty.

\emph{(b), all $x$.} We have $\Delta_k(x)=(t_k/t_K)^2\sum_{G_0}c^{[k]}_{G_0}(x)$. By the localised form of (P) with $j^*=k$, and with $d(\Lambda)$ in scale-$k$ units,
\[
\abs{\Delta_k(x;\Lambda)}\le C_P\Big(\frac{t_k}{t_K}\Big)^2\abs e^\alpha\lambda^\beta t_K^o\Big(\frac{t_k}{t_K}\Big)^{o-2}\mathrm{e}^{-\delta_0d(\Lambda)}=C_Pe(t)^\alpha\lambda(t)^\beta\mathrm{e}^{-\delta_0d(\Lambda)}.
\]
Fix $\square(p)\ni x$. Every cube of $\Lambda$ lies within $d(\Lambda)+c_d$ of $\square(p)$. The number of cubes of side $1$ or $2$ with corners in the unit lattice within distance $D+1+c_d$ of $\square(p)$ is at most $c'_d(D+2)^d$, for a constant $c'_d\ge1$ depending only on $d$ and $c_d$. Hence
\[
\#\{\Lambda:d(\Lambda)\in[D,D+1)\}\le(c'_d(D+2)^d)^{N_s},\qquad \sum_\Lambda\mathrm{e}^{-\delta_0d(\Lambda)}\le\sum_{D\ge0}(c'_d(D+2)^d)^{N_s}\mathrm{e}^{-\delta_0D}\le S_0 .
\]
Since $\sum_{G_0}e(t)^\alpha\lambda(t)^\beta\le N_{\rm ct}P_m(t)$, this bounds $\Delta_k$, and a fortiori $\Delta^{\rm bad}_k$, which is a sub-sum of absolute values.

\emph{(b), $x\in\mathcal D_k$.} By \cite[(2.8)]{Ba2}, $\Lambda_8^{(k)c}$ is the union of the large blocks within distance $r(t)$ of $\Lambda_7^{(k)c}$. Hence $\dist(x,W^c)>r(t)-c_d/2$.

A bad $\Lambda$ has a tree that joins $\square(p)\ni x$ to a cube meeting $W^c$. Both cubes have diameter at most $c_d/2$, so $d(\Lambda)\ge\dist(x,W^c)-c_d\ge r(t)-2c_d$. Then
\[
\mathrm{e}^{-\delta_0d(\Lambda)}\le\mathrm{e}^{-\delta_0(r(t)-2c_d)/2}\mathrm{e}^{-\delta_0d(\Lambda)/2}=\mathrm{e}^{\delta_0c_d}\mathrm{e}^{-\delta_0r(t)/2}\mathrm{e}^{-\delta_0d(\Lambda)/2}.
\]
The count above, now with rate $\delta_0/2$, gives the factor $S_0$.

\emph{Vertex form.} On $W$, $\abs\varphi^2\le c_1^2p(t)^2\lambda(t)^{-1/2}$. Since $\Lambda_8^{(k)}\subset\Lambda_7^{(k)}$, we have $W=\mathcal L_k\sqcup\mathcal D_k$. Apply the first bound on $\mathcal L_k$ and the second on $\mathcal D_k$. Finally, $\eta^d$ times the number of sites of a set is its volume in scale-$k$ units.
\end{proof}

\begin{proposition}[size of the bare mass counterterm]\label{app:K:L14-4}\label{app:G1:Cdelta}
Let $X\in\{\eps\Z^3,T_\eps\}$ with $\ell_T\ge\ell_0$. Let $\int$ denote the Brillouin-zone integral, resp.\ $\ell_T^{-3}\sum_{p\in T^*}$. Put $\hat p_\mu=s_\mu(p):=\frac2\eps\sin\frac{\eps p_\mu}2$, $\hat p^2=\sum_\mu\hat p_\mu^2$, $\hat G=(\hat p^2+M^2)^{-1}$ and $\hat D=(\hat p^2+\mu_*^2)^{-1}$, and
\[
D_X=\textstyle\int\hat D,\quad G_X(0)=\int\hat G,\quad X_X(0)=\int\hat D\hat G,\quad\delta_X=\int\hat D^2,\quad\tau_X=\int\hat G^2,\quad J_X=M^2\int\hat D\hat G .
\]
Let $G_X(z)$, $D_X(z)$ be the position-space kernels. Assume:
\begin{itemize}
\item[(Bd)] for $\eps\le\eps_{\max}$: $\eps D_X\le\bar\omega$, $\eps G_X(0)\le\bar\omega_M$, $X_X(0)\le\bar\psi(M)$, $\delta_X\le\bar\delta$ and $\tau_X\le\bar\tau(M)$, with constants independent of $\eps$ and of $\ell_T\ge\ell_0$;
\item[(Ct)] the coefficients of King's counterterm \cite{K1} through order $4$ are
\begin{itemize}
\item $\delta m^2_\lambda=-16\lambda G_X(0)$;
\item $\delta m^2_{e^2}=-2e^2D_X-e^2J_X$;
\item $\delta m^2_{\lambda^2}=128\lambda^2\sum_z\eps^3G_X(z)^3$;
\item $\delta m^2_{e^2\lambda}=16\lambda\int\hat G^2\tilde\Pi$;
\item $\delta m^2_{e^4}=\tfrac34\eps^2e^4D_X^2+(A)+(C)+(D1)+(D2)+(D3)$.
\end{itemize}
Here $\tilde\Pi(p)=\Sigma(p)-\Sigma(0)$, with $\Sigma(p)/e^2=D_X\sum_\mu\cos(\eps p_\mu)-\int_{p'}\hat D(p-p')\widehat{(p+p')}^2\hat G(p')$, and
\begin{itemize}
\item $(A)=e^2M^2\int\hat D\hat G^2\tilde\Pi$;
\item $(C)=-\frac{\eps^2e^4D_X}{2}(D_X-J_X)$;
\item $(D1)=\frac{e^4}2\sum_\mu\sum_z\eps^3[D_X(z)^2+D_X(z+\eps e_\mu)^2][G_X(z)+G_X(z+\eps e_\mu)]$;
\item $(D2)=-e^4\sum_\mu\int_ks_\mu(k)\hat G(k)\hat D(k)S_\mu(k)$, with $S_\mu(k)=\int_P[s_\mu(k)+s_\mu(2P+k)]\hat G(P)\hat D(P-k)$;
\item $(D3)=e^2\int\hat D^2\sum_\mu\Pi_{\mu\mu}$, with $\Pi_{\mu\mu}(k)/e^2=2G_X(\eps e_\mu)-\int_q\widehat{(2q+k)}_\mu^2\hat G(q)\hat G(q+k)$.
\end{itemize}
\end{itemize}
Let $\delta m^2_{K,X}(M)$ be the sum of these five coefficients. Then $\eps\abs{\delta m^2_{K,X}(M)}\le C_\delta(M)$ for every dyadic $\eps\le\eps_{\max}$, where
\[
C_\delta(M)=(16\lambda+3e^2)\max(\bar\omega,\bar\omega_M)+128\lambda^2\bar\omega_M\bar\tau(M)+16\lambda e^2[10\bar\omega\bar\tau(M)+2\bar\psi(M)\bar\omega_M]+e^4C_4(M),
\]
\[
C_4(M)=\tfrac54\eps_{\max}\bar\omega^2+16\bar\omega_M\bar\delta+14\bar\psi(M)\bar\omega+4\bar\delta\bar\omega .
\]
Consequently $\lambda(\eps)s^2+\frac12\delta m^2_{K,X}(M)\eps^2s\ge-C_\delta(M)^2\eps/(16\lambda)$ for all $s\ge0$.
\end{proposition}

\begin{proof}
\emph{Elementary facts.}
\begin{itemize}
\item The kernel of $(-\Delta+m^2)^{-1}$ on $X$, $m>0$, is pointwise non-negative: it is the Neumann series of $(2d\eps^{-2}+m^2)^{-1}$ times a non-negative hopping matrix. It is also positive definite and translation invariant. Hence $0\le G_X(z)=\int\hat G\,\mathrm{e}^{\mathrm{i}pz}\le G_X(0)$, and the same holds for $D_X$.
\item A convolution $\int_{p'}\hat D(p-p')\hat G(p')$ is the transform of $D_X(z)G_X(z)\ge0$. It is therefore maximal at $p=0$, where it equals $X_X(0)$.
\item All integrands are periodic with period $2\pi/\eps$ in each component, and the set of momenta is a group modulo $2\pi/\eps$, invariant under $P\mapsto-P$. So $\int_{p'}\hat D(p-p')=D_X$, and reflections and translations are allowed; $\hat G$ and $\hat D$ are even.
\item $\hat p^2\hat G\le1$ and $M^2\hat G\le1$.
\item $\widehat{(a+b)}_\mu^2\le2\hat a_\mu^2+2\hat b_\mu^2$, because $\abs{\sin(x+y)}\le\abs{\sin x}+\abs{\sin y}$.
\item By Parseval, $\sum_z\eps^3D_X(z)^2=\delta_X$ and $\sum_z\eps^3G_X(z)^2=\tau_X$.
\end{itemize}

\emph{Self-energy.} By these facts,
\[
\abs{\Sigma(p)}\le e^2\Big[3D_X+\int\hat D(p-p')(2\hat p'^2+2\hat p^2)\hat G(p')\Big]\le e^2[5D_X+2\hat p^2X_X(0)],
\]
hence $\abs{\tilde\Pi(p)}\le e^2[10D_X+2\hat p^2X_X(0)]$.

\emph{Orders $\lambda$, $e^2$, $\lambda^2$, $e^2\lambda$.}
\begin{itemize}
\item One loop: $0\le J_X\le D_X$, so $\abs{\delta m^2_\lambda}+\abs{\delta m^2_{e^2}}\le16\lambda G_X(0)+3e^2D_X$. Times $\eps$, this is at most $(16\lambda+3e^2)\max(\bar\omega,\bar\omega_M)$.
\item Order $\lambda^2$: $\sum_z\eps^3G_X^3\le G_X(0)\tau_X$, which gives $128\lambda^2\bar\omega_M\bar\tau(M)$ after multiplying by $\eps$.
\item Order $e^2\lambda$: $\abs{\delta m^2_{e^2\lambda}}\le16\lambda e^2[10D_X\tau_X+2X_X(0)\int\hat G^2\hat p^2]$, and $\int\hat G^2\hat p^2\le G_X(0)$. Times $\eps$, this is at most $16\lambda e^2[10\bar\omega\bar\tau(M)+2\bar\psi(M)\bar\omega_M]$.
\end{itemize}

\emph{Order $e^4$.}
\begin{itemize}
\item The constant term and (C). Since $0\le D_X-J_X\le D_X$,
\[
\tfrac34\eps^2e^4D_X^2+\abs{(C)}\le\tfrac54e^4\eps^2D_X^2 .
\]
Times $\eps$, this is at most $\frac54e^4\eps_{\max}\bar\omega^2$.
\item (A): $\abs{(A)}\le e^4\big[10D_X\int\hat D(M^2\hat G)\hat G+2X_X(0)\int\hat D(M^2\hat G)(\hat p^2\hat G)\big]\le12e^4X_X(0)D_X$.
\item (D1): $0\le(D1)\le\frac{e^4}2\cdot3\cdot2\delta_X\cdot2G_X(0)=6e^4\delta_XG_X(0)$.
\item (D3): $\sum_\mu2G_X(\eps e_\mu)\in[0,6G_X(0)]$. The bubble $\sum_\mu\int_q\widehat{(2q+k)}_\mu^2\hat G(q)\hat G(q+k)$ lies in $[0,4G_X(0)]$, by $\sum_\mu\widehat{(2q+k)}_\mu^2\le2\hat q^2+2\widehat{(q+k)}^2$ and $\hat p^2\hat G\le1$. Hence $\abs{(D3)}\le6e^4\delta_XG_X(0)\le10e^4\delta_XG_X(0)$.
\item (D2): First,
\[
s_\mu(2P+k)=\tfrac2\eps\sin(\eps P_\mu)\cos\tfrac{\eps k_\mu}2+s_\mu(k)\cos(\eps P_\mu).
\]
Symmetrising under $P\mapsto-P$ gives
\begin{multline*}
S_\mu(k)=s_\mu(k)\!\int\!\hat G\hat D(\cdot-k)\\
+\tfrac12\!\int\!\hat G(P)\Big\{\tfrac2\eps\sin(\eps P_\mu)\cos\tfrac{\eps k_\mu}2[\hat D(P-k)-\hat D(P+k)]+s_\mu(k)\cos(\eps P_\mu)[\hat D(P-k)+\hat D(P+k)]\Big\}.
\end{multline*}
The terms containing $s_\mu(k)$ are at most $2\abs{s_\mu(k)}X_X(0)$ in absolute value.

Next, $\hat D(P-k)-\hat D(P+k)=[\widehat{(P+k)}^2-\widehat{(P-k)}^2]\hat D(P-k)\hat D(P+k)$, and
\[
\widehat{(P+k)}^2-\widehat{(P-k)}^2=\sum_\nu\tfrac4{\eps^2}\sin(\eps P_\nu)\sin(\eps k_\nu).
\]
Since $\abs{\frac2\eps\sin\eps x}\le2\abs{\hat x}$, this difference is at most $4\hat P\hat k$. Using $\sum_\mu\abs{s_\mu(k)}\abs{\hat P_\mu}\le\hat k\hat P$, $\hat P^2\hat G\le1$, and $\int\hat D(P-k)\hat D(P+k)\le\delta_X$ (Cauchy--Schwarz), we obtain
\[
\sum_\mu\abs{s_\mu(k)}\abs{S_\mu(k)}\le\hat k^2(2X_X(0)+4\delta_X).
\]
With $\hat k^2\hat G\le1$ this gives $\abs{(D2)}\le e^4(2X_X(0)+4\delta_X)D_X$.
\end{itemize}

\emph{Conclusion.} Multiply by $\eps$ and insert (Bd). The $e^4$ coefficients add up to:
\begin{itemize}
\item $\frac34+\frac12=\frac54$ for $\eps_{\max}\bar\omega^2$;
\item $6+10=16$ for $\bar\omega_M\bar\delta$;
\item $12+2=14$ for $\bar\psi(M)\bar\omega$;
\item $4$ for $\bar\delta\bar\omega$.
\end{itemize}
This is $C_4(M)$, and the lower-order terms give the rest of $C_\delta(M)$.

For the last claim, $\min_{s\ge0}(as^2+bs)\ge-b^2/(4a)$, with $a=\lambda(\eps)=\lambda\eps$ and $b=\frac12\delta m^2_{K,X}(M)\eps^2$. This gives the lower bound $-(\delta m^2_{K,X}(M))^2\eps^3/(16\lambda)\ge-C_\delta(M)^2\eps/(16\lambda)$.
\end{proof}

\paragraph{Families with $D$-lines.}

\begin{lemma}[families with $D$-lines]\label{app:K:L14-5}\label{app:G1:Dlines}
Work in organisation $(\alpha)$ at scale $n$ ($\eta=L^{-n}$), and put $c_n:=\kappa\mu_*^2L^2t_n^2$. A $D$-line is $uc_nD^\eta_n$ with $D^\eta_n=\sum_{j<n}(L^j\eta)^2G^\eta_{(j)}$.
\begin{itemize}
\item[(a)] Every connected coloured graph with an external leg and $m\ge1$ $D$-lines has superficial degree $D^\natural\ge1$.
\item[(b)] Every divergent proper subgraph of a coloured graph is $C$-only. Hence a family with an external leg and $m\ge1$ is convergent. Its own divergent subgraphs are renormalised exactly as at $u=0$, by $\delta m^2_{K,T}$ and by the Ward--Takahashi and $\varphi\varphi A$ partners, and the family itself receives no counterterm.
\item[(c)] Let $\square_0$ be a scale-$n$ unit block, and let $\mathcal M_n$ be the sum of all families with two scalar external legs and $m\ge1$ over the localisations whose first cube is $\square_0$. Then
\[
\abs{\mathcal M_n}\le C_{U1}e^2\abs u\kappa\mu_*^2L^2t_n^3\sup\norm{h\Phi_{\rm ext}}_{1,\alpha_0},
\]
where
\[
C_{U1}=O(1)\,N_F\,N_\ell\,S_0\,(1+\bar c)^{N_\ell-1}\sum_{(\alpha,\beta)}\abs e^{\alpha-2}\lambda^\beta\eps_0^{o-1}.
\]
The sum runs over $\alpha\ge2$ and $2\le\alpha+2\beta\le2\nbar$, with $o=\alpha/2+\beta$. Here:
\begin{itemize}
\item $O(1)$ is the constant of Theorem~\ref{app:G1:WP-WP1}(c);
\item $N_F$ is the number of connected uncoloured two-scalar-leg graphs of order $\le\nbar$, and $N_\ell$ is their maximal number of vector lines;
\item $\bar c:=\abs u\kappa\mu_*^2L^2\eps_0^2$, and $S_0$ is as in Lemma~\ref{app:K:L14-3}.
\end{itemize}
The same bound holds for the share of $\mathcal M_n$ anchored in any sub-block of $\square_0$.
\item[(d)] Assume the weighted-field structure of Lemma~\ref{app:K:L14-6}. We use it in the following form: in the $u$-model, the step-$i$ vector fluctuation has covariance $w_iC^{(i)}$, with $w_i=1+u\delta_i$ and $0\le\delta_i\le c_n$ for $i<n$; and it enters all later expressions only through linear maps that do not depend on $w$.

Decompose any vector line $\mathcal L$ used at scale $n$ by step of origin, $\mathcal L=\sum_i\mathcal L_i$. The line may be a composed propagator, a domain-restricted step covariance, a Dirichlet-to-free difference, a localised piece, a cut line, or a difference of such objects. Then the $u$-model line is $\sum_iw_i\mathcal L_i$, and its $D$-part is $u\sum_i\delta_i\mathcal L_i$. Every per-index bound used for $\mathcal L_i$ holds for the $D$-part with the extra factor $\abs u\delta_i\le\abs uc_n$.

Consequently, take a localised family of order $\le\nbar$ with at most $N_\ell$ vector lines. Its $u$-model value is its $u=0$ value plus a sum of coloured families with $m\ge1$. That sum is bounded by the right side of \cite[(1.33)]{Ba3}, with the constant of Theorem~\ref{app:G1:WP-WP1}(c), times
\[
(1+\abs uc_n)^{N_\ell}-1\le N_\ell\abs uc_n(1+\abs uc_n)^{N_\ell-1}.
\]
This statement concerns lines only; it does not cover non-perturbative corrections to the precision of the fluctuation integral.
\end{itemize}
\end{lemma}

\begin{proof}
(a) and the first claim of (b) are Lemma~\ref{app:G1:WP-deg}(c),(d). A proper subgraph of a connected graph has an external leg, so by (a) every subgraph containing a $D$-line is convergent. Renormalisation \cite[pp.~434--445]{Ba3} acts only on divergent subgraphs. These are $C$-only, and their partners are present by Lemma~\ref{app:G1:WP-fam}.

(c) By Theorem~\ref{app:G1:WP-WP1}(c), each localised coloured family with $m$ $D$-lines and order $(\alpha,\beta)$ is bounded by
\[
O(1)(\abs uc_n)^m\abs e^\alpha\lambda^\beta t_n^{o}\mathrm{e}^{-\delta_0d(\Lambda)}\norm{h\Phi_{\rm ext}}_{1,\alpha_0}.
\]
With the first cube fixed, the localisation count in the proof of Lemma~\ref{app:K:L14-3}(b) gives $\sum_\Lambda\mathrm{e}^{-\delta_0d(\Lambda)}\le S_0$.

Families are disjoint non-empty sets of coloured graphs. Hence the number of coloured families with exactly $m$ $D$-lines is at most $N_F\binom{N_\ell}m$.

A $D$-line is a vector line, so $\alpha\ge2$; also $o\ge1$ and $t_n\le\eps_0$. Therefore
\[
\abs e^\alpha\lambda^\beta t_n^o\le e^2t_n\,\abs e^{\alpha-2}\lambda^\beta\eps_0^{o-1},
\]
and the last factor is at most the sum over $(\alpha,\beta)$ in $C_{U1}$. Since $\abs uc_n\le\bar c$,
\[
\sum_{m\ge1}\binom{N_\ell}m(\abs uc_n)^m=(1+\abs uc_n)^{N_\ell}-1\le N_\ell\abs uc_n(1+\bar c)^{N_\ell-1}.
\]
Finally $e^2t_n\cdot\abs u\kappa\mu_*^2L^2t_n^2=e^2\abs u\kappa\mu_*^2L^2t_n^3$. A sub-block anchors a subset of the localisations, and every bound above is a sum of absolute values.

(d) A line is a covariance pushed forward by $w$-independent maps, so $M(w_iC^{(i)})M'^*=w_i\,MC^{(i)}M'^*$. Restriction to domains, Dirichlet-to-free replacement, localisation cut-offs and differences are linear in the covariance and act index by index. Hence $\mathcal L(u)=\sum_iw_i\mathcal L_i$ and $\mathcal L(u)-\mathcal L=u\sum_i\delta_i\mathcal L_i$. The kernel bounds for vector lines are per index \cite[(2.6), (2.10)--(2.12), pp.~424--426]{Ba3}, so they transfer to the $D$-part with the factor $\abs u\delta_i\le\abs uc_n$.

Expand the lines multilinearly into $C$- and $D$-parts. The colouring with $m=0$ reproduces the $u=0$ family together with its counterterm partners. By Lemma~\ref{app:G1:WP-fam} and Theorem~\ref{app:G1:WP-WP1}(b), the colourings with $m\ge1$ are coloured families: their $D$-containing subgraphs are convergent, and their $C$-only divergent subgraphs keep their partners. Each such family is bounded by Theorem~\ref{app:G1:WP-WP1}(c), with the factor $(\abs uc_n)^m$.

There are at most $\binom{N_\ell}m$ colourings with $m$ $D$-parts. Summing, and using $(1+x)^N-1\le Nx(1+x)^{N-1}$ for $x\ge0$ (mean value theorem), gives the claim.
\end{proof}

\paragraph{The anchor blocks.}

\begin{lemma}[weighted external field]\label{app:K:L14-7}\label{app:G1:FAext}
At zero background, for two-variable lattice fields $g$ put
\[
\norm g_{1,\alpha_0}:=\sup\abs g+\sup\abs{\partial g}+\sup_{y\ne y'}\abs{\partial g(y)-\partial g(y')}/\abs{y-y'}^{\alpha_0}.
\]
Here $\partial$ ranges over the lattice differences in either variable, in the units of the scale at hand. Suppose $f$ depends on the first variable only and $\abs{\partial^\gamma f}\le c$ for $\abs\gamma\le2$. Then $\norm{h(f\otimes1)gh'}_{1,\alpha_0}\le C_Lc\norm{hgh'}_{1,\alpha_0}$, with $C_L=2+8\sqrt d$.
\end{lemma}

\begin{proof}
Put $F=f\otimes1$ and $N:=\norm{hgh'}_{1,\alpha_0}$. Since multiplication operators commute, $h(Fg)h'=F(hgh')$, so we may replace $g$ by $hgh'$.

\emph{Sup terms.} By the discrete Leibniz rule,
\[
\partial(Fg)(y)=(\partial F)(y)g(y+e)+F(y)\partial g(y),
\]
and $\partial F=0$ for differences in the second variable. Hence $\sup\abs{Fg}+\sup\abs{\partial(Fg)}\le2cN$.

\emph{H\"older term.} Write $\partial(Fg)(y)-\partial(Fg)(y')$ as the sum of four differences:
\[
[\partial F(y)-\partial F(y')]g(y+e)+\partial F(y')[g(y+e)-g(y'+e)]+[F(y)-F(y')]\partial g(y)+F(y')[\partial g(y)-\partial g(y')].
\]
Put $r=\abs{y-y'}$. Since $F$ varies only in the first variable, joining the first variables by a lattice path of $\ell^1$-length at most $\sqrt dr$ gives
\[
\abs{\partial F(y)-\partial F(y')}\le\min(\sqrt dcr,2c)\le2\sqrt dc\,r^{\alpha_0}.
\]
To see the last inequality, consider $r\le1$ and $r\ge1$ separately. The same bound holds for $F$. For the field $g$ of two variables,
\[
\abs{g(z)-g(z')}\le\min(\sqrt{2d}\sup\abs{\partial g}\abs{z-z'},2\sup\abs g)\le2\sqrt{2d}\,N\abs{z-z'}^{\alpha_0}.
\]
The four terms are therefore at most $2\sqrt dcN$, $2\sqrt{2d}cN$, $2\sqrt dcN$ and $cN$, each times $r^{\alpha_0}$. So the H\"older quotient is at most $\big(1+(4+2\sqrt2)\sqrt d\big)cN$.

Adding the sup terms, $\norm{F(hgh')}_{1,\alpha_0}\le\big(3+(4+2\sqrt2)\sqrt d\big)cN\le(2+8\sqrt d)cN$ for $d\ge1$.
\end{proof}

\begin{lemma}[collar weights: bound per anchor block]\label{app:K:L14-8}\label{app:G1:FAanchor}
Let $0\le j^*<K$, and let $F$ be the set of terms of $E_{K,j^*}(u)-E_{K,j^*+1}(u)$ with a common top line $l$. By Lemma~\ref{app:G1:cut}, cutting $l$ turns the members of $F$ into one renormalised two-leg class $G'_{\rm ren}$. Its two legs are labelled once for the whole class, with leg~1 at the vertex $v_1$, the anchor end. The legs are transported by the partner operations, and ties are broken by one rule compatible with the class.

Put $\Psi_F[\Phi]:=\sum_\Lambda E(G'_{\rm ren},\Lambda,\Phi)$ for two-variable fields $\Phi$, so that the value of $F$ is $\Psi_F[K_l]$. Here $K_l=(uc_{j^*})^{m_l}\tilde K_l$, with $m_l\in\{0,1\}$ and $c_{j^*}=\kappa\mu_*^2L^2t_{j^*}^2$; $\tilde K_l$ is the index-$j^*$ kernel, and $\norm{h\tilde K_lh'}_{1,\alpha_0}\le C_K$. Let $F$ have orders $(\alpha,\beta)$ and $m$ $D$-lines, $l$ included.

Let $\{\chi_a\}_{a\in\mathfrak B_{j^*}}$ be a partition of unity over the scale-$j^*$ blocks, with $\supp\chi_a$ within $t_{j^*}/2$ of $a$ and $\abs{\partial^\gamma\chi_a}\le c_\chi t_{j^*}^{-\abs\gamma}$ for $\abs\gamma\le2$. Let the localising function $h$ of \cite[p.~420]{Ba3} vanish outside the $c_h$-neighbourhood of $\square(v_1)$, and let the localisation cubes be as in Lemma~\ref{app:K:L14-3}. Then for every $a\in\mathfrak B_{j^*}$
\[
\big|\Psi_F[(\chi_a\otimes1)K_l]\big|\le C_A^{\rm FA}A_{j^*}(F),\qquad A_{j^*}(F):=\abs{e(t_{j^*})}^\alpha\lambda(t_{j^*})^\beta(\abs u\kappa\mu_*^2L^2t_{j^*}^2)^m,
\]
where
\[
C_A^{\rm FA}:=2^{\nbar}4^{N_\ell}C_{1.33}C_Lc_\chi C_K\cdot2(2c_h+5)^d\cdot S(N_s,d,\delta_0)+C_{\rm im},\qquad S(N_s,d,\delta):=\sum_{D\ge0}(c''_d(D+2)^d)^{N_s}\mathrm{e}^{-\delta D}.
\]
Here:
\begin{itemize}
\item $C_{1.33}$ is the constant of Theorem~\ref{app:G1:WP-WP1}(c), applied to the cut class with a two-variable external field \cite[p.~419]{Ba3};
\item $C_L$ is as in Lemma~\ref{app:K:L14-7};
\item $C_{\rm im}$ is the constant for the image terms of the torus lines in the proof of Proposition~\ref{app:G1:I9};
\item $c''_d$ depends only on $d$ and on the localisation cubes, and the factor $2(2c_h+5)^d$ bounds the number of possible anchor cubes $\square(v_1)$.
\end{itemize}
$C_A^{\rm FA}$ is independent of $\eps$, $K$, $j^*$, $a$, $\ell_T\ge\ell_0$ and the domains, and it is finite at $u=-\eta$ and $M=m_1$.
\end{lemma}

\begin{proof}
Work in scale-$j^*$ units and put $\Phi_a:=(\chi_a\otimes1)K_l$. By the labelling at the level of the class, the first argument of the external field is the same summation variable in every member, since partners never re-attach the anchor end. So $\Psi_F[\Phi_a]$ is the class functional of the single field $\Phi_a$.

\emph{The bridge case.} If $l$ is a bridge, both components of $G'_{\rm ren}$ vanish identically by Lemma~\ref{app:G1:cutclasses}(c), for every two-variable external field.

\emph{The bound per localisation.} Otherwise $G'_{\rm ren}$ is connected and non-vacuum, and $d(\Lambda)$ joins all its cubes, including $\square(v_1)$ and $\square(v_2)$. Its lines have index $\le j^*$. Apply Theorem~\ref{app:G1:WP-WP1}(c),(d) at scale $j^*+1$ to the sub-sum with largest index $j^*$, and rewrite in scale-$j^*$ units. Since $e(t_{j^*+1})^\alpha\lambda(t_{j^*+1})^\beta=2^oe(t_{j^*})^\alpha\lambda(t_{j^*})^\beta$ and $c_{j^*+1}=4c_{j^*}$, this gives
\[
\abs{E(G'_{\rm ren},\Lambda,\Phi_a)}\le2^{o}4^{m-m_l}C_{1.33}\abs{e(t_{j^*})}^\alpha\lambda(t_{j^*})^\beta(\abs uc_{j^*})^{m-m_l}\mathrm{e}^{-\delta_0d(\Lambda)}\norm{h\Phi_ah'}_{1,\alpha_0}.
\]
Moreover $2^o4^{m-m_l}\le2^{\nbar}4^{N_\ell}$.

\emph{The external field.} In scale-$j^*$ units $\abs{\partial^\gamma\chi_a}\le c_\chi$ for $\abs\gamma\le2$. So Lemma~\ref{app:K:L14-7} gives
\[
\norm{h\Phi_ah'}_{1,\alpha_0}\le C_Lc_\chi(\abs uc_{j^*})^{m_l}C_K .
\]
Since $c_{j^*}=\kappa\mu_*^2L^2t_{j^*}^2$, the coupling factors and the weight factors combine into $A_{j^*}(F)$.

\emph{The anchor cubes.} $h\Phi_a=0$ unless $\dist_\infty(\square(v_1),a)\le c_h+\frac12$. For each side $s\in\{1,2\}$, at most $2c_h+3+s$ unit-lattice positions per coordinate satisfy this. So there are at most $2(2c_h+5)^d$ possible cubes $\square(v_1)$.

\emph{The other cubes.} Given $\square(v_1)$, every other cube of $\Lambda$ lies within $d(\Lambda)+c_d$ of it. Hence
\[
\#\{\Lambda:d(\Lambda)\in[D,D+1)\}\le(c''_d(D+2)^d)^{N_s},\qquad \sum_\Lambda\mathrm{e}^{-\delta_0d(\Lambda)}\le S(N_s,d,\delta_0).
\]
There is no factor for the choice of the anchor vertex, because $F$ fixes it.

\emph{Conclusion.} Multiplying these factors bounds the part without image terms by the first summand of $C_A^{\rm FA}$ times $A_{j^*}(F)$. The image terms of the torus lines are bounded as in the proof of Proposition~\ref{app:G1:I9}; this gives the summand $C_{\rm im}A_{j^*}(F)$. None of the factors depends on $\eps$, $K$, $j^*$, $a$, $\ell_T\ge\ell_0$ or the domains, and all of them are finite at $u=-\eta$ and $M=m_1$.
\end{proof}

\subsection{Ba{\l}aban's Lemma~2.7 and the normalisations $Z^{(j)}$}\label{app:G1:L27Z}

This subsection proves two statements used in the proof of Lemma~\ref{app:G1:S05-SC}. The first is the smallness of the scalar fluctuation \cite[Lemma~2.7]{Ba2}. The second is the expansion of the normalisation factors $Z^{(j)}$ of \cite[(2.100)--(2.101)]{Ba2}. Here $j$ and $k$ are scale indices. Unit blocks are written $\square$, and $B(x)$ is the unit block of the unit lattice point $x$. For a finite set $\mathcal X$ and fields with values in $\R^n$ we put $\norm V_{\ell^1}:=\sum_{x,y\in\mathcal X}\norm{V(x,y)}$.

\begin{lemma}[smallness of the scalar fluctuation; {\cite[Lemma~2.7]{Ba2}}]\label{app:K:L13-5}\label{app:G1:L27}
Work in scale-$k$ units: $t=t_k$, $\eta=L^{-k}$, $e_t=\abs e\,t^{1/2}$, $\lambda(t)=\lambda t$, $p=p(t)$, $r=r(t)$. Put $\Omega=B^k(\Lambda_2^{(k)})$ and $\hat B=B^{(k+1),\eta}$. Write $Q_k=Q_k(\hat B)$ and $Q=Q(\hat B)$, so that $Q_kQ_k^*=I$ and $QQ^*=I$. Let
\[
G_k=\big(-\Delta_{\hat B,\Omega}+m^2t^2+a_kQ_k^*Q_k\big)^{-1}\ \text{(Neumann on $\partial\Omega$)},\quad \Delta=a_k-a_k^2Q_kG_kQ_k^*,\quad \mathsf M=\Delta+aL^{-2}Q^*Q,
\]
and $C=(\Lambda_4^{(k)}\mathsf M\Lambda_4^{(k)})^{-1}$. The sets $\Lambda_5\subset\Lambda_4\subset\Lambda_2$ (superscripts $(k)$ dropped) are unions of blocks of scale $k+1$, with $\dist(\Lambda_5,\Lambda_4^c)>r$ and $\dist(\Lambda_4,\Lambda_2^c)>2r$ \cite[(2.8)]{Ba2}. Blocks of scale $k+1$ are written $y$, and $y(z)$, $\square(z)$ are the blocks containing $z$.

Let $\psi$ be a field on the blocks $y\subset\Lambda_2$. Put $u=Q_k^*Q^*\psi$ and $g=Q^*\psi=Q_ku$. Thus $u(z)=U(\hat B(\Gamma_z))\psi(y(z))$, where $\Gamma_z$ is the composite contour from the base point of $y(z)$ to $z$. For a fine bond $b=(b_-,b_+)\subset\Omega$ put $d_bf=f(b_+)-U(\hat B_b)f(b_-)$ and $\gamma_b=\Gamma_{b_-}+b-\Gamma_{b_+}$. Let $\tilde c(\delta)=\sum_{z\in\Z^d}\mathrm{e}^{-\delta\abs z}$ and $S(\delta)=\sup_{x'}\sum_{\square}\mathrm{e}^{-\delta\dist(\square,B(x'))}$. Assume:
\begin{itemize}
\item[(F1)] $\hat B$ is real and regular on $\Omega$ in the sense of \cite[Prop.~2.1]{Ba1}, and $\abs{\partial^\eta_\mu\hat B_\nu}\le c_Bp$ on $\Omega$.
\item[(F2)] Let $f$ be a function on $\Omega$ and $x'\in\Lambda_2$, and let $z\in\Omega$ and the fine bond $b\subset\Omega$ be at distance $\ge R_0$ from $\Omega^c$. Then $\abs{(G_kf)(z)}\le c_0\mathrm{e}^{-\delta_0\dist(z,\supp f)}\norm f_\infty$ and $\abs{\eta^{-1}d_b(G_kQ_k^*\delta_{x'})}\le c_0\mathrm{e}^{-\delta_0\dist(b,B(x'))}$. Moreover $\abs{C(x,x')}\le c_0\mathrm{e}^{-\delta_0\abs{x-x'}}$ for $x,x'\in\Lambda_4$.
\item[(F3)] $\abs{\psi}\le2\lambda(t)^{-1/4}p$ on $\Lambda_2$. If $y(b_-)\ne y(b_+)$, then $\abs{\psi(y(b_+))-U(\hat B(\gamma_b))\psi(y(b_-))}\le c_\psi p$. Moreover $\abs{\varphi(x)-(Q^*\psi)(x)}\le c_\varphi p$ for $x\in\Lambda_5$.
\item[(F4)] If $y(b_-)=y(b_+)$, then $\gamma_b$ is closed and bounds a surface of fine plaquettes of area $\le c_\Gamma L^2$. The area is $\le c_\Gamma\eta$ if $\square(b_-)=\square(b_+)$.
\item[(F5)] $r\ge2R_0+2$.
\end{itemize}
Then for $x\in\Lambda_5$
\[
aL^{-2}(CQ^*\psi)(x)=(Q^*\psi)(x)+R(x),\qquad\abs{R(x)}\le\big(K_0+\omega(t)\big)p,
\]
where $K_0=2d\,\mathrm{e}^{\delta_0}a_kc_0^2\,\tilde c(\delta_0)S(\delta_0)\,c_\psi$ and
\[
\omega(t)=C_1\lambda(t)^{-1/4}\big[\mathrm{e}^{-\delta_0(r-1-\sqrt d)/2}+m^2t^2+e_tp\big],
\]
with $C_1$ depending only on $d,a,L,c_0,\delta_0,c_B,c_\Gamma$. If $\omega(t)\le1$, then $\varphi'=\varphi-aL^{-2}CQ^*\psi$ satisfies $\abs{\varphi'(x)}\le(c_\varphi+K_0+1)p$ on $\Lambda_5$.
\end{lemma}

\begin{remark}\label{app:G1:L27-rem}
In the application, (F1) is \cite[Lemma~2.3, (2.60), p.~571]{Ba2} at the step $k+1$. (F2) is \cite[(1.10), (1.16)]{Ba4} together with \cite[(2.58)]{Ba2}, for $e_t$ small. (F3) is not literally \cite[(2.16)]{Ba2}: there the transports are by the block field. Replacing them by transports by $\hat B$ along the composite contours changes the left-hand sides by Stokes terms, and (F3) is assumed in the form stated. Write $\ell(t)=1+\log t^{-1}$, so that $p=b_0\ell(t)^p$ and $r=R\,\ell(t)^{\varrho}$ with $\varrho>1$. Since $e_t\lambda(t)^{-1/4}p=\abs e\lambda^{-1/4}t^{1/4}p$ and $\lambda(t)^{-1/4}m^2t^2=m^2\lambda^{-1/4}t^{7/4}$, we have $\omega(t)\to0$ as $t\to0$.
\end{remark}

\begin{proof}[Proof of Lemma~\ref{app:K:L13-5}]
\emph{Step 1 (exact decomposition).} Since $QQ^*=I$ we have $Q^*Qg=g$, and $\Lambda_4$ commutes with $Q^*Q$. Hence $\Lambda_4\mathsf M\Lambda_4g=\Lambda_4\Delta g-\Lambda_4\Delta\Lambda_4^cg+aL^{-2}\Lambda_4g$. Applying $C$ and using $C\Lambda_4\mathsf M\Lambda_4=\Lambda_4$ gives
\[
aL^{-2}C\Lambda_4g=\Lambda_4g+C\Lambda_4\Delta\Lambda_4^cg-C\Lambda_4\Delta g .
\]
So $R(x)=(C\Lambda_4\Delta\Lambda_4^cg)(x)-(C\Lambda_4\Delta g)(x)$ for $x\in\Lambda_5$.

\emph{Step 2 (collar term).} Let $x'\in\Lambda_4$ and $x''\in\Lambda_2\setminus\Lambda_4$. Then $\Delta(x',x'')=-a_k^2(Q_kG_kQ_k^*)(x',x'')$, and the points of $B(x')$ are at distance $>2r\ge R_0$ from $\Omega^c$. By (F2) with $f=Q_k^*\delta_{x''}$ we get $\abs{\Delta(x',x'')}\le a_k^2c_0\mathrm{e}^{\delta_0\sqrt d}\mathrm{e}^{-\delta_0\abs{x'-x''}}$. Since $\abs{x-x'}+\abs{x'-x''}>r$ and $\abs g\le\abs\psi$, halving the exponents gives
\[
\abs{(C\Lambda_4\Delta\Lambda_4^cg)(x)}\le a_k^2c_0^2\mathrm{e}^{\delta_0\sqrt d}\tilde c(\delta_0/2)^2\mathrm{e}^{-\delta_0r/2}\,2\lambda(t)^{-1/4}p .
\]

\emph{Step 3 (cut-off identity).} Put $\zeta(\square)=\min\{1,\max\{0,2\dist(\square,\Omega^c)/r-1\}\}$, extended to the fine lattice as a function constant on unit blocks. Then $\abs{\zeta(\square)-\zeta(\square')}\le2/r$ for neighbours, and $\zeta=1$ on $\Lambda_4$. Hence $Q_k^*(\zeta g)=\zeta u$ and $Q_k^*Q_k(\zeta u)=\zeta u$. Put $H=-\Delta_{\hat B,\Omega}+m^2t^2$, so that $G_k(H+a_kQ_k^*Q_k)=I$. Then
\[
a_k^2Q_kG_kQ_k^*(\zeta g)=a_k\zeta g-a_kQ_kG_kH(\zeta u).
\]
For $x'\in\Lambda_4$ let $w=G_kQ_k^*\delta_{x'}$, and use the pairing $\ip fh=\sum_z\eta^d\bar fh$. Since $G_k$ is self-adjoint and the Neumann form is a sum over the bonds of $\Omega$,
\[
(\Delta g)(x')=a_k\sum_{b\subset\Omega}\eta^d\,\overline{\eta^{-1}d_bw}\,\eta^{-1}d_b(\zeta u)+a_km^2t^2\ip w{\zeta u}-a_k^2\big(Q_kG_kQ_k^*((1-\zeta)g)\big)(x'),
\]
where $d_b(\zeta u)=\zeta(b_+)d_bu+(\zeta(b_+)-\zeta(b_-))U(\hat B_b)u(b_-)$. Every bond with $d_b(\zeta u)\ne0$, and every $z$ with $\zeta(z)\ne0$, lies at distance $\ge r/2-1\ge R_0$ from $\Omega^c$ by (F5). So (F2) applies at all of them.

\emph{Step 4 (the bond sums).} The transports are abelian, so $\abs{d_bu}=\abs{\psi(y(b_+))-U(\hat B(\gamma_b))\psi(y(b_-))}$. If $y(b_-)=y(b_+)$, Stokes' theorem, (F1), (F4) and $\abs{1-\mathrm{e}^{i\theta}}\le\abs\theta$ give $\abs{d_bu}\le2c_B\,e_t\,p\,\mathrm{area}(\gamma_b)\abs\psi$. A unit face is crossed by $\eta^{1-d}$ fine bonds, and a unit block contains at most $d\eta^{-d}$ fine bonds. For $b_-\in\square$ we have $\mathrm{e}^{-\delta_0\dist(b,B(x'))}\le\mathrm{e}^{\delta_0}\mathrm{e}^{-\delta_0\dist(\square,B(x'))}$. Together with (F2) and (F3) this bounds the contributions as follows.
\begin{itemize}
\item Bonds between different blocks $y$: $\le2dc_0\mathrm{e}^{\delta_0}S(\delta_0)c_\psi p$.
\item Bonds crossing unit faces inside one block $y$: $\le2dc_0\mathrm{e}^{\delta_0}S(\delta_0)\,2c_Bc_\Gamma L^2e_tp\cdot2\lambda(t)^{-1/4}p$.
\item Bonds inside a unit block: $\le dc_0\mathrm{e}^{\delta_0}S(\delta_0)\,2c_Bc_\Gamma e_tp\cdot2\lambda(t)^{-1/4}p$. The factor $\eta^{-1}$ of the difference is cancelled by the area bound $c_\Gamma\eta$.
\item Jumps of $\zeta$: they occur only at unit blocks with $\dist(\square,\Omega^c)\in[r/2-1,r+1]$, hence with $\dist(\square,B(x'))\ge r-1-\sqrt d$. They give $\le8dr^{-1}c_0\mathrm{e}^{\delta_0}S(\delta_0/2)\mathrm{e}^{-\delta_0(r-1-\sqrt d)/2}\lambda(t)^{-1/4}p$.
\item The mass term: $\abs{w}\le c_0\mathrm{e}^{-\delta_0\dist(\cdot,B(x'))}$ where $\zeta\ne0$, so this term is $\le m^2t^2c_0S(\delta_0)\,2\lambda(t)^{-1/4}p$.
\item The tail term: $\supp Q_k^*((1-\zeta)g)$ is at distance $\ge r-\sqrt d$ from $B(x')$, so this term is $\le a_k^2c_0\mathrm{e}^{-\delta_0(r-\sqrt d)}\,2\lambda(t)^{-1/4}p$.
\end{itemize}
Hence $\sup_{x'\in\Lambda_4}\abs{(\Delta g)(x')}\le a_k\big[2dc_0\mathrm{e}^{\delta_0}S(\delta_0)c_\psi+\omega_1(t)\big]p$, where $\omega_1(t)$ is a sum of terms of the three types in $\omega(t)$. Their constants depend only on $d,a,L,c_0,\delta_0,c_B,c_\Gamma$; note $r\ge2$.

\emph{Step 5 (assembly).} By (F2), $\abs{(C\Lambda_4\Delta g)(x)}\le c_0\tilde c(\delta_0)\sup_{\Lambda_4}\abs{\Delta g}$. Together with Step 2 this gives the bound on $R$, because $\mathrm{e}^{-\delta_0r/2}\le\mathrm{e}^{-\delta_0(r-1-\sqrt d)/2}$ and, since $r\ge2$, $\mathrm{e}^{-\delta_0(r-\sqrt d)}\le\mathrm{e}^{\delta_0\max\{0,\sqrt d-3\}/2}\mathrm{e}^{-\delta_0(r-1-\sqrt d)/2}$. Finally, by (F3), $\abs{\varphi'}\le\abs{\varphi-Q^*\psi}+\abs R\le(c_\varphi+K_0+\omega)p$.
\end{proof}

\paragraph{Setting for the normalisations.}
Fix $j$, write $\eta=L^{-j}$ and $e_j=e(t_j)=\abs e\,t_j^{1/2}$. Let $\Omega\subset\eta\Z^d$ be a union of unit blocks, and let $\mathcal X$ be the set of unit lattice points $x$ with $B(x)\subset\Omega$. Fields take values in $\R^n$, and $q$ is the antisymmetric charge matrix. Let $\tilde B$ and $\cAc$ be real vector fields on $\Omega$, and put $\varepsilon=\norm q\,e_j\sup_b\abs{\cAc_b}$. For $\sigma\in[0,1]$ put $U_b(\sigma)=\mathrm{e}^{\eta e_jq(\sigma\cAc_b+\tilde B_b)}$ and $(D_\sigma\phi)(b)=\eta^{-1}(U_b(\sigma)\phi(b_+)-\phi(b_-))$. Further put $(Q_j(\sigma)\phi)(y)=\sum_{x\in B(y)}\eta^dU(\sigma,\Gamma_{y,x})\phi(x)$, with $U(\sigma,\Gamma)=\mathrm{e}^{\eta e_jq\sum_{b\in\Gamma}(\sigma\cAc_b+\tilde B_b)}$. Then
\[
H(\sigma)=D_\sigma^*D_\sigma+m_j^2+a_jQ_j(\sigma)^*Q_j(\sigma),\qquad G(\sigma)=H(\sigma)^{-1}\ \text{(Neumann on $\partial\Omega$)},
\]
with a constant $m_j^2\ge0$. Put $\Delta(\sigma)=a_j-a_j^2Q_jGQ_j^*$ and $P(\sigma)=Q(\sigma)^*Q(\sigma)$, where $Q(\sigma)$ is the next-scale averaging on $\mathcal X$, with transports of the same form. Put
\[
W(\sigma)=\Delta(0)+aL^{-2}P(0)-\Delta(\sigma)-aL^{-2}P(\sigma).
\]
This is the operator of \cite[(2.99)]{Ba2} at the field $\sigma\cAc+\tilde B$. All operators are real, since $\mathrm{e}^{sq}$ is orthogonal. We use:
\begin{itemize}
\item[(Z.C)] every contour of $Q_j$ or $Q$ lies in its block and has length $\eta\,\#\Gamma\le\ell_\Gamma$; put $\ell_1=\max\{1,\ell_\Gamma\}$;
\item[(Z.H)] $H(\sigma)\ge\gamma_G>0$ for $\sigma\in[0,1]$ \cite[(1.8)]{Ba4}, and $a_j\le\bar a$.
\end{itemize}
For an operator $O$ put $\norm O_{\square\square'}=\norm{1_\square O1_{\square'}}$. On bonds, $1_\square$ selects the bonds with $b_-\in\square$; on $\mathcal X$, the point $x$ is identified with $B(x)$. Let $\tilde c(\delta)=\sup_\square\sum_{\square'}\mathrm{e}^{-\delta\dist(\square,\square')}$. Inserting $\sum_\square1_\square=I$, and using (Z.H) ($H\ge D^*D$, $H\ge a_jQ_j^*Q_j$), we get
\begin{equation}\label{app:G1:Z-chain}
\norm{OO'}_{\square\square'}\le\sum_{\square''}\norm O_{\square\square''}\norm{O'}_{\square''\square'},
\end{equation}
\begin{equation}\label{app:G1:Z-energy}
\norm{DG^{1/2}}\le1,\qquad\norm{a_j^{1/2}Q_jG^{1/2}}\le1,\qquad\norm{G^{1/2}}\le\gamma_G^{-1/2}.
\end{equation}
Let $\cF=\{1,D,a_j^{1/2}Q_j\}$.

\begin{lemma}[localised decay]\label{app:G1:Z-CT}
Assume (Z.H) at $\sigma$. Put $c_d=\sqrt d\,\mathrm{e}^{\sqrt d}$, $\delta_1=\min\{1,\gamma_G^{1/2}/(5(1+a_j^{1/2})c_d)\}$ and $c_M=\max\{1,\gamma_G^{-1/2}\}+\frac15$. Then for all $\Xi,\Upsilon\in\cF$ and all unit blocks $\square,\square'$
\[
\norm{\Xi G\Upsilon^*}_{\square\square'}\le2c_M^2\mathrm{e}^{-\delta_1\dist(\square,\square')}.
\]
\end{lemma}

\begin{proof}
Let $\rho$ be $1$-Lipschitz and $0<\delta\le\delta_1$. Let $E=\mathrm{e}^{\delta\rho}$ act on sites by $\rho(x)$, on bonds by $\rho(b_-)$, and on $\mathcal X$ by $\rho(y)$. Then $EDE^{-1}=D+J_1$ and $E^{-1}DE=D+J_2$, where $(J_1\phi)(b)=\eta^{-1}(\mathrm{e}^{\delta(\rho(b_-)-\rho(b_+))}-1)U_b\phi(b_+)$. Since $\eta^{-1}(\mathrm{e}^{\delta\eta}-1)\le\delta\mathrm{e}^\delta$ and each site is $b_+$ of at most $d$ bonds, $\norm{J_i}\le c_d\delta$. Likewise $E^{\pm1}Q_jE^{\mp1}=Q_j+J_{3,4}$ with $\norm{J_{3,4}}\le\mathrm{e}^{\delta\sqrt d}-1\le c_d\delta$, by the Cauchy--Schwarz inequality on each block and $\abs{\rho(y)-\rho(x)}\le\sqrt d$ on $B(y)$. Hence $EHE^{-1}=H+V$ with
\[
V=D^*J_1+J_2^*D+J_2^*J_1+a_j(Q_j^*J_3+J_4^*Q_j+J_4^*J_3).
\]
Put $u=(1+a_j^{1/2})c_d\delta\gamma_G^{-1/2}\le\frac15$. By \eqref{app:G1:Z-energy}, $K=G^{1/2}VG^{1/2}$ obeys $\norm K\le2u+u^2\le\frac12$. Thus $EGE^{-1}=G^{1/2}(I+K)^{-1}G^{1/2}$, and
\[
E\Xi G\Upsilon^*E^{-1}=(\Xi+J_\Xi)G^{1/2}(I+K)^{-1}G^{1/2}(\Upsilon+J'_\Upsilon)^*.
\]
Here $J_\Xi,J'_\Upsilon$ are the corresponding $J_i$, or $0$ when the operator is $1$. By \eqref{app:G1:Z-energy}, $\norm{(\Xi+J_\Xi)G^{1/2}}\le c_M$, so $\norm{E\Xi G\Upsilon^*E^{-1}}\le2c_M^2$. Now take $\rho=\dist(\cdot,\square')$ and $\delta=\delta_1$. Then $E^{-1}\le\mathrm{e}^{-\delta_1\dist(\square,\square')}$ on the range of $1_\square$, and $E=1$ on the range of $1_{\square'}$. Write $1_\square\Xi G\Upsilon^*1_{\square'}=1_\square E^{-1}(E\Xi G\Upsilon^*E^{-1})E1_{\square'}$.
\end{proof}

\begin{lemma}[domain change]\label{app:G1:Z-DC}
Let $\Omega\subset\Omega_0$ be unions of unit blocks with (Z.H) on both, and write $G_\Omega,G_{\Omega_0}$ for the propagators at the same field. Let $\square,\square'\subset\Omega$ and $\rho_0=\min\{\dist(\square,\Omega^c),\dist(\square',\Omega^c)\}$. Then for $\Xi,\Upsilon\in\cF$
\[
\norm{\Xi(G_\Omega-G_{\Omega_0})\Upsilon^*}_{\square\square'}\le c_3\mathrm{e}^{-\delta_3(\dist(\square,\square')+\rho_0)},\qquad\delta_3=\delta_1/6,
\]
with $c_3$ depending only on $d,\gamma_G,\bar a$.
\end{lemma}

\begin{proof}
Put $\rho_*=6\sqrt d+6$. If $\rho_0\le\rho_*$, Lemma~\ref{app:G1:Z-CT} applied to both terms gives the claim with $c_3\ge4c_M^2\mathrm{e}^{\delta_3\rho_*}$. Let $\rho_0>\rho_*$. Let $\chi=\phi(\dist(\cdot,\Omega^c))$, where $\phi$ vanishes on $[0,\rho_0/3]$, equals $1$ on $[2\rho_0/3,\infty)$ and is linear in between. Then $\chi=1$ within distance $2\sqrt d$ of $\square\cup\square'$. Every bond and block meeting $\supp\chi$ lies in $\Omega$, so $H_\Omega\chi=H_{\Omega_0}\chi$. Hence $G_\Omega\chi-\chi G_{\Omega_0}=G_\Omega[\chi,H_{\Omega_0}]G_{\Omega_0}$. By the locality of $\Xi,\Upsilon$,
\[
1_\square\Xi(G_\Omega-G_{\Omega_0})\Upsilon^*1_{\square'}=1_\square\Xi G_\Omega[\chi,H_{\Omega_0}]G_{\Omega_0}\Upsilon^*1_{\square'}.
\]
Now $[\chi,H]=[\chi,D^*]D+D^*[\chi,D]+a_j([\chi,Q_j^*]Q_j+Q_j^*[\chi,Q_j])$, with $\norm{[\chi,D]}\le3\sqrt d/\rho_0\le1$ and $\norm{[\chi,Q_j]}\le1$. Both commutators are local and supported in the blocks $\square''$ with $\rho_0/3-2\sqrt d\le\dist(\square'',\Omega^c)\le2\rho_0/3+2\sqrt d$. So each of the four terms has the form $\Xi G\alpha^*O\beta G\Upsilon^*$ with $\alpha,\beta\in\cF$ and $O$ local of norm $\le\max\{1,a_j^{1/2}\}$. By \eqref{app:G1:Z-chain} and Lemma~\ref{app:G1:Z-CT}, the left side is bounded by a sum of $c\,\mathrm{e}^{-\delta_1s}$ over such $\square''$ and blocks $\square'''$ adjacent to $\square''$, where $s=\dist(\square,\square'')+\dist(\square''',\square')$. We have $s\ge\dist(\square,\square')-3\sqrt d$ and $s\ge2\rho_0/3-6\sqrt d$, hence $s\ge(\dist(\square,\square')+\rho_0)/3-5\sqrt d$. Use $\mathrm{e}^{-\delta_1s}\le\mathrm{e}^{-\delta_1\dist(\square,\square'')/2}\mathrm{e}^{-\delta_1s/2}$, and sum over $\square''$ (a factor $\tilde c(\delta_1/2)$) and over $\square'''$ (a factor $3^d$).
\end{proof}

\begin{proposition}[expansion of $W$]\label{app:G1:Z-W}
Assume (Z.C) and (Z.H). Let $S_1$ be the set of unit blocks within distance $2L\sqrt d$ of $\supp\cAc$. There is $C_*=C_*(d,L,a,\bar a,\gamma_G,\ell_\Gamma)$ such that, with $w=C_*\varepsilon$, every $\sigma_0\in[0,1]$ and every $q\ge1$ satisfy
\[
\norm{W_q(\sigma_0)}\le w^q,\qquad\norm{W_q(\sigma_0)}_{\ell^1}\le w^q\abs{S_1},\qquad W_q(\sigma_0):=\frac{W^{(q)}(\sigma_0)}{q!}.
\]
Consequently $W(\tau)=\sum_{q=1}^{\nbar}\tau^qW_q+W_>(\tau)$ with $W_q=W_q(0)$. For $\tau\in[0,1]$ the operator $W_>(\tau)$ obeys both bounds with $(w\tau)^{\nbar+1}$ in place of $w^q$. All these operators are real and symmetric.
\end{proposition}

\begin{proof}
\emph{Vertices.} For $s\ge1$ put $D_s=\partial_\sigma^sD/s!$. Then $(D_s\phi)(b)=\eta^{-1}\frac{(\eta e_jq\cAc_b)^s}{s!}U_b\phi(b_+)$, so $\norm{D_s}\le\sqrt d\,\varepsilon^s/s!$. The entries of $Q_{j,s}=\partial_\sigma^sQ_j/s!$ are $\eta^d\frac{(\eta e_jq\cAc(\Gamma))^s}{s!}U$, so $\norm{Q_{j,s}}\le(\ell_1\varepsilon)^s/s!$ by (Z.C) and the Cauchy--Schwarz inequality on each block; the same holds for $Q_s$. All of these vanish outside the blocks meeting $\supp\cAc$. The operators $D_s$, $a_j^{1/2}Q_{j,s}$ ($s\ge1$) are called vertices of order $s$. Each is local, supported in $S_1$, and has norm $\le c_v(\ell_1\varepsilon)^s$ with $c_v=\max\{1,\sqrt d,\bar a^{1/2}\}$.

\emph{Coefficients.} $H$ is smooth and invertible on $[0,1]$. Write $O_r=O^{(r)}(\sigma_0)/r!$, the index $0$ denoting operators at $\sigma_0$. The Leibniz rule applied to $GH=I$ gives $G_q=-\sum_{r=1}^qG_{q-r}H_rG$, so
\[
G_q=\sum_{l=1}^q(-1)^l\sum_{r_1+\dots+r_l=q,\ r_i\ge1}GH_{r_1}G\cdots H_{r_l}G,\qquad H_r=\sum_{r_1+r_2=r}\big(D_{r_1}^*D_{r_2}+a_jQ_{j,r_1}^*Q_{j,r_2}\big),
\]
\[
W_q=a_j^2\sum_{s+q'+s'=q}Q_{j,s}G_{q'}Q_{j,s'}^*-aL^{-2}\sum_{s+s'=q}Q_s^*Q_{s'}.
\]
Expanding, the first sum consists of at most $(q+1)^22^q4^q\le32^q$ terms. Each has the form $a_j(\Xi_0G\Upsilon_1^*)(\Xi_1G\Upsilon_2^*)\cdots(\Xi_lG\Upsilon_{l+1}^*)$ with $l\le q$. Each $\Xi_i,\Upsilon_i$ is in $\cF\setminus\{1\}$ or is a vertex. The vertex orders sum to $q$, so there is at least one vertex and at most $2q+2$.

\emph{Operator norm.} Put $g_*=\max\{1,\gamma_G^{-1/2}\}$ and $a_*=\max\{1,\bar a\}$. By \eqref{app:G1:Z-energy}, $\norm{\Xi G\Upsilon^*}\le\norm{\Xi G^{1/2}}\norm{G^{1/2}\Upsilon^*}$. Here $\norm{\Xi G^{1/2}}\le1$ for $\Xi\in\cF\setminus\{1\}$, and $\norm{\Xi G^{1/2}}\le\gamma_G^{-1/2}\norm\Xi$ for a vertex. So each term is at most $a_jg_*^{2q+2}c_v^{2q+2}(\ell_1\varepsilon)^q\le(a_*(g_*c_v)^4\ell_1\varepsilon)^q$.

\emph{$\ell^1$ norm.} Since $1_xQ_{j,s}=1_xQ_{j,s}1_{B(x)}$, the block $W_q(x,y)$ equals $1_xW_q1_y$. By the locality of the vertices and Lemma~\ref{app:G1:Z-CT}, a factor containing vertices obeys
\[
\norm{\Xi G\Upsilon^*}_{\square\square'}\le c_4\prod(\text{vertex norms})\,\mathrm{e}^{-\delta_1\dist(\square,\square')},\qquad c_4=9^d2c_M^2\mathrm{e}^{2\delta_1\sqrt d}.
\]
The block next to a vertex lies in $S_1$. By \eqref{app:G1:Z-chain}, $\sum_{x,y}\norm{1_x(\cdots)1_y}$ is bounded by a sum over chains of blocks. Fix the block of one vertex in $S_1$ and sum the other $l+1$ blocks, gaining $\tilde c(\delta_1)$ each. This gives at most $\abs{S_1}a_j(c_4\tilde c(\delta_1))^{q+1}c_v^{2q+2}(\ell_1\varepsilon)^q$ per term.

\emph{The $P$ term.} It is local inside $L$-blocks, and its nonzero blocks lie in $S_1$. So its operator norm is $\le aL^{-2}(2\ell_1\varepsilon)^q$, and its $\ell^1$ norm is at most $L^d\abs{S_1}$ times that.

\emph{Constant.} Put $c_5=\max\{g_*,c_4\tilde c(\delta_1)\}$ and $C_*=64\max\{a_*,L^daL^{-2}\}(c_5c_v)^4\ell_1$. Both parts are then $\le(w/2)^q$ (times $\abs{S_1}$ for the $\ell^1$ norm), so the sum obeys the claimed bounds. These hold at every $\sigma_0\in[0,1]$, by (Z.H) there and the unitarity of the transports.

\emph{Remainder.} Taylor's formula gives $W_>(\tau)=(\nbar+1)\tau^{\nbar+1}\int_0^1(1-s)^{\nbar}W_{\nbar+1}(s\tau)\,ds$. So both norms of $W_>(\tau)$ are at most $\tau^{\nbar+1}$ times the supremum over $\sigma\in[0,1]$ of the norms of $W_{\nbar+1}(\sigma)$. Finally, $W$ is real and symmetric because $\Delta$ and $P$ are, and $W(0)=0$.
\end{proof}

\begin{lemma}[determinant factors]\label{app:G1:Z-det}
Let $\mathcal Y$ be a finite set, and let $\abs{\mathcal X}>0$ be a real number (for instance the cardinality of a finite set $\mathcal X$). Let $C$ be symmetric and positive on $\bigoplus_{x\in\mathcal Y}\R^n$ with $\norm C\le\gamma_0^{-1}$. All operators below act on this space, and $\norm\cdot_{\ell^1}$ is taken over $\mathcal Y$.
\begin{itemize}
\item[(a)] Let $W$ be symmetric with $\norm W\le\omega$ and $\norm W_{\ell^1}\le\omega\abs{\mathcal X}$. Put $K=C^{1/2}WC^{1/2}$ and assume $\bar k:=\gamma_0^{-1}\omega\le\frac12$. Then $\big|-\frac12\log\det(I-K)-\sum_{l=1}^{\nbar}\frac1{2l}\tr K^l\big|\le\frac{\bar k^{\nbar+1}}{\nbar+1}\,n\abs{\mathcal X}$.
\item[(b)] Let $\tau\mapsto W(\tau)$ be smooth and symmetric on $[0,1]$, with $W(\tau)=\sum_{q=1}^{\nbar}\tau^qW_q+W_>(\tau)$. Assume that $W_q$ obeys the two bounds of (a) with $w^q$ in place of $\omega$, and $W_>(\tau)$ with $(w\tau)^{\nbar+1}$ in place of $\omega$. Assume $4w\le\gamma_0\le1$, and put $f(\tau)=-\frac12\log\det(I-C^{1/2}W(\tau)C^{1/2})$. Then for $1\le m'\le\nbar$
\[
\frac{f^{(m')}(0)}{m'!}=\sum_{l=1}^{m'}\frac1{2l}\sum_{q_1+\dots+q_l=m'}\tr(CW_{q_1}\cdots CW_{q_l}),
\]
and $\big|f(1)-\sum_{m'=1}^{\nbar}f^{(m')}(0)/m'!\big|\le C_Z\,n\,w^{\nbar+1}\abs{\mathcal X}$, where $C_Z=(2/\gamma_0)^{\nbar+1}+\sum_{l\le\nbar}\frac{(\nbar+1)^l}{2l\gamma_0^l}$.
\end{itemize}
\end{lemma}

\begin{proof}
(a) We have $\norm K\le\norm{C^{1/2}}^2\norm W\le\bar k\le\frac12$, so $I-K>0$. Every eigenvalue $\mu$ of $K$ has $\abs\mu\le\frac12$, and
\[
\Big|-\tfrac12\log(1-\mu)-\sum_{l\le\nbar}\frac{\mu^l}{2l}\Big|\le\sum_{l>\nbar}\frac{\abs\mu^l}{2l}\le\frac{\abs\mu^{\nbar+1}}{\nbar+1}.
\]
Summing over the eigenvalues gives at most $(\nbar+1)^{-1}\bar k^{\nbar}\norm K_1$. Next, $\norm K_1\le\norm C\norm W_1\le\gamma_0^{-1}n\norm W_{\ell^1}$, because an $n\times n$ block has trace norm at most $n$ times its operator norm.

(b) For $\tau\in[0,1]$ we have $\norm{W(\tau)}\le\sum_q(w\tau)^q+(w\tau)^{\nbar+1}\le2w\tau$, and likewise $\norm{W(\tau)}_{\ell^1}\le2w\tau\abs{\mathcal X}$. So (a) applies with $\bar k=2w\tau/\gamma_0\le\frac12$, and the truncation error is $\le(2w\tau/\gamma_0)^{\nbar+1}n\abs{\mathcal X}$. Now expand $\tr K^l=\tr(CW(\tau))^l$, $l\le\nbar$, into at most $(\nbar+1)^l$ traces $\tr(CV_1\cdots CV_l)$. Let $P(\tau)$ be the sum of the terms with every $V_i=W_{q_i}$ and $\sum q_i\le\nbar$; it is a polynomial of degree $\le\nbar$ with $P(0)=0$. Each remaining term carries a size factor at most $(w\tau)^{\nbar+1}$. Using cyclicity to put the factor measured in $\ell^1$ first, it is bounded by $\norm{CV_1}_1\prod_{i\ge2}\norm{CV_i}\le\gamma_0^{-l}n\abs{\mathcal X}(w\tau)^{\nbar+1}$. Hence $\abs{f(\tau)-P(\tau)}\le C_Zn(w\tau)^{\nbar+1}\abs{\mathcal X}$. Since $f$ is smooth and $f-P=O(\tau^{\nbar+1})$ as $\tau\downarrow0$, $P$ is the Taylor polynomial of $f$ of order $\nbar$. This gives both claims, the second at $\tau=1$.
\end{proof}

\begin{lemma}[re-localisation of cyclic traces]\label{app:G1:Z-R}
Let $T=\tr_{\mathcal X}(O_1\cdots O_N)$, and let $U$ be a union of unit blocks. Assume that $\norm{O_i}_{\square\square'}\le y_i\mathrm{e}^{-\delta\dist(\square,\square')}$ for all $i$. Assume that some factor $O_{i_0}$ is not replaced and satisfies $1_\square O_{i_0}1_{\square'}=0$ unless $\square\in S_1$. Let $T'$ be obtained by replacing some other factors $O_i$ by $O_i'$, where $O_i'$ has the same decay and
\[
\norm{O_i-O_i'}_{\square\square'}\le y_i\mathrm{e}^{-\delta(\dist(\square,\square')+\dist(\square,U^c))}.
\]
Then, with $\rho=\dist(S_1,U^c)$,
\[
\abs{T-T'}\le nN\Big(\prod_iy_i\Big)\tilde c(\delta/2)^{N-1}\mathrm{e}^{\delta(N+2)\sqrt d/2}\mathrm{e}^{-\delta\rho/2}\abs{S_1}.
\]
\end{lemma}

\begin{proof}
Telescope: $T-T'=\sum_i\tr(O'_1\cdots O'_{i-1}(O_i-O'_i)O_{i+1}\cdots O_N)$, the sum running over the replaced indices. Use $\abs{\tr M}\le n\norm M$ for the $n\times n$ blocks at the points of $\mathcal X$, together with \eqref{app:G1:Z-chain}. In the $i$-th term the cycle runs through blocks $\square_0,\dots,\square_{N-1}$, with $P:=\square_{i_0-1}\in S_1$. The hops $h$ along the cycle from $P$ to $\square_{i-1}$ satisfy $\sum h\ge\dist(P,\square_{i-1})-N\sqrt d$. Moreover $\dist(\square_{i-1},U^c)\ge\rho-\dist(P,\square_{i-1})-2\sqrt d$. Hence the product of the decay factors is at most $\prod\mathrm{e}^{-\delta h/2}\cdot\mathrm{e}^{\delta(N+2)\sqrt d/2}\mathrm{e}^{-\delta\rho/2}$. Fix $P$, sum the other $N-1$ blocks along the chain with the factors $\mathrm{e}^{-\delta h/2}$, bound the closing hop by $1$, and sum $P$ over $S_1$.
\end{proof}

\begin{proposition}[expansion of the normalisations $Z^{(j)}$]\label{app:K:L13-7}\label{app:G1:Zexp}
Work in the setting of Lemma~\ref{app:K:L13-6}, with $j<k$ and $m=k-j$. Let $\cAc=\theta_{k+1}A'^{(k),L^{-j}}$ in scale-$j$ units, $\Omega=B^j(\Lambda_2^{(j)})$ and $\Omega_0=B^k(\Lambda_0^{(k)})$ \cite[p.~578]{Ba2}. Let $W^{(j)}(\tau)$ be the operator $W(\tau)$ above, i.e.\ \cite[(2.99)]{Ba2} at $\tau\cAc+\tilde B$. Put $C=C^{(j)}_{\Lambda_5^{(j)}}(\Omega,\tilde B)$ and $C'=C^{(j)}_{\Lambda_5^{(j)}}(\Omega_0,\tilde B)$, and let $X_j$ be the set of scale-$j$ points of $B^k(\Lambda_2^{(k)})$, so that $\abs{X_j}=L^{dm}\abs{\Lambda_2^{(k)}}$. Let $U$ be the union of the blocks $B(x)$, $x\in\Lambda_5^{(j)}$, intersected with $\Omega\cap\Omega_0$. We use the following hypotheses, provided by Lemma~\ref{app:K:L13-6} and Proposition~\ref{app:K:L13-8}:
\begin{itemize}
\item[(Z.C), (Z.H)] as above, on $\Omega$ for $\sigma\in[0,1]$, and on $\Omega_0$ and $\Omega\cap\Omega_0$ at $\sigma=0$;
\item[(Z.A)] $\sup\abs{\cAc}\le c_Ap(t_j)$;
\item[(Z.K)] $C,C'$ are symmetric and positive, with $\norm C\le\gamma_0^{-1}$ \cite[(1.15)]{Ba4}. For $x,x'$ in their domain, $\abs{C(x,x')},\abs{C'(x,x')}\le c_0\mathrm{e}^{-\delta_0\abs{x-x'}}$ \cite[(1.16)]{Ba4}, and $\abs{(C-C')(x,x')}\le2c_0\mathrm{e}^{-\delta_0(\abs{x-x'}+\dist(x,U^c))}$ \cite[(1.18), (1.20)]{Ba4};
\item[(Z.S)] $\abs{S_1}\le c_S\abs{X_j}$ and $\rho:=\dist(S_1,U^c)\ge L^mr(t_k)-c$ (Lemma~\ref{app:K:L13-6}(iii), (iv)).
\end{itemize}
Here $t_j=L^{-m}t_k$, $t_k\le1$, $\ell(t)=1+\log t^{-1}$, $p(t)=b_0\ell(t)^p$, $r(t)=R\,\ell(t)^{\varrho}$ with $\varrho>1$, and $e(t)=\abs et^{1/2}$. Then:
\begin{itemize}
\item[(a)] $W^{(j)}(\tau)=\sum_{q=1}^{\nbar}\tau^qW_q+W_>(\tau)$ on $[0,1]$ with real symmetric $W_q$, $\norm{W_q}\le w_j^q$ and $\norm{W_q}_{\ell^1}\le c_Sw_j^q\abs{X_j}$. The same bounds hold for $W_>(\tau)$ with $(w_j\tau)^{\nbar+1}$ in place of $w_j^q$. Here $w_j=C_*\varepsilon\le C_We(t_j)p(t_j)$ with $C_W=C_*c_A\norm q$.
\item[(b)] Put $c^{(j)}_{m'}=\sum_{l=1}^{m'}\frac1{2l}\sum_{q_1+\dots+q_l=m'}\tr(CW_{q_1}\cdots CW_{q_l})$. If $4w_j\le\gamma_0\le1$, then
\[
\Big|\log\frac{Z^{(j)}(\theta_{k+1}A'+\tilde B)}{Z^{(j)}(\tilde B)}-\sum_{m'=1}^{\nbar}c^{(j)}_{m'}\Big|\le C_Zc_Snw_j^{\nbar+1}\abs{X_j}.
\]
\item[(c)] Let $c'^{(j)}_{m'}$ be $c^{(j)}_{m'}$ with every $G_j(\Omega,\tilde B)$ replaced by $G_j(\Omega_0,\tilde B)$ and $C$ by $C'$. Then $\abs{c^{(j)}_{m'}-c'^{(j)}_{m'}}\le nK(\nbar)(C'_*w_j)^{m'}\mathrm{e}^{-\delta_5\rho}\abs{S_1}$, where $\delta_5=\min\{\delta_0,\delta_3\}/4$.
\item[(d)] Assume $4w_j\le\gamma_0\le1$ for all $j<k$, $\frac{\nbar+1}2>d$ and $\delta_5r(t_k)\ge d\log L+\log2$. Then
\[
\sum_{j<k}\Big|\log\frac{Z^{(j)}(\theta_{k+1}A'+\tilde B)}{Z^{(j)}(\tilde B)}-\sum_{m'\le\nbar}c^{(j)}_{m'}\Big|\le Ct_k^{(\nbar+1)/2}\ell(t_k)^{p(\nbar+1)}\abs{\Lambda_2^{(k)}},
\]
and $\sum_{j<k}\sum_{m'\le\nbar}\abs{c^{(j)}_{m'}-c'^{(j)}_{m'}}\le C_Nt_k^N\abs{\Lambda_2^{(k)}}$ for every $N$.
\end{itemize}
The constants $C_*,C'_*,K(\nbar)$ depend only on $d,L,a,\bar a,\gamma_G,\ell_\Gamma,\nbar,c_0,\delta_0$. The constant $C$ depends in addition on $n,c_S,c_A,\norm q,\abs e,b_0,p,\gamma_0$, and $C_N$ on $N,n,c_S,R,\varrho,c$ as well.
\end{proposition}

\begin{proof}
(a) This is Proposition~\ref{app:G1:Z-W}, with $\abs{S_1}\le c_S\abs{X_j}$ by (Z.S) and $\varepsilon\le\norm q\,e(t_j)c_Ap(t_j)$ by (Z.A).

(b) By \cite[(2.100)--(2.101)]{Ba2}, $\log\big(Z^{(j)}(\tau\theta_{k+1}A'+\tilde B)/Z^{(j)}(\tilde B)\big)=f(\tau)$ with $f$ as in Lemma~\ref{app:G1:Z-det}(b). Only $\Lambda_5^{(j)}W\Lambda_5^{(j)}$ enters $C^{1/2}WC^{1/2}$, and restriction does not increase the norms. Since $\tr(C^{1/2}WC^{1/2})^l=\tr(CW)^l$, Lemma~\ref{app:G1:Z-det}(b) with $\mathcal Y=\Lambda_5^{(j)}$ and $\abs{\mathcal X}=c_S\abs{X_j}$, together with (a) and (Z.K), gives the claim.

(c) Insert into each trace the expansion of each $W_q$ from the proof of Proposition~\ref{app:G1:Z-W}. In every factor $\Xi G\Upsilon^*$ of that proof, write a vertex $\Xi$ as $\Xi\cdot1$ and a vertex $\Upsilon$ as $1\cdot\Upsilon$, and treat the vertex, or its adjoint, as a separate factor. Then $c^{(j)}_{m'}$ is a sum of at most $N(\nbar)$ terms $\frac1{2l}\tr(O_1\cdots O_N)$ with $N\le N(\nbar)$. Each $O_i$ is $C$, a factor $\Xi G\Upsilon^*$ with $\Xi,\Upsilon\in\cF$, a vertex or the adjoint of a vertex, or a $P$-vertex. Every term contains a factor of one of the last two kinds. Such a factor is not replaced, and $1_\square O_i1_{\square'}=0$ unless $\square\in S_1$. With $y_i$ taken from (Z.K), Lemma~\ref{app:G1:Z-CT} and the vertex norms, $\prod_iy_i\le(C''\varepsilon)^{m'}$. The replaced factors obey the hypothesis of Lemma~\ref{app:G1:Z-R} with $\delta=\min\{\delta_0,\delta_3\}/2$. For $C$ this is (Z.K). For the propagators, apply Lemma~\ref{app:G1:Z-DC} to the pairs $(\Omega\cap\Omega_0,\Omega)$ and $(\Omega\cap\Omega_0,\Omega_0)$. Since $U\subset\Omega\cap\Omega_0$, we have $\rho_0\ge\dist(\square,U^c)-\dist(\square,\square')-\sqrt d$. Hence
\[
\dist(\square,\square')+\rho_0\ge\tfrac12\big(\dist(\square,\square')+\dist(\square,U^c)\big)-\tfrac{\sqrt d}2 .
\]
Lemma~\ref{app:G1:Z-R} and the sum over the at most $N(\nbar)$ traces give the claim with $C'_*=C''/C_*$.

(d) By (b), the $j$-th remainder is at most $C_Znw_j^{\nbar+1}L^{dm}\abs{\Lambda_2^{(k)}}$. We have $w_j\le C_W\abs eb_0\,t_j^{1/2}\ell(t_j)^p$ and $t_j=L^{-m}t_k$. Since $\ell(t_k)\ge1$, also $\ell(t_j)=\ell(t_k)+m\log L\le\ell(t_k)(1+m\log L)$. So the sum over $j$ is at most
\[
C_Zn(C_W\abs eb_0)^{\nbar+1}t_k^{(\nbar+1)/2}\ell(t_k)^{p(\nbar+1)}\sum_{m\ge1}L^{-m((\nbar+1)/2-d)}(1+m\log L)^{p(\nbar+1)}\abs{\Lambda_2^{(k)}}.
\]
The series converges since $\frac{\nbar+1}2>d$.

For the re-localisation, (c) gives at scale $j$ at most $\nbar nK(\nbar)\max\{1,C'_*\}^{\nbar}\mathrm{e}^{\delta_5c}\mathrm{e}^{-\delta_5L^mr(t_k)}L^{dm}\abs{\Lambda_2^{(k)}}$. Here we used $w_j\le1$, $\abs{S_1}\le\abs{X_j}$ and (Z.S). Since $L^m\ge1+m$ and $\delta_5r(t_k)\ge d\log L+\log2$,
\[
\sum_{m\ge1}L^{dm}\mathrm{e}^{-\delta_5L^mr(t_k)}\le\mathrm{e}^{-\delta_5r(t_k)}\sum_{m\ge1}\mathrm{e}^{m(d\log L-\delta_5r(t_k))}\le\mathrm{e}^{-\delta_5r(t_k)}.
\]
Since $\varrho>1$, the number $C'_N:=\exp\sup_{x\ge1}(Nx-\delta_5Rx^{\varrho})$ is finite. Then $\mathrm{e}^{-\delta_5r(t)}=\mathrm{e}^{-N\ell(t)}\mathrm{e}^{N\ell(t)-\delta_5R\ell(t)^{\varrho}}\le C'_N\mathrm{e}^{-N}t^N$.
\end{proof}

\subsection{Scalar kernels on the torus, Proposition~3.1 of \cite{Ba2} at a real background, and the background as an external field}\label{app:G1:part-SK}

\paragraph{Conventions.} Fix $k\ge1$ and work in scale-$k$ units. The fine lattice is the torus $T_h$ with spacing $h:=L^{-k}$. Since $L=2$, we have $h\le\frac12$. The unit blocks are $B(y)$, $y\in T^{(k)}$, and the big blocks have side $M_b$. We put $c_M:=2\sqrt dM_b$. Distances $\abs{x-y}$ are the $\ell^\infty$ torus distances of \cite[(1.3)]{Ba1}; this is the metric of \cite[(2.8), (2.44)]{Ba2}. Every point of $B(y)$ lies within distance $\sqrt d$ of $y$.

Norms are $\ell^2(h^d)$ norms on fine functions and counting norms on $T^{(k)}$. We put $t:=t_k$, $e_k:=et^{1/2}$, $\ell(t):=1+\log t^{-1}$, $r(t)=R\ell(t)^2$ and $p(t)=b_0\ell(t)^6$.

For a union of blocks $\Omega$ and a real vector field $A$ on its bonds, let
\[
q_{A,\Omega}(\psi)=\sum_{b\subset\Omega}h^{d-2}\abs{\psi(b_-)-U(A_b)\psi(b_+)}^2
\]
be the form of the Neumann operator $-\Delta^N_{A,\Omega}$, and put
\[
X(\Omega,A)=-\Delta^N_{A,\Omega}+m^2t^2+a_kP_k(A),\qquad G_k=X^{-1},\qquad\Delta^{(k)}(\Omega,A)=a_k-a_k^2Q_kG_kQ_k^*.
\]
Here $Q_k=Q_k(A)$ is the averaging operator of \cite[(1.4)]{Ba4}, so that $Q_kQ_k^*=I$, and $P_k=Q_k^*Q_k$. The operator $Q_k$ averages over $B(y)$ with unitary parallel transporters. Hence the kernel of $P_k$ vanishes unless $x,x'$ lie in one block, it has modulus at most $1$, and $P_k=P_k^*\ge0$.

Two standing assumptions are used:
\begin{itemize}
\item[(N0)] $-\Delta^N_{A,\Omega}+aP_k(A)\ge\gamma_0$ for every union of big blocks $\Omega$ and every $A$ that is regular in the sense of \cite[(1.7)]{Ba4}. This is \cite[(1.8), p.~573]{Ba4}.
\item[(NW)] The side of $T^{(k)}$ exceeds $3r(t)$. For $t_k\le\eps_0$ this holds under (T.nw) and (E.r), by Lemma~\ref{app:G1:nowrap}.
\end{itemize}
We assume $m^2t^2\le1$. Since $a(1-L^{-2})\le a_k\le a$ and $-\Delta^N\ge0$, (N0) gives
\[
X(\Omega,A)\ge\frac{a_k}{a}\big(-\Delta^N_{A,\Omega}+aP_k\big)\ge\gamma:=(1-L^{-2})\gamma_0 .
\]
We fix $0<\delta\le1$, depending only on $(d,a,\gamma_0)$, such that $c(\delta):=d\delta^2\cosh\delta+a(\mathrm{e}^{\sqrt d\delta}-1)\le\gamma/2$.

\begin{lemma}[weighted resolvent bound]\label{app:G1:SK-CT}
Let $X=q+v+a'P$ on $\ell^2(\Omega,h^d)$, $h\le\frac12$, where:
\begin{itemize}
\item $q$ is the form of a set of nearest-neighbour bonds with unitary transporters, and $v\ge0$ is a multiplication operator;
\item $0<a'\le a$ and $P=P^*\ge0$, where $P(x,x')=0$ unless $x,x'$ lie in one block, and $\abs{P(x,x')}\le1$;
\item $X\ge\gamma$.
\end{itemize}
Let $\rho$ be $1$-Lipschitz. Then $w:=\mathrm{e}^{\delta\rho}X^{-1}\mathrm{e}^{-\delta\rho}F$ satisfies $\norm w\le2\gamma^{-1}\norm F$ and $q(w)\le4\gamma^{-1}\norm F^2$.
\end{lemma}
\begin{proof}
We have $X^\rho w=F$ with $X^\rho:=\mathrm{e}^{\delta\rho}X\mathrm{e}^{-\delta\rho}$.

\emph{Bond terms.} Take a bond $b=(x,x')$ and put $s:=\delta(\rho(x)-\rho(x'))$, so $\abs s\le\delta h$. The bond-$b$ part of $\rRe\ip w{X^\rho w}$ is
\[
h^{d-2}\big[\abs{w_x}^2+\abs{w_{x'}}^2-2\cosh s\,\rRe(\bar w_xU_bw_{x'})\big]\ge q_b(w)-(\cosh s-1)h^{d-2}\big(\abs{w_x}^2+\abs{w_{x'}}^2\big).
\]
We have $\cosh s-1\le\frac12\delta^2h^2\cosh\delta$, and every site lies in at most $2d$ bonds. Hence the bond terms sum to at least $q(w)-d\delta^2\cosh\delta\norm w^2$.

\emph{The other terms.} The operator $v$ commutes with $\mathrm{e}^{\delta\rho}$. The kernel of $P^\rho-P$ is $(\mathrm{e}^{\delta(\rho(x)-\rho(x'))}-1)P(x,x')$, with $\abs{\rho(x)-\rho(x')}\le\sqrt d$ inside a block, and every block has volume $1$. So the row and column sums of $\abs{P^\rho-P}$ are at most $\mathrm{e}^{\sqrt d\delta}-1$, and therefore $\norm{P^\rho-P}\le \mathrm{e}^{\sqrt d\delta}-1$.

\emph{Conclusion.} Combining,
\[
\rRe\ip w{X^\rho w}\ge\ip w{Xw}-c(\delta)\norm w^2\ge\tfrac\gamma2\norm w^2 .
\]
Since $\rRe\ip w{X^\rho w}=\rRe\ip wF\le\norm w\norm F$, the first bound follows. Since $v,P\ge0$,
\[
q(w)\le\ip w{Xw}\le\norm w\norm F+c(\delta)\norm w^2\le4\gamma^{-1}\norm F^2. \qedhere
\]
\end{proof}

\begin{lemma}[locality]\label{app:G1:SK-loc}
Let $X_i$ be as in Lemma~\ref{app:G1:SK-CT} on $\Omega_i$, $i=1,2$. Let $S\subset\Omega_1\cap\Omega_2$ be a union of blocks with the following properties:
\begin{itemize}
\item every bond of $\Omega_1$ or $\Omega_2$ with an endpoint in $S$ is a bond of both, with the same transporter;
\item $v_1=v_2$ on $S$, and $P_1=P_2$ on the blocks of $S$.
\end{itemize}
Put $D(y):=\dist(y,(\Omega_1\cup\Omega_2)\setminus S)$. Let $D(y),D(y')\ge3+\sqrt d$, and let $f,g$ be supported in $B(y),B(y')$ with $\norm f,\norm g\le1$. Then
\[
\abs{\ip f{(X_1^{-1}-X_2^{-1})g}}\le C\,\mathrm{e}^{-\frac\delta2(\abs{y-y'}+D(y)+D(y'))},\qquad C=C(d,a,\gamma).
\]
\end{lemma}
\begin{proof}
\emph{The identity.} Let $\chi:=\min\{1,\frac12\dist(\cdot,(\Omega_1\cup\Omega_2)\setminus S)\}$. Then $\chi$ is real and $\frac12$-Lipschitz, and $\chi=0$ off $S$. Moreover $\chi=1$ on $B(y)\cup B(y')$, because points there are at distance $\ge D-\sqrt d\ge2$ from the complement.

Put $u=X_1^{-1}g$. The function $\chi u$ is supported in $S$. By the hypotheses, $X_1(\chi u)$ and $X_2(\chi u)$ coincide, both being supported in $S$ and its bond neighbours. Since $\chi g=g$,
\[
X_2(\chi u)=X_1(\chi u)=g+[X_1,\chi]u .
\]
We also have $\ip f{\chi u}=\ip fu$. Using the self-adjointness of $X_2$ we get
\[
\ip f{(X_1^{-1}-X_2^{-1})g}=\ip{\alpha}{[X_1,\chi]\beta},\qquad\alpha=X_2^{-1}f,\quad\beta=X_1^{-1}g .
\]

\emph{The commutator.} Write $d_bc=c(x)-U_bc(x')$. Then $d_b(\chi c)=\chi(x)d_bc+(\chi(x)-\chi(x'))U_bc(x')$, with $\abs{\chi(x)-\chi(x')}\le h/2$. Next, $[P,\chi](x,x')=P(x,x')(\chi(x')-\chi(x))$, with $\abs{\chi(x')-\chi(x)}\le\sqrt d/2$ inside a block. The operator $v$ commutes with $\chi$. Let $K$ be the set of sites within one bond or one block of $\{0<\chi<1\}$. Then, by the Cauchy--Schwarz inequality on each block,
\[
\abs{\ip\alpha{[X_1,\chi]\beta}}\le\sum_{b\cap K\ne\emptyset}h^{d-1}\big(\abs{d_b\alpha}\abs{\beta(x')}+\abs{\alpha(x')}\abs{d_b\beta}\big)+a\sqrt d\sum_{B\cap K\ne\emptyset}\norm\alpha_B\norm\beta_B .
\]
A bond on which $\chi$ is not constant has an endpoint in $\{0<\chi<1\}$, since $\chi$ is $\frac12$-Lipschitz and $h<2$; so the bond sum may be restricted to these bonds. Since $\{0<\chi<1\}\subset S$, they have an endpoint in $S$ and are bonds of both domains.

\emph{The bound.} Write $\alpha=\mathrm{e}^{-\delta\rho_y}w$ and $\beta=\mathrm{e}^{-\delta\rho_{y'}}w'$, with $\rho_y=\abs{\cdot-y}$. Since $\delta h\le\frac12$, we have $\abs{\mathrm{e}^s-1}\le2\abs s$ for $\abs s\le\delta h$. Hence
\[
h^{-1}\abs{d_b\alpha}\le \mathrm{e}^{-\delta\rho_y(x)}\big(h^{-1}\abs{d_bw}+2\delta\abs{w(x')}\big).
\]
Lemma~\ref{app:G1:SK-CT}, applied with $F=\mathrm{e}^{\delta\rho_y}f$ and $\norm F\le \mathrm{e}^{\sqrt d\delta}$, gives $q(w)+\norm w^2\le C$; likewise for $w'$. Shifts inside one bond or block change the weights by at most $\mathrm{e}^{2\sqrt d\delta}$. Hence, by the Cauchy--Schwarz inequality, the right-hand side is at most $C\sup_{x\in K}\mathrm{e}^{-\delta(\rho_y(x)+\rho_{y'}(x))}$.

\emph{Geometry.} On $K$ we have $\dist(x,(\Omega_1\cup\Omega_2)\setminus S)\le2+\sqrt d$. So $\rho_y(x)\ge D(y)-2-\sqrt d$, and likewise for $y'$, while $\rho_y(x)+\rho_{y'}(x)\ge\abs{y-y'}$. Therefore
\[
\rho_y(x)+\rho_{y'}(x)\ge\tfrac12\big(\abs{y-y'}+D(y)+D(y')\big)-(2+\sqrt d),
\]
which gives the claim. If $K=\emptyset$, the difference vanishes.
\end{proof}

\begin{proposition}[scalar kernels on the torus]\label{app:K:L13-8}\label{app:G1:SK-kernels}
Let $\Omega\subset T_h$ be a union of big blocks, and let $A$ be a real field on $\Omega$ that is regular in the sense of \cite[(1.7)]{Ba4}, with $e_k$ small as there. Assume (N0) and (NW). Then the following hold.
\begin{itemize}
\item[(i)] \cite[(1.9)--(1.12), (1.15)--(1.20)]{Ba4} hold on $\Omega$ with torus distances. Their constants do not depend on $A$ (in particular not on its constant part), on $\Omega$, $k$, $h$, or the side of the torus.
\item[(ii)] $\abs{\Delta^{(k)}(\Omega,A;y,y')-a_k\delta_{yy'}}\le c_B'\mathrm{e}^{-\delta\abs{y-y'}}$, where $c_B'=2a_k^2\gamma^{-1}\mathrm{e}^{2\sqrt d\delta}$.
\item[(iii)] Let $(\Omega_i,A_i)$, $i=1,2$, be as above and let $S\subset\Omega_1\cap\Omega_2$ be a union of blocks with two properties: every bond of $\Omega_1$ or $\Omega_2$ with an endpoint in $S$ is a bond of both, and $A_1=A_2$ on these bonds. Put $D(y):=\dist(y,(\Omega_1\cup\Omega_2)\setminus S)$. If $D(y),D(y')\ge3+\sqrt d$, then
\[
\abs{\Delta^{(k)}(\Omega_1,A_1;y,y')-\Delta^{(k)}(\Omega_2,A_2;y,y')}\le a_k^2C\mathrm{e}^{-\frac\delta2(\abs{y-y'}+D(y)+D(y'))}.
\]
\item[(iv)] Let $\Omega:=B^k(\Lambda_2^{(k-1)\prime})$, $\Omega_{\rm in}:=B^k(\Lambda_2^{(k)})$ and $\Omega_{\rm out}:=B^k(\Lambda_2^{(k-1)\prime}\cap\Lambda_5^{(k)c})$. Let $\tilde B$ be regular on $\Omega$ with $\tilde B=B^{(k+1),\eta}$ on $\Omega_{\rm in}$ (the latter holds by Proposition~\ref{app:K:L13-11}(ii)). Put
\[
U=\Lambda_6^{(k-1)\prime}\cap\Lambda_3^{(k)c},\quad V=\Lambda_3^{(k)},\quad W=V\setminus\Lambda_4^{(k)},\quad\Delta_{\rm out}=\Delta^{(k)}(\Omega_{\rm out},\tilde B),\quad\Delta_{\rm in}=\Delta^{(k)}(\Omega_{\rm in},B^{(k+1),\eta}).
\]
Let $H_k$ be the self-adjoint operator on $\ell^2(\Lambda_6^{(k-1)\prime})$ defined by
\[
H_k:=\Delta^{(k)}(\Omega,\tilde B)-U\Delta_{\rm out}U-\big(U\Delta_{\rm out}W+W\Delta_{\rm out}U\big)-V\Delta_{\rm in}V,
\]
where a set stands for multiplication by its indicator. Equivalently, $H_k$ is the unique self-adjoint operator for which \cite[(2.108)]{Ba2} holds with the cross term read as its real part and with the factor $\frac12$ in front of $\ip{\Lambda_3^{(k)}\varphi}{\Delta_{\rm in}\Lambda_3^{(k)}\varphi}$. Let $h_k$ be the kernel of $H_k$. If $r(t)\ge4+3\sqrt dM_b$, then for $y,y'\in\Lambda_6^{(k-1)\prime}$
\[
\abs{h_k(y,y')}\le C_H\mathrm{e}^{-\frac\delta2r(t)}\mathrm{e}^{-\frac\delta2\abs{y-y'}},
\]
and $\mathrm{e}^{-\frac\delta2r(t)}\le C_\kappa t^\kappa$ for every $\kappa>0$.
\item[(v)] If $\abs\varphi\le\bar\varphi$ on $\Lambda_6^{(k-1)\prime}$, then $\abs{\ip\varphi{H_k\varphi}}\le C_H\mathrm{e}^{-\frac\delta2r(t)}S_\delta\bar\varphi^2\abs{\Lambda_6^{(k-1)\prime}}$, where $S_\delta=\sum_{z\in\Z^d}\mathrm{e}^{-\delta\abs z/2}$.
\end{itemize}
Here $\gamma,\delta,C$ depend only on $(d,a,\gamma_0)$, and $C_H$ depends only on $(d,a,\gamma_0,M_b)$.
\end{proposition}
\begin{proof}
(i) The setting of \cite{Ba4} includes subsets of a torus $T_h$ with periodic conditions \cite[p.~572]{Ba4}. The proofs there localise in cubes of side $O(M_b)$, and by (NW) these do not wrap. A constant field $c$ has $\partial^h_\mu c=0$, so $c+A''$ satisfies \cite[(1.7)]{Ba4} with the constant of $A''$. Every constant there depends on $A$ only through (1.7). On the torus, $\abs{\cdot-y'}$ is $1$-Lipschitz, so the decay estimates hold with torus distances.

(ii) Put $\rho=\abs{\cdot-y'}$. Then
\[
\Delta^{(k)}(y,y')-a_k\delta_{yy'}=-a_k^2\ip{\mathrm{e}^{-\delta\rho}Q_k^*\delta_y}{w},\qquad w=\mathrm{e}^{\delta\rho}G_k\mathrm{e}^{-\delta\rho}F,\quad F=\mathrm{e}^{\delta\rho}Q_k^*\delta_{y'} .
\]
We have $\norm{Q_k^*\delta_y}=1$, so $\norm F\le \mathrm{e}^{\sqrt d\delta}$ and $\norm{\mathrm{e}^{-\delta\rho}Q_k^*\delta_y}\le \mathrm{e}^{-\delta(\abs{y-y'}-\sqrt d)}$. The operator $X(\Omega,A)$ is of the form of Lemma~\ref{app:G1:SK-CT}, with $v=m^2t^2$, $a'=a_k$ and $P=P_k$, and $X\ge\gamma$. So $\norm w\le2\gamma^{-1}\norm F$, and the bound follows.

(iii) The constants $a_k\delta_{yy'}$ cancel. Since $B(y)\subset S$, and $Q_k$ reads the field only on contours inside a block, the vector $f=Q_k^*\delta_y$ is the same for both fields; likewise $g=Q_k^*\delta_{y'}$. Both have norm $1$. By the hypothesis, the bonds, the transporters, $v=m^2t^2$ and the $P_k$ on the blocks of $S$ agree. Apply Lemma~\ref{app:G1:SK-loc} to $X(\Omega_i,A_i)$.

(iv) Since $M_b\ge1$, the hypothesis gives $r(t)\ge4+\sqrt d+c_M$. By \cite[(2.8)]{Ba2}, $\Lambda_{i+1}^{(k)}$ is obtained from $\Lambda_i^{(k)}$ by removing a layer of big blocks of thickness $r(t)$. Allowing for the rounding to blocks, points of $\Lambda_i^{(k)c}$ and of $\Lambda_{i+1}^{(k)}$ (or of $B^k(\Lambda_{i+1}^{(k)})$) are at distance $\ge r-c_M$.

The kernel $h_k(y,y')$, and its symmetric counterpart, is of one of four kinds.
\begin{itemize}
\item[(a)] $y,y'\in U$: $\Delta^{(k)}(\Omega,\tilde B)-\Delta_{\rm out}$.
\item[(b)] $y\in U$, $y'\in W$: the same difference.
\item[(c)] $y\in U$, $y'\in V\cap\Lambda_4^{(k)}$: $\Delta^{(k)}(\Omega,\tilde B)(y,y')$ itself.
\item[(d)] $y,y'\in V$: $\Delta^{(k)}(\Omega,\tilde B)-\Delta_{\rm in}$.
\end{itemize}

In (a) and (b), apply (iii) with $(\Omega_1,A_1)=(\Omega,\tilde B)$ and $(\Omega_2,A_2)=(\Omega_{\rm out},\tilde B)$. Take for $S$ the set $\Omega_{\rm out}$ minus the blocks within distance $1$ of $B^k(\Lambda_5^{(k)})$. A bond of $\Omega$ with an endpoint in $S$ then lies in $\Omega_{\rm out}$. Since $U\subset\Lambda_3^{(k)c}$ and $W\subset\Lambda_4^{(k)c}$, we get $D(y),D(y')\ge r-c_M-1\ge3+\sqrt d$. Hence
\[
\abs{h_k(y,y')}\le a^2C\mathrm{e}^{\delta(c_M+1)}\mathrm{e}^{-\delta r}\mathrm{e}^{-\frac\delta2\abs{y-y'}}.
\]

In (c) the diagonal term vanishes, and $\abs{y-y'}\ge r-c_M$. By (ii),
\[
\abs{h_k}\le c_B'\mathrm{e}^{-\frac\delta2(r-c_M)}\mathrm{e}^{-\frac\delta2\abs{y-y'}}.
\]

In (d), $\tilde B=B^{(k+1),\eta}$ on $\Omega_{\rm in}$, and a Neumann operator and the averaging on $\Omega_{\rm in}$ read the field only on bonds of $\Omega_{\rm in}$. Hence $\Delta_{\rm in}=\Delta^{(k)}(\Omega_{\rm in},\tilde B)$. Apply (iii) with $S=\Omega_{\rm in}$ minus its boundary layer of blocks. Since $y,y'\in\Lambda_3^{(k)}$, we get $D(y),D(y')\ge r-c_M-1$, and the bound of (a) holds.

Since $c_B'\le2a^2\gamma^{-1}\mathrm{e}^{2\sqrt d\delta}$, (iv) holds with $C_H:=(a^2C+2a^2\gamma^{-1}\mathrm{e}^{2\sqrt d\delta})\mathrm{e}^{\delta(c_M+1)}$. Finally, $\mathrm{e}^{-\delta R\ell(t)^2/2}\le \mathrm{e}^{-\delta R(\log t^{-1})^2/2}\le C_\kappa t^\kappa$ for $0<t\le1$.

(v) We have
\[
\abs{\ip\varphi{H_k\varphi}}\le\bar\varphi^2\sum_{y\in\Lambda_6^{(k-1)\prime}}\sum_{y'}\abs{h_k(y,y')} .
\]
Each point of the torus has a representative $z\in\Z^d$ with $\abs z$ equal to its torus distance from $y$. So $\sum_{y'}\mathrm{e}^{-\delta\abs{y-y'}/2}\le S_\delta$, and (iv) gives the claim.
\end{proof}

\begin{proposition}[{\cite[Prop.~3.1]{Ba2}} at a real background]\label{app:K:L13-9}\label{app:G1:SK-Prop31}
Let $S_k:=\Lambda_5^{(k-1)\prime}\cap\Lambda_5^{(k)c}$ for $k\ge1$, and $S_0:=\Lambda_5^{(0)c}$. Let $\Phi=(\varphi_k|_{S_k})_k$, and define $\ip\Phi{\Delta(\tilde A)\Phi}$ by \cite[(3.23), (3.25)]{Ba2}. Assume:
\begin{itemize}
\item[(F1)] $\abs{\partial\tilde A}\le c_Ap(t_k)$ on each $B^k(S_k)$, in scale-$k$ units;
\item[(F2)] $\abs{\varphi_k(x)}\le c_1\lambda(t_{k-1})^{-1/4}p(t_{k-1})$ for $x\in S_k$, in scale-$(k-1)$ units;
\item[(F3)] $m^2\eps_0^2\le1$, $\eps_0\le\min(1,e^{-4})$ with $e$ the charge, $e\eps_0^{1/2}\le e_*(\frac1{10})$, where $e_*(\alpha)$ is the bound on the coupling under which \cite[(1.22)]{Ba4} holds with the factor $e^{2-\alpha}$, and the conditions (T.nw) and (E.r).
\end{itemize}
Let $\gamma_0>0$ be a constant for which both (N0) and \cite[(1.22)]{Ba4} hold; both bounds remain true when $\gamma_0$ is decreased. Put $\gamma_0'=\min(1,\gamma_0)$ and $\kappa_0=\frac15$. Then
\[
\ip\Phi{\Delta(\tilde A)\Phi}\ge\gamma_0'\sum_{k=0}^K\Big[\sum_{\langle x,x'\rangle\subset S_k}t_k^{d-2}\abs{U(\tilde A(\langle x,x'\rangle))\varphi_k(x')-\varphi_k(x)}^2+\sum_{x\in S_k}t_k^dm^2\abs{\varphi_k(x)}^2\Big]-\sum_{k=1}^KC(\kappa_0)t_k^{\kappa_0}\abs{S_k}.
\]
Here $C(\kappa_0)$ depends on $d,a,\gamma_0,c_A,c_1,b_0,\lambda,e$ and $L=2$. This is \cite[(3.26)]{Ba2}.
\end{proposition}
\begin{proof}
\emph{Step 1 (marginal and minimum).} Let $F(\xi,\zeta)$ be a positive-definite quadratic form. Completing the square gives $\int d\xi\,\mathrm{e}^{-F/2}=c\,\mathrm{e}^{-\frac12\min_\xi F}$, with $c$ independent of $\zeta$. Hence $\ip\Phi{\Delta\Phi}=\min\mathcal F(\varphi_0;\Phi)$, where $\mathcal F$ is the exponent of \cite[(3.23)]{Ba2} and the minimum is over the integrated variable $\varphi_0|_{\Lambda_5^{(0)}}$. Positive definiteness holds because $m^2>0$.

\emph{Step 2 (Neumann decoupling).} Delete from $\mathcal F$ the non-negative bond terms $\eps^d\abs{(D_{\tilde A}\varphi_0)(b)}^2$ of the bonds $b$ that join different regions $B^k(S_k)$; this lowers the minimum. The averaging $Q_k(\tilde A)\varphi_0$ on $S_k$ reads $\varphi_0$ only in $B^k(S_k)$, so the reduced problem splits over $k$. Each piece has the form $\ip\psi{Y\psi}+a\norm{\phi-Q\psi}^2$, where $Y$ collects the bond and mass terms of the piece and $a$ is the coefficient of its averaging term. Minimising over $\psi$ gives $(Y+aQ^*Q)\psi=aQ^*\phi$, and hence
\[
\inf_\psi\big[\ip\psi{Y\psi}+a\norm{\phi-Q\psi}^2\big]=\ip\phi{\big(a-a^2Q(Y+aQ^*Q)^{-1}Q^*\big)\phi}.
\]
This yields \cite[(3.27)--(3.28)]{Ba2}:
\[
\ip\Phi{\Delta\Phi}\ge\sum_k\ip{\varphi_k}{\Delta^{(k)}(B^k(S_k),\tilde A)\varphi_k}.
\]
For $k=0$ the term is $\ip{\varphi_0}{(-\Delta^N_{\tilde A}+m^2)\varphi_0}$ on $S_0$. This is the $k=0$ bracket, and it is at least $\gamma_0'$ times that bracket.

\emph{Step 3 (each shell).} Pass to scale-$k$ units, $\varphi^s=t^{(d-2)/2}\varphi_k$, with $t=t_k\le\eps_0$, rescaled mass $m^2t^2\le1$ and coupling $e_k\le e_*$. We first show $1+\log e_k^{-1}\ge\frac14\ell(t)$.
\begin{itemize}
\item If $e\le1$, then $\log e_k^{-1}\ge\frac12\log t^{-1}$.
\item If $e>1$, then $\log t^{-1}\ge\log\eps_0^{-1}\ge4\log e$ by (F3). Hence $\log e_k^{-1}=\frac12\log t^{-1}-\log e\ge\frac14\log t^{-1}$.
\end{itemize}
So $p(t)\le4^6(b_0/a_0)p_{\rm B}(e_k)$, where $p_{\rm B}(e'):=a_0(1+\log e'^{-1})^6$ is the function of \cite[(1.21)]{Ba4} with exponent $6$. Hence (F1) implies \cite[(1.21)]{Ba4} with a constant independent of $e$. By (T.nw), (E.r) and Lemma~\ref{app:G1:nowrap}, (NW) holds. Then \cite[Prop.~3.1$'$, (1.22)]{Ba4}, applied on the union of unit blocks $B^k(S_k)$ of the torus (subsets of a torus are in the setting of \cite[p.~572]{Ba4}), with the constant $\gamma_0$ of the hypothesis, gives, with $\alpha=\frac1{10}$ and $C_\alpha$ the constant $O(1)$ of \cite[(1.22)]{Ba4},
\[
\ip{\varphi^s}{\Delta^{(k)}\varphi^s}\ge\gamma_0\Big(\sum_{\langle x,x'\rangle\subset S_k}\abs{U\varphi^s(x')-\varphi^s(x)}^2+m^2t^2\sum_{x\in S_k}\abs{\varphi^s}^2\Big)-C_\alpha e_k^{2-\alpha}\sum_{x\in S_k}\abs{\varphi^s(x)}^2.
\]
Undoing the rescaling gives the $k$-th bracket.

\emph{Step 4 (the error).} The scale-$k$ value of the bound (F2) is at most $L^{(d-2)/2}c_1\lambda(t/L)^{-1/4}p(t/L)$. With $d=3$, $\lambda(t)=\lambda t$ and $p(t/L)\le(1+\log L)^6p(t)$, this gives
\[
\abs{\varphi^s}^2\le L^{3/2}(1+\log L)^{12}c_1^2\lambda^{-1/2}t^{-1/2}p(t)^2.
\]
Hence
\[
C_\alpha e_k^{2-\alpha}\abs{\varphi^s}^2\le C\,t^{(1-\alpha)/2}p(t)^2=C\,t^{1/5}\,t^{1/4}b_0^2\ell(t)^{12}\le C(\kappa_0)\,t^{\kappa_0},
\]
since $\sup_{0<t\le1}t^{1/4}\ell(t)^{12}<\infty$. Summing over $S_k$ and over $k$ gives the claim.
\end{proof}

\begin{proposition}[the background as an external field]\label{app:K:L13-11}\label{app:G1:SK-bg}
Put $s_k:=(1-\theta_{k+1})\theta_k$ and $\tilde B:=s_kA^{(k)}+\theta_{k+1}B^{(k+1),\eta}$ \cite[p.~582]{Ba2}.
\begin{itemize}
\item[(i)] Let $A'$ be defined by the translation \cite[(2.80)]{Ba2}, and consider the step-$k$ integral over $A'|_{\Lambda_5^{(k)}}$. The vector variable enters a scalar object only through
\[
\theta_kA^{(k)}=s_kA^{(k)}+\theta_{k+1}A'^{(k)}+\mathcal F(B).
\]
Here $s_kA^{(k)}$ depends on $A'$ only through $A'|_{\Lambda_3^{(k)c}}$, and $\mathcal F(B)=\theta_{k+1}(aL^{-2}C^{(k)}_{\Lambda_0^{(k)}}Q^*B)^{(k)}$. The scalar objects are the normalisations $Z^{(j)}$, the averaging kernel, the basic scalar form, $\mathcal P^{(k)}$ and the scalar characteristic functions.

Moreover, let $\mathcal R$ be a class of real external fields and $\mathcal A$ a set of vector configurations with $\theta_kA^{(k)}\in\mathcal R$ for $A\in\mathcal A$. Suppose $\int d\varphi\,F(A,b,\varphi)\le\mathcal B(A,b)$ for all $A\in\mathcal A$ and $b\in\mathcal R$, with $F\ge0$. Then
\[
\int_{\mathcal A}dA\int d\varphi\,F(A,\theta_kA^{(k)},\varphi)\le\int_{\mathcal A}dA\,\mathcal B(A,\theta_kA^{(k)}).
\]
\item[(ii)] $\tilde B=B^{(k+1),\eta}$ on $B^k(\Lambda_2^{(k)})$. Hence $\Delta^{(k)}$, $G_k$ and $C^{(k)}_{\Lambda_5}$ on that domain, and $Q(\tilde B)$ on blocks inside it, agree at $\tilde B$ and at $B^{(k+1),\eta}$.
\item[(iii)] Let $L=2$, $r(\eps_0)\ge4\sqrt dM_b$ and $2R\log2\,\ell(\eps_0)\ge2\sqrt dM_b$. Then:
 \begin{itemize}
 \item the $\zeta^{(l)}$-balls around $\supp\theta_l$ and around the $\theta_{l+1}$-slices lie in $\Lambda_{-1}^{(l-1)\prime}$, the domain of \cite[(2.55)]{Ba2};
 \item the split \cite[(2.89)]{Ba2} holds as printed;
 \item \cite[Lemma~2.3, (2.93)--(2.98)]{Ba2} hold from the restrictions in force; in particular $\abs{\partial\tilde B}\le O(1)p(t)$ on $B^k(\Lambda_2^{(k-1)\prime})$.
 \end{itemize}
\item[(iv)] Let $w_j\in[w_-,w_+]$ and $\sigma:=\max(w_+,w_-^{-1})^{1/2}$. Assume that on the slice line the vector characteristic functions admit exactly the configurations obeying the vector restrictions of \cite[(2.55)]{Ba2} with $c_1$ replaced by $\sigma c_1$. (By Lemma~\ref{app:K:L13-10}(iii) this holds for $s_k=\sigma^{-1}$ when $\chi_k^{\rm v}$ is the indicator of $\mathcal R_k(c_1)$. Corollary~\ref{app:G1:restr-slice-use} gives the conclusion under the weaker property of Lemma~\ref{app:K:L13-10}(ii).) Then (i)--(iii) and Propositions~\ref{app:K:L13-8} and~\ref{app:K:L13-9} hold on the slice line, with $c_1$ replaced by $\sigma c_1$ in the vector restrictions. For $w_j=1+u\delta_j$ with $\abs u\le1$ and $0\le\delta_j\le\frac1{12}$, we have $\sigma\le(12/11)^{1/2}<\sqrt2$.
\end{itemize}
\end{proposition}
\begin{proof}
(i) \emph{Locality of $s_kA^{(k)}$.} By \cite[(2.44)]{Ba2}, the value of $A^{(k)}$ at $x$ reads $A_k$ only at points $y$ with $\zeta^{(k)}(x,y)\ne0$, and these satisfy $\abs{x-y}<r(t)-2M_b$. Also $\theta_{k+1}=1$ on $B^k(\Lambda_2^{(k)})$ \cite[p.~567]{Ba2}. So if $x\in\supp s_k$, then $x\notin B^k(\Lambda_2^{(k)})$. For such $x$, the big block of $y$ is within distance $<r$ of $\Lambda_2^{(k)c}$; the margin $2M_b$ covers the rounding to big blocks. By \cite[(2.8)]{Ba2} that block is not in $\Lambda_3^{(k)}$. For $j\ne k$, $A^{(j)}$ is linear in $A_j$ alone, since $G_j$ is the vector propagator at zero background.

\emph{The decomposition.} We have $\theta_k\theta_{k+1}=\theta_{k+1}$, for three reasons: $\theta_{k+1}$ is supported within distance $M_b$ of $B^k(\Lambda_2^{(k)})$; $\Lambda_0^{(k)}\subset\Lambda_7^{(k-1)\prime}$ \cite[p.~566]{Ba2}; and $\theta_k=1$ on $B^{k-1}(\Lambda_2^{(k-1)})$ \cite[p.~567]{Ba2}. Hence $\theta_kA^{(k)}=s_kA^{(k)}+\theta_{k+1}A^{(k)}$. Since $A^{(k)}$ is linear in $A=A'+aL^{-2}C^{(k)}_{\Lambda_0^{(k)}}Q^*B$, we get $\theta_{k+1}A^{(k)}=\theta_{k+1}A'^{(k)}+\mathcal F(B)$. By the locality above, $s_kA^{(k)}$ reads $A'$ only on $\Lambda_3^{(k)c}\subset\Lambda_5^{(k)c}$.

\emph{The list of places.} The places where the vector variable enters a scalar object are those of \cite[(2.53), pp.~576, 579]{Ba2}, together with the scalar characteristic functions. The latter are bounded by $1$ for upper bounds, or are retained where they read no integrated variable.

\emph{The transfer.} The integrand is non-negative, so the scalar variables may be integrated first at each fixed vector configuration $A\in\mathcal A$. At that configuration the external field $b=\theta_kA^{(k)}$ is fixed and lies in $\mathcal R$. Hence the inner integral is at most $\mathcal B(A,\theta_kA^{(k)})$, and integrating over $\mathcal A$ gives the claim.

(ii) On $B^k(\Lambda_2^{(k)})$ we have $\theta_{k+1}=1$ and $s_k=0$, so $\tilde B=B^{(k+1),\eta}$ there. A Neumann operator on a domain reads the field only on the bonds of that domain, and $Q$ reads it only on contours inside a block.

(iii) \emph{The $\zeta$-balls.} In scale-$l$ units, a $\zeta^{(l)}$-ball around $\supp\theta_l$ reaches at most $r(t_l)-M_b+c_M$ beyond $B^{l-1}(\Lambda_2^{(l-1)})$. By three applications of \cite[(2.8)]{Ba2}, $\Lambda_{-1}^{(l-1)\prime}$ contains the neighbourhood of $\Lambda_2^{(l-1)\prime}$ of radius $(3\lfloor r(t_{l-1})\rfloor+1)/L$ in scale-$l$ units. At $L=2$ the inclusion follows from
\[
2r(t_l)-2M_b+2c_M\le3r(t_{l-1})-2M_b\le3\lfloor r(t_{l-1})\rfloor+1 .
\]
This uses three facts: $r(t_l)\le r(t_{l-1})$, because $t_{l-1}=t_l/L$ and $r$ decreases; $r(t_{l-1})\ge r(\eps_0)\ge2c_M$; and $M_b\ge1$. Call this inclusion (G1). The same count one scale up handles the $\theta_{l+1}$-slices.

\emph{The split.} The split \cite[(2.89)]{Ba2} as printed differs from the exact one by
\[
a_{k+1}L^{-2}\theta_{k+1}\zeta^{(k+1)}G_{k+1}Q_{k+1}^*\Lambda_0^{(k)\prime c}B .
\]
This term vanishes when $Lr(t_{k+1})+Lc_M\le2r(t_k)$. Indeed, the $\zeta^{(k+1)}$-balls around $\supp\theta_{k+1}$ then stay inside $\Lambda_0^{(k)\prime}$, by two applications of \cite[(2.8)]{Ba2}. At $L=2$ the condition reads $r(t_k)-r(t_{k+1})\ge c_M$. Since $\ell(t_k)=\ell(t_{k+1})+\log2$ and $t_{k+1}\le\eps_0$,
\[
r(t_k)-r(t_{k+1})=R\log2\,\big(2\ell(t_{k+1})+\log2\big)\ge2R\log2\,\ell(\eps_0)\ge c_M .
\]

\emph{The regularity bounds.} The proofs of \cite[Lemma~2.3, (2.93)--(2.98)]{Ba2} use the restrictions \cite[(2.55)]{Ba2} only at the block of $x$ and at the points $y'$ with $\zeta(x,y')\ne0$ \cite[(2.61)--(2.64)]{Ba2}. By (G1) these points lie in $\Lambda_{-1}^{(l-1)\prime}$, where the restrictions are in force.

(iv) Every ingredient of (i)--(iii) and of Propositions~\ref{app:K:L13-8} and~\ref{app:K:L13-9} is of one of three kinds: geometric; about scalar operators at a fixed real vector configuration; or an identity for integrals of non-negative functions or for normalisations. None of them involves the photon covariance.

In particular, the averaging background may be chosen freely. Indeed, for fixed $(A,\varphi)$ and any measurable $b(A)$, the integral $\int d\psi\,\mathcal N_a^{-1}\mathrm{e}^{-\frac a2\abs{\psi-Q(b(A))\varphi}^2}$ equals $1$, where $\mathcal N_a=\int d\psi\,\mathrm{e}^{-\frac a2\abs\psi^2}$, because the integrand is a translate of a normalised Gaussian. So one may take the background of \cite[(2.44)]{Ba2}. The vector configuration therefore enters only through the class of admissible backgrounds, and every constant depends on that class only through the restriction constants. Under the assumption these are those of \cite[(2.55)]{Ba2} with $\sigma c_1$ in place of $c_1$.

Finally, if $w_j\in[1-\frac1{12},1+\frac1{12}]$, then $w_+\le\frac{13}{12}$ and $w_-^{-1}\le\frac{12}{11}$. So $\sigma\le(12/11)^{1/2}<\sqrt2$.
\end{proof}

\subsection{Bounds on the background at step $k$}\label{app:G1:bg-bounds}

Throughout this subsection $1\le k\le K-1$, and we work in scale-$k$ units: $\eta=L^{-k}$, $t=t_k$, and $\partial^\eta_\mu$ is the forward lattice derivative on the $\eta$-lattice. Put $\Omega=B^k(\Lambda_0^{(k)})$. We use the fields of \cite[pp.~570--576]{Ba2}. The photon mass $\mu_0$ of \cite{Ba2} is $\mu_*$. $A^{(k)}$ is the field of \cite[(2.56), (2.61)]{Ba2}, and $A'$ is defined by the translation \cite[(2.80)]{Ba2}. The fields $A'^{(k)}=a_k\zeta^{(k)}G_kQ_k^*A'$, $\tilde B^{(k)}$ and $B^{(k+1),\eta}$ are defined by the decomposition \cite[(2.89), p.~575]{Ba2}:
\begin{equation}\label{app:G1:bg-289}
\theta_{k+1}A^{(k)}=\theta_{k+1}A'^{(k)}+\tilde B^{(k)}+\theta_{k+1}B^{(k+1),\eta}.
\end{equation}
Put $\tilde A:=\theta_{k+1}A'^{(k)}+\tilde B^{(k)}$ and $\tilde B:=(1-\theta_{k+1})\theta_kA^{(k)}+\theta_{k+1}B^{(k+1),\eta}$. The cut-off functions take values in $[0,1]$. The product of a point function $\theta$ with a vector field $f$ is taken in a fixed convention $(\theta f)_b:=(\lambda\theta(b_-)+(1-\lambda)\theta(b_+))f_b$ with a fixed $\lambda\in[0,1]$; $\lambda=1$ is the printed convention (proof of Lemma~\ref{app:G1:Rthetab}(ii)). This convention is linear in $\theta$, sends $1$ to the identity, has factor in $[0,1]$, and has factor $1$ when $\theta=1$ at both end points of $b$. $g_k$ is the localisation function of \cite[p.~579]{Ba2}: a smooth function with $\supp g_k\subset B^k(\Lambda_7^{(k)})$ and $g_k=1$ on $B^k(\Lambda_8^{(k)})$. As for the function $g_k$ of \cite[p.~412]{Ba3} (``bounded by $1$ together with derivatives''), we fix it with $0\le g_k\le1$ and $\abs{\partial^\eta g_k}\le1$. Such a choice exists: by Lemma~\ref{app:G1:Rnest}(v), $D:=\dist(B^k(\Lambda_8^{(k)}),B^k(\Lambda_7^{(k)})^c)\ge\lfloor r_k\rfloor+1\ge2$ in scale-$k$ units, and $\min\{1,\dist(x,B^k(\Lambda_7^{(k)})^c)/(D-1)\}$ has these properties.

\begin{lemma}[bounds on the background]\label{app:K:L13-3}
Let $L=2$, $1\le k\le K-1$, and assume (IH$_k$) of \eqref{app:G1:eq-IH}. In particular the restrictions \cite[(2.55), p.~570]{Ba2} hold on $\Lambda_{-1}^{(k-1)\prime}$, and the restrictions on $(B,\psi,A,\varphi)$ on $\Lambda_7^{(k-1)\prime}$ under which \cite[(2.86)--(2.91), p.~575]{Ba2} are stated also hold. Let $C_{59},C_{60},C_{90},C_{91}$ be the constants in \cite[(2.59), (2.60), (2.90) with $\kappa=1$, (2.91)]{Ba2}. Let $C_{90}',C_{91}'$ be the constants of the derivative clauses printed with them \cite[p.~575]{Ba2}: after (2.90), ``and similar estimates for the derivative'' (taken with $\kappa=1$), and after (2.91), ``similarly for the derivative''. Put
\[
C:=L(1+\log L)^p+\frac{C_{59}+C_{90}+C_{91}}{c_1},\qquad C_B:=C_{60}+C_{90}'+C_{91}'+C_{91}.
\]
Then the following hold.
\begin{itemize}
\item[(a)] On $\Omega$ we have $\abs{A^{(k)}}\le C c_1p(t)/(\mu_*t)$, $\abs{\tilde A}\le(C_{90}+C_{91})p(t)$, and $\abs{\tilde B},\abs{\theta_{k+1}B^{(k+1),\eta}}\le Cc_1p(t)/(\mu_*t)$. On $\Omega\cap\{\theta_{k+1}=1\}$, which contains every bond of $\Omega\cap B^k(\Lambda_2^{(k)})$, we have $\abs{B^{(k+1),\eta}}\le Cc_1p(t)/(\mu_*t)$.
\item[(b)] On every bond of $\Omega$ we have $\abs{\partial^\eta\tilde B}\le C_Bp(t)$, $\abs{\partial^\eta(\theta_{k+1}A'^{(k)})}\le C_Bp(t)$ and $\abs{\partial^\eta(g_kA'^{(k)})}\le C_Bp(t)$. Consequently the field strengths of $\tilde B$, $\theta_{k+1}A'^{(k)}$ and $g_kA'^{(k)}$ are bounded by $2C_Bp(t)$ on the plaquettes of $\Omega$.
\end{itemize}
$C$ and $C_B$ depend only on the constants of \cite{Ba2} named above, on $c_1$, $L$ and $p$. In particular they do not depend on $\eps$, $k$, the torus or the configuration. For $k=0$ every use of these bounds involves the empty product $\prod_{j<0}$ (or the empty sum $\sum_{j<0}$), so no bound is needed there.
\end{lemma}

\begin{lemma}[elementary comparisons]\label{app:G1:bg-elem}
Let $L=2$, $\eps=2^{-K}\eps_0$ and $t_k=L^k\eps$. For $1\le k\le K-1$ we have $t_k\le1$, $\mu_*t_k\le1$, $p(t_k)\ge1$ and $p(t_{k-1})\le(1+\log L)^pp(t_k)$.
\end{lemma}
\begin{proof}
First, $t_k=2^{k-K}\eps_0<t_K=\eps_0$. We have $\eps_0=2^{-n_0}\le1$ by (T.div), and $\eps_0\le1/\mu_*$ by (E.fin). This gives the first two bounds. Hence $\ell(t_k)=1+\log t_k^{-1}\ge1$, and $p(t_k)=b_0\ell(t_k)^p\ge b_0\ge1$ by the lower bound on $b_0$ in Table~\ref{p:tab:A3}. Finally, since $t_{k-1}=t_k/L$,
\[
\ell(t_{k-1})=\ell(t_k)+\log L\le(1+\log L)\,\ell(t_k)
\]
because $\ell(t_k)\ge1$. Raise this to the power $p$ and multiply by $b_0$.
\end{proof}

\begin{lemma}[the background is the field minus the small field]\label{app:G1:bg-id}
Let $\Omega_\eta$ be the set of bonds whose end points lie within distance $2\eta$ of $\Omega$. On every bond of $\Omega_\eta$ we have $\tilde B=A^{(k)}-\tilde A$, and $\theta_{k+1}B^{(k+1),\eta}=\theta_{k+1}A^{(k)}-\tilde A$.
\end{lemma}
\begin{proof}
By the definition of $\Lambda_0^{(k)}$ \cite[p.~570]{Ba2}, $\Lambda_0^{(k)}$ is contained in $\Lambda_7^{(k-1)\prime}$. We have $\Lambda_7^{(k-1)\prime}\subset\Lambda_2^{(k-1)\prime}\subset\Lambda_0^{(k-1)\prime}$ by Lemma~\ref{app:G1:Rnest}(i), and $\Lambda_0^{(k-1)\prime}\subset\Lambda_{-1}^{(k-1)\prime}$ because $\Lambda_{-1}$ is the set of points at distance $<r$ from $\Lambda_0$ \cite[p.~558]{Ba2}. The function $\theta_k$ equals $1$ on $B^{k-1}(\Lambda_2^{(k-1)})$ \cite[p.~567]{Ba2}, and $B^{k-1}(\Lambda_2^{(k-1)})=B^k(\Lambda_2^{(k-1)\prime})\supset B^k(\Lambda_7^{(k-1)\prime})\supset\Omega$ by Lemma~\ref{app:G1:Rnest}(ii),(iii). By Lemma~\ref{app:G1:Rnest}(v), $\dist(B^k(\Lambda_7^{(k-1)\prime}),B^k(\Lambda_2^{(k-1)\prime})^c)>(5r_{k-1}-4)/L-1>2\eta$. So $\theta_k=1$ at both end points of every bond of $\Omega_\eta$, and these bonds lie in $B^k(\Lambda_2^{(k-1)\prime})$. By the convention for products, $(1-\theta_{k+1})\theta_kA^{(k)}=(1-\theta_{k+1})A^{(k)}$ on these bonds.

The convention is linear in the point function and sends $1$ to the identity, so $A^{(k)}=(1-\theta_{k+1})A^{(k)}+\theta_{k+1}A^{(k)}$; this is the splitting used in \cite[p.~575]{Ba2}. Insert \eqref{app:G1:bg-289}. This gives
\[
A^{(k)}=(1-\theta_{k+1})A^{(k)}+\theta_{k+1}B^{(k+1),\eta}+\theta_{k+1}A'^{(k)}+\tilde B^{(k)}=\tilde B+\tilde A
\]
on the bonds of $\Omega_\eta$. The second identity is \eqref{app:G1:bg-289} rearranged.
\end{proof}

\begin{proof}[Proof of Lemma~\ref{app:K:L13-3}]
\emph{(a).} Let $x\in\Omega$ and let $y$ be the unit-lattice point with $x\in B^k(y)$. Then $y\in\Lambda_2^{(k-1)\prime}\subset\Lambda_{-1}^{(k-1)\prime}$, as in the proof of Lemma~\ref{app:G1:bg-id}. By \cite[(2.59), p.~571]{Ba2}, by the restriction $\abs{A(y)}\le c_1p(t_{k-1})/(\mu_*t_{k-1})$ of \cite[(2.55)]{Ba2}, and by Lemma~\ref{app:G1:bg-elem} together with $t_{k-1}=t/L$,
\[
\abs{A^{(k)}(x)}\le\frac{c_1p(t_{k-1})}{\mu_*t_{k-1}}+C_{59}p(t)\le\Big[L(1+\log L)^p+\frac{C_{59}}{c_1}\Big]\frac{c_1p(t)}{\mu_*t}.
\]
In the last step we used $\mu_*t\le1$. Next, by \cite[(2.91)]{Ba2} and by \cite[(2.90)]{Ba2} with $\kappa=1$,
\[
\abs{\tilde A}\le\abs{\theta_{k+1}A'^{(k)}}+\abs{\tilde B^{(k)}}\le C_{91}p(t)+C_{90}t\le(C_{90}+C_{91})p(t),
\]
since $t\le1\le p(t)$. Lemma~\ref{app:G1:bg-id} and $\abs{\theta_{k+1}A^{(k)}}\le\abs{A^{(k)}}$ then give
\[
\abs{\tilde B},\ \abs{\theta_{k+1}B^{(k+1),\eta}}\le\abs{A^{(k)}}+\abs{\tilde A}\le Cc_1p(t)/(\mu_*t),
\]
again using $\mu_*t\le1$. On bonds where $\theta_{k+1}=1$ at both end points, $\theta_{k+1}B^{(k+1),\eta}=B^{(k+1),\eta}$ by the convention. Every bond of $B^k(\Lambda_2^{(k)})$ is such a bond, by Lemma~\ref{app:G1:Rthetab}(i).

\emph{(b).} By the derivative clauses of \cite[(2.90) with $\kappa=1$, (2.91), p.~575]{Ba2}, $\abs{\partial^\eta\tilde B^{(k)}}\le C_{90}'t\le C_{90}'p(t)$ and $\abs{\partial^\eta(\theta_{k+1}A'^{(k)})}\le C_{91}'p(t)\le C_Bp(t)$. Hence $\abs{\partial^\eta\tilde A}\le(C_{90}'+C_{91}')p(t)$. By Lemma~\ref{app:G1:bg-id} on $\Omega_\eta$, and by \cite[(2.60), p.~571]{Ba2}, which holds on $B^k(\Lambda_2^{(k-1)\prime})$ and so on the bonds of $\Omega_\eta$ (proof of Lemma~\ref{app:G1:bg-id}),
\[
\abs{\partial^\eta\tilde B}\le\abs{\partial^\eta A^{(k)}}+\abs{\partial^\eta\tilde A}\le(C_{60}+C_{90}'+C_{91}')p(t)\le C_Bp(t).
\]

Now consider $g_kA'^{(k)}$, and let $(\partial^\eta_\nu(g_kA'^{(k)}))(x)$ be a value that is not zero. Then $g_k$ does not vanish at some point within distance $2\eta$ of $x$. So every point and bond entering this difference lies within distance $3\eta$ of $\supp g_k\subset B^k(\Lambda_7^{(k)})$. Hence it lies in $B^k(\Lambda_2^{(k)})$, because $\dist(B^k(\Lambda_7^{(k)}),B^k(\Lambda_2^{(k)})^c)\ge5\lfloor r_k\rfloor>3\eta$ (Lemma~\ref{app:G1:Rnest}(v)); there $\theta_{k+1}=1$ (Lemma~\ref{app:G1:Rthetab}(i)). Therefore $g_kA'^{(k)}=g_k\cdot(\theta_{k+1}A'^{(k)})$ at all these bonds. The discrete Leibniz rule
\[
\partial^\eta_\nu(gf)(x)=g(x+\eta e_\nu)\,(\partial^\eta_\nu f)(x)+(\partial^\eta_\nu g)(x)\,f(x)
\]
holds for the start-point convention ($\lambda=1$). For general $\lambda$, write $(g f)_\mu(x)=c_\mu(x)f_\mu(x)$ with $c_\mu(x)=\lambda g(x)+(1-\lambda)g(x+\eta e_\mu)$. Then $\partial^\eta_\nu(gf)_\mu(x)=c_\mu(x+\eta e_\nu)(\partial^\eta_\nu f_\mu)(x)+(\partial^\eta_\nu c_\mu)(x)f_\mu(x)$, with $0\le c_\mu\le1$ and $\abs{\partial^\eta_\nu c_\mu}\le\lambda\abs{\partial^\eta_\nu g(x)}+(1-\lambda)\abs{\partial^\eta_\nu g(x+\eta e_\mu)}\le\sup\abs{\partial^\eta g}$. With $0\le g_k\le1$ and $\abs{\partial^\eta g_k}\le1$ (the choice of $g_k$ above), and with \cite[(2.91)]{Ba2} and its derivative clause,
\[
\abs{\partial^\eta(g_kA'^{(k)})}\le C_{91}'p(t)+C_{91}p(t)\le C_Bp(t).
\]
Finally, $F_{\mu\nu}(f)=\partial^\eta_\mu f_\nu-\partial^\eta_\nu f_\mu$, so $\abs{F(f)}\le2\max_{\mu,\nu}\sup\abs{\partial^\eta_\mu f_\nu}$.
\end{proof}

\begin{remark}\label{app:G1:bg-rem}
(i) Lemma~\ref{app:G1:bg-id} makes the slice analysis \cite[(2.95)--(2.98), pp.~576--577]{Ba2} unnecessary for the bound on $\partial^\eta\tilde B$ at scale $k$. The bound is consistent with \cite[(2.94), (2.98)]{Ba2}. Indeed, rescaled to scale-$k$ units, the factor $(L^j\eta)^{d/2}$ of \cite[(2.94)]{Ba2} is exactly the scaling of a derivative of a vector field.

(ii) Rescaled to scale $j<k$, a vector field acquires the factor $L^{-(k-j)(d-2)/2}$ and its derivative acquires $L^{-(k-j)d/2}$. Hence Lemma~\ref{app:K:L13-3} gives the bounds on the fields at scale $j$ that are used in Lemma~\ref{app:K:L13-4}.
\end{remark}

\subsection{Collars of the polynomial $\mathcal Q^{(k)}$}\label{app:G1:Qcoll}

\emph{Conventions.} We work in scale-$k$ units, so $t=t_k$ and $r=r(t)=R\,\ell(t)^2$ with $\ell(t)=1+\log t^{-1}$. Large blocks \cite[p.~607]{Ba1} are unions of scale-$k$ unit blocks and have side at most $s_b:=LM_b$; only this upper bound is used. Each set $\Lambda_n=\Lambda_n^{(k)}$ is a union of closed large blocks, and $\Lambda_n^c$ denotes the union of the remaining large blocks. Both are closed, and $\Lambda_n\cup\Lambda_n^c$ is the whole torus. By \cite[(2.8), p.~558]{Ba2}, applied at scale $k$ with $r(\eps)$ replaced by $r(t)$ \cite[p.~566]{Ba2}, $\Lambda_{n+1}^c$ is the union of the large blocks whose distance from $\Lambda_n^c$ is at most $r$. Distances are Euclidean torus distances.

We have $\Omega=B^k(\Lambda_0)$, as in \cite[(2.106), p.~579]{Ba2}. By \cite[p.~567]{Ba2}, $\theta_{k+1}=1$ on $B^k(\Lambda_2)$, and $\theta_{k+1}$ vanishes at distance $\ge M_\theta$ from $\Lambda_2$, where $M_\theta\le s_b$. By \cite[p.~579]{Ba2}, $g_k=1$ on $\Lambda_8$ and $\supp g_k\subset\Lambda_7$. Put $a:=\theta_{k+1}A'^{(k)}$ and $a':=g_kA'^{(k)}$, and write $a_j,a'_j$ for these fields in scale-$j$ units. Put $\mathcal C:=\Lambda_2\cap\Lambda_7^c$. Here $\abs{\cdot}$ counts scale-$k$ unit blocks.

The sites $y$ at scale $j<k$ lie in the scale-$k$ unit blocks; each unit block contains $L^{3(k-j)}$ of them. By the setting of Subsection~\ref{app:G1:cumcoll},
\[
q_{j,y}(a)=-\tfrac12\sum_{n=1}^{\nbar}\tfrac1{n!}D^nf_y(\tilde B_j)[a_j,\dots,a_j].
\]

\begin{lemma}[collar geometry]\label{app:G1:Qcoll-geo}
Assume $r\ge 4M_b+4$. Put $S:=\{x\notin\Lambda_2:\dist(x,\Lambda_2)<M_\theta\}$, let $\mathcal D$ be the closed $1$-neighbourhood of $(\Lambda_2\setminus\Lambda_8)\cup S$, and put $\rho_1:=r+\sqrt3\,s_b+1$.
\begin{itemize}
\item[(a)] $\Lambda_{n+1}\subset\Lambda_n$, and $\dist(\Lambda_{n+i},\Lambda_n^c)>ir$ for all $n\ge0$ and $i\ge1$.
\item[(b)] Every point of $\mathcal D$ lies within distance $\rho_1$ of $\mathcal C$. Every point of $\mathcal D$ lies within distance $1+M_\theta$ of $\Lambda_2$. In particular $\mathcal D=\emptyset$ if $\mathcal C=\emptyset$.
\item[(c)] For $u\ge0$, the number of scale-$k$ unit blocks that meet $\{x:\dist(x,\mathcal D)\le u\}$ is at most $(2u+2\rho_1+2\sqrt3+1)^3\abs{\mathcal C}$.
\end{itemize}
\end{lemma}

\begin{proof}
(a) A block of $\Lambda_n^c$ has distance $0$ from $\Lambda_n^c$, so it lies in $\Lambda_{n+1}^c$; hence $\Lambda_{n+1}\subset\Lambda_n$. The case $i=1$ is the definition. For the induction step, let $x\in\Lambda_{n+i+1}$ and $z\in\Lambda_n^c$, and suppose $\abs{x-z}\le(i+1)r$. If $\abs{x-z}\le r$, this contradicts the case $i=1$, because $x\in\Lambda_{n+1}$. Otherwise, choose $w$ on a shortest path from $x$ to $z$ with $\abs{w-z}=r$, so that $\abs{x-w}\le ir$. Every block containing $w$ lies within $r$ of $\Lambda_n^c$, hence in $\Lambda_{n+1}^c$. Therefore $\dist(x,\Lambda_{n+1}^c)\le ir$, which contradicts the induction hypothesis at $(n+1,i)$.

(b) Let $x\in\mathcal D$, and choose $x'\in(\Lambda_2\setminus\Lambda_8)\cup S$ with $\abs{x-x'}\le1$.

First case: $x'\in\Lambda_2\setminus\Lambda_8$. Then $x'$ lies in a block $\square\subset\Lambda_2\cap\Lambda_8^c$. If $\square\subset\Lambda_7^c$, then $x'\in\mathcal C$. Otherwise $\square\subset\Lambda_7$, and $\dist(\square,\Lambda_7^c)\le r$ by the definition of $\Lambda_8$. Let $z\in\Lambda_7^c$ be a nearest point to $x'$. Then $\abs{x'-z}\le r+\sqrt3s_b$, and $z\in\Lambda_7$, since points of a shortest path from $x'$ to $z$ that precede $z$ lie in $\Lambda_7$. Let $\square'\subset\Lambda_7^c$ be a block containing $z$. By (a), $\dist(\square',\Lambda_2^c)\ge\dist(\Lambda_7,\Lambda_2^c)-\sqrt3s_b>5r-\sqrt3s_b>0$, so $\square'\subset\Lambda_2$ and $z\in\mathcal C$.

Second case: $x'\in S$. Let $z\in\Lambda_2$ be a nearest point to $x'$, so $\abs{x'-z}<M_\theta$. Then $z$ also lies in $\Lambda_2^c$. A block $\square'\subset\Lambda_2$ containing $z$ therefore has distance $0$ from $\Lambda_2^c$, so $\square'\subset\Lambda_3^c\subset\Lambda_7^c$ by (a), and $z\in\mathcal C$.

In both cases $\dist(x,\mathcal C)\le1+\max(r+\sqrt3s_b,M_\theta)=\rho_1$, since $M_\theta\le s_b$. The second claim is clear.

(c) Let $b$ be a unit block containing a point $p$ with $\dist(p,\mathcal D)\le u$. In the proof of (b), the point of $\mathcal C$ that was found ($x'\in\square$ or $z\in\square'$) lies in a large block contained in $\mathcal C$, and large blocks are unions of unit blocks. Hence there is a point $q$ in a unit block $b_0\subset\mathcal C$ with $\abs{p-q}\le u+\rho_1$. The centres of $b$ and $b_0$ are then within $s:=u+\rho_1+\sqrt3$ of each other. At most $(2s+1)^3$ unit blocks have their centre within $s$ of a given centre.
\end{proof}

\begin{lemma}[collars of $\mathcal Q^{(k)}$]\label{app:K:L13-4}
Let $L=2$, $0\le k\le K-1$ and $t=t_k\le\eps_0$, with $\eps_0$ obeying (E.ksp1) and (E.ksp2) and with the regulator data of Theorem~\ref{thm:T6} (in particular $R\ge R^{**}$). Assume (IH$_k$). Use the setting, the hypotheses (Q0)--(Q2) and (Fd1)--(Fd3), and the site decomposition $\mathcal Q^{(k)}=\sum_{j<k}\sum_yq_{j,y}$ of Subsection~\ref{app:G1:cumcoll}, with $\Omega=B^k(\Lambda_0^{(k)})$. Suppose that $\Omega^c\ne\emptyset$ (or that $R_y$ is capped) and that $R_y\le r(t)+M_b<\ell_T/t$. Then, uniformly in $k$,
\[
\abs{\mathcal Q^{(k)}(\Omega;\tilde B,\theta_{k+1}A')-\mathcal Q^{(k)}(\Omega;\tilde B,g_kA')}\le\Big[\frac{C_{\rm coll}e^2b_0^2\,t\,\ell(t)^{2p}(r(t)+M_b)^3}{L-1}+C_Nt^N\Big]\abs{\Lambda_2^{(k)}\cap\Lambda_7^{(k)c}}
\]
for every $N$. Here $C_{\rm coll}$ depends only on $(d,L,a,M_b,c_1,\gamma_0,\gamma_0',c_M,\delta,\delta_0,\nbar)$, and $C_N$ depends in addition on $N$ and the parameters.
\end{lemma}

\begin{proof}
Since $\ell(t)\ge1$, we have $r=R\ell(t)^2\ge R\ge R^{**}\ge32M_b+25$, so Lemma~\ref{app:G1:Qcoll-geo} applies.

\emph{Step 1 (support of the difference).} The values of $a$ and $a'$ on a scale-$j$ bond $b$ are determined by $A'^{(k)}(b)$ and by the values of $\theta_{k+1}$ and $g_k$ at the endpoints of $b$, which lie within $\frac12$ of $b_-$. Both functions equal $1$ on $B^k(\Lambda_8)$, since $\Lambda_8\subset\Lambda_2$. Both vanish off $\Lambda_2\cup S$, since $\Lambda_7\subset\Lambda_2$. Hence $a_j(b)\ne a'_j(b)$ implies $b_-\in\mathcal D$, so
\[
\norm{a_j-a'_j}_y\le2P_1(t)\,\mathrm{e}^{-\delta_1L^{k-j}\dist(y,\mathcal D)}
\]
in the norm of Lemma~\ref{app:G1:Q-S}. Here we used $\sup\abs{a_j},\sup\abs{a'_j}\le P_1(t)$ (paragraph~\ref{app:G1:Q-fields}). If $\mathcal C=\emptyset$, then $\mathcal D=\emptyset$ and the left side vanishes, so assume $\mathcal C\ne\emptyset$.

\emph{Step 2 (near sites).} Let $Y$ be the set of scale-$k$ unit blocks that meet $\{x:\dist(x,\mathcal D)\le r/4\}$. By Lemma~\ref{app:G1:Qcoll-geo}(b) and (a), every point of $Y$ lies within $r/4+\sqrt3+1+M_\theta$ of $\Lambda_2$, and $\dist(\Lambda_2,\Lambda_0^c)>2r$. Hence
\[
\dist(Y,\Omega^c)>\tfrac74r-\sqrt3-1-s_b\ge\tfrac r2 .
\]
So Corollary~\ref{app:G1:Q-sum} applies to $a$ and to $a'$, both of which are fields of paragraph~\ref{app:G1:Q-fields}. With $\vartheta^2=e^2t\,p(t)^2=e^2b_0^2t\ell(t)^{2p}$, it gives
\[
\sum_{j<k}\sum_{y\in B^{k-j}(Y)}\abs{q_{j,y}(a)-q_{j,y}(a')}\le2\Big[\frac{8C_Qc_{AX}^2C_B^2\vartheta^2}{L-1}+C_Nt^N\Big]\abs Y .
\]
Lemma~\ref{app:G1:Qcoll-geo}(c) with $u=r/4$ gives $2u+2\rho_1+2\sqrt3+1\le\frac52r+8M_b+9\le15(r+M_b)$, so $\abs Y\le15^3(r+M_b)^3\abs{\mathcal C}$.

\emph{Step 3 (far sites).} Let $y$ be a scale-$j$ site, $m:=k-j\ge1$, with $y\notin B^{m}(Y)$; then $\dist(y,\mathcal D)>r/4$. The forms $D^nf_y(\tilde B_j)$ are symmetric and multilinear, so
\[
D^nf_y[a_j^{\otimes n}]-D^nf_y[a_j'^{\otimes n}]=\sum_{l=1}^nD^nf_y[a'_j,\dots,a'_j,a_j-a'_j,a_j,\dots,a_j],
\]
with $a'_j$ in the first $l-1$ arguments. Apply Lemma~\ref{app:G1:Q-S} with the base point $\tilde B_j$, which is real and regular, and raw legs, using $\norm{a_j}_y,\norm{a'_j}_y\le P_1(t)$ and Step~1. This gives
\[
\abs{q_{j,y}(a)-q_{j,y}(a')}\le C_F\hat P^{\nbar}\mathrm{e}^{-\delta_1L^m\dist(y,\mathcal D)},
\]
with $C_F:=2\sum_{n\le\nbar}C_n/(n-1)!$ and $\hat P:=\max(1,P_1(t))$. The factor $2$ makes the bound hold whether the normalisation $-\frac12$ of the setting sits inside $q_{j,y}$ or in front of $\sum_j\sum_y$. In either case the left side of Lemma~\ref{app:K:L13-4} is at most $\sum_{j,y}\abs{q_{j,y}(a)-q_{j,y}(a')}$, and Step~2 is unaffected. By Lemma~\ref{app:G1:Qcoll-geo}(c), the number of scale-$j$ sites with $r/4+n<\dist(y,\mathcal D)\le r/4+n+1$ is at most $L^{3m}\,15^3(r+M_b)^3(n+1)^3\abs{\mathcal C}$. Moreover $\mathrm{e}^{-\delta_1L^mr/4}\le\mathrm{e}^{-\delta_1r/4}\mathrm{e}^{-\delta_1(L^m-1)/4}$, because $r\ge1$. The far sites therefore contribute at most
\[
C_F\hat P^{\nbar}15^3(r+M_b)^3\,c_L\,c_3\,\mathrm{e}^{-\delta_1r/4}\abs{\mathcal C},
\]
where $c_L:=\sum_{m\ge1}L^{3m}\mathrm{e}^{-\delta_1(L^m-1)/4}$ and $c_3:=\sum_{n\ge0}(n+1)^3\mathrm{e}^{-\delta_1n}$.

\emph{Step 4 (the $t^N$ bound).} We have $\ell(t)\le t^{-1}$, $p(t)=b_0\ell(t)^p$ and $P_1(t)=\abs e c_1C\mu_*^{-1}t^{-1/2}p(t)$ (Lemma~\ref{app:G1:RW-KSP}). Hence $\hat P^{\nbar}(r+M_b)^3\le C't^{-\nbar(p+1)-6}$. Also, for $N_1\ge0$, $\sup_{\ell\ge1}(N_1\ell-\frac{\delta_1R}4\ell^2)\le N_1^2/(\delta_1R)$ and $\mathrm{e}^{-N_1\ell(t)}\le t^{N_1}$, so $\mathrm{e}^{-\delta_1r/4}\le\mathrm{e}^{N_1^2/(\delta_1R)}t^{N_1}$. Taking $N_1=N+\nbar(p+1)+6$ bounds Step~3 by $C_Nt^N\abs{\mathcal C}$. In Step~2, apply Corollary~\ref{app:G1:Q-sum} with $N+6$ in place of $N$; since $(r+M_b)^3\le Ct^{-6}$, we get $2\cdot15^3(r+M_b)^3C_{N+6}t^{N+6}\le C'_Nt^N$.

\emph{Step 5.} The sites of $\Omega$ split into those of Step~2 and those of Step~3. Adding the two bounds gives the claim with $C_{\rm coll}:=16\cdot15^3C_Qc_{AX}^2C_B^2$, which depends only on the constants of Corollary~\ref{app:G1:Q-sum}. No constant depends on $k$: the sums over $j<k$ are bounded by sums over $m\ge1$.
\end{proof}

\subsection{Restriction constants on the slice line}\label{app:G1:ssec-restr-slice}

Fix $1\le k\le K$ and put $t:=L^{k-1}\eps$, in the units of \cite[(2.55), p.~570]{Ba2}. The vector mass, written $\mu_0$ in \cite{Ba2}, is $\mu_*$ here. For $c>0$ let $\mathcal R_k(c)$ be the set of real vector configurations $A$ on $T_1^{(k-1)}$ such that
\[
\abs{(\partial A)(b)}\le c\,p(t),\qquad \abs{A(x)}\le\frac{c}{\mu_*t}\,p(t)\qquad\text{for all bonds } b\subset\Lambda_{-1}^{(k-1)\prime}\text{ and all } x\in\Lambda_{-1}^{(k-1)\prime}.
\]
These are the vector restrictions of \cite[(2.55)]{Ba2}. There they are imposed by the vector factors of the characteristic functions $\chi_k$ of \cite[(2.43)--(2.45), pp.~566--567]{Ba2}. Those factors are defined as in \cite[(2.4), p.~557]{Ba2}, with $\eps$ replaced by $L^{k-1}\eps$ \cite[p.~567]{Ba2}. Write $\chi_k^{\rm v}$ for the product of these vector factors. By \cite[(2.4), p.~557]{Ba2} it is a product of indicators of conditions $N\le m$ or $N>m$, where $m>0$ is fixed and $N$ is one of $\abs{B(y)-(QA)(y)}$, $\abs{(\partial A)(b)}$, $\abs{A(x)}$. Each $N$ is absolutely homogeneous of degree one in the vector variables it reads; these are $A$ and, through $B-QA$, the block field $B$. Below, $A$ stands for this tuple of vector variables, and $sA$ means that every entry is multiplied by $s$. By \cite[(2.55)]{Ba2}, $\chi_k^{\rm v}(A)\ne0$ implies that the $A$-component lies in $\mathcal R_k(c_1)$; we write this as $A\in\mathcal R_k(c_1)$.

In the same way, for $c>0$ let $\mathcal R'_k(c)$ be the set of fluctuations $A'$ with $\abs{A'(x)}\le c\,p(L^k\eps)$ for $x\in\Lambda_5^{(k)}$. The characteristic function $\chi(A')$ of \cite[p.~581]{Ba2} imposes the restriction \cite[(2.86), p.~575]{Ba2}; that is, it is the indicator of $\mathcal R'_k(c_A)$, where $c_A$ is the constant $O(1)$ of (2.86).

On the slice line, the step-$j$ vector fluctuation has covariance $w_jC^{(j)}_{\Lambda_5^{(j)}}$ with $w_j=1+u\delta_j$, as in Lemma~\ref{app:G1:S05-SC}(C6). The vector characteristic functions there are Balaban's functions evaluated at rescaled arguments,
\[
\chi_k^{\rm v,sl}(A):=\chi_k^{\rm v}(s_kA),\qquad \chi^{\rm sl}(A'):=\chi(s_k'A'),
\]
with numbers $s_k,s_k'>0$. The lemma below holds for every choice with $s_k,s_k'\in[\sigma^{-1},\sigma]$. The convention of Lemma~\ref{app:G1:S05-SC}(C6), where $\chi(A')$ appears unrescaled, is $s_k=s_k'=1$: Balaban's functions on Balaban's variables \cite[(2.80), p.~574]{Ba2}. Another admissible choice is $s_k'=w_k^{-1/2}$ (and $s_k=w_{k-1}^{-1/2}$): Balaban's functions on the variable $X=w_k^{-1/2}A'$, whose law is Balaban's (part (iv) below). These two choices are different restrictions, not the same one in two coordinates.

\begin{lemma}[restriction constants on the slice line]\label{app:K:L13-10}\label{app:G1:restr-slice}
Let $0<w_-\le w_j\le w_+$ for all $j$ and $\sigma:=\max(w_+,w_-^{-1})^{1/2}$. Let $s_k,s_k'\in[\sigma^{-1},\sigma]$ for every $k$.
\begin{itemize}
\item[(i)] $\sigma\ge1$, and $1,w_j^{1/2},w_j^{-1/2}\in[\sigma^{-1},\sigma]$ for every $j$. In particular both choices above are admissible.
\item[(ii)] (Admitted only, not exactly.) If $\chi_k^{\rm v,sl}(A)\ne0$, then $A\in\mathcal R_k(c_1/s_k)\subset\mathcal R_k(\sigma c_1)$. If $\chi^{\rm sl}(A')\ne0$, then $A'\in\mathcal R'_k(c_A/s'_k)\subset\mathcal R'_k(\sigma c_A)$. So every configuration admitted by the slice-line vector characteristic functions obeys the vector restrictions of \cite[(2.55)]{Ba2} with $c_1$ replaced by $\sigma c_1$, and the restriction \cite[(2.86)]{Ba2} with its constant multiplied by $\sigma$.
\item[(iii)] Suppose $\chi_k^{\rm v}$ is the indicator of $\mathcal R_k(c_1)$. Then the set admitted by $\chi_k^{\rm v,sl}$ is exactly $\mathcal R_k(c_1/s_k)$, and
\[
\mathcal R_k(c_1/\sigma)\subset\mathcal R_k(c_1/s_k)\subset\mathcal R_k(\sigma c_1).
\]
The same holds for $\chi^{\rm sl}$, with $\mathcal R'_k$ and $c_A$ in place of $\mathcal R_k$ and $c_1$. For $s_k=\sigma^{-1}$ the admitted set is exactly $\mathcal R_k(\sigma c_1)$.
\item[(iv)] For every measurable $F\ge0$, every $s>0$ and every $w>0$,
\[
\int d\mu_{wC}(A')\,\chi(sA')F(A')=\int d\mu_C(X)\,\chi(sw^{1/2}X)\,F(w^{1/2}X),\qquad C=C^{(k)}_{\Lambda_5^{(k)}} .
\]
\item[(v)] If $w_j=1+u\delta_j$ with $\abs u\le1$ and $0\le\delta_j\le\frac1{12}$, then $\sigma\le(12/11)^{1/2}<\sqrt2$. On the slice line of Theorem~\ref{thm:T6} these hypotheses hold. Indeed $u\in\{1,-\eta\}$ with $0<\eta<1$, and $\delta_j\le\kappa\mu_*^2\eps_0^2\le\frac1{12}$ by \eqref{eq:T6-GammaB} and condition (E.$\delta$) of Appendix~\ref{app:B}.
\end{itemize}
\end{lemma}

\begin{proof}
(i) For $w\in[w_-,w_+]$ we have $w^{1/2}\le w_+^{1/2}\le\sigma$ and $w^{-1/2}\le w_-^{-1/2}\le\sigma$. Hence $w^{\pm1/2}\in[\sigma^{-1},\sigma]$. If $w_+\ge1$, then $\sigma\ge1$ directly. If $w_+<1$, then $w_-^{-1}\ge w_+^{-1}>1$, so again $\sigma\ge1$. Thus $1\in[\sigma^{-1},\sigma]$.

(ii) Each condition defining $\mathcal R_k(c)$ has the form $N(A)\le c\,m$, where $m>0$ is a fixed number and $N$ is absolutely homogeneous of degree one. There are two kinds of such $N$. The first is $N(A)=\abs{(\partial A)(b)}$; here $(\partial A)(b)$ is a real-linear function of the values of $A$. The second is $N(A)=\abs{A(x)}$, the Euclidean norm of the vector $A(x)$. For $s>0$ we have $N(sA)=sN(A)$, so $N(sA)\le c\,m$ holds if and only if $N(A)\le(c/s)m$. Hence
\[
\{A:\ sA\in\mathcal R_k(c)\}=\mathcal R_k(c/s).
\]
Now let $\chi_k^{\rm v,sl}(A)\ne0$. Then $\chi_k^{\rm v}(s_kA)\ne0$, so $s_kA\in\mathcal R_k(c_1)$, that is, $A\in\mathcal R_k(c_1/s_k)$. Since $s_k\ge\sigma^{-1}$ we have $c_1/s_k\le\sigma c_1$. The sets $\mathcal R_k(c)$ increase with $c$, which gives the inclusion. The argument for $A'$ is the same, with $N(A')=\abs{A'(x)}$ for $x\in\Lambda_5^{(k)}$.

(iii) If $\chi_k^{\rm v}=1_{\mathcal R_k(c_1)}$, then $\chi_k^{\rm v,sl}=1_{\{A:\,s_kA\in\mathcal R_k(c_1)\}}=1_{\mathcal R_k(c_1/s_k)}$, by the identity in (ii). Moreover $c_1/\sigma\le c_1/s_k\le\sigma c_1$, and monotonicity in $c$ gives both inclusions. For $s_k=\sigma^{-1}$ we get $c_1/s_k=\sigma c_1$.

(iv) If $X$ is centred Gaussian with covariance $C$, then $w^{1/2}X$ is centred Gaussian with covariance $wC$. So $d\mu_{wC}$ is the image of $d\mu_C$ under $X\mapsto w^{1/2}X$. Apply this to the non-negative function $A'\mapsto\chi(sA')F(A')$.

(v) We have $w_j\in[\frac{11}{12},\frac{13}{12}]$, so $w_+\le\frac{13}{12}$ and $w_-^{-1}\le\frac{12}{11}$. Since $13\cdot11=143<144$, we get $\frac{13}{12}<\frac{12}{11}$, and therefore $\sigma^2\le\frac{12}{11}<2$. The two bounds $\abs u\le1$ and $\delta_j\le\kappa\mu_*^2t_{j+1}^2\le\kappa\mu_*^2\eps_0^2$ hold by definition, using $t_{j+1}\le\eps_0$.
\end{proof}

\begin{corollary}[the use in Proposition~\ref{app:K:L13-11}(iv)]\label{app:G1:restr-slice-use}
Let $s_k,s_k'\in[\sigma^{-1},\sigma]$. The conclusion of Proposition~\ref{app:K:L13-11}(iv) holds with its hypothesis ``admit exactly'' replaced by ``admit only''. The latter is what part (ii) of Lemma~\ref{app:G1:restr-slice} provides.
\end{corollary}
\begin{proof}
By (ii), every vector configuration at which a slice-line integrand is non-zero lies in $\mathcal R_k(\sigma c_1)$, and its fluctuation lies in $\mathcal R'_k(\sigma c_A)$. The statements in question are the class $\mathcal R$ of Proposition~\ref{app:K:L13-11}(i), the regularity bounds of part (iii) of that proposition, and Propositions~\ref{app:K:L13-8} and~\ref{app:K:L13-9}. Each is proved for every configuration obeying pointwise restrictions with a given constant, and its constants depend on the configuration only through that constant. So each holds, with the constant $\sigma c_1$, at every configuration in $\mathcal R_k(\sigma c_1)$, and hence on the admitted subset. The proof of Proposition~\ref{app:K:L13-11}(iv) uses its hypothesis only in this way, through ``the class of admissible backgrounds''. The scalar characteristic functions are not rescaled: they read the vector configuration only as an external field, namely the configuration itself \cite[p.~568]{Ba2}.
\end{proof}
\begin{remark}
The reverse inclusion of part (iii) holds only under its indicator hypothesis, and it is not used. On the large-field factors $\chi(\{N>m\})$ evaluated at $s_kA$, the condition reads $N(A)>m/s_k\ge m/\sigma$, so the gain \cite[(2.3)]{Ba2} becomes $\exp(-c_0p^2/\sigma^2)$. Part (iv) does not identify the choices $s_k'=1$ and $s_k'=w_k^{-1/2}$. It rewrites each in Balaban's law: under $d\mu_C(X)$ they become $\chi(w_k^{1/2}X)$ and $\chi(X)$, that is, restriction constants $w_k^{-1/2}c_A$ and $c_A$.
\end{remark}

\subsection{The vector lines of a step}\label{app:G1:VL}
In this subsection $0\le k\le K-1$ is fixed. At step $i$ we use scale-$i$ units, and $C^{(i)}$, $C^{(i)}_\Lambda$, $G_i$, $Q_i$, $a_i$ denote Ba{\l}aban's operators at $u=0$. In the construction of Lemma~\ref{app:G1:S05-PH}, which we call the weighted chain, the vector fluctuation of step $i$ has the law $N(0,w_iC^{(i)}_{\Lambda_5^{(i)}})$ (Lemma~\ref{app:G1:S05-PH}(a), applied at every step $i\le k$). Ba{\l}aban's construction has $N(0,C^{(i)}_{\Lambda_5^{(i)}})$ instead. The translations, the lifts, the cut-off functions and the restriction $\chi(A')$ are those of Ba{\l}aban's construction. By a \emph{vector leg} we mean a linear functional of a vector fluctuation that occurs as a factor in a vertex. By a \emph{vector line} we mean the kernel by which a pair of vector legs of the step fluctuations ($A'$ of the current step, $A'_j$ of \cite[(1.1), (1.2)]{Ba3}) is contracted, or a kernel into which such a kernel is transformed. The propagators inside the counterterm coefficients $\delta m^2$, $E_1$ of \cite[(1.4), (1.26)--(1.28)]{Ba3} are not vector lines in this sense. They belong to the fixed, $u$-independent counterterm (Lemma~\ref{app:G1:S05-PH}(c)), stay at their $u=0$ value in the weighted chain, and are not covered by Lemma~\ref{app:K:L13-14}.

\begin{lemma}[three Gaussian facts]\label{app:G1:VL-gauss}
(i) Let $q$ be a positive definite quadratic form on $X\times Y$, let $s>0$, and let $r$ be a function on $Y$ such that $\exp(-\frac12sq(x,y)-r(y))$ is integrable. Under the normalised density proportional to this function, the conditional law of $x$ given $y$ is Gaussian. Its mean $m(y)$ does not depend on $s$, and its covariance is $s^{-1}$ times the covariance at $s=1$.
(ii) Let $HA=\sum_{i<K}w_i^{-1}I_i(A_i,A_{i+1})+D_K(A_K)$ be the form of Lemma~\ref{app:G1:S05-T}, with $D_K$ positive definite. Under the density proportional to $\mathrm{e}^{-\frac12HA}$, the variables $\hat S_i:=A_i-\mathcal H_iA_{i+1}$, $i<K$, and $A_K$ are independent, and $\hat S_i\sim N(0,w_iC^{(i)})$.
(iii) Let $Z=\sum_iN_iZ_i$, where the $Z_i$ are independent centred Gaussian vectors with covariances $\Gamma_i$ and the $N_i$ are linear maps. For linear functionals $\ell,\ell'$, $\E[\ell(Z)\ell'(Z)]=\sum_i\ell N_i\Gamma_iN_i^*\ell'^*$. The moments of $\ell_1(Z),\dots,\ell_{2n}(Z)$ are sums over pairings of products of such kernels, and the cumulants are the sums over the connected pairings.
\end{lemma}
\begin{proof}
(i) Write $q(x,y)=\ip{x}{Px}+2\ip{x}{Ry}+\ip{y}{Sy}$ with $P>0$. Completing the square gives $q=\ip{x-m(y)}{P(x-m(y))}+q'(y)$, where $m(y)=-P^{-1}Ry$. The conditional density of $x$ is therefore proportional to $\exp(-\frac s2\ip{x-m(y)}{P(x-m(y))})$: its mean is $m(y)$ and its covariance is $(sP)^{-1}$.
(ii) The map $(A_0,\dots,A_K)\mapsto(\hat S_0,\dots,\hat S_{K-1},A_K)$ is linear and unit triangular, so its Jacobian is $1$. By Lemma~\ref{app:G1:S05-T}, $I_i(A_i,A_{i+1})=\ip{\hat S_i}{(C^{(i)})^{-1}\hat S_i}$ (complete the square in $D_i(A_i)+\beta_i\norm{QA_i-A_{i+1}}^2$). Hence the density is $\prod_{i<K}\exp(-\frac12w_i^{-1}\ip{\hat S_i}{(C^{(i)})^{-1}\hat S_i})\cdot\exp(-\frac12D_K(A_K))$, which is a product.
(iii) By independence, $\E[\ell(Z)\ell'(Z)]=\sum_{i,i'}\ell N_i\E[Z_iZ_{i'}^*]N_{i'}^*\ell'^*$, and $\E[Z_iZ_{i'}^*]=\delta_{ii'}\Gamma_i$. The rest is Wick's rule for the Gaussian vector $(\ell_a(Z))_a$, together with the representation of cumulants by connected pairings (Lemma~\ref{app:G1:I5-lc}).
\end{proof}

\begin{definition}[step-graded kernels]\label{app:G1:VL-def}
A \emph{step-graded kernel} is a family $\mathcal L=(\mathcal L_i)_{0\le i\le k}$ of real matrices indexed by vector-field components. Its value at the weights $w=(w_i)$ is $\mathcal L[w]:=\sum_iw_i\mathcal L_i$, and $\mathcal L[1]$ is its value at $u=0$. The \emph{elementary} families are $\mathcal E^{(i)}_\Lambda:=(\delta_{ij}C^{(i)}_\Lambda)_{0\le j\le k}$ with $\Lambda=\Lambda_5^{(i)}$ or $\Lambda=T^{(i)}$ (then $C^{(i)}_{T^{(i)}}=C^{(i)}$). The \emph{admissible operations} are:
\begin{itemize}
\item[(O1)] push-forward $(\mathcal L_i)\mapsto(N_i\mathcal L_iN_i'^*)$, where $N_i,N_i'$ are compositions of the following $w$-free linear maps: the linear parts of the translation \cite[(2.80)]{Ba2} and of the lifts \cite[(2.88), (2.89)]{Ba2}, \cite[(1.2)]{Ba3}; the conditional means; the maps $A\mapsto A(b)$, $A\mapsto A(\Gamma)$ (contour sums), $Q_i$ and the lattice derivatives, which build the vector legs of the vertices \cite[(1.8)--(1.15), (1.17)]{Ba3}; and the rescalings between lattices;
\item[(O2)] sums $\mathcal L+\mathcal L'$. The composition \cite[(2.42), p.~612]{Ba1}, $G_{k+1}=a_k^2(L^k\varepsilon)^{-4}G_kQ_k^*C^{(k)}Q_kG_k+G_k$, is (O1) applied to $\mathcal E^{(k)}$ with the $w$-free maps $a_k(L^k\varepsilon)^{-2}G_kQ_k^*$ of Ba{\l}aban's $u=0$ operator $G_k$, followed by (O2). The same holds for its vector form (the equality after \cite[(2.6), p.~424]{Ba3}; \cite[(I.2.43), cited on p.~414]{Ba3});
\item[(O3)] restriction to domains, $\mathcal L_i\mapsto1_X\mathcal L_i1_Y$;
\item[(O4)] Dirichlet-to-free replacement, i.e.\ replacing an elementary family $\mathcal E^{(i)}_{\Lambda_5^{(i)}}$ by $\mathcal E^{(i)}_{T^{(i)}}$, or by the difference $\mathcal E^{(i)}_{\Lambda_5^{(i)}}-\mathcal E^{(i)}_{T^{(i)}}$;
\item[(O5)] localisation cut-offs, $\mathcal L_i(x,y)\mapsto g(x)\mathcal L_i(x,y)g'(y)$ with $w$-free functions $g,g'$ (the functions $\theta_{i+1}$, $\zeta^{(i)}$ of \cite[(2.88), (2.89)]{Ba2} and $g_i$ of \cite[p.~419]{Ba3});
\item[(O6)] differences, and splitting by the index, $\mathcal L\mapsto(1\{i\in I\}\mathcal L_i)_i$.
\end{itemize}
\end{definition}

\begin{lemma}[weights pass through the operations]\label{app:G1:VL-alg}
Let $\mathcal L$ be obtained from elementary families by finitely many admissible operations. If the same operations are applied to the weighted elementary families $(\delta_{ij}w_iC^{(i)}_\Lambda)_j$, the result is $\mathcal L[w]$, i.e.\ its component $i$ is $w_i\mathcal L_i$.
\end{lemma}
\begin{proof}
We use induction on the number of operations. For elementary families the claim is the definition. Each of (O1)--(O3), (O5) and (O6) acts on the component $i$ by a $w$-free linear map, so it commutes with multiplication of that component by $w_i$. (O4) replaces $w_iC^{(i)}_{\Lambda_5^{(i)}}$ by $w_iC^{(i)}$, a component with the same index and the same weight.
\end{proof}

\begin{lemma}[the vector operations of a step]\label{app:K:L13-14}
Every vector line occurring in the step-$k$ construction of \cite[Sect.~2C]{Ba2} and \cite{Ba3} is obtained from the step covariances $C^{(i)}_{\Lambda_5^{(i)}}$, $i\le k$, by finitely many of the following operations: push-forward by the $w$-free linear maps of the translations, conditional means and lifts (together with the leg maps and rescalings of (O1)); composition \cite[(2.42)]{Ba1}; restriction to domains; Dirichlet-to-free replacement; localisation cut-offs; differences. More precisely:
\begin{itemize}
\item[(a)] Every vector line of the step-$k$ construction is $\mathcal L[1]$ for some step-graded $\mathcal L$ obtained from elementary families by finitely many admissible operations.
\item[(b)] In the weighted chain the line at the same place is $\mathcal L[w]=\sum_iw_i\mathcal L_i$.
\item[(c)] After the decomposition by indices \cite[(2.6), p.~424; (3.4), p.~434]{Ba3}, every vector line carries one index $j_l$, and in the weighted chain it equals $w_{j_l}$ times its value at $u=0$. Hence every indexed expression $E(G(\mathbf j))$ of the weighted chain equals $\prod_lw_{j_l}$ times Ba{\l}aban's, where the product runs over the vector lines of $G$. Every identity or bound proved for a single indexed expression $E(G(\mathbf j))$ at $u=0$ therefore holds in the weighted chain after multiplication by this product. This does not apply to identities between sums over a renormalised class, such as \cite[(3.5)]{Ba3}, whose members have different sets of vector lines.
\end{itemize}
\end{lemma}

\begin{proof}
\emph{Step 1: the vector legs.} After the translation \cite[(2.80), p.~574]{Ba2}, $A=A'+aL^{-2}C^{(k)}_{\Lambda_0^{(k)}}Q^*B$, the quadratic form in $A'$ restricted to $\Lambda_5^{(k)}$ is $\frac12\ip{\Lambda_5^{(k)}A'}{(C^{(k)}_{\Lambda_5^{(k)}})^{-1}\Lambda_5^{(k)}A'}$. The remaining terms of the action depend on $A'$ only through the field $A^{(k)}$ \cite[(2.87) and the text after it, p.~575]{Ba2}. By \cite[(2.88), (2.89), p.~575]{Ba2}, $A^{(k)}$ is affine in $A'$. In particular $\theta_{k+1}A^{(k)}=a_k\theta_{k+1}\zeta^{(k)}G_kQ_k^*A'+\tilde B^{(k)}+\theta_{k+1}B^{(k+1),\eta}$, where only the first term contains $A'$. The expansion is made with respect to the small field $\theta_{k+1}A'^{(k)}+\tilde B^{(k)}$ \cite[p.~576]{Ba2}, \cite[p.~412]{Ba3}. After all transformations of the step, the fluctuation integral of step $k$ is
\[
\int d\mu_{C^{(k)}_{\Lambda_5^{(k)}}}(A')\int d\mu(\varphi')\,\chi(A')\chi(\varphi')\exp V^{(k)}(\Lambda_7^{(k)},B^{(k+1),\eta},\psi,g_kA'^{(k)},\varphi'),
\]
and all the dependence on $A'|_{\Lambda_5^{(k)}}$ lies inside it \cite[(2.115), p.~581]{Ba2}. The argument $g_kA'^{(k)}$ is that of the display at the top of \cite[p.~581]{Ba2}; (2.115) prints $A'^{(k)}$. The $A'$-integration in (2.115) is conditional on $\Lambda_5^{(k)c}$, and its conditional means are those of the quadratic form (2.87). Hence every vector leg created at step $k$ has the form $\ell(A')=N\,g_k\,a_k\zeta^{(k)}G_kQ_k^*A'$, where $N$ is one of the leg maps of (O1) \cite[(1.8)--(1.15), pp.~413--414]{Ba3}. Every factor of this map is $w$-free.

The vector field of the earlier steps enters in the same way. By \cite[(1.1), (1.2), p.~412]{Ba3}, the field at scale $k$ is $A=\sum_{j<k}A'^{(j),\eta}+A^{(k)}$ with $A'^{(j),\eta}=a_j(L^j\eta)^{-2}G^\eta_jQ_j^*A'_j$, and $E_k$ in \cite[(1.4)]{Ba3} depends on the $A'_j$ only through $\mathrm{e}'g_kA'+A^{(k)}$, under the product measure $\prod_jd\mu_{C^{(j),L^j\eta}}(A'_j)$. So every vector leg of \cite[(1.5)]{Ba3} is $\ell(A')=\sum_jN\,g_k\,M_jA'_j$ with $M_j:=a_j(L^j\eta)^{-2}G^\eta_jQ_j^*$, which is again $w$-free.

\emph{Step 2: contraction.} The restriction $\chi(A')$ is removed in the cumulant bound, whose cumulants are taken for the Gaussian measure without $\chi$ (Theorem~\ref{app:K:L13-1}). In \cite[(1.4)]{Ba3} no restriction is present. By Lemma~\ref{app:G1:VL-gauss}(iii), the pair $\ell,\ell'$ is contracted at step $k$ by $\ell C^{(k)}_{\Lambda_5^{(k)}}\ell'^*$. In \cite[(1.4)]{Ba3} it is contracted by $\sum_j\ell_jC^{(j)}\ell_j'^*$ with $\ell_j=Ng_kM_j$. By \cite[(2.6), p.~424]{Ba3} and the sentence after it, the latter sum is the composed vector propagator $G_k=\sum_jG_{(j)}$. Ba{\l}aban states this on \cite[p.~414]{Ba3}: the $A'_j$ are independent with covariances $C^{(j),L^j\eta}$, the field $A'$ has the covariance $G_k$, and all $A'$-legs are contracted in pairs by this propagator. Both kernels are $\mathcal L[1]$ for families obtained from elementary families by (O1) and (O2).

\emph{Step 3: the transformations in print.} In \cite[Sect.~2C]{Ba2} the vector lines are transformed by the following operations.
\begin{itemize}
\item The interaction is removed from $B^k(\Lambda_7^{(k-1)\prime}\cap\Lambda_7^{(k)c})$ \cite[Prop.~2.1, p.~570]{Ba2}, and the vertex summations are restricted to $B^k(\Lambda_7^{(k)})$ or cut off by $g_k$ \cite[p.~581]{Ba2}. These are (O3) and (O5).
\item The terms of order higher than $\nbar$ are estimated. This deletes terms and creates no kernel.
\item The covariances with Dirichlet boundary conditions are replaced by the free ones, and some of them are composed by \cite[(2.42)]{Ba1} \cite[p.~582]{Ba2}. These are (O4) and (O2).
\item The result is rescaled \cite[p.~582]{Ba2}. This is (O1).
\end{itemize}
In \cite{Ba3}, before the decomposition by indices, the following operations occur.
\begin{itemize}
\item The localisation by $g_k$ on the vector legs \cite[p.~419]{Ba3}. This is (O5).
\item The decomposition of $G_K$ and $G_k$ in the counterterms \cite[(1.27), (1.28), p.~418]{Ba3}, the extraction of the terms with index $k$ \cite[p.~422]{Ba3}, and the treatment of their index-$k$ lines as external fields \cite[p.~422]{Ba3} act only on counterterm coefficients. These contain no vector line in the sense above (they are $u$-independent in the weighted chain), so no vector line is changed.
\item The lines treated as external fields in \cite[(2.5), p.~424]{Ba3} are $\delta G_k$ and \cite[(1.16)]{Ba3}, which are scalar.
\item The changes of domain $\Omega\to\Omega_2$ and $\Omega\to\Pi$, the operators $\delta G_k$ and \cite[(1.16)]{Ba3}, the expansion $B=B_0+B'$, the gauge transformation removing $B_0$ and the replacement $G_k(\Pi,0)\to G_k(0)$ \cite[p.~414; p.~433]{Ba3} act on scalar propagators only. The vector propagator is $G_k$ throughout \cite[(1.17), p.~415]{Ba3}, and it does not depend on the background.
\end{itemize}
Each listed operation maps a kernel of the form $\mathcal L[1]$ to one of the same form, with $\mathcal L$ changed by an admissible operation. Together with Steps~1 and~2 this proves (a).

\emph{Step 4: the weighted chain.} By Lemma~\ref{app:G1:S05-PH}(a), the weighted chain differs from Ba{\l}aban's construction in the vector sector only in the laws $N(0,w_iC^{(i)}_{\Lambda_5^{(i)}})$. By Lemma~\ref{app:G1:S05-PH}(a), the vector quadratic form of step $i$ in the weighted chain is $w_i^{-1}$ times Ba{\l}aban's form \cite[(2.87)]{Ba2}. Hence, by Lemma~\ref{app:G1:VL-gauss}(i) with $s=w_i^{-1}$, its minimiser, which is the translation (2.80), and the conditional means of the conditioning on $\Lambda_5^{(i)c}$ in (2.115) are unchanged. In the form of \cite[(1.4)]{Ba3} the measure is, by construction, the product of the step measures. In the weighted chain the factor of step $j$ is $d\mu_{w_jC^{(j)}}$, by Lemma~\ref{app:G1:S05-PH}(a) at step $j$ followed by (O4). (Lemma~\ref{app:G1:VL-gauss}(ii) gives the same product structure for the full law $\mathrm{e}^{-\frac12HA}$.) So the operations of Steps~1--3 are applied, unchanged, to the weighted elementary families, and Lemma~\ref{app:G1:VL-alg} gives (b).

\emph{Step 5: indexed expressions.} After \cite[(2.6)]{Ba3}, applied in \cite[(3.4), p.~434]{Ba3} to the vector propagator $G_k$, each vector line $l$ of $G(\mathbf j)$ is $N_lG_{(j_l)}N_l'^*$, a family supported on the single index $j_l$. By (b), its weighted value is $w_{j_l}N_lG_{(j_l)}N_l'^*$. The value $E(G(\mathbf j))$ is a sum, over the positions of the vertices, of products in which every line of $G$ occurs exactly once as a factor, so it is linear in each line separately. Multiplying the line $l$ by the number $w_{j_l}$ therefore multiplies $E(G(\mathbf j))$ by $w_{j_l}$. This gives the first claim of (c). Any rewriting of $E(G(\mathbf j))$ used in \cite[Sects.~2--3]{Ba3} is an identity between numbers, for instance the sums over orderings, Taylor's formula \cite[(3.10)]{Ba3} or the integration by parts \cite[(2.8)]{Ba3}. Such an identity remains true after both sides are multiplied by the constant $\prod_lw_{j_l}$, and so does every bound. This gives the second claim of (c). The class identities \cite[(3.5)]{Ba3} compare graphs with counterterm partners, whose lines differ, and the counterterm coefficients are $u$-independent. They are not covered here; the paper treats them by the colour expansion in Lemma~\ref{app:G1:S05-PH}(c).
\end{proof}

\noindent In particular, the line $\mathcal L=\sum_i\mathcal L_i$ of Lemma~\ref{app:G1:PH-WL} becomes $\sum_iw_i\mathcal L_i=\mathcal L+u\sum_i\delta_i\mathcal L_i$ in the weighted chain. This holds because $w_i=1+u\delta_i$ with $\delta_i=\kappa\mu_*^2t_{i+1}^2$, and $c_k=\kappa\mu_*^2L^2t_k^2=\kappa\mu_*^2t_{k+1}^2$. So $\delta_i\le c_k$ for all $i\le k$, with equality at $i=k$ (the definitions at the start of the proof of Lemma~\ref{app:G1:S05-PH}(b),(c),(e) and in Lemma~\ref{app:G1:S05-PH}(c)). Taking all indices $\le k$ in Lemma~\ref{app:K:L13-14}(b) gives the new photon line $\sum_{j\le k}w_jG_{(j)}$ used in the proof of Lemma~\ref{app:G1:S05-PH}(c).

\subsection{The weighted-field structure of the slice model}\label{app:G1:WFS}

Throughout this subsection $t_j=L^j\eps$. For $j<K$ put $\delta_j=\kappa\mu_*^2t_{j+1}^2$ and $w_j=1+u\delta_j$, and put $w_K:=1$. Let $\omega_j=w_j^{-1}-w_{j-1}^{-1}$ be as in Lemma~\ref{app:G1:S05-T}, and $c_n=\kappa\mu_*^2L^2t_n^2$. The photon covariance of the slice model is $C_B(u)=C_{\mu_*}+u\Gamma^B$ with $\Gamma^B=\sum_{j<K}\delta_jH_j$ \eqref{eq:T6-GammaB}. Here $H_j=G_{(j)}$ are King's slices \cite[(2.17), p.~653]{K1}, taken for the vector field ($m^2=\mu_*^2$, $N=d$, $A=0$). The vector objects are $\Delta^{(0)}:=-\Delta_\eps+\mu_*^2$, the block actions $\Delta^{(j)}$, $\beta_j:=at_{j+1}^{-2}$, and $C^{(j)}=(\Delta^{(j)}+\beta_jQ^*Q)^{-1}$ \cite[(2.14), (2.16), p.~653]{K1}. The lifts are $F_0:=1$ and $F_j:=a_jt_j^{-2}G_jQ_j^*$ \cite[(2.18), p.~654]{K1}, so that $H_j=F_jC^{(j)}F_j^*$. For a set $\Lambda$, $C^{(j)}_\Lambda$ denotes the covariance of \cite[(2.87), p.~575]{Ba2}. That is, the $\Lambda$-block of the quadratic form in the translated field is $\frac12\ip{\Lambda A'}{(C^{(j)}_\Lambda)^{-1}\Lambda A'}$, and $C^{(j)}$ is its free version.

\begin{lemma}[Gaussian integration with a weight]\label{app:G1:WFS-gauss}
Let $P>0$ on $\R^p$, let $J:\R^q\to\R^p$ be linear and $R$ symmetric on $\R^q$, and put $q(x,y)=\ip{x}{Px}-2\ip{x}{Jy}+\ip{y}{Ry}$. Put $m(y):=P^{-1}Jy$ and $q_*(y):=\min_xq(x,y)=\ip{y}{(R-J^*P^{-1}J)y}$. Then for $w>0$ and every measurable $F\ge0$,
\[
\int F(x,y)\,\exp\big(-\tfrac1{2w}q(x,y)\big)\,dx=(2\pi w)^{p/2}(\det P)^{-1/2}\exp\big(-\tfrac1{2w}q_*(y)\big)\int F(m(y)+x',y)\,d\mu_{wP^{-1}}(x').
\]
So the minimiser $m(y)$ and the minimised form $q_*$ do not depend on $w$, and the fluctuation $x'$ has covariance $wP^{-1}$. If a later expression depends on $x'$ only through $Mx'$ and $M'x'$ for fixed linear maps $M,M'$, then $\int(Mx')(M'x')^T\,d\mu_{wP^{-1}}=w\,MP^{-1}M'^*$.
\end{lemma}

\begin{proof}
We have $q(x,y)=\ip{x-m(y)}{P(x-m(y))}+q_*(y)$. Substitute $x=m(y)+x'$ and use $\int\exp(-\frac1{2w}\ip{x'}{Px'})dx'=(2\pi w)^{p/2}(\det P)^{-1/2}$. The last claim is the covariance of $N(0,wP^{-1})$ pushed forward by $M$ and $M'$.
\end{proof}

\begin{lemma}[realisation of the slice photon]\label{app:G1:WFS-real}
Let $S_j\sim N(0,w_jC^{(j)})$ for $j<K$, independent of each other and of a field $A_K$ distributed as King's average field of scale $K$ at $u=0$ \cite[(2.18), p.~654]{K1}. Then $A:=\sum_{j<K}F_jS_j+F_KA_K$ has covariance $C_B(u)$. For $n\le K$, the part $\sum_{j<n}F_jS_j$ has covariance $\sum_{j<n}w_jH_j=G_n+u\sum_{j<n}\delta_jG_{(j)}$.
\end{lemma}

\begin{proof}
At $u=0$, \cite[(2.17)--(2.18), pp.~653--654]{K1} write the vector field of covariance $C_{\mu_*}$ as the sum of the independent fields $F_jA'_j$ ($j<K$) and $F_KA_K$, with $\mathrm{Cov}(F_jA'_j)=G_{(j)}=H_j$. Hence $F_K\mathrm{Cov}(A_K)F_K^*=C_{\mu_*}-\sum_{j<K}H_j$. For general $u$, independence gives $\mathrm{Cov}(A)=\sum_jw_jH_j+C_{\mu_*}-\sum_jH_j=C_{\mu_*}+u\Gamma^B$. The second claim follows from the same computation together with $\sum_{j<n}G_{(j)}=G_n$ \cite[(2.17)]{K1}.
\end{proof}

\begin{lemma}[the vector part at every stage]\label{app:G1:WFS-skel}
Fix $(u,M)$, $0\le k\le K-1$ and an admissible family $\boldsymbol\Lambda_k$, and put $\Lambda_6:=\Lambda_6^{(k-1)\prime}$. Hold all variables other than $A_k$ fixed. As a function of $A_k$, the integrand of \eqref{app:G1:eq-IH} at stage $k$ equals
\[
\exp\Big(-\tfrac12w_k^{-1}\big[\ip{\Lambda_6A_k}{\Delta^{(k)}\Lambda_6A_k}+\beta_k\norm{QA_k-A_{k+1}}^2\big]\Big)\,\chi_k(A_k)\,\Psi_k .
\]
Here $\chi_k$ is a product of characteristic functions that are the same functions of the fields as at $u=0$. The factor $\Psi_k$ depends on $A_k$ only through $\Lambda_7^{(k-1)\prime c}A_k$ and through the configuration $A^{(k)}$ of \cite[(2.44), p.~566]{Ba2}.
\end{lemma}

\begin{proof}
We use induction on $k$.

\emph{Base case $k=0$.} We have $\Lambda_6^{(-1)\prime}=T_\eps$. By Lemmas~\ref{app:G1:S05-T} and~\ref{app:G1:S05-B0}, the integrand is $\exp(-\frac12w_0^{-1}\ip{A_0}{\Delta^{(0)}A_0})\prod_{j\ge0}\mathcal K^w_j$ times the scalar factor. The kernel $\mathcal K^w_0$ supplies $w_0^{-1}\beta_0\norm{QA_0-A_1}^2$. Its spectator $\omega_1\ip{A_1}{\Delta^{(1)}A_1}$ does not depend on $A_0$, and neither do the kernels $\mathcal K^w_j$ with $j\ge1$. The scalar factor depends on $A_0$ through the covariant Laplacian, that is, through $A^{(0)}=A_0$.

\emph{Inductive step $k\to k+1$.} Write $B:=A_{k+1}$, $\Lambda_6':=\Lambda_6^{(k)\prime}$ and $\Lambda_7':=\Lambda_7^{(k)\prime}$. The step consists of the operations of \cite[pp.~569--582]{Ba2}, in that order.
\begin{enumerate}
\item[(i)] Since $\Lambda_0^{(k)}\subset\Lambda_7^{(k-1)\prime}$ \cite[p.~566]{Ba2}, the part of the exponent that depends on $\Lambda_0^{(k)}A_k$ is $w_k^{-1}$ times the left side of \cite[(2.87), p.~575]{Ba2}. Every other factor depends on $\Lambda_0^{(k)}A_k$ only through $A^{(k)}$ \cite[p.~575, after (2.87)]{Ba2} and through $\chi_k$.
\item[(ii)] By Lemma~\ref{app:G1:WFS-gauss}, the translation \cite[(2.80), p.~574]{Ba2}, which is the minimiser, is the same as at $u=0$, and \cite[(2.87)]{Ba2} holds multiplied by $w_k^{-1}$. The configuration $A^{(k)}$ is expressed through $A'$ and $B$ by the $u$-independent linear maps \cite[(2.88)--(2.89), p.~575]{Ba2}.
\item[(iii)] Next comes the conditional integration over $\Lambda_5^{(k)}A'$ with conditioning on $\Lambda_5^{(k)c}$ \cite[(2.115), p.~581]{Ba2}. By Lemma~\ref{app:G1:WFS-gauss} it has measure $d\mu_{w_kC^{(k)}_{\Lambda_5^{(k)}}}$, a $u$-independent mean, and a remaining Gaussian form equal to $w_k^{-1}$ times the form of the case $u=0$.
\item[(iv)] At $u=0$, the subsequent operations turn this form into $\frac12\ip{\Lambda_6'B}{\Delta^{(k+1)}\Lambda_6'B}$ plus terms that depend on $B$ only through $\Lambda_7'^cB$ and $A^{(k+1)}$ \cite[(2.47)--(2.50), pp.~567--568]{Ba2}, applied at $k+1$, and \cite[p.~582]{Ba2}. By \cite[(2.50)]{Ba2}, the $A^{(k+1)}$-dependence enters only through the scalar factors, which carry no weight. The operations on the vector form rewrite a quadratic form and estimate parts of it, separately from the scalar forms. Applied to $w_k^{-1}$ times the form, they give $w_k^{-1}$ times the same terms, with $w_k^{-1}\le\sigma_+$ times the printed errors; this is the cost statement of Lemma~\ref{app:G1:S05-PH}(d).
\item[(v)] The fluctuation integral gives a function of $B$ through $B^{(k+1),\eta}$ only \cite[p.~582]{Ba2}.
\item[(vi)] Add the spectator $\frac12\omega_{k+1}\ip{B}{\Delta^{(k+1)}B}$ of $\mathcal K^w_k$. Split
\[
\ip{B}{\Delta^{(k+1)}B}=\ip{\Lambda_6'B}{\Delta^{(k+1)}\Lambda_6'B}+2X_{k+1}(B)+\Sigma_{k+1}(B),
\]
where $X_{k+1}(B)=\ip{\Lambda_7'B}{\Delta^{(k+1)}\Lambda_6'^cB}$. The remainder $\Sigma_{k+1}$ collects the terms with an argument in $\Lambda_6'^c$ or in $\Lambda_6'\setminus\Lambda_7'$, so it depends on $B$ only through $\Lambda_7'^cB$. Since $w_k^{-1}+\omega_{k+1}=w_{k+1}^{-1}$, the $\Lambda_6'$-block carries the coefficient $w_{k+1}^{-1}$.
\item[(vii)] The cross term $\omega_{k+1}X_{k+1}$ is removed by Lemma~\ref{app:G1:S05-PH}(d). This costs a $c$-number and a factor that depends on $B$ only through $\Lambda_6'^cB\subset\Lambda_7'^cB$.
\end{enumerate}
Finally, $\mathcal K^w_{k+1}$ supplies $w_{k+1}^{-1}\beta_{k+1}\norm{QB-A_{k+2}}^2$. The characteristic functions of step $k+1$ are defined by conditions on the fields \cite[(2.43), p.~566]{Ba2}, so they do not depend on $w$. This is the claim at $k+1$.
\end{proof}

\begin{lemma}[the weighted-field structure of the slice model]\label{app:K:L14-6}\label{app:G1:WFS-main}
Let $0<\eta<1$, $u\in[-\eta,1]$ and let $n\le K$ be a scale. In the slice model at $(u,M)$, with the construction of the proof of Proposition~\ref{app:G1:assembly}, the following holds for every step $i<n$.
\begin{itemize}
\item[(a)] $0\le\delta_i\le\kappa\mu_*^2t_n^2=L^{-2}c_n\le c_n$. In scale-$n$ units (lattice spacing $\eta_n:=L^{-n}$; this is the $\eta$ of Lemma~\ref{app:K:L14-5}, not the deformation parameter $\eta$ of $u\in[-\eta,1]$), $\delta_i=c_n(L^i\eta_n)^2=c_n(t_i/t_n)^2$. Moreover $w_i\ge1-\eta\kappa\mu_*^2\eps_0^2\ge\frac{11}{12}>0$.
\item[(b)] The translation of step $i$, its conditional means, the lifts $F_j$, the configuration $A^{(i)}$ as a function of $(A',B)$, and the characteristic function $\chi(A')$ \cite[(2.86), p.~575; p.~581]{Ba2} are the maps of the case $u=0$. The step-$i$ vector fluctuation on its integration set has law $N(0,w_iC^{(i)}_{\Lambda_5^{(i)}})$, where $C^{(i)}_{\Lambda_5^{(i)}}$ is the covariance of the case $u=0$. Its free version is $w_iC^{(i)}$.
\item[(c)] Let $\mathcal L$ be a vector line used at scale $n$: a composed propagator, a domain-restricted step covariance, a Dirichlet-to-free difference, a localised piece, a cut line, or a difference of such objects. Let $\mathcal L=\sum_{i<n}\mathcal L_i$ be its decomposition at $u=0$ by step of origin \cite[(2.6), p.~424]{Ba3}. Then the slice-model line at the same place is $\sum_iw_i\mathcal L_i$, and its $D$-part is $u\sum_i\delta_i\mathcal L_i$. In particular the composed line is $G_n+uc_nD^{\eta_n}_n$, with $D^{\eta_n}_n=\sum_{j<n}(L^j\eta_n)^2G^{\eta_n}_{(j)}$, $\eta_n=L^{-n}$; this is the $D$-line $uc_nD^\eta_n$ of Lemma~\ref{app:K:L14-5}.
\item[(d)] The bounds \cite[(2.6), (2.10)--(2.12), pp.~424--426]{Ba3} for $\mathcal L_i$ hold with the constants of the case $u=0$. The index-$i$ part $w_i\mathcal L_i$ obeys them with the extra factor $w_i$, and the $D$-part $u\delta_i\mathcal L_i$ obeys them with the extra factor $\abs u\delta_i\le\abs uc_n$.
\end{itemize}
\end{lemma}

\begin{proof}
(a) For $i<n\le K$ we have $t_{i+1}\le t_n\le t_K=\eps_0$, so $0\le\delta_i=\kappa\mu_*^2t_{i+1}^2\le\kappa\mu_*^2t_n^2=L^{-2}c_n$. In scale-$n$ units $L^i\eta_n=L^{i-n}=t_i/t_n$, so $c_n(L^i\eta_n)^2=\kappa\mu_*^2L^2t_i^2=\kappa\mu_*^2t_{i+1}^2=\delta_i$. Also $\delta_i\le\kappa\mu_*^2\eps_0^2\le\frac1{12}$ by condition (E.$\delta$) of Definition~\ref{app:G1:rowlist}. Since $\delta_i\ge0$ and $u\ge-\eta$ with $0<\eta<1$, $w_i=1+u\delta_i\ge1-\eta\delta_i\ge1-\eta\kappa\mu_*^2\eps_0^2\ge1-\frac1{12}$.

(b) This is steps (i)--(iii) of the proof of Lemma~\ref{app:G1:WFS-skel} at $k=i$, given that lemma at stage $i$. The function $\chi(A')$ is defined by the bounds \cite[(2.86)]{Ba2} on $\Lambda_5^{(i)}$, which are conditions on the field \cite[p.~581]{Ba2}. The lifts are fixed operators \cite[(2.18)]{K1}, and $F_j\mathcal H_j=F_{j+1}$ by \cite[(3.42), p.~619]{Ba1}. The free version of the step-$i$ precision is $w_i^{-1}(\Delta^{(i)}+\beta_iQ^*Q)$. Its inverse is $w_iC^{(i)}$, and its $\Lambda_5^{(i)}$-block has inverse $w_iC^{(i)}_{\Lambda_5^{(i)}}$, since by Lemma~\ref{app:G1:WFS-skel} the whole step-$i$ form is $w_i^{-1}$ times the form of the case $u=0$.

(c) At $u=0$, a vector line at scale $n$ arises in two stages. First, a contraction of two step-$i$ vector legs ($i<n$) in the perturbative part of the cumulant expansion of the fluctuation integral \cite[p.~582]{Ba2}. Such a contraction is a Wick pairing for the Gaussian measure, of the form $M_1C^{(i)}_{\Lambda}M_2^*$. Here $M_1,M_2$ are the maps through which $A'$ enters the vertices: $A'^{(i)}$ by \cite[(2.88)--(2.89)]{Ba2}, followed by the lifts. Second, the following operations act on these contractions: replacement of Dirichlet covariances by free ones, composition by \cite[(I.2.42)]{Ba2} (that is, summation over the index), and the decomposition \cite[(2.6), p.~424]{Ba3} ("the similar equality for the vector field propagator"); then localisation by the functions $h,h'$, cutting of a line into an external field, and differences of such objects \cite[pp.~419--426]{Ba3}. Each of these operations is linear in every line, and each acts on the single-index pieces separately. In the slice model the same maps occur, by (b). By Lemma~\ref{app:G1:WFS-gauss}, each step-$i$ contraction equals $w_iM_1C^{(i)}_\Lambda M_2^*$, and its free version equals $w_iM_1C^{(i)}M_2^*$; the Dirichlet-to-free difference is $w_i$ times that of the case $u=0$. Linearity and the index-wise action give $\mathcal L(u)=\sum_iw_i\mathcal L_i$, hence $\mathcal L(u)-\mathcal L=u\sum_i\delta_i\mathcal L_i$. For the composed line, Lemma~\ref{app:G1:WFS-real} and (a) give, in scale-$n$ units, $\sum_{j<n}w_jG_{(j)}=G_n+uc_n\sum_{j<n}(L^j\eta_n)^2G^{\eta_n}_{(j)}$. The $\chi$-corrections and the terms of order $>\nbar$ are not lines \cite[p.~582]{Ba2}.

(d) $\mathcal L_i$ is the index-$i$ part of the case $u=0$, so \cite[(2.10)--(2.12)]{Ba3} hold for it as printed ("the same inequalities hold for vector field propagators", p.~426). By (c), the slice-model index-$i$ part is $w_i\mathcal L_i$ and its $D$-part is $u\delta_i\mathcal L_i$. These are scalar multiples, so the bounds hold with the factors $w_i$ and $\abs u\delta_i\le\abs uc_n$, the last by (a).
\end{proof}

\subsection{The inputs of the expansion of the normalisations}\label{app:G1:ssec-Zinputs}

Throughout this subsection, $0\le j<k\le K-1$ and $m=k-j$. Distances are those of the torus sup-metric \cite[(1.3)]{Ba1}. A length of $1$ in scale-$k$ units is a length of $L^m$ in scale-$j$ units. Blocks are $B(y)=\{x:y_\mu\le x_\mu<y_\mu+L\eta\}$ \cite[(1.17), p.~606]{Ba1}, and $B^i(y)$ is the iterated block. We write $r_i=r(t_i)$ and $M_b$ for the side of a big block. We use the following facts of the construction.
\begin{itemize}
\item[(N1)] \cite[pp.~566--567]{Ba2}: $\theta_{k+1}$ is defined on $T_\eta$, $\eta=L^{-k}$. It equals $1$ on $B^k(\Lambda_2^{(k)})$, takes values in $[0,1]$, and vanishes at scale-$k$ distance $\ge M_b$ from $B^k(\Lambda_2^{(k)})$.
\item[(N2)] \cite[(2.89), p.~575]{Ba2}: $\theta_{k+1}A'^{(k)}=a_k\theta_{k+1}\zeta^{(k)}G_kQ_k^*A'$. That is,
\[
(\theta_{k+1}A'^{(k)})(x)=a_k\theta_{k+1}(x)\sum_{y\in T_1^{(k)}}\zeta^{(k)}(x,y)\,(G_kQ_k^*\delta_y)(x)\,A'(y).
\]
Here $\abs{(\partial^\eta_x\zeta^{(k)})(b,y)}\le1$, $\zeta^{(k)}(x,y)=1$ if $\abs{x-y}\le\frac12r_k$, and $\zeta^{(k)}(x,y)=0$ if $\abs{x-y}\ge r_k-2M_b$ \cite[p.~566]{Ba2}. The vector field is treated as an $\R^d$-valued field at zero background \cite[p.~608]{Ba1}. $G_k$ is the scale-$k$ vector propagator on $\Omega=T_\eta$, and \cite[(2.25)]{Ba1} holds for it at every $x$ (Proposition~\ref{prop:UV-BaProp21}; no restriction on $x$, since $\Omega=T_\eta$ \cite[p.~611]{Ba1}). Its constants $c_0,\delta_0$ do not depend on $k$, $\eps$ or the side of the torus (Lemma~\ref{app:G1:torusuniform}). Moreover $a_k=a(1-L^{-2})/(1-L^{-2k})\le a$ \cite[(2.15), p.~609]{Ba1}.
\item[(N3)] Lemma~\ref{app:G1:S05-PH}(a): $\abs{A'(y)}\le c_{A'}p(t_k)$ for $y\in\Lambda_1^{(k)}$, with $c_{A'}$ independent of $k,\eps,\ell_T$.
\item[(N4)] Lemma~\ref{app:G1:Rnest}: for $y\in\Lambda_{i+s}^{(l)}$ we have $\dist(y,\Lambda_i^{(l)c})\ge s\lfloor r_l\rfloor+1$. Also $\Lambda_{i+1}^{(l)}\subset\Lambda_i^{(l)}$, $\Lambda_0^{(l+1)}\subset\Lambda_7^{(l)\prime}$, and $B^{l+1}(\Lambda')=B^l(\Lambda)$ for unions $\Lambda$ of big blocks.
\end{itemize}
In scale-$j$ units, let $\cAc(x):=L^{-m(d-2)/2}(\theta_{k+1}A'^{(k)})(L^{-m}x)$ be the rescaled field. Its transporters are $\exp(\eta e(t_j)q\cAc_b)$, $\eta=L^{-j}$, $e(t)=\abs et^{1/2}$ \cite[(3.14), p.~614]{Ba1}. Put $\Omega=B^j(\Lambda_2^{(j)})$ and $\Omega_0=B^k(\Lambda_0^{(k)})$, and let $U$ be the union of the blocks $B^j(x)$, $x\in\Lambda_5^{(j)}$, intersected with $\Omega\cap\Omega_0$. For $c_1\ge0$, let $S_1(c_1)$ be the set of scale-$j$ unit blocks at scale-$j$ distance $\le c_1$ from $\supp\theta_{k+1}$.

\begin{lemma}[inputs of the expansion of the normalisations]\label{app:K:L13-6}\label{app:G1:Zinputs}
Assume (IH$_k$) and $R\ge R^{**}$, so that \eqref{eq:eppp} holds and $r(t)\ge R$. Then:
\begin{itemize}
\item[(i)] $\sup\abs{\cAc}\le c_Ap(t_j)$, with
\[
c_A=a\,c_0\,c_{A'}\,S,\qquad S=\sum_{n\ge0}(2n+3)^d\big(1+d(n+1)\big)\mathrm{e}^{-\delta_0(n-1)} .
\]
\item[(ii)] In scale-$j$ units, every transporter in $Q_j(A)$, in $Q(A)$ and in $W^{(j)}$ of \cite[(2.99)]{Ba2} runs along one bond, along a contour $\Gamma^{(j)}_{y,x}\subset B^j(y)$ of \cite[(2.2)]{Ba1}, or along a contour $\Gamma_{y',y}\subset B(y')$ of \cite[(2.1)]{Ba1}, with $y'$ on the scale-$(j+1)$ lattice. Its length $\eta\,\#\Gamma$ is at most $\ell:=\max\{1,d(L-1)\}$.
\item[(iii)] $U=\Omega_0$, and $\dist(S_1(c_1),U^c)\ge L^mr(t_k)-c_1-1$. For $c_1=L$, every vertex built from $\cAc$ in the proof of Proposition~\ref{app:G1:Z-W} vanishes outside the blocks of $S_1(L)$.
\item[(iv)] $\abs{S_1(c_1)}\le c_S\abs{X_j}$ with $c_S=(2M_b+2c_1+7)^d$.
\end{itemize}
The constants $c_A,\ell,c_1+1,c_S$ do not depend on $j,k,\eps,\ell_T$.
\end{lemma}

\begin{proof}
(i) \emph{Support.} Let $\theta_{k+1}(x)\ne0$ and $\zeta^{(k)}(x,y)\ne0$. By (N1) there are $x'\in B^k(z)$, $z\in\Lambda_2^{(k)}$, with $\abs{x-x'}<M_b$, and $\abs{x'-z}\le1-\eta$. By (N2), $\abs{x-y}<r_k-2M_b$. Hence $\abs{y-z}<r_k-M_b+1\le r_k$. Since $y,z$ are unit-lattice points, $\abs{y-z}\le\lfloor r_k\rfloor$. By (N4) with $s=1$, $y\in\Lambda_1^{(k)}$, so $\abs{A'(y)}\le c_{A'}p(t_k)$ by (N3).

\emph{The cut-off.} Join $y$ to $x$ by a coordinate path of at most $d\abs{x-y}/\eta$ bonds of $T_\eta$. Each bond changes $\zeta^{(k)}(\cdot,y)$ by at most $\eta$, and $\zeta^{(k)}(y,y)=1$. Hence $\abs{\zeta^{(k)}(x,y)}\le1+d\abs{x-y}$.

\emph{Kernel.} Apply \cite[(2.25)]{Ba1} (with $N=d$) to the $\R^d$-valued $f=Q_k^*(\delta_yA'(y))$, which has $\supp f=B^k(y)$ and $\norm f_\infty=\abs{A'(y)}$. Since $\dist(x,B^k(y))\ge\abs{x-y}-1$, we get $\abs{(G_kQ_k^*\delta_y)(x)A'(y)}\le c_0\mathrm{e}^{-\delta_0(\abs{x-y}-1)}\abs{A'(y)}$. There are at most $(2n+3)^d$ points $y$ with $n\le\abs{x-y}<n+1$. With $0\le\theta_{k+1}\le1$ and $a_k\le a$, this gives $\abs{(\theta_{k+1}A'^{(k)})(x)}\le a c_0Sc_{A'}p(t_k)$.

\emph{Rescaling.} The factor in the definition of $\cAc$ is $L^{-m(d-2)/2}\le1$. Moreover $t_j=L^{-m}t_k<t_k\le1$, and $p$ is decreasing, so $p(t_k)\le p(t_j)$.

(ii) By \cite[(2.99), p.~577]{Ba2}, $W^{(j)}$ is a combination of $Q_j(\cdot)$, $Q(\cdot)$, $D_B$ and $G_j$, and of the operators $F_2$, $F_{2,j}$, $V_j$. By \cite[(3.14)--(3.16), pp.~614--615]{Ba1} and \cite[(3.44)--(3.45), p.~620]{Ba1}, these last operators contain transporters only along single bonds (the factors $F_{1,j}$) and along the contours of $Q_j$ and $Q$.

\emph{Bonds.} A bond has length $\eta=L^{-j}\le1$.

\emph{The contours of $Q_j$.} By \cite[(2.2)]{Ba1},
\[
\Gamma^{(j)}_{y,x}=\Gamma_{y,x_{j-1}}\cup\dots\cup\Gamma_{x_1,x},
\]
with $x_0=x$ and $x_i\in B(x_{i+1})$ on the lattice of spacing $L^{i-j}$. By \cite[(2.1)]{Ba1}, $\Gamma_{x_{i+1},x_i}$ consists of $d$ coordinate segments. All their points have coordinates between those of $x_{i+1}$ and $x_i$. So the contour lies in the block of $x_{i+1}$, which lies in $B^j(y)$. By \cite[(1.17)]{Ba1}, $\abs{(x_i-x_{i+1})_\mu}\le(L-1)L^{i-j}$. Hence the length is at most
\[
d(L-1)\sum_{i=0}^{j-1}L^{i-j}=d(1-L^{-j})<d .
\]

\emph{The contours of $Q$.} For $Q$ (unit lattice to scale $j+1$), $\Gamma_{y',y}\subset B(y')$, and its $d$ segments each have length $\le L-1$.

(iii) \emph{The set $U$.} For $j\le l<k$, (N4) gives
\[
B^{l+1}(\Lambda_0^{(l+1)})\subset B^{l+1}(\Lambda_7^{(l)\prime})=B^l(\Lambda_7^{(l)})\subset B^l(\Lambda_0^{(l)}).
\]
Iterating from $l=k-1$ down to $l=j$ gives $\Omega_0\subset B^j(\Lambda_7^{(j)})\subset B^j(\Lambda_5^{(j)})\subset\Omega$. Hence $U=\Omega_0$.

\emph{Distance.} Let $x\in\supp\theta_{k+1}$ and $w\notin\Omega_0$. Then $w\in B^k(y')$ with $y'\notin\Lambda_0^{(k)}$. With $x'$ and $z$ as in (i), (N4) with $s=2$ gives $\abs{z-y'}\ge2\lfloor r_k\rfloor+1$. So, in scale-$k$ units,
\[
\abs{x-w}\ge\abs{z-y'}-\abs{z-x'}-\abs{w-y'}-\abs{x-x'}>2r_k-3-M_b\ge r_k ,
\]
since $r_k\ge R\ge32M_b+25$ by \eqref{eq:eppp}. In scale-$j$ units, $\dist(\supp\theta_{k+1},U^c)\ge L^mr_k$. A block of $S_1(c_1)$ has sup-diameter $<1$, so each of its points lies within $c_1+1$ of $\supp\theta_{k+1}$, and the bound follows.

\emph{Vertices.} The vertices of Proposition~\ref{app:G1:Z-W} are of three kinds. A vertex $D_s$ lives on bonds $b$ with $b_-\in\supp\cAc$. A nonzero entry $(y,x)$ of $Q_{j,s}$ needs a bond of $\Gamma^{(j)}_{y,x}\subset B^j(y)$ in $\supp\cAc$. A nonzero entry of $Q_s$ needs one in $\Gamma_{y',y}\subset B(y')$, whose sup-diameter is $<L$. Finally, $\supp\cAc\subset\supp\theta_{k+1}$.

(iv) Every block of $S_1(c_1)$ lies within scale-$k$ distance $s=M_b+c_1+1$ of $B^k(\Lambda_2^{(k)})$, hence within $s+1$ of some $z\in\Lambda_2^{(k)}$. It is therefore contained in a scale-$k$ block $B^k(y'')$ with $\abs{y''-z}\le s+2$. There are at most $(2s+5)^d$ such $y''$ for each $z$. Each scale-$k$ block contains $L^{dm}$ scale-$j$ blocks, and $\abs{X_j}=L^{dm}\abs{\Lambda_2^{(k)}}$.
\end{proof}

\begin{remark}\label{app:G1:Zinputs-rem}
In Proposition~\ref{app:G1:Zexp}, (Z.A) is part (i). The hypothesis (Z.C) is part (ii), with $\ell=d(L-1)=3$ at $L=2$, $d=3$. The bound of (Z.S) on $\rho$ is part (iii) with $S_1=S_1(c_1)$, $c_1=2L\sqrt d$, and $c=2L\sqrt d+1$. This set contains $S_1(L)$, hence every vertex, and it contains every block within Euclidean distance $2L\sqrt d$ of $\supp\cAc$, hence the set of blocks counted in the proof of Proposition~\ref{app:G1:Z-W}. The bound $\abs{S_1}\le c_S\abs{X_j}$ of (Z.S) is part (iv), with $c_S=(2M_b+4L\sqrt d+7)^d\ge1$. An inclusion $S_1\subset X_j$ fails in general, since $\theta_{k+1}$ reaches up to $M_b$ beyond $B^k(\Lambda_2^{(k)})$; it is not used. The factor $c_S$ appears explicitly in (a) and (b), enters (c) through $\abs{S_1}$, and enters the constants $C$ and $C_N$ of (d). It depends only on $d,L,M_b$. The quantity $w_j$ does not contain $c_S$, so the hypothesis $4w_j\le\gamma_0$ is a smallness condition on $\eps_0$.
\end{remark}

\section{The counterterm gaps}\label{app:G1:part-D8}

This appendix proves the bounds behind Theorem~\ref{thm:ch-W}(a)--(c): the two-loop gap bound (Proposition~\ref{app:G1:G2}), the derivative bound (Proposition~\ref{app:G1:R}) and the coefficient condition (Corollary~\ref{app:G1:CorC}), with their numerical evaluations (Corollaries~\ref{app:G1:Wa}, \ref{app:G1:Wb}, \ref{app:G1:Wc}). They rest on the exact two-loop form of King's counterterm and on slope bounds for each coupling order (Lemmas~\ref{app:G1:poly} and \ref{app:G1:dens}), which are proved in Subsection~\ref{app:G1:part-Tpoly}.

\emph{Units and notation.} By the scaling step of the proof of Theorem~\ref{thm:ch-W}, we work in units $m=\mu_*=1$. ``Every $\eps\le1$'' then means every $\eps\le1/m$, and a torus floor $\ell_0$ means $\ell_0/m$. The numerical constants below are computed in these units, for $\eps\le1$, and so are valid for every $\eps\le1/m$.
$T$ is a periodic torus of side $\ell_T\ge\ell_0$ with $\ell_T/\eps\in2\N$, and $T^*$ is its dual lattice. For $k\in T^*$ put $\hat k^2:=\sum_\mu(2/\eps)^2\sin^2(\eps k_\mu/2)$, $\hat G_M(k):=(\hat k^2+M^2)^{-1}$ and $\hat D(k):=(\hat k^2+\mu_*^2)^{-1}$, and write $\int_T f:=\abs T^{-1}\sum_{k\in T^*}f(k)$. Let $G_M(x):=\abs T^{-1}\sum_{k\in T^*}\mathrm{e}^{ik\cdot x}\hat G_M(k)$ be the lattice Green's function on $T$. Define
\[
C_T(M):=\int_T\hat G_M,\quad D_T:=\int_T\hat D,\quad J_T(M):=M^2\int_T\hat D\hat G_M,\quad S_T(M):=\sum_{x\in T}\eps^3G_M(x)^3.
\]
Let $\widetilde\Pi_T(p;M)$ be the subtracted order-$e^2$ self-energy of \cite[(3.32)--(3.33), p.~659]{K1} on $T$ (photon propagator $(-\Delta_\eps+\mu_*^2)^{-1}$), and put $F_T(M):=\mathrm{e}^{-2}\int_T\hat G_M^2\,\widetilde\Pi_T(\cdot;M)$. Write $Q_T(M):=\mathrm{e}^{-4}\delta m^2_{e^4,T}(\eps;M)$ for the order-$e^4$ coefficient of the counterterm of Definition~\ref{def:counterterm}.
Finally, the functions
\begin{gather*}
\tau(x):=1+\tfrac{12}{\pi x}+\tfrac{24}{\pi x^2}+\tfrac{64}{\pi^2x^3},\\
\theta(x):=\tfrac{3\pi}{16x}+\tfrac{3\pi}{8x^2}+\tfrac1{x^3},\\
A_R(x):=768\big[\tfrac{c_N}{8\pi}+\theta(x)\big]^2,\\
c_N:=\tfrac{\pi^3}{8}
\end{gather*}
are all decreasing in $x>0$.

\paragraph{Two lemmas.}\label{app:G1:ssec-recinputs}
\begin{lemma}[exact two-loop form]\label{app:G1:poly}
For every $K$ and every torus $T$,
\[
a_{K,T}(M)=M^2-16\lambda C_T(M)-2e^2D_T-e^2J_T(M)+128\lambda^2S_T(M)+16e^2\lambda F_T(M)+e^4Q_T(M),
\]
with no further terms. $Q_T=\tfrac34\eps^2D_T^2+(A)+(C)+(D1)+(D2)+(D3)$ is the Ward form of Corollary~\ref{app:G1:PD-e4}. Its first term does not depend on $M$. Moreover $M\mapsto a_{K,T}(M)$ is $C^\infty$ on $(0,\infty)$.
\end{lemma}
\noindent The polynomial structure, i.e.\ only the monomials $e^2,\lambda,e^4,e^2\lambda,\lambda^2$, is that of Definition~\ref{def:counterterm}: King's (3.32) sums over $2\le\alpha+2\beta\le4$, and the odd orders in $e$ vanish.

\emph{Smoothness (elementary).} At fixed $\eps$ and $T$ each term is a finite sum, over $T^*$ or over the sites of $T$, of products of factors $\hat G_M(k)$ and of $M$-independent factors. Since $\hat k^2+M^2\ge M^2>0$, each such term is a rational function of $M^2$ without poles on $(0,\infty)$.
Hence, for $0<M_1<M_2$, $a_{K,T}(M_2)-a_{K,T}(M_1)=\int_{M_1}^{M_2}\partial_Ma_{K,T}\,dM$, and the same identity holds for each order separately.

\begin{lemma}[slope bounds for each coupling order]\label{app:G1:dens}
Let $\eps\le1$, let $T$ have side $\ell_T\ge\ell_0$, and let $M>0$. Let $\bar\tau(M)\ge\int_T\hat G_M^2$, $\bar\psi(M)\ge\int_T\hat G_M\hat D$, $\bar\delta,\bar\Theta,\bar J_s,\bar J'_s,\bar\Lambda_i$ be the one-loop majorants of Lemma~\ref{app:G1:PD-oneloop}, all evaluated at $\ell_0$; here $\bar\tau=\mathcal I_{2,0}$ is the Riemann form, not an image form, which would need $\eps M\le1$. Let $\bar\omega\ge\sup\eps D_T$ be the constant $w+\bar E(\mu_*,\ell_0)$, $w=0.252731$, or its sharper form with $\min\{\bar E,\tilde E\}$ (Lemma~\ref{app:G1:PD-img}(c)). All of these majorants are non-negative, independent of the couplings and of $\eps$, and non-increasing in $M$ and in $\ell_0$. Then:
\begin{itemize}
\item[(o-$\lambda$)] $\partial_MC_T=-2M\int_T\hat G_M^2\le0$. Moreover
\[
C_T(M_1)-C_T(M_2)\ge h_1(M_1,M_2):=(2\pi^2)^{-1}[M_2\arctan(\pi/M_2)-M_1\arctan(\pi/M_1)]\ge0.
\]
\item[(o-$e^2$)] $\partial_MJ_T=2M\int_T\hat D\hat k^2\hat G_M^2$, and $0\le\partial_MJ_T\le j(M):=2M\min\{\bar\tau(M),\bar\psi(M)\}$. Also $\partial_MJ_T\le\frac{\pi^2}{32}\tau(M\ell_0)$.
\item[(o-$\lambda^2$)] $0\le-128\,\partial_MS_T\le768M\bar\tau(M)^2\le A_R(M\ell_0)/M$.
\item[(o-$e^2\lambda$)] $\abs{\partial_MF_T}\le\varphi_3(M):=4M\big[2\sqrt2\,\bar\Theta\bar J_{3/2}+\abs{\mu_*^2-2M^2}\bar\psi\bar J_0+3\sqrt3\,\bar\omega\bar J_1\big]+\bar\beta\bar\tau$, where $\bar\beta:=2M[2\bar\tau+8\bar\psi]$.
\item[(o-$e^4$)] $\partial_MQ_T\ge-b_4(M)$, where $b_4:=\varphi_A+28M\bar\delta\bar\tau+\varphi_{D2}$, $\varphi_A:=6M\big[2\sqrt2\,\bar\Theta\bar J'_{3/2}+\abs{\mu_*^2-2M^2}\bar\psi\bar J'_0+3\sqrt3\,\bar\omega\bar J'_1\big]+\bar\beta\bar\psi$, and $\varphi_{D2}:=2M\bar\psi(2\bar\psi+4\bar\delta)+\min\{8M\bar\psi(6\bar\psi+2\bar\tau+4\bar\delta),\,2^{5/2}M(2\bar\Lambda_1+4\bar\Lambda_2)\bar\Lambda_0\}$.
\end{itemize}
\end{lemma}
\noindent The elementary parts are the following.
\begin{itemize}
\item From $\partial_M\hat G_M=-2M\hat G_M^2$ we get $\partial_MC_T\le0$. We also get $\partial_MJ_T=2M\int_T\hat D\hat G_M(1-M^2\hat G_M)=2M\int_T\hat D\hat k^2\hat G_M^2$, which is $\ge0$.
\item Since $\hat D\hat k^2\le1$, $\partial_MJ_T\le2M\bar\tau$; since $\hat k^2\hat G_M\le1$, $\partial_MJ_T\le2M\bar\psi$.
\item For the order $\lambda^2$: $\partial_MG_M=-2M\,G_M*G_M$, so $-\partial_MS_T=6M\sum_x\eps^3G_M(x)^2(G_M*G_M)(x)$. Then $(G_M*G_M)(x)\le(G_M*G_M)(0)=\int_T\hat G_M^2$ (the Fourier coefficients $\hat G_M^2$ are non-negative), and $\sum_x\eps^3G_M^2=\int_T\hat G_M^2$ (Parseval). This gives $-\partial_MS_T\le6M(\int_T\hat G_M^2)^2$.
\item The non-negativity in (o-$\lambda^2$) uses $G_M\ge0$ (Lemma~\ref{app:G1:PD-pos}(a)).
\item The last bounds in (o-$e^2$) and (o-$\lambda^2$) use $\int_T\hat G_M^2\le(\pi^2/(64M))\tau(M\ell_T)$ and $M\bar\tau\le c_N/(8\pi)+\theta(M\ell_T)$ (Lemma~\ref{app:G1:PD-riem}), together with the monotonicity of $\tau$, $\theta$ and $\ell_T\ge\ell_0$.
\end{itemize}

\paragraph{The two-loop gap and statement (a).}
\begin{proposition}[two-loop gap between anchor masses on tori]\label{app:G1:G2}
Let $0<M_1<M_2$, $\mu_*\in(0,1]$ and $\ell_0>0$. For every $K\ge0$ (that is, $\eps\le1$) and every torus of side $\ell_T\ge\ell_0$,
\[
a_{K,T}(M_2)-a_{K,T}(M_1)\ \ge\ P_{\rm gap}:=(M_2^2-M_1^2)+16\lambda h_1(M_1,M_2)-e^2\bar j-\lambda^2\bar s-e^2\lambda\bar B_3-e^4\bar B_4^-,
\]
where $\bar j:=\int_{M_1}^{M_2}j$, $\bar s:=\int_{M_1}^{M_2}768M\bar\tau^2$, $\bar B_3:=16\int_{M_1}^{M_2}\varphi_3$ and $\bar B_4^-:=\int_{M_1}^{M_2}b_4$. All four are non-negative and independent of the couplings.
\end{proposition}
\noindent With Corollary~\ref{app:G1:Wa} this gives Theorem~\ref{thm:ch-W}(a).
\begin{proof}
Fix $\eps$ and $T$, and take the difference of Lemma~\ref{app:G1:poly} at $M_2$ and $M_1$. The terms $-2e^2D_T$ and $\frac34\eps^2e^4D_T^2$ do not depend on $M$ and cancel. Write each remaining order as $\int_{M_1}^{M_2}\partial_M$ (Lemma~\ref{app:G1:poly}, smoothness):
\begin{itemize}
\item order $\lambda$: $16\lambda[C_T(M_1)-C_T(M_2)]\ge16\lambda h_1$, by (o-$\lambda$);
\item order $e^2$: $-e^2[J_T(M_2)-J_T(M_1)]\ge-e^2\bar j$, by (o-$e^2$);
\item order $\lambda^2$: $128\lambda^2[S_T(M_2)-S_T(M_1)]\ge-\lambda^2\bar s$, by (o-$\lambda^2$);
\item order $e^2\lambda$: $16e^2\lambda[F_T(M_2)-F_T(M_1)]\ge-16e^2\lambda\int\abs{\partial_MF_T}\ge-e^2\lambda\bar B_3$, by (o-$e^2\lambda$);
\item order $e^4$: $e^4[Q_T(M_2)-Q_T(M_1)]\ge-e^4\bar B_4^-$, by (o-$e^4$).
\end{itemize}
Each majorant is evaluated at $\ell_0$ and does not depend on $\eps\le1$ or on $\ell_T\ge\ell_0$, so the bound holds for the infimum over $K$ and $T$.
\end{proof}

\begin{corollary}[statement (a) of Theorem~\ref{thm:ch-W}]\label{app:G1:Wa}
Take $\mu_*=1$, $(M_1,M_2)=(0.7,1)$ and $\ell_0=96$. On the closed box $0\le\lambda\le0.1$, $0\le e^2\le0.1$, for every $\eps\le1$ and every torus of side $\ge96$, we have $a_{K,T}(1)-a_{K,T}(0.7)\ge0.3190043>0.319$.
\end{corollary}

\begin{proof}
\emph{Cell upper sums.} Fix $\mu_*$ and $\ell_0$, and let $\phi$ be one of the functions $j$, $768M\bar\tau^2$, $16\varphi_3$, $b_4$ of Lemma~\ref{app:G1:dens}. Each of them is built from the majorants $\bar\tau,\bar\psi,\bar\delta,\bar\Theta,\bar J_s,\bar J'_s,\bar\Lambda_i$ and the constant $\bar\omega$ by sums, products and minima with non-negative numerical coefficients, multiplied by the prefactors $M$ and $\abs{\mu_*^2-2M^2}$. The majorants are non-negative and non-increasing in $M$ (Lemma~\ref{app:G1:dens}), and $\bar\omega$ does not depend on $M$. The prefactor $M$ is increasing. The prefactor $\abs{\mu_*^2-2M^2}$ is monotone on each side of its zero $M=\mu_*/\sqrt2$, so its maximum on an interval is attained at an endpoint. For an interval $[a,b]\subset(0,\infty)$ let $\phi^{\rm cell}[a,b]$ be the number obtained from the formula for $\phi$ by evaluating every majorant at $a$ and every prefactor at its maximum on $[a,b]$. Sums, products and minima of non-negative numbers are non-decreasing in each argument, so $\phi(M)\le\phi^{\rm cell}[a,b]$ for every $M\in[a,b]$. Hence, for every partition $a=M^{(0)}<M^{(1)}<\dots<M^{(n)}=b$,
\begin{equation}\label{app:G1:eq-cell}
\int_a^b\phi\,dM\le\sum_{i=0}^{n-1}\big(M^{(i+1)}-M^{(i)}\big)\,\phi^{\rm cell}[M^{(i)},M^{(i+1)}],\qquad \sup_{[a,b]}\phi\le\max_{0\le i<n}\phi^{\rm cell}[M^{(i)},M^{(i+1)}].
\end{equation}
Take $[a,b]=[0.7,1]$, the uniform partition of $[0.7,1]$ into $400$ cells, and all majorants evaluated at $\ell_0=96$. Then the upper sums~\eqref{app:G1:eq-cell} give
\[
\bar j\le0.0886293,\quad \bar s\le7.179064,\quad \bar B_3\le9.330427,\quad \bar B_4^-\le1.703784.
\]
These four upper sums are evaluated in floating-point arithmetic with $40$ significant digits. In this evaluation the constant $w$ in $\bar\omega$ is replaced by the upper bound $0.25273101$, every term that is subtracted in $P_{\rm gap}$ is multiplied by $1+10^{-25}$, and each result is rounded upward in the last digit shown.

Next, $h_1\ge0$, because $x\mapsto x\arctan(\pi/x)$ is increasing: its derivative is $\arctan y-y/(1+y^2)>0$ with $y=\pi/x>0$. Drop the term $16\lambda h_1$. Each remaining negative term is non-increasing in $\lambda\ge0$ and in $e^2\ge0$. So on the box $P_{\rm gap}$ is at least its value at the corner $(\lambda,e^2)=(0.1,0.1)$:
\[
P_{\rm gap}\ \ge\ 0.51-0.1\cdot0.0886293-0.01\,(7.179064+9.330427+1.703784)=0.31900432.
\]
The majorants are non-increasing in $\ell_0$, so the bound at $\ell_0=96$ holds on every torus of side $\ge96$. With the scaling in the proof of Theorem~\ref{thm:ch-W}, this gives Theorem~\ref{thm:ch-W}(a): $a_{K,T}(m)-a_{K,T}(0.7m)\ge0.319\,m^2$ for $\eps\le1/m$ and $\ell_T\ge96/m$.
\end{proof}
\begin{remark}
The proof of Theorem~\ref{thm:ch-W}(a) uses the rounded values $\bar j\le0.088629$ and $\bar B_3\le9.3304$; they are slightly below the computed upper sums $0.08862928\ldots$ and $9.330426\ldots$. With the outward-rounded values above the corner value is $0.3190043$ instead of $0.3190041$, so the conclusion is unchanged.
\end{remark}

\paragraph{The derivative bound and statement (b).}
\begin{proposition}[derivative bound and gap between anchor masses above the target mass]\label{app:G1:R}
Let $\lambda\ge0$, $e\in\R$, $\mu_*>0$, $\ell_0>0$ and $0<m_1<m_2$. Put
\[
p(M):=e^2\tfrac{\pi^2}{32}\tau(M\ell_0)+\lambda^2\frac{A_R(M\ell_0)}{M}+16e^2\lambda\,\varphi_3(M)+e^4b_4(M).
\]
Then for every $K\ge0$ ($\eps\le1$) and every torus of side $\ell_T\ge\ell_0$:
\begin{itemize}
\item[(i)] $\partial_Ma_{K,T}(M)\ge2M-p(M)$ for every $M>0$;
\item[(ii)] $a_{K,T}(m_2)-a_{K,T}(m_1)\ge(m_2-m_1)(m_2+m_1-S)$ for every $S\ge\sup_{[m_1,m_2]}p$;
\item[(iii)] $\sup_{M\ge m_1}p(M)\le e^2s_1+\lambda^2s_2+e^2\lambda s_3+e^4s_4$, where $s_1=\frac{\pi^2}{32}\tau(m_1\ell_0)$, $s_2=A_R(m_1\ell_0)/m_1$, $s_3=16\sup_{M\ge m_1}\varphi_3$ and $s_4=\sup_{M\ge m_1}b_4$.
\end{itemize}
\end{proposition}
\noindent With Corollary~\ref{app:G1:Wb} this gives Theorem~\ref{thm:ch-W}(b).
\begin{proof}
(i): Differentiate Lemma~\ref{app:G1:poly}. The $M$-free terms drop out. The order-$\lambda$ term contributes $-16\lambda\,\partial_MC_T\ge0$ by (o-$\lambda$), and we drop it. The order-$e^2$ term contributes $-e^2\partial_MJ_T\ge-e^2\frac{\pi^2}{32}\tau(M\ell_0)$. The order-$\lambda^2$ term contributes $128\lambda^2\partial_MS_T\ge-\lambda^2A_R(M\ell_0)/M$. The order-$e^2\lambda$ term is at least $-16e^2\lambda\varphi_3$, and the order-$e^4$ term is at least $-e^4b_4$, by Lemma~\ref{app:G1:dens}.
(ii): By smoothness (Lemma~\ref{app:G1:poly}), $a(m_2)-a(m_1)=\int_{m_1}^{m_2}\partial_Ma\,dM\ge\int_{m_1}^{m_2}(2M-S)\,dM=(m_2^2-m_1^2)-(m_2-m_1)S$. The bound does not depend on $\eps\le1$ or on $\ell_T\ge\ell_0$.
(iii): The supremum of a sum is at most the sum of the suprema. Since $\tau$ and $x\mapsto A_R(x\ell_0)/x$ are decreasing, the first two suprema are attained at $M=m_1$.
\end{proof}
\begin{corollary}[statement (b) of Theorem~\ref{thm:ch-W}]\label{app:G1:Wb}
Take $\mu_*=1$ and $\ell_0=64$. On the box $0\le\lambda,e^2\le0.1$, for $m_2>m_1>2$, every $\eps\le1$ and every torus of side $\ge64$ (in particular $\ge96$), we have $a_{K,T}(m_2)-a_{K,T}(m_1)\ge(m_2-m_1)(m_2+m_1-0.3622)\ge(m_2-m_1)(m_2+m_1-0.3631)>0$.
\end{corollary}
\noindent Here $s_1$ and $s_2$ are closed forms. The numbers $s_3$ and $s_4$ are bounded by upper sums over the cells of a partition of a bounded interval of masses, together with a bound on the tail of large masses.
\begin{proof}
At $m_1=2$, $\ell_0=64$: $s_1=\frac{\pi^2}{32}\tau(128)=0.317774\ldots\le0.3178$ and $s_2=A_R(128)/2=9.694097\ldots\le9.6941$.

For $s_3=16\sup_{M\ge2}\varphi_3$ and $s_4=\sup_{M\ge2}b_4$ choose $M_c\ge2$ and a partition of $[2,M_c]$. On $[2,M_c]$ the supremum is bounded by the maximum of the cell values, by the second inequality in~\eqref{app:G1:eq-cell}. On $[M_c,\infty)$ it is bounded by the monomial bound of Lemma~\ref{app:K:L15-1} at $M_c$, because $M_c\ge2>\mu_*$. The resulting upper bounds are evaluated numerically, in floating-point arithmetic, from the explicit majorants of Lemma~\ref{app:G1:dens} at $\ell_0=64$ and from the monomials of Lemma~\ref{app:K:L15-1}. Rounded to nearest at four decimals, the values are $17.3775$ and $5.9678$. Hence $s_3\le17.37755$ and $s_4\le5.96785$.

Each term of $p$ has a non-negative coefficient times a non-negative monomial in $(e^2,\lambda)$, so $p$ is non-decreasing in $\lambda,e^2\ge0$. Hence, on the box,
\[
\sup_{M\ge2}p\le0.1\cdot0.3178+0.01\,(9.6941+17.37755+5.96785)=0.362175<0.3622.
\]
For $m_1>2$ we have $[m_1,m_2]\subset[2,\infty)$, so $S:=0.3622$ is admissible in (ii). Finally, $m_2+m_1-0.3631>4-0.3631>0$.
\end{proof}

\paragraph{The coefficient condition and statement (c).}
\begin{corollary}[the coefficient condition]\label{app:G1:CorC}
Let $m>0$, $c:=2H\ge0$, and let $S(m)\ge\sup_{M\ge m}p(M)$. For every $\eps\le1$ and every torus of side $\ge\ell_0$, and every $m_1\ge m$,
\[
a_{K,T}(m_1)-a_{K,T}(m)\ge(m_1-m)(m_1+m-S(m)).
\]
Hence $a_{K,T}(m_1)\ge a_{K,T}(m)+2H$ whenever $m_1\ge m_1^c:=\frac12\big[S(m)+\big((S(m)-2m)^2+4c\big)^{1/2}\big]$. Moreover $m_1^c\ge m$, and $m_1^c$ is non-decreasing in $S(m)$.
\end{corollary}
\noindent This is Proposition~\ref{app:G1:R} with the anchors $(m,m_1)$, written as an inequality between King's c-numbers; with Corollary~\ref{app:G1:Wc} it gives Theorem~\ref{thm:ch-W}(c). In Section~\ref{sec:chain} the bare coefficient is $b_T=a_{K,T}(m)$ (Definition~\ref{def:counterterm}), which may be used on every torus by Lemma~\ref{app:G1:P1}(b).
\begin{proof}
The inequality is (ii) of Proposition~\ref{app:G1:R} with $(m_1,m_2)$ replaced by $(m,m_1)$ and $S=S(m)$ (it is trivial if $m_1=m$). Put $q(x):=(x-m)(x+m-S)-c=x^2-Sx+(Sm-m^2-c)$. Its roots are $\frac12[S\pm(S^2-4Sm+4m^2+4c)^{1/2}]=\frac12[S\pm((S-2m)^2+4c)^{1/2}]$. Since the leading coefficient is positive, $q(x)\ge0$ for $x\ge m_1^c$. Also $((S-2m)^2+4c)^{1/2}\ge\abs{S-2m}\ge2m-S$, so $m_1^c\ge m$. Finally,
\[
\partial_Sm_1^c=\tfrac12\big[1+(S-2m)\big((S-2m)^2+4c\big)^{-1/2}\big]\ge0.
\]
\end{proof}
\begin{corollary}[statement (c) of Theorem~\ref{thm:ch-W}]\label{app:G1:Wc}
Take $\mu_*=m=1$, $H=1$ and $\ell_0=64$. On the box $0\le\lambda,e^2\le0.1$ we have $m_1^c\le1.882<2$. Hence $a_{K,T}(m_1)\ge a_{K,T}(1)+2H$ for every $m_1>2$, every $\eps\le1$ and every torus of side $\ge64$.
\end{corollary}

\begin{proof}
Bound $\sup_{M\ge1}p(M)$ at the corner $(0.1,0.1)$ as follows. The terms $e^2\frac{\pi^2}{32}\tau(M\ell_0)$ and $\lambda^2A_R(M\ell_0)/M$ are decreasing in $M$, so their suprema are their values at $M=1$. For $16e^2\lambda\varphi_3+e^4b_4$, choose $M_c\ge1$ and use the cell values of~\eqref{app:G1:eq-cell} on a partition of $[1,M_c]$, and the monomial bound of Lemma~\ref{app:K:L15-1} on $[M_c,\infty)$ (here $M_c\ge1=\mu_*$). This upper bound is evaluated numerically, in floating-point arithmetic, from the explicit majorants of Lemma~\ref{app:G1:dens} at $\ell_0=64$. Rounded to nearest at four decimals it is $0.6119$, so $S(1)\le0.61195\le0.612$. Since $S$ is non-decreasing in $\lambda,e^2\ge0$, $S(1)\le0.612$ on the box. By the monotonicity in Corollary~\ref{app:G1:CorC},
\[
m_1^c\le\tfrac12\big[0.612+(1.388^2+8)^{1/2}\big]=1.88132\ldots\le1.882.
\]
\end{proof}

\subsection{\texorpdfstring{Proofs of Lemmas~\ref{app:G1:poly} and \ref{app:G1:dens}}{Proofs of the two-loop form and of the slope bounds}}\label{app:G1:part-Tpoly}
Notation is that of the beginning of this appendix, together with the following. $\mathsf X\in\{T,f\}$, where $f$ stands for $\eps\Z^3$; $\mathrm{BZ}_\eps:=[-\pi/\eps,\pi/\eps)^3$ and $\int_f:=\int_{\mathrm{BZ}_\eps}d^3k/(2\pi)^3$. $G_{M,\mathsf X}(x)$ and $D_{\mathsf X}(x)$ are the position kernels of $\hat G_M$ and $\hat D$, so $G_{M,T}=G_M$, $C_{\mathsf X}(M)=G_{M,\mathsf X}(0)$ and $D_{\mathsf X}:=D_{\mathsf X}(0)=\int_{\mathsf X}\hat D$. $\mu:=\mu_*$. Products of kernels are convolutions: $(K*K')(x):=\sum_y\eps^3K(x-y)K'(y)$. Finally $E_{\mathsf X}(M)$ (for $\mathsf X=T$) is $E_T(M):=C_T(M)-C_f(M)$, and $w:=\eps\int_f\hat k^{-2}=0.252731$, which does not depend on $\eps$.

\paragraph{Lattice tools.}
\begin{lemma}[density of states]\label{app:G1:PD-dos}
For every $\eps>0$ and every non-increasing $F:[0,\infty)\to[0,\infty)$,
\[
\int_{\mathrm{BZ}_\eps}F(\hat k^2)\,d^3k\le c_N\int_{\R^3}F(\abs k^2)\,d^3k,\qquad c_N=\tfrac{\pi^3}{8}.
\]
\end{lemma}
\begin{proof}
On $[-\pi,\pi]$ we have $\abs{\sin(y/2)}\ge\abs y/\pi$, so $\hat k^2\ge(4/\pi^2)\abs k^2$ on $\mathrm{BZ}_\eps$. Hence $\{\hat k^2\le E\}\subset\{\abs k\le\frac\pi2E^{1/2}\}$, and this set has volume at most $\frac{4\pi}3(\frac{\pi^2E}4)^{3/2}=c_N\cdot\frac{4\pi}3E^{3/2}$.

For the integral, use the layer-cake formula: $\int_{\mathrm{BZ}}F(\hat k^2)=\int_0^\infty\abs{\{F(\hat k^2)>t\}}\,dt$. Since $F$ is non-increasing, $\{F(\hat k^2)>t\}$ is $\{\hat k^2<E_t\}$ or $\{\hat k^2\le E_t\}$. The level sets of $\hat k^2$ are null, so its measure is at most $c_N\abs{\{k\in\R^3:F(\abs k^2)>t\}}$. If $E_t=\infty$ the right side is infinite. Integrate in $t$.
\end{proof}

\begin{lemma}[torus Riemann sums]\label{app:G1:PD-riem}
Let $F$ be as in Lemma~\ref{app:G1:PD-dos}, and let $T$ have side $\ell=N\eps$ with $N$ even. Then
\[
\int_TF(\hat k^2)\le c_N\!\int_{\R^3}\!F(\abs k^2)\tfrac{d^3k}{(2\pi)^3}+\tfrac{3\pi}{16\ell}\int_0^\infty\!F(u)\,du+\tfrac3{4\ell^2}\int_\R\!F(q^2)\,dq+\tfrac{F(0)}{\ell^3}.
\]
In particular $\int_T\hat G_M^2\le M^{-1}[c_N/(8\pi)+\theta(M\ell)]=\frac{\pi^2}{64M}\tau(M\ell)$.
\end{lemma}
\begin{proof}
\emph{One coordinate.} Let $g$ be even and non-increasing in $\abs k$ on $[-\pi/\eps,\pi/\eps]$, and put $h=2\pi/\ell$. For $1\le j\le N/2$ we have $h\,g(jh)\le\int_{(j-1)h}^{jh}g$, and symmetrically for $j<0$. Summing, $\ell^{-1}\sum_jg(jh)\le\int g\,dk/(2\pi)+g(0)/\ell$; the edge point $j=N/2$ is included.

\emph{Three coordinates.} $\hat k^2$ is increasing in each $\abs{k_\mu}$ on the Brillouin zone. So $F(\hat k^2)$, and every partial sum or integral of it, stays even and non-increasing in each remaining coordinate. Iterating over the three coordinates gives $\sum_{S\subset\{1,2,3\}}\ell^{-\abs S}\int F(\hat k^2)|_{k_S=0}$.

\emph{Bounding the pieces.} The term $S=\emptyset$ is bounded by Lemma~\ref{app:G1:PD-dos}. In the lower-dimensional integrals use $\hat k_\mu^2\ge(4/\pi^2)k_\mu^2$ and extend to $\R^{d}$. This gives $\int_{\R^2}F(\frac4{\pi^2}\abs k^2)\frac{d^2k}{(2\pi)^2}=\frac\pi{16}\int_0^\infty F$ and $\int_\R F(\frac4{\pi^2}k^2)\frac{dk}{2\pi}=\frac14\int F(q^2)dq$.

\emph{The instance.} For $F=(u+M^2)^{-2}$ the four terms are $\frac{c_N}{8\pi M}$, $\frac{3\pi}{16\ell M^2}$, $\frac{3\pi}{8\ell^2M^3}$ and $\frac1{\ell^3M^4}$. Their sum is $M^{-1}[c_N/(8\pi)+\theta(M\ell)]$. With $c_N=\pi^3/8$ this equals $\frac{\pi^2}{64M}\tau(M\ell)$, since $\frac{\pi^2}{64}\tau(x)=\frac{\pi^2}{64}+\theta(x)$.
\end{proof}

\begin{lemma}[positivity and images]\label{app:G1:PD-pos}
Let $M,\mu>0$ and $\mathsf X\in\{T,f\}$.
\begin{itemize}
\item[(a)] $G_{M,\mathsf X}(x)>0$ for every $x$, and $-\partial_{M^2}G_{M,\mathsf X}=G_{M,\mathsf X}*G_{M,\mathsf X}\ge0$ pointwise.
\item[(b)] $G_{M,T}(x)=\sum_{n\in\Z^3}G_{M,f}(x+n\ell)$, and likewise for $D$ and for $G*G$. In particular $E_T(M)=\sum_{n\ne0}G_{M,f}(n\ell)\ge0$ and $D_T-D_f=E_T(\mu)$.
\item[(c)] For $s,t>0$ the position kernels of $\hat G^s$ and $\hat D^t$ are pointwise positive. Hence, for every $q$ (with $q\in T^*$ on a torus),
\[
\int_{\mathsf X,P}\hat G(P)^s\hat D(P\pm q)^t\le\int_{\mathsf X}\hat G^s\hat D^t,\quad \int_{\mathsf X,P}\hat G(P)\hat G(P+q)\le\int_{\mathsf X}\hat G^2,\quad \int_{\mathsf X,P}\hat D(P-q)\hat D(P+q)\le\int_{\mathsf X}\hat D^2.
\]
Also $\abs{(K*K')(x)}\le(K*K')(0)$ whenever $\hat K,\hat K'\ge0$.
\end{itemize}
\end{lemma}
\begin{proof}
(a) Write $-\Delta_\eps+M^2=(6/\eps^2+M^2)\big(1-P/(6+\eps^2M^2)\big)$, where $P$ is the nearest-neighbour adjacency matrix: its entries are $\ge0$ and its norm is $6$. The Neumann series therefore has non-negative terms, its zeroth term is positive on the diagonal, and its $n$-th term reaches distance $n$. Moreover $-\partial_{M^2}(-\Delta+M^2)^{-1}=(-\Delta+M^2)^{-2}$.

(b) Both sides solve $(-\Delta+M^2)G=\eps^{-3}\delta$ on $T$, and that solution is unique. The image sum converges absolutely by Lemma~\ref{app:G1:PD-img}(b) below.

(c) Positivity follows from subordination, $(-\Delta_\eps+M^2)^{-s}=\Gamma(s)^{-1}\int_0^\infty\tau^{s-1}\mathrm{e}^{-\tau(-\Delta_\eps+M^2)}d\tau$: the lattice heat kernel is positive, and on $T$ one uses images. Then Parseval gives $\int_P\hat K(P)\hat K'(P-q)=\sum_x\eps^3K(x)K'(x)\mathrm{e}^{-iqx}\le\sum_x\eps^3K(x)K'(x)$. For the last claim, $K*K'$ has non-negative Fourier transform $\hat K\hat K'$.
\end{proof}

\begin{lemma}[splitting and two lattice identities]\label{app:G1:PD-split}
Write $\hat x:=(2/\eps)\sin(\eps x/2)$ componentwise.
\begin{itemize}
\item[(a)] For $q\ne0$ and $0<\theta\le\min(s,t)$: $\int_{\mathsf X,P}\hat G(P)^s\hat D(P-q)^t\le(2/\hat q)^{2\theta}\big[\int_{\mathsf X}\hat G^{s-\theta}\hat D^t+\int_{\mathsf X}\hat G^s\hat D^{t-\theta}\big]$.
\item[(b)] $\abs{\widehat{(a+b)}_\nu}\le\abs{\hat a_\nu}+\abs{\hat b_\nu}$, hence $\widehat{(a+b)}^2\le2\hat a^2+2\hat b^2$. Also $\widehat{(p+p')}^2\le2\widehat{(p-p')}^2+8\hat p'^2$.
\item[(c)] $\abs{\hat D(P-q)-\hat D(P+q)}\le4\hat P\hat q\,\hat D(P-q)\hat D(P+q)$.
\item[(d)] Componentwise, $\widehat{(a+b)}_\mu^2+\widehat{(a-b)}_\mu^2=2\hat a_\mu^2+2\hat b_\mu^2-\eps^2\hat a_\mu^2\hat b_\mu^2$.
\end{itemize}
\end{lemma}
\begin{proof}
(a) By $\abs{\sin(x+y)}\le\abs{\sin x}+\abs{\sin y}$ componentwise and Minkowski, $\hat q\le\hat P+\widehat{(q-P)}$. So either $\hat P\ge\hat q/2$, and then $\hat G(P)^\theta\le(2/\hat q)^{2\theta}$; or $\widehat{(P-q)}\ge\hat q/2$, and then $\hat D(P-q)^\theta\le(2/\hat q)^{2\theta}$. Bound the integrand by the sum of the two cases and apply Lemma~\ref{app:G1:PD-pos}(c). When $s=\theta$, the factor $\hat G^0$ is the $\delta$-kernel.

(b) The first claim is the same subadditivity. The second uses $\abs{\sin2x}\le2\abs{\sin x}$, so $\widehat{2v}\le2\hat v$, followed by $(a+b)^2\le2a^2+2b^2$.

(c) Use $\sin^2A-\sin^2B=\sin(A+B)\sin(A-B)$, then $\abs{(2/\eps)\sin\eps x}\le2\abs{\hat x}$ and Cauchy--Schwarz.

(d) This is $4-4\cos a\cos b=2(2-2\cos a)+2(2-2\cos b)-(2-2\cos a)(2-2\cos b)$, divided by $\eps^2$.
\end{proof}

\begin{lemma}[torus images; the constant $\bar\omega$]\label{app:G1:PD-img}
\begin{itemize}
\item[(a)] For every $\eps$ with $\ell/\eps\in\N$ and every $M>0$,
\[
E_T(M)\le\bar E(M,\ell):=\tfrac{7.37}{\ell}+\tfrac{3\pi}{4\ell}E_1\!\big(\tfrac{M^2\ell^2}8\big)+\tfrac{12\pi\,\mathrm{erfc}(M\ell/(2\sqrt2))}{M\ell^2}+\tfrac{64\mathrm{e}^{-M^2\ell^2/8}}{M^2\ell^3},
\]
and $\bar E$ is decreasing in $\ell$.
\item[(b)] Let $z\in\eps\Z^3\setminus\{0\}$, $r:=\abs z_\infty$, and $\kappa(x):=\ln(1+x/(2\sqrt2))/(x/(2\sqrt2))$. Then $G_{M,f}(z)\le\frac{c_N}{\pi(2r-\eps)}(1+\frac{\eps M}{2\sqrt2})^{1-2r/\eps}$. Hence, for $\eps\le1$,
\[
E_T(M)\le\tilde E(M,\ell):=\sum_{j\ge1}(24j^2+2)\frac{c_N(1+M/(2\sqrt2))}{\pi(2j\ell-1)}\mathrm{e}^{-\kappa(M)Mj\ell/\sqrt2},
\]
which is decreasing in $\ell$.
\item[(c)] Put $\hat E:=\min\{\bar E,\tilde E\}$. For $\eps\le1$ and $\ell_T\ge\ell_0$: $\eps D_T\le\bar\omega:=w+\hat E(\mu,\ell_0)$.
\end{itemize}
\end{lemma}
\noindent In (c), ``$\eps\le1$'' is in the units in which $\bar E$ and $\tilde E$ are evaluated; see the units paragraph of Appendix~\ref{app:G1:part-D8}.
\begin{proof}
(a) This is Lemma~\ref{app:G1:WP-img}, with $7.364\le7.37$. Its hypothesis $\ell/\eps\ge2$ holds on the tori of this paper, where $\ell_T/\eps\in2\N$.

(b) First, $G_M=\int_{M^2}^\infty(G_s*G_s)\,ds$ by Lemma~\ref{app:G1:PD-pos}(a). Let $z_1=r>0$. Shift $k_1\to k_1+i\eta$ in $(G_s*G_s)(z)=\int_f\mathrm{e}^{ikz}\hat G_s(k)^2$; the integrand is periodic, so the vertical sides cancel. Since $\rRe\,\frac4{\eps^2}\sin^2\frac{\eps(k_1+i\eta)}2\ge\hat k_1^2-\frac4{\eps^2}\sinh^2\frac{\eps\eta}2$, the choice $\frac4{\eps^2}\sinh^2\frac{\eps\eta_s}2=\frac s2$ gives $\abs{\hat G_s}\le(\hat k^2+s/2)^{-1}$. Lemma~\ref{app:G1:PD-dos} then gives $\abs{(G_s*G_s)(z)}\le \mathrm{e}^{-\eta_sr}c_N/(8\pi(s/2)^{1/2})$, with $\eta_s\ge(2/\eps)\ln(1+\eps\sqrt s/(2\sqrt2))$. Integrating over $s$ gives the first bound.

For $\tilde E$: use the images of Lemma~\ref{app:G1:PD-pos}(b), $\#\{n:\abs n_\infty=j\}=24j^2+2$, and $r=j\ell$. Each factor of the first bound is non-decreasing in $\eps$ ($\kappa$ is decreasing), so $\eps\le1$ gives the value at $\eps=1$.

(c) $D_T=D_f+E_T(\mu)$ by Lemma~\ref{app:G1:PD-pos}(b). Also $\eps D_f\le\eps\int_f\hat k^{-2}=w$. Finally $\eps E_T(\mu)\le\hat E(\mu,\ell_T)\le\hat E(\mu,\ell_0)$, using $\eps\le1$ and monotonicity in $\ell$.
\end{proof}

\begin{lemma}[one-loop majorants]\label{app:G1:PD-oneloop}
Let $\alpha,\beta\ge0$ with $g:=\alpha+\beta>3/2$, let $a,b>0$, let $\nu:=\min(a,b)$, and put $F(u):=(u+a^2)^{-\alpha}(u+b^2)^{-\beta}$. Define
\[
\mathcal I_{\alpha,\beta}(a,b;\ell):=c_NC_3+\tfrac{3\pi}{16\ell}C_2+\tfrac3{4\ell^2}C_1+\tfrac{a^{-2\alpha}b^{-2\beta}}{\ell^3},
\]
where each $C_d$ is the minimum of the following admissible closed forms.
\begin{itemize}
\item $C_3$:
\begin{itemize}
\item $\Gamma(g-\frac32)\nu^{3-2g}/((4\pi)^{3/2}\Gamma(g))$;
\item if $\beta<\frac32<2g$: $\frac1{4\pi^2}B(c,\alpha-c)a^{2c-2\alpha}$ with $c=\frac{3-2\beta}2$, and the same with $(a,\alpha)\leftrightarrow(b,\beta)$;
\item the exact values $\frac1{4\pi(a+b)}$ at $(1,1)$ and $\frac1{8\pi a(a+b)^2}$ at $(2,1)$.
\end{itemize}
\item $C_2$: $\nu^{2-2g}/(g-1)$; or, if $0<\beta<1$, $a^{2(1-g)}B(1-\beta,g-1)$, and symmetrically.
\item $C_1$: $\sqrt\pi\Gamma(g-\frac12)\nu^{1-2g}/\Gamma(g)$; or, if $0<\beta<\frac12$, $a^{1-2g}B(\frac12-\beta,g-\frac12)$, and symmetrically.
\item For every $d$, the H\"older form $C_d(g,a;0)^{\alpha/g}C_d(g,b;0)^{\beta/g}$.
\end{itemize}
Then:
\begin{itemize}
\item[(a)] For every $\eps$ and every torus of side $\ell$: $\int_TF(\hat k^2)\le\mathcal I_{\alpha,\beta}(a,b;\ell)$. On $\eps\Z^3$ the same holds with $\ell=\infty$.
\item[(b)] $\mathcal I_{\alpha,\beta}$ is decreasing in $\ell$.
\item[(c)] At $a=M$, $b=\mu$, every admissible form is non-increasing in $M$, and hence so is $\mathcal I_{\alpha,\beta}(M,\mu;\ell)$.
\end{itemize}
A factor $\hat p^s$ with $0\le s<2\alpha$ is absorbed by $\hat p^s\le(\hat p^2+a^2)^{s/2}$, i.e.\ $\alpha\to\alpha-s/2$. The majorants used in Lemma~\ref{app:G1:dens} are the following, all at $(M,\mu;\ell_0)$:
\begin{center}
\begin{tabular}{lll}
symbol & bounds & majorant\\\hline
$\bar\tau(M)$ & $\tau_{\mathsf X}:=\int_{\mathsf X}\hat G^2$ & $\mathcal I_{2,0}$\\
$\bar\psi(M)$ & $\Xi_{\mathsf X}(0):=\int_{\mathsf X}\hat D\hat G$ & $\mathcal I_{1,1}$\\
$\bar\delta$ & $\delta_{\mathsf X}:=\int_{\mathsf X}\hat D^2$ & $\mathcal I_{0,2}$\\
$\bar\Theta(M)$ & $\Theta_{\mathsf X}:=\int_{\mathsf X}\hat D^{3/4}\hat G+\int_{\mathsf X}\hat D\hat G^{3/4}$ & $\mathcal I_{1,3/4}+\mathcal I_{3/4,1}$\\
$\bar J_s(M)$ & $J_s:=\int_{\mathsf X}\hat G^3\hat p^s$, $s\in\{0,1,\frac32\}$ & $\mathcal I_{3-s/2,0}$\\
$\bar J'_s(M)$ & $J'_s:=\int_{\mathsf X}\hat D\hat G^2\hat p^s$ & $\mathcal I_{2-s/2,1}$\\
$\bar\Lambda_0$ & $\int\hat k^{1/2}\hat G\hat D$ & $\mathcal I_{3/4,1}$\\
$\bar\Lambda_1$ & $\int\hat G^{5/4}\hat D+\int\hat G^2\hat D^{1/4}$ & $\mathcal I_{5/4,1}+\mathcal I_{2,1/4}$\\
$\bar\Lambda_2$ & $\int\hat G^{1/4}\hat D^2+\int\hat G\hat D^{5/4}$ & $\mathcal I_{1/4,2}+\mathcal I_{1,5/4}$
\end{tabular}
\end{center}
In particular $\bar\tau(M)=\mathcal I_{2,0}(M;\ell_0)$ is the Riemann form $M^{-1}[c_N/(8\pi)+\theta(M\ell_0)]$ of Lemma~\ref{app:G1:PD-riem}. It is \emph{not} an image form, which would need $\eps M\le1$.
\end{lemma}
\begin{proof}
(a) Apply Lemma~\ref{app:G1:PD-riem}. Each closed form is the integral of a pointwise larger function: $F\le(u+\nu^2)^{-g}$, or $(k^2+b^2)^{-\beta}\le k^{-2\beta}$, or H\"older with exponents $g/\alpha$ and $g/\beta$. Each is evaluated by $\int_0^\infty t^{c-1}(t+a^2)^{-x}dt=a^{2(c-x)}B(c,x-c)$.

(b) The $\ell$-terms are positive multiples of $\ell^{-1},\ell^{-2},\ell^{-3}$.

(c) Check each form in turn:
\begin{itemize}
\item the $\nu$-forms: $\nu=\min(M,\mu)$ is non-decreasing and the exponent $d-2g$ is negative;
\item the Beta forms: $M^{2c-2\alpha}$ with $c<\alpha$, or independent of $M$;
\item the exact forms are decreasing;
\item the H\"older forms are $M^{(d-2g)\alpha/g}$;
\item the zero-mode term is $M^{-2\alpha}\mu^{-2\beta}$.
\end{itemize}
A minimum of non-increasing functions is non-increasing.
\end{proof}

\paragraph{The low orders and the order $e^2\lambda$.}
Write $\psi:=(\varphi_1+i\varphi_2)/\sqrt2$, $U_\mu(x)=\mathrm{e}^{i\eps eA_\mu(x)}$, $s_\mu(q):=(2/\eps)\sin(\eps q_\mu/2)$. King's hopping term $\frac12\sum_{x,\mu}\eps\abs{U_\mu(x)\varphi(x+\eps e_\mu)-\varphi(x)}^2$ expands into:
\begin{itemize}
\item the free part $\sum_x\eps^3\bar\psi(-\Delta_\eps+M^2)\psi$, so $\ev{\psi\bar\psi}=\hat G_M$;
\item the vertices $V_1=\sum_x\eps^3(2e/\eps)\sum_\mu A_\mu\,\mathrm{Im}(\bar\psi\psi')$ and $V_2=\sum_x\eps^3e^2\sum_\mu A_\mu^2\,\rRe(\bar\psi\psi')$;
\item $V_3=-(\eps e^3/3)\sum_x\eps^3\sum_\mu A_\mu^3\,\mathrm{Im}(\bar\psi\psi')$ and $V_4=-(\eps^2e^4/12)\sum_x\eps^3\sum_\mu A_\mu^4\,\rRe(\bar\psi\psi')$,
\end{itemize}
where $\psi'=\psi(x+\eps e_\mu)$. The quartic term is $\lambda\abs\varphi^4=4\lambda(\bar\psi\psi)^2$. A self-energy $\Sigma$ is a shift of the coefficient of $\abs\psi^2$, and $\delta m^2=-\Sigma(0)$ order by order. Every $V_n$ carries $e^n$ and $n$ photon legs, and all photon legs are contracted, so only even powers of $e$ occur.

\begin{lemma}[orders $\lambda$, $e^2$, $\lambda^2$, $e^2\lambda$]\label{app:G1:PD-low}
On $\mathsf X\in\{T,f\}$:
\begin{itemize}
\item[(i)] $\delta m^2_{\lambda,\mathsf X}=-16\lambda C_{\mathsf X}(M)$.
\item[(ii)] The order-$e^2$ self-energy is
\[
\Sigma_{e^2,\mathsf X}(p;M)=e^2D_{\mathsf X}\sum_\mu\cos(\eps p_\mu)-e^2\int_{\mathsf X,p'}\hat D(p-p')\widehat{(p+p')}^2\hat G_M(p'),
\]
and $\delta m^2_{e^2,\mathsf X}=-2e^2D_{\mathsf X}-e^2J_{\mathsf X}(M)$: the seagull gives $-3e^2D_{\mathsf X}$ and the exchange gives $e^2D_{\mathsf X}-e^2J_{\mathsf X}$.
\item[(iii)] $\delta m^2_{\lambda^2,\mathsf X}=+128\lambda^2S_{\mathsf X}(M)$.
\item[(iv)] $\delta m^2_{e^2\lambda,\mathsf X}=16\lambda\int_{\mathsf X}\hat G_M^2\tilde\Pi_{\mathsf X}(\cdot;M)=16e^2\lambda F_{\mathsf X}(M)$, where $\tilde\Pi_{\mathsf X}(p):=\Sigma_{e^2,\mathsf X}(p)-\Sigma_{e^2,\mathsf X}(0)$. Explicitly,
\[
\tilde\Pi_{\mathsf X}(p;M)=-e^2D_{\mathsf X}\sum_\mu(1-\cos\eps p_\mu)-e^2\int_{\mathsf X,p'}\big[\hat D(p-p')\widehat{(p+p')}^2-\hat D(p')\hat p'^2\big]\hat G_M(p').
\]
\end{itemize}
\end{lemma}
\noindent The proof of (iv) uses Lemma~\ref{app:G1:PD-list21} below, which lists all one-particle-irreducible graphs of order $e^2\lambda$.
\begin{proof}
(i) Contracting one pair in $4\lambda(\bar\psi\psi)^2$ gives $16\lambda C\bar\psi\psi$.

(ii) The seagull with its photon pair contracted gives $e^2D_{\mathsf X}\sum_\mu\cos(\eps p_\mu)$, which is $3e^2D_{\mathsf X}$ at $p=0$. The second-order term $-\frac12\ev{V_1V_1}$, with one photon and one scalar contraction, has vertex factors whose product is $e^2\widehat{(p+p')}_\mu^2$. At $p=0$, $\int\hat D\hat p'^2\hat G=D_{\mathsf X}-J_{\mathsf X}$, because $\hat p'^2\hat G=1-M^2\hat G$.

(iii) Only the sunset survives: the cactus is cancelled exactly by the $\delta m^2_\lambda$ insertion, since both contain the same $C_{\mathsf X}$. The Wick count is $2\cdot2\cdot2$ per orientation, times $2$ orientations, times $\frac12(4\lambda)^2$, i.e.\ $32(N+2)=128$ for $N=2$.

(iv) By definition, the counterterm $\delta m^2_{e^2\lambda,\mathsf X}$ is minus the zero-momentum value of the sum of the 1PI graphs of order $e^2\lambda$, the vertex $c_{2,1}$ itself excluded (Corollary~\ref{app:G1:PD-irr}(c) below). By Lemma~\ref{app:G1:PD-list21} these graphs are (a$_1$)--(a$_3$), (b$_1$), (b$_2$) and (c). The argument goes graph by graph, so it holds for $\mathsf X\in\{T,f\}$.
\begin{itemize}
\item The graphs (a$_1$), (a$_2$), (a$_3$) are the order-$\lambda$ tadpole $16\lambda\int_{\mathsf X}\hat G$ with one of three insertions on its loop line: the seagull, the exchange, or the counterterm $\delta m^2_{e^2,\mathsf X}$. Their sum is the tadpole with $\hat G$ replaced by $-\hat G(\Sigma_{e^2,\mathsf X}+\delta m^2_{e^2,\mathsf X})\hat G=-\hat G^2\tilde\Pi_{\mathsf X}$, by (ii) and $\delta m^2_{e^2,\mathsf X}=-\Sigma_{e^2,\mathsf X}(0)$. So their contribution to the self-energy is $-16\lambda\int_{\mathsf X}\hat G^2\tilde\Pi_{\mathsf X}$.
\item The graph (b$_1$) is the exchange graph with the insertion $16\lambda C_{\mathsf X}$ on its internal scalar line, and (b$_2$) is the same graph with the insertion $\delta m^2_{\lambda,\mathsf X}=-16\lambda C_{\mathsf X}$. Their sum is $0$.
\item Each graph (c) vanishes before the sum over the position of its $\lambda$-vertex. Let $w$ be that position, and let $V_1'$ sit on the bond $(z,z+\eps e_\mu)$; both of its scalar lines end at $w$. By~\eqref{app:G1:eq-Vn}, the sum over the two orientations of $V_1'$ is
\[
c_1\big[G_M(z+\eps e_\mu-w)\,G_M(z-w)-G_M(z-w)\,G_M(z+\eps e_\mu-w)\big]=0,
\]
because $G_M$ is even.
\end{itemize}
Hence $\delta m^2_{e^2\lambda,\mathsf X}=16\lambda\int_{\mathsf X}\hat G_M^2\tilde\Pi_{\mathsf X}$.
\end{proof}

\paragraph{Graphs, irreducibility and one-photon parts}
In this paragraph $T$ is a torus of side $\ell_T$ with $\ell_T/\eps\in2\N$, $\psi=\varphi/\sqrt2$, and
\[
(-\Delta_A\psi)(x)=\eps^{-2}\sum_\mu\big[2\psi(x)-\mathrm{e}^{i\eps eA_\mu(x)}\psi(x+\eps e_\mu)-\mathrm{e}^{-i\eps eA_\mu(x-\eps e_\mu)}\psi(x-\eps e_\mu)\big],
\]
so that the hopping term of the action of \cite{K1} is $\ip{\psi}{-\Delta_A\psi}$. Let $V_n$ be the part of $-\Delta_A-(-\Delta_\eps)$ of degree $n$ in $A$:
\begin{equation}\label{app:G1:eq-Vn}
(V_n\psi)(x)=c_n\sum_\mu\Big[A_\mu(x)^n\psi(x+\eps e_\mu)+(-1)^nA_\mu(x-\eps e_\mu)^n\psi(x-\eps e_\mu)\Big],\qquad c_n:=-\eps^{-2}\frac{(i\eps e)^n}{n!}.
\end{equation}
The quadratic forms $\ip{\psi}{V_n\psi}$, $n=1,\dots,4$, are the vertices $V_1,\dots,V_4$ written before Lemma~\ref{app:G1:PD-low}. We read~\eqref{app:G1:eq-Vn} as follows. A vertex of type $V_n$ sits on a bond $b=(x,x+\eps e_\mu)$ and carries the field $A_\mu(x)^n$ of that bond. Its factor is $c_n$ if its $\bar\psi$-leg is at the lower end $x$ and its $\psi$-leg at the upper end $x+\eps e_\mu$, and $(-1)^nc_n$ in the opposite case.

Let $\mathbb E_0$ be the Gaussian expectation under which $\psi$ is a complex field with $\mathbb E_0[\psi(x)\bar\psi(y)]=G_M(x-y)$ and $\mathbb E_0[\psi\psi]=0$, and $A$ is an independent real field with $\mathbb E_0[A_\mu(x)A_\nu(y)]=\delta_{\mu\nu}D_T(x-y)$. The interaction is
\[
\mathcal V:=\ip{\psi}{(-\Delta_A+\Delta_\eps)\psi}+4\lambda\sum_x\eps^3\abs{\psi(x)}^4+\delta m^2\sum_x\eps^3\abs{\psi(x)}^2,\qquad \delta m^2=\sum e^\alpha\lambda^\beta\,\delta m^2(\alpha,\beta).
\]
The full two-point function $G_{\rm full}(x-y):=\mathbb E_0[\psi(x)\bar\psi(y)\mathrm{e}^{-\mathcal V}]/\mathbb E_0[\mathrm{e}^{-\mathcal V}]$ is a formal power series in $(e,\lambda)$, and each coefficient is a finite sum of Gaussian moments.

By Wick's theorem each coefficient is a finite sum of values of \emph{graphs}. A graph has vertices of three kinds. A vertex of type $V_n$ has one $\psi$-leg, one $\bar\psi$-leg and $n$ photon legs. A vertex $\lambda$ sits at a site and has two $\psi$-legs and two $\bar\psi$-legs. A vertex $c_{\alpha,\beta}$ sits at a site, has one $\psi$-leg and one $\bar\psi$-leg, and has the factor $\delta m^2(\alpha,\beta)$. A graph has two external scalar legs. All its other legs are paired into lines. A scalar line joins a $\psi$-leg at $u$ with a $\bar\psi$-leg at $v$ and has the value $G_M(u-v)$. A photon line joins legs on the bonds $(x,\mu)$ and $(y,\nu)$ and has the value $\delta_{\mu\nu}D_T(x-y)$. The value of a graph is the sum, over the positions of its vertices, of the product of its line values and vertex factors. Only connected graphs contribute to $G_{\rm full}$.

A connected two-point graph is \emph{one-particle irreducible} (1PI) if it stays connected after the removal of any one line. It is \emph{scalar-irreducible} if it stays connected after the removal of any one scalar line. A \emph{one-photon part} of a graph is a connected set $\gamma$ of vertices that carries no external leg and is joined to the other vertices by exactly one line, this line being a photon line. By \cite[p.~659]{K1}, the sum of the 1PI two-point graphs is $-\delta m^2+\Gamma^\eps$, where $-\delta m^2$ is the single counterterm vertex. So the function $\Gamma^\eps$ in (3.33) there is, up to the sign convention, the sum of the 1PI graphs other than the single counterterm vertices.

\begin{lemma}[scalar parity and lattice inversion]\label{app:G1:PD-parity}
\begin{itemize}
\item[(a)] For every set $S$ of vertices of a graph, the number of scalar legs at vertices of $S$ that are not paired with a scalar leg at a vertex of $S$ is even.
\item[(b)] Let $\gamma$ be a connected configuration of vertices and lines in which all scalar legs are paired and exactly one photon leg is unpaired. This leg belongs to a vertex $v_0$ of type $V_n$. For a bond $(z,\mu)$ let $w_\mu(z)$ be the sum, over all positions of the vertices of $\gamma$ with $v_0$ on $(z,\mu)$, of the product of the line values and vertex factors of $\gamma$. Then $w_\mu(z)=0$.
\item[(c)] Every graph that contains a one-photon part has value $0$.
\end{itemize}
\end{lemma}
\begin{proof}
(a) Every vertex has two or four scalar legs, and every scalar line with both ends in $S$ uses two of them. So the number of unpaired legs is an even number minus an even number.

(b) The value $w_\mu(z)$ does not depend on $z$, because all line values are translation invariant. Let $\iota$ be the inversion of $T$: $\iota(x)=-x$ on sites, and $\iota$ maps the bond $(x,x+\eps e_\mu)$ to the bond $(-x-\eps e_\mu,-x)$. Then $\iota$ maps the lower end of a bond to the upper end of its image and vice versa, and it preserves the direction $\mu$. Apply $\iota$ to every position in a position assignment of $\gamma$, carrying each leg to the image of its site. This is a bijection from the assignments with $v_0$ on $(z,\mu)$ onto those with $v_0$ on $\iota(z,\mu)$. Under it:
\begin{itemize}
\item a scalar line value $G_M(u-v)$ becomes $G_M(v-u)=G_M(u-v)$;
\item a photon line value $\delta_{\mu\nu}D_T(x-y)$ becomes $\delta_{\mu\nu}D_T(y-x)$, which is the same;
\item the factors of the vertices $\lambda$ and $c_{\alpha,\beta}$ do not change;
\item by~\eqref{app:G1:eq-Vn}, the factor of each vertex of type $V_{n_v}$ is multiplied by $(-1)^{n_v}$, because its $\bar\psi$- and $\psi$-legs exchange the roles of lower and upper end.
\end{itemize}
Every photon line uses two photon legs and one photon leg is unpaired, so $\sum_vn_v$ is odd. Hence $w_\mu(z)=-w_\mu(\iota(z,\mu))=-w_\mu(z)$.

(c) Let $\gamma$ be a one-photon part, joined to the rest by the photon line from the leg of $v_0\in\gamma$ on $(z,\mu)$ to a leg on $(y,\mu)$. By definition, every other line at a vertex of $\gamma$ has both ends in $\gamma$. Fix $z$ and the positions outside $\gamma$, and sum first over the positions of the vertices of $\gamma$ other than $v_0$. This inner sum is $w_\mu(z)$ times factors that do not involve $\gamma$, and $w_\mu(z)=0$ by (b). Here $\gamma$, with the joining line removed, satisfies the hypotheses of (b): its scalar legs are all paired, by (a) applied to $\gamma$ and because $\gamma$ carries no external leg.
\end{proof}

\begin{corollary}[1PI and scalar-irreducible self-energies]\label{app:G1:PD-irr}
\begin{itemize}
\item[(a)] Let $\Gamma$ be a connected two-point graph. Every scalar line whose removal disconnects $\Gamma$ separates the two external legs.
\item[(b)] A scalar-irreducible two-point graph that is not 1PI contains a one-photon part, and its value is $0$.
\item[(c)] Let $\Sigma_s(p)$ be the sum of the amputated scalar-irreducible two-point graphs other than the single counterterm vertices, with the sign convention of Lemma~\ref{app:G1:PD-low}. Then, as formal power series,
\[
\hat G_{\rm full}(p)^{-1}=\hat p^2+M^2+\delta m^2+\Sigma_s(p),
\]
and at every order $\Sigma_s(0)$ is the zero-momentum value $\Sigma_{\rm 1PI}(0)$ of the sum of the amputated 1PI graphs other than the single counterterm vertices, i.e.\ of the function $\Gamma^\eps$ of \cite[(3.33)]{K1}, up to the sign convention. Hence the renormalisation condition of \cite[(3.33)]{K1}, $\delta m^2(\alpha,\beta)+\Sigma_{\rm 1PI}(0;\alpha,\beta)=0$ for $2\le\alpha+2\beta\le4$, where $\Sigma_{\rm 1PI}(0;\alpha,\beta)$ is the coefficient of $e^\alpha\lambda^\beta$, holds if and only if the coefficients of $e^\alpha\lambda^\beta$, $2\le\alpha+2\beta\le4$, of $\hat G_{\rm full}(0)-M^{-2}$ vanish.
\end{itemize}
\end{corollary}
\begin{proof}
(a) Removing such a line leaves two components $H_1,H_2$, and each contains one end of the line. By Lemma~\ref{app:G1:PD-parity}(a), the number of external legs in $H_i$ plus one is even. So each $H_i$ contains exactly one external leg.

(b) Let a photon line disconnect $\Gamma$ into $H_1,H_2$. By Lemma~\ref{app:G1:PD-parity}(a) the number of external (scalar) legs in each $H_i$ is even, so one of them, say $H_1$, contains none. Then $H_1$ is a one-photon part, and Lemma~\ref{app:G1:PD-parity}(c) applies.

(c) By (a), the scalar bridges of a connected two-point graph lie on every path between its external legs. So the amputated graph is a chain of scalar-irreducible two-point parts joined by single scalar lines. For a complex field the Wick weight of such a chain is the product of the weights of its parts, and summing the chains gives the geometric series $\hat G\sum_{n\ge0}\big(-(\delta m^2+\Sigma_s)\hat G\big)^n$. By (b), $\Sigma_s$ and the sum of the 1PI graphs have the same value at every order. The last statement follows by induction on the order, because $\hat G(0)=M^{-2}$.
\end{proof}

\begin{lemma}[the graphs of order $e^2\lambda$]\label{app:G1:PD-list21}
Apart from the single vertex $c_{2,1}$, the vertex sets of order $e^2\lambda$ whose vertex factors do not vanish are
\[
S_1=\{\lambda,V_2\},\ S_2=\{\lambda,V_1,V_1'\},\ S_3=\{\lambda,c_{2,0}\},\ S_4=\{V_2,c_{0,1}\},\ S_5=\{V_1,V_1',c_{0,1}\},\ S_6=\{c_{2,0},c_{0,1}\}.
\]
The 1PI two-point graphs on these sets are exactly the following six:
\begin{itemize}
\item[(a$_1$)] on $S_1$: both external legs at $\lambda$, two scalar lines $\lambda$--$V_2$, and the photon line from $V_2$ to itself;
\item[(a$_2$)] on $S_2$: both external legs at $\lambda$, scalar lines $\lambda$--$V_1$, $V_1$--$V_1'$, $V_1'$--$\lambda$, and the photon line $V_1$--$V_1'$;
\item[(a$_3$)] on $S_3$: both external legs at $\lambda$, and two scalar lines $\lambda$--$c_{2,0}$;
\item[(b$_1$)] on $S_2$: external legs at $V_1$ and $V_1'$, scalar lines $V_1$--$\lambda$, $\lambda$--$V_1'$ and a scalar line from $\lambda$ to itself, and the photon line $V_1$--$V_1'$;
\item[(b$_2$)] on $S_5$: external legs at $V_1$ and $V_1'$, scalar lines $V_1$--$c_{0,1}$, $c_{0,1}$--$V_1'$, and the photon line $V_1$--$V_1'$;
\item[(c)] on $S_2$: external legs at $\lambda$ and at $V_1$ (or at $V_1'$), one scalar line $\lambda$--$V_1$, two scalar lines $\lambda$--$V_1'$, and the photon line $V_1$--$V_1'$ (and the same with $V_1$ and $V_1'$ exchanged).
\end{itemize}
Every other graph on $S_1,\dots,S_6$ is disconnected or has a bridge.
\end{lemma}
\begin{proof}
The counterterms of odd order in $e$ vanish (Definition~\ref{def:counterterm}). So the power $\lambda^1$ comes from exactly one vertex $\lambda$ or $c_{0,1}$, unless the graph is the single vertex $c_{2,1}$. The power $e^2$ comes from $V_2$, from $V_1V_1$, or from $c_{2,0}$, because $V_n$ carries $e^n$. This gives $S_1,\dots,S_6$.

On each set the photon legs can be paired in only one way: the two legs of $V_2$ with each other, or the legs of $V_1,V_1'$ with each other. Once the external legs are placed, the scalar lines form a multigraph whose degrees are the numbers of free scalar legs. We list all such multigraphs, writing $uu$ for a scalar line from $u$ to itself.
\begin{itemize}
\item $S_1$. External legs both at $\lambda$, degrees $(2,2)$: either $\{\lambda V_2,\lambda V_2\}$, which is (a$_1$), or $\{\lambda\lambda,V_2V_2\}$, which is disconnected. Both at $V_2$, degrees $(4,0)$: $\lambda$ is isolated. One at each, degrees $(3,1)$: $\{\lambda\lambda,\lambda V_2\}$, and $\lambda V_2$ is a bridge.
\item $S_3$. As for $S_1$, with $c_{2,0}$ in place of $V_2$ and no photon line. This gives (a$_3$), a disconnected graph, or a graph with a bridge.
\item $S_4$ and $S_6$. There is one scalar line. If the external legs are at different vertices this line is a bridge; otherwise the other vertex is isolated.
\item $S_2$, with degrees listed in the order $(\lambda,V_1,V_1')$.
 \begin{itemize}
 \item External legs $(\lambda,\lambda)$, degrees $(2,2,2)$. The triangle is (a$_2$). In $\{\lambda V_1,\lambda V_1,V_1'V_1'\}$ and $\{\lambda V_1',\lambda V_1',V_1V_1\}$ the photon line is a bridge. The graphs $\{\lambda\lambda,V_1V_1',V_1V_1'\}$ and $\{\lambda\lambda,V_1V_1,V_1'V_1'\}$ are disconnected.
 \item External legs $(V_1,V_1')$, degrees $(4,1,1)$. Either $\{\lambda\lambda,\lambda V_1,\lambda V_1'\}$, which is (b$_1$), or $\{\lambda\lambda,\lambda\lambda,V_1V_1'\}$, which is disconnected.
 \item External legs $(V_1,V_1)$, degrees $(4,0,2)$. Either $\{\lambda\lambda,\lambda V_1',\lambda V_1'\}$, where the photon line is a bridge, or $\{\lambda\lambda,\lambda\lambda,V_1'V_1'\}$, which is disconnected. The case $(V_1',V_1')$ is the same.
 \item External legs $(\lambda,V_1)$, degrees $(3,1,2)$. The graph $\{\lambda V_1,\lambda V_1',\lambda V_1'\}$ is (c). In $\{\lambda\lambda,\lambda V_1,V_1'V_1'\}$ the line $\lambda V_1$ and the photon line are bridges. In $\{\lambda\lambda,\lambda V_1',V_1'V_1\}$ the line $\lambda V_1'$ is a bridge. There are no other cases: $\lambda$ has odd degree, so it has one or three lines to the other vertices, at most one of them to $V_1$ and at most two to $V_1'$. The case $(\lambda,V_1')$ is the same.
 \end{itemize}
\item $S_5$, with degrees listed in the order $(c_{0,1},V_1,V_1')$ and $c:=c_{0,1}$.
 \begin{itemize}
 \item External legs $(V_1,V_1')$: either $\{V_1c,cV_1'\}$, which is (b$_2$), or $\{V_1V_1',cc\}$, which is disconnected.
 \item External legs $(V_1,V_1)$: either $\{cV_1',cV_1'\}$, where the photon line is a bridge, or $\{cc,V_1'V_1'\}$, which is disconnected.
 \item External legs $(c,c)$: $c$ is isolated.
 \item External legs $(c,V_1)$: either $\{cV_1',V_1'V_1\}$ or $\{cV_1,V_1'V_1'\}$, and in both the line at $c$ is a bridge.
 \item The cases $(V_1',V_1')$ and $(c,V_1')$ are the same, with $V_1$ and $V_1'$ exchanged.
 \end{itemize}
\end{itemize}
\end{proof}

\begin{lemma}[structure of $\tilde\Pi$]\label{app:G1:PD-Pi}
Put $\Xi_{\mathsf X}(p):=\int_{\mathsf X,p'}\hat D(p-p')\hat G_M(p')\in[0,\Xi_{\mathsf X}(0)]$ and $Y_\mu(p):=\int_{\mathsf X,p'}\hat D(p-p')\hat p_\mu'^2\hat G_M(p')\in[0,D_{\mathsf X}]$.
\begin{itemize}
\item[(a)] $\tilde\Pi_{\mathsf X}(p;M)=-e^2\big\{2\hat p^2\Xi_{\mathsf X}(p)+(\mu^2-2M^2)[\Xi_{\mathsf X}(p)-\Xi_{\mathsf X}(0)]-\eps^2\sum_\mu\hat p_\mu^2Y_\mu(p)+D_{\mathsf X}\tfrac{\eps^2}2\hat p^2\big\}$.
\item[(b)] $\abs{\tilde\Pi_{\mathsf X}(p;M)}/e^2\le2\sqrt2\,\Theta_{\mathsf X}(M)\hat p^{3/2}+\abs{\mu^2-2M^2}\,\Xi_{\mathsf X}(0)+3\sqrt3\,(\eps D_{\mathsf X})\,\hat p$.
\item[(c)] $0\le\partial_M\Sigma_{e^2,\mathsf X}(p;M)\le e^2\beta_{\mathsf X}(M)$ with $\beta_{\mathsf X}:=2M[2\tau_{\mathsf X}+8\Xi_{\mathsf X}(0)]$. Hence $\abs{\partial_M\tilde\Pi_{\mathsf X}(p;M)}\le e^2\beta_{\mathsf X}(M)$ for all $p$.
\end{itemize}
\end{lemma}
\noindent In (b) the seagull and $Y$ terms could be sharpened from $\frac32\eps^2D\hat p^2$ to $\frac12\eps^2D\hat p^2$; this is not used.
\begin{proof}
(a) Apply Lemma~\ref{app:G1:PD-split}(d) to $\widehat{(p+p')}^2$ in Lemma~\ref{app:G1:PD-low}(iv), and use $\hat k^2\hat D=1-\mu^2\hat D$ and $\hat p'^2\hat G=1-M^2\hat G$. The linearly divergent constants $2D_{\mathsf X}-C_{\mathsf X}$ cancel in $\Sigma(p)-\Sigma(0)$. The identity is algebraic, so it holds on $T$ and on $\eps\Z^3$. The ranges of $\Xi_{\mathsf X}$ and $Y_\mu$ follow from Lemma~\ref{app:G1:PD-pos}(c) and $\hat p_\mu'^2\hat G\le1$.

(b) Lemma~\ref{app:G1:PD-split}(a) with $s=t=1$ and $\theta=\frac14$ gives $\Xi_{\mathsf X}(p)\le\sqrt2\hat p^{-1/2}\Theta_{\mathsf X}$, which bounds $2\hat p^2\Xi_{\mathsf X}(p)$. Next, $\abs{\Xi_{\mathsf X}(p)-\Xi_{\mathsf X}(0)}\le\Xi_{\mathsf X}(0)$. The seagull and $Y$ terms together are at most $\frac32\eps^2D_{\mathsf X}\hat p^2$, and $\eps\hat p\le2\sqrt3$ turns this into $3\sqrt3(\eps D_{\mathsf X})\hat p$.

(c) We have $\partial_M\Sigma_{\rm ex}(p)=2Me^2\int\hat D(p-p')\widehat{(p+p')}^2\hat G(p')^2\ge0$. Bound it with Lemma~\ref{app:G1:PD-split}(b), $\hat D\hat k^2\le1$, $\hat p'^2\hat G\le1$ and Lemma~\ref{app:G1:PD-pos}(c). The seagull and the counterterm do not depend on $M$, or do not depend on $p$, and so drop out of $\partial_M\tilde\Pi$. Then $\partial_M\Sigma_{\rm ex}(p)$ and $\partial_M\Sigma_{\rm ex}(0)$ both lie in $[0,e^2\beta_{\mathsf X}]$, so their difference is at most $e^2\beta_{\mathsf X}$ in absolute value.
\end{proof}

\paragraph{The order $e^4$: Ward form.}
\begin{proposition}[Ward form of King's counterterm]\label{app:G1:PD-ward}
Take $\lambda=0$, a torus $T$ and any $\eps$. Let $a=M^2+\delta m^2$ and $\mathcal G_{A,a}:=(-\Delta_A+a)^{-1}$. Let $\ev\cdot$ be the photon expectation of the model at $\lambda=0$: the Feynman Gaussian with mass $\mu$, times $\det(-\Delta_A+a)^{-1}$. Put $\hat G_{\rm full}:=(\ev{\mathcal G_{A,a}})^\wedge$, $W_A:=-\Delta_A\mathbf1$ and $R_A:=W_A+ie\,\mathrm{div}A$. Then, as formal power series in $e$,
\[
\delta m^2=e^2D_T-e^2M^2\int_T\hat D\,\hat G_{\rm full}-M^2\ev{(\mathcal G_{A,a}R_A)(x)}.
\]
King's $\delta m^2$, truncated at $\alpha+2\beta\le4$, satisfies this identity through order $e^4$, which is all that is used. Moreover $R_A=R^{(2)}+R^{(3)}+R^{(4)}+O(e^5)$ with
\begin{itemize}
\item $R^{(2)}=\frac{e^2}2\sum_\mu[A_\mu(x)^2+A_\mu(x-\eps e_\mu)^2]$,
\item $R^{(3)}=\frac{i\eps^2e^3}6\sum_\mu\nabla^-_\mu(A_\mu^3)(x)$,
\item $R^{(4)}=-\frac{\eps^2e^4}{24}\sum_\mu[A_\mu(x)^4+A_\mu(x-\eps e_\mu)^4]$.
\end{itemize}
\end{proposition}
\begin{proof}
\emph{Step 1.} King's condition \cite[(3.33), p.~659]{K1} reads $\delta m^2(\alpha,\beta)+\Sigma_{\rm 1PI}(0;\alpha,\beta)=0$ at each order. By Dyson at $p=0$, $\hat G_{\rm full}(0)^{-1}=M^2+\delta m^2+\Sigma_{\rm 1PI}(0)$. So King's condition is equivalent to $\hat G_{\rm full}(0)=1/M^2$, order by order.

\emph{Step 2.} $(-\Delta_A+a)\mathbf1=a+W_A$, so $a\,\mathcal G_{A,a}\mathbf1=1-\mathcal G_{A,a}W_A$. Averaging gives $\delta m^2=-M^2\ev{(\mathcal G_{A,a}W_A)(x)}$.

\emph{Step 3.} Gauge covariance: $\mathcal G_{A+\nabla^+\chi,a}(x,y)=\mathrm{e}^{-ie(\chi(x)-\chi(y))}\mathcal G_{A,a}(x,y)$, and the determinant is invariant.

\emph{Step 4.} Gaussian integration by parts with the Feynman covariance, together with Step 3 (the determinant contributes $0$), gives $\ev{\mathrm{div}A(y)F}=\ev{\frac d{dt}F(A+t\nabla^+\chi_y)|_0}$ with $\chi_y(z)=-D_T(y-z)$.

\emph{Step 5.} Applied to $F=\mathcal G_{A,a}(x,y)$: $\ev{\mathcal G(x,y)\mathrm{div}A(y)}=-ie(D_T(0)-D_T(x-y))G_{\rm full}(x-y)$. Multiplying by $-ie\eps^3$, summing over $y$, and using Parseval and Step 1 gives $-M^2\ev{\mathcal G\cdot(-ie\,\mathrm{div}A)}=e^2D_T-e^2M^2\int\hat D\hat G_{\rm full}$.

\emph{Step 6.} Add Steps 2 and 5. The expansion of $R_A$ is that of $1-\cos\vartheta$ and $\sin\vartheta/\eps-eA$, with $\vartheta=\eps eA$.
\end{proof}

\begin{corollary}[the $e^4$ coefficient]\label{app:G1:PD-e4}
On $\mathsf X\in\{T,\eps\Z^3\}$, $\delta m^2_{e^4,\mathsf X}(M)=e^4Q_{\mathsf X}(M)$ with
\[
Q_{\mathsf X}=\tfrac34\eps^2D_{\mathsf X}^2+(A)+(C)+(D1)+(D2)+(D3),
\]
where, with $s_\mu(k):=(2/\eps)\sin(\eps k_\mu/2)=\hat k_\mu$:
\begin{itemize}
\item $(A)=M^2\int_{\mathsf X}\hat D\hat G^2\tilde\Pi_{\mathsf X}/e^2$;
\item $(C)=-\frac{\eps^2D_{\mathsf X}}2(D_{\mathsf X}-J_{\mathsf X}(M))$;
\item $(D1)=\frac12\sum_\mu\sum_z\eps^3[D_{\mathsf X}(z)^2+D_{\mathsf X}(z+\eps e_\mu)^2][G_{M,\mathsf X}(z)+G_{M,\mathsf X}(z+\eps e_\mu)]$;
\item $(D2)=-\sum_\mu\int_k\int_Ps_\mu(k)[s_\mu(k)+s_\mu(2P+k)]\hat G(k)\hat G(P)\hat D(k)\hat D(P-k)$;
\item $(D3)=\int_{\mathsf X}\hat D(k)^2\sum_\mu\Pi_{\mu\mu}(k;M)/e^2$, where $\Pi_{\mu\mu}(k;M)/e^2=2G_{M,\mathsf X}(\eps e_\mu)-\int_{\mathsf X,q}\widehat{(2q+k)}_\mu^2\hat G(q)\hat G(q+k)$.
\end{itemize}
The first term does not depend on $M$.
\end{corollary}
\noindent The proof expands Proposition~\ref{app:G1:PD-ward} to order $e^4$ on $T$; the case $\eps\Z^3$ follows by a limit.
\begin{proof}
We work on $\mathsf X=T$, with the notation of Proposition~\ref{app:G1:PD-ward} (the resolvent $(-\Delta_A+a)^{-1}$, written $G_{A,a}$ below, $R_A$ and its parts $R^{(n)}$ of order $e^n$) and of the paragraph before Lemma~\ref{app:G1:PD-parity}. Fourier transforms are $\hat f(k)=\sum_x\eps^3\mathrm{e}^{-ik\cdot x}f(x)$ and $f(x)=\int_T\mathrm{e}^{ik\cdot x}\hat f(k)$. The photon field is expanded at the lower ends of the bonds, $A_\mu(x)=\int_T\mathrm{e}^{ik\cdot x}\hat A_\mu(k)$, so all momentum-space factors below are periodic. By~\eqref{app:G1:eq-Vn}, the kernel of $V_1$ between an outgoing $\bar\psi$-momentum $p$ and an incoming $\psi$-momentum $p'$, with photon momentum $p-p'$ entering, is
\begin{equation}\label{app:G1:eq-v1}
v_1(p,p')=-\tfrac{ie}{\eps}\big(\mathrm{e}^{i\eps p'_\mu}-\mathrm{e}^{-i\eps p_\mu}\big),\qquad v_1(p',p)=\overline{v_1(p,p')},\qquad \abs{v_1(p,p')}^2=e^2s_\mu(p+p')^2.
\end{equation}
Write $G:=(-\Delta_\eps+M^2)^{-1}$, so that $G\mathbf1=M^{-2}\mathbf1$ and $\hat G(0)=M^{-2}$, and write $(KR)(x):=\sum_y\eps^3K(x,y)R(y)$.

\emph{Step 0 (the order-$e^4$ terms).} Expand Proposition~\ref{app:G1:PD-ward} as a formal series in $e$, at $\lambda=0$, with $\delta m^2=e^2d_2+e^4d_4$. The Neumann series gives $G_{A,a}=G^{(0)}+G^{(1)}+G^{(2)}+O(e^3)$ with
\[
G^{(0)}=G,\qquad G^{(1)}=-GV_1G,\qquad G^{(2)}=-G(V_2+e^2d_2)G+GV_1GV_1G.
\]
The weight is $\det(-\Delta_A+a)^{-1}=\det(-\Delta_\eps+a)^{-1}\mathrm{e}^{-U}$ with $U=\operatorname{Tr}\log(1+(-\Delta_\eps+a)^{-1}(-\Delta_A+\Delta_\eps))=U_1+U_2+U_3+O(e^4)$, where $U_m$ has order $e^m$ and
\[
U_1=\operatorname{Tr}(GV_1)=0,\qquad U_2=\operatorname{Tr}(GV_2)-\tfrac12\operatorname{Tr}(GV_1GV_1).
\]
Here $U_1=0$ because, by~\eqref{app:G1:eq-Vn}, $\operatorname{Tr}(GV_1)=c_1\sum_{x,\mu}\eps^3A_\mu(x)\big[G_M(\eps e_\mu)-G_M(-\eps e_\mu)\big]$ and $G_M$ is even. Hence $\ev F=\mathbb E_0F-\operatorname{Cov}_0(F,U_2)-\operatorname{Cov}_0(F,U_3)+\dots$, and $\operatorname{Cov}_0(F,U_m)$ has order at least (order of $F$)$+m$. The function $R_A=R^{(2)}+R^{(3)}+R^{(4)}+O(e^5)$ starts at order $e^2$. So the order-$e^4$ part of $\ev{(G_{A,a}R_A)(x)}$ is
\begin{equation}\label{app:G1:eq-e4terms}
\mathbb E_0\big[(GR^{(4)})(x)\big]+\mathbb E_0\big[(G^{(1)}R^{(3)})(x)\big]+\mathbb E_0\big[(G^{(2)}R^{(2)})(x)\big]-\operatorname{Cov}_0\big((GR^{(2)})(x),U_2\big).
\end{equation}
Similarly $[\hat G_{\rm full}]_{e^2}$ is the Fourier transform of $\mathbb E_0G^{(2)}$, because $\operatorname{Cov}_0(G,U_2)=0$ for the $A$-independent $G$. By Lemma~\ref{app:G1:PD-low}(ii) and $e^2d_2=\delta m^2_{e^2,T}=-\Sigma_{e^2,T}(0)$, this transform is $-\hat G^2\tilde\Pi_T$. All expectations are invariant under the translations of $T$. Since $\sum_x\eps^3G(x-z)=M^{-2}$, every term of the form $-M^2\,\mathbb E[(GY)(x)]$ equals $-\mathbb E[Y(0)]$; we use this repeatedly.

\emph{Step 1 (the term $-e^2M^2\int_T\hat D\hat G_{\rm full}$).} Its order-$e^4$ part is $e^2M^2\int_T\hat D\hat G^2\tilde\Pi_T$, which is $e^4(A)$.

\emph{Step 2 (the constant).} $-M^2\mathbb E_0[(GR^{(4)})(x)]=-\mathbb E_0R^{(4)}(0)=\frac{\eps^2e^4}{24}\sum_\mu2\cdot3D_T^2=\frac34\eps^2e^4D_T^2$, because $\mathbb E_0A_\mu^4=3D_T^2$.

\emph{Step 3 (the term $(C)$).} For $L$ linear in $A$, $\mathbb E_0[L\,A_\mu(z)^3]=3D_T\,\mathbb E_0[L\,A_\mu(z)]$. So in $\mathbb E_0[G^{(1)}R^{(3)}]$ the factor $R^{(3)}=\frac{i\eps^2e^3}6\sum_\mu\nabla^-_\mu(A_\mu^3)$ may be replaced by $\frac{i\eps^2e^3D_T}2\,\mathrm{div}A=-\frac{\eps^2e^2D_T}2(-ie\,\mathrm{div}A)$. Steps~3 and~4 of the proof of Proposition~\ref{app:G1:PD-ward}, applied to the Gaussian $\mathbb E_0$ and to $G^{(1)}(x,y)$, which is linear in $A$, give
\[
\mathbb E_0\big[G^{(1)}(x,y)\,\mathrm{div}A(y)\big]=-ie\big(D_T(0)-D_T(x-y)\big)G_M(x-y).
\]
Hence $-M^2\mathbb E_0[(G^{(1)}(-ie\,\mathrm{div}A))(x)]=e^2M^2\big(D_T\hat G(0)-\int_T\hat D\hat G\big)=e^2(D_T-J_T)$, and the term is $-\frac{\eps^2e^4D_T}2(D_T-J_T)=e^4(C)$.

\emph{Step 4 (the term $\mathbb E_0[G^{(2)}R^{(2)}]$).} The two photon fields of $R^{(2)}$ are paired either with each other or with photon fields of the vertices in $G^{(2)}$.
\begin{itemize}
\item \emph{Pairing with each other.} $\mathbb E_0R^{(2)}=3e^2D_T$ is constant, so this part is $3e^2D_T\sum_y\eps^3\mathbb E_0G^{(2)}(x,y)=3e^2D_T[\hat G_{\rm full}]_{e^2}(0)$. It vanishes, because $\hat G_{\rm full}(0)=M^{-2}$ at order $e^2$ (Corollary~\ref{app:G1:PD-irr}(c)).
\item \emph{Both fields paired with $V_2$.} This gives $e^4(D1)$. Indeed, by~\eqref{app:G1:eq-Vn}, $(V_2f)(0)=\frac{e^2}2\sum_\mu[A_\mu(0)^2f(\eps e_\mu)+A_\mu(-\eps e_\mu)^2f(-\eps e_\mu)]$, and the connected moment is $\mathbb E_0[A_\mu(u)^2R^{(2)}(y)]_{\rm c}=e^2[D_T(u-y)^2+D_T(u-y+\eps e_\mu)^2]$. Therefore, after the shift $z=y-\eps e_\mu$ in the first term and with $G_M(-z)=G_M(z)$,
\begin{multline*}
-M^2\,\mathbb E_0\big[(-GV_2GR^{(2)})(x)\big]_{\rm c}=\mathbb E_0\big[(V_2GR^{(2)})(0)\big]_{\rm c}\\
=\frac{e^4}2\sum_\mu\sum_z\eps^3\big[D_T(z)^2+D_T(z+\eps e_\mu)^2\big]\big[G_M(z)+G_M(z+\eps e_\mu)\big].
\end{multline*}
\item \emph{One field paired with each $V_1$.} This gives $e^4(D2)$. In $\mathbb E_0[(V_1GV_1GR^{(2)})(0)]$, let $P$ be the momentum of the right propagator and $k$ that of the left one. The left vertex has factor $v_1(0,k)$ and the right vertex $v_1(k,P)$. The photon lines carry $k$ and $P-k$, with values $\hat D(k)$ and $\hat D(P-k)$. The field $R^{(2)}=\frac{e^2}2\sum_\mu[A_\mu(y)^2+A_\mu(y-\eps e_\mu)^2]$, with total photon momentum $P$ and both fields paired, contributes $e^2(1+\mathrm{e}^{-i\eps P_\mu})$. From $\mathrm{e}^{i\eps k_\mu}-1=i\eps\,\mathrm{e}^{i\eps k_\mu/2}s_\mu(k)$ and $\mathrm{e}^{i\eps P_\mu}+1-\mathrm{e}^{-i\eps k_\mu}-\mathrm{e}^{-i\eps(k+P)_\mu}=i\eps\,\mathrm{e}^{-i\eps k_\mu/2}[s_\mu(2P+k)+s_\mu(k)]$ we get
\[
(\mathrm{e}^{i\eps k_\mu}-1)\big(\mathrm{e}^{i\eps P_\mu}-\mathrm{e}^{-i\eps k_\mu}\big)\big(1+\mathrm{e}^{-i\eps P_\mu}\big)=(i\eps)^2\,s_\mu(k)\big[s_\mu(k)+s_\mu(2P+k)\big],
\]
so the product of the three factors is $e^4s_\mu(k)[s_\mu(k)+s_\mu(2P+k)]$. Hence
\[
-M^2\,\mathbb E_0\big[(GV_1GV_1GR^{(2)})(x)\big]_{\rm c}=-e^4\sum_\mu\int_{T,k}\int_{T,P}s_\mu(k)\big[s_\mu(k)+s_\mu(2P+k)\big]\hat G(k)\hat G(P)\hat D(k)\hat D(P-k).
\]
\end{itemize}

\emph{Step 5 (the measure term, $(D3)$).} $U_2$ is a quadratic form in $A$, $U_2=\frac12\ip{A}{\Pi A}$. By~\eqref{app:G1:eq-Vn}, $\operatorname{Tr}(GV_2)=\sum_{x,\mu}\eps^3e^2G_M(\eps e_\mu)A_\mu(x)^2$. By~\eqref{app:G1:eq-v1}, the diagonal symbol of $-\frac12\operatorname{Tr}(GV_1GV_1)$ is $-\frac{e^2}2\int_{T,q}s_\mu(2q+k)^2\hat G(q)\hat G(q+k)$ per component. So
\[
\Pi_{\mu\mu}(k)=2e^2G_M(\eps e_\mu)-e^2\int_{T,q}\widehat{(2q+k)}_\mu^{\,2}\hat G(q)\hat G(q+k).
\]
For a Gaussian, $\operatorname{Cov}_0(A_\mu(y)^2,\frac12\ip{A}{\Pi A})=(D\Pi D)_{\mu\mu}(y,y)=\int_T\hat D^2\Pi_{\mu\mu}$, and the same holds for $A_\mu(y-\eps e_\mu)^2$. Only the diagonal of $\Pi$ enters, because the covariance is $\delta_{\mu\nu}D_T$. Therefore
\[
M^2\operatorname{Cov}_0\big((GR^{(2)})(x),U_2\big)=\operatorname{Cov}_0\big(R^{(2)}(0),U_2\big)=e^2\sum_\mu\int_T\hat D(k)^2\Pi_{\mu\mu}(k),
\]
and the term $-M^2\big(-\operatorname{Cov}_0((GR^{(2)})(x),U_2)\big)$ is $e^4(D3)$.

\emph{Step 6 (conclusion).} The term $e^2D_T$ of Proposition~\ref{app:G1:PD-ward} has order $e^2$. By~\eqref{app:G1:eq-e4terms}, the order-$e^4$ part of $-M^2\ev{(G_{A,a}R_A)(x)}$ is the sum of Steps~2--5. Adding Step~1 gives $\delta m^2_{e^4,T}=e^4\big[\frac34\eps^2D_T^2+(A)+(C)+(D1)+(D2)+(D3)\big]$. The first term does not depend on $M$. The identity of Proposition~\ref{app:G1:PD-ward} holds through order $e^4$ for the counterterm of Definition~\ref{def:counterterm}: by Corollary~\ref{app:G1:PD-irr}(c), the renormalisation condition is equivalent to $\hat G_{\rm full}(0)=M^{-2}$ through that order, and this fixes $d_2,d_4$.

\emph{The lattice $\eps\Z^3$.} Lemma~\ref{app:G1:poly}, and through it Propositions~\ref{app:G1:G2} and~\ref{app:G1:R} and Corollary~\ref{app:G1:CorC}, use the present corollary only for $\mathsf X=T$. For $\mathsf X=\eps\Z^3$ the coefficient is defined graph by graph, with $\int_f$ in place of $\int_T$. Then both sides of the identity are finite sums of torus Riemann sums $\abs T^{-L}\sum_{(T^*)^L}$ of continuous periodic functions on $\mathrm{BZ}_\eps^L$, and all denominators are at least $\min(M^2,\mu_*^2)$. As $\ell_T\to\infty$ at fixed $\eps$, these Riemann sums converge to the corresponding integrals, and the identity passes to the limit.
\end{proof}

\paragraph{Proof of Lemma~\ref{app:G1:poly}.}
\begin{proof}
By \cite[(3.32), p.~659]{K1} the counterterm has the orders $2\le\alpha+2\beta\le4$, and the odd powers of $e$ vanish (above, or Definition~\ref{def:counterterm}). So $\delta m^2_{K,T}$ is exactly the sum of its five orders $\lambda$, $e^2$, $\lambda^2$, $e^2\lambda$, $e^4$, with no remainder.

Lemma~\ref{app:G1:PD-low} gives the first four orders. Corollary~\ref{app:G1:PD-e4} gives the order $e^4$: the coefficient of $e^4\lambda^0$ comes from graphs and insertions without a $\lambda$-vertex, so it is the $e^4$ coefficient of the $\lambda=0$ model. Hence
\[
a_{K,T}=M^2-16\lambda C_T-2e^2D_T-e^2J_T+128\lambda^2S_T+16e^2\lambda F_T+e^4Q_T.
\]
\emph{Smoothness.} At fixed $\eps$ and $T$ each term is a finite sum of products of factors $\hat G_M(k)$ and $M$-independent factors; $G_M(x)$ is itself such a finite sum. Since $\hat k^2+M^2\ge M^2$, each term is a rational function of $M^2$ with no pole on $(0,\infty)$.
\end{proof}

\paragraph{Proof of Lemma~\ref{app:G1:dens}.}
\begin{proof}
Throughout, the barred majorants are those of Lemma~\ref{app:G1:PD-oneloop} at $\ell_0$, and $\bar\omega$ is that of Lemma~\ref{app:G1:PD-img}(c). Every actual one-loop sum on $T$ ($\ell_T\ge\ell_0$) is at most its majorant, by Lemma~\ref{app:G1:PD-oneloop}(a),(b). The bound $\eps D_T\le\bar\omega$ is the only place where $\eps\le1$ enters, apart from $h_\eps\ge h_1$. No step needs $\eps M\le1$.

\emph{(o-$\lambda$).} $\partial_M\hat G_M=-2M\hat G_M^2$. By Lemma~\ref{app:G1:PD-pos}(b),(c), $C_T(M_1)-C_T(M_2)=(M_2^2-M_1^2)(G_{M_1}*G_{M_2})_T(0)$, which is at least its value on $\eps\Z^3$. On $\mathrm{BZ}_\eps$ we have $\hat k^2\le\abs k^2$, and $\mathrm{BZ}_\eps$ contains the ball of radius $R=\pi/\eps$. So
\[
C_f(M_1)-C_f(M_2)\ge\tfrac1{2\pi^2}\int_0^R\Big[\tfrac{M_2^2}{k^2+M_2^2}-\tfrac{M_1^2}{k^2+M_1^2}\Big]dk=:h_\eps(M_1,M_2).
\]
The integrand is $\ge0$, so $h_\eps$ is non-decreasing in $R$, and $h_\eps\ge h_1$ for $\eps\le1$.

\emph{(o-$e^2$).} $\partial_MJ_T=2M\int\hat D\hat G(1-M^2\hat G)=2M\int\hat D\hat k^2\hat G^2\ge0$. With $\hat D\hat k^2\le1$ this is at most $2M\tau_T\le2M\bar\tau$; with $\hat k^2\hat G\le1$ it is at most $2M\Xi_T(0)\le2M\bar\psi$. Lemma~\ref{app:G1:PD-riem} gives $2M\tau_T\le\frac{\pi^2}{32}\tau(M\ell_T)\le\frac{\pi^2}{32}\tau(M\ell_0)$.

\emph{(o-$\lambda^2$).} By Lemma~\ref{app:G1:PD-pos}(a), $-\partial_MS_T=3\sum_x\eps^3G_M^2(-\partial_MG_M)=6M\sum_x\eps^3G_M(x)^2(G_M*G_M)(x)\ge0$. Next, $(G*G)(x)\le(G*G)(0)=\tau_T$ by Lemma~\ref{app:G1:PD-pos}(c), and $\sum_x\eps^3G_M^2=\tau_T$ by Parseval. Hence $-128\partial_MS_T\le768M\tau_T^2\le768M\bar\tau^2=A_R(M\ell_0)/M$.

\emph{(o-$e^2\lambda$).} We have $\partial_MF_T=-4M\int\hat G^3\tilde\Pi/e^2+\int\hat G^2\partial_M\tilde\Pi/e^2$. Lemma~\ref{app:G1:PD-Pi}(b) bounds the first term by $4M[2\sqrt2\Theta_TJ_{3/2}+\abs{\mu^2-2M^2}\Xi_T(0)J_0+3\sqrt3(\eps D_T)J_1]$, and Lemma~\ref{app:G1:PD-Pi}(c) bounds the second by $\beta_T\tau_T$. Replace each factor by its majorant and $\eps D_T$ by $\bar\omega$; this gives $\varphi_3$.

\emph{(o-$e^4$).} By Corollary~\ref{app:G1:PD-e4}, $\partial_MQ_T$ is the sum of the following five derivatives.
\begin{itemize}
\item[(C)] $\partial_M(C)=\frac{\eps^2D_T}2\partial_MJ_T\ge0$ by (o-$e^2$). It is dropped.
\item[(A)] $(A)=M^2H$ with $H:=\int\hat D\hat G^2\tilde\Pi/e^2$, so $\partial_M(M^2H)=2MH-4M^3\int\hat D\hat G^3\tilde\Pi/e^2+M^2\int\hat D\hat G^2\partial_M\tilde\Pi/e^2$. Use $M^2\hat G\le1$ twice and Lemma~\ref{app:G1:PD-Pi}(b),(c): $\abs{\partial_M(A)}\le6M\int\hat D\hat G^2\abs{\tilde\Pi}/e^2+\beta_T\Xi_T(0)\le\varphi_A$, with $J'_s\le\bar J'_s$.
\item[(D1)] $\partial_MG(z)=-2M(G*G)(z)\in[-2M\tau_T,0]$ by Lemma~\ref{app:G1:PD-pos}(a),(c), and $\sum_z\eps^3D(z)^2=\sum_z\eps^3D(z+\eps e_\mu)^2=\delta_T$. Over three directions this gives $-12M\bar\delta\bar\tau\le\partial_M(D1)\le0$.
\item[(D3)] $\sum_\mu\partial_M[2G(\eps e_\mu)]\in[-12M\tau,0]$. The bubble part has derivative $2M\int\widehat{(2q+k)}_\mu^2(\hat G^2\hat G'+\hat G\hat G'^2)\ge0$, where $\hat G':=\hat G(q+k)$. Summed over $\mu$, by Lemma~\ref{app:G1:PD-split}(b) (with $a=q$, $b=q+k$), $\hat q^2\hat G\le1$ and Lemma~\ref{app:G1:PD-pos}(c), it is at most $2M[4\int\hat G\hat G'+4\tau]\le16M\tau$. Since $\int\hat D^2=\delta$, $\abs{\partial_M(D3)}\le16M\bar\delta\bar\tau$.
\item[(D2)] Write $(D2)=-\sum_\mu\int_ks_\mu(k)\hat G(k)\hat D(k)S_\mu(k)$, where $S_\mu(k):=\int_P[s_\mu(k)+s_\mu(2P+k)]\hat G(P)\hat D(P-k)$. Decompose $s_\mu(2P+k)=(2/\eps)[\sin(\eps P_\mu)\cos(\eps k_\mu/2)+\cos(\eps P_\mu)\sin(\eps k_\mu/2)]$.
\begin{itemize}
\item The first piece is odd in $P$ except through $\hat D(P-k)$. Symmetrising $P\to-P$ and applying Lemma~\ref{app:G1:PD-split}(c) bounds it by $\int2\abs{\hat P_\mu}\cdot2\hat P\hat k\,\hat G\hat D(P-k)\hat D(P+k)$.
\item The other two pieces are each at most $\abs{s_\mu(k)}\int\hat G\hat D(\cdot-k)$.
\end{itemize}
With $\sum_\mu\abs{\hat k_\mu}\abs{\hat P_\mu}\le\hat k\hat P$, $\hat P^2\hat G\le1$ and Lemma~\ref{app:G1:PD-pos}(c):
\[
\sum_\mu\abs{s_\mu}\abs{S_\mu}\le\hat k^2(2\Xi(0)+4\delta).
\]
Likewise, using also $\hat D\hat D'\le\frac12(\hat D^2+\hat D'^2)$ and Lemma~\ref{app:G1:PD-split}(a) with $\theta\in(\frac12,1]$:
\[
\sum_\mu\abs{s_\mu}\abs{\partial_MS_\mu}\le2M\hat k^2(2/\hat k)^{2\theta}\big[2(\textstyle\int\hat G^{2-\theta}\hat D+\int\hat G^2\hat D^{1-\theta})+4(\int\hat G^{1-\theta}\hat D^2+\int\hat G\hat D^{2-\theta})\big].
\]
When $\partial_M$ hits $\hat G(k)$, use $\hat k^2\hat G^2\hat D\le\hat G\hat D$: the contribution is at most $2M\Xi(0)(2\Xi(0)+4\delta)$. When $\partial_M$ hits $\hat G(P)$: at $\theta=1$ the contribution is at most $8M\Xi(0)(6\Xi(0)+2\tau+4\delta)$, and at $\theta=\frac34$ it is at most $2^{5/2}M(2\Lambda_1+4\Lambda_2)\Lambda_0$. Together these give $\abs{\partial_M(D2)}\le\varphi_{D2}$.
\end{itemize}
Summing, $\partial_MQ_T\ge-[\varphi_A+(12+16)M\bar\delta\bar\tau+\varphi_{D2}]=-b_4$.

\emph{Properties of the majorants.} All of them are non-negative and coupling-free. The $\mathcal I$'s and $\bar\omega$ do not depend on $\eps$, and they are non-increasing in $M$ and in $\ell_0$ (Lemmas~\ref{app:G1:PD-oneloop}(b),(c) and~\ref{app:G1:PD-img}). The densities $\varphi_3$ and $b_4$ contain the increasing prefactors $M$ and $\abs{\mu^2-2M^2}$; these are handled by the upper sums over cells in Corollaries~\ref{app:G1:Wa} and \ref{app:G1:Wb}.
\end{proof}

\subsection{Monomial majorants for large masses}\label{app:G1:ssec-monomial}

In this subsection $\mu:=\mu_*$. For $\ell_0=\infty$ we put $\ell_0^{-1}:=0$; by Lemma~\ref{app:G1:PD-oneloop}(a) every majorant $\mathcal I_{\alpha,\beta}(M,\mu;\infty)$ then reduces to its term $c_NC_3$. For $g>1/2$ put
\[
\gamma_3(g):=\frac{\Gamma(g-\frac32)}{(4\pi)^{3/2}\Gamma(g)}\ (g>\tfrac32),\qquad \gamma_2(g):=\frac1{g-1}\ (g>1),\qquad \gamma_1(g):=\frac{\sqrt\pi\,\Gamma(g-\frac12)}{\Gamma(g)} .
\]
These are the closed forms of $\int_{\R^3}(\abs k^2+a^2)^{-g}\frac{d^3k}{(2\pi)^3}$, $\int_0^\infty(u+a^2)^{-g}du$ and $\int_\R(q^2+a^2)^{-g}dq$ at $a=1$. The integral at mass $a$ is $\gamma_d(g)a^{d-2g}$.

\begin{lemma}[one one-loop majorant]\label{app:G1:mono-oneloop}
Let $\alpha,\beta\ge0$ with $g:=\alpha+\beta>3/2$, let $\mu>0$ and $\ell_0\in(0,\infty]$. Put $q_d:=(d-2g)\alpha/g$ for $d=1,2$ and $q_0:=-2\alpha$. Put $q_3:=3-2g$ if $\beta<3/2$, and $q_3:=(3-2g)\alpha/g$ if $\beta\ge3/2$. Let $q_{\alpha,\beta}:=\max\{q_0,q_1,q_2,q_3\}$. Then $q_{\alpha,\beta}\le0$, and $q_{\alpha,\beta}<0$ if $\alpha>0$. For every $M\ge\mu$,
\[
\mathcal I_{\alpha,\beta}(M,\mu;\ell_0)\le c_Nk_3M^{q_3}+\tfrac{3\pi}{16\ell_0}k_2M^{q_2}+\tfrac3{4\ell_0^2}k_1M^{q_1}+\tfrac{\mu^{-2\beta}}{\ell_0^3}M^{q_0}\le K_{\alpha,\beta}M^{q_{\alpha,\beta}} .
\]
Here $k_d:=\gamma_d(g)\mu^{(d-2g)\beta/g}$ for $d=1,2$. Further, $k_3:=(4\pi^2)^{-1}B(\frac32-\beta,g-\frac32)$ if $\beta<3/2$, and $k_3:=\gamma_3(g)\mu^{(3-2g)\beta/g}$ otherwise. Finally,
\[
K_{\alpha,\beta}:=c_Nk_3\mu^{q_3-q_{\alpha,\beta}}+\tfrac{3\pi}{16\ell_0}k_2\mu^{q_2-q_{\alpha,\beta}}+\tfrac3{4\ell_0^2}k_1\mu^{q_1-q_{\alpha,\beta}}+\tfrac{\mu^{-2\beta}}{\ell_0^3}\mu^{q_0-q_{\alpha,\beta}} .
\]
This constant depends only on $\alpha,\beta,\mu,\ell_0$.
\end{lemma}

\begin{proof}
By Lemma~\ref{app:G1:PD-oneloop}, at $(a,b)=(M,\mu)$,
\[
\mathcal I_{\alpha,\beta}=c_NC_3+\tfrac{3\pi}{16\ell_0}C_2+\tfrac3{4\ell_0^2}C_1+M^{-2\alpha}\mu^{-2\beta}\ell_0^{-3},
\]
where each $C_d$ is the minimum of its admissible closed forms. So $C_d$ is at most any one of these forms.

Let $d=3$ and $\beta<3/2$. Then $\beta<\frac32<2g$, so the Beta form with $c=\frac{3-2\beta}2$ is admissible. Here $\alpha-c=g-\frac32>0$, and the form equals $(4\pi^2)^{-1}B(c,\alpha-c)M^{2c-2\alpha}=k_3M^{3-2g}$.

In every other case use the H\"older form $C_d(g,M;0)^{\alpha/g}C_d(g,\mu;0)^{\beta/g}$. Its factor $C_d(g,a;0)$ is the closed form $\gamma_d(g)a^{d-2g}$ of the integral of $(u+a^2)^{-g}$. The form is the bound that H\"older's inequality with exponents $g/\alpha$, $g/\beta$ gives for $(u+M^2)^{-\alpha}(u+\mu^2)^{-\beta}$, and it equals $\gamma_d(g)\mu^{(d-2g)\beta/g}M^{(d-2g)\alpha/g}$. A factor with exponent $0$ is read as $1$. This proves the first inequality.

Since $g>3/2$ and $d\le3$, we have $d-2g<0$ and $3-2g<0$. Hence every $q_d\le0$, and every $q_d<0$ when $\alpha>0$. In the case $\alpha=0$ we have $\beta=g>3/2$, so all $q_d=0$.

For $M\ge\mu$ and $q_d\le q_{\alpha,\beta}$ we have $M^{q_d}=M^{q_{\alpha,\beta}}(M/\mu)^{q_d-q_{\alpha,\beta}}\mu^{q_d-q_{\alpha,\beta}}\le\mu^{q_d-q_{\alpha,\beta}}M^{q_{\alpha,\beta}}$. This gives the second inequality.
\end{proof}

\begin{lemma}[monomial majorants for large masses]\label{app:K:L15-1}\label{app:G1:mono}
Let $\mu_*>0$ and $\ell_0\in(0,\infty]$, where $\ell_0=\infty$ stands for the lattice $\eps\Z^3$. Let $\varphi_3$ and $b_4$ be the functions of Lemma~\ref{app:G1:dens}, formed with the majorants $\bar\tau,\bar\psi,\bar\delta,\bar\Theta,\bar J_s,\bar J'_s,\bar\Lambda_i$ and the constant $\bar\omega$ evaluated at $\ell_0$. Then there are $n\in\N$ and numbers $c_1,\dots,c_n\ge0$, $c'_1,\dots,c'_n\ge0$, $p_1,\dots,p_n\le0$, $p'_1,\dots,p'_n\le0$, depending only on $\mu_*$ and $\ell_0$, such that for every $M\ge\mu_*$
\[
\varphi_3(M)\le\sum_{i=1}^nc_iM^{p_i},\qquad b_4(M)\le\sum_{i=1}^nc'_iM^{p'_i}.
\]
These monomials are obtained as follows: each majorant is replaced by an admissible majorant of the form (constant)$\times M^{q}$, $\abs{\mu_*^2-2M^2}$ is replaced by $2M^2$, and the minimum in $\varphi_{D2}$ is replaced by its first branch. One may take $n=5$, $c_1=c$ (the constant defined in the proof), $p_1=-1$, $c_i=0$ and $p_i=0$ for $i\ge2$, and the pairs $(c'_i,p'_i)$ of~\eqref{app:G1:eq-mono-b4}. In particular, for every $M_c\ge\mu_*$,
\[
\sup_{M\ge M_c}\varphi_3(M)\le\sum_ic_iM_c^{p_i},\qquad \sup_{M\ge M_c}b_4(M)\le\sum_ic'_iM_c^{p'_i}.
\]
\end{lemma}

\begin{proof}
Let $M\ge\mu=\mu_*$. The majorants are the instances of $\mathcal I_{\alpha,\beta}(M,\mu;\ell_0)$ listed in Lemma~\ref{app:G1:PD-oneloop}. By Lemma~\ref{app:G1:mono-oneloop}, each of them is at most $K_XM^{q_X}$, where $K_X$ is the corresponding constant $K_{\alpha,\beta}$, or the sum of two such constants for $\bar\Theta$. The exponents $(q_3,q_2,q_1,q_0)$ of Lemma~\ref{app:G1:mono-oneloop} and their maximum $q_X$ are:
\begin{center}
\begin{tabular}{llll}
majorant & instance & $(q_3,q_2,q_1,q_0)$ & $q_X$\\\hline
$\bar\tau$ & $\mathcal I_{2,0}$ & $(-1,-2,-3,-4)$ & $-1$\\
$\bar\psi$ & $\mathcal I_{1,1}$ & $(-1,-1,-\frac32,-2)$ & $-1$\\
$\bar\delta$ & $\mathcal I_{0,2}$ & $(0,0,0,0)$ & $0$\\
$\bar\Theta$ & $\mathcal I_{1,3/4}$, $\mathcal I_{3/4,1}$ & $(-\frac12,-\frac67,-\frac{10}7,-2)$, $(-\frac12,-\frac9{14},-\frac{15}{14},-\frac32)$ & $-\frac12$\\
$\bar J_0,\bar J_1,\bar J_{3/2}$ & $\mathcal I_{3,0},\mathcal I_{5/2,0},\mathcal I_{9/4,0}$ & $(d-2g)_{d=3,2,1,0}$ & $-3,-2,-\frac32$\\
$\bar J'_0$ & $\mathcal I_{2,1}$ & $(-3,-\frac83,-\frac{10}3,-4)$ & $-\frac83$\\
$\bar J'_1$ & $\mathcal I_{3/2,1}$ & $(-2,-\frac95,-\frac{12}5,-3)$ & $-\frac95$\\
$\bar J'_{3/2}$ & $\mathcal I_{5/4,1}$ & $(-\frac32,-\frac{25}{18},-\frac{35}{18},-\frac52)$ & $-\frac{25}{18}$
\end{tabular}
\end{center}
Further, $\bar\omega\ge0$ is a number that does not depend on $M$ (Lemma~\ref{app:G1:PD-img}(c); for $\ell_0=\infty$ one takes $\bar\omega=w$). For $M\ge\mu$ we have $2M^2-\mu^2\ge M^2\ge0$, hence $\abs{\mu^2-2M^2}\le2M^2$. Finally $\min\{x,y\}\le x$, so the majorants $\bar\Lambda_i$ are not needed.

\emph{The function $\varphi_3$.} With $\bar\beta=4M\bar\tau+16M\bar\psi$,
\[
\varphi_3(M)\le8\sqrt2M\bar\Theta\bar J_{3/2}+8M^3\bar\psi\bar J_0+12\sqrt3\,\bar\omega M\bar J_1+4M\bar\tau^2+16M\bar\psi\bar\tau .
\]
The exponents of the five terms are $1-\frac12-\frac32$, $3-1-3$, $1-2$, $1-2$ and $1-2$, all equal to $-1$. Hence $\varphi_3(M)\le cM^{-1}$, with
\[
c:=8\sqrt2K_{\Theta}K_{J_{3/2}}+8K_\psi K_{J_0}+12\sqrt3\,\bar\omega K_{J_1}+4K_\tau^2+16K_\psi K_\tau .
\]

\emph{The function $b_4$.} We have $\varphi_{D2}\le2M\bar\psi(2\bar\psi+4\bar\delta)+8M\bar\psi(6\bar\psi+2\bar\tau+4\bar\delta)=52M\bar\psi^2+16M\bar\psi\bar\tau+40M\bar\psi\bar\delta$. Also
\[
\varphi_A\le12\sqrt2M\bar\Theta\bar J'_{3/2}+12M^3\bar\psi\bar J'_0+18\sqrt3\,\bar\omega M\bar J'_1+4M\bar\tau\bar\psi+16M\bar\psi^2 .
\]
The exponents are $1-\frac12-\frac{25}{18}=-\frac89$, $3-1-\frac83=-\frac23$, $1-\frac95=-\frac45$, $-1$ and $-1$. The terms $28M\bar\delta\bar\tau$ and $40M\bar\psi\bar\delta$ have exponent $0$, and the terms $52M\bar\psi^2$ and $16M\bar\psi\bar\tau$ have exponent $-1$. Hence $b_4(M)\le\sum_{i=1}^5c'_iM^{p'_i}$ with
\begin{equation}\label{app:G1:eq-mono-b4}
\begin{aligned}
&(c'_1,p'_1)=(28K_\delta K_\tau+40K_\psi K_\delta,\,0),\quad (c'_2,p'_2)=(12K_\psi K_{J'_0},\,-\tfrac23),\quad (c'_3,p'_3)=(18\sqrt3\,\bar\omega K_{J'_1},\,-\tfrac45),\\
&(c'_4,p'_4)=(12\sqrt2K_\Theta K_{J'_{3/2}},\,-\tfrac89),\quad (c'_5,p'_5)=(20K_\tau K_\psi+68K_\psi^2,\,-1).
\end{aligned}
\end{equation}

All the constants are non-negative and depend only on $\mu_*$ and $\ell_0$. For $M\ge M_c\ge\mu_*$ and $p\le0$ we have $M^p\le M_c^p$, which gives the last claim.
\end{proof}

\begin{remark}
The conclusion of Lemma~\ref{app:G1:mono}, including the bound on $\sup_{M\ge M_c}$, holds for every choice of the following kind. Each slice $C_d$ of each instance $\mathcal I_{\alpha,\beta}$ in the table is bounded on $[\mu_*,\infty)$ by $\kappa M^{q}$, where $\kappa\ge0$, $q\le q_X$, and $\kappa M^q$ is an admissible form of Lemma~\ref{app:G1:PD-oneloop} or an upper bound for one there (for example $\frac1{4\pi(M+\mu)}\le\frac1{4\pi M}$). The zero-mode term is $\mu^{-2\beta}\ell_0^{-3}M^{-2\alpha}$, and the slices may be kept as separate monomials. Proof: a slice is at most any of its admissible forms, so each majorant $X$ is at most a finite sum of non-negative multiples of powers $M^q$ with $q\le q_X$. Each term of $\varphi_3$ and $b_4$ is a non-negative prefactor $M^r$ times a product of such majorants. Exponents add, so every monomial in its expansion has an exponent at most the exponent computed in the proof above, and that exponent is $\le0$. The bound on $\sup_{M\ge M_c}$ then follows from $M^p\le M_c^p$ for $p\le0$. In particular this covers choosing, for each slice, the pure-power option with the smallest exponent: the Beta and H\"older forms used above are themselves such options, so the smallest exponent is at most theirs.
\end{remark}

\section{The c-number bounds}\label{app:G1:part-D5}

This appendix proves the c-number bounds used in Section~\ref{sec:chain}: the normalisation identities behind Propositions~\ref{prop:ch-CN} and~\ref{prop:ch-T5} (Lemma~\ref{app:G1:vacnorm}), the bound on c-number differences of Proposition~\ref{prop:ch-CN} (Proposition~\ref{app:G1:CN}), and the two bounds \eqref{p:ch-E2B} and \eqref{p:ch-CNu} of Proposition~\ref{prop:ch-T5} (Lemma~\ref{app:G1:E2B} and Proposition~\ref{app:G1:CNuB}). The inputs from power counting, the weighted-line form of \cite[Prop.~1]{Ba3} and the kernel bounds are in Subsection~\ref{app:G1:part-TWP1}.

\emph{Notation.} Every symbol of Section~\ref{sec:chain} keeps its meaning. In addition: $\hat p^2:=\sum_\mu(4/\eps^2)\sin^2(\eps p_\mu/2)$; $\hat G_M(p):=(\hat p^2+M^2)^{-1}$ is the free scalar propagator per real component; $T^*:=(2\pi/\ell_T)\Z^3\bmod(2\pi/\eps)\Z^3$ is the set of torus momenta, and $\int_p$ means $\abs T^{-1}\sum_{p\in T^*}$ on $T$ and $\int_{[-\pi/\eps,\pi/\eps)^3}d^3p/(2\pi)^3$ on $\eps\Z^3$. Write $C^T(s):=\abs T^{-1}\sum_{p\in T^*}\hat G_{\sqrt s}(p)$ and $C^\infty(s):=\int_{\rm BZ}\hat G_{\sqrt s}$. $X^T_\eps(M)$ is the order-$e^2$ coefficient of $\langle r_0^2\rangle$ on $T$ at Higgs mass $M$ with King's torus counterterm, and $X_\eps(M)$ is the same on $\eps\Z^3$, so $z_\eps=e^2X_\eps(m)$. $[F]_{\alpha\beta}$ is the coefficient of $e^\alpha\lambda^\beta$ in a polynomial $F$ in $(e,\lambda)$, and $T_{\nbar}$ is the Taylor polynomial of total order $\le\nbar$ in $(e,\lambda)$ at $0$, with $\delta m^2_{K,T}(M)$ expanded as well. $Z^F_T(a):=\int \mathrm{e}^{-S_{\mu_*,a}}\prod_bdA_b\prod_xd^2\varphi(x)$ is the Lebesgue integral of the model \eqref{eq:latticeaction} at photon mass $\mu_*$. Note that the photon measure is \emph{not} normalised here.

\paragraph{Normalisation identities.}

\begin{lemma}[vacuum normalisations; identity (3.3)]\label{app:G1:vacnorm}
For every torus $T=T_j$, every spacing of the family and every $M>0$, put $\gamma^\circ_T(M):=\abs T^{-1}[E_0(M)+E_1(M)]+\frac12a_{K,T}(M)(W_\eps+z_\eps)$. Then:
\begin{enumerate}
\item[(a)] $f_T(a_{K,T}(M))=\abs T^{-1}\log Z_K(T;M)+\gamma^\circ_T(M)$; in particular the $\gamma_T(M)$ of Proposition~\ref{prop:ch-CN} equals $\gamma^\circ_T(M)$;
\item[(b)] $\abs T^{-1}\log Z_K(T;M)=\abs T^{-1}\Phi_M(0)+e(\eps;0;M)$, with $\Phi_M$ and $e$ as in Subsection~\ref{ssec:T6-statement};
\item[(c)] $e(\eps;0;M')-e(\eps;0;M)+\frac12\big(a_{K,T}(M)-a_{K,T}(M')\big)(W_\eps+z_\eps)=\gamma_T(M)-\gamma_T(M')$ for all $M,M'>0$.
\end{enumerate}
Consequently, the hypothesis ``Theorem~\ref{thm:UV-K0minus} at $m_1$ in the form $\abs T^{-1}\Phi_{m_1}(0)+e(\eps;0;m_1)\ge-E_-(m_1)$'' of Proposition~\ref{prop:ch-T5} is exactly Theorem~\ref{thm:UV-K0minus} at $M=m_1$.
\end{lemma}
\noindent The lemma uses Proposition~\ref{prop:UV-B1e}. It gives the identity used in Proposition~\ref{prop:ch-CN} and the form of the lower vacuum bound used in Proposition~\ref{prop:ch-T5}.

\begin{proof}
Write $a:=a_{K,T}(M)$ and $N_{\rm ph}:=\int\exp(-\frac12\ip{A}{(-\Delta_\eps+\mu_*^2)A})\prod_bdA_b$. This normalisation does not depend on $M$, $e$ or $\lambda$.
(i) By the definition of $Z(\Gamma;a)$ in Subsection~\ref{ssec:T6-statement} with $\Gamma=C_B(0)=C_{\mu_*}$, $Z^F_T(a)=N_{\rm ph}\,Z(C_{\mu_*};a)$, so $\log Z^F_T(a)=\log N_{\rm ph}+\Phi_M(0)$.
(ii) King's vacuum constant is the Taylor polynomial $E_0(M)+E_1(M)=T_{\nbar}[\log Z^F_T(a)]$ \cite[(2.6), (3.35)]{K1}. By Proposition~\ref{prop:UV-B1e}, i.e.\ Lemma~\ref{app:G1:vac}(a),(b), it is Ba{\l}aban's $E_0+E_1$ of \cite[(1.12)--(1.13), p.~606]{Ba1}; the pairing with the normalised photon is Lemma~\ref{app:G1:vac}(c). There $\exp(E_0)=\exp(E_{0,v})\exp(E_{0,s})$ is the free Lebesgue normalisation, with $E_{0,v}=\log N_{\rm ph}$, and $E_1$ is the sum of the Taylor coefficients of orders $1\le\alpha+\beta\le\nbar$ of $\log\int\!\int\exp(-S+E)$. Since $\log N_{\rm ph}$ is independent of $(e,\lambda)$, $T_{\nbar}[\log N_{\rm ph}+\Phi_M(0)]=\log N_{\rm ph}+T_{\nbar}[\Phi_M(0)]$, and the order-zero term of $T_{\nbar}[\Phi_M(0)]$ is $E_{0,s}(M)$.
(iii) By definition $Z_K(T;M)=Z^F_T(a)\,\mathrm{e}^{-E_0(M)-E_1(M)}$ (Subsection~\ref{ssec:UV-K0}), so by (i)--(ii)
\[
\log Z_K(T;M)=\Phi_M(0)-T_{\nbar}[\Phi_M(0)]=\Phi_M(0)+\abs T\,e(\eps;0;M),
\]
which is (b). (iv) $f_T(a)=\abs T^{-1}\log Z^F_T(a)+\frac12a(W_\eps+z_\eps)=\abs T^{-1}[\log Z_K(T;M)+E_0(M)+E_1(M)]+\frac12a(W_\eps+z_\eps)$, which is (a). (v) By (ii), $e(\eps;0;M)=-\abs T^{-1}[E_0(M)+E_1(M)]+\abs T^{-1}E_{0,v}(\mu_*)$. The last term is free of $M$ and cancels in the difference in (c); what remains is $\gamma^\circ_T(M)-\gamma^\circ_T(M')$. The final sentence of the lemma follows from (b).
\end{proof}

\paragraph{Lattice power counting in three dimensions.}

\emph{Reisz's theorems in $d=3$} are Proposition~\ref{app:G1:reisz3d}, item~3 of Subsection~\ref{ssec:UV-other}.

\emph{Three inputs.} The proof of Proposition~\ref{app:G1:CN} uses three inputs, which we name.
(q1) The transfer from $d=4$ to $d=3$ of Proposition~\ref{app:G1:reisz3d}(a)--(b) (Lemma~\ref{app:G1:WP-R3transfer}).
(q2) The identification of King's counterterm expansion with a Reisz forest sum (Lemmas~\ref{app:G1:forest} and \ref{app:G1:WP-q2}), by the standard counterterm $\Leftrightarrow$ forest combinatorics (Bogoliubov--Zimmermann).
(q3) The constant $z_\eps$ is the graphical one, $z_\eps=e^2X_\eps(m)$ (Subsection~\ref{ssec:observables}); this is the scheme of Subsection~\ref{ssec:ch-scheme}.

\paragraph{The c-number differences.}
Fix anchor masses $0<\underline M<\overline M$ with $\eps\overline M\le1$; for the spacings of the construction $\eps\le1/m_2$ by Lemma~\ref{lem:window}. Put $\hat\Delta:=\hat G_{\underline M}-\hat G_{\overline M}=(\overline M^2-\underline M^2)\hat G_{\underline M}\hat G_{\overline M}$ and $C_t:=\hat G_{\overline M}+t\hat\Delta$ for $t\in[0,1]$. Let $\delta m^2_t$ be King's recursion \cite[(3.33)]{K1} with scalar lines $C_t$ and photon lines $C_{\mu_*}$, and write $\delta\dot m^2_t:=\partial_t\delta m^2_t$. Let $\langle\cdot\rangle_t$ be the model with scalar Gaussian $C_t\otimes\mathbf 1_2$, photon Gaussian $C_{\mu_*}$ and interaction $V_t$ (the $e$-dependent hopping terms, $\lambda\sum\eps^3r^4$ and $\frac12\delta m^2_t\sum\eps^3r^2$). Write $\mathcal G_t:=\langle\varphi\varphi^{\rm T}\rangle_t$, $\mathcal A_t:=C_t^{-1}(\mathcal G_t-C_t)C_t^{-1}$, $\mathcal P_t:=\int_p\operatorname{tr}[\hat G_{\underline M}\hat G_{\overline M}\mathcal A_t]-z_\eps$ and $\rho_t:=\langle r_0^2\rangle_t-W_\eps-z_\eps$.

\begin{lemma}[covariance path]\label{app:G1:covpath}
For every $(\alpha,\beta)\ne(0,0)$,
\[
[\gamma_T(\underline M)-\gamma_T(\overline M)]_{\alpha\beta}=\int_0^1\Big\{\tfrac12(\overline M^2-\underline M^2)[\mathcal P_t]_{\alpha\beta}-\tfrac12[\delta\dot m^2_t\,\rho_t]_{\alpha\beta}\Big\}dt.
\]
Every coefficient of $\mathcal P_t,\rho_t,\delta\dot m^2_t$ is a polynomial in $t$. The $(0,0)$ coefficient equals $\int_{\underline M^2}^{\overline M^2}[C^T(s)-C^\infty(s)]\,ds+\int_{\underline M^2}^{\overline M^2}[C^\infty(s)-C^\infty(m^2)]\,ds$.
\end{lemma}

\begin{proof}
Put $s(t):=\overline M^2+t(\underline M^2-\overline M^2)$ and $\Gamma(t):=\abs T^{-1}T^{\ge1}\log\langle \mathrm{e}^{-V_t}\rangle_{C_t}+\frac12s(t)z_\eps+\frac12\delta m^2_t(W_\eps+z_\eps)$, where $T^{\ge1}$ keeps the orders $1\le\alpha+\beta\le\nbar$. At $t=0,1$ the Gaussian is the free measure at $\overline M$, resp.\ $\underline M$, so $\abs T^{-1}T^{\ge1}\log\langle \mathrm{e}^{-V}\rangle=\abs T^{-1}E_1$, and the c-numbers equal $\frac12a_{K,T}(M)(W_\eps+z_\eps)-\frac12M^2W_\eps$. Hence $\Gamma(1)-\Gamma(0)$ is the order-$\ge1$ part of $\gamma_T(\underline M)-\gamma_T(\overline M)$. Each coefficient of $\Gamma$ is a polynomial in $t$ (a finite Wick sum), so $\Gamma(1)-\Gamma(0)=\int_0^1\Gamma'$ coefficientwise. $\Gamma'$ has three parts.
(i) The covariance: $\frac{d}{dt}\log\int d\mu_{C_t}F=\frac12\operatorname{tr}[\dot C_tC_t^{-1}(\mathcal G-C_t)C_t^{-1}]$ (differentiate the density $\propto\det C^{-1/2}\mathrm{e}^{-\frac12\ip\varphi{C^{-1}\varphi}}$ in finite dimension). With $\dot C_t=\hat\Delta$ this gives, per volume, $\frac12(\overline M^2-\underline M^2)\int_p\operatorname{tr}[\hat G_{\underline M}\hat G_{\overline M}\mathcal A_t]$.
(ii) The vertex: $-\frac12\delta\dot m^2_t\langle r_0^2\rangle_t$.
(iii) The c-numbers: $\frac12(\underline M^2-\overline M^2)z_\eps+\frac12\delta\dot m^2_t(W_\eps+z_\eps)$.
Summing (i)--(iii) gives the identity.
For the $(0,0)$ coefficient: $\abs T^{-1}E_0(M)=-\int_p\log(\hat p^2+M^2)+{}$(an $M$-free term), so $\partial_{M^2}\abs T^{-1}E_0=-C^T(M^2)$. Also $\frac12(\underline M^2-\overline M^2)W_\eps=-\int_{\underline M^2}^{\overline M^2}C^\infty(m^2)\,ds$.
\end{proof}

\begin{lemma}[graph bookkeeping on the torus]\label{app:G1:forest}
Each $t^j e^\alpha\lambda^\beta$-coefficient of $\delta m^2_t-\delta m^2_0$, $\delta\dot m^2_t$, $\mathcal P_t$ and $\rho_t$ (minus its $W_\eps$, $z_\eps$ terms) is a finite rational combination of terms. Each term is a product of three kinds of factor: lower-order coefficients of $\delta m^2_t-\delta m^2_0$; Ward defects $\Pi^T_c(0)$; and a torus sum $S_T(\mathcal S):=\abs T^{-L}\sum_{k\in(T^*)^L}R_{\mathcal S}(k)$ of a Reisz forest integrand of a reduced graph $\mathcal S$. Here:
\begin{itemize}
\item a reduced graph is a graph with lines $\hat G_{\overline M}$ or $\hat\Delta$ and $C_{\mu_*}$ in which pairwise disjoint subgraphs of class R3 of Lemma~\ref{app:G1:CN-list} (two-photon parts of order $o\in\{1,2\}$) are contracted to constant vertices;
\item $\Pi^T_c(0)$ is the zero-momentum value of the \emph{full} order-$o$ Taylor coefficient of the $\hat\Delta$-free, $\delta m^2$-difference-free 1PI two-photon function of the uniform-mass-$\overline M$ model, with King's counterterm $\delta m^2_0$ inside;
\item the forests subtract with $\delta(\gamma)=\omega(\gamma)$ every 1PI proper part $\gamma$ that contains no $\hat\Delta$-line, no $\delta m^2$-difference vertex and not the insertion, where $\omega=3-\frac12(E_\varphi+E_A)-o_V-2v_2$ is the degree of Lemma~\ref{app:G1:CN-list}, with $v_2$ the number of two-point vertices of $\gamma$. By Lemma~\ref{app:G1:CN-list}, only the parts of the classes R2--R6 are actually subtracted.
\end{itemize}
Every forest of a reduced graph is Reisz-admissible. Every reduced graph has $\omega\le-1$, apart from the overall non-negative pieces listed in step~(iv) of Proposition~\ref{app:G1:CN}.
\end{lemma}
\noindent The lemma uses Lemma~\ref{app:G1:WP-q2a}, Lemma~\ref{app:G1:WP-q2a-disj} and Lemma~\ref{app:G1:CN-list}, i.e.\ step~(a) of the proof of Lemma~\ref{app:G1:WP-q2} in its torus form. It also uses the torus Ward identity, Proposition~\ref{app:K:L16-1}.
\begin{proof}
By Lemma~\ref{app:G1:WP-q2a}, which holds on $T$ with sums over $T^*$, the insertion of King's counterterm $\delta m^2_0$ is exactly the sum over the $t^0$-forests of the blocks. By the remark after that lemma and by Lemma~\ref{app:G1:WP-q2a-disj}, the blocks are exactly the scalar self-energy parts that contain no $\hat\Delta$-line, no $f$-vertex and not the insertion and have order at most $2$, i.e.\ the parts of class R2 of Lemma~\ref{app:G1:CN-list}. Also, the forests of blocks are Reisz's forests restricted to these parts. So King's counterterm expansion on $T$ is the forest sum of \cite{Rei88b} with $t^0$ on the R2 parts. We now add the subtractions of the remaining divergent parts, together with the $t^1$ subtractions of the R2 parts.
On $T$ the $t^1-t^0$ terms of R2 parts, and the $t^0$ terms of R4, R5, R6 parts, are odd under $k\mapsto-k$ of the internal momenta. Since $T^*=-T^*$, their torus sums vanish exactly. Adding them does not change the sum. Graphs with a one-photon part vanish identically on $T$, as in step~(d) of the proof of Lemma~\ref{app:G1:WP-q2}. So the R6 parts that contain a $\hat\Delta$-line, an $f$-vertex or the insertion (Lemma~\ref{app:G1:CN-list}(a)) play no role.
For R3 this cancellation fails: a constant photon field on $T$ is a twist, not a gauge transformation, so the class sum $\Pi^T_c(0)$ need not vanish.
For a graph $G$ let $\mathrm{Phys}(G)$ be the torus sum of its integrand with the subtractions described so far, none of them on R3 parts. Let $\mathrm{Full}(G)$ be the same sum with, in addition, every R3 part subtracted with $\delta=\omega$. Here the R3 graphs at a given place are summed over their class, so a subtracted R3 element of class $c$ contributes the factor $-\Pi^T_c(0)$. For a set $P$ of such elements put $\Pi^P:=\prod_{\gamma\in P}\Pi^T_{c(\gamma)}(0)$, where $c(\gamma)$ is the class of $\gamma$, and $(-\Pi)^P:=(-1)^{\abs P}\Pi^P$.
Group the forests by their set $P$ of R3 elements. These are pairwise disjoint, since at $o\le2$ no R3 part contains another. This gives $\mathrm{Full}(G)=\sum_P(-\Pi)^P\,\mathrm{Phys}(G/P)$. The positions form a Boolean lattice, so M\"obius inversion gives $\mathrm{Phys}(G)=\sum_P\Pi^P\,\mathrm{Full}(G/P)$.
Admissibility: $\omega(\tau/P)=\omega(\tau)-\sum_{\gamma\in P,\gamma\subset\tau}\omega(\gamma)\le\delta(\tau)$. This is an over-subtraction, which \cite[(4-3)]{Rei88b} allows. Condition (4-4) holds because the dropped elements have $\delta\ge0$.
The Ward identity of Lemma~\ref{app:G1:CNward}, i.e.\ Proposition~\ref{app:K:L16-1} with $a=\overline M^2+\delta m^2_0$, holds only for the full class coefficient of the gauge-invariant model of uniform mass. The sum above runs over exactly that class.
The assertions on $\omega$ follow from Lemma~\ref{app:G1:CN-list}(b).
\end{proof}

\begin{lemma}[torus sums]\label{app:G1:CNtorus}
Let $R$ be a forest integrand as in Proposition~\ref{app:G1:reisz3d}(b)--(c) with $L$ loops, $3h+\deg_HR\le-1$ for all $H$, and bounded external momenta. Put $\mathfrak s(x):=\sum_{n\in\Z^3\setminus0}\mathrm{e}^{-x\abs n}$. For $\ell_T\ge\ell_0$ and $\eps<a_0$, $\abs{S_T-I_\infty}\le C'[(1+\mathfrak s(\eta_0\ell_T))^L-1]$, where $I_\infty:=\int_{{\rm BZ}^L}R$. Hence $\sup_{\eps,\ell_T\ge\ell_0}\abs{S_T}<\infty$.
\end{lemma}
\noindent The lemma uses Proposition~\ref{app:G1:reisz3d}.
\begin{proof}
Poisson summation on the lattice torus gives $\abs T^{-L}\sum_{(T^*)^L}f=\sum_{n\in\Z^{3L}}\int_{{\rm BZ}^L}f(k)\mathrm{e}^{i\ell_Tn\cdot k}$; check it on the characters $\mathrm{e}^{-ik\cdot x}$.
For $n\ne0$, shift $k_j\mapsto k_j+i\eta_0n_j/\abs{n_j}$. This is allowed because $R$ is analytic in the strip and periodic, so the vertical sides of the contour cancel. The shift gives $\abs{\int R\mathrm{e}^{i\ell_Tn\cdot k}}\le \mathrm{e}^{-\eta_0\ell_T\sum_j\abs{n_j}}C'$ by Proposition~\ref{app:G1:reisz3d}(c).
Finally $\sum_{n\ne0}\mathrm{e}^{-x\sum_j\abs{n_j}}=(1+\mathfrak s(x))^L-1$, and $\abs{I_\infty}\le C$ by Proposition~\ref{app:G1:reisz3d}(a).
\end{proof}

\begin{lemma}[torus Ward identity and its defect]\label{app:G1:CNward}
(a) For $o\in\{1,2\}$, the order-$o$ 1PI two-photon coefficient $\Pi^{(o)}$ of the uniform-mass model on $T$ satisfies $\sum_\nu\Pi^{(o)}_{\mu\nu}(q)\tilde q_\nu=0$ for all $q\in T^*$, with $\tilde q_\nu:=(\mathrm{e}^{i\eps q_\nu}-1)/\eps$. In particular $\Pi^{(o)}_{\mu\nu}(q_1e_\nu)=0$ for $q_1:=2\pi/\ell_T$. The value $\Pi^{(o)}(0)$ is not constrained.
(b) For $\overline M\in\{m,m_1,m_2\}$, $\sup_{\eps<a_0,\ell_T\ge\ell_0}\abs{\Pi^T_c(0)}<\infty$.
\end{lemma}
\noindent Part (a) is unconditional; part (b) uses Proposition~\ref{app:G1:reisz3d}.
\begin{proof}
(a) Every coefficient is a Gaussian moment on the finite torus. The change of variables $\delta\varphi=e\chi q\varphi$, $\delta A=-d\chi$ ($\chi$ periodic) leaves the matter action invariant (Subsection~\ref{ssec:action}) and has Jacobian $1$. It changes $\frac12\ip A{(-\Delta_\eps+\mu_*^2)A}$ by a term \emph{linear} in $A$. Gaussian integration by parts then shows that $\Gamma_{\rm int}$ (the Legendre transform minus the free quadratic form) is invariant to first order under $A\mapsto A-d\chi$, so $\Pi\,d\chi=0$. The gradients $d\chi$ span exactly the non-constant modes.
(b) By (a), $\Pi^T(0)=-[\Pi^T(q_1e_\nu)-\Pi^T(0)]$. The bracket is the torus sum of the $(1-t^\omega)$-subtracted class integrand at the fixed external momentum $q_1e_\nu$; for $o=1$ the extra $t^1-t^0$ part is odd and sums to $0$. This integrand is Reisz-admissible with $3h+\deg_H\le-1$, so Lemma~\ref{app:G1:CNtorus} applies.
\end{proof}

\begin{lemma}[one propagator]\label{app:G1:CNprop}
For $\eps M\le1$ and $x\in\eps\Z^3\setminus0$, $\abs{G^\eps_M(x)}\le77.1\,M^{-3}\abs x^{-4}$, where $G^\eps_M(x):=\int_{\rm BZ}\hat G_M\mathrm{e}^{ik\cdot x}$. Hence $\abs{C^T(M^2)-C^\infty(M^2)}\le77.1\,\zeta_3(4)M^{-3}\ell_T^{-4}$, where $\zeta_3(4):=\sum_{n\in\Z^3\setminus0}\abs n^{-4}\le16.6$.
\end{lemma}

\begin{proof}
$\abs x^4\abs{G(x)}\le\int_{\rm BZ}\abs{\Delta_k^2\hat G_M}$, by two periodic integrations by parts. Put $h=\hat k^2$ and $s=h+M^2$. Then $\abs{\partial_\mu h}\le2h_\mu^{1/2}$, $\abs{\partial^2_\mu h}\le2$, $\abs{\partial^3_\mu h}\le2\eps$, $\abs{\partial^4_\mu h}\le2\eps^2$, and $\eps^2\le24/s$. Fa\`a di Bruno gives $\abs{\Delta^2\hat G}\le2000s^{-3}$. Since $h\ge(4/\pi^2)\abs k^2$, $\int_{\rm BZ}s^{-3}\le\pi^2/(256M^3)$, and $2000\pi^2/256\le77.1$. The torus bound follows from $C^T-C^\infty=\sum_{n\ne0}G^\eps_M(n\ell_T)$.
\end{proof}

\begin{lemma}[the order $e^2$ of $\ev{\abs{\varphi(0)}^2}$ on the torus]\label{app:G1:CNH1}
For $\eps M\le1$ and $\ell_T\ge\ell_0$, $\abs{X^T_\eps(M)-X_\eps(M)}\le C_{z}(M,\mu_*)\,\zeta_6(8)\,\ell_T^{-8}$, where $\zeta_6(8):=\sum_{N\in\Z^6\setminus0}\abs N^{-8}\le22.1$.
\end{lemma}
\noindent The lemma uses Proposition~\ref{app:G1:reisz3d}(a)--(b); the constant is Reisz's and is not explicit.
\begin{proof}
Let $\sigma(k,q)$ be the one-loop $e^2$ self-energy integrand (exchange plus seagull). Then $X^T_\eps=c_X\abs T^{-2}\sum_{k,q\in T^*}\hat G_M(k)^2[\sigma(k,q)-\sigma(0,q)]$, and $X_\eps$ is the same with integrals.
Since $\sigma(-k,-q)=\sigma(k,q)$ and $T^*=-T^*$, $\sigma$ may be replaced by $\tilde\sigma:=\frac12[\sigma(k,q)+\sigma(k,-q)]$, which has $\partial_k\tilde\sigma(0,q)=0$.
The line set $\{k,q,k+q,k-q\}$ of the symmetrised integrand is \emph{not} natural, so Proposition~\ref{app:G1:reisz3d} is applied to the two halves $\frac12\hat G_M(k)^2(1-t^1_k)\sigma(k,\pm q)$ separately. Each half has the natural line set $\{k,q,k\pm q\}$ and $\delta(\gamma)=\omega(\gamma)=1$, and the two halves sum to the symmetrised integrand pointwise.
Each derivative lowers the overall degree $0$ and raises none, so every term of $(\Delta_k+\Delta_q)^4$ applied to each half has $3h+\deg_H\le-1$, and its $L^1$ norm is at most $C_4$.
Poisson summation in $(k,q)$ and four integrations by parts give $\abs{\int\!\int(\cdot)\mathrm{e}^{i\ell_TN\cdot(k,q)}}\le C_4(\ell_T\abs N)^{-8}$. Sum over $N\ne0$.
\end{proof}

\begin{proposition}[c-number differences on tori]\label{app:G1:CN}
For every pair of masses $0<\underline M<\overline M\le m_2$ with $\overline M\in\{m,m_1,m_2\}$, and every $(\alpha,\beta)$ with $\alpha+\beta\le\nbar$, $\sup_{K,\,\ell_T\ge\ell_0}\abs{[\gamma_T(\overline M)-\gamma_T(\underline M)]_{\alpha\beta}}<\infty$. Hence $\sup_{K,\ell_T\ge\ell_0}\abs{\gamma_T(M)-\gamma_T(M')}<\infty$ for every pair of distinct anchor masses $M,M'\in\{M_1,m,m_1,m_2\}$: since $M_1<m<m_1<m_2$, the larger mass of the pair lies in $\{m,m_1,m_2\}$.
\end{proposition}
\noindent The proposition uses the following results:
\begin{itemize}
\item Proposition~\ref{app:G1:reisz3d};
\item Lemma~\ref{app:G1:WP-q2}, through Lemma~\ref{app:G1:forest}; its step~(a) is Lemma~\ref{app:G1:WP-q2a};
\item the Ward identities of Proposition~\ref{app:K:L16-1} and Corollary~\ref{app:K:L16-2};
\item the Dyson chains of Lemma~\ref{app:K:L16-3};
\item the choice $z_\eps=e^2X_\eps(m)$ of the scheme (Subsection~\ref{ssec:ch-scheme}).
\end{itemize}
It gives Proposition~\ref{prop:ch-CN}. By Lemma~\ref{app:G1:vacnorm}(a), the $\gamma_T$ of that proposition is the $\gamma_T$ here.
\begin{proof}
Lemma~\ref{app:G1:covpath} reduces the claim to the $t^j$-coefficients of $\mathcal P_t,\rho_t,\delta\dot m^2_t$, since $\int_0^1t^jdt\le1$ and the bound is two-sided. The coefficients of $\mathcal P_t$ and $\rho_t$ are built from the scalar two-point function $G^s$, so by Lemma~\ref{app:K:L16-3}(a) they are sums of weights of connected graphs. The coefficients of $\delta\dot m^2_t$ are sums of weights of connected graphs by King's recursion \cite[(3.33)]{K1}; see (iii) below. By Lemma~\ref{app:G1:forest}, each such coefficient is a finite sum of products of the three kinds of factor listed there. These factors and the overall non-negative pieces are bounded as follows.
(i) Torus sums of reduced graphs with $\omega\le-1$: by Lemma~\ref{app:G1:CNtorus}.
(ii) Ward defects: by Lemma~\ref{app:G1:CNward}(b), at the uniform mass $\overline M$.
(iii) The coefficients of $\delta m^2_t-\delta m^2_0$ and of $\delta\dot m^2_t$ are zero-momentum self-energies. King's recursion \cite[(3.33)]{K1} defines them through the irreducible two-point pieces, and these are the 1PI pieces by Lemma~\ref{app:K:L16-3}(c), in its second case (scalar covariance $C_t\otimes\mathbf 1_2$ and counterterm $\delta m^2_t$). They contain at least one $\hat\Delta$-line or $\delta m^2$-difference vertex, so $\omega\le-1$ by Lemma~\ref{app:G1:CN-list}(b). They are bounded by (i)--(ii), by induction on the order.
(iv) The overall non-negative pieces:
\begin{itemize}
\item $\rho_t$ at order $0$ equals $2[C^T(\overline M^2)-C^\infty(\overline M^2)]+2t[S_T-I_\infty](\hat\Delta)+2\int(C_t-\hat G_m)$. We bound it by three results: Lemma~\ref{app:G1:CNprop}; Lemma~\ref{app:G1:CNtorus}, since $\hat\Delta$ has $\omega=-1$; and the free bound $\abs{C^\infty(M^2)-C^\infty(m^2)}\le\abs{M-m}\pi^2/32$ (Lemma~\ref{app:G1:WP-freeC}), which holds uniformly in $\eps$.
\item The coefficients of $e^0\lambda^1$ of $\rho_t$ and $\mathcal P_t$ vanish identically on $T$, for every $t\in[0,1]$. At $e=0$ the photon decouples. By Lemma~\ref{app:K:L16-3}(b), to first order in $\lambda$ we have $\mathcal G_t=C_t-C_t\Sigma_tC_t$. Here $\Sigma_t$ is the order-$\lambda$ self-energy: the $\lambda$-tadpole plus the order-$\lambda$ part of $\delta m^2_t$. By translation invariance the tadpole is a multiple of $C_t(x,x)=\abs T^{-1}\sum_{p\in T^*}C_t(p)$, so it is independent of the site and of the momentum. King's recursion \cite[(3.33)]{K1}, computed with the same lines $C_t$, fixes the order-$\lambda$ part of $\delta m^2_t$ so that the order-$\lambda$ self-energy vanishes at zero momentum. Hence $\Sigma_t=0$ identically, so $[\mathcal A_t]_{01}=0$ and $[\ev{r_0^2}_t]_{01}=0$. Finally, $W_\eps$ has order $0$ and $z_\eps=e^2X_\eps(m)$ has order $e^2$.
\item At order $e^2$, the graphs that contain a $\hat\Delta$-line or an $f$-vertex have $\omega\le-1$ (Lemma~\ref{app:G1:CN-list}(b)), and (i) and (iii) cover them. The remaining part of $\rho_t$ is $e^2[X^T_\eps(\overline M)-X_\eps(m)]=e^2[X^T_\eps(\overline M)-X^T_\eps(m)]+e^2[X^T_\eps(m)-X_\eps(m)]$. Lemma~\ref{app:G1:CNH1} bounds the second bracket.

For the first bracket, write, as in the proof of Lemma~\ref{app:G1:CNH1}, $X^T_\eps(M)=c_X\,S_kS_q\,\hat G_M(k)^2[\sigma_M(k,q)-\sigma_M(0,q)]$, where $\sigma_M$ is the one-loop $e^2$ self-energy integrand at mass $M$ and $S:=\abs T^{-1}\sum_{T^*}$. Put $\Delta':=\hat G_{\overline M}-\hat G_m=(m^2-\overline M^2)\hat G_{\overline M}\hat G_m$ and $\Delta\sigma:=\sigma_{\overline M}-\sigma_m$. Then
\begin{multline*}
\frac{X^T_\eps(\overline M)-X^T_\eps(m)}{c_X}=S_kS_q\,\Delta'(\hat G_{\overline M}+\hat G_m)(k)\big[\sigma_m(k,q)-\sigma_m(0,q)\big]\\
+S_kS_q\,\hat G_{\overline M}(k)^2\Delta\sigma(k,q)-\big[S_k\hat G_{\overline M}(k)^2\big]\big[S_q\Delta\sigma(0,q)\big].
\end{multline*}
In the first term the self-energy subgraph is subtracted at zero momentum. Since $\sigma_m(-k,-q)=\sigma_m(k,q)$ and $T^*=-T^*$, we have $S_q\partial_k\sigma_m(0,q)=0$, so $t^0$ may be replaced by $t^1$. The integrand is then a Reisz forest integrand with $\delta(\gamma)=\omega(\gamma)=1$ and overall degree $-2$. In the second term, $\Delta\sigma$ contains the line $\Delta'(k+q)$. The Zimmermann subspaces have degrees $-1,-3,-1$ (with $q$, $k$, resp.\ $k+q$ varying) and the overall degree is $-2$, so no subtraction is needed. The third term is a product of two one-loop sums of degree $-1$. All line sets are $\{k,q,k+q\}$, which is natural, and there is no two-photon subgraph. Lemma~\ref{app:G1:CNtorus} therefore bounds each term uniformly in $\eps<a_0$ and $\ell_T\ge\ell_0$. The same decomposition treats $\mathcal P_t$, with $\hat G_{\underline M}\hat G_{\overline M}-\hat G_m^2=(\hat G_{\underline M}-\hat G_m)\hat G_{\overline M}+\hat G_m(\hat G_{\overline M}-\hat G_m)$ in place of $\hat G_{\overline M}^2-\hat G_m^2$.
\end{itemize}
(v) The $(0,0)$ coefficient: by Lemma~\ref{app:G1:covpath}, Lemma~\ref{app:G1:CNprop} and the same free bound.
(vi) There are only finitely many coarse spacings $\eps\ge\eps_1$. There every quantity is a finite lattice sum of a bounded continuous integrand whose denominators are at least $\min(\underline M^2,\mu_*^2)$, and $\abs T^{-L}\sum\abs{R}\le\sup\abs R$ uniformly in $\ell_T$.
Since $\gamma_T(\overline M)-\gamma_T(\underline M)$ is a polynomial in $(e,\lambda)$ of order at most $\nbar$, the coefficient bounds give the bound on the difference.
\end{proof}

\paragraph{The two c-number bounds of Proposition~\ref{prop:ch-T5}.}
In this subsection $M\in\{m_1,m_2\}$. The slice direction is $\Gamma^B=\kappa\mu_*^2D=\sum_j\delta_jH_j$, with $D:=\sum_{j<K}t_{j+1}^2H_j$, and $\mathcal E(u;M)=-\abs T\,e(\eps;u;M)$ is the total (not per-volume) counterterm polynomial. Put $\bar\delta:=\kappa\mu_*^2\eps_0^2$. Write $[\mathcal E]_{1,\alpha\beta}$ for the coefficient of $u\,e^\alpha\lambda^\beta$ in $\mathcal E$, $o:=\alpha/2+\beta$, and $t_K:=L^K\eps\in(\eps_0/L,\eps_0]$ for the final scale. We use the condition (E.$\delta$), $\bar\delta\le1/12$. It gives $C_B(u)>0$ and slice weights $1+u\delta_j\in[1-\eta,2]$ for $u\in[-\eta,1]$, and it implies $\bar\delta\le1$ and $\eta\bar\delta\le\frac12$.

\begin{lemma}[the second derivative in $u$]\label{app:G1:E2B}
Assume the hypotheses of Theorem~\ref{thm:T6}. In particular, the conditions (E.fin) and (T.per) of Definition~\ref{app:G1:rowlist} hold. Let $M\in\{m_1,m_2\}$ and $\eps=2^{-K}\eps_0\le\eps_0\le1$. Let $\nbar>12$ satisfy $2\nbar+1\le3R$, where $R$ is the parameter of \eqref{eq:UV-r}. For $f\in\{s,v\}$ put $M_s:=M$ and $M_v:=\mu_*$, with the photon taken componentwise (Lemma~\ref{app:G1:VSR-V0}). In scale-$k$ units let $G_{k,f}:=(-\Delta^{\eta_k}+M_f^2t_k^2+a_kP_k)^{-1}$ for $1\le k\le K$. Assume that these kernels and their derivatives of order at most $2$ satisfy (H$_\infty$) of Subsection~\ref{ssec:T6-IMG}, with the constants $C_0,\delta_0$ given there. Then, on every torus $T=T_j$ of side $\ell_T\ge\ell_0$,
\[
\sup_\eps\sup_{u\in[-\eta,1]}\abs{\partial_u^2e(\eps;u;M)}\le E_2(M),
\]
where
\begin{multline*}
E_2(M):=C_W\sum_{m'=2}^{N_v}m'^2\sum_{\alpha\ge2m',\,\alpha+\beta\le\nbar}e^\alpha\lambda^\beta(\kappa\mu_*^2L^2)^{m'}\eps_0^{\alpha/2+\beta+2m'-3}\\
=4C_We^4\kappa^2\mu_*^4L^4\eps_0^3\,[1+O(e^2\eps_0+\lambda\eps_0+e^2\kappa\mu_*^2L^2\eps_0^3)].
\end{multline*}
Here $N_v$ is the largest number of slice lines in a graph of order $\le\nbar$. The constant $C_W=C_W(\nbar,\alpha_0,L,a,M_b,\delta_0,M\eps_0,\mu_*\eps_0,\ell_0/\eps_0)$ does not depend on $\eps$, $u$, $\kappa$ or $\ell_T\ge\ell_0$. The bound is two-sided and uniform in $u\in[-\eta,1]$, not only at the two points of Theorem~\ref{thm:T6}.
\end{lemma}
\noindent The lemma uses \cite{Ba3} in the organisation of \cite[(1.12)--(1.13)]{Ba1}, Theorem~\ref{app:K:L17-1}, and Proposition~\ref{prop:UV-B1e}. Theorem~\ref{app:K:L17-1} is the zero-background form of Theorem~\ref{app:G1:WP-WP1}(c) for the vacuum families with at least two slice lines. The lemma is the bound \eqref{p:ch-E2B} of Proposition~\ref{prop:ch-T5}.
\begin{proof}
(1) $\partial_u^2\mathcal E=T_{\nbar}[\Phi_M''(u)]$. Expand $C_B(u)=C_{\mu_*}+u\Gamma^B$ multilinearly. Then $T_{\nbar}[\Phi_M(u)]=\sum_Gu^{m'(G)}\mathrm{val}(G)$, summed over the connected coloured vacuum graphs of order $\le\nbar$, where $m'(G)$ is the number of $\Gamma^B$-lines. Hence $\partial^2_u\mathcal E=\sum_{m'(G)\ge2}m'(m'-1)u^{m'-2}\mathrm{val}(G)$.

(2) By \cite[(1.25)--(1.26)]{Ba3}, $C_{\mu_*}=G_K+(C_{\mu_*}-G_K)$. The second part is regular and is treated as an external field. The slice part $\Gamma^B$ is exactly $c_K$ times the rescaled slice kernel, with $c_K=\kappa\mu_*^2L^2t_K^2$.

(3) Rescale to unit $t_K$.

(4) The graphs fall into the families of Lemma~\ref{app:G1:WP-fam}. Each family has a fixed number $m'$ of $\Gamma^B$-lines and a fixed order $e^\alpha\lambda^\beta$, where the orders of the counterterm vertices are included. No family with $m'\ge2$ needs a vacuum counterterm, because its degree is at least $3$ (Lemma~\ref{app:G1:WP-deg}(c)). By Theorem~\ref{app:K:L17-1}, the coefficient of $u^{m'}$ of such a family is at most $C_Fe^\alpha\lambda^\beta(\kappa\mu_*^2L^2)^{m'}\eps_0^{\alpha/2+\beta+2m'-3}\abs T$. Each slice line has two $e$-legs, so $\alpha\ge2m'$, and the exponent of $\eps_0$ is at least $3m'-3\ge3$.

(5) Multiply by $m'(m'-1)\abs u^{m'-2}\le m'^2$ (because $\abs u\le1$), sum over the finitely many families, and divide by $\abs T$. The constant $C_W$ is the largest $C_F$ times the number of families. It depends only on the data listed in Theorem~\ref{app:K:L17-1}.

The bare coefficient is $a_{K,T}(M)$. By Proposition~\ref{prop:UV-B1e}, it is the reference counterterm of the organisation \cite[(1.12)--(1.13)]{Ba1} at mass $M$ on $T$.

\emph{The inputs of step (4).} Step (4) uses Theorem~\ref{app:K:L17-1} in place of Theorem~\ref{app:G1:WP-WP1}(c). The sketched items I1--I5 listed after Theorem~\ref{app:G1:WP-WP1} enter only in the following way.
\begin{itemize}
\item The expansion of $\mathcal E(u;M)$ has no external fields, zero background and no block-spin transformation. Its vertices are those of the lattice action \cite[(1.6)--(1.8), (1.10)]{Ba3} and the mass-counterterm vertices \cite[(1.7)]{Ba3}. The vertices \cite[(1.9), (1.11)]{Ba3} are Taylor remainders, so they do not occur in a Taylor polynomial. The vertices \cite[(1.12)--(1.15)]{Ba3} arise only from a block-spin transformation. Hence the H\"older norms of the background (I5) do not occur. Since $\mathcal E(u;M)$ is one number on $T$, the relocalisation of the vacuum counterterms between scales (I2) does not occur either, and no vacuum counterterm is attached (Theorem~\ref{app:G1:WP-WP1}(b)).
\item The line bounds (I4) for King's slices $H_{j,f}=G_{j+1,f}-G_{j,f}$ are given by Corollary~\ref{app:G1:VSR-VSphys}. It applies because, by (E.fin), $t_{j+1}\le t_K\le\eps_0\le\min(1/M,1/\mu_*,L)$. By Lemma~\ref{app:G1:WP-Kslices}, King's slices are the slices of \cite[(2.6)]{Ba3}. The remainder lines $C_f-G_{K,f}$ obey Proposition~\ref{app:G1:VSR-R}.
\item The finite counterterms (I3) are bounded by Lemma~\ref{app:K:L17-2}.
\item The combinatorics of the families (I1; step S10 of Table~\ref{tab:prop1}) enters only through Theorem~\ref{app:K:L17-1}. On the torus, the counterterm partners of the classes are compared with their form on $\eta_K\Z^3$ by Lemma~\ref{app:K:L17-3} and Statement~\ref{app:K:L17-4}. These comparisons take place on cubes of side $2\nbar+1$ that do not wrap around the torus, by the hypothesis $2\nbar+1\le3R$, (T.per) and Lemma~\ref{lem:IMG}.
\end{itemize}
\end{proof}

\begin{lemma}[the one-loop mass difference of the term linear in $u$]\label{app:G1:CU1}
Let $\Pi_M(k)_{\mu\nu}:=e^2\big[2\delta_{\mu\nu}S_q\cos(\eps q_\mu)G(q)-S_qV_\mu V_\nu G(q)G(q+k)\big]$, with $V_\mu:=(2/\eps)\sin(\eps(2q+k)_\mu/2)$, $G(q):=(\hat q^2+M^2)^{-1}$ and $S_q:=\abs T^{-1}\sum_{q\in T^*}$. This is the order-$e^2$ vacuum polarisation of the complex scalar on $T_\eps$ with King's hopping, and $[\partial_ue(\eps;0;M)]_{20}=\frac12\abs T^{-1}\operatorname{Tr}(\Gamma^B\Pi_M)$. For $0<m_1<m_2$,
\begin{multline*}
\tfrac12\abs T^{-1}\big\lvert\operatorname{Tr}\big(\Gamma^B(\Pi_{m_1}-\Pi_{m_2})\big)\big\rvert\\
\le10e^2(m_2-m_1)\sup_{M\in[m_1,m_2]}[M\,C_T(M,\ell_0)]\cdot3\kappa\mu_*^2L^2\Big(c_1\eps+\frac{C_H\eps_0}{L-1}\Big).
\end{multline*}
Here $C_T(M,\ell):=a_*^3\sqrt\pi/M+3a_*^2b/M^2+\frac{3\sqrt\pi}2a_*b^2/M^3+b^3/M^4$, with $a_*:=\sqrt\pi/4$ and $b:=4/\ell$. The constants $c_1$, $C_H$ are those of the diagonal bound on the slice kernels (Lemma~\ref{app:G1:WP-RBdiag}).
\end{lemma}
\noindent The lemma uses the diagonal bound of Lemma~\ref{app:G1:WP-RBdiag} and the heat-kernel bound of Corollary~\ref{app:G1:WP-heat}.
\begin{proof}
(1) For $\Gamma\ge0$ and bounded self-adjoint $B$, $\abs{\operatorname{Tr}(\Gamma B)}\le\norm B\operatorname{Tr}\Gamma$. Here $B=\Pi_{m_1}-\Pi_{m_2}$ is a Fourier multiplier, so $\norm B=\sup_k\norm{B(k)}_{\rm op}$.
(2) $\abs T^{-1}\operatorname{Tr}\Gamma^B\le3\kappa\mu_*^2\max_xD(x,x)\le3\kappa\mu_*^2L^2(c_1\eps+C_H\eps_0/(L-1))$ (Lemma~\ref{app:G1:WP-RBdiag}, $H_j(x,x)\le C_H/t_j$).
(3) Differentiating, $\partial_M\Pi_M=e^2[-4M\delta_{\mu\nu}S_q\cos(\eps q_\mu)G^2+2MS_qV_\mu V_\nu(G(q)^2G(q+k)+G(q)G(q+k)^2)]$. Write $\mathsf A:=\hat q$ and $\mathsf B:=\widehat{q+k}$. The sine addition formula gives $\abs{V_\mu}\le\abs{\mathsf A_\mu}+\abs{\mathsf B_\mu}$, so $\abs V^2\le2(\mathsf A^2+\mathsf B^2)$. Since $\mathsf A^2G_{\mathsf A},\mathsf B^2G_{\mathsf B}\le1$, $\abs V^2(G_{\mathsf A}^2G_{\mathsf B}+G_{\mathsf A}G_{\mathsf B}^2)\le4(G_{\mathsf A}^2+G_{\mathsf B}^2)$. By shift invariance of $S_q$, $\norm{\partial_M\Pi_M(k)}_{\rm op}\le20Me^2S_qG(q)^2$.
(4) $S_qG^2=\int_0^\infty t\mathrm{e}^{-tM^2}S_q\mathrm{e}^{-t\hat q^2}dt$, with $S_q\mathrm{e}^{-t\hat q^2}\le(a_*/\sqrt t+4/\ell_T)^3$ (Corollary~\ref{app:G1:WP-heat}, which does not depend on $\eps$). Integrating term by term gives $S_qG^2\le C_T(M,\ell_T)\le C_T(M,\ell_0)$.
(5) Write $\Pi_{m_1}-\Pi_{m_2}=-\int_{m_1}^{m_2}\partial_M\Pi_M\,dM$, bound by (3)--(4), and combine with (1), (2) and the factor $\frac12$.
\end{proof}

\begin{proposition}[the term linear in $u$]\label{app:G1:CNuB}
Assume the hypotheses of Lemma~\ref{app:G1:E2B} at $M=m_1$ and at $M=m_2$. Then $\bar\Gamma^{(1)}_{21}<\infty$. More precisely,
\begin{multline*}
\bar\Gamma^{(1)}_{21}\le\sup_{K,\ell_T\ge\ell_0}[\gamma_T(m_2)-\gamma_T(m_1)]+\sum_{\substack{(\alpha,\beta)\ne(2,0),\,\alpha\ge2\\ \alpha+\beta\le\nbar}}e^\alpha\lambda^\beta\big(C_{1,\alpha\beta}(m_1)+C_{1,\alpha\beta}(m_2)\big)\kappa\mu_*^2L^2\eps_0^{o-1}\\
+\mathfrak C_{\rm CU1}+\tfrac12\big(E_2(m_1)+E_2(m_2)\big).
\end{multline*}
Here $o:=\alpha/2+\beta$, $\mathfrak C_{\rm CU1}$ is the right side of Lemma~\ref{app:G1:CU1}, taken at $\eps\le\eps_0$, and the $C_{1,\alpha\beta}$ are the constants of the bound on the vacuum families linear in $u$ (Corollary~\ref{app:G1:WP-AC11}).
\end{proposition}
\noindent The proposition uses \cite{Ba3} in the organisation \cite[(1.12)--(1.13)]{Ba1}, Proposition~\ref{app:G1:CN}, Lemma~\ref{app:G1:E2B}, Lemma~\ref{app:G1:CU1} and Proposition~\ref{prop:UV-B1e}. It is the bound \eqref{p:ch-CNu} of Proposition~\ref{prop:ch-T5}.
\begin{proof}
Write $a_i:=a_{K,T}(m_i)$. By Lemma~\ref{app:G1:edeg}, $u\mapsto e(\eps;u;M)$ is a polynomial, so
\[
e(\eps;1;M)=e(\eps;0;M)+\partial_ue(\eps;0;M)+\int_0^1(1-u)\,\partial_u^2e(\eps;u;M)\,du.
\]
Insert this into the bracket that defines $\bar\Gamma^{(1)}_{21}$ and treat the three parts in turn.

\emph{Order $u^0$.} By Lemma~\ref{app:G1:vacnorm}(c) with $(M,M')=(m_2,m_1)$, the $u^0$ part is $\gamma_T(m_2)-\gamma_T(m_1)$. Its supremum over the family is finite by Proposition~\ref{app:G1:CN} with the pair $\{m_1,m_2\}$.

\emph{Remainder.} By Lemma~\ref{app:G1:E2B}, the remainder is at most $\frac12E_2(M)$ for each $M$.

\emph{Order $u^1$.} We bound $\abs{\partial_ue(\eps;0;m_1)-\partial_ue(\eps;0;m_2)}$ coefficient by coefficient.
\begin{itemize}
\item Odd $\alpha$: these coefficients vanish. The map $(e,A)\mapsto(-e,-A)$ leaves the action and the centred Gaussian invariant, so $\Phi_M$ is even in $e$.
\item $\alpha=0$: these coefficients vanish. At $e=0$ the matter integral does not depend on $A$, so $\Phi_M(u)|_{e=0}$ does not depend on $u$.
\item $\alpha\ge2$ with $(\alpha,\beta)\ne(2,0)$: then $o\ge2$. The bound $\abs{[\mathcal E]_{1,\alpha\beta}}\le C_{1,\alpha\beta}\kappa\mu_*^2L^2\abs T\eps_0^{o-1}$ (Corollary~\ref{app:G1:WP-AC11}, an instance of Theorem~\ref{app:G1:WP-WP1}(c)) bounds each mass separately, per unit volume.
\item $(\alpha,\beta)=(2,0)$: this is the one-loop graph. It diverges logarithmically at each $M$, and Lemma~\ref{app:G1:CU1} bounds its mass difference. As in that lemma, the Taylor coefficient is taken with the scalar lines at the bare mass $M^2$; the counterterms enter as vertices of higher order.
\end{itemize}
All bounds are uniform in $\eps$ and in $\ell_T\ge\ell_0$, and all sums are finite.

The constants $c_1$ and $C_H$ in Lemma~\ref{app:G1:CU1} come from step~(2) of its proof, which is Lemma~\ref{app:G1:WP-RBdiag}. That lemma rests on two inputs: the block Poincar\'e inequality (Lemma~\ref{app:G1:DOM-poinc}), and the bound $C^{(j)}\le c_Ct_j^2$ on King's fluctuation covariance (Lemma~\ref{app:G1:Cj-bound}). The second follows from the variational form of $\Delta^{(j)}$ (Lemma~\ref{app:G1:Dk-var}). By Lemma~\ref{app:G1:diagC}, one may take $c_1=4.42+\eps_0/(\ell_0^3\mu_*^2)$. Hence $\mathfrak C_{\rm CU1}$ is finite and does not depend on $\eps\le\eps_0$ or $\ell_T\ge\ell_0$.
\end{proof}

\subsection{Power counting, the weighted-line form of Proposition~1, and kernel bounds}\label{app:G1:part-TWP1}

\emph{Notation.} Notation is as at the beginning of this appendix. In addition: $t_k:=L^k\eps$; $\eta:=L^{-k}$ is the lattice spacing in scale-$k$ units; $c_k:=\kappa\mu_*^2L^2t_k^2$; $\delta_j:=\kappa\mu_*^2t_{j+1}^2$, so that $\Gamma^B=\kappa\mu_*^2D=\sum_{j<K}\delta_jH_j$ with $D:=\sum_{j<K}t_{j+1}^2H_j$. $H_j$ is King's photon slice \cite[(2.17)]{K1}. $G^\eta_{(j)}$ is Ba{\l}aban's slice \cite[(2.6), p.~424]{Ba3}. For a polynomial $F$ in $(u,e,\lambda)$, $[F]_{m,\alpha\beta}$ is its $u^me^\alpha\lambda^\beta$ coefficient, and $o:=\alpha/2+\beta$. We write $\mathcal E(u;M)=-\abs T\,e(\eps;u;M)=T_{\nbar}[\Phi_M(u)]$ (Lemma~\ref{app:G1:vacnorm}(b)). All bounds below are on absolute values, so signs play no role.

\paragraph{Lattice power counting: the transfer from $d=4$ to $d=3$.}

\begin{lemma}[where the dimension enters Reisz's proofs]\label{app:G1:WP-R3transfer}
Let $\mathsf d$ be the space-time dimension. In the proofs of \cite[Thm.~1, Thm.~2, Thm.~3]{Rei88} (Sects.~4--9, App.~A--D) and of the Renormalization Theorem \cite[p.~92]{Rei88b} (Sects.~3--6, App.~A--C), the value $\mathsf d=4$ is used at exactly the following places and nowhere else:
\begin{enumerate}
\item[(i)] the degree definitions, each of the form ``$4\times$(dimension of the loop space)'': \cite[(3.8), (5.9), (7.9), (7.14), (9.3), (9.29)]{Rei88} and \cite[(2-28), (4-11)--(4-16), (5-5), (5-10), (C-1)--(C-2)]{Rei88b};
\item[(ii)] the single scaling step \cite[(9.16)--(9.17a)]{Rei88}, where the inner variables are rescaled by $\norm{t_1}$;
\item[(iii)] two constants that use $\sqrt{\mathsf d}\le2$: $\delta$ in \cite[(7.8)]{Rei88}, and step~1 of the proof of \cite[Lemma~D.1]{Rei88} (p.~125).
\end{enumerate}
Replacing $4$ by $\mathsf d$ in (i), the proofs remain valid for every $\mathsf d\le4$, and in particular for $\mathsf d=3$. They then give Proposition~\ref{app:G1:reisz3d}(a)--(b) in the bound (absolute-value) form.
\end{lemma}
\noindent The proofs of \cite{Rei88}, Sects.~4--9 and App.~A--D, and of \cite{Rei88b}, Sects.~3--6 and App.~A, C, are used as printed; the absolute-value form is read off \cite[(7.4)--(7.5)]{Rei88}. Reisz's printed theorems are for $\mathsf d=4$ (``The following considerations can easily be extended to other dimensions'', \cite[p.~86]{Rei88}). For $\mathsf d=3$ the lemma is a reading of a published proof, not a new proof.

\begin{proof}
We go through the items.
(1) \emph{(i): the degree definitions.} $\deg_H\hat I=4\dim H+\deg F$ \cite[(3.8)]{Rei88}; $\omega(H)=4\dim H+\delta(H)-\deg E$ (5.9); the counts (7.9), (7.14); $\omega_0=4m+\Delta-2\sum n_j$ (9.3); and (9.29). In \cite{Rei88b}: $\omega(\gamma)=\sum\omega(\Delta_L)+\sum\omega(V_B)+4m(\gamma)$ (2-28)/(2-31); (4-11)--(4-16); $M_U(\gamma)=4\sum m(\tau(U))$ (5-10). Each is ``$4\times$ the number of 4-dimensional integration variables'', and we replace $4$ by $\mathsf d$.
(2) \emph{App.~C.} Reisz states and uses a \emph{lower} bound: $M_U(\Gamma)\ge4h$ (C-1) and $M_U(\gamma)\ge4\,\mathrm{rank}$ (C-2) \cite[pp.~97, 106]{Rei88b}. Its proof is the count $M_U=4\cdot\#\mathcal T^*\ge4\,\mathrm{rank}$, which is linear in the $4$. Hence $M_U(\Gamma)\ge\mathsf dh$, and this is what gives $\mathsf dh+\deg_HR_\Gamma\le-1$ in (5-13).
(3) \emph{(ii): the scaling step.} The sector $X_\xi$ (9.8) has outer variable $t_1\in\R^{\mathsf d}$ with $\norm{t_1}\ge1$; the inner variables $(t_2,\dots,t_m)\in\R^{\mathsf d(m-1)}$ are rescaled as $t_j=\norm{t_1}t'_j$ (9.16). The Jacobian contributes $\norm{t_1}^{\mathsf d(m-1)}$. The numerator contributes at most $\norm{t_1}^\Delta$ (by (9.19), $r_i-p_i\le\Delta$), and the homogeneous denominator $\norm{t_1}^{-2\sum n_j}$. So (9.17a) becomes $\int d^{\mathsf d}t_1\norm{t_1}^{\omega_0-\mathsf d}\,\mathcal J^{\rm int}$ with $\omega_0:=\mathsf dm+\Delta-2\sum n_j$. Lemma~9.2 applies the induction hypothesis to the inner integral, with $\omega(\tilde H)=\omega(H)$ (9.29).
Reisz does not print the outer integral (``the remaining integration can be done without problems'', p.~119). In polar coordinates, $r^{\mathsf d-1}r^{\omega_0-\mathsf d}=r^{\omega_0-1}$, and $\int_1^{\lambda+K}r^{\omega_0-1}f(\lambda')\,dr$ with $\lambda'=b(\lambda+d)/r+e$ (Lemma 9.1) does not depend on $\mathsf d$:
\begin{itemize}
\item $f=1$, $\omega_0<0$: the integral is bounded;
\item $f=\lambda'^{-1}\log^{m-1}\lambda'$, $\omega_0\le-1$: it is at most $\lambda^{-1}\log^m\lambda$;
\item $f=\lambda'^{\bar\omega}\log^{m-1}\lambda'$: it is at most $\lambda^{\max(\omega_0,\bar\omega)}\log^m\lambda$.
\end{itemize}
This is \cite[(9.30)]{Rei88} for every $\mathsf d$.
(4) \emph{(iii): the two constants.} In (7.8), $\delta=\max(\pi\eps,4\pi(1+\norm{\delta_i}/2))$ must make $\kappa_J$ contain the translated sector. From (4.7), $\norm{k_j}\le(\pi/a)(\sqrt{\mathsf d}+2\norm{\delta_j})$, which suffices if $\sqrt{\mathsf d}\le4$. Lemma~D.1, step~1, uses $\norm{k_j}\le2\pi/a=\sqrt4\,\pi/a$. Both hold a fortiori for $\mathsf d=3$; for $\mathsf d>4$ they fail, which is why the lemma is restricted to $\mathsf d\le4$.
(5) \emph{Dimension-free steps.} The following involve no $\mathsf d$:
\begin{itemize}
\item the function classes and propagator class, Defs.~2.1--2.5 and Lemma~2.1 (vanishing orders in $w$; King's $\hat l^2=\sum_\mu(4/a^2)\sin^2(al_\mu/2)$ is in Def.~2.5 for $\mathsf d=3$);
\item naturalness, Def.~3.1 and App.~D (linear algebra over $\Z$, coordinatewise);
\item the sector decomposition (4.1)--(4.9) and the propagator bounds (4.8)--(4.9);
\item Theorem~3 and Lemmas~6.1--6.3 (Taylor's formula in groups $v_i\in\R^{m_i}$);
\item the powers of $a$ in (7.10)--(7.17) (twice the number of propagators far from a pole);
\item Lemmas~8.1--8.10 and App.~A--B (interpolation, set inclusions, linear changes of variables, Fubini);
\item in \cite{Rei88b}: Lemma~3.1 and App.~A (subtraction operators; $t^0$, $t^1$ with $\sin(\eps q)/\eps$ are of type (A-1)), Lemma~5.1 (Zimmermann's complete forests), and Sect.~6, the index distributions (finite index sets).
\end{itemize}
Masses enter only through lower bounds (all $\mu_i^2>0$), so $\mathsf d=3$ needs no extra infrared condition. Two further items are dimension-free: \cite[Sect.~6]{Rei88b}, and complex-valued numerators. The latter are harmless because $\abs V\le\abs{\rRe V}+\abs{\operatorname{Im}V}$ and real and imaginary parts have vanishing orders at least those of $V$.
(6) \emph{Absolute-value form.} (4.3) is a partition into finitely many sectors by non-negative $\Theta$'s. (7.4) and (7.5a--b) bound each sector by integrals of non-negative majorants, and the integrands of Theorem~2 (5.2) are non-negative. So the proof bounds $\int_{\rm BZ}\abs F$, uniformly for bounded external momenta (p.~96, p.~102). Reducing to $V\in\mathcal C^c_{m_0}$ is harmless, because $\abs V\le\sum\abs{V_i}$.
\end{proof}

\paragraph{The counterterm--forest bookkeeping.}

\begin{lemma}[King's expansion as Reisz forest integrals]\label{app:G1:WP-q2}
On $\eps\Z^3$ take the $C_t$-model of the paragraph on the c-number differences: scalar lines $C_t=\hat G_{\overline M}+t\hat\Delta$, and $\delta m^2_t=\delta m^2_0+f_t$ with $f_t:=\delta m^2_t-\delta m^2_0$, a polynomial in $t$ with $f_0=0$. Then for every $(\alpha,\beta,j)$, the $e^\alpha\lambda^\beta t^j$ coefficient of $\mathcal P_t$, of $\rho_t$ (minus its $W_\eps$, $z_\eps$ terms) and of $f_t$ is a finite sum $\sum_Gc_G\,\prod(\text{lower-order coefficients of }f)\int_{{\rm BZ}^L}R_G$, with rational $c_G$. Here:
\begin{itemize}
\item $G$ runs over 1PI graphs with scalar lines $\hat G_{\overline M}$ or $\hat\Delta$, photon lines, the vertices of King's action, the insertion, and two-point ``$f$-vertices''. There are no $\delta m^2_0$ vertices.
\item $R_G$ is Reisz's forest integrand with $\delta(\gamma)=\omega(\gamma)$ for every 1PI $\gamma\subsetneq G$. This means $t^1$ (in $\sin(\eps q)/\eps$) if $\omega=1$, $t^0$ if $\omega=0$, and no subtraction if $\omega<0$. There is no overall subtraction.
\end{itemize}
The identity holds exactly at every $\eps$. Its torus form, with the Ward-defect correction, is Lemma~\ref{app:G1:forest}.
\end{lemma}
\noindent Step~(a) of the proof is Lemma~\ref{app:G1:WP-q2a} below. It identifies the insertion of King's constant counterterms with the sum over $t^0$-forests (Bogoliubov's $R$-operation for zero-momentum counterterms, in Zimmermann's forest form; cf.\ \cite[pp.~79--80, 94--95]{Rei88b}). The regrouping of the two-photon parts in step~(c) rests on the Ward identity on $\eps\Z^3$, Corollary~\ref{app:K:L16-2}.

\begin{lemma}[the one-particle-irreducible parts with $\omega\ge0$]\label{app:G1:CN-list}
Let $\gamma$ be a connected 1PI graph with a loop. It is built from the following vertices:
\begin{itemize}
\item King's vertices $V_n$, $n\ge1$, of Lemma~\ref{app:G1:WP-q2a};
\item the $\lambda$-vertex;
\item two-point vertices: the $f$-vertices, the vertex inside each $\hat\Delta$-line, and at most one insertion.
\end{itemize}
It contains no $\delta m^2_0$-vertex. Let $E_\varphi$ and $E_A$ be its numbers of external scalar and photon legs, and put $E:=E_\varphi+E_A$. Let $k$ be its number of $\hat\Delta$-lines, $v_f$ its number of $f$-vertices and $v_X\in\{0,1\}$ its number of insertions. Put $o:=o_V(\gamma):=\sum_n\frac n2v_n+v_\lambda$. Its degree is
\[
\omega(\gamma)=3-\tfrac12E-o-2(k+v_f+v_X).
\]
\begin{enumerate}
\item[(a)] If $\omega(\gamma)\ge0$, then $\gamma$ belongs to exactly one of the following classes:
\begin{itemize}
\item[R0] vacuum graphs without insertion, $E=0$, $v_X=0$: $\omega=3-o-2(k+v_f)$, so $o\le3$;
\item[R1] vacuum graphs with one insertion and $k=v_f=0$, as the whole graph only: $o=0$ ($\omega=1$) or $o=1$ ($\omega=0$);
\item[R2] scalar self-energies, $E_\varphi=2$, $E_A=0$: $o=1$ ($\omega=1$) or $o=2$ ($\omega=0$);
\item[R3] photon self-energies, $E_A=2$, $E_\varphi=0$: $o=1$ ($\omega=1$) or $o=2$ ($\omega=0$);
\item[R4] $\varphi\varphi A$ parts, $o=\frac32$, $\omega=0$;
\item[R5] three-photon parts, $o=\frac32$, $\omega=0$;
\item[R6] one-photon parts, $E_A=1$, $E_\varphi=0$.
\end{itemize}
In R2--R5 one has $k=v_f=v_X=0$. In R6 one has $o\in\{\frac12,\frac32,\frac52\}$ and $k+v_f+v_X\le1$, with equality only if $o=\frac12$ (and then $\omega=0$).
\item[(b)] The degree formula holds for every connected graph built from these vertices. A graph of $\rho_t$ or $\mathcal P_t$ has $\omega=1-o-2(k+v_f)$. A zero-momentum self-energy graph with $k+v_f\ge1$ (a graph of $f_t$ or of $\delta\dot m^2_t$) has $\omega=2-o-2(k+v_f)\le-1$.
\end{enumerate}
\end{lemma}
\begin{proof}
Each line has degree $-2$. The vertex function of $V_n$ is $e^n\eps^{n-2}$ times a trigonometric polynomial in $\eps$ times the momenta, so its degree is at most $2-n$. The $\lambda$-vertex and every two-point vertex have degree $0$. A $\hat\Delta$-line is two lines with the same momentum joined by a two-point vertex. Reisz's degree \cite[(2-28)]{Rei88b}, with $3$ in place of $4$, is $\omega=3L-2I+\sum_v\deg(v)$. With $L=I-V+1$ and $2I+E=\sum_v\mathrm{legs}(v)$ this gives the displayed formula: the vertex contributions $\frac12\mathrm{legs}(v)-3+\deg(v)$ are $-\frac n2$, $-1$ and $-2$, respectively.

Every vertex has $2$ or $4$ scalar legs, so $E_\varphi$ is even. Every photon line has two photon ends, so $\sum_nnv_n\equiv E_A\pmod2$; hence $o\in\frac12+\N$ if and only if $E_A$ is odd. With $o=0$ only two-point vertices occur, and they form no loop with external legs. A single $V_1$ has three legs. So a loop with $E\ge1$ has $o\ge1$, or $o=\frac12$ with $E=1$.
$E=0$: $\gamma$ is a whole graph. This gives R0 or R1, since $\omega=1-o-2(k+v_f)$ when $v_X=1$.
$E=1$: then $E_A=1$, which gives R6. There $o\in\frac12+\N$ and $\omega=\frac52-o-2(k+v_f+v_X)\ge0$, which gives the bounds stated for R6.
$E=2$: $(E_\varphi,E_A)\in\{(2,0),(0,2)\}$, so $o$ is an integer $\ge1$. Then $\omega=2-o-2(k+v_f+v_X)\ge0$ forces $k=v_f=v_X=0$ and $o\in\{1,2\}$. This gives R2 and R3.
$E=3$: $(E_\varphi,E_A)\in\{(2,1),(0,3)\}$, so $o\in\{\frac12,\frac32,\dots\}$. A loop needs $o\ge\frac32$, and $\omega=\frac32-o-2(k+v_f+v_X)\ge0$ forces $o=\frac32$ and no two-point vertices. This gives R4 and R5.
$E=4$: $\omega\ge0$ forces $o\le1$ and no two-point vertices. The vertex sets $\{V_1,V_1\}$, $\{V_2\}$, $\{\lambda\}$ with four external legs are trees.
$E\ge5$: $\omega\le\frac12-o<0$ for $o\ge1$, and $o=\frac12$ gives a tree.
(b) A graph of $\rho_t$ or $\mathcal P_t$ is a vacuum graph with one insertion. A zero-momentum self-energy has $E_\varphi=2$, $E_A=0$ and $v_X=0$, and it has integer $o\ge1$ by the above.
\end{proof}

\begin{lemma}[zero-momentum counterterms as forests]\label{app:G1:WP-q2a}
Work in the $C_t$-model on $T$ or on $\eps\Z^3$. Let $J_0:=\begin{pmatrix}0&-1\\1&0\end{pmatrix}$ act on the two real components of $\varphi$, and let $U(A):=\mathrm{e}^{e\eps AJ_0}$. The vertices are of the following kinds:
\begin{itemize}
\item King's vertices $V_n(A,\varphi)=-\frac{e^n\eps^{n-2}}{n!}\sum_{x,\mu}\eps^3A_{x,\mu}^n\,\rRe[i^n\bar\varphi(x)\varphi(x+\eps e_\mu)]$, $n\ge1$, of order $o(V_n):=\frac n2$, and the vertex $\lambda\sum_x\eps^3\abs\varphi^4$, of order $1$;
\item two-legged counterterm vertices $c_o$, $o\in\{1,2\}$, with coupling $-\frac12\delta m^{2,(o)}_0$, where $\delta m^{2,(o)}_0$ is the sum of the $e^\alpha\lambda^\beta$ terms of $\delta m^2_0$ with $\frac\alpha2+\beta=o$;
\item two-legged $f$-vertices;
\item an insertion vertex $X$. For $\rho_t$, $X$ is $r_0^2$. For $\mathcal P_t$, $X$ is the pair $\hat G_{\underline M}\hat G_{\overline M}$, joined to the rest by two scalar legs of equal momentum.
\end{itemize}
For a finite multiset $\nu$ of kinds put $\nu!:=\prod_\kappa\nu_\kappa!$. A graph of content $\nu$ consists of two pieces of data:
\begin{itemize}
\item a perfect matching of the legs of the labelled vertices $(\kappa,k)$, $k\le\nu_\kappa$;
\item scalar lines coloured $\hat G_{\overline M}$ or $\hat\Delta$, and photon lines $C_{\mu_*}$.
\end{itemize}
$A(G)=\int I_G$ is the product of the couplings times the lattice sum, resp.\ Brillouin-zone integral, of its line and vertex factors. Here $I_G$ is the integrand, and $\int$ is $\abs T^{-L}\sum_{(T^*)^L}$ on $T$ and $\int_{{\rm BZ}^L}$ on $\eps\Z^3$.

A \emph{block} of $G$ is a set $S$ of vertices of the kinds $V_n$ and $\lambda$ with the following properties:
\begin{enumerate}
\item every line joining two vertices of $S$ is a $\hat G_{\overline M}$ or $C_{\mu_*}$ line;
\item the induced graph $G[S]$ is one-particle irreducible;
\item exactly two legs at $S$ are not matched inside $S$, and both are scalar;
\item $o(S):=\sum_{v\in S}o(v)\le2$.
\end{enumerate}
A \emph{forest} is a set of blocks any two of which are nested or disjoint. For a forest $F$ put $A_F(G):=\int\prod_{\gamma\in F}(-t^0_\gamma)I_G$, where $t^0_\gamma$ sets the external momentum of $\gamma$ to $0$ and the operators of nested elements act from the inside out. For a two-point graph $\gamma$, $\hat A_F(\gamma)(0)$ is the $2\times2$ matrix obtained in the same way at external momentum $0$. Let $\mathcal C$ be one of the following classes:
\begin{itemize}
\item connected graphs containing $X$ once (for $\rho_t$, $\mathcal P_t$);
\item one-particle irreducible amputated two-point graphs at zero momentum with at least one $\hat\Delta$-line or $f$-vertex (for $f_t$).
\end{itemize}
Let $\mathcal C_0\subset\mathcal C$ be the graphs without $c_o$-vertices. Then, for every order,
\begin{equation}\label{eq:q2a}
\sum_\mu\frac1{\mu!}\sum_{H\in\mathcal C(\mu)}A(H)=\sum_\nu\frac1{\nu!}\sum_{G\in\mathcal C_0(\nu)}\sum_{F}A_F(G),
\end{equation}
where $F$ runs over the forests of $G$.
\end{lemma}

\begin{proof}
\emph{Step 1 (the counterterm is scalar).} Let $\Gamma^{(o)}$ be King's order-$o$ amputated one-particle-irreducible two-point function of the uniform-mass model, excluding the single counterterm vertex. Its zero-momentum value $\hat\Gamma^{(o)}(0)$ is a real $2\times2$ matrix. The action and the Gaussian are invariant under two maps:
\begin{itemize}
\item $\varphi\mapsto\mathrm{e}^{\theta J_0}\varphi$;
\item $\varphi\mapsto P\varphi$, $A\mapsto-A$, with $P=\mathrm{diag}(1,-1)$, since $PU(A)P=U(-A)$.
\end{itemize}
So $\hat\Gamma^{(o)}(0)$ commutes with $\mathrm{e}^{\theta J_0}$, hence equals $a\mathbf 1+bJ_0$. It also commutes with $P$, hence $b=0$. By \cite[(3.33)]{K1}, $\hat\Gamma^{(o)}(0)=\delta m^{2,(o)}_0\mathbf 1$.

\emph{Step 2 (factorisation).} Let $S$ be a block. Its kernel $K_S(z,z')$ between the sites of its two free legs is translation invariant. Hence $(t^0_SI_G)$ is the integrand of $G/S$, which has at the place of $S$ a two-legged vertex of value $\hat K_S(0)=\sum_{z'}\eps^3K_S(z,z')$. After Reisz's routing $S_\gamma$ \cite[(2-23)]{Rei88b}, the internal momenta of $S$ are independent variables, and the integrand is bounded on the compact zone. By Fubini, $\int(-t^0_S)I_G=-\hat K_S(0)\cdot A(G/S)$. The same holds recursively for nested blocks. On $T$ the same argument works with sums over $T^*\ni0$.

\emph{Step 3 (substitution).} Let $H\in\mathcal C(\mu)$ have counterterm vertices $v_1,\dots,v_r$, with $r_o$ of them of order $o$. For each $v_i$ choose a counterterm graph $\gamma_i$ and a forest $F_i$ of proper blocks of $\gamma_i$. Here a counterterm graph is a one-particle-irreducible uniform-mass two-point graph of content $\kappa_i$ and order $o(v_i)$ without $c$-vertices. Delete $v_i$ and match the free legs $x,y$ of $\gamma_i$ with the partners of legs $1,2$ of $v_i$. Then distribute the labels: there are $\prod_\kappa\nu_\kappa!/(\mu_\kappa!\prod_i\kappa_{i,\kappa}!)$ ways, with $\nu=\mu_K+\sum_i\kappa_i$. Each way gives a pair $(G,F)$ with $G\in\mathcal C_0(\nu)$ and $F=\{S_i\}\cup\bigcup_iF_i$, where $S_i$ is the vertex set of $\gamma_i$. Indeed, $G[S_i]=\gamma_i$. Contracting a connected one-particle-irreducible part preserves connectedness, one-particle irreducibility and the presence of $X$, $\hat\Delta$ or $f$. And the $S_i$ are the maximal elements of $F$. Conversely, contracting the maximal elements of a forest recovers $(H,(\gamma_i,F_i))$. This uses three facts:
\begin{itemize}
\item $o(S_i)\in\{1,2\}$, because $E_A=0$ forces the order to be an integer;
\item a proper sub-block has strictly smaller order;
\item the blocks of $G$ inside $S_i$ are exactly the blocks of $\gamma_i$.
\end{itemize}
This reconstruction has $2^r\prod_or_o!$ preimages, coming from the names $x,y$ of the free legs of each block and from the labels of the $c_o$-vertices. By Step 2, replacing the coupling $-\frac12\delta m^{2,(o_i)}_0\mathbf 1$ of $v_i$ by $-\frac12\hat A_{F_i}(\gamma_i)(0)$ turns $A(H)$ into $2^{-r}A_F(G)$. Since $\mu!=\prod_\kappa\mu_\kappa!\prod_or_o!$,
\[
\frac1{\mu!}\prod_i\frac1{\kappa_i!}\cdot2^{-r}\cdot\frac{2^r\prod_or_o!}{\prod_\kappa\nu_\kappa!/(\mu_\kappa!\prod_i\kappa_{i,\kappa}!)}=\frac1{\nu!}.
\]

\emph{Step 4 (induction on $o$).} Claim: $\delta m^{2,(o)}_0\mathbf 1=\sum_\kappa\frac1{\kappa!}\sum_{\gamma}\sum_{F\ \text{proper}}\hat A_F(\gamma)(0)$, where the sum runs over the counterterm graphs $\gamma$ of order $o$. For $o=1$ the graphs of $\Gamma^{(1)}$ contain no counterterm vertex and no proper block. For $o=2$ the graphs of $\Gamma^{(2)}$ contain only vertices $c_1$. Insert the case $o=1$ and apply Step 3 to the class of $\Gamma^{(2)}$; Step 1 converts the matrix sum into $\delta m^{2,(2)}_0\mathbf 1$. Insert the claim into every $c$-vertex of every $H\in\mathcal C(\mu)$ and apply Step 3 once more. This gives \eqref{eq:q2a}.

\emph{Step 5 ($f_t$).} By \cite[(3.33)]{K1} with lines $C_t$, $\delta m^{2,(o)}_t=\hat\Gamma^{(o)}_t(0)$. Expand $C_t=\hat G_{\overline M}+t\hat\Delta$ multilinearly and $\delta m^2_t=\delta m^2_0+f_t$ at lower orders. The terms with only $\hat G_{\overline M}$-lines and only $\delta m^2_0$-vertices form the right side of King's recursion at $\overline M$, which is $\delta m^{2,(o)}_0$. So $f^{(o)}_t$ is the sum over the second class $\mathcal C$, to which \eqref{eq:q2a} applies.
\end{proof}

\begin{remark}
For a scalar two-point part with no two-point vertices, $\omega=2-o$ (Lemma~\ref{app:G1:CN-list}). Hence the blocks are exactly the proper one-particle-irreducible scalar two-point parts with $\omega\ge0$ that contain no $\hat\Delta$-line, no $f$-vertex and not the insertion; these are the R2 parts. A subdiagram with two external scalar legs whose vertex set is $S$ contains every line of $G[S]$, since an omitted line adds two external legs. So the forests above are Reisz's forests \cite[p.~87]{Rei88b} restricted to these parts, with $t^0$ (Lemma~\ref{app:G1:WP-q2a-disj}). The right side of \eqref{eq:q2a} is $\sum_\nu\frac1{\nu!}\sum_G\int R^{(a)}_G$, where $R^{(a)}_G$ is the forest integrand \cite[(4-2)]{Rei88b}.
\end{remark}

\begin{lemma}[line-disjoint blocks are vertex-disjoint]\label{app:G1:WP-q2a-disj}
In the setting of Lemma~\ref{app:G1:WP-q2a}, let $\gamma_1,\gamma_2$ be one-particle-irreducible subgraphs of $G$ with the following properties: each has exactly two external legs, both scalar; neither contains the insertion; and they have no line in common. Then $\gamma_1$ and $\gamma_2$ have no vertex in common. In particular, the forests of Lemma~\ref{app:G1:WP-q2a} (nested or vertex-disjoint blocks) are the same as the forests of two-point parts in the sense of \cite[p.~87]{Rei88b} (nested or line-disjoint), restricted to the blocks.
\end{lemma}
\begin{proof}
Suppose $\gamma_1,\gamma_2$ have a common vertex $v$, of degree $d$. Let $a_i$ be the number of legs at $v$ that lie on lines of $\gamma_i$. Since $\gamma_i$ is one-particle irreducible, $a_i\ge2$. The $d-a_i$ remaining legs at $v$ are external to $\gamma_i$, so $d-a_i\le2$. Line-disjointness gives $a_1+a_2\le d$. Hence $d=4$ and $a_1=a_2=2$. Then every external leg of $\gamma_1$ lies on a line of $\gamma_2$, and conversely. So $\gamma_1\cup\gamma_2$ is a connected component without external legs, and therefore it is all of $G$. This is impossible: every graph considered has either two external legs or the insertion, and the insertion lies in neither $\gamma_i$.
\end{proof}

\begin{proof}[Proof of Lemma~\ref{app:G1:WP-q2}]
The degrees are $\omega(\gamma)=3-\frac12(E_\varphi+E_A)-o_V-2v_2$, where $v_2=k+v_f+v_X$ is the number of two-point vertices (Lemma~\ref{app:G1:CN-list}). Lemma~\ref{app:G1:CN-list} lists the parts with $\omega\ge0$ in the classes R0--R6. Those that occur as proper parts are:
\begin{itemize}
\item R2, the scalar self-energies;
\item R3, the two-photon parts of order $o\in\{1,2\}$;
\item the types R4 and R5, whose zero-order forest terms are odd at zero external momentum;
\item R6, the parts with one photon leg.
\end{itemize}
None of the parts of R2--R5 contains a $\hat\Delta$-line, an $f$-vertex or the insertion. An R6 part of order $\frac12$ may contain one of them, but graphs with an R6 part vanish identically, see (d).
(a) \emph{$\delta m^2_0$ and the forests of the R2 parts.} King's recursion \cite[(3.33)]{K1} sets $\delta m^2(\alpha,\beta)$ equal to the zero-momentum value of the full 1PI two-point function at that order. Expand it in the lines and in $\delta m^2_t=\delta m^2_0+f_t$. The terms with only $\hat G_{\overline M}$ lines and only $\delta m^2_0$ insertions reproduce King's counterterm at mass $\overline M$, which is $\delta m^2_0$. All other terms (those with at least one $\hat\Delta$ line or one $f$-insertion) make up $f_t$. By Lemma~\ref{app:G1:WP-q2a}, inserting the vertices $\delta m^2_0$ is the same as summing over the $t^0$-forests of the scalar self-energy parts that contain no $\hat\Delta$-line, no $f$-vertex and not the insertion and have order at most $2$. These forests are Reisz's forests of these parts (Lemma~\ref{app:G1:WP-q2a-disj} and the remark after Lemma~\ref{app:G1:WP-q2a}). Self-energy parts that contain the insertion, or have order $o\ge3$, are not subtracted by King's counterterm. The parts that contain a $\hat\Delta$ line or an $f$-vertex are not subtracted either: their counterterm is the explicit $f$-vertex. All these parts have $\omega\le-1$ (Lemma~\ref{app:G1:CN-list}).
(b) \emph{$t^0\to t^1$ for the R2 parts with $\omega=1$.} The added term $\eps^{-1}\sin(\eps q_\mu)\partial_{q_\mu}R_\gamma|_{q=0}$ integrates to zero over the internal momenta of $\gamma$, which are independent integration variables after $S_\gamma$ \cite[(2-23)]{Rei88b}. Fubini applies because the integrand is bounded on the compact BZ: all denominators are at least $\min(\underline M^2,\mu_*^2)$.
(c) \emph{R3--R5.} For R4 and R5 the forest terms vanish graph by graph, by oddness at $q=0$. For the two-photon parts R3, consider the zero-order terms summed over all R3 graphs of a given order at the same place. They add up to $\Pi^{(o)}_{\overline M}(0)\times(\text{reduced integrand})$, which is $0$ by the Ward identity on $\eps\Z^3$, Corollary~\ref{app:K:L16-2}. This needs the R3 parts to be free of $\hat\Delta$ and $f$ (uniform mass $\overline M$), which is why one expands in $\hat\Delta$ rather than in the two masses. For periodicity, $A$ is labelled by the lower end point. In this labelling the linear Taylor term of an R3 part does \emph{not} vanish graph by graph: it acquires $(i\eps/2)(q_\mu-q_\nu)\Pi^{\rm mid}_\gamma(0)$. It vanishes only after the same per-place regrouping, since $\Pi^{\rm low}(q)=\mathrm{e}^{i\eps(q_\mu-q_\nu)/2}\Pi^{\rm mid}(q)$ with $\Pi^{\rm mid}$ even and $\Pi^{\rm mid}(0)=0$.
(d) \emph{R6.} Graphs with a one-photon part vanish identically.
(e) \emph{Completeness.} Take a 1PI $\gamma\subsetneq G$ with $\omega(\gamma)\ge0$. By Lemma~\ref{app:G1:CN-list}, it lies in R2--R6, and unless it lies in R6 it contains no $\hat\Delta$-line, no $f$-vertex and not the insertion. It is subtracted with $\delta=\omega$: by (a)--(b) for R2, by (c) for R3--R5, and by (d) for R6. Every other $\gamma$ has $\omega<0$ and $t=0$. Reisz's conditions (4-3)--(4-4) with $\delta=\omega$, and $t:=0$ when $\omega<0$ (including for the whole graph), are ``automatically satisfied'' \cite[p.~92]{Rei88b}.
\end{proof}

\paragraph{The weighted-line form of Ba{\l}aban's Proposition~1.}

In scale-$k$ units, a $\Gamma^B$-line (``$D$-line'') is $u\,c_kD^\eta_k$ with $D^\eta_k:=\sum_{j<k}(L^j\eta)^2G^\eta_{(j)}$, because $\delta_j=c_k(L^j\eta)^2$. At $k=K$ we have $\kappa\mu_*^2D=c_KD^\eta_K$ exactly, with no coarse remainder. Expanding every vector line $G_k(u)=G_k+uc_kD^\eta_k$ multilinearly, the value of a graph of the $u$-model is $\sum_{\rm colourings}(uc_k)^{m}\times$(value of the coloured graph). Here every internal vector line is coloured $C$ (kernel $G_k$) or $D$, and $m$ is the number of $D$-lines. Write $D^\natural$ for Ba{\l}aban's degree \cite[(2.1)--(2.3), pp.~422--423]{Ba3} with each $D$-leg of dimension $(-\mathsf d+4)/2$; a subgraph is \emph{divergent} if $D^\natural\le0$.
\emph{Organisation $(\alpha)$:} the mass counterterm is the fixed $\delta m^2_{\rm ref}$ \cite[(1.23), (1.26)--(1.29)]{Ba3}, built from $C$-lines only, and only $C$-only divergent subgraphs are renormalised (by $\delta m^2_{\rm ref}$, Ward--Takahashi partners, and $\varphi\varphi A$ partners).

\begin{lemma}[$D$-lines and degrees]\label{app:G1:WP-deg}
Let $\mathsf d=3$.
(a) $D^\eta_k$ has the representation \cite[(2.6)]{Ba3}. Its slice kernels $(L^j\eta)^2G^\eta_{(j)}$ obey \cite[(2.10)--(2.12)]{Ba3} with $(L^j\eta)^{-\mathsf d+2}$ replaced by $(L^j\eta)^{-\mathsf d+4}$ and with the same exponential factors, and $D^\eta_k$ obeys \cite[(3.1)]{Ba3}. Hence $D^\natural(G)=D(G)+2m(G)$, where $D$ is Ba{\l}aban's degree with the colours forgotten.
(b) Let $G$ be a connected graph without the vertices (1.9), (1.11), (1.13)--(1.15) and without finite counterterms. Then $D(G)=I_v+E'_v+\sum_vn'_v+d_s+\frac12E_s+E_\partial-3$. Here $I_v$ is the number of internal vector lines, $E'_v$ the number of external $A'$ legs, $n'_v$ the number of background legs, $d_s$ the number of $\lambda$-vertices, $E_s$ the number of external scalar legs, and $E_\partial$ the number of differentiations on external legs.
(c) If $G$ has an external leg and $m\ge1$, then $D^\natural(G)\ge1$. For a vacuum graph, $D^\natural=I_v+d_s-3+2m$, plus the non-negative terms of (b) when (b) applies. So $m=1$ gives $D^\natural\ge0$, with equality only at $(\alpha,\beta)=(2,0)$, and $m\ge2$ gives $D^\natural\ge3$.
(d) Every divergent proper subgraph of a coloured graph is $C$-only and belongs to one of Ba{\l}aban's classes \cite[pp.~429--432]{Ba3}. A vacuum graph with $m\ge2$, or with $m=1$ and $o\ge2$, is convergent.
\end{lemma}
\noindent The lemma uses the bounds \cite[(2.10)--(2.12), p.~426]{Ba3} for the vector slices. The vector lines do not depend on the background. At zero background, these bounds for King's slices $H_j=G_{j+1}-G_j$ are given by Corollary~\ref{app:G1:VSR-VSphys}. By Lemma~\ref{app:G1:WP-Kslices} below, King's slices are Ba{\l}aban's slices $G^\eta_{(j)}$ of \cite[(2.6)]{Ba3}, including $j=0$: King's $H_0=G_1$ and Ba{\l}aban's $G_{(0)}=C^{(0),\eta}$.
\begin{proof}
(a) Multiply \cite[(2.10)--(2.12)]{Ba3} by $(L^j\eta)^2\le1$. Both legs of a line carry the same index $j$, so splitting $(L^j\eta)^{-\mathsf d+4}$ as $(L^j\eta)^{(-\mathsf d+4)/2}$ per leg is legitimate in \cite[(2.14)]{Ba3}. For localisations at distance at least $1$, $\sum_j(L^j\eta)^{2-\mathsf d+2-1-\alpha}\mathrm{e}^{-\delta_1\dist/(L^j\eta)}\le O(1)\mathrm{e}^{-\delta_0\dist}$, because each term is at most $O(1)\mathrm{e}^{-\delta_1\dist/2}(L^j\eta)^{-\alpha}\mathrm{e}^{-\delta_1/(2L^j\eta)}$, which is summable.
(b) Per vertex, $D_G(v)$ is the number of $\eta$-factors, plus the dimensions of the legs on internal lines, minus the differentiations on internal lines \cite[p.~423]{Ba3}. The $\eta$-counts are: $\mathsf d$ for (1.6) and (1.7); $\mathsf d+n+n'-1$ for (1.8) (one differentiation, on a scalar leg); $\mathsf d+n+n'-2$ for (1.10) (no differentiation). Sum over vertices using $\sum_vn_v=2I_v+E'_v$ and $\mathsf d=3$. For a counterterm vertex, \cite[(2.3)]{Ba3} gives $D(G)=D(G')+E_\partial(G_0)\ge D(G')$, with equality for the counterterms of \cite[(1.23)]{Ba3}. In full generality the formula carries the extra terms $+2V_7+\sum\deg(\mathrm{ct})$, both non-negative.
(c) If $G$ contains (1.13)--(1.15), an R-vertex or a finite counterterm, then $D>0$ by \cite[Cor.~2.3, p.~429]{Ba3}, so $D^\natural>2$. Otherwise $m\ge1$ gives $I_v\ge1$, and an external leg means $E'_v\ge1$, or $E_s\ge2$ ($E_s$ is even), or $n'\ge1$. So $D\ge-1$ and $D^\natural\ge1$. For vacuum graphs, $E'_v=E_s=E_\partial=n'=0$ and $I_v\ge m$.
(d) A proper subgraph of a connected graph has an external leg, so use (c). The $(2,0)$ one-$D$-line vacuum graphs (bubble, seagull, and the $(1.8)$, $n=2$ self-contractions) are the only ones with $D^\natural=0$.
Ba{\l}aban's graph (2.4) is the tree with one internal \emph{scalar} line ($I_v=0$, $E'_v=2$, $E_s=2$, $D=0$). There are no ``coloured (2.4) graphs''. A (2.4) subgraph whose external vector legs are $D$-lines is integrated by parts \cite[(2.8)]{Ba3} using the derivative bound of (a).
\end{proof}

\emph{Block lattices.} Let $T$ and $\Lambda=T_\eps$ be as in Definition~\ref{app:G1:DOM-set}, with $\ip fg=\eps^3\sum_{x\in\Lambda}f(x)g(x)$ and with the blocks $B^k(y)$ of Definition~\ref{app:G1:DOM-set}(a). For $0\le k\le K$ let $\Lambda^{(k)}:=T\cap t_k\Z^3$, and define
\[
\ip{B}{B'}_k:=t_k^3\sum_{y\in\Lambda^{(k)}}B(y)B'(y),\qquad \norm B_k^2:=\ip BB_k .
\]
Define $Q_k:\ell^2(\Lambda)\to\ell^2(\Lambda^{(k)})$ by $(Q_kA)(y):=L^{-3k}\sum_{x\in B^k(y)}A(x)$ \cite[(2.10), p.~653]{K1}. For $k<K$ define $Q:\ell^2(\Lambda^{(k)})\to\ell^2(\Lambda^{(k+1)})$ by $(QB)(z):=L^{-3}\sum_{y\in\Lambda^{(k)}\cap B^{k+1}(z)}B(y)$. Adjoints are taken with respect to these inner products.

\begin{lemma}[block averages]\label{app:G1:blockops}
\begin{enumerate}
\item[(a)] For $0\le k\le K$: $Q_k^*$ is the extension by constants on the blocks $B^k(y)$; $Q_kQ_k^*=1$; $Q_k^*Q_k=P_k$; and $\norm B_k=\norm{Q_k^*B}$.
\item[(b)] For $0\le k<K$: $Q^*$ is the extension by constants from $\Lambda^{(k+1)}$ to $\Lambda^{(k)}$; $QQ^*=1$; $Q^*Q$ is an orthogonal projection; $\norm Q\le1$; and $QQ_k=Q_{k+1}$.
\item[(c)] For $0\le k<K$: $Q_k^*(1-Q^*Q)Q_k=P_k-P_{k+1}$.
\end{enumerate}
\end{lemma}
\begin{proof}
By (F2) of Definition~\ref{app:G1:DOM-set}, for $k\le K$ the blocks $B^k(y)$, $y\in\Lambda^{(k)}$, partition $\Lambda$ into cubes of $L^{3k}$ sites. Each $B^{k+1}(z)$ is the disjoint union of the $L^3$ blocks $B^k(y)$ with $y\in\Lambda^{(k)}\cap B^{k+1}(z)$.

(a) Let $y(x)$ be the corner of the block that contains $x$. Since $t_k^3L^{-3k}=\eps^3$,
\[
\ip{Q_kA}{B}_k=\eps^3\sum_xA(x)B(y(x)),
\]
so $(Q_k^*B)(x)=B(y(x))$. Averaging a block-constant function returns its value, so $Q_kQ_k^*=1$. The operator $Q_k^*Q_k$ replaces $A$ on each block by its block average, which is the orthogonal projection $P_k$. Finally, $\norm{Q_k^*B}^2=\ip{B}{Q_kQ_k^*B}_k=\norm B_k^2$.

(b) The same computation, with $t_{k+1}^3L^{-3}=t_k^3$, gives the first two claims. Hence $(Q^*Q)^2=Q^*Q=(Q^*Q)^*$, and $\norm Q^2=\norm{Q^*Q}\le1$. The blocks are nested, so the average over $B^{k+1}(z)$ is the average of the $L^3$ averages over the blocks $B^k(y)\subset B^{k+1}(z)$. Thus $QQ_k=Q_{k+1}$.

(c) Use (a) at $k$ and at $k+1\le K$, together with (b):
\[
Q_k^*Q_k-Q_k^*Q^*QQ_k=P_k-(QQ_k)^*(QQ_k)=P_k-P_{k+1}. \qedhere
\]
\end{proof}

Let $M=-\Delta_\eps+\mu_*^2$, $r_k=a_kt_k^{-2}$ and $G_k=(M+r_kP_k)^{-1}$ be as in Definition~\ref{app:G1:DOM-set}. For $1\le k\le K$, King's effective block operator \cite[(2.14), p.~653]{K1} and his lift \cite[(2.18), p.~654]{K1} are
\[
\Delta^{(k)}:=r_k-r_k^2\,Q_kG_kQ_k^*\ \text{ on }\ell^2(\Lambda^{(k)}),\qquad F_k:=r_kG_kQ_k^* .
\]
For $1\le k<K$, the fluctuation covariance is $C^{(k)}:=(\Delta^{(k)}+\beta_kQ^*Q)^{-1}$ with $\beta_k:=a\,t_{k+1}^{-2}$ \cite[(2.16), p.~653]{K1}. The constant here is $a$, not $a_{k+1}$.

\begin{lemma}[variational form of $\Delta^{(k)}$]\label{app:G1:Dk-var}
Let $1\le k\le K$ and $B\in\ell^2(\Lambda^{(k)})$, and put $J_B(A):=\ip{A}{MA}+r_k\norm{Q_kA-B}_k^2$ for $A\in\ell^2(\Lambda)$. Then, for all $A$,
\[
J_B(A)=\ip{A-F_kB}{G_k^{-1}(A-F_kB)}+\ip{B}{\Delta^{(k)}B}_k .
\]
Hence $F_kB$ is the unique minimiser of $J_B$, $\min_AJ_B(A)=\ip B{\Delta^{(k)}B}_k$, and $0<\Delta^{(k)}\le r_k$. In particular, for $k<K$ and every $\beta\ge0$ the operator $\Delta^{(k)}+\beta Q^*Q$ is positive definite, so $C^{(k)}$ is well defined.
\end{lemma}
\begin{proof}
By Lemma~\ref{app:G1:blockops}(a), $\norm{Q_kA-B}_k^2=\ip{A}{P_kA}-2\ip{A}{Q_k^*B}+\norm B_k^2$, so $J_B(A)=\ip{A}{G_k^{-1}A}-2r_k\ip{A}{Q_k^*B}+r_k\norm B_k^2$. Let $A_0:=F_kB$, so that $G_k^{-1}A_0=r_kQ_k^*B$. Then
\[
\ip{A-A_0}{G_k^{-1}(A-A_0)}=\ip{A}{G_k^{-1}A}-2r_k\ip{A}{Q_k^*B}+r_k^2\ip{B}{Q_kG_kQ_k^*B}_k .
\]
Subtracting, and using $r_k\norm B_k^2-r_k^2\ip{B}{Q_kG_kQ_k^*B}_k=\ip{B}{\Delta^{(k)}B}_k$, gives the identity. Since $G_k^{-1}\ge M\ge\mu_*^2>0$, the first term is positive unless $A=A_0$. So $A_0$ is the unique minimiser, and $\ip B{\Delta^{(k)}B}_k=J_B(A_0)$.

Taking $A=0$ gives $\ip B{\Delta^{(k)}B}_k\le J_B(0)=r_k\norm B_k^2$. Suppose $\ip B{\Delta^{(k)}B}_k=0$. Then $\ip{A_0}{MA_0}=0$, so $A_0=0$, and therefore $r_k\norm B_k^2=J_B(0)=J_B(A_0)=0$, that is, $B=0$. Hence $\Delta^{(k)}>0$. The last claim follows because $Q^*Q\ge0$.
\end{proof}

\begin{remark}\label{app:G1:Dk-var-rem}
In King's notation, the lemma says that the average field $A^{(k)}=F_kA_k$ of \cite[(2.18), p.~654]{K1} minimises the fine action under the soft constraint $Q_kA\approx A_k$. The same completion of the square, carried out one scale up and in the block field, is Ba{\l}aban's translation \cite[(3.41), p.~619]{Ba1}.
\end{remark}

\begin{lemma}[King's slices satisfy Ba{\l}aban's recursion]\label{app:G1:WP-Kslices}
Let $a_k=a(1-L^{-2})(1-L^{-2k})^{-1}$ be King's constants \cite[(2.13), p.~653]{K1}. Let $H_j=G_{j+1}-G_j$, $0\le j<K$, with $G_0:=0$, be King's slices (Subsection~\ref{ssec:T6-statement}).
\begin{enumerate}
\item[(a)] $r_{j+1}^{-1}=r_j^{-1}+\beta_j^{-1}$ for $1\le j<K$.
\item[(b)] $H_j=F_jC^{(j)}F_j^*=r_j^2\,G_jQ_j^*C^{(j)}Q_jG_j$ for $1\le j<K$.
\item[(c)] $H_0=G_1=(M+a\,t_1^{-2}P_1)^{-1}$.
\end{enumerate}
\end{lemma}
\begin{proof}
(a) Since $t_j=L^{-1}t_{j+1}$,
\[
\frac{t_j^2}{a_j}+\frac{t_{j+1}^2}{a}=\frac{t_{j+1}^2}{a(1-L^{-2})}\big[L^{-2}(1-L^{-2j})+1-L^{-2}\big]=\frac{t_{j+1}^2}{a_{j+1}}.
\]

(b) On $\ell^2(\Lambda)\oplus\ell^2(\Lambda^{(j)})$, with the inner products $\ip\cdot\cdot$ and $\ip\cdot\cdot_j$, consider the quadratic form
\[
S(A,B):=\ip{A}{MA}+r_j\norm{Q_jA-B}_j^2+\beta_j\norm{QB}_{j+1}^2 .
\]
First, by Lemma~\ref{app:G1:Dk-var} and $\norm{QB}_{j+1}^2=\ip{B}{Q^*QB}_j$,
\[
S(A,B)=\ip{A-F_jB}{G_j^{-1}(A-F_jB)}+\ip{B}{(C^{(j)})^{-1}B}_j .
\]
Thus $S=\Theta^*\operatorname{diag}(G_j^{-1},(C^{(j)})^{-1})\Theta$, where $\Theta(A,B):=(A-F_jB,B)$ is invertible with $\Theta^{-1}(A',B)=(A'+F_jB,B)$. So $S$ is positive definite, and the $AA$-block of $S^{-1}=\Theta^{-1}\operatorname{diag}(G_j,C^{(j)})\Theta^{-*}$ is $G_j+F_jC^{(j)}F_j^*$.

Second, fix $A$ and put $c:=Q_jA$. Write $B=Q^*z+w$ with $z:=QB$ and $w:=(1-Q^*Q)B$. Since $QQ^*=1$ (Lemma~\ref{app:G1:blockops}(b)), $\norm{Q^*v}_j=\norm v_{j+1}$, so
\[
\norm{c-B}_j^2=\norm{Qc-z}_{j+1}^2+\norm{(1-Q^*Q)c-w}_j^2 .
\]
Hence
\[
\min_B\big[r_j\norm{c-B}_j^2+\beta_j\norm{QB}_{j+1}^2\big]=\min_z\big[r_j\norm{Qc-z}^2_{j+1}+\beta_j\norm z^2_{j+1}\big]=\frac{r_j\beta_j}{r_j+\beta_j}\norm{Qc}_{j+1}^2 .
\]
By (a), $r_j\beta_j/(r_j+\beta_j)=r_{j+1}$. By Lemma~\ref{app:G1:blockops}(a),(b), $\norm{Qc}_{j+1}^2=\norm{Q_{j+1}A}_{j+1}^2=\ip{A}{P_{j+1}A}$. So $\min_BS(A,B)=\ip{A}{G_{j+1}^{-1}A}$.

For a positive definite block form, the minimum over the second variable is the Schur complement, whose inverse is the $AA$-block of $S^{-1}$. Hence $G_{j+1}=G_j+F_jC^{(j)}F_j^*$. Finally, $F_j^*=r_jQ_jG_j$.

(c) $G_0=0$ gives $H_0=G_1=(M+r_1P_1)^{-1}$, and $a_1=a$.
\end{proof}
\noindent Part (b) is the recursion \cite[(2.42)]{Ba1} at zero background, written in King's notation, and part (c) is the identity $G_1=C^{(0)}$ that precedes \cite[(2.43)]{Ba1}. Summing (b) and (c) over $j<k$ gives
\[
G_k=G_1+\sum_{j=1}^{k-1}a_j^2t_j^{-4}G_jQ_j^*C^{(j)}Q_jG_j .
\]
In scale-$k$ units ($t_j/t_k=L^j\eta$) this is the decomposition \cite[(2.6)]{Ba3}$=$\cite[(2.43)]{Ba1} at $\tilde B=0$. Thus, when the zero-background operators $G_j,Q_j,C^{(j)}$ and the constants $a_j$ of \cite[Sect.~2]{Ba1} are written in King's normalisation, the $H_j$ are the slices $G^\eta_{(j)}$ of \cite[(2.6)]{Ba3} for $0\le j<k$; for $j=0$, $H_0=G_1=C^{(0)}$. The bounds use only the operators $H_j$, through Corollary~\ref{app:G1:VSR-VSphys}, which is proved for them directly.

\begin{lemma}[families]\label{app:G1:WP-fam}
Call two coloured graphs \emph{partners} if one arises from the other by one of Ba{\l}aban's partner operations applied to a $C$-only subgraph, with the colours of all other lines unchanged. The operations are: (i) replace a $C$-only divergent scalar self-energy $G_0$ by the vertex (1.7) with $\delta m^2_{G_0}$, or back; (ii) exchange $C$-only two-vector-leg graphs of the same order \cite[p.~431]{Ba3}; (iii) the $\varphi\varphi A$ partners \cite[(2.18)--(2.22)]{Ba3}. Then partners have the same $(m,d_v,d_s)$. Every partner a coloured graph needs occurs in the expansion with the right coefficient, and no counterterm vertex is left without its partner. Hence the coloured expansion is the disjoint union of the partner classes (the \emph{families}), and for $m=0$ these are Ba{\l}aban's.
\end{lemma}
\noindent The lemma uses Ba{\l}aban's construction of the families for $u=0$ \cite[pp.~420, 429--432]{Ba3}.
\begin{proof}
The operations act inside $C$-only parts and preserve the coupling orders (\cite[p.~420]{Ba3}: ``all graphs in the family $G_{\rm ren}$ have the same orders''). The counterterm $\delta m^2_{\rm ref}=\sum_{G_0}\delta m^2_{G_0}$ runs over all $C$-only self-energy graphs of order $\le4$, including those with lower-order insertions \cite[p.~417]{Ba3}. The $C$-only graphs of every order are present as the $u^0$ part of the expansion. The partner carries the same index assignment as $G_0$ \cite[(3.9), p.~435]{Ba3}, and colouring the complement commutes with the operations, so the combinatorial factors are Ba{\l}aban's. $D$-containing subgraphs need no partner, by Lemma~\ref{app:G1:WP-deg}(d). The partition is the orbit decomposition of the partner relation.
\end{proof}

\begin{theorem}[the weighted-line form of {\cite[Prop.~1]{Ba3}}]\label{app:G1:WP-WP1}
Fix $\nbar>12$, $0<\alpha_0<1$, $k\le K$ and $\eta=L^{-k}$. Let $\Omega,\Omega_1,\Omega_2,\tilde B$ satisfy the hypotheses of \cite[Prop.~1, pp.~412, 419--421]{Ba3} on Ba{\l}aban's periodic lattice torus. Work in organisation~$(\alpha)$.
\begin{enumerate}
\item[(a)] (Degrees.) Lemma~\ref{app:G1:WP-deg}(c).
\item[(b)] (Families.) The coloured expansion of the perturbative expressions of the $u$-model, with $\delta m^2_{\rm ref}$, is a finite disjoint sum of the families of Lemma~\ref{app:G1:WP-fam}, each with fixed $(m,d_v,d_s)$. No vacuum counterterm is attached to the vacuum polynomial $T_{\nbar}[\Phi(u)]$ itself, since it is the quantity being bounded. Inside the per-step expressions, by contrast, the vacuum families carry their pieces of $E_k(u)$.
\item[(c)] (Bound.) Let $F$ be a non-vacuum family, or a vacuum family with $m\ge2$, or one with $m=1$ and $o\ge2$, whose $D$-lines carry $D^\eta_k$. Then
\[
\abs{E(F,\{\square(v)\},\Phi_{\rm ext},A_{\rm ext})}\le O(1)\,(e(L^k\eps))^{d_v}(\lambda(L^k\eps))^{d_s}\,\mathrm{e}^{-\delta_0d(\{\square(v)\})}\,\norm{h\Phi_{\rm ext}}_{1,\alpha_0}\norm{h'A_{\rm ext}}_{1,\alpha_0}.
\]
Here $\delta_0>0$ is Ba{\l}aban's tree rate of \cite[(1.33)]{Ba3}, and $O(1)=O(\alpha_0,\nbar)$ (together with his lattice data $L$, $a$) does not depend on $\eps$, $k$, $\Omega$, $\Omega_1$, $\Omega_2$, $\tilde B$, $u$ or $\kappa$. In physical normalisation the family carries the explicit factor $(uc_k)^m$.
\item[(d)] (Per step.) The same bound holds for the sub-sum whose largest index is $k-1$.
\end{enumerate}
\end{theorem}
\noindent The theorem is stated on the torus. The statement that the vacuum family with $m=1$, $o=1$ at a fixed largest index is $O(1)$ is not used. King's cube with Neumann and Dirichlet conditions is not used.
\begin{proof}
Every step transfers a step of \cite[pp.~422--445]{Ba3}.
\emph{What transfers.} Vertices, scalar lines and external fields are unchanged by colouring \cite[pp.~411--422]{Ba3}. So are the parity cancellation $\tr q^{2n+1}=0$ (p.~434), the localisation and background expansion (3.2)--(3.3) (pp.~432--433), and gauge covariance (vector lines do not depend on $\tilde B$ and are gauge-neutral). Ba{\l}aban's Ward--Takahashi identities (2.23)--(2.28) hold at every fixed internal vector field, hence for any Gaussian law of the internal photons, coefficientwise in $s$ for the covariance $G+sD$; only the $s^0$ coefficient is used. The classification \cite[Cor.~2.3, pp.~429--432]{Ba3} and the renormalisation identities \cite[(3.6)--(3.38)]{Ba3} are needed only for $C$-only subgraphs, by Lemma~\ref{app:G1:WP-deg}(d), and there they are Ba{\l}aban's own. $D$-lines enter them only as external legs, through (2.10)--(2.12).
The steps where Ba{\l}aban replaces an internal \emph{photon} line by $C^\xi$ are (3.24), (3.31)--(3.32) and (3.37)--(3.38). Steps (3.16) and (3.29) are not among them: (3.16) replaces only a scalar line, and (3.29) has no internal photon. With a weighted photon, $\sum_{j<j_0}w_jG_{(j)}=G_{j_0}+uc_kD^\eta_{j_0}$: Ba{\l}aban's identity applies to the $C$ part, and the $D$ part has degree at least $2$ more.
\emph{Transversality.} A $C$-only vacuum-polarisation subgraph $\Pi$ has negative degree. It is renormalised by Ba{\l}aban's real-space Ward step \cite[(3.25)--(3.32), pp.~439--442]{Ba3}: Taylor expansion of the external legs (3.26); summation of the internal indices below $j_0$; the replacement by $C^\xi$; vanishing of the symmetric term (3.27)--(3.28); and the identity (2.26) at (3.29). It acts on the class with arbitrary regular external legs (p.~440), and so in particular with $D$-line legs, which need only the H\"older bound of Lemma~\ref{app:G1:WP-deg}(a). A $\Pi$ with an internal $D$-line has $D^\natural\ge2$ and needs nothing.
\emph{The bound (c).}
(1) Lines between localisations at distance at least $1$, whether $C$ or $D$, become pairs of external fields by \cite[(3.1)]{Ba3}. For a $D$-line this lowers $D^\natural$ by $1$, rather than raising it; the conclusion survives. If $l$ is a $D$-line, $G\setminus l$ has the same subgraphs as the line subsets of $G$ not containing $l$, with the same degrees. Its $D$-containing components have $D^\natural\ge1$. Its $C$-only divergent components are $C$-only divergent subgraphs of $G$, already renormalised in the family, or (2.4)-type trees, or odd graphs, which vanish.
(2) Localise and remove the background as in \cite[(3.2)--(3.3)]{Ba3}.
(3) Decompose every line by \cite[(2.6)]{Ba3}, the $D$-lines by Lemma~\ref{app:G1:WP-deg}(a), and sum over orderings (3.4).
(4) In a fixed ordering, the $C$-only components are treated by Ba{\l}aban's construction (3.5), unchanged. The $D$-containing components have $D^\natural\ge1$. Ba{\l}aban's transformations never lower the degree of a larger subgraph (pp.~425, 445). So every component of every subgraph in every admissible ordering has positive degree.
(5) Apply \cite[Prop.~2.2, p.~428]{Ba3} (``lines with the same exponential factors but with arbitrary dimensions instead of $-d+2$''), with the $D$-lines of dimension $-\mathsf d+4$. This gives (2.16), hence (3.3), (3.2) and the displayed bound, exactly as at \cite[pp.~432--434]{Ba3}. Each $D$-line has two $A'$-legs, both counted in $d_v$.
(6) The $u$-model value is $(uc_k)^m$ times this, by the multilinear expansion.
\emph{(d).} In (2.16) the last sum runs over the largest index. Fixing it gives a sub-sum of a majorant sum.
\emph{(a), (b).} Lemmas~\ref{app:G1:WP-deg} and~\ref{app:G1:WP-fam}.
\end{proof}

The theorem covers only organisation~$(\alpha)$. In the other organisation (weighted counterterms, with a finite mass vertex depending on $u$ at every step), the per-step statements are not proved here and are not used.
The following sketched items of \cite{Ba3} are used as in the proof of \cite[Prop.~1]{Ba3}; the colouring does not change them:
\begin{itemize}
\item I1: family combinatorics, ``inspecting all possible cases'' (p.~432);
\item I2: bad localisations of the vacuum counterterms, ``even simpler and we omit them'' (p.~422); this does not affect $E_2$ or Corollary~\ref{app:G1:WP-AC11};
\item I3: convergence of $\delta m^2_{\rm fin}$ and $E_{\rm fin}$, and of the regular remainders treated as external fields (p.~418);
\item I4: \cite[Props.~2.1--2.3]{Ba1}, cited by Ba{\l}aban, behind (2.10)--(2.12);
\item I5: the background H\"older norms (p.~412).
\end{itemize}
In addition, \cite[Prop.~2.2, p.~428]{Ba3} is a remark promoted to a proposition, without a separate proof.
The proof of Theorem~\ref{app:G1:WP-WP1} is a transfer of \cite[pp.~422--445]{Ba3}: only the new ingredients are written. The tree rate given by the proof depends on $\nbar$, through the line-moving rates of \cite[p.~427]{Ba3}, and not only on $\mathsf d$; only the existence of some $\delta_0>0$ is used (Proposition~\ref{app:G1:VSR-CAP}).

\emph{The sketched items and Table~\ref{tab:prop1}.} By page, the items correspond to Table~\ref{tab:prop1} as follows.
\begin{itemize}
\item I1 is step S10 (p.~432). The classification \cite[Cor.~2.3]{Ba3} behind it is step S7.
\item I4 is step S3 (pp.~424, 426). At zero background, the slice bounds for King's slices are Corollary~\ref{app:G1:VSR-VSphys}, and these slices are Ba{\l}aban's by Lemma~\ref{app:G1:WP-Kslices}. The vector lines do not depend on $\tilde B$, so the bounds for the $D$-lines are always of this kind (Lemma~\ref{app:G1:WP-deg}(a)).
\item I2, I3 and I5 concern pp.~422, 418 and 412, which carry no step of Table~\ref{tab:prop1}. Their depth is at most that of Proposition~\ref{prop:prop1}. For I3, the bounds on $\delta m^2_{\rm fin}$ and $E_{\rm fin}$ uniform in the torus are Propositions~\ref{app:G1:I3b} and~\ref{app:G1:Efin}, and the last remainder line is Proposition~\ref{app:G1:VSR-R}. The bad localisations of the moved vacuum pieces (I2) are treated in Proposition~\ref{app:G1:R9vac}.
\end{itemize}
The proof of \cite[Prop.~2.1, pp.~424--428]{Ba3} uses of the lines only the representation \cite[(2.6)]{Ba3} and the bounds \cite[(2.10)--(2.12)]{Ba3}. The leg dimensions enter only through \cite[(2.14)]{Ba3}. Ba{\l}aban's \cite[Prop.~2.2]{Ba3} is this proof with arbitrary line dimensions. In the form used in step (5) of the proof of Theorem~\ref{app:G1:WP-WP1} it reads as follows: for a connected generalised graph all of whose connected components in a fixed ordering have positive degree, \cite[(2.16)]{Ba3} holds. This is Lemma~\ref{app:G1:VSR-A}(i), and Lemma~\ref{app:G1:VSR-Btop} gives the version with the largest index fixed. Subgraphs of type \cite[(2.4)]{Ba3} are treated beforehand by summation by parts \cite[(2.8)]{Ba3}, which is step S4 of Table~\ref{tab:prop1}. The $D$-lines are lines of this generalised kind, of dimension $-\mathsf d+4$, by Lemma~\ref{app:G1:WP-deg}(a).
For the vacuum families with $m\ge2$ that are used in Lemma~\ref{app:G1:E2B}, part (c) is replaced there by Theorem~\ref{app:K:L17-1}, which is the corresponding bound at zero background. For the other uses of part (c), namely Corollary~\ref{app:G1:WP-AC11} and the non-vacuum families, the items I1--I5 are used as described above.

\begin{remark}[the $E_2$ constant]\label{app:G1:WP-E2}
The constant $4C_W$ of Lemma~\ref{app:G1:E2B} is conservative: $m(m-1)\le m^2$ costs a factor $2$, and $c_K$ overcounts the top $D$-slice by $L^2$. The constant $C_W$ depends on $\mu_*$ only through $\mu_*\eps_0$. Such factors enter through $C_{\mu_*}-G_K$, which is treated as an external field (I3). They do not depend on $\eps$, so they do not affect $\sup_\eps$.
\end{remark}

\paragraph{The vacuum coefficients linear in $u$.}

\begin{corollary}[the vacuum coefficients linear in $u$ on the torus]\label{app:G1:WP-AC11}
For every $(\alpha,\beta)\ne(2,0)$ with $\alpha+\beta\le\nbar$ (so $o\ge2$ whenever the coefficient is non-zero, since $\alpha\ge2$ is even), and every $M\in\{m_1,m_2\}$,
\[
\abs{[\mathcal E(\cdot;M)]_{1,\alpha\beta}}\le C_{1,\alpha\beta}(M)\,\kappa\mu_*^2L^2\,\abs T\,\eps_0^{\,o-1},
\]
uniformly in $\eps=2^{-K}\eps_0$ and in the torus. Here $C_{1,\alpha\beta}$ is Ba{\l}aban's $O(1)(\nbar,\alpha_0,\mathsf d,L,a,M)$ times a graph count, independent of $\eps$, $u$, $\kappa$ and $\ell_T$.
\end{corollary}
\noindent This is an instance of Theorem~\ref{app:G1:WP-WP1}(c). The excluded coefficient $(2,0)$ ($D^\natural=0$, logarithmically divergent) is treated by its mass difference in Lemma~\ref{app:G1:CU1}.

\begin{proof}
The covariance $C_{\mu_*}+u\Gamma^B$ is linear in $u$. So the $u^1$ part of $T_{\nbar}[\Phi_M(u)]$ is the sum of the coloured vacuum graphs with exactly one $D$-line. Each such graph is $e^\alpha\lambda^\beta$ times a number, the $\delta m^2_{\rm ref}$ vertices carrying their own orders. By Lemma~\ref{app:G1:WP-fam}, the families partition the set of graphs with fixed $(m,\alpha,\beta)=(1,\alpha,\beta)$, so $[\mathcal E]_{1,\alpha\beta}=\pm\sum_F c_K\,\mathrm{val}(F)$ in scale-$K$ units.
For $o\ge2$, every such family is covered by Theorem~\ref{app:G1:WP-WP1}(c) ($D^\natural\ge1$, Lemma~\ref{app:G1:WP-deg}(c)). Its $C$-only photon self-energy closed by the single $D$-line is renormalised by Ba{\l}aban's Ward step \cite[(3.25)--(3.32)]{Ba3}, with the two ends of the $D$-line as external legs. No vacuum counterterm is attached (Theorem~\ref{app:G1:WP-WP1}(b)).
Counting: fix one vertex in a unit block of scale $K$. The other localisations sum to $O(1)$ by $\mathrm{e}^{-\delta_0d}$, and there are $\abs T/t_K^3$ blocks. A family therefore contributes at most
\[
O(1)(e^2t_K)^{\alpha/2}(\lambda t_K)^\beta c_K\,t_K^{-3}\abs T=O(1)\,e^\alpha\lambda^\beta\kappa\mu_*^2L^2\,t_K^{\,o-1}\abs T.
\]
Since $o-1\ge1$ and $t_K\le\eps_0$ (only $t_K\le\eps_0$ is used), $t_K^{o-1}\le\eps_0^{o-1}$. Sum over the finitely many families.
\end{proof}

\paragraph{The diagonal of King's slices.}

\emph{King's objects.} These are as in \cite[(2.14)--(2.17)]{K1}. $Q_j$ is the block average to the scale-$j$ block lattice, and $P_j$ is the orthogonal projection onto fields that are constant on scale-$j$ blocks. $a_j$ are King's block constants, $r_j:=a_jt_j^{-2}$, and $G_j$ is King's scale-$j$ Green's function. $C^{(j)}=(\Delta^{(j)}+at_{j+1}^{-2}Q^*Q)^{-1}$ is the fluctuation covariance. $H_j=F_jC^{(j)}F_j^*$ with $F_j=r_jG_jQ_j^*$. Per component, $C_{\mu_*}=(-\Delta_\eps+\mu_*^2)^{-1}$. Put $a_\flat:=\inf_ja_j=a(1-L^{-2})$; we use $a_j\le a$.

\begin{lemma}[King's slice at a point]\label{app:G1:WP-RBdiag}
On the periodic lattice torus, for $1\le j<K$ and every site $x$ (per component),
\[
H_j(x,x)\le C_H/t_j,\qquad C_H:=c_C\,a^2\,(m_a+3^{1/4}C_{\rm osc})^2,\quad c_C:=\frac4{a_\flat}+\frac{L^2}8+\frac{L^2}a,\quad m_a:=\max\Big(\frac18,\frac1{a_\flat}\Big),
\]
and $H_0(x,x)\le C_{\mu_*}(x,x)$. Here $C_{\rm osc}$ is the constant of the lattice Green's-function estimate
\[
\text{(G-osc)}\qquad\ip{\delta_x-\delta_y}{C_{\mu_*}^2(\delta_x-\delta_y)}\le C_{\rm osc}^2\abs{x-y},
\]
which holds on the periodic lattice with $C_{\rm osc}^2\le8.1$, uniformly in $\eps$, $\ell_T$ and $\mu_*$. Consequently
\[
\max_xD(x,x)\le L^2\eps\,c_1+C_HL^2\eps_0/(L-1),
\]
where $c_1$ is such that $\eps\,C_{\mu_*}(x,x)\le c_1$ (see Remark~\ref{app:G1:WP-zeromode} and Lemma~\ref{app:G1:diagC}).
\end{lemma}
\noindent On the periodic lattice, (G-osc) is proved below, so the lemma holds on the torus without further conditions. The proof uses the block Poincar\'e inequality (Lemma~\ref{app:G1:DOM-poinc}) and the variational form of $\Delta^{(j)}$ (Lemma~\ref{app:G1:Dk-var}), through Lemma~\ref{app:G1:Cj-bound}, which is step~(a) of the proof. The lemma gives step~(2) of Lemma~\ref{app:G1:CU1}.
\begin{proof}
(a) \emph{$C^{(j)}\le c_Ct_j^2$ on the scale-$j$ block lattice.} By the variational form, $\ip{B}{\Delta^{(j)}B}_j=\min_A[\ip{A}{(-\Delta_\eps+\mu_*^2)A}+r_j\norm{Q_jA-B}_j^2]$. Let $A^*$ be the minimiser and $B':=Q_jA^*$. Then $\norm{(1-Q^*Q)B'}=\norm{(P_j-P_{j+1})A^*}\le\norm{(1-P_{j+1})A^*}\le(t_{j+1}^2/8)^{1/2}\ip{A^*}{(-\Delta_\eps+\mu_*^2)A^*}^{1/2}$, by the block Poincar\'e inequality. Here $Q_{j+1}=QQ_j$ and $Q_jQ_j^*=1$; the factor $1/8$ comes from the lower bound $8t_{j+1}^{-2}$ on the smallest non-zero Neumann eigenvalue of a block of side $t_{j+1}$ with at least $2$ sites per side. Also $\norm{QB'}\le\norm{QB}+\norm{B-B'}$. Hence $\norm B\le2\norm{B-B'}+(t_{j+1}^2/8)^{1/2}\ip{A^*}{(-\Delta_\eps+\mu_*^2)A^*}^{1/2}+\norm{QB}$. Cauchy--Schwarz with weights $(r_j,1,\beta)$, $\beta=at_{j+1}^{-2}$, gives $\norm B^2\le(4/r_j+t_{j+1}^2/8+1/\beta)\ip{B}{C^{(j)-1}B}=t_j^2(4/a_j+L^2/8+L^2/a)\ip{B}{C^{(j)-1}B}\le c_Ct_j^2\ip{B}{C^{(j)-1}B}$.
(b) Hence $H_j=F_jC^{(j)}F_j^*\le c_Ct_j^2F_jF_j^*$, and $H_j(x,x)\le c_Ct_j^2r_j^2\norm{P_jG_j\delta_x}^2$ with $\delta_x:=\eps^{-3}1_{\{x\}}$.
(c) Split $\delta_x=\rho_x+v$ with $\rho_x:=P_j\delta_x=t_j^{-3}1_{\square_x}$ and $P_jv=0$. First, $\norm{P_jG_j\rho_x}\le\norm{G_j}\norm{\rho_x}\le m_at_j^2\,t_j^{-3/2}$, using $\norm{G_j}\le t_j^2\max(1/8,1/a_\flat)$. Second, $G_j=(-\Delta_\eps+\mu_*^2)^{-1}-G_jr_jP_j(-\Delta_\eps+\mu_*^2)^{-1}$ gives $P_jG_jv=(1-r_jP_jG_jP_j)P_j(-\Delta_\eps+\mu_*^2)^{-1}v$, with $0\le r_jP_jG_jP_j\le1$ on $\operatorname{ran}P_j$. So $\norm{P_jG_jv}\le\sup_{\norm w\le1}[\psi_w(x)-\mathrm{avg}_{\square_x}\psi_w]$ with $\psi_w:=C_{\mu_*}w$. Moreover $\abs{\psi_w(x)-\psi_w(y)}\le\ip{\delta_x-\delta_y}{C_{\mu_*}^2(\delta_x-\delta_y)}^{1/2}\norm w\le C_{\rm osc}(\sqrt3t_j)^{1/2}$ for $y\in\square_x$. Hence $\norm{P_jG_j\delta_x}\le(m_a+3^{1/4}C_{\rm osc})t_j^{1/2}$, and $H_j(x,x)\le c_Ca_j^2t_j^{-2}(m_a+3^{1/4}C_{\rm osc})^2t_j\le C_H/t_j$.
\emph{(G-osc) on the periodic lattice.} Write $s:=\abs{x-y}$ and $\lambda(k)=\sum_i(4/\eps^2)\sin^2(\eps k_i/2)\ge(4/\pi^2)\abs k^2$. The left side equals $\ell_T^{-3}\sum_{k}\abs{\mathrm{e}^{ik\cdot x}-\mathrm{e}^{ik\cdot y}}^2(\lambda(k)+\mu_*^2)^{-2}$; the $k=0$ term vanishes. The summand is at most $(\pi^4/16)\min(4\abs k^{-4},s^2\abs k^{-2})$, which is radially non-increasing. Every $k'$ in the cell of side $2\pi/\ell_T$ around $k\ne0$ has $0.134\abs k\le\abs{k'}\le1.866\abs k$. This gives the Riemann bound by $(1.87^3/(8\pi^3))\int f$ (dropping the BZ restriction only adds positive terms), and $\int(\pi^4/16)\abs q^{-4}\min(4,\abs q^2s^2)\,d^3q=\pi^5s$. So the left side is at most $8.07\,s\le8.1\,s$.
\emph{The display.} We have $D=\sum_{j<K}t_{j+1}^2H_j$ and $t_1^2H_0(x,x)\le L^2\eps^2C_{\mu_*}(x,x)\le L^2\eps c_1$. For $j\ge1$, $\sum_{j\ge1}t_{j+1}^2/t_j=L^2\sum_{j\ge1}t_j\le L^2\eps_0/(L-1)$.
\end{proof}

\begin{lemma}[the fluctuation covariance]\label{app:G1:Cj-bound}
For $1\le j<K$, $C^{(j)}\le c_Ct_j^2$ on $\ell^2(\Lambda^{(j)})$, with $c_C$ as in Lemma~\ref{app:G1:WP-RBdiag}.
\end{lemma}
\begin{proof}
Let $B\in\ell^2(\Lambda^{(j)})$, let $A^*:=F_jB$ be the minimiser of Lemma~\ref{app:G1:Dk-var}, and put $B':=Q_jA^*$. Then
\[
r_j\norm{B-B'}_j^2+\ip{A^*}{MA^*}=\ip B{\Delta^{(j)}B}_j .
\]
Decompose $B=(B-B')+(1-Q^*Q)B'+Q^*QB'$.

The operators $1-Q^*Q$ and $P_j-P_{j+1}$ are orthogonal projections. So Lemma~\ref{app:G1:blockops}(a)--(c) gives $\norm{(1-Q^*Q)B'}_j^2=\ip{A^*}{(P_j-P_{j+1})A^*}=\norm{(P_j-P_{j+1})A^*}^2$. Since $P_j-P_{j+1}\le1-P_{j+1}$, Lemma~\ref{app:G1:DOM-poinc}(a) at scale $j+1\le K$ gives
\[
\norm{(P_j-P_{j+1})A^*}\le\norm{(1-P_{j+1})A^*}\le(t_{j+1}^2/8)^{1/2}\ip{A^*}{MA^*}^{1/2}.
\]
By Lemma~\ref{app:G1:blockops}(b), $\norm{Q^*QB'}_j=\norm{QB'}_{j+1}\le\norm{QB}_{j+1}+\norm{B'-B}_j$. Hence
\[
\norm B_j\le2\norm{B-B'}_j+(t_{j+1}^2/8)^{1/2}\ip{A^*}{MA^*}^{1/2}+\norm{QB}_{j+1}.
\]
Apply the Cauchy--Schwarz inequality with weights $(r_j,1,\beta_j)$ and use $\norm{QB}_{j+1}^2=\ip{B}{Q^*QB}_j$:
\[
\norm B_j^2\le\Big(\frac4{r_j}+\frac{t_{j+1}^2}8+\frac1{\beta_j}\Big)\ip{B}{(\Delta^{(j)}+\beta_jQ^*Q)B}_j=t_j^2\Big(\frac4{a_j}+\frac{L^2}8+\frac{L^2}a\Big)\ip{B}{(C^{(j)})^{-1}B}_j .
\]
By (F1) of Definition~\ref{app:G1:DOM-set}, $a_j>a_\flat$, so the bracket is at most $c_C$. Thus $(C^{(j)})^{-1}\ge(c_Ct_j^2)^{-1}$, that is, $C^{(j)}\le c_Ct_j^2$.
\end{proof}

\begin{lemma}[the diagonal of $C_{\mu_*}$ on the torus]\label{app:G1:diagC}
In the setting of Definition~\ref{app:G1:DOM-set}, for every $x\in T_\eps$,
\[
\eps\,C_{\mu_*}(x,x)\le\frac{\sqrt3\,\pi}{8}\Big(1+\frac{\sqrt3}2\Big)^3+\frac{\eps}{\ell_T^3\mu_*^2}<4.42+\frac{\eps}{\ell_T^3\mu_*^2}.
\]
In particular, the constant $c_1$ of Lemmas~\ref{app:G1:WP-RBdiag} and~\ref{app:G1:CU1} may be taken as $c_1:=4.42+\eps_0/(\ell_0^3\mu_*^2)$, for all $\eps\le\eps_0$ and $\ell_T\ge\ell_0$.
\end{lemma}
\begin{proof}
By (F2), $\ell_T/\eps\ge4$, so $-\Delta_\eps$ has the symbol $\lambda(k)=\sum_i(4/\eps^2)\sin^2(\eps k_i/2)$ on $T^*:=(2\pi/\ell_T)\Z^3\cap(-\pi/\eps,\pi/\eps]^3$. Hence $C_{\mu_*}(x,x)=\ell_T^{-3}\sum_{k\in T^*}(\lambda(k)+\mu_*^2)^{-1}$ for every $x$. The term $k=0$ is $\ell_T^{-3}\mu_*^{-2}$.

For $k\ne0$, the inequality $\sin(t/2)\ge t/\pi$ on $[0,\pi]$ gives $\lambda(k)\ge(4/\pi^2)\abs k^2$. Every $k\in T^*$ satisfies $\abs k\le R':=\sqrt3\pi/\eps$. Put $c:=1+\sqrt3/2$ and $f(q):=\abs q^{-2}\mathbf 1_{\{\abs q\le R'\}}$. The cubes $Z_k:=k+[-\pi/\ell_T,\pi/\ell_T]^3$ have disjoint interiors and volume $(2\pi/\ell_T)^3$. For $k\ne0$ and $k'\in Z_k$ we have $\abs{k'-k}\le\sqrt3\pi/\ell_T\le(\sqrt3/2)\abs k$, so $0<\abs{k'}\le c\abs k$. Hence $\abs{k'/c}\le\abs k\le R'$ and $f(k)\le f(k'/c)$. Averaging over $Z_k$ and summing over $k$ gives
\[
\sum_{k\in T^*\setminus\{0\}}\abs k^{-2}\le\Big(\frac{\ell_T}{2\pi}\Big)^3\int_{\R^3}f(q/c)\,d^3q=\Big(\frac{\ell_T}{2\pi}\Big)^3c^3\,4\pi R' .
\]
Therefore
\[
\eps\,\ell_T^{-3}\sum_{k\ne0}(\lambda(k)+\mu_*^2)^{-1}\le\eps\,\frac{\pi^2}{4}\,\frac{4\pi c^3R'}{(2\pi)^3}=\frac{\sqrt3\pi}{8}\,c^3=4.4195\ldots
\]
Adding the term $k=0$ proves the bound. The right side increases in $\eps$ and decreases in $\ell_T$, which gives the statement about $c_1$.
\end{proof}
\noindent With Lemma~\ref{app:G1:diagC}, the display of Lemma~\ref{app:G1:WP-RBdiag} holds in the form
\[
\max_xD(x,x)\le L^2\eps\,c_1+\frac{C_HL^2\eps_0}{L-1},\qquad c_1=4.42+\frac{\eps_0}{\ell_0^3\mu_*^2},
\]
uniformly in $\eps\le\eps_0$ and $\ell_T\ge\ell_0$. This is the form used in step~(2) of Lemma~\ref{app:G1:CU1}.

\begin{remark}[the torus zero mode]\label{app:G1:WP-zeromode}
On King's Neumann cube one has $C_{\mu_*}(x,x)\le c_1/\eps+1/(\ell^3\mu_*^2)$ with $c_1\le36$. On the periodic torus, the Fourier sum $C_{\mu_*}(x,x)=\ell_T^{-3}\sum_{k}(\lambda(k)+\mu_*^2)^{-1}$ contains the zero-mode term $\ell_T^{-3}\mu_*^{-2}$, and the Riemann argument of the proof bounds the rest by $c_1/\eps$. So on $T$ one should read $c_1$ as $c_1+\eps/(\ell_0^3\mu_*^2)$. This adds at most $L^2\eps^2/(\ell_0^3\mu_*^2)$ to $\max_xD(x,x)$, which is bounded and tends to $0$ with $\eps$, so no conclusion of Lemma~\ref{app:G1:CU1} changes.
\end{remark}

\paragraph{Torus images and the heat kernel.}

For $\Lambda\in\{T,\eps\Z^3\}$ write $C_{M,\Lambda}:=\int_\Lambda(\hat k^2+M^2)^{-1}$, so that $C_{M,T}=C^T(M^2)$ and $C_{M,\eps\Z^3}=C^\infty(M^2)$ in the notation of the beginning of this appendix. Write $G^\eps_M(x)=\int_{\rm BZ}\hat G_M\mathrm{e}^{ik\cdot x}$, as there.

\begin{lemma}[torus images, free of $\eps$]\label{app:G1:WP-img}
For every $\eps$ with $\ell_T/\eps\in\N$, $\ell_T/\eps\ge2$, and every $M>0$,
\[
E_T(M):=C_{M,T}-C_{M,\eps\Z^3}=\sum_{n\in\Z^3\setminus0}G^\eps_M(n\ell_T)\le\bar E(M,\ell_T),
\]
\[
\bar E(M,\ell):=\frac{7.37}{\ell}+\frac{3\pi}{4\ell}E_1\Big(\frac{M^2\ell^2}8\Big)+\frac{12\pi\,\mathrm{erfc}(M\ell/(2\sqrt2))}{M\ell^2}+\frac{64\,\mathrm{e}^{-M^2\ell^2/8}}{M^2\ell^3},
\]
where $E_1$ is the exponential integral. Moreover $\bar E$ is decreasing in $\ell$. The constant is $7.37\ge3\pi/4+768\sqrt\pi\,(2/5)\,8^{-5/2}+2=7.364$.
\end{lemma}
\noindent The bound is only polynomial in $1/\ell_T$, whereas the true size is expected to be exponentially small. For the photon mass, $E_T(\mu_*)\ge1/(\mu_*^2\ell_T^3)$ because of the zero mode.

\begin{proof}
(1) $G^\eps_M(x)=\int_0^\infty \mathrm{e}^{-M^2t}p_t(x)\,dt$, where $p_t=\prod_\mu p^{(1)}_t(x_\mu)$ and $p^{(1)}_t(x)=\eps^{-1}\mathbb P(X_t=x)$ for the continuous-time walk on $\eps\Z$ with jumps $\pm\eps$ at rate $\eps^{-2}$ each. Its Fourier symbol is $(4/\eps^2)\sin^2(\eps k/2)$, $\operatorname{Var}X_t=2t$, and $p^{(1)}_t$ is non-increasing in $\abs x$ ($I_n(u)$ decreases in $n$).
(2) By monotonicity, and since $(n\ell-\ell/2,n\ell]$ contains at least $\ell/(4\eps)$ sites when $\ell/\eps\ge2$, we have $p^{(1)}_t(n\ell)\le(4\eps/\ell)\sum_{x\in(n\ell-\ell/2,n\ell]}p^{(1)}_t(x)$ for $n\ge1$, and symmetrically for $n\le-1$. Hence $\delta(t):=\sum_{n\ne0}p^{(1)}_t(n\ell)\le(4/\ell)\mathbb P(\abs{X_t}\ge\ell/2)\le(4/\ell)\min(1,8t/\ell^2)$, by Chebyshev.
(3) $p^0:=p^{(1)}_t(0)=\eps^{-1}\mathrm{e}^{-v}I_0(v)$ with $v=2t/\eps^2$. Since $\mathrm{e}^{-v}I_0(v)=\pi^{-1}\int_0^\pi \mathrm{e}^{-v(1-\cos\theta)}d\theta$ and $1-\cos\theta\ge2\theta^2/\pi^2$ on $[0,\pi]$, $\mathrm{e}^{-v}I_0(v)\le\sqrt{\pi/(8v)}$. Hence $p^0\le a_*/\sqrt t$ with $a_*:=\sqrt\pi/4$, with no $\eps$.
(4) $\sum_{n\in\Z^3\setminus0}p_t(n\ell)=(p^0+\delta)^3-(p^0)^3=3(p^0)^2\delta+3p^0\delta^2+\delta^3$. Integrate against $\mathrm{e}^{-M^2t}$, using $\mathrm{e}^{-M^2t}\le1$ on $t\le\ell^2/8$ and the two forms of $\min(1,8t/\ell^2)$. The resulting closed forms give $\bar E$.
\end{proof}

\begin{corollary}[the torus heat kernel; the form used in Lemma~\ref{app:G1:CU1}]\label{app:G1:WP-heat}
Under the hypotheses of Lemma~\ref{app:G1:WP-img}, with $S_q=\abs T^{-1}\sum_{q\in T^*}$ and $b:=4/\ell_T$,
\[
S_q\,\mathrm{e}^{-t\hat q^2}\le\big(a_*/\sqrt t+b\big)^3\quad(t>0),\qquad S_q\,\hat G_M(q)^2\le C_T(M,\ell_T)\le C_T(M,\ell_0)\quad(\ell_T\ge\ell_0),
\]
with $C_T(M,\ell):=a_*^3\sqrt\pi/M+3a_*^2b/M^2+\frac{3\sqrt\pi}2a_*b^2/M^3+b^3/M^4$, where $b=4/\ell$.
\end{corollary}
\noindent This is steps (2)--(3) of the proof of Lemma~\ref{app:G1:WP-img}.
\begin{proof}
By Poisson summation on the lattice torus (as in Lemma~\ref{app:G1:CNtorus}), $S_q\mathrm{e}^{-t\hat q^2}=\sum_{n\in\Z^3}p_t(n\ell_T)=(p^0+\delta)^3$, using the product structure of $p_t$. By steps (2)--(3), this is at most $(a_*/\sqrt t+4/\ell_T)^3$. Then $S_q\hat G_M^2=\int_0^\infty t\mathrm{e}^{-tM^2}S_q\mathrm{e}^{-t\hat q^2}dt$. Integrating term by term, $\int_0^\infty t^{1/2+j/2-3/2+1-1}\cdots$ gives the four terms: $a_*^3\Gamma(\tfrac12)M^{-1}$, $3a_*^2b\,M^{-2}$, $3a_*b^2\Gamma(\tfrac32)M^{-3}$ and $b^3M^{-4}$. Finally $C_T$ is increasing in $b$, hence decreasing in $\ell$.
\end{proof}

\begin{lemma}[the free mass difference]\label{app:G1:WP-freeC}
For $M,m>0$ and every $\eps$, $\abs{C^\infty(M^2)-C^\infty(m^2)}\le\abs{M-m}\,\pi^2/32$.
\end{lemma}
\noindent This is the free bound used in step~(iv) of the proof of Proposition~\ref{app:G1:CN}.
\begin{proof}
Take $M<m$; the other case follows by symmetry. Then $C^\infty(M^2)-C^\infty(m^2)=(m^2-M^2)\int_{\rm BZ}\frac{d^3k}{(2\pi)^3}\,[(\hat k^2+M^2)(\hat k^2+m^2)]^{-1}\ge0$. Since $\sin x\ge2x/\pi$ on $[0,\pi/2]$, $\hat k^2\ge(4/\pi^2)\abs k^2$ on the BZ. Substitute $q=(2/\pi)k$ and extend to $\R^3$: the integral is at most $(\pi/2)^3\int_{\R^3}\frac{d^3q}{(2\pi)^3}[(q^2+M^2)(q^2+m^2)]^{-1}=(\pi^3/8)/(4\pi(M+m))$. Multiplying by $m^2-M^2$ gives $(m-M)\pi^2/32$.
\end{proof}

\subsection{The Ward identities and the Dyson chain}\label{app:G1:ssec-ward-dyson}

In this subsection $T$ is a torus of side $\ell_T$ with $\ell_T/\eps\in\N$. Its sites are the points of $\eps\Z^3$ modulo $\ell_T\Z^3$, its bonds are the pairs $b=(x,\mu)$, $\abs T:=\ell_T^3$, and $T^*:=(2\pi/\ell_T)\Z^3$ modulo $(2\pi/\eps)\Z^3$. We write $\hat p^2:=\sum_\mu4\eps^{-2}\sin^2(\eps p_\mu/2)$ for the symbol of $-\Delta_\eps$ and $\tilde q_\nu:=(\mathrm{e}^{i\eps q_\nu}-1)/\eps$. For a translation-invariant operator $Y$ with kernel $Y(x,y)$ with respect to $\sum_y\eps^3$, put $\hat Y(p):=\sum_z\eps^3\mathrm{e}^{-ip\cdot z}Y(z,0)$. The exponential series of a formal power series $W$ is written $\exp(W)$.

\paragraph{Formal perturbation theory.}
Let $I=I_s\sqcup I_A$ be a finite set. Let $\E_0$ be the centred real Gaussian expectation of variables $(\Phi_i)_{i\in I}$ with covariance $C=C^s\oplus C^A$, where $C^s$ (on $I_s$) and $C^A$ (on $I_A$) are positive definite. The variables with $i\in I_s$ are called scalar and are written $\varphi_a$; those with $i\in I_A$ are called photon variables and are written $A_b$. Let $g=(g_1,\dots,g_p)$ be couplings, and let $V=\sum_{r\in\mathcal R}g^{(r)}m_r$. Here each $m_r$ is a product of variables, its factors, listed in a fixed order, with an even number of scalar factors. Each $g^{(r)}$ is a real formal power series in $g$ without constant term, and for each $d$ only finitely many $g^{(r)}$ have a nonzero term of degree at most $d$. For a polynomial $F$ in the variables put $\ev F:=\E_0[F\exp(-V)]/\E_0[\exp(-V)]$. This is a formal power series in $g$, and each of its coefficients is a finite sum of Gaussian moments.

A labelled graph $\gamma$ with legs $F=\Phi_{i_1}\cdots\Phi_{i_p}$ consists of three data:
\begin{itemize}
\item a number $n\ge0$;
\item labels $r_1,\dots,r_n\in\mathcal R$, where vertex $j$ carries $m_{r_j}$;
\item a pairing $\pi$ of all factors of $Fm_{r_1}\cdots m_{r_n}$ that pairs scalar factors with scalar factors and photon factors with photon factors.
\end{itemize}
Its weight is
\[
w(\gamma):=\frac{(-1)^n}{n!}\prod_{j=1}^ng^{(r_j)}\prod_{\{\Phi_i,\Phi_{i'}\}\in\pi}\E_0[\Phi_i\Phi_{i'}] .
\]
The pairs are the lines of $\gamma$. The legs are univalent external vertices, and a line joins the vertices or legs that carry its two factors. Since $\E_0[\varphi_aA_b]=0$, Wick's theorem gives, coefficientwise, $\E_0[F\exp(-V)]=\sum_\gamma w(\gamma)$ over the labelled graphs with legs $F$. At degree $d$ only graphs with $n\le d$, and with labels from a finite set, contribute.

An amputated two-leg piece (a piece, for short) is defined in the same way, with $n\ge1$, except that two distinct factors of $m_{r_1}\cdots m_{r_n}$, the in-leg and the out-leg, are left unpaired. Its weight is given by the same formula. The vertices carrying the in-leg and the out-leg are the in-vertex and the out-vertex; they may coincide. A piece is irreducible if it is connected and no line separates its in-vertex from its out-vertex. It is one-particle irreducible (1PI) if it is connected and stays connected after the removal of any one line.

Let $G$ be a graph of content $\nu$ ($\nu_r$ vertices with label $r$) in which, for each $r$, the vertices with label $r$ are numbered $1,\dots,\nu_r$, and put $\nu!:=\prod_r\nu_r!$. The labelled graphs inducing $G$ correspond bijectively to the $n!/\nu!$ sequences $(r_1,\dots,r_n)$ of content $\nu$, so the sum of their weights is $A(G)/\nu!$ with $A(G)=(-1)^n\prod_jg^{(r_j)}\prod\E_0[\cdot]$. This is the convention of Lemma~\ref{app:G1:WP-q2a}, so every sum below has the same value in either convention.

\begin{lemma}[parity of the scalar factors]\label{app:G1:parity}
Let $\gamma$ be a connected labelled graph or a connected piece, and let $e$ be a line whose removal disconnects $\gamma$. Let $X$ be a component of $\gamma\setminus e$, and let $s(X)$ be the number of scalar factors among the legs of $\gamma$ in $X$ (for a piece: among its in-leg and out-leg factors at vertices of $X$). Then the end of $e$ in $X$ is scalar if and only if $s(X)$ is odd. In particular:
\begin{enumerate}
\item[(i)] in a graph with legs $\varphi_a,\varphi_b$, every line separating $a$ from $b$ is scalar; in a graph with legs $A_b,A_{b'}$, every such line is a photon line;
\item[(ii)] in a piece whose two legs are of the same type, every line separating the in-vertex from the out-vertex is of that type, and every disconnecting line that leaves both legs on one side is a photon line;
\item[(iii)] there is no labelled graph with the single leg $\varphi_a$.
\end{enumerate}
\end{lemma}

\begin{proof}
The vertices in $X$ carry an even number of scalar factors each. The lines inside $X$ pair the factors at $X$ in pairs of equal type. These factors are the vertex factors in $X$ and, for a graph, the legs in $X$. The only factors at $X$ not paired inside $X$ are the end of $e$ and, for a piece, the in-leg and out-leg factors at $X$. Counting scalar factors modulo $2$ gives the first claim. For (i), the component of $a$ contains exactly one leg. For (ii), the component of the in-vertex contains exactly one leg in the first case, and the other side contains no leg in the second case. For (iii), the number of scalar factors is odd, so no pairing exists.
\end{proof}

\begin{lemma}[linked clusters and Dyson chains]\label{app:K:L16-3}\label{app:G1:dyson}
In the setting above, put $G^s_{ab}:=\ev{\varphi_a\varphi_b}$ for $a,b\in I_s$, and $G^A_{bb'}:=\ev{A_bA_{b'}}-\ev{A_b}\ev{A_{b'}}$ for $b,b'\in I_A$.
\begin{enumerate}
\item[(a)] Coefficientwise, $G^s_{ab}$ and $G^A_{bb'}$ are the sums of the weights of the connected labelled graphs with legs $\varphi_a\varphi_b$, resp.\ $A_bA_{b'}$.
\item[(b)] For $\sharp\in\{s,A\}$ and $i,j\in I_\sharp$, let $K^\sharp_{ij}$ be the sum of the weights of the irreducible pieces with in-leg a factor $\Phi_i$ and out-leg a factor $\Phi_j$. By Lemma~\ref{app:G1:parity}(ii), a connected piece of this kind is irreducible if and only if no line of the leg type separates its two legs. Then, as matrices over the formal power series in $g$,
\[
G^\sharp=\sum_{k\ge0}C^\sharp(K^\sharp C^\sharp)^k,\qquad (G^\sharp)^{-1}=(C^\sharp)^{-1}-K^\sharp ,
\]
and the series is finite in each order.
\item[(c)] Let $I_s$ be the pairs $a=(x,i)$, with $x$ a site of $T$ and $i\in\{1,2\}$, and let $I_A$ be the bonds of $T$. For a site $z$ define the substitutions
\[
\tau_z:\ \varphi_i(x)\mapsto\varphi_i(x+z),\ A_{x,\mu}\mapsto A_{x+z,\mu};\qquad \iota:\ \varphi_i(x)\mapsto\varphi_i(-x),\ A_{x,\mu}\mapsto-A_{-x-\eps e_\mu,\mu}.
\]
Assume that each of them leaves the Gaussian law $\E_0$ invariant and leaves $V$ invariant, coefficientwise in $g$. Then $K^s$ and $K^A$ are the sums over the 1PI pieces only. These assumptions hold in either of the following two cases.
\begin{itemize}
\item $C^s=(\eps^3(-\Delta_\eps+M^2))^{-1}\otimes\mathbf 1_2$ and $C^A=(\eps^3(-\Delta_\eps+\mu_*^2))^{-1}$ on the bond functions, componentwise. The interaction $V$ consists of the hopping vertices of \eqref{eq:matteraction}, the quartic vertex $\lambda\sum_x\eps^3\abs{\varphi(x)}^4$ and a constant counterterm $\frac12c\sum_x\eps^3\abs{\varphi(x)}^2$; these are the vertices of Lemma~\ref{app:G1:WP-q2a}. This is the uniform-mass model with any constant counterterm $c$.
\item The same, with the scalar covariance $C_t\otimes\mathbf 1_2$ and the counterterm $\delta m^2_t$ of the proof of Proposition~\ref{app:G1:CN}.
\end{itemize}
\item[(d)] In the setting of (c), let $C^s=(-\Delta_\eps+M^2)^{-1}\otimes\mathbf 1_2$ with respect to $\sum_y\eps^3$, that is, $C^s=(\eps^3(-\Delta_\eps+M^2))^{-1}\otimes\mathbf 1_2$ as a matrix. Let $V$ contain the term $\frac12c\sum_x\eps^3\abs{\varphi(x)}^2$, where $c$ is a formal power series without constant term. Let $\mathcal G$ be the operator with kernel $G^s$ with respect to $\sum_y\eps^3$. Let $K^s_{\rm ct}$ be the sum of the weights of the one-vertex pieces whose vertex is a term of this counterterm, and put $\Sigma:=-\eps^{-3}(K^s-K^s_{\rm ct})$. Then $K^s_{\rm ct}=-c\,\eps^3\mathbf 1$ and $\mathcal G^{-1}=-\Delta_\eps+M^2+c+\Sigma$. If $C^s$ and $V$ are invariant under every $\tau_z$, then $\hat{\mathcal G}(p)^{-1}=\hat p^2+M^2+c+\hat\Sigma(p)$ for $p\in T^*$.
\end{enumerate}
\end{lemma}

\begin{proof}
(a) Let $\gamma$ have legs $F$ and $n$ vertices. Let $S\subset\{1,\dots,n\}$ be the set of vertices joined to some leg by a path, and put $k:=\abs S$. Renumbering $S$ and its complement increasingly gives a bijection $\gamma\mapsto(S,\gamma_S,\gamma_{S^c})$ onto the triples formed by:
\begin{itemize}
\item a $k$-subset $S$;
\item a graph $\gamma_S$ with legs $F$ and $k$ vertices, each joined to a leg;
\item a vacuum graph $\gamma_{S^c}$ with $n-k$ vertices.
\end{itemize}
Moreover $w(\gamma)=\binom nk^{-1}w(\gamma_S)w(\gamma_{S^c})$. Summing over the $\binom nk$ sets $S$ gives $\E_0[F\exp(-V)]=L_F\,\E_0[\exp(-V)]$, where $L_F$ is the sum over the graphs all of whose vertices are joined to a leg. The series $\E_0[\exp(-V)]$ has constant term $1$, so $\ev F=L_F$.

For $F=\Phi_i\Phi_j$, a graph counted in $L_F$ is either connected or has exactly two components, one containing $i$ and one containing $j$. The same splitting, with $S$ the vertex set of the component of $i$, shows that the graphs of the second kind sum to $L_{\Phi_i}L_{\Phi_j}=\ev{\Phi_i}\ev{\Phi_j}$. This proves the claim for $G^A$. By Lemma~\ref{app:G1:parity}(iii), $\ev{\varphi_a}=0$, which proves the claim for $G^s$.

(b) The graph with $n=0$ is the single line $i$--$j$, of weight $C_{ij}$. Let $\gamma$ be connected with legs $\Phi_i,\Phi_j$ and $n\ge1$. Let $P$ be a simple path from $i$ to $j$. Every separating line lies on $P$. The lines $e_0$ at $i$ and $e_k$ at $j$ are separating and distinct: since $n\ge1$ and $\gamma$ is connected, no line joins the two legs. Order the separating lines along $P$ as $e_0,\dots,e_k$, so $k\ge1$. Every component of $\gamma\setminus\{e_0,\dots,e_k\}$ meets $P$: a shortest path from a vertex to $P$ uses no $e_l$, since every $e_l$ has both ends on $P$. Two different segments of $P$ between consecutive $e_l$ cannot lie in one component. Otherwise $i$ and $j$ would remain connected after the removal of the $e_l$ between them.

So the components are $\{i\},Q_1,\dots,Q_k,\{j\}$, where $Q_l$ contains the segment between $e_{l-1}$ and $e_l$ and has at least one vertex. Take the ends of $e_{l-1}$ and $e_l$ in $Q_l$ as in-leg and out-leg. Then $Q_l$ is connected, and it is irreducible for the following reason. A simple path from $i$ to $j$ crosses every separating line exactly once. Between $e_{l-1}$ and $e_l$ it runs inside $Q_l$, from the in-vertex to the out-vertex. A line of $Q_l$ separating these two vertices would therefore separate $i$ from $j$, but it is not among the $e_l$. By Lemma~\ref{app:G1:parity}(i), every $e_l$ is of the leg type.

Conversely, chaining irreducible pieces by lines of the leg type gives a connected graph whose separating lines are exactly the chaining lines. Indeed, a path from the in-vertex to the out-vertex inside each piece avoids any given line of that piece. Hence $\gamma$ corresponds bijectively to the following data: $k$, an ordered partition $(S_1,\dots,S_k)$ of $\{1,\dots,n\}$ into non-empty blocks, and an irreducible piece on each $S_l$, renumbered increasingly. Moreover
\[
w(\gamma)=\frac{\prod_ln_l!}{n!}\prod_lw(Q_l)\prod_{l=0}^k\E_0[\text{ends of }e_l],\qquad n_l=\abs{S_l},
\]
and there are $n!/\prod_ln_l!$ ordered partitions with block sizes $(n_l)$. Summing, with $\E_0[\Phi_i\Phi_x]=C_{ix}$, $\E_0[\Phi_y\Phi_{x'}]=C_{yx'}$ and $\E_0[\Phi_y\Phi_j]=C_{yj}$, gives $G=C+\sum_{k\ge1}C(KC)^k$. Every piece has $n\ge1$, so $K$ has no term of degree $0$. Hence the series is finite in each order, $\mathbf 1-KC$ is invertible, $G=C(\mathbf 1-KC)^{-1}$, and $G^{-1}=C^{-1}-K$.

(c) Reduction. Whether a piece is irreducible or 1PI depends only on $n$, on the monomials $m_{r_j}$ and on the pairing. The weight is multilinear in the coefficients $g^{(r_j)}$. So $K^\sharp$ does not change if all terms of $V$ with the same monomial are merged. We may therefore assume that distinct labels carry distinct monomials, and write $v_m$ for the coefficient of the monomial $m$.

Hanging parts. Let $Q$ be irreducible but not 1PI, and let $e$ be a line whose removal disconnects $Q$. Since $Q$ is irreducible, both legs lie on one side of $e$, and the other side $F_e$ contains neither the in-vertex nor the out-vertex. Call $F_e$ a hanging part. By Lemma~\ref{app:G1:parity}(ii), $e$ is a photon line.

Two hanging parts $F_e,F_f$ are nested or disjoint. If $f$ has both ends in $F_e$, then the component of the legs in $Q\setminus f$ contains the complement of $F_e$, which is connected in $Q\setminus e$; so $F_f\subset F_e$. If $f$ has both ends outside $F_e$, then $F_e$ together with the outer end of $e$ is connected in $Q\setminus f$; so $F_e\subset F_f$ or $F_e\cap F_f=\emptyset$.

Let $v_0$ be the smallest vertex number lying in some hanging part. The hanging parts containing $v_0$ form a chain, so there is a largest one; call it $H(Q)$, with line $e$ and vertex set $S$. Let $Q'$ arise from $Q$ by replacing the labels on $S$, the pairing among the factors of the vertices of $S$, and the factor of $S$ paired with the outer end $A_b$ of $e$, by any other choice for which the part on $S$ is connected.

We claim that $Q'$ is again an irreducible piece, with the same $v_0$ and with $H(Q')$ on the same set $S$. The set $S$ is attached to the rest of $Q'$ only by $e$, and the rest is connected and contains both legs. Hence:
\begin{itemize}
\item $Q'$ is connected, and no line inside $S\cup\{e\}$ separates the legs;
\item the lines inside $S$ cut off only subsets of $S$;
\item every line outside $S\cup\{e\}$ cuts off the same vertex set as in $Q$.
\end{itemize}
So $Q'$ is irreducible, and the union of the hanging parts is unchanged; in particular $v_0$ is unchanged. A hanging part of $Q'$ strictly containing $S$ would be cut off by a line outside $S$, hence would also be a hanging part of $Q$, which contradicts the maximality of $S$. So $H(Q')$ lies on $S$.

Classes. Thus the irreducible, non-1PI pieces fall into classes. In each class, everything outside $S$ is fixed, the in-leg and out-leg are fixed, and the part on $S$ is arbitrary. The sum of the weights over a class is a common factor times $\sum_{b'}C^A_{bb'}\hat h_{b'}$, where
\[
\hat h_{b'}:=\sum\prod_{j\in S}v_{m_j}\prod_{\text{lines}}\E_0[\cdot],
\]
summed over the connected labelled graphs on the vertex set $S$ in which exactly one factor, a factor $A_{b'}$, is left unpaired. It suffices to show $\hat h\equiv0$.

Symmetry. Write a substitution $\sigma\in\{\tau_z,\iota\}$ as $\Phi_i\mapsto s(i)\Phi_{\varsigma(i)}$, where $\varsigma$ is a bijection of $I$ and $s(i)\in\{\pm1\}$. Here $s\equiv1$ for $\tau_z$, while for $\iota$ we have $s=1$ on scalar indices and $s=-1$ on photon indices. Invariance of $\E_0$ means $s(i)s(i')C_{\varsigma(i)\varsigma(i')}=C_{ii'}$. Invariance of $V$ means $v_{m^\varsigma}=s(m)v_m$, where $m^\varsigma$ is obtained from $m$ by replacing every factor $\Phi_i$ by $\Phi_{\varsigma(i)}$, and $s(m)$ is the product of $s$ over the factors of $m$.

Replace every vertex monomial by its image and keep the pairing of the corresponding factors. This is a bijection between the graphs counted in $\hat h_{b'}$ and those counted in $\hat h_{\varsigma(b')}$, and it preserves connectedness. It multiplies the weight by $\prod_js(m_j)\prod_{\text{lines}}s(i)s(i')$. Every factor except the free one occurs once in each product, so this ratio equals $s(b')$. Hence $\hat h_{\varsigma(b')}=s(b')\hat h_{b'}$. For $\tau_z$ this shows that $\hat h_{(x,\mu)}=:\hat h_\mu$ does not depend on $x$. For $\iota$ it gives $\hat h_{(-x-\eps e_\mu,\mu)}=-\hat h_{(x,\mu)}$. Together, $\hat h_\mu=-\hat h_\mu=0$.

Finally, an irreducible piece without a hanging part is 1PI, since a disconnecting line would either separate the legs or cut off a hanging part. Conversely, every 1PI piece is irreducible.

The models. Every covariance of the form $c(x-y)\delta_{ij}$ (scalar) or $c'(x-y)\delta_{\mu\nu}$ (photon), with $c(-z)=c(z)$ and $c'(-z)=c'(z)$, is $\tau_z$-invariant. It is also $\iota$-invariant: the shifts $\eps e_\mu$ cancel in $(-x-\eps e_\mu)-(-y-\eps e_\mu)$, and the photon signs enter squared. The kernels of $(-\Delta_\eps+M^2)^{-1}$, of $(-\Delta_\eps+\mu_*^2)^{-1}$ and of $C_t=(1-t)\hat G_{\overline M}+t\hat G_{\underline M}$ are of this form, because their symbols depend on $p$ only through $\hat p^2$, which is even in $p$. The single-site terms of $V$ depend only on $\abs{\varphi(x)}$. By Lemma~\ref{app:G1:K1form}(c), the hopping part of $V$ is $-\sum_{x,\mu}\eps\,\rRe\big(\overline{\varphi(x)}(u_{x,\mu}-1)\varphi(x+\eps e_\mu)\big)$, with $u_{x,\mu}=\mathrm{e}^{ie\eps A_{x,\mu}}$. Put $y=-x-\eps e_\mu$. Then $\iota$ maps $\rRe\big(\overline{\varphi(x)}(u_{x,\mu}-1)\varphi(x+\eps e_\mu)\big)$ to
\[
\rRe\big(\overline{\varphi(y+\eps e_\mu)}(\bar u_{y,\mu}-1)\varphi(y)\big)=\rRe\big(\overline{\varphi(y)}(u_{y,\mu}-1)\varphi(y+\eps e_\mu)\big),
\]
and $y$ runs over all sites when $x$ does. Translations clearly preserve $V$.

(d) The single counterterm vertex $\frac12c\,\eps^3\varphi_i(x)^2$ has two ordered choices of in-leg and out-leg. So $K^s_{\rm ct}=-2\cdot\frac12c\,\eps^3\mathbf 1=-c\,\eps^3\mathbf 1$. The matrix of $\mathcal G$ is $\eps^3G^s$. By (b),
\[
\mathcal G^{-1}=\eps^{-3}(G^s)^{-1}=\eps^{-3}(C^s)^{-1}-\eps^{-3}K^s=-\Delta_\eps+M^2+c+\Sigma .
\]
If $C^s$ and $V$ are $\tau_z$-invariant, then so are $\ev\cdot$, $G^s$ and $K^s=(C^s)^{-1}-(G^s)^{-1}$. Translation-invariant operators compose to the products of their symbols, which gives the Fourier form.
\end{proof}

\begin{proposition}[Ward identity on the torus]\label{app:K:L16-1}\label{app:G1:wardT}
Let $\mu_*>0$, let $a$ be a real polynomial in $(e,\lambda)$ with $a(0,0)>0$, and put $\overline M:=a(0,0)^{1/2}$. Write $\ip fh:=\sum_b\eps^3f_bh_b$ on bond functions and $B:=-\Delta_\eps+\mu_*^2$, acting componentwise (Lemma~\ref{app:G1:X}). On $T$ consider the probability measure
\[
Z^{-1}\exp\big(-S^{\mathrm m}_a(A,\varphi)-\tfrac12\ip A{BA}\big)\prod_bdA_b\prod_xd^2\varphi(x),
\]
where $S^{\mathrm m}_a$ is the matter action \eqref{eq:matteraction} with uniform mass term $\frac12a\sum_x\eps^3\abs\varphi^2$. Let $\mathcal G$ be the photon covariance, $\ip f{\mathcal Gh}:=\ev{\ip fA\ip hA}-\ev{\ip fA}\ev{\ip hA}$.
\begin{enumerate}
\item[(i)] There is a neighbourhood $U=(-e_0,e_0)\times[0,\lambda_0)$ of $(0,0)$ in $\R\times[0,\infty)$ on which $\mathcal G$ is positive definite and $C^\infty$ in $(e,\lambda)$, one-sided in $\lambda$ at $\lambda=0$.
\item[(ii)] $\Pi:=B-\mathcal G^{-1}$ is translation invariant on $U$, and $\Pi\,d\psi=0$ for every real function $\psi$ on the sites, where $(d\psi)_{x,\mu}:=\eps^{-1}(\psi(x+\eps e_\mu)-\psi(x))$.
\item[(iii)] Let $\Pi^{\alpha\beta}$ be the $e^\alpha\lambda^\beta$ Taylor coefficient of $\Pi$ at $(0,0)$. It equals the $e^\alpha\lambda^\beta$ coefficient of the 1PI two-photon function $\eps^{-3}K^A$, where $K^A$ is taken in Lemma~\ref{app:G1:dyson}(b) for the following expansion: $C^s=(\eps^3(-\Delta_\eps+\overline M^2))^{-1}\otimes\mathbf 1_2$, $C^A=(\eps^3B)^{-1}$, and $V:=S^{\mathrm m}_a-\frac12\sum_x\eps^3\bar\varphi(-\Delta_\eps+\overline M^2)\varphi$. By Lemma~\ref{app:G1:dyson}(c), $K^A$ is the sum of the 1PI two-photon pieces with the counterterm $a-\overline M^2$ inserted.
\item[(iv)] $\sum_\nu\hat\Pi^{\alpha\beta}_{\mu\nu}(q)\tilde q_\nu=0$ for all $q\in T^*$, all $\mu$ and all $(\alpha,\beta)$.
\item[(v)] If $\ell_T/\eps\ge2$, then $\hat\Pi^{\alpha\beta}_{\mu\nu}(q_1e_\nu)=0$ for $q_1:=2\pi/\ell_T$ and all $\mu,\nu$.
\end{enumerate}
\end{proposition}

\begin{proof}
Gauge invariance. Let $\psi$ be real and $y=x+\eps e_\mu$. Under $A\mapsto A+d\psi$ and $\varphi\mapsto\mathrm{e}^{-ie\psi}\varphi$, the expression $\mathrm{e}^{ie\eps A_{x,\mu}}\varphi(y)-\varphi(x)$ becomes $\mathrm{e}^{-ie\psi(x)}(\mathrm{e}^{ie\eps A_{x,\mu}}\varphi(y)-\varphi(x))$, also on the bonds that wrap around $T$; the single-site terms are unchanged. So $S^{\mathrm m}_a(A+d\psi,\mathrm{e}^{-ie\psi}\varphi)=S^{\mathrm m}_a(A,\varphi)$ (compare Lemma~\ref{app:G1:gauge}). The rotation of each $\varphi(x)$ has Jacobian $1$. Hence $Z^{\mathrm m}(A):=\int\exp(-S^{\mathrm m}_a(A,\varphi))\prod_xd^2\varphi(x)$ satisfies $Z^{\mathrm m}(A+d\psi)=Z^{\mathrm m}(A)$.

Regularity. Choose $e_0,\lambda_0\in(0,\infty)$ such that $a\ge\overline M^2/2$ on $U$, and write $u_{x,\mu}:=\mathrm{e}^{ie\eps A_{x,\mu}}$. By Lemma~\ref{app:G1:K1form}(c), the hopping part of $S^{\mathrm m}_a$ equals $\frac12\sum_{x,\mu}\eps\abs{u_{x,\mu}\varphi(x+\eps e_\mu)-\varphi(x)}^2\ge0$, and the quartic term is non-negative for $\lambda\ge0$. Hence $\exp(-S^{\mathrm m}_a)\le\exp(-\frac14\overline M^2\sum_x\eps^3\abs\varphi^2)$ on $U$, and $0<Z^{\mathrm m}(A)\le C<\infty$ uniformly in $A$. Every $(e,\lambda)$-derivative of $P\exp(-S^{\mathrm m}_a-\frac12\ip A{BA})$, with $P$ a polynomial, is a polynomial in $(A,\varphi,\bar\varphi,u,\bar u)$ times this integrand, with coefficients bounded on $U$. Since $\abs u=1$, it is dominated on $U$ by a polynomial in $(A,\varphi)$ times $\exp(-\frac14\overline M^2\sum\eps^3\abs\varphi^2-\frac12\ip A{BA})$. By dominated convergence, every $\ev P$ is $C^\infty$ on $U$.

Taylor series. The integrand equals $\exp(-V)$ times the Gaussian density of $\E_0$ (up to a constant), and $V$ is entire in $e$ and polynomial in $\lambda$. Therefore the Taylor coefficients of $\E_0[P\exp(-V)]$ at $(0,0)$ are the Gaussian moments of the Taylor coefficients of $\exp(-V)$. Taylor expansion is a ring homomorphism, so the Taylor series of $\ev P$ is the formal series $\ev P$ of Lemma~\ref{app:G1:dyson}. The hypotheses of that lemma hold: every term of $V$ has an even number of scalar factors, and its coefficient has no constant term, because $u-1=O(e)$ and $a-\overline M^2$ vanishes at $(0,0)$.

The identity. Fix real $f$ and $\psi$ and put $v:=s\,d\psi$. Substitute $A\mapsto A+v$ in the integrals $\int\ip fAZ^{\mathrm m}(A)\mathrm{e}^{-\frac12\ip A{BA}}dA$ and $\int Z^{\mathrm m}(A)\mathrm{e}^{-\frac12\ip A{BA}}dA$, and use $Z^{\mathrm m}(A+v)=Z^{\mathrm m}(A)$. This shows that neither integral changes when $\mathrm{e}^{-\frac12\ip A{BA}}$ is replaced by $\mathrm{e}^{-\frac12\ip{A+v}{B(A+v)}}$ and $\ip fA$ by $\ip f{A+v}$. For $\abs s\le1$ we have $\ip{A+v}{B(A+v)}\ge\frac12\ip A{BA}-\ip v{Bv}$, so we may differentiate at $s=0$. This gives
\[
\ev{\ip{B\,d\psi}A}=0,\qquad \ip f{d\psi}=\ev{\ip fA\ip{B\,d\psi}A},
\]
hence $\mathcal G\,B\,d\psi=d\psi$.

Positivity. The law of $A$ has the everywhere positive density $Z^{-1}Z^{\mathrm m}(A)\mathrm{e}^{-\frac12\ip A{BA}}$. So for $f\ne0$ the variable $\ip fA$ is not almost surely constant, and $\ip f{\mathcal Gf}>0$. The entries of $\mathcal G$ are smooth, so $\mathcal G^{-1}$ is smooth (Cramer's rule). This proves (i), and $\Pi\,d\psi=B\,d\psi-\mathcal G^{-1}d\psi=0$.

Translation invariance. The measure is invariant under the translations of $T$, so $\mathcal G$ commutes with them; hence so do $\mathcal G^{-1}$ and $\Pi$. This proves (ii).

(iii) The matrix of $\mathcal G$ is $\eps^3G^A$, and $C^A=(\eps^3B)^{-1}$. By the Taylor-series step and Lemma~\ref{app:G1:dyson}(b), the Taylor series of $\mathcal G^{-1}$ is $\eps^{-3}(G^A)^{-1}=B-\eps^{-3}K^A$, so the Taylor series of $\Pi$ is $\eps^{-3}K^A$. The model satisfies the assumptions of Lemma~\ref{app:G1:dyson}(c) (its first case, with $M=\overline M$ and $c=a-\overline M^2$), so $K^A$ is the sum over the 1PI pieces.

(iv) Every Taylor coefficient of $\Pi\,d\psi=0$ vanishes, so $\Pi^{\alpha\beta}d\psi=0$. By linearity this also holds for complex $\psi$. By translation invariance, $(\Pi^{\alpha\beta}h)_{x,\mu}=\sum_{y,\nu}\eps^3\pi_{\mu\nu}(x-y)h_{y,\nu}$. For $q\in T^*$, the function $\psi(y)=\mathrm{e}^{iq\cdot y}$ is well defined on $T$, and $(d\psi)_{y,\nu}=\tilde q_\nu\mathrm{e}^{iq\cdot y}$. Then $(\Pi^{\alpha\beta}d\psi)_{x,\mu}=\mathrm{e}^{iq\cdot x}\sum_\nu\hat\Pi^{\alpha\beta}_{\mu\nu}(q)\tilde q_\nu$, which vanishes.

(v) Take $q=q_1e_\nu$. Then $\tilde q_\rho=0$ for $\rho\ne\nu$, and $\tilde q_\nu\ne0$ because $\eps q_1=2\pi\eps/\ell_T\in(0,\pi]$. So (iv) gives $\hat\Pi^{\alpha\beta}_{\mu\nu}(q_1e_\nu)=0$.
\end{proof}

\noindent At $q=0$ we have $\tilde q=0$, so (iv) imposes no condition on $\hat\Pi^{\alpha\beta}(0)$. Nor does the identity $\Pi\,d\psi=0$: a constant bond function $\mathbf 1_\mu$ satisfies $\ip{\mathbf 1_\mu}{d\psi}=\ip{d^*\mathbf 1_\mu}\psi=0$, so it is not of the form $d\psi$.

\paragraph{Momentum representation.}
Fix $\eps$, $\overline M>0$, $\mu_*>0$ and a real polynomial $\delta m^2_\infty$ in $(e,\lambda)$ vanishing at $(0,0)$. Consider the vertices $V$ of Proposition~\ref{app:G1:wardT}(iii) with $a=\overline M^2+\delta m^2_\infty$: the hopping vertices of order $e^n$, $n\ge1$, the vertex $\lambda\abs\varphi^4$ and the mass vertex $\frac12\delta m^2_\infty\abs\varphi^2$. Each term of $V$ is a coefficient times a product of variables at the sites $x+\eps w_f$, with $w_f\in\{0,e_1,e_2,e_3\}$, summed over a base site $x$; the variable $A_{x,\mu}$ is attached to the site $x$.

A class $\gamma$ of order $(\alpha,\beta)$ consists of the following data:
\begin{itemize}
\item $n$, and the types of the $n$ vertices (term of $V$, including the chosen monomial in $(e,\lambda)$ of its coefficient);
\item an in-leg factor $A_\mu$ and an out-leg factor $A_\nu$;
\item a pairing of the remaining factors, such that the resulting piece is 1PI.
\end{itemize}
These data do not depend on the torus. On $T$, the labelled 1PI pieces of order $(\alpha,\beta)$ correspond bijectively to the pairs (class, base sites $x_1,\dots,x_n$).

Write each line, joining factors at the sites $y$ and $y'$, as $\abs T^{-1}\sum_{k\in T^*}\mathrm{e}^{ik\cdot(y-y')}\hat c(k)$ times $\delta_{ij}$, resp.\ $\delta_{\mu\nu}$, with $\hat c(k)=(\hat k^2+\overline M^2)^{-1}$ for a scalar line and $\hat c(k)=(\hat k^2+\mu_*^2)^{-1}$ for a photon line. Average the leg site over $T$ using translation invariance, and sum freely over the base sites. Each site sum gives $\eps^{-3}\abs T$ times a Kronecker constraint on the momenta at the vertex. The graph is connected and the external momenta are $q$ and $-q$, so the constraints are dependent; choosing a spanning tree, the momenta of its lines are determined by $q$ and by the momenta of the remaining $L_\gamma=N_\gamma-n+1$ lines, where $N_\gamma$ is the number of lines. This gives the finite sum
\[
\hat\Pi^{\alpha\beta,T}(q)=\sum_\gamma c_\gamma\abs T^{-L_\gamma}\sum_{k\in(T^*)^{L_\gamma}}I_\gamma(k,q),\qquad q\in T^*.
\]
Here $\hat\Pi^{\alpha\beta,T}$ is the coefficient of Proposition~\ref{app:G1:wardT}(iii) on $T$ with $a=\overline M^2+\delta m^2_\infty$. The integrand $I_\gamma$ is the product of the line symbols at the line momenta (integer combinations of $k$ and $q$) and of the phases $\mathrm{e}^{i\eps\kappa\cdot w_f}$. The number $c_\gamma$ (the sign, $1/n!$, the coefficients and the powers of $\eps$) depends on $\gamma$, $\eps$ and the coefficients of $\delta m^2_\infty$, but not on $T$. The 1PI two-photon coefficient on $\eps\Z^3$ is, by definition, the same sum with the weights of the graphs and their integrands integrated over the Brillouin zone ${\rm BZ}=[-\pi/\eps,\pi/\eps)^3$:
\[
\Pi^\infty(q):=\sum_\gamma c_\gamma\int_{{\rm BZ}^{L_\gamma}}I_\gamma(k,q)\,\frac{d^{3L_\gamma}k}{(2\pi)^{3L_\gamma}},\qquad q\in\R^3 .
\]

\begin{corollary}[Ward identity on $\eps\Z^3$]\label{app:K:L16-2}\label{app:G1:CN-wardZ}
In the setting of the preceding paragraph (scalar lines $\hat G_{\overline M}$, photon lines $C_{\mu_*}$, the vertices of Lemma~\ref{app:G1:WP-q2a} and the mass counterterm $\frac12\delta m^2_\infty\sum_x\eps^3\abs\varphi^2$), for every $(\alpha,\beta)$ we have $\sum_\nu\Pi^\infty_{\mu\nu}(q)\tilde q_\nu=0$ for all $q\in\R^3$ and all $\mu$. In particular $\Pi^\infty_{\mu\nu}(0)=0$.
\end{corollary}

\begin{proof}
Let $\ell_n:=2^n\ell_0$ with $\ell_0/\eps\in\N$, and let $T_n$ be the torus of side $\ell_n$. Put $D:=\bigcup_n(2\pi/\ell_n)\Z^3$; every $q\in D$ lies in $T_n^*$ for all large $n$.

All denominators are at least $\min(\overline M^2,\mu_*^2)>0$. So $I_\gamma$ is continuous in $(k,q)$ and $(2\pi/\eps)$-periodic in each loop momentum. The points of $(T_n^*)^{L}$ form a grid of mesh $2\pi/\ell_n$ on the period cell, and $(2\pi)^{-3L}(2\pi/\ell_n)^{3L}=\abs{T_n}^{-L}$. Hence the Riemann sums $\abs{T_n}^{-L}\sum_{(T_n^*)^L}I_\gamma(\cdot,q)$ converge to $\int_{{\rm BZ}^L}I_\gamma(k,q)(2\pi)^{-3L}d^{3L}k$, by uniform continuity. By the momentum representation, $\hat\Pi^{\alpha\beta,T_n}(q)\to\Pi^\infty(q)$ for every $q\in D$.

By Proposition~\ref{app:G1:wardT}(iv), applied with the polynomial $a=\overline M^2+\delta m^2_\infty$, for which $a(0,0)=\overline M^2>0$, we have $\sum_\nu\hat\Pi^{\alpha\beta,T_n}_{\mu\nu}(q)\tilde q_\nu=0$ for $q\in T_n^*$. Hence $\sum_\nu\Pi^\infty_{\mu\nu}(q)\tilde q_\nu=0$ on $D$.

The set $D$ is dense in $\R^3$. $\Pi^\infty$ is continuous, as a finite sum of integrals over a compact set of integrands continuous in $(k,q)$, and $q\mapsto\tilde q$ is continuous. So the identity holds for all $q\in\R^3$. For $q=se_\nu$ with $0<\abs s<2\pi/\eps$ we have $\tilde q_\nu\ne0$ and $\tilde q_\rho=0$ for $\rho\ne\nu$, so $\Pi^\infty_{\mu\nu}(se_\nu)=0$. Letting $s\to0$ gives $\Pi^\infty_{\mu\nu}(0)=0$.
\end{proof}

\subsection{Zero-background vacuum families, finite vertices at the top scale, and torus partners}\label{app:G1:ssec-V0}

\emph{Conventions.} In this subsection the hypotheses of Theorem~\ref{thm:T6} hold and $T=T_j$. As in Subsection~\ref{ssec:UV-torus}, $L=2$ and $\eps=2^{-K}\eps_0$, so $t_K=L^K\eps=\eps_0$. The indices $j,k$ are scale indices. We work in scale-$K$ units. The lattice is $\eta_K\Z^3/S\Z^3$, where $\eta_K:=L^{-K}$ and $S:=\ell_T/\eps_0$; $S$ is an integer by Lemma~\ref{lem:sizes}(ii). A kernel $h$ becomes $(X,Y)\mapsto\eps_0h(\eps_0X,\eps_0Y)$. A kernel $g_k$ written in scale-$k$ units becomes $L^{K-k}g_k(L^{K-k}X,L^{K-k}Y)$. By (T.per) \eqref{eq:GIMG} and $\ell_T\ge\ell_0$,
\begin{equation}\label{app:G1:V0-S}
S\ge\ell_0/\eps_0\ge\max\{6\sqrt3\,r(\eps_0)+2,\ 2/\delta_0\}.
\end{equation}
Write $f\in\{s,v\}$, with $M_s:=M\in\{m_1,m_2\}$ and $M_v:=\mu_*$; the photon is taken componentwise (Lemma~\ref{app:G1:VSR-V0}). For $1\le k\le K$ put $G_{k,f}:=(-\Delta^{\eta_k}+M_f^2t_k^2+a_kP_k)^{-1}$ in scale-$k$ units, and put $G_{0,f}:=0$. Whenever these kernels and their derivatives with at most one lattice difference in each of the two variables are assumed to satisfy (H$_\infty$) of Subsection~\ref{ssec:T6-IMG} with the constants $C_0,\delta_0$ given there, Lemma~\ref{lem:IMG}(i)--(iii) hold for them. Indeed, the proofs of (i)--(iii) use only (H$_\infty$) and the invariance under $\Gamma_k$.

For a kernel $h$ and sets $A,A'$, let $\norm h_{A\times A'}$ be the maximum over $x\in A$, $y\in A'$ of the following quantities: $\abs h$, its first differences in $x$ and in $y$, its mixed second difference, and the $\alpha_0$-H\"older quotients of the first and mixed differences.

\begin{lemma}[torus counterterms in class form]\label{app:K:L17-3}\label{app:G1:IMGv}
Assume (H$_\infty$) of Subsection~\ref{ssec:T6-IMG}, the condition (E.r) ($\eps_0<\mathrm{e}^{1-\varrho}$) and (T.per) \eqref{eq:GIMG}. Let $k\le K$, and let $\square\subset T_{\eta_k}$ be a cube of side $s\le3r(t_k)$ in scale-$k$ units. Let $S_k$, $\Gamma_k$, $\mathcal K$ and $C_1$ be as in Subsection~\ref{ssec:T6-IMG} and Lemma~\ref{lem:IMG}.
\begin{enumerate}
\item[(a)] There is a cube $\tilde\square\subset\eta_k\Z^3$ on which the projection $\pi$ to the torus is a bijection onto $\square$. For $x,x'\in\square$ with lifts $\tilde x,\tilde x'\in\tilde\square$, and for every $n\in\Gamma_k\setminus\{0\}$, $\abs{\tilde x-\tilde x'-n}\ge\max\{\abs{\tilde x-\tilde x'},S_k/2\}$. In particular, the torus distance of $x,x'$ is $\abs{\tilde x-\tilde x'}$.
\item[(b)] For every $\mathsf K\in\mathcal K$ and $x,x'\in\square$, $\mathsf K_T(x,x')=\mathsf K_\infty(\tilde x,\tilde x')+\delta\mathsf K(x,x')$, with $\norm{\delta\mathsf K}_{1,\alpha;\square\times\square}\le C_1\mathrm{e}^{-\delta_0S_k/2}\le C_1\mathrm{e}^{-3\delta_0r(t_k)}$.
\item[(c)] Transform a counterterm graph of \cite[(1.26)--(1.31)]{Ba3} in the way the classes are transformed in \cite[pp.~433, 445]{Ba3}: localise its vertices, replace the far lines by external fields, and regroup. Let $\square$ be the cube of side $3r(t_k)$ of \cite[p.~433]{Ba3} attached to one of its connected components. Insert (b) for every scalar line of that component and expand the product. Then:
\begin{itemize}
\item under the identification $\square\cong\tilde\square$, the term in which every scalar line is $G_k(0)$ is the counterterm in the form used in \cite[(3.7)--(3.9), (3.19), pp.~444--445]{Ba3};
\item every other term is an expression $E(G'_{\rm ren},\{\square(v)\},\Phi_{\rm ext},A_{\rm ext})$ of the type in \cite[(3.3)]{Ba3}, in which $a\ge1$ of the external scalar fields are restrictions of $\delta G_k$ to $\square\times\square$. Its modulus is therefore at most $O(1)(C_1\mathrm{e}^{-3\delta_0r(t_k)})^a$ times the norms of the remaining external fields.
\end{itemize}
\end{enumerate}
In particular, take $k=K$ and $\tilde B=0$, and let $Q$ be a cube of side $\le3r(\eps_0)$ that contains the localisation cubes of a cluster. Each partner of a class supported in $Q$ is then the sum of two parts. The first is its form on $\eta_K\Z^3$. The second consists of at most $2^{m(\nbar)}-1$ expressions of the type in \cite[(3.3)]{Ba3} of the same order, where $m(\nbar)$ is the largest number of lines of a graph of order $\le\nbar$. Each of these expressions contains at least one additional external two-variable field of norm $\le C_1\mathrm{e}^{-3\delta_0r(\eps_0)}$, and every graph in them has at most $2\nbar$ vertices.
\end{lemma}

\begin{proof}
By Lemma~\ref{lem:B:rows}(b), for every $k\le K$,
\[
S_k\ge\max\{6\sqrt3\,r(t_k)+2,\ 2/\delta_0\}.
\]
(a) Choose $\tilde\square$ to be a lift of $\square$ such that $\abs{\tilde x_\mu-\tilde x'_\mu}\le s$ for all $\tilde x,\tilde x'\in\tilde\square$. We have $s\le3r(t_k)<S_k/2$. Let $n=\sum_\mu n_\mu e_\mu\in\Gamma_k$. If $n_\mu\ne0$, then $\abs{n_\mu}\ge S_k$, so
\[
\abs{\tilde x_\mu-\tilde x'_\mu-n_\mu}\ge S_k-s>S_k/2>\abs{\tilde x_\mu-\tilde x'_\mu}.
\]
Coordinates with $n_\mu=0$ are unchanged. Hence, for $n\ne0$, $\abs{\tilde x-\tilde x'-n}\ge\abs{\tilde x-\tilde x'}$ and $\abs{\tilde x-\tilde x'-n}\ge S_k/2$. If $\pi\tilde x=\pi\tilde x'$, then $\tilde x-\tilde x'\in\Gamma_k$ has coordinates of modulus $\le s<S_k$, so $\tilde x=\tilde x'$. The torus distance is $\min_n\abs{\tilde x-\tilde x'-n}$, and this minimum is attained at $n=0$.

(b) By Lemma~\ref{lem:IMG}(i), $\mathsf K_T(x,x')=\sum_{n\in\Gamma_k}\mathsf K_\infty(\tilde x,\tilde x'+n)$. Hence $\delta\mathsf K=\sum_{n\ne0}\mathsf K_\infty(\tilde x,\tilde x'+n)$. The cube $\tilde\square$ has diameter $D=\sqrt3\,s\le3\sqrt3\,r(t_k)$, and $S_k\ge\max\{2D,2/\delta_0\}$. So Lemma~\ref{lem:IMG}(iii) gives $\norm{\delta\mathsf K}_{1,\alpha}\le C_1\mathrm{e}^{-\delta_0S_k/2}$. The derivatives in the norm are themselves kernels of $\mathcal K$, and they are periodised by Lemma~\ref{lem:IMG}(i). The bound therefore needs only points of $\square$; no margin around $\square$ is required. Finally, $\delta_0S_k/2\ge\delta_0(3\sqrt3\,r(t_k)+1)\ge3\delta_0r(t_k)$.

(c) We argue by induction on the order of the counterterm graph. A counterterm of lower order inserted in it is treated first, by the same statement. This is possible because the insertions of \cite[p.~417]{Ba3} have smaller order.

At zero background the counterterms are c-numbers. Their scalar lines are the torus kernels $G_k(T_{\eta_k},0)$ and their slices, all of which belong to $\mathcal K$. In the class, \cite[p.~433]{Ba3} makes three substitutions, in this order:
\begin{itemize}
\item $G_k(\Omega,\tilde B)=G_k(\square,\tilde B)+\delta G_k(\Omega,\square,\tilde B)$;
\item the expansion in $\tilde B'=\tilde B-B_0$, and the gauge transformation that removes $B_0$;
\item $G_k(\square,0)=G_k(0)+\delta G_k(\square,\eta_k\Z^3,0)$.
\end{itemize}
The last substitution regards $\square$ as a subset of $\eta_k\Z^3$. By (a), this is the identification $\square\cong\tilde\square$. The identification preserves distances. It also preserves the nearest-neighbour differences at the vertices, which lie in the central cube $\square_1$ of side $r(t_k)$, at distance $r(t_k)$ from $\partial\square$. Finally, it preserves the values of the external fields, since these are functions on $T_{\eta_k}$ composed with $\pi$.

\emph{The term with only $G_k(0)$.} In the counterterm, (b) replaces each near scalar line by $G_k(0)(\tilde x,\tilde x')+\delta G_k(x,x')$. Taking $G_k(0)$ on every near scalar line gives the near scalar lines of the class, evaluated at the same lifted points. The substitution does not change the vertices or the far lines; the far lines became external fields in the same way on both sides. Vector lines are not substituted: they are the same periodic kernels in the class and in the counterterm. Hence this term is exactly the counterterm in the form of \cite[(3.7)--(3.9), (3.19), pp.~444--445]{Ba3}.

\emph{The other terms.} There are at most $2^{m(\nbar)}-1$ of them. Fix one, and let $A\neq\emptyset$ be the set of lines that carry $\delta G_k$. Drop the lines in $A$ and put the two-variable external fields $\delta G_k|_{\square\times\square}$ in their place. This is the operation that \cite[p.~433]{Ba3} performs with $\delta G_k(\square,\eta_k\Z^3,0)$, on the same lines and with kernels on the same set $\square\times\square$. The graphs obtained, and their decomposition into renormalised classes, depend only on which lines are dropped, and not on the values of the fields \cite[p.~432]{Ba3}. Hence the term is an expression of the type in \cite[(3.3)]{Ba3}. That estimate holds for arbitrary external fields of finite $(1,\alpha_0)$-norm and is multilinear in their norms. By (b), the term is therefore bounded by $O(1)(C_1\mathrm{e}^{-3\delta_0r(t_k)})^{\abs A}$ times the other norms.

\emph{The final statement.} Apply (a)--(c) at $k=K$ with the cube $Q$; parts (a) and (b) need only side $\le3r(\eps_0)$. At $\tilde B=0$ the expansion in the background and the gauge transformation are void, since $B_0=0$ and $\tilde B'=0$. The number of terms was counted in (c). Every vertex other than a finite mass vertex has order at least $\frac12$, and every finite mass vertex has order at least $1$. The vertices of a partner graph of order $o$ carry its total order $o$, so they number at most $2o$. Hence every graph has at most $2\nbar$ vertices.
\end{proof}

\begin{lemma}[finite mass vertices at the top scale]\label{app:K:L17-2}\label{app:G1:Ffin}
Assume the hypotheses of Theorem~\ref{thm:T6}. Let $M\in\{m_1,m_2\}$, $\eps=2^{-K}\eps_0$ and $T=T_j$. Put $C_f:=(-\Delta_\eps+M_f^2)^{-1}$, with $M_s=M$ and $M_v=\mu_*$, and put $R_f:=C_f-G_{K,f}(0)$. Following \cite[(1.26), p.~418]{Ba3}, let $\delta m^2_{\rm fin}(x)$ be the part of the mass counterterm contained in $a_{K,T}(M)$ (Proposition~\ref{prop:UV-B1e}) that is given by the graphs with at least one line $R_f$. Write $\delta m^2_{\rm fin}(x)=\sum_{\alpha,\beta}e^\alpha\lambda^\beta c_{\alpha\beta}(x)$, and put $e(t):=e\,t^{1/2}$ and $\lambda(t):=\lambda t$.
\begin{enumerate}
\item[(a)] $c_{\alpha\beta}\ne0$ only if $\alpha$ is even and $o:=\alpha/2+\beta\in\{1,2\}$.
\item[(b)] For every $x\in T$, $\eps_0^2\,\abs{e^\alpha\lambda^\beta c_{\alpha\beta}(x)}\le C_{\rm fin}\,e(t_K)^\alpha\lambda(t_K)^\beta$. Here, with $\rho_0:=\min(M\eps_0,\mu_*\eps_0)$ and $\bar Z:=2(M\eps_0)^{-1}\rho_0^{-2}(\eps_0/\ell_0)^3$,
\[
C_{\rm fin}:=5\Big[C_b\Big(\tfrac{L^2}{L-1}+2+\log_L\tfrac1{M\eps_0}\Big)+(L+1)\,C'\Big(1+\rho_0^{-2}+\big(\tfrac{8}{\min(\rho_0,\kappa_G)}\big)^3\Big)^5(1+\bar Z)^5\Big].
\]
The constants $C_b$ and $C'=C(\alpha_0,\nbar,d,L,a,M_b)$ are those of Proposition~\ref{app:G1:I3b}, and $\kappa_G$ is the decay rate of Proposition~\ref{app:G1:VSR-R}. The constant $C_{\rm fin}$ does not depend on $\eps$, $x$, $u$, $\kappa$ or $\ell_T\ge\ell_0$.
\end{enumerate}
\end{lemma}

\begin{proof}
(a) The charge $e$ enters only through two kinds of vertex: the vertices \cite[(1.8)]{Ba3}, with one vector leg and factor $e$, and \cite[(1.10)]{Ba3}, with two vector legs and factor $e^2$. The counterterm graphs are scalar self-energies, so all vector legs are internal and are paired by vector lines. Hence $\alpha$ equals the number of vector legs, and it is even. By the definition of the counterterm (\cite[(1.23)]{Ba3}; Lemma~\ref{app:G1:WP-fam}), the graphs have order $2\le\alpha+2\beta\le4$. Hence $o\in\{1,2\}$.

(b) \emph{Step 1 (the hypotheses).} By Lemma~\ref{lem:T6-rows}(b) and $M\le m_2$, we have $\eps_0\le\min(1/m_2,1/\mu_*,L)\le\bar\ell(M):=\min(1/M,1/\mu_*,L)$. Moreover $\ell_T\ge\ell_0\ge L^3$. So \eqref{app:G1:eq-hyp} holds (Lemma~\ref{app:G1:N9prem}), and Proposition~\ref{app:G1:I3b} applies at this $M$. Its hypotheses and constants do not involve $(e,\lambda)$. Nor do the coefficients $c_{\alpha\beta}(x)$, which are built from free lines. Let $\bar K$, $\rho_{\min}$ and $Z_f=t_{\bar K}/(M_f^2\ell_T^3)$ be as in Proposition~\ref{app:G1:I3b}. Write $P_o$ for the part of order $o$. By (a), $\delta m^2_{\rm fin}=P_1+P_2$, and for all $e,\lambda>0$,
\[
\abs{P_1+P_2}\le(e^2+\lambda)A_1+(e^4+e^2\lambda+\lambda^2)B_2,
\]
where
\[
A_1:=C_b\tfrac{L^2}{(L-1)\eps_0}+C_b^R\max(LM,L\mu_*,1)(1+Z_s+Z_v)^2,\qquad
B_2:=C_b\big(2+\log_L^+(\bar\ell/\eps_0)\big)+C_b^R(1+Z_s+Z_v)^5 .
\]
Here $C_b^R=C'\big(1+\rho_{\min}^{-2}+(8/\min(\rho_{\min},\kappa_G))^3\big)^5$.

\emph{Step 2 (order by order).} Replace $(e,\lambda)$ by $(s^{1/2}e,s\lambda)$; this multiplies $P_o$ by $s^o$.
\begin{itemize}
\item Dividing by $s$ and letting $s\downarrow0$ gives $\abs{P_1}\le(e^2+\lambda)A_1$.
\item Dividing by $s^2$ and letting $s\to\infty$ gives $\abs{P_2}\le(e^4+e^2\lambda+\lambda^2)B_2$.
\item Letting $\lambda\downarrow0$ and then $e\downarrow0$ gives $\abs{c_{20}},\abs{c_{01}}\le A_1$ and $\abs{c_{40}},\abs{c_{02}}\le B_2$.
\item Taking $e=\lambda=1$ then gives $\abs{c_{21}}\le3B_2+2B_2=5B_2$.
\end{itemize}

\emph{Step 3 (the constants).}
\begin{itemize}
\item Since $\bar K\ge K$, we have $t_{\bar K}\ge t_K=\eps_0$, so $\rho_{\min}\ge\rho_0$. The constant $C_b^R$ is non-increasing in $\rho_{\min}$ (Remark~\ref{app:G1:N9-supply}(iii)), so $C_b^R\le C'\big(1+\rho_0^{-2}+(8/\min(\rho_0,\kappa_G))^3\big)^5$.
\item Since $t_{\bar K}\le\bar\ell\le1/M$ and $\ell_T\ge\ell_0$, we get $Z_f\le(M\eps_0)^{-1}(M_f\eps_0)^{-2}(\eps_0/\ell_0)^3$, so $Z_s+Z_v\le\bar Z$.
\item $\bar\ell/\eps_0\le1/(M\eps_0)$, and $M\eps_0\le1$.
\item $\eps_0\max(LM,L\mu_*,1)\le L$, because $M\eps_0,\mu_*\eps_0\le1$ and $\eps_0=2^{-n_0}\le1$ by (T.div).
\end{itemize}

\emph{Step 4 (rescaling).} By \cite[(1.29)--(1.30)]{Ba3}, the term $\frac12\eps^3\delta m^2\abs{\varphi(x)}^2$ becomes $\frac12\eta_K^3\,t_K^2\delta m^2\abs{\varphi'(x')}^2$. So the scale-$K$ coefficient is $\eps_0^2e^\alpha\lambda^\beta c_{\alpha\beta}=\eps_0^{2-o}c_{\alpha\beta}\,e(t_K)^\alpha\lambda(t_K)^\beta$.
\begin{itemize}
\item For $o=1$: by Steps 2 and 3, $\eps_0\abs{c_{\alpha\beta}}\le\eps_0A_1\le C_bL^2/(L-1)+L\,C_b^R(1+\bar Z)^5$.
\item For $o=2$: $\abs{c_{\alpha\beta}}\le5B_2\le5[C_b(2+\log_L(1/(M\eps_0)))+C_b^R(1+\bar Z)^5]$.
\end{itemize}
By the bound on $C_b^R$ in Step 3, both are at most $C_{\rm fin}$. None of these constants involves $\eps$, $x$, $u$, $\kappa$ or $\ell_T\ge\ell_0$.
\end{proof}

Put $H_{j,f}:=G_{j+1,f}-G_{j,f}$, $R_f:=C_f-G_{K,f}$, $\rho_s:=M\eps_0$ and $\rho_v:=\mu_*\eps_0$. Both $\rho_f$ lie in $(0,1]$ by (E.fin). Put also $\kappa_{R,f}:=\min(\rho_f,\kappa_G)/8$ and $\zeta_f:=(\rho_f^2S^3)^{-1}\le\rho_f^{-2}(\eps_0/\ell_0)^3$. Since $\delta_j=\kappa\mu_*^2t_{j+1}^2$,
\begin{equation}\label{app:G1:V0-Gamma}
\Gamma^B=\sum_{j<K}\delta_jH_{j,v}=c_K\sum_{j<K}L^{-2(K-j)}H_{j,v},\qquad c_K:=\kappa\mu_*^2L^2t_K^2 .
\end{equation}

\begin{lemma}[the lines]\label{app:G1:V0-lines}
Assume (H$_\infty$) for the kernels $G_{k,f}$, $1\le k\le K$, and their derivatives of order at most $2$. Let $\Delta,\Delta'$ be unit cubes of $T$, and put $c_L:=\sum_{n\ge0}L^{4n}\mathrm{e}^{-(L^n-1)}$, so that $c_2<24$. There is $C_*=C_*(a,L,M_b,\delta_0)$ such that:
\begin{itemize}
\item[(a)] if $\dist(\Delta,\Delta')\ge1$, then $\norm{G_{K,f}}_{\Delta\times\Delta'}\le C_*\mathrm{e}^{-\delta_0\dist(\Delta,\Delta')/2}$ and $\norm{\Gamma^B}_{\Delta\times\Delta'}\le C_*c_K\mathrm{e}^{-\delta_0\dist(\Delta,\Delta')/2}$;
\item[(b)] $\norm{R_f}_{\Delta\times\Delta'}\le B^\sharp\big[\mathrm{e}^{-\kappa_{R,f}\dist_1(\Delta,\Delta')}+\zeta_f\big]$, with $B^\sharp=B^\sharp(a,L,\alpha_0)$. Hence
\[
\sup_\Delta\sum_{\Delta'}\norm{R_f}_{\Delta\times\Delta'}\le N_f:=B^\sharp\big[\mathrm{e}^{3\kappa_{R,f}}(1+2/\kappa_{R,f})^3+\rho_f^{-2}\big]\quad\text{and}\quad
\norm{R_f}_{\Delta\times\Delta'}\le B^\sharp(1+\zeta_f);
\]
\item[(c)] let $Q\subset T$ be a cube of Euclidean diameter $D$ with $S\ge2D$. Then $Q$ lifts isometrically to a cube $\tilde Q\subset\eta_K\Z^3$. On $Q\times Q$ we have $G_{K,f}=G^\infty_{K,f}+\delta G_{K,f}$ and $\Gamma^B=\Gamma^{B,\infty}+\delta\Gamma^B$, where $G^\infty_{K,f}$ and $\Gamma^{B,\infty}=c_K\sum_{j<K}L^{-2(K-j)}H^\infty_{j,v}$ are the kernels on $\eta_K\Z^3$ evaluated at the lifts. Moreover $\norm{\delta G_{K,f}}_{Q\times Q}\le C_1$ and $\norm{\delta\Gamma^B}_{Q\times Q}\le2c_LC_1c_K$.
\end{itemize}
\end{lemma}

\begin{proof}
(a) By Lemma~\ref{lem:IMG}(i), $G_{K,f}(x,x')=\sum_nG^\infty_{K,f}(\tilde x,\tilde x'+nS)$, and the same holds for every derivative. Choose the lifts so that $\abs{\tilde x_\mu-\tilde x'_\mu}\le S/2$ for each $\mu$, and put $\Delta x:=\tilde x-\tilde x'$. For $n\ne0$ and $n_\mu\neq0$, $\abs{\Delta x_\mu-n_\mu S}\ge\abs{n_\mu}S/2\ge\abs{\Delta x_\mu}$. Hence $\abs{\Delta x-nS}\ge\abs{\Delta x}\ge\dist(\Delta,\Delta')\ge1$ and $\abs{\Delta x-nS}\ge\abs nS/2$, so that $\abs{\Delta x-nS}\ge\frac12\abs{\Delta x}+\frac14\abs nS$. By (H$_\infty$),
\[
\norm{G_{K,f}}_{\Delta\times\Delta'}\le C_0\mathrm{e}^{-\delta_0\dist(\Delta,\Delta')/2}\sum_n\mathrm{e}^{-\delta_0\abs nS/4},
\]
and the last sum is at most $\sum_n\mathrm{e}^{-\abs n/2}$ because $\delta_0S\ge2$ by \eqref{app:G1:V0-S}.

For $\Gamma^B$, write $H_{j,v}=G_{j+1,v}-G_{j,v}$ and treat each $G_{k,v}$, $1\le k\le K$, in its own scale-$k$ units. There distances are multiplied by $L^{K-k}$, and the torus side is $S_k=L^{K-k}S$. Going back to scale-$K$ units:
\begin{itemize}
\item the kernel is multiplied by $L^{K-k}$;
\item each difference costs a factor $L^{K-k}$;
\item a H\"older quotient costs a factor $L^{(K-k)\alpha_0}$.
\end{itemize}
So the norm costs at most $L^{4(K-k)}$. Each $G_{k,v}$ occurs at most twice in \eqref{app:G1:V0-Gamma}, with $L^{-2(K-j)}\le1$. Applying the previous estimate at scale $k$, with $C:=C_0\sum_n\mathrm{e}^{-\abs n/2}$ and $\dist:=\dist(\Delta,\Delta')\ge1$, we get
\[
\norm{\Gamma^B}_{\Delta\times\Delta'}\le2c_KC\sum_{n\ge0}L^{4n}\mathrm{e}^{-\delta_0L^n\dist/2}\le2c_KC\,\mathrm{e}^{-\delta_0\dist/2}\sum_{n\ge0}L^{4n}\mathrm{e}^{-\delta_0(L^n-1)/2}.
\]

(b) At the index $K$ the torus side in Proposition~\ref{app:G1:VSR-R} is $S$ and the mass parameter is $\rho_f$. The pointwise bound and the derivative and H\"older bounds of Proposition~\ref{app:G1:VSR-R} give the first claim, with $B^\sharp$ the largest of their constants. Unit cubes whose offset is $n\in\Z^3$ satisfy $\dist_1\ge\abs n_1-3$. Moreover $\sum_{n\in\Z^3}\mathrm{e}^{-\kappa\abs n_1}=\coth(\kappa/2)^3\le(1+2/\kappa)^3$. Finally, there are $S^3$ unit cubes and $S^3\zeta_f=\rho_f^{-2}$.

(c) The side of $Q$ is less than $S/2$, so its projection is injective, and for $x,x'\in Q$ the nearest image is $n=0$; hence the lift is isometric. For $1\le k\le K$, the lift of $Q$ in scale-$k$ units has diameter $L^{K-k}D\le S_k/2$, and $S_k\ge S\ge2/\delta_0$. Lemma~\ref{lem:IMG}(iii) therefore bounds the image part of $G_{k,f}$ by $C_1\mathrm{e}^{-\delta_0S_k/2}$ in scale-$k$ units. For $k=K$ this gives $\norm{\delta G_{K,f}}\le C_1$. For $\Gamma^B$, sum as in (a), using $\mathrm{e}^{-\delta_0L^nS/2}\le\mathrm{e}^{-\delta_0S/2}\mathrm{e}^{-(L^n-1)}$; this gives $\norm{\delta\Gamma^B}\le2c_KC_1c_L$. The formula for $\Gamma^{B,\infty}$ follows from \eqref{app:G1:V0-Gamma} and Lemma~\ref{lem:IMG}(i) applied to each $G_{k,v}$. Finally, $c_2=\sum_n16^n\mathrm{e}^{-(2^n-1)}<23.4$.
\end{proof}

\begin{lemma}[degrees]\label{app:G1:V0-deg}
Let $G$ be a connected graph at zero background. Its vertices are those of \cite[(1.6)--(1.11)]{Ba3}, among them vertices (1.7) that carry a partner, and $V_f$ finite mass vertices. Its lines, $m$ of which are $\Gamma^B$-lines, are $C$-lines or $\Gamma^B$-lines. Let $E=E_s+E_v$ be its number of external legs, $E_s$ of them scalar and $E_v$ vector, $E_\partial$ the number of differentiations on external legs, and $o'$ the order carried by the vertices other than the finite mass vertices. Let $D(G)$ be the degree of \cite[(2.1)--(2.3)]{Ba3}, and $D^\natural(G)$ the degree with the $\Gamma^B$-lines counted as in Lemma~\ref{app:G1:WP-deg}. If $G$ contains no vertex (1.9), (1.11), (1.13)--(1.15), then
\[
D^\natural(G)=o'+2V_f+\tfrac12E+E_\partial-3+2m .
\]
Consequently:
\begin{itemize}
\item[(a)] if $G$ is a vacuum graph and $m\ge2$, then $D^\natural(G)\ge3m-3\ge3$;
\item[(b)] if $G$ has an external leg and $V_f+m\ge1$, then $D^\natural(G)>0$. The only exceptions are a bare vertex, and the one-loop graph through one finite mass vertex and one vertex (1.8) with $n=1$ with a single external photon leg; that graph vanishes \cite[p.~434]{Ba3};
\item[(c)] if $l$ is a $C$-line of $G$ that is not a bridge, then $D(G\setminus l)=D(G)+1+(\text{number of differentiations on }l)$;
\item[(d)] if $G$ is divergent, has an external leg, is not one of the graphs whose value vanishes by \cite[p.~434]{Ba3}, and $l$ is a $C$-line of $G$ that is not a bridge, then $D(G\setminus l)\ge0$. Equality holds only if $G\setminus l$ is the tree \cite[(2.4)]{Ba3} or a bare vertex.
\end{itemize}
Graphs that contain a vertex (1.9), (1.11), (1.13)--(1.15) have $D>0$ \cite[Cor.~2.3]{Ba3}.
\end{lemma}

\begin{proof}
By \cite[(2.1)--(2.3)]{Ba3}, an external leg contributes $0$ to $D_G(v)$ instead of $-\frac12$, and a differentiation on an external leg contributes $0$ instead of $-1$. Hence
\[
D(G)=\sum_vD(v)+\tfrac12E+E_\partial-3+\sum\deg .
\]
Here $D(v)=o(v)$ for the vertices (1.6), (1.8) and (1.10). A vertex (1.7) has $D(v)=2$. If it carries a partner $\delta m^2_{G_0}$, then $D(v)+\deg=2+D(G_0)=o(G_0)$, by induction on the order, since $G_0$ has two scalar legs and $E_\partial(G_0)=0$. If it carries a finite counterterm, then $\deg=0$. A $\Gamma^B$-line adds $+1$ per leg (Lemma~\ref{app:G1:WP-deg}(a)). This proves the identity.

(a) Each $\Gamma^B$-line ends on two photon legs of coupling vertices, and each photon leg carries order $\frac12$. Hence $o'\ge m$, and with $E=E_\partial=0$ the identity gives $D^\natural\ge3m-3$.

(b) If $m\ge1$ and $V_f=0$, this is Lemma~\ref{app:G1:WP-deg}(c). Now let $V_f\ge1$.
\begin{itemize}
\item If $E\ge2$, then $D^\natural\ge o'+2V_f-2\ge0$. Equality forces $o'=0$, $V_f=1$, $E=2$ and $m=0$; then $G$ is a single finite mass vertex with both legs external, that is, a bare vertex. If $o'=0$ and $V_f\ge2$, then $D^\natural\ge2$.
\item If $E=1$, the leg is a photon leg, since $E_s$ is even. Then $o'\ge\frac12$, so $D^\natural\ge o'-\frac12\ge0$. Equality forces $o'=\frac12$, $V_f=1$ and $m=0$, which is the graph named in (b).
\end{itemize}

(c) Cutting $l$ makes its two legs external, each contributing $+\frac12$, and each differentiation on $l$ contributes $+1$. Since $G\setminus l$ is connected, the term $-3$ is unchanged.

(d) Since $G$ has the line $l$, it is not a bare vertex. Since its value does not vanish, (b) and Lemma~\ref{app:G1:WP-deg}(c) give $m=V_f=0$. Let $e^\alpha\lambda^\beta$ be the order of $G$, $o:=\alpha/2+\beta$, and let $I_v$ be the number of internal vector lines of $G$ and of the graphs of its partners. Then $D(G)=o+\frac12E+E_\partial-3$, and $\alpha=2I_v+E_v$, since the partners have scalar legs only. Since $E_s$ is even, the order $o$ is an integer when $E_v$ is even. The case $E=1$ is excluded, because such graphs vanish \cite[p.~434]{Ba3}. If $E=2$, then $o\ge1$; if $E\ge3$, then $o\ge\frac12$. Hence $D(G)\ge-1$. Moreover $D(G)=-1$ forces $E=2$, $o=1$ and $E_\partial=0$: the only other candidate, $E=3$ with $o=\frac12$, is a bare vertex, which has no line. By (c), $D(G\setminus l)\ge D(G)+1\ge0$. Equality forces $D(G)=-1$ and no differentiation on $l$, so $G\setminus l$ has $E=4$, $o=1$ and $E_\partial=0$. Three cases remain.
\begin{itemize}
\item $(E_s,E_v)=(4,0)$. A graph with a line, $(E_s,E_v)=(4,0)$ and $o=1$ consists of two vertices (1.8) with $n=1$ joined by a vector line. Both of its differentiations sit on external legs, so it has $E_\partial=2$. Hence $G\setminus l$ is the bare vertex (1.6).
\item $(E_s,E_v)=(0,4)$. This needs $\alpha\ge4$, that is, $o\ge2$.
\item $(E_s,E_v)=(2,2)$. Then $\alpha=2I_v+2=2$ gives $I_v=0$ and $\beta=0$. So $G\setminus l$ is either the bare vertex (1.10) with $n=2$, or two vertices (1.8) with $n=1$ joined by a scalar line that carries both differentiations, which is the tree \cite[(2.4)]{Ba3}.
\end{itemize}
The last sentence of the lemma is \cite[Cor.~2.3]{Ba3}.
\end{proof}

The proof of the next theorem uses three facts.
\begin{itemize}
\item[(CB)] The class bound: \cite[(3.3)]{Ba3} with its proof \cite[pp.~433--445]{Ba3}, in the coloured form of Theorem~\ref{app:G1:WP-WP1}. Let $G_{\rm ren}$ be a connected renormalised class on $\eta_K\Z^3$, with lines $G^\infty_{K,f}$ and $\Gamma^{B,\infty}$, with partners in the form of \cite[(3.7)--(3.9), (3.19)]{Ba3}, and with finite mass vertices whose coefficients have modulus at most $1$. Let its vertices be localised in unit cubes, and let $\Phi$ be external fields of one or two vertex positions. Suppose $G_{\rm ren}$ is not a vacuum class with at most one $\Gamma^B$-line. Then, after the couplings are extracted, $\abs{E(G_{\rm ren},\{\square(v)\},\Phi)}\le C_B\,c_K^{m(G_{\rm ren})}\prod\norm\Phi$, with $C_B=C_B(\nbar,\alpha_0,L,a)$.
\item[(PS)] The partner splitting: Statement~\ref{app:K:L17-4} at scale $K$ and zero background. Each torus partner equals its form on $\eta_K\Z^3$ plus a finite mass vertex of the same order. The scale-$K$ coefficient of that vertex is at most $C_P$ times the rescaled couplings, with $C_P=C_P(\nbar,\alpha_0,L,a,M_b,\delta_0)$.
\item[(FV)] The finite-vertex bound: by Lemma~\ref{app:G1:Ffin}, each finite mass vertex of order $e^{\alpha_v}\lambda^{\beta_v}$ has scale-$K$ coefficient at most $C_{\rm fin}\,e(t_K)^{\alpha_v}\lambda(t_K)^{\beta_v}$.
\end{itemize}

\begin{theorem}[vacuum families at zero background]\label{app:K:L17-1}\label{app:G1:V0-thm}
Assume the hypotheses of Theorem~\ref{thm:T6}, in particular (E.fin) and (T.per). Assume $\nbar>12$ and $2\nbar+1\le3R$, with $R$ as in \eqref{eq:UV-r}. Let $M\in\{m_1,m_2\}$ and $u\in[-\eta,1]$. For $f\in\{s,v\}$ let $G_{k,f}=(-\Delta^{\eta_k}+M_f^2t_k^2+a_kP_k)^{-1}$ in scale-$k$ units, with $M_s=M$, $M_v=\mu_*$, and the photon taken componentwise. Assume that these kernels and their derivatives of order at most $2$ satisfy (H$_\infty$) of Subsection~\ref{ssec:T6-IMG}, with the constants $C_0,\delta_0$ given there. Let $F$ be a family (Lemma~\ref{app:G1:WP-fam}) of connected coloured vacuum graphs that occur in $T_{\nbar}[\Phi_M(u)]$ on a torus $T=T_j$ of side $\ell_T\ge\ell_0$. Let $F$ have order $e^\alpha\lambda^\beta$, where the orders of the finite counterterm vertices are included, and $m'\ge2$ lines $\Gamma^B$. Let $\operatorname{val}(F)$ be its coefficient of $u^{m'}$. Then
\[
\abs{\operatorname{val}(F)}\le C_F\,e^\alpha\lambda^\beta(\kappa\mu_*^2L^2)^{m'}\eps_0^{\alpha/2+\beta+2m'-3}\abs T,
\]
with $C_F=C_F(\nbar,\alpha_0,L,a,M_b,\delta_0,M\eps_0,\mu_*\eps_0,\ell_0/\eps_0)$ independent of $\eps=2^{-K}\eps_0$, $K$, $u$, $\kappa$ and $\ell_T\ge\ell_0$.
\end{theorem}

\begin{proof}
\emph{Step 0 (rescaling).} The rescaling \cite[(1.29)--(1.30)]{Ba3} preserves values and replaces $e,\lambda$ by $e(t_K)=e\eps_0^{1/2}$ and $\lambda(t_K)=\lambda\eps_0$. Write the vertex sums as sums over unit cubes, and let $\mathcal V(\Delta_0)$ be the sum with the cube $\Delta_0$ of a root vertex fixed. There are $S^3=\abs T/\eps_0^3$ unit cubes, so $\abs{\operatorname{val}(F)}\le S^3\sup_{\Delta_0}\abs{\mathcal V(\Delta_0)}$.

\emph{Step 1 (lines and vertices).} Write each $C$-line as $C_f=G_{K,f}+R_f$ \cite[(1.25)]{Ba3}, keep $\Gamma^B$ as in \eqref{app:G1:V0-Gamma}, and expand multilinearly; this gives at most $2^{m(\nbar)}$ terms. Every graph has at most $2\nbar$ vertices. Indeed, each vertex other than a finite mass vertex has $o(v)\ge\frac12$, each finite mass vertex has order at least $1$, and $\alpha/2+\beta\le\nbar$.

\emph{Step 2 (cuts on $T$).} Fix the localisations. Every $R_f$-line is treated as a two-variable external field. This is legitimate because (CB) is multilinear in the external fields, and Lemma~\ref{app:G1:V0-lines}(b) bounds their norms. Every $G_{K,f}$- or $\Gamma^B$-line joining cubes at distance at least $1$ is also treated as an external field \cite[(3.1), p.~432]{Ba3}; its norm is bounded by Lemma~\ref{app:G1:V0-lines}(a). The remaining lines join cubes at distance $0$. The connected components under these remaining lines are called \emph{clusters}. By \cite[p.~432]{Ba3} and step~(1) of the proof of Theorem~\ref{app:G1:WP-WP1}, the expression is a sum of products of renormalised classes, one product for each decomposition into clusters.

\emph{Step 3 (non-wrapping cubes).} Let a cluster occupy $N\le2\nbar$ unit cubes, each touching another. Adding a touching cube enlarges the bounding box by at most $1$ in each coordinate, so the cluster lies in a cube of side $N$. The vertices have range $\eta_K\le1$, so all points at which the vertices of the cluster couple lie in a cube $Q$ of side $2\nbar+1$, of Euclidean diameter $D=\sqrt3(2\nbar+1)$. We have $r(\eps_0)\ge R$ by \eqref{eq:UV-r}, since $\eps_0\le1$, and so $2\nbar+1\le3R\le3r(\eps_0)$. Then \eqref{app:G1:V0-S} gives $2D\le6\sqrt3\,r(\eps_0)<S$. So Lemma~\ref{app:G1:V0-lines}(c) applies to $Q$.

\emph{Step 4 (lifting).} On each cluster, split every uncut line by Lemma~\ref{app:G1:V0-lines}(c), and split every partner by (PS). The terms $\delta G_{K,f}$ and $\delta\Gamma^B$ are two-variable external fields, as in \cite[p.~433]{Ba3}, and the image parts of the partners are finite mass vertices. At $\tilde B=0$ the background expansion of \cite[p.~433]{Ba3} is void, because $B_0=0$ and $\tilde B'=0$. We regroup as in \cite[p.~433]{Ba3}. Each cluster then contributes a sum of products of connected classes on $\eta_K\Z^3$, with lines $G^\infty_{K,f}$ and $\Gamma^{B,\infty}$, and with external fields coming from cut lines. The number of terms is bounded in terms of $\nbar$ alone.

\emph{Step 5 (degrees).} Consider one resulting component $G'$.
\begin{itemize}
\item If no line is cut on $G'$, then $G'$ is a whole vacuum graph with $m'\ge2$, so $D^\natural\ge3$ by Lemma~\ref{app:G1:V0-deg}(a). In particular it needs no vacuum counterterm.
\item Otherwise $G'$ has an external leg. If it contains a $\Gamma^B$-line or a finite mass vertex, then $D^\natural>0$ by Lemma~\ref{app:G1:V0-deg}(b). The exceptions are bare vertices, which carry no sum, and the graph named there, which vanishes.
\item If $G'$ has an external leg and contains neither, it is a class with $C$-lines only, renormalised by its partners on $\eta_K\Z^3$ (Lemma~\ref{app:G1:WP-fam}). A cut line inside a divergent subgraph is then covered by Lemma~\ref{app:G1:V0-deg}(c),(d) and \cite[(2.8)]{Ba3}. The divergent subgraphs excluded in Lemma~\ref{app:G1:V0-deg}(d) are graphs whose value vanishes by \cite[p.~434]{Ba3}, and the terms that contain them vanish.
\end{itemize}
So (CB) applies to every component.

\emph{Step 6 (one localisation).} Use (CB) for each component, and (FV) and (PS) for the finite mass vertices, after their coefficients are extracted. The number of $\Gamma^B$-lines is $m'$ in every term, whether they stay inside a class, are cut, or are image lines. Since (CB) is multilinear in the external fields, the value at fixed localisations is at most
\[
C_B^{2\nbar}(1+C_{\rm fin}+C_P)^{2\nbar}\,e(t_K)^\alpha\lambda(t_K)^\beta c_K^{m'}\prod_{\text{cut }l}\norm l'_{\Delta_l\times\Delta'_l}.
\]
Here $\norm l'$ is the norm of the line divided by $c_K$ if $l$ is a $\Gamma^B$- or $\delta\Gamma^B$-line, and the norm of the line otherwise. An image line has $\norm l'\le2c_LC_1$ or $\le C_1$.

\emph{Step 7 (sum over localisations).} Fix a decomposition into clusters; there are finitely many, their number depending on $\nbar$ only. Inside a cluster with $N$ cubes, once one cube is chosen, each of the others has at most $(2N+1)^3$ choices. Let $\Gamma$ be the graph whose vertices are the clusters and whose edges are the cut lines between different clusters; it is connected because $G$ is. Choose a spanning tree of $\Gamma$, and sum from its leaves towards the cluster of $\Delta_0$.
\begin{itemize}
\item A tree edge costs $\sup_\Delta\sum_{\Delta'}\norm l'$. By Lemma~\ref{app:G1:V0-lines}(a),(b), this is at most $C_*\sum_{n\in\Z^3}\mathrm{e}^{-\delta_0(\abs n_\infty-1)/2}$ for $G_{K,f}$- and $\Gamma^B$-lines, and at most $N_f$ for $R_f$-lines.
\item Every other cut line costs at most $\max\{C_*,2c_LC_1,B^\sharp(1+\rho_f^{-2}(\eps_0/\ell_0)^3)\}$.
\end{itemize}
There are at most $m(\nbar)$ lines. Hence $\sup_{\Delta_0}\abs{\mathcal V(\Delta_0)}\le C_F\,e(t_K)^\alpha\lambda(t_K)^\beta c_K^{m'}$. The constant $C_F$ depends on $M\eps_0$, $\mu_*\eps_0$ and $\ell_0/\eps_0$ only through $N_f$, $\zeta_f$ and $C_{\rm fin}$, and otherwise only on $\nbar,\alpha_0,L,a,M_b,\delta_0$. It is uniform in $\ell_T\ge\ell_0$ because $\zeta_f\le\rho_f^{-2}(\eps_0/\ell_0)^3$, and it does not involve $\kappa$ or $u$.

\emph{Step 8.} Multiply by $S^3=\abs T\eps_0^{-3}$ and insert $c_K=\kappa\mu_*^2L^2\eps_0^2$, $e(t_K)^\alpha=e^\alpha\eps_0^{\alpha/2}$ and $\lambda(t_K)^\beta=\lambda^\beta\eps_0^\beta$. This gives the claim.
\end{proof}

\section{The negative-source bound and the choice of parameters}\label{app:G1:part-D6}
We prove Proposition~\ref{prop:ch-Tflat} (Subsection~\ref{p:appI:Tflat}), the domination of the infrared part of the massless photon by the slice covariance $\Gamma^B$ (Subsection~\ref{app:G1:part-Tdom}), and the statements on the auxiliary parameters and on the hypotheses of Theorem~\ref{thm:ch-T9} (Subsection~\ref{app:G1:part-D7}).

\subsection{\texorpdfstring{The proof of Proposition~\ref{prop:ch-Tflat}}{The proof of the negative-source bound}}\label{p:appI:Tflat}

Section~\ref{sec:chain} uses three facts that are proved here: the definition of $\Gamma_\ge$ and of the comparison covariance $C^\flat_{B,s}$, with the domination $C_{\mu_0}\le C^\flat_{B,s}$ (Subsection~\ref{ssec:ch-scheme}); the translation-averaging lemma (step~(1) of the proof of Proposition~\ref{prop:ch-Tflat}); and the determinant comparison for the infrared part with the constants $J_{\rm IR}$, $Z_T$ (step~(3) there). With them, Proposition~\ref{app:G1:Tflat} below is the full proof of Proposition~\ref{prop:ch-Tflat}.

\emph{Standing conventions.} Fix $j\ge1$, $T=T_j$, $\ell_T=\ell_j$ (the index $i$ labels King's slices $H_i$ below; $j$ is the torus index), and $\eps=L^{-K}\eps_0$ with $K\ge1$. By Lemma~\ref{lem:sizes}(ii) with $k=K$, $\ell_T/(M_b\eps_0)\in2\N$; hence the blocks of every scale $t_k=L^k\eps$, $0\le k\le K$, tile $T_\eps$, which is what the definitions of $G_k$, $H_i$ and $\Gamma^B$ in Subsection~\ref{ssec:T6-statement} require.
Covariances are operators on real bond functions, with the inner product of Subsection~\ref{ssec:action}. $-\Delta_\eps$ acts on each component; its eigenvalues are $\hat k^2:=\sum_\mu4\eps^{-2}\sin^2(\eps k_\mu/2)$, where $k\in(2\pi/\ell_T)\Z^3$ has components in $[-\pi/\eps,\pi/\eps)$; $k=0$ corresponds to $\operatorname{ran}\Pi_0$.

For a site function $b(\cdot)$ (a \emph{profile}; a constant $b$ is a special case), let $S^{\mathrm m}_{b(\cdot)}$ be the matter action of Definition~\ref{def:algebra} with single-site term $\frac12b(x)\eps^3r_x^2$ at $x$. Put $Z_\varphi(A;b):=\int \mathrm{e}^{-S^{\mathrm m}_{b}(A,\varphi)}\prod_xd^2\varphi(x)$ and $\Psi^b_F(A):=\int F\,\mathrm{e}^{-S^{\mathrm m}_{b}(A,\varphi)}\prod_xd^2\varphi(x)$. For a covariance $\Gamma\ge0$ put
\[
\ev{F}_{\Gamma;b}:=\frac{\E_{N(0,\Gamma)}[\Psi^b_F]}{\E_{N(0,\Gamma)}[Z_\varphi(\cdot;b)]},\qquad Z(\Gamma;b)=\E_{N(0,\Gamma)}[Z_\varphi(\cdot;b)],
\]
as in Subsection~\ref{ssec:T6-statement}. If $\Gamma=\Gamma^\perp+\infty\cdot\Pi_0$ we mean the product of $N(0,\Gamma^\perp)$ on $\operatorname{ran}\Pi^\perp$ with the uniform probability measure on $[0,P)^3$ for $c$; we call this a \emph{Haar zero mode}. Thus $\ev{\cdot}_{\eps,T,\mu_0}=\ev{\cdot}_{C_{\mu_0};b_T}$ and $\ev{\cdot}_{\eps,T,\fl}=\ev{\cdot}_{(-\Delta_\eps)^{-1}\Pi^\perp+\infty\Pi_0;\,b_T}$ (Definition~\ref{def:states}). Finally put $\rho_\Gamma(y;b):=\ev{r_y^2-W_\eps-z_\eps}_{\Gamma;b}$ and $\bar\rho_\Gamma(b):=\abs T^{-1}\sum_y\eps^3\rho_\Gamma(y;b)$.

\begin{definition}[infrared part and slice-modified comparison covariance]\label{app:G1:Tflatdef}
Let
\[
D^<:=\int_0^{\eps_0^2}\mathrm{e}^{t\Delta_\eps}(1-\mathrm{e}^{-t\mu_*^2})\,dt\;\Pi^\perp,\qquad \Gamma_\ge:=\int_{\eps_0^2}^\infty \mathrm{e}^{t\Delta_\eps}(1-\mathrm{e}^{-t\mu_*^2})\,dt\;\Pi^\perp,
\]
let $\gamma_\ge(u):=\int_{\eps_0^2}^\infty \mathrm{e}^{-tu}(1-\mathrm{e}^{-t\mu_*^2})\,dt$ for $u>0$, so that $\Gamma_\ge=\gamma_\ge(-\Delta_\eps)\Pi^\perp$, and put
\begin{gather*}
C^\flat_{B,s}:=(C_{\mu_*}+\Gamma^B+\Gamma_\ge)\Pi^\perp+s\abs T\,\Pi_0\quad(0<s<\infty),\\
C^\flat_B:=C^\flat_{B,\infty}\ \text{(Haar zero mode)},\\
C_1^B:=C_{\mu_*}+\Gamma^B.
\end{gather*}
\end{definition}
Since $\int_0^\infty \mathrm{e}^{-tu}(1-\mathrm{e}^{-t\mu_*^2})\,dt=u^{-1}-(u+\mu_*^2)^{-1}$ for $u>0$, and $-\Delta_\eps>0$ on $\operatorname{ran}\Pi^\perp$,
\begin{equation}\label{app:G1:eq-Xsplit}
(-\Delta_\eps)^{-1}\Pi^\perp=C_{\mu_*}\Pi^\perp+D^<+\Gamma_\ge,\qquad D^<\ge0,\quad\Gamma_\ge\ge0.
\end{equation}

\begin{lemma}[structure]\label{app:G1:struct}
(a) $G_k$, $H_i$ and $\Gamma^B$ commute with $\Pi_0$, and $\Gamma^B\ge0$. On $\operatorname{ran}\Pi_0$, $\Gamma^B$ acts as a number $\gamma_0^B\ge0$. Hence, with $s_1:=\abs T^{-1}(\mu_*^{-2}+\gamma^B_0)\ge s(\mu_*)$,
\[
C_1^B=(C_{\mu_*}+\Gamma^B)\Pi^\perp+s_1\abs T\,\Pi_0,\qquad C^\flat_{B,s}-C_1^B=\Gamma_\ge+(s-s_1)\abs T\,\Pi_0.
\]
(b) \emph{(Domination.)} If the constant $\kappa$ in \eqref{eq:T6-GammaB} satisfies $\kappa\ge\kappa_1\kappa_3$, with $\kappa_1,\kappa_3$ as in Theorem~\ref{app:G1:DOM-nat}, then $D^<\le\Gamma^B\Pi^\perp$.
\end{lemma}
\begin{proof}
(a) $\Delta_\eps\Pi_0=0=\Pi_0\Delta_\eps$. The block average $P_k$ is self-adjoint and fixes the constants, so $P_k\Pi_0=\Pi_0=\Pi_0P_k$. Hence $G_k^{-1}=-\Delta_\eps+\mu_*^2+a_kt_k^{-2}P_k$ commutes with $\Pi_0$, and therefore so do $G_k$, $H_i=G_{i+1}-G_i$ and $\Gamma^B$. On a constant bond function, $G_k\mathbf 1_\mu=(\mu_*^2+a_kt_k^{-2})^{-1}\mathbf 1_\mu$, the same number in every direction; hence $\Gamma^B$ acts on $\operatorname{ran}\Pi_0$ as a number $\gamma_0^B$. Each $H_i$ is a covariance, King's fluctuation covariance \cite[(2.17)]{K1}, so $H_i\ge0$; hence $\Gamma^B=\sum_i\delta_iH_i\ge0$ and $\gamma^B_0\ge0$. On $\operatorname{ran}\Pi_0$, $C_{\mu_*}$ acts as $\mu_*^{-2}=s(\mu_*)\abs T$. This gives the formula for $C_1^B$. The difference formula then follows from Definition~\ref{app:G1:Tflatdef}, since $\Gamma_\ge=\Gamma_\ge\Pi^\perp$.
(b) Theorem~\ref{app:G1:DOM-nat} states $\int_0^{\eps_0^2}\mathrm{e}^{t\Delta_\eps}(1-\mathrm{e}^{-t\mu_*^2})\,dt\le\kappa\mu_*^2\sum_{i<K}t_{i+1}^2H_i=\Gamma^B$ on the whole space, zero mode included. Both sides commute with $\Pi_0$, by (a), so compressing to $\operatorname{ran}\Pi^\perp$ gives the claim.
\end{proof}

\begin{lemma}[domination of the true photon]\label{app:G1:HBdom}
For every $\mu_0\in(0,\mu_*]$ and every $s\ge s(\mu_0)$, $C_{\mu_0}\le C^\flat_{B,s}$; also $C^\flat_s\le C^\flat_{B,s}$ for every $s>0$. Consequently, for every $h\ge0$, every profile $b$ and every $s\ge s(\mu_0)$,
\[
\ev{\mathrm{e}^{-\XR(h)}}_{C_{\mu_0};b}\le\ev{\mathrm{e}^{-\XR(h)}}_{C^\flat_{B,s};b}\xrightarrow[s\to\infty]{}\ev{\mathrm{e}^{-\XR(h)}}_{C^\flat_B;b}.
\]
This is the domination $C_{\mu_0}\le C^\flat_{B,s}$ stated in Subsection~\ref{ssec:ch-scheme}.
\end{lemma}
\begin{proof}
All operators involved commute with $\Pi_0$ (Lemma~\ref{app:G1:struct}(a)), so it suffices to compare them on $\operatorname{ran}\Pi^\perp$ and on $\operatorname{ran}\Pi_0$ separately. On $\operatorname{ran}\Pi^\perp$, by \eqref{app:G1:eq-Xsplit} and Lemma~\ref{app:G1:struct}(b),
\[
(-\Delta_\eps+\mu_0^2)^{-1}\Pi^\perp\le (-\Delta_\eps)^{-1}\Pi^\perp=C_{\mu_*}\Pi^\perp+D^<+\Gamma_\ge\le(C_{\mu_*}+\Gamma^B+\Gamma_\ge)\Pi^\perp.
\]
On $\operatorname{ran}\Pi_0$, $C_{\mu_0}$ acts as $\mu_0^{-2}=s(\mu_0)\abs T$, which is at most $s\abs T$ for $s\ge s(\mu_0)$. The middle inequality above gives $C^\flat_s\le C^\flat_{B,s}$ as well. Next, $\mathrm{e}^{-\XR(h)}=\mathrm{e}^{\norm h_{L^1,\eps}(W_\eps+z_\eps)}\mathrm{e}^{-r^2(\eps^3h)}$ with $\eps^3h\ge0$, a positive constant times an element of the class $\cD$ of Lemma~\ref{lem:C-cov}. By that lemma (\cite[Cor.~6.3(1), p.~281]{BFS1} on the torus graph) its expectation does not decrease when the covariance increases. The limit $s\to\infty$ is Lemma~\ref{app:G1:haar}, applied with the bounded observable $\mathrm{e}^{-\XR(h)}$.
\end{proof}

\begin{lemma}[Haar limit of the zero mode]\label{app:G1:haar}
Let $\Gamma^\perp\ge0$ act on $\operatorname{ran}\Pi^\perp$ and let $\Gamma_s:=\Gamma^\perp+s\abs T\,\Pi_0$. Fix a profile $b$. Let $G$ be either a bounded element of $\cA_\eps$ or a polynomial in the $r_x^2$ with nonnegative coefficients. Then $\ev{G}_{\Gamma_s;b}\to\ev{G}_{\Gamma_\infty;b}$ as $s\to\infty$.
\end{lemma}
\begin{proof}
By Lemma~\ref{lem:C-current} (last sentence) and Lemma~\ref{lem:holonomy}, $c\mapsto\Psi^b_G(A^\perp+c)$ is $P$-periodic in each $c_\mu$. It is also bounded. For bounded $G$, $\abs{\Psi^b_G}\le\norm G_\infty Z_\varphi(0;b)$. For $G=\prod_xr_x^{2n_x}$, the proof of Lemma~\ref{lem:postype} applies verbatim with the radial weight $v_x(r)$ replaced by $r^{2n_x}v_x(r)$, which is again nonnegative and integrable since $\lambda>0$. This gives $\Psi^b_G(A)=\sum_nw_G(n)\mathrm{e}^{ie\eps\sum_bn_bA_b}$ with $w_G\ge0$, hence $\abs{\Psi^b_G(A)}\le\Psi^b_G(0)<\infty$. Under $N(0,\Gamma_s)$, $A^\perp$ and $c$ are independent, the law of $A^\perp$ does not depend on $s$, and $c_\mu$ has variance $s$. Only $c\bmod P$ enters. Its law has density $P^{-1}\sum_{\nu\in\Z}\mathrm{e}^{-2\pi^2\nu^2s/P^2}\cos(2\pi\nu x/P)$ in each coordinate, which converges uniformly to $1/P$, hence in total variation to the uniform law (as in the proof of Lemma~\ref{lem:C-flatlimit}). By bounded convergence the numerator and the denominator $\E[Z_\varphi(\cdot;b)]$ converge, and the limit of the denominator is positive.
\end{proof}

\begin{lemma}[one-point reduction]\label{app:G1:onept}
Let $\Gamma>0$ be a covariance on bond functions, $a(\cdot)$ a profile, and $h\ge0$. Then for every profile $a'\ge a+2h$ (sitewise)
\[
\log\ev{\mathrm{e}^{-\XR(h)}}_{\Gamma;a}=-\int_0^1\sum_y\eps^3h(y)\,\rho_\Gamma(y;a+2th)\,dt\le-\sum_y\eps^3h(y)\,\rho_\Gamma(y;a').
\]
\end{lemma}
\begin{proof}
Let $f(t):=\log\ev{\mathrm{e}^{-t\XR(h)}}_{\Gamma;a}$ for $t\in[0,1]$; then $f(0)=0$. Since $\mathrm{e}^{-t\XR(h)}$ is a c-number times $\exp(-t\sum_x\eps^3h(x)r_x^2)$ and the single-site term is $\frac12a(x)\eps^3r_x^2$, the tilted measure is the model with profile $a+2th$; the c-numbers cancel in the normalisation. Differentiation under the integral is justified because $\mathrm{e}^{-t\sum\eps^3hr^2}\le1$ and $r_x^2$ is integrable ($\lambda>0$). Hence $f'(t)=-\sum_y\eps^3h(y)\rho_\Gamma(y;a+2th)$. Moreover $\partial\ev{r_y^2}_{\Gamma;b}/\partial b(x)=-\frac12\eps^3\ev{r_y^2;r_x^2}_{\Gamma;b}\le0$. This is \cite[Thm.~6.1, p.~276]{BFS1} with $P=r_y^2$, $Q=r_x^2$, $n=m=0$ and zero linear functionals, on the torus graph: Lemma~\ref{app:G1:ginibre} in the instance of Proposition~\ref{app:G1:BFStorus}, with $b_T$ replaced by the profile, which Definition~\ref{app:G1:BFSsetting} allows because the single-site measures $\nu_x$ may depend on $x$. It holds for any centred Gaussian $\Gamma>0$. So $\rho_\Gamma(y;\cdot)$ does not increase along the straight path from $a+2th$ to $a'\ge a+2th$. With $h\ge0$ this gives $f'(t)\le-\sum_y\eps^3h(y)\rho_\Gamma(y;a')$; integrate over $t\in[0,1]$.
\end{proof}

\begin{lemma}[translation averaging of the source]\label{app:G1:TA}
Let $0\le h\le H$, $\mu_0\in(0,\mu_*]$, and let $a'\ge b_T+2H$ be a constant. Then
\[
\log\ev{\mathrm{e}^{-\XR(h)}}_{\eps,T,\mu_0}\le-\norm h_{L^1,\eps}\,\bar\rho^B(a'),\qquad\bar\rho^B:=\bar\rho_{C^\flat_B}.
\]
\end{lemma}
\begin{proof}
Let $\mathcal T$ be the group of the $n:=\abs T/\eps^3$ lattice translations of $T_\eps$, and put $(\tau h)(y):=h(y-\tau)$.

(i) The state $\ev{\cdot}_{\eps,T,\mu_0}$ is invariant under $\mathcal T$ (Lemma~\ref{lem:C-sym}(i)), and $\mathrm{e}^{-\XR(h)}$ composed with the translation by $\tau$ is $\mathrm{e}^{-\XR(\tau h)}$; the c-number $\norm{\tau h}_{L^1,\eps}(W_\eps+z_\eps)=\norm h_{L^1,\eps}(W_\eps+z_\eps)$ is unchanged. Hence $\ev{\mathrm{e}^{-\XR(h)}}_{\eps,T,\mu_0}=\ev{\mathrm{e}^{-\XR(\tau h)}}_{\eps,T,\mu_0}$ for every $\tau\in\mathcal T$.

(ii) Fix $s\ge s(\mu_0)$. By Lemma~\ref{app:G1:HBdom} applied to $\tau h\ge0$, and then Lemma~\ref{app:G1:onept} with $\Gamma=C^\flat_{B,s}>0$, $a\equiv b_T$ and $a'\ge b_T+2H\ge b_T+2\tau h$,
\[
\log\ev{\mathrm{e}^{-\XR(h)}}_{\eps,T,\mu_0}\le\log\ev{\mathrm{e}^{-\XR(\tau h)}}_{C^\flat_{B,s};b_T}\le-\sum_y\eps^3h(y-\tau)\,\rho_{C^\flat_{B,s}}(y;a').
\]

(iii) The left side does not depend on $\tau$. Average the right side over $\tau\in\mathcal T$, using $n^{-1}\sum_{\tau}h(y-\tau)=n^{-1}\sum_xh(x)=\norm h_{L^1,\eps}/\abs T$ for every $y$ (as $h\ge0$):
\[
\log\ev{\mathrm{e}^{-\XR(h)}}_{\eps,T,\mu_0}\le-\frac{\norm h_{L^1,\eps}}{\abs T}\sum_y\eps^3\rho_{C^\flat_{B,s}}(y;a')=-\norm h_{L^1,\eps}\,\bar\rho_{C^\flat_{B,s}}(a').
\]
(iv) Let $s\to\infty$. By Lemma~\ref{app:G1:haar}, with $G=r_y^2$ and $G=1$ at the constant profile $a'$, $\rho_{C^\flat_{B,s}}(y;a')\to\rho^B(y;a')$.
\end{proof}
\begin{remark}\label{app:G1:TArem}
One may also average over $\eps\Z^3/\eps_0\Z^3$ only, and then use that $y\mapsto\rho^B(y;b)$ is $\eps_0\Z^3$-periodic. This holds because $\Gamma^B$ commutes with the translations by $\eps_0\Z^3=t_K\Z^3$ (not with those by $\eps\Z^3$ in general). The proof above averages over all of $\mathcal T$ and does not need this periodicity; only the $\eps\Z^3$-invariance of the true model is used. That invariance fails on a cube with boundary conditions, which is why the lemma is a torus statement. The choice $a'=\hat a_1$, which Lemma~\ref{app:G1:onept} allows, makes a separate step for the monotonicity of $\rho^B$ in the constant coefficient unnecessary.
\end{remark}

\begin{lemma}[determinant comparison for the infrared part; $J_{\rm IR}$ and $Z_T$]\label{app:G1:IR}
Put $\delta_P:=2\sum_{\nu\ge1}\mathrm{e}^{-2\pi^2\nu^2}\ (<5.4\cdot10^{-9})$, and
\begin{gather*}
J_T:=\frac32\abs T^{-1}\sum_{k\ne0}(\hat k^2+\mu_*^2)\,\gamma_\ge(\hat k^2),\\
Z_T:=\abs T^{-1}\Big[\tfrac32\log_+\frac{4\pi^2\mu_*^2\ell_T}{e^2}+3\delta_P\Big],
\end{gather*}
$J_{\rm IR}:=\sup_{K\ge1,\,j\ge1}J_{T_j}$ and $\bar Z:=\sup_{j\ge1}Z_{T_j}$. For every constant $a$:
\begin{enumerate}
\item[(i)] (density ratio) if $0<\Gamma\le\Gamma'$ on $\R^n$ and $f\ge0$, then $\E_{N(0,\Gamma')}[f]\ge\det(\Gamma'\Gamma^{-1})^{-1/2}\E_{N(0,\Gamma)}[f]$;
\item[(ii)] (upper) $Z(C^\flat_B;a)\le Z(C_1^B;a)$;
\item[(iii)] (lower) $\log Z(C^\flat_B;a)\ge\log Z(C_1^B;a)-(J_T+Z_T)\abs T$;
\item[(iv)] with $\alpha:=\sqrt\pi/4$,
\[
J_{\rm IR}\le\tfrac32(\alpha/\eps_0+1/\ell_0)^3+3\mu_*^2(\alpha+1)^3\big(1/\eps_0+1/(32\ell_0)\big),\qquad \bar Z\le\ell_0^{-3}\big[\tfrac32\log_+(4\pi^2\mu_*^2\ell_0/e^2)+0.19+3\delta_P\big].
\] So $J_{\rm IR}$ depends on $\mu_*,\eps_0,\ell_0$ only, and $\bar Z$ on $\mu_*,\ell_0,e$ only.
\end{enumerate}
\end{lemma}
\begin{proof}
(i) $\frac{dN(0,\Gamma')}{dN(0,\Gamma)}(x)=(\det\Gamma/\det\Gamma')^{1/2}\exp(\frac12x\cdot(\Gamma^{-1}-\Gamma'^{-1})x)\ge(\det\Gamma/\det\Gamma')^{1/2}$, since $\Gamma^{-1}\ge\Gamma'^{-1}$.

(ii) By Lemma~\ref{lem:postype}, $Z_\varphi(A;a)=\sum_nw(n)\mathrm{e}^{ie\eps\sum_bn_bA_b}$ with $w\ge0$ and $\sum_nw(n)<\infty$. By Fubini, $Z(\Gamma;a)=\sum_nw(n)\exp(-\frac12e^2\eps^2\operatorname{Var}_\Gamma(\sum_bn_bA_b))$ for every centred Gaussian $\Gamma$, and this is non-increasing in $\Gamma$ in form order. By Lemma~\ref{app:G1:struct}(a), $C^\flat_{B,s}\ge C_1^B$ for $s\ge s_1$, so $Z(C^\flat_{B,s};a)\le Z(C_1^B;a)$. Let $s\to\infty$ (Lemma~\ref{app:G1:haar} with $G=1$; the unnormalised integrals converge by the same bounded-convergence argument).

(iii) \emph{Zero mode.} Let $F(c):=\E_{N(0,(C_{\mu_*}+\Gamma^B+\Gamma_\ge)\Pi^\perp)}[Z_\varphi(A^\perp+c;a)]\ge0$. It is $P$-periodic in each $c_\mu$ (Lemma~\ref{lem:postype}). Put $s':=\max(P^2,s_1)$. For variance $s'\ge P^2$ the wrapped normal density on a circle of length $P$ deviates from $1/P$ by at most $P^{-1}\cdot2\sum_{\nu\ge1}\mathrm{e}^{-2\pi^2\nu^2s'/P^2}\le\delta_P/P$. Hence $\E_{N(0,s')^{\otimes3}}[F]\le(1+\delta_P)^3P^{-3}\int_{[0,P)^3}F$, that is, $Z(C^\flat_B;a)\ge(1+\delta_P)^{-3}Z(C^\flat_{B,s'};a)$. By (i) on the three zero-mode coordinates, $Z(C^\flat_{B,s'};a)\ge(s_1/s')^{3/2}Z(C^\flat_{B,s_1};a)$. Since $\log(1+\delta_P)\le\delta_P$, $\log(s'/s_1)=\log_+(P^2/s_1)$, $s_1\ge s(\mu_*)$ (Lemma~\ref{app:G1:struct}(a)) and $P^2/s(\mu_*)=4\pi^2\mu_*^2\ell_T/e^2$ (with $P=2\pi/(\abs e\ell_T)$ and $s(\mu_*)=(\mu_*^2\ell_T^3)^{-1}$),
\[
\log Z(C^\flat_B;a)\ge\log Z(C^\flat_{B,s_1};a)-Z_T\abs T.
\]
\emph{Infrared part.} By Lemma~\ref{app:G1:struct}(a), $C^\flat_{B,s_1}=C_1^B+\Gamma_\ge$, with $C_1^B\ge C_{\mu_*}>0$. By (i) with $f=Z_\varphi(\cdot;a)\ge0$, and $\log\det(1+K)\le\tr K$ for $K\ge0$,
\[
\log Z(C^\flat_{B,s_1};a)\ge\log Z(C_1^B;a)-\tfrac12\tr\big((C_1^B)^{-1}\Gamma_\ge\big)\ge\log Z(C_1^B;a)-\tfrac12\tr\big(C_{\mu_*}^{-1}\Gamma_\ge\big).
\]
The last step uses $(C_1^B)^{-1}\le C_{\mu_*}^{-1}$ (the inverse is operator monotone) and $\tr\big((C_1^B)^{-1}\Gamma_\ge\big)=\tr\big(\Gamma_\ge^{1/2}(C_1^B)^{-1}\Gamma_\ge^{1/2}\big)$. Finally, $C_{\mu_*}^{-1}\Gamma_\ge=(-\Delta_\eps+\mu_*^2)\gamma_\ge(-\Delta_\eps)\Pi^\perp$ has the eigenvalues $(\hat k^2+\mu_*^2)\gamma_\ge(\hat k^2)$, $k\ne0$, each three times (three components), so $\frac12\tr(C_{\mu_*}^{-1}\Gamma_\ge)=J_T\abs T$. Combining the two displays gives (iii).

(iv) Since $1-\mathrm{e}^{-t\mu_*^2}\le1$, $\gamma_\ge(u)\le \mathrm{e}^{-\eps_0^2u}/u$, so $(u+\mu_*^2)\gamma_\ge(u)\le \mathrm{e}^{-\eps_0^2u}+\mu_*^2\mathrm{e}^{-\eps_0^2u}/u$. With $q^T(t):=\abs T^{-1}\sum_k\mathrm{e}^{-t\hat k^2}$,
\[
J_T\le\tfrac32\big[q^T(\eps_0^2)-\abs T^{-1}\big]+\tfrac32\mu_*^2\int_{\eps_0^2}^\infty\big[q^T(t)-\abs T^{-1}\big]\,dt.
\]
Now $q^T(t)=q_1(t)^3$, with $q_1(t)=\ell_T^{-1}\sum_{k_1}\mathrm{e}^{-4t\eps^{-2}\sin^2(\eps k_1/2)}$. Since $\abs{\sin x}\ge2\abs x/\pi$ for $\abs x\le\pi/2$, we have $4\eps^{-2}\sin^2(\eps k_1/2)\ge4k_1^2/\pi^2$, which equals $16\nu^2/\ell_T^2$ for $k_1=2\pi\nu/\ell_T$. Hence
\[
q_1(t)\le\ell_T^{-1}\sum_{\nu\in\Z}\mathrm{e}^{-16t\nu^2/\ell_T^2}\le\ell_T^{-1}+\ell_T^{-1}\int_\R \mathrm{e}^{-16tx^2/\ell_T^2}dx=\ell_T^{-1}+\alpha t^{-1/2}.
\]
The same inequality gives the spectral gap $\min_{k\ne0}\hat k^2\ge16/\ell_T^2$, so $q^T(t)-\abs T^{-1}\le \mathrm{e}^{-16(t-\ell_T^2)/\ell_T^2}\big(q^T(\ell_T^2)-\abs T^{-1}\big)$ for $t\ge\ell_T^2$. For $t\le\ell_T^2$, $q^T(t)\le(\alpha+1)^3t^{-3/2}$. Therefore (using $\ell_T\ge\ell_0\ge\eps_0$)
\[
\int_{\eps_0^2}^\infty\big[q^T-\abs T^{-1}\big]dt\le\int_{\eps_0^2}^\infty(\alpha+1)^3t^{-3/2}dt+\frac{\ell_T^2}{16}\,(\alpha+1)^3\ell_T^{-3}\le(\alpha+1)^3\Big(\frac2{\eps_0}+\frac1{16\ell_0}\Big),
\]
and $q^T(\eps_0^2)-\abs T^{-1}\le(\alpha/\eps_0+1/\ell_0)^3$. Neither bound depends on $\eps$, $K$ or $j$, which proves the bound on $J_{\rm IR}$.

For $\bar Z$: since $\ell_{T_j}\ge\ell_0$, $\bar Z\le\sup_{\ell_T\ge\ell_0}Z_T$, and Lemma~\ref{app:G1:ZT} gives the bound, because $1/(2\mathrm{e})<0.19$; the estimate there uses the subadditivity of $\log_+$ and $\max_{y\ge1}y^{-3}\log y=1/(3\mathrm{e})$.
\end{proof}

\begin{proposition}[Proposition~\ref{prop:ch-Tflat}, full proof]\label{app:G1:Tflat}
Assume the slice vacuum bound $\Delta_1^B<\infty$ and the anchor condition \eqref{p:ch-anchor}: $b_T+2H\le \hat a_1$ on every torus of $\mathfrak C$, and $g_{21}>0$ (Theorem~\ref{thm:ch-W}(b),(c)). Then for every $0\le h\le H$, every $\mu_0\in(0,\mu_*]\cup\{\flat\}$ and every $(\eps,T_j)\in\mathfrak C$,
\[
\log\ev{\mathrm{e}^{-\XR(h)}}_{\eps,T_j,\mu_0}\le2\norm h_{L^1,\eps}\,S^B_\flat,\qquad S^B_\flat:=[\Delta_1^B+J_{\rm IR}+\bar Z]_+/g_{21},
\]
with $J_{\rm IR}$, $\bar Z$ as in Lemma~\ref{app:G1:IR}. Since $Z_{T_j}\le\bar Z$, this is Proposition~\ref{prop:ch-Tflat}, with $Z_T$ there read as $\bar Z$, so that $S^B_\flat$ does not depend on $T$.
\end{proposition}
\begin{proof}
Write $\hat a_i:=a_{K,T}(m_i)$ and $\rho^B:=\rho_{C^\flat_B}$, $\bar\rho^B:=\bar\rho_{C^\flat_B}$.

\emph{Case $\mu_0\in(0,\mu_*]$.} (1) By Lemma~\ref{app:G1:TA} with $a':=\hat a_1\ge b_T+2H$, $\log\ev{\mathrm{e}^{-\XR(h)}}_{\eps,T,\mu_0}\le-\norm h_{L^1,\eps}\,\bar\rho^B(\hat a_1)$.

(2) Let $f^B(b):=\abs T^{-1}\log Z(C^\flat_B;b)+\frac12b(W_\eps+z_\eps)$ for constant $b$. Since $\partial_bZ_\varphi(A;b)=-\frac12\sum_x\eps^3\Psi^b_{r_x^2}(A)$, and on $b\ge b_{\min}$ we have $\abs{\Psi^b_{r_x^2}(A)}\le\Psi^{b_{\min}}_{r_x^2}(0)$ (positive type, as in the proof of Lemma~\ref{app:G1:haar}), we may differentiate under $\E_{C^\flat_B}$, twice. Using $\sum_x\eps^3=\abs T$,
\begin{gather*}
f^{B\prime}(b)=-\tfrac12\abs T^{-1}\sum_x\eps^3\big(\ev{r_x^2}_{C^\flat_B;b}-W_\eps-z_\eps\big)=-\tfrac12\bar\rho^B(b),\\
f^{B\prime\prime}(b)=\tfrac14\abs T^{-1}\operatorname{Var}_{C^\flat_B;b}\Big(\sum_x\eps^3r_x^2\Big)\ge0.
\end{gather*}
So $f^B$ is convex, and since $\hat a_2-\hat a_1\ge g_{21}>0$, $-\bar\rho^B(\hat a_1)=2f^{B\prime}(\hat a_1)\le2[f^B(\hat a_2)-f^B(\hat a_1)]/(\hat a_2-\hat a_1)$.

(3) By Lemma~\ref{app:G1:IR}(ii) at $\hat a_2$ and (iii) at $\hat a_1$, and since $J_T\le J_{\rm IR}$, $Z_T\le\bar Z$ and by the definition of $\Delta_1^B$,
\[
f^B(\hat a_2)-f^B(\hat a_1)\le\abs T^{-1}\log\frac{Z(C_1^B;\hat a_2)}{Z(C_1^B;\hat a_1)}+\tfrac12(\hat a_2-\hat a_1)(W_\eps+z_\eps)+J_{\rm IR}+\bar Z\le\Delta_1^B+J_{\rm IR}+\bar Z.
\]

(4) Combine (1)--(3) with $\norm h_{L^1,\eps}\ge0$:
\[
\log\ev{\mathrm{e}^{-\XR(h)}}_{\eps,T,\mu_0}\le\frac{2\norm h_{L^1,\eps}[\Delta_1^B+J_{\rm IR}+\bar Z]}{\hat a_2-\hat a_1}\le\frac{2\norm h_{L^1,\eps}[\Delta_1^B+J_{\rm IR}+\bar Z]_+}{g_{21}}.
\]

\emph{Case $\flat$.} $G=\mathrm{e}^{-\XR(h)}$ with $h\ge0$ is an elementary observable of Lemma~\ref{lem:C-flatlimit}, so $\ev{G}_{\eps,T,\flat}=\lim_{\mu_0\downarrow0}\ev{G}_{\eps,T,\mu_0}$. The bound just proved does not depend on $\mu_0$, so it passes to the limit.

\emph{Uniformity.} $\Delta_1^B$ is a supremum over $\mathfrak C$ (Subsection~\ref{ssec:ch-minus}). By Lemma~\ref{app:G1:IR}(iv), $J_{\rm IR}$ and $\bar Z$ do not depend on $\eps$, $j$, $\mu_0$ or $h$, and $g_{21}$ is an infimum over $\eps$ and $T$. No renormalisation-group statement carries a source; the bound is linear in $\norm h_{L^1,\eps}$.
\end{proof}

\subsection{The domination of the infrared part}\label{app:G1:part-Tdom}
This subsection proves the domination of the infrared part $D^<$ by the slice covariance $\Gamma^B$ of \eqref{eq:T6-GammaB}. Lemma~\ref{app:G1:struct}(b) and the constant $\kappa$ of Subsection~\ref{ssec:T6-statement} use it. Only the free (Gaussian) photon enters: there is no Higgs field and no observable.

\begin{definition}[setting]\label{app:G1:DOM-set}
Fix $j\ge1$, $K\ge1$, $\eps=L^{-K}\eps_0$ and $t_k:=L^k\eps$ for $0\le k\le K$, so that $t_K=\eps_0$. Let $\Lambda:=T_{j,\eps}$, with $\ip fg:=\eps^3\sum_{x\in\Lambda}f(x)g(x)$ on real functions on $\Lambda$. $-\Delta_\eps$ is the periodic lattice Laplacian, $\ip\varphi{(-\Delta_\eps)\varphi}=\eps\sum_{\{x,y\}}(\varphi(x)-\varphi(y))^2$, where the sum runs over all nearest-neighbour pairs of $T_{j,\eps}$. Let $M:=-\Delta_\eps+\mu_*^2$, so that $C_{\mu_*}=M^{-1}$.
\begin{enumerate}
\item[(a)] \emph{Blocks.} For $0\le k\le K$ let $\mathcal P_k:=\{B^k(y):y\in t_k\Z^3\cap T_j\}$, where $B^k(y):=\{x\in\Lambda:y_i\le x_i<y_i+t_k\}$ \cite[p.~653]{K1}. Let $\cB_k$ be the functions constant on each block of $\mathcal P_k$, and let $P_k$ be the orthogonal projection onto $\cB_k$, i.e.\ the block average $Q_k^*Q_k$ of \cite[(2.10)]{K1}. Thus $P_0=1$; this $P_0$ is not the projection $\Pi_0$ onto the constants. Moreover $P_{k+1}\le P_k$ and $P_kP_{k+1}=P_{k+1}P_k=P_{k+1}$.
\item[(b)] \emph{Coefficients.} $a_k:=a(1-L^{-2})(1-L^{-2k})^{-1}$ \cite[p.~653]{K1}, $a_\flat:=a(1-L^{-2})$ and $r_k:=a_kt_k^{-2}$.
\item[(c)] \emph{Covariances.} $G_k:=(M+r_kP_k)^{-1}$ for $1\le k\le K$ and $G_0:=0$; $H_j:=G_{j+1}-G_j$; $D:=\sum_{j<K}t_{j+1}^2H_j$, so that $\Gamma^B=\kappa\mu_*^2D$ in \eqref{eq:T6-GammaB}; and
\[
D^{\natural<}:=\int_0^{\eps_0^2}\mathrm{e}^{t\Delta_\eps}\big(1-\mathrm{e}^{-t\mu_*^2}\big)\,dt,\qquad\text{so that } D^<=D^{\natural<}\Pi^\perp\ \text{(Remark~\ref{rem:T6-routes}).}
\]
\item[(d)] \emph{Constants.}
\[
\kappa_1:=\max\Big\{\frac18,\ \frac{L^2/8+1/a_\flat}{L^2-1},\ \frac{L^4}{(L^2-1)^2a_\flat}\Big\},\qquad \kappa_3:=2(1+\mu_*^2\eps_0^2)^2+a_K^2.
\]
\end{enumerate}
$-\Delta_\eps$ and $P_k$ act on each component of a bond function separately, and so does every operator above. An inequality between such operators is therefore the same inequality in each component, and we work with one component throughout. $A\le B$ means $\ip\varphi{A\varphi}\le\ip\varphi{B\varphi}$ for all $\varphi$.
\end{definition}

\emph{Elementary facts.}
(F1) We have $a_1=a$ and $a_{k+1}/a_k=(1-L^{-2k})/(1-L^{-2k-2})<1$. Hence $a\ge a_k>a_\flat$, and $r_m/r_K=L^{2(K-m)}a_m/a_K\ge L^{2(K-m)}$ for $m\le K$; in particular $r_{K-1}-r_K\ge r_{K-1}(1-L^{-2})$.
(F2) For $k\le K$, Lemma~\ref{lem:sizes}(ii) gives $\ell_j/(M_bt_k)\in2\N$. So $\mathcal P_k$ is a partition of $\Lambda$ into lattice cubes with $L^k$ sites per side, and each $\mathcal P_{k+1}$-block is a union of $\mathcal P_k$-blocks. Since $\ell_j\ge2t_k$, no block contains both ends of a wrap-around bond of the torus, so the nearest-neighbour pairs inside a block are exactly those of the cube.
(F3) The block-Neumann forms satisfy
\[
\ip\varphi{(-\Delta_\eps)\varphi}\ \ge\ \sum_{\square\in\mathcal P_k}\eps\sum_{\{x,y\}\subset\square}(\varphi(x)-\varphi(y))^2.
\]
The left side sums over all pairs and the right side only over pairs inside blocks.
(F4) $G_K\le M^{-1}$, since $G_K^{-1}=M+r_KP_K\ge M$.

\begin{lemma}[block Poincar\'e inequality]\label{app:G1:DOM-poinc}
Let $1\le k\le K$. For all $\varphi,v$:
\begin{enumerate}
\item[(a)] $\ip\varphi{(-\Delta_\eps)\varphi}\ge(8/t_k^2)\norm{(1-P_k)\varphi}^2$;
\item[(b)] if $P_kv=0$, then $\ip v{M^{-1}v}\le(t_k^2/8)\norm v^2$;
\item[(c)] $\ip v{G_kv}\le(t_k^2/8)\norm{(1-P_k)v}^2+r_k^{-1}\norm{P_kv}^2$.
\end{enumerate}
\end{lemma}
\begin{proof}
(a) Let $\square$ be a cube with $b=L^k\ge2$ sites per side. With respect to $\ip\cdot\cdot$, the form $\eps\sum_{\{x,y\}\subset\square}(\varphi(x)-\varphi(y))^2$ belongs to the operator $-\Delta_{N,\square}$, $\eps^{-2}$ times the graph Laplacian of $\square$. Its eigenvalues are $\eps^{-2}\sum_{i=1}^3 4\sin^2(\pi n_i/(2b))$ with $0\le n_i<b$. Its kernel consists of the constants on $\square$, and its first nonzero eigenvalue is $(4/\eps^2)\sin^2(\pi/(2b))$. The map $b\mapsto4b^2\sin^2(\pi/(2b))=\pi^2(\sin y/y)^2$, $y=\pi/(2b)$, increases in $b$ and equals $8$ at $b=2$. Hence that eigenvalue is at least $8/(b\eps)^2=8/t_k^2$. The constant part of $\varphi$ on $\square$ is $(P_k\varphi)|_\square$. Summing over $\square\in\mathcal P_k$ and using (F3) gives (a).
(b), (c) For $B>0$ we have $\ip v{B^{-1}v}=\sup_\varphi[2\ip v\varphi-\ip\varphi{B\varphi}]$.
For (b): $\ip v\varphi=\ip v{(1-P_k)\varphi}\le\norm v\,(t_k^2/8)^{1/2}\ip\varphi{(-\Delta_\eps)\varphi}^{1/2}\le\norm v\,(t_k^2/8)^{1/2}\ip\varphi{M\varphi}^{1/2}$ by (a), and $\sup_{s\ge0}(2\alpha s-s^2)=\alpha^2$.
For (c), take $B=M+r_kP_k$. Then $\ip v\varphi=\ip{(1-P_k)v}{(1-P_k)\varphi}+\ip{P_kv}{P_k\varphi}\le\alpha s+\beta\sigma$, with $\alpha:=(t_k^2/8)^{1/2}\norm{(1-P_k)v}$, $s:=\ip\varphi{M\varphi}^{1/2}$ (using (a)), $\beta:=\norm{P_kv}$ and $\sigma:=\norm{P_k\varphi}$. Since $\ip\varphi{B\varphi}=s^2+r_k\sigma^2$ and $\sup_{s,\sigma}[2\alpha s+2\beta\sigma-s^2-r_k\sigma^2]=\alpha^2+\beta^2/r_k$, (c) follows.
\end{proof}

\begin{lemma}[Abel summation]\label{app:G1:DOM-abel}
$D=\sum_{m=0}^{K-1}w_m(G_K-G_m)$, where $w_0:=t_1^2$ and $w_m:=t_{m+1}^2-t_m^2=(L^2-1)t_m^2$ for $m\ge1$. Moreover $G_K-G_m\ge0$ for $0\le m\le K$.
\end{lemma}
\begin{proof}
Since $\sum_{m=0}^{j}w_m=t_{j+1}^2$,
\[
D=\sum_{j<K}\sum_{m\le j}w_m(G_{j+1}-G_j)=\sum_{m<K}w_m\sum_{j=m}^{K-1}(G_{j+1}-G_j)=\sum_{m<K}w_m(G_K-G_m).
\]
For $m=0$, $G_K\ge0$. For $1\le m<K$, $G_m^{-1}-G_K^{-1}=r_mP_m-r_KP_K=r_m(P_m-P_K)+(r_m-r_K)P_K\ge0$, by $P_K\le P_m$ and (F1). Inversion is operator monotone.
\end{proof}

\begin{lemma}[one summand, dual cost]\label{app:G1:DOM-wood}
Let $Y:=G_K^{-1}=M+r_KP_K$. For $1\le m\le K-1$, every $v\in\cB_m$ and every $\varphi$,
\[
\ip\varphi{G_Kv}^2\le\ip\varphi{(G_K-G_m)\varphi}\,c_m(v),\qquad c_m(v):=\ip v{G_Kv}+r_m^{-1}\norm{(P_m-P_K)v}^2+(r_m-r_K)^{-1}\norm{P_Kv}^2.
\]
For $m=0$ and every $v$: $\ip\varphi{G_Kv}^2\le\ip\varphi{G_K\varphi}\,c_0(v)$, with $c_0(v):=\ip v{G_Kv}$.
\end{lemma}
\begin{proof}
The case $m=0$ is the Cauchy--Schwarz inequality for the form of $G_K$. Let $m\ge1$ and put $E_m:=r_m(P_m-P_K)+(r_m-r_K)P_K\ge0$. Then $G_m^{-1}=Y+E_m$, and $E_m$ acts on $\cB_m=\operatorname{ran}P_m$ as $r_m$ on $\operatorname{ran}(P_m-P_K)$ and as $r_m-r_K>0$ on $\operatorname{ran}P_K$, and vanishes on $\cB_m^\perp$. For $\eta>0$ let $E_\eta:=E_m+\eta$ and $B_\eta:=Y^{-1}-(Y+E_\eta)^{-1}=Y^{-1}E_\eta(Y+E_\eta)^{-1}>0$. Its inverse is $(Y+E_\eta)E_\eta^{-1}Y=Y+YE_\eta^{-1}Y$. With $f:=G_Kv=Y^{-1}v$ this gives $\ip f{B_\eta^{-1}f}=\ip v{Y^{-1}v}+\ip v{E_\eta^{-1}v}$, and Cauchy--Schwarz in the $B_\eta$-product gives $\ip\varphi f^2\le\ip\varphi{B_\eta\varphi}\ip f{B_\eta^{-1}f}$. As $\eta\downarrow0$ (finite dimension), $B_\eta\to Y^{-1}-G_m=G_K-G_m$. For $v\in\cB_m$,
\[
\ip v{E_\eta^{-1}v}=(r_m+\eta)^{-1}\norm{(P_m-P_K)v}^2+(r_m-r_K+\eta)^{-1}\norm{P_Kv}^2,
\]
which tends to the last two terms of $c_m(v)$.
\end{proof}

\begin{lemma}[block lemma]\label{app:G1:DOM-block}
$G_K^2\le\kappa_1D$.
\end{lemma}
\begin{proof}
\emph{Case $K\ge2$.} Fix $\varphi$ and $u$. Decompose $u$ into the mutually orthogonal pieces
\[
\begin{gathered}
v_0:=(1-P_1)u,\qquad v_m:=(P_m-P_{m+1})u\ \ (1\le m\le K-2),\\
v_{K-1}:=P_{K-1}u=v'+v'',\qquad v':=(P_{K-1}-P_K)u,\quad v'':=P_Ku,
\end{gathered}
\]
so that $v_m\in\cB_m$ and $\norm{v_0}^2+\sum_{1\le m\le K-2}\norm{v_m}^2+\norm{v'}^2+\norm{v''}^2=\norm u^2$. By Lemma~\ref{app:G1:DOM-wood}, Cauchy--Schwarz over $m$, and Lemma~\ref{app:G1:DOM-abel},
\begin{multline*}
\ip\varphi{G_Ku}=\sum_{m<K}\ip\varphi{G_Kv_m}\le\sum_{m<K}\big(w_m\ip\varphi{(G_K-G_m)\varphi}\big)^{1/2}\Big(\frac{c_m(v_m)}{w_m}\Big)^{1/2}\\
\le\ip\varphi{D\varphi}^{1/2}\Big(\sum_{m<K}\frac{c_m(v_m)}{w_m}\Big)^{1/2}.
\end{multline*}
We bound the costs using Lemma~\ref{app:G1:DOM-poinc}, (F1) and (F4).
\begin{itemize}
\item $m=0$: $P_1v_0=0$, so $c_0\le\ip{v_0}{M^{-1}v_0}\le(t_1^2/8)\norm{v_0}^2$ and $c_0/w_0\le\frac18\norm{v_0}^2$.
\item $1\le m\le K-2$: $P_{m+1}v_m=0=P_Kv_m$, so $c_m\le(t_{m+1}^2/8+t_m^2/a_m)\norm{v_m}^2$ and $c_m/w_m\le\frac{L^2/8+1/a_\flat}{L^2-1}\norm{v_m}^2$.
\item $m=K-1$: $(P_{K-1}-P_K)v_{K-1}=v'$ and $P_Kv_{K-1}=v''$. Lemma~\ref{app:G1:DOM-poinc}(c) with $k=K$ gives $\ip{v_{K-1}}{G_Kv_{K-1}}\le(t_K^2/8)\norm{v'}^2+(t_K^2/a_K)\norm{v''}^2$. With $r_{K-1}^{-1}=t_{K-1}^2/a_{K-1}$ and $(r_{K-1}-r_K)^{-1}\le L^2t_{K-1}^2/((L^2-1)a_{K-1})$,
\[
\frac{c_{K-1}}{w_{K-1}}\le\frac{L^2/8+1/a_\flat}{L^2-1}\norm{v'}^2+\frac{L^2/a_\flat+L^2/((L^2-1)a_\flat)}{L^2-1}\norm{v''}^2,
\]
and the last coefficient equals $L^4/((L^2-1)^2a_\flat)$.
\end{itemize}
Hence $\sum_mc_m(v_m)/w_m\le\kappa_1\norm u^2$, i.e.\ $\ip{G_K\varphi}u^2\le\kappa_1\ip\varphi{D\varphi}\norm u^2$ for all $u$. Taking $u=G_K\varphi$ gives $\norm{G_K\varphi}^2\le\kappa_1\ip\varphi{D\varphi}$, i.e.\ $G_K^2\le\kappa_1D$.

\emph{Case $K=1$.} Here $D=t_1^2G_1$, and Lemma~\ref{app:G1:DOM-poinc}(c) with $a_1=a$ gives $G_1\le t_1^2\max(1/8,1/a)$. Hence $G_1^2\le\norm{G_1}G_1\le\max(1/8,1/a)D\le\kappa_1D$, because $L^4/((L^2-1)^2a_\flat)=L^6/((L^2-1)^3a)\ge1/a$.
\end{proof}

\begin{lemma}[block geometry]\label{app:G1:DOM-geom}
In the setting of Definition~\ref{app:G1:DOM-set}, for $0\le k\le K$:
\begin{enumerate}
\item[(a)] $\ell_j/t_k\in2M_b\N$. In particular $\ell_j\ge2t_k$, and $N:=\ell_j/\eps\ge4$;
\item[(b)] $\mathcal P_k$ is a partition of $\Lambda$ into cubes with $L^k$ sites per side, and for $k<K$ each block of $\mathcal P_{k+1}$ is a union of $L^3$ blocks of $\mathcal P_k$. Hence $\cB_{k+1}\subset\cB_k$, $P_{k+1}\le P_k$ and $P_kP_{k+1}=P_{k+1}P_k=P_{k+1}$;
\item[(c)] for every block $\square\in\mathcal P_k$, a bond of $T_{j,\eps}$ has both endpoints in $\square$ if and only if it joins two points of $\square$ whose coordinates differ by $\eps$ in exactly one direction, without wrap-around. Thus the bonds inside $\square$ are exactly the edges of the cube graph $\{0,\dots,L^k-1\}^3$;
\item[(d)] $\ip\varphi{-\Delta_\eps\varphi}\ge\sum_{\square\in\mathcal P_k}\eps\sum_{b\subset\square}(\varphi(b_+)-\varphi(b_-))^2$, where $b\subset\square$ means that both endpoints of the bond $b=(b_-,b_+)$ lie in $\square$.
\end{enumerate}
\end{lemma}
\begin{proof}
(a) We have $t_k=L^k\eps_K$ and $k\le K\le K+n_0$, so Lemma~\ref{lem:sizes}(ii) gives $\ell_j/(M_bt_k)\in2\N$. Hence $\ell_j/t_k\in2M_b\N$ and $\ell_j\ge2M_bt_k\ge2t_k$. Taking $k=K$ gives $N=\ell_j/\eps=L^K(\ell_j/t_K)\ge2L\ge4$.

(b) By (a), $n_k:=\ell_j/t_k$ is an integer. In each direction the intervals $[y_i,y_i+t_k)$ with $y_i\in t_k\{0,\dots,n_k-1\}$ partition $[0,\ell_j)$, and each contains $t_k/\eps=L^k$ lattice points. Taking products gives the partition of $\Lambda$. For $y\in t_{k+1}\Z^3$ we have $[y_i,y_i+Lt_k)=\bigcup_{m=0}^{L-1}[y_i+mt_k,y_i+(m+1)t_k)$ with $y_i+mt_k\in t_k\Z$, which gives the nesting. A function that is constant on the coarser blocks is constant on the finer ones, and the statements about the projections follow.

(c) Let $\square=B^k(y)$. In direction $i$ the points of $\square$ have coordinates in $[y_i,y_i+t_k)\subset[0,\ell_j)$, so two of its points differ in the $i$th coordinate by at most $t_k-\eps$. A bond in direction $i$ joins points $x,x'$ that agree off coordinate $i$. Either $x'_i=x_i+\eps$ (an interior bond), or $x_i=\ell_j-\eps$ and $x'_i=0$ (a wrap-around bond). Both endpoints of a wrap-around bond lie in $\square$ only if $\ell_j-\eps\le t_k-\eps$, that is $\ell_j\le t_k$, and (a) excludes this. Conversely, two points of $\square$ that differ by $\eps$ in one coordinate and agree in the others are the endpoints of an interior bond. Since $N\ge4$, distinct bonds have distinct pairs of endpoints.

(d) The blocks of $\mathcal P_k$ are disjoint, so each bond lies inside at most one block. The left side is $\eps$ times the sum of $(\varphi(b_+)-\varphi(b_-))^2$ over all bonds. The right side is the same sum restricted to the bonds that lie inside some block.
\end{proof}

\begin{theorem}[domination of the slice-modified photon]\label{app:G1:DOM-nat}
Put $\bar\kappa:=\kappa_1\big(2(1+\mu_*^2\eps_0^2)^2+a^2\big)$. Then on all of $\ell^2(\Lambda)$, zero mode included,
\[
D^{\natural<}\le\kappa_1\kappa_3\,\mu_*^2\,D\le\bar\kappa\,\mu_*^2\,D .
\]
Consequently $D^{\natural<}\le\kappa\mu_*^2D$ for every $\kappa\ge\bar\kappa$, for every $K\ge1$. The constant $\bar\kappa$ depends only on $L$, $a$ and $\mu_*\eps_0$. It does not depend on $\eps$, $K$ or $\ell_j$. If $L=2$, $a=1$ and $0<\mu_*\eps_0\le\frac15$, then for every $K\ge1$
\[
\frac{164}{27}<\kappa_1\kappa_3\le\bar\kappa\le\frac{42176}{5625}<\frac{15}2 .
\]
\end{theorem}
\begin{proof}
Put $\vartheta:=\eps_0^2=t_K^2$ and $g^\natural(\lambda):=\mu_*^{-2}\int_0^\vartheta \mathrm{e}^{-s\lambda}(1-\mathrm{e}^{-s\mu_*^2})\,ds$ for $\lambda\ge0$, so that $D^{\natural<}=\mu_*^2g^\natural(-\Delta_\eps)$.

\emph{Step 1: $G_K^{-2}\le h(-\Delta_\eps)$, with $h(\lambda):=2(\lambda+\mu_*^2)^2+2r_K^2$.} From $(M-r_KP_K)^2\ge0$ and $P_K^2=P_K$ we get $r_K(MP_K+P_KM)\le M^2+r_K^2P_K$. Hence $G_K^{-2}=(M+r_KP_K)^2\le2M^2+2r_K^2P_K\le2M^2+2r_K^2=h(-\Delta_\eps)$. This uses only the inequality $AB+BA\le A^2+B^2$ for self-adjoint $A,B$; it never treats $G_K$ as a function of $-\Delta_\eps$. Since inversion is operator monotone, $G_K^2\ge h(-\Delta_\eps)^{-1}$.

\emph{Step 2: $g^\natural h\le\kappa_3$ on $[0,\infty)$.} Since $1-\mathrm{e}^{-y}\le y$, we have $g^\natural(\lambda)\le g_0(\lambda):=\int_0^\vartheta s\mathrm{e}^{-s\lambda}\,ds$. Moreover $g_0\le\vartheta^2/2$, and $g_0(\lambda)\le\int_0^\infty s\mathrm{e}^{-s\lambda}ds=\lambda^{-2}$ for $\lambda>0$. With $m':=\max(\lambda,1/\vartheta)$ this gives $g_0(\lambda)m'^2\le1$. Also $\lambda+\mu_*^2\le(1+\mu_*^2\vartheta)m'$, since $\vartheta m'\ge1$. As $r_K=a_K/\vartheta$,
\[
g^\natural(\lambda)h(\lambda)\le2(1+\mu_*^2\vartheta)^2+2\frac{a_K^2}{\vartheta^2}\cdot\frac{\vartheta^2}2=2(1+\mu_*^2\eps_0^2)^2+a_K^2=\kappa_3.
\]
This includes $\lambda=0$, the constants.

\emph{Step 3.} $g^\natural(-\Delta_\eps)$ and $h(-\Delta_\eps)$ are commuting functions of $-\Delta_\eps$, and $h>0$. So Step 2 gives $g^\natural(-\Delta_\eps)\le\kappa_3h(-\Delta_\eps)^{-1}$, and Step 1 gives $\kappa_3h(-\Delta_\eps)^{-1}\le\kappa_3G_K^2$. Lemma~\ref{app:G1:DOM-block} then gives $D^{\natural<}=\mu_*^2g^\natural(-\Delta_\eps)\le\kappa_3\mu_*^2G_K^2\le\kappa_1\kappa_3\mu_*^2D$. Lemma~\ref{app:G1:DOM-block} rests on the block Poincar\'e inequality, Lemma~\ref{app:G1:DOM-poinc}. That inequality uses the cube structure of the blocks, which is Lemma~\ref{app:G1:DOM-geom}(c),(d): no block contains both endpoints of a wrap-around bond, because $\ell_j\ge2t_k$ for every $k\le K$. By (F1), $a_K\le a_1=a$, so $\kappa_3\le2(1+\mu_*^2\eps_0^2)^2+a^2$ and $\kappa_1\kappa_3\le\bar\kappa$. Finally $D\ge0$ (Lemma~\ref{app:G1:DOM-abel}), so the bound with $\bar\kappa$ implies the bound with every $\kappa\ge\bar\kappa$.

\emph{Step 4: the values at $L=2$, $a=1$.} Here $a_\flat=\frac34$. The three numbers in the definition of $\kappa_1$ are $\frac18$, $(\frac12+\frac43)/3=\frac{11}{18}$ and $16/(9\cdot\frac34)=\frac{64}{27}$, so $\kappa_1=\frac{64}{27}$. By (F1), $\frac34<a_K\le1$. Put $x:=\mu_*\eps_0\in(0,\frac15]$. Then $\kappa_3=2(1+x^2)^2+a_K^2$ increases in $x$ and in $a_K$, so $\kappa_3\le2(1+x^2)^2+1\le2\cdot\frac{676}{625}+1=\frac{1977}{625}$. Hence $\kappa_1\kappa_3\le\bar\kappa\le\frac{64}{27}\cdot\frac{1977}{625}=\frac{42176}{5625}=7.4979\ldots<\frac{15}2$. On the other side, $\kappa_3>2+a_\flat^2=\frac{41}{16}$, so $\kappa_1\kappa_3>\frac{64}{27}\cdot\frac{41}{16}=\frac{164}{27}$. The value $\frac{164}{27}$ is the limit of $\kappa_1\kappa_3$ as $\mu_*\eps_0\to0$ and $K\to\infty$, and it is not attained. So this argument gives no constant $\kappa\le\frac{164}{27}$. The constant of \eqref{eq:T6-GammaB} must be a fixed number, independent of $K$. The fixed value $\kappa=\frac{15}2$ satisfies $\kappa\ge\bar\kappa$ under the condition (E.dom) $\mu_*\eps_0\le0.2$ (Remark~\ref{app:G1:DOM-kappa}).
\end{proof}

\begin{corollary}[torus form]\label{app:G1:DOM-torus}
Let $\kappa\ge\kappa_1\kappa_3$, and let $\Gamma^B$ and $\bar\delta$ be as in \eqref{eq:T6-GammaB}. Then:
\begin{enumerate}
\item[(a)] $D^<=D^{\natural<}\Pi^\perp\le\Gamma^B\Pi^\perp$;
\item[(b)] $D\le\eps_0^2G_K\le\eps_0^2C_{\mu_*}$, and hence $\Gamma^B\le\kappa\mu_*^2\eps_0^2G_K\le\bar\delta\,C_{\mu_*}$.
\end{enumerate}
\end{corollary}
\begin{proof}
(a) $P_k$ is self-adjoint and fixes the constants, so $P_k\Pi_0=\Pi_0=\Pi_0P_k$. Hence $G_k$, $H_j$, $D$ and $\Gamma^B$ commute with $\Pi_0$, and so does $\mathrm{e}^{t\Delta_\eps}$. Compressing the inequality $D^{\natural<}\le\kappa\mu_*^2D=\Gamma^B$ of Theorem~\ref{app:G1:DOM-nat} by $\Pi^\perp$ gives $\Pi^\perp D^{\natural<}\Pi^\perp\le\Pi^\perp\Gamma^B\Pi^\perp$, and the two sides are $D^<$ and $\Gamma^B\Pi^\perp$.
(b) $H_j\ge0$ by Lemma~\ref{app:G1:DOM-abel}, since $H_j=(G_K-G_j)-(G_K-G_{j+1})$ and $G_{j+1}^{-1}\le G_j^{-1}$ by the same computation. Also $t_{j+1}\le t_K=\eps_0$ for $j<K$ and $\sum_{j<K}H_j=G_K-G_0=G_K$. Hence $D\le\eps_0^2G_K$, and $G_K\le C_{\mu_*}$ by (F4). Multiply by $\kappa\mu_*^2$.
\end{proof}

\begin{remark}[the value of $\kappa$]\label{app:G1:DOM-kappa}
At $L=2$ and $a=1$ (stages~1--2 of Table~\ref{p:tab:stages}) we have $a_\flat=3/4$ and $\kappa_1=\max\{1/8,\,11/18,\,64/27\}=64/27$. Under the condition (E.dom), $\mu_*\eps_0\le0.2$, imposed at stage~10a, and with $a_K\le a=1$ by (F1),
\[
\kappa_1\kappa_3\le\frac{64}{27}\Big(2\cdot\frac{676}{625}+1\Big)=\frac{126528}{16875}<7.498<7.5\qquad\text{for every }K\ge1.
\]
As $\mu_*\eps_0\to0$ and $K\to\infty$ (so $a_K\to a_\flat$), $\kappa_1\kappa_3\to\frac{64}{27}\cdot\frac{41}{16}=\frac{164}{27}\approx6.07$. This limit is not attained: $a_K>a_\flat$ by (F1), so $\kappa_3>2+a_\flat^2=\frac{41}{16}$ and $\kappa_1\kappa_3>\frac{164}{27}>6.07$ for every $K$ and every $\mu_*\eps_0$. Hence Theorem~\ref{app:G1:DOM-nat} does not give the value $\kappa=6.07$, and the condition (E.$\delta$) is used with $\kappa=7.5$ (Lemma~\ref{app:G1:RW-B17}). By Theorem~\ref{app:G1:DOM-nat} and Corollary~\ref{app:G1:DOM-torus}, any fixed $\kappa\in[126528/16875,\,\infty)$, in particular $\kappa=7.5$, can serve as the $K$-independent domination constant in \eqref{eq:T6-GammaB}; it is fixed at stage~0, under the a~priori bound (E.dom). Note that the inequality $\Gamma^B\le\kappa\mu_*^2\eps_0^2G_K\le\bar\delta C_{\mu_*}$ of Corollary~\ref{app:G1:DOM-torus}(b) holds for every $\kappa>0$. The role of $\kappa$ is the domination $D^{\natural<}\le\kappa\mu_*^2D$.
\end{remark}

\subsection{The parameters and the hypotheses of \texorpdfstring{Theorem~\ref{thm:ch-T9}}{the tube bound}}\label{app:G1:part-D7}

This subsection contains four statements used in the proof of Theorem~\ref{thm:main}: the choice of the auxiliary parameters $(\eta,m_1,m_2)$ in Theorem~\ref{thm:ch-T9} (Lemma~\ref{app:G1:window-apply}); the list of the estimates of \cite{Ba2,Ba3} that are not written out in print and are used in the proofs of Theorems~\ref{thm:UV-K0plus} and~\ref{thm:T6}, with the places where they are proved (Lemma~\ref{app:G1:items}); the transfer of these estimates to the masses at which Theorem~\ref{thm:UV-K0plus} is used (Lemma~\ref{app:G1:RMG}); and the verification of the hypotheses of Theorem~\ref{thm:ch-T9} on every torus of the family $\mathfrak C$ (Proposition~\ref{app:G1:VACT}).

Throughout, $\mathbb M_4=\{M_1,m,m_1,m_2\}$, the stages 0--11 are those of Table~\ref{p:tab:stages}, and, as in Section~\ref{sec:chain}, $\mathfrak C$ is the family of pairs $(\eps,T_j)$ with $\eps=2^{-K}\eps_0$, $K\ge K_0=\max(K_0',n_0+1)$, $j\ge1$. A Higgs-mass entry of a constant or of a condition is \emph{of type $(\alpha)$} if the mass $M$ enters only through the rescaled mass $M^2t^2$, $t\le\eps_0$, and the constant is uniform for $0\le M^2t^2\le1$. It is \emph{of type $(\beta)$} if $M$ enters through positive powers or logarithms of $M$, or through quantities finite at every fixed $M>0$, and only inside downward-closed upper bounds on $\eps_0$, monotone constants of the stages 8--9, or the floor on $\ell_0$. It is \emph{of type $(\gamma)$} if it orders $M$ against $\mu_*$ or uses $m_1>2m$.

\paragraph{The auxiliary parameters.}

\begin{lemma}[the choice of $(\eta,m_1,m_2)$]\label{app:G1:window-apply}
Let $m>0$, $0<\lambda<0.1m$ and $0<e^2<0.1m$, as in Theorem~\ref{thm:main}. Put $\pi(e,\lambda,m):=(e,\lambda,m,\tfrac12,3m,4m)$. Then:
\begin{enumerate}
\item[(a)] $\pi(e,\lambda,m)\in\Wwin$, and at this point Theorem~\ref{thm:ch-W}(a)--(d) give $2\sigma_+\ge0.319m^2$, $g_{21}\ge6.6369\,m^2$, $a_{K,T}(3m)\ge a_{K,T}(m)+2H$ and $3m\ge\sqrt2\mu_*$, uniformly in $\eps\le1/m$ and in tori of side $\ge96/m$;
\item[(b)] Lemma~\ref{lem:window} at $\pi(e,\lambda,m)$ provides $R$, $b_0$, $n_0$ and $\ell_0\in2M_b\N$ with $\ell_0\ge183/m$ that meet every condition of Appendix~\ref{app:B}, and every $(\eps,T_j)\in\mathfrak C$ has $\eps\le1/(4m)$ and $\ell_j\ge96/m$;
\item[(c)] with these data, the constants of Theorem~\ref{thm:main}(i) are $\sigma_+:=g/2$, with $g$ as in Proposition~\ref{prop:ch-Pplus}, and $C:=\mathsf K=\max(\kappa_+,2S^B_\flat)$, with $\kappa_+$ from Proposition~\ref{prop:ch-Pplus} and $S^B_\flat$ from Proposition~\ref{prop:ch-Tflat} at $(\eta,m_1,m_2)=(\tfrac12,3m,4m)$, with $Z_T$ read as $\bar Z$ (Proposition~\ref{app:G1:Tflat}, Corollary~\ref{app:G1:T9const}); the threshold is $K_0:=\max(K_0',n_0+1)$.
\end{enumerate}
Any other $\eta\in(0,1)$ and $2m<m_1<m_2$ may be used instead; the constants depend on the choice.
\end{lemma}
\begin{proof}
(a) $0<\tfrac12<1$ and $2m<3m<4m$, and the conditions on $(e,\lambda,m)$ are those of $\Wwin$ in \eqref{eq:window}, so $\pi(e,\lambda,m)\in\Wwin$. Theorem~\ref{thm:ch-W}(b) gives $g_{21}\ge(4m-3m)(4m+3m-0.3631m)=6.6369\,m^2$. Theorem~\ref{thm:ch-W}(c) holds because $3m>2m>1.882m\ge m_1^c$, and $3m\ge\sqrt2\,m=\sqrt2\mu_*$.
(b) This is the conclusion of Lemma~\ref{lem:window} at $\pi(e,\lambda,m)$. The condition (E.m2), $m_2^2\eps_0^2\le1$, gives $\eps\le\eps_0\le1/m_2=1/(4m)$. The torus floor (rule~(iii) of Subsection~\ref{p:appA:window}) is $\ell_0^\natural=\max\{8,128/(3m),128/(0.7m)\}=182.857\ldots/m$, so the choice $\ell_0\ge183/m$ meets it.
(c) Theorem~\ref{thm:ch-T9} gives the bound of Theorem~\ref{thm:main}(i) with the tube $\mathcal D_{\sigma_+,H}$, $\sigma_+=g/2$, and the constant $\mathsf K$. The spacing threshold is that of Theorem~\ref{thm:ch-T9}.
\end{proof}

\paragraph{The estimates of \texorpdfstring{\cite{Ba2,Ba3}}{Balaban} at the points of Theorem~\ref{thm:T6}.}

\paragraph{The estimates of the proof of Theorem~\ref{thm:T6} at its two points.}
Throughout this paragraph the hypotheses of Theorem~\ref{thm:T6} are in force. We take $\eps=2^{-K}\eps_0$ with $K\ge K_0=\max(K_0',n_0+1)$, $T=T_j$ with $j\ge1$, and $t_k:=L^k\eps$, so that $t_K=\eps_0$. For $X\subset T$ we write $\abs X_k$ for the number of points of $X$ in the lattice $t_k\Z^3$. For a set $\Lambda$ of scale $k-1$ we write $\Lambda':=\Lambda\cap t_k\Z^3$ \cite[(1.16), p.~606]{Ba1}. Inside this paragraph, $i$ denotes a scale index.

\begin{definition}[the omitted estimates]\label{app:G1:itemlist}
Let $\mathcal I$ be the following collection of statements. Each one is followed by the place where it is proved.
\begin{enumerate}
\item[(a)] The ingredients of Table~\ref{tab:T6-inputs}:
\begin{itemize}
\item the moved vacuum pieces at every scale (Proposition~\ref{app:G1:I1});
\item the Ward--Takahashi step on the torus (Proposition~\ref{app:G1:I2});
\item the conditions of \cite[Prop.~1]{Ba3} as used in \cite{Ba2} (Proposition~\ref{app:G1:I3});
\item the restriction to real gauge backgrounds near Neumann faces (Proposition~\ref{app:G1:I4});
\item \cite[Prop.~1]{Ba3} at the order cap $\Nbar$ (Proposition~\ref{app:G1:I5});
\item the region entropy, that is, the estimate \cite[(3.42)--(3.54), pp.~592--594]{Ba2} on the torus (Proposition~\ref{app:G1:I6});
\item the mass reading of \cite[p.~589]{Ba2} (Proposition~\ref{app:G1:I7});
\item the torus-uniform final c-numbers (Proposition~\ref{app:G1:I8});
\item the vacuum classes at the last scale (Proposition~\ref{app:G1:I9}).
\end{itemize}
\item[(b)] The statements that make up Proposition~\ref{app:G1:I3}, which are the entries of Table~\ref{app:G1:I3tab}:
\begin{itemize}
\item the one-variable remainder-kernel bound (entry 1, Lemma~\ref{app:G1:I3r-U1});
\item the two-variable bound on $T_\eta$ (entry 2, Lemma~\ref{app:G1:I3r-U1m}) and on Neumann domains $\Omega$ (entry 3, Lemma~\ref{app:G1:I3r-U1mOmega});
\item the absence of mass insertions other than $\delta m^2$ (entry 4, Lemma~\ref{app:G1:I3r-ME});
\item the abbreviated steps of \cite{Ba2,Ba3} that need more than the printed text (entries 5, $5'$, Lemma~\ref{app:G1:I3r-Pw});
\item the separation statement for unions of big blocks (entry 6, Lemma~\ref{app:G1:I3r-SEP});
\item the cube remainder of \cite[p.~433]{Ba3} (entry 7, Lemma~\ref{app:G1:I3r-cube});
\item the two-domain bounds for $\delta G_k$ (entry 8, Lemma~\ref{app:G1:I3r-DGM});
\item the inclusions of the cube $\square$ in the outer domains and the retirement of $\square\subset\Omega_2$ (entries 9, 10, Lemmas~\ref{app:G1:I3r-Wsq}, \ref{app:G1:I3r-W2R});
\item the end-point statements for $\Omega_2$-lines in \cite[Chaps.~2, 3]{Ba2} (entries 11, 12, Lemmas~\ref{app:G1:I3r-HOm2}, \ref{app:G1:I3r-HOm2ch3});
\item the nesting $\Omega\supset\Omega_1\supset\Omega_2$ (entry 13, Lemma~\ref{app:G1:Rnest} and Remark~\ref{app:G1:Rnest-scope});
\item the support hypothesis in the form \eqref{eq:R4ppp} (entry 14, Corollary~\ref{cor:prop1R4});
\item the conventions R1--R3 on gating, nesting and metric (entry 15, Lemma~\ref{app:G1:I3r-R123});
\item the completeness of the lists of uses of \cite[Prop.~1]{Ba3} (entry $15'$, Lemma~\ref{app:G1:I3r-Rexh});
\item the local $A\cdot A$ terms (entry 16, Lemmas~\ref{app:G1:AA}, \ref{app:G1:AA4} and~\ref{app:G1:I3r-AA});
\item the cuts beyond the order $\nbar$ of \cite[p.~579]{Ba2} (entry 17, Proposition~\ref{app:G1:T579} and Lemma~\ref{app:G1:I3r-T579});
\item the two-domain bounds at the spacings $0<\eta\le\frac12$ and $\eta=1$ (entry 8, Proposition~\ref{app:G1:twodomain} and Lemma~\ref{app:G1:I3a-eta1}), and the statements of the last row of the table (Lemmas~\ref{app:G1:I3a-onebg}, \ref{app:G1:I3a-Bloc}, \ref{app:G1:I3a-margins} and~\ref{app:G1:I3a-rate}).
\end{itemize}
\item[(c)] The statements on domains, supports and localisation widths: Corollary~\ref{cor:prop1R4}, which gives the form \eqref{eq:R4ppp} of \cite[p.~412]{Ba3}, and Lemmas~\ref{app:G1:Rnest}, \ref{app:G1:Rthetab} and~\ref{app:G1:Rloc7}.
\item[(d)] The order cap $\Nbar=\nbar^2=169$ for \cite[(1.33)]{Ba3} (Proposition~\ref{app:G1:I5}).
\item[(e)] The final step of the upper bound, \cite[(2.116)--(2.118), p.~582]{Ba2} (Lemma~\ref{app:G1:N10} below).
\item[(f)] \cite[Lemma~2.4]{Ba2} at edge points with the cap \eqref{eq:lem24-cap}, that is, condition (E.edge) (Lemma~\ref{lem:lem24}); the periodisation lemma with condition (T.per) (Lemma~\ref{lem:IMG}); and the depth of each sketched step of \cite[Prop.~1]{Ba3} (Proposition~\ref{prop:prop1} and Lemma~\ref{lem:Omegafree}).
\end{enumerate}
\end{definition}

\begin{lemma}[the final step of the upper bound]\label{app:G1:N10}
Let $\rho^{(K)}$ be the density of \cite[(2.43), p.~566]{Ba2} at $k=K$, so that, in the notation of \cite[pp.~566--582]{Ba2},
\begin{equation}\label{app:G1:eq-243K}
\begin{split}
Z^\eps\le{}&\int\!dB\!\int\!d\psi\sum_{\Lambda_0^{(0)},\dots,\Lambda_0^{(K-1)}}\rho^{(K)}\exp\Big(\mathcal P^{(K)}\big(\Lambda_7^{(K-1)\prime},\theta_KB^{(K),\eps},\psi\big)-E_0\\
&+\sum_{i=0}^{K-1}O(t_i^{\kappa_0})\abs{\Lambda_7^{(i-1)\prime}\cap\Lambda_7^{(i)c}}\Big)\exp\Big(\sum_{i<K}O(t_i^{\kappa})\abs{T_1^{(i)}}\Big),
\end{split}
\end{equation}
with $\Lambda_7^{(-1)\prime}:=T$. Here $\kappa>0$ is the exponent of \cite[Prop.~2.2, (2.67)]{Ba2}; it is unrelated to the slice constant of \eqref{eq:T6-GammaB}. Let $\kappa_0>0$ and $C\ge0$, and assume that on the support of $\rho^{(K)}$
\begin{equation}\label{app:G1:eq-F}
\text{(F)}\qquad \mathcal P^{(K)}\big(\Lambda_7^{(K-1)\prime},\theta_KB^{(K),\eps},\psi\big)\le C\,t_K^{\kappa_0}\,\abs{\Lambda_7^{(K-1)\prime}}_K .
\end{equation}
Then:
\begin{enumerate}
\item[(a)] $\mathcal P^{(K)}(\Lambda_7^{(K-1)\prime},\theta_KB^{(K),\eps},\psi)\le C\eps_0^{\kappa_0-3}\abs T$ on the support of $\rho^{(K)}$;
\item[(b)] with $c$ the constant of \cite[(2.118), p.~582]{Ba2},
\[
Z^\eps\le \mathrm{e}^{(C\eps_0^{\kappa_0-3}+c)\abs T}\int\!dB\!\int\!d\psi\sum_{\Lambda_0^{(0)},\dots,\Lambda_0^{(K-1)}}\rho^{(K)}\mathrm{e}^{-E_0}\exp\Big(\sum_{i=0}^{K-1}O(t_i^{\kappa_0})\abs{\Lambda_7^{(i-1)\prime}\cap\Lambda_7^{(i)c}}\Big).
\]
Moreover, if the sizes in \cite[(2.117)]{Ba2} are read as $\abs\cdot_{K-1}$, then the right side of \eqref{app:G1:eq-243K} is at most $\mathrm{e}^{c\abs T}$ times the left side of \cite[(2.117)]{Ba2}, with last constant $L^{\kappa_0}C$;
\item[(c)] (F) follows from the printed bound \cite[(2.116), p.~582]{Ba2}, $\mathcal P^{(K)}\le O(t_K^{\kappa_0})\abs{\Lambda_7^{(K)}}$, under each of the following readings of its right side: $\Lambda_7^{(K)}=\emptyset$; or $\Lambda_7^{(K)}$ is any set obtained by one further step of the construction of \cite[p.~566]{Ba2}. In both readings $\abs\cdot$ may be either the count $\abs\cdot_K$ or the volume $t_K^3\abs\cdot_K$;
\item[(d)] (a)--(c) hold unchanged with $\rho^{(K)}$ replaced by any non-negative density. In particular they hold for the density of the slice model at $(u,M)$, $u\in[-\eta,1]$, whose photon weights $1+u\delta_i$ are positive.
\end{enumerate}
The constants added by (a)--(c) are $1$ and $L^{\kappa_0}$. They do not depend on $\eps$, $K$, $T$, $u$ or $M$.
\end{lemma}
\begin{proof}
(a) We have $\Lambda_7^{(K-1)\prime}\subset T$, so $\abs{\Lambda_7^{(K-1)\prime}}_K\le\abs T_K=\abs T\,t_K^{-3}=\abs T\eps_0^{-3}$.

(b) In \eqref{app:G1:eq-243K} the factor $\rho^{(K)}$ is the density of a positive measure: it is a product of positive kernels and characteristic functions. The two exponential factors are positive, and $x\mapsto \mathrm{e}^x$ is increasing. So replacing $\mathcal P^{(K)}$ by the right side of (F), which bounds it on the support of $\rho^{(K)}$, does not decrease the right side of \eqref{app:G1:eq-243K}. By (a) this replacement contributes at most $\mathrm{e}^{C\eps_0^{\kappa_0-3}\abs T}$. By \cite[(2.118)]{Ba2}, $\sum_{i<K}O(t_i^\kappa)\abs{T_1^{(i)}}=\sum_{i<K}O(1)t_i^{\kappa_0}\abs T\le c\abs T$. This gives the first display.

For the second statement, note that $t_K^{\kappa_0}=L^{\kappa_0}t_{K-1}^{\kappa_0}$. As point sets, $\Lambda_7^{(K-1)\prime}\subset\Lambda_7^{(K-1)}$, and $t_K\Z^3\subset t_{K-1}\Z^3$ because $L$ is an integer. Hence $\abs{\Lambda_7^{(K-1)\prime}}_K\le\abs{\Lambda_7^{(K-1)}}_{K-1}$, and so $Ct_K^{\kappa_0}\abs{\Lambda_7^{(K-1)\prime}}_K\le L^{\kappa_0}C\,t_{K-1}^{\kappa_0}\abs{\Lambda_7^{(K-1)}}_{K-1}$. This is the last term of the exponent in \cite[(2.117)]{Ba2}. The other terms of that exponent are the collar terms $i=0,\dots,K-1$ of \eqref{app:G1:eq-243K}.

(c) If $\Lambda_7^{(K)}=\emptyset$, the printed bound reads $\mathcal P^{(K)}\le0$, which implies (F). Now let $\Lambda_7^{(K)}$ be produced by one further step. By \cite[p.~566]{Ba2} the sets of each scale are nested, so $\Lambda_7^{(K)}\subset\Lambda_0^{(K)}$, and an admissible set satisfies $\Lambda_0^{(K)}\subset\Lambda_7^{(K-1)\prime}$. Hence $\abs{\Lambda_7^{(K)}}_K\le\abs{\Lambda_7^{(K-1)\prime}}_K$, and (F) holds with the printed constant. In the volume convention, $t_K^3\abs X_K\le\abs X_K$ because $t_K=\eps_0\le1$. So the volume form of the printed bound implies its count form, and hence (F).

(d) The proof of (b) used only $\rho^{(K)}\ge0$, and (a), (c) do not involve the density.
\end{proof}

\noindent Lemma~\ref{app:G1:N10} does not prove (F). On the torus, (F) is the final perturbative bound of the closing step, part (S5) of Proposition~\ref{app:G1:assembly}, and it is obtained there from \cite[Prop.~1]{Ba3}.

\begin{remark}[King's cube terms and field restrictions]\label{app:G1:N10-King}
(i) On the periodic torus $T_\eps$ the set $\partial T^{(k)}$ is empty, and we use the convention $\dist(x,\emptyset)=+\infty$. Consider the terms that \cite[Sect.~5.2, pp.~343--344, (5.41)--(5.49)]{K2} adds for Dirichlet scalar boundary conditions on a cube. The product (5.41) runs over $x\in\Lambda_0^{(k)}\cap\partial T^{(k)}$, so it is empty and equals $1$. The hypothesis $\dist(x,\partial T^{(k)})<p(L^k\eps)$ of (5.44)--(5.45) holds for no $x$. (5.42)--(5.43) are \cite[(2.110)]{Ba2} and its consequence. The hypothesis $\dist(x,\partial T^{(k)})\ge R$ of (5.48) holds at every $x$, and the operator $G_k(A^{(k)})$ there is the periodic one. So (5.46)--(5.49) reduce to the quartic term of \cite[p.~570]{Ba2} combined with \cite[Lemma~2.4, p.~572]{Ba2} (Lemma~\ref{lem:lem24}). Therefore on $T_\eps$ the cube terms add nothing to the estimates of \cite[Chap.~2]{Ba2}.

(ii) The small-field restrictions in the upper bound \eqref{eq:T6} are those of \cite[(2.55), p.~570]{Ba2}. King's restrictions \cite[(3.2), p.~655]{K1} and his modified lower bound \cite[(A.1), p.~676]{K1} enter only the lower bound, Theorem~\ref{thm:UV-K0minus}. The condition $\mu_0t\le1$ of \cite[(A.4)--(A.5), p.~676]{K1} is implied by (E.dom). The bound \eqref{eq:T6} contains no characteristic function of the fields, so none of its thresholds has to match those of the lower bound.
\end{remark}

\begin{lemma}[the omitted estimates of \cite{Ba2,Ba3}]\label{app:G1:items}
Assume the hypotheses of Theorem~\ref{thm:T6}. Then every statement of $\mathcal I$ (Definition~\ref{app:G1:itemlist}) holds at $(u,M)\in\{(1,m_2),(-\eta,m_1)\}$, for every $K\ge K_0$ and every $j\ge1$. Its constants are independent of $\eps$ and $K$. They depend on the torus $T_j$ at most through non-increasing functions of $\ell_{T_j}$, so they are bounded by their values at the torus floor $\ell_0^\natural$ of (T.floor). They are finite at $u=-\eta$ and at $M=m_1$.
\end{lemma}
\begin{proof}
\emph{Step 1: item (a).}
\begin{itemize}
\item The statement of Proposition~\ref{app:G1:I1} includes bounds uniform in $\eps$, $K$ and $\ell_T\ge\ell_0$, and the finiteness at $u=-\eta$, $M=m_1$. Every torus $T_j$ has $\ell_{T_j}\ge\ell_0$. So its constants may be taken to be these uniform bounds, which do not depend on $T_j$.
\item Proposition~\ref{app:G1:I2} holds under condition (T.nw) (Lemma~\ref{app:G1:nowrap}) and introduces no constant.
\item Proposition~\ref{app:G1:I3} is treated in Step 2.
\item Proposition~\ref{app:G1:I4} is a statement on the backgrounds at which bounds are used, and it introduces no constant.
\item The constants $O(\alpha_0,\Nbar;d,L,a,M_b)$ and $\delta(\Nbar)>0$ of Proposition~\ref{app:G1:I5} are independent of $\eps$, $k$, the domains and the background. They are uniform over $\ell_T\ge\ell_0$, with values taken at the floor where $\delta m^2_{\rm fin}$ vertices occur. Its colouring factor is at most $(1+L^2)^{N_v}\le5^{84}$ for $u\in[-\eta,1]$.
\item Proposition~\ref{app:G1:I6} is uniform in $K$, $k$ and the torus.
\item Proposition~\ref{app:G1:I7} holds under (E.m2) and (E.$\Phi$) (stage 10a of Appendix~\ref{app:B}), $\mu_*\eps_0\le1$ (by (E.dom)) and $\eps_0\le\eps_1\le1$. Its constants depend on none of $\eps$, $\eps_0$, $k$, $\ell_T$, $u$, $M$, and the proposition holds at every $M\in(0,m_2]$.
\item Proposition~\ref{app:G1:I8} holds under (E.fin), $\eps_0\le\min(1/m_2,1/\mu_*,L)$, which the hypotheses imply (Lemma~\ref{lem:T6-rows}(b)), and because $\eps_0$ is a power of $L$ (Lemma~\ref{app:G1:N9prem}). It is uniform in $\eps$, $x$ and $j$. The torus enters only through non-increasing functions of $\ell_T$, which are bounded by their values at $\ell_0^\natural$.
\item The constant $C_9$ of Proposition~\ref{app:G1:I9} is independent of $\eps$, $k$, the domains and $\ell_T$. Its image term satisfies $\sum_{k<K}\abs{\cV_k^{\rm img}}\le C\abs T\eps_0^{-3}\mathrm{e}^{-\delta_0\ell_T/\eps_0}\le C\abs T\eps_0^{-3}\mathrm{e}^{-\delta_0\ell_0/\eps_0}$ for $\ell_T\ge\ell_0$, with $C$ independent of $\ell_T$, so per unit volume it is bounded by its value at $\ell_0$. In its proof, Lemma~\ref{app:G1:VS} is applied with $W=B^k(\Lambda_7^{(k)})$, which is Lemma~\ref{app:K:L18-2}.
\end{itemize}

\emph{Step 2: item (b).} The entries of Table~\ref{app:G1:I3tab} are proved in the results listed in Definition~\ref{app:G1:itemlist}(b),(c): in Subsection~\ref{app:G1:part-TI3}, in Lemmas~\ref{app:G1:Rnest}, \ref{app:G1:Rthetab}, \ref{app:G1:Rloc7}, Corollary~\ref{cor:prop1R4}, Propositions~\ref{app:G1:twodomain}, \ref{app:G1:T579}, Lemmas~\ref{app:G1:AA}, \ref{app:G1:AA4}, and Lemmas~\ref{app:G1:I3a-onebg}, \ref{app:G1:I3a-eta1}, \ref{app:G1:I3a-Bloc}, \ref{app:G1:I3a-margins}, \ref{app:G1:I3a-rate}. Proposition~\ref{app:G1:I3} composes them. At the first step $k=1$, the two-domain bounds of Lemma~\ref{app:G1:I3r-DGM} are used with the margins of Lemma~\ref{app:K:L18-1}. Proposition~\ref{app:G1:T579} and Lemma~\ref{app:G1:AA4} are used under the hypotheses stated there. The constants these statements introduce are $\ell_{\Nbar}=2V(\Nbar)$, $C_{\rm loc}(\Nbar)$, $C_{\Nbar}$, the constant of the two-domain bound at margin $R_0+12$, $\beta_\square$, $C_{\rm old}$ and $C_>$ (Definition~\ref{app:G1:rowlist}), the tree rate $\beta_P$ of Proposition~\ref{app:G1:I3}, and the constant of (N2) in Lemma~\ref{app:G1:I3r-AA}. Each of the first seven, except $C_>$, is a function of $d,a,L,M_b,R_0,\delta_0,\Nbar$ only. The constant $C_>$ is a supremum over $t\in(0,1]$, evaluated at $\ell_T=\ell_0^\natural$. The rate $\beta_P$ depends only on the quantities listed at the end of Proposition~\ref{app:G1:I3}. None of these depends on $\eps$, $K$, $u$ or $M$. The constant of (N2) is taken at the two points $(u,M)$ of Theorem~\ref{thm:T6}; it is independent of $\eps$, $k$ and $\ell_T$, and finite there.

\emph{Step 3: items (c)--(f).} The form \eqref{eq:R4ppp} holds at every use under \eqref{eq:eppp} (Corollary~\ref{cor:prop1R4}). Lemmas~\ref{app:G1:Rnest}, \ref{app:G1:Rthetab} and~\ref{app:G1:Rloc7} are statements about domains, supports and localisation widths, and they contain no constant that depends on $\eps$, $K$, $T$, $u$ or $M$. The cap $\Nbar=169$ is part of Proposition~\ref{app:G1:I5}. The final step is Lemma~\ref{app:G1:N10}, whose added constants are $1$ and $L^{\kappa_0}$. It is applied to the non-negative density of the slice model (Lemma~\ref{app:G1:N10}(d)), with (F) supplied by part (S5) of Proposition~\ref{app:G1:assembly}. Lemma~\ref{lem:lem24} holds under its cap (E.edge). Lemma~\ref{lem:IMG} has constants $C_1=c_dC_0$ and $\delta_0$, where $C_0,\delta_0$ are the constants of (H$_\infty$) in Subsection~\ref{ssec:T6-IMG}; these do not depend on $\eps$, $K$, $T$, $u$ or $M$. Under (E.r) and (T.per) its image terms are at most $C_1\mathrm{e}^{-3\delta_0r(t_k)}$. Proposition~\ref{prop:prop1} and Lemma~\ref{lem:Omegafree} concern depths and the dependence on $\Omega$. They introduce no constant beyond those of Proposition~\ref{prop:UV-BaProp21}.

\emph{Step 4: the two points.} For $u\in[-\eta,1]$ the photon weights satisfy $1+u\delta_i\in[1-\eta\bar\delta,1+\bar\delta]$, since $0\le\delta_i\le\bar\delta$ by \eqref{eq:T6-GammaB}. By (E.$\delta$) (Lemma~\ref{lem:T6-rows}(a)), $\bar\delta\le\frac1{12}$. Hence $1+u\delta_i\in[1-\eta/12,\,13/12]\subset[1-\eta,2]$. Every weighted photon propagator is therefore at most $\frac{13}{12}$ times the unweighted one, and $C_B(u)\ge(1-\eta/12)C_{\mu_*}>0$. So no constant of Steps 1--3 blows up at $u=-\eta$.

Consider now the Higgs mass. Wherever it enters from above, it does so through $M^2t^2\le m_2^2\eps_0^2\le1$, by (E.m2). Wherever it enters from below, it does so through quantities that are finite at every fixed $M>0$: $\gamma_1M^2t^2$, $\bar\ell(M)=\min(1/M,1/\mu_*,L)$, $\log(1/(\min(M,\mu_*)\eps_0))$ and the zero-mode gap $M^2(1+4M^2\eps_0^2/(3a))^{-1}$. These are finite at $M=m_1>0$.

\emph{What the cited results use.} The results cited in Steps 1--3 quote, in their proofs, the following statements: Lemmas~\ref{app:G1:S05-TS} and~\ref{app:G1:S05-RF} (in Proposition~\ref{app:G1:I1}); Proposition~\ref{app:G1:I145-prop1base} (in Propositions~\ref{app:G1:I1}, \ref{app:G1:I5} and~\ref{app:G1:I9}); Lemmas~\ref{app:G1:I3r-AA} and~\ref{app:G1:I3r-T579}, beyond Proposition~\ref{app:G1:T579} and Lemmas~\ref{app:G1:AA}, \ref{app:G1:AA4}; Propositions~\ref{app:G1:I3b}, \ref{app:G1:Efin} and \ref{app:K:L04-4}, with the factor $\bar\sigma_N(u)$ (in Proposition~\ref{app:G1:I8}); and Lemma~\ref{app:G1:cutclasses} (in Proposition~\ref{app:G1:I9}). The lemma is proved from the results cited in Steps 1--3 and uses no other input. Combining Steps 1--4 proves the lemma.
\end{proof}

\paragraph{Transfer to the masses of \texorpdfstring{Theorem~\ref{thm:UV-K0plus}}{the upper vacuum bound}.}

\begin{lemma}[transfer of the estimates to $u=0$ and to the masses $M_1,m_1,m_2$]\label{app:G1:RMG}
Let $\pi\in\Wwin$, and let the regulator data be chosen as in Lemma~\ref{lem:window}, with the four-mass rule and the torus floor of Subsection~\ref{p:appA:window}. Assume Statement~\ref{app:K:SL}. Let $M\in\{M_1,m_1,m_2\}$. Then every estimate of $\mathcal I$ (Definition~\ref{app:G1:itemlist}) holds at $(u,M)=(0,M)$, in Ba{\l}aban's model with photon mass $\mu_*=m$. Consequently, Theorem~\ref{thm:UV-K0plus} holds at $M_1=0.7m$, $m_1$ and $m_2$ on every torus of $\mathfrak C$, with its printed proof completed by these estimates.
\end{lemma}
\noindent The proof transfers the dependence on $u$ and on the Higgs mass, as it can be read off the proof of each estimate; no estimate is proved again.
\begin{proof}
\emph{(i) The point $u=0$.} Since $0\in[-\eta,1]$, the point $u=0$ lies in the interval on which the estimates are proved. The weights $1+u\delta_i$ act on the photon fluctuations only, and the scalar side carries no $u$-weight. At $u=0$ every weight equals $1$. Every factor that contains $\abs u$ in the constants of these estimates, including $C_>$ through $(1+\abs u c_k)^{N_v}$, is at $u=0$ at most its value at $u=1$ and at $u=-\eta$. The model at $u=0$ is Ba{\l}aban's own model, so the estimates are statements on his objects.

\emph{(ii) The Higgs mass, estimate by estimate.} The proof of each estimate shows how the Higgs mass enters it:
\begin{itemize}
\item \emph{Type $(\alpha)$:} the remainder bounds of Lemmas~\ref{app:G1:I3r-U1}, \ref{app:G1:I3r-U1m} and~\ref{app:G1:I3r-U1mOmega}, where $m^2$ enters the form $\norm{Df}^2+m^2t^2\norm f^2+a\norm{Qf}^2$ with a $+$ sign; Lemma~\ref{app:G1:I3r-SEP}; the two-domain bounds of Lemma~\ref{app:G1:I3r-DGM} (both routes), with the margins of Lemma~\ref{app:K:L18-1}; Lemmas~\ref{app:G1:I3r-Wsq} and~\ref{app:G1:I3r-W2R}; the statements of item~(c) and Lemmas~\ref{app:G1:I3r-R123}, \ref{app:G1:I3r-Rexh} (either no mass or type $(\alpha)$); the $A\cdot A$ local terms; the region entropy of Proposition~\ref{app:G1:I6}; and Proposition~\ref{app:G1:I7} (the condition (E.md)).
\item \emph{Type $(\beta)$:} Proposition~\ref{app:G1:T579} (through the $\delta m^2$ vertices); Proposition~\ref{app:G1:I4} (only through $-\Delta_A+m^2>0$); Proposition~\ref{app:G1:I8} (through $\bar\ell(M)$, $\rho_{\min}(M)$, $Z_f$ and $\log(1/(\min(M,\mu_*)\eps_0))$); the steps of Subsection~\ref{app:G1:part-TS05} (through the condition (E.5) and Lemma~\ref{app:G1:RW-MASS}); the closing step (Lemmas~\ref{app:G1:S05-CL} and~\ref{app:G1:RW-ZMs}, with the gap $\ge M^2(1+4M^2\eps_0^2/(3a))^{-1}$); the constant $C_\delta(M)$, which is finite at every $M>0$; $C_b'$; and $C_{F1}$ and $\mathfrak E_k$ of Proposition~\ref{app:G1:assembly}.
\item \emph{No mass entry:} the cube remainder of \cite[p.~433]{Ba3}; the order cap $\Nbar=169$; the Ward--Takahashi step on the torus, Proposition~\ref{app:G1:I2} (King's counterterm at $M_1$ on $M_1$-lines; Lemma~\ref{app:G1:ctid} holds at every $M>0$); and the final step of the upper bound, Lemma~\ref{app:G1:N10}. That lemma's added constants are $1$ and $L^{\kappa_0}$, and it holds for every non-negative density (Lemma~\ref{app:G1:N10}(d)), in particular for that of Ba{\l}aban's model at $u=0$ and Higgs mass $M$.
\end{itemize}
Entries of type $(\alpha)$ hold at $M\le m_2$, because (E.m2) and $t\le\eps_0$ give $M^2t^2\le m_2^2\eps_0^2\le1$. The left-hand sides of the upper bounds (E.m2) and (E.$\Phi$) increase in $M$, and their constants do not depend on $M$. So imposing them at $m_2$ covers every $M\le m_2$, in particular $M_1$. Entries of type $(\beta)$ are finite at each fixed $M>0$, and they enter only conditions that the data of Lemma~\ref{lem:window} meet at every $M\in\mathbb M_4$. Consider an upper bound on $\eps_0$ in which the mass enters from below: (E.5), the $\eps_0$-window of the closing step, Lemma~\ref{app:G1:RW-ZMs}, or a condition of Lemma~\ref{app:G1:RW-logp} containing $\log_+(1/(M_1\eps_0))$. Each is a supremum, over $t\le\eps_0$, of a quantity that tends to $0$ as $\eps_0\downarrow0$ at fixed $M_1>0$. Each is therefore a downward-closed upper bound, and the minimum of four positive thresholds is positive. These bounds bind first at the smallest mass $M_1$.

The constants of the stages 8--9 that carry $M$ are $c_7\ni C_\delta(M)^2/(16\lambda)$, $c_1^I$ and the lower bounds on $b_0$. They are evaluated at the four masses and combined by the maximum; this is the four-mass rule. Their boundedness for $\eps_0\in(0,\eps_1]$ is Statement~\ref{app:K:SL}. The constants that depend on the torus are taken at $\ell_0^\natural=\max\{8,128/m_1,128/M_1\}=128/M_1$; this is the torus floor. They are non-increasing in $\ell_T$, so their values at this floor dominate their values on every torus of $\mathfrak C$.

\emph{(iii) No entry of type $(\gamma)$.} No estimate of $\mathcal I$ orders $M$ against $\mu_*$ or uses $m_1>2m$. Such orderings occur, for instance, in Theorem~\ref{thm:ch-W} and in the anchor condition \eqref{p:ch-anchor}. Those are used at $m_1,m_2$ and are not estimates of $\mathcal I$. The order of the two field masses flips in only one place: Proposition~\ref{app:G1:I8} at $M_1<\mu_*$ (by Lemma~\ref{lem:B:massrange}(a); see (v)), where the scalar becomes the light field. The construction there is symmetric. Both fields are sliced to one common index $\bar K$ with $t_{\bar K}\le\bar\ell(M)=\min(1/M,1/\mu_*,L)$, and the mass ratio enters only through $\rho_{\min}(M)$. At $(M_1,\mu_*)=(0.7,1)m$ and $L=2$, in units $m=1$, we get $\bar\ell(M_1)=1$ and $\rho_{\min}(M_1)\ge0.7\cdot\min(1,\tfrac12)=0.35>0$. The hypothesis $\eps_0\le\bar\ell(M_1)$ follows from (E.fin), and $\eps_0$ is a power of $L$ by (T.div). Theorem (1.3) of \cite[p.~556]{Ba2} asks for no relation between the Higgs mass and $\mu_0$. It is stated for the model of \cite[p.~555]{Ba2}, at the masses $m^2$, $\mu_0^2$ and coupling $\lambda$ there, with constants $E_\pm$ independent of $\eps$ and $T_\eps$.

\emph{(iv) Domains, supports and structure.} The statements of item~(c), the end-point statements of Lemmas~\ref{app:G1:I3r-HOm2} and~\ref{app:G1:I3r-HOm2ch3}, and Lemma~\ref{app:G1:I3r-ME} are statements about domains, supports, localisation widths and the insertion structure of the vertices. They contain no constant, and the value of $M$ does not occur in them. For each of its items, Lemma~\ref{app:G1:I3r-Pw} names the place of the argument and the conditions used. The places are entries of Table~\ref{app:G1:I3tab}, Lemmas~\ref{app:G1:M4D}, \ref{lem:lem24} and~\ref{app:G1:RW-KSP}, and Propositions~\ref{prop:prop1} and~\ref{app:K:L10-2}. The conditions are (E.md), (R.edge), (E.edge) and those of Lemma~\ref{app:G1:RW-KSP}, all of which belong to Definition~\ref{app:G1:rowlist}. Lemma~\ref{app:G1:I3r-Pw} therefore transfers together with the statements at these places.

\emph{(v) Statement ranges.} Three ingredients are stated at the masses of Theorem~\ref{thm:T6}: Proposition~\ref{app:G1:I8} (at $M\in\{m_1,m_2\}$), the $A\cdot A$ local terms ($M\in[m_1,m_2]$) and Lemma~\ref{app:G1:RW-MASS} (in the regime of $m_1$). By their proofs they hold at every $M\in(0,m_2]$. Proposition~\ref{app:G1:T579} is stated at every $M\in(0,m_2]$. Together this is Lemma~\ref{lem:B:massrange}(a)--(d). No proof step is redone.

\emph{(vi) The estimates of item~(f).} For them the transfer can be read off their statements. The cap \eqref{eq:lem24-cap}, condition (E.edge), does not depend on the Higgs mass: in the proof of Lemma~\ref{lem:lem24} the mass occurs only together with the coefficients $a_k$, and the cap involves only $a_k$, which do not depend on $M$. So it is compatible with the minimum over $\mathbb M_4$, and Lemma~\ref{lem:lem24} holds at $(0,M)$ for each $M\in\{M_1,m_1,m_2\}$ under the same cap. Lemma~\ref{lem:IMG} has constants $C_1=c_dC_0$ and $\delta_0$, where $C_0,\delta_0$ are the constants of (H$_\infty$) in Subsection~\ref{ssec:T6-IMG}. For the photon kernels these are the constants at the photon mass $\mu_*=m$, which is the same at every $(u,M)$. For the scalar kernels they depend only on $a$, $L$ and $\alpha$, at every $M\in(0,m_2]$ with $M\eps_0\le1$ (Proposition~\ref{app:K:L19-1}(ii)), and so at $M_1$, $m_1$ and $m_2$ by (E.m2). They do not depend on $u$. The torus enters only through its side, and under (E.r) and (T.per) the image terms are at most $C_1\mathrm{e}^{-3\delta_0r(t_k)}$ at every $(u,M)$. Proposition~\ref{prop:prop1} and Lemma~\ref{lem:Omegafree} concern depths and the dependence on $\Omega$. They introduce no constant beyond those of Proposition~\ref{prop:UV-BaProp21}, which is stated for Ba{\l}aban's model, that is, at $u=0$. In the written steps of \cite[Prop.~1]{Ba3} the only mass that occurs is the term $+m^2$ of the operator $-\Delta^{\eta,N}+m^2+a_kP_k$, which is of type $(\alpha)$. Hence the three estimates of item~(f) hold at $(0,M)$ for every $M\in\{M_1,m_1,m_2\}$.

\emph{Conclusion.} By Lemma~\ref{app:G1:items}, the estimates of $\mathcal I$ hold at $(1,m_2)$ and $(-\eta,m_1)$. By (i)--(vi) they hold at $(0,M)$ for $M\in\{M_1,m_1,m_2\}$, under the four-mass rule and the torus floor. Both rules only shrink $\eps_0$ or raise $\ell_0$, so Lemma~\ref{lem:window} can meet them. These estimates are those of Remark~\ref{rem:UV-blockD}, so Theorem~\ref{thm:UV-K0plus} holds at these masses.
\end{proof}

\begin{remark}\label{app:G1:RMG-scope}
Proposition~\ref{prop:ch-K0} states the upper bound only at $M\in\{M_1,m_1,m_2\}$, and no upper bound at $M=m$ is used. The chain uses only $E_+(M_1)$, in Proposition~\ref{prop:ch-Pplus}. Only $(0,M_1)$ lies on the path to Theorem~\ref{thm:ch-T9}; the masses $m_1,m_2$ are anchors off that path.
\end{remark}

\paragraph{The hypotheses of Theorem~\ref{thm:ch-T9}.}

\begin{proposition}[the hypotheses of Theorem~\ref{thm:ch-T9} on $\mathfrak C$]\label{app:G1:VACT}
Let $(e,\lambda,m)$ be as in Theorem~\ref{thm:main}, take $\pi=\pi(e,\lambda,m)$ and the data of Lemma~\ref{app:G1:window-apply}, and let $(\eps,T_j)\in\mathfrak C$. Then every hypothesis of Theorem~\ref{thm:ch-T9} holds:
\begin{enumerate}
\item[(1)] (the hypotheses of Theorem~\ref{thm:T6}) the conditions of Appendix~\ref{app:B}, including \eqref{eq:lem24-cap} and \eqref{eq:GIMG}, hold, and $\eta\bar\delta<1/12$; hence \eqref{eq:T6} holds at $(1,m_2)$ and $(-\eta,m_1)$ on every torus of $\mathfrak C$;
\item[(2)] (the slice vacuum bound) $\Delta_1^B<\infty$, the supremum being taken over $\mathfrak C$, with
$\Delta_1^B\le C_+^B(1;m_2)+\eta^{-1}C_+^B(-\eta;m_1)+(1+\eta^{-1})E_-(m_1)+\tfrac12(1+\eta)E_2(m_1)+\bar\Gamma^{(1)}_{21}$;
\item[(3)] (the anchor condition \eqref{p:ch-anchor}) $b_T+2H\le a_{K,T}(m_1)$ and $g_{21}>0$ on every torus of $\mathfrak C$; hence Proposition~\ref{prop:ch-Tflat} applies;
\item[(4)] (the ``$+$'' side) the hypotheses of Proposition~\ref{prop:ch-Pplus} hold, with $g\ge0.319m^2$; hence Propositions~\ref{prop:ch-Pplus} and \ref{prop:ch-T7} hold for every $j$ and every $K$ of $\mathfrak C$;
\item[(5)] (the curvature bound) Lemma~\ref{lem:C-Fsector} holds on every torus of $\mathfrak C$ at the coefficient $b_{T_j}$.
\end{enumerate}
Consequently Theorem~\ref{thm:ch-T9} holds on $\mathfrak C$ with $\sigma_+=g/2$, $H=m^2$ and $\mathsf K=\max(\kappa_+,2S^B_\flat)$.
\end{proposition}
\begin{proof}
(1) Lemma~\ref{app:G1:window-apply}(b) gives the conditions of Appendix~\ref{app:B}. Theorem~\ref{thm:T6} also uses $\eta\bar\delta<1$. It follows from the condition (E.$\delta$), $\kappa\mu_*^2\eps_0^2\le1/12$, imposed at stage~10b, that is $\bar\delta\le1/12$; hence $\eta\bar\delta<1/12$ for $\eta<1$. (Alternatively, (E.dom), $\mu_*\eps_0\le0.2$, together with $\kappa\le7.5$ gives $\bar\delta\le0.3$.) With $\Gamma^B\le\bar\delta\,C_{\mu_*}$ from \eqref{eq:T6-GammaB} this gives $C_B(-\eta)=C_{\mu_*}-\eta\Gamma^B\ge(1-\eta\bar\delta)C_{\mu_*}>0$, so the slice model at $u=-\eta$ is well defined and $\Phi_{m_1}$ is convex on $[-\eta,1]$. The spacings of $\mathfrak C$ satisfy $K\ge K_0=\max(K_0',n_0+1)$, as Theorem~\ref{thm:T6} requires.

(2) Apply Proposition~\ref{prop:ch-T5}. Its hypotheses hold:
\begin{itemize}
\item Theorem~\ref{thm:T6} at $(1,m_2)$ and $(-\eta,m_1)$: by (1).
\item Theorem~\ref{thm:UV-K0minus} at $m_1$, in the form $\abs T^{-1}\Phi_{m_1}(0)+e(\eps;0;m_1)\ge-E_-(m_1)$: this is the lower bound of Proposition~\ref{prop:ch-K0} at $m_1$, by Lemma~\ref{app:G1:vacnorm}(b), read with the order-zero term $E_{0,s}(m_1)$ of Remark~\ref{rem:T6-textfixes}(i). It needs none of the estimates of Lemma~\ref{app:G1:items}.
\item The bound \eqref{p:ch-E2B} at $m_1,m_2$ (Lemma~\ref{app:G1:E2B}). The hypotheses of that lemma hold at $M=m_1$ and at $M=m_2$. Those of Theorem~\ref{thm:T6} hold by (1). The condition $2\nbar+1\le3R$ holds because Appendix~\ref{app:B} fixes $\nbar=13$ and $R\ge107$. (H$_\infty$) holds for the kernels at the masses $m_1$, $m_2$ and $\mu_*$ by Proposition~\ref{app:G1:Hinf-M}, with the photon taken componentwise (Lemma~\ref{app:G1:VSR-V0}). The bound is a supremum over the whole interval $u\in[-\eta,1]$, not only over the two points. What it inherits from Propositions~\ref{app:G1:I1}--\ref{app:G1:I5} holds at $m_1,m_2$ uniformly over the weight class $\mathbb W_0=[1-\eta,2]$ of the slice weights $w_j=1+u\delta_j$, i.e.\ for all $u\in[-\eta,1]$.
\item The bound \eqref{p:ch-CNu} at $(m_1,m_2)$: Proposition~\ref{app:G1:CNuB}, given Proposition~\ref{prop:ch-CN} and \cite{Ba3}. Its hypotheses are those of Lemma~\ref{app:G1:E2B} at $m_1$ and at $m_2$, which were checked in the previous item.
\end{itemize}
Both instances of \eqref{eq:T6} and both c-number bounds are used at the masses $m_1,m_2$ of Theorem~\ref{thm:T6}, so no transfer to other masses is needed on this side. Proposition~\ref{prop:ch-T5} then gives the bound in (2).
The conversion between King's Lebesgue-normalised photon in $Z_K$ and the normalised photon in $\Phi_M(0)$ is Lemma~\ref{app:G1:vacnorm}(b).

(3) Every $(\eps,T_j)\in\mathfrak C$ has $\eps\le1/m$ and $\ell_j\ge96/m$, by Lemma~\ref{app:G1:window-apply}(b), so Theorem~\ref{thm:ch-W}(b),(c) apply. Theorem~\ref{thm:ch-W}(c) reads $a_{K,T}(m_1)\ge a_{K,T}(m)+2H=b_T+2H$. Theorem~\ref{thm:ch-W}(b) gives $g_{21}\ge6.6369\,m^2>0$. With (2), these are the hypotheses of Proposition~\ref{prop:ch-Tflat}.

(4) Take $M_2=m$ and the anchor $M_1=0.7m$. Then:
\begin{itemize}
\item $g\ge2\sigma_+\ge0.319m^2$ by Theorem~\ref{thm:ch-W}(a).
\item The upper bound of Theorem~\ref{thm:UV-K0plus} at $M_1$ holds on every $T_j$ by Lemma~\ref{app:G1:RMG}, under Statement~\ref{app:K:SL}.
\item The lower bound of Theorem~\ref{thm:UV-K0minus} at $m$ holds by that theorem.
\item $\bar\Gamma(M_1,m)<\infty$ by Proposition~\ref{prop:ch-CN}.
\item The coefficient $b_T=a_{K,T}(m)$ is the one at which Ba{\l}aban's theorems are stated on the torus, by Proposition~\ref{prop:UV-B1e}.
\item Each $T_j$ is a torus of the form \cite[(1.2), p.~604]{Ba1} (Subsection~\ref{ssec:UV-torus}).
\end{itemize}
Theorems~\ref{thm:UV-K0minus} and~\ref{thm:UV-K0plus} are stated for Ba{\l}aban's step count $K'-n_0\ge1$, which holds on $\mathfrak C$ because $K\ge n_0+1$. Proposition~\ref{prop:ch-Pplus} follows, and with it Proposition~\ref{prop:ch-T7}: the states of $\mathfrak C$ have the coefficient $b_{T_j}$, independent of $\mu_0$; $N$ is even; and the reflection positivity of Proposition~\ref{prop:C-RP} holds.

(5) The curvature bound of Lemma~\ref{lem:C-Fsector} holds on every torus, for every real source of $\wick{\varphi^2}$ and every bare coefficient. In particular it holds at $b_{T_j}$.

Steps (1)--(6) of the proof of Theorem~\ref{thm:ch-T9} use exactly Proposition~\ref{prop:ch-T7} at $q=2k_+$, Proposition~\ref{prop:ch-Tflat} at $h=2k_-$ and Lemma~\ref{lem:C-Fsector}, each on one torus at the coefficient $b_{T_j}$. By (2)--(5) all three hold on $\mathfrak C$. The constants $\sigma_+$, $\kappa_+$, $S^B_\flat$ do not depend on $\eps$, $j$ or $\mu_0$.
\end{proof}

\section{Measure-theoretic facts}\label{app:A}
The following three facts are used in Section~\ref{sec:OS4}: a theorem of Lusin and Souslin on injective Borel maps, the density of a generating algebra of bounded functions in $L^2$, and the reflection positivity of a Gaussian measure whose covariance is reflection positive.

\begin{lemma}[Lusin--Souslin]\label{lem:A:LS}
Let $X$ be Polish and $\iota:X\to\R^\N$ continuous and injective. Then $\iota$ is a Borel isomorphism onto its image, which is Borel; in particular the Borel $\sigma$-algebra of $X$ is generated by the coordinates of $\iota$.
\end{lemma}
\begin{proof}
\cite[Thm.~15.1 and Cor.~15.2]{Kec95}.
\end{proof}

\begin{lemma}[density of generating algebras]\label{lem:A:mc}
Let $m$ be a probability measure on a measurable space, and $\cAc$ an algebra of bounded complex functions, closed under complex conjugation and containing the constants. Then $\cAc$ is dense in $L^2(\sigma(\cAc),m)$. The same holds for a countable $\Q[i]$-algebra $\cAc_\Q$, closed under conjugation and containing $1$, whose complex span is such an algebra.
\end{lemma}
\begin{proof}
Let $\mathcal K$ be the set of bounded real $\sigma(\cAc)$-measurable functions lying in the $L^2(m)$-closure of the real parts of elements of $\cAc$. It is a vector space containing the constants and $\operatorname{Re}\cAc$ (which is an algebra, being the real algebra generated by real and imaginary parts), and it is closed under bounded monotone limits by dominated convergence. By the functional monotone class theorem \cite[Ch.~I, Thm.~21]{DM78} it contains all bounded real $\sigma(\cAc)$-measurable functions, which are dense in $L^2(\sigma(\cAc),m)$. Complex functions are treated by real and imaginary parts, and $\cAc_\Q$ is $L^2$-dense in its complex span.
\end{proof}

\begin{lemma}[reflection positivity of Gaussian laws]\label{lem:A:gaussRP}
Let $\mu$ be a centred Gaussian measure on $X$ with covariance $C$, invariant under a reflection $\theta$ with $\theta(\Sig(\Hp{j}{0}))=\Sig(\Hm{j}{0})$, and suppose $\ip{\theta f}{Cf}\ge0$ for all real $f$ supported in $\Hp{j}{0}$. Then $\mu$ satisfies \eqref{eq:D:RP} for $c=0$.
\end{lemma}
\begin{proof}
For real $f_1,\dots,f_N$ supported in $\Hp{j}{0}$ and $\Theta_k=\mathrm{e}^{i\Phi(f_k)}$,
\[
\E\big[\overline{\Theta_k\circ\theta}\,\Theta_l\big]=\mathrm{e}^{-\frac12\ip{f_k}{Cf_k}}\,\mathrm{e}^{-\frac12\ip{f_l}{Cf_l}}\,\mathrm{e}^{\ip{\theta f_k}{Cf_l}},
\]
using $\ip{\theta f}{C\theta f}=\ip{f}{Cf}$. The kernel $(k,l)\mapsto\ip{\theta f_k}{Cf_l}$ is positive semidefinite (it is the form of the hypothesis, polarised, on the span of the $f_k$), so by the Schur product theorem its entrywise exponential is positive semidefinite, and so is the displayed matrix. Hence \eqref{eq:D:RP} holds on the span of such $\Theta_k$, and by Lemma~\ref{lem:A:mc} on all of $L^2(\Sig(\Hp{j}{0}),\mu)$.
\end{proof}
\section{Statements used without proof in this version}\label{app:K}

The proofs of Theorems~\ref{thm:main} and \ref{thm:T6} use the statements collected in this appendix. Their proofs are not included in this version. Each statement is in the notation of the place where it is used, and is followed by a list of these places. Statements that are printed with a proof in the literature are cited where they are used and are not repeated here.

\begin{kstatement}[the per-step constants of the large-field estimate are free of $\eps_0$]\label{app:K:SL}
Let $\pi\in\fW$ and let $\eps_1$ be as in Lemma~\ref{app:G1:L10asup}. Each of the following estimates holds at every scale $k\le K$, given the inductive hypothesis at scale $k$, with a constant per site of at most $F(t_k)$: the replacements \cite[(2.47)--(2.48)]{Ba2} and their inverse \cite[(3.1)]{Ba2}; \cite[Prop.~2.1, (2.57)]{Ba2}, with the vacuum part on the collar bounded by \cite[Prop.~1]{Ba3}; the mass-counterterm pieces of step $k$ and the layer pieces of Lemma~\ref{app:G1:S05-RF}; \cite[Prop.~3.1, (3.26)]{Ba2} in the form of Lemma~\ref{app:G1:M4B}; and the normalisations and the rescaling \cite[(3.32)--(3.36)]{Ba2}. Here $F:(0,1]\to[0,\infty)$ satisfies $\sup_{0<t\le\eps_1}F(t)<\infty$, and it depends on the structural data and on $\pi$ but not on $\eps_0$, $K$, $k$, the torus, the region sequence, the value of $u$ or the mass $M\in\{M_1,m,m_1,m_2\}$. In particular the constants $c_5$ and $c_7$ of \cite[(3.42)]{Ba2} satisfy $\sup_{0<\eps_0\le\eps_1}c_5<\infty$ and $\sup_{0<\eps_0\le\eps_1}c_7<\infty$.
\end{kstatement}
\noindent This statement is used in Proposition~\ref{app:G1:I6}, Lemma~\ref{app:G1:RW-BM1} (the condition (E.ent) of stage 10b) and Lemma~\ref{app:G1:L10asup}. It is the form in which \cite{Ba2} states its per-step bounds (per block scale, with constants independent of the scale, for small running couplings, for instance \cite[Prop.~2.2, p.~570; Prop.~3.1, p.~589]{Ba2}); it cannot be derived from the printed text alone, because the proofs of \cite[Prop.~2.1, (2.57)]{Ba2} and of \cite[Prop.~3.1]{Ba2} at a nonzero background are not written there with the dependence of the constants on $\eps_0$.

\paragraph{The final c-numbers.}

\begin{kstatement}[the vacuum output for Ba{\l}aban's scalar staging]\label{app:K:L04-6}
Consider Ba{\l}aban's staging of the scalar integrals \cite[(1.25), (1.27)]{Ba3}, in which the fluctuation covariance of step $k$ depends on the gauge background, with the truncation of the effective action at order $\nbar$ in the couplings. Then the part of order $\le\nbar$ of the output of step $k$ at zero block fields and zero background equals $\cV_k$ of Definition~\ref{app:G1:vacgrade}, the slices being evaluated at zero background.
\end{kstatement}
\noindent This statement is used in the proof of Lemma~\ref{app:G1:S05-TS}.

\paragraph{The coefficients of the $A\cdot A$ vertices.}

\begin{kstatement}[inner renormalisation]\label{app:K:L05-3}
Let $0\le m^2_{j_0}\le1$, $(\alpha,\beta)\in\{(4,0),(2,1)\}$, let the photon lines be as in Lemma~\ref{app:G1:AA4}, and let $\Lambda$ be a label configuration. For a member $T$ of $\Pi^\Lambda_G$ and $\emptyset\ne S$, let the critical subgraphs of $T^S$, their types and the sets $\mathcal K(T)$ be as in Definition~\ref{app:G1:Rm-crit}; chains are as in Proposition~\ref{app:G1:Rm-chain}. Then there is a constant $C_N$ with the following properties, uniformly in $\Lambda$, the bond $b$, $\xi$, $j_0$, $k$, $\eps$, $\ell_T$ and $m^2_{j_0}\in[0,1]$:
\begin{itemize}
\item[(N.2)] If $T$ is not a chain and $T^S$ has a critical subgraph of type 2 (exactly two external legs, both scalar, no difference of a vertex of the subgraph acting on them), then $\bigl|\sum_{T'\in\mathcal K(T)}\sum_{b'}\xi^3T'^S(b,b')\bigr|\le C_N$.
\item[(N.3)] The same bound holds if $T$ is not a chain and $T^S$ has a critical subgraph of type 3 (exactly two external legs, both scalar, a difference acting on exactly one of them).
\item[(N.5)] The same bound holds if $T$ is not a chain and $T^S$ has a critical subgraph of type 5 (exactly three external legs, two scalar and one vector).
\end{itemize}
$C_N$ depends only on the order, $d$, $a$, $L$, the block size, the constants of the slice bounds of the scalar and photon lines (the latter independent of $\mu_*$, with weights in $[0,2]$), and $\bar c_{\rm fin}$.
\end{kstatement}
\noindent This statement is used in the proof of Theorem~\ref{app:K:L05-2}, which is used in the proof of Lemma~\ref{app:G1:AA4}.

\paragraph{The tail beyond the order $\nbar$.}

\begin{kstatement}[exceptional vertices with zeroth-order legs at the first step]\label{app:K:X0}
At $k=0$, the bound (X) of Lemma~\ref{app:G1:T579-po} holds for the vertices of coupling order one whose two scalar legs are external fields and neither of which is a chain.
\end{kstatement}
\noindent This statement is used in Proposition~\ref{app:G1:T579} and in Step 3a of its proof; for $1\le k\le K-1$ the same bound is Proposition~\ref{app:K:L06-3} together with Statement~\ref{app:K:L06-4}.

\begin{kstatement}[exceptional vertices with zeroth-order legs: contour-type partner]\label{app:K:L06-4}
Assume the setting and the hypotheses (a)--(c) and (iii) of Proposition~\ref{app:K:L06-3}. Let $F$ be a connected localised family of $\mathcal P^{(k)}_w$ with at least two vertices containing a vertex $v$ of coupling order one whose scalar legs are $\Phi_0$ or $\phi(y)$, and whose vector leg is contracted with a contour leg $A'(\Gamma^{(k)}_{y',x'})$ of a factor \cite[(1.14)]{Ba3} or \cite[(1.15)]{Ba3} of a product \cite[(1.18)]{Ba3}. Then $F$ obeys the bound \cite[(1.33)]{Ba3}, with the factor $(\abs uc_k)^m$ and with some decay rate $\delta_1>0$ independent of $\eps$, $k$, $\ell_T$, in which the factor $e(t)\norm{h\Phi_1}_{1,\alpha_0}\norm{h\Phi_2}_{1,\alpha_0}$ of $v$ is replaced by $C\,C_{\rm ex}t^{1/4}p^2$, with $C$ independent of $\eps$, $k$ and $\ell_T$.
\end{kstatement}
\noindent This statement is assumed in Proposition~\ref{app:G1:T579} and used in Step 3a of its proof.

\paragraph{The conditions of Ba{\l}aban's Proposition~1 at its applications.}

\begin{kstatement}[the remaining abbreviated steps]\label{app:K:L10-1}
Assume the hypotheses of Theorem~\ref{thm:T6}, and let the regulator data satisfy the conditions of Definition~\ref{app:G1:rowlist}. Let $s$ be an estimate, identity or property with the following three features: it is asserted on \cite[pp.~555--594]{Ba2}, on \cite[pp.~411--445]{Ba3}, or on a page of \cite{Ba1} invoked on those pages; its argument is not printed there; and it is not one of the items (a)--(i) of Lemma~\ref{app:G1:I3r-Pw}. Then $s$ holds at every scale $k$ and at every configuration at which it is asserted, with constants independent of $\eps$, $k$ and $\ell_T$.
\end{kstatement}
\noindent This statement is used in Remark~\ref{app:G1:I3a-complete} (the completeness of the list of applications in Definition~\ref{app:G1:I3a-uses}, on which Proposition~\ref{app:G1:I3} is stated, and of the list in Lemma~\ref{app:G1:I3r-Pw}), and through the printed steps of \cite{Ba2,Ba3} in the proof of Theorem~\ref{thm:T6}.

\paragraph{The c-number bounds: the weighted-line form.}

\begin{kstatement}[pointwise partner bound]\label{app:K:L17-4}
Assume the hypotheses of Theorem~\ref{app:K:L17-1}. Work at the top scale $k=K$, in scale-$K$ units and at zero background. Let $G_0$ be a $C$-only self-energy graph of \cite[(1.21)--(1.23)]{Ba3}, with two scalar legs and order $e^\alpha\lambda^\beta$, whose lines are the kernels $G_{K,s}(0)$, $G_{K,v}(0)$. For $x\in T$, let $\delta m^2_{G_0,T}(x)$ be its counterterm coefficient on $T$ (the expression with leg $1$ fixed at $x$ and leg $2$ summed), and let $\delta m^2_{G_0,\infty}$ be the same expression on $\eta_K\Z^3$. Then, for every $x\in T$,
\[
\abs{\delta m^2_{G_0,T}(x)-\delta m^2_{G_0,\infty}}\le C_P\,e(t_K)^\alpha\lambda(t_K)^\beta,
\]
with $e(t)=e\,t^{1/2}$, $\lambda(t)=\lambda t$ and $C_P=C_P(\nbar,\alpha_0,L,a,M_b,\delta_0)$ independent of $\eps$, $x$ and $\ell_T\ge\ell_0$. Consequently, each torus partner of a class equals its form on $\eta_K\Z^3$ plus a finite mass vertex of the same order whose coefficient obeys this bound.
\end{kstatement}
\noindent This statement is used in the proof of Lemma~\ref{app:G1:E2B}, through Theorem~\ref{app:K:L17-1}, for the counterterm partners on the torus.

\begin{thebibliography}{GJ87}

\bibitem{Ba1} T.~Ba{\l}aban, (Higgs)$_{2,3}$ quantum fields in a finite volume. I. A lower bound, Commun. Math. Phys. \textbf{85} (1982), 603--626.
\bibitem{Ba2} T.~Ba{\l}aban, (Higgs)$_{2,3}$ quantum fields in a finite volume. II. An upper bound, Commun. Math. Phys. \textbf{86} (1982), 555--594.
\bibitem{Ba3} T.~Ba{\l}aban, (Higgs)$_{2,3}$ quantum fields in a finite volume. III. Renormalization, Commun. Math. Phys. \textbf{88} (1983), 411--445.
\bibitem{Ba4} T.~Ba{\l}aban, Regularity and decay of lattice Green's functions, Commun. Math. Phys. \textbf{89} (1983), 571--597.
\bibitem{BBIJ} T.~Ba{\l}aban, D.~Brydges, J.~Imbrie, A.~Jaffe, The mass gap for Higgs models on a unit lattice, Ann. Physics \textbf{158} (1984), 281--319.
\bibitem{BCGNOPS} G.~Benfatto, M.~Cassandro, G.~Gallavotti, F.~Nicol\`o, E.~Olivieri, E.~Presutti, E.~Scacciatelli, Some probabilistic techniques in field theory, Comm. Math. Phys. \textbf{59} (1978), 143--166.
\bibitem{BIJ} T.~Ba{\l}aban, J.~Imbrie, A.~Jaffe, Effective action and cluster properties of the abelian Higgs model, Commun. Math. Phys. \textbf{114} (1988), 257--315.
\bibitem{BFS1} D.~Brydges, J.~Fr\"ohlich, E.~Seiler, On the construction of quantized gauge fields. I. General results, Ann. Physics \textbf{121} (1979), 227--284.
\bibitem{BFS2} D.~Brydges, J.~Fr\"ohlich, E.~Seiler, Construction of quantised gauge fields. II. Convergence of the lattice approximation, Commun. Math. Phys. \textbf{71} (1980), 159--205.
\bibitem{BFS3} D.~Brydges, J.~Fr\"ohlich, E.~Seiler, On the construction of quantized gauge fields. III. The two-dimensional abelian Higgs model without cutoffs, Commun. Math. Phys. \textbf{79} (1981), 353--399.
\bibitem{CCHS} A.~Chandra, I.~Chevyrev, M.~Hairer, H.~Shen, Stochastic quantisation of Yang--Mills--Higgs in 3D, Invent. Math. \textbf{237} (2024), 541--696; arXiv:2201.03487.
\bibitem{CS25} I.~Chevyrev, H.~Shen, Uniqueness of gauge covariant renormalisation of stochastic 3D Yang--Mills--Higgs, arXiv:2503.03060.
\bibitem{DM78} C.~Dellacherie, P.-A.~Meyer, \emph{Probabilities and Potential}, North-Holland Mathematics Studies \textbf{29}, North-Holland, Amsterdam, 1978.
\bibitem{Di15} J.~Dimock, Nonperturbative renormalization of scalar quantum electrodynamics in $d=3$, J. Math. Phys. \textbf{56} (2015), 102304; arXiv:1502.02946.
\bibitem{Di18} J.~Dimock, Ultraviolet regularity for QED in $d=3$, J. Math. Phys. \textbf{59} (2018), 012301; arXiv:1512.04373.
\bibitem{Dr87} B.~K.~Driver, Convergence of the $U(1)_4$ lattice gauge theory to its continuum limit, Commun. Math. Phys. \textbf{110} (1987), 479--501.
\bibitem{EIC26} Euclidean invariance and clustering for the curvature field of the three-dimensional lattice abelian Higgs model, preprint (2026).
\bibitem{FILS78} J.~Fr\"ohlich, R.~Israel, E.~H.~Lieb, B.~Simon, Phase transitions and reflection positivity. I. General theory and long range lattice models, Commun. Math. Phys. \textbf{62} (1978), 1--34.
\bibitem{FSS76} J.~Fr\"ohlich, B.~Simon, T.~Spencer, Infrared bounds, phase transitions and continuous symmetry breaking, Commun. Math. Phys. \textbf{50} (1976), 79--95.
\bibitem{GJ} J.~Glimm, A.~Jaffe, \emph{Quantum Physics. A Functional Integral Point of View}, 2nd ed., Springer, New York, 1987.
\bibitem{GM82} M.~G\"opfert, G.~Mack, Proof of confinement of static quarks in 3-dimensional $U(1)$ lattice gauge theory for all values of the coupling constant, Commun. Math. Phys. \textbf{82} (1982), 545--606.
\bibitem{Gr83} L.~Gross, Convergence of $U(1)_3$ lattice gauge theory to its continuum limit, Commun. Math. Phys. \textbf{92} (1983), 137--162.
\bibitem{Jan97} S.~Janson, \emph{Gaussian Hilbert Spaces}, Cambridge Tracts in Mathematics \textbf{129}, Cambridge University Press, 1997.
\bibitem{Kec95} A.~S.~Kechris, \emph{Classical Descriptive Set Theory}, Graduate Texts in Mathematics \textbf{156}, Springer, New York, 1995.
\bibitem{K1} C.~King, The $U(1)$ Higgs model. I. The continuum limit, Commun. Math. Phys. \textbf{102} (1986), 649--677.
\bibitem{K2} C.~King, The $U(1)$ Higgs model. II. The infinite volume limit, Commun. Math. Phys. \textbf{103} (1986), 323--349.
\bibitem{New80} C.~M.~Newman, Normal fluctuations and the FKG inequalities, Commun. Math. Phys. \textbf{74} (1980), 119--128.
\bibitem{OS73} K.~Osterwalder, R.~Schrader, Axioms for Euclidean Green's functions, Commun. Math. Phys. \textbf{31} (1973), 83--112.
\bibitem{OS75} K.~Osterwalder, R.~Schrader, Axioms for Euclidean Green's functions. II, Commun. Math. Phys. \textbf{42} (1975), 281--305.
\bibitem{OS78} K.~Osterwalder, E.~Seiler, Gauge field theories on a lattice, Ann. Physics \textbf{110} (1978), 440--471.
\bibitem{Pre74} C.~J.~Preston, A generalization of the FKG inequalities, Commun. Math. Phys. \textbf{36} (1974), 233--241.
\bibitem{RS1} M.~Reed, B.~Simon, \emph{Methods of Modern Mathematical Physics. I. Functional Analysis}, revised and enlarged ed., Academic Press, New York, 1980.
\bibitem{Rei88} T.~Reisz, A power counting theorem for Feynman integrals on the lattice, Commun. Math. Phys. \textbf{116} (1988), 81--126.
\bibitem{Rei88b} T.~Reisz, Renormalization of Feynman integrals on the lattice, Commun. Math. Phys. \textbf{117} (1988), 79--108.
\bibitem{Sei82} E.~Seiler, \emph{Gauge Theories as a Problem of Constructive Quantum Field Theory and Statistical Mechanics}, Lecture Notes in Physics \textbf{159}, Springer, Berlin, 1982.
\bibitem{War83} F.~W.~Warner, \emph{Foundations of Differentiable Manifolds and Lie Groups}, Graduate Texts in Mathematics \textbf{94}, Springer, New York, 1983.

\end{thebibliography}

\end{document}

